% AUTO-GENERATED with python3 -m tools.reference.gen_valhall_spec_tex, DO NOT EDIT
\documentclass[10pt,letterpaper]{report}
\usepackage[margin=20mm]{geometry}
\usepackage{charter}
\usepackage[scaled=0.88]{helvet}
\usepackage[T1]{fontenc}
\usepackage{amsmath,amssymb}
\usepackage{booktabs}
\usepackage{longtable}
\usepackage{tabularx}
\usepackage{array}
\usepackage{xcolor}
\usepackage{fancyhdr}
\usepackage{microtype}
\usepackage{hyperref}

\definecolor{primary}{RGB}{36, 52, 118}
\definecolor{subtle}{RGB}{88, 94, 110}
\definecolor{codebg}{RGB}{246, 247, 251}
% Distinct instruction syntax colors (Mnemonic, Modifier, Destination, Source, Immediate)
\definecolor{cMnem}{RGB}{20, 58, 148}
\definecolor{cMod}{RGB}{128, 34, 138}
\definecolor{cDst}{RGB}{16, 118, 60}
\definecolor{cSrc}{RGB}{178, 56, 22}
\definecolor{cImm}{RGB}{0, 114, 128}

\newcommand{\mnem}[1]{\textcolor{cMnem}{\textbf{#1}}}
\newcommand{\imod}[1]{\textcolor{cMod}{#1}}
\newcommand{\dst}[1]{\textcolor{cDst}{\textbf{#1}}}
\newcommand{\src}[1]{\textcolor{cSrc}{#1}}
\newcommand{\imm}[1]{\textcolor{cImm}{#1}}

\hypersetup{
  colorlinks=true,
  linkcolor=primary,
  urlcolor=primary,
  citecolor=primary,
  pdftitle={Valhall Instruction Set Reference (v9--v14)},
  pdfauthor={3-Workstream Ground-Truth Audited Specification}
}

\setlength{\parskip}{5pt plus 2pt minus 1pt}
\setlength{\parindent}{0pt}
\setlength{\headheight}{14pt}

\pagestyle{fancy}
\fancyhf{}
\fancyhead[L]{\small\color{subtle}\textsf{Valhall Instruction Set Reference (v9--v14)}}
\fancyhead[R]{\small\color{subtle}\textsf{\leftmark}}
\fancyfoot[C]{\small\thepage}
\renewcommand{\headrulewidth}{0.4pt}

\newcommand{\codebox}[1]{%
  \par\vspace{2pt}\noindent\colorbox{codebg}{%
    \parbox{\dimexpr\linewidth-2\fboxsep\relax}{\ttfamily\footnotesize #1}%
  }\par\vspace{3pt}%
}

\newenvironment{tightitemize}{%
  \begin{itemize}\setlength{\itemsep}{1pt}\setlength{\parskip}{0pt}\setlength{\parsep}{0pt}%
}{\end{itemize}}

\sloppy
\begin{document}

% ============================================================================
% Title Page
% ============================================================================
\begin{titlepage}
\centering
\vspace*{4.5cm}
{\Huge\bfseries\textsf{Valhall Instruction Set Reference}\par}
\vspace{0.7cm}
{\Large\color{primary}\textsf{Architecture Versions v9 through v14}\par}
\vfill
\begin{minipage}{0.84\textwidth}
\centering\small\color{subtle}
\url{https://github.com/leegao/valhall2vulkan/valhall-isa-reference.pdf}\\[1.2em]
{\bfseries Revision \today}
\end{minipage}
\vspace*{1.8cm}
\end{titlepage}

% ============================================================================
% Provenance, Syntax Color Legend & Change Log
% ============================================================================
\thispagestyle{plain}
\section*{Assembly Syntax Color Legend}
All Valhall instruction syntax blocks and worked examples in this manual are color-coded by operand role:
\codebox{\mnem{FMA.f32}\imod{.clamp\_0\_1} \dst{r0}\imod{.h1}, \src{r1\textasciicircum}\imod{.abs}, \src{r2}\imod{.neg}, \src{u0.w0} \quad\textnormal{\textit{and}}\quad \mnem{SHADD.u64} \dst{[r0:r1]}, \src{[r2:r3]}, \src{r4}\imod{.w0}, \imm{\#2}}
\begin{itemize}
  \item \texttt{\mnem{Blue Bold}}: Canonical mnemonic and type variant (\texttt{\mnem{FMA.f32}}, \texttt{\mnem{LD\_CHECKED.i32}})
  \item \texttt{\imod{Purple}}: Instruction flags, pre-source modifiers (\texttt{\imod{.abs}}, \texttt{\imod{.neg}}, \texttt{\imod{.h1}}), and post-destination modifiers (\texttt{\imod{.h01}})
  \item \texttt{\dst{Green Bold}}: Destination GPR or write staging window (\texttt{\dst{dest}}, \texttt{\dst{r0}}, \texttt{\dst{@r0:r3}})
  \item \texttt{\src{Orange-Red}}: Source GPRs, uniforms, FAU constants, or read staging windows (\texttt{\src{src0}}, \texttt{\src{r1\textasciicircum}}, \texttt{\src{u0.w0}}, \texttt{\src{k1f}})
  \item \texttt{\imm{Teal}}: Inline immediates, offsets, and shift counts (\texttt{\imm{\#0x10}}, \texttt{\imm{\#2}})
\end{itemize}

\newpage
\tableofcontents
\newpage

% ============================================================================
% Chapter 1: Architecture Overview
% ============================================================================
\chapter{Architecture Overview}

\section{GPU Generation Matrix (v9--v14)}
Valhall is a dynamically scheduled, 64-bit fixed-width superscalar GPU ISA introduced in architecture major version 9. It supports 16-thread warps, SIMD-within-a-word 8/16/32-bit operations, and explicit scoreboard slots (\texttt{slot0..slot7}).

\begin{table}[!htbp]
\centering
\caption{Arm Mali Valhall Architecture Generations (v9--v14)}
\begin{tabularx}{\linewidth}{@{}lllX@{}}
\toprule
Arch & Representative GPUs & Codenames & Key ISA Milestones \\
\midrule
v9 & Mali-G77, G57, G78, G68 & Trym, Natt, Borr, Ottr & Initial Valhall ISA; sole version with \texttt{TEX\_WEIGHTS}; includes \texttt{V2F32\_TO\_V2F16}, \texttt{HADD}, \texttt{ILOGB}/\texttt{LOGB}, \texttt{v4*8} ALU. \\
v10 & Mali-G710, G610, G510 & Odinn, Turse-pre & Removes \texttt{TEX\_WEIGHTS}; adds \texttt{source\_format} on \texttt{LD\_VAR\_BUF*} / \texttt{VAR\_TEX\_BUF*}; software cubemap projection. \\
v11 & Mali-G715, Immortalis-G715 & Turse & Adds \texttt{F32\_TO\_F16}, \texttt{AND/OR/XOR\_IMM.i32}, \texttt{MMUL}, hardware \texttt{TEX\_SINGLE.cube.project}, \texttt{CLPER} \texttt{low/prefix}; removes 16 legacy ops (\texttt{V2F32\_TO\_V2F16}, \texttt{HADD}, \texttt{v4*8} ALU, narrow converts). \\
v12 & Mali-G720, Immortalis-G720 & Titan & Adds native FP16 transcendentals (\texttt{FEXP/FLOG/FPOW.f16}) and hardware bounds-checked \texttt{LD\_CHECKED*}/\texttt{ST\_CHECKED*}. \\
v13 & Immortalis-G925 & Krake (Avalon) & Retains v12 shader ISA; CSF v13 queue/descriptor extensions. \\
v14 & Next-Gen Valhall & Avalon+ & Adds \texttt{FEXPF.f32}, \texttt{FLOG.f32}, \texttt{FSINCOS.f32}, \texttt{FSQRT.f32}, \texttt{CLPER\_OR}, \texttt{LD\_MULTI}/\texttt{ST\_MULTI}, \texttt{LD\_VAR\_BUF\_FLAT*}, \texttt{MMUL.f16}, mixed-sign \texttt{IDPADD}/\texttt{MMUL.v4s8}, \texttt{oob\_mode}; removes \texttt{NOT\_OLD}. \\
\bottomrule
\end{tabularx}
\end{table}

\section{Instruction Word, Operands \& Execution Pipeline}
Every Valhall instruction executes in three strictly ordered stages:
\begin{enumerate}
  \item \textbf{Stage A (Pre-Source):} Swizzle/widen/lane extraction runs \emph{before} \texttt{abs}/\texttt{neg}/\texttt{not} (\texttt{kraid/ir.rs:882}). Exception: \texttt{BRANCH}/\texttt{JUMP}/\texttt{JUMP\_EX} \texttt{.not} negates the boolean branch condition, not \texttt{src0} bits.
  \item \textbf{Stage B (Core Op \& Post-Op):} Evaluates the core operation. On \texttt{FMIN}/\texttt{FMAX}, \texttt{sem} runs \emph{before} \texttt{clamp}. On \texttt{V2F32\_TO\_V2F16} and \texttt{F32\_TO\_F16}, \texttt{clamp} runs \emph{before} FP32$\to$FP16 rounding.
  \item \textbf{Stage C (Post-Destination):} \texttt{dest\_width} (\texttt{h0}/\texttt{h1}) is a \textbf{pure lane write-mask} (\texttt{h1} writes \texttt{val32[31:16]} to \texttt{dst[31:16]} without shifting); \texttt{dest\_width\_narrow} (\texttt{h0}/\texttt{h1}/\texttt{h01}) narrows FP32$\to$FP16 using the instruction's \texttt{round} mode; \texttt{dest\_width\_replicate} places/replicates scalar FP16 outputs.
\end{enumerate}

\section{Small Constants Table (FAU Page 0)}
\begin{table}[!htbp]
\centering
\caption{32 Hardware Small Constants (FAU Page 0, Indices \texttt{0xC0..0xDF})}
\begin{tabularx}{\linewidth}{@{}lll|lll@{}}
\toprule
Idx & Value & Meaning & Idx & Value & Meaning \\
\midrule
0 & \texttt{0x00000000} & Zero (\texttt{k0}) & 16 & \texttt{0x3F800000} & Float $1.0$ (\texttt{k1f}) \\
1 & \texttt{0xFFFFFFFF} & All ones / $-1$ (\texttt{k1n}) & 17 & \texttt{0x3DCCCCCD} & Float $0.1$ \\
2 & \texttt{0x7FFFFFFF} & \texttt{INT32\_MAX} / qNaN & 18 & \texttt{0x3EA2F983} & Float $1/\pi$ \\
3 & \texttt{0xFAFCFDFE} & Bytes $(-2,-3,-4,-5)$ & 19 & \texttt{0x3F317218} & Float $\ln 2$ \\
4 & \texttt{0x01000000} & 16-bit $2^8$ & 20 & \texttt{0x40490FDB} & Float $\pi$ \\
5 & \texttt{0x80002000} & Bytes $(0, 32, 0, 128)$ & 21 & \texttt{0x00000000} & Float $0.0$ \\
6 & \texttt{0x70605030} & Bytes $(48, 80, 96, 112)$ & 22 & \texttt{0x477FFF00} & Float $65535.0$ \\
7 & \texttt{0xC0B0A090} & Bytes $(144, 160, 176, 192)$ & 23 & \texttt{0x5C005BF8} & Half $(255.0, 256.0)$ \\
8 & \texttt{0x03020100} & Bytes $(0, 1, 2, 3)$ & 24 & \texttt{0x2E660000} & Half $0.1$ (hi16) \\
9 & \texttt{0x07060504} & Bytes $(4, 5, 6, 7)$ & 25 & \texttt{0x34000000} & Half $0.25$ (hi16) \\
10 & \texttt{0x0B0A0908} & Bytes $(8, 9, 10, 11)$ & 26 & \texttt{0x38000000} & Half $0.5$ (hi16) \\
11 & \texttt{0x0F0E0D0C} & Bytes $(12, 13, 14, 15)$ & 27 & \texttt{0x3C000000} & Half $1.0$ (hi16) \\
12 & \texttt{0x13121110} & Bytes $(16, 17, 18, 19)$ & 28 & \texttt{0x40000000} & Half $2.0$ (hi16) \\
13 & \texttt{0x17161514} & Bytes $(20, 21, 22, 23)$ & 29 & \texttt{0x44000000} & Half $4.0$ (hi16) \\
14 & \texttt{0x1B1A1918} & Bytes $(24, 25, 26, 27)$ & 30 & \texttt{0x48000000} & Half $8.0$ (hi16) \\
15 & \texttt{0x1F1E1D1C} & Bytes $(28, 29, 30, 31)$ & 31 & \texttt{0x42480000} & Half $\pi$ (hi16) \\
\bottomrule
\end{tabularx}
\end{table}

\section{CSF Execution Context \& Descriptor Graph ABI}
\label{sec:csf-abi}

\subsection{ABI}
\label{sec:csf-modes}
In Arm's Mali Command Stream Frontend (CSF, architecture v10--v14), Mesa Panfrost (\texttt{genxml/v10.xml}--\texttt{v14.xml}, \texttt{panvk}), and Arm's proprietary DDK (\texttt{libGLES\_mali.so}), the per-dispatch/per-draw register state and memory-resident descriptor hierarchy is formally the \textbf{CSF Execution Context \& Descriptor Graph ABI} (also called the \textbf{CSF Staging Register \& Descriptor Table ABI}).

Once a CSF command stream populates the per-queue 96$\times$32-bit register file (\texttt{r0..r63} staging window; \texttt{r64..r95} CSHWIF scratch) and executes \texttt{RUN\_COMPUTE}, \texttt{RUN\_IDVS}, or \texttt{RUN\_FRAGMENT}, every ABI object is consumed in one or more of \textbf{three hardware modes}:
\begin{enumerate}
  \item \textbf{Implicit Pre-Execution (Warp Launch \& Preload):} Before instruction \texttt{0x00} executes, hardware reads \texttt{SPD\_k}, \texttt{TSD\_k}, \texttt{WG\_SIZE}/\texttt{JOB\_OFFSET}/\texttt{JOB\_SIZE}, or \texttt{INDEX\_BUFFER}/\texttt{VERTEX\_OFFSET} to configure warp occupancy (32 vs 64 GPRs), FP rounding/FTZ/NaN/Inf modes, TLS/WLS scratch pools, and pre-populated stage GPRs.
  \item \textbf{Explicit Valhall ISA Instruction or Source Operand:} A 64-bit Valhall instruction directly references the ABI object---either via a message opcode (\texttt{LD\_ATTR\_IMM}, \texttt{LEA\_BUF\_IMM}, \texttt{LD\_PKA}, \texttt{LD\_VAR\_IMM}, \texttt{TEX\_SINGLE}, \texttt{VAR\_TEX\_SINGLE}, \texttt{STORE.istream/estream}, \texttt{ATEST}, \texttt{ZS\_EMIT}, \texttt{BLEND.end}) or via an 8-bit source operand (\texttt{mode=2} \texttt{FAU} uniform \texttt{u0..u31}; \texttt{mode=3} \texttt{FauSpecial} such as \texttt{pc}, \texttt{lane\_id}, \texttt{tls\_ptr}, \texttt{wls\_ptr}, \texttt{blend\_descriptor\_0..7}).
  \item \textbf{Implicit Inter-Stage / Post-Execution (Fixed-Function):} Hardware units running between or after shader warps consume ABI descriptors without shader instructions touching them (Index Fetcher, Tiler clipping/culling/16$\times$16 binning, Early-Z, on-chip Tile SRAM clear/preload, and Morton 16$\times$16 U-interleaved / AFBC writeback).
\end{enumerate}

\codebox{%
COMPUTE TIMELINE (Pilot + Main Compute):\\
~~[IMPLICIT] SPD\_0, WG\_SIZE=(1,1,1), JOB\_SIZE=(1,1,1) -> 1 Pilot warp (r0:r1=fau\_dst, r2:r3=srt\_va)\\
~~[EXPLICIT] Pilot: LOAD [r2:r3+off] (SRT\_0->Resource->Buffer[]) -> STORE [r0:r1+off] (populates FAU\_0)\\
~~[SYNC]     CSWAIT scoreboard sync\\
~~[IMPLICIT] SPD\_1, TSD\_0, WG\_SIZE, JOB\_OFFSET, JOB\_SIZE -> preloads r55..r57=WG\_ID, r60..r62=Local/GlobalID\\
~~[EXPLICIT] Main: src=u0..u31 (FAU\_0), src=pc (embedded LUT), LOAD/STORE/ATOM via 64-bit BDA\\
\\
GRAPHICS TIMELINE (RUN\_IDVS + RUN\_FRAGMENT):\\
~~[IMPLICIT] INDEX\_BUFFER, VERTEX/INSTANCE/GLOBAL\_ATTR\_OFF, SPD\_0, TSD\_0 -> r59/r60=VertID, r61=InstID\\
~~[EXPLICIT] Position EBIN: LD\_ATTR\_IMM (SRT\_0) + FAU\_0 -> STORE.istream.end (TILER\_CTX::geometry\_buffer)\\
~~[IMPLICIT] Fixed-Function Tiler: TILER\_CTX, TILER\_FLAGS, DCD0, SCISSOR\_BOX, DEPTH\_CLAMP -> 16x16 bins\\
~~[EXPLICIT] Varying EBIN (visible verts only): LD\_ATTR\_IMM -> STORE.estream.end (VARY\_SIZE buffer)\\
~~[IMPLICIT] FINISH\_TILING -> RUN\_FRAGMENT: FBD\_POINTER, BBOX, ZSD+DCD2 (Early-Z), SPD\_2 -> r59=z, r60=mask\\
~~[EXPLICIT] Fragment EBIN: LD\_VAR\_IMM / VAR\_TEX\_SINGLE (VARY\_SIZE + SRT\_2) -> ATEST/ZS\_EMIT -> BLEND.end\\
~~[IMPLICIT] RenderTarget[0..N-1] (FBD +0xC0): retiles 16x16 Tile SRAM -> Morton 16x16 writeback\_buffer%
}

\subsection{Complete CSF Staging Register Map (\texttt{r0..r63})}
\label{sec:csf-regs}
Each CSF Command Stream Interface (CSI) queue owns an isolated 96$\times$32-bit register file (\texttt{CsfQueueState::regs[96]}). Registers \texttt{r0..r63} form the staging window re-interpreted by \texttt{RUN\_COMPUTE}, \texttt{RUN\_IDVS}, and \texttt{RUN\_FRAGMENT}; \texttt{r64..r95} are CSHWIF scratch/overflow registers. Tagged 64-bit registers pack counts or flags in low/high bits that \textbf{must} be stripped before dereferencing.

\begin{center}\footnotesize
\begin{longtable}{@{}l l l l >{\raggedright\arraybackslash}p{0.34\linewidth}@{}}
\toprule
Regs & \texttt{COMPUTESR} & \texttt{IDVSSR} & \texttt{FRAGMENTSR} & Bitfield Packing, Alignment \& Conditions \\
\midrule\endhead
\texttt{r0:r7} & \texttt{SRT\_0..3} & \texttt{SRT\_0..3} & \texttt{SRT\_0..3} & \texttt{[47:6]} 64B-aligned VA (\texttt{\& \textasciitilde{}0x3FULL}); \texttt{[5:0]} entry count (max 63). \\
\texttt{r8:r15} & \texttt{FAU\_0..3} & \texttt{FAU\_0..3} & \texttt{FAU\_0..3} & \texttt{[47:0]} 48-bit VA; \texttt{[63:56]} 64-bit slot count ($\times 2$ for 32-bit words). \\
\texttt{r16:r23} & \texttt{SPD\_0..3} & \texttt{SPD\_0..2} & \texttt{SPD\_0..3} & \texttt{[47:0]} 32B-aligned VA to \texttt{ShaderProgram} (IDVS: 0=Pos, 1=Var, 2=Frag). \\
\texttt{r24:r31} & \texttt{TSD\_0..3} & \texttt{TSD\_0..3} & \texttt{TSD\_0..3} & \texttt{[47:0]} 64B-aligned VA to 32B \texttt{LocalStorage} (TLS/WLS). \\
\texttt{r32} & \texttt{GLOBAL\_ATTR\_OFF} & \texttt{GLOBAL\_ATTR\_OFF} & --- & \texttt{[31:0]} Global vertex/attribute index offset. \\
\texttt{r33} & \texttt{WG\_SIZE} & \texttt{INDEX\_COUNT} & --- & Comp: \texttt{[9:0]=x-1, [19:10]=y-1, [29:20]=z-1}; IDVS: vertex/index count. \\
\texttt{r34..r36} & \texttt{JOB\_OFFSET\_X/Y/Z} & \texttt{INST\_CNT, IDX\_OFF, VTX\_OFF} & --- & Comp: workgroup ID base offset; IDVS: instance count, index \& vertex offset. \\
\texttt{r37..r39} & \texttt{JOB\_SIZE\_X/Y/Z} & \texttt{INST\_OFF, DCD2, IDX\_BUF\_SZ} & --- & Comp: dispatch grid; IDVS: instance offset, \texttt{DCD2} (ZS read/write), index bytes. \\
\texttt{r40:r41} & --- & \texttt{TILER\_CTX} & \texttt{FBD\_POINTER} & IDVS: 64B-aligned TCD VA; Frag: \texttt{[0]=type, [1]=zs\_crc, [5:2]=rt\_cnt-1, [63:6]=VA}. \\
\texttt{r42:r43} & --- & \texttt{SCISSOR\_BOX} & \texttt{BBOX\_MIN/MAX} & Packed 16-bit \texttt{(X, Y)} min/max inclusive coordinates (\texttt{+1} for extent). \\
\texttt{r44:r45} & --- & \texttt{LOW/HIGH\_DEPTH\_CLAMP} & \texttt{TEM\_POINTER} & IDVS: IEEE-754 \texttt{f32} depth clamp bounds; Frag: Tile Emission Map VA. \\
\texttt{r46:r47} & --- & \texttt{OQ} & \texttt{TEM\_ROW\_STRIDE} & IDVS: occlusion query counter VA; Frag: TEM row stride. \\
\texttt{r48:r49} & --- & \texttt{VARY\_SIZE} & --- & Per-vertex varying buffer allocation size in bytes. \\
\texttt{r50:r51} & --- & \texttt{BLEND\_DESC} & --- & \texttt{[47:4]} 16B-aligned VA (\texttt{\& \textasciitilde{}0xFULL}); \texttt{[3:0]} blend descriptor count. \\
\texttt{r52:r53} & --- & \texttt{ZSD} & --- & \texttt{[47:0]} 32B-aligned VA to 32B \texttt{Depth/stencil} descriptor. \\
\texttt{r54:r55} & --- & \texttt{INDEX\_BUFFER} & --- & \texttt{[47:0]} Index buffer base VA (stride from \texttt{PrimitiveFlags::index\_type}). \\
\texttt{r56} & --- & \texttt{TILER\_FLAGS} & --- & Merged via \texttt{r56 | CSRUNIDVS::flags\_override} into \texttt{PrimitiveFlags}. \\
\texttt{r57..r58} & --- & \texttt{DCD0, DCD1} & --- & \texttt{DCD0}: cull/front-face/kill/A2C; \texttt{DCD1}: sample mask \& RT mask. \\
\texttt{r60..r61} & --- & \texttt{PRIMITIVE\_SIZE, TILER\_FLAGS2} & --- & \texttt{f32} point size / line width and secondary tiler flags. \\
\texttt{r62:r63} & \texttt{ITER\_TRACE\_ID0/1} & \texttt{ITER\_TRACE\_ID0/1} & \texttt{ITER\_TRACE\_ID0/1} & Hardware performance/iterator trace correlation IDs. \\
\bottomrule
\end{longtable}
\end{center}

\subsection{Pointer-Chasing Descriptor Graph \& Memory Hierarchy}
\label{sec:csf-chains}
Every descriptor table indirection follows strict alignment and bounds invariants verified in \texttt{specs/csf/v11.xml}, \texttt{panvk}, and \texttt{valhall\_csf\_pipeline.cpp}:
\codebox{%
1. Vertex Attribute Pulling (LD\_ATTR\_IMM table=T, index=I):\\
~~~IDVSSR::SRT\_0 (va = raw \& \textasciitilde{}0x3F, cnt = raw \& 0x3F)  [Precond: T < cnt]\\
~~~  -> Resource[T] (16B @ va + T*16: address, size, contains\_descriptors=1)\\
~~~  -> Attribute[I] (32B @ Resource[T].address + I*32: format, table=B\_T, buf\_idx=B\_I, off, stride)\\
~~~  -> Resource[B\_T] (16B) -> Buffer[B\_I] (32B @ Resource[B\_T].address + B\_I*32: address, size)\\
~~~  -> Vertex RAM @ Buffer[B\_I].address + Attribute[I].offset + VertexIndex * Attribute[I].stride\\
\\
2. Texture \& Sampler Lookup (TEX\_SINGLE / VAR\_TEX\_SINGLE):\\
~~~IDVSSR::SRT\_2 (64B-aligned VA) -> Resource[T] -> Texture[tex\_idx] (32B, type=2)\\
~~~  -> GenericPlane[mip] (32B @ Texture.surfaces + mip*32, type=10: pointer, row\_stride, Morton 16x16)\\
~~~  +  Sampler[samp\_idx] (32B, type=1: filter, wrap, LOD clamp/bias, compare\_function)\\
\\
3. Framebuffer Contiguous Memory Block (RUN\_FRAGMENT):\\
~~~FRAGMENTSR::FBD\_POINTER (va = raw \& \textasciitilde{}0x3F, type=bit0, zs\_crc=bit1, rt\_cnt=(bits[5:2])+1)\\
~~~  +0x00: FramebufferParameters (72B fields + 56B pad = 128B)\\
~~~  +0x80: ZSCRCExtension (64B, present iff zs\_crc == 1)\\
~~~  +0x80/+0xC0: RenderTarget[0..rt\_cnt-1] (64B each; RGB writeback\_buffer @ +0x20, clear @ +0x30)%
}

\subsection{Descriptor Memory \& Table Entry Layouts (\texttt{specs/csf/v11.xml})}
\label{sec:csf-layouts}

\subsubsection{\texttt{ShaderProgram} (\texttt{SPD}, 32B, align 32) \& Stage Preload Table}
\label{sec:csf-spd}
\begin{center}\footnotesize
\begin{tabularx}{\linewidth}{@{}l l l X@{}}
\toprule
Word:Bits & Width & Field & Architectural Role \\
\midrule
\texttt{0:0..3} & 4b & \texttt{Type} & \texttt{DescriptorType::Shader} (\texttt{8}). \\
\texttt{0:4..7} & 4b & \texttt{Stage} & \texttt{1}=Compute, \texttt{2}=Fragment, \texttt{3}=Vertex. \\
\texttt{0:16..19} & 4b & \texttt{FP Modes} & \texttt{Suppress NaN} (b16), \texttt{Flush to zero mode} (b17:18), \texttt{Suppress Inf} (b19). \\
\texttt{0:28..31} & 4b & \texttt{Exec Flags} & \texttt{Requires helper threads} (b28), \texttt{Shader contains JUMP\_EX} (b29), \texttt{Register allocation} (b30:31: \texttt{0}=64 GPRs, \texttt{2}=32 GPRs). \\
\texttt{1:0..15} & 16b & \texttt{Preload} & Stage-specific bitmask controlling which GPRs hardware initializes at PC \texttt{0x00}. \\
\texttt{2:0..3:31} & 64b & \texttt{Binary} & 48-bit GPU VA of 64-bit Valhall shader object code (\texttt{EBIN}). \\
\bottomrule
\end{tabularx}
\end{center}
\textbf{Implicit Hardware Preloads at PC \texttt{0x00}:} Pilot: \texttt{r0:r1 = pilot\_fau\_dst\_va}, \texttt{r2:r3 = srt\_base\_va}; Compute: \texttt{r55..r57 = WorkGroupID.xyz + JOB\_OFFSET}, \texttt{r60..r62 = LocalInvocationID.xyz} (or \texttt{GlobalInvocationID.xyz}); Vertex (Pos \& Var): \texttt{r59, r60 = VertexIndex}, \texttt{r61 = InstanceIndex}; Fragment: \texttt{r59 = floatBitsToUint(FragCoord.z)}, \texttt{r60 = sample coverage mask}.

\subsubsection{\texttt{Resource} Table Entry (\texttt{SRT}, 16B, align 64) \& \texttt{DescriptorType}}
\label{sec:csf-srt}
An \texttt{SRT\_k} register points to an array of \texttt{count} 16-byte \texttt{Resource} entries. Each 32-byte child descriptor carries \texttt{DescriptorType} in \texttt{word0[3:0]}: \texttt{0}=Null, \texttt{1}=Sampler, \texttt{2}=Texture, \texttt{5}=Attribute, \texttt{7}=Depth/stencil, \texttt{8}=Shader, \texttt{9}=Buffer, \texttt{10}=Plane.
\begin{center}\footnotesize
\begin{tabularx}{\linewidth}{@{}l l l X@{}}
\toprule
Word:Bits & Width & Field & Architectural Role \\
\midrule
\texttt{0:0..1:23} & 56b & \texttt{Address} & Base GPU VA of the child 32-byte descriptor array (or raw data table). \\
\texttt{1:24} & 1b & \texttt{Contains descriptors} & \texttt{1} if array holds 32B hardware descriptors (\texttt{Buffer}/\texttt{Attribute}/\texttt{Texture}/\texttt{Sampler}). \\
\texttt{2:0..3:31} & 64b & \texttt{Size} & Table byte size (\texttt{Size / 32} = number of 32B descriptors; bounds-checked by HW). \\
\bottomrule
\end{tabularx}
\end{center}

\subsubsection{\texttt{Buffer} (32B, \texttt{type=9}) \& \texttt{Attribute} (32B, \texttt{type=5})}
\label{sec:csf-buf-attr}
\begin{center}\footnotesize
\begin{tabularx}{\linewidth}{@{}l l l X@{}}
\toprule
Struct & Word:Bits & Field & Architectural Role \\
\midrule
\texttt{Buffer} & \texttt{0:0..7}, \texttt{1:0..31}, \texttt{5:0..31} & \texttt{Type=9}, \texttt{Buffer type}, \texttt{Size} & 64-bit byte size (\texttt{Size | Size\_hi<<32}) checked by \texttt{LD\_CHECKED*}/\texttt{ST\_CHECKED*}. \\
\texttt{Buffer} & \texttt{2:0..3:31}, \texttt{4:0..31} & \texttt{Address} (64b), \texttt{Stride} & Base GPU VA of SSBO/UBO/vertex buffer; optional packet/attribute stride. \\
\texttt{Attribute} & \texttt{0:0..31} & \texttt{Type=5}, \texttt{Offset en}, \texttt{Format} (22b) & Hardware \texttt{PixelFormat} (e.g., \texttt{RGB32F}, \texttt{RGBA8\_UNORM}) converted on load. \\
\texttt{Attribute} & \texttt{1:0..7}, \texttt{1:24..29} & \texttt{Table} (6b), \texttt{Frequency}, \texttt{Divisor R/E} & Selects \texttt{Resource[Table]} holding target \texttt{Buffer[]} and per-vertex/instanced rate. \\
\texttt{Attribute} & \texttt{2:0..5:31} & \texttt{Offset}, \texttt{Buffer index}, \texttt{Stride}, \texttt{Divisor D} & Effective VA = \texttt{Buffer[idx].address + Offset + VertexIdx * Stride}. \\
\bottomrule
\end{tabularx}
\end{center}

\subsubsection{\texttt{Texture} (32B, \texttt{type=2}), \texttt{GenericPlane} (32B, \texttt{type=10}) \& \texttt{Sampler} (32B, \texttt{type=1})}
\label{sec:csf-tex-samp}
\begin{center}\footnotesize
\begin{tabularx}{\linewidth}{@{}l l l X@{}}
\toprule
Struct & Word:Bits & Field & Architectural Role \\
\midrule
\texttt{Texture} & \texttt{0:0..31}, \texttt{1:0..31} & \texttt{Type=2}, \texttt{Dim}(3b), \texttt{Format}(22b), \texttt{W/H}(16b) & 1D/2D/3D/Cube dimensions, texel format, and base width/height. \\
\texttt{Texture} & \texttt{2:0..3:31}, \texttt{6..7} & \texttt{Swizzle}, \texttt{Levels}, \texttt{Min/Max LOD}, \texttt{Array/Depth} & Component swizzle, mipmap count, MSAA sample count, array layers, 3D depth. \\
\texttt{Texture} & \texttt{4:0..5:31} & \texttt{Surfaces} (64b VA) & Pointer to contiguous array of 32B \texttt{GenericPlane} descriptors (1 per mip/plane). \\
\texttt{GenericPlane} & \texttt{0:8..31}, \texttt{2:0..3:15} & \texttt{Clump ordering/format}, \texttt{Pointer} (48b) & Surface tiling (\texttt{Tiled\_U\_Interleaved} Morton 16$\times$16, linear, AFBC) \& texel VA. \\
\texttt{GenericPlane} & \texttt{1:0..31}, \texttt{4:0..7:31} & \texttt{Size}, \texttt{Row/Slice stride}, \texttt{W/H} & Byte size, row/slice pitch, and plane dimensions. \\
\texttt{Sampler} & \texttt{0:4..31}, \texttt{1..2}, \texttt{4..7} & \texttt{Wrap S/T/R}, \texttt{Filter}, \texttt{LOD/Compare}, \texttt{Border} & Minify/magnify/mipmap mode, LOD clamp \& bias, anisotropy, shadow compare, RGBA border. \\
\bottomrule
\end{tabularx}
\end{center}

\subsubsection{\texttt{LocalStorage} (\texttt{TSD}, 32B, align 64), \texttt{FAU} \& \texttt{ValhallSidebandBuffer}}
\label{sec:csf-tsd-fau}
\begin{center}\footnotesize
\begin{tabularx}{\linewidth}{@{}l l l X@{}}
\toprule
Word:Bits & Width & Field (\texttt{LocalStorage} / \texttt{TSD}) & Architectural Role \\
\midrule
\texttt{0:0..4} & 5b & \texttt{TLS Size} & Log2 per-thread scratch stack size (spills \& \texttt{Private} arrays). \\
\texttt{1:0..12} & 13b & \texttt{WLS Instances}, \texttt{WLS Size Base/Scale} & Workgroup Local Storage (\texttt{shared} memory) pool sizing per workgroup. \\
\texttt{2:0..3:31} & 48b+4b & \texttt{TLS Base Pointer}, \texttt{TLS Address Mode} & Base VA returned via \texttt{FauSpecial::tls\_ptr} for \texttt{SEG\_ADD}/\texttt{LOAD}/\texttt{STORE}. \\
\texttt{4:0..5:31} & 64b & \texttt{WLS Base Pointer} & Base VA returned via \texttt{FauSpecial::wls\_ptr} for workgroup shared memory. \\
\bottomrule
\end{tabularx}
\end{center}
\textbf{FAU \& 8-Bit Source Multiplexing:} Every 8-bit Valhall source operand splits \texttt{mode = raw >> 6}: \texttt{0/1}=GPR \texttt{r0..r63} (\texttt{1}=last-use \texttt{\textasciicircum}); \texttt{2}=\texttt{FauUniform} reading \texttt{FAU\_k} (\texttt{u0..u31} on Page 0; 64-bit slot count in \texttt{FAU\_k[63:56]}); \texttt{3, bit5=0}=\texttt{SmallConstant} (Table~1.2); \texttt{3, bit5=1}=\texttt{FauSpecial} (\texttt{pc}, \texttt{lane\_id}, \texttt{lane\_width=16}, \texttt{warp\_id}, \texttt{framebuffer\_size}, \texttt{atest\_datum}, \texttt{sample}, \texttt{blend\_descriptor\_0..7}, \texttt{tls\_ptr}, \texttt{wls\_ptr}). In \texttt{valhall2spirv}, the 344-byte \texttt{ValhallSidebandBuffer} (Set~0, Binding~0) bridges \texttt{fau\_words[64]} (\texttt{+0}), \texttt{shader\_pc\_base} (\texttt{+256}), \texttt{srt\_base\_va} (\texttt{+264}), \texttt{pilot\_fau\_dst\_va} (\texttt{+272}), \texttt{blend\_descs[4]} (\texttt{+280}), \texttt{local\_size\_xyz+lane\_width} (\texttt{+312}), and \texttt{fb\_info} (\texttt{+328}).

\subsubsection{\texttt{TilerContext} (\texttt{TCD}, 128B), \texttt{PrimitiveFlags} \& \texttt{DCD0..2}}
\label{sec:csf-tiler}
\begin{center}\footnotesize
\begin{tabularx}{\linewidth}{@{}l l l X@{}}
\toprule
Struct / Reg & Word:Bits & Key Fields & Fixed-Function / Shader Role \\
\midrule
\texttt{TilerContext} & \texttt{0..1}, \texttt{6..7}, \texttt{8..9} & \texttt{Polygon List}, \texttt{Heap}, \texttt{Geometry buffer} & 16$\times$16 tile polygon lists, heap VA, and Position FIFO written by \texttt{STORE.istream}. \\
\texttt{TilerContext} & \texttt{2..4} & \texttt{Hierarchy Mask}, \texttt{FB W/H}, \texttt{Layers} & Hierarchical binning mask, framebuffer dimensions, and provoking vertex rule. \\
\texttt{PrimitiveFlags} (\texttt{r56}) & \texttt{0:0..31} & \texttt{Draw mode}, \texttt{Index type}, \texttt{Secondary Shader}, \texttt{Depth Cull} & OR'd with \texttt{CSRUNIDVS::flags\_override}; enables Varying EBIN (\texttt{b18}). \\
\texttt{DCD0} (\texttt{r57}) & \texttt{0:0..26} & \texttt{Front Face CCW}, \texttt{Cull Front/Back}, \texttt{A2C}, \texttt{Cull Zero Area} & Rasterization culling, alpha-to-coverage, \texttt{Shader Modifies Coverage} (b20). \\
\texttt{DCD1} (\texttt{r58}) & \texttt{0:0..27} & \texttt{Sample Mask} (16b), \texttt{RT Mask} (11b) & Active MSAA sample mask and color render-target enable mask. \\
\texttt{DCD2} (\texttt{r38}) & \texttt{0:0..25} & \texttt{Read/Write Mask}, \texttt{No Shader Depth/Stencil Read} & Controls Early-Z (\texttt{b25=1}) vs Late-Z (\texttt{b25=0}) with \texttt{ZS\_EMIT}/\texttt{ATEST}. \\
\bottomrule
\end{tabularx}
\end{center}

\subsubsection{\texttt{Depth/stencil} (\texttt{ZSD}, 32B, \texttt{type=7}) \& \texttt{Blend} (16B, align 64)}
\label{sec:csf-zsd-blend}
\begin{center}\footnotesize
\begin{tabularx}{\linewidth}{@{}l l l X@{}}
\toprule
Struct & Word:Bits & Key Fields & Hardware Role \\
\midrule
\texttt{Depth/stencil} & \texttt{0:0..31}, \texttt{1..2} & \texttt{Front/Back Stencil Ops}, \texttt{Stencil from shader}(b30), \texttt{Masks/Ref} & Early/Late stencil test; \texttt{b30=1} accepts \texttt{ZS\_EMIT.stencil}. \\
\texttt{Depth/stencil} & \texttt{4:22..31}, \texttt{5..7} & \texttt{Depth write}(b27), \texttt{Depth func}(b29:31), \texttt{Units/Factor/Clamp} & Depth compare function, write enable, and polygon offset bias. \\
\texttt{Blend} & \texttt{0:0..31}, \texttt{1:0..31} & \texttt{Load Dest}, \texttt{Enable}(b9), \texttt{sRGB}, \texttt{Constant}, \texttt{Equation} & Fixed-function blend enable, constant alpha, and RGB/A blend equation. \\
\texttt{Blend} & \texttt{2:0..3:31} & \texttt{Internal} (\texttt{Mode} b0:1: Fixed-Function vs \texttt{Shader} PC) & Consumed by \texttt{BLEND.end} via \texttt{FauSpecial::blend\_descriptor\_0..7}. \\
\bottomrule
\end{tabularx}
\end{center}

\subsubsection{\texttt{FramebufferParameters} (128B), \texttt{ZSCRCExtension} (64B) \& \texttt{RenderTarget} (64B)}
\label{sec:csf-fbd}
\begin{center}\footnotesize
\begin{tabularx}{\linewidth}{@{}l l l X@{}}
\toprule
Struct & Word:Bits & Key Fields & Hardware Role \\
\midrule
\texttt{FramebufferParams} & \texttt{8..10}, \texttt{11:0..31} & \texttt{W/H}, \texttt{Bound Min/Max}, \texttt{Sample Cnt}, \texttt{Tile Size}, \texttt{RT Cnt} & Tile grid dimensions (16$\times$16 Tile SRAM), MSAA count, and color buffer allocation. \\
\texttt{FramebufferParams} & \texttt{12:0..31}, \texttt{13}, \texttt{14..15} & \texttt{S/Z Clear/Preload/Unload}, \texttt{Has ZS CRC}(b21), \texttt{Z Clear}, \texttt{Tiler} & Initializes Tile SRAM depth/stencil clear values and links Tiler polygon lists. \\
\texttt{ZSCRCExtension} & \texttt{0..15} (64B) & \texttt{CRC}, \texttt{ZS} / \texttt{AFBC ZS}, \texttt{S} / \texttt{AFBC S} & Present at \texttt{FBD +0x80} when \texttt{FBD\_POINTER[1]==1}; backs \texttt{LD\_TILE.z\_stencil}. \\
\texttt{RenderTarget[i]} & \texttt{0..1}, \texttt{8..11}, \texttt{12..15} & \texttt{Internal/Writeback Format}, \texttt{Block Format}, \texttt{Writeback buffer}, \texttt{Clear} & 64B per RT; \texttt{Writeback buffer} (\texttt{+0x20}) holds destination VA \& \texttt{row\_stride}; \texttt{Clear} (\texttt{+0x30}) holds 128-bit clear color. \\
\bottomrule
\end{tabularx}
\end{center}
\textbf{Morton 16$\times$16 \texttt{Tiled\_U\_Interleaved} Bit Interleaving:} Within each 16$\times$16 tile, texel $(t_x, t_y)$ maps to linear tile index $b_7 b_6 b_5 b_4 b_3 b_2 b_1 b_0$ where $b_0 = t_{x0} \oplus t_{y0}$, $b_1 = t_{y0}$, $b_2 = t_{x1} \oplus t_{y1}$, $b_3 = t_{y1}$, $b_4 = t_{x2}$, $b_5 = t_{y2}$, $b_6 = t_{x3}$, $b_7 = t_{y3}$.

\subsection{Worked End-to-End ABI Traces (Compute \& Graphics)}
\label{sec:csf-traces}

\subsubsection*{Trace A: End-to-End Compute Shader Execution (\texttt{sample.comp}, Mali-G615, 256 Threads)}
Verified bit-exact across 256/256 output elements and 5/5 atomic debug counters (\texttt{tests/test\_moc\_roundtrip\_compute.cpp}):
\begin{enumerate}
  \item \textbf{Phase 1 --- Pilot Shader (\texttt{EBIN\_44}, 1 thread):} The command stream populates \texttt{COMPUTESR::WG\_SIZE = 0x0} ($1\times 1\times 1$), \texttt{JOB\_SIZE = (1, 1, 1)}, \texttt{SRT\_0 = 0x500000010000}, and \texttt{FAU\_0 = 0x500000020000}, then fires \texttt{RUN\_COMPUTE (spd=0, srt=0, fau=0)}. Hardware preloads \texttt{r0:r1 = 0x500000020000} (\texttt{pilot\_fau\_dst\_va}) and \texttt{r2:r3 = 0x500000010000} (\texttt{srt\_base\_va}).
  \codebox{%
  0x00: LOAD.i64  @r4:r5, [r2:r3], \#+0   ; reads SRT\_0::Resource[0].address (Buffer[] table VA)\\
  0x08: LOAD.i64  @r6:r7, [r4:r5], \#+8   ; reads Buffer[0].address (input SSBO VA)\\
  0x10: STORE.i64 @r6:r7, [r0:r1], \#+0   ; writes input SSBO VA into FAU\_0 slot u0\\
  ...   (repeats for Buffer[1] -> u1 @ +8, Buffer[2] -> u2 @ +16, and writes 16u -> u3.w0 @ +24)\\
  0x48: NOP.end                           ; signals scoreboard; CSWAIT unblocks Main Compute%
  }
  \item \textbf{Phase 2 --- Main Compute (\texttt{EBIN\_42}, 256 threads + Embedded LUT at \texttt{+0x100}):} The command stream sets \texttt{SPD\_1} (64 GPRs/thread, PC \texttt{0x00 = EBIN\_42}), \texttt{WG\_SIZE = 0x0000000f} ($16\times 1\times 1$), \texttt{JOB\_OFFSET = (0,0,0)}, and \texttt{JOB\_SIZE = (16,1,1)}, spawning 16 workgroups. Hardware preloads \texttt{r55..r57 = WorkGroupID.xyz} and \texttt{r60..r62 = LocalInvocationID.xyz}.
  \codebox{%
  ; 1. Global thread ID using implicit preloads r55/r60 + Pilot-populated FAU u3.w0 (16):\\
  IMAD.u32     r0, r55, u3.w0, r60           ; gid = WorkGroupID.x * 16 + LocalInvocationID.x\\
  ; 2. Direct 64-bit BDA load from input SSBO (FAU u0):\\
  SHADD.u64    [r2:r3], u0, r0.w0, \#2       ; \&input[gid] = u0 + (gid << 2)\\
  LOAD.i32     @r4, [r2:r3], \#0             ; reads input[gid]\\
  ; 3. Embedded constant LUT read via FauSpecial::pc (+0x100 past .end inside shader BO):\\
  IADD\_IMM.u64 [r6:r7], pc, \#+0x100         ; \&packed\_trit\_data\_0[0]\\
  LOAD.i8      @r8, [r6:r7], \#0             ; reads LUT byte\\
  ; 4. Atomic debug counters (FAU u2) \& output store (FAU u1):\\
  ATOM1.i32.ainc.slot0 [u2 + off]            ; increments 5 debug counters -> 256\\
  STORE.i32.slot0.end @r10, [u1 + gid*4]     ; writes output[gid] and terminates warp%
  }
\end{enumerate}

\subsubsection*{Trace B: End-to-End Graphics Execution (\texttt{vkcube\_frame1}, Pixel~8 \texttt{libr54p3\_vulkan\_icd.so}, 256$\times$256)}
Every raw 64-bit word and address below comes verbatim from \texttt{tests/data/vkcube\_frame1/}:
\begin{enumerate}
  \item \textbf{Stage 1 --- Position EBIN (\texttt{pos\_shader.bin}, 54 insns, \texttt{SPD\_0 = 0x10050580}):} \texttt{IDVSSR::SRT\_0 = 0x000078682c5fc1c8} strips \texttt{\& \textasciitilde{}0x3F} $\to$ base VA \texttt{0x78682c5fc1c0}, \texttt{count = 8}. \texttt{FAU\_0 = 0x010078682c5fc240} gives \texttt{slots64 = 1} (\texttt{u0.w0 = 128.0f}, \texttt{u0.w1 = 256.0f}). The Index Fetcher reads 36 \texttt{uint16} indices from \texttt{INDEX\_BUFFER} (\texttt{0x78685bd17000}) and preloads \texttt{r59 = r60 = VertexIndex}, \texttt{r61 = InstanceIndex (0)}.
  \codebox{%
  0x000: 0x0066800622017d7c  LD\_ATTR\_IMM @r0:r2, r60, r61, \#0, \#0\\
  ~~-> Chases SRT\_0 -> Resource[1] -> Attribute[0] (RGB32F, off=0, stride=24, table=2)\\
  ~~-> Resource[2] -> Buffer[0] (0x78685bd19000) -> loads vec3 position into @r0:r2\\
  0x008..0x020:              LD\_PKA @r4..r7, u\_buf + 0x00..0x30 (flow=wait0)\\
  ~~-> Chases Resource[4] -> Buffer[1] (64B @ 0x78685bd18000) -> loads 4x4 MVP matrix\\
  0x028..0x1a0:              48x FMA.f32 / FMUL.f32 reading FAU\_0 u0.w0 (128.0f), u0.w1 (256.0f)\\
  0x1a8: 0x7861480839000040  STORE.istream.end @r0:r3 (flow=0xf = END)\\
  ~~-> Writes 16B clip-space gl\_Position into TILER\_CTX::geometry\_buffer FIFO%
  }
  \item \textbf{Stage 2 --- Fixed-Function Tiler (100\% Implicit):} Zero shader instructions execute. The Tiler merges \texttt{r56 | flags\_override = 0x00470208} (\texttt{Triangles}, \texttt{UINT16}, \texttt{secondary\_shader=1}), applies \texttt{DCD0 = 0x04010040} (CCW front face, zero-area cull), clips against \texttt{SCISSOR\_BOX = (0,0)..(255,255)} and \texttt{DEPTH\_CLAMP = [-3.77e-5f, 1.0f]}, and bins surviving triangles into 16$\times$16 tile polygon lists in \texttt{TILER\_CTX} (\texttt{0x78682c63e400}).
  \item \textbf{Stage 3 --- Varying EBIN (\texttt{var\_shader.bin}, 3 insns, \texttt{SPD\_1 = 0x10050800}):} Runs \emph{only} for vertices of surviving visible triangles; hardware reloads \texttt{r59..r61}.
  \codebox{%
  0x000: 0x0066800622117d7c  LD\_ATTR\_IMM @r0:r2, r60, r61, \#0, \#1\\
  ~~-> Chases Resource[1] -> Attribute[1] (RGB32F, off=12, stride=24, table=2) -> vec3 color\\
  0x010: 0x7861400632000044  STORE.estream.end @r0:r2 -> writes 12B into VARY\_SIZE (0x0c) buffer%
  }
  \item \textbf{Stage 4 --- Fragment EBIN (\texttt{frag\_shader.bin}, 3 insns, \texttt{SPD\_2 = 0x10050900}) \& Tile Writeback:} \texttt{FBD\_POINTER = 0x000078682c63e483} unpacks \texttt{type=1, zs\_crc=1, rt\_count=1}, base VA \texttt{0x78682c63e480}, clearing 16$\times$16 Tile SRAM to \texttt{z\_clear = 1.0f} and \texttt{RTD[0]::clear\_color = 0xff2e241f}. \texttt{ZSD} (\texttt{0x78685bcfc0e0}) + \texttt{DCD2} (\texttt{no\_shader\_depth\_read=1}) run implicit Early-Z (\texttt{Less}) before spawning warps with \texttt{r59 = FragCoord.z}, \texttt{r60 = mask}.
  \codebox{%
  0x000: 0x685c8006220000c0  LD\_VAR\_IMM.f32.v3 @r0:r2, slot=0 (flow=0xd wait/reconverge)\\
  ~~-> Hardware barycentric interpolator reads 12B vec3 color from VARY\_SIZE into r0..r2\\
  0x008: 0x4910c33f800000c0  MOV\_IMM r3, \#0x3f800000 (sets alpha = 1.0f)\\
  0x010: 0x787f4008327c00f0  BLEND.end @r0:r3, blend\_descriptor\_0 (flow=0xf = END)\\
  ~~-> IDVSSR::BLEND\_DESC (0x78685bcfc0c1 \& \textasciitilde{}0xF -> 0x78685bcfc0c0, count=1) -> Blend[0] -> Tile SRAM%
  }
  Once all triangles in a 16$\times$16 tile complete, fixed-function hardware retiles on-chip Tile SRAM into Morton 16$\times$16 \texttt{Tiled\_U\_Interleaved} layout at \texttt{RTD[0]} (\texttt{FBD +0xC0}, \texttt{writeback\_buffer.base = 0x78685bd1b000}, \texttt{row\_stride = 16384}). In textured pipelines (\texttt{test\_moc\_roundtrip\_graphics.cpp}), a single \texttt{VAR\_TEX\_SINGLE} fuses \texttt{LD\_VAR\_IMM} (\texttt{VARY\_SIZE}) with \texttt{SRT\_2} $\to$ \texttt{Texture[0]} $\to$ \texttt{GenericPlane} (Morton detile) + \texttt{Sampler[0]} directly into \texttt{@r0:r3} for \texttt{BLEND.end}.
\end{enumerate}

% ============================================================================
% Chapter 2: Enumerations & Modifiers (Every Entry + Worked Example, All 128 Enums)
% ============================================================================
\chapter{Enumerations \& Modifiers}
This chapter documents all \textbf{128} Valhall modifier enumerations directly from \texttt{specs/kraid/isa-v9-v14.xml}.

\section{Flow Control (flow\_control\_m, Bits 59:62)}
\label{enum:flow_control_m}
Embedded scheduling, scoreboard waits, warp reconvergence, and shader termination.
\par\textbf{CSF ABI Interaction:} Scoreboard waits (\texttt{.wait0..wait}) synchronize async \texttt{SRT}/\texttt{LOAD}/\texttt{TEX}/\texttt{LD\_VAR} messages; \texttt{.end} signals warp completion to \hyperref[sec:csf-modes]{CSF \texttt{RUN\_COMPUTE}/\texttt{RUN\_IDVS}/\texttt{RUN\_FRAGMENT}} and triggers \texttt{CSWAIT} or Tiler/Tile-SRAM writeback.
\begin{center}\footnotesize
\begin{longtable}{@{}rll>{\raggedright\arraybackslash}p{0.62\linewidth}@{}}
\toprule
Enc & Value & Arch & Description \& Worked Example \\
\midrule\endhead
\texttt{0} & \texttt{\imod{none}} & v9 & No scoreboard wait or warp control action; continue immediately to next instruction. \newline {\scriptsize\ttfamily \mnem{FADD.f32} \dst{r0}, \src{r1}, \src{r2}} \\[2pt]
\texttt{1} & \texttt{\imod{wait0}} & v9 & Stall warp until all asynchronous messages in scoreboard slot 0 complete. \newline {\scriptsize\ttfamily \mnem{FADD.f32}\imod{.wait0} \dst{r0}, \src{r1}, \src{r2}} \\[2pt]
\texttt{2} & \texttt{\imod{wait1}} & v9 & Stall warp until all asynchronous messages in scoreboard slot 1 complete. \newline {\scriptsize\ttfamily \mnem{FADD.f32}\imod{.wait1} \dst{r0}, \src{r1}, \src{r2}} \\[2pt]
\texttt{3} & \texttt{\imod{wait01}} & v9 & Stall warp until both scoreboard slots 0 and 1 drain completely. \newline {\scriptsize\ttfamily \mnem{FMUL.f32}\imod{.wait01} \dst{r0}, \src{r1}, \src{r2}} \\[2pt]
\texttt{4} & \texttt{\imod{wait2}} & v9 & Stall warp until all asynchronous messages in scoreboard slot 2 complete. \newline {\scriptsize\ttfamily \mnem{IADD.s32}\imod{.wait2} \dst{r0}, \src{r1}, \src{r2}} \\[2pt]
\texttt{5} & \texttt{\imod{wait02}} & v9 & Stall warp until both scoreboard slots 0 and 2 drain completely. \newline {\scriptsize\ttfamily \mnem{FMA.f32}\imod{.wait02} \dst{r0}, \src{r1}, \src{r2}, \src{r3}} \\[2pt]
\texttt{6} & \texttt{\imod{wait12}} & v9 & Stall warp until both scoreboard slots 1 and 2 drain completely. \newline {\scriptsize\ttfamily \mnem{FMA.f32}\imod{.wait12} \dst{r0}, \src{r1}, \src{r2}, \src{r3}} \\[2pt]
\texttt{7} & \texttt{\imod{wait012}} & v9 & Stall warp until scoreboard slots 0, 1, and 2 all complete. \newline {\scriptsize\ttfamily \mnem{MOV.i32}\imod{.wait012} \dst{r0}, \src{r1}} \\[2pt]
\texttt{8} & \texttt{\imod{wait0126}} & v9 & Stall warp until slots 0, 1, 2, and tile/blend slot 6 all complete. \newline {\scriptsize\ttfamily \mnem{NOP}\imod{.wait0126}} \\[2pt]
\texttt{9} & \texttt{\imod{wait}} & v9 & Full scoreboard barrier: stall until all in-flight slots (0--7) complete. \newline {\scriptsize\ttfamily \mnem{NOP}\imod{.wait}} \\[2pt]
\texttt{10} & \texttt{\imod{reconverge}} & v9 & Signal hardware warp reconvergence point after divergent control flow. \newline {\scriptsize\ttfamily \mnem{FADD.f32}\imod{.reconverge} \dst{r0}, \src{r1}, \src{r2}} \\[2pt]
\texttt{12} & \texttt{\imod{wait\_resource}} & v10--v13 & Wait for hardware resource availability (v10--v13 barrier/queue sync). \newline {\scriptsize\ttfamily \mnem{NOP}\imod{.wait\_resource}} \\[2pt]
\texttt{13} & \texttt{\imod{discard}} & v9 & Mark active fragment lanes as helper/discarded (suppresses side effects \& writes). \newline {\scriptsize\ttfamily \mnem{NOP}\imod{.discard}} \\[2pt]
\texttt{14} & \texttt{\imod{continue}} & v14 & Continue hardware-managed loop/warp execution (v14+). \newline {\scriptsize\ttfamily \mnem{NOP}\imod{.continue}} \\[2pt]
\texttt{15} & \texttt{\imod{end}} & v9 & Terminate shader warp execution and release register/tile resources. \newline {\scriptsize\ttfamily \mnem{NOP}\imod{.end}} \\[2pt]
\bottomrule
\end{longtable}
\end{center}
\par\textbf{Combined Walkthroughs:}
\par\textbf{Example 1 (Wait on slot0 + warp reconvergence):}
\codebox{\mnem{FADD.f32}\imod{.wait0.reconverge} \dst{r0}, \src{r1}, \src{r2}}
Stalls until the asynchronous message in \texttt{slot0} writes \texttt{r1}, executes \texttt{FADD.f32}, and marks the hardware reconvergence point after a divergent branch.
\par\textbf{Example 2 (Full scoreboard drain + shader end):}
\codebox{\mnem{NOP}\imod{.wait.end}}
Drains all in-flight asynchronous slots (\texttt{slot0..slot7}) so all SSBO/tile writes complete before releasing the warp (\texttt{.end}).

\section{FP32 Source Widen (widen\_m)}
\label{enum:widen_m}
Widens an FP16 lane to exact IEEE-754 FP32 before applying \texttt{abs}/\texttt{neg}. Default: \texttt{none}.

\begin{center}\footnotesize
\begin{longtable}{@{}rll>{\raggedright\arraybackslash}p{0.62\linewidth}@{}}
\toprule
Enc & Value & Arch & Description \& Worked Example \\
\midrule\endhead
\texttt{0} & \texttt{\imod{none}} & v9--v14 & Use full 32-bit operand without narrowing/widening. \newline {\scriptsize\ttfamily \mnem{ATEST} \dst{r0}, \src{r1}\imod{}} \\[2pt]
\texttt{1} & \texttt{\imod{h0}} & v9--v14 & Select low 16-bit halfword \texttt{[15:0]} (widen to 32-bit if FP/int widen). \newline {\scriptsize\ttfamily \mnem{ATEST} \dst{r0}, \src{r1}\imod{.h0}} \\[2pt]
\texttt{2} & \texttt{\imod{h1}} & v9--v14 & Select high 16-bit halfword \texttt{[31:16]} (widen to 32-bit if FP/int widen). \newline {\scriptsize\ttfamily \mnem{ATEST} \dst{r0}, \src{r1}\imod{.h1}} \\[2pt]
\bottomrule
\end{longtable}
\end{center}
\par\textbf{Combined Walkthroughs:}
\par\textbf{Example 1 (Widen upper FP16 lane then negate):}
\codebox{\mnem{FADD.f32} \dst{r0}, \src{r1}\imod{.h1.neg}, \src{r2}}
With \texttt{r1 = 0x3C000000} (\texttt{h1 = 1.0h}), \texttt{r2 = 5.0f}: Stage A widens \texttt{h1 -\textgreater{} 1.0f} first, then applies \texttt{.neg -\textgreater{} -1.0f}, then adds \texttt{r2} -\textgreater{} \texttt{r0 = 4.0f}.
\par\textbf{Example 2 (Widen low FP16 subnormal + abs):}
\codebox{\mnem{FMA.f32} \dst{r0}, \src{r1}\imod{.h0.abs}, \src{r2}, \src{r3}}
With \texttt{r1.h0 = 0x8001} (\texttt{-2\textasciicircum{}-24h}): Stage A promotes to exact FP32 \texttt{-2\textasciicircum{}-24f}, \texttt{.abs} clears sign -\textgreater{} \texttt{+2\textasciicircum{}-24f}, then multiplies by \texttt{r2} and adds \texttt{r3}.

\section{FP16 Vector Swizzle (swiz\_m)}
\label{enum:swiz_m}
Halfword swizzle for \texttt{v2f16} operands (\texttt{hXY} places half \texttt{X} in \texttt{[15:0]}, \texttt{Y} in \texttt{[31:16]}). Default: \texttt{h01}.

\begin{center}\footnotesize
\begin{longtable}{@{}rll>{\raggedright\arraybackslash}p{0.62\linewidth}@{}}
\toprule
Enc & Value & Arch & Description \& Worked Example \\
\midrule\endhead
\texttt{0} & \texttt{\imod{h00}} & v9--v14 & 16-bit halfword swizzle \texttt{h00}: lane 0 $\leftarrow$ h0, lane 1 $\leftarrow$ h0. \newline {\scriptsize\ttfamily \mnem{CLZ} \dst{r0}, \src{r1}\imod{.h00}, \src{r2}} \\[2pt]
\texttt{1} & \texttt{\imod{h10}} & v9--v14 & 16-bit halfword swizzle \texttt{h10}: lane 0 $\leftarrow$ h1, lane 1 $\leftarrow$ h0. \newline {\scriptsize\ttfamily \mnem{CLZ} \dst{r0}, \src{r1}\imod{.h10}, \src{r2}} \\[2pt]
\texttt{2} & \texttt{\imod{h01}} & v9--v14 & 16-bit halfword swizzle \texttt{h01}: lane 0 $\leftarrow$ h0, lane 1 $\leftarrow$ h1. \newline {\scriptsize\ttfamily \mnem{CLZ} \dst{r0}, \src{r1}\imod{.h01}, \src{r2}} \\[2pt]
\texttt{3} & \texttt{\imod{h11}} & v9--v14 & 16-bit halfword swizzle \texttt{h11}: lane 0 $\leftarrow$ h1, lane 1 $\leftarrow$ h1. \newline {\scriptsize\ttfamily \mnem{CLZ} \dst{r0}, \src{r1}\imod{.h11}, \src{r2}} \\[2pt]
\bottomrule
\end{longtable}
\end{center}
\par\textbf{Combined Walkthroughs:}
\par\textbf{Example 1 (Swap + broadcast in v2f16 add):}
\codebox{\mnem{FADD.v2f16} \dst{r0}, \src{r1}\imod{.h10}, \src{r2}\imod{.h00}}
With \texttt{r1 = [2.0h, 1.0h]} (\texttt{0x40003C00}), \texttt{r2 = [4.0h, 0.5h]} (\texttt{0x44003800}): \texttt{r1.h10 = [1.0h, 2.0h]}, \texttt{r2.h00 = [0.5h, 0.5h]} -\textgreater{} \texttt{r0 = [1.5h, 2.5h]}.
\par\textbf{Example 2 (Broadcast high lane with negation):}
\codebox{\mnem{FMUL.v2f16} \dst{r0}, \src{r1}\imod{.h01}, \src{r2}\imod{.h11.neg}}
With \texttt{r1 = [3.0h, 2.0h]}, \texttt{r2 = [4.0h, 1.0h]}: \texttt{r2.h11} broadcasts \texttt{4.0h} and \texttt{.neg} gives \texttt{[-4.0h, -4.0h]} -\textgreater{} \texttt{r0 = [-12.0h, -8.0h]}.

\section{Integer Source Widen (widen\_int\_m)}
\label{enum:widen_int_m}
Sign-extends (\texttt{s32}/\texttt{s64}) or zero-extends (\texttt{u32}/\texttt{u64}) narrow byte/half/word lanes. On 64-bit ops, byte/half swizzles fold low 32 bits then sign-extend to 64. Default: \texttt{none}.

\begin{center}\footnotesize
\begin{longtable}{@{}rll>{\raggedright\arraybackslash}p{0.62\linewidth}@{}}
\toprule
Enc & Value & Arch & Description \& Worked Example \\
\midrule\endhead
\texttt{0} & \texttt{\imod{none}} & v9--v14 & Use full 32-bit operand without narrowing/widening. \newline {\scriptsize\ttfamily \mnem{AND} \dst{r0}, \src{r1}\imod{}} \\[2pt]
\texttt{2} & \texttt{\imod{h0}} & v9--v14 & Select low 16-bit halfword \texttt{[15:0]} (widen to 32-bit if FP/int widen). \newline {\scriptsize\ttfamily \mnem{AND} \dst{r0}, \src{r1}\imod{.h0}} \\[2pt]
\texttt{3} & \texttt{\imod{h1}} & v9--v14 & Select high 16-bit halfword \texttt{[31:16]} (widen to 32-bit if FP/int widen). \newline {\scriptsize\ttfamily \mnem{AND} \dst{r0}, \src{r1}\imod{.h1}} \\[2pt]
\texttt{4} & \texttt{\imod{b0}} & v9--v14 & Select byte 0 \texttt{[7:0]} (zero/sign-extended to lane width). \newline {\scriptsize\ttfamily \mnem{AND} \dst{r0}, \src{r1}\imod{.b0}} \\[2pt]
\texttt{5} & \texttt{\imod{b1}} & v9--v14 & Select byte 1 \texttt{[15:8]} (zero/sign-extended to lane width). \newline {\scriptsize\ttfamily \mnem{AND} \dst{r0}, \src{r1}\imod{.b1}} \\[2pt]
\texttt{6} & \texttt{\imod{b2}} & v9--v14 & Select byte 2 \texttt{[23:16]} (zero/sign-extended to lane width). \newline {\scriptsize\ttfamily \mnem{AND} \dst{r0}, \src{r1}\imod{.b2}} \\[2pt]
\texttt{7} & \texttt{\imod{b3}} & v9--v14 & Select byte 3 \texttt{[31:24]} (zero/sign-extended to lane width). \newline {\scriptsize\ttfamily \mnem{AND} \dst{r0}, \src{r1}\imod{.b3}} \\[2pt]
\texttt{8} & \texttt{\imod{w0}} & v9--v14 & Select 32-bit word 0 \texttt{[31:0]} (extend to 64-bit). \newline {\scriptsize\ttfamily \mnem{AND} \dst{r0}, \src{r1}\imod{.w0}} \\[2pt]
\texttt{9} & \texttt{\imod{w1}} & v9--v14 & Select 32-bit word 1 \texttt{[63:32]} (extend to 64-bit). \newline {\scriptsize\ttfamily \mnem{AND} \dst{r0}, \src{r1}\imod{.w1}} \\[2pt]
\bottomrule
\end{longtable}
\end{center}
\par\textbf{Combined Walkthroughs:}
\par\textbf{Example 1 (Signed byte widen on IADD.s32):}
\codebox{\mnem{IADD.s32} \dst{r0}, \src{r1}\imod{.b0}, \src{r2}}
With \texttt{r1 = 0x000000FE} (\texttt{-2} in \texttt{b0}), \texttt{r2 = 10}: Stage A sign-extends \texttt{b0 -\textgreater{} 0xFFFFFFFE}; Stage B computes \texttt{r0 = 8}.
\par\textbf{Example 2 (Unsigned word widen on IMUL.u64):}
\codebox{\mnem{IMUL.u64} \dst{[r0:r1]}, \src{[r2:r3].w0}, \src{[r4:r5].w0}}
With \texttt{r2 = 0x80000000, r4 = 4}: zero-extends \texttt{.w0} to \texttt{0x00000000\_80000000} (not sign-extended) -\textgreater{} \texttt{[r0:r1] = 0x00000002\_00000000}.

\section{16-Bit Integer Swizzle (swiz\_int\_m)}
\label{enum:swiz_int_m}
Halfword swizzle or 8-to-16 byte widening (\texttt{bXY}: sign-extends on \texttt{v2s16}, zero-extends on \texttt{v2u16}). Default: \texttt{h01}.

\begin{center}\footnotesize
\begin{longtable}{@{}rll>{\raggedright\arraybackslash}p{0.62\linewidth}@{}}
\toprule
Enc & Value & Arch & Description \& Worked Example \\
\midrule\endhead
\texttt{0} & \texttt{\imod{h00}} & v9--v14 & 16-bit halfword swizzle \texttt{h00}: lane 0 $\leftarrow$ h0, lane 1 $\leftarrow$ h0. \newline {\scriptsize\ttfamily \mnem{AND} \dst{r0}, \src{r1}\imod{.h00}, \src{r2}} \\[2pt]
\texttt{1} & \texttt{\imod{h10}} & v9--v14 & 16-bit halfword swizzle \texttt{h10}: lane 0 $\leftarrow$ h1, lane 1 $\leftarrow$ h0. \newline {\scriptsize\ttfamily \mnem{AND} \dst{r0}, \src{r1}\imod{.h10}, \src{r2}} \\[2pt]
\texttt{2} & \texttt{\imod{h01}} & v9--v14 & 16-bit halfword swizzle \texttt{h01}: lane 0 $\leftarrow$ h0, lane 1 $\leftarrow$ h1. \newline {\scriptsize\ttfamily \mnem{AND} \dst{r0}, \src{r1}\imod{.h01}, \src{r2}} \\[2pt]
\texttt{3} & \texttt{\imod{h11}} & v9--v14 & 16-bit halfword swizzle \texttt{h11}: lane 0 $\leftarrow$ h1, lane 1 $\leftarrow$ h1. \newline {\scriptsize\ttfamily \mnem{AND} \dst{r0}, \src{r1}\imod{.h11}, \src{r2}} \\[2pt]
\texttt{4} & \texttt{\imod{b00}} & v9--v14 & 2-byte swizzle \texttt{b00}: low byte $\leftarrow$ b0, high byte $\leftarrow$ b0. \newline {\scriptsize\ttfamily \mnem{AND} \dst{r0}, \src{r1}\imod{.b00}, \src{r2}} \\[2pt]
\texttt{5} & \texttt{\imod{b20}} & v9--v14 & 2-byte swizzle \texttt{b20}: low byte $\leftarrow$ b2, high byte $\leftarrow$ b0. \newline {\scriptsize\ttfamily \mnem{AND} \dst{r0}, \src{r1}\imod{.b20}, \src{r2}} \\[2pt]
\texttt{6} & \texttt{\imod{b02}} & v9--v14 & 2-byte swizzle \texttt{b02}: low byte $\leftarrow$ b0, high byte $\leftarrow$ b2. \newline {\scriptsize\ttfamily \mnem{AND} \dst{r0}, \src{r1}\imod{.b02}, \src{r2}} \\[2pt]
\texttt{7} & \texttt{\imod{b22}} & v9--v14 & 2-byte swizzle \texttt{b22}: low byte $\leftarrow$ b2, high byte $\leftarrow$ b2. \newline {\scriptsize\ttfamily \mnem{AND} \dst{r0}, \src{r1}\imod{.b22}, \src{r2}} \\[2pt]
\texttt{8} & \texttt{\imod{b11}} & v9--v14 & 2-byte swizzle \texttt{b11}: low byte $\leftarrow$ b1, high byte $\leftarrow$ b1. \newline {\scriptsize\ttfamily \mnem{AND} \dst{r0}, \src{r1}\imod{.b11}, \src{r2}} \\[2pt]
\texttt{9} & \texttt{\imod{b31}} & v9--v14 & 2-byte swizzle \texttt{b31}: low byte $\leftarrow$ b3, high byte $\leftarrow$ b1. \newline {\scriptsize\ttfamily \mnem{AND} \dst{r0}, \src{r1}\imod{.b31}, \src{r2}} \\[2pt]
\texttt{10} & \texttt{\imod{b13}} & v9--v14 & 2-byte swizzle \texttt{b13}: low byte $\leftarrow$ b1, high byte $\leftarrow$ b3. \newline {\scriptsize\ttfamily \mnem{AND} \dst{r0}, \src{r1}\imod{.b13}, \src{r2}} \\[2pt]
\texttt{11} & \texttt{\imod{b33}} & v9--v14 & 2-byte swizzle \texttt{b33}: low byte $\leftarrow$ b3, high byte $\leftarrow$ b3. \newline {\scriptsize\ttfamily \mnem{AND} \dst{r0}, \src{r1}\imod{.b33}, \src{r2}} \\[2pt]
\texttt{12} & \texttt{\imod{b01}} & v9--v14 & 2-byte swizzle \texttt{b01}: low byte $\leftarrow$ b0, high byte $\leftarrow$ b1. \newline {\scriptsize\ttfamily \mnem{AND} \dst{r0}, \src{r1}\imod{.b01}, \src{r2}} \\[2pt]
\texttt{13} & \texttt{\imod{b23}} & v9--v14 & 2-byte swizzle \texttt{b23}: low byte $\leftarrow$ b2, high byte $\leftarrow$ b3. \newline {\scriptsize\ttfamily \mnem{AND} \dst{r0}, \src{r1}\imod{.b23}, \src{r2}} \\[2pt]
\bottomrule
\end{longtable}
\end{center}
\par\textbf{Combined Walkthroughs:}
\par\textbf{Example 1 (Signed vs unsigned byte-pair widen (`b02`)):}
\codebox{\mnem{IADD.v2s16} \dst{r0}, \src{r1}\imod{.b02}, \src{k0}}
With \texttt{r1 = 0x00FE00FF} (\texttt{b0 = -1, b2 = -2}): \texttt{v2s16} sign-extends both bytes -\textgreater{} \texttt{r0 = 0xFFFEFFFF} (whereas \texttt{v2u16} yields \texttt{0x00FE00FF}).
\par\textbf{Example 2 (Swap halfwords (`h10`) before subtraction):}
\codebox{\mnem{ISUB.v2s16} \dst{r0}, \src{r1}\imod{.h10}, \src{r2}\imod{.h01}}
With \texttt{r1 = [10, 3]} (\texttt{0x000A0003}), \texttt{r2 = [2, 1]} (\texttt{0x00020001}): \texttt{r1.h10 = [3, 10]}, subtracting \texttt{[2, 1]} -\textgreater{} \texttt{r0 = [1, 9]} (\texttt{0x00010009}).

\section{2-Byte Full Swizzle (swiz8\_2\_full\_m)}
\label{enum:swiz8_2_full_m}
Extracts any 2 of 4 bytes into \texttt{[7:0]} and \texttt{[23:16]}. Default \texttt{b02} on \texttt{V2*8\_TO\_*}; virtual fixed \texttt{b01} on \texttt{MKVEC.v2i8}.

\begin{center}\footnotesize
\begin{longtable}{@{}rll>{\raggedright\arraybackslash}p{0.62\linewidth}@{}}
\toprule
Enc & Value & Arch & Description \& Worked Example \\
\midrule\endhead
\texttt{0} & \texttt{\imod{b00}} & v9--v14 & 2-byte swizzle \texttt{b00}: low byte $\leftarrow$ b0, high byte $\leftarrow$ b0. \newline {\scriptsize\ttfamily \mnem{MKVEC} \dst{r0}, \src{r1}\imod{.b00}, \src{r2}} \\[2pt]
\texttt{1} & \texttt{\imod{b10}} & v9--v14 & 2-byte swizzle \texttt{b10}: low byte $\leftarrow$ b1, high byte $\leftarrow$ b0. \newline {\scriptsize\ttfamily \mnem{MKVEC} \dst{r0}, \src{r1}\imod{.b10}, \src{r2}} \\[2pt]
\texttt{2} & \texttt{\imod{b20}} & v9--v14 & 2-byte swizzle \texttt{b20}: low byte $\leftarrow$ b2, high byte $\leftarrow$ b0. \newline {\scriptsize\ttfamily \mnem{MKVEC} \dst{r0}, \src{r1}\imod{.b20}, \src{r2}} \\[2pt]
\texttt{3} & \texttt{\imod{b30}} & v9--v14 & 2-byte swizzle \texttt{b30}: low byte $\leftarrow$ b3, high byte $\leftarrow$ b0. \newline {\scriptsize\ttfamily \mnem{MKVEC} \dst{r0}, \src{r1}\imod{.b30}, \src{r2}} \\[2pt]
\texttt{4} & \texttt{\imod{b01}} & v9--v14 & 2-byte swizzle \texttt{b01}: low byte $\leftarrow$ b0, high byte $\leftarrow$ b1. \newline {\scriptsize\ttfamily \mnem{MKVEC} \dst{r0}, \src{r1}\imod{.b01}, \src{r2}} \\[2pt]
\texttt{5} & \texttt{\imod{b11}} & v9--v14 & 2-byte swizzle \texttt{b11}: low byte $\leftarrow$ b1, high byte $\leftarrow$ b1. \newline {\scriptsize\ttfamily \mnem{MKVEC} \dst{r0}, \src{r1}\imod{.b11}, \src{r2}} \\[2pt]
\texttt{6} & \texttt{\imod{b21}} & v9--v14 & 2-byte swizzle \texttt{b21}: low byte $\leftarrow$ b2, high byte $\leftarrow$ b1. \newline {\scriptsize\ttfamily \mnem{MKVEC} \dst{r0}, \src{r1}\imod{.b21}, \src{r2}} \\[2pt]
\texttt{7} & \texttt{\imod{b31}} & v9--v14 & 2-byte swizzle \texttt{b31}: low byte $\leftarrow$ b3, high byte $\leftarrow$ b1. \newline {\scriptsize\ttfamily \mnem{MKVEC} \dst{r0}, \src{r1}\imod{.b31}, \src{r2}} \\[2pt]
\texttt{8} & \texttt{\imod{b02}} & v9--v14 & 2-byte swizzle \texttt{b02}: low byte $\leftarrow$ b0, high byte $\leftarrow$ b2. \newline {\scriptsize\ttfamily \mnem{MKVEC} \dst{r0}, \src{r1}\imod{.b02}, \src{r2}} \\[2pt]
\texttt{9} & \texttt{\imod{b12}} & v9--v14 & 2-byte swizzle \texttt{b12}: low byte $\leftarrow$ b1, high byte $\leftarrow$ b2. \newline {\scriptsize\ttfamily \mnem{MKVEC} \dst{r0}, \src{r1}\imod{.b12}, \src{r2}} \\[2pt]
\texttt{10} & \texttt{\imod{b22}} & v9--v14 & 2-byte swizzle \texttt{b22}: low byte $\leftarrow$ b2, high byte $\leftarrow$ b2. \newline {\scriptsize\ttfamily \mnem{MKVEC} \dst{r0}, \src{r1}\imod{.b22}, \src{r2}} \\[2pt]
\texttt{11} & \texttt{\imod{b32}} & v9--v14 & 2-byte swizzle \texttt{b32}: low byte $\leftarrow$ b3, high byte $\leftarrow$ b2. \newline {\scriptsize\ttfamily \mnem{MKVEC} \dst{r0}, \src{r1}\imod{.b32}, \src{r2}} \\[2pt]
\texttt{12} & \texttt{\imod{b03}} & v9--v14 & 2-byte swizzle \texttt{b03}: low byte $\leftarrow$ b0, high byte $\leftarrow$ b3. \newline {\scriptsize\ttfamily \mnem{MKVEC} \dst{r0}, \src{r1}\imod{.b03}, \src{r2}} \\[2pt]
\texttt{13} & \texttt{\imod{b13}} & v9--v14 & 2-byte swizzle \texttt{b13}: low byte $\leftarrow$ b1, high byte $\leftarrow$ b3. \newline {\scriptsize\ttfamily \mnem{MKVEC} \dst{r0}, \src{r1}\imod{.b13}, \src{r2}} \\[2pt]
\texttt{14} & \texttt{\imod{b23}} & v9--v14 & 2-byte swizzle \texttt{b23}: low byte $\leftarrow$ b2, high byte $\leftarrow$ b3. \newline {\scriptsize\ttfamily \mnem{MKVEC} \dst{r0}, \src{r1}\imod{.b23}, \src{r2}} \\[2pt]
\texttt{15} & \texttt{\imod{b33}} & v9--v14 & 2-byte swizzle \texttt{b33}: low byte $\leftarrow$ b3, high byte $\leftarrow$ b3. \newline {\scriptsize\ttfamily \mnem{MKVEC} \dst{r0}, \src{r1}\imod{.b33}, \src{r2}} \\[2pt]
\bottomrule
\end{longtable}
\end{center}
\par\textbf{Combined Walkthroughs:}
\par\textbf{Example 1 (Cross-lane byte swap (`b20`) on signed widen):}
\codebox{\mnem{V2S8\_TO\_V2S16} \dst{r0}, \src{r1}\imod{.b20}}
With \texttt{r1 = 0x000500FE}: \texttt{b20} puts \texttt{b2 = +5} in lane 0 and \texttt{b0 = -2} in lane 1 -\textgreater{} \texttt{r0 = 0xFFFE0005}.
\par\textbf{Example 2 (Duplicate high byte (`b33`) on unsigned widen):}
\codebox{\mnem{V2U8\_TO\_V2U16} \dst{r0}, \src{r1}\imod{.b33}}
With \texttt{r1 = 0xAB123456}: \texttt{b33} zero-extends byte 3 (\texttt{0xAB}) into both lanes -\textgreater{} \texttt{r0 = 0x00AB00AB}.

\section{8-Bit Vector Lanes (lanes\_int\_m)}
\label{enum:lanes_int_m}
4-byte permutation for \texttt{v4s8}/\texttt{v4u8}/\texttt{v4i8} (arch v9--v10). Default: \texttt{b0123}.

\begin{center}\footnotesize
\begin{longtable}{@{}rll>{\raggedright\arraybackslash}p{0.62\linewidth}@{}}
\toprule
Enc & Value & Arch & Description \& Worked Example \\
\midrule\endhead
\texttt{0} & \texttt{\imod{b0123}} & v9--v14 & 4-byte swizzle \texttt{b0123}: output bytes [0..3] $\leftarrow$ input bytes [0, 1, 2, 3]. \newline {\scriptsize\ttfamily \mnem{AND} \dst{r0}, \src{r1}\imod{.b0123}, \src{r2}} \\[2pt]
\texttt{1} & \texttt{\imod{b3210}} & v9--v14 & 4-byte swizzle \texttt{b3210}: output bytes [0..3] $\leftarrow$ input bytes [3, 2, 1, 0]. \newline {\scriptsize\ttfamily \mnem{AND} \dst{r0}, \src{r1}\imod{.b3210}, \src{r2}} \\[2pt]
\texttt{2} & \texttt{\imod{b0101}} & v9--v14 & 4-byte swizzle \texttt{b0101}: output bytes [0..3] $\leftarrow$ input bytes [0, 1, 0, 1]. \newline {\scriptsize\ttfamily \mnem{AND} \dst{r0}, \src{r1}\imod{.b0101}, \src{r2}} \\[2pt]
\texttt{3} & \texttt{\imod{b2323}} & v9--v14 & 4-byte swizzle \texttt{b2323}: output bytes [0..3] $\leftarrow$ input bytes [2, 3, 2, 3]. \newline {\scriptsize\ttfamily \mnem{AND} \dst{r0}, \src{r1}\imod{.b2323}, \src{r2}} \\[2pt]
\texttt{4} & \texttt{\imod{b0000}} & v9--v14 & 4-byte swizzle \texttt{b0000}: output bytes [0..3] $\leftarrow$ input bytes [0, 0, 0, 0]. \newline {\scriptsize\ttfamily \mnem{AND} \dst{r0}, \src{r1}\imod{.b0000}, \src{r2}} \\[2pt]
\texttt{5} & \texttt{\imod{b1111}} & v9--v14 & 4-byte swizzle \texttt{b1111}: output bytes [0..3] $\leftarrow$ input bytes [1, 1, 1, 1]. \newline {\scriptsize\ttfamily \mnem{AND} \dst{r0}, \src{r1}\imod{.b1111}, \src{r2}} \\[2pt]
\texttt{6} & \texttt{\imod{b2222}} & v9--v14 & 4-byte swizzle \texttt{b2222}: output bytes [0..3] $\leftarrow$ input bytes [2, 2, 2, 2]. \newline {\scriptsize\ttfamily \mnem{AND} \dst{r0}, \src{r1}\imod{.b2222}, \src{r2}} \\[2pt]
\texttt{7} & \texttt{\imod{b3333}} & v9--v14 & 4-byte swizzle \texttt{b3333}: output bytes [0..3] $\leftarrow$ input bytes [3, 3, 3, 3]. \newline {\scriptsize\ttfamily \mnem{AND} \dst{r0}, \src{r1}\imod{.b3333}, \src{r2}} \\[2pt]
\texttt{8} & \texttt{\imod{b2301}} & v9--v14 & 4-byte swizzle \texttt{b2301}: output bytes [0..3] $\leftarrow$ input bytes [2, 3, 0, 1]. \newline {\scriptsize\ttfamily \mnem{AND} \dst{r0}, \src{r1}\imod{.b2301}, \src{r2}} \\[2pt]
\texttt{9} & \texttt{\imod{b1032}} & v9--v14 & 4-byte swizzle \texttt{b1032}: output bytes [0..3] $\leftarrow$ input bytes [1, 0, 3, 2]. \newline {\scriptsize\ttfamily \mnem{AND} \dst{r0}, \src{r1}\imod{.b1032}, \src{r2}} \\[2pt]
\texttt{10} & \texttt{\imod{b0011}} & v9--v14 & 4-byte swizzle \texttt{b0011}: output bytes [0..3] $\leftarrow$ input bytes [0, 0, 1, 1]. \newline {\scriptsize\ttfamily \mnem{AND} \dst{r0}, \src{r1}\imod{.b0011}, \src{r2}} \\[2pt]
\texttt{11} & \texttt{\imod{b2233}} & v9--v14 & 4-byte swizzle \texttt{b2233}: output bytes [0..3] $\leftarrow$ input bytes [2, 2, 3, 3]. \newline {\scriptsize\ttfamily \mnem{AND} \dst{r0}, \src{r1}\imod{.b2233}, \src{r2}} \\[2pt]
\bottomrule
\end{longtable}
\end{center}
\par\textbf{Combined Walkthroughs:}
\par\textbf{Example 1 (Endianness byte reversal (`b3210`)):}
\codebox{\mnem{SWZ.v4i8} \dst{r0}, \src{r1}\imod{.b3210}}
With \texttt{r1 = 0x11223344}: \texttt{b3210} reverses all 4 bytes -\textgreater{} \texttt{r0 = 0x44332211}.
\par\textbf{Example 2 (Byte broadcast (`b0000`) in packed add):}
\codebox{\mnem{IADD.v4u8}\imod{.sat} \dst{r0}, \src{r1}\imod{.b0123}, \src{r2}\imod{.b0000}}
With \texttt{r1 = 0xF0201000}, \texttt{r2.b0 = 0x20}: broadcasts \texttt{0x20} to all 4 bytes, saturating byte 3 (\texttt{0xF0+0x20 -\textgreater{} 0xFF}) -\textgreater{} \texttt{r0 = 0xFF403020}.

\section{8-Bit Single Lane (lane8\_m)}
\label{enum:lane8_m}
Extracts a single byte \texttt{b0..b3}. On \texttt{v4*8} replicates across 4 lanes; on \texttt{v2*16} widens and replicates to 2x16. Default: \texttt{b0}.

\begin{center}\footnotesize
\begin{longtable}{@{}rll>{\raggedright\arraybackslash}p{0.62\linewidth}@{}}
\toprule
Enc & Value & Arch & Description \& Worked Example \\
\midrule\endhead
\texttt{0} & \texttt{\imod{b0}} & v9--v14 & Select byte 0 \texttt{[7:0]} (zero/sign-extended to lane width). \newline {\scriptsize\ttfamily \mnem{MKVEC} \dst{r0}, \src{r1}\imod{.b0}} \\[2pt]
\texttt{1} & \texttt{\imod{b1}} & v9--v14 & Select byte 1 \texttt{[15:8]} (zero/sign-extended to lane width). \newline {\scriptsize\ttfamily \mnem{MKVEC} \dst{r0}, \src{r1}\imod{.b1}} \\[2pt]
\texttt{2} & \texttt{\imod{b2}} & v9--v14 & Select byte 2 \texttt{[23:16]} (zero/sign-extended to lane width). \newline {\scriptsize\ttfamily \mnem{MKVEC} \dst{r0}, \src{r1}\imod{.b2}} \\[2pt]
\texttt{3} & \texttt{\imod{b3}} & v9--v14 & Select byte 3 \texttt{[31:24]} (zero/sign-extended to lane width). \newline {\scriptsize\ttfamily \mnem{MKVEC} \dst{r0}, \src{r1}\imod{.b3}} \\[2pt]
\bottomrule
\end{longtable}
\end{center}
\par\textbf{Combined Walkthroughs:}
\par\textbf{Example 1 (Extract byte 2 and zero-extend):}
\codebox{\mnem{U8\_TO\_U32} \dst{r0}, \src{r1}\imod{.b2}}
With \texttt{r1 = 0x00AB0000}: extracts byte 2 (\texttt{0xAB}) -\textgreater{} \texttt{r0 = 0x000000AB}.
\par\textbf{Example 2 (Extract byte 3 and sign-extend):}
\codebox{\mnem{S8\_TO\_S32} \dst{r0}, \src{r1}\imod{.b3}}
With \texttt{r1 = 0xFE000000}: extracts byte 3 (\texttt{-2}) and sign-extends -\textgreater{} \texttt{r0 = 0xFFFFFFFE}.

\section{16-Bit Single Lane (lane16\_m)}
\label{enum:lane16_m}
Extracts a single 16-bit halfword \texttt{h0} or \texttt{h1}. Default: \texttt{h0}.

\begin{center}\footnotesize
\begin{longtable}{@{}rll>{\raggedright\arraybackslash}p{0.62\linewidth}@{}}
\toprule
Enc & Value & Arch & Description \& Worked Example \\
\midrule\endhead
\texttt{0} & \texttt{\imod{h0}} & v9--v14 & Select low 16-bit halfword \texttt{[15:0]} (widen to 32-bit if FP/int widen). \newline {\scriptsize\ttfamily \mnem{F16\_TO\_F32} \dst{r0}, \src{r1}\imod{.h0}} \\[2pt]
\texttt{1} & \texttt{\imod{h1}} & v9--v14 & Select high 16-bit halfword \texttt{[31:16]} (widen to 32-bit if FP/int widen). \newline {\scriptsize\ttfamily \mnem{F16\_TO\_F32} \dst{r0}, \src{r1}\imod{.h1}} \\[2pt]
\bottomrule
\end{longtable}
\end{center}
\par\textbf{Combined Walkthroughs:}
\par\textbf{Example 1 (Extract upper FP16 halfword (`h1`)):}
\codebox{\mnem{F16\_TO\_F32} \dst{r0}, \src{r1}\imod{.h1}}
With \texttt{r1 = 0x40003C00}: extracts \texttt{h1 = 0x4000} (\texttt{2.0h}) -\textgreater{} \texttt{r0 = 0x40000000} (\texttt{2.0f}).
\par\textbf{Example 2 (Extract low signed 16-bit halfword (`h0`)):}
\codebox{\mnem{S16\_TO\_S32} \dst{r0}, \src{r1}\imod{.h0}}
With \texttt{r1 = 0x1234FFFC}: extracts \texttt{h0 = 0xFFFC} (\texttt{-4}) and sign-extends -\textgreater{} \texttt{r0 = 0xFFFFFFFC}.

\section{32/64-Bit Shift Amount Lane (shift\_bitwise\_i32\_lanes\_m)}
\label{enum:shift_bitwise_i32_lanes_m}
Selects unsigned 8-bit shift count byte from \texttt{src1} (masked modulo lane width: 32 or 64). Default: \texttt{b0}.

\begin{center}\footnotesize
\begin{longtable}{@{}rll>{\raggedright\arraybackslash}p{0.62\linewidth}@{}}
\toprule
Enc & Value & Arch & Description \& Worked Example \\
\midrule\endhead
\texttt{4} & \texttt{\imod{b0}} & v9--v14 & Single-byte extract \texttt{b0}: select byte 0 zero/sign-extended to 32 bits. \newline {\scriptsize\ttfamily \mnem{AND} \dst{r0}, \src{r1}\imod{.b0}, \src{r2}} \\[2pt]
\texttt{5} & \texttt{\imod{b1}} & v9--v14 & Single-byte extract \texttt{b1}: select byte 1 zero/sign-extended to 32 bits. \newline {\scriptsize\ttfamily \mnem{AND} \dst{r0}, \src{r1}\imod{.b1}, \src{r2}} \\[2pt]
\texttt{6} & \texttt{\imod{b2}} & v9--v14 & Single-byte extract \texttt{b2}: select byte 2 zero/sign-extended to 32 bits. \newline {\scriptsize\ttfamily \mnem{AND} \dst{r0}, \src{r1}\imod{.b2}, \src{r2}} \\[2pt]
\texttt{7} & \texttt{\imod{b3}} & v9--v14 & Single-byte extract \texttt{b3}: select byte 3 zero/sign-extended to 32 bits. \newline {\scriptsize\ttfamily \mnem{AND} \dst{r0}, \src{r1}\imod{.b3}, \src{r2}} \\[2pt]
\bottomrule
\end{longtable}
\end{center}
\par\textbf{Combined Walkthroughs:}
\par\textbf{Example 1 (Read shift amount from byte 1 (`b1`)):}
\codebox{\mnem{LSHIFT.i32} \dst{r0}, \src{r1}, \src{r2}\imod{.b1}}
With \texttt{r1 = 1, r2 = 0x00000300}: reads \texttt{b1 = 3} -\textgreater{} \texttt{r0 = 1 \textless{}\textless{} 3 = 8}.
\par\textbf{Example 2 (Modulo-32 wrap on shift count \textgreater{}= 32):}
\codebox{\mnem{RSHIFT.i32} \dst{r0}, \src{r1}, \src{r2}\imod{.b0}}
With \texttt{r1 = 0x80000000, r2.b0 = 33}: hardware masks \texttt{33 \& 31 = 1} -\textgreater{} \texttt{r0 = 0x40000000} (NOT 0).

\section{16-Bit Shift Amount Lanes (shift\_bitwise\_i16\_lanes\_m)}
\label{enum:shift_bitwise_i16_lanes_m}
Selects two unsigned 8-bit shift counts (lane 0 \& lane 1, masked mod 16) on \texttt{v2i16} shifts. Default: \texttt{b02}.

\begin{center}\footnotesize
\begin{longtable}{@{}rll>{\raggedright\arraybackslash}p{0.62\linewidth}@{}}
\toprule
Enc & Value & Arch & Description \& Worked Example \\
\midrule\endhead
\texttt{0} & \texttt{\imod{b02}} & v9--v14 & 2-byte swizzle \texttt{b02}: low byte $\leftarrow$ b0, high byte $\leftarrow$ b2. \newline {\scriptsize\ttfamily \mnem{AND} \dst{r0}, \src{r1}\imod{.b02}, \src{r2}} \\[2pt]
\texttt{4} & \texttt{\imod{b00}} & v9--v14 & 2-byte swizzle \texttt{b00}: low byte $\leftarrow$ b0, high byte $\leftarrow$ b0. \newline {\scriptsize\ttfamily \mnem{AND} \dst{r0}, \src{r1}\imod{.b00}, \src{r2}} \\[2pt]
\texttt{5} & \texttt{\imod{b11}} & v9--v14 & 2-byte swizzle \texttt{b11}: low byte $\leftarrow$ b1, high byte $\leftarrow$ b1. \newline {\scriptsize\ttfamily \mnem{AND} \dst{r0}, \src{r1}\imod{.b11}, \src{r2}} \\[2pt]
\texttt{6} & \texttt{\imod{b22}} & v9--v14 & 2-byte swizzle \texttt{b22}: low byte $\leftarrow$ b2, high byte $\leftarrow$ b2. \newline {\scriptsize\ttfamily \mnem{AND} \dst{r0}, \src{r1}\imod{.b22}, \src{r2}} \\[2pt]
\texttt{7} & \texttt{\imod{b33}} & v9--v14 & 2-byte swizzle \texttt{b33}: low byte $\leftarrow$ b3, high byte $\leftarrow$ b3. \newline {\scriptsize\ttfamily \mnem{AND} \dst{r0}, \src{r1}\imod{.b33}, \src{r2}} \\[2pt]
\texttt{10} & \texttt{\imod{b01}} & v9--v14 & 2-byte swizzle \texttt{b01}: low byte $\leftarrow$ b0, high byte $\leftarrow$ b1. \newline {\scriptsize\ttfamily \mnem{AND} \dst{r0}, \src{r1}\imod{.b01}, \src{r2}} \\[2pt]
\texttt{11} & \texttt{\imod{b23}} & v9--v14 & 2-byte swizzle \texttt{b23}: low byte $\leftarrow$ b2, high byte $\leftarrow$ b3. \newline {\scriptsize\ttfamily \mnem{AND} \dst{r0}, \src{r1}\imod{.b23}, \src{r2}} \\[2pt]
\bottomrule
\end{longtable}
\end{center}
\par\textbf{Combined Walkthroughs:}
\par\textbf{Example 1 (Independent per-lane shift counts (`b02`)):}
\codebox{\mnem{LSHIFT.v2i16} \dst{r0}, \src{r1}, \src{r2}\imod{.b02}}
With \texttt{r1 = 0x00010001, r2 = 0x00110010} (\texttt{b0=16 -\textgreater{} 0}, \texttt{b2=17 -\textgreater{} 1}): lane 0 shifts by 0, lane 1 shifts by 1 -\textgreater{} \texttt{r0 = 0x00020001}.
\par\textbf{Example 2 (Shared shift count (`b00`)):}
\codebox{\mnem{RSHIFT.v2i16} \dst{r0}, \src{r1}, \src{r2}\imod{.b00}}
With \texttt{r1 = 0x00800040, r2.b0 = 2}: shifts both 16-bit lanes right by 2 -\textgreater{} \texttt{r0 = 0x00200010}.

\section{Pre-Source Floating-Point Abs (abs\_m)}
\label{enum:abs_m}
Clears IEEE sign bit(s) after swizzle (\texttt{\& 0x7FFFFFFF} f32, \texttt{\& 0x7FFF7FFF} f16/v2f16). Default: \texttt{none}.

\begin{center}\footnotesize
\begin{longtable}{@{}rll>{\raggedright\arraybackslash}p{0.62\linewidth}@{}}
\toprule
Enc & Value & Arch & Description \& Worked Example \\
\midrule\endhead
\texttt{0} & \texttt{\imod{none}} & v9--v14 & Default behavior (flag disabled for \texttt{abs\_m}). \newline {\scriptsize\ttfamily \mnem{CUBEFACE1}\imod{} \dst{r0}, \src{r1}} \\[2pt]
\texttt{1} & \texttt{\imod{abs}} & v9--v14 & Clear sign bit ($|x|$) before operation (after widen, before neg). \newline {\scriptsize\ttfamily \mnem{CUBEFACE1}\imod{.abs} \dst{r0}, \src{r1}} \\[2pt]
\bottomrule
\end{longtable}
\end{center}
\par\textbf{Combined Walkthroughs:}
\par\textbf{Example 1 (Pre-source `.abs` on FP32 add):}
\codebox{\mnem{FADD.f32} \dst{r0}, \src{r1}\imod{.abs}, \src{r2}}
With \texttt{r1 = -3.0f, r2 = 1.0f}: Stage A clears sign bit of \texttt{r1} (\texttt{+3.0f}) -\textgreater{} \texttt{r0 = 4.0f}.
\par\textbf{Example 2 (Combined `.abs.neg` (always negative magnitude)):}
\codebox{\mnem{FADD.f32} \dst{r0}, \src{r1}\imod{.abs.neg}, \src{r2}}
With \texttt{r1 = +3.0f, r2 = 10.0f}: \texttt{.abs} runs before \texttt{.neg}, forcing \texttt{-\textbar{}r1\textbar{} = -3.0f} -\textgreater{} \texttt{r0 = 7.0f}.

\section{Pre-Source Floating-Point Neg (neg\_m)}
\label{enum:neg_m}
Flips IEEE sign bit(s) after swizzle (\texttt{\textasciicircum{} 0x80000000} f32, \texttt{\textasciicircum{} 0x80008000} f16/v2f16). Default: \texttt{none}.

\begin{center}\footnotesize
\begin{longtable}{@{}rll>{\raggedright\arraybackslash}p{0.62\linewidth}@{}}
\toprule
Enc & Value & Arch & Description \& Worked Example \\
\midrule\endhead
\texttt{0} & \texttt{\imod{none}} & v9--v14 & Default behavior (flag disabled for \texttt{neg\_m}). \newline {\scriptsize\ttfamily \mnem{CUBEFACE1}\imod{} \dst{r0}, \src{r1}} \\[2pt]
\texttt{1} & \texttt{\imod{neg}} & v9--v14 & Flip sign bit ($-x$) before operation (after abs). \newline {\scriptsize\ttfamily \mnem{CUBEFACE1}\imod{.neg} \dst{r0}, \src{r1}} \\[2pt]
\bottomrule
\end{longtable}
\end{center}
\par\textbf{Combined Walkthroughs:}
\par\textbf{Example 1 (Floating-point subtraction via `src1.neg`):}
\codebox{\mnem{FADD.f32} \dst{r0}, \src{r1}, \src{r2}\imod{.neg}}
With \texttt{r1 = 5.0f, r2 = 2.0f}: Stage A negates \texttt{r2} (\texttt{-2.0f}) -\textgreater{} \texttt{r0 = 3.0f} (Valhall has no separate \texttt{FSUB} opcode).
\par\textbf{Example 2 (FMA negated product (`-a*b + c`)):}
\codebox{\mnem{FMA.f32} \dst{r0}, \src{r1}\imod{.neg}, \src{r2}, \src{r3}}
With \texttt{r1 = 3.0f, r2 = 4.0f, r3 = 20.0f}: Stage A negates \texttt{r1} to \texttt{-3.0f} -\textgreater{} \texttt{r0 = -12.0 + 20.0 = 8.0f}.

\section{Bitwise NOT Modifier (not\_m)}
\label{enum:not_m}
Bitwise complement \texttt{\textasciitilde{}x} on \texttt{src2} (\texttt{not2}) or result (\texttt{not\_result}) of shift-logic ops. Default: \texttt{none}.

\begin{center}\footnotesize
\begin{longtable}{@{}rll>{\raggedright\arraybackslash}p{0.62\linewidth}@{}}
\toprule
Enc & Value & Arch & Description \& Worked Example \\
\midrule\endhead
\texttt{0} & \texttt{\imod{none}} & v9--v14 & Default behavior (flag disabled for \texttt{not\_m}). \newline {\scriptsize\ttfamily \mnem{AND}\imod{} \dst{r0}, \src{r1}} \\[2pt]
\texttt{1} & \texttt{\imod{not}} & v9--v14 & Bitwise invert operand or negate branch predicate ($\sim x$ / $!p$). \newline {\scriptsize\ttfamily \mnem{AND}\imod{.not} \dst{r0}, \src{r1}} \\[2pt]
\bottomrule
\end{longtable}
\end{center}
\par\textbf{Combined Walkthroughs:}
\par\textbf{Example 1 (Combined `src2.not` and `.not\_result`):}
\codebox{\mnem{AND.i32}\imod{.not\_result} \dst{r0}, \src{r1}, \src{r2}\imod{.not}}
With \texttt{r1 = 0x0F, r2 = 0x03}: \texttt{\textasciitilde{}(0x0F \& \textasciitilde{}0x03) = \textasciitilde{}(0x0C) = 0xFFFFFFF3}.
\par\textbf{Example 2 (Bitwise BIC (`a \& \textasciitilde{}b`) without `.not\_result`):}
\codebox{\mnem{AND.i32} \dst{r0}, \src{r1}, \src{r2}\imod{.not}}
With \texttt{r1 = 0xFFFF00FF, r2 = 0x00FF000F}: clears masked bits -\textgreater{} \texttt{r0 = 0xFF0000F0}.

\section{Floating-Point Clamp (clamp\_m)}
\label{enum:clamp_m}
Post-op IEEE-754 clamp with deterministic \texttt{-0.0 \textless{} +0.0} and NaN suppression (\texttt{NaN -\textgreater{} lo}). Default: \texttt{none}.

\begin{center}\footnotesize
\begin{longtable}{@{}rll>{\raggedright\arraybackslash}p{0.62\linewidth}@{}}
\toprule
Enc & Value & Arch & Description \& Worked Example \\
\midrule\endhead
\texttt{0} & \texttt{\imod{none}} & v9--v14 & No output clamping; pass IEEE-754 result through unchanged. \newline {\scriptsize\ttfamily \mnem{FADD.f32} \dst{r0}, \src{r1}, \src{r2}} \\[2pt]
\texttt{1} & \texttt{\imod{clamp\_0\_inf}} & v9--v14 & Clamp to $[+0.0, +\infty]$; negative values, $-0.0$, and NaN become $+0.0$. \newline {\scriptsize\ttfamily \mnem{FADD.f32}\imod{.clamp\_0\_inf} \dst{r0}, \src{r1}, \src{r2}} \\[2pt]
\texttt{2} & \texttt{\imod{clamp\_m1\_1}} & v9--v14 & Clamp to $[-1.0, +1.0]$; NaN becomes $-1.0$ (\texttt{0xBF800000}) per DDK. \newline {\scriptsize\ttfamily \mnem{FMA.f32}\imod{.clamp\_m1\_1} \dst{r0}, \src{r1}, \src{r2}, \src{r3}} \\[2pt]
\texttt{3} & \texttt{\imod{clamp\_0\_1}} & v9--v14 & Clamp to $[+0.0, +1.0]$ (saturate); $-0.0$ and NaN become $+0.0$. \newline {\scriptsize\ttfamily \mnem{FMA.f32}\imod{.clamp\_0\_1} \dst{r0}, \src{r1}, \src{r2}, \src{r3}} \\[2pt]
\bottomrule
\end{longtable}
\end{center}
\par\textbf{Combined Walkthroughs:}
\par\textbf{Example 1 (`.clamp\_0\_1` on `-0.0` and `NaN`):}
\codebox{\mnem{FMA.f32}\imod{.clamp\_0\_1} \dst{r0}, \src{r1}, \src{r2}, \src{r3}}
With result \texttt{-0.0f} or \texttt{NaN}: \texttt{.clamp\_0\_1} clamps to \texttt{+0.0f} (\texttt{0x00000000}), never \texttt{-0.0f} or \texttt{NaN}.
\par\textbf{Example 2 (`.clamp\_m1\_1` NaN-to-low-bound rule):}
\codebox{\mnem{FADD.f32}\imod{.clamp\_m1\_1} \dst{r0}, \src{r1}, \src{r2}}
With result \texttt{NaN} or \texttt{-2.5f}: \texttt{.clamp\_m1\_1} returns \texttt{-1.0f} (\texttt{0xBF800000}); with \texttt{+3.0f}, returns \texttt{+1.0f} (\texttt{0x3F800000}).

\section{Floating-Point Rounding (round\_m)}
\label{enum:round_m}
Standard 4 IEEE-754 rounding modes for FP ALU and Stage C \texttt{dest\_width\_narrow}. Default: \texttt{none} (RNE).

\begin{center}\footnotesize
\begin{longtable}{@{}rll>{\raggedright\arraybackslash}p{0.62\linewidth}@{}}
\toprule
Enc & Value & Arch & Description \& Worked Example \\
\midrule\endhead
\texttt{0} & \texttt{\imod{none}} & v9--v14 & Round to nearest ties-to-even (RTE, default IEEE-754 mode). \newline {\scriptsize\ttfamily \mnem{FADD.f32} \dst{r0}, \src{r1}, \src{r2}} \\[2pt]
\texttt{1} & \texttt{\imod{round\_up}} & v9--v14 & Round toward $+\infty$ (RTP / ceiling). \newline {\scriptsize\ttfamily \mnem{FADD.f32}\imod{.round\_up} \dst{r0}, \src{r1}, \src{r2}} \\[2pt]
\texttt{2} & \texttt{\imod{round\_down}} & v9--v14 & Round toward $-\infty$ (RTN / floor). \newline {\scriptsize\ttfamily \mnem{FADD.f32}\imod{.round\_down} \dst{r0}, \src{r1}, \src{r2}} \\[2pt]
\texttt{3} & \texttt{\imod{round\_zero}} & v9--v14 & Round toward zero (RTZ / truncate). \newline {\scriptsize\ttfamily \mnem{FADD.f32}\imod{.round\_zero} \dst{r0}, \src{r1}, \src{r2}} \\[2pt]
\bottomrule
\end{longtable}
\end{center}
\par\textbf{Combined Walkthroughs:}
\par\textbf{Example 1 (Stage C FP32 -\textgreater{} FP16 narrowing with `.round\_zero`):}
\codebox{\mnem{FADD.f32}\imod{.round\_zero} \dst{r0}\imod{.h0}, \src{r1}, \src{r2}}
With \texttt{1.0f + 2\textasciicircum{}-11}: RNE rounds up to \texttt{0x3C02} in FP16, whereas \texttt{.round\_zero} truncates to \texttt{0x3C01}.
\par\textbf{Example 2 (Directed rounding `.round\_up` vs `.round\_down`):}
\codebox{\mnem{FADD.f32}\imod{.round\_up} \dst{r0}, \src{r1}, \src{r2}}
With \texttt{1.0f + 2\textasciicircum{}-25}: RNE returns \texttt{1.0f} (\texttt{0x3F800000}), whereas \texttt{.round\_up} rounds toward \texttt{+Inf} -\textgreater{} \texttt{0x3F800001}.

\section{Integer / Extended Rounding (round\_integer\_m)}
\label{enum:round_integer_m}
Adds \texttt{round\_na} (ties away from zero) for \texttt{FROUND} and float-to-int conversions. Default: \texttt{none} (RNE).

\begin{center}\footnotesize
\begin{longtable}{@{}rll>{\raggedright\arraybackslash}p{0.62\linewidth}@{}}
\toprule
Enc & Value & Arch & Description \& Worked Example \\
\midrule\endhead
\texttt{0} & \texttt{\imod{none}} & v9--v14 & Round float to integer using nearest ties-to-even (RTE). \newline {\scriptsize\ttfamily \mnem{F32\_TO\_S32} \dst{r0}, \src{r1}} \\[2pt]
\texttt{1} & \texttt{\imod{round\_up}} & v9--v14 & Round float to integer toward $+\infty$ (\texttt{ceil}). \newline {\scriptsize\ttfamily \mnem{F32\_TO\_S32}\imod{.round\_up} \dst{r0}, \src{r1}} \\[2pt]
\texttt{2} & \texttt{\imod{round\_down}} & v9--v14 & Round float to integer toward $-\infty$ (\texttt{floor}). \newline {\scriptsize\ttfamily \mnem{F32\_TO\_S32}\imod{.round\_down} \dst{r0}, \src{r1}} \\[2pt]
\texttt{3} & \texttt{\imod{round\_zero}} & v9--v14 & Truncate fractional part toward zero (GLSL \texttt{int(x)}). \newline {\scriptsize\ttfamily \mnem{F32\_TO\_S32}\imod{.round\_zero} \dst{r0}, \src{r1}} \\[2pt]
\texttt{4} & \texttt{\imod{round\_na}} & v9--v14 & Round to nearest integer, ties away from zero. \newline {\scriptsize\ttfamily \mnem{F32\_TO\_S32}\imod{.round\_na} \dst{r0}, \src{r1}} \\[2pt]
\bottomrule
\end{longtable}
\end{center}
\par\textbf{Combined Walkthroughs:}
\par\textbf{Example 1 (Halfway tie: RNE vs `.round\_na`):}
\codebox{\mnem{FROUND.f32}\imod{.round\_na} \dst{r0}, \src{r1}}
With \texttt{r1 = 2.5f}: RNE returns \texttt{2.0f} (even), whereas \texttt{.round\_na} rounds ties away from zero -\textgreater{} \texttt{3.0f} (\texttt{-2.5f -\textgreater{} -3.0f}).
\par\textbf{Example 2 (GLSL `int(x)` truncation with `.round\_zero`):}
\codebox{\mnem{F32\_TO\_S32}\imod{.round\_zero} \dst{r0}, \src{r1}}
With \texttt{r1 = -3.9f}: \texttt{.round\_zero} truncates to \texttt{-3}, whereas \texttt{.round\_down} (\texttt{floor}) would yield \texttt{-4}.

\section{Halving-Add Rounding (hadd\_round\_m)}
\label{enum:hadd_round_m}
Selects floor (\texttt{round\_down}) vs ceil (\texttt{round\_up}) for v9--v10 \texttt{HADD}. Default: \texttt{round\_down}.

\begin{center}\footnotesize
\begin{longtable}{@{}rll>{\raggedright\arraybackslash}p{0.62\linewidth}@{}}
\toprule
Enc & Value & Arch & Description \& Worked Example \\
\midrule\endhead
\texttt{0} & \texttt{\imod{round\_down}} & v9--v10 & Halving add floor: $(a + b) \gg 1$. \newline {\scriptsize\ttfamily \mnem{HADD}\imod{.round\_down} \dst{r0}, \src{r1}} \\[2pt]
\texttt{1} & \texttt{\imod{round\_up}} & v9--v10 & Halving add ceil: $(a + b + 1) \gg 1$. \newline {\scriptsize\ttfamily \mnem{HADD}\imod{.round\_up} \dst{r0}, \src{r1}} \\[2pt]
\bottomrule
\end{longtable}
\end{center}
\par\textbf{Combined Walkthroughs:}
\par\textbf{Example 1 (Odd sum floor vs ceil):}
\codebox{\mnem{HADD.u32}\imod{.round\_up} \dst{r0}, \src{r1}, \src{r2}}
With \texttt{r1 = 1, r2 = 2}: \texttt{round\_down} gives \texttt{(1+2)\textgreater{}\textgreater{}1 = 1}; \texttt{round\_up} gives \texttt{(1+2+1)\textgreater{}\textgreater{}1 = 2}.
\par\textbf{Example 2 (Signed negative odd sum):}
\codebox{\mnem{HADD.s32}\imod{.round\_down} \dst{r0}, \src{r1}, \src{r2}}
With \texttt{r1 = -3, r2 = -2}: arithmetic shift \texttt{(-5) \textgreater{}\textgreater{} 1 = -3} under \texttt{.round\_down}, vs \texttt{(-5 + 1) \textgreater{}\textgreater{} 1 = -2} under \texttt{.round\_up}.

\section{Integer Saturation (saturate\_m)}
\label{enum:saturate_m}
Clamps integer overflow/underflow to \texttt{[MIN\_N, MAX\_N]} or \texttt{[0, 2\textasciicircum{}N - 1]}. Default: \texttt{none}.

\begin{center}\footnotesize
\begin{longtable}{@{}rll>{\raggedright\arraybackslash}p{0.62\linewidth}@{}}
\toprule
Enc & Value & Arch & Description \& Worked Example \\
\midrule\endhead
\texttt{0} & \texttt{\imod{none}} & v9--v14 & Default behavior (flag disabled for \texttt{saturate\_m}). \newline {\scriptsize\ttfamily \mnem{IADD}\imod{} \dst{r0}, \src{r1}} \\[2pt]
\texttt{1} & \texttt{\imod{sat}} & v9--v14 & Clamp integer result to signed/unsigned type range instead of wrapping. \newline {\scriptsize\ttfamily \mnem{IADD}\imod{.sat} \dst{r0}, \src{r1}} \\[2pt]
\bottomrule
\end{longtable}
\end{center}
\par\textbf{Combined Walkthroughs:}
\par\textbf{Example 1 (Signed positive overflow clamp):}
\codebox{\mnem{IADD.s32}\imod{.sat} \dst{r0}, \src{r1}, \src{r2}}
With \texttt{r1 = 0x7FFFFFFF, r2 = 1}: wraps to \texttt{0x80000000} without \texttt{.sat}; saturates to \texttt{0x7FFFFFFF} with \texttt{.sat}.
\par\textbf{Example 2 (Unsigned subtraction underflow clamp):}
\codebox{\mnem{ISUB.u32}\imod{.sat} \dst{r0}, \src{r1}, \src{r2}}
With \texttt{r1 = 10, r2 = 25}: \texttt{.sat} clamps underflow to \texttt{r0 = 0} instead of wrapping to \texttt{0xFFFFFFF1}.

\section{FMIN/FMAX NaN Semantics (sem\_m)}
\label{enum:sem_m}
Selects NaN suppression (default) vs payload-preserving NaN propagation (\texttt{\textbar{} qbit}). Applied before \texttt{clamp}.

\begin{center}\footnotesize
\begin{longtable}{@{}rll>{\raggedright\arraybackslash}p{0.62\linewidth}@{}}
\toprule
Enc & Value & Arch & Description \& Worked Example \\
\midrule\endhead
\texttt{0} & \texttt{\imod{nan\_suppress}} & v9--v14 & IEEE-754 minNum/maxNum: non-NaN wins over NaN; $-0.0 < +0.0$ deterministically. \newline {\scriptsize\ttfamily \mnem{FMIN.f32}\imod{.nan\_suppress} \dst{r0}, \src{r1}, \src{r2}} \\[2pt]
\texttt{1} & \texttt{\imod{nan\_propagate}} & v9--v14 & Propagate NaN: if either input is NaN, return quietened winning NaN preserving payload. \newline {\scriptsize\ttfamily \mnem{FMAX.f32}\imod{.nan\_propagate} \dst{r0}, \src{r1}, \src{r2}} \\[2pt]
\bottomrule
\end{longtable}
\end{center}
\par\textbf{Combined Walkthroughs:}
\par\textbf{Example 1 (Payload-preserving NaN propagation):}
\codebox{\mnem{FMIN.f32}\imod{.nan\_propagate} \dst{r0}, \src{r1}, \src{r2}}
With \texttt{r1 = 0xFF800001} (sNaN), \texttt{r2 = 1.0f}: \texttt{nan\_suppress} returns \texttt{1.0f}; \texttt{nan\_propagate} sets bit 22 and returns \texttt{0xFFC00001}.
\par\textbf{Example 2 (Deterministic signed zero regardless of `sem`):}
\codebox{\mnem{FMAX.f32}\imod{.nan\_suppress} \dst{r0}, \src{r1}, \src{r2}}
With \texttt{r1 = -0.0f, r2 = +0.0f}: both \texttt{nan\_suppress} and \texttt{nan\_propagate} deterministically return \texttt{+0.0f} (\texttt{FMIN} returns \texttt{-0.0f}).

\section{FLUSH Denormal Mode (flush\_ftz\_m)}
\label{enum:flush_ftz_m}
Flushes subnormal inputs to signed zero in \texttt{FLUSH}. Default: \texttt{none}.

\begin{center}\footnotesize
\begin{longtable}{@{}rll>{\raggedright\arraybackslash}p{0.62\linewidth}@{}}
\toprule
Enc & Value & Arch & Description \& Worked Example \\
\midrule\endhead
\texttt{0} & \texttt{\imod{none}} & v9--v14 & Default behavior (flag disabled for \texttt{flush\_ftz\_m}). \newline {\scriptsize\ttfamily \mnem{FLUSH}\imod{} \dst{r0}, \src{r1}} \\[2pt]
\texttt{1} & \texttt{\imod{ftz}} & v9--v14 & Flush subnormal float inputs/outputs to signed zero ($\pm 0.0$). \newline {\scriptsize\ttfamily \mnem{FLUSH}\imod{.ftz} \dst{r0}, \src{r1}} \\[2pt]
\bottomrule
\end{longtable}
\end{center}
\par\textbf{Combined Walkthroughs:}
\par\textbf{Example 1 (Negative subnormal FTZ):}
\codebox{\mnem{FLUSH.f32}\imod{.ftz} \dst{r0}, \src{r1}}
With \texttt{r1 = 0x80000001} (negative subnormal): \texttt{.ftz} flushes to \texttt{-0.0f} (\texttt{0x80000000}), preserving sign.
\par\textbf{Example 2 (Normal value passes unchanged):}
\codebox{\mnem{FLUSH.f32}\imod{.ftz} \dst{r0}, \src{r1}}
With \texttt{r1 = 0x00800000} (\texttt{FLT\_MIN}, smallest normal): exponent is non-zero, so \texttt{r0 = 0x00800000} unchanged.

\section{FLUSH Infinity Mode (flush\_inf\_m)}
\label{enum:flush_inf_m}
Clamps \texttt{$\pm$Inf} to \texttt{$\pm$MAX\_FINITE} in \texttt{FLUSH}. Default: \texttt{none}.

\begin{center}\footnotesize
\begin{longtable}{@{}rll>{\raggedright\arraybackslash}p{0.62\linewidth}@{}}
\toprule
Enc & Value & Arch & Description \& Worked Example \\
\midrule\endhead
\texttt{0} & \texttt{\imod{none}} & v9--v14 & Default behavior (flag disabled for \texttt{flush\_inf\_m}). \newline {\scriptsize\ttfamily \mnem{FLUSH}\imod{} \dst{r0}, \src{r1}} \\[2pt]
\texttt{1} & \texttt{\imod{flush\_inf}} & v9--v14 & Flush $\pm\infty$ to finite clamp/zero per \texttt{FLUSH} mode. \newline {\scriptsize\ttfamily \mnem{FLUSH}\imod{.flush\_inf} \dst{r0}, \src{r1}} \\[2pt]
\bottomrule
\end{longtable}
\end{center}
\par\textbf{Combined Walkthroughs:}
\par\textbf{Example 1 (Clamp `+Inf` to `+FLT\_MAX`):}
\codebox{\mnem{FLUSH.f32}\imod{.flush\_inf} \dst{r0}, \src{r1}}
With \texttt{r1 = 0x7F800000} (\texttt{+Inf}): \texttt{.flush\_inf} clamps to \texttt{0x7F7FFFFF} (\texttt{+3.4028235e38f}).
\par\textbf{Example 2 (Clamp `-Inf` to `-FLT\_MAX`):}
\codebox{\mnem{FLUSH.f32}\imod{.flush\_inf} \dst{r0}, \src{r1}}
With \texttt{r1 = 0xFF800000} (\texttt{-Inf}): \texttt{.flush\_inf} clamps to \texttt{0xFF7FFFFF} (\texttt{-3.4028235e38f}).

\section{FLUSH NaN Mode (flush\_nan\_m)}
\label{enum:flush_nan_m}
Flushes NaN to \texttt{+0.0} (\texttt{flush\_nan}) or quietens sNaN (\texttt{quiet\_nan}). Default: \texttt{none}.

\begin{center}\footnotesize
\begin{longtable}{@{}rll>{\raggedright\arraybackslash}p{0.62\linewidth}@{}}
\toprule
Enc & Value & Arch & Description \& Worked Example \\
\midrule\endhead
\texttt{0} & \texttt{\imod{none}} & v9--v14 & Default behavior (flag disabled for \texttt{flush\_nan\_m}). \newline {\scriptsize\ttfamily \mnem{FLUSH}\imod{} \dst{r0}, \src{r1}} \\[2pt]
\texttt{1} & \texttt{\imod{flush\_nan}} & v9--v14 & Flush NaN inputs to $+0.0$. \newline {\scriptsize\ttfamily \mnem{FLUSH}\imod{.flush\_nan} \dst{r0}, \src{r1}} \\[2pt]
\texttt{2} & \texttt{\imod{quiet\_nan}} & v9--v14 & Convert signaling NaNs (sNaN) to quiet NaNs (qNaN) by setting bit 22. \newline {\scriptsize\ttfamily \mnem{FLUSH}\imod{.quiet\_nan} \dst{r0}, \src{r1}} \\[2pt]
\bottomrule
\end{longtable}
\end{center}
\par\textbf{Combined Walkthroughs:}
\par\textbf{Example 1 (Flush NaN to `+0.0`):}
\codebox{\mnem{FLUSH.f32}\imod{.flush\_nan} \dst{r0}, \src{r1}}
With \texttt{r1 = 0x7FC00000} (qNaN) or \texttt{0xFF800001} (sNaN): \texttt{.flush\_nan} returns \texttt{+0.0f} (\texttt{0x00000000}).
\par\textbf{Example 2 (Quieten sNaN preserving payload (`.quiet\_nan`)):}
\codebox{\mnem{FLUSH.f32}\imod{.quiet\_nan} \dst{r0}, \src{r1}}
With \texttt{r1 = 0x7F800042} (sNaN): \texttt{.quiet\_nan} sets quiet bit 22 -\textgreater{} \texttt{r0 = 0x7FC00042}.

\section{FMA\_RSCALE Special Modes (fma\_rscale\_special\_32\_m)}
\label{enum:fma_rscale_special_32_m}
Controls \texttt{0 * Inf} and \texttt{Inf} addend overrides in \texttt{FMA\_RSCALE}/\texttt{FADD\_RSCALE}/\texttt{FMUL\_RSCALE}. Default: \texttt{none}.

\begin{center}\footnotesize
\begin{longtable}{@{}rll>{\raggedright\arraybackslash}p{0.62\linewidth}@{}}
\toprule
Enc & Value & Arch & Description \& Worked Example \\
\midrule\endhead
\texttt{0} & \texttt{\imod{none}} & v9--v14 & Standard FP64-widened $2^e$-scaled FMA/FADD/FMUL. \newline {\scriptsize\ttfamily \mnem{FMA\_RSCALE.f32} \dst{r0}, \src{r1}, \src{r2}, \src{r3}, \src{r4}} \\[2pt]
\texttt{1} & \texttt{\imod{n}} & v9--v14 & If $|a|=0$ or $|b|=0$, force product to $\pm 0.0$ instead of $0\cdot\infty=\text{NaN}$. \newline {\scriptsize\ttfamily \mnem{FMA\_RSCALE.f32}\imod{.n} \dst{r0}, \src{r1}, \src{r2}, \src{r3}, \src{r4}} \\[2pt]
\texttt{2} & \texttt{\imod{left}} & v9--v14 & If $0\cdot\infty$ or $\infty\cdot 0$, return left operand $a$ unmodified. \newline {\scriptsize\ttfamily \mnem{FMUL\_RSCALE.f32}\imod{.left} \dst{r0}, \src{r1}, \src{r2}, \src{r3}} \\[2pt]
\texttt{3} & \texttt{\imod{n\_add}} & v9--v14 & Suppress $\infty + (-\infty)$ NaN during range reduction (valid on FMA/FADD\_RSCALE only). \newline {\scriptsize\ttfamily \mnem{FADD\_RSCALE.f32}\imod{.n\_add} \dst{r0}, \src{r1}, \src{r2}, \src{r3}} \\[2pt]
\bottomrule
\end{longtable}
\end{center}
\par\textbf{Combined Walkthroughs:}
\par\textbf{Example 1 (Suppress `0 * Inf` NaN with `.n`):}
\codebox{\mnem{FMA\_RSCALE.f32}\imod{.n} \dst{r0}, \src{k0}, \src{r\_inf}, \src{k1f}, \src{k0}}
With \texttt{0 * Inf + 1.0}: default yields \texttt{NaN}; \texttt{.n} forces product to \texttt{0.0}, yielding \texttt{1.0f}.
\par\textbf{Example 2 (Preserve left operand with `.left` on `FMUL\_RSCALE`):}
\codebox{\mnem{FMUL\_RSCALE.f32}\imod{.left} \dst{r0}, \src{r1}, \src{r\_inf}, \src{k0}}
With \texttt{r1 = -0.0f, r\_inf = +Inf}: \texttt{.left} returns \texttt{r1} (\texttt{-0.0f}) instead of \texttt{NaN}.

\section{FREXP Normalization Mode (frexp\_special\_m)}
\label{enum:frexp_special_m}
Selects significand interval: \texttt{none} \texttt{[0.5, 1)}, \texttt{sqrt} \texttt{[0.25, 1)} (\texttt{4\textasciicircum{}E}), \texttt{log} \texttt{[0.75, 1.5)}. Default: \texttt{none}.

\begin{center}\footnotesize
\begin{longtable}{@{}rll>{\raggedright\arraybackslash}p{0.62\linewidth}@{}}
\toprule
Enc & Value & Arch & Description \& Worked Example \\
\midrule\endhead
\texttt{0} & \texttt{\imod{none}} & v9--v14 & Standard IEEE \texttt{frexp}: mantissa in $[0.5, 1.0)$, exponent $e - 126$. \newline {\scriptsize\ttfamily \mnem{FREXPM.f32} \dst{r0}, \src{r1}} \\[2pt]
\texttt{1} & \texttt{\imod{sqrt}} & v9--v14 & Square-root decomposition: mantissa in $[0.25, 1.0)$, even exponent $(e-125)\gg 1$. \newline {\scriptsize\ttfamily \mnem{FREXPE.f32}\imod{.sqrt} \dst{r0}, \src{r1}} \\[2pt]
\texttt{2} & \texttt{\imod{log}} & v9--v14 & Logarithm decomposition: mantissa in $[0.75, 1.5)$; $\pm\infty \to \pm 1.0$ on \texttt{FREXPM}. \newline {\scriptsize\ttfamily \mnem{FREXPM.f32}\imod{.log} \dst{r0}, \src{r1}} \\[2pt]
\bottomrule
\end{longtable}
\end{center}
\par\textbf{Combined Walkthroughs:}
\par\textbf{Example 1 (Infinity override under `.log`):}
\codebox{\mnem{FREXPM.f32}\imod{.log} \dst{r0}, \src{r\_inf}}
On \texttt{+Inf}: \texttt{none} and \texttt{sqrt} return \texttt{+Inf}, whereas \texttt{.log} returns \texttt{+1.0f} (\texttt{0x3F800000}).
\par\textbf{Example 2 (Even-exponent `.sqrt` decomposition):}
\codebox{\mnem{FREXPM.f32}\imod{.sqrt} \dst{r0}, \src{r1}}
With \texttt{r1 = 2.0f}: \texttt{.sqrt} normalizes into \texttt{[0.25, 1.0)} -\textgreater{} \texttt{0.5f} with \texttt{FREXPE.sqrt = 1} (\texttt{0.5 * 4\textasciicircum{}1 = 2.0}).

\section{LDEXP Special Threshold Mode (ldexp\_inf\_m)}
\label{enum:ldexp_inf_m}
Overrides \texttt{$\pm$Inf}/\texttt{$\pm$0.0} at extreme exponents \texttt{exp \textless{} -16384} (\texttt{.inf}) and \texttt{exp \textgreater{}= +16384} (\texttt{.inf0}). Default: \texttt{none}.

\begin{center}\footnotesize
\begin{longtable}{@{}rll>{\raggedright\arraybackslash}p{0.62\linewidth}@{}}
\toprule
Enc & Value & Arch & Description \& Worked Example \\
\midrule\endhead
\texttt{0} & \texttt{\imod{none}} & v9--v14 & Standard \texttt{ldexp(x, e)} without extreme-exponent $\pm\infty/\pm 0$ overrides. \newline {\scriptsize\ttfamily \mnem{LDEXP.f32} \dst{r0}, \src{r1}, \src{r2}} \\[2pt]
\texttt{1} & \texttt{\imod{inf}} & v9--v14 & For $e < -16384$: $\pm\infty \to \pm 1.0$ and finite $\to \pm 0.0$. \newline {\scriptsize\ttfamily \mnem{LDEXP.f32}\imod{.inf} \dst{r0}, \src{r1}, \src{r2}} \\[2pt]
\texttt{2} & \texttt{\imod{inf0}} & v9--v14 & Also for $e \ge +16384$: $\pm 0.0 \to \pm 1.0$ and non-zero $\to \pm\infty$. \newline {\scriptsize\ttfamily \mnem{LDEXP.f32}\imod{.inf0} \dst{r0}, \src{r1}, \src{r2}} \\[2pt]
\bottomrule
\end{longtable}
\end{center}
\par\textbf{Combined Walkthroughs:}
\par\textbf{Example 1 (Negative extreme exponent (`.inf`)):}
\codebox{\mnem{LDEXP.f32}\imod{.inf} \dst{r0}, \src{r\_inf}, \src{r\_neg20000}}
With \texttt{r\_inf = +Inf, exp = -20000 (\textless{} -16384)}: \texttt{.inf} returns \texttt{+1.0f} (used in \texttt{atan2} denominator scaling).
\par\textbf{Example 2 (Positive extreme exponent (`.inf0`)):}
\codebox{\mnem{LDEXP.f32}\imod{.inf0} \dst{r0}, \src{k0}, \src{r\_pos20000}}
With \texttt{k0 = +0.0f, exp = +20000 (\textgreater{}= 16384)}: \texttt{.inf0} overrides \texttt{0 * 2\textasciicircum{}20000} to \texttt{+1.0f} (\texttt{0x3F800000}).

\section{Comparison Condition (cmpf\_m)}
\label{enum:cmpf_m}
Condition code for \texttt{FCMP*}, \texttt{ICMP*}, \texttt{CSEL}, and \texttt{DISCARD.f32}. FP \texttt{ne} is unordered; \texttt{gtlt}/\texttt{total} are FP-only.

\begin{center}\footnotesize
\begin{longtable}{@{}rll>{\raggedright\arraybackslash}p{0.62\linewidth}@{}}
\toprule
Enc & Value & Arch & Description \& Worked Example \\
\midrule\endhead
\texttt{0} & \texttt{\imod{eq}} & v9--v14 & Equal ($a == b$); false if either float input is NaN. \newline {\scriptsize\ttfamily \mnem{FCMP.f32}\imod{.eq.i1} \dst{r0}, \src{r1}, \src{r2}} \\[2pt]
\texttt{1} & \texttt{\imod{gt}} & v9--v14 & Greater than ($a > b$); false on NaN. \newline {\scriptsize\ttfamily \mnem{ICMP.s32}\imod{.gt.i1} \dst{r0}, \src{r1}, \src{r2}} \\[2pt]
\texttt{2} & \texttt{\imod{ge}} & v9--v14 & Greater than or equal ($a \ge b$); false on NaN. \newline {\scriptsize\ttfamily \mnem{FCMP.f32}\imod{.ge.i1} \dst{r0}, \src{r1}, \src{r2}} \\[2pt]
\texttt{3} & \texttt{\imod{ne}} & v9--v14 & Not equal ($a \ne b$); true if either float input is NaN (unordered). \newline {\scriptsize\ttfamily \mnem{FCMP.f32}\imod{.ne.i1} \dst{r0}, \src{r1}, \src{r2}} \\[2pt]
\texttt{4} & \texttt{\imod{lt}} & v9--v14 & Less than ($a < b$); false on NaN. \newline {\scriptsize\ttfamily \mnem{ICMP.u32}\imod{.lt.i1} \dst{r0}, \src{r1}, \src{r2}} \\[2pt]
\texttt{5} & \texttt{\imod{le}} & v9--v14 & Less than or equal ($a \le b$); false on NaN. \newline {\scriptsize\ttfamily \mnem{FCMP.f32}\imod{.le.i1} \dst{r0}, \src{r1}, \src{r2}} \\[2pt]
\texttt{6} & \texttt{\imod{gtlt}} & v9--v14 & Ordered not-equal ($a < b \lor a > b$); false if either float is NaN. \newline {\scriptsize\ttfamily \mnem{FCMP.f32}\imod{.gtlt.i1} \dst{r0}, \src{r1}, \src{r2}} \\[2pt]
\texttt{7} & \texttt{\imod{total}} & v9--v14 & IEEE-754 totalOrder predicate (orders $-0.0 < +0.0$ and NaN bit patterns). \newline {\scriptsize\ttfamily \mnem{FCMP.f32}\imod{.total.i1} \dst{r0}, \src{r1}, \src{r2}} \\[2pt]
\bottomrule
\end{longtable}
\end{center}
\par\textbf{Combined Walkthroughs:}
\par\textbf{Example 1 (Ordered `.gtlt` vs unordered `.ne` on NaN):}
\codebox{\mnem{FCMP.f32}\imod{.gtlt.i1} \dst{r0}, \src{r\_nan}, \src{k0}}
On \texttt{(NaN, 0.0)}: \texttt{.ne} returns \texttt{1} (unordered), whereas \texttt{.gtlt} returns \texttt{0} (ordered not-equal).
\par\textbf{Example 2 (IEEE `.total` ordering of `-0.0` vs `+0.0`):}
\codebox{\mnem{FCMP.f32}\imod{.total.i1} \dst{r0}, \src{r\_neg0}, \src{k0}}
Standard \texttt{.lt} says \texttt{-0.0 \textless{} +0.0} is \texttt{0} (false), whereas \texttt{.total} (IEEE \texttt{totalOrder}) returns \texttt{1} (\texttt{-0.0 \textless{} +0.0}).

\section{Comparison Result Format (cmp\_result\_type\_m)}
\label{enum:cmp_result_type_m}
Formats boolean comparison output (\texttt{i1}: 1, \texttt{f1}: 1.0, \texttt{m1}: all-ones, \texttt{c}: 3-way \texttt{\{-1,0,+1\}} on \texttt{ICMP\_MULTI.u32} ONLY). Default: \texttt{i1}.

\begin{center}\footnotesize
\begin{longtable}{@{}rll>{\raggedright\arraybackslash}p{0.62\linewidth}@{}}
\toprule
Enc & Value & Arch & Description \& Worked Example \\
\midrule\endhead
\texttt{0} & \texttt{\imod{i1}} & v9--v14 & Integer boolean result: \texttt{1} if true, \texttt{0} if false. \newline {\scriptsize\ttfamily \mnem{ICMP.u32}\imod{.eq.i1} \dst{r0}, \src{r1}, \src{r2}} \\[2pt]
\texttt{1} & \texttt{\imod{f1}} & v9--v14 & Floating-point boolean result: \texttt{1.0} (\texttt{0x3F800000}) if true, \texttt{0.0} if false. \newline {\scriptsize\ttfamily \mnem{FCMP.f32}\imod{.lt.f1} \dst{r0}, \src{r1}, \src{r2}} \\[2pt]
\texttt{2} & \texttt{\imod{m1}} & v9--v14 & Bitmask boolean result: all-ones (\texttt{0xFFFFFFFF} / \texttt{-1}) if true, \texttt{0} if false. \newline {\scriptsize\ttfamily \mnem{ICMP.s32}\imod{.ge.m1} \dst{r0}, \src{r1}, \src{r2}} \\[2pt]
\texttt{3} & \texttt{\imod{c}} & v9--v14 & Carry/chained comparison result (legal on \texttt{ICMP\_MULTI} only per DDK \texttt{ModInfo}). \newline {\scriptsize\ttfamily \mnem{ICMP\_MULTI.u32}\imod{.eq.c} \dst{r0}, \src{r1}, \src{r2}, \src{r3}} \\[2pt]
\bottomrule
\end{longtable}
\end{center}
\par\textbf{Combined Walkthroughs:}
\par\textbf{Example 1 (Three boolean encodings (`.i1`, `.f1`, `.m1`)):}
\codebox{\mnem{ICMP.u32}\imod{.eq.m1} \dst{r0}, \src{r1}, \src{r1}}
On equal operands: \texttt{.i1} writes \texttt{0x00000001}, \texttt{.f1} writes \texttt{0x3F800000} (\texttt{1.0f}), \texttt{.m1} writes \texttt{0xFFFFFFFF} (\texttt{-1}).
\par\textbf{Example 2 (3-way chained carry (`.c`) on `ICMP\_MULTI.u32`):}
\codebox{\mnem{ICMP\_MULTI.u32}\imod{.lt.c} \dst{r0}, \src{r1}, \src{r2}, \src{k0}}
Outputs \texttt{-1} (\texttt{0xFFFFFFFF}) if \texttt{r1 \textless{} r2}, \texttt{0} if \texttt{r1 == r2}, or \texttt{+1} if \texttt{r1 \textgreater{} r2} for 64-bit chained comparison.

\section{Mux Condition (mux\_cmpf\_m)}
\label{enum:mux_cmpf_m}
Condition tested on \texttt{src2} selector in \texttt{MUX}. Default: \texttt{int\_zero} (encoding \texttt{1}, NOT \texttt{bit}!).

\begin{center}\footnotesize
\begin{longtable}{@{}rll>{\raggedright\arraybackslash}p{0.62\linewidth}@{}}
\toprule
Enc & Value & Arch & Description \& Worked Example \\
\midrule\endhead
\texttt{0} & \texttt{\imod{neg}} & v9--v14 & Select \texttt{src0} if sign bit (bit 31) of \texttt{src2} is 1, else \texttt{src1}. \newline {\scriptsize\ttfamily \mnem{MUX.i32}\imod{.neg} \dst{r0}, \src{r1}, \src{r2}, \src{r3}} \\[2pt]
\texttt{1} & \texttt{\imod{int\_zero}} & v9--v14 & Select \texttt{src0} if integer \texttt{src2 == 0}, else \texttt{src1}. \newline {\scriptsize\ttfamily \mnem{MUX.i32}\imod{.int\_zero} \dst{r0}, \src{r1}, \src{r2}, \src{r3}} \\[2pt]
\texttt{2} & \texttt{\imod{fp\_zero}} & v9--v14 & Select \texttt{src0} if float \texttt{src2 == ±0.0}, else \texttt{src1}. \newline {\scriptsize\ttfamily \mnem{MUX.i32}\imod{.fp\_zero} \dst{r0}, \src{r1}, \src{r2}, \src{r3}} \\[2pt]
\texttt{3} & \texttt{\imod{bit}} & v9--v14 & Per-bit blend: \texttt{(src0 \& src2) | (src1 \& ~src2)}. \newline {\scriptsize\ttfamily \mnem{MUX.i32}\imod{.bit} \dst{r0}, \src{r1}, \src{r2}, \src{r3}} \\[2pt]
\bottomrule
\end{longtable}
\end{center}
\par\textbf{Combined Walkthroughs:}
\par\textbf{Example 1 (Default `.int\_zero` selection):}
\codebox{\mnem{MUX.i32}\imod{.int\_zero} \dst{r0}, \src{r1}, \src{r2}, \src{r3}}
With \texttt{r1 = 0xAAAA, r2 = 0xBBBB, r3 = 0}: \texttt{int\_zero} sees \texttt{r3 == 0} -\textgreater{} selects \texttt{r1 = 0xAAAA} (if \texttt{r3 != 0}, selects \texttt{r2}).
\par\textbf{Example 2 (Per-bit blend (`.bit`) vs sign-bit select (`.neg`)):}
\codebox{\mnem{MUX.i32}\imod{.bit} \dst{r0}, \src{r1}, \src{r2}, \src{r3}}
With \texttt{r1 = 0x0000FFFF, r2 = 0xFFFF0000, r3 = 0x0F0F0F0F}: blends per bit -\textgreater{} \texttt{r0 = 0xF0F00F0F}. Under \texttt{.neg}, bit 31 of \texttt{r3} chooses the entire word.

\section{Destination Lane Write-Mask (dest\_width\_m)}
\label{enum:dest_width_m}
Pure 16-bit lane write-mask on 32-bit/packed ops (NO bit shift!). Default: \texttt{none} (full 32-bit).

\begin{center}\footnotesize
\begin{longtable}{@{}rll>{\raggedright\arraybackslash}p{0.62\linewidth}@{}}
\toprule
Enc & Value & Arch & Description \& Worked Example \\
\midrule\endhead
\texttt{1} & \texttt{\imod{h0}} & v9--v14 & Write low 16-bit halfword \texttt{dst[15:0]}; zero-fill upper \texttt{dst[31:16]}. \newline {\scriptsize\ttfamily \mnem{ADR.f16} \dst{r0}\imod{.h0}, \src{r1}} \\[2pt]
\texttt{2} & \texttt{\imod{h1}} & v9--v14 & Write high 16-bit halfword \texttt{dst[31:16]}; preserve/mask low halfword per lane-mask rule. \newline {\scriptsize\ttfamily \mnem{ADR.f16} \dst{r0}\imod{.h1}, \src{r1}} \\[2pt]
\texttt{3} & \texttt{\imod{none}} & v9--v14 & Write full 32-bit destination register (both halfwords). \newline {\scriptsize\ttfamily \mnem{ADR.f32} \dst{r0}, \src{r1}} \\[2pt]
\bottomrule
\end{longtable}
\end{center}
\par\textbf{Combined Walkthroughs:}
\par\textbf{Example 1 (Upper halfword write-mask (`r0.h1`)):}
\codebox{\mnem{FADD.v2f16} \dst{r0}\imod{.h1}, \src{r1}, \src{r2}}
Computes both FP16 lanes; \texttt{.h1} writes ONLY lane 1 (\texttt{val32[31:16]}) into \texttt{r0[31:16]}, preserving \texttt{r0[15:0]}.
\par\textbf{Example 2 (Low halfword write-mask (`r0.h0`) on `MOV\_IMM`):}
\codebox{\mnem{MOV\_IMM.i32} \dst{r0}\imod{.h0}, \imm{\#0xAAAA1234}}
With \texttt{r0\_old = 0xDEAD0000}: \texttt{.h0} writes \texttt{0x1234} to \texttt{r0[15:0]} and preserves \texttt{r0[31:16]} -\textgreater{} \texttt{r0 = 0xDEAD1234}.

\section{Narrowing Destination Writeback (dest\_width\_narrow\_m)}
\label{enum:dest_width_narrow_m}
Fused FP32-\textgreater{}FP16 narrowing using the instruction's \texttt{round} mode into \texttt{h0}, \texttt{h1}, or replicated \texttt{h01}. Default: \texttt{none}.

\begin{center}\footnotesize
\begin{longtable}{@{}rll>{\raggedright\arraybackslash}p{0.62\linewidth}@{}}
\toprule
Enc & Value & Arch & Description \& Worked Example \\
\midrule\endhead
\texttt{0} & \texttt{\imod{h01}} & v9--v14 & Replicate/pack 16-bit result into both \texttt{dst[15:0]} and \texttt{dst[31:16]}. \newline {\scriptsize\ttfamily \mnem{F32\_TO\_F16.f16} \dst{r0}\imod{.h01}, \src{r1}} \\[2pt]
\texttt{1} & \texttt{\imod{h0}} & v9--v14 & Write low 16-bit halfword \texttt{dst[15:0]}; zero-fill upper \texttt{dst[31:16]}. \newline {\scriptsize\ttfamily \mnem{F32\_TO\_F16.f16} \dst{r0}\imod{.h0}, \src{r1}} \\[2pt]
\texttt{2} & \texttt{\imod{h1}} & v9--v14 & Write high 16-bit halfword \texttt{dst[31:16]}; preserve/mask low halfword per lane-mask rule. \newline {\scriptsize\ttfamily \mnem{F32\_TO\_F16.f16} \dst{r0}\imod{.h1}, \src{r1}} \\[2pt]
\texttt{3} & \texttt{\imod{none}} & v9--v14 & Write full 32-bit destination register (both halfwords). \newline {\scriptsize\ttfamily \mnem{F32\_TO\_F16.f32} \dst{r0}, \src{r1}} \\[2pt]
\bottomrule
\end{longtable}
\end{center}
\par\textbf{Combined Walkthroughs:}
\par\textbf{Example 1 (Narrow + broadcast to both halfwords (`.h01`)):}
\codebox{\mnem{FADD.f32}\imod{.round\_zero} \dst{r0}\imod{.h01}, \src{r1}, \src{r2}}
With \texttt{r1 = 1.0f, r2 = 1.0f}: narrows \texttt{2.0f -\textgreater{} 0x4000} and \texttt{.h01} replicates to both halves -\textgreater{} \texttt{r0 = 0x40004000}.
\par\textbf{Example 2 (Narrow into upper half preserving low half (`.h1`)):}
\codebox{\mnem{FMA.f32} \dst{r0}\imod{.h1}, \src{r1}, \src{r2}, \src{r3}}
With \texttt{r0\_old = 0x00003C00} and FP32 result \texttt{4.0f}: narrows to \texttt{0x4400} in \texttt{r0[31:16]} -\textgreater{} \texttt{r0 = 0x44003C00}.

\section{Replicating Scalar FP16 Destination (dest\_width\_replicate\_m)}
\label{enum:dest_width_replicate_m}
Places scalar FP16 output (\texttt{FRCP/FRSQ/FEXP/FLOG/FPOW.f16}) into \texttt{h0}, \texttt{h1}, or replicated \texttt{h01}. Default: \texttt{none}.

\begin{center}\footnotesize
\begin{longtable}{@{}rll>{\raggedright\arraybackslash}p{0.62\linewidth}@{}}
\toprule
Enc & Value & Arch & Description \& Worked Example \\
\midrule\endhead
\texttt{0} & \texttt{\imod{h01}} & v11--v14 & Replicate/pack 16-bit result into both \texttt{dst[15:0]} and \texttt{dst[31:16]}. \newline {\scriptsize\ttfamily \mnem{FEXP.f16} \dst{r0}\imod{.h01}, \src{r1}} \\[2pt]
\texttt{1} & \texttt{\imod{h0}} & v11--v14 & Write low 16-bit halfword \texttt{dst[15:0]}; zero-fill upper \texttt{dst[31:16]}. \newline {\scriptsize\ttfamily \mnem{FEXP.f16} \dst{r0}\imod{.h0}, \src{r1}} \\[2pt]
\texttt{2} & \texttt{\imod{h1}} & v11--v14 & Write high 16-bit halfword \texttt{dst[31:16]}; preserve/mask low halfword per lane-mask rule. \newline {\scriptsize\ttfamily \mnem{FEXP.f16} \dst{r0}\imod{.h1}, \src{r1}} \\[2pt]
\texttt{3} & \texttt{\imod{none}} & v11--v14 & Write full 32-bit destination register (both halfwords). \newline {\scriptsize\ttfamily \mnem{FEXP.f32} \dst{r0}, \src{r1}} \\[2pt]
\bottomrule
\end{longtable}
\end{center}
\par\textbf{Combined Walkthroughs:}
\par\textbf{Example 1 (Broadcast scalar FP16 reciprocal (`.h01`)):}
\codebox{\mnem{FRCP.f16} \dst{r0}\imod{.h01}, \src{r1}\imod{.h0}}
With \texttt{r1.h0 = 2.0h}: \texttt{1/2.0 = 0.5h} (\texttt{0x3800}), replicated to \texttt{r0 = 0x38003800}.
\par\textbf{Example 2 (Write scalar FP16 into upper half only (`.h1`)):}
\codebox{\mnem{FLOG.f16} \dst{r0}\imod{.h1}, \src{r1}\imod{.h1}}
With \texttt{r1.h1 = 4.0h}, \texttt{r0\_old = 0x00003C00}: writes \texttt{log2(4.0h) = 2.0h} (\texttt{0x4000}) to \texttt{r0[31:16]} -\textgreater{} \texttt{r0 = 0x40003C00}.

\section{Branch Condition Combine (branch\_combine\_m)}
\label{enum:branch_combine_m}
Predicate reduction on \texttt{src0} before optional \texttt{.not} boolean negation. Default: \texttt{none} (\texttt{!= 0}).

\begin{center}\footnotesize
\begin{longtable}{@{}rll>{\raggedright\arraybackslash}p{0.62\linewidth}@{}}
\toprule
Enc & Value & Arch & Description \& Worked Example \\
\midrule\endhead
\texttt{0} & \texttt{\imod{none}} & v9--v14 & Test full 32-bit predicate operand \texttt{src0 != 0}. \newline {\scriptsize\ttfamily \mnem{BRANCH} \src{r0}, \imm{\#+4}} \\[2pt]
\texttt{1} & \texttt{\imod{h0}} & v9--v14 & Test low 16-bit halfword \texttt{src0[15:0] != 0}. \newline {\scriptsize\ttfamily \mnem{BRANCH}\imod{.h0} \src{r0}, \imm{\#+4}} \\[2pt]
\texttt{2} & \texttt{\imod{h1}} & v9--v14 & Test high 16-bit halfword \texttt{src0[31:16] != 0}. \newline {\scriptsize\ttfamily \mnem{BRANCH}\imod{.h1} \src{r0}, \imm{\#+4}} \\[2pt]
\texttt{3} & \texttt{\imod{and}} & v9--v14 & Branch if both halfwords are non-zero (\texttt{h0 != 0 \&\& h1 != 0}). \newline {\scriptsize\ttfamily \mnem{BRANCH}\imod{.and} \src{r0}, \imm{\#+4}} \\[2pt]
\texttt{4} & \texttt{\imod{lowbits}} & v9--v14 & Test low bit(s) of predicate lane(s) for warp condition. \newline {\scriptsize\ttfamily \mnem{BRANCH}\imod{.lowbits} \src{r0}, \imm{\#+4}} \\[2pt]
\bottomrule
\end{longtable}
\end{center}
\par\textbf{Combined Walkthroughs:}
\par\textbf{Example 1 (Both halfwords non-zero (`.and`)):}
\codebox{\mnem{BRANCH} \src{r4}\imod{.and}, \imm{\#0x20}}
With \texttt{r4 = 0x00010001}: both halfwords non-zero -\textgreater{} taken; with \texttt{r4 = 0x0000FFFF}: upper half is 0 -\textgreater{} not taken.
\par\textbf{Example 2 (Single halfword predicate (`.h1`)):}
\codebox{\mnem{BRANCH} \src{r4}\imod{.h1}, \imm{\#0x10}}
Tests \texttt{r4[31:16] != 0} while ignoring \texttt{r4[15:0]} (useful after \texttt{v2*16} comparisons).

\section{Branch Condition Negate (branch\_negate\_m)}
\label{enum:branch_negate_m}
Negates the boolean result of \texttt{combine(src0)} (NOT a bitwise NOT on \texttt{src0}). Default: \texttt{none}.

\begin{center}\footnotesize
\begin{longtable}{@{}rll>{\raggedright\arraybackslash}p{0.62\linewidth}@{}}
\toprule
Enc & Value & Arch & Description \& Worked Example \\
\midrule\endhead
\texttt{0} & \texttt{\imod{none}} & v9--v14 & Default behavior (flag disabled for \texttt{branch\_negate\_m}). \newline {\scriptsize\ttfamily \mnem{BRANCH}\imod{} \dst{r0}, \src{r1}} \\[2pt]
\texttt{1} & \texttt{\imod{not}} & v9--v14 & Bitwise invert operand or negate branch predicate ($\sim x$ / $!p$). \newline {\scriptsize\ttfamily \mnem{BRANCH}\imod{.not} \dst{r0}, \src{r1}} \\[2pt]
\bottomrule
\end{longtable}
\end{center}
\par\textbf{Combined Walkthroughs:}
\par\textbf{Example 1 (Panfrost unconditional branch (`.not k0`)):}
\codebox{\mnem{BRANCH}\imod{.not} \src{k0}, \imm{\#0x10}}
\texttt{k0 != 0} is \texttt{False}; \texttt{.not} flips it to \texttt{True} -\textgreater{} always branches to \texttt{PC + 8 + 0x10}.
\par\textbf{Example 2 (Branch-if-zero (`if (r0 == 0)`)):}
\codebox{\mnem{BRANCH}\imod{.not} \src{r0}, \imm{\#0x28}}
Taken when \texttt{r0 == 0} (\texttt{combine} = False, \texttt{.not} = True); falls through for any non-zero \texttt{r0}.

\section{Jump Address Mode (jump\_address\_mode\_m)}
\label{enum:jump_address_mode_m}
Selects PC-relative vs absolute target address on \texttt{JUMP} (v11+ fixes \texttt{absolute}).

\begin{center}\footnotesize
\begin{longtable}{@{}rll>{\raggedright\arraybackslash}p{0.62\linewidth}@{}}
\toprule
Enc & Value & Arch & Description \& Worked Example \\
\midrule\endhead
\texttt{0} & \texttt{\imod{relative}} & v9--v10 & PC-relative jump target offset. \newline {\scriptsize\ttfamily \mnem{JUMP}\imod{.relative} \dst{r0}, \src{r1}} \\[2pt]
\texttt{1} & \texttt{\imod{absolute}} & v9--v10 & Absolute 64-bit GPU virtual address jump target. \newline {\scriptsize\ttfamily \mnem{JUMP}\imod{.absolute} \dst{r0}, \src{r1}} \\[2pt]
\bottomrule
\end{longtable}
\end{center}
\par\textbf{Combined Walkthroughs:}
\par\textbf{Example 1 (Absolute blend-shader jump (`.absolute`)):}
\codebox{\mnem{JUMP}\imod{.absolute} \src{k1n}, \src{r54}}
Jumps directly to 32-bit shader PC in \texttt{r54} (\texttt{r54 == 0} terminates the shader cleanly).
\par\textbf{Example 2 (PC-relative jump table (v9--v10)):}
\codebox{\mnem{JUMP} \src{k1n}, \src{r2}}
On v9--v10 with relative mode (\texttt{0}): adds signed offset \texttt{r2} to \texttt{PC + 8}.

\section{Subgroup Size (clper\_subgroup\_size\_m)}
\label{enum:clper_subgroup_size_m}
Partitions 16-lane warp into 2-, 4-, 8-, or 16-lane subgroups for \texttt{CLPER*} and \texttt{WMASK}.
\par\textbf{CSF ABI Interaction:} Operates within Mali's 16-lane hardware warp (\texttt{FauSpecial::lane\_width=16}, \hyperref[sec:csf-tsd-fau]{Section~\ref{sec:csf-tsd-fau}}).
\begin{center}\footnotesize
\begin{longtable}{@{}rll>{\raggedright\arraybackslash}p{0.62\linewidth}@{}}
\toprule
Enc & Value & Arch & Description \& Worked Example \\
\midrule\endhead
\texttt{0} & \texttt{\imod{subgroup2}} & v9--v14 & Partition warp into independent 2-lane subgroups for shuffle/reduction. \newline {\scriptsize\ttfamily \mnem{CLPER}\imod{.subgroup2} \dst{r0}, \src{r1}, \src{r2}} \\[2pt]
\texttt{1} & \texttt{\imod{subgroup4}} & v9--v14 & Partition warp into independent 4-lane subgroups for shuffle/reduction. \newline {\scriptsize\ttfamily \mnem{CLPER}\imod{.subgroup4} \dst{r0}, \src{r1}, \src{r2}} \\[2pt]
\texttt{2} & \texttt{\imod{subgroup8}} & v9--v14 & Partition warp into independent 8-lane subgroups for shuffle/reduction. \newline {\scriptsize\ttfamily \mnem{CLPER}\imod{.subgroup8} \dst{r0}, \src{r1}, \src{r2}} \\[2pt]
\texttt{3} & \texttt{\imod{subgroup16}} & v9--v14 & Partition warp into independent 16-lane subgroups for shuffle/reduction. \newline {\scriptsize\ttfamily \mnem{CLPER}\imod{.subgroup16} \dst{r0}, \src{r1}, \src{r2}} \\[2pt]
\bottomrule
\end{longtable}
\end{center}
\par\textbf{Combined Walkthroughs:}
\par\textbf{Example 1 (4-lane fragment quad derivative (`.subgroup4`)):}
\codebox{\mnem{CLPER}\imod{.subgroup4.xor.zero} \dst{r0}, \src{r1}, \src{k01000000}\imod{.b3}}
Exchanges \texttt{r1} with horizontal quad neighbor (\texttt{lane \textasciicircum{} 1} inside each 4-lane quad, used for \texttt{dFdx}).
\par\textbf{Example 2 (Full 16-lane warp reduction (`.subgroup16`)):}
\codebox{\mnem{WMASK}\imod{.subgroup16} \dst{r0}, \src{r1}}
Collects predicate bits from all 16 lanes of the warp into a single 16-bit ballot mask.

\section{CLPER Lane Operation (clper\_lane\_op\_m)}
\label{enum:clper_lane_op_m}
Computes target lane from \texttt{(lane\_id, src1.b0)} (\texttt{none}, \texttt{xor}, \texttt{accumulate}, \texttt{shift}, \texttt{rotate}, v11+ \texttt{low}/\texttt{prefix}).

\begin{center}\footnotesize
\begin{longtable}{@{}rll>{\raggedright\arraybackslash}p{0.62\linewidth}@{}}
\toprule
Enc & Value & Arch & Description \& Worked Example \\
\midrule\endhead
\texttt{0} & \texttt{\imod{none}} & v9--v13 & Direct lane index lookup: read from lane 	exttt{src1 \& (N-1)}. \newline {\scriptsize\ttfamily \mnem{CLPER}\imod{.none} \dst{r0}, \src{r1}, \src{r2}} \\[2pt]
\texttt{1} & \texttt{\imod{xor}} & v9--v13 & Butterfly XOR shuffle: read from lane 	exttt{my\_lane 	extasciicircum{} src1}. \newline {\scriptsize\ttfamily \mnem{CLPER}\imod{.xor} \dst{r0}, \src{r1}, \src{r2}} \\[2pt]
\texttt{2} & \texttt{\imod{accumulate}} & v9--v13 & Tree-reduction accumulate helper across subgroup lanes. \newline {\scriptsize\ttfamily \mnem{CLPER}\imod{.accumulate} \dst{r0}, \src{r1}, \src{r2}} \\[2pt]
\texttt{3} & \texttt{\imod{shift}} & v9--v13 & Subgroup shift up/down: read from 	exttt{my\_lane - src1} (clamped/inactive if OOB). \newline {\scriptsize\ttfamily \mnem{CLPER}\imod{.shift} \dst{r0}, \src{r1}, \src{r2}} \\[2pt]
\texttt{4} & \texttt{\imod{rotate}} & v9--v13 & Cyclic subgroup rotate: read from 	exttt{(my\_lane + src1) mod N}. \newline {\scriptsize\ttfamily \mnem{CLPER}\imod{.rotate} \dst{r0}, \src{r1}, \src{r2}} \\[2pt]
\texttt{5} & \texttt{\imod{low}} & v14 & Broadcast lowest active lane index in the subgroup. \newline {\scriptsize\ttfamily \mnem{CLPER}\imod{.low} \dst{r0}, \src{r1}, \src{r2}} \\[2pt]
\texttt{6} & \texttt{\imod{low\_alt}} & v14 & Alternate low-lane broadcast mode for quad/clustered ops. \newline {\scriptsize\ttfamily \mnem{CLPER}\imod{.low\_alt} \dst{r0}, \src{r1}, \src{r2}} \\[2pt]
\texttt{7} & \texttt{\imod{prefix}} & v14 & Inclusive/exclusive prefix-scan lane index computation. \newline {\scriptsize\ttfamily \mnem{CLPER}\imod{.prefix} \dst{r0}, \src{r1}, \src{r2}} \\[2pt]
\bottomrule
\end{longtable}
\end{center}
\par\textbf{Combined Walkthroughs:}
\par\textbf{Example 1 (Butterfly reduction (`.accumulate`)):}
\codebox{\mnem{CLPER}\imod{.subgroup16.accumulate.zero} \dst{r0}, \src{r1}, \src{r2}\imod{.b0}}
Reads partner lane in a binary-tree butterfly reduction for \texttt{subgroupAdd} / \texttt{subgroupMax}.
\par\textbf{Example 2 (Cyclic subgroup rotate (`.rotate`)):}
\codebox{\mnem{CLPER}\imod{.subgroup16.rotate.zero} \dst{r0}, \src{r1}, \src{r2}\imod{.b0}}
With \texttt{r2.b0 = 1}: lane \texttt{L} reads \texttt{r1} from lane \texttt{(L + 1) \& 15}, wrapping lane 15 back to lane 0.

\section{CLPER Inactive Lane Result (clper\_inactive\_result\_m)}
\label{enum:clper_inactive_result_m}
16 hardware constants returned when a \texttt{CLPER} target lane is inactive or out of bounds. Default: \texttt{zero}.

\begin{center}\footnotesize
\begin{longtable}{@{}rll>{\raggedright\arraybackslash}p{0.62\linewidth}@{}}
\toprule
Enc & Value & Arch & Description \& Worked Example \\
\midrule\endhead
\texttt{0} & \texttt{\imod{zero}} & v9--v14 & Return 0x00000000 for inactive source lanes. \newline {\scriptsize\ttfamily \mnem{CLPER}\imod{.zero} \dst{r0}, \src{r1}, \src{r2}} \\[2pt]
\texttt{1} & \texttt{\imod{umax}} & v9--v14 & Return 0xFFFFFFFF (UINT32\_MAX) for inactive lanes. \newline {\scriptsize\ttfamily \mnem{CLPER}\imod{.umax} \dst{r0}, \src{r1}, \src{r2}} \\[2pt]
\texttt{2} & \texttt{\imod{i1}} & v9--v14 & Return 0x00000001 (int32 1) for inactive lanes. \newline {\scriptsize\ttfamily \mnem{CLPER}\imod{.i1} \dst{r0}, \src{r1}, \src{r2}} \\[2pt]
\texttt{3} & \texttt{\imod{v2i1}} & v9--v14 & Return 0x00010001 (v2i16 {1, 1}) for inactive lanes. \newline {\scriptsize\ttfamily \mnem{CLPER}\imod{.v2i1} \dst{r0}, \src{r1}, \src{r2}} \\[2pt]
\texttt{4} & \texttt{\imod{smin}} & v9--v14 & Return 0x80000000 (INT32\_MIN) for inactive lanes. \newline {\scriptsize\ttfamily \mnem{CLPER}\imod{.smin} \dst{r0}, \src{r1}, \src{r2}} \\[2pt]
\texttt{5} & \texttt{\imod{smax}} & v9--v14 & Return 0x7FFFFFFF (INT32\_MAX) for inactive lanes. \newline {\scriptsize\ttfamily \mnem{CLPER}\imod{.smax} \dst{r0}, \src{r1}, \src{r2}} \\[2pt]
\texttt{6} & \texttt{\imod{v2smin}} & v9--v14 & Return 0x80008000 (v2s16 {INT16\_MIN, INT16\_MIN}). \newline {\scriptsize\ttfamily \mnem{CLPER}\imod{.v2smin} \dst{r0}, \src{r1}, \src{r2}} \\[2pt]
\texttt{7} & \texttt{\imod{v2smax}} & v9--v14 & Return 0x7FFF7FFF (v2s16 {INT16\_MAX, INT16\_MAX}). \newline {\scriptsize\ttfamily \mnem{CLPER}\imod{.v2smax} \dst{r0}, \src{r1}, \src{r2}} \\[2pt]
\texttt{8} & \texttt{\imod{v4smin}} & v9--v14 & Return 0x80808080 (v4s8 {-128, -128, -128, -128}). \newline {\scriptsize\ttfamily \mnem{CLPER}\imod{.v4smin} \dst{r0}, \src{r1}, \src{r2}} \\[2pt]
\texttt{9} & \texttt{\imod{v4smax}} & v9--v14 & Return 0x7F7F7F7F (v4s8 {127, 127, 127, 127}). \newline {\scriptsize\ttfamily \mnem{CLPER}\imod{.v4smax} \dst{r0}, \src{r1}, \src{r2}} \\[2pt]
\texttt{10} & \texttt{\imod{f1}} & v9--v14 & Return 0x3F800000 (FP32 +1.0) for inactive lanes. \newline {\scriptsize\ttfamily \mnem{CLPER}\imod{.f1} \dst{r0}, \src{r1}, \src{r2}} \\[2pt]
\texttt{11} & \texttt{\imod{v2f1}} & v9--v14 & Return 0x3C003C00 (v2f16 {+1.0, +1.0}) for inactive lanes. \newline {\scriptsize\ttfamily \mnem{CLPER}\imod{.v2f1} \dst{r0}, \src{r1}, \src{r2}} \\[2pt]
\texttt{12} & \texttt{\imod{infn}} & v9--v14 & Modifier option \texttt{infn} (encoding 12). \newline {\scriptsize\ttfamily \mnem{CLPER}\imod{.infn} \dst{r0}, \src{r1}} \\[2pt]
\texttt{13} & \texttt{\imod{inf}} & v9--v14 & Return 0x7F800000 (FP32 +Inf) for inactive lanes (min identity). \newline {\scriptsize\ttfamily \mnem{CLPER}\imod{.inf} \dst{r0}, \src{r1}, \src{r2}} \\[2pt]
\texttt{14} & \texttt{\imod{v2infn}} & v9--v14 & Modifier option \texttt{v2infn} (encoding 14). \newline {\scriptsize\ttfamily \mnem{CLPER}\imod{.v2infn} \dst{r0}, \src{r1}} \\[2pt]
\texttt{15} & \texttt{\imod{v2inf}} & v9--v14 & Return 0x7C007C00 (v2f16 {+Inf, +Inf}) for inactive lanes. \newline {\scriptsize\ttfamily \mnem{CLPER}\imod{.v2inf} \dst{r0}, \src{r1}, \src{r2}} \\[2pt]
\bottomrule
\end{longtable}
\end{center}
\par\textbf{Combined Walkthroughs:}
\par\textbf{Example 1 (Additive identity (`.zero`) under divergence):}
\codebox{\mnem{CLPER}\imod{.subgroup16.xor.zero} \dst{r0}, \src{r1}, \src{r2}\imod{.b0}}
If target lane \texttt{L \textasciicircum{} r2.b0} is masked off by divergence, returns \texttt{0} so \texttt{IADD}/\texttt{FADD} reductions stay exact.
\par\textbf{Example 2 (Unsigned min identity (`.umax`) \& FP max identity (`.ninf`)):}
\codebox{\mnem{CLPER}\imod{.subgroup16.xor.umax} \dst{r0}, \src{r1}, \src{r2}\imod{.b0}}
Substitutes \texttt{0xFFFFFFFF} for inactive lanes in \texttt{subgroupMin(uint)} (or \texttt{-Inf} with \texttt{.ninf} for \texttt{subgroupMax(float)}).

\section{Load Extension Mode (sign\_extend\_or\_none\_m)}
\label{enum:sign_extend_or_none_m}
Sign- or zero-extends narrow \texttt{LOAD.i8}/\texttt{i16} into 32-bit staging register. Default: \texttt{none}/\texttt{zext\_or\_none}.

\begin{center}\footnotesize
\begin{longtable}{@{}rll>{\raggedright\arraybackslash}p{0.62\linewidth}@{}}
\toprule
Enc & Value & Arch & Description \& Worked Example \\
\midrule\endhead
\texttt{0} & \texttt{\imod{sext}} & v9--v14 & Sign-extend narrow loaded value (8/16-bit) to 32 bits. \newline {\scriptsize\ttfamily \mnem{LD\_PKA}\imod{.sext} \dst{r0}, \src{r1}} \\[2pt]
\texttt{1} & \texttt{\imod{zext}} & v9--v14 & Zero-extend narrow loaded value (8/16-bit) to 32 bits. \newline {\scriptsize\ttfamily \mnem{LD\_PKA}\imod{.zext} \dst{r0}, \src{r1}} \\[2pt]
\texttt{2} & \texttt{\imod{none}} & v9--v14 & Default behavior (flag disabled for \texttt{sign\_extend\_or\_none\_m}). \newline {\scriptsize\ttfamily \mnem{LD\_PKA}\imod{} \dst{r0}, \src{r1}} \\[2pt]
\bottomrule
\end{longtable}
\end{center}
\par\textbf{Combined Walkthroughs:}
\par\textbf{Example 1 (Sign-extended 8-bit load (`.sext`)):}
\codebox{\mnem{LOAD.i8}\imod{.sext.slot0} \dst{@r1}, \src{[r2:r3]}, \imm{\#0}}
Loading byte \texttt{0x80} (\texttt{-128}): \texttt{.sext} writes \texttt{0xFFFFFF80} into \texttt{r1}.
\par\textbf{Example 2 (Zero-extended 16-bit load (default)):}
\codebox{\mnem{LOAD.i16}\imod{.slot0} \dst{@r1}, \src{[r2:r3]}, \imm{\#0x2}}
Loading halfword \texttt{0x8001}: default \texttt{zext\_or\_none} writes \texttt{0x00008001} (unsigned \texttt{32769}).

\section{Memory Cache Access Hint (access\_load\_m)}
\label{enum:access_load_m}
Cache allocation/streaming policy (\texttt{none}, \texttt{istream}, \texttt{estream}, \texttt{force}). Default: \texttt{none}.
\par\textbf{CSF ABI Interaction:} Controls L1/L2 cache allocation on direct 64-bit VA loads (backed by \hyperref[sec:csf-tsd-fau]{\texttt{FAU\_k}} buffer VAs or \texttt{TSD\_k} TLS/WLS pools).
\begin{center}\footnotesize
\begin{longtable}{@{}rll>{\raggedright\arraybackslash}p{0.62\linewidth}@{}}
\toprule
Enc & Value & Arch & Description \& Worked Example \\
\midrule\endhead
\texttt{0} & \texttt{\imod{none}} & v9--v14 & Default L1/L2 cached load policy. \newline {\scriptsize\ttfamily \mnem{LOAD.i32} \dst{@r0}, \src{r2}, \imm{\#0}} \\[2pt]
\texttt{1} & \texttt{\imod{istream}} & v9--v14 & Inner-streaming hint: bypass/deprioritize L1 cache allocation. \newline {\scriptsize\ttfamily \mnem{LOAD.i32}\imod{.istream} \dst{@r0}, \src{r2}, \imm{\#0}} \\[2pt]
\texttt{2} & \texttt{\imod{estream}} & v9--v14 & Outer-streaming hint: streaming access through L2/outer caches. \newline {\scriptsize\ttfamily \mnem{LOAD.i32}\imod{.estream} \dst{@r0}, \src{r2}, \imm{\#0}} \\[2pt]
\texttt{3} & \texttt{\imod{force}} & v9--v14 & Force cache coherence / bypass stale cache lines (volatile/coherent read). \newline {\scriptsize\ttfamily \mnem{LOAD.i32}\imod{.force} \dst{@r0}, \src{r2}, \imm{\#0}} \\[2pt]
\bottomrule
\end{longtable}
\end{center}
\par\textbf{Combined Walkthroughs:}
\par\textbf{Example 1 (Outer-streaming sequential read (`.estream`)):}
\codebox{\mnem{LOAD.i128}\imod{.estream.slot0} \dst{@r0:r3}, \src{[r4:r5]}, \imm{\#0}}
Hints evict-first streaming cache behavior for large sequential buffer scans to avoid thrashing L2.
\par\textbf{Example 2 (Coherent volatile read (`.force`)):}
\codebox{\mnem{LOAD.i32}\imod{.force.slot0} \dst{@r0}, \src{[r4:r5]}, \imm{\#0}}
Bypasses stale L1 lines to observe writes from other workgroups (\texttt{coherent} SSBO).

\section{Checked Load OOB Mode (oob\_mode\_m, v14)}
\label{enum:oob_mode_m}
Controls out-of-bounds behavior on v14 \texttt{LD\_CHECKED*} (\texttt{none}: return 0; \texttt{oob\_ignore}: disable OOB zeroing).
\par\textbf{CSF ABI Interaction:} Works with \hyperref[sec:csf-buf-attr]{\texttt{Buffer::size}} (32B descriptor in \texttt{SRT\_k}) for hardware-checked \texttt{robustBufferAccess} bounds enforcement.
\begin{center}\footnotesize
\begin{longtable}{@{}rll>{\raggedright\arraybackslash}p{0.62\linewidth}@{}}
\toprule
Enc & Value & Arch & Description \& Worked Example \\
\midrule\endhead
\texttt{0} & \texttt{\imod{none}} & v14 & Default behavior (flag disabled for \texttt{oob\_mode\_m}). \newline {\scriptsize\ttfamily \mnem{LD\_CHECKED}\imod{} \dst{r0}, \src{r1}} \\[2pt]
\texttt{1} & \texttt{\imod{oob\_ignore}} & v14 & Silently return 0 on OOB load or drop OOB store without faulting. \newline {\scriptsize\ttfamily \mnem{LD\_CHECKED}\imod{.oob\_ignore} \dst{r0}, \src{r1}} \\[2pt]
\bottomrule
\end{longtable}
\end{center}
\par\textbf{Combined Walkthroughs:}
\par\textbf{Example 1 (Standard robust OOB zeroing (`none`)):}
\codebox{\mnem{LD\_CHECKED.i32}\imod{.slot0} \dst{@r0}, \src{u3}, \src{r1}, \src{r2}}
With \texttt{r1 \textgreater{}= r2}: suppresses the memory read and writes \texttt{r0 = 0} per Vulkan \texttt{robustBufferAccess}.
\par\textbf{Example 2 (Suppress OOB zeroing (`.oob\_ignore`, v14)):}
\codebox{\mnem{LD\_CHECKED.i32}\imod{.oob\_ignore.slot0} \dst{@r0}, \src{u3}, \src{r1}, \src{r2}}
Leaves destination unmodified on OOB when the shader tracks bounds separately.

\section{Binary Atomic Operation (atom\_operation\_m)}
\label{enum:atom_operation_m}
Selects RMW operation for \texttt{ATOM} / \texttt{ATOM\_RETURN} (\texttt{aadd}, \texttt{asmin}, \texttt{asmax}, \texttt{aumin}, \texttt{aumax}, \texttt{aand}, \texttt{aor}, \texttt{axor}, \texttt{axchg}, \texttt{aaddu}, \texttt{aadds}).

\begin{center}\footnotesize
\begin{longtable}{@{}rll>{\raggedright\arraybackslash}p{0.62\linewidth}@{}}
\toprule
Enc & Value & Arch & Description \& Worked Example \\
\midrule\endhead
\texttt{0} & \texttt{\imod{aaddu}} & v9--v14 & Atomic unsigned add (saturating/wrapping per opcode): \texttt{*p += v}. \newline {\scriptsize\ttfamily \mnem{ATOM.i32}\imod{.aaddu} \src{@r0}, \src{r2}, \imm{\#0}} \\[2pt]
\texttt{1} & \texttt{\imod{aadds}} & v9--v14 & Atomic signed add: \texttt{*p += (int)v}. \newline {\scriptsize\ttfamily \mnem{ATOM.i32}\imod{.aadds} \src{@r0}, \src{r2}, \imm{\#0}} \\[2pt]
\texttt{2} & \texttt{\imod{aadd}} & v9--v14 & Atomic integer add (wrapping two's complement): \texttt{*p += v}. \newline {\scriptsize\ttfamily \mnem{ATOM.i32}\imod{.aadd} \src{@r0}, \src{r2}, \imm{\#0}} \\[2pt]
\texttt{8} & \texttt{\imod{asmin}} & v9--v14 & Atomic signed minimum: \texttt{*p = min((int)*p, (int)v)}. \newline {\scriptsize\ttfamily \mnem{ATOM.i32}\imod{.asmin} \src{@r0}, \src{r2}, \imm{\#0}} \\[2pt]
\texttt{9} & \texttt{\imod{asmax}} & v9--v14 & Atomic signed maximum: \texttt{*p = max((int)*p, (int)v)}. \newline {\scriptsize\ttfamily \mnem{ATOM.i32}\imod{.asmax} \src{@r0}, \src{r2}, \imm{\#0}} \\[2pt]
\texttt{10} & \texttt{\imod{aumin}} & v9--v14 & Atomic unsigned minimum: \texttt{*p = min((uint)*p, (uint)v)}. \newline {\scriptsize\ttfamily \mnem{ATOM.i32}\imod{.aumin} \src{@r0}, \src{r2}, \imm{\#0}} \\[2pt]
\texttt{11} & \texttt{\imod{aumax}} & v9--v14 & Atomic unsigned maximum: \texttt{*p = max((uint)*p, (uint)v)}. \newline {\scriptsize\ttfamily \mnem{ATOM.i32}\imod{.aumax} \src{@r0}, \src{r2}, \imm{\#0}} \\[2pt]
\texttt{12} & \texttt{\imod{aand}} & v9--v14 & Atomic bitwise AND: \texttt{*p \&= v}. \newline {\scriptsize\ttfamily \mnem{ATOM.i32}\imod{.aand} \src{@r0}, \src{r2}, \imm{\#0}} \\[2pt]
\texttt{13} & \texttt{\imod{aor}} & v9--v14 & Atomic bitwise OR: \texttt{*p |= v}. \newline {\scriptsize\ttfamily \mnem{ATOM.i32}\imod{.aor} \src{@r0}, \src{r2}, \imm{\#0}} \\[2pt]
\texttt{14} & \texttt{\imod{axor}} & v9--v14 & Atomic bitwise XOR: \texttt{*p \^= v}. \newline {\scriptsize\ttfamily \mnem{ATOM.i32}\imod{.axor} \src{@r0}, \src{r2}, \imm{\#0}} \\[2pt]
\texttt{15} & \texttt{\imod{axchg}} & v9--v14 & Atomic exchange: \texttt{old = *p; *p = v}. \newline {\scriptsize\ttfamily \mnem{ATOM\_RETURN.i32}\imod{.axchg} \dst{@r0}, \src{r2}, \imm{\#0}} \\[2pt]
\bottomrule
\end{longtable}
\end{center}
\par\textbf{Combined Walkthroughs:}
\par\textbf{Example 1 (Signed max returning previous value (`.asmax`)):}
\codebox{\mnem{ATOM\_RETURN.i32}\imod{.asmax.slot0} \dst{@r0}, \src{[r2:r3]}, \imm{\#0}}
With \texttt{*addr = -5, r0 = 3}: updates \texttt{*addr = 3} and returns previous value \texttt{@r0 = -5}.
\par\textbf{Example 2 (Bitwise XOR without return (`ATOM` + `.axor`)):}
\codebox{\mnem{ATOM.i32}\imod{.axor.slot7} \dst{@r0}, \src{[r2:r3]}, \imm{\#0x4}}
Atomically toggles bits \texttt{*([r2:r3] + 4) \textasciicircum{}= r0} on non-blocking \texttt{slot7}.

\section{Unary Atomic Operation (atom1\_operation\_m)}
\label{enum:atom1_operation_m}
Implicit-operand RMW opcode for \texttt{ATOM1} / \texttt{ATOM1\_RETURN} (\texttt{ainc}, \texttt{adec}, \texttt{asmax1}, \texttt{aumax1}, \texttt{aor1}).

\begin{center}\footnotesize
\begin{longtable}{@{}rll>{\raggedright\arraybackslash}p{0.62\linewidth}@{}}
\toprule
Enc & Value & Arch & Description \& Worked Example \\
\midrule\endhead
\texttt{0} & \texttt{\imod{ainc}} & v9--v14 & Unary atomic increment: \texttt{*p += 1}. \newline {\scriptsize\ttfamily \mnem{ATOM1.i32}\imod{.ainc} \src{r2}, \imm{\#0}} \\[2pt]
\texttt{1} & \texttt{\imod{adec}} & v9--v14 & Unary atomic decrement: \texttt{*p -= 1}. \newline {\scriptsize\ttfamily \mnem{ATOM1.i32}\imod{.adec} \src{r2}, \imm{\#0}} \\[2pt]
\texttt{2} & \texttt{\imod{aumax1}} & v9--v14 & Unary atomic unsigned max with 1: \texttt{*p = max(*p, 1u)}. \newline {\scriptsize\ttfamily \mnem{ATOM1.i32}\imod{.aumax1} \src{r2}, \imm{\#0}} \\[2pt]
\texttt{3} & \texttt{\imod{asmax1}} & v9--v14 & Unary atomic signed max with 1: \texttt{*p = max((int)*p, 1)}. \newline {\scriptsize\ttfamily \mnem{ATOM1.i32}\imod{.asmax1} \src{r2}, \imm{\#0}} \\[2pt]
\texttt{4} & \texttt{\imod{aor1}} & v9--v14 & Unary atomic bitwise OR with 1: \texttt{*p |= 1}. \newline {\scriptsize\ttfamily \mnem{ATOM1.i32}\imod{.aor1} \src{r2}, \imm{\#0}} \\[2pt]
\bottomrule
\end{longtable}
\end{center}
\par\textbf{Combined Walkthroughs:}
\par\textbf{Example 1 (Work-queue ticket increment (`.ainc`)):}
\codebox{\mnem{ATOM1\_RETURN.i32}\imod{.ainc.slot0} \dst{@r0}, \src{[r2:r3]}, \imm{\#0}}
Atomically increments \texttt{*addr} by \texttt{1} and returns the old ticket in \texttt{@r0} without needing an input value register.
\par\textbf{Example 2 (Set dirty flag bit 0 (`.aor1`)):}
\codebox{\mnem{ATOM1.i32}\imod{.aor1.slot7} \src{[r2:r3]}, \imm{\#0}}
Atomically sets bit 0 (\texttt{*addr \textbar{}= 1}) with zero value-register traffic.

\section{Register Conversion Format (register\_file\_format\_general\_m)}
\label{enum:register_file_format_general_m}
Staging register data format for \texttt{LD\_CVT}, \texttt{ST\_CVT}, \texttt{LD\_ATTR*}, \texttt{LD\_VAR*}, \texttt{LD\_TILE}, \texttt{BLEND}.
\par\textbf{CSF ABI Interaction:} Governs staging register conversion against \hyperref[sec:csf-buf-attr]{\texttt{Attribute::format}}, \hyperref[sec:csf-fbd]{\texttt{RenderTarget::internal\_format}}, and \hyperref[sec:csf-zsd-blend]{\texttt{Blend[rt]}}.
\begin{center}\footnotesize
\begin{longtable}{@{}rll>{\raggedright\arraybackslash}p{0.62\linewidth}@{}}
\toprule
Enc & Value & Arch & Description \& Worked Example \\
\midrule\endhead
\texttt{0} & \texttt{\imod{auto32}} & v9--v14 & Format determined by descriptor; write 32-bit GPRs. \newline {\scriptsize\ttfamily \mnem{LD\_ATTR}\imod{.auto32.v4} \dst{@r0:r3}, \src{r4}} \\[2pt]
\texttt{2} & \texttt{\imod{f32}} & v9--v14 & Convert descriptor data to 32-bit IEEE-754 float in GPRs. \newline {\scriptsize\ttfamily \mnem{LD\_ATTR}\imod{.f32.v4} \dst{@r0:r3}, \src{r4}} \\[2pt]
\texttt{3} & \texttt{\imod{f16}} & v9--v14 & Convert descriptor data to packed FP16 pairs in GPRs. \newline {\scriptsize\ttfamily \mnem{LD\_ATTR}\imod{.f16.v4} \dst{@r0:r1}, \src{r4}} \\[2pt]
\texttt{4} & \texttt{\imod{s32}} & v9--v14 & Sign-extend integer attribute/tile data to 32-bit signed int. \newline {\scriptsize\ttfamily \mnem{LD\_ATTR}\imod{.s32.v4} \dst{@r0:r3}, \src{r4}} \\[2pt]
\texttt{5} & \texttt{\imod{s16}} & v9--v14 & Convert integer attribute/tile data to packed signed 16-bit. \newline {\scriptsize\ttfamily \mnem{LD\_ATTR}\imod{.s16.v4} \dst{@r0:r1}, \src{r4}} \\[2pt]
\texttt{6} & \texttt{\imod{u32}} & v9--v14 & Zero-extend integer attribute/tile data to 32-bit unsigned int. \newline {\scriptsize\ttfamily \mnem{LD\_ATTR}\imod{.u32.v4} \dst{@r0:r3}, \src{r4}} \\[2pt]
\texttt{7} & \texttt{\imod{u16}} & v9--v14 & Convert integer attribute/tile data to packed unsigned 16-bit. \newline {\scriptsize\ttfamily \mnem{LD\_ATTR}\imod{.u16.v4} \dst{@r0:r1}, \src{r4}} \\[2pt]
\bottomrule
\end{longtable}
\end{center}
\par\textbf{Combined Walkthroughs:}
\par\textbf{Example 1 (FP32 vertex attribute unpack (`.f32`)):}
\codebox{\mnem{LD\_ATTR\_IMM}\imod{.f32.v4.slot0} \dst{@r0:r3}, \src{r60}, \src{r61}, \imm{\#0}, \imm{\#0}}
Converts e.g. \texttt{R8G8B8A8\_UNORM} vertex buffer data into 4x IEEE FP32 in \texttt{@r0:r3}.
\par\textbf{Example 2 (Packed FP16 tilebuffer fetch (`.f16`)):}
\codebox{\mnem{LD\_TILE}\imod{.f16.v4.slot0} \dst{@r0:r1}, \src{r59}, \src{r60}, \src{u0}}
Packs 4x converted FP16 color channels into just 2 GPRs \texttt{@r0:r1}.

\section{Vector Size (vecsize\_varying\_m)}
\label{enum:vecsize_varying_m}
Number of channels (\texttt{v1}/\texttt{none}, \texttt{v2}, \texttt{v3}, \texttt{v4}) transferred by varying/attribute/tile/blend messages.

\begin{center}\footnotesize
\begin{longtable}{@{}rll>{\raggedright\arraybackslash}p{0.62\linewidth}@{}}
\toprule
Enc & Value & Arch & Description \& Worked Example \\
\midrule\endhead
\texttt{0} & \texttt{\imod{none}} & v9--v14 & 1 component (scalar). \newline {\scriptsize\ttfamily \mnem{BLEND}\imod{} \dst{@r0:r3}, \src{r4}} \\[2pt]
\texttt{1} & \texttt{\imod{v2}} & v9--v14 & 2 components (vec2). \newline {\scriptsize\ttfamily \mnem{BLEND}\imod{.v2} \dst{@r0:r3}, \src{r4}} \\[2pt]
\texttt{2} & \texttt{\imod{v3}} & v9--v14 & 3 components (vec3). \newline {\scriptsize\ttfamily \mnem{BLEND}\imod{.v3} \dst{@r0:r3}, \src{r4}} \\[2pt]
\texttt{3} & \texttt{\imod{v4}} & v9--v14 & 4 components (vec4). \newline {\scriptsize\ttfamily \mnem{BLEND}\imod{.v4} \dst{@r0:r3}, \src{r4}} \\[2pt]
\bottomrule
\end{longtable}
\end{center}
\par\textbf{Combined Walkthroughs:}
\par\textbf{Example 1 (2-channel UV varying load (`.v2`)):}
\codebox{\mnem{LD\_VAR\_IMM}\imod{.f32.v2.center.slot0} \dst{@r0:r1}, \src{r59}, \imm{\#0}, \imm{\#0}}
Interpolates 2 FP32 channels (\texttt{vec2}) into staging pair \texttt{@r0:r1}.
\par\textbf{Example 2 (3-channel normal attribute load (`.v3`)):}
\codebox{\mnem{LD\_ATTR\_IMM}\imod{.f32.v3.slot0} \dst{@r0:r2}, \src{r60}, \src{r61}, \imm{\#0}, \imm{\#1}}
Loads 3 FP32 channels (\texttt{vec3}) into \texttt{@r0:r2}, leaving \texttt{r3} untouched.

\section{Varying Sample Position (sample\_position\_m)}
\label{enum:sample_position_m}
Selects pixel center, centroid, per-sample, or explicit offset interpolation.
\par\textbf{CSF ABI Interaction:} Reads sample coordinates configured by \hyperref[sec:csf-fbd]{\texttt{FramebufferParameters::sample\_locations}} and \hyperref[sec:csf-tiler]{\texttt{DCD0::multisample\_enable}} / \texttt{evaluate\_per\_sample}.
\begin{center}\footnotesize
\begin{longtable}{@{}rll>{\raggedright\arraybackslash}p{0.62\linewidth}@{}}
\toprule
Enc & Value & Arch & Description \& Worked Example \\
\midrule\endhead
\texttt{0} & \texttt{\imod{center}} & v9--v13 & Evaluate varying at pixel center $(0.5, 0.5)$. \newline {\scriptsize\ttfamily \mnem{LD\_VAR\_SPECIAL}\imod{.center} \dst{@r0:r3}, \src{r4}} \\[2pt]
\texttt{1} & \texttt{\imod{centroid}} & v9--v13 & Evaluate varying at MSAA sample centroid (avoids extrapolation outside primitive). \newline {\scriptsize\ttfamily \mnem{LD\_VAR\_SPECIAL}\imod{.centroid} \dst{@r0:r3}, \src{r4}} \\[2pt]
\texttt{2} & \texttt{\imod{sample}} & v9--v13 & Evaluate varying at current MSAA sample position (\texttt{interpolateAtSample}). \newline {\scriptsize\ttfamily \mnem{LD\_VAR\_SPECIAL}\imod{.sample} \dst{@r0:r3}, \src{r4}} \\[2pt]
\texttt{3} & \texttt{\imod{explicit}} & v9--v13 & Evaluate varying at explicit $(x, y)$ offset register (\texttt{interpolateAtOffset}). \newline {\scriptsize\ttfamily \mnem{LD\_VAR\_SPECIAL}\imod{.explicit} \dst{@r0:r3}, \src{r4}} \\[2pt]
\bottomrule
\end{longtable}
\end{center}
\par\textbf{Combined Walkthroughs:}
\par\textbf{Example 1 (MSAA centroid interpolation (`.centroid`)):}
\codebox{\mnem{LD\_VAR\_IMM}\imod{.f32.v4.centroid.store.slot0} \dst{@r0:r3}, \src{r59}, \imm{\#0}, \imm{\#0}}
Evaluates \texttt{centroid in vec4} inside the covered portion of the primitive to avoid extrapolation artifacts.
\par\textbf{Example 2 (Explicit sub-pixel offset (`.explicit`)):}
\codebox{\mnem{LD\_VAR\_IMM}\imod{.f32.v2.explicit.clobber.slot0} \dst{@r0:r1}, \src{r4}, \imm{\#0}, \imm{\#0}}
Implements GLSL \texttt{interpolateAtOffset(v, ofs)} using caller-supplied offset \texttt{r4}.

\section{Varying Cache Update Mode (update\_mode\_m)}
\label{enum:update_mode_m}
Controls hidden barycentric interpolation cache (\texttt{store}, \texttt{retrieve}, \texttt{clobber}).
\par\textbf{CSF ABI Interaction:} Controls the hardware barycentric cache backed by \hyperref[sec:csf-tiler]{\texttt{IDVSSR::VARY\_SIZE}} (\texttt{r48:r49}).
\begin{center}\footnotesize
\begin{longtable}{@{}rll>{\raggedright\arraybackslash}p{0.62\linewidth}@{}}
\toprule
Enc & Value & Arch & Description \& Worked Example \\
\midrule\endhead
\texttt{0} & \texttt{\imod{store}} & v9--v14 & Compute barycentric/varying coordinates and store in hardware varying cache. \newline {\scriptsize\ttfamily \mnem{LD\_VAR}\imod{.store} \dst{@r0:r3}, \src{r4}} \\[2pt]
\texttt{1} & \texttt{\imod{retrieve}} & v9--v14 & Reuse previously cached barycentric coordinates (saves ALU/interpolation power). \newline {\scriptsize\ttfamily \mnem{LD\_VAR}\imod{.retrieve} \dst{@r0:r3}, \src{r4}} \\[2pt]
\texttt{3} & \texttt{\imod{clobber}} & v9--v14 & Compute fresh coordinates and overwrite hardware varying cache. \newline {\scriptsize\ttfamily \mnem{LD\_VAR}\imod{.clobber} \dst{@r0:r3}, \src{r4}} \\[2pt]
\bottomrule
\end{longtable}
\end{center}
\par\textbf{Combined Walkthroughs:}
\par\textbf{Example 1 (First varying saves barycentrics (`.store`), second reuses (`.retrieve`)):}
\codebox{\mnem{LD\_VAR\_IMM}\imod{.f32.v4.center.retrieve.slot0} \dst{@r4:r7}, \src{r59}, \imm{\#0}, \imm{\#1}}
Reuses barycentric weights saved by a prior \texttt{.store} varying load at the same sample position, saving interpolation cycles.
\par\textbf{Example 2 (One-off special varying without polluting cache (`.clobber`)):}
\codebox{\mnem{LD\_VAR\_SPECIAL}\imod{.frag\_w.center.clobber.slot0} \dst{@r0}, \src{r59}}
Computes \texttt{frag\_w} directly when no subsequent \texttt{.retrieve} depends on the cache.

\section{Texture Dimensionality (tex\_dimensionality\_m)}
\label{enum:tex_dimensionality_m}
Selects 1D, 2D, 3D, or Cube texture addressing.
\par\textbf{CSF ABI Interaction:} Must match \hyperref[sec:csf-tex-samp]{\texttt{Texture::dimension}} (word 0 bits 4:6) fetched via \texttt{SRT\_2}/\texttt{SRT\_0}.
\begin{center}\footnotesize
\begin{longtable}{@{}rll>{\raggedright\arraybackslash}p{0.62\linewidth}@{}}
\toprule
Enc & Value & Arch & Description \& Worked Example \\
\midrule\endhead
\texttt{0} & \texttt{\imod{tex1d}} & v9--v14 & 1D texture coordinate $(u)$. \newline {\scriptsize\ttfamily \mnem{TEX\_SINGLE}\imod{.tex1d} \dst{@r0:r3}, \src{@r4}} \\[2pt]
\texttt{1} & \texttt{\imod{tex2d}} & v9--v14 & 2D texture coordinates $(u, v)$. \newline {\scriptsize\ttfamily \mnem{TEX\_SINGLE}\imod{.tex2d} \dst{@r0:r3}, \src{@r4:r5}} \\[2pt]
\texttt{2} & \texttt{\imod{tex3d}} & v9--v14 & 3D volumetric texture coordinates $(u, v, w)$. \newline {\scriptsize\ttfamily \mnem{TEX\_SINGLE}\imod{.tex3d} \dst{@r0:r3}, \src{@r4:r6}} \\[2pt]
\texttt{3} & \texttt{\imod{cube}} & v9--v14 & Cubemap direction/face-projected coordinates. \newline {\scriptsize\ttfamily \mnem{TEX\_SINGLE}\imod{.cube} \dst{@r0:r3}, \src{@r4:r6}} \\[2pt]
\bottomrule
\end{longtable}
\end{center}
\par\textbf{Combined Walkthroughs:}
\par\textbf{Example 1 (Hardware-projected cubemap (v11+ `.cube.project`)):}
\codebox{\mnem{TEX\_SINGLE}\imod{.cube.project.dst32.rgba.slot0} \dst{@r0:r3}, \src{@r4:r6}, \src{u0}}
Samples a cubemap from 3D direction \texttt{@r4:r6} directly in hardware (eliminates v9--v10 \texttt{CUBEFACE1/2} + \texttt{CUBE\_SSEL/TSEL}).
\par\textbf{Example 2 (3D volumetric texture (`.tex3d`)):}
\codebox{\mnem{TEX\_SINGLE}\imod{.tex3d.computed.dst32.rgba.slot0} \dst{@r0:r3}, \src{@r4:r6}, \src{u0}}
Samples a 3D LUT using trilinear \texttt{(u, v, w)} coordinates in \texttt{@r4:r6}.

\section{Texture LOD Mode (tex\_lod\_mode\_m)}
\label{enum:tex_lod_mode_m}
Controls mipmap level-of-detail selection (\texttt{computed}, \texttt{explicit}, \texttt{zero}, \texttt{computed\_bias}, etc.).
\par\textbf{CSF ABI Interaction:} Interacts with \hyperref[sec:csf-tex-samp]{\texttt{Sampler}} LOD clamp (\texttt{minimum\_lod}/\texttt{maximum\_lod}), \texttt{lod\_bias}, and \texttt{Texture::levels} via \texttt{SRT\_2}.
\begin{center}\footnotesize
\begin{longtable}{@{}rll>{\raggedright\arraybackslash}p{0.62\linewidth}@{}}
\toprule
Enc & Value & Arch & Description \& Worked Example \\
\midrule\endhead
\texttt{0} & \texttt{\imod{none}} & v9 & Implicit zero LOD (base mip level 0, no derivatives). \newline {\scriptsize\ttfamily \mnem{TEX\_SINGLE}\imod{.tex2d} \dst{@r0:r3}, \src{@r4:r5}} \\[2pt]
\texttt{1} & \texttt{\imod{computed}} & v9 & Hardware 2x2 quad finite-difference LOD (\texttt{texture()}). \newline {\scriptsize\ttfamily \mnem{TEX\_SINGLE}\imod{.tex2d.computed} \dst{@r0:r3}, \src{@r4:r5}} \\[2pt]
\texttt{3} & \texttt{\imod{computed\_force\_delta}} & v10--v14 & Computed LOD forcing quad delta evaluation even in helper lanes. \newline {\scriptsize\ttfamily \mnem{TEX\_SINGLE}\imod{.tex2d.computed\_force\_delta} \dst{@r0:r3}, \src{@r4:r5}} \\[2pt]
\texttt{4} & \texttt{\imod{explicit}} & v9 & Explicit LOD supplied in staging register (\texttt{textureLod()}). \newline {\scriptsize\ttfamily \mnem{TEX\_SINGLE}\imod{.tex2d.explicit} \dst{@r0:r3}, \src{@r4:r6}} \\[2pt]
\texttt{5} & \texttt{\imod{computed\_bias}} & v9 & Computed quad LOD plus shader bias (\texttt{texture(..., bias)}). \newline {\scriptsize\ttfamily \mnem{TEX\_SINGLE}\imod{.tex2d.computed\_bias} \dst{@r0:r3}, \src{@r4:r6}} \\[2pt]
\texttt{6} & \texttt{\imod{grdesc}} & v9 & Explicit $dP/dx, dP/dy$ gradients in staging regs (\texttt{textureGrad()}). \newline {\scriptsize\ttfamily \mnem{TEX\_SINGLE}\imod{.tex2d.grdesc} \dst{@r0:r3}, \src{@r4:r9}} \\[2pt]
\texttt{7} & \texttt{\imod{computed\_bias\_force\_delta}} & v10--v14 & Computed LOD + bias with forced quad delta evaluation. \newline {\scriptsize\ttfamily \mnem{TEX\_SINGLE}\imod{.tex2d.computed\_bias\_force\_delta} \dst{@r0:r3}, \src{@r4:r6}} \\[2pt]
\bottomrule
\end{longtable}
\end{center}
\par\textbf{Combined Walkthroughs:}
\par\textbf{Example 1 (Explicit mipmap LOD (`.explicit`)):}
\codebox{\mnem{TEX\_SINGLE}\imod{.tex2d.explicit.dst32.rgba.slot0} \dst{@r0:r3}, \src{@r4:r6}, \src{u0}}
Reads explicit 8.8 fixed-point LOD from \texttt{@r6} (implements GLSL \texttt{textureLod(samp, uv, lod)}).
\par\textbf{Example 2 (Computed quad LOD + shader bias (`.computed\_bias`)):}
\codebox{\mnem{TEX\_SINGLE}\imod{.tex2d.computed\_bias.dst32.rgba.slot0} \dst{@r0:r3}, \src{@r4:r6}, \src{u0}}
Computes 2x2 finite-difference LOD and adds bias \texttt{@r6} (implements \texttt{texture(samp, uv, bias)}).

\section{Texture Gather Component (tex\_gather\_component\_m)}
\label{enum:tex_gather_component_m}
Selects channel (\texttt{gather4\_r}, \texttt{gather4\_g}, \texttt{gather4\_b}, \texttt{gather4\_a}) for \texttt{TEX\_GATHER}.
\par\textbf{CSF ABI Interaction:} Gathers the selected channel from \hyperref[sec:csf-tex-samp]{\texttt{Texture}} $\to$ \texttt{GenericPlane} (Morton 16$\times$16 detiled) via \texttt{SRT\_2}.
\begin{center}\footnotesize
\begin{longtable}{@{}rll>{\raggedright\arraybackslash}p{0.62\linewidth}@{}}
\toprule
Enc & Value & Arch & Description \& Worked Example \\
\midrule\endhead
\texttt{0} & \texttt{\imod{gather4\_r}} & v9--v14 & Gather 2x2 footprint of channel R (\texttt{textureGather(..., 0)}). \newline {\scriptsize\ttfamily \mnem{TEX\_GATHER}\imod{.gather4\_r} \dst{@r0:r3}, \src{@r4:r5}} \\[2pt]
\texttt{1} & \texttt{\imod{gather4\_g}} & v9--v14 & Gather 2x2 footprint of channel G (\texttt{textureGather(..., 1)}). \newline {\scriptsize\ttfamily \mnem{TEX\_GATHER}\imod{.gather4\_g} \dst{@r0:r3}, \src{@r4:r5}} \\[2pt]
\texttt{2} & \texttt{\imod{gather4\_b}} & v9--v14 & Gather 2x2 footprint of channel B (\texttt{textureGather(..., 2)}). \newline {\scriptsize\ttfamily \mnem{TEX\_GATHER}\imod{.gather4\_b} \dst{@r0:r3}, \src{@r4:r5}} \\[2pt]
\texttt{3} & \texttt{\imod{gather4\_a}} & v9--v14 & Gather 2x2 footprint of channel A (\texttt{textureGather(..., 3)}). \newline {\scriptsize\ttfamily \mnem{TEX\_GATHER}\imod{.gather4\_a} \dst{@r0:r3}, \src{@r4:r5}} \\[2pt]
\bottomrule
\end{longtable}
\end{center}
\par\textbf{Combined Walkthroughs:}
\par\textbf{Example 1 (Green-channel 2x2 gather (`.gather4\_g`)):}
\codebox{\mnem{TEX\_GATHER}\imod{.tex2d.gather4\_g.dst32.rgba.slot0} \dst{@r0:r3}, \src{@r4:r5}, \src{u0}}
Implements GLSL \texttt{textureGather(samp, uv, 1)}, returning the 4 green texels in counter-clockwise order.
\par\textbf{Example 2 (Alpha-channel 2x2 gather (`.gather4\_a`)):}
\codebox{\mnem{TEX\_GATHER}\imod{.tex2d.gather4\_a.dst32.rgba.slot0} \dst{@r0:r3}, \src{@r4:r5}, \src{u0}}
Implements GLSL \texttt{textureGather(samp, uv, 3)} for custom alpha-cutout filtering.


\section{Additional Specialized Modifier Enumerations (78 Enums)}
The summary table below indexes all remaining secondary modifier enumerations in \texttt{isa-v9-v14.xml}; each enumeration is walked through entry-by-entry with a concrete worked example in the subsections immediately following.
\begin{center}\footnotesize
\begin{longtable}{@{}p{0.36\linewidth}>{\raggedright\arraybackslash}p{0.60\linewidth}@{}}
\toprule
Enum Name (Jump) & Values (Encoding) \\
\midrule\endhead
\hyperref[enum:access_store_m]{\texttt{access\_store\_m}} & \texttt{none(0), istream(1), estream(2), force(3)} \\
\hyperref[enum:branch_conservative_m]{\texttt{branch\_conservative\_m}} & \texttt{none(0), conservative(1)} \\
\hyperref[enum:clper_or_lane_select_m]{\texttt{clper\_or\_lane\_select\_m}} & \texttt{lanes04(0), lanes15(1), lanes26(2), lanes37(3), lanes8C(4), lanes9D(5), lanesAE(6), lanesBF(7)} \\
\hyperref[enum:clz_mask_m]{\texttt{clz\_mask\_m}} & \texttt{none(0), mask(1)} \\
\hyperref[enum:dest_width_single_lane_16_bit_m]{\texttt{dest\_width\_single\_lane\_16\_bit\_m}} & \texttt{h0(1), h1(2), none(3)} \\
\hyperref[enum:encoded_sample_and_update_m]{\texttt{encoded\_sample\_and\_update\_m}} & \texttt{center\_store(0), centroid\_store(1), sample\_store(2), explicit\_store(3), center\_clobber(4), sample\_clobber(6), retrieve(7)} \\
\hyperref[enum:encoded_sample_position_none_m]{\texttt{encoded\_sample\_position\_none\_m}} & \texttt{center\_or\_none(0), centroid(1), sample(2), explicit(3)} \\
\hyperref[enum:encoded_update_mode_special_m]{\texttt{encoded\_update\_mode\_special\_m}} & \texttt{store(0), retrieve(1), clobber\_or\_none(3)} \\
\hyperref[enum:fadd_sticky_m]{\texttt{fadd\_sticky\_m}} & \texttt{none(0), sticky(1)} \\
\hyperref[enum:fexp_adj_m]{\texttt{fexp\_adj\_m}} & \texttt{none(0), low(2)} \\
\hyperref[enum:fma_rscale_round_m]{\texttt{fma\_rscale\_round\_m}} & \texttt{none(0), round\_zero(1)} \\
\hyperref[enum:fpow_sc_det_function_m]{\texttt{fpow\_sc\_det\_function\_m}} & \texttt{pow(0), powr(1)} \\
\hyperref[enum:iaddc_return_carry_m]{\texttt{iaddc\_return\_carry\_m}} & \texttt{none(0), c(1)} \\
\hyperref[enum:jump_divergence_mode_m]{\texttt{jump\_divergence\_mode\_m}} & \texttt{none(0), diverge(1)} \\
\hyperref[enum:jump_stack_mode_m]{\texttt{jump\_stack\_mode\_m}} & \texttt{return(0), call(1), none(2), replace(3)} \\
\hyperref[enum:lane_all_m]{\texttt{lane\_all\_m}} & \texttt{b0(0), b1(1), b2(2), b3(3), h0(4), h1(5), w0(6), d0(7)} \\
\hyperref[enum:ld_gclk_source_m]{\texttt{ld\_gclk\_source\_m}} & \texttt{cycle\_counter(6), system\_timestamp(7)} \\
\hyperref[enum:ld_var_buf_encoded_update_mode_none_m]{\texttt{ld\_var\_buf\_encoded\_update\_mode\_none\_m}} & \texttt{store(0), retrieve\_or\_none(1), clobber(3)} \\
\hyperref[enum:ld_var_buf_flat_format_m]{\texttt{ld\_var\_buf\_flat\_format\_m}} & \texttt{flat32(2), flat16(3)} \\
\hyperref[enum:ld_var_buf_source_format_m]{\texttt{ld\_var\_buf\_source\_format\_m}} & \texttt{src\_flat32(0), src\_flat16(1), src\_f32(2), src\_f16(3)} \\
\hyperref[enum:ldexp_round_m]{\texttt{ldexp\_round\_m}} & \texttt{none(0), round\_up(1), round\_down(2), round\_zero(3), round\_na(4)} \\
\hyperref[enum:load_lane_16_m]{\texttt{load\_lane\_16\_m}} & \texttt{h0(0), h1(1), w0(2), d0(3)} \\
\hyperref[enum:load_lane_32_m]{\texttt{load\_lane\_32\_m}} & \texttt{none(0), d0(1)} \\
\hyperref[enum:ls_multi_sr_count_m]{\texttt{ls\_multi\_sr\_count\_m}} & \texttt{i32(0), i64(1), i96(2), i128(3), i160(4), i192(5), i224(6), i256(7), ... (64 total)} \\
\hyperref[enum:ls_z_last_read_m]{\texttt{ls\_z\_last\_read\_m}} & \texttt{none(0), last\_read(1)} \\
\hyperref[enum:ls_z_stencil_m]{\texttt{ls\_z\_stencil\_m}} & \texttt{none(0), z\_stencil(1)} \\
\hyperref[enum:message_slot_index_m]{\texttt{message\_slot\_index\_m}} & \texttt{slot0(0), slot1(1), slot2(2), slot7(7)} \\
\hyperref[enum:round_extended_m]{\texttt{round\_extended\_m}} & \texttt{none(0), round\_up(1), round\_down(2), round\_zero(3), round\_na(4)} \\
\hyperref[enum:sample_position_none_m]{\texttt{sample\_position\_none\_m}} & \texttt{center(0), centroid(1), sample(2), explicit(3), none(4)} \\
\hyperref[enum:segadd_preserve_null_m]{\texttt{segadd\_preserve\_null\_m}} & \texttt{none(0), preserve\_null(1)} \\
\hyperref[enum:sign_extend_or_none_compressed_encoding_m]{\texttt{sign\_extend\_or\_none\_compressed\_encoding\_m}} & \texttt{sext(0), zext\_or\_none(1)} \\
\hyperref[enum:stencil_m]{\texttt{stencil\_m}} & \texttt{none(0), stencil(1)} \\
\hyperref[enum:submatrix_sel_m]{\texttt{submatrix\_sel\_m}} & \texttt{f0(0), f1(1), none(2)} \\
\hyperref[enum:table_flog_mode_m]{\texttt{table\_flog\_mode\_m}} & \texttt{red(0), base2(1), natural(2)} \\
\hyperref[enum:table_flog_precision_m]{\texttt{table\_flog\_precision\_m}} & \texttt{none(0), high(1), low(2)} \\
\hyperref[enum:table_sincos_offset_m]{\texttt{table\_sincos\_offset\_m}} & \texttt{none(0), offset(1)} \\
\hyperref[enum:table_sincos_operation_m]{\texttt{table\_sincos\_operation\_m}} & \texttt{cos(0), sin(1)} \\
\hyperref[enum:table_sincos_scale_m]{\texttt{table\_sincos\_scale\_m}} & \texttt{none(0), scale(1)} \\
\hyperref[enum:tex_array_enable_m]{\texttt{tex\_array\_enable\_m}} & \texttt{none(0), array(1)} \\
\hyperref[enum:tex_compare_enable_m]{\texttt{tex\_compare\_enable\_m}} & \texttt{none(0), compare(1)} \\
\hyperref[enum:tex_coordinate_mode_m]{\texttt{tex\_coordinate\_mode\_m}} & \texttt{float\_coordinates(0), integer\_coordinates(1)} \\
\hyperref[enum:tex_coordinate_or_derivative_m]{\texttt{tex\_coordinate\_or\_derivative\_m}} & \texttt{none(0), force\_delta(1), derivative(2)} \\
\hyperref[enum:tex_derivative_enable_m]{\texttt{tex\_derivative\_enable\_m}} & \texttt{none(0), derivative(1)} \\
\hyperref[enum:tex_dual_array_enable_m]{\texttt{tex\_dual\_array\_enable\_m}} & \texttt{none(0), array(1)} \\
\hyperref[enum:tex_dual_lod_mode_m]{\texttt{tex\_dual\_lod\_mode\_m}} & \texttt{none(0), computed(1), computed\_force\_delta(3), explicit(4), computed\_bias(5), grdesc(6), computed\_bias\_force\_delta(7)} \\
\hyperref[enum:tex_dual_primary_width_m]{\texttt{tex\_dual\_primary\_width\_m}} & \texttt{primary\_dst16(0), primary\_f16(0), primary\_s16(0), primary\_u16(0), primary\_dst32(1), primary\_f32(1), primary\_s32(1), primary\_u32(1)} \\
\hyperref[enum:tex_dual_primary_write_mask_m]{\texttt{tex\_dual\_primary\_write\_mask\_m}} & \texttt{primary\_r(1), primary\_g(2), primary\_rg(3), primary\_b(4), primary\_rb(5), primary\_gb(6), primary\_rgb(7), primary\_a(8), ... (15 total)} \\
\hyperref[enum:tex_dual_secondary_width_m]{\texttt{tex\_dual\_secondary\_width\_m}} & \texttt{secondary\_dst16(0), secondary\_f16(0), secondary\_s16(0), secondary\_u16(0), secondary\_dst32(1), secondary\_f32(1), secondary\_s32(1), secondary\_u32(1)} \\
\hyperref[enum:tex_dual_secondary_write_mask_m]{\texttt{tex\_dual\_secondary\_write\_mask\_m}} & \texttt{secondary\_r(1), secondary\_rg(3), secondary\_rgb(7), secondary\_rgba(15)} \\
\hyperref[enum:tex_footprint_m]{\texttt{tex\_footprint\_m}} & \texttt{footprint\_2x2(0)} \\
\hyperref[enum:tex_force_delta_enable_m]{\texttt{tex\_force\_delta\_enable\_m}} & \texttt{none(0), force\_delta(1)} \\
\hyperref[enum:tex_gradient_skip_m]{\texttt{tex\_gradient\_skip\_m}} & \texttt{none(0), skip(1)} \\
\hyperref[enum:tex_gradient_width_m]{\texttt{tex\_gradient\_width\_m}} & \texttt{dst16(0), f16(0), s16(0), u16(0), dst32(1), f32(1), s32(1), u32(1)} \\
\hyperref[enum:tex_gradient_write_mask_m]{\texttt{tex\_gradient\_write\_mask\_m}} & \texttt{r(1), g(2), rg(3), b(4), rb(5), gb(6), rgb(7), a(8), ... (15 total)} \\
\hyperref[enum:tex_lod_bias_disable_m]{\texttt{tex\_lod\_bias\_disable\_m}} & \texttt{none(0), disable\_lod\_bias(1)} \\
\hyperref[enum:tex_lod_clamp_disable_m]{\texttt{tex\_lod\_clamp\_disable\_m}} & \texttt{none(0), disable\_lod\_clamp(1)} \\
\hyperref[enum:tex_projection_enable_m]{\texttt{tex\_projection\_enable\_m}} & \texttt{none(0), project(1)} \\
\hyperref[enum:tex_skip_m]{\texttt{tex\_skip\_m}} & \texttt{none(0), skip(1)} \\
\hyperref[enum:tex_texel_offset_m]{\texttt{tex\_texel\_offset\_m}} & \texttt{none(0), texel\_offset(1)} \\
\hyperref[enum:tex_texel_offset_si_m]{\texttt{tex\_texel\_offset\_si\_m}} & \texttt{none(0), texel\_offset(1)} \\
\hyperref[enum:tex_wide_indices_enable_m]{\texttt{tex\_wide\_indices\_enable\_m}} & \texttt{none(0), wide(1)} \\
\hyperref[enum:tex_width_m]{\texttt{tex\_width\_m}} & \texttt{dst16(0), f16(0), s16(0), u16(0), dst32(1), f32(1), s32(1), u32(1)} \\
\hyperref[enum:tex_write_mask_m]{\texttt{tex\_write\_mask\_m}} & \texttt{r(1), g(2), rg(3), b(4), rb(5), gb(6), rgb(7), a(8), ... (15 total)} \\
\hyperref[enum:update_mode_none_m]{\texttt{update\_mode\_none\_m}} & \texttt{store(0), retrieve(1), clobber(3), none(4)} \\
\hyperref[enum:update_mode_special_m]{\texttt{update\_mode\_special\_m}} & \texttt{store(0), retrieve(1), clobber(3), none(4)} \\
\hyperref[enum:v4u8_m]{\texttt{v4u8\_m}} & \texttt{none(0), v4u8(1)} \\
\hyperref[enum:var_tex_array_enable_m]{\texttt{var\_tex\_array\_enable\_m}} & \texttt{none(0), array(1)} \\
\hyperref[enum:var_tex_buf_source_format_m]{\texttt{var\_tex\_buf\_source\_format\_m}} & \texttt{src\_f32(0), src\_f16(1)} \\
\hyperref[enum:var_tex_compare_enable_m]{\texttt{var\_tex\_compare\_enable\_m}} & \texttt{none(0), compare(1)} \\
\hyperref[enum:var_tex_dimensionality_m]{\texttt{var\_tex\_dimensionality\_m}} & \texttt{tex1d(0), tex2d(1), tex3d(2), cube(3)} \\
\hyperref[enum:var_tex_dual_array_enable_m]{\texttt{var\_tex\_dual\_array\_enable\_m}} & \texttt{none(0), array(1)} \\
\hyperref[enum:var_tex_dual_lod_mode_m]{\texttt{var\_tex\_dual\_lod\_mode\_m}} & \texttt{none(0), computed(1)} \\
\hyperref[enum:var_tex_lod_mode_m]{\texttt{var\_tex\_lod\_mode\_m}} & \texttt{none(0), computed(1)} \\
\hyperref[enum:var_tex_projection_enable_m]{\texttt{var\_tex\_projection\_enable\_m}} & \texttt{none(0), project(1)} \\
\hyperref[enum:varying_name_m]{\texttt{varying\_name\_m}} & \texttt{point(0), bary(1), frag\_w(2), frag\_z(3)} \\
\hyperref[enum:vertex_m]{\texttt{vertex\_m}} & \texttt{none(0), vertex0(1), vertex1(2), vertex2(3)} \\
\hyperref[enum:vertex_type_m]{\texttt{vertex\_type\_m}} & \texttt{none(0), explicit\_vertex(1)} \\
\hyperref[enum:z_m]{\texttt{z\_m}} & \texttt{none(0), z(1)} \\
\bottomrule
\end{longtable}
\end{center}

\subsection{Access Store (\texttt{access\_store\_m})}
\label{enum:access_store_m}
Used by \texttt{STORE} and related instructions (4 encoding values):
\par\textbf{CSF ABI Interaction:} In IDVS graphics, \texttt{STORE.istream} routes \texttt{gl\_Position} to \hyperref[sec:csf-tiler]{\texttt{TILER\_CTX::geometry\_buffer}} (\texttt{r40:r41}) and \texttt{STORE.estream} writes per-vertex varyings into \hyperref[sec:csf-tiler]{\texttt{IDVSSR::VARY\_SIZE}} (\texttt{r48:r49}).
\begin{center}\footnotesize
\begin{longtable}{@{}rll>{\raggedright\arraybackslash}p{0.62\linewidth}@{}}
\toprule
Enc & Value & Arch & Description \& Worked Example \\
\midrule\endhead
\texttt{0} & \texttt{\imod{none}} & v9--v14 & Default write-back cached store policy. \newline {\scriptsize\ttfamily \mnem{STORE.i32} \src{@r0}, \src{r2}, \imm{\#0}} \\[2pt]
\texttt{1} & \texttt{\imod{istream}} & v9--v14 & Inner-streaming store hint: avoid polluting L1 cache. \newline {\scriptsize\ttfamily \mnem{STORE.i32}\imod{.istream} \src{@r0}, \src{r2}, \imm{\#0}} \\[2pt]
\texttt{2} & \texttt{\imod{estream}} & v9--v14 & Outer-streaming store hint: stream past L2 toward memory. \newline {\scriptsize\ttfamily \mnem{STORE.i32}\imod{.estream} \src{@r0}, \src{r2}, \imm{\#0}} \\[2pt]
\texttt{3} & \texttt{\imod{force}} & v9--v14 & Force write-through/coherent visibility (coherent SSBO/image store). \newline {\scriptsize\ttfamily \mnem{STORE.i32}\imod{.force} \src{@r0}, \src{r2}, \imm{\#0}} \\[2pt]
\bottomrule
\end{longtable}
\end{center}
\par\textbf{Worked Example:} \emph{Coherent SSBO store vs streaming store}
\codebox{\mnem{STORE.i32}\imod{.force.slot7} \dst{@r0}, \src{[r2:r3]}, \imm{\#0}}
With \texttt{.force}, bypasses write-back caching so other workgroups immediately observe the store; \texttt{.estream} instead streams bulk writes past L2.

\subsection{Branch Conservative (\texttt{branch\_conservative\_m})}
\label{enum:branch_conservative_m}
Used by \texttt{BRANCH} and related instructions (2 encoding values):

\begin{center}\footnotesize
\begin{longtable}{@{}rll>{\raggedright\arraybackslash}p{0.62\linewidth}@{}}
\toprule
Enc & Value & Arch & Description \& Worked Example \\
\midrule\endhead
\texttt{0} & \texttt{\imod{none}} & v9--v14 & Default behavior (flag disabled for \texttt{branch\_conservative\_m}). \newline {\scriptsize\ttfamily \mnem{BRANCH}\imod{} \dst{r0}, \src{r1}} \\[2pt]
\texttt{1} & \texttt{\imod{conservative}} & v9--v14 & Conservative warp branch hint when lanes may diverge. \newline {\scriptsize\ttfamily \mnem{BRANCH}\imod{.conservative} \dst{r0}, \src{r1}} \\[2pt]
\bottomrule
\end{longtable}
\end{center}
\par\textbf{Worked Example:} \emph{Conservative loop back-edge branch}
\codebox{\mnem{BRANCH}\imod{.conservative} \src{r0}, \imm{\#-0x18}}
Hints to the warp scheduler that divergence across lanes is non-trivial so reconvergence tracking stays conservative.

\subsection{Clper Or Lane Select (\texttt{clper\_or\_lane\_select\_m})}
\label{enum:clper_or_lane_select_m}
Used by \texttt{CLPER\_OR} and related instructions (8 encoding values):

\begin{center}\footnotesize
\begin{longtable}{@{}rll>{\raggedright\arraybackslash}p{0.62\linewidth}@{}}
\toprule
Enc & Value & Arch & Description \& Worked Example \\
\midrule\endhead
\texttt{0} & \texttt{\imod{lanes04}} & v14 & Select lane pair \texttt{lanes04} for cross-lane bitwise OR reduction. \newline {\scriptsize\ttfamily \mnem{CLPER\_OR}\imod{.lanes04} \dst{r0}, \src{r1}, \src{r2}} \\[2pt]
\texttt{1} & \texttt{\imod{lanes15}} & v14 & Select lane pair \texttt{lanes15} for cross-lane bitwise OR reduction. \newline {\scriptsize\ttfamily \mnem{CLPER\_OR}\imod{.lanes15} \dst{r0}, \src{r1}, \src{r2}} \\[2pt]
\texttt{2} & \texttt{\imod{lanes26}} & v14 & Select lane pair \texttt{lanes26} for cross-lane bitwise OR reduction. \newline {\scriptsize\ttfamily \mnem{CLPER\_OR}\imod{.lanes26} \dst{r0}, \src{r1}, \src{r2}} \\[2pt]
\texttt{3} & \texttt{\imod{lanes37}} & v14 & Select lane pair \texttt{lanes37} for cross-lane bitwise OR reduction. \newline {\scriptsize\ttfamily \mnem{CLPER\_OR}\imod{.lanes37} \dst{r0}, \src{r1}, \src{r2}} \\[2pt]
\texttt{4} & \texttt{\imod{lanes8C}} & v14 & Select lane pair \texttt{lanes8C} for cross-lane bitwise OR reduction. \newline {\scriptsize\ttfamily \mnem{CLPER\_OR}\imod{.lanes8C} \dst{r0}, \src{r1}, \src{r2}} \\[2pt]
\texttt{5} & \texttt{\imod{lanes9D}} & v14 & Select lane pair \texttt{lanes9D} for cross-lane bitwise OR reduction. \newline {\scriptsize\ttfamily \mnem{CLPER\_OR}\imod{.lanes9D} \dst{r0}, \src{r1}, \src{r2}} \\[2pt]
\texttt{6} & \texttt{\imod{lanesAE}} & v14 & Select lane pair \texttt{lanesAE} for cross-lane bitwise OR reduction. \newline {\scriptsize\ttfamily \mnem{CLPER\_OR}\imod{.lanesAE} \dst{r0}, \src{r1}, \src{r2}} \\[2pt]
\texttt{7} & \texttt{\imod{lanesBF}} & v14 & Select lane pair \texttt{lanesBF} for cross-lane bitwise OR reduction. \newline {\scriptsize\ttfamily \mnem{CLPER\_OR}\imod{.lanesBF} \dst{r0}, \src{r1}, \src{r2}} \\[2pt]
\bottomrule
\end{longtable}
\end{center}
\par\textbf{Worked Example:} \emph{Fused cross-lane shuffle + OR (v14)}
\codebox{\mnem{CLPER\_OR}\imod{.subgroup16.lanes04} \dst{r0}, \src{r1}, \src{r2}\imod{.b0}, \src{r3}}
Selects the \texttt{lanes04} pair routing in v14 \texttt{CLPER\_OR}, permutes \texttt{r1}, and bitwise-ORs with \texttt{r3} in a single cycle.

\subsection{Clz Mask (\texttt{clz\_mask\_m})}
\label{enum:clz_mask_m}
Used by \texttt{CLZ} and related instructions (2 encoding values):

\begin{center}\footnotesize
\begin{longtable}{@{}rll>{\raggedright\arraybackslash}p{0.62\linewidth}@{}}
\toprule
Enc & Value & Arch & Description \& Worked Example \\
\midrule\endhead
\texttt{0} & \texttt{\imod{none}} & v9--v14 & Default behavior (flag disabled for \texttt{clz\_mask\_m}). \newline {\scriptsize\ttfamily \mnem{CLZ}\imod{} \dst{r0}, \src{r1}} \\[2pt]
\texttt{1} & \texttt{\imod{mask}} & v9--v14 & Apply lane/bitmask variant of count-leading-zeros. \newline {\scriptsize\ttfamily \mnem{CLZ}\imod{.mask} \dst{r0}, \src{r1}} \\[2pt]
\bottomrule
\end{longtable}
\end{center}
\par\textbf{Worked Example:} \emph{GLSL `findMSB(0) == -1` via `.mask`}
\codebox{\mnem{CLZ.u32}\imod{.mask} \dst{r0}, \src{r1}}
When \texttt{r1 = 0}: default \texttt{none} returns \texttt{32}, whereas \texttt{.mask} returns \texttt{0xFFFFFFFF} (\texttt{-1}). For \texttt{r1 = 0x00010000}, both return \texttt{15}.

\subsection{Dest Width Single Lane 16 Bit (\texttt{dest\_width\_single\_lane\_16\_bit\_m})}
\label{enum:dest_width_single_lane_16_bit_m}
Used by \texttt{FRCP} and related instructions (3 encoding values):

\begin{center}\footnotesize
\begin{longtable}{@{}rll>{\raggedright\arraybackslash}p{0.62\linewidth}@{}}
\toprule
Enc & Value & Arch & Description \& Worked Example \\
\midrule\endhead
\texttt{1} & \texttt{\imod{h0}} & v9--v10 & Write low 16-bit halfword \texttt{dst[15:0]}; zero-fill upper \texttt{dst[31:16]}. \newline {\scriptsize\ttfamily \mnem{FRCP.f16} \dst{r0}\imod{.h0}, \src{r1}} \\[2pt]
\texttt{2} & \texttt{\imod{h1}} & v9--v10 & Write high 16-bit halfword \texttt{dst[31:16]}; preserve/mask low halfword per lane-mask rule. \newline {\scriptsize\ttfamily \mnem{FRCP.f16} \dst{r0}\imod{.h1}, \src{r1}} \\[2pt]
\texttt{3} & \texttt{\imod{none}} & v9--v10 & Write full 32-bit destination register (both halfwords). \newline {\scriptsize\ttfamily \mnem{FRCP.f32} \dst{r0}, \src{r1}} \\[2pt]
\bottomrule
\end{longtable}
\end{center}
\par\textbf{Worked Example:} \emph{Single-lane 16-bit writeback preserving peer lane}
\codebox{\mnem{F32\_TO\_F16} \dst{r0}\imod{.h1}, \src{r1}}
With \texttt{r1 = 2.0f}, \texttt{r0\_old = 0x00003C00}: \texttt{.h1} writes \texttt{0x4000} to \texttt{r0[31:16]} while preserving \texttt{r0[15:0]} -\textgreater{} \texttt{r0 = 0x40003C00}.

\subsection{Encoded Sample And Update (\texttt{encoded\_sample\_and\_update\_m})}
\label{enum:encoded_sample_and_update_m}
Used by \texttt{VAR\_TEX\_BUF\_DUAL} and related instructions (7 encoding values):

\begin{center}\footnotesize
\begin{longtable}{@{}rll>{\raggedright\arraybackslash}p{0.62\linewidth}@{}}
\toprule
Enc & Value & Arch & Description \& Worked Example \\
\midrule\endhead
\texttt{0} & \texttt{\imod{center\_store}} & v9--v14 & Evaluate at pixel center and store barycentrics to varying cache. \newline {\scriptsize\ttfamily \mnem{VAR\_TEX\_BUF\_DUAL}\imod{.center\_store} \dst{@r0:r3}, \src{r4}} \\[2pt]
\texttt{1} & \texttt{\imod{centroid\_store}} & v9--v14 & Evaluate at centroid and store barycentrics to varying cache. \newline {\scriptsize\ttfamily \mnem{VAR\_TEX\_BUF\_DUAL}\imod{.centroid\_store} \dst{@r0:r3}, \src{r4}} \\[2pt]
\texttt{2} & \texttt{\imod{sample\_store}} & v9--v14 & Evaluate at per-sample position and store barycentrics. \newline {\scriptsize\ttfamily \mnem{VAR\_TEX\_BUF\_DUAL}\imod{.sample\_store} \dst{@r0:r3}, \src{r4}} \\[2pt]
\texttt{3} & \texttt{\imod{explicit\_store}} & v9--v14 & Evaluate at explicit offset and store barycentrics. \newline {\scriptsize\ttfamily \mnem{VAR\_TEX\_BUF\_DUAL}\imod{.explicit\_store} \dst{@r0:r3}, \src{r4}} \\[2pt]
\texttt{4} & \texttt{\imod{center\_clobber}} & v9--v14 & Evaluate at pixel center, clobbering existing cached barycentrics. \newline {\scriptsize\ttfamily \mnem{VAR\_TEX\_BUF\_DUAL}\imod{.center\_clobber} \dst{@r0:r3}, \src{r4}} \\[2pt]
\texttt{6} & \texttt{\imod{sample\_clobber}} & v9--v14 & Evaluate at sample position, clobbering cached barycentrics. \newline {\scriptsize\ttfamily \mnem{VAR\_TEX\_BUF\_DUAL}\imod{.sample\_clobber} \dst{@r0:r3}, \src{r4}} \\[2pt]
\texttt{7} & \texttt{\imod{retrieve}} & v9--v14 & Reuse previously stored barycentrics from the varying cache without recomputing. \newline {\scriptsize\ttfamily \mnem{VAR\_TEX\_BUF\_DUAL}\imod{.retrieve} \dst{@r0:r3}, \src{r4}} \\[2pt]
\bottomrule
\end{longtable}
\end{center}
\par\textbf{Worked Example:} \emph{Joint sample-position + cache-store encoding}
\codebox{\mnem{VAR\_TEX\_BUF\_SINGLE}\imod{.center\_store.tex2d.slot0} \dst{@r0:r3}, \src{u0}, \src{r59}, \src{r4}}
Interpolates \texttt{vUV} at pixel center and stores the barycentric weights in hardware cache for subsequent \texttt{.retrieve} varying loads.

\subsection{Encoded Sample Position None (\texttt{encoded\_sample\_position\_none\_m})}
\label{enum:encoded_sample_position_none_m}
Used by \texttt{LD\_VAR} and related instructions (4 encoding values):

\begin{center}\footnotesize
\begin{longtable}{@{}rll>{\raggedright\arraybackslash}p{0.62\linewidth}@{}}
\toprule
Enc & Value & Arch & Description \& Worked Example \\
\midrule\endhead
\texttt{0} & \texttt{\imod{center\_or\_none}} & v9--v14 & Evaluate at pixel center (default / flat varying). \newline {\scriptsize\ttfamily \mnem{LD\_VAR}\imod{.center\_or\_none} \dst{@r0:r3}, \src{r4}} \\[2pt]
\texttt{1} & \texttt{\imod{centroid}} & v9--v14 & Evaluate varying at MSAA sample centroid (avoids extrapolation outside primitive). \newline {\scriptsize\ttfamily \mnem{LD\_VAR}\imod{.centroid} \dst{@r0:r3}, \src{r4}} \\[2pt]
\texttt{2} & \texttt{\imod{sample}} & v9--v14 & Evaluate varying at current MSAA sample position (\texttt{interpolateAtSample}). \newline {\scriptsize\ttfamily \mnem{LD\_VAR}\imod{.sample} \dst{@r0:r3}, \src{r4}} \\[2pt]
\texttt{3} & \texttt{\imod{explicit}} & v9--v14 & Evaluate varying at explicit $(x, y)$ offset register (\texttt{interpolateAtOffset}). \newline {\scriptsize\ttfamily \mnem{LD\_VAR}\imod{.explicit} \dst{@r0:r3}, \src{r4}} \\[2pt]
\bottomrule
\end{longtable}
\end{center}
\par\textbf{Worked Example:} \emph{Per-sample MSAA varying load}
\codebox{\mnem{LD\_VAR\_IMM}\imod{.f32.v4.sample.slot0} \dst{@r0:r3}, \src{r59}, \imm{\#0}, \imm{\#0}}
Encodes \texttt{sample} (\texttt{2}) instead of default \texttt{center\_or\_none} (\texttt{0}) to interpolate at each MSAA sample location.

\subsection{Encoded Update Mode Special (\texttt{encoded\_update\_mode\_special\_m})}
\label{enum:encoded_update_mode_special_m}
Used by \texttt{LD\_VAR\_SPECIAL} and related instructions (3 encoding values):

\begin{center}\footnotesize
\begin{longtable}{@{}rll>{\raggedright\arraybackslash}p{0.62\linewidth}@{}}
\toprule
Enc & Value & Arch & Description \& Worked Example \\
\midrule\endhead
\texttt{0} & \texttt{\imod{store}} & v14 & Compute barycentric/varying coordinates and store in hardware varying cache. \newline {\scriptsize\ttfamily \mnem{LD\_VAR\_SPECIAL}\imod{.store} \dst{@r0:r3}, \src{r4}} \\[2pt]
\texttt{1} & \texttt{\imod{retrieve}} & v14 & Reuse previously cached barycentric coordinates (saves ALU/interpolation power). \newline {\scriptsize\ttfamily \mnem{LD\_VAR\_SPECIAL}\imod{.retrieve} \dst{@r0:r3}, \src{r4}} \\[2pt]
\texttt{3} & \texttt{\imod{clobber\_or\_none}} & v9--v13 & Overwrite varying cache (default for special varyings). \newline {\scriptsize\ttfamily \mnem{LD\_VAR\_SPECIAL}\imod{.clobber\_or\_none} \dst{@r0:r3}, \src{r4}} \\[2pt]
\bottomrule
\end{longtable}
\end{center}
\par\textbf{Worked Example:} \emph{Caching barycentrics from `LD\_VAR\_SPECIAL`}
\codebox{\mnem{LD\_VAR\_SPECIAL}\imod{.bary.store.slot0} \dst{@r0:r1}, \src{r59}}
Uses \texttt{.store} (\texttt{1}) instead of default \texttt{clobber\_or\_none} (\texttt{0}) so later varying loads can \texttt{.retrieve} the same weights.

\subsection{Fadd Sticky (\texttt{fadd\_sticky\_m})}
\label{enum:fadd_sticky_m}
Used by \texttt{F32\_TO\_F16} and related instructions (2 encoding values):

\begin{center}\footnotesize
\begin{longtable}{@{}rll>{\raggedright\arraybackslash}p{0.62\linewidth}@{}}
\toprule
Enc & Value & Arch & Description \& Worked Example \\
\midrule\endhead
\texttt{0} & \texttt{\imod{none}} & v9--v14 & Default behavior (flag disabled for \texttt{fadd\_sticky\_m}). \newline {\scriptsize\ttfamily \mnem{F32\_TO\_F16}\imod{} \dst{r0}, \src{r1}} \\[2pt]
\texttt{1} & \texttt{\imod{sticky}} & v9--v14 & OR-reduce shifted-out bits into LSB (round-to-odd / sticky bit for FP64 emulation). \newline {\scriptsize\ttfamily \mnem{F32\_TO\_F16}\imod{.sticky} \dst{r0}, \src{r1}} \\[2pt]
\bottomrule
\end{longtable}
\end{center}
\par\textbf{Worked Example:} \emph{Round-to-odd sticky bit for emulated FP64}
\codebox{\mnem{FADD.f32}\imod{.sticky} \dst{r0}, \src{r1}, \src{r2}}
OR-reduces shifted-out bits into bit 0 of the mantissa so a subsequent rounding step avoids double-rounding error.

\subsection{Fexp Adj (\texttt{fexp\_adj\_m})}
\label{enum:fexp_adj_m}
Used by \texttt{FEXP\_TABLE} and related instructions (2 encoding values):

\begin{center}\footnotesize
\begin{longtable}{@{}rll>{\raggedright\arraybackslash}p{0.62\linewidth}@{}}
\toprule
Enc & Value & Arch & Description \& Worked Example \\
\midrule\endhead
\texttt{0} & \texttt{\imod{none}} & v9--v14 & Default behavior (flag disabled for \texttt{fexp\_adj\_m}). \newline {\scriptsize\ttfamily \mnem{FEXP\_TABLE}\imod{} \dst{r0}, \src{r1}} \\[2pt]
\texttt{2} & \texttt{\imod{low}} & v9--v14 & Select low half of exponential lookup table entry. \newline {\scriptsize\ttfamily \mnem{FEXP\_TABLE}\imod{.low} \dst{r0}, \src{r1}} \\[2pt]
\bottomrule
\end{longtable}
\end{center}
\par\textbf{Worked Example:} \emph{Low-word correction in legacy `FEXP\_TABLE`}
\codebox{\mnem{FEXP\_TABLE.u4}\imod{.low} \dst{r0}, \src{r1}}
Selects the \texttt{.low} table bank to fetch the fine-precision correction term paired with the main \texttt{2\textasciicircum{}(idx/16)} entry.

\subsection{Fma Rscale Round (\texttt{fma\_rscale\_round\_m})}
\label{enum:fma_rscale_round_m}
Used by \texttt{FADD\_RSCALE} and related instructions (2 encoding values):

\begin{center}\footnotesize
\begin{longtable}{@{}rll>{\raggedright\arraybackslash}p{0.62\linewidth}@{}}
\toprule
Enc & Value & Arch & Description \& Worked Example \\
\midrule\endhead
\texttt{0} & \texttt{\imod{none}} & v9--v14 & Default behavior (flag disabled for \texttt{fma\_rscale\_round\_m}). \newline {\scriptsize\ttfamily \mnem{FADD\_RSCALE}\imod{} \dst{r0}, \src{r1}} \\[2pt]
\texttt{1} & \texttt{\imod{round\_zero}} & v9--v14 & Round toward zero (RTZ) instead of default RTE. \newline {\scriptsize\ttfamily \mnem{FADD\_RSCALE}\imod{.round\_zero} \dst{r0}, \src{r1}} \\[2pt]
\bottomrule
\end{longtable}
\end{center}
\par\textbf{Worked Example:} \emph{Truncated scaled FMA (`.round\_zero`)}
\codebox{\mnem{FMA\_RSCALE.f32}\imod{.round\_zero} \dst{r0}, \src{r1}, \src{r2}, \src{r3}, \src{r4}}
Evaluates \texttt{fma(r1, r2, r3) * 2\textasciicircum{}r4} in FP64 and rounds toward zero instead of default RNE.

\subsection{Fpow Sc Det Function (\texttt{fpow\_sc\_det\_function\_m})}
\label{enum:fpow_sc_det_function_m}
Used by \texttt{FPOW\_SC\_DET} and related instructions (2 encoding values):

\begin{center}\footnotesize
\begin{longtable}{@{}rll>{\raggedright\arraybackslash}p{0.62\linewidth}@{}}
\toprule
Enc & Value & Arch & Description \& Worked Example \\
\midrule\endhead
\texttt{0} & \texttt{\imod{pow}} & v9--v14 & Evaluate \texttt{pow(x, y)} special-case detection (negative base allowed for integer $y$). \newline {\scriptsize\ttfamily \mnem{FPOW\_SC\_DET}\imod{.pow} \dst{r0}, \src{r1}} \\[2pt]
\texttt{1} & \texttt{\imod{powr}} & v9--v14 & Evaluate OpenCL \texttt{powr(x, y)} special-case detection (negative base $\to$ NaN). \newline {\scriptsize\ttfamily \mnem{FPOW\_SC\_DET}\imod{.powr} \dst{r0}, \src{r1}} \\[2pt]
\bottomrule
\end{longtable}
\end{center}
\par\textbf{Worked Example:} \emph{`pow(0,0) = 1` vs `powr(0,0) = NaN`}
\codebox{\mnem{FPOW\_SC\_DET.f32}\imod{.powr} \dst{r0}, \src{r1}, \src{r2}}
With \texttt{r1 = 0.0f, r2 = 0.0f}: \texttt{.pow} signals \texttt{1.0f} override, whereas \texttt{.powr} signals \texttt{NaN} per OpenCL \texttt{powr} semantics.

\subsection{Iaddc Return Carry (\texttt{iaddc\_return\_carry\_m})}
\label{enum:iaddc_return_carry_m}
Used by \texttt{IADDC} and related instructions (2 encoding values):

\begin{center}\footnotesize
\begin{longtable}{@{}rll>{\raggedright\arraybackslash}p{0.62\linewidth}@{}}
\toprule
Enc & Value & Arch & Description \& Worked Example \\
\midrule\endhead
\texttt{0} & \texttt{\imod{none}} & v9--v14 & Default behavior (flag disabled for \texttt{iaddc\_return\_carry\_m}). \newline {\scriptsize\ttfamily \mnem{IADDC}\imod{} \dst{r0}, \src{r1}} \\[2pt]
\texttt{1} & \texttt{\imod{c}} & v9--v14 & Return carry-out / borrow-out bit (0 or 1) instead of sum/difference. \newline {\scriptsize\ttfamily \mnem{IADDC}\imod{.c} \dst{r0}, \src{r1}} \\[2pt]
\bottomrule
\end{longtable}
\end{center}
\par\textbf{Worked Example:} \emph{Extracting 1-bit carry-out for 64-bit add}
\codebox{\mnem{IADDC.u32}\imod{.c} \dst{r0}, \src{r1}, \src{r2}, \src{k0}}
With \texttt{r1 = 0xFFFFFFFF, r2 = 2}: \texttt{none} returns sum \texttt{1}, whereas \texttt{.c} returns carry-out \texttt{1} to feed the high-word \texttt{IADDC}.

\subsection{Jump Divergence Mode (\texttt{jump\_divergence\_mode\_m})}
\label{enum:jump_divergence_mode_m}
Used by \texttt{JUMP\_EX} and related instructions (2 encoding values):

\begin{center}\footnotesize
\begin{longtable}{@{}rll>{\raggedright\arraybackslash}p{0.62\linewidth}@{}}
\toprule
Enc & Value & Arch & Description \& Worked Example \\
\midrule\endhead
\texttt{0} & \texttt{\imod{none}} & v9--v14 & Default behavior (flag disabled for \texttt{jump\_divergence\_mode\_m}). \newline {\scriptsize\ttfamily \mnem{JUMP\_EX}\imod{} \dst{r0}, \src{r1}} \\[2pt]
\texttt{1} & \texttt{\imod{diverge}} & v9--v14 & Mark indirect jump as potentially divergent across warp lanes. \newline {\scriptsize\ttfamily \mnem{JUMP\_EX}\imod{.diverge} \dst{r0}, \src{r1}} \\[2pt]
\bottomrule
\end{longtable}
\end{center}
\par\textbf{Worked Example:} \emph{Divergent indirect function call}
\codebox{\mnem{JUMP\_EX}\imod{.call.diverge} \src{r0}, \src{r1}, \src{r2}}
Marks \texttt{.diverge} when lanes in the warp may jump to different target addresses via a function pointer.

\subsection{Jump Stack Mode (\texttt{jump\_stack\_mode\_m})}
\label{enum:jump_stack_mode_m}
Used by \texttt{JUMP\_EX} and related instructions (4 encoding values):
\par\textbf{CSF ABI Interaction:} Requires \hyperref[sec:csf-spd]{\texttt{SPD\_k::shader\_contains\_jump\_ex=1}} and \hyperref[sec:csf-tsd-fau]{\texttt{TSD\_k}} call-stack backing.
\begin{center}\footnotesize
\begin{longtable}{@{}rll>{\raggedright\arraybackslash}p{0.62\linewidth}@{}}
\toprule
Enc & Value & Arch & Description \& Worked Example \\
\midrule\endhead
\texttt{0} & \texttt{\imod{return}} & v9--v14 & Pop return address from hardware call stack and branch. \newline {\scriptsize\ttfamily \mnem{JUMP\_EX}\imod{.return} \src{r0}} \\[2pt]
\texttt{1} & \texttt{\imod{call}} & v9--v14 & Push next PC onto hardware call stack and jump to target. \newline {\scriptsize\ttfamily \mnem{JUMP\_EX}\imod{.call} \src{r0}, \imm{\#+16}} \\[2pt]
\texttt{2} & \texttt{\imod{none}} & v9--v14 & Indirect jump without modifying hardware call stack. \newline {\scriptsize\ttfamily \mnem{JUMP\_EX} \src{r0}, \src{r1}} \\[2pt]
\texttt{3} & \texttt{\imod{replace}} & v9--v14 & Tail-call: replace top-of-stack frame without pushing. \newline {\scriptsize\ttfamily \mnem{JUMP\_EX}\imod{.replace} \src{r0}, \src{r1}} \\[2pt]
\bottomrule
\end{longtable}
\end{center}
\par\textbf{Worked Example:} \emph{Hardware call-stack push and tail-call replace}
\codebox{\mnem{JUMP\_EX}\imod{.replace} \src{k1n}, \src{r1}, \src{r2}}
\texttt{.call} pushes \texttt{PC+8}, \texttt{.return} pops it, and \texttt{.replace} performs a tail call without growing the hardware stack.

\subsection{Lane All (\texttt{lane\_all\_m})}
\label{enum:lane_all_m}
Used by \texttt{LD\_CHECKED} and related instructions (8 encoding values):

\begin{center}\footnotesize
\begin{longtable}{@{}rll>{\raggedright\arraybackslash}p{0.62\linewidth}@{}}
\toprule
Enc & Value & Arch & Description \& Worked Example \\
\midrule\endhead
\texttt{0} & \texttt{\imod{b0}} & v9--v14 & Select byte 0 \texttt{[7:0]} (zero/sign-extended to lane width). \newline {\scriptsize\ttfamily \mnem{LD\_CHECKED} \dst{r0}, \src{r1}\imod{.b0}} \\[2pt]
\texttt{1} & \texttt{\imod{b1}} & v9--v14 & Select byte 1 \texttt{[15:8]} (zero/sign-extended to lane width). \newline {\scriptsize\ttfamily \mnem{LD\_CHECKED} \dst{r0}, \src{r1}\imod{.b1}} \\[2pt]
\texttt{2} & \texttt{\imod{b2}} & v9--v14 & Select byte 2 \texttt{[23:16]} (zero/sign-extended to lane width). \newline {\scriptsize\ttfamily \mnem{LD\_CHECKED} \dst{r0}, \src{r1}\imod{.b2}} \\[2pt]
\texttt{3} & \texttt{\imod{b3}} & v9--v14 & Select byte 3 \texttt{[31:24]} (zero/sign-extended to lane width). \newline {\scriptsize\ttfamily \mnem{LD\_CHECKED} \dst{r0}, \src{r1}\imod{.b3}} \\[2pt]
\texttt{4} & \texttt{\imod{h0}} & v9--v14 & Select low 16-bit halfword \texttt{[15:0]} (widen to 32-bit if FP/int widen). \newline {\scriptsize\ttfamily \mnem{LD\_CHECKED} \dst{r0}, \src{r1}\imod{.h0}} \\[2pt]
\texttt{5} & \texttt{\imod{h1}} & v9--v14 & Select high 16-bit halfword \texttt{[31:16]} (widen to 32-bit if FP/int widen). \newline {\scriptsize\ttfamily \mnem{LD\_CHECKED} \dst{r0}, \src{r1}\imod{.h1}} \\[2pt]
\texttt{6} & \texttt{\imod{w0}} & v9--v14 & Select 32-bit word 0 \texttt{[31:0]} (extend to 64-bit). \newline {\scriptsize\ttfamily \mnem{LD\_CHECKED} \dst{r0}, \src{r1}\imod{.w0}} \\[2pt]
\texttt{7} & \texttt{\imod{d0}} & v9--v14 & Select full 64-bit doubleword \texttt{[63:0]}. \newline {\scriptsize\ttfamily \mnem{LD\_CHECKED} \dst{r0}, \src{r1}\imod{.d0}} \\[2pt]
\bottomrule
\end{longtable}
\end{center}
\par\textbf{Worked Example:} \emph{Selecting sub-word lane in checked/narrow load}
\codebox{\mnem{LOAD.i8}\imod{.b2.slot0} \dst{@r0}, \src{[r2:r3]}, \imm{\#0}}
Loads 1 byte from memory and inserts it into byte lane 2 (\texttt{r0[23:16]}), leaving other bytes of \texttt{r0} intact.

\subsection{Ld Gclk Source (\texttt{ld\_gclk\_source\_m})}
\label{enum:ld_gclk_source_m}
Used by \texttt{LD\_GCLK} and related instructions (2 encoding values):

\begin{center}\footnotesize
\begin{longtable}{@{}rll>{\raggedright\arraybackslash}p{0.62\linewidth}@{}}
\toprule
Enc & Value & Arch & Description \& Worked Example \\
\midrule\endhead
\texttt{6} & \texttt{\imod{cycle\_counter}} & v9--v14 & Read 64-bit GPU shader-core cycle counter. \newline {\scriptsize\ttfamily \mnem{LD\_GCLK}\imod{.cycle\_counter} \dst{r0}, \src{r1}} \\[2pt]
\texttt{7} & \texttt{\imod{system\_timestamp}} & v9--v14 & Read 64-bit global system timestamp counter. \newline {\scriptsize\ttfamily \mnem{LD\_GCLK}\imod{.system\_timestamp} \dst{r0}, \src{r1}} \\[2pt]
\bottomrule
\end{longtable}
\end{center}
\par\textbf{Worked Example:} \emph{Reading GPU cycle counter vs system clock}
\codebox{\mnem{LD\_GCLK}\imod{.system\_timestamp.slot0} \dst{@r0:r1}}
\texttt{.cycle\_counter} reads per-core shader cycles; \texttt{.system\_timestamp} reads the cross-core monotonic wall clock.

\subsection{Ld Var Buf Encoded Update Mode None (\texttt{ld\_var\_buf\_encoded\_update\_mode\_none\_m})}
\label{enum:ld_var_buf_encoded_update_mode_none_m}
Used by \texttt{LD\_VAR\_BUF} and related instructions (3 encoding values):

\begin{center}\footnotesize
\begin{longtable}{@{}rll>{\raggedright\arraybackslash}p{0.62\linewidth}@{}}
\toprule
Enc & Value & Arch & Description \& Worked Example \\
\midrule\endhead
\texttt{0} & \texttt{\imod{store}} & v9--v13 & Compute barycentric/varying coordinates and store in hardware varying cache. \newline {\scriptsize\ttfamily \mnem{LD\_VAR\_BUF}\imod{.store} \dst{@r0:r3}, \src{r4}} \\[2pt]
\texttt{1} & \texttt{\imod{retrieve\_or\_none}} & v9--v13 & Reuse cached barycentrics if available (default). \newline {\scriptsize\ttfamily \mnem{LD\_VAR\_BUF}\imod{.retrieve\_or\_none} \dst{@r0:r3}, \src{r4}} \\[2pt]
\texttt{3} & \texttt{\imod{clobber}} & v9--v13 & Compute fresh coordinates and overwrite hardware varying cache. \newline {\scriptsize\ttfamily \mnem{LD\_VAR\_BUF}\imod{.clobber} \dst{@r0:r3}, \src{r4}} \\[2pt]
\bottomrule
\end{longtable}
\end{center}
\par\textbf{Worked Example:} \emph{Default `.retrieve\_or\_none` in varying buffer}
\codebox{\mnem{LD\_VAR\_BUF\_IMM.f32}\imod{.v4.center.store.slot0} \dst{@r0:r3}, \src{r59}, \imm{\#0}}
First varying uses \texttt{.store} (\texttt{0}); subsequent varyings default to \texttt{retrieve\_or\_none} (\texttt{1}) to reuse cached barycentrics.

\subsection{Ld Var Buf Flat Format (\texttt{ld\_var\_buf\_flat\_format\_m})}
\label{enum:ld_var_buf_flat_format_m}
Used by \texttt{LD\_VAR\_BUF\_FLAT} and related instructions (2 encoding values):

\begin{center}\footnotesize
\begin{longtable}{@{}rll>{\raggedright\arraybackslash}p{0.62\linewidth}@{}}
\toprule
Enc & Value & Arch & Description \& Worked Example \\
\midrule\endhead
\texttt{2} & \texttt{\imod{flat32}} & v14 & Flat (uninterpolated) 32-bit varying read from provoking vertex. \newline {\scriptsize\ttfamily \mnem{LD\_VAR\_BUF\_FLAT}\imod{.flat32} \dst{r0}, \src{r1}} \\[2pt]
\texttt{3} & \texttt{\imod{flat16}} & v14 & Flat (uninterpolated) 16-bit varying read from provoking vertex. \newline {\scriptsize\ttfamily \mnem{LD\_VAR\_BUF\_FLAT}\imod{.flat16} \dst{r0}, \src{r1}} \\[2pt]
\bottomrule
\end{longtable}
\end{center}
\par\textbf{Worked Example:} \emph{Flat 16-bit vs 32-bit varying buffer read (v14)}
\codebox{\mnem{LD\_VAR\_BUF\_FLAT\_IMM}\imod{.flat16.slot0} \dst{@r0}, \imm{\#0x04}}
Reads a 16-bit flat varying (\texttt{flat16}) instead of a 32-bit word (\texttt{flat32}) from varying buffer offset \texttt{0x04}.

\subsection{Ld Var Buf Source Format (\texttt{ld\_var\_buf\_source\_format\_m})}
\label{enum:ld_var_buf_source_format_m}
Used by \texttt{LD\_VAR\_BUF} and related instructions (4 encoding values):

\begin{center}\footnotesize
\begin{longtable}{@{}rll>{\raggedright\arraybackslash}p{0.62\linewidth}@{}}
\toprule
Enc & Value & Arch & Description \& Worked Example \\
\midrule\endhead
\texttt{0} & \texttt{\imod{src\_flat32}} & v9--v13 & Source buffer holds flat 32-bit varying values. \newline {\scriptsize\ttfamily \mnem{LD\_VAR\_BUF}\imod{.src\_flat32} \dst{r0}, \src{r1}} \\[2pt]
\texttt{1} & \texttt{\imod{src\_flat16}} & v9--v13 & Source buffer holds flat 16-bit varying values. \newline {\scriptsize\ttfamily \mnem{LD\_VAR\_BUF}\imod{.src\_flat16} \dst{r0}, \src{r1}} \\[2pt]
\texttt{2} & \texttt{\imod{src\_f32}} & v9--v13 & Source buffer holds FP32 varying values. \newline {\scriptsize\ttfamily \mnem{LD\_VAR\_BUF}\imod{.src\_f32} \dst{r0}, \src{r1}} \\[2pt]
\texttt{3} & \texttt{\imod{src\_f16}} & v9--v13 & Source buffer holds FP16 varying values. \newline {\scriptsize\ttfamily \mnem{LD\_VAR\_BUF}\imod{.src\_f16} \dst{r0}, \src{r1}} \\[2pt]
\bottomrule
\end{longtable}
\end{center}
\par\textbf{Worked Example:} \emph{FP16 varying buffer widened to FP32 registers}
\codebox{\mnem{LD\_VAR\_BUF.f32}\imod{.src\_f16.v2.center.slot0} \dst{@r0:r1}, \src{r59}, \src{r4}}
Declares that the varying buffer stores \texttt{f16} (\texttt{src\_f16}) while delivering interpolated \texttt{f32} results in \texttt{@r0:r1}.

\subsection{Ldexp Round (\texttt{ldexp\_round\_m})}
\label{enum:ldexp_round_m}
Used by \texttt{LDEXP} and related instructions (5 encoding values):

\begin{center}\footnotesize
\begin{longtable}{@{}rll>{\raggedright\arraybackslash}p{0.62\linewidth}@{}}
\toprule
Enc & Value & Arch & Description \& Worked Example \\
\midrule\endhead
\texttt{0} & \texttt{\imod{none}} & v9--v13 & Scale by $2^e$ with round-to-nearest ties-to-even (RTE). \newline {\scriptsize\ttfamily \mnem{LDEXP.f32} \dst{r0}, \src{r1}, \src{r2}} \\[2pt]
\texttt{1} & \texttt{\imod{round\_up}} & v9--v13 & Scale by $2^e$ rounding toward $+\infty$. \newline {\scriptsize\ttfamily \mnem{LDEXP.f32}\imod{.round\_up} \dst{r0}, \src{r1}, \src{r2}} \\[2pt]
\texttt{2} & \texttt{\imod{round\_down}} & v9--v13 & Scale by $2^e$ rounding toward $-\infty$. \newline {\scriptsize\ttfamily \mnem{LDEXP.f32}\imod{.round\_down} \dst{r0}, \src{r1}, \src{r2}} \\[2pt]
\texttt{3} & \texttt{\imod{round\_zero}} & v9--v13 & Scale by $2^e$ rounding toward zero. \newline {\scriptsize\ttfamily \mnem{LDEXP.f32}\imod{.round\_zero} \dst{r0}, \src{r1}, \src{r2}} \\[2pt]
\texttt{4} & \texttt{\imod{round\_na}} & v9--v13 & Scale by $2^e$ rounding ties away from zero. \newline {\scriptsize\ttfamily \mnem{LDEXP.f32}\imod{.round\_na} \dst{r0}, \src{r1}, \src{r2}} \\[2pt]
\bottomrule
\end{longtable}
\end{center}
\par\textbf{Worked Example:} \emph{Rounding under subnormal underflow in `LDEXP`}
\codebox{\mnem{LDEXP.f32}\imod{.round\_down} \dst{r0}, \src{r1}, \src{r2}}
With \texttt{r1 = -1.5f, r2 = -150}: rounds underflowing negative subnormal toward \texttt{-Inf} (\texttt{0x80000001}) instead of flushing to \texttt{-0.0}.

\subsection{Load Lane 16 (\texttt{load\_lane\_16\_m})}
\label{enum:load_lane_16_m}
Used by \texttt{LD\_CHECKED} and related instructions (4 encoding values):

\begin{center}\footnotesize
\begin{longtable}{@{}rll>{\raggedright\arraybackslash}p{0.62\linewidth}@{}}
\toprule
Enc & Value & Arch & Description \& Worked Example \\
\midrule\endhead
\texttt{0} & \texttt{\imod{h0}} & v9--v14 & Select low 16-bit halfword \texttt{[15:0]} (widen to 32-bit if FP/int widen). \newline {\scriptsize\ttfamily \mnem{LD\_CHECKED} \dst{r0}, \src{r1}\imod{.h0}} \\[2pt]
\texttt{1} & \texttt{\imod{h1}} & v9--v14 & Select high 16-bit halfword \texttt{[31:16]} (widen to 32-bit if FP/int widen). \newline {\scriptsize\ttfamily \mnem{LD\_CHECKED} \dst{r0}, \src{r1}\imod{.h1}} \\[2pt]
\texttt{2} & \texttt{\imod{w0}} & v9--v14 & Select 32-bit word 0 \texttt{[31:0]} (extend to 64-bit). \newline {\scriptsize\ttfamily \mnem{LD\_CHECKED} \dst{r0}, \src{r1}\imod{.w0}} \\[2pt]
\texttt{3} & \texttt{\imod{d0}} & v9--v14 & Select full 64-bit doubleword \texttt{[63:0]}. \newline {\scriptsize\ttfamily \mnem{LD\_CHECKED} \dst{r0}, \src{r1}\imod{.d0}} \\[2pt]
\bottomrule
\end{longtable}
\end{center}
\par\textbf{Worked Example:} \emph{Inserting a 16-bit load into `r0.h1`}
\codebox{\mnem{LOAD.i16}\imod{.h1.slot0} \dst{@r0}, \src{[r2:r3]}, \imm{\#0x2}}
Loads 16 bits from \texttt{[r2:r3] + 2} directly into \texttt{r0[31:16]} without overwriting \texttt{r0[15:0]}.

\subsection{Load Lane 32 (\texttt{load\_lane\_32\_m})}
\label{enum:load_lane_32_m}
Used by \texttt{LD\_PKA} and related instructions (2 encoding values):

\begin{center}\footnotesize
\begin{longtable}{@{}rll>{\raggedright\arraybackslash}p{0.62\linewidth}@{}}
\toprule
Enc & Value & Arch & Description \& Worked Example \\
\midrule\endhead
\texttt{0} & \texttt{\imod{none}} & v9--v14 & Use full 32-bit operand without narrowing/widening. \newline {\scriptsize\ttfamily \mnem{LD\_PKA} \dst{r0}, \src{r1}\imod{}} \\[2pt]
\texttt{1} & \texttt{\imod{d0}} & v9--v14 & Select full 64-bit doubleword \texttt{[63:0]}. \newline {\scriptsize\ttfamily \mnem{LD\_PKA} \dst{r0}, \src{r1}\imod{.d0}} \\[2pt]
\bottomrule
\end{longtable}
\end{center}
\par\textbf{Worked Example:} \emph{32-bit vs 64-bit (`d0`) load lane mode}
\codebox{\mnem{LOAD.i32}\imod{.slot0} \dst{@r0}, \src{[r2:r3]}, \imm{\#0}}
Default \texttt{none} writes a 32-bit word into \texttt{@r0}; \texttt{d0} selects full 64-bit doubleword lane addressing.

\subsection{Ls Multi Sr Count (\texttt{ls\_multi\_sr\_count\_m})}
\label{enum:ls_multi_sr_count_m}
Used by \texttt{LD\_MULTI} and related instructions (64 encoding values):

\begin{center}\footnotesize
\begin{longtable}{@{}rll>{\raggedright\arraybackslash}p{0.62\linewidth}@{}}
\toprule
Enc & Value & Arch & Description \& Worked Example \\
\midrule\endhead
\texttt{0} & \texttt{\imod{i32}} & v14 & Transfer 1 consecutive 32-bit staging register(s) (32 bits total). \newline {\scriptsize\ttfamily \mnem{LD\_MULTI}\imod{.i32} \dst{@r0}, \src{r60}, \imm{\#0}} \\[2pt]
\texttt{1} & \texttt{\imod{i64}} & v14 & Transfer 2 consecutive 32-bit staging register(s) (64 bits total). \newline {\scriptsize\ttfamily \mnem{LD\_MULTI}\imod{.i64} \dst{@r0:r1}, \src{r60}, \imm{\#0}} \\[2pt]
\texttt{2} & \texttt{\imod{i96}} & v14 & Transfer 3 consecutive 32-bit staging register(s) (96 bits total). \newline {\scriptsize\ttfamily \mnem{LD\_MULTI}\imod{.i96} \dst{@r0:r2}, \src{r60}, \imm{\#0}} \\[2pt]
\texttt{3} & \texttt{\imod{i128}} & v14 & Transfer 4 consecutive 32-bit staging register(s) (128 bits total). \newline {\scriptsize\ttfamily \mnem{LD\_MULTI}\imod{.i128} \dst{@r0:r3}, \src{r60}, \imm{\#0}} \\[2pt]
\texttt{4} & \texttt{\imod{i160}} & v14 & Transfer 5 consecutive 32-bit staging register(s) (160 bits total). \newline {\scriptsize\ttfamily \mnem{LD\_MULTI}\imod{.i160} \dst{@r0:r4}, \src{r60}, \imm{\#0}} \\[2pt]
\texttt{5} & \texttt{\imod{i192}} & v14 & Transfer 6 consecutive 32-bit staging register(s) (192 bits total). \newline {\scriptsize\ttfamily \mnem{LD\_MULTI}\imod{.i192} \dst{@r0:r5}, \src{r60}, \imm{\#0}} \\[2pt]
\texttt{6} & \texttt{\imod{i224}} & v14 & Transfer 7 consecutive 32-bit staging register(s) (224 bits total). \newline {\scriptsize\ttfamily \mnem{LD\_MULTI}\imod{.i224} \dst{@r0:r6}, \src{r60}, \imm{\#0}} \\[2pt]
\texttt{7} & \texttt{\imod{i256}} & v14 & Transfer 8 consecutive 32-bit staging register(s) (256 bits total). \newline {\scriptsize\ttfamily \mnem{LD\_MULTI}\imod{.i256} \dst{@r0:r7}, \src{r60}, \imm{\#0}} \\[2pt]
\texttt{8} & \texttt{\imod{i288}} & v14 & Transfer 9 consecutive 32-bit staging register(s) (288 bits total). \newline {\scriptsize\ttfamily \mnem{LD\_MULTI}\imod{.i288} \dst{@r0:r8}, \src{r60}, \imm{\#0}} \\[2pt]
\texttt{9} & \texttt{\imod{i320}} & v14 & Transfer 10 consecutive 32-bit staging register(s) (320 bits total). \newline {\scriptsize\ttfamily \mnem{LD\_MULTI}\imod{.i320} \dst{@r0:r9}, \src{r60}, \imm{\#0}} \\[2pt]
\texttt{10} & \texttt{\imod{i352}} & v14 & Transfer 11 consecutive 32-bit staging register(s) (352 bits total). \newline {\scriptsize\ttfamily \mnem{LD\_MULTI}\imod{.i352} \dst{@r0:r10}, \src{r60}, \imm{\#0}} \\[2pt]
\texttt{11} & \texttt{\imod{i384}} & v14 & Transfer 12 consecutive 32-bit staging register(s) (384 bits total). \newline {\scriptsize\ttfamily \mnem{LD\_MULTI}\imod{.i384} \dst{@r0:r11}, \src{r60}, \imm{\#0}} \\[2pt]
\texttt{12} & \texttt{\imod{i416}} & v14 & Transfer 13 consecutive 32-bit staging register(s) (416 bits total). \newline {\scriptsize\ttfamily \mnem{LD\_MULTI}\imod{.i416} \dst{@r0:r12}, \src{r60}, \imm{\#0}} \\[2pt]
\texttt{13} & \texttt{\imod{i448}} & v14 & Transfer 14 consecutive 32-bit staging register(s) (448 bits total). \newline {\scriptsize\ttfamily \mnem{LD\_MULTI}\imod{.i448} \dst{@r0:r13}, \src{r60}, \imm{\#0}} \\[2pt]
\texttt{14} & \texttt{\imod{i480}} & v14 & Transfer 15 consecutive 32-bit staging register(s) (480 bits total). \newline {\scriptsize\ttfamily \mnem{LD\_MULTI}\imod{.i480} \dst{@r0:r14}, \src{r60}, \imm{\#0}} \\[2pt]
\texttt{15} & \texttt{\imod{i512}} & v14 & Transfer 16 consecutive 32-bit staging register(s) (512 bits total). \newline {\scriptsize\ttfamily \mnem{LD\_MULTI}\imod{.i512} \dst{@r0:r15}, \src{r60}, \imm{\#0}} \\[2pt]
\texttt{16} & \texttt{\imod{i544}} & v14 & Transfer 17 consecutive 32-bit staging register(s) (544 bits total). \newline {\scriptsize\ttfamily \mnem{LD\_MULTI}\imod{.i544} \dst{@r0:r16}, \src{r60}, \imm{\#0}} \\[2pt]
\texttt{17} & \texttt{\imod{i576}} & v14 & Transfer 18 consecutive 32-bit staging register(s) (576 bits total). \newline {\scriptsize\ttfamily \mnem{LD\_MULTI}\imod{.i576} \dst{@r0:r17}, \src{r60}, \imm{\#0}} \\[2pt]
\texttt{18} & \texttt{\imod{i608}} & v14 & Transfer 19 consecutive 32-bit staging register(s) (608 bits total). \newline {\scriptsize\ttfamily \mnem{LD\_MULTI}\imod{.i608} \dst{@r0:r18}, \src{r60}, \imm{\#0}} \\[2pt]
\texttt{19} & \texttt{\imod{i640}} & v14 & Transfer 20 consecutive 32-bit staging register(s) (640 bits total). \newline {\scriptsize\ttfamily \mnem{LD\_MULTI}\imod{.i640} \dst{@r0:r19}, \src{r60}, \imm{\#0}} \\[2pt]
\texttt{20} & \texttt{\imod{i672}} & v14 & Transfer 21 consecutive 32-bit staging register(s) (672 bits total). \newline {\scriptsize\ttfamily \mnem{LD\_MULTI}\imod{.i672} \dst{@r0:r20}, \src{r60}, \imm{\#0}} \\[2pt]
\texttt{21} & \texttt{\imod{i704}} & v14 & Transfer 22 consecutive 32-bit staging register(s) (704 bits total). \newline {\scriptsize\ttfamily \mnem{LD\_MULTI}\imod{.i704} \dst{@r0:r21}, \src{r60}, \imm{\#0}} \\[2pt]
\texttt{22} & \texttt{\imod{i736}} & v14 & Transfer 23 consecutive 32-bit staging register(s) (736 bits total). \newline {\scriptsize\ttfamily \mnem{LD\_MULTI}\imod{.i736} \dst{@r0:r22}, \src{r60}, \imm{\#0}} \\[2pt]
\texttt{23} & \texttt{\imod{i768}} & v14 & Transfer 24 consecutive 32-bit staging register(s) (768 bits total). \newline {\scriptsize\ttfamily \mnem{LD\_MULTI}\imod{.i768} \dst{@r0:r23}, \src{r60}, \imm{\#0}} \\[2pt]
\texttt{24} & \texttt{\imod{i800}} & v14 & Transfer 25 consecutive 32-bit staging register(s) (800 bits total). \newline {\scriptsize\ttfamily \mnem{LD\_MULTI}\imod{.i800} \dst{@r0:r24}, \src{r60}, \imm{\#0}} \\[2pt]
\texttt{25} & \texttt{\imod{i832}} & v14 & Transfer 26 consecutive 32-bit staging register(s) (832 bits total). \newline {\scriptsize\ttfamily \mnem{LD\_MULTI}\imod{.i832} \dst{@r0:r25}, \src{r60}, \imm{\#0}} \\[2pt]
\texttt{26} & \texttt{\imod{i864}} & v14 & Transfer 27 consecutive 32-bit staging register(s) (864 bits total). \newline {\scriptsize\ttfamily \mnem{LD\_MULTI}\imod{.i864} \dst{@r0:r26}, \src{r60}, \imm{\#0}} \\[2pt]
\texttt{27} & \texttt{\imod{i896}} & v14 & Transfer 28 consecutive 32-bit staging register(s) (896 bits total). \newline {\scriptsize\ttfamily \mnem{LD\_MULTI}\imod{.i896} \dst{@r0:r27}, \src{r60}, \imm{\#0}} \\[2pt]
\texttt{28} & \texttt{\imod{i928}} & v14 & Transfer 29 consecutive 32-bit staging register(s) (928 bits total). \newline {\scriptsize\ttfamily \mnem{LD\_MULTI}\imod{.i928} \dst{@r0:r28}, \src{r60}, \imm{\#0}} \\[2pt]
\texttt{29} & \texttt{\imod{i960}} & v14 & Transfer 30 consecutive 32-bit staging register(s) (960 bits total). \newline {\scriptsize\ttfamily \mnem{LD\_MULTI}\imod{.i960} \dst{@r0:r29}, \src{r60}, \imm{\#0}} \\[2pt]
\texttt{30} & \texttt{\imod{i992}} & v14 & Transfer 31 consecutive 32-bit staging register(s) (992 bits total). \newline {\scriptsize\ttfamily \mnem{LD\_MULTI}\imod{.i992} \dst{@r0:r30}, \src{r60}, \imm{\#0}} \\[2pt]
\texttt{31} & \texttt{\imod{i1024}} & v14 & Transfer 32 consecutive 32-bit staging register(s) (1024 bits total). \newline {\scriptsize\ttfamily \mnem{LD\_MULTI}\imod{.i1024} \dst{@r0:r31}, \src{r60}, \imm{\#0}} \\[2pt]
\texttt{32} & \texttt{\imod{i1056}} & v14 & Transfer 33 consecutive 32-bit staging register(s) (1056 bits total). \newline {\scriptsize\ttfamily \mnem{LD\_MULTI}\imod{.i1056} \dst{@r0:r32}, \src{r60}, \imm{\#0}} \\[2pt]
\texttt{33} & \texttt{\imod{i1088}} & v14 & Transfer 34 consecutive 32-bit staging register(s) (1088 bits total). \newline {\scriptsize\ttfamily \mnem{LD\_MULTI}\imod{.i1088} \dst{@r0:r33}, \src{r60}, \imm{\#0}} \\[2pt]
\texttt{34} & \texttt{\imod{i1120}} & v14 & Transfer 35 consecutive 32-bit staging register(s) (1120 bits total). \newline {\scriptsize\ttfamily \mnem{LD\_MULTI}\imod{.i1120} \dst{@r0:r34}, \src{r60}, \imm{\#0}} \\[2pt]
\texttt{35} & \texttt{\imod{i1152}} & v14 & Transfer 36 consecutive 32-bit staging register(s) (1152 bits total). \newline {\scriptsize\ttfamily \mnem{LD\_MULTI}\imod{.i1152} \dst{@r0:r35}, \src{r60}, \imm{\#0}} \\[2pt]
\texttt{36} & \texttt{\imod{i1184}} & v14 & Transfer 37 consecutive 32-bit staging register(s) (1184 bits total). \newline {\scriptsize\ttfamily \mnem{LD\_MULTI}\imod{.i1184} \dst{@r0:r36}, \src{r60}, \imm{\#0}} \\[2pt]
\texttt{37} & \texttt{\imod{i1216}} & v14 & Transfer 38 consecutive 32-bit staging register(s) (1216 bits total). \newline {\scriptsize\ttfamily \mnem{LD\_MULTI}\imod{.i1216} \dst{@r0:r37}, \src{r60}, \imm{\#0}} \\[2pt]
\texttt{38} & \texttt{\imod{i1248}} & v14 & Transfer 39 consecutive 32-bit staging register(s) (1248 bits total). \newline {\scriptsize\ttfamily \mnem{LD\_MULTI}\imod{.i1248} \dst{@r0:r38}, \src{r60}, \imm{\#0}} \\[2pt]
\texttt{39} & \texttt{\imod{i1280}} & v14 & Transfer 40 consecutive 32-bit staging register(s) (1280 bits total). \newline {\scriptsize\ttfamily \mnem{LD\_MULTI}\imod{.i1280} \dst{@r0:r39}, \src{r60}, \imm{\#0}} \\[2pt]
\texttt{40} & \texttt{\imod{i1312}} & v14 & Transfer 41 consecutive 32-bit staging register(s) (1312 bits total). \newline {\scriptsize\ttfamily \mnem{LD\_MULTI}\imod{.i1312} \dst{@r0:r40}, \src{r60}, \imm{\#0}} \\[2pt]
\texttt{41} & \texttt{\imod{i1344}} & v14 & Transfer 42 consecutive 32-bit staging register(s) (1344 bits total). \newline {\scriptsize\ttfamily \mnem{LD\_MULTI}\imod{.i1344} \dst{@r0:r41}, \src{r60}, \imm{\#0}} \\[2pt]
\texttt{42} & \texttt{\imod{i1376}} & v14 & Transfer 43 consecutive 32-bit staging register(s) (1376 bits total). \newline {\scriptsize\ttfamily \mnem{LD\_MULTI}\imod{.i1376} \dst{@r0:r42}, \src{r60}, \imm{\#0}} \\[2pt]
\texttt{43} & \texttt{\imod{i1408}} & v14 & Transfer 44 consecutive 32-bit staging register(s) (1408 bits total). \newline {\scriptsize\ttfamily \mnem{LD\_MULTI}\imod{.i1408} \dst{@r0:r43}, \src{r60}, \imm{\#0}} \\[2pt]
\texttt{44} & \texttt{\imod{i1440}} & v14 & Transfer 45 consecutive 32-bit staging register(s) (1440 bits total). \newline {\scriptsize\ttfamily \mnem{LD\_MULTI}\imod{.i1440} \dst{@r0:r44}, \src{r60}, \imm{\#0}} \\[2pt]
\texttt{45} & \texttt{\imod{i1472}} & v14 & Transfer 46 consecutive 32-bit staging register(s) (1472 bits total). \newline {\scriptsize\ttfamily \mnem{LD\_MULTI}\imod{.i1472} \dst{@r0:r45}, \src{r60}, \imm{\#0}} \\[2pt]
\texttt{46} & \texttt{\imod{i1504}} & v14 & Transfer 47 consecutive 32-bit staging register(s) (1504 bits total). \newline {\scriptsize\ttfamily \mnem{LD\_MULTI}\imod{.i1504} \dst{@r0:r46}, \src{r60}, \imm{\#0}} \\[2pt]
\texttt{47} & \texttt{\imod{i1536}} & v14 & Transfer 48 consecutive 32-bit staging register(s) (1536 bits total). \newline {\scriptsize\ttfamily \mnem{LD\_MULTI}\imod{.i1536} \dst{@r0:r47}, \src{r60}, \imm{\#0}} \\[2pt]
\texttt{48} & \texttt{\imod{i1568}} & v14 & Transfer 49 consecutive 32-bit staging register(s) (1568 bits total). \newline {\scriptsize\ttfamily \mnem{LD\_MULTI}\imod{.i1568} \dst{@r0:r48}, \src{r60}, \imm{\#0}} \\[2pt]
\texttt{49} & \texttt{\imod{i1600}} & v14 & Transfer 50 consecutive 32-bit staging register(s) (1600 bits total). \newline {\scriptsize\ttfamily \mnem{LD\_MULTI}\imod{.i1600} \dst{@r0:r49}, \src{r60}, \imm{\#0}} \\[2pt]
\texttt{50} & \texttt{\imod{i1632}} & v14 & Transfer 51 consecutive 32-bit staging register(s) (1632 bits total). \newline {\scriptsize\ttfamily \mnem{LD\_MULTI}\imod{.i1632} \dst{@r0:r50}, \src{r60}, \imm{\#0}} \\[2pt]
\texttt{51} & \texttt{\imod{i1664}} & v14 & Transfer 52 consecutive 32-bit staging register(s) (1664 bits total). \newline {\scriptsize\ttfamily \mnem{LD\_MULTI}\imod{.i1664} \dst{@r0:r51}, \src{r60}, \imm{\#0}} \\[2pt]
\texttt{52} & \texttt{\imod{i1696}} & v14 & Transfer 53 consecutive 32-bit staging register(s) (1696 bits total). \newline {\scriptsize\ttfamily \mnem{LD\_MULTI}\imod{.i1696} \dst{@r0:r52}, \src{r60}, \imm{\#0}} \\[2pt]
\texttt{53} & \texttt{\imod{i1728}} & v14 & Transfer 54 consecutive 32-bit staging register(s) (1728 bits total). \newline {\scriptsize\ttfamily \mnem{LD\_MULTI}\imod{.i1728} \dst{@r0:r53}, \src{r60}, \imm{\#0}} \\[2pt]
\texttt{54} & \texttt{\imod{i1760}} & v14 & Transfer 55 consecutive 32-bit staging register(s) (1760 bits total). \newline {\scriptsize\ttfamily \mnem{LD\_MULTI}\imod{.i1760} \dst{@r0:r54}, \src{r60}, \imm{\#0}} \\[2pt]
\texttt{55} & \texttt{\imod{i1792}} & v14 & Transfer 56 consecutive 32-bit staging register(s) (1792 bits total). \newline {\scriptsize\ttfamily \mnem{LD\_MULTI}\imod{.i1792} \dst{@r0:r55}, \src{r60}, \imm{\#0}} \\[2pt]
\texttt{56} & \texttt{\imod{i1824}} & v14 & Transfer 57 consecutive 32-bit staging register(s) (1824 bits total). \newline {\scriptsize\ttfamily \mnem{LD\_MULTI}\imod{.i1824} \dst{@r0:r56}, \src{r60}, \imm{\#0}} \\[2pt]
\texttt{57} & \texttt{\imod{i1856}} & v14 & Transfer 58 consecutive 32-bit staging register(s) (1856 bits total). \newline {\scriptsize\ttfamily \mnem{LD\_MULTI}\imod{.i1856} \dst{@r0:r57}, \src{r60}, \imm{\#0}} \\[2pt]
\texttt{58} & \texttt{\imod{i1888}} & v14 & Transfer 59 consecutive 32-bit staging register(s) (1888 bits total). \newline {\scriptsize\ttfamily \mnem{LD\_MULTI}\imod{.i1888} \dst{@r0:r58}, \src{r60}, \imm{\#0}} \\[2pt]
\texttt{59} & \texttt{\imod{i1920}} & v14 & Transfer 60 consecutive 32-bit staging register(s) (1920 bits total). \newline {\scriptsize\ttfamily \mnem{LD\_MULTI}\imod{.i1920} \dst{@r0:r59}, \src{r60}, \imm{\#0}} \\[2pt]
\texttt{60} & \texttt{\imod{i1952}} & v14 & Transfer 61 consecutive 32-bit staging register(s) (1952 bits total). \newline {\scriptsize\ttfamily \mnem{LD\_MULTI}\imod{.i1952} \dst{@r0:r60}, \src{r60}, \imm{\#0}} \\[2pt]
\texttt{61} & \texttt{\imod{i1984}} & v14 & Transfer 62 consecutive 32-bit staging register(s) (1984 bits total). \newline {\scriptsize\ttfamily \mnem{LD\_MULTI}\imod{.i1984} \dst{@r0:r61}, \src{r60}, \imm{\#0}} \\[2pt]
\texttt{62} & \texttt{\imod{i2016}} & v14 & Transfer 63 consecutive 32-bit staging register(s) (2016 bits total). \newline {\scriptsize\ttfamily \mnem{LD\_MULTI}\imod{.i2016} \dst{@r0:r62}, \src{r60}, \imm{\#0}} \\[2pt]
\texttt{63} & \texttt{\imod{i2048}} & v14 & Transfer 64 consecutive 32-bit staging register(s) (2048 bits total). \newline {\scriptsize\ttfamily \mnem{LD\_MULTI}\imod{.i2048} \dst{@r0:r63}, \src{r60}, \imm{\#0}} \\[2pt]
\bottomrule
\end{longtable}
\end{center}
\par\textbf{Worked Example:} \emph{Bulk 8-register (256-bit) transfer on v14}
\codebox{\mnem{LD\_MULTI}\imod{.i256.slot0} \dst{@r0:r7}, \src{[r8:r9]}, \imm{\#0}}
Encoding \texttt{7} (\texttt{.i256}) transfers \texttt{(7 + 1) * 32 = 256} bits (8 consecutive staging registers \texttt{@r0:r7}) in one message.

\subsection{Ls Z Last Read (\texttt{ls\_z\_last\_read\_m})}
\label{enum:ls_z_last_read_m}
Used by \texttt{LD\_TILE} and related instructions (2 encoding values):
\par\textbf{CSF ABI Interaction:} Signals final read of \hyperref[sec:csf-fbd]{\texttt{RenderTarget[rt]}} on-chip 16$\times$16 Tile SRAM.
\begin{center}\footnotesize
\begin{longtable}{@{}rll>{\raggedright\arraybackslash}p{0.62\linewidth}@{}}
\toprule
Enc & Value & Arch & Description \& Worked Example \\
\midrule\endhead
\texttt{0} & \texttt{\imod{none}} & v11--v14 & Default behavior (flag disabled for \texttt{ls\_z\_last\_read\_m}). \newline {\scriptsize\ttfamily \mnem{LD\_TILE}\imod{} \dst{r0}, \src{r1}} \\[2pt]
\texttt{1} & \texttt{\imod{last\_read}} & v11--v14 & Signal final tile-buffer read so hardware may release tile memory. \newline {\scriptsize\ttfamily \mnem{LD\_TILE}\imod{.last\_read} \dst{r0}, \src{r1}} \\[2pt]
\bottomrule
\end{longtable}
\end{center}
\par\textbf{Worked Example:} \emph{Reclaiming tile storage on final `LD\_TILE`}
\codebox{\mnem{LD\_TILE}\imod{.f32.v4.last\_read.slot0} \dst{@r0:r3}, \src{r59}, \src{r60}, \src{u0}}
\texttt{.last\_read} signals that no further tilebuffer reads will occur for this pixel so hardware can release tile resources early.

\subsection{Ls Z Stencil (\texttt{ls\_z\_stencil\_m})}
\label{enum:ls_z_stencil_m}
Used by \texttt{LD\_TILE} and related instructions (2 encoding values):
\par\textbf{CSF ABI Interaction:} Targets \hyperref[sec:csf-fbd]{\texttt{FRAGMENTSR::FBD\_POINTER}} (\texttt{zs\_crc\_ext=1}) $\to$ \texttt{ZSCRCExtension} (64B) instead of \texttt{RenderTarget[rt]}.
\begin{center}\footnotesize
\begin{longtable}{@{}rll>{\raggedright\arraybackslash}p{0.62\linewidth}@{}}
\toprule
Enc & Value & Arch & Description \& Worked Example \\
\midrule\endhead
\texttt{0} & \texttt{\imod{none}} & v9--v14 & Default behavior (flag disabled for \texttt{ls\_z\_stencil\_m}). \newline {\scriptsize\ttfamily \mnem{LD\_TILE}\imod{} \dst{r0}, \src{r1}} \\[2pt]
\texttt{1} & \texttt{\imod{z\_stencil}} & v9--v14 & Access combined depth/stencil tile buffer instead of color tile. \newline {\scriptsize\ttfamily \mnem{LD\_TILE}\imod{.z\_stencil} \dst{r0}, \src{r1}} \\[2pt]
\bottomrule
\end{longtable}
\end{center}
\par\textbf{Worked Example:} \emph{Reading depth/stencil tilebuffer instead of color}
\codebox{\mnem{LD\_TILE}\imod{.f32.z\_stencil.slot0} \dst{@r0}, \src{r59}, \src{r60}, \src{u1}}
With \texttt{.z\_stencil}, targets the on-chip depth/stencil attachment rather than a color render target.

\subsection{Message Slot Index (\texttt{message\_slot\_index\_m})}
\label{enum:message_slot_index_m}
Used by \texttt{ACMPXCHG} and related instructions (4 encoding values):

\begin{center}\footnotesize
\begin{longtable}{@{}rll>{\raggedright\arraybackslash}p{0.62\linewidth}@{}}
\toprule
Enc & Value & Arch & Description \& Worked Example \\
\midrule\endhead
\texttt{0} & \texttt{\imod{slot0}} & v9--v14 & Track async message completion in scoreboard slot 0. \newline {\scriptsize\ttfamily \mnem{LOAD.i32}\imod{.slot0} \dst{@r0}, \src{r2}, \imm{\#0}} \\[2pt]
\texttt{1} & \texttt{\imod{slot1}} & v9--v14 & Track async message completion in scoreboard slot 1. \newline {\scriptsize\ttfamily \mnem{TEX\_SINGLE}\imod{.slot1} \dst{@r0:r3}, \src{@r4:r5}} \\[2pt]
\texttt{2} & \texttt{\imod{slot2}} & v9--v14 & Track async message completion in scoreboard slot 2. \newline {\scriptsize\ttfamily \mnem{LD\_VAR\_BUF.f32}\imod{.slot2} \dst{@r0:r3}, \src{r4}} \\[2pt]
\texttt{7} & \texttt{\imod{slot7}} & v9--v14 & Track async message in slot 7 (non-blocking / fire-and-forget write). \newline {\scriptsize\ttfamily \mnem{STORE.i32}\imod{.slot7} \src{@r0}, \src{r2}, \imm{\#0}} \\[2pt]
\bottomrule
\end{longtable}
\end{center}
\par\textbf{Worked Example:} \emph{Overlapping async load (`slot0`) and texture (`slot1`)}
\codebox{\mnem{LOAD.i32}\imod{.slot0} \dst{@r0}, \src{[r2:r3]}, \imm{\#0}}
Assigning independent messages to \texttt{slot0} and \texttt{slot1} lets the shader wait on \texttt{.wait0} while \texttt{slot1} remains in flight (\texttt{slot7} is non-blocking).

\subsection{Round Extended (\texttt{round\_extended\_m})}
\label{enum:round_extended_m}
Used by \texttt{LDEXP} and related instructions (5 encoding values):

\begin{center}\footnotesize
\begin{longtable}{@{}rll>{\raggedright\arraybackslash}p{0.62\linewidth}@{}}
\toprule
Enc & Value & Arch & Description \& Worked Example \\
\midrule\endhead
\texttt{0} & \texttt{\imod{none}} & v9--v14 & Round to nearest ties-to-even (RTE). \newline {\scriptsize\ttfamily \mnem{S32\_TO\_F32} \dst{r0}, \src{r1}} \\[2pt]
\texttt{1} & \texttt{\imod{round\_up}} & v9--v14 & Round toward $+\infty$ (RTP). \newline {\scriptsize\ttfamily \mnem{S32\_TO\_F32}\imod{.round\_up} \dst{r0}, \src{r1}} \\[2pt]
\texttt{2} & \texttt{\imod{round\_down}} & v9--v14 & Round toward $-\infty$ (RTN). \newline {\scriptsize\ttfamily \mnem{S32\_TO\_F32}\imod{.round\_down} \dst{r0}, \src{r1}} \\[2pt]
\texttt{3} & \texttt{\imod{round\_zero}} & v9--v14 & Round toward zero (RTZ). \newline {\scriptsize\ttfamily \mnem{S32\_TO\_F32}\imod{.round\_zero} \dst{r0}, \src{r1}} \\[2pt]
\texttt{4} & \texttt{\imod{round\_na}} & v9--v14 & Round to nearest ties-away-from-zero (RNA). \newline {\scriptsize\ttfamily \mnem{S32\_TO\_F32}\imod{.round\_na} \dst{r0}, \src{r1}} \\[2pt]
\bottomrule
\end{longtable}
\end{center}
\par\textbf{Worked Example:} \emph{Ties-away-from-zero int-to-float conversion}
\codebox{\mnem{S32\_TO\_F32}\imod{.round\_na} \dst{r0}, \src{r1}}
For \texttt{r1 = 0x01000001} (halfway between two FP32 values): \texttt{.round\_na} rounds away from zero while default RNE rounds to even.

\subsection{Sample Position None (\texttt{sample\_position\_none\_m})}
\label{enum:sample_position_none_m}
Used by \texttt{LD\_VAR} and related instructions (5 encoding values):

\begin{center}\footnotesize
\begin{longtable}{@{}rll>{\raggedright\arraybackslash}p{0.62\linewidth}@{}}
\toprule
Enc & Value & Arch & Description \& Worked Example \\
\midrule\endhead
\texttt{0} & \texttt{\imod{center}} & v9--v14 & Evaluate varying at pixel center $(0.5, 0.5)$. \newline {\scriptsize\ttfamily \mnem{LD\_VAR}\imod{.center} \dst{@r0:r3}, \src{r4}} \\[2pt]
\texttt{1} & \texttt{\imod{centroid}} & v9--v14 & Evaluate varying at MSAA sample centroid (avoids extrapolation outside primitive). \newline {\scriptsize\ttfamily \mnem{LD\_VAR}\imod{.centroid} \dst{@r0:r3}, \src{r4}} \\[2pt]
\texttt{2} & \texttt{\imod{sample}} & v9--v14 & Evaluate varying at current MSAA sample position (\texttt{interpolateAtSample}). \newline {\scriptsize\ttfamily \mnem{LD\_VAR}\imod{.sample} \dst{@r0:r3}, \src{r4}} \\[2pt]
\texttt{3} & \texttt{\imod{explicit}} & v9--v14 & Evaluate varying at explicit $(x, y)$ offset register (\texttt{interpolateAtOffset}). \newline {\scriptsize\ttfamily \mnem{LD\_VAR}\imod{.explicit} \dst{@r0:r3}, \src{r4}} \\[2pt]
\texttt{4} & \texttt{\imod{none}} & v9--v14 & No sample position override (flat or retrieve-cached). \newline {\scriptsize\ttfamily \mnem{LD\_VAR}\imod{} \dst{@r0:r3}, \src{r4}} \\[2pt]
\bottomrule
\end{longtable}
\end{center}
\par\textbf{Worked Example:} \emph{Centroid vs uninterpolated `none` sample position}
\codebox{\mnem{LD\_VAR\_BUF.f32}\imod{.v4.centroid.slot0} \dst{@r0:r3}, \src{r59}, \src{r4}}
Selects \texttt{centroid} for MSAA-safe interpolation, or \texttt{none} when interpolation position is unused.

\subsection{Segadd Preserve Null (\texttt{segadd\_preserve\_null\_m})}
\label{enum:segadd_preserve_null_m}
Used by \texttt{SEG\_ADD} and related instructions (2 encoding values):
\par\textbf{CSF ABI Interaction:} Preserves null pointers during \hyperref[sec:csf-tsd-fau]{\texttt{TSD\_k}} TLS/WLS segmented address arithmetic.
\begin{center}\footnotesize
\begin{longtable}{@{}rll>{\raggedright\arraybackslash}p{0.62\linewidth}@{}}
\toprule
Enc & Value & Arch & Description \& Worked Example \\
\midrule\endhead
\texttt{0} & \texttt{\imod{none}} & v9--v14 & Default behavior (flag disabled for \texttt{segadd\_preserve\_null\_m}). \newline {\scriptsize\ttfamily \mnem{SEG\_ADD}\imod{} \dst{r0}, \src{r1}} \\[2pt]
\texttt{1} & \texttt{\imod{preserve\_null}} & v9--v14 & If base address is 0 (NULL), return 0 instead of adding offset. \newline {\scriptsize\ttfamily \mnem{SEG\_ADD}\imod{.preserve\_null} \dst{r0}, \src{r1}} \\[2pt]
\bottomrule
\end{longtable}
\end{center}
\par\textbf{Worked Example:} \emph{Keeping NULL pointer zero on offset add}
\codebox{\mnem{SEG\_ADD.u64}\imod{.preserve\_null} \dst{[r0:r1]}, \src{[r2:r3]}, \src{r4}\imod{.w0}}
With \texttt{[r2:r3] = 0, r4 = 16}: \texttt{.preserve\_null} keeps \texttt{[r0:r1] = 0}; with \texttt{none}, computes \texttt{[r0:r1] = 16}.

\subsection{Sign Extend Or None Compressed Encoding (\texttt{sign\_extend\_or\_none\_compressed\_encoding\_m})}
\label{enum:sign_extend_or_none_compressed_encoding_m}
Used by \texttt{LD\_PKA} and related instructions (2 encoding values):

\begin{center}\footnotesize
\begin{longtable}{@{}rll>{\raggedright\arraybackslash}p{0.62\linewidth}@{}}
\toprule
Enc & Value & Arch & Description \& Worked Example \\
\midrule\endhead
\texttt{0} & \texttt{\imod{sext}} & v9--v14 & Sign-extend narrow loaded value (8/16-bit) to 32 bits. \newline {\scriptsize\ttfamily \mnem{LD\_PKA}\imod{.sext} \dst{r0}, \src{r1}} \\[2pt]
\texttt{1} & \texttt{\imod{zext\_or\_none}} & v9--v14 & Zero-extend narrow load (or no extension for 32-bit+ loads). \newline {\scriptsize\ttfamily \mnem{LD\_PKA}\imod{.zext\_or\_none} \dst{r0}, \src{r1}} \\[2pt]
\bottomrule
\end{longtable}
\end{center}
\par\textbf{Worked Example:} \emph{Compressed 1-bit sign-extend flag on narrow `LOAD`}
\codebox{\mnem{LOAD.i8}\imod{.sext.slot0} \dst{@r0}, \src{[r2:r3]}, \imm{\#0}}
Encoding \texttt{0} (\texttt{.sext}) sign-extends \texttt{0xFE -\textgreater{} 0xFFFFFFFE}; encoding \texttt{1} (\texttt{zext\_or\_none}) zero-extends \texttt{0xFE -\textgreater{} 0x000000FE}.

\subsection{Stencil (\texttt{stencil\_m})}
\label{enum:stencil_m}
Used by \texttt{ZS\_EMIT} and related instructions (2 encoding values):
\par\textbf{CSF ABI Interaction:} Exports shader stencil ref to \hyperref[sec:csf-zsd-blend]{\texttt{IDVSSR::ZSD}} (\texttt{stencil\_from\_shader=1}, word 0 bit 30).
\begin{center}\footnotesize
\begin{longtable}{@{}rll>{\raggedright\arraybackslash}p{0.62\linewidth}@{}}
\toprule
Enc & Value & Arch & Description \& Worked Example \\
\midrule\endhead
\texttt{0} & \texttt{\imod{none}} & v9--v14 & Default behavior (flag disabled for \texttt{stencil\_m}). \newline {\scriptsize\ttfamily \mnem{ZS\_EMIT}\imod{} \dst{r0}, \src{r1}} \\[2pt]
\texttt{1} & \texttt{\imod{stencil}} & v9--v14 & Emit fragment stencil output to tile depth/stencil unit. \newline {\scriptsize\ttfamily \mnem{ZS\_EMIT}\imod{.stencil} \dst{r0}, \src{r1}} \\[2pt]
\bottomrule
\end{longtable}
\end{center}
\par\textbf{Worked Example:} \emph{Exporting programmatic fragment stencil reference}
\codebox{\mnem{ZS\_EMIT}\imod{.stencil.slot0} \dst{@r60}, \src{k0}, \src{r1}, \src{r60}}
With \texttt{.stencil}, writes \texttt{r1[7:0]} as \texttt{gl\_FragStencilRefARB} and returns updated coverage in \texttt{@r60}.

\subsection{Submatrix Sel (\texttt{submatrix\_sel\_m})}
\label{enum:submatrix_sel_m}
Used by \texttt{MMUL} and related instructions (3 encoding values):

\begin{center}\footnotesize
\begin{longtable}{@{}rll>{\raggedright\arraybackslash}p{0.62\linewidth}@{}}
\toprule
Enc & Value & Arch & Description \& Worked Example \\
\midrule\endhead
\texttt{0} & \texttt{\imod{f0}} & v11--v14 & Select submatrix half 0 for matrix multiply. \newline {\scriptsize\ttfamily \mnem{MMUL}\imod{.f0} \dst{r0}, \src{r1}} \\[2pt]
\texttt{1} & \texttt{\imod{f1}} & v11--v14 & Select submatrix half 1 for matrix multiply. \newline {\scriptsize\ttfamily \mnem{MMUL}\imod{.f1} \dst{r0}, \src{r1}} \\[2pt]
\texttt{2} & \texttt{\imod{none}} & v11--v14 & Default behavior (flag disabled for \texttt{submatrix\_sel\_m}). \newline {\scriptsize\ttfamily \mnem{MMUL}\imod{} \dst{r0}, \src{r1}} \\[2pt]
\bottomrule
\end{longtable}
\end{center}
\par\textbf{Worked Example:} \emph{Selecting 4x4 submatrix halves in `MMUL.v2f16`}
\codebox{\mnem{MMUL.v2f16} \dst{r0}, \src{r1}\imod{.f0}, \src{r2}\imod{.f1}, \src{r3}}
\texttt{.f0} selects the first 4x4 FP16 half of matrix A and \texttt{.f1} selects the second 4x4 half of matrix B.

\subsection{Table Flog Mode (\texttt{table\_flog\_mode\_m})}
\label{enum:table_flog_mode_m}
Used by \texttt{FLOG\_TABLE} and related instructions (3 encoding values):

\begin{center}\footnotesize
\begin{longtable}{@{}rll>{\raggedright\arraybackslash}p{0.62\linewidth}@{}}
\toprule
Enc & Value & Arch & Description \& Worked Example \\
\midrule\endhead
\texttt{0} & \texttt{\imod{red}} & v9--v14 & Range-reduced logarithm table lookup. \newline {\scriptsize\ttfamily \mnem{FLOG\_TABLE}\imod{.red} \dst{r0}, \src{r1}} \\[2pt]
\texttt{1} & \texttt{\imod{base2}} & v9--v14 & Base-2 logarithm ($\log_2$) lookup table. \newline {\scriptsize\ttfamily \mnem{FLOG\_TABLE}\imod{.base2} \dst{r0}, \src{r1}} \\[2pt]
\texttt{2} & \texttt{\imod{natural}} & v9--v14 & Natural logarithm ($\ln$) lookup table. \newline {\scriptsize\ttfamily \mnem{FLOG\_TABLE}\imod{.natural} \dst{r0}, \src{r1}} \\[2pt]
\bottomrule
\end{longtable}
\end{center}
\par\textbf{Worked Example:} \emph{Selecting natural vs base-2 log polynomial}
\codebox{\mnem{FLOG\_TABLE.f32}\imod{.natural} \dst{r0}, \src{r1}}
\texttt{.base2} fetches \texttt{log2} LUT terms, \texttt{.natural} fetches \texttt{ln} terms, and \texttt{.red} fetches reduced-range coefficients.

\subsection{Table Flog Precision (\texttt{table\_flog\_precision\_m})}
\label{enum:table_flog_precision_m}
Used by \texttt{FLOG\_TABLE} and related instructions (3 encoding values):

\begin{center}\footnotesize
\begin{longtable}{@{}rll>{\raggedright\arraybackslash}p{0.62\linewidth}@{}}
\toprule
Enc & Value & Arch & Description \& Worked Example \\
\midrule\endhead
\texttt{0} & \texttt{\imod{none}} & v9--v14 & Default behavior (flag disabled for \texttt{table\_flog\_precision\_m}). \newline {\scriptsize\ttfamily \mnem{FLOG\_TABLE}\imod{} \dst{r0}, \src{r1}} \\[2pt]
\texttt{1} & \texttt{\imod{high}} & v9--v14 & High-precision logarithm table mode. \newline {\scriptsize\ttfamily \mnem{FLOG\_TABLE}\imod{.high} \dst{r0}, \src{r1}} \\[2pt]
\texttt{2} & \texttt{\imod{low}} & v9--v14 & Select low half of exponential lookup table entry. \newline {\scriptsize\ttfamily \mnem{FLOG\_TABLE}\imod{.low} \dst{r0}, \src{r1}} \\[2pt]
\bottomrule
\end{longtable}
\end{center}
\par\textbf{Worked Example:} \emph{High vs low precision `FLOG\_TABLE` lookup}
\codebox{\mnem{FLOG\_TABLE.f32}\imod{.base2.high} \dst{r0}, \src{r1}}
\texttt{.high} returns the upper-order polynomial coefficient; \texttt{.low} returns the fine refinement coefficient.

\subsection{Table Sincos Offset (\texttt{table\_sincos\_offset\_m})}
\label{enum:table_sincos_offset_m}
Used by \texttt{FCOS\_TABLE} and related instructions (2 encoding values):

\begin{center}\footnotesize
\begin{longtable}{@{}rll>{\raggedright\arraybackslash}p{0.62\linewidth}@{}}
\toprule
Enc & Value & Arch & Description \& Worked Example \\
\midrule\endhead
\texttt{0} & \texttt{\imod{none}} & v9--v14 & Default behavior (flag disabled for \texttt{table\_sincos\_offset\_m}). \newline {\scriptsize\ttfamily \mnem{FCOS\_TABLE}\imod{} \dst{r0}, \src{r1}} \\[2pt]
\texttt{1} & \texttt{\imod{offset}} & v9--v14 & Apply phase offset in sin/cos table lookup. \newline {\scriptsize\ttfamily \mnem{FCOS\_TABLE}\imod{.offset} \dst{r0}, \src{r1}} \\[2pt]
\bottomrule
\end{longtable}
\end{center}
\par\textbf{Worked Example:} \emph{Applying quadrant angle offset in `FSINCOS`}
\codebox{\mnem{FSINCOS.f32}\imod{.sin.offset} \dst{r0}, \src{r1}, \src{r2}}
Enables the quadrant phase offset adjustment during v14 hardware \texttt{FSINCOS} evaluation.

\subsection{Table Sincos Operation (\texttt{table\_sincos\_operation\_m})}
\label{enum:table_sincos_operation_m}
Used by \texttt{FSINCOS} and related instructions (2 encoding values):

\begin{center}\footnotesize
\begin{longtable}{@{}rll>{\raggedright\arraybackslash}p{0.62\linewidth}@{}}
\toprule
Enc & Value & Arch & Description \& Worked Example \\
\midrule\endhead
\texttt{0} & \texttt{\imod{cos}} & v14 & Compute cosine in \texttt{FSINCOS}. \newline {\scriptsize\ttfamily \mnem{FSINCOS}\imod{.cos} \dst{r0}, \src{r1}} \\[2pt]
\texttt{1} & \texttt{\imod{sin}} & v14 & Compute sine in \texttt{FSINCOS}. \newline {\scriptsize\ttfamily \mnem{FSINCOS}\imod{.sin} \dst{r0}, \src{r1}} \\[2pt]
\bottomrule
\end{longtable}
\end{center}
\par\textbf{Worked Example:} \emph{Selecting `sin` vs `cos` output on `FSINCOS`}
\codebox{\mnem{FSINCOS.f32}\imod{.cos} \dst{r0}, \src{r1}, \src{r2}}
Encoding \texttt{0} computes \texttt{cos(x)}; encoding \texttt{1} computes \texttt{sin(x)} from the same range-reduced inputs.

\subsection{Table Sincos Scale (\texttt{table\_sincos\_scale\_m})}
\label{enum:table_sincos_scale_m}
Used by \texttt{FSINCOS\_OFFSET} and related instructions (2 encoding values):

\begin{center}\footnotesize
\begin{longtable}{@{}rll>{\raggedright\arraybackslash}p{0.62\linewidth}@{}}
\toprule
Enc & Value & Arch & Description \& Worked Example \\
\midrule\endhead
\texttt{0} & \texttt{\imod{none}} & v9--v14 & Default behavior (flag disabled for \texttt{table\_sincos\_scale\_m}). \newline {\scriptsize\ttfamily \mnem{FSINCOS\_OFFSET}\imod{} \dst{r0}, \src{r1}} \\[2pt]
\texttt{1} & \texttt{\imod{scale}} & v9--v14 & Apply $\pi$-scaling in \texttt{FSINCOS\_OFFSET}. \newline {\scriptsize\ttfamily \mnem{FSINCOS\_OFFSET}\imod{.scale} \dst{r0}, \src{r1}} \\[2pt]
\bottomrule
\end{longtable}
\end{center}
\par\textbf{Worked Example:} \emph{Adding `+24` exponent scale in `FSINCOS\_OFFSET`}
\codebox{\mnem{FSINCOS\_OFFSET.u6}\imod{.scale} \dst{r0}, \src{r1}}
\texttt{.scale} boosts the returned FP32 offset exponent by \texttt{+24} to align with fixed-point range reduction.

\subsection{Tex Array Enable (\texttt{tex\_array\_enable\_m})}
\label{enum:tex_array_enable_m}
Used by \texttt{TEX\_FETCH} and related instructions (2 encoding values):

\begin{center}\footnotesize
\begin{longtable}{@{}rll>{\raggedright\arraybackslash}p{0.62\linewidth}@{}}
\toprule
Enc & Value & Arch & Description \& Worked Example \\
\midrule\endhead
\texttt{0} & \texttt{\imod{none}} & v9--v14 & Default behavior (flag disabled for \texttt{tex\_array\_enable\_m}). \newline {\scriptsize\ttfamily \mnem{TEX\_FETCH}\imod{} \dst{r0}, \src{r1}} \\[2pt]
\texttt{1} & \texttt{\imod{array}} & v9--v14 & Enable texture array layer index coordinate (\texttt{sampler2DArray}). \newline {\scriptsize\ttfamily \mnem{TEX\_FETCH}\imod{.array} \dst{r0}, \src{r1}} \\[2pt]
\bottomrule
\end{longtable}
\end{center}
\par\textbf{Worked Example:} \emph{Sampling a `sampler2DArray` layer}
\codebox{\mnem{TEX\_SINGLE}\imod{.tex2d.array.computed.slot0} \dst{@r0:r3}, \src{@r4:r6}, \src{u0}}
\texttt{.array} tells the texture unit to read the array layer index from staging register \texttt{@r6} after \texttt{(u, v)}.

\subsection{Tex Compare Enable (\texttt{tex\_compare\_enable\_m})}
\label{enum:tex_compare_enable_m}
Used by \texttt{TEX\_DUAL} and related instructions (2 encoding values):

\begin{center}\footnotesize
\begin{longtable}{@{}rll>{\raggedright\arraybackslash}p{0.62\linewidth}@{}}
\toprule
Enc & Value & Arch & Description \& Worked Example \\
\midrule\endhead
\texttt{0} & \texttt{\imod{none}} & v9--v14 & Default behavior (flag disabled for \texttt{tex\_compare\_enable\_m}). \newline {\scriptsize\ttfamily \mnem{TEX\_DUAL}\imod{} \dst{r0}, \src{r1}} \\[2pt]
\texttt{1} & \texttt{\imod{compare}} & v9--v14 & Enable shadow depth comparison reference coordinate (\texttt{sampler2DShadow}). \newline {\scriptsize\ttfamily \mnem{TEX\_DUAL}\imod{.compare} \dst{r0}, \src{r1}} \\[2pt]
\bottomrule
\end{longtable}
\end{center}
\par\textbf{Worked Example:} \emph{Hardware shadow-map PCF comparison}
\codebox{\mnem{TEX\_SINGLE}\imod{.tex2d.compare.computed.slot0} \dst{@r0}, \src{@r4:r6}, \src{u0}}
\texttt{.compare} reads depth reference \texttt{Dref} from \texttt{@r6} and returns filtered shadow factor in \texttt{[0.0, 1.0]}.

\subsection{Tex Coordinate Mode (\texttt{tex\_coordinate\_mode\_m})}
\label{enum:tex_coordinate_mode_m}
Used by \texttt{TEX\_GATHER} and related instructions (2 encoding values):

\begin{center}\footnotesize
\begin{longtable}{@{}rll>{\raggedright\arraybackslash}p{0.62\linewidth}@{}}
\toprule
Enc & Value & Arch & Description \& Worked Example \\
\midrule\endhead
\texttt{0} & \texttt{\imod{float\_coordinates}} & v9--v14 & Use normalized floating-point texture coordinates in $[0, 1]$. \newline {\scriptsize\ttfamily \mnem{TEX\_GATHER}\imod{.float\_coordinates} \dst{r0}, \src{r1}} \\[2pt]
\texttt{1} & \texttt{\imod{integer\_coordinates}} & v9--v14 & Use unnormalized integer texel coordinates (\texttt{texelFetch}). \newline {\scriptsize\ttfamily \mnem{TEX\_GATHER}\imod{.integer\_coordinates} \dst{r0}, \src{r1}} \\[2pt]
\bottomrule
\end{longtable}
\end{center}
\par\textbf{Worked Example:} \emph{Unnormalized integer coordinates for `texelFetch`}
\codebox{\mnem{TEX\_FETCH}\imod{.tex2d.integer\_coordinates.slot0} \dst{@r0:r3}, \src{@r4:r6}, \src{u0}}
\texttt{.integer\_coordinates} interprets \texttt{@r4:r5} as raw integer texel indices \texttt{(x, y)} rather than normalized \texttt{[0, 1]} floats.

\subsection{Tex Coordinate Or Derivative (\texttt{tex\_coordinate\_or\_derivative\_m})}
\label{enum:tex_coordinate_or_derivative_m}
Used by \texttt{TEX\_GRADIENT} and related instructions (3 encoding values):

\begin{center}\footnotesize
\begin{longtable}{@{}rll>{\raggedright\arraybackslash}p{0.62\linewidth}@{}}
\toprule
Enc & Value & Arch & Description \& Worked Example \\
\midrule\endhead
\texttt{0} & \texttt{\imod{none}} & v9 & Default behavior (flag disabled for \texttt{tex\_coordinate\_or\_derivative\_m}). \newline {\scriptsize\ttfamily \mnem{TEX\_GRADIENT}\imod{} \dst{r0}, \src{r1}} \\[2pt]
\texttt{1} & \texttt{\imod{force\_delta}} & v10--v14 & Force 2x2 quad derivative calculation even on inactive/helper lanes. \newline {\scriptsize\ttfamily \mnem{TEX\_GRADIENT}\imod{.force\_delta} \dst{r0}, \src{r1}} \\[2pt]
\texttt{2} & \texttt{\imod{derivative}} & v9 & Return screen-space coordinate derivatives instead of filtered texels. \newline {\scriptsize\ttfamily \mnem{TEX\_GRADIENT}\imod{.derivative} \dst{r0}, \src{r1}} \\[2pt]
\bottomrule
\end{longtable}
\end{center}
\par\textbf{Worked Example:} \emph{Requesting raw screen-space derivatives in `TEX\_GRADIENT`}
\codebox{\mnem{TEX\_GRADIENT}\imod{.tex2d.derivative.slot0} \dst{@r0:r3}, \src{@r4:r5}, \src{u0}}
\texttt{.derivative} outputs \texttt{(dU/dx, dV/dx, dU/dy, dV/dy)} instead of scalar LOD; \texttt{.force\_delta} forces quad deltas across helper lanes.

\subsection{Tex Derivative Enable (\texttt{tex\_derivative\_enable\_m})}
\label{enum:tex_derivative_enable_m}
Used by \texttt{TEX\_GRADIENT} and related instructions (2 encoding values):

\begin{center}\footnotesize
\begin{longtable}{@{}rll>{\raggedright\arraybackslash}p{0.62\linewidth}@{}}
\toprule
Enc & Value & Arch & Description \& Worked Example \\
\midrule\endhead
\texttt{0} & \texttt{\imod{none}} & v9--v14 & Default behavior (flag disabled for \texttt{tex\_derivative\_enable\_m}). \newline {\scriptsize\ttfamily \mnem{TEX\_GRADIENT}\imod{} \dst{r0}, \src{r1}} \\[2pt]
\texttt{1} & \texttt{\imod{derivative}} & v9--v14 & Return screen-space coordinate derivatives instead of filtered texels. \newline {\scriptsize\ttfamily \mnem{TEX\_GRADIENT}\imod{.derivative} \dst{r0}, \src{r1}} \\[2pt]
\bottomrule
\end{longtable}
\end{center}
\par\textbf{Worked Example:} \emph{Returning derivatives from `VAR\_TEX\_GRADIENT`}
\codebox{\mnem{VAR\_TEX\_GRADIENT}\imod{.tex2d.derivative.center.slot0} \dst{@r0:r3}, \src{u0}, \src{r59}, \src{u1}}
Enables raw derivative output directly from interpolated varying coordinates.

\subsection{Tex Dual Array Enable (\texttt{tex\_dual\_array\_enable\_m})}
\label{enum:tex_dual_array_enable_m}
Used by \texttt{TEX\_DUAL} and related instructions (2 encoding values):

\begin{center}\footnotesize
\begin{longtable}{@{}rll>{\raggedright\arraybackslash}p{0.62\linewidth}@{}}
\toprule
Enc & Value & Arch & Description \& Worked Example \\
\midrule\endhead
\texttt{0} & \texttt{\imod{none}} & v9--v14 & Default behavior (flag disabled for \texttt{tex\_dual\_array\_enable\_m}). \newline {\scriptsize\ttfamily \mnem{TEX\_DUAL}\imod{} \dst{r0}, \src{r1}} \\[2pt]
\texttt{1} & \texttt{\imod{array}} & v9--v14 & Enable texture array layer index coordinate (\texttt{sampler2DArray}). \newline {\scriptsize\ttfamily \mnem{TEX\_DUAL}\imod{.array} \dst{r0}, \src{r1}} \\[2pt]
\bottomrule
\end{longtable}
\end{center}
\par\textbf{Worked Example:} \emph{Arrayed dual texture sample}
\codebox{\mnem{TEX\_DUAL}\imod{.tex2d.array.computed.slot0} \dst{@r0:r7}, \src{@r8:r10}, \src{u0}}
Passes shared \texttt{(u, v, layer)} in \texttt{@r8:r10} to both paired texture lookups simultaneously.

\subsection{Tex Dual Lod Mode (\texttt{tex\_dual\_lod\_mode\_m})}
\label{enum:tex_dual_lod_mode_m}
Used by \texttt{TEX\_DUAL} and related instructions (7 encoding values):

\begin{center}\footnotesize
\begin{longtable}{@{}rll>{\raggedright\arraybackslash}p{0.62\linewidth}@{}}
\toprule
Enc & Value & Arch & Description \& Worked Example \\
\midrule\endhead
\texttt{0} & \texttt{\imod{none}} & v9 & Implicit zero LOD (base mip level 0, no derivatives). \newline {\scriptsize\ttfamily \mnem{TEX\_DUAL}\imod{} \dst{@r0:r3}, \src{@r4:r5}} \\[2pt]
\texttt{1} & \texttt{\imod{computed}} & v9 & Hardware 2x2 quad finite-difference LOD (\texttt{texture()}). \newline {\scriptsize\ttfamily \mnem{TEX\_DUAL}\imod{.computed} \dst{@r0:r3}, \src{@r4:r5}} \\[2pt]
\texttt{3} & \texttt{\imod{computed\_force\_delta}} & v10--v14 & Computed LOD forcing quad delta evaluation even in helper lanes. \newline {\scriptsize\ttfamily \mnem{TEX\_DUAL}\imod{.computed\_force\_delta} \dst{@r0:r3}, \src{@r4:r5}} \\[2pt]
\texttt{4} & \texttt{\imod{explicit}} & v9 & Explicit LOD supplied in staging register (\texttt{textureLod()}). \newline {\scriptsize\ttfamily \mnem{TEX\_DUAL}\imod{.explicit} \dst{@r0:r3}, \src{@r4:r5}} \\[2pt]
\texttt{5} & \texttt{\imod{computed\_bias}} & v9 & Computed quad LOD plus shader bias (\texttt{texture(..., bias)}). \newline {\scriptsize\ttfamily \mnem{TEX\_DUAL}\imod{.computed\_bias} \dst{@r0:r3}, \src{@r4:r5}} \\[2pt]
\texttt{6} & \texttt{\imod{grdesc}} & v9 & Explicit $dP/dx, dP/dy$ gradients in staging regs (\texttt{textureGrad()}). \newline {\scriptsize\ttfamily \mnem{TEX\_DUAL}\imod{.grdesc} \dst{@r0:r3}, \src{@r4:r5}} \\[2pt]
\texttt{7} & \texttt{\imod{computed\_bias\_force\_delta}} & v10--v14 & Computed LOD + bias with forced quad delta evaluation. \newline {\scriptsize\ttfamily \mnem{TEX\_DUAL}\imod{.computed\_bias\_force\_delta} \dst{@r0:r3}, \src{@r4:r5}} \\[2pt]
\bottomrule
\end{longtable}
\end{center}
\par\textbf{Worked Example:} \emph{Explicit LOD on dual texture sample}
\codebox{\mnem{TEX\_DUAL}\imod{.tex2d.explicit.slot0} \dst{@r0:r7}, \src{@r8:r10}, \src{u0}}
Applies explicit LOD \texttt{@r10} to both paired texture fetches in \texttt{TEX\_DUAL}.

\subsection{Tex Dual Primary Width (\texttt{tex\_dual\_primary\_width\_m})}
\label{enum:tex_dual_primary_width_m}
Used by \texttt{TEX\_DUAL} and related instructions (8 encoding values):

\begin{center}\footnotesize
\begin{longtable}{@{}rll>{\raggedright\arraybackslash}p{0.62\linewidth}@{}}
\toprule
Enc & Value & Arch & Description \& Worked Example \\
\midrule\endhead
\texttt{0} & \texttt{\imod{primary\_dst16}} & v9--v14 & Primary texture return format \texttt{primary\_dst16} (16-bit packed staging register elements). \newline {\scriptsize\ttfamily \mnem{TEX\_DUAL}\imod{.primary\_dst16} \dst{@r0:r3}, \src{@r4:r5}} \\[2pt]
\texttt{0} & \texttt{\imod{primary\_f16}} & v9--v14 & Primary texture return format \texttt{primary\_f16} (16-bit packed staging register elements). \newline {\scriptsize\ttfamily \mnem{TEX\_DUAL}\imod{.primary\_f16} \dst{@r0:r3}, \src{@r4:r5}} \\[2pt]
\texttt{0} & \texttt{\imod{primary\_s16}} & v9--v14 & Primary texture return format \texttt{primary\_s16} (16-bit packed staging register elements). \newline {\scriptsize\ttfamily \mnem{TEX\_DUAL}\imod{.primary\_s16} \dst{@r0:r3}, \src{@r4:r5}} \\[2pt]
\texttt{0} & \texttt{\imod{primary\_u16}} & v9--v14 & Primary texture return format \texttt{primary\_u16} (16-bit packed staging register elements). \newline {\scriptsize\ttfamily \mnem{TEX\_DUAL}\imod{.primary\_u16} \dst{@r0:r3}, \src{@r4:r5}} \\[2pt]
\texttt{1} & \texttt{\imod{primary\_dst32}} & v9--v14 & Primary texture return format \texttt{primary\_dst32} (32-bit staging register elements). \newline {\scriptsize\ttfamily \mnem{TEX\_DUAL}\imod{.primary\_dst32} \dst{@r0:r3}, \src{@r4:r5}} \\[2pt]
\texttt{1} & \texttt{\imod{primary\_f32}} & v9--v14 & Primary texture return format \texttt{primary\_f32} (32-bit staging register elements). \newline {\scriptsize\ttfamily \mnem{TEX\_DUAL}\imod{.primary\_f32} \dst{@r0:r3}, \src{@r4:r5}} \\[2pt]
\texttt{1} & \texttt{\imod{primary\_s32}} & v9--v14 & Primary texture return format \texttt{primary\_s32} (32-bit staging register elements). \newline {\scriptsize\ttfamily \mnem{TEX\_DUAL}\imod{.primary\_s32} \dst{@r0:r3}, \src{@r4:r5}} \\[2pt]
\texttt{1} & \texttt{\imod{primary\_u32}} & v9--v14 & Primary texture return format \texttt{primary\_u32} (32-bit staging register elements). \newline {\scriptsize\ttfamily \mnem{TEX\_DUAL}\imod{.primary\_u32} \dst{@r0:r3}, \src{@r4:r5}} \\[2pt]
\bottomrule
\end{longtable}
\end{center}
\par\textbf{Worked Example:} \emph{FP16 primary output width in `TEX\_DUAL`}
\codebox{\mnem{TEX\_DUAL}\imod{.tex2d.primary\_f16.secondary\_f32.slot0} \dst{@r0:r5}, \src{@r6:r7}, \src{u0}}
Packs the primary texture's RGBA into 2 FP16 registers (\texttt{@r0:r1}) while keeping the secondary texture in 4 FP32 registers (\texttt{@r2:r5}).

\subsection{Tex Dual Primary Write Mask (\texttt{tex\_dual\_primary\_write\_mask\_m})}
\label{enum:tex_dual_primary_write_mask_m}
Used by \texttt{TEX\_DUAL} and related instructions (15 encoding values):

\begin{center}\footnotesize
\begin{longtable}{@{}rll>{\raggedright\arraybackslash}p{0.62\linewidth}@{}}
\toprule
Enc & Value & Arch & Description \& Worked Example \\
\midrule\endhead
\texttt{1} & \texttt{\imod{primary\_r}} & v9--v14 & Primary channel write mask: write only R component(s) to consecutive staging GPRs. \newline {\scriptsize\ttfamily \mnem{TEX\_DUAL}\imod{.primary\_r} \dst{@r0:r3}, \src{@r4:r5}} \\[2pt]
\texttt{2} & \texttt{\imod{primary\_g}} & v9--v14 & Primary channel write mask: write only G component(s) to consecutive staging GPRs. \newline {\scriptsize\ttfamily \mnem{TEX\_DUAL}\imod{.primary\_g} \dst{@r0:r3}, \src{@r4:r5}} \\[2pt]
\texttt{3} & \texttt{\imod{primary\_rg}} & v9--v14 & Primary channel write mask: write only RG component(s) to consecutive staging GPRs. \newline {\scriptsize\ttfamily \mnem{TEX\_DUAL}\imod{.primary\_rg} \dst{@r0:r3}, \src{@r4:r5}} \\[2pt]
\texttt{4} & \texttt{\imod{primary\_b}} & v9--v14 & Primary channel write mask: write only B component(s) to consecutive staging GPRs. \newline {\scriptsize\ttfamily \mnem{TEX\_DUAL}\imod{.primary\_b} \dst{@r0:r3}, \src{@r4:r5}} \\[2pt]
\texttt{5} & \texttt{\imod{primary\_rb}} & v9--v14 & Primary channel write mask: write only RB component(s) to consecutive staging GPRs. \newline {\scriptsize\ttfamily \mnem{TEX\_DUAL}\imod{.primary\_rb} \dst{@r0:r3}, \src{@r4:r5}} \\[2pt]
\texttt{6} & \texttt{\imod{primary\_gb}} & v9--v14 & Primary channel write mask: write only GB component(s) to consecutive staging GPRs. \newline {\scriptsize\ttfamily \mnem{TEX\_DUAL}\imod{.primary\_gb} \dst{@r0:r3}, \src{@r4:r5}} \\[2pt]
\texttt{7} & \texttt{\imod{primary\_rgb}} & v9--v14 & Primary channel write mask: write only RGB component(s) to consecutive staging GPRs. \newline {\scriptsize\ttfamily \mnem{TEX\_DUAL}\imod{.primary\_rgb} \dst{@r0:r3}, \src{@r4:r5}} \\[2pt]
\texttt{8} & \texttt{\imod{primary\_a}} & v9--v14 & Primary channel write mask: write only A component(s) to consecutive staging GPRs. \newline {\scriptsize\ttfamily \mnem{TEX\_DUAL}\imod{.primary\_a} \dst{@r0:r3}, \src{@r4:r5}} \\[2pt]
\texttt{9} & \texttt{\imod{primary\_ra}} & v9--v14 & Primary channel write mask: write only RA component(s) to consecutive staging GPRs. \newline {\scriptsize\ttfamily \mnem{TEX\_DUAL}\imod{.primary\_ra} \dst{@r0:r3}, \src{@r4:r5}} \\[2pt]
\texttt{10} & \texttt{\imod{primary\_ga}} & v9--v14 & Primary channel write mask: write only GA component(s) to consecutive staging GPRs. \newline {\scriptsize\ttfamily \mnem{TEX\_DUAL}\imod{.primary\_ga} \dst{@r0:r3}, \src{@r4:r5}} \\[2pt]
\texttt{11} & \texttt{\imod{primary\_rga}} & v9--v14 & Primary channel write mask: write only RGA component(s) to consecutive staging GPRs. \newline {\scriptsize\ttfamily \mnem{TEX\_DUAL}\imod{.primary\_rga} \dst{@r0:r3}, \src{@r4:r5}} \\[2pt]
\texttt{12} & \texttt{\imod{primary\_ba}} & v9--v14 & Primary channel write mask: write only BA component(s) to consecutive staging GPRs. \newline {\scriptsize\ttfamily \mnem{TEX\_DUAL}\imod{.primary\_ba} \dst{@r0:r3}, \src{@r4:r5}} \\[2pt]
\texttt{13} & \texttt{\imod{primary\_rba}} & v9--v14 & Primary channel write mask: write only RBA component(s) to consecutive staging GPRs. \newline {\scriptsize\ttfamily \mnem{TEX\_DUAL}\imod{.primary\_rba} \dst{@r0:r3}, \src{@r4:r5}} \\[2pt]
\texttt{14} & \texttt{\imod{primary\_gba}} & v9--v14 & Primary channel write mask: write only GBA component(s) to consecutive staging GPRs. \newline {\scriptsize\ttfamily \mnem{TEX\_DUAL}\imod{.primary\_gba} \dst{@r0:r3}, \src{@r4:r5}} \\[2pt]
\texttt{15} & \texttt{\imod{primary\_rgba}} & v9--v14 & Primary channel write mask: write only RGBA component(s) to consecutive staging GPRs. \newline {\scriptsize\ttfamily \mnem{TEX\_DUAL}\imod{.primary\_rgba} \dst{@r0:r3}, \src{@r4:r5}} \\[2pt]
\bottomrule
\end{longtable}
\end{center}
\par\textbf{Worked Example:} \emph{RGB-only primary write mask in `TEX\_DUAL`}
\codebox{\mnem{TEX\_DUAL}\imod{.tex2d.primary\_rgb.secondary\_r.slot0} \dst{@r0:r3}, \src{@r4:r5}, \src{u0}}
Suppresses unused primary alpha writeback, writing only \texttt{.rgb} to \texttt{@r0:r2}.

\subsection{Tex Dual Secondary Width (\texttt{tex\_dual\_secondary\_width\_m})}
\label{enum:tex_dual_secondary_width_m}
Used by \texttt{TEX\_DUAL} and related instructions (8 encoding values):

\begin{center}\footnotesize
\begin{longtable}{@{}rll>{\raggedright\arraybackslash}p{0.62\linewidth}@{}}
\toprule
Enc & Value & Arch & Description \& Worked Example \\
\midrule\endhead
\texttt{0} & \texttt{\imod{secondary\_dst16}} & v9--v14 & Secondary texture return format \texttt{secondary\_dst16} (16-bit packed staging register elements). \newline {\scriptsize\ttfamily \mnem{TEX\_DUAL}\imod{.secondary\_dst16} \dst{@r0:r3}, \src{@r4:r5}} \\[2pt]
\texttt{0} & \texttt{\imod{secondary\_f16}} & v9--v14 & Secondary texture return format \texttt{secondary\_f16} (16-bit packed staging register elements). \newline {\scriptsize\ttfamily \mnem{TEX\_DUAL}\imod{.secondary\_f16} \dst{@r0:r3}, \src{@r4:r5}} \\[2pt]
\texttt{0} & \texttt{\imod{secondary\_s16}} & v9--v14 & Secondary texture return format \texttt{secondary\_s16} (16-bit packed staging register elements). \newline {\scriptsize\ttfamily \mnem{TEX\_DUAL}\imod{.secondary\_s16} \dst{@r0:r3}, \src{@r4:r5}} \\[2pt]
\texttt{0} & \texttt{\imod{secondary\_u16}} & v9--v14 & Secondary texture return format \texttt{secondary\_u16} (16-bit packed staging register elements). \newline {\scriptsize\ttfamily \mnem{TEX\_DUAL}\imod{.secondary\_u16} \dst{@r0:r3}, \src{@r4:r5}} \\[2pt]
\texttt{1} & \texttt{\imod{secondary\_dst32}} & v9--v14 & Secondary texture return format \texttt{secondary\_dst32} (32-bit staging register elements). \newline {\scriptsize\ttfamily \mnem{TEX\_DUAL}\imod{.secondary\_dst32} \dst{@r0:r3}, \src{@r4:r5}} \\[2pt]
\texttt{1} & \texttt{\imod{secondary\_f32}} & v9--v14 & Secondary texture return format \texttt{secondary\_f32} (32-bit staging register elements). \newline {\scriptsize\ttfamily \mnem{TEX\_DUAL}\imod{.secondary\_f32} \dst{@r0:r3}, \src{@r4:r5}} \\[2pt]
\texttt{1} & \texttt{\imod{secondary\_s32}} & v9--v14 & Secondary texture return format \texttt{secondary\_s32} (32-bit staging register elements). \newline {\scriptsize\ttfamily \mnem{TEX\_DUAL}\imod{.secondary\_s32} \dst{@r0:r3}, \src{@r4:r5}} \\[2pt]
\texttt{1} & \texttt{\imod{secondary\_u32}} & v9--v14 & Secondary texture return format \texttt{secondary\_u32} (32-bit staging register elements). \newline {\scriptsize\ttfamily \mnem{TEX\_DUAL}\imod{.secondary\_u32} \dst{@r0:r3}, \src{@r4:r5}} \\[2pt]
\bottomrule
\end{longtable}
\end{center}
\par\textbf{Worked Example:} \emph{FP16 secondary output width in `TEX\_DUAL`}
\codebox{\mnem{TEX\_DUAL}\imod{.tex2d.primary\_f32.secondary\_f16.slot0} \dst{@r0:r5}, \src{@r6:r7}, \src{u0}}
Returns secondary texture channels as packed FP16 pairs (\texttt{@r4:r5}) to cut register pressure.

\subsection{Tex Dual Secondary Write Mask (\texttt{tex\_dual\_secondary\_write\_mask\_m})}
\label{enum:tex_dual_secondary_write_mask_m}
Used by \texttt{TEX\_DUAL} and related instructions (4 encoding values):

\begin{center}\footnotesize
\begin{longtable}{@{}rll>{\raggedright\arraybackslash}p{0.62\linewidth}@{}}
\toprule
Enc & Value & Arch & Description \& Worked Example \\
\midrule\endhead
\texttt{1} & \texttt{\imod{secondary\_r}} & v9--v14 & Secondary channel write mask: write only R component(s) to consecutive staging GPRs. \newline {\scriptsize\ttfamily \mnem{TEX\_DUAL}\imod{.secondary\_r} \dst{@r0:r3}, \src{@r4:r5}} \\[2pt]
\texttt{3} & \texttt{\imod{secondary\_rg}} & v9--v14 & Secondary channel write mask: write only RG component(s) to consecutive staging GPRs. \newline {\scriptsize\ttfamily \mnem{TEX\_DUAL}\imod{.secondary\_rg} \dst{@r0:r3}, \src{@r4:r5}} \\[2pt]
\texttt{7} & \texttt{\imod{secondary\_rgb}} & v9--v14 & Secondary channel write mask: write only RGB component(s) to consecutive staging GPRs. \newline {\scriptsize\ttfamily \mnem{TEX\_DUAL}\imod{.secondary\_rgb} \dst{@r0:r3}, \src{@r4:r5}} \\[2pt]
\texttt{15} & \texttt{\imod{secondary\_rgba}} & v9--v14 & Secondary channel write mask: write only RGBA component(s) to consecutive staging GPRs. \newline {\scriptsize\ttfamily \mnem{TEX\_DUAL}\imod{.secondary\_rgba} \dst{@r0:r3}, \src{@r4:r5}} \\[2pt]
\bottomrule
\end{longtable}
\end{center}
\par\textbf{Worked Example:} \emph{Single-channel secondary mask in `TEX\_DUAL`}
\codebox{\mnem{TEX\_DUAL}\imod{.tex2d.primary\_rgba.secondary\_r.slot0} \dst{@r0:r4}, \src{@r5:r6}, \src{u0}}
Writes only the red channel of the second texture into \texttt{@r4} (e.g., ambient occlusion map).

\subsection{Tex Footprint (\texttt{tex\_footprint\_m})}
\label{enum:tex_footprint_m}
Used by \texttt{TEX\_WEIGHTS} and related instructions (1 encoding values):

\begin{center}\footnotesize
\begin{longtable}{@{}rll>{\raggedright\arraybackslash}p{0.62\linewidth}@{}}
\toprule
Enc & Value & Arch & Description \& Worked Example \\
\midrule\endhead
\texttt{0} & \texttt{\imod{footprint\_2x2}} & v9 & Evaluate 2x2 bilinear texel weight footprint. \newline {\scriptsize\ttfamily \mnem{TEX\_WEIGHTS}\imod{.footprint\_2x2} \dst{r0}, \src{r1}} \\[2pt]
\bottomrule
\end{longtable}
\end{center}
\par\textbf{Worked Example:} \emph{2x2 bilinear footprint weights query (v9)}
\codebox{\mnem{TEX\_WEIGHTS}\imod{.tex2d.footprint\_2x2.slot0} \dst{@r0:r3}, \src{@r4:r5}, \src{u0}}
Returns the 4 normalized bilinear weights for the 2x2 footprint surrounding \texttt{@r4:r5}.

\subsection{Tex Force Delta Enable (\texttt{tex\_force\_delta\_enable\_m})}
\label{enum:tex_force_delta_enable_m}
Used by \texttt{TEX\_GRADIENT} and related instructions (2 encoding values):

\begin{center}\footnotesize
\begin{longtable}{@{}rll>{\raggedright\arraybackslash}p{0.62\linewidth}@{}}
\toprule
Enc & Value & Arch & Description \& Worked Example \\
\midrule\endhead
\texttt{0} & \texttt{\imod{none}} & v10--v14 & Default behavior (flag disabled for \texttt{tex\_force\_delta\_enable\_m}). \newline {\scriptsize\ttfamily \mnem{TEX\_GRADIENT}\imod{} \dst{r0}, \src{r1}} \\[2pt]
\texttt{1} & \texttt{\imod{force\_delta}} & v10--v14 & Force 2x2 quad derivative calculation even on inactive/helper lanes. \newline {\scriptsize\ttfamily \mnem{TEX\_GRADIENT}\imod{.force\_delta} \dst{r0}, \src{r1}} \\[2pt]
\bottomrule
\end{longtable}
\end{center}
\par\textbf{Worked Example:} \emph{Forcing quad delta evaluation on helper lanes}
\codebox{\mnem{VAR\_TEX\_BUF\_GRADIENT}\imod{.tex2d.force\_delta.slot0} \dst{@r0:r1}, \src{u0}, \src{r59}, \src{r4}}
Ensures valid 2x2 finite differences even when neighboring lanes in the quad are helper/discarded pixels.

\subsection{Tex Gradient Skip (\texttt{tex\_gradient\_skip\_m})}
\label{enum:tex_gradient_skip_m}
Used by \texttt{TEX\_GRADIENT} and related instructions (2 encoding values):

\begin{center}\footnotesize
\begin{longtable}{@{}rll>{\raggedright\arraybackslash}p{0.62\linewidth}@{}}
\toprule
Enc & Value & Arch & Description \& Worked Example \\
\midrule\endhead
\texttt{0} & \texttt{\imod{none}} & v9--v14 & Default behavior (flag disabled for \texttt{tex\_gradient\_skip\_m}). \newline {\scriptsize\ttfamily \mnem{TEX\_GRADIENT}\imod{} \dst{r0}, \src{r1}} \\[2pt]
\texttt{1} & \texttt{\imod{skip}} & v9--v14 & Skip texture fetch for helper/discarded fragment lanes. \newline {\scriptsize\ttfamily \mnem{TEX\_GRADIENT}\imod{.skip} \dst{r0}, \src{r1}} \\[2pt]
\bottomrule
\end{longtable}
\end{center}
\par\textbf{Worked Example:} \emph{Skipping helper lanes on `TEX\_GRADIENT`}
\codebox{\mnem{TEX\_GRADIENT}\imod{.tex2d.skip.slot0} \dst{@r0:r3}, \src{@r4:r9}, \src{u0}}
Skips texture memory fetches on discarded/helper lanes while preserving quad LOD math.

\subsection{Tex Gradient Width (\texttt{tex\_gradient\_width\_m})}
\label{enum:tex_gradient_width_m}
Used by \texttt{TEX\_GRADIENT} and related instructions (8 encoding values):

\begin{center}\footnotesize
\begin{longtable}{@{}rll>{\raggedright\arraybackslash}p{0.62\linewidth}@{}}
\toprule
Enc & Value & Arch & Description \& Worked Example \\
\midrule\endhead
\texttt{0} & \texttt{\imod{dst16}} & v9--v14 & texture return format \texttt{dst16} (16-bit packed staging register elements). \newline {\scriptsize\ttfamily \mnem{TEX\_GRADIENT}\imod{.dst16} \dst{@r0:r3}, \src{@r4:r5}} \\[2pt]
\texttt{0} & \texttt{\imod{f16}} & v9--v14 & texture return format \texttt{f16} (16-bit packed staging register elements). \newline {\scriptsize\ttfamily \mnem{TEX\_GRADIENT}\imod{.f16} \dst{@r0:r3}, \src{@r4:r5}} \\[2pt]
\texttt{0} & \texttt{\imod{s16}} & v9--v14 & texture return format \texttt{s16} (16-bit packed staging register elements). \newline {\scriptsize\ttfamily \mnem{TEX\_GRADIENT}\imod{.s16} \dst{@r0:r3}, \src{@r4:r5}} \\[2pt]
\texttt{0} & \texttt{\imod{u16}} & v9--v14 & texture return format \texttt{u16} (16-bit packed staging register elements). \newline {\scriptsize\ttfamily \mnem{TEX\_GRADIENT}\imod{.u16} \dst{@r0:r3}, \src{@r4:r5}} \\[2pt]
\texttt{1} & \texttt{\imod{dst32}} & v9--v14 & texture return format \texttt{dst32} (32-bit staging register elements). \newline {\scriptsize\ttfamily \mnem{TEX\_GRADIENT}\imod{.dst32} \dst{@r0:r3}, \src{@r4:r5}} \\[2pt]
\texttt{1} & \texttt{\imod{f32}} & v9--v14 & texture return format \texttt{f32} (32-bit staging register elements). \newline {\scriptsize\ttfamily \mnem{TEX\_GRADIENT}\imod{.f32} \dst{@r0:r3}, \src{@r4:r5}} \\[2pt]
\texttt{1} & \texttt{\imod{s32}} & v9--v14 & texture return format \texttt{s32} (32-bit staging register elements). \newline {\scriptsize\ttfamily \mnem{TEX\_GRADIENT}\imod{.s32} \dst{@r0:r3}, \src{@r4:r5}} \\[2pt]
\texttt{1} & \texttt{\imod{u32}} & v9--v14 & texture return format \texttt{u32} (32-bit staging register elements). \newline {\scriptsize\ttfamily \mnem{TEX\_GRADIENT}\imod{.u32} \dst{@r0:r3}, \src{@r4:r5}} \\[2pt]
\bottomrule
\end{longtable}
\end{center}
\par\textbf{Worked Example:} \emph{FP16 return format on `TEX\_GRADIENT`}
\codebox{\mnem{TEX\_GRADIENT}\imod{.tex2d.f16.rgba.slot0} \dst{@r0:r1}, \src{@r4:r9}, \src{u0}}
Returns the filtered RGBA result packed as 4x FP16 in \texttt{@r0:r1} instead of 4x FP32.

\subsection{Tex Gradient Write Mask (\texttt{tex\_gradient\_write\_mask\_m})}
\label{enum:tex_gradient_write_mask_m}
Used by \texttt{TEX\_GRADIENT} and related instructions (15 encoding values):

\begin{center}\footnotesize
\begin{longtable}{@{}rll>{\raggedright\arraybackslash}p{0.62\linewidth}@{}}
\toprule
Enc & Value & Arch & Description \& Worked Example \\
\midrule\endhead
\texttt{1} & \texttt{\imod{r}} & v9--v14 & channel write mask: write only R component(s) to consecutive staging GPRs. \newline {\scriptsize\ttfamily \mnem{TEX\_GRADIENT}\imod{.r} \dst{@r0:r3}, \src{@r4:r5}} \\[2pt]
\texttt{2} & \texttt{\imod{g}} & v9--v14 & channel write mask: write only G component(s) to consecutive staging GPRs. \newline {\scriptsize\ttfamily \mnem{TEX\_GRADIENT}\imod{.g} \dst{@r0:r3}, \src{@r4:r5}} \\[2pt]
\texttt{3} & \texttt{\imod{rg}} & v9--v14 & channel write mask: write only RG component(s) to consecutive staging GPRs. \newline {\scriptsize\ttfamily \mnem{TEX\_GRADIENT}\imod{.rg} \dst{@r0:r3}, \src{@r4:r5}} \\[2pt]
\texttt{4} & \texttt{\imod{b}} & v9--v14 & channel write mask: write only B component(s) to consecutive staging GPRs. \newline {\scriptsize\ttfamily \mnem{TEX\_GRADIENT}\imod{.b} \dst{@r0:r3}, \src{@r4:r5}} \\[2pt]
\texttt{5} & \texttt{\imod{rb}} & v9--v14 & channel write mask: write only RB component(s) to consecutive staging GPRs. \newline {\scriptsize\ttfamily \mnem{TEX\_GRADIENT}\imod{.rb} \dst{@r0:r3}, \src{@r4:r5}} \\[2pt]
\texttt{6} & \texttt{\imod{gb}} & v9--v14 & channel write mask: write only GB component(s) to consecutive staging GPRs. \newline {\scriptsize\ttfamily \mnem{TEX\_GRADIENT}\imod{.gb} \dst{@r0:r3}, \src{@r4:r5}} \\[2pt]
\texttt{7} & \texttt{\imod{rgb}} & v9--v14 & channel write mask: write only RGB component(s) to consecutive staging GPRs. \newline {\scriptsize\ttfamily \mnem{TEX\_GRADIENT}\imod{.rgb} \dst{@r0:r3}, \src{@r4:r5}} \\[2pt]
\texttt{8} & \texttt{\imod{a}} & v9--v14 & channel write mask: write only A component(s) to consecutive staging GPRs. \newline {\scriptsize\ttfamily \mnem{TEX\_GRADIENT}\imod{.a} \dst{@r0:r3}, \src{@r4:r5}} \\[2pt]
\texttt{9} & \texttt{\imod{ra}} & v9--v14 & channel write mask: write only RA component(s) to consecutive staging GPRs. \newline {\scriptsize\ttfamily \mnem{TEX\_GRADIENT}\imod{.ra} \dst{@r0:r3}, \src{@r4:r5}} \\[2pt]
\texttt{10} & \texttt{\imod{ga}} & v9--v14 & channel write mask: write only GA component(s) to consecutive staging GPRs. \newline {\scriptsize\ttfamily \mnem{TEX\_GRADIENT}\imod{.ga} \dst{@r0:r3}, \src{@r4:r5}} \\[2pt]
\texttt{11} & \texttt{\imod{rga}} & v9--v14 & channel write mask: write only RGA component(s) to consecutive staging GPRs. \newline {\scriptsize\ttfamily \mnem{TEX\_GRADIENT}\imod{.rga} \dst{@r0:r3}, \src{@r4:r5}} \\[2pt]
\texttt{12} & \texttt{\imod{ba}} & v9--v14 & channel write mask: write only BA component(s) to consecutive staging GPRs. \newline {\scriptsize\ttfamily \mnem{TEX\_GRADIENT}\imod{.ba} \dst{@r0:r3}, \src{@r4:r5}} \\[2pt]
\texttt{13} & \texttt{\imod{rba}} & v9--v14 & channel write mask: write only RBA component(s) to consecutive staging GPRs. \newline {\scriptsize\ttfamily \mnem{TEX\_GRADIENT}\imod{.rba} \dst{@r0:r3}, \src{@r4:r5}} \\[2pt]
\texttt{14} & \texttt{\imod{gba}} & v9--v14 & channel write mask: write only GBA component(s) to consecutive staging GPRs. \newline {\scriptsize\ttfamily \mnem{TEX\_GRADIENT}\imod{.gba} \dst{@r0:r3}, \src{@r4:r5}} \\[2pt]
\texttt{15} & \texttt{\imod{rgba}} & v9--v14 & channel write mask: write only RGBA component(s) to consecutive staging GPRs. \newline {\scriptsize\ttfamily \mnem{TEX\_GRADIENT}\imod{.rgba} \dst{@r0:r3}, \src{@r4:r5}} \\[2pt]
\bottomrule
\end{longtable}
\end{center}
\par\textbf{Worked Example:} \emph{Two-channel `.rg` mask on `TEX\_GRADIENT`}
\codebox{\mnem{TEX\_GRADIENT}\imod{.tex2d.f32.rg.slot0} \dst{@r0:r1}, \src{@r4:r9}, \src{u0}}
Writes only \texttt{.rg} components into \texttt{@r0:r1} (e.g., 2-channel normal map or \texttt{textureQueryLod}).

\subsection{Tex Lod Bias Disable (\texttt{tex\_lod\_bias\_disable\_m})}
\label{enum:tex_lod_bias_disable_m}
Used by \texttt{TEX\_GRADIENT} and related instructions (2 encoding values):

\begin{center}\footnotesize
\begin{longtable}{@{}rll>{\raggedright\arraybackslash}p{0.62\linewidth}@{}}
\toprule
Enc & Value & Arch & Description \& Worked Example \\
\midrule\endhead
\texttt{0} & \texttt{\imod{none}} & v9--v14 & Default behavior (flag disabled for \texttt{tex\_lod\_bias\_disable\_m}). \newline {\scriptsize\ttfamily \mnem{TEX\_GRADIENT}\imod{} \dst{r0}, \src{r1}} \\[2pt]
\texttt{1} & \texttt{\imod{disable\_lod\_bias}} & v9--v14 & Ignore sampler/shader LOD bias during gradient calculation. \newline {\scriptsize\ttfamily \mnem{TEX\_GRADIENT}\imod{.disable\_lod\_bias} \dst{r0}, \src{r1}} \\[2pt]
\bottomrule
\end{longtable}
\end{center}
\par\textbf{Worked Example:} \emph{Ignoring sampler LOD bias for raw `textureQueryLod`}
\codebox{\mnem{TEX\_GRADIENT}\imod{.tex2d.disable\_lod\_bias.slot0} \dst{@r0:r1}, \src{@r4:r5}, \src{u0}}
Suppresses sampler descriptor LOD bias when querying unbiased mipmap level.

\subsection{Tex Lod Clamp Disable (\texttt{tex\_lod\_clamp\_disable\_m})}
\label{enum:tex_lod_clamp_disable_m}
Used by \texttt{TEX\_GRADIENT} and related instructions (2 encoding values):

\begin{center}\footnotesize
\begin{longtable}{@{}rll>{\raggedright\arraybackslash}p{0.62\linewidth}@{}}
\toprule
Enc & Value & Arch & Description \& Worked Example \\
\midrule\endhead
\texttt{0} & \texttt{\imod{none}} & v9--v14 & Default behavior (flag disabled for \texttt{tex\_lod\_clamp\_disable\_m}). \newline {\scriptsize\ttfamily \mnem{TEX\_GRADIENT}\imod{} \dst{r0}, \src{r1}} \\[2pt]
\texttt{1} & \texttt{\imod{disable\_lod\_clamp}} & v9--v14 & Ignore sampler min/max LOD clamp bounds. \newline {\scriptsize\ttfamily \mnem{TEX\_GRADIENT}\imod{.disable\_lod\_clamp} \dst{r0}, \src{r1}} \\[2pt]
\bottomrule
\end{longtable}
\end{center}
\par\textbf{Worked Example:} \emph{Ignoring sampler min/max LOD clamp}
\codebox{\mnem{TEX\_GRADIENT}\imod{.tex2d.disable\_lod\_clamp.slot0} \dst{@r0:r1}, \src{@r4:r5}, \src{u0}}
Returns the unclamped mipmap level \texttt{lambda} required by the second component of GLSL \texttt{textureQueryLod}.

\subsection{Tex Projection Enable (\texttt{tex\_projection\_enable\_m})}
\label{enum:tex_projection_enable_m}
Used by \texttt{TEX\_DUAL} and related instructions (2 encoding values):

\begin{center}\footnotesize
\begin{longtable}{@{}rll>{\raggedright\arraybackslash}p{0.62\linewidth}@{}}
\toprule
Enc & Value & Arch & Description \& Worked Example \\
\midrule\endhead
\texttt{0} & \texttt{\imod{none}} & v11--v14 & Default behavior (flag disabled for \texttt{tex\_projection\_enable\_m}). \newline {\scriptsize\ttfamily \mnem{TEX\_DUAL}\imod{} \dst{r0}, \src{r1}} \\[2pt]
\texttt{1} & \texttt{\imod{project}} & v11--v14 & Divide texture coordinates by homogeneous $q$ (\texttt{textureProj}). \newline {\scriptsize\ttfamily \mnem{TEX\_DUAL}\imod{.project} \dst{r0}, \src{r1}} \\[2pt]
\bottomrule
\end{longtable}
\end{center}
\par\textbf{Worked Example:} \emph{Homogeneous projective divide (`textureProj`)}
\codebox{\mnem{TEX\_SINGLE}\imod{.tex2d.project.computed.slot0} \dst{@r0:r3}, \src{@r4:r6}, \src{u0}}
Divides \texttt{(u, v)} by homogeneous coordinate \texttt{q = @r6} inside the texture unit before filtering.

\subsection{Tex Skip (\texttt{tex\_skip\_m})}
\label{enum:tex_skip_m}
Used by \texttt{TEX\_DUAL} and related instructions (2 encoding values):

\begin{center}\footnotesize
\begin{longtable}{@{}rll>{\raggedright\arraybackslash}p{0.62\linewidth}@{}}
\toprule
Enc & Value & Arch & Description \& Worked Example \\
\midrule\endhead
\texttt{0} & \texttt{\imod{none}} & v9--v14 & Default behavior (flag disabled for \texttt{tex\_skip\_m}). \newline {\scriptsize\ttfamily \mnem{TEX\_DUAL}\imod{} \dst{r0}, \src{r1}} \\[2pt]
\texttt{1} & \texttt{\imod{skip}} & v9--v14 & Skip texture fetch for helper/discarded fragment lanes. \newline {\scriptsize\ttfamily \mnem{TEX\_DUAL}\imod{.skip} \dst{r0}, \src{r1}} \\[2pt]
\bottomrule
\end{longtable}
\end{center}
\par\textbf{Worked Example:} \emph{Skipping helper-lane memory traffic (`.skip`)}
\codebox{\mnem{TEX\_SINGLE}\imod{.tex2d.computed.skip.slot0} \dst{@r0:r3}, \src{@r4:r5}, \src{u0}}
Helper lanes participate in 2x2 quad LOD derivation, then \texttt{.skip} suppresses their actual cache fetch.

\subsection{Tex Texel Offset (\texttt{tex\_texel\_offset\_m})}
\label{enum:tex_texel_offset_m}
Used by \texttt{TEX\_DUAL} and related instructions (2 encoding values):

\begin{center}\footnotesize
\begin{longtable}{@{}rll>{\raggedright\arraybackslash}p{0.62\linewidth}@{}}
\toprule
Enc & Value & Arch & Description \& Worked Example \\
\midrule\endhead
\texttt{0} & \texttt{\imod{none}} & v9--v14 & Default behavior (flag disabled for \texttt{tex\_texel\_offset\_m}). \newline {\scriptsize\ttfamily \mnem{TEX\_DUAL}\imod{} \dst{r0}, \src{r1}} \\[2pt]
\texttt{1} & \texttt{\imod{texel\_offset}} & v9--v14 & Add signed 4-bit texel offset $(dx, dy, dz)$ (\texttt{textureOffset}). \newline {\scriptsize\ttfamily \mnem{TEX\_DUAL}\imod{.texel\_offset} \dst{r0}, \src{r1}} \\[2pt]
\bottomrule
\end{longtable}
\end{center}
\par\textbf{Worked Example:} \emph{Signed 4-bit texel offset (`textureOffset`)}
\codebox{\mnem{TEX\_SINGLE}\imod{.tex2d.texel\_offset.computed.slot0} \dst{@r0:r3}, \src{@r4:r6}, \src{u0}}
Applies packed signed 4-bit \texttt{(dx, dy)} offset in \texttt{@r6} (range \texttt{[-8, +7]}) before sampling.

\subsection{Tex Texel Offset Si (\texttt{tex\_texel\_offset\_si\_m})}
\label{enum:tex_texel_offset_si_m}
Used by \texttt{TEX\_FETCH} and related instructions (2 encoding values):

\begin{center}\footnotesize
\begin{longtable}{@{}rll>{\raggedright\arraybackslash}p{0.62\linewidth}@{}}
\toprule
Enc & Value & Arch & Description \& Worked Example \\
\midrule\endhead
\texttt{0} & \texttt{\imod{none}} & v9--v14 & Default behavior (flag disabled for \texttt{tex\_texel\_offset\_si\_m}). \newline {\scriptsize\ttfamily \mnem{TEX\_FETCH}\imod{} \dst{r0}, \src{r1}} \\[2pt]
\texttt{1} & \texttt{\imod{texel\_offset}} & v9--v14 & Add signed 4-bit texel offset $(dx, dy, dz)$ (\texttt{textureOffset}). \newline {\scriptsize\ttfamily \mnem{TEX\_FETCH}\imod{.texel\_offset} \dst{r0}, \src{r1}} \\[2pt]
\bottomrule
\end{longtable}
\end{center}
\par\textbf{Worked Example:} \emph{Immediate-style texel offset on `TEX\_FETCH`/`TEX\_GATHER`}
\codebox{\mnem{TEX\_GATHER}\imod{.tex2d.gather4\_r.texel\_offset.slot0} \dst{@r0:r3}, \src{@r4:r6}, \src{u0}}
Offsets the 2x2 gather footprint by signed \texttt{(dx, dy)} for GLSL \texttt{textureGatherOffset}.

\subsection{Tex Wide Indices Enable (\texttt{tex\_wide\_indices\_enable\_m})}
\label{enum:tex_wide_indices_enable_m}
Used by \texttt{TEX\_DUAL} and related instructions (2 encoding values):

\begin{center}\footnotesize
\begin{longtable}{@{}rll>{\raggedright\arraybackslash}p{0.62\linewidth}@{}}
\toprule
Enc & Value & Arch & Description \& Worked Example \\
\midrule\endhead
\texttt{0} & \texttt{\imod{none}} & v9--v14 & Default behavior (flag disabled for \texttt{tex\_wide\_indices\_enable\_m}). \newline {\scriptsize\ttfamily \mnem{TEX\_DUAL}\imod{} \dst{r0}, \src{r1}} \\[2pt]
\texttt{1} & \texttt{\imod{wide}} & v9--v14 & Use 32-bit wide texture/sampler descriptor indices. \newline {\scriptsize\ttfamily \mnem{TEX\_DUAL}\imod{.wide} \dst{r0}, \src{r1}} \\[2pt]
\bottomrule
\end{longtable}
\end{center}
\par\textbf{Worked Example:} \emph{32-bit bindless descriptor indices (`.wide`)}
\codebox{\mnem{TEX\_SINGLE}\imod{.tex2d.wide.computed.slot0} \dst{@r0:r3}, \src{@r4:r5}, \src{r6}}
Allows full 32-bit texture/sampler table indices in \texttt{r6} for large Vulkan bindless descriptor arrays.

\subsection{Tex Width (\texttt{tex\_width\_m})}
\label{enum:tex_width_m}
Used by \texttt{TEX\_FETCH} and related instructions (8 encoding values):

\begin{center}\footnotesize
\begin{longtable}{@{}rll>{\raggedright\arraybackslash}p{0.62\linewidth}@{}}
\toprule
Enc & Value & Arch & Description \& Worked Example \\
\midrule\endhead
\texttt{0} & \texttt{\imod{dst16}} & v9--v14 & texture return format \texttt{dst16} (16-bit packed staging register elements). \newline {\scriptsize\ttfamily \mnem{TEX\_FETCH}\imod{.dst16} \dst{@r0:r3}, \src{@r4:r5}} \\[2pt]
\texttt{0} & \texttt{\imod{f16}} & v9--v14 & texture return format \texttt{f16} (16-bit packed staging register elements). \newline {\scriptsize\ttfamily \mnem{TEX\_FETCH}\imod{.f16} \dst{@r0:r3}, \src{@r4:r5}} \\[2pt]
\texttt{0} & \texttt{\imod{s16}} & v9--v14 & texture return format \texttt{s16} (16-bit packed staging register elements). \newline {\scriptsize\ttfamily \mnem{TEX\_FETCH}\imod{.s16} \dst{@r0:r3}, \src{@r4:r5}} \\[2pt]
\texttt{0} & \texttt{\imod{u16}} & v9--v14 & texture return format \texttt{u16} (16-bit packed staging register elements). \newline {\scriptsize\ttfamily \mnem{TEX\_FETCH}\imod{.u16} \dst{@r0:r3}, \src{@r4:r5}} \\[2pt]
\texttt{1} & \texttt{\imod{dst32}} & v9--v14 & texture return format \texttt{dst32} (32-bit staging register elements). \newline {\scriptsize\ttfamily \mnem{TEX\_FETCH}\imod{.dst32} \dst{@r0:r3}, \src{@r4:r5}} \\[2pt]
\texttt{1} & \texttt{\imod{f32}} & v9--v14 & texture return format \texttt{f32} (32-bit staging register elements). \newline {\scriptsize\ttfamily \mnem{TEX\_FETCH}\imod{.f32} \dst{@r0:r3}, \src{@r4:r5}} \\[2pt]
\texttt{1} & \texttt{\imod{s32}} & v9--v14 & texture return format \texttt{s32} (32-bit staging register elements). \newline {\scriptsize\ttfamily \mnem{TEX\_FETCH}\imod{.s32} \dst{@r0:r3}, \src{@r4:r5}} \\[2pt]
\texttt{1} & \texttt{\imod{u32}} & v9--v14 & texture return format \texttt{u32} (32-bit staging register elements). \newline {\scriptsize\ttfamily \mnem{TEX\_FETCH}\imod{.u32} \dst{@r0:r3}, \src{@r4:r5}} \\[2pt]
\bottomrule
\end{longtable}
\end{center}
\par\textbf{Worked Example:} \emph{Signed integer (`s32`) vs FP16 (`f16`) texture return}
\codebox{\mnem{TEX\_SINGLE}\imod{.tex2d.f16.rgba.slot0} \dst{@r0:r1}, \src{@r4:r5}, \src{u0}}
\texttt{.f16} returns 4x FP16 in \texttt{@r0:r1}; \texttt{.s32}/\texttt{.u32} returns integer texels for \texttt{isampler2D}/\texttt{usampler2D}.

\subsection{Tex Write Mask (\texttt{tex\_write\_mask\_m})}
\label{enum:tex_write_mask_m}
Used by \texttt{TEX\_FETCH} and related instructions (15 encoding values):

\begin{center}\footnotesize
\begin{longtable}{@{}rll>{\raggedright\arraybackslash}p{0.62\linewidth}@{}}
\toprule
Enc & Value & Arch & Description \& Worked Example \\
\midrule\endhead
\texttt{1} & \texttt{\imod{r}} & v9--v14 & channel write mask: write only R component(s) to consecutive staging GPRs. \newline {\scriptsize\ttfamily \mnem{TEX\_FETCH}\imod{.r} \dst{@r0:r3}, \src{@r4:r5}} \\[2pt]
\texttt{2} & \texttt{\imod{g}} & v9--v14 & channel write mask: write only G component(s) to consecutive staging GPRs. \newline {\scriptsize\ttfamily \mnem{TEX\_FETCH}\imod{.g} \dst{@r0:r3}, \src{@r4:r5}} \\[2pt]
\texttt{3} & \texttt{\imod{rg}} & v9--v14 & channel write mask: write only RG component(s) to consecutive staging GPRs. \newline {\scriptsize\ttfamily \mnem{TEX\_FETCH}\imod{.rg} \dst{@r0:r3}, \src{@r4:r5}} \\[2pt]
\texttt{4} & \texttt{\imod{b}} & v9--v14 & channel write mask: write only B component(s) to consecutive staging GPRs. \newline {\scriptsize\ttfamily \mnem{TEX\_FETCH}\imod{.b} \dst{@r0:r3}, \src{@r4:r5}} \\[2pt]
\texttt{5} & \texttt{\imod{rb}} & v9--v14 & channel write mask: write only RB component(s) to consecutive staging GPRs. \newline {\scriptsize\ttfamily \mnem{TEX\_FETCH}\imod{.rb} \dst{@r0:r3}, \src{@r4:r5}} \\[2pt]
\texttt{6} & \texttt{\imod{gb}} & v9--v14 & channel write mask: write only GB component(s) to consecutive staging GPRs. \newline {\scriptsize\ttfamily \mnem{TEX\_FETCH}\imod{.gb} \dst{@r0:r3}, \src{@r4:r5}} \\[2pt]
\texttt{7} & \texttt{\imod{rgb}} & v9--v14 & channel write mask: write only RGB component(s) to consecutive staging GPRs. \newline {\scriptsize\ttfamily \mnem{TEX\_FETCH}\imod{.rgb} \dst{@r0:r3}, \src{@r4:r5}} \\[2pt]
\texttt{8} & \texttt{\imod{a}} & v9--v14 & channel write mask: write only A component(s) to consecutive staging GPRs. \newline {\scriptsize\ttfamily \mnem{TEX\_FETCH}\imod{.a} \dst{@r0:r3}, \src{@r4:r5}} \\[2pt]
\texttt{9} & \texttt{\imod{ra}} & v9--v14 & channel write mask: write only RA component(s) to consecutive staging GPRs. \newline {\scriptsize\ttfamily \mnem{TEX\_FETCH}\imod{.ra} \dst{@r0:r3}, \src{@r4:r5}} \\[2pt]
\texttt{10} & \texttt{\imod{ga}} & v9--v14 & channel write mask: write only GA component(s) to consecutive staging GPRs. \newline {\scriptsize\ttfamily \mnem{TEX\_FETCH}\imod{.ga} \dst{@r0:r3}, \src{@r4:r5}} \\[2pt]
\texttt{11} & \texttt{\imod{rga}} & v9--v14 & channel write mask: write only RGA component(s) to consecutive staging GPRs. \newline {\scriptsize\ttfamily \mnem{TEX\_FETCH}\imod{.rga} \dst{@r0:r3}, \src{@r4:r5}} \\[2pt]
\texttt{12} & \texttt{\imod{ba}} & v9--v14 & channel write mask: write only BA component(s) to consecutive staging GPRs. \newline {\scriptsize\ttfamily \mnem{TEX\_FETCH}\imod{.ba} \dst{@r0:r3}, \src{@r4:r5}} \\[2pt]
\texttt{13} & \texttt{\imod{rba}} & v9--v14 & channel write mask: write only RBA component(s) to consecutive staging GPRs. \newline {\scriptsize\ttfamily \mnem{TEX\_FETCH}\imod{.rba} \dst{@r0:r3}, \src{@r4:r5}} \\[2pt]
\texttt{14} & \texttt{\imod{gba}} & v9--v14 & channel write mask: write only GBA component(s) to consecutive staging GPRs. \newline {\scriptsize\ttfamily \mnem{TEX\_FETCH}\imod{.gba} \dst{@r0:r3}, \src{@r4:r5}} \\[2pt]
\texttt{15} & \texttt{\imod{rgba}} & v9--v14 & channel write mask: write only RGBA component(s) to consecutive staging GPRs. \newline {\scriptsize\ttfamily \mnem{TEX\_FETCH}\imod{.rgba} \dst{@r0:r3}, \src{@r4:r5}} \\[2pt]
\bottomrule
\end{longtable}
\end{center}
\par\textbf{Worked Example:} \emph{Single-channel `.r` mask saving staging GPRs}
\codebox{\mnem{TEX\_SINGLE}\imod{.tex2d.f32.r.slot0} \dst{@r0}, \src{@r4:r5}, \src{u0}}
Writes only the red channel to \texttt{@r0} instead of allocating 4 staging registers \texttt{@r0:r3}.

\subsection{Update Mode None (\texttt{update\_mode\_none\_m})}
\label{enum:update_mode_none_m}
Used by \texttt{LD\_VAR\_BUF} and related instructions (4 encoding values):

\begin{center}\footnotesize
\begin{longtable}{@{}rll>{\raggedright\arraybackslash}p{0.62\linewidth}@{}}
\toprule
Enc & Value & Arch & Description \& Worked Example \\
\midrule\endhead
\texttt{0} & \texttt{\imod{store}} & v9--v13 & Compute barycentric/varying coordinates and store in hardware varying cache. \newline {\scriptsize\ttfamily \mnem{LD\_VAR\_BUF}\imod{.store} \dst{@r0:r3}, \src{r4}} \\[2pt]
\texttt{1} & \texttt{\imod{retrieve}} & v9--v13 & Reuse previously cached barycentric coordinates (saves ALU/interpolation power). \newline {\scriptsize\ttfamily \mnem{LD\_VAR\_BUF}\imod{.retrieve} \dst{@r0:r3}, \src{r4}} \\[2pt]
\texttt{3} & \texttt{\imod{clobber}} & v9--v13 & Compute fresh coordinates and overwrite hardware varying cache. \newline {\scriptsize\ttfamily \mnem{LD\_VAR\_BUF}\imod{.clobber} \dst{@r0:r3}, \src{r4}} \\[2pt]
\texttt{4} & \texttt{\imod{none}} & v9--v13 & Do not read or modify the hardware varying cache. \newline {\scriptsize\ttfamily \mnem{LD\_VAR\_BUF}\imod{} \dst{@r0:r3}, \src{r4}} \\[2pt]
\bottomrule
\end{longtable}
\end{center}
\par\textbf{Worked Example:} \emph{Varying cache lifecycle (`store` -\textgreater{} `retrieve` -\textgreater{} `none`)}
\codebox{\mnem{LD\_VAR\_IMM}\imod{.f32.v4.center.store.slot0} \dst{@r0:r3}, \src{r59}, \imm{\#0}, \imm{\#0}}
\texttt{.store} populates the barycentric cache; \texttt{.retrieve} reuses it; \texttt{.none} leaves cache state untouched.

\subsection{Update Mode Special (\texttt{update\_mode\_special\_m})}
\label{enum:update_mode_special_m}
Used by \texttt{LD\_VAR\_SPECIAL} and related instructions (4 encoding values):

\begin{center}\footnotesize
\begin{longtable}{@{}rll>{\raggedright\arraybackslash}p{0.62\linewidth}@{}}
\toprule
Enc & Value & Arch & Description \& Worked Example \\
\midrule\endhead
\texttt{0} & \texttt{\imod{store}} & v14 & Compute barycentric/varying coordinates and store in hardware varying cache. \newline {\scriptsize\ttfamily \mnem{LD\_VAR\_SPECIAL}\imod{.store} \dst{@r0:r3}, \src{r4}} \\[2pt]
\texttt{1} & \texttt{\imod{retrieve}} & v14 & Reuse previously cached barycentric coordinates (saves ALU/interpolation power). \newline {\scriptsize\ttfamily \mnem{LD\_VAR\_SPECIAL}\imod{.retrieve} \dst{@r0:r3}, \src{r4}} \\[2pt]
\texttt{3} & \texttt{\imod{clobber}} & v9--v13 & Compute fresh coordinates and overwrite hardware varying cache. \newline {\scriptsize\ttfamily \mnem{LD\_VAR\_SPECIAL}\imod{.clobber} \dst{@r0:r3}, \src{r4}} \\[2pt]
\texttt{4} & \texttt{\imod{none}} & v9--v13 & Do not read or modify the hardware varying cache. \newline {\scriptsize\ttfamily \mnem{LD\_VAR\_SPECIAL}\imod{} \dst{@r0:r3}, \src{r4}} \\[2pt]
\bottomrule
\end{longtable}
\end{center}
\par\textbf{Worked Example:} \emph{Special varying cache update mode}
\codebox{\mnem{LD\_VAR\_SPECIAL}\imod{.frag\_z.center.clobber.slot0} \dst{@r0}, \src{r59}}
Uses \texttt{.clobber} when interpolating \texttt{gl\_FragCoord.z} independently of user varyings.

\subsection{V4U8 (\texttt{v4u8\_m})}
\label{enum:v4u8_m}
Used by \texttt{IDPADD} and related instructions (2 encoding values):

\begin{center}\footnotesize
\begin{longtable}{@{}rll>{\raggedright\arraybackslash}p{0.62\linewidth}@{}}
\toprule
Enc & Value & Arch & Description \& Worked Example \\
\midrule\endhead
\texttt{0} & \texttt{\imod{none}} & v14 & Default behavior (flag disabled for \texttt{v4u8\_m}). \newline {\scriptsize\ttfamily \mnem{IDPADD}\imod{} \dst{r0}, \src{r1}} \\[2pt]
\texttt{1} & \texttt{\imod{v4u8}} & v14 & Interpret 32-bit operands as 4 packed unsigned 8-bit lanes. \newline {\scriptsize\ttfamily \mnem{IDPADD}\imod{.v4u8} \dst{r0}, \src{r1}} \\[2pt]
\bottomrule
\end{longtable}
\end{center}
\par\textbf{Worked Example:} \emph{Unsigned 4x8-bit lane interpretation on `IDPADD`/`MMUL`}
\codebox{\mnem{IDPADD.v4s8}\imod{.v4u8.sat} \dst{r0}, \src{r1}, \src{r2}, \src{r3}}
Overrides a source operand from signed \texttt{s8} to unsigned \texttt{u8} for mixed-sign \texttt{OpSUDotKHR} dot products (v14).

\subsection{Var Tex Array Enable (\texttt{var\_tex\_array\_enable\_m})}
\label{enum:var_tex_array_enable_m}
Used by \texttt{VAR\_TEX\_BUF\_GATHER} and related instructions (2 encoding values):

\begin{center}\footnotesize
\begin{longtable}{@{}rll>{\raggedright\arraybackslash}p{0.62\linewidth}@{}}
\toprule
Enc & Value & Arch & Description \& Worked Example \\
\midrule\endhead
\texttt{0} & \texttt{\imod{none}} & v9--v14 & Default behavior (flag disabled for \texttt{var\_tex\_array\_enable\_m}). \newline {\scriptsize\ttfamily \mnem{VAR\_TEX\_BUF\_GATHER}\imod{} \dst{r0}, \src{r1}} \\[2pt]
\texttt{1} & \texttt{\imod{array}} & v9--v14 & Enable texture array layer index coordinate (\texttt{sampler2DArray}). \newline {\scriptsize\ttfamily \mnem{VAR\_TEX\_BUF\_GATHER}\imod{.array} \dst{r0}, \src{r1}} \\[2pt]
\bottomrule
\end{longtable}
\end{center}
\par\textbf{Worked Example:} \emph{Fused arrayed varying-texture sample}
\codebox{\mnem{VAR\_TEX\_SINGLE}\imod{.tex2d.array.computed.center.slot0} \dst{@r0:r3}, \src{u0}, \src{r59}, \src{u1}}
Interpolates \texttt{(u, v, layer)} directly from varying \texttt{u1} and samples a \texttt{sampler2DArray} in one instruction.

\subsection{Var Tex Buf Source Format (\texttt{var\_tex\_buf\_source\_format\_m})}
\label{enum:var_tex_buf_source_format_m}
Used by \texttt{VAR\_TEX\_BUF\_DUAL} and related instructions (2 encoding values):

\begin{center}\footnotesize
\begin{longtable}{@{}rll>{\raggedright\arraybackslash}p{0.62\linewidth}@{}}
\toprule
Enc & Value & Arch & Description \& Worked Example \\
\midrule\endhead
\texttt{0} & \texttt{\imod{src\_f32}} & v9 & Source buffer holds FP32 varying values. \newline {\scriptsize\ttfamily \mnem{VAR\_TEX\_BUF\_DUAL}\imod{.src\_f32} \dst{r0}, \src{r1}} \\[2pt]
\texttt{1} & \texttt{\imod{src\_f16}} & v10--v14 & Source buffer holds FP16 varying values. \newline {\scriptsize\ttfamily \mnem{VAR\_TEX\_BUF\_DUAL}\imod{.src\_f16} \dst{r0}, \src{r1}} \\[2pt]
\bottomrule
\end{longtable}
\end{center}
\par\textbf{Worked Example:} \emph{FP16 UV varying buffer fused into texture sample}
\codebox{\mnem{VAR\_TEX\_BUF\_SINGLE}\imod{.tex2d.src\_f16.computed.slot0} \dst{@r0:r3}, \src{u0}, \src{r59}, \src{r4}}
Reads \texttt{mediump} FP16 UVs from varying buffer offset \texttt{r4} and samples texture \texttt{u0} without intermediate GPRs.

\subsection{Var Tex Compare Enable (\texttt{var\_tex\_compare\_enable\_m})}
\label{enum:var_tex_compare_enable_m}
Used by \texttt{VAR\_TEX\_BUF\_DUAL} and related instructions (2 encoding values):

\begin{center}\footnotesize
\begin{longtable}{@{}rll>{\raggedright\arraybackslash}p{0.62\linewidth}@{}}
\toprule
Enc & Value & Arch & Description \& Worked Example \\
\midrule\endhead
\texttt{0} & \texttt{\imod{none}} & v9--v14 & Default behavior (flag disabled for \texttt{var\_tex\_compare\_enable\_m}). \newline {\scriptsize\ttfamily \mnem{VAR\_TEX\_BUF\_DUAL}\imod{} \dst{r0}, \src{r1}} \\[2pt]
\texttt{1} & \texttt{\imod{compare}} & v9--v14 & Enable shadow depth comparison reference coordinate (\texttt{sampler2DShadow}). \newline {\scriptsize\ttfamily \mnem{VAR\_TEX\_BUF\_DUAL}\imod{.compare} \dst{r0}, \src{r1}} \\[2pt]
\bottomrule
\end{longtable}
\end{center}
\par\textbf{Worked Example:} \emph{Fused varying interpolation + shadow comparison}
\codebox{\mnem{VAR\_TEX\_BUF\_SINGLE}\imod{.tex2d.compare.computed.slot0} \dst{@r0}, \src{u0}, \src{r59}, \src{r4}}
Interpolates shadow-map \texttt{(u, v, refZ)} from varying buffer and performs hardware PCF depth comparison.

\subsection{Var Tex Dimensionality (\texttt{var\_tex\_dimensionality\_m})}
\label{enum:var_tex_dimensionality_m}
Used by \texttt{VAR\_TEX\_BUF\_DUAL} and related instructions (4 encoding values):

\begin{center}\footnotesize
\begin{longtable}{@{}rll>{\raggedright\arraybackslash}p{0.62\linewidth}@{}}
\toprule
Enc & Value & Arch & Description \& Worked Example \\
\midrule\endhead
\texttt{0} & \texttt{\imod{tex1d}} & v9--v10 & 1D texture coordinate $(u)$. \newline {\scriptsize\ttfamily \mnem{VAR\_TEX\_BUF\_DUAL}\imod{.tex1d} \dst{@r0:r3}, \src{@r4:r5}} \\[2pt]
\texttt{1} & \texttt{\imod{tex2d}} & v9--v10 & 2D texture coordinates $(u, v)$. \newline {\scriptsize\ttfamily \mnem{VAR\_TEX\_BUF\_DUAL}\imod{.tex2d} \dst{@r0:r3}, \src{@r4:r5}} \\[2pt]
\texttt{2} & \texttt{\imod{tex3d}} & v9--v10 & 3D volumetric texture coordinates $(u, v, w)$. \newline {\scriptsize\ttfamily \mnem{VAR\_TEX\_BUF\_DUAL}\imod{.tex3d} \dst{@r0:r3}, \src{@r4:r5}} \\[2pt]
\texttt{3} & \texttt{\imod{cube}} & v11--v14 & Cubemap direction/face-projected coordinates. \newline {\scriptsize\ttfamily \mnem{VAR\_TEX\_BUF\_DUAL}\imod{.cube} \dst{@r0:r3}, \src{@r4:r5}} \\[2pt]
\bottomrule
\end{longtable}
\end{center}
\par\textbf{Worked Example:} \emph{Fused 3D varying-texture lookup (`.tex3d`)}
\codebox{\mnem{VAR\_TEX\_SINGLE}\imod{.tex3d.computed.center.slot0} \dst{@r0:r3}, \src{u0}, \src{r59}, \src{u1}}
Interpolates 3D \texttt{(u, v, w)} varying coordinates directly into a volumetric texture sample.

\subsection{Var Tex Dual Array Enable (\texttt{var\_tex\_dual\_array\_enable\_m})}
\label{enum:var_tex_dual_array_enable_m}
Used by \texttt{VAR\_TEX\_BUF\_DUAL} and related instructions (2 encoding values):

\begin{center}\footnotesize
\begin{longtable}{@{}rll>{\raggedright\arraybackslash}p{0.62\linewidth}@{}}
\toprule
Enc & Value & Arch & Description \& Worked Example \\
\midrule\endhead
\texttt{0} & \texttt{\imod{none}} & v9--v14 & Default behavior (flag disabled for \texttt{var\_tex\_dual\_array\_enable\_m}). \newline {\scriptsize\ttfamily \mnem{VAR\_TEX\_BUF\_DUAL}\imod{} \dst{r0}, \src{r1}} \\[2pt]
\texttt{1} & \texttt{\imod{array}} & v9--v14 & Enable texture array layer index coordinate (\texttt{sampler2DArray}). \newline {\scriptsize\ttfamily \mnem{VAR\_TEX\_BUF\_DUAL}\imod{.array} \dst{r0}, \src{r1}} \\[2pt]
\bottomrule
\end{longtable}
\end{center}
\par\textbf{Worked Example:} \emph{Fused arrayed dual texture sample}
\codebox{\mnem{VAR\_TEX\_BUF\_DUAL}\imod{.tex2d.array.computed.slot0} \dst{@r0:r7}, \src{u0}, \src{r59}, \src{r4}}
Interpolates \texttt{(u, v, layer)} once and feeds two paired \texttt{sampler2DArray} lookups.

\subsection{Var Tex Dual Lod Mode (\texttt{var\_tex\_dual\_lod\_mode\_m})}
\label{enum:var_tex_dual_lod_mode_m}
Used by \texttt{VAR\_TEX\_BUF\_DUAL} and related instructions (2 encoding values):

\begin{center}\footnotesize
\begin{longtable}{@{}rll>{\raggedright\arraybackslash}p{0.62\linewidth}@{}}
\toprule
Enc & Value & Arch & Description \& Worked Example \\
\midrule\endhead
\texttt{0} & \texttt{\imod{none}} & v9--v14 & Implicit zero LOD (base mip level 0, no derivatives). \newline {\scriptsize\ttfamily \mnem{VAR\_TEX\_BUF\_DUAL}\imod{} \dst{@r0:r3}, \src{@r4:r5}} \\[2pt]
\texttt{1} & \texttt{\imod{computed}} & v9--v14 & Hardware 2x2 quad finite-difference LOD (\texttt{texture()}). \newline {\scriptsize\ttfamily \mnem{VAR\_TEX\_BUF\_DUAL}\imod{.computed} \dst{@r0:r3}, \src{@r4:r5}} \\[2pt]
\bottomrule
\end{longtable}
\end{center}
\par\textbf{Worked Example:} \emph{Computed quad LOD vs base-mip (`none`) on `VAR\_TEX\_DUAL`}
\codebox{\mnem{VAR\_TEX\_DUAL}\imod{.tex2d.computed.center.slot0} \dst{@r0:r7}, \src{u0}, \src{r59}, \src{u1}}
\texttt{.computed} enables mipmapping via 2x2 quad derivatives; \texttt{none} forces mip level 0.

\subsection{Var Tex Lod Mode (\texttt{var\_tex\_lod\_mode\_m})}
\label{enum:var_tex_lod_mode_m}
Used by \texttt{VAR\_TEX\_BUF\_SINGLE} and related instructions (2 encoding values):

\begin{center}\footnotesize
\begin{longtable}{@{}rll>{\raggedright\arraybackslash}p{0.62\linewidth}@{}}
\toprule
Enc & Value & Arch & Description \& Worked Example \\
\midrule\endhead
\texttt{0} & \texttt{\imod{none}} & v9--v14 & Implicit zero LOD (base mip level 0, no derivatives). \newline {\scriptsize\ttfamily \mnem{VAR\_TEX\_BUF\_SINGLE}\imod{} \dst{@r0:r3}, \src{@r4:r5}} \\[2pt]
\texttt{1} & \texttt{\imod{computed}} & v9--v14 & Hardware 2x2 quad finite-difference LOD (\texttt{texture()}). \newline {\scriptsize\ttfamily \mnem{VAR\_TEX\_BUF\_SINGLE}\imod{.computed} \dst{@r0:r3}, \src{@r4:r5}} \\[2pt]
\bottomrule
\end{longtable}
\end{center}
\par\textbf{Worked Example:} \emph{Computed quad LOD vs mip-0 on `VAR\_TEX\_SINGLE`}
\codebox{\mnem{VAR\_TEX\_BUF\_SINGLE}\imod{.tex2d.computed.slot0} \dst{@r0:r3}, \src{u0}, \src{r59}, \src{r4}}
Uses \texttt{.computed} for standard \texttt{texture(s, vUV)} or \texttt{none} when sampling a single-level texture.

\subsection{Var Tex Projection Enable (\texttt{var\_tex\_projection\_enable\_m})}
\label{enum:var_tex_projection_enable_m}
Used by \texttt{VAR\_TEX\_BUF\_DUAL} and related instructions (2 encoding values):

\begin{center}\footnotesize
\begin{longtable}{@{}rll>{\raggedright\arraybackslash}p{0.62\linewidth}@{}}
\toprule
Enc & Value & Arch & Description \& Worked Example \\
\midrule\endhead
\texttt{0} & \texttt{\imod{none}} & v11--v14 & Default behavior (flag disabled for \texttt{var\_tex\_projection\_enable\_m}). \newline {\scriptsize\ttfamily \mnem{VAR\_TEX\_BUF\_DUAL}\imod{} \dst{r0}, \src{r1}} \\[2pt]
\texttt{1} & \texttt{\imod{project}} & v11--v14 & Divide texture coordinates by homogeneous $q$ (\texttt{textureProj}). \newline {\scriptsize\ttfamily \mnem{VAR\_TEX\_BUF\_DUAL}\imod{.project} \dst{r0}, \src{r1}} \\[2pt]
\bottomrule
\end{longtable}
\end{center}
\par\textbf{Worked Example:} \emph{Fused projective varying-texture sample}
\codebox{\mnem{VAR\_TEX\_SINGLE}\imod{.tex2d.project.computed.center.slot0} \dst{@r0:r3}, \src{u0}, \src{r59}, \src{u1}}
Interpolates homogeneous \texttt{(u, v, q)} and divides by \texttt{q} before filtering (\texttt{textureProj(s, vUVQ)}).

\subsection{Varying Name (\texttt{varying\_name\_m})}
\label{enum:varying_name_m}
Used by \texttt{LD\_VAR\_SPECIAL} and related instructions (4 encoding values):

\begin{center}\footnotesize
\begin{longtable}{@{}rll>{\raggedright\arraybackslash}p{0.62\linewidth}@{}}
\toprule
Enc & Value & Arch & Description \& Worked Example \\
\midrule\endhead
\texttt{0} & \texttt{\imod{point}} & v9--v13 & Read \texttt{gl\_PointCoord} $(s, t)$ inside point sprite. \newline {\scriptsize\ttfamily \mnem{LD\_VAR\_SPECIAL}\imod{.point} \dst{@r0:r1}, \src{r2}} \\[2pt]
\texttt{1} & \texttt{\imod{bary}} & v14 & Read perspective/noperspective barycentric weights $(i, j)$. \newline {\scriptsize\ttfamily \mnem{LD\_VAR\_SPECIAL}\imod{.bary} \dst{@r0:r1}, \src{r2}} \\[2pt]
\texttt{2} & \texttt{\imod{frag\_w}} & v9--v13 & Read reciprocal clip-space $1/w$ (\texttt{gl\_FragCoord.w}). \newline {\scriptsize\ttfamily \mnem{LD\_VAR\_SPECIAL}\imod{.frag\_w} \dst{@r0}, \src{r2}} \\[2pt]
\texttt{3} & \texttt{\imod{frag\_z}} & v9--v13 & Read interpolated window-space depth (\texttt{gl\_FragCoord.z}). \newline {\scriptsize\ttfamily \mnem{LD\_VAR\_SPECIAL}\imod{.frag\_z} \dst{@r0}, \src{r2}} \\[2pt]
\bottomrule
\end{longtable}
\end{center}
\par\textbf{Worked Example:} \emph{Reading fixed-function `gl\_PointCoord` vs `gl\_FragCoord.w`}
\codebox{\mnem{LD\_VAR\_SPECIAL}\imod{.point.center.clobber.slot0} \dst{@r0:r1}, \src{r59}}
\texttt{.point} reads point-sprite \texttt{(s, t)}, \texttt{.bary} reads barycentrics, \texttt{.frag\_w} reads \texttt{1/w}, and \texttt{.frag\_z} reads depth \texttt{z}.

\subsection{Vertex (\texttt{vertex\_m})}
\label{enum:vertex_m}
Used by \texttt{LD\_VAR\_BUF\_FLAT\_IMM} and related instructions (4 encoding values):

\begin{center}\footnotesize
\begin{longtable}{@{}rll>{\raggedright\arraybackslash}p{0.62\linewidth}@{}}
\toprule
Enc & Value & Arch & Description \& Worked Example \\
\midrule\endhead
\texttt{0} & \texttt{\imod{none}} & v14 & Default provoking vertex. \newline {\scriptsize\ttfamily \mnem{LD\_VAR\_BUF\_FLAT\_IMM}\imod{} \dst{@r0:r3}, \src{r4}} \\[2pt]
\texttt{1} & \texttt{\imod{vertex0}} & v14 & Read flat varying from triangle vertex 0. \newline {\scriptsize\ttfamily \mnem{LD\_VAR\_BUF\_FLAT\_IMM}\imod{.vertex0} \dst{@r0:r3}, \src{r4}} \\[2pt]
\texttt{2} & \texttt{\imod{vertex1}} & v14 & Read flat varying from triangle vertex 1. \newline {\scriptsize\ttfamily \mnem{LD\_VAR\_BUF\_FLAT\_IMM}\imod{.vertex1} \dst{@r0:r3}, \src{r4}} \\[2pt]
\texttt{3} & \texttt{\imod{vertex2}} & v14 & Read flat varying from triangle vertex 2. \newline {\scriptsize\ttfamily \mnem{LD\_VAR\_BUF\_FLAT\_IMM}\imod{.vertex2} \dst{@r0:r3}, \src{r4}} \\[2pt]
\bottomrule
\end{longtable}
\end{center}
\par\textbf{Worked Example:} \emph{Selecting triangle vertex 1 for `pervertexEXT`}
\codebox{\mnem{LD\_VAR\_FLAT\_IMM}\imod{.f32.v4.vertex1.slot0} \dst{@r0:r3}, \imm{\#0}, \imm{\#0}}
Reads vertex 1's uninterpolated attribute directly instead of the default provoking vertex (\texttt{none}/\texttt{vertex0}).

\subsection{Vertex Type (\texttt{vertex\_type\_m})}
\label{enum:vertex_type_m}
Used by \texttt{LD\_VAR\_BUF\_FLAT} and related instructions (2 encoding values):

\begin{center}\footnotesize
\begin{longtable}{@{}rll>{\raggedright\arraybackslash}p{0.62\linewidth}@{}}
\toprule
Enc & Value & Arch & Description \& Worked Example \\
\midrule\endhead
\texttt{0} & \texttt{\imod{none}} & v14 & Default behavior (flag disabled for \texttt{vertex\_type\_m}). \newline {\scriptsize\ttfamily \mnem{LD\_VAR\_BUF\_FLAT}\imod{} \dst{r0}, \src{r1}} \\[2pt]
\texttt{1} & \texttt{\imod{explicit\_vertex}} & v14 & Select flat varying vertex index explicitly via register. \newline {\scriptsize\ttfamily \mnem{LD\_VAR\_BUF\_FLAT}\imod{.explicit\_vertex} \dst{r0}, \src{r1}} \\[2pt]
\bottomrule
\end{longtable}
\end{center}
\par\textbf{Worked Example:} \emph{Dynamic register-indexed triangle vertex (v14)}
\codebox{\mnem{LD\_VAR\_BUF\_FLAT}\imod{.f32.v4.explicit\_vertex.slot0} \dst{@r0:r3}, \src{r4}}
With \texttt{.explicit\_vertex}, reads the triangle vertex index (\texttt{0..2}) dynamically from a register instead of fixing the provoking vertex.

\subsection{Z (\texttt{z\_m})}
\label{enum:z_m}
Used by \texttt{ZS\_EMIT} and related instructions (2 encoding values):
\par\textbf{CSF ABI Interaction:} Exports shader depth to \hyperref[sec:csf-zsd-blend]{\texttt{IDVSSR::ZSD}} (\texttt{r52:r53}) and requires \hyperref[sec:csf-tiler]{\texttt{DCD2::no\_shader\_depth\_read=0}} (Late-Z).
\begin{center}\footnotesize
\begin{longtable}{@{}rll>{\raggedright\arraybackslash}p{0.62\linewidth}@{}}
\toprule
Enc & Value & Arch & Description \& Worked Example \\
\midrule\endhead
\texttt{0} & \texttt{\imod{none}} & v9--v14 & Default behavior (flag disabled for \texttt{z\_m}). \newline {\scriptsize\ttfamily \mnem{ZS\_EMIT}\imod{} \dst{r0}, \src{r1}} \\[2pt]
\texttt{1} & \texttt{\imod{z}} & v9--v14 & Emit fragment depth ($z$) output to tile depth/stencil unit. \newline {\scriptsize\ttfamily \mnem{ZS\_EMIT}\imod{.z} \dst{r0}, \src{r1}} \\[2pt]
\bottomrule
\end{longtable}
\end{center}
\par\textbf{Worked Example:} \emph{Exporting programmatic fragment depth (`gl\_FragDepth`)}
\codebox{\mnem{ZS\_EMIT}\imod{.z.slot0} \dst{@r60}, \src{r0}, \src{k0}, \src{r60}}
With \texttt{.z}, submits \texttt{r0} as the fragment's depth value to the late-Z test and returns updated coverage in \texttt{@r60}.


% ============================================================================
% Chapter 3: Complete Instruction Reference (All 201 Mnemonics / 627 Specs)
% ============================================================================
\chapter{Instruction Reference}
This chapter documents all \textbf{201} canonical Valhall mnemonics (\textbf{627} \texttt{<instruction>} XML specifications across v9--v14). Each instruction starts on its own page and provides:
\begin{itemize}
  \item Operational semantics, structured bulleted notes, and color-coded syntax (\mnem{Mnemonic}, \imod{Modifiers}, \dst{Destination}, \src{Sources}, \imm{Immediates}).
  \item Closest SPIR-V opcode, syntax-aligned Equivalent GLSL operational semantics (with flag variants), bulleted Execution Context (cross-linked to Section~\ref{sec:csf-abi}), and multiple concrete worked examples.
  \item Variant/architecture/opcode table, representative 64-bit encoding bit-layout table (with clickable backlinks from modifier field names to their Chapter~2 definitions), and explicit explanations for every instruction field.
\end{itemize}
\clearpage

\section{Integer Arithmetic (16 mnemonics)}
\subsection*{BITREV}
\addcontentsline{toc}{subsection}{BITREV}
Reverses the bit order of the 32-bit unsigned integer \texttt{src0} and writes the result to \texttt{dest}: for every bit index $k \in \{0, \dots, 31\}$, sets \texttt{dest[k] = src0[31 - k]}.
\par\textbf{Notes:}
\begin{tightitemize}
  \item Exact bit-for-bit reversal (\texttt{hw\_tests.rs:1265-1272}).
  \item \textbf{Source references:} \texttt{kraid/ops.rs:331-355}, \texttt{encode\_v9.rs:811-828}, \texttt{nir.rs:1255-1261}, \texttt{hw\_tests.rs:1265-1272}.
\end{tightitemize}
\codebox{\mnem{BITREV}\imod{\{.flow\}} \dst{dest}\imod{\{.dest\_width\}}, \src{src0}}
\textbf{Operands:}
\begin{tightitemize}\footnotesize
  \item \texttt{\dst{dest}} (\texttt{dest\_register\_index}, \texttt{dest\_width}): 32-bit GPR receiving the bit-reversed word.
  \item \texttt{\src{src0}}: 32-bit input integer whose bit order is reversed.
\end{tightitemize}
\par\textbf{Closest SPIR-V Op:} \texttt{OpBitReverse}
\par\textbf{Pseudo-GLSL (Operational Semantics):}
\codebox{uint out\_bits = 0u;\\for (int k = 0; k \textless{} 32; ++k) out\_bits \textbar{}= ((src0 \textgreater{}\textgreater{} k) \& 1u) \textless{}\textless{} (31 - k);\\dest = out\_bits; // equivalent to bitfieldReverse(src0)}
\textbf{Execution Context:}
\begin{tightitemize}\footnotesize
  \item \textbf{Compute stage preloads (\hyperref[sec:csf-spd]{\texttt{SPD\_k}}, \texttt{r16..r23}):} Hardware pre-fills \texttt{r55..r57} with \texttt{WorkGroupID.xyz + JOB\_OFFSET} (from \hyperref[sec:csf-regs]{\texttt{COMPUTESR::JOB\_OFFSET\_X/Y/Z}} \texttt{r34..r36}) and \texttt{r60..r62} with \texttt{LocalInvocationID.xyz} or \texttt{GlobalInvocationID.xyz} (from \texttt{WG\_SIZE} \texttt{r33} \& \texttt{JOB\_SIZE} \texttt{r37..r39}, selected by \texttt{ShaderProgram::preload}) so compute shaders can derive thread/element indices without memory loads.
  \item \textbf{Vertex stage preloads (\hyperref[sec:csf-spd]{\texttt{SPD\_0}/\texttt{SPD\_1}}, \texttt{r16..r19}):} The hardware Index Fetcher reads \hyperref[sec:csf-regs]{\texttt{IDVSSR::INDEX\_BUFFER}} (\texttt{r54:r55}) and pre-fills \texttt{r59, r60} with \texttt{VertexIndex} (\texttt{+ VERTEX\_OFFSET + GLOBAL\_ATTR\_OFF}) and \texttt{r61} with \texttt{InstanceIndex} (\texttt{+ INSTANCE\_OFFSET}) so Position and Varying shaders can pull vertex attributes directly.
  \item \textbf{Inline FAU \& Special Operands (\hyperref[sec:csf-tsd-fau]{\texttt{FAU\_k}}, \texttt{r8..r15}):} Any 8-bit source with \texttt{mode=2} reads zero-latency uniform \texttt{u0..u31} from \texttt{FAU\_k} (\texttt{slots64 = bits[63:56]}); \texttt{mode=3} reads \texttt{SmallConstant} (Table~1.2) or \texttt{FauSpecial} (\texttt{pc}, \texttt{lane\_id}, \texttt{lane\_width}, \texttt{tls\_ptr}, \texttt{wls\_ptr}, \texttt{blend\_descriptor\_0..7}).
\end{tightitemize}
\par\textbf{Worked Examples:}
\par\textbf{Example 1 (Canonical 32-bit reversal):}
\codebox{\mnem{BITREV.i32} \dst{r0}, \src{r1}}
With \texttt{r1 = 0x12345678} (\texttt{0001\_0010\_0011\_0100\_0101\_0110\_0111\_1000}): bit \texttt{k} moves to bit \texttt{31 - k}, producing \texttt{r0 = 0x1E6A2C48}.
\par\textbf{Example 2 (Palindrome \& MSB/LSB swap):}
\codebox{\mnem{BITREV.i32} \dst{r0}, \src{r1}}
With \texttt{r1 = 0x80000001}: bit 31 and bit 0 swap cleanly -\textgreater{} \texttt{r0 = 0x80000001}; for \texttt{r1 = 0x00000003} -\textgreater{} \texttt{r0 = 0xC0000000}.

\begin{center}\footnotesize
\begin{tabularx}{\linewidth}{@{}lllX@{}}
\toprule
Variant & Arch & Class & Opcode Fields \\
\midrule
\texttt{(none)} & v9--v14 & \texttt{R1} & \texttt{opcode2=BITREV, opcode1=MOV} \\
\bottomrule
\end{tabularx}
\end{center}
\begin{center}\footnotesize
\begin{tabularx}{0.88\linewidth}{@{}l l X@{}}
\toprule
Field Name & Bits & Type / Fixed Value \\
\midrule
\texttt{src0} & \texttt{0:7} & \texttt{SourceEncoding} \\
\textit{Reserved} & \texttt{8:15} & \texttt{zero} \\
\texttt{opcode2} & \texttt{16:20} & \texttt{exact=BITREV} \\
\textit{Reserved} & \texttt{21:39} & \texttt{zero} \\
\texttt{dest\_register\_index} & \texttt{40:45} & \texttt{uint} \\
\hyperref[enum:dest_width_m]{\texttt{dest\_width}} & \texttt{46:47} & \texttt{dest\_width\_m} \\
\texttt{opcode1} & \texttt{48:51} & \texttt{exact=MOV} \\
\texttt{instruction\_class} & \texttt{52:56} & \texttt{exact=R1} \\
\texttt{fau\_page\_index} & \texttt{57:58} & \texttt{uint} \\
\hyperref[enum:flow_control_m]{\texttt{flow}} & \texttt{59:62} & \texttt{flow\_control\_m} \\
\textit{Reserved} & \texttt{63} & \texttt{zero} \\
\bottomrule
\end{tabularx}
\end{center}
\textbf{Fields Explained:}
\begin{tightitemize}\footnotesize
  \item \texttt{src0}: Primary source operand (\texttt{r0..r63}, \texttt{u0..u63} uniform, \texttt{k*} FAU constant, or \texttt{t0..t2} temporary; \texttt{\textasciicircum{}} marks discard/last-use). \emph{Example:} {\scriptsize\ttfamily \mnem{FADD.f32} \dst{r0}, \src{r1}, \src{r2}} (\textbf{SPIR-V:} \texttt{OpFAdd (Operand 0)}; \textbf{GLSL:} \texttt{float c = a + b;})
  \item \texttt{opcode2}: Secondary sub-opcode further specializing instructions sharing \texttt{instruction\_class} and \texttt{opcode1} (e.g., conversions, SFU tables, moves). \emph{Example:} {\scriptsize\ttfamily \mnem{F32\_TO\_S32}\imod{.round\_zero} \dst{r0}, \src{r1}} (\textbf{SPIR-V:} \texttt{OpConvertFToS}; \textbf{GLSL:} \texttt{int y = int(x);})
  \item \texttt{dest\_register\_index}: 6-bit destination GPR index (\texttt{r0..r63}) receiving the Stage C writeback result (even-aligned pair \texttt{[rN:rN+1]} for 64-bit ops). \emph{Example:} {\scriptsize\ttfamily \mnem{FMA.f32} \dst{r0}, \src{r1}, \src{r2}, \src{r3}} (\textbf{SPIR-V:} \texttt{SSA Result ID}; \textbf{GLSL:} \texttt{float r0 = fma(r1, r2, r3);})
  \item \hyperref[enum:dest_width_m]{\texttt{dest\_width}}: Stage C destination lane write-mask (\texttt{none} = full 32-bit, \texttt{h0} = low 16-bit, \texttt{h1} = high 16-bit without bit-shifting). \emph{Example:} {\scriptsize\ttfamily \mnem{FADD.v2f16} \dst{r0}\imod{.h1}, \src{r1}, \src{r2}} (\textbf{SPIR-V:} \texttt{OpVectorInsertDynamic / WriteMask}; \textbf{GLSL:} \texttt{dst.y = a.y + b.y;})
  \item \texttt{opcode1}: Primary sub-opcode within the instruction class uniquely identifying the operation family. \emph{Example:} {\scriptsize\ttfamily \mnem{IADD.s32} \dst{r0}, \src{r1}, \src{r2}} (\textbf{SPIR-V:} \texttt{OpIAdd (operation selector)}; \textbf{GLSL:} \texttt{// Primary opcode bits})
  \item \texttt{instruction\_class}: 4-bit major instruction class selector (\texttt{bits 48:51}) partitioning the 64-bit Valhall opcode space (e.g., \texttt{FMA}, \texttt{CVT}, \texttt{SFU}, \texttt{LS}, \texttt{TEX}, \texttt{BRANCH}). \emph{Example:} {\scriptsize\ttfamily \mnem{FADD.f32} \dst{r0}, \src{r1}, \src{r2}} (\textbf{SPIR-V:} \texttt{N/A (encoding class)}; \textbf{GLSL:} \texttt{// Hardware execution pipe / decode class})
  \item \texttt{fau\_page\_index}: Selects the FAU (Fast Access Uniform) page (\texttt{0..3}) when a source operand references hardware small constants, thread IDs, or lane constants. \emph{Example:} {\scriptsize\ttfamily \mnem{FADD.f32} \dst{r0}, \src{r1}, \src{k1f}} (\textbf{SPIR-V:} \texttt{OpConstant / BuiltIn}; \textbf{GLSL:} \texttt{float y = x + 1.0;})
  \item \hyperref[enum:flow_control_m]{\texttt{flow}}: 4-bit embedded flow-control \& scoreboard wait modifier (\texttt{bits 59:62}: \texttt{none}, \texttt{wait0..wait0126}, \texttt{wait}, \texttt{reconverge}, \texttt{discard}, \texttt{end}). \emph{Example:} {\scriptsize\ttfamily \mnem{FADD.f32}\imod{.wait0.reconverge} \dst{r0}, \src{r1}, \src{r2}} (\textbf{SPIR-V:} \texttt{OpControlBarrier / OpReturn}; \textbf{GLSL:} \texttt{// Scoreboard wait + warp reconvergence})
\end{tightitemize}
\clearpage

\subsection*{CLZ}
\addcontentsline{toc}{subsection}{CLZ}
Counts leading zero bits independently in each SIMD element $i \in \{0, \dots, N-1\}$ of \texttt{src0} (where $N=1$ for \texttt{u32}, $N=2$ for \texttt{v2u16}, and $N=4$ for \texttt{v4u8} on v9--v10, with lane width $W = 32/N$). When \texttt{src0[i] != 0}, writes \texttt{leading\_zeros\_W(src0[i])} to \texttt{dest[i]}; when \texttt{src0[i] == 0}, writes \texttt{W} by default or all-ones (\texttt{\textasciitilde{}0\_W}) if \texttt{.mask} is set.
\par\textbf{Notes:}
\begin{tightitemize}
  \item When \texttt{src0[i] == 0}: returns \texttt{bits} (32/16/8) if \texttt{!mask}, or \texttt{0xFFFFFFFF}/\texttt{0xFFFF}/\texttt{0xFF} if \texttt{mask} (\texttt{hw\_tests.rs:1279-1282}).
  \item \textbf{Source references:} \texttt{kraid/ops.rs:696-732}, \texttt{encode\_v9.rs:902-920}, \texttt{nir.rs:1240-1279}, \texttt{hw\_tests.rs:1275-1315}.
\end{tightitemize}
\codebox{\mnem{CLZ}\imod{\{.flow\}}\imod{\{.mask\}} \dst{dest}\imod{\{.dest\_width\}}, \src{src0}\\\mnem{CLZ}\imod{\{.flow\}}\imod{\{.mask\}} \dst{dest}\imod{\{.dest\_width\}}, \src{src0}\imod{\{.swz0\}}}
\textbf{Operands:}
\begin{tightitemize}\footnotesize
  \item \texttt{\dst{dest}} (\texttt{dest\_register\_index}, \texttt{dest\_width}): Destination GPR receiving per-lane leading-zero counts (\texttt{W} or \texttt{\textasciitilde{}0} on zero if \texttt{.mask}).
  \item \texttt{\src{src0}}: Input integer (\texttt{u32}, \texttt{v2u16}, or \texttt{v4u8}).
\end{tightitemize}
\par\textbf{Closest SPIR-V Op:} \texttt{GLSL.std.450 FindUMsb (31 - clz)}
\par\textbf{Pseudo-GLSL (Operational Semantics):}
\codebox{for (int i = 0; i \textless{} N; ++i) \{\\    if (src0[i] == 0u) dest[i] = flag\_mask ? \textasciitilde{}0u : uint(W);\\    else               dest[i] = uint(W - 1 - findMSB(src0[i]));\\\}}
\textbf{Execution Context:}
\begin{tightitemize}\footnotesize
  \item \textbf{Compute stage preloads (\hyperref[sec:csf-spd]{\texttt{SPD\_k}}, \texttt{r16..r23}):} Hardware pre-fills \texttt{r55..r57} with \texttt{WorkGroupID.xyz + JOB\_OFFSET} (from \hyperref[sec:csf-regs]{\texttt{COMPUTESR::JOB\_OFFSET\_X/Y/Z}} \texttt{r34..r36}) and \texttt{r60..r62} with \texttt{LocalInvocationID.xyz} or \texttt{GlobalInvocationID.xyz} (from \texttt{WG\_SIZE} \texttt{r33} \& \texttt{JOB\_SIZE} \texttt{r37..r39}, selected by \texttt{ShaderProgram::preload}) so compute shaders can derive thread/element indices without memory loads.
  \item \textbf{Vertex stage preloads (\hyperref[sec:csf-spd]{\texttt{SPD\_0}/\texttt{SPD\_1}}, \texttt{r16..r19}):} The hardware Index Fetcher reads \hyperref[sec:csf-regs]{\texttt{IDVSSR::INDEX\_BUFFER}} (\texttt{r54:r55}) and pre-fills \texttt{r59, r60} with \texttt{VertexIndex} (\texttt{+ VERTEX\_OFFSET + GLOBAL\_ATTR\_OFF}) and \texttt{r61} with \texttt{InstanceIndex} (\texttt{+ INSTANCE\_OFFSET}) so Position and Varying shaders can pull vertex attributes directly.
  \item \textbf{Inline FAU \& Special Operands (\hyperref[sec:csf-tsd-fau]{\texttt{FAU\_k}}, \texttt{r8..r15}):} Any 8-bit source with \texttt{mode=2} reads zero-latency uniform \texttt{u0..u31} from \texttt{FAU\_k} (\texttt{slots64 = bits[63:56]}); \texttt{mode=3} reads \texttt{SmallConstant} (Table~1.2) or \texttt{FauSpecial} (\texttt{pc}, \texttt{lane\_id}, \texttt{lane\_width}, \texttt{tls\_ptr}, \texttt{wls\_ptr}, \texttt{blend\_descriptor\_0..7}).
\end{tightitemize}
\par\textbf{Worked Examples:}
\par\textbf{Example 1 (Non-zero input):}
\codebox{\mnem{CLZ.u32} \dst{r0}, \src{r1}}
With \texttt{r1 = 0x00100000} (bit 20 set): 11 zero bits precede bit 20 -\textgreater{} \texttt{r0 = 11}. GLSL \texttt{findMSB(r1)} is then \texttt{31 - 11 = 20}.
\par\textbf{Example 2 (Zero input with .mask vs default):}
\codebox{\mnem{CLZ.u32}\imod{.mask} \dst{r0}, \src{r1}}
With \texttt{r1 = 0x00000000}: without \texttt{.mask}, hardware returns \texttt{r0 = 32} (lane width); with \texttt{.mask}, hardware returns \texttt{r0 = 0xFFFFFFFF} (\texttt{-1}, matching GLSL \texttt{findMSB(0) == -1}).

\begin{center}\footnotesize
\begin{tabularx}{\linewidth}{@{}lllX@{}}
\toprule
Variant & Arch & Class & Opcode Fields \\
\midrule
\texttt{(none)} & v9--v14 & \texttt{R1} & \texttt{opcode2=CLZ\_U32, opcode1=MOV} \\
\texttt{(none)} & v9--v14 & \texttt{R1} & \texttt{opcode2=CLZ\_V2U16, opcode1=MOV} \\
\texttt{(none)} & v9--v14 & \texttt{R1} & \texttt{opcode2=CLZ\_V4U8, opcode1=MOV} \\
\bottomrule
\end{tabularx}
\end{center}
\begin{center}\footnotesize
\begin{tabularx}{0.88\linewidth}{@{}l l X@{}}
\toprule
Field Name & Bits & Type / Fixed Value \\
\midrule
\texttt{src0} & \texttt{0:7} & \texttt{SourceEncoding} \\
\textit{Reserved} & \texttt{8:15} & \texttt{zero} \\
\texttt{opcode2} & \texttt{16:20} & \texttt{exact=CLZ\_U32} \\
\textit{Reserved} & \texttt{21:29} & \texttt{zero} \\
\hyperref[enum:clz_mask_m]{\texttt{mask}} & \texttt{30} & \texttt{clz\_mask\_m} \\
\textit{Reserved} & \texttt{31:39} & \texttt{zero} \\
\texttt{dest\_register\_index} & \texttt{40:45} & \texttt{uint} \\
\hyperref[enum:dest_width_m]{\texttt{dest\_width}} & \texttt{46:47} & \texttt{dest\_width\_m} \\
\texttt{opcode1} & \texttt{48:51} & \texttt{exact=MOV} \\
\texttt{instruction\_class} & \texttt{52:56} & \texttt{exact=R1} \\
\texttt{fau\_page\_index} & \texttt{57:58} & \texttt{uint} \\
\hyperref[enum:flow_control_m]{\texttt{flow}} & \texttt{59:62} & \texttt{flow\_control\_m} \\
\textit{Reserved} & \texttt{63} & \texttt{zero} \\
\bottomrule
\end{tabularx}
\end{center}
\textbf{Fields Explained:}
\begin{tightitemize}\footnotesize
  \item \texttt{src0}: Primary source operand (\texttt{r0..r63}, \texttt{u0..u63} uniform, \texttt{k*} FAU constant, or \texttt{t0..t2} temporary; \texttt{\textasciicircum{}} marks discard/last-use). \emph{Example:} {\scriptsize\ttfamily \mnem{FADD.f32} \dst{r0}, \src{r1}, \src{r2}} (\textbf{SPIR-V:} \texttt{OpFAdd (Operand 0)}; \textbf{GLSL:} \texttt{float c = a + b;})
  \item \texttt{opcode2}: Secondary sub-opcode further specializing instructions sharing \texttt{instruction\_class} and \texttt{opcode1} (e.g., conversions, SFU tables, moves). \emph{Example:} {\scriptsize\ttfamily \mnem{F32\_TO\_S32}\imod{.round\_zero} \dst{r0}, \src{r1}} (\textbf{SPIR-V:} \texttt{OpConvertFToS}; \textbf{GLSL:} \texttt{int y = int(x);})
  \item \hyperref[enum:clz_mask_m]{\texttt{mask}}: Enables bitmask-mode leading-zero count in \texttt{CLZ}. \emph{Example:} {\scriptsize\ttfamily \mnem{CLZ.u32}\imod{.mask} \dst{r0}, \src{r1}} (\textbf{SPIR-V:} \texttt{GLSL.std.450 FindUMsb}; \textbf{GLSL:} \texttt{int msb = findMSB(x);})
  \item \texttt{dest\_register\_index}: 6-bit destination GPR index (\texttt{r0..r63}) receiving the Stage C writeback result (even-aligned pair \texttt{[rN:rN+1]} for 64-bit ops). \emph{Example:} {\scriptsize\ttfamily \mnem{FMA.f32} \dst{r0}, \src{r1}, \src{r2}, \src{r3}} (\textbf{SPIR-V:} \texttt{SSA Result ID}; \textbf{GLSL:} \texttt{float r0 = fma(r1, r2, r3);})
  \item \hyperref[enum:dest_width_m]{\texttt{dest\_width}}: Stage C destination lane write-mask (\texttt{none} = full 32-bit, \texttt{h0} = low 16-bit, \texttt{h1} = high 16-bit without bit-shifting). \emph{Example:} {\scriptsize\ttfamily \mnem{FADD.v2f16} \dst{r0}\imod{.h1}, \src{r1}, \src{r2}} (\textbf{SPIR-V:} \texttt{OpVectorInsertDynamic / WriteMask}; \textbf{GLSL:} \texttt{dst.y = a.y + b.y;})
  \item \texttt{opcode1}: Primary sub-opcode within the instruction class uniquely identifying the operation family. \emph{Example:} {\scriptsize\ttfamily \mnem{IADD.s32} \dst{r0}, \src{r1}, \src{r2}} (\textbf{SPIR-V:} \texttt{OpIAdd (operation selector)}; \textbf{GLSL:} \texttt{// Primary opcode bits})
  \item \texttt{instruction\_class}: 4-bit major instruction class selector (\texttt{bits 48:51}) partitioning the 64-bit Valhall opcode space (e.g., \texttt{FMA}, \texttt{CVT}, \texttt{SFU}, \texttt{LS}, \texttt{TEX}, \texttt{BRANCH}). \emph{Example:} {\scriptsize\ttfamily \mnem{FADD.f32} \dst{r0}, \src{r1}, \src{r2}} (\textbf{SPIR-V:} \texttt{N/A (encoding class)}; \textbf{GLSL:} \texttt{// Hardware execution pipe / decode class})
  \item \texttt{fau\_page\_index}: Selects the FAU (Fast Access Uniform) page (\texttt{0..3}) when a source operand references hardware small constants, thread IDs, or lane constants. \emph{Example:} {\scriptsize\ttfamily \mnem{FADD.f32} \dst{r0}, \src{r1}, \src{k1f}} (\textbf{SPIR-V:} \texttt{OpConstant / BuiltIn}; \textbf{GLSL:} \texttt{float y = x + 1.0;})
  \item \hyperref[enum:flow_control_m]{\texttt{flow}}: 4-bit embedded flow-control \& scoreboard wait modifier (\texttt{bits 59:62}: \texttt{none}, \texttt{wait0..wait0126}, \texttt{wait}, \texttt{reconverge}, \texttt{discard}, \texttt{end}). \emph{Example:} {\scriptsize\ttfamily \mnem{FADD.f32}\imod{.wait0.reconverge} \dst{r0}, \src{r1}, \src{r2}} (\textbf{SPIR-V:} \texttt{OpControlBarrier / OpReturn}; \textbf{GLSL:} \texttt{// Scoreboard wait + warp reconvergence})
\end{tightitemize}
\clearpage

\subsection*{HADD}
\addcontentsline{toc}{subsection}{HADD}
Computes the overflow-free halving average independently for each SIMD element $i \in \{0, \dots, N-1\}$ (where $N=1$ for \texttt{s32}/\texttt{u32}, $N=2$ for \texttt{v2s16}/\texttt{v2u16}, $N=4$ for \texttt{v4s8}/\texttt{v4u8}; arch \textless{}= v10): \texttt{dest[i] = (ext(src0[i]) + ext(src1[i]) + (round == round\_up ? 1 : 0)) \textgreater{}\textgreater{} 1} using an exact $(W+1)$-bit intermediate sum.
\par\textbf{Notes:}
\begin{tightitemize}
  \item Removed in v11 (\texttt{hw\_tests.rs:1768}).
  \item Arithmetic shift right for signed, logical shift right for unsigned.
  \item \textbf{Source references:} \texttt{kraid/ops.rs:2157-2203}, \texttt{encode\_v9.rs:1710-1734}, \texttt{nir.rs:1071-1090}, \texttt{hw\_tests.rs:1765-1812}.
\end{tightitemize}
\codebox{\mnem{HADD}\imod{\{.flow\}}\imod{\{.round\}} \dst{dest}\imod{\{.dest\_width\}}, \src{src0}\imod{\{.widen0\}}, \src{src1}\imod{\{.widen1\}}\\\mnem{HADD}\imod{\{.flow\}}\imod{\{.round\}} \dst{dest}\imod{\{.dest\_width\}}, \src{src0}\imod{\{.swz0\}}, \src{src1}\imod{\{.swz1\}}\\\mnem{HADD}\imod{\{.flow\}}\imod{\{.round\}} \dst{dest}\imod{\{.dest\_width\}}, \src{src0}\imod{\{.lanes0\}}, \src{src1}\imod{\{.lanes1\}}}
\textbf{Operands:}
\begin{tightitemize}\footnotesize
  \item \texttt{\dst{dest}} (\texttt{dest\_register\_index}, \texttt{dest\_width}): Destination GPR receiving per-lane halving averages \texttt{(src0 + src1 + round) \textgreater{}\textgreater{} 1}.
  \item \texttt{\src{src0}}: First integer operand.
  \item \texttt{\src{src1}}: Second integer operand.
\end{tightitemize}
\par\textbf{Closest SPIR-V Op:} \texttt{None (lowered from `(a+b)\textgreater{}\textgreater{}1` / `findMSB(int)`)}
\par\textbf{Pseudo-GLSL (Operational Semantics):}
\codebox{for (int i = 0; i \textless{} N; ++i) \{\\    int64\_t sum = ext(src0[i]) + ext(src1[i]) + (round == round\_up ? 1 : 0);\\    dest[i] = T(sum \textgreater{}\textgreater{} 1); // exact (W+1)-bit sum, no overflow\\\}}
\textbf{Execution Context:}
\begin{tightitemize}\footnotesize
  \item \textbf{Compute stage preloads (\hyperref[sec:csf-spd]{\texttt{SPD\_k}}, \texttt{r16..r23}):} Hardware pre-fills \texttt{r55..r57} with \texttt{WorkGroupID.xyz + JOB\_OFFSET} (from \hyperref[sec:csf-regs]{\texttt{COMPUTESR::JOB\_OFFSET\_X/Y/Z}} \texttt{r34..r36}) and \texttt{r60..r62} with \texttt{LocalInvocationID.xyz} or \texttt{GlobalInvocationID.xyz} (from \texttt{WG\_SIZE} \texttt{r33} \& \texttt{JOB\_SIZE} \texttt{r37..r39}, selected by \texttt{ShaderProgram::preload}) so compute shaders can derive thread/element indices without memory loads.
  \item \textbf{Vertex stage preloads (\hyperref[sec:csf-spd]{\texttt{SPD\_0}/\texttt{SPD\_1}}, \texttt{r16..r19}):} The hardware Index Fetcher reads \hyperref[sec:csf-regs]{\texttt{IDVSSR::INDEX\_BUFFER}} (\texttt{r54:r55}) and pre-fills \texttt{r59, r60} with \texttt{VertexIndex} (\texttt{+ VERTEX\_OFFSET + GLOBAL\_ATTR\_OFF}) and \texttt{r61} with \texttt{InstanceIndex} (\texttt{+ INSTANCE\_OFFSET}) so Position and Varying shaders can pull vertex attributes directly.
  \item \textbf{Inline FAU \& Special Operands (\hyperref[sec:csf-tsd-fau]{\texttt{FAU\_k}}, \texttt{r8..r15}):} Any 8-bit source with \texttt{mode=2} reads zero-latency uniform \texttt{u0..u31} from \texttt{FAU\_k} (\texttt{slots64 = bits[63:56]}); \texttt{mode=3} reads \texttt{SmallConstant} (Table~1.2) or \texttt{FauSpecial} (\texttt{pc}, \texttt{lane\_id}, \texttt{lane\_width}, \texttt{tls\_ptr}, \texttt{wls\_ptr}, \texttt{blend\_descriptor\_0..7}).
\end{tightitemize}
\par\textbf{Worked Examples:}
\par\textbf{Example 1 (Overflow-free unsigned ceil average):}
\codebox{\mnem{HADD.u32}\imod{.round\_up} \dst{r0}, \src{r1}, \src{r2}}
With \texttt{r1 = 0xFFFFFFFF, r2 = 0xFFFFFFFF}: exact 33-bit sum \texttt{(2\textasciicircum{}32-1 + 2\textasciicircum{}32-1 + 1) \textgreater{}\textgreater{} 1 = 0xFFFFFFFF} with zero overflow.
\par\textbf{Example 2 (Signed floor vs ceil rounding):}
\codebox{\mnem{HADD.s32}\imod{.round\_down} \dst{r0}, \src{r1}, \src{r2}}
With \texttt{r1 = 1, r2 = 2}: \texttt{.round\_down} computes \texttt{(1 + 2) \textgreater{}\textgreater{} 1 = 1}, whereas \texttt{.round\_up} computes \texttt{(1 + 2 + 1) \textgreater{}\textgreater{} 1 = 2} (v9--v10 only).
\par\textbf{Mined SPIR-V $\to$ Valhall ISA Example (Halving signed vertex-index divide-by-2 --- Vertex shader, Mali-G710 (v10)):}
\codebox{\textbf{; SPIR-V (input)}\\\%v = OpLoad \%int \%gl\_VertexIndex\\\%half = OpShiftRightArithmetic \%int \%v \%int\_1\\\%f = OpConvertSToF \%float \%half\\[3pt]\textbf{; Valhall ISA (malioc output)}\\HADD.s32.round\_down r0, r60\textasciicircum{}, k0\\S32\_TO\_F32 r0, r0\textasciicircum{}}
Instead of emitting a shift-unit \texttt{ARSHIFT} instruction, \texttt{malioc} lowers \texttt{(gl\_VertexIndex \textgreater{}\textgreater{} 1)} on v10 to \texttt{HADD.s32.round\_down r0, r60\textasciicircum{}, k0}, computing \texttt{(r60 + 0) \textgreater{}\textgreater{} 1} on the integer add pipe with zero extra constant setup (since \texttt{k0} is the inline zero constant and \texttt{r60} is the preloaded \texttt{VertexIndex}).
\begin{center}\footnotesize
\begin{tabularx}{\linewidth}{@{}lllX@{}}
\toprule
Variant & Arch & Class & Opcode Fields \\
\midrule
\texttt{s32} & v9--v14 & \texttt{R2} & \texttt{opcode2=HADD} \\
\texttt{u32} & v9--v14 & \texttt{R2} & \texttt{opcode2=HADD} \\
\texttt{v2s16} & v9--v14 & \texttt{R2} & \texttt{opcode2=HADD} \\
\texttt{v2u16} & v9--v14 & \texttt{R2} & \texttt{opcode2=HADD} \\
\texttt{v4s8} & v9--v14 & \texttt{R2} & \texttt{opcode2=HADD} \\
\texttt{v4u8} & v9--v14 & \texttt{R2} & \texttt{opcode2=HADD} \\
\bottomrule
\end{tabularx}
\end{center}
\begin{center}\footnotesize
\begin{tabularx}{0.88\linewidth}{@{}l l X@{}}
\toprule
Field Name & Bits & Type / Fixed Value \\
\midrule
\texttt{src0} & \texttt{0:7} & \texttt{SourceEncoding} \\
\texttt{src1} & \texttt{8:15} & \texttt{SourceEncoding} \\
\texttt{opcode2} & \texttt{16:19} & \texttt{exact=HADD} \\
\textit{Reserved} & \texttt{20:25} & \texttt{zero} \\
\hyperref[enum:widen_int_m]{\texttt{widen1}} & \texttt{26:29} & \texttt{widen\_int\_m (def none)} \\
\hyperref[enum:hadd_round_m]{\texttt{round}} & \texttt{30} & \texttt{hadd\_round\_m (def round\_down)} \\
\textit{Reserved} & \texttt{31:35} & \texttt{zero} \\
\hyperref[enum:widen_int_m]{\texttt{widen0}} & \texttt{36:39} & \texttt{widen\_int\_m (def none)} \\
\texttt{dest\_register\_index} & \texttt{40:45} & \texttt{uint} \\
\hyperref[enum:dest_width_m]{\texttt{dest\_width}} & \texttt{46:47} & \texttt{dest\_width\_m} \\
\texttt{datatype} & \texttt{48:51} & \texttt{exact=s32} \\
\texttt{instruction\_class} & \texttt{52:56} & \texttt{exact=R2} \\
\texttt{fau\_page\_index} & \texttt{57:58} & \texttt{uint} \\
\hyperref[enum:flow_control_m]{\texttt{flow}} & \texttt{59:62} & \texttt{flow\_control\_m} \\
\textit{Reserved} & \texttt{63} & \texttt{zero} \\
\bottomrule
\end{tabularx}
\end{center}
\textbf{Fields Explained:}
\begin{tightitemize}\footnotesize
  \item \texttt{src0}: Primary source operand (\texttt{r0..r63}, \texttt{u0..u63} uniform, \texttt{k*} FAU constant, or \texttt{t0..t2} temporary; \texttt{\textasciicircum{}} marks discard/last-use). \emph{Example:} {\scriptsize\ttfamily \mnem{FADD.f32} \dst{r0}, \src{r1}, \src{r2}} (\textbf{SPIR-V:} \texttt{OpFAdd (Operand 0)}; \textbf{GLSL:} \texttt{float c = a + b;})
  \item \texttt{src1}: Secondary source operand (register, uniform, or FAU constant; in shifts, supplies the shift count byte). \emph{Example:} {\scriptsize\ttfamily \mnem{LSHIFT.i32} \dst{r0}, \src{r1}, \src{r2}\imod{.b0}} (\textbf{SPIR-V:} \texttt{OpShiftLeftLogical (Shift)}; \textbf{GLSL:} \texttt{uint y = a \textless{}\textless{} s;})
  \item \texttt{opcode2}: Secondary sub-opcode further specializing instructions sharing \texttt{instruction\_class} and \texttt{opcode1} (e.g., conversions, SFU tables, moves). \emph{Example:} {\scriptsize\ttfamily \mnem{F32\_TO\_S32}\imod{.round\_zero} \dst{r0}, \src{r1}} (\textbf{SPIR-V:} \texttt{OpConvertFToS}; \textbf{GLSL:} \texttt{int y = int(x);})
  \item \hyperref[enum:widen_int_m]{\texttt{widen1}}: Stage A widening extractor on \texttt{src1} (promotes narrow byte/half/word lane before Stage B). \emph{Example:} {\scriptsize\ttfamily \mnem{IMUL.u64} \dst{[r0:r1]}, \src{[r2:r3]}\imod{.w0}, \src{[r4:r5]}\imod{.w0}} (\textbf{SPIR-V:} \texttt{OpUConvert / OpUMulExtended}; \textbf{GLSL:} \texttt{uint64\_t p = uint64\_t(a) * uint64\_t(b);})
  \item \hyperref[enum:hadd_round_m]{\texttt{round}}: Rounding mode (\texttt{none}/RNE, \texttt{round\_up}, \texttt{round\_down}, \texttt{round\_zero}, \texttt{round\_na}) applied in FP ALU, conversions, \texttt{HADD}, \texttt{LDEXP}, and Stage C narrowing. \emph{Example:} {\scriptsize\ttfamily \mnem{FROUND.f32}\imod{.round\_down} \dst{r0}, \src{r1}} (\textbf{SPIR-V:} \texttt{GLSL.std.450 Floor / Ceil / Trunc}; \textbf{GLSL:} \texttt{float y = floor(x);})
  \item \hyperref[enum:widen_int_m]{\texttt{widen0}}: Stage A widening extractor on \texttt{src0} (promotes FP16 \texttt{h0}/\texttt{h1} to FP32, or sign/zero-extends \texttt{b0..b3}/\texttt{h0..h1}/\texttt{w0..w1} to 32/64-bit before \texttt{abs}/\texttt{neg}). \emph{Example:} {\scriptsize\ttfamily \mnem{IADD.s32} \dst{r0}, \src{r1}\imod{.h0}, \src{r2}} (\textbf{SPIR-V:} \texttt{OpSConvert / OpFConvert}; \textbf{GLSL:} \texttt{int c = int(int16\_t(a)) + b;})
  \item \texttt{dest\_register\_index}: 6-bit destination GPR index (\texttt{r0..r63}) receiving the Stage C writeback result (even-aligned pair \texttt{[rN:rN+1]} for 64-bit ops). \emph{Example:} {\scriptsize\ttfamily \mnem{FMA.f32} \dst{r0}, \src{r1}, \src{r2}, \src{r3}} (\textbf{SPIR-V:} \texttt{SSA Result ID}; \textbf{GLSL:} \texttt{float r0 = fma(r1, r2, r3);})
  \item \hyperref[enum:dest_width_m]{\texttt{dest\_width}}: Stage C destination lane write-mask (\texttt{none} = full 32-bit, \texttt{h0} = low 16-bit, \texttt{h1} = high 16-bit without bit-shifting). \emph{Example:} {\scriptsize\ttfamily \mnem{FADD.v2f16} \dst{r0}\imod{.h1}, \src{r1}, \src{r2}} (\textbf{SPIR-V:} \texttt{OpVectorInsertDynamic / WriteMask}; \textbf{GLSL:} \texttt{dst.y = a.y + b.y;})
  \item \texttt{datatype}: Selects the lane data type (\texttt{f32}, \texttt{f16}, \texttt{v2f16}, \texttt{s32}, \texttt{u32}, \texttt{v2s16}, \texttt{v2u16}, \texttt{v4s8}, \texttt{v4u8}, \texttt{s64}, \texttt{u64}, etc.). \emph{Example:} {\scriptsize\ttfamily \mnem{IADD.v2s16} \dst{r0}, \src{r1}, \src{r2}} (\textbf{SPIR-V:} \texttt{OpTypeInt / OpTypeFloat}; \textbf{GLSL:} \texttt{i16vec2 c = a + b;})
  \item \texttt{instruction\_class}: 4-bit major instruction class selector (\texttt{bits 48:51}) partitioning the 64-bit Valhall opcode space (e.g., \texttt{FMA}, \texttt{CVT}, \texttt{SFU}, \texttt{LS}, \texttt{TEX}, \texttt{BRANCH}). \emph{Example:} {\scriptsize\ttfamily \mnem{FADD.f32} \dst{r0}, \src{r1}, \src{r2}} (\textbf{SPIR-V:} \texttt{N/A (encoding class)}; \textbf{GLSL:} \texttt{// Hardware execution pipe / decode class})
  \item \texttt{fau\_page\_index}: Selects the FAU (Fast Access Uniform) page (\texttt{0..3}) when a source operand references hardware small constants, thread IDs, or lane constants. \emph{Example:} {\scriptsize\ttfamily \mnem{FADD.f32} \dst{r0}, \src{r1}, \src{k1f}} (\textbf{SPIR-V:} \texttt{OpConstant / BuiltIn}; \textbf{GLSL:} \texttt{float y = x + 1.0;})
  \item \hyperref[enum:flow_control_m]{\texttt{flow}}: 4-bit embedded flow-control \& scoreboard wait modifier (\texttt{bits 59:62}: \texttt{none}, \texttt{wait0..wait0126}, \texttt{wait}, \texttt{reconverge}, \texttt{discard}, \texttt{end}). \emph{Example:} {\scriptsize\ttfamily \mnem{FADD.f32}\imod{.wait0.reconverge} \dst{r0}, \src{r1}, \src{r2}} (\textbf{SPIR-V:} \texttt{OpControlBarrier / OpReturn}; \textbf{GLSL:} \texttt{// Scoreboard wait + warp reconvergence})
\end{tightitemize}
\clearpage

\subsection*{IABS}
\addcontentsline{toc}{subsection}{IABS}
Computes the wrapping signed absolute value for each SIMD element $i \in \{0, \dots, N-1\}$ of \texttt{src0} ($N=1$ for \texttt{s32}, $N=2$ for \texttt{v2s16}, $N=4$ for \texttt{v4s8} on v9--v10): \texttt{dest[i] = (src0[i] \textless{} 0) ? -src0[i] : src0[i]} (wrapping on \texttt{INT\_MIN}).
\par\textbf{Notes:}
\begin{tightitemize}
  \item For \texttt{INT\_MIN} (\texttt{-2\textasciicircum{}(N-1)}), \texttt{unsigned\_abs()} returns \texttt{2\textasciicircum{}(N-1)} (identical bit pattern to \texttt{INT\_MIN}, wrapping, non-saturating; verified in \texttt{hw\_tests.rs:1815-1845}).
  \item \textbf{Source references:} \texttt{kraid/ops.rs:2211-2240}, \texttt{encode\_v9.rs:1736-1753}, \texttt{nir.rs:1064-1070}, \texttt{hw\_tests.rs:1815-1845}.
\end{tightitemize}
\codebox{\mnem{IABS}\imod{\{.flow\}} \dst{dest}\imod{\{.dest\_width\}}, \src{src0}\imod{\{.widen0\}}\\\mnem{IABS}\imod{\{.flow\}} \dst{dest}\imod{\{.dest\_width\}}, \src{src0}\imod{\{.swz0\}}\\\mnem{IABS}\imod{\{.flow\}} \dst{dest}\imod{\{.dest\_width\}}, \src{src0}\imod{\{.lanes0\}}}
\textbf{Operands:}
\begin{tightitemize}\footnotesize
  \item \texttt{\dst{dest}} (\texttt{dest\_register\_index}, \texttt{dest\_width}): Destination GPR receiving per-lane wrapping absolute values.
  \item \texttt{\src{src0}}: Signed integer input (\texttt{s32}, \texttt{v2s16}, or \texttt{v4s8}).
\end{tightitemize}
\par\textbf{Closest SPIR-V Op:} \texttt{GLSL.std.450 SAbs}
\par\textbf{Pseudo-GLSL (Operational Semantics):}
\codebox{for (int i = 0; i \textless{} N; ++i)\\    dest[i] = (src0[i] \textless{} 0) ? uint(-src0[i]) : uint(src0[i]); // INT\_MIN wraps}
\textbf{Execution Context:}
\begin{tightitemize}\footnotesize
  \item \textbf{Compute stage preloads (\hyperref[sec:csf-spd]{\texttt{SPD\_k}}, \texttt{r16..r23}):} Hardware pre-fills \texttt{r55..r57} with \texttt{WorkGroupID.xyz + JOB\_OFFSET} (from \hyperref[sec:csf-regs]{\texttt{COMPUTESR::JOB\_OFFSET\_X/Y/Z}} \texttt{r34..r36}) and \texttt{r60..r62} with \texttt{LocalInvocationID.xyz} or \texttt{GlobalInvocationID.xyz} (from \texttt{WG\_SIZE} \texttt{r33} \& \texttt{JOB\_SIZE} \texttt{r37..r39}, selected by \texttt{ShaderProgram::preload}) so compute shaders can derive thread/element indices without memory loads.
  \item \textbf{Vertex stage preloads (\hyperref[sec:csf-spd]{\texttt{SPD\_0}/\texttt{SPD\_1}}, \texttt{r16..r19}):} The hardware Index Fetcher reads \hyperref[sec:csf-regs]{\texttt{IDVSSR::INDEX\_BUFFER}} (\texttt{r54:r55}) and pre-fills \texttt{r59, r60} with \texttt{VertexIndex} (\texttt{+ VERTEX\_OFFSET + GLOBAL\_ATTR\_OFF}) and \texttt{r61} with \texttt{InstanceIndex} (\texttt{+ INSTANCE\_OFFSET}) so Position and Varying shaders can pull vertex attributes directly.
  \item \textbf{Inline FAU \& Special Operands (\hyperref[sec:csf-tsd-fau]{\texttt{FAU\_k}}, \texttt{r8..r15}):} Any 8-bit source with \texttt{mode=2} reads zero-latency uniform \texttt{u0..u31} from \texttt{FAU\_k} (\texttt{slots64 = bits[63:56]}); \texttt{mode=3} reads \texttt{SmallConstant} (Table~1.2) or \texttt{FauSpecial} (\texttt{pc}, \texttt{lane\_id}, \texttt{lane\_width}, \texttt{tls\_ptr}, \texttt{wls\_ptr}, \texttt{blend\_descriptor\_0..7}).
\end{tightitemize}
\par\textbf{Worked Examples:}
\par\textbf{Example 1 (Negative 32-bit signed integer):}
\codebox{\mnem{IABS.s32} \dst{r0}, \src{r1}}
With \texttt{r1 = 0xFFFFFFFB} (\texttt{-5}): negates to \texttt{r0 = 5} (\texttt{0x00000005}). Positive inputs pass through unchanged.
\par\textbf{Example 2 (INT32\_MIN wrapping edge case):}
\codebox{\mnem{IABS.s32} \dst{r0}, \src{r1}}
With \texttt{r1 = 0x80000000} (\texttt{-2147483648}): \texttt{IABS} is non-saturating, so \texttt{unsigned\_abs()} returns \texttt{0x80000000} (identical bit pattern, NOT \texttt{0x7FFFFFFF}).
\par\textbf{Mined SPIR-V $\to$ Valhall ISA Example (Signed integer absolute value before float conversion --- Vertex shader, Mali-G710 (v10)):}
\codebox{\textbf{; SPIR-V (input)}\\\%a = OpExtInst \%int \%glsl SAbs \%idx\\\%f = OpConvertUToF \%float \%a\\[3pt]\textbf{; Valhall ISA (malioc output)}\\IABS.s32 r1, r0\\U32\_TO\_F32 r2, r1\\FMUL.f32.clamp\_0\_inf r2, r2\textasciicircum{}, u4.w0}
\texttt{GLSL.std.450 SAbs} maps 1:1 to \texttt{IABS.s32}. Because the result is guaranteed non-negative (except \texttt{INT32\_MIN}, which wraps to \texttt{0x80000000}), \texttt{malioc} feeds \texttt{r1} directly into \texttt{U32\_TO\_F32}.
\begin{center}\footnotesize
\begin{tabularx}{\linewidth}{@{}lllX@{}}
\toprule
Variant & Arch & Class & Opcode Fields \\
\midrule
\texttt{(none)} & v9--v14 & \texttt{R1} & \texttt{opcode2=IABS\_S32, opcode1=MOV} \\
\texttt{(none)} & v9--v14 & \texttt{R1} & \texttt{opcode2=IABS\_V2S16, opcode1=MOV} \\
\texttt{(none)} & v9--v14 & \texttt{R1} & \texttt{opcode2=IABS\_V4S8, opcode1=MOV} \\
\bottomrule
\end{tabularx}
\end{center}
\begin{center}\footnotesize
\begin{tabularx}{0.88\linewidth}{@{}l l X@{}}
\toprule
Field Name & Bits & Type / Fixed Value \\
\midrule
\texttt{src0} & \texttt{0:7} & \texttt{SourceEncoding} \\
\textit{Reserved} & \texttt{8:15} & \texttt{zero} \\
\texttt{opcode2} & \texttt{16:20} & \texttt{exact=IABS\_S32} \\
\textit{Reserved} & \texttt{21:35} & \texttt{zero} \\
\hyperref[enum:widen_int_m]{\texttt{widen0}} & \texttt{36:39} & \texttt{widen\_int\_m (def none)} \\
\texttt{dest\_register\_index} & \texttt{40:45} & \texttt{uint} \\
\hyperref[enum:dest_width_m]{\texttt{dest\_width}} & \texttt{46:47} & \texttt{dest\_width\_m} \\
\texttt{opcode1} & \texttt{48:51} & \texttt{exact=MOV} \\
\texttt{instruction\_class} & \texttt{52:56} & \texttt{exact=R1} \\
\texttt{fau\_page\_index} & \texttt{57:58} & \texttt{uint} \\
\hyperref[enum:flow_control_m]{\texttt{flow}} & \texttt{59:62} & \texttt{flow\_control\_m} \\
\textit{Reserved} & \texttt{63} & \texttt{zero} \\
\bottomrule
\end{tabularx}
\end{center}
\textbf{Fields Explained:}
\begin{tightitemize}\footnotesize
  \item \texttt{src0}: Primary source operand (\texttt{r0..r63}, \texttt{u0..u63} uniform, \texttt{k*} FAU constant, or \texttt{t0..t2} temporary; \texttt{\textasciicircum{}} marks discard/last-use). \emph{Example:} {\scriptsize\ttfamily \mnem{FADD.f32} \dst{r0}, \src{r1}, \src{r2}} (\textbf{SPIR-V:} \texttt{OpFAdd (Operand 0)}; \textbf{GLSL:} \texttt{float c = a + b;})
  \item \texttt{opcode2}: Secondary sub-opcode further specializing instructions sharing \texttt{instruction\_class} and \texttt{opcode1} (e.g., conversions, SFU tables, moves). \emph{Example:} {\scriptsize\ttfamily \mnem{F32\_TO\_S32}\imod{.round\_zero} \dst{r0}, \src{r1}} (\textbf{SPIR-V:} \texttt{OpConvertFToS}; \textbf{GLSL:} \texttt{int y = int(x);})
  \item \hyperref[enum:widen_int_m]{\texttt{widen0}}: Stage A widening extractor on \texttt{src0} (promotes FP16 \texttt{h0}/\texttt{h1} to FP32, or sign/zero-extends \texttt{b0..b3}/\texttt{h0..h1}/\texttt{w0..w1} to 32/64-bit before \texttt{abs}/\texttt{neg}). \emph{Example:} {\scriptsize\ttfamily \mnem{IADD.s32} \dst{r0}, \src{r1}\imod{.h0}, \src{r2}} (\textbf{SPIR-V:} \texttt{OpSConvert / OpFConvert}; \textbf{GLSL:} \texttt{int c = int(int16\_t(a)) + b;})
  \item \texttt{dest\_register\_index}: 6-bit destination GPR index (\texttt{r0..r63}) receiving the Stage C writeback result (even-aligned pair \texttt{[rN:rN+1]} for 64-bit ops). \emph{Example:} {\scriptsize\ttfamily \mnem{FMA.f32} \dst{r0}, \src{r1}, \src{r2}, \src{r3}} (\textbf{SPIR-V:} \texttt{SSA Result ID}; \textbf{GLSL:} \texttt{float r0 = fma(r1, r2, r3);})
  \item \hyperref[enum:dest_width_m]{\texttt{dest\_width}}: Stage C destination lane write-mask (\texttt{none} = full 32-bit, \texttt{h0} = low 16-bit, \texttt{h1} = high 16-bit without bit-shifting). \emph{Example:} {\scriptsize\ttfamily \mnem{FADD.v2f16} \dst{r0}\imod{.h1}, \src{r1}, \src{r2}} (\textbf{SPIR-V:} \texttt{OpVectorInsertDynamic / WriteMask}; \textbf{GLSL:} \texttt{dst.y = a.y + b.y;})
  \item \texttt{opcode1}: Primary sub-opcode within the instruction class uniquely identifying the operation family. \emph{Example:} {\scriptsize\ttfamily \mnem{IADD.s32} \dst{r0}, \src{r1}, \src{r2}} (\textbf{SPIR-V:} \texttt{OpIAdd (operation selector)}; \textbf{GLSL:} \texttt{// Primary opcode bits})
  \item \texttt{instruction\_class}: 4-bit major instruction class selector (\texttt{bits 48:51}) partitioning the 64-bit Valhall opcode space (e.g., \texttt{FMA}, \texttt{CVT}, \texttt{SFU}, \texttt{LS}, \texttt{TEX}, \texttt{BRANCH}). \emph{Example:} {\scriptsize\ttfamily \mnem{FADD.f32} \dst{r0}, \src{r1}, \src{r2}} (\textbf{SPIR-V:} \texttt{N/A (encoding class)}; \textbf{GLSL:} \texttt{// Hardware execution pipe / decode class})
  \item \texttt{fau\_page\_index}: Selects the FAU (Fast Access Uniform) page (\texttt{0..3}) when a source operand references hardware small constants, thread IDs, or lane constants. \emph{Example:} {\scriptsize\ttfamily \mnem{FADD.f32} \dst{r0}, \src{r1}, \src{k1f}} (\textbf{SPIR-V:} \texttt{OpConstant / BuiltIn}; \textbf{GLSL:} \texttt{float y = x + 1.0;})
  \item \hyperref[enum:flow_control_m]{\texttt{flow}}: 4-bit embedded flow-control \& scoreboard wait modifier (\texttt{bits 59:62}: \texttt{none}, \texttt{wait0..wait0126}, \texttt{wait}, \texttt{reconverge}, \texttt{discard}, \texttt{end}). \emph{Example:} {\scriptsize\ttfamily \mnem{FADD.f32}\imod{.wait0.reconverge} \dst{r0}, \src{r1}, \src{r2}} (\textbf{SPIR-V:} \texttt{OpControlBarrier / OpReturn}; \textbf{GLSL:} \texttt{// Scoreboard wait + warp reconvergence})
\end{tightitemize}
\clearpage

\subsection*{IADD}
\addcontentsline{toc}{subsection}{IADD}
Adds Stage A widened/swizzled \texttt{src0} and \texttt{src1} independently for each SIMD element $i \in \{0, \dots, N-1\}$ ($N=1$ for \texttt{s32}/\texttt{u32}/\texttt{s64}/\texttt{u64}, $N=2$ for \texttt{v2s16}/\texttt{v2u16}, $N=4$ for \texttt{v4s8}/\texttt{v4u8} on v9--v10): \texttt{dest[i] = sat(widen0(src0)[i] + widen1(src1)[i])} when \texttt{.sat} is set, or wraps modulo $2^W$ otherwise.
\par\textbf{Notes:}
\begin{tightitemize}
  \item 64-bit \texttt{IADD.sat} is NOT supported by HW (\texttt{kraid/ops.rs:2281-2284}, \texttt{hw\_tests.rs:1885}).
  \item \texttt{v4s8}/\texttt{v4u8} only exist on arch \textless{}= v10 (\texttt{hw\_tests.rs:1876}).
  \item \textbf{Source references:} \texttt{kraid/ops.rs:2250-2298}, \texttt{encode\_v9.rs:1755-1793}, \texttt{nir.rs:1091-1105}, \texttt{hw\_tests.rs:1848-1902}.
\end{tightitemize}
\codebox{\mnem{IADD}\imod{\{.flow\}}\imod{\{.saturate\}} \dst{dest}\imod{\{.dest\_width\}}, \src{src0}\imod{\{.widen0\}}, \src{src1}\imod{\{.widen1\}}\\\mnem{IADD}\imod{\{.flow\}}\imod{\{.saturate\}} \dst{dest}, \src{src0}\imod{\{.widen0\}}, \src{src1}\imod{\{.widen1\}}\\\mnem{IADD}\imod{\{.flow\}} \dst{dest}, \src{src0}\imod{\{.widen0\}}, \src{src1}\imod{\{.widen1\}}\\\mnem{IADD}\imod{\{.flow\}}\imod{\{.saturate\}} \dst{dest}\imod{\{.dest\_width\}}, \src{src0}\imod{\{.swz0\}}, \src{src1}\imod{\{.swz1\}}\\\mnem{IADD}\imod{\{.flow\}}\imod{\{.saturate\}} \dst{dest}\imod{\{.dest\_width\}}, \src{src0}\imod{\{.lanes0\}}, \src{src1}\imod{\{.lanes1\}}}
\textbf{Operands:}
\begin{tightitemize}\footnotesize
  \item \texttt{\dst{dest}} (\texttt{dest\_register\_index}, \texttt{dest\_width}): Destination GPR (or register pair for \texttt{s64}/\texttt{u64}) receiving the sum.
  \item \texttt{\src{src0}}: First addend (supports Stage A widening/swizzle).
  \item \texttt{\src{src1}}: Second addend (supports Stage A widening/swizzle).
\end{tightitemize}
\par\textbf{Closest SPIR-V Op:} \texttt{OpIAdd}
\par\textbf{Pseudo-GLSL (Operational Semantics):}
\codebox{for (int i = 0; i \textless{} N; ++i) \{\\    auto sum = ext(widen0(src0)[i]) + ext(widen1(src1)[i]);\\    dest[i] = flag\_sat ? clamp(sum, MIN\_W, MAX\_W) : wrap\_W(sum);\\\}}
\textbf{Execution Context:}
\begin{tightitemize}\footnotesize
  \item \textbf{Compute stage preloads (\hyperref[sec:csf-spd]{\texttt{SPD\_k}}, \texttt{r16..r23}):} Hardware pre-fills \texttt{r55..r57} with \texttt{WorkGroupID.xyz + JOB\_OFFSET} (from \hyperref[sec:csf-regs]{\texttt{COMPUTESR::JOB\_OFFSET\_X/Y/Z}} \texttt{r34..r36}) and \texttt{r60..r62} with \texttt{LocalInvocationID.xyz} or \texttt{GlobalInvocationID.xyz} (from \texttt{WG\_SIZE} \texttt{r33} \& \texttt{JOB\_SIZE} \texttt{r37..r39}, selected by \texttt{ShaderProgram::preload}) so compute shaders can derive thread/element indices without memory loads.
  \item \textbf{Vertex stage preloads (\hyperref[sec:csf-spd]{\texttt{SPD\_0}/\texttt{SPD\_1}}, \texttt{r16..r19}):} The hardware Index Fetcher reads \hyperref[sec:csf-regs]{\texttt{IDVSSR::INDEX\_BUFFER}} (\texttt{r54:r55}) and pre-fills \texttt{r59, r60} with \texttt{VertexIndex} (\texttt{+ VERTEX\_OFFSET + GLOBAL\_ATTR\_OFF}) and \texttt{r61} with \texttt{InstanceIndex} (\texttt{+ INSTANCE\_OFFSET}) so Position and Varying shaders can pull vertex attributes directly.
  \item \textbf{Inline FAU \& Special Operands (\hyperref[sec:csf-tsd-fau]{\texttt{FAU\_k}}, \texttt{r8..r15}):} Any 8-bit source with \texttt{mode=2} reads zero-latency uniform \texttt{u0..u31} from \texttt{FAU\_k} (\texttt{slots64 = bits[63:56]}); \texttt{mode=3} reads \texttt{SmallConstant} (Table~1.2) or \texttt{FauSpecial} (\texttt{pc}, \texttt{lane\_id}, \texttt{lane\_width}, \texttt{tls\_ptr}, \texttt{wls\_ptr}, \texttt{blend\_descriptor\_0..7}).
\end{tightitemize}
\par\textbf{Worked Examples:}
\par\textbf{Example 1 (Widening + saturating 32-bit signed add):}
\codebox{\mnem{IADD.s32}\imod{.sat} \dst{r0}, \src{r1}\imod{.h0}, \src{r2}}
With \texttt{r1 = 0x00007FFF} (\texttt{+32767}), \texttt{r2 = 0x7FFFFFF0}: Stage A sign-extends \texttt{r1.h0} to \texttt{0x00007FFF}; Stage B clamps overflow to \texttt{0x7FFFFFFF}.
\par\textbf{Example 2 (Packed SIMD v2s16 with lane-mask writeback):}
\codebox{\mnem{IADD.v2s16} \dst{r0}\imod{.h1}, \src{r1}\imod{.h01}, \src{r2}\imod{.h01}}
With \texttt{r1 = [10, 20]} (\texttt{0x0014000A}), \texttt{r2 = [5, -3]} (\texttt{0xFFFD0005}), \texttt{r0\_old = 0x0000BEEF}: lane 1 computes \texttt{20 + (-3) = 17} (\texttt{0x0011}); Stage C \texttt{.h1} writes \texttt{r0 = 0x0011BEEF}.
\par\textbf{Mined SPIR-V $\to$ Valhall ISA Example (Buffer byte-offset calculation from vertex attribute and UBO base --- Vertex shader, Mali-G710 (v10)):}
\codebox{\textbf{; SPIR-V (input)}\\\%idx = OpLoad \%uint \%attr\_in\\\%off = OpIAdd \%uint \%ubo\_base \%idx\\\%byte\_off = OpIMul \%uint \%off \%uint\_16\\[3pt]\textbf{; Valhall ISA (malioc output)}\\IADD.u32 r0, u3.w0, r0\textasciicircum{}\\IMUL.u32 r0, r0\textasciicircum{}, k13121110.b0\\LD\_PKA.i32.slot0.wait0 @r0, r0\textasciicircum{}, u12.w1}
\texttt{OpIAdd} lowers to \texttt{IADD.u32} reading the uniform base directly from FAU (\texttt{u3.w0}) and killing the last use of \texttt{r0} (\texttt{r0\textasciicircum{}}). The resulting index is scaled by 16 (\texttt{k13121110.b0 = 0x10}) for \texttt{LD\_PKA}.
\begin{center}\footnotesize
\begin{tabularx}{\linewidth}{@{}lllX@{}}
\toprule
Variant & Arch & Class & Opcode Fields \\
\midrule
\texttt{s32} & v9--v14 & \texttt{R2} & \texttt{opcode2=ADD} \\
\texttt{s64} & v9--v14 & \texttt{D2} & \texttt{opcode2=ADD} \\
\texttt{s64} & v9--v14 & \texttt{D2} & \texttt{opcode2=ADD} \\
\texttt{u32} & v9--v14 & \texttt{R2} & \texttt{opcode2=ADD} \\
\texttt{u64} & v9--v14 & \texttt{D2} & \texttt{opcode2=ADD} \\
\texttt{u64} & v9--v14 & \texttt{D2} & \texttt{opcode2=ADD} \\
\texttt{v2s16} & v9--v14 & \texttt{R2} & \texttt{opcode2=ADD} \\
\texttt{v2u16} & v9--v14 & \texttt{R2} & \texttt{opcode2=ADD} \\
\texttt{v4s8} & v9--v14 & \texttt{R2} & \texttt{opcode2=ADD} \\
\texttt{v4u8} & v9--v14 & \texttt{R2} & \texttt{opcode2=ADD} \\
\bottomrule
\end{tabularx}
\end{center}
\begin{center}\footnotesize
\begin{tabularx}{0.88\linewidth}{@{}l l X@{}}
\toprule
Field Name & Bits & Type / Fixed Value \\
\midrule
\texttt{src0} & \texttt{0:7} & \texttt{SourceEncoding} \\
\texttt{src1} & \texttt{8:15} & \texttt{SourceEncoding} \\
\texttt{opcode2} & \texttt{16:19} & \texttt{exact=ADD} \\
\textit{Reserved} & \texttt{20:25} & \texttt{zero} \\
\hyperref[enum:widen_int_m]{\texttt{widen1}} & \texttt{26:29} & \texttt{widen\_int\_m (def none)} \\
\hyperref[enum:saturate_m]{\texttt{saturate}} & \texttt{30} & \texttt{saturate\_m (def none)} \\
\textit{Reserved} & \texttt{31:35} & \texttt{zero} \\
\hyperref[enum:widen_int_m]{\texttt{widen0}} & \texttt{36:39} & \texttt{widen\_int\_m (def none)} \\
\texttt{dest\_register\_index} & \texttt{40:45} & \texttt{uint} \\
\hyperref[enum:dest_width_m]{\texttt{dest\_width}} & \texttt{46:47} & \texttt{dest\_width\_m} \\
\texttt{datatype} & \texttt{48:51} & \texttt{exact=s32} \\
\texttt{instruction\_class} & \texttt{52:56} & \texttt{exact=R2} \\
\texttt{fau\_page\_index} & \texttt{57:58} & \texttt{uint} \\
\hyperref[enum:flow_control_m]{\texttt{flow}} & \texttt{59:62} & \texttt{flow\_control\_m} \\
\textit{Reserved} & \texttt{63} & \texttt{zero} \\
\bottomrule
\end{tabularx}
\end{center}
\textbf{Fields Explained:}
\begin{tightitemize}\footnotesize
  \item \texttt{src0}: Primary source operand (\texttt{r0..r63}, \texttt{u0..u63} uniform, \texttt{k*} FAU constant, or \texttt{t0..t2} temporary; \texttt{\textasciicircum{}} marks discard/last-use). \emph{Example:} {\scriptsize\ttfamily \mnem{FADD.f32} \dst{r0}, \src{r1}, \src{r2}} (\textbf{SPIR-V:} \texttt{OpFAdd (Operand 0)}; \textbf{GLSL:} \texttt{float c = a + b;})
  \item \texttt{src1}: Secondary source operand (register, uniform, or FAU constant; in shifts, supplies the shift count byte). \emph{Example:} {\scriptsize\ttfamily \mnem{LSHIFT.i32} \dst{r0}, \src{r1}, \src{r2}\imod{.b0}} (\textbf{SPIR-V:} \texttt{OpShiftLeftLogical (Shift)}; \textbf{GLSL:} \texttt{uint y = a \textless{}\textless{} s;})
  \item \texttt{opcode2}: Secondary sub-opcode further specializing instructions sharing \texttt{instruction\_class} and \texttt{opcode1} (e.g., conversions, SFU tables, moves). \emph{Example:} {\scriptsize\ttfamily \mnem{F32\_TO\_S32}\imod{.round\_zero} \dst{r0}, \src{r1}} (\textbf{SPIR-V:} \texttt{OpConvertFToS}; \textbf{GLSL:} \texttt{int y = int(x);})
  \item \hyperref[enum:widen_int_m]{\texttt{widen1}}: Stage A widening extractor on \texttt{src1} (promotes narrow byte/half/word lane before Stage B). \emph{Example:} {\scriptsize\ttfamily \mnem{IMUL.u64} \dst{[r0:r1]}, \src{[r2:r3]}\imod{.w0}, \src{[r4:r5]}\imod{.w0}} (\textbf{SPIR-V:} \texttt{OpUConvert / OpUMulExtended}; \textbf{GLSL:} \texttt{uint64\_t p = uint64\_t(a) * uint64\_t(b);})
  \item \hyperref[enum:saturate_m]{\texttt{saturate}}: Stage B integer saturation flag (\texttt{.sat}): clamps signed/unsigned overflow or underflow to \texttt{[MIN\_N, MAX\_N]} / \texttt{[0, 2\textasciicircum{}N - 1]} instead of wrapping. \emph{Example:} {\scriptsize\ttfamily \mnem{IADD.s32}\imod{.sat} \dst{r0}, \src{r1}, \src{r2}} (\textbf{SPIR-V:} \texttt{OpSDotAccSat / saturating add}; \textbf{GLSL:} \texttt{int c = clamp(int64\_t(a) + int64\_t(b), INT\_MIN, INT\_MAX);})
  \item \hyperref[enum:widen_int_m]{\texttt{widen0}}: Stage A widening extractor on \texttt{src0} (promotes FP16 \texttt{h0}/\texttt{h1} to FP32, or sign/zero-extends \texttt{b0..b3}/\texttt{h0..h1}/\texttt{w0..w1} to 32/64-bit before \texttt{abs}/\texttt{neg}). \emph{Example:} {\scriptsize\ttfamily \mnem{IADD.s32} \dst{r0}, \src{r1}\imod{.h0}, \src{r2}} (\textbf{SPIR-V:} \texttt{OpSConvert / OpFConvert}; \textbf{GLSL:} \texttt{int c = int(int16\_t(a)) + b;})
  \item \texttt{dest\_register\_index}: 6-bit destination GPR index (\texttt{r0..r63}) receiving the Stage C writeback result (even-aligned pair \texttt{[rN:rN+1]} for 64-bit ops). \emph{Example:} {\scriptsize\ttfamily \mnem{FMA.f32} \dst{r0}, \src{r1}, \src{r2}, \src{r3}} (\textbf{SPIR-V:} \texttt{SSA Result ID}; \textbf{GLSL:} \texttt{float r0 = fma(r1, r2, r3);})
  \item \hyperref[enum:dest_width_m]{\texttt{dest\_width}}: Stage C destination lane write-mask (\texttt{none} = full 32-bit, \texttt{h0} = low 16-bit, \texttt{h1} = high 16-bit without bit-shifting). \emph{Example:} {\scriptsize\ttfamily \mnem{FADD.v2f16} \dst{r0}\imod{.h1}, \src{r1}, \src{r2}} (\textbf{SPIR-V:} \texttt{OpVectorInsertDynamic / WriteMask}; \textbf{GLSL:} \texttt{dst.y = a.y + b.y;})
  \item \texttt{datatype}: Selects the lane data type (\texttt{f32}, \texttt{f16}, \texttt{v2f16}, \texttt{s32}, \texttt{u32}, \texttt{v2s16}, \texttt{v2u16}, \texttt{v4s8}, \texttt{v4u8}, \texttt{s64}, \texttt{u64}, etc.). \emph{Example:} {\scriptsize\ttfamily \mnem{IADD.v2s16} \dst{r0}, \src{r1}, \src{r2}} (\textbf{SPIR-V:} \texttt{OpTypeInt / OpTypeFloat}; \textbf{GLSL:} \texttt{i16vec2 c = a + b;})
  \item \texttt{instruction\_class}: 4-bit major instruction class selector (\texttt{bits 48:51}) partitioning the 64-bit Valhall opcode space (e.g., \texttt{FMA}, \texttt{CVT}, \texttt{SFU}, \texttt{LS}, \texttt{TEX}, \texttt{BRANCH}). \emph{Example:} {\scriptsize\ttfamily \mnem{FADD.f32} \dst{r0}, \src{r1}, \src{r2}} (\textbf{SPIR-V:} \texttt{N/A (encoding class)}; \textbf{GLSL:} \texttt{// Hardware execution pipe / decode class})
  \item \texttt{fau\_page\_index}: Selects the FAU (Fast Access Uniform) page (\texttt{0..3}) when a source operand references hardware small constants, thread IDs, or lane constants. \emph{Example:} {\scriptsize\ttfamily \mnem{FADD.f32} \dst{r0}, \src{r1}, \src{k1f}} (\textbf{SPIR-V:} \texttt{OpConstant / BuiltIn}; \textbf{GLSL:} \texttt{float y = x + 1.0;})
  \item \hyperref[enum:flow_control_m]{\texttt{flow}}: 4-bit embedded flow-control \& scoreboard wait modifier (\texttt{bits 59:62}: \texttt{none}, \texttt{wait0..wait0126}, \texttt{wait}, \texttt{reconverge}, \texttt{discard}, \texttt{end}). \emph{Example:} {\scriptsize\ttfamily \mnem{FADD.f32}\imod{.wait0.reconverge} \dst{r0}, \src{r1}, \src{r2}} (\textbf{SPIR-V:} \texttt{OpControlBarrier / OpReturn}; \textbf{GLSL:} \texttt{// Scoreboard wait + warp reconvergence})
\end{tightitemize}
\clearpage

\subsection*{IADDC}
\addcontentsline{toc}{subsection}{IADDC}
Computes 32-bit unsigned addition with 1-bit carry-in: \texttt{sum64 = uint64(src0) + uint64(src1) + uint64(src2 \& 1)}. By default writes the low 32-bit sum \texttt{sum64[31:0]} to \texttt{dest}; when \texttt{.c} (\texttt{return\_carry}) is set, writes the 1-bit carry-out \texttt{(sum64 \textgreater{}\textgreater{} 32) \& 1} to \texttt{dest} instead.
\par\textbf{Notes:}
\begin{tightitemize}
  \item Unsigned 32-bit only.
  \item Carry-in uses LSB of \texttt{src2}.
  \item Used on Bifrost/Valhall as FMA-unit \texttt{IADD} when \texttt{src2=0} (\texttt{bi\_schedule.c:537}).
  \item \textbf{Source references:} \texttt{isa-v9-v14.xml}, \texttt{bifrost/bifrost/ISA.xml:1474}, \texttt{bi\_schedule.c:537-544}.
\end{tightitemize}
\codebox{\mnem{IADDC}\imod{\{.flow\}}\imod{\{.return\_carry\}} \dst{dest}\imod{\{.dest\_width\}}, \src{src0}, \src{src1}, \src{src2}}
\textbf{Operands:}
\begin{tightitemize}\footnotesize
  \item \texttt{\dst{dest}} (\texttt{dest\_register\_index}, \texttt{dest\_width}): 32-bit destination GPR receiving the low sum word (default) or carry-out bit (\texttt{.c}).
  \item \texttt{\src{src0}}: First 32-bit unsigned addend.
  \item \texttt{\src{src1}}: Second 32-bit unsigned addend.
  \item \texttt{\src{src2}}: Carry-in operand (only bit 0, \texttt{src2 \& 1}, is added).
\end{tightitemize}
\par\textbf{Closest SPIR-V Op:} \texttt{OpIAddCarry}
\par\textbf{Pseudo-GLSL (Operational Semantics):}
\codebox{uint64\_t sum = uint64\_t(src0) + uint64\_t(src1) + uint64\_t(src2 \& 1u);\\dest = (return\_carry == c) ? uint(sum \textgreater{}\textgreater{} 32) : uint(sum \& 0xFFFFFFFFul);}
\textbf{Execution Context:}
\begin{tightitemize}\footnotesize
  \item \textbf{Compute stage preloads (\hyperref[sec:csf-spd]{\texttt{SPD\_k}}, \texttt{r16..r23}):} Hardware pre-fills \texttt{r55..r57} with \texttt{WorkGroupID.xyz + JOB\_OFFSET} (from \hyperref[sec:csf-regs]{\texttt{COMPUTESR::JOB\_OFFSET\_X/Y/Z}} \texttt{r34..r36}) and \texttt{r60..r62} with \texttt{LocalInvocationID.xyz} or \texttt{GlobalInvocationID.xyz} (from \texttt{WG\_SIZE} \texttt{r33} \& \texttt{JOB\_SIZE} \texttt{r37..r39}, selected by \texttt{ShaderProgram::preload}) so compute shaders can derive thread/element indices without memory loads.
  \item \textbf{Vertex stage preloads (\hyperref[sec:csf-spd]{\texttt{SPD\_0}/\texttt{SPD\_1}}, \texttt{r16..r19}):} The hardware Index Fetcher reads \hyperref[sec:csf-regs]{\texttt{IDVSSR::INDEX\_BUFFER}} (\texttt{r54:r55}) and pre-fills \texttt{r59, r60} with \texttt{VertexIndex} (\texttt{+ VERTEX\_OFFSET + GLOBAL\_ATTR\_OFF}) and \texttt{r61} with \texttt{InstanceIndex} (\texttt{+ INSTANCE\_OFFSET}) so Position and Varying shaders can pull vertex attributes directly.
  \item \textbf{Inline FAU \& Special Operands (\hyperref[sec:csf-tsd-fau]{\texttt{FAU\_k}}, \texttt{r8..r15}):} Any 8-bit source with \texttt{mode=2} reads zero-latency uniform \texttt{u0..u31} from \texttt{FAU\_k} (\texttt{slots64 = bits[63:56]}); \texttt{mode=3} reads \texttt{SmallConstant} (Table~1.2) or \texttt{FauSpecial} (\texttt{pc}, \texttt{lane\_id}, \texttt{lane\_width}, \texttt{tls\_ptr}, \texttt{wls\_ptr}, \texttt{blend\_descriptor\_0..7}).
\end{tightitemize}
\par\textbf{Worked Examples:}
\par\textbf{Example 1 (Carry-out extraction (.c)):}
\codebox{\mnem{IADDC.u32}\imod{.c} \dst{r0}, \src{r1}, \src{r2}, \src{r3}}
With \texttt{r1 = 0xFFFFFFFF, r2 = 0, r3 = 1}: 64-bit sum is \texttt{0x1\_00000000}; \texttt{.c} returns the carry-out bit \texttt{r0 = 1}.
\par\textbf{Example 2 (32-bit sum with carry-in (no .c)):}
\codebox{\mnem{IADDC.u32} \dst{r0}, \src{r1}, \src{r2}, \src{r3}}
With \texttt{r1 = 0xFFFFFFFF, r2 = 5, r3 = 0x00000003} (LSB=1): \texttt{0xFFFFFFFF + 5 + 1 = 0x1\_00000005}; without \texttt{.c}, returns low 32 bits \texttt{r0 = 5}.

\begin{center}\footnotesize
\begin{tabularx}{\linewidth}{@{}lllX@{}}
\toprule
Variant & Arch & Class & Opcode Fields \\
\midrule
\texttt{(none)} & v9--v14 & \texttt{R3} & \texttt{opcode3=IADDC, opcode1=MULTI} \\
\bottomrule
\end{tabularx}
\end{center}
\begin{center}\footnotesize
\begin{tabularx}{0.88\linewidth}{@{}l l X@{}}
\toprule
Field Name & Bits & Type / Fixed Value \\
\midrule
\texttt{src0} & \texttt{0:7} & \texttt{SourceEncoding} \\
\texttt{src1} & \texttt{8:15} & \texttt{SourceEncoding} \\
\texttt{src2} & \texttt{16:23} & \texttt{SourceEncoding} \\
\texttt{opcode3} & \texttt{24:25} & \texttt{exact=IADDC} \\
\textit{Reserved} & \texttt{26:29} & \texttt{zero} \\
\hyperref[enum:iaddc_return_carry_m]{\texttt{return\_carry}} & \texttt{30} & \texttt{iaddc\_return\_carry\_m (def none)} \\
\textit{Reserved} & \texttt{31:39} & \texttt{zero} \\
\texttt{dest\_register\_index} & \texttt{40:45} & \texttt{uint} \\
\hyperref[enum:dest_width_m]{\texttt{dest\_width}} & \texttt{46:47} & \texttt{dest\_width\_m} \\
\texttt{opcode1} & \texttt{48:51} & \texttt{exact=MULTI} \\
\texttt{instruction\_class} & \texttt{52:56} & \texttt{exact=R3} \\
\texttt{fau\_page\_index} & \texttt{57:58} & \texttt{uint} \\
\hyperref[enum:flow_control_m]{\texttt{flow}} & \texttt{59:62} & \texttt{flow\_control\_m} \\
\textit{Reserved} & \texttt{63} & \texttt{zero} \\
\bottomrule
\end{tabularx}
\end{center}
\textbf{Fields Explained:}
\begin{tightitemize}\footnotesize
  \item \texttt{src0}: Primary source operand (\texttt{r0..r63}, \texttt{u0..u63} uniform, \texttt{k*} FAU constant, or \texttt{t0..t2} temporary; \texttt{\textasciicircum{}} marks discard/last-use). \emph{Example:} {\scriptsize\ttfamily \mnem{FADD.f32} \dst{r0}, \src{r1}, \src{r2}} (\textbf{SPIR-V:} \texttt{OpFAdd (Operand 0)}; \textbf{GLSL:} \texttt{float c = a + b;})
  \item \texttt{src1}: Secondary source operand (register, uniform, or FAU constant; in shifts, supplies the shift count byte). \emph{Example:} {\scriptsize\ttfamily \mnem{LSHIFT.i32} \dst{r0}, \src{r1}, \src{r2}\imod{.b0}} (\textbf{SPIR-V:} \texttt{OpShiftLeftLogical (Shift)}; \textbf{GLSL:} \texttt{uint y = a \textless{}\textless{} s;})
  \item \texttt{src2}: Third source operand (FMA addend, \texttt{MUX}/\texttt{CSEL} selector, \texttt{IADDC} carry-in, or bitwise mask in fused shift-logic). \emph{Example:} {\scriptsize\ttfamily \mnem{FMA.f32} \dst{r0}, \src{r1}, \src{r2}, \src{r3}} (\textbf{SPIR-V:} \texttt{GLSL.std.450 Fma (Addend)}; \textbf{GLSL:} \texttt{float d = fma(a, b, c);})
  \item \texttt{opcode3}: Tertiary sub-opcode disambiguating specialized multi-op families (\texttt{FPOW\_SC\_*}, SFU lookup tables, \texttt{ICMP\_MULTI}). \emph{Example:} {\scriptsize\ttfamily \mnem{FPOW\_SC\_DET.f32}\imod{.pow} \dst{r0}, \src{r1}, \src{r2}} (\textbf{SPIR-V:} \texttt{GLSL.std.450 Pow}; \textbf{GLSL:} \texttt{float z = pow(x, y);})
  \item \hyperref[enum:iaddc_return_carry_m]{\texttt{return\_carry}}: When \texttt{.c} is set on \texttt{IADDC}/\texttt{ISUBB}, returns the 1-bit carry/borrow instead of the 32-bit sum/diff. \emph{Example:} {\scriptsize\ttfamily \mnem{IADDC.u32}\imod{.c} \dst{r0}, \src{r1}, \src{r2}} (\textbf{SPIR-V:} \texttt{OpIAddCarry}; \textbf{GLSL:} \texttt{uint carry; uint s = uaddCarry(a, b, carry);})
  \item \texttt{dest\_register\_index}: 6-bit destination GPR index (\texttt{r0..r63}) receiving the Stage C writeback result (even-aligned pair \texttt{[rN:rN+1]} for 64-bit ops). \emph{Example:} {\scriptsize\ttfamily \mnem{FMA.f32} \dst{r0}, \src{r1}, \src{r2}, \src{r3}} (\textbf{SPIR-V:} \texttt{SSA Result ID}; \textbf{GLSL:} \texttt{float r0 = fma(r1, r2, r3);})
  \item \hyperref[enum:dest_width_m]{\texttt{dest\_width}}: Stage C destination lane write-mask (\texttt{none} = full 32-bit, \texttt{h0} = low 16-bit, \texttt{h1} = high 16-bit without bit-shifting). \emph{Example:} {\scriptsize\ttfamily \mnem{FADD.v2f16} \dst{r0}\imod{.h1}, \src{r1}, \src{r2}} (\textbf{SPIR-V:} \texttt{OpVectorInsertDynamic / WriteMask}; \textbf{GLSL:} \texttt{dst.y = a.y + b.y;})
  \item \texttt{opcode1}: Primary sub-opcode within the instruction class uniquely identifying the operation family. \emph{Example:} {\scriptsize\ttfamily \mnem{IADD.s32} \dst{r0}, \src{r1}, \src{r2}} (\textbf{SPIR-V:} \texttt{OpIAdd (operation selector)}; \textbf{GLSL:} \texttt{// Primary opcode bits})
  \item \texttt{instruction\_class}: 4-bit major instruction class selector (\texttt{bits 48:51}) partitioning the 64-bit Valhall opcode space (e.g., \texttt{FMA}, \texttt{CVT}, \texttt{SFU}, \texttt{LS}, \texttt{TEX}, \texttt{BRANCH}). \emph{Example:} {\scriptsize\ttfamily \mnem{FADD.f32} \dst{r0}, \src{r1}, \src{r2}} (\textbf{SPIR-V:} \texttt{N/A (encoding class)}; \textbf{GLSL:} \texttt{// Hardware execution pipe / decode class})
  \item \texttt{fau\_page\_index}: Selects the FAU (Fast Access Uniform) page (\texttt{0..3}) when a source operand references hardware small constants, thread IDs, or lane constants. \emph{Example:} {\scriptsize\ttfamily \mnem{FADD.f32} \dst{r0}, \src{r1}, \src{k1f}} (\textbf{SPIR-V:} \texttt{OpConstant / BuiltIn}; \textbf{GLSL:} \texttt{float y = x + 1.0;})
  \item \hyperref[enum:flow_control_m]{\texttt{flow}}: 4-bit embedded flow-control \& scoreboard wait modifier (\texttt{bits 59:62}: \texttt{none}, \texttt{wait0..wait0126}, \texttt{wait}, \texttt{reconverge}, \texttt{discard}, \texttt{end}). \emph{Example:} {\scriptsize\ttfamily \mnem{FADD.f32}\imod{.wait0.reconverge} \dst{r0}, \src{r1}, \src{r2}} (\textbf{SPIR-V:} \texttt{OpControlBarrier / OpReturn}; \textbf{GLSL:} \texttt{// Scoreboard wait + warp reconvergence})
\end{tightitemize}
\clearpage

\subsection*{IADD\_IMM}
\addcontentsline{toc}{subsection}{IADD\_IMM}
Adds the inline 32-bit immediate \texttt{imm1w} to \texttt{src0} with wrapping arithmetic for each SIMD element $i \in \{0, \dots, N-1\}$ ($N=1$ for \texttt{i32}, $N=2$ for \texttt{v2i16}, $N=4$ for \texttt{v4i8} on v9--v10): \texttt{dest[i] = src0[i] + imm1w[i]}.
\par\textbf{Notes:}
\begin{tightitemize}
  \item With \texttt{src0 = 0}, canonical immediate move (\texttt{ISA.xml:62}).
  \item \texttt{v4i8} only arch \textless{}= v10.
  \item \textbf{Source references:} \texttt{kraid/encode\_v9.rs:1766-1782}, \texttt{bifrost/valhall/ISA.xml}.
\end{tightitemize}
\codebox{\mnem{IADD\_IMM}\imod{\{.flow\}} \dst{dest}\imod{\{.dest\_width\}}, \src{src0}, \imm{\#imm1w}}
\textbf{Operands:}
\begin{tightitemize}\footnotesize
  \item \texttt{\dst{dest}} (\texttt{dest\_register\_index}, \texttt{dest\_width}): Destination GPR receiving \texttt{src0 + \#imm1w}.
  \item \texttt{\src{src0}}: Input register (\texttt{k0} encodes the \texttt{MOV\_IMM} alias).
  \item \texttt{\imm{\#imm1w}} (\texttt{imm1w}): Inline 32-bit immediate addend (interpreted per lane for \texttt{v2i16}/\texttt{v4i8}).
\end{tightitemize}
\par\textbf{Closest SPIR-V Op:} \texttt{OpIAdd (constant RHS)}
\par\textbf{Pseudo-GLSL (Operational Semantics):}
\codebox{for (int i = 0; i \textless{} N; ++i)\\    dest[i] = wrap\_W(src0[i] + \#imm1w[i]); // MOV\_IMM when src0 == k0}
\textbf{Execution Context:}
\begin{tightitemize}\footnotesize
  \item \textbf{Compute stage preloads (\hyperref[sec:csf-spd]{\texttt{SPD\_k}}, \texttt{r16..r23}):} Hardware pre-fills \texttt{r55..r57} with \texttt{WorkGroupID.xyz + JOB\_OFFSET} (from \hyperref[sec:csf-regs]{\texttt{COMPUTESR::JOB\_OFFSET\_X/Y/Z}} \texttt{r34..r36}) and \texttt{r60..r62} with \texttt{LocalInvocationID.xyz} or \texttt{GlobalInvocationID.xyz} (from \texttt{WG\_SIZE} \texttt{r33} \& \texttt{JOB\_SIZE} \texttt{r37..r39}, selected by \texttt{ShaderProgram::preload}) so compute shaders can derive thread/element indices without memory loads.
  \item \textbf{Vertex stage preloads (\hyperref[sec:csf-spd]{\texttt{SPD\_0}/\texttt{SPD\_1}}, \texttt{r16..r19}):} The hardware Index Fetcher reads \hyperref[sec:csf-regs]{\texttt{IDVSSR::INDEX\_BUFFER}} (\texttt{r54:r55}) and pre-fills \texttt{r59, r60} with \texttt{VertexIndex} (\texttt{+ VERTEX\_OFFSET + GLOBAL\_ATTR\_OFF}) and \texttt{r61} with \texttt{InstanceIndex} (\texttt{+ INSTANCE\_OFFSET}) so Position and Varying shaders can pull vertex attributes directly.
  \item \textbf{Inline FAU \& Special Operands (\hyperref[sec:csf-tsd-fau]{\texttt{FAU\_k}}, \texttt{r8..r15}):} Any 8-bit source with \texttt{mode=2} reads zero-latency uniform \texttt{u0..u31} from \texttt{FAU\_k} (\texttt{slots64 = bits[63:56]}); \texttt{mode=3} reads \texttt{SmallConstant} (Table~1.2) or \texttt{FauSpecial} (\texttt{pc}, \texttt{lane\_id}, \texttt{lane\_width}, \texttt{tls\_ptr}, \texttt{wls\_ptr}, \texttt{blend\_descriptor\_0..7}).
\end{tightitemize}
\par\textbf{Worked Examples:}
\par\textbf{Example 1 (32-bit inline constant add):}
\codebox{\mnem{IADD\_IMM.i32} \dst{r0}, \src{r1}, \imm{\#0x100}}
With \texttt{r1 = 0x00000040}: adds 32-bit inline immediate -\textgreater{} \texttt{r0 = 0x00000140} (wraps mod \texttt{2\textasciicircum{}32} on overflow).
\par\textbf{Example 2 (Packed v2i16 per-lane immediate add):}
\codebox{\mnem{IADD\_IMM.v2i16} \dst{r0}, \src{r1}, \imm{\#0x00020001}}
With \texttt{r1 = 0x00100020} (\texttt{[16, 32]}): adds \texttt{1} to lane 0 and \texttt{2} to lane 1 independently -\textgreater{} \texttt{r0 = 0x00120021}.
\par\textbf{Mined SPIR-V $\to$ Valhall ISA Example (Materializing 32-bit integer constants via MOV\_IMM alias --- Compute shader, Mali-G710 (v10)):}
\codebox{\textbf{; SPIR-V (input)}\\OpStore \%out\_x \%uint\_1000\\OpStore \%out\_y \%uint\_4294966296 ; -1000\\OpStore \%out\_z \%uint\_0\\OpStore \%out\_w \%uint\_0\\[3pt]\textbf{; Valhall ISA (malioc output)}\\MOV\_IMM.i32 r2, \#0x0\\MOV.i32 r3, r2\\MOV\_IMM.i32 r1, \#0xfffffc18\\MOV\_IMM.i32 r0, \#0x3e8\\STORE.i128.slot0.end @r0:r1:r2:r3, u0, \#0x0}
Valhall encodes 32-bit immediate loads (\texttt{MOV\_IMM.i32}) as \texttt{IADD\_IMM.i32} with \texttt{src0 = k0} (0). Here four scalar stores to a 16-byte SSBO struct are coalesced into one \texttt{STORE.i128.end} after loading \texttt{\#0x3e8} (\texttt{1000}) and \texttt{\#0xfffffc18} (\texttt{-1000}).
\begin{center}\footnotesize
\begin{tabularx}{\linewidth}{@{}lllX@{}}
\toprule
Variant & Arch & Class & Opcode Fields \\
\midrule
\texttt{(none)} & v9--v14 & \texttt{R1I} & \texttt{opcode1=ADD\_IMM\_I32} \\
\texttt{(none)} & v9--v14 & \texttt{R1I} & \texttt{opcode1=ADD\_IMM\_V2I16} \\
\texttt{(none)} & v9--v14 & \texttt{R1I} & \texttt{opcode1=ADD\_IMM\_V4I8} \\
\bottomrule
\end{tabularx}
\end{center}
\begin{center}\footnotesize
\begin{tabularx}{0.88\linewidth}{@{}l l X@{}}
\toprule
Field Name & Bits & Type / Fixed Value \\
\midrule
\texttt{src0} & \texttt{0:7} & \texttt{SourceEncoding} \\
\texttt{imm1w} & \texttt{8:39} & \texttt{uint} \\
\texttt{dest\_register\_index} & \texttt{40:45} & \texttt{uint} \\
\hyperref[enum:dest_width_m]{\texttt{dest\_width}} & \texttt{46:47} & \texttt{dest\_width\_m} \\
\texttt{opcode1} & \texttt{48:51} & \texttt{exact=ADD\_IMM\_I32} \\
\texttt{instruction\_class} & \texttt{52:56} & \texttt{exact=R1I} \\
\texttt{fau\_page\_index} & \texttt{57:58} & \texttt{uint} \\
\hyperref[enum:flow_control_m]{\texttt{flow}} & \texttt{59:62} & \texttt{flow\_control\_m} \\
\textit{Reserved} & \texttt{63} & \texttt{zero} \\
\bottomrule
\end{tabularx}
\end{center}
\textbf{Fields Explained:}
\begin{tightitemize}\footnotesize
  \item \texttt{src0}: Primary source operand (\texttt{r0..r63}, \texttt{u0..u63} uniform, \texttt{k*} FAU constant, or \texttt{t0..t2} temporary; \texttt{\textasciicircum{}} marks discard/last-use). \emph{Example:} {\scriptsize\ttfamily \mnem{FADD.f32} \dst{r0}, \src{r1}, \src{r2}} (\textbf{SPIR-V:} \texttt{OpFAdd (Operand 0)}; \textbf{GLSL:} \texttt{float c = a + b;})
  \item \texttt{imm1w}: 32-bit inline immediate operand encoded directly in the instruction word (bits 8:39). \emph{Example:} {\scriptsize\ttfamily \mnem{IADD\_IMM.i32} \dst{r0}, \src{r1}, \imm{\#0x10}} (\textbf{SPIR-V:} \texttt{OpIAdd (constant)}; \textbf{GLSL:} \texttt{int y = x + 16;})
  \item \texttt{dest\_register\_index}: 6-bit destination GPR index (\texttt{r0..r63}) receiving the Stage C writeback result (even-aligned pair \texttt{[rN:rN+1]} for 64-bit ops). \emph{Example:} {\scriptsize\ttfamily \mnem{FMA.f32} \dst{r0}, \src{r1}, \src{r2}, \src{r3}} (\textbf{SPIR-V:} \texttt{SSA Result ID}; \textbf{GLSL:} \texttt{float r0 = fma(r1, r2, r3);})
  \item \hyperref[enum:dest_width_m]{\texttt{dest\_width}}: Stage C destination lane write-mask (\texttt{none} = full 32-bit, \texttt{h0} = low 16-bit, \texttt{h1} = high 16-bit without bit-shifting). \emph{Example:} {\scriptsize\ttfamily \mnem{FADD.v2f16} \dst{r0}\imod{.h1}, \src{r1}, \src{r2}} (\textbf{SPIR-V:} \texttt{OpVectorInsertDynamic / WriteMask}; \textbf{GLSL:} \texttt{dst.y = a.y + b.y;})
  \item \texttt{opcode1}: Primary sub-opcode within the instruction class uniquely identifying the operation family. \emph{Example:} {\scriptsize\ttfamily \mnem{IADD.s32} \dst{r0}, \src{r1}, \src{r2}} (\textbf{SPIR-V:} \texttt{OpIAdd (operation selector)}; \textbf{GLSL:} \texttt{// Primary opcode bits})
  \item \texttt{instruction\_class}: 4-bit major instruction class selector (\texttt{bits 48:51}) partitioning the 64-bit Valhall opcode space (e.g., \texttt{FMA}, \texttt{CVT}, \texttt{SFU}, \texttt{LS}, \texttt{TEX}, \texttt{BRANCH}). \emph{Example:} {\scriptsize\ttfamily \mnem{FADD.f32} \dst{r0}, \src{r1}, \src{r2}} (\textbf{SPIR-V:} \texttt{N/A (encoding class)}; \textbf{GLSL:} \texttt{// Hardware execution pipe / decode class})
  \item \texttt{fau\_page\_index}: Selects the FAU (Fast Access Uniform) page (\texttt{0..3}) when a source operand references hardware small constants, thread IDs, or lane constants. \emph{Example:} {\scriptsize\ttfamily \mnem{FADD.f32} \dst{r0}, \src{r1}, \src{k1f}} (\textbf{SPIR-V:} \texttt{OpConstant / BuiltIn}; \textbf{GLSL:} \texttt{float y = x + 1.0;})
  \item \hyperref[enum:flow_control_m]{\texttt{flow}}: 4-bit embedded flow-control \& scoreboard wait modifier (\texttt{bits 59:62}: \texttt{none}, \texttt{wait0..wait0126}, \texttt{wait}, \texttt{reconverge}, \texttt{discard}, \texttt{end}). \emph{Example:} {\scriptsize\ttfamily \mnem{FADD.f32}\imod{.wait0.reconverge} \dst{r0}, \src{r1}, \src{r2}} (\textbf{SPIR-V:} \texttt{OpControlBarrier / OpReturn}; \textbf{GLSL:} \texttt{// Scoreboard wait + warp reconvergence})
\end{tightitemize}
\clearpage

\subsection*{IDPADD}
\addcontentsline{toc}{subsection}{IDPADD}
Computes the 4-element 8-bit integer dot product of bytes $k \in \{0, 1, 2, 3\}$ of \texttt{src0} and \texttt{src1} and accumulates into 32-bit \texttt{src2}: \texttt{dot = sum\_\{k=0..3\}(ext(src0.b[k]) * ext(src1.b[k]))}; writes \texttt{sat(src2 + dot)} if \texttt{.sat} is set, or wrapping \texttt{src2 + dot} otherwise.
\par\textbf{Notes:}
\begin{tightitemize}
  \item On \texttt{v4s8} (v9-v13), both inputs are signed (\texttt{V4S8}), dst is \texttt{S32}.
  \item On v14, \texttt{unsigned0}/\texttt{unsigned1} override per-src signedness while dst remains \texttt{S32}.
  \item On \texttt{v4u8}, all unsigned.
  \item \textbf{Source references:} \texttt{kraid/ops.rs:2354-2423}, \texttt{encode\_v9.rs:1820-1865}, \texttt{hw\_tests.rs:1952-1989}.
\end{tightitemize}
\codebox{\mnem{IDPADD}\imod{\{.flow\}}\imod{\{.saturate\}} \dst{dest}\imod{\{.dest\_width\}}, \src{src0}, \src{src1}, \src{src2}\\\mnem{IDPADD}\imod{\{.flow\}}\imod{\{.saturate\}} \dst{dest}\imod{\{.dest\_width\}}, \src{src0}\imod{\{.unsigned0\}}, \src{src1}\imod{\{.unsigned1\}}, \src{src2}}
\textbf{Operands:}
\begin{tightitemize}\footnotesize
  \item \texttt{\dst{dest}} (\texttt{dest\_register\_index}, \texttt{dest\_width}): 32-bit destination GPR receiving \texttt{src2 + dot4x8(src0, src1)}.
  \item \texttt{\src{src0}}: Packed 4x8-bit integer vector A (\texttt{b0..b3}).
  \item \texttt{\src{src1}}: Packed 4x8-bit integer vector B (\texttt{b0..b3}).
  \item \texttt{\src{src2}}: 32-bit accumulator addend.
\end{tightitemize}
\par\textbf{Closest SPIR-V Op:} \texttt{OpSDotAccSat / OpUDotAccSat}
\par\textbf{Pseudo-GLSL (Operational Semantics):}
\codebox{int64\_t acc = int64\_t(src2);\\for (int k = 0; k \textless{} 4; ++k) acc += ext(src0.b[k]) * ext(src1.b[k]);\\dest = flag\_sat ? clamp(acc, MIN\_32, MAX\_32) : wrap\_32(acc);}
\textbf{Execution Context:}
\begin{tightitemize}\footnotesize
  \item \textbf{Compute stage preloads (\hyperref[sec:csf-spd]{\texttt{SPD\_k}}, \texttt{r16..r23}):} Hardware pre-fills \texttt{r55..r57} with \texttt{WorkGroupID.xyz + JOB\_OFFSET} (from \hyperref[sec:csf-regs]{\texttt{COMPUTESR::JOB\_OFFSET\_X/Y/Z}} \texttt{r34..r36}) and \texttt{r60..r62} with \texttt{LocalInvocationID.xyz} or \texttt{GlobalInvocationID.xyz} (from \texttt{WG\_SIZE} \texttt{r33} \& \texttt{JOB\_SIZE} \texttt{r37..r39}, selected by \texttt{ShaderProgram::preload}) so compute shaders can derive thread/element indices without memory loads.
  \item \textbf{Vertex stage preloads (\hyperref[sec:csf-spd]{\texttt{SPD\_0}/\texttt{SPD\_1}}, \texttt{r16..r19}):} The hardware Index Fetcher reads \hyperref[sec:csf-regs]{\texttt{IDVSSR::INDEX\_BUFFER}} (\texttt{r54:r55}) and pre-fills \texttt{r59, r60} with \texttt{VertexIndex} (\texttt{+ VERTEX\_OFFSET + GLOBAL\_ATTR\_OFF}) and \texttt{r61} with \texttt{InstanceIndex} (\texttt{+ INSTANCE\_OFFSET}) so Position and Varying shaders can pull vertex attributes directly.
  \item \textbf{Inline FAU \& Special Operands (\hyperref[sec:csf-tsd-fau]{\texttt{FAU\_k}}, \texttt{r8..r15}):} Any 8-bit source with \texttt{mode=2} reads zero-latency uniform \texttt{u0..u31} from \texttt{FAU\_k} (\texttt{slots64 = bits[63:56]}); \texttt{mode=3} reads \texttt{SmallConstant} (Table~1.2) or \texttt{FauSpecial} (\texttt{pc}, \texttt{lane\_id}, \texttt{lane\_width}, \texttt{tls\_ptr}, \texttt{wls\_ptr}, \texttt{blend\_descriptor\_0..7}).
\end{tightitemize}
\par\textbf{Worked Examples:}
\par\textbf{Example 1 (Signed 4x8 dot product with saturation):}
\codebox{\mnem{IDPADD.v4s8}\imod{.sat} \dst{r0}, \src{r1}, \src{r2}, \src{r3}}
With \texttt{r1 = 0x01020304}, \texttt{r2 = 0x01010101}, \texttt{r3 = 10}: dot \texttt{1*1 + 2*1 + 3*1 + 4*1 = 10}, accumulated with \texttt{r3} -\textgreater{} \texttt{r0 = 20}.
\par\textbf{Example 2 (Negative signed lanes \& saturation clamp):}
\codebox{\mnem{IDPADD.v4s8}\imod{.sat} \dst{r0}, \src{r1}, \src{r2}, \src{r3}}
With \texttt{r1 = 0x7F7F7F7F} (\texttt{[127,127,127,127]}), \texttt{r2 = 0x80808080} (\texttt{[-128,...]}), \texttt{r3 = 0x80000000} (\texttt{INT32\_MIN}): dot is \texttt{-65024}; \texttt{.sat} clamps underflow to \texttt{r0 = 0x80000000}.

\begin{center}\footnotesize
\begin{tabularx}{\linewidth}{@{}lllX@{}}
\toprule
Variant & Arch & Class & Opcode Fields \\
\midrule
\texttt{(none)} & v9--v14 & \texttt{R3A} & \texttt{opcode1=IDPADD\_V4I8} \\
\texttt{(none)} & v9--v14 & \texttt{R3A} & \texttt{opcode1=IDPADD\_V4I8} \\
\texttt{(none)} & v9--v14 & \texttt{R3A} & \texttt{opcode1=IDPADD\_V4I8} \\
\bottomrule
\end{tabularx}
\end{center}
\begin{center}\footnotesize
\begin{tabularx}{0.88\linewidth}{@{}l l X@{}}
\toprule
Field Name & Bits & Type / Fixed Value \\
\midrule
\texttt{src0} & \texttt{0:7} & \texttt{SourceEncoding} \\
\texttt{src1} & \texttt{8:15} & \texttt{SourceEncoding} \\
\texttt{src2} & \texttt{16:23} & \texttt{SourceEncoding} \\
\texttt{signedness} & \texttt{24} & \texttt{exact=op\_signed} \\
\textit{Reserved} & \texttt{25:29} & \texttt{zero} \\
\hyperref[enum:saturate_m]{\texttt{saturate}} & \texttt{30} & \texttt{saturate\_m (def none)} \\
\textit{Reserved} & \texttt{31:39} & \texttt{zero} \\
\texttt{dest\_register\_index} & \texttt{40:45} & \texttt{uint} \\
\hyperref[enum:dest_width_m]{\texttt{dest\_width}} & \texttt{46:47} & \texttt{dest\_width\_m} \\
\texttt{opcode1} & \texttt{48:51} & \texttt{exact=IDPADD\_V4I8} \\
\texttt{instruction\_class} & \texttt{52:56} & \texttt{exact=R3A} \\
\texttt{fau\_page\_index} & \texttt{57:58} & \texttt{uint} \\
\hyperref[enum:flow_control_m]{\texttt{flow}} & \texttt{59:62} & \texttt{flow\_control\_m} \\
\textit{Reserved} & \texttt{63} & \texttt{zero} \\
\bottomrule
\end{tabularx}
\end{center}
\textbf{Fields Explained:}
\begin{tightitemize}\footnotesize
  \item \texttt{src0}: Primary source operand (\texttt{r0..r63}, \texttt{u0..u63} uniform, \texttt{k*} FAU constant, or \texttt{t0..t2} temporary; \texttt{\textasciicircum{}} marks discard/last-use). \emph{Example:} {\scriptsize\ttfamily \mnem{FADD.f32} \dst{r0}, \src{r1}, \src{r2}} (\textbf{SPIR-V:} \texttt{OpFAdd (Operand 0)}; \textbf{GLSL:} \texttt{float c = a + b;})
  \item \texttt{src1}: Secondary source operand (register, uniform, or FAU constant; in shifts, supplies the shift count byte). \emph{Example:} {\scriptsize\ttfamily \mnem{LSHIFT.i32} \dst{r0}, \src{r1}, \src{r2}\imod{.b0}} (\textbf{SPIR-V:} \texttt{OpShiftLeftLogical (Shift)}; \textbf{GLSL:} \texttt{uint y = a \textless{}\textless{} s;})
  \item \texttt{src2}: Third source operand (FMA addend, \texttt{MUX}/\texttt{CSEL} selector, \texttt{IADDC} carry-in, or bitwise mask in fused shift-logic). \emph{Example:} {\scriptsize\ttfamily \mnem{FMA.f32} \dst{r0}, \src{r1}, \src{r2}, \src{r3}} (\textbf{SPIR-V:} \texttt{GLSL.std.450 Fma (Addend)}; \textbf{GLSL:} \texttt{float d = fma(a, b, c);})
  \item \texttt{signedness}: Selects signed vs unsigned interpretation for integer memory/comparison/conversion. \emph{Example:} {\scriptsize\ttfamily \mnem{LD\_ATTR}\imod{.s32.v4} \dst{@r0:r3}, \src{r4}} (\textbf{SPIR-V:} \texttt{OpLoad (ivec4)}; \textbf{GLSL:} \texttt{in ivec4 a\_idx;})
  \item \hyperref[enum:saturate_m]{\texttt{saturate}}: Stage B integer saturation flag (\texttt{.sat}): clamps signed/unsigned overflow or underflow to \texttt{[MIN\_N, MAX\_N]} / \texttt{[0, 2\textasciicircum{}N - 1]} instead of wrapping. \emph{Example:} {\scriptsize\ttfamily \mnem{IADD.s32}\imod{.sat} \dst{r0}, \src{r1}, \src{r2}} (\textbf{SPIR-V:} \texttt{OpSDotAccSat / saturating add}; \textbf{GLSL:} \texttt{int c = clamp(int64\_t(a) + int64\_t(b), INT\_MIN, INT\_MAX);})
  \item \texttt{dest\_register\_index}: 6-bit destination GPR index (\texttt{r0..r63}) receiving the Stage C writeback result (even-aligned pair \texttt{[rN:rN+1]} for 64-bit ops). \emph{Example:} {\scriptsize\ttfamily \mnem{FMA.f32} \dst{r0}, \src{r1}, \src{r2}, \src{r3}} (\textbf{SPIR-V:} \texttt{SSA Result ID}; \textbf{GLSL:} \texttt{float r0 = fma(r1, r2, r3);})
  \item \hyperref[enum:dest_width_m]{\texttt{dest\_width}}: Stage C destination lane write-mask (\texttt{none} = full 32-bit, \texttt{h0} = low 16-bit, \texttt{h1} = high 16-bit without bit-shifting). \emph{Example:} {\scriptsize\ttfamily \mnem{FADD.v2f16} \dst{r0}\imod{.h1}, \src{r1}, \src{r2}} (\textbf{SPIR-V:} \texttt{OpVectorInsertDynamic / WriteMask}; \textbf{GLSL:} \texttt{dst.y = a.y + b.y;})
  \item \texttt{opcode1}: Primary sub-opcode within the instruction class uniquely identifying the operation family. \emph{Example:} {\scriptsize\ttfamily \mnem{IADD.s32} \dst{r0}, \src{r1}, \src{r2}} (\textbf{SPIR-V:} \texttt{OpIAdd (operation selector)}; \textbf{GLSL:} \texttt{// Primary opcode bits})
  \item \texttt{instruction\_class}: 4-bit major instruction class selector (\texttt{bits 48:51}) partitioning the 64-bit Valhall opcode space (e.g., \texttt{FMA}, \texttt{CVT}, \texttt{SFU}, \texttt{LS}, \texttt{TEX}, \texttt{BRANCH}). \emph{Example:} {\scriptsize\ttfamily \mnem{FADD.f32} \dst{r0}, \src{r1}, \src{r2}} (\textbf{SPIR-V:} \texttt{N/A (encoding class)}; \textbf{GLSL:} \texttt{// Hardware execution pipe / decode class})
  \item \texttt{fau\_page\_index}: Selects the FAU (Fast Access Uniform) page (\texttt{0..3}) when a source operand references hardware small constants, thread IDs, or lane constants. \emph{Example:} {\scriptsize\ttfamily \mnem{FADD.f32} \dst{r0}, \src{r1}, \src{k1f}} (\textbf{SPIR-V:} \texttt{OpConstant / BuiltIn}; \textbf{GLSL:} \texttt{float y = x + 1.0;})
  \item \hyperref[enum:flow_control_m]{\texttt{flow}}: 4-bit embedded flow-control \& scoreboard wait modifier (\texttt{bits 59:62}: \texttt{none}, \texttt{wait0..wait0126}, \texttt{wait}, \texttt{reconverge}, \texttt{discard}, \texttt{end}). \emph{Example:} {\scriptsize\ttfamily \mnem{FADD.f32}\imod{.wait0.reconverge} \dst{r0}, \src{r1}, \src{r2}} (\textbf{SPIR-V:} \texttt{OpControlBarrier / OpReturn}; \textbf{GLSL:} \texttt{// Scoreboard wait + warp reconvergence})
\end{tightitemize}
\clearpage

\subsection*{IMUL}
\addcontentsline{toc}{subsection}{IMUL}
Multiplies Stage A widened/swizzled \texttt{src0} and \texttt{src1} independently for each SIMD element $i \in \{0, \dots, N-1\}$ ($N=1$ for \texttt{s32}/\texttt{u32}/\texttt{s64}/\texttt{u64}, $N=2$ for \texttt{v2s16}/\texttt{v2u16}, $N=4$ for \texttt{v4s8}/\texttt{v4u8} on v9--v10): \texttt{dest[i] = sat(widen0(src0)[i] * widen1(src1)[i])} under \texttt{.sat}, or low $W$ bits otherwise. For 64-bit \texttt{s64}/\texttt{u64}, performs a $32 \times 32 \to 64$ widening multiply requiring \texttt{.w0}/\texttt{.w1} on both sources.
\par\textbf{Notes:}
\begin{tightitemize}
  \item 64-bit \texttt{IMUL.sat} NOT supported (\texttt{ops.rs:2465-2468}).
  \item \texttt{IMUL.s64/u64} requires widening swizzle (\texttt{None} swizzle unsupported, \texttt{hw\_tests.rs:2022-2027}).
  \item \textbf{Source references:} \texttt{kraid/ops.rs:2433-2483}, \texttt{encode\_v9.rs:1867-1896}, \texttt{nir.rs:1106-1136}, \texttt{hw\_tests.rs:1992-2060}.
\end{tightitemize}
\codebox{\mnem{IMUL}\imod{\{.flow\}}\imod{\{.saturate\}} \dst{dest}\imod{\{.dest\_width\}}, \src{src0}\imod{\{.widen0\}}, \src{src1}\imod{\{.widen1\}}\\\mnem{IMUL}\imod{\{.flow\}} \dst{dest}, \src{src0}\imod{\{.widen0\}}, \src{src1}\imod{\{.widen1\}}\\\mnem{IMUL}\imod{\{.flow\}}\imod{\{.saturate\}} \dst{dest}\imod{\{.dest\_width\}}, \src{src0}\imod{\{.swz0\}}, \src{src1}\imod{\{.swz1\}}\\\mnem{IMUL}\imod{\{.flow\}}\imod{\{.saturate\}} \dst{dest}\imod{\{.dest\_width\}}, \src{src0}\imod{\{.lanes0\}}, \src{src1}\imod{\{.lanes1\}}}
\textbf{Operands:}
\begin{tightitemize}\footnotesize
  \item \texttt{\dst{dest}} (\texttt{dest\_register\_index}, \texttt{dest\_width}): Destination GPR (or 64-bit pair for \texttt{s64}/\texttt{u64} widening multiply) receiving the product.
  \item \texttt{\src{src0}}: First multiplicand (must use \texttt{.w0}/\texttt{.w1} for \texttt{s64}/\texttt{u64}).
  \item \texttt{\src{src1}}: Second multiplicand (must use \texttt{.w0}/\texttt{.w1} for \texttt{s64}/\texttt{u64}).
\end{tightitemize}
\par\textbf{Closest SPIR-V Op:} \texttt{OpIMul / OpUMulExtended / OpSMulExtended}
\par\textbf{Pseudo-GLSL (Operational Semantics):}
\codebox{if (is\_64bit) dest = ext64(widen0(src0)) * ext64(widen1(src1)); // 32x32-\textgreater{}64\\else for (int i = 0; i \textless{} N; ++i) \{\\    auto prod = ext(widen0(src0)[i]) * ext(widen1(src1)[i]);\\    dest[i] = flag\_sat ? clamp(prod, MIN\_W, MAX\_W) : wrap\_W(prod);\\\}}
\textbf{Execution Context:}
\begin{tightitemize}\footnotesize
  \item \textbf{Compute stage preloads (\hyperref[sec:csf-spd]{\texttt{SPD\_k}}, \texttt{r16..r23}):} Hardware pre-fills \texttt{r55..r57} with \texttt{WorkGroupID.xyz + JOB\_OFFSET} (from \hyperref[sec:csf-regs]{\texttt{COMPUTESR::JOB\_OFFSET\_X/Y/Z}} \texttt{r34..r36}) and \texttt{r60..r62} with \texttt{LocalInvocationID.xyz} or \texttt{GlobalInvocationID.xyz} (from \texttt{WG\_SIZE} \texttt{r33} \& \texttt{JOB\_SIZE} \texttt{r37..r39}, selected by \texttt{ShaderProgram::preload}) so compute shaders can derive thread/element indices without memory loads.
  \item \textbf{Vertex stage preloads (\hyperref[sec:csf-spd]{\texttt{SPD\_0}/\texttt{SPD\_1}}, \texttt{r16..r19}):} The hardware Index Fetcher reads \hyperref[sec:csf-regs]{\texttt{IDVSSR::INDEX\_BUFFER}} (\texttt{r54:r55}) and pre-fills \texttt{r59, r60} with \texttt{VertexIndex} (\texttt{+ VERTEX\_OFFSET + GLOBAL\_ATTR\_OFF}) and \texttt{r61} with \texttt{InstanceIndex} (\texttt{+ INSTANCE\_OFFSET}) so Position and Varying shaders can pull vertex attributes directly.
  \item \textbf{Inline FAU \& Special Operands (\hyperref[sec:csf-tsd-fau]{\texttt{FAU\_k}}, \texttt{r8..r15}):} Any 8-bit source with \texttt{mode=2} reads zero-latency uniform \texttt{u0..u31} from \texttt{FAU\_k} (\texttt{slots64 = bits[63:56]}); \texttt{mode=3} reads \texttt{SmallConstant} (Table~1.2) or \texttt{FauSpecial} (\texttt{pc}, \texttt{lane\_id}, \texttt{lane\_width}, \texttt{tls\_ptr}, \texttt{wls\_ptr}, \texttt{blend\_descriptor\_0..7}).
\end{tightitemize}
\par\textbf{Worked Examples:}
\par\textbf{Example 1 (32x32 -\textgreater{} 64-bit widening unsigned multiply):}
\codebox{\mnem{IMUL.u64} \dst{[r0:r1]}, \src{[r2:r3].w0}, \src{[r4:r5].w0}}
With \texttt{r2 = 0x80000000, r4 = 4}: Stage A zero-extends \texttt{.w0} to 64-bit; Stage B computes \texttt{0x00000002\_00000000} in \texttt{[r0:r1]} (\texttt{r0 = 0}, \texttt{r1 = 2}).
\par\textbf{Example 2 (32-bit saturating signed multiply):}
\codebox{\mnem{IMUL.s32}\imod{.sat} \dst{r0}, \src{r1}, \src{r2}}
With \texttt{r1 = 0x00010000} (\texttt{65536}), \texttt{r2 = 0x00010000} (\texttt{65536}): exact product \texttt{2\textasciicircum{}32} overflows \texttt{s32}; \texttt{.sat} clamps to \texttt{r0 = 0x7FFFFFFF}.
\par\textbf{Mined SPIR-V $\to$ Valhall ISA Example (Struct stride multiplication using FAU byte constant --- Vertex shader, Mali-G710 (v10)):}
\codebox{\textbf{; SPIR-V (input)}\\\%sum = OpIAdd \%uint \%base \%attr\_idx\\\%stride\_off = OpIMul \%uint \%sum \%uint\_16\\[3pt]\textbf{; Valhall ISA (malioc output)}\\IADD.u32 r0, u3.w0, r0\textasciicircum{}\\IMUL.u32 r0, r0\textasciicircum{}, k13121110.b0}
\texttt{OpIMul} lowers to \texttt{IMUL.u32} with inline FAU small constant \texttt{k13121110.b0} (byte 0 = \texttt{0x10} = 16), avoiding a separate immediate-load instruction.
\begin{center}\footnotesize
\begin{tabularx}{\linewidth}{@{}lllX@{}}
\toprule
Variant & Arch & Class & Opcode Fields \\
\midrule
\texttt{s32} & v9--v14 & \texttt{R2} & \texttt{opcode2=MUL} \\
\texttt{s64} & v9--v14 & \texttt{D2} & \texttt{opcode2=MUL} \\
\texttt{u32} & v9--v14 & \texttt{R2} & \texttt{opcode2=MUL} \\
\texttt{u64} & v9--v14 & \texttt{D2} & \texttt{opcode2=MUL} \\
\texttt{v2s16} & v9--v14 & \texttt{R2} & \texttt{opcode2=MUL} \\
\texttt{v2u16} & v9--v14 & \texttt{R2} & \texttt{opcode2=MUL} \\
\texttt{v4s8} & v9--v14 & \texttt{R2} & \texttt{opcode2=MUL} \\
\texttt{v4u8} & v9--v14 & \texttt{R2} & \texttt{opcode2=MUL} \\
\bottomrule
\end{tabularx}
\end{center}
\begin{center}\footnotesize
\begin{tabularx}{0.88\linewidth}{@{}l l X@{}}
\toprule
Field Name & Bits & Type / Fixed Value \\
\midrule
\texttt{src0} & \texttt{0:7} & \texttt{SourceEncoding} \\
\texttt{src1} & \texttt{8:15} & \texttt{SourceEncoding} \\
\texttt{opcode2} & \texttt{16:19} & \texttt{exact=MUL} \\
\textit{Reserved} & \texttt{20:25} & \texttt{zero} \\
\hyperref[enum:widen_int_m]{\texttt{widen1}} & \texttt{26:29} & \texttt{widen\_int\_m (def none)} \\
\hyperref[enum:saturate_m]{\texttt{saturate}} & \texttt{30} & \texttt{saturate\_m (def none)} \\
\textit{Reserved} & \texttt{31:35} & \texttt{zero} \\
\hyperref[enum:widen_int_m]{\texttt{widen0}} & \texttt{36:39} & \texttt{widen\_int\_m (def none)} \\
\texttt{dest\_register\_index} & \texttt{40:45} & \texttt{uint} \\
\hyperref[enum:dest_width_m]{\texttt{dest\_width}} & \texttt{46:47} & \texttt{dest\_width\_m} \\
\texttt{datatype} & \texttt{48:51} & \texttt{exact=s32} \\
\texttt{instruction\_class} & \texttt{52:56} & \texttt{exact=R2} \\
\texttt{fau\_page\_index} & \texttt{57:58} & \texttt{uint} \\
\hyperref[enum:flow_control_m]{\texttt{flow}} & \texttt{59:62} & \texttt{flow\_control\_m} \\
\textit{Reserved} & \texttt{63} & \texttt{zero} \\
\bottomrule
\end{tabularx}
\end{center}
\textbf{Fields Explained:}
\begin{tightitemize}\footnotesize
  \item \texttt{src0}: Primary source operand (\texttt{r0..r63}, \texttt{u0..u63} uniform, \texttt{k*} FAU constant, or \texttt{t0..t2} temporary; \texttt{\textasciicircum{}} marks discard/last-use). \emph{Example:} {\scriptsize\ttfamily \mnem{FADD.f32} \dst{r0}, \src{r1}, \src{r2}} (\textbf{SPIR-V:} \texttt{OpFAdd (Operand 0)}; \textbf{GLSL:} \texttt{float c = a + b;})
  \item \texttt{src1}: Secondary source operand (register, uniform, or FAU constant; in shifts, supplies the shift count byte). \emph{Example:} {\scriptsize\ttfamily \mnem{LSHIFT.i32} \dst{r0}, \src{r1}, \src{r2}\imod{.b0}} (\textbf{SPIR-V:} \texttt{OpShiftLeftLogical (Shift)}; \textbf{GLSL:} \texttt{uint y = a \textless{}\textless{} s;})
  \item \texttt{opcode2}: Secondary sub-opcode further specializing instructions sharing \texttt{instruction\_class} and \texttt{opcode1} (e.g., conversions, SFU tables, moves). \emph{Example:} {\scriptsize\ttfamily \mnem{F32\_TO\_S32}\imod{.round\_zero} \dst{r0}, \src{r1}} (\textbf{SPIR-V:} \texttt{OpConvertFToS}; \textbf{GLSL:} \texttt{int y = int(x);})
  \item \hyperref[enum:widen_int_m]{\texttt{widen1}}: Stage A widening extractor on \texttt{src1} (promotes narrow byte/half/word lane before Stage B). \emph{Example:} {\scriptsize\ttfamily \mnem{IMUL.u64} \dst{[r0:r1]}, \src{[r2:r3]}\imod{.w0}, \src{[r4:r5]}\imod{.w0}} (\textbf{SPIR-V:} \texttt{OpUConvert / OpUMulExtended}; \textbf{GLSL:} \texttt{uint64\_t p = uint64\_t(a) * uint64\_t(b);})
  \item \hyperref[enum:saturate_m]{\texttt{saturate}}: Stage B integer saturation flag (\texttt{.sat}): clamps signed/unsigned overflow or underflow to \texttt{[MIN\_N, MAX\_N]} / \texttt{[0, 2\textasciicircum{}N - 1]} instead of wrapping. \emph{Example:} {\scriptsize\ttfamily \mnem{IADD.s32}\imod{.sat} \dst{r0}, \src{r1}, \src{r2}} (\textbf{SPIR-V:} \texttt{OpSDotAccSat / saturating add}; \textbf{GLSL:} \texttt{int c = clamp(int64\_t(a) + int64\_t(b), INT\_MIN, INT\_MAX);})
  \item \hyperref[enum:widen_int_m]{\texttt{widen0}}: Stage A widening extractor on \texttt{src0} (promotes FP16 \texttt{h0}/\texttt{h1} to FP32, or sign/zero-extends \texttt{b0..b3}/\texttt{h0..h1}/\texttt{w0..w1} to 32/64-bit before \texttt{abs}/\texttt{neg}). \emph{Example:} {\scriptsize\ttfamily \mnem{IADD.s32} \dst{r0}, \src{r1}\imod{.h0}, \src{r2}} (\textbf{SPIR-V:} \texttt{OpSConvert / OpFConvert}; \textbf{GLSL:} \texttt{int c = int(int16\_t(a)) + b;})
  \item \texttt{dest\_register\_index}: 6-bit destination GPR index (\texttt{r0..r63}) receiving the Stage C writeback result (even-aligned pair \texttt{[rN:rN+1]} for 64-bit ops). \emph{Example:} {\scriptsize\ttfamily \mnem{FMA.f32} \dst{r0}, \src{r1}, \src{r2}, \src{r3}} (\textbf{SPIR-V:} \texttt{SSA Result ID}; \textbf{GLSL:} \texttt{float r0 = fma(r1, r2, r3);})
  \item \hyperref[enum:dest_width_m]{\texttt{dest\_width}}: Stage C destination lane write-mask (\texttt{none} = full 32-bit, \texttt{h0} = low 16-bit, \texttt{h1} = high 16-bit without bit-shifting). \emph{Example:} {\scriptsize\ttfamily \mnem{FADD.v2f16} \dst{r0}\imod{.h1}, \src{r1}, \src{r2}} (\textbf{SPIR-V:} \texttt{OpVectorInsertDynamic / WriteMask}; \textbf{GLSL:} \texttt{dst.y = a.y + b.y;})
  \item \texttt{datatype}: Selects the lane data type (\texttt{f32}, \texttt{f16}, \texttt{v2f16}, \texttt{s32}, \texttt{u32}, \texttt{v2s16}, \texttt{v2u16}, \texttt{v4s8}, \texttt{v4u8}, \texttt{s64}, \texttt{u64}, etc.). \emph{Example:} {\scriptsize\ttfamily \mnem{IADD.v2s16} \dst{r0}, \src{r1}, \src{r2}} (\textbf{SPIR-V:} \texttt{OpTypeInt / OpTypeFloat}; \textbf{GLSL:} \texttt{i16vec2 c = a + b;})
  \item \texttt{instruction\_class}: 4-bit major instruction class selector (\texttt{bits 48:51}) partitioning the 64-bit Valhall opcode space (e.g., \texttt{FMA}, \texttt{CVT}, \texttt{SFU}, \texttt{LS}, \texttt{TEX}, \texttt{BRANCH}). \emph{Example:} {\scriptsize\ttfamily \mnem{FADD.f32} \dst{r0}, \src{r1}, \src{r2}} (\textbf{SPIR-V:} \texttt{N/A (encoding class)}; \textbf{GLSL:} \texttt{// Hardware execution pipe / decode class})
  \item \texttt{fau\_page\_index}: Selects the FAU (Fast Access Uniform) page (\texttt{0..3}) when a source operand references hardware small constants, thread IDs, or lane constants. \emph{Example:} {\scriptsize\ttfamily \mnem{FADD.f32} \dst{r0}, \src{r1}, \src{k1f}} (\textbf{SPIR-V:} \texttt{OpConstant / BuiltIn}; \textbf{GLSL:} \texttt{float y = x + 1.0;})
  \item \hyperref[enum:flow_control_m]{\texttt{flow}}: 4-bit embedded flow-control \& scoreboard wait modifier (\texttt{bits 59:62}: \texttt{none}, \texttt{wait0..wait0126}, \texttt{wait}, \texttt{reconverge}, \texttt{discard}, \texttt{end}). \emph{Example:} {\scriptsize\ttfamily \mnem{FADD.f32}\imod{.wait0.reconverge} \dst{r0}, \src{r1}, \src{r2}} (\textbf{SPIR-V:} \texttt{OpControlBarrier / OpReturn}; \textbf{GLSL:} \texttt{// Scoreboard wait + warp reconvergence})
\end{tightitemize}
\clearpage

\subsection*{ISUB}
\addcontentsline{toc}{subsection}{ISUB}
Subtracts Stage A widened/swizzled \texttt{src1} from \texttt{src0} independently for each SIMD element $i \in \{0, \dots, N-1\}$ ($N=1$ for \texttt{s32}/\texttt{u32}/\texttt{s64}/\texttt{u64}, $N=2$ for \texttt{v2s16}/\texttt{v2u16}, $N=4$ for \texttt{v4s8}/\texttt{v4u8} on v9--v10): \texttt{dest[i] = sat(widen0(src0)[i] - widen1(src1)[i])} when \texttt{.sat} is set, or wraps modulo $2^W$ otherwise.
\par\textbf{Notes:}
\begin{tightitemize}
  \item 64-bit \texttt{ISUB.sat} NOT supported by HW (\texttt{kraid/ops.rs:2523-2526}, \texttt{hw\_tests.rs:2093}).
  \item Unsigned sat clamps underflow to 0.
  \item \texttt{v4s8}/\texttt{v4u8} arch \textless{}= v10.
  \item \textbf{Source references:} \texttt{kraid/ops.rs:2492-2540}, \texttt{encode\_v9.rs:1898-1918}, \texttt{nir.rs:1137-1159}, \texttt{hw\_tests.rs:2063-2110}.
\end{tightitemize}
\codebox{\mnem{ISUB}\imod{\{.flow\}}\imod{\{.saturate\}} \dst{dest}\imod{\{.dest\_width\}}, \src{src0}\imod{\{.widen0\}}, \src{src1}\imod{\{.widen1\}}\\\mnem{ISUB}\imod{\{.flow\}}\imod{\{.saturate\}} \dst{dest}, \src{src0}\imod{\{.widen0\}}, \src{src1}\imod{\{.widen1\}}\\\mnem{ISUB}\imod{\{.flow\}} \dst{dest}, \src{src0}\imod{\{.widen0\}}, \src{src1}\imod{\{.widen1\}}\\\mnem{ISUB}\imod{\{.flow\}}\imod{\{.saturate\}} \dst{dest}\imod{\{.dest\_width\}}, \src{src0}\imod{\{.swz0\}}, \src{src1}\imod{\{.swz1\}}\\\mnem{ISUB}\imod{\{.flow\}}\imod{\{.saturate\}} \dst{dest}\imod{\{.dest\_width\}}, \src{src0}\imod{\{.lanes0\}}, \src{src1}\imod{\{.lanes1\}}}
\textbf{Operands:}
\begin{tightitemize}\footnotesize
  \item \texttt{\dst{dest}} (\texttt{dest\_register\_index}, \texttt{dest\_width}): Destination GPR (or 64-bit pair) receiving \texttt{src0 - src1}.
  \item \texttt{\src{src0}}: Minuend (supports Stage A widening/swizzle).
  \item \texttt{\src{src1}}: Subtrahend (supports Stage A widening/swizzle).
\end{tightitemize}
\par\textbf{Closest SPIR-V Op:} \texttt{OpISub}
\par\textbf{Pseudo-GLSL (Operational Semantics):}
\codebox{for (int i = 0; i \textless{} N; ++i) \{\\    auto diff = ext(widen0(src0)[i]) - ext(widen1(src1)[i]);\\    dest[i] = flag\_sat ? clamp(diff, MIN\_W, MAX\_W) : wrap\_W(diff);\\\}}
\textbf{Execution Context:}
\begin{tightitemize}\footnotesize
  \item \textbf{Compute stage preloads (\hyperref[sec:csf-spd]{\texttt{SPD\_k}}, \texttt{r16..r23}):} Hardware pre-fills \texttt{r55..r57} with \texttt{WorkGroupID.xyz + JOB\_OFFSET} (from \hyperref[sec:csf-regs]{\texttt{COMPUTESR::JOB\_OFFSET\_X/Y/Z}} \texttt{r34..r36}) and \texttt{r60..r62} with \texttt{LocalInvocationID.xyz} or \texttt{GlobalInvocationID.xyz} (from \texttt{WG\_SIZE} \texttt{r33} \& \texttt{JOB\_SIZE} \texttt{r37..r39}, selected by \texttt{ShaderProgram::preload}) so compute shaders can derive thread/element indices without memory loads.
  \item \textbf{Vertex stage preloads (\hyperref[sec:csf-spd]{\texttt{SPD\_0}/\texttt{SPD\_1}}, \texttt{r16..r19}):} The hardware Index Fetcher reads \hyperref[sec:csf-regs]{\texttt{IDVSSR::INDEX\_BUFFER}} (\texttt{r54:r55}) and pre-fills \texttt{r59, r60} with \texttt{VertexIndex} (\texttt{+ VERTEX\_OFFSET + GLOBAL\_ATTR\_OFF}) and \texttt{r61} with \texttt{InstanceIndex} (\texttt{+ INSTANCE\_OFFSET}) so Position and Varying shaders can pull vertex attributes directly.
  \item \textbf{Inline FAU \& Special Operands (\hyperref[sec:csf-tsd-fau]{\texttt{FAU\_k}}, \texttt{r8..r15}):} Any 8-bit source with \texttt{mode=2} reads zero-latency uniform \texttt{u0..u31} from \texttt{FAU\_k} (\texttt{slots64 = bits[63:56]}); \texttt{mode=3} reads \texttt{SmallConstant} (Table~1.2) or \texttt{FauSpecial} (\texttt{pc}, \texttt{lane\_id}, \texttt{lane\_width}, \texttt{tls\_ptr}, \texttt{wls\_ptr}, \texttt{blend\_descriptor\_0..7}).
\end{tightitemize}
\par\textbf{Worked Examples:}
\par\textbf{Example 1 (Unsigned saturating underflow clamp):}
\codebox{\mnem{ISUB.u32}\imod{.sat} \dst{r0}, \src{r1}, \src{r2}}
With \texttt{r1 = 5, r2 = 12}: unsigned underflow clamps to \texttt{r0 = 0} under \texttt{.sat} (without \texttt{.sat}, wraps to \texttt{0xFFFFFFF9}).
\par\textbf{Example 2 (Signed saturating overflow on subtracting negative):}
\codebox{\mnem{ISUB.s32}\imod{.sat} \dst{r0}, \src{r1}, \src{r2}}
With \texttt{r1 = 0x7FFFFFF0, r2 = -32} (\texttt{0xFFFFFFE0}): exact \texttt{0x7FFFFFF0 - (-32) = 0x80000010} overflows \texttt{s32} -\textgreater{} \texttt{.sat} clamps to \texttt{r0 = 0x7FFFFFFF}.
\par\textbf{Mined SPIR-V $\to$ Valhall ISA Example (Subtracting BaseInstance from InstanceIndex --- Vertex shader, Mali-G710 (v10)):}
\codebox{\textbf{; SPIR-V (input)}\\\%inst = OpLoad \%int \%gl\_InstanceIndex\\\%base = OpLoad \%int \%ubo\_base\_inst\\\%rel  = OpISub \%int \%inst \%base\\\%f    = OpConvertSToF \%float \%rel\\[3pt]\textbf{; Valhall ISA (malioc output)}\\LD\_PKA.i32.slot0 @r0, u15.w0, u13.w1\\ISUB.u32 r0, r61\textasciicircum{}, r0\textasciicircum{}\\S32\_TO\_F32 r0, r0\textasciicircum{}}
\texttt{r61} holds hardware-preloaded \texttt{gl\_InstanceIndex}. After \texttt{LD\_PKA} fetches the base offset into \texttt{r0}, \texttt{ISUB.u32} computes the relative instance index (\texttt{r61 - r0}) before converting to float.
\begin{center}\footnotesize
\begin{tabularx}{\linewidth}{@{}lllX@{}}
\toprule
Variant & Arch & Class & Opcode Fields \\
\midrule
\texttt{s32} & v9--v14 & \texttt{R2} & \texttt{opcode2=SUB} \\
\texttt{s64} & v9--v14 & \texttt{D2} & \texttt{opcode2=SUB} \\
\texttt{s64} & v9--v14 & \texttt{D2} & \texttt{opcode2=SUB} \\
\texttt{u32} & v9--v14 & \texttt{R2} & \texttt{opcode2=SUB} \\
\texttt{u64} & v9--v14 & \texttt{D2} & \texttt{opcode2=SUB} \\
\texttt{u64} & v9--v14 & \texttt{D2} & \texttt{opcode2=SUB} \\
\texttt{v2s16} & v9--v14 & \texttt{R2} & \texttt{opcode2=SUB} \\
\texttt{v2u16} & v9--v14 & \texttt{R2} & \texttt{opcode2=SUB} \\
\texttt{v4s8} & v9--v14 & \texttt{R2} & \texttt{opcode2=SUB} \\
\texttt{v4u8} & v9--v14 & \texttt{R2} & \texttt{opcode2=SUB} \\
\bottomrule
\end{tabularx}
\end{center}
\begin{center}\footnotesize
\begin{tabularx}{0.88\linewidth}{@{}l l X@{}}
\toprule
Field Name & Bits & Type / Fixed Value \\
\midrule
\texttt{src0} & \texttt{0:7} & \texttt{SourceEncoding} \\
\texttt{src1} & \texttt{8:15} & \texttt{SourceEncoding} \\
\texttt{opcode2} & \texttt{16:19} & \texttt{exact=SUB} \\
\textit{Reserved} & \texttt{20:25} & \texttt{zero} \\
\hyperref[enum:widen_int_m]{\texttt{widen1}} & \texttt{26:29} & \texttt{widen\_int\_m (def none)} \\
\hyperref[enum:saturate_m]{\texttt{saturate}} & \texttt{30} & \texttt{saturate\_m (def none)} \\
\textit{Reserved} & \texttt{31:35} & \texttt{zero} \\
\hyperref[enum:widen_int_m]{\texttt{widen0}} & \texttt{36:39} & \texttt{widen\_int\_m (def none)} \\
\texttt{dest\_register\_index} & \texttt{40:45} & \texttt{uint} \\
\hyperref[enum:dest_width_m]{\texttt{dest\_width}} & \texttt{46:47} & \texttt{dest\_width\_m} \\
\texttt{datatype} & \texttt{48:51} & \texttt{exact=s32} \\
\texttt{instruction\_class} & \texttt{52:56} & \texttt{exact=R2} \\
\texttt{fau\_page\_index} & \texttt{57:58} & \texttt{uint} \\
\hyperref[enum:flow_control_m]{\texttt{flow}} & \texttt{59:62} & \texttt{flow\_control\_m} \\
\textit{Reserved} & \texttt{63} & \texttt{zero} \\
\bottomrule
\end{tabularx}
\end{center}
\textbf{Fields Explained:}
\begin{tightitemize}\footnotesize
  \item \texttt{src0}: Primary source operand (\texttt{r0..r63}, \texttt{u0..u63} uniform, \texttt{k*} FAU constant, or \texttt{t0..t2} temporary; \texttt{\textasciicircum{}} marks discard/last-use). \emph{Example:} {\scriptsize\ttfamily \mnem{FADD.f32} \dst{r0}, \src{r1}, \src{r2}} (\textbf{SPIR-V:} \texttt{OpFAdd (Operand 0)}; \textbf{GLSL:} \texttt{float c = a + b;})
  \item \texttt{src1}: Secondary source operand (register, uniform, or FAU constant; in shifts, supplies the shift count byte). \emph{Example:} {\scriptsize\ttfamily \mnem{LSHIFT.i32} \dst{r0}, \src{r1}, \src{r2}\imod{.b0}} (\textbf{SPIR-V:} \texttt{OpShiftLeftLogical (Shift)}; \textbf{GLSL:} \texttt{uint y = a \textless{}\textless{} s;})
  \item \texttt{opcode2}: Secondary sub-opcode further specializing instructions sharing \texttt{instruction\_class} and \texttt{opcode1} (e.g., conversions, SFU tables, moves). \emph{Example:} {\scriptsize\ttfamily \mnem{F32\_TO\_S32}\imod{.round\_zero} \dst{r0}, \src{r1}} (\textbf{SPIR-V:} \texttt{OpConvertFToS}; \textbf{GLSL:} \texttt{int y = int(x);})
  \item \hyperref[enum:widen_int_m]{\texttt{widen1}}: Stage A widening extractor on \texttt{src1} (promotes narrow byte/half/word lane before Stage B). \emph{Example:} {\scriptsize\ttfamily \mnem{IMUL.u64} \dst{[r0:r1]}, \src{[r2:r3]}\imod{.w0}, \src{[r4:r5]}\imod{.w0}} (\textbf{SPIR-V:} \texttt{OpUConvert / OpUMulExtended}; \textbf{GLSL:} \texttt{uint64\_t p = uint64\_t(a) * uint64\_t(b);})
  \item \hyperref[enum:saturate_m]{\texttt{saturate}}: Stage B integer saturation flag (\texttt{.sat}): clamps signed/unsigned overflow or underflow to \texttt{[MIN\_N, MAX\_N]} / \texttt{[0, 2\textasciicircum{}N - 1]} instead of wrapping. \emph{Example:} {\scriptsize\ttfamily \mnem{IADD.s32}\imod{.sat} \dst{r0}, \src{r1}, \src{r2}} (\textbf{SPIR-V:} \texttt{OpSDotAccSat / saturating add}; \textbf{GLSL:} \texttt{int c = clamp(int64\_t(a) + int64\_t(b), INT\_MIN, INT\_MAX);})
  \item \hyperref[enum:widen_int_m]{\texttt{widen0}}: Stage A widening extractor on \texttt{src0} (promotes FP16 \texttt{h0}/\texttt{h1} to FP32, or sign/zero-extends \texttt{b0..b3}/\texttt{h0..h1}/\texttt{w0..w1} to 32/64-bit before \texttt{abs}/\texttt{neg}). \emph{Example:} {\scriptsize\ttfamily \mnem{IADD.s32} \dst{r0}, \src{r1}\imod{.h0}, \src{r2}} (\textbf{SPIR-V:} \texttt{OpSConvert / OpFConvert}; \textbf{GLSL:} \texttt{int c = int(int16\_t(a)) + b;})
  \item \texttt{dest\_register\_index}: 6-bit destination GPR index (\texttt{r0..r63}) receiving the Stage C writeback result (even-aligned pair \texttt{[rN:rN+1]} for 64-bit ops). \emph{Example:} {\scriptsize\ttfamily \mnem{FMA.f32} \dst{r0}, \src{r1}, \src{r2}, \src{r3}} (\textbf{SPIR-V:} \texttt{SSA Result ID}; \textbf{GLSL:} \texttt{float r0 = fma(r1, r2, r3);})
  \item \hyperref[enum:dest_width_m]{\texttt{dest\_width}}: Stage C destination lane write-mask (\texttt{none} = full 32-bit, \texttt{h0} = low 16-bit, \texttt{h1} = high 16-bit without bit-shifting). \emph{Example:} {\scriptsize\ttfamily \mnem{FADD.v2f16} \dst{r0}\imod{.h1}, \src{r1}, \src{r2}} (\textbf{SPIR-V:} \texttt{OpVectorInsertDynamic / WriteMask}; \textbf{GLSL:} \texttt{dst.y = a.y + b.y;})
  \item \texttt{datatype}: Selects the lane data type (\texttt{f32}, \texttt{f16}, \texttt{v2f16}, \texttt{s32}, \texttt{u32}, \texttt{v2s16}, \texttt{v2u16}, \texttt{v4s8}, \texttt{v4u8}, \texttt{s64}, \texttt{u64}, etc.). \emph{Example:} {\scriptsize\ttfamily \mnem{IADD.v2s16} \dst{r0}, \src{r1}, \src{r2}} (\textbf{SPIR-V:} \texttt{OpTypeInt / OpTypeFloat}; \textbf{GLSL:} \texttt{i16vec2 c = a + b;})
  \item \texttt{instruction\_class}: 4-bit major instruction class selector (\texttt{bits 48:51}) partitioning the 64-bit Valhall opcode space (e.g., \texttt{FMA}, \texttt{CVT}, \texttt{SFU}, \texttt{LS}, \texttt{TEX}, \texttt{BRANCH}). \emph{Example:} {\scriptsize\ttfamily \mnem{FADD.f32} \dst{r0}, \src{r1}, \src{r2}} (\textbf{SPIR-V:} \texttt{N/A (encoding class)}; \textbf{GLSL:} \texttt{// Hardware execution pipe / decode class})
  \item \texttt{fau\_page\_index}: Selects the FAU (Fast Access Uniform) page (\texttt{0..3}) when a source operand references hardware small constants, thread IDs, or lane constants. \emph{Example:} {\scriptsize\ttfamily \mnem{FADD.f32} \dst{r0}, \src{r1}, \src{k1f}} (\textbf{SPIR-V:} \texttt{OpConstant / BuiltIn}; \textbf{GLSL:} \texttt{float y = x + 1.0;})
  \item \hyperref[enum:flow_control_m]{\texttt{flow}}: 4-bit embedded flow-control \& scoreboard wait modifier (\texttt{bits 59:62}: \texttt{none}, \texttt{wait0..wait0126}, \texttt{wait}, \texttt{reconverge}, \texttt{discard}, \texttt{end}). \emph{Example:} {\scriptsize\ttfamily \mnem{FADD.f32}\imod{.wait0.reconverge} \dst{r0}, \src{r1}, \src{r2}} (\textbf{SPIR-V:} \texttt{OpControlBarrier / OpReturn}; \textbf{GLSL:} \texttt{// Scoreboard wait + warp reconvergence})
\end{tightitemize}
\clearpage

\subsection*{ISUBB}
\addcontentsline{toc}{subsection}{ISUBB}
32-bit unsigned subtract with borrow-in: \texttt{diff64 = zext64(src0) - zext64(src1) - (src2 \& 1)}; if \texttt{return\_carry == c}, returns borrow-out (\texttt{src0 \textless{} src1 + bin ? 1 : 0}), else \texttt{diff64[31:0]}.
\par\textbf{Notes:}
\begin{tightitemize}
  \item Unsigned 32-bit only.
  \item Borrow-in uses LSB of \texttt{src2}.
  \item \textbf{Source references:} \texttt{isa-v9-v14.xml}, \texttt{bifrost/bifrost/ISA.xml:1615}.
\end{tightitemize}
\codebox{\mnem{ISUBB}\imod{\{.flow\}}\imod{\{.return\_carry\}} \dst{dest}\imod{\{.dest\_width\}}, \src{src0}, \src{src1}, \src{src2}}
\textbf{Operands:}
\begin{tightitemize}\footnotesize
  \item \texttt{\dst{dest}} (\texttt{dest\_register\_index}, \texttt{dest\_width}): 32-bit destination GPR receiving the low difference word (default) or borrow-out bit (\texttt{.c}).
  \item \texttt{\src{src0}}: 32-bit unsigned minuend.
  \item \texttt{\src{src1}}: 32-bit unsigned subtrahend.
  \item \texttt{\src{src2}}: Borrow-in operand (only bit 0, \texttt{src2 \& 1}, is subtracted).
\end{tightitemize}
\par\textbf{Closest SPIR-V Op:} \texttt{OpISubBorrow}
\par\textbf{Pseudo-GLSL (Operational Semantics):}
\codebox{uint64\_t rhs = uint64\_t(src1) + uint64\_t(src2 \& 1u);\\if (return\_carry == c) dest = (uint64\_t(src0) \textless{} rhs) ? 1u : 0u;\\else                   dest = uint((uint64\_t(src0) - rhs) \& 0xFFFFFFFFul);}
\textbf{Execution Context:}
\begin{tightitemize}\footnotesize
  \item \textbf{Compute stage preloads (\hyperref[sec:csf-spd]{\texttt{SPD\_k}}, \texttt{r16..r23}):} Hardware pre-fills \texttt{r55..r57} with \texttt{WorkGroupID.xyz + JOB\_OFFSET} (from \hyperref[sec:csf-regs]{\texttt{COMPUTESR::JOB\_OFFSET\_X/Y/Z}} \texttt{r34..r36}) and \texttt{r60..r62} with \texttt{LocalInvocationID.xyz} or \texttt{GlobalInvocationID.xyz} (from \texttt{WG\_SIZE} \texttt{r33} \& \texttt{JOB\_SIZE} \texttt{r37..r39}, selected by \texttt{ShaderProgram::preload}) so compute shaders can derive thread/element indices without memory loads.
  \item \textbf{Vertex stage preloads (\hyperref[sec:csf-spd]{\texttt{SPD\_0}/\texttt{SPD\_1}}, \texttt{r16..r19}):} The hardware Index Fetcher reads \hyperref[sec:csf-regs]{\texttt{IDVSSR::INDEX\_BUFFER}} (\texttt{r54:r55}) and pre-fills \texttt{r59, r60} with \texttt{VertexIndex} (\texttt{+ VERTEX\_OFFSET + GLOBAL\_ATTR\_OFF}) and \texttt{r61} with \texttt{InstanceIndex} (\texttt{+ INSTANCE\_OFFSET}) so Position and Varying shaders can pull vertex attributes directly.
  \item \textbf{Inline FAU \& Special Operands (\hyperref[sec:csf-tsd-fau]{\texttt{FAU\_k}}, \texttt{r8..r15}):} Any 8-bit source with \texttt{mode=2} reads zero-latency uniform \texttt{u0..u31} from \texttt{FAU\_k} (\texttt{slots64 = bits[63:56]}); \texttt{mode=3} reads \texttt{SmallConstant} (Table~1.2) or \texttt{FauSpecial} (\texttt{pc}, \texttt{lane\_id}, \texttt{lane\_width}, \texttt{tls\_ptr}, \texttt{wls\_ptr}, \texttt{blend\_descriptor\_0..7}).
\end{tightitemize}
\par\textbf{Worked Examples:}
\par\textbf{Example 1 (Borrow-out extraction (.c)):}
\codebox{\mnem{ISUBB.u32}\imod{.c} \dst{r0}, \src{r1}, \src{r2}, \src{r3}}
With \texttt{r1 = 3, r2 = 3, r3 = 1} (borrow-in): \texttt{3 \textless{} 3 + 1}, so \texttt{.c} returns borrow-out \texttt{r0 = 1}.
\par\textbf{Example 2 (32-bit difference with borrow-in (no .c)):}
\codebox{\mnem{ISUBB.u32} \dst{r0}, \src{r1}, \src{r2}, \src{r3}}
With \texttt{r1 = 10, r2 = 4, r3 = 1}: computes \texttt{10 - 4 - 1 = 5} -\textgreater{} \texttt{r0 = 5} (multi-word 64/128-bit subtraction).

\begin{center}\footnotesize
\begin{tabularx}{\linewidth}{@{}lllX@{}}
\toprule
Variant & Arch & Class & Opcode Fields \\
\midrule
\texttt{(none)} & v9--v14 & \texttt{R3} & \texttt{opcode3=ISUBB, opcode1=MULTI} \\
\bottomrule
\end{tabularx}
\end{center}
\begin{center}\footnotesize
\begin{tabularx}{0.88\linewidth}{@{}l l X@{}}
\toprule
Field Name & Bits & Type / Fixed Value \\
\midrule
\texttt{src0} & \texttt{0:7} & \texttt{SourceEncoding} \\
\texttt{src1} & \texttt{8:15} & \texttt{SourceEncoding} \\
\texttt{src2} & \texttt{16:23} & \texttt{SourceEncoding} \\
\texttt{opcode3} & \texttt{24:25} & \texttt{exact=ISUBB} \\
\textit{Reserved} & \texttt{26:29} & \texttt{zero} \\
\hyperref[enum:iaddc_return_carry_m]{\texttt{return\_carry}} & \texttt{30} & \texttt{iaddc\_return\_carry\_m (def none)} \\
\textit{Reserved} & \texttt{31:39} & \texttt{zero} \\
\texttt{dest\_register\_index} & \texttt{40:45} & \texttt{uint} \\
\hyperref[enum:dest_width_m]{\texttt{dest\_width}} & \texttt{46:47} & \texttt{dest\_width\_m} \\
\texttt{opcode1} & \texttt{48:51} & \texttt{exact=MULTI} \\
\texttt{instruction\_class} & \texttt{52:56} & \texttt{exact=R3} \\
\texttt{fau\_page\_index} & \texttt{57:58} & \texttt{uint} \\
\hyperref[enum:flow_control_m]{\texttt{flow}} & \texttt{59:62} & \texttt{flow\_control\_m} \\
\textit{Reserved} & \texttt{63} & \texttt{zero} \\
\bottomrule
\end{tabularx}
\end{center}
\textbf{Fields Explained:}
\begin{tightitemize}\footnotesize
  \item \texttt{src0}: Primary source operand (\texttt{r0..r63}, \texttt{u0..u63} uniform, \texttt{k*} FAU constant, or \texttt{t0..t2} temporary; \texttt{\textasciicircum{}} marks discard/last-use). \emph{Example:} {\scriptsize\ttfamily \mnem{FADD.f32} \dst{r0}, \src{r1}, \src{r2}} (\textbf{SPIR-V:} \texttt{OpFAdd (Operand 0)}; \textbf{GLSL:} \texttt{float c = a + b;})
  \item \texttt{src1}: Secondary source operand (register, uniform, or FAU constant; in shifts, supplies the shift count byte). \emph{Example:} {\scriptsize\ttfamily \mnem{LSHIFT.i32} \dst{r0}, \src{r1}, \src{r2}\imod{.b0}} (\textbf{SPIR-V:} \texttt{OpShiftLeftLogical (Shift)}; \textbf{GLSL:} \texttt{uint y = a \textless{}\textless{} s;})
  \item \texttt{src2}: Third source operand (FMA addend, \texttt{MUX}/\texttt{CSEL} selector, \texttt{IADDC} carry-in, or bitwise mask in fused shift-logic). \emph{Example:} {\scriptsize\ttfamily \mnem{FMA.f32} \dst{r0}, \src{r1}, \src{r2}, \src{r3}} (\textbf{SPIR-V:} \texttt{GLSL.std.450 Fma (Addend)}; \textbf{GLSL:} \texttt{float d = fma(a, b, c);})
  \item \texttt{opcode3}: Tertiary sub-opcode disambiguating specialized multi-op families (\texttt{FPOW\_SC\_*}, SFU lookup tables, \texttt{ICMP\_MULTI}). \emph{Example:} {\scriptsize\ttfamily \mnem{FPOW\_SC\_DET.f32}\imod{.pow} \dst{r0}, \src{r1}, \src{r2}} (\textbf{SPIR-V:} \texttt{GLSL.std.450 Pow}; \textbf{GLSL:} \texttt{float z = pow(x, y);})
  \item \hyperref[enum:iaddc_return_carry_m]{\texttt{return\_carry}}: When \texttt{.c} is set on \texttt{IADDC}/\texttt{ISUBB}, returns the 1-bit carry/borrow instead of the 32-bit sum/diff. \emph{Example:} {\scriptsize\ttfamily \mnem{IADDC.u32}\imod{.c} \dst{r0}, \src{r1}, \src{r2}} (\textbf{SPIR-V:} \texttt{OpIAddCarry}; \textbf{GLSL:} \texttt{uint carry; uint s = uaddCarry(a, b, carry);})
  \item \texttt{dest\_register\_index}: 6-bit destination GPR index (\texttt{r0..r63}) receiving the Stage C writeback result (even-aligned pair \texttt{[rN:rN+1]} for 64-bit ops). \emph{Example:} {\scriptsize\ttfamily \mnem{FMA.f32} \dst{r0}, \src{r1}, \src{r2}, \src{r3}} (\textbf{SPIR-V:} \texttt{SSA Result ID}; \textbf{GLSL:} \texttt{float r0 = fma(r1, r2, r3);})
  \item \hyperref[enum:dest_width_m]{\texttt{dest\_width}}: Stage C destination lane write-mask (\texttt{none} = full 32-bit, \texttt{h0} = low 16-bit, \texttt{h1} = high 16-bit without bit-shifting). \emph{Example:} {\scriptsize\ttfamily \mnem{FADD.v2f16} \dst{r0}\imod{.h1}, \src{r1}, \src{r2}} (\textbf{SPIR-V:} \texttt{OpVectorInsertDynamic / WriteMask}; \textbf{GLSL:} \texttt{dst.y = a.y + b.y;})
  \item \texttt{opcode1}: Primary sub-opcode within the instruction class uniquely identifying the operation family. \emph{Example:} {\scriptsize\ttfamily \mnem{IADD.s32} \dst{r0}, \src{r1}, \src{r2}} (\textbf{SPIR-V:} \texttt{OpIAdd (operation selector)}; \textbf{GLSL:} \texttt{// Primary opcode bits})
  \item \texttt{instruction\_class}: 4-bit major instruction class selector (\texttt{bits 48:51}) partitioning the 64-bit Valhall opcode space (e.g., \texttt{FMA}, \texttt{CVT}, \texttt{SFU}, \texttt{LS}, \texttt{TEX}, \texttt{BRANCH}). \emph{Example:} {\scriptsize\ttfamily \mnem{FADD.f32} \dst{r0}, \src{r1}, \src{r2}} (\textbf{SPIR-V:} \texttt{N/A (encoding class)}; \textbf{GLSL:} \texttt{// Hardware execution pipe / decode class})
  \item \texttt{fau\_page\_index}: Selects the FAU (Fast Access Uniform) page (\texttt{0..3}) when a source operand references hardware small constants, thread IDs, or lane constants. \emph{Example:} {\scriptsize\ttfamily \mnem{FADD.f32} \dst{r0}, \src{r1}, \src{k1f}} (\textbf{SPIR-V:} \texttt{OpConstant / BuiltIn}; \textbf{GLSL:} \texttt{float y = x + 1.0;})
  \item \hyperref[enum:flow_control_m]{\texttt{flow}}: 4-bit embedded flow-control \& scoreboard wait modifier (\texttt{bits 59:62}: \texttt{none}, \texttt{wait0..wait0126}, \texttt{wait}, \texttt{reconverge}, \texttt{discard}, \texttt{end}). \emph{Example:} {\scriptsize\ttfamily \mnem{FADD.f32}\imod{.wait0.reconverge} \dst{r0}, \src{r1}, \src{r2}} (\textbf{SPIR-V:} \texttt{OpControlBarrier / OpReturn}; \textbf{GLSL:} \texttt{// Scoreboard wait + warp reconvergence})
\end{tightitemize}
\clearpage

\subsection*{MMUL}
\addcontentsline{toc}{subsection}{MMUL}
Subgroup/quad cooperative matrix multiply-accumulate \texttt{D = A * B + C} across lanes: \texttt{f32} (4x4*4x4-\textgreater{}4x4), \texttt{v2f16} (4x4 submat of 4x8 -\textgreater{} 4x4 f32), \texttt{f16} (4x4 submat of 4x8 * 4x8 -\textgreater{} 4x8 f16, v14), \texttt{v4s8}/\texttt{v4u8} (4x16*16x4 -\textgreater{} 4x4 s32/u32).
\par\textbf{Notes:}
\begin{tightitemize}
  \item Arch \textgreater{}= v11 (\texttt{f16} variant and mixed-sign/sat \texttt{v4s8} on v14).
  \item Cross-lane subgroup operation.
  \item \textbf{Source references:} \texttt{kraid/ops.rs:3411-3550}, \texttt{encode\_v9.rs:2723-2819}, \texttt{nir.rs:1907-1957}.
\end{tightitemize}
\codebox{\mnem{MMUL}\imod{\{.flow\}} \dst{dest}\imod{\{.dest\_width\}}, \src{src0}\imod{\{.sub0\}}, \src{src1}, \src{src2}\\\mnem{MMUL}\imod{\{.flow\}} \dst{dest}\imod{\{.dest\_width\}}, \src{src0}, \src{src1}, \src{src2}\\\mnem{MMUL}\imod{\{.flow\}} \dst{dest}\imod{\{.dest\_width\}}, \src{src0}\imod{\{.sub0\}}, \src{src1}\imod{\{.sub1\}}, \src{src2}\\\mnem{MMUL}\imod{\{.flow\}} \dst{dest}\imod{\{.dest\_width\}}, \src{src0}\imod{\{.unsigned0\}}, \src{src1}\imod{\{.unsigned1\}}, \src{src2}\\\mnem{MMUL}\imod{\{.flow\}}\imod{\{.saturate\}} \dst{dest}\imod{\{.dest\_width\}}, \src{src0}, \src{src1}, \src{src2}}
\textbf{Operands:}
\begin{tightitemize}\footnotesize
  \item \texttt{\dst{dest}} (\texttt{dest\_register\_index}, \texttt{dest\_width}): Per-lane slice of cooperative output matrix tile \texttt{D = A * B + C}.
  \item \texttt{\src{src0}}: Per-lane slice of cooperative input matrix tile \texttt{A}.
  \item \texttt{\src{src1}}: Per-lane slice of cooperative input matrix tile \texttt{B}.
  \item \texttt{\src{src2}}: Per-lane slice of cooperative accumulator matrix tile \texttt{C}.
\end{tightitemize}
\par\textbf{Closest SPIR-V Op:} \texttt{OpCooperativeMatrixMulAddKHR}
\par\textbf{Pseudo-GLSL (Operational Semantics):}
\codebox{// Cooperative subgroup matrix multiply-accumulate across warp lanes:\\mat\_tile D = mat\_tile(src2) + submat(src0, sub0) * submat(src1, sub1);\\dest = flag\_sat ? saturate(D)[lane\_slice] : D[lane\_slice];}
\textbf{Execution Context:}
\begin{tightitemize}\footnotesize
  \item \textbf{Compute stage preloads (\hyperref[sec:csf-spd]{\texttt{SPD\_k}}, \texttt{r16..r23}):} Hardware pre-fills \texttt{r55..r57} with \texttt{WorkGroupID.xyz + JOB\_OFFSET} (from \hyperref[sec:csf-regs]{\texttt{COMPUTESR::JOB\_OFFSET\_X/Y/Z}} \texttt{r34..r36}) and \texttt{r60..r62} with \texttt{LocalInvocationID.xyz} or \texttt{GlobalInvocationID.xyz} (from \texttt{WG\_SIZE} \texttt{r33} \& \texttt{JOB\_SIZE} \texttt{r37..r39}, selected by \texttt{ShaderProgram::preload}) so compute shaders can derive thread/element indices without memory loads.
  \item \textbf{Vertex stage preloads (\hyperref[sec:csf-spd]{\texttt{SPD\_0}/\texttt{SPD\_1}}, \texttt{r16..r19}):} The hardware Index Fetcher reads \hyperref[sec:csf-regs]{\texttt{IDVSSR::INDEX\_BUFFER}} (\texttt{r54:r55}) and pre-fills \texttt{r59, r60} with \texttt{VertexIndex} (\texttt{+ VERTEX\_OFFSET + GLOBAL\_ATTR\_OFF}) and \texttt{r61} with \texttt{InstanceIndex} (\texttt{+ INSTANCE\_OFFSET}) so Position and Varying shaders can pull vertex attributes directly.
  \item \textbf{Inline FAU \& Special Operands (\hyperref[sec:csf-tsd-fau]{\texttt{FAU\_k}}, \texttt{r8..r15}):} Any 8-bit source with \texttt{mode=2} reads zero-latency uniform \texttt{u0..u31} from \texttt{FAU\_k} (\texttt{slots64 = bits[63:56]}); \texttt{mode=3} reads \texttt{SmallConstant} (Table~1.2) or \texttt{FauSpecial} (\texttt{pc}, \texttt{lane\_id}, \texttt{lane\_width}, \texttt{tls\_ptr}, \texttt{wls\_ptr}, \texttt{blend\_descriptor\_0..7}).
\end{tightitemize}
\par\textbf{Worked Examples:}
\par\textbf{Example 1 (FP16 submatrix multiply-accumulate into FP32):}
\codebox{\mnem{MMUL.v2f16} \dst{r0}, \src{r1}\imod{.f0}, \src{r2}\imod{.f1}, \src{r3}}
Each lane in a 4x4 subgroup computes one output element of \texttt{D = A[f0] * B[f1] + C} using 4x4 FP16 submatrices accumulated in FP32 \texttt{r0}.
\par\textbf{Example 2 (v14 mixed-sign INT8 saturating matrix multiply):}
\codebox{\mnem{MMUL.v4s8}\imod{.sat} \dst{r0}, \src{r1}, \src{r2}, \src{r3}}
Multiplies 4x16 \texttt{v4s8} tile \texttt{A} by 16x4 \texttt{B}, accumulates into signed 32-bit \texttt{r3}, and saturates to \texttt{[INT32\_MIN, INT32\_MAX]} (v14+).

\begin{center}\footnotesize
\begin{tabularx}{\linewidth}{@{}lllX@{}}
\toprule
Variant & Arch & Class & Opcode Fields \\
\midrule
\texttt{(none)} & v9--v14 & \texttt{M3} & \texttt{opcode1=MMUL\_F16} \\
\texttt{(none)} & v9--v14 & \texttt{M3} & \texttt{opcode1=MMUL\_F32} \\
\texttt{(none)} & v9--v14 & \texttt{M3} & \texttt{opcode1=MMUL\_V2F16} \\
\texttt{(none)} & v9--v14 & \texttt{M3} & \texttt{opcode1=MMUL\_V4S8} \\
\texttt{(none)} & v9--v14 & \texttt{M3} & \texttt{opcode1=MMUL\_V4S8} \\
\texttt{(none)} & v9--v14 & \texttt{M3} & \texttt{opcode1=MMUL\_V4U8} \\
\texttt{(none)} & v9--v14 & \texttt{M3} & \texttt{opcode1=MMUL\_V4U8} \\
\bottomrule
\end{tabularx}
\end{center}
\begin{center}\footnotesize
\begin{tabularx}{0.88\linewidth}{@{}l l X@{}}
\toprule
Field Name & Bits & Type / Fixed Value \\
\midrule
\texttt{src0} & \texttt{0:7} & \texttt{SourceEncoding} \\
\texttt{src1} & \texttt{8:15} & \texttt{SourceEncoding} \\
\texttt{src2} & \texttt{16:23} & \texttt{SourceEncoding} \\
\textit{Reserved} & \texttt{24:25} & \texttt{zero} \\
\hyperref[enum:submatrix_sel_m]{\texttt{sub1}} & \texttt{26:27} & \texttt{exact=none} \\
\hyperref[enum:submatrix_sel_m]{\texttt{sub0}} & \texttt{28:29} & \texttt{submatrix\_sel\_m} \\
\textit{Reserved} & \texttt{30:39} & \texttt{zero} \\
\texttt{dest\_register\_index} & \texttt{40:45} & \texttt{uint} \\
\hyperref[enum:dest_width_m]{\texttt{dest\_width}} & \texttt{46:47} & \texttt{exact=none} \\
\texttt{opcode1} & \texttt{48:51} & \texttt{exact=MMUL\_F16} \\
\texttt{instruction\_class} & \texttt{52:56} & \texttt{exact=M3} \\
\texttt{fau\_page\_index} & \texttt{57:58} & \texttt{exact=0} \\
\hyperref[enum:flow_control_m]{\texttt{flow}} & \texttt{59:62} & \texttt{flow\_control\_m} \\
\textit{Reserved} & \texttt{63} & \texttt{zero} \\
\bottomrule
\end{tabularx}
\end{center}
\textbf{Fields Explained:}
\begin{tightitemize}\footnotesize
  \item \texttt{src0}: Primary source operand (\texttt{r0..r63}, \texttt{u0..u63} uniform, \texttt{k*} FAU constant, or \texttt{t0..t2} temporary; \texttt{\textasciicircum{}} marks discard/last-use). \emph{Example:} {\scriptsize\ttfamily \mnem{FADD.f32} \dst{r0}, \src{r1}, \src{r2}} (\textbf{SPIR-V:} \texttt{OpFAdd (Operand 0)}; \textbf{GLSL:} \texttt{float c = a + b;})
  \item \texttt{src1}: Secondary source operand (register, uniform, or FAU constant; in shifts, supplies the shift count byte). \emph{Example:} {\scriptsize\ttfamily \mnem{LSHIFT.i32} \dst{r0}, \src{r1}, \src{r2}\imod{.b0}} (\textbf{SPIR-V:} \texttt{OpShiftLeftLogical (Shift)}; \textbf{GLSL:} \texttt{uint y = a \textless{}\textless{} s;})
  \item \texttt{src2}: Third source operand (FMA addend, \texttt{MUX}/\texttt{CSEL} selector, \texttt{IADDC} carry-in, or bitwise mask in fused shift-logic). \emph{Example:} {\scriptsize\ttfamily \mnem{FMA.f32} \dst{r0}, \src{r1}, \src{r2}, \src{r3}} (\textbf{SPIR-V:} \texttt{GLSL.std.450 Fma (Addend)}; \textbf{GLSL:} \texttt{float d = fma(a, b, c);})
  \item \hyperref[enum:submatrix_sel_m]{\texttt{sub1}}: Selects submatrix half (\texttt{f0}, \texttt{f1}, \texttt{none}) of source matrix B in \texttt{MMUL}. \emph{Example:} {\scriptsize\ttfamily \mnem{MMUL.v2f16}\imod{.f1} \dst{r0}, \src{r1}, \src{r2}, \src{r3}} (\textbf{SPIR-V:} \texttt{OpCooperativeMatrixMulAddKHR}; \textbf{GLSL:} \texttt{// 4x4/8x8 cooperative matrix multiply})
  \item \hyperref[enum:submatrix_sel_m]{\texttt{sub0}}: Selects submatrix half (\texttt{f0}, \texttt{f1}, \texttt{none}) of source matrix A in \texttt{MMUL}. \emph{Example:} {\scriptsize\ttfamily \mnem{MMUL.v2f16}\imod{.f0} \dst{r0}, \src{r1}, \src{r2}, \src{r3}} (\textbf{SPIR-V:} \texttt{OpCooperativeMatrixMulAddKHR}; \textbf{GLSL:} \texttt{// 4x4/8x8 cooperative matrix multiply})
  \item \texttt{dest\_register\_index}: 6-bit destination GPR index (\texttt{r0..r63}) receiving the Stage C writeback result (even-aligned pair \texttt{[rN:rN+1]} for 64-bit ops). \emph{Example:} {\scriptsize\ttfamily \mnem{FMA.f32} \dst{r0}, \src{r1}, \src{r2}, \src{r3}} (\textbf{SPIR-V:} \texttt{SSA Result ID}; \textbf{GLSL:} \texttt{float r0 = fma(r1, r2, r3);})
  \item \hyperref[enum:dest_width_m]{\texttt{dest\_width}}: Stage C destination lane write-mask (\texttt{none} = full 32-bit, \texttt{h0} = low 16-bit, \texttt{h1} = high 16-bit without bit-shifting). \emph{Example:} {\scriptsize\ttfamily \mnem{FADD.v2f16} \dst{r0}\imod{.h1}, \src{r1}, \src{r2}} (\textbf{SPIR-V:} \texttt{OpVectorInsertDynamic / WriteMask}; \textbf{GLSL:} \texttt{dst.y = a.y + b.y;})
  \item \texttt{opcode1}: Primary sub-opcode within the instruction class uniquely identifying the operation family. \emph{Example:} {\scriptsize\ttfamily \mnem{IADD.s32} \dst{r0}, \src{r1}, \src{r2}} (\textbf{SPIR-V:} \texttt{OpIAdd (operation selector)}; \textbf{GLSL:} \texttt{// Primary opcode bits})
  \item \texttt{instruction\_class}: 4-bit major instruction class selector (\texttt{bits 48:51}) partitioning the 64-bit Valhall opcode space (e.g., \texttt{FMA}, \texttt{CVT}, \texttt{SFU}, \texttt{LS}, \texttt{TEX}, \texttt{BRANCH}). \emph{Example:} {\scriptsize\ttfamily \mnem{FADD.f32} \dst{r0}, \src{r1}, \src{r2}} (\textbf{SPIR-V:} \texttt{N/A (encoding class)}; \textbf{GLSL:} \texttt{// Hardware execution pipe / decode class})
  \item \texttt{fau\_page\_index}: Selects the FAU (Fast Access Uniform) page (\texttt{0..3}) when a source operand references hardware small constants, thread IDs, or lane constants. \emph{Example:} {\scriptsize\ttfamily \mnem{FADD.f32} \dst{r0}, \src{r1}, \src{k1f}} (\textbf{SPIR-V:} \texttt{OpConstant / BuiltIn}; \textbf{GLSL:} \texttt{float y = x + 1.0;})
  \item \hyperref[enum:flow_control_m]{\texttt{flow}}: 4-bit embedded flow-control \& scoreboard wait modifier (\texttt{bits 59:62}: \texttt{none}, \texttt{wait0..wait0126}, \texttt{wait}, \texttt{reconverge}, \texttt{discard}, \texttt{end}). \emph{Example:} {\scriptsize\ttfamily \mnem{FADD.f32}\imod{.wait0.reconverge} \dst{r0}, \src{r1}, \src{r2}} (\textbf{SPIR-V:} \texttt{OpControlBarrier / OpReturn}; \textbf{GLSL:} \texttt{// Scoreboard wait + warp reconvergence})
\end{tightitemize}
\clearpage

\subsection*{POPCOUNT}
\addcontentsline{toc}{subsection}{POPCOUNT}
Counts the total number of \texttt{1} bits in the 32-bit word \texttt{src0} and writes the exact unsigned count (in the range $0 \dots 32$) to \texttt{dest}.
\par\textbf{Notes:}
\begin{tightitemize}
  \item Exact 32-bit popcount (\texttt{hw\_tests.rs:2138-2145}).
  \item \textbf{Source references:} \texttt{kraid/ops.rs:3738-3762}, \texttt{encode\_v9.rs:2897-2914}, \texttt{nir.rs:1248-1254}, \texttt{hw\_tests.rs:2138-2145}.
\end{tightitemize}
\codebox{\mnem{POPCOUNT}\imod{\{.flow\}} \dst{dest}\imod{\{.dest\_width\}}, \src{src0}}
\textbf{Operands:}
\begin{tightitemize}\footnotesize
  \item \texttt{\dst{dest}} (\texttt{dest\_register\_index}, \texttt{dest\_width}): 32-bit destination GPR receiving the number of set bits in \texttt{src0} (\texttt{0..32}).
  \item \texttt{\src{src0}}: 32-bit input bitmask.
\end{tightitemize}
\par\textbf{Closest SPIR-V Op:} \texttt{OpBitCount}
\par\textbf{Pseudo-GLSL (Operational Semantics):}
\codebox{dest = uint(bitCount(src0)); // exact count of 1-bits in 0..32}
\textbf{Execution Context:}
\begin{tightitemize}\footnotesize
  \item \textbf{Compute stage preloads (\hyperref[sec:csf-spd]{\texttt{SPD\_k}}, \texttt{r16..r23}):} Hardware pre-fills \texttt{r55..r57} with \texttt{WorkGroupID.xyz + JOB\_OFFSET} (from \hyperref[sec:csf-regs]{\texttt{COMPUTESR::JOB\_OFFSET\_X/Y/Z}} \texttt{r34..r36}) and \texttt{r60..r62} with \texttt{LocalInvocationID.xyz} or \texttt{GlobalInvocationID.xyz} (from \texttt{WG\_SIZE} \texttt{r33} \& \texttt{JOB\_SIZE} \texttt{r37..r39}, selected by \texttt{ShaderProgram::preload}) so compute shaders can derive thread/element indices without memory loads.
  \item \textbf{Vertex stage preloads (\hyperref[sec:csf-spd]{\texttt{SPD\_0}/\texttt{SPD\_1}}, \texttt{r16..r19}):} The hardware Index Fetcher reads \hyperref[sec:csf-regs]{\texttt{IDVSSR::INDEX\_BUFFER}} (\texttt{r54:r55}) and pre-fills \texttt{r59, r60} with \texttt{VertexIndex} (\texttt{+ VERTEX\_OFFSET + GLOBAL\_ATTR\_OFF}) and \texttt{r61} with \texttt{InstanceIndex} (\texttt{+ INSTANCE\_OFFSET}) so Position and Varying shaders can pull vertex attributes directly.
  \item \textbf{Inline FAU \& Special Operands (\hyperref[sec:csf-tsd-fau]{\texttt{FAU\_k}}, \texttt{r8..r15}):} Any 8-bit source with \texttt{mode=2} reads zero-latency uniform \texttt{u0..u31} from \texttt{FAU\_k} (\texttt{slots64 = bits[63:56]}); \texttt{mode=3} reads \texttt{SmallConstant} (Table~1.2) or \texttt{FauSpecial} (\texttt{pc}, \texttt{lane\_id}, \texttt{lane\_width}, \texttt{tls\_ptr}, \texttt{wls\_ptr}, \texttt{blend\_descriptor\_0..7}).
\end{tightitemize}
\par\textbf{Worked Examples:}
\par\textbf{Example 1 (Mixed bit pattern):}
\codebox{\mnem{POPCOUNT.i32} \dst{r0}, \src{r1}}
With \texttt{r1 = 0xF0F00003}: counts set bits \texttt{4 + 4 + 2 = 10} -\textgreater{} \texttt{r0 = 10}.
\par\textbf{Example 2 (Boundary values (0 and all-ones)):}
\codebox{\mnem{POPCOUNT.i32} \dst{r0}, \src{r1}}
With \texttt{r1 = 0x00000000} -\textgreater{} \texttt{r0 = 0}; with \texttt{r1 = 0xFFFFFFFF} -\textgreater{} \texttt{r0 = 32}.

\begin{center}\footnotesize
\begin{tabularx}{\linewidth}{@{}lllX@{}}
\toprule
Variant & Arch & Class & Opcode Fields \\
\midrule
\texttt{(none)} & v9--v14 & \texttt{R1} & \texttt{opcode2=POPCOUNT, opcode1=MOV} \\
\bottomrule
\end{tabularx}
\end{center}
\begin{center}\footnotesize
\begin{tabularx}{0.88\linewidth}{@{}l l X@{}}
\toprule
Field Name & Bits & Type / Fixed Value \\
\midrule
\texttt{src0} & \texttt{0:7} & \texttt{SourceEncoding} \\
\textit{Reserved} & \texttt{8:15} & \texttt{zero} \\
\texttt{opcode2} & \texttt{16:20} & \texttt{exact=POPCOUNT} \\
\textit{Reserved} & \texttt{21:39} & \texttt{zero} \\
\texttt{dest\_register\_index} & \texttt{40:45} & \texttt{uint} \\
\hyperref[enum:dest_width_m]{\texttt{dest\_width}} & \texttt{46:47} & \texttt{dest\_width\_m} \\
\texttt{opcode1} & \texttt{48:51} & \texttt{exact=MOV} \\
\texttt{instruction\_class} & \texttt{52:56} & \texttt{exact=R1} \\
\texttt{fau\_page\_index} & \texttt{57:58} & \texttt{uint} \\
\hyperref[enum:flow_control_m]{\texttt{flow}} & \texttt{59:62} & \texttt{flow\_control\_m} \\
\textit{Reserved} & \texttt{63} & \texttt{zero} \\
\bottomrule
\end{tabularx}
\end{center}
\textbf{Fields Explained:}
\begin{tightitemize}\footnotesize
  \item \texttt{src0}: Primary source operand (\texttt{r0..r63}, \texttt{u0..u63} uniform, \texttt{k*} FAU constant, or \texttt{t0..t2} temporary; \texttt{\textasciicircum{}} marks discard/last-use). \emph{Example:} {\scriptsize\ttfamily \mnem{FADD.f32} \dst{r0}, \src{r1}, \src{r2}} (\textbf{SPIR-V:} \texttt{OpFAdd (Operand 0)}; \textbf{GLSL:} \texttt{float c = a + b;})
  \item \texttt{opcode2}: Secondary sub-opcode further specializing instructions sharing \texttt{instruction\_class} and \texttt{opcode1} (e.g., conversions, SFU tables, moves). \emph{Example:} {\scriptsize\ttfamily \mnem{F32\_TO\_S32}\imod{.round\_zero} \dst{r0}, \src{r1}} (\textbf{SPIR-V:} \texttt{OpConvertFToS}; \textbf{GLSL:} \texttt{int y = int(x);})
  \item \texttt{dest\_register\_index}: 6-bit destination GPR index (\texttt{r0..r63}) receiving the Stage C writeback result (even-aligned pair \texttt{[rN:rN+1]} for 64-bit ops). \emph{Example:} {\scriptsize\ttfamily \mnem{FMA.f32} \dst{r0}, \src{r1}, \src{r2}, \src{r3}} (\textbf{SPIR-V:} \texttt{SSA Result ID}; \textbf{GLSL:} \texttt{float r0 = fma(r1, r2, r3);})
  \item \hyperref[enum:dest_width_m]{\texttt{dest\_width}}: Stage C destination lane write-mask (\texttt{none} = full 32-bit, \texttt{h0} = low 16-bit, \texttt{h1} = high 16-bit without bit-shifting). \emph{Example:} {\scriptsize\ttfamily \mnem{FADD.v2f16} \dst{r0}\imod{.h1}, \src{r1}, \src{r2}} (\textbf{SPIR-V:} \texttt{OpVectorInsertDynamic / WriteMask}; \textbf{GLSL:} \texttt{dst.y = a.y + b.y;})
  \item \texttt{opcode1}: Primary sub-opcode within the instruction class uniquely identifying the operation family. \emph{Example:} {\scriptsize\ttfamily \mnem{IADD.s32} \dst{r0}, \src{r1}, \src{r2}} (\textbf{SPIR-V:} \texttt{OpIAdd (operation selector)}; \textbf{GLSL:} \texttt{// Primary opcode bits})
  \item \texttt{instruction\_class}: 4-bit major instruction class selector (\texttt{bits 48:51}) partitioning the 64-bit Valhall opcode space (e.g., \texttt{FMA}, \texttt{CVT}, \texttt{SFU}, \texttt{LS}, \texttt{TEX}, \texttt{BRANCH}). \emph{Example:} {\scriptsize\ttfamily \mnem{FADD.f32} \dst{r0}, \src{r1}, \src{r2}} (\textbf{SPIR-V:} \texttt{N/A (encoding class)}; \textbf{GLSL:} \texttt{// Hardware execution pipe / decode class})
  \item \texttt{fau\_page\_index}: Selects the FAU (Fast Access Uniform) page (\texttt{0..3}) when a source operand references hardware small constants, thread IDs, or lane constants. \emph{Example:} {\scriptsize\ttfamily \mnem{FADD.f32} \dst{r0}, \src{r1}, \src{k1f}} (\textbf{SPIR-V:} \texttt{OpConstant / BuiltIn}; \textbf{GLSL:} \texttt{float y = x + 1.0;})
  \item \hyperref[enum:flow_control_m]{\texttt{flow}}: 4-bit embedded flow-control \& scoreboard wait modifier (\texttt{bits 59:62}: \texttt{none}, \texttt{wait0..wait0126}, \texttt{wait}, \texttt{reconverge}, \texttt{discard}, \texttt{end}). \emph{Example:} {\scriptsize\ttfamily \mnem{FADD.f32}\imod{.wait0.reconverge} \dst{r0}, \src{r1}, \src{r2}} (\textbf{SPIR-V:} \texttt{OpControlBarrier / OpReturn}; \textbf{GLSL:} \texttt{// Scoreboard wait + warp reconvergence})
\end{tightitemize}
\clearpage

\subsection*{SEG\_ADD}
\addcontentsline{toc}{subsection}{SEG\_ADD}
Computes 64-bit segment pointer addition: widens \texttt{src1} to 64 bits and adds it to the 64-bit base pointer \texttt{src0} (\texttt{dest = src0 + widen64(src1)}), wrapping modulo $2^{64}$. When \texttt{.preserve\_null} is set and \texttt{src0 == 0}, returns \texttt{0} so null pointers remain null.
\par\textbf{Notes:}
\begin{tightitemize}
  \item Used for WLS/TLS segment pointer arithmetic.
  \item \texttt{preserve\_null} keeps null pointer \texttt{0} as \texttt{0}.
  \item \textbf{Source references:} \texttt{bifrost/valhall/ISA.xml:2666-2684}, \texttt{bifrost/compiler.h:720-722}, \texttt{isa-v9-v14.xml}.
\end{tightitemize}
\codebox{\mnem{SEG\_ADD}\imod{\{.flow\}}\imod{\{.preserve\_null\}} \dst{dest}, \src{src0}, \src{src1}\imod{\{.widen1\}}}
\textbf{Operands:}
\begin{tightitemize}\footnotesize
  \item \texttt{\dst{dest}} (\texttt{dest\_register\_index}): 64-bit destination register pair receiving the offset segment pointer.
  \item \texttt{\src{src0}}: 64-bit base segment pointer (preserved as \texttt{0} when \texttt{.preserve\_null} is set).
  \item \texttt{\src{src1}}: Signed/unsigned offset widened to 64 bits and added to \texttt{src0}.
\end{tightitemize}
\par\textbf{Closest SPIR-V Op:} \texttt{OpPtrAccessChain (TLS/WLS segment)}
\par\textbf{Pseudo-GLSL (Operational Semantics):}
\codebox{if (flag\_preserve\_null \&\& src0 == 0ul) dest = 0ul;\\else                                   dest = src0 + widen64(src1);}
\textbf{Execution Context:}
\begin{tightitemize}\footnotesize
  \item \textbf{Thread/Workgroup Local Storage (\hyperref[sec:csf-tsd-fau]{\texttt{TSD\_k}}, \texttt{r24..r31}):} 32B \texttt{LocalStorage} descriptor (\texttt{tls\_size}, \texttt{wls\_instances}, \texttt{wls\_size\_base/scale}, \texttt{tls\_base\_pointer}, \texttt{wls\_base\_pointer}) exposed inline via \texttt{FauSpecial::tls\_ptr}/\texttt{wls\_ptr}.
  \item \textbf{Inline FAU \& Special Operands (\hyperref[sec:csf-tsd-fau]{\texttt{FAU\_k}}, \texttt{r8..r15}):} Any 8-bit source with \texttt{mode=2} reads zero-latency uniform \texttt{u0..u31} from \texttt{FAU\_k} (\texttt{slots64 = bits[63:56]}); \texttt{mode=3} reads \texttt{SmallConstant} (Table~1.2) or \texttt{FauSpecial} (\texttt{pc}, \texttt{lane\_id}, \texttt{lane\_width}, \texttt{tls\_ptr}, \texttt{wls\_ptr}, \texttt{blend\_descriptor\_0..7}).
\end{tightitemize}
\par\textbf{Worked Examples:}
\par\textbf{Example 1 (Null-preserving pointer offset):}
\codebox{\mnem{SEG\_ADD.u64}\imod{.preserve\_null} \dst{[r0:r1]}, \src{[r2:r3]}, \src{r4}\imod{.w0}}
With \texttt{[r2:r3] = 0x0}, \texttt{r4 = 0x20}: \texttt{.preserve\_null} detects null base and keeps \texttt{[r0:r1] = 0x0} instead of \texttt{0x20}.
\par\textbf{Example 2 (Non-null WLS pointer advance):}
\codebox{\mnem{SEG\_ADD.u64}\imod{.preserve\_null} \dst{[r0:r1]}, \src{[r2:r3]}, \src{r4}\imod{.w0}}
With \texttt{[r2:r3] = 0x00000001\_00001000}, \texttt{r4 = 0x40}: zero-extends \texttt{r4.w0} and adds -\textgreater{} \texttt{[r0:r1] = 0x00000001\_00001040}.

\begin{center}\footnotesize
\begin{tabularx}{\linewidth}{@{}lllX@{}}
\toprule
Variant & Arch & Class & Opcode Fields \\
\midrule
\texttt{u64} & v9--v14 & \texttt{D2} & \texttt{opcode2=SEG\_ADD} \\
\bottomrule
\end{tabularx}
\end{center}
\begin{center}\footnotesize
\begin{tabularx}{0.88\linewidth}{@{}l l X@{}}
\toprule
Field Name & Bits & Type / Fixed Value \\
\midrule
\texttt{src0} & \texttt{0:7} & \texttt{SourceEncoding64} \\
\texttt{src1} & \texttt{8:15} & \texttt{SourceEncoding64} \\
\texttt{opcode2} & \texttt{16:19} & \texttt{exact=SEG\_ADD} \\
\textit{Reserved} & \texttt{20:25} & \texttt{zero} \\
\hyperref[enum:widen_int_m]{\texttt{widen1}} & \texttt{26:29} & \texttt{widen\_int\_m (def none)} \\
\textit{Reserved} & \texttt{30:37} & \texttt{zero} \\
\hyperref[enum:neg_m]{\texttt{neg}} & \texttt{38} & \texttt{exact=none} \\
\hyperref[enum:segadd_preserve_null_m]{\texttt{preserve\_null}} & \texttt{39} & \texttt{segadd\_preserve\_null\_m (def none)} \\
\texttt{dest\_register\_index} & \texttt{40:45} & \texttt{uint} \\
\texttt{double\_width\_instruction} & \texttt{46:47} & \texttt{exact=3} \\
\texttt{datatype} & \texttt{48:51} & \texttt{exact=u64} \\
\texttt{instruction\_class} & \texttt{52:56} & \texttt{exact=D2} \\
\texttt{fau\_page\_index} & \texttt{57:58} & \texttt{uint} \\
\hyperref[enum:flow_control_m]{\texttt{flow}} & \texttt{59:62} & \texttt{flow\_control\_m} \\
\textit{Reserved} & \texttt{63} & \texttt{zero} \\
\bottomrule
\end{tabularx}
\end{center}
\textbf{Fields Explained:}
\begin{tightitemize}\footnotesize
  \item \texttt{src0}: Primary source operand (\texttt{r0..r63}, \texttt{u0..u63} uniform, \texttt{k*} FAU constant, or \texttt{t0..t2} temporary; \texttt{\textasciicircum{}} marks discard/last-use). \emph{Example:} {\scriptsize\ttfamily \mnem{FADD.f32} \dst{r0}, \src{r1}, \src{r2}} (\textbf{SPIR-V:} \texttt{OpFAdd (Operand 0)}; \textbf{GLSL:} \texttt{float c = a + b;})
  \item \texttt{src1}: Secondary source operand (register, uniform, or FAU constant; in shifts, supplies the shift count byte). \emph{Example:} {\scriptsize\ttfamily \mnem{LSHIFT.i32} \dst{r0}, \src{r1}, \src{r2}\imod{.b0}} (\textbf{SPIR-V:} \texttt{OpShiftLeftLogical (Shift)}; \textbf{GLSL:} \texttt{uint y = a \textless{}\textless{} s;})
  \item \texttt{opcode2}: Secondary sub-opcode further specializing instructions sharing \texttt{instruction\_class} and \texttt{opcode1} (e.g., conversions, SFU tables, moves). \emph{Example:} {\scriptsize\ttfamily \mnem{F32\_TO\_S32}\imod{.round\_zero} \dst{r0}, \src{r1}} (\textbf{SPIR-V:} \texttt{OpConvertFToS}; \textbf{GLSL:} \texttt{int y = int(x);})
  \item \hyperref[enum:widen_int_m]{\texttt{widen1}}: Stage A widening extractor on \texttt{src1} (promotes narrow byte/half/word lane before Stage B). \emph{Example:} {\scriptsize\ttfamily \mnem{IMUL.u64} \dst{[r0:r1]}, \src{[r2:r3]}\imod{.w0}, \src{[r4:r5]}\imod{.w0}} (\textbf{SPIR-V:} \texttt{OpUConvert / OpUMulExtended}; \textbf{GLSL:} \texttt{uint64\_t p = uint64\_t(a) * uint64\_t(b);})
  \item \hyperref[enum:neg_m]{\texttt{neg}}: Negation flag on segmented pointer arithmetic (\texttt{SEG\_ADD} vs \texttt{SEG\_SUB}: subtracts widened \texttt{src1} when set) or unary FP negation. \emph{Example:} {\scriptsize\ttfamily \mnem{SEG\_SUB.u64} \dst{[r0:r1]}, \src{[r2:r3]}, \src{r4}\imod{.w0}} (\textbf{SPIR-V:} \texttt{OpPtrAccessChain (negative index)}; \textbf{GLSL:} \texttt{shared\_arr[idx - off]})
  \item \hyperref[enum:segadd_preserve_null_m]{\texttt{preserve\_null}}: Preserves NULL (0) pointer on segmented address add/sub instead of adding segment base. \emph{Example:} {\scriptsize\ttfamily \mnem{SEG\_ADD}\imod{.preserve\_null} \dst{[r0:r1]}, \src{r2}} (\textbf{SPIR-V:} \texttt{OpPtrAccessChain (PhysicalStorageBuffer)}; \textbf{GLSL:} \texttt{if (ptr != null) val = ptr.x;})
  \item \texttt{dest\_register\_index}: 6-bit destination GPR index (\texttt{r0..r63}) receiving the Stage C writeback result (even-aligned pair \texttt{[rN:rN+1]} for 64-bit ops). \emph{Example:} {\scriptsize\ttfamily \mnem{FMA.f32} \dst{r0}, \src{r1}, \src{r2}, \src{r3}} (\textbf{SPIR-V:} \texttt{SSA Result ID}; \textbf{GLSL:} \texttt{float r0 = fma(r1, r2, r3);})
  \item \texttt{double\_width\_instruction}: Indicates 64-bit / dual-slot wide instruction encoding variant. \emph{Example:} {\scriptsize\ttfamily \mnem{IADD.u64} \dst{[r0:r1]}, \src{[r2:r3]}, \src{[r4:r5]}} (\textbf{SPIR-V:} \texttt{OpIAdd (64-bit)}; \textbf{GLSL:} \texttt{uint64\_t z = x + y;})
  \item \texttt{datatype}: Selects the lane data type (\texttt{f32}, \texttt{f16}, \texttt{v2f16}, \texttt{s32}, \texttt{u32}, \texttt{v2s16}, \texttt{v2u16}, \texttt{v4s8}, \texttt{v4u8}, \texttt{s64}, \texttt{u64}, etc.). \emph{Example:} {\scriptsize\ttfamily \mnem{IADD.v2s16} \dst{r0}, \src{r1}, \src{r2}} (\textbf{SPIR-V:} \texttt{OpTypeInt / OpTypeFloat}; \textbf{GLSL:} \texttt{i16vec2 c = a + b;})
  \item \texttt{instruction\_class}: 4-bit major instruction class selector (\texttt{bits 48:51}) partitioning the 64-bit Valhall opcode space (e.g., \texttt{FMA}, \texttt{CVT}, \texttt{SFU}, \texttt{LS}, \texttt{TEX}, \texttt{BRANCH}). \emph{Example:} {\scriptsize\ttfamily \mnem{FADD.f32} \dst{r0}, \src{r1}, \src{r2}} (\textbf{SPIR-V:} \texttt{N/A (encoding class)}; \textbf{GLSL:} \texttt{// Hardware execution pipe / decode class})
  \item \texttt{fau\_page\_index}: Selects the FAU (Fast Access Uniform) page (\texttt{0..3}) when a source operand references hardware small constants, thread IDs, or lane constants. \emph{Example:} {\scriptsize\ttfamily \mnem{FADD.f32} \dst{r0}, \src{r1}, \src{k1f}} (\textbf{SPIR-V:} \texttt{OpConstant / BuiltIn}; \textbf{GLSL:} \texttt{float y = x + 1.0;})
  \item \hyperref[enum:flow_control_m]{\texttt{flow}}: 4-bit embedded flow-control \& scoreboard wait modifier (\texttt{bits 59:62}: \texttt{none}, \texttt{wait0..wait0126}, \texttt{wait}, \texttt{reconverge}, \texttt{discard}, \texttt{end}). \emph{Example:} {\scriptsize\ttfamily \mnem{FADD.f32}\imod{.wait0.reconverge} \dst{r0}, \src{r1}, \src{r2}} (\textbf{SPIR-V:} \texttt{OpControlBarrier / OpReturn}; \textbf{GLSL:} \texttt{// Scoreboard wait + warp reconvergence})
\end{tightitemize}
\clearpage

\subsection*{SEG\_SUB}
\addcontentsline{toc}{subsection}{SEG\_SUB}
Computes 64-bit segment pointer subtraction: widens \texttt{src1} to 64 bits and subtracts it from the 64-bit base pointer \texttt{src0} (\texttt{dest = src0 - widen64(src1)}), wrapping modulo $2^{64}$. When \texttt{.preserve\_null} is set and \texttt{src0 == 0}, returns \texttt{0} so null pointers remain null.
\par\textbf{Notes:}
\begin{tightitemize}
  \item Same null-preservation rule as \texttt{SEG\_ADD}.
  \item \textbf{Source references:} \texttt{bifrost/valhall/ISA.xml:2666-2684}, \texttt{isa-v9-v14.xml}.
\end{tightitemize}
\codebox{\mnem{SEG\_SUB}\imod{\{.flow\}}\imod{\{.preserve\_null\}} \dst{dest}, \src{src0}, \src{src1}\imod{\{.widen1\}}}
\textbf{Operands:}
\begin{tightitemize}\footnotesize
  \item \texttt{\dst{dest}} (\texttt{dest\_register\_index}): 64-bit destination register pair receiving the offset segment pointer.
  \item \texttt{\src{src0}}: 64-bit base segment pointer (preserved as \texttt{0} when \texttt{.preserve\_null} is set).
  \item \texttt{\src{src1}}: Offset widened to 64 bits and subtracted from \texttt{src0}.
\end{tightitemize}
\par\textbf{Closest SPIR-V Op:} \texttt{OpPtrAccessChain (negative offset)}
\par\textbf{Pseudo-GLSL (Operational Semantics):}
\codebox{if (flag\_preserve\_null \&\& src0 == 0ul) dest = 0ul;\\else                                   dest = src0 - widen64(src1);}
\textbf{Execution Context:}
\begin{tightitemize}\footnotesize
  \item \textbf{Thread/Workgroup Local Storage (\hyperref[sec:csf-tsd-fau]{\texttt{TSD\_k}}, \texttt{r24..r31}):} 32B \texttt{LocalStorage} descriptor (\texttt{tls\_size}, \texttt{wls\_instances}, \texttt{wls\_size\_base/scale}, \texttt{tls\_base\_pointer}, \texttt{wls\_base\_pointer}) exposed inline via \texttt{FauSpecial::tls\_ptr}/\texttt{wls\_ptr}.
  \item \textbf{Inline FAU \& Special Operands (\hyperref[sec:csf-tsd-fau]{\texttt{FAU\_k}}, \texttt{r8..r15}):} Any 8-bit source with \texttt{mode=2} reads zero-latency uniform \texttt{u0..u31} from \texttt{FAU\_k} (\texttt{slots64 = bits[63:56]}); \texttt{mode=3} reads \texttt{SmallConstant} (Table~1.2) or \texttt{FauSpecial} (\texttt{pc}, \texttt{lane\_id}, \texttt{lane\_width}, \texttt{tls\_ptr}, \texttt{wls\_ptr}, \texttt{blend\_descriptor\_0..7}).
\end{tightitemize}
\par\textbf{Worked Examples:}
\par\textbf{Example 1 (Non-null pointer decrement):}
\codebox{\mnem{SEG\_SUB.u64}\imod{.preserve\_null} \dst{[r0:r1]}, \src{[r2:r3]}, \src{r4}\imod{.w0}}
With \texttt{[r2:r3] = 0x1000}, \texttt{r4 = 0x20}: subtracts widened \texttt{r4} -\textgreater{} \texttt{[r0:r1] = 0x0FE0}.
\par\textbf{Example 2 (Null pointer preservation):}
\codebox{\mnem{SEG\_SUB.u64}\imod{.preserve\_null} \dst{[r0:r1]}, \src{[r2:r3]}, \src{r4}\imod{.w0}}
With \texttt{[r2:r3] = 0}, \texttt{r4 = 0x10}: \texttt{.preserve\_null} keeps \texttt{[r0:r1] = 0} rather than wrapping to \texttt{0xFFFFFFFF\_FFFFFFF0}.

\begin{center}\footnotesize
\begin{tabularx}{\linewidth}{@{}lllX@{}}
\toprule
Variant & Arch & Class & Opcode Fields \\
\midrule
\texttt{u64} & v9--v14 & \texttt{D2} & \texttt{opcode2=SEG\_ADD} \\
\bottomrule
\end{tabularx}
\end{center}
\begin{center}\footnotesize
\begin{tabularx}{0.88\linewidth}{@{}l l X@{}}
\toprule
Field Name & Bits & Type / Fixed Value \\
\midrule
\texttt{src0} & \texttt{0:7} & \texttt{SourceEncoding64} \\
\texttt{src1} & \texttt{8:15} & \texttt{SourceEncoding64} \\
\texttt{opcode2} & \texttt{16:19} & \texttt{exact=SEG\_ADD} \\
\textit{Reserved} & \texttt{20:25} & \texttt{zero} \\
\hyperref[enum:widen_int_m]{\texttt{widen1}} & \texttt{26:29} & \texttt{widen\_int\_m (def none)} \\
\textit{Reserved} & \texttt{30:37} & \texttt{zero} \\
\hyperref[enum:neg_m]{\texttt{neg}} & \texttt{38} & \texttt{exact=neg} \\
\hyperref[enum:segadd_preserve_null_m]{\texttt{preserve\_null}} & \texttt{39} & \texttt{segadd\_preserve\_null\_m (def none)} \\
\texttt{dest\_register\_index} & \texttt{40:45} & \texttt{uint} \\
\texttt{double\_width\_instruction} & \texttt{46:47} & \texttt{exact=3} \\
\texttt{datatype} & \texttt{48:51} & \texttt{exact=u64} \\
\texttt{instruction\_class} & \texttt{52:56} & \texttt{exact=D2} \\
\texttt{fau\_page\_index} & \texttt{57:58} & \texttt{uint} \\
\hyperref[enum:flow_control_m]{\texttt{flow}} & \texttt{59:62} & \texttt{flow\_control\_m} \\
\textit{Reserved} & \texttt{63} & \texttt{zero} \\
\bottomrule
\end{tabularx}
\end{center}
\textbf{Fields Explained:}
\begin{tightitemize}\footnotesize
  \item \texttt{src0}: Primary source operand (\texttt{r0..r63}, \texttt{u0..u63} uniform, \texttt{k*} FAU constant, or \texttt{t0..t2} temporary; \texttt{\textasciicircum{}} marks discard/last-use). \emph{Example:} {\scriptsize\ttfamily \mnem{FADD.f32} \dst{r0}, \src{r1}, \src{r2}} (\textbf{SPIR-V:} \texttt{OpFAdd (Operand 0)}; \textbf{GLSL:} \texttt{float c = a + b;})
  \item \texttt{src1}: Secondary source operand (register, uniform, or FAU constant; in shifts, supplies the shift count byte). \emph{Example:} {\scriptsize\ttfamily \mnem{LSHIFT.i32} \dst{r0}, \src{r1}, \src{r2}\imod{.b0}} (\textbf{SPIR-V:} \texttt{OpShiftLeftLogical (Shift)}; \textbf{GLSL:} \texttt{uint y = a \textless{}\textless{} s;})
  \item \texttt{opcode2}: Secondary sub-opcode further specializing instructions sharing \texttt{instruction\_class} and \texttt{opcode1} (e.g., conversions, SFU tables, moves). \emph{Example:} {\scriptsize\ttfamily \mnem{F32\_TO\_S32}\imod{.round\_zero} \dst{r0}, \src{r1}} (\textbf{SPIR-V:} \texttt{OpConvertFToS}; \textbf{GLSL:} \texttt{int y = int(x);})
  \item \hyperref[enum:widen_int_m]{\texttt{widen1}}: Stage A widening extractor on \texttt{src1} (promotes narrow byte/half/word lane before Stage B). \emph{Example:} {\scriptsize\ttfamily \mnem{IMUL.u64} \dst{[r0:r1]}, \src{[r2:r3]}\imod{.w0}, \src{[r4:r5]}\imod{.w0}} (\textbf{SPIR-V:} \texttt{OpUConvert / OpUMulExtended}; \textbf{GLSL:} \texttt{uint64\_t p = uint64\_t(a) * uint64\_t(b);})
  \item \hyperref[enum:neg_m]{\texttt{neg}}: Negation flag on segmented pointer arithmetic (\texttt{SEG\_ADD} vs \texttt{SEG\_SUB}: subtracts widened \texttt{src1} when set) or unary FP negation. \emph{Example:} {\scriptsize\ttfamily \mnem{SEG\_SUB.u64} \dst{[r0:r1]}, \src{[r2:r3]}, \src{r4}\imod{.w0}} (\textbf{SPIR-V:} \texttt{OpPtrAccessChain (negative index)}; \textbf{GLSL:} \texttt{shared\_arr[idx - off]})
  \item \hyperref[enum:segadd_preserve_null_m]{\texttt{preserve\_null}}: Preserves NULL (0) pointer on segmented address add/sub instead of adding segment base. \emph{Example:} {\scriptsize\ttfamily \mnem{SEG\_ADD}\imod{.preserve\_null} \dst{[r0:r1]}, \src{r2}} (\textbf{SPIR-V:} \texttt{OpPtrAccessChain (PhysicalStorageBuffer)}; \textbf{GLSL:} \texttt{if (ptr != null) val = ptr.x;})
  \item \texttt{dest\_register\_index}: 6-bit destination GPR index (\texttt{r0..r63}) receiving the Stage C writeback result (even-aligned pair \texttt{[rN:rN+1]} for 64-bit ops). \emph{Example:} {\scriptsize\ttfamily \mnem{FMA.f32} \dst{r0}, \src{r1}, \src{r2}, \src{r3}} (\textbf{SPIR-V:} \texttt{SSA Result ID}; \textbf{GLSL:} \texttt{float r0 = fma(r1, r2, r3);})
  \item \texttt{double\_width\_instruction}: Indicates 64-bit / dual-slot wide instruction encoding variant. \emph{Example:} {\scriptsize\ttfamily \mnem{IADD.u64} \dst{[r0:r1]}, \src{[r2:r3]}, \src{[r4:r5]}} (\textbf{SPIR-V:} \texttt{OpIAdd (64-bit)}; \textbf{GLSL:} \texttt{uint64\_t z = x + y;})
  \item \texttt{datatype}: Selects the lane data type (\texttt{f32}, \texttt{f16}, \texttt{v2f16}, \texttt{s32}, \texttt{u32}, \texttt{v2s16}, \texttt{v2u16}, \texttt{v4s8}, \texttt{v4u8}, \texttt{s64}, \texttt{u64}, etc.). \emph{Example:} {\scriptsize\ttfamily \mnem{IADD.v2s16} \dst{r0}, \src{r1}, \src{r2}} (\textbf{SPIR-V:} \texttt{OpTypeInt / OpTypeFloat}; \textbf{GLSL:} \texttt{i16vec2 c = a + b;})
  \item \texttt{instruction\_class}: 4-bit major instruction class selector (\texttt{bits 48:51}) partitioning the 64-bit Valhall opcode space (e.g., \texttt{FMA}, \texttt{CVT}, \texttt{SFU}, \texttt{LS}, \texttt{TEX}, \texttt{BRANCH}). \emph{Example:} {\scriptsize\ttfamily \mnem{FADD.f32} \dst{r0}, \src{r1}, \src{r2}} (\textbf{SPIR-V:} \texttt{N/A (encoding class)}; \textbf{GLSL:} \texttt{// Hardware execution pipe / decode class})
  \item \texttt{fau\_page\_index}: Selects the FAU (Fast Access Uniform) page (\texttt{0..3}) when a source operand references hardware small constants, thread IDs, or lane constants. \emph{Example:} {\scriptsize\ttfamily \mnem{FADD.f32} \dst{r0}, \src{r1}, \src{k1f}} (\textbf{SPIR-V:} \texttt{OpConstant / BuiltIn}; \textbf{GLSL:} \texttt{float y = x + 1.0;})
  \item \hyperref[enum:flow_control_m]{\texttt{flow}}: 4-bit embedded flow-control \& scoreboard wait modifier (\texttt{bits 59:62}: \texttt{none}, \texttt{wait0..wait0126}, \texttt{wait}, \texttt{reconverge}, \texttt{discard}, \texttt{end}). \emph{Example:} {\scriptsize\ttfamily \mnem{FADD.f32}\imod{.wait0.reconverge} \dst{r0}, \src{r1}, \src{r2}} (\textbf{SPIR-V:} \texttt{OpControlBarrier / OpReturn}; \textbf{GLSL:} \texttt{// Scoreboard wait + warp reconvergence})
\end{tightitemize}
\clearpage

\subsection*{SHADD}
\addcontentsline{toc}{subsection}{SHADD}
Computes 64-bit shift-and-add address arithmetic (\texttt{s64} or \texttt{u64}): widens \texttt{src1} to 64 bits, shifts it left by the 3-bit immediate \texttt{\#shift} ($0 \dots 7$), and adds it to the 64-bit base \texttt{src0} (\texttt{dest = src0 + (widen64(src1) \textless{}\textless{} \#shift)}), wrapping modulo $2^{64}$.
\par\textbf{Notes:}
\begin{tightitemize}
  \item Only \texttt{s64} and \texttt{u64} variants.
  \item Wraps modulo \texttt{2\textasciicircum{}64}.
  \item \textbf{Source references:} \texttt{bifrost/valhall/ISA.xml:2686-2707}, \texttt{bifrost/bir.c:80,128-130}, \texttt{isa-v9-v14.xml}.
\end{tightitemize}
\codebox{\mnem{SHADD}\imod{\{.flow\}} \dst{dest}, \src{src0}, \src{src1}\imod{\{.widen1\}}, \imm{\#shift}}
\textbf{Operands:}
\begin{tightitemize}\footnotesize
  \item \texttt{\dst{dest}} (\texttt{dest\_register\_index}): 64-bit destination register pair receiving \texttt{src0 + (widen64(src1) \textless{}\textless{} \#shift)}.
  \item \texttt{\src{src0}}: 64-bit base address or integer.
  \item \texttt{\src{src1}}: Index/offset widened to 64 bits before shifting.
  \item \texttt{\imm{\#shift}} (\texttt{shift}): 3-bit unsigned left-shift amount (\texttt{0..7}, e.g. \texttt{\#2} for 4-byte stride).
\end{tightitemize}
\par\textbf{Closest SPIR-V Op:} \texttt{OpIAdd + OpShiftLeftLogical (64-bit addr)}
\par\textbf{Pseudo-GLSL (Operational Semantics):}
\codebox{dest = src0 + (widen64(src1) \textless{}\textless{} \#shift); // \#shift in 0..7, wraps mod 2\textasciicircum{}64}
\textbf{Execution Context:}
\begin{tightitemize}\footnotesize
  \item \textbf{Compute stage preloads (\hyperref[sec:csf-spd]{\texttt{SPD\_k}}, \texttt{r16..r23}):} Hardware pre-fills \texttt{r55..r57} with \texttt{WorkGroupID.xyz + JOB\_OFFSET} (from \hyperref[sec:csf-regs]{\texttt{COMPUTESR::JOB\_OFFSET\_X/Y/Z}} \texttt{r34..r36}) and \texttt{r60..r62} with \texttt{LocalInvocationID.xyz} or \texttt{GlobalInvocationID.xyz} (from \texttt{WG\_SIZE} \texttt{r33} \& \texttt{JOB\_SIZE} \texttt{r37..r39}, selected by \texttt{ShaderProgram::preload}) so compute shaders can derive thread/element indices without memory loads.
  \item \textbf{Vertex stage preloads (\hyperref[sec:csf-spd]{\texttt{SPD\_0}/\texttt{SPD\_1}}, \texttt{r16..r19}):} The hardware Index Fetcher reads \hyperref[sec:csf-regs]{\texttt{IDVSSR::INDEX\_BUFFER}} (\texttt{r54:r55}) and pre-fills \texttt{r59, r60} with \texttt{VertexIndex} (\texttt{+ VERTEX\_OFFSET + GLOBAL\_ATTR\_OFF}) and \texttt{r61} with \texttt{InstanceIndex} (\texttt{+ INSTANCE\_OFFSET}) so Position and Varying shaders can pull vertex attributes directly.
  \item \textbf{Inline FAU \& Special Operands (\hyperref[sec:csf-tsd-fau]{\texttt{FAU\_k}}, \texttt{r8..r15}):} Any 8-bit source with \texttt{mode=2} reads zero-latency uniform \texttt{u0..u31} from \texttt{FAU\_k} (\texttt{slots64 = bits[63:56]}); \texttt{mode=3} reads \texttt{SmallConstant} (Table~1.2) or \texttt{FauSpecial} (\texttt{pc}, \texttt{lane\_id}, \texttt{lane\_width}, \texttt{tls\_ptr}, \texttt{wls\_ptr}, \texttt{blend\_descriptor\_0..7}).
\end{tightitemize}
\par\textbf{Worked Examples:}
\par\textbf{Example 1 (Unsigned array stride addressing):}
\codebox{\mnem{SHADD.u64} \dst{[r0:r1]}, \src{[r2:r3]}, \src{r4}\imod{.w0}, \imm{\#2}}
With \texttt{[r2:r3] = 0x1000}, \texttt{r4 = 5}, \texttt{\#shift = 2}: \texttt{0x1000 + (5 \textless{}\textless{} 2) = 0x1014} (indexes 4-byte elements).
\par\textbf{Example 2 (Signed negative index stride):}
\codebox{\mnem{SHADD.s64} \dst{[r0:r1]}, \src{[r2:r3]}, \src{r4}\imod{.w0}, \imm{\#3}}
With \texttt{[r2:r3] = 0x2000}, \texttt{r4 = 0xFFFFFFFE} (\texttt{-2}), \texttt{\#shift = 3}: sign-extends \texttt{-2} to 64-bit, shifts \texttt{\textless{}\textless{} 3} (\texttt{-16}), adds -\textgreater{} \texttt{[r0:r1] = 0x1FF0}.
\par\textbf{Mined SPIR-V $\to$ Valhall ISA Example (Scaled 64-bit buffer address calculation for 32-bit element load --- Fragment shader, Mali-G710 (v10)):}
\codebox{\textbf{; SPIR-V (input)}\\\%ptr = OpAccessChain \%\_ptr\_Uniform\_uint \%ssbo \%int\_0 \%idx\\\%val = OpLoad \%uint \%ptr\\[3pt]\textbf{; Valhall ISA (malioc output)}\\SHADD.u64 [r0:r1], sample, [r0\textasciicircum{}:r1\textasciicircum{}].w0, \#0x2\\LOAD.i32.slot0.wait0 @r0, [r0\textasciicircum{}:r1\textasciicircum{}], \#0x0}
\texttt{SHADD.u64} fuses \texttt{base\_64 + (uint64(index\_32) \textless{}\textless{} 2)} into a single 64-bit address instruction, zero-extending the 32-bit element index via \texttt{.w0} and scaling by \texttt{sizeof(uint) == 4} (\texttt{\#0x2}).
\begin{center}\footnotesize
\begin{tabularx}{\linewidth}{@{}lllX@{}}
\toprule
Variant & Arch & Class & Opcode Fields \\
\midrule
\texttt{s64} & v9--v14 & \texttt{D2} & \texttt{opcode2=SHADD} \\
\texttt{u64} & v9--v14 & \texttt{D2} & \texttt{opcode2=SHADD} \\
\bottomrule
\end{tabularx}
\end{center}
\begin{center}\footnotesize
\begin{tabularx}{0.88\linewidth}{@{}l l X@{}}
\toprule
Field Name & Bits & Type / Fixed Value \\
\midrule
\texttt{src0} & \texttt{0:7} & \texttt{SourceEncoding64} \\
\texttt{src1} & \texttt{8:15} & \texttt{SourceEncoding64} \\
\texttt{opcode2} & \texttt{16:19} & \texttt{exact=SHADD} \\
\texttt{shift} & \texttt{20:22} & \texttt{uint} \\
\textit{Reserved} & \texttt{23:25} & \texttt{zero} \\
\hyperref[enum:widen_int_m]{\texttt{widen1}} & \texttt{26:29} & \texttt{widen\_int\_m (def none)} \\
\textit{Reserved} & \texttt{30:39} & \texttt{zero} \\
\texttt{dest\_register\_index} & \texttt{40:45} & \texttt{uint} \\
\texttt{double\_width\_instruction} & \texttt{46:47} & \texttt{exact=3} \\
\texttt{datatype} & \texttt{48:51} & \texttt{exact=s64} \\
\texttt{instruction\_class} & \texttt{52:56} & \texttt{exact=D2} \\
\texttt{fau\_page\_index} & \texttt{57:58} & \texttt{uint} \\
\hyperref[enum:flow_control_m]{\texttt{flow}} & \texttt{59:62} & \texttt{flow\_control\_m} \\
\textit{Reserved} & \texttt{63} & \texttt{zero} \\
\bottomrule
\end{tabularx}
\end{center}
\textbf{Fields Explained:}
\begin{tightitemize}\footnotesize
  \item \texttt{src0}: Primary source operand (\texttt{r0..r63}, \texttt{u0..u63} uniform, \texttt{k*} FAU constant, or \texttt{t0..t2} temporary; \texttt{\textasciicircum{}} marks discard/last-use). \emph{Example:} {\scriptsize\ttfamily \mnem{FADD.f32} \dst{r0}, \src{r1}, \src{r2}} (\textbf{SPIR-V:} \texttt{OpFAdd (Operand 0)}; \textbf{GLSL:} \texttt{float c = a + b;})
  \item \texttt{src1}: Secondary source operand (register, uniform, or FAU constant; in shifts, supplies the shift count byte). \emph{Example:} {\scriptsize\ttfamily \mnem{LSHIFT.i32} \dst{r0}, \src{r1}, \src{r2}\imod{.b0}} (\textbf{SPIR-V:} \texttt{OpShiftLeftLogical (Shift)}; \textbf{GLSL:} \texttt{uint y = a \textless{}\textless{} s;})
  \item \texttt{opcode2}: Secondary sub-opcode further specializing instructions sharing \texttt{instruction\_class} and \texttt{opcode1} (e.g., conversions, SFU tables, moves). \emph{Example:} {\scriptsize\ttfamily \mnem{F32\_TO\_S32}\imod{.round\_zero} \dst{r0}, \src{r1}} (\textbf{SPIR-V:} \texttt{OpConvertFToS}; \textbf{GLSL:} \texttt{int y = int(x);})
  \item \texttt{shift}: Inline shift count (0--31) in \texttt{SHADD} / shift-add instructions. \emph{Example:} {\scriptsize\ttfamily \mnem{SHADD.u64} \dst{[r0:r1]}, \src{[r2:r3]}, \src{r4}, \imm{\#2}} (\textbf{SPIR-V:} \texttt{OpPtrAccessChain}; \textbf{GLSL:} \texttt{uint x = ssbo.arr[i]; // base + (i\textless{}\textless{}2)})
  \item \hyperref[enum:widen_int_m]{\texttt{widen1}}: Stage A widening extractor on \texttt{src1} (promotes narrow byte/half/word lane before Stage B). \emph{Example:} {\scriptsize\ttfamily \mnem{IMUL.u64} \dst{[r0:r1]}, \src{[r2:r3]}\imod{.w0}, \src{[r4:r5]}\imod{.w0}} (\textbf{SPIR-V:} \texttt{OpUConvert / OpUMulExtended}; \textbf{GLSL:} \texttt{uint64\_t p = uint64\_t(a) * uint64\_t(b);})
  \item \texttt{dest\_register\_index}: 6-bit destination GPR index (\texttt{r0..r63}) receiving the Stage C writeback result (even-aligned pair \texttt{[rN:rN+1]} for 64-bit ops). \emph{Example:} {\scriptsize\ttfamily \mnem{FMA.f32} \dst{r0}, \src{r1}, \src{r2}, \src{r3}} (\textbf{SPIR-V:} \texttt{SSA Result ID}; \textbf{GLSL:} \texttt{float r0 = fma(r1, r2, r3);})
  \item \texttt{double\_width\_instruction}: Indicates 64-bit / dual-slot wide instruction encoding variant. \emph{Example:} {\scriptsize\ttfamily \mnem{IADD.u64} \dst{[r0:r1]}, \src{[r2:r3]}, \src{[r4:r5]}} (\textbf{SPIR-V:} \texttt{OpIAdd (64-bit)}; \textbf{GLSL:} \texttt{uint64\_t z = x + y;})
  \item \texttt{datatype}: Selects the lane data type (\texttt{f32}, \texttt{f16}, \texttt{v2f16}, \texttt{s32}, \texttt{u32}, \texttt{v2s16}, \texttt{v2u16}, \texttt{v4s8}, \texttt{v4u8}, \texttt{s64}, \texttt{u64}, etc.). \emph{Example:} {\scriptsize\ttfamily \mnem{IADD.v2s16} \dst{r0}, \src{r1}, \src{r2}} (\textbf{SPIR-V:} \texttt{OpTypeInt / OpTypeFloat}; \textbf{GLSL:} \texttt{i16vec2 c = a + b;})
  \item \texttt{instruction\_class}: 4-bit major instruction class selector (\texttt{bits 48:51}) partitioning the 64-bit Valhall opcode space (e.g., \texttt{FMA}, \texttt{CVT}, \texttt{SFU}, \texttt{LS}, \texttt{TEX}, \texttt{BRANCH}). \emph{Example:} {\scriptsize\ttfamily \mnem{FADD.f32} \dst{r0}, \src{r1}, \src{r2}} (\textbf{SPIR-V:} \texttt{N/A (encoding class)}; \textbf{GLSL:} \texttt{// Hardware execution pipe / decode class})
  \item \texttt{fau\_page\_index}: Selects the FAU (Fast Access Uniform) page (\texttt{0..3}) when a source operand references hardware small constants, thread IDs, or lane constants. \emph{Example:} {\scriptsize\ttfamily \mnem{FADD.f32} \dst{r0}, \src{r1}, \src{k1f}} (\textbf{SPIR-V:} \texttt{OpConstant / BuiltIn}; \textbf{GLSL:} \texttt{float y = x + 1.0;})
  \item \hyperref[enum:flow_control_m]{\texttt{flow}}: 4-bit embedded flow-control \& scoreboard wait modifier (\texttt{bits 59:62}: \texttt{none}, \texttt{wait0..wait0126}, \texttt{wait}, \texttt{reconverge}, \texttt{discard}, \texttt{end}). \emph{Example:} {\scriptsize\ttfamily \mnem{FADD.f32}\imod{.wait0.reconverge} \dst{r0}, \src{r1}, \src{r2}} (\textbf{SPIR-V:} \texttt{OpControlBarrier / OpReturn}; \textbf{GLSL:} \texttt{// Scoreboard wait + warp reconvergence})
\end{tightitemize}
\clearpage

\section{Bitwise \& Shift-Logic (28 mnemonics)}
\subsection*{AND}
\addcontentsline{toc}{subsection}{AND}
Computes bitwise AND independently across each SIMD element $i \in \{0, \dots, N-1\}$ ($N=1$ for \texttt{i32}/\texttt{i64}, $N=2$ for \texttt{v2i16}, $N=4$ for \texttt{v4i8} on v9--v10): \texttt{res = widen0(src0)[i] \& (not2 ? \textasciitilde{}src2[i] : src2[i])}, then writes \texttt{not\_result ? \textasciitilde{}res : res} to \texttt{dest[i]} (\texttt{src1} unused/fixed 0).
\par\textbf{Notes:}
\begin{tightitemize}
  \item Second operand is \texttt{src2} (NOT \texttt{src1}) in XML encoding (\texttt{encode\_v9.rs:2982-2985}).
  \item Supports \texttt{i32}, \texttt{i64}, \texttt{v2i16}, \texttt{v4i8}.
  \item \textbf{Source references:} \texttt{kraid/ops.rs:3944-4025}, \texttt{encode\_v9.rs:2916-3093}, \texttt{hw\_tests.rs:2148-2198}.
\end{tightitemize}
\codebox{\mnem{AND}\imod{\{.flow\}}\imod{\{.not\_result\}} \dst{dest}\imod{\{.dest\_width\}}, \src{src0}\imod{\{.widen0\}}, \src{src2}\imod{\{.not2\}}\\\mnem{AND}\imod{\{.flow\}}\imod{\{.not\_result\}} \dst{dest}, \src{src0}\imod{\{.widen0\}}, \src{src2}\imod{\{.not2\}}\\\mnem{AND}\imod{\{.flow\}}\imod{\{.not\_result\}} \dst{dest}\imod{\{.dest\_width\}}, \src{src0}\imod{\{.swz0\}}, \src{src2}\imod{\{.not2\}}\\\mnem{AND}\imod{\{.flow\}}\imod{\{.not\_result\}} \dst{dest}\imod{\{.dest\_width\}}, \src{src0}\imod{\{.lanes0\}}, \src{src2}\imod{\{.not2\}}}
\textbf{Operands:}
\begin{tightitemize}\footnotesize
  \item \texttt{\dst{dest}} (\texttt{dest\_register\_index}, \texttt{dest\_width}): Destination GPR receiving \texttt{src0 \& src2} (or \texttt{\textasciitilde{}(src0 \& src2)} with \texttt{.not\_result}).
  \item \texttt{\src{src0}}: First bitwise operand.
  \item \texttt{\src{src2}}: Second bitwise operand (note: encoded in \texttt{src2}, not \texttt{src1}; supports \texttt{.not2}).
\end{tightitemize}
\par\textbf{Closest SPIR-V Op:} \texttt{OpBitwiseAnd}
\par\textbf{Pseudo-GLSL (Operational Semantics):}
\codebox{for (int i = 0; i \textless{} N; ++i) \{\\    auto b = flag\_not2 ? \textasciitilde{}src2[i] : src2[i];\\    auto r = widen0(src0)[i] \& b;\\    dest[i] = flag\_not\_result ? \textasciitilde{}r : r;\\\}}
\textbf{Execution Context:}
\begin{tightitemize}\footnotesize
  \item \textbf{Compute stage preloads (\hyperref[sec:csf-spd]{\texttt{SPD\_k}}, \texttt{r16..r23}):} Hardware pre-fills \texttt{r55..r57} with \texttt{WorkGroupID.xyz + JOB\_OFFSET} (from \hyperref[sec:csf-regs]{\texttt{COMPUTESR::JOB\_OFFSET\_X/Y/Z}} \texttt{r34..r36}) and \texttt{r60..r62} with \texttt{LocalInvocationID.xyz} or \texttt{GlobalInvocationID.xyz} (from \texttt{WG\_SIZE} \texttt{r33} \& \texttt{JOB\_SIZE} \texttt{r37..r39}, selected by \texttt{ShaderProgram::preload}) so compute shaders can derive thread/element indices without memory loads.
  \item \textbf{Vertex stage preloads (\hyperref[sec:csf-spd]{\texttt{SPD\_0}/\texttt{SPD\_1}}, \texttt{r16..r19}):} The hardware Index Fetcher reads \hyperref[sec:csf-regs]{\texttt{IDVSSR::INDEX\_BUFFER}} (\texttt{r54:r55}) and pre-fills \texttt{r59, r60} with \texttt{VertexIndex} (\texttt{+ VERTEX\_OFFSET + GLOBAL\_ATTR\_OFF}) and \texttt{r61} with \texttt{InstanceIndex} (\texttt{+ INSTANCE\_OFFSET}) so Position and Varying shaders can pull vertex attributes directly.
  \item \textbf{Inline FAU \& Special Operands (\hyperref[sec:csf-tsd-fau]{\texttt{FAU\_k}}, \texttt{r8..r15}):} Any 8-bit source with \texttt{mode=2} reads zero-latency uniform \texttt{u0..u31} from \texttt{FAU\_k} (\texttt{slots64 = bits[63:56]}); \texttt{mode=3} reads \texttt{SmallConstant} (Table~1.2) or \texttt{FauSpecial} (\texttt{pc}, \texttt{lane\_id}, \texttt{lane\_width}, \texttt{tls\_ptr}, \texttt{wls\_ptr}, \texttt{blend\_descriptor\_0..7}).
\end{tightitemize}
\par\textbf{Worked Examples:}
\par\textbf{Example 1 (Bit-clear using src2.not):}
\codebox{\mnem{AND.i32} \dst{r0}, \src{r1}, \src{r2}\imod{.not}}
With \texttt{r1 = 0xFF, r2 = 0x0F}: Stage A inverts \texttt{src2} (\texttt{\textasciitilde{}0x0F = 0xFFFFFFF0}); Stage B computes \texttt{0xFF \& 0xFFFFFFF0 = 0xF0}.
\par\textbf{Example 2 (NAND via .not\_result):}
\codebox{\mnem{AND.i32}\imod{.not\_result} \dst{r0}, \src{r1}, \src{r2}}
With \texttt{r1 = 0x0000FFFF, r2 = 0x00FF00FF}: \texttt{r1 \& r2 = 0x000000FF}, then \texttt{.not\_result} complements -\textgreater{} \texttt{r0 = 0xFFFFFF00}.

\begin{center}\footnotesize
\begin{tabularx}{\linewidth}{@{}lllX@{}}
\toprule
Variant & Arch & Class & Opcode Fields \\
\midrule
\texttt{(none)} & v9--v14 & \texttt{R3} & \texttt{opcode1=SHIFT\_BITWISE\_I32} \\
\texttt{(none)} & v9--v14 & \texttt{R3} & \texttt{opcode1=SHIFT\_BITWISE\_I32} \\
\texttt{(none)} & v9--v14 & \texttt{R3} & \texttt{opcode1=LOGIC\_SHIFT\_I32} \\
\texttt{(none)} & v9--v14 & \texttt{D3} & \texttt{opcode1=SHIFT\_BITWISE\_I64} \\
\texttt{(none)} & v9--v14 & \texttt{D3} & \texttt{opcode1=LOGIC\_SHIFT\_I64} \\
\texttt{(none)} & v9--v14 & \texttt{R3} & \texttt{opcode1=SHIFT\_BITWISE\_V2I16} \\
\texttt{(none)} & v9--v14 & \texttt{R3} & \texttt{opcode1=LOGIC\_SHIFT\_V2I16} \\
\texttt{(none)} & v9--v14 & \texttt{R3} & \texttt{opcode1=SHIFT\_BITWISE\_V4I8} \\
\texttt{(none)} & v9--v14 & \texttt{R3} & \texttt{opcode1=LOGIC\_SHIFT\_V4I8} \\
\bottomrule
\end{tabularx}
\end{center}
\begin{center}\footnotesize
\begin{tabularx}{0.88\linewidth}{@{}l l X@{}}
\toprule
Field Name & Bits & Type / Fixed Value \\
\midrule
\texttt{src0} & \texttt{0:7} & \texttt{SourceEncoding} \\
\texttt{src1} & \texttt{8:15} & \texttt{exact=192} \\
\texttt{src2} & \texttt{16:23} & \texttt{SourceEncoding} \\
\texttt{bitwise\_operation} & \texttt{24:25} & \texttt{exact=and} \\
\hyperref[enum:shift_bitwise_i32_lanes_m]{\texttt{lanes1}} & \texttt{26:29} & \texttt{exact=b0} \\
\hyperref[enum:not_m]{\texttt{not\_result}} & \texttt{30} & \texttt{not\_m (def none)} \\
\textit{Reserved} & \texttt{31} & \texttt{zero} \\
\texttt{shift\_operation} & \texttt{32:34} & \texttt{exact=lshift} \\
\hyperref[enum:not_m]{\texttt{not2}} & \texttt{35} & \texttt{not\_m (def none)} \\
\hyperref[enum:widen_int_m]{\texttt{widen0}} & \texttt{36:39} & \texttt{widen\_int\_m (def none)} \\
\texttt{dest\_register\_index} & \texttt{40:45} & \texttt{uint} \\
\hyperref[enum:dest_width_m]{\texttt{dest\_width}} & \texttt{46:47} & \texttt{dest\_width\_m} \\
\texttt{opcode1} & \texttt{48:51} & \texttt{exact=SHIFT\_BITWISE\_I32} \\
\texttt{instruction\_class} & \texttt{52:56} & \texttt{exact=R3} \\
\texttt{fau\_page\_index} & \texttt{57:58} & \texttt{uint} \\
\hyperref[enum:flow_control_m]{\texttt{flow}} & \texttt{59:62} & \texttt{flow\_control\_m} \\
\textit{Reserved} & \texttt{63} & \texttt{zero} \\
\bottomrule
\end{tabularx}
\end{center}
\textbf{Fields Explained:}
\begin{tightitemize}\footnotesize
  \item \texttt{src0}: Primary source operand (\texttt{r0..r63}, \texttt{u0..u63} uniform, \texttt{k*} FAU constant, or \texttt{t0..t2} temporary; \texttt{\textasciicircum{}} marks discard/last-use). \emph{Example:} {\scriptsize\ttfamily \mnem{FADD.f32} \dst{r0}, \src{r1}, \src{r2}} (\textbf{SPIR-V:} \texttt{OpFAdd (Operand 0)}; \textbf{GLSL:} \texttt{float c = a + b;})
  \item \texttt{src1}: Secondary source operand (register, uniform, or FAU constant; in shifts, supplies the shift count byte). \emph{Example:} {\scriptsize\ttfamily \mnem{LSHIFT.i32} \dst{r0}, \src{r1}, \src{r2}\imod{.b0}} (\textbf{SPIR-V:} \texttt{OpShiftLeftLogical (Shift)}; \textbf{GLSL:} \texttt{uint y = a \textless{}\textless{} s;})
  \item \texttt{src2}: Third source operand (FMA addend, \texttt{MUX}/\texttt{CSEL} selector, \texttt{IADDC} carry-in, or bitwise mask in fused shift-logic). \emph{Example:} {\scriptsize\ttfamily \mnem{FMA.f32} \dst{r0}, \src{r1}, \src{r2}, \src{r3}} (\textbf{SPIR-V:} \texttt{GLSL.std.450 Fma (Addend)}; \textbf{GLSL:} \texttt{float d = fma(a, b, c);})
  \item \texttt{bitwise\_operation}: Selects the bitwise ALU op (\texttt{and}, \texttt{or}, \texttt{xor}, \texttt{lshift\_and}, etc.) fused with shifts. \emph{Example:} {\scriptsize\ttfamily \mnem{LSHIFT\_OR.i32} \dst{r0}, \src{r1}, \src{r2}, \src{r3}} (\textbf{SPIR-V:} \texttt{OpShiftLeftLogical + OpBitwiseOr}; \textbf{GLSL:} \texttt{uint y = (a \textless{}\textless{} s) \textbar{} b;})
  \item \hyperref[enum:shift_bitwise_i32_lanes_m]{\texttt{lanes1}}: Stage A byte-lane selector on \texttt{src1} (selects shift-count byte \texttt{b0..b3} on 32/64-bit shifts, \texttt{bXY} on \texttt{v2i16} shifts, or 4-byte swizzle on \texttt{v4*8}). \emph{Example:} {\scriptsize\ttfamily \mnem{LSHIFT.i32} \dst{r0}, \src{r1}, \src{r2}\imod{.b1}} (\textbf{SPIR-V:} \texttt{OpBitFieldUExtract + OpShiftLeftLogical}; \textbf{GLSL:} \texttt{uint y = a \textless{}\textless{} ((s \textgreater{}\textgreater{} 8) \& 0xFFu);})
  \item \hyperref[enum:not_m]{\texttt{not\_result}}: Bitwise-inverts the final output of a bitwise/shift-bitwise instruction (\textasciitilde{}result). \emph{Example:} {\scriptsize\ttfamily \mnem{AND.i32}\imod{.not\_result} \dst{r0}, \src{r1}, \src{r2}} (\textbf{SPIR-V:} \texttt{OpBitwiseAnd + OpNot}; \textbf{GLSL:} \texttt{uint y = \textasciitilde{}(a \& b);})
  \item \texttt{shift\_operation}: Selects shift direction/type (\texttt{lshift}, \texttt{rshift}, \texttt{arshift}) in fused shift-bitwise instructions. \emph{Example:} {\scriptsize\ttfamily \mnem{RSHIFT\_AND.i32} \dst{r0}, \src{r1}, \src{r2}, \src{r3}} (\textbf{SPIR-V:} \texttt{OpShiftRightLogical + OpBitwiseAnd}; \textbf{GLSL:} \texttt{uint y = (a \textgreater{}\textgreater{} s) \& m;})
  \item \hyperref[enum:not_m]{\texttt{not2}}: Stage A bitwise complement (\texttt{\textasciitilde{}src2}) applied to \texttt{src2} before the bitwise/shift-logic operation. \emph{Example:} {\scriptsize\ttfamily \mnem{AND.i32} \dst{r0}, \src{r1}, \src{r2}\imod{.not}} (\textbf{SPIR-V:} \texttt{OpNot + OpBitwiseAnd}; \textbf{GLSL:} \texttt{uint y = a \& \textasciitilde{}b;})
  \item \hyperref[enum:widen_int_m]{\texttt{widen0}}: Stage A widening extractor on \texttt{src0} (promotes FP16 \texttt{h0}/\texttt{h1} to FP32, or sign/zero-extends \texttt{b0..b3}/\texttt{h0..h1}/\texttt{w0..w1} to 32/64-bit before \texttt{abs}/\texttt{neg}). \emph{Example:} {\scriptsize\ttfamily \mnem{IADD.s32} \dst{r0}, \src{r1}\imod{.h0}, \src{r2}} (\textbf{SPIR-V:} \texttt{OpSConvert / OpFConvert}; \textbf{GLSL:} \texttt{int c = int(int16\_t(a)) + b;})
  \item \texttt{dest\_register\_index}: 6-bit destination GPR index (\texttt{r0..r63}) receiving the Stage C writeback result (even-aligned pair \texttt{[rN:rN+1]} for 64-bit ops). \emph{Example:} {\scriptsize\ttfamily \mnem{FMA.f32} \dst{r0}, \src{r1}, \src{r2}, \src{r3}} (\textbf{SPIR-V:} \texttt{SSA Result ID}; \textbf{GLSL:} \texttt{float r0 = fma(r1, r2, r3);})
  \item \hyperref[enum:dest_width_m]{\texttt{dest\_width}}: Stage C destination lane write-mask (\texttt{none} = full 32-bit, \texttt{h0} = low 16-bit, \texttt{h1} = high 16-bit without bit-shifting). \emph{Example:} {\scriptsize\ttfamily \mnem{FADD.v2f16} \dst{r0}\imod{.h1}, \src{r1}, \src{r2}} (\textbf{SPIR-V:} \texttt{OpVectorInsertDynamic / WriteMask}; \textbf{GLSL:} \texttt{dst.y = a.y + b.y;})
  \item \texttt{opcode1}: Primary sub-opcode within the instruction class uniquely identifying the operation family. \emph{Example:} {\scriptsize\ttfamily \mnem{IADD.s32} \dst{r0}, \src{r1}, \src{r2}} (\textbf{SPIR-V:} \texttt{OpIAdd (operation selector)}; \textbf{GLSL:} \texttt{// Primary opcode bits})
  \item \texttt{instruction\_class}: 4-bit major instruction class selector (\texttt{bits 48:51}) partitioning the 64-bit Valhall opcode space (e.g., \texttt{FMA}, \texttt{CVT}, \texttt{SFU}, \texttt{LS}, \texttt{TEX}, \texttt{BRANCH}). \emph{Example:} {\scriptsize\ttfamily \mnem{FADD.f32} \dst{r0}, \src{r1}, \src{r2}} (\textbf{SPIR-V:} \texttt{N/A (encoding class)}; \textbf{GLSL:} \texttt{// Hardware execution pipe / decode class})
  \item \texttt{fau\_page\_index}: Selects the FAU (Fast Access Uniform) page (\texttt{0..3}) when a source operand references hardware small constants, thread IDs, or lane constants. \emph{Example:} {\scriptsize\ttfamily \mnem{FADD.f32} \dst{r0}, \src{r1}, \src{k1f}} (\textbf{SPIR-V:} \texttt{OpConstant / BuiltIn}; \textbf{GLSL:} \texttt{float y = x + 1.0;})
  \item \hyperref[enum:flow_control_m]{\texttt{flow}}: 4-bit embedded flow-control \& scoreboard wait modifier (\texttt{bits 59:62}: \texttt{none}, \texttt{wait0..wait0126}, \texttt{wait}, \texttt{reconverge}, \texttt{discard}, \texttt{end}). \emph{Example:} {\scriptsize\ttfamily \mnem{FADD.f32}\imod{.wait0.reconverge} \dst{r0}, \src{r1}, \src{r2}} (\textbf{SPIR-V:} \texttt{OpControlBarrier / OpReturn}; \textbf{GLSL:} \texttt{// Scoreboard wait + warp reconvergence})
\end{tightitemize}
\clearpage

\subsection*{AND\_IMM}
\addcontentsline{toc}{subsection}{AND\_IMM}
Computes a 32-bit bitwise AND between \texttt{src0} and the inline 32-bit immediate \texttt{\#imm1w} (\texttt{dest = src0 \& \#imm1w}) on Valhall v11+. Does not support source swizzles or \texttt{.not\_result}; supports \texttt{dest\_width} halfword write masking.
\par\textbf{Notes:}
\begin{tightitemize}
  \item No \texttt{not\_result} or \texttt{src0} swizzle supported (\texttt{encode\_v9.rs:3043}).
  \item \textbf{Source references:} \texttt{kraid/encode\_v9.rs:2919-2926, 3036-3054}.
\end{tightitemize}
\codebox{\mnem{AND\_IMM}\imod{\{.flow\}} \dst{dest}\imod{\{.dest\_width\}}, \src{src0}, \imm{\#imm1w}}
\textbf{Operands:}
\begin{tightitemize}\footnotesize
  \item \texttt{\dst{dest}} (\texttt{dest\_register\_index}, \texttt{dest\_width}): 32-bit destination GPR receiving \texttt{src0 \& \#imm1w}.
  \item \texttt{\src{src0}}: 32-bit input operand.
  \item \texttt{\imm{\#imm1w}} (\texttt{imm1w}): Inline 32-bit bitwise AND mask.
\end{tightitemize}
\par\textbf{Closest SPIR-V Op:} \texttt{OpBitwiseAnd (constant)}
\par\textbf{Pseudo-GLSL (Operational Semantics):}
\codebox{dest = maskWrite(dest\_width, src0 \& \#imm1w);}
\textbf{Execution Context:}
\begin{tightitemize}\footnotesize
  \item \textbf{Compute stage preloads (\hyperref[sec:csf-spd]{\texttt{SPD\_k}}, \texttt{r16..r23}):} Hardware pre-fills \texttt{r55..r57} with \texttt{WorkGroupID.xyz + JOB\_OFFSET} (from \hyperref[sec:csf-regs]{\texttt{COMPUTESR::JOB\_OFFSET\_X/Y/Z}} \texttt{r34..r36}) and \texttt{r60..r62} with \texttt{LocalInvocationID.xyz} or \texttt{GlobalInvocationID.xyz} (from \texttt{WG\_SIZE} \texttt{r33} \& \texttt{JOB\_SIZE} \texttt{r37..r39}, selected by \texttt{ShaderProgram::preload}) so compute shaders can derive thread/element indices without memory loads.
  \item \textbf{Vertex stage preloads (\hyperref[sec:csf-spd]{\texttt{SPD\_0}/\texttt{SPD\_1}}, \texttt{r16..r19}):} The hardware Index Fetcher reads \hyperref[sec:csf-regs]{\texttt{IDVSSR::INDEX\_BUFFER}} (\texttt{r54:r55}) and pre-fills \texttt{r59, r60} with \texttt{VertexIndex} (\texttt{+ VERTEX\_OFFSET + GLOBAL\_ATTR\_OFF}) and \texttt{r61} with \texttt{InstanceIndex} (\texttt{+ INSTANCE\_OFFSET}) so Position and Varying shaders can pull vertex attributes directly.
  \item \textbf{Inline FAU \& Special Operands (\hyperref[sec:csf-tsd-fau]{\texttt{FAU\_k}}, \texttt{r8..r15}):} Any 8-bit source with \texttt{mode=2} reads zero-latency uniform \texttt{u0..u31} from \texttt{FAU\_k} (\texttt{slots64 = bits[63:56]}); \texttt{mode=3} reads \texttt{SmallConstant} (Table~1.2) or \texttt{FauSpecial} (\texttt{pc}, \texttt{lane\_id}, \texttt{lane\_width}, \texttt{tls\_ptr}, \texttt{wls\_ptr}, \texttt{blend\_descriptor\_0..7}).
\end{tightitemize}
\par\textbf{Worked Examples:}
\par\textbf{Example 1 (32-bit inline bitmask (v11+)):}
\codebox{\mnem{AND\_IMM.i32} \dst{r0}, \src{r1}, \imm{\#0x00FF00FF}}
With \texttt{r1 = 0x12345678}: masks with 32-bit inline immediate -\textgreater{} \texttt{r0 = 0x00340078} without consuming an FAU slot.
\par\textbf{Example 2 (Sign-bit extraction):}
\codebox{\mnem{AND\_IMM.i32} \dst{r0}, \src{r1}, \imm{\#0x80000000}}
With \texttt{r1 = 0xBF800000} (\texttt{-1.0f}): isolates IEEE sign bit -\textgreater{} \texttt{r0 = 0x80000000}.
\par\textbf{Mined SPIR-V $\to$ Valhall ISA Example (Inline 32-bit bitmask clear on v11+ (Immortalis-G925) --- Vertex shader, Immortalis-G925 (v14)):}
\codebox{\textbf{; SPIR-V (input)}\\\%masked = OpBitwiseAnd \%uint \%flags \%uint\_0xffffffdf\\\%eq     = OpIEqual \%bool \%masked \%expected\\[3pt]\textbf{; Valhall ISA (malioc output)}\\AND\_IMM.i32 r0, r0\textasciicircum{}, \#0xffffffdf\\CSEL.u32.eq r0, r0\textasciicircum{}, u12.w1, r2\textasciicircum{}, u12.w0}
On Valhall v11+, \texttt{OpBitwiseAnd} with a 32-bit constant mask that does not fit in FAU small constants lowers directly to a single \texttt{AND\_IMM.i32} instruction (clearing bit 5 via \texttt{\#0xffffffdf}) rather than \texttt{MOV\_IMM} + \texttt{AND}.
\begin{center}\footnotesize
\begin{tabularx}{\linewidth}{@{}lllX@{}}
\toprule
Variant & Arch & Class & Opcode Fields \\
\midrule
\texttt{(none)} & v9--v14 & \texttt{R1I} & \texttt{opcode1=AND\_IMM\_I32} \\
\bottomrule
\end{tabularx}
\end{center}
\begin{center}\footnotesize
\begin{tabularx}{0.88\linewidth}{@{}l l X@{}}
\toprule
Field Name & Bits & Type / Fixed Value \\
\midrule
\texttt{src0} & \texttt{0:7} & \texttt{SourceEncoding} \\
\texttt{imm1w} & \texttt{8:39} & \texttt{uint} \\
\texttt{dest\_register\_index} & \texttt{40:45} & \texttt{uint} \\
\hyperref[enum:dest_width_m]{\texttt{dest\_width}} & \texttt{46:47} & \texttt{dest\_width\_m} \\
\texttt{opcode1} & \texttt{48:51} & \texttt{exact=AND\_IMM\_I32} \\
\texttt{instruction\_class} & \texttt{52:56} & \texttt{exact=R1I} \\
\texttt{fau\_page\_index} & \texttt{57:58} & \texttt{uint} \\
\hyperref[enum:flow_control_m]{\texttt{flow}} & \texttt{59:62} & \texttt{flow\_control\_m} \\
\textit{Reserved} & \texttt{63} & \texttt{zero} \\
\bottomrule
\end{tabularx}
\end{center}
\textbf{Fields Explained:}
\begin{tightitemize}\footnotesize
  \item \texttt{src0}: Primary source operand (\texttt{r0..r63}, \texttt{u0..u63} uniform, \texttt{k*} FAU constant, or \texttt{t0..t2} temporary; \texttt{\textasciicircum{}} marks discard/last-use). \emph{Example:} {\scriptsize\ttfamily \mnem{FADD.f32} \dst{r0}, \src{r1}, \src{r2}} (\textbf{SPIR-V:} \texttt{OpFAdd (Operand 0)}; \textbf{GLSL:} \texttt{float c = a + b;})
  \item \texttt{imm1w}: 32-bit inline immediate operand encoded directly in the instruction word (bits 8:39). \emph{Example:} {\scriptsize\ttfamily \mnem{IADD\_IMM.i32} \dst{r0}, \src{r1}, \imm{\#0x10}} (\textbf{SPIR-V:} \texttt{OpIAdd (constant)}; \textbf{GLSL:} \texttt{int y = x + 16;})
  \item \texttt{dest\_register\_index}: 6-bit destination GPR index (\texttt{r0..r63}) receiving the Stage C writeback result (even-aligned pair \texttt{[rN:rN+1]} for 64-bit ops). \emph{Example:} {\scriptsize\ttfamily \mnem{FMA.f32} \dst{r0}, \src{r1}, \src{r2}, \src{r3}} (\textbf{SPIR-V:} \texttt{SSA Result ID}; \textbf{GLSL:} \texttt{float r0 = fma(r1, r2, r3);})
  \item \hyperref[enum:dest_width_m]{\texttt{dest\_width}}: Stage C destination lane write-mask (\texttt{none} = full 32-bit, \texttt{h0} = low 16-bit, \texttt{h1} = high 16-bit without bit-shifting). \emph{Example:} {\scriptsize\ttfamily \mnem{FADD.v2f16} \dst{r0}\imod{.h1}, \src{r1}, \src{r2}} (\textbf{SPIR-V:} \texttt{OpVectorInsertDynamic / WriteMask}; \textbf{GLSL:} \texttt{dst.y = a.y + b.y;})
  \item \texttt{opcode1}: Primary sub-opcode within the instruction class uniquely identifying the operation family. \emph{Example:} {\scriptsize\ttfamily \mnem{IADD.s32} \dst{r0}, \src{r1}, \src{r2}} (\textbf{SPIR-V:} \texttt{OpIAdd (operation selector)}; \textbf{GLSL:} \texttt{// Primary opcode bits})
  \item \texttt{instruction\_class}: 4-bit major instruction class selector (\texttt{bits 48:51}) partitioning the 64-bit Valhall opcode space (e.g., \texttt{FMA}, \texttt{CVT}, \texttt{SFU}, \texttt{LS}, \texttt{TEX}, \texttt{BRANCH}). \emph{Example:} {\scriptsize\ttfamily \mnem{FADD.f32} \dst{r0}, \src{r1}, \src{r2}} (\textbf{SPIR-V:} \texttt{N/A (encoding class)}; \textbf{GLSL:} \texttt{// Hardware execution pipe / decode class})
  \item \texttt{fau\_page\_index}: Selects the FAU (Fast Access Uniform) page (\texttt{0..3}) when a source operand references hardware small constants, thread IDs, or lane constants. \emph{Example:} {\scriptsize\ttfamily \mnem{FADD.f32} \dst{r0}, \src{r1}, \src{k1f}} (\textbf{SPIR-V:} \texttt{OpConstant / BuiltIn}; \textbf{GLSL:} \texttt{float y = x + 1.0;})
  \item \hyperref[enum:flow_control_m]{\texttt{flow}}: 4-bit embedded flow-control \& scoreboard wait modifier (\texttt{bits 59:62}: \texttt{none}, \texttt{wait0..wait0126}, \texttt{wait}, \texttt{reconverge}, \texttt{discard}, \texttt{end}). \emph{Example:} {\scriptsize\ttfamily \mnem{FADD.f32}\imod{.wait0.reconverge} \dst{r0}, \src{r1}, \src{r2}} (\textbf{SPIR-V:} \texttt{OpControlBarrier / OpReturn}; \textbf{GLSL:} \texttt{// Scoreboard wait + warp reconvergence})
\end{tightitemize}
\clearpage

\subsection*{ARSHIFT}
\addcontentsline{toc}{subsection}{ARSHIFT}
For each SIMD element $i \in \{0, \dots, N-1\}$ ($N=1$ for \texttt{i32}/\texttt{i64}, $N=2$ for \texttt{v2i16}, $N=4$ for \texttt{v4i8} on v9--v10, lane width $W$), computes arithmetic right shift \texttt{s = (int\_W(val) \textgreater{}\textgreater{} (amt \& (W - 1)))} where \texttt{val = src0[i]} and the 8-bit shift count \texttt{amt} is extracted from \texttt{src1} via \texttt{lanes1}; writes \texttt{not\_result ? \textasciitilde{}s : s} to \texttt{dest[i]}.
\par\textbf{Notes:}
\begin{tightitemize}
  \item Shift amount is masked modulo lane width (\texttt{wrapping\_shl}/\texttt{wrapping\_shr}/\texttt{rotate\_*}, \texttt{kraid/ops.rs:3984-4009}).
  \item On \texttt{v2i16}, \texttt{src1} is \texttt{V2I8} (\texttt{ops.rs:3953}) with byte lanes selected by \texttt{lanes1} (default \texttt{b02}: byte 0 -\textgreater{} lane 0, byte 2 -\textgreater{} lane 1).
  \item \textbf{Source references:} \texttt{kraid/ops.rs:3878-4025}, \texttt{encode\_v9.rs:2940-3093}, \texttt{hw\_tests.rs:2148-2198}.
\end{tightitemize}
\codebox{\mnem{ARSHIFT}\imod{\{.flow\}}\imod{\{.not\_result\}} \dst{dest}\imod{\{.dest\_width\}}, \src{src0}\imod{\{.widen0\}}, \src{src1}\imod{\{.lanes1\}}\\\mnem{ARSHIFT}\imod{\{.flow\}}\imod{\{.not\_result\}} \dst{dest}, \src{src0}\imod{\{.widen0\}}, \src{src1}\imod{\{.lanes1\}}\\\mnem{ARSHIFT}\imod{\{.flow\}}\imod{\{.not\_result\}} \dst{dest}\imod{\{.dest\_width\}}, \src{src0}\imod{\{.swz0\}}, \src{src1}\imod{\{.lanes1\}}\\\mnem{ARSHIFT}\imod{\{.flow\}}\imod{\{.not\_result\}} \dst{dest}\imod{\{.dest\_width\}}, \src{src0}\imod{\{.lanes0\}}, \src{src1}\imod{\{.lanes1\}}}
\textbf{Operands:}
\begin{tightitemize}\footnotesize
  \item \texttt{\dst{dest}} (\texttt{dest\_register\_index}, \texttt{dest\_width}): Destination GPR receiving the arithmetic right shift (inverted if \texttt{.not\_result}).
  \item \texttt{\src{src0}}: Input integer value to shift/rotate.
  \item \texttt{\src{src1}}: 8-bit shift/rotate count per lane (selected via \texttt{lanes1}, masked modulo lane width \texttt{W}).
\end{tightitemize}
\par\textbf{Closest SPIR-V Op:} \texttt{OpShiftRightArithmetic}
\par\textbf{Pseudo-GLSL (Operational Semantics):}
\codebox{for (int i = 0; i \textless{} N; ++i) \{\\    uint s = extractByte(src1, lanes1, i) \& (W - 1);\\    auto r = (int(src0[i]) \textgreater{}\textgreater{} s);\\    dest[i] = flag\_not\_result ? \textasciitilde{}r : r;\\\}}
\textbf{Execution Context:}
\begin{tightitemize}\footnotesize
  \item \textbf{Compute stage preloads (\hyperref[sec:csf-spd]{\texttt{SPD\_k}}, \texttt{r16..r23}):} Hardware pre-fills \texttt{r55..r57} with \texttt{WorkGroupID.xyz + JOB\_OFFSET} (from \hyperref[sec:csf-regs]{\texttt{COMPUTESR::JOB\_OFFSET\_X/Y/Z}} \texttt{r34..r36}) and \texttt{r60..r62} with \texttt{LocalInvocationID.xyz} or \texttt{GlobalInvocationID.xyz} (from \texttt{WG\_SIZE} \texttt{r33} \& \texttt{JOB\_SIZE} \texttt{r37..r39}, selected by \texttt{ShaderProgram::preload}) so compute shaders can derive thread/element indices without memory loads.
  \item \textbf{Vertex stage preloads (\hyperref[sec:csf-spd]{\texttt{SPD\_0}/\texttt{SPD\_1}}, \texttt{r16..r19}):} The hardware Index Fetcher reads \hyperref[sec:csf-regs]{\texttt{IDVSSR::INDEX\_BUFFER}} (\texttt{r54:r55}) and pre-fills \texttt{r59, r60} with \texttt{VertexIndex} (\texttt{+ VERTEX\_OFFSET + GLOBAL\_ATTR\_OFF}) and \texttt{r61} with \texttt{InstanceIndex} (\texttt{+ INSTANCE\_OFFSET}) so Position and Varying shaders can pull vertex attributes directly.
  \item \textbf{Inline FAU \& Special Operands (\hyperref[sec:csf-tsd-fau]{\texttt{FAU\_k}}, \texttt{r8..r15}):} Any 8-bit source with \texttt{mode=2} reads zero-latency uniform \texttt{u0..u31} from \texttt{FAU\_k} (\texttt{slots64 = bits[63:56]}); \texttt{mode=3} reads \texttt{SmallConstant} (Table~1.2) or \texttt{FauSpecial} (\texttt{pc}, \texttt{lane\_id}, \texttt{lane\_width}, \texttt{tls\_ptr}, \texttt{wls\_ptr}, \texttt{blend\_descriptor\_0..7}).
\end{tightitemize}
\par\textbf{Worked Examples:}
\par\textbf{Example 1 (Canonical 32-bit shift/rotate):}
\codebox{\mnem{ARSHIFT.i32} \dst{r0}, \src{r1}, \src{r2}\imod{.b0}}
With \texttt{r1 = 0x80000000}, \texttt{r2.b0 = 4}: yields \texttt{r0 = 0xF8000000}.
\par\textbf{Example 2 (Byte-lane selector \& modulo-32 masking):}
\codebox{\mnem{ARSHIFT.i32} \dst{r0}, \src{r1}, \src{r2}\imod{.b1}}
With \texttt{r1 = 0x80000000}, \texttt{r2.b1 = 31} (shift count masked mod 32 -\textgreater{} \texttt{31}): yields \texttt{r0 = 0xFFFFFFFF}.

\begin{center}\footnotesize
\begin{tabularx}{\linewidth}{@{}lllX@{}}
\toprule
Variant & Arch & Class & Opcode Fields \\
\midrule
\texttt{(none)} & v9--v14 & \texttt{R3} & \texttt{opcode1=SHIFT\_BITWISE\_I32} \\
\texttt{(none)} & v9--v14 & \texttt{R3} & \texttt{opcode1=SHIFT\_BITWISE\_I32} \\
\texttt{(none)} & v9--v14 & \texttt{R3} & \texttt{opcode1=LOGIC\_SHIFT\_I32} \\
\texttt{(none)} & v9--v14 & \texttt{D3} & \texttt{opcode1=SHIFT\_BITWISE\_I64} \\
\texttt{(none)} & v9--v14 & \texttt{D3} & \texttt{opcode1=LOGIC\_SHIFT\_I64} \\
\texttt{(none)} & v9--v14 & \texttt{R3} & \texttt{opcode1=SHIFT\_BITWISE\_V2I16} \\
\texttt{(none)} & v9--v14 & \texttt{R3} & \texttt{opcode1=LOGIC\_SHIFT\_V2I16} \\
\texttt{(none)} & v9--v14 & \texttt{R3} & \texttt{opcode1=SHIFT\_BITWISE\_V4I8} \\
\texttt{(none)} & v9--v14 & \texttt{R3} & \texttt{opcode1=LOGIC\_SHIFT\_V4I8} \\
\bottomrule
\end{tabularx}
\end{center}
\begin{center}\footnotesize
\begin{tabularx}{0.88\linewidth}{@{}l l X@{}}
\toprule
Field Name & Bits & Type / Fixed Value \\
\midrule
\texttt{src0} & \texttt{0:7} & \texttt{SourceEncoding} \\
\texttt{src1} & \texttt{8:15} & \texttt{SourceEncoding} \\
\texttt{src2} & \texttt{16:23} & \texttt{exact=192} \\
\texttt{bitwise\_operation} & \texttt{24:25} & \texttt{exact=or} \\
\hyperref[enum:shift_bitwise_i32_lanes_m]{\texttt{lanes1}} & \texttt{26:29} & \texttt{shift\_bitwise\_i32\_lanes\_m (def b0)} \\
\hyperref[enum:not_m]{\texttt{not\_result}} & \texttt{30} & \texttt{not\_m (def none)} \\
\textit{Reserved} & \texttt{31} & \texttt{zero} \\
\texttt{shift\_operation} & \texttt{32:34} & \texttt{exact=arshift} \\
\hyperref[enum:not_m]{\texttt{not2}} & \texttt{35} & \texttt{exact=none} \\
\hyperref[enum:widen_int_m]{\texttt{widen0}} & \texttt{36:39} & \texttt{widen\_int\_m (def none)} \\
\texttt{dest\_register\_index} & \texttt{40:45} & \texttt{uint} \\
\hyperref[enum:dest_width_m]{\texttt{dest\_width}} & \texttt{46:47} & \texttt{dest\_width\_m} \\
\texttt{opcode1} & \texttt{48:51} & \texttt{exact=SHIFT\_BITWISE\_I32} \\
\texttt{instruction\_class} & \texttt{52:56} & \texttt{exact=R3} \\
\texttt{fau\_page\_index} & \texttt{57:58} & \texttt{uint} \\
\hyperref[enum:flow_control_m]{\texttt{flow}} & \texttt{59:62} & \texttt{flow\_control\_m} \\
\textit{Reserved} & \texttt{63} & \texttt{zero} \\
\bottomrule
\end{tabularx}
\end{center}
\textbf{Fields Explained:}
\begin{tightitemize}\footnotesize
  \item \texttt{src0}: Primary source operand (\texttt{r0..r63}, \texttt{u0..u63} uniform, \texttt{k*} FAU constant, or \texttt{t0..t2} temporary; \texttt{\textasciicircum{}} marks discard/last-use). \emph{Example:} {\scriptsize\ttfamily \mnem{FADD.f32} \dst{r0}, \src{r1}, \src{r2}} (\textbf{SPIR-V:} \texttt{OpFAdd (Operand 0)}; \textbf{GLSL:} \texttt{float c = a + b;})
  \item \texttt{src1}: Secondary source operand (register, uniform, or FAU constant; in shifts, supplies the shift count byte). \emph{Example:} {\scriptsize\ttfamily \mnem{LSHIFT.i32} \dst{r0}, \src{r1}, \src{r2}\imod{.b0}} (\textbf{SPIR-V:} \texttt{OpShiftLeftLogical (Shift)}; \textbf{GLSL:} \texttt{uint y = a \textless{}\textless{} s;})
  \item \texttt{src2}: Third source operand (FMA addend, \texttt{MUX}/\texttt{CSEL} selector, \texttt{IADDC} carry-in, or bitwise mask in fused shift-logic). \emph{Example:} {\scriptsize\ttfamily \mnem{FMA.f32} \dst{r0}, \src{r1}, \src{r2}, \src{r3}} (\textbf{SPIR-V:} \texttt{GLSL.std.450 Fma (Addend)}; \textbf{GLSL:} \texttt{float d = fma(a, b, c);})
  \item \texttt{bitwise\_operation}: Selects the bitwise ALU op (\texttt{and}, \texttt{or}, \texttt{xor}, \texttt{lshift\_and}, etc.) fused with shifts. \emph{Example:} {\scriptsize\ttfamily \mnem{LSHIFT\_OR.i32} \dst{r0}, \src{r1}, \src{r2}, \src{r3}} (\textbf{SPIR-V:} \texttt{OpShiftLeftLogical + OpBitwiseOr}; \textbf{GLSL:} \texttt{uint y = (a \textless{}\textless{} s) \textbar{} b;})
  \item \hyperref[enum:shift_bitwise_i32_lanes_m]{\texttt{lanes1}}: Stage A byte-lane selector on \texttt{src1} (selects shift-count byte \texttt{b0..b3} on 32/64-bit shifts, \texttt{bXY} on \texttt{v2i16} shifts, or 4-byte swizzle on \texttt{v4*8}). \emph{Example:} {\scriptsize\ttfamily \mnem{LSHIFT.i32} \dst{r0}, \src{r1}, \src{r2}\imod{.b1}} (\textbf{SPIR-V:} \texttt{OpBitFieldUExtract + OpShiftLeftLogical}; \textbf{GLSL:} \texttt{uint y = a \textless{}\textless{} ((s \textgreater{}\textgreater{} 8) \& 0xFFu);})
  \item \hyperref[enum:not_m]{\texttt{not\_result}}: Bitwise-inverts the final output of a bitwise/shift-bitwise instruction (\textasciitilde{}result). \emph{Example:} {\scriptsize\ttfamily \mnem{AND.i32}\imod{.not\_result} \dst{r0}, \src{r1}, \src{r2}} (\textbf{SPIR-V:} \texttt{OpBitwiseAnd + OpNot}; \textbf{GLSL:} \texttt{uint y = \textasciitilde{}(a \& b);})
  \item \texttt{shift\_operation}: Selects shift direction/type (\texttt{lshift}, \texttt{rshift}, \texttt{arshift}) in fused shift-bitwise instructions. \emph{Example:} {\scriptsize\ttfamily \mnem{RSHIFT\_AND.i32} \dst{r0}, \src{r1}, \src{r2}, \src{r3}} (\textbf{SPIR-V:} \texttt{OpShiftRightLogical + OpBitwiseAnd}; \textbf{GLSL:} \texttt{uint y = (a \textgreater{}\textgreater{} s) \& m;})
  \item \hyperref[enum:not_m]{\texttt{not2}}: Stage A bitwise complement (\texttt{\textasciitilde{}src2}) applied to \texttt{src2} before the bitwise/shift-logic operation. \emph{Example:} {\scriptsize\ttfamily \mnem{AND.i32} \dst{r0}, \src{r1}, \src{r2}\imod{.not}} (\textbf{SPIR-V:} \texttt{OpNot + OpBitwiseAnd}; \textbf{GLSL:} \texttt{uint y = a \& \textasciitilde{}b;})
  \item \hyperref[enum:widen_int_m]{\texttt{widen0}}: Stage A widening extractor on \texttt{src0} (promotes FP16 \texttt{h0}/\texttt{h1} to FP32, or sign/zero-extends \texttt{b0..b3}/\texttt{h0..h1}/\texttt{w0..w1} to 32/64-bit before \texttt{abs}/\texttt{neg}). \emph{Example:} {\scriptsize\ttfamily \mnem{IADD.s32} \dst{r0}, \src{r1}\imod{.h0}, \src{r2}} (\textbf{SPIR-V:} \texttt{OpSConvert / OpFConvert}; \textbf{GLSL:} \texttt{int c = int(int16\_t(a)) + b;})
  \item \texttt{dest\_register\_index}: 6-bit destination GPR index (\texttt{r0..r63}) receiving the Stage C writeback result (even-aligned pair \texttt{[rN:rN+1]} for 64-bit ops). \emph{Example:} {\scriptsize\ttfamily \mnem{FMA.f32} \dst{r0}, \src{r1}, \src{r2}, \src{r3}} (\textbf{SPIR-V:} \texttt{SSA Result ID}; \textbf{GLSL:} \texttt{float r0 = fma(r1, r2, r3);})
  \item \hyperref[enum:dest_width_m]{\texttt{dest\_width}}: Stage C destination lane write-mask (\texttt{none} = full 32-bit, \texttt{h0} = low 16-bit, \texttt{h1} = high 16-bit without bit-shifting). \emph{Example:} {\scriptsize\ttfamily \mnem{FADD.v2f16} \dst{r0}\imod{.h1}, \src{r1}, \src{r2}} (\textbf{SPIR-V:} \texttt{OpVectorInsertDynamic / WriteMask}; \textbf{GLSL:} \texttt{dst.y = a.y + b.y;})
  \item \texttt{opcode1}: Primary sub-opcode within the instruction class uniquely identifying the operation family. \emph{Example:} {\scriptsize\ttfamily \mnem{IADD.s32} \dst{r0}, \src{r1}, \src{r2}} (\textbf{SPIR-V:} \texttt{OpIAdd (operation selector)}; \textbf{GLSL:} \texttt{// Primary opcode bits})
  \item \texttt{instruction\_class}: 4-bit major instruction class selector (\texttt{bits 48:51}) partitioning the 64-bit Valhall opcode space (e.g., \texttt{FMA}, \texttt{CVT}, \texttt{SFU}, \texttt{LS}, \texttt{TEX}, \texttt{BRANCH}). \emph{Example:} {\scriptsize\ttfamily \mnem{FADD.f32} \dst{r0}, \src{r1}, \src{r2}} (\textbf{SPIR-V:} \texttt{N/A (encoding class)}; \textbf{GLSL:} \texttt{// Hardware execution pipe / decode class})
  \item \texttt{fau\_page\_index}: Selects the FAU (Fast Access Uniform) page (\texttt{0..3}) when a source operand references hardware small constants, thread IDs, or lane constants. \emph{Example:} {\scriptsize\ttfamily \mnem{FADD.f32} \dst{r0}, \src{r1}, \src{k1f}} (\textbf{SPIR-V:} \texttt{OpConstant / BuiltIn}; \textbf{GLSL:} \texttt{float y = x + 1.0;})
  \item \hyperref[enum:flow_control_m]{\texttt{flow}}: 4-bit embedded flow-control \& scoreboard wait modifier (\texttt{bits 59:62}: \texttt{none}, \texttt{wait0..wait0126}, \texttt{wait}, \texttt{reconverge}, \texttt{discard}, \texttt{end}). \emph{Example:} {\scriptsize\ttfamily \mnem{FADD.f32}\imod{.wait0.reconverge} \dst{r0}, \src{r1}, \src{r2}} (\textbf{SPIR-V:} \texttt{OpControlBarrier / OpReturn}; \textbf{GLSL:} \texttt{// Scoreboard wait + warp reconvergence})
\end{tightitemize}
\clearpage

\subsection*{ARSHIFT\_AND}
\addcontentsline{toc}{subsection}{ARSHIFT\_AND}
For each SIMD element $i \in \{0, \dots, N-1\}$ ($N=1$ for \texttt{i32}/\texttt{i64}, $N=2$ for \texttt{v2i16}, $N=4$ for \texttt{v4i8} on v9--v10, lane width $W$), computes arithmetic right shift \texttt{s = (int\_W(val) \textgreater{}\textgreater{} (amt \& (W - 1)))} where \texttt{val = src0[i]} and the 8-bit shift count \texttt{amt} is extracted from \texttt{src1} via \texttt{lanes1}, then combines with \texttt{src2[i]} via bitwise AND \texttt{res = s \& (not2 ? \textasciitilde{}src2[i] : src2[i])} and writes \texttt{not\_result ? \textasciitilde{}res : res} to \texttt{dest[i]}.
\par\textbf{Notes:}
\begin{tightitemize}
  \item Shift amount is masked modulo lane width (\texttt{wrapping\_shl}/\texttt{wrapping\_shr}/\texttt{rotate\_*}, \texttt{kraid/ops.rs:3984-4009}).
  \item On \texttt{v2i16}, \texttt{src1} is \texttt{V2I8} (\texttt{ops.rs:3953}) with byte lanes selected by \texttt{lanes1} (default \texttt{b02}: byte 0 -\textgreater{} lane 0, byte 2 -\textgreater{} lane 1).
  \item \textbf{Source references:} \texttt{kraid/ops.rs:3878-4025}, \texttt{encode\_v9.rs:2940-3093}, \texttt{hw\_tests.rs:2148-2198}.
\end{tightitemize}
\codebox{\mnem{ARSHIFT\_AND}\imod{\{.flow\}}\imod{\{.not\_result\}} \dst{dest}\imod{\{.dest\_width\}}, \src{src0}\imod{\{.widen0\}}, \src{src1}\imod{\{.lanes1\}}, \src{src2}\imod{\{.not2\}}\\\mnem{ARSHIFT\_AND}\imod{\{.flow\}}\imod{\{.not\_result\}} \dst{dest}, \src{src0}\imod{\{.widen0\}}, \src{src1}\imod{\{.lanes1\}}, \src{src2}\imod{\{.not2\}}\\\mnem{ARSHIFT\_AND}\imod{\{.flow\}}\imod{\{.not\_result\}} \dst{dest}\imod{\{.dest\_width\}}, \src{src0}\imod{\{.swz0\}}, \src{src1}\imod{\{.lanes1\}}, \src{src2}\imod{\{.not2\}}\\\mnem{ARSHIFT\_AND}\imod{\{.flow\}}\imod{\{.not\_result\}} \dst{dest}\imod{\{.dest\_width\}}, \src{src0}\imod{\{.lanes0\}}, \src{src1}\imod{\{.lanes1\}}, \src{src2}\imod{\{.not2\}}}
\textbf{Operands:}
\begin{tightitemize}\footnotesize
  \item \texttt{\dst{dest}} (\texttt{dest\_register\_index}, \texttt{dest\_width}): Destination GPR receiving the arithmetic right shift combined via bitwise AND (inverted if \texttt{.not\_result}).
  \item \texttt{\src{src0}}: Input integer value to shift/rotate.
  \item \texttt{\src{src1}}: 8-bit shift/rotate count per lane (selected via \texttt{lanes1}, masked modulo lane width \texttt{W}).
  \item \texttt{\src{src2}}: Second operand combined via bitwise AND after shifting (inverted if \texttt{.not2}).
\end{tightitemize}
\par\textbf{Closest SPIR-V Op:} \texttt{OpShiftRightArithmetic + OpBitwiseAnd}
\par\textbf{Pseudo-GLSL (Operational Semantics):}
\codebox{for (int i = 0; i \textless{} N; ++i) \{\\    uint s = extractByte(src1, lanes1, i) \& (W - 1);\\    auto r = (int(src0[i]) \textgreater{}\textgreater{} s) \& (flag\_not2 ? \textasciitilde{}src2[i] : src2[i]);\\    dest[i] = flag\_not\_result ? \textasciitilde{}r : r;\\\}}
\textbf{Execution Context:}
\begin{tightitemize}\footnotesize
  \item \textbf{Compute stage preloads (\hyperref[sec:csf-spd]{\texttt{SPD\_k}}, \texttt{r16..r23}):} Hardware pre-fills \texttt{r55..r57} with \texttt{WorkGroupID.xyz + JOB\_OFFSET} (from \hyperref[sec:csf-regs]{\texttt{COMPUTESR::JOB\_OFFSET\_X/Y/Z}} \texttt{r34..r36}) and \texttt{r60..r62} with \texttt{LocalInvocationID.xyz} or \texttt{GlobalInvocationID.xyz} (from \texttt{WG\_SIZE} \texttt{r33} \& \texttt{JOB\_SIZE} \texttt{r37..r39}, selected by \texttt{ShaderProgram::preload}) so compute shaders can derive thread/element indices without memory loads.
  \item \textbf{Vertex stage preloads (\hyperref[sec:csf-spd]{\texttt{SPD\_0}/\texttt{SPD\_1}}, \texttt{r16..r19}):} The hardware Index Fetcher reads \hyperref[sec:csf-regs]{\texttt{IDVSSR::INDEX\_BUFFER}} (\texttt{r54:r55}) and pre-fills \texttt{r59, r60} with \texttt{VertexIndex} (\texttt{+ VERTEX\_OFFSET + GLOBAL\_ATTR\_OFF}) and \texttt{r61} with \texttt{InstanceIndex} (\texttt{+ INSTANCE\_OFFSET}) so Position and Varying shaders can pull vertex attributes directly.
  \item \textbf{Inline FAU \& Special Operands (\hyperref[sec:csf-tsd-fau]{\texttt{FAU\_k}}, \texttt{r8..r15}):} Any 8-bit source with \texttt{mode=2} reads zero-latency uniform \texttt{u0..u31} from \texttt{FAU\_k} (\texttt{slots64 = bits[63:56]}); \texttt{mode=3} reads \texttt{SmallConstant} (Table~1.2) or \texttt{FauSpecial} (\texttt{pc}, \texttt{lane\_id}, \texttt{lane\_width}, \texttt{tls\_ptr}, \texttt{wls\_ptr}, \texttt{blend\_descriptor\_0..7}).
\end{tightitemize}
\par\textbf{Worked Examples:}
\par\textbf{Example 1 (Standard fused shift + bitwise op):}
\codebox{\mnem{ARSHIFT\_AND.i32} \dst{r0}, \src{r1}, \src{r2}\imod{.b0}, \src{r3}}
With \texttt{r1 = 0x80000000}, \texttt{r2.b0 = 4} (masked mod 32): shift yields \texttt{0xF8000000}, then \texttt{\& 0x0FFFFFFF} yields \texttt{r0 = 0x8000000}.
\par\textbf{Example 2 (Modulo-32 shift count \& .not\_result):}
\codebox{\mnem{ARSHIFT\_AND.i32}\imod{.not\_result} \dst{r0}, \src{r1}, \src{r2}\imod{.b0}, \src{r3}}
With \texttt{r1 = 0x80000000}, \texttt{r2.b0 = 31} (effective shift \texttt{31}): shift yields \texttt{0xFFFFFFFF}, \texttt{\& 0x0FFFFFFF} yields \texttt{0x0FFFFFFF}, and \texttt{.not\_result} inverts -\textgreater{} \texttt{r0 = 0xF0000000}.

\begin{center}\footnotesize
\begin{tabularx}{\linewidth}{@{}lllX@{}}
\toprule
Variant & Arch & Class & Opcode Fields \\
\midrule
\texttt{(none)} & v9--v14 & \texttt{R3} & \texttt{opcode1=SHIFT\_BITWISE\_I32} \\
\texttt{(none)} & v9--v14 & \texttt{R3} & \texttt{opcode1=SHIFT\_BITWISE\_I32} \\
\texttt{(none)} & v9--v14 & \texttt{R3} & \texttt{opcode1=LOGIC\_SHIFT\_I32} \\
\texttt{(none)} & v9--v14 & \texttt{D3} & \texttt{opcode1=SHIFT\_BITWISE\_I64} \\
\texttt{(none)} & v9--v14 & \texttt{D3} & \texttt{opcode1=LOGIC\_SHIFT\_I64} \\
\texttt{(none)} & v9--v14 & \texttt{R3} & \texttt{opcode1=SHIFT\_BITWISE\_V2I16} \\
\texttt{(none)} & v9--v14 & \texttt{R3} & \texttt{opcode1=LOGIC\_SHIFT\_V2I16} \\
\texttt{(none)} & v9--v14 & \texttt{R3} & \texttt{opcode1=SHIFT\_BITWISE\_V4I8} \\
\texttt{(none)} & v9--v14 & \texttt{R3} & \texttt{opcode1=LOGIC\_SHIFT\_V4I8} \\
\bottomrule
\end{tabularx}
\end{center}
\begin{center}\footnotesize
\begin{tabularx}{0.88\linewidth}{@{}l l X@{}}
\toprule
Field Name & Bits & Type / Fixed Value \\
\midrule
\texttt{src0} & \texttt{0:7} & \texttt{SourceEncoding} \\
\texttt{src1} & \texttt{8:15} & \texttt{SourceEncoding} \\
\texttt{src2} & \texttt{16:23} & \texttt{SourceEncoding} \\
\texttt{bitwise\_operation} & \texttt{24:25} & \texttt{exact=and} \\
\hyperref[enum:shift_bitwise_i32_lanes_m]{\texttt{lanes1}} & \texttt{26:29} & \texttt{shift\_bitwise\_i32\_lanes\_m (def b0)} \\
\hyperref[enum:not_m]{\texttt{not\_result}} & \texttt{30} & \texttt{not\_m (def none)} \\
\textit{Reserved} & \texttt{31} & \texttt{zero} \\
\texttt{shift\_operation} & \texttt{32:34} & \texttt{exact=arshift} \\
\hyperref[enum:not_m]{\texttt{not2}} & \texttt{35} & \texttt{not\_m (def none)} \\
\hyperref[enum:widen_int_m]{\texttt{widen0}} & \texttt{36:39} & \texttt{widen\_int\_m (def none)} \\
\texttt{dest\_register\_index} & \texttt{40:45} & \texttt{uint} \\
\hyperref[enum:dest_width_m]{\texttt{dest\_width}} & \texttt{46:47} & \texttt{dest\_width\_m} \\
\texttt{opcode1} & \texttt{48:51} & \texttt{exact=SHIFT\_BITWISE\_I32} \\
\texttt{instruction\_class} & \texttt{52:56} & \texttt{exact=R3} \\
\texttt{fau\_page\_index} & \texttt{57:58} & \texttt{uint} \\
\hyperref[enum:flow_control_m]{\texttt{flow}} & \texttt{59:62} & \texttt{flow\_control\_m} \\
\textit{Reserved} & \texttt{63} & \texttt{zero} \\
\bottomrule
\end{tabularx}
\end{center}
\textbf{Fields Explained:}
\begin{tightitemize}\footnotesize
  \item \texttt{src0}: Primary source operand (\texttt{r0..r63}, \texttt{u0..u63} uniform, \texttt{k*} FAU constant, or \texttt{t0..t2} temporary; \texttt{\textasciicircum{}} marks discard/last-use). \emph{Example:} {\scriptsize\ttfamily \mnem{FADD.f32} \dst{r0}, \src{r1}, \src{r2}} (\textbf{SPIR-V:} \texttt{OpFAdd (Operand 0)}; \textbf{GLSL:} \texttt{float c = a + b;})
  \item \texttt{src1}: Secondary source operand (register, uniform, or FAU constant; in shifts, supplies the shift count byte). \emph{Example:} {\scriptsize\ttfamily \mnem{LSHIFT.i32} \dst{r0}, \src{r1}, \src{r2}\imod{.b0}} (\textbf{SPIR-V:} \texttt{OpShiftLeftLogical (Shift)}; \textbf{GLSL:} \texttt{uint y = a \textless{}\textless{} s;})
  \item \texttt{src2}: Third source operand (FMA addend, \texttt{MUX}/\texttt{CSEL} selector, \texttt{IADDC} carry-in, or bitwise mask in fused shift-logic). \emph{Example:} {\scriptsize\ttfamily \mnem{FMA.f32} \dst{r0}, \src{r1}, \src{r2}, \src{r3}} (\textbf{SPIR-V:} \texttt{GLSL.std.450 Fma (Addend)}; \textbf{GLSL:} \texttt{float d = fma(a, b, c);})
  \item \texttt{bitwise\_operation}: Selects the bitwise ALU op (\texttt{and}, \texttt{or}, \texttt{xor}, \texttt{lshift\_and}, etc.) fused with shifts. \emph{Example:} {\scriptsize\ttfamily \mnem{LSHIFT\_OR.i32} \dst{r0}, \src{r1}, \src{r2}, \src{r3}} (\textbf{SPIR-V:} \texttt{OpShiftLeftLogical + OpBitwiseOr}; \textbf{GLSL:} \texttt{uint y = (a \textless{}\textless{} s) \textbar{} b;})
  \item \hyperref[enum:shift_bitwise_i32_lanes_m]{\texttt{lanes1}}: Stage A byte-lane selector on \texttt{src1} (selects shift-count byte \texttt{b0..b3} on 32/64-bit shifts, \texttt{bXY} on \texttt{v2i16} shifts, or 4-byte swizzle on \texttt{v4*8}). \emph{Example:} {\scriptsize\ttfamily \mnem{LSHIFT.i32} \dst{r0}, \src{r1}, \src{r2}\imod{.b1}} (\textbf{SPIR-V:} \texttt{OpBitFieldUExtract + OpShiftLeftLogical}; \textbf{GLSL:} \texttt{uint y = a \textless{}\textless{} ((s \textgreater{}\textgreater{} 8) \& 0xFFu);})
  \item \hyperref[enum:not_m]{\texttt{not\_result}}: Bitwise-inverts the final output of a bitwise/shift-bitwise instruction (\textasciitilde{}result). \emph{Example:} {\scriptsize\ttfamily \mnem{AND.i32}\imod{.not\_result} \dst{r0}, \src{r1}, \src{r2}} (\textbf{SPIR-V:} \texttt{OpBitwiseAnd + OpNot}; \textbf{GLSL:} \texttt{uint y = \textasciitilde{}(a \& b);})
  \item \texttt{shift\_operation}: Selects shift direction/type (\texttt{lshift}, \texttt{rshift}, \texttt{arshift}) in fused shift-bitwise instructions. \emph{Example:} {\scriptsize\ttfamily \mnem{RSHIFT\_AND.i32} \dst{r0}, \src{r1}, \src{r2}, \src{r3}} (\textbf{SPIR-V:} \texttt{OpShiftRightLogical + OpBitwiseAnd}; \textbf{GLSL:} \texttt{uint y = (a \textgreater{}\textgreater{} s) \& m;})
  \item \hyperref[enum:not_m]{\texttt{not2}}: Stage A bitwise complement (\texttt{\textasciitilde{}src2}) applied to \texttt{src2} before the bitwise/shift-logic operation. \emph{Example:} {\scriptsize\ttfamily \mnem{AND.i32} \dst{r0}, \src{r1}, \src{r2}\imod{.not}} (\textbf{SPIR-V:} \texttt{OpNot + OpBitwiseAnd}; \textbf{GLSL:} \texttt{uint y = a \& \textasciitilde{}b;})
  \item \hyperref[enum:widen_int_m]{\texttt{widen0}}: Stage A widening extractor on \texttt{src0} (promotes FP16 \texttt{h0}/\texttt{h1} to FP32, or sign/zero-extends \texttt{b0..b3}/\texttt{h0..h1}/\texttt{w0..w1} to 32/64-bit before \texttt{abs}/\texttt{neg}). \emph{Example:} {\scriptsize\ttfamily \mnem{IADD.s32} \dst{r0}, \src{r1}\imod{.h0}, \src{r2}} (\textbf{SPIR-V:} \texttt{OpSConvert / OpFConvert}; \textbf{GLSL:} \texttt{int c = int(int16\_t(a)) + b;})
  \item \texttt{dest\_register\_index}: 6-bit destination GPR index (\texttt{r0..r63}) receiving the Stage C writeback result (even-aligned pair \texttt{[rN:rN+1]} for 64-bit ops). \emph{Example:} {\scriptsize\ttfamily \mnem{FMA.f32} \dst{r0}, \src{r1}, \src{r2}, \src{r3}} (\textbf{SPIR-V:} \texttt{SSA Result ID}; \textbf{GLSL:} \texttt{float r0 = fma(r1, r2, r3);})
  \item \hyperref[enum:dest_width_m]{\texttt{dest\_width}}: Stage C destination lane write-mask (\texttt{none} = full 32-bit, \texttt{h0} = low 16-bit, \texttt{h1} = high 16-bit without bit-shifting). \emph{Example:} {\scriptsize\ttfamily \mnem{FADD.v2f16} \dst{r0}\imod{.h1}, \src{r1}, \src{r2}} (\textbf{SPIR-V:} \texttt{OpVectorInsertDynamic / WriteMask}; \textbf{GLSL:} \texttt{dst.y = a.y + b.y;})
  \item \texttt{opcode1}: Primary sub-opcode within the instruction class uniquely identifying the operation family. \emph{Example:} {\scriptsize\ttfamily \mnem{IADD.s32} \dst{r0}, \src{r1}, \src{r2}} (\textbf{SPIR-V:} \texttt{OpIAdd (operation selector)}; \textbf{GLSL:} \texttt{// Primary opcode bits})
  \item \texttt{instruction\_class}: 4-bit major instruction class selector (\texttt{bits 48:51}) partitioning the 64-bit Valhall opcode space (e.g., \texttt{FMA}, \texttt{CVT}, \texttt{SFU}, \texttt{LS}, \texttt{TEX}, \texttt{BRANCH}). \emph{Example:} {\scriptsize\ttfamily \mnem{FADD.f32} \dst{r0}, \src{r1}, \src{r2}} (\textbf{SPIR-V:} \texttt{N/A (encoding class)}; \textbf{GLSL:} \texttt{// Hardware execution pipe / decode class})
  \item \texttt{fau\_page\_index}: Selects the FAU (Fast Access Uniform) page (\texttt{0..3}) when a source operand references hardware small constants, thread IDs, or lane constants. \emph{Example:} {\scriptsize\ttfamily \mnem{FADD.f32} \dst{r0}, \src{r1}, \src{k1f}} (\textbf{SPIR-V:} \texttt{OpConstant / BuiltIn}; \textbf{GLSL:} \texttt{float y = x + 1.0;})
  \item \hyperref[enum:flow_control_m]{\texttt{flow}}: 4-bit embedded flow-control \& scoreboard wait modifier (\texttt{bits 59:62}: \texttt{none}, \texttt{wait0..wait0126}, \texttt{wait}, \texttt{reconverge}, \texttt{discard}, \texttt{end}). \emph{Example:} {\scriptsize\ttfamily \mnem{FADD.f32}\imod{.wait0.reconverge} \dst{r0}, \src{r1}, \src{r2}} (\textbf{SPIR-V:} \texttt{OpControlBarrier / OpReturn}; \textbf{GLSL:} \texttt{// Scoreboard wait + warp reconvergence})
\end{tightitemize}
\clearpage

\subsection*{ARSHIFT\_OR}
\addcontentsline{toc}{subsection}{ARSHIFT\_OR}
For each SIMD element $i \in \{0, \dots, N-1\}$ ($N=1$ for \texttt{i32}/\texttt{i64}, $N=2$ for \texttt{v2i16}, $N=4$ for \texttt{v4i8} on v9--v10, lane width $W$), computes arithmetic right shift \texttt{s = (int\_W(val) \textgreater{}\textgreater{} (amt \& (W - 1)))} where \texttt{val = src0[i]} and the 8-bit shift count \texttt{amt} is extracted from \texttt{src1} via \texttt{lanes1}, then combines with \texttt{src2[i]} via bitwise OR \texttt{res = s \textbar{} (not2 ? \textasciitilde{}src2[i] : src2[i])} and writes \texttt{not\_result ? \textasciitilde{}res : res} to \texttt{dest[i]}.
\par\textbf{Notes:}
\begin{tightitemize}
  \item Shift amount is masked modulo lane width (\texttt{wrapping\_shl}/\texttt{wrapping\_shr}/\texttt{rotate\_*}, \texttt{kraid/ops.rs:3984-4009}).
  \item On \texttt{v2i16}, \texttt{src1} is \texttt{V2I8} (\texttt{ops.rs:3953}) with byte lanes selected by \texttt{lanes1} (default \texttt{b02}: byte 0 -\textgreater{} lane 0, byte 2 -\textgreater{} lane 1).
  \item \textbf{Source references:} \texttt{kraid/ops.rs:3878-4025}, \texttt{encode\_v9.rs:2940-3093}, \texttt{hw\_tests.rs:2148-2198}.
\end{tightitemize}
\codebox{\mnem{ARSHIFT\_OR}\imod{\{.flow\}}\imod{\{.not\_result\}} \dst{dest}\imod{\{.dest\_width\}}, \src{src0}\imod{\{.widen0\}}, \src{src1}\imod{\{.lanes1\}}, \src{src2}\imod{\{.not2\}}\\\mnem{ARSHIFT\_OR}\imod{\{.flow\}}\imod{\{.not\_result\}} \dst{dest}, \src{src0}\imod{\{.widen0\}}, \src{src1}\imod{\{.lanes1\}}, \src{src2}\imod{\{.not2\}}\\\mnem{ARSHIFT\_OR}\imod{\{.flow\}}\imod{\{.not\_result\}} \dst{dest}\imod{\{.dest\_width\}}, \src{src0}\imod{\{.swz0\}}, \src{src1}\imod{\{.lanes1\}}, \src{src2}\imod{\{.not2\}}\\\mnem{ARSHIFT\_OR}\imod{\{.flow\}}\imod{\{.not\_result\}} \dst{dest}\imod{\{.dest\_width\}}, \src{src0}\imod{\{.lanes0\}}, \src{src1}\imod{\{.lanes1\}}, \src{src2}\imod{\{.not2\}}}
\textbf{Operands:}
\begin{tightitemize}\footnotesize
  \item \texttt{\dst{dest}} (\texttt{dest\_register\_index}, \texttt{dest\_width}): Destination GPR receiving the arithmetic right shift combined via bitwise OR (inverted if \texttt{.not\_result}).
  \item \texttt{\src{src0}}: Input integer value to shift/rotate.
  \item \texttt{\src{src1}}: 8-bit shift/rotate count per lane (selected via \texttt{lanes1}, masked modulo lane width \texttt{W}).
  \item \texttt{\src{src2}}: Second operand combined via bitwise OR after shifting (inverted if \texttt{.not2}).
\end{tightitemize}
\par\textbf{Closest SPIR-V Op:} \texttt{OpShiftRightArithmetic + OpBitwiseOr}
\par\textbf{Pseudo-GLSL (Operational Semantics):}
\codebox{for (int i = 0; i \textless{} N; ++i) \{\\    uint s = extractByte(src1, lanes1, i) \& (W - 1);\\    auto r = (int(src0[i]) \textgreater{}\textgreater{} s) \textbar{} (flag\_not2 ? \textasciitilde{}src2[i] : src2[i]);\\    dest[i] = flag\_not\_result ? \textasciitilde{}r : r;\\\}}
\textbf{Execution Context:}
\begin{tightitemize}\footnotesize
  \item \textbf{Compute stage preloads (\hyperref[sec:csf-spd]{\texttt{SPD\_k}}, \texttt{r16..r23}):} Hardware pre-fills \texttt{r55..r57} with \texttt{WorkGroupID.xyz + JOB\_OFFSET} (from \hyperref[sec:csf-regs]{\texttt{COMPUTESR::JOB\_OFFSET\_X/Y/Z}} \texttt{r34..r36}) and \texttt{r60..r62} with \texttt{LocalInvocationID.xyz} or \texttt{GlobalInvocationID.xyz} (from \texttt{WG\_SIZE} \texttt{r33} \& \texttt{JOB\_SIZE} \texttt{r37..r39}, selected by \texttt{ShaderProgram::preload}) so compute shaders can derive thread/element indices without memory loads.
  \item \textbf{Vertex stage preloads (\hyperref[sec:csf-spd]{\texttt{SPD\_0}/\texttt{SPD\_1}}, \texttt{r16..r19}):} The hardware Index Fetcher reads \hyperref[sec:csf-regs]{\texttt{IDVSSR::INDEX\_BUFFER}} (\texttt{r54:r55}) and pre-fills \texttt{r59, r60} with \texttt{VertexIndex} (\texttt{+ VERTEX\_OFFSET + GLOBAL\_ATTR\_OFF}) and \texttt{r61} with \texttt{InstanceIndex} (\texttt{+ INSTANCE\_OFFSET}) so Position and Varying shaders can pull vertex attributes directly.
  \item \textbf{Inline FAU \& Special Operands (\hyperref[sec:csf-tsd-fau]{\texttt{FAU\_k}}, \texttt{r8..r15}):} Any 8-bit source with \texttt{mode=2} reads zero-latency uniform \texttt{u0..u31} from \texttt{FAU\_k} (\texttt{slots64 = bits[63:56]}); \texttt{mode=3} reads \texttt{SmallConstant} (Table~1.2) or \texttt{FauSpecial} (\texttt{pc}, \texttt{lane\_id}, \texttt{lane\_width}, \texttt{tls\_ptr}, \texttt{wls\_ptr}, \texttt{blend\_descriptor\_0..7}).
\end{tightitemize}
\par\textbf{Worked Examples:}
\par\textbf{Example 1 (Standard fused shift + bitwise op):}
\codebox{\mnem{ARSHIFT\_OR.i32} \dst{r0}, \src{r1}, \src{r2}\imod{.b0}, \src{r3}}
With \texttt{r1 = 0x80000000}, \texttt{r2.b0 = 4} (masked mod 32): shift yields \texttt{0xF8000000}, then \texttt{\textbar{} 0x0000000F} yields \texttt{r0 = 0xF800000F}.
\par\textbf{Example 2 (Modulo-32 shift count \& .not\_result):}
\codebox{\mnem{ARSHIFT\_OR.i32}\imod{.not\_result} \dst{r0}, \src{r1}, \src{r2}\imod{.b0}, \src{r3}}
With \texttt{r1 = 0x80000000}, \texttt{r2.b0 = 31} (effective shift \texttt{31}): shift yields \texttt{0xFFFFFFFF}, \texttt{\textbar{} 0x0000000F} yields \texttt{0xFFFFFFFF}, and \texttt{.not\_result} inverts -\textgreater{} \texttt{r0 = 0x00000000}.
\par\textbf{Mined SPIR-V $\to$ Valhall ISA Example (Signed arithmetic right shift of bitfield value --- Vertex shader, Mali-G710 (v10)):}
\codebox{\textbf{; SPIR-V (input)}\\\%shr = OpShiftRightArithmetic \%int \%val \%int\_21\\[3pt]\textbf{; Valhall ISA (malioc output)}\\ARSHIFT.i32 r4, r0, k17161514.b1}
Standalone \texttt{ARSHIFT.i32} is encoded as \texttt{ARSHIFT\_OR} with \texttt{src2 = k0} (0). Here \texttt{k17161514.b1} supplies the inline shift count \texttt{0x15} (21) directly from the FAU small-constant table.
\begin{center}\footnotesize
\begin{tabularx}{\linewidth}{@{}lllX@{}}
\toprule
Variant & Arch & Class & Opcode Fields \\
\midrule
\texttt{(none)} & v9--v14 & \texttt{R3} & \texttt{opcode1=SHIFT\_BITWISE\_I32} \\
\texttt{(none)} & v9--v14 & \texttt{R3} & \texttt{opcode1=SHIFT\_BITWISE\_I32} \\
\texttt{(none)} & v9--v14 & \texttt{R3} & \texttt{opcode1=LOGIC\_SHIFT\_I32} \\
\texttt{(none)} & v9--v14 & \texttt{D3} & \texttt{opcode1=SHIFT\_BITWISE\_I64} \\
\texttt{(none)} & v9--v14 & \texttt{D3} & \texttt{opcode1=LOGIC\_SHIFT\_I64} \\
\texttt{(none)} & v9--v14 & \texttt{R3} & \texttt{opcode1=SHIFT\_BITWISE\_V2I16} \\
\texttt{(none)} & v9--v14 & \texttt{R3} & \texttt{opcode1=LOGIC\_SHIFT\_V2I16} \\
\texttt{(none)} & v9--v14 & \texttt{R3} & \texttt{opcode1=SHIFT\_BITWISE\_V4I8} \\
\texttt{(none)} & v9--v14 & \texttt{R3} & \texttt{opcode1=LOGIC\_SHIFT\_V4I8} \\
\bottomrule
\end{tabularx}
\end{center}
\begin{center}\footnotesize
\begin{tabularx}{0.88\linewidth}{@{}l l X@{}}
\toprule
Field Name & Bits & Type / Fixed Value \\
\midrule
\texttt{src0} & \texttt{0:7} & \texttt{SourceEncoding} \\
\texttt{src1} & \texttt{8:15} & \texttt{SourceEncoding} \\
\texttt{src2} & \texttt{16:23} & \texttt{SourceEncoding} \\
\texttt{bitwise\_operation} & \texttt{24:25} & \texttt{exact=or} \\
\hyperref[enum:shift_bitwise_i32_lanes_m]{\texttt{lanes1}} & \texttt{26:29} & \texttt{shift\_bitwise\_i32\_lanes\_m (def b0)} \\
\hyperref[enum:not_m]{\texttt{not\_result}} & \texttt{30} & \texttt{not\_m (def none)} \\
\textit{Reserved} & \texttt{31} & \texttt{zero} \\
\texttt{shift\_operation} & \texttt{32:34} & \texttt{exact=arshift} \\
\hyperref[enum:not_m]{\texttt{not2}} & \texttt{35} & \texttt{not\_m (def none)} \\
\hyperref[enum:widen_int_m]{\texttt{widen0}} & \texttt{36:39} & \texttt{widen\_int\_m (def none)} \\
\texttt{dest\_register\_index} & \texttt{40:45} & \texttt{uint} \\
\hyperref[enum:dest_width_m]{\texttt{dest\_width}} & \texttt{46:47} & \texttt{dest\_width\_m} \\
\texttt{opcode1} & \texttt{48:51} & \texttt{exact=SHIFT\_BITWISE\_I32} \\
\texttt{instruction\_class} & \texttt{52:56} & \texttt{exact=R3} \\
\texttt{fau\_page\_index} & \texttt{57:58} & \texttt{uint} \\
\hyperref[enum:flow_control_m]{\texttt{flow}} & \texttt{59:62} & \texttt{flow\_control\_m} \\
\textit{Reserved} & \texttt{63} & \texttt{zero} \\
\bottomrule
\end{tabularx}
\end{center}
\textbf{Fields Explained:}
\begin{tightitemize}\footnotesize
  \item \texttt{src0}: Primary source operand (\texttt{r0..r63}, \texttt{u0..u63} uniform, \texttt{k*} FAU constant, or \texttt{t0..t2} temporary; \texttt{\textasciicircum{}} marks discard/last-use). \emph{Example:} {\scriptsize\ttfamily \mnem{FADD.f32} \dst{r0}, \src{r1}, \src{r2}} (\textbf{SPIR-V:} \texttt{OpFAdd (Operand 0)}; \textbf{GLSL:} \texttt{float c = a + b;})
  \item \texttt{src1}: Secondary source operand (register, uniform, or FAU constant; in shifts, supplies the shift count byte). \emph{Example:} {\scriptsize\ttfamily \mnem{LSHIFT.i32} \dst{r0}, \src{r1}, \src{r2}\imod{.b0}} (\textbf{SPIR-V:} \texttt{OpShiftLeftLogical (Shift)}; \textbf{GLSL:} \texttt{uint y = a \textless{}\textless{} s;})
  \item \texttt{src2}: Third source operand (FMA addend, \texttt{MUX}/\texttt{CSEL} selector, \texttt{IADDC} carry-in, or bitwise mask in fused shift-logic). \emph{Example:} {\scriptsize\ttfamily \mnem{FMA.f32} \dst{r0}, \src{r1}, \src{r2}, \src{r3}} (\textbf{SPIR-V:} \texttt{GLSL.std.450 Fma (Addend)}; \textbf{GLSL:} \texttt{float d = fma(a, b, c);})
  \item \texttt{bitwise\_operation}: Selects the bitwise ALU op (\texttt{and}, \texttt{or}, \texttt{xor}, \texttt{lshift\_and}, etc.) fused with shifts. \emph{Example:} {\scriptsize\ttfamily \mnem{LSHIFT\_OR.i32} \dst{r0}, \src{r1}, \src{r2}, \src{r3}} (\textbf{SPIR-V:} \texttt{OpShiftLeftLogical + OpBitwiseOr}; \textbf{GLSL:} \texttt{uint y = (a \textless{}\textless{} s) \textbar{} b;})
  \item \hyperref[enum:shift_bitwise_i32_lanes_m]{\texttt{lanes1}}: Stage A byte-lane selector on \texttt{src1} (selects shift-count byte \texttt{b0..b3} on 32/64-bit shifts, \texttt{bXY} on \texttt{v2i16} shifts, or 4-byte swizzle on \texttt{v4*8}). \emph{Example:} {\scriptsize\ttfamily \mnem{LSHIFT.i32} \dst{r0}, \src{r1}, \src{r2}\imod{.b1}} (\textbf{SPIR-V:} \texttt{OpBitFieldUExtract + OpShiftLeftLogical}; \textbf{GLSL:} \texttt{uint y = a \textless{}\textless{} ((s \textgreater{}\textgreater{} 8) \& 0xFFu);})
  \item \hyperref[enum:not_m]{\texttt{not\_result}}: Bitwise-inverts the final output of a bitwise/shift-bitwise instruction (\textasciitilde{}result). \emph{Example:} {\scriptsize\ttfamily \mnem{AND.i32}\imod{.not\_result} \dst{r0}, \src{r1}, \src{r2}} (\textbf{SPIR-V:} \texttt{OpBitwiseAnd + OpNot}; \textbf{GLSL:} \texttt{uint y = \textasciitilde{}(a \& b);})
  \item \texttt{shift\_operation}: Selects shift direction/type (\texttt{lshift}, \texttt{rshift}, \texttt{arshift}) in fused shift-bitwise instructions. \emph{Example:} {\scriptsize\ttfamily \mnem{RSHIFT\_AND.i32} \dst{r0}, \src{r1}, \src{r2}, \src{r3}} (\textbf{SPIR-V:} \texttt{OpShiftRightLogical + OpBitwiseAnd}; \textbf{GLSL:} \texttt{uint y = (a \textgreater{}\textgreater{} s) \& m;})
  \item \hyperref[enum:not_m]{\texttt{not2}}: Stage A bitwise complement (\texttt{\textasciitilde{}src2}) applied to \texttt{src2} before the bitwise/shift-logic operation. \emph{Example:} {\scriptsize\ttfamily \mnem{AND.i32} \dst{r0}, \src{r1}, \src{r2}\imod{.not}} (\textbf{SPIR-V:} \texttt{OpNot + OpBitwiseAnd}; \textbf{GLSL:} \texttt{uint y = a \& \textasciitilde{}b;})
  \item \hyperref[enum:widen_int_m]{\texttt{widen0}}: Stage A widening extractor on \texttt{src0} (promotes FP16 \texttt{h0}/\texttt{h1} to FP32, or sign/zero-extends \texttt{b0..b3}/\texttt{h0..h1}/\texttt{w0..w1} to 32/64-bit before \texttt{abs}/\texttt{neg}). \emph{Example:} {\scriptsize\ttfamily \mnem{IADD.s32} \dst{r0}, \src{r1}\imod{.h0}, \src{r2}} (\textbf{SPIR-V:} \texttt{OpSConvert / OpFConvert}; \textbf{GLSL:} \texttt{int c = int(int16\_t(a)) + b;})
  \item \texttt{dest\_register\_index}: 6-bit destination GPR index (\texttt{r0..r63}) receiving the Stage C writeback result (even-aligned pair \texttt{[rN:rN+1]} for 64-bit ops). \emph{Example:} {\scriptsize\ttfamily \mnem{FMA.f32} \dst{r0}, \src{r1}, \src{r2}, \src{r3}} (\textbf{SPIR-V:} \texttt{SSA Result ID}; \textbf{GLSL:} \texttt{float r0 = fma(r1, r2, r3);})
  \item \hyperref[enum:dest_width_m]{\texttt{dest\_width}}: Stage C destination lane write-mask (\texttt{none} = full 32-bit, \texttt{h0} = low 16-bit, \texttt{h1} = high 16-bit without bit-shifting). \emph{Example:} {\scriptsize\ttfamily \mnem{FADD.v2f16} \dst{r0}\imod{.h1}, \src{r1}, \src{r2}} (\textbf{SPIR-V:} \texttt{OpVectorInsertDynamic / WriteMask}; \textbf{GLSL:} \texttt{dst.y = a.y + b.y;})
  \item \texttt{opcode1}: Primary sub-opcode within the instruction class uniquely identifying the operation family. \emph{Example:} {\scriptsize\ttfamily \mnem{IADD.s32} \dst{r0}, \src{r1}, \src{r2}} (\textbf{SPIR-V:} \texttt{OpIAdd (operation selector)}; \textbf{GLSL:} \texttt{// Primary opcode bits})
  \item \texttt{instruction\_class}: 4-bit major instruction class selector (\texttt{bits 48:51}) partitioning the 64-bit Valhall opcode space (e.g., \texttt{FMA}, \texttt{CVT}, \texttt{SFU}, \texttt{LS}, \texttt{TEX}, \texttt{BRANCH}). \emph{Example:} {\scriptsize\ttfamily \mnem{FADD.f32} \dst{r0}, \src{r1}, \src{r2}} (\textbf{SPIR-V:} \texttt{N/A (encoding class)}; \textbf{GLSL:} \texttt{// Hardware execution pipe / decode class})
  \item \texttt{fau\_page\_index}: Selects the FAU (Fast Access Uniform) page (\texttt{0..3}) when a source operand references hardware small constants, thread IDs, or lane constants. \emph{Example:} {\scriptsize\ttfamily \mnem{FADD.f32} \dst{r0}, \src{r1}, \src{k1f}} (\textbf{SPIR-V:} \texttt{OpConstant / BuiltIn}; \textbf{GLSL:} \texttt{float y = x + 1.0;})
  \item \hyperref[enum:flow_control_m]{\texttt{flow}}: 4-bit embedded flow-control \& scoreboard wait modifier (\texttt{bits 59:62}: \texttt{none}, \texttt{wait0..wait0126}, \texttt{wait}, \texttt{reconverge}, \texttt{discard}, \texttt{end}). \emph{Example:} {\scriptsize\ttfamily \mnem{FADD.f32}\imod{.wait0.reconverge} \dst{r0}, \src{r1}, \src{r2}} (\textbf{SPIR-V:} \texttt{OpControlBarrier / OpReturn}; \textbf{GLSL:} \texttt{// Scoreboard wait + warp reconvergence})
\end{tightitemize}
\clearpage

\subsection*{ARSHIFT\_XOR}
\addcontentsline{toc}{subsection}{ARSHIFT\_XOR}
For each SIMD element $i \in \{0, \dots, N-1\}$ ($N=1$ for \texttt{i32}/\texttt{i64}, $N=2$ for \texttt{v2i16}, $N=4$ for \texttt{v4i8} on v9--v10, lane width $W$), computes arithmetic right shift \texttt{s = (int\_W(val) \textgreater{}\textgreater{} (amt \& (W - 1)))} where \texttt{val = src0[i]} and the 8-bit shift count \texttt{amt} is extracted from \texttt{src1} via \texttt{lanes1}, then combines with \texttt{src2[i]} via bitwise XOR \texttt{res = s \textasciicircum{} (not2 ? \textasciitilde{}src2[i] : src2[i])} and writes \texttt{not\_result ? \textasciitilde{}res : res} to \texttt{dest[i]}.
\par\textbf{Notes:}
\begin{tightitemize}
  \item Shift amount is masked modulo lane width (\texttt{wrapping\_shl}/\texttt{wrapping\_shr}/\texttt{rotate\_*}, \texttt{kraid/ops.rs:3984-4009}).
  \item On \texttt{v2i16}, \texttt{src1} is \texttt{V2I8} (\texttt{ops.rs:3953}) with byte lanes selected by \texttt{lanes1} (default \texttt{b02}: byte 0 -\textgreater{} lane 0, byte 2 -\textgreater{} lane 1).
  \item \textbf{Source references:} \texttt{kraid/ops.rs:3878-4025}, \texttt{encode\_v9.rs:2940-3093}, \texttt{hw\_tests.rs:2148-2198}.
\end{tightitemize}
\codebox{\mnem{ARSHIFT\_XOR}\imod{\{.flow\}}\imod{\{.not\_result\}} \dst{dest}\imod{\{.dest\_width\}}, \src{src0}\imod{\{.widen0\}}, \src{src1}\imod{\{.lanes1\}}, \src{src2}\imod{\{.not2\}}\\\mnem{ARSHIFT\_XOR}\imod{\{.flow\}}\imod{\{.not\_result\}} \dst{dest}, \src{src0}\imod{\{.widen0\}}, \src{src1}\imod{\{.lanes1\}}, \src{src2}\imod{\{.not2\}}\\\mnem{ARSHIFT\_XOR}\imod{\{.flow\}}\imod{\{.not\_result\}} \dst{dest}\imod{\{.dest\_width\}}, \src{src0}\imod{\{.swz0\}}, \src{src1}\imod{\{.lanes1\}}, \src{src2}\imod{\{.not2\}}\\\mnem{ARSHIFT\_XOR}\imod{\{.flow\}}\imod{\{.not\_result\}} \dst{dest}\imod{\{.dest\_width\}}, \src{src0}\imod{\{.lanes0\}}, \src{src1}\imod{\{.lanes1\}}, \src{src2}\imod{\{.not2\}}}
\textbf{Operands:}
\begin{tightitemize}\footnotesize
  \item \texttt{\dst{dest}} (\texttt{dest\_register\_index}, \texttt{dest\_width}): Destination GPR receiving the arithmetic right shift combined via bitwise XOR (inverted if \texttt{.not\_result}).
  \item \texttt{\src{src0}}: Input integer value to shift/rotate.
  \item \texttt{\src{src1}}: 8-bit shift/rotate count per lane (selected via \texttt{lanes1}, masked modulo lane width \texttt{W}).
  \item \texttt{\src{src2}}: Second operand combined via bitwise XOR after shifting (inverted if \texttt{.not2}).
\end{tightitemize}
\par\textbf{Closest SPIR-V Op:} \texttt{OpShiftRightArithmetic + OpBitwiseXor}
\par\textbf{Pseudo-GLSL (Operational Semantics):}
\codebox{for (int i = 0; i \textless{} N; ++i) \{\\    uint s = extractByte(src1, lanes1, i) \& (W - 1);\\    auto r = (int(src0[i]) \textgreater{}\textgreater{} s) \textasciicircum{} (flag\_not2 ? \textasciitilde{}src2[i] : src2[i]);\\    dest[i] = flag\_not\_result ? \textasciitilde{}r : r;\\\}}
\textbf{Execution Context:}
\begin{tightitemize}\footnotesize
  \item \textbf{Compute stage preloads (\hyperref[sec:csf-spd]{\texttt{SPD\_k}}, \texttt{r16..r23}):} Hardware pre-fills \texttt{r55..r57} with \texttt{WorkGroupID.xyz + JOB\_OFFSET} (from \hyperref[sec:csf-regs]{\texttt{COMPUTESR::JOB\_OFFSET\_X/Y/Z}} \texttt{r34..r36}) and \texttt{r60..r62} with \texttt{LocalInvocationID.xyz} or \texttt{GlobalInvocationID.xyz} (from \texttt{WG\_SIZE} \texttt{r33} \& \texttt{JOB\_SIZE} \texttt{r37..r39}, selected by \texttt{ShaderProgram::preload}) so compute shaders can derive thread/element indices without memory loads.
  \item \textbf{Vertex stage preloads (\hyperref[sec:csf-spd]{\texttt{SPD\_0}/\texttt{SPD\_1}}, \texttt{r16..r19}):} The hardware Index Fetcher reads \hyperref[sec:csf-regs]{\texttt{IDVSSR::INDEX\_BUFFER}} (\texttt{r54:r55}) and pre-fills \texttt{r59, r60} with \texttt{VertexIndex} (\texttt{+ VERTEX\_OFFSET + GLOBAL\_ATTR\_OFF}) and \texttt{r61} with \texttt{InstanceIndex} (\texttt{+ INSTANCE\_OFFSET}) so Position and Varying shaders can pull vertex attributes directly.
  \item \textbf{Inline FAU \& Special Operands (\hyperref[sec:csf-tsd-fau]{\texttt{FAU\_k}}, \texttt{r8..r15}):} Any 8-bit source with \texttt{mode=2} reads zero-latency uniform \texttt{u0..u31} from \texttt{FAU\_k} (\texttt{slots64 = bits[63:56]}); \texttt{mode=3} reads \texttt{SmallConstant} (Table~1.2) or \texttt{FauSpecial} (\texttt{pc}, \texttt{lane\_id}, \texttt{lane\_width}, \texttt{tls\_ptr}, \texttt{wls\_ptr}, \texttt{blend\_descriptor\_0..7}).
\end{tightitemize}
\par\textbf{Worked Examples:}
\par\textbf{Example 1 (Standard fused shift + bitwise op):}
\codebox{\mnem{ARSHIFT\_XOR.i32} \dst{r0}, \src{r1}, \src{r2}\imod{.b0}, \src{r3}}
With \texttt{r1 = 0x80000000}, \texttt{r2.b0 = 4} (masked mod 32): shift yields \texttt{0xF8000000}, then \texttt{\textasciicircum{} 0x000000FF} yields \texttt{r0 = 0xF80000FF}.
\par\textbf{Example 2 (Modulo-32 shift count \& .not\_result):}
\codebox{\mnem{ARSHIFT\_XOR.i32}\imod{.not\_result} \dst{r0}, \src{r1}, \src{r2}\imod{.b0}, \src{r3}}
With \texttt{r1 = 0x80000000}, \texttt{r2.b0 = 31} (effective shift \texttt{31}): shift yields \texttt{0xFFFFFFFF}, \texttt{\textasciicircum{} 0x000000FF} yields \texttt{0xFFFFFF00}, and \texttt{.not\_result} inverts -\textgreater{} \texttt{r0 = 0x000000FF}.

\begin{center}\footnotesize
\begin{tabularx}{\linewidth}{@{}lllX@{}}
\toprule
Variant & Arch & Class & Opcode Fields \\
\midrule
\texttt{(none)} & v9--v14 & \texttt{R3} & \texttt{opcode1=SHIFT\_BITWISE\_I32} \\
\texttt{(none)} & v9--v14 & \texttt{R3} & \texttt{opcode1=SHIFT\_BITWISE\_I32} \\
\texttt{(none)} & v9--v14 & \texttt{R3} & \texttt{opcode1=LOGIC\_SHIFT\_I32} \\
\texttt{(none)} & v9--v14 & \texttt{D3} & \texttt{opcode1=SHIFT\_BITWISE\_I64} \\
\texttt{(none)} & v9--v14 & \texttt{D3} & \texttt{opcode1=LOGIC\_SHIFT\_I64} \\
\texttt{(none)} & v9--v14 & \texttt{R3} & \texttt{opcode1=SHIFT\_BITWISE\_V2I16} \\
\texttt{(none)} & v9--v14 & \texttt{R3} & \texttt{opcode1=LOGIC\_SHIFT\_V2I16} \\
\texttt{(none)} & v9--v14 & \texttt{R3} & \texttt{opcode1=SHIFT\_BITWISE\_V4I8} \\
\texttt{(none)} & v9--v14 & \texttt{R3} & \texttt{opcode1=LOGIC\_SHIFT\_V4I8} \\
\bottomrule
\end{tabularx}
\end{center}
\begin{center}\footnotesize
\begin{tabularx}{0.88\linewidth}{@{}l l X@{}}
\toprule
Field Name & Bits & Type / Fixed Value \\
\midrule
\texttt{src0} & \texttt{0:7} & \texttt{SourceEncoding} \\
\texttt{src1} & \texttt{8:15} & \texttt{SourceEncoding} \\
\texttt{src2} & \texttt{16:23} & \texttt{SourceEncoding} \\
\texttt{bitwise\_operation} & \texttt{24:25} & \texttt{exact=xor} \\
\hyperref[enum:shift_bitwise_i32_lanes_m]{\texttt{lanes1}} & \texttt{26:29} & \texttt{shift\_bitwise\_i32\_lanes\_m (def b0)} \\
\hyperref[enum:not_m]{\texttt{not\_result}} & \texttt{30} & \texttt{not\_m (def none)} \\
\textit{Reserved} & \texttt{31} & \texttt{zero} \\
\texttt{shift\_operation} & \texttt{32:34} & \texttt{exact=arshift} \\
\hyperref[enum:not_m]{\texttt{not2}} & \texttt{35} & \texttt{not\_m (def none)} \\
\hyperref[enum:widen_int_m]{\texttt{widen0}} & \texttt{36:39} & \texttt{widen\_int\_m (def none)} \\
\texttt{dest\_register\_index} & \texttt{40:45} & \texttt{uint} \\
\hyperref[enum:dest_width_m]{\texttt{dest\_width}} & \texttt{46:47} & \texttt{dest\_width\_m} \\
\texttt{opcode1} & \texttt{48:51} & \texttt{exact=SHIFT\_BITWISE\_I32} \\
\texttt{instruction\_class} & \texttt{52:56} & \texttt{exact=R3} \\
\texttt{fau\_page\_index} & \texttt{57:58} & \texttt{uint} \\
\hyperref[enum:flow_control_m]{\texttt{flow}} & \texttt{59:62} & \texttt{flow\_control\_m} \\
\textit{Reserved} & \texttt{63} & \texttt{zero} \\
\bottomrule
\end{tabularx}
\end{center}
\textbf{Fields Explained:}
\begin{tightitemize}\footnotesize
  \item \texttt{src0}: Primary source operand (\texttt{r0..r63}, \texttt{u0..u63} uniform, \texttt{k*} FAU constant, or \texttt{t0..t2} temporary; \texttt{\textasciicircum{}} marks discard/last-use). \emph{Example:} {\scriptsize\ttfamily \mnem{FADD.f32} \dst{r0}, \src{r1}, \src{r2}} (\textbf{SPIR-V:} \texttt{OpFAdd (Operand 0)}; \textbf{GLSL:} \texttt{float c = a + b;})
  \item \texttt{src1}: Secondary source operand (register, uniform, or FAU constant; in shifts, supplies the shift count byte). \emph{Example:} {\scriptsize\ttfamily \mnem{LSHIFT.i32} \dst{r0}, \src{r1}, \src{r2}\imod{.b0}} (\textbf{SPIR-V:} \texttt{OpShiftLeftLogical (Shift)}; \textbf{GLSL:} \texttt{uint y = a \textless{}\textless{} s;})
  \item \texttt{src2}: Third source operand (FMA addend, \texttt{MUX}/\texttt{CSEL} selector, \texttt{IADDC} carry-in, or bitwise mask in fused shift-logic). \emph{Example:} {\scriptsize\ttfamily \mnem{FMA.f32} \dst{r0}, \src{r1}, \src{r2}, \src{r3}} (\textbf{SPIR-V:} \texttt{GLSL.std.450 Fma (Addend)}; \textbf{GLSL:} \texttt{float d = fma(a, b, c);})
  \item \texttt{bitwise\_operation}: Selects the bitwise ALU op (\texttt{and}, \texttt{or}, \texttt{xor}, \texttt{lshift\_and}, etc.) fused with shifts. \emph{Example:} {\scriptsize\ttfamily \mnem{LSHIFT\_OR.i32} \dst{r0}, \src{r1}, \src{r2}, \src{r3}} (\textbf{SPIR-V:} \texttt{OpShiftLeftLogical + OpBitwiseOr}; \textbf{GLSL:} \texttt{uint y = (a \textless{}\textless{} s) \textbar{} b;})
  \item \hyperref[enum:shift_bitwise_i32_lanes_m]{\texttt{lanes1}}: Stage A byte-lane selector on \texttt{src1} (selects shift-count byte \texttt{b0..b3} on 32/64-bit shifts, \texttt{bXY} on \texttt{v2i16} shifts, or 4-byte swizzle on \texttt{v4*8}). \emph{Example:} {\scriptsize\ttfamily \mnem{LSHIFT.i32} \dst{r0}, \src{r1}, \src{r2}\imod{.b1}} (\textbf{SPIR-V:} \texttt{OpBitFieldUExtract + OpShiftLeftLogical}; \textbf{GLSL:} \texttt{uint y = a \textless{}\textless{} ((s \textgreater{}\textgreater{} 8) \& 0xFFu);})
  \item \hyperref[enum:not_m]{\texttt{not\_result}}: Bitwise-inverts the final output of a bitwise/shift-bitwise instruction (\textasciitilde{}result). \emph{Example:} {\scriptsize\ttfamily \mnem{AND.i32}\imod{.not\_result} \dst{r0}, \src{r1}, \src{r2}} (\textbf{SPIR-V:} \texttt{OpBitwiseAnd + OpNot}; \textbf{GLSL:} \texttt{uint y = \textasciitilde{}(a \& b);})
  \item \texttt{shift\_operation}: Selects shift direction/type (\texttt{lshift}, \texttt{rshift}, \texttt{arshift}) in fused shift-bitwise instructions. \emph{Example:} {\scriptsize\ttfamily \mnem{RSHIFT\_AND.i32} \dst{r0}, \src{r1}, \src{r2}, \src{r3}} (\textbf{SPIR-V:} \texttt{OpShiftRightLogical + OpBitwiseAnd}; \textbf{GLSL:} \texttt{uint y = (a \textgreater{}\textgreater{} s) \& m;})
  \item \hyperref[enum:not_m]{\texttt{not2}}: Stage A bitwise complement (\texttt{\textasciitilde{}src2}) applied to \texttt{src2} before the bitwise/shift-logic operation. \emph{Example:} {\scriptsize\ttfamily \mnem{AND.i32} \dst{r0}, \src{r1}, \src{r2}\imod{.not}} (\textbf{SPIR-V:} \texttt{OpNot + OpBitwiseAnd}; \textbf{GLSL:} \texttt{uint y = a \& \textasciitilde{}b;})
  \item \hyperref[enum:widen_int_m]{\texttt{widen0}}: Stage A widening extractor on \texttt{src0} (promotes FP16 \texttt{h0}/\texttt{h1} to FP32, or sign/zero-extends \texttt{b0..b3}/\texttt{h0..h1}/\texttt{w0..w1} to 32/64-bit before \texttt{abs}/\texttt{neg}). \emph{Example:} {\scriptsize\ttfamily \mnem{IADD.s32} \dst{r0}, \src{r1}\imod{.h0}, \src{r2}} (\textbf{SPIR-V:} \texttt{OpSConvert / OpFConvert}; \textbf{GLSL:} \texttt{int c = int(int16\_t(a)) + b;})
  \item \texttt{dest\_register\_index}: 6-bit destination GPR index (\texttt{r0..r63}) receiving the Stage C writeback result (even-aligned pair \texttt{[rN:rN+1]} for 64-bit ops). \emph{Example:} {\scriptsize\ttfamily \mnem{FMA.f32} \dst{r0}, \src{r1}, \src{r2}, \src{r3}} (\textbf{SPIR-V:} \texttt{SSA Result ID}; \textbf{GLSL:} \texttt{float r0 = fma(r1, r2, r3);})
  \item \hyperref[enum:dest_width_m]{\texttt{dest\_width}}: Stage C destination lane write-mask (\texttt{none} = full 32-bit, \texttt{h0} = low 16-bit, \texttt{h1} = high 16-bit without bit-shifting). \emph{Example:} {\scriptsize\ttfamily \mnem{FADD.v2f16} \dst{r0}\imod{.h1}, \src{r1}, \src{r2}} (\textbf{SPIR-V:} \texttt{OpVectorInsertDynamic / WriteMask}; \textbf{GLSL:} \texttt{dst.y = a.y + b.y;})
  \item \texttt{opcode1}: Primary sub-opcode within the instruction class uniquely identifying the operation family. \emph{Example:} {\scriptsize\ttfamily \mnem{IADD.s32} \dst{r0}, \src{r1}, \src{r2}} (\textbf{SPIR-V:} \texttt{OpIAdd (operation selector)}; \textbf{GLSL:} \texttt{// Primary opcode bits})
  \item \texttt{instruction\_class}: 4-bit major instruction class selector (\texttt{bits 48:51}) partitioning the 64-bit Valhall opcode space (e.g., \texttt{FMA}, \texttt{CVT}, \texttt{SFU}, \texttt{LS}, \texttt{TEX}, \texttt{BRANCH}). \emph{Example:} {\scriptsize\ttfamily \mnem{FADD.f32} \dst{r0}, \src{r1}, \src{r2}} (\textbf{SPIR-V:} \texttt{N/A (encoding class)}; \textbf{GLSL:} \texttt{// Hardware execution pipe / decode class})
  \item \texttt{fau\_page\_index}: Selects the FAU (Fast Access Uniform) page (\texttt{0..3}) when a source operand references hardware small constants, thread IDs, or lane constants. \emph{Example:} {\scriptsize\ttfamily \mnem{FADD.f32} \dst{r0}, \src{r1}, \src{k1f}} (\textbf{SPIR-V:} \texttt{OpConstant / BuiltIn}; \textbf{GLSL:} \texttt{float y = x + 1.0;})
  \item \hyperref[enum:flow_control_m]{\texttt{flow}}: 4-bit embedded flow-control \& scoreboard wait modifier (\texttt{bits 59:62}: \texttt{none}, \texttt{wait0..wait0126}, \texttt{wait}, \texttt{reconverge}, \texttt{discard}, \texttt{end}). \emph{Example:} {\scriptsize\ttfamily \mnem{FADD.f32}\imod{.wait0.reconverge} \dst{r0}, \src{r1}, \src{r2}} (\textbf{SPIR-V:} \texttt{OpControlBarrier / OpReturn}; \textbf{GLSL:} \texttt{// Scoreboard wait + warp reconvergence})
\end{tightitemize}
\clearpage

\subsection*{LROT}
\addcontentsline{toc}{subsection}{LROT}
For each SIMD element $i \in \{0, \dots, N-1\}$ ($N=1$ for \texttt{i32}/\texttt{i64}, $N=2$ for \texttt{v2i16}, $N=4$ for \texttt{v4i8} on v9--v10, lane width $W$), computes rotate left \texttt{s = rotate\_left(val, amt \& (W - 1))} where \texttt{val = src0[i]} and the 8-bit shift count \texttt{amt} is extracted from \texttt{src1} via \texttt{lanes1}; writes \texttt{not\_result ? \textasciitilde{}s : s} to \texttt{dest[i]}.
\par\textbf{Notes:}
\begin{tightitemize}
  \item Shift amount is masked modulo lane width (\texttt{wrapping\_shl}/\texttt{wrapping\_shr}/\texttt{rotate\_*}, \texttt{kraid/ops.rs:3984-4009}).
  \item On \texttt{v2i16}, \texttt{src1} is \texttt{V2I8} (\texttt{ops.rs:3953}) with byte lanes selected by \texttt{lanes1} (default \texttt{b02}: byte 0 -\textgreater{} lane 0, byte 2 -\textgreater{} lane 1).
  \item \textbf{Source references:} \texttt{kraid/ops.rs:3878-4025}, \texttt{encode\_v9.rs:2940-3093}, \texttt{hw\_tests.rs:2148-2198}.
\end{tightitemize}
\codebox{\mnem{LROT}\imod{\{.flow\}}\imod{\{.not\_result\}} \dst{dest}\imod{\{.dest\_width\}}, \src{src0}\imod{\{.widen0\}}, \src{src1}\imod{\{.lanes1\}}\\\mnem{LROT}\imod{\{.flow\}}\imod{\{.not\_result\}} \dst{dest}, \src{src0}\imod{\{.widen0\}}, \src{src1}\imod{\{.lanes1\}}\\\mnem{LROT}\imod{\{.flow\}}\imod{\{.not\_result\}} \dst{dest}\imod{\{.dest\_width\}}, \src{src0}\imod{\{.swz0\}}, \src{src1}\imod{\{.lanes1\}}\\\mnem{LROT}\imod{\{.flow\}}\imod{\{.not\_result\}} \dst{dest}\imod{\{.dest\_width\}}, \src{src0}\imod{\{.lanes0\}}, \src{src1}\imod{\{.lanes1\}}}
\textbf{Operands:}
\begin{tightitemize}\footnotesize
  \item \texttt{\dst{dest}} (\texttt{dest\_register\_index}, \texttt{dest\_width}): Destination GPR receiving the rotate left (inverted if \texttt{.not\_result}).
  \item \texttt{\src{src0}}: Input integer value to shift/rotate.
  \item \texttt{\src{src1}}: 8-bit shift/rotate count per lane (selected via \texttt{lanes1}, masked modulo lane width \texttt{W}).
\end{tightitemize}
\par\textbf{Closest SPIR-V Op:} \texttt{None (lowered from `(x\textless{}\textless{}s)\textbar{}(x\textgreater{}\textgreater{}(32-s))`)}
\par\textbf{Pseudo-GLSL (Operational Semantics):}
\codebox{for (int i = 0; i \textless{} N; ++i) \{\\    uint s = extractByte(src1, lanes1, i) \& (W - 1);\\    auto r = ((src0[i] \textless{}\textless{} s) \textbar{} (src0[i] \textgreater{}\textgreater{} ((W - s) \& (W - 1))));\\    dest[i] = flag\_not\_result ? \textasciitilde{}r : r;\\\}}
\textbf{Execution Context:}
\begin{tightitemize}\footnotesize
  \item \textbf{Compute stage preloads (\hyperref[sec:csf-spd]{\texttt{SPD\_k}}, \texttt{r16..r23}):} Hardware pre-fills \texttt{r55..r57} with \texttt{WorkGroupID.xyz + JOB\_OFFSET} (from \hyperref[sec:csf-regs]{\texttt{COMPUTESR::JOB\_OFFSET\_X/Y/Z}} \texttt{r34..r36}) and \texttt{r60..r62} with \texttt{LocalInvocationID.xyz} or \texttt{GlobalInvocationID.xyz} (from \texttt{WG\_SIZE} \texttt{r33} \& \texttt{JOB\_SIZE} \texttt{r37..r39}, selected by \texttt{ShaderProgram::preload}) so compute shaders can derive thread/element indices without memory loads.
  \item \textbf{Vertex stage preloads (\hyperref[sec:csf-spd]{\texttt{SPD\_0}/\texttt{SPD\_1}}, \texttt{r16..r19}):} The hardware Index Fetcher reads \hyperref[sec:csf-regs]{\texttt{IDVSSR::INDEX\_BUFFER}} (\texttt{r54:r55}) and pre-fills \texttt{r59, r60} with \texttt{VertexIndex} (\texttt{+ VERTEX\_OFFSET + GLOBAL\_ATTR\_OFF}) and \texttt{r61} with \texttt{InstanceIndex} (\texttt{+ INSTANCE\_OFFSET}) so Position and Varying shaders can pull vertex attributes directly.
  \item \textbf{Inline FAU \& Special Operands (\hyperref[sec:csf-tsd-fau]{\texttt{FAU\_k}}, \texttt{r8..r15}):} Any 8-bit source with \texttt{mode=2} reads zero-latency uniform \texttt{u0..u31} from \texttt{FAU\_k} (\texttt{slots64 = bits[63:56]}); \texttt{mode=3} reads \texttt{SmallConstant} (Table~1.2) or \texttt{FauSpecial} (\texttt{pc}, \texttt{lane\_id}, \texttt{lane\_width}, \texttt{tls\_ptr}, \texttt{wls\_ptr}, \texttt{blend\_descriptor\_0..7}).
\end{tightitemize}
\par\textbf{Worked Examples:}
\par\textbf{Example 1 (Canonical 32-bit shift/rotate):}
\codebox{\mnem{LROT.i32} \dst{r0}, \src{r1}, \src{r2}\imod{.b0}}
With \texttt{r1 = 0x80000001}, \texttt{r2.b0 = 1}: yields \texttt{r0 = 0x00000003}.
\par\textbf{Example 2 (Byte-lane selector \& modulo-32 masking):}
\codebox{\mnem{LROT.i32} \dst{r0}, \src{r1}, \src{r2}\imod{.b1}}
With \texttt{r1 = 0x12345678}, \texttt{r2.b1 = 8} (shift count masked mod 32 -\textgreater{} \texttt{8}): yields \texttt{r0 = 0x34567812}.

\begin{center}\footnotesize
\begin{tabularx}{\linewidth}{@{}lllX@{}}
\toprule
Variant & Arch & Class & Opcode Fields \\
\midrule
\texttt{(none)} & v9--v14 & \texttt{R3} & \texttt{opcode1=SHIFT\_BITWISE\_I32} \\
\texttt{(none)} & v9--v14 & \texttt{R3} & \texttt{opcode1=LOGIC\_SHIFT\_I32} \\
\texttt{(none)} & v9--v14 & \texttt{D3} & \texttt{opcode1=SHIFT\_BITWISE\_I64} \\
\texttt{(none)} & v9--v14 & \texttt{D3} & \texttt{opcode1=LOGIC\_SHIFT\_I64} \\
\texttt{(none)} & v9--v14 & \texttt{R3} & \texttt{opcode1=SHIFT\_BITWISE\_V2I16} \\
\texttt{(none)} & v9--v14 & \texttt{R3} & \texttt{opcode1=LOGIC\_SHIFT\_V2I16} \\
\texttt{(none)} & v9--v14 & \texttt{R3} & \texttt{opcode1=SHIFT\_BITWISE\_V4I8} \\
\texttt{(none)} & v9--v14 & \texttt{R3} & \texttt{opcode1=LOGIC\_SHIFT\_V4I8} \\
\bottomrule
\end{tabularx}
\end{center}
\begin{center}\footnotesize
\begin{tabularx}{0.88\linewidth}{@{}l l X@{}}
\toprule
Field Name & Bits & Type / Fixed Value \\
\midrule
\texttt{src0} & \texttt{0:7} & \texttt{SourceEncoding} \\
\texttt{src1} & \texttt{8:15} & \texttt{SourceEncoding} \\
\texttt{src2} & \texttt{16:23} & \texttt{exact=192} \\
\texttt{bitwise\_operation} & \texttt{24:25} & \texttt{exact=or} \\
\hyperref[enum:shift_bitwise_i32_lanes_m]{\texttt{lanes1}} & \texttt{26:29} & \texttt{shift\_bitwise\_i32\_lanes\_m (def b0)} \\
\hyperref[enum:not_m]{\texttt{not\_result}} & \texttt{30} & \texttt{not\_m (def none)} \\
\textit{Reserved} & \texttt{31} & \texttt{zero} \\
\texttt{shift\_operation} & \texttt{32:34} & \texttt{exact=lrot} \\
\hyperref[enum:not_m]{\texttt{not2}} & \texttt{35} & \texttt{exact=none} \\
\hyperref[enum:widen_int_m]{\texttt{widen0}} & \texttt{36:39} & \texttt{widen\_int\_m (def none)} \\
\texttt{dest\_register\_index} & \texttt{40:45} & \texttt{uint} \\
\hyperref[enum:dest_width_m]{\texttt{dest\_width}} & \texttt{46:47} & \texttt{dest\_width\_m} \\
\texttt{opcode1} & \texttt{48:51} & \texttt{exact=SHIFT\_BITWISE\_I32} \\
\texttt{instruction\_class} & \texttt{52:56} & \texttt{exact=R3} \\
\texttt{fau\_page\_index} & \texttt{57:58} & \texttt{uint} \\
\hyperref[enum:flow_control_m]{\texttt{flow}} & \texttt{59:62} & \texttt{flow\_control\_m} \\
\textit{Reserved} & \texttt{63} & \texttt{zero} \\
\bottomrule
\end{tabularx}
\end{center}
\textbf{Fields Explained:}
\begin{tightitemize}\footnotesize
  \item \texttt{src0}: Primary source operand (\texttt{r0..r63}, \texttt{u0..u63} uniform, \texttt{k*} FAU constant, or \texttt{t0..t2} temporary; \texttt{\textasciicircum{}} marks discard/last-use). \emph{Example:} {\scriptsize\ttfamily \mnem{FADD.f32} \dst{r0}, \src{r1}, \src{r2}} (\textbf{SPIR-V:} \texttt{OpFAdd (Operand 0)}; \textbf{GLSL:} \texttt{float c = a + b;})
  \item \texttt{src1}: Secondary source operand (register, uniform, or FAU constant; in shifts, supplies the shift count byte). \emph{Example:} {\scriptsize\ttfamily \mnem{LSHIFT.i32} \dst{r0}, \src{r1}, \src{r2}\imod{.b0}} (\textbf{SPIR-V:} \texttt{OpShiftLeftLogical (Shift)}; \textbf{GLSL:} \texttt{uint y = a \textless{}\textless{} s;})
  \item \texttt{src2}: Third source operand (FMA addend, \texttt{MUX}/\texttt{CSEL} selector, \texttt{IADDC} carry-in, or bitwise mask in fused shift-logic). \emph{Example:} {\scriptsize\ttfamily \mnem{FMA.f32} \dst{r0}, \src{r1}, \src{r2}, \src{r3}} (\textbf{SPIR-V:} \texttt{GLSL.std.450 Fma (Addend)}; \textbf{GLSL:} \texttt{float d = fma(a, b, c);})
  \item \texttt{bitwise\_operation}: Selects the bitwise ALU op (\texttt{and}, \texttt{or}, \texttt{xor}, \texttt{lshift\_and}, etc.) fused with shifts. \emph{Example:} {\scriptsize\ttfamily \mnem{LSHIFT\_OR.i32} \dst{r0}, \src{r1}, \src{r2}, \src{r3}} (\textbf{SPIR-V:} \texttt{OpShiftLeftLogical + OpBitwiseOr}; \textbf{GLSL:} \texttt{uint y = (a \textless{}\textless{} s) \textbar{} b;})
  \item \hyperref[enum:shift_bitwise_i32_lanes_m]{\texttt{lanes1}}: Stage A byte-lane selector on \texttt{src1} (selects shift-count byte \texttt{b0..b3} on 32/64-bit shifts, \texttt{bXY} on \texttt{v2i16} shifts, or 4-byte swizzle on \texttt{v4*8}). \emph{Example:} {\scriptsize\ttfamily \mnem{LSHIFT.i32} \dst{r0}, \src{r1}, \src{r2}\imod{.b1}} (\textbf{SPIR-V:} \texttt{OpBitFieldUExtract + OpShiftLeftLogical}; \textbf{GLSL:} \texttt{uint y = a \textless{}\textless{} ((s \textgreater{}\textgreater{} 8) \& 0xFFu);})
  \item \hyperref[enum:not_m]{\texttt{not\_result}}: Bitwise-inverts the final output of a bitwise/shift-bitwise instruction (\textasciitilde{}result). \emph{Example:} {\scriptsize\ttfamily \mnem{AND.i32}\imod{.not\_result} \dst{r0}, \src{r1}, \src{r2}} (\textbf{SPIR-V:} \texttt{OpBitwiseAnd + OpNot}; \textbf{GLSL:} \texttt{uint y = \textasciitilde{}(a \& b);})
  \item \texttt{shift\_operation}: Selects shift direction/type (\texttt{lshift}, \texttt{rshift}, \texttt{arshift}) in fused shift-bitwise instructions. \emph{Example:} {\scriptsize\ttfamily \mnem{RSHIFT\_AND.i32} \dst{r0}, \src{r1}, \src{r2}, \src{r3}} (\textbf{SPIR-V:} \texttt{OpShiftRightLogical + OpBitwiseAnd}; \textbf{GLSL:} \texttt{uint y = (a \textgreater{}\textgreater{} s) \& m;})
  \item \hyperref[enum:not_m]{\texttt{not2}}: Stage A bitwise complement (\texttt{\textasciitilde{}src2}) applied to \texttt{src2} before the bitwise/shift-logic operation. \emph{Example:} {\scriptsize\ttfamily \mnem{AND.i32} \dst{r0}, \src{r1}, \src{r2}\imod{.not}} (\textbf{SPIR-V:} \texttt{OpNot + OpBitwiseAnd}; \textbf{GLSL:} \texttt{uint y = a \& \textasciitilde{}b;})
  \item \hyperref[enum:widen_int_m]{\texttt{widen0}}: Stage A widening extractor on \texttt{src0} (promotes FP16 \texttt{h0}/\texttt{h1} to FP32, or sign/zero-extends \texttt{b0..b3}/\texttt{h0..h1}/\texttt{w0..w1} to 32/64-bit before \texttt{abs}/\texttt{neg}). \emph{Example:} {\scriptsize\ttfamily \mnem{IADD.s32} \dst{r0}, \src{r1}\imod{.h0}, \src{r2}} (\textbf{SPIR-V:} \texttt{OpSConvert / OpFConvert}; \textbf{GLSL:} \texttt{int c = int(int16\_t(a)) + b;})
  \item \texttt{dest\_register\_index}: 6-bit destination GPR index (\texttt{r0..r63}) receiving the Stage C writeback result (even-aligned pair \texttt{[rN:rN+1]} for 64-bit ops). \emph{Example:} {\scriptsize\ttfamily \mnem{FMA.f32} \dst{r0}, \src{r1}, \src{r2}, \src{r3}} (\textbf{SPIR-V:} \texttt{SSA Result ID}; \textbf{GLSL:} \texttt{float r0 = fma(r1, r2, r3);})
  \item \hyperref[enum:dest_width_m]{\texttt{dest\_width}}: Stage C destination lane write-mask (\texttt{none} = full 32-bit, \texttt{h0} = low 16-bit, \texttt{h1} = high 16-bit without bit-shifting). \emph{Example:} {\scriptsize\ttfamily \mnem{FADD.v2f16} \dst{r0}\imod{.h1}, \src{r1}, \src{r2}} (\textbf{SPIR-V:} \texttt{OpVectorInsertDynamic / WriteMask}; \textbf{GLSL:} \texttt{dst.y = a.y + b.y;})
  \item \texttt{opcode1}: Primary sub-opcode within the instruction class uniquely identifying the operation family. \emph{Example:} {\scriptsize\ttfamily \mnem{IADD.s32} \dst{r0}, \src{r1}, \src{r2}} (\textbf{SPIR-V:} \texttt{OpIAdd (operation selector)}; \textbf{GLSL:} \texttt{// Primary opcode bits})
  \item \texttt{instruction\_class}: 4-bit major instruction class selector (\texttt{bits 48:51}) partitioning the 64-bit Valhall opcode space (e.g., \texttt{FMA}, \texttt{CVT}, \texttt{SFU}, \texttt{LS}, \texttt{TEX}, \texttt{BRANCH}). \emph{Example:} {\scriptsize\ttfamily \mnem{FADD.f32} \dst{r0}, \src{r1}, \src{r2}} (\textbf{SPIR-V:} \texttt{N/A (encoding class)}; \textbf{GLSL:} \texttt{// Hardware execution pipe / decode class})
  \item \texttt{fau\_page\_index}: Selects the FAU (Fast Access Uniform) page (\texttt{0..3}) when a source operand references hardware small constants, thread IDs, or lane constants. \emph{Example:} {\scriptsize\ttfamily \mnem{FADD.f32} \dst{r0}, \src{r1}, \src{k1f}} (\textbf{SPIR-V:} \texttt{OpConstant / BuiltIn}; \textbf{GLSL:} \texttt{float y = x + 1.0;})
  \item \hyperref[enum:flow_control_m]{\texttt{flow}}: 4-bit embedded flow-control \& scoreboard wait modifier (\texttt{bits 59:62}: \texttt{none}, \texttt{wait0..wait0126}, \texttt{wait}, \texttt{reconverge}, \texttt{discard}, \texttt{end}). \emph{Example:} {\scriptsize\ttfamily \mnem{FADD.f32}\imod{.wait0.reconverge} \dst{r0}, \src{r1}, \src{r2}} (\textbf{SPIR-V:} \texttt{OpControlBarrier / OpReturn}; \textbf{GLSL:} \texttt{// Scoreboard wait + warp reconvergence})
\end{tightitemize}
\clearpage

\subsection*{LROT\_AND}
\addcontentsline{toc}{subsection}{LROT\_AND}
For each SIMD element $i \in \{0, \dots, N-1\}$ ($N=1$ for \texttt{i32}/\texttt{i64}, $N=2$ for \texttt{v2i16}, $N=4$ for \texttt{v4i8} on v9--v10, lane width $W$), computes rotate left \texttt{s = rotate\_left(val, amt \& (W - 1))} where \texttt{val = src0[i]} and the 8-bit shift count \texttt{amt} is extracted from \texttt{src1} via \texttt{lanes1}, then combines with \texttt{src2[i]} via bitwise AND \texttt{res = s \& (not2 ? \textasciitilde{}src2[i] : src2[i])} and writes \texttt{not\_result ? \textasciitilde{}res : res} to \texttt{dest[i]}.
\par\textbf{Notes:}
\begin{tightitemize}
  \item Shift amount is masked modulo lane width (\texttt{wrapping\_shl}/\texttt{wrapping\_shr}/\texttt{rotate\_*}, \texttt{kraid/ops.rs:3984-4009}).
  \item On \texttt{v2i16}, \texttt{src1} is \texttt{V2I8} (\texttt{ops.rs:3953}) with byte lanes selected by \texttt{lanes1} (default \texttt{b02}: byte 0 -\textgreater{} lane 0, byte 2 -\textgreater{} lane 1).
  \item \textbf{Source references:} \texttt{kraid/ops.rs:3878-4025}, \texttt{encode\_v9.rs:2940-3093}, \texttt{hw\_tests.rs:2148-2198}.
\end{tightitemize}
\codebox{\mnem{LROT\_AND}\imod{\{.flow\}}\imod{\{.not\_result\}} \dst{dest}\imod{\{.dest\_width\}}, \src{src0}\imod{\{.widen0\}}, \src{src1}\imod{\{.lanes1\}}, \src{src2}\imod{\{.not2\}}\\\mnem{LROT\_AND}\imod{\{.flow\}}\imod{\{.not\_result\}} \dst{dest}, \src{src0}\imod{\{.widen0\}}, \src{src1}\imod{\{.lanes1\}}, \src{src2}\imod{\{.not2\}}\\\mnem{LROT\_AND}\imod{\{.flow\}}\imod{\{.not\_result\}} \dst{dest}\imod{\{.dest\_width\}}, \src{src0}\imod{\{.swz0\}}, \src{src1}\imod{\{.lanes1\}}, \src{src2}\imod{\{.not2\}}\\\mnem{LROT\_AND}\imod{\{.flow\}}\imod{\{.not\_result\}} \dst{dest}\imod{\{.dest\_width\}}, \src{src0}\imod{\{.lanes0\}}, \src{src1}\imod{\{.lanes1\}}, \src{src2}\imod{\{.not2\}}}
\textbf{Operands:}
\begin{tightitemize}\footnotesize
  \item \texttt{\dst{dest}} (\texttt{dest\_register\_index}, \texttt{dest\_width}): Destination GPR receiving the rotate left combined via bitwise AND (inverted if \texttt{.not\_result}).
  \item \texttt{\src{src0}}: Input integer value to shift/rotate.
  \item \texttt{\src{src1}}: 8-bit shift/rotate count per lane (selected via \texttt{lanes1}, masked modulo lane width \texttt{W}).
  \item \texttt{\src{src2}}: Second operand combined via bitwise AND after shifting (inverted if \texttt{.not2}).
\end{tightitemize}
\par\textbf{Closest SPIR-V Op:} \texttt{None (lowered from `(x\textless{}\textless{}s)\textbar{}(x\textgreater{}\textgreater{}(32-s))`) + OpBitwiseAnd}
\par\textbf{Pseudo-GLSL (Operational Semantics):}
\codebox{for (int i = 0; i \textless{} N; ++i) \{\\    uint s = extractByte(src1, lanes1, i) \& (W - 1);\\    auto r = ((src0[i] \textless{}\textless{} s) \textbar{} (src0[i] \textgreater{}\textgreater{} ((W - s) \& (W - 1)))) \& (flag\_not2 ? \textasciitilde{}src2[i] : src2[i]);\\    dest[i] = flag\_not\_result ? \textasciitilde{}r : r;\\\}}
\textbf{Execution Context:}
\begin{tightitemize}\footnotesize
  \item \textbf{Compute stage preloads (\hyperref[sec:csf-spd]{\texttt{SPD\_k}}, \texttt{r16..r23}):} Hardware pre-fills \texttt{r55..r57} with \texttt{WorkGroupID.xyz + JOB\_OFFSET} (from \hyperref[sec:csf-regs]{\texttt{COMPUTESR::JOB\_OFFSET\_X/Y/Z}} \texttt{r34..r36}) and \texttt{r60..r62} with \texttt{LocalInvocationID.xyz} or \texttt{GlobalInvocationID.xyz} (from \texttt{WG\_SIZE} \texttt{r33} \& \texttt{JOB\_SIZE} \texttt{r37..r39}, selected by \texttt{ShaderProgram::preload}) so compute shaders can derive thread/element indices without memory loads.
  \item \textbf{Vertex stage preloads (\hyperref[sec:csf-spd]{\texttt{SPD\_0}/\texttt{SPD\_1}}, \texttt{r16..r19}):} The hardware Index Fetcher reads \hyperref[sec:csf-regs]{\texttt{IDVSSR::INDEX\_BUFFER}} (\texttt{r54:r55}) and pre-fills \texttt{r59, r60} with \texttt{VertexIndex} (\texttt{+ VERTEX\_OFFSET + GLOBAL\_ATTR\_OFF}) and \texttt{r61} with \texttt{InstanceIndex} (\texttt{+ INSTANCE\_OFFSET}) so Position and Varying shaders can pull vertex attributes directly.
  \item \textbf{Inline FAU \& Special Operands (\hyperref[sec:csf-tsd-fau]{\texttt{FAU\_k}}, \texttt{r8..r15}):} Any 8-bit source with \texttt{mode=2} reads zero-latency uniform \texttt{u0..u31} from \texttt{FAU\_k} (\texttt{slots64 = bits[63:56]}); \texttt{mode=3} reads \texttt{SmallConstant} (Table~1.2) or \texttt{FauSpecial} (\texttt{pc}, \texttt{lane\_id}, \texttt{lane\_width}, \texttt{tls\_ptr}, \texttt{wls\_ptr}, \texttt{blend\_descriptor\_0..7}).
\end{tightitemize}
\par\textbf{Worked Examples:}
\par\textbf{Example 1 (Standard fused shift + bitwise op):}
\codebox{\mnem{LROT\_AND.i32} \dst{r0}, \src{r1}, \src{r2}\imod{.b0}, \src{r3}}
With \texttt{r1 = 0x80000001}, \texttt{r2.b0 = 1} (masked mod 32): shift yields \texttt{0x00000003}, then \texttt{\& 0x0FFFFFFF} yields \texttt{r0 = 0x3}.
\par\textbf{Example 2 (Modulo-32 shift count \& .not\_result):}
\codebox{\mnem{LROT\_AND.i32}\imod{.not\_result} \dst{r0}, \src{r1}, \src{r2}\imod{.b0}, \src{r3}}
With \texttt{r1 = 0x12345678}, \texttt{r2.b0 = 8} (effective shift \texttt{8}): shift yields \texttt{0x34567812}, \texttt{\& 0x0FFFFFFF} yields \texttt{0x04567812}, and \texttt{.not\_result} inverts -\textgreater{} \texttt{r0 = 0xFBA987ED}.

\begin{center}\footnotesize
\begin{tabularx}{\linewidth}{@{}lllX@{}}
\toprule
Variant & Arch & Class & Opcode Fields \\
\midrule
\texttt{(none)} & v9--v14 & \texttt{R3} & \texttt{opcode1=SHIFT\_BITWISE\_I32} \\
\texttt{(none)} & v9--v14 & \texttt{R3} & \texttt{opcode1=LOGIC\_SHIFT\_I32} \\
\texttt{(none)} & v9--v14 & \texttt{D3} & \texttt{opcode1=SHIFT\_BITWISE\_I64} \\
\texttt{(none)} & v9--v14 & \texttt{D3} & \texttt{opcode1=LOGIC\_SHIFT\_I64} \\
\texttt{(none)} & v9--v14 & \texttt{R3} & \texttt{opcode1=SHIFT\_BITWISE\_V2I16} \\
\texttt{(none)} & v9--v14 & \texttt{R3} & \texttt{opcode1=LOGIC\_SHIFT\_V2I16} \\
\texttt{(none)} & v9--v14 & \texttt{R3} & \texttt{opcode1=SHIFT\_BITWISE\_V4I8} \\
\texttt{(none)} & v9--v14 & \texttt{R3} & \texttt{opcode1=LOGIC\_SHIFT\_V4I8} \\
\bottomrule
\end{tabularx}
\end{center}
\begin{center}\footnotesize
\begin{tabularx}{0.88\linewidth}{@{}l l X@{}}
\toprule
Field Name & Bits & Type / Fixed Value \\
\midrule
\texttt{src0} & \texttt{0:7} & \texttt{SourceEncoding} \\
\texttt{src1} & \texttt{8:15} & \texttt{SourceEncoding} \\
\texttt{src2} & \texttt{16:23} & \texttt{SourceEncoding} \\
\texttt{bitwise\_operation} & \texttt{24:25} & \texttt{exact=and} \\
\hyperref[enum:shift_bitwise_i32_lanes_m]{\texttt{lanes1}} & \texttt{26:29} & \texttt{shift\_bitwise\_i32\_lanes\_m (def b0)} \\
\hyperref[enum:not_m]{\texttt{not\_result}} & \texttt{30} & \texttt{not\_m (def none)} \\
\textit{Reserved} & \texttt{31} & \texttt{zero} \\
\texttt{shift\_operation} & \texttt{32:34} & \texttt{exact=lrot} \\
\hyperref[enum:not_m]{\texttt{not2}} & \texttt{35} & \texttt{not\_m (def none)} \\
\hyperref[enum:widen_int_m]{\texttt{widen0}} & \texttt{36:39} & \texttt{widen\_int\_m (def none)} \\
\texttt{dest\_register\_index} & \texttt{40:45} & \texttt{uint} \\
\hyperref[enum:dest_width_m]{\texttt{dest\_width}} & \texttt{46:47} & \texttt{dest\_width\_m} \\
\texttt{opcode1} & \texttt{48:51} & \texttt{exact=SHIFT\_BITWISE\_I32} \\
\texttt{instruction\_class} & \texttt{52:56} & \texttt{exact=R3} \\
\texttt{fau\_page\_index} & \texttt{57:58} & \texttt{uint} \\
\hyperref[enum:flow_control_m]{\texttt{flow}} & \texttt{59:62} & \texttt{flow\_control\_m} \\
\textit{Reserved} & \texttt{63} & \texttt{zero} \\
\bottomrule
\end{tabularx}
\end{center}
\textbf{Fields Explained:}
\begin{tightitemize}\footnotesize
  \item \texttt{src0}: Primary source operand (\texttt{r0..r63}, \texttt{u0..u63} uniform, \texttt{k*} FAU constant, or \texttt{t0..t2} temporary; \texttt{\textasciicircum{}} marks discard/last-use). \emph{Example:} {\scriptsize\ttfamily \mnem{FADD.f32} \dst{r0}, \src{r1}, \src{r2}} (\textbf{SPIR-V:} \texttt{OpFAdd (Operand 0)}; \textbf{GLSL:} \texttt{float c = a + b;})
  \item \texttt{src1}: Secondary source operand (register, uniform, or FAU constant; in shifts, supplies the shift count byte). \emph{Example:} {\scriptsize\ttfamily \mnem{LSHIFT.i32} \dst{r0}, \src{r1}, \src{r2}\imod{.b0}} (\textbf{SPIR-V:} \texttt{OpShiftLeftLogical (Shift)}; \textbf{GLSL:} \texttt{uint y = a \textless{}\textless{} s;})
  \item \texttt{src2}: Third source operand (FMA addend, \texttt{MUX}/\texttt{CSEL} selector, \texttt{IADDC} carry-in, or bitwise mask in fused shift-logic). \emph{Example:} {\scriptsize\ttfamily \mnem{FMA.f32} \dst{r0}, \src{r1}, \src{r2}, \src{r3}} (\textbf{SPIR-V:} \texttt{GLSL.std.450 Fma (Addend)}; \textbf{GLSL:} \texttt{float d = fma(a, b, c);})
  \item \texttt{bitwise\_operation}: Selects the bitwise ALU op (\texttt{and}, \texttt{or}, \texttt{xor}, \texttt{lshift\_and}, etc.) fused with shifts. \emph{Example:} {\scriptsize\ttfamily \mnem{LSHIFT\_OR.i32} \dst{r0}, \src{r1}, \src{r2}, \src{r3}} (\textbf{SPIR-V:} \texttt{OpShiftLeftLogical + OpBitwiseOr}; \textbf{GLSL:} \texttt{uint y = (a \textless{}\textless{} s) \textbar{} b;})
  \item \hyperref[enum:shift_bitwise_i32_lanes_m]{\texttt{lanes1}}: Stage A byte-lane selector on \texttt{src1} (selects shift-count byte \texttt{b0..b3} on 32/64-bit shifts, \texttt{bXY} on \texttt{v2i16} shifts, or 4-byte swizzle on \texttt{v4*8}). \emph{Example:} {\scriptsize\ttfamily \mnem{LSHIFT.i32} \dst{r0}, \src{r1}, \src{r2}\imod{.b1}} (\textbf{SPIR-V:} \texttt{OpBitFieldUExtract + OpShiftLeftLogical}; \textbf{GLSL:} \texttt{uint y = a \textless{}\textless{} ((s \textgreater{}\textgreater{} 8) \& 0xFFu);})
  \item \hyperref[enum:not_m]{\texttt{not\_result}}: Bitwise-inverts the final output of a bitwise/shift-bitwise instruction (\textasciitilde{}result). \emph{Example:} {\scriptsize\ttfamily \mnem{AND.i32}\imod{.not\_result} \dst{r0}, \src{r1}, \src{r2}} (\textbf{SPIR-V:} \texttt{OpBitwiseAnd + OpNot}; \textbf{GLSL:} \texttt{uint y = \textasciitilde{}(a \& b);})
  \item \texttt{shift\_operation}: Selects shift direction/type (\texttt{lshift}, \texttt{rshift}, \texttt{arshift}) in fused shift-bitwise instructions. \emph{Example:} {\scriptsize\ttfamily \mnem{RSHIFT\_AND.i32} \dst{r0}, \src{r1}, \src{r2}, \src{r3}} (\textbf{SPIR-V:} \texttt{OpShiftRightLogical + OpBitwiseAnd}; \textbf{GLSL:} \texttt{uint y = (a \textgreater{}\textgreater{} s) \& m;})
  \item \hyperref[enum:not_m]{\texttt{not2}}: Stage A bitwise complement (\texttt{\textasciitilde{}src2}) applied to \texttt{src2} before the bitwise/shift-logic operation. \emph{Example:} {\scriptsize\ttfamily \mnem{AND.i32} \dst{r0}, \src{r1}, \src{r2}\imod{.not}} (\textbf{SPIR-V:} \texttt{OpNot + OpBitwiseAnd}; \textbf{GLSL:} \texttt{uint y = a \& \textasciitilde{}b;})
  \item \hyperref[enum:widen_int_m]{\texttt{widen0}}: Stage A widening extractor on \texttt{src0} (promotes FP16 \texttt{h0}/\texttt{h1} to FP32, or sign/zero-extends \texttt{b0..b3}/\texttt{h0..h1}/\texttt{w0..w1} to 32/64-bit before \texttt{abs}/\texttt{neg}). \emph{Example:} {\scriptsize\ttfamily \mnem{IADD.s32} \dst{r0}, \src{r1}\imod{.h0}, \src{r2}} (\textbf{SPIR-V:} \texttt{OpSConvert / OpFConvert}; \textbf{GLSL:} \texttt{int c = int(int16\_t(a)) + b;})
  \item \texttt{dest\_register\_index}: 6-bit destination GPR index (\texttt{r0..r63}) receiving the Stage C writeback result (even-aligned pair \texttt{[rN:rN+1]} for 64-bit ops). \emph{Example:} {\scriptsize\ttfamily \mnem{FMA.f32} \dst{r0}, \src{r1}, \src{r2}, \src{r3}} (\textbf{SPIR-V:} \texttt{SSA Result ID}; \textbf{GLSL:} \texttt{float r0 = fma(r1, r2, r3);})
  \item \hyperref[enum:dest_width_m]{\texttt{dest\_width}}: Stage C destination lane write-mask (\texttt{none} = full 32-bit, \texttt{h0} = low 16-bit, \texttt{h1} = high 16-bit without bit-shifting). \emph{Example:} {\scriptsize\ttfamily \mnem{FADD.v2f16} \dst{r0}\imod{.h1}, \src{r1}, \src{r2}} (\textbf{SPIR-V:} \texttt{OpVectorInsertDynamic / WriteMask}; \textbf{GLSL:} \texttt{dst.y = a.y + b.y;})
  \item \texttt{opcode1}: Primary sub-opcode within the instruction class uniquely identifying the operation family. \emph{Example:} {\scriptsize\ttfamily \mnem{IADD.s32} \dst{r0}, \src{r1}, \src{r2}} (\textbf{SPIR-V:} \texttt{OpIAdd (operation selector)}; \textbf{GLSL:} \texttt{// Primary opcode bits})
  \item \texttt{instruction\_class}: 4-bit major instruction class selector (\texttt{bits 48:51}) partitioning the 64-bit Valhall opcode space (e.g., \texttt{FMA}, \texttt{CVT}, \texttt{SFU}, \texttt{LS}, \texttt{TEX}, \texttt{BRANCH}). \emph{Example:} {\scriptsize\ttfamily \mnem{FADD.f32} \dst{r0}, \src{r1}, \src{r2}} (\textbf{SPIR-V:} \texttt{N/A (encoding class)}; \textbf{GLSL:} \texttt{// Hardware execution pipe / decode class})
  \item \texttt{fau\_page\_index}: Selects the FAU (Fast Access Uniform) page (\texttt{0..3}) when a source operand references hardware small constants, thread IDs, or lane constants. \emph{Example:} {\scriptsize\ttfamily \mnem{FADD.f32} \dst{r0}, \src{r1}, \src{k1f}} (\textbf{SPIR-V:} \texttt{OpConstant / BuiltIn}; \textbf{GLSL:} \texttt{float y = x + 1.0;})
  \item \hyperref[enum:flow_control_m]{\texttt{flow}}: 4-bit embedded flow-control \& scoreboard wait modifier (\texttt{bits 59:62}: \texttt{none}, \texttt{wait0..wait0126}, \texttt{wait}, \texttt{reconverge}, \texttt{discard}, \texttt{end}). \emph{Example:} {\scriptsize\ttfamily \mnem{FADD.f32}\imod{.wait0.reconverge} \dst{r0}, \src{r1}, \src{r2}} (\textbf{SPIR-V:} \texttt{OpControlBarrier / OpReturn}; \textbf{GLSL:} \texttt{// Scoreboard wait + warp reconvergence})
\end{tightitemize}
\clearpage

\subsection*{LROT\_OR}
\addcontentsline{toc}{subsection}{LROT\_OR}
For each SIMD element $i \in \{0, \dots, N-1\}$ ($N=1$ for \texttt{i32}/\texttt{i64}, $N=2$ for \texttt{v2i16}, $N=4$ for \texttt{v4i8} on v9--v10, lane width $W$), computes rotate left \texttt{s = rotate\_left(val, amt \& (W - 1))} where \texttt{val = src0[i]} and the 8-bit shift count \texttt{amt} is extracted from \texttt{src1} via \texttt{lanes1}, then combines with \texttt{src2[i]} via bitwise OR \texttt{res = s \textbar{} (not2 ? \textasciitilde{}src2[i] : src2[i])} and writes \texttt{not\_result ? \textasciitilde{}res : res} to \texttt{dest[i]}.
\par\textbf{Notes:}
\begin{tightitemize}
  \item Shift amount is masked modulo lane width (\texttt{wrapping\_shl}/\texttt{wrapping\_shr}/\texttt{rotate\_*}, \texttt{kraid/ops.rs:3984-4009}).
  \item On \texttt{v2i16}, \texttt{src1} is \texttt{V2I8} (\texttt{ops.rs:3953}) with byte lanes selected by \texttt{lanes1} (default \texttt{b02}: byte 0 -\textgreater{} lane 0, byte 2 -\textgreater{} lane 1).
  \item \textbf{Source references:} \texttt{kraid/ops.rs:3878-4025}, \texttt{encode\_v9.rs:2940-3093}, \texttt{hw\_tests.rs:2148-2198}.
\end{tightitemize}
\codebox{\mnem{LROT\_OR}\imod{\{.flow\}}\imod{\{.not\_result\}} \dst{dest}\imod{\{.dest\_width\}}, \src{src0}\imod{\{.widen0\}}, \src{src1}\imod{\{.lanes1\}}, \src{src2}\imod{\{.not2\}}\\\mnem{LROT\_OR}\imod{\{.flow\}}\imod{\{.not\_result\}} \dst{dest}, \src{src0}\imod{\{.widen0\}}, \src{src1}\imod{\{.lanes1\}}, \src{src2}\imod{\{.not2\}}\\\mnem{LROT\_OR}\imod{\{.flow\}}\imod{\{.not\_result\}} \dst{dest}\imod{\{.dest\_width\}}, \src{src0}\imod{\{.swz0\}}, \src{src1}\imod{\{.lanes1\}}, \src{src2}\imod{\{.not2\}}\\\mnem{LROT\_OR}\imod{\{.flow\}}\imod{\{.not\_result\}} \dst{dest}\imod{\{.dest\_width\}}, \src{src0}\imod{\{.lanes0\}}, \src{src1}\imod{\{.lanes1\}}, \src{src2}\imod{\{.not2\}}}
\textbf{Operands:}
\begin{tightitemize}\footnotesize
  \item \texttt{\dst{dest}} (\texttt{dest\_register\_index}, \texttt{dest\_width}): Destination GPR receiving the rotate left combined via bitwise OR (inverted if \texttt{.not\_result}).
  \item \texttt{\src{src0}}: Input integer value to shift/rotate.
  \item \texttt{\src{src1}}: 8-bit shift/rotate count per lane (selected via \texttt{lanes1}, masked modulo lane width \texttt{W}).
  \item \texttt{\src{src2}}: Second operand combined via bitwise OR after shifting (inverted if \texttt{.not2}).
\end{tightitemize}
\par\textbf{Closest SPIR-V Op:} \texttt{None (lowered from `(x\textless{}\textless{}s)\textbar{}(x\textgreater{}\textgreater{}(32-s))`) + OpBitwiseOr}
\par\textbf{Pseudo-GLSL (Operational Semantics):}
\codebox{for (int i = 0; i \textless{} N; ++i) \{\\    uint s = extractByte(src1, lanes1, i) \& (W - 1);\\    auto r = ((src0[i] \textless{}\textless{} s) \textbar{} (src0[i] \textgreater{}\textgreater{} ((W - s) \& (W - 1)))) \textbar{} (flag\_not2 ? \textasciitilde{}src2[i] : src2[i]);\\    dest[i] = flag\_not\_result ? \textasciitilde{}r : r;\\\}}
\textbf{Execution Context:}
\begin{tightitemize}\footnotesize
  \item \textbf{Compute stage preloads (\hyperref[sec:csf-spd]{\texttt{SPD\_k}}, \texttt{r16..r23}):} Hardware pre-fills \texttt{r55..r57} with \texttt{WorkGroupID.xyz + JOB\_OFFSET} (from \hyperref[sec:csf-regs]{\texttt{COMPUTESR::JOB\_OFFSET\_X/Y/Z}} \texttt{r34..r36}) and \texttt{r60..r62} with \texttt{LocalInvocationID.xyz} or \texttt{GlobalInvocationID.xyz} (from \texttt{WG\_SIZE} \texttt{r33} \& \texttt{JOB\_SIZE} \texttt{r37..r39}, selected by \texttt{ShaderProgram::preload}) so compute shaders can derive thread/element indices without memory loads.
  \item \textbf{Vertex stage preloads (\hyperref[sec:csf-spd]{\texttt{SPD\_0}/\texttt{SPD\_1}}, \texttt{r16..r19}):} The hardware Index Fetcher reads \hyperref[sec:csf-regs]{\texttt{IDVSSR::INDEX\_BUFFER}} (\texttt{r54:r55}) and pre-fills \texttt{r59, r60} with \texttt{VertexIndex} (\texttt{+ VERTEX\_OFFSET + GLOBAL\_ATTR\_OFF}) and \texttt{r61} with \texttt{InstanceIndex} (\texttt{+ INSTANCE\_OFFSET}) so Position and Varying shaders can pull vertex attributes directly.
  \item \textbf{Inline FAU \& Special Operands (\hyperref[sec:csf-tsd-fau]{\texttt{FAU\_k}}, \texttt{r8..r15}):} Any 8-bit source with \texttt{mode=2} reads zero-latency uniform \texttt{u0..u31} from \texttt{FAU\_k} (\texttt{slots64 = bits[63:56]}); \texttt{mode=3} reads \texttt{SmallConstant} (Table~1.2) or \texttt{FauSpecial} (\texttt{pc}, \texttt{lane\_id}, \texttt{lane\_width}, \texttt{tls\_ptr}, \texttt{wls\_ptr}, \texttt{blend\_descriptor\_0..7}).
\end{tightitemize}
\par\textbf{Worked Examples:}
\par\textbf{Example 1 (Standard fused shift + bitwise op):}
\codebox{\mnem{LROT\_OR.i32} \dst{r0}, \src{r1}, \src{r2}\imod{.b0}, \src{r3}}
With \texttt{r1 = 0x80000001}, \texttt{r2.b0 = 1} (masked mod 32): shift yields \texttt{0x00000003}, then \texttt{\textbar{} 0x0000000F} yields \texttt{r0 = 0xF}.
\par\textbf{Example 2 (Modulo-32 shift count \& .not\_result):}
\codebox{\mnem{LROT\_OR.i32}\imod{.not\_result} \dst{r0}, \src{r1}, \src{r2}\imod{.b0}, \src{r3}}
With \texttt{r1 = 0x12345678}, \texttt{r2.b0 = 8} (effective shift \texttt{8}): shift yields \texttt{0x34567812}, \texttt{\textbar{} 0x0000000F} yields \texttt{0x3456781F}, and \texttt{.not\_result} inverts -\textgreater{} \texttt{r0 = 0xCBA987E0}.

\begin{center}\footnotesize
\begin{tabularx}{\linewidth}{@{}lllX@{}}
\toprule
Variant & Arch & Class & Opcode Fields \\
\midrule
\texttt{(none)} & v9--v14 & \texttt{R3} & \texttt{opcode1=SHIFT\_BITWISE\_I32} \\
\texttt{(none)} & v9--v14 & \texttt{R3} & \texttt{opcode1=LOGIC\_SHIFT\_I32} \\
\texttt{(none)} & v9--v14 & \texttt{D3} & \texttt{opcode1=SHIFT\_BITWISE\_I64} \\
\texttt{(none)} & v9--v14 & \texttt{D3} & \texttt{opcode1=LOGIC\_SHIFT\_I64} \\
\texttt{(none)} & v9--v14 & \texttt{R3} & \texttt{opcode1=SHIFT\_BITWISE\_V2I16} \\
\texttt{(none)} & v9--v14 & \texttt{R3} & \texttt{opcode1=LOGIC\_SHIFT\_V2I16} \\
\texttt{(none)} & v9--v14 & \texttt{R3} & \texttt{opcode1=SHIFT\_BITWISE\_V4I8} \\
\texttt{(none)} & v9--v14 & \texttt{R3} & \texttt{opcode1=LOGIC\_SHIFT\_V4I8} \\
\bottomrule
\end{tabularx}
\end{center}
\begin{center}\footnotesize
\begin{tabularx}{0.88\linewidth}{@{}l l X@{}}
\toprule
Field Name & Bits & Type / Fixed Value \\
\midrule
\texttt{src0} & \texttt{0:7} & \texttt{SourceEncoding} \\
\texttt{src1} & \texttt{8:15} & \texttt{SourceEncoding} \\
\texttt{src2} & \texttt{16:23} & \texttt{SourceEncoding} \\
\texttt{bitwise\_operation} & \texttt{24:25} & \texttt{exact=or} \\
\hyperref[enum:shift_bitwise_i32_lanes_m]{\texttt{lanes1}} & \texttt{26:29} & \texttt{shift\_bitwise\_i32\_lanes\_m (def b0)} \\
\hyperref[enum:not_m]{\texttt{not\_result}} & \texttt{30} & \texttt{not\_m (def none)} \\
\textit{Reserved} & \texttt{31} & \texttt{zero} \\
\texttt{shift\_operation} & \texttt{32:34} & \texttt{exact=lrot} \\
\hyperref[enum:not_m]{\texttt{not2}} & \texttt{35} & \texttt{not\_m (def none)} \\
\hyperref[enum:widen_int_m]{\texttt{widen0}} & \texttt{36:39} & \texttt{widen\_int\_m (def none)} \\
\texttt{dest\_register\_index} & \texttt{40:45} & \texttt{uint} \\
\hyperref[enum:dest_width_m]{\texttt{dest\_width}} & \texttt{46:47} & \texttt{dest\_width\_m} \\
\texttt{opcode1} & \texttt{48:51} & \texttt{exact=SHIFT\_BITWISE\_I32} \\
\texttt{instruction\_class} & \texttt{52:56} & \texttt{exact=R3} \\
\texttt{fau\_page\_index} & \texttt{57:58} & \texttt{uint} \\
\hyperref[enum:flow_control_m]{\texttt{flow}} & \texttt{59:62} & \texttt{flow\_control\_m} \\
\textit{Reserved} & \texttt{63} & \texttt{zero} \\
\bottomrule
\end{tabularx}
\end{center}
\textbf{Fields Explained:}
\begin{tightitemize}\footnotesize
  \item \texttt{src0}: Primary source operand (\texttt{r0..r63}, \texttt{u0..u63} uniform, \texttt{k*} FAU constant, or \texttt{t0..t2} temporary; \texttt{\textasciicircum{}} marks discard/last-use). \emph{Example:} {\scriptsize\ttfamily \mnem{FADD.f32} \dst{r0}, \src{r1}, \src{r2}} (\textbf{SPIR-V:} \texttt{OpFAdd (Operand 0)}; \textbf{GLSL:} \texttt{float c = a + b;})
  \item \texttt{src1}: Secondary source operand (register, uniform, or FAU constant; in shifts, supplies the shift count byte). \emph{Example:} {\scriptsize\ttfamily \mnem{LSHIFT.i32} \dst{r0}, \src{r1}, \src{r2}\imod{.b0}} (\textbf{SPIR-V:} \texttt{OpShiftLeftLogical (Shift)}; \textbf{GLSL:} \texttt{uint y = a \textless{}\textless{} s;})
  \item \texttt{src2}: Third source operand (FMA addend, \texttt{MUX}/\texttt{CSEL} selector, \texttt{IADDC} carry-in, or bitwise mask in fused shift-logic). \emph{Example:} {\scriptsize\ttfamily \mnem{FMA.f32} \dst{r0}, \src{r1}, \src{r2}, \src{r3}} (\textbf{SPIR-V:} \texttt{GLSL.std.450 Fma (Addend)}; \textbf{GLSL:} \texttt{float d = fma(a, b, c);})
  \item \texttt{bitwise\_operation}: Selects the bitwise ALU op (\texttt{and}, \texttt{or}, \texttt{xor}, \texttt{lshift\_and}, etc.) fused with shifts. \emph{Example:} {\scriptsize\ttfamily \mnem{LSHIFT\_OR.i32} \dst{r0}, \src{r1}, \src{r2}, \src{r3}} (\textbf{SPIR-V:} \texttt{OpShiftLeftLogical + OpBitwiseOr}; \textbf{GLSL:} \texttt{uint y = (a \textless{}\textless{} s) \textbar{} b;})
  \item \hyperref[enum:shift_bitwise_i32_lanes_m]{\texttt{lanes1}}: Stage A byte-lane selector on \texttt{src1} (selects shift-count byte \texttt{b0..b3} on 32/64-bit shifts, \texttt{bXY} on \texttt{v2i16} shifts, or 4-byte swizzle on \texttt{v4*8}). \emph{Example:} {\scriptsize\ttfamily \mnem{LSHIFT.i32} \dst{r0}, \src{r1}, \src{r2}\imod{.b1}} (\textbf{SPIR-V:} \texttt{OpBitFieldUExtract + OpShiftLeftLogical}; \textbf{GLSL:} \texttt{uint y = a \textless{}\textless{} ((s \textgreater{}\textgreater{} 8) \& 0xFFu);})
  \item \hyperref[enum:not_m]{\texttt{not\_result}}: Bitwise-inverts the final output of a bitwise/shift-bitwise instruction (\textasciitilde{}result). \emph{Example:} {\scriptsize\ttfamily \mnem{AND.i32}\imod{.not\_result} \dst{r0}, \src{r1}, \src{r2}} (\textbf{SPIR-V:} \texttt{OpBitwiseAnd + OpNot}; \textbf{GLSL:} \texttt{uint y = \textasciitilde{}(a \& b);})
  \item \texttt{shift\_operation}: Selects shift direction/type (\texttt{lshift}, \texttt{rshift}, \texttt{arshift}) in fused shift-bitwise instructions. \emph{Example:} {\scriptsize\ttfamily \mnem{RSHIFT\_AND.i32} \dst{r0}, \src{r1}, \src{r2}, \src{r3}} (\textbf{SPIR-V:} \texttt{OpShiftRightLogical + OpBitwiseAnd}; \textbf{GLSL:} \texttt{uint y = (a \textgreater{}\textgreater{} s) \& m;})
  \item \hyperref[enum:not_m]{\texttt{not2}}: Stage A bitwise complement (\texttt{\textasciitilde{}src2}) applied to \texttt{src2} before the bitwise/shift-logic operation. \emph{Example:} {\scriptsize\ttfamily \mnem{AND.i32} \dst{r0}, \src{r1}, \src{r2}\imod{.not}} (\textbf{SPIR-V:} \texttt{OpNot + OpBitwiseAnd}; \textbf{GLSL:} \texttt{uint y = a \& \textasciitilde{}b;})
  \item \hyperref[enum:widen_int_m]{\texttt{widen0}}: Stage A widening extractor on \texttt{src0} (promotes FP16 \texttt{h0}/\texttt{h1} to FP32, or sign/zero-extends \texttt{b0..b3}/\texttt{h0..h1}/\texttt{w0..w1} to 32/64-bit before \texttt{abs}/\texttt{neg}). \emph{Example:} {\scriptsize\ttfamily \mnem{IADD.s32} \dst{r0}, \src{r1}\imod{.h0}, \src{r2}} (\textbf{SPIR-V:} \texttt{OpSConvert / OpFConvert}; \textbf{GLSL:} \texttt{int c = int(int16\_t(a)) + b;})
  \item \texttt{dest\_register\_index}: 6-bit destination GPR index (\texttt{r0..r63}) receiving the Stage C writeback result (even-aligned pair \texttt{[rN:rN+1]} for 64-bit ops). \emph{Example:} {\scriptsize\ttfamily \mnem{FMA.f32} \dst{r0}, \src{r1}, \src{r2}, \src{r3}} (\textbf{SPIR-V:} \texttt{SSA Result ID}; \textbf{GLSL:} \texttt{float r0 = fma(r1, r2, r3);})
  \item \hyperref[enum:dest_width_m]{\texttt{dest\_width}}: Stage C destination lane write-mask (\texttt{none} = full 32-bit, \texttt{h0} = low 16-bit, \texttt{h1} = high 16-bit without bit-shifting). \emph{Example:} {\scriptsize\ttfamily \mnem{FADD.v2f16} \dst{r0}\imod{.h1}, \src{r1}, \src{r2}} (\textbf{SPIR-V:} \texttt{OpVectorInsertDynamic / WriteMask}; \textbf{GLSL:} \texttt{dst.y = a.y + b.y;})
  \item \texttt{opcode1}: Primary sub-opcode within the instruction class uniquely identifying the operation family. \emph{Example:} {\scriptsize\ttfamily \mnem{IADD.s32} \dst{r0}, \src{r1}, \src{r2}} (\textbf{SPIR-V:} \texttt{OpIAdd (operation selector)}; \textbf{GLSL:} \texttt{// Primary opcode bits})
  \item \texttt{instruction\_class}: 4-bit major instruction class selector (\texttt{bits 48:51}) partitioning the 64-bit Valhall opcode space (e.g., \texttt{FMA}, \texttt{CVT}, \texttt{SFU}, \texttt{LS}, \texttt{TEX}, \texttt{BRANCH}). \emph{Example:} {\scriptsize\ttfamily \mnem{FADD.f32} \dst{r0}, \src{r1}, \src{r2}} (\textbf{SPIR-V:} \texttt{N/A (encoding class)}; \textbf{GLSL:} \texttt{// Hardware execution pipe / decode class})
  \item \texttt{fau\_page\_index}: Selects the FAU (Fast Access Uniform) page (\texttt{0..3}) when a source operand references hardware small constants, thread IDs, or lane constants. \emph{Example:} {\scriptsize\ttfamily \mnem{FADD.f32} \dst{r0}, \src{r1}, \src{k1f}} (\textbf{SPIR-V:} \texttt{OpConstant / BuiltIn}; \textbf{GLSL:} \texttt{float y = x + 1.0;})
  \item \hyperref[enum:flow_control_m]{\texttt{flow}}: 4-bit embedded flow-control \& scoreboard wait modifier (\texttt{bits 59:62}: \texttt{none}, \texttt{wait0..wait0126}, \texttt{wait}, \texttt{reconverge}, \texttt{discard}, \texttt{end}). \emph{Example:} {\scriptsize\ttfamily \mnem{FADD.f32}\imod{.wait0.reconverge} \dst{r0}, \src{r1}, \src{r2}} (\textbf{SPIR-V:} \texttt{OpControlBarrier / OpReturn}; \textbf{GLSL:} \texttt{// Scoreboard wait + warp reconvergence})
\end{tightitemize}
\clearpage

\subsection*{LROT\_XOR}
\addcontentsline{toc}{subsection}{LROT\_XOR}
For each SIMD element $i \in \{0, \dots, N-1\}$ ($N=1$ for \texttt{i32}/\texttt{i64}, $N=2$ for \texttt{v2i16}, $N=4$ for \texttt{v4i8} on v9--v10, lane width $W$), computes rotate left \texttt{s = rotate\_left(val, amt \& (W - 1))} where \texttt{val = src0[i]} and the 8-bit shift count \texttt{amt} is extracted from \texttt{src1} via \texttt{lanes1}, then combines with \texttt{src2[i]} via bitwise XOR \texttt{res = s \textasciicircum{} (not2 ? \textasciitilde{}src2[i] : src2[i])} and writes \texttt{not\_result ? \textasciitilde{}res : res} to \texttt{dest[i]}.
\par\textbf{Notes:}
\begin{tightitemize}
  \item Shift amount is masked modulo lane width (\texttt{wrapping\_shl}/\texttt{wrapping\_shr}/\texttt{rotate\_*}, \texttt{kraid/ops.rs:3984-4009}).
  \item On \texttt{v2i16}, \texttt{src1} is \texttt{V2I8} (\texttt{ops.rs:3953}) with byte lanes selected by \texttt{lanes1} (default \texttt{b02}: byte 0 -\textgreater{} lane 0, byte 2 -\textgreater{} lane 1).
  \item \textbf{Source references:} \texttt{kraid/ops.rs:3878-4025}, \texttt{encode\_v9.rs:2940-3093}, \texttt{hw\_tests.rs:2148-2198}.
\end{tightitemize}
\codebox{\mnem{LROT\_XOR}\imod{\{.flow\}}\imod{\{.not\_result\}} \dst{dest}\imod{\{.dest\_width\}}, \src{src0}\imod{\{.widen0\}}, \src{src1}\imod{\{.lanes1\}}, \src{src2}\imod{\{.not2\}}\\\mnem{LROT\_XOR}\imod{\{.flow\}}\imod{\{.not\_result\}} \dst{dest}, \src{src0}\imod{\{.widen0\}}, \src{src1}\imod{\{.lanes1\}}, \src{src2}\imod{\{.not2\}}\\\mnem{LROT\_XOR}\imod{\{.flow\}}\imod{\{.not\_result\}} \dst{dest}\imod{\{.dest\_width\}}, \src{src0}\imod{\{.swz0\}}, \src{src1}\imod{\{.lanes1\}}, \src{src2}\imod{\{.not2\}}\\\mnem{LROT\_XOR}\imod{\{.flow\}}\imod{\{.not\_result\}} \dst{dest}\imod{\{.dest\_width\}}, \src{src0}\imod{\{.lanes0\}}, \src{src1}\imod{\{.lanes1\}}, \src{src2}\imod{\{.not2\}}}
\textbf{Operands:}
\begin{tightitemize}\footnotesize
  \item \texttt{\dst{dest}} (\texttt{dest\_register\_index}, \texttt{dest\_width}): Destination GPR receiving the rotate left combined via bitwise XOR (inverted if \texttt{.not\_result}).
  \item \texttt{\src{src0}}: Input integer value to shift/rotate.
  \item \texttt{\src{src1}}: 8-bit shift/rotate count per lane (selected via \texttt{lanes1}, masked modulo lane width \texttt{W}).
  \item \texttt{\src{src2}}: Second operand combined via bitwise XOR after shifting (inverted if \texttt{.not2}).
\end{tightitemize}
\par\textbf{Closest SPIR-V Op:} \texttt{None (lowered from `(x\textless{}\textless{}s)\textbar{}(x\textgreater{}\textgreater{}(32-s))`) + OpBitwiseXor}
\par\textbf{Pseudo-GLSL (Operational Semantics):}
\codebox{for (int i = 0; i \textless{} N; ++i) \{\\    uint s = extractByte(src1, lanes1, i) \& (W - 1);\\    auto r = ((src0[i] \textless{}\textless{} s) \textbar{} (src0[i] \textgreater{}\textgreater{} ((W - s) \& (W - 1)))) \textasciicircum{} (flag\_not2 ? \textasciitilde{}src2[i] : src2[i]);\\    dest[i] = flag\_not\_result ? \textasciitilde{}r : r;\\\}}
\textbf{Execution Context:}
\begin{tightitemize}\footnotesize
  \item \textbf{Compute stage preloads (\hyperref[sec:csf-spd]{\texttt{SPD\_k}}, \texttt{r16..r23}):} Hardware pre-fills \texttt{r55..r57} with \texttt{WorkGroupID.xyz + JOB\_OFFSET} (from \hyperref[sec:csf-regs]{\texttt{COMPUTESR::JOB\_OFFSET\_X/Y/Z}} \texttt{r34..r36}) and \texttt{r60..r62} with \texttt{LocalInvocationID.xyz} or \texttt{GlobalInvocationID.xyz} (from \texttt{WG\_SIZE} \texttt{r33} \& \texttt{JOB\_SIZE} \texttt{r37..r39}, selected by \texttt{ShaderProgram::preload}) so compute shaders can derive thread/element indices without memory loads.
  \item \textbf{Vertex stage preloads (\hyperref[sec:csf-spd]{\texttt{SPD\_0}/\texttt{SPD\_1}}, \texttt{r16..r19}):} The hardware Index Fetcher reads \hyperref[sec:csf-regs]{\texttt{IDVSSR::INDEX\_BUFFER}} (\texttt{r54:r55}) and pre-fills \texttt{r59, r60} with \texttt{VertexIndex} (\texttt{+ VERTEX\_OFFSET + GLOBAL\_ATTR\_OFF}) and \texttt{r61} with \texttt{InstanceIndex} (\texttt{+ INSTANCE\_OFFSET}) so Position and Varying shaders can pull vertex attributes directly.
  \item \textbf{Inline FAU \& Special Operands (\hyperref[sec:csf-tsd-fau]{\texttt{FAU\_k}}, \texttt{r8..r15}):} Any 8-bit source with \texttt{mode=2} reads zero-latency uniform \texttt{u0..u31} from \texttt{FAU\_k} (\texttt{slots64 = bits[63:56]}); \texttt{mode=3} reads \texttt{SmallConstant} (Table~1.2) or \texttt{FauSpecial} (\texttt{pc}, \texttt{lane\_id}, \texttt{lane\_width}, \texttt{tls\_ptr}, \texttt{wls\_ptr}, \texttt{blend\_descriptor\_0..7}).
\end{tightitemize}
\par\textbf{Worked Examples:}
\par\textbf{Example 1 (Standard fused shift + bitwise op):}
\codebox{\mnem{LROT\_XOR.i32} \dst{r0}, \src{r1}, \src{r2}\imod{.b0}, \src{r3}}
With \texttt{r1 = 0x80000001}, \texttt{r2.b0 = 1} (masked mod 32): shift yields \texttt{0x00000003}, then \texttt{\textasciicircum{} 0x000000FF} yields \texttt{r0 = 0xFC}.
\par\textbf{Example 2 (Modulo-32 shift count \& .not\_result):}
\codebox{\mnem{LROT\_XOR.i32}\imod{.not\_result} \dst{r0}, \src{r1}, \src{r2}\imod{.b0}, \src{r3}}
With \texttt{r1 = 0x12345678}, \texttt{r2.b0 = 8} (effective shift \texttt{8}): shift yields \texttt{0x34567812}, \texttt{\textasciicircum{} 0x000000FF} yields \texttt{0x345678ED}, and \texttt{.not\_result} inverts -\textgreater{} \texttt{r0 = 0xCBA98712}.

\begin{center}\footnotesize
\begin{tabularx}{\linewidth}{@{}lllX@{}}
\toprule
Variant & Arch & Class & Opcode Fields \\
\midrule
\texttt{(none)} & v9--v14 & \texttt{R3} & \texttt{opcode1=SHIFT\_BITWISE\_I32} \\
\texttt{(none)} & v9--v14 & \texttt{R3} & \texttt{opcode1=LOGIC\_SHIFT\_I32} \\
\texttt{(none)} & v9--v14 & \texttt{D3} & \texttt{opcode1=SHIFT\_BITWISE\_I64} \\
\texttt{(none)} & v9--v14 & \texttt{D3} & \texttt{opcode1=LOGIC\_SHIFT\_I64} \\
\texttt{(none)} & v9--v14 & \texttt{R3} & \texttt{opcode1=SHIFT\_BITWISE\_V2I16} \\
\texttt{(none)} & v9--v14 & \texttt{R3} & \texttt{opcode1=LOGIC\_SHIFT\_V2I16} \\
\texttt{(none)} & v9--v14 & \texttt{R3} & \texttt{opcode1=SHIFT\_BITWISE\_V4I8} \\
\texttt{(none)} & v9--v14 & \texttt{R3} & \texttt{opcode1=LOGIC\_SHIFT\_V4I8} \\
\bottomrule
\end{tabularx}
\end{center}
\begin{center}\footnotesize
\begin{tabularx}{0.88\linewidth}{@{}l l X@{}}
\toprule
Field Name & Bits & Type / Fixed Value \\
\midrule
\texttt{src0} & \texttt{0:7} & \texttt{SourceEncoding} \\
\texttt{src1} & \texttt{8:15} & \texttt{SourceEncoding} \\
\texttt{src2} & \texttt{16:23} & \texttt{SourceEncoding} \\
\texttt{bitwise\_operation} & \texttt{24:25} & \texttt{exact=xor} \\
\hyperref[enum:shift_bitwise_i32_lanes_m]{\texttt{lanes1}} & \texttt{26:29} & \texttt{shift\_bitwise\_i32\_lanes\_m (def b0)} \\
\hyperref[enum:not_m]{\texttt{not\_result}} & \texttt{30} & \texttt{not\_m (def none)} \\
\textit{Reserved} & \texttt{31} & \texttt{zero} \\
\texttt{shift\_operation} & \texttt{32:34} & \texttt{exact=lrot} \\
\hyperref[enum:not_m]{\texttt{not2}} & \texttt{35} & \texttt{not\_m (def none)} \\
\hyperref[enum:widen_int_m]{\texttt{widen0}} & \texttt{36:39} & \texttt{widen\_int\_m (def none)} \\
\texttt{dest\_register\_index} & \texttt{40:45} & \texttt{uint} \\
\hyperref[enum:dest_width_m]{\texttt{dest\_width}} & \texttt{46:47} & \texttt{dest\_width\_m} \\
\texttt{opcode1} & \texttt{48:51} & \texttt{exact=SHIFT\_BITWISE\_I32} \\
\texttt{instruction\_class} & \texttt{52:56} & \texttt{exact=R3} \\
\texttt{fau\_page\_index} & \texttt{57:58} & \texttt{uint} \\
\hyperref[enum:flow_control_m]{\texttt{flow}} & \texttt{59:62} & \texttt{flow\_control\_m} \\
\textit{Reserved} & \texttt{63} & \texttt{zero} \\
\bottomrule
\end{tabularx}
\end{center}
\textbf{Fields Explained:}
\begin{tightitemize}\footnotesize
  \item \texttt{src0}: Primary source operand (\texttt{r0..r63}, \texttt{u0..u63} uniform, \texttt{k*} FAU constant, or \texttt{t0..t2} temporary; \texttt{\textasciicircum{}} marks discard/last-use). \emph{Example:} {\scriptsize\ttfamily \mnem{FADD.f32} \dst{r0}, \src{r1}, \src{r2}} (\textbf{SPIR-V:} \texttt{OpFAdd (Operand 0)}; \textbf{GLSL:} \texttt{float c = a + b;})
  \item \texttt{src1}: Secondary source operand (register, uniform, or FAU constant; in shifts, supplies the shift count byte). \emph{Example:} {\scriptsize\ttfamily \mnem{LSHIFT.i32} \dst{r0}, \src{r1}, \src{r2}\imod{.b0}} (\textbf{SPIR-V:} \texttt{OpShiftLeftLogical (Shift)}; \textbf{GLSL:} \texttt{uint y = a \textless{}\textless{} s;})
  \item \texttt{src2}: Third source operand (FMA addend, \texttt{MUX}/\texttt{CSEL} selector, \texttt{IADDC} carry-in, or bitwise mask in fused shift-logic). \emph{Example:} {\scriptsize\ttfamily \mnem{FMA.f32} \dst{r0}, \src{r1}, \src{r2}, \src{r3}} (\textbf{SPIR-V:} \texttt{GLSL.std.450 Fma (Addend)}; \textbf{GLSL:} \texttt{float d = fma(a, b, c);})
  \item \texttt{bitwise\_operation}: Selects the bitwise ALU op (\texttt{and}, \texttt{or}, \texttt{xor}, \texttt{lshift\_and}, etc.) fused with shifts. \emph{Example:} {\scriptsize\ttfamily \mnem{LSHIFT\_OR.i32} \dst{r0}, \src{r1}, \src{r2}, \src{r3}} (\textbf{SPIR-V:} \texttt{OpShiftLeftLogical + OpBitwiseOr}; \textbf{GLSL:} \texttt{uint y = (a \textless{}\textless{} s) \textbar{} b;})
  \item \hyperref[enum:shift_bitwise_i32_lanes_m]{\texttt{lanes1}}: Stage A byte-lane selector on \texttt{src1} (selects shift-count byte \texttt{b0..b3} on 32/64-bit shifts, \texttt{bXY} on \texttt{v2i16} shifts, or 4-byte swizzle on \texttt{v4*8}). \emph{Example:} {\scriptsize\ttfamily \mnem{LSHIFT.i32} \dst{r0}, \src{r1}, \src{r2}\imod{.b1}} (\textbf{SPIR-V:} \texttt{OpBitFieldUExtract + OpShiftLeftLogical}; \textbf{GLSL:} \texttt{uint y = a \textless{}\textless{} ((s \textgreater{}\textgreater{} 8) \& 0xFFu);})
  \item \hyperref[enum:not_m]{\texttt{not\_result}}: Bitwise-inverts the final output of a bitwise/shift-bitwise instruction (\textasciitilde{}result). \emph{Example:} {\scriptsize\ttfamily \mnem{AND.i32}\imod{.not\_result} \dst{r0}, \src{r1}, \src{r2}} (\textbf{SPIR-V:} \texttt{OpBitwiseAnd + OpNot}; \textbf{GLSL:} \texttt{uint y = \textasciitilde{}(a \& b);})
  \item \texttt{shift\_operation}: Selects shift direction/type (\texttt{lshift}, \texttt{rshift}, \texttt{arshift}) in fused shift-bitwise instructions. \emph{Example:} {\scriptsize\ttfamily \mnem{RSHIFT\_AND.i32} \dst{r0}, \src{r1}, \src{r2}, \src{r3}} (\textbf{SPIR-V:} \texttt{OpShiftRightLogical + OpBitwiseAnd}; \textbf{GLSL:} \texttt{uint y = (a \textgreater{}\textgreater{} s) \& m;})
  \item \hyperref[enum:not_m]{\texttt{not2}}: Stage A bitwise complement (\texttt{\textasciitilde{}src2}) applied to \texttt{src2} before the bitwise/shift-logic operation. \emph{Example:} {\scriptsize\ttfamily \mnem{AND.i32} \dst{r0}, \src{r1}, \src{r2}\imod{.not}} (\textbf{SPIR-V:} \texttt{OpNot + OpBitwiseAnd}; \textbf{GLSL:} \texttt{uint y = a \& \textasciitilde{}b;})
  \item \hyperref[enum:widen_int_m]{\texttt{widen0}}: Stage A widening extractor on \texttt{src0} (promotes FP16 \texttt{h0}/\texttt{h1} to FP32, or sign/zero-extends \texttt{b0..b3}/\texttt{h0..h1}/\texttt{w0..w1} to 32/64-bit before \texttt{abs}/\texttt{neg}). \emph{Example:} {\scriptsize\ttfamily \mnem{IADD.s32} \dst{r0}, \src{r1}\imod{.h0}, \src{r2}} (\textbf{SPIR-V:} \texttt{OpSConvert / OpFConvert}; \textbf{GLSL:} \texttt{int c = int(int16\_t(a)) + b;})
  \item \texttt{dest\_register\_index}: 6-bit destination GPR index (\texttt{r0..r63}) receiving the Stage C writeback result (even-aligned pair \texttt{[rN:rN+1]} for 64-bit ops). \emph{Example:} {\scriptsize\ttfamily \mnem{FMA.f32} \dst{r0}, \src{r1}, \src{r2}, \src{r3}} (\textbf{SPIR-V:} \texttt{SSA Result ID}; \textbf{GLSL:} \texttt{float r0 = fma(r1, r2, r3);})
  \item \hyperref[enum:dest_width_m]{\texttt{dest\_width}}: Stage C destination lane write-mask (\texttt{none} = full 32-bit, \texttt{h0} = low 16-bit, \texttt{h1} = high 16-bit without bit-shifting). \emph{Example:} {\scriptsize\ttfamily \mnem{FADD.v2f16} \dst{r0}\imod{.h1}, \src{r1}, \src{r2}} (\textbf{SPIR-V:} \texttt{OpVectorInsertDynamic / WriteMask}; \textbf{GLSL:} \texttt{dst.y = a.y + b.y;})
  \item \texttt{opcode1}: Primary sub-opcode within the instruction class uniquely identifying the operation family. \emph{Example:} {\scriptsize\ttfamily \mnem{IADD.s32} \dst{r0}, \src{r1}, \src{r2}} (\textbf{SPIR-V:} \texttt{OpIAdd (operation selector)}; \textbf{GLSL:} \texttt{// Primary opcode bits})
  \item \texttt{instruction\_class}: 4-bit major instruction class selector (\texttt{bits 48:51}) partitioning the 64-bit Valhall opcode space (e.g., \texttt{FMA}, \texttt{CVT}, \texttt{SFU}, \texttt{LS}, \texttt{TEX}, \texttt{BRANCH}). \emph{Example:} {\scriptsize\ttfamily \mnem{FADD.f32} \dst{r0}, \src{r1}, \src{r2}} (\textbf{SPIR-V:} \texttt{N/A (encoding class)}; \textbf{GLSL:} \texttt{// Hardware execution pipe / decode class})
  \item \texttt{fau\_page\_index}: Selects the FAU (Fast Access Uniform) page (\texttt{0..3}) when a source operand references hardware small constants, thread IDs, or lane constants. \emph{Example:} {\scriptsize\ttfamily \mnem{FADD.f32} \dst{r0}, \src{r1}, \src{k1f}} (\textbf{SPIR-V:} \texttt{OpConstant / BuiltIn}; \textbf{GLSL:} \texttt{float y = x + 1.0;})
  \item \hyperref[enum:flow_control_m]{\texttt{flow}}: 4-bit embedded flow-control \& scoreboard wait modifier (\texttt{bits 59:62}: \texttt{none}, \texttt{wait0..wait0126}, \texttt{wait}, \texttt{reconverge}, \texttt{discard}, \texttt{end}). \emph{Example:} {\scriptsize\ttfamily \mnem{FADD.f32}\imod{.wait0.reconverge} \dst{r0}, \src{r1}, \src{r2}} (\textbf{SPIR-V:} \texttt{OpControlBarrier / OpReturn}; \textbf{GLSL:} \texttt{// Scoreboard wait + warp reconvergence})
\end{tightitemize}
\clearpage

\subsection*{LSHIFT}
\addcontentsline{toc}{subsection}{LSHIFT}
For each SIMD element $i \in \{0, \dots, N-1\}$ ($N=1$ for \texttt{i32}/\texttt{i64}, $N=2$ for \texttt{v2i16}, $N=4$ for \texttt{v4i8} on v9--v10, lane width $W$), computes logical left shift \texttt{s = (val \textless{}\textless{} (amt \& (W - 1)))} where \texttt{val = src0[i]} and the 8-bit shift count \texttt{amt} is extracted from \texttt{src1} via \texttt{lanes1}; writes \texttt{not\_result ? \textasciitilde{}s : s} to \texttt{dest[i]}.
\par\textbf{Notes:}
\begin{tightitemize}
  \item Shift amount is masked modulo lane width (\texttt{wrapping\_shl}/\texttt{wrapping\_shr}/\texttt{rotate\_*}, \texttt{kraid/ops.rs:3984-4009}).
  \item On \texttt{v2i16}, \texttt{src1} is \texttt{V2I8} (\texttt{ops.rs:3953}) with byte lanes selected by \texttt{lanes1} (default \texttt{b02}: byte 0 -\textgreater{} lane 0, byte 2 -\textgreater{} lane 1).
  \item \textbf{Source references:} \texttt{kraid/ops.rs:3878-4025}, \texttt{encode\_v9.rs:2940-3093}, \texttt{hw\_tests.rs:2148-2198}.
\end{tightitemize}
\codebox{\mnem{LSHIFT}\imod{\{.flow\}}\imod{\{.not\_result\}} \dst{dest}\imod{\{.dest\_width\}}, \src{src0}\imod{\{.widen0\}}, \src{src1}\imod{\{.lanes1\}}\\\mnem{LSHIFT}\imod{\{.flow\}}\imod{\{.not\_result\}} \dst{dest}, \src{src0}\imod{\{.widen0\}}, \src{src1}\imod{\{.lanes1\}}\\\mnem{LSHIFT}\imod{\{.flow\}}\imod{\{.not\_result\}} \dst{dest}\imod{\{.dest\_width\}}, \src{src0}\imod{\{.swz0\}}, \src{src1}\imod{\{.lanes1\}}\\\mnem{LSHIFT}\imod{\{.flow\}}\imod{\{.not\_result\}} \dst{dest}\imod{\{.dest\_width\}}, \src{src0}\imod{\{.lanes0\}}, \src{src1}\imod{\{.lanes1\}}}
\textbf{Operands:}
\begin{tightitemize}\footnotesize
  \item \texttt{\dst{dest}} (\texttt{dest\_register\_index}, \texttt{dest\_width}): Destination GPR receiving the logical left shift (inverted if \texttt{.not\_result}).
  \item \texttt{\src{src0}}: Input integer value to shift/rotate.
  \item \texttt{\src{src1}}: 8-bit shift/rotate count per lane (selected via \texttt{lanes1}, masked modulo lane width \texttt{W}).
\end{tightitemize}
\par\textbf{Closest SPIR-V Op:} \texttt{OpShiftLeftLogical}
\par\textbf{Pseudo-GLSL (Operational Semantics):}
\codebox{for (int i = 0; i \textless{} N; ++i) \{\\    uint s = extractByte(src1, lanes1, i) \& (W - 1);\\    auto r = (src0[i] \textless{}\textless{} s);\\    dest[i] = flag\_not\_result ? \textasciitilde{}r : r;\\\}}
\textbf{Execution Context:}
\begin{tightitemize}\footnotesize
  \item \textbf{Compute stage preloads (\hyperref[sec:csf-spd]{\texttt{SPD\_k}}, \texttt{r16..r23}):} Hardware pre-fills \texttt{r55..r57} with \texttt{WorkGroupID.xyz + JOB\_OFFSET} (from \hyperref[sec:csf-regs]{\texttt{COMPUTESR::JOB\_OFFSET\_X/Y/Z}} \texttt{r34..r36}) and \texttt{r60..r62} with \texttt{LocalInvocationID.xyz} or \texttt{GlobalInvocationID.xyz} (from \texttt{WG\_SIZE} \texttt{r33} \& \texttt{JOB\_SIZE} \texttt{r37..r39}, selected by \texttt{ShaderProgram::preload}) so compute shaders can derive thread/element indices without memory loads.
  \item \textbf{Vertex stage preloads (\hyperref[sec:csf-spd]{\texttt{SPD\_0}/\texttt{SPD\_1}}, \texttt{r16..r19}):} The hardware Index Fetcher reads \hyperref[sec:csf-regs]{\texttt{IDVSSR::INDEX\_BUFFER}} (\texttt{r54:r55}) and pre-fills \texttt{r59, r60} with \texttt{VertexIndex} (\texttt{+ VERTEX\_OFFSET + GLOBAL\_ATTR\_OFF}) and \texttt{r61} with \texttt{InstanceIndex} (\texttt{+ INSTANCE\_OFFSET}) so Position and Varying shaders can pull vertex attributes directly.
  \item \textbf{Inline FAU \& Special Operands (\hyperref[sec:csf-tsd-fau]{\texttt{FAU\_k}}, \texttt{r8..r15}):} Any 8-bit source with \texttt{mode=2} reads zero-latency uniform \texttt{u0..u31} from \texttt{FAU\_k} (\texttt{slots64 = bits[63:56]}); \texttt{mode=3} reads \texttt{SmallConstant} (Table~1.2) or \texttt{FauSpecial} (\texttt{pc}, \texttt{lane\_id}, \texttt{lane\_width}, \texttt{tls\_ptr}, \texttt{wls\_ptr}, \texttt{blend\_descriptor\_0..7}).
\end{tightitemize}
\par\textbf{Worked Examples:}
\par\textbf{Example 1 (Canonical 32-bit shift/rotate):}
\codebox{\mnem{LSHIFT.i32} \dst{r0}, \src{r1}, \src{r2}\imod{.b0}}
With \texttt{r1 = 0x00000010}, \texttt{r2.b0 = 4}: yields \texttt{r0 = 0x00000100}.
\par\textbf{Example 2 (Byte-lane selector \& modulo-32 masking):}
\codebox{\mnem{LSHIFT.i32} \dst{r0}, \src{r1}, \src{r2}\imod{.b1}}
With \texttt{r1 = 0x00000001}, \texttt{r2.b1 = 33} (shift count masked mod 32 -\textgreater{} \texttt{1}): yields \texttt{r0 = 0x00000002}.

\begin{center}\footnotesize
\begin{tabularx}{\linewidth}{@{}lllX@{}}
\toprule
Variant & Arch & Class & Opcode Fields \\
\midrule
\texttt{(none)} & v9--v14 & \texttt{R3} & \texttt{opcode1=SHIFT\_BITWISE\_I32} \\
\texttt{(none)} & v9--v14 & \texttt{R3} & \texttt{opcode1=SHIFT\_BITWISE\_I32} \\
\texttt{(none)} & v9--v14 & \texttt{R3} & \texttt{opcode1=LOGIC\_SHIFT\_I32} \\
\texttt{(none)} & v9--v14 & \texttt{D3} & \texttt{opcode1=SHIFT\_BITWISE\_I64} \\
\texttt{(none)} & v9--v14 & \texttt{D3} & \texttt{opcode1=LOGIC\_SHIFT\_I64} \\
\texttt{(none)} & v9--v14 & \texttt{R3} & \texttt{opcode1=SHIFT\_BITWISE\_V2I16} \\
\texttt{(none)} & v9--v14 & \texttt{R3} & \texttt{opcode1=LOGIC\_SHIFT\_V2I16} \\
\texttt{(none)} & v9--v14 & \texttt{R3} & \texttt{opcode1=SHIFT\_BITWISE\_V4I8} \\
\texttt{(none)} & v9--v14 & \texttt{R3} & \texttt{opcode1=LOGIC\_SHIFT\_V4I8} \\
\bottomrule
\end{tabularx}
\end{center}
\begin{center}\footnotesize
\begin{tabularx}{0.88\linewidth}{@{}l l X@{}}
\toprule
Field Name & Bits & Type / Fixed Value \\
\midrule
\texttt{src0} & \texttt{0:7} & \texttt{SourceEncoding} \\
\texttt{src1} & \texttt{8:15} & \texttt{SourceEncoding} \\
\texttt{src2} & \texttt{16:23} & \texttt{exact=192} \\
\texttt{bitwise\_operation} & \texttt{24:25} & \texttt{exact=or} \\
\hyperref[enum:shift_bitwise_i32_lanes_m]{\texttt{lanes1}} & \texttt{26:29} & \texttt{shift\_bitwise\_i32\_lanes\_m (def b0)} \\
\hyperref[enum:not_m]{\texttt{not\_result}} & \texttt{30} & \texttt{not\_m (def none)} \\
\textit{Reserved} & \texttt{31} & \texttt{zero} \\
\texttt{shift\_operation} & \texttt{32:34} & \texttt{exact=lshift} \\
\hyperref[enum:not_m]{\texttt{not2}} & \texttt{35} & \texttt{exact=none} \\
\hyperref[enum:widen_int_m]{\texttt{widen0}} & \texttt{36:39} & \texttt{widen\_int\_m (def none)} \\
\texttt{dest\_register\_index} & \texttt{40:45} & \texttt{uint} \\
\hyperref[enum:dest_width_m]{\texttt{dest\_width}} & \texttt{46:47} & \texttt{dest\_width\_m} \\
\texttt{opcode1} & \texttt{48:51} & \texttt{exact=SHIFT\_BITWISE\_I32} \\
\texttt{instruction\_class} & \texttt{52:56} & \texttt{exact=R3} \\
\texttt{fau\_page\_index} & \texttt{57:58} & \texttt{uint} \\
\hyperref[enum:flow_control_m]{\texttt{flow}} & \texttt{59:62} & \texttt{flow\_control\_m} \\
\textit{Reserved} & \texttt{63} & \texttt{zero} \\
\bottomrule
\end{tabularx}
\end{center}
\textbf{Fields Explained:}
\begin{tightitemize}\footnotesize
  \item \texttt{src0}: Primary source operand (\texttt{r0..r63}, \texttt{u0..u63} uniform, \texttt{k*} FAU constant, or \texttt{t0..t2} temporary; \texttt{\textasciicircum{}} marks discard/last-use). \emph{Example:} {\scriptsize\ttfamily \mnem{FADD.f32} \dst{r0}, \src{r1}, \src{r2}} (\textbf{SPIR-V:} \texttt{OpFAdd (Operand 0)}; \textbf{GLSL:} \texttt{float c = a + b;})
  \item \texttt{src1}: Secondary source operand (register, uniform, or FAU constant; in shifts, supplies the shift count byte). \emph{Example:} {\scriptsize\ttfamily \mnem{LSHIFT.i32} \dst{r0}, \src{r1}, \src{r2}\imod{.b0}} (\textbf{SPIR-V:} \texttt{OpShiftLeftLogical (Shift)}; \textbf{GLSL:} \texttt{uint y = a \textless{}\textless{} s;})
  \item \texttt{src2}: Third source operand (FMA addend, \texttt{MUX}/\texttt{CSEL} selector, \texttt{IADDC} carry-in, or bitwise mask in fused shift-logic). \emph{Example:} {\scriptsize\ttfamily \mnem{FMA.f32} \dst{r0}, \src{r1}, \src{r2}, \src{r3}} (\textbf{SPIR-V:} \texttt{GLSL.std.450 Fma (Addend)}; \textbf{GLSL:} \texttt{float d = fma(a, b, c);})
  \item \texttt{bitwise\_operation}: Selects the bitwise ALU op (\texttt{and}, \texttt{or}, \texttt{xor}, \texttt{lshift\_and}, etc.) fused with shifts. \emph{Example:} {\scriptsize\ttfamily \mnem{LSHIFT\_OR.i32} \dst{r0}, \src{r1}, \src{r2}, \src{r3}} (\textbf{SPIR-V:} \texttt{OpShiftLeftLogical + OpBitwiseOr}; \textbf{GLSL:} \texttt{uint y = (a \textless{}\textless{} s) \textbar{} b;})
  \item \hyperref[enum:shift_bitwise_i32_lanes_m]{\texttt{lanes1}}: Stage A byte-lane selector on \texttt{src1} (selects shift-count byte \texttt{b0..b3} on 32/64-bit shifts, \texttt{bXY} on \texttt{v2i16} shifts, or 4-byte swizzle on \texttt{v4*8}). \emph{Example:} {\scriptsize\ttfamily \mnem{LSHIFT.i32} \dst{r0}, \src{r1}, \src{r2}\imod{.b1}} (\textbf{SPIR-V:} \texttt{OpBitFieldUExtract + OpShiftLeftLogical}; \textbf{GLSL:} \texttt{uint y = a \textless{}\textless{} ((s \textgreater{}\textgreater{} 8) \& 0xFFu);})
  \item \hyperref[enum:not_m]{\texttt{not\_result}}: Bitwise-inverts the final output of a bitwise/shift-bitwise instruction (\textasciitilde{}result). \emph{Example:} {\scriptsize\ttfamily \mnem{AND.i32}\imod{.not\_result} \dst{r0}, \src{r1}, \src{r2}} (\textbf{SPIR-V:} \texttt{OpBitwiseAnd + OpNot}; \textbf{GLSL:} \texttt{uint y = \textasciitilde{}(a \& b);})
  \item \texttt{shift\_operation}: Selects shift direction/type (\texttt{lshift}, \texttt{rshift}, \texttt{arshift}) in fused shift-bitwise instructions. \emph{Example:} {\scriptsize\ttfamily \mnem{RSHIFT\_AND.i32} \dst{r0}, \src{r1}, \src{r2}, \src{r3}} (\textbf{SPIR-V:} \texttt{OpShiftRightLogical + OpBitwiseAnd}; \textbf{GLSL:} \texttt{uint y = (a \textgreater{}\textgreater{} s) \& m;})
  \item \hyperref[enum:not_m]{\texttt{not2}}: Stage A bitwise complement (\texttt{\textasciitilde{}src2}) applied to \texttt{src2} before the bitwise/shift-logic operation. \emph{Example:} {\scriptsize\ttfamily \mnem{AND.i32} \dst{r0}, \src{r1}, \src{r2}\imod{.not}} (\textbf{SPIR-V:} \texttt{OpNot + OpBitwiseAnd}; \textbf{GLSL:} \texttt{uint y = a \& \textasciitilde{}b;})
  \item \hyperref[enum:widen_int_m]{\texttt{widen0}}: Stage A widening extractor on \texttt{src0} (promotes FP16 \texttt{h0}/\texttt{h1} to FP32, or sign/zero-extends \texttt{b0..b3}/\texttt{h0..h1}/\texttt{w0..w1} to 32/64-bit before \texttt{abs}/\texttt{neg}). \emph{Example:} {\scriptsize\ttfamily \mnem{IADD.s32} \dst{r0}, \src{r1}\imod{.h0}, \src{r2}} (\textbf{SPIR-V:} \texttt{OpSConvert / OpFConvert}; \textbf{GLSL:} \texttt{int c = int(int16\_t(a)) + b;})
  \item \texttt{dest\_register\_index}: 6-bit destination GPR index (\texttt{r0..r63}) receiving the Stage C writeback result (even-aligned pair \texttt{[rN:rN+1]} for 64-bit ops). \emph{Example:} {\scriptsize\ttfamily \mnem{FMA.f32} \dst{r0}, \src{r1}, \src{r2}, \src{r3}} (\textbf{SPIR-V:} \texttt{SSA Result ID}; \textbf{GLSL:} \texttt{float r0 = fma(r1, r2, r3);})
  \item \hyperref[enum:dest_width_m]{\texttt{dest\_width}}: Stage C destination lane write-mask (\texttt{none} = full 32-bit, \texttt{h0} = low 16-bit, \texttt{h1} = high 16-bit without bit-shifting). \emph{Example:} {\scriptsize\ttfamily \mnem{FADD.v2f16} \dst{r0}\imod{.h1}, \src{r1}, \src{r2}} (\textbf{SPIR-V:} \texttt{OpVectorInsertDynamic / WriteMask}; \textbf{GLSL:} \texttt{dst.y = a.y + b.y;})
  \item \texttt{opcode1}: Primary sub-opcode within the instruction class uniquely identifying the operation family. \emph{Example:} {\scriptsize\ttfamily \mnem{IADD.s32} \dst{r0}, \src{r1}, \src{r2}} (\textbf{SPIR-V:} \texttt{OpIAdd (operation selector)}; \textbf{GLSL:} \texttt{// Primary opcode bits})
  \item \texttt{instruction\_class}: 4-bit major instruction class selector (\texttt{bits 48:51}) partitioning the 64-bit Valhall opcode space (e.g., \texttt{FMA}, \texttt{CVT}, \texttt{SFU}, \texttt{LS}, \texttt{TEX}, \texttt{BRANCH}). \emph{Example:} {\scriptsize\ttfamily \mnem{FADD.f32} \dst{r0}, \src{r1}, \src{r2}} (\textbf{SPIR-V:} \texttt{N/A (encoding class)}; \textbf{GLSL:} \texttt{// Hardware execution pipe / decode class})
  \item \texttt{fau\_page\_index}: Selects the FAU (Fast Access Uniform) page (\texttt{0..3}) when a source operand references hardware small constants, thread IDs, or lane constants. \emph{Example:} {\scriptsize\ttfamily \mnem{FADD.f32} \dst{r0}, \src{r1}, \src{k1f}} (\textbf{SPIR-V:} \texttt{OpConstant / BuiltIn}; \textbf{GLSL:} \texttt{float y = x + 1.0;})
  \item \hyperref[enum:flow_control_m]{\texttt{flow}}: 4-bit embedded flow-control \& scoreboard wait modifier (\texttt{bits 59:62}: \texttt{none}, \texttt{wait0..wait0126}, \texttt{wait}, \texttt{reconverge}, \texttt{discard}, \texttt{end}). \emph{Example:} {\scriptsize\ttfamily \mnem{FADD.f32}\imod{.wait0.reconverge} \dst{r0}, \src{r1}, \src{r2}} (\textbf{SPIR-V:} \texttt{OpControlBarrier / OpReturn}; \textbf{GLSL:} \texttt{// Scoreboard wait + warp reconvergence})
\end{tightitemize}
\clearpage

\subsection*{LSHIFT\_AND}
\addcontentsline{toc}{subsection}{LSHIFT\_AND}
For each SIMD element $i \in \{0, \dots, N-1\}$ ($N=1$ for \texttt{i32}/\texttt{i64}, $N=2$ for \texttt{v2i16}, $N=4$ for \texttt{v4i8} on v9--v10, lane width $W$), computes logical left shift \texttt{s = (val \textless{}\textless{} (amt \& (W - 1)))} where \texttt{val = src0[i]} and the 8-bit shift count \texttt{amt} is extracted from \texttt{src1} via \texttt{lanes1}, then combines with \texttt{src2[i]} via bitwise AND \texttt{res = s \& (not2 ? \textasciitilde{}src2[i] : src2[i])} and writes \texttt{not\_result ? \textasciitilde{}res : res} to \texttt{dest[i]}.
\par\textbf{Notes:}
\begin{tightitemize}
  \item Shift amount is masked modulo lane width (\texttt{wrapping\_shl}/\texttt{wrapping\_shr}/\texttt{rotate\_*}, \texttt{kraid/ops.rs:3984-4009}).
  \item On \texttt{v2i16}, \texttt{src1} is \texttt{V2I8} (\texttt{ops.rs:3953}) with byte lanes selected by \texttt{lanes1} (default \texttt{b02}: byte 0 -\textgreater{} lane 0, byte 2 -\textgreater{} lane 1).
  \item \textbf{Source references:} \texttt{kraid/ops.rs:3878-4025}, \texttt{encode\_v9.rs:2940-3093}, \texttt{hw\_tests.rs:2148-2198}.
\end{tightitemize}
\codebox{\mnem{LSHIFT\_AND}\imod{\{.flow\}}\imod{\{.not\_result\}} \dst{dest}\imod{\{.dest\_width\}}, \src{src0}\imod{\{.widen0\}}, \src{src1}\imod{\{.lanes1\}}, \src{src2}\imod{\{.not2\}}\\\mnem{LSHIFT\_AND}\imod{\{.flow\}}\imod{\{.not\_result\}} \dst{dest}, \src{src0}\imod{\{.widen0\}}, \src{src1}\imod{\{.lanes1\}}, \src{src2}\imod{\{.not2\}}\\\mnem{LSHIFT\_AND}\imod{\{.flow\}}\imod{\{.not\_result\}} \dst{dest}\imod{\{.dest\_width\}}, \src{src0}\imod{\{.swz0\}}, \src{src1}\imod{\{.lanes1\}}, \src{src2}\imod{\{.not2\}}\\\mnem{LSHIFT\_AND}\imod{\{.flow\}}\imod{\{.not\_result\}} \dst{dest}\imod{\{.dest\_width\}}, \src{src0}\imod{\{.lanes0\}}, \src{src1}\imod{\{.lanes1\}}, \src{src2}\imod{\{.not2\}}}
\textbf{Operands:}
\begin{tightitemize}\footnotesize
  \item \texttt{\dst{dest}} (\texttt{dest\_register\_index}, \texttt{dest\_width}): Destination GPR receiving the logical left shift combined via bitwise AND (inverted if \texttt{.not\_result}).
  \item \texttt{\src{src0}}: Input integer value to shift/rotate.
  \item \texttt{\src{src1}}: 8-bit shift/rotate count per lane (selected via \texttt{lanes1}, masked modulo lane width \texttt{W}).
  \item \texttt{\src{src2}}: Second operand combined via bitwise AND after shifting (inverted if \texttt{.not2}).
\end{tightitemize}
\par\textbf{Closest SPIR-V Op:} \texttt{OpShiftLeftLogical + OpBitwiseAnd}
\par\textbf{Pseudo-GLSL (Operational Semantics):}
\codebox{for (int i = 0; i \textless{} N; ++i) \{\\    uint s = extractByte(src1, lanes1, i) \& (W - 1);\\    auto r = (src0[i] \textless{}\textless{} s) \& (flag\_not2 ? \textasciitilde{}src2[i] : src2[i]);\\    dest[i] = flag\_not\_result ? \textasciitilde{}r : r;\\\}}
\textbf{Execution Context:}
\begin{tightitemize}\footnotesize
  \item \textbf{Compute stage preloads (\hyperref[sec:csf-spd]{\texttt{SPD\_k}}, \texttt{r16..r23}):} Hardware pre-fills \texttt{r55..r57} with \texttt{WorkGroupID.xyz + JOB\_OFFSET} (from \hyperref[sec:csf-regs]{\texttt{COMPUTESR::JOB\_OFFSET\_X/Y/Z}} \texttt{r34..r36}) and \texttt{r60..r62} with \texttt{LocalInvocationID.xyz} or \texttt{GlobalInvocationID.xyz} (from \texttt{WG\_SIZE} \texttt{r33} \& \texttt{JOB\_SIZE} \texttt{r37..r39}, selected by \texttt{ShaderProgram::preload}) so compute shaders can derive thread/element indices without memory loads.
  \item \textbf{Vertex stage preloads (\hyperref[sec:csf-spd]{\texttt{SPD\_0}/\texttt{SPD\_1}}, \texttt{r16..r19}):} The hardware Index Fetcher reads \hyperref[sec:csf-regs]{\texttt{IDVSSR::INDEX\_BUFFER}} (\texttt{r54:r55}) and pre-fills \texttt{r59, r60} with \texttt{VertexIndex} (\texttt{+ VERTEX\_OFFSET + GLOBAL\_ATTR\_OFF}) and \texttt{r61} with \texttt{InstanceIndex} (\texttt{+ INSTANCE\_OFFSET}) so Position and Varying shaders can pull vertex attributes directly.
  \item \textbf{Inline FAU \& Special Operands (\hyperref[sec:csf-tsd-fau]{\texttt{FAU\_k}}, \texttt{r8..r15}):} Any 8-bit source with \texttt{mode=2} reads zero-latency uniform \texttt{u0..u31} from \texttt{FAU\_k} (\texttt{slots64 = bits[63:56]}); \texttt{mode=3} reads \texttt{SmallConstant} (Table~1.2) or \texttt{FauSpecial} (\texttt{pc}, \texttt{lane\_id}, \texttt{lane\_width}, \texttt{tls\_ptr}, \texttt{wls\_ptr}, \texttt{blend\_descriptor\_0..7}).
\end{tightitemize}
\par\textbf{Worked Examples:}
\par\textbf{Example 1 (Standard fused shift + bitwise op):}
\codebox{\mnem{LSHIFT\_AND.i32} \dst{r0}, \src{r1}, \src{r2}\imod{.b0}, \src{r3}}
With \texttt{r1 = 0x00000010}, \texttt{r2.b0 = 4} (masked mod 32): shift yields \texttt{0x00000100}, then \texttt{\& 0x0FFFFFFF} yields \texttt{r0 = 0x100}.
\par\textbf{Example 2 (Modulo-32 shift count \& .not\_result):}
\codebox{\mnem{LSHIFT\_AND.i32}\imod{.not\_result} \dst{r0}, \src{r1}, \src{r2}\imod{.b0}, \src{r3}}
With \texttt{r1 = 0x00000001}, \texttt{r2.b0 = 33} (effective shift \texttt{1}): shift yields \texttt{0x00000002}, \texttt{\& 0x0FFFFFFF} yields \texttt{0x00000002}, and \texttt{.not\_result} inverts -\textgreater{} \texttt{r0 = 0xFFFFFFFD}.
\par\textbf{Mined SPIR-V $\to$ Valhall ISA Example (Bitmasking vertex attribute bits via AND alias --- Vertex shader, Mali-G710 (v10)):}
\codebox{\textbf{; SPIR-V (input)}\\\%bit1 = OpBitwiseAnd \%int \%bits \%int\_2\\\%bit0 = OpBitwiseAnd \%int \%bits \%int\_1\\[3pt]\textbf{; Valhall ISA (malioc output)}\\AND.i32 r1, k03020100.b2, r0\\AND.i32 r0, k01000000.b3, r0\textasciicircum{}}
Standalone \texttt{AND.i32} encodes as \texttt{LSHIFT\_AND} with shift amount 0. Notice \texttt{malioc} pulls small bitmasks (\texttt{2} from \texttt{k03020100.b2}, \texttt{1} from \texttt{k01000000.b3}) directly from FAU byte lanes.
\begin{center}\footnotesize
\begin{tabularx}{\linewidth}{@{}lllX@{}}
\toprule
Variant & Arch & Class & Opcode Fields \\
\midrule
\texttt{(none)} & v9--v14 & \texttt{R3} & \texttt{opcode1=SHIFT\_BITWISE\_I32} \\
\texttt{(none)} & v9--v14 & \texttt{R3} & \texttt{opcode1=SHIFT\_BITWISE\_I32} \\
\texttt{(none)} & v9--v14 & \texttt{R3} & \texttt{opcode1=LOGIC\_SHIFT\_I32} \\
\texttt{(none)} & v9--v14 & \texttt{D3} & \texttt{opcode1=SHIFT\_BITWISE\_I64} \\
\texttt{(none)} & v9--v14 & \texttt{D3} & \texttt{opcode1=LOGIC\_SHIFT\_I64} \\
\texttt{(none)} & v9--v14 & \texttt{R3} & \texttt{opcode1=SHIFT\_BITWISE\_V2I16} \\
\texttt{(none)} & v9--v14 & \texttt{R3} & \texttt{opcode1=LOGIC\_SHIFT\_V2I16} \\
\texttt{(none)} & v9--v14 & \texttt{R3} & \texttt{opcode1=SHIFT\_BITWISE\_V4I8} \\
\texttt{(none)} & v9--v14 & \texttt{R3} & \texttt{opcode1=LOGIC\_SHIFT\_V4I8} \\
\bottomrule
\end{tabularx}
\end{center}
\begin{center}\footnotesize
\begin{tabularx}{0.88\linewidth}{@{}l l X@{}}
\toprule
Field Name & Bits & Type / Fixed Value \\
\midrule
\texttt{src0} & \texttt{0:7} & \texttt{SourceEncoding} \\
\texttt{src1} & \texttt{8:15} & \texttt{SourceEncoding} \\
\texttt{src2} & \texttt{16:23} & \texttt{SourceEncoding} \\
\texttt{bitwise\_operation} & \texttt{24:25} & \texttt{exact=and} \\
\hyperref[enum:shift_bitwise_i32_lanes_m]{\texttt{lanes1}} & \texttt{26:29} & \texttt{shift\_bitwise\_i32\_lanes\_m (def b0)} \\
\hyperref[enum:not_m]{\texttt{not\_result}} & \texttt{30} & \texttt{not\_m (def none)} \\
\textit{Reserved} & \texttt{31} & \texttt{zero} \\
\texttt{shift\_operation} & \texttt{32:34} & \texttt{exact=lshift} \\
\hyperref[enum:not_m]{\texttt{not2}} & \texttt{35} & \texttt{not\_m (def none)} \\
\hyperref[enum:widen_int_m]{\texttt{widen0}} & \texttt{36:39} & \texttt{widen\_int\_m (def none)} \\
\texttt{dest\_register\_index} & \texttt{40:45} & \texttt{uint} \\
\hyperref[enum:dest_width_m]{\texttt{dest\_width}} & \texttt{46:47} & \texttt{dest\_width\_m} \\
\texttt{opcode1} & \texttt{48:51} & \texttt{exact=SHIFT\_BITWISE\_I32} \\
\texttt{instruction\_class} & \texttt{52:56} & \texttt{exact=R3} \\
\texttt{fau\_page\_index} & \texttt{57:58} & \texttt{uint} \\
\hyperref[enum:flow_control_m]{\texttt{flow}} & \texttt{59:62} & \texttt{flow\_control\_m} \\
\textit{Reserved} & \texttt{63} & \texttt{zero} \\
\bottomrule
\end{tabularx}
\end{center}
\textbf{Fields Explained:}
\begin{tightitemize}\footnotesize
  \item \texttt{src0}: Primary source operand (\texttt{r0..r63}, \texttt{u0..u63} uniform, \texttt{k*} FAU constant, or \texttt{t0..t2} temporary; \texttt{\textasciicircum{}} marks discard/last-use). \emph{Example:} {\scriptsize\ttfamily \mnem{FADD.f32} \dst{r0}, \src{r1}, \src{r2}} (\textbf{SPIR-V:} \texttt{OpFAdd (Operand 0)}; \textbf{GLSL:} \texttt{float c = a + b;})
  \item \texttt{src1}: Secondary source operand (register, uniform, or FAU constant; in shifts, supplies the shift count byte). \emph{Example:} {\scriptsize\ttfamily \mnem{LSHIFT.i32} \dst{r0}, \src{r1}, \src{r2}\imod{.b0}} (\textbf{SPIR-V:} \texttt{OpShiftLeftLogical (Shift)}; \textbf{GLSL:} \texttt{uint y = a \textless{}\textless{} s;})
  \item \texttt{src2}: Third source operand (FMA addend, \texttt{MUX}/\texttt{CSEL} selector, \texttt{IADDC} carry-in, or bitwise mask in fused shift-logic). \emph{Example:} {\scriptsize\ttfamily \mnem{FMA.f32} \dst{r0}, \src{r1}, \src{r2}, \src{r3}} (\textbf{SPIR-V:} \texttt{GLSL.std.450 Fma (Addend)}; \textbf{GLSL:} \texttt{float d = fma(a, b, c);})
  \item \texttt{bitwise\_operation}: Selects the bitwise ALU op (\texttt{and}, \texttt{or}, \texttt{xor}, \texttt{lshift\_and}, etc.) fused with shifts. \emph{Example:} {\scriptsize\ttfamily \mnem{LSHIFT\_OR.i32} \dst{r0}, \src{r1}, \src{r2}, \src{r3}} (\textbf{SPIR-V:} \texttt{OpShiftLeftLogical + OpBitwiseOr}; \textbf{GLSL:} \texttt{uint y = (a \textless{}\textless{} s) \textbar{} b;})
  \item \hyperref[enum:shift_bitwise_i32_lanes_m]{\texttt{lanes1}}: Stage A byte-lane selector on \texttt{src1} (selects shift-count byte \texttt{b0..b3} on 32/64-bit shifts, \texttt{bXY} on \texttt{v2i16} shifts, or 4-byte swizzle on \texttt{v4*8}). \emph{Example:} {\scriptsize\ttfamily \mnem{LSHIFT.i32} \dst{r0}, \src{r1}, \src{r2}\imod{.b1}} (\textbf{SPIR-V:} \texttt{OpBitFieldUExtract + OpShiftLeftLogical}; \textbf{GLSL:} \texttt{uint y = a \textless{}\textless{} ((s \textgreater{}\textgreater{} 8) \& 0xFFu);})
  \item \hyperref[enum:not_m]{\texttt{not\_result}}: Bitwise-inverts the final output of a bitwise/shift-bitwise instruction (\textasciitilde{}result). \emph{Example:} {\scriptsize\ttfamily \mnem{AND.i32}\imod{.not\_result} \dst{r0}, \src{r1}, \src{r2}} (\textbf{SPIR-V:} \texttt{OpBitwiseAnd + OpNot}; \textbf{GLSL:} \texttt{uint y = \textasciitilde{}(a \& b);})
  \item \texttt{shift\_operation}: Selects shift direction/type (\texttt{lshift}, \texttt{rshift}, \texttt{arshift}) in fused shift-bitwise instructions. \emph{Example:} {\scriptsize\ttfamily \mnem{RSHIFT\_AND.i32} \dst{r0}, \src{r1}, \src{r2}, \src{r3}} (\textbf{SPIR-V:} \texttt{OpShiftRightLogical + OpBitwiseAnd}; \textbf{GLSL:} \texttt{uint y = (a \textgreater{}\textgreater{} s) \& m;})
  \item \hyperref[enum:not_m]{\texttt{not2}}: Stage A bitwise complement (\texttt{\textasciitilde{}src2}) applied to \texttt{src2} before the bitwise/shift-logic operation. \emph{Example:} {\scriptsize\ttfamily \mnem{AND.i32} \dst{r0}, \src{r1}, \src{r2}\imod{.not}} (\textbf{SPIR-V:} \texttt{OpNot + OpBitwiseAnd}; \textbf{GLSL:} \texttt{uint y = a \& \textasciitilde{}b;})
  \item \hyperref[enum:widen_int_m]{\texttt{widen0}}: Stage A widening extractor on \texttt{src0} (promotes FP16 \texttt{h0}/\texttt{h1} to FP32, or sign/zero-extends \texttt{b0..b3}/\texttt{h0..h1}/\texttt{w0..w1} to 32/64-bit before \texttt{abs}/\texttt{neg}). \emph{Example:} {\scriptsize\ttfamily \mnem{IADD.s32} \dst{r0}, \src{r1}\imod{.h0}, \src{r2}} (\textbf{SPIR-V:} \texttt{OpSConvert / OpFConvert}; \textbf{GLSL:} \texttt{int c = int(int16\_t(a)) + b;})
  \item \texttt{dest\_register\_index}: 6-bit destination GPR index (\texttt{r0..r63}) receiving the Stage C writeback result (even-aligned pair \texttt{[rN:rN+1]} for 64-bit ops). \emph{Example:} {\scriptsize\ttfamily \mnem{FMA.f32} \dst{r0}, \src{r1}, \src{r2}, \src{r3}} (\textbf{SPIR-V:} \texttt{SSA Result ID}; \textbf{GLSL:} \texttt{float r0 = fma(r1, r2, r3);})
  \item \hyperref[enum:dest_width_m]{\texttt{dest\_width}}: Stage C destination lane write-mask (\texttt{none} = full 32-bit, \texttt{h0} = low 16-bit, \texttt{h1} = high 16-bit without bit-shifting). \emph{Example:} {\scriptsize\ttfamily \mnem{FADD.v2f16} \dst{r0}\imod{.h1}, \src{r1}, \src{r2}} (\textbf{SPIR-V:} \texttt{OpVectorInsertDynamic / WriteMask}; \textbf{GLSL:} \texttt{dst.y = a.y + b.y;})
  \item \texttt{opcode1}: Primary sub-opcode within the instruction class uniquely identifying the operation family. \emph{Example:} {\scriptsize\ttfamily \mnem{IADD.s32} \dst{r0}, \src{r1}, \src{r2}} (\textbf{SPIR-V:} \texttt{OpIAdd (operation selector)}; \textbf{GLSL:} \texttt{// Primary opcode bits})
  \item \texttt{instruction\_class}: 4-bit major instruction class selector (\texttt{bits 48:51}) partitioning the 64-bit Valhall opcode space (e.g., \texttt{FMA}, \texttt{CVT}, \texttt{SFU}, \texttt{LS}, \texttt{TEX}, \texttt{BRANCH}). \emph{Example:} {\scriptsize\ttfamily \mnem{FADD.f32} \dst{r0}, \src{r1}, \src{r2}} (\textbf{SPIR-V:} \texttt{N/A (encoding class)}; \textbf{GLSL:} \texttt{// Hardware execution pipe / decode class})
  \item \texttt{fau\_page\_index}: Selects the FAU (Fast Access Uniform) page (\texttt{0..3}) when a source operand references hardware small constants, thread IDs, or lane constants. \emph{Example:} {\scriptsize\ttfamily \mnem{FADD.f32} \dst{r0}, \src{r1}, \src{k1f}} (\textbf{SPIR-V:} \texttt{OpConstant / BuiltIn}; \textbf{GLSL:} \texttt{float y = x + 1.0;})
  \item \hyperref[enum:flow_control_m]{\texttt{flow}}: 4-bit embedded flow-control \& scoreboard wait modifier (\texttt{bits 59:62}: \texttt{none}, \texttt{wait0..wait0126}, \texttt{wait}, \texttt{reconverge}, \texttt{discard}, \texttt{end}). \emph{Example:} {\scriptsize\ttfamily \mnem{FADD.f32}\imod{.wait0.reconverge} \dst{r0}, \src{r1}, \src{r2}} (\textbf{SPIR-V:} \texttt{OpControlBarrier / OpReturn}; \textbf{GLSL:} \texttt{// Scoreboard wait + warp reconvergence})
\end{tightitemize}
\clearpage

\subsection*{LSHIFT\_OR}
\addcontentsline{toc}{subsection}{LSHIFT\_OR}
For each SIMD element $i \in \{0, \dots, N-1\}$ ($N=1$ for \texttt{i32}/\texttt{i64}, $N=2$ for \texttt{v2i16}, $N=4$ for \texttt{v4i8} on v9--v10, lane width $W$), computes logical left shift \texttt{s = (val \textless{}\textless{} (amt \& (W - 1)))} where \texttt{val = src0[i]} and the 8-bit shift count \texttt{amt} is extracted from \texttt{src1} via \texttt{lanes1}, then combines with \texttt{src2[i]} via bitwise OR \texttt{res = s \textbar{} (not2 ? \textasciitilde{}src2[i] : src2[i])} and writes \texttt{not\_result ? \textasciitilde{}res : res} to \texttt{dest[i]}.
\par\textbf{Notes:}
\begin{tightitemize}
  \item Shift amount is masked modulo lane width (\texttt{wrapping\_shl}/\texttt{wrapping\_shr}/\texttt{rotate\_*}, \texttt{kraid/ops.rs:3984-4009}).
  \item On \texttt{v2i16}, \texttt{src1} is \texttt{V2I8} (\texttt{ops.rs:3953}) with byte lanes selected by \texttt{lanes1} (default \texttt{b02}: byte 0 -\textgreater{} lane 0, byte 2 -\textgreater{} lane 1).
  \item \textbf{Source references:} \texttt{kraid/ops.rs:3878-4025}, \texttt{encode\_v9.rs:2940-3093}, \texttt{hw\_tests.rs:2148-2198}.
\end{tightitemize}
\codebox{\mnem{LSHIFT\_OR}\imod{\{.flow\}}\imod{\{.not\_result\}} \dst{dest}\imod{\{.dest\_width\}}, \src{src0}\imod{\{.widen0\}}, \src{src1}\imod{\{.lanes1\}}, \src{src2}\imod{\{.not2\}}\\\mnem{LSHIFT\_OR}\imod{\{.flow\}}\imod{\{.not\_result\}} \dst{dest}, \src{src0}\imod{\{.widen0\}}, \src{src1}\imod{\{.lanes1\}}, \src{src2}\imod{\{.not2\}}\\\mnem{LSHIFT\_OR}\imod{\{.flow\}}\imod{\{.not\_result\}} \dst{dest}\imod{\{.dest\_width\}}, \src{src0}\imod{\{.swz0\}}, \src{src1}\imod{\{.lanes1\}}, \src{src2}\imod{\{.not2\}}\\\mnem{LSHIFT\_OR}\imod{\{.flow\}}\imod{\{.not\_result\}} \dst{dest}\imod{\{.dest\_width\}}, \src{src0}\imod{\{.lanes0\}}, \src{src1}\imod{\{.lanes1\}}, \src{src2}\imod{\{.not2\}}}
\textbf{Operands:}
\begin{tightitemize}\footnotesize
  \item \texttt{\dst{dest}} (\texttt{dest\_register\_index}, \texttt{dest\_width}): Destination GPR receiving the logical left shift combined via bitwise OR (inverted if \texttt{.not\_result}).
  \item \texttt{\src{src0}}: Input integer value to shift/rotate.
  \item \texttt{\src{src1}}: 8-bit shift/rotate count per lane (selected via \texttt{lanes1}, masked modulo lane width \texttt{W}).
  \item \texttt{\src{src2}}: Second operand combined via bitwise OR after shifting (inverted if \texttt{.not2}).
\end{tightitemize}
\par\textbf{Closest SPIR-V Op:} \texttt{OpShiftLeftLogical + OpBitwiseOr}
\par\textbf{Pseudo-GLSL (Operational Semantics):}
\codebox{for (int i = 0; i \textless{} N; ++i) \{\\    uint s = extractByte(src1, lanes1, i) \& (W - 1);\\    auto r = (src0[i] \textless{}\textless{} s) \textbar{} (flag\_not2 ? \textasciitilde{}src2[i] : src2[i]);\\    dest[i] = flag\_not\_result ? \textasciitilde{}r : r;\\\}}
\textbf{Execution Context:}
\begin{tightitemize}\footnotesize
  \item \textbf{Compute stage preloads (\hyperref[sec:csf-spd]{\texttt{SPD\_k}}, \texttt{r16..r23}):} Hardware pre-fills \texttt{r55..r57} with \texttt{WorkGroupID.xyz + JOB\_OFFSET} (from \hyperref[sec:csf-regs]{\texttt{COMPUTESR::JOB\_OFFSET\_X/Y/Z}} \texttt{r34..r36}) and \texttt{r60..r62} with \texttt{LocalInvocationID.xyz} or \texttt{GlobalInvocationID.xyz} (from \texttt{WG\_SIZE} \texttt{r33} \& \texttt{JOB\_SIZE} \texttt{r37..r39}, selected by \texttt{ShaderProgram::preload}) so compute shaders can derive thread/element indices without memory loads.
  \item \textbf{Vertex stage preloads (\hyperref[sec:csf-spd]{\texttt{SPD\_0}/\texttt{SPD\_1}}, \texttt{r16..r19}):} The hardware Index Fetcher reads \hyperref[sec:csf-regs]{\texttt{IDVSSR::INDEX\_BUFFER}} (\texttt{r54:r55}) and pre-fills \texttt{r59, r60} with \texttt{VertexIndex} (\texttt{+ VERTEX\_OFFSET + GLOBAL\_ATTR\_OFF}) and \texttt{r61} with \texttt{InstanceIndex} (\texttt{+ INSTANCE\_OFFSET}) so Position and Varying shaders can pull vertex attributes directly.
  \item \textbf{Inline FAU \& Special Operands (\hyperref[sec:csf-tsd-fau]{\texttt{FAU\_k}}, \texttt{r8..r15}):} Any 8-bit source with \texttt{mode=2} reads zero-latency uniform \texttt{u0..u31} from \texttt{FAU\_k} (\texttt{slots64 = bits[63:56]}); \texttt{mode=3} reads \texttt{SmallConstant} (Table~1.2) or \texttt{FauSpecial} (\texttt{pc}, \texttt{lane\_id}, \texttt{lane\_width}, \texttt{tls\_ptr}, \texttt{wls\_ptr}, \texttt{blend\_descriptor\_0..7}).
\end{tightitemize}
\par\textbf{Worked Examples:}
\par\textbf{Example 1 (Standard fused shift + bitwise op):}
\codebox{\mnem{LSHIFT\_OR.i32} \dst{r0}, \src{r1}, \src{r2}\imod{.b0}, \src{r3}}
With \texttt{r1 = 0x00000010}, \texttt{r2.b0 = 4} (masked mod 32): shift yields \texttt{0x00000100}, then \texttt{\textbar{} 0x0000000F} yields \texttt{r0 = 0x10F}.
\par\textbf{Example 2 (Modulo-32 shift count \& .not\_result):}
\codebox{\mnem{LSHIFT\_OR.i32}\imod{.not\_result} \dst{r0}, \src{r1}, \src{r2}\imod{.b0}, \src{r3}}
With \texttt{r1 = 0x00000001}, \texttt{r2.b0 = 33} (effective shift \texttt{1}): shift yields \texttt{0x00000002}, \texttt{\textbar{} 0x0000000F} yields \texttt{0x0000000F}, and \texttt{.not\_result} inverts -\textgreater{} \texttt{r0 = 0xFFFFFFF0}.
\par\textbf{Mined SPIR-V $\to$ Valhall ISA Example (Combining two comparison masks for fragment discard --- Fragment shader, Mali-G710 (v10)):}
\codebox{\textbf{; SPIR-V (input)}\\\%is\_nz  = OpFOrdNotEqual \%bool \%v \%float\_0\\\%is\_nan = OpFUnordNotEqual \%bool \%v \%v\\\%either = OpLogicalOr \%bool \%is\_nz \%is\_nan\\[3pt]\textbf{; Valhall ISA (malioc output)}\\FCMP.f32.discard.m1.gtlt r0, r2, k0\\FCMP.f32.m1.ne r1, r2, r2\\OR.i32 r0, r0\textasciicircum{}, r1\textasciicircum{}}
Two \texttt{.m1} (\texttt{0xFFFFFFFF} / \texttt{0x00000000}) comparison masks are combined with \texttt{OR.i32} (encoded as \texttt{LSHIFT\_OR} with shift 0).
\begin{center}\footnotesize
\begin{tabularx}{\linewidth}{@{}lllX@{}}
\toprule
Variant & Arch & Class & Opcode Fields \\
\midrule
\texttt{(none)} & v9--v14 & \texttt{R3} & \texttt{opcode1=SHIFT\_BITWISE\_I32} \\
\texttt{(none)} & v9--v14 & \texttt{R3} & \texttt{opcode1=SHIFT\_BITWISE\_I32} \\
\texttt{(none)} & v9--v14 & \texttt{R3} & \texttt{opcode1=LOGIC\_SHIFT\_I32} \\
\texttt{(none)} & v9--v14 & \texttt{D3} & \texttt{opcode1=SHIFT\_BITWISE\_I64} \\
\texttt{(none)} & v9--v14 & \texttt{D3} & \texttt{opcode1=LOGIC\_SHIFT\_I64} \\
\texttt{(none)} & v9--v14 & \texttt{R3} & \texttt{opcode1=SHIFT\_BITWISE\_V2I16} \\
\texttt{(none)} & v9--v14 & \texttt{R3} & \texttt{opcode1=LOGIC\_SHIFT\_V2I16} \\
\texttt{(none)} & v9--v14 & \texttt{R3} & \texttt{opcode1=SHIFT\_BITWISE\_V4I8} \\
\texttt{(none)} & v9--v14 & \texttt{R3} & \texttt{opcode1=LOGIC\_SHIFT\_V4I8} \\
\bottomrule
\end{tabularx}
\end{center}
\begin{center}\footnotesize
\begin{tabularx}{0.88\linewidth}{@{}l l X@{}}
\toprule
Field Name & Bits & Type / Fixed Value \\
\midrule
\texttt{src0} & \texttt{0:7} & \texttt{SourceEncoding} \\
\texttt{src1} & \texttt{8:15} & \texttt{SourceEncoding} \\
\texttt{src2} & \texttt{16:23} & \texttt{SourceEncoding} \\
\texttt{bitwise\_operation} & \texttt{24:25} & \texttt{exact=or} \\
\hyperref[enum:shift_bitwise_i32_lanes_m]{\texttt{lanes1}} & \texttt{26:29} & \texttt{shift\_bitwise\_i32\_lanes\_m (def b0)} \\
\hyperref[enum:not_m]{\texttt{not\_result}} & \texttt{30} & \texttt{not\_m (def none)} \\
\textit{Reserved} & \texttt{31} & \texttt{zero} \\
\texttt{shift\_operation} & \texttt{32:34} & \texttt{exact=lshift} \\
\hyperref[enum:not_m]{\texttt{not2}} & \texttt{35} & \texttt{not\_m (def none)} \\
\hyperref[enum:widen_int_m]{\texttt{widen0}} & \texttt{36:39} & \texttt{widen\_int\_m (def none)} \\
\texttt{dest\_register\_index} & \texttt{40:45} & \texttt{uint} \\
\hyperref[enum:dest_width_m]{\texttt{dest\_width}} & \texttt{46:47} & \texttt{dest\_width\_m} \\
\texttt{opcode1} & \texttt{48:51} & \texttt{exact=SHIFT\_BITWISE\_I32} \\
\texttt{instruction\_class} & \texttt{52:56} & \texttt{exact=R3} \\
\texttt{fau\_page\_index} & \texttt{57:58} & \texttt{uint} \\
\hyperref[enum:flow_control_m]{\texttt{flow}} & \texttt{59:62} & \texttt{flow\_control\_m} \\
\textit{Reserved} & \texttt{63} & \texttt{zero} \\
\bottomrule
\end{tabularx}
\end{center}
\textbf{Fields Explained:}
\begin{tightitemize}\footnotesize
  \item \texttt{src0}: Primary source operand (\texttt{r0..r63}, \texttt{u0..u63} uniform, \texttt{k*} FAU constant, or \texttt{t0..t2} temporary; \texttt{\textasciicircum{}} marks discard/last-use). \emph{Example:} {\scriptsize\ttfamily \mnem{FADD.f32} \dst{r0}, \src{r1}, \src{r2}} (\textbf{SPIR-V:} \texttt{OpFAdd (Operand 0)}; \textbf{GLSL:} \texttt{float c = a + b;})
  \item \texttt{src1}: Secondary source operand (register, uniform, or FAU constant; in shifts, supplies the shift count byte). \emph{Example:} {\scriptsize\ttfamily \mnem{LSHIFT.i32} \dst{r0}, \src{r1}, \src{r2}\imod{.b0}} (\textbf{SPIR-V:} \texttt{OpShiftLeftLogical (Shift)}; \textbf{GLSL:} \texttt{uint y = a \textless{}\textless{} s;})
  \item \texttt{src2}: Third source operand (FMA addend, \texttt{MUX}/\texttt{CSEL} selector, \texttt{IADDC} carry-in, or bitwise mask in fused shift-logic). \emph{Example:} {\scriptsize\ttfamily \mnem{FMA.f32} \dst{r0}, \src{r1}, \src{r2}, \src{r3}} (\textbf{SPIR-V:} \texttt{GLSL.std.450 Fma (Addend)}; \textbf{GLSL:} \texttt{float d = fma(a, b, c);})
  \item \texttt{bitwise\_operation}: Selects the bitwise ALU op (\texttt{and}, \texttt{or}, \texttt{xor}, \texttt{lshift\_and}, etc.) fused with shifts. \emph{Example:} {\scriptsize\ttfamily \mnem{LSHIFT\_OR.i32} \dst{r0}, \src{r1}, \src{r2}, \src{r3}} (\textbf{SPIR-V:} \texttt{OpShiftLeftLogical + OpBitwiseOr}; \textbf{GLSL:} \texttt{uint y = (a \textless{}\textless{} s) \textbar{} b;})
  \item \hyperref[enum:shift_bitwise_i32_lanes_m]{\texttt{lanes1}}: Stage A byte-lane selector on \texttt{src1} (selects shift-count byte \texttt{b0..b3} on 32/64-bit shifts, \texttt{bXY} on \texttt{v2i16} shifts, or 4-byte swizzle on \texttt{v4*8}). \emph{Example:} {\scriptsize\ttfamily \mnem{LSHIFT.i32} \dst{r0}, \src{r1}, \src{r2}\imod{.b1}} (\textbf{SPIR-V:} \texttt{OpBitFieldUExtract + OpShiftLeftLogical}; \textbf{GLSL:} \texttt{uint y = a \textless{}\textless{} ((s \textgreater{}\textgreater{} 8) \& 0xFFu);})
  \item \hyperref[enum:not_m]{\texttt{not\_result}}: Bitwise-inverts the final output of a bitwise/shift-bitwise instruction (\textasciitilde{}result). \emph{Example:} {\scriptsize\ttfamily \mnem{AND.i32}\imod{.not\_result} \dst{r0}, \src{r1}, \src{r2}} (\textbf{SPIR-V:} \texttt{OpBitwiseAnd + OpNot}; \textbf{GLSL:} \texttt{uint y = \textasciitilde{}(a \& b);})
  \item \texttt{shift\_operation}: Selects shift direction/type (\texttt{lshift}, \texttt{rshift}, \texttt{arshift}) in fused shift-bitwise instructions. \emph{Example:} {\scriptsize\ttfamily \mnem{RSHIFT\_AND.i32} \dst{r0}, \src{r1}, \src{r2}, \src{r3}} (\textbf{SPIR-V:} \texttt{OpShiftRightLogical + OpBitwiseAnd}; \textbf{GLSL:} \texttt{uint y = (a \textgreater{}\textgreater{} s) \& m;})
  \item \hyperref[enum:not_m]{\texttt{not2}}: Stage A bitwise complement (\texttt{\textasciitilde{}src2}) applied to \texttt{src2} before the bitwise/shift-logic operation. \emph{Example:} {\scriptsize\ttfamily \mnem{AND.i32} \dst{r0}, \src{r1}, \src{r2}\imod{.not}} (\textbf{SPIR-V:} \texttt{OpNot + OpBitwiseAnd}; \textbf{GLSL:} \texttt{uint y = a \& \textasciitilde{}b;})
  \item \hyperref[enum:widen_int_m]{\texttt{widen0}}: Stage A widening extractor on \texttt{src0} (promotes FP16 \texttt{h0}/\texttt{h1} to FP32, or sign/zero-extends \texttt{b0..b3}/\texttt{h0..h1}/\texttt{w0..w1} to 32/64-bit before \texttt{abs}/\texttt{neg}). \emph{Example:} {\scriptsize\ttfamily \mnem{IADD.s32} \dst{r0}, \src{r1}\imod{.h0}, \src{r2}} (\textbf{SPIR-V:} \texttt{OpSConvert / OpFConvert}; \textbf{GLSL:} \texttt{int c = int(int16\_t(a)) + b;})
  \item \texttt{dest\_register\_index}: 6-bit destination GPR index (\texttt{r0..r63}) receiving the Stage C writeback result (even-aligned pair \texttt{[rN:rN+1]} for 64-bit ops). \emph{Example:} {\scriptsize\ttfamily \mnem{FMA.f32} \dst{r0}, \src{r1}, \src{r2}, \src{r3}} (\textbf{SPIR-V:} \texttt{SSA Result ID}; \textbf{GLSL:} \texttt{float r0 = fma(r1, r2, r3);})
  \item \hyperref[enum:dest_width_m]{\texttt{dest\_width}}: Stage C destination lane write-mask (\texttt{none} = full 32-bit, \texttt{h0} = low 16-bit, \texttt{h1} = high 16-bit without bit-shifting). \emph{Example:} {\scriptsize\ttfamily \mnem{FADD.v2f16} \dst{r0}\imod{.h1}, \src{r1}, \src{r2}} (\textbf{SPIR-V:} \texttt{OpVectorInsertDynamic / WriteMask}; \textbf{GLSL:} \texttt{dst.y = a.y + b.y;})
  \item \texttt{opcode1}: Primary sub-opcode within the instruction class uniquely identifying the operation family. \emph{Example:} {\scriptsize\ttfamily \mnem{IADD.s32} \dst{r0}, \src{r1}, \src{r2}} (\textbf{SPIR-V:} \texttt{OpIAdd (operation selector)}; \textbf{GLSL:} \texttt{// Primary opcode bits})
  \item \texttt{instruction\_class}: 4-bit major instruction class selector (\texttt{bits 48:51}) partitioning the 64-bit Valhall opcode space (e.g., \texttt{FMA}, \texttt{CVT}, \texttt{SFU}, \texttt{LS}, \texttt{TEX}, \texttt{BRANCH}). \emph{Example:} {\scriptsize\ttfamily \mnem{FADD.f32} \dst{r0}, \src{r1}, \src{r2}} (\textbf{SPIR-V:} \texttt{N/A (encoding class)}; \textbf{GLSL:} \texttt{// Hardware execution pipe / decode class})
  \item \texttt{fau\_page\_index}: Selects the FAU (Fast Access Uniform) page (\texttt{0..3}) when a source operand references hardware small constants, thread IDs, or lane constants. \emph{Example:} {\scriptsize\ttfamily \mnem{FADD.f32} \dst{r0}, \src{r1}, \src{k1f}} (\textbf{SPIR-V:} \texttt{OpConstant / BuiltIn}; \textbf{GLSL:} \texttt{float y = x + 1.0;})
  \item \hyperref[enum:flow_control_m]{\texttt{flow}}: 4-bit embedded flow-control \& scoreboard wait modifier (\texttt{bits 59:62}: \texttt{none}, \texttt{wait0..wait0126}, \texttt{wait}, \texttt{reconverge}, \texttt{discard}, \texttt{end}). \emph{Example:} {\scriptsize\ttfamily \mnem{FADD.f32}\imod{.wait0.reconverge} \dst{r0}, \src{r1}, \src{r2}} (\textbf{SPIR-V:} \texttt{OpControlBarrier / OpReturn}; \textbf{GLSL:} \texttt{// Scoreboard wait + warp reconvergence})
\end{tightitemize}
\clearpage

\subsection*{LSHIFT\_XOR}
\addcontentsline{toc}{subsection}{LSHIFT\_XOR}
For each SIMD element $i \in \{0, \dots, N-1\}$ ($N=1$ for \texttt{i32}/\texttt{i64}, $N=2$ for \texttt{v2i16}, $N=4$ for \texttt{v4i8} on v9--v10, lane width $W$), computes logical left shift \texttt{s = (val \textless{}\textless{} (amt \& (W - 1)))} where \texttt{val = src0[i]} and the 8-bit shift count \texttt{amt} is extracted from \texttt{src1} via \texttt{lanes1}, then combines with \texttt{src2[i]} via bitwise XOR \texttt{res = s \textasciicircum{} (not2 ? \textasciitilde{}src2[i] : src2[i])} and writes \texttt{not\_result ? \textasciitilde{}res : res} to \texttt{dest[i]}.
\par\textbf{Notes:}
\begin{tightitemize}
  \item Shift amount is masked modulo lane width (\texttt{wrapping\_shl}/\texttt{wrapping\_shr}/\texttt{rotate\_*}, \texttt{kraid/ops.rs:3984-4009}).
  \item On \texttt{v2i16}, \texttt{src1} is \texttt{V2I8} (\texttt{ops.rs:3953}) with byte lanes selected by \texttt{lanes1} (default \texttt{b02}: byte 0 -\textgreater{} lane 0, byte 2 -\textgreater{} lane 1).
  \item \textbf{Source references:} \texttt{kraid/ops.rs:3878-4025}, \texttt{encode\_v9.rs:2940-3093}, \texttt{hw\_tests.rs:2148-2198}.
\end{tightitemize}
\codebox{\mnem{LSHIFT\_XOR}\imod{\{.flow\}}\imod{\{.not\_result\}} \dst{dest}\imod{\{.dest\_width\}}, \src{src0}\imod{\{.widen0\}}, \src{src1}\imod{\{.lanes1\}}, \src{src2}\imod{\{.not2\}}\\\mnem{LSHIFT\_XOR}\imod{\{.flow\}}\imod{\{.not\_result\}} \dst{dest}, \src{src0}\imod{\{.widen0\}}, \src{src1}\imod{\{.lanes1\}}, \src{src2}\imod{\{.not2\}}\\\mnem{LSHIFT\_XOR}\imod{\{.flow\}}\imod{\{.not\_result\}} \dst{dest}\imod{\{.dest\_width\}}, \src{src0}\imod{\{.swz0\}}, \src{src1}\imod{\{.lanes1\}}, \src{src2}\imod{\{.not2\}}\\\mnem{LSHIFT\_XOR}\imod{\{.flow\}}\imod{\{.not\_result\}} \dst{dest}\imod{\{.dest\_width\}}, \src{src0}\imod{\{.lanes0\}}, \src{src1}\imod{\{.lanes1\}}, \src{src2}\imod{\{.not2\}}}
\textbf{Operands:}
\begin{tightitemize}\footnotesize
  \item \texttt{\dst{dest}} (\texttt{dest\_register\_index}, \texttt{dest\_width}): Destination GPR receiving the logical left shift combined via bitwise XOR (inverted if \texttt{.not\_result}).
  \item \texttt{\src{src0}}: Input integer value to shift/rotate.
  \item \texttt{\src{src1}}: 8-bit shift/rotate count per lane (selected via \texttt{lanes1}, masked modulo lane width \texttt{W}).
  \item \texttt{\src{src2}}: Second operand combined via bitwise XOR after shifting (inverted if \texttt{.not2}).
\end{tightitemize}
\par\textbf{Closest SPIR-V Op:} \texttt{OpShiftLeftLogical + OpBitwiseXor}
\par\textbf{Pseudo-GLSL (Operational Semantics):}
\codebox{for (int i = 0; i \textless{} N; ++i) \{\\    uint s = extractByte(src1, lanes1, i) \& (W - 1);\\    auto r = (src0[i] \textless{}\textless{} s) \textasciicircum{} (flag\_not2 ? \textasciitilde{}src2[i] : src2[i]);\\    dest[i] = flag\_not\_result ? \textasciitilde{}r : r;\\\}}
\textbf{Execution Context:}
\begin{tightitemize}\footnotesize
  \item \textbf{Compute stage preloads (\hyperref[sec:csf-spd]{\texttt{SPD\_k}}, \texttt{r16..r23}):} Hardware pre-fills \texttt{r55..r57} with \texttt{WorkGroupID.xyz + JOB\_OFFSET} (from \hyperref[sec:csf-regs]{\texttt{COMPUTESR::JOB\_OFFSET\_X/Y/Z}} \texttt{r34..r36}) and \texttt{r60..r62} with \texttt{LocalInvocationID.xyz} or \texttt{GlobalInvocationID.xyz} (from \texttt{WG\_SIZE} \texttt{r33} \& \texttt{JOB\_SIZE} \texttt{r37..r39}, selected by \texttt{ShaderProgram::preload}) so compute shaders can derive thread/element indices without memory loads.
  \item \textbf{Vertex stage preloads (\hyperref[sec:csf-spd]{\texttt{SPD\_0}/\texttt{SPD\_1}}, \texttt{r16..r19}):} The hardware Index Fetcher reads \hyperref[sec:csf-regs]{\texttt{IDVSSR::INDEX\_BUFFER}} (\texttt{r54:r55}) and pre-fills \texttt{r59, r60} with \texttt{VertexIndex} (\texttt{+ VERTEX\_OFFSET + GLOBAL\_ATTR\_OFF}) and \texttt{r61} with \texttt{InstanceIndex} (\texttt{+ INSTANCE\_OFFSET}) so Position and Varying shaders can pull vertex attributes directly.
  \item \textbf{Inline FAU \& Special Operands (\hyperref[sec:csf-tsd-fau]{\texttt{FAU\_k}}, \texttt{r8..r15}):} Any 8-bit source with \texttt{mode=2} reads zero-latency uniform \texttt{u0..u31} from \texttt{FAU\_k} (\texttt{slots64 = bits[63:56]}); \texttt{mode=3} reads \texttt{SmallConstant} (Table~1.2) or \texttt{FauSpecial} (\texttt{pc}, \texttt{lane\_id}, \texttt{lane\_width}, \texttt{tls\_ptr}, \texttt{wls\_ptr}, \texttt{blend\_descriptor\_0..7}).
\end{tightitemize}
\par\textbf{Worked Examples:}
\par\textbf{Example 1 (Standard fused shift + bitwise op):}
\codebox{\mnem{LSHIFT\_XOR.i32} \dst{r0}, \src{r1}, \src{r2}\imod{.b0}, \src{r3}}
With \texttt{r1 = 0x00000010}, \texttt{r2.b0 = 4} (masked mod 32): shift yields \texttt{0x00000100}, then \texttt{\textasciicircum{} 0x000000FF} yields \texttt{r0 = 0x1FF}.
\par\textbf{Example 2 (Modulo-32 shift count \& .not\_result):}
\codebox{\mnem{LSHIFT\_XOR.i32}\imod{.not\_result} \dst{r0}, \src{r1}, \src{r2}\imod{.b0}, \src{r3}}
With \texttt{r1 = 0x00000001}, \texttt{r2.b0 = 33} (effective shift \texttt{1}): shift yields \texttt{0x00000002}, \texttt{\textasciicircum{} 0x000000FF} yields \texttt{0x000000FD}, and \texttt{.not\_result} inverts -\textgreater{} \texttt{r0 = 0xFFFFFF02}.
\par\textbf{Mined SPIR-V $\to$ Valhall ISA Example (Inverting a boolean mask via NOT alias --- Fragment shader, Mali-G710 (v10)):}
\codebox{\textbf{; SPIR-V (input)}\\\%keep = OpLogicalNot \%bool \%cond\\OpBranchConditional \%keep \%kill\_blk \%cont\_blk\\[3pt]\textbf{; Valhall ISA (malioc output)}\\NOT.i32 r1, r0\\DISCARD.f32.wait0126.ne r1\textasciicircum{}, k0}
Because \texttt{.m1} booleans use all-ones (\texttt{0xFFFFFFFF}) for true and \texttt{0} for false, \texttt{OpLogicalNot} lowers to \texttt{NOT.i32} (encoded via shift-logic with \texttt{.not\_result}) before testing against \texttt{k0} in \texttt{DISCARD}.
\begin{center}\footnotesize
\begin{tabularx}{\linewidth}{@{}lllX@{}}
\toprule
Variant & Arch & Class & Opcode Fields \\
\midrule
\texttt{(none)} & v9--v14 & \texttt{R3} & \texttt{opcode1=SHIFT\_BITWISE\_I32} \\
\texttt{(none)} & v9--v14 & \texttt{R3} & \texttt{opcode1=SHIFT\_BITWISE\_I32} \\
\texttt{(none)} & v9--v14 & \texttt{R3} & \texttt{opcode1=LOGIC\_SHIFT\_I32} \\
\texttt{(none)} & v9--v14 & \texttt{D3} & \texttt{opcode1=SHIFT\_BITWISE\_I64} \\
\texttt{(none)} & v9--v14 & \texttt{D3} & \texttt{opcode1=LOGIC\_SHIFT\_I64} \\
\texttt{(none)} & v9--v14 & \texttt{R3} & \texttt{opcode1=SHIFT\_BITWISE\_V2I16} \\
\texttt{(none)} & v9--v14 & \texttt{R3} & \texttt{opcode1=LOGIC\_SHIFT\_V2I16} \\
\texttt{(none)} & v9--v14 & \texttt{R3} & \texttt{opcode1=SHIFT\_BITWISE\_V4I8} \\
\texttt{(none)} & v9--v14 & \texttt{R3} & \texttt{opcode1=LOGIC\_SHIFT\_V4I8} \\
\bottomrule
\end{tabularx}
\end{center}
\begin{center}\footnotesize
\begin{tabularx}{0.88\linewidth}{@{}l l X@{}}
\toprule
Field Name & Bits & Type / Fixed Value \\
\midrule
\texttt{src0} & \texttt{0:7} & \texttt{SourceEncoding} \\
\texttt{src1} & \texttt{8:15} & \texttt{SourceEncoding} \\
\texttt{src2} & \texttt{16:23} & \texttt{SourceEncoding} \\
\texttt{bitwise\_operation} & \texttt{24:25} & \texttt{exact=xor} \\
\hyperref[enum:shift_bitwise_i32_lanes_m]{\texttt{lanes1}} & \texttt{26:29} & \texttt{shift\_bitwise\_i32\_lanes\_m (def b0)} \\
\hyperref[enum:not_m]{\texttt{not\_result}} & \texttt{30} & \texttt{not\_m (def none)} \\
\textit{Reserved} & \texttt{31} & \texttt{zero} \\
\texttt{shift\_operation} & \texttt{32:34} & \texttt{exact=lshift} \\
\hyperref[enum:not_m]{\texttt{not2}} & \texttt{35} & \texttt{not\_m (def none)} \\
\hyperref[enum:widen_int_m]{\texttt{widen0}} & \texttt{36:39} & \texttt{widen\_int\_m (def none)} \\
\texttt{dest\_register\_index} & \texttt{40:45} & \texttt{uint} \\
\hyperref[enum:dest_width_m]{\texttt{dest\_width}} & \texttt{46:47} & \texttt{dest\_width\_m} \\
\texttt{opcode1} & \texttt{48:51} & \texttt{exact=SHIFT\_BITWISE\_I32} \\
\texttt{instruction\_class} & \texttt{52:56} & \texttt{exact=R3} \\
\texttt{fau\_page\_index} & \texttt{57:58} & \texttt{uint} \\
\hyperref[enum:flow_control_m]{\texttt{flow}} & \texttt{59:62} & \texttt{flow\_control\_m} \\
\textit{Reserved} & \texttt{63} & \texttt{zero} \\
\bottomrule
\end{tabularx}
\end{center}
\textbf{Fields Explained:}
\begin{tightitemize}\footnotesize
  \item \texttt{src0}: Primary source operand (\texttt{r0..r63}, \texttt{u0..u63} uniform, \texttt{k*} FAU constant, or \texttt{t0..t2} temporary; \texttt{\textasciicircum{}} marks discard/last-use). \emph{Example:} {\scriptsize\ttfamily \mnem{FADD.f32} \dst{r0}, \src{r1}, \src{r2}} (\textbf{SPIR-V:} \texttt{OpFAdd (Operand 0)}; \textbf{GLSL:} \texttt{float c = a + b;})
  \item \texttt{src1}: Secondary source operand (register, uniform, or FAU constant; in shifts, supplies the shift count byte). \emph{Example:} {\scriptsize\ttfamily \mnem{LSHIFT.i32} \dst{r0}, \src{r1}, \src{r2}\imod{.b0}} (\textbf{SPIR-V:} \texttt{OpShiftLeftLogical (Shift)}; \textbf{GLSL:} \texttt{uint y = a \textless{}\textless{} s;})
  \item \texttt{src2}: Third source operand (FMA addend, \texttt{MUX}/\texttt{CSEL} selector, \texttt{IADDC} carry-in, or bitwise mask in fused shift-logic). \emph{Example:} {\scriptsize\ttfamily \mnem{FMA.f32} \dst{r0}, \src{r1}, \src{r2}, \src{r3}} (\textbf{SPIR-V:} \texttt{GLSL.std.450 Fma (Addend)}; \textbf{GLSL:} \texttt{float d = fma(a, b, c);})
  \item \texttt{bitwise\_operation}: Selects the bitwise ALU op (\texttt{and}, \texttt{or}, \texttt{xor}, \texttt{lshift\_and}, etc.) fused with shifts. \emph{Example:} {\scriptsize\ttfamily \mnem{LSHIFT\_OR.i32} \dst{r0}, \src{r1}, \src{r2}, \src{r3}} (\textbf{SPIR-V:} \texttt{OpShiftLeftLogical + OpBitwiseOr}; \textbf{GLSL:} \texttt{uint y = (a \textless{}\textless{} s) \textbar{} b;})
  \item \hyperref[enum:shift_bitwise_i32_lanes_m]{\texttt{lanes1}}: Stage A byte-lane selector on \texttt{src1} (selects shift-count byte \texttt{b0..b3} on 32/64-bit shifts, \texttt{bXY} on \texttt{v2i16} shifts, or 4-byte swizzle on \texttt{v4*8}). \emph{Example:} {\scriptsize\ttfamily \mnem{LSHIFT.i32} \dst{r0}, \src{r1}, \src{r2}\imod{.b1}} (\textbf{SPIR-V:} \texttt{OpBitFieldUExtract + OpShiftLeftLogical}; \textbf{GLSL:} \texttt{uint y = a \textless{}\textless{} ((s \textgreater{}\textgreater{} 8) \& 0xFFu);})
  \item \hyperref[enum:not_m]{\texttt{not\_result}}: Bitwise-inverts the final output of a bitwise/shift-bitwise instruction (\textasciitilde{}result). \emph{Example:} {\scriptsize\ttfamily \mnem{AND.i32}\imod{.not\_result} \dst{r0}, \src{r1}, \src{r2}} (\textbf{SPIR-V:} \texttt{OpBitwiseAnd + OpNot}; \textbf{GLSL:} \texttt{uint y = \textasciitilde{}(a \& b);})
  \item \texttt{shift\_operation}: Selects shift direction/type (\texttt{lshift}, \texttt{rshift}, \texttt{arshift}) in fused shift-bitwise instructions. \emph{Example:} {\scriptsize\ttfamily \mnem{RSHIFT\_AND.i32} \dst{r0}, \src{r1}, \src{r2}, \src{r3}} (\textbf{SPIR-V:} \texttt{OpShiftRightLogical + OpBitwiseAnd}; \textbf{GLSL:} \texttt{uint y = (a \textgreater{}\textgreater{} s) \& m;})
  \item \hyperref[enum:not_m]{\texttt{not2}}: Stage A bitwise complement (\texttt{\textasciitilde{}src2}) applied to \texttt{src2} before the bitwise/shift-logic operation. \emph{Example:} {\scriptsize\ttfamily \mnem{AND.i32} \dst{r0}, \src{r1}, \src{r2}\imod{.not}} (\textbf{SPIR-V:} \texttt{OpNot + OpBitwiseAnd}; \textbf{GLSL:} \texttt{uint y = a \& \textasciitilde{}b;})
  \item \hyperref[enum:widen_int_m]{\texttt{widen0}}: Stage A widening extractor on \texttt{src0} (promotes FP16 \texttt{h0}/\texttt{h1} to FP32, or sign/zero-extends \texttt{b0..b3}/\texttt{h0..h1}/\texttt{w0..w1} to 32/64-bit before \texttt{abs}/\texttt{neg}). \emph{Example:} {\scriptsize\ttfamily \mnem{IADD.s32} \dst{r0}, \src{r1}\imod{.h0}, \src{r2}} (\textbf{SPIR-V:} \texttt{OpSConvert / OpFConvert}; \textbf{GLSL:} \texttt{int c = int(int16\_t(a)) + b;})
  \item \texttt{dest\_register\_index}: 6-bit destination GPR index (\texttt{r0..r63}) receiving the Stage C writeback result (even-aligned pair \texttt{[rN:rN+1]} for 64-bit ops). \emph{Example:} {\scriptsize\ttfamily \mnem{FMA.f32} \dst{r0}, \src{r1}, \src{r2}, \src{r3}} (\textbf{SPIR-V:} \texttt{SSA Result ID}; \textbf{GLSL:} \texttt{float r0 = fma(r1, r2, r3);})
  \item \hyperref[enum:dest_width_m]{\texttt{dest\_width}}: Stage C destination lane write-mask (\texttt{none} = full 32-bit, \texttt{h0} = low 16-bit, \texttt{h1} = high 16-bit without bit-shifting). \emph{Example:} {\scriptsize\ttfamily \mnem{FADD.v2f16} \dst{r0}\imod{.h1}, \src{r1}, \src{r2}} (\textbf{SPIR-V:} \texttt{OpVectorInsertDynamic / WriteMask}; \textbf{GLSL:} \texttt{dst.y = a.y + b.y;})
  \item \texttt{opcode1}: Primary sub-opcode within the instruction class uniquely identifying the operation family. \emph{Example:} {\scriptsize\ttfamily \mnem{IADD.s32} \dst{r0}, \src{r1}, \src{r2}} (\textbf{SPIR-V:} \texttt{OpIAdd (operation selector)}; \textbf{GLSL:} \texttt{// Primary opcode bits})
  \item \texttt{instruction\_class}: 4-bit major instruction class selector (\texttt{bits 48:51}) partitioning the 64-bit Valhall opcode space (e.g., \texttt{FMA}, \texttt{CVT}, \texttt{SFU}, \texttt{LS}, \texttt{TEX}, \texttt{BRANCH}). \emph{Example:} {\scriptsize\ttfamily \mnem{FADD.f32} \dst{r0}, \src{r1}, \src{r2}} (\textbf{SPIR-V:} \texttt{N/A (encoding class)}; \textbf{GLSL:} \texttt{// Hardware execution pipe / decode class})
  \item \texttt{fau\_page\_index}: Selects the FAU (Fast Access Uniform) page (\texttt{0..3}) when a source operand references hardware small constants, thread IDs, or lane constants. \emph{Example:} {\scriptsize\ttfamily \mnem{FADD.f32} \dst{r0}, \src{r1}, \src{k1f}} (\textbf{SPIR-V:} \texttt{OpConstant / BuiltIn}; \textbf{GLSL:} \texttt{float y = x + 1.0;})
  \item \hyperref[enum:flow_control_m]{\texttt{flow}}: 4-bit embedded flow-control \& scoreboard wait modifier (\texttt{bits 59:62}: \texttt{none}, \texttt{wait0..wait0126}, \texttt{wait}, \texttt{reconverge}, \texttt{discard}, \texttt{end}). \emph{Example:} {\scriptsize\ttfamily \mnem{FADD.f32}\imod{.wait0.reconverge} \dst{r0}, \src{r1}, \src{r2}} (\textbf{SPIR-V:} \texttt{OpControlBarrier / OpReturn}; \textbf{GLSL:} \texttt{// Scoreboard wait + warp reconvergence})
\end{tightitemize}
\clearpage

\subsection*{NOT}
\addcontentsline{toc}{subsection}{NOT}
Computes the bitwise complement of \texttt{src0} (\texttt{dest = \textasciitilde{}widen(src0)}) across all lanes. Encoded as an alias of \texttt{OR} with \texttt{src2 = 0} and \texttt{.not\_result} enabled.
\par\textbf{Notes:}
\begin{tightitemize}
  \item Bitwise NOT across all bits of destination width.
  \item \textbf{Source references:} \texttt{kraid/nir.rs:1160-1171}, \texttt{encode\_v9.rs:3060-3066}.
\end{tightitemize}
\codebox{\mnem{NOT}\imod{\{.flow\}} \dst{dest}\imod{\{.dest\_width\}}, \src{src0}\imod{\{.widen0\}}\\\mnem{NOT}\imod{\{.flow\}} \dst{dest}, \src{src0}\imod{\{.widen0\}}\\\mnem{NOT}\imod{\{.flow\}} \dst{dest}\imod{\{.dest\_width\}}, \src{src0}\imod{\{.swz0\}}\\\mnem{NOT}\imod{\{.flow\}} \dst{dest}\imod{\{.dest\_width\}}, \src{src0}\imod{\{.lanes0\}}}
\textbf{Operands:}
\begin{tightitemize}\footnotesize
  \item \texttt{\dst{dest}} (\texttt{dest\_register\_index}, \texttt{dest\_width}): Destination GPR receiving \texttt{\textasciitilde{}src0}.
  \item \texttt{\src{src0}}: Input operand to invert bitwise.
\end{tightitemize}
\par\textbf{Closest SPIR-V Op:} \texttt{OpNot}
\par\textbf{Pseudo-GLSL (Operational Semantics):}
\codebox{dest = \textasciitilde{}widen0(src0); // alias of OR(src0, 0) with .not\_result}
\textbf{Execution Context:}
\begin{tightitemize}\footnotesize
  \item \textbf{Compute stage preloads (\hyperref[sec:csf-spd]{\texttt{SPD\_k}}, \texttt{r16..r23}):} Hardware pre-fills \texttt{r55..r57} with \texttt{WorkGroupID.xyz + JOB\_OFFSET} (from \hyperref[sec:csf-regs]{\texttt{COMPUTESR::JOB\_OFFSET\_X/Y/Z}} \texttt{r34..r36}) and \texttt{r60..r62} with \texttt{LocalInvocationID.xyz} or \texttt{GlobalInvocationID.xyz} (from \texttt{WG\_SIZE} \texttt{r33} \& \texttt{JOB\_SIZE} \texttt{r37..r39}, selected by \texttt{ShaderProgram::preload}) so compute shaders can derive thread/element indices without memory loads.
  \item \textbf{Vertex stage preloads (\hyperref[sec:csf-spd]{\texttt{SPD\_0}/\texttt{SPD\_1}}, \texttt{r16..r19}):} The hardware Index Fetcher reads \hyperref[sec:csf-regs]{\texttt{IDVSSR::INDEX\_BUFFER}} (\texttt{r54:r55}) and pre-fills \texttt{r59, r60} with \texttt{VertexIndex} (\texttt{+ VERTEX\_OFFSET + GLOBAL\_ATTR\_OFF}) and \texttt{r61} with \texttt{InstanceIndex} (\texttt{+ INSTANCE\_OFFSET}) so Position and Varying shaders can pull vertex attributes directly.
  \item \textbf{Inline FAU \& Special Operands (\hyperref[sec:csf-tsd-fau]{\texttt{FAU\_k}}, \texttt{r8..r15}):} Any 8-bit source with \texttt{mode=2} reads zero-latency uniform \texttt{u0..u31} from \texttt{FAU\_k} (\texttt{slots64 = bits[63:56]}); \texttt{mode=3} reads \texttt{SmallConstant} (Table~1.2) or \texttt{FauSpecial} (\texttt{pc}, \texttt{lane\_id}, \texttt{lane\_width}, \texttt{tls\_ptr}, \texttt{wls\_ptr}, \texttt{blend\_descriptor\_0..7}).
\end{tightitemize}
\par\textbf{Worked Examples:}
\par\textbf{Example 1 (32-bit bitwise complement):}
\codebox{\mnem{NOT.i32} \dst{r0}, \src{r1}}
With \texttt{r1 = 0x0000FFFF}: flips every bit -\textgreater{} \texttt{r0 = 0xFFFF0000}.
\par\textbf{Example 2 (Zero and all-ones complement):}
\codebox{\mnem{NOT.i32} \dst{r0}, \src{r1}}
With \texttt{r1 = 0x00000000} -\textgreater{} \texttt{r0 = 0xFFFFFFFF}; with \texttt{r1 = 0xFFFFFFFF} -\textgreater{} \texttt{r0 = 0x00000000}.

\begin{center}\footnotesize
\begin{tabularx}{\linewidth}{@{}lllX@{}}
\toprule
Variant & Arch & Class & Opcode Fields \\
\midrule
\texttt{(none)} & v9--v14 & \texttt{R3} & \texttt{opcode1=SHIFT\_BITWISE\_I32} \\
\texttt{(none)} & v9--v14 & \texttt{R3} & \texttt{opcode1=SHIFT\_BITWISE\_I32} \\
\texttt{(none)} & v9--v14 & \texttt{R3} & \texttt{opcode1=LOGIC\_SHIFT\_I32} \\
\texttt{(none)} & v9--v14 & \texttt{D3} & \texttt{opcode1=SHIFT\_BITWISE\_I64} \\
\texttt{(none)} & v9--v14 & \texttt{D3} & \texttt{opcode1=LOGIC\_SHIFT\_I64} \\
\texttt{(none)} & v9--v14 & \texttt{R3} & \texttt{opcode1=SHIFT\_BITWISE\_V2I16} \\
\texttt{(none)} & v9--v14 & \texttt{R3} & \texttt{opcode1=LOGIC\_SHIFT\_V2I16} \\
\texttt{(none)} & v9--v14 & \texttt{R3} & \texttt{opcode1=SHIFT\_BITWISE\_V4I8} \\
\texttt{(none)} & v9--v14 & \texttt{R3} & \texttt{opcode1=LOGIC\_SHIFT\_V4I8} \\
\bottomrule
\end{tabularx}
\end{center}
\begin{center}\footnotesize
\begin{tabularx}{0.88\linewidth}{@{}l l X@{}}
\toprule
Field Name & Bits & Type / Fixed Value \\
\midrule
\texttt{src0} & \texttt{0:7} & \texttt{SourceEncoding} \\
\texttt{src1} & \texttt{8:15} & \texttt{exact=192} \\
\texttt{src2} & \texttt{16:23} & \texttt{exact=192} \\
\texttt{bitwise\_operation} & \texttt{24:25} & \texttt{exact=xor} \\
\hyperref[enum:shift_bitwise_i32_lanes_m]{\texttt{lanes1}} & \texttt{26:29} & \texttt{exact=b0} \\
\hyperref[enum:not_m]{\texttt{not\_result}} & \texttt{30} & \texttt{exact=none} \\
\textit{Reserved} & \texttt{31} & \texttt{zero} \\
\texttt{shift\_operation} & \texttt{32:34} & \texttt{exact=lshift} \\
\hyperref[enum:not_m]{\texttt{not2}} & \texttt{35} & \texttt{exact=not} \\
\hyperref[enum:widen_int_m]{\texttt{widen0}} & \texttt{36:39} & \texttt{widen\_int\_m (def none)} \\
\texttt{dest\_register\_index} & \texttt{40:45} & \texttt{uint} \\
\hyperref[enum:dest_width_m]{\texttt{dest\_width}} & \texttt{46:47} & \texttt{dest\_width\_m} \\
\texttt{opcode1} & \texttt{48:51} & \texttt{exact=SHIFT\_BITWISE\_I32} \\
\texttt{instruction\_class} & \texttt{52:56} & \texttt{exact=R3} \\
\texttt{fau\_page\_index} & \texttt{57:58} & \texttt{uint} \\
\hyperref[enum:flow_control_m]{\texttt{flow}} & \texttt{59:62} & \texttt{flow\_control\_m} \\
\textit{Reserved} & \texttt{63} & \texttt{zero} \\
\bottomrule
\end{tabularx}
\end{center}
\textbf{Fields Explained:}
\begin{tightitemize}\footnotesize
  \item \texttt{src0}: Primary source operand (\texttt{r0..r63}, \texttt{u0..u63} uniform, \texttt{k*} FAU constant, or \texttt{t0..t2} temporary; \texttt{\textasciicircum{}} marks discard/last-use). \emph{Example:} {\scriptsize\ttfamily \mnem{FADD.f32} \dst{r0}, \src{r1}, \src{r2}} (\textbf{SPIR-V:} \texttt{OpFAdd (Operand 0)}; \textbf{GLSL:} \texttt{float c = a + b;})
  \item \texttt{src1}: Secondary source operand (register, uniform, or FAU constant; in shifts, supplies the shift count byte). \emph{Example:} {\scriptsize\ttfamily \mnem{LSHIFT.i32} \dst{r0}, \src{r1}, \src{r2}\imod{.b0}} (\textbf{SPIR-V:} \texttt{OpShiftLeftLogical (Shift)}; \textbf{GLSL:} \texttt{uint y = a \textless{}\textless{} s;})
  \item \texttt{src2}: Third source operand (FMA addend, \texttt{MUX}/\texttt{CSEL} selector, \texttt{IADDC} carry-in, or bitwise mask in fused shift-logic). \emph{Example:} {\scriptsize\ttfamily \mnem{FMA.f32} \dst{r0}, \src{r1}, \src{r2}, \src{r3}} (\textbf{SPIR-V:} \texttt{GLSL.std.450 Fma (Addend)}; \textbf{GLSL:} \texttt{float d = fma(a, b, c);})
  \item \texttt{bitwise\_operation}: Selects the bitwise ALU op (\texttt{and}, \texttt{or}, \texttt{xor}, \texttt{lshift\_and}, etc.) fused with shifts. \emph{Example:} {\scriptsize\ttfamily \mnem{LSHIFT\_OR.i32} \dst{r0}, \src{r1}, \src{r2}, \src{r3}} (\textbf{SPIR-V:} \texttt{OpShiftLeftLogical + OpBitwiseOr}; \textbf{GLSL:} \texttt{uint y = (a \textless{}\textless{} s) \textbar{} b;})
  \item \hyperref[enum:shift_bitwise_i32_lanes_m]{\texttt{lanes1}}: Stage A byte-lane selector on \texttt{src1} (selects shift-count byte \texttt{b0..b3} on 32/64-bit shifts, \texttt{bXY} on \texttt{v2i16} shifts, or 4-byte swizzle on \texttt{v4*8}). \emph{Example:} {\scriptsize\ttfamily \mnem{LSHIFT.i32} \dst{r0}, \src{r1}, \src{r2}\imod{.b1}} (\textbf{SPIR-V:} \texttt{OpBitFieldUExtract + OpShiftLeftLogical}; \textbf{GLSL:} \texttt{uint y = a \textless{}\textless{} ((s \textgreater{}\textgreater{} 8) \& 0xFFu);})
  \item \hyperref[enum:not_m]{\texttt{not\_result}}: Bitwise-inverts the final output of a bitwise/shift-bitwise instruction (\textasciitilde{}result). \emph{Example:} {\scriptsize\ttfamily \mnem{AND.i32}\imod{.not\_result} \dst{r0}, \src{r1}, \src{r2}} (\textbf{SPIR-V:} \texttt{OpBitwiseAnd + OpNot}; \textbf{GLSL:} \texttt{uint y = \textasciitilde{}(a \& b);})
  \item \texttt{shift\_operation}: Selects shift direction/type (\texttt{lshift}, \texttt{rshift}, \texttt{arshift}) in fused shift-bitwise instructions. \emph{Example:} {\scriptsize\ttfamily \mnem{RSHIFT\_AND.i32} \dst{r0}, \src{r1}, \src{r2}, \src{r3}} (\textbf{SPIR-V:} \texttt{OpShiftRightLogical + OpBitwiseAnd}; \textbf{GLSL:} \texttt{uint y = (a \textgreater{}\textgreater{} s) \& m;})
  \item \hyperref[enum:not_m]{\texttt{not2}}: Stage A bitwise complement (\texttt{\textasciitilde{}src2}) applied to \texttt{src2} before the bitwise/shift-logic operation. \emph{Example:} {\scriptsize\ttfamily \mnem{AND.i32} \dst{r0}, \src{r1}, \src{r2}\imod{.not}} (\textbf{SPIR-V:} \texttt{OpNot + OpBitwiseAnd}; \textbf{GLSL:} \texttt{uint y = a \& \textasciitilde{}b;})
  \item \hyperref[enum:widen_int_m]{\texttt{widen0}}: Stage A widening extractor on \texttt{src0} (promotes FP16 \texttt{h0}/\texttt{h1} to FP32, or sign/zero-extends \texttt{b0..b3}/\texttt{h0..h1}/\texttt{w0..w1} to 32/64-bit before \texttt{abs}/\texttt{neg}). \emph{Example:} {\scriptsize\ttfamily \mnem{IADD.s32} \dst{r0}, \src{r1}\imod{.h0}, \src{r2}} (\textbf{SPIR-V:} \texttt{OpSConvert / OpFConvert}; \textbf{GLSL:} \texttt{int c = int(int16\_t(a)) + b;})
  \item \texttt{dest\_register\_index}: 6-bit destination GPR index (\texttt{r0..r63}) receiving the Stage C writeback result (even-aligned pair \texttt{[rN:rN+1]} for 64-bit ops). \emph{Example:} {\scriptsize\ttfamily \mnem{FMA.f32} \dst{r0}, \src{r1}, \src{r2}, \src{r3}} (\textbf{SPIR-V:} \texttt{SSA Result ID}; \textbf{GLSL:} \texttt{float r0 = fma(r1, r2, r3);})
  \item \hyperref[enum:dest_width_m]{\texttt{dest\_width}}: Stage C destination lane write-mask (\texttt{none} = full 32-bit, \texttt{h0} = low 16-bit, \texttt{h1} = high 16-bit without bit-shifting). \emph{Example:} {\scriptsize\ttfamily \mnem{FADD.v2f16} \dst{r0}\imod{.h1}, \src{r1}, \src{r2}} (\textbf{SPIR-V:} \texttt{OpVectorInsertDynamic / WriteMask}; \textbf{GLSL:} \texttt{dst.y = a.y + b.y;})
  \item \texttt{opcode1}: Primary sub-opcode within the instruction class uniquely identifying the operation family. \emph{Example:} {\scriptsize\ttfamily \mnem{IADD.s32} \dst{r0}, \src{r1}, \src{r2}} (\textbf{SPIR-V:} \texttt{OpIAdd (operation selector)}; \textbf{GLSL:} \texttt{// Primary opcode bits})
  \item \texttt{instruction\_class}: 4-bit major instruction class selector (\texttt{bits 48:51}) partitioning the 64-bit Valhall opcode space (e.g., \texttt{FMA}, \texttt{CVT}, \texttt{SFU}, \texttt{LS}, \texttt{TEX}, \texttt{BRANCH}). \emph{Example:} {\scriptsize\ttfamily \mnem{FADD.f32} \dst{r0}, \src{r1}, \src{r2}} (\textbf{SPIR-V:} \texttt{N/A (encoding class)}; \textbf{GLSL:} \texttt{// Hardware execution pipe / decode class})
  \item \texttt{fau\_page\_index}: Selects the FAU (Fast Access Uniform) page (\texttt{0..3}) when a source operand references hardware small constants, thread IDs, or lane constants. \emph{Example:} {\scriptsize\ttfamily \mnem{FADD.f32} \dst{r0}, \src{r1}, \src{k1f}} (\textbf{SPIR-V:} \texttt{OpConstant / BuiltIn}; \textbf{GLSL:} \texttt{float y = x + 1.0;})
  \item \hyperref[enum:flow_control_m]{\texttt{flow}}: 4-bit embedded flow-control \& scoreboard wait modifier (\texttt{bits 59:62}: \texttt{none}, \texttt{wait0..wait0126}, \texttt{wait}, \texttt{reconverge}, \texttt{discard}, \texttt{end}). \emph{Example:} {\scriptsize\ttfamily \mnem{FADD.f32}\imod{.wait0.reconverge} \dst{r0}, \src{r1}, \src{r2}} (\textbf{SPIR-V:} \texttt{OpControlBarrier / OpReturn}; \textbf{GLSL:} \texttt{// Scoreboard wait + warp reconvergence})
\end{tightitemize}
\clearpage

\subsection*{NOT\_OLD}
\addcontentsline{toc}{subsection}{NOT\_OLD}
Legacy v9--v13 unary bitwise complement on \texttt{i32} or \texttt{i64}: inverts every bit of \texttt{src0} and writes \texttt{dest = \textasciitilde{}src0}. Superseded by shift-logic \texttt{NOT} on newer architectures.
\par\textbf{Notes:}
\begin{tightitemize}
  \item Superseded by shift-logic \texttt{NOT} (\texttt{bifrost/valhall/ISA.xml:137}).
  \item \textbf{Source references:} \texttt{bifrost/valhall/ISA.xml:137-138}, \texttt{isa-v9-v14.xml}.
\end{tightitemize}
\codebox{\mnem{NOT\_OLD}\imod{\{.flow\}} \dst{dest}\imod{\{.dest\_width\}}, \src{src0}\\\mnem{NOT\_OLD}\imod{\{.flow\}} \dst{dest}, \src{src0}}
\textbf{Operands:}
\begin{tightitemize}\footnotesize
  \item \texttt{\dst{dest}} (\texttt{dest\_register\_index}, \texttt{dest\_width}): Destination GPR receiving \texttt{\textasciitilde{}src0}.
  \item \texttt{\src{src0}}: Input operand to invert bitwise.
\end{tightitemize}
\par\textbf{Closest SPIR-V Op:} \texttt{OpNot (legacy v9--v13)}
\par\textbf{Pseudo-GLSL (Operational Semantics):}
\codebox{dest = \textasciitilde{}src0; // legacy v9--v13 bitwise complement}
\textbf{Execution Context:}
\begin{tightitemize}\footnotesize
  \item \textbf{Compute stage preloads (\hyperref[sec:csf-spd]{\texttt{SPD\_k}}, \texttt{r16..r23}):} Hardware pre-fills \texttt{r55..r57} with \texttt{WorkGroupID.xyz + JOB\_OFFSET} (from \hyperref[sec:csf-regs]{\texttt{COMPUTESR::JOB\_OFFSET\_X/Y/Z}} \texttt{r34..r36}) and \texttt{r60..r62} with \texttt{LocalInvocationID.xyz} or \texttt{GlobalInvocationID.xyz} (from \texttt{WG\_SIZE} \texttt{r33} \& \texttt{JOB\_SIZE} \texttt{r37..r39}, selected by \texttt{ShaderProgram::preload}) so compute shaders can derive thread/element indices without memory loads.
  \item \textbf{Vertex stage preloads (\hyperref[sec:csf-spd]{\texttt{SPD\_0}/\texttt{SPD\_1}}, \texttt{r16..r19}):} The hardware Index Fetcher reads \hyperref[sec:csf-regs]{\texttt{IDVSSR::INDEX\_BUFFER}} (\texttt{r54:r55}) and pre-fills \texttt{r59, r60} with \texttt{VertexIndex} (\texttt{+ VERTEX\_OFFSET + GLOBAL\_ATTR\_OFF}) and \texttt{r61} with \texttt{InstanceIndex} (\texttt{+ INSTANCE\_OFFSET}) so Position and Varying shaders can pull vertex attributes directly.
  \item \textbf{Inline FAU \& Special Operands (\hyperref[sec:csf-tsd-fau]{\texttt{FAU\_k}}, \texttt{r8..r15}):} Any 8-bit source with \texttt{mode=2} reads zero-latency uniform \texttt{u0..u31} from \texttt{FAU\_k} (\texttt{slots64 = bits[63:56]}); \texttt{mode=3} reads \texttt{SmallConstant} (Table~1.2) or \texttt{FauSpecial} (\texttt{pc}, \texttt{lane\_id}, \texttt{lane\_width}, \texttt{tls\_ptr}, \texttt{wls\_ptr}, \texttt{blend\_descriptor\_0..7}).
\end{tightitemize}
\par\textbf{Worked Examples:}
\par\textbf{Example 1 (Legacy SFU bitwise complement (v9--v13)):}
\codebox{\mnem{NOT\_OLD.i32} \dst{r0}, \src{r1}}
With \texttt{r1 = 0x12345678}: bitwise complement -\textgreater{} \texttt{r0 = 0xEDCBA987} (removed in v14 in favor of \texttt{NOT}).
\par\textbf{Example 2 (Packed v2i16 complement):}
\codebox{\mnem{NOT\_OLD.i32} \dst{r0}\imod{.h0}, \src{r1}}
With \texttt{r1 = 0x123400FF}, \texttt{r0\_old = 0xAAAA0000}: complements low halfword to \texttt{0xFF00} and \texttt{.h0} writes \texttt{r0 = 0xAAAAFF00}.

\begin{center}\footnotesize
\begin{tabularx}{\linewidth}{@{}lllX@{}}
\toprule
Variant & Arch & Class & Opcode Fields \\
\midrule
\texttt{(none)} & v9--v14 & \texttt{R1} & \texttt{opcode2=NOT\_OLD, opcode1=MOV} \\
\texttt{(none)} & v9--v14 & \texttt{D1} & \texttt{opcode2=NOT\_OLD, opcode1=MOV} \\
\bottomrule
\end{tabularx}
\end{center}
\begin{center}\footnotesize
\begin{tabularx}{0.88\linewidth}{@{}l l X@{}}
\toprule
Field Name & Bits & Type / Fixed Value \\
\midrule
\texttt{src0} & \texttt{0:7} & \texttt{SourceEncoding} \\
\textit{Reserved} & \texttt{8:15} & \texttt{zero} \\
\texttt{opcode2} & \texttt{16:20} & \texttt{exact=NOT\_OLD} \\
\textit{Reserved} & \texttt{21:39} & \texttt{zero} \\
\texttt{dest\_register\_index} & \texttt{40:45} & \texttt{uint} \\
\hyperref[enum:dest_width_m]{\texttt{dest\_width}} & \texttt{46:47} & \texttt{dest\_width\_m} \\
\texttt{opcode1} & \texttt{48:51} & \texttt{exact=MOV} \\
\texttt{instruction\_class} & \texttt{52:56} & \texttt{exact=R1} \\
\texttt{fau\_page\_index} & \texttt{57:58} & \texttt{uint} \\
\hyperref[enum:flow_control_m]{\texttt{flow}} & \texttt{59:62} & \texttt{flow\_control\_m} \\
\textit{Reserved} & \texttt{63} & \texttt{zero} \\
\bottomrule
\end{tabularx}
\end{center}
\textbf{Fields Explained:}
\begin{tightitemize}\footnotesize
  \item \texttt{src0}: Primary source operand (\texttt{r0..r63}, \texttt{u0..u63} uniform, \texttt{k*} FAU constant, or \texttt{t0..t2} temporary; \texttt{\textasciicircum{}} marks discard/last-use). \emph{Example:} {\scriptsize\ttfamily \mnem{FADD.f32} \dst{r0}, \src{r1}, \src{r2}} (\textbf{SPIR-V:} \texttt{OpFAdd (Operand 0)}; \textbf{GLSL:} \texttt{float c = a + b;})
  \item \texttt{opcode2}: Secondary sub-opcode further specializing instructions sharing \texttt{instruction\_class} and \texttt{opcode1} (e.g., conversions, SFU tables, moves). \emph{Example:} {\scriptsize\ttfamily \mnem{F32\_TO\_S32}\imod{.round\_zero} \dst{r0}, \src{r1}} (\textbf{SPIR-V:} \texttt{OpConvertFToS}; \textbf{GLSL:} \texttt{int y = int(x);})
  \item \texttt{dest\_register\_index}: 6-bit destination GPR index (\texttt{r0..r63}) receiving the Stage C writeback result (even-aligned pair \texttt{[rN:rN+1]} for 64-bit ops). \emph{Example:} {\scriptsize\ttfamily \mnem{FMA.f32} \dst{r0}, \src{r1}, \src{r2}, \src{r3}} (\textbf{SPIR-V:} \texttt{SSA Result ID}; \textbf{GLSL:} \texttt{float r0 = fma(r1, r2, r3);})
  \item \hyperref[enum:dest_width_m]{\texttt{dest\_width}}: Stage C destination lane write-mask (\texttt{none} = full 32-bit, \texttt{h0} = low 16-bit, \texttt{h1} = high 16-bit without bit-shifting). \emph{Example:} {\scriptsize\ttfamily \mnem{FADD.v2f16} \dst{r0}\imod{.h1}, \src{r1}, \src{r2}} (\textbf{SPIR-V:} \texttt{OpVectorInsertDynamic / WriteMask}; \textbf{GLSL:} \texttt{dst.y = a.y + b.y;})
  \item \texttt{opcode1}: Primary sub-opcode within the instruction class uniquely identifying the operation family. \emph{Example:} {\scriptsize\ttfamily \mnem{IADD.s32} \dst{r0}, \src{r1}, \src{r2}} (\textbf{SPIR-V:} \texttt{OpIAdd (operation selector)}; \textbf{GLSL:} \texttt{// Primary opcode bits})
  \item \texttt{instruction\_class}: 4-bit major instruction class selector (\texttt{bits 48:51}) partitioning the 64-bit Valhall opcode space (e.g., \texttt{FMA}, \texttt{CVT}, \texttt{SFU}, \texttt{LS}, \texttt{TEX}, \texttt{BRANCH}). \emph{Example:} {\scriptsize\ttfamily \mnem{FADD.f32} \dst{r0}, \src{r1}, \src{r2}} (\textbf{SPIR-V:} \texttt{N/A (encoding class)}; \textbf{GLSL:} \texttt{// Hardware execution pipe / decode class})
  \item \texttt{fau\_page\_index}: Selects the FAU (Fast Access Uniform) page (\texttt{0..3}) when a source operand references hardware small constants, thread IDs, or lane constants. \emph{Example:} {\scriptsize\ttfamily \mnem{FADD.f32} \dst{r0}, \src{r1}, \src{k1f}} (\textbf{SPIR-V:} \texttt{OpConstant / BuiltIn}; \textbf{GLSL:} \texttt{float y = x + 1.0;})
  \item \hyperref[enum:flow_control_m]{\texttt{flow}}: 4-bit embedded flow-control \& scoreboard wait modifier (\texttt{bits 59:62}: \texttt{none}, \texttt{wait0..wait0126}, \texttt{wait}, \texttt{reconverge}, \texttt{discard}, \texttt{end}). \emph{Example:} {\scriptsize\ttfamily \mnem{FADD.f32}\imod{.wait0.reconverge} \dst{r0}, \src{r1}, \src{r2}} (\textbf{SPIR-V:} \texttt{OpControlBarrier / OpReturn}; \textbf{GLSL:} \texttt{// Scoreboard wait + warp reconvergence})
\end{tightitemize}
\clearpage

\subsection*{OR}
\addcontentsline{toc}{subsection}{OR}
Computes bitwise OR independently across each SIMD element $i \in \{0, \dots, N-1\}$ ($N=1$ for \texttt{i32}/\texttt{i64}, $N=2$ for \texttt{v2i16}, $N=4$ for \texttt{v4i8} on v9--v10): \texttt{res = widen0(src0)[i] \textbar{} (not2 ? \textasciitilde{}src2[i] : src2[i])}, then writes \texttt{not\_result ? \textasciitilde{}res : res} to \texttt{dest[i]}.
\par\textbf{Notes:}
\begin{tightitemize}
  \item Second operand is \texttt{src2}.
  \item \texttt{OpShiftLop(None, None)} encodes as \texttt{OR src0, 0} (\texttt{encode\_v9.rs:3060-3066}).
  \item \textbf{Source references:} \texttt{kraid/ops.rs:3944-4025}, \texttt{encode\_v9.rs:2916-3093}, \texttt{hw\_tests.rs:2148-2198}.
\end{tightitemize}
\codebox{\mnem{OR}\imod{\{.flow\}}\imod{\{.not\_result\}} \dst{dest}\imod{\{.dest\_width\}}, \src{src0}\imod{\{.widen0\}}, \src{src2}\imod{\{.not2\}}\\\mnem{OR}\imod{\{.flow\}}\imod{\{.not\_result\}} \dst{dest}, \src{src0}\imod{\{.widen0\}}, \src{src2}\imod{\{.not2\}}\\\mnem{OR}\imod{\{.flow\}}\imod{\{.not\_result\}} \dst{dest}\imod{\{.dest\_width\}}, \src{src0}\imod{\{.swz0\}}, \src{src2}\imod{\{.not2\}}\\\mnem{OR}\imod{\{.flow\}}\imod{\{.not\_result\}} \dst{dest}\imod{\{.dest\_width\}}, \src{src0}\imod{\{.lanes0\}}, \src{src2}\imod{\{.not2\}}}
\textbf{Operands:}
\begin{tightitemize}\footnotesize
  \item \texttt{\dst{dest}} (\texttt{dest\_register\_index}, \texttt{dest\_width}): Destination GPR receiving \texttt{src0 \textbar{} src2} (or NOR with \texttt{.not\_result}).
  \item \texttt{\src{src0}}: First bitwise operand.
  \item \texttt{\src{src2}}: Second bitwise operand (encoded in \texttt{src2}, not \texttt{src1}; supports \texttt{.not2}).
\end{tightitemize}
\par\textbf{Closest SPIR-V Op:} \texttt{OpBitwiseOr}
\par\textbf{Pseudo-GLSL (Operational Semantics):}
\codebox{for (int i = 0; i \textless{} N; ++i) \{\\    auto b = flag\_not2 ? \textasciitilde{}src2[i] : src2[i];\\    auto r = widen0(src0)[i] \textbar{} b;\\    dest[i] = flag\_not\_result ? \textasciitilde{}r : r; // NOR when .not\_result\\\}}
\textbf{Execution Context:}
\begin{tightitemize}\footnotesize
  \item \textbf{Compute stage preloads (\hyperref[sec:csf-spd]{\texttt{SPD\_k}}, \texttt{r16..r23}):} Hardware pre-fills \texttt{r55..r57} with \texttt{WorkGroupID.xyz + JOB\_OFFSET} (from \hyperref[sec:csf-regs]{\texttt{COMPUTESR::JOB\_OFFSET\_X/Y/Z}} \texttt{r34..r36}) and \texttt{r60..r62} with \texttt{LocalInvocationID.xyz} or \texttt{GlobalInvocationID.xyz} (from \texttt{WG\_SIZE} \texttt{r33} \& \texttt{JOB\_SIZE} \texttt{r37..r39}, selected by \texttt{ShaderProgram::preload}) so compute shaders can derive thread/element indices without memory loads.
  \item \textbf{Vertex stage preloads (\hyperref[sec:csf-spd]{\texttt{SPD\_0}/\texttt{SPD\_1}}, \texttt{r16..r19}):} The hardware Index Fetcher reads \hyperref[sec:csf-regs]{\texttt{IDVSSR::INDEX\_BUFFER}} (\texttt{r54:r55}) and pre-fills \texttt{r59, r60} with \texttt{VertexIndex} (\texttt{+ VERTEX\_OFFSET + GLOBAL\_ATTR\_OFF}) and \texttt{r61} with \texttt{InstanceIndex} (\texttt{+ INSTANCE\_OFFSET}) so Position and Varying shaders can pull vertex attributes directly.
  \item \textbf{Inline FAU \& Special Operands (\hyperref[sec:csf-tsd-fau]{\texttt{FAU\_k}}, \texttt{r8..r15}):} Any 8-bit source with \texttt{mode=2} reads zero-latency uniform \texttt{u0..u31} from \texttt{FAU\_k} (\texttt{slots64 = bits[63:56]}); \texttt{mode=3} reads \texttt{SmallConstant} (Table~1.2) or \texttt{FauSpecial} (\texttt{pc}, \texttt{lane\_id}, \texttt{lane\_width}, \texttt{tls\_ptr}, \texttt{wls\_ptr}, \texttt{blend\_descriptor\_0..7}).
\end{tightitemize}
\par\textbf{Worked Examples:}
\par\textbf{Example 1 (NOR via .not\_result):}
\codebox{\mnem{OR.i32}\imod{.not\_result} \dst{r0}, \src{r1}, \src{r2}}
With \texttt{r1 = 0xF0, r2 = 0x0A}: \texttt{0xF0 \textbar{} 0x0A = 0xFA}, then \texttt{.not\_result} inverts -\textgreater{} \texttt{r0 = 0xFFFFFF05}.
\par\textbf{Example 2 (Bitwise OR with inverted second operand):}
\codebox{\mnem{OR.i32} \dst{r0}, \src{r1}, \src{r2}\imod{.not}}
With \texttt{r1 = 0x0000000F, r2 = 0xFFFFFFF0}: \texttt{\textasciitilde{}r2 = 0x0000000F}, so \texttt{0xF \textbar{} 0xF = 0x0000000F}.

\begin{center}\footnotesize
\begin{tabularx}{\linewidth}{@{}lllX@{}}
\toprule
Variant & Arch & Class & Opcode Fields \\
\midrule
\texttt{(none)} & v9--v14 & \texttt{R3} & \texttt{opcode1=SHIFT\_BITWISE\_I32} \\
\texttt{(none)} & v9--v14 & \texttt{R3} & \texttt{opcode1=SHIFT\_BITWISE\_I32} \\
\texttt{(none)} & v9--v14 & \texttt{R3} & \texttt{opcode1=LOGIC\_SHIFT\_I32} \\
\texttt{(none)} & v9--v14 & \texttt{D3} & \texttt{opcode1=SHIFT\_BITWISE\_I64} \\
\texttt{(none)} & v9--v14 & \texttt{D3} & \texttt{opcode1=LOGIC\_SHIFT\_I64} \\
\texttt{(none)} & v9--v14 & \texttt{R3} & \texttt{opcode1=SHIFT\_BITWISE\_V2I16} \\
\texttt{(none)} & v9--v14 & \texttt{R3} & \texttt{opcode1=LOGIC\_SHIFT\_V2I16} \\
\texttt{(none)} & v9--v14 & \texttt{R3} & \texttt{opcode1=SHIFT\_BITWISE\_V4I8} \\
\texttt{(none)} & v9--v14 & \texttt{R3} & \texttt{opcode1=LOGIC\_SHIFT\_V4I8} \\
\bottomrule
\end{tabularx}
\end{center}
\begin{center}\footnotesize
\begin{tabularx}{0.88\linewidth}{@{}l l X@{}}
\toprule
Field Name & Bits & Type / Fixed Value \\
\midrule
\texttt{src0} & \texttt{0:7} & \texttt{SourceEncoding} \\
\texttt{src1} & \texttt{8:15} & \texttt{exact=192} \\
\texttt{src2} & \texttt{16:23} & \texttt{SourceEncoding} \\
\texttt{bitwise\_operation} & \texttt{24:25} & \texttt{exact=or} \\
\hyperref[enum:shift_bitwise_i32_lanes_m]{\texttt{lanes1}} & \texttt{26:29} & \texttt{exact=b0} \\
\hyperref[enum:not_m]{\texttt{not\_result}} & \texttt{30} & \texttt{not\_m (def none)} \\
\textit{Reserved} & \texttt{31} & \texttt{zero} \\
\texttt{shift\_operation} & \texttt{32:34} & \texttt{exact=lshift} \\
\hyperref[enum:not_m]{\texttt{not2}} & \texttt{35} & \texttt{not\_m (def none)} \\
\hyperref[enum:widen_int_m]{\texttt{widen0}} & \texttt{36:39} & \texttt{widen\_int\_m (def none)} \\
\texttt{dest\_register\_index} & \texttt{40:45} & \texttt{uint} \\
\hyperref[enum:dest_width_m]{\texttt{dest\_width}} & \texttt{46:47} & \texttt{dest\_width\_m} \\
\texttt{opcode1} & \texttt{48:51} & \texttt{exact=SHIFT\_BITWISE\_I32} \\
\texttt{instruction\_class} & \texttt{52:56} & \texttt{exact=R3} \\
\texttt{fau\_page\_index} & \texttt{57:58} & \texttt{uint} \\
\hyperref[enum:flow_control_m]{\texttt{flow}} & \texttt{59:62} & \texttt{flow\_control\_m} \\
\textit{Reserved} & \texttt{63} & \texttt{zero} \\
\bottomrule
\end{tabularx}
\end{center}
\textbf{Fields Explained:}
\begin{tightitemize}\footnotesize
  \item \texttt{src0}: Primary source operand (\texttt{r0..r63}, \texttt{u0..u63} uniform, \texttt{k*} FAU constant, or \texttt{t0..t2} temporary; \texttt{\textasciicircum{}} marks discard/last-use). \emph{Example:} {\scriptsize\ttfamily \mnem{FADD.f32} \dst{r0}, \src{r1}, \src{r2}} (\textbf{SPIR-V:} \texttt{OpFAdd (Operand 0)}; \textbf{GLSL:} \texttt{float c = a + b;})
  \item \texttt{src1}: Secondary source operand (register, uniform, or FAU constant; in shifts, supplies the shift count byte). \emph{Example:} {\scriptsize\ttfamily \mnem{LSHIFT.i32} \dst{r0}, \src{r1}, \src{r2}\imod{.b0}} (\textbf{SPIR-V:} \texttt{OpShiftLeftLogical (Shift)}; \textbf{GLSL:} \texttt{uint y = a \textless{}\textless{} s;})
  \item \texttt{src2}: Third source operand (FMA addend, \texttt{MUX}/\texttt{CSEL} selector, \texttt{IADDC} carry-in, or bitwise mask in fused shift-logic). \emph{Example:} {\scriptsize\ttfamily \mnem{FMA.f32} \dst{r0}, \src{r1}, \src{r2}, \src{r3}} (\textbf{SPIR-V:} \texttt{GLSL.std.450 Fma (Addend)}; \textbf{GLSL:} \texttt{float d = fma(a, b, c);})
  \item \texttt{bitwise\_operation}: Selects the bitwise ALU op (\texttt{and}, \texttt{or}, \texttt{xor}, \texttt{lshift\_and}, etc.) fused with shifts. \emph{Example:} {\scriptsize\ttfamily \mnem{LSHIFT\_OR.i32} \dst{r0}, \src{r1}, \src{r2}, \src{r3}} (\textbf{SPIR-V:} \texttt{OpShiftLeftLogical + OpBitwiseOr}; \textbf{GLSL:} \texttt{uint y = (a \textless{}\textless{} s) \textbar{} b;})
  \item \hyperref[enum:shift_bitwise_i32_lanes_m]{\texttt{lanes1}}: Stage A byte-lane selector on \texttt{src1} (selects shift-count byte \texttt{b0..b3} on 32/64-bit shifts, \texttt{bXY} on \texttt{v2i16} shifts, or 4-byte swizzle on \texttt{v4*8}). \emph{Example:} {\scriptsize\ttfamily \mnem{LSHIFT.i32} \dst{r0}, \src{r1}, \src{r2}\imod{.b1}} (\textbf{SPIR-V:} \texttt{OpBitFieldUExtract + OpShiftLeftLogical}; \textbf{GLSL:} \texttt{uint y = a \textless{}\textless{} ((s \textgreater{}\textgreater{} 8) \& 0xFFu);})
  \item \hyperref[enum:not_m]{\texttt{not\_result}}: Bitwise-inverts the final output of a bitwise/shift-bitwise instruction (\textasciitilde{}result). \emph{Example:} {\scriptsize\ttfamily \mnem{AND.i32}\imod{.not\_result} \dst{r0}, \src{r1}, \src{r2}} (\textbf{SPIR-V:} \texttt{OpBitwiseAnd + OpNot}; \textbf{GLSL:} \texttt{uint y = \textasciitilde{}(a \& b);})
  \item \texttt{shift\_operation}: Selects shift direction/type (\texttt{lshift}, \texttt{rshift}, \texttt{arshift}) in fused shift-bitwise instructions. \emph{Example:} {\scriptsize\ttfamily \mnem{RSHIFT\_AND.i32} \dst{r0}, \src{r1}, \src{r2}, \src{r3}} (\textbf{SPIR-V:} \texttt{OpShiftRightLogical + OpBitwiseAnd}; \textbf{GLSL:} \texttt{uint y = (a \textgreater{}\textgreater{} s) \& m;})
  \item \hyperref[enum:not_m]{\texttt{not2}}: Stage A bitwise complement (\texttt{\textasciitilde{}src2}) applied to \texttt{src2} before the bitwise/shift-logic operation. \emph{Example:} {\scriptsize\ttfamily \mnem{AND.i32} \dst{r0}, \src{r1}, \src{r2}\imod{.not}} (\textbf{SPIR-V:} \texttt{OpNot + OpBitwiseAnd}; \textbf{GLSL:} \texttt{uint y = a \& \textasciitilde{}b;})
  \item \hyperref[enum:widen_int_m]{\texttt{widen0}}: Stage A widening extractor on \texttt{src0} (promotes FP16 \texttt{h0}/\texttt{h1} to FP32, or sign/zero-extends \texttt{b0..b3}/\texttt{h0..h1}/\texttt{w0..w1} to 32/64-bit before \texttt{abs}/\texttt{neg}). \emph{Example:} {\scriptsize\ttfamily \mnem{IADD.s32} \dst{r0}, \src{r1}\imod{.h0}, \src{r2}} (\textbf{SPIR-V:} \texttt{OpSConvert / OpFConvert}; \textbf{GLSL:} \texttt{int c = int(int16\_t(a)) + b;})
  \item \texttt{dest\_register\_index}: 6-bit destination GPR index (\texttt{r0..r63}) receiving the Stage C writeback result (even-aligned pair \texttt{[rN:rN+1]} for 64-bit ops). \emph{Example:} {\scriptsize\ttfamily \mnem{FMA.f32} \dst{r0}, \src{r1}, \src{r2}, \src{r3}} (\textbf{SPIR-V:} \texttt{SSA Result ID}; \textbf{GLSL:} \texttt{float r0 = fma(r1, r2, r3);})
  \item \hyperref[enum:dest_width_m]{\texttt{dest\_width}}: Stage C destination lane write-mask (\texttt{none} = full 32-bit, \texttt{h0} = low 16-bit, \texttt{h1} = high 16-bit without bit-shifting). \emph{Example:} {\scriptsize\ttfamily \mnem{FADD.v2f16} \dst{r0}\imod{.h1}, \src{r1}, \src{r2}} (\textbf{SPIR-V:} \texttt{OpVectorInsertDynamic / WriteMask}; \textbf{GLSL:} \texttt{dst.y = a.y + b.y;})
  \item \texttt{opcode1}: Primary sub-opcode within the instruction class uniquely identifying the operation family. \emph{Example:} {\scriptsize\ttfamily \mnem{IADD.s32} \dst{r0}, \src{r1}, \src{r2}} (\textbf{SPIR-V:} \texttt{OpIAdd (operation selector)}; \textbf{GLSL:} \texttt{// Primary opcode bits})
  \item \texttt{instruction\_class}: 4-bit major instruction class selector (\texttt{bits 48:51}) partitioning the 64-bit Valhall opcode space (e.g., \texttt{FMA}, \texttt{CVT}, \texttt{SFU}, \texttt{LS}, \texttt{TEX}, \texttt{BRANCH}). \emph{Example:} {\scriptsize\ttfamily \mnem{FADD.f32} \dst{r0}, \src{r1}, \src{r2}} (\textbf{SPIR-V:} \texttt{N/A (encoding class)}; \textbf{GLSL:} \texttt{// Hardware execution pipe / decode class})
  \item \texttt{fau\_page\_index}: Selects the FAU (Fast Access Uniform) page (\texttt{0..3}) when a source operand references hardware small constants, thread IDs, or lane constants. \emph{Example:} {\scriptsize\ttfamily \mnem{FADD.f32} \dst{r0}, \src{r1}, \src{k1f}} (\textbf{SPIR-V:} \texttt{OpConstant / BuiltIn}; \textbf{GLSL:} \texttt{float y = x + 1.0;})
  \item \hyperref[enum:flow_control_m]{\texttt{flow}}: 4-bit embedded flow-control \& scoreboard wait modifier (\texttt{bits 59:62}: \texttt{none}, \texttt{wait0..wait0126}, \texttt{wait}, \texttt{reconverge}, \texttt{discard}, \texttt{end}). \emph{Example:} {\scriptsize\ttfamily \mnem{FADD.f32}\imod{.wait0.reconverge} \dst{r0}, \src{r1}, \src{r2}} (\textbf{SPIR-V:} \texttt{OpControlBarrier / OpReturn}; \textbf{GLSL:} \texttt{// Scoreboard wait + warp reconvergence})
\end{tightitemize}
\clearpage

\subsection*{OR\_IMM}
\addcontentsline{toc}{subsection}{OR\_IMM}
Computes a 32-bit bitwise OR between \texttt{src0} and the inline 32-bit immediate \texttt{\#imm1w} (\texttt{dest = src0 \textbar{} \#imm1w}) on Valhall v11+. Does not support source swizzles or \texttt{.not\_result}; supports \texttt{dest\_width} halfword write masking.
\par\textbf{Notes:}
\begin{tightitemize}
  \item No \texttt{not\_result} or \texttt{src0} swizzle supported.
  \item \textbf{Source references:} \texttt{kraid/encode\_v9.rs:2919-2926, 3036-3054}.
\end{tightitemize}
\codebox{\mnem{OR\_IMM}\imod{\{.flow\}} \dst{dest}\imod{\{.dest\_width\}}, \src{src0}, \imm{\#imm1w}}
\textbf{Operands:}
\begin{tightitemize}\footnotesize
  \item \texttt{\dst{dest}} (\texttt{dest\_register\_index}, \texttt{dest\_width}): 32-bit destination GPR receiving \texttt{src0 \textbar{} \#imm1w}.
  \item \texttt{\src{src0}}: 32-bit input operand.
  \item \texttt{\imm{\#imm1w}} (\texttt{imm1w}): Inline 32-bit bitwise OR mask.
\end{tightitemize}
\par\textbf{Closest SPIR-V Op:} \texttt{OpBitwiseOr (constant)}
\par\textbf{Pseudo-GLSL (Operational Semantics):}
\codebox{dest = maskWrite(dest\_width, src0 \textbar{} \#imm1w);}
\textbf{Execution Context:}
\begin{tightitemize}\footnotesize
  \item \textbf{Compute stage preloads (\hyperref[sec:csf-spd]{\texttt{SPD\_k}}, \texttt{r16..r23}):} Hardware pre-fills \texttt{r55..r57} with \texttt{WorkGroupID.xyz + JOB\_OFFSET} (from \hyperref[sec:csf-regs]{\texttt{COMPUTESR::JOB\_OFFSET\_X/Y/Z}} \texttt{r34..r36}) and \texttt{r60..r62} with \texttt{LocalInvocationID.xyz} or \texttt{GlobalInvocationID.xyz} (from \texttt{WG\_SIZE} \texttt{r33} \& \texttt{JOB\_SIZE} \texttt{r37..r39}, selected by \texttt{ShaderProgram::preload}) so compute shaders can derive thread/element indices without memory loads.
  \item \textbf{Vertex stage preloads (\hyperref[sec:csf-spd]{\texttt{SPD\_0}/\texttt{SPD\_1}}, \texttt{r16..r19}):} The hardware Index Fetcher reads \hyperref[sec:csf-regs]{\texttt{IDVSSR::INDEX\_BUFFER}} (\texttt{r54:r55}) and pre-fills \texttt{r59, r60} with \texttt{VertexIndex} (\texttt{+ VERTEX\_OFFSET + GLOBAL\_ATTR\_OFF}) and \texttt{r61} with \texttt{InstanceIndex} (\texttt{+ INSTANCE\_OFFSET}) so Position and Varying shaders can pull vertex attributes directly.
  \item \textbf{Inline FAU \& Special Operands (\hyperref[sec:csf-tsd-fau]{\texttt{FAU\_k}}, \texttt{r8..r15}):} Any 8-bit source with \texttt{mode=2} reads zero-latency uniform \texttt{u0..u31} from \texttt{FAU\_k} (\texttt{slots64 = bits[63:56]}); \texttt{mode=3} reads \texttt{SmallConstant} (Table~1.2) or \texttt{FauSpecial} (\texttt{pc}, \texttt{lane\_id}, \texttt{lane\_width}, \texttt{tls\_ptr}, \texttt{wls\_ptr}, \texttt{blend\_descriptor\_0..7}).
\end{tightitemize}
\par\textbf{Worked Examples:}
\par\textbf{Example 1 (Force negative sign bit (v11+)):}
\codebox{\mnem{OR\_IMM.i32} \dst{r0}, \src{r1}, \imm{\#0x80000000}}
With \texttt{r1 = 0x3F800000} (\texttt{+1.0f}): sets bit 31 -\textgreater{} \texttt{r0 = 0xBF800000} (\texttt{-1.0f}).
\par\textbf{Example 2 (Set flag bits in packed bitfield):}
\codebox{\mnem{OR\_IMM.i32} \dst{r0}, \src{r1}, \imm{\#0x00000005}}
With \texttt{r1 = 0x00000010}: sets bits 0 and 2 -\textgreater{} \texttt{r0 = 0x00000015}.

\begin{center}\footnotesize
\begin{tabularx}{\linewidth}{@{}lllX@{}}
\toprule
Variant & Arch & Class & Opcode Fields \\
\midrule
\texttt{(none)} & v9--v14 & \texttt{R1I} & \texttt{opcode1=OR\_IMM\_I32} \\
\bottomrule
\end{tabularx}
\end{center}
\begin{center}\footnotesize
\begin{tabularx}{0.88\linewidth}{@{}l l X@{}}
\toprule
Field Name & Bits & Type / Fixed Value \\
\midrule
\texttt{src0} & \texttt{0:7} & \texttt{SourceEncoding} \\
\texttt{imm1w} & \texttt{8:39} & \texttt{uint} \\
\texttt{dest\_register\_index} & \texttt{40:45} & \texttt{uint} \\
\hyperref[enum:dest_width_m]{\texttt{dest\_width}} & \texttt{46:47} & \texttt{dest\_width\_m} \\
\texttt{opcode1} & \texttt{48:51} & \texttt{exact=OR\_IMM\_I32} \\
\texttt{instruction\_class} & \texttt{52:56} & \texttt{exact=R1I} \\
\texttt{fau\_page\_index} & \texttt{57:58} & \texttt{uint} \\
\hyperref[enum:flow_control_m]{\texttt{flow}} & \texttt{59:62} & \texttt{flow\_control\_m} \\
\textit{Reserved} & \texttt{63} & \texttt{zero} \\
\bottomrule
\end{tabularx}
\end{center}
\textbf{Fields Explained:}
\begin{tightitemize}\footnotesize
  \item \texttt{src0}: Primary source operand (\texttt{r0..r63}, \texttt{u0..u63} uniform, \texttt{k*} FAU constant, or \texttt{t0..t2} temporary; \texttt{\textasciicircum{}} marks discard/last-use). \emph{Example:} {\scriptsize\ttfamily \mnem{FADD.f32} \dst{r0}, \src{r1}, \src{r2}} (\textbf{SPIR-V:} \texttt{OpFAdd (Operand 0)}; \textbf{GLSL:} \texttt{float c = a + b;})
  \item \texttt{imm1w}: 32-bit inline immediate operand encoded directly in the instruction word (bits 8:39). \emph{Example:} {\scriptsize\ttfamily \mnem{IADD\_IMM.i32} \dst{r0}, \src{r1}, \imm{\#0x10}} (\textbf{SPIR-V:} \texttt{OpIAdd (constant)}; \textbf{GLSL:} \texttt{int y = x + 16;})
  \item \texttt{dest\_register\_index}: 6-bit destination GPR index (\texttt{r0..r63}) receiving the Stage C writeback result (even-aligned pair \texttt{[rN:rN+1]} for 64-bit ops). \emph{Example:} {\scriptsize\ttfamily \mnem{FMA.f32} \dst{r0}, \src{r1}, \src{r2}, \src{r3}} (\textbf{SPIR-V:} \texttt{SSA Result ID}; \textbf{GLSL:} \texttt{float r0 = fma(r1, r2, r3);})
  \item \hyperref[enum:dest_width_m]{\texttt{dest\_width}}: Stage C destination lane write-mask (\texttt{none} = full 32-bit, \texttt{h0} = low 16-bit, \texttt{h1} = high 16-bit without bit-shifting). \emph{Example:} {\scriptsize\ttfamily \mnem{FADD.v2f16} \dst{r0}\imod{.h1}, \src{r1}, \src{r2}} (\textbf{SPIR-V:} \texttt{OpVectorInsertDynamic / WriteMask}; \textbf{GLSL:} \texttt{dst.y = a.y + b.y;})
  \item \texttt{opcode1}: Primary sub-opcode within the instruction class uniquely identifying the operation family. \emph{Example:} {\scriptsize\ttfamily \mnem{IADD.s32} \dst{r0}, \src{r1}, \src{r2}} (\textbf{SPIR-V:} \texttt{OpIAdd (operation selector)}; \textbf{GLSL:} \texttt{// Primary opcode bits})
  \item \texttt{instruction\_class}: 4-bit major instruction class selector (\texttt{bits 48:51}) partitioning the 64-bit Valhall opcode space (e.g., \texttt{FMA}, \texttt{CVT}, \texttt{SFU}, \texttt{LS}, \texttt{TEX}, \texttt{BRANCH}). \emph{Example:} {\scriptsize\ttfamily \mnem{FADD.f32} \dst{r0}, \src{r1}, \src{r2}} (\textbf{SPIR-V:} \texttt{N/A (encoding class)}; \textbf{GLSL:} \texttt{// Hardware execution pipe / decode class})
  \item \texttt{fau\_page\_index}: Selects the FAU (Fast Access Uniform) page (\texttt{0..3}) when a source operand references hardware small constants, thread IDs, or lane constants. \emph{Example:} {\scriptsize\ttfamily \mnem{FADD.f32} \dst{r0}, \src{r1}, \src{k1f}} (\textbf{SPIR-V:} \texttt{OpConstant / BuiltIn}; \textbf{GLSL:} \texttt{float y = x + 1.0;})
  \item \hyperref[enum:flow_control_m]{\texttt{flow}}: 4-bit embedded flow-control \& scoreboard wait modifier (\texttt{bits 59:62}: \texttt{none}, \texttt{wait0..wait0126}, \texttt{wait}, \texttt{reconverge}, \texttt{discard}, \texttt{end}). \emph{Example:} {\scriptsize\ttfamily \mnem{FADD.f32}\imod{.wait0.reconverge} \dst{r0}, \src{r1}, \src{r2}} (\textbf{SPIR-V:} \texttt{OpControlBarrier / OpReturn}; \textbf{GLSL:} \texttt{// Scoreboard wait + warp reconvergence})
\end{tightitemize}
\clearpage

\subsection*{RROT}
\addcontentsline{toc}{subsection}{RROT}
For each SIMD element $i \in \{0, \dots, N-1\}$ ($N=1$ for \texttt{i32}/\texttt{i64}, $N=2$ for \texttt{v2i16}, $N=4$ for \texttt{v4i8} on v9--v10, lane width $W$), computes rotate right \texttt{s = rotate\_right(val, amt \& (W - 1))} where \texttt{val = src0[i]} and the 8-bit shift count \texttt{amt} is extracted from \texttt{src1} via \texttt{lanes1}; writes \texttt{not\_result ? \textasciitilde{}s : s} to \texttt{dest[i]}.
\par\textbf{Notes:}
\begin{tightitemize}
  \item Shift amount is masked modulo lane width (\texttt{wrapping\_shl}/\texttt{wrapping\_shr}/\texttt{rotate\_*}, \texttt{kraid/ops.rs:3984-4009}).
  \item On \texttt{v2i16}, \texttt{src1} is \texttt{V2I8} (\texttt{ops.rs:3953}) with byte lanes selected by \texttt{lanes1} (default \texttt{b02}: byte 0 -\textgreater{} lane 0, byte 2 -\textgreater{} lane 1).
  \item \textbf{Source references:} \texttt{kraid/ops.rs:3878-4025}, \texttt{encode\_v9.rs:2940-3093}, \texttt{hw\_tests.rs:2148-2198}.
\end{tightitemize}
\codebox{\mnem{RROT}\imod{\{.flow\}}\imod{\{.not\_result\}} \dst{dest}\imod{\{.dest\_width\}}, \src{src0}\imod{\{.widen0\}}, \src{src1}\imod{\{.lanes1\}}\\\mnem{RROT}\imod{\{.flow\}}\imod{\{.not\_result\}} \dst{dest}, \src{src0}\imod{\{.widen0\}}, \src{src1}\imod{\{.lanes1\}}\\\mnem{RROT}\imod{\{.flow\}}\imod{\{.not\_result\}} \dst{dest}\imod{\{.dest\_width\}}, \src{src0}\imod{\{.swz0\}}, \src{src1}\imod{\{.lanes1\}}\\\mnem{RROT}\imod{\{.flow\}}\imod{\{.not\_result\}} \dst{dest}\imod{\{.dest\_width\}}, \src{src0}\imod{\{.lanes0\}}, \src{src1}\imod{\{.lanes1\}}}
\textbf{Operands:}
\begin{tightitemize}\footnotesize
  \item \texttt{\dst{dest}} (\texttt{dest\_register\_index}, \texttt{dest\_width}): Destination GPR receiving the rotate right (inverted if \texttt{.not\_result}).
  \item \texttt{\src{src0}}: Input integer value to shift/rotate.
  \item \texttt{\src{src1}}: 8-bit shift/rotate count per lane (selected via \texttt{lanes1}, masked modulo lane width \texttt{W}).
\end{tightitemize}
\par\textbf{Closest SPIR-V Op:} \texttt{None (lowered from `(x\textgreater{}\textgreater{}s)\textbar{}(x\textless{}\textless{}(32-s))`)}
\par\textbf{Pseudo-GLSL (Operational Semantics):}
\codebox{for (int i = 0; i \textless{} N; ++i) \{\\    uint s = extractByte(src1, lanes1, i) \& (W - 1);\\    auto r = ((src0[i] \textgreater{}\textgreater{} s) \textbar{} (src0[i] \textless{}\textless{} ((W - s) \& (W - 1))));\\    dest[i] = flag\_not\_result ? \textasciitilde{}r : r;\\\}}
\textbf{Execution Context:}
\begin{tightitemize}\footnotesize
  \item \textbf{Compute stage preloads (\hyperref[sec:csf-spd]{\texttt{SPD\_k}}, \texttt{r16..r23}):} Hardware pre-fills \texttt{r55..r57} with \texttt{WorkGroupID.xyz + JOB\_OFFSET} (from \hyperref[sec:csf-regs]{\texttt{COMPUTESR::JOB\_OFFSET\_X/Y/Z}} \texttt{r34..r36}) and \texttt{r60..r62} with \texttt{LocalInvocationID.xyz} or \texttt{GlobalInvocationID.xyz} (from \texttt{WG\_SIZE} \texttt{r33} \& \texttt{JOB\_SIZE} \texttt{r37..r39}, selected by \texttt{ShaderProgram::preload}) so compute shaders can derive thread/element indices without memory loads.
  \item \textbf{Vertex stage preloads (\hyperref[sec:csf-spd]{\texttt{SPD\_0}/\texttt{SPD\_1}}, \texttt{r16..r19}):} The hardware Index Fetcher reads \hyperref[sec:csf-regs]{\texttt{IDVSSR::INDEX\_BUFFER}} (\texttt{r54:r55}) and pre-fills \texttt{r59, r60} with \texttt{VertexIndex} (\texttt{+ VERTEX\_OFFSET + GLOBAL\_ATTR\_OFF}) and \texttt{r61} with \texttt{InstanceIndex} (\texttt{+ INSTANCE\_OFFSET}) so Position and Varying shaders can pull vertex attributes directly.
  \item \textbf{Inline FAU \& Special Operands (\hyperref[sec:csf-tsd-fau]{\texttt{FAU\_k}}, \texttt{r8..r15}):} Any 8-bit source with \texttt{mode=2} reads zero-latency uniform \texttt{u0..u31} from \texttt{FAU\_k} (\texttt{slots64 = bits[63:56]}); \texttt{mode=3} reads \texttt{SmallConstant} (Table~1.2) or \texttt{FauSpecial} (\texttt{pc}, \texttt{lane\_id}, \texttt{lane\_width}, \texttt{tls\_ptr}, \texttt{wls\_ptr}, \texttt{blend\_descriptor\_0..7}).
\end{tightitemize}
\par\textbf{Worked Examples:}
\par\textbf{Example 1 (Canonical 32-bit shift/rotate):}
\codebox{\mnem{RROT.i32} \dst{r0}, \src{r1}, \src{r2}\imod{.b0}}
With \texttt{r1 = 0x00000003}, \texttt{r2.b0 = 1}: yields \texttt{r0 = 0x80000001}.
\par\textbf{Example 2 (Byte-lane selector \& modulo-32 masking):}
\codebox{\mnem{RROT.i32} \dst{r0}, \src{r1}, \src{r2}\imod{.b1}}
With \texttt{r1 = 0x12345678}, \texttt{r2.b1 = 8} (shift count masked mod 32 -\textgreater{} \texttt{8}): yields \texttt{r0 = 0x78123456}.

\begin{center}\footnotesize
\begin{tabularx}{\linewidth}{@{}lllX@{}}
\toprule
Variant & Arch & Class & Opcode Fields \\
\midrule
\texttt{(none)} & v9--v14 & \texttt{R3} & \texttt{opcode1=SHIFT\_BITWISE\_I32} \\
\texttt{(none)} & v9--v14 & \texttt{R3} & \texttt{opcode1=LOGIC\_SHIFT\_I32} \\
\texttt{(none)} & v9--v14 & \texttt{D3} & \texttt{opcode1=SHIFT\_BITWISE\_I64} \\
\texttt{(none)} & v9--v14 & \texttt{D3} & \texttt{opcode1=LOGIC\_SHIFT\_I64} \\
\texttt{(none)} & v9--v14 & \texttt{R3} & \texttt{opcode1=SHIFT\_BITWISE\_V2I16} \\
\texttt{(none)} & v9--v14 & \texttt{R3} & \texttt{opcode1=LOGIC\_SHIFT\_V2I16} \\
\texttt{(none)} & v9--v14 & \texttt{R3} & \texttt{opcode1=SHIFT\_BITWISE\_V4I8} \\
\texttt{(none)} & v9--v14 & \texttt{R3} & \texttt{opcode1=LOGIC\_SHIFT\_V4I8} \\
\bottomrule
\end{tabularx}
\end{center}
\begin{center}\footnotesize
\begin{tabularx}{0.88\linewidth}{@{}l l X@{}}
\toprule
Field Name & Bits & Type / Fixed Value \\
\midrule
\texttt{src0} & \texttt{0:7} & \texttt{SourceEncoding} \\
\texttt{src1} & \texttt{8:15} & \texttt{SourceEncoding} \\
\texttt{src2} & \texttt{16:23} & \texttt{exact=192} \\
\texttt{bitwise\_operation} & \texttt{24:25} & \texttt{exact=or} \\
\hyperref[enum:shift_bitwise_i32_lanes_m]{\texttt{lanes1}} & \texttt{26:29} & \texttt{shift\_bitwise\_i32\_lanes\_m (def b0)} \\
\hyperref[enum:not_m]{\texttt{not\_result}} & \texttt{30} & \texttt{not\_m (def none)} \\
\textit{Reserved} & \texttt{31} & \texttt{zero} \\
\texttt{shift\_operation} & \texttt{32:34} & \texttt{exact=rrot} \\
\hyperref[enum:not_m]{\texttt{not2}} & \texttt{35} & \texttt{exact=none} \\
\hyperref[enum:widen_int_m]{\texttt{widen0}} & \texttt{36:39} & \texttt{widen\_int\_m (def none)} \\
\texttt{dest\_register\_index} & \texttt{40:45} & \texttt{uint} \\
\hyperref[enum:dest_width_m]{\texttt{dest\_width}} & \texttt{46:47} & \texttt{dest\_width\_m} \\
\texttt{opcode1} & \texttt{48:51} & \texttt{exact=SHIFT\_BITWISE\_I32} \\
\texttt{instruction\_class} & \texttt{52:56} & \texttt{exact=R3} \\
\texttt{fau\_page\_index} & \texttt{57:58} & \texttt{uint} \\
\hyperref[enum:flow_control_m]{\texttt{flow}} & \texttt{59:62} & \texttt{flow\_control\_m} \\
\textit{Reserved} & \texttt{63} & \texttt{zero} \\
\bottomrule
\end{tabularx}
\end{center}
\textbf{Fields Explained:}
\begin{tightitemize}\footnotesize
  \item \texttt{src0}: Primary source operand (\texttt{r0..r63}, \texttt{u0..u63} uniform, \texttt{k*} FAU constant, or \texttt{t0..t2} temporary; \texttt{\textasciicircum{}} marks discard/last-use). \emph{Example:} {\scriptsize\ttfamily \mnem{FADD.f32} \dst{r0}, \src{r1}, \src{r2}} (\textbf{SPIR-V:} \texttt{OpFAdd (Operand 0)}; \textbf{GLSL:} \texttt{float c = a + b;})
  \item \texttt{src1}: Secondary source operand (register, uniform, or FAU constant; in shifts, supplies the shift count byte). \emph{Example:} {\scriptsize\ttfamily \mnem{LSHIFT.i32} \dst{r0}, \src{r1}, \src{r2}\imod{.b0}} (\textbf{SPIR-V:} \texttt{OpShiftLeftLogical (Shift)}; \textbf{GLSL:} \texttt{uint y = a \textless{}\textless{} s;})
  \item \texttt{src2}: Third source operand (FMA addend, \texttt{MUX}/\texttt{CSEL} selector, \texttt{IADDC} carry-in, or bitwise mask in fused shift-logic). \emph{Example:} {\scriptsize\ttfamily \mnem{FMA.f32} \dst{r0}, \src{r1}, \src{r2}, \src{r3}} (\textbf{SPIR-V:} \texttt{GLSL.std.450 Fma (Addend)}; \textbf{GLSL:} \texttt{float d = fma(a, b, c);})
  \item \texttt{bitwise\_operation}: Selects the bitwise ALU op (\texttt{and}, \texttt{or}, \texttt{xor}, \texttt{lshift\_and}, etc.) fused with shifts. \emph{Example:} {\scriptsize\ttfamily \mnem{LSHIFT\_OR.i32} \dst{r0}, \src{r1}, \src{r2}, \src{r3}} (\textbf{SPIR-V:} \texttt{OpShiftLeftLogical + OpBitwiseOr}; \textbf{GLSL:} \texttt{uint y = (a \textless{}\textless{} s) \textbar{} b;})
  \item \hyperref[enum:shift_bitwise_i32_lanes_m]{\texttt{lanes1}}: Stage A byte-lane selector on \texttt{src1} (selects shift-count byte \texttt{b0..b3} on 32/64-bit shifts, \texttt{bXY} on \texttt{v2i16} shifts, or 4-byte swizzle on \texttt{v4*8}). \emph{Example:} {\scriptsize\ttfamily \mnem{LSHIFT.i32} \dst{r0}, \src{r1}, \src{r2}\imod{.b1}} (\textbf{SPIR-V:} \texttt{OpBitFieldUExtract + OpShiftLeftLogical}; \textbf{GLSL:} \texttt{uint y = a \textless{}\textless{} ((s \textgreater{}\textgreater{} 8) \& 0xFFu);})
  \item \hyperref[enum:not_m]{\texttt{not\_result}}: Bitwise-inverts the final output of a bitwise/shift-bitwise instruction (\textasciitilde{}result). \emph{Example:} {\scriptsize\ttfamily \mnem{AND.i32}\imod{.not\_result} \dst{r0}, \src{r1}, \src{r2}} (\textbf{SPIR-V:} \texttt{OpBitwiseAnd + OpNot}; \textbf{GLSL:} \texttt{uint y = \textasciitilde{}(a \& b);})
  \item \texttt{shift\_operation}: Selects shift direction/type (\texttt{lshift}, \texttt{rshift}, \texttt{arshift}) in fused shift-bitwise instructions. \emph{Example:} {\scriptsize\ttfamily \mnem{RSHIFT\_AND.i32} \dst{r0}, \src{r1}, \src{r2}, \src{r3}} (\textbf{SPIR-V:} \texttt{OpShiftRightLogical + OpBitwiseAnd}; \textbf{GLSL:} \texttt{uint y = (a \textgreater{}\textgreater{} s) \& m;})
  \item \hyperref[enum:not_m]{\texttt{not2}}: Stage A bitwise complement (\texttt{\textasciitilde{}src2}) applied to \texttt{src2} before the bitwise/shift-logic operation. \emph{Example:} {\scriptsize\ttfamily \mnem{AND.i32} \dst{r0}, \src{r1}, \src{r2}\imod{.not}} (\textbf{SPIR-V:} \texttt{OpNot + OpBitwiseAnd}; \textbf{GLSL:} \texttt{uint y = a \& \textasciitilde{}b;})
  \item \hyperref[enum:widen_int_m]{\texttt{widen0}}: Stage A widening extractor on \texttt{src0} (promotes FP16 \texttt{h0}/\texttt{h1} to FP32, or sign/zero-extends \texttt{b0..b3}/\texttt{h0..h1}/\texttt{w0..w1} to 32/64-bit before \texttt{abs}/\texttt{neg}). \emph{Example:} {\scriptsize\ttfamily \mnem{IADD.s32} \dst{r0}, \src{r1}\imod{.h0}, \src{r2}} (\textbf{SPIR-V:} \texttt{OpSConvert / OpFConvert}; \textbf{GLSL:} \texttt{int c = int(int16\_t(a)) + b;})
  \item \texttt{dest\_register\_index}: 6-bit destination GPR index (\texttt{r0..r63}) receiving the Stage C writeback result (even-aligned pair \texttt{[rN:rN+1]} for 64-bit ops). \emph{Example:} {\scriptsize\ttfamily \mnem{FMA.f32} \dst{r0}, \src{r1}, \src{r2}, \src{r3}} (\textbf{SPIR-V:} \texttt{SSA Result ID}; \textbf{GLSL:} \texttt{float r0 = fma(r1, r2, r3);})
  \item \hyperref[enum:dest_width_m]{\texttt{dest\_width}}: Stage C destination lane write-mask (\texttt{none} = full 32-bit, \texttt{h0} = low 16-bit, \texttt{h1} = high 16-bit without bit-shifting). \emph{Example:} {\scriptsize\ttfamily \mnem{FADD.v2f16} \dst{r0}\imod{.h1}, \src{r1}, \src{r2}} (\textbf{SPIR-V:} \texttt{OpVectorInsertDynamic / WriteMask}; \textbf{GLSL:} \texttt{dst.y = a.y + b.y;})
  \item \texttt{opcode1}: Primary sub-opcode within the instruction class uniquely identifying the operation family. \emph{Example:} {\scriptsize\ttfamily \mnem{IADD.s32} \dst{r0}, \src{r1}, \src{r2}} (\textbf{SPIR-V:} \texttt{OpIAdd (operation selector)}; \textbf{GLSL:} \texttt{// Primary opcode bits})
  \item \texttt{instruction\_class}: 4-bit major instruction class selector (\texttt{bits 48:51}) partitioning the 64-bit Valhall opcode space (e.g., \texttt{FMA}, \texttt{CVT}, \texttt{SFU}, \texttt{LS}, \texttt{TEX}, \texttt{BRANCH}). \emph{Example:} {\scriptsize\ttfamily \mnem{FADD.f32} \dst{r0}, \src{r1}, \src{r2}} (\textbf{SPIR-V:} \texttt{N/A (encoding class)}; \textbf{GLSL:} \texttt{// Hardware execution pipe / decode class})
  \item \texttt{fau\_page\_index}: Selects the FAU (Fast Access Uniform) page (\texttt{0..3}) when a source operand references hardware small constants, thread IDs, or lane constants. \emph{Example:} {\scriptsize\ttfamily \mnem{FADD.f32} \dst{r0}, \src{r1}, \src{k1f}} (\textbf{SPIR-V:} \texttt{OpConstant / BuiltIn}; \textbf{GLSL:} \texttt{float y = x + 1.0;})
  \item \hyperref[enum:flow_control_m]{\texttt{flow}}: 4-bit embedded flow-control \& scoreboard wait modifier (\texttt{bits 59:62}: \texttt{none}, \texttt{wait0..wait0126}, \texttt{wait}, \texttt{reconverge}, \texttt{discard}, \texttt{end}). \emph{Example:} {\scriptsize\ttfamily \mnem{FADD.f32}\imod{.wait0.reconverge} \dst{r0}, \src{r1}, \src{r2}} (\textbf{SPIR-V:} \texttt{OpControlBarrier / OpReturn}; \textbf{GLSL:} \texttt{// Scoreboard wait + warp reconvergence})
\end{tightitemize}
\clearpage

\subsection*{RROT\_AND}
\addcontentsline{toc}{subsection}{RROT\_AND}
For each SIMD element $i \in \{0, \dots, N-1\}$ ($N=1$ for \texttt{i32}/\texttt{i64}, $N=2$ for \texttt{v2i16}, $N=4$ for \texttt{v4i8} on v9--v10, lane width $W$), computes rotate right \texttt{s = rotate\_right(val, amt \& (W - 1))} where \texttt{val = src0[i]} and the 8-bit shift count \texttt{amt} is extracted from \texttt{src1} via \texttt{lanes1}, then combines with \texttt{src2[i]} via bitwise AND \texttt{res = s \& (not2 ? \textasciitilde{}src2[i] : src2[i])} and writes \texttt{not\_result ? \textasciitilde{}res : res} to \texttt{dest[i]}.
\par\textbf{Notes:}
\begin{tightitemize}
  \item Shift amount is masked modulo lane width (\texttt{wrapping\_shl}/\texttt{wrapping\_shr}/\texttt{rotate\_*}, \texttt{kraid/ops.rs:3984-4009}).
  \item On \texttt{v2i16}, \texttt{src1} is \texttt{V2I8} (\texttt{ops.rs:3953}) with byte lanes selected by \texttt{lanes1} (default \texttt{b02}: byte 0 -\textgreater{} lane 0, byte 2 -\textgreater{} lane 1).
  \item \textbf{Source references:} \texttt{kraid/ops.rs:3878-4025}, \texttt{encode\_v9.rs:2940-3093}, \texttt{hw\_tests.rs:2148-2198}.
\end{tightitemize}
\codebox{\mnem{RROT\_AND}\imod{\{.flow\}}\imod{\{.not\_result\}} \dst{dest}\imod{\{.dest\_width\}}, \src{src0}\imod{\{.widen0\}}, \src{src1}\imod{\{.lanes1\}}, \src{src2}\imod{\{.not2\}}\\\mnem{RROT\_AND}\imod{\{.flow\}}\imod{\{.not\_result\}} \dst{dest}, \src{src0}\imod{\{.widen0\}}, \src{src1}\imod{\{.lanes1\}}, \src{src2}\imod{\{.not2\}}\\\mnem{RROT\_AND}\imod{\{.flow\}}\imod{\{.not\_result\}} \dst{dest}\imod{\{.dest\_width\}}, \src{src0}\imod{\{.swz0\}}, \src{src1}\imod{\{.lanes1\}}, \src{src2}\imod{\{.not2\}}\\\mnem{RROT\_AND}\imod{\{.flow\}}\imod{\{.not\_result\}} \dst{dest}\imod{\{.dest\_width\}}, \src{src0}\imod{\{.lanes0\}}, \src{src1}\imod{\{.lanes1\}}, \src{src2}\imod{\{.not2\}}}
\textbf{Operands:}
\begin{tightitemize}\footnotesize
  \item \texttt{\dst{dest}} (\texttt{dest\_register\_index}, \texttt{dest\_width}): Destination GPR receiving the rotate right combined via bitwise AND (inverted if \texttt{.not\_result}).
  \item \texttt{\src{src0}}: Input integer value to shift/rotate.
  \item \texttt{\src{src1}}: 8-bit shift/rotate count per lane (selected via \texttt{lanes1}, masked modulo lane width \texttt{W}).
  \item \texttt{\src{src2}}: Second operand combined via bitwise AND after shifting (inverted if \texttt{.not2}).
\end{tightitemize}
\par\textbf{Closest SPIR-V Op:} \texttt{None (lowered from `(x\textgreater{}\textgreater{}s)\textbar{}(x\textless{}\textless{}(32-s))`) + OpBitwiseAnd}
\par\textbf{Pseudo-GLSL (Operational Semantics):}
\codebox{for (int i = 0; i \textless{} N; ++i) \{\\    uint s = extractByte(src1, lanes1, i) \& (W - 1);\\    auto r = ((src0[i] \textgreater{}\textgreater{} s) \textbar{} (src0[i] \textless{}\textless{} ((W - s) \& (W - 1)))) \& (flag\_not2 ? \textasciitilde{}src2[i] : src2[i]);\\    dest[i] = flag\_not\_result ? \textasciitilde{}r : r;\\\}}
\textbf{Execution Context:}
\begin{tightitemize}\footnotesize
  \item \textbf{Compute stage preloads (\hyperref[sec:csf-spd]{\texttt{SPD\_k}}, \texttt{r16..r23}):} Hardware pre-fills \texttt{r55..r57} with \texttt{WorkGroupID.xyz + JOB\_OFFSET} (from \hyperref[sec:csf-regs]{\texttt{COMPUTESR::JOB\_OFFSET\_X/Y/Z}} \texttt{r34..r36}) and \texttt{r60..r62} with \texttt{LocalInvocationID.xyz} or \texttt{GlobalInvocationID.xyz} (from \texttt{WG\_SIZE} \texttt{r33} \& \texttt{JOB\_SIZE} \texttt{r37..r39}, selected by \texttt{ShaderProgram::preload}) so compute shaders can derive thread/element indices without memory loads.
  \item \textbf{Vertex stage preloads (\hyperref[sec:csf-spd]{\texttt{SPD\_0}/\texttt{SPD\_1}}, \texttt{r16..r19}):} The hardware Index Fetcher reads \hyperref[sec:csf-regs]{\texttt{IDVSSR::INDEX\_BUFFER}} (\texttt{r54:r55}) and pre-fills \texttt{r59, r60} with \texttt{VertexIndex} (\texttt{+ VERTEX\_OFFSET + GLOBAL\_ATTR\_OFF}) and \texttt{r61} with \texttt{InstanceIndex} (\texttt{+ INSTANCE\_OFFSET}) so Position and Varying shaders can pull vertex attributes directly.
  \item \textbf{Inline FAU \& Special Operands (\hyperref[sec:csf-tsd-fau]{\texttt{FAU\_k}}, \texttt{r8..r15}):} Any 8-bit source with \texttt{mode=2} reads zero-latency uniform \texttt{u0..u31} from \texttt{FAU\_k} (\texttt{slots64 = bits[63:56]}); \texttt{mode=3} reads \texttt{SmallConstant} (Table~1.2) or \texttt{FauSpecial} (\texttt{pc}, \texttt{lane\_id}, \texttt{lane\_width}, \texttt{tls\_ptr}, \texttt{wls\_ptr}, \texttt{blend\_descriptor\_0..7}).
\end{tightitemize}
\par\textbf{Worked Examples:}
\par\textbf{Example 1 (Standard fused shift + bitwise op):}
\codebox{\mnem{RROT\_AND.i32} \dst{r0}, \src{r1}, \src{r2}\imod{.b0}, \src{r3}}
With \texttt{r1 = 0x00000003}, \texttt{r2.b0 = 1} (masked mod 32): shift yields \texttt{0x80000001}, then \texttt{\& 0x0FFFFFFF} yields \texttt{r0 = 0x1}.
\par\textbf{Example 2 (Modulo-32 shift count \& .not\_result):}
\codebox{\mnem{RROT\_AND.i32}\imod{.not\_result} \dst{r0}, \src{r1}, \src{r2}\imod{.b0}, \src{r3}}
With \texttt{r1 = 0x12345678}, \texttt{r2.b0 = 8} (effective shift \texttt{8}): shift yields \texttt{0x78123456}, \texttt{\& 0x0FFFFFFF} yields \texttt{0x08123456}, and \texttt{.not\_result} inverts -\textgreater{} \texttt{r0 = 0xF7EDCBA9}.

\begin{center}\footnotesize
\begin{tabularx}{\linewidth}{@{}lllX@{}}
\toprule
Variant & Arch & Class & Opcode Fields \\
\midrule
\texttt{(none)} & v9--v14 & \texttt{R3} & \texttt{opcode1=SHIFT\_BITWISE\_I32} \\
\texttt{(none)} & v9--v14 & \texttt{R3} & \texttt{opcode1=LOGIC\_SHIFT\_I32} \\
\texttt{(none)} & v9--v14 & \texttt{D3} & \texttt{opcode1=SHIFT\_BITWISE\_I64} \\
\texttt{(none)} & v9--v14 & \texttt{D3} & \texttt{opcode1=LOGIC\_SHIFT\_I64} \\
\texttt{(none)} & v9--v14 & \texttt{R3} & \texttt{opcode1=SHIFT\_BITWISE\_V2I16} \\
\texttt{(none)} & v9--v14 & \texttt{R3} & \texttt{opcode1=LOGIC\_SHIFT\_V2I16} \\
\texttt{(none)} & v9--v14 & \texttt{R3} & \texttt{opcode1=SHIFT\_BITWISE\_V4I8} \\
\texttt{(none)} & v9--v14 & \texttt{R3} & \texttt{opcode1=LOGIC\_SHIFT\_V4I8} \\
\bottomrule
\end{tabularx}
\end{center}
\begin{center}\footnotesize
\begin{tabularx}{0.88\linewidth}{@{}l l X@{}}
\toprule
Field Name & Bits & Type / Fixed Value \\
\midrule
\texttt{src0} & \texttt{0:7} & \texttt{SourceEncoding} \\
\texttt{src1} & \texttt{8:15} & \texttt{SourceEncoding} \\
\texttt{src2} & \texttt{16:23} & \texttt{SourceEncoding} \\
\texttt{bitwise\_operation} & \texttt{24:25} & \texttt{exact=and} \\
\hyperref[enum:shift_bitwise_i32_lanes_m]{\texttt{lanes1}} & \texttt{26:29} & \texttt{shift\_bitwise\_i32\_lanes\_m (def b0)} \\
\hyperref[enum:not_m]{\texttt{not\_result}} & \texttt{30} & \texttt{not\_m (def none)} \\
\textit{Reserved} & \texttt{31} & \texttt{zero} \\
\texttt{shift\_operation} & \texttt{32:34} & \texttt{exact=rrot} \\
\hyperref[enum:not_m]{\texttt{not2}} & \texttt{35} & \texttt{not\_m (def none)} \\
\hyperref[enum:widen_int_m]{\texttt{widen0}} & \texttt{36:39} & \texttt{widen\_int\_m (def none)} \\
\texttt{dest\_register\_index} & \texttt{40:45} & \texttt{uint} \\
\hyperref[enum:dest_width_m]{\texttt{dest\_width}} & \texttt{46:47} & \texttt{dest\_width\_m} \\
\texttt{opcode1} & \texttt{48:51} & \texttt{exact=SHIFT\_BITWISE\_I32} \\
\texttt{instruction\_class} & \texttt{52:56} & \texttt{exact=R3} \\
\texttt{fau\_page\_index} & \texttt{57:58} & \texttt{uint} \\
\hyperref[enum:flow_control_m]{\texttt{flow}} & \texttt{59:62} & \texttt{flow\_control\_m} \\
\textit{Reserved} & \texttt{63} & \texttt{zero} \\
\bottomrule
\end{tabularx}
\end{center}
\textbf{Fields Explained:}
\begin{tightitemize}\footnotesize
  \item \texttt{src0}: Primary source operand (\texttt{r0..r63}, \texttt{u0..u63} uniform, \texttt{k*} FAU constant, or \texttt{t0..t2} temporary; \texttt{\textasciicircum{}} marks discard/last-use). \emph{Example:} {\scriptsize\ttfamily \mnem{FADD.f32} \dst{r0}, \src{r1}, \src{r2}} (\textbf{SPIR-V:} \texttt{OpFAdd (Operand 0)}; \textbf{GLSL:} \texttt{float c = a + b;})
  \item \texttt{src1}: Secondary source operand (register, uniform, or FAU constant; in shifts, supplies the shift count byte). \emph{Example:} {\scriptsize\ttfamily \mnem{LSHIFT.i32} \dst{r0}, \src{r1}, \src{r2}\imod{.b0}} (\textbf{SPIR-V:} \texttt{OpShiftLeftLogical (Shift)}; \textbf{GLSL:} \texttt{uint y = a \textless{}\textless{} s;})
  \item \texttt{src2}: Third source operand (FMA addend, \texttt{MUX}/\texttt{CSEL} selector, \texttt{IADDC} carry-in, or bitwise mask in fused shift-logic). \emph{Example:} {\scriptsize\ttfamily \mnem{FMA.f32} \dst{r0}, \src{r1}, \src{r2}, \src{r3}} (\textbf{SPIR-V:} \texttt{GLSL.std.450 Fma (Addend)}; \textbf{GLSL:} \texttt{float d = fma(a, b, c);})
  \item \texttt{bitwise\_operation}: Selects the bitwise ALU op (\texttt{and}, \texttt{or}, \texttt{xor}, \texttt{lshift\_and}, etc.) fused with shifts. \emph{Example:} {\scriptsize\ttfamily \mnem{LSHIFT\_OR.i32} \dst{r0}, \src{r1}, \src{r2}, \src{r3}} (\textbf{SPIR-V:} \texttt{OpShiftLeftLogical + OpBitwiseOr}; \textbf{GLSL:} \texttt{uint y = (a \textless{}\textless{} s) \textbar{} b;})
  \item \hyperref[enum:shift_bitwise_i32_lanes_m]{\texttt{lanes1}}: Stage A byte-lane selector on \texttt{src1} (selects shift-count byte \texttt{b0..b3} on 32/64-bit shifts, \texttt{bXY} on \texttt{v2i16} shifts, or 4-byte swizzle on \texttt{v4*8}). \emph{Example:} {\scriptsize\ttfamily \mnem{LSHIFT.i32} \dst{r0}, \src{r1}, \src{r2}\imod{.b1}} (\textbf{SPIR-V:} \texttt{OpBitFieldUExtract + OpShiftLeftLogical}; \textbf{GLSL:} \texttt{uint y = a \textless{}\textless{} ((s \textgreater{}\textgreater{} 8) \& 0xFFu);})
  \item \hyperref[enum:not_m]{\texttt{not\_result}}: Bitwise-inverts the final output of a bitwise/shift-bitwise instruction (\textasciitilde{}result). \emph{Example:} {\scriptsize\ttfamily \mnem{AND.i32}\imod{.not\_result} \dst{r0}, \src{r1}, \src{r2}} (\textbf{SPIR-V:} \texttt{OpBitwiseAnd + OpNot}; \textbf{GLSL:} \texttt{uint y = \textasciitilde{}(a \& b);})
  \item \texttt{shift\_operation}: Selects shift direction/type (\texttt{lshift}, \texttt{rshift}, \texttt{arshift}) in fused shift-bitwise instructions. \emph{Example:} {\scriptsize\ttfamily \mnem{RSHIFT\_AND.i32} \dst{r0}, \src{r1}, \src{r2}, \src{r3}} (\textbf{SPIR-V:} \texttt{OpShiftRightLogical + OpBitwiseAnd}; \textbf{GLSL:} \texttt{uint y = (a \textgreater{}\textgreater{} s) \& m;})
  \item \hyperref[enum:not_m]{\texttt{not2}}: Stage A bitwise complement (\texttt{\textasciitilde{}src2}) applied to \texttt{src2} before the bitwise/shift-logic operation. \emph{Example:} {\scriptsize\ttfamily \mnem{AND.i32} \dst{r0}, \src{r1}, \src{r2}\imod{.not}} (\textbf{SPIR-V:} \texttt{OpNot + OpBitwiseAnd}; \textbf{GLSL:} \texttt{uint y = a \& \textasciitilde{}b;})
  \item \hyperref[enum:widen_int_m]{\texttt{widen0}}: Stage A widening extractor on \texttt{src0} (promotes FP16 \texttt{h0}/\texttt{h1} to FP32, or sign/zero-extends \texttt{b0..b3}/\texttt{h0..h1}/\texttt{w0..w1} to 32/64-bit before \texttt{abs}/\texttt{neg}). \emph{Example:} {\scriptsize\ttfamily \mnem{IADD.s32} \dst{r0}, \src{r1}\imod{.h0}, \src{r2}} (\textbf{SPIR-V:} \texttt{OpSConvert / OpFConvert}; \textbf{GLSL:} \texttt{int c = int(int16\_t(a)) + b;})
  \item \texttt{dest\_register\_index}: 6-bit destination GPR index (\texttt{r0..r63}) receiving the Stage C writeback result (even-aligned pair \texttt{[rN:rN+1]} for 64-bit ops). \emph{Example:} {\scriptsize\ttfamily \mnem{FMA.f32} \dst{r0}, \src{r1}, \src{r2}, \src{r3}} (\textbf{SPIR-V:} \texttt{SSA Result ID}; \textbf{GLSL:} \texttt{float r0 = fma(r1, r2, r3);})
  \item \hyperref[enum:dest_width_m]{\texttt{dest\_width}}: Stage C destination lane write-mask (\texttt{none} = full 32-bit, \texttt{h0} = low 16-bit, \texttt{h1} = high 16-bit without bit-shifting). \emph{Example:} {\scriptsize\ttfamily \mnem{FADD.v2f16} \dst{r0}\imod{.h1}, \src{r1}, \src{r2}} (\textbf{SPIR-V:} \texttt{OpVectorInsertDynamic / WriteMask}; \textbf{GLSL:} \texttt{dst.y = a.y + b.y;})
  \item \texttt{opcode1}: Primary sub-opcode within the instruction class uniquely identifying the operation family. \emph{Example:} {\scriptsize\ttfamily \mnem{IADD.s32} \dst{r0}, \src{r1}, \src{r2}} (\textbf{SPIR-V:} \texttt{OpIAdd (operation selector)}; \textbf{GLSL:} \texttt{// Primary opcode bits})
  \item \texttt{instruction\_class}: 4-bit major instruction class selector (\texttt{bits 48:51}) partitioning the 64-bit Valhall opcode space (e.g., \texttt{FMA}, \texttt{CVT}, \texttt{SFU}, \texttt{LS}, \texttt{TEX}, \texttt{BRANCH}). \emph{Example:} {\scriptsize\ttfamily \mnem{FADD.f32} \dst{r0}, \src{r1}, \src{r2}} (\textbf{SPIR-V:} \texttt{N/A (encoding class)}; \textbf{GLSL:} \texttt{// Hardware execution pipe / decode class})
  \item \texttt{fau\_page\_index}: Selects the FAU (Fast Access Uniform) page (\texttt{0..3}) when a source operand references hardware small constants, thread IDs, or lane constants. \emph{Example:} {\scriptsize\ttfamily \mnem{FADD.f32} \dst{r0}, \src{r1}, \src{k1f}} (\textbf{SPIR-V:} \texttt{OpConstant / BuiltIn}; \textbf{GLSL:} \texttt{float y = x + 1.0;})
  \item \hyperref[enum:flow_control_m]{\texttt{flow}}: 4-bit embedded flow-control \& scoreboard wait modifier (\texttt{bits 59:62}: \texttt{none}, \texttt{wait0..wait0126}, \texttt{wait}, \texttt{reconverge}, \texttt{discard}, \texttt{end}). \emph{Example:} {\scriptsize\ttfamily \mnem{FADD.f32}\imod{.wait0.reconverge} \dst{r0}, \src{r1}, \src{r2}} (\textbf{SPIR-V:} \texttt{OpControlBarrier / OpReturn}; \textbf{GLSL:} \texttt{// Scoreboard wait + warp reconvergence})
\end{tightitemize}
\clearpage

\subsection*{RROT\_OR}
\addcontentsline{toc}{subsection}{RROT\_OR}
For each SIMD element $i \in \{0, \dots, N-1\}$ ($N=1$ for \texttt{i32}/\texttt{i64}, $N=2$ for \texttt{v2i16}, $N=4$ for \texttt{v4i8} on v9--v10, lane width $W$), computes rotate right \texttt{s = rotate\_right(val, amt \& (W - 1))} where \texttt{val = src0[i]} and the 8-bit shift count \texttt{amt} is extracted from \texttt{src1} via \texttt{lanes1}, then combines with \texttt{src2[i]} via bitwise OR \texttt{res = s \textbar{} (not2 ? \textasciitilde{}src2[i] : src2[i])} and writes \texttt{not\_result ? \textasciitilde{}res : res} to \texttt{dest[i]}.
\par\textbf{Notes:}
\begin{tightitemize}
  \item Shift amount is masked modulo lane width (\texttt{wrapping\_shl}/\texttt{wrapping\_shr}/\texttt{rotate\_*}, \texttt{kraid/ops.rs:3984-4009}).
  \item On \texttt{v2i16}, \texttt{src1} is \texttt{V2I8} (\texttt{ops.rs:3953}) with byte lanes selected by \texttt{lanes1} (default \texttt{b02}: byte 0 -\textgreater{} lane 0, byte 2 -\textgreater{} lane 1).
  \item \textbf{Source references:} \texttt{kraid/ops.rs:3878-4025}, \texttt{encode\_v9.rs:2940-3093}, \texttt{hw\_tests.rs:2148-2198}.
\end{tightitemize}
\codebox{\mnem{RROT\_OR}\imod{\{.flow\}}\imod{\{.not\_result\}} \dst{dest}\imod{\{.dest\_width\}}, \src{src0}\imod{\{.widen0\}}, \src{src1}\imod{\{.lanes1\}}, \src{src2}\imod{\{.not2\}}\\\mnem{RROT\_OR}\imod{\{.flow\}}\imod{\{.not\_result\}} \dst{dest}, \src{src0}\imod{\{.widen0\}}, \src{src1}\imod{\{.lanes1\}}, \src{src2}\imod{\{.not2\}}\\\mnem{RROT\_OR}\imod{\{.flow\}}\imod{\{.not\_result\}} \dst{dest}\imod{\{.dest\_width\}}, \src{src0}\imod{\{.swz0\}}, \src{src1}\imod{\{.lanes1\}}, \src{src2}\imod{\{.not2\}}\\\mnem{RROT\_OR}\imod{\{.flow\}}\imod{\{.not\_result\}} \dst{dest}\imod{\{.dest\_width\}}, \src{src0}\imod{\{.lanes0\}}, \src{src1}\imod{\{.lanes1\}}, \src{src2}\imod{\{.not2\}}}
\textbf{Operands:}
\begin{tightitemize}\footnotesize
  \item \texttt{\dst{dest}} (\texttt{dest\_register\_index}, \texttt{dest\_width}): Destination GPR receiving the rotate right combined via bitwise OR (inverted if \texttt{.not\_result}).
  \item \texttt{\src{src0}}: Input integer value to shift/rotate.
  \item \texttt{\src{src1}}: 8-bit shift/rotate count per lane (selected via \texttt{lanes1}, masked modulo lane width \texttt{W}).
  \item \texttt{\src{src2}}: Second operand combined via bitwise OR after shifting (inverted if \texttt{.not2}).
\end{tightitemize}
\par\textbf{Closest SPIR-V Op:} \texttt{None (lowered from `(x\textgreater{}\textgreater{}s)\textbar{}(x\textless{}\textless{}(32-s))`) + OpBitwiseOr}
\par\textbf{Pseudo-GLSL (Operational Semantics):}
\codebox{for (int i = 0; i \textless{} N; ++i) \{\\    uint s = extractByte(src1, lanes1, i) \& (W - 1);\\    auto r = ((src0[i] \textgreater{}\textgreater{} s) \textbar{} (src0[i] \textless{}\textless{} ((W - s) \& (W - 1)))) \textbar{} (flag\_not2 ? \textasciitilde{}src2[i] : src2[i]);\\    dest[i] = flag\_not\_result ? \textasciitilde{}r : r;\\\}}
\textbf{Execution Context:}
\begin{tightitemize}\footnotesize
  \item \textbf{Compute stage preloads (\hyperref[sec:csf-spd]{\texttt{SPD\_k}}, \texttt{r16..r23}):} Hardware pre-fills \texttt{r55..r57} with \texttt{WorkGroupID.xyz + JOB\_OFFSET} (from \hyperref[sec:csf-regs]{\texttt{COMPUTESR::JOB\_OFFSET\_X/Y/Z}} \texttt{r34..r36}) and \texttt{r60..r62} with \texttt{LocalInvocationID.xyz} or \texttt{GlobalInvocationID.xyz} (from \texttt{WG\_SIZE} \texttt{r33} \& \texttt{JOB\_SIZE} \texttt{r37..r39}, selected by \texttt{ShaderProgram::preload}) so compute shaders can derive thread/element indices without memory loads.
  \item \textbf{Vertex stage preloads (\hyperref[sec:csf-spd]{\texttt{SPD\_0}/\texttt{SPD\_1}}, \texttt{r16..r19}):} The hardware Index Fetcher reads \hyperref[sec:csf-regs]{\texttt{IDVSSR::INDEX\_BUFFER}} (\texttt{r54:r55}) and pre-fills \texttt{r59, r60} with \texttt{VertexIndex} (\texttt{+ VERTEX\_OFFSET + GLOBAL\_ATTR\_OFF}) and \texttt{r61} with \texttt{InstanceIndex} (\texttt{+ INSTANCE\_OFFSET}) so Position and Varying shaders can pull vertex attributes directly.
  \item \textbf{Inline FAU \& Special Operands (\hyperref[sec:csf-tsd-fau]{\texttt{FAU\_k}}, \texttt{r8..r15}):} Any 8-bit source with \texttt{mode=2} reads zero-latency uniform \texttt{u0..u31} from \texttt{FAU\_k} (\texttt{slots64 = bits[63:56]}); \texttt{mode=3} reads \texttt{SmallConstant} (Table~1.2) or \texttt{FauSpecial} (\texttt{pc}, \texttt{lane\_id}, \texttt{lane\_width}, \texttt{tls\_ptr}, \texttt{wls\_ptr}, \texttt{blend\_descriptor\_0..7}).
\end{tightitemize}
\par\textbf{Worked Examples:}
\par\textbf{Example 1 (Standard fused shift + bitwise op):}
\codebox{\mnem{RROT\_OR.i32} \dst{r0}, \src{r1}, \src{r2}\imod{.b0}, \src{r3}}
With \texttt{r1 = 0x00000003}, \texttt{r2.b0 = 1} (masked mod 32): shift yields \texttt{0x80000001}, then \texttt{\textbar{} 0x0000000F} yields \texttt{r0 = 0x8000000F}.
\par\textbf{Example 2 (Modulo-32 shift count \& .not\_result):}
\codebox{\mnem{RROT\_OR.i32}\imod{.not\_result} \dst{r0}, \src{r1}, \src{r2}\imod{.b0}, \src{r3}}
With \texttt{r1 = 0x12345678}, \texttt{r2.b0 = 8} (effective shift \texttt{8}): shift yields \texttt{0x78123456}, \texttt{\textbar{} 0x0000000F} yields \texttt{0x7812345F}, and \texttt{.not\_result} inverts -\textgreater{} \texttt{r0 = 0x87EDCBA0}.

\begin{center}\footnotesize
\begin{tabularx}{\linewidth}{@{}lllX@{}}
\toprule
Variant & Arch & Class & Opcode Fields \\
\midrule
\texttt{(none)} & v9--v14 & \texttt{R3} & \texttt{opcode1=SHIFT\_BITWISE\_I32} \\
\texttt{(none)} & v9--v14 & \texttt{R3} & \texttt{opcode1=LOGIC\_SHIFT\_I32} \\
\texttt{(none)} & v9--v14 & \texttt{D3} & \texttt{opcode1=SHIFT\_BITWISE\_I64} \\
\texttt{(none)} & v9--v14 & \texttt{D3} & \texttt{opcode1=LOGIC\_SHIFT\_I64} \\
\texttt{(none)} & v9--v14 & \texttt{R3} & \texttt{opcode1=SHIFT\_BITWISE\_V2I16} \\
\texttt{(none)} & v9--v14 & \texttt{R3} & \texttt{opcode1=LOGIC\_SHIFT\_V2I16} \\
\texttt{(none)} & v9--v14 & \texttt{R3} & \texttt{opcode1=SHIFT\_BITWISE\_V4I8} \\
\texttt{(none)} & v9--v14 & \texttt{R3} & \texttt{opcode1=LOGIC\_SHIFT\_V4I8} \\
\bottomrule
\end{tabularx}
\end{center}
\begin{center}\footnotesize
\begin{tabularx}{0.88\linewidth}{@{}l l X@{}}
\toprule
Field Name & Bits & Type / Fixed Value \\
\midrule
\texttt{src0} & \texttt{0:7} & \texttt{SourceEncoding} \\
\texttt{src1} & \texttt{8:15} & \texttt{SourceEncoding} \\
\texttt{src2} & \texttt{16:23} & \texttt{SourceEncoding} \\
\texttt{bitwise\_operation} & \texttt{24:25} & \texttt{exact=or} \\
\hyperref[enum:shift_bitwise_i32_lanes_m]{\texttt{lanes1}} & \texttt{26:29} & \texttt{shift\_bitwise\_i32\_lanes\_m (def b0)} \\
\hyperref[enum:not_m]{\texttt{not\_result}} & \texttt{30} & \texttt{not\_m (def none)} \\
\textit{Reserved} & \texttt{31} & \texttt{zero} \\
\texttt{shift\_operation} & \texttt{32:34} & \texttt{exact=rrot} \\
\hyperref[enum:not_m]{\texttt{not2}} & \texttt{35} & \texttt{not\_m (def none)} \\
\hyperref[enum:widen_int_m]{\texttt{widen0}} & \texttt{36:39} & \texttt{widen\_int\_m (def none)} \\
\texttt{dest\_register\_index} & \texttt{40:45} & \texttt{uint} \\
\hyperref[enum:dest_width_m]{\texttt{dest\_width}} & \texttt{46:47} & \texttt{dest\_width\_m} \\
\texttt{opcode1} & \texttt{48:51} & \texttt{exact=SHIFT\_BITWISE\_I32} \\
\texttt{instruction\_class} & \texttt{52:56} & \texttt{exact=R3} \\
\texttt{fau\_page\_index} & \texttt{57:58} & \texttt{uint} \\
\hyperref[enum:flow_control_m]{\texttt{flow}} & \texttt{59:62} & \texttt{flow\_control\_m} \\
\textit{Reserved} & \texttt{63} & \texttt{zero} \\
\bottomrule
\end{tabularx}
\end{center}
\textbf{Fields Explained:}
\begin{tightitemize}\footnotesize
  \item \texttt{src0}: Primary source operand (\texttt{r0..r63}, \texttt{u0..u63} uniform, \texttt{k*} FAU constant, or \texttt{t0..t2} temporary; \texttt{\textasciicircum{}} marks discard/last-use). \emph{Example:} {\scriptsize\ttfamily \mnem{FADD.f32} \dst{r0}, \src{r1}, \src{r2}} (\textbf{SPIR-V:} \texttt{OpFAdd (Operand 0)}; \textbf{GLSL:} \texttt{float c = a + b;})
  \item \texttt{src1}: Secondary source operand (register, uniform, or FAU constant; in shifts, supplies the shift count byte). \emph{Example:} {\scriptsize\ttfamily \mnem{LSHIFT.i32} \dst{r0}, \src{r1}, \src{r2}\imod{.b0}} (\textbf{SPIR-V:} \texttt{OpShiftLeftLogical (Shift)}; \textbf{GLSL:} \texttt{uint y = a \textless{}\textless{} s;})
  \item \texttt{src2}: Third source operand (FMA addend, \texttt{MUX}/\texttt{CSEL} selector, \texttt{IADDC} carry-in, or bitwise mask in fused shift-logic). \emph{Example:} {\scriptsize\ttfamily \mnem{FMA.f32} \dst{r0}, \src{r1}, \src{r2}, \src{r3}} (\textbf{SPIR-V:} \texttt{GLSL.std.450 Fma (Addend)}; \textbf{GLSL:} \texttt{float d = fma(a, b, c);})
  \item \texttt{bitwise\_operation}: Selects the bitwise ALU op (\texttt{and}, \texttt{or}, \texttt{xor}, \texttt{lshift\_and}, etc.) fused with shifts. \emph{Example:} {\scriptsize\ttfamily \mnem{LSHIFT\_OR.i32} \dst{r0}, \src{r1}, \src{r2}, \src{r3}} (\textbf{SPIR-V:} \texttt{OpShiftLeftLogical + OpBitwiseOr}; \textbf{GLSL:} \texttt{uint y = (a \textless{}\textless{} s) \textbar{} b;})
  \item \hyperref[enum:shift_bitwise_i32_lanes_m]{\texttt{lanes1}}: Stage A byte-lane selector on \texttt{src1} (selects shift-count byte \texttt{b0..b3} on 32/64-bit shifts, \texttt{bXY} on \texttt{v2i16} shifts, or 4-byte swizzle on \texttt{v4*8}). \emph{Example:} {\scriptsize\ttfamily \mnem{LSHIFT.i32} \dst{r0}, \src{r1}, \src{r2}\imod{.b1}} (\textbf{SPIR-V:} \texttt{OpBitFieldUExtract + OpShiftLeftLogical}; \textbf{GLSL:} \texttt{uint y = a \textless{}\textless{} ((s \textgreater{}\textgreater{} 8) \& 0xFFu);})
  \item \hyperref[enum:not_m]{\texttt{not\_result}}: Bitwise-inverts the final output of a bitwise/shift-bitwise instruction (\textasciitilde{}result). \emph{Example:} {\scriptsize\ttfamily \mnem{AND.i32}\imod{.not\_result} \dst{r0}, \src{r1}, \src{r2}} (\textbf{SPIR-V:} \texttt{OpBitwiseAnd + OpNot}; \textbf{GLSL:} \texttt{uint y = \textasciitilde{}(a \& b);})
  \item \texttt{shift\_operation}: Selects shift direction/type (\texttt{lshift}, \texttt{rshift}, \texttt{arshift}) in fused shift-bitwise instructions. \emph{Example:} {\scriptsize\ttfamily \mnem{RSHIFT\_AND.i32} \dst{r0}, \src{r1}, \src{r2}, \src{r3}} (\textbf{SPIR-V:} \texttt{OpShiftRightLogical + OpBitwiseAnd}; \textbf{GLSL:} \texttt{uint y = (a \textgreater{}\textgreater{} s) \& m;})
  \item \hyperref[enum:not_m]{\texttt{not2}}: Stage A bitwise complement (\texttt{\textasciitilde{}src2}) applied to \texttt{src2} before the bitwise/shift-logic operation. \emph{Example:} {\scriptsize\ttfamily \mnem{AND.i32} \dst{r0}, \src{r1}, \src{r2}\imod{.not}} (\textbf{SPIR-V:} \texttt{OpNot + OpBitwiseAnd}; \textbf{GLSL:} \texttt{uint y = a \& \textasciitilde{}b;})
  \item \hyperref[enum:widen_int_m]{\texttt{widen0}}: Stage A widening extractor on \texttt{src0} (promotes FP16 \texttt{h0}/\texttt{h1} to FP32, or sign/zero-extends \texttt{b0..b3}/\texttt{h0..h1}/\texttt{w0..w1} to 32/64-bit before \texttt{abs}/\texttt{neg}). \emph{Example:} {\scriptsize\ttfamily \mnem{IADD.s32} \dst{r0}, \src{r1}\imod{.h0}, \src{r2}} (\textbf{SPIR-V:} \texttt{OpSConvert / OpFConvert}; \textbf{GLSL:} \texttt{int c = int(int16\_t(a)) + b;})
  \item \texttt{dest\_register\_index}: 6-bit destination GPR index (\texttt{r0..r63}) receiving the Stage C writeback result (even-aligned pair \texttt{[rN:rN+1]} for 64-bit ops). \emph{Example:} {\scriptsize\ttfamily \mnem{FMA.f32} \dst{r0}, \src{r1}, \src{r2}, \src{r3}} (\textbf{SPIR-V:} \texttt{SSA Result ID}; \textbf{GLSL:} \texttt{float r0 = fma(r1, r2, r3);})
  \item \hyperref[enum:dest_width_m]{\texttt{dest\_width}}: Stage C destination lane write-mask (\texttt{none} = full 32-bit, \texttt{h0} = low 16-bit, \texttt{h1} = high 16-bit without bit-shifting). \emph{Example:} {\scriptsize\ttfamily \mnem{FADD.v2f16} \dst{r0}\imod{.h1}, \src{r1}, \src{r2}} (\textbf{SPIR-V:} \texttt{OpVectorInsertDynamic / WriteMask}; \textbf{GLSL:} \texttt{dst.y = a.y + b.y;})
  \item \texttt{opcode1}: Primary sub-opcode within the instruction class uniquely identifying the operation family. \emph{Example:} {\scriptsize\ttfamily \mnem{IADD.s32} \dst{r0}, \src{r1}, \src{r2}} (\textbf{SPIR-V:} \texttt{OpIAdd (operation selector)}; \textbf{GLSL:} \texttt{// Primary opcode bits})
  \item \texttt{instruction\_class}: 4-bit major instruction class selector (\texttt{bits 48:51}) partitioning the 64-bit Valhall opcode space (e.g., \texttt{FMA}, \texttt{CVT}, \texttt{SFU}, \texttt{LS}, \texttt{TEX}, \texttt{BRANCH}). \emph{Example:} {\scriptsize\ttfamily \mnem{FADD.f32} \dst{r0}, \src{r1}, \src{r2}} (\textbf{SPIR-V:} \texttt{N/A (encoding class)}; \textbf{GLSL:} \texttt{// Hardware execution pipe / decode class})
  \item \texttt{fau\_page\_index}: Selects the FAU (Fast Access Uniform) page (\texttt{0..3}) when a source operand references hardware small constants, thread IDs, or lane constants. \emph{Example:} {\scriptsize\ttfamily \mnem{FADD.f32} \dst{r0}, \src{r1}, \src{k1f}} (\textbf{SPIR-V:} \texttt{OpConstant / BuiltIn}; \textbf{GLSL:} \texttt{float y = x + 1.0;})
  \item \hyperref[enum:flow_control_m]{\texttt{flow}}: 4-bit embedded flow-control \& scoreboard wait modifier (\texttt{bits 59:62}: \texttt{none}, \texttt{wait0..wait0126}, \texttt{wait}, \texttt{reconverge}, \texttt{discard}, \texttt{end}). \emph{Example:} {\scriptsize\ttfamily \mnem{FADD.f32}\imod{.wait0.reconverge} \dst{r0}, \src{r1}, \src{r2}} (\textbf{SPIR-V:} \texttt{OpControlBarrier / OpReturn}; \textbf{GLSL:} \texttt{// Scoreboard wait + warp reconvergence})
\end{tightitemize}
\clearpage

\subsection*{RROT\_XOR}
\addcontentsline{toc}{subsection}{RROT\_XOR}
For each SIMD element $i \in \{0, \dots, N-1\}$ ($N=1$ for \texttt{i32}/\texttt{i64}, $N=2$ for \texttt{v2i16}, $N=4$ for \texttt{v4i8} on v9--v10, lane width $W$), computes rotate right \texttt{s = rotate\_right(val, amt \& (W - 1))} where \texttt{val = src0[i]} and the 8-bit shift count \texttt{amt} is extracted from \texttt{src1} via \texttt{lanes1}, then combines with \texttt{src2[i]} via bitwise XOR \texttt{res = s \textasciicircum{} (not2 ? \textasciitilde{}src2[i] : src2[i])} and writes \texttt{not\_result ? \textasciitilde{}res : res} to \texttt{dest[i]}.
\par\textbf{Notes:}
\begin{tightitemize}
  \item Shift amount is masked modulo lane width (\texttt{wrapping\_shl}/\texttt{wrapping\_shr}/\texttt{rotate\_*}, \texttt{kraid/ops.rs:3984-4009}).
  \item On \texttt{v2i16}, \texttt{src1} is \texttt{V2I8} (\texttt{ops.rs:3953}) with byte lanes selected by \texttt{lanes1} (default \texttt{b02}: byte 0 -\textgreater{} lane 0, byte 2 -\textgreater{} lane 1).
  \item \textbf{Source references:} \texttt{kraid/ops.rs:3878-4025}, \texttt{encode\_v9.rs:2940-3093}, \texttt{hw\_tests.rs:2148-2198}.
\end{tightitemize}
\codebox{\mnem{RROT\_XOR}\imod{\{.flow\}}\imod{\{.not\_result\}} \dst{dest}\imod{\{.dest\_width\}}, \src{src0}\imod{\{.widen0\}}, \src{src1}\imod{\{.lanes1\}}, \src{src2}\imod{\{.not2\}}\\\mnem{RROT\_XOR}\imod{\{.flow\}}\imod{\{.not\_result\}} \dst{dest}, \src{src0}\imod{\{.widen0\}}, \src{src1}\imod{\{.lanes1\}}, \src{src2}\imod{\{.not2\}}\\\mnem{RROT\_XOR}\imod{\{.flow\}}\imod{\{.not\_result\}} \dst{dest}\imod{\{.dest\_width\}}, \src{src0}\imod{\{.swz0\}}, \src{src1}\imod{\{.lanes1\}}, \src{src2}\imod{\{.not2\}}\\\mnem{RROT\_XOR}\imod{\{.flow\}}\imod{\{.not\_result\}} \dst{dest}\imod{\{.dest\_width\}}, \src{src0}\imod{\{.lanes0\}}, \src{src1}\imod{\{.lanes1\}}, \src{src2}\imod{\{.not2\}}}
\textbf{Operands:}
\begin{tightitemize}\footnotesize
  \item \texttt{\dst{dest}} (\texttt{dest\_register\_index}, \texttt{dest\_width}): Destination GPR receiving the rotate right combined via bitwise XOR (inverted if \texttt{.not\_result}).
  \item \texttt{\src{src0}}: Input integer value to shift/rotate.
  \item \texttt{\src{src1}}: 8-bit shift/rotate count per lane (selected via \texttt{lanes1}, masked modulo lane width \texttt{W}).
  \item \texttt{\src{src2}}: Second operand combined via bitwise XOR after shifting (inverted if \texttt{.not2}).
\end{tightitemize}
\par\textbf{Closest SPIR-V Op:} \texttt{None (lowered from `(x\textgreater{}\textgreater{}s)\textbar{}(x\textless{}\textless{}(32-s))`) + OpBitwiseXor}
\par\textbf{Pseudo-GLSL (Operational Semantics):}
\codebox{for (int i = 0; i \textless{} N; ++i) \{\\    uint s = extractByte(src1, lanes1, i) \& (W - 1);\\    auto r = ((src0[i] \textgreater{}\textgreater{} s) \textbar{} (src0[i] \textless{}\textless{} ((W - s) \& (W - 1)))) \textasciicircum{} (flag\_not2 ? \textasciitilde{}src2[i] : src2[i]);\\    dest[i] = flag\_not\_result ? \textasciitilde{}r : r;\\\}}
\textbf{Execution Context:}
\begin{tightitemize}\footnotesize
  \item \textbf{Compute stage preloads (\hyperref[sec:csf-spd]{\texttt{SPD\_k}}, \texttt{r16..r23}):} Hardware pre-fills \texttt{r55..r57} with \texttt{WorkGroupID.xyz + JOB\_OFFSET} (from \hyperref[sec:csf-regs]{\texttt{COMPUTESR::JOB\_OFFSET\_X/Y/Z}} \texttt{r34..r36}) and \texttt{r60..r62} with \texttt{LocalInvocationID.xyz} or \texttt{GlobalInvocationID.xyz} (from \texttt{WG\_SIZE} \texttt{r33} \& \texttt{JOB\_SIZE} \texttt{r37..r39}, selected by \texttt{ShaderProgram::preload}) so compute shaders can derive thread/element indices without memory loads.
  \item \textbf{Vertex stage preloads (\hyperref[sec:csf-spd]{\texttt{SPD\_0}/\texttt{SPD\_1}}, \texttt{r16..r19}):} The hardware Index Fetcher reads \hyperref[sec:csf-regs]{\texttt{IDVSSR::INDEX\_BUFFER}} (\texttt{r54:r55}) and pre-fills \texttt{r59, r60} with \texttt{VertexIndex} (\texttt{+ VERTEX\_OFFSET + GLOBAL\_ATTR\_OFF}) and \texttt{r61} with \texttt{InstanceIndex} (\texttt{+ INSTANCE\_OFFSET}) so Position and Varying shaders can pull vertex attributes directly.
  \item \textbf{Inline FAU \& Special Operands (\hyperref[sec:csf-tsd-fau]{\texttt{FAU\_k}}, \texttt{r8..r15}):} Any 8-bit source with \texttt{mode=2} reads zero-latency uniform \texttt{u0..u31} from \texttt{FAU\_k} (\texttt{slots64 = bits[63:56]}); \texttt{mode=3} reads \texttt{SmallConstant} (Table~1.2) or \texttt{FauSpecial} (\texttt{pc}, \texttt{lane\_id}, \texttt{lane\_width}, \texttt{tls\_ptr}, \texttt{wls\_ptr}, \texttt{blend\_descriptor\_0..7}).
\end{tightitemize}
\par\textbf{Worked Examples:}
\par\textbf{Example 1 (Standard fused shift + bitwise op):}
\codebox{\mnem{RROT\_XOR.i32} \dst{r0}, \src{r1}, \src{r2}\imod{.b0}, \src{r3}}
With \texttt{r1 = 0x00000003}, \texttt{r2.b0 = 1} (masked mod 32): shift yields \texttt{0x80000001}, then \texttt{\textasciicircum{} 0x000000FF} yields \texttt{r0 = 0x800000FE}.
\par\textbf{Example 2 (Modulo-32 shift count \& .not\_result):}
\codebox{\mnem{RROT\_XOR.i32}\imod{.not\_result} \dst{r0}, \src{r1}, \src{r2}\imod{.b0}, \src{r3}}
With \texttt{r1 = 0x12345678}, \texttt{r2.b0 = 8} (effective shift \texttt{8}): shift yields \texttt{0x78123456}, \texttt{\textasciicircum{} 0x000000FF} yields \texttt{0x781234A9}, and \texttt{.not\_result} inverts -\textgreater{} \texttt{r0 = 0x87EDCB56}.

\begin{center}\footnotesize
\begin{tabularx}{\linewidth}{@{}lllX@{}}
\toprule
Variant & Arch & Class & Opcode Fields \\
\midrule
\texttt{(none)} & v9--v14 & \texttt{R3} & \texttt{opcode1=SHIFT\_BITWISE\_I32} \\
\texttt{(none)} & v9--v14 & \texttt{R3} & \texttt{opcode1=LOGIC\_SHIFT\_I32} \\
\texttt{(none)} & v9--v14 & \texttt{D3} & \texttt{opcode1=SHIFT\_BITWISE\_I64} \\
\texttt{(none)} & v9--v14 & \texttt{D3} & \texttt{opcode1=LOGIC\_SHIFT\_I64} \\
\texttt{(none)} & v9--v14 & \texttt{R3} & \texttt{opcode1=SHIFT\_BITWISE\_V2I16} \\
\texttt{(none)} & v9--v14 & \texttt{R3} & \texttt{opcode1=LOGIC\_SHIFT\_V2I16} \\
\texttt{(none)} & v9--v14 & \texttt{R3} & \texttt{opcode1=SHIFT\_BITWISE\_V4I8} \\
\texttt{(none)} & v9--v14 & \texttt{R3} & \texttt{opcode1=LOGIC\_SHIFT\_V4I8} \\
\bottomrule
\end{tabularx}
\end{center}
\begin{center}\footnotesize
\begin{tabularx}{0.88\linewidth}{@{}l l X@{}}
\toprule
Field Name & Bits & Type / Fixed Value \\
\midrule
\texttt{src0} & \texttt{0:7} & \texttt{SourceEncoding} \\
\texttt{src1} & \texttt{8:15} & \texttt{SourceEncoding} \\
\texttt{src2} & \texttt{16:23} & \texttt{SourceEncoding} \\
\texttt{bitwise\_operation} & \texttt{24:25} & \texttt{exact=xor} \\
\hyperref[enum:shift_bitwise_i32_lanes_m]{\texttt{lanes1}} & \texttt{26:29} & \texttt{shift\_bitwise\_i32\_lanes\_m (def b0)} \\
\hyperref[enum:not_m]{\texttt{not\_result}} & \texttt{30} & \texttt{not\_m (def none)} \\
\textit{Reserved} & \texttt{31} & \texttt{zero} \\
\texttt{shift\_operation} & \texttt{32:34} & \texttt{exact=rrot} \\
\hyperref[enum:not_m]{\texttt{not2}} & \texttt{35} & \texttt{not\_m (def none)} \\
\hyperref[enum:widen_int_m]{\texttt{widen0}} & \texttt{36:39} & \texttt{widen\_int\_m (def none)} \\
\texttt{dest\_register\_index} & \texttt{40:45} & \texttt{uint} \\
\hyperref[enum:dest_width_m]{\texttt{dest\_width}} & \texttt{46:47} & \texttt{dest\_width\_m} \\
\texttt{opcode1} & \texttt{48:51} & \texttt{exact=SHIFT\_BITWISE\_I32} \\
\texttt{instruction\_class} & \texttt{52:56} & \texttt{exact=R3} \\
\texttt{fau\_page\_index} & \texttt{57:58} & \texttt{uint} \\
\hyperref[enum:flow_control_m]{\texttt{flow}} & \texttt{59:62} & \texttt{flow\_control\_m} \\
\textit{Reserved} & \texttt{63} & \texttt{zero} \\
\bottomrule
\end{tabularx}
\end{center}
\textbf{Fields Explained:}
\begin{tightitemize}\footnotesize
  \item \texttt{src0}: Primary source operand (\texttt{r0..r63}, \texttt{u0..u63} uniform, \texttt{k*} FAU constant, or \texttt{t0..t2} temporary; \texttt{\textasciicircum{}} marks discard/last-use). \emph{Example:} {\scriptsize\ttfamily \mnem{FADD.f32} \dst{r0}, \src{r1}, \src{r2}} (\textbf{SPIR-V:} \texttt{OpFAdd (Operand 0)}; \textbf{GLSL:} \texttt{float c = a + b;})
  \item \texttt{src1}: Secondary source operand (register, uniform, or FAU constant; in shifts, supplies the shift count byte). \emph{Example:} {\scriptsize\ttfamily \mnem{LSHIFT.i32} \dst{r0}, \src{r1}, \src{r2}\imod{.b0}} (\textbf{SPIR-V:} \texttt{OpShiftLeftLogical (Shift)}; \textbf{GLSL:} \texttt{uint y = a \textless{}\textless{} s;})
  \item \texttt{src2}: Third source operand (FMA addend, \texttt{MUX}/\texttt{CSEL} selector, \texttt{IADDC} carry-in, or bitwise mask in fused shift-logic). \emph{Example:} {\scriptsize\ttfamily \mnem{FMA.f32} \dst{r0}, \src{r1}, \src{r2}, \src{r3}} (\textbf{SPIR-V:} \texttt{GLSL.std.450 Fma (Addend)}; \textbf{GLSL:} \texttt{float d = fma(a, b, c);})
  \item \texttt{bitwise\_operation}: Selects the bitwise ALU op (\texttt{and}, \texttt{or}, \texttt{xor}, \texttt{lshift\_and}, etc.) fused with shifts. \emph{Example:} {\scriptsize\ttfamily \mnem{LSHIFT\_OR.i32} \dst{r0}, \src{r1}, \src{r2}, \src{r3}} (\textbf{SPIR-V:} \texttt{OpShiftLeftLogical + OpBitwiseOr}; \textbf{GLSL:} \texttt{uint y = (a \textless{}\textless{} s) \textbar{} b;})
  \item \hyperref[enum:shift_bitwise_i32_lanes_m]{\texttt{lanes1}}: Stage A byte-lane selector on \texttt{src1} (selects shift-count byte \texttt{b0..b3} on 32/64-bit shifts, \texttt{bXY} on \texttt{v2i16} shifts, or 4-byte swizzle on \texttt{v4*8}). \emph{Example:} {\scriptsize\ttfamily \mnem{LSHIFT.i32} \dst{r0}, \src{r1}, \src{r2}\imod{.b1}} (\textbf{SPIR-V:} \texttt{OpBitFieldUExtract + OpShiftLeftLogical}; \textbf{GLSL:} \texttt{uint y = a \textless{}\textless{} ((s \textgreater{}\textgreater{} 8) \& 0xFFu);})
  \item \hyperref[enum:not_m]{\texttt{not\_result}}: Bitwise-inverts the final output of a bitwise/shift-bitwise instruction (\textasciitilde{}result). \emph{Example:} {\scriptsize\ttfamily \mnem{AND.i32}\imod{.not\_result} \dst{r0}, \src{r1}, \src{r2}} (\textbf{SPIR-V:} \texttt{OpBitwiseAnd + OpNot}; \textbf{GLSL:} \texttt{uint y = \textasciitilde{}(a \& b);})
  \item \texttt{shift\_operation}: Selects shift direction/type (\texttt{lshift}, \texttt{rshift}, \texttt{arshift}) in fused shift-bitwise instructions. \emph{Example:} {\scriptsize\ttfamily \mnem{RSHIFT\_AND.i32} \dst{r0}, \src{r1}, \src{r2}, \src{r3}} (\textbf{SPIR-V:} \texttt{OpShiftRightLogical + OpBitwiseAnd}; \textbf{GLSL:} \texttt{uint y = (a \textgreater{}\textgreater{} s) \& m;})
  \item \hyperref[enum:not_m]{\texttt{not2}}: Stage A bitwise complement (\texttt{\textasciitilde{}src2}) applied to \texttt{src2} before the bitwise/shift-logic operation. \emph{Example:} {\scriptsize\ttfamily \mnem{AND.i32} \dst{r0}, \src{r1}, \src{r2}\imod{.not}} (\textbf{SPIR-V:} \texttt{OpNot + OpBitwiseAnd}; \textbf{GLSL:} \texttt{uint y = a \& \textasciitilde{}b;})
  \item \hyperref[enum:widen_int_m]{\texttt{widen0}}: Stage A widening extractor on \texttt{src0} (promotes FP16 \texttt{h0}/\texttt{h1} to FP32, or sign/zero-extends \texttt{b0..b3}/\texttt{h0..h1}/\texttt{w0..w1} to 32/64-bit before \texttt{abs}/\texttt{neg}). \emph{Example:} {\scriptsize\ttfamily \mnem{IADD.s32} \dst{r0}, \src{r1}\imod{.h0}, \src{r2}} (\textbf{SPIR-V:} \texttt{OpSConvert / OpFConvert}; \textbf{GLSL:} \texttt{int c = int(int16\_t(a)) + b;})
  \item \texttt{dest\_register\_index}: 6-bit destination GPR index (\texttt{r0..r63}) receiving the Stage C writeback result (even-aligned pair \texttt{[rN:rN+1]} for 64-bit ops). \emph{Example:} {\scriptsize\ttfamily \mnem{FMA.f32} \dst{r0}, \src{r1}, \src{r2}, \src{r3}} (\textbf{SPIR-V:} \texttt{SSA Result ID}; \textbf{GLSL:} \texttt{float r0 = fma(r1, r2, r3);})
  \item \hyperref[enum:dest_width_m]{\texttt{dest\_width}}: Stage C destination lane write-mask (\texttt{none} = full 32-bit, \texttt{h0} = low 16-bit, \texttt{h1} = high 16-bit without bit-shifting). \emph{Example:} {\scriptsize\ttfamily \mnem{FADD.v2f16} \dst{r0}\imod{.h1}, \src{r1}, \src{r2}} (\textbf{SPIR-V:} \texttt{OpVectorInsertDynamic / WriteMask}; \textbf{GLSL:} \texttt{dst.y = a.y + b.y;})
  \item \texttt{opcode1}: Primary sub-opcode within the instruction class uniquely identifying the operation family. \emph{Example:} {\scriptsize\ttfamily \mnem{IADD.s32} \dst{r0}, \src{r1}, \src{r2}} (\textbf{SPIR-V:} \texttt{OpIAdd (operation selector)}; \textbf{GLSL:} \texttt{// Primary opcode bits})
  \item \texttt{instruction\_class}: 4-bit major instruction class selector (\texttt{bits 48:51}) partitioning the 64-bit Valhall opcode space (e.g., \texttt{FMA}, \texttt{CVT}, \texttt{SFU}, \texttt{LS}, \texttt{TEX}, \texttt{BRANCH}). \emph{Example:} {\scriptsize\ttfamily \mnem{FADD.f32} \dst{r0}, \src{r1}, \src{r2}} (\textbf{SPIR-V:} \texttt{N/A (encoding class)}; \textbf{GLSL:} \texttt{// Hardware execution pipe / decode class})
  \item \texttt{fau\_page\_index}: Selects the FAU (Fast Access Uniform) page (\texttt{0..3}) when a source operand references hardware small constants, thread IDs, or lane constants. \emph{Example:} {\scriptsize\ttfamily \mnem{FADD.f32} \dst{r0}, \src{r1}, \src{k1f}} (\textbf{SPIR-V:} \texttt{OpConstant / BuiltIn}; \textbf{GLSL:} \texttt{float y = x + 1.0;})
  \item \hyperref[enum:flow_control_m]{\texttt{flow}}: 4-bit embedded flow-control \& scoreboard wait modifier (\texttt{bits 59:62}: \texttt{none}, \texttt{wait0..wait0126}, \texttt{wait}, \texttt{reconverge}, \texttt{discard}, \texttt{end}). \emph{Example:} {\scriptsize\ttfamily \mnem{FADD.f32}\imod{.wait0.reconverge} \dst{r0}, \src{r1}, \src{r2}} (\textbf{SPIR-V:} \texttt{OpControlBarrier / OpReturn}; \textbf{GLSL:} \texttt{// Scoreboard wait + warp reconvergence})
\end{tightitemize}
\clearpage

\subsection*{RSHIFT}
\addcontentsline{toc}{subsection}{RSHIFT}
For each SIMD element $i \in \{0, \dots, N-1\}$ ($N=1$ for \texttt{i32}/\texttt{i64}, $N=2$ for \texttt{v2i16}, $N=4$ for \texttt{v4i8} on v9--v10, lane width $W$), computes logical right shift \texttt{s = (uint\_W(val) \textgreater{}\textgreater{} (amt \& (W - 1)))} where \texttt{val = src0[i]} and the 8-bit shift count \texttt{amt} is extracted from \texttt{src1} via \texttt{lanes1}; writes \texttt{not\_result ? \textasciitilde{}s : s} to \texttt{dest[i]}.
\par\textbf{Notes:}
\begin{tightitemize}
  \item Shift amount is masked modulo lane width (\texttt{wrapping\_shl}/\texttt{wrapping\_shr}/\texttt{rotate\_*}, \texttt{kraid/ops.rs:3984-4009}).
  \item On \texttt{v2i16}, \texttt{src1} is \texttt{V2I8} (\texttt{ops.rs:3953}) with byte lanes selected by \texttt{lanes1} (default \texttt{b02}: byte 0 -\textgreater{} lane 0, byte 2 -\textgreater{} lane 1).
  \item \textbf{Source references:} \texttt{kraid/ops.rs:3878-4025}, \texttt{encode\_v9.rs:2940-3093}, \texttt{hw\_tests.rs:2148-2198}.
\end{tightitemize}
\codebox{\mnem{RSHIFT}\imod{\{.flow\}}\imod{\{.not\_result\}} \dst{dest}\imod{\{.dest\_width\}}, \src{src0}\imod{\{.widen0\}}, \src{src1}\imod{\{.lanes1\}}\\\mnem{RSHIFT}\imod{\{.flow\}}\imod{\{.not\_result\}} \dst{dest}, \src{src0}\imod{\{.widen0\}}, \src{src1}\imod{\{.lanes1\}}\\\mnem{RSHIFT}\imod{\{.flow\}}\imod{\{.not\_result\}} \dst{dest}\imod{\{.dest\_width\}}, \src{src0}\imod{\{.swz0\}}, \src{src1}\imod{\{.lanes1\}}\\\mnem{RSHIFT}\imod{\{.flow\}}\imod{\{.not\_result\}} \dst{dest}\imod{\{.dest\_width\}}, \src{src0}\imod{\{.lanes0\}}, \src{src1}\imod{\{.lanes1\}}}
\textbf{Operands:}
\begin{tightitemize}\footnotesize
  \item \texttt{\dst{dest}} (\texttt{dest\_register\_index}, \texttt{dest\_width}): Destination GPR receiving the logical right shift (inverted if \texttt{.not\_result}).
  \item \texttt{\src{src0}}: Input integer value to shift/rotate.
  \item \texttt{\src{src1}}: 8-bit shift/rotate count per lane (selected via \texttt{lanes1}, masked modulo lane width \texttt{W}).
\end{tightitemize}
\par\textbf{Closest SPIR-V Op:} \texttt{OpShiftRightLogical}
\par\textbf{Pseudo-GLSL (Operational Semantics):}
\codebox{for (int i = 0; i \textless{} N; ++i) \{\\    uint s = extractByte(src1, lanes1, i) \& (W - 1);\\    auto r = (uint(src0[i]) \textgreater{}\textgreater{} s);\\    dest[i] = flag\_not\_result ? \textasciitilde{}r : r;\\\}}
\textbf{Execution Context:}
\begin{tightitemize}\footnotesize
  \item \textbf{Compute stage preloads (\hyperref[sec:csf-spd]{\texttt{SPD\_k}}, \texttt{r16..r23}):} Hardware pre-fills \texttt{r55..r57} with \texttt{WorkGroupID.xyz + JOB\_OFFSET} (from \hyperref[sec:csf-regs]{\texttt{COMPUTESR::JOB\_OFFSET\_X/Y/Z}} \texttt{r34..r36}) and \texttt{r60..r62} with \texttt{LocalInvocationID.xyz} or \texttt{GlobalInvocationID.xyz} (from \texttt{WG\_SIZE} \texttt{r33} \& \texttt{JOB\_SIZE} \texttt{r37..r39}, selected by \texttt{ShaderProgram::preload}) so compute shaders can derive thread/element indices without memory loads.
  \item \textbf{Vertex stage preloads (\hyperref[sec:csf-spd]{\texttt{SPD\_0}/\texttt{SPD\_1}}, \texttt{r16..r19}):} The hardware Index Fetcher reads \hyperref[sec:csf-regs]{\texttt{IDVSSR::INDEX\_BUFFER}} (\texttt{r54:r55}) and pre-fills \texttt{r59, r60} with \texttt{VertexIndex} (\texttt{+ VERTEX\_OFFSET + GLOBAL\_ATTR\_OFF}) and \texttt{r61} with \texttt{InstanceIndex} (\texttt{+ INSTANCE\_OFFSET}) so Position and Varying shaders can pull vertex attributes directly.
  \item \textbf{Inline FAU \& Special Operands (\hyperref[sec:csf-tsd-fau]{\texttt{FAU\_k}}, \texttt{r8..r15}):} Any 8-bit source with \texttt{mode=2} reads zero-latency uniform \texttt{u0..u31} from \texttt{FAU\_k} (\texttt{slots64 = bits[63:56]}); \texttt{mode=3} reads \texttt{SmallConstant} (Table~1.2) or \texttt{FauSpecial} (\texttt{pc}, \texttt{lane\_id}, \texttt{lane\_width}, \texttt{tls\_ptr}, \texttt{wls\_ptr}, \texttt{blend\_descriptor\_0..7}).
\end{tightitemize}
\par\textbf{Worked Examples:}
\par\textbf{Example 1 (Canonical 32-bit shift/rotate):}
\codebox{\mnem{RSHIFT.i32} \dst{r0}, \src{r1}, \src{r2}\imod{.b0}}
With \texttt{r1 = 0x80000000}, \texttt{r2.b0 = 4}: yields \texttt{r0 = 0x08000000}.
\par\textbf{Example 2 (Byte-lane selector \& modulo-32 masking):}
\codebox{\mnem{RSHIFT.i32} \dst{r0}, \src{r1}, \src{r2}\imod{.b1}}
With \texttt{r1 = 0x80000000}, \texttt{r2.b1 = 32} (shift count masked mod 32 -\textgreater{} \texttt{0}): yields \texttt{r0 = 0x80000000}.

\begin{center}\footnotesize
\begin{tabularx}{\linewidth}{@{}lllX@{}}
\toprule
Variant & Arch & Class & Opcode Fields \\
\midrule
\texttt{(none)} & v9--v14 & \texttt{R3} & \texttt{opcode1=SHIFT\_BITWISE\_I32} \\
\texttt{(none)} & v9--v14 & \texttt{R3} & \texttt{opcode1=SHIFT\_BITWISE\_I32} \\
\texttt{(none)} & v9--v14 & \texttt{R3} & \texttt{opcode1=LOGIC\_SHIFT\_I32} \\
\texttt{(none)} & v9--v14 & \texttt{D3} & \texttt{opcode1=SHIFT\_BITWISE\_I64} \\
\texttt{(none)} & v9--v14 & \texttt{D3} & \texttt{opcode1=LOGIC\_SHIFT\_I64} \\
\texttt{(none)} & v9--v14 & \texttt{R3} & \texttt{opcode1=SHIFT\_BITWISE\_V2I16} \\
\texttt{(none)} & v9--v14 & \texttt{R3} & \texttt{opcode1=LOGIC\_SHIFT\_V2I16} \\
\texttt{(none)} & v9--v14 & \texttt{R3} & \texttt{opcode1=SHIFT\_BITWISE\_V4I8} \\
\texttt{(none)} & v9--v14 & \texttt{R3} & \texttt{opcode1=LOGIC\_SHIFT\_V4I8} \\
\bottomrule
\end{tabularx}
\end{center}
\begin{center}\footnotesize
\begin{tabularx}{0.88\linewidth}{@{}l l X@{}}
\toprule
Field Name & Bits & Type / Fixed Value \\
\midrule
\texttt{src0} & \texttt{0:7} & \texttt{SourceEncoding} \\
\texttt{src1} & \texttt{8:15} & \texttt{SourceEncoding} \\
\texttt{src2} & \texttt{16:23} & \texttt{exact=192} \\
\texttt{bitwise\_operation} & \texttt{24:25} & \texttt{exact=or} \\
\hyperref[enum:shift_bitwise_i32_lanes_m]{\texttt{lanes1}} & \texttt{26:29} & \texttt{shift\_bitwise\_i32\_lanes\_m (def b0)} \\
\hyperref[enum:not_m]{\texttt{not\_result}} & \texttt{30} & \texttt{not\_m (def none)} \\
\textit{Reserved} & \texttt{31} & \texttt{zero} \\
\texttt{shift\_operation} & \texttt{32:34} & \texttt{exact=rshift} \\
\hyperref[enum:not_m]{\texttt{not2}} & \texttt{35} & \texttt{exact=none} \\
\hyperref[enum:widen_int_m]{\texttt{widen0}} & \texttt{36:39} & \texttt{widen\_int\_m (def none)} \\
\texttt{dest\_register\_index} & \texttt{40:45} & \texttt{uint} \\
\hyperref[enum:dest_width_m]{\texttt{dest\_width}} & \texttt{46:47} & \texttt{dest\_width\_m} \\
\texttt{opcode1} & \texttt{48:51} & \texttt{exact=SHIFT\_BITWISE\_I32} \\
\texttt{instruction\_class} & \texttt{52:56} & \texttt{exact=R3} \\
\texttt{fau\_page\_index} & \texttt{57:58} & \texttt{uint} \\
\hyperref[enum:flow_control_m]{\texttt{flow}} & \texttt{59:62} & \texttt{flow\_control\_m} \\
\textit{Reserved} & \texttt{63} & \texttt{zero} \\
\bottomrule
\end{tabularx}
\end{center}
\textbf{Fields Explained:}
\begin{tightitemize}\footnotesize
  \item \texttt{src0}: Primary source operand (\texttt{r0..r63}, \texttt{u0..u63} uniform, \texttt{k*} FAU constant, or \texttt{t0..t2} temporary; \texttt{\textasciicircum{}} marks discard/last-use). \emph{Example:} {\scriptsize\ttfamily \mnem{FADD.f32} \dst{r0}, \src{r1}, \src{r2}} (\textbf{SPIR-V:} \texttt{OpFAdd (Operand 0)}; \textbf{GLSL:} \texttt{float c = a + b;})
  \item \texttt{src1}: Secondary source operand (register, uniform, or FAU constant; in shifts, supplies the shift count byte). \emph{Example:} {\scriptsize\ttfamily \mnem{LSHIFT.i32} \dst{r0}, \src{r1}, \src{r2}\imod{.b0}} (\textbf{SPIR-V:} \texttt{OpShiftLeftLogical (Shift)}; \textbf{GLSL:} \texttt{uint y = a \textless{}\textless{} s;})
  \item \texttt{src2}: Third source operand (FMA addend, \texttt{MUX}/\texttt{CSEL} selector, \texttt{IADDC} carry-in, or bitwise mask in fused shift-logic). \emph{Example:} {\scriptsize\ttfamily \mnem{FMA.f32} \dst{r0}, \src{r1}, \src{r2}, \src{r3}} (\textbf{SPIR-V:} \texttt{GLSL.std.450 Fma (Addend)}; \textbf{GLSL:} \texttt{float d = fma(a, b, c);})
  \item \texttt{bitwise\_operation}: Selects the bitwise ALU op (\texttt{and}, \texttt{or}, \texttt{xor}, \texttt{lshift\_and}, etc.) fused with shifts. \emph{Example:} {\scriptsize\ttfamily \mnem{LSHIFT\_OR.i32} \dst{r0}, \src{r1}, \src{r2}, \src{r3}} (\textbf{SPIR-V:} \texttt{OpShiftLeftLogical + OpBitwiseOr}; \textbf{GLSL:} \texttt{uint y = (a \textless{}\textless{} s) \textbar{} b;})
  \item \hyperref[enum:shift_bitwise_i32_lanes_m]{\texttt{lanes1}}: Stage A byte-lane selector on \texttt{src1} (selects shift-count byte \texttt{b0..b3} on 32/64-bit shifts, \texttt{bXY} on \texttt{v2i16} shifts, or 4-byte swizzle on \texttt{v4*8}). \emph{Example:} {\scriptsize\ttfamily \mnem{LSHIFT.i32} \dst{r0}, \src{r1}, \src{r2}\imod{.b1}} (\textbf{SPIR-V:} \texttt{OpBitFieldUExtract + OpShiftLeftLogical}; \textbf{GLSL:} \texttt{uint y = a \textless{}\textless{} ((s \textgreater{}\textgreater{} 8) \& 0xFFu);})
  \item \hyperref[enum:not_m]{\texttt{not\_result}}: Bitwise-inverts the final output of a bitwise/shift-bitwise instruction (\textasciitilde{}result). \emph{Example:} {\scriptsize\ttfamily \mnem{AND.i32}\imod{.not\_result} \dst{r0}, \src{r1}, \src{r2}} (\textbf{SPIR-V:} \texttt{OpBitwiseAnd + OpNot}; \textbf{GLSL:} \texttt{uint y = \textasciitilde{}(a \& b);})
  \item \texttt{shift\_operation}: Selects shift direction/type (\texttt{lshift}, \texttt{rshift}, \texttt{arshift}) in fused shift-bitwise instructions. \emph{Example:} {\scriptsize\ttfamily \mnem{RSHIFT\_AND.i32} \dst{r0}, \src{r1}, \src{r2}, \src{r3}} (\textbf{SPIR-V:} \texttt{OpShiftRightLogical + OpBitwiseAnd}; \textbf{GLSL:} \texttt{uint y = (a \textgreater{}\textgreater{} s) \& m;})
  \item \hyperref[enum:not_m]{\texttt{not2}}: Stage A bitwise complement (\texttt{\textasciitilde{}src2}) applied to \texttt{src2} before the bitwise/shift-logic operation. \emph{Example:} {\scriptsize\ttfamily \mnem{AND.i32} \dst{r0}, \src{r1}, \src{r2}\imod{.not}} (\textbf{SPIR-V:} \texttt{OpNot + OpBitwiseAnd}; \textbf{GLSL:} \texttt{uint y = a \& \textasciitilde{}b;})
  \item \hyperref[enum:widen_int_m]{\texttt{widen0}}: Stage A widening extractor on \texttt{src0} (promotes FP16 \texttt{h0}/\texttt{h1} to FP32, or sign/zero-extends \texttt{b0..b3}/\texttt{h0..h1}/\texttt{w0..w1} to 32/64-bit before \texttt{abs}/\texttt{neg}). \emph{Example:} {\scriptsize\ttfamily \mnem{IADD.s32} \dst{r0}, \src{r1}\imod{.h0}, \src{r2}} (\textbf{SPIR-V:} \texttt{OpSConvert / OpFConvert}; \textbf{GLSL:} \texttt{int c = int(int16\_t(a)) + b;})
  \item \texttt{dest\_register\_index}: 6-bit destination GPR index (\texttt{r0..r63}) receiving the Stage C writeback result (even-aligned pair \texttt{[rN:rN+1]} for 64-bit ops). \emph{Example:} {\scriptsize\ttfamily \mnem{FMA.f32} \dst{r0}, \src{r1}, \src{r2}, \src{r3}} (\textbf{SPIR-V:} \texttt{SSA Result ID}; \textbf{GLSL:} \texttt{float r0 = fma(r1, r2, r3);})
  \item \hyperref[enum:dest_width_m]{\texttt{dest\_width}}: Stage C destination lane write-mask (\texttt{none} = full 32-bit, \texttt{h0} = low 16-bit, \texttt{h1} = high 16-bit without bit-shifting). \emph{Example:} {\scriptsize\ttfamily \mnem{FADD.v2f16} \dst{r0}\imod{.h1}, \src{r1}, \src{r2}} (\textbf{SPIR-V:} \texttt{OpVectorInsertDynamic / WriteMask}; \textbf{GLSL:} \texttt{dst.y = a.y + b.y;})
  \item \texttt{opcode1}: Primary sub-opcode within the instruction class uniquely identifying the operation family. \emph{Example:} {\scriptsize\ttfamily \mnem{IADD.s32} \dst{r0}, \src{r1}, \src{r2}} (\textbf{SPIR-V:} \texttt{OpIAdd (operation selector)}; \textbf{GLSL:} \texttt{// Primary opcode bits})
  \item \texttt{instruction\_class}: 4-bit major instruction class selector (\texttt{bits 48:51}) partitioning the 64-bit Valhall opcode space (e.g., \texttt{FMA}, \texttt{CVT}, \texttt{SFU}, \texttt{LS}, \texttt{TEX}, \texttt{BRANCH}). \emph{Example:} {\scriptsize\ttfamily \mnem{FADD.f32} \dst{r0}, \src{r1}, \src{r2}} (\textbf{SPIR-V:} \texttt{N/A (encoding class)}; \textbf{GLSL:} \texttt{// Hardware execution pipe / decode class})
  \item \texttt{fau\_page\_index}: Selects the FAU (Fast Access Uniform) page (\texttt{0..3}) when a source operand references hardware small constants, thread IDs, or lane constants. \emph{Example:} {\scriptsize\ttfamily \mnem{FADD.f32} \dst{r0}, \src{r1}, \src{k1f}} (\textbf{SPIR-V:} \texttt{OpConstant / BuiltIn}; \textbf{GLSL:} \texttt{float y = x + 1.0;})
  \item \hyperref[enum:flow_control_m]{\texttt{flow}}: 4-bit embedded flow-control \& scoreboard wait modifier (\texttt{bits 59:62}: \texttt{none}, \texttt{wait0..wait0126}, \texttt{wait}, \texttt{reconverge}, \texttt{discard}, \texttt{end}). \emph{Example:} {\scriptsize\ttfamily \mnem{FADD.f32}\imod{.wait0.reconverge} \dst{r0}, \src{r1}, \src{r2}} (\textbf{SPIR-V:} \texttt{OpControlBarrier / OpReturn}; \textbf{GLSL:} \texttt{// Scoreboard wait + warp reconvergence})
\end{tightitemize}
\clearpage

\subsection*{RSHIFT\_AND}
\addcontentsline{toc}{subsection}{RSHIFT\_AND}
For each SIMD element $i \in \{0, \dots, N-1\}$ ($N=1$ for \texttt{i32}/\texttt{i64}, $N=2$ for \texttt{v2i16}, $N=4$ for \texttt{v4i8} on v9--v10, lane width $W$), computes logical right shift \texttt{s = (uint\_W(val) \textgreater{}\textgreater{} (amt \& (W - 1)))} where \texttt{val = src0[i]} and the 8-bit shift count \texttt{amt} is extracted from \texttt{src1} via \texttt{lanes1}, then combines with \texttt{src2[i]} via bitwise AND \texttt{res = s \& (not2 ? \textasciitilde{}src2[i] : src2[i])} and writes \texttt{not\_result ? \textasciitilde{}res : res} to \texttt{dest[i]}.
\par\textbf{Notes:}
\begin{tightitemize}
  \item Shift amount is masked modulo lane width (\texttt{wrapping\_shl}/\texttt{wrapping\_shr}/\texttt{rotate\_*}, \texttt{kraid/ops.rs:3984-4009}).
  \item On \texttt{v2i16}, \texttt{src1} is \texttt{V2I8} (\texttt{ops.rs:3953}) with byte lanes selected by \texttt{lanes1} (default \texttt{b02}: byte 0 -\textgreater{} lane 0, byte 2 -\textgreater{} lane 1).
  \item \textbf{Source references:} \texttt{kraid/ops.rs:3878-4025}, \texttt{encode\_v9.rs:2940-3093}, \texttt{hw\_tests.rs:2148-2198}.
\end{tightitemize}
\codebox{\mnem{RSHIFT\_AND}\imod{\{.flow\}}\imod{\{.not\_result\}} \dst{dest}\imod{\{.dest\_width\}}, \src{src0}\imod{\{.widen0\}}, \src{src1}\imod{\{.lanes1\}}, \src{src2}\imod{\{.not2\}}\\\mnem{RSHIFT\_AND}\imod{\{.flow\}}\imod{\{.not\_result\}} \dst{dest}, \src{src0}\imod{\{.widen0\}}, \src{src1}\imod{\{.lanes1\}}, \src{src2}\imod{\{.not2\}}\\\mnem{RSHIFT\_AND}\imod{\{.flow\}}\imod{\{.not\_result\}} \dst{dest}\imod{\{.dest\_width\}}, \src{src0}\imod{\{.swz0\}}, \src{src1}\imod{\{.lanes1\}}, \src{src2}\imod{\{.not2\}}\\\mnem{RSHIFT\_AND}\imod{\{.flow\}}\imod{\{.not\_result\}} \dst{dest}\imod{\{.dest\_width\}}, \src{src0}\imod{\{.lanes0\}}, \src{src1}\imod{\{.lanes1\}}, \src{src2}\imod{\{.not2\}}}
\textbf{Operands:}
\begin{tightitemize}\footnotesize
  \item \texttt{\dst{dest}} (\texttt{dest\_register\_index}, \texttt{dest\_width}): Destination GPR receiving the logical right shift combined via bitwise AND (inverted if \texttt{.not\_result}).
  \item \texttt{\src{src0}}: Input integer value to shift/rotate.
  \item \texttt{\src{src1}}: 8-bit shift/rotate count per lane (selected via \texttt{lanes1}, masked modulo lane width \texttt{W}).
  \item \texttt{\src{src2}}: Second operand combined via bitwise AND after shifting (inverted if \texttt{.not2}).
\end{tightitemize}
\par\textbf{Closest SPIR-V Op:} \texttt{OpShiftRightLogical + OpBitwiseAnd}
\par\textbf{Pseudo-GLSL (Operational Semantics):}
\codebox{for (int i = 0; i \textless{} N; ++i) \{\\    uint s = extractByte(src1, lanes1, i) \& (W - 1);\\    auto r = (uint(src0[i]) \textgreater{}\textgreater{} s) \& (flag\_not2 ? \textasciitilde{}src2[i] : src2[i]);\\    dest[i] = flag\_not\_result ? \textasciitilde{}r : r;\\\}}
\textbf{Execution Context:}
\begin{tightitemize}\footnotesize
  \item \textbf{Compute stage preloads (\hyperref[sec:csf-spd]{\texttt{SPD\_k}}, \texttt{r16..r23}):} Hardware pre-fills \texttt{r55..r57} with \texttt{WorkGroupID.xyz + JOB\_OFFSET} (from \hyperref[sec:csf-regs]{\texttt{COMPUTESR::JOB\_OFFSET\_X/Y/Z}} \texttt{r34..r36}) and \texttt{r60..r62} with \texttt{LocalInvocationID.xyz} or \texttt{GlobalInvocationID.xyz} (from \texttt{WG\_SIZE} \texttt{r33} \& \texttt{JOB\_SIZE} \texttt{r37..r39}, selected by \texttt{ShaderProgram::preload}) so compute shaders can derive thread/element indices without memory loads.
  \item \textbf{Vertex stage preloads (\hyperref[sec:csf-spd]{\texttt{SPD\_0}/\texttt{SPD\_1}}, \texttt{r16..r19}):} The hardware Index Fetcher reads \hyperref[sec:csf-regs]{\texttt{IDVSSR::INDEX\_BUFFER}} (\texttt{r54:r55}) and pre-fills \texttt{r59, r60} with \texttt{VertexIndex} (\texttt{+ VERTEX\_OFFSET + GLOBAL\_ATTR\_OFF}) and \texttt{r61} with \texttt{InstanceIndex} (\texttt{+ INSTANCE\_OFFSET}) so Position and Varying shaders can pull vertex attributes directly.
  \item \textbf{Inline FAU \& Special Operands (\hyperref[sec:csf-tsd-fau]{\texttt{FAU\_k}}, \texttt{r8..r15}):} Any 8-bit source with \texttt{mode=2} reads zero-latency uniform \texttt{u0..u31} from \texttt{FAU\_k} (\texttt{slots64 = bits[63:56]}); \texttt{mode=3} reads \texttt{SmallConstant} (Table~1.2) or \texttt{FauSpecial} (\texttt{pc}, \texttt{lane\_id}, \texttt{lane\_width}, \texttt{tls\_ptr}, \texttt{wls\_ptr}, \texttt{blend\_descriptor\_0..7}).
\end{tightitemize}
\par\textbf{Worked Examples:}
\par\textbf{Example 1 (Standard fused shift + bitwise op):}
\codebox{\mnem{RSHIFT\_AND.i32} \dst{r0}, \src{r1}, \src{r2}\imod{.b0}, \src{r3}}
With \texttt{r1 = 0x80000000}, \texttt{r2.b0 = 4} (masked mod 32): shift yields \texttt{0x08000000}, then \texttt{\& 0x0FFFFFFF} yields \texttt{r0 = 0x8000000}.
\par\textbf{Example 2 (Modulo-32 shift count \& .not\_result):}
\codebox{\mnem{RSHIFT\_AND.i32}\imod{.not\_result} \dst{r0}, \src{r1}, \src{r2}\imod{.b0}, \src{r3}}
With \texttt{r1 = 0x80000000}, \texttt{r2.b0 = 32} (effective shift \texttt{0}): shift yields \texttt{0x80000000}, \texttt{\& 0x0FFFFFFF} yields \texttt{0x00000000}, and \texttt{.not\_result} inverts -\textgreater{} \texttt{r0 = 0xFFFFFFFF}.
\par\textbf{Mined SPIR-V $\to$ Valhall ISA Example (Fused dynamic bit-extraction: (bits \textgreater{}\textgreater{} shift) \& mask --- Vertex shader, Mali-G710 (v10)):}
\codebox{\textbf{; SPIR-V (input)}\\\%s = OpShiftRightLogical \%uint \%word \%bit\_idx\\\%m = OpBitwiseAnd \%uint \%s \%ubo\_mask\\[3pt]\textbf{; Valhall ISA (malioc output)}\\RSHIFT\_AND.i32 r6, r11\textasciicircum{}, r6\textasciicircum{}.b0, u13.w0\\CSEL.u32.eq r6, r6\textasciicircum{}, k0, u13.w1, r7\textasciicircum{}}
\texttt{malioc} fuses \texttt{OpShiftRightLogical} + \texttt{OpBitwiseAnd} into a single \texttt{RSHIFT\_AND.i32} instruction: \texttt{r11} is logically right-shifted by byte 0 of \texttt{r6} (\texttt{r6\textasciicircum{}.b0}) and masked with \texttt{u13.w0} in one cycle.
\begin{center}\footnotesize
\begin{tabularx}{\linewidth}{@{}lllX@{}}
\toprule
Variant & Arch & Class & Opcode Fields \\
\midrule
\texttt{(none)} & v9--v14 & \texttt{R3} & \texttt{opcode1=SHIFT\_BITWISE\_I32} \\
\texttt{(none)} & v9--v14 & \texttt{R3} & \texttt{opcode1=SHIFT\_BITWISE\_I32} \\
\texttt{(none)} & v9--v14 & \texttt{R3} & \texttt{opcode1=LOGIC\_SHIFT\_I32} \\
\texttt{(none)} & v9--v14 & \texttt{D3} & \texttt{opcode1=SHIFT\_BITWISE\_I64} \\
\texttt{(none)} & v9--v14 & \texttt{D3} & \texttt{opcode1=LOGIC\_SHIFT\_I64} \\
\texttt{(none)} & v9--v14 & \texttt{R3} & \texttt{opcode1=SHIFT\_BITWISE\_V2I16} \\
\texttt{(none)} & v9--v14 & \texttt{R3} & \texttt{opcode1=LOGIC\_SHIFT\_V2I16} \\
\texttt{(none)} & v9--v14 & \texttt{R3} & \texttt{opcode1=SHIFT\_BITWISE\_V4I8} \\
\texttt{(none)} & v9--v14 & \texttt{R3} & \texttt{opcode1=LOGIC\_SHIFT\_V4I8} \\
\bottomrule
\end{tabularx}
\end{center}
\begin{center}\footnotesize
\begin{tabularx}{0.88\linewidth}{@{}l l X@{}}
\toprule
Field Name & Bits & Type / Fixed Value \\
\midrule
\texttt{src0} & \texttt{0:7} & \texttt{SourceEncoding} \\
\texttt{src1} & \texttt{8:15} & \texttt{SourceEncoding} \\
\texttt{src2} & \texttt{16:23} & \texttt{SourceEncoding} \\
\texttt{bitwise\_operation} & \texttt{24:25} & \texttt{exact=and} \\
\hyperref[enum:shift_bitwise_i32_lanes_m]{\texttt{lanes1}} & \texttt{26:29} & \texttt{shift\_bitwise\_i32\_lanes\_m (def b0)} \\
\hyperref[enum:not_m]{\texttt{not\_result}} & \texttt{30} & \texttt{not\_m (def none)} \\
\textit{Reserved} & \texttt{31} & \texttt{zero} \\
\texttt{shift\_operation} & \texttt{32:34} & \texttt{exact=rshift} \\
\hyperref[enum:not_m]{\texttt{not2}} & \texttt{35} & \texttt{not\_m (def none)} \\
\hyperref[enum:widen_int_m]{\texttt{widen0}} & \texttt{36:39} & \texttt{widen\_int\_m (def none)} \\
\texttt{dest\_register\_index} & \texttt{40:45} & \texttt{uint} \\
\hyperref[enum:dest_width_m]{\texttt{dest\_width}} & \texttt{46:47} & \texttt{dest\_width\_m} \\
\texttt{opcode1} & \texttt{48:51} & \texttt{exact=SHIFT\_BITWISE\_I32} \\
\texttt{instruction\_class} & \texttt{52:56} & \texttt{exact=R3} \\
\texttt{fau\_page\_index} & \texttt{57:58} & \texttt{uint} \\
\hyperref[enum:flow_control_m]{\texttt{flow}} & \texttt{59:62} & \texttt{flow\_control\_m} \\
\textit{Reserved} & \texttt{63} & \texttt{zero} \\
\bottomrule
\end{tabularx}
\end{center}
\textbf{Fields Explained:}
\begin{tightitemize}\footnotesize
  \item \texttt{src0}: Primary source operand (\texttt{r0..r63}, \texttt{u0..u63} uniform, \texttt{k*} FAU constant, or \texttt{t0..t2} temporary; \texttt{\textasciicircum{}} marks discard/last-use). \emph{Example:} {\scriptsize\ttfamily \mnem{FADD.f32} \dst{r0}, \src{r1}, \src{r2}} (\textbf{SPIR-V:} \texttt{OpFAdd (Operand 0)}; \textbf{GLSL:} \texttt{float c = a + b;})
  \item \texttt{src1}: Secondary source operand (register, uniform, or FAU constant; in shifts, supplies the shift count byte). \emph{Example:} {\scriptsize\ttfamily \mnem{LSHIFT.i32} \dst{r0}, \src{r1}, \src{r2}\imod{.b0}} (\textbf{SPIR-V:} \texttt{OpShiftLeftLogical (Shift)}; \textbf{GLSL:} \texttt{uint y = a \textless{}\textless{} s;})
  \item \texttt{src2}: Third source operand (FMA addend, \texttt{MUX}/\texttt{CSEL} selector, \texttt{IADDC} carry-in, or bitwise mask in fused shift-logic). \emph{Example:} {\scriptsize\ttfamily \mnem{FMA.f32} \dst{r0}, \src{r1}, \src{r2}, \src{r3}} (\textbf{SPIR-V:} \texttt{GLSL.std.450 Fma (Addend)}; \textbf{GLSL:} \texttt{float d = fma(a, b, c);})
  \item \texttt{bitwise\_operation}: Selects the bitwise ALU op (\texttt{and}, \texttt{or}, \texttt{xor}, \texttt{lshift\_and}, etc.) fused with shifts. \emph{Example:} {\scriptsize\ttfamily \mnem{LSHIFT\_OR.i32} \dst{r0}, \src{r1}, \src{r2}, \src{r3}} (\textbf{SPIR-V:} \texttt{OpShiftLeftLogical + OpBitwiseOr}; \textbf{GLSL:} \texttt{uint y = (a \textless{}\textless{} s) \textbar{} b;})
  \item \hyperref[enum:shift_bitwise_i32_lanes_m]{\texttt{lanes1}}: Stage A byte-lane selector on \texttt{src1} (selects shift-count byte \texttt{b0..b3} on 32/64-bit shifts, \texttt{bXY} on \texttt{v2i16} shifts, or 4-byte swizzle on \texttt{v4*8}). \emph{Example:} {\scriptsize\ttfamily \mnem{LSHIFT.i32} \dst{r0}, \src{r1}, \src{r2}\imod{.b1}} (\textbf{SPIR-V:} \texttt{OpBitFieldUExtract + OpShiftLeftLogical}; \textbf{GLSL:} \texttt{uint y = a \textless{}\textless{} ((s \textgreater{}\textgreater{} 8) \& 0xFFu);})
  \item \hyperref[enum:not_m]{\texttt{not\_result}}: Bitwise-inverts the final output of a bitwise/shift-bitwise instruction (\textasciitilde{}result). \emph{Example:} {\scriptsize\ttfamily \mnem{AND.i32}\imod{.not\_result} \dst{r0}, \src{r1}, \src{r2}} (\textbf{SPIR-V:} \texttt{OpBitwiseAnd + OpNot}; \textbf{GLSL:} \texttt{uint y = \textasciitilde{}(a \& b);})
  \item \texttt{shift\_operation}: Selects shift direction/type (\texttt{lshift}, \texttt{rshift}, \texttt{arshift}) in fused shift-bitwise instructions. \emph{Example:} {\scriptsize\ttfamily \mnem{RSHIFT\_AND.i32} \dst{r0}, \src{r1}, \src{r2}, \src{r3}} (\textbf{SPIR-V:} \texttt{OpShiftRightLogical + OpBitwiseAnd}; \textbf{GLSL:} \texttt{uint y = (a \textgreater{}\textgreater{} s) \& m;})
  \item \hyperref[enum:not_m]{\texttt{not2}}: Stage A bitwise complement (\texttt{\textasciitilde{}src2}) applied to \texttt{src2} before the bitwise/shift-logic operation. \emph{Example:} {\scriptsize\ttfamily \mnem{AND.i32} \dst{r0}, \src{r1}, \src{r2}\imod{.not}} (\textbf{SPIR-V:} \texttt{OpNot + OpBitwiseAnd}; \textbf{GLSL:} \texttt{uint y = a \& \textasciitilde{}b;})
  \item \hyperref[enum:widen_int_m]{\texttt{widen0}}: Stage A widening extractor on \texttt{src0} (promotes FP16 \texttt{h0}/\texttt{h1} to FP32, or sign/zero-extends \texttt{b0..b3}/\texttt{h0..h1}/\texttt{w0..w1} to 32/64-bit before \texttt{abs}/\texttt{neg}). \emph{Example:} {\scriptsize\ttfamily \mnem{IADD.s32} \dst{r0}, \src{r1}\imod{.h0}, \src{r2}} (\textbf{SPIR-V:} \texttt{OpSConvert / OpFConvert}; \textbf{GLSL:} \texttt{int c = int(int16\_t(a)) + b;})
  \item \texttt{dest\_register\_index}: 6-bit destination GPR index (\texttt{r0..r63}) receiving the Stage C writeback result (even-aligned pair \texttt{[rN:rN+1]} for 64-bit ops). \emph{Example:} {\scriptsize\ttfamily \mnem{FMA.f32} \dst{r0}, \src{r1}, \src{r2}, \src{r3}} (\textbf{SPIR-V:} \texttt{SSA Result ID}; \textbf{GLSL:} \texttt{float r0 = fma(r1, r2, r3);})
  \item \hyperref[enum:dest_width_m]{\texttt{dest\_width}}: Stage C destination lane write-mask (\texttt{none} = full 32-bit, \texttt{h0} = low 16-bit, \texttt{h1} = high 16-bit without bit-shifting). \emph{Example:} {\scriptsize\ttfamily \mnem{FADD.v2f16} \dst{r0}\imod{.h1}, \src{r1}, \src{r2}} (\textbf{SPIR-V:} \texttt{OpVectorInsertDynamic / WriteMask}; \textbf{GLSL:} \texttt{dst.y = a.y + b.y;})
  \item \texttt{opcode1}: Primary sub-opcode within the instruction class uniquely identifying the operation family. \emph{Example:} {\scriptsize\ttfamily \mnem{IADD.s32} \dst{r0}, \src{r1}, \src{r2}} (\textbf{SPIR-V:} \texttt{OpIAdd (operation selector)}; \textbf{GLSL:} \texttt{// Primary opcode bits})
  \item \texttt{instruction\_class}: 4-bit major instruction class selector (\texttt{bits 48:51}) partitioning the 64-bit Valhall opcode space (e.g., \texttt{FMA}, \texttt{CVT}, \texttt{SFU}, \texttt{LS}, \texttt{TEX}, \texttt{BRANCH}). \emph{Example:} {\scriptsize\ttfamily \mnem{FADD.f32} \dst{r0}, \src{r1}, \src{r2}} (\textbf{SPIR-V:} \texttt{N/A (encoding class)}; \textbf{GLSL:} \texttt{// Hardware execution pipe / decode class})
  \item \texttt{fau\_page\_index}: Selects the FAU (Fast Access Uniform) page (\texttt{0..3}) when a source operand references hardware small constants, thread IDs, or lane constants. \emph{Example:} {\scriptsize\ttfamily \mnem{FADD.f32} \dst{r0}, \src{r1}, \src{k1f}} (\textbf{SPIR-V:} \texttt{OpConstant / BuiltIn}; \textbf{GLSL:} \texttt{float y = x + 1.0;})
  \item \hyperref[enum:flow_control_m]{\texttt{flow}}: 4-bit embedded flow-control \& scoreboard wait modifier (\texttt{bits 59:62}: \texttt{none}, \texttt{wait0..wait0126}, \texttt{wait}, \texttt{reconverge}, \texttt{discard}, \texttt{end}). \emph{Example:} {\scriptsize\ttfamily \mnem{FADD.f32}\imod{.wait0.reconverge} \dst{r0}, \src{r1}, \src{r2}} (\textbf{SPIR-V:} \texttt{OpControlBarrier / OpReturn}; \textbf{GLSL:} \texttt{// Scoreboard wait + warp reconvergence})
\end{tightitemize}
\clearpage

\subsection*{RSHIFT\_OR}
\addcontentsline{toc}{subsection}{RSHIFT\_OR}
For each SIMD element $i \in \{0, \dots, N-1\}$ ($N=1$ for \texttt{i32}/\texttt{i64}, $N=2$ for \texttt{v2i16}, $N=4$ for \texttt{v4i8} on v9--v10, lane width $W$), computes logical right shift \texttt{s = (uint\_W(val) \textgreater{}\textgreater{} (amt \& (W - 1)))} where \texttt{val = src0[i]} and the 8-bit shift count \texttt{amt} is extracted from \texttt{src1} via \texttt{lanes1}, then combines with \texttt{src2[i]} via bitwise OR \texttt{res = s \textbar{} (not2 ? \textasciitilde{}src2[i] : src2[i])} and writes \texttt{not\_result ? \textasciitilde{}res : res} to \texttt{dest[i]}.
\par\textbf{Notes:}
\begin{tightitemize}
  \item Shift amount is masked modulo lane width (\texttt{wrapping\_shl}/\texttt{wrapping\_shr}/\texttt{rotate\_*}, \texttt{kraid/ops.rs:3984-4009}).
  \item On \texttt{v2i16}, \texttt{src1} is \texttt{V2I8} (\texttt{ops.rs:3953}) with byte lanes selected by \texttt{lanes1} (default \texttt{b02}: byte 0 -\textgreater{} lane 0, byte 2 -\textgreater{} lane 1).
  \item \textbf{Source references:} \texttt{kraid/ops.rs:3878-4025}, \texttt{encode\_v9.rs:2940-3093}, \texttt{hw\_tests.rs:2148-2198}.
\end{tightitemize}
\codebox{\mnem{RSHIFT\_OR}\imod{\{.flow\}}\imod{\{.not\_result\}} \dst{dest}\imod{\{.dest\_width\}}, \src{src0}\imod{\{.widen0\}}, \src{src1}\imod{\{.lanes1\}}, \src{src2}\imod{\{.not2\}}\\\mnem{RSHIFT\_OR}\imod{\{.flow\}}\imod{\{.not\_result\}} \dst{dest}, \src{src0}\imod{\{.widen0\}}, \src{src1}\imod{\{.lanes1\}}, \src{src2}\imod{\{.not2\}}\\\mnem{RSHIFT\_OR}\imod{\{.flow\}}\imod{\{.not\_result\}} \dst{dest}\imod{\{.dest\_width\}}, \src{src0}\imod{\{.swz0\}}, \src{src1}\imod{\{.lanes1\}}, \src{src2}\imod{\{.not2\}}\\\mnem{RSHIFT\_OR}\imod{\{.flow\}}\imod{\{.not\_result\}} \dst{dest}\imod{\{.dest\_width\}}, \src{src0}\imod{\{.lanes0\}}, \src{src1}\imod{\{.lanes1\}}, \src{src2}\imod{\{.not2\}}}
\textbf{Operands:}
\begin{tightitemize}\footnotesize
  \item \texttt{\dst{dest}} (\texttt{dest\_register\_index}, \texttt{dest\_width}): Destination GPR receiving the logical right shift combined via bitwise OR (inverted if \texttt{.not\_result}).
  \item \texttt{\src{src0}}: Input integer value to shift/rotate.
  \item \texttt{\src{src1}}: 8-bit shift/rotate count per lane (selected via \texttt{lanes1}, masked modulo lane width \texttt{W}).
  \item \texttt{\src{src2}}: Second operand combined via bitwise OR after shifting (inverted if \texttt{.not2}).
\end{tightitemize}
\par\textbf{Closest SPIR-V Op:} \texttt{OpShiftRightLogical + OpBitwiseOr}
\par\textbf{Pseudo-GLSL (Operational Semantics):}
\codebox{for (int i = 0; i \textless{} N; ++i) \{\\    uint s = extractByte(src1, lanes1, i) \& (W - 1);\\    auto r = (uint(src0[i]) \textgreater{}\textgreater{} s) \textbar{} (flag\_not2 ? \textasciitilde{}src2[i] : src2[i]);\\    dest[i] = flag\_not\_result ? \textasciitilde{}r : r;\\\}}
\textbf{Execution Context:}
\begin{tightitemize}\footnotesize
  \item \textbf{Compute stage preloads (\hyperref[sec:csf-spd]{\texttt{SPD\_k}}, \texttt{r16..r23}):} Hardware pre-fills \texttt{r55..r57} with \texttt{WorkGroupID.xyz + JOB\_OFFSET} (from \hyperref[sec:csf-regs]{\texttt{COMPUTESR::JOB\_OFFSET\_X/Y/Z}} \texttt{r34..r36}) and \texttt{r60..r62} with \texttt{LocalInvocationID.xyz} or \texttt{GlobalInvocationID.xyz} (from \texttt{WG\_SIZE} \texttt{r33} \& \texttt{JOB\_SIZE} \texttt{r37..r39}, selected by \texttt{ShaderProgram::preload}) so compute shaders can derive thread/element indices without memory loads.
  \item \textbf{Vertex stage preloads (\hyperref[sec:csf-spd]{\texttt{SPD\_0}/\texttt{SPD\_1}}, \texttt{r16..r19}):} The hardware Index Fetcher reads \hyperref[sec:csf-regs]{\texttt{IDVSSR::INDEX\_BUFFER}} (\texttt{r54:r55}) and pre-fills \texttt{r59, r60} with \texttt{VertexIndex} (\texttt{+ VERTEX\_OFFSET + GLOBAL\_ATTR\_OFF}) and \texttt{r61} with \texttt{InstanceIndex} (\texttt{+ INSTANCE\_OFFSET}) so Position and Varying shaders can pull vertex attributes directly.
  \item \textbf{Inline FAU \& Special Operands (\hyperref[sec:csf-tsd-fau]{\texttt{FAU\_k}}, \texttt{r8..r15}):} Any 8-bit source with \texttt{mode=2} reads zero-latency uniform \texttt{u0..u31} from \texttt{FAU\_k} (\texttt{slots64 = bits[63:56]}); \texttt{mode=3} reads \texttt{SmallConstant} (Table~1.2) or \texttt{FauSpecial} (\texttt{pc}, \texttt{lane\_id}, \texttt{lane\_width}, \texttt{tls\_ptr}, \texttt{wls\_ptr}, \texttt{blend\_descriptor\_0..7}).
\end{tightitemize}
\par\textbf{Worked Examples:}
\par\textbf{Example 1 (Standard fused shift + bitwise op):}
\codebox{\mnem{RSHIFT\_OR.i32} \dst{r0}, \src{r1}, \src{r2}\imod{.b0}, \src{r3}}
With \texttt{r1 = 0x80000000}, \texttt{r2.b0 = 4} (masked mod 32): shift yields \texttt{0x08000000}, then \texttt{\textbar{} 0x0000000F} yields \texttt{r0 = 0x800000F}.
\par\textbf{Example 2 (Modulo-32 shift count \& .not\_result):}
\codebox{\mnem{RSHIFT\_OR.i32}\imod{.not\_result} \dst{r0}, \src{r1}, \src{r2}\imod{.b0}, \src{r3}}
With \texttt{r1 = 0x80000000}, \texttt{r2.b0 = 32} (effective shift \texttt{0}): shift yields \texttt{0x80000000}, \texttt{\textbar{} 0x0000000F} yields \texttt{0x8000000F}, and \texttt{.not\_result} inverts -\textgreater{} \texttt{r0 = 0x7FFFFFF0}.
\par\textbf{Mined SPIR-V $\to$ Valhall ISA Example (Logical right shift for index division by 32 --- Vertex shader, Mali-G710 (v10)):}
\codebox{\textbf{; SPIR-V (input)}\\\%word\_idx = OpShiftRightLogical \%uint \%bit\_idx \%uint\_5\\\%addr\_idx = OpIAdd \%uint \%word\_idx \%base\\[3pt]\textbf{; Valhall ISA (malioc output)}\\RSHIFT.i32 r11, r6, k07060504.b1\\IADD.u32 r11, r11\textasciicircum{}, u3.w0}
\texttt{RSHIFT.i32} (encoded as \texttt{RSHIFT\_OR} with \texttt{src2 = k0}) shifts \texttt{r6} right by 5 (\texttt{k07060504.b1 = 0x05}) to compute a 32-bit word index.
\begin{center}\footnotesize
\begin{tabularx}{\linewidth}{@{}lllX@{}}
\toprule
Variant & Arch & Class & Opcode Fields \\
\midrule
\texttt{(none)} & v9--v14 & \texttt{R3} & \texttt{opcode1=SHIFT\_BITWISE\_I32} \\
\texttt{(none)} & v9--v14 & \texttt{R3} & \texttt{opcode1=SHIFT\_BITWISE\_I32} \\
\texttt{(none)} & v9--v14 & \texttt{R3} & \texttt{opcode1=LOGIC\_SHIFT\_I32} \\
\texttt{(none)} & v9--v14 & \texttt{D3} & \texttt{opcode1=SHIFT\_BITWISE\_I64} \\
\texttt{(none)} & v9--v14 & \texttt{D3} & \texttt{opcode1=LOGIC\_SHIFT\_I64} \\
\texttt{(none)} & v9--v14 & \texttt{R3} & \texttt{opcode1=SHIFT\_BITWISE\_V2I16} \\
\texttt{(none)} & v9--v14 & \texttt{R3} & \texttt{opcode1=LOGIC\_SHIFT\_V2I16} \\
\texttt{(none)} & v9--v14 & \texttt{R3} & \texttt{opcode1=SHIFT\_BITWISE\_V4I8} \\
\texttt{(none)} & v9--v14 & \texttt{R3} & \texttt{opcode1=LOGIC\_SHIFT\_V4I8} \\
\bottomrule
\end{tabularx}
\end{center}
\begin{center}\footnotesize
\begin{tabularx}{0.88\linewidth}{@{}l l X@{}}
\toprule
Field Name & Bits & Type / Fixed Value \\
\midrule
\texttt{src0} & \texttt{0:7} & \texttt{SourceEncoding} \\
\texttt{src1} & \texttt{8:15} & \texttt{SourceEncoding} \\
\texttt{src2} & \texttt{16:23} & \texttt{SourceEncoding} \\
\texttt{bitwise\_operation} & \texttt{24:25} & \texttt{exact=or} \\
\hyperref[enum:shift_bitwise_i32_lanes_m]{\texttt{lanes1}} & \texttt{26:29} & \texttt{shift\_bitwise\_i32\_lanes\_m (def b0)} \\
\hyperref[enum:not_m]{\texttt{not\_result}} & \texttt{30} & \texttt{not\_m (def none)} \\
\textit{Reserved} & \texttt{31} & \texttt{zero} \\
\texttt{shift\_operation} & \texttt{32:34} & \texttt{exact=rshift} \\
\hyperref[enum:not_m]{\texttt{not2}} & \texttt{35} & \texttt{not\_m (def none)} \\
\hyperref[enum:widen_int_m]{\texttt{widen0}} & \texttt{36:39} & \texttt{widen\_int\_m (def none)} \\
\texttt{dest\_register\_index} & \texttt{40:45} & \texttt{uint} \\
\hyperref[enum:dest_width_m]{\texttt{dest\_width}} & \texttt{46:47} & \texttt{dest\_width\_m} \\
\texttt{opcode1} & \texttt{48:51} & \texttt{exact=SHIFT\_BITWISE\_I32} \\
\texttt{instruction\_class} & \texttt{52:56} & \texttt{exact=R3} \\
\texttt{fau\_page\_index} & \texttt{57:58} & \texttt{uint} \\
\hyperref[enum:flow_control_m]{\texttt{flow}} & \texttt{59:62} & \texttt{flow\_control\_m} \\
\textit{Reserved} & \texttt{63} & \texttt{zero} \\
\bottomrule
\end{tabularx}
\end{center}
\textbf{Fields Explained:}
\begin{tightitemize}\footnotesize
  \item \texttt{src0}: Primary source operand (\texttt{r0..r63}, \texttt{u0..u63} uniform, \texttt{k*} FAU constant, or \texttt{t0..t2} temporary; \texttt{\textasciicircum{}} marks discard/last-use). \emph{Example:} {\scriptsize\ttfamily \mnem{FADD.f32} \dst{r0}, \src{r1}, \src{r2}} (\textbf{SPIR-V:} \texttt{OpFAdd (Operand 0)}; \textbf{GLSL:} \texttt{float c = a + b;})
  \item \texttt{src1}: Secondary source operand (register, uniform, or FAU constant; in shifts, supplies the shift count byte). \emph{Example:} {\scriptsize\ttfamily \mnem{LSHIFT.i32} \dst{r0}, \src{r1}, \src{r2}\imod{.b0}} (\textbf{SPIR-V:} \texttt{OpShiftLeftLogical (Shift)}; \textbf{GLSL:} \texttt{uint y = a \textless{}\textless{} s;})
  \item \texttt{src2}: Third source operand (FMA addend, \texttt{MUX}/\texttt{CSEL} selector, \texttt{IADDC} carry-in, or bitwise mask in fused shift-logic). \emph{Example:} {\scriptsize\ttfamily \mnem{FMA.f32} \dst{r0}, \src{r1}, \src{r2}, \src{r3}} (\textbf{SPIR-V:} \texttt{GLSL.std.450 Fma (Addend)}; \textbf{GLSL:} \texttt{float d = fma(a, b, c);})
  \item \texttt{bitwise\_operation}: Selects the bitwise ALU op (\texttt{and}, \texttt{or}, \texttt{xor}, \texttt{lshift\_and}, etc.) fused with shifts. \emph{Example:} {\scriptsize\ttfamily \mnem{LSHIFT\_OR.i32} \dst{r0}, \src{r1}, \src{r2}, \src{r3}} (\textbf{SPIR-V:} \texttt{OpShiftLeftLogical + OpBitwiseOr}; \textbf{GLSL:} \texttt{uint y = (a \textless{}\textless{} s) \textbar{} b;})
  \item \hyperref[enum:shift_bitwise_i32_lanes_m]{\texttt{lanes1}}: Stage A byte-lane selector on \texttt{src1} (selects shift-count byte \texttt{b0..b3} on 32/64-bit shifts, \texttt{bXY} on \texttt{v2i16} shifts, or 4-byte swizzle on \texttt{v4*8}). \emph{Example:} {\scriptsize\ttfamily \mnem{LSHIFT.i32} \dst{r0}, \src{r1}, \src{r2}\imod{.b1}} (\textbf{SPIR-V:} \texttt{OpBitFieldUExtract + OpShiftLeftLogical}; \textbf{GLSL:} \texttt{uint y = a \textless{}\textless{} ((s \textgreater{}\textgreater{} 8) \& 0xFFu);})
  \item \hyperref[enum:not_m]{\texttt{not\_result}}: Bitwise-inverts the final output of a bitwise/shift-bitwise instruction (\textasciitilde{}result). \emph{Example:} {\scriptsize\ttfamily \mnem{AND.i32}\imod{.not\_result} \dst{r0}, \src{r1}, \src{r2}} (\textbf{SPIR-V:} \texttt{OpBitwiseAnd + OpNot}; \textbf{GLSL:} \texttt{uint y = \textasciitilde{}(a \& b);})
  \item \texttt{shift\_operation}: Selects shift direction/type (\texttt{lshift}, \texttt{rshift}, \texttt{arshift}) in fused shift-bitwise instructions. \emph{Example:} {\scriptsize\ttfamily \mnem{RSHIFT\_AND.i32} \dst{r0}, \src{r1}, \src{r2}, \src{r3}} (\textbf{SPIR-V:} \texttt{OpShiftRightLogical + OpBitwiseAnd}; \textbf{GLSL:} \texttt{uint y = (a \textgreater{}\textgreater{} s) \& m;})
  \item \hyperref[enum:not_m]{\texttt{not2}}: Stage A bitwise complement (\texttt{\textasciitilde{}src2}) applied to \texttt{src2} before the bitwise/shift-logic operation. \emph{Example:} {\scriptsize\ttfamily \mnem{AND.i32} \dst{r0}, \src{r1}, \src{r2}\imod{.not}} (\textbf{SPIR-V:} \texttt{OpNot + OpBitwiseAnd}; \textbf{GLSL:} \texttt{uint y = a \& \textasciitilde{}b;})
  \item \hyperref[enum:widen_int_m]{\texttt{widen0}}: Stage A widening extractor on \texttt{src0} (promotes FP16 \texttt{h0}/\texttt{h1} to FP32, or sign/zero-extends \texttt{b0..b3}/\texttt{h0..h1}/\texttt{w0..w1} to 32/64-bit before \texttt{abs}/\texttt{neg}). \emph{Example:} {\scriptsize\ttfamily \mnem{IADD.s32} \dst{r0}, \src{r1}\imod{.h0}, \src{r2}} (\textbf{SPIR-V:} \texttt{OpSConvert / OpFConvert}; \textbf{GLSL:} \texttt{int c = int(int16\_t(a)) + b;})
  \item \texttt{dest\_register\_index}: 6-bit destination GPR index (\texttt{r0..r63}) receiving the Stage C writeback result (even-aligned pair \texttt{[rN:rN+1]} for 64-bit ops). \emph{Example:} {\scriptsize\ttfamily \mnem{FMA.f32} \dst{r0}, \src{r1}, \src{r2}, \src{r3}} (\textbf{SPIR-V:} \texttt{SSA Result ID}; \textbf{GLSL:} \texttt{float r0 = fma(r1, r2, r3);})
  \item \hyperref[enum:dest_width_m]{\texttt{dest\_width}}: Stage C destination lane write-mask (\texttt{none} = full 32-bit, \texttt{h0} = low 16-bit, \texttt{h1} = high 16-bit without bit-shifting). \emph{Example:} {\scriptsize\ttfamily \mnem{FADD.v2f16} \dst{r0}\imod{.h1}, \src{r1}, \src{r2}} (\textbf{SPIR-V:} \texttt{OpVectorInsertDynamic / WriteMask}; \textbf{GLSL:} \texttt{dst.y = a.y + b.y;})
  \item \texttt{opcode1}: Primary sub-opcode within the instruction class uniquely identifying the operation family. \emph{Example:} {\scriptsize\ttfamily \mnem{IADD.s32} \dst{r0}, \src{r1}, \src{r2}} (\textbf{SPIR-V:} \texttt{OpIAdd (operation selector)}; \textbf{GLSL:} \texttt{// Primary opcode bits})
  \item \texttt{instruction\_class}: 4-bit major instruction class selector (\texttt{bits 48:51}) partitioning the 64-bit Valhall opcode space (e.g., \texttt{FMA}, \texttt{CVT}, \texttt{SFU}, \texttt{LS}, \texttt{TEX}, \texttt{BRANCH}). \emph{Example:} {\scriptsize\ttfamily \mnem{FADD.f32} \dst{r0}, \src{r1}, \src{r2}} (\textbf{SPIR-V:} \texttt{N/A (encoding class)}; \textbf{GLSL:} \texttt{// Hardware execution pipe / decode class})
  \item \texttt{fau\_page\_index}: Selects the FAU (Fast Access Uniform) page (\texttt{0..3}) when a source operand references hardware small constants, thread IDs, or lane constants. \emph{Example:} {\scriptsize\ttfamily \mnem{FADD.f32} \dst{r0}, \src{r1}, \src{k1f}} (\textbf{SPIR-V:} \texttt{OpConstant / BuiltIn}; \textbf{GLSL:} \texttt{float y = x + 1.0;})
  \item \hyperref[enum:flow_control_m]{\texttt{flow}}: 4-bit embedded flow-control \& scoreboard wait modifier (\texttt{bits 59:62}: \texttt{none}, \texttt{wait0..wait0126}, \texttt{wait}, \texttt{reconverge}, \texttt{discard}, \texttt{end}). \emph{Example:} {\scriptsize\ttfamily \mnem{FADD.f32}\imod{.wait0.reconverge} \dst{r0}, \src{r1}, \src{r2}} (\textbf{SPIR-V:} \texttt{OpControlBarrier / OpReturn}; \textbf{GLSL:} \texttt{// Scoreboard wait + warp reconvergence})
\end{tightitemize}
\clearpage

\subsection*{RSHIFT\_XOR}
\addcontentsline{toc}{subsection}{RSHIFT\_XOR}
For each SIMD element $i \in \{0, \dots, N-1\}$ ($N=1$ for \texttt{i32}/\texttt{i64}, $N=2$ for \texttt{v2i16}, $N=4$ for \texttt{v4i8} on v9--v10, lane width $W$), computes logical right shift \texttt{s = (uint\_W(val) \textgreater{}\textgreater{} (amt \& (W - 1)))} where \texttt{val = src0[i]} and the 8-bit shift count \texttt{amt} is extracted from \texttt{src1} via \texttt{lanes1}, then combines with \texttt{src2[i]} via bitwise XOR \texttt{res = s \textasciicircum{} (not2 ? \textasciitilde{}src2[i] : src2[i])} and writes \texttt{not\_result ? \textasciitilde{}res : res} to \texttt{dest[i]}.
\par\textbf{Notes:}
\begin{tightitemize}
  \item Shift amount is masked modulo lane width (\texttt{wrapping\_shl}/\texttt{wrapping\_shr}/\texttt{rotate\_*}, \texttt{kraid/ops.rs:3984-4009}).
  \item On \texttt{v2i16}, \texttt{src1} is \texttt{V2I8} (\texttt{ops.rs:3953}) with byte lanes selected by \texttt{lanes1} (default \texttt{b02}: byte 0 -\textgreater{} lane 0, byte 2 -\textgreater{} lane 1).
  \item \textbf{Source references:} \texttt{kraid/ops.rs:3878-4025}, \texttt{encode\_v9.rs:2940-3093}, \texttt{hw\_tests.rs:2148-2198}.
\end{tightitemize}
\codebox{\mnem{RSHIFT\_XOR}\imod{\{.flow\}}\imod{\{.not\_result\}} \dst{dest}\imod{\{.dest\_width\}}, \src{src0}\imod{\{.widen0\}}, \src{src1}\imod{\{.lanes1\}}, \src{src2}\imod{\{.not2\}}\\\mnem{RSHIFT\_XOR}\imod{\{.flow\}}\imod{\{.not\_result\}} \dst{dest}, \src{src0}\imod{\{.widen0\}}, \src{src1}\imod{\{.lanes1\}}, \src{src2}\imod{\{.not2\}}\\\mnem{RSHIFT\_XOR}\imod{\{.flow\}}\imod{\{.not\_result\}} \dst{dest}\imod{\{.dest\_width\}}, \src{src0}\imod{\{.swz0\}}, \src{src1}\imod{\{.lanes1\}}, \src{src2}\imod{\{.not2\}}\\\mnem{RSHIFT\_XOR}\imod{\{.flow\}}\imod{\{.not\_result\}} \dst{dest}\imod{\{.dest\_width\}}, \src{src0}\imod{\{.lanes0\}}, \src{src1}\imod{\{.lanes1\}}, \src{src2}\imod{\{.not2\}}}
\textbf{Operands:}
\begin{tightitemize}\footnotesize
  \item \texttt{\dst{dest}} (\texttt{dest\_register\_index}, \texttt{dest\_width}): Destination GPR receiving the logical right shift combined via bitwise XOR (inverted if \texttt{.not\_result}).
  \item \texttt{\src{src0}}: Input integer value to shift/rotate.
  \item \texttt{\src{src1}}: 8-bit shift/rotate count per lane (selected via \texttt{lanes1}, masked modulo lane width \texttt{W}).
  \item \texttt{\src{src2}}: Second operand combined via bitwise XOR after shifting (inverted if \texttt{.not2}).
\end{tightitemize}
\par\textbf{Closest SPIR-V Op:} \texttt{OpShiftRightLogical + OpBitwiseXor}
\par\textbf{Pseudo-GLSL (Operational Semantics):}
\codebox{for (int i = 0; i \textless{} N; ++i) \{\\    uint s = extractByte(src1, lanes1, i) \& (W - 1);\\    auto r = (uint(src0[i]) \textgreater{}\textgreater{} s) \textasciicircum{} (flag\_not2 ? \textasciitilde{}src2[i] : src2[i]);\\    dest[i] = flag\_not\_result ? \textasciitilde{}r : r;\\\}}
\textbf{Execution Context:}
\begin{tightitemize}\footnotesize
  \item \textbf{Compute stage preloads (\hyperref[sec:csf-spd]{\texttt{SPD\_k}}, \texttt{r16..r23}):} Hardware pre-fills \texttt{r55..r57} with \texttt{WorkGroupID.xyz + JOB\_OFFSET} (from \hyperref[sec:csf-regs]{\texttt{COMPUTESR::JOB\_OFFSET\_X/Y/Z}} \texttt{r34..r36}) and \texttt{r60..r62} with \texttt{LocalInvocationID.xyz} or \texttt{GlobalInvocationID.xyz} (from \texttt{WG\_SIZE} \texttt{r33} \& \texttt{JOB\_SIZE} \texttt{r37..r39}, selected by \texttt{ShaderProgram::preload}) so compute shaders can derive thread/element indices without memory loads.
  \item \textbf{Vertex stage preloads (\hyperref[sec:csf-spd]{\texttt{SPD\_0}/\texttt{SPD\_1}}, \texttt{r16..r19}):} The hardware Index Fetcher reads \hyperref[sec:csf-regs]{\texttt{IDVSSR::INDEX\_BUFFER}} (\texttt{r54:r55}) and pre-fills \texttt{r59, r60} with \texttt{VertexIndex} (\texttt{+ VERTEX\_OFFSET + GLOBAL\_ATTR\_OFF}) and \texttt{r61} with \texttt{InstanceIndex} (\texttt{+ INSTANCE\_OFFSET}) so Position and Varying shaders can pull vertex attributes directly.
  \item \textbf{Inline FAU \& Special Operands (\hyperref[sec:csf-tsd-fau]{\texttt{FAU\_k}}, \texttt{r8..r15}):} Any 8-bit source with \texttt{mode=2} reads zero-latency uniform \texttt{u0..u31} from \texttt{FAU\_k} (\texttt{slots64 = bits[63:56]}); \texttt{mode=3} reads \texttt{SmallConstant} (Table~1.2) or \texttt{FauSpecial} (\texttt{pc}, \texttt{lane\_id}, \texttt{lane\_width}, \texttt{tls\_ptr}, \texttt{wls\_ptr}, \texttt{blend\_descriptor\_0..7}).
\end{tightitemize}
\par\textbf{Worked Examples:}
\par\textbf{Example 1 (Standard fused shift + bitwise op):}
\codebox{\mnem{RSHIFT\_XOR.i32} \dst{r0}, \src{r1}, \src{r2}\imod{.b0}, \src{r3}}
With \texttt{r1 = 0x80000000}, \texttt{r2.b0 = 4} (masked mod 32): shift yields \texttt{0x08000000}, then \texttt{\textasciicircum{} 0x000000FF} yields \texttt{r0 = 0x80000FF}.
\par\textbf{Example 2 (Modulo-32 shift count \& .not\_result):}
\codebox{\mnem{RSHIFT\_XOR.i32}\imod{.not\_result} \dst{r0}, \src{r1}, \src{r2}\imod{.b0}, \src{r3}}
With \texttt{r1 = 0x80000000}, \texttt{r2.b0 = 32} (effective shift \texttt{0}): shift yields \texttt{0x80000000}, \texttt{\textasciicircum{} 0x000000FF} yields \texttt{0x800000FF}, and \texttt{.not\_result} inverts -\textgreater{} \texttt{r0 = 0x7FFFFF00}.

\begin{center}\footnotesize
\begin{tabularx}{\linewidth}{@{}lllX@{}}
\toprule
Variant & Arch & Class & Opcode Fields \\
\midrule
\texttt{(none)} & v9--v14 & \texttt{R3} & \texttt{opcode1=SHIFT\_BITWISE\_I32} \\
\texttt{(none)} & v9--v14 & \texttt{R3} & \texttt{opcode1=SHIFT\_BITWISE\_I32} \\
\texttt{(none)} & v9--v14 & \texttt{R3} & \texttt{opcode1=LOGIC\_SHIFT\_I32} \\
\texttt{(none)} & v9--v14 & \texttt{D3} & \texttt{opcode1=SHIFT\_BITWISE\_I64} \\
\texttt{(none)} & v9--v14 & \texttt{D3} & \texttt{opcode1=LOGIC\_SHIFT\_I64} \\
\texttt{(none)} & v9--v14 & \texttt{R3} & \texttt{opcode1=SHIFT\_BITWISE\_V2I16} \\
\texttt{(none)} & v9--v14 & \texttt{R3} & \texttt{opcode1=LOGIC\_SHIFT\_V2I16} \\
\texttt{(none)} & v9--v14 & \texttt{R3} & \texttt{opcode1=SHIFT\_BITWISE\_V4I8} \\
\texttt{(none)} & v9--v14 & \texttt{R3} & \texttt{opcode1=LOGIC\_SHIFT\_V4I8} \\
\bottomrule
\end{tabularx}
\end{center}
\begin{center}\footnotesize
\begin{tabularx}{0.88\linewidth}{@{}l l X@{}}
\toprule
Field Name & Bits & Type / Fixed Value \\
\midrule
\texttt{src0} & \texttt{0:7} & \texttt{SourceEncoding} \\
\texttt{src1} & \texttt{8:15} & \texttt{SourceEncoding} \\
\texttt{src2} & \texttt{16:23} & \texttt{SourceEncoding} \\
\texttt{bitwise\_operation} & \texttt{24:25} & \texttt{exact=xor} \\
\hyperref[enum:shift_bitwise_i32_lanes_m]{\texttt{lanes1}} & \texttt{26:29} & \texttt{shift\_bitwise\_i32\_lanes\_m (def b0)} \\
\hyperref[enum:not_m]{\texttt{not\_result}} & \texttt{30} & \texttt{not\_m (def none)} \\
\textit{Reserved} & \texttt{31} & \texttt{zero} \\
\texttt{shift\_operation} & \texttt{32:34} & \texttt{exact=rshift} \\
\hyperref[enum:not_m]{\texttt{not2}} & \texttt{35} & \texttt{not\_m (def none)} \\
\hyperref[enum:widen_int_m]{\texttt{widen0}} & \texttt{36:39} & \texttt{widen\_int\_m (def none)} \\
\texttt{dest\_register\_index} & \texttt{40:45} & \texttt{uint} \\
\hyperref[enum:dest_width_m]{\texttt{dest\_width}} & \texttt{46:47} & \texttt{dest\_width\_m} \\
\texttt{opcode1} & \texttt{48:51} & \texttt{exact=SHIFT\_BITWISE\_I32} \\
\texttt{instruction\_class} & \texttt{52:56} & \texttt{exact=R3} \\
\texttt{fau\_page\_index} & \texttt{57:58} & \texttt{uint} \\
\hyperref[enum:flow_control_m]{\texttt{flow}} & \texttt{59:62} & \texttt{flow\_control\_m} \\
\textit{Reserved} & \texttt{63} & \texttt{zero} \\
\bottomrule
\end{tabularx}
\end{center}
\textbf{Fields Explained:}
\begin{tightitemize}\footnotesize
  \item \texttt{src0}: Primary source operand (\texttt{r0..r63}, \texttt{u0..u63} uniform, \texttt{k*} FAU constant, or \texttt{t0..t2} temporary; \texttt{\textasciicircum{}} marks discard/last-use). \emph{Example:} {\scriptsize\ttfamily \mnem{FADD.f32} \dst{r0}, \src{r1}, \src{r2}} (\textbf{SPIR-V:} \texttt{OpFAdd (Operand 0)}; \textbf{GLSL:} \texttt{float c = a + b;})
  \item \texttt{src1}: Secondary source operand (register, uniform, or FAU constant; in shifts, supplies the shift count byte). \emph{Example:} {\scriptsize\ttfamily \mnem{LSHIFT.i32} \dst{r0}, \src{r1}, \src{r2}\imod{.b0}} (\textbf{SPIR-V:} \texttt{OpShiftLeftLogical (Shift)}; \textbf{GLSL:} \texttt{uint y = a \textless{}\textless{} s;})
  \item \texttt{src2}: Third source operand (FMA addend, \texttt{MUX}/\texttt{CSEL} selector, \texttt{IADDC} carry-in, or bitwise mask in fused shift-logic). \emph{Example:} {\scriptsize\ttfamily \mnem{FMA.f32} \dst{r0}, \src{r1}, \src{r2}, \src{r3}} (\textbf{SPIR-V:} \texttt{GLSL.std.450 Fma (Addend)}; \textbf{GLSL:} \texttt{float d = fma(a, b, c);})
  \item \texttt{bitwise\_operation}: Selects the bitwise ALU op (\texttt{and}, \texttt{or}, \texttt{xor}, \texttt{lshift\_and}, etc.) fused with shifts. \emph{Example:} {\scriptsize\ttfamily \mnem{LSHIFT\_OR.i32} \dst{r0}, \src{r1}, \src{r2}, \src{r3}} (\textbf{SPIR-V:} \texttt{OpShiftLeftLogical + OpBitwiseOr}; \textbf{GLSL:} \texttt{uint y = (a \textless{}\textless{} s) \textbar{} b;})
  \item \hyperref[enum:shift_bitwise_i32_lanes_m]{\texttt{lanes1}}: Stage A byte-lane selector on \texttt{src1} (selects shift-count byte \texttt{b0..b3} on 32/64-bit shifts, \texttt{bXY} on \texttt{v2i16} shifts, or 4-byte swizzle on \texttt{v4*8}). \emph{Example:} {\scriptsize\ttfamily \mnem{LSHIFT.i32} \dst{r0}, \src{r1}, \src{r2}\imod{.b1}} (\textbf{SPIR-V:} \texttt{OpBitFieldUExtract + OpShiftLeftLogical}; \textbf{GLSL:} \texttt{uint y = a \textless{}\textless{} ((s \textgreater{}\textgreater{} 8) \& 0xFFu);})
  \item \hyperref[enum:not_m]{\texttt{not\_result}}: Bitwise-inverts the final output of a bitwise/shift-bitwise instruction (\textasciitilde{}result). \emph{Example:} {\scriptsize\ttfamily \mnem{AND.i32}\imod{.not\_result} \dst{r0}, \src{r1}, \src{r2}} (\textbf{SPIR-V:} \texttt{OpBitwiseAnd + OpNot}; \textbf{GLSL:} \texttt{uint y = \textasciitilde{}(a \& b);})
  \item \texttt{shift\_operation}: Selects shift direction/type (\texttt{lshift}, \texttt{rshift}, \texttt{arshift}) in fused shift-bitwise instructions. \emph{Example:} {\scriptsize\ttfamily \mnem{RSHIFT\_AND.i32} \dst{r0}, \src{r1}, \src{r2}, \src{r3}} (\textbf{SPIR-V:} \texttt{OpShiftRightLogical + OpBitwiseAnd}; \textbf{GLSL:} \texttt{uint y = (a \textgreater{}\textgreater{} s) \& m;})
  \item \hyperref[enum:not_m]{\texttt{not2}}: Stage A bitwise complement (\texttt{\textasciitilde{}src2}) applied to \texttt{src2} before the bitwise/shift-logic operation. \emph{Example:} {\scriptsize\ttfamily \mnem{AND.i32} \dst{r0}, \src{r1}, \src{r2}\imod{.not}} (\textbf{SPIR-V:} \texttt{OpNot + OpBitwiseAnd}; \textbf{GLSL:} \texttt{uint y = a \& \textasciitilde{}b;})
  \item \hyperref[enum:widen_int_m]{\texttt{widen0}}: Stage A widening extractor on \texttt{src0} (promotes FP16 \texttt{h0}/\texttt{h1} to FP32, or sign/zero-extends \texttt{b0..b3}/\texttt{h0..h1}/\texttt{w0..w1} to 32/64-bit before \texttt{abs}/\texttt{neg}). \emph{Example:} {\scriptsize\ttfamily \mnem{IADD.s32} \dst{r0}, \src{r1}\imod{.h0}, \src{r2}} (\textbf{SPIR-V:} \texttt{OpSConvert / OpFConvert}; \textbf{GLSL:} \texttt{int c = int(int16\_t(a)) + b;})
  \item \texttt{dest\_register\_index}: 6-bit destination GPR index (\texttt{r0..r63}) receiving the Stage C writeback result (even-aligned pair \texttt{[rN:rN+1]} for 64-bit ops). \emph{Example:} {\scriptsize\ttfamily \mnem{FMA.f32} \dst{r0}, \src{r1}, \src{r2}, \src{r3}} (\textbf{SPIR-V:} \texttt{SSA Result ID}; \textbf{GLSL:} \texttt{float r0 = fma(r1, r2, r3);})
  \item \hyperref[enum:dest_width_m]{\texttt{dest\_width}}: Stage C destination lane write-mask (\texttt{none} = full 32-bit, \texttt{h0} = low 16-bit, \texttt{h1} = high 16-bit without bit-shifting). \emph{Example:} {\scriptsize\ttfamily \mnem{FADD.v2f16} \dst{r0}\imod{.h1}, \src{r1}, \src{r2}} (\textbf{SPIR-V:} \texttt{OpVectorInsertDynamic / WriteMask}; \textbf{GLSL:} \texttt{dst.y = a.y + b.y;})
  \item \texttt{opcode1}: Primary sub-opcode within the instruction class uniquely identifying the operation family. \emph{Example:} {\scriptsize\ttfamily \mnem{IADD.s32} \dst{r0}, \src{r1}, \src{r2}} (\textbf{SPIR-V:} \texttt{OpIAdd (operation selector)}; \textbf{GLSL:} \texttt{// Primary opcode bits})
  \item \texttt{instruction\_class}: 4-bit major instruction class selector (\texttt{bits 48:51}) partitioning the 64-bit Valhall opcode space (e.g., \texttt{FMA}, \texttt{CVT}, \texttt{SFU}, \texttt{LS}, \texttt{TEX}, \texttt{BRANCH}). \emph{Example:} {\scriptsize\ttfamily \mnem{FADD.f32} \dst{r0}, \src{r1}, \src{r2}} (\textbf{SPIR-V:} \texttt{N/A (encoding class)}; \textbf{GLSL:} \texttt{// Hardware execution pipe / decode class})
  \item \texttt{fau\_page\_index}: Selects the FAU (Fast Access Uniform) page (\texttt{0..3}) when a source operand references hardware small constants, thread IDs, or lane constants. \emph{Example:} {\scriptsize\ttfamily \mnem{FADD.f32} \dst{r0}, \src{r1}, \src{k1f}} (\textbf{SPIR-V:} \texttt{OpConstant / BuiltIn}; \textbf{GLSL:} \texttt{float y = x + 1.0;})
  \item \hyperref[enum:flow_control_m]{\texttt{flow}}: 4-bit embedded flow-control \& scoreboard wait modifier (\texttt{bits 59:62}: \texttt{none}, \texttt{wait0..wait0126}, \texttt{wait}, \texttt{reconverge}, \texttt{discard}, \texttt{end}). \emph{Example:} {\scriptsize\ttfamily \mnem{FADD.f32}\imod{.wait0.reconverge} \dst{r0}, \src{r1}, \src{r2}} (\textbf{SPIR-V:} \texttt{OpControlBarrier / OpReturn}; \textbf{GLSL:} \texttt{// Scoreboard wait + warp reconvergence})
\end{tightitemize}
\clearpage

\subsection*{XOR}
\addcontentsline{toc}{subsection}{XOR}
Computes bitwise XOR independently across each SIMD element $i \in \{0, \dots, N-1\}$ ($N=1$ for \texttt{i32}/\texttt{i64}, $N=2$ for \texttt{v2i16}, $N=4$ for \texttt{v4i8} on v9--v10): \texttt{res = widen0(src0)[i] \textasciicircum{} (not2 ? \textasciitilde{}src2[i] : src2[i])}, then writes \texttt{not\_result ? \textasciitilde{}res : res} to \texttt{dest[i]}.
\par\textbf{Notes:}
\begin{tightitemize}
  \item Second operand is \texttt{src2}.
  \item \textbf{Source references:} \texttt{kraid/ops.rs:3944-4025}, \texttt{encode\_v9.rs:2916-3093}, \texttt{hw\_tests.rs:2148-2198}.
\end{tightitemize}
\codebox{\mnem{XOR}\imod{\{.flow\}}\imod{\{.not\_result\}} \dst{dest}\imod{\{.dest\_width\}}, \src{src0}\imod{\{.widen0\}}, \src{src2}\imod{\{.not2\}}\\\mnem{XOR}\imod{\{.flow\}}\imod{\{.not\_result\}} \dst{dest}, \src{src0}\imod{\{.widen0\}}, \src{src2}\imod{\{.not2\}}\\\mnem{XOR}\imod{\{.flow\}}\imod{\{.not\_result\}} \dst{dest}\imod{\{.dest\_width\}}, \src{src0}\imod{\{.swz0\}}, \src{src2}\imod{\{.not2\}}\\\mnem{XOR}\imod{\{.flow\}}\imod{\{.not\_result\}} \dst{dest}\imod{\{.dest\_width\}}, \src{src0}\imod{\{.lanes0\}}, \src{src2}\imod{\{.not2\}}}
\textbf{Operands:}
\begin{tightitemize}\footnotesize
  \item \texttt{\dst{dest}} (\texttt{dest\_register\_index}, \texttt{dest\_width}): Destination GPR receiving \texttt{src0 \textasciicircum{} src2} (or XNOR with \texttt{.not\_result}).
  \item \texttt{\src{src0}}: First bitwise operand.
  \item \texttt{\src{src2}}: Second bitwise operand (encoded in \texttt{src2}, not \texttt{src1}; supports \texttt{.not2}).
\end{tightitemize}
\par\textbf{Closest SPIR-V Op:} \texttt{OpBitwiseXor}
\par\textbf{Pseudo-GLSL (Operational Semantics):}
\codebox{for (int i = 0; i \textless{} N; ++i) \{\\    auto b = flag\_not2 ? \textasciitilde{}src2[i] : src2[i];\\    auto r = widen0(src0)[i] \textasciicircum{} b;\\    dest[i] = flag\_not\_result ? \textasciitilde{}r : r; // XNOR when .not\_result\\\}}
\textbf{Execution Context:}
\begin{tightitemize}\footnotesize
  \item \textbf{Compute stage preloads (\hyperref[sec:csf-spd]{\texttt{SPD\_k}}, \texttt{r16..r23}):} Hardware pre-fills \texttt{r55..r57} with \texttt{WorkGroupID.xyz + JOB\_OFFSET} (from \hyperref[sec:csf-regs]{\texttt{COMPUTESR::JOB\_OFFSET\_X/Y/Z}} \texttt{r34..r36}) and \texttt{r60..r62} with \texttt{LocalInvocationID.xyz} or \texttt{GlobalInvocationID.xyz} (from \texttt{WG\_SIZE} \texttt{r33} \& \texttt{JOB\_SIZE} \texttt{r37..r39}, selected by \texttt{ShaderProgram::preload}) so compute shaders can derive thread/element indices without memory loads.
  \item \textbf{Vertex stage preloads (\hyperref[sec:csf-spd]{\texttt{SPD\_0}/\texttt{SPD\_1}}, \texttt{r16..r19}):} The hardware Index Fetcher reads \hyperref[sec:csf-regs]{\texttt{IDVSSR::INDEX\_BUFFER}} (\texttt{r54:r55}) and pre-fills \texttt{r59, r60} with \texttt{VertexIndex} (\texttt{+ VERTEX\_OFFSET + GLOBAL\_ATTR\_OFF}) and \texttt{r61} with \texttt{InstanceIndex} (\texttt{+ INSTANCE\_OFFSET}) so Position and Varying shaders can pull vertex attributes directly.
  \item \textbf{Inline FAU \& Special Operands (\hyperref[sec:csf-tsd-fau]{\texttt{FAU\_k}}, \texttt{r8..r15}):} Any 8-bit source with \texttt{mode=2} reads zero-latency uniform \texttt{u0..u31} from \texttt{FAU\_k} (\texttt{slots64 = bits[63:56]}); \texttt{mode=3} reads \texttt{SmallConstant} (Table~1.2) or \texttt{FauSpecial} (\texttt{pc}, \texttt{lane\_id}, \texttt{lane\_width}, \texttt{tls\_ptr}, \texttt{wls\_ptr}, \texttt{blend\_descriptor\_0..7}).
\end{tightitemize}
\par\textbf{Worked Examples:}
\par\textbf{Example 1 (Bitwise toggle):}
\codebox{\mnem{XOR.i32} \dst{r0}, \src{r1}, \src{r2}}
With \texttt{r1 = 0xAAAA5555, r2 = 0xFFFFFFFF}: flips all 32 bits -\textgreater{} \texttt{r0 = 0x5555AAAA}.
\par\textbf{Example 2 (XNOR via .not\_result):}
\codebox{\mnem{XOR.i32}\imod{.not\_result} \dst{r0}, \src{r1}, \src{r2}}
With \texttt{r1 = 0x000000FF, r2 = 0x000000FF}: \texttt{r1 \textasciicircum{} r2 = 0}, then \texttt{.not\_result} yields \texttt{r0 = 0xFFFFFFFF} (bitwise equality mask).

\begin{center}\footnotesize
\begin{tabularx}{\linewidth}{@{}lllX@{}}
\toprule
Variant & Arch & Class & Opcode Fields \\
\midrule
\texttt{(none)} & v9--v14 & \texttt{R3} & \texttt{opcode1=SHIFT\_BITWISE\_I32} \\
\texttt{(none)} & v9--v14 & \texttt{R3} & \texttt{opcode1=SHIFT\_BITWISE\_I32} \\
\texttt{(none)} & v9--v14 & \texttt{R3} & \texttt{opcode1=LOGIC\_SHIFT\_I32} \\
\texttt{(none)} & v9--v14 & \texttt{D3} & \texttt{opcode1=SHIFT\_BITWISE\_I64} \\
\texttt{(none)} & v9--v14 & \texttt{D3} & \texttt{opcode1=LOGIC\_SHIFT\_I64} \\
\texttt{(none)} & v9--v14 & \texttt{R3} & \texttt{opcode1=SHIFT\_BITWISE\_V2I16} \\
\texttt{(none)} & v9--v14 & \texttt{R3} & \texttt{opcode1=LOGIC\_SHIFT\_V2I16} \\
\texttt{(none)} & v9--v14 & \texttt{R3} & \texttt{opcode1=SHIFT\_BITWISE\_V4I8} \\
\texttt{(none)} & v9--v14 & \texttt{R3} & \texttt{opcode1=LOGIC\_SHIFT\_V4I8} \\
\bottomrule
\end{tabularx}
\end{center}
\begin{center}\footnotesize
\begin{tabularx}{0.88\linewidth}{@{}l l X@{}}
\toprule
Field Name & Bits & Type / Fixed Value \\
\midrule
\texttt{src0} & \texttt{0:7} & \texttt{SourceEncoding} \\
\texttt{src1} & \texttt{8:15} & \texttt{exact=192} \\
\texttt{src2} & \texttt{16:23} & \texttt{SourceEncoding} \\
\texttt{bitwise\_operation} & \texttt{24:25} & \texttt{exact=xor} \\
\hyperref[enum:shift_bitwise_i32_lanes_m]{\texttt{lanes1}} & \texttt{26:29} & \texttt{exact=b0} \\
\hyperref[enum:not_m]{\texttt{not\_result}} & \texttt{30} & \texttt{not\_m (def none)} \\
\textit{Reserved} & \texttt{31} & \texttt{zero} \\
\texttt{shift\_operation} & \texttt{32:34} & \texttt{exact=lshift} \\
\hyperref[enum:not_m]{\texttt{not2}} & \texttt{35} & \texttt{not\_m (def none)} \\
\hyperref[enum:widen_int_m]{\texttt{widen0}} & \texttt{36:39} & \texttt{widen\_int\_m (def none)} \\
\texttt{dest\_register\_index} & \texttt{40:45} & \texttt{uint} \\
\hyperref[enum:dest_width_m]{\texttt{dest\_width}} & \texttt{46:47} & \texttt{dest\_width\_m} \\
\texttt{opcode1} & \texttt{48:51} & \texttt{exact=SHIFT\_BITWISE\_I32} \\
\texttt{instruction\_class} & \texttt{52:56} & \texttt{exact=R3} \\
\texttt{fau\_page\_index} & \texttt{57:58} & \texttt{uint} \\
\hyperref[enum:flow_control_m]{\texttt{flow}} & \texttt{59:62} & \texttt{flow\_control\_m} \\
\textit{Reserved} & \texttt{63} & \texttt{zero} \\
\bottomrule
\end{tabularx}
\end{center}
\textbf{Fields Explained:}
\begin{tightitemize}\footnotesize
  \item \texttt{src0}: Primary source operand (\texttt{r0..r63}, \texttt{u0..u63} uniform, \texttt{k*} FAU constant, or \texttt{t0..t2} temporary; \texttt{\textasciicircum{}} marks discard/last-use). \emph{Example:} {\scriptsize\ttfamily \mnem{FADD.f32} \dst{r0}, \src{r1}, \src{r2}} (\textbf{SPIR-V:} \texttt{OpFAdd (Operand 0)}; \textbf{GLSL:} \texttt{float c = a + b;})
  \item \texttt{src1}: Secondary source operand (register, uniform, or FAU constant; in shifts, supplies the shift count byte). \emph{Example:} {\scriptsize\ttfamily \mnem{LSHIFT.i32} \dst{r0}, \src{r1}, \src{r2}\imod{.b0}} (\textbf{SPIR-V:} \texttt{OpShiftLeftLogical (Shift)}; \textbf{GLSL:} \texttt{uint y = a \textless{}\textless{} s;})
  \item \texttt{src2}: Third source operand (FMA addend, \texttt{MUX}/\texttt{CSEL} selector, \texttt{IADDC} carry-in, or bitwise mask in fused shift-logic). \emph{Example:} {\scriptsize\ttfamily \mnem{FMA.f32} \dst{r0}, \src{r1}, \src{r2}, \src{r3}} (\textbf{SPIR-V:} \texttt{GLSL.std.450 Fma (Addend)}; \textbf{GLSL:} \texttt{float d = fma(a, b, c);})
  \item \texttt{bitwise\_operation}: Selects the bitwise ALU op (\texttt{and}, \texttt{or}, \texttt{xor}, \texttt{lshift\_and}, etc.) fused with shifts. \emph{Example:} {\scriptsize\ttfamily \mnem{LSHIFT\_OR.i32} \dst{r0}, \src{r1}, \src{r2}, \src{r3}} (\textbf{SPIR-V:} \texttt{OpShiftLeftLogical + OpBitwiseOr}; \textbf{GLSL:} \texttt{uint y = (a \textless{}\textless{} s) \textbar{} b;})
  \item \hyperref[enum:shift_bitwise_i32_lanes_m]{\texttt{lanes1}}: Stage A byte-lane selector on \texttt{src1} (selects shift-count byte \texttt{b0..b3} on 32/64-bit shifts, \texttt{bXY} on \texttt{v2i16} shifts, or 4-byte swizzle on \texttt{v4*8}). \emph{Example:} {\scriptsize\ttfamily \mnem{LSHIFT.i32} \dst{r0}, \src{r1}, \src{r2}\imod{.b1}} (\textbf{SPIR-V:} \texttt{OpBitFieldUExtract + OpShiftLeftLogical}; \textbf{GLSL:} \texttt{uint y = a \textless{}\textless{} ((s \textgreater{}\textgreater{} 8) \& 0xFFu);})
  \item \hyperref[enum:not_m]{\texttt{not\_result}}: Bitwise-inverts the final output of a bitwise/shift-bitwise instruction (\textasciitilde{}result). \emph{Example:} {\scriptsize\ttfamily \mnem{AND.i32}\imod{.not\_result} \dst{r0}, \src{r1}, \src{r2}} (\textbf{SPIR-V:} \texttt{OpBitwiseAnd + OpNot}; \textbf{GLSL:} \texttt{uint y = \textasciitilde{}(a \& b);})
  \item \texttt{shift\_operation}: Selects shift direction/type (\texttt{lshift}, \texttt{rshift}, \texttt{arshift}) in fused shift-bitwise instructions. \emph{Example:} {\scriptsize\ttfamily \mnem{RSHIFT\_AND.i32} \dst{r0}, \src{r1}, \src{r2}, \src{r3}} (\textbf{SPIR-V:} \texttt{OpShiftRightLogical + OpBitwiseAnd}; \textbf{GLSL:} \texttt{uint y = (a \textgreater{}\textgreater{} s) \& m;})
  \item \hyperref[enum:not_m]{\texttt{not2}}: Stage A bitwise complement (\texttt{\textasciitilde{}src2}) applied to \texttt{src2} before the bitwise/shift-logic operation. \emph{Example:} {\scriptsize\ttfamily \mnem{AND.i32} \dst{r0}, \src{r1}, \src{r2}\imod{.not}} (\textbf{SPIR-V:} \texttt{OpNot + OpBitwiseAnd}; \textbf{GLSL:} \texttt{uint y = a \& \textasciitilde{}b;})
  \item \hyperref[enum:widen_int_m]{\texttt{widen0}}: Stage A widening extractor on \texttt{src0} (promotes FP16 \texttt{h0}/\texttt{h1} to FP32, or sign/zero-extends \texttt{b0..b3}/\texttt{h0..h1}/\texttt{w0..w1} to 32/64-bit before \texttt{abs}/\texttt{neg}). \emph{Example:} {\scriptsize\ttfamily \mnem{IADD.s32} \dst{r0}, \src{r1}\imod{.h0}, \src{r2}} (\textbf{SPIR-V:} \texttt{OpSConvert / OpFConvert}; \textbf{GLSL:} \texttt{int c = int(int16\_t(a)) + b;})
  \item \texttt{dest\_register\_index}: 6-bit destination GPR index (\texttt{r0..r63}) receiving the Stage C writeback result (even-aligned pair \texttt{[rN:rN+1]} for 64-bit ops). \emph{Example:} {\scriptsize\ttfamily \mnem{FMA.f32} \dst{r0}, \src{r1}, \src{r2}, \src{r3}} (\textbf{SPIR-V:} \texttt{SSA Result ID}; \textbf{GLSL:} \texttt{float r0 = fma(r1, r2, r3);})
  \item \hyperref[enum:dest_width_m]{\texttt{dest\_width}}: Stage C destination lane write-mask (\texttt{none} = full 32-bit, \texttt{h0} = low 16-bit, \texttt{h1} = high 16-bit without bit-shifting). \emph{Example:} {\scriptsize\ttfamily \mnem{FADD.v2f16} \dst{r0}\imod{.h1}, \src{r1}, \src{r2}} (\textbf{SPIR-V:} \texttt{OpVectorInsertDynamic / WriteMask}; \textbf{GLSL:} \texttt{dst.y = a.y + b.y;})
  \item \texttt{opcode1}: Primary sub-opcode within the instruction class uniquely identifying the operation family. \emph{Example:} {\scriptsize\ttfamily \mnem{IADD.s32} \dst{r0}, \src{r1}, \src{r2}} (\textbf{SPIR-V:} \texttt{OpIAdd (operation selector)}; \textbf{GLSL:} \texttt{// Primary opcode bits})
  \item \texttt{instruction\_class}: 4-bit major instruction class selector (\texttt{bits 48:51}) partitioning the 64-bit Valhall opcode space (e.g., \texttt{FMA}, \texttt{CVT}, \texttt{SFU}, \texttt{LS}, \texttt{TEX}, \texttt{BRANCH}). \emph{Example:} {\scriptsize\ttfamily \mnem{FADD.f32} \dst{r0}, \src{r1}, \src{r2}} (\textbf{SPIR-V:} \texttt{N/A (encoding class)}; \textbf{GLSL:} \texttt{// Hardware execution pipe / decode class})
  \item \texttt{fau\_page\_index}: Selects the FAU (Fast Access Uniform) page (\texttt{0..3}) when a source operand references hardware small constants, thread IDs, or lane constants. \emph{Example:} {\scriptsize\ttfamily \mnem{FADD.f32} \dst{r0}, \src{r1}, \src{k1f}} (\textbf{SPIR-V:} \texttt{OpConstant / BuiltIn}; \textbf{GLSL:} \texttt{float y = x + 1.0;})
  \item \hyperref[enum:flow_control_m]{\texttt{flow}}: 4-bit embedded flow-control \& scoreboard wait modifier (\texttt{bits 59:62}: \texttt{none}, \texttt{wait0..wait0126}, \texttt{wait}, \texttt{reconverge}, \texttt{discard}, \texttt{end}). \emph{Example:} {\scriptsize\ttfamily \mnem{FADD.f32}\imod{.wait0.reconverge} \dst{r0}, \src{r1}, \src{r2}} (\textbf{SPIR-V:} \texttt{OpControlBarrier / OpReturn}; \textbf{GLSL:} \texttt{// Scoreboard wait + warp reconvergence})
\end{tightitemize}
\clearpage

\subsection*{XOR\_IMM}
\addcontentsline{toc}{subsection}{XOR\_IMM}
Computes a 32-bit bitwise XOR between \texttt{src0} and the inline 32-bit immediate \texttt{\#imm1w} (\texttt{dest = src0 \textasciicircum{} \#imm1w}) on Valhall v11+. Does not support source swizzles or \texttt{.not\_result}; supports \texttt{dest\_width} halfword write masking.
\par\textbf{Notes:}
\begin{tightitemize}
  \item No \texttt{not\_result} or \texttt{src0} swizzle supported.
  \item \textbf{Source references:} \texttt{kraid/encode\_v9.rs:2919-2926, 3036-3054}.
\end{tightitemize}
\codebox{\mnem{XOR\_IMM}\imod{\{.flow\}} \dst{dest}\imod{\{.dest\_width\}}, \src{src0}, \imm{\#imm1w}}
\textbf{Operands:}
\begin{tightitemize}\footnotesize
  \item \texttt{\dst{dest}} (\texttt{dest\_register\_index}, \texttt{dest\_width}): 32-bit destination GPR receiving \texttt{src0 \textasciicircum{} \#imm1w}.
  \item \texttt{\src{src0}}: 32-bit input operand.
  \item \texttt{\imm{\#imm1w}} (\texttt{imm1w}): Inline 32-bit bitwise XOR mask.
\end{tightitemize}
\par\textbf{Closest SPIR-V Op:} \texttt{OpBitwiseXor (constant)}
\par\textbf{Pseudo-GLSL (Operational Semantics):}
\codebox{dest = maskWrite(dest\_width, src0 \textasciicircum{} \#imm1w);}
\textbf{Execution Context:}
\begin{tightitemize}\footnotesize
  \item \textbf{Compute stage preloads (\hyperref[sec:csf-spd]{\texttt{SPD\_k}}, \texttt{r16..r23}):} Hardware pre-fills \texttt{r55..r57} with \texttt{WorkGroupID.xyz + JOB\_OFFSET} (from \hyperref[sec:csf-regs]{\texttt{COMPUTESR::JOB\_OFFSET\_X/Y/Z}} \texttt{r34..r36}) and \texttt{r60..r62} with \texttt{LocalInvocationID.xyz} or \texttt{GlobalInvocationID.xyz} (from \texttt{WG\_SIZE} \texttt{r33} \& \texttt{JOB\_SIZE} \texttt{r37..r39}, selected by \texttt{ShaderProgram::preload}) so compute shaders can derive thread/element indices without memory loads.
  \item \textbf{Vertex stage preloads (\hyperref[sec:csf-spd]{\texttt{SPD\_0}/\texttt{SPD\_1}}, \texttt{r16..r19}):} The hardware Index Fetcher reads \hyperref[sec:csf-regs]{\texttt{IDVSSR::INDEX\_BUFFER}} (\texttt{r54:r55}) and pre-fills \texttt{r59, r60} with \texttt{VertexIndex} (\texttt{+ VERTEX\_OFFSET + GLOBAL\_ATTR\_OFF}) and \texttt{r61} with \texttt{InstanceIndex} (\texttt{+ INSTANCE\_OFFSET}) so Position and Varying shaders can pull vertex attributes directly.
  \item \textbf{Inline FAU \& Special Operands (\hyperref[sec:csf-tsd-fau]{\texttt{FAU\_k}}, \texttt{r8..r15}):} Any 8-bit source with \texttt{mode=2} reads zero-latency uniform \texttt{u0..u31} from \texttt{FAU\_k} (\texttt{slots64 = bits[63:56]}); \texttt{mode=3} reads \texttt{SmallConstant} (Table~1.2) or \texttt{FauSpecial} (\texttt{pc}, \texttt{lane\_id}, \texttt{lane\_width}, \texttt{tls\_ptr}, \texttt{wls\_ptr}, \texttt{blend\_descriptor\_0..7}).
\end{tightitemize}
\par\textbf{Worked Examples:}
\par\textbf{Example 1 (Flip FP32 sign bit (v11+)):}
\codebox{\mnem{XOR\_IMM.i32} \dst{r0}, \src{r1}, \imm{\#0x80000000}}
With \texttt{r1 = 0xBF800000} (\texttt{-1.0f}): toggles sign bit -\textgreater{} \texttt{r0 = 0x3F800000} (\texttt{+1.0f}).
\par\textbf{Example 2 (Toggle dual FP16 sign bits):}
\codebox{\mnem{XOR\_IMM.i32} \dst{r0}, \src{r1}, \imm{\#0x80008000}}
With \texttt{r1 = 0x3C00BC00} (\texttt{[+1.0h, -1.0h]}): toggles both halfword sign bits -\textgreater{} \texttt{r0 = 0xBC003C00} (\texttt{[-1.0h, +1.0h]}).
\par\textbf{Mined SPIR-V $\to$ Valhall ISA Example (Inverting a comparison mask with inline 32-bit XOR on v14 --- Compute shader, Immortalis-G925 (v14)):}
\codebox{\textbf{; SPIR-V (input)}\\\%oob = OpLogicalOr \%bool \%lt\_min \%gt\_max\\\%in\_bounds = OpLogicalNot \%bool \%oob\\OpBranchConditional \%in\_bounds \%body \%exit\\[3pt]\textbf{; Valhall ISA (malioc output)}\\FCMP.f32.m1.lt r12, r15, u2.w0\\OR.i32 r12, r12\textasciicircum{}, r50\textasciicircum{}\\XOR\_IMM.i32 r12, r12\textasciicircum{}, \#0xffffffff\\BRANCH.not r12, 0x248}
On Immortalis-G925 (v14), \texttt{malioc} emits \texttt{XOR\_IMM.i32 r12, r12\textasciicircum{}, \#0xffffffff} to flip all 32 bits of the \texttt{.m1} predicate mask before branching.
\begin{center}\footnotesize
\begin{tabularx}{\linewidth}{@{}lllX@{}}
\toprule
Variant & Arch & Class & Opcode Fields \\
\midrule
\texttt{(none)} & v9--v14 & \texttt{R1I} & \texttt{opcode1=XOR\_IMM\_I32} \\
\bottomrule
\end{tabularx}
\end{center}
\begin{center}\footnotesize
\begin{tabularx}{0.88\linewidth}{@{}l l X@{}}
\toprule
Field Name & Bits & Type / Fixed Value \\
\midrule
\texttt{src0} & \texttt{0:7} & \texttt{SourceEncoding} \\
\texttt{imm1w} & \texttt{8:39} & \texttt{uint} \\
\texttt{dest\_register\_index} & \texttt{40:45} & \texttt{uint} \\
\hyperref[enum:dest_width_m]{\texttt{dest\_width}} & \texttt{46:47} & \texttt{dest\_width\_m} \\
\texttt{opcode1} & \texttt{48:51} & \texttt{exact=XOR\_IMM\_I32} \\
\texttt{instruction\_class} & \texttt{52:56} & \texttt{exact=R1I} \\
\texttt{fau\_page\_index} & \texttt{57:58} & \texttt{uint} \\
\hyperref[enum:flow_control_m]{\texttt{flow}} & \texttt{59:62} & \texttt{flow\_control\_m} \\
\textit{Reserved} & \texttt{63} & \texttt{zero} \\
\bottomrule
\end{tabularx}
\end{center}
\textbf{Fields Explained:}
\begin{tightitemize}\footnotesize
  \item \texttt{src0}: Primary source operand (\texttt{r0..r63}, \texttt{u0..u63} uniform, \texttt{k*} FAU constant, or \texttt{t0..t2} temporary; \texttt{\textasciicircum{}} marks discard/last-use). \emph{Example:} {\scriptsize\ttfamily \mnem{FADD.f32} \dst{r0}, \src{r1}, \src{r2}} (\textbf{SPIR-V:} \texttt{OpFAdd (Operand 0)}; \textbf{GLSL:} \texttt{float c = a + b;})
  \item \texttt{imm1w}: 32-bit inline immediate operand encoded directly in the instruction word (bits 8:39). \emph{Example:} {\scriptsize\ttfamily \mnem{IADD\_IMM.i32} \dst{r0}, \src{r1}, \imm{\#0x10}} (\textbf{SPIR-V:} \texttt{OpIAdd (constant)}; \textbf{GLSL:} \texttt{int y = x + 16;})
  \item \texttt{dest\_register\_index}: 6-bit destination GPR index (\texttt{r0..r63}) receiving the Stage C writeback result (even-aligned pair \texttt{[rN:rN+1]} for 64-bit ops). \emph{Example:} {\scriptsize\ttfamily \mnem{FMA.f32} \dst{r0}, \src{r1}, \src{r2}, \src{r3}} (\textbf{SPIR-V:} \texttt{SSA Result ID}; \textbf{GLSL:} \texttt{float r0 = fma(r1, r2, r3);})
  \item \hyperref[enum:dest_width_m]{\texttt{dest\_width}}: Stage C destination lane write-mask (\texttt{none} = full 32-bit, \texttt{h0} = low 16-bit, \texttt{h1} = high 16-bit without bit-shifting). \emph{Example:} {\scriptsize\ttfamily \mnem{FADD.v2f16} \dst{r0}\imod{.h1}, \src{r1}, \src{r2}} (\textbf{SPIR-V:} \texttt{OpVectorInsertDynamic / WriteMask}; \textbf{GLSL:} \texttt{dst.y = a.y + b.y;})
  \item \texttt{opcode1}: Primary sub-opcode within the instruction class uniquely identifying the operation family. \emph{Example:} {\scriptsize\ttfamily \mnem{IADD.s32} \dst{r0}, \src{r1}, \src{r2}} (\textbf{SPIR-V:} \texttt{OpIAdd (operation selector)}; \textbf{GLSL:} \texttt{// Primary opcode bits})
  \item \texttt{instruction\_class}: 4-bit major instruction class selector (\texttt{bits 48:51}) partitioning the 64-bit Valhall opcode space (e.g., \texttt{FMA}, \texttt{CVT}, \texttt{SFU}, \texttt{LS}, \texttt{TEX}, \texttt{BRANCH}). \emph{Example:} {\scriptsize\ttfamily \mnem{FADD.f32} \dst{r0}, \src{r1}, \src{r2}} (\textbf{SPIR-V:} \texttt{N/A (encoding class)}; \textbf{GLSL:} \texttt{// Hardware execution pipe / decode class})
  \item \texttt{fau\_page\_index}: Selects the FAU (Fast Access Uniform) page (\texttt{0..3}) when a source operand references hardware small constants, thread IDs, or lane constants. \emph{Example:} {\scriptsize\ttfamily \mnem{FADD.f32} \dst{r0}, \src{r1}, \src{k1f}} (\textbf{SPIR-V:} \texttt{OpConstant / BuiltIn}; \textbf{GLSL:} \texttt{float y = x + 1.0;})
  \item \hyperref[enum:flow_control_m]{\texttt{flow}}: 4-bit embedded flow-control \& scoreboard wait modifier (\texttt{bits 59:62}: \texttt{none}, \texttt{wait0..wait0126}, \texttt{wait}, \texttt{reconverge}, \texttt{discard}, \texttt{end}). \emph{Example:} {\scriptsize\ttfamily \mnem{FADD.f32}\imod{.wait0.reconverge} \dst{r0}, \src{r1}, \src{r2}} (\textbf{SPIR-V:} \texttt{OpControlBarrier / OpReturn}; \textbf{GLSL:} \texttt{// Scoreboard wait + warp reconvergence})
\end{tightitemize}
\clearpage

\section{Floating-Point ALU \& Canonicalization (20 mnemonics)}
\subsection*{FABS}
\addcontentsline{toc}{subsection}{FABS}
Computes the floating-point absolute value of \texttt{src0} by clearing its IEEE-754 sign bit (\texttt{dest = \textbar{}src0\textbar{}}), then applies optional \texttt{round}, \texttt{clamp}, and \texttt{dest\_width\_narrow} modifiers. Encoded as an alias of \texttt{FADD(abs(src0), -0.0)}.
\par\textbf{Notes:}
\begin{tightitemize}
  \item Lowered in \texttt{kraid/nir.rs:766-775} as \texttt{OpFAdd(-0.0, src0.fabs())}.
  \item Because it is an FADD alias, adding \texttt{-0.0} canonicalizes sNaN-\textgreater{}qNaN and applies FTZ if global FTZ is enabled.
  \item \textbf{Source references:} \texttt{kraid/nir.rs:766-775}, \texttt{isa-v9-v14.xml}.
\end{tightitemize}
\codebox{\mnem{FABS}\imod{\{.flow\}}\imod{\{.clamp\}} \dst{dest}\imod{\{.dest\_width\_narrow\}}, \src{src0}\imod{\{.widen0\}}\\\mnem{FABS}\imod{\{.flow\}}\imod{\{.clamp\}} \dst{dest}\imod{\{.dest\_width\}}, \src{src0}\imod{\{.swz0\}}}
\textbf{Operands:}
\begin{tightitemize}\footnotesize
  \item \texttt{\dst{dest}} (\texttt{dest\_register\_index}, \texttt{dest\_width}, \texttt{dest\_width\_narrow}): Destination GPR receiving \texttt{\textbar{}src0\textbar{}} (sign bit cleared).
  \item \texttt{\src{src0}}: Floating-point input.
\end{tightitemize}
\par\textbf{Closest SPIR-V Op:} \texttt{GLSL.std.450 FAbs}
\par\textbf{Pseudo-GLSL (Operational Semantics):}
\codebox{float v = uintBitsToFloat(floatBitsToUint(src0) \& 0x7FFFFFFFu);\\dest = narrowOrMask(dest\_width\_narrow, clamp\_mode(round\_rm(v)));}
\textbf{Execution Context:}
\begin{tightitemize}\footnotesize
  \item \textbf{Compute stage preloads (\hyperref[sec:csf-spd]{\texttt{SPD\_k}}, \texttt{r16..r23}):} Hardware pre-fills \texttt{r55..r57} with \texttt{WorkGroupID.xyz + JOB\_OFFSET} (from \hyperref[sec:csf-regs]{\texttt{COMPUTESR::JOB\_OFFSET\_X/Y/Z}} \texttt{r34..r36}) and \texttt{r60..r62} with \texttt{LocalInvocationID.xyz} or \texttt{GlobalInvocationID.xyz} (from \texttt{WG\_SIZE} \texttt{r33} \& \texttt{JOB\_SIZE} \texttt{r37..r39}, selected by \texttt{ShaderProgram::preload}) so compute shaders can derive thread/element indices without memory loads.
  \item \textbf{Vertex stage preloads (\hyperref[sec:csf-spd]{\texttt{SPD\_0}/\texttt{SPD\_1}}, \texttt{r16..r19}):} The hardware Index Fetcher reads \hyperref[sec:csf-regs]{\texttt{IDVSSR::INDEX\_BUFFER}} (\texttt{r54:r55}) and pre-fills \texttt{r59, r60} with \texttt{VertexIndex} (\texttt{+ VERTEX\_OFFSET + GLOBAL\_ATTR\_OFF}) and \texttt{r61} with \texttt{InstanceIndex} (\texttt{+ INSTANCE\_OFFSET}) so Position and Varying shaders can pull vertex attributes directly.
  \item \textbf{Fragment stage preloads \& FP modes (\hyperref[sec:csf-spd]{\texttt{SPD\_2}}, \texttt{r20:r21}):} Hardware pre-fills \texttt{r59} with \texttt{floatBitsToUint(FragCoord.z)} and \texttt{r60} with the MSAA sample coverage mask (\texttt{gl\_SampleMaskIn[0]}), while \texttt{ShaderProgram} word~0 latches 32/64 GPR allocation and warp-wide \texttt{flush\_to\_zero\_mode}, \texttt{suppress\_na\_n}, and \texttt{suppress\_inf}.
  \item \textbf{Inline FAU \& Special Operands (\hyperref[sec:csf-tsd-fau]{\texttt{FAU\_k}}, \texttt{r8..r15}):} Any 8-bit source with \texttt{mode=2} reads zero-latency uniform \texttt{u0..u31} from \texttt{FAU\_k} (\texttt{slots64 = bits[63:56]}); \texttt{mode=3} reads \texttt{SmallConstant} (Table~1.2) or \texttt{FauSpecial} (\texttt{pc}, \texttt{lane\_id}, \texttt{lane\_width}, \texttt{tls\_ptr}, \texttt{wls\_ptr}, \texttt{blend\_descriptor\_0..7}).
\end{tightitemize}
\par\textbf{Worked Examples:}
\par\textbf{Example 1 (Clear negative sign bit):}
\codebox{\mnem{FABS.f32} \dst{r0}, \src{r1}}
With \texttt{r1 = 0xBF800000} (\texttt{-1.0f}): clears bit 31 -\textgreater{} \texttt{r0 = 0x3F800000} (\texttt{+1.0f}).
\par\textbf{Example 2 (Negative zero \& negative NaN):}
\codebox{\mnem{FABS.f32} \dst{r0}, \src{r1}}
With \texttt{r1 = 0x80000000} (\texttt{-0.0f}) -\textgreater{} \texttt{r0 = 0x00000000} (\texttt{+0.0f}); with \texttt{r1 = 0xFFC00001} (\texttt{-qNaN}) -\textgreater{} clears sign bit.

\begin{center}\footnotesize
\begin{tabularx}{\linewidth}{@{}lllX@{}}
\toprule
Variant & Arch & Class & Opcode Fields \\
\midrule
\texttt{f32} & v9--v14 & \texttt{R2} & \texttt{opcode2=ADD} \\
\texttt{v2f16} & v9--v14 & \texttt{R2} & \texttt{opcode2=ADD} \\
\bottomrule
\end{tabularx}
\end{center}
\begin{center}\footnotesize
\begin{tabularx}{0.88\linewidth}{@{}l l X@{}}
\toprule
Field Name & Bits & Type / Fixed Value \\
\midrule
\texttt{src0} & \texttt{0:7} & \texttt{SourceEncoding} \\
\texttt{src1} & \texttt{8:15} & \texttt{exact=192} \\
\texttt{opcode2} & \texttt{16:19} & \texttt{exact=ADD} \\
\textit{Reserved} & \texttt{20:25} & \texttt{zero} \\
\hyperref[enum:widen_m]{\texttt{widen1}} & \texttt{26:27} & \texttt{exact=none} \\
\hyperref[enum:widen_m]{\texttt{widen0}} & \texttt{28:29} & \texttt{widen\_m (def none)} \\
\hyperref[enum:round_m]{\texttt{round}} & \texttt{30:31} & \texttt{exact=none} \\
\hyperref[enum:clamp_m]{\texttt{clamp}} & \texttt{32:33} & \texttt{clamp\_m (def none)} \\
\hyperref[enum:fadd_sticky_m]{\texttt{sticky}} & \texttt{34} & \texttt{exact=none} \\
\textit{Reserved} & \texttt{35} & \texttt{zero} \\
\hyperref[enum:neg_m]{\texttt{neg1}} & \texttt{36} & \texttt{exact=neg} \\
\hyperref[enum:abs_m]{\texttt{abs1}} & \texttt{37} & \texttt{exact=none} \\
\hyperref[enum:neg_m]{\texttt{neg0}} & \texttt{38} & \texttt{exact=none} \\
\hyperref[enum:abs_m]{\texttt{abs0}} & \texttt{39} & \texttt{exact=abs} \\
\texttt{dest\_register\_index} & \texttt{40:45} & \texttt{uint} \\
\hyperref[enum:dest_width_narrow_m]{\texttt{dest\_width\_narrow}} & \texttt{46:47} & \texttt{dest\_width\_narrow\_m} \\
\texttt{datatype} & \texttt{48:51} & \texttt{exact=f32} \\
\texttt{instruction\_class} & \texttt{52:56} & \texttt{exact=R2} \\
\texttt{fau\_page\_index} & \texttt{57:58} & \texttt{uint} \\
\hyperref[enum:flow_control_m]{\texttt{flow}} & \texttt{59:62} & \texttt{flow\_control\_m} \\
\textit{Reserved} & \texttt{63} & \texttt{zero} \\
\bottomrule
\end{tabularx}
\end{center}
\textbf{Fields Explained:}
\begin{tightitemize}\footnotesize
  \item \texttt{src0}: Primary source operand (\texttt{r0..r63}, \texttt{u0..u63} uniform, \texttt{k*} FAU constant, or \texttt{t0..t2} temporary; \texttt{\textasciicircum{}} marks discard/last-use). \emph{Example:} {\scriptsize\ttfamily \mnem{FADD.f32} \dst{r0}, \src{r1}, \src{r2}} (\textbf{SPIR-V:} \texttt{OpFAdd (Operand 0)}; \textbf{GLSL:} \texttt{float c = a + b;})
  \item \texttt{src1}: Secondary source operand (register, uniform, or FAU constant; in shifts, supplies the shift count byte). \emph{Example:} {\scriptsize\ttfamily \mnem{LSHIFT.i32} \dst{r0}, \src{r1}, \src{r2}\imod{.b0}} (\textbf{SPIR-V:} \texttt{OpShiftLeftLogical (Shift)}; \textbf{GLSL:} \texttt{uint y = a \textless{}\textless{} s;})
  \item \texttt{opcode2}: Secondary sub-opcode further specializing instructions sharing \texttt{instruction\_class} and \texttt{opcode1} (e.g., conversions, SFU tables, moves). \emph{Example:} {\scriptsize\ttfamily \mnem{F32\_TO\_S32}\imod{.round\_zero} \dst{r0}, \src{r1}} (\textbf{SPIR-V:} \texttt{OpConvertFToS}; \textbf{GLSL:} \texttt{int y = int(x);})
  \item \hyperref[enum:widen_m]{\texttt{widen1}}: Stage A widening extractor on \texttt{src1} (promotes narrow byte/half/word lane before Stage B). \emph{Example:} {\scriptsize\ttfamily \mnem{IMUL.u64} \dst{[r0:r1]}, \src{[r2:r3]}\imod{.w0}, \src{[r4:r5]}\imod{.w0}} (\textbf{SPIR-V:} \texttt{OpUConvert / OpUMulExtended}; \textbf{GLSL:} \texttt{uint64\_t p = uint64\_t(a) * uint64\_t(b);})
  \item \hyperref[enum:widen_m]{\texttt{widen0}}: Stage A widening extractor on \texttt{src0} (promotes FP16 \texttt{h0}/\texttt{h1} to FP32, or sign/zero-extends \texttt{b0..b3}/\texttt{h0..h1}/\texttt{w0..w1} to 32/64-bit before \texttt{abs}/\texttt{neg}). \emph{Example:} {\scriptsize\ttfamily \mnem{IADD.s32} \dst{r0}, \src{r1}\imod{.h0}, \src{r2}} (\textbf{SPIR-V:} \texttt{OpSConvert / OpFConvert}; \textbf{GLSL:} \texttt{int c = int(int16\_t(a)) + b;})
  \item \hyperref[enum:round_m]{\texttt{round}}: Rounding mode (\texttt{none}/RNE, \texttt{round\_up}, \texttt{round\_down}, \texttt{round\_zero}, \texttt{round\_na}) applied in FP ALU, conversions, \texttt{HADD}, \texttt{LDEXP}, and Stage C narrowing. \emph{Example:} {\scriptsize\ttfamily \mnem{FROUND.f32}\imod{.round\_down} \dst{r0}, \src{r1}} (\textbf{SPIR-V:} \texttt{GLSL.std.450 Floor / Ceil / Trunc}; \textbf{GLSL:} \texttt{float y = floor(x);})
  \item \hyperref[enum:clamp_m]{\texttt{clamp}}: Stage B post-op floating-point clamp (\texttt{none}, \texttt{clamp\_0\_inf}, \texttt{clamp\_m1\_1}, \texttt{clamp\_0\_1}) with \texttt{-0.0 \textless{} +0.0} and NaN suppression (\texttt{NaN -\textgreater{} lo}). \emph{Example:} {\scriptsize\ttfamily \mnem{FMA.f32}\imod{.clamp\_0\_1} \dst{r0}, \src{r1}, \src{r2}, \src{r3}} (\textbf{SPIR-V:} \texttt{GLSL.std.450 FClamp}; \textbf{GLSL:} \texttt{float y = clamp(fma(a, b, c), 0.0, 1.0);})
  \item \hyperref[enum:fadd_sticky_m]{\texttt{sticky}}: OR-reduces shifted-out rounding bits into LSB (round-to-odd for double-precision emulation). \emph{Example:} {\scriptsize\ttfamily \mnem{FADD.f32}\imod{.sticky} \dst{r0}, \src{r1}, \src{r2}} (\textbf{SPIR-V:} \texttt{OpFAdd (emulated FP64)}; \textbf{GLSL:} \texttt{double c = a + b; // lowered})
  \item \hyperref[enum:neg_m]{\texttt{neg1}}: Stage A floating-point negation on \texttt{src1} (flips IEEE sign bit after swizzle/widen and \texttt{abs1}). \emph{Example:} {\scriptsize\ttfamily \mnem{FMA.f32} \dst{r0}, \src{r1}, \src{r2}\imod{.neg}, \src{r3}} (\textbf{SPIR-V:} \texttt{OpFNegate}; \textbf{GLSL:} \texttt{float d = fma(a, -b, c);})
  \item \hyperref[enum:abs_m]{\texttt{abs1}}: Stage A floating-point absolute value on \texttt{src1} (clears IEEE sign bit after swizzle/widen, before \texttt{neg1}). \emph{Example:} {\scriptsize\ttfamily \mnem{FMA.f32} \dst{r0}, \src{r1}, \src{r2}\imod{.abs}, \src{r3}} (\textbf{SPIR-V:} \texttt{GLSL.std.450 FAbs}; \textbf{GLSL:} \texttt{float d = fma(a, abs(b), c);})
  \item \hyperref[enum:neg_m]{\texttt{neg0}}: Stage A floating-point negation on \texttt{src0} (flips IEEE sign bit after swizzle/widen and \texttt{abs0}). \emph{Example:} {\scriptsize\ttfamily \mnem{FADD.f32} \dst{r0}, \src{r1}\imod{.neg}, \src{r2}} (\textbf{SPIR-V:} \texttt{OpFNegate}; \textbf{GLSL:} \texttt{float c = -a + b;})
  \item \hyperref[enum:abs_m]{\texttt{abs0}}: Stage A floating-point absolute value on \texttt{src0} (clears IEEE sign bit after swizzle/widen, before \texttt{neg0}). \emph{Example:} {\scriptsize\ttfamily \mnem{FADD.f32} \dst{r0}, \src{r1}\imod{.abs}, \src{r2}} (\textbf{SPIR-V:} \texttt{GLSL.std.450 FAbs}; \textbf{GLSL:} \texttt{float c = abs(a) + b;})
  \item \texttt{dest\_register\_index}: 6-bit destination GPR index (\texttt{r0..r63}) receiving the Stage C writeback result (even-aligned pair \texttt{[rN:rN+1]} for 64-bit ops). \emph{Example:} {\scriptsize\ttfamily \mnem{FMA.f32} \dst{r0}, \src{r1}, \src{r2}, \src{r3}} (\textbf{SPIR-V:} \texttt{SSA Result ID}; \textbf{GLSL:} \texttt{float r0 = fma(r1, r2, r3);})
  \item \hyperref[enum:dest_width_narrow_m]{\texttt{dest\_width\_narrow}}: Destination halfword lane-mask / replication mode (\texttt{h0}, \texttt{h1}, \texttt{h01}, \texttt{none}). \emph{Example:} {\scriptsize\ttfamily \mnem{F32\_TO\_F16} \dst{r0}\imod{.h01}, \src{r1}} (\textbf{SPIR-V:} \texttt{OpFConvert (f32 -\textgreater{} f16)}; \textbf{GLSL:} \texttt{float16\_t h = float16\_t(x);})
  \item \texttt{datatype}: Selects the lane data type (\texttt{f32}, \texttt{f16}, \texttt{v2f16}, \texttt{s32}, \texttt{u32}, \texttt{v2s16}, \texttt{v2u16}, \texttt{v4s8}, \texttt{v4u8}, \texttt{s64}, \texttt{u64}, etc.). \emph{Example:} {\scriptsize\ttfamily \mnem{IADD.v2s16} \dst{r0}, \src{r1}, \src{r2}} (\textbf{SPIR-V:} \texttt{OpTypeInt / OpTypeFloat}; \textbf{GLSL:} \texttt{i16vec2 c = a + b;})
  \item \texttt{instruction\_class}: 4-bit major instruction class selector (\texttt{bits 48:51}) partitioning the 64-bit Valhall opcode space (e.g., \texttt{FMA}, \texttt{CVT}, \texttt{SFU}, \texttt{LS}, \texttt{TEX}, \texttt{BRANCH}). \emph{Example:} {\scriptsize\ttfamily \mnem{FADD.f32} \dst{r0}, \src{r1}, \src{r2}} (\textbf{SPIR-V:} \texttt{N/A (encoding class)}; \textbf{GLSL:} \texttt{// Hardware execution pipe / decode class})
  \item \texttt{fau\_page\_index}: Selects the FAU (Fast Access Uniform) page (\texttt{0..3}) when a source operand references hardware small constants, thread IDs, or lane constants. \emph{Example:} {\scriptsize\ttfamily \mnem{FADD.f32} \dst{r0}, \src{r1}, \src{k1f}} (\textbf{SPIR-V:} \texttt{OpConstant / BuiltIn}; \textbf{GLSL:} \texttt{float y = x + 1.0;})
  \item \hyperref[enum:flow_control_m]{\texttt{flow}}: 4-bit embedded flow-control \& scoreboard wait modifier (\texttt{bits 59:62}: \texttt{none}, \texttt{wait0..wait0126}, \texttt{wait}, \texttt{reconverge}, \texttt{discard}, \texttt{end}). \emph{Example:} {\scriptsize\ttfamily \mnem{FADD.f32}\imod{.wait0.reconverge} \dst{r0}, \src{r1}, \src{r2}} (\textbf{SPIR-V:} \texttt{OpControlBarrier / OpReturn}; \textbf{GLSL:} \texttt{// Scoreboard wait + warp reconvergence})
\end{tightitemize}
\clearpage

\subsection*{FADD}
\addcontentsline{toc}{subsection}{FADD}
Adds floating-point operands \texttt{src0} and \texttt{src1} independently for each SIMD element $i \in \{0, \dots, N-1\}$ ($N=1$ for \texttt{f32}, $N=2$ for packed \texttt{v2f16}), rounding via \texttt{round} and applying optional post-op \texttt{clamp}: \texttt{dest[i] = clamp(round\_rm(src0[i] + src1[i]))}.
\par\textbf{Notes:}
\begin{tightitemize}
  \item Clamp suppresses NaN (\texttt{x.ieee\_max(lo).ieee\_min(hi)}, \texttt{ops.rs:1136-1143}).
  \item \texttt{sticky} (\texttt{fadd\_sticky\_m}) enables IEEE-754 round-to-odd (jamming LSB to 1 on inexact sums; \texttt{\_mali\_sf64\_to\_sf32\_sticky} \texttt{0x2e297a0}).
  \item \texttt{dest\_width\_narrow} narrows f32-\textgreater{}f16 using \texttt{round} (\texttt{dst\_mod\_prop.rs:226-229}, \texttt{builder.rs:50-65}).
  \item \textbf{Source references:} \texttt{kraid/ops.rs:1267-1302}, \texttt{encode\_v9.rs:1271-1315}, \texttt{dst\_mod\_prop.rs:199-236}, \texttt{hw\_tests.rs:1486-1519}.
\end{tightitemize}
\codebox{\mnem{FADD}\imod{\{.flow\}}\imod{\{.round\}}\imod{\{.sticky\}}\imod{\{.clamp\}} \dst{dest}\imod{\{.dest\_width\_narrow\}}, \src{src0}\imod{\{.widen0\}}\imod{\{.abs0\}}\imod{\{.neg0\}}, \src{src1}\imod{\{.widen1\}}\imod{\{.abs1\}}\imod{\{.neg1\}}\\\mnem{FADD}\imod{\{.flow\}}\imod{\{.round\}}\imod{\{.sticky\}}\imod{\{.clamp\}} \dst{dest}\imod{\{.dest\_width\}}, \src{src0}\imod{\{.swz0\}}\imod{\{.abs0\}}\imod{\{.neg0\}}, \src{src1}\imod{\{.swz1\}}\imod{\{.abs1\}}\imod{\{.neg1\}}}
\textbf{Operands:}
\begin{tightitemize}\footnotesize
  \item \texttt{\dst{dest}} (\texttt{dest\_register\_index}, \texttt{dest\_width}, \texttt{dest\_width\_narrow}): Destination GPR receiving \texttt{src0 + src1} (optionally narrowed via \texttt{dest\_width\_narrow}).
  \item \texttt{\src{src0}}: First floating-point addend (supports \texttt{.abs}/\texttt{.neg}/swizzle).
  \item \texttt{\src{src1}}: Second floating-point addend (supports \texttt{.abs}/\texttt{.neg}/swizzle).
\end{tightitemize}
\par\textbf{Closest SPIR-V Op:} \texttt{OpFAdd}
\par\textbf{Pseudo-GLSL (Operational Semantics):}
\codebox{for (int i = 0; i \textless{} N; ++i) \{\\    auto r = clamp\_mode(round\_rm(mod(src0)[i] + mod(src1)[i]));\\    dest[i] = narrowOrMask(dest\_width\_narrow, r);\\\}}
\textbf{Execution Context:}
\begin{tightitemize}\footnotesize
  \item \textbf{Compute stage preloads (\hyperref[sec:csf-spd]{\texttt{SPD\_k}}, \texttt{r16..r23}):} Hardware pre-fills \texttt{r55..r57} with \texttt{WorkGroupID.xyz + JOB\_OFFSET} (from \hyperref[sec:csf-regs]{\texttt{COMPUTESR::JOB\_OFFSET\_X/Y/Z}} \texttt{r34..r36}) and \texttt{r60..r62} with \texttt{LocalInvocationID.xyz} or \texttt{GlobalInvocationID.xyz} (from \texttt{WG\_SIZE} \texttt{r33} \& \texttt{JOB\_SIZE} \texttt{r37..r39}, selected by \texttt{ShaderProgram::preload}) so compute shaders can derive thread/element indices without memory loads.
  \item \textbf{Vertex stage preloads (\hyperref[sec:csf-spd]{\texttt{SPD\_0}/\texttt{SPD\_1}}, \texttt{r16..r19}):} The hardware Index Fetcher reads \hyperref[sec:csf-regs]{\texttt{IDVSSR::INDEX\_BUFFER}} (\texttt{r54:r55}) and pre-fills \texttt{r59, r60} with \texttt{VertexIndex} (\texttt{+ VERTEX\_OFFSET + GLOBAL\_ATTR\_OFF}) and \texttt{r61} with \texttt{InstanceIndex} (\texttt{+ INSTANCE\_OFFSET}) so Position and Varying shaders can pull vertex attributes directly.
  \item \textbf{Fragment stage preloads \& FP modes (\hyperref[sec:csf-spd]{\texttt{SPD\_2}}, \texttt{r20:r21}):} Hardware pre-fills \texttt{r59} with \texttt{floatBitsToUint(FragCoord.z)} and \texttt{r60} with the MSAA sample coverage mask (\texttt{gl\_SampleMaskIn[0]}), while \texttt{ShaderProgram} word~0 latches 32/64 GPR allocation and warp-wide \texttt{flush\_to\_zero\_mode}, \texttt{suppress\_na\_n}, and \texttt{suppress\_inf}.
  \item \textbf{Inline FAU \& Special Operands (\hyperref[sec:csf-tsd-fau]{\texttt{FAU\_k}}, \texttt{r8..r15}):} Any 8-bit source with \texttt{mode=2} reads zero-latency uniform \texttt{u0..u31} from \texttt{FAU\_k} (\texttt{slots64 = bits[63:56]}); \texttt{mode=3} reads \texttt{SmallConstant} (Table~1.2) or \texttt{FauSpecial} (\texttt{pc}, \texttt{lane\_id}, \texttt{lane\_width}, \texttt{tls\_ptr}, \texttt{wls\_ptr}, \texttt{blend\_descriptor\_0..7}).
\end{tightitemize}
\par\textbf{Worked Examples:}
\par\textbf{Example 1 (Clamped FP32 add with fused Stage C FP16 narrowing):}
\codebox{\mnem{FADD.f32}\imod{.clamp\_0\_1} \dst{r0}\imod{.h1}, \src{r1}, \src{r2}}
With \texttt{r1 = 1.0f, r2 = 0.5f, r0\_old = 0x0000BEEF}: Stage B clamps \texttt{1.5f} to \texttt{1.0f}; Stage C \texttt{dest\_width\_narrow=h1} narrows to \texttt{0x3C00} in \texttt{r0[31:16]} -\textgreater{} \texttt{r0 = 0x3C00BEEF}.
\par\textbf{Example 2 (FP16 widening source + rounding mode):}
\codebox{\mnem{FADD.f32}\imod{.round\_down} \dst{r0}, \src{r1}\imod{.h0.neg}, \src{r2}}
With \texttt{r1.h0 = 2.0h} (\texttt{0x4000}), \texttt{r2 = 5.0f}: Stage A widens \texttt{2.0h -\textgreater{} 2.0f} then negates to \texttt{-2.0f}; Stage B computes \texttt{5.0f + (-2.0f) = 3.0f} under RTN.
\par\textbf{Mined SPIR-V $\to$ Valhall ISA Example (Clamping a varying channel to [0, 1] using FADD.clamp\_0\_1 --- Fragment shader, Mali-G710 (v10)):}
\codebox{\textbf{; SPIR-V (input)}\\\%c = OpExtInst \%float \%glsl FClamp \%alpha \%float\_0 \%float\_1\\\%out = OpFMul \%float \%c \%fade\\[3pt]\textbf{; Valhall ISA (malioc output)}\\FADD.f32.wait1.clamp\_0\_1 r1, r7\textasciicircum{}, k0.neg\\FMUL.f32.wait.clamp\_0\_1 r7, r1\textasciicircum{}, r0\textasciicircum{}}
\texttt{FClamp(x, 0.0, 1.0)} lowers to \texttt{FADD.f32.clamp\_0\_1 r1, r7\textasciicircum{}, k0.neg} (adding \texttt{-0.0} preserves the sign of \texttt{0.0} while applying the free hardware \texttt{[0, 1]} output clamp).
\begin{center}\footnotesize
\begin{tabularx}{\linewidth}{@{}lllX@{}}
\toprule
Variant & Arch & Class & Opcode Fields \\
\midrule
\texttt{f32} & v9--v14 & \texttt{R2} & \texttt{opcode2=ADD} \\
\texttt{v2f16} & v9--v14 & \texttt{R2} & \texttt{opcode2=ADD} \\
\bottomrule
\end{tabularx}
\end{center}
\begin{center}\footnotesize
\begin{tabularx}{0.88\linewidth}{@{}l l X@{}}
\toprule
Field Name & Bits & Type / Fixed Value \\
\midrule
\texttt{src0} & \texttt{0:7} & \texttt{SourceEncoding} \\
\texttt{src1} & \texttt{8:15} & \texttt{SourceEncoding} \\
\texttt{opcode2} & \texttt{16:19} & \texttt{exact=ADD} \\
\textit{Reserved} & \texttt{20:25} & \texttt{zero} \\
\hyperref[enum:widen_m]{\texttt{widen1}} & \texttt{26:27} & \texttt{widen\_m (def none)} \\
\hyperref[enum:widen_m]{\texttt{widen0}} & \texttt{28:29} & \texttt{widen\_m (def none)} \\
\hyperref[enum:round_m]{\texttt{round}} & \texttt{30:31} & \texttt{round\_m (def none)} \\
\hyperref[enum:clamp_m]{\texttt{clamp}} & \texttt{32:33} & \texttt{clamp\_m (def none)} \\
\hyperref[enum:fadd_sticky_m]{\texttt{sticky}} & \texttt{34} & \texttt{fadd\_sticky\_m (def none)} \\
\textit{Reserved} & \texttt{35} & \texttt{zero} \\
\hyperref[enum:neg_m]{\texttt{neg1}} & \texttt{36} & \texttt{neg\_m (def none)} \\
\hyperref[enum:abs_m]{\texttt{abs1}} & \texttt{37} & \texttt{abs\_m (def none)} \\
\hyperref[enum:neg_m]{\texttt{neg0}} & \texttt{38} & \texttt{neg\_m (def none)} \\
\hyperref[enum:abs_m]{\texttt{abs0}} & \texttt{39} & \texttt{abs\_m (def none)} \\
\texttt{dest\_register\_index} & \texttt{40:45} & \texttt{uint} \\
\hyperref[enum:dest_width_narrow_m]{\texttt{dest\_width\_narrow}} & \texttt{46:47} & \texttt{dest\_width\_narrow\_m} \\
\texttt{datatype} & \texttt{48:51} & \texttt{exact=f32} \\
\texttt{instruction\_class} & \texttt{52:56} & \texttt{exact=R2} \\
\texttt{fau\_page\_index} & \texttt{57:58} & \texttt{uint} \\
\hyperref[enum:flow_control_m]{\texttt{flow}} & \texttt{59:62} & \texttt{flow\_control\_m} \\
\textit{Reserved} & \texttt{63} & \texttt{zero} \\
\bottomrule
\end{tabularx}
\end{center}
\textbf{Fields Explained:}
\begin{tightitemize}\footnotesize
  \item \texttt{src0}: Primary source operand (\texttt{r0..r63}, \texttt{u0..u63} uniform, \texttt{k*} FAU constant, or \texttt{t0..t2} temporary; \texttt{\textasciicircum{}} marks discard/last-use). \emph{Example:} {\scriptsize\ttfamily \mnem{FADD.f32} \dst{r0}, \src{r1}, \src{r2}} (\textbf{SPIR-V:} \texttt{OpFAdd (Operand 0)}; \textbf{GLSL:} \texttt{float c = a + b;})
  \item \texttt{src1}: Secondary source operand (register, uniform, or FAU constant; in shifts, supplies the shift count byte). \emph{Example:} {\scriptsize\ttfamily \mnem{LSHIFT.i32} \dst{r0}, \src{r1}, \src{r2}\imod{.b0}} (\textbf{SPIR-V:} \texttt{OpShiftLeftLogical (Shift)}; \textbf{GLSL:} \texttt{uint y = a \textless{}\textless{} s;})
  \item \texttt{opcode2}: Secondary sub-opcode further specializing instructions sharing \texttt{instruction\_class} and \texttt{opcode1} (e.g., conversions, SFU tables, moves). \emph{Example:} {\scriptsize\ttfamily \mnem{F32\_TO\_S32}\imod{.round\_zero} \dst{r0}, \src{r1}} (\textbf{SPIR-V:} \texttt{OpConvertFToS}; \textbf{GLSL:} \texttt{int y = int(x);})
  \item \hyperref[enum:widen_m]{\texttt{widen1}}: Stage A widening extractor on \texttt{src1} (promotes narrow byte/half/word lane before Stage B). \emph{Example:} {\scriptsize\ttfamily \mnem{IMUL.u64} \dst{[r0:r1]}, \src{[r2:r3]}\imod{.w0}, \src{[r4:r5]}\imod{.w0}} (\textbf{SPIR-V:} \texttt{OpUConvert / OpUMulExtended}; \textbf{GLSL:} \texttt{uint64\_t p = uint64\_t(a) * uint64\_t(b);})
  \item \hyperref[enum:widen_m]{\texttt{widen0}}: Stage A widening extractor on \texttt{src0} (promotes FP16 \texttt{h0}/\texttt{h1} to FP32, or sign/zero-extends \texttt{b0..b3}/\texttt{h0..h1}/\texttt{w0..w1} to 32/64-bit before \texttt{abs}/\texttt{neg}). \emph{Example:} {\scriptsize\ttfamily \mnem{IADD.s32} \dst{r0}, \src{r1}\imod{.h0}, \src{r2}} (\textbf{SPIR-V:} \texttt{OpSConvert / OpFConvert}; \textbf{GLSL:} \texttt{int c = int(int16\_t(a)) + b;})
  \item \hyperref[enum:round_m]{\texttt{round}}: Rounding mode (\texttt{none}/RNE, \texttt{round\_up}, \texttt{round\_down}, \texttt{round\_zero}, \texttt{round\_na}) applied in FP ALU, conversions, \texttt{HADD}, \texttt{LDEXP}, and Stage C narrowing. \emph{Example:} {\scriptsize\ttfamily \mnem{FROUND.f32}\imod{.round\_down} \dst{r0}, \src{r1}} (\textbf{SPIR-V:} \texttt{GLSL.std.450 Floor / Ceil / Trunc}; \textbf{GLSL:} \texttt{float y = floor(x);})
  \item \hyperref[enum:clamp_m]{\texttt{clamp}}: Stage B post-op floating-point clamp (\texttt{none}, \texttt{clamp\_0\_inf}, \texttt{clamp\_m1\_1}, \texttt{clamp\_0\_1}) with \texttt{-0.0 \textless{} +0.0} and NaN suppression (\texttt{NaN -\textgreater{} lo}). \emph{Example:} {\scriptsize\ttfamily \mnem{FMA.f32}\imod{.clamp\_0\_1} \dst{r0}, \src{r1}, \src{r2}, \src{r3}} (\textbf{SPIR-V:} \texttt{GLSL.std.450 FClamp}; \textbf{GLSL:} \texttt{float y = clamp(fma(a, b, c), 0.0, 1.0);})
  \item \hyperref[enum:fadd_sticky_m]{\texttt{sticky}}: OR-reduces shifted-out rounding bits into LSB (round-to-odd for double-precision emulation). \emph{Example:} {\scriptsize\ttfamily \mnem{FADD.f32}\imod{.sticky} \dst{r0}, \src{r1}, \src{r2}} (\textbf{SPIR-V:} \texttt{OpFAdd (emulated FP64)}; \textbf{GLSL:} \texttt{double c = a + b; // lowered})
  \item \hyperref[enum:neg_m]{\texttt{neg1}}: Stage A floating-point negation on \texttt{src1} (flips IEEE sign bit after swizzle/widen and \texttt{abs1}). \emph{Example:} {\scriptsize\ttfamily \mnem{FMA.f32} \dst{r0}, \src{r1}, \src{r2}\imod{.neg}, \src{r3}} (\textbf{SPIR-V:} \texttt{OpFNegate}; \textbf{GLSL:} \texttt{float d = fma(a, -b, c);})
  \item \hyperref[enum:abs_m]{\texttt{abs1}}: Stage A floating-point absolute value on \texttt{src1} (clears IEEE sign bit after swizzle/widen, before \texttt{neg1}). \emph{Example:} {\scriptsize\ttfamily \mnem{FMA.f32} \dst{r0}, \src{r1}, \src{r2}\imod{.abs}, \src{r3}} (\textbf{SPIR-V:} \texttt{GLSL.std.450 FAbs}; \textbf{GLSL:} \texttt{float d = fma(a, abs(b), c);})
  \item \hyperref[enum:neg_m]{\texttt{neg0}}: Stage A floating-point negation on \texttt{src0} (flips IEEE sign bit after swizzle/widen and \texttt{abs0}). \emph{Example:} {\scriptsize\ttfamily \mnem{FADD.f32} \dst{r0}, \src{r1}\imod{.neg}, \src{r2}} (\textbf{SPIR-V:} \texttt{OpFNegate}; \textbf{GLSL:} \texttt{float c = -a + b;})
  \item \hyperref[enum:abs_m]{\texttt{abs0}}: Stage A floating-point absolute value on \texttt{src0} (clears IEEE sign bit after swizzle/widen, before \texttt{neg0}). \emph{Example:} {\scriptsize\ttfamily \mnem{FADD.f32} \dst{r0}, \src{r1}\imod{.abs}, \src{r2}} (\textbf{SPIR-V:} \texttt{GLSL.std.450 FAbs}; \textbf{GLSL:} \texttt{float c = abs(a) + b;})
  \item \texttt{dest\_register\_index}: 6-bit destination GPR index (\texttt{r0..r63}) receiving the Stage C writeback result (even-aligned pair \texttt{[rN:rN+1]} for 64-bit ops). \emph{Example:} {\scriptsize\ttfamily \mnem{FMA.f32} \dst{r0}, \src{r1}, \src{r2}, \src{r3}} (\textbf{SPIR-V:} \texttt{SSA Result ID}; \textbf{GLSL:} \texttt{float r0 = fma(r1, r2, r3);})
  \item \hyperref[enum:dest_width_narrow_m]{\texttt{dest\_width\_narrow}}: Destination halfword lane-mask / replication mode (\texttt{h0}, \texttt{h1}, \texttt{h01}, \texttt{none}). \emph{Example:} {\scriptsize\ttfamily \mnem{F32\_TO\_F16} \dst{r0}\imod{.h01}, \src{r1}} (\textbf{SPIR-V:} \texttt{OpFConvert (f32 -\textgreater{} f16)}; \textbf{GLSL:} \texttt{float16\_t h = float16\_t(x);})
  \item \texttt{datatype}: Selects the lane data type (\texttt{f32}, \texttt{f16}, \texttt{v2f16}, \texttt{s32}, \texttt{u32}, \texttt{v2s16}, \texttt{v2u16}, \texttt{v4s8}, \texttt{v4u8}, \texttt{s64}, \texttt{u64}, etc.). \emph{Example:} {\scriptsize\ttfamily \mnem{IADD.v2s16} \dst{r0}, \src{r1}, \src{r2}} (\textbf{SPIR-V:} \texttt{OpTypeInt / OpTypeFloat}; \textbf{GLSL:} \texttt{i16vec2 c = a + b;})
  \item \texttt{instruction\_class}: 4-bit major instruction class selector (\texttt{bits 48:51}) partitioning the 64-bit Valhall opcode space (e.g., \texttt{FMA}, \texttt{CVT}, \texttt{SFU}, \texttt{LS}, \texttt{TEX}, \texttt{BRANCH}). \emph{Example:} {\scriptsize\ttfamily \mnem{FADD.f32} \dst{r0}, \src{r1}, \src{r2}} (\textbf{SPIR-V:} \texttt{N/A (encoding class)}; \textbf{GLSL:} \texttt{// Hardware execution pipe / decode class})
  \item \texttt{fau\_page\_index}: Selects the FAU (Fast Access Uniform) page (\texttt{0..3}) when a source operand references hardware small constants, thread IDs, or lane constants. \emph{Example:} {\scriptsize\ttfamily \mnem{FADD.f32} \dst{r0}, \src{r1}, \src{k1f}} (\textbf{SPIR-V:} \texttt{OpConstant / BuiltIn}; \textbf{GLSL:} \texttt{float y = x + 1.0;})
  \item \hyperref[enum:flow_control_m]{\texttt{flow}}: 4-bit embedded flow-control \& scoreboard wait modifier (\texttt{bits 59:62}: \texttt{none}, \texttt{wait0..wait0126}, \texttt{wait}, \texttt{reconverge}, \texttt{discard}, \texttt{end}). \emph{Example:} {\scriptsize\ttfamily \mnem{FADD.f32}\imod{.wait0.reconverge} \dst{r0}, \src{r1}, \src{r2}} (\textbf{SPIR-V:} \texttt{OpControlBarrier / OpReturn}; \textbf{GLSL:} \texttt{// Scoreboard wait + warp reconvergence})
\end{tightitemize}
\clearpage

\subsection*{FADD\_IMM}
\addcontentsline{toc}{subsection}{FADD\_IMM}
Adds the inline 32-bit floating-point immediate \texttt{imm1w} to \texttt{src0} under IEEE-754 RNE for each SIMD element $i \in \{0, \dots, N-1\}$ ($N=1$ for \texttt{f32}, $N=2$ for \texttt{v2f16}): \texttt{dest[i] = round\_rne(src0[i] + imm1w[i])}.
\par\textbf{Notes:}
\begin{tightitemize}
  \item Always RNE, unclamped (\texttt{encode\_v9.rs:1282-1302}).
  \item \textbf{Source references:} \texttt{kraid/encode\_v9.rs:1282-1302}, \texttt{bifrost/valhall/ISA.xml:40-42}.
\end{tightitemize}
\codebox{\mnem{FADD\_IMM}\imod{\{.flow\}} \dst{dest}\imod{\{.dest\_width\}}, \src{src0}, \imm{\#imm1w}}
\textbf{Operands:}
\begin{tightitemize}\footnotesize
  \item \texttt{\dst{dest}} (\texttt{dest\_register\_index}, \texttt{dest\_width}): Destination GPR receiving \texttt{src0 + \#imm1w}.
  \item \texttt{\src{src0}}: Floating-point input.
  \item \texttt{\imm{\#imm1w}} (\texttt{imm1w}): Inline 32-bit IEEE-754 float immediate (or two packed FP16 values for \texttt{v2f16}).
\end{tightitemize}
\par\textbf{Closest SPIR-V Op:} \texttt{OpFAdd (constant)}
\par\textbf{Pseudo-GLSL (Operational Semantics):}
\codebox{for (int i = 0; i \textless{} N; ++i)\\    dest[i] = round\_rne(src0[i] + \#imm1w[i]);}
\textbf{Execution Context:}
\begin{tightitemize}\footnotesize
  \item \textbf{Compute stage preloads (\hyperref[sec:csf-spd]{\texttt{SPD\_k}}, \texttt{r16..r23}):} Hardware pre-fills \texttt{r55..r57} with \texttt{WorkGroupID.xyz + JOB\_OFFSET} (from \hyperref[sec:csf-regs]{\texttt{COMPUTESR::JOB\_OFFSET\_X/Y/Z}} \texttt{r34..r36}) and \texttt{r60..r62} with \texttt{LocalInvocationID.xyz} or \texttt{GlobalInvocationID.xyz} (from \texttt{WG\_SIZE} \texttt{r33} \& \texttt{JOB\_SIZE} \texttt{r37..r39}, selected by \texttt{ShaderProgram::preload}) so compute shaders can derive thread/element indices without memory loads.
  \item \textbf{Vertex stage preloads (\hyperref[sec:csf-spd]{\texttt{SPD\_0}/\texttt{SPD\_1}}, \texttt{r16..r19}):} The hardware Index Fetcher reads \hyperref[sec:csf-regs]{\texttt{IDVSSR::INDEX\_BUFFER}} (\texttt{r54:r55}) and pre-fills \texttt{r59, r60} with \texttt{VertexIndex} (\texttt{+ VERTEX\_OFFSET + GLOBAL\_ATTR\_OFF}) and \texttt{r61} with \texttt{InstanceIndex} (\texttt{+ INSTANCE\_OFFSET}) so Position and Varying shaders can pull vertex attributes directly.
  \item \textbf{Fragment stage preloads \& FP modes (\hyperref[sec:csf-spd]{\texttt{SPD\_2}}, \texttt{r20:r21}):} Hardware pre-fills \texttt{r59} with \texttt{floatBitsToUint(FragCoord.z)} and \texttt{r60} with the MSAA sample coverage mask (\texttt{gl\_SampleMaskIn[0]}), while \texttt{ShaderProgram} word~0 latches 32/64 GPR allocation and warp-wide \texttt{flush\_to\_zero\_mode}, \texttt{suppress\_na\_n}, and \texttt{suppress\_inf}.
  \item \textbf{Inline FAU \& Special Operands (\hyperref[sec:csf-tsd-fau]{\texttt{FAU\_k}}, \texttt{r8..r15}):} Any 8-bit source with \texttt{mode=2} reads zero-latency uniform \texttt{u0..u31} from \texttt{FAU\_k} (\texttt{slots64 = bits[63:56]}); \texttt{mode=3} reads \texttt{SmallConstant} (Table~1.2) or \texttt{FauSpecial} (\texttt{pc}, \texttt{lane\_id}, \texttt{lane\_width}, \texttt{tls\_ptr}, \texttt{wls\_ptr}, \texttt{blend\_descriptor\_0..7}).
\end{tightitemize}
\par\textbf{Worked Examples:}
\par\textbf{Example 1 (Inline FP32 constant addition):}
\codebox{\mnem{FADD\_IMM.f32} \dst{r0}, \src{r1}, \imm{\#0x3FC00000}}
With \texttt{r1 = 2.5f}: adds inline \texttt{1.5f} (\texttt{0x3FC00000}) under RNE -\textgreater{} \texttt{r0 = 4.0f} (\texttt{0x40800000}).
\par\textbf{Example 2 (Signed zero rule (`-0.0 + +0.0`)):}
\codebox{\mnem{FADD\_IMM.f32} \dst{r0}, \src{r1}, \imm{\#0x00000000}}
With \texttt{r1 = -0.0f} (\texttt{0x80000000}) and inline \texttt{+0.0f}: IEEE RNE yields \texttt{+0.0f} (\texttt{0x00000000}), canonicalizing \texttt{-0.0}.
\par\textbf{Mined SPIR-V $\to$ Valhall ISA Example (Adding a non-FAU float epsilon constant inline --- Fragment shader, Mali-G710 (v10)):}
\codebox{\textbf{; SPIR-V (input)}\\\%denom = OpFAdd \%float \%val \%float\_0\_01 ; 0x3c23d70a\\\%inv   = OpFDiv \%float \%float\_1 \%denom\\[3pt]\textbf{; Valhall ISA (malioc output)}\\FADD\_IMM.f32 r63, r61, \#0x3c23d70a\\FRCP.f32 r63, r63\textasciicircum{}}
Because \texttt{0.01f} (\texttt{0x3c23d70a}) is not in the FAU small-constant table, \texttt{malioc} uses \texttt{FADD\_IMM.f32} to embed the IEEE-754 bit pattern directly in the instruction word before \texttt{FRCP.f32}.
\begin{center}\footnotesize
\begin{tabularx}{\linewidth}{@{}lllX@{}}
\toprule
Variant & Arch & Class & Opcode Fields \\
\midrule
\texttt{(none)} & v9--v14 & \texttt{R1I} & \texttt{opcode1=ADD\_IMM\_F32} \\
\texttt{(none)} & v9--v14 & \texttt{R1I} & \texttt{opcode1=ADD\_IMM\_V2F16} \\
\bottomrule
\end{tabularx}
\end{center}
\begin{center}\footnotesize
\begin{tabularx}{0.88\linewidth}{@{}l l X@{}}
\toprule
Field Name & Bits & Type / Fixed Value \\
\midrule
\texttt{src0} & \texttt{0:7} & \texttt{SourceEncoding} \\
\texttt{imm1w} & \texttt{8:39} & \texttt{float} \\
\texttt{dest\_register\_index} & \texttt{40:45} & \texttt{uint} \\
\hyperref[enum:dest_width_m]{\texttt{dest\_width}} & \texttt{46:47} & \texttt{dest\_width\_m} \\
\texttt{opcode1} & \texttt{48:51} & \texttt{exact=ADD\_IMM\_F32} \\
\texttt{instruction\_class} & \texttt{52:56} & \texttt{exact=R1I} \\
\texttt{fau\_page\_index} & \texttt{57:58} & \texttt{uint} \\
\hyperref[enum:flow_control_m]{\texttt{flow}} & \texttt{59:62} & \texttt{flow\_control\_m} \\
\textit{Reserved} & \texttt{63} & \texttt{zero} \\
\bottomrule
\end{tabularx}
\end{center}
\textbf{Fields Explained:}
\begin{tightitemize}\footnotesize
  \item \texttt{src0}: Primary source operand (\texttt{r0..r63}, \texttt{u0..u63} uniform, \texttt{k*} FAU constant, or \texttt{t0..t2} temporary; \texttt{\textasciicircum{}} marks discard/last-use). \emph{Example:} {\scriptsize\ttfamily \mnem{FADD.f32} \dst{r0}, \src{r1}, \src{r2}} (\textbf{SPIR-V:} \texttt{OpFAdd (Operand 0)}; \textbf{GLSL:} \texttt{float c = a + b;})
  \item \texttt{imm1w}: 32-bit inline immediate operand encoded directly in the instruction word (bits 8:39). \emph{Example:} {\scriptsize\ttfamily \mnem{IADD\_IMM.i32} \dst{r0}, \src{r1}, \imm{\#0x10}} (\textbf{SPIR-V:} \texttt{OpIAdd (constant)}; \textbf{GLSL:} \texttt{int y = x + 16;})
  \item \texttt{dest\_register\_index}: 6-bit destination GPR index (\texttt{r0..r63}) receiving the Stage C writeback result (even-aligned pair \texttt{[rN:rN+1]} for 64-bit ops). \emph{Example:} {\scriptsize\ttfamily \mnem{FMA.f32} \dst{r0}, \src{r1}, \src{r2}, \src{r3}} (\textbf{SPIR-V:} \texttt{SSA Result ID}; \textbf{GLSL:} \texttt{float r0 = fma(r1, r2, r3);})
  \item \hyperref[enum:dest_width_m]{\texttt{dest\_width}}: Stage C destination lane write-mask (\texttt{none} = full 32-bit, \texttt{h0} = low 16-bit, \texttt{h1} = high 16-bit without bit-shifting). \emph{Example:} {\scriptsize\ttfamily \mnem{FADD.v2f16} \dst{r0}\imod{.h1}, \src{r1}, \src{r2}} (\textbf{SPIR-V:} \texttt{OpVectorInsertDynamic / WriteMask}; \textbf{GLSL:} \texttt{dst.y = a.y + b.y;})
  \item \texttt{opcode1}: Primary sub-opcode within the instruction class uniquely identifying the operation family. \emph{Example:} {\scriptsize\ttfamily \mnem{IADD.s32} \dst{r0}, \src{r1}, \src{r2}} (\textbf{SPIR-V:} \texttt{OpIAdd (operation selector)}; \textbf{GLSL:} \texttt{// Primary opcode bits})
  \item \texttt{instruction\_class}: 4-bit major instruction class selector (\texttt{bits 48:51}) partitioning the 64-bit Valhall opcode space (e.g., \texttt{FMA}, \texttt{CVT}, \texttt{SFU}, \texttt{LS}, \texttt{TEX}, \texttt{BRANCH}). \emph{Example:} {\scriptsize\ttfamily \mnem{FADD.f32} \dst{r0}, \src{r1}, \src{r2}} (\textbf{SPIR-V:} \texttt{N/A (encoding class)}; \textbf{GLSL:} \texttt{// Hardware execution pipe / decode class})
  \item \texttt{fau\_page\_index}: Selects the FAU (Fast Access Uniform) page (\texttt{0..3}) when a source operand references hardware small constants, thread IDs, or lane constants. \emph{Example:} {\scriptsize\ttfamily \mnem{FADD.f32} \dst{r0}, \src{r1}, \src{k1f}} (\textbf{SPIR-V:} \texttt{OpConstant / BuiltIn}; \textbf{GLSL:} \texttt{float y = x + 1.0;})
  \item \hyperref[enum:flow_control_m]{\texttt{flow}}: 4-bit embedded flow-control \& scoreboard wait modifier (\texttt{bits 59:62}: \texttt{none}, \texttt{wait0..wait0126}, \texttt{wait}, \texttt{reconverge}, \texttt{discard}, \texttt{end}). \emph{Example:} {\scriptsize\ttfamily \mnem{FADD.f32}\imod{.wait0.reconverge} \dst{r0}, \src{r1}, \src{r2}} (\textbf{SPIR-V:} \texttt{OpControlBarrier / OpReturn}; \textbf{GLSL:} \texttt{// Scoreboard wait + warp reconvergence})
\end{tightitemize}
\clearpage

\subsection*{FEXP}
\addcontentsline{toc}{subsection}{FEXP}
Computes base-2 exponential $2^x$. For \texttt{FEXP.f32}, \texttt{src0} supplies a signed Q8.24 fixed-point exponent (\texttt{int(round(x * 2\textasciicircum{}24))}) and \texttt{src1} supplies the original FP32 \texttt{x} for overflow/underflow/NaN propagation; subnormal outputs flush to zero. For \texttt{FEXP.f16} (v12+), computes $2^{\text{src1}}$ directly from the FP16 lane in \texttt{src1} (\texttt{src0} is unused) and supports \texttt{dest\_width\_replicate}.
\par\textbf{Notes:}
\begin{tightitemize}
  \item \texttt{ ULP (}hw\_tests.rs:2297`).
  \item \textbf{Source references:} \texttt{kraid/ops.rs:1589-1635}, \texttt{encode\_v9.rs:1361-1400}, \texttt{builder.rs:353-384}, \texttt{hw\_tests.rs:2295-2345}.
\end{tightitemize}
\codebox{\mnem{FEXP}\imod{\{.flow\}} \dst{dest}\imod{\{.dest\_width\_replicate\}}, \src{src1}\imod{\{.lane1\}}\imod{\{.abs1\}}\imod{\{.neg1\}}\\\mnem{FEXP}\imod{\{.flow\}} \dst{dest}\imod{\{.dest\_width\}}, \src{src0}, \src{src1}\imod{\{.widen1\}}\imod{\{.abs1\}}\imod{\{.neg1\}}}
\textbf{Operands:}
\begin{tightitemize}\footnotesize
  \item \texttt{\dst{dest}} (\texttt{dest\_register\_index}, \texttt{dest\_width}, \texttt{dest\_width\_replicate}): Destination GPR receiving \texttt{2\textasciicircum{}x}.
  \item \texttt{\src{src0}}: Signed Q8.24 fixed-point exponent (\texttt{x * 2\textasciicircum{}24}) for \texttt{FEXP.f32} (unused on \texttt{FEXP.f16}).
  \item \texttt{\src{src1}}: Original FP32 input for overflow/underflow/NaN edge cases on \texttt{FEXP.f32} (or direct FP16 exponent on \texttt{FEXP.f16}).
\end{tightitemize}
\par\textbf{Closest SPIR-V Op:} \texttt{GLSL.std.450 Exp2}
\par\textbf{Pseudo-GLSL (Operational Semantics):}
\codebox{if (is\_f32) dest = ftz\_out(specialCaseFexp(src1, exp2(float(int(src0)) / 16777216.0)));\\else        dest = replicate(dest\_width\_replicate, exp2(lane1(src1))); // v12+ f16}
\textbf{Execution Context:}
\begin{tightitemize}\footnotesize
  \item \textbf{Compute stage preloads (\hyperref[sec:csf-spd]{\texttt{SPD\_k}}, \texttt{r16..r23}):} Hardware pre-fills \texttt{r55..r57} with \texttt{WorkGroupID.xyz + JOB\_OFFSET} (from \hyperref[sec:csf-regs]{\texttt{COMPUTESR::JOB\_OFFSET\_X/Y/Z}} \texttt{r34..r36}) and \texttt{r60..r62} with \texttt{LocalInvocationID.xyz} or \texttt{GlobalInvocationID.xyz} (from \texttt{WG\_SIZE} \texttt{r33} \& \texttt{JOB\_SIZE} \texttt{r37..r39}, selected by \texttt{ShaderProgram::preload}) so compute shaders can derive thread/element indices without memory loads.
  \item \textbf{Vertex stage preloads (\hyperref[sec:csf-spd]{\texttt{SPD\_0}/\texttt{SPD\_1}}, \texttt{r16..r19}):} The hardware Index Fetcher reads \hyperref[sec:csf-regs]{\texttt{IDVSSR::INDEX\_BUFFER}} (\texttt{r54:r55}) and pre-fills \texttt{r59, r60} with \texttt{VertexIndex} (\texttt{+ VERTEX\_OFFSET + GLOBAL\_ATTR\_OFF}) and \texttt{r61} with \texttt{InstanceIndex} (\texttt{+ INSTANCE\_OFFSET}) so Position and Varying shaders can pull vertex attributes directly.
  \item \textbf{Fragment stage preloads \& FP modes (\hyperref[sec:csf-spd]{\texttt{SPD\_2}}, \texttt{r20:r21}):} Hardware pre-fills \texttt{r59} with \texttt{floatBitsToUint(FragCoord.z)} and \texttt{r60} with the MSAA sample coverage mask (\texttt{gl\_SampleMaskIn[0]}), while \texttt{ShaderProgram} word~0 latches 32/64 GPR allocation and warp-wide \texttt{flush\_to\_zero\_mode}, \texttt{suppress\_na\_n}, and \texttt{suppress\_inf}.
  \item \textbf{Inline FAU \& Special Operands (\hyperref[sec:csf-tsd-fau]{\texttt{FAU\_k}}, \texttt{r8..r15}):} Any 8-bit source with \texttt{mode=2} reads zero-latency uniform \texttt{u0..u31} from \texttt{FAU\_k} (\texttt{slots64 = bits[63:56]}); \texttt{mode=3} reads \texttt{SmallConstant} (Table~1.2) or \texttt{FauSpecial} (\texttt{pc}, \texttt{lane\_id}, \texttt{lane\_width}, \texttt{tls\_ptr}, \texttt{wls\_ptr}, \texttt{blend\_descriptor\_0..7}).
\end{tightitemize}
\par\textbf{Worked Examples:}
\par\textbf{Example 1 (Q8.24 fixed-point + FP32 input):}
\codebox{\mnem{FEXP.f32} \dst{r0}, \src{r1}, \src{r2}}
With \texttt{r1 = 0x02000000} (\texttt{+2.0} in Q8.24), \texttt{r2 = 2.0f}: computes \texttt{2\textasciicircum{}2.0 = 4.0f} (\texttt{0x40800000}).
\par\textbf{Example 2 (Negative exponent \& underflow):}
\codebox{\mnem{FEXP.f32} \dst{r0}, \src{r1}, \src{r2}}
With \texttt{r1 = 0xFE000000} (\texttt{-2.0} in Q8.24), \texttt{r2 = -2.0f}: computes \texttt{2\textasciicircum{}-2.0 = 0.25f} (\texttt{0x3E800000}); inputs \texttt{\textless{}= -127.0} flush to \texttt{+0.0f}.
\par\textbf{Mined SPIR-V $\to$ Valhall ISA Example (Lowering GLSL.std.450 Exp2 via Q8.24 fixed-point pre-scaling --- Fragment shader, Mali-G710 (v10)):}
\codebox{\textbf{; SPIR-V (input)}\\\%e = OpExtInst \%float \%glsl Exp2 \%x\\[3pt]\textbf{; Valhall ISA (malioc output)}\\LDEXP.f32 r55, r4, u9.w1          ; scale r4 by 2\textasciicircum{}24\\F32\_TO\_S32 r55, r55\textasciicircum{}              ; convert to signed Q8.24\\FEXP.f32 r55, r55\textasciicircum{}, r4            ; 2\textasciicircum{}(r55 / 2\textasciicircum{}24), range/NaN from r4}
\texttt{FEXP.f32} does not take raw FP32 directly: \texttt{malioc} scales \texttt{x} by $2^{24}$ (\texttt{LDEXP} or \texttt{FMUL\_RSCALE}), converts to signed Q8.24 fixed-point (\texttt{F32\_TO\_S32}), and passes both the fixed-point value (\texttt{src0 = r55\textasciicircum{}}) and original float (\texttt{src1 = r4}, used for overflow/underflow/NaN edge cases) into \texttt{FEXP.f32}.
\begin{center}\footnotesize
\begin{tabularx}{\linewidth}{@{}lllX@{}}
\toprule
Variant & Arch & Class & Opcode Fields \\
\midrule
\texttt{f16} & v9--v14 & \texttt{R2} & \texttt{opcode2=FPOW, opcode3=FPOW} \\
\texttt{f32} & v9--v14 & \texttt{R2} & \texttt{opcode2=FEXP} \\
\bottomrule
\end{tabularx}
\end{center}
\begin{center}\footnotesize
\begin{tabularx}{0.88\linewidth}{@{}l l X@{}}
\toprule
Field Name & Bits & Type / Fixed Value \\
\midrule
\texttt{src0} & \texttt{0:7} & \texttt{exact=220} \\
\texttt{src1} & \texttt{8:15} & \texttt{SourceEncoding} \\
\texttt{opcode2} & \texttt{16:19} & \texttt{exact=FPOW} \\
\textit{Reserved} & \texttt{20:23} & \texttt{zero} \\
\texttt{opcode3} & \texttt{24:25} & \texttt{exact=FPOW} \\
\hyperref[enum:lane16_m]{\texttt{lane1}} & \texttt{26} & \texttt{lane16\_m (def h0)} \\
\textit{Reserved} & \texttt{27} & \texttt{zero} \\
\hyperref[enum:lane16_m]{\texttt{lane0}} & \texttt{28} & \texttt{exact=h1} \\
\textit{Reserved} & \texttt{29:35} & \texttt{zero} \\
\hyperref[enum:neg_m]{\texttt{neg1}} & \texttt{36} & \texttt{neg\_m (def none)} \\
\hyperref[enum:abs_m]{\texttt{abs1}} & \texttt{37} & \texttt{abs\_m (def none)} \\
\hyperref[enum:neg_m]{\texttt{neg0}} & \texttt{38} & \texttt{exact=none} \\
\hyperref[enum:abs_m]{\texttt{abs0}} & \texttt{39} & \texttt{exact=none} \\
\texttt{dest\_register\_index} & \texttt{40:45} & \texttt{uint} \\
\hyperref[enum:dest_width_replicate_m]{\texttt{dest\_width\_replicate}} & \texttt{46:47} & \texttt{dest\_width\_replicate\_m} \\
\texttt{datatype} & \texttt{48:51} & \texttt{exact=f16} \\
\texttt{instruction\_class} & \texttt{52:56} & \texttt{exact=R2} \\
\texttt{fau\_page\_index} & \texttt{57:58} & \texttt{uint} \\
\hyperref[enum:flow_control_m]{\texttt{flow}} & \texttt{59:62} & \texttt{flow\_control\_m} \\
\textit{Reserved} & \texttt{63} & \texttt{zero} \\
\bottomrule
\end{tabularx}
\end{center}
\textbf{Fields Explained:}
\begin{tightitemize}\footnotesize
  \item \texttt{src0}: Primary source operand (\texttt{r0..r63}, \texttt{u0..u63} uniform, \texttt{k*} FAU constant, or \texttt{t0..t2} temporary; \texttt{\textasciicircum{}} marks discard/last-use). \emph{Example:} {\scriptsize\ttfamily \mnem{FADD.f32} \dst{r0}, \src{r1}, \src{r2}} (\textbf{SPIR-V:} \texttt{OpFAdd (Operand 0)}; \textbf{GLSL:} \texttt{float c = a + b;})
  \item \texttt{src1}: Secondary source operand (register, uniform, or FAU constant; in shifts, supplies the shift count byte). \emph{Example:} {\scriptsize\ttfamily \mnem{LSHIFT.i32} \dst{r0}, \src{r1}, \src{r2}\imod{.b0}} (\textbf{SPIR-V:} \texttt{OpShiftLeftLogical (Shift)}; \textbf{GLSL:} \texttt{uint y = a \textless{}\textless{} s;})
  \item \texttt{opcode2}: Secondary sub-opcode further specializing instructions sharing \texttt{instruction\_class} and \texttt{opcode1} (e.g., conversions, SFU tables, moves). \emph{Example:} {\scriptsize\ttfamily \mnem{F32\_TO\_S32}\imod{.round\_zero} \dst{r0}, \src{r1}} (\textbf{SPIR-V:} \texttt{OpConvertFToS}; \textbf{GLSL:} \texttt{int y = int(x);})
  \item \texttt{opcode3}: Tertiary sub-opcode disambiguating specialized multi-op families (\texttt{FPOW\_SC\_*}, SFU lookup tables, \texttt{ICMP\_MULTI}). \emph{Example:} {\scriptsize\ttfamily \mnem{FPOW\_SC\_DET.f32}\imod{.pow} \dst{r0}, \src{r1}, \src{r2}} (\textbf{SPIR-V:} \texttt{GLSL.std.450 Pow}; \textbf{GLSL:} \texttt{float z = pow(x, y);})
  \item \hyperref[enum:lane16_m]{\texttt{lane1}}: Stage A single-lane extractor on \texttt{src1} (\texttt{b0..b3} or \texttt{h0}/\texttt{h1}) for \texttt{MKVEC} packing and narrow binary ops. \emph{Example:} {\scriptsize\ttfamily \mnem{MKVEC.v2i8} \dst{r0}, \src{r1}\imod{.b0}, \src{r2}\imod{.b1}, \src{r3}\imod{.b01}} (\textbf{SPIR-V:} \texttt{OpCompositeConstruct}; \textbf{GLSL:} \texttt{u8vec4 v = u8vec4(a.x, b.y, c.x, c.y);})
  \item \hyperref[enum:lane16_m]{\texttt{lane0}}: Stage A single-lane extractor on \texttt{src0} (\texttt{b0..b3} for 8-bit or \texttt{h0}/\texttt{h1} for 16-bit narrow conversions and \texttt{MKVEC}). \emph{Example:} {\scriptsize\ttfamily \mnem{F16\_TO\_F32} \dst{r0}, \src{r1}\imod{.h1}} (\textbf{SPIR-V:} \texttt{OpCompositeExtract}; \textbf{GLSL:} \texttt{float y = float(x.y);})
  \item \hyperref[enum:neg_m]{\texttt{neg1}}: Stage A floating-point negation on \texttt{src1} (flips IEEE sign bit after swizzle/widen and \texttt{abs1}). \emph{Example:} {\scriptsize\ttfamily \mnem{FMA.f32} \dst{r0}, \src{r1}, \src{r2}\imod{.neg}, \src{r3}} (\textbf{SPIR-V:} \texttt{OpFNegate}; \textbf{GLSL:} \texttt{float d = fma(a, -b, c);})
  \item \hyperref[enum:abs_m]{\texttt{abs1}}: Stage A floating-point absolute value on \texttt{src1} (clears IEEE sign bit after swizzle/widen, before \texttt{neg1}). \emph{Example:} {\scriptsize\ttfamily \mnem{FMA.f32} \dst{r0}, \src{r1}, \src{r2}\imod{.abs}, \src{r3}} (\textbf{SPIR-V:} \texttt{GLSL.std.450 FAbs}; \textbf{GLSL:} \texttt{float d = fma(a, abs(b), c);})
  \item \hyperref[enum:neg_m]{\texttt{neg0}}: Stage A floating-point negation on \texttt{src0} (flips IEEE sign bit after swizzle/widen and \texttt{abs0}). \emph{Example:} {\scriptsize\ttfamily \mnem{FADD.f32} \dst{r0}, \src{r1}\imod{.neg}, \src{r2}} (\textbf{SPIR-V:} \texttt{OpFNegate}; \textbf{GLSL:} \texttt{float c = -a + b;})
  \item \hyperref[enum:abs_m]{\texttt{abs0}}: Stage A floating-point absolute value on \texttt{src0} (clears IEEE sign bit after swizzle/widen, before \texttt{neg0}). \emph{Example:} {\scriptsize\ttfamily \mnem{FADD.f32} \dst{r0}, \src{r1}\imod{.abs}, \src{r2}} (\textbf{SPIR-V:} \texttt{GLSL.std.450 FAbs}; \textbf{GLSL:} \texttt{float c = abs(a) + b;})
  \item \texttt{dest\_register\_index}: 6-bit destination GPR index (\texttt{r0..r63}) receiving the Stage C writeback result (even-aligned pair \texttt{[rN:rN+1]} for 64-bit ops). \emph{Example:} {\scriptsize\ttfamily \mnem{FMA.f32} \dst{r0}, \src{r1}, \src{r2}, \src{r3}} (\textbf{SPIR-V:} \texttt{SSA Result ID}; \textbf{GLSL:} \texttt{float r0 = fma(r1, r2, r3);})
  \item \hyperref[enum:dest_width_replicate_m]{\texttt{dest\_width\_replicate}}: Destination halfword lane-mask / replication mode (\texttt{h0}, \texttt{h1}, \texttt{h01}, \texttt{none}). \emph{Example:} {\scriptsize\ttfamily \mnem{F32\_TO\_F16} \dst{r0}\imod{.h01}, \src{r1}} (\textbf{SPIR-V:} \texttt{OpFConvert (f32 -\textgreater{} f16)}; \textbf{GLSL:} \texttt{float16\_t h = float16\_t(x);})
  \item \texttt{datatype}: Selects the lane data type (\texttt{f32}, \texttt{f16}, \texttt{v2f16}, \texttt{s32}, \texttt{u32}, \texttt{v2s16}, \texttt{v2u16}, \texttt{v4s8}, \texttt{v4u8}, \texttt{s64}, \texttt{u64}, etc.). \emph{Example:} {\scriptsize\ttfamily \mnem{IADD.v2s16} \dst{r0}, \src{r1}, \src{r2}} (\textbf{SPIR-V:} \texttt{OpTypeInt / OpTypeFloat}; \textbf{GLSL:} \texttt{i16vec2 c = a + b;})
  \item \texttt{instruction\_class}: 4-bit major instruction class selector (\texttt{bits 48:51}) partitioning the 64-bit Valhall opcode space (e.g., \texttt{FMA}, \texttt{CVT}, \texttt{SFU}, \texttt{LS}, \texttt{TEX}, \texttt{BRANCH}). \emph{Example:} {\scriptsize\ttfamily \mnem{FADD.f32} \dst{r0}, \src{r1}, \src{r2}} (\textbf{SPIR-V:} \texttt{N/A (encoding class)}; \textbf{GLSL:} \texttt{// Hardware execution pipe / decode class})
  \item \texttt{fau\_page\_index}: Selects the FAU (Fast Access Uniform) page (\texttt{0..3}) when a source operand references hardware small constants, thread IDs, or lane constants. \emph{Example:} {\scriptsize\ttfamily \mnem{FADD.f32} \dst{r0}, \src{r1}, \src{k1f}} (\textbf{SPIR-V:} \texttt{OpConstant / BuiltIn}; \textbf{GLSL:} \texttt{float y = x + 1.0;})
  \item \hyperref[enum:flow_control_m]{\texttt{flow}}: 4-bit embedded flow-control \& scoreboard wait modifier (\texttt{bits 59:62}: \texttt{none}, \texttt{wait0..wait0126}, \texttt{wait}, \texttt{reconverge}, \texttt{discard}, \texttt{end}). \emph{Example:} {\scriptsize\ttfamily \mnem{FADD.f32}\imod{.wait0.reconverge} \dst{r0}, \src{r1}, \src{r2}} (\textbf{SPIR-V:} \texttt{OpControlBarrier / OpReturn}; \textbf{GLSL:} \texttt{// Scoreboard wait + warp reconvergence})
\end{tightitemize}
\clearpage

\subsection*{FEXPF}
\addcontentsline{toc}{subsection}{FEXPF}
Directly computes the single-precision base-2 exponential \texttt{dest = exp2(src0)} on Valhall v14+ without requiring Q8.24 fixed-point pre-scaling in software. Handles standard IEEE-754 overflow to \texttt{+Inf}, underflow to \texttt{+0.0}, and NaN propagation.
\par\textbf{Notes:}
\begin{tightitemize}
  \item Replaces the 3-instruction \texttt{FMA\_RSCALE} + \texttt{F32\_TO\_S32} + \texttt{FEXP.f32} sequence on v14+.
  \item \textbf{Source references:} \texttt{isa-v9-v14.xml}.
\end{tightitemize}
\codebox{\mnem{FEXPF}\imod{\{.flow\}} \dst{dest}\imod{\{.dest\_width\}}, \src{src0}\imod{\{.widen0\}}\imod{\{.abs0\}}\imod{\{.neg0\}}}
\textbf{Operands:}
\begin{tightitemize}\footnotesize
  \item \texttt{\dst{dest}} (\texttt{dest\_register\_index}, \texttt{dest\_width}): 32-bit destination GPR receiving \texttt{exp2(src0)}.
  \item \texttt{\src{src0}}: Direct FP32 exponent input (no Q8.24 pre-scaling).
\end{tightitemize}
\par\textbf{Closest SPIR-V Op:} \texttt{GLSL.std.450 Exp2 (v14+)}
\par\textbf{Pseudo-GLSL (Operational Semantics):}
\codebox{dest = exp2(src0); // direct FP32 exp2 (v14+)}
\textbf{Execution Context:}
\begin{tightitemize}\footnotesize
  \item \textbf{Compute stage preloads (\hyperref[sec:csf-spd]{\texttt{SPD\_k}}, \texttt{r16..r23}):} Hardware pre-fills \texttt{r55..r57} with \texttt{WorkGroupID.xyz + JOB\_OFFSET} (from \hyperref[sec:csf-regs]{\texttt{COMPUTESR::JOB\_OFFSET\_X/Y/Z}} \texttt{r34..r36}) and \texttt{r60..r62} with \texttt{LocalInvocationID.xyz} or \texttt{GlobalInvocationID.xyz} (from \texttt{WG\_SIZE} \texttt{r33} \& \texttt{JOB\_SIZE} \texttt{r37..r39}, selected by \texttt{ShaderProgram::preload}) so compute shaders can derive thread/element indices without memory loads.
  \item \textbf{Vertex stage preloads (\hyperref[sec:csf-spd]{\texttt{SPD\_0}/\texttt{SPD\_1}}, \texttt{r16..r19}):} The hardware Index Fetcher reads \hyperref[sec:csf-regs]{\texttt{IDVSSR::INDEX\_BUFFER}} (\texttt{r54:r55}) and pre-fills \texttt{r59, r60} with \texttt{VertexIndex} (\texttt{+ VERTEX\_OFFSET + GLOBAL\_ATTR\_OFF}) and \texttt{r61} with \texttt{InstanceIndex} (\texttt{+ INSTANCE\_OFFSET}) so Position and Varying shaders can pull vertex attributes directly.
  \item \textbf{Fragment stage preloads \& FP modes (\hyperref[sec:csf-spd]{\texttt{SPD\_2}}, \texttt{r20:r21}):} Hardware pre-fills \texttt{r59} with \texttt{floatBitsToUint(FragCoord.z)} and \texttt{r60} with the MSAA sample coverage mask (\texttt{gl\_SampleMaskIn[0]}), while \texttt{ShaderProgram} word~0 latches 32/64 GPR allocation and warp-wide \texttt{flush\_to\_zero\_mode}, \texttt{suppress\_na\_n}, and \texttt{suppress\_inf}.
  \item \textbf{Inline FAU \& Special Operands (\hyperref[sec:csf-tsd-fau]{\texttt{FAU\_k}}, \texttt{r8..r15}):} Any 8-bit source with \texttt{mode=2} reads zero-latency uniform \texttt{u0..u31} from \texttt{FAU\_k} (\texttt{slots64 = bits[63:56]}); \texttt{mode=3} reads \texttt{SmallConstant} (Table~1.2) or \texttt{FauSpecial} (\texttt{pc}, \texttt{lane\_id}, \texttt{lane\_width}, \texttt{tls\_ptr}, \texttt{wls\_ptr}, \texttt{blend\_descriptor\_0..7}).
\end{tightitemize}
\par\textbf{Worked Examples:}
\par\textbf{Example 1 (Direct FP32 exp2 (v14+)):}
\codebox{\mnem{FEXPF.f32} \dst{r0}, \src{r1}}
With \texttt{r1 = 3.0f}: direct floating-point \texttt{exp2(3.0) = 8.0f} (\texttt{0x41000000}) without needing Q8.24 pre-scaling.
\par\textbf{Example 2 (Special values (`-Inf`, `+Inf`, `NaN`)):}
\codebox{\mnem{FEXPF.f32} \dst{r0}, \src{r1}}
With \texttt{r1 = -Inf} -\textgreater{} \texttt{+0.0f}; with \texttt{r1 = +Inf} (or \texttt{\textgreater{}= 128.0f}) -\textgreater{} \texttt{+Inf}; with \texttt{r1 = NaN} -\textgreater{} \texttt{qNaN}.

\begin{center}\footnotesize
\begin{tabularx}{\linewidth}{@{}lllX@{}}
\toprule
Variant & Arch & Class & Opcode Fields \\
\midrule
\texttt{(none)} & v9--v14 & \texttt{R1} & \texttt{opcode2=FEXPF\_F32, opcode1=TABLE} \\
\bottomrule
\end{tabularx}
\end{center}
\begin{center}\footnotesize
\begin{tabularx}{0.88\linewidth}{@{}l l X@{}}
\toprule
Field Name & Bits & Type / Fixed Value \\
\midrule
\texttt{src0} & \texttt{0:7} & \texttt{SourceEncoding} \\
\textit{Reserved} & \texttt{8:15} & \texttt{zero} \\
\texttt{opcode2} & \texttt{16:20} & \texttt{exact=FEXPF\_F32} \\
\textit{Reserved} & \texttt{21:27} & \texttt{zero} \\
\hyperref[enum:widen_m]{\texttt{widen0}} & \texttt{28:29} & \texttt{widen\_m (def none)} \\
\textit{Reserved} & \texttt{30:37} & \texttt{zero} \\
\hyperref[enum:neg_m]{\texttt{neg0}} & \texttt{38} & \texttt{neg\_m (def none)} \\
\hyperref[enum:abs_m]{\texttt{abs0}} & \texttt{39} & \texttt{abs\_m (def none)} \\
\texttt{dest\_register\_index} & \texttt{40:45} & \texttt{uint} \\
\hyperref[enum:dest_width_m]{\texttt{dest\_width}} & \texttt{46:47} & \texttt{dest\_width\_m} \\
\texttt{opcode1} & \texttt{48:51} & \texttt{exact=TABLE} \\
\texttt{instruction\_class} & \texttt{52:56} & \texttt{exact=R1} \\
\texttt{fau\_page\_index} & \texttt{57:58} & \texttt{uint} \\
\hyperref[enum:flow_control_m]{\texttt{flow}} & \texttt{59:62} & \texttt{flow\_control\_m} \\
\textit{Reserved} & \texttt{63} & \texttt{zero} \\
\bottomrule
\end{tabularx}
\end{center}
\textbf{Fields Explained:}
\begin{tightitemize}\footnotesize
  \item \texttt{src0}: Primary source operand (\texttt{r0..r63}, \texttt{u0..u63} uniform, \texttt{k*} FAU constant, or \texttt{t0..t2} temporary; \texttt{\textasciicircum{}} marks discard/last-use). \emph{Example:} {\scriptsize\ttfamily \mnem{FADD.f32} \dst{r0}, \src{r1}, \src{r2}} (\textbf{SPIR-V:} \texttt{OpFAdd (Operand 0)}; \textbf{GLSL:} \texttt{float c = a + b;})
  \item \texttt{opcode2}: Secondary sub-opcode further specializing instructions sharing \texttt{instruction\_class} and \texttt{opcode1} (e.g., conversions, SFU tables, moves). \emph{Example:} {\scriptsize\ttfamily \mnem{F32\_TO\_S32}\imod{.round\_zero} \dst{r0}, \src{r1}} (\textbf{SPIR-V:} \texttt{OpConvertFToS}; \textbf{GLSL:} \texttt{int y = int(x);})
  \item \hyperref[enum:widen_m]{\texttt{widen0}}: Stage A widening extractor on \texttt{src0} (promotes FP16 \texttt{h0}/\texttt{h1} to FP32, or sign/zero-extends \texttt{b0..b3}/\texttt{h0..h1}/\texttt{w0..w1} to 32/64-bit before \texttt{abs}/\texttt{neg}). \emph{Example:} {\scriptsize\ttfamily \mnem{IADD.s32} \dst{r0}, \src{r1}\imod{.h0}, \src{r2}} (\textbf{SPIR-V:} \texttt{OpSConvert / OpFConvert}; \textbf{GLSL:} \texttt{int c = int(int16\_t(a)) + b;})
  \item \hyperref[enum:neg_m]{\texttt{neg0}}: Stage A floating-point negation on \texttt{src0} (flips IEEE sign bit after swizzle/widen and \texttt{abs0}). \emph{Example:} {\scriptsize\ttfamily \mnem{FADD.f32} \dst{r0}, \src{r1}\imod{.neg}, \src{r2}} (\textbf{SPIR-V:} \texttt{OpFNegate}; \textbf{GLSL:} \texttt{float c = -a + b;})
  \item \hyperref[enum:abs_m]{\texttt{abs0}}: Stage A floating-point absolute value on \texttt{src0} (clears IEEE sign bit after swizzle/widen, before \texttt{neg0}). \emph{Example:} {\scriptsize\ttfamily \mnem{FADD.f32} \dst{r0}, \src{r1}\imod{.abs}, \src{r2}} (\textbf{SPIR-V:} \texttt{GLSL.std.450 FAbs}; \textbf{GLSL:} \texttt{float c = abs(a) + b;})
  \item \texttt{dest\_register\_index}: 6-bit destination GPR index (\texttt{r0..r63}) receiving the Stage C writeback result (even-aligned pair \texttt{[rN:rN+1]} for 64-bit ops). \emph{Example:} {\scriptsize\ttfamily \mnem{FMA.f32} \dst{r0}, \src{r1}, \src{r2}, \src{r3}} (\textbf{SPIR-V:} \texttt{SSA Result ID}; \textbf{GLSL:} \texttt{float r0 = fma(r1, r2, r3);})
  \item \hyperref[enum:dest_width_m]{\texttt{dest\_width}}: Stage C destination lane write-mask (\texttt{none} = full 32-bit, \texttt{h0} = low 16-bit, \texttt{h1} = high 16-bit without bit-shifting). \emph{Example:} {\scriptsize\ttfamily \mnem{FADD.v2f16} \dst{r0}\imod{.h1}, \src{r1}, \src{r2}} (\textbf{SPIR-V:} \texttt{OpVectorInsertDynamic / WriteMask}; \textbf{GLSL:} \texttt{dst.y = a.y + b.y;})
  \item \texttt{opcode1}: Primary sub-opcode within the instruction class uniquely identifying the operation family. \emph{Example:} {\scriptsize\ttfamily \mnem{IADD.s32} \dst{r0}, \src{r1}, \src{r2}} (\textbf{SPIR-V:} \texttt{OpIAdd (operation selector)}; \textbf{GLSL:} \texttt{// Primary opcode bits})
  \item \texttt{instruction\_class}: 4-bit major instruction class selector (\texttt{bits 48:51}) partitioning the 64-bit Valhall opcode space (e.g., \texttt{FMA}, \texttt{CVT}, \texttt{SFU}, \texttt{LS}, \texttt{TEX}, \texttt{BRANCH}). \emph{Example:} {\scriptsize\ttfamily \mnem{FADD.f32} \dst{r0}, \src{r1}, \src{r2}} (\textbf{SPIR-V:} \texttt{N/A (encoding class)}; \textbf{GLSL:} \texttt{// Hardware execution pipe / decode class})
  \item \texttt{fau\_page\_index}: Selects the FAU (Fast Access Uniform) page (\texttt{0..3}) when a source operand references hardware small constants, thread IDs, or lane constants. \emph{Example:} {\scriptsize\ttfamily \mnem{FADD.f32} \dst{r0}, \src{r1}, \src{k1f}} (\textbf{SPIR-V:} \texttt{OpConstant / BuiltIn}; \textbf{GLSL:} \texttt{float y = x + 1.0;})
  \item \hyperref[enum:flow_control_m]{\texttt{flow}}: 4-bit embedded flow-control \& scoreboard wait modifier (\texttt{bits 59:62}: \texttt{none}, \texttt{wait0..wait0126}, \texttt{wait}, \texttt{reconverge}, \texttt{discard}, \texttt{end}). \emph{Example:} {\scriptsize\ttfamily \mnem{FADD.f32}\imod{.wait0.reconverge} \dst{r0}, \src{r1}, \src{r2}} (\textbf{SPIR-V:} \texttt{OpControlBarrier / OpReturn}; \textbf{GLSL:} \texttt{// Scoreboard wait + warp reconvergence})
\end{tightitemize}
\clearpage

\subsection*{FLOG}
\addcontentsline{toc}{subsection}{FLOG}
Directly computes the base-2 logarithm \texttt{dest = log2(src0)} (\texttt{f16} on v12+, \texttt{f32} on v14+) in a single SFU instruction. Returns \texttt{-Inf} for \texttt{$\pm$0.0}, \texttt{NaN} for negative inputs, and \texttt{+Inf} for \texttt{+Inf}.
\par\textbf{Notes:}
\begin{tightitemize}
  \item Negative -\textgreater{} NaN, \texttt{0} -\textgreater{} \texttt{-Inf}, \texttt{+Inf} -\textgreater{} \texttt{+Inf}.
  \item \textbf{Source references:} \texttt{kraid/ops.rs:1640-1656}, \texttt{encode\_v9.rs:1402-1419}, \texttt{nir.rs:931-942}, \texttt{isa-v9-v14.xml}.
\end{tightitemize}
\codebox{\mnem{FLOG}\imod{\{.flow\}} \dst{dest}\imod{\{.dest\_width\_replicate\}}, \src{src0}\imod{\{.lane0\}}\imod{\{.abs0\}}\imod{\{.neg0\}}\\\mnem{FLOG}\imod{\{.flow\}} \dst{dest}\imod{\{.dest\_width\}}, \src{src0}\imod{\{.widen0\}}\imod{\{.abs0\}}\imod{\{.neg0\}}}
\textbf{Operands:}
\begin{tightitemize}\footnotesize
  \item \texttt{\dst{dest}} (\texttt{dest\_register\_index}, \texttt{dest\_width}, \texttt{dest\_width\_replicate}): Destination GPR receiving \texttt{log2(src0)}.
  \item \texttt{\src{src0}}: Floating-point input (\texttt{f16} on v12+, \texttt{f32} on v14+).
\end{tightitemize}
\par\textbf{Closest SPIR-V Op:} \texttt{GLSL.std.450 Log2 (v12+ f16 / v14+ f32)}
\par\textbf{Pseudo-GLSL (Operational Semantics):}
\codebox{dest = (src0 == 0.0) ? -INF : ((src0 \textless{} 0.0) ? NaN : log2(src0));}
\textbf{Execution Context:}
\begin{tightitemize}\footnotesize
  \item \textbf{Compute stage preloads (\hyperref[sec:csf-spd]{\texttt{SPD\_k}}, \texttt{r16..r23}):} Hardware pre-fills \texttt{r55..r57} with \texttt{WorkGroupID.xyz + JOB\_OFFSET} (from \hyperref[sec:csf-regs]{\texttt{COMPUTESR::JOB\_OFFSET\_X/Y/Z}} \texttt{r34..r36}) and \texttt{r60..r62} with \texttt{LocalInvocationID.xyz} or \texttt{GlobalInvocationID.xyz} (from \texttt{WG\_SIZE} \texttt{r33} \& \texttt{JOB\_SIZE} \texttt{r37..r39}, selected by \texttt{ShaderProgram::preload}) so compute shaders can derive thread/element indices without memory loads.
  \item \textbf{Vertex stage preloads (\hyperref[sec:csf-spd]{\texttt{SPD\_0}/\texttt{SPD\_1}}, \texttt{r16..r19}):} The hardware Index Fetcher reads \hyperref[sec:csf-regs]{\texttt{IDVSSR::INDEX\_BUFFER}} (\texttt{r54:r55}) and pre-fills \texttt{r59, r60} with \texttt{VertexIndex} (\texttt{+ VERTEX\_OFFSET + GLOBAL\_ATTR\_OFF}) and \texttt{r61} with \texttt{InstanceIndex} (\texttt{+ INSTANCE\_OFFSET}) so Position and Varying shaders can pull vertex attributes directly.
  \item \textbf{Fragment stage preloads \& FP modes (\hyperref[sec:csf-spd]{\texttt{SPD\_2}}, \texttt{r20:r21}):} Hardware pre-fills \texttt{r59} with \texttt{floatBitsToUint(FragCoord.z)} and \texttt{r60} with the MSAA sample coverage mask (\texttt{gl\_SampleMaskIn[0]}), while \texttt{ShaderProgram} word~0 latches 32/64 GPR allocation and warp-wide \texttt{flush\_to\_zero\_mode}, \texttt{suppress\_na\_n}, and \texttt{suppress\_inf}.
  \item \textbf{Inline FAU \& Special Operands (\hyperref[sec:csf-tsd-fau]{\texttt{FAU\_k}}, \texttt{r8..r15}):} Any 8-bit source with \texttt{mode=2} reads zero-latency uniform \texttt{u0..u31} from \texttt{FAU\_k} (\texttt{slots64 = bits[63:56]}); \texttt{mode=3} reads \texttt{SmallConstant} (Table~1.2) or \texttt{FauSpecial} (\texttt{pc}, \texttt{lane\_id}, \texttt{lane\_width}, \texttt{tls\_ptr}, \texttt{wls\_ptr}, \texttt{blend\_descriptor\_0..7}).
\end{tightitemize}
\par\textbf{Worked Examples:}
\par\textbf{Example 1 (Exact power of two):}
\codebox{\mnem{FLOG.f32} \dst{r0}, \src{r1}}
With \texttt{r1 = 8.0f}: computes \texttt{log2(8.0) = 3.0f} (\texttt{0x40400000}); with \texttt{r1 = 0.5f} -\textgreater{} \texttt{-1.0f}.
\par\textbf{Example 2 (Domain boundary values):}
\codebox{\mnem{FLOG.f32} \dst{r0}, \src{r1}}
With \texttt{r1 = $\pm$0.0f} -\textgreater{} \texttt{-Inf} (\texttt{0xFF800000}); with \texttt{r1 = -1.0f} -\textgreater{} \texttt{qNaN}; with \texttt{r1 = +Inf} -\textgreater{} \texttt{+Inf}.

\begin{center}\footnotesize
\begin{tabularx}{\linewidth}{@{}lllX@{}}
\toprule
Variant & Arch & Class & Opcode Fields \\
\midrule
\texttt{(none)} & v9--v14 & \texttt{R1} & \texttt{opcode2=FLOG\_F16, opcode1=TABLE} \\
\texttt{(none)} & v9--v14 & \texttt{R1} & \texttt{opcode2=FLOG\_F32, opcode1=TABLE} \\
\bottomrule
\end{tabularx}
\end{center}
\begin{center}\footnotesize
\begin{tabularx}{0.88\linewidth}{@{}l l X@{}}
\toprule
Field Name & Bits & Type / Fixed Value \\
\midrule
\texttt{src0} & \texttt{0:7} & \texttt{SourceEncoding} \\
\textit{Reserved} & \texttt{8:15} & \texttt{zero} \\
\texttt{opcode2} & \texttt{16:20} & \texttt{exact=FLOG\_F16} \\
\textit{Reserved} & \texttt{21:27} & \texttt{zero} \\
\hyperref[enum:lane16_m]{\texttt{lane0}} & \texttt{28} & \texttt{lane16\_m (def h0)} \\
\textit{Reserved} & \texttt{29:37} & \texttt{zero} \\
\hyperref[enum:neg_m]{\texttt{neg0}} & \texttt{38} & \texttt{neg\_m (def none)} \\
\hyperref[enum:abs_m]{\texttt{abs0}} & \texttt{39} & \texttt{abs\_m (def none)} \\
\texttt{dest\_register\_index} & \texttt{40:45} & \texttt{uint} \\
\hyperref[enum:dest_width_replicate_m]{\texttt{dest\_width\_replicate}} & \texttt{46:47} & \texttt{dest\_width\_replicate\_m} \\
\texttt{opcode1} & \texttt{48:51} & \texttt{exact=TABLE} \\
\texttt{instruction\_class} & \texttt{52:56} & \texttt{exact=R1} \\
\texttt{fau\_page\_index} & \texttt{57:58} & \texttt{uint} \\
\hyperref[enum:flow_control_m]{\texttt{flow}} & \texttt{59:62} & \texttt{flow\_control\_m} \\
\textit{Reserved} & \texttt{63} & \texttt{zero} \\
\bottomrule
\end{tabularx}
\end{center}
\textbf{Fields Explained:}
\begin{tightitemize}\footnotesize
  \item \texttt{src0}: Primary source operand (\texttt{r0..r63}, \texttt{u0..u63} uniform, \texttt{k*} FAU constant, or \texttt{t0..t2} temporary; \texttt{\textasciicircum{}} marks discard/last-use). \emph{Example:} {\scriptsize\ttfamily \mnem{FADD.f32} \dst{r0}, \src{r1}, \src{r2}} (\textbf{SPIR-V:} \texttt{OpFAdd (Operand 0)}; \textbf{GLSL:} \texttt{float c = a + b;})
  \item \texttt{opcode2}: Secondary sub-opcode further specializing instructions sharing \texttt{instruction\_class} and \texttt{opcode1} (e.g., conversions, SFU tables, moves). \emph{Example:} {\scriptsize\ttfamily \mnem{F32\_TO\_S32}\imod{.round\_zero} \dst{r0}, \src{r1}} (\textbf{SPIR-V:} \texttt{OpConvertFToS}; \textbf{GLSL:} \texttt{int y = int(x);})
  \item \hyperref[enum:lane16_m]{\texttt{lane0}}: Stage A single-lane extractor on \texttt{src0} (\texttt{b0..b3} for 8-bit or \texttt{h0}/\texttt{h1} for 16-bit narrow conversions and \texttt{MKVEC}). \emph{Example:} {\scriptsize\ttfamily \mnem{F16\_TO\_F32} \dst{r0}, \src{r1}\imod{.h1}} (\textbf{SPIR-V:} \texttt{OpCompositeExtract}; \textbf{GLSL:} \texttt{float y = float(x.y);})
  \item \hyperref[enum:neg_m]{\texttt{neg0}}: Stage A floating-point negation on \texttt{src0} (flips IEEE sign bit after swizzle/widen and \texttt{abs0}). \emph{Example:} {\scriptsize\ttfamily \mnem{FADD.f32} \dst{r0}, \src{r1}\imod{.neg}, \src{r2}} (\textbf{SPIR-V:} \texttt{OpFNegate}; \textbf{GLSL:} \texttt{float c = -a + b;})
  \item \hyperref[enum:abs_m]{\texttt{abs0}}: Stage A floating-point absolute value on \texttt{src0} (clears IEEE sign bit after swizzle/widen, before \texttt{neg0}). \emph{Example:} {\scriptsize\ttfamily \mnem{FADD.f32} \dst{r0}, \src{r1}\imod{.abs}, \src{r2}} (\textbf{SPIR-V:} \texttt{GLSL.std.450 FAbs}; \textbf{GLSL:} \texttt{float c = abs(a) + b;})
  \item \texttt{dest\_register\_index}: 6-bit destination GPR index (\texttt{r0..r63}) receiving the Stage C writeback result (even-aligned pair \texttt{[rN:rN+1]} for 64-bit ops). \emph{Example:} {\scriptsize\ttfamily \mnem{FMA.f32} \dst{r0}, \src{r1}, \src{r2}, \src{r3}} (\textbf{SPIR-V:} \texttt{SSA Result ID}; \textbf{GLSL:} \texttt{float r0 = fma(r1, r2, r3);})
  \item \hyperref[enum:dest_width_replicate_m]{\texttt{dest\_width\_replicate}}: Destination halfword lane-mask / replication mode (\texttt{h0}, \texttt{h1}, \texttt{h01}, \texttt{none}). \emph{Example:} {\scriptsize\ttfamily \mnem{F32\_TO\_F16} \dst{r0}\imod{.h01}, \src{r1}} (\textbf{SPIR-V:} \texttt{OpFConvert (f32 -\textgreater{} f16)}; \textbf{GLSL:} \texttt{float16\_t h = float16\_t(x);})
  \item \texttt{opcode1}: Primary sub-opcode within the instruction class uniquely identifying the operation family. \emph{Example:} {\scriptsize\ttfamily \mnem{IADD.s32} \dst{r0}, \src{r1}, \src{r2}} (\textbf{SPIR-V:} \texttt{OpIAdd (operation selector)}; \textbf{GLSL:} \texttt{// Primary opcode bits})
  \item \texttt{instruction\_class}: 4-bit major instruction class selector (\texttt{bits 48:51}) partitioning the 64-bit Valhall opcode space (e.g., \texttt{FMA}, \texttt{CVT}, \texttt{SFU}, \texttt{LS}, \texttt{TEX}, \texttt{BRANCH}). \emph{Example:} {\scriptsize\ttfamily \mnem{FADD.f32} \dst{r0}, \src{r1}, \src{r2}} (\textbf{SPIR-V:} \texttt{N/A (encoding class)}; \textbf{GLSL:} \texttt{// Hardware execution pipe / decode class})
  \item \texttt{fau\_page\_index}: Selects the FAU (Fast Access Uniform) page (\texttt{0..3}) when a source operand references hardware small constants, thread IDs, or lane constants. \emph{Example:} {\scriptsize\ttfamily \mnem{FADD.f32} \dst{r0}, \src{r1}, \src{k1f}} (\textbf{SPIR-V:} \texttt{OpConstant / BuiltIn}; \textbf{GLSL:} \texttt{float y = x + 1.0;})
  \item \hyperref[enum:flow_control_m]{\texttt{flow}}: 4-bit embedded flow-control \& scoreboard wait modifier (\texttt{bits 59:62}: \texttt{none}, \texttt{wait0..wait0126}, \texttt{wait}, \texttt{reconverge}, \texttt{discard}, \texttt{end}). \emph{Example:} {\scriptsize\ttfamily \mnem{FADD.f32}\imod{.wait0.reconverge} \dst{r0}, \src{r1}, \src{r2}} (\textbf{SPIR-V:} \texttt{OpControlBarrier / OpReturn}; \textbf{GLSL:} \texttt{// Scoreboard wait + warp reconvergence})
\end{tightitemize}
\clearpage

\subsection*{FLUSH}
\addcontentsline{toc}{subsection}{FLUSH}
Canonicalizes IEEE-754 special values independently for each floating-point element $i \in \{0, \dots, N-1\}$ of \texttt{src0} ($N=1$ \texttt{f32}, $N=2$ \texttt{v2f16}): NaN maps to \texttt{+0.0} (\texttt{.flush\_nan}) or quietened NaN (\texttt{.quiet\_nan}); \texttt{$\pm$Inf} clamps to \texttt{$\pm$MAX\_FINITE} (\texttt{.flush\_inf}); subnormals flush to \texttt{$\pm$0.0} (\texttt{.ftz}); normal finite values pass through unchanged.
\par\textbf{Notes:}
\begin{tightitemize}
  \item \texttt{qnan\_bit} is \texttt{0x00400000} (f32) / \texttt{0x0200} (f16).
  \item \texttt{MAX\_FINITE} is \texttt{0x7F7FFFFF} / \texttt{0x7BFF}.
  \item Exact bit-for-bit verified (\texttt{hw\_tests.rs:1570-1596}).
  \item \textbf{Source references:} \texttt{kraid/ops.rs:1681-1783}, \texttt{encode\_v9.rs:1440-1470}, \texttt{hw\_tests.rs:1570-1596}.
\end{tightitemize}
\codebox{\mnem{FLUSH}\imod{\{.flow\}}\imod{\{.flush\_to\_zero\_mode\}}\imod{\{.inf\_mode\}}\imod{\{.nan\_mode\}} \dst{dest}\imod{\{.dest\_width\}}, \src{src0}\imod{\{.widen0\}}\imod{\{.abs0\}}\imod{\{.neg0\}}\\\mnem{FLUSH}\imod{\{.flow\}}\imod{\{.flush\_to\_zero\_mode\}}\imod{\{.inf\_mode\}}\imod{\{.nan\_mode\}} \dst{dest}\imod{\{.dest\_width\}}, \src{src0}\imod{\{.swz0\}}\imod{\{.abs0\}}\imod{\{.neg0\}}}
\textbf{Operands:}
\begin{tightitemize}\footnotesize
  \item \texttt{\dst{dest}} (\texttt{dest\_register\_index}, \texttt{dest\_width}): Destination GPR receiving canonicalized floating-point output.
  \item \texttt{\src{src0}}: Floating-point input subject to \texttt{.ftz}, \texttt{.flush\_inf}, \texttt{.flush\_nan}, or \texttt{.quiet\_nan}.
\end{tightitemize}
\par\textbf{Closest SPIR-V Op:} \texttt{None (FP canonicalization / FTZ)}
\par\textbf{Pseudo-GLSL (Operational Semantics):}
\codebox{for (int i = 0; i \textless{} N; ++i) \{\\    auto x = src0[i];\\    if (flag\_ftz \&\& isSubnormal(x))   x = copysign(0.0, x);\\    if (flag\_flush\_inf \&\& isinf(x))   x = copysign(MAX\_FINITE, x);\\    if (flag\_flush\_nan \&\& isnan(x))   x = +0.0;\\    else if (flag\_quiet\_nan \&\& isnan(x)) x = quieten(x);\\    dest[i] = x;\\\}}
\textbf{Execution Context:}
\begin{tightitemize}\footnotesize
  \item \textbf{Compute stage preloads (\hyperref[sec:csf-spd]{\texttt{SPD\_k}}, \texttt{r16..r23}):} Hardware pre-fills \texttt{r55..r57} with \texttt{WorkGroupID.xyz + JOB\_OFFSET} (from \hyperref[sec:csf-regs]{\texttt{COMPUTESR::JOB\_OFFSET\_X/Y/Z}} \texttt{r34..r36}) and \texttt{r60..r62} with \texttt{LocalInvocationID.xyz} or \texttt{GlobalInvocationID.xyz} (from \texttt{WG\_SIZE} \texttt{r33} \& \texttt{JOB\_SIZE} \texttt{r37..r39}, selected by \texttt{ShaderProgram::preload}) so compute shaders can derive thread/element indices without memory loads.
  \item \textbf{Vertex stage preloads (\hyperref[sec:csf-spd]{\texttt{SPD\_0}/\texttt{SPD\_1}}, \texttt{r16..r19}):} The hardware Index Fetcher reads \hyperref[sec:csf-regs]{\texttt{IDVSSR::INDEX\_BUFFER}} (\texttt{r54:r55}) and pre-fills \texttt{r59, r60} with \texttt{VertexIndex} (\texttt{+ VERTEX\_OFFSET + GLOBAL\_ATTR\_OFF}) and \texttt{r61} with \texttt{InstanceIndex} (\texttt{+ INSTANCE\_OFFSET}) so Position and Varying shaders can pull vertex attributes directly.
  \item \textbf{Fragment stage preloads \& FP modes (\hyperref[sec:csf-spd]{\texttt{SPD\_2}}, \texttt{r20:r21}):} Hardware pre-fills \texttt{r59} with \texttt{floatBitsToUint(FragCoord.z)} and \texttt{r60} with the MSAA sample coverage mask (\texttt{gl\_SampleMaskIn[0]}), while \texttt{ShaderProgram} word~0 latches 32/64 GPR allocation and warp-wide \texttt{flush\_to\_zero\_mode}, \texttt{suppress\_na\_n}, and \texttt{suppress\_inf}.
  \item \textbf{Inline FAU \& Special Operands (\hyperref[sec:csf-tsd-fau]{\texttt{FAU\_k}}, \texttt{r8..r15}):} Any 8-bit source with \texttt{mode=2} reads zero-latency uniform \texttt{u0..u31} from \texttt{FAU\_k} (\texttt{slots64 = bits[63:56]}); \texttt{mode=3} reads \texttt{SmallConstant} (Table~1.2) or \texttt{FauSpecial} (\texttt{pc}, \texttt{lane\_id}, \texttt{lane\_width}, \texttt{tls\_ptr}, \texttt{wls\_ptr}, \texttt{blend\_descriptor\_0..7}).
\end{tightitemize}
\par\textbf{Worked Examples:}
\par\textbf{Example 1 (Infinity clamping + sNaN quietening):}
\codebox{\mnem{FLUSH.f32}\imod{.ftz.flush\_inf.quiet\_nan} \dst{r0}, \src{r1}}
With \texttt{r1 = 0xFF800000} (\texttt{-Inf}): \texttt{.flush\_inf} clamps to \texttt{-MAX\_FINITE} (\texttt{0xFF7FFFFF}); an sNaN \texttt{0x7F800001} becomes \texttt{0x7FC00001}.
\par\textbf{Example 2 (Subnormal FTZ preserving sign):}
\codebox{\mnem{FLUSH.f32}\imod{.ftz} \dst{r0}, \src{r1}}
With \texttt{r1 = 0x80000001} (smallest negative FP32 subnormal): \texttt{.ftz} flushes to \texttt{-0.0f} (\texttt{0x80000000}), preserving sign bit.

\begin{center}\footnotesize
\begin{tabularx}{\linewidth}{@{}lllX@{}}
\toprule
Variant & Arch & Class & Opcode Fields \\
\midrule
\texttt{(none)} & v9--v14 & \texttt{R1} & \texttt{opcode2=FLUSH\_F32, opcode1=FLOAT} \\
\texttt{(none)} & v9--v14 & \texttt{R1} & \texttt{opcode2=FLUSH\_V2F16, opcode1=FLOAT} \\
\bottomrule
\end{tabularx}
\end{center}
\begin{center}\footnotesize
\begin{tabularx}{0.88\linewidth}{@{}l l X@{}}
\toprule
Field Name & Bits & Type / Fixed Value \\
\midrule
\texttt{src0} & \texttt{0:7} & \texttt{SourceEncoding} \\
\hyperref[enum:flush_nan_m]{\texttt{nan\_mode}} & \texttt{8:9} & \texttt{flush\_nan\_m} \\
\hyperref[enum:flush_ftz_m]{\texttt{flush\_to\_zero\_mode}} & \texttt{10} & \texttt{flush\_ftz\_m} \\
\hyperref[enum:flush_inf_m]{\texttt{inf\_mode}} & \texttt{11} & \texttt{flush\_inf\_m} \\
\textit{Reserved} & \texttt{12:15} & \texttt{zero} \\
\texttt{opcode2} & \texttt{16:20} & \texttt{exact=FLUSH\_F32} \\
\textit{Reserved} & \texttt{21:27} & \texttt{zero} \\
\hyperref[enum:widen_m]{\texttt{widen0}} & \texttt{28:29} & \texttt{widen\_m (def none)} \\
\textit{Reserved} & \texttt{30:37} & \texttt{zero} \\
\hyperref[enum:neg_m]{\texttt{neg0}} & \texttt{38} & \texttt{neg\_m (def none)} \\
\hyperref[enum:abs_m]{\texttt{abs0}} & \texttt{39} & \texttt{abs\_m (def none)} \\
\texttt{dest\_register\_index} & \texttt{40:45} & \texttt{uint} \\
\hyperref[enum:dest_width_m]{\texttt{dest\_width}} & \texttt{46:47} & \texttt{dest\_width\_m} \\
\texttt{opcode1} & \texttt{48:51} & \texttt{exact=FLOAT} \\
\texttt{instruction\_class} & \texttt{52:56} & \texttt{exact=R1} \\
\texttt{fau\_page\_index} & \texttt{57:58} & \texttt{uint} \\
\hyperref[enum:flow_control_m]{\texttt{flow}} & \texttt{59:62} & \texttt{flow\_control\_m} \\
\textit{Reserved} & \texttt{63} & \texttt{zero} \\
\bottomrule
\end{tabularx}
\end{center}
\textbf{Fields Explained:}
\begin{tightitemize}\footnotesize
  \item \texttt{src0}: Primary source operand (\texttt{r0..r63}, \texttt{u0..u63} uniform, \texttt{k*} FAU constant, or \texttt{t0..t2} temporary; \texttt{\textasciicircum{}} marks discard/last-use). \emph{Example:} {\scriptsize\ttfamily \mnem{FADD.f32} \dst{r0}, \src{r1}, \src{r2}} (\textbf{SPIR-V:} \texttt{OpFAdd (Operand 0)}; \textbf{GLSL:} \texttt{float c = a + b;})
  \item \hyperref[enum:flush_nan_m]{\texttt{nan\_mode}}: Controls NaN canonicalization in \texttt{FLUSH} (\texttt{none}, \texttt{flush\_nan}, \texttt{quiet\_nan}). \emph{Example:} {\scriptsize\ttfamily \mnem{FLUSH.f32}\imod{.quiet\_nan} \dst{r0}, \src{r1}} (\textbf{SPIR-V:} \texttt{None (HW sNaN quietening)}; \textbf{GLSL:} \texttt{// Quieten signaling NaN})
  \item \hyperref[enum:flush_ftz_m]{\texttt{flush\_to\_zero\_mode}}: Flushes subnormal values to signed zero in \texttt{FLUSH}. \emph{Example:} {\scriptsize\ttfamily \mnem{FLUSH.f32}\imod{.ftz} \dst{r0}, \src{r1}} (\textbf{SPIR-V:} \texttt{ExecutionMode DenormFlushToZero}; \textbf{GLSL:} \texttt{// FP32 FTZ canonicalization})
  \item \hyperref[enum:flush_inf_m]{\texttt{inf\_mode}}: Controls whether $\pm\infty$ is flushed by \texttt{FLUSH}. \emph{Example:} {\scriptsize\ttfamily \mnem{FLUSH.f32}\imod{.flush\_inf} \dst{r0}, \src{r1}} (\textbf{SPIR-V:} \texttt{None (HW canonicalization)}; \textbf{GLSL:} \texttt{// Clamp/flush Inf in special math})
  \item \texttt{opcode2}: Secondary sub-opcode further specializing instructions sharing \texttt{instruction\_class} and \texttt{opcode1} (e.g., conversions, SFU tables, moves). \emph{Example:} {\scriptsize\ttfamily \mnem{F32\_TO\_S32}\imod{.round\_zero} \dst{r0}, \src{r1}} (\textbf{SPIR-V:} \texttt{OpConvertFToS}; \textbf{GLSL:} \texttt{int y = int(x);})
  \item \hyperref[enum:widen_m]{\texttt{widen0}}: Stage A widening extractor on \texttt{src0} (promotes FP16 \texttt{h0}/\texttt{h1} to FP32, or sign/zero-extends \texttt{b0..b3}/\texttt{h0..h1}/\texttt{w0..w1} to 32/64-bit before \texttt{abs}/\texttt{neg}). \emph{Example:} {\scriptsize\ttfamily \mnem{IADD.s32} \dst{r0}, \src{r1}\imod{.h0}, \src{r2}} (\textbf{SPIR-V:} \texttt{OpSConvert / OpFConvert}; \textbf{GLSL:} \texttt{int c = int(int16\_t(a)) + b;})
  \item \hyperref[enum:neg_m]{\texttt{neg0}}: Stage A floating-point negation on \texttt{src0} (flips IEEE sign bit after swizzle/widen and \texttt{abs0}). \emph{Example:} {\scriptsize\ttfamily \mnem{FADD.f32} \dst{r0}, \src{r1}\imod{.neg}, \src{r2}} (\textbf{SPIR-V:} \texttt{OpFNegate}; \textbf{GLSL:} \texttt{float c = -a + b;})
  \item \hyperref[enum:abs_m]{\texttt{abs0}}: Stage A floating-point absolute value on \texttt{src0} (clears IEEE sign bit after swizzle/widen, before \texttt{neg0}). \emph{Example:} {\scriptsize\ttfamily \mnem{FADD.f32} \dst{r0}, \src{r1}\imod{.abs}, \src{r2}} (\textbf{SPIR-V:} \texttt{GLSL.std.450 FAbs}; \textbf{GLSL:} \texttt{float c = abs(a) + b;})
  \item \texttt{dest\_register\_index}: 6-bit destination GPR index (\texttt{r0..r63}) receiving the Stage C writeback result (even-aligned pair \texttt{[rN:rN+1]} for 64-bit ops). \emph{Example:} {\scriptsize\ttfamily \mnem{FMA.f32} \dst{r0}, \src{r1}, \src{r2}, \src{r3}} (\textbf{SPIR-V:} \texttt{SSA Result ID}; \textbf{GLSL:} \texttt{float r0 = fma(r1, r2, r3);})
  \item \hyperref[enum:dest_width_m]{\texttt{dest\_width}}: Stage C destination lane write-mask (\texttt{none} = full 32-bit, \texttt{h0} = low 16-bit, \texttt{h1} = high 16-bit without bit-shifting). \emph{Example:} {\scriptsize\ttfamily \mnem{FADD.v2f16} \dst{r0}\imod{.h1}, \src{r1}, \src{r2}} (\textbf{SPIR-V:} \texttt{OpVectorInsertDynamic / WriteMask}; \textbf{GLSL:} \texttt{dst.y = a.y + b.y;})
  \item \texttt{opcode1}: Primary sub-opcode within the instruction class uniquely identifying the operation family. \emph{Example:} {\scriptsize\ttfamily \mnem{IADD.s32} \dst{r0}, \src{r1}, \src{r2}} (\textbf{SPIR-V:} \texttt{OpIAdd (operation selector)}; \textbf{GLSL:} \texttt{// Primary opcode bits})
  \item \texttt{instruction\_class}: 4-bit major instruction class selector (\texttt{bits 48:51}) partitioning the 64-bit Valhall opcode space (e.g., \texttt{FMA}, \texttt{CVT}, \texttt{SFU}, \texttt{LS}, \texttt{TEX}, \texttt{BRANCH}). \emph{Example:} {\scriptsize\ttfamily \mnem{FADD.f32} \dst{r0}, \src{r1}, \src{r2}} (\textbf{SPIR-V:} \texttt{N/A (encoding class)}; \textbf{GLSL:} \texttt{// Hardware execution pipe / decode class})
  \item \texttt{fau\_page\_index}: Selects the FAU (Fast Access Uniform) page (\texttt{0..3}) when a source operand references hardware small constants, thread IDs, or lane constants. \emph{Example:} {\scriptsize\ttfamily \mnem{FADD.f32} \dst{r0}, \src{r1}, \src{k1f}} (\textbf{SPIR-V:} \texttt{OpConstant / BuiltIn}; \textbf{GLSL:} \texttt{float y = x + 1.0;})
  \item \hyperref[enum:flow_control_m]{\texttt{flow}}: 4-bit embedded flow-control \& scoreboard wait modifier (\texttt{bits 59:62}: \texttt{none}, \texttt{wait0..wait0126}, \texttt{wait}, \texttt{reconverge}, \texttt{discard}, \texttt{end}). \emph{Example:} {\scriptsize\ttfamily \mnem{FADD.f32}\imod{.wait0.reconverge} \dst{r0}, \src{r1}, \src{r2}} (\textbf{SPIR-V:} \texttt{OpControlBarrier / OpReturn}; \textbf{GLSL:} \texttt{// Scoreboard wait + warp reconvergence})
\end{tightitemize}
\clearpage

\subsection*{FMA}
\addcontentsline{toc}{subsection}{FMA}
Computes single-rounded fused multiply-add independently for each floating-point element $i \in \{0, \dots, N-1\}$ ($N=1$ for \texttt{f32}, $N=2$ for \texttt{v2f16}): \texttt{dest[i] = clamp(round\_rm(fma(src0[i], src1[i], src2[i])))}.
\par\textbf{Notes:}
\begin{tightitemize}
  \item Single rounding for \texttt{a*b+c}.
  \item \texttt{dest\_width\_narrow} narrows f32-\textgreater{}f16 using \texttt{round} (\texttt{dst\_mod\_prop.rs:566-569}).
  \item \textbf{Source references:} \texttt{kraid/ops.rs:1788-1824}, \texttt{encode\_v9.rs:1472-1495}, \texttt{hw\_tests.rs:1599-1631}.
\end{tightitemize}
\codebox{\mnem{FMA}\imod{\{.flow\}}\imod{\{.round\}}\imod{\{.clamp\}} \dst{dest}\imod{\{.dest\_width\_narrow\}}, \src{src0}\imod{\{.widen0\}}\imod{\{.abs0\}}\imod{\{.neg0\}}, \src{src1}\imod{\{.widen1\}}\imod{\{.abs1\}}\imod{\{.neg1\}}, \src{src2}\imod{\{.widen2\}}\imod{\{.abs2\}}\imod{\{.neg2\}}\\\mnem{FMA}\imod{\{.flow\}}\imod{\{.round\}}\imod{\{.clamp\}} \dst{dest}\imod{\{.dest\_width\}}, \src{src0}\imod{\{.swz0\}}\imod{\{.abs0\}}\imod{\{.neg0\}}, \src{src1}\imod{\{.swz1\}}\imod{\{.abs1\}}\imod{\{.neg1\}}, \src{src2}\imod{\{.swz2\}}\imod{\{.abs2\}}\imod{\{.neg2\}}}
\textbf{Operands:}
\begin{tightitemize}\footnotesize
  \item \texttt{\dst{dest}} (\texttt{dest\_register\_index}, \texttt{dest\_width}, \texttt{dest\_width\_narrow}): Destination GPR receiving \texttt{fma(src0, src1, src2)} (with optional clamp/narrowing).
  \item \texttt{\src{src0}}: First multiplicand.
  \item \texttt{\src{src1}}: Second multiplicand.
  \item \texttt{\src{src2}}: Addend accumulated with single rounding.
\end{tightitemize}
\par\textbf{Closest SPIR-V Op:} \texttt{GLSL.std.450 Fma}
\par\textbf{Pseudo-GLSL (Operational Semantics):}
\codebox{for (int i = 0; i \textless{} N; ++i) \{\\    auto r = clamp\_mode(round\_rm(fma(mod(src0)[i], mod(src1)[i], mod(src2)[i])));\\    dest[i] = narrowOrMask(dest\_width\_narrow, r);\\\}}
\textbf{Execution Context:}
\begin{tightitemize}\footnotesize
  \item \textbf{Compute stage preloads (\hyperref[sec:csf-spd]{\texttt{SPD\_k}}, \texttt{r16..r23}):} Hardware pre-fills \texttt{r55..r57} with \texttt{WorkGroupID.xyz + JOB\_OFFSET} (from \hyperref[sec:csf-regs]{\texttt{COMPUTESR::JOB\_OFFSET\_X/Y/Z}} \texttt{r34..r36}) and \texttt{r60..r62} with \texttt{LocalInvocationID.xyz} or \texttt{GlobalInvocationID.xyz} (from \texttt{WG\_SIZE} \texttt{r33} \& \texttt{JOB\_SIZE} \texttt{r37..r39}, selected by \texttt{ShaderProgram::preload}) so compute shaders can derive thread/element indices without memory loads.
  \item \textbf{Vertex stage preloads (\hyperref[sec:csf-spd]{\texttt{SPD\_0}/\texttt{SPD\_1}}, \texttt{r16..r19}):} The hardware Index Fetcher reads \hyperref[sec:csf-regs]{\texttt{IDVSSR::INDEX\_BUFFER}} (\texttt{r54:r55}) and pre-fills \texttt{r59, r60} with \texttt{VertexIndex} (\texttt{+ VERTEX\_OFFSET + GLOBAL\_ATTR\_OFF}) and \texttt{r61} with \texttt{InstanceIndex} (\texttt{+ INSTANCE\_OFFSET}) so Position and Varying shaders can pull vertex attributes directly.
  \item \textbf{Fragment stage preloads \& FP modes (\hyperref[sec:csf-spd]{\texttt{SPD\_2}}, \texttt{r20:r21}):} Hardware pre-fills \texttt{r59} with \texttt{floatBitsToUint(FragCoord.z)} and \texttt{r60} with the MSAA sample coverage mask (\texttt{gl\_SampleMaskIn[0]}), while \texttt{ShaderProgram} word~0 latches 32/64 GPR allocation and warp-wide \texttt{flush\_to\_zero\_mode}, \texttt{suppress\_na\_n}, and \texttt{suppress\_inf}.
  \item \textbf{Inline FAU \& Special Operands (\hyperref[sec:csf-tsd-fau]{\texttt{FAU\_k}}, \texttt{r8..r15}):} Any 8-bit source with \texttt{mode=2} reads zero-latency uniform \texttt{u0..u31} from \texttt{FAU\_k} (\texttt{slots64 = bits[63:56]}); \texttt{mode=3} reads \texttt{SmallConstant} (Table~1.2) or \texttt{FauSpecial} (\texttt{pc}, \texttt{lane\_id}, \texttt{lane\_width}, \texttt{tls\_ptr}, \texttt{wls\_ptr}, \texttt{blend\_descriptor\_0..7}).
\end{tightitemize}
\par\textbf{Worked Examples:}
\par\textbf{Example 1 (Single-rounded multiply-subtract):}
\codebox{\mnem{FMA.f32} \dst{r0}, \src{r1}, \src{r2}\imod{.neg}, \src{r3}}
With \texttt{r1 = 2.0f, r2 = 3.0f, r3 = 10.0f}: Stage A negates \texttt{r2} to \texttt{-3.0f}; single-rounded \texttt{2.0 * (-3.0) + 10.0 = 4.0f}.
\par\textbf{Example 2 (Clamped FMA with fused FP16 replication (.h01)):}
\codebox{\mnem{FMA.f32}\imod{.clamp\_0\_1} \dst{r0}\imod{.h01}, \src{r1}, \src{r2}, \src{r3}}
With \texttt{r1 = 0.5f, r2 = 1.0f, r3 = 0.75f}: exact \texttt{1.25f} clamps to \texttt{1.0f}, narrows to \texttt{0x3C00}, and \texttt{.h01} replicates -\textgreater{} \texttt{r0 = 0x3C003C00}.
\par\textbf{Mined SPIR-V $\to$ Valhall ISA Example (3D vector normalization and dot product with fused multiply-add --- Fragment shader, Mali-G710 (v10)):}
\codebox{\textbf{; SPIR-V (input)}\\\%len2 = OpDot \%float \%n \%n\\\%inv  = OpExtInst \%float \%glsl InverseSqrt \%len2\\\%ndotl = OpExtInst \%float \%glsl Fma \%ny \%ly \%nx\_lx\\\%half\_lambert = OpExtInst \%float \%glsl FClamp (fma(\%ndotl, 0.5, 0.5)) 0.0 1.0\\[3pt]\textbf{; Valhall ISA (malioc output)}\\FMUL.f32 r3, r0, r0\\FMA.f32 r3, r1, r1, r3\textasciicircum{}\\FMA.f32 r3, r2, r2, r3\textasciicircum{}\\FRSQ.f32 r3, r3\textasciicircum{}\\FMA.f32.wait1.clamp\_0\_1 r0, r0\textasciicircum{}, k0fr2f.h1, k0fr2f.h1}
\texttt{OpDot} lowers to \texttt{FMUL} + two \texttt{FMA} instructions; \texttt{fma(d, 0.5, 0.5)} clamped to \texttt{[0, 1]} fuses into a single \texttt{FMA.f32.clamp\_0\_1} using \texttt{k0fr2f.h1} (\texttt{0.5f}).
\begin{center}\footnotesize
\begin{tabularx}{\linewidth}{@{}lllX@{}}
\toprule
Variant & Arch & Class & Opcode Fields \\
\midrule
\texttt{(none)} & v9--v14 & \texttt{R3} & \texttt{opcode1=FMA\_F32} \\
\texttt{(none)} & v9--v14 & \texttt{R3} & \texttt{opcode1=FMA\_V2F16} \\
\bottomrule
\end{tabularx}
\end{center}
\begin{center}\footnotesize
\begin{tabularx}{0.88\linewidth}{@{}l l X@{}}
\toprule
Field Name & Bits & Type / Fixed Value \\
\midrule
\texttt{src0} & \texttt{0:7} & \texttt{SourceEncoding} \\
\texttt{src1} & \texttt{8:15} & \texttt{SourceEncoding} \\
\texttt{src2} & \texttt{16:23} & \texttt{SourceEncoding} \\
\hyperref[enum:widen_m]{\texttt{widen2}} & \texttt{24:25} & \texttt{widen\_m (def none)} \\
\hyperref[enum:widen_m]{\texttt{widen1}} & \texttt{26:27} & \texttt{widen\_m (def none)} \\
\hyperref[enum:widen_m]{\texttt{widen0}} & \texttt{28:29} & \texttt{widen\_m (def none)} \\
\hyperref[enum:round_m]{\texttt{round}} & \texttt{30:31} & \texttt{round\_m (def none)} \\
\hyperref[enum:clamp_m]{\texttt{clamp}} & \texttt{32:33} & \texttt{clamp\_m (def none)} \\
\hyperref[enum:neg_m]{\texttt{neg2}} & \texttt{34} & \texttt{neg\_m (def none)} \\
\hyperref[enum:abs_m]{\texttt{abs2}} & \texttt{35} & \texttt{abs\_m (def none)} \\
\hyperref[enum:neg_m]{\texttt{neg1}} & \texttt{36} & \texttt{neg\_m (def none)} \\
\hyperref[enum:abs_m]{\texttt{abs1}} & \texttt{37} & \texttt{abs\_m (def none)} \\
\hyperref[enum:neg_m]{\texttt{neg0}} & \texttt{38} & \texttt{neg\_m (def none)} \\
\hyperref[enum:abs_m]{\texttt{abs0}} & \texttt{39} & \texttt{abs\_m (def none)} \\
\texttt{dest\_register\_index} & \texttt{40:45} & \texttt{uint} \\
\hyperref[enum:dest_width_narrow_m]{\texttt{dest\_width\_narrow}} & \texttt{46:47} & \texttt{dest\_width\_narrow\_m} \\
\texttt{opcode1} & \texttt{48:51} & \texttt{exact=FMA\_F32} \\
\texttt{instruction\_class} & \texttt{52:56} & \texttt{exact=R3} \\
\texttt{fau\_page\_index} & \texttt{57:58} & \texttt{uint} \\
\hyperref[enum:flow_control_m]{\texttt{flow}} & \texttt{59:62} & \texttt{flow\_control\_m} \\
\textit{Reserved} & \texttt{63} & \texttt{zero} \\
\bottomrule
\end{tabularx}
\end{center}
\textbf{Fields Explained:}
\begin{tightitemize}\footnotesize
  \item \texttt{src0}: Primary source operand (\texttt{r0..r63}, \texttt{u0..u63} uniform, \texttt{k*} FAU constant, or \texttt{t0..t2} temporary; \texttt{\textasciicircum{}} marks discard/last-use). \emph{Example:} {\scriptsize\ttfamily \mnem{FADD.f32} \dst{r0}, \src{r1}, \src{r2}} (\textbf{SPIR-V:} \texttt{OpFAdd (Operand 0)}; \textbf{GLSL:} \texttt{float c = a + b;})
  \item \texttt{src1}: Secondary source operand (register, uniform, or FAU constant; in shifts, supplies the shift count byte). \emph{Example:} {\scriptsize\ttfamily \mnem{LSHIFT.i32} \dst{r0}, \src{r1}, \src{r2}\imod{.b0}} (\textbf{SPIR-V:} \texttt{OpShiftLeftLogical (Shift)}; \textbf{GLSL:} \texttt{uint y = a \textless{}\textless{} s;})
  \item \texttt{src2}: Third source operand (FMA addend, \texttt{MUX}/\texttt{CSEL} selector, \texttt{IADDC} carry-in, or bitwise mask in fused shift-logic). \emph{Example:} {\scriptsize\ttfamily \mnem{FMA.f32} \dst{r0}, \src{r1}, \src{r2}, \src{r3}} (\textbf{SPIR-V:} \texttt{GLSL.std.450 Fma (Addend)}; \textbf{GLSL:} \texttt{float d = fma(a, b, c);})
  \item \hyperref[enum:widen_m]{\texttt{widen2}}: Stage A widening extractor on \texttt{src2} (e.g., promoting FP16 addend \texttt{h0}/\texttt{h1} to FP32 in \texttt{FMA.f32}). \emph{Example:} {\scriptsize\ttfamily \mnem{FMA.f32} \dst{r0}, \src{r1}, \src{r2}, \src{r3}\imod{.h0}} (\textbf{SPIR-V:} \texttt{OpFConvert + GLSL.std.450 Fma}; \textbf{GLSL:} \texttt{float d = fma(a, b, float(c16));})
  \item \hyperref[enum:widen_m]{\texttt{widen1}}: Stage A widening extractor on \texttt{src1} (promotes narrow byte/half/word lane before Stage B). \emph{Example:} {\scriptsize\ttfamily \mnem{IMUL.u64} \dst{[r0:r1]}, \src{[r2:r3]}\imod{.w0}, \src{[r4:r5]}\imod{.w0}} (\textbf{SPIR-V:} \texttt{OpUConvert / OpUMulExtended}; \textbf{GLSL:} \texttt{uint64\_t p = uint64\_t(a) * uint64\_t(b);})
  \item \hyperref[enum:widen_m]{\texttt{widen0}}: Stage A widening extractor on \texttt{src0} (promotes FP16 \texttt{h0}/\texttt{h1} to FP32, or sign/zero-extends \texttt{b0..b3}/\texttt{h0..h1}/\texttt{w0..w1} to 32/64-bit before \texttt{abs}/\texttt{neg}). \emph{Example:} {\scriptsize\ttfamily \mnem{IADD.s32} \dst{r0}, \src{r1}\imod{.h0}, \src{r2}} (\textbf{SPIR-V:} \texttt{OpSConvert / OpFConvert}; \textbf{GLSL:} \texttt{int c = int(int16\_t(a)) + b;})
  \item \hyperref[enum:round_m]{\texttt{round}}: Rounding mode (\texttt{none}/RNE, \texttt{round\_up}, \texttt{round\_down}, \texttt{round\_zero}, \texttt{round\_na}) applied in FP ALU, conversions, \texttt{HADD}, \texttt{LDEXP}, and Stage C narrowing. \emph{Example:} {\scriptsize\ttfamily \mnem{FROUND.f32}\imod{.round\_down} \dst{r0}, \src{r1}} (\textbf{SPIR-V:} \texttt{GLSL.std.450 Floor / Ceil / Trunc}; \textbf{GLSL:} \texttt{float y = floor(x);})
  \item \hyperref[enum:clamp_m]{\texttt{clamp}}: Stage B post-op floating-point clamp (\texttt{none}, \texttt{clamp\_0\_inf}, \texttt{clamp\_m1\_1}, \texttt{clamp\_0\_1}) with \texttt{-0.0 \textless{} +0.0} and NaN suppression (\texttt{NaN -\textgreater{} lo}). \emph{Example:} {\scriptsize\ttfamily \mnem{FMA.f32}\imod{.clamp\_0\_1} \dst{r0}, \src{r1}, \src{r2}, \src{r3}} (\textbf{SPIR-V:} \texttt{GLSL.std.450 FClamp}; \textbf{GLSL:} \texttt{float y = clamp(fma(a, b, c), 0.0, 1.0);})
  \item \hyperref[enum:neg_m]{\texttt{neg2}}: Stage A floating-point negation on \texttt{src2} (flips IEEE sign bit of FMA addend after swizzle/widen and \texttt{abs2}). \emph{Example:} {\scriptsize\ttfamily \mnem{FMA.f32} \dst{r0}, \src{r1}, \src{r2}, \src{r3}\imod{.neg}} (\textbf{SPIR-V:} \texttt{OpFNegate}; \textbf{GLSL:} \texttt{float d = fma(a, b, -c);})
  \item \hyperref[enum:abs_m]{\texttt{abs2}}: Stage A floating-point absolute value on \texttt{src2} (clears IEEE sign bit of FMA addend after swizzle/widen). \emph{Example:} {\scriptsize\ttfamily \mnem{FMA.f32} \dst{r0}, \src{r1}, \src{r2}, \src{r3}\imod{.abs}} (\textbf{SPIR-V:} \texttt{GLSL.std.450 FAbs}; \textbf{GLSL:} \texttt{float d = fma(a, b, abs(c));})
  \item \hyperref[enum:neg_m]{\texttt{neg1}}: Stage A floating-point negation on \texttt{src1} (flips IEEE sign bit after swizzle/widen and \texttt{abs1}). \emph{Example:} {\scriptsize\ttfamily \mnem{FMA.f32} \dst{r0}, \src{r1}, \src{r2}\imod{.neg}, \src{r3}} (\textbf{SPIR-V:} \texttt{OpFNegate}; \textbf{GLSL:} \texttt{float d = fma(a, -b, c);})
  \item \hyperref[enum:abs_m]{\texttt{abs1}}: Stage A floating-point absolute value on \texttt{src1} (clears IEEE sign bit after swizzle/widen, before \texttt{neg1}). \emph{Example:} {\scriptsize\ttfamily \mnem{FMA.f32} \dst{r0}, \src{r1}, \src{r2}\imod{.abs}, \src{r3}} (\textbf{SPIR-V:} \texttt{GLSL.std.450 FAbs}; \textbf{GLSL:} \texttt{float d = fma(a, abs(b), c);})
  \item \hyperref[enum:neg_m]{\texttt{neg0}}: Stage A floating-point negation on \texttt{src0} (flips IEEE sign bit after swizzle/widen and \texttt{abs0}). \emph{Example:} {\scriptsize\ttfamily \mnem{FADD.f32} \dst{r0}, \src{r1}\imod{.neg}, \src{r2}} (\textbf{SPIR-V:} \texttt{OpFNegate}; \textbf{GLSL:} \texttt{float c = -a + b;})
  \item \hyperref[enum:abs_m]{\texttt{abs0}}: Stage A floating-point absolute value on \texttt{src0} (clears IEEE sign bit after swizzle/widen, before \texttt{neg0}). \emph{Example:} {\scriptsize\ttfamily \mnem{FADD.f32} \dst{r0}, \src{r1}\imod{.abs}, \src{r2}} (\textbf{SPIR-V:} \texttt{GLSL.std.450 FAbs}; \textbf{GLSL:} \texttt{float c = abs(a) + b;})
  \item \texttt{dest\_register\_index}: 6-bit destination GPR index (\texttt{r0..r63}) receiving the Stage C writeback result (even-aligned pair \texttt{[rN:rN+1]} for 64-bit ops). \emph{Example:} {\scriptsize\ttfamily \mnem{FMA.f32} \dst{r0}, \src{r1}, \src{r2}, \src{r3}} (\textbf{SPIR-V:} \texttt{SSA Result ID}; \textbf{GLSL:} \texttt{float r0 = fma(r1, r2, r3);})
  \item \hyperref[enum:dest_width_narrow_m]{\texttt{dest\_width\_narrow}}: Destination halfword lane-mask / replication mode (\texttt{h0}, \texttt{h1}, \texttt{h01}, \texttt{none}). \emph{Example:} {\scriptsize\ttfamily \mnem{F32\_TO\_F16} \dst{r0}\imod{.h01}, \src{r1}} (\textbf{SPIR-V:} \texttt{OpFConvert (f32 -\textgreater{} f16)}; \textbf{GLSL:} \texttt{float16\_t h = float16\_t(x);})
  \item \texttt{opcode1}: Primary sub-opcode within the instruction class uniquely identifying the operation family. \emph{Example:} {\scriptsize\ttfamily \mnem{IADD.s32} \dst{r0}, \src{r1}, \src{r2}} (\textbf{SPIR-V:} \texttt{OpIAdd (operation selector)}; \textbf{GLSL:} \texttt{// Primary opcode bits})
  \item \texttt{instruction\_class}: 4-bit major instruction class selector (\texttt{bits 48:51}) partitioning the 64-bit Valhall opcode space (e.g., \texttt{FMA}, \texttt{CVT}, \texttt{SFU}, \texttt{LS}, \texttt{TEX}, \texttt{BRANCH}). \emph{Example:} {\scriptsize\ttfamily \mnem{FADD.f32} \dst{r0}, \src{r1}, \src{r2}} (\textbf{SPIR-V:} \texttt{N/A (encoding class)}; \textbf{GLSL:} \texttt{// Hardware execution pipe / decode class})
  \item \texttt{fau\_page\_index}: Selects the FAU (Fast Access Uniform) page (\texttt{0..3}) when a source operand references hardware small constants, thread IDs, or lane constants. \emph{Example:} {\scriptsize\ttfamily \mnem{FADD.f32} \dst{r0}, \src{r1}, \src{k1f}} (\textbf{SPIR-V:} \texttt{OpConstant / BuiltIn}; \textbf{GLSL:} \texttt{float y = x + 1.0;})
  \item \hyperref[enum:flow_control_m]{\texttt{flow}}: 4-bit embedded flow-control \& scoreboard wait modifier (\texttt{bits 59:62}: \texttt{none}, \texttt{wait0..wait0126}, \texttt{wait}, \texttt{reconverge}, \texttt{discard}, \texttt{end}). \emph{Example:} {\scriptsize\ttfamily \mnem{FADD.f32}\imod{.wait0.reconverge} \dst{r0}, \src{r1}, \src{r2}} (\textbf{SPIR-V:} \texttt{OpControlBarrier / OpReturn}; \textbf{GLSL:} \texttt{// Scoreboard wait + warp reconvergence})
\end{tightitemize}
\clearpage

\subsection*{FMAX}
\addcontentsline{toc}{subsection}{FMAX}
Computes IEEE-754 maximum independently for each floating-point element $i \in \{0, \dots, N-1\}$ ($N=1$ \texttt{f32}, $N=2$ \texttt{v2f16}) with deterministic \texttt{-0.0 \textless{} +0.0}. Under \texttt{nan\_suppress} (default), a non-NaN wins over NaN; under \texttt{nan\_propagate}, returns the quietened winning NaN preserving sign and payload (\texttt{\textbar{} qbit}). Optional \texttt{clamp} runs after \texttt{sem}.
\par\textbf{Note:} Verified in \texttt{libGLES\_mali.so} \texttt{\_mali\_max*\_sf32} (\texttt{0x2e28300, 0x2e28590}) / \texttt{\_mali\_max\_sf16} (\texttt{0x2e28450}) and \texttt{kraid/ops.rs:1862-1902}: when both inputs are NaN, both \texttt{nan\_suppress} and \texttt{nan\_propagate} quieten both (\texttt{\textbar{} qbit}) and pick the larger NaN under unsigned total-order key.
\codebox{\mnem{FMAX}\imod{\{.flow\}}\imod{\{.sem\}}\imod{\{.clamp\}} \dst{dest}\imod{\{.dest\_width\}}, \src{src0}\imod{\{.widen0\}}\imod{\{.abs0\}}\imod{\{.neg0\}}, \src{src1}\imod{\{.widen1\}}\imod{\{.abs1\}}\imod{\{.neg1\}}\\\mnem{FMAX}\imod{\{.flow\}}\imod{\{.sem\}}\imod{\{.clamp\}} \dst{dest}\imod{\{.dest\_width\}}, \src{src0}\imod{\{.swz0\}}\imod{\{.abs0\}}\imod{\{.neg0\}}, \src{src1}\imod{\{.swz1\}}\imod{\{.abs1\}}\imod{\{.neg1\}}}
\textbf{Operands:}
\begin{tightitemize}\footnotesize
  \item \texttt{\dst{dest}} (\texttt{dest\_register\_index}, \texttt{dest\_width}): Destination GPR receiving \texttt{max(src0, src1)} per NaN handling mode.
  \item \texttt{\src{src0}}: First floating-point operand.
  \item \texttt{\src{src1}}: Second floating-point operand.
\end{tightitemize}
\par\textbf{Closest SPIR-V Op:} \texttt{GLSL.std.450 NMax / FMax}
\par\textbf{Pseudo-GLSL (Operational Semantics):}
\codebox{for (int i = 0; i \textless{} N; ++i) \{\\    auto a = mod(src0)[i], b = mod(src1)[i];\\    auto r = flag\_nan\_suppress ? (isnan(a) ? b : (isnan(b) ? a : maxTotal(a, b)))\\                               : ((isnan(a) \textbar{}\textbar{} isnan(b)) ? quieten(winNaN(a, b)) : maxTotal(a, b));\\    dest[i] = clamp\_mode(r);\\\}}
\textbf{Execution Context:}
\begin{tightitemize}\footnotesize
  \item \textbf{Compute stage preloads (\hyperref[sec:csf-spd]{\texttt{SPD\_k}}, \texttt{r16..r23}):} Hardware pre-fills \texttt{r55..r57} with \texttt{WorkGroupID.xyz + JOB\_OFFSET} (from \hyperref[sec:csf-regs]{\texttt{COMPUTESR::JOB\_OFFSET\_X/Y/Z}} \texttt{r34..r36}) and \texttt{r60..r62} with \texttt{LocalInvocationID.xyz} or \texttt{GlobalInvocationID.xyz} (from \texttt{WG\_SIZE} \texttt{r33} \& \texttt{JOB\_SIZE} \texttt{r37..r39}, selected by \texttt{ShaderProgram::preload}) so compute shaders can derive thread/element indices without memory loads.
  \item \textbf{Vertex stage preloads (\hyperref[sec:csf-spd]{\texttt{SPD\_0}/\texttt{SPD\_1}}, \texttt{r16..r19}):} The hardware Index Fetcher reads \hyperref[sec:csf-regs]{\texttt{IDVSSR::INDEX\_BUFFER}} (\texttt{r54:r55}) and pre-fills \texttt{r59, r60} with \texttt{VertexIndex} (\texttt{+ VERTEX\_OFFSET + GLOBAL\_ATTR\_OFF}) and \texttt{r61} with \texttt{InstanceIndex} (\texttt{+ INSTANCE\_OFFSET}) so Position and Varying shaders can pull vertex attributes directly.
  \item \textbf{Fragment stage preloads \& FP modes (\hyperref[sec:csf-spd]{\texttt{SPD\_2}}, \texttt{r20:r21}):} Hardware pre-fills \texttt{r59} with \texttt{floatBitsToUint(FragCoord.z)} and \texttt{r60} with the MSAA sample coverage mask (\texttt{gl\_SampleMaskIn[0]}), while \texttt{ShaderProgram} word~0 latches 32/64 GPR allocation and warp-wide \texttt{flush\_to\_zero\_mode}, \texttt{suppress\_na\_n}, and \texttt{suppress\_inf}.
  \item \textbf{Inline FAU \& Special Operands (\hyperref[sec:csf-tsd-fau]{\texttt{FAU\_k}}, \texttt{r8..r15}):} Any 8-bit source with \texttt{mode=2} reads zero-latency uniform \texttt{u0..u31} from \texttt{FAU\_k} (\texttt{slots64 = bits[63:56]}); \texttt{mode=3} reads \texttt{SmallConstant} (Table~1.2) or \texttt{FauSpecial} (\texttt{pc}, \texttt{lane\_id}, \texttt{lane\_width}, \texttt{tls\_ptr}, \texttt{wls\_ptr}, \texttt{blend\_descriptor\_0..7}).
\end{tightitemize}
\par\textbf{Worked Examples:}
\par\textbf{Example 1 (Signed zero ordering \& NaN propagation):}
\codebox{\mnem{FMAX.f32}\imod{.nan\_propagate} \dst{r0}, \src{r1}, \src{r2}}
With \texttt{r1 = -0.0f, r2 = +0.0f}: deterministically returns \texttt{+0.0f}. With \texttt{r1 = 0xFF800001} (sNaN), \texttt{r2 = 1.0f}: \texttt{.nan\_propagate} returns quietened \texttt{0xFFC00001}.
\par\textbf{Example 2 (NaN suppression (IEEE maxNum)):}
\codebox{\mnem{FMAX.f32}\imod{.nan\_suppress} \dst{r0}, \src{r1}, \src{r2}}
With \texttt{r1 = NaN, r2 = -5.0f}: \texttt{.nan\_suppress} ignores NaN and returns \texttt{-5.0f} (\texttt{0xC0A00000}).
\par\textbf{Mined SPIR-V $\to$ Valhall ISA Example (2x2 textureGather max reduction with NaN suppression --- Fragment shader, Mali-G710 (v10)):}
\codebox{\textbf{; SPIR-V (input)}\\\%g = OpImageGather \%v4float \%tex \%uv \%int\_0 None\\\%m0 = OpExtInst \%float \%glsl FMax \%g.y \%g.x\\\%m1 = OpExtInst \%float \%glsl FMax \%g.z \%g.w\\\%res = OpExtInst \%float \%glsl FMax \%m0 \%m1\\[3pt]\textbf{; Valhall ISA (malioc output)}\\TEX\_GATHER.slot0.wait0.skip.tex2d.dst32.rgba.wide.float\_coordinates.gather4\_r @r0:r1:r2:r3, @r0:r1, u0\\FMAX.f32.nan\_suppress r0, r1\textasciicircum{}, r0\textasciicircum{}\\FMAX.f32.nan\_suppress r2, r2\textasciicircum{}, r3\textasciicircum{}\\FMAX.f32.wait.nan\_suppress r4, r0\textasciicircum{}, r2\textasciicircum{}}
Vulkan/GLSL \texttt{FMax} specifies that if one operand is NaN, the non-NaN operand is returned; \texttt{malioc} maps this directly to \texttt{FMAX.f32.nan\_suppress}.
\begin{center}\footnotesize
\begin{tabularx}{\linewidth}{@{}lllX@{}}
\toprule
Variant & Arch & Class & Opcode Fields \\
\midrule
\texttt{f32} & v9--v14 & \texttt{R2} & \texttt{opcode2=MAX} \\
\texttt{v2f16} & v9--v14 & \texttt{R2} & \texttt{opcode2=MAX} \\
\bottomrule
\end{tabularx}
\end{center}
\begin{center}\footnotesize
\begin{tabularx}{0.88\linewidth}{@{}l l X@{}}
\toprule
Field Name & Bits & Type / Fixed Value \\
\midrule
\texttt{src0} & \texttt{0:7} & \texttt{SourceEncoding} \\
\texttt{src1} & \texttt{8:15} & \texttt{SourceEncoding} \\
\texttt{opcode2} & \texttt{16:19} & \texttt{exact=MAX} \\
\textit{Reserved} & \texttt{20:25} & \texttt{zero} \\
\hyperref[enum:widen_m]{\texttt{widen1}} & \texttt{26:27} & \texttt{widen\_m (def none)} \\
\hyperref[enum:widen_m]{\texttt{widen0}} & \texttt{28:29} & \texttt{widen\_m (def none)} \\
\hyperref[enum:sem_m]{\texttt{sem}} & \texttt{30:31} & \texttt{sem\_m (def nan\_suppress)} \\
\hyperref[enum:clamp_m]{\texttt{clamp}} & \texttt{32:33} & \texttt{clamp\_m (def none)} \\
\textit{Reserved} & \texttt{34:35} & \texttt{zero} \\
\hyperref[enum:neg_m]{\texttt{neg1}} & \texttt{36} & \texttt{neg\_m (def none)} \\
\hyperref[enum:abs_m]{\texttt{abs1}} & \texttt{37} & \texttt{abs\_m (def none)} \\
\hyperref[enum:neg_m]{\texttt{neg0}} & \texttt{38} & \texttt{neg\_m (def none)} \\
\hyperref[enum:abs_m]{\texttt{abs0}} & \texttt{39} & \texttt{abs\_m (def none)} \\
\texttt{dest\_register\_index} & \texttt{40:45} & \texttt{uint} \\
\hyperref[enum:dest_width_m]{\texttt{dest\_width}} & \texttt{46:47} & \texttt{dest\_width\_m} \\
\texttt{datatype} & \texttt{48:51} & \texttt{exact=f32} \\
\texttt{instruction\_class} & \texttt{52:56} & \texttt{exact=R2} \\
\texttt{fau\_page\_index} & \texttt{57:58} & \texttt{uint} \\
\hyperref[enum:flow_control_m]{\texttt{flow}} & \texttt{59:62} & \texttt{flow\_control\_m} \\
\textit{Reserved} & \texttt{63} & \texttt{zero} \\
\bottomrule
\end{tabularx}
\end{center}
\textbf{Fields Explained:}
\begin{tightitemize}\footnotesize
  \item \texttt{src0}: Primary source operand (\texttt{r0..r63}, \texttt{u0..u63} uniform, \texttt{k*} FAU constant, or \texttt{t0..t2} temporary; \texttt{\textasciicircum{}} marks discard/last-use). \emph{Example:} {\scriptsize\ttfamily \mnem{FADD.f32} \dst{r0}, \src{r1}, \src{r2}} (\textbf{SPIR-V:} \texttt{OpFAdd (Operand 0)}; \textbf{GLSL:} \texttt{float c = a + b;})
  \item \texttt{src1}: Secondary source operand (register, uniform, or FAU constant; in shifts, supplies the shift count byte). \emph{Example:} {\scriptsize\ttfamily \mnem{LSHIFT.i32} \dst{r0}, \src{r1}, \src{r2}\imod{.b0}} (\textbf{SPIR-V:} \texttt{OpShiftLeftLogical (Shift)}; \textbf{GLSL:} \texttt{uint y = a \textless{}\textless{} s;})
  \item \texttt{opcode2}: Secondary sub-opcode further specializing instructions sharing \texttt{instruction\_class} and \texttt{opcode1} (e.g., conversions, SFU tables, moves). \emph{Example:} {\scriptsize\ttfamily \mnem{F32\_TO\_S32}\imod{.round\_zero} \dst{r0}, \src{r1}} (\textbf{SPIR-V:} \texttt{OpConvertFToS}; \textbf{GLSL:} \texttt{int y = int(x);})
  \item \hyperref[enum:widen_m]{\texttt{widen1}}: Stage A widening extractor on \texttt{src1} (promotes narrow byte/half/word lane before Stage B). \emph{Example:} {\scriptsize\ttfamily \mnem{IMUL.u64} \dst{[r0:r1]}, \src{[r2:r3]}\imod{.w0}, \src{[r4:r5]}\imod{.w0}} (\textbf{SPIR-V:} \texttt{OpUConvert / OpUMulExtended}; \textbf{GLSL:} \texttt{uint64\_t p = uint64\_t(a) * uint64\_t(b);})
  \item \hyperref[enum:widen_m]{\texttt{widen0}}: Stage A widening extractor on \texttt{src0} (promotes FP16 \texttt{h0}/\texttt{h1} to FP32, or sign/zero-extends \texttt{b0..b3}/\texttt{h0..h1}/\texttt{w0..w1} to 32/64-bit before \texttt{abs}/\texttt{neg}). \emph{Example:} {\scriptsize\ttfamily \mnem{IADD.s32} \dst{r0}, \src{r1}\imod{.h0}, \src{r2}} (\textbf{SPIR-V:} \texttt{OpSConvert / OpFConvert}; \textbf{GLSL:} \texttt{int c = int(int16\_t(a)) + b;})
  \item \hyperref[enum:sem_m]{\texttt{sem}}: IEEE-754 NaN handling on \texttt{FMIN}/\texttt{FMAX}: \texttt{nan\_suppress} (minNum/maxNum) vs \texttt{nan\_propagate} (payload-preserving). \emph{Example:} {\scriptsize\ttfamily \mnem{FMAX.f32}\imod{.nan\_propagate} \dst{r0}, \src{r1}, \src{r2}} (\textbf{SPIR-V:} \texttt{GLSL.std.450 NMax / FMax}; \textbf{GLSL:} \texttt{float m = max(a, b);})
  \item \hyperref[enum:clamp_m]{\texttt{clamp}}: Stage B post-op floating-point clamp (\texttt{none}, \texttt{clamp\_0\_inf}, \texttt{clamp\_m1\_1}, \texttt{clamp\_0\_1}) with \texttt{-0.0 \textless{} +0.0} and NaN suppression (\texttt{NaN -\textgreater{} lo}). \emph{Example:} {\scriptsize\ttfamily \mnem{FMA.f32}\imod{.clamp\_0\_1} \dst{r0}, \src{r1}, \src{r2}, \src{r3}} (\textbf{SPIR-V:} \texttt{GLSL.std.450 FClamp}; \textbf{GLSL:} \texttt{float y = clamp(fma(a, b, c), 0.0, 1.0);})
  \item \hyperref[enum:neg_m]{\texttt{neg1}}: Stage A floating-point negation on \texttt{src1} (flips IEEE sign bit after swizzle/widen and \texttt{abs1}). \emph{Example:} {\scriptsize\ttfamily \mnem{FMA.f32} \dst{r0}, \src{r1}, \src{r2}\imod{.neg}, \src{r3}} (\textbf{SPIR-V:} \texttt{OpFNegate}; \textbf{GLSL:} \texttt{float d = fma(a, -b, c);})
  \item \hyperref[enum:abs_m]{\texttt{abs1}}: Stage A floating-point absolute value on \texttt{src1} (clears IEEE sign bit after swizzle/widen, before \texttt{neg1}). \emph{Example:} {\scriptsize\ttfamily \mnem{FMA.f32} \dst{r0}, \src{r1}, \src{r2}\imod{.abs}, \src{r3}} (\textbf{SPIR-V:} \texttt{GLSL.std.450 FAbs}; \textbf{GLSL:} \texttt{float d = fma(a, abs(b), c);})
  \item \hyperref[enum:neg_m]{\texttt{neg0}}: Stage A floating-point negation on \texttt{src0} (flips IEEE sign bit after swizzle/widen and \texttt{abs0}). \emph{Example:} {\scriptsize\ttfamily \mnem{FADD.f32} \dst{r0}, \src{r1}\imod{.neg}, \src{r2}} (\textbf{SPIR-V:} \texttt{OpFNegate}; \textbf{GLSL:} \texttt{float c = -a + b;})
  \item \hyperref[enum:abs_m]{\texttt{abs0}}: Stage A floating-point absolute value on \texttt{src0} (clears IEEE sign bit after swizzle/widen, before \texttt{neg0}). \emph{Example:} {\scriptsize\ttfamily \mnem{FADD.f32} \dst{r0}, \src{r1}\imod{.abs}, \src{r2}} (\textbf{SPIR-V:} \texttt{GLSL.std.450 FAbs}; \textbf{GLSL:} \texttt{float c = abs(a) + b;})
  \item \texttt{dest\_register\_index}: 6-bit destination GPR index (\texttt{r0..r63}) receiving the Stage C writeback result (even-aligned pair \texttt{[rN:rN+1]} for 64-bit ops). \emph{Example:} {\scriptsize\ttfamily \mnem{FMA.f32} \dst{r0}, \src{r1}, \src{r2}, \src{r3}} (\textbf{SPIR-V:} \texttt{SSA Result ID}; \textbf{GLSL:} \texttt{float r0 = fma(r1, r2, r3);})
  \item \hyperref[enum:dest_width_m]{\texttt{dest\_width}}: Stage C destination lane write-mask (\texttt{none} = full 32-bit, \texttt{h0} = low 16-bit, \texttt{h1} = high 16-bit without bit-shifting). \emph{Example:} {\scriptsize\ttfamily \mnem{FADD.v2f16} \dst{r0}\imod{.h1}, \src{r1}, \src{r2}} (\textbf{SPIR-V:} \texttt{OpVectorInsertDynamic / WriteMask}; \textbf{GLSL:} \texttt{dst.y = a.y + b.y;})
  \item \texttt{datatype}: Selects the lane data type (\texttt{f32}, \texttt{f16}, \texttt{v2f16}, \texttt{s32}, \texttt{u32}, \texttt{v2s16}, \texttt{v2u16}, \texttt{v4s8}, \texttt{v4u8}, \texttt{s64}, \texttt{u64}, etc.). \emph{Example:} {\scriptsize\ttfamily \mnem{IADD.v2s16} \dst{r0}, \src{r1}, \src{r2}} (\textbf{SPIR-V:} \texttt{OpTypeInt / OpTypeFloat}; \textbf{GLSL:} \texttt{i16vec2 c = a + b;})
  \item \texttt{instruction\_class}: 4-bit major instruction class selector (\texttt{bits 48:51}) partitioning the 64-bit Valhall opcode space (e.g., \texttt{FMA}, \texttt{CVT}, \texttt{SFU}, \texttt{LS}, \texttt{TEX}, \texttt{BRANCH}). \emph{Example:} {\scriptsize\ttfamily \mnem{FADD.f32} \dst{r0}, \src{r1}, \src{r2}} (\textbf{SPIR-V:} \texttt{N/A (encoding class)}; \textbf{GLSL:} \texttt{// Hardware execution pipe / decode class})
  \item \texttt{fau\_page\_index}: Selects the FAU (Fast Access Uniform) page (\texttt{0..3}) when a source operand references hardware small constants, thread IDs, or lane constants. \emph{Example:} {\scriptsize\ttfamily \mnem{FADD.f32} \dst{r0}, \src{r1}, \src{k1f}} (\textbf{SPIR-V:} \texttt{OpConstant / BuiltIn}; \textbf{GLSL:} \texttt{float y = x + 1.0;})
  \item \hyperref[enum:flow_control_m]{\texttt{flow}}: 4-bit embedded flow-control \& scoreboard wait modifier (\texttt{bits 59:62}: \texttt{none}, \texttt{wait0..wait0126}, \texttt{wait}, \texttt{reconverge}, \texttt{discard}, \texttt{end}). \emph{Example:} {\scriptsize\ttfamily \mnem{FADD.f32}\imod{.wait0.reconverge} \dst{r0}, \src{r1}, \src{r2}} (\textbf{SPIR-V:} \texttt{OpControlBarrier / OpReturn}; \textbf{GLSL:} \texttt{// Scoreboard wait + warp reconvergence})
\end{tightitemize}
\clearpage

\subsection*{FMIN}
\addcontentsline{toc}{subsection}{FMIN}
Computes IEEE-754 minimum independently for each floating-point element $i \in \{0, \dots, N-1\}$ ($N=1$ \texttt{f32}, $N=2$ \texttt{v2f16}) with deterministic \texttt{-0.0 \textless{} +0.0}. Under \texttt{nan\_suppress} (default), a non-NaN wins over NaN; under \texttt{nan\_propagate}, returns the quietened winning NaN preserving sign and payload (\texttt{\textbar{} qbit}). Optional \texttt{clamp} runs after \texttt{sem}.
\par\textbf{Note:} Verified in \texttt{libGLES\_mali.so} \texttt{\_mali\_min\_sf32} (\texttt{0x2e28500}) / \texttt{\_mali\_min\_sf16} (\texttt{0x2e283a0}) and \texttt{kraid/ops.rs:1907-1947}: when both inputs are NaN, even \texttt{nan\_suppress} quietens both (\texttt{\textbar{} qbit}) and picks the larger NaN under unsigned total-order key.
\codebox{\mnem{FMIN}\imod{\{.flow\}}\imod{\{.sem\}}\imod{\{.clamp\}} \dst{dest}\imod{\{.dest\_width\}}, \src{src0}\imod{\{.widen0\}}\imod{\{.abs0\}}\imod{\{.neg0\}}, \src{src1}\imod{\{.widen1\}}\imod{\{.abs1\}}\imod{\{.neg1\}}\\\mnem{FMIN}\imod{\{.flow\}}\imod{\{.sem\}}\imod{\{.clamp\}} \dst{dest}\imod{\{.dest\_width\}}, \src{src0}\imod{\{.swz0\}}\imod{\{.abs0\}}\imod{\{.neg0\}}, \src{src1}\imod{\{.swz1\}}\imod{\{.abs1\}}\imod{\{.neg1\}}}
\textbf{Operands:}
\begin{tightitemize}\footnotesize
  \item \texttt{\dst{dest}} (\texttt{dest\_register\_index}, \texttt{dest\_width}): Destination GPR receiving \texttt{min(src0, src1)} per NaN handling mode.
  \item \texttt{\src{src0}}: First floating-point operand.
  \item \texttt{\src{src1}}: Second floating-point operand.
\end{tightitemize}
\par\textbf{Closest SPIR-V Op:} \texttt{GLSL.std.450 NMin / FMin}
\par\textbf{Pseudo-GLSL (Operational Semantics):}
\codebox{for (int i = 0; i \textless{} N; ++i) \{\\    auto a = mod(src0)[i], b = mod(src1)[i];\\    auto r = flag\_nan\_suppress ? (isnan(a) ? b : (isnan(b) ? a : minTotal(a, b)))\\                               : ((isnan(a) \textbar{}\textbar{} isnan(b)) ? quieten(winNaN(a, b)) : minTotal(a, b));\\    dest[i] = clamp\_mode(r);\\\}}
\textbf{Execution Context:}
\begin{tightitemize}\footnotesize
  \item \textbf{Compute stage preloads (\hyperref[sec:csf-spd]{\texttt{SPD\_k}}, \texttt{r16..r23}):} Hardware pre-fills \texttt{r55..r57} with \texttt{WorkGroupID.xyz + JOB\_OFFSET} (from \hyperref[sec:csf-regs]{\texttt{COMPUTESR::JOB\_OFFSET\_X/Y/Z}} \texttt{r34..r36}) and \texttt{r60..r62} with \texttt{LocalInvocationID.xyz} or \texttt{GlobalInvocationID.xyz} (from \texttt{WG\_SIZE} \texttt{r33} \& \texttt{JOB\_SIZE} \texttt{r37..r39}, selected by \texttt{ShaderProgram::preload}) so compute shaders can derive thread/element indices without memory loads.
  \item \textbf{Vertex stage preloads (\hyperref[sec:csf-spd]{\texttt{SPD\_0}/\texttt{SPD\_1}}, \texttt{r16..r19}):} The hardware Index Fetcher reads \hyperref[sec:csf-regs]{\texttt{IDVSSR::INDEX\_BUFFER}} (\texttt{r54:r55}) and pre-fills \texttt{r59, r60} with \texttt{VertexIndex} (\texttt{+ VERTEX\_OFFSET + GLOBAL\_ATTR\_OFF}) and \texttt{r61} with \texttt{InstanceIndex} (\texttt{+ INSTANCE\_OFFSET}) so Position and Varying shaders can pull vertex attributes directly.
  \item \textbf{Fragment stage preloads \& FP modes (\hyperref[sec:csf-spd]{\texttt{SPD\_2}}, \texttt{r20:r21}):} Hardware pre-fills \texttt{r59} with \texttt{floatBitsToUint(FragCoord.z)} and \texttt{r60} with the MSAA sample coverage mask (\texttt{gl\_SampleMaskIn[0]}), while \texttt{ShaderProgram} word~0 latches 32/64 GPR allocation and warp-wide \texttt{flush\_to\_zero\_mode}, \texttt{suppress\_na\_n}, and \texttt{suppress\_inf}.
  \item \textbf{Inline FAU \& Special Operands (\hyperref[sec:csf-tsd-fau]{\texttt{FAU\_k}}, \texttt{r8..r15}):} Any 8-bit source with \texttt{mode=2} reads zero-latency uniform \texttt{u0..u31} from \texttt{FAU\_k} (\texttt{slots64 = bits[63:56]}); \texttt{mode=3} reads \texttt{SmallConstant} (Table~1.2) or \texttt{FauSpecial} (\texttt{pc}, \texttt{lane\_id}, \texttt{lane\_width}, \texttt{tls\_ptr}, \texttt{wls\_ptr}, \texttt{blend\_descriptor\_0..7}).
\end{tightitemize}
\par\textbf{Worked Examples:}
\par\textbf{Example 1 (Deterministic `-0.0 \textless{} +0.0` \& NaN suppression):}
\codebox{\mnem{FMIN.f32}\imod{.nan\_suppress} \dst{r0}, \src{r1}, \src{r2}}
With \texttt{r1 = -0.0f, r2 = +0.0f}: deterministically returns \texttt{-0.0f} (\texttt{0x80000000}). With \texttt{r1 = NaN, r2 = 1.0f}: returns \texttt{1.0f}.
\par\textbf{Example 2 (NaN propagation + post-op clamp order):}
\codebox{\mnem{FMIN.f32}\imod{.nan\_propagate.clamp\_0\_1} \dst{r0}, \src{r1}, \src{r2}}
With \texttt{r1 = NaN, r2 = 0.5f}: \texttt{sem} runs BEFORE \texttt{clamp} (propagating NaN), then \texttt{.clamp\_0\_1} maps \texttt{NaN -\textgreater{} +0.0f} (\texttt{0x00000000}).
\par\textbf{Mined SPIR-V $\to$ Valhall ISA Example (2x2 textureGather min reduction with NaN suppression --- Fragment shader, Mali-G710 (v10)):}
\codebox{\textbf{; SPIR-V (input)}\\\%g = OpImageGather \%v4float \%tex \%uv \%int\_0 None\\\%m0 = OpExtInst \%float \%glsl FMin \%g.y \%g.x\\\%m1 = OpExtInst \%float \%glsl FMin \%g.z \%g.w\\\%res = OpExtInst \%float \%glsl FMin \%m0 \%m1\\[3pt]\textbf{; Valhall ISA (malioc output)}\\TEX\_GATHER.slot0.wait0.skip.tex2d.dst32.rgba.wide.float\_coordinates.gather4\_r @r0:r1:r2:r3, @r0:r1, u0\\FMIN.f32.nan\_suppress r0, r1\textasciicircum{}, r0\textasciicircum{}\\FMIN.f32.nan\_suppress r2, r2\textasciicircum{}, r3\textasciicircum{}\\FMIN.f32.nan\_suppress r0, r0\textasciicircum{}, r2\textasciicircum{}}
Just like \texttt{FMAX}, \texttt{GLSL.std.450 FMin} uses \texttt{.nan\_suppress} so a single NaN texel does not poison the 2x2 footprint minimum.
\begin{center}\footnotesize
\begin{tabularx}{\linewidth}{@{}lllX@{}}
\toprule
Variant & Arch & Class & Opcode Fields \\
\midrule
\texttt{f32} & v9--v14 & \texttt{R2} & \texttt{opcode2=MIN} \\
\texttt{v2f16} & v9--v14 & \texttt{R2} & \texttt{opcode2=MIN} \\
\bottomrule
\end{tabularx}
\end{center}
\begin{center}\footnotesize
\begin{tabularx}{0.88\linewidth}{@{}l l X@{}}
\toprule
Field Name & Bits & Type / Fixed Value \\
\midrule
\texttt{src0} & \texttt{0:7} & \texttt{SourceEncoding} \\
\texttt{src1} & \texttt{8:15} & \texttt{SourceEncoding} \\
\texttt{opcode2} & \texttt{16:19} & \texttt{exact=MIN} \\
\textit{Reserved} & \texttt{20:25} & \texttt{zero} \\
\hyperref[enum:widen_m]{\texttt{widen1}} & \texttt{26:27} & \texttt{widen\_m (def none)} \\
\hyperref[enum:widen_m]{\texttt{widen0}} & \texttt{28:29} & \texttt{widen\_m (def none)} \\
\hyperref[enum:sem_m]{\texttt{sem}} & \texttt{30:31} & \texttt{sem\_m (def nan\_suppress)} \\
\hyperref[enum:clamp_m]{\texttt{clamp}} & \texttt{32:33} & \texttt{clamp\_m (def none)} \\
\textit{Reserved} & \texttt{34:35} & \texttt{zero} \\
\hyperref[enum:neg_m]{\texttt{neg1}} & \texttt{36} & \texttt{neg\_m (def none)} \\
\hyperref[enum:abs_m]{\texttt{abs1}} & \texttt{37} & \texttt{abs\_m (def none)} \\
\hyperref[enum:neg_m]{\texttt{neg0}} & \texttt{38} & \texttt{neg\_m (def none)} \\
\hyperref[enum:abs_m]{\texttt{abs0}} & \texttt{39} & \texttt{abs\_m (def none)} \\
\texttt{dest\_register\_index} & \texttt{40:45} & \texttt{uint} \\
\hyperref[enum:dest_width_m]{\texttt{dest\_width}} & \texttt{46:47} & \texttt{dest\_width\_m} \\
\texttt{datatype} & \texttt{48:51} & \texttt{exact=f32} \\
\texttt{instruction\_class} & \texttt{52:56} & \texttt{exact=R2} \\
\texttt{fau\_page\_index} & \texttt{57:58} & \texttt{uint} \\
\hyperref[enum:flow_control_m]{\texttt{flow}} & \texttt{59:62} & \texttt{flow\_control\_m} \\
\textit{Reserved} & \texttt{63} & \texttt{zero} \\
\bottomrule
\end{tabularx}
\end{center}
\textbf{Fields Explained:}
\begin{tightitemize}\footnotesize
  \item \texttt{src0}: Primary source operand (\texttt{r0..r63}, \texttt{u0..u63} uniform, \texttt{k*} FAU constant, or \texttt{t0..t2} temporary; \texttt{\textasciicircum{}} marks discard/last-use). \emph{Example:} {\scriptsize\ttfamily \mnem{FADD.f32} \dst{r0}, \src{r1}, \src{r2}} (\textbf{SPIR-V:} \texttt{OpFAdd (Operand 0)}; \textbf{GLSL:} \texttt{float c = a + b;})
  \item \texttt{src1}: Secondary source operand (register, uniform, or FAU constant; in shifts, supplies the shift count byte). \emph{Example:} {\scriptsize\ttfamily \mnem{LSHIFT.i32} \dst{r0}, \src{r1}, \src{r2}\imod{.b0}} (\textbf{SPIR-V:} \texttt{OpShiftLeftLogical (Shift)}; \textbf{GLSL:} \texttt{uint y = a \textless{}\textless{} s;})
  \item \texttt{opcode2}: Secondary sub-opcode further specializing instructions sharing \texttt{instruction\_class} and \texttt{opcode1} (e.g., conversions, SFU tables, moves). \emph{Example:} {\scriptsize\ttfamily \mnem{F32\_TO\_S32}\imod{.round\_zero} \dst{r0}, \src{r1}} (\textbf{SPIR-V:} \texttt{OpConvertFToS}; \textbf{GLSL:} \texttt{int y = int(x);})
  \item \hyperref[enum:widen_m]{\texttt{widen1}}: Stage A widening extractor on \texttt{src1} (promotes narrow byte/half/word lane before Stage B). \emph{Example:} {\scriptsize\ttfamily \mnem{IMUL.u64} \dst{[r0:r1]}, \src{[r2:r3]}\imod{.w0}, \src{[r4:r5]}\imod{.w0}} (\textbf{SPIR-V:} \texttt{OpUConvert / OpUMulExtended}; \textbf{GLSL:} \texttt{uint64\_t p = uint64\_t(a) * uint64\_t(b);})
  \item \hyperref[enum:widen_m]{\texttt{widen0}}: Stage A widening extractor on \texttt{src0} (promotes FP16 \texttt{h0}/\texttt{h1} to FP32, or sign/zero-extends \texttt{b0..b3}/\texttt{h0..h1}/\texttt{w0..w1} to 32/64-bit before \texttt{abs}/\texttt{neg}). \emph{Example:} {\scriptsize\ttfamily \mnem{IADD.s32} \dst{r0}, \src{r1}\imod{.h0}, \src{r2}} (\textbf{SPIR-V:} \texttt{OpSConvert / OpFConvert}; \textbf{GLSL:} \texttt{int c = int(int16\_t(a)) + b;})
  \item \hyperref[enum:sem_m]{\texttt{sem}}: IEEE-754 NaN handling on \texttt{FMIN}/\texttt{FMAX}: \texttt{nan\_suppress} (minNum/maxNum) vs \texttt{nan\_propagate} (payload-preserving). \emph{Example:} {\scriptsize\ttfamily \mnem{FMAX.f32}\imod{.nan\_propagate} \dst{r0}, \src{r1}, \src{r2}} (\textbf{SPIR-V:} \texttt{GLSL.std.450 NMax / FMax}; \textbf{GLSL:} \texttt{float m = max(a, b);})
  \item \hyperref[enum:clamp_m]{\texttt{clamp}}: Stage B post-op floating-point clamp (\texttt{none}, \texttt{clamp\_0\_inf}, \texttt{clamp\_m1\_1}, \texttt{clamp\_0\_1}) with \texttt{-0.0 \textless{} +0.0} and NaN suppression (\texttt{NaN -\textgreater{} lo}). \emph{Example:} {\scriptsize\ttfamily \mnem{FMA.f32}\imod{.clamp\_0\_1} \dst{r0}, \src{r1}, \src{r2}, \src{r3}} (\textbf{SPIR-V:} \texttt{GLSL.std.450 FClamp}; \textbf{GLSL:} \texttt{float y = clamp(fma(a, b, c), 0.0, 1.0);})
  \item \hyperref[enum:neg_m]{\texttt{neg1}}: Stage A floating-point negation on \texttt{src1} (flips IEEE sign bit after swizzle/widen and \texttt{abs1}). \emph{Example:} {\scriptsize\ttfamily \mnem{FMA.f32} \dst{r0}, \src{r1}, \src{r2}\imod{.neg}, \src{r3}} (\textbf{SPIR-V:} \texttt{OpFNegate}; \textbf{GLSL:} \texttt{float d = fma(a, -b, c);})
  \item \hyperref[enum:abs_m]{\texttt{abs1}}: Stage A floating-point absolute value on \texttt{src1} (clears IEEE sign bit after swizzle/widen, before \texttt{neg1}). \emph{Example:} {\scriptsize\ttfamily \mnem{FMA.f32} \dst{r0}, \src{r1}, \src{r2}\imod{.abs}, \src{r3}} (\textbf{SPIR-V:} \texttt{GLSL.std.450 FAbs}; \textbf{GLSL:} \texttt{float d = fma(a, abs(b), c);})
  \item \hyperref[enum:neg_m]{\texttt{neg0}}: Stage A floating-point negation on \texttt{src0} (flips IEEE sign bit after swizzle/widen and \texttt{abs0}). \emph{Example:} {\scriptsize\ttfamily \mnem{FADD.f32} \dst{r0}, \src{r1}\imod{.neg}, \src{r2}} (\textbf{SPIR-V:} \texttt{OpFNegate}; \textbf{GLSL:} \texttt{float c = -a + b;})
  \item \hyperref[enum:abs_m]{\texttt{abs0}}: Stage A floating-point absolute value on \texttt{src0} (clears IEEE sign bit after swizzle/widen, before \texttt{neg0}). \emph{Example:} {\scriptsize\ttfamily \mnem{FADD.f32} \dst{r0}, \src{r1}\imod{.abs}, \src{r2}} (\textbf{SPIR-V:} \texttt{GLSL.std.450 FAbs}; \textbf{GLSL:} \texttt{float c = abs(a) + b;})
  \item \texttt{dest\_register\_index}: 6-bit destination GPR index (\texttt{r0..r63}) receiving the Stage C writeback result (even-aligned pair \texttt{[rN:rN+1]} for 64-bit ops). \emph{Example:} {\scriptsize\ttfamily \mnem{FMA.f32} \dst{r0}, \src{r1}, \src{r2}, \src{r3}} (\textbf{SPIR-V:} \texttt{SSA Result ID}; \textbf{GLSL:} \texttt{float r0 = fma(r1, r2, r3);})
  \item \hyperref[enum:dest_width_m]{\texttt{dest\_width}}: Stage C destination lane write-mask (\texttt{none} = full 32-bit, \texttt{h0} = low 16-bit, \texttt{h1} = high 16-bit without bit-shifting). \emph{Example:} {\scriptsize\ttfamily \mnem{FADD.v2f16} \dst{r0}\imod{.h1}, \src{r1}, \src{r2}} (\textbf{SPIR-V:} \texttt{OpVectorInsertDynamic / WriteMask}; \textbf{GLSL:} \texttt{dst.y = a.y + b.y;})
  \item \texttt{datatype}: Selects the lane data type (\texttt{f32}, \texttt{f16}, \texttt{v2f16}, \texttt{s32}, \texttt{u32}, \texttt{v2s16}, \texttt{v2u16}, \texttt{v4s8}, \texttt{v4u8}, \texttt{s64}, \texttt{u64}, etc.). \emph{Example:} {\scriptsize\ttfamily \mnem{IADD.v2s16} \dst{r0}, \src{r1}, \src{r2}} (\textbf{SPIR-V:} \texttt{OpTypeInt / OpTypeFloat}; \textbf{GLSL:} \texttt{i16vec2 c = a + b;})
  \item \texttt{instruction\_class}: 4-bit major instruction class selector (\texttt{bits 48:51}) partitioning the 64-bit Valhall opcode space (e.g., \texttt{FMA}, \texttt{CVT}, \texttt{SFU}, \texttt{LS}, \texttt{TEX}, \texttt{BRANCH}). \emph{Example:} {\scriptsize\ttfamily \mnem{FADD.f32} \dst{r0}, \src{r1}, \src{r2}} (\textbf{SPIR-V:} \texttt{N/A (encoding class)}; \textbf{GLSL:} \texttt{// Hardware execution pipe / decode class})
  \item \texttt{fau\_page\_index}: Selects the FAU (Fast Access Uniform) page (\texttt{0..3}) when a source operand references hardware small constants, thread IDs, or lane constants. \emph{Example:} {\scriptsize\ttfamily \mnem{FADD.f32} \dst{r0}, \src{r1}, \src{k1f}} (\textbf{SPIR-V:} \texttt{OpConstant / BuiltIn}; \textbf{GLSL:} \texttt{float y = x + 1.0;})
  \item \hyperref[enum:flow_control_m]{\texttt{flow}}: 4-bit embedded flow-control \& scoreboard wait modifier (\texttt{bits 59:62}: \texttt{none}, \texttt{wait0..wait0126}, \texttt{wait}, \texttt{reconverge}, \texttt{discard}, \texttt{end}). \emph{Example:} {\scriptsize\ttfamily \mnem{FADD.f32}\imod{.wait0.reconverge} \dst{r0}, \src{r1}, \src{r2}} (\textbf{SPIR-V:} \texttt{OpControlBarrier / OpReturn}; \textbf{GLSL:} \texttt{// Scoreboard wait + warp reconvergence})
\end{tightitemize}
\clearpage

\subsection*{FMUL}
\addcontentsline{toc}{subsection}{FMUL}
Multiplies floating-point \texttt{src0} and \texttt{src1} independently for each SIMD element $i \in \{0, \dots, N-1\}$ ($N=1$ \texttt{f32}, $N=2$ \texttt{v2f16}): \texttt{dest[i] = clamp(round\_rm(src0[i] * src1[i]))} (hardware alias of \texttt{FMA(src0, src1, -0.0)} so \texttt{(+0.0)*(-0.0)} preserves \texttt{-0.0}).
\par\textbf{Notes:}
\begin{tightitemize}
  \item In \texttt{kraid/nir.rs:896-904} and \texttt{malioc} v10 (\texttt{neg2=neg}), \texttt{FMUL} encodes \texttt{FMA(a, b, -0.0)} so \texttt{a*b + (-0.0)} preserves \texttt{-0.0} under RNE/RTP/RTZ.
  \item Explicit \texttt{neg2=none} adds \texttt{+0.0}, flushing \texttt{-0.0} to \texttt{+0.0}.
  \item \textbf{Source references:} \texttt{kraid/nir.rs:896-904}, \texttt{isa-v9-v14.xml}.
\end{tightitemize}
\codebox{\mnem{FMUL}\imod{\{.flow\}}\imod{\{.round\}}\imod{\{.clamp\}} \dst{dest}\imod{\{.dest\_width\_narrow\}}, \src{src0}\imod{\{.widen0\}}\imod{\{.abs0\}}\imod{\{.neg0\}}, \src{src1}\imod{\{.widen1\}}\imod{\{.abs1\}}\imod{\{.neg1\}}\\\mnem{FMUL}\imod{\{.flow\}}\imod{\{.round\}}\imod{\{.clamp\}} \dst{dest}\imod{\{.dest\_width\}}, \src{src0}\imod{\{.swz0\}}\imod{\{.abs0\}}\imod{\{.neg0\}}, \src{src1}\imod{\{.swz1\}}\imod{\{.abs1\}}\imod{\{.neg1\}}}
\textbf{Operands:}
\begin{tightitemize}\footnotesize
  \item \texttt{\dst{dest}} (\texttt{dest\_register\_index}, \texttt{dest\_width}, \texttt{dest\_width\_narrow}): Destination GPR receiving \texttt{src0 * src1}.
  \item \texttt{\src{src0}}: First multiplicand.
  \item \texttt{\src{src1}}: Second multiplicand.
\end{tightitemize}
\par\textbf{Closest SPIR-V Op:} \texttt{OpFMul}
\par\textbf{Pseudo-GLSL (Operational Semantics):}
\codebox{for (int i = 0; i \textless{} N; ++i)\\    dest[i] = clamp\_mode(round\_rm(fma(mod(src0)[i], mod(src1)[i], -0.0)));}
\textbf{Execution Context:}
\begin{tightitemize}\footnotesize
  \item \textbf{Compute stage preloads (\hyperref[sec:csf-spd]{\texttt{SPD\_k}}, \texttt{r16..r23}):} Hardware pre-fills \texttt{r55..r57} with \texttt{WorkGroupID.xyz + JOB\_OFFSET} (from \hyperref[sec:csf-regs]{\texttt{COMPUTESR::JOB\_OFFSET\_X/Y/Z}} \texttt{r34..r36}) and \texttt{r60..r62} with \texttt{LocalInvocationID.xyz} or \texttt{GlobalInvocationID.xyz} (from \texttt{WG\_SIZE} \texttt{r33} \& \texttt{JOB\_SIZE} \texttt{r37..r39}, selected by \texttt{ShaderProgram::preload}) so compute shaders can derive thread/element indices without memory loads.
  \item \textbf{Vertex stage preloads (\hyperref[sec:csf-spd]{\texttt{SPD\_0}/\texttt{SPD\_1}}, \texttt{r16..r19}):} The hardware Index Fetcher reads \hyperref[sec:csf-regs]{\texttt{IDVSSR::INDEX\_BUFFER}} (\texttt{r54:r55}) and pre-fills \texttt{r59, r60} with \texttt{VertexIndex} (\texttt{+ VERTEX\_OFFSET + GLOBAL\_ATTR\_OFF}) and \texttt{r61} with \texttt{InstanceIndex} (\texttt{+ INSTANCE\_OFFSET}) so Position and Varying shaders can pull vertex attributes directly.
  \item \textbf{Fragment stage preloads \& FP modes (\hyperref[sec:csf-spd]{\texttt{SPD\_2}}, \texttt{r20:r21}):} Hardware pre-fills \texttt{r59} with \texttt{floatBitsToUint(FragCoord.z)} and \texttt{r60} with the MSAA sample coverage mask (\texttt{gl\_SampleMaskIn[0]}), while \texttt{ShaderProgram} word~0 latches 32/64 GPR allocation and warp-wide \texttt{flush\_to\_zero\_mode}, \texttt{suppress\_na\_n}, and \texttt{suppress\_inf}.
  \item \textbf{Inline FAU \& Special Operands (\hyperref[sec:csf-tsd-fau]{\texttt{FAU\_k}}, \texttt{r8..r15}):} Any 8-bit source with \texttt{mode=2} reads zero-latency uniform \texttt{u0..u31} from \texttt{FAU\_k} (\texttt{slots64 = bits[63:56]}); \texttt{mode=3} reads \texttt{SmallConstant} (Table~1.2) or \texttt{FauSpecial} (\texttt{pc}, \texttt{lane\_id}, \texttt{lane\_width}, \texttt{tls\_ptr}, \texttt{wls\_ptr}, \texttt{blend\_descriptor\_0..7}).
\end{tightitemize}
\par\textbf{Worked Examples:}
\par\textbf{Example 1 (Signed-zero preservation (FMA alias)):}
\codebox{\mnem{FMUL.f32} \dst{r0}, \src{r1}, \src{r2}}
Encoded as \texttt{FMA(r1, r2, -0.0)}: with \texttt{r1 = +0.0f, r2 = -0.0f}, produces \texttt{-0.0f} (\texttt{0x80000000}).
\par\textbf{Example 2 (Packed v2f16 multiply with swizzle):}
\codebox{\mnem{FMUL.v2f16} \dst{r0}, \src{r1}\imod{.h01}, \src{r2}\imod{.h00}}
With \texttt{r1 = [3.0h, 2.0h]}, \texttt{r2 = [0.0h, 4.0h]} (\texttt{h0 = 4.0h}): broadcasts \texttt{r2.h0} to both lanes -\textgreater{} \texttt{r0 = [12.0h, 8.0h]}.
\par\textbf{Mined SPIR-V $\to$ Valhall ISA Example (Perspective divide preparation and clamped alpha multiply --- Fragment shader, Mali-G710 (v10)):}
\codebox{\textbf{; SPIR-V (input)}\\\%scaled = OpFMul \%float \%u \%inv\_w\\\%clamped = OpExtInst \%float \%glsl FClamp \%scaled \%float\_0 \%float\_1\\[3pt]\textbf{; Valhall ISA (malioc output)}\\FMUL.f32.wait.clamp\_0\_1 r7, r1\textasciicircum{}, r0\textasciicircum{}}
\texttt{FMUL.f32} is an encoding alias of \texttt{FMA.f32} with addend \texttt{-0.0}, supporting full source modifiers (\texttt{.abs}, \texttt{.neg}) and output \texttt{.clamp\_0\_1}.
\begin{center}\footnotesize
\begin{tabularx}{\linewidth}{@{}lllX@{}}
\toprule
Variant & Arch & Class & Opcode Fields \\
\midrule
\texttt{(none)} & v9--v14 & \texttt{R3} & \texttt{opcode1=FMA\_F32} \\
\texttt{(none)} & v9--v14 & \texttt{R3} & \texttt{opcode1=FMA\_F32} \\
\texttt{(none)} & v9--v14 & \texttt{R3} & \texttt{opcode1=FMA\_V2F16} \\
\texttt{(none)} & v9--v14 & \texttt{R3} & \texttt{opcode1=FMA\_V2F16} \\
\bottomrule
\end{tabularx}
\end{center}
\begin{center}\footnotesize
\begin{tabularx}{0.88\linewidth}{@{}l l X@{}}
\toprule
Field Name & Bits & Type / Fixed Value \\
\midrule
\texttt{src0} & \texttt{0:7} & \texttt{SourceEncoding} \\
\texttt{src1} & \texttt{8:15} & \texttt{SourceEncoding} \\
\texttt{src2} & \texttt{16:23} & \texttt{exact=192} \\
\hyperref[enum:widen_m]{\texttt{widen2}} & \texttt{24:25} & \texttt{exact=none} \\
\hyperref[enum:widen_m]{\texttt{widen1}} & \texttt{26:27} & \texttt{widen\_m (def none)} \\
\hyperref[enum:widen_m]{\texttt{widen0}} & \texttt{28:29} & \texttt{widen\_m (def none)} \\
\hyperref[enum:round_m]{\texttt{round}} & \texttt{30:31} & \texttt{round\_m (def none)} \\
\hyperref[enum:clamp_m]{\texttt{clamp}} & \texttt{32:33} & \texttt{clamp\_m (def none)} \\
\hyperref[enum:neg_m]{\texttt{neg2}} & \texttt{34} & \texttt{neg\_m (def none)} \\
\hyperref[enum:abs_m]{\texttt{abs2}} & \texttt{35} & \texttt{exact=none} \\
\hyperref[enum:neg_m]{\texttt{neg1}} & \texttt{36} & \texttt{neg\_m (def none)} \\
\hyperref[enum:abs_m]{\texttt{abs1}} & \texttt{37} & \texttt{abs\_m (def none)} \\
\hyperref[enum:neg_m]{\texttt{neg0}} & \texttt{38} & \texttt{neg\_m (def none)} \\
\hyperref[enum:abs_m]{\texttt{abs0}} & \texttt{39} & \texttt{abs\_m (def none)} \\
\texttt{dest\_register\_index} & \texttt{40:45} & \texttt{uint} \\
\hyperref[enum:dest_width_narrow_m]{\texttt{dest\_width\_narrow}} & \texttt{46:47} & \texttt{dest\_width\_narrow\_m} \\
\texttt{opcode1} & \texttt{48:51} & \texttt{exact=FMA\_F32} \\
\texttt{instruction\_class} & \texttt{52:56} & \texttt{exact=R3} \\
\texttt{fau\_page\_index} & \texttt{57:58} & \texttt{uint} \\
\hyperref[enum:flow_control_m]{\texttt{flow}} & \texttt{59:62} & \texttt{flow\_control\_m} \\
\textit{Reserved} & \texttt{63} & \texttt{zero} \\
\bottomrule
\end{tabularx}
\end{center}
\textbf{Fields Explained:}
\begin{tightitemize}\footnotesize
  \item \texttt{src0}: Primary source operand (\texttt{r0..r63}, \texttt{u0..u63} uniform, \texttt{k*} FAU constant, or \texttt{t0..t2} temporary; \texttt{\textasciicircum{}} marks discard/last-use). \emph{Example:} {\scriptsize\ttfamily \mnem{FADD.f32} \dst{r0}, \src{r1}, \src{r2}} (\textbf{SPIR-V:} \texttt{OpFAdd (Operand 0)}; \textbf{GLSL:} \texttt{float c = a + b;})
  \item \texttt{src1}: Secondary source operand (register, uniform, or FAU constant; in shifts, supplies the shift count byte). \emph{Example:} {\scriptsize\ttfamily \mnem{LSHIFT.i32} \dst{r0}, \src{r1}, \src{r2}\imod{.b0}} (\textbf{SPIR-V:} \texttt{OpShiftLeftLogical (Shift)}; \textbf{GLSL:} \texttt{uint y = a \textless{}\textless{} s;})
  \item \texttt{src2}: Third source operand (FMA addend, \texttt{MUX}/\texttt{CSEL} selector, \texttt{IADDC} carry-in, or bitwise mask in fused shift-logic). \emph{Example:} {\scriptsize\ttfamily \mnem{FMA.f32} \dst{r0}, \src{r1}, \src{r2}, \src{r3}} (\textbf{SPIR-V:} \texttt{GLSL.std.450 Fma (Addend)}; \textbf{GLSL:} \texttt{float d = fma(a, b, c);})
  \item \hyperref[enum:widen_m]{\texttt{widen2}}: Stage A widening extractor on \texttt{src2} (e.g., promoting FP16 addend \texttt{h0}/\texttt{h1} to FP32 in \texttt{FMA.f32}). \emph{Example:} {\scriptsize\ttfamily \mnem{FMA.f32} \dst{r0}, \src{r1}, \src{r2}, \src{r3}\imod{.h0}} (\textbf{SPIR-V:} \texttt{OpFConvert + GLSL.std.450 Fma}; \textbf{GLSL:} \texttt{float d = fma(a, b, float(c16));})
  \item \hyperref[enum:widen_m]{\texttt{widen1}}: Stage A widening extractor on \texttt{src1} (promotes narrow byte/half/word lane before Stage B). \emph{Example:} {\scriptsize\ttfamily \mnem{IMUL.u64} \dst{[r0:r1]}, \src{[r2:r3]}\imod{.w0}, \src{[r4:r5]}\imod{.w0}} (\textbf{SPIR-V:} \texttt{OpUConvert / OpUMulExtended}; \textbf{GLSL:} \texttt{uint64\_t p = uint64\_t(a) * uint64\_t(b);})
  \item \hyperref[enum:widen_m]{\texttt{widen0}}: Stage A widening extractor on \texttt{src0} (promotes FP16 \texttt{h0}/\texttt{h1} to FP32, or sign/zero-extends \texttt{b0..b3}/\texttt{h0..h1}/\texttt{w0..w1} to 32/64-bit before \texttt{abs}/\texttt{neg}). \emph{Example:} {\scriptsize\ttfamily \mnem{IADD.s32} \dst{r0}, \src{r1}\imod{.h0}, \src{r2}} (\textbf{SPIR-V:} \texttt{OpSConvert / OpFConvert}; \textbf{GLSL:} \texttt{int c = int(int16\_t(a)) + b;})
  \item \hyperref[enum:round_m]{\texttt{round}}: Rounding mode (\texttt{none}/RNE, \texttt{round\_up}, \texttt{round\_down}, \texttt{round\_zero}, \texttt{round\_na}) applied in FP ALU, conversions, \texttt{HADD}, \texttt{LDEXP}, and Stage C narrowing. \emph{Example:} {\scriptsize\ttfamily \mnem{FROUND.f32}\imod{.round\_down} \dst{r0}, \src{r1}} (\textbf{SPIR-V:} \texttt{GLSL.std.450 Floor / Ceil / Trunc}; \textbf{GLSL:} \texttt{float y = floor(x);})
  \item \hyperref[enum:clamp_m]{\texttt{clamp}}: Stage B post-op floating-point clamp (\texttt{none}, \texttt{clamp\_0\_inf}, \texttt{clamp\_m1\_1}, \texttt{clamp\_0\_1}) with \texttt{-0.0 \textless{} +0.0} and NaN suppression (\texttt{NaN -\textgreater{} lo}). \emph{Example:} {\scriptsize\ttfamily \mnem{FMA.f32}\imod{.clamp\_0\_1} \dst{r0}, \src{r1}, \src{r2}, \src{r3}} (\textbf{SPIR-V:} \texttt{GLSL.std.450 FClamp}; \textbf{GLSL:} \texttt{float y = clamp(fma(a, b, c), 0.0, 1.0);})
  \item \hyperref[enum:neg_m]{\texttt{neg2}}: Stage A floating-point negation on \texttt{src2} (flips IEEE sign bit of FMA addend after swizzle/widen and \texttt{abs2}). \emph{Example:} {\scriptsize\ttfamily \mnem{FMA.f32} \dst{r0}, \src{r1}, \src{r2}, \src{r3}\imod{.neg}} (\textbf{SPIR-V:} \texttt{OpFNegate}; \textbf{GLSL:} \texttt{float d = fma(a, b, -c);})
  \item \hyperref[enum:abs_m]{\texttt{abs2}}: Stage A floating-point absolute value on \texttt{src2} (clears IEEE sign bit of FMA addend after swizzle/widen). \emph{Example:} {\scriptsize\ttfamily \mnem{FMA.f32} \dst{r0}, \src{r1}, \src{r2}, \src{r3}\imod{.abs}} (\textbf{SPIR-V:} \texttt{GLSL.std.450 FAbs}; \textbf{GLSL:} \texttt{float d = fma(a, b, abs(c));})
  \item \hyperref[enum:neg_m]{\texttt{neg1}}: Stage A floating-point negation on \texttt{src1} (flips IEEE sign bit after swizzle/widen and \texttt{abs1}). \emph{Example:} {\scriptsize\ttfamily \mnem{FMA.f32} \dst{r0}, \src{r1}, \src{r2}\imod{.neg}, \src{r3}} (\textbf{SPIR-V:} \texttt{OpFNegate}; \textbf{GLSL:} \texttt{float d = fma(a, -b, c);})
  \item \hyperref[enum:abs_m]{\texttt{abs1}}: Stage A floating-point absolute value on \texttt{src1} (clears IEEE sign bit after swizzle/widen, before \texttt{neg1}). \emph{Example:} {\scriptsize\ttfamily \mnem{FMA.f32} \dst{r0}, \src{r1}, \src{r2}\imod{.abs}, \src{r3}} (\textbf{SPIR-V:} \texttt{GLSL.std.450 FAbs}; \textbf{GLSL:} \texttt{float d = fma(a, abs(b), c);})
  \item \hyperref[enum:neg_m]{\texttt{neg0}}: Stage A floating-point negation on \texttt{src0} (flips IEEE sign bit after swizzle/widen and \texttt{abs0}). \emph{Example:} {\scriptsize\ttfamily \mnem{FADD.f32} \dst{r0}, \src{r1}\imod{.neg}, \src{r2}} (\textbf{SPIR-V:} \texttt{OpFNegate}; \textbf{GLSL:} \texttt{float c = -a + b;})
  \item \hyperref[enum:abs_m]{\texttt{abs0}}: Stage A floating-point absolute value on \texttt{src0} (clears IEEE sign bit after swizzle/widen, before \texttt{neg0}). \emph{Example:} {\scriptsize\ttfamily \mnem{FADD.f32} \dst{r0}, \src{r1}\imod{.abs}, \src{r2}} (\textbf{SPIR-V:} \texttt{GLSL.std.450 FAbs}; \textbf{GLSL:} \texttt{float c = abs(a) + b;})
  \item \texttt{dest\_register\_index}: 6-bit destination GPR index (\texttt{r0..r63}) receiving the Stage C writeback result (even-aligned pair \texttt{[rN:rN+1]} for 64-bit ops). \emph{Example:} {\scriptsize\ttfamily \mnem{FMA.f32} \dst{r0}, \src{r1}, \src{r2}, \src{r3}} (\textbf{SPIR-V:} \texttt{SSA Result ID}; \textbf{GLSL:} \texttt{float r0 = fma(r1, r2, r3);})
  \item \hyperref[enum:dest_width_narrow_m]{\texttt{dest\_width\_narrow}}: Destination halfword lane-mask / replication mode (\texttt{h0}, \texttt{h1}, \texttt{h01}, \texttt{none}). \emph{Example:} {\scriptsize\ttfamily \mnem{F32\_TO\_F16} \dst{r0}\imod{.h01}, \src{r1}} (\textbf{SPIR-V:} \texttt{OpFConvert (f32 -\textgreater{} f16)}; \textbf{GLSL:} \texttt{float16\_t h = float16\_t(x);})
  \item \texttt{opcode1}: Primary sub-opcode within the instruction class uniquely identifying the operation family. \emph{Example:} {\scriptsize\ttfamily \mnem{IADD.s32} \dst{r0}, \src{r1}, \src{r2}} (\textbf{SPIR-V:} \texttt{OpIAdd (operation selector)}; \textbf{GLSL:} \texttt{// Primary opcode bits})
  \item \texttt{instruction\_class}: 4-bit major instruction class selector (\texttt{bits 48:51}) partitioning the 64-bit Valhall opcode space (e.g., \texttt{FMA}, \texttt{CVT}, \texttt{SFU}, \texttt{LS}, \texttt{TEX}, \texttt{BRANCH}). \emph{Example:} {\scriptsize\ttfamily \mnem{FADD.f32} \dst{r0}, \src{r1}, \src{r2}} (\textbf{SPIR-V:} \texttt{N/A (encoding class)}; \textbf{GLSL:} \texttt{// Hardware execution pipe / decode class})
  \item \texttt{fau\_page\_index}: Selects the FAU (Fast Access Uniform) page (\texttt{0..3}) when a source operand references hardware small constants, thread IDs, or lane constants. \emph{Example:} {\scriptsize\ttfamily \mnem{FADD.f32} \dst{r0}, \src{r1}, \src{k1f}} (\textbf{SPIR-V:} \texttt{OpConstant / BuiltIn}; \textbf{GLSL:} \texttt{float y = x + 1.0;})
  \item \hyperref[enum:flow_control_m]{\texttt{flow}}: 4-bit embedded flow-control \& scoreboard wait modifier (\texttt{bits 59:62}: \texttt{none}, \texttt{wait0..wait0126}, \texttt{wait}, \texttt{reconverge}, \texttt{discard}, \texttt{end}). \emph{Example:} {\scriptsize\ttfamily \mnem{FADD.f32}\imod{.wait0.reconverge} \dst{r0}, \src{r1}, \src{r2}} (\textbf{SPIR-V:} \texttt{OpControlBarrier / OpReturn}; \textbf{GLSL:} \texttt{// Scoreboard wait + warp reconvergence})
\end{tightitemize}
\clearpage

\subsection*{FNEG}
\addcontentsline{toc}{subsection}{FNEG}
Negates floating-point \texttt{src0} by flipping its IEEE-754 sign bit (\texttt{dest = -src0}), then applies optional \texttt{round}, \texttt{clamp}, and \texttt{dest\_width\_narrow} modifiers. Encoded as an alias of \texttt{FADD(-src0, -0.0)}.
\par\textbf{Notes:}
\begin{tightitemize}
  \item Lowered in \texttt{kraid/nir.rs:776-784} as \texttt{OpFAdd(-0.0, src0.fneg())}.
  \item \textbf{Source references:} \texttt{kraid/nir.rs:776-784}, \texttt{isa-v9-v14.xml}.
\end{tightitemize}
\codebox{\mnem{FNEG}\imod{\{.flow\}}\imod{\{.clamp\}} \dst{dest}\imod{\{.dest\_width\_narrow\}}, \src{src0}\imod{\{.widen0\}}\\\mnem{FNEG}\imod{\{.flow\}}\imod{\{.clamp\}} \dst{dest}\imod{\{.dest\_width\}}, \src{src0}\imod{\{.swz0\}}}
\textbf{Operands:}
\begin{tightitemize}\footnotesize
  \item \texttt{\dst{dest}} (\texttt{dest\_register\_index}, \texttt{dest\_width}, \texttt{dest\_width\_narrow}): Destination GPR receiving \texttt{-src0} (sign bit flipped).
  \item \texttt{\src{src0}}: Floating-point input.
\end{tightitemize}
\par\textbf{Closest SPIR-V Op:} \texttt{OpFNegate}
\par\textbf{Pseudo-GLSL (Operational Semantics):}
\codebox{float v = uintBitsToFloat(floatBitsToUint(src0) \textasciicircum{} 0x80000000u);\\dest = narrowOrMask(dest\_width\_narrow, clamp\_mode(round\_rm(v)));}
\textbf{Execution Context:}
\begin{tightitemize}\footnotesize
  \item \textbf{Compute stage preloads (\hyperref[sec:csf-spd]{\texttt{SPD\_k}}, \texttt{r16..r23}):} Hardware pre-fills \texttt{r55..r57} with \texttt{WorkGroupID.xyz + JOB\_OFFSET} (from \hyperref[sec:csf-regs]{\texttt{COMPUTESR::JOB\_OFFSET\_X/Y/Z}} \texttt{r34..r36}) and \texttt{r60..r62} with \texttt{LocalInvocationID.xyz} or \texttt{GlobalInvocationID.xyz} (from \texttt{WG\_SIZE} \texttt{r33} \& \texttt{JOB\_SIZE} \texttt{r37..r39}, selected by \texttt{ShaderProgram::preload}) so compute shaders can derive thread/element indices without memory loads.
  \item \textbf{Vertex stage preloads (\hyperref[sec:csf-spd]{\texttt{SPD\_0}/\texttt{SPD\_1}}, \texttt{r16..r19}):} The hardware Index Fetcher reads \hyperref[sec:csf-regs]{\texttt{IDVSSR::INDEX\_BUFFER}} (\texttt{r54:r55}) and pre-fills \texttt{r59, r60} with \texttt{VertexIndex} (\texttt{+ VERTEX\_OFFSET + GLOBAL\_ATTR\_OFF}) and \texttt{r61} with \texttt{InstanceIndex} (\texttt{+ INSTANCE\_OFFSET}) so Position and Varying shaders can pull vertex attributes directly.
  \item \textbf{Fragment stage preloads \& FP modes (\hyperref[sec:csf-spd]{\texttt{SPD\_2}}, \texttt{r20:r21}):} Hardware pre-fills \texttt{r59} with \texttt{floatBitsToUint(FragCoord.z)} and \texttt{r60} with the MSAA sample coverage mask (\texttt{gl\_SampleMaskIn[0]}), while \texttt{ShaderProgram} word~0 latches 32/64 GPR allocation and warp-wide \texttt{flush\_to\_zero\_mode}, \texttt{suppress\_na\_n}, and \texttt{suppress\_inf}.
  \item \textbf{Inline FAU \& Special Operands (\hyperref[sec:csf-tsd-fau]{\texttt{FAU\_k}}, \texttt{r8..r15}):} Any 8-bit source with \texttt{mode=2} reads zero-latency uniform \texttt{u0..u31} from \texttt{FAU\_k} (\texttt{slots64 = bits[63:56]}); \texttt{mode=3} reads \texttt{SmallConstant} (Table~1.2) or \texttt{FauSpecial} (\texttt{pc}, \texttt{lane\_id}, \texttt{lane\_width}, \texttt{tls\_ptr}, \texttt{wls\_ptr}, \texttt{blend\_descriptor\_0..7}).
\end{tightitemize}
\par\textbf{Worked Examples:}
\par\textbf{Example 1 (Negate positive float):}
\codebox{\mnem{FNEG.f32} \dst{r0}, \src{r1}}
With \texttt{r1 = 0x3F800000} (\texttt{+1.0f}): flips sign bit -\textgreater{} \texttt{r0 = 0xBF800000} (\texttt{-1.0f}).
\par\textbf{Example 2 (Negate zero and infinity):}
\codebox{\mnem{FNEG.f32} \dst{r0}, \src{r1}}
With \texttt{r1 = +0.0f} (\texttt{0x00000000}) -\textgreater{} \texttt{-0.0f} (\texttt{0x80000000}); with \texttt{r1 = -Inf} (\texttt{0xFF800000}) -\textgreater{} \texttt{+Inf} (\texttt{0x7F800000}).

\begin{center}\footnotesize
\begin{tabularx}{\linewidth}{@{}lllX@{}}
\toprule
Variant & Arch & Class & Opcode Fields \\
\midrule
\texttt{f32} & v9--v14 & \texttt{R2} & \texttt{opcode2=ADD} \\
\texttt{v2f16} & v9--v14 & \texttt{R2} & \texttt{opcode2=ADD} \\
\bottomrule
\end{tabularx}
\end{center}
\begin{center}\footnotesize
\begin{tabularx}{0.88\linewidth}{@{}l l X@{}}
\toprule
Field Name & Bits & Type / Fixed Value \\
\midrule
\texttt{src0} & \texttt{0:7} & \texttt{SourceEncoding} \\
\texttt{src1} & \texttt{8:15} & \texttt{exact=192} \\
\texttt{opcode2} & \texttt{16:19} & \texttt{exact=ADD} \\
\textit{Reserved} & \texttt{20:25} & \texttt{zero} \\
\hyperref[enum:widen_m]{\texttt{widen1}} & \texttt{26:27} & \texttt{exact=none} \\
\hyperref[enum:widen_m]{\texttt{widen0}} & \texttt{28:29} & \texttt{widen\_m (def none)} \\
\hyperref[enum:round_m]{\texttt{round}} & \texttt{30:31} & \texttt{exact=none} \\
\hyperref[enum:clamp_m]{\texttt{clamp}} & \texttt{32:33} & \texttt{clamp\_m (def none)} \\
\hyperref[enum:fadd_sticky_m]{\texttt{sticky}} & \texttt{34} & \texttt{exact=none} \\
\textit{Reserved} & \texttt{35} & \texttt{zero} \\
\hyperref[enum:neg_m]{\texttt{neg1}} & \texttt{36} & \texttt{exact=neg} \\
\hyperref[enum:abs_m]{\texttt{abs1}} & \texttt{37} & \texttt{exact=none} \\
\hyperref[enum:neg_m]{\texttt{neg0}} & \texttt{38} & \texttt{exact=neg} \\
\hyperref[enum:abs_m]{\texttt{abs0}} & \texttt{39} & \texttt{exact=none} \\
\texttt{dest\_register\_index} & \texttt{40:45} & \texttt{uint} \\
\hyperref[enum:dest_width_narrow_m]{\texttt{dest\_width\_narrow}} & \texttt{46:47} & \texttt{dest\_width\_narrow\_m} \\
\texttt{datatype} & \texttt{48:51} & \texttt{exact=f32} \\
\texttt{instruction\_class} & \texttt{52:56} & \texttt{exact=R2} \\
\texttt{fau\_page\_index} & \texttt{57:58} & \texttt{uint} \\
\hyperref[enum:flow_control_m]{\texttt{flow}} & \texttt{59:62} & \texttt{flow\_control\_m} \\
\textit{Reserved} & \texttt{63} & \texttt{zero} \\
\bottomrule
\end{tabularx}
\end{center}
\textbf{Fields Explained:}
\begin{tightitemize}\footnotesize
  \item \texttt{src0}: Primary source operand (\texttt{r0..r63}, \texttt{u0..u63} uniform, \texttt{k*} FAU constant, or \texttt{t0..t2} temporary; \texttt{\textasciicircum{}} marks discard/last-use). \emph{Example:} {\scriptsize\ttfamily \mnem{FADD.f32} \dst{r0}, \src{r1}, \src{r2}} (\textbf{SPIR-V:} \texttt{OpFAdd (Operand 0)}; \textbf{GLSL:} \texttt{float c = a + b;})
  \item \texttt{src1}: Secondary source operand (register, uniform, or FAU constant; in shifts, supplies the shift count byte). \emph{Example:} {\scriptsize\ttfamily \mnem{LSHIFT.i32} \dst{r0}, \src{r1}, \src{r2}\imod{.b0}} (\textbf{SPIR-V:} \texttt{OpShiftLeftLogical (Shift)}; \textbf{GLSL:} \texttt{uint y = a \textless{}\textless{} s;})
  \item \texttt{opcode2}: Secondary sub-opcode further specializing instructions sharing \texttt{instruction\_class} and \texttt{opcode1} (e.g., conversions, SFU tables, moves). \emph{Example:} {\scriptsize\ttfamily \mnem{F32\_TO\_S32}\imod{.round\_zero} \dst{r0}, \src{r1}} (\textbf{SPIR-V:} \texttt{OpConvertFToS}; \textbf{GLSL:} \texttt{int y = int(x);})
  \item \hyperref[enum:widen_m]{\texttt{widen1}}: Stage A widening extractor on \texttt{src1} (promotes narrow byte/half/word lane before Stage B). \emph{Example:} {\scriptsize\ttfamily \mnem{IMUL.u64} \dst{[r0:r1]}, \src{[r2:r3]}\imod{.w0}, \src{[r4:r5]}\imod{.w0}} (\textbf{SPIR-V:} \texttt{OpUConvert / OpUMulExtended}; \textbf{GLSL:} \texttt{uint64\_t p = uint64\_t(a) * uint64\_t(b);})
  \item \hyperref[enum:widen_m]{\texttt{widen0}}: Stage A widening extractor on \texttt{src0} (promotes FP16 \texttt{h0}/\texttt{h1} to FP32, or sign/zero-extends \texttt{b0..b3}/\texttt{h0..h1}/\texttt{w0..w1} to 32/64-bit before \texttt{abs}/\texttt{neg}). \emph{Example:} {\scriptsize\ttfamily \mnem{IADD.s32} \dst{r0}, \src{r1}\imod{.h0}, \src{r2}} (\textbf{SPIR-V:} \texttt{OpSConvert / OpFConvert}; \textbf{GLSL:} \texttt{int c = int(int16\_t(a)) + b;})
  \item \hyperref[enum:round_m]{\texttt{round}}: Rounding mode (\texttt{none}/RNE, \texttt{round\_up}, \texttt{round\_down}, \texttt{round\_zero}, \texttt{round\_na}) applied in FP ALU, conversions, \texttt{HADD}, \texttt{LDEXP}, and Stage C narrowing. \emph{Example:} {\scriptsize\ttfamily \mnem{FROUND.f32}\imod{.round\_down} \dst{r0}, \src{r1}} (\textbf{SPIR-V:} \texttt{GLSL.std.450 Floor / Ceil / Trunc}; \textbf{GLSL:} \texttt{float y = floor(x);})
  \item \hyperref[enum:clamp_m]{\texttt{clamp}}: Stage B post-op floating-point clamp (\texttt{none}, \texttt{clamp\_0\_inf}, \texttt{clamp\_m1\_1}, \texttt{clamp\_0\_1}) with \texttt{-0.0 \textless{} +0.0} and NaN suppression (\texttt{NaN -\textgreater{} lo}). \emph{Example:} {\scriptsize\ttfamily \mnem{FMA.f32}\imod{.clamp\_0\_1} \dst{r0}, \src{r1}, \src{r2}, \src{r3}} (\textbf{SPIR-V:} \texttt{GLSL.std.450 FClamp}; \textbf{GLSL:} \texttt{float y = clamp(fma(a, b, c), 0.0, 1.0);})
  \item \hyperref[enum:fadd_sticky_m]{\texttt{sticky}}: OR-reduces shifted-out rounding bits into LSB (round-to-odd for double-precision emulation). \emph{Example:} {\scriptsize\ttfamily \mnem{FADD.f32}\imod{.sticky} \dst{r0}, \src{r1}, \src{r2}} (\textbf{SPIR-V:} \texttt{OpFAdd (emulated FP64)}; \textbf{GLSL:} \texttt{double c = a + b; // lowered})
  \item \hyperref[enum:neg_m]{\texttt{neg1}}: Stage A floating-point negation on \texttt{src1} (flips IEEE sign bit after swizzle/widen and \texttt{abs1}). \emph{Example:} {\scriptsize\ttfamily \mnem{FMA.f32} \dst{r0}, \src{r1}, \src{r2}\imod{.neg}, \src{r3}} (\textbf{SPIR-V:} \texttt{OpFNegate}; \textbf{GLSL:} \texttt{float d = fma(a, -b, c);})
  \item \hyperref[enum:abs_m]{\texttt{abs1}}: Stage A floating-point absolute value on \texttt{src1} (clears IEEE sign bit after swizzle/widen, before \texttt{neg1}). \emph{Example:} {\scriptsize\ttfamily \mnem{FMA.f32} \dst{r0}, \src{r1}, \src{r2}\imod{.abs}, \src{r3}} (\textbf{SPIR-V:} \texttt{GLSL.std.450 FAbs}; \textbf{GLSL:} \texttt{float d = fma(a, abs(b), c);})
  \item \hyperref[enum:neg_m]{\texttt{neg0}}: Stage A floating-point negation on \texttt{src0} (flips IEEE sign bit after swizzle/widen and \texttt{abs0}). \emph{Example:} {\scriptsize\ttfamily \mnem{FADD.f32} \dst{r0}, \src{r1}\imod{.neg}, \src{r2}} (\textbf{SPIR-V:} \texttt{OpFNegate}; \textbf{GLSL:} \texttt{float c = -a + b;})
  \item \hyperref[enum:abs_m]{\texttt{abs0}}: Stage A floating-point absolute value on \texttt{src0} (clears IEEE sign bit after swizzle/widen, before \texttt{neg0}). \emph{Example:} {\scriptsize\ttfamily \mnem{FADD.f32} \dst{r0}, \src{r1}\imod{.abs}, \src{r2}} (\textbf{SPIR-V:} \texttt{GLSL.std.450 FAbs}; \textbf{GLSL:} \texttt{float c = abs(a) + b;})
  \item \texttt{dest\_register\_index}: 6-bit destination GPR index (\texttt{r0..r63}) receiving the Stage C writeback result (even-aligned pair \texttt{[rN:rN+1]} for 64-bit ops). \emph{Example:} {\scriptsize\ttfamily \mnem{FMA.f32} \dst{r0}, \src{r1}, \src{r2}, \src{r3}} (\textbf{SPIR-V:} \texttt{SSA Result ID}; \textbf{GLSL:} \texttt{float r0 = fma(r1, r2, r3);})
  \item \hyperref[enum:dest_width_narrow_m]{\texttt{dest\_width\_narrow}}: Destination halfword lane-mask / replication mode (\texttt{h0}, \texttt{h1}, \texttt{h01}, \texttt{none}). \emph{Example:} {\scriptsize\ttfamily \mnem{F32\_TO\_F16} \dst{r0}\imod{.h01}, \src{r1}} (\textbf{SPIR-V:} \texttt{OpFConvert (f32 -\textgreater{} f16)}; \textbf{GLSL:} \texttt{float16\_t h = float16\_t(x);})
  \item \texttt{datatype}: Selects the lane data type (\texttt{f32}, \texttt{f16}, \texttt{v2f16}, \texttt{s32}, \texttt{u32}, \texttt{v2s16}, \texttt{v2u16}, \texttt{v4s8}, \texttt{v4u8}, \texttt{s64}, \texttt{u64}, etc.). \emph{Example:} {\scriptsize\ttfamily \mnem{IADD.v2s16} \dst{r0}, \src{r1}, \src{r2}} (\textbf{SPIR-V:} \texttt{OpTypeInt / OpTypeFloat}; \textbf{GLSL:} \texttt{i16vec2 c = a + b;})
  \item \texttt{instruction\_class}: 4-bit major instruction class selector (\texttt{bits 48:51}) partitioning the 64-bit Valhall opcode space (e.g., \texttt{FMA}, \texttt{CVT}, \texttt{SFU}, \texttt{LS}, \texttt{TEX}, \texttt{BRANCH}). \emph{Example:} {\scriptsize\ttfamily \mnem{FADD.f32} \dst{r0}, \src{r1}, \src{r2}} (\textbf{SPIR-V:} \texttt{N/A (encoding class)}; \textbf{GLSL:} \texttt{// Hardware execution pipe / decode class})
  \item \texttt{fau\_page\_index}: Selects the FAU (Fast Access Uniform) page (\texttt{0..3}) when a source operand references hardware small constants, thread IDs, or lane constants. \emph{Example:} {\scriptsize\ttfamily \mnem{FADD.f32} \dst{r0}, \src{r1}, \src{k1f}} (\textbf{SPIR-V:} \texttt{OpConstant / BuiltIn}; \textbf{GLSL:} \texttt{float y = x + 1.0;})
  \item \hyperref[enum:flow_control_m]{\texttt{flow}}: 4-bit embedded flow-control \& scoreboard wait modifier (\texttt{bits 59:62}: \texttt{none}, \texttt{wait0..wait0126}, \texttt{wait}, \texttt{reconverge}, \texttt{discard}, \texttt{end}). \emph{Example:} {\scriptsize\ttfamily \mnem{FADD.f32}\imod{.wait0.reconverge} \dst{r0}, \src{r1}, \src{r2}} (\textbf{SPIR-V:} \texttt{OpControlBarrier / OpReturn}; \textbf{GLSL:} \texttt{// Scoreboard wait + warp reconvergence})
\end{tightitemize}
\clearpage

\subsection*{FPCLASS}
\addcontentsline{toc}{subsection}{FPCLASS}
Classifies each floating-point element of \texttt{src0} (\texttt{f32} or \texttt{f16}) into an IEEE-754 class bitmask written to \texttt{dest}, distinguishing sign, zero, subnormal, normal, infinity, and quiet/signaling NaN.
\par\textbf{Notes:}
\begin{tightitemize}
  \item SFU table unit instruction.
  \item Also used in Bifrost/Valhall transcendental special-case detection.
  \item \textbf{Source references:} \texttt{bifrost/valhall/ISA.xml:2286-2297}, \texttt{bifrost/bi\_lower\_swizzle.c:182}.
\end{tightitemize}
\codebox{\mnem{FPCLASS}\imod{\{.flow\}} \dst{dest}\imod{\{.dest\_width\}}, \src{src0}\imod{\{.lane0\}}\imod{\{.abs0\}}\imod{\{.neg0\}}\\\mnem{FPCLASS}\imod{\{.flow\}} \dst{dest}\imod{\{.dest\_width\}}, \src{src0}\imod{\{.abs0\}}\imod{\{.neg0\}}}
\textbf{Operands:}
\begin{tightitemize}\footnotesize
  \item \texttt{\dst{dest}} (\texttt{dest\_register\_index}, \texttt{dest\_width}): Destination GPR receiving IEEE-754 classification bitmask.
  \item \texttt{\src{src0}}: Floating-point value to classify.
\end{tightitemize}
\par\textbf{Closest SPIR-V Op:} \texttt{OpIsNan / OpIsInf}
\par\textbf{Pseudo-GLSL (Operational Semantics):}
\codebox{dest = classifyIEEE754Bitmask(src0); // bits: sign, zero, subnormal, normal, inf, qnan, snan}
\textbf{Execution Context:}
\begin{tightitemize}\footnotesize
  \item \textbf{Compute stage preloads (\hyperref[sec:csf-spd]{\texttt{SPD\_k}}, \texttt{r16..r23}):} Hardware pre-fills \texttt{r55..r57} with \texttt{WorkGroupID.xyz + JOB\_OFFSET} (from \hyperref[sec:csf-regs]{\texttt{COMPUTESR::JOB\_OFFSET\_X/Y/Z}} \texttt{r34..r36}) and \texttt{r60..r62} with \texttt{LocalInvocationID.xyz} or \texttt{GlobalInvocationID.xyz} (from \texttt{WG\_SIZE} \texttt{r33} \& \texttt{JOB\_SIZE} \texttt{r37..r39}, selected by \texttt{ShaderProgram::preload}) so compute shaders can derive thread/element indices without memory loads.
  \item \textbf{Vertex stage preloads (\hyperref[sec:csf-spd]{\texttt{SPD\_0}/\texttt{SPD\_1}}, \texttt{r16..r19}):} The hardware Index Fetcher reads \hyperref[sec:csf-regs]{\texttt{IDVSSR::INDEX\_BUFFER}} (\texttt{r54:r55}) and pre-fills \texttt{r59, r60} with \texttt{VertexIndex} (\texttt{+ VERTEX\_OFFSET + GLOBAL\_ATTR\_OFF}) and \texttt{r61} with \texttt{InstanceIndex} (\texttt{+ INSTANCE\_OFFSET}) so Position and Varying shaders can pull vertex attributes directly.
  \item \textbf{Fragment stage preloads \& FP modes (\hyperref[sec:csf-spd]{\texttt{SPD\_2}}, \texttt{r20:r21}):} Hardware pre-fills \texttt{r59} with \texttt{floatBitsToUint(FragCoord.z)} and \texttt{r60} with the MSAA sample coverage mask (\texttt{gl\_SampleMaskIn[0]}), while \texttt{ShaderProgram} word~0 latches 32/64 GPR allocation and warp-wide \texttt{flush\_to\_zero\_mode}, \texttt{suppress\_na\_n}, and \texttt{suppress\_inf}.
  \item \textbf{Inline FAU \& Special Operands (\hyperref[sec:csf-tsd-fau]{\texttt{FAU\_k}}, \texttt{r8..r15}):} Any 8-bit source with \texttt{mode=2} reads zero-latency uniform \texttt{u0..u31} from \texttt{FAU\_k} (\texttt{slots64 = bits[63:56]}); \texttt{mode=3} reads \texttt{SmallConstant} (Table~1.2) or \texttt{FauSpecial} (\texttt{pc}, \texttt{lane\_id}, \texttt{lane\_width}, \texttt{tls\_ptr}, \texttt{wls\_ptr}, \texttt{blend\_descriptor\_0..7}).
\end{tightitemize}
\par\textbf{Worked Examples:}
\par\textbf{Example 1 (Classify infinity / NaN):}
\codebox{\mnem{FPCLASS.f32} \dst{r0}, \src{r1}}
With \texttt{r1 = 0x7F800000} (\texttt{+Inf}): sets the hardware positive-infinity class bit in \texttt{r0}.
\par\textbf{Example 2 (Classify subnormal vs zero):}
\codebox{\mnem{FPCLASS.f32} \dst{r0}, \src{r1}}
With \texttt{r1 = 0x80000000} (\texttt{-0.0f}): sets the negative-zero bit distinct from negative-subnormal (\texttt{0x80000001}).

\begin{center}\footnotesize
\begin{tabularx}{\linewidth}{@{}lllX@{}}
\toprule
Variant & Arch & Class & Opcode Fields \\
\midrule
\texttt{(none)} & v9--v14 & \texttt{R1} & \texttt{opcode2=FPCLASS\_F16, opcode1=TABLE} \\
\texttt{(none)} & v9--v14 & \texttt{R1} & \texttt{opcode2=FPCLASS\_F32, opcode1=TABLE} \\
\bottomrule
\end{tabularx}
\end{center}
\begin{center}\footnotesize
\begin{tabularx}{0.88\linewidth}{@{}l l X@{}}
\toprule
Field Name & Bits & Type / Fixed Value \\
\midrule
\texttt{src0} & \texttt{0:7} & \texttt{SourceEncoding} \\
\textit{Reserved} & \texttt{8:15} & \texttt{zero} \\
\texttt{opcode2} & \texttt{16:20} & \texttt{exact=FPCLASS\_F16} \\
\textit{Reserved} & \texttt{21:27} & \texttt{zero} \\
\hyperref[enum:lane16_m]{\texttt{lane0}} & \texttt{28} & \texttt{lane16\_m (def h0)} \\
\textit{Reserved} & \texttt{29:37} & \texttt{zero} \\
\hyperref[enum:neg_m]{\texttt{neg0}} & \texttt{38} & \texttt{neg\_m (def none)} \\
\hyperref[enum:abs_m]{\texttt{abs0}} & \texttt{39} & \texttt{abs\_m (def none)} \\
\texttt{dest\_register\_index} & \texttt{40:45} & \texttt{uint} \\
\hyperref[enum:dest_width_m]{\texttt{dest\_width}} & \texttt{46:47} & \texttt{dest\_width\_m} \\
\texttt{opcode1} & \texttt{48:51} & \texttt{exact=TABLE} \\
\texttt{instruction\_class} & \texttt{52:56} & \texttt{exact=R1} \\
\texttt{fau\_page\_index} & \texttt{57:58} & \texttt{uint} \\
\hyperref[enum:flow_control_m]{\texttt{flow}} & \texttt{59:62} & \texttt{flow\_control\_m} \\
\textit{Reserved} & \texttt{63} & \texttt{zero} \\
\bottomrule
\end{tabularx}
\end{center}
\textbf{Fields Explained:}
\begin{tightitemize}\footnotesize
  \item \texttt{src0}: Primary source operand (\texttt{r0..r63}, \texttt{u0..u63} uniform, \texttt{k*} FAU constant, or \texttt{t0..t2} temporary; \texttt{\textasciicircum{}} marks discard/last-use). \emph{Example:} {\scriptsize\ttfamily \mnem{FADD.f32} \dst{r0}, \src{r1}, \src{r2}} (\textbf{SPIR-V:} \texttt{OpFAdd (Operand 0)}; \textbf{GLSL:} \texttt{float c = a + b;})
  \item \texttt{opcode2}: Secondary sub-opcode further specializing instructions sharing \texttt{instruction\_class} and \texttt{opcode1} (e.g., conversions, SFU tables, moves). \emph{Example:} {\scriptsize\ttfamily \mnem{F32\_TO\_S32}\imod{.round\_zero} \dst{r0}, \src{r1}} (\textbf{SPIR-V:} \texttt{OpConvertFToS}; \textbf{GLSL:} \texttt{int y = int(x);})
  \item \hyperref[enum:lane16_m]{\texttt{lane0}}: Stage A single-lane extractor on \texttt{src0} (\texttt{b0..b3} for 8-bit or \texttt{h0}/\texttt{h1} for 16-bit narrow conversions and \texttt{MKVEC}). \emph{Example:} {\scriptsize\ttfamily \mnem{F16\_TO\_F32} \dst{r0}, \src{r1}\imod{.h1}} (\textbf{SPIR-V:} \texttt{OpCompositeExtract}; \textbf{GLSL:} \texttt{float y = float(x.y);})
  \item \hyperref[enum:neg_m]{\texttt{neg0}}: Stage A floating-point negation on \texttt{src0} (flips IEEE sign bit after swizzle/widen and \texttt{abs0}). \emph{Example:} {\scriptsize\ttfamily \mnem{FADD.f32} \dst{r0}, \src{r1}\imod{.neg}, \src{r2}} (\textbf{SPIR-V:} \texttt{OpFNegate}; \textbf{GLSL:} \texttt{float c = -a + b;})
  \item \hyperref[enum:abs_m]{\texttt{abs0}}: Stage A floating-point absolute value on \texttt{src0} (clears IEEE sign bit after swizzle/widen, before \texttt{neg0}). \emph{Example:} {\scriptsize\ttfamily \mnem{FADD.f32} \dst{r0}, \src{r1}\imod{.abs}, \src{r2}} (\textbf{SPIR-V:} \texttt{GLSL.std.450 FAbs}; \textbf{GLSL:} \texttt{float c = abs(a) + b;})
  \item \texttt{dest\_register\_index}: 6-bit destination GPR index (\texttt{r0..r63}) receiving the Stage C writeback result (even-aligned pair \texttt{[rN:rN+1]} for 64-bit ops). \emph{Example:} {\scriptsize\ttfamily \mnem{FMA.f32} \dst{r0}, \src{r1}, \src{r2}, \src{r3}} (\textbf{SPIR-V:} \texttt{SSA Result ID}; \textbf{GLSL:} \texttt{float r0 = fma(r1, r2, r3);})
  \item \hyperref[enum:dest_width_m]{\texttt{dest\_width}}: Stage C destination lane write-mask (\texttt{none} = full 32-bit, \texttt{h0} = low 16-bit, \texttt{h1} = high 16-bit without bit-shifting). \emph{Example:} {\scriptsize\ttfamily \mnem{FADD.v2f16} \dst{r0}\imod{.h1}, \src{r1}, \src{r2}} (\textbf{SPIR-V:} \texttt{OpVectorInsertDynamic / WriteMask}; \textbf{GLSL:} \texttt{dst.y = a.y + b.y;})
  \item \texttt{opcode1}: Primary sub-opcode within the instruction class uniquely identifying the operation family. \emph{Example:} {\scriptsize\ttfamily \mnem{IADD.s32} \dst{r0}, \src{r1}, \src{r2}} (\textbf{SPIR-V:} \texttt{OpIAdd (operation selector)}; \textbf{GLSL:} \texttt{// Primary opcode bits})
  \item \texttt{instruction\_class}: 4-bit major instruction class selector (\texttt{bits 48:51}) partitioning the 64-bit Valhall opcode space (e.g., \texttt{FMA}, \texttt{CVT}, \texttt{SFU}, \texttt{LS}, \texttt{TEX}, \texttt{BRANCH}). \emph{Example:} {\scriptsize\ttfamily \mnem{FADD.f32} \dst{r0}, \src{r1}, \src{r2}} (\textbf{SPIR-V:} \texttt{N/A (encoding class)}; \textbf{GLSL:} \texttt{// Hardware execution pipe / decode class})
  \item \texttt{fau\_page\_index}: Selects the FAU (Fast Access Uniform) page (\texttt{0..3}) when a source operand references hardware small constants, thread IDs, or lane constants. \emph{Example:} {\scriptsize\ttfamily \mnem{FADD.f32} \dst{r0}, \src{r1}, \src{k1f}} (\textbf{SPIR-V:} \texttt{OpConstant / BuiltIn}; \textbf{GLSL:} \texttt{float y = x + 1.0;})
  \item \hyperref[enum:flow_control_m]{\texttt{flow}}: 4-bit embedded flow-control \& scoreboard wait modifier (\texttt{bits 59:62}: \texttt{none}, \texttt{wait0..wait0126}, \texttt{wait}, \texttt{reconverge}, \texttt{discard}, \texttt{end}). \emph{Example:} {\scriptsize\ttfamily \mnem{FADD.f32}\imod{.wait0.reconverge} \dst{r0}, \src{r1}, \src{r2}} (\textbf{SPIR-V:} \texttt{OpControlBarrier / OpReturn}; \textbf{GLSL:} \texttt{// Scoreboard wait + warp reconvergence})
\end{tightitemize}
\clearpage

\subsection*{FPOW}
\addcontentsline{toc}{subsection}{FPOW}
Computes the FP16 power function \texttt{dest = pow(src0, src1)} directly on Valhall v12+ for scalar half-precision inputs selected by \texttt{lane0}/\texttt{lane1}, writing to \texttt{dest} per \texttt{dest\_width\_replicate}.
\par\textbf{Notes:}
\begin{tightitemize}
  \item Arch v12-v14 only.
  \item \texttt{f32} pow is lowered via \texttt{flog2\_32} + \texttt{fexp\_32\_to} (\texttt{kraid/nir.rs:943-954}).
  \item \textbf{Source references:} \texttt{isa-v9-v14.xml}, \texttt{kraid/nir.rs:943-954}.
\end{tightitemize}
\codebox{\mnem{FPOW}\imod{\{.flow\}} \dst{dest}\imod{\{.dest\_width\_replicate\}}, \src{src0}\imod{\{.lane0\}}\imod{\{.abs0\}}\imod{\{.neg0\}}, \src{src1}\imod{\{.lane1\}}\imod{\{.abs1\}}\imod{\{.neg1\}}}
\textbf{Operands:}
\begin{tightitemize}\footnotesize
  \item \texttt{\dst{dest}} (\texttt{dest\_register\_index}, \texttt{dest\_width\_replicate}): Destination GPR receiving \texttt{pow(src0, src1)} in FP16.
  \item \texttt{\src{src0}}: FP16 base operand.
  \item \texttt{\src{src1}}: FP16 exponent operand.
\end{tightitemize}
\par\textbf{Closest SPIR-V Op:} \texttt{GLSL.std.450 Pow (f16, v12+)}
\par\textbf{Pseudo-GLSL (Operational Semantics):}
\codebox{dest = replicate(dest\_width\_replicate, pow(lane0(src0), lane1(src1))); // v12+ f16}
\textbf{Execution Context:}
\begin{tightitemize}\footnotesize
  \item \textbf{Compute stage preloads (\hyperref[sec:csf-spd]{\texttt{SPD\_k}}, \texttt{r16..r23}):} Hardware pre-fills \texttt{r55..r57} with \texttt{WorkGroupID.xyz + JOB\_OFFSET} (from \hyperref[sec:csf-regs]{\texttt{COMPUTESR::JOB\_OFFSET\_X/Y/Z}} \texttt{r34..r36}) and \texttt{r60..r62} with \texttt{LocalInvocationID.xyz} or \texttt{GlobalInvocationID.xyz} (from \texttt{WG\_SIZE} \texttt{r33} \& \texttt{JOB\_SIZE} \texttt{r37..r39}, selected by \texttt{ShaderProgram::preload}) so compute shaders can derive thread/element indices without memory loads.
  \item \textbf{Vertex stage preloads (\hyperref[sec:csf-spd]{\texttt{SPD\_0}/\texttt{SPD\_1}}, \texttt{r16..r19}):} The hardware Index Fetcher reads \hyperref[sec:csf-regs]{\texttt{IDVSSR::INDEX\_BUFFER}} (\texttt{r54:r55}) and pre-fills \texttt{r59, r60} with \texttt{VertexIndex} (\texttt{+ VERTEX\_OFFSET + GLOBAL\_ATTR\_OFF}) and \texttt{r61} with \texttt{InstanceIndex} (\texttt{+ INSTANCE\_OFFSET}) so Position and Varying shaders can pull vertex attributes directly.
  \item \textbf{Fragment stage preloads \& FP modes (\hyperref[sec:csf-spd]{\texttt{SPD\_2}}, \texttt{r20:r21}):} Hardware pre-fills \texttt{r59} with \texttt{floatBitsToUint(FragCoord.z)} and \texttt{r60} with the MSAA sample coverage mask (\texttt{gl\_SampleMaskIn[0]}), while \texttt{ShaderProgram} word~0 latches 32/64 GPR allocation and warp-wide \texttt{flush\_to\_zero\_mode}, \texttt{suppress\_na\_n}, and \texttt{suppress\_inf}.
  \item \textbf{Inline FAU \& Special Operands (\hyperref[sec:csf-tsd-fau]{\texttt{FAU\_k}}, \texttt{r8..r15}):} Any 8-bit source with \texttt{mode=2} reads zero-latency uniform \texttt{u0..u31} from \texttt{FAU\_k} (\texttt{slots64 = bits[63:56]}); \texttt{mode=3} reads \texttt{SmallConstant} (Table~1.2) or \texttt{FauSpecial} (\texttt{pc}, \texttt{lane\_id}, \texttt{lane\_width}, \texttt{tls\_ptr}, \texttt{wls\_ptr}, \texttt{blend\_descriptor\_0..7}).
\end{tightitemize}
\par\textbf{Worked Examples:}
\par\textbf{Example 1 (Scalar FP16 power with .h01 replication):}
\codebox{\mnem{FPOW.f16} \dst{r0}\imod{.h01}, \src{r1}\imod{.h0}, \src{r2}\imod{.h0}}
With \texttt{r1.h0 = 2.0h, r2.h0 = 3.0h}: computes \texttt{2\textasciicircum{}3 = 8.0h} (\texttt{0x4800}) and \texttt{.h01} replicates -\textgreater{} \texttt{r0 = 0x48004800}.
\par\textbf{Example 2 (Zero and negative base edge cases):}
\codebox{\mnem{FPOW.f16} \dst{r0}\imod{.h0}, \src{r1}\imod{.h0}, \src{r2}\imod{.h0}}
With \texttt{r1.h0 = 4.0h, r2.h0 = 0.5h}: returns \texttt{2.0h} (\texttt{0x4000}); negative base with non-integer exponent returns \texttt{qNaN}.

\begin{center}\footnotesize
\begin{tabularx}{\linewidth}{@{}lllX@{}}
\toprule
Variant & Arch & Class & Opcode Fields \\
\midrule
\texttt{f16} & v9--v14 & \texttt{R2} & \texttt{opcode2=FPOW, opcode3=FPOW} \\
\bottomrule
\end{tabularx}
\end{center}
\begin{center}\footnotesize
\begin{tabularx}{0.88\linewidth}{@{}l l X@{}}
\toprule
Field Name & Bits & Type / Fixed Value \\
\midrule
\texttt{src0} & \texttt{0:7} & \texttt{SourceEncoding} \\
\texttt{src1} & \texttt{8:15} & \texttt{SourceEncoding} \\
\texttt{opcode2} & \texttt{16:19} & \texttt{exact=FPOW} \\
\textit{Reserved} & \texttt{20:23} & \texttt{zero} \\
\texttt{opcode3} & \texttt{24:25} & \texttt{exact=FPOW} \\
\hyperref[enum:lane16_m]{\texttt{lane1}} & \texttt{26} & \texttt{lane16\_m (def h0)} \\
\textit{Reserved} & \texttt{27} & \texttt{zero} \\
\hyperref[enum:lane16_m]{\texttt{lane0}} & \texttt{28} & \texttt{lane16\_m (def h0)} \\
\textit{Reserved} & \texttt{29:35} & \texttt{zero} \\
\hyperref[enum:neg_m]{\texttt{neg1}} & \texttt{36} & \texttt{neg\_m (def none)} \\
\hyperref[enum:abs_m]{\texttt{abs1}} & \texttt{37} & \texttt{abs\_m (def none)} \\
\hyperref[enum:neg_m]{\texttt{neg0}} & \texttt{38} & \texttt{neg\_m (def none)} \\
\hyperref[enum:abs_m]{\texttt{abs0}} & \texttt{39} & \texttt{abs\_m (def none)} \\
\texttt{dest\_register\_index} & \texttt{40:45} & \texttt{uint} \\
\hyperref[enum:dest_width_replicate_m]{\texttt{dest\_width\_replicate}} & \texttt{46:47} & \texttt{dest\_width\_replicate\_m} \\
\texttt{datatype} & \texttt{48:51} & \texttt{exact=f16} \\
\texttt{instruction\_class} & \texttt{52:56} & \texttt{exact=R2} \\
\texttt{fau\_page\_index} & \texttt{57:58} & \texttt{uint} \\
\hyperref[enum:flow_control_m]{\texttt{flow}} & \texttt{59:62} & \texttt{flow\_control\_m} \\
\textit{Reserved} & \texttt{63} & \texttt{zero} \\
\bottomrule
\end{tabularx}
\end{center}
\textbf{Fields Explained:}
\begin{tightitemize}\footnotesize
  \item \texttt{src0}: Primary source operand (\texttt{r0..r63}, \texttt{u0..u63} uniform, \texttt{k*} FAU constant, or \texttt{t0..t2} temporary; \texttt{\textasciicircum{}} marks discard/last-use). \emph{Example:} {\scriptsize\ttfamily \mnem{FADD.f32} \dst{r0}, \src{r1}, \src{r2}} (\textbf{SPIR-V:} \texttt{OpFAdd (Operand 0)}; \textbf{GLSL:} \texttt{float c = a + b;})
  \item \texttt{src1}: Secondary source operand (register, uniform, or FAU constant; in shifts, supplies the shift count byte). \emph{Example:} {\scriptsize\ttfamily \mnem{LSHIFT.i32} \dst{r0}, \src{r1}, \src{r2}\imod{.b0}} (\textbf{SPIR-V:} \texttt{OpShiftLeftLogical (Shift)}; \textbf{GLSL:} \texttt{uint y = a \textless{}\textless{} s;})
  \item \texttt{opcode2}: Secondary sub-opcode further specializing instructions sharing \texttt{instruction\_class} and \texttt{opcode1} (e.g., conversions, SFU tables, moves). \emph{Example:} {\scriptsize\ttfamily \mnem{F32\_TO\_S32}\imod{.round\_zero} \dst{r0}, \src{r1}} (\textbf{SPIR-V:} \texttt{OpConvertFToS}; \textbf{GLSL:} \texttt{int y = int(x);})
  \item \texttt{opcode3}: Tertiary sub-opcode disambiguating specialized multi-op families (\texttt{FPOW\_SC\_*}, SFU lookup tables, \texttt{ICMP\_MULTI}). \emph{Example:} {\scriptsize\ttfamily \mnem{FPOW\_SC\_DET.f32}\imod{.pow} \dst{r0}, \src{r1}, \src{r2}} (\textbf{SPIR-V:} \texttt{GLSL.std.450 Pow}; \textbf{GLSL:} \texttt{float z = pow(x, y);})
  \item \hyperref[enum:lane16_m]{\texttt{lane1}}: Stage A single-lane extractor on \texttt{src1} (\texttt{b0..b3} or \texttt{h0}/\texttt{h1}) for \texttt{MKVEC} packing and narrow binary ops. \emph{Example:} {\scriptsize\ttfamily \mnem{MKVEC.v2i8} \dst{r0}, \src{r1}\imod{.b0}, \src{r2}\imod{.b1}, \src{r3}\imod{.b01}} (\textbf{SPIR-V:} \texttt{OpCompositeConstruct}; \textbf{GLSL:} \texttt{u8vec4 v = u8vec4(a.x, b.y, c.x, c.y);})
  \item \hyperref[enum:lane16_m]{\texttt{lane0}}: Stage A single-lane extractor on \texttt{src0} (\texttt{b0..b3} for 8-bit or \texttt{h0}/\texttt{h1} for 16-bit narrow conversions and \texttt{MKVEC}). \emph{Example:} {\scriptsize\ttfamily \mnem{F16\_TO\_F32} \dst{r0}, \src{r1}\imod{.h1}} (\textbf{SPIR-V:} \texttt{OpCompositeExtract}; \textbf{GLSL:} \texttt{float y = float(x.y);})
  \item \hyperref[enum:neg_m]{\texttt{neg1}}: Stage A floating-point negation on \texttt{src1} (flips IEEE sign bit after swizzle/widen and \texttt{abs1}). \emph{Example:} {\scriptsize\ttfamily \mnem{FMA.f32} \dst{r0}, \src{r1}, \src{r2}\imod{.neg}, \src{r3}} (\textbf{SPIR-V:} \texttt{OpFNegate}; \textbf{GLSL:} \texttt{float d = fma(a, -b, c);})
  \item \hyperref[enum:abs_m]{\texttt{abs1}}: Stage A floating-point absolute value on \texttt{src1} (clears IEEE sign bit after swizzle/widen, before \texttt{neg1}). \emph{Example:} {\scriptsize\ttfamily \mnem{FMA.f32} \dst{r0}, \src{r1}, \src{r2}\imod{.abs}, \src{r3}} (\textbf{SPIR-V:} \texttt{GLSL.std.450 FAbs}; \textbf{GLSL:} \texttt{float d = fma(a, abs(b), c);})
  \item \hyperref[enum:neg_m]{\texttt{neg0}}: Stage A floating-point negation on \texttt{src0} (flips IEEE sign bit after swizzle/widen and \texttt{abs0}). \emph{Example:} {\scriptsize\ttfamily \mnem{FADD.f32} \dst{r0}, \src{r1}\imod{.neg}, \src{r2}} (\textbf{SPIR-V:} \texttt{OpFNegate}; \textbf{GLSL:} \texttt{float c = -a + b;})
  \item \hyperref[enum:abs_m]{\texttt{abs0}}: Stage A floating-point absolute value on \texttt{src0} (clears IEEE sign bit after swizzle/widen, before \texttt{neg0}). \emph{Example:} {\scriptsize\ttfamily \mnem{FADD.f32} \dst{r0}, \src{r1}\imod{.abs}, \src{r2}} (\textbf{SPIR-V:} \texttt{GLSL.std.450 FAbs}; \textbf{GLSL:} \texttt{float c = abs(a) + b;})
  \item \texttt{dest\_register\_index}: 6-bit destination GPR index (\texttt{r0..r63}) receiving the Stage C writeback result (even-aligned pair \texttt{[rN:rN+1]} for 64-bit ops). \emph{Example:} {\scriptsize\ttfamily \mnem{FMA.f32} \dst{r0}, \src{r1}, \src{r2}, \src{r3}} (\textbf{SPIR-V:} \texttt{SSA Result ID}; \textbf{GLSL:} \texttt{float r0 = fma(r1, r2, r3);})
  \item \hyperref[enum:dest_width_replicate_m]{\texttt{dest\_width\_replicate}}: Destination halfword lane-mask / replication mode (\texttt{h0}, \texttt{h1}, \texttt{h01}, \texttt{none}). \emph{Example:} {\scriptsize\ttfamily \mnem{F32\_TO\_F16} \dst{r0}\imod{.h01}, \src{r1}} (\textbf{SPIR-V:} \texttt{OpFConvert (f32 -\textgreater{} f16)}; \textbf{GLSL:} \texttt{float16\_t h = float16\_t(x);})
  \item \texttt{datatype}: Selects the lane data type (\texttt{f32}, \texttt{f16}, \texttt{v2f16}, \texttt{s32}, \texttt{u32}, \texttt{v2s16}, \texttt{v2u16}, \texttt{v4s8}, \texttt{v4u8}, \texttt{s64}, \texttt{u64}, etc.). \emph{Example:} {\scriptsize\ttfamily \mnem{IADD.v2s16} \dst{r0}, \src{r1}, \src{r2}} (\textbf{SPIR-V:} \texttt{OpTypeInt / OpTypeFloat}; \textbf{GLSL:} \texttt{i16vec2 c = a + b;})
  \item \texttt{instruction\_class}: 4-bit major instruction class selector (\texttt{bits 48:51}) partitioning the 64-bit Valhall opcode space (e.g., \texttt{FMA}, \texttt{CVT}, \texttt{SFU}, \texttt{LS}, \texttt{TEX}, \texttt{BRANCH}). \emph{Example:} {\scriptsize\ttfamily \mnem{FADD.f32} \dst{r0}, \src{r1}, \src{r2}} (\textbf{SPIR-V:} \texttt{N/A (encoding class)}; \textbf{GLSL:} \texttt{// Hardware execution pipe / decode class})
  \item \texttt{fau\_page\_index}: Selects the FAU (Fast Access Uniform) page (\texttt{0..3}) when a source operand references hardware small constants, thread IDs, or lane constants. \emph{Example:} {\scriptsize\ttfamily \mnem{FADD.f32} \dst{r0}, \src{r1}, \src{k1f}} (\textbf{SPIR-V:} \texttt{OpConstant / BuiltIn}; \textbf{GLSL:} \texttt{float y = x + 1.0;})
  \item \hyperref[enum:flow_control_m]{\texttt{flow}}: 4-bit embedded flow-control \& scoreboard wait modifier (\texttt{bits 59:62}: \texttt{none}, \texttt{wait0..wait0126}, \texttt{wait}, \texttt{reconverge}, \texttt{discard}, \texttt{end}). \emph{Example:} {\scriptsize\ttfamily \mnem{FADD.f32}\imod{.wait0.reconverge} \dst{r0}, \src{r1}, \src{r2}} (\textbf{SPIR-V:} \texttt{OpControlBarrier / OpReturn}; \textbf{GLSL:} \texttt{// Scoreboard wait + warp reconvergence})
\end{tightitemize}
\clearpage

\subsection*{FRCP}
\addcontentsline{toc}{subsection}{FRCP}
Computes the floating-point reciprocal \texttt{dest = 1.0 / src0} (\texttt{f32} or \texttt{f16}) on the SFU with \texttt{\textless{} 1 ULP} accuracy. On \texttt{f32}, subnormal inputs and subnormal outputs are flushed to signed zero (\texttt{$\pm$0.0}); on \texttt{f16}, subnormals are preserved.
\par\textbf{Notes:}
\begin{tightitemize}
  \item Hardware precision is \texttt{\textless{} 1 ULP} (\texttt{hw\_tests.rs:1697}).
  \item No \texttt{round} or \texttt{clamp} modifier.
  \item \textbf{Source references:} \texttt{kraid/ops.rs:1952-1976}, \texttt{encode\_v9.rs:1581-1598}, \texttt{hw\_tests.rs:1688-1700}.
\end{tightitemize}
\codebox{\mnem{FRCP}\imod{\{.flow\}} \dst{dest}\imod{\{.dest\_width\}}, \src{src0}\imod{\{.lane0\}}\imod{\{.abs0\}}\imod{\{.neg0\}}\\\mnem{FRCP}\imod{\{.flow\}} \dst{dest}\imod{\{.dest\_width\_replicate\}}, \src{src0}\imod{\{.lane0\}}\imod{\{.abs0\}}\imod{\{.neg0\}}\\\mnem{FRCP}\imod{\{.flow\}} \dst{dest}\imod{\{.dest\_width\}}, \src{src0}\imod{\{.widen0\}}\imod{\{.abs0\}}\imod{\{.neg0\}}}
\textbf{Operands:}
\begin{tightitemize}\footnotesize
  \item \texttt{\dst{dest}} (\texttt{dest\_register\_index}, \texttt{dest\_width}, \texttt{dest\_width\_replicate}): Destination GPR receiving reciprocal \texttt{1.0 / src0}.
  \item \texttt{\src{src0}}: Floating-point divisor (subnormals flush to \texttt{0.0} on \texttt{f32}).
\end{tightitemize}
\par\textbf{Closest SPIR-V Op:} \texttt{OpFDiv (`1.0 / x`)}
\par\textbf{Pseudo-GLSL (Operational Semantics):}
\codebox{if (is\_f32) dest = ftz(1.0 / ftz(src0)); // \textless{} 1 ULP, flushes subnormals\\else        dest = float16\_t(1.0h / lane0(src0)); // preserves f16 subnormals}
\textbf{Execution Context:}
\begin{tightitemize}\footnotesize
  \item \textbf{Compute stage preloads (\hyperref[sec:csf-spd]{\texttt{SPD\_k}}, \texttt{r16..r23}):} Hardware pre-fills \texttt{r55..r57} with \texttt{WorkGroupID.xyz + JOB\_OFFSET} (from \hyperref[sec:csf-regs]{\texttt{COMPUTESR::JOB\_OFFSET\_X/Y/Z}} \texttt{r34..r36}) and \texttt{r60..r62} with \texttt{LocalInvocationID.xyz} or \texttt{GlobalInvocationID.xyz} (from \texttt{WG\_SIZE} \texttt{r33} \& \texttt{JOB\_SIZE} \texttt{r37..r39}, selected by \texttt{ShaderProgram::preload}) so compute shaders can derive thread/element indices without memory loads.
  \item \textbf{Vertex stage preloads (\hyperref[sec:csf-spd]{\texttt{SPD\_0}/\texttt{SPD\_1}}, \texttt{r16..r19}):} The hardware Index Fetcher reads \hyperref[sec:csf-regs]{\texttt{IDVSSR::INDEX\_BUFFER}} (\texttt{r54:r55}) and pre-fills \texttt{r59, r60} with \texttt{VertexIndex} (\texttt{+ VERTEX\_OFFSET + GLOBAL\_ATTR\_OFF}) and \texttt{r61} with \texttt{InstanceIndex} (\texttt{+ INSTANCE\_OFFSET}) so Position and Varying shaders can pull vertex attributes directly.
  \item \textbf{Fragment stage preloads \& FP modes (\hyperref[sec:csf-spd]{\texttt{SPD\_2}}, \texttt{r20:r21}):} Hardware pre-fills \texttt{r59} with \texttt{floatBitsToUint(FragCoord.z)} and \texttt{r60} with the MSAA sample coverage mask (\texttt{gl\_SampleMaskIn[0]}), while \texttt{ShaderProgram} word~0 latches 32/64 GPR allocation and warp-wide \texttt{flush\_to\_zero\_mode}, \texttt{suppress\_na\_n}, and \texttt{suppress\_inf}.
  \item \textbf{Inline FAU \& Special Operands (\hyperref[sec:csf-tsd-fau]{\texttt{FAU\_k}}, \texttt{r8..r15}):} Any 8-bit source with \texttt{mode=2} reads zero-latency uniform \texttt{u0..u31} from \texttt{FAU\_k} (\texttt{slots64 = bits[63:56]}); \texttt{mode=3} reads \texttt{SmallConstant} (Table~1.2) or \texttt{FauSpecial} (\texttt{pc}, \texttt{lane\_id}, \texttt{lane\_width}, \texttt{tls\_ptr}, \texttt{wls\_ptr}, \texttt{blend\_descriptor\_0..7}).
\end{tightitemize}
\par\textbf{Worked Examples:}
\par\textbf{Example 1 (Normal FP32 reciprocal):}
\codebox{\mnem{FRCP.f32} \dst{r0}, \src{r1}}
With \texttt{r1 = 4.0f}: returns \texttt{0.25f} (\texttt{0x3E800000}, \texttt{\textless{} 1 ULP}).
\par\textbf{Example 2 (Subnormal FTZ \& signed zero):}
\codebox{\mnem{FRCP.f32} \dst{r0}, \src{r1}}
With \texttt{r1 = -0.0f} (or any negative subnormal, flushed to \texttt{-0.0f} on input): returns \texttt{-Inf} (\texttt{0xFF800000}). With \texttt{r1 = +Inf}: returns \texttt{+0.0f}.
\par\textbf{Mined SPIR-V $\to$ Valhall ISA Example (Perspective division by interpolated 1/w --- Fragment shader, Mali-G710 (v10)):}
\codebox{\textbf{; SPIR-V (input)}\\\%w    = OpLoad \%float \%frag\_w\\\%invw = OpFDiv \%float \%float\_1 \%w\\\%u    = OpFMul \%float \%u\_times\_w \%invw\\[3pt]\textbf{; Valhall ISA (malioc output)}\\LD\_VAR\_SPECIAL.slot0.wait0.frag\_w.f32.sample.clobber @r2, r61\textasciicircum{}\\FRCP.f32 r3, r2\\FMUL.f32 r1, r3, r1\textasciicircum{}}
Floating-point reciprocal \texttt{1.0 / w} lowers directly to a single SFU \texttt{FRCP.f32} instruction (\textless{} 1 ULP accuracy, with FTZ on \texttt{f32}).
\begin{center}\footnotesize
\begin{tabularx}{\linewidth}{@{}lllX@{}}
\toprule
Variant & Arch & Class & Opcode Fields \\
\midrule
\texttt{(none)} & v9--v14 & \texttt{R1} & \texttt{opcode2=FRCP\_F16, opcode1=TABLE} \\
\texttt{(none)} & v9--v14 & \texttt{R1} & \texttt{opcode2=FRCP\_F16, opcode1=TABLE} \\
\texttt{(none)} & v9--v14 & \texttt{R1} & \texttt{opcode2=FRCP\_F32, opcode1=TABLE} \\
\bottomrule
\end{tabularx}
\end{center}
\begin{center}\footnotesize
\begin{tabularx}{0.88\linewidth}{@{}l l X@{}}
\toprule
Field Name & Bits & Type / Fixed Value \\
\midrule
\texttt{src0} & \texttt{0:7} & \texttt{SourceEncoding} \\
\textit{Reserved} & \texttt{8:15} & \texttt{zero} \\
\texttt{opcode2} & \texttt{16:20} & \texttt{exact=FRCP\_F16} \\
\textit{Reserved} & \texttt{21:27} & \texttt{zero} \\
\hyperref[enum:lane16_m]{\texttt{lane0}} & \texttt{28} & \texttt{lane16\_m (def h0)} \\
\textit{Reserved} & \texttt{29:37} & \texttt{zero} \\
\hyperref[enum:neg_m]{\texttt{neg0}} & \texttt{38} & \texttt{neg\_m (def none)} \\
\hyperref[enum:abs_m]{\texttt{abs0}} & \texttt{39} & \texttt{abs\_m (def none)} \\
\texttt{dest\_register\_index} & \texttt{40:45} & \texttt{uint} \\
\hyperref[enum:dest_width_single_lane_16_bit_m]{\texttt{dest\_width}} & \texttt{46:47} & \texttt{dest\_width\_single\_lane\_16\_bit\_m} \\
\texttt{opcode1} & \texttt{48:51} & \texttt{exact=TABLE} \\
\texttt{instruction\_class} & \texttt{52:56} & \texttt{exact=R1} \\
\texttt{fau\_page\_index} & \texttt{57:58} & \texttt{uint} \\
\hyperref[enum:flow_control_m]{\texttt{flow}} & \texttt{59:62} & \texttt{flow\_control\_m} \\
\textit{Reserved} & \texttt{63} & \texttt{zero} \\
\bottomrule
\end{tabularx}
\end{center}
\textbf{Fields Explained:}
\begin{tightitemize}\footnotesize
  \item \texttt{src0}: Primary source operand (\texttt{r0..r63}, \texttt{u0..u63} uniform, \texttt{k*} FAU constant, or \texttt{t0..t2} temporary; \texttt{\textasciicircum{}} marks discard/last-use). \emph{Example:} {\scriptsize\ttfamily \mnem{FADD.f32} \dst{r0}, \src{r1}, \src{r2}} (\textbf{SPIR-V:} \texttt{OpFAdd (Operand 0)}; \textbf{GLSL:} \texttt{float c = a + b;})
  \item \texttt{opcode2}: Secondary sub-opcode further specializing instructions sharing \texttt{instruction\_class} and \texttt{opcode1} (e.g., conversions, SFU tables, moves). \emph{Example:} {\scriptsize\ttfamily \mnem{F32\_TO\_S32}\imod{.round\_zero} \dst{r0}, \src{r1}} (\textbf{SPIR-V:} \texttt{OpConvertFToS}; \textbf{GLSL:} \texttt{int y = int(x);})
  \item \hyperref[enum:lane16_m]{\texttt{lane0}}: Stage A single-lane extractor on \texttt{src0} (\texttt{b0..b3} for 8-bit or \texttt{h0}/\texttt{h1} for 16-bit narrow conversions and \texttt{MKVEC}). \emph{Example:} {\scriptsize\ttfamily \mnem{F16\_TO\_F32} \dst{r0}, \src{r1}\imod{.h1}} (\textbf{SPIR-V:} \texttt{OpCompositeExtract}; \textbf{GLSL:} \texttt{float y = float(x.y);})
  \item \hyperref[enum:neg_m]{\texttt{neg0}}: Stage A floating-point negation on \texttt{src0} (flips IEEE sign bit after swizzle/widen and \texttt{abs0}). \emph{Example:} {\scriptsize\ttfamily \mnem{FADD.f32} \dst{r0}, \src{r1}\imod{.neg}, \src{r2}} (\textbf{SPIR-V:} \texttt{OpFNegate}; \textbf{GLSL:} \texttt{float c = -a + b;})
  \item \hyperref[enum:abs_m]{\texttt{abs0}}: Stage A floating-point absolute value on \texttt{src0} (clears IEEE sign bit after swizzle/widen, before \texttt{neg0}). \emph{Example:} {\scriptsize\ttfamily \mnem{FADD.f32} \dst{r0}, \src{r1}\imod{.abs}, \src{r2}} (\textbf{SPIR-V:} \texttt{GLSL.std.450 FAbs}; \textbf{GLSL:} \texttt{float c = abs(a) + b;})
  \item \texttt{dest\_register\_index}: 6-bit destination GPR index (\texttt{r0..r63}) receiving the Stage C writeback result (even-aligned pair \texttt{[rN:rN+1]} for 64-bit ops). \emph{Example:} {\scriptsize\ttfamily \mnem{FMA.f32} \dst{r0}, \src{r1}, \src{r2}, \src{r3}} (\textbf{SPIR-V:} \texttt{SSA Result ID}; \textbf{GLSL:} \texttt{float r0 = fma(r1, r2, r3);})
  \item \hyperref[enum:dest_width_single_lane_16_bit_m]{\texttt{dest\_width}}: Stage C destination lane write-mask (\texttt{none} = full 32-bit, \texttt{h0} = low 16-bit, \texttt{h1} = high 16-bit without bit-shifting). \emph{Example:} {\scriptsize\ttfamily \mnem{FADD.v2f16} \dst{r0}\imod{.h1}, \src{r1}, \src{r2}} (\textbf{SPIR-V:} \texttt{OpVectorInsertDynamic / WriteMask}; \textbf{GLSL:} \texttt{dst.y = a.y + b.y;})
  \item \texttt{opcode1}: Primary sub-opcode within the instruction class uniquely identifying the operation family. \emph{Example:} {\scriptsize\ttfamily \mnem{IADD.s32} \dst{r0}, \src{r1}, \src{r2}} (\textbf{SPIR-V:} \texttt{OpIAdd (operation selector)}; \textbf{GLSL:} \texttt{// Primary opcode bits})
  \item \texttt{instruction\_class}: 4-bit major instruction class selector (\texttt{bits 48:51}) partitioning the 64-bit Valhall opcode space (e.g., \texttt{FMA}, \texttt{CVT}, \texttt{SFU}, \texttt{LS}, \texttt{TEX}, \texttt{BRANCH}). \emph{Example:} {\scriptsize\ttfamily \mnem{FADD.f32} \dst{r0}, \src{r1}, \src{r2}} (\textbf{SPIR-V:} \texttt{N/A (encoding class)}; \textbf{GLSL:} \texttt{// Hardware execution pipe / decode class})
  \item \texttt{fau\_page\_index}: Selects the FAU (Fast Access Uniform) page (\texttt{0..3}) when a source operand references hardware small constants, thread IDs, or lane constants. \emph{Example:} {\scriptsize\ttfamily \mnem{FADD.f32} \dst{r0}, \src{r1}, \src{k1f}} (\textbf{SPIR-V:} \texttt{OpConstant / BuiltIn}; \textbf{GLSL:} \texttt{float y = x + 1.0;})
  \item \hyperref[enum:flow_control_m]{\texttt{flow}}: 4-bit embedded flow-control \& scoreboard wait modifier (\texttt{bits 59:62}: \texttt{none}, \texttt{wait0..wait0126}, \texttt{wait}, \texttt{reconverge}, \texttt{discard}, \texttt{end}). \emph{Example:} {\scriptsize\ttfamily \mnem{FADD.f32}\imod{.wait0.reconverge} \dst{r0}, \src{r1}, \src{r2}} (\textbf{SPIR-V:} \texttt{OpControlBarrier / OpReturn}; \textbf{GLSL:} \texttt{// Scoreboard wait + warp reconvergence})
\end{tightitemize}
\clearpage

\subsection*{FROUND}
\addcontentsline{toc}{subsection}{FROUND}
Rounds each floating-point element $i \in \{0, \dots, N-1\}$ of \texttt{src0} ($N=1$ \texttt{f32}, $N=2$ \texttt{v2f16}) to an integral floating-point value using \texttt{round} (\texttt{round\_integer\_m}: \texttt{none}=RNE, \texttt{round\_up}=ceil, \texttt{round\_down}=floor, \texttt{round\_zero}=trunc, \texttt{round\_na}=ties-away-from-zero).
\par\textbf{Notes:}
\begin{tightitemize}
  \item \texttt{v2f16} only on arch \textless{}= v10 (\texttt{hw\_tests.rs:1716-1719}).
  \item On arch \textgreater{}= v11, \texttt{round\_up}/\texttt{round\_down} does NOT auto-FTZ inputs (\texttt{kraid/nir.rs:849-868}).
  \item \textbf{Source references:} \texttt{kraid/ops.rs:2053-2077}, \texttt{encode\_v9.rs:1651-1669}, \texttt{nir.rs:839-876}, \texttt{hw\_tests.rs:1703-1730}.
\end{tightitemize}
\codebox{\mnem{FROUND}\imod{\{.flow\}}\imod{\{.round\}} \dst{dest}\imod{\{.dest\_width\}}, \src{src0}\imod{\{.widen0\}}\imod{\{.abs0\}}\imod{\{.neg0\}}\\\mnem{FROUND}\imod{\{.flow\}}\imod{\{.round\}} \dst{dest}\imod{\{.dest\_width\}}, \src{src0}\imod{\{.abs0\}}\imod{\{.neg0\}}\\\mnem{FROUND}\imod{\{.flow\}}\imod{\{.round\}} \dst{dest}\imod{\{.dest\_width\}}, \src{src0}\imod{\{.swz0\}}\imod{\{.abs0\}}\imod{\{.neg0\}}}
\textbf{Operands:}
\begin{tightitemize}\footnotesize
  \item \texttt{\dst{dest}} (\texttt{dest\_register\_index}, \texttt{dest\_width}): Destination GPR receiving rounded integral float value.
  \item \texttt{\src{src0}}: Floating-point input rounded according to \texttt{round} modifier.
\end{tightitemize}
\par\textbf{Closest SPIR-V Op:} \texttt{GLSL.std.450 RoundEven / Floor / Ceil / Trunc}
\par\textbf{Pseudo-GLSL (Operational Semantics):}
\codebox{for (int i = 0; i \textless{} N; ++i) \{\\    if      (round == round\_down) dest[i] = floor(src0[i]);\\    else if (round == round\_up)   dest[i] = ceil(src0[i]);\\    else if (round == round\_zero) dest[i] = trunc(src0[i]);\\    else if (round == round\_na)   dest[i] = sign(src0[i]) * floor(abs(src0[i]) + 0.5);\\    else                          dest[i] = roundEven(src0[i]);\\\}}
\textbf{Execution Context:}
\begin{tightitemize}\footnotesize
  \item \textbf{Compute stage preloads (\hyperref[sec:csf-spd]{\texttt{SPD\_k}}, \texttt{r16..r23}):} Hardware pre-fills \texttt{r55..r57} with \texttt{WorkGroupID.xyz + JOB\_OFFSET} (from \hyperref[sec:csf-regs]{\texttt{COMPUTESR::JOB\_OFFSET\_X/Y/Z}} \texttt{r34..r36}) and \texttt{r60..r62} with \texttt{LocalInvocationID.xyz} or \texttt{GlobalInvocationID.xyz} (from \texttt{WG\_SIZE} \texttt{r33} \& \texttt{JOB\_SIZE} \texttt{r37..r39}, selected by \texttt{ShaderProgram::preload}) so compute shaders can derive thread/element indices without memory loads.
  \item \textbf{Vertex stage preloads (\hyperref[sec:csf-spd]{\texttt{SPD\_0}/\texttt{SPD\_1}}, \texttt{r16..r19}):} The hardware Index Fetcher reads \hyperref[sec:csf-regs]{\texttt{IDVSSR::INDEX\_BUFFER}} (\texttt{r54:r55}) and pre-fills \texttt{r59, r60} with \texttt{VertexIndex} (\texttt{+ VERTEX\_OFFSET + GLOBAL\_ATTR\_OFF}) and \texttt{r61} with \texttt{InstanceIndex} (\texttt{+ INSTANCE\_OFFSET}) so Position and Varying shaders can pull vertex attributes directly.
  \item \textbf{Fragment stage preloads \& FP modes (\hyperref[sec:csf-spd]{\texttt{SPD\_2}}, \texttt{r20:r21}):} Hardware pre-fills \texttt{r59} with \texttt{floatBitsToUint(FragCoord.z)} and \texttt{r60} with the MSAA sample coverage mask (\texttt{gl\_SampleMaskIn[0]}), while \texttt{ShaderProgram} word~0 latches 32/64 GPR allocation and warp-wide \texttt{flush\_to\_zero\_mode}, \texttt{suppress\_na\_n}, and \texttt{suppress\_inf}.
  \item \textbf{Inline FAU \& Special Operands (\hyperref[sec:csf-tsd-fau]{\texttt{FAU\_k}}, \texttt{r8..r15}):} Any 8-bit source with \texttt{mode=2} reads zero-latency uniform \texttt{u0..u31} from \texttt{FAU\_k} (\texttt{slots64 = bits[63:56]}); \texttt{mode=3} reads \texttt{SmallConstant} (Table~1.2) or \texttt{FauSpecial} (\texttt{pc}, \texttt{lane\_id}, \texttt{lane\_width}, \texttt{tls\_ptr}, \texttt{wls\_ptr}, \texttt{blend\_descriptor\_0..7}).
\end{tightitemize}
\par\textbf{Worked Examples:}
\par\textbf{Example 1 (Floor of negative fractional float):}
\codebox{\mnem{FROUND.f32}\imod{.round\_down} \dst{r0}, \src{r1}}
With \texttt{r1 = -1.25f}: \texttt{.round\_down} returns \texttt{-2.0f} (\texttt{0xC0000000}); \texttt{-0.0f} preserves \texttt{-0.0f}.
\par\textbf{Example 2 (Ties-to-even (RNE) vs ties-away (.round\_na)):}
\codebox{\mnem{FROUND.f32}\imod{.round\_na} \dst{r0}, \src{r1}}
With \texttt{r1 = 2.5f}: default RNE returns \texttt{2.0f} (even), whereas \texttt{.round\_na} rounds ties away from zero -\textgreater{} \texttt{3.0f}.
\par\textbf{Mined SPIR-V $\to$ Valhall ISA Example (Ceiling a shifted LOD level with .round\_up --- Fragment shader, Mali-G710 (v10)):}
\codebox{\textbf{; SPIR-V (input)}\\\%shifted = OpFSub \%float \%lod \%float\_0\_5\\\%ceiled  = OpExtInst \%float \%glsl Ceil \%shifted\\[3pt]\textbf{; Valhall ISA (malioc output)}\\FADD.f32 r3, r2, k0fr2f.h1.neg\\FROUND.f32.round\_up r3, r3\textasciicircum{}}
\texttt{GLSL.std.450 Ceil} maps directly to \texttt{FROUND.f32.round\_up} (while \texttt{Floor}, \texttt{Trunc}, and \texttt{RoundEven} select \texttt{.round\_down}, \texttt{.round\_zero}, and default RNE respectively).
\begin{center}\footnotesize
\begin{tabularx}{\linewidth}{@{}lllX@{}}
\toprule
Variant & Arch & Class & Opcode Fields \\
\midrule
\texttt{(none)} & v9--v14 & \texttt{R1} & \texttt{opcode2=FROUND\_F32, opcode1=CONVERT} \\
\texttt{(none)} & v9--v14 & \texttt{R1} & \texttt{opcode2=FROUND\_F32, opcode1=CONVERT} \\
\texttt{(none)} & v9--v14 & \texttt{R1} & \texttt{opcode2=FROUND\_F32, opcode1=CONVERT} \\
\texttt{(none)} & v9--v14 & \texttt{R1} & \texttt{opcode2=FROUND\_V2F16, opcode1=CONVERT} \\
\bottomrule
\end{tabularx}
\end{center}
\begin{center}\footnotesize
\begin{tabularx}{0.88\linewidth}{@{}l l X@{}}
\toprule
Field Name & Bits & Type / Fixed Value \\
\midrule
\texttt{src0} & \texttt{0:7} & \texttt{SourceEncoding} \\
\textit{Reserved} & \texttt{8:15} & \texttt{zero} \\
\texttt{opcode2} & \texttt{16:20} & \texttt{exact=FROUND\_F32} \\
\textit{Reserved} & \texttt{21:27} & \texttt{zero} \\
\hyperref[enum:widen_m]{\texttt{widen0}} & \texttt{28:29} & \texttt{widen\_m (def none)} \\
\hyperref[enum:round_integer_m]{\texttt{round}} & \texttt{30:32} & \texttt{round\_integer\_m (def none)} \\
\textit{Reserved} & \texttt{33:37} & \texttt{zero} \\
\hyperref[enum:neg_m]{\texttt{neg0}} & \texttt{38} & \texttt{neg\_m (def none)} \\
\hyperref[enum:abs_m]{\texttt{abs0}} & \texttt{39} & \texttt{abs\_m (def none)} \\
\texttt{dest\_register\_index} & \texttt{40:45} & \texttt{uint} \\
\hyperref[enum:dest_width_m]{\texttt{dest\_width}} & \texttt{46:47} & \texttt{dest\_width\_m} \\
\texttt{opcode1} & \texttt{48:51} & \texttt{exact=CONVERT} \\
\texttt{instruction\_class} & \texttt{52:56} & \texttt{exact=R1} \\
\texttt{fau\_page\_index} & \texttt{57:58} & \texttt{uint} \\
\hyperref[enum:flow_control_m]{\texttt{flow}} & \texttt{59:62} & \texttt{flow\_control\_m} \\
\textit{Reserved} & \texttt{63} & \texttt{zero} \\
\bottomrule
\end{tabularx}
\end{center}
\textbf{Fields Explained:}
\begin{tightitemize}\footnotesize
  \item \texttt{src0}: Primary source operand (\texttt{r0..r63}, \texttt{u0..u63} uniform, \texttt{k*} FAU constant, or \texttt{t0..t2} temporary; \texttt{\textasciicircum{}} marks discard/last-use). \emph{Example:} {\scriptsize\ttfamily \mnem{FADD.f32} \dst{r0}, \src{r1}, \src{r2}} (\textbf{SPIR-V:} \texttt{OpFAdd (Operand 0)}; \textbf{GLSL:} \texttt{float c = a + b;})
  \item \texttt{opcode2}: Secondary sub-opcode further specializing instructions sharing \texttt{instruction\_class} and \texttt{opcode1} (e.g., conversions, SFU tables, moves). \emph{Example:} {\scriptsize\ttfamily \mnem{F32\_TO\_S32}\imod{.round\_zero} \dst{r0}, \src{r1}} (\textbf{SPIR-V:} \texttt{OpConvertFToS}; \textbf{GLSL:} \texttt{int y = int(x);})
  \item \hyperref[enum:widen_m]{\texttt{widen0}}: Stage A widening extractor on \texttt{src0} (promotes FP16 \texttt{h0}/\texttt{h1} to FP32, or sign/zero-extends \texttt{b0..b3}/\texttt{h0..h1}/\texttt{w0..w1} to 32/64-bit before \texttt{abs}/\texttt{neg}). \emph{Example:} {\scriptsize\ttfamily \mnem{IADD.s32} \dst{r0}, \src{r1}\imod{.h0}, \src{r2}} (\textbf{SPIR-V:} \texttt{OpSConvert / OpFConvert}; \textbf{GLSL:} \texttt{int c = int(int16\_t(a)) + b;})
  \item \hyperref[enum:round_integer_m]{\texttt{round}}: Rounding mode (\texttt{none}/RNE, \texttt{round\_up}, \texttt{round\_down}, \texttt{round\_zero}, \texttt{round\_na}) applied in FP ALU, conversions, \texttt{HADD}, \texttt{LDEXP}, and Stage C narrowing. \emph{Example:} {\scriptsize\ttfamily \mnem{FROUND.f32}\imod{.round\_down} \dst{r0}, \src{r1}} (\textbf{SPIR-V:} \texttt{GLSL.std.450 Floor / Ceil / Trunc}; \textbf{GLSL:} \texttt{float y = floor(x);})
  \item \hyperref[enum:neg_m]{\texttt{neg0}}: Stage A floating-point negation on \texttt{src0} (flips IEEE sign bit after swizzle/widen and \texttt{abs0}). \emph{Example:} {\scriptsize\ttfamily \mnem{FADD.f32} \dst{r0}, \src{r1}\imod{.neg}, \src{r2}} (\textbf{SPIR-V:} \texttt{OpFNegate}; \textbf{GLSL:} \texttt{float c = -a + b;})
  \item \hyperref[enum:abs_m]{\texttt{abs0}}: Stage A floating-point absolute value on \texttt{src0} (clears IEEE sign bit after swizzle/widen, before \texttt{neg0}). \emph{Example:} {\scriptsize\ttfamily \mnem{FADD.f32} \dst{r0}, \src{r1}\imod{.abs}, \src{r2}} (\textbf{SPIR-V:} \texttt{GLSL.std.450 FAbs}; \textbf{GLSL:} \texttt{float c = abs(a) + b;})
  \item \texttt{dest\_register\_index}: 6-bit destination GPR index (\texttt{r0..r63}) receiving the Stage C writeback result (even-aligned pair \texttt{[rN:rN+1]} for 64-bit ops). \emph{Example:} {\scriptsize\ttfamily \mnem{FMA.f32} \dst{r0}, \src{r1}, \src{r2}, \src{r3}} (\textbf{SPIR-V:} \texttt{SSA Result ID}; \textbf{GLSL:} \texttt{float r0 = fma(r1, r2, r3);})
  \item \hyperref[enum:dest_width_m]{\texttt{dest\_width}}: Stage C destination lane write-mask (\texttt{none} = full 32-bit, \texttt{h0} = low 16-bit, \texttt{h1} = high 16-bit without bit-shifting). \emph{Example:} {\scriptsize\ttfamily \mnem{FADD.v2f16} \dst{r0}\imod{.h1}, \src{r1}, \src{r2}} (\textbf{SPIR-V:} \texttt{OpVectorInsertDynamic / WriteMask}; \textbf{GLSL:} \texttt{dst.y = a.y + b.y;})
  \item \texttt{opcode1}: Primary sub-opcode within the instruction class uniquely identifying the operation family. \emph{Example:} {\scriptsize\ttfamily \mnem{IADD.s32} \dst{r0}, \src{r1}, \src{r2}} (\textbf{SPIR-V:} \texttt{OpIAdd (operation selector)}; \textbf{GLSL:} \texttt{// Primary opcode bits})
  \item \texttt{instruction\_class}: 4-bit major instruction class selector (\texttt{bits 48:51}) partitioning the 64-bit Valhall opcode space (e.g., \texttt{FMA}, \texttt{CVT}, \texttt{SFU}, \texttt{LS}, \texttt{TEX}, \texttt{BRANCH}). \emph{Example:} {\scriptsize\ttfamily \mnem{FADD.f32} \dst{r0}, \src{r1}, \src{r2}} (\textbf{SPIR-V:} \texttt{N/A (encoding class)}; \textbf{GLSL:} \texttt{// Hardware execution pipe / decode class})
  \item \texttt{fau\_page\_index}: Selects the FAU (Fast Access Uniform) page (\texttt{0..3}) when a source operand references hardware small constants, thread IDs, or lane constants. \emph{Example:} {\scriptsize\ttfamily \mnem{FADD.f32} \dst{r0}, \src{r1}, \src{k1f}} (\textbf{SPIR-V:} \texttt{OpConstant / BuiltIn}; \textbf{GLSL:} \texttt{float y = x + 1.0;})
  \item \hyperref[enum:flow_control_m]{\texttt{flow}}: 4-bit embedded flow-control \& scoreboard wait modifier (\texttt{bits 59:62}: \texttt{none}, \texttt{wait0..wait0126}, \texttt{wait}, \texttt{reconverge}, \texttt{discard}, \texttt{end}). \emph{Example:} {\scriptsize\ttfamily \mnem{FADD.f32}\imod{.wait0.reconverge} \dst{r0}, \src{r1}, \src{r2}} (\textbf{SPIR-V:} \texttt{OpControlBarrier / OpReturn}; \textbf{GLSL:} \texttt{// Scoreboard wait + warp reconvergence})
\end{tightitemize}
\clearpage

\subsection*{FRSQ}
\addcontentsline{toc}{subsection}{FRSQ}
Computes the floating-point reciprocal square root \texttt{dest = 1.0 / sqrt(src0)} (\texttt{f32} or \texttt{f16}) on the SFU with \texttt{\textless{}= 1 ULP} accuracy. Returns \texttt{+Inf} for \texttt{+0.0}, \texttt{-Inf} for \texttt{-0.0}, and \texttt{NaN} for negative inputs; on \texttt{f32}, subnormal inputs are flushed to zero.
\par\textbf{Notes:}
\begin{tightitemize}
  \item Hardware precision is \texttt{1 ULP} (\texttt{hw\_tests.rs:1742}).
  \item Negative non-zero -\textgreater{} NaN, \texttt{+0.0} -\textgreater{} \texttt{+Inf}, \texttt{-0.0} -\textgreater{} \texttt{-Inf}.
  \item \textbf{Source references:} \texttt{kraid/ops.rs:2082-2103}, \texttt{encode\_v9.rs:1671-1688}, \texttt{hw\_tests.rs:1733-1745}.
\end{tightitemize}
\codebox{\mnem{FRSQ}\imod{\{.flow\}} \dst{dest}\imod{\{.dest\_width\}}, \src{src0}\imod{\{.lane0\}}\imod{\{.abs0\}}\imod{\{.neg0\}}\\\mnem{FRSQ}\imod{\{.flow\}} \dst{dest}\imod{\{.dest\_width\_replicate\}}, \src{src0}\imod{\{.lane0\}}\imod{\{.abs0\}}\imod{\{.neg0\}}\\\mnem{FRSQ}\imod{\{.flow\}} \dst{dest}\imod{\{.dest\_width\}}, \src{src0}\imod{\{.widen0\}}\imod{\{.abs0\}}\imod{\{.neg0\}}}
\textbf{Operands:}
\begin{tightitemize}\footnotesize
  \item \texttt{\dst{dest}} (\texttt{dest\_register\_index}, \texttt{dest\_width}, \texttt{dest\_width\_replicate}): Destination GPR receiving reciprocal square root \texttt{1.0 / sqrt(src0)}.
  \item \texttt{\src{src0}}: Floating-point input.
\end{tightitemize}
\par\textbf{Closest SPIR-V Op:} \texttt{GLSL.std.450 InverseSqrt}
\par\textbf{Pseudo-GLSL (Operational Semantics):}
\codebox{float x = is\_f32 ? ftz(src0) : float(lane0(src0));\\dest = (x \textless{} 0.0) ? NaN : ((x == 0.0) ? copysign(INF, x) : inversesqrt(x));}
\textbf{Execution Context:}
\begin{tightitemize}\footnotesize
  \item \textbf{Compute stage preloads (\hyperref[sec:csf-spd]{\texttt{SPD\_k}}, \texttt{r16..r23}):} Hardware pre-fills \texttt{r55..r57} with \texttt{WorkGroupID.xyz + JOB\_OFFSET} (from \hyperref[sec:csf-regs]{\texttt{COMPUTESR::JOB\_OFFSET\_X/Y/Z}} \texttt{r34..r36}) and \texttt{r60..r62} with \texttt{LocalInvocationID.xyz} or \texttt{GlobalInvocationID.xyz} (from \texttt{WG\_SIZE} \texttt{r33} \& \texttt{JOB\_SIZE} \texttt{r37..r39}, selected by \texttt{ShaderProgram::preload}) so compute shaders can derive thread/element indices without memory loads.
  \item \textbf{Vertex stage preloads (\hyperref[sec:csf-spd]{\texttt{SPD\_0}/\texttt{SPD\_1}}, \texttt{r16..r19}):} The hardware Index Fetcher reads \hyperref[sec:csf-regs]{\texttt{IDVSSR::INDEX\_BUFFER}} (\texttt{r54:r55}) and pre-fills \texttt{r59, r60} with \texttt{VertexIndex} (\texttt{+ VERTEX\_OFFSET + GLOBAL\_ATTR\_OFF}) and \texttt{r61} with \texttt{InstanceIndex} (\texttt{+ INSTANCE\_OFFSET}) so Position and Varying shaders can pull vertex attributes directly.
  \item \textbf{Fragment stage preloads \& FP modes (\hyperref[sec:csf-spd]{\texttt{SPD\_2}}, \texttt{r20:r21}):} Hardware pre-fills \texttt{r59} with \texttt{floatBitsToUint(FragCoord.z)} and \texttt{r60} with the MSAA sample coverage mask (\texttt{gl\_SampleMaskIn[0]}), while \texttt{ShaderProgram} word~0 latches 32/64 GPR allocation and warp-wide \texttt{flush\_to\_zero\_mode}, \texttt{suppress\_na\_n}, and \texttt{suppress\_inf}.
  \item \textbf{Inline FAU \& Special Operands (\hyperref[sec:csf-tsd-fau]{\texttt{FAU\_k}}, \texttt{r8..r15}):} Any 8-bit source with \texttt{mode=2} reads zero-latency uniform \texttt{u0..u31} from \texttt{FAU\_k} (\texttt{slots64 = bits[63:56]}); \texttt{mode=3} reads \texttt{SmallConstant} (Table~1.2) or \texttt{FauSpecial} (\texttt{pc}, \texttt{lane\_id}, \texttt{lane\_width}, \texttt{tls\_ptr}, \texttt{wls\_ptr}, \texttt{blend\_descriptor\_0..7}).
\end{tightitemize}
\par\textbf{Worked Examples:}
\par\textbf{Example 1 (Exact square reciprocal):}
\codebox{\mnem{FRSQ.f32} \dst{r0}, \src{r1}}
With \texttt{r1 = 4.0f}: returns \texttt{0.5f} (\texttt{0x3F000000}, \texttt{\textless{}= 1 ULP}).
\par\textbf{Example 2 (Zero and negative inputs):}
\codebox{\mnem{FRSQ.f32} \dst{r0}, \src{r1}}
With \texttt{r1 = -0.0f} -\textgreater{} \texttt{-Inf} (\texttt{0xFF800000}); with \texttt{r1 = -1.0f} -\textgreater{} \texttt{qNaN}; with \texttt{r1 = +Inf} -\textgreater{} \texttt{+0.0f}.
\par\textbf{Mined SPIR-V $\to$ Valhall ISA Example (Vector normalization via SFU reciprocal square root --- Fragment shader, Mali-G710 (v10)):}
\codebox{\textbf{; SPIR-V (input)}\\\%dot  = OpDot \%float \%n \%n\\\%rinv = OpExtInst \%float \%glsl InverseSqrt \%dot\\\%nx   = OpFMul \%float \%n.x \%rinv\\[3pt]\textbf{; Valhall ISA (malioc output)}\\FMUL.f32 r3, r0, r0\\FMA.f32 r3, r1, r1, r3\textasciicircum{}\\FMA.f32 r3, r2, r2, r3\textasciicircum{}\\FRSQ.f32 r3, r3\textasciicircum{}\\FMUL.f32 r0, r0\textasciicircum{}, r3}
\texttt{GLSL.std.450 InverseSqrt} lowers directly to a single \texttt{FRSQ.f32} instruction on the SFU pipe.
\begin{center}\footnotesize
\begin{tabularx}{\linewidth}{@{}lllX@{}}
\toprule
Variant & Arch & Class & Opcode Fields \\
\midrule
\texttt{(none)} & v9--v14 & \texttt{R1} & \texttt{opcode2=FRSQ\_F16, opcode1=TABLE} \\
\texttt{(none)} & v9--v14 & \texttt{R1} & \texttt{opcode2=FRSQ\_F16, opcode1=TABLE} \\
\texttt{(none)} & v9--v14 & \texttt{R1} & \texttt{opcode2=FRSQ\_F32, opcode1=TABLE} \\
\bottomrule
\end{tabularx}
\end{center}
\begin{center}\footnotesize
\begin{tabularx}{0.88\linewidth}{@{}l l X@{}}
\toprule
Field Name & Bits & Type / Fixed Value \\
\midrule
\texttt{src0} & \texttt{0:7} & \texttt{SourceEncoding} \\
\textit{Reserved} & \texttt{8:15} & \texttt{zero} \\
\texttt{opcode2} & \texttt{16:20} & \texttt{exact=FRSQ\_F16} \\
\textit{Reserved} & \texttt{21:27} & \texttt{zero} \\
\hyperref[enum:lane16_m]{\texttt{lane0}} & \texttt{28} & \texttt{lane16\_m (def h0)} \\
\textit{Reserved} & \texttt{29:37} & \texttt{zero} \\
\hyperref[enum:neg_m]{\texttt{neg0}} & \texttt{38} & \texttt{neg\_m (def none)} \\
\hyperref[enum:abs_m]{\texttt{abs0}} & \texttt{39} & \texttt{abs\_m (def none)} \\
\texttt{dest\_register\_index} & \texttt{40:45} & \texttt{uint} \\
\hyperref[enum:dest_width_single_lane_16_bit_m]{\texttt{dest\_width}} & \texttt{46:47} & \texttt{dest\_width\_single\_lane\_16\_bit\_m} \\
\texttt{opcode1} & \texttt{48:51} & \texttt{exact=TABLE} \\
\texttt{instruction\_class} & \texttt{52:56} & \texttt{exact=R1} \\
\texttt{fau\_page\_index} & \texttt{57:58} & \texttt{uint} \\
\hyperref[enum:flow_control_m]{\texttt{flow}} & \texttt{59:62} & \texttt{flow\_control\_m} \\
\textit{Reserved} & \texttt{63} & \texttt{zero} \\
\bottomrule
\end{tabularx}
\end{center}
\textbf{Fields Explained:}
\begin{tightitemize}\footnotesize
  \item \texttt{src0}: Primary source operand (\texttt{r0..r63}, \texttt{u0..u63} uniform, \texttt{k*} FAU constant, or \texttt{t0..t2} temporary; \texttt{\textasciicircum{}} marks discard/last-use). \emph{Example:} {\scriptsize\ttfamily \mnem{FADD.f32} \dst{r0}, \src{r1}, \src{r2}} (\textbf{SPIR-V:} \texttt{OpFAdd (Operand 0)}; \textbf{GLSL:} \texttt{float c = a + b;})
  \item \texttt{opcode2}: Secondary sub-opcode further specializing instructions sharing \texttt{instruction\_class} and \texttt{opcode1} (e.g., conversions, SFU tables, moves). \emph{Example:} {\scriptsize\ttfamily \mnem{F32\_TO\_S32}\imod{.round\_zero} \dst{r0}, \src{r1}} (\textbf{SPIR-V:} \texttt{OpConvertFToS}; \textbf{GLSL:} \texttt{int y = int(x);})
  \item \hyperref[enum:lane16_m]{\texttt{lane0}}: Stage A single-lane extractor on \texttt{src0} (\texttt{b0..b3} for 8-bit or \texttt{h0}/\texttt{h1} for 16-bit narrow conversions and \texttt{MKVEC}). \emph{Example:} {\scriptsize\ttfamily \mnem{F16\_TO\_F32} \dst{r0}, \src{r1}\imod{.h1}} (\textbf{SPIR-V:} \texttt{OpCompositeExtract}; \textbf{GLSL:} \texttt{float y = float(x.y);})
  \item \hyperref[enum:neg_m]{\texttt{neg0}}: Stage A floating-point negation on \texttt{src0} (flips IEEE sign bit after swizzle/widen and \texttt{abs0}). \emph{Example:} {\scriptsize\ttfamily \mnem{FADD.f32} \dst{r0}, \src{r1}\imod{.neg}, \src{r2}} (\textbf{SPIR-V:} \texttt{OpFNegate}; \textbf{GLSL:} \texttt{float c = -a + b;})
  \item \hyperref[enum:abs_m]{\texttt{abs0}}: Stage A floating-point absolute value on \texttt{src0} (clears IEEE sign bit after swizzle/widen, before \texttt{neg0}). \emph{Example:} {\scriptsize\ttfamily \mnem{FADD.f32} \dst{r0}, \src{r1}\imod{.abs}, \src{r2}} (\textbf{SPIR-V:} \texttt{GLSL.std.450 FAbs}; \textbf{GLSL:} \texttt{float c = abs(a) + b;})
  \item \texttt{dest\_register\_index}: 6-bit destination GPR index (\texttt{r0..r63}) receiving the Stage C writeback result (even-aligned pair \texttt{[rN:rN+1]} for 64-bit ops). \emph{Example:} {\scriptsize\ttfamily \mnem{FMA.f32} \dst{r0}, \src{r1}, \src{r2}, \src{r3}} (\textbf{SPIR-V:} \texttt{SSA Result ID}; \textbf{GLSL:} \texttt{float r0 = fma(r1, r2, r3);})
  \item \hyperref[enum:dest_width_single_lane_16_bit_m]{\texttt{dest\_width}}: Stage C destination lane write-mask (\texttt{none} = full 32-bit, \texttt{h0} = low 16-bit, \texttt{h1} = high 16-bit without bit-shifting). \emph{Example:} {\scriptsize\ttfamily \mnem{FADD.v2f16} \dst{r0}\imod{.h1}, \src{r1}, \src{r2}} (\textbf{SPIR-V:} \texttt{OpVectorInsertDynamic / WriteMask}; \textbf{GLSL:} \texttt{dst.y = a.y + b.y;})
  \item \texttt{opcode1}: Primary sub-opcode within the instruction class uniquely identifying the operation family. \emph{Example:} {\scriptsize\ttfamily \mnem{IADD.s32} \dst{r0}, \src{r1}, \src{r2}} (\textbf{SPIR-V:} \texttt{OpIAdd (operation selector)}; \textbf{GLSL:} \texttt{// Primary opcode bits})
  \item \texttt{instruction\_class}: 4-bit major instruction class selector (\texttt{bits 48:51}) partitioning the 64-bit Valhall opcode space (e.g., \texttt{FMA}, \texttt{CVT}, \texttt{SFU}, \texttt{LS}, \texttt{TEX}, \texttt{BRANCH}). \emph{Example:} {\scriptsize\ttfamily \mnem{FADD.f32} \dst{r0}, \src{r1}, \src{r2}} (\textbf{SPIR-V:} \texttt{N/A (encoding class)}; \textbf{GLSL:} \texttt{// Hardware execution pipe / decode class})
  \item \texttt{fau\_page\_index}: Selects the FAU (Fast Access Uniform) page (\texttt{0..3}) when a source operand references hardware small constants, thread IDs, or lane constants. \emph{Example:} {\scriptsize\ttfamily \mnem{FADD.f32} \dst{r0}, \src{r1}, \src{k1f}} (\textbf{SPIR-V:} \texttt{OpConstant / BuiltIn}; \textbf{GLSL:} \texttt{float y = x + 1.0;})
  \item \hyperref[enum:flow_control_m]{\texttt{flow}}: 4-bit embedded flow-control \& scoreboard wait modifier (\texttt{bits 59:62}: \texttt{none}, \texttt{wait0..wait0126}, \texttt{wait}, \texttt{reconverge}, \texttt{discard}, \texttt{end}). \emph{Example:} {\scriptsize\ttfamily \mnem{FADD.f32}\imod{.wait0.reconverge} \dst{r0}, \src{r1}, \src{r2}} (\textbf{SPIR-V:} \texttt{OpControlBarrier / OpReturn}; \textbf{GLSL:} \texttt{// Scoreboard wait + warp reconvergence})
\end{tightitemize}
\clearpage

\subsection*{FSINCOS}
\addcontentsline{toc}{subsection}{FSINCOS}
\texttt{FSINCOS.f32} (v14+): computes \texttt{sin} or \texttt{cos} (selected by \texttt{operation} modifier: \texttt{cos}/\texttt{sin}) from range-reduced inputs \texttt{src0}, \texttt{src1}.
\par\textbf{Notes:}
\begin{tightitemize}
  \item Distinct \texttt{cos} vs \texttt{sin} behavior controlled by \texttt{operation} modifier.
  \item \textbf{Source references:} \texttt{isa-v9-v14.xml}.
\end{tightitemize}
\codebox{\mnem{FSINCOS}\imod{\{.flow\}}\imod{\{.operation\}} \dst{dest}\imod{\{.dest\_width\}}, \src{src0}, \src{src1}}
\textbf{Operands:}
\begin{tightitemize}\footnotesize
  \item \texttt{\dst{dest}} (\texttt{dest\_register\_index}, \texttt{dest\_width}): Destination GPR receiving \texttt{sin} or \texttt{cos} of the reconstructed angle (v14+).
  \item \texttt{\src{src0}}: Primary reduced angle operand.
  \item \texttt{\src{src1}}: Auxiliary quadrant/range-reduction operand.
\end{tightitemize}
\par\textbf{Closest SPIR-V Op:} \texttt{GLSL.std.450 Sin / Cos (v14+)}
\par\textbf{Pseudo-GLSL (Operational Semantics):}
\codebox{float theta = reconstructAngle(src0, src1);\\dest = (func == sin) ? sin(theta) : cos(theta); // v14+}
\textbf{Execution Context:}
\begin{tightitemize}\footnotesize
  \item \textbf{Compute stage preloads (\hyperref[sec:csf-spd]{\texttt{SPD\_k}}, \texttt{r16..r23}):} Hardware pre-fills \texttt{r55..r57} with \texttt{WorkGroupID.xyz + JOB\_OFFSET} (from \hyperref[sec:csf-regs]{\texttt{COMPUTESR::JOB\_OFFSET\_X/Y/Z}} \texttt{r34..r36}) and \texttt{r60..r62} with \texttt{LocalInvocationID.xyz} or \texttt{GlobalInvocationID.xyz} (from \texttt{WG\_SIZE} \texttt{r33} \& \texttt{JOB\_SIZE} \texttt{r37..r39}, selected by \texttt{ShaderProgram::preload}) so compute shaders can derive thread/element indices without memory loads.
  \item \textbf{Vertex stage preloads (\hyperref[sec:csf-spd]{\texttt{SPD\_0}/\texttt{SPD\_1}}, \texttt{r16..r19}):} The hardware Index Fetcher reads \hyperref[sec:csf-regs]{\texttt{IDVSSR::INDEX\_BUFFER}} (\texttt{r54:r55}) and pre-fills \texttt{r59, r60} with \texttt{VertexIndex} (\texttt{+ VERTEX\_OFFSET + GLOBAL\_ATTR\_OFF}) and \texttt{r61} with \texttt{InstanceIndex} (\texttt{+ INSTANCE\_OFFSET}) so Position and Varying shaders can pull vertex attributes directly.
  \item \textbf{Fragment stage preloads \& FP modes (\hyperref[sec:csf-spd]{\texttt{SPD\_2}}, \texttt{r20:r21}):} Hardware pre-fills \texttt{r59} with \texttt{floatBitsToUint(FragCoord.z)} and \texttt{r60} with the MSAA sample coverage mask (\texttt{gl\_SampleMaskIn[0]}), while \texttt{ShaderProgram} word~0 latches 32/64 GPR allocation and warp-wide \texttt{flush\_to\_zero\_mode}, \texttt{suppress\_na\_n}, and \texttt{suppress\_inf}.
  \item \textbf{Inline FAU \& Special Operands (\hyperref[sec:csf-tsd-fau]{\texttt{FAU\_k}}, \texttt{r8..r15}):} Any 8-bit source with \texttt{mode=2} reads zero-latency uniform \texttt{u0..u31} from \texttt{FAU\_k} (\texttt{slots64 = bits[63:56]}); \texttt{mode=3} reads \texttt{SmallConstant} (Table~1.2) or \texttt{FauSpecial} (\texttt{pc}, \texttt{lane\_id}, \texttt{lane\_width}, \texttt{tls\_ptr}, \texttt{wls\_ptr}, \texttt{blend\_descriptor\_0..7}).
\end{tightitemize}
\par\textbf{Worked Examples:}
\par\textbf{Example 1 (Hardware sine evaluation (v14+)):}
\codebox{\mnem{FSINCOS.f32}\imod{.sin} \dst{r0}, \src{r1}, \src{r2}}
Evaluates hardware \texttt{sin} from range-reduced quadrant/angle inputs \texttt{r1, r2}: at \texttt{angle = pi/2} returns \texttt{+1.0f} (\texttt{0x3F800000}).
\par\textbf{Example 2 (Hardware cosine mode (.cos)):}
\codebox{\mnem{FSINCOS.f32}\imod{.cos} \dst{r0}, \src{r1}, \src{r2}}
With \texttt{.cos} on \texttt{angle = 0.0f}: returns \texttt{+1.0f}; at \texttt{angle = pi} returns \texttt{-1.0f} (\texttt{0xBF800000}).

\begin{center}\footnotesize
\begin{tabularx}{\linewidth}{@{}lllX@{}}
\toprule
Variant & Arch & Class & Opcode Fields \\
\midrule
\texttt{f32} & v9--v14 & \texttt{R2} & \texttt{opcode2=TABLE, opcode3=FSINCOS} \\
\bottomrule
\end{tabularx}
\end{center}
\begin{center}\footnotesize
\begin{tabularx}{0.88\linewidth}{@{}l l X@{}}
\toprule
Field Name & Bits & Type / Fixed Value \\
\midrule
\texttt{src0} & \texttt{0:7} & \texttt{SourceEncoding} \\
\texttt{src1} & \texttt{8:15} & \texttt{SourceEncoding} \\
\texttt{opcode2} & \texttt{16:19} & \texttt{exact=TABLE} \\
\textit{Reserved} & \texttt{20:23} & \texttt{zero} \\
\texttt{opcode3} & \texttt{24:25} & \texttt{exact=FSINCOS} \\
\textit{Reserved} & \texttt{26:29} & \texttt{zero} \\
\hyperref[enum:table_sincos_operation_m]{\texttt{operation}} & \texttt{30} & \texttt{table\_sincos\_operation\_m} \\
\textit{Reserved} & \texttt{31:39} & \texttt{zero} \\
\texttt{dest\_register\_index} & \texttt{40:45} & \texttt{uint} \\
\hyperref[enum:dest_width_m]{\texttt{dest\_width}} & \texttt{46:47} & \texttt{dest\_width\_m} \\
\texttt{datatype} & \texttt{48:51} & \texttt{exact=f32} \\
\texttt{instruction\_class} & \texttt{52:56} & \texttt{exact=R2} \\
\texttt{fau\_page\_index} & \texttt{57:58} & \texttt{uint} \\
\hyperref[enum:flow_control_m]{\texttt{flow}} & \texttt{59:62} & \texttt{flow\_control\_m} \\
\textit{Reserved} & \texttt{63} & \texttt{zero} \\
\bottomrule
\end{tabularx}
\end{center}
\textbf{Fields Explained:}
\begin{tightitemize}\footnotesize
  \item \texttt{src0}: Primary source operand (\texttt{r0..r63}, \texttt{u0..u63} uniform, \texttt{k*} FAU constant, or \texttt{t0..t2} temporary; \texttt{\textasciicircum{}} marks discard/last-use). \emph{Example:} {\scriptsize\ttfamily \mnem{FADD.f32} \dst{r0}, \src{r1}, \src{r2}} (\textbf{SPIR-V:} \texttt{OpFAdd (Operand 0)}; \textbf{GLSL:} \texttt{float c = a + b;})
  \item \texttt{src1}: Secondary source operand (register, uniform, or FAU constant; in shifts, supplies the shift count byte). \emph{Example:} {\scriptsize\ttfamily \mnem{LSHIFT.i32} \dst{r0}, \src{r1}, \src{r2}\imod{.b0}} (\textbf{SPIR-V:} \texttt{OpShiftLeftLogical (Shift)}; \textbf{GLSL:} \texttt{uint y = a \textless{}\textless{} s;})
  \item \texttt{opcode2}: Secondary sub-opcode further specializing instructions sharing \texttt{instruction\_class} and \texttt{opcode1} (e.g., conversions, SFU tables, moves). \emph{Example:} {\scriptsize\ttfamily \mnem{F32\_TO\_S32}\imod{.round\_zero} \dst{r0}, \src{r1}} (\textbf{SPIR-V:} \texttt{OpConvertFToS}; \textbf{GLSL:} \texttt{int y = int(x);})
  \item \texttt{opcode3}: Tertiary sub-opcode disambiguating specialized multi-op families (\texttt{FPOW\_SC\_*}, SFU lookup tables, \texttt{ICMP\_MULTI}). \emph{Example:} {\scriptsize\ttfamily \mnem{FPOW\_SC\_DET.f32}\imod{.pow} \dst{r0}, \src{r1}, \src{r2}} (\textbf{SPIR-V:} \texttt{GLSL.std.450 Pow}; \textbf{GLSL:} \texttt{float z = pow(x, y);})
  \item \hyperref[enum:table_sincos_operation_m]{\texttt{operation}}: Selects \texttt{sin} vs \texttt{cos} output in \texttt{FSINCOS}. \emph{Example:} {\scriptsize\ttfamily \mnem{FSINCOS}\imod{.sin} \dst{r0}, \src{r1}} (\textbf{SPIR-V:} \texttt{GLSL.std.450 Sin}; \textbf{GLSL:} \texttt{float y = sin(x);})
  \item \texttt{dest\_register\_index}: 6-bit destination GPR index (\texttt{r0..r63}) receiving the Stage C writeback result (even-aligned pair \texttt{[rN:rN+1]} for 64-bit ops). \emph{Example:} {\scriptsize\ttfamily \mnem{FMA.f32} \dst{r0}, \src{r1}, \src{r2}, \src{r3}} (\textbf{SPIR-V:} \texttt{SSA Result ID}; \textbf{GLSL:} \texttt{float r0 = fma(r1, r2, r3);})
  \item \hyperref[enum:dest_width_m]{\texttt{dest\_width}}: Stage C destination lane write-mask (\texttt{none} = full 32-bit, \texttt{h0} = low 16-bit, \texttt{h1} = high 16-bit without bit-shifting). \emph{Example:} {\scriptsize\ttfamily \mnem{FADD.v2f16} \dst{r0}\imod{.h1}, \src{r1}, \src{r2}} (\textbf{SPIR-V:} \texttt{OpVectorInsertDynamic / WriteMask}; \textbf{GLSL:} \texttt{dst.y = a.y + b.y;})
  \item \texttt{datatype}: Selects the lane data type (\texttt{f32}, \texttt{f16}, \texttt{v2f16}, \texttt{s32}, \texttt{u32}, \texttt{v2s16}, \texttt{v2u16}, \texttt{v4s8}, \texttt{v4u8}, \texttt{s64}, \texttt{u64}, etc.). \emph{Example:} {\scriptsize\ttfamily \mnem{IADD.v2s16} \dst{r0}, \src{r1}, \src{r2}} (\textbf{SPIR-V:} \texttt{OpTypeInt / OpTypeFloat}; \textbf{GLSL:} \texttt{i16vec2 c = a + b;})
  \item \texttt{instruction\_class}: 4-bit major instruction class selector (\texttt{bits 48:51}) partitioning the 64-bit Valhall opcode space (e.g., \texttt{FMA}, \texttt{CVT}, \texttt{SFU}, \texttt{LS}, \texttt{TEX}, \texttt{BRANCH}). \emph{Example:} {\scriptsize\ttfamily \mnem{FADD.f32} \dst{r0}, \src{r1}, \src{r2}} (\textbf{SPIR-V:} \texttt{N/A (encoding class)}; \textbf{GLSL:} \texttt{// Hardware execution pipe / decode class})
  \item \texttt{fau\_page\_index}: Selects the FAU (Fast Access Uniform) page (\texttt{0..3}) when a source operand references hardware small constants, thread IDs, or lane constants. \emph{Example:} {\scriptsize\ttfamily \mnem{FADD.f32} \dst{r0}, \src{r1}, \src{k1f}} (\textbf{SPIR-V:} \texttt{OpConstant / BuiltIn}; \textbf{GLSL:} \texttt{float y = x + 1.0;})
  \item \hyperref[enum:flow_control_m]{\texttt{flow}}: 4-bit embedded flow-control \& scoreboard wait modifier (\texttt{bits 59:62}: \texttt{none}, \texttt{wait0..wait0126}, \texttt{wait}, \texttt{reconverge}, \texttt{discard}, \texttt{end}). \emph{Example:} {\scriptsize\ttfamily \mnem{FADD.f32}\imod{.wait0.reconverge} \dst{r0}, \src{r1}, \src{r2}} (\textbf{SPIR-V:} \texttt{OpControlBarrier / OpReturn}; \textbf{GLSL:} \texttt{// Scoreboard wait + warp reconvergence})
\end{tightitemize}
\clearpage

\subsection*{FSQRT}
\addcontentsline{toc}{subsection}{FSQRT}
Computes the single-precision IEEE-754 square root \texttt{dest = sqrt(src0)} directly on Valhall v14+. Preserves \texttt{-0.0 -\textgreater{} -0.0}, returns \texttt{+Inf} for \texttt{+Inf}, and produces \texttt{NaN} for strictly negative inputs.
\par\textbf{Notes:}
\begin{tightitemize}
  \item IEEE-754 \texttt{sqrt}: negative non-zero -\textgreater{} NaN, \texttt{-0.0} -\textgreater{} \texttt{-0.0}, \texttt{+Inf} -\textgreater{} \texttt{+Inf}.
  \item \textbf{Source references:} \texttt{isa-v9-v14.xml}.
\end{tightitemize}
\codebox{\mnem{FSQRT}\imod{\{.flow\}} \dst{dest}\imod{\{.dest\_width\}}, \src{src0}\imod{\{.widen0\}}\imod{\{.abs0\}}\imod{\{.neg0\}}}
\textbf{Operands:}
\begin{tightitemize}\footnotesize
  \item \texttt{\dst{dest}} (\texttt{dest\_register\_index}, \texttt{dest\_width}): 32-bit destination GPR receiving \texttt{sqrt(src0)}.
  \item \texttt{\src{src0}}: FP32 input value.
\end{tightitemize}
\par\textbf{Closest SPIR-V Op:} \texttt{GLSL.std.450 Sqrt (v14+)}
\par\textbf{Pseudo-GLSL (Operational Semantics):}
\codebox{dest = (src0 \textless{} 0.0) ? NaN : sqrt(src0); // preserves -0.0 -\textgreater{} -0.0 (v14+)}
\textbf{Execution Context:}
\begin{tightitemize}\footnotesize
  \item \textbf{Compute stage preloads (\hyperref[sec:csf-spd]{\texttt{SPD\_k}}, \texttt{r16..r23}):} Hardware pre-fills \texttt{r55..r57} with \texttt{WorkGroupID.xyz + JOB\_OFFSET} (from \hyperref[sec:csf-regs]{\texttt{COMPUTESR::JOB\_OFFSET\_X/Y/Z}} \texttt{r34..r36}) and \texttt{r60..r62} with \texttt{LocalInvocationID.xyz} or \texttt{GlobalInvocationID.xyz} (from \texttt{WG\_SIZE} \texttt{r33} \& \texttt{JOB\_SIZE} \texttt{r37..r39}, selected by \texttt{ShaderProgram::preload}) so compute shaders can derive thread/element indices without memory loads.
  \item \textbf{Vertex stage preloads (\hyperref[sec:csf-spd]{\texttt{SPD\_0}/\texttt{SPD\_1}}, \texttt{r16..r19}):} The hardware Index Fetcher reads \hyperref[sec:csf-regs]{\texttt{IDVSSR::INDEX\_BUFFER}} (\texttt{r54:r55}) and pre-fills \texttt{r59, r60} with \texttt{VertexIndex} (\texttt{+ VERTEX\_OFFSET + GLOBAL\_ATTR\_OFF}) and \texttt{r61} with \texttt{InstanceIndex} (\texttt{+ INSTANCE\_OFFSET}) so Position and Varying shaders can pull vertex attributes directly.
  \item \textbf{Fragment stage preloads \& FP modes (\hyperref[sec:csf-spd]{\texttt{SPD\_2}}, \texttt{r20:r21}):} Hardware pre-fills \texttt{r59} with \texttt{floatBitsToUint(FragCoord.z)} and \texttt{r60} with the MSAA sample coverage mask (\texttt{gl\_SampleMaskIn[0]}), while \texttt{ShaderProgram} word~0 latches 32/64 GPR allocation and warp-wide \texttt{flush\_to\_zero\_mode}, \texttt{suppress\_na\_n}, and \texttt{suppress\_inf}.
  \item \textbf{Inline FAU \& Special Operands (\hyperref[sec:csf-tsd-fau]{\texttt{FAU\_k}}, \texttt{r8..r15}):} Any 8-bit source with \texttt{mode=2} reads zero-latency uniform \texttt{u0..u31} from \texttt{FAU\_k} (\texttt{slots64 = bits[63:56]}); \texttt{mode=3} reads \texttt{SmallConstant} (Table~1.2) or \texttt{FauSpecial} (\texttt{pc}, \texttt{lane\_id}, \texttt{lane\_width}, \texttt{tls\_ptr}, \texttt{wls\_ptr}, \texttt{blend\_descriptor\_0..7}).
\end{tightitemize}
\par\textbf{Worked Examples:}
\par\textbf{Example 1 (Positive normal square root (v14+)):}
\codebox{\mnem{FSQRT.f32} \dst{r0}, \src{r1}}
With \texttt{r1 = 9.0f}: returns \texttt{3.0f} (\texttt{0x40400000}); v9--v13 lower \texttt{sqrt(x)} to \texttt{FRSQ} + \texttt{FRCP}.
\par\textbf{Example 2 (Signed zero \& negative domain):}
\codebox{\mnem{FSQRT.f32} \dst{r0}, \src{r1}}
With \texttt{r1 = -0.0f}: IEEE-754 requires \texttt{sqrt(-0.0) = -0.0f} (\texttt{0x80000000}). With \texttt{r1 = -4.0f} -\textgreater{} \texttt{qNaN}.

\begin{center}\footnotesize
\begin{tabularx}{\linewidth}{@{}lllX@{}}
\toprule
Variant & Arch & Class & Opcode Fields \\
\midrule
\texttt{(none)} & v9--v14 & \texttt{R1} & \texttt{opcode2=FSQRT\_F32, opcode1=TABLE} \\
\bottomrule
\end{tabularx}
\end{center}
\begin{center}\footnotesize
\begin{tabularx}{0.88\linewidth}{@{}l l X@{}}
\toprule
Field Name & Bits & Type / Fixed Value \\
\midrule
\texttt{src0} & \texttt{0:7} & \texttt{SourceEncoding} \\
\textit{Reserved} & \texttt{8:15} & \texttt{zero} \\
\texttt{opcode2} & \texttt{16:20} & \texttt{exact=FSQRT\_F32} \\
\textit{Reserved} & \texttt{21:27} & \texttt{zero} \\
\hyperref[enum:widen_m]{\texttt{widen0}} & \texttt{28:29} & \texttt{widen\_m (def none)} \\
\textit{Reserved} & \texttt{30:37} & \texttt{zero} \\
\hyperref[enum:neg_m]{\texttt{neg0}} & \texttt{38} & \texttt{neg\_m (def none)} \\
\hyperref[enum:abs_m]{\texttt{abs0}} & \texttt{39} & \texttt{abs\_m (def none)} \\
\texttt{dest\_register\_index} & \texttt{40:45} & \texttt{uint} \\
\hyperref[enum:dest_width_m]{\texttt{dest\_width}} & \texttt{46:47} & \texttt{dest\_width\_m} \\
\texttt{opcode1} & \texttt{48:51} & \texttt{exact=TABLE} \\
\texttt{instruction\_class} & \texttt{52:56} & \texttt{exact=R1} \\
\texttt{fau\_page\_index} & \texttt{57:58} & \texttt{uint} \\
\hyperref[enum:flow_control_m]{\texttt{flow}} & \texttt{59:62} & \texttt{flow\_control\_m} \\
\textit{Reserved} & \texttt{63} & \texttt{zero} \\
\bottomrule
\end{tabularx}
\end{center}
\textbf{Fields Explained:}
\begin{tightitemize}\footnotesize
  \item \texttt{src0}: Primary source operand (\texttt{r0..r63}, \texttt{u0..u63} uniform, \texttt{k*} FAU constant, or \texttt{t0..t2} temporary; \texttt{\textasciicircum{}} marks discard/last-use). \emph{Example:} {\scriptsize\ttfamily \mnem{FADD.f32} \dst{r0}, \src{r1}, \src{r2}} (\textbf{SPIR-V:} \texttt{OpFAdd (Operand 0)}; \textbf{GLSL:} \texttt{float c = a + b;})
  \item \texttt{opcode2}: Secondary sub-opcode further specializing instructions sharing \texttt{instruction\_class} and \texttt{opcode1} (e.g., conversions, SFU tables, moves). \emph{Example:} {\scriptsize\ttfamily \mnem{F32\_TO\_S32}\imod{.round\_zero} \dst{r0}, \src{r1}} (\textbf{SPIR-V:} \texttt{OpConvertFToS}; \textbf{GLSL:} \texttt{int y = int(x);})
  \item \hyperref[enum:widen_m]{\texttt{widen0}}: Stage A widening extractor on \texttt{src0} (promotes FP16 \texttt{h0}/\texttt{h1} to FP32, or sign/zero-extends \texttt{b0..b3}/\texttt{h0..h1}/\texttt{w0..w1} to 32/64-bit before \texttt{abs}/\texttt{neg}). \emph{Example:} {\scriptsize\ttfamily \mnem{IADD.s32} \dst{r0}, \src{r1}\imod{.h0}, \src{r2}} (\textbf{SPIR-V:} \texttt{OpSConvert / OpFConvert}; \textbf{GLSL:} \texttt{int c = int(int16\_t(a)) + b;})
  \item \hyperref[enum:neg_m]{\texttt{neg0}}: Stage A floating-point negation on \texttt{src0} (flips IEEE sign bit after swizzle/widen and \texttt{abs0}). \emph{Example:} {\scriptsize\ttfamily \mnem{FADD.f32} \dst{r0}, \src{r1}\imod{.neg}, \src{r2}} (\textbf{SPIR-V:} \texttt{OpFNegate}; \textbf{GLSL:} \texttt{float c = -a + b;})
  \item \hyperref[enum:abs_m]{\texttt{abs0}}: Stage A floating-point absolute value on \texttt{src0} (clears IEEE sign bit after swizzle/widen, before \texttt{neg0}). \emph{Example:} {\scriptsize\ttfamily \mnem{FADD.f32} \dst{r0}, \src{r1}\imod{.abs}, \src{r2}} (\textbf{SPIR-V:} \texttt{GLSL.std.450 FAbs}; \textbf{GLSL:} \texttt{float c = abs(a) + b;})
  \item \texttt{dest\_register\_index}: 6-bit destination GPR index (\texttt{r0..r63}) receiving the Stage C writeback result (even-aligned pair \texttt{[rN:rN+1]} for 64-bit ops). \emph{Example:} {\scriptsize\ttfamily \mnem{FMA.f32} \dst{r0}, \src{r1}, \src{r2}, \src{r3}} (\textbf{SPIR-V:} \texttt{SSA Result ID}; \textbf{GLSL:} \texttt{float r0 = fma(r1, r2, r3);})
  \item \hyperref[enum:dest_width_m]{\texttt{dest\_width}}: Stage C destination lane write-mask (\texttt{none} = full 32-bit, \texttt{h0} = low 16-bit, \texttt{h1} = high 16-bit without bit-shifting). \emph{Example:} {\scriptsize\ttfamily \mnem{FADD.v2f16} \dst{r0}\imod{.h1}, \src{r1}, \src{r2}} (\textbf{SPIR-V:} \texttt{OpVectorInsertDynamic / WriteMask}; \textbf{GLSL:} \texttt{dst.y = a.y + b.y;})
  \item \texttt{opcode1}: Primary sub-opcode within the instruction class uniquely identifying the operation family. \emph{Example:} {\scriptsize\ttfamily \mnem{IADD.s32} \dst{r0}, \src{r1}, \src{r2}} (\textbf{SPIR-V:} \texttt{OpIAdd (operation selector)}; \textbf{GLSL:} \texttt{// Primary opcode bits})
  \item \texttt{instruction\_class}: 4-bit major instruction class selector (\texttt{bits 48:51}) partitioning the 64-bit Valhall opcode space (e.g., \texttt{FMA}, \texttt{CVT}, \texttt{SFU}, \texttt{LS}, \texttt{TEX}, \texttt{BRANCH}). \emph{Example:} {\scriptsize\ttfamily \mnem{FADD.f32} \dst{r0}, \src{r1}, \src{r2}} (\textbf{SPIR-V:} \texttt{N/A (encoding class)}; \textbf{GLSL:} \texttt{// Hardware execution pipe / decode class})
  \item \texttt{fau\_page\_index}: Selects the FAU (Fast Access Uniform) page (\texttt{0..3}) when a source operand references hardware small constants, thread IDs, or lane constants. \emph{Example:} {\scriptsize\ttfamily \mnem{FADD.f32} \dst{r0}, \src{r1}, \src{k1f}} (\textbf{SPIR-V:} \texttt{OpConstant / BuiltIn}; \textbf{GLSL:} \texttt{float y = x + 1.0;})
  \item \hyperref[enum:flow_control_m]{\texttt{flow}}: 4-bit embedded flow-control \& scoreboard wait modifier (\texttt{bits 59:62}: \texttt{none}, \texttt{wait0..wait0126}, \texttt{wait}, \texttt{reconverge}, \texttt{discard}, \texttt{end}). \emph{Example:} {\scriptsize\ttfamily \mnem{FADD.f32}\imod{.wait0.reconverge} \dst{r0}, \src{r1}, \src{r2}} (\textbf{SPIR-V:} \texttt{OpControlBarrier / OpReturn}; \textbf{GLSL:} \texttt{// Scoreboard wait + warp reconvergence})
\end{tightitemize}
\clearpage

\subsection*{LDEXP}
\addcontentsline{toc}{subsection}{LDEXP}
Scales each floating-point element $i \in \{0, \dots, N-1\}$ of \texttt{src0} by $2^{\text{src1}[i]}$ ($N=1$ \texttt{f32}, $N=2$ \texttt{v2f16}; \texttt{src1[i]} signed integer), rounding via \texttt{round}. Under \texttt{.inf}, \texttt{exp \textless{} -16384} maps \texttt{$\pm$Inf -\textgreater{} $\pm$1.0} and finite \texttt{-\textgreater{} $\pm$0.0}; under \texttt{.inf0}, \texttt{exp \textgreater{}= +16384} also maps \texttt{$\pm$0.0 -\textgreater{} $\pm$1.0} and non-zero \texttt{-\textgreater{} $\pm$Inf}.
\par\textbf{Note:} Verified in \texttt{libGLES\_mali.so} \texttt{\_mali\_ldexp\_\{sf32,inf\_sf32,inf0\_sf32\}} (\texttt{0x2e288c0..0x2e28b60}).
\codebox{\mnem{LDEXP}\imod{\{.flow\}}\imod{\{.round\}}\imod{\{.inf\}} \dst{dest}\imod{\{.dest\_width\}}, \src{src0}\imod{\{.abs0\}}\imod{\{.neg0\}}, \src{src1}}
\textbf{Operands:}
\begin{tightitemize}\footnotesize
  \item \texttt{\dst{dest}} (\texttt{dest\_register\_index}, \texttt{dest\_width}): Destination GPR receiving \texttt{src0 * 2\textasciicircum{}src1}.
  \item \texttt{\src{src0}}: Floating-point significand/value.
  \item \texttt{\src{src1}}: Signed integer power-of-two exponent adjustment.
\end{tightitemize}
\par\textbf{Closest SPIR-V Op:} \texttt{GLSL.std.450 Ldexp}
\par\textbf{Pseudo-GLSL (Operational Semantics):}
\codebox{for (int i = 0; i \textless{} N; ++i) \{\\    int e = int(src1[i]);\\    if (flag\_inf\_override \&\& ((isinf(src0[i]) \&\& e \textless{} -16384) \textbar{}\textbar{} (src0[i] == 0.0 \&\& e \textgreater{}= 16384)))\\        dest[i] = copysign(1.0, src0[i]);\\    else dest[i] = round\_rm(ldexp(src0[i], e));\\\}}
\textbf{Execution Context:}
\begin{tightitemize}\footnotesize
  \item \textbf{Compute stage preloads (\hyperref[sec:csf-spd]{\texttt{SPD\_k}}, \texttt{r16..r23}):} Hardware pre-fills \texttt{r55..r57} with \texttt{WorkGroupID.xyz + JOB\_OFFSET} (from \hyperref[sec:csf-regs]{\texttt{COMPUTESR::JOB\_OFFSET\_X/Y/Z}} \texttt{r34..r36}) and \texttt{r60..r62} with \texttt{LocalInvocationID.xyz} or \texttt{GlobalInvocationID.xyz} (from \texttt{WG\_SIZE} \texttt{r33} \& \texttt{JOB\_SIZE} \texttt{r37..r39}, selected by \texttt{ShaderProgram::preload}) so compute shaders can derive thread/element indices without memory loads.
  \item \textbf{Vertex stage preloads (\hyperref[sec:csf-spd]{\texttt{SPD\_0}/\texttt{SPD\_1}}, \texttt{r16..r19}):} The hardware Index Fetcher reads \hyperref[sec:csf-regs]{\texttt{IDVSSR::INDEX\_BUFFER}} (\texttt{r54:r55}) and pre-fills \texttt{r59, r60} with \texttt{VertexIndex} (\texttt{+ VERTEX\_OFFSET + GLOBAL\_ATTR\_OFF}) and \texttt{r61} with \texttt{InstanceIndex} (\texttt{+ INSTANCE\_OFFSET}) so Position and Varying shaders can pull vertex attributes directly.
  \item \textbf{Fragment stage preloads \& FP modes (\hyperref[sec:csf-spd]{\texttt{SPD\_2}}, \texttt{r20:r21}):} Hardware pre-fills \texttt{r59} with \texttt{floatBitsToUint(FragCoord.z)} and \texttt{r60} with the MSAA sample coverage mask (\texttt{gl\_SampleMaskIn[0]}), while \texttt{ShaderProgram} word~0 latches 32/64 GPR allocation and warp-wide \texttt{flush\_to\_zero\_mode}, \texttt{suppress\_na\_n}, and \texttt{suppress\_inf}.
  \item \textbf{Inline FAU \& Special Operands (\hyperref[sec:csf-tsd-fau]{\texttt{FAU\_k}}, \texttt{r8..r15}):} Any 8-bit source with \texttt{mode=2} reads zero-latency uniform \texttt{u0..u31} from \texttt{FAU\_k} (\texttt{slots64 = bits[63:56]}); \texttt{mode=3} reads \texttt{SmallConstant} (Table~1.2) or \texttt{FauSpecial} (\texttt{pc}, \texttt{lane\_id}, \texttt{lane\_width}, \texttt{tls\_ptr}, \texttt{wls\_ptr}, \texttt{blend\_descriptor\_0..7}).
\end{tightitemize}
\par\textbf{Worked Examples:}
\par\textbf{Example 1 (Standard power-of-two scaling):}
\codebox{\mnem{LDEXP.f32} \dst{r0}, \src{r1}, \src{r2}}
With \texttt{r1 = 1.5f, r2 = 3}: adds \texttt{3} to the exponent -\textgreater{} \texttt{1.5 * 2\textasciicircum{}3 = 12.0f} (\texttt{0x41400000}).
\par\textbf{Example 2 (Extreme-exponent override (.inf0)):}
\codebox{\mnem{LDEXP.f32}\imod{.inf0} \dst{r0}, \src{r1}, \src{r2}}
With \texttt{r1 = +0.0f, r2 = 20000 (\textgreater{}= 16384)}: normal \texttt{ldexp(0, 20000)} is \texttt{0.0f}, but \texttt{.inf0} overrides it to \texttt{+1.0f} (\texttt{0x3F800000}).
\par\textbf{Mined SPIR-V $\to$ Valhall ISA Example (Power-of-two exponent scaling before FEXP.f32 --- Fragment shader, Mali-G710 (v10)):}
\codebox{\textbf{; SPIR-V (input)}\\\%exp\_res = OpExtInst \%float \%glsl Exp2 \%x\\[3pt]\textbf{; Valhall ISA (malioc output)}\\LDEXP.f32 r55, r4, u9.w1          ; r55 = r4 * 2\textasciicircum{}24\\F32\_TO\_S32 r55, r55\textasciicircum{}\\FEXP.f32 r55, r55\textasciicircum{}, r4}
\texttt{malioc} uses \texttt{LDEXP.f32} with exponent \texttt{24} (from uniform \texttt{u9.w1}) to shift the binary point of \texttt{r4} left by 24 bits before converting to Q8.24 fixed-point for \texttt{FEXP.f32}.
\begin{center}\footnotesize
\begin{tabularx}{\linewidth}{@{}lllX@{}}
\toprule
Variant & Arch & Class & Opcode Fields \\
\midrule
\texttt{f32} & v9--v14 & \texttt{R2} & \texttt{opcode2=LDEXP} \\
\texttt{f32} & v9--v14 & \texttt{R2} & \texttt{opcode2=LDEXP} \\
\texttt{v2f16} & v9--v14 & \texttt{R2} & \texttt{opcode2=LDEXP} \\
\texttt{v2f16} & v9--v14 & \texttt{R2} & \texttt{opcode2=LDEXP} \\
\bottomrule
\end{tabularx}
\end{center}
\begin{center}\footnotesize
\begin{tabularx}{0.88\linewidth}{@{}l l X@{}}
\toprule
Field Name & Bits & Type / Fixed Value \\
\midrule
\texttt{src0} & \texttt{0:7} & \texttt{SourceEncoding} \\
\texttt{src1} & \texttt{8:15} & \texttt{SourceEncoding} \\
\texttt{opcode2} & \texttt{16:19} & \texttt{exact=LDEXP} \\
\textit{Reserved} & \texttt{20:27} & \texttt{zero} \\
\hyperref[enum:ldexp_inf_m]{\texttt{inf}} & \texttt{28:29} & \texttt{ldexp\_inf\_m (def none)} \\
\hyperref[enum:ldexp_round_m]{\texttt{round}} & \texttt{30:32} & \texttt{ldexp\_round\_m (def none)} \\
\textit{Reserved} & \texttt{33:37} & \texttt{zero} \\
\hyperref[enum:neg_m]{\texttt{neg0}} & \texttt{38} & \texttt{neg\_m (def none)} \\
\hyperref[enum:abs_m]{\texttt{abs0}} & \texttt{39} & \texttt{abs\_m (def none)} \\
\texttt{dest\_register\_index} & \texttt{40:45} & \texttt{uint} \\
\hyperref[enum:dest_width_m]{\texttt{dest\_width}} & \texttt{46:47} & \texttt{dest\_width\_m} \\
\texttt{datatype} & \texttt{48:51} & \texttt{exact=f32} \\
\texttt{instruction\_class} & \texttt{52:56} & \texttt{exact=R2} \\
\texttt{fau\_page\_index} & \texttt{57:58} & \texttt{uint} \\
\hyperref[enum:flow_control_m]{\texttt{flow}} & \texttt{59:62} & \texttt{flow\_control\_m} \\
\textit{Reserved} & \texttt{63} & \texttt{zero} \\
\bottomrule
\end{tabularx}
\end{center}
\textbf{Fields Explained:}
\begin{tightitemize}\footnotesize
  \item \texttt{src0}: Primary source operand (\texttt{r0..r63}, \texttt{u0..u63} uniform, \texttt{k*} FAU constant, or \texttt{t0..t2} temporary; \texttt{\textasciicircum{}} marks discard/last-use). \emph{Example:} {\scriptsize\ttfamily \mnem{FADD.f32} \dst{r0}, \src{r1}, \src{r2}} (\textbf{SPIR-V:} \texttt{OpFAdd (Operand 0)}; \textbf{GLSL:} \texttt{float c = a + b;})
  \item \texttt{src1}: Secondary source operand (register, uniform, or FAU constant; in shifts, supplies the shift count byte). \emph{Example:} {\scriptsize\ttfamily \mnem{LSHIFT.i32} \dst{r0}, \src{r1}, \src{r2}\imod{.b0}} (\textbf{SPIR-V:} \texttt{OpShiftLeftLogical (Shift)}; \textbf{GLSL:} \texttt{uint y = a \textless{}\textless{} s;})
  \item \texttt{opcode2}: Secondary sub-opcode further specializing instructions sharing \texttt{instruction\_class} and \texttt{opcode1} (e.g., conversions, SFU tables, moves). \emph{Example:} {\scriptsize\ttfamily \mnem{F32\_TO\_S32}\imod{.round\_zero} \dst{r0}, \src{r1}} (\textbf{SPIR-V:} \texttt{OpConvertFToS}; \textbf{GLSL:} \texttt{int y = int(x);})
  \item \hyperref[enum:ldexp_inf_m]{\texttt{inf}}: Handles extreme \texttt{LDEXP} exponents ($|e| \ge 16384$): \texttt{.inf} maps $\pm\infty \to \pm 1.0$; \texttt{.inf0} also maps $\pm 0 \to \pm 1.0$. \emph{Example:} {\scriptsize\ttfamily \mnem{LDEXP.f32}\imod{.inf0} \dst{r0}, \src{r1}, \src{r2}} (\textbf{SPIR-V:} \texttt{GLSL.std.450 Ldexp}; \textbf{GLSL:} \texttt{float y = ldexp(x, e);})
  \item \hyperref[enum:ldexp_round_m]{\texttt{round}}: Rounding mode (\texttt{none}/RNE, \texttt{round\_up}, \texttt{round\_down}, \texttt{round\_zero}, \texttt{round\_na}) applied in FP ALU, conversions, \texttt{HADD}, \texttt{LDEXP}, and Stage C narrowing. \emph{Example:} {\scriptsize\ttfamily \mnem{FROUND.f32}\imod{.round\_down} \dst{r0}, \src{r1}} (\textbf{SPIR-V:} \texttt{GLSL.std.450 Floor / Ceil / Trunc}; \textbf{GLSL:} \texttt{float y = floor(x);})
  \item \hyperref[enum:neg_m]{\texttt{neg0}}: Stage A floating-point negation on \texttt{src0} (flips IEEE sign bit after swizzle/widen and \texttt{abs0}). \emph{Example:} {\scriptsize\ttfamily \mnem{FADD.f32} \dst{r0}, \src{r1}\imod{.neg}, \src{r2}} (\textbf{SPIR-V:} \texttt{OpFNegate}; \textbf{GLSL:} \texttt{float c = -a + b;})
  \item \hyperref[enum:abs_m]{\texttt{abs0}}: Stage A floating-point absolute value on \texttt{src0} (clears IEEE sign bit after swizzle/widen, before \texttt{neg0}). \emph{Example:} {\scriptsize\ttfamily \mnem{FADD.f32} \dst{r0}, \src{r1}\imod{.abs}, \src{r2}} (\textbf{SPIR-V:} \texttt{GLSL.std.450 FAbs}; \textbf{GLSL:} \texttt{float c = abs(a) + b;})
  \item \texttt{dest\_register\_index}: 6-bit destination GPR index (\texttt{r0..r63}) receiving the Stage C writeback result (even-aligned pair \texttt{[rN:rN+1]} for 64-bit ops). \emph{Example:} {\scriptsize\ttfamily \mnem{FMA.f32} \dst{r0}, \src{r1}, \src{r2}, \src{r3}} (\textbf{SPIR-V:} \texttt{SSA Result ID}; \textbf{GLSL:} \texttt{float r0 = fma(r1, r2, r3);})
  \item \hyperref[enum:dest_width_m]{\texttt{dest\_width}}: Stage C destination lane write-mask (\texttt{none} = full 32-bit, \texttt{h0} = low 16-bit, \texttt{h1} = high 16-bit without bit-shifting). \emph{Example:} {\scriptsize\ttfamily \mnem{FADD.v2f16} \dst{r0}\imod{.h1}, \src{r1}, \src{r2}} (\textbf{SPIR-V:} \texttt{OpVectorInsertDynamic / WriteMask}; \textbf{GLSL:} \texttt{dst.y = a.y + b.y;})
  \item \texttt{datatype}: Selects the lane data type (\texttt{f32}, \texttt{f16}, \texttt{v2f16}, \texttt{s32}, \texttt{u32}, \texttt{v2s16}, \texttt{v2u16}, \texttt{v4s8}, \texttt{v4u8}, \texttt{s64}, \texttt{u64}, etc.). \emph{Example:} {\scriptsize\ttfamily \mnem{IADD.v2s16} \dst{r0}, \src{r1}, \src{r2}} (\textbf{SPIR-V:} \texttt{OpTypeInt / OpTypeFloat}; \textbf{GLSL:} \texttt{i16vec2 c = a + b;})
  \item \texttt{instruction\_class}: 4-bit major instruction class selector (\texttt{bits 48:51}) partitioning the 64-bit Valhall opcode space (e.g., \texttt{FMA}, \texttt{CVT}, \texttt{SFU}, \texttt{LS}, \texttt{TEX}, \texttt{BRANCH}). \emph{Example:} {\scriptsize\ttfamily \mnem{FADD.f32} \dst{r0}, \src{r1}, \src{r2}} (\textbf{SPIR-V:} \texttt{N/A (encoding class)}; \textbf{GLSL:} \texttt{// Hardware execution pipe / decode class})
  \item \texttt{fau\_page\_index}: Selects the FAU (Fast Access Uniform) page (\texttt{0..3}) when a source operand references hardware small constants, thread IDs, or lane constants. \emph{Example:} {\scriptsize\ttfamily \mnem{FADD.f32} \dst{r0}, \src{r1}, \src{k1f}} (\textbf{SPIR-V:} \texttt{OpConstant / BuiltIn}; \textbf{GLSL:} \texttt{float y = x + 1.0;})
  \item \hyperref[enum:flow_control_m]{\texttt{flow}}: 4-bit embedded flow-control \& scoreboard wait modifier (\texttt{bits 59:62}: \texttt{none}, \texttt{wait0..wait0126}, \texttt{wait}, \texttt{reconverge}, \texttt{discard}, \texttt{end}). \emph{Example:} {\scriptsize\ttfamily \mnem{FADD.f32}\imod{.wait0.reconverge} \dst{r0}, \src{r1}, \src{r2}} (\textbf{SPIR-V:} \texttt{OpControlBarrier / OpReturn}; \textbf{GLSL:} \texttt{// Scoreboard wait + warp reconvergence})
\end{tightitemize}
\clearpage

\section{SFU Range-Reduction \& Helpers (25 mnemonics)}
\subsection*{FADD\_LSCALE}
\addcontentsline{toc}{subsection}{FADD\_LSCALE}
Reduces normal FP32 \texttt{src1} into the logarithm mantissa interval $[0.75, 1.5)$ (forcing exponent to 0 so $m \in [1.0, 2.0)$, then halving $m$ if $|m| \ge 1.5$; non-normals pass unchanged) and computes \texttt{dest = clamp(round\_rm(src0 + reduce\_log(src1)))}. Paired with \texttt{src0 = -1.0} to compute $m - 1$ for \texttt{log2}.
\par\textbf{Notes:}
\begin{tightitemize}
  \item Used in \texttt{flog2\_32\_to} with \texttt{src0 = -1.0} (\texttt{builder.rs:410-415}).
  \item Supports \texttt{sticky} round-to-odd (\texttt{fadd\_sticky\_m}).
  \item Verified in \texttt{hw\_tests.rs:1522-1550}.
  \item \textbf{Source references:} \texttt{kraid/ops.rs:1308-1354}, \texttt{encode\_v9.rs:1317-1339}, \texttt{builder.rs:409-415}, \texttt{hw\_tests.rs:1522-1550}.
\end{tightitemize}
\codebox{\mnem{FADD\_LSCALE}\imod{\{.flow\}}\imod{\{.round\}}\imod{\{.sticky\}}\imod{\{.clamp\}} \dst{dest}\imod{\{.dest\_width\_narrow\}}, \src{src0}\imod{\{.widen0\}}\imod{\{.abs0\}}\imod{\{.neg0\}}, \src{src1}\imod{\{.widen1\}}\imod{\{.abs1\}}\imod{\{.neg1\}}}
\textbf{Operands:}
\begin{tightitemize}\footnotesize
  \item \texttt{\dst{dest}} (\texttt{dest\_register\_index}, \texttt{dest\_width\_narrow}): FP32 destination GPR receiving \texttt{src0 + reduce\_log(src1)} (typically \texttt{m - 1} when \texttt{src0 = -1.0}).
  \item \texttt{\src{src0}}: FP32 addend (usually \texttt{k1f.neg} = \texttt{-1.0}).
  \item \texttt{\src{src1}}: FP32 input whose normal mantissa is reduced into \texttt{[0.75, 1.5)}.
\end{tightitemize}
\par\textbf{Closest SPIR-V Op:} \texttt{None (lowered from GLSL.std.450 Log2)}
\par\textbf{Pseudo-GLSL (Operational Semantics):}
\codebox{float m = isNormal(src1) ? reduceTo\_0\_75\_to\_1\_5(src1) : src1;\\dest = clamp\_mode(round\_rm(mod(src0) + m)); // src0=-1.0 yields m - 1}
\textbf{Execution Context:}
\begin{tightitemize}\footnotesize
  \item \textbf{Compute stage preloads (\hyperref[sec:csf-spd]{\texttt{SPD\_k}}, \texttt{r16..r23}):} Hardware pre-fills \texttt{r55..r57} with \texttt{WorkGroupID.xyz + JOB\_OFFSET} (from \hyperref[sec:csf-regs]{\texttt{COMPUTESR::JOB\_OFFSET\_X/Y/Z}} \texttt{r34..r36}) and \texttt{r60..r62} with \texttt{LocalInvocationID.xyz} or \texttt{GlobalInvocationID.xyz} (from \texttt{WG\_SIZE} \texttt{r33} \& \texttt{JOB\_SIZE} \texttt{r37..r39}, selected by \texttt{ShaderProgram::preload}) so compute shaders can derive thread/element indices without memory loads.
  \item \textbf{Vertex stage preloads (\hyperref[sec:csf-spd]{\texttt{SPD\_0}/\texttt{SPD\_1}}, \texttt{r16..r19}):} The hardware Index Fetcher reads \hyperref[sec:csf-regs]{\texttt{IDVSSR::INDEX\_BUFFER}} (\texttt{r54:r55}) and pre-fills \texttt{r59, r60} with \texttt{VertexIndex} (\texttt{+ VERTEX\_OFFSET + GLOBAL\_ATTR\_OFF}) and \texttt{r61} with \texttt{InstanceIndex} (\texttt{+ INSTANCE\_OFFSET}) so Position and Varying shaders can pull vertex attributes directly.
  \item \textbf{Fragment stage preloads \& FP modes (\hyperref[sec:csf-spd]{\texttt{SPD\_2}}, \texttt{r20:r21}):} Hardware pre-fills \texttt{r59} with \texttt{floatBitsToUint(FragCoord.z)} and \texttt{r60} with the MSAA sample coverage mask (\texttt{gl\_SampleMaskIn[0]}), while \texttt{ShaderProgram} word~0 latches 32/64 GPR allocation and warp-wide \texttt{flush\_to\_zero\_mode}, \texttt{suppress\_na\_n}, and \texttt{suppress\_inf}.
  \item \textbf{Inline FAU \& Special Operands (\hyperref[sec:csf-tsd-fau]{\texttt{FAU\_k}}, \texttt{r8..r15}):} Any 8-bit source with \texttt{mode=2} reads zero-latency uniform \texttt{u0..u31} from \texttt{FAU\_k} (\texttt{slots64 = bits[63:56]}); \texttt{mode=3} reads \texttt{SmallConstant} (Table~1.2) or \texttt{FauSpecial} (\texttt{pc}, \texttt{lane\_id}, \texttt{lane\_width}, \texttt{tls\_ptr}, \texttt{wls\_ptr}, \texttt{blend\_descriptor\_0..7}).
\end{tightitemize}
\par\textbf{Worked Examples:}
\par\textbf{Example 1 (Log2 mantissa reduction (`m - 1`)):}
\codebox{\mnem{FADD\_LSCALE.f32} \dst{r0}, \src{k1f}\imod{.neg}, \src{r1}}
With \texttt{r1 = 6.0f} (\texttt{1.5 * 2\textasciicircum{}2}): reduces \texttt{r1} into \texttt{[0.75, 1.5)} -\textgreater{} \texttt{0.75f}, adds \texttt{-1.0f} -\textgreater{} \texttt{r0 = -0.25f} (\texttt{0xBE800000}).
\par\textbf{Example 2 (Input already in `[0.75, 1.5)`):}
\codebox{\mnem{FADD\_LSCALE.f32} \dst{r0}, \src{k1f}\imod{.neg}, \src{r1}}
With \texttt{r1 = 1.25f}: stays \texttt{1.25f}, adds \texttt{-1.0f} -\textgreater{} \texttt{r0 = +0.25f} (\texttt{0x3E800000}), feeding the \texttt{log2} polynomial.
\par\textbf{Mined SPIR-V $\to$ Valhall ISA Example (Complete 5-instruction GLSL Log2 lowering sequence --- Fragment shader, Mali-G710 (v10)):}
\codebox{\textbf{; SPIR-V (input)}\\\%y = OpExtInst \%float \%glsl Log2 \%x\\[3pt]\textbf{; Valhall ISA (malioc output)}\\FREXPE.f32.log r4, r2             ; unbiased exponent e (for m in [0.75, 1.5))\\S32\_TO\_F32 r4, r4\textasciicircum{}                ; float(e)\\FLOGD.f32 r6, r2                  ; log2(m) / (m - 1)\\FADD\_LSCALE.f32 r2, k1f.neg, r2\textasciicircum{}  ; m - 1 = (-1.0) + reduce\_log(r2)\\FMA.f32 r2, r6\textasciicircum{}, r2\textasciicircum{}, r4\textasciicircum{}         ; log2(x) = (log2(m)/(m-1)) * (m-1) + e}
This real \texttt{malioc} sequence shows why Valhall has three dedicated log helpers: \texttt{FREXPE.log} extracts the integer exponent $e$, \texttt{FLOGD} evaluates the smooth divided difference $\log_2(m)/(m-1)$ on $m \in [0.75, 1.5)$, \texttt{FADD\_LSCALE} computes $m - 1$ by adding $-1.0$ (\texttt{k1f.neg}) to the reduced mantissa, and a final \texttt{FMA} reconstructs $\log_2(x) = \frac{\log_2(m)}{m-1}(m-1) + e$.
\begin{center}\footnotesize
\begin{tabularx}{\linewidth}{@{}lllX@{}}
\toprule
Variant & Arch & Class & Opcode Fields \\
\midrule
\texttt{f32} & v9--v14 & \texttt{R2} & \texttt{opcode2=FADD\_LSCALE} \\
\bottomrule
\end{tabularx}
\end{center}
\begin{center}\footnotesize
\begin{tabularx}{0.88\linewidth}{@{}l l X@{}}
\toprule
Field Name & Bits & Type / Fixed Value \\
\midrule
\texttt{src0} & \texttt{0:7} & \texttt{SourceEncoding} \\
\texttt{src1} & \texttt{8:15} & \texttt{SourceEncoding} \\
\texttt{opcode2} & \texttt{16:19} & \texttt{exact=FADD\_LSCALE} \\
\textit{Reserved} & \texttt{20:25} & \texttt{zero} \\
\hyperref[enum:widen_m]{\texttt{widen1}} & \texttt{26:27} & \texttt{widen\_m (def none)} \\
\hyperref[enum:widen_m]{\texttt{widen0}} & \texttt{28:29} & \texttt{widen\_m (def none)} \\
\hyperref[enum:round_m]{\texttt{round}} & \texttt{30:31} & \texttt{round\_m (def none)} \\
\hyperref[enum:clamp_m]{\texttt{clamp}} & \texttt{32:33} & \texttt{clamp\_m (def none)} \\
\hyperref[enum:fadd_sticky_m]{\texttt{sticky}} & \texttt{34} & \texttt{fadd\_sticky\_m (def none)} \\
\textit{Reserved} & \texttt{35} & \texttt{zero} \\
\hyperref[enum:neg_m]{\texttt{neg1}} & \texttt{36} & \texttt{neg\_m (def none)} \\
\hyperref[enum:abs_m]{\texttt{abs1}} & \texttt{37} & \texttt{abs\_m (def none)} \\
\hyperref[enum:neg_m]{\texttt{neg0}} & \texttt{38} & \texttt{neg\_m (def none)} \\
\hyperref[enum:abs_m]{\texttt{abs0}} & \texttt{39} & \texttt{abs\_m (def none)} \\
\texttt{dest\_register\_index} & \texttt{40:45} & \texttt{uint} \\
\hyperref[enum:dest_width_narrow_m]{\texttt{dest\_width\_narrow}} & \texttt{46:47} & \texttt{dest\_width\_narrow\_m} \\
\texttt{datatype} & \texttt{48:51} & \texttt{exact=f32} \\
\texttt{instruction\_class} & \texttt{52:56} & \texttt{exact=R2} \\
\texttt{fau\_page\_index} & \texttt{57:58} & \texttt{uint} \\
\hyperref[enum:flow_control_m]{\texttt{flow}} & \texttt{59:62} & \texttt{flow\_control\_m} \\
\textit{Reserved} & \texttt{63} & \texttt{zero} \\
\bottomrule
\end{tabularx}
\end{center}
\textbf{Fields Explained:}
\begin{tightitemize}\footnotesize
  \item \texttt{src0}: Primary source operand (\texttt{r0..r63}, \texttt{u0..u63} uniform, \texttt{k*} FAU constant, or \texttt{t0..t2} temporary; \texttt{\textasciicircum{}} marks discard/last-use). \emph{Example:} {\scriptsize\ttfamily \mnem{FADD.f32} \dst{r0}, \src{r1}, \src{r2}} (\textbf{SPIR-V:} \texttt{OpFAdd (Operand 0)}; \textbf{GLSL:} \texttt{float c = a + b;})
  \item \texttt{src1}: Secondary source operand (register, uniform, or FAU constant; in shifts, supplies the shift count byte). \emph{Example:} {\scriptsize\ttfamily \mnem{LSHIFT.i32} \dst{r0}, \src{r1}, \src{r2}\imod{.b0}} (\textbf{SPIR-V:} \texttt{OpShiftLeftLogical (Shift)}; \textbf{GLSL:} \texttt{uint y = a \textless{}\textless{} s;})
  \item \texttt{opcode2}: Secondary sub-opcode further specializing instructions sharing \texttt{instruction\_class} and \texttt{opcode1} (e.g., conversions, SFU tables, moves). \emph{Example:} {\scriptsize\ttfamily \mnem{F32\_TO\_S32}\imod{.round\_zero} \dst{r0}, \src{r1}} (\textbf{SPIR-V:} \texttt{OpConvertFToS}; \textbf{GLSL:} \texttt{int y = int(x);})
  \item \hyperref[enum:widen_m]{\texttt{widen1}}: Stage A widening extractor on \texttt{src1} (promotes narrow byte/half/word lane before Stage B). \emph{Example:} {\scriptsize\ttfamily \mnem{IMUL.u64} \dst{[r0:r1]}, \src{[r2:r3]}\imod{.w0}, \src{[r4:r5]}\imod{.w0}} (\textbf{SPIR-V:} \texttt{OpUConvert / OpUMulExtended}; \textbf{GLSL:} \texttt{uint64\_t p = uint64\_t(a) * uint64\_t(b);})
  \item \hyperref[enum:widen_m]{\texttt{widen0}}: Stage A widening extractor on \texttt{src0} (promotes FP16 \texttt{h0}/\texttt{h1} to FP32, or sign/zero-extends \texttt{b0..b3}/\texttt{h0..h1}/\texttt{w0..w1} to 32/64-bit before \texttt{abs}/\texttt{neg}). \emph{Example:} {\scriptsize\ttfamily \mnem{IADD.s32} \dst{r0}, \src{r1}\imod{.h0}, \src{r2}} (\textbf{SPIR-V:} \texttt{OpSConvert / OpFConvert}; \textbf{GLSL:} \texttt{int c = int(int16\_t(a)) + b;})
  \item \hyperref[enum:round_m]{\texttt{round}}: Rounding mode (\texttt{none}/RNE, \texttt{round\_up}, \texttt{round\_down}, \texttt{round\_zero}, \texttt{round\_na}) applied in FP ALU, conversions, \texttt{HADD}, \texttt{LDEXP}, and Stage C narrowing. \emph{Example:} {\scriptsize\ttfamily \mnem{FROUND.f32}\imod{.round\_down} \dst{r0}, \src{r1}} (\textbf{SPIR-V:} \texttt{GLSL.std.450 Floor / Ceil / Trunc}; \textbf{GLSL:} \texttt{float y = floor(x);})
  \item \hyperref[enum:clamp_m]{\texttt{clamp}}: Stage B post-op floating-point clamp (\texttt{none}, \texttt{clamp\_0\_inf}, \texttt{clamp\_m1\_1}, \texttt{clamp\_0\_1}) with \texttt{-0.0 \textless{} +0.0} and NaN suppression (\texttt{NaN -\textgreater{} lo}). \emph{Example:} {\scriptsize\ttfamily \mnem{FMA.f32}\imod{.clamp\_0\_1} \dst{r0}, \src{r1}, \src{r2}, \src{r3}} (\textbf{SPIR-V:} \texttt{GLSL.std.450 FClamp}; \textbf{GLSL:} \texttt{float y = clamp(fma(a, b, c), 0.0, 1.0);})
  \item \hyperref[enum:fadd_sticky_m]{\texttt{sticky}}: OR-reduces shifted-out rounding bits into LSB (round-to-odd for double-precision emulation). \emph{Example:} {\scriptsize\ttfamily \mnem{FADD.f32}\imod{.sticky} \dst{r0}, \src{r1}, \src{r2}} (\textbf{SPIR-V:} \texttt{OpFAdd (emulated FP64)}; \textbf{GLSL:} \texttt{double c = a + b; // lowered})
  \item \hyperref[enum:neg_m]{\texttt{neg1}}: Stage A floating-point negation on \texttt{src1} (flips IEEE sign bit after swizzle/widen and \texttt{abs1}). \emph{Example:} {\scriptsize\ttfamily \mnem{FMA.f32} \dst{r0}, \src{r1}, \src{r2}\imod{.neg}, \src{r3}} (\textbf{SPIR-V:} \texttt{OpFNegate}; \textbf{GLSL:} \texttt{float d = fma(a, -b, c);})
  \item \hyperref[enum:abs_m]{\texttt{abs1}}: Stage A floating-point absolute value on \texttt{src1} (clears IEEE sign bit after swizzle/widen, before \texttt{neg1}). \emph{Example:} {\scriptsize\ttfamily \mnem{FMA.f32} \dst{r0}, \src{r1}, \src{r2}\imod{.abs}, \src{r3}} (\textbf{SPIR-V:} \texttt{GLSL.std.450 FAbs}; \textbf{GLSL:} \texttt{float d = fma(a, abs(b), c);})
  \item \hyperref[enum:neg_m]{\texttt{neg0}}: Stage A floating-point negation on \texttt{src0} (flips IEEE sign bit after swizzle/widen and \texttt{abs0}). \emph{Example:} {\scriptsize\ttfamily \mnem{FADD.f32} \dst{r0}, \src{r1}\imod{.neg}, \src{r2}} (\textbf{SPIR-V:} \texttt{OpFNegate}; \textbf{GLSL:} \texttt{float c = -a + b;})
  \item \hyperref[enum:abs_m]{\texttt{abs0}}: Stage A floating-point absolute value on \texttt{src0} (clears IEEE sign bit after swizzle/widen, before \texttt{neg0}). \emph{Example:} {\scriptsize\ttfamily \mnem{FADD.f32} \dst{r0}, \src{r1}\imod{.abs}, \src{r2}} (\textbf{SPIR-V:} \texttt{GLSL.std.450 FAbs}; \textbf{GLSL:} \texttt{float c = abs(a) + b;})
  \item \texttt{dest\_register\_index}: 6-bit destination GPR index (\texttt{r0..r63}) receiving the Stage C writeback result (even-aligned pair \texttt{[rN:rN+1]} for 64-bit ops). \emph{Example:} {\scriptsize\ttfamily \mnem{FMA.f32} \dst{r0}, \src{r1}, \src{r2}, \src{r3}} (\textbf{SPIR-V:} \texttt{SSA Result ID}; \textbf{GLSL:} \texttt{float r0 = fma(r1, r2, r3);})
  \item \hyperref[enum:dest_width_narrow_m]{\texttt{dest\_width\_narrow}}: Destination halfword lane-mask / replication mode (\texttt{h0}, \texttt{h1}, \texttt{h01}, \texttt{none}). \emph{Example:} {\scriptsize\ttfamily \mnem{F32\_TO\_F16} \dst{r0}\imod{.h01}, \src{r1}} (\textbf{SPIR-V:} \texttt{OpFConvert (f32 -\textgreater{} f16)}; \textbf{GLSL:} \texttt{float16\_t h = float16\_t(x);})
  \item \texttt{datatype}: Selects the lane data type (\texttt{f32}, \texttt{f16}, \texttt{v2f16}, \texttt{s32}, \texttt{u32}, \texttt{v2s16}, \texttt{v2u16}, \texttt{v4s8}, \texttt{v4u8}, \texttt{s64}, \texttt{u64}, etc.). \emph{Example:} {\scriptsize\ttfamily \mnem{IADD.v2s16} \dst{r0}, \src{r1}, \src{r2}} (\textbf{SPIR-V:} \texttt{OpTypeInt / OpTypeFloat}; \textbf{GLSL:} \texttt{i16vec2 c = a + b;})
  \item \texttt{instruction\_class}: 4-bit major instruction class selector (\texttt{bits 48:51}) partitioning the 64-bit Valhall opcode space (e.g., \texttt{FMA}, \texttt{CVT}, \texttt{SFU}, \texttt{LS}, \texttt{TEX}, \texttt{BRANCH}). \emph{Example:} {\scriptsize\ttfamily \mnem{FADD.f32} \dst{r0}, \src{r1}, \src{r2}} (\textbf{SPIR-V:} \texttt{N/A (encoding class)}; \textbf{GLSL:} \texttt{// Hardware execution pipe / decode class})
  \item \texttt{fau\_page\_index}: Selects the FAU (Fast Access Uniform) page (\texttt{0..3}) when a source operand references hardware small constants, thread IDs, or lane constants. \emph{Example:} {\scriptsize\ttfamily \mnem{FADD.f32} \dst{r0}, \src{r1}, \src{k1f}} (\textbf{SPIR-V:} \texttt{OpConstant / BuiltIn}; \textbf{GLSL:} \texttt{float y = x + 1.0;})
  \item \hyperref[enum:flow_control_m]{\texttt{flow}}: 4-bit embedded flow-control \& scoreboard wait modifier (\texttt{bits 59:62}: \texttt{none}, \texttt{wait0..wait0126}, \texttt{wait}, \texttt{reconverge}, \texttt{discard}, \texttt{end}). \emph{Example:} {\scriptsize\ttfamily \mnem{FADD.f32}\imod{.wait0.reconverge} \dst{r0}, \src{r1}, \src{r2}} (\textbf{SPIR-V:} \texttt{OpControlBarrier / OpReturn}; \textbf{GLSL:} \texttt{// Scoreboard wait + warp reconvergence})
\end{tightitemize}
\clearpage

\subsection*{FADD\_RSCALE}
\addcontentsline{toc}{subsection}{FADD\_RSCALE}
Computes \texttt{dest = clamp(round\_rm((src0 + src2) * 2\textasciicircum{}clamp(src3, -512, 512)))} in exact FP64 with a single rounding to FP32 (encoded as \texttt{FMA\_RSCALE} with \texttt{src1 = 1.0}). Note that the second addend is \texttt{src2} and the signed integer exponent scale is \texttt{src3}.
\par\textbf{Note:} \texttt{libGLES\_mali.so} \texttt{ModInfo} allows \texttt{.special=n\_add} only on \texttt{FADD\_RSCALE} (\texttt{0x2e28e80}).
\codebox{\mnem{FADD\_RSCALE}\imod{\{.flow\}}\imod{\{.special\}}\imod{\{.round\}}\imod{\{.clamp\}} \dst{dest}\imod{\{.dest\_width\}}, \src{src0}\imod{\{.abs0\}}\imod{\{.neg0\}}, \src{src2}\imod{\{.abs2\}}\imod{\{.neg2\}}, \src{src3}}
\textbf{Operands:}
\begin{tightitemize}\footnotesize
  \item \texttt{\dst{dest}} (\texttt{dest\_register\_index}, \texttt{dest\_width}): FP32 destination GPR receiving \texttt{(src0 + src2) * 2\textasciicircum{}clamp(src3, -512, 512)}.
  \item \texttt{\src{src0}}: First FP32 addend.
  \item \texttt{\src{src2}}: Second FP32 addend (encoded in \texttt{src2}; \texttt{src1} is fixed to \texttt{1.0}).
  \item \texttt{\src{src3}}: Signed integer exponent scale clamped to \texttt{[-512, 512]}.
\end{tightitemize}
\par\textbf{Closest SPIR-V Op:} \texttt{None (scaled FP add)}
\par\textbf{Pseudo-GLSL (Operational Semantics):}
\codebox{double s = double(src0) + double(src2);\\if (special == n\_add \&\& isinf(src0) \&\& isinf(src2) \&\& sign(src0) != sign(src2)) s = 0.0;\\dest = clamp\_mode(round\_rm(float(s * exp2(double(clamp(int(src3), -512, 512))))));}
\textbf{Execution Context:}
\begin{tightitemize}\footnotesize
  \item \textbf{Compute stage preloads (\hyperref[sec:csf-spd]{\texttt{SPD\_k}}, \texttt{r16..r23}):} Hardware pre-fills \texttt{r55..r57} with \texttt{WorkGroupID.xyz + JOB\_OFFSET} (from \hyperref[sec:csf-regs]{\texttt{COMPUTESR::JOB\_OFFSET\_X/Y/Z}} \texttt{r34..r36}) and \texttt{r60..r62} with \texttt{LocalInvocationID.xyz} or \texttt{GlobalInvocationID.xyz} (from \texttt{WG\_SIZE} \texttt{r33} \& \texttt{JOB\_SIZE} \texttt{r37..r39}, selected by \texttt{ShaderProgram::preload}) so compute shaders can derive thread/element indices without memory loads.
  \item \textbf{Vertex stage preloads (\hyperref[sec:csf-spd]{\texttt{SPD\_0}/\texttt{SPD\_1}}, \texttt{r16..r19}):} The hardware Index Fetcher reads \hyperref[sec:csf-regs]{\texttt{IDVSSR::INDEX\_BUFFER}} (\texttt{r54:r55}) and pre-fills \texttt{r59, r60} with \texttt{VertexIndex} (\texttt{+ VERTEX\_OFFSET + GLOBAL\_ATTR\_OFF}) and \texttt{r61} with \texttt{InstanceIndex} (\texttt{+ INSTANCE\_OFFSET}) so Position and Varying shaders can pull vertex attributes directly.
  \item \textbf{Fragment stage preloads \& FP modes (\hyperref[sec:csf-spd]{\texttt{SPD\_2}}, \texttt{r20:r21}):} Hardware pre-fills \texttt{r59} with \texttt{floatBitsToUint(FragCoord.z)} and \texttt{r60} with the MSAA sample coverage mask (\texttt{gl\_SampleMaskIn[0]}), while \texttt{ShaderProgram} word~0 latches 32/64 GPR allocation and warp-wide \texttt{flush\_to\_zero\_mode}, \texttt{suppress\_na\_n}, and \texttt{suppress\_inf}.
  \item \textbf{Inline FAU \& Special Operands (\hyperref[sec:csf-tsd-fau]{\texttt{FAU\_k}}, \texttt{r8..r15}):} Any 8-bit source with \texttt{mode=2} reads zero-latency uniform \texttt{u0..u31} from \texttt{FAU\_k} (\texttt{slots64 = bits[63:56]}); \texttt{mode=3} reads \texttt{SmallConstant} (Table~1.2) or \texttt{FauSpecial} (\texttt{pc}, \texttt{lane\_id}, \texttt{lane\_width}, \texttt{tls\_ptr}, \texttt{wls\_ptr}, \texttt{blend\_descriptor\_0..7}).
\end{tightitemize}
\par\textbf{Worked Examples:}
\par\textbf{Example 1 (Exact FP64-widened midpoint):}
\codebox{\mnem{FADD\_RSCALE.f32} \dst{r0}, \src{r1}, \src{r2}, \src{k1n}}
With \texttt{r1 = 3.0f, r2 = 5.0f, scale = -1}: evaluates \texttt{(3 + 5) * 2\textasciicircum{}-1 = 4.0f} (\texttt{0x40800000}) without intermediate FP32 overflow.
\par\textbf{Example 2 (Infinity cancellation override (.n\_add)):}
\codebox{\mnem{FADD\_RSCALE.f32}\imod{.n\_add} \dst{r0}, \src{r1}, \src{r2}, \src{k0}}
Used in \texttt{expm1}: prevents \texttt{+Inf + (-Inf)} from producing NaN during range reduction when \texttt{.n\_add} is set.

\begin{center}\footnotesize
\begin{tabularx}{\linewidth}{@{}lllX@{}}
\toprule
Variant & Arch & Class & Opcode Fields \\
\midrule
\texttt{(none)} & v9--v14 & \texttt{R4\_FMAR} & \texttt{opcode1=FMA\_RSCALE\_F32} \\
\bottomrule
\end{tabularx}
\end{center}
\begin{center}\footnotesize
\begin{tabularx}{0.88\linewidth}{@{}l l X@{}}
\toprule
Field Name & Bits & Type / Fixed Value \\
\midrule
\texttt{src0} & \texttt{0:7} & \texttt{SourceEncoding} \\
\texttt{src1} & \texttt{8:15} & \texttt{exact=208} \\
\texttt{src2} & \texttt{16:23} & \texttt{SourceEncoding} \\
\texttt{src3} & \texttt{24:31} & \texttt{SourceEncoding} \\
\hyperref[enum:clamp_m]{\texttt{clamp}} & \texttt{32:33} & \texttt{clamp\_m (def none)} \\
\hyperref[enum:neg_m]{\texttt{neg2}} & \texttt{34} & \texttt{neg\_m (def none)} \\
\hyperref[enum:abs_m]{\texttt{abs2}} & \texttt{35} & \texttt{abs\_m (def none)} \\
\hyperref[enum:neg_m]{\texttt{neg1}} & \texttt{36} & \texttt{exact=none} \\
\hyperref[enum:abs_m]{\texttt{abs1}} & \texttt{37} & \texttt{exact=none} \\
\hyperref[enum:neg_m]{\texttt{neg0}} & \texttt{38} & \texttt{neg\_m (def none)} \\
\hyperref[enum:abs_m]{\texttt{abs0}} & \texttt{39} & \texttt{abs\_m (def none)} \\
\texttt{dest\_register\_index} & \texttt{40:45} & \texttt{uint} \\
\hyperref[enum:dest_width_m]{\texttt{dest\_width}} & \texttt{46:47} & \texttt{dest\_width\_m} \\
\hyperref[enum:fma_rscale_special_32_m]{\texttt{special}} & \texttt{48:49} & \texttt{fma\_rscale\_special\_32\_m} \\
\hyperref[enum:fma_rscale_round_m]{\texttt{round}} & \texttt{50} & \texttt{fma\_rscale\_round\_m} \\
\texttt{opcode1} & \texttt{51} & \texttt{exact=FMA\_RSCALE\_F32} \\
\texttt{instruction\_class} & \texttt{52:56} & \texttt{exact=R4\_FMAR} \\
\texttt{fau\_page\_index} & \texttt{57:58} & \texttt{uint} \\
\hyperref[enum:flow_control_m]{\texttt{flow}} & \texttt{59:62} & \texttt{flow\_control\_m} \\
\textit{Reserved} & \texttt{63} & \texttt{zero} \\
\bottomrule
\end{tabularx}
\end{center}
\textbf{Fields Explained:}
\begin{tightitemize}\footnotesize
  \item \texttt{src0}: Primary source operand (\texttt{r0..r63}, \texttt{u0..u63} uniform, \texttt{k*} FAU constant, or \texttt{t0..t2} temporary; \texttt{\textasciicircum{}} marks discard/last-use). \emph{Example:} {\scriptsize\ttfamily \mnem{FADD.f32} \dst{r0}, \src{r1}, \src{r2}} (\textbf{SPIR-V:} \texttt{OpFAdd (Operand 0)}; \textbf{GLSL:} \texttt{float c = a + b;})
  \item \texttt{src1}: Secondary source operand (register, uniform, or FAU constant; in shifts, supplies the shift count byte). \emph{Example:} {\scriptsize\ttfamily \mnem{LSHIFT.i32} \dst{r0}, \src{r1}, \src{r2}\imod{.b0}} (\textbf{SPIR-V:} \texttt{OpShiftLeftLogical (Shift)}; \textbf{GLSL:} \texttt{uint y = a \textless{}\textless{} s;})
  \item \texttt{src2}: Third source operand (FMA addend, \texttt{MUX}/\texttt{CSEL} selector, \texttt{IADDC} carry-in, or bitwise mask in fused shift-logic). \emph{Example:} {\scriptsize\ttfamily \mnem{FMA.f32} \dst{r0}, \src{r1}, \src{r2}, \src{r3}} (\textbf{SPIR-V:} \texttt{GLSL.std.450 Fma (Addend)}; \textbf{GLSL:} \texttt{float d = fma(a, b, c);})
  \item \texttt{src3}: Fourth source operand (false-branch value in \texttt{CSEL}, exponent scale in \texttt{FMA\_RSCALE}, or fourth matrix/select input). \emph{Example:} {\scriptsize\ttfamily \mnem{CSEL.f32}\imod{.lt} \dst{r0}, \src{r1}, \src{r2}, \src{r3}, \src{r4}} (\textbf{SPIR-V:} \texttt{OpSelect (Object 2)}; \textbf{GLSL:} \texttt{float r = (a \textless{} b) ? x : y;})
  \item \hyperref[enum:clamp_m]{\texttt{clamp}}: Stage B post-op floating-point clamp (\texttt{none}, \texttt{clamp\_0\_inf}, \texttt{clamp\_m1\_1}, \texttt{clamp\_0\_1}) with \texttt{-0.0 \textless{} +0.0} and NaN suppression (\texttt{NaN -\textgreater{} lo}). \emph{Example:} {\scriptsize\ttfamily \mnem{FMA.f32}\imod{.clamp\_0\_1} \dst{r0}, \src{r1}, \src{r2}, \src{r3}} (\textbf{SPIR-V:} \texttt{GLSL.std.450 FClamp}; \textbf{GLSL:} \texttt{float y = clamp(fma(a, b, c), 0.0, 1.0);})
  \item \hyperref[enum:neg_m]{\texttt{neg2}}: Stage A floating-point negation on \texttt{src2} (flips IEEE sign bit of FMA addend after swizzle/widen and \texttt{abs2}). \emph{Example:} {\scriptsize\ttfamily \mnem{FMA.f32} \dst{r0}, \src{r1}, \src{r2}, \src{r3}\imod{.neg}} (\textbf{SPIR-V:} \texttt{OpFNegate}; \textbf{GLSL:} \texttt{float d = fma(a, b, -c);})
  \item \hyperref[enum:abs_m]{\texttt{abs2}}: Stage A floating-point absolute value on \texttt{src2} (clears IEEE sign bit of FMA addend after swizzle/widen). \emph{Example:} {\scriptsize\ttfamily \mnem{FMA.f32} \dst{r0}, \src{r1}, \src{r2}, \src{r3}\imod{.abs}} (\textbf{SPIR-V:} \texttt{GLSL.std.450 FAbs}; \textbf{GLSL:} \texttt{float d = fma(a, b, abs(c));})
  \item \hyperref[enum:neg_m]{\texttt{neg1}}: Stage A floating-point negation on \texttt{src1} (flips IEEE sign bit after swizzle/widen and \texttt{abs1}). \emph{Example:} {\scriptsize\ttfamily \mnem{FMA.f32} \dst{r0}, \src{r1}, \src{r2}\imod{.neg}, \src{r3}} (\textbf{SPIR-V:} \texttt{OpFNegate}; \textbf{GLSL:} \texttt{float d = fma(a, -b, c);})
  \item \hyperref[enum:abs_m]{\texttt{abs1}}: Stage A floating-point absolute value on \texttt{src1} (clears IEEE sign bit after swizzle/widen, before \texttt{neg1}). \emph{Example:} {\scriptsize\ttfamily \mnem{FMA.f32} \dst{r0}, \src{r1}, \src{r2}\imod{.abs}, \src{r3}} (\textbf{SPIR-V:} \texttt{GLSL.std.450 FAbs}; \textbf{GLSL:} \texttt{float d = fma(a, abs(b), c);})
  \item \hyperref[enum:neg_m]{\texttt{neg0}}: Stage A floating-point negation on \texttt{src0} (flips IEEE sign bit after swizzle/widen and \texttt{abs0}). \emph{Example:} {\scriptsize\ttfamily \mnem{FADD.f32} \dst{r0}, \src{r1}\imod{.neg}, \src{r2}} (\textbf{SPIR-V:} \texttt{OpFNegate}; \textbf{GLSL:} \texttt{float c = -a + b;})
  \item \hyperref[enum:abs_m]{\texttt{abs0}}: Stage A floating-point absolute value on \texttt{src0} (clears IEEE sign bit after swizzle/widen, before \texttt{neg0}). \emph{Example:} {\scriptsize\ttfamily \mnem{FADD.f32} \dst{r0}, \src{r1}\imod{.abs}, \src{r2}} (\textbf{SPIR-V:} \texttt{GLSL.std.450 FAbs}; \textbf{GLSL:} \texttt{float c = abs(a) + b;})
  \item \texttt{dest\_register\_index}: 6-bit destination GPR index (\texttt{r0..r63}) receiving the Stage C writeback result (even-aligned pair \texttt{[rN:rN+1]} for 64-bit ops). \emph{Example:} {\scriptsize\ttfamily \mnem{FMA.f32} \dst{r0}, \src{r1}, \src{r2}, \src{r3}} (\textbf{SPIR-V:} \texttt{SSA Result ID}; \textbf{GLSL:} \texttt{float r0 = fma(r1, r2, r3);})
  \item \hyperref[enum:dest_width_m]{\texttt{dest\_width}}: Stage C destination lane write-mask (\texttt{none} = full 32-bit, \texttt{h0} = low 16-bit, \texttt{h1} = high 16-bit without bit-shifting). \emph{Example:} {\scriptsize\ttfamily \mnem{FADD.v2f16} \dst{r0}\imod{.h1}, \src{r1}, \src{r2}} (\textbf{SPIR-V:} \texttt{OpVectorInsertDynamic / WriteMask}; \textbf{GLSL:} \texttt{dst.y = a.y + b.y;})
  \item \hyperref[enum:fma_rscale_special_32_m]{\texttt{special}}: Selects \texttt{FMA\_RSCALE}/\texttt{FREXP} special-case override (\texttt{n}, \texttt{left}, \texttt{n\_add}, \texttt{sqrt}, \texttt{log}). \emph{Example:} {\scriptsize\ttfamily \mnem{FREXPM.f32}\imod{.log} \dst{r0}, \src{r1}} (\textbf{SPIR-V:} \texttt{GLSL.std.450 Log2 (range reduction)}; \textbf{GLSL:} \texttt{float y = log2(x);})
  \item \hyperref[enum:fma_rscale_round_m]{\texttt{round}}: Rounding mode (\texttt{none}/RNE, \texttt{round\_up}, \texttt{round\_down}, \texttt{round\_zero}, \texttt{round\_na}) applied in FP ALU, conversions, \texttt{HADD}, \texttt{LDEXP}, and Stage C narrowing. \emph{Example:} {\scriptsize\ttfamily \mnem{FROUND.f32}\imod{.round\_down} \dst{r0}, \src{r1}} (\textbf{SPIR-V:} \texttt{GLSL.std.450 Floor / Ceil / Trunc}; \textbf{GLSL:} \texttt{float y = floor(x);})
  \item \texttt{opcode1}: Primary sub-opcode within the instruction class uniquely identifying the operation family. \emph{Example:} {\scriptsize\ttfamily \mnem{IADD.s32} \dst{r0}, \src{r1}, \src{r2}} (\textbf{SPIR-V:} \texttt{OpIAdd (operation selector)}; \textbf{GLSL:} \texttt{// Primary opcode bits})
  \item \texttt{instruction\_class}: 4-bit major instruction class selector (\texttt{bits 48:51}) partitioning the 64-bit Valhall opcode space (e.g., \texttt{FMA}, \texttt{CVT}, \texttt{SFU}, \texttt{LS}, \texttt{TEX}, \texttt{BRANCH}). \emph{Example:} {\scriptsize\ttfamily \mnem{FADD.f32} \dst{r0}, \src{r1}, \src{r2}} (\textbf{SPIR-V:} \texttt{N/A (encoding class)}; \textbf{GLSL:} \texttt{// Hardware execution pipe / decode class})
  \item \texttt{fau\_page\_index}: Selects the FAU (Fast Access Uniform) page (\texttt{0..3}) when a source operand references hardware small constants, thread IDs, or lane constants. \emph{Example:} {\scriptsize\ttfamily \mnem{FADD.f32} \dst{r0}, \src{r1}, \src{k1f}} (\textbf{SPIR-V:} \texttt{OpConstant / BuiltIn}; \textbf{GLSL:} \texttt{float y = x + 1.0;})
  \item \hyperref[enum:flow_control_m]{\texttt{flow}}: 4-bit embedded flow-control \& scoreboard wait modifier (\texttt{bits 59:62}: \texttt{none}, \texttt{wait0..wait0126}, \texttt{wait}, \texttt{reconverge}, \texttt{discard}, \texttt{end}). \emph{Example:} {\scriptsize\ttfamily \mnem{FADD.f32}\imod{.wait0.reconverge} \dst{r0}, \src{r1}, \src{r2}} (\textbf{SPIR-V:} \texttt{OpControlBarrier / OpReturn}; \textbf{GLSL:} \texttt{// Scoreboard wait + warp reconvergence})
\end{tightitemize}
\clearpage

\subsection*{FATAN\_ASSIST}
\addcontentsline{toc}{subsection}{FATAN\_ASSIST}
First-stage SFU helper for \texttt{atan2(y, x)} (\texttt{src0 = y}, \texttt{src1 = x}): determines the octant and ratio exponent and packs them into \texttt{dest} as \texttt{(i16 octant/exp, f16 seed)} for subsequent polynomial evaluation and \texttt{FATAN\_TABLE}.
\par\textbf{Notes:}
\begin{tightitemize}
  \item Paired with \texttt{FATAN\_TABLE.f32}.
  \item \textbf{Source references:} \texttt{bifrost/valhall/ISA.xml:2490-2498}, \texttt{isa-v9-v14.xml}.
\end{tightitemize}
\codebox{\mnem{FATAN\_ASSIST}\imod{\{.flow\}} \dst{dest}\imod{\{.dest\_width\}}, \src{src0}, \src{src1}}
\textbf{Operands:}
\begin{tightitemize}\footnotesize
  \item \texttt{\dst{dest}} (\texttt{dest\_register\_index}, \texttt{dest\_width}): 32-bit destination GPR receiving packed \texttt{(i16 octant/exp, f16 seed)} assist word.
  \item \texttt{\src{src0}}: FP32 \texttt{y} argument of \texttt{atan2(y, x)}.
  \item \texttt{\src{src1}}: FP32 \texttt{x} argument of \texttt{atan2(y, x)}.
\end{tightitemize}
\par\textbf{Closest SPIR-V Op:} \texttt{None (lowered from GLSL.std.450 Atan/Atan2)}
\par\textbf{Pseudo-GLSL (Operational Semantics):}
\codebox{dest = packAtanOctantAssist(src0 /*y*/, src1 /*x*/); // (i16 octant/exp, f16 seed)}
\textbf{Execution Context:}
\begin{tightitemize}\footnotesize
  \item \textbf{Compute stage preloads (\hyperref[sec:csf-spd]{\texttt{SPD\_k}}, \texttt{r16..r23}):} Hardware pre-fills \texttt{r55..r57} with \texttt{WorkGroupID.xyz + JOB\_OFFSET} (from \hyperref[sec:csf-regs]{\texttt{COMPUTESR::JOB\_OFFSET\_X/Y/Z}} \texttt{r34..r36}) and \texttt{r60..r62} with \texttt{LocalInvocationID.xyz} or \texttt{GlobalInvocationID.xyz} (from \texttt{WG\_SIZE} \texttt{r33} \& \texttt{JOB\_SIZE} \texttt{r37..r39}, selected by \texttt{ShaderProgram::preload}) so compute shaders can derive thread/element indices without memory loads.
  \item \textbf{Vertex stage preloads (\hyperref[sec:csf-spd]{\texttt{SPD\_0}/\texttt{SPD\_1}}, \texttt{r16..r19}):} The hardware Index Fetcher reads \hyperref[sec:csf-regs]{\texttt{IDVSSR::INDEX\_BUFFER}} (\texttt{r54:r55}) and pre-fills \texttt{r59, r60} with \texttt{VertexIndex} (\texttt{+ VERTEX\_OFFSET + GLOBAL\_ATTR\_OFF}) and \texttt{r61} with \texttt{InstanceIndex} (\texttt{+ INSTANCE\_OFFSET}) so Position and Varying shaders can pull vertex attributes directly.
  \item \textbf{Fragment stage preloads \& FP modes (\hyperref[sec:csf-spd]{\texttt{SPD\_2}}, \texttt{r20:r21}):} Hardware pre-fills \texttt{r59} with \texttt{floatBitsToUint(FragCoord.z)} and \texttt{r60} with the MSAA sample coverage mask (\texttt{gl\_SampleMaskIn[0]}), while \texttt{ShaderProgram} word~0 latches 32/64 GPR allocation and warp-wide \texttt{flush\_to\_zero\_mode}, \texttt{suppress\_na\_n}, and \texttt{suppress\_inf}.
  \item \textbf{Inline FAU \& Special Operands (\hyperref[sec:csf-tsd-fau]{\texttt{FAU\_k}}, \texttt{r8..r15}):} Any 8-bit source with \texttt{mode=2} reads zero-latency uniform \texttt{u0..u31} from \texttt{FAU\_k} (\texttt{slots64 = bits[63:56]}); \texttt{mode=3} reads \texttt{SmallConstant} (Table~1.2) or \texttt{FauSpecial} (\texttt{pc}, \texttt{lane\_id}, \texttt{lane\_width}, \texttt{tls\_ptr}, \texttt{wls\_ptr}, \texttt{blend\_descriptor\_0..7}).
\end{tightitemize}
\par\textbf{Worked Examples:}
\par\textbf{Example 1 (First-quadrant atan2 assist):}
\codebox{\mnem{FATAN\_ASSIST.f32} \dst{r0}, \src{r1}, \src{r2}}
With \texttt{y = r1 = 1.0f, x = r2 = 2.0f}: packs octant/exponent assist metadata \texttt{(i16, f16)} into \texttt{r0} for \texttt{FATAN\_TABLE}.
\par\textbf{Example 2 (Steep slope (`\textbar{}y\textbar{} \textgreater{} \textbar{}x\textbar{}`) octant swap):}
\codebox{\mnem{FATAN\_ASSIST.f32} \dst{r0}, \src{r1}, \src{r2}}
With \texttt{y = r1 = 4.0f, x = r2 = 1.0f}: records octant swap flag so \texttt{FATAN\_TABLE} applies the \texttt{pi/2 - atan(x/y)} identity.

\begin{center}\footnotesize
\begin{tabularx}{\linewidth}{@{}lllX@{}}
\toprule
Variant & Arch & Class & Opcode Fields \\
\midrule
\texttt{f32} & v9--v14 & \texttt{R2} & \texttt{opcode2=TABLE, opcode3=FATAN\_ASSIST} \\
\bottomrule
\end{tabularx}
\end{center}
\begin{center}\footnotesize
\begin{tabularx}{0.88\linewidth}{@{}l l X@{}}
\toprule
Field Name & Bits & Type / Fixed Value \\
\midrule
\texttt{src0} & \texttt{0:7} & \texttt{SourceEncoding} \\
\texttt{src1} & \texttt{8:15} & \texttt{SourceEncoding} \\
\texttt{opcode2} & \texttt{16:19} & \texttt{exact=TABLE} \\
\textit{Reserved} & \texttt{20:23} & \texttt{zero} \\
\texttt{opcode3} & \texttt{24:25} & \texttt{exact=FATAN\_ASSIST} \\
\textit{Reserved} & \texttt{26:39} & \texttt{zero} \\
\texttt{dest\_register\_index} & \texttt{40:45} & \texttt{uint} \\
\hyperref[enum:dest_width_m]{\texttt{dest\_width}} & \texttt{46:47} & \texttt{dest\_width\_m} \\
\texttt{datatype} & \texttt{48:51} & \texttt{exact=f32} \\
\texttt{instruction\_class} & \texttt{52:56} & \texttt{exact=R2} \\
\texttt{fau\_page\_index} & \texttt{57:58} & \texttt{uint} \\
\hyperref[enum:flow_control_m]{\texttt{flow}} & \texttt{59:62} & \texttt{flow\_control\_m} \\
\textit{Reserved} & \texttt{63} & \texttt{zero} \\
\bottomrule
\end{tabularx}
\end{center}
\textbf{Fields Explained:}
\begin{tightitemize}\footnotesize
  \item \texttt{src0}: Primary source operand (\texttt{r0..r63}, \texttt{u0..u63} uniform, \texttt{k*} FAU constant, or \texttt{t0..t2} temporary; \texttt{\textasciicircum{}} marks discard/last-use). \emph{Example:} {\scriptsize\ttfamily \mnem{FADD.f32} \dst{r0}, \src{r1}, \src{r2}} (\textbf{SPIR-V:} \texttt{OpFAdd (Operand 0)}; \textbf{GLSL:} \texttt{float c = a + b;})
  \item \texttt{src1}: Secondary source operand (register, uniform, or FAU constant; in shifts, supplies the shift count byte). \emph{Example:} {\scriptsize\ttfamily \mnem{LSHIFT.i32} \dst{r0}, \src{r1}, \src{r2}\imod{.b0}} (\textbf{SPIR-V:} \texttt{OpShiftLeftLogical (Shift)}; \textbf{GLSL:} \texttt{uint y = a \textless{}\textless{} s;})
  \item \texttt{opcode2}: Secondary sub-opcode further specializing instructions sharing \texttt{instruction\_class} and \texttt{opcode1} (e.g., conversions, SFU tables, moves). \emph{Example:} {\scriptsize\ttfamily \mnem{F32\_TO\_S32}\imod{.round\_zero} \dst{r0}, \src{r1}} (\textbf{SPIR-V:} \texttt{OpConvertFToS}; \textbf{GLSL:} \texttt{int y = int(x);})
  \item \texttt{opcode3}: Tertiary sub-opcode disambiguating specialized multi-op families (\texttt{FPOW\_SC\_*}, SFU lookup tables, \texttt{ICMP\_MULTI}). \emph{Example:} {\scriptsize\ttfamily \mnem{FPOW\_SC\_DET.f32}\imod{.pow} \dst{r0}, \src{r1}, \src{r2}} (\textbf{SPIR-V:} \texttt{GLSL.std.450 Pow}; \textbf{GLSL:} \texttt{float z = pow(x, y);})
  \item \texttt{dest\_register\_index}: 6-bit destination GPR index (\texttt{r0..r63}) receiving the Stage C writeback result (even-aligned pair \texttt{[rN:rN+1]} for 64-bit ops). \emph{Example:} {\scriptsize\ttfamily \mnem{FMA.f32} \dst{r0}, \src{r1}, \src{r2}, \src{r3}} (\textbf{SPIR-V:} \texttt{SSA Result ID}; \textbf{GLSL:} \texttt{float r0 = fma(r1, r2, r3);})
  \item \hyperref[enum:dest_width_m]{\texttt{dest\_width}}: Stage C destination lane write-mask (\texttt{none} = full 32-bit, \texttt{h0} = low 16-bit, \texttt{h1} = high 16-bit without bit-shifting). \emph{Example:} {\scriptsize\ttfamily \mnem{FADD.v2f16} \dst{r0}\imod{.h1}, \src{r1}, \src{r2}} (\textbf{SPIR-V:} \texttt{OpVectorInsertDynamic / WriteMask}; \textbf{GLSL:} \texttt{dst.y = a.y + b.y;})
  \item \texttt{datatype}: Selects the lane data type (\texttt{f32}, \texttt{f16}, \texttt{v2f16}, \texttt{s32}, \texttt{u32}, \texttt{v2s16}, \texttt{v2u16}, \texttt{v4s8}, \texttt{v4u8}, \texttt{s64}, \texttt{u64}, etc.). \emph{Example:} {\scriptsize\ttfamily \mnem{IADD.v2s16} \dst{r0}, \src{r1}, \src{r2}} (\textbf{SPIR-V:} \texttt{OpTypeInt / OpTypeFloat}; \textbf{GLSL:} \texttt{i16vec2 c = a + b;})
  \item \texttt{instruction\_class}: 4-bit major instruction class selector (\texttt{bits 48:51}) partitioning the 64-bit Valhall opcode space (e.g., \texttt{FMA}, \texttt{CVT}, \texttt{SFU}, \texttt{LS}, \texttt{TEX}, \texttt{BRANCH}). \emph{Example:} {\scriptsize\ttfamily \mnem{FADD.f32} \dst{r0}, \src{r1}, \src{r2}} (\textbf{SPIR-V:} \texttt{N/A (encoding class)}; \textbf{GLSL:} \texttt{// Hardware execution pipe / decode class})
  \item \texttt{fau\_page\_index}: Selects the FAU (Fast Access Uniform) page (\texttt{0..3}) when a source operand references hardware small constants, thread IDs, or lane constants. \emph{Example:} {\scriptsize\ttfamily \mnem{FADD.f32} \dst{r0}, \src{r1}, \src{k1f}} (\textbf{SPIR-V:} \texttt{OpConstant / BuiltIn}; \textbf{GLSL:} \texttt{float y = x + 1.0;})
  \item \hyperref[enum:flow_control_m]{\texttt{flow}}: 4-bit embedded flow-control \& scoreboard wait modifier (\texttt{bits 59:62}: \texttt{none}, \texttt{wait0..wait0126}, \texttt{wait}, \texttt{reconverge}, \texttt{discard}, \texttt{end}). \emph{Example:} {\scriptsize\ttfamily \mnem{FADD.f32}\imod{.wait0.reconverge} \dst{r0}, \src{r1}, \src{r2}} (\textbf{SPIR-V:} \texttt{OpControlBarrier / OpReturn}; \textbf{GLSL:} \texttt{// Scoreboard wait + warp reconvergence})
\end{tightitemize}
\clearpage

\subsection*{FATAN\_TABLE}
\addcontentsline{toc}{subsection}{FATAN\_TABLE}
Second-stage SFU lookup table for \texttt{atan2}: indexes the 8-entry octant offset table using the low 3 bits of \texttt{src0} and writes the FP32 offset constant (\texttt{0}, \texttt{$\pm$pi/2}, or \texttt{$\pm$pi}) to \texttt{dest}.
\par\textbf{Notes:}
\begin{tightitemize}
  \item Paired with \texttt{FATAN\_ASSIST.f32}.
  \item \textbf{Source references:} \texttt{bifrost/valhall/ISA.xml:2498}, \texttt{isa-v9-v14.xml}.
\end{tightitemize}
\codebox{\mnem{FATAN\_TABLE}\imod{\{.flow\}} \dst{dest}\imod{\{.dest\_width\}}, \src{src0}, \src{src1}}
\textbf{Operands:}
\begin{tightitemize}\footnotesize
  \item \texttt{\dst{dest}} (\texttt{dest\_register\_index}, \texttt{dest\_width}): FP32 destination GPR receiving the octant offset constant (\texttt{0}, \texttt{$\pm$pi/2}, \texttt{$\pm$pi}).
  \item \texttt{\src{src0}}: 32-bit assist/index word produced by \texttt{FATAN\_ASSIST}.
  \item \texttt{\src{src1}}: Source operand \texttt{src1} (GPR, FAU uniform, or inline constant).
\end{tightitemize}
\par\textbf{Closest SPIR-V Op:} \texttt{None (lowered from GLSL.std.450 Atan/Atan2)}
\par\textbf{Pseudo-GLSL (Operational Semantics):}
\codebox{dest = ATAN\_OCTANT\_OFFSET\_LUT[src0 \& 0x7u]; // 0, +-pi/2, +-pi}
\textbf{Execution Context:}
\begin{tightitemize}\footnotesize
  \item \textbf{Compute stage preloads (\hyperref[sec:csf-spd]{\texttt{SPD\_k}}, \texttt{r16..r23}):} Hardware pre-fills \texttt{r55..r57} with \texttt{WorkGroupID.xyz + JOB\_OFFSET} (from \hyperref[sec:csf-regs]{\texttt{COMPUTESR::JOB\_OFFSET\_X/Y/Z}} \texttt{r34..r36}) and \texttt{r60..r62} with \texttt{LocalInvocationID.xyz} or \texttt{GlobalInvocationID.xyz} (from \texttt{WG\_SIZE} \texttt{r33} \& \texttt{JOB\_SIZE} \texttt{r37..r39}, selected by \texttt{ShaderProgram::preload}) so compute shaders can derive thread/element indices without memory loads.
  \item \textbf{Vertex stage preloads (\hyperref[sec:csf-spd]{\texttt{SPD\_0}/\texttt{SPD\_1}}, \texttt{r16..r19}):} The hardware Index Fetcher reads \hyperref[sec:csf-regs]{\texttt{IDVSSR::INDEX\_BUFFER}} (\texttt{r54:r55}) and pre-fills \texttt{r59, r60} with \texttt{VertexIndex} (\texttt{+ VERTEX\_OFFSET + GLOBAL\_ATTR\_OFF}) and \texttt{r61} with \texttt{InstanceIndex} (\texttt{+ INSTANCE\_OFFSET}) so Position and Varying shaders can pull vertex attributes directly.
  \item \textbf{Fragment stage preloads \& FP modes (\hyperref[sec:csf-spd]{\texttt{SPD\_2}}, \texttt{r20:r21}):} Hardware pre-fills \texttt{r59} with \texttt{floatBitsToUint(FragCoord.z)} and \texttt{r60} with the MSAA sample coverage mask (\texttt{gl\_SampleMaskIn[0]}), while \texttt{ShaderProgram} word~0 latches 32/64 GPR allocation and warp-wide \texttt{flush\_to\_zero\_mode}, \texttt{suppress\_na\_n}, and \texttt{suppress\_inf}.
  \item \textbf{Inline FAU \& Special Operands (\hyperref[sec:csf-tsd-fau]{\texttt{FAU\_k}}, \texttt{r8..r15}):} Any 8-bit source with \texttt{mode=2} reads zero-latency uniform \texttt{u0..u31} from \texttt{FAU\_k} (\texttt{slots64 = bits[63:56]}); \texttt{mode=3} reads \texttt{SmallConstant} (Table~1.2) or \texttt{FauSpecial} (\texttt{pc}, \texttt{lane\_id}, \texttt{lane\_width}, \texttt{tls\_ptr}, \texttt{wls\_ptr}, \texttt{blend\_descriptor\_0..7}).
\end{tightitemize}
\par\textbf{Worked Examples:}
\par\textbf{Example 1 (Octant 0 base offset lookup):}
\codebox{\mnem{FATAN\_TABLE.f32} \dst{r0}, \src{r1}}
Given \texttt{r1} from \texttt{FATAN\_ASSIST} in octant 0 (\texttt{0 \textless{}= y \textless{}= x}): returns base angle offset \texttt{0.0f}.
\par\textbf{Example 2 (Octant 1/2 (`pi/2` or `pi`) offset lookup):}
\codebox{\mnem{FATAN\_TABLE.f32} \dst{r0}, \src{r1}}
Given \texttt{r1} in octant with \texttt{x \textless{} 0, y \textgreater{} 0}: returns \texttt{+pi} (\texttt{0x40490FDB}) or \texttt{+pi/2} (\texttt{0x3FC90FDB}) for quadrant reconstruction.

\begin{center}\footnotesize
\begin{tabularx}{\linewidth}{@{}lllX@{}}
\toprule
Variant & Arch & Class & Opcode Fields \\
\midrule
\texttt{f32} & v9--v14 & \texttt{R2} & \texttt{opcode2=TABLE, opcode3=FATAN\_TABLE} \\
\bottomrule
\end{tabularx}
\end{center}
\begin{center}\footnotesize
\begin{tabularx}{0.88\linewidth}{@{}l l X@{}}
\toprule
Field Name & Bits & Type / Fixed Value \\
\midrule
\texttt{src0} & \texttt{0:7} & \texttt{SourceEncoding} \\
\texttt{src1} & \texttt{8:15} & \texttt{SourceEncoding} \\
\texttt{opcode2} & \texttt{16:19} & \texttt{exact=TABLE} \\
\textit{Reserved} & \texttt{20:23} & \texttt{zero} \\
\texttt{opcode3} & \texttt{24:25} & \texttt{exact=FATAN\_TABLE} \\
\textit{Reserved} & \texttt{26:39} & \texttt{zero} \\
\texttt{dest\_register\_index} & \texttt{40:45} & \texttt{uint} \\
\hyperref[enum:dest_width_m]{\texttt{dest\_width}} & \texttt{46:47} & \texttt{dest\_width\_m} \\
\texttt{datatype} & \texttt{48:51} & \texttt{exact=f32} \\
\texttt{instruction\_class} & \texttt{52:56} & \texttt{exact=R2} \\
\texttt{fau\_page\_index} & \texttt{57:58} & \texttt{uint} \\
\hyperref[enum:flow_control_m]{\texttt{flow}} & \texttt{59:62} & \texttt{flow\_control\_m} \\
\textit{Reserved} & \texttt{63} & \texttt{zero} \\
\bottomrule
\end{tabularx}
\end{center}
\textbf{Fields Explained:}
\begin{tightitemize}\footnotesize
  \item \texttt{src0}: Primary source operand (\texttt{r0..r63}, \texttt{u0..u63} uniform, \texttt{k*} FAU constant, or \texttt{t0..t2} temporary; \texttt{\textasciicircum{}} marks discard/last-use). \emph{Example:} {\scriptsize\ttfamily \mnem{FADD.f32} \dst{r0}, \src{r1}, \src{r2}} (\textbf{SPIR-V:} \texttt{OpFAdd (Operand 0)}; \textbf{GLSL:} \texttt{float c = a + b;})
  \item \texttt{src1}: Secondary source operand (register, uniform, or FAU constant; in shifts, supplies the shift count byte). \emph{Example:} {\scriptsize\ttfamily \mnem{LSHIFT.i32} \dst{r0}, \src{r1}, \src{r2}\imod{.b0}} (\textbf{SPIR-V:} \texttt{OpShiftLeftLogical (Shift)}; \textbf{GLSL:} \texttt{uint y = a \textless{}\textless{} s;})
  \item \texttt{opcode2}: Secondary sub-opcode further specializing instructions sharing \texttt{instruction\_class} and \texttt{opcode1} (e.g., conversions, SFU tables, moves). \emph{Example:} {\scriptsize\ttfamily \mnem{F32\_TO\_S32}\imod{.round\_zero} \dst{r0}, \src{r1}} (\textbf{SPIR-V:} \texttt{OpConvertFToS}; \textbf{GLSL:} \texttt{int y = int(x);})
  \item \texttt{opcode3}: Tertiary sub-opcode disambiguating specialized multi-op families (\texttt{FPOW\_SC\_*}, SFU lookup tables, \texttt{ICMP\_MULTI}). \emph{Example:} {\scriptsize\ttfamily \mnem{FPOW\_SC\_DET.f32}\imod{.pow} \dst{r0}, \src{r1}, \src{r2}} (\textbf{SPIR-V:} \texttt{GLSL.std.450 Pow}; \textbf{GLSL:} \texttt{float z = pow(x, y);})
  \item \texttt{dest\_register\_index}: 6-bit destination GPR index (\texttt{r0..r63}) receiving the Stage C writeback result (even-aligned pair \texttt{[rN:rN+1]} for 64-bit ops). \emph{Example:} {\scriptsize\ttfamily \mnem{FMA.f32} \dst{r0}, \src{r1}, \src{r2}, \src{r3}} (\textbf{SPIR-V:} \texttt{SSA Result ID}; \textbf{GLSL:} \texttt{float r0 = fma(r1, r2, r3);})
  \item \hyperref[enum:dest_width_m]{\texttt{dest\_width}}: Stage C destination lane write-mask (\texttt{none} = full 32-bit, \texttt{h0} = low 16-bit, \texttt{h1} = high 16-bit without bit-shifting). \emph{Example:} {\scriptsize\ttfamily \mnem{FADD.v2f16} \dst{r0}\imod{.h1}, \src{r1}, \src{r2}} (\textbf{SPIR-V:} \texttt{OpVectorInsertDynamic / WriteMask}; \textbf{GLSL:} \texttt{dst.y = a.y + b.y;})
  \item \texttt{datatype}: Selects the lane data type (\texttt{f32}, \texttt{f16}, \texttt{v2f16}, \texttt{s32}, \texttt{u32}, \texttt{v2s16}, \texttt{v2u16}, \texttt{v4s8}, \texttt{v4u8}, \texttt{s64}, \texttt{u64}, etc.). \emph{Example:} {\scriptsize\ttfamily \mnem{IADD.v2s16} \dst{r0}, \src{r1}, \src{r2}} (\textbf{SPIR-V:} \texttt{OpTypeInt / OpTypeFloat}; \textbf{GLSL:} \texttt{i16vec2 c = a + b;})
  \item \texttt{instruction\_class}: 4-bit major instruction class selector (\texttt{bits 48:51}) partitioning the 64-bit Valhall opcode space (e.g., \texttt{FMA}, \texttt{CVT}, \texttt{SFU}, \texttt{LS}, \texttt{TEX}, \texttt{BRANCH}). \emph{Example:} {\scriptsize\ttfamily \mnem{FADD.f32} \dst{r0}, \src{r1}, \src{r2}} (\textbf{SPIR-V:} \texttt{N/A (encoding class)}; \textbf{GLSL:} \texttt{// Hardware execution pipe / decode class})
  \item \texttt{fau\_page\_index}: Selects the FAU (Fast Access Uniform) page (\texttt{0..3}) when a source operand references hardware small constants, thread IDs, or lane constants. \emph{Example:} {\scriptsize\ttfamily \mnem{FADD.f32} \dst{r0}, \src{r1}, \src{k1f}} (\textbf{SPIR-V:} \texttt{OpConstant / BuiltIn}; \textbf{GLSL:} \texttt{float y = x + 1.0;})
  \item \hyperref[enum:flow_control_m]{\texttt{flow}}: 4-bit embedded flow-control \& scoreboard wait modifier (\texttt{bits 59:62}: \texttt{none}, \texttt{wait0..wait0126}, \texttt{wait}, \texttt{reconverge}, \texttt{discard}, \texttt{end}). \emph{Example:} {\scriptsize\ttfamily \mnem{FADD.f32}\imod{.wait0.reconverge} \dst{r0}, \src{r1}, \src{r2}} (\textbf{SPIR-V:} \texttt{OpControlBarrier / OpReturn}; \textbf{GLSL:} \texttt{// Scoreboard wait + warp reconvergence})
\end{tightitemize}
\clearpage

\subsection*{FCOS\_TABLE}
\addcontentsline{toc}{subsection}{FCOS\_TABLE}
64-entry cosine lookup table for trigonometric range reduction: if \texttt{src0} is NaN returns \texttt{src0}; if \texttt{src0} is Inf returns \texttt{NaN}; if \texttt{(sign, exp) != (0, 146)} returns \texttt{-0.0}; otherwise indexes \texttt{idx = bits(src0) \& 0x3F} and writes \texttt{cos(idx * pi / 32)} to \texttt{dest} (with exact \texttt{+0.0} at \texttt{idx = 16} and \texttt{-0.0} at \texttt{idx = 48}).
\par\textbf{Notes:}
\begin{tightitemize}
  \item Exact hardware match verified in \texttt{hw\_tests.rs:1553-1567} (with \texttt{offset=false}).
  \item \textbf{Source references:} \texttt{kraid/ops.rs:1526-1584}, \texttt{encode\_v9.rs:1341-1359}, \texttt{nir.rs:1996-2003}, \texttt{hw\_tests.rs:1553-1567}.
\end{tightitemize}
\codebox{\mnem{FCOS\_TABLE}\imod{\{.flow\}}\imod{\{.offset\}} \dst{dest}\imod{\{.dest\_width\}}, \src{src0}}
\textbf{Operands:}
\begin{tightitemize}\footnotesize
  \item \texttt{\dst{dest}} (\texttt{dest\_register\_index}, \texttt{dest\_width}): FP32 destination GPR receiving \texttt{cos(idx * pi / 32)}.
  \item \texttt{\src{src0}}: Fixed-point angle word whose low 6 bits select the 64-entry cosine LUT entry.
\end{tightitemize}
\par\textbf{Closest SPIR-V Op:} \texttt{None (lowered from GLSL.std.450 Cos/Sin)}
\par\textbf{Pseudo-GLSL (Operational Semantics):}
\codebox{if (isnan(src0))      dest = src0;\\else if (isinf(src0)) dest = NaN;\\else if (signExp(src0) != (0, 146)) dest = -0.0;\\else \{ uint idx = floatBitsToUint(src0) \& 0x3Fu; dest = cos(float(idx) * PI / 32.0); \}}
\textbf{Execution Context:}
\begin{tightitemize}\footnotesize
  \item \textbf{Compute stage preloads (\hyperref[sec:csf-spd]{\texttt{SPD\_k}}, \texttt{r16..r23}):} Hardware pre-fills \texttt{r55..r57} with \texttt{WorkGroupID.xyz + JOB\_OFFSET} (from \hyperref[sec:csf-regs]{\texttt{COMPUTESR::JOB\_OFFSET\_X/Y/Z}} \texttt{r34..r36}) and \texttt{r60..r62} with \texttt{LocalInvocationID.xyz} or \texttt{GlobalInvocationID.xyz} (from \texttt{WG\_SIZE} \texttt{r33} \& \texttt{JOB\_SIZE} \texttt{r37..r39}, selected by \texttt{ShaderProgram::preload}) so compute shaders can derive thread/element indices without memory loads.
  \item \textbf{Vertex stage preloads (\hyperref[sec:csf-spd]{\texttt{SPD\_0}/\texttt{SPD\_1}}, \texttt{r16..r19}):} The hardware Index Fetcher reads \hyperref[sec:csf-regs]{\texttt{IDVSSR::INDEX\_BUFFER}} (\texttt{r54:r55}) and pre-fills \texttt{r59, r60} with \texttt{VertexIndex} (\texttt{+ VERTEX\_OFFSET + GLOBAL\_ATTR\_OFF}) and \texttt{r61} with \texttt{InstanceIndex} (\texttt{+ INSTANCE\_OFFSET}) so Position and Varying shaders can pull vertex attributes directly.
  \item \textbf{Fragment stage preloads \& FP modes (\hyperref[sec:csf-spd]{\texttt{SPD\_2}}, \texttt{r20:r21}):} Hardware pre-fills \texttt{r59} with \texttt{floatBitsToUint(FragCoord.z)} and \texttt{r60} with the MSAA sample coverage mask (\texttt{gl\_SampleMaskIn[0]}), while \texttt{ShaderProgram} word~0 latches 32/64 GPR allocation and warp-wide \texttt{flush\_to\_zero\_mode}, \texttt{suppress\_na\_n}, and \texttt{suppress\_inf}.
  \item \textbf{Inline FAU \& Special Operands (\hyperref[sec:csf-tsd-fau]{\texttt{FAU\_k}}, \texttt{r8..r15}):} Any 8-bit source with \texttt{mode=2} reads zero-latency uniform \texttt{u0..u31} from \texttt{FAU\_k} (\texttt{slots64 = bits[63:56]}); \texttt{mode=3} reads \texttt{SmallConstant} (Table~1.2) or \texttt{FauSpecial} (\texttt{pc}, \texttt{lane\_id}, \texttt{lane\_width}, \texttt{tls\_ptr}, \texttt{wls\_ptr}, \texttt{blend\_descriptor\_0..7}).
\end{tightitemize}
\par\textbf{Worked Examples:}
\par\textbf{Example 1 (Table index 0 (`cos(0) = 1.0`)):}
\codebox{\mnem{FCOS\_TABLE.u6} \dst{r0}, \src{r1}}
With \texttt{r1 = 0x49400000} (\texttt{idx = 0}, exp \texttt{127+19}): returns \texttt{cos(0) = +1.0f} (\texttt{0x3F800000}).
\par\textbf{Example 2 (Table index 16 (`cos(pi/2) = 0.0`)):}
\codebox{\mnem{FCOS\_TABLE.u6} \dst{r0}, \src{r1}}
With \texttt{r1 = 0x49400010} (\texttt{idx = 16}): returns \texttt{cos(16*pi/32) = +0.0f} (\texttt{0x00000000}).
\par\textbf{Mined SPIR-V $\to$ Valhall ISA Example (Cosine table lookup in GLSL Cos range-reduction --- Vertex shader, Mali-G710 (v10)):}
\codebox{\textbf{; SPIR-V (input)}\\\%c = OpExtInst \%float \%glsl Cos \%angle\\[3pt]\textbf{; Valhall ISA (malioc output)}\\FMUL\_RSCALE.f32 r7, r2, r2, k1n   ; r7 = (r2 * r2) * 2\textasciicircum{}(-1) = r2\textasciicircum{}2 / 2\\FSIN\_TABLE.u6 r63, r5             ; sin(idx * pi / 32)\\FCOS\_TABLE.u6 r5, r5\textasciicircum{}             ; cos(idx * pi / 32)\\FMUL.f32 r7, r7\textasciicircum{}, r5.neg          ; -(r2\textasciicircum{}2 / 2) * cos(idx * pi / 32)\\FMA.f32.clamp\_m1\_1 r2, r2\textasciicircum{}, r63\textasciicircum{}.neg, r7\textasciicircum{}}
Valhall splits $\cos(\theta_0 + \delta) \approx \cos(\theta_0) - \delta\sin(\theta_0) - \frac{\delta^2}{2}\cos(\theta_0)$ using 64-entry LUTs \texttt{FCOS\_TABLE.u6} and \texttt{FSIN\_TABLE.u6} indexed by the 6-bit quadrant/step index \texttt{r5}, finishing with \texttt{.clamp\_m1\_1} to keep the result within $[-1, 1]$.
\begin{center}\footnotesize
\begin{tabularx}{\linewidth}{@{}lllX@{}}
\toprule
Variant & Arch & Class & Opcode Fields \\
\midrule
\texttt{(none)} & v9--v14 & \texttt{R1} & \texttt{opcode2=FCOS\_TABLE, opcode1=TABLE} \\
\bottomrule
\end{tabularx}
\end{center}
\begin{center}\footnotesize
\begin{tabularx}{0.88\linewidth}{@{}l l X@{}}
\toprule
Field Name & Bits & Type / Fixed Value \\
\midrule
\texttt{src0} & \texttt{0:7} & \texttt{SourceEncoding} \\
\textit{Reserved} & \texttt{8:15} & \texttt{zero} \\
\texttt{opcode2} & \texttt{16:20} & \texttt{exact=FCOS\_TABLE} \\
\textit{Reserved} & \texttt{21:29} & \texttt{zero} \\
\hyperref[enum:table_sincos_offset_m]{\texttt{offset}} & \texttt{30} & \texttt{table\_sincos\_offset\_m} \\
\textit{Reserved} & \texttt{31:39} & \texttt{zero} \\
\texttt{dest\_register\_index} & \texttt{40:45} & \texttt{uint} \\
\hyperref[enum:dest_width_m]{\texttt{dest\_width}} & \texttt{46:47} & \texttt{dest\_width\_m} \\
\texttt{opcode1} & \texttt{48:51} & \texttt{exact=TABLE} \\
\texttt{instruction\_class} & \texttt{52:56} & \texttt{exact=R1} \\
\texttt{fau\_page\_index} & \texttt{57:58} & \texttt{uint} \\
\hyperref[enum:flow_control_m]{\texttt{flow}} & \texttt{59:62} & \texttt{flow\_control\_m} \\
\textit{Reserved} & \texttt{63} & \texttt{zero} \\
\bottomrule
\end{tabularx}
\end{center}
\textbf{Fields Explained:}
\begin{tightitemize}\footnotesize
  \item \texttt{src0}: Primary source operand (\texttt{r0..r63}, \texttt{u0..u63} uniform, \texttt{k*} FAU constant, or \texttt{t0..t2} temporary; \texttt{\textasciicircum{}} marks discard/last-use). \emph{Example:} {\scriptsize\ttfamily \mnem{FADD.f32} \dst{r0}, \src{r1}, \src{r2}} (\textbf{SPIR-V:} \texttt{OpFAdd (Operand 0)}; \textbf{GLSL:} \texttt{float c = a + b;})
  \item \texttt{opcode2}: Secondary sub-opcode further specializing instructions sharing \texttt{instruction\_class} and \texttt{opcode1} (e.g., conversions, SFU tables, moves). \emph{Example:} {\scriptsize\ttfamily \mnem{F32\_TO\_S32}\imod{.round\_zero} \dst{r0}, \src{r1}} (\textbf{SPIR-V:} \texttt{OpConvertFToS}; \textbf{GLSL:} \texttt{int y = int(x);})
  \item \hyperref[enum:table_sincos_offset_m]{\texttt{offset}}: Signed byte offset added to base address register in memory load/store/atomic instructions. \emph{Example:} {\scriptsize\ttfamily \mnem{LOAD.i32} \dst{@r0}, \src{[r2:r3]}, \imm{\#16}} (\textbf{SPIR-V:} \texttt{OpAccessChain + OpLoad}; \textbf{GLSL:} \texttt{uint y = ssbo.arr[4];})
  \item \texttt{dest\_register\_index}: 6-bit destination GPR index (\texttt{r0..r63}) receiving the Stage C writeback result (even-aligned pair \texttt{[rN:rN+1]} for 64-bit ops). \emph{Example:} {\scriptsize\ttfamily \mnem{FMA.f32} \dst{r0}, \src{r1}, \src{r2}, \src{r3}} (\textbf{SPIR-V:} \texttt{SSA Result ID}; \textbf{GLSL:} \texttt{float r0 = fma(r1, r2, r3);})
  \item \hyperref[enum:dest_width_m]{\texttt{dest\_width}}: Stage C destination lane write-mask (\texttt{none} = full 32-bit, \texttt{h0} = low 16-bit, \texttt{h1} = high 16-bit without bit-shifting). \emph{Example:} {\scriptsize\ttfamily \mnem{FADD.v2f16} \dst{r0}\imod{.h1}, \src{r1}, \src{r2}} (\textbf{SPIR-V:} \texttt{OpVectorInsertDynamic / WriteMask}; \textbf{GLSL:} \texttt{dst.y = a.y + b.y;})
  \item \texttt{opcode1}: Primary sub-opcode within the instruction class uniquely identifying the operation family. \emph{Example:} {\scriptsize\ttfamily \mnem{IADD.s32} \dst{r0}, \src{r1}, \src{r2}} (\textbf{SPIR-V:} \texttt{OpIAdd (operation selector)}; \textbf{GLSL:} \texttt{// Primary opcode bits})
  \item \texttt{instruction\_class}: 4-bit major instruction class selector (\texttt{bits 48:51}) partitioning the 64-bit Valhall opcode space (e.g., \texttt{FMA}, \texttt{CVT}, \texttt{SFU}, \texttt{LS}, \texttt{TEX}, \texttt{BRANCH}). \emph{Example:} {\scriptsize\ttfamily \mnem{FADD.f32} \dst{r0}, \src{r1}, \src{r2}} (\textbf{SPIR-V:} \texttt{N/A (encoding class)}; \textbf{GLSL:} \texttt{// Hardware execution pipe / decode class})
  \item \texttt{fau\_page\_index}: Selects the FAU (Fast Access Uniform) page (\texttt{0..3}) when a source operand references hardware small constants, thread IDs, or lane constants. \emph{Example:} {\scriptsize\ttfamily \mnem{FADD.f32} \dst{r0}, \src{r1}, \src{k1f}} (\textbf{SPIR-V:} \texttt{OpConstant / BuiltIn}; \textbf{GLSL:} \texttt{float y = x + 1.0;})
  \item \hyperref[enum:flow_control_m]{\texttt{flow}}: 4-bit embedded flow-control \& scoreboard wait modifier (\texttt{bits 59:62}: \texttt{none}, \texttt{wait0..wait0126}, \texttt{wait}, \texttt{reconverge}, \texttt{discard}, \texttt{end}). \emph{Example:} {\scriptsize\ttfamily \mnem{FADD.f32}\imod{.wait0.reconverge} \dst{r0}, \src{r1}, \src{r2}} (\textbf{SPIR-V:} \texttt{OpControlBarrier / OpReturn}; \textbf{GLSL:} \texttt{// Scoreboard wait + warp reconvergence})
\end{tightitemize}
\clearpage

\subsection*{FEXP\_TABLE}
\addcontentsline{toc}{subsection}{FEXP\_TABLE}
16-entry exponential lookup table indexed by the low 4 bits (\texttt{u4}) of \texttt{src0}: writes \texttt{2\textasciicircum{}(idx / 16)} in FP32 to \texttt{dest} by default, or the low-order correction table entry when \texttt{.low} (\texttt{adj}) is set.
\par\textbf{Notes:}
\begin{tightitemize}
  \item Legacy SFU table helper (\texttt{bifrost/valhall/ISA.xml:2346}, \texttt{compiler.h:684}).
  \item \textbf{Source references:} \texttt{bifrost/valhall/ISA.xml:2346}, \texttt{bifrost/compiler.h:684}, \texttt{isa-v9-v14.xml}.
\end{tightitemize}
\codebox{\mnem{FEXP\_TABLE}\imod{\{.flow\}}\imod{\{.adj\}} \dst{dest}\imod{\{.dest\_width\}}, \src{src0}}
\textbf{Operands:}
\begin{tightitemize}\footnotesize
  \item \texttt{\dst{dest}} (\texttt{dest\_register\_index}, \texttt{dest\_width}): FP32 destination GPR receiving \texttt{2\textasciicircum{}(idx / 16)} (or low-order correction with \texttt{.low}).
  \item \texttt{\src{src0}}: Input whose low 4 bits (\texttt{u4}) index the 16-entry exponential LUT.
\end{tightitemize}
\par\textbf{Closest SPIR-V Op:} \texttt{None (legacy SFU exp2 table)}
\par\textbf{Pseudo-GLSL (Operational Semantics):}
\codebox{uint idx = src0 \& 0xFu;\\dest = (adj == low) ? EXP2\_LOW\_CORRECTION\_LUT[idx] : exp2(float(idx) / 16.0);}
\textbf{Execution Context:}
\begin{tightitemize}\footnotesize
  \item \textbf{Compute stage preloads (\hyperref[sec:csf-spd]{\texttt{SPD\_k}}, \texttt{r16..r23}):} Hardware pre-fills \texttt{r55..r57} with \texttt{WorkGroupID.xyz + JOB\_OFFSET} (from \hyperref[sec:csf-regs]{\texttt{COMPUTESR::JOB\_OFFSET\_X/Y/Z}} \texttt{r34..r36}) and \texttt{r60..r62} with \texttt{LocalInvocationID.xyz} or \texttt{GlobalInvocationID.xyz} (from \texttt{WG\_SIZE} \texttt{r33} \& \texttt{JOB\_SIZE} \texttt{r37..r39}, selected by \texttt{ShaderProgram::preload}) so compute shaders can derive thread/element indices without memory loads.
  \item \textbf{Vertex stage preloads (\hyperref[sec:csf-spd]{\texttt{SPD\_0}/\texttt{SPD\_1}}, \texttt{r16..r19}):} The hardware Index Fetcher reads \hyperref[sec:csf-regs]{\texttt{IDVSSR::INDEX\_BUFFER}} (\texttt{r54:r55}) and pre-fills \texttt{r59, r60} with \texttt{VertexIndex} (\texttt{+ VERTEX\_OFFSET + GLOBAL\_ATTR\_OFF}) and \texttt{r61} with \texttt{InstanceIndex} (\texttt{+ INSTANCE\_OFFSET}) so Position and Varying shaders can pull vertex attributes directly.
  \item \textbf{Fragment stage preloads \& FP modes (\hyperref[sec:csf-spd]{\texttt{SPD\_2}}, \texttt{r20:r21}):} Hardware pre-fills \texttt{r59} with \texttt{floatBitsToUint(FragCoord.z)} and \texttt{r60} with the MSAA sample coverage mask (\texttt{gl\_SampleMaskIn[0]}), while \texttt{ShaderProgram} word~0 latches 32/64 GPR allocation and warp-wide \texttt{flush\_to\_zero\_mode}, \texttt{suppress\_na\_n}, and \texttt{suppress\_inf}.
  \item \textbf{Inline FAU \& Special Operands (\hyperref[sec:csf-tsd-fau]{\texttt{FAU\_k}}, \texttt{r8..r15}):} Any 8-bit source with \texttt{mode=2} reads zero-latency uniform \texttt{u0..u31} from \texttt{FAU\_k} (\texttt{slots64 = bits[63:56]}); \texttt{mode=3} reads \texttt{SmallConstant} (Table~1.2) or \texttt{FauSpecial} (\texttt{pc}, \texttt{lane\_id}, \texttt{lane\_width}, \texttt{tls\_ptr}, \texttt{wls\_ptr}, \texttt{blend\_descriptor\_0..7}).
\end{tightitemize}
\par\textbf{Worked Examples:}
\par\textbf{Example 1 (4-bit fractional `2\textasciicircum{}(idx/16)` lookup):}
\codebox{\mnem{FEXP\_TABLE.u4} \dst{r0}, \src{r1}}
With \texttt{r1 \& 0xF = 8}: looks up \texttt{2\textasciicircum{}(8/16) = sqrt(2) \textasciitilde{}= 1.4142135f} (\texttt{0x3FB504F3}).
\par\textbf{Example 2 (Adjusted low-precision entry (.low)):}
\codebox{\mnem{FEXP\_TABLE.u4}\imod{.low} \dst{r0}, \src{r1}}
With \texttt{.low} and \texttt{r1 \& 0xF = 0}: returns the low-word correction term for extended-precision \texttt{exp2} reconstruction.

\begin{center}\footnotesize
\begin{tabularx}{\linewidth}{@{}lllX@{}}
\toprule
Variant & Arch & Class & Opcode Fields \\
\midrule
\texttt{(none)} & v9--v14 & \texttt{R1} & \texttt{opcode2=FEXP\_TABLE, opcode1=TABLE} \\
\bottomrule
\end{tabularx}
\end{center}
\begin{center}\footnotesize
\begin{tabularx}{0.88\linewidth}{@{}l l X@{}}
\toprule
Field Name & Bits & Type / Fixed Value \\
\midrule
\texttt{src0} & \texttt{0:7} & \texttt{SourceEncoding} \\
\textit{Reserved} & \texttt{8:15} & \texttt{zero} \\
\texttt{opcode2} & \texttt{16:20} & \texttt{exact=FEXP\_TABLE} \\
\textit{Reserved} & \texttt{21:30} & \texttt{zero} \\
\hyperref[enum:fexp_adj_m]{\texttt{adj}} & \texttt{31:32} & \texttt{fexp\_adj\_m (def none)} \\
\textit{Reserved} & \texttt{33:39} & \texttt{zero} \\
\texttt{dest\_register\_index} & \texttt{40:45} & \texttt{uint} \\
\hyperref[enum:dest_width_m]{\texttt{dest\_width}} & \texttt{46:47} & \texttt{dest\_width\_m} \\
\texttt{opcode1} & \texttt{48:51} & \texttt{exact=TABLE} \\
\texttt{instruction\_class} & \texttt{52:56} & \texttt{exact=R1} \\
\texttt{fau\_page\_index} & \texttt{57:58} & \texttt{uint} \\
\hyperref[enum:flow_control_m]{\texttt{flow}} & \texttt{59:62} & \texttt{flow\_control\_m} \\
\textit{Reserved} & \texttt{63} & \texttt{zero} \\
\bottomrule
\end{tabularx}
\end{center}
\textbf{Fields Explained:}
\begin{tightitemize}\footnotesize
  \item \texttt{src0}: Primary source operand (\texttt{r0..r63}, \texttt{u0..u63} uniform, \texttt{k*} FAU constant, or \texttt{t0..t2} temporary; \texttt{\textasciicircum{}} marks discard/last-use). \emph{Example:} {\scriptsize\ttfamily \mnem{FADD.f32} \dst{r0}, \src{r1}, \src{r2}} (\textbf{SPIR-V:} \texttt{OpFAdd (Operand 0)}; \textbf{GLSL:} \texttt{float c = a + b;})
  \item \texttt{opcode2}: Secondary sub-opcode further specializing instructions sharing \texttt{instruction\_class} and \texttt{opcode1} (e.g., conversions, SFU tables, moves). \emph{Example:} {\scriptsize\ttfamily \mnem{F32\_TO\_S32}\imod{.round\_zero} \dst{r0}, \src{r1}} (\textbf{SPIR-V:} \texttt{OpConvertFToS}; \textbf{GLSL:} \texttt{int y = int(x);})
  \item \hyperref[enum:fexp_adj_m]{\texttt{adj}}: Selects low/adjusted table entry in \texttt{FEXP\_TABLE}. \emph{Example:} {\scriptsize\ttfamily \mnem{FEXP\_TABLE}\imod{.low} \dst{r0}, \src{r1}} (\textbf{SPIR-V:} \texttt{GLSL.std.450 Exp2}; \textbf{GLSL:} \texttt{float y = exp2(x);})
  \item \texttt{dest\_register\_index}: 6-bit destination GPR index (\texttt{r0..r63}) receiving the Stage C writeback result (even-aligned pair \texttt{[rN:rN+1]} for 64-bit ops). \emph{Example:} {\scriptsize\ttfamily \mnem{FMA.f32} \dst{r0}, \src{r1}, \src{r2}, \src{r3}} (\textbf{SPIR-V:} \texttt{SSA Result ID}; \textbf{GLSL:} \texttt{float r0 = fma(r1, r2, r3);})
  \item \hyperref[enum:dest_width_m]{\texttt{dest\_width}}: Stage C destination lane write-mask (\texttt{none} = full 32-bit, \texttt{h0} = low 16-bit, \texttt{h1} = high 16-bit without bit-shifting). \emph{Example:} {\scriptsize\ttfamily \mnem{FADD.v2f16} \dst{r0}\imod{.h1}, \src{r1}, \src{r2}} (\textbf{SPIR-V:} \texttt{OpVectorInsertDynamic / WriteMask}; \textbf{GLSL:} \texttt{dst.y = a.y + b.y;})
  \item \texttt{opcode1}: Primary sub-opcode within the instruction class uniquely identifying the operation family. \emph{Example:} {\scriptsize\ttfamily \mnem{IADD.s32} \dst{r0}, \src{r1}, \src{r2}} (\textbf{SPIR-V:} \texttt{OpIAdd (operation selector)}; \textbf{GLSL:} \texttt{// Primary opcode bits})
  \item \texttt{instruction\_class}: 4-bit major instruction class selector (\texttt{bits 48:51}) partitioning the 64-bit Valhall opcode space (e.g., \texttt{FMA}, \texttt{CVT}, \texttt{SFU}, \texttt{LS}, \texttt{TEX}, \texttt{BRANCH}). \emph{Example:} {\scriptsize\ttfamily \mnem{FADD.f32} \dst{r0}, \src{r1}, \src{r2}} (\textbf{SPIR-V:} \texttt{N/A (encoding class)}; \textbf{GLSL:} \texttt{// Hardware execution pipe / decode class})
  \item \texttt{fau\_page\_index}: Selects the FAU (Fast Access Uniform) page (\texttt{0..3}) when a source operand references hardware small constants, thread IDs, or lane constants. \emph{Example:} {\scriptsize\ttfamily \mnem{FADD.f32} \dst{r0}, \src{r1}, \src{k1f}} (\textbf{SPIR-V:} \texttt{OpConstant / BuiltIn}; \textbf{GLSL:} \texttt{float y = x + 1.0;})
  \item \hyperref[enum:flow_control_m]{\texttt{flow}}: 4-bit embedded flow-control \& scoreboard wait modifier (\texttt{bits 59:62}: \texttt{none}, \texttt{wait0..wait0126}, \texttt{wait}, \texttt{reconverge}, \texttt{discard}, \texttt{end}). \emph{Example:} {\scriptsize\ttfamily \mnem{FADD.f32}\imod{.wait0.reconverge} \dst{r0}, \src{r1}, \src{r2}} (\textbf{SPIR-V:} \texttt{OpControlBarrier / OpReturn}; \textbf{GLSL:} \texttt{// Scoreboard wait + warp reconvergence})
\end{tightitemize}
\clearpage

\subsection*{FLOGD}
\addcontentsline{toc}{subsection}{FLOGD}
Logarithm divided-difference helper (\texttt{f32}): reduces \texttt{src0} to its mantissa $m \in [0.75, 1.5)$ and writes \texttt{dest = log2(m) / (m - 1)} (with the removable singularity at $m = 1$ evaluating to $1 / \ln 2$). Multiplied by $(m - 1)$ from \texttt{FADD\_LSCALE} and added to \texttt{FREXPE.log} to compute \texttt{log2(x)} without cancellation near 1.
\par\textbf{Notes:}
\begin{tightitemize}
  \item Paired with \texttt{FREXPE.log} + \texttt{FADD\_LSCALE(-1.0, x)} + \texttt{FMA} to compute \texttt{log2(x)} to \textless{}= 3 ULP (\texttt{builder.rs:387-418}, \texttt{hw\_tests.rs:2348-2393}).
  \item \textbf{Source references:} \texttt{kraid/ops.rs:1658-1677}, \texttt{encode\_v9.rs:1421-1438}, \texttt{builder.rs:404-417}, \texttt{hw\_tests.rs:2348-2393}.
\end{tightitemize}
\codebox{\mnem{FLOGD}\imod{\{.flow\}} \dst{dest}\imod{\{.dest\_width\}}, \src{src0}\imod{\{.widen0\}}\imod{\{.abs0\}}\imod{\{.neg0\}}}
\textbf{Operands:}
\begin{tightitemize}\footnotesize
  \item \texttt{\dst{dest}} (\texttt{dest\_register\_index}, \texttt{dest\_width}): FP32 destination GPR receiving \texttt{log2(m) / (m - 1)} for reduced mantissa \texttt{m}.
  \item \texttt{\src{src0}}: FP32 input whose mantissa is reduced into \texttt{[0.75, 1.5)}.
\end{tightitemize}
\par\textbf{Closest SPIR-V Op:} \texttt{None (lowered from GLSL.std.450 Log2)}
\par\textbf{Pseudo-GLSL (Operational Semantics):}
\codebox{float m = reduceTo\_0\_75\_to\_1\_5(src0);\\dest = (m == 1.0) ? (1.0 / log(2.0)) : (log2(m) / (m - 1.0));}
\textbf{Execution Context:}
\begin{tightitemize}\footnotesize
  \item \textbf{Compute stage preloads (\hyperref[sec:csf-spd]{\texttt{SPD\_k}}, \texttt{r16..r23}):} Hardware pre-fills \texttt{r55..r57} with \texttt{WorkGroupID.xyz + JOB\_OFFSET} (from \hyperref[sec:csf-regs]{\texttt{COMPUTESR::JOB\_OFFSET\_X/Y/Z}} \texttt{r34..r36}) and \texttt{r60..r62} with \texttt{LocalInvocationID.xyz} or \texttt{GlobalInvocationID.xyz} (from \texttt{WG\_SIZE} \texttt{r33} \& \texttt{JOB\_SIZE} \texttt{r37..r39}, selected by \texttt{ShaderProgram::preload}) so compute shaders can derive thread/element indices without memory loads.
  \item \textbf{Vertex stage preloads (\hyperref[sec:csf-spd]{\texttt{SPD\_0}/\texttt{SPD\_1}}, \texttt{r16..r19}):} The hardware Index Fetcher reads \hyperref[sec:csf-regs]{\texttt{IDVSSR::INDEX\_BUFFER}} (\texttt{r54:r55}) and pre-fills \texttt{r59, r60} with \texttt{VertexIndex} (\texttt{+ VERTEX\_OFFSET + GLOBAL\_ATTR\_OFF}) and \texttt{r61} with \texttt{InstanceIndex} (\texttt{+ INSTANCE\_OFFSET}) so Position and Varying shaders can pull vertex attributes directly.
  \item \textbf{Fragment stage preloads \& FP modes (\hyperref[sec:csf-spd]{\texttt{SPD\_2}}, \texttt{r20:r21}):} Hardware pre-fills \texttt{r59} with \texttt{floatBitsToUint(FragCoord.z)} and \texttt{r60} with the MSAA sample coverage mask (\texttt{gl\_SampleMaskIn[0]}), while \texttt{ShaderProgram} word~0 latches 32/64 GPR allocation and warp-wide \texttt{flush\_to\_zero\_mode}, \texttt{suppress\_na\_n}, and \texttt{suppress\_inf}.
  \item \textbf{Inline FAU \& Special Operands (\hyperref[sec:csf-tsd-fau]{\texttt{FAU\_k}}, \texttt{r8..r15}):} Any 8-bit source with \texttt{mode=2} reads zero-latency uniform \texttt{u0..u31} from \texttt{FAU\_k} (\texttt{slots64 = bits[63:56]}); \texttt{mode=3} reads \texttt{SmallConstant} (Table~1.2) or \texttt{FauSpecial} (\texttt{pc}, \texttt{lane\_id}, \texttt{lane\_width}, \texttt{tls\_ptr}, \texttt{wls\_ptr}, \texttt{blend\_descriptor\_0..7}).
\end{tightitemize}
\par\textbf{Worked Examples:}
\par\textbf{Example 1 (Reduced mantissa `m = 1.0f`):}
\codebox{\mnem{FLOGD.f32} \dst{r0}, \src{r1}}
With \texttt{r1 = 1.0f}: evaluates \texttt{lim\_\{m-\textgreater{}1\} log2(m)/(m - 1) = 1/ln(2) \textasciitilde{}= 1.442695f} (\texttt{0x3FB8AA3B}).
\par\textbf{Example 2 (Upper interval edge `m = 1.5f`):}
\codebox{\mnem{FLOGD.f32} \dst{r0}, \src{r1}}
With \texttt{r1 = 1.5f} (or \texttt{6.0f}, reduced to \texttt{0.75f}): evaluates 16-entry cubic LUT \texttt{log2(m)/(m - 1)} for \texttt{FMA} with \texttt{FADD\_LSCALE}.
\par\textbf{Mined SPIR-V $\to$ Valhall ISA Example (Logarithm divided-difference helper in GLSL Log2 --- Fragment shader, Mali-G710 (v10)):}
\codebox{\textbf{; SPIR-V (input)}\\\%y = OpExtInst \%float \%glsl Log2 \%x\\[3pt]\textbf{; Valhall ISA (malioc output)}\\FREXPE.f32.log r4, r2\\S32\_TO\_F32 r4, r4\textasciicircum{}\\FLOGD.f32 r6, r2\\FADD\_LSCALE.f32 r2, k1f.neg, r2\textasciicircum{}\\FMA.f32 r2, r6\textasciicircum{}, r2\textasciicircum{}, r4\textasciicircum{}}
\texttt{FLOGD.f32} evaluates $\log_2(m)/(m-1)$ for the mantissa $m \in [0.75, 1.5)$ of \texttt{r2}, avoiding catastrophic cancellation near $m = 1$ when multiplied by $(m - 1)$ from \texttt{FADD\_LSCALE}.
\begin{center}\footnotesize
\begin{tabularx}{\linewidth}{@{}lllX@{}}
\toprule
Variant & Arch & Class & Opcode Fields \\
\midrule
\texttt{(none)} & v9--v14 & \texttt{R1} & \texttt{opcode2=FLOGD\_F32, opcode1=TABLE} \\
\bottomrule
\end{tabularx}
\end{center}
\begin{center}\footnotesize
\begin{tabularx}{0.88\linewidth}{@{}l l X@{}}
\toprule
Field Name & Bits & Type / Fixed Value \\
\midrule
\texttt{src0} & \texttt{0:7} & \texttt{SourceEncoding} \\
\textit{Reserved} & \texttt{8:15} & \texttt{zero} \\
\texttt{opcode2} & \texttt{16:20} & \texttt{exact=FLOGD\_F32} \\
\textit{Reserved} & \texttt{21:27} & \texttt{zero} \\
\hyperref[enum:widen_m]{\texttt{widen0}} & \texttt{28:29} & \texttt{widen\_m (def none)} \\
\textit{Reserved} & \texttt{30:37} & \texttt{zero} \\
\hyperref[enum:neg_m]{\texttt{neg0}} & \texttt{38} & \texttt{neg\_m (def none)} \\
\hyperref[enum:abs_m]{\texttt{abs0}} & \texttt{39} & \texttt{abs\_m (def none)} \\
\texttt{dest\_register\_index} & \texttt{40:45} & \texttt{uint} \\
\hyperref[enum:dest_width_m]{\texttt{dest\_width}} & \texttt{46:47} & \texttt{dest\_width\_m} \\
\texttt{opcode1} & \texttt{48:51} & \texttt{exact=TABLE} \\
\texttt{instruction\_class} & \texttt{52:56} & \texttt{exact=R1} \\
\texttt{fau\_page\_index} & \texttt{57:58} & \texttt{uint} \\
\hyperref[enum:flow_control_m]{\texttt{flow}} & \texttt{59:62} & \texttt{flow\_control\_m} \\
\textit{Reserved} & \texttt{63} & \texttt{zero} \\
\bottomrule
\end{tabularx}
\end{center}
\textbf{Fields Explained:}
\begin{tightitemize}\footnotesize
  \item \texttt{src0}: Primary source operand (\texttt{r0..r63}, \texttt{u0..u63} uniform, \texttt{k*} FAU constant, or \texttt{t0..t2} temporary; \texttt{\textasciicircum{}} marks discard/last-use). \emph{Example:} {\scriptsize\ttfamily \mnem{FADD.f32} \dst{r0}, \src{r1}, \src{r2}} (\textbf{SPIR-V:} \texttt{OpFAdd (Operand 0)}; \textbf{GLSL:} \texttt{float c = a + b;})
  \item \texttt{opcode2}: Secondary sub-opcode further specializing instructions sharing \texttt{instruction\_class} and \texttt{opcode1} (e.g., conversions, SFU tables, moves). \emph{Example:} {\scriptsize\ttfamily \mnem{F32\_TO\_S32}\imod{.round\_zero} \dst{r0}, \src{r1}} (\textbf{SPIR-V:} \texttt{OpConvertFToS}; \textbf{GLSL:} \texttt{int y = int(x);})
  \item \hyperref[enum:widen_m]{\texttt{widen0}}: Stage A widening extractor on \texttt{src0} (promotes FP16 \texttt{h0}/\texttt{h1} to FP32, or sign/zero-extends \texttt{b0..b3}/\texttt{h0..h1}/\texttt{w0..w1} to 32/64-bit before \texttt{abs}/\texttt{neg}). \emph{Example:} {\scriptsize\ttfamily \mnem{IADD.s32} \dst{r0}, \src{r1}\imod{.h0}, \src{r2}} (\textbf{SPIR-V:} \texttt{OpSConvert / OpFConvert}; \textbf{GLSL:} \texttt{int c = int(int16\_t(a)) + b;})
  \item \hyperref[enum:neg_m]{\texttt{neg0}}: Stage A floating-point negation on \texttt{src0} (flips IEEE sign bit after swizzle/widen and \texttt{abs0}). \emph{Example:} {\scriptsize\ttfamily \mnem{FADD.f32} \dst{r0}, \src{r1}\imod{.neg}, \src{r2}} (\textbf{SPIR-V:} \texttt{OpFNegate}; \textbf{GLSL:} \texttt{float c = -a + b;})
  \item \hyperref[enum:abs_m]{\texttt{abs0}}: Stage A floating-point absolute value on \texttt{src0} (clears IEEE sign bit after swizzle/widen, before \texttt{neg0}). \emph{Example:} {\scriptsize\ttfamily \mnem{FADD.f32} \dst{r0}, \src{r1}\imod{.abs}, \src{r2}} (\textbf{SPIR-V:} \texttt{GLSL.std.450 FAbs}; \textbf{GLSL:} \texttt{float c = abs(a) + b;})
  \item \texttt{dest\_register\_index}: 6-bit destination GPR index (\texttt{r0..r63}) receiving the Stage C writeback result (even-aligned pair \texttt{[rN:rN+1]} for 64-bit ops). \emph{Example:} {\scriptsize\ttfamily \mnem{FMA.f32} \dst{r0}, \src{r1}, \src{r2}, \src{r3}} (\textbf{SPIR-V:} \texttt{SSA Result ID}; \textbf{GLSL:} \texttt{float r0 = fma(r1, r2, r3);})
  \item \hyperref[enum:dest_width_m]{\texttt{dest\_width}}: Stage C destination lane write-mask (\texttt{none} = full 32-bit, \texttt{h0} = low 16-bit, \texttt{h1} = high 16-bit without bit-shifting). \emph{Example:} {\scriptsize\ttfamily \mnem{FADD.v2f16} \dst{r0}\imod{.h1}, \src{r1}, \src{r2}} (\textbf{SPIR-V:} \texttt{OpVectorInsertDynamic / WriteMask}; \textbf{GLSL:} \texttt{dst.y = a.y + b.y;})
  \item \texttt{opcode1}: Primary sub-opcode within the instruction class uniquely identifying the operation family. \emph{Example:} {\scriptsize\ttfamily \mnem{IADD.s32} \dst{r0}, \src{r1}, \src{r2}} (\textbf{SPIR-V:} \texttt{OpIAdd (operation selector)}; \textbf{GLSL:} \texttt{// Primary opcode bits})
  \item \texttt{instruction\_class}: 4-bit major instruction class selector (\texttt{bits 48:51}) partitioning the 64-bit Valhall opcode space (e.g., \texttt{FMA}, \texttt{CVT}, \texttt{SFU}, \texttt{LS}, \texttt{TEX}, \texttt{BRANCH}). \emph{Example:} {\scriptsize\ttfamily \mnem{FADD.f32} \dst{r0}, \src{r1}, \src{r2}} (\textbf{SPIR-V:} \texttt{N/A (encoding class)}; \textbf{GLSL:} \texttt{// Hardware execution pipe / decode class})
  \item \texttt{fau\_page\_index}: Selects the FAU (Fast Access Uniform) page (\texttt{0..3}) when a source operand references hardware small constants, thread IDs, or lane constants. \emph{Example:} {\scriptsize\ttfamily \mnem{FADD.f32} \dst{r0}, \src{r1}, \src{k1f}} (\textbf{SPIR-V:} \texttt{OpConstant / BuiltIn}; \textbf{GLSL:} \texttt{float y = x + 1.0;})
  \item \hyperref[enum:flow_control_m]{\texttt{flow}}: 4-bit embedded flow-control \& scoreboard wait modifier (\texttt{bits 59:62}: \texttt{none}, \texttt{wait0..wait0126}, \texttt{wait}, \texttt{reconverge}, \texttt{discard}, \texttt{end}). \emph{Example:} {\scriptsize\ttfamily \mnem{FADD.f32}\imod{.wait0.reconverge} \dst{r0}, \src{r1}, \src{r2}} (\textbf{SPIR-V:} \texttt{OpControlBarrier / OpReturn}; \textbf{GLSL:} \texttt{// Scoreboard wait + warp reconvergence})
\end{tightitemize}
\clearpage

\subsection*{FLOG\_TABLE}
\addcontentsline{toc}{subsection}{FLOG\_TABLE}
Logarithm lookup table (\texttt{f32}): returns table coefficient for \texttt{log2}/\texttt{ln}/reduced log based on \texttt{mode} (\texttt{red}/\texttt{base2}/\texttt{natural}) and \texttt{precision} (\texttt{none}/\texttt{high}/\texttt{low}).
\par\textbf{Notes:}
\begin{tightitemize}
  \item Used in Bifrost/Valhall legacy log sequences (\texttt{bi\_opt\_mod\_props.c:42,61}).
  \item \textbf{Source references:} \texttt{bifrost/valhall/ISA.xml:2298}, \texttt{bifrost/compiler.h:781-782}, \texttt{isa-v9-v14.xml}.
\end{tightitemize}
\codebox{\mnem{FLOG\_TABLE}\imod{\{.flow\}}\imod{\{.mode\}}\imod{\{.precision\}} \dst{dest}\imod{\{.dest\_width\}}, \src{src0}\imod{\{.widen0\}}\imod{\{.abs0\}}\imod{\{.neg0\}}}
\textbf{Operands:}
\begin{tightitemize}\footnotesize
  \item \texttt{\dst{dest}} (\texttt{dest\_register\_index}, \texttt{dest\_width}): FP32 destination GPR receiving the polynomial coefficient from the log LUT.
  \item \texttt{\src{src0}}: FP32 input whose top mantissa bits index the table.
\end{tightitemize}
\par\textbf{Closest SPIR-V Op:} \texttt{None (legacy SFU log table)}
\par\textbf{Pseudo-GLSL (Operational Semantics):}
\codebox{uint idx = topMantissaBits(src0);\\dest = (base == natural) ? LN\_POLY\_LUT[precision][idx] : LOG2\_POLY\_LUT[precision][idx];}
\textbf{Execution Context:}
\begin{tightitemize}\footnotesize
  \item \textbf{Compute stage preloads (\hyperref[sec:csf-spd]{\texttt{SPD\_k}}, \texttt{r16..r23}):} Hardware pre-fills \texttt{r55..r57} with \texttt{WorkGroupID.xyz + JOB\_OFFSET} (from \hyperref[sec:csf-regs]{\texttt{COMPUTESR::JOB\_OFFSET\_X/Y/Z}} \texttt{r34..r36}) and \texttt{r60..r62} with \texttt{LocalInvocationID.xyz} or \texttt{GlobalInvocationID.xyz} (from \texttt{WG\_SIZE} \texttt{r33} \& \texttt{JOB\_SIZE} \texttt{r37..r39}, selected by \texttt{ShaderProgram::preload}) so compute shaders can derive thread/element indices without memory loads.
  \item \textbf{Vertex stage preloads (\hyperref[sec:csf-spd]{\texttt{SPD\_0}/\texttt{SPD\_1}}, \texttt{r16..r19}):} The hardware Index Fetcher reads \hyperref[sec:csf-regs]{\texttt{IDVSSR::INDEX\_BUFFER}} (\texttt{r54:r55}) and pre-fills \texttt{r59, r60} with \texttt{VertexIndex} (\texttt{+ VERTEX\_OFFSET + GLOBAL\_ATTR\_OFF}) and \texttt{r61} with \texttt{InstanceIndex} (\texttt{+ INSTANCE\_OFFSET}) so Position and Varying shaders can pull vertex attributes directly.
  \item \textbf{Fragment stage preloads \& FP modes (\hyperref[sec:csf-spd]{\texttt{SPD\_2}}, \texttt{r20:r21}):} Hardware pre-fills \texttt{r59} with \texttt{floatBitsToUint(FragCoord.z)} and \texttt{r60} with the MSAA sample coverage mask (\texttt{gl\_SampleMaskIn[0]}), while \texttt{ShaderProgram} word~0 latches 32/64 GPR allocation and warp-wide \texttt{flush\_to\_zero\_mode}, \texttt{suppress\_na\_n}, and \texttt{suppress\_inf}.
  \item \textbf{Inline FAU \& Special Operands (\hyperref[sec:csf-tsd-fau]{\texttt{FAU\_k}}, \texttt{r8..r15}):} Any 8-bit source with \texttt{mode=2} reads zero-latency uniform \texttt{u0..u31} from \texttt{FAU\_k} (\texttt{slots64 = bits[63:56]}); \texttt{mode=3} reads \texttt{SmallConstant} (Table~1.2) or \texttt{FauSpecial} (\texttt{pc}, \texttt{lane\_id}, \texttt{lane\_width}, \texttt{tls\_ptr}, \texttt{wls\_ptr}, \texttt{blend\_descriptor\_0..7}).
\end{tightitemize}
\par\textbf{Worked Examples:}
\par\textbf{Example 1 (Base-2 log LUT entry):}
\codebox{\mnem{FLOG\_TABLE.f32}\imod{.base2} \dst{r0}, \src{r1}}
Indexes the base-2 logarithm coefficient table using the top mantissa bits of \texttt{r1}.
\par\textbf{Example 2 (Natural log high-precision LUT entry):}
\codebox{\mnem{FLOG\_TABLE.f32}\imod{.natural.high} \dst{r0}, \src{r1}}
Selects the high-precision natural-log (\texttt{ln}) table coefficient for \texttt{log(x)} lowering.

\begin{center}\footnotesize
\begin{tabularx}{\linewidth}{@{}lllX@{}}
\toprule
Variant & Arch & Class & Opcode Fields \\
\midrule
\texttt{(none)} & v9--v14 & \texttt{R1} & \texttt{opcode2=FLOG\_TABLE, opcode1=TABLE} \\
\bottomrule
\end{tabularx}
\end{center}
\begin{center}\footnotesize
\begin{tabularx}{0.88\linewidth}{@{}l l X@{}}
\toprule
Field Name & Bits & Type / Fixed Value \\
\midrule
\texttt{src0} & \texttt{0:7} & \texttt{SourceEncoding} \\
\textit{Reserved} & \texttt{8:15} & \texttt{zero} \\
\texttt{opcode2} & \texttt{16:20} & \texttt{exact=FLOG\_TABLE} \\
\textit{Reserved} & \texttt{21:27} & \texttt{zero} \\
\hyperref[enum:widen_m]{\texttt{widen0}} & \texttt{28:29} & \texttt{widen\_m (def none)} \\
\textit{Reserved} & \texttt{30} & \texttt{zero} \\
\hyperref[enum:table_flog_mode_m]{\texttt{mode}} & \texttt{31:32} & \texttt{table\_flog\_mode\_m} \\
\hyperref[enum:table_flog_precision_m]{\texttt{precision}} & \texttt{33:34} & \texttt{table\_flog\_precision\_m} \\
\textit{Reserved} & \texttt{35:37} & \texttt{zero} \\
\hyperref[enum:neg_m]{\texttt{neg0}} & \texttt{38} & \texttt{neg\_m (def none)} \\
\hyperref[enum:abs_m]{\texttt{abs0}} & \texttt{39} & \texttt{abs\_m (def none)} \\
\texttt{dest\_register\_index} & \texttt{40:45} & \texttt{uint} \\
\hyperref[enum:dest_width_m]{\texttt{dest\_width}} & \texttt{46:47} & \texttt{dest\_width\_m} \\
\texttt{opcode1} & \texttt{48:51} & \texttt{exact=TABLE} \\
\texttt{instruction\_class} & \texttt{52:56} & \texttt{exact=R1} \\
\texttt{fau\_page\_index} & \texttt{57:58} & \texttt{uint} \\
\hyperref[enum:flow_control_m]{\texttt{flow}} & \texttt{59:62} & \texttt{flow\_control\_m} \\
\textit{Reserved} & \texttt{63} & \texttt{zero} \\
\bottomrule
\end{tabularx}
\end{center}
\textbf{Fields Explained:}
\begin{tightitemize}\footnotesize
  \item \texttt{src0}: Primary source operand (\texttt{r0..r63}, \texttt{u0..u63} uniform, \texttt{k*} FAU constant, or \texttt{t0..t2} temporary; \texttt{\textasciicircum{}} marks discard/last-use). \emph{Example:} {\scriptsize\ttfamily \mnem{FADD.f32} \dst{r0}, \src{r1}, \src{r2}} (\textbf{SPIR-V:} \texttt{OpFAdd (Operand 0)}; \textbf{GLSL:} \texttt{float c = a + b;})
  \item \texttt{opcode2}: Secondary sub-opcode further specializing instructions sharing \texttt{instruction\_class} and \texttt{opcode1} (e.g., conversions, SFU tables, moves). \emph{Example:} {\scriptsize\ttfamily \mnem{F32\_TO\_S32}\imod{.round\_zero} \dst{r0}, \src{r1}} (\textbf{SPIR-V:} \texttt{OpConvertFToS}; \textbf{GLSL:} \texttt{int y = int(x);})
  \item \hyperref[enum:widen_m]{\texttt{widen0}}: Stage A widening extractor on \texttt{src0} (promotes FP16 \texttt{h0}/\texttt{h1} to FP32, or sign/zero-extends \texttt{b0..b3}/\texttt{h0..h1}/\texttt{w0..w1} to 32/64-bit before \texttt{abs}/\texttt{neg}). \emph{Example:} {\scriptsize\ttfamily \mnem{IADD.s32} \dst{r0}, \src{r1}\imod{.h0}, \src{r2}} (\textbf{SPIR-V:} \texttt{OpSConvert / OpFConvert}; \textbf{GLSL:} \texttt{int c = int(int16\_t(a)) + b;})
  \item \hyperref[enum:table_flog_mode_m]{\texttt{mode}}: Selects logarithm table base (\texttt{red}, \texttt{base2}, \texttt{natural}) in \texttt{FLOG\_TABLE}. \emph{Example:} {\scriptsize\ttfamily \mnem{FLOG\_TABLE}\imod{.base2} \dst{r0}, \src{r1}} (\textbf{SPIR-V:} \texttt{GLSL.std.450 Log2}; \textbf{GLSL:} \texttt{float y = log2(x);})
  \item \hyperref[enum:table_flog_precision_m]{\texttt{precision}}: Selects high vs low precision table lookup in \texttt{FLOG\_TABLE}. \emph{Example:} {\scriptsize\ttfamily \mnem{FLOG\_TABLE}\imod{.high} \dst{r0}, \src{r1}} (\textbf{SPIR-V:} \texttt{GLSL.std.450 Log}; \textbf{GLSL:} \texttt{float y = log(x);})
  \item \hyperref[enum:neg_m]{\texttt{neg0}}: Stage A floating-point negation on \texttt{src0} (flips IEEE sign bit after swizzle/widen and \texttt{abs0}). \emph{Example:} {\scriptsize\ttfamily \mnem{FADD.f32} \dst{r0}, \src{r1}\imod{.neg}, \src{r2}} (\textbf{SPIR-V:} \texttt{OpFNegate}; \textbf{GLSL:} \texttt{float c = -a + b;})
  \item \hyperref[enum:abs_m]{\texttt{abs0}}: Stage A floating-point absolute value on \texttt{src0} (clears IEEE sign bit after swizzle/widen, before \texttt{neg0}). \emph{Example:} {\scriptsize\ttfamily \mnem{FADD.f32} \dst{r0}, \src{r1}\imod{.abs}, \src{r2}} (\textbf{SPIR-V:} \texttt{GLSL.std.450 FAbs}; \textbf{GLSL:} \texttt{float c = abs(a) + b;})
  \item \texttt{dest\_register\_index}: 6-bit destination GPR index (\texttt{r0..r63}) receiving the Stage C writeback result (even-aligned pair \texttt{[rN:rN+1]} for 64-bit ops). \emph{Example:} {\scriptsize\ttfamily \mnem{FMA.f32} \dst{r0}, \src{r1}, \src{r2}, \src{r3}} (\textbf{SPIR-V:} \texttt{SSA Result ID}; \textbf{GLSL:} \texttt{float r0 = fma(r1, r2, r3);})
  \item \hyperref[enum:dest_width_m]{\texttt{dest\_width}}: Stage C destination lane write-mask (\texttt{none} = full 32-bit, \texttt{h0} = low 16-bit, \texttt{h1} = high 16-bit without bit-shifting). \emph{Example:} {\scriptsize\ttfamily \mnem{FADD.v2f16} \dst{r0}\imod{.h1}, \src{r1}, \src{r2}} (\textbf{SPIR-V:} \texttt{OpVectorInsertDynamic / WriteMask}; \textbf{GLSL:} \texttt{dst.y = a.y + b.y;})
  \item \texttt{opcode1}: Primary sub-opcode within the instruction class uniquely identifying the operation family. \emph{Example:} {\scriptsize\ttfamily \mnem{IADD.s32} \dst{r0}, \src{r1}, \src{r2}} (\textbf{SPIR-V:} \texttt{OpIAdd (operation selector)}; \textbf{GLSL:} \texttt{// Primary opcode bits})
  \item \texttt{instruction\_class}: 4-bit major instruction class selector (\texttt{bits 48:51}) partitioning the 64-bit Valhall opcode space (e.g., \texttt{FMA}, \texttt{CVT}, \texttt{SFU}, \texttt{LS}, \texttt{TEX}, \texttt{BRANCH}). \emph{Example:} {\scriptsize\ttfamily \mnem{FADD.f32} \dst{r0}, \src{r1}, \src{r2}} (\textbf{SPIR-V:} \texttt{N/A (encoding class)}; \textbf{GLSL:} \texttt{// Hardware execution pipe / decode class})
  \item \texttt{fau\_page\_index}: Selects the FAU (Fast Access Uniform) page (\texttt{0..3}) when a source operand references hardware small constants, thread IDs, or lane constants. \emph{Example:} {\scriptsize\ttfamily \mnem{FADD.f32} \dst{r0}, \src{r1}, \src{k1f}} (\textbf{SPIR-V:} \texttt{OpConstant / BuiltIn}; \textbf{GLSL:} \texttt{float y = x + 1.0;})
  \item \hyperref[enum:flow_control_m]{\texttt{flow}}: 4-bit embedded flow-control \& scoreboard wait modifier (\texttt{bits 59:62}: \texttt{none}, \texttt{wait0..wait0126}, \texttt{wait}, \texttt{reconverge}, \texttt{discard}, \texttt{end}). \emph{Example:} {\scriptsize\ttfamily \mnem{FADD.f32}\imod{.wait0.reconverge} \dst{r0}, \src{r1}, \src{r2}} (\textbf{SPIR-V:} \texttt{OpControlBarrier / OpReturn}; \textbf{GLSL:} \texttt{// Scoreboard wait + warp reconvergence})
\end{tightitemize}
\clearpage

\subsection*{FMA\_RSCALE}
\addcontentsline{toc}{subsection}{FMA\_RSCALE}
Computes \texttt{dest = clamp(round\_rm((src0 * src1 + src2) * 2\textasciicircum{}clamp(src3, -512, 512)))} in exact FP64 with a single rounding to FP32, preventing intermediate overflow or underflow during transcendental range reduction. The \texttt{.special} modifier controls zero-times-infinity handling: \texttt{.n} forces \texttt{0 * anything -\textgreater{} $\pm$0}; \texttt{.left} returns \texttt{src0} on \texttt{0 * Inf}; \texttt{.n\_add} suppresses \texttt{(+Inf) + (-Inf)} NaNs.
\par\textbf{Notes:}
\begin{tightitemize}
  \item Only RNE (\texttt{none}) and RTZ (\texttt{round\_zero}) rounding.
  \item Verified in \texttt{libGLES\_mali.so} \texttt{\_mali\_fma\_po2\_scale\_sf32} (\texttt{0x2e28c20}) and \texttt{kraid/ops.rs:1828-1857}.
\end{tightitemize}
\codebox{\mnem{FMA\_RSCALE}\imod{\{.flow\}}\imod{\{.special\}}\imod{\{.round\}}\imod{\{.clamp\}} \dst{dest}\imod{\{.dest\_width\}}, \src{src0}\imod{\{.abs0\}}\imod{\{.neg0\}}, \src{src1}\imod{\{.abs1\}}\imod{\{.neg1\}}, \src{src2}\imod{\{.abs2\}}\imod{\{.neg2\}}, \src{src3}}
\textbf{Operands:}
\begin{tightitemize}\footnotesize
  \item \texttt{\dst{dest}} (\texttt{dest\_register\_index}, \texttt{dest\_width}): FP32 destination GPR receiving \texttt{(src0 * src1 + src2) * 2\textasciicircum{}clamp(src3, -512, 512)}.
  \item \texttt{\src{src0}}: First FP32 multiplicand.
  \item \texttt{\src{src1}}: Second FP32 multiplicand.
  \item \texttt{\src{src2}}: FP32 addend.
  \item \texttt{\src{src3}}: Signed integer exponent scale clamped to \texttt{[-512, 512]}.
\end{tightitemize}
\par\textbf{Closest SPIR-V Op:} \texttt{None (scaled FMA for Exp2/Sin/Cube)}
\par\textbf{Pseudo-GLSL (Operational Semantics):}
\codebox{double p = (special == n \&\& (src0 == 0.0 \textbar{}\textbar{} src1 == 0.0)) ? copysign(0.0, src0 * src1)\\                                                         : double(src0) * double(src1);\\int e = clamp(int(src3), -512, 512);\\dest = (special == left \&\& isZeroTimesInf(src0, src1)) ? src0\\                                                       : clamp\_mode(round\_rm(float((p + double(src2)) * exp2(double(e)))));}
\textbf{Execution Context:}
\begin{tightitemize}\footnotesize
  \item \textbf{Compute stage preloads (\hyperref[sec:csf-spd]{\texttt{SPD\_k}}, \texttt{r16..r23}):} Hardware pre-fills \texttt{r55..r57} with \texttt{WorkGroupID.xyz + JOB\_OFFSET} (from \hyperref[sec:csf-regs]{\texttt{COMPUTESR::JOB\_OFFSET\_X/Y/Z}} \texttt{r34..r36}) and \texttt{r60..r62} with \texttt{LocalInvocationID.xyz} or \texttt{GlobalInvocationID.xyz} (from \texttt{WG\_SIZE} \texttt{r33} \& \texttt{JOB\_SIZE} \texttt{r37..r39}, selected by \texttt{ShaderProgram::preload}) so compute shaders can derive thread/element indices without memory loads.
  \item \textbf{Vertex stage preloads (\hyperref[sec:csf-spd]{\texttt{SPD\_0}/\texttt{SPD\_1}}, \texttt{r16..r19}):} The hardware Index Fetcher reads \hyperref[sec:csf-regs]{\texttt{IDVSSR::INDEX\_BUFFER}} (\texttt{r54:r55}) and pre-fills \texttt{r59, r60} with \texttt{VertexIndex} (\texttt{+ VERTEX\_OFFSET + GLOBAL\_ATTR\_OFF}) and \texttt{r61} with \texttt{InstanceIndex} (\texttt{+ INSTANCE\_OFFSET}) so Position and Varying shaders can pull vertex attributes directly.
  \item \textbf{Fragment stage preloads \& FP modes (\hyperref[sec:csf-spd]{\texttt{SPD\_2}}, \texttt{r20:r21}):} Hardware pre-fills \texttt{r59} with \texttt{floatBitsToUint(FragCoord.z)} and \texttt{r60} with the MSAA sample coverage mask (\texttt{gl\_SampleMaskIn[0]}), while \texttt{ShaderProgram} word~0 latches 32/64 GPR allocation and warp-wide \texttt{flush\_to\_zero\_mode}, \texttt{suppress\_na\_n}, and \texttt{suppress\_inf}.
  \item \textbf{Inline FAU \& Special Operands (\hyperref[sec:csf-tsd-fau]{\texttt{FAU\_k}}, \texttt{r8..r15}):} Any 8-bit source with \texttt{mode=2} reads zero-latency uniform \texttt{u0..u31} from \texttt{FAU\_k} (\texttt{slots64 = bits[63:56]}); \texttt{mode=3} reads \texttt{SmallConstant} (Table~1.2) or \texttt{FauSpecial} (\texttt{pc}, \texttt{lane\_id}, \texttt{lane\_width}, \texttt{tls\_ptr}, \texttt{wls\_ptr}, \texttt{blend\_descriptor\_0..7}).
\end{tightitemize}
\par\textbf{Worked Examples:}
\par\textbf{Example 1 (Overflow-free intermediate product scaling):}
\codebox{\mnem{FMA\_RSCALE.f32} \dst{r0}, \src{r1}, \src{r2}, \src{r3}, \src{r4}}
With \texttt{r1 = 2\textasciicircum{}100, r2 = 2\textasciicircum{}50, r3 = 0, r4 = -150}: single-rounded FP64 evaluates \texttt{2\textasciicircum{}150 * 2\textasciicircum{}-150 = 1.0f} without intermediate FP32 overflow.
\par\textbf{Example 2 (Special `0 * Inf` override (.n)):}
\codebox{\mnem{FMA\_RSCALE.f32}\imod{.n} \dst{r0}, \src{k0}, \src{r\_inf}, \src{k1f}, \src{k0}}
With \texttt{0.0 * +Inf + 1.0}: standard FMA yields \texttt{NaN}, whereas \texttt{.n} forces the product to \texttt{0.0}, returning \texttt{1.0f}.
\par\textbf{Mined SPIR-V $\to$ Valhall ISA Example (FP64-intermediate multiply with power-of-two scaling (FMUL\_RSCALE alias) --- Fragment shader, Mali-G710 (v10)):}
\codebox{\textbf{; SPIR-V (input)}\\\%y = OpExtInst \%float \%glsl Exp2 \%x\\[3pt]\textbf{; Valhall ISA (malioc output)}\\FMUL\_RSCALE.f32 r57, r57\textasciicircum{}, u3.w1, r4\\F32\_TO\_S32 r5, r57\\FEXP.f32 r57, r5\textasciicircum{}, r57\textasciicircum{}}
\texttt{FMUL\_RSCALE.f32} (an alias of \texttt{FMA\_RSCALE} with \texttt{src2 = -0.0}) multiplies \texttt{r57 * u3.w1} and scales by $2^{r4}$ in exact FP64 with a single rounding to FP32 --- avoiding intermediate overflow/underflow before \texttt{FEXP.f32}.
\begin{center}\footnotesize
\begin{tabularx}{\linewidth}{@{}lllX@{}}
\toprule
Variant & Arch & Class & Opcode Fields \\
\midrule
\texttt{(none)} & v9--v14 & \texttt{R4\_FMAR} & \texttt{opcode1=FMA\_RSCALE\_F32} \\
\bottomrule
\end{tabularx}
\end{center}
\begin{center}\footnotesize
\begin{tabularx}{0.88\linewidth}{@{}l l X@{}}
\toprule
Field Name & Bits & Type / Fixed Value \\
\midrule
\texttt{src0} & \texttt{0:7} & \texttt{SourceEncoding} \\
\texttt{src1} & \texttt{8:15} & \texttt{SourceEncoding} \\
\texttt{src2} & \texttt{16:23} & \texttt{SourceEncoding} \\
\texttt{src3} & \texttt{24:31} & \texttt{SourceEncoding} \\
\hyperref[enum:clamp_m]{\texttt{clamp}} & \texttt{32:33} & \texttt{clamp\_m (def none)} \\
\hyperref[enum:neg_m]{\texttt{neg2}} & \texttt{34} & \texttt{neg\_m (def none)} \\
\hyperref[enum:abs_m]{\texttt{abs2}} & \texttt{35} & \texttt{abs\_m (def none)} \\
\hyperref[enum:neg_m]{\texttt{neg1}} & \texttt{36} & \texttt{neg\_m (def none)} \\
\hyperref[enum:abs_m]{\texttt{abs1}} & \texttt{37} & \texttt{abs\_m (def none)} \\
\hyperref[enum:neg_m]{\texttt{neg0}} & \texttt{38} & \texttt{neg\_m (def none)} \\
\hyperref[enum:abs_m]{\texttt{abs0}} & \texttt{39} & \texttt{abs\_m (def none)} \\
\texttt{dest\_register\_index} & \texttt{40:45} & \texttt{uint} \\
\hyperref[enum:dest_width_m]{\texttt{dest\_width}} & \texttt{46:47} & \texttt{dest\_width\_m} \\
\hyperref[enum:fma_rscale_special_32_m]{\texttt{special}} & \texttt{48:49} & \texttt{fma\_rscale\_special\_32\_m} \\
\hyperref[enum:fma_rscale_round_m]{\texttt{round}} & \texttt{50} & \texttt{fma\_rscale\_round\_m} \\
\texttt{opcode1} & \texttt{51} & \texttt{exact=FMA\_RSCALE\_F32} \\
\texttt{instruction\_class} & \texttt{52:56} & \texttt{exact=R4\_FMAR} \\
\texttt{fau\_page\_index} & \texttt{57:58} & \texttt{uint} \\
\hyperref[enum:flow_control_m]{\texttt{flow}} & \texttt{59:62} & \texttt{flow\_control\_m} \\
\textit{Reserved} & \texttt{63} & \texttt{zero} \\
\bottomrule
\end{tabularx}
\end{center}
\textbf{Fields Explained:}
\begin{tightitemize}\footnotesize
  \item \texttt{src0}: Primary source operand (\texttt{r0..r63}, \texttt{u0..u63} uniform, \texttt{k*} FAU constant, or \texttt{t0..t2} temporary; \texttt{\textasciicircum{}} marks discard/last-use). \emph{Example:} {\scriptsize\ttfamily \mnem{FADD.f32} \dst{r0}, \src{r1}, \src{r2}} (\textbf{SPIR-V:} \texttt{OpFAdd (Operand 0)}; \textbf{GLSL:} \texttt{float c = a + b;})
  \item \texttt{src1}: Secondary source operand (register, uniform, or FAU constant; in shifts, supplies the shift count byte). \emph{Example:} {\scriptsize\ttfamily \mnem{LSHIFT.i32} \dst{r0}, \src{r1}, \src{r2}\imod{.b0}} (\textbf{SPIR-V:} \texttt{OpShiftLeftLogical (Shift)}; \textbf{GLSL:} \texttt{uint y = a \textless{}\textless{} s;})
  \item \texttt{src2}: Third source operand (FMA addend, \texttt{MUX}/\texttt{CSEL} selector, \texttt{IADDC} carry-in, or bitwise mask in fused shift-logic). \emph{Example:} {\scriptsize\ttfamily \mnem{FMA.f32} \dst{r0}, \src{r1}, \src{r2}, \src{r3}} (\textbf{SPIR-V:} \texttt{GLSL.std.450 Fma (Addend)}; \textbf{GLSL:} \texttt{float d = fma(a, b, c);})
  \item \texttt{src3}: Fourth source operand (false-branch value in \texttt{CSEL}, exponent scale in \texttt{FMA\_RSCALE}, or fourth matrix/select input). \emph{Example:} {\scriptsize\ttfamily \mnem{CSEL.f32}\imod{.lt} \dst{r0}, \src{r1}, \src{r2}, \src{r3}, \src{r4}} (\textbf{SPIR-V:} \texttt{OpSelect (Object 2)}; \textbf{GLSL:} \texttt{float r = (a \textless{} b) ? x : y;})
  \item \hyperref[enum:clamp_m]{\texttt{clamp}}: Stage B post-op floating-point clamp (\texttt{none}, \texttt{clamp\_0\_inf}, \texttt{clamp\_m1\_1}, \texttt{clamp\_0\_1}) with \texttt{-0.0 \textless{} +0.0} and NaN suppression (\texttt{NaN -\textgreater{} lo}). \emph{Example:} {\scriptsize\ttfamily \mnem{FMA.f32}\imod{.clamp\_0\_1} \dst{r0}, \src{r1}, \src{r2}, \src{r3}} (\textbf{SPIR-V:} \texttt{GLSL.std.450 FClamp}; \textbf{GLSL:} \texttt{float y = clamp(fma(a, b, c), 0.0, 1.0);})
  \item \hyperref[enum:neg_m]{\texttt{neg2}}: Stage A floating-point negation on \texttt{src2} (flips IEEE sign bit of FMA addend after swizzle/widen and \texttt{abs2}). \emph{Example:} {\scriptsize\ttfamily \mnem{FMA.f32} \dst{r0}, \src{r1}, \src{r2}, \src{r3}\imod{.neg}} (\textbf{SPIR-V:} \texttt{OpFNegate}; \textbf{GLSL:} \texttt{float d = fma(a, b, -c);})
  \item \hyperref[enum:abs_m]{\texttt{abs2}}: Stage A floating-point absolute value on \texttt{src2} (clears IEEE sign bit of FMA addend after swizzle/widen). \emph{Example:} {\scriptsize\ttfamily \mnem{FMA.f32} \dst{r0}, \src{r1}, \src{r2}, \src{r3}\imod{.abs}} (\textbf{SPIR-V:} \texttt{GLSL.std.450 FAbs}; \textbf{GLSL:} \texttt{float d = fma(a, b, abs(c));})
  \item \hyperref[enum:neg_m]{\texttt{neg1}}: Stage A floating-point negation on \texttt{src1} (flips IEEE sign bit after swizzle/widen and \texttt{abs1}). \emph{Example:} {\scriptsize\ttfamily \mnem{FMA.f32} \dst{r0}, \src{r1}, \src{r2}\imod{.neg}, \src{r3}} (\textbf{SPIR-V:} \texttt{OpFNegate}; \textbf{GLSL:} \texttt{float d = fma(a, -b, c);})
  \item \hyperref[enum:abs_m]{\texttt{abs1}}: Stage A floating-point absolute value on \texttt{src1} (clears IEEE sign bit after swizzle/widen, before \texttt{neg1}). \emph{Example:} {\scriptsize\ttfamily \mnem{FMA.f32} \dst{r0}, \src{r1}, \src{r2}\imod{.abs}, \src{r3}} (\textbf{SPIR-V:} \texttt{GLSL.std.450 FAbs}; \textbf{GLSL:} \texttt{float d = fma(a, abs(b), c);})
  \item \hyperref[enum:neg_m]{\texttt{neg0}}: Stage A floating-point negation on \texttt{src0} (flips IEEE sign bit after swizzle/widen and \texttt{abs0}). \emph{Example:} {\scriptsize\ttfamily \mnem{FADD.f32} \dst{r0}, \src{r1}\imod{.neg}, \src{r2}} (\textbf{SPIR-V:} \texttt{OpFNegate}; \textbf{GLSL:} \texttt{float c = -a + b;})
  \item \hyperref[enum:abs_m]{\texttt{abs0}}: Stage A floating-point absolute value on \texttt{src0} (clears IEEE sign bit after swizzle/widen, before \texttt{neg0}). \emph{Example:} {\scriptsize\ttfamily \mnem{FADD.f32} \dst{r0}, \src{r1}\imod{.abs}, \src{r2}} (\textbf{SPIR-V:} \texttt{GLSL.std.450 FAbs}; \textbf{GLSL:} \texttt{float c = abs(a) + b;})
  \item \texttt{dest\_register\_index}: 6-bit destination GPR index (\texttt{r0..r63}) receiving the Stage C writeback result (even-aligned pair \texttt{[rN:rN+1]} for 64-bit ops). \emph{Example:} {\scriptsize\ttfamily \mnem{FMA.f32} \dst{r0}, \src{r1}, \src{r2}, \src{r3}} (\textbf{SPIR-V:} \texttt{SSA Result ID}; \textbf{GLSL:} \texttt{float r0 = fma(r1, r2, r3);})
  \item \hyperref[enum:dest_width_m]{\texttt{dest\_width}}: Stage C destination lane write-mask (\texttt{none} = full 32-bit, \texttt{h0} = low 16-bit, \texttt{h1} = high 16-bit without bit-shifting). \emph{Example:} {\scriptsize\ttfamily \mnem{FADD.v2f16} \dst{r0}\imod{.h1}, \src{r1}, \src{r2}} (\textbf{SPIR-V:} \texttt{OpVectorInsertDynamic / WriteMask}; \textbf{GLSL:} \texttt{dst.y = a.y + b.y;})
  \item \hyperref[enum:fma_rscale_special_32_m]{\texttt{special}}: Selects \texttt{FMA\_RSCALE}/\texttt{FREXP} special-case override (\texttt{n}, \texttt{left}, \texttt{n\_add}, \texttt{sqrt}, \texttt{log}). \emph{Example:} {\scriptsize\ttfamily \mnem{FREXPM.f32}\imod{.log} \dst{r0}, \src{r1}} (\textbf{SPIR-V:} \texttt{GLSL.std.450 Log2 (range reduction)}; \textbf{GLSL:} \texttt{float y = log2(x);})
  \item \hyperref[enum:fma_rscale_round_m]{\texttt{round}}: Rounding mode (\texttt{none}/RNE, \texttt{round\_up}, \texttt{round\_down}, \texttt{round\_zero}, \texttt{round\_na}) applied in FP ALU, conversions, \texttt{HADD}, \texttt{LDEXP}, and Stage C narrowing. \emph{Example:} {\scriptsize\ttfamily \mnem{FROUND.f32}\imod{.round\_down} \dst{r0}, \src{r1}} (\textbf{SPIR-V:} \texttt{GLSL.std.450 Floor / Ceil / Trunc}; \textbf{GLSL:} \texttt{float y = floor(x);})
  \item \texttt{opcode1}: Primary sub-opcode within the instruction class uniquely identifying the operation family. \emph{Example:} {\scriptsize\ttfamily \mnem{IADD.s32} \dst{r0}, \src{r1}, \src{r2}} (\textbf{SPIR-V:} \texttt{OpIAdd (operation selector)}; \textbf{GLSL:} \texttt{// Primary opcode bits})
  \item \texttt{instruction\_class}: 4-bit major instruction class selector (\texttt{bits 48:51}) partitioning the 64-bit Valhall opcode space (e.g., \texttt{FMA}, \texttt{CVT}, \texttt{SFU}, \texttt{LS}, \texttt{TEX}, \texttt{BRANCH}). \emph{Example:} {\scriptsize\ttfamily \mnem{FADD.f32} \dst{r0}, \src{r1}, \src{r2}} (\textbf{SPIR-V:} \texttt{N/A (encoding class)}; \textbf{GLSL:} \texttt{// Hardware execution pipe / decode class})
  \item \texttt{fau\_page\_index}: Selects the FAU (Fast Access Uniform) page (\texttt{0..3}) when a source operand references hardware small constants, thread IDs, or lane constants. \emph{Example:} {\scriptsize\ttfamily \mnem{FADD.f32} \dst{r0}, \src{r1}, \src{k1f}} (\textbf{SPIR-V:} \texttt{OpConstant / BuiltIn}; \textbf{GLSL:} \texttt{float y = x + 1.0;})
  \item \hyperref[enum:flow_control_m]{\texttt{flow}}: 4-bit embedded flow-control \& scoreboard wait modifier (\texttt{bits 59:62}: \texttt{none}, \texttt{wait0..wait0126}, \texttt{wait}, \texttt{reconverge}, \texttt{discard}, \texttt{end}). \emph{Example:} {\scriptsize\ttfamily \mnem{FADD.f32}\imod{.wait0.reconverge} \dst{r0}, \src{r1}, \src{r2}} (\textbf{SPIR-V:} \texttt{OpControlBarrier / OpReturn}; \textbf{GLSL:} \texttt{// Scoreboard wait + warp reconvergence})
\end{tightitemize}
\clearpage

\subsection*{FMUL\_CSLICE}
\addcontentsline{toc}{subsection}{FMUL\_CSLICE}
Extended-precision cosine range-reduction slice multiply helper (\texttt{f32}): multiplies \texttt{src0} (supporting \texttt{lane0}, \texttt{.abs}, \texttt{.neg}) against the Payne-Hanek slice constant selected by \texttt{src1} and writes the partial product to \texttt{dest}.
\par\textbf{Notes:}
\begin{tightitemize}
  \item SFU helper opcode.
  \item Not emitted by kraid.
  \item \textbf{Source references:} \texttt{isa-v9-v14.xml}.
\end{tightitemize}
\codebox{\mnem{FMUL\_CSLICE}\imod{\{.flow\}} \dst{dest}\imod{\{.dest\_width\}}, \src{src0}\imod{\{.lane0\}}\imod{\{.abs0\}}\imod{\{.neg0\}}, \src{src1}}
\textbf{Operands:}
\begin{tightitemize}\footnotesize
  \item \texttt{\dst{dest}} (\texttt{dest\_register\_index}, \texttt{dest\_width}): FP32 destination GPR receiving the cosine range-reduction slice product.
  \item \texttt{\src{src0}}: FP32 angle input (supports \texttt{lane0}/\texttt{.abs}/\texttt{.neg}).
  \item \texttt{\src{src1}}: Slice selector / extended-precision constant operand.
\end{tightitemize}
\par\textbf{Closest SPIR-V Op:} \texttt{None (SFU trig range reduction)}
\par\textbf{Pseudo-GLSL (Operational Semantics):}
\codebox{dest = payneHanekCosSlice(mod(lane0(src0)), src1);}
\textbf{Execution Context:}
\begin{tightitemize}\footnotesize
  \item \textbf{Compute stage preloads (\hyperref[sec:csf-spd]{\texttt{SPD\_k}}, \texttt{r16..r23}):} Hardware pre-fills \texttt{r55..r57} with \texttt{WorkGroupID.xyz + JOB\_OFFSET} (from \hyperref[sec:csf-regs]{\texttt{COMPUTESR::JOB\_OFFSET\_X/Y/Z}} \texttt{r34..r36}) and \texttt{r60..r62} with \texttt{LocalInvocationID.xyz} or \texttt{GlobalInvocationID.xyz} (from \texttt{WG\_SIZE} \texttt{r33} \& \texttt{JOB\_SIZE} \texttt{r37..r39}, selected by \texttt{ShaderProgram::preload}) so compute shaders can derive thread/element indices without memory loads.
  \item \textbf{Vertex stage preloads (\hyperref[sec:csf-spd]{\texttt{SPD\_0}/\texttt{SPD\_1}}, \texttt{r16..r19}):} The hardware Index Fetcher reads \hyperref[sec:csf-regs]{\texttt{IDVSSR::INDEX\_BUFFER}} (\texttt{r54:r55}) and pre-fills \texttt{r59, r60} with \texttt{VertexIndex} (\texttt{+ VERTEX\_OFFSET + GLOBAL\_ATTR\_OFF}) and \texttt{r61} with \texttt{InstanceIndex} (\texttt{+ INSTANCE\_OFFSET}) so Position and Varying shaders can pull vertex attributes directly.
  \item \textbf{Fragment stage preloads \& FP modes (\hyperref[sec:csf-spd]{\texttt{SPD\_2}}, \texttt{r20:r21}):} Hardware pre-fills \texttt{r59} with \texttt{floatBitsToUint(FragCoord.z)} and \texttt{r60} with the MSAA sample coverage mask (\texttt{gl\_SampleMaskIn[0]}), while \texttt{ShaderProgram} word~0 latches 32/64 GPR allocation and warp-wide \texttt{flush\_to\_zero\_mode}, \texttt{suppress\_na\_n}, and \texttt{suppress\_inf}.
  \item \textbf{Inline FAU \& Special Operands (\hyperref[sec:csf-tsd-fau]{\texttt{FAU\_k}}, \texttt{r8..r15}):} Any 8-bit source with \texttt{mode=2} reads zero-latency uniform \texttt{u0..u31} from \texttt{FAU\_k} (\texttt{slots64 = bits[63:56]}); \texttt{mode=3} reads \texttt{SmallConstant} (Table~1.2) or \texttt{FauSpecial} (\texttt{pc}, \texttt{lane\_id}, \texttt{lane\_width}, \texttt{tls\_ptr}, \texttt{wls\_ptr}, \texttt{blend\_descriptor\_0..7}).
\end{tightitemize}
\par\textbf{Worked Examples:}
\par\textbf{Example 1 (Cosine Payne-Hanek slice product):}
\codebox{\mnem{FMUL\_CSLICE} \dst{r0}, \src{r1}\imod{.h0}, \src{r2}}
Computes the high-order cosine range-reduction slice product of \texttt{r1.h0} and \texttt{r2} for extended-precision trig.
\par\textbf{Example 2 (Upper halfword slice selection):}
\codebox{\mnem{FMUL\_CSLICE} \dst{r0}, \src{r1}\imod{.h1}, \src{r2}}
Selects \texttt{r1.h1} as the second chunk in a multi-slice \texttt{2/pi} reduction sequence.

\begin{center}\footnotesize
\begin{tabularx}{\linewidth}{@{}lllX@{}}
\toprule
Variant & Arch & Class & Opcode Fields \\
\midrule
\texttt{v2f16} & v9--v14 & \texttt{R2} & \texttt{opcode2=MUL} \\
\bottomrule
\end{tabularx}
\end{center}
\begin{center}\footnotesize
\begin{tabularx}{0.88\linewidth}{@{}l l X@{}}
\toprule
Field Name & Bits & Type / Fixed Value \\
\midrule
\texttt{src0} & \texttt{0:7} & \texttt{SourceEncoding} \\
\texttt{src1} & \texttt{8:15} & \texttt{SourceEncoding} \\
\texttt{opcode2} & \texttt{16:19} & \texttt{exact=MUL} \\
\textit{Reserved} & \texttt{20:27} & \texttt{zero} \\
\hyperref[enum:lane16_m]{\texttt{lane0}} & \texttt{28} & \texttt{lane16\_m (def h0)} \\
\textit{Reserved} & \texttt{29:37} & \texttt{zero} \\
\hyperref[enum:neg_m]{\texttt{neg0}} & \texttt{38} & \texttt{neg\_m (def none)} \\
\hyperref[enum:abs_m]{\texttt{abs0}} & \texttt{39} & \texttt{abs\_m (def none)} \\
\texttt{dest\_register\_index} & \texttt{40:45} & \texttt{uint} \\
\hyperref[enum:dest_width_m]{\texttt{dest\_width}} & \texttt{46:47} & \texttt{dest\_width\_m} \\
\texttt{datatype} & \texttt{48:51} & \texttt{exact=v2f16} \\
\texttt{instruction\_class} & \texttt{52:56} & \texttt{exact=R2} \\
\texttt{fau\_page\_index} & \texttt{57:58} & \texttt{uint} \\
\hyperref[enum:flow_control_m]{\texttt{flow}} & \texttt{59:62} & \texttt{flow\_control\_m} \\
\textit{Reserved} & \texttt{63} & \texttt{zero} \\
\bottomrule
\end{tabularx}
\end{center}
\textbf{Fields Explained:}
\begin{tightitemize}\footnotesize
  \item \texttt{src0}: Primary source operand (\texttt{r0..r63}, \texttt{u0..u63} uniform, \texttt{k*} FAU constant, or \texttt{t0..t2} temporary; \texttt{\textasciicircum{}} marks discard/last-use). \emph{Example:} {\scriptsize\ttfamily \mnem{FADD.f32} \dst{r0}, \src{r1}, \src{r2}} (\textbf{SPIR-V:} \texttt{OpFAdd (Operand 0)}; \textbf{GLSL:} \texttt{float c = a + b;})
  \item \texttt{src1}: Secondary source operand (register, uniform, or FAU constant; in shifts, supplies the shift count byte). \emph{Example:} {\scriptsize\ttfamily \mnem{LSHIFT.i32} \dst{r0}, \src{r1}, \src{r2}\imod{.b0}} (\textbf{SPIR-V:} \texttt{OpShiftLeftLogical (Shift)}; \textbf{GLSL:} \texttt{uint y = a \textless{}\textless{} s;})
  \item \texttt{opcode2}: Secondary sub-opcode further specializing instructions sharing \texttt{instruction\_class} and \texttt{opcode1} (e.g., conversions, SFU tables, moves). \emph{Example:} {\scriptsize\ttfamily \mnem{F32\_TO\_S32}\imod{.round\_zero} \dst{r0}, \src{r1}} (\textbf{SPIR-V:} \texttt{OpConvertFToS}; \textbf{GLSL:} \texttt{int y = int(x);})
  \item \hyperref[enum:lane16_m]{\texttt{lane0}}: Stage A single-lane extractor on \texttt{src0} (\texttt{b0..b3} for 8-bit or \texttt{h0}/\texttt{h1} for 16-bit narrow conversions and \texttt{MKVEC}). \emph{Example:} {\scriptsize\ttfamily \mnem{F16\_TO\_F32} \dst{r0}, \src{r1}\imod{.h1}} (\textbf{SPIR-V:} \texttt{OpCompositeExtract}; \textbf{GLSL:} \texttt{float y = float(x.y);})
  \item \hyperref[enum:neg_m]{\texttt{neg0}}: Stage A floating-point negation on \texttt{src0} (flips IEEE sign bit after swizzle/widen and \texttt{abs0}). \emph{Example:} {\scriptsize\ttfamily \mnem{FADD.f32} \dst{r0}, \src{r1}\imod{.neg}, \src{r2}} (\textbf{SPIR-V:} \texttt{OpFNegate}; \textbf{GLSL:} \texttt{float c = -a + b;})
  \item \hyperref[enum:abs_m]{\texttt{abs0}}: Stage A floating-point absolute value on \texttt{src0} (clears IEEE sign bit after swizzle/widen, before \texttt{neg0}). \emph{Example:} {\scriptsize\ttfamily \mnem{FADD.f32} \dst{r0}, \src{r1}\imod{.abs}, \src{r2}} (\textbf{SPIR-V:} \texttt{GLSL.std.450 FAbs}; \textbf{GLSL:} \texttt{float c = abs(a) + b;})
  \item \texttt{dest\_register\_index}: 6-bit destination GPR index (\texttt{r0..r63}) receiving the Stage C writeback result (even-aligned pair \texttt{[rN:rN+1]} for 64-bit ops). \emph{Example:} {\scriptsize\ttfamily \mnem{FMA.f32} \dst{r0}, \src{r1}, \src{r2}, \src{r3}} (\textbf{SPIR-V:} \texttt{SSA Result ID}; \textbf{GLSL:} \texttt{float r0 = fma(r1, r2, r3);})
  \item \hyperref[enum:dest_width_m]{\texttt{dest\_width}}: Stage C destination lane write-mask (\texttt{none} = full 32-bit, \texttt{h0} = low 16-bit, \texttt{h1} = high 16-bit without bit-shifting). \emph{Example:} {\scriptsize\ttfamily \mnem{FADD.v2f16} \dst{r0}\imod{.h1}, \src{r1}, \src{r2}} (\textbf{SPIR-V:} \texttt{OpVectorInsertDynamic / WriteMask}; \textbf{GLSL:} \texttt{dst.y = a.y + b.y;})
  \item \texttt{datatype}: Selects the lane data type (\texttt{f32}, \texttt{f16}, \texttt{v2f16}, \texttt{s32}, \texttt{u32}, \texttt{v2s16}, \texttt{v2u16}, \texttt{v4s8}, \texttt{v4u8}, \texttt{s64}, \texttt{u64}, etc.). \emph{Example:} {\scriptsize\ttfamily \mnem{IADD.v2s16} \dst{r0}, \src{r1}, \src{r2}} (\textbf{SPIR-V:} \texttt{OpTypeInt / OpTypeFloat}; \textbf{GLSL:} \texttt{i16vec2 c = a + b;})
  \item \texttt{instruction\_class}: 4-bit major instruction class selector (\texttt{bits 48:51}) partitioning the 64-bit Valhall opcode space (e.g., \texttt{FMA}, \texttt{CVT}, \texttt{SFU}, \texttt{LS}, \texttt{TEX}, \texttt{BRANCH}). \emph{Example:} {\scriptsize\ttfamily \mnem{FADD.f32} \dst{r0}, \src{r1}, \src{r2}} (\textbf{SPIR-V:} \texttt{N/A (encoding class)}; \textbf{GLSL:} \texttt{// Hardware execution pipe / decode class})
  \item \texttt{fau\_page\_index}: Selects the FAU (Fast Access Uniform) page (\texttt{0..3}) when a source operand references hardware small constants, thread IDs, or lane constants. \emph{Example:} {\scriptsize\ttfamily \mnem{FADD.f32} \dst{r0}, \src{r1}, \src{k1f}} (\textbf{SPIR-V:} \texttt{OpConstant / BuiltIn}; \textbf{GLSL:} \texttt{float y = x + 1.0;})
  \item \hyperref[enum:flow_control_m]{\texttt{flow}}: 4-bit embedded flow-control \& scoreboard wait modifier (\texttt{bits 59:62}: \texttt{none}, \texttt{wait0..wait0126}, \texttt{wait}, \texttt{reconverge}, \texttt{discard}, \texttt{end}). \emph{Example:} {\scriptsize\ttfamily \mnem{FADD.f32}\imod{.wait0.reconverge} \dst{r0}, \src{r1}, \src{r2}} (\textbf{SPIR-V:} \texttt{OpControlBarrier / OpReturn}; \textbf{GLSL:} \texttt{// Scoreboard wait + warp reconvergence})
\end{tightitemize}
\clearpage

\subsection*{FMUL\_RSCALE}
\addcontentsline{toc}{subsection}{FMUL\_RSCALE}
Computes \texttt{dest = clamp(round\_rm((src0 * src1) * 2\textasciicircum{}clamp(src3, -512, 512)))} in exact FP64 with a single rounding to FP32 (encoded as \texttt{FMA\_RSCALE} with \texttt{src2 = -0.0}). The \texttt{.special} modifier selects \texttt{.n} (\texttt{0 * anything -\textgreater{} $\pm$0}) or \texttt{.left} (\texttt{0 * Inf -\textgreater{} src0} unscaled).
\par\textbf{Note:} \texttt{libGLES\_mali.so} \texttt{ModInfo} allows \texttt{n}/\texttt{left} only (forbids \texttt{n\_add}; \texttt{\_mali\_mul\_left\_po2\_scale\_sf32} \texttt{0x2e29550}).
\codebox{\mnem{FMUL\_RSCALE}\imod{\{.flow\}}\imod{\{.special\}}\imod{\{.round\}}\imod{\{.clamp\}} \dst{dest}\imod{\{.dest\_width\}}, \src{src0}\imod{\{.abs0\}}\imod{\{.neg0\}}, \src{src1}\imod{\{.abs1\}}\imod{\{.neg1\}}, \src{src3}}
\textbf{Operands:}
\begin{tightitemize}\footnotesize
  \item \texttt{\dst{dest}} (\texttt{dest\_register\_index}, \texttt{dest\_width}): FP32 destination GPR receiving \texttt{(src0 * src1) * 2\textasciicircum{}clamp(src3, -512, 512)}.
  \item \texttt{\src{src0}}: First FP32 multiplicand.
  \item \texttt{\src{src1}}: Second FP32 multiplicand (\texttt{src2} is fixed to \texttt{-0.0}).
  \item \texttt{\src{src3}}: Signed integer exponent scale clamped to \texttt{[-512, 512]}.
\end{tightitemize}
\par\textbf{Closest SPIR-V Op:} \texttt{None (scaled FP mul for Sin/Pow)}
\par\textbf{Pseudo-GLSL (Operational Semantics):}
\codebox{int e = clamp(int(src3), -512, 512);\\if (special == left \&\& isZeroTimesInf(src0, src1)) dest = src0;\\else dest = clamp\_mode(round\_rm(float(double(src0) * double(src1) * exp2(double(e)))));}
\textbf{Execution Context:}
\begin{tightitemize}\footnotesize
  \item \textbf{Compute stage preloads (\hyperref[sec:csf-spd]{\texttt{SPD\_k}}, \texttt{r16..r23}):} Hardware pre-fills \texttt{r55..r57} with \texttt{WorkGroupID.xyz + JOB\_OFFSET} (from \hyperref[sec:csf-regs]{\texttt{COMPUTESR::JOB\_OFFSET\_X/Y/Z}} \texttt{r34..r36}) and \texttt{r60..r62} with \texttt{LocalInvocationID.xyz} or \texttt{GlobalInvocationID.xyz} (from \texttt{WG\_SIZE} \texttt{r33} \& \texttt{JOB\_SIZE} \texttt{r37..r39}, selected by \texttt{ShaderProgram::preload}) so compute shaders can derive thread/element indices without memory loads.
  \item \textbf{Vertex stage preloads (\hyperref[sec:csf-spd]{\texttt{SPD\_0}/\texttt{SPD\_1}}, \texttt{r16..r19}):} The hardware Index Fetcher reads \hyperref[sec:csf-regs]{\texttt{IDVSSR::INDEX\_BUFFER}} (\texttt{r54:r55}) and pre-fills \texttt{r59, r60} with \texttt{VertexIndex} (\texttt{+ VERTEX\_OFFSET + GLOBAL\_ATTR\_OFF}) and \texttt{r61} with \texttt{InstanceIndex} (\texttt{+ INSTANCE\_OFFSET}) so Position and Varying shaders can pull vertex attributes directly.
  \item \textbf{Fragment stage preloads \& FP modes (\hyperref[sec:csf-spd]{\texttt{SPD\_2}}, \texttt{r20:r21}):} Hardware pre-fills \texttt{r59} with \texttt{floatBitsToUint(FragCoord.z)} and \texttt{r60} with the MSAA sample coverage mask (\texttt{gl\_SampleMaskIn[0]}), while \texttt{ShaderProgram} word~0 latches 32/64 GPR allocation and warp-wide \texttt{flush\_to\_zero\_mode}, \texttt{suppress\_na\_n}, and \texttt{suppress\_inf}.
  \item \textbf{Inline FAU \& Special Operands (\hyperref[sec:csf-tsd-fau]{\texttt{FAU\_k}}, \texttt{r8..r15}):} Any 8-bit source with \texttt{mode=2} reads zero-latency uniform \texttt{u0..u31} from \texttt{FAU\_k} (\texttt{slots64 = bits[63:56]}); \texttt{mode=3} reads \texttt{SmallConstant} (Table~1.2) or \texttt{FauSpecial} (\texttt{pc}, \texttt{lane\_id}, \texttt{lane\_width}, \texttt{tls\_ptr}, \texttt{wls\_ptr}, \texttt{blend\_descriptor\_0..7}).
\end{tightitemize}
\par\textbf{Worked Examples:}
\par\textbf{Example 1 (Exact FP64-widened scaled multiply):}
\codebox{\mnem{FMUL\_RSCALE.f32} \dst{r0}, \src{r1}, \src{r2}, \src{k1n}}
With \texttt{r1 = 2\textasciicircum{}80, r2 = 2\textasciicircum{}60, scale = -140}: produces exact \texttt{1.0f} (\texttt{0x3F800000}) despite \texttt{r1 * r2 = 2\textasciicircum{}140 \textgreater{} FLT\_MAX}.
\par\textbf{Example 2 (Left-operand preservation on `0 * Inf` (.left)):}
\codebox{\mnem{FMUL\_RSCALE.f32}\imod{.left} \dst{r0}, \src{r1}, \src{r2}, \src{k1n}}
With \texttt{r1 = 0.0f, r2 = +Inf}: \texttt{.left} returns \texttt{r1} (\texttt{0.0f}) instead of \texttt{NaN}.
\par\textbf{Mined SPIR-V $\to$ Valhall ISA Example (Computing delta\textasciicircum{}2 / 2 during trig polynomial evaluation --- Vertex shader, Mali-G710 (v10)):}
\codebox{\textbf{; SPIR-V (input)}\\\%c = OpExtInst \%float \%glsl Cos \%angle\\[3pt]\textbf{; Valhall ISA (malioc output)}\\FMUL\_RSCALE.f32 r7, r2, r2, k1n   ; (r2 * r2) * 2\textasciicircum{}(-1)\\FSIN\_TABLE.u6 r63, r5\\FCOS\_TABLE.u6 r5, r5\textasciicircum{}}
Passing \texttt{k1n} (\texttt{-1}) as the exponent scale \texttt{src3} divides \texttt{r2 * r2} by 2 in exact FP64 without needing a separate \texttt{0.5f} constant multiply.
\begin{center}\footnotesize
\begin{tabularx}{\linewidth}{@{}lllX@{}}
\toprule
Variant & Arch & Class & Opcode Fields \\
\midrule
\texttt{(none)} & v9--v14 & \texttt{R4\_FMAR} & \texttt{opcode1=FMA\_RSCALE\_F32} \\
\bottomrule
\end{tabularx}
\end{center}
\begin{center}\footnotesize
\begin{tabularx}{0.88\linewidth}{@{}l l X@{}}
\toprule
Field Name & Bits & Type / Fixed Value \\
\midrule
\texttt{src0} & \texttt{0:7} & \texttt{SourceEncoding} \\
\texttt{src1} & \texttt{8:15} & \texttt{SourceEncoding} \\
\texttt{src2} & \texttt{16:23} & \texttt{exact=192} \\
\texttt{src3} & \texttt{24:31} & \texttt{SourceEncoding} \\
\hyperref[enum:clamp_m]{\texttt{clamp}} & \texttt{32:33} & \texttt{clamp\_m (def none)} \\
\hyperref[enum:neg_m]{\texttt{neg2}} & \texttt{34} & \texttt{exact=neg} \\
\hyperref[enum:abs_m]{\texttt{abs2}} & \texttt{35} & \texttt{exact=none} \\
\hyperref[enum:neg_m]{\texttt{neg1}} & \texttt{36} & \texttt{neg\_m (def none)} \\
\hyperref[enum:abs_m]{\texttt{abs1}} & \texttt{37} & \texttt{abs\_m (def none)} \\
\hyperref[enum:neg_m]{\texttt{neg0}} & \texttt{38} & \texttt{neg\_m (def none)} \\
\hyperref[enum:abs_m]{\texttt{abs0}} & \texttt{39} & \texttt{abs\_m (def none)} \\
\texttt{dest\_register\_index} & \texttt{40:45} & \texttt{uint} \\
\hyperref[enum:dest_width_m]{\texttt{dest\_width}} & \texttt{46:47} & \texttt{dest\_width\_m} \\
\hyperref[enum:fma_rscale_special_32_m]{\texttt{special}} & \texttt{48:49} & \texttt{fma\_rscale\_special\_32\_m} \\
\hyperref[enum:fma_rscale_round_m]{\texttt{round}} & \texttt{50} & \texttt{fma\_rscale\_round\_m} \\
\texttt{opcode1} & \texttt{51} & \texttt{exact=FMA\_RSCALE\_F32} \\
\texttt{instruction\_class} & \texttt{52:56} & \texttt{exact=R4\_FMAR} \\
\texttt{fau\_page\_index} & \texttt{57:58} & \texttt{uint} \\
\hyperref[enum:flow_control_m]{\texttt{flow}} & \texttt{59:62} & \texttt{flow\_control\_m} \\
\textit{Reserved} & \texttt{63} & \texttt{zero} \\
\bottomrule
\end{tabularx}
\end{center}
\textbf{Fields Explained:}
\begin{tightitemize}\footnotesize
  \item \texttt{src0}: Primary source operand (\texttt{r0..r63}, \texttt{u0..u63} uniform, \texttt{k*} FAU constant, or \texttt{t0..t2} temporary; \texttt{\textasciicircum{}} marks discard/last-use). \emph{Example:} {\scriptsize\ttfamily \mnem{FADD.f32} \dst{r0}, \src{r1}, \src{r2}} (\textbf{SPIR-V:} \texttt{OpFAdd (Operand 0)}; \textbf{GLSL:} \texttt{float c = a + b;})
  \item \texttt{src1}: Secondary source operand (register, uniform, or FAU constant; in shifts, supplies the shift count byte). \emph{Example:} {\scriptsize\ttfamily \mnem{LSHIFT.i32} \dst{r0}, \src{r1}, \src{r2}\imod{.b0}} (\textbf{SPIR-V:} \texttt{OpShiftLeftLogical (Shift)}; \textbf{GLSL:} \texttt{uint y = a \textless{}\textless{} s;})
  \item \texttt{src2}: Third source operand (FMA addend, \texttt{MUX}/\texttt{CSEL} selector, \texttt{IADDC} carry-in, or bitwise mask in fused shift-logic). \emph{Example:} {\scriptsize\ttfamily \mnem{FMA.f32} \dst{r0}, \src{r1}, \src{r2}, \src{r3}} (\textbf{SPIR-V:} \texttt{GLSL.std.450 Fma (Addend)}; \textbf{GLSL:} \texttt{float d = fma(a, b, c);})
  \item \texttt{src3}: Fourth source operand (false-branch value in \texttt{CSEL}, exponent scale in \texttt{FMA\_RSCALE}, or fourth matrix/select input). \emph{Example:} {\scriptsize\ttfamily \mnem{CSEL.f32}\imod{.lt} \dst{r0}, \src{r1}, \src{r2}, \src{r3}, \src{r4}} (\textbf{SPIR-V:} \texttt{OpSelect (Object 2)}; \textbf{GLSL:} \texttt{float r = (a \textless{} b) ? x : y;})
  \item \hyperref[enum:clamp_m]{\texttt{clamp}}: Stage B post-op floating-point clamp (\texttt{none}, \texttt{clamp\_0\_inf}, \texttt{clamp\_m1\_1}, \texttt{clamp\_0\_1}) with \texttt{-0.0 \textless{} +0.0} and NaN suppression (\texttt{NaN -\textgreater{} lo}). \emph{Example:} {\scriptsize\ttfamily \mnem{FMA.f32}\imod{.clamp\_0\_1} \dst{r0}, \src{r1}, \src{r2}, \src{r3}} (\textbf{SPIR-V:} \texttt{GLSL.std.450 FClamp}; \textbf{GLSL:} \texttt{float y = clamp(fma(a, b, c), 0.0, 1.0);})
  \item \hyperref[enum:neg_m]{\texttt{neg2}}: Stage A floating-point negation on \texttt{src2} (flips IEEE sign bit of FMA addend after swizzle/widen and \texttt{abs2}). \emph{Example:} {\scriptsize\ttfamily \mnem{FMA.f32} \dst{r0}, \src{r1}, \src{r2}, \src{r3}\imod{.neg}} (\textbf{SPIR-V:} \texttt{OpFNegate}; \textbf{GLSL:} \texttt{float d = fma(a, b, -c);})
  \item \hyperref[enum:abs_m]{\texttt{abs2}}: Stage A floating-point absolute value on \texttt{src2} (clears IEEE sign bit of FMA addend after swizzle/widen). \emph{Example:} {\scriptsize\ttfamily \mnem{FMA.f32} \dst{r0}, \src{r1}, \src{r2}, \src{r3}\imod{.abs}} (\textbf{SPIR-V:} \texttt{GLSL.std.450 FAbs}; \textbf{GLSL:} \texttt{float d = fma(a, b, abs(c));})
  \item \hyperref[enum:neg_m]{\texttt{neg1}}: Stage A floating-point negation on \texttt{src1} (flips IEEE sign bit after swizzle/widen and \texttt{abs1}). \emph{Example:} {\scriptsize\ttfamily \mnem{FMA.f32} \dst{r0}, \src{r1}, \src{r2}\imod{.neg}, \src{r3}} (\textbf{SPIR-V:} \texttt{OpFNegate}; \textbf{GLSL:} \texttt{float d = fma(a, -b, c);})
  \item \hyperref[enum:abs_m]{\texttt{abs1}}: Stage A floating-point absolute value on \texttt{src1} (clears IEEE sign bit after swizzle/widen, before \texttt{neg1}). \emph{Example:} {\scriptsize\ttfamily \mnem{FMA.f32} \dst{r0}, \src{r1}, \src{r2}\imod{.abs}, \src{r3}} (\textbf{SPIR-V:} \texttt{GLSL.std.450 FAbs}; \textbf{GLSL:} \texttt{float d = fma(a, abs(b), c);})
  \item \hyperref[enum:neg_m]{\texttt{neg0}}: Stage A floating-point negation on \texttt{src0} (flips IEEE sign bit after swizzle/widen and \texttt{abs0}). \emph{Example:} {\scriptsize\ttfamily \mnem{FADD.f32} \dst{r0}, \src{r1}\imod{.neg}, \src{r2}} (\textbf{SPIR-V:} \texttt{OpFNegate}; \textbf{GLSL:} \texttt{float c = -a + b;})
  \item \hyperref[enum:abs_m]{\texttt{abs0}}: Stage A floating-point absolute value on \texttt{src0} (clears IEEE sign bit after swizzle/widen, before \texttt{neg0}). \emph{Example:} {\scriptsize\ttfamily \mnem{FADD.f32} \dst{r0}, \src{r1}\imod{.abs}, \src{r2}} (\textbf{SPIR-V:} \texttt{GLSL.std.450 FAbs}; \textbf{GLSL:} \texttt{float c = abs(a) + b;})
  \item \texttt{dest\_register\_index}: 6-bit destination GPR index (\texttt{r0..r63}) receiving the Stage C writeback result (even-aligned pair \texttt{[rN:rN+1]} for 64-bit ops). \emph{Example:} {\scriptsize\ttfamily \mnem{FMA.f32} \dst{r0}, \src{r1}, \src{r2}, \src{r3}} (\textbf{SPIR-V:} \texttt{SSA Result ID}; \textbf{GLSL:} \texttt{float r0 = fma(r1, r2, r3);})
  \item \hyperref[enum:dest_width_m]{\texttt{dest\_width}}: Stage C destination lane write-mask (\texttt{none} = full 32-bit, \texttt{h0} = low 16-bit, \texttt{h1} = high 16-bit without bit-shifting). \emph{Example:} {\scriptsize\ttfamily \mnem{FADD.v2f16} \dst{r0}\imod{.h1}, \src{r1}, \src{r2}} (\textbf{SPIR-V:} \texttt{OpVectorInsertDynamic / WriteMask}; \textbf{GLSL:} \texttt{dst.y = a.y + b.y;})
  \item \hyperref[enum:fma_rscale_special_32_m]{\texttt{special}}: Selects \texttt{FMA\_RSCALE}/\texttt{FREXP} special-case override (\texttt{n}, \texttt{left}, \texttt{n\_add}, \texttt{sqrt}, \texttt{log}). \emph{Example:} {\scriptsize\ttfamily \mnem{FREXPM.f32}\imod{.log} \dst{r0}, \src{r1}} (\textbf{SPIR-V:} \texttt{GLSL.std.450 Log2 (range reduction)}; \textbf{GLSL:} \texttt{float y = log2(x);})
  \item \hyperref[enum:fma_rscale_round_m]{\texttt{round}}: Rounding mode (\texttt{none}/RNE, \texttt{round\_up}, \texttt{round\_down}, \texttt{round\_zero}, \texttt{round\_na}) applied in FP ALU, conversions, \texttt{HADD}, \texttt{LDEXP}, and Stage C narrowing. \emph{Example:} {\scriptsize\ttfamily \mnem{FROUND.f32}\imod{.round\_down} \dst{r0}, \src{r1}} (\textbf{SPIR-V:} \texttt{GLSL.std.450 Floor / Ceil / Trunc}; \textbf{GLSL:} \texttt{float y = floor(x);})
  \item \texttt{opcode1}: Primary sub-opcode within the instruction class uniquely identifying the operation family. \emph{Example:} {\scriptsize\ttfamily \mnem{IADD.s32} \dst{r0}, \src{r1}, \src{r2}} (\textbf{SPIR-V:} \texttt{OpIAdd (operation selector)}; \textbf{GLSL:} \texttt{// Primary opcode bits})
  \item \texttt{instruction\_class}: 4-bit major instruction class selector (\texttt{bits 48:51}) partitioning the 64-bit Valhall opcode space (e.g., \texttt{FMA}, \texttt{CVT}, \texttt{SFU}, \texttt{LS}, \texttt{TEX}, \texttt{BRANCH}). \emph{Example:} {\scriptsize\ttfamily \mnem{FADD.f32} \dst{r0}, \src{r1}, \src{r2}} (\textbf{SPIR-V:} \texttt{N/A (encoding class)}; \textbf{GLSL:} \texttt{// Hardware execution pipe / decode class})
  \item \texttt{fau\_page\_index}: Selects the FAU (Fast Access Uniform) page (\texttt{0..3}) when a source operand references hardware small constants, thread IDs, or lane constants. \emph{Example:} {\scriptsize\ttfamily \mnem{FADD.f32} \dst{r0}, \src{r1}, \src{k1f}} (\textbf{SPIR-V:} \texttt{OpConstant / BuiltIn}; \textbf{GLSL:} \texttt{float y = x + 1.0;})
  \item \hyperref[enum:flow_control_m]{\texttt{flow}}: 4-bit embedded flow-control \& scoreboard wait modifier (\texttt{bits 59:62}: \texttt{none}, \texttt{wait0..wait0126}, \texttt{wait}, \texttt{reconverge}, \texttt{discard}, \texttt{end}). \emph{Example:} {\scriptsize\ttfamily \mnem{FADD.f32}\imod{.wait0.reconverge} \dst{r0}, \src{r1}, \src{r2}} (\textbf{SPIR-V:} \texttt{OpControlBarrier / OpReturn}; \textbf{GLSL:} \texttt{// Scoreboard wait + warp reconvergence})
\end{tightitemize}
\clearpage

\subsection*{FMUL\_SLICE}
\addcontentsline{toc}{subsection}{FMUL\_SLICE}
Trigonometric range-reduction multiply helper (\texttt{f32}): computes sliced product for Payne-Hanek / Cody-Waite \texttt{sin}/\texttt{cos} reduction.
\par\textbf{Notes:}
\begin{tightitemize}
  \item SFU helper opcode.
  \item Not emitted by kraid.
  \item \textbf{Source references:} \texttt{isa-v9-v14.xml}, \texttt{bifrost/bifrost/ISA.xml:1290}.
\end{tightitemize}
\codebox{\mnem{FMUL\_SLICE}\imod{\{.flow\}} \dst{dest}\imod{\{.dest\_width\}}, \src{src0}, \src{src1}}
\textbf{Operands:}
\begin{tightitemize}\footnotesize
  \item \texttt{\dst{dest}} (\texttt{dest\_register\_index}, \texttt{dest\_width}): FP32 destination GPR receiving the sliced \texttt{2/pi} range-reduction product.
  \item \texttt{\src{src0}}: FP32 input angle.
  \item \texttt{\src{src1}}: Slice index / reduction operand.
\end{tightitemize}
\par\textbf{Closest SPIR-V Op:} \texttt{None (SFU trig range reduction)}
\par\textbf{Pseudo-GLSL (Operational Semantics):}
\codebox{dest = payneHanekTwoOverPiSlice(src0, src1);}
\textbf{Execution Context:}
\begin{tightitemize}\footnotesize
  \item \textbf{Compute stage preloads (\hyperref[sec:csf-spd]{\texttt{SPD\_k}}, \texttt{r16..r23}):} Hardware pre-fills \texttt{r55..r57} with \texttt{WorkGroupID.xyz + JOB\_OFFSET} (from \hyperref[sec:csf-regs]{\texttt{COMPUTESR::JOB\_OFFSET\_X/Y/Z}} \texttt{r34..r36}) and \texttt{r60..r62} with \texttt{LocalInvocationID.xyz} or \texttt{GlobalInvocationID.xyz} (from \texttt{WG\_SIZE} \texttt{r33} \& \texttt{JOB\_SIZE} \texttt{r37..r39}, selected by \texttt{ShaderProgram::preload}) so compute shaders can derive thread/element indices without memory loads.
  \item \textbf{Vertex stage preloads (\hyperref[sec:csf-spd]{\texttt{SPD\_0}/\texttt{SPD\_1}}, \texttt{r16..r19}):} The hardware Index Fetcher reads \hyperref[sec:csf-regs]{\texttt{IDVSSR::INDEX\_BUFFER}} (\texttt{r54:r55}) and pre-fills \texttt{r59, r60} with \texttt{VertexIndex} (\texttt{+ VERTEX\_OFFSET + GLOBAL\_ATTR\_OFF}) and \texttt{r61} with \texttt{InstanceIndex} (\texttt{+ INSTANCE\_OFFSET}) so Position and Varying shaders can pull vertex attributes directly.
  \item \textbf{Fragment stage preloads \& FP modes (\hyperref[sec:csf-spd]{\texttt{SPD\_2}}, \texttt{r20:r21}):} Hardware pre-fills \texttt{r59} with \texttt{floatBitsToUint(FragCoord.z)} and \texttt{r60} with the MSAA sample coverage mask (\texttt{gl\_SampleMaskIn[0]}), while \texttt{ShaderProgram} word~0 latches 32/64 GPR allocation and warp-wide \texttt{flush\_to\_zero\_mode}, \texttt{suppress\_na\_n}, and \texttt{suppress\_inf}.
  \item \textbf{Inline FAU \& Special Operands (\hyperref[sec:csf-tsd-fau]{\texttt{FAU\_k}}, \texttt{r8..r15}):} Any 8-bit source with \texttt{mode=2} reads zero-latency uniform \texttt{u0..u31} from \texttt{FAU\_k} (\texttt{slots64 = bits[63:56]}); \texttt{mode=3} reads \texttt{SmallConstant} (Table~1.2) or \texttt{FauSpecial} (\texttt{pc}, \texttt{lane\_id}, \texttt{lane\_width}, \texttt{tls\_ptr}, \texttt{wls\_ptr}, \texttt{blend\_descriptor\_0..7}).
\end{tightitemize}
\par\textbf{Worked Examples:}
\par\textbf{Example 1 (Payne-Hanek `2/pi` slice multiply):}
\codebox{\mnem{FMUL\_SLICE.f32} \dst{r0}, \src{r1}, \src{r2}}
Extracts a partial product slice of \texttt{x * (2/pi)} for large-angle \texttt{sin(x)} / \texttt{cos(x)} range reduction.
\par\textbf{Example 2 (Negated slice accumulation):}
\codebox{\mnem{FMUL\_SLICE.f32} \dst{r0}, \src{r1}, \src{r2}\imod{.neg}}
Negates \texttt{r2} in Stage A to compute the residual subtraction term during multi-word trig reduction.

\begin{center}\footnotesize
\begin{tabularx}{\linewidth}{@{}lllX@{}}
\toprule
Variant & Arch & Class & Opcode Fields \\
\midrule
\texttt{f32} & v9--v14 & \texttt{R2} & \texttt{opcode2=MUL} \\
\bottomrule
\end{tabularx}
\end{center}
\begin{center}\footnotesize
\begin{tabularx}{0.88\linewidth}{@{}l l X@{}}
\toprule
Field Name & Bits & Type / Fixed Value \\
\midrule
\texttt{src0} & \texttt{0:7} & \texttt{SourceEncoding} \\
\texttt{src1} & \texttt{8:15} & \texttt{SourceEncoding} \\
\texttt{opcode2} & \texttt{16:19} & \texttt{exact=MUL} \\
\textit{Reserved} & \texttt{20:39} & \texttt{zero} \\
\texttt{dest\_register\_index} & \texttt{40:45} & \texttt{uint} \\
\hyperref[enum:dest_width_m]{\texttt{dest\_width}} & \texttt{46:47} & \texttt{dest\_width\_m} \\
\texttt{datatype} & \texttt{48:51} & \texttt{exact=f32} \\
\texttt{instruction\_class} & \texttt{52:56} & \texttt{exact=R2} \\
\texttt{fau\_page\_index} & \texttt{57:58} & \texttt{uint} \\
\hyperref[enum:flow_control_m]{\texttt{flow}} & \texttt{59:62} & \texttt{flow\_control\_m} \\
\textit{Reserved} & \texttt{63} & \texttt{zero} \\
\bottomrule
\end{tabularx}
\end{center}
\textbf{Fields Explained:}
\begin{tightitemize}\footnotesize
  \item \texttt{src0}: Primary source operand (\texttt{r0..r63}, \texttt{u0..u63} uniform, \texttt{k*} FAU constant, or \texttt{t0..t2} temporary; \texttt{\textasciicircum{}} marks discard/last-use). \emph{Example:} {\scriptsize\ttfamily \mnem{FADD.f32} \dst{r0}, \src{r1}, \src{r2}} (\textbf{SPIR-V:} \texttt{OpFAdd (Operand 0)}; \textbf{GLSL:} \texttt{float c = a + b;})
  \item \texttt{src1}: Secondary source operand (register, uniform, or FAU constant; in shifts, supplies the shift count byte). \emph{Example:} {\scriptsize\ttfamily \mnem{LSHIFT.i32} \dst{r0}, \src{r1}, \src{r2}\imod{.b0}} (\textbf{SPIR-V:} \texttt{OpShiftLeftLogical (Shift)}; \textbf{GLSL:} \texttt{uint y = a \textless{}\textless{} s;})
  \item \texttt{opcode2}: Secondary sub-opcode further specializing instructions sharing \texttt{instruction\_class} and \texttt{opcode1} (e.g., conversions, SFU tables, moves). \emph{Example:} {\scriptsize\ttfamily \mnem{F32\_TO\_S32}\imod{.round\_zero} \dst{r0}, \src{r1}} (\textbf{SPIR-V:} \texttt{OpConvertFToS}; \textbf{GLSL:} \texttt{int y = int(x);})
  \item \texttt{dest\_register\_index}: 6-bit destination GPR index (\texttt{r0..r63}) receiving the Stage C writeback result (even-aligned pair \texttt{[rN:rN+1]} for 64-bit ops). \emph{Example:} {\scriptsize\ttfamily \mnem{FMA.f32} \dst{r0}, \src{r1}, \src{r2}, \src{r3}} (\textbf{SPIR-V:} \texttt{SSA Result ID}; \textbf{GLSL:} \texttt{float r0 = fma(r1, r2, r3);})
  \item \hyperref[enum:dest_width_m]{\texttt{dest\_width}}: Stage C destination lane write-mask (\texttt{none} = full 32-bit, \texttt{h0} = low 16-bit, \texttt{h1} = high 16-bit without bit-shifting). \emph{Example:} {\scriptsize\ttfamily \mnem{FADD.v2f16} \dst{r0}\imod{.h1}, \src{r1}, \src{r2}} (\textbf{SPIR-V:} \texttt{OpVectorInsertDynamic / WriteMask}; \textbf{GLSL:} \texttt{dst.y = a.y + b.y;})
  \item \texttt{datatype}: Selects the lane data type (\texttt{f32}, \texttt{f16}, \texttt{v2f16}, \texttt{s32}, \texttt{u32}, \texttt{v2s16}, \texttt{v2u16}, \texttt{v4s8}, \texttt{v4u8}, \texttt{s64}, \texttt{u64}, etc.). \emph{Example:} {\scriptsize\ttfamily \mnem{IADD.v2s16} \dst{r0}, \src{r1}, \src{r2}} (\textbf{SPIR-V:} \texttt{OpTypeInt / OpTypeFloat}; \textbf{GLSL:} \texttt{i16vec2 c = a + b;})
  \item \texttt{instruction\_class}: 4-bit major instruction class selector (\texttt{bits 48:51}) partitioning the 64-bit Valhall opcode space (e.g., \texttt{FMA}, \texttt{CVT}, \texttt{SFU}, \texttt{LS}, \texttt{TEX}, \texttt{BRANCH}). \emph{Example:} {\scriptsize\ttfamily \mnem{FADD.f32} \dst{r0}, \src{r1}, \src{r2}} (\textbf{SPIR-V:} \texttt{N/A (encoding class)}; \textbf{GLSL:} \texttt{// Hardware execution pipe / decode class})
  \item \texttt{fau\_page\_index}: Selects the FAU (Fast Access Uniform) page (\texttt{0..3}) when a source operand references hardware small constants, thread IDs, or lane constants. \emph{Example:} {\scriptsize\ttfamily \mnem{FADD.f32} \dst{r0}, \src{r1}, \src{k1f}} (\textbf{SPIR-V:} \texttt{OpConstant / BuiltIn}; \textbf{GLSL:} \texttt{float y = x + 1.0;})
  \item \hyperref[enum:flow_control_m]{\texttt{flow}}: 4-bit embedded flow-control \& scoreboard wait modifier (\texttt{bits 59:62}: \texttt{none}, \texttt{wait0..wait0126}, \texttt{wait}, \texttt{reconverge}, \texttt{discard}, \texttt{end}). \emph{Example:} {\scriptsize\ttfamily \mnem{FADD.f32}\imod{.wait0.reconverge} \dst{r0}, \src{r1}, \src{r2}} (\textbf{SPIR-V:} \texttt{OpControlBarrier / OpReturn}; \textbf{GLSL:} \texttt{// Scoreboard wait + warp reconvergence})
\end{tightitemize}
\clearpage

\subsection*{FPOW\_SC\_APPLY}
\addcontentsline{toc}{subsection}{FPOW\_SC\_APPLY}
Power special-case applicator (\texttt{f32}): inspects the special-case tag word in \texttt{src1} (produced by \texttt{FPOW\_SC\_DET}); if a special-case override is flagged, writes the decoded IEEE-754 override value to \texttt{dest}, otherwise passes \texttt{src0} through unchanged.
\par\textbf{Notes:}
\begin{tightitemize}
  \item Paired with \texttt{FPOW\_SC\_DET}.
  \item \textbf{Source references:} \texttt{bifrost/bifrost/ISA.xml:4634}, \texttt{isa-v9-v14.xml}.
\end{tightitemize}
\codebox{\mnem{FPOW\_SC\_APPLY}\imod{\{.flow\}} \dst{dest}\imod{\{.dest\_width\}}, \src{src0}, \src{src1}}
\textbf{Operands:}
\begin{tightitemize}\footnotesize
  \item \texttt{\dst{dest}} (\texttt{dest\_register\_index}, \texttt{dest\_width}): FP32 destination GPR receiving either the special-case override or \texttt{src0}.
  \item \texttt{\src{src0}}: Computed general-case power result \texttt{exp2(y * log2(\textbar{}x\textbar{}))}.
  \item \texttt{\src{src1}}: Special-case tag word produced by \texttt{FPOW\_SC\_DET}.
\end{tightitemize}
\par\textbf{Closest SPIR-V Op:} \texttt{None (lowered from GLSL.std.450 Pow)}
\par\textbf{Pseudo-GLSL (Operational Semantics):}
\codebox{dest = hasSpecialOverride(src1) ? decodeSpecialOverride(src1) : src0;}
\textbf{Execution Context:}
\begin{tightitemize}\footnotesize
  \item \textbf{Compute stage preloads (\hyperref[sec:csf-spd]{\texttt{SPD\_k}}, \texttt{r16..r23}):} Hardware pre-fills \texttt{r55..r57} with \texttt{WorkGroupID.xyz + JOB\_OFFSET} (from \hyperref[sec:csf-regs]{\texttt{COMPUTESR::JOB\_OFFSET\_X/Y/Z}} \texttt{r34..r36}) and \texttt{r60..r62} with \texttt{LocalInvocationID.xyz} or \texttt{GlobalInvocationID.xyz} (from \texttt{WG\_SIZE} \texttt{r33} \& \texttt{JOB\_SIZE} \texttt{r37..r39}, selected by \texttt{ShaderProgram::preload}) so compute shaders can derive thread/element indices without memory loads.
  \item \textbf{Vertex stage preloads (\hyperref[sec:csf-spd]{\texttt{SPD\_0}/\texttt{SPD\_1}}, \texttt{r16..r19}):} The hardware Index Fetcher reads \hyperref[sec:csf-regs]{\texttt{IDVSSR::INDEX\_BUFFER}} (\texttt{r54:r55}) and pre-fills \texttt{r59, r60} with \texttt{VertexIndex} (\texttt{+ VERTEX\_OFFSET + GLOBAL\_ATTR\_OFF}) and \texttt{r61} with \texttt{InstanceIndex} (\texttt{+ INSTANCE\_OFFSET}) so Position and Varying shaders can pull vertex attributes directly.
  \item \textbf{Fragment stage preloads \& FP modes (\hyperref[sec:csf-spd]{\texttt{SPD\_2}}, \texttt{r20:r21}):} Hardware pre-fills \texttt{r59} with \texttt{floatBitsToUint(FragCoord.z)} and \texttt{r60} with the MSAA sample coverage mask (\texttt{gl\_SampleMaskIn[0]}), while \texttt{ShaderProgram} word~0 latches 32/64 GPR allocation and warp-wide \texttt{flush\_to\_zero\_mode}, \texttt{suppress\_na\_n}, and \texttt{suppress\_inf}.
  \item \textbf{Inline FAU \& Special Operands (\hyperref[sec:csf-tsd-fau]{\texttt{FAU\_k}}, \texttt{r8..r15}):} Any 8-bit source with \texttt{mode=2} reads zero-latency uniform \texttt{u0..u31} from \texttt{FAU\_k} (\texttt{slots64 = bits[63:56]}); \texttt{mode=3} reads \texttt{SmallConstant} (Table~1.2) or \texttt{FauSpecial} (\texttt{pc}, \texttt{lane\_id}, \texttt{lane\_width}, \texttt{tls\_ptr}, \texttt{wls\_ptr}, \texttt{blend\_descriptor\_0..7}).
\end{tightitemize}
\par\textbf{Worked Examples:}
\par\textbf{Example 1 (Normal path passthrough):}
\codebox{\mnem{FPOW\_SC\_APPLY} \dst{r0}, \src{r1}, \src{r2}}
When \texttt{r2} from \texttt{FPOW\_SC\_DET} indicates no special case: passes computed \texttt{exp2(y * log2(x))} in \texttt{r1} straight through to \texttt{r0}.
\par\textbf{Example 2 (Special-case override (`pow(0.0, 0.0) = 1.0`)):}
\codebox{\mnem{FPOW\_SC\_APPLY} \dst{r0}, \src{r1}, \src{r2}}
When \texttt{FPOW\_SC\_DET} flags \texttt{x == 0, y == 0}: overrides \texttt{r1} (\texttt{NaN}) with \texttt{r0 = 1.0f} (\texttt{0x3F800000}).

\begin{center}\footnotesize
\begin{tabularx}{\linewidth}{@{}lllX@{}}
\toprule
Variant & Arch & Class & Opcode Fields \\
\midrule
\texttt{f32} & v9--v14 & \texttt{R2} & \texttt{opcode2=FPOW, opcode3=FPOW\_SC\_APPLY} \\
\bottomrule
\end{tabularx}
\end{center}
\begin{center}\footnotesize
\begin{tabularx}{0.88\linewidth}{@{}l l X@{}}
\toprule
Field Name & Bits & Type / Fixed Value \\
\midrule
\texttt{src0} & \texttt{0:7} & \texttt{SourceEncoding} \\
\texttt{src1} & \texttt{8:15} & \texttt{SourceEncoding} \\
\texttt{opcode2} & \texttt{16:19} & \texttt{exact=FPOW} \\
\textit{Reserved} & \texttt{20:23} & \texttt{zero} \\
\texttt{opcode3} & \texttt{24:25} & \texttt{exact=FPOW\_SC\_APPLY} \\
\textit{Reserved} & \texttt{26:39} & \texttt{zero} \\
\texttt{dest\_register\_index} & \texttt{40:45} & \texttt{uint} \\
\hyperref[enum:dest_width_m]{\texttt{dest\_width}} & \texttt{46:47} & \texttt{dest\_width\_m} \\
\texttt{datatype} & \texttt{48:51} & \texttt{exact=f32} \\
\texttt{instruction\_class} & \texttt{52:56} & \texttt{exact=R2} \\
\texttt{fau\_page\_index} & \texttt{57:58} & \texttt{uint} \\
\hyperref[enum:flow_control_m]{\texttt{flow}} & \texttt{59:62} & \texttt{flow\_control\_m} \\
\textit{Reserved} & \texttt{63} & \texttt{zero} \\
\bottomrule
\end{tabularx}
\end{center}
\textbf{Fields Explained:}
\begin{tightitemize}\footnotesize
  \item \texttt{src0}: Primary source operand (\texttt{r0..r63}, \texttt{u0..u63} uniform, \texttt{k*} FAU constant, or \texttt{t0..t2} temporary; \texttt{\textasciicircum{}} marks discard/last-use). \emph{Example:} {\scriptsize\ttfamily \mnem{FADD.f32} \dst{r0}, \src{r1}, \src{r2}} (\textbf{SPIR-V:} \texttt{OpFAdd (Operand 0)}; \textbf{GLSL:} \texttt{float c = a + b;})
  \item \texttt{src1}: Secondary source operand (register, uniform, or FAU constant; in shifts, supplies the shift count byte). \emph{Example:} {\scriptsize\ttfamily \mnem{LSHIFT.i32} \dst{r0}, \src{r1}, \src{r2}\imod{.b0}} (\textbf{SPIR-V:} \texttt{OpShiftLeftLogical (Shift)}; \textbf{GLSL:} \texttt{uint y = a \textless{}\textless{} s;})
  \item \texttt{opcode2}: Secondary sub-opcode further specializing instructions sharing \texttt{instruction\_class} and \texttt{opcode1} (e.g., conversions, SFU tables, moves). \emph{Example:} {\scriptsize\ttfamily \mnem{F32\_TO\_S32}\imod{.round\_zero} \dst{r0}, \src{r1}} (\textbf{SPIR-V:} \texttt{OpConvertFToS}; \textbf{GLSL:} \texttt{int y = int(x);})
  \item \texttt{opcode3}: Tertiary sub-opcode disambiguating specialized multi-op families (\texttt{FPOW\_SC\_*}, SFU lookup tables, \texttt{ICMP\_MULTI}). \emph{Example:} {\scriptsize\ttfamily \mnem{FPOW\_SC\_DET.f32}\imod{.pow} \dst{r0}, \src{r1}, \src{r2}} (\textbf{SPIR-V:} \texttt{GLSL.std.450 Pow}; \textbf{GLSL:} \texttt{float z = pow(x, y);})
  \item \texttt{dest\_register\_index}: 6-bit destination GPR index (\texttt{r0..r63}) receiving the Stage C writeback result (even-aligned pair \texttt{[rN:rN+1]} for 64-bit ops). \emph{Example:} {\scriptsize\ttfamily \mnem{FMA.f32} \dst{r0}, \src{r1}, \src{r2}, \src{r3}} (\textbf{SPIR-V:} \texttt{SSA Result ID}; \textbf{GLSL:} \texttt{float r0 = fma(r1, r2, r3);})
  \item \hyperref[enum:dest_width_m]{\texttt{dest\_width}}: Stage C destination lane write-mask (\texttt{none} = full 32-bit, \texttt{h0} = low 16-bit, \texttt{h1} = high 16-bit without bit-shifting). \emph{Example:} {\scriptsize\ttfamily \mnem{FADD.v2f16} \dst{r0}\imod{.h1}, \src{r1}, \src{r2}} (\textbf{SPIR-V:} \texttt{OpVectorInsertDynamic / WriteMask}; \textbf{GLSL:} \texttt{dst.y = a.y + b.y;})
  \item \texttt{datatype}: Selects the lane data type (\texttt{f32}, \texttt{f16}, \texttt{v2f16}, \texttt{s32}, \texttt{u32}, \texttt{v2s16}, \texttt{v2u16}, \texttt{v4s8}, \texttt{v4u8}, \texttt{s64}, \texttt{u64}, etc.). \emph{Example:} {\scriptsize\ttfamily \mnem{IADD.v2s16} \dst{r0}, \src{r1}, \src{r2}} (\textbf{SPIR-V:} \texttt{OpTypeInt / OpTypeFloat}; \textbf{GLSL:} \texttt{i16vec2 c = a + b;})
  \item \texttt{instruction\_class}: 4-bit major instruction class selector (\texttt{bits 48:51}) partitioning the 64-bit Valhall opcode space (e.g., \texttt{FMA}, \texttt{CVT}, \texttt{SFU}, \texttt{LS}, \texttt{TEX}, \texttt{BRANCH}). \emph{Example:} {\scriptsize\ttfamily \mnem{FADD.f32} \dst{r0}, \src{r1}, \src{r2}} (\textbf{SPIR-V:} \texttt{N/A (encoding class)}; \textbf{GLSL:} \texttt{// Hardware execution pipe / decode class})
  \item \texttt{fau\_page\_index}: Selects the FAU (Fast Access Uniform) page (\texttt{0..3}) when a source operand references hardware small constants, thread IDs, or lane constants. \emph{Example:} {\scriptsize\ttfamily \mnem{FADD.f32} \dst{r0}, \src{r1}, \src{k1f}} (\textbf{SPIR-V:} \texttt{OpConstant / BuiltIn}; \textbf{GLSL:} \texttt{float y = x + 1.0;})
  \item \hyperref[enum:flow_control_m]{\texttt{flow}}: 4-bit embedded flow-control \& scoreboard wait modifier (\texttt{bits 59:62}: \texttt{none}, \texttt{wait0..wait0126}, \texttt{wait}, \texttt{reconverge}, \texttt{discard}, \texttt{end}). \emph{Example:} {\scriptsize\ttfamily \mnem{FADD.f32}\imod{.wait0.reconverge} \dst{r0}, \src{r1}, \src{r2}} (\textbf{SPIR-V:} \texttt{OpControlBarrier / OpReturn}; \textbf{GLSL:} \texttt{// Scoreboard wait + warp reconvergence})
\end{tightitemize}
\clearpage

\subsection*{FPOW\_SC\_DET}
\addcontentsline{toc}{subsection}{FPOW\_SC\_DET}
Power special-case detector (\texttt{f32}): inspects base \texttt{src0} and exponent \texttt{src1} for edge cases (\texttt{0}, \texttt{$\pm$1}, \texttt{$\pm$Inf}, \texttt{NaN}, and negative bases) according to \texttt{func} (\texttt{.pow} vs \texttt{.powr}) and writes a packed override tag to \texttt{dest} for \texttt{FPOW\_SC\_APPLY}.
\par\textbf{Notes:}
\begin{tightitemize}
  \item Paired with \texttt{FPOW\_SC\_APPLY} (\texttt{bifrost/compiler.h:686}).
  \item \textbf{Source references:} \texttt{bifrost/compiler.h:686}, \texttt{bifrost/bifrost/ISA.xml:4689}, \texttt{isa-v9-v14.xml}.
\end{tightitemize}
\codebox{\mnem{FPOW\_SC\_DET}\imod{\{.flow\}}\imod{\{.func\}} \dst{dest}\imod{\{.dest\_width\}}, \src{src0}, \src{src1}}
\textbf{Operands:}
\begin{tightitemize}\footnotesize
  \item \texttt{\dst{dest}} (\texttt{dest\_register\_index}, \texttt{dest\_width}): 32-bit destination GPR receiving the encoded special-case tag for \texttt{FPOW\_SC\_APPLY}.
  \item \texttt{\src{src0}}: FP32 base \texttt{x} of \texttt{pow(x, y)} / \texttt{powr(x, y)}.
  \item \texttt{\src{src1}}: FP32 exponent \texttt{y} of \texttt{pow(x, y)} / \texttt{powr(x, y)}.
\end{tightitemize}
\par\textbf{Closest SPIR-V Op:} \texttt{None (lowered from GLSL.std.450 Pow)}
\par\textbf{Pseudo-GLSL (Operational Semantics):}
\codebox{dest = (func == powr) ? detectPowrSpecialCase(src0, src1) : detectPowSpecialCase(src0, src1);}
\textbf{Execution Context:}
\begin{tightitemize}\footnotesize
  \item \textbf{Compute stage preloads (\hyperref[sec:csf-spd]{\texttt{SPD\_k}}, \texttt{r16..r23}):} Hardware pre-fills \texttt{r55..r57} with \texttt{WorkGroupID.xyz + JOB\_OFFSET} (from \hyperref[sec:csf-regs]{\texttt{COMPUTESR::JOB\_OFFSET\_X/Y/Z}} \texttt{r34..r36}) and \texttt{r60..r62} with \texttt{LocalInvocationID.xyz} or \texttt{GlobalInvocationID.xyz} (from \texttt{WG\_SIZE} \texttt{r33} \& \texttt{JOB\_SIZE} \texttt{r37..r39}, selected by \texttt{ShaderProgram::preload}) so compute shaders can derive thread/element indices without memory loads.
  \item \textbf{Vertex stage preloads (\hyperref[sec:csf-spd]{\texttt{SPD\_0}/\texttt{SPD\_1}}, \texttt{r16..r19}):} The hardware Index Fetcher reads \hyperref[sec:csf-regs]{\texttt{IDVSSR::INDEX\_BUFFER}} (\texttt{r54:r55}) and pre-fills \texttt{r59, r60} with \texttt{VertexIndex} (\texttt{+ VERTEX\_OFFSET + GLOBAL\_ATTR\_OFF}) and \texttt{r61} with \texttt{InstanceIndex} (\texttt{+ INSTANCE\_OFFSET}) so Position and Varying shaders can pull vertex attributes directly.
  \item \textbf{Fragment stage preloads \& FP modes (\hyperref[sec:csf-spd]{\texttt{SPD\_2}}, \texttt{r20:r21}):} Hardware pre-fills \texttt{r59} with \texttt{floatBitsToUint(FragCoord.z)} and \texttt{r60} with the MSAA sample coverage mask (\texttt{gl\_SampleMaskIn[0]}), while \texttt{ShaderProgram} word~0 latches 32/64 GPR allocation and warp-wide \texttt{flush\_to\_zero\_mode}, \texttt{suppress\_na\_n}, and \texttt{suppress\_inf}.
  \item \textbf{Inline FAU \& Special Operands (\hyperref[sec:csf-tsd-fau]{\texttt{FAU\_k}}, \texttt{r8..r15}):} Any 8-bit source with \texttt{mode=2} reads zero-latency uniform \texttt{u0..u31} from \texttt{FAU\_k} (\texttt{slots64 = bits[63:56]}); \texttt{mode=3} reads \texttt{SmallConstant} (Table~1.2) or \texttt{FauSpecial} (\texttt{pc}, \texttt{lane\_id}, \texttt{lane\_width}, \texttt{tls\_ptr}, \texttt{wls\_ptr}, \texttt{blend\_descriptor\_0..7}).
\end{tightitemize}
\par\textbf{Worked Examples:}
\par\textbf{Example 1 (GLSL/OpenCL `pow` mode (.pow)):}
\codebox{\mnem{FPOW\_SC\_DET.f32}\imod{.pow} \dst{r0}, \src{r1}, \src{r2}}
With \texttt{r1 = -2.0f, r2 = 3.0f} (odd integer): detects negative-base odd-integer sign flip for \texttt{FPOW\_SC\_APPLY}.
\par\textbf{Example 2 (OpenCL `powr` mode (.powr)):}
\codebox{\mnem{FPOW\_SC\_DET.f32}\imod{.powr} \dst{r0}, \src{r1}, \src{r2}}
With \texttt{r1 = 0.0f, r2 = 0.0f}: under \texttt{.powr}, flags \texttt{NaN} (unlike \texttt{.pow}, which returns \texttt{1.0f}).

\begin{center}\footnotesize
\begin{tabularx}{\linewidth}{@{}lllX@{}}
\toprule
Variant & Arch & Class & Opcode Fields \\
\midrule
\texttt{f32} & v9--v14 & \texttt{R2} & \texttt{opcode2=FPOW, opcode3=FPOW\_SC\_DET} \\
\bottomrule
\end{tabularx}
\end{center}
\begin{center}\footnotesize
\begin{tabularx}{0.88\linewidth}{@{}l l X@{}}
\toprule
Field Name & Bits & Type / Fixed Value \\
\midrule
\texttt{src0} & \texttt{0:7} & \texttt{SourceEncoding} \\
\texttt{src1} & \texttt{8:15} & \texttt{SourceEncoding} \\
\texttt{opcode2} & \texttt{16:19} & \texttt{exact=FPOW} \\
\hyperref[enum:fpow_sc_det_function_m]{\texttt{func}} & \texttt{20:21} & \texttt{fpow\_sc\_det\_function\_m} \\
\textit{Reserved} & \texttt{22:23} & \texttt{zero} \\
\texttt{opcode3} & \texttt{24:25} & \texttt{exact=FPOW\_SC\_DET} \\
\textit{Reserved} & \texttt{26:39} & \texttt{zero} \\
\texttt{dest\_register\_index} & \texttt{40:45} & \texttt{uint} \\
\hyperref[enum:dest_width_m]{\texttt{dest\_width}} & \texttt{46:47} & \texttt{dest\_width\_m} \\
\texttt{datatype} & \texttt{48:51} & \texttt{exact=f32} \\
\texttt{instruction\_class} & \texttt{52:56} & \texttt{exact=R2} \\
\texttt{fau\_page\_index} & \texttt{57:58} & \texttt{uint} \\
\hyperref[enum:flow_control_m]{\texttt{flow}} & \texttt{59:62} & \texttt{flow\_control\_m} \\
\textit{Reserved} & \texttt{63} & \texttt{zero} \\
\bottomrule
\end{tabularx}
\end{center}
\textbf{Fields Explained:}
\begin{tightitemize}\footnotesize
  \item \texttt{src0}: Primary source operand (\texttt{r0..r63}, \texttt{u0..u63} uniform, \texttt{k*} FAU constant, or \texttt{t0..t2} temporary; \texttt{\textasciicircum{}} marks discard/last-use). \emph{Example:} {\scriptsize\ttfamily \mnem{FADD.f32} \dst{r0}, \src{r1}, \src{r2}} (\textbf{SPIR-V:} \texttt{OpFAdd (Operand 0)}; \textbf{GLSL:} \texttt{float c = a + b;})
  \item \texttt{src1}: Secondary source operand (register, uniform, or FAU constant; in shifts, supplies the shift count byte). \emph{Example:} {\scriptsize\ttfamily \mnem{LSHIFT.i32} \dst{r0}, \src{r1}, \src{r2}\imod{.b0}} (\textbf{SPIR-V:} \texttt{OpShiftLeftLogical (Shift)}; \textbf{GLSL:} \texttt{uint y = a \textless{}\textless{} s;})
  \item \texttt{opcode2}: Secondary sub-opcode further specializing instructions sharing \texttt{instruction\_class} and \texttt{opcode1} (e.g., conversions, SFU tables, moves). \emph{Example:} {\scriptsize\ttfamily \mnem{F32\_TO\_S32}\imod{.round\_zero} \dst{r0}, \src{r1}} (\textbf{SPIR-V:} \texttt{OpConvertFToS}; \textbf{GLSL:} \texttt{int y = int(x);})
  \item \hyperref[enum:fpow_sc_det_function_m]{\texttt{func}}: Selects \texttt{pow} vs \texttt{powr} special-case detection in \texttt{FPOW\_SC\_DET}. \emph{Example:} {\scriptsize\ttfamily \mnem{FPOW\_SC\_DET}\imod{.pow} \dst{r0}, \src{r1}, \src{r2}} (\textbf{SPIR-V:} \texttt{GLSL.std.450 Pow}; \textbf{GLSL:} \texttt{float y = pow(x, p);})
  \item \texttt{opcode3}: Tertiary sub-opcode disambiguating specialized multi-op families (\texttt{FPOW\_SC\_*}, SFU lookup tables, \texttt{ICMP\_MULTI}). \emph{Example:} {\scriptsize\ttfamily \mnem{FPOW\_SC\_DET.f32}\imod{.pow} \dst{r0}, \src{r1}, \src{r2}} (\textbf{SPIR-V:} \texttt{GLSL.std.450 Pow}; \textbf{GLSL:} \texttt{float z = pow(x, y);})
  \item \texttt{dest\_register\_index}: 6-bit destination GPR index (\texttt{r0..r63}) receiving the Stage C writeback result (even-aligned pair \texttt{[rN:rN+1]} for 64-bit ops). \emph{Example:} {\scriptsize\ttfamily \mnem{FMA.f32} \dst{r0}, \src{r1}, \src{r2}, \src{r3}} (\textbf{SPIR-V:} \texttt{SSA Result ID}; \textbf{GLSL:} \texttt{float r0 = fma(r1, r2, r3);})
  \item \hyperref[enum:dest_width_m]{\texttt{dest\_width}}: Stage C destination lane write-mask (\texttt{none} = full 32-bit, \texttt{h0} = low 16-bit, \texttt{h1} = high 16-bit without bit-shifting). \emph{Example:} {\scriptsize\ttfamily \mnem{FADD.v2f16} \dst{r0}\imod{.h1}, \src{r1}, \src{r2}} (\textbf{SPIR-V:} \texttt{OpVectorInsertDynamic / WriteMask}; \textbf{GLSL:} \texttt{dst.y = a.y + b.y;})
  \item \texttt{datatype}: Selects the lane data type (\texttt{f32}, \texttt{f16}, \texttt{v2f16}, \texttt{s32}, \texttt{u32}, \texttt{v2s16}, \texttt{v2u16}, \texttt{v4s8}, \texttt{v4u8}, \texttt{s64}, \texttt{u64}, etc.). \emph{Example:} {\scriptsize\ttfamily \mnem{IADD.v2s16} \dst{r0}, \src{r1}, \src{r2}} (\textbf{SPIR-V:} \texttt{OpTypeInt / OpTypeFloat}; \textbf{GLSL:} \texttt{i16vec2 c = a + b;})
  \item \texttt{instruction\_class}: 4-bit major instruction class selector (\texttt{bits 48:51}) partitioning the 64-bit Valhall opcode space (e.g., \texttt{FMA}, \texttt{CVT}, \texttt{SFU}, \texttt{LS}, \texttt{TEX}, \texttt{BRANCH}). \emph{Example:} {\scriptsize\ttfamily \mnem{FADD.f32} \dst{r0}, \src{r1}, \src{r2}} (\textbf{SPIR-V:} \texttt{N/A (encoding class)}; \textbf{GLSL:} \texttt{// Hardware execution pipe / decode class})
  \item \texttt{fau\_page\_index}: Selects the FAU (Fast Access Uniform) page (\texttt{0..3}) when a source operand references hardware small constants, thread IDs, or lane constants. \emph{Example:} {\scriptsize\ttfamily \mnem{FADD.f32} \dst{r0}, \src{r1}, \src{k1f}} (\textbf{SPIR-V:} \texttt{OpConstant / BuiltIn}; \textbf{GLSL:} \texttt{float y = x + 1.0;})
  \item \hyperref[enum:flow_control_m]{\texttt{flow}}: 4-bit embedded flow-control \& scoreboard wait modifier (\texttt{bits 59:62}: \texttt{none}, \texttt{wait0..wait0126}, \texttt{wait}, \texttt{reconverge}, \texttt{discard}, \texttt{end}). \emph{Example:} {\scriptsize\ttfamily \mnem{FADD.f32}\imod{.wait0.reconverge} \dst{r0}, \src{r1}, \src{r2}} (\textbf{SPIR-V:} \texttt{OpControlBarrier / OpReturn}; \textbf{GLSL:} \texttt{// Scoreboard wait + warp reconvergence})
\end{tightitemize}
\clearpage

\subsection*{FRCBRT\_APPROX\_A}
\addcontentsline{toc}{subsection}{FRCBRT\_APPROX\_A}
Stage-A reciprocal cube-root (\texttt{x\textasciicircum{}(-1/3)}) SFU helper (\texttt{f32}): computes the initial low-precision seed approximation of \texttt{src0\textasciicircum{}(-1/3)} in \texttt{dest} for the multi-stage cube-root sequence.
\par\textbf{Notes:}
\begin{tightitemize}
  \item SFU table seed instruction (\texttt{bifrost/ISA.xml:4700}).
  \item \textbf{Source references:} \texttt{bifrost/bifrost/ISA.xml:4700}, \texttt{isa-v9-v14.xml}.
\end{tightitemize}
\codebox{\mnem{FRCBRT\_APPROX\_A}\imod{\{.flow\}} \dst{dest}\imod{\{.dest\_width\}}, \src{src0}\imod{\{.widen0\}}\imod{\{.abs0\}}\imod{\{.neg0\}}}
\textbf{Operands:}
\begin{tightitemize}\footnotesize
  \item \texttt{\dst{dest}} (\texttt{dest\_register\_index}, \texttt{dest\_width}): FP32 destination GPR receiving the initial \texttt{x\textasciicircum{}(-1/3)} approximation seed.
  \item \texttt{\src{src0}}: FP32 input \texttt{x}.
\end{tightitemize}
\par\textbf{Closest SPIR-V Op:} \texttt{None (SFU cube-root seed A)}
\par\textbf{Pseudo-GLSL (Operational Semantics):}
\codebox{dest = approxInvCbrtStageA(src0); // initial x\textasciicircum{}(-1/3) seed}
\textbf{Execution Context:}
\begin{tightitemize}\footnotesize
  \item \textbf{Compute stage preloads (\hyperref[sec:csf-spd]{\texttt{SPD\_k}}, \texttt{r16..r23}):} Hardware pre-fills \texttt{r55..r57} with \texttt{WorkGroupID.xyz + JOB\_OFFSET} (from \hyperref[sec:csf-regs]{\texttt{COMPUTESR::JOB\_OFFSET\_X/Y/Z}} \texttt{r34..r36}) and \texttt{r60..r62} with \texttt{LocalInvocationID.xyz} or \texttt{GlobalInvocationID.xyz} (from \texttt{WG\_SIZE} \texttt{r33} \& \texttt{JOB\_SIZE} \texttt{r37..r39}, selected by \texttt{ShaderProgram::preload}) so compute shaders can derive thread/element indices without memory loads.
  \item \textbf{Vertex stage preloads (\hyperref[sec:csf-spd]{\texttt{SPD\_0}/\texttt{SPD\_1}}, \texttt{r16..r19}):} The hardware Index Fetcher reads \hyperref[sec:csf-regs]{\texttt{IDVSSR::INDEX\_BUFFER}} (\texttt{r54:r55}) and pre-fills \texttt{r59, r60} with \texttt{VertexIndex} (\texttt{+ VERTEX\_OFFSET + GLOBAL\_ATTR\_OFF}) and \texttt{r61} with \texttt{InstanceIndex} (\texttt{+ INSTANCE\_OFFSET}) so Position and Varying shaders can pull vertex attributes directly.
  \item \textbf{Fragment stage preloads \& FP modes (\hyperref[sec:csf-spd]{\texttt{SPD\_2}}, \texttt{r20:r21}):} Hardware pre-fills \texttt{r59} with \texttt{floatBitsToUint(FragCoord.z)} and \texttt{r60} with the MSAA sample coverage mask (\texttt{gl\_SampleMaskIn[0]}), while \texttt{ShaderProgram} word~0 latches 32/64 GPR allocation and warp-wide \texttt{flush\_to\_zero\_mode}, \texttt{suppress\_na\_n}, and \texttt{suppress\_inf}.
  \item \textbf{Inline FAU \& Special Operands (\hyperref[sec:csf-tsd-fau]{\texttt{FAU\_k}}, \texttt{r8..r15}):} Any 8-bit source with \texttt{mode=2} reads zero-latency uniform \texttt{u0..u31} from \texttt{FAU\_k} (\texttt{slots64 = bits[63:56]}); \texttt{mode=3} reads \texttt{SmallConstant} (Table~1.2) or \texttt{FauSpecial} (\texttt{pc}, \texttt{lane\_id}, \texttt{lane\_width}, \texttt{tls\_ptr}, \texttt{wls\_ptr}, \texttt{blend\_descriptor\_0..7}).
\end{tightitemize}
\par\textbf{Worked Examples:}
\par\textbf{Example 1 (Initial `x\textasciicircum{}(-1/3)` seed for `x = 8.0f`):}
\codebox{\mnem{FRCBRT\_APPROX\_A.f32} \dst{r0}, \src{r1}}
With \texttt{r1 = 8.0f}: returns initial reciprocal cube-root seed \texttt{\textasciitilde{}0.5f} for Newton-Raphson refinement.
\par\textbf{Example 2 (Negative input preservation):}
\codebox{\mnem{FRCBRT\_APPROX\_A.f32} \dst{r0}, \src{r1}}
With \texttt{r1 = -27.0f}: produces negative seed \texttt{\textasciitilde{}-0.333f} preserving sign of the cube root.

\begin{center}\footnotesize
\begin{tabularx}{\linewidth}{@{}lllX@{}}
\toprule
Variant & Arch & Class & Opcode Fields \\
\midrule
\texttt{(none)} & v9--v14 & \texttt{R1} & \texttt{opcode2=FRCBRT\_APPROX\_A, opcode1=FRCBRT} \\
\bottomrule
\end{tabularx}
\end{center}
\begin{center}\footnotesize
\begin{tabularx}{0.88\linewidth}{@{}l l X@{}}
\toprule
Field Name & Bits & Type / Fixed Value \\
\midrule
\texttt{src0} & \texttt{0:7} & \texttt{SourceEncoding} \\
\textit{Reserved} & \texttt{8:15} & \texttt{zero} \\
\texttt{opcode2} & \texttt{16:20} & \texttt{exact=FRCBRT\_APPROX\_A} \\
\textit{Reserved} & \texttt{21:27} & \texttt{zero} \\
\hyperref[enum:widen_m]{\texttt{widen0}} & \texttt{28:29} & \texttt{widen\_m (def none)} \\
\textit{Reserved} & \texttt{30:37} & \texttt{zero} \\
\hyperref[enum:neg_m]{\texttt{neg0}} & \texttt{38} & \texttt{neg\_m (def none)} \\
\hyperref[enum:abs_m]{\texttt{abs0}} & \texttt{39} & \texttt{abs\_m (def none)} \\
\texttt{dest\_register\_index} & \texttt{40:45} & \texttt{uint} \\
\hyperref[enum:dest_width_m]{\texttt{dest\_width}} & \texttt{46:47} & \texttt{dest\_width\_m} \\
\texttt{opcode1} & \texttt{48:51} & \texttt{exact=FRCBRT} \\
\texttt{instruction\_class} & \texttt{52:56} & \texttt{exact=R1} \\
\texttt{fau\_page\_index} & \texttt{57:58} & \texttt{uint} \\
\hyperref[enum:flow_control_m]{\texttt{flow}} & \texttt{59:62} & \texttt{flow\_control\_m} \\
\textit{Reserved} & \texttt{63} & \texttt{zero} \\
\bottomrule
\end{tabularx}
\end{center}
\textbf{Fields Explained:}
\begin{tightitemize}\footnotesize
  \item \texttt{src0}: Primary source operand (\texttt{r0..r63}, \texttt{u0..u63} uniform, \texttt{k*} FAU constant, or \texttt{t0..t2} temporary; \texttt{\textasciicircum{}} marks discard/last-use). \emph{Example:} {\scriptsize\ttfamily \mnem{FADD.f32} \dst{r0}, \src{r1}, \src{r2}} (\textbf{SPIR-V:} \texttt{OpFAdd (Operand 0)}; \textbf{GLSL:} \texttt{float c = a + b;})
  \item \texttt{opcode2}: Secondary sub-opcode further specializing instructions sharing \texttt{instruction\_class} and \texttt{opcode1} (e.g., conversions, SFU tables, moves). \emph{Example:} {\scriptsize\ttfamily \mnem{F32\_TO\_S32}\imod{.round\_zero} \dst{r0}, \src{r1}} (\textbf{SPIR-V:} \texttt{OpConvertFToS}; \textbf{GLSL:} \texttt{int y = int(x);})
  \item \hyperref[enum:widen_m]{\texttt{widen0}}: Stage A widening extractor on \texttt{src0} (promotes FP16 \texttt{h0}/\texttt{h1} to FP32, or sign/zero-extends \texttt{b0..b3}/\texttt{h0..h1}/\texttt{w0..w1} to 32/64-bit before \texttt{abs}/\texttt{neg}). \emph{Example:} {\scriptsize\ttfamily \mnem{IADD.s32} \dst{r0}, \src{r1}\imod{.h0}, \src{r2}} (\textbf{SPIR-V:} \texttt{OpSConvert / OpFConvert}; \textbf{GLSL:} \texttt{int c = int(int16\_t(a)) + b;})
  \item \hyperref[enum:neg_m]{\texttt{neg0}}: Stage A floating-point negation on \texttt{src0} (flips IEEE sign bit after swizzle/widen and \texttt{abs0}). \emph{Example:} {\scriptsize\ttfamily \mnem{FADD.f32} \dst{r0}, \src{r1}\imod{.neg}, \src{r2}} (\textbf{SPIR-V:} \texttt{OpFNegate}; \textbf{GLSL:} \texttt{float c = -a + b;})
  \item \hyperref[enum:abs_m]{\texttt{abs0}}: Stage A floating-point absolute value on \texttt{src0} (clears IEEE sign bit after swizzle/widen, before \texttt{neg0}). \emph{Example:} {\scriptsize\ttfamily \mnem{FADD.f32} \dst{r0}, \src{r1}\imod{.abs}, \src{r2}} (\textbf{SPIR-V:} \texttt{GLSL.std.450 FAbs}; \textbf{GLSL:} \texttt{float c = abs(a) + b;})
  \item \texttt{dest\_register\_index}: 6-bit destination GPR index (\texttt{r0..r63}) receiving the Stage C writeback result (even-aligned pair \texttt{[rN:rN+1]} for 64-bit ops). \emph{Example:} {\scriptsize\ttfamily \mnem{FMA.f32} \dst{r0}, \src{r1}, \src{r2}, \src{r3}} (\textbf{SPIR-V:} \texttt{SSA Result ID}; \textbf{GLSL:} \texttt{float r0 = fma(r1, r2, r3);})
  \item \hyperref[enum:dest_width_m]{\texttt{dest\_width}}: Stage C destination lane write-mask (\texttt{none} = full 32-bit, \texttt{h0} = low 16-bit, \texttt{h1} = high 16-bit without bit-shifting). \emph{Example:} {\scriptsize\ttfamily \mnem{FADD.v2f16} \dst{r0}\imod{.h1}, \src{r1}, \src{r2}} (\textbf{SPIR-V:} \texttt{OpVectorInsertDynamic / WriteMask}; \textbf{GLSL:} \texttt{dst.y = a.y + b.y;})
  \item \texttt{opcode1}: Primary sub-opcode within the instruction class uniquely identifying the operation family. \emph{Example:} {\scriptsize\ttfamily \mnem{IADD.s32} \dst{r0}, \src{r1}, \src{r2}} (\textbf{SPIR-V:} \texttt{OpIAdd (operation selector)}; \textbf{GLSL:} \texttt{// Primary opcode bits})
  \item \texttt{instruction\_class}: 4-bit major instruction class selector (\texttt{bits 48:51}) partitioning the 64-bit Valhall opcode space (e.g., \texttt{FMA}, \texttt{CVT}, \texttt{SFU}, \texttt{LS}, \texttt{TEX}, \texttt{BRANCH}). \emph{Example:} {\scriptsize\ttfamily \mnem{FADD.f32} \dst{r0}, \src{r1}, \src{r2}} (\textbf{SPIR-V:} \texttt{N/A (encoding class)}; \textbf{GLSL:} \texttt{// Hardware execution pipe / decode class})
  \item \texttt{fau\_page\_index}: Selects the FAU (Fast Access Uniform) page (\texttt{0..3}) when a source operand references hardware small constants, thread IDs, or lane constants. \emph{Example:} {\scriptsize\ttfamily \mnem{FADD.f32} \dst{r0}, \src{r1}, \src{k1f}} (\textbf{SPIR-V:} \texttt{OpConstant / BuiltIn}; \textbf{GLSL:} \texttt{float y = x + 1.0;})
  \item \hyperref[enum:flow_control_m]{\texttt{flow}}: 4-bit embedded flow-control \& scoreboard wait modifier (\texttt{bits 59:62}: \texttt{none}, \texttt{wait0..wait0126}, \texttt{wait}, \texttt{reconverge}, \texttt{discard}, \texttt{end}). \emph{Example:} {\scriptsize\ttfamily \mnem{FADD.f32}\imod{.wait0.reconverge} \dst{r0}, \src{r1}, \src{r2}} (\textbf{SPIR-V:} \texttt{OpControlBarrier / OpReturn}; \textbf{GLSL:} \texttt{// Scoreboard wait + warp reconvergence})
\end{tightitemize}
\clearpage

\subsection*{FRCBRT\_APPROX\_B}
\addcontentsline{toc}{subsection}{FRCBRT\_APPROX\_B}
Stage-B reciprocal cube-root (\texttt{x\textasciicircum{}(-1/3)}) SFU helper (\texttt{f32}): applies the second-stage polynomial refinement to the intermediate value \texttt{src0} and writes the refined estimate to \texttt{dest}.
\par\textbf{Notes:}
\begin{tightitemize}
  \item SFU table seed instruction (\texttt{bifrost/ISA.xml:4722}).
  \item \textbf{Source references:} \texttt{bifrost/bifrost/ISA.xml:4722}, \texttt{isa-v9-v14.xml}.
\end{tightitemize}
\codebox{\mnem{FRCBRT\_APPROX\_B}\imod{\{.flow\}} \dst{dest}\imod{\{.dest\_width\}}, \src{src0}}
\textbf{Operands:}
\begin{tightitemize}\footnotesize
  \item \texttt{\dst{dest}} (\texttt{dest\_register\_index}, \texttt{dest\_width}): FP32 destination GPR receiving the stage-B \texttt{x\textasciicircum{}(-1/3)} refinement value.
  \item \texttt{\src{src0}}: Intermediate FP32 value from stage A.
\end{tightitemize}
\par\textbf{Closest SPIR-V Op:} \texttt{None (SFU cube-root seed B)}
\par\textbf{Pseudo-GLSL (Operational Semantics):}
\codebox{dest = approxInvCbrtStageB(src0); // stage-B refinement}
\textbf{Execution Context:}
\begin{tightitemize}\footnotesize
  \item \textbf{Compute stage preloads (\hyperref[sec:csf-spd]{\texttt{SPD\_k}}, \texttt{r16..r23}):} Hardware pre-fills \texttt{r55..r57} with \texttt{WorkGroupID.xyz + JOB\_OFFSET} (from \hyperref[sec:csf-regs]{\texttt{COMPUTESR::JOB\_OFFSET\_X/Y/Z}} \texttt{r34..r36}) and \texttt{r60..r62} with \texttt{LocalInvocationID.xyz} or \texttt{GlobalInvocationID.xyz} (from \texttt{WG\_SIZE} \texttt{r33} \& \texttt{JOB\_SIZE} \texttt{r37..r39}, selected by \texttt{ShaderProgram::preload}) so compute shaders can derive thread/element indices without memory loads.
  \item \textbf{Vertex stage preloads (\hyperref[sec:csf-spd]{\texttt{SPD\_0}/\texttt{SPD\_1}}, \texttt{r16..r19}):} The hardware Index Fetcher reads \hyperref[sec:csf-regs]{\texttt{IDVSSR::INDEX\_BUFFER}} (\texttt{r54:r55}) and pre-fills \texttt{r59, r60} with \texttt{VertexIndex} (\texttt{+ VERTEX\_OFFSET + GLOBAL\_ATTR\_OFF}) and \texttt{r61} with \texttt{InstanceIndex} (\texttt{+ INSTANCE\_OFFSET}) so Position and Varying shaders can pull vertex attributes directly.
  \item \textbf{Fragment stage preloads \& FP modes (\hyperref[sec:csf-spd]{\texttt{SPD\_2}}, \texttt{r20:r21}):} Hardware pre-fills \texttt{r59} with \texttt{floatBitsToUint(FragCoord.z)} and \texttt{r60} with the MSAA sample coverage mask (\texttt{gl\_SampleMaskIn[0]}), while \texttt{ShaderProgram} word~0 latches 32/64 GPR allocation and warp-wide \texttt{flush\_to\_zero\_mode}, \texttt{suppress\_na\_n}, and \texttt{suppress\_inf}.
  \item \textbf{Inline FAU \& Special Operands (\hyperref[sec:csf-tsd-fau]{\texttt{FAU\_k}}, \texttt{r8..r15}):} Any 8-bit source with \texttt{mode=2} reads zero-latency uniform \texttt{u0..u31} from \texttt{FAU\_k} (\texttt{slots64 = bits[63:56]}); \texttt{mode=3} reads \texttt{SmallConstant} (Table~1.2) or \texttt{FauSpecial} (\texttt{pc}, \texttt{lane\_id}, \texttt{lane\_width}, \texttt{tls\_ptr}, \texttt{wls\_ptr}, \texttt{blend\_descriptor\_0..7}).
\end{tightitemize}
\par\textbf{Worked Examples:}
\par\textbf{Example 1 (Second-stage cbrt polynomial refinement):}
\codebox{\mnem{FRCBRT\_APPROX\_B.f32} \dst{r0}, \src{r1}}
Refines the exponent/mantissa seed from \texttt{FRCBRT\_APPROX\_A} to increase bit accuracy before final FMA steps.
\par\textbf{Example 2 (Subnormal/edge handling):}
\codebox{\mnem{FRCBRT\_APPROX\_B.f32} \dst{r0}, \src{r1}}
Handles normalized intermediate from stage A across all finite FP32 exponents.

\begin{center}\footnotesize
\begin{tabularx}{\linewidth}{@{}lllX@{}}
\toprule
Variant & Arch & Class & Opcode Fields \\
\midrule
\texttt{(none)} & v9--v14 & \texttt{R1} & \texttt{opcode2=FRCBRT\_APPROX\_B, opcode1=FRCBRT} \\
\bottomrule
\end{tabularx}
\end{center}
\begin{center}\footnotesize
\begin{tabularx}{0.88\linewidth}{@{}l l X@{}}
\toprule
Field Name & Bits & Type / Fixed Value \\
\midrule
\texttt{src0} & \texttt{0:7} & \texttt{SourceEncoding} \\
\textit{Reserved} & \texttt{8:15} & \texttt{zero} \\
\texttt{opcode2} & \texttt{16:20} & \texttt{exact=FRCBRT\_APPROX\_B} \\
\textit{Reserved} & \texttt{21:39} & \texttt{zero} \\
\texttt{dest\_register\_index} & \texttt{40:45} & \texttt{uint} \\
\hyperref[enum:dest_width_m]{\texttt{dest\_width}} & \texttt{46:47} & \texttt{dest\_width\_m} \\
\texttt{opcode1} & \texttt{48:51} & \texttt{exact=FRCBRT} \\
\texttt{instruction\_class} & \texttt{52:56} & \texttt{exact=R1} \\
\texttt{fau\_page\_index} & \texttt{57:58} & \texttt{uint} \\
\hyperref[enum:flow_control_m]{\texttt{flow}} & \texttt{59:62} & \texttt{flow\_control\_m} \\
\textit{Reserved} & \texttt{63} & \texttt{zero} \\
\bottomrule
\end{tabularx}
\end{center}
\textbf{Fields Explained:}
\begin{tightitemize}\footnotesize
  \item \texttt{src0}: Primary source operand (\texttt{r0..r63}, \texttt{u0..u63} uniform, \texttt{k*} FAU constant, or \texttt{t0..t2} temporary; \texttt{\textasciicircum{}} marks discard/last-use). \emph{Example:} {\scriptsize\ttfamily \mnem{FADD.f32} \dst{r0}, \src{r1}, \src{r2}} (\textbf{SPIR-V:} \texttt{OpFAdd (Operand 0)}; \textbf{GLSL:} \texttt{float c = a + b;})
  \item \texttt{opcode2}: Secondary sub-opcode further specializing instructions sharing \texttt{instruction\_class} and \texttt{opcode1} (e.g., conversions, SFU tables, moves). \emph{Example:} {\scriptsize\ttfamily \mnem{F32\_TO\_S32}\imod{.round\_zero} \dst{r0}, \src{r1}} (\textbf{SPIR-V:} \texttt{OpConvertFToS}; \textbf{GLSL:} \texttt{int y = int(x);})
  \item \texttt{dest\_register\_index}: 6-bit destination GPR index (\texttt{r0..r63}) receiving the Stage C writeback result (even-aligned pair \texttt{[rN:rN+1]} for 64-bit ops). \emph{Example:} {\scriptsize\ttfamily \mnem{FMA.f32} \dst{r0}, \src{r1}, \src{r2}, \src{r3}} (\textbf{SPIR-V:} \texttt{SSA Result ID}; \textbf{GLSL:} \texttt{float r0 = fma(r1, r2, r3);})
  \item \hyperref[enum:dest_width_m]{\texttt{dest\_width}}: Stage C destination lane write-mask (\texttt{none} = full 32-bit, \texttt{h0} = low 16-bit, \texttt{h1} = high 16-bit without bit-shifting). \emph{Example:} {\scriptsize\ttfamily \mnem{FADD.v2f16} \dst{r0}\imod{.h1}, \src{r1}, \src{r2}} (\textbf{SPIR-V:} \texttt{OpVectorInsertDynamic / WriteMask}; \textbf{GLSL:} \texttt{dst.y = a.y + b.y;})
  \item \texttt{opcode1}: Primary sub-opcode within the instruction class uniquely identifying the operation family. \emph{Example:} {\scriptsize\ttfamily \mnem{IADD.s32} \dst{r0}, \src{r1}, \src{r2}} (\textbf{SPIR-V:} \texttt{OpIAdd (operation selector)}; \textbf{GLSL:} \texttt{// Primary opcode bits})
  \item \texttt{instruction\_class}: 4-bit major instruction class selector (\texttt{bits 48:51}) partitioning the 64-bit Valhall opcode space (e.g., \texttt{FMA}, \texttt{CVT}, \texttt{SFU}, \texttt{LS}, \texttt{TEX}, \texttt{BRANCH}). \emph{Example:} {\scriptsize\ttfamily \mnem{FADD.f32} \dst{r0}, \src{r1}, \src{r2}} (\textbf{SPIR-V:} \texttt{N/A (encoding class)}; \textbf{GLSL:} \texttt{// Hardware execution pipe / decode class})
  \item \texttt{fau\_page\_index}: Selects the FAU (Fast Access Uniform) page (\texttt{0..3}) when a source operand references hardware small constants, thread IDs, or lane constants. \emph{Example:} {\scriptsize\ttfamily \mnem{FADD.f32} \dst{r0}, \src{r1}, \src{k1f}} (\textbf{SPIR-V:} \texttt{OpConstant / BuiltIn}; \textbf{GLSL:} \texttt{float y = x + 1.0;})
  \item \hyperref[enum:flow_control_m]{\texttt{flow}}: 4-bit embedded flow-control \& scoreboard wait modifier (\texttt{bits 59:62}: \texttt{none}, \texttt{wait0..wait0126}, \texttt{wait}, \texttt{reconverge}, \texttt{discard}, \texttt{end}). \emph{Example:} {\scriptsize\ttfamily \mnem{FADD.f32}\imod{.wait0.reconverge} \dst{r0}, \src{r1}, \src{r2}} (\textbf{SPIR-V:} \texttt{OpControlBarrier / OpReturn}; \textbf{GLSL:} \texttt{// Scoreboard wait + warp reconvergence})
\end{tightitemize}
\clearpage

\subsection*{FRCBRT\_APPROX\_C}
\addcontentsline{toc}{subsection}{FRCBRT\_APPROX\_C}
Stage-C reciprocal cube-root (\texttt{x\textasciicircum{}(-1/3)}) SFU helper (\texttt{f32}): applies the final refinement step to \texttt{src0} and writes the resulting \texttt{x\textasciicircum{}(-1/3)} approximation to \texttt{dest}.
\par\textbf{Notes:}
\begin{tightitemize}
  \item SFU table seed instruction (\texttt{bifrost/ISA.xml:4726}).
  \item \textbf{Source references:} \texttt{bifrost/bifrost/ISA.xml:4726}, \texttt{isa-v9-v14.xml}.
\end{tightitemize}
\codebox{\mnem{FRCBRT\_APPROX\_C}\imod{\{.flow\}} \dst{dest}\imod{\{.dest\_width\}}, \src{src0}}
\textbf{Operands:}
\begin{tightitemize}\footnotesize
  \item \texttt{\dst{dest}} (\texttt{dest\_register\_index}, \texttt{dest\_width}): FP32 destination GPR receiving the stage-C \texttt{x\textasciicircum{}(-1/3)} refinement value.
  \item \texttt{\src{src0}}: Intermediate FP32 value from stage B.
\end{tightitemize}
\par\textbf{Closest SPIR-V Op:} \texttt{None (SFU cube-root seed C)}
\par\textbf{Pseudo-GLSL (Operational Semantics):}
\codebox{dest = approxInvCbrtStageC(src0); // stage-C refinement}
\textbf{Execution Context:}
\begin{tightitemize}\footnotesize
  \item \textbf{Compute stage preloads (\hyperref[sec:csf-spd]{\texttt{SPD\_k}}, \texttt{r16..r23}):} Hardware pre-fills \texttt{r55..r57} with \texttt{WorkGroupID.xyz + JOB\_OFFSET} (from \hyperref[sec:csf-regs]{\texttt{COMPUTESR::JOB\_OFFSET\_X/Y/Z}} \texttt{r34..r36}) and \texttt{r60..r62} with \texttt{LocalInvocationID.xyz} or \texttt{GlobalInvocationID.xyz} (from \texttt{WG\_SIZE} \texttt{r33} \& \texttt{JOB\_SIZE} \texttt{r37..r39}, selected by \texttt{ShaderProgram::preload}) so compute shaders can derive thread/element indices without memory loads.
  \item \textbf{Vertex stage preloads (\hyperref[sec:csf-spd]{\texttt{SPD\_0}/\texttt{SPD\_1}}, \texttt{r16..r19}):} The hardware Index Fetcher reads \hyperref[sec:csf-regs]{\texttt{IDVSSR::INDEX\_BUFFER}} (\texttt{r54:r55}) and pre-fills \texttt{r59, r60} with \texttt{VertexIndex} (\texttt{+ VERTEX\_OFFSET + GLOBAL\_ATTR\_OFF}) and \texttt{r61} with \texttt{InstanceIndex} (\texttt{+ INSTANCE\_OFFSET}) so Position and Varying shaders can pull vertex attributes directly.
  \item \textbf{Fragment stage preloads \& FP modes (\hyperref[sec:csf-spd]{\texttt{SPD\_2}}, \texttt{r20:r21}):} Hardware pre-fills \texttt{r59} with \texttt{floatBitsToUint(FragCoord.z)} and \texttt{r60} with the MSAA sample coverage mask (\texttt{gl\_SampleMaskIn[0]}), while \texttt{ShaderProgram} word~0 latches 32/64 GPR allocation and warp-wide \texttt{flush\_to\_zero\_mode}, \texttt{suppress\_na\_n}, and \texttt{suppress\_inf}.
  \item \textbf{Inline FAU \& Special Operands (\hyperref[sec:csf-tsd-fau]{\texttt{FAU\_k}}, \texttt{r8..r15}):} Any 8-bit source with \texttt{mode=2} reads zero-latency uniform \texttt{u0..u31} from \texttt{FAU\_k} (\texttt{slots64 = bits[63:56]}); \texttt{mode=3} reads \texttt{SmallConstant} (Table~1.2) or \texttt{FauSpecial} (\texttt{pc}, \texttt{lane\_id}, \texttt{lane\_width}, \texttt{tls\_ptr}, \texttt{wls\_ptr}, \texttt{blend\_descriptor\_0..7}).
\end{tightitemize}
\par\textbf{Worked Examples:}
\par\textbf{Example 1 (Third-stage cbrt refinement):}
\codebox{\mnem{FRCBRT\_APPROX\_C.f32} \dst{r0}, \src{r1}}
Completes the 3-stage hardware reciprocal cube-root seed pipeline (\texttt{A -\textgreater{} B -\textgreater{} C}).
\par\textbf{Example 2 (Final FMA pairing):}
\codebox{\mnem{FRCBRT\_APPROX\_C.f32} \dst{r0}, \src{r1}}
Outputs the final refined reciprocal cube-root approximation ready for \texttt{FMUL} to recover \texttt{cbrt(x) = x * x\textasciicircum{}(-2/3)}.

\begin{center}\footnotesize
\begin{tabularx}{\linewidth}{@{}lllX@{}}
\toprule
Variant & Arch & Class & Opcode Fields \\
\midrule
\texttt{(none)} & v9--v14 & \texttt{R1} & \texttt{opcode2=FRCBRT\_APPROX\_C, opcode1=FRCBRT} \\
\bottomrule
\end{tabularx}
\end{center}
\begin{center}\footnotesize
\begin{tabularx}{0.88\linewidth}{@{}l l X@{}}
\toprule
Field Name & Bits & Type / Fixed Value \\
\midrule
\texttt{src0} & \texttt{0:7} & \texttt{SourceEncoding} \\
\textit{Reserved} & \texttt{8:15} & \texttt{zero} \\
\texttt{opcode2} & \texttt{16:20} & \texttt{exact=FRCBRT\_APPROX\_C} \\
\textit{Reserved} & \texttt{21:39} & \texttt{zero} \\
\texttt{dest\_register\_index} & \texttt{40:45} & \texttt{uint} \\
\hyperref[enum:dest_width_m]{\texttt{dest\_width}} & \texttt{46:47} & \texttt{dest\_width\_m} \\
\texttt{opcode1} & \texttt{48:51} & \texttt{exact=FRCBRT} \\
\texttt{instruction\_class} & \texttt{52:56} & \texttt{exact=R1} \\
\texttt{fau\_page\_index} & \texttt{57:58} & \texttt{uint} \\
\hyperref[enum:flow_control_m]{\texttt{flow}} & \texttt{59:62} & \texttt{flow\_control\_m} \\
\textit{Reserved} & \texttt{63} & \texttt{zero} \\
\bottomrule
\end{tabularx}
\end{center}
\textbf{Fields Explained:}
\begin{tightitemize}\footnotesize
  \item \texttt{src0}: Primary source operand (\texttt{r0..r63}, \texttt{u0..u63} uniform, \texttt{k*} FAU constant, or \texttt{t0..t2} temporary; \texttt{\textasciicircum{}} marks discard/last-use). \emph{Example:} {\scriptsize\ttfamily \mnem{FADD.f32} \dst{r0}, \src{r1}, \src{r2}} (\textbf{SPIR-V:} \texttt{OpFAdd (Operand 0)}; \textbf{GLSL:} \texttt{float c = a + b;})
  \item \texttt{opcode2}: Secondary sub-opcode further specializing instructions sharing \texttt{instruction\_class} and \texttt{opcode1} (e.g., conversions, SFU tables, moves). \emph{Example:} {\scriptsize\ttfamily \mnem{F32\_TO\_S32}\imod{.round\_zero} \dst{r0}, \src{r1}} (\textbf{SPIR-V:} \texttt{OpConvertFToS}; \textbf{GLSL:} \texttt{int y = int(x);})
  \item \texttt{dest\_register\_index}: 6-bit destination GPR index (\texttt{r0..r63}) receiving the Stage C writeback result (even-aligned pair \texttt{[rN:rN+1]} for 64-bit ops). \emph{Example:} {\scriptsize\ttfamily \mnem{FMA.f32} \dst{r0}, \src{r1}, \src{r2}, \src{r3}} (\textbf{SPIR-V:} \texttt{SSA Result ID}; \textbf{GLSL:} \texttt{float r0 = fma(r1, r2, r3);})
  \item \hyperref[enum:dest_width_m]{\texttt{dest\_width}}: Stage C destination lane write-mask (\texttt{none} = full 32-bit, \texttt{h0} = low 16-bit, \texttt{h1} = high 16-bit without bit-shifting). \emph{Example:} {\scriptsize\ttfamily \mnem{FADD.v2f16} \dst{r0}\imod{.h1}, \src{r1}, \src{r2}} (\textbf{SPIR-V:} \texttt{OpVectorInsertDynamic / WriteMask}; \textbf{GLSL:} \texttt{dst.y = a.y + b.y;})
  \item \texttt{opcode1}: Primary sub-opcode within the instruction class uniquely identifying the operation family. \emph{Example:} {\scriptsize\ttfamily \mnem{IADD.s32} \dst{r0}, \src{r1}, \src{r2}} (\textbf{SPIR-V:} \texttt{OpIAdd (operation selector)}; \textbf{GLSL:} \texttt{// Primary opcode bits})
  \item \texttt{instruction\_class}: 4-bit major instruction class selector (\texttt{bits 48:51}) partitioning the 64-bit Valhall opcode space (e.g., \texttt{FMA}, \texttt{CVT}, \texttt{SFU}, \texttt{LS}, \texttt{TEX}, \texttt{BRANCH}). \emph{Example:} {\scriptsize\ttfamily \mnem{FADD.f32} \dst{r0}, \src{r1}, \src{r2}} (\textbf{SPIR-V:} \texttt{N/A (encoding class)}; \textbf{GLSL:} \texttt{// Hardware execution pipe / decode class})
  \item \texttt{fau\_page\_index}: Selects the FAU (Fast Access Uniform) page (\texttt{0..3}) when a source operand references hardware small constants, thread IDs, or lane constants. \emph{Example:} {\scriptsize\ttfamily \mnem{FADD.f32} \dst{r0}, \src{r1}, \src{k1f}} (\textbf{SPIR-V:} \texttt{OpConstant / BuiltIn}; \textbf{GLSL:} \texttt{float y = x + 1.0;})
  \item \hyperref[enum:flow_control_m]{\texttt{flow}}: 4-bit embedded flow-control \& scoreboard wait modifier (\texttt{bits 59:62}: \texttt{none}, \texttt{wait0..wait0126}, \texttt{wait}, \texttt{reconverge}, \texttt{discard}, \texttt{end}). \emph{Example:} {\scriptsize\ttfamily \mnem{FADD.f32}\imod{.wait0.reconverge} \dst{r0}, \src{r1}, \src{r2}} (\textbf{SPIR-V:} \texttt{OpControlBarrier / OpReturn}; \textbf{GLSL:} \texttt{// Scoreboard wait + warp reconvergence})
\end{tightitemize}
\clearpage

\subsection*{FRCP\_APPROX}
\addcontentsline{toc}{subsection}{FRCP\_APPROX}
Low-precision reciprocal seed (\texttt{f32}): computes an initial \textasciitilde{}12-bit approximation of \texttt{1.0 / src0} in \texttt{dest}, intended as the starting estimate for Newton-Raphson refinement.
\par\textbf{Notes:}
\begin{tightitemize}
  \item SFU table seed instruction.
  \item \textbf{Source references:} \texttt{bifrost/valhall/ISA.xml}, \texttt{isa-v9-v14.xml}.
\end{tightitemize}
\codebox{\mnem{FRCP\_APPROX}\imod{\{.flow\}} \dst{dest}\imod{\{.dest\_width\}}, \src{src0}\imod{\{.widen0\}}\imod{\{.abs0\}}\imod{\{.neg0\}}}
\textbf{Operands:}
\begin{tightitemize}\footnotesize
  \item \texttt{\dst{dest}} (\texttt{dest\_register\_index}, \texttt{dest\_width}): FP32 destination GPR receiving the low-precision \texttt{1/x} seed for Newton-Raphson.
  \item \texttt{\src{src0}}: FP32 input \texttt{x}.
\end{tightitemize}
\par\textbf{Closest SPIR-V Op:} \texttt{None (SFU reciprocal seed)}
\par\textbf{Pseudo-GLSL (Operational Semantics):}
\codebox{dest = approxRcp12Bit(src0); // Newton-Raphson 1/x seed}
\textbf{Execution Context:}
\begin{tightitemize}\footnotesize
  \item \textbf{Compute stage preloads (\hyperref[sec:csf-spd]{\texttt{SPD\_k}}, \texttt{r16..r23}):} Hardware pre-fills \texttt{r55..r57} with \texttt{WorkGroupID.xyz + JOB\_OFFSET} (from \hyperref[sec:csf-regs]{\texttt{COMPUTESR::JOB\_OFFSET\_X/Y/Z}} \texttt{r34..r36}) and \texttt{r60..r62} with \texttt{LocalInvocationID.xyz} or \texttt{GlobalInvocationID.xyz} (from \texttt{WG\_SIZE} \texttt{r33} \& \texttt{JOB\_SIZE} \texttt{r37..r39}, selected by \texttt{ShaderProgram::preload}) so compute shaders can derive thread/element indices without memory loads.
  \item \textbf{Vertex stage preloads (\hyperref[sec:csf-spd]{\texttt{SPD\_0}/\texttt{SPD\_1}}, \texttt{r16..r19}):} The hardware Index Fetcher reads \hyperref[sec:csf-regs]{\texttt{IDVSSR::INDEX\_BUFFER}} (\texttt{r54:r55}) and pre-fills \texttt{r59, r60} with \texttt{VertexIndex} (\texttt{+ VERTEX\_OFFSET + GLOBAL\_ATTR\_OFF}) and \texttt{r61} with \texttt{InstanceIndex} (\texttt{+ INSTANCE\_OFFSET}) so Position and Varying shaders can pull vertex attributes directly.
  \item \textbf{Fragment stage preloads \& FP modes (\hyperref[sec:csf-spd]{\texttt{SPD\_2}}, \texttt{r20:r21}):} Hardware pre-fills \texttt{r59} with \texttt{floatBitsToUint(FragCoord.z)} and \texttt{r60} with the MSAA sample coverage mask (\texttt{gl\_SampleMaskIn[0]}), while \texttt{ShaderProgram} word~0 latches 32/64 GPR allocation and warp-wide \texttt{flush\_to\_zero\_mode}, \texttt{suppress\_na\_n}, and \texttt{suppress\_inf}.
  \item \textbf{Inline FAU \& Special Operands (\hyperref[sec:csf-tsd-fau]{\texttt{FAU\_k}}, \texttt{r8..r15}):} Any 8-bit source with \texttt{mode=2} reads zero-latency uniform \texttt{u0..u31} from \texttt{FAU\_k} (\texttt{slots64 = bits[63:56]}); \texttt{mode=3} reads \texttt{SmallConstant} (Table~1.2) or \texttt{FauSpecial} (\texttt{pc}, \texttt{lane\_id}, \texttt{lane\_width}, \texttt{tls\_ptr}, \texttt{wls\_ptr}, \texttt{blend\_descriptor\_0..7}).
\end{tightitemize}
\par\textbf{Worked Examples:}
\par\textbf{Example 1 (Initial Newton-Raphson seed for `1/4.0`):}
\codebox{\mnem{FRCP\_APPROX.f32} \dst{r0}, \src{r1}}
With \texttt{r1 = 4.0f}: returns fast low-precision estimate \texttt{\textasciitilde{}0.25f} used as \texttt{y0} in \texttt{y1 = y0 * (2 - x * y0)}.
\par\textbf{Example 2 (FP64 division lowering):}
\codebox{\mnem{FRCP\_APPROX.f32} \dst{r0}, \src{r1}}
Provides the \textasciitilde{}12-bit seed for emulated FP64 reciprocal before 3 Newton-Raphson iterations.

\begin{center}\footnotesize
\begin{tabularx}{\linewidth}{@{}lllX@{}}
\toprule
Variant & Arch & Class & Opcode Fields \\
\midrule
\texttt{(none)} & v9--v14 & \texttt{R1} & \texttt{opcode2=FRCP\_APPROX, opcode1=TABLE} \\
\bottomrule
\end{tabularx}
\end{center}
\begin{center}\footnotesize
\begin{tabularx}{0.88\linewidth}{@{}l l X@{}}
\toprule
Field Name & Bits & Type / Fixed Value \\
\midrule
\texttt{src0} & \texttt{0:7} & \texttt{SourceEncoding} \\
\textit{Reserved} & \texttt{8:15} & \texttt{zero} \\
\texttt{opcode2} & \texttt{16:20} & \texttt{exact=FRCP\_APPROX} \\
\textit{Reserved} & \texttt{21:27} & \texttt{zero} \\
\hyperref[enum:widen_m]{\texttt{widen0}} & \texttt{28:29} & \texttt{widen\_m (def none)} \\
\textit{Reserved} & \texttt{30:37} & \texttt{zero} \\
\hyperref[enum:neg_m]{\texttt{neg0}} & \texttt{38} & \texttt{neg\_m (def none)} \\
\hyperref[enum:abs_m]{\texttt{abs0}} & \texttt{39} & \texttt{abs\_m (def none)} \\
\texttt{dest\_register\_index} & \texttt{40:45} & \texttt{uint} \\
\hyperref[enum:dest_width_m]{\texttt{dest\_width}} & \texttt{46:47} & \texttt{dest\_width\_m} \\
\texttt{opcode1} & \texttt{48:51} & \texttt{exact=TABLE} \\
\texttt{instruction\_class} & \texttt{52:56} & \texttt{exact=R1} \\
\texttt{fau\_page\_index} & \texttt{57:58} & \texttt{uint} \\
\hyperref[enum:flow_control_m]{\texttt{flow}} & \texttt{59:62} & \texttt{flow\_control\_m} \\
\textit{Reserved} & \texttt{63} & \texttt{zero} \\
\bottomrule
\end{tabularx}
\end{center}
\textbf{Fields Explained:}
\begin{tightitemize}\footnotesize
  \item \texttt{src0}: Primary source operand (\texttt{r0..r63}, \texttt{u0..u63} uniform, \texttt{k*} FAU constant, or \texttt{t0..t2} temporary; \texttt{\textasciicircum{}} marks discard/last-use). \emph{Example:} {\scriptsize\ttfamily \mnem{FADD.f32} \dst{r0}, \src{r1}, \src{r2}} (\textbf{SPIR-V:} \texttt{OpFAdd (Operand 0)}; \textbf{GLSL:} \texttt{float c = a + b;})
  \item \texttt{opcode2}: Secondary sub-opcode further specializing instructions sharing \texttt{instruction\_class} and \texttt{opcode1} (e.g., conversions, SFU tables, moves). \emph{Example:} {\scriptsize\ttfamily \mnem{F32\_TO\_S32}\imod{.round\_zero} \dst{r0}, \src{r1}} (\textbf{SPIR-V:} \texttt{OpConvertFToS}; \textbf{GLSL:} \texttt{int y = int(x);})
  \item \hyperref[enum:widen_m]{\texttt{widen0}}: Stage A widening extractor on \texttt{src0} (promotes FP16 \texttt{h0}/\texttt{h1} to FP32, or sign/zero-extends \texttt{b0..b3}/\texttt{h0..h1}/\texttt{w0..w1} to 32/64-bit before \texttt{abs}/\texttt{neg}). \emph{Example:} {\scriptsize\ttfamily \mnem{IADD.s32} \dst{r0}, \src{r1}\imod{.h0}, \src{r2}} (\textbf{SPIR-V:} \texttt{OpSConvert / OpFConvert}; \textbf{GLSL:} \texttt{int c = int(int16\_t(a)) + b;})
  \item \hyperref[enum:neg_m]{\texttt{neg0}}: Stage A floating-point negation on \texttt{src0} (flips IEEE sign bit after swizzle/widen and \texttt{abs0}). \emph{Example:} {\scriptsize\ttfamily \mnem{FADD.f32} \dst{r0}, \src{r1}\imod{.neg}, \src{r2}} (\textbf{SPIR-V:} \texttt{OpFNegate}; \textbf{GLSL:} \texttt{float c = -a + b;})
  \item \hyperref[enum:abs_m]{\texttt{abs0}}: Stage A floating-point absolute value on \texttt{src0} (clears IEEE sign bit after swizzle/widen, before \texttt{neg0}). \emph{Example:} {\scriptsize\ttfamily \mnem{FADD.f32} \dst{r0}, \src{r1}\imod{.abs}, \src{r2}} (\textbf{SPIR-V:} \texttt{GLSL.std.450 FAbs}; \textbf{GLSL:} \texttt{float c = abs(a) + b;})
  \item \texttt{dest\_register\_index}: 6-bit destination GPR index (\texttt{r0..r63}) receiving the Stage C writeback result (even-aligned pair \texttt{[rN:rN+1]} for 64-bit ops). \emph{Example:} {\scriptsize\ttfamily \mnem{FMA.f32} \dst{r0}, \src{r1}, \src{r2}, \src{r3}} (\textbf{SPIR-V:} \texttt{SSA Result ID}; \textbf{GLSL:} \texttt{float r0 = fma(r1, r2, r3);})
  \item \hyperref[enum:dest_width_m]{\texttt{dest\_width}}: Stage C destination lane write-mask (\texttt{none} = full 32-bit, \texttt{h0} = low 16-bit, \texttt{h1} = high 16-bit without bit-shifting). \emph{Example:} {\scriptsize\ttfamily \mnem{FADD.v2f16} \dst{r0}\imod{.h1}, \src{r1}, \src{r2}} (\textbf{SPIR-V:} \texttt{OpVectorInsertDynamic / WriteMask}; \textbf{GLSL:} \texttt{dst.y = a.y + b.y;})
  \item \texttt{opcode1}: Primary sub-opcode within the instruction class uniquely identifying the operation family. \emph{Example:} {\scriptsize\ttfamily \mnem{IADD.s32} \dst{r0}, \src{r1}, \src{r2}} (\textbf{SPIR-V:} \texttt{OpIAdd (operation selector)}; \textbf{GLSL:} \texttt{// Primary opcode bits})
  \item \texttt{instruction\_class}: 4-bit major instruction class selector (\texttt{bits 48:51}) partitioning the 64-bit Valhall opcode space (e.g., \texttt{FMA}, \texttt{CVT}, \texttt{SFU}, \texttt{LS}, \texttt{TEX}, \texttt{BRANCH}). \emph{Example:} {\scriptsize\ttfamily \mnem{FADD.f32} \dst{r0}, \src{r1}, \src{r2}} (\textbf{SPIR-V:} \texttt{N/A (encoding class)}; \textbf{GLSL:} \texttt{// Hardware execution pipe / decode class})
  \item \texttt{fau\_page\_index}: Selects the FAU (Fast Access Uniform) page (\texttt{0..3}) when a source operand references hardware small constants, thread IDs, or lane constants. \emph{Example:} {\scriptsize\ttfamily \mnem{FADD.f32} \dst{r0}, \src{r1}, \src{k1f}} (\textbf{SPIR-V:} \texttt{OpConstant / BuiltIn}; \textbf{GLSL:} \texttt{float y = x + 1.0;})
  \item \hyperref[enum:flow_control_m]{\texttt{flow}}: 4-bit embedded flow-control \& scoreboard wait modifier (\texttt{bits 59:62}: \texttt{none}, \texttt{wait0..wait0126}, \texttt{wait}, \texttt{reconverge}, \texttt{discard}, \texttt{end}). \emph{Example:} {\scriptsize\ttfamily \mnem{FADD.f32}\imod{.wait0.reconverge} \dst{r0}, \src{r1}, \src{r2}} (\textbf{SPIR-V:} \texttt{OpControlBarrier / OpReturn}; \textbf{GLSL:} \texttt{// Scoreboard wait + warp reconvergence})
\end{tightitemize}
\clearpage

\subsection*{FREXPE}
\addcontentsline{toc}{subsection}{FREXPE}
Extracts the signed integer exponent $E$ of each floating-point element $i \in \{0, \dots, N-1\}$ of \texttt{src0} ($N=1$ \texttt{f32}, $N=2$ \texttt{v2f16}), returning \texttt{0} for \texttt{$\pm$0}, \texttt{$\pm$Inf}, and \texttt{NaN}: \texttt{none} gives \texttt{e - 126} (\texttt{[0.5, 1.0)}), \texttt{sqrt} gives \texttt{(e - 125) \textgreater{}\textgreater{} 1} (\texttt{[0.25, 1.0)}), and \texttt{log} gives \texttt{(e - 127) + top\_mant\_bit} (\texttt{[0.75, 1.5)}). Negates the result when \texttt{neg\_result=neg}.
\par\textbf{Notes:}
\begin{tightitemize}
  \item Bit-exact verified in \texttt{libGLES\_mali.so} \texttt{\_mali\_frexpe\_*} (\texttt{0x2e291d0..0x2e293c8}, \texttt{\_mali\_frexpe\_sf16} \texttt{0x2e294b0}).
  \item \texttt{v2f16} arch v9--v10 only.
\end{tightitemize}
\codebox{\mnem{FREXPE}\imod{\{.flow\}}\imod{\{.special\}}\imod{\{.neg\_result\}} \dst{dest}\imod{\{.dest\_width\}}, \src{src0}\imod{\{.widen0\}}\\\mnem{FREXPE}\imod{\{.flow\}}\imod{\{.special\}}\imod{\{.neg\_result\}} \dst{dest}\imod{\{.dest\_width\}}, \src{src0}\imod{\{.swz0\}}}
\textbf{Operands:}
\begin{tightitemize}\footnotesize
  \item \texttt{\dst{dest}} (\texttt{dest\_register\_index}, \texttt{dest\_width}): Signed integer destination GPR receiving the unbiased exponent \texttt{e} (adjusted for \texttt{.sqrt}/\texttt{.log}).
  \item \texttt{\src{src0}}: Floating-point input.
\end{tightitemize}
\par\textbf{Closest SPIR-V Op:} \texttt{GLSL.std.450 Frexp (exponent)}
\par\textbf{Pseudo-GLSL (Operational Semantics):}
\codebox{for (int i = 0; i \textless{} N; ++i) \{\\    int e; frexp(src0[i], e);\\    if (mode == sqrt) e = (e - 125) \textgreater{}\textgreater{} 1;\\    else if (mode == log) e = (e - 127) + (topMantBit(src0[i]) ? 1 : 0);\\    dest[i] = flag\_neg\_result ? -e : e;\\\}}
\textbf{Execution Context:}
\begin{tightitemize}\footnotesize
  \item \textbf{Compute stage preloads (\hyperref[sec:csf-spd]{\texttt{SPD\_k}}, \texttt{r16..r23}):} Hardware pre-fills \texttt{r55..r57} with \texttt{WorkGroupID.xyz + JOB\_OFFSET} (from \hyperref[sec:csf-regs]{\texttt{COMPUTESR::JOB\_OFFSET\_X/Y/Z}} \texttt{r34..r36}) and \texttt{r60..r62} with \texttt{LocalInvocationID.xyz} or \texttt{GlobalInvocationID.xyz} (from \texttt{WG\_SIZE} \texttt{r33} \& \texttt{JOB\_SIZE} \texttt{r37..r39}, selected by \texttt{ShaderProgram::preload}) so compute shaders can derive thread/element indices without memory loads.
  \item \textbf{Vertex stage preloads (\hyperref[sec:csf-spd]{\texttt{SPD\_0}/\texttt{SPD\_1}}, \texttt{r16..r19}):} The hardware Index Fetcher reads \hyperref[sec:csf-regs]{\texttt{IDVSSR::INDEX\_BUFFER}} (\texttt{r54:r55}) and pre-fills \texttt{r59, r60} with \texttt{VertexIndex} (\texttt{+ VERTEX\_OFFSET + GLOBAL\_ATTR\_OFF}) and \texttt{r61} with \texttt{InstanceIndex} (\texttt{+ INSTANCE\_OFFSET}) so Position and Varying shaders can pull vertex attributes directly.
  \item \textbf{Fragment stage preloads \& FP modes (\hyperref[sec:csf-spd]{\texttt{SPD\_2}}, \texttt{r20:r21}):} Hardware pre-fills \texttt{r59} with \texttt{floatBitsToUint(FragCoord.z)} and \texttt{r60} with the MSAA sample coverage mask (\texttt{gl\_SampleMaskIn[0]}), while \texttt{ShaderProgram} word~0 latches 32/64 GPR allocation and warp-wide \texttt{flush\_to\_zero\_mode}, \texttt{suppress\_na\_n}, and \texttt{suppress\_inf}.
  \item \textbf{Inline FAU \& Special Operands (\hyperref[sec:csf-tsd-fau]{\texttt{FAU\_k}}, \texttt{r8..r15}):} Any 8-bit source with \texttt{mode=2} reads zero-latency uniform \texttt{u0..u31} from \texttt{FAU\_k} (\texttt{slots64 = bits[63:56]}); \texttt{mode=3} reads \texttt{SmallConstant} (Table~1.2) or \texttt{FauSpecial} (\texttt{pc}, \texttt{lane\_id}, \texttt{lane\_width}, \texttt{tls\_ptr}, \texttt{wls\_ptr}, \texttt{blend\_descriptor\_0..7}).
\end{tightitemize}
\par\textbf{Worked Examples:}
\par\textbf{Example 1 (Standard vs `.log` exponent extraction):}
\codebox{\mnem{FREXPE.f32}\imod{.log} \dst{r0}, \src{r1}}
With \texttt{r1 = 3.0f} (\texttt{0x40400000}, biased \texttt{e = 128}, mantissa \texttt{1.5}): default returns \texttt{2} (\texttt{0.75 * 2\textasciicircum{}2}); \texttt{.log} returns \texttt{2} (\texttt{0.75 * 2\textasciicircum{}2} in \texttt{[0.75, 1.5)}).
\par\textbf{Example 2 (Square-root even exponent with `.neg`):}
\codebox{\mnem{FREXPE.f32}\imod{.sqrt.neg} \dst{r0}, \src{r1}}
With \texttt{r1 = 16.0f} (\texttt{2\textasciicircum{}4}): \texttt{.sqrt} halves exponent to \texttt{2} (\texttt{4\textasciicircum{}2}), and \texttt{.neg} negates -\textgreater{} \texttt{r0 = -2} (for \texttt{inversesqrt}).
\par\textbf{Mined SPIR-V $\to$ Valhall ISA Example (Extracting the .log-biased exponent for Log2 --- Fragment shader, Mali-G710 (v10)):}
\codebox{\textbf{; SPIR-V (input)}\\\%y = OpExtInst \%float \%glsl Log2 \%x\\[3pt]\textbf{; Valhall ISA (malioc output)}\\FREXPE.f32.log r4, r2\\S32\_TO\_F32 r4, r4\textasciicircum{}\\FLOGD.f32 r6, r2\\FADD\_LSCALE.f32 r2, k1f.neg, r2\textasciicircum{}\\FMA.f32 r2, r6\textasciicircum{}, r2\textasciicircum{}, r4\textasciicircum{}}
With \texttt{.log}, \texttt{FREXPE} increments the unbiased exponent by 1 whenever the mantissa is $\ge 1.5$ so that the reduced mantissa $m$ lies in $[0.75, 1.5)$ instead of $[1.0, 2.0)$, matching \texttt{FLOGD} and \texttt{FADD\_LSCALE}.
\begin{center}\footnotesize
\begin{tabularx}{\linewidth}{@{}lllX@{}}
\toprule
Variant & Arch & Class & Opcode Fields \\
\midrule
\texttt{(none)} & v9--v14 & \texttt{R1} & \texttt{opcode2=FREXPE\_F32, opcode1=FREXP} \\
\texttt{(none)} & v9--v14 & \texttt{R1} & \texttt{opcode2=FREXPE\_V2F16, opcode1=FREXP} \\
\bottomrule
\end{tabularx}
\end{center}
\begin{center}\footnotesize
\begin{tabularx}{0.88\linewidth}{@{}l l X@{}}
\toprule
Field Name & Bits & Type / Fixed Value \\
\midrule
\texttt{src0} & \texttt{0:7} & \texttt{SourceEncoding} \\
\textit{Reserved} & \texttt{8:15} & \texttt{zero} \\
\texttt{opcode2} & \texttt{16:20} & \texttt{exact=FREXPE\_F32} \\
\textit{Reserved} & \texttt{21:23} & \texttt{zero} \\
\hyperref[enum:frexp_special_m]{\texttt{special}} & \texttt{24:25} & \texttt{frexp\_special\_m (def none)} \\
\textit{Reserved} & \texttt{26:27} & \texttt{zero} \\
\hyperref[enum:widen_m]{\texttt{widen0}} & \texttt{28:29} & \texttt{widen\_m (def none)} \\
\hyperref[enum:neg_m]{\texttt{neg\_result}} & \texttt{30} & \texttt{neg\_m (def none)} \\
\textit{Reserved} & \texttt{31:39} & \texttt{zero} \\
\texttt{dest\_register\_index} & \texttt{40:45} & \texttt{uint} \\
\hyperref[enum:dest_width_m]{\texttt{dest\_width}} & \texttt{46:47} & \texttt{dest\_width\_m} \\
\texttt{opcode1} & \texttt{48:51} & \texttt{exact=FREXP} \\
\texttt{instruction\_class} & \texttt{52:56} & \texttt{exact=R1} \\
\texttt{fau\_page\_index} & \texttt{57:58} & \texttt{uint} \\
\hyperref[enum:flow_control_m]{\texttt{flow}} & \texttt{59:62} & \texttt{flow\_control\_m} \\
\textit{Reserved} & \texttt{63} & \texttt{zero} \\
\bottomrule
\end{tabularx}
\end{center}
\textbf{Fields Explained:}
\begin{tightitemize}\footnotesize
  \item \texttt{src0}: Primary source operand (\texttt{r0..r63}, \texttt{u0..u63} uniform, \texttt{k*} FAU constant, or \texttt{t0..t2} temporary; \texttt{\textasciicircum{}} marks discard/last-use). \emph{Example:} {\scriptsize\ttfamily \mnem{FADD.f32} \dst{r0}, \src{r1}, \src{r2}} (\textbf{SPIR-V:} \texttt{OpFAdd (Operand 0)}; \textbf{GLSL:} \texttt{float c = a + b;})
  \item \texttt{opcode2}: Secondary sub-opcode further specializing instructions sharing \texttt{instruction\_class} and \texttt{opcode1} (e.g., conversions, SFU tables, moves). \emph{Example:} {\scriptsize\ttfamily \mnem{F32\_TO\_S32}\imod{.round\_zero} \dst{r0}, \src{r1}} (\textbf{SPIR-V:} \texttt{OpConvertFToS}; \textbf{GLSL:} \texttt{int y = int(x);})
  \item \hyperref[enum:frexp_special_m]{\texttt{special}}: Selects \texttt{FMA\_RSCALE}/\texttt{FREXP} special-case override (\texttt{n}, \texttt{left}, \texttt{n\_add}, \texttt{sqrt}, \texttt{log}). \emph{Example:} {\scriptsize\ttfamily \mnem{FREXPM.f32}\imod{.log} \dst{r0}, \src{r1}} (\textbf{SPIR-V:} \texttt{GLSL.std.450 Log2 (range reduction)}; \textbf{GLSL:} \texttt{float y = log2(x);})
  \item \hyperref[enum:widen_m]{\texttt{widen0}}: Stage A widening extractor on \texttt{src0} (promotes FP16 \texttt{h0}/\texttt{h1} to FP32, or sign/zero-extends \texttt{b0..b3}/\texttt{h0..h1}/\texttt{w0..w1} to 32/64-bit before \texttt{abs}/\texttt{neg}). \emph{Example:} {\scriptsize\ttfamily \mnem{IADD.s32} \dst{r0}, \src{r1}\imod{.h0}, \src{r2}} (\textbf{SPIR-V:} \texttt{OpSConvert / OpFConvert}; \textbf{GLSL:} \texttt{int c = int(int16\_t(a)) + b;})
  \item \hyperref[enum:neg_m]{\texttt{neg\_result}}: Negates the final output of \texttt{FREXPE} (used for reciprocal-sqrt / inverse exponent scaling). \emph{Example:} {\scriptsize\ttfamily \mnem{FREXPE.f32}\imod{.sqrt.neg} \dst{r0}, \src{r1}} (\textbf{SPIR-V:} \texttt{GLSL.std.450 InverseSqrt}; \textbf{GLSL:} \texttt{float y = inversesqrt(x);})
  \item \texttt{dest\_register\_index}: 6-bit destination GPR index (\texttt{r0..r63}) receiving the Stage C writeback result (even-aligned pair \texttt{[rN:rN+1]} for 64-bit ops). \emph{Example:} {\scriptsize\ttfamily \mnem{FMA.f32} \dst{r0}, \src{r1}, \src{r2}, \src{r3}} (\textbf{SPIR-V:} \texttt{SSA Result ID}; \textbf{GLSL:} \texttt{float r0 = fma(r1, r2, r3);})
  \item \hyperref[enum:dest_width_m]{\texttt{dest\_width}}: Stage C destination lane write-mask (\texttt{none} = full 32-bit, \texttt{h0} = low 16-bit, \texttt{h1} = high 16-bit without bit-shifting). \emph{Example:} {\scriptsize\ttfamily \mnem{FADD.v2f16} \dst{r0}\imod{.h1}, \src{r1}, \src{r2}} (\textbf{SPIR-V:} \texttt{OpVectorInsertDynamic / WriteMask}; \textbf{GLSL:} \texttt{dst.y = a.y + b.y;})
  \item \texttt{opcode1}: Primary sub-opcode within the instruction class uniquely identifying the operation family. \emph{Example:} {\scriptsize\ttfamily \mnem{IADD.s32} \dst{r0}, \src{r1}, \src{r2}} (\textbf{SPIR-V:} \texttt{OpIAdd (operation selector)}; \textbf{GLSL:} \texttt{// Primary opcode bits})
  \item \texttt{instruction\_class}: 4-bit major instruction class selector (\texttt{bits 48:51}) partitioning the 64-bit Valhall opcode space (e.g., \texttt{FMA}, \texttt{CVT}, \texttt{SFU}, \texttt{LS}, \texttt{TEX}, \texttt{BRANCH}). \emph{Example:} {\scriptsize\ttfamily \mnem{FADD.f32} \dst{r0}, \src{r1}, \src{r2}} (\textbf{SPIR-V:} \texttt{N/A (encoding class)}; \textbf{GLSL:} \texttt{// Hardware execution pipe / decode class})
  \item \texttt{fau\_page\_index}: Selects the FAU (Fast Access Uniform) page (\texttt{0..3}) when a source operand references hardware small constants, thread IDs, or lane constants. \emph{Example:} {\scriptsize\ttfamily \mnem{FADD.f32} \dst{r0}, \src{r1}, \src{k1f}} (\textbf{SPIR-V:} \texttt{OpConstant / BuiltIn}; \textbf{GLSL:} \texttt{float y = x + 1.0;})
  \item \hyperref[enum:flow_control_m]{\texttt{flow}}: 4-bit embedded flow-control \& scoreboard wait modifier (\texttt{bits 59:62}: \texttt{none}, \texttt{wait0..wait0126}, \texttt{wait}, \texttt{reconverge}, \texttt{discard}, \texttt{end}). \emph{Example:} {\scriptsize\ttfamily \mnem{FADD.f32}\imod{.wait0.reconverge} \dst{r0}, \src{r1}, \src{r2}} (\textbf{SPIR-V:} \texttt{OpControlBarrier / OpReturn}; \textbf{GLSL:} \texttt{// Scoreboard wait + warp reconvergence})
\end{tightitemize}
\clearpage

\subsection*{FREXPM}
\addcontentsline{toc}{subsection}{FREXPM}
Extracts the normalized significand $M$ of each floating-point element $i \in \{0, \dots, N-1\}$ of \texttt{src0} ($N=1$ \texttt{f32}, $N=2$ \texttt{v2f16}): passes \texttt{$\pm$0} unchanged and quietens \texttt{NaN}; maps finite values into \texttt{[0.5, 1.0)} (\texttt{none}), \texttt{[0.25, 1.0)} (\texttt{sqrt}), or \texttt{[0.75, 1.5)} (\texttt{log}). On \texttt{$\pm$Inf}, \texttt{none}/\texttt{sqrt} return \texttt{$\pm$Inf} whereas \texttt{log} returns \texttt{$\pm$1.0} (\texttt{x \& 0xbf800000}).
\par\textbf{Notes:}
\begin{tightitemize}
  \item Bit-exact verified in \texttt{libGLES\_mali.so} \texttt{\_mali\_frexpm\_*} (\texttt{0x2e28fb0..0x2e2914c}, \texttt{\_mali\_frexpm\_sf16} \texttt{0x2e293d0}).
  \item \texttt{v2f16} arch v9--v10 only.
\end{tightitemize}
\codebox{\mnem{FREXPM}\imod{\{.flow\}}\imod{\{.special\}} \dst{dest}\imod{\{.dest\_width\}}, \src{src0}\imod{\{.widen0\}}\imod{\{.abs0\}}\imod{\{.neg0\}}\\\mnem{FREXPM}\imod{\{.flow\}}\imod{\{.special\}} \dst{dest}\imod{\{.dest\_width\}}, \src{src0}\imod{\{.swz0\}}\imod{\{.abs0\}}\imod{\{.neg0\}}}
\textbf{Operands:}
\begin{tightitemize}\footnotesize
  \item \texttt{\dst{dest}} (\texttt{dest\_register\_index}, \texttt{dest\_width}): Floating-point destination GPR receiving the normalized significand (\texttt{[0.5, 1.0)} or \texttt{[0.75, 1.5)} with \texttt{.log}).
  \item \texttt{\src{src0}}: Floating-point input.
\end{tightitemize}
\par\textbf{Closest SPIR-V Op:} \texttt{GLSL.std.450 Frexp (mantissa)}
\par\textbf{Pseudo-GLSL (Operational Semantics):}
\codebox{for (int i = 0; i \textless{} N; ++i) \{\\    if (mode == log) dest[i] = isinf(src0[i]) ? copysign(1.0, src0[i]) : reduceTo\_0\_75\_to\_1\_5(src0[i]);\\    else \{ int e; dest[i] = frexp(src0[i], e); \} // in [0.5, 1.0)\\\}}
\textbf{Execution Context:}
\begin{tightitemize}\footnotesize
  \item \textbf{Compute stage preloads (\hyperref[sec:csf-spd]{\texttt{SPD\_k}}, \texttt{r16..r23}):} Hardware pre-fills \texttt{r55..r57} with \texttt{WorkGroupID.xyz + JOB\_OFFSET} (from \hyperref[sec:csf-regs]{\texttt{COMPUTESR::JOB\_OFFSET\_X/Y/Z}} \texttt{r34..r36}) and \texttt{r60..r62} with \texttt{LocalInvocationID.xyz} or \texttt{GlobalInvocationID.xyz} (from \texttt{WG\_SIZE} \texttt{r33} \& \texttt{JOB\_SIZE} \texttt{r37..r39}, selected by \texttt{ShaderProgram::preload}) so compute shaders can derive thread/element indices without memory loads.
  \item \textbf{Vertex stage preloads (\hyperref[sec:csf-spd]{\texttt{SPD\_0}/\texttt{SPD\_1}}, \texttt{r16..r19}):} The hardware Index Fetcher reads \hyperref[sec:csf-regs]{\texttt{IDVSSR::INDEX\_BUFFER}} (\texttt{r54:r55}) and pre-fills \texttt{r59, r60} with \texttt{VertexIndex} (\texttt{+ VERTEX\_OFFSET + GLOBAL\_ATTR\_OFF}) and \texttt{r61} with \texttt{InstanceIndex} (\texttt{+ INSTANCE\_OFFSET}) so Position and Varying shaders can pull vertex attributes directly.
  \item \textbf{Fragment stage preloads \& FP modes (\hyperref[sec:csf-spd]{\texttt{SPD\_2}}, \texttt{r20:r21}):} Hardware pre-fills \texttt{r59} with \texttt{floatBitsToUint(FragCoord.z)} and \texttt{r60} with the MSAA sample coverage mask (\texttt{gl\_SampleMaskIn[0]}), while \texttt{ShaderProgram} word~0 latches 32/64 GPR allocation and warp-wide \texttt{flush\_to\_zero\_mode}, \texttt{suppress\_na\_n}, and \texttt{suppress\_inf}.
  \item \textbf{Inline FAU \& Special Operands (\hyperref[sec:csf-tsd-fau]{\texttt{FAU\_k}}, \texttt{r8..r15}):} Any 8-bit source with \texttt{mode=2} reads zero-latency uniform \texttt{u0..u31} from \texttt{FAU\_k} (\texttt{slots64 = bits[63:56]}); \texttt{mode=3} reads \texttt{SmallConstant} (Table~1.2) or \texttt{FauSpecial} (\texttt{pc}, \texttt{lane\_id}, \texttt{lane\_width}, \texttt{tls\_ptr}, \texttt{wls\_ptr}, \texttt{blend\_descriptor\_0..7}).
\end{tightitemize}
\par\textbf{Worked Examples:}
\par\textbf{Example 1 (Standard `[0.5, 1.0)` mantissa):}
\codebox{\mnem{FREXPM.f32} \dst{r0}, \src{r1}}
With \texttt{r1 = 3.0f}: returns \texttt{r0 = 0.75f} (\texttt{0x3F400000}), paired with \texttt{FREXPE = 2} since \texttt{0.75 * 2\textasciicircum{}2 = 3.0}.
\par\textbf{Example 2 (DDK `.log` infinity override):}
\codebox{\mnem{FREXPM.f32}\imod{.log} \dst{r0}, \src{r1}}
With \texttt{r1 = +Inf} (\texttt{0x7F800000}): default returns \texttt{+Inf}, whereas \texttt{.log} returns \texttt{+1.0f} (\texttt{0x3F800000}) and \texttt{-Inf -\textgreater{} -1.0f}.

\begin{center}\footnotesize
\begin{tabularx}{\linewidth}{@{}lllX@{}}
\toprule
Variant & Arch & Class & Opcode Fields \\
\midrule
\texttt{(none)} & v9--v14 & \texttt{R1} & \texttt{opcode2=FREXPM\_F32, opcode1=FREXP} \\
\texttt{(none)} & v9--v14 & \texttt{R1} & \texttt{opcode2=FREXPM\_V2F16, opcode1=FREXP} \\
\bottomrule
\end{tabularx}
\end{center}
\begin{center}\footnotesize
\begin{tabularx}{0.88\linewidth}{@{}l l X@{}}
\toprule
Field Name & Bits & Type / Fixed Value \\
\midrule
\texttt{src0} & \texttt{0:7} & \texttt{SourceEncoding} \\
\textit{Reserved} & \texttt{8:15} & \texttt{zero} \\
\texttt{opcode2} & \texttt{16:20} & \texttt{exact=FREXPM\_F32} \\
\textit{Reserved} & \texttt{21:23} & \texttt{zero} \\
\hyperref[enum:frexp_special_m]{\texttt{special}} & \texttt{24:25} & \texttt{frexp\_special\_m (def none)} \\
\textit{Reserved} & \texttt{26:27} & \texttt{zero} \\
\hyperref[enum:widen_m]{\texttt{widen0}} & \texttt{28:29} & \texttt{widen\_m (def none)} \\
\textit{Reserved} & \texttt{30:37} & \texttt{zero} \\
\hyperref[enum:neg_m]{\texttt{neg0}} & \texttt{38} & \texttt{neg\_m (def none)} \\
\hyperref[enum:abs_m]{\texttt{abs0}} & \texttt{39} & \texttt{abs\_m (def none)} \\
\texttt{dest\_register\_index} & \texttt{40:45} & \texttt{uint} \\
\hyperref[enum:dest_width_m]{\texttt{dest\_width}} & \texttt{46:47} & \texttt{dest\_width\_m} \\
\texttt{opcode1} & \texttt{48:51} & \texttt{exact=FREXP} \\
\texttt{instruction\_class} & \texttt{52:56} & \texttt{exact=R1} \\
\texttt{fau\_page\_index} & \texttt{57:58} & \texttt{uint} \\
\hyperref[enum:flow_control_m]{\texttt{flow}} & \texttt{59:62} & \texttt{flow\_control\_m} \\
\textit{Reserved} & \texttt{63} & \texttt{zero} \\
\bottomrule
\end{tabularx}
\end{center}
\textbf{Fields Explained:}
\begin{tightitemize}\footnotesize
  \item \texttt{src0}: Primary source operand (\texttt{r0..r63}, \texttt{u0..u63} uniform, \texttt{k*} FAU constant, or \texttt{t0..t2} temporary; \texttt{\textasciicircum{}} marks discard/last-use). \emph{Example:} {\scriptsize\ttfamily \mnem{FADD.f32} \dst{r0}, \src{r1}, \src{r2}} (\textbf{SPIR-V:} \texttt{OpFAdd (Operand 0)}; \textbf{GLSL:} \texttt{float c = a + b;})
  \item \texttt{opcode2}: Secondary sub-opcode further specializing instructions sharing \texttt{instruction\_class} and \texttt{opcode1} (e.g., conversions, SFU tables, moves). \emph{Example:} {\scriptsize\ttfamily \mnem{F32\_TO\_S32}\imod{.round\_zero} \dst{r0}, \src{r1}} (\textbf{SPIR-V:} \texttt{OpConvertFToS}; \textbf{GLSL:} \texttt{int y = int(x);})
  \item \hyperref[enum:frexp_special_m]{\texttt{special}}: Selects \texttt{FMA\_RSCALE}/\texttt{FREXP} special-case override (\texttt{n}, \texttt{left}, \texttt{n\_add}, \texttt{sqrt}, \texttt{log}). \emph{Example:} {\scriptsize\ttfamily \mnem{FREXPM.f32}\imod{.log} \dst{r0}, \src{r1}} (\textbf{SPIR-V:} \texttt{GLSL.std.450 Log2 (range reduction)}; \textbf{GLSL:} \texttt{float y = log2(x);})
  \item \hyperref[enum:widen_m]{\texttt{widen0}}: Stage A widening extractor on \texttt{src0} (promotes FP16 \texttt{h0}/\texttt{h1} to FP32, or sign/zero-extends \texttt{b0..b3}/\texttt{h0..h1}/\texttt{w0..w1} to 32/64-bit before \texttt{abs}/\texttt{neg}). \emph{Example:} {\scriptsize\ttfamily \mnem{IADD.s32} \dst{r0}, \src{r1}\imod{.h0}, \src{r2}} (\textbf{SPIR-V:} \texttt{OpSConvert / OpFConvert}; \textbf{GLSL:} \texttt{int c = int(int16\_t(a)) + b;})
  \item \hyperref[enum:neg_m]{\texttt{neg0}}: Stage A floating-point negation on \texttt{src0} (flips IEEE sign bit after swizzle/widen and \texttt{abs0}). \emph{Example:} {\scriptsize\ttfamily \mnem{FADD.f32} \dst{r0}, \src{r1}\imod{.neg}, \src{r2}} (\textbf{SPIR-V:} \texttt{OpFNegate}; \textbf{GLSL:} \texttt{float c = -a + b;})
  \item \hyperref[enum:abs_m]{\texttt{abs0}}: Stage A floating-point absolute value on \texttt{src0} (clears IEEE sign bit after swizzle/widen, before \texttt{neg0}). \emph{Example:} {\scriptsize\ttfamily \mnem{FADD.f32} \dst{r0}, \src{r1}\imod{.abs}, \src{r2}} (\textbf{SPIR-V:} \texttt{GLSL.std.450 FAbs}; \textbf{GLSL:} \texttt{float c = abs(a) + b;})
  \item \texttt{dest\_register\_index}: 6-bit destination GPR index (\texttt{r0..r63}) receiving the Stage C writeback result (even-aligned pair \texttt{[rN:rN+1]} for 64-bit ops). \emph{Example:} {\scriptsize\ttfamily \mnem{FMA.f32} \dst{r0}, \src{r1}, \src{r2}, \src{r3}} (\textbf{SPIR-V:} \texttt{SSA Result ID}; \textbf{GLSL:} \texttt{float r0 = fma(r1, r2, r3);})
  \item \hyperref[enum:dest_width_m]{\texttt{dest\_width}}: Stage C destination lane write-mask (\texttt{none} = full 32-bit, \texttt{h0} = low 16-bit, \texttt{h1} = high 16-bit without bit-shifting). \emph{Example:} {\scriptsize\ttfamily \mnem{FADD.v2f16} \dst{r0}\imod{.h1}, \src{r1}, \src{r2}} (\textbf{SPIR-V:} \texttt{OpVectorInsertDynamic / WriteMask}; \textbf{GLSL:} \texttt{dst.y = a.y + b.y;})
  \item \texttt{opcode1}: Primary sub-opcode within the instruction class uniquely identifying the operation family. \emph{Example:} {\scriptsize\ttfamily \mnem{IADD.s32} \dst{r0}, \src{r1}, \src{r2}} (\textbf{SPIR-V:} \texttt{OpIAdd (operation selector)}; \textbf{GLSL:} \texttt{// Primary opcode bits})
  \item \texttt{instruction\_class}: 4-bit major instruction class selector (\texttt{bits 48:51}) partitioning the 64-bit Valhall opcode space (e.g., \texttt{FMA}, \texttt{CVT}, \texttt{SFU}, \texttt{LS}, \texttt{TEX}, \texttt{BRANCH}). \emph{Example:} {\scriptsize\ttfamily \mnem{FADD.f32} \dst{r0}, \src{r1}, \src{r2}} (\textbf{SPIR-V:} \texttt{N/A (encoding class)}; \textbf{GLSL:} \texttt{// Hardware execution pipe / decode class})
  \item \texttt{fau\_page\_index}: Selects the FAU (Fast Access Uniform) page (\texttt{0..3}) when a source operand references hardware small constants, thread IDs, or lane constants. \emph{Example:} {\scriptsize\ttfamily \mnem{FADD.f32} \dst{r0}, \src{r1}, \src{k1f}} (\textbf{SPIR-V:} \texttt{OpConstant / BuiltIn}; \textbf{GLSL:} \texttt{float y = x + 1.0;})
  \item \hyperref[enum:flow_control_m]{\texttt{flow}}: 4-bit embedded flow-control \& scoreboard wait modifier (\texttt{bits 59:62}: \texttt{none}, \texttt{wait0..wait0126}, \texttt{wait}, \texttt{reconverge}, \texttt{discard}, \texttt{end}). \emph{Example:} {\scriptsize\ttfamily \mnem{FADD.f32}\imod{.wait0.reconverge} \dst{r0}, \src{r1}, \src{r2}} (\textbf{SPIR-V:} \texttt{OpControlBarrier / OpReturn}; \textbf{GLSL:} \texttt{// Scoreboard wait + warp reconvergence})
\end{tightitemize}
\clearpage

\subsection*{FRSQ\_APPROX}
\addcontentsline{toc}{subsection}{FRSQ\_APPROX}
Low-precision reciprocal square-root seed (\texttt{f32}): computes an initial \textasciitilde{}12-bit approximation of \texttt{1.0 / sqrt(src0)} in \texttt{dest}, intended as the starting estimate for Newton-Raphson refinement.
\par\textbf{Notes:}
\begin{tightitemize}
  \item SFU table seed instruction.
  \item \textbf{Source references:} \texttt{bifrost/valhall/ISA.xml}, \texttt{isa-v9-v14.xml}.
\end{tightitemize}
\codebox{\mnem{FRSQ\_APPROX}\imod{\{.flow\}} \dst{dest}\imod{\{.dest\_width\}}, \src{src0}\imod{\{.widen0\}}\imod{\{.abs0\}}\imod{\{.neg0\}}}
\textbf{Operands:}
\begin{tightitemize}\footnotesize
  \item \texttt{\dst{dest}} (\texttt{dest\_register\_index}, \texttt{dest\_width}): FP32 destination GPR receiving the low-precision \texttt{1/sqrt(x)} seed for Newton-Raphson.
  \item \texttt{\src{src0}}: FP32 input \texttt{x}.
\end{tightitemize}
\par\textbf{Closest SPIR-V Op:} \texttt{None (SFU rsqrt seed)}
\par\textbf{Pseudo-GLSL (Operational Semantics):}
\codebox{dest = approxRsqrt12Bit(src0); // Newton-Raphson 1/sqrt(x) seed}
\textbf{Execution Context:}
\begin{tightitemize}\footnotesize
  \item \textbf{Compute stage preloads (\hyperref[sec:csf-spd]{\texttt{SPD\_k}}, \texttt{r16..r23}):} Hardware pre-fills \texttt{r55..r57} with \texttt{WorkGroupID.xyz + JOB\_OFFSET} (from \hyperref[sec:csf-regs]{\texttt{COMPUTESR::JOB\_OFFSET\_X/Y/Z}} \texttt{r34..r36}) and \texttt{r60..r62} with \texttt{LocalInvocationID.xyz} or \texttt{GlobalInvocationID.xyz} (from \texttt{WG\_SIZE} \texttt{r33} \& \texttt{JOB\_SIZE} \texttt{r37..r39}, selected by \texttt{ShaderProgram::preload}) so compute shaders can derive thread/element indices without memory loads.
  \item \textbf{Vertex stage preloads (\hyperref[sec:csf-spd]{\texttt{SPD\_0}/\texttt{SPD\_1}}, \texttt{r16..r19}):} The hardware Index Fetcher reads \hyperref[sec:csf-regs]{\texttt{IDVSSR::INDEX\_BUFFER}} (\texttt{r54:r55}) and pre-fills \texttt{r59, r60} with \texttt{VertexIndex} (\texttt{+ VERTEX\_OFFSET + GLOBAL\_ATTR\_OFF}) and \texttt{r61} with \texttt{InstanceIndex} (\texttt{+ INSTANCE\_OFFSET}) so Position and Varying shaders can pull vertex attributes directly.
  \item \textbf{Fragment stage preloads \& FP modes (\hyperref[sec:csf-spd]{\texttt{SPD\_2}}, \texttt{r20:r21}):} Hardware pre-fills \texttt{r59} with \texttt{floatBitsToUint(FragCoord.z)} and \texttt{r60} with the MSAA sample coverage mask (\texttt{gl\_SampleMaskIn[0]}), while \texttt{ShaderProgram} word~0 latches 32/64 GPR allocation and warp-wide \texttt{flush\_to\_zero\_mode}, \texttt{suppress\_na\_n}, and \texttt{suppress\_inf}.
  \item \textbf{Inline FAU \& Special Operands (\hyperref[sec:csf-tsd-fau]{\texttt{FAU\_k}}, \texttt{r8..r15}):} Any 8-bit source with \texttt{mode=2} reads zero-latency uniform \texttt{u0..u31} from \texttt{FAU\_k} (\texttt{slots64 = bits[63:56]}); \texttt{mode=3} reads \texttt{SmallConstant} (Table~1.2) or \texttt{FauSpecial} (\texttt{pc}, \texttt{lane\_id}, \texttt{lane\_width}, \texttt{tls\_ptr}, \texttt{wls\_ptr}, \texttt{blend\_descriptor\_0..7}).
\end{tightitemize}
\par\textbf{Worked Examples:}
\par\textbf{Example 1 (Initial `1/sqrt(4.0)` seed):}
\codebox{\mnem{FRSQ\_APPROX.f32} \dst{r0}, \src{r1}}
With \texttt{r1 = 4.0f}: returns low-precision seed \texttt{\textasciitilde{}0.5f} for Newton-Raphson \texttt{y1 = y0 * (1.5 - 0.5 * x * y0\textasciicircum{}2)}.
\par\textbf{Example 2 (Vector normalization fast path):}
\codebox{\mnem{FRSQ\_APPROX.f32} \dst{r0}, \src{r1}}
Used as the initial estimate for software-refined FP64 \texttt{inversesqrt}.

\begin{center}\footnotesize
\begin{tabularx}{\linewidth}{@{}lllX@{}}
\toprule
Variant & Arch & Class & Opcode Fields \\
\midrule
\texttt{(none)} & v9--v14 & \texttt{R1} & \texttt{opcode2=FRSQ\_APPROX, opcode1=TABLE} \\
\bottomrule
\end{tabularx}
\end{center}
\begin{center}\footnotesize
\begin{tabularx}{0.88\linewidth}{@{}l l X@{}}
\toprule
Field Name & Bits & Type / Fixed Value \\
\midrule
\texttt{src0} & \texttt{0:7} & \texttt{SourceEncoding} \\
\textit{Reserved} & \texttt{8:15} & \texttt{zero} \\
\texttt{opcode2} & \texttt{16:20} & \texttt{exact=FRSQ\_APPROX} \\
\textit{Reserved} & \texttt{21:27} & \texttt{zero} \\
\hyperref[enum:widen_m]{\texttt{widen0}} & \texttt{28:29} & \texttt{widen\_m (def none)} \\
\textit{Reserved} & \texttt{30:37} & \texttt{zero} \\
\hyperref[enum:neg_m]{\texttt{neg0}} & \texttt{38} & \texttt{neg\_m (def none)} \\
\hyperref[enum:abs_m]{\texttt{abs0}} & \texttt{39} & \texttt{abs\_m (def none)} \\
\texttt{dest\_register\_index} & \texttt{40:45} & \texttt{uint} \\
\hyperref[enum:dest_width_m]{\texttt{dest\_width}} & \texttt{46:47} & \texttt{dest\_width\_m} \\
\texttt{opcode1} & \texttt{48:51} & \texttt{exact=TABLE} \\
\texttt{instruction\_class} & \texttt{52:56} & \texttt{exact=R1} \\
\texttt{fau\_page\_index} & \texttt{57:58} & \texttt{uint} \\
\hyperref[enum:flow_control_m]{\texttt{flow}} & \texttt{59:62} & \texttt{flow\_control\_m} \\
\textit{Reserved} & \texttt{63} & \texttt{zero} \\
\bottomrule
\end{tabularx}
\end{center}
\textbf{Fields Explained:}
\begin{tightitemize}\footnotesize
  \item \texttt{src0}: Primary source operand (\texttt{r0..r63}, \texttt{u0..u63} uniform, \texttt{k*} FAU constant, or \texttt{t0..t2} temporary; \texttt{\textasciicircum{}} marks discard/last-use). \emph{Example:} {\scriptsize\ttfamily \mnem{FADD.f32} \dst{r0}, \src{r1}, \src{r2}} (\textbf{SPIR-V:} \texttt{OpFAdd (Operand 0)}; \textbf{GLSL:} \texttt{float c = a + b;})
  \item \texttt{opcode2}: Secondary sub-opcode further specializing instructions sharing \texttt{instruction\_class} and \texttt{opcode1} (e.g., conversions, SFU tables, moves). \emph{Example:} {\scriptsize\ttfamily \mnem{F32\_TO\_S32}\imod{.round\_zero} \dst{r0}, \src{r1}} (\textbf{SPIR-V:} \texttt{OpConvertFToS}; \textbf{GLSL:} \texttt{int y = int(x);})
  \item \hyperref[enum:widen_m]{\texttt{widen0}}: Stage A widening extractor on \texttt{src0} (promotes FP16 \texttt{h0}/\texttt{h1} to FP32, or sign/zero-extends \texttt{b0..b3}/\texttt{h0..h1}/\texttt{w0..w1} to 32/64-bit before \texttt{abs}/\texttt{neg}). \emph{Example:} {\scriptsize\ttfamily \mnem{IADD.s32} \dst{r0}, \src{r1}\imod{.h0}, \src{r2}} (\textbf{SPIR-V:} \texttt{OpSConvert / OpFConvert}; \textbf{GLSL:} \texttt{int c = int(int16\_t(a)) + b;})
  \item \hyperref[enum:neg_m]{\texttt{neg0}}: Stage A floating-point negation on \texttt{src0} (flips IEEE sign bit after swizzle/widen and \texttt{abs0}). \emph{Example:} {\scriptsize\ttfamily \mnem{FADD.f32} \dst{r0}, \src{r1}\imod{.neg}, \src{r2}} (\textbf{SPIR-V:} \texttt{OpFNegate}; \textbf{GLSL:} \texttt{float c = -a + b;})
  \item \hyperref[enum:abs_m]{\texttt{abs0}}: Stage A floating-point absolute value on \texttt{src0} (clears IEEE sign bit after swizzle/widen, before \texttt{neg0}). \emph{Example:} {\scriptsize\ttfamily \mnem{FADD.f32} \dst{r0}, \src{r1}\imod{.abs}, \src{r2}} (\textbf{SPIR-V:} \texttt{GLSL.std.450 FAbs}; \textbf{GLSL:} \texttt{float c = abs(a) + b;})
  \item \texttt{dest\_register\_index}: 6-bit destination GPR index (\texttt{r0..r63}) receiving the Stage C writeback result (even-aligned pair \texttt{[rN:rN+1]} for 64-bit ops). \emph{Example:} {\scriptsize\ttfamily \mnem{FMA.f32} \dst{r0}, \src{r1}, \src{r2}, \src{r3}} (\textbf{SPIR-V:} \texttt{SSA Result ID}; \textbf{GLSL:} \texttt{float r0 = fma(r1, r2, r3);})
  \item \hyperref[enum:dest_width_m]{\texttt{dest\_width}}: Stage C destination lane write-mask (\texttt{none} = full 32-bit, \texttt{h0} = low 16-bit, \texttt{h1} = high 16-bit without bit-shifting). \emph{Example:} {\scriptsize\ttfamily \mnem{FADD.v2f16} \dst{r0}\imod{.h1}, \src{r1}, \src{r2}} (\textbf{SPIR-V:} \texttt{OpVectorInsertDynamic / WriteMask}; \textbf{GLSL:} \texttt{dst.y = a.y + b.y;})
  \item \texttt{opcode1}: Primary sub-opcode within the instruction class uniquely identifying the operation family. \emph{Example:} {\scriptsize\ttfamily \mnem{IADD.s32} \dst{r0}, \src{r1}, \src{r2}} (\textbf{SPIR-V:} \texttt{OpIAdd (operation selector)}; \textbf{GLSL:} \texttt{// Primary opcode bits})
  \item \texttt{instruction\_class}: 4-bit major instruction class selector (\texttt{bits 48:51}) partitioning the 64-bit Valhall opcode space (e.g., \texttt{FMA}, \texttt{CVT}, \texttt{SFU}, \texttt{LS}, \texttt{TEX}, \texttt{BRANCH}). \emph{Example:} {\scriptsize\ttfamily \mnem{FADD.f32} \dst{r0}, \src{r1}, \src{r2}} (\textbf{SPIR-V:} \texttt{N/A (encoding class)}; \textbf{GLSL:} \texttt{// Hardware execution pipe / decode class})
  \item \texttt{fau\_page\_index}: Selects the FAU (Fast Access Uniform) page (\texttt{0..3}) when a source operand references hardware small constants, thread IDs, or lane constants. \emph{Example:} {\scriptsize\ttfamily \mnem{FADD.f32} \dst{r0}, \src{r1}, \src{k1f}} (\textbf{SPIR-V:} \texttt{OpConstant / BuiltIn}; \textbf{GLSL:} \texttt{float y = x + 1.0;})
  \item \hyperref[enum:flow_control_m]{\texttt{flow}}: 4-bit embedded flow-control \& scoreboard wait modifier (\texttt{bits 59:62}: \texttt{none}, \texttt{wait0..wait0126}, \texttt{wait}, \texttt{reconverge}, \texttt{discard}, \texttt{end}). \emph{Example:} {\scriptsize\ttfamily \mnem{FADD.f32}\imod{.wait0.reconverge} \dst{r0}, \src{r1}, \src{r2}} (\textbf{SPIR-V:} \texttt{OpControlBarrier / OpReturn}; \textbf{GLSL:} \texttt{// Scoreboard wait + warp reconvergence})
\end{tightitemize}
\clearpage

\subsection*{FSINCOS\_OFFSET}
\addcontentsline{toc}{subsection}{FSINCOS\_OFFSET}
Trigonometric angle-offset lookup table (\texttt{u6}): indexes the 8-entry quadrant offset table using the low bits of \texttt{src0} and writes the FP32 offset to \texttt{dest} (scaled by $2^{24}$ when \texttt{.scale} is set).
\par\textbf{Notes:}
\begin{tightitemize}
  \item SFU table helper (\texttt{bifrost/valhall/ISA.xml:2340}, \texttt{compiler.h:691}).
  \item \textbf{Source references:} \texttt{bifrost/valhall/ISA.xml:2340}, \texttt{bifrost/compiler.h:691}, \texttt{isa-v9-v14.xml}.
\end{tightitemize}
\codebox{\mnem{FSINCOS\_OFFSET}\imod{\{.flow\}}\imod{\{.scale\}} \dst{dest}\imod{\{.dest\_width\}}, \src{src0}}
\textbf{Operands:}
\begin{tightitemize}\footnotesize
  \item \texttt{\dst{dest}} (\texttt{dest\_register\_index}, \texttt{dest\_width}): FP32 destination GPR receiving the trig quadrant angle offset (scaled by \texttt{2\textasciicircum{}24} if \texttt{.scale}).
  \item \texttt{\src{src0}}: Input whose low bits select the trig offset LUT entry.
\end{tightitemize}
\par\textbf{Closest SPIR-V Op:} \texttt{None (SFU trig offset LUT)}
\par\textbf{Pseudo-GLSL (Operational Semantics):}
\codebox{float off = SINCOS\_OFFSET\_LUT[src0 \& 0x7u];\\dest = (scale == scale\_24) ? ldexp(off, 24) : off;}
\textbf{Execution Context:}
\begin{tightitemize}\footnotesize
  \item \textbf{Compute stage preloads (\hyperref[sec:csf-spd]{\texttt{SPD\_k}}, \texttt{r16..r23}):} Hardware pre-fills \texttt{r55..r57} with \texttt{WorkGroupID.xyz + JOB\_OFFSET} (from \hyperref[sec:csf-regs]{\texttt{COMPUTESR::JOB\_OFFSET\_X/Y/Z}} \texttt{r34..r36}) and \texttt{r60..r62} with \texttt{LocalInvocationID.xyz} or \texttt{GlobalInvocationID.xyz} (from \texttt{WG\_SIZE} \texttt{r33} \& \texttt{JOB\_SIZE} \texttt{r37..r39}, selected by \texttt{ShaderProgram::preload}) so compute shaders can derive thread/element indices without memory loads.
  \item \textbf{Vertex stage preloads (\hyperref[sec:csf-spd]{\texttt{SPD\_0}/\texttt{SPD\_1}}, \texttt{r16..r19}):} The hardware Index Fetcher reads \hyperref[sec:csf-regs]{\texttt{IDVSSR::INDEX\_BUFFER}} (\texttt{r54:r55}) and pre-fills \texttt{r59, r60} with \texttt{VertexIndex} (\texttt{+ VERTEX\_OFFSET + GLOBAL\_ATTR\_OFF}) and \texttt{r61} with \texttt{InstanceIndex} (\texttt{+ INSTANCE\_OFFSET}) so Position and Varying shaders can pull vertex attributes directly.
  \item \textbf{Fragment stage preloads \& FP modes (\hyperref[sec:csf-spd]{\texttt{SPD\_2}}, \texttt{r20:r21}):} Hardware pre-fills \texttt{r59} with \texttt{floatBitsToUint(FragCoord.z)} and \texttt{r60} with the MSAA sample coverage mask (\texttt{gl\_SampleMaskIn[0]}), while \texttt{ShaderProgram} word~0 latches 32/64 GPR allocation and warp-wide \texttt{flush\_to\_zero\_mode}, \texttt{suppress\_na\_n}, and \texttt{suppress\_inf}.
  \item \textbf{Inline FAU \& Special Operands (\hyperref[sec:csf-tsd-fau]{\texttt{FAU\_k}}, \texttt{r8..r15}):} Any 8-bit source with \texttt{mode=2} reads zero-latency uniform \texttt{u0..u31} from \texttt{FAU\_k} (\texttt{slots64 = bits[63:56]}); \texttt{mode=3} reads \texttt{SmallConstant} (Table~1.2) or \texttt{FauSpecial} (\texttt{pc}, \texttt{lane\_id}, \texttt{lane\_width}, \texttt{tls\_ptr}, \texttt{wls\_ptr}, \texttt{blend\_descriptor\_0..7}).
\end{tightitemize}
\par\textbf{Worked Examples:}
\par\textbf{Example 1 (Scaled quadrant offset (.scale)):}
\codebox{\mnem{FSINCOS\_OFFSET.u6}\imod{.scale} \dst{r0}, \src{r1}}
Reads 8-entry quadrant-symmetric angle offset LUT for \texttt{r1 \& 0x3F}; \texttt{.scale} adds \texttt{+24} to the FP32 exponent.
\par\textbf{Example 2 (Unscaled base offset):}
\codebox{\mnem{FSINCOS\_OFFSET.u6} \dst{r0}, \src{r1}}
Returns the unscaled quadrant offset constant directly for \texttt{FMA} angle reconstruction.

\begin{center}\footnotesize
\begin{tabularx}{\linewidth}{@{}lllX@{}}
\toprule
Variant & Arch & Class & Opcode Fields \\
\midrule
\texttt{(none)} & v9--v14 & \texttt{R1} & \texttt{opcode2=FSINCOS\_OFFSET, opcode1=TABLE} \\
\bottomrule
\end{tabularx}
\end{center}
\begin{center}\footnotesize
\begin{tabularx}{0.88\linewidth}{@{}l l X@{}}
\toprule
Field Name & Bits & Type / Fixed Value \\
\midrule
\texttt{src0} & \texttt{0:7} & \texttt{SourceEncoding} \\
\textit{Reserved} & \texttt{8:15} & \texttt{zero} \\
\texttt{opcode2} & \texttt{16:20} & \texttt{exact=FSINCOS\_OFFSET} \\
\textit{Reserved} & \texttt{21:29} & \texttt{zero} \\
\hyperref[enum:table_sincos_scale_m]{\texttt{scale}} & \texttt{30} & \texttt{table\_sincos\_scale\_m} \\
\textit{Reserved} & \texttt{31:39} & \texttt{zero} \\
\texttt{dest\_register\_index} & \texttt{40:45} & \texttt{uint} \\
\hyperref[enum:dest_width_m]{\texttt{dest\_width}} & \texttt{46:47} & \texttt{dest\_width\_m} \\
\texttt{opcode1} & \texttt{48:51} & \texttt{exact=TABLE} \\
\texttt{instruction\_class} & \texttt{52:56} & \texttt{exact=R1} \\
\texttt{fau\_page\_index} & \texttt{57:58} & \texttt{uint} \\
\hyperref[enum:flow_control_m]{\texttt{flow}} & \texttt{59:62} & \texttt{flow\_control\_m} \\
\textit{Reserved} & \texttt{63} & \texttt{zero} \\
\bottomrule
\end{tabularx}
\end{center}
\textbf{Fields Explained:}
\begin{tightitemize}\footnotesize
  \item \texttt{src0}: Primary source operand (\texttt{r0..r63}, \texttt{u0..u63} uniform, \texttt{k*} FAU constant, or \texttt{t0..t2} temporary; \texttt{\textasciicircum{}} marks discard/last-use). \emph{Example:} {\scriptsize\ttfamily \mnem{FADD.f32} \dst{r0}, \src{r1}, \src{r2}} (\textbf{SPIR-V:} \texttt{OpFAdd (Operand 0)}; \textbf{GLSL:} \texttt{float c = a + b;})
  \item \texttt{opcode2}: Secondary sub-opcode further specializing instructions sharing \texttt{instruction\_class} and \texttt{opcode1} (e.g., conversions, SFU tables, moves). \emph{Example:} {\scriptsize\ttfamily \mnem{F32\_TO\_S32}\imod{.round\_zero} \dst{r0}, \src{r1}} (\textbf{SPIR-V:} \texttt{OpConvertFToS}; \textbf{GLSL:} \texttt{int y = int(x);})
  \item \hyperref[enum:table_sincos_scale_m]{\texttt{scale}}: Enables $\pi$-scaling in \texttt{FSINCOS\_OFFSET}. \emph{Example:} {\scriptsize\ttfamily \mnem{FSINCOS\_OFFSET}\imod{.scale} \dst{r0}, \src{r1}} (\textbf{SPIR-V:} \texttt{GLSL.std.450 Sin}; \textbf{GLSL:} \texttt{float y = sin(x);})
  \item \texttt{dest\_register\_index}: 6-bit destination GPR index (\texttt{r0..r63}) receiving the Stage C writeback result (even-aligned pair \texttt{[rN:rN+1]} for 64-bit ops). \emph{Example:} {\scriptsize\ttfamily \mnem{FMA.f32} \dst{r0}, \src{r1}, \src{r2}, \src{r3}} (\textbf{SPIR-V:} \texttt{SSA Result ID}; \textbf{GLSL:} \texttt{float r0 = fma(r1, r2, r3);})
  \item \hyperref[enum:dest_width_m]{\texttt{dest\_width}}: Stage C destination lane write-mask (\texttt{none} = full 32-bit, \texttt{h0} = low 16-bit, \texttt{h1} = high 16-bit without bit-shifting). \emph{Example:} {\scriptsize\ttfamily \mnem{FADD.v2f16} \dst{r0}\imod{.h1}, \src{r1}, \src{r2}} (\textbf{SPIR-V:} \texttt{OpVectorInsertDynamic / WriteMask}; \textbf{GLSL:} \texttt{dst.y = a.y + b.y;})
  \item \texttt{opcode1}: Primary sub-opcode within the instruction class uniquely identifying the operation family. \emph{Example:} {\scriptsize\ttfamily \mnem{IADD.s32} \dst{r0}, \src{r1}, \src{r2}} (\textbf{SPIR-V:} \texttt{OpIAdd (operation selector)}; \textbf{GLSL:} \texttt{// Primary opcode bits})
  \item \texttt{instruction\_class}: 4-bit major instruction class selector (\texttt{bits 48:51}) partitioning the 64-bit Valhall opcode space (e.g., \texttt{FMA}, \texttt{CVT}, \texttt{SFU}, \texttt{LS}, \texttt{TEX}, \texttt{BRANCH}). \emph{Example:} {\scriptsize\ttfamily \mnem{FADD.f32} \dst{r0}, \src{r1}, \src{r2}} (\textbf{SPIR-V:} \texttt{N/A (encoding class)}; \textbf{GLSL:} \texttt{// Hardware execution pipe / decode class})
  \item \texttt{fau\_page\_index}: Selects the FAU (Fast Access Uniform) page (\texttt{0..3}) when a source operand references hardware small constants, thread IDs, or lane constants. \emph{Example:} {\scriptsize\ttfamily \mnem{FADD.f32} \dst{r0}, \src{r1}, \src{k1f}} (\textbf{SPIR-V:} \texttt{OpConstant / BuiltIn}; \textbf{GLSL:} \texttt{float y = x + 1.0;})
  \item \hyperref[enum:flow_control_m]{\texttt{flow}}: 4-bit embedded flow-control \& scoreboard wait modifier (\texttt{bits 59:62}: \texttt{none}, \texttt{wait0..wait0126}, \texttt{wait}, \texttt{reconverge}, \texttt{discard}, \texttt{end}). \emph{Example:} {\scriptsize\ttfamily \mnem{FADD.f32}\imod{.wait0.reconverge} \dst{r0}, \src{r1}, \src{r2}} (\textbf{SPIR-V:} \texttt{OpControlBarrier / OpReturn}; \textbf{GLSL:} \texttt{// Scoreboard wait + warp reconvergence})
\end{tightitemize}
\clearpage

\subsection*{FSIN\_TABLE}
\addcontentsline{toc}{subsection}{FSIN\_TABLE}
64-entry sine lookup table for trigonometric range reduction: if \texttt{src0} is NaN returns \texttt{src0}; if \texttt{src0} is Inf returns \texttt{NaN}; if \texttt{(sign, exp) != (0, 146)} returns \texttt{-0.0}; otherwise indexes \texttt{idx = bits(src0) \& 0x3F} and writes \texttt{sin(idx * pi / 32)} to \texttt{dest} (with exact \texttt{-0.0} at \texttt{idx = 0} and \texttt{+0.0} at \texttt{idx = 32}).
\par\textbf{Notes:}
\begin{tightitemize}
  \item Exact hardware match verified in \texttt{hw\_tests.rs:1748-1762} (with \texttt{offset=false}).
  \item \textbf{Source references:} \texttt{kraid/ops.rs:2107-2147}, \texttt{encode\_v9.rs:1690-1708}, \texttt{nir.rs:1988-1995}, \texttt{hw\_tests.rs:1748-1762}.
\end{tightitemize}
\codebox{\mnem{FSIN\_TABLE}\imod{\{.flow\}}\imod{\{.offset\}} \dst{dest}\imod{\{.dest\_width\}}, \src{src0}}
\textbf{Operands:}
\begin{tightitemize}\footnotesize
  \item \texttt{\dst{dest}} (\texttt{dest\_register\_index}, \texttt{dest\_width}): FP32 destination GPR receiving \texttt{sin(idx * pi / 32)}.
  \item \texttt{\src{src0}}: Fixed-point angle word whose low 6 bits select the 64-entry sine LUT entry.
\end{tightitemize}
\par\textbf{Closest SPIR-V Op:} \texttt{None (lowered from GLSL.std.450 Sin/Cos)}
\par\textbf{Pseudo-GLSL (Operational Semantics):}
\codebox{if (isnan(src0))      dest = src0;\\else if (isinf(src0)) dest = NaN;\\else if (signExp(src0) != (0, 146)) dest = -0.0;\\else \{ uint idx = floatBitsToUint(src0) \& 0x3Fu; dest = (idx == 0u) ? -0.0 : sin(float(idx) * PI / 32.0); \}}
\textbf{Execution Context:}
\begin{tightitemize}\footnotesize
  \item \textbf{Compute stage preloads (\hyperref[sec:csf-spd]{\texttt{SPD\_k}}, \texttt{r16..r23}):} Hardware pre-fills \texttt{r55..r57} with \texttt{WorkGroupID.xyz + JOB\_OFFSET} (from \hyperref[sec:csf-regs]{\texttt{COMPUTESR::JOB\_OFFSET\_X/Y/Z}} \texttt{r34..r36}) and \texttt{r60..r62} with \texttt{LocalInvocationID.xyz} or \texttt{GlobalInvocationID.xyz} (from \texttt{WG\_SIZE} \texttt{r33} \& \texttt{JOB\_SIZE} \texttt{r37..r39}, selected by \texttt{ShaderProgram::preload}) so compute shaders can derive thread/element indices without memory loads.
  \item \textbf{Vertex stage preloads (\hyperref[sec:csf-spd]{\texttt{SPD\_0}/\texttt{SPD\_1}}, \texttt{r16..r19}):} The hardware Index Fetcher reads \hyperref[sec:csf-regs]{\texttt{IDVSSR::INDEX\_BUFFER}} (\texttt{r54:r55}) and pre-fills \texttt{r59, r60} with \texttt{VertexIndex} (\texttt{+ VERTEX\_OFFSET + GLOBAL\_ATTR\_OFF}) and \texttt{r61} with \texttt{InstanceIndex} (\texttt{+ INSTANCE\_OFFSET}) so Position and Varying shaders can pull vertex attributes directly.
  \item \textbf{Fragment stage preloads \& FP modes (\hyperref[sec:csf-spd]{\texttt{SPD\_2}}, \texttt{r20:r21}):} Hardware pre-fills \texttt{r59} with \texttt{floatBitsToUint(FragCoord.z)} and \texttt{r60} with the MSAA sample coverage mask (\texttt{gl\_SampleMaskIn[0]}), while \texttt{ShaderProgram} word~0 latches 32/64 GPR allocation and warp-wide \texttt{flush\_to\_zero\_mode}, \texttt{suppress\_na\_n}, and \texttt{suppress\_inf}.
  \item \textbf{Inline FAU \& Special Operands (\hyperref[sec:csf-tsd-fau]{\texttt{FAU\_k}}, \texttt{r8..r15}):} Any 8-bit source with \texttt{mode=2} reads zero-latency uniform \texttt{u0..u31} from \texttt{FAU\_k} (\texttt{slots64 = bits[63:56]}); \texttt{mode=3} reads \texttt{SmallConstant} (Table~1.2) or \texttt{FauSpecial} (\texttt{pc}, \texttt{lane\_id}, \texttt{lane\_width}, \texttt{tls\_ptr}, \texttt{wls\_ptr}, \texttt{blend\_descriptor\_0..7}).
\end{tightitemize}
\par\textbf{Worked Examples:}
\par\textbf{Example 1 (Index 16 (`sin(pi/2) = +1.0`)):}
\codebox{\mnem{FSIN\_TABLE.u6} \dst{r0}, \src{r1}}
With \texttt{r1 = 0x49400010} (\texttt{idx = 16}, exp \texttt{127+19}): returns \texttt{sin(16*pi/32) = +1.0f} (\texttt{0x3F800000}).
\par\textbf{Example 2 (Index 0 signed-zero edge case):}
\codebox{\mnem{FSIN\_TABLE.u6} \dst{r0}, \src{r1}}
With \texttt{r1 = 0x49400000} (\texttt{idx = 0}): returns \texttt{-0.0f} (\texttt{0x80000000}) so \texttt{FMA} preserves \texttt{sin(-0.0) = -0.0}.
\par\textbf{Mined SPIR-V $\to$ Valhall ISA Example (Sine table lookup paired with FCOS\_TABLE in trig lowering --- Vertex shader, Mali-G710 (v10)):}
\codebox{\textbf{; SPIR-V (input)}\\\%c = OpExtInst \%float \%glsl Cos \%angle\\[3pt]\textbf{; Valhall ISA (malioc output)}\\FSIN\_TABLE.u6 r63, r5             ; sin(idx * pi / 32)\\FCOS\_TABLE.u6 r5, r5\textasciicircum{}             ; cos(idx * pi / 32)\\FMUL.f32 r7, r7\textasciicircum{}, r5.neg\\FMA.f32.clamp\_m1\_1 r2, r2\textasciicircum{}, r63\textasciicircum{}.neg, r7\textasciicircum{}}
\texttt{FSIN\_TABLE.u6} reads the 6-bit index in \texttt{r5} (low 6 bits of the range-reduced fixed-point angle) and returns $\sin(\text{idx} \cdot \pi / 32)$ for the linear Taylor term $-\delta \sin(\theta_0)$.
\begin{center}\footnotesize
\begin{tabularx}{\linewidth}{@{}lllX@{}}
\toprule
Variant & Arch & Class & Opcode Fields \\
\midrule
\texttt{(none)} & v9--v14 & \texttt{R1} & \texttt{opcode2=FSIN\_TABLE, opcode1=TABLE} \\
\bottomrule
\end{tabularx}
\end{center}
\begin{center}\footnotesize
\begin{tabularx}{0.88\linewidth}{@{}l l X@{}}
\toprule
Field Name & Bits & Type / Fixed Value \\
\midrule
\texttt{src0} & \texttt{0:7} & \texttt{SourceEncoding} \\
\textit{Reserved} & \texttt{8:15} & \texttt{zero} \\
\texttt{opcode2} & \texttt{16:20} & \texttt{exact=FSIN\_TABLE} \\
\textit{Reserved} & \texttt{21:29} & \texttt{zero} \\
\hyperref[enum:table_sincos_offset_m]{\texttt{offset}} & \texttt{30} & \texttt{table\_sincos\_offset\_m} \\
\textit{Reserved} & \texttt{31:39} & \texttt{zero} \\
\texttt{dest\_register\_index} & \texttt{40:45} & \texttt{uint} \\
\hyperref[enum:dest_width_m]{\texttt{dest\_width}} & \texttt{46:47} & \texttt{dest\_width\_m} \\
\texttt{opcode1} & \texttt{48:51} & \texttt{exact=TABLE} \\
\texttt{instruction\_class} & \texttt{52:56} & \texttt{exact=R1} \\
\texttt{fau\_page\_index} & \texttt{57:58} & \texttt{uint} \\
\hyperref[enum:flow_control_m]{\texttt{flow}} & \texttt{59:62} & \texttt{flow\_control\_m} \\
\textit{Reserved} & \texttt{63} & \texttt{zero} \\
\bottomrule
\end{tabularx}
\end{center}
\textbf{Fields Explained:}
\begin{tightitemize}\footnotesize
  \item \texttt{src0}: Primary source operand (\texttt{r0..r63}, \texttt{u0..u63} uniform, \texttt{k*} FAU constant, or \texttt{t0..t2} temporary; \texttt{\textasciicircum{}} marks discard/last-use). \emph{Example:} {\scriptsize\ttfamily \mnem{FADD.f32} \dst{r0}, \src{r1}, \src{r2}} (\textbf{SPIR-V:} \texttt{OpFAdd (Operand 0)}; \textbf{GLSL:} \texttt{float c = a + b;})
  \item \texttt{opcode2}: Secondary sub-opcode further specializing instructions sharing \texttt{instruction\_class} and \texttt{opcode1} (e.g., conversions, SFU tables, moves). \emph{Example:} {\scriptsize\ttfamily \mnem{F32\_TO\_S32}\imod{.round\_zero} \dst{r0}, \src{r1}} (\textbf{SPIR-V:} \texttt{OpConvertFToS}; \textbf{GLSL:} \texttt{int y = int(x);})
  \item \hyperref[enum:table_sincos_offset_m]{\texttt{offset}}: Signed byte offset added to base address register in memory load/store/atomic instructions. \emph{Example:} {\scriptsize\ttfamily \mnem{LOAD.i32} \dst{@r0}, \src{[r2:r3]}, \imm{\#16}} (\textbf{SPIR-V:} \texttt{OpAccessChain + OpLoad}; \textbf{GLSL:} \texttt{uint y = ssbo.arr[4];})
  \item \texttt{dest\_register\_index}: 6-bit destination GPR index (\texttt{r0..r63}) receiving the Stage C writeback result (even-aligned pair \texttt{[rN:rN+1]} for 64-bit ops). \emph{Example:} {\scriptsize\ttfamily \mnem{FMA.f32} \dst{r0}, \src{r1}, \src{r2}, \src{r3}} (\textbf{SPIR-V:} \texttt{SSA Result ID}; \textbf{GLSL:} \texttt{float r0 = fma(r1, r2, r3);})
  \item \hyperref[enum:dest_width_m]{\texttt{dest\_width}}: Stage C destination lane write-mask (\texttt{none} = full 32-bit, \texttt{h0} = low 16-bit, \texttt{h1} = high 16-bit without bit-shifting). \emph{Example:} {\scriptsize\ttfamily \mnem{FADD.v2f16} \dst{r0}\imod{.h1}, \src{r1}, \src{r2}} (\textbf{SPIR-V:} \texttt{OpVectorInsertDynamic / WriteMask}; \textbf{GLSL:} \texttt{dst.y = a.y + b.y;})
  \item \texttt{opcode1}: Primary sub-opcode within the instruction class uniquely identifying the operation family. \emph{Example:} {\scriptsize\ttfamily \mnem{IADD.s32} \dst{r0}, \src{r1}, \src{r2}} (\textbf{SPIR-V:} \texttt{OpIAdd (operation selector)}; \textbf{GLSL:} \texttt{// Primary opcode bits})
  \item \texttt{instruction\_class}: 4-bit major instruction class selector (\texttt{bits 48:51}) partitioning the 64-bit Valhall opcode space (e.g., \texttt{FMA}, \texttt{CVT}, \texttt{SFU}, \texttt{LS}, \texttt{TEX}, \texttt{BRANCH}). \emph{Example:} {\scriptsize\ttfamily \mnem{FADD.f32} \dst{r0}, \src{r1}, \src{r2}} (\textbf{SPIR-V:} \texttt{N/A (encoding class)}; \textbf{GLSL:} \texttt{// Hardware execution pipe / decode class})
  \item \texttt{fau\_page\_index}: Selects the FAU (Fast Access Uniform) page (\texttt{0..3}) when a source operand references hardware small constants, thread IDs, or lane constants. \emph{Example:} {\scriptsize\ttfamily \mnem{FADD.f32} \dst{r0}, \src{r1}, \src{k1f}} (\textbf{SPIR-V:} \texttt{OpConstant / BuiltIn}; \textbf{GLSL:} \texttt{float y = x + 1.0;})
  \item \hyperref[enum:flow_control_m]{\texttt{flow}}: 4-bit embedded flow-control \& scoreboard wait modifier (\texttt{bits 59:62}: \texttt{none}, \texttt{wait0..wait0126}, \texttt{wait}, \texttt{reconverge}, \texttt{discard}, \texttt{end}). \emph{Example:} {\scriptsize\ttfamily \mnem{FADD.f32}\imod{.wait0.reconverge} \dst{r0}, \src{r1}, \src{r2}} (\textbf{SPIR-V:} \texttt{OpControlBarrier / OpReturn}; \textbf{GLSL:} \texttt{// Scoreboard wait + warp reconvergence})
\end{tightitemize}
\clearpage

\subsection*{ILOGB}
\addcontentsline{toc}{subsection}{ILOGB}
Extracts the unbiased base-2 integer exponent \texttt{floor(log2(\textbar{}src0[i]\textbar{}))} for each floating-point element $i \in \{0, \dots, N-1\}$ ($N=1$ \texttt{f32 -\textgreater{} s32}, $N=2$ \texttt{v2f16 -\textgreater{} v2s16}; arch \textless{}= v10).
\par\textbf{Notes:}
\begin{tightitemize}
  \item IEEE-754 \texttt{ilogb} special cases: \texttt{0 -\textgreater{} FP\_ILOGB0}, \texttt{Inf -\textgreater{} INT\_MAX}, \texttt{NaN -\textgreater{} FP\_ILOGBNAN}.
  \item \textbf{Source references:} \texttt{bifrost/bifrost/ISA.xml:5778-5788}, \texttt{isa-v9-v14.xml}.
\end{tightitemize}
\codebox{\mnem{ILOGB}\imod{\{.flow\}} \dst{dest}\imod{\{.dest\_width\}}, \src{src0}\imod{\{.widen0\}}\\\mnem{ILOGB}\imod{\{.flow\}} \dst{dest}\imod{\{.dest\_width\}}, \src{src0}\imod{\{.swz0\}}}
\textbf{Operands:}
\begin{tightitemize}\footnotesize
  \item \texttt{\dst{dest}} (\texttt{dest\_register\_index}, \texttt{dest\_width}): Signed integer destination GPR receiving \texttt{floor(log2(\textbar{}src0\textbar{}))} (\texttt{INT32\_MIN} on 0, \texttt{INT32\_MAX} on Inf).
  \item \texttt{\src{src0}}: Floating-point input.
\end{tightitemize}
\par\textbf{Closest SPIR-V Op:} \texttt{None (OpenCL `ilogb`, v9--v10)}
\par\textbf{Pseudo-GLSL (Operational Semantics):}
\codebox{for (int i = 0; i \textless{} N; ++i) \{\\    if (src0[i] == 0.0)    dest[i] = INT32\_MIN;\\    else if (isinf(src0[i])) dest[i] = INT32\_MAX;\\    else                   dest[i] = int(floor(log2(abs(src0[i]))));\\\}}
\textbf{Execution Context:}
\begin{tightitemize}\footnotesize
  \item \textbf{Compute stage preloads (\hyperref[sec:csf-spd]{\texttt{SPD\_k}}, \texttt{r16..r23}):} Hardware pre-fills \texttt{r55..r57} with \texttt{WorkGroupID.xyz + JOB\_OFFSET} (from \hyperref[sec:csf-regs]{\texttt{COMPUTESR::JOB\_OFFSET\_X/Y/Z}} \texttt{r34..r36}) and \texttt{r60..r62} with \texttt{LocalInvocationID.xyz} or \texttt{GlobalInvocationID.xyz} (from \texttt{WG\_SIZE} \texttt{r33} \& \texttt{JOB\_SIZE} \texttt{r37..r39}, selected by \texttt{ShaderProgram::preload}) so compute shaders can derive thread/element indices without memory loads.
  \item \textbf{Vertex stage preloads (\hyperref[sec:csf-spd]{\texttt{SPD\_0}/\texttt{SPD\_1}}, \texttt{r16..r19}):} The hardware Index Fetcher reads \hyperref[sec:csf-regs]{\texttt{IDVSSR::INDEX\_BUFFER}} (\texttt{r54:r55}) and pre-fills \texttt{r59, r60} with \texttt{VertexIndex} (\texttt{+ VERTEX\_OFFSET + GLOBAL\_ATTR\_OFF}) and \texttt{r61} with \texttt{InstanceIndex} (\texttt{+ INSTANCE\_OFFSET}) so Position and Varying shaders can pull vertex attributes directly.
  \item \textbf{Fragment stage preloads \& FP modes (\hyperref[sec:csf-spd]{\texttt{SPD\_2}}, \texttt{r20:r21}):} Hardware pre-fills \texttt{r59} with \texttt{floatBitsToUint(FragCoord.z)} and \texttt{r60} with the MSAA sample coverage mask (\texttt{gl\_SampleMaskIn[0]}), while \texttt{ShaderProgram} word~0 latches 32/64 GPR allocation and warp-wide \texttt{flush\_to\_zero\_mode}, \texttt{suppress\_na\_n}, and \texttt{suppress\_inf}.
  \item \textbf{Inline FAU \& Special Operands (\hyperref[sec:csf-tsd-fau]{\texttt{FAU\_k}}, \texttt{r8..r15}):} Any 8-bit source with \texttt{mode=2} reads zero-latency uniform \texttt{u0..u31} from \texttt{FAU\_k} (\texttt{slots64 = bits[63:56]}); \texttt{mode=3} reads \texttt{SmallConstant} (Table~1.2) or \texttt{FauSpecial} (\texttt{pc}, \texttt{lane\_id}, \texttt{lane\_width}, \texttt{tls\_ptr}, \texttt{wls\_ptr}, \texttt{blend\_descriptor\_0..7}).
\end{tightitemize}
\par\textbf{Worked Examples:}
\par\textbf{Example 1 (Normal positive float (v9--v10)):}
\codebox{\mnem{ILOGB.f32} \dst{r0}, \src{r1}}
With \texttt{r1 = 6.0f} (\texttt{1.5 * 2\textasciicircum{}2}): returns integer unbiased exponent \texttt{floor(log2(6.0)) = 2}.
\par\textbf{Example 2 (Special values (`0`, `Inf`, `NaN`)):}
\codebox{\mnem{ILOGB.f32} \dst{r0}, \src{r1}}
With \texttt{r1 = $\pm$0.0f} -\textgreater{} \texttt{FP\_ILOGB0} (\texttt{0x80000000}); with \texttt{r1 = $\pm$Inf} -\textgreater{} \texttt{INT32\_MAX} (\texttt{0x7FFFFFFF}); removed in v11.

\begin{center}\footnotesize
\begin{tabularx}{\linewidth}{@{}lllX@{}}
\toprule
Variant & Arch & Class & Opcode Fields \\
\midrule
\texttt{(none)} & v9--v14 & \texttt{R1} & \texttt{opcode2=ILOGB\_F32, opcode1=FLOAT} \\
\texttt{(none)} & v9--v14 & \texttt{R1} & \texttt{opcode2=ILOGB\_V2F16, opcode1=FLOAT} \\
\bottomrule
\end{tabularx}
\end{center}
\begin{center}\footnotesize
\begin{tabularx}{0.88\linewidth}{@{}l l X@{}}
\toprule
Field Name & Bits & Type / Fixed Value \\
\midrule
\texttt{src0} & \texttt{0:7} & \texttt{SourceEncoding} \\
\textit{Reserved} & \texttt{8:15} & \texttt{zero} \\
\texttt{opcode2} & \texttt{16:20} & \texttt{exact=ILOGB\_F32} \\
\textit{Reserved} & \texttt{21:27} & \texttt{zero} \\
\hyperref[enum:widen_m]{\texttt{widen0}} & \texttt{28:29} & \texttt{widen\_m (def none)} \\
\textit{Reserved} & \texttt{30:39} & \texttt{zero} \\
\texttt{dest\_register\_index} & \texttt{40:45} & \texttt{uint} \\
\hyperref[enum:dest_width_m]{\texttt{dest\_width}} & \texttt{46:47} & \texttt{dest\_width\_m} \\
\texttt{opcode1} & \texttt{48:51} & \texttt{exact=FLOAT} \\
\texttt{instruction\_class} & \texttt{52:56} & \texttt{exact=R1} \\
\texttt{fau\_page\_index} & \texttt{57:58} & \texttt{uint} \\
\hyperref[enum:flow_control_m]{\texttt{flow}} & \texttt{59:62} & \texttt{flow\_control\_m} \\
\textit{Reserved} & \texttt{63} & \texttt{zero} \\
\bottomrule
\end{tabularx}
\end{center}
\textbf{Fields Explained:}
\begin{tightitemize}\footnotesize
  \item \texttt{src0}: Primary source operand (\texttt{r0..r63}, \texttt{u0..u63} uniform, \texttt{k*} FAU constant, or \texttt{t0..t2} temporary; \texttt{\textasciicircum{}} marks discard/last-use). \emph{Example:} {\scriptsize\ttfamily \mnem{FADD.f32} \dst{r0}, \src{r1}, \src{r2}} (\textbf{SPIR-V:} \texttt{OpFAdd (Operand 0)}; \textbf{GLSL:} \texttt{float c = a + b;})
  \item \texttt{opcode2}: Secondary sub-opcode further specializing instructions sharing \texttt{instruction\_class} and \texttt{opcode1} (e.g., conversions, SFU tables, moves). \emph{Example:} {\scriptsize\ttfamily \mnem{F32\_TO\_S32}\imod{.round\_zero} \dst{r0}, \src{r1}} (\textbf{SPIR-V:} \texttt{OpConvertFToS}; \textbf{GLSL:} \texttt{int y = int(x);})
  \item \hyperref[enum:widen_m]{\texttt{widen0}}: Stage A widening extractor on \texttt{src0} (promotes FP16 \texttt{h0}/\texttt{h1} to FP32, or sign/zero-extends \texttt{b0..b3}/\texttt{h0..h1}/\texttt{w0..w1} to 32/64-bit before \texttt{abs}/\texttt{neg}). \emph{Example:} {\scriptsize\ttfamily \mnem{IADD.s32} \dst{r0}, \src{r1}\imod{.h0}, \src{r2}} (\textbf{SPIR-V:} \texttt{OpSConvert / OpFConvert}; \textbf{GLSL:} \texttt{int c = int(int16\_t(a)) + b;})
  \item \texttt{dest\_register\_index}: 6-bit destination GPR index (\texttt{r0..r63}) receiving the Stage C writeback result (even-aligned pair \texttt{[rN:rN+1]} for 64-bit ops). \emph{Example:} {\scriptsize\ttfamily \mnem{FMA.f32} \dst{r0}, \src{r1}, \src{r2}, \src{r3}} (\textbf{SPIR-V:} \texttt{SSA Result ID}; \textbf{GLSL:} \texttt{float r0 = fma(r1, r2, r3);})
  \item \hyperref[enum:dest_width_m]{\texttt{dest\_width}}: Stage C destination lane write-mask (\texttt{none} = full 32-bit, \texttt{h0} = low 16-bit, \texttt{h1} = high 16-bit without bit-shifting). \emph{Example:} {\scriptsize\ttfamily \mnem{FADD.v2f16} \dst{r0}\imod{.h1}, \src{r1}, \src{r2}} (\textbf{SPIR-V:} \texttt{OpVectorInsertDynamic / WriteMask}; \textbf{GLSL:} \texttt{dst.y = a.y + b.y;})
  \item \texttt{opcode1}: Primary sub-opcode within the instruction class uniquely identifying the operation family. \emph{Example:} {\scriptsize\ttfamily \mnem{IADD.s32} \dst{r0}, \src{r1}, \src{r2}} (\textbf{SPIR-V:} \texttt{OpIAdd (operation selector)}; \textbf{GLSL:} \texttt{// Primary opcode bits})
  \item \texttt{instruction\_class}: 4-bit major instruction class selector (\texttt{bits 48:51}) partitioning the 64-bit Valhall opcode space (e.g., \texttt{FMA}, \texttt{CVT}, \texttt{SFU}, \texttt{LS}, \texttt{TEX}, \texttt{BRANCH}). \emph{Example:} {\scriptsize\ttfamily \mnem{FADD.f32} \dst{r0}, \src{r1}, \src{r2}} (\textbf{SPIR-V:} \texttt{N/A (encoding class)}; \textbf{GLSL:} \texttt{// Hardware execution pipe / decode class})
  \item \texttt{fau\_page\_index}: Selects the FAU (Fast Access Uniform) page (\texttt{0..3}) when a source operand references hardware small constants, thread IDs, or lane constants. \emph{Example:} {\scriptsize\ttfamily \mnem{FADD.f32} \dst{r0}, \src{r1}, \src{k1f}} (\textbf{SPIR-V:} \texttt{OpConstant / BuiltIn}; \textbf{GLSL:} \texttt{float y = x + 1.0;})
  \item \hyperref[enum:flow_control_m]{\texttt{flow}}: 4-bit embedded flow-control \& scoreboard wait modifier (\texttt{bits 59:62}: \texttt{none}, \texttt{wait0..wait0126}, \texttt{wait}, \texttt{reconverge}, \texttt{discard}, \texttt{end}). \emph{Example:} {\scriptsize\ttfamily \mnem{FADD.f32}\imod{.wait0.reconverge} \dst{r0}, \src{r1}, \src{r2}} (\textbf{SPIR-V:} \texttt{OpControlBarrier / OpReturn}; \textbf{GLSL:} \texttt{// Scoreboard wait + warp reconvergence})
\end{tightitemize}
\clearpage

\subsection*{LOGB}
\addcontentsline{toc}{subsection}{LOGB}
Extracts the unbiased base-2 floating-point exponent \texttt{floor(log2(\textbar{}src0[i]\textbar{}))} for each floating-point element $i \in \{0, \dots, N-1\}$ ($N=1$ \texttt{f32}, $N=2$ \texttt{v2f16}; arch \textless{}= v10).
\par\textbf{Notes:}
\begin{tightitemize}
  \item IEEE-754 \texttt{logb} special cases: \texttt{+-0 -\textgreater{} -Inf}, \texttt{+-Inf -\textgreater{} +Inf}, \texttt{NaN -\textgreater{} NaN}.
  \item \textbf{Source references:} \texttt{bifrost/bifrost/ISA.xml:7386-7396}, \texttt{isa-v9-v14.xml}.
\end{tightitemize}
\codebox{\mnem{LOGB}\imod{\{.flow\}} \dst{dest}\imod{\{.dest\_width\}}, \src{src0}\imod{\{.widen0\}}\\\mnem{LOGB}\imod{\{.flow\}} \dst{dest}\imod{\{.dest\_width\}}, \src{src0}\imod{\{.swz0\}}}
\textbf{Operands:}
\begin{tightitemize}\footnotesize
  \item \texttt{\dst{dest}} (\texttt{dest\_register\_index}, \texttt{dest\_width}): Floating-point destination GPR receiving \texttt{floor(log2(\textbar{}src0\textbar{}))} (\texttt{-Inf} on 0, \texttt{+Inf} on Inf).
  \item \texttt{\src{src0}}: Floating-point input.
\end{tightitemize}
\par\textbf{Closest SPIR-V Op:} \texttt{None (OpenCL `logb`, v9--v10)}
\par\textbf{Pseudo-GLSL (Operational Semantics):}
\codebox{for (int i = 0; i \textless{} N; ++i)\\    dest[i] = (src0[i] == 0.0) ? -INF : (isinf(src0[i]) ? +INF : floor(log2(abs(src0[i]))));}
\textbf{Execution Context:}
\begin{tightitemize}\footnotesize
  \item \textbf{Compute stage preloads (\hyperref[sec:csf-spd]{\texttt{SPD\_k}}, \texttt{r16..r23}):} Hardware pre-fills \texttt{r55..r57} with \texttt{WorkGroupID.xyz + JOB\_OFFSET} (from \hyperref[sec:csf-regs]{\texttt{COMPUTESR::JOB\_OFFSET\_X/Y/Z}} \texttt{r34..r36}) and \texttt{r60..r62} with \texttt{LocalInvocationID.xyz} or \texttt{GlobalInvocationID.xyz} (from \texttt{WG\_SIZE} \texttt{r33} \& \texttt{JOB\_SIZE} \texttt{r37..r39}, selected by \texttt{ShaderProgram::preload}) so compute shaders can derive thread/element indices without memory loads.
  \item \textbf{Vertex stage preloads (\hyperref[sec:csf-spd]{\texttt{SPD\_0}/\texttt{SPD\_1}}, \texttt{r16..r19}):} The hardware Index Fetcher reads \hyperref[sec:csf-regs]{\texttt{IDVSSR::INDEX\_BUFFER}} (\texttt{r54:r55}) and pre-fills \texttt{r59, r60} with \texttt{VertexIndex} (\texttt{+ VERTEX\_OFFSET + GLOBAL\_ATTR\_OFF}) and \texttt{r61} with \texttt{InstanceIndex} (\texttt{+ INSTANCE\_OFFSET}) so Position and Varying shaders can pull vertex attributes directly.
  \item \textbf{Fragment stage preloads \& FP modes (\hyperref[sec:csf-spd]{\texttt{SPD\_2}}, \texttt{r20:r21}):} Hardware pre-fills \texttt{r59} with \texttt{floatBitsToUint(FragCoord.z)} and \texttt{r60} with the MSAA sample coverage mask (\texttt{gl\_SampleMaskIn[0]}), while \texttt{ShaderProgram} word~0 latches 32/64 GPR allocation and warp-wide \texttt{flush\_to\_zero\_mode}, \texttt{suppress\_na\_n}, and \texttt{suppress\_inf}.
  \item \textbf{Inline FAU \& Special Operands (\hyperref[sec:csf-tsd-fau]{\texttt{FAU\_k}}, \texttt{r8..r15}):} Any 8-bit source with \texttt{mode=2} reads zero-latency uniform \texttt{u0..u31} from \texttt{FAU\_k} (\texttt{slots64 = bits[63:56]}); \texttt{mode=3} reads \texttt{SmallConstant} (Table~1.2) or \texttt{FauSpecial} (\texttt{pc}, \texttt{lane\_id}, \texttt{lane\_width}, \texttt{tls\_ptr}, \texttt{wls\_ptr}, \texttt{blend\_descriptor\_0..7}).
\end{tightitemize}
\par\textbf{Worked Examples:}
\par\textbf{Example 1 (Normal float unbiased exponent (v9--v10)):}
\codebox{\mnem{LOGB.f32} \dst{r0}, \src{r1}}
With \texttt{r1 = 6.0f}: returns float unbiased exponent \texttt{2.0f} (\texttt{0x40000000}).
\par\textbf{Example 2 (Zero and infinity):}
\codebox{\mnem{LOGB.f32} \dst{r0}, \src{r1}}
With \texttt{r1 = $\pm$0.0f} -\textgreater{} \texttt{-Inf} (\texttt{0xFF800000}); with \texttt{r1 = $\pm$Inf} -\textgreater{} \texttt{+Inf} (\texttt{0x7F800000}); removed in v11.

\begin{center}\footnotesize
\begin{tabularx}{\linewidth}{@{}lllX@{}}
\toprule
Variant & Arch & Class & Opcode Fields \\
\midrule
\texttt{(none)} & v9--v14 & \texttt{R1} & \texttt{opcode2=LOGB\_F32, opcode1=FLOAT} \\
\texttt{(none)} & v9--v14 & \texttt{R1} & \texttt{opcode2=LOGB\_V2F16, opcode1=FLOAT} \\
\bottomrule
\end{tabularx}
\end{center}
\begin{center}\footnotesize
\begin{tabularx}{0.88\linewidth}{@{}l l X@{}}
\toprule
Field Name & Bits & Type / Fixed Value \\
\midrule
\texttt{src0} & \texttt{0:7} & \texttt{SourceEncoding} \\
\textit{Reserved} & \texttt{8:15} & \texttt{zero} \\
\texttt{opcode2} & \texttt{16:20} & \texttt{exact=LOGB\_F32} \\
\textit{Reserved} & \texttt{21:27} & \texttt{zero} \\
\hyperref[enum:widen_m]{\texttt{widen0}} & \texttt{28:29} & \texttt{widen\_m (def none)} \\
\textit{Reserved} & \texttt{30:39} & \texttt{zero} \\
\texttt{dest\_register\_index} & \texttt{40:45} & \texttt{uint} \\
\hyperref[enum:dest_width_m]{\texttt{dest\_width}} & \texttt{46:47} & \texttt{dest\_width\_m} \\
\texttt{opcode1} & \texttt{48:51} & \texttt{exact=FLOAT} \\
\texttt{instruction\_class} & \texttt{52:56} & \texttt{exact=R1} \\
\texttt{fau\_page\_index} & \texttt{57:58} & \texttt{uint} \\
\hyperref[enum:flow_control_m]{\texttt{flow}} & \texttt{59:62} & \texttt{flow\_control\_m} \\
\textit{Reserved} & \texttt{63} & \texttt{zero} \\
\bottomrule
\end{tabularx}
\end{center}
\textbf{Fields Explained:}
\begin{tightitemize}\footnotesize
  \item \texttt{src0}: Primary source operand (\texttt{r0..r63}, \texttt{u0..u63} uniform, \texttt{k*} FAU constant, or \texttt{t0..t2} temporary; \texttt{\textasciicircum{}} marks discard/last-use). \emph{Example:} {\scriptsize\ttfamily \mnem{FADD.f32} \dst{r0}, \src{r1}, \src{r2}} (\textbf{SPIR-V:} \texttt{OpFAdd (Operand 0)}; \textbf{GLSL:} \texttt{float c = a + b;})
  \item \texttt{opcode2}: Secondary sub-opcode further specializing instructions sharing \texttt{instruction\_class} and \texttt{opcode1} (e.g., conversions, SFU tables, moves). \emph{Example:} {\scriptsize\ttfamily \mnem{F32\_TO\_S32}\imod{.round\_zero} \dst{r0}, \src{r1}} (\textbf{SPIR-V:} \texttt{OpConvertFToS}; \textbf{GLSL:} \texttt{int y = int(x);})
  \item \hyperref[enum:widen_m]{\texttt{widen0}}: Stage A widening extractor on \texttt{src0} (promotes FP16 \texttt{h0}/\texttt{h1} to FP32, or sign/zero-extends \texttt{b0..b3}/\texttt{h0..h1}/\texttt{w0..w1} to 32/64-bit before \texttt{abs}/\texttt{neg}). \emph{Example:} {\scriptsize\ttfamily \mnem{IADD.s32} \dst{r0}, \src{r1}\imod{.h0}, \src{r2}} (\textbf{SPIR-V:} \texttt{OpSConvert / OpFConvert}; \textbf{GLSL:} \texttt{int c = int(int16\_t(a)) + b;})
  \item \texttt{dest\_register\_index}: 6-bit destination GPR index (\texttt{r0..r63}) receiving the Stage C writeback result (even-aligned pair \texttt{[rN:rN+1]} for 64-bit ops). \emph{Example:} {\scriptsize\ttfamily \mnem{FMA.f32} \dst{r0}, \src{r1}, \src{r2}, \src{r3}} (\textbf{SPIR-V:} \texttt{SSA Result ID}; \textbf{GLSL:} \texttt{float r0 = fma(r1, r2, r3);})
  \item \hyperref[enum:dest_width_m]{\texttt{dest\_width}}: Stage C destination lane write-mask (\texttt{none} = full 32-bit, \texttt{h0} = low 16-bit, \texttt{h1} = high 16-bit without bit-shifting). \emph{Example:} {\scriptsize\ttfamily \mnem{FADD.v2f16} \dst{r0}\imod{.h1}, \src{r1}, \src{r2}} (\textbf{SPIR-V:} \texttt{OpVectorInsertDynamic / WriteMask}; \textbf{GLSL:} \texttt{dst.y = a.y + b.y;})
  \item \texttt{opcode1}: Primary sub-opcode within the instruction class uniquely identifying the operation family. \emph{Example:} {\scriptsize\ttfamily \mnem{IADD.s32} \dst{r0}, \src{r1}, \src{r2}} (\textbf{SPIR-V:} \texttt{OpIAdd (operation selector)}; \textbf{GLSL:} \texttt{// Primary opcode bits})
  \item \texttt{instruction\_class}: 4-bit major instruction class selector (\texttt{bits 48:51}) partitioning the 64-bit Valhall opcode space (e.g., \texttt{FMA}, \texttt{CVT}, \texttt{SFU}, \texttt{LS}, \texttt{TEX}, \texttt{BRANCH}). \emph{Example:} {\scriptsize\ttfamily \mnem{FADD.f32} \dst{r0}, \src{r1}, \src{r2}} (\textbf{SPIR-V:} \texttt{N/A (encoding class)}; \textbf{GLSL:} \texttt{// Hardware execution pipe / decode class})
  \item \texttt{fau\_page\_index}: Selects the FAU (Fast Access Uniform) page (\texttt{0..3}) when a source operand references hardware small constants, thread IDs, or lane constants. \emph{Example:} {\scriptsize\ttfamily \mnem{FADD.f32} \dst{r0}, \src{r1}, \src{k1f}} (\textbf{SPIR-V:} \texttt{OpConstant / BuiltIn}; \textbf{GLSL:} \texttt{float y = x + 1.0;})
  \item \hyperref[enum:flow_control_m]{\texttt{flow}}: 4-bit embedded flow-control \& scoreboard wait modifier (\texttt{bits 59:62}: \texttt{none}, \texttt{wait0..wait0126}, \texttt{wait}, \texttt{reconverge}, \texttt{discard}, \texttt{end}). \emph{Example:} {\scriptsize\ttfamily \mnem{FADD.f32}\imod{.wait0.reconverge} \dst{r0}, \src{r1}, \src{r2}} (\textbf{SPIR-V:} \texttt{OpControlBarrier / OpReturn}; \textbf{GLSL:} \texttt{// Scoreboard wait + warp reconvergence})
\end{tightitemize}
\clearpage

\section{Cubemap Helpers (4 mnemonics)}
\subsection*{CUBEFACE1}
\addcontentsline{toc}{subsection}{CUBEFACE1}
Computes the cubemap major-axis magnitude from 3D direction \texttt{(x=src0, y=src1, z=src2)}: \texttt{m = max(\textbar{}x\textbar{}, \textbar{}y\textbar{}, \textbar{}z\textbar{})}. Writes \texttt{NaN} if any input is NaN, \texttt{1.0} if \texttt{m} is infinite, \texttt{2.0} if \texttt{m} is zero or subnormal (preventing divide-by-zero), and \texttt{m} otherwise to \texttt{dest}.
\par\textbf{Notes:}
\begin{tightitemize}
  \item Verified in \texttt{hw\_tests.rs:1443-1450} (\texttt{Ulp(1)}).
  \item \textbf{Source references:} \texttt{kraid/ops.rs:948-994}, \texttt{encode\_v9.rs:1074-1094}, \texttt{nir.rs:1958-1971}, \texttt{hw\_tests.rs:1443-1450}.
\end{tightitemize}
\codebox{\mnem{CUBEFACE1}\imod{\{.flow\}} \dst{dest}\imod{\{.dest\_width\}}, \src{src0}\imod{\{.widen0\}}\imod{\{.abs0\}}\imod{\{.neg0\}}, \src{src1}\imod{\{.widen1\}}\imod{\{.abs1\}}\imod{\{.neg1\}}, \src{src2}\imod{\{.widen2\}}\imod{\{.abs2\}}\imod{\{.neg2\}}}
\textbf{Operands:}
\begin{tightitemize}\footnotesize
  \item \texttt{\dst{dest}} (\texttt{dest\_register\_index}, \texttt{dest\_width}): FP32 destination GPR receiving major-axis magnitude \texttt{max(\textbar{}x\textbar{}, \textbar{}y\textbar{}, \textbar{}z\textbar{})} (\texttt{1.0} on Inf, \texttt{2.0} on zero/subnormal).
  \item \texttt{\src{src0}}: FP32 cubemap direction \texttt{x} component.
  \item \texttt{\src{src1}}: FP32 cubemap direction \texttt{y} component.
  \item \texttt{\src{src2}}: FP32 cubemap direction \texttt{z} component.
\end{tightitemize}
\par\textbf{Closest SPIR-V Op:} \texttt{None (lowered from OpImageSample* Cube, v9--v10)}
\par\textbf{Pseudo-GLSL (Operational Semantics):}
\codebox{float m = max(abs(src0), max(abs(src1), abs(src2)));\\if (isnan(src0) \textbar{}\textbar{} isnan(src1) \textbar{}\textbar{} isnan(src2)) dest = NaN;\\else if (isinf(m)) dest = 1.0;\\else if (m \textless{} FLT\_MIN) dest = 2.0;\\else dest = m;}
\textbf{Execution Context:}
\begin{tightitemize}\footnotesize
  \item \textbf{Compute stage preloads (\hyperref[sec:csf-spd]{\texttt{SPD\_k}}, \texttt{r16..r23}):} Hardware pre-fills \texttt{r55..r57} with \texttt{WorkGroupID.xyz + JOB\_OFFSET} (from \hyperref[sec:csf-regs]{\texttt{COMPUTESR::JOB\_OFFSET\_X/Y/Z}} \texttt{r34..r36}) and \texttt{r60..r62} with \texttt{LocalInvocationID.xyz} or \texttt{GlobalInvocationID.xyz} (from \texttt{WG\_SIZE} \texttt{r33} \& \texttt{JOB\_SIZE} \texttt{r37..r39}, selected by \texttt{ShaderProgram::preload}) so compute shaders can derive thread/element indices without memory loads.
  \item \textbf{Vertex stage preloads (\hyperref[sec:csf-spd]{\texttt{SPD\_0}/\texttt{SPD\_1}}, \texttt{r16..r19}):} The hardware Index Fetcher reads \hyperref[sec:csf-regs]{\texttt{IDVSSR::INDEX\_BUFFER}} (\texttt{r54:r55}) and pre-fills \texttt{r59, r60} with \texttt{VertexIndex} (\texttt{+ VERTEX\_OFFSET + GLOBAL\_ATTR\_OFF}) and \texttt{r61} with \texttt{InstanceIndex} (\texttt{+ INSTANCE\_OFFSET}) so Position and Varying shaders can pull vertex attributes directly.
  \item \textbf{Fragment stage preloads \& FP modes (\hyperref[sec:csf-spd]{\texttt{SPD\_2}}, \texttt{r20:r21}):} Hardware pre-fills \texttt{r59} with \texttt{floatBitsToUint(FragCoord.z)} and \texttt{r60} with the MSAA sample coverage mask (\texttt{gl\_SampleMaskIn[0]}), while \texttt{ShaderProgram} word~0 latches 32/64 GPR allocation and warp-wide \texttt{flush\_to\_zero\_mode}, \texttt{suppress\_na\_n}, and \texttt{suppress\_inf}.
  \item \textbf{Inline FAU \& Special Operands (\hyperref[sec:csf-tsd-fau]{\texttt{FAU\_k}}, \texttt{r8..r15}):} Any 8-bit source with \texttt{mode=2} reads zero-latency uniform \texttt{u0..u31} from \texttt{FAU\_k} (\texttt{slots64 = bits[63:56]}); \texttt{mode=3} reads \texttt{SmallConstant} (Table~1.2) or \texttt{FauSpecial} (\texttt{pc}, \texttt{lane\_id}, \texttt{lane\_width}, \texttt{tls\_ptr}, \texttt{wls\_ptr}, \texttt{blend\_descriptor\_0..7}).
\end{tightitemize}
\par\textbf{Worked Examples:}
\par\textbf{Example 1 (Major-axis magnitude extraction):}
\codebox{\mnem{CUBEFACE1} \dst{r0}, \src{r1}, \src{r2}, \src{r3}}
With \texttt{(x, y, z) = (1.0, -4.0, 2.0)}: returns \texttt{max(\textbar{}x\textbar{}, \textbar{}y\textbar{}, \textbar{}z\textbar{}) = 4.0f} (\texttt{0x40800000}) for \texttt{FRCP} projection.
\par\textbf{Example 2 (Infinity and zero clamping):}
\codebox{\mnem{CUBEFACE1} \dst{r0}, \src{r1}, \src{r2}, \src{r3}}
If any component is \texttt{$\pm$Inf}, returns \texttt{1.0f}; if all components are \texttt{$\pm$0.0}/subnormal, returns \texttt{2.0f} to avoid division by zero.
\par\textbf{Mined SPIR-V $\to$ Valhall ISA Example (Full 10-instruction SPIR-V Cube texture sampling lowering --- Fragment shader, Mali-G710 (v10)):}
\codebox{\textbf{; SPIR-V (input)}\\\%cube\_img = OpTypeImage \%float Cube 0 0 0 1 Unknown\\\%color    = OpImageSampleImplicitLod \%v4float \%sampled\_cube \%dir\_xyz None\\[3pt]\textbf{; Valhall ISA (malioc output)}\\CUBEFACE2 r6, r5, r4, r3          ; face\_idx\textless{}\textless{}29 \textbar{} inf\_flag\textless{}\textless{}16 from (x=r5,y=r4,z=r3)\\CUBE\_SSEL r7, r3, r5, r6          ; select face S numerator from (z=r3, x=r5, face=r6)\\CUBEFACE1 r5, r5\textasciicircum{}, r4, r3         ; major axis \textbar{}M\textbar{} = max(\textbar{}x\textbar{},\textbar{}y\textbar{},\textbar{}z\textbar{})\\FRCP.f32 r5, r5\textasciicircum{}                  ; 1 / \textbar{}M\textbar{}\\FMUL.f32 r5, r5\textasciicircum{}, k0fr2f.h1       ; 0.5 / \textbar{}M\textbar{}\\FMA.f32.clamp\_0\_1 r7, r5, r7\textasciicircum{}, k0fr2f.h1 ; s = clamp(S/(2\textbar{}M\textbar{}) + 0.5, 0, 1)\\MUX.i32.bit r8, r7\textasciicircum{}, r6, u1.w1    ; pack 3-bit face\_idx into s\\CUBE\_TSEL r3, r4\textasciicircum{}, r3\textasciicircum{}, r6\textasciicircum{}       ; select face T numerator from (y=r4, z=r3, face=r6)\\FMA.f32.clamp\_0\_1 r9, r5\textasciicircum{}, r3\textasciicircum{}, k0fr2f.h1 ; t = clamp(T/(2\textbar{}M\textbar{}) + 0.5, 0, 1)\\TEX\_SINGLE.slot0.reconverge.skip.cube.dst32.rgba.wide.computed @r4:r5:r6:r7, @r8:r9, u2}
Valhall does not accept raw 3D \texttt{(x, y, z)} vectors in \texttt{TEX\_SINGLE.cube}. Instead, \texttt{malioc} lowers \texttt{OpImageSampleImplicitLod} on a \texttt{Cube} image using all four cubemap helpers: \texttt{CUBEFACE1} computes the major-axis magnitude $|M|$ (returning \texttt{1.0} on Inf and \texttt{2.0} on 0 to prevent NaN division), \texttt{FRCP} + \texttt{FMUL 0.5} computes $0.5/|M|$, and \texttt{FMA.clamp\_0\_1} maps both \texttt{S} and \texttt{T} into $[0, 1]$ before \texttt{TEX\_SINGLE.cube}.
\begin{center}\footnotesize
\begin{tabularx}{\linewidth}{@{}lllX@{}}
\toprule
Variant & Arch & Class & Opcode Fields \\
\midrule
\texttt{(none)} & v9--v14 & \texttt{R3A} & \texttt{opcode1=CUBEFACE1} \\
\bottomrule
\end{tabularx}
\end{center}
\begin{center}\footnotesize
\begin{tabularx}{0.88\linewidth}{@{}l l X@{}}
\toprule
Field Name & Bits & Type / Fixed Value \\
\midrule
\texttt{src0} & \texttt{0:7} & \texttt{SourceEncoding} \\
\texttt{src1} & \texttt{8:15} & \texttt{SourceEncoding} \\
\texttt{src2} & \texttt{16:23} & \texttt{SourceEncoding} \\
\hyperref[enum:widen_m]{\texttt{widen2}} & \texttt{24:25} & \texttt{widen\_m (def none)} \\
\hyperref[enum:widen_m]{\texttt{widen1}} & \texttt{26:27} & \texttt{widen\_m (def none)} \\
\hyperref[enum:widen_m]{\texttt{widen0}} & \texttt{28:29} & \texttt{widen\_m (def none)} \\
\textit{Reserved} & \texttt{30:33} & \texttt{zero} \\
\hyperref[enum:neg_m]{\texttt{neg2}} & \texttt{34} & \texttt{neg\_m (def none)} \\
\hyperref[enum:abs_m]{\texttt{abs2}} & \texttt{35} & \texttt{abs\_m (def none)} \\
\hyperref[enum:neg_m]{\texttt{neg1}} & \texttt{36} & \texttt{neg\_m (def none)} \\
\hyperref[enum:abs_m]{\texttt{abs1}} & \texttt{37} & \texttt{abs\_m (def none)} \\
\hyperref[enum:neg_m]{\texttt{neg0}} & \texttt{38} & \texttt{neg\_m (def none)} \\
\hyperref[enum:abs_m]{\texttt{abs0}} & \texttt{39} & \texttt{abs\_m (def none)} \\
\texttt{dest\_register\_index} & \texttt{40:45} & \texttt{uint} \\
\hyperref[enum:dest_width_m]{\texttt{dest\_width}} & \texttt{46:47} & \texttt{dest\_width\_m} \\
\texttt{opcode1} & \texttt{48:51} & \texttt{exact=CUBEFACE1} \\
\texttt{instruction\_class} & \texttt{52:56} & \texttt{exact=R3A} \\
\texttt{fau\_page\_index} & \texttt{57:58} & \texttt{uint} \\
\hyperref[enum:flow_control_m]{\texttt{flow}} & \texttt{59:62} & \texttt{flow\_control\_m} \\
\textit{Reserved} & \texttt{63} & \texttt{zero} \\
\bottomrule
\end{tabularx}
\end{center}
\textbf{Fields Explained:}
\begin{tightitemize}\footnotesize
  \item \texttt{src0}: Primary source operand (\texttt{r0..r63}, \texttt{u0..u63} uniform, \texttt{k*} FAU constant, or \texttt{t0..t2} temporary; \texttt{\textasciicircum{}} marks discard/last-use). \emph{Example:} {\scriptsize\ttfamily \mnem{FADD.f32} \dst{r0}, \src{r1}, \src{r2}} (\textbf{SPIR-V:} \texttt{OpFAdd (Operand 0)}; \textbf{GLSL:} \texttt{float c = a + b;})
  \item \texttt{src1}: Secondary source operand (register, uniform, or FAU constant; in shifts, supplies the shift count byte). \emph{Example:} {\scriptsize\ttfamily \mnem{LSHIFT.i32} \dst{r0}, \src{r1}, \src{r2}\imod{.b0}} (\textbf{SPIR-V:} \texttt{OpShiftLeftLogical (Shift)}; \textbf{GLSL:} \texttt{uint y = a \textless{}\textless{} s;})
  \item \texttt{src2}: Third source operand (FMA addend, \texttt{MUX}/\texttt{CSEL} selector, \texttt{IADDC} carry-in, or bitwise mask in fused shift-logic). \emph{Example:} {\scriptsize\ttfamily \mnem{FMA.f32} \dst{r0}, \src{r1}, \src{r2}, \src{r3}} (\textbf{SPIR-V:} \texttt{GLSL.std.450 Fma (Addend)}; \textbf{GLSL:} \texttt{float d = fma(a, b, c);})
  \item \hyperref[enum:widen_m]{\texttt{widen2}}: Stage A widening extractor on \texttt{src2} (e.g., promoting FP16 addend \texttt{h0}/\texttt{h1} to FP32 in \texttt{FMA.f32}). \emph{Example:} {\scriptsize\ttfamily \mnem{FMA.f32} \dst{r0}, \src{r1}, \src{r2}, \src{r3}\imod{.h0}} (\textbf{SPIR-V:} \texttt{OpFConvert + GLSL.std.450 Fma}; \textbf{GLSL:} \texttt{float d = fma(a, b, float(c16));})
  \item \hyperref[enum:widen_m]{\texttt{widen1}}: Stage A widening extractor on \texttt{src1} (promotes narrow byte/half/word lane before Stage B). \emph{Example:} {\scriptsize\ttfamily \mnem{IMUL.u64} \dst{[r0:r1]}, \src{[r2:r3]}\imod{.w0}, \src{[r4:r5]}\imod{.w0}} (\textbf{SPIR-V:} \texttt{OpUConvert / OpUMulExtended}; \textbf{GLSL:} \texttt{uint64\_t p = uint64\_t(a) * uint64\_t(b);})
  \item \hyperref[enum:widen_m]{\texttt{widen0}}: Stage A widening extractor on \texttt{src0} (promotes FP16 \texttt{h0}/\texttt{h1} to FP32, or sign/zero-extends \texttt{b0..b3}/\texttt{h0..h1}/\texttt{w0..w1} to 32/64-bit before \texttt{abs}/\texttt{neg}). \emph{Example:} {\scriptsize\ttfamily \mnem{IADD.s32} \dst{r0}, \src{r1}\imod{.h0}, \src{r2}} (\textbf{SPIR-V:} \texttt{OpSConvert / OpFConvert}; \textbf{GLSL:} \texttt{int c = int(int16\_t(a)) + b;})
  \item \hyperref[enum:neg_m]{\texttt{neg2}}: Stage A floating-point negation on \texttt{src2} (flips IEEE sign bit of FMA addend after swizzle/widen and \texttt{abs2}). \emph{Example:} {\scriptsize\ttfamily \mnem{FMA.f32} \dst{r0}, \src{r1}, \src{r2}, \src{r3}\imod{.neg}} (\textbf{SPIR-V:} \texttt{OpFNegate}; \textbf{GLSL:} \texttt{float d = fma(a, b, -c);})
  \item \hyperref[enum:abs_m]{\texttt{abs2}}: Stage A floating-point absolute value on \texttt{src2} (clears IEEE sign bit of FMA addend after swizzle/widen). \emph{Example:} {\scriptsize\ttfamily \mnem{FMA.f32} \dst{r0}, \src{r1}, \src{r2}, \src{r3}\imod{.abs}} (\textbf{SPIR-V:} \texttt{GLSL.std.450 FAbs}; \textbf{GLSL:} \texttt{float d = fma(a, b, abs(c));})
  \item \hyperref[enum:neg_m]{\texttt{neg1}}: Stage A floating-point negation on \texttt{src1} (flips IEEE sign bit after swizzle/widen and \texttt{abs1}). \emph{Example:} {\scriptsize\ttfamily \mnem{FMA.f32} \dst{r0}, \src{r1}, \src{r2}\imod{.neg}, \src{r3}} (\textbf{SPIR-V:} \texttt{OpFNegate}; \textbf{GLSL:} \texttt{float d = fma(a, -b, c);})
  \item \hyperref[enum:abs_m]{\texttt{abs1}}: Stage A floating-point absolute value on \texttt{src1} (clears IEEE sign bit after swizzle/widen, before \texttt{neg1}). \emph{Example:} {\scriptsize\ttfamily \mnem{FMA.f32} \dst{r0}, \src{r1}, \src{r2}\imod{.abs}, \src{r3}} (\textbf{SPIR-V:} \texttt{GLSL.std.450 FAbs}; \textbf{GLSL:} \texttt{float d = fma(a, abs(b), c);})
  \item \hyperref[enum:neg_m]{\texttt{neg0}}: Stage A floating-point negation on \texttt{src0} (flips IEEE sign bit after swizzle/widen and \texttt{abs0}). \emph{Example:} {\scriptsize\ttfamily \mnem{FADD.f32} \dst{r0}, \src{r1}\imod{.neg}, \src{r2}} (\textbf{SPIR-V:} \texttt{OpFNegate}; \textbf{GLSL:} \texttt{float c = -a + b;})
  \item \hyperref[enum:abs_m]{\texttt{abs0}}: Stage A floating-point absolute value on \texttt{src0} (clears IEEE sign bit after swizzle/widen, before \texttt{neg0}). \emph{Example:} {\scriptsize\ttfamily \mnem{FADD.f32} \dst{r0}, \src{r1}\imod{.abs}, \src{r2}} (\textbf{SPIR-V:} \texttt{GLSL.std.450 FAbs}; \textbf{GLSL:} \texttt{float c = abs(a) + b;})
  \item \texttt{dest\_register\_index}: 6-bit destination GPR index (\texttt{r0..r63}) receiving the Stage C writeback result (even-aligned pair \texttt{[rN:rN+1]} for 64-bit ops). \emph{Example:} {\scriptsize\ttfamily \mnem{FMA.f32} \dst{r0}, \src{r1}, \src{r2}, \src{r3}} (\textbf{SPIR-V:} \texttt{SSA Result ID}; \textbf{GLSL:} \texttt{float r0 = fma(r1, r2, r3);})
  \item \hyperref[enum:dest_width_m]{\texttt{dest\_width}}: Stage C destination lane write-mask (\texttt{none} = full 32-bit, \texttt{h0} = low 16-bit, \texttt{h1} = high 16-bit without bit-shifting). \emph{Example:} {\scriptsize\ttfamily \mnem{FADD.v2f16} \dst{r0}\imod{.h1}, \src{r1}, \src{r2}} (\textbf{SPIR-V:} \texttt{OpVectorInsertDynamic / WriteMask}; \textbf{GLSL:} \texttt{dst.y = a.y + b.y;})
  \item \texttt{opcode1}: Primary sub-opcode within the instruction class uniquely identifying the operation family. \emph{Example:} {\scriptsize\ttfamily \mnem{IADD.s32} \dst{r0}, \src{r1}, \src{r2}} (\textbf{SPIR-V:} \texttt{OpIAdd (operation selector)}; \textbf{GLSL:} \texttt{// Primary opcode bits})
  \item \texttt{instruction\_class}: 4-bit major instruction class selector (\texttt{bits 48:51}) partitioning the 64-bit Valhall opcode space (e.g., \texttt{FMA}, \texttt{CVT}, \texttt{SFU}, \texttt{LS}, \texttt{TEX}, \texttt{BRANCH}). \emph{Example:} {\scriptsize\ttfamily \mnem{FADD.f32} \dst{r0}, \src{r1}, \src{r2}} (\textbf{SPIR-V:} \texttt{N/A (encoding class)}; \textbf{GLSL:} \texttt{// Hardware execution pipe / decode class})
  \item \texttt{fau\_page\_index}: Selects the FAU (Fast Access Uniform) page (\texttt{0..3}) when a source operand references hardware small constants, thread IDs, or lane constants. \emph{Example:} {\scriptsize\ttfamily \mnem{FADD.f32} \dst{r0}, \src{r1}, \src{k1f}} (\textbf{SPIR-V:} \texttt{OpConstant / BuiltIn}; \textbf{GLSL:} \texttt{float y = x + 1.0;})
  \item \hyperref[enum:flow_control_m]{\texttt{flow}}: 4-bit embedded flow-control \& scoreboard wait modifier (\texttt{bits 59:62}: \texttt{none}, \texttt{wait0..wait0126}, \texttt{wait}, \texttt{reconverge}, \texttt{discard}, \texttt{end}). \emph{Example:} {\scriptsize\ttfamily \mnem{FADD.f32}\imod{.wait0.reconverge} \dst{r0}, \src{r1}, \src{r2}} (\textbf{SPIR-V:} \texttt{OpControlBarrier / OpReturn}; \textbf{GLSL:} \texttt{// Scoreboard wait + warp reconvergence})
\end{tightitemize}
\clearpage

\subsection*{CUBEFACE2}
\addcontentsline{toc}{subsection}{CUBEFACE2}
Determines the dominant cubemap face index (\texttt{0=+X, 1=-X, 2=+Y, 3=-Y, 4=+Z, 5=-Z}) from 3D direction \texttt{(x=src0, y=src1, z=src2)} with tie-breaking priority $|z| \ge |y| \ge |x|$ (NaN defaults to face \texttt{4}). Packs the 3-bit face index into bits \texttt{[31:29]} and sets bit \texttt{16} if any finite/infinite component is \texttt{$\pm$Inf} (and no component is NaN): \texttt{dest = (face \textless{}\textless{} 29) \textbar{} (is\_inf \textless{}\textless{} 16)}.
\par\textbf{Notes:}
\begin{tightitemize}
  \item Tie-breaking priority: \texttt{\textbar{}z\textbar{} \textgreater{}= \textbar{}y\textbar{} \textgreater{}= \textbar{}x\textbar{}}.
  \item Exact match verified in \texttt{hw\_tests.rs:1433-1440}.
  \item \textbf{Source references:} \texttt{kraid/ops.rs:898-944}, \texttt{encode\_v9.rs:1052-1072}, \texttt{nir.rs:1958-1971}, \texttt{hw\_tests.rs:1433-1440}.
\end{tightitemize}
\codebox{\mnem{CUBEFACE2}\imod{\{.flow\}} \dst{dest}\imod{\{.dest\_width\}}, \src{src0}\imod{\{.widen0\}}\imod{\{.abs0\}}\imod{\{.neg0\}}, \src{src1}\imod{\{.widen1\}}\imod{\{.abs1\}}\imod{\{.neg1\}}, \src{src2}\imod{\{.widen2\}}\imod{\{.abs2\}}\imod{\{.neg2\}}}
\textbf{Operands:}
\begin{tightitemize}\footnotesize
  \item \texttt{\dst{dest}} (\texttt{dest\_register\_index}, \texttt{dest\_width}): 32-bit destination GPR receiving \texttt{(face\_idx \textless{}\textless{} 29) \textbar{} (any\_inf \textless{}\textless{} 16)}.
  \item \texttt{\src{src0}}: FP32 cubemap direction \texttt{x} component.
  \item \texttt{\src{src1}}: FP32 cubemap direction \texttt{y} component.
  \item \texttt{\src{src2}}: FP32 cubemap direction \texttt{z} component.
\end{tightitemize}
\par\textbf{Closest SPIR-V Op:} \texttt{None (lowered from OpImageSample* Cube, v9--v10)}
\par\textbf{Pseudo-GLSL (Operational Semantics):}
\codebox{uint face = majorAxisFaceIndex(src0, src1, src2); // \textbar{}z\textbar{} \textgreater{}= \textbar{}y\textbar{} \textgreater{}= \textbar{}x\textbar{}; NaN -\textgreater{} 4\\uint inf\_bit = (anyInf(src0, src1, src2) \&\& !anyNaN(src0, src1, src2)) ? 0x10000u : 0u;\\dest = (face \textless{}\textless{} 29) \textbar{} inf\_bit;}
\textbf{Execution Context:}
\begin{tightitemize}\footnotesize
  \item \textbf{Compute stage preloads (\hyperref[sec:csf-spd]{\texttt{SPD\_k}}, \texttt{r16..r23}):} Hardware pre-fills \texttt{r55..r57} with \texttt{WorkGroupID.xyz + JOB\_OFFSET} (from \hyperref[sec:csf-regs]{\texttt{COMPUTESR::JOB\_OFFSET\_X/Y/Z}} \texttt{r34..r36}) and \texttt{r60..r62} with \texttt{LocalInvocationID.xyz} or \texttt{GlobalInvocationID.xyz} (from \texttt{WG\_SIZE} \texttt{r33} \& \texttt{JOB\_SIZE} \texttt{r37..r39}, selected by \texttt{ShaderProgram::preload}) so compute shaders can derive thread/element indices without memory loads.
  \item \textbf{Vertex stage preloads (\hyperref[sec:csf-spd]{\texttt{SPD\_0}/\texttt{SPD\_1}}, \texttt{r16..r19}):} The hardware Index Fetcher reads \hyperref[sec:csf-regs]{\texttt{IDVSSR::INDEX\_BUFFER}} (\texttt{r54:r55}) and pre-fills \texttt{r59, r60} with \texttt{VertexIndex} (\texttt{+ VERTEX\_OFFSET + GLOBAL\_ATTR\_OFF}) and \texttt{r61} with \texttt{InstanceIndex} (\texttt{+ INSTANCE\_OFFSET}) so Position and Varying shaders can pull vertex attributes directly.
  \item \textbf{Fragment stage preloads \& FP modes (\hyperref[sec:csf-spd]{\texttt{SPD\_2}}, \texttt{r20:r21}):} Hardware pre-fills \texttt{r59} with \texttt{floatBitsToUint(FragCoord.z)} and \texttt{r60} with the MSAA sample coverage mask (\texttt{gl\_SampleMaskIn[0]}), while \texttt{ShaderProgram} word~0 latches 32/64 GPR allocation and warp-wide \texttt{flush\_to\_zero\_mode}, \texttt{suppress\_na\_n}, and \texttt{suppress\_inf}.
  \item \textbf{Inline FAU \& Special Operands (\hyperref[sec:csf-tsd-fau]{\texttt{FAU\_k}}, \texttt{r8..r15}):} Any 8-bit source with \texttt{mode=2} reads zero-latency uniform \texttt{u0..u31} from \texttt{FAU\_k} (\texttt{slots64 = bits[63:56]}); \texttt{mode=3} reads \texttt{SmallConstant} (Table~1.2) or \texttt{FauSpecial} (\texttt{pc}, \texttt{lane\_id}, \texttt{lane\_width}, \texttt{tls\_ptr}, \texttt{wls\_ptr}, \texttt{blend\_descriptor\_0..7}).
\end{tightitemize}
\par\textbf{Worked Examples:}
\par\textbf{Example 1 (`-Z` face index packing):}
\codebox{\mnem{CUBEFACE2} \dst{r0}, \src{r1}, \src{r2}, \src{r3}}
With \texttt{(0.0, 0.0, -1.0)}: selects face 5 (\texttt{-Z}) -\textgreater{} \texttt{r0 = (5 \textless{}\textless{} 29) = 0xA0000000}.
\par\textbf{Example 2 (`+X` face with infinity flag):}
\codebox{\mnem{CUBEFACE2} \dst{r0}, \src{r1}, \src{r2}, \src{r3}}
With \texttt{(+Inf, 1.0, 0.0)}: selects face 0 (\texttt{+X}) and sets bit 16 (\texttt{0x00010000}) to signal an infinite major axis.
\par\textbf{Mined SPIR-V $\to$ Valhall ISA Example (Encoding the 3-bit cubemap face index in bits [31:29] --- Fragment shader, Mali-G710 (v10)):}
\codebox{\textbf{; SPIR-V (input)}\\\%color = OpImageSampleImplicitLod \%v4float \%sampled\_cube \%dir\_xyz None\\[3pt]\textbf{; Valhall ISA (malioc output)}\\CUBEFACE2 r6, r5, r4, r3          ; r6[31:29] = face (0..5), r6[16] = any\_inf\\CUBE\_SSEL r7, r3, r5, r6\\MUX.i32.bit r8, r7\textasciicircum{}, r6, u1.w1    ; merge r6[31:29] into s coordinate\\CUBE\_TSEL r3, r4\textasciicircum{}, r3\textasciicircum{}, r6\textasciicircum{}}
\texttt{CUBEFACE2} places the 3-bit face ID in bits \texttt{[31:29]} so it can be consumed directly by \texttt{CUBE\_SSEL}/\texttt{CUBE\_TSEL} and bit-blended into the texture unit's packed cube \texttt{s} coordinate via \texttt{MUX.i32.bit}.
\begin{center}\footnotesize
\begin{tabularx}{\linewidth}{@{}lllX@{}}
\toprule
Variant & Arch & Class & Opcode Fields \\
\midrule
\texttt{(none)} & v9--v14 & \texttt{R3A} & \texttt{opcode1=CUBEFACE2} \\
\bottomrule
\end{tabularx}
\end{center}
\begin{center}\footnotesize
\begin{tabularx}{0.88\linewidth}{@{}l l X@{}}
\toprule
Field Name & Bits & Type / Fixed Value \\
\midrule
\texttt{src0} & \texttt{0:7} & \texttt{SourceEncoding} \\
\texttt{src1} & \texttt{8:15} & \texttt{SourceEncoding} \\
\texttt{src2} & \texttt{16:23} & \texttt{SourceEncoding} \\
\hyperref[enum:widen_m]{\texttt{widen2}} & \texttt{24:25} & \texttt{widen\_m (def none)} \\
\hyperref[enum:widen_m]{\texttt{widen1}} & \texttt{26:27} & \texttt{widen\_m (def none)} \\
\hyperref[enum:widen_m]{\texttt{widen0}} & \texttt{28:29} & \texttt{widen\_m (def none)} \\
\textit{Reserved} & \texttt{30:33} & \texttt{zero} \\
\hyperref[enum:neg_m]{\texttt{neg2}} & \texttt{34} & \texttt{neg\_m (def none)} \\
\hyperref[enum:abs_m]{\texttt{abs2}} & \texttt{35} & \texttt{abs\_m (def none)} \\
\hyperref[enum:neg_m]{\texttt{neg1}} & \texttt{36} & \texttt{neg\_m (def none)} \\
\hyperref[enum:abs_m]{\texttt{abs1}} & \texttt{37} & \texttt{abs\_m (def none)} \\
\hyperref[enum:neg_m]{\texttt{neg0}} & \texttt{38} & \texttt{neg\_m (def none)} \\
\hyperref[enum:abs_m]{\texttt{abs0}} & \texttt{39} & \texttt{abs\_m (def none)} \\
\texttt{dest\_register\_index} & \texttt{40:45} & \texttt{uint} \\
\hyperref[enum:dest_width_m]{\texttt{dest\_width}} & \texttt{46:47} & \texttt{dest\_width\_m} \\
\texttt{opcode1} & \texttt{48:51} & \texttt{exact=CUBEFACE2} \\
\texttt{instruction\_class} & \texttt{52:56} & \texttt{exact=R3A} \\
\texttt{fau\_page\_index} & \texttt{57:58} & \texttt{uint} \\
\hyperref[enum:flow_control_m]{\texttt{flow}} & \texttt{59:62} & \texttt{flow\_control\_m} \\
\textit{Reserved} & \texttt{63} & \texttt{zero} \\
\bottomrule
\end{tabularx}
\end{center}
\textbf{Fields Explained:}
\begin{tightitemize}\footnotesize
  \item \texttt{src0}: Primary source operand (\texttt{r0..r63}, \texttt{u0..u63} uniform, \texttt{k*} FAU constant, or \texttt{t0..t2} temporary; \texttt{\textasciicircum{}} marks discard/last-use). \emph{Example:} {\scriptsize\ttfamily \mnem{FADD.f32} \dst{r0}, \src{r1}, \src{r2}} (\textbf{SPIR-V:} \texttt{OpFAdd (Operand 0)}; \textbf{GLSL:} \texttt{float c = a + b;})
  \item \texttt{src1}: Secondary source operand (register, uniform, or FAU constant; in shifts, supplies the shift count byte). \emph{Example:} {\scriptsize\ttfamily \mnem{LSHIFT.i32} \dst{r0}, \src{r1}, \src{r2}\imod{.b0}} (\textbf{SPIR-V:} \texttt{OpShiftLeftLogical (Shift)}; \textbf{GLSL:} \texttt{uint y = a \textless{}\textless{} s;})
  \item \texttt{src2}: Third source operand (FMA addend, \texttt{MUX}/\texttt{CSEL} selector, \texttt{IADDC} carry-in, or bitwise mask in fused shift-logic). \emph{Example:} {\scriptsize\ttfamily \mnem{FMA.f32} \dst{r0}, \src{r1}, \src{r2}, \src{r3}} (\textbf{SPIR-V:} \texttt{GLSL.std.450 Fma (Addend)}; \textbf{GLSL:} \texttt{float d = fma(a, b, c);})
  \item \hyperref[enum:widen_m]{\texttt{widen2}}: Stage A widening extractor on \texttt{src2} (e.g., promoting FP16 addend \texttt{h0}/\texttt{h1} to FP32 in \texttt{FMA.f32}). \emph{Example:} {\scriptsize\ttfamily \mnem{FMA.f32} \dst{r0}, \src{r1}, \src{r2}, \src{r3}\imod{.h0}} (\textbf{SPIR-V:} \texttt{OpFConvert + GLSL.std.450 Fma}; \textbf{GLSL:} \texttt{float d = fma(a, b, float(c16));})
  \item \hyperref[enum:widen_m]{\texttt{widen1}}: Stage A widening extractor on \texttt{src1} (promotes narrow byte/half/word lane before Stage B). \emph{Example:} {\scriptsize\ttfamily \mnem{IMUL.u64} \dst{[r0:r1]}, \src{[r2:r3]}\imod{.w0}, \src{[r4:r5]}\imod{.w0}} (\textbf{SPIR-V:} \texttt{OpUConvert / OpUMulExtended}; \textbf{GLSL:} \texttt{uint64\_t p = uint64\_t(a) * uint64\_t(b);})
  \item \hyperref[enum:widen_m]{\texttt{widen0}}: Stage A widening extractor on \texttt{src0} (promotes FP16 \texttt{h0}/\texttt{h1} to FP32, or sign/zero-extends \texttt{b0..b3}/\texttt{h0..h1}/\texttt{w0..w1} to 32/64-bit before \texttt{abs}/\texttt{neg}). \emph{Example:} {\scriptsize\ttfamily \mnem{IADD.s32} \dst{r0}, \src{r1}\imod{.h0}, \src{r2}} (\textbf{SPIR-V:} \texttt{OpSConvert / OpFConvert}; \textbf{GLSL:} \texttt{int c = int(int16\_t(a)) + b;})
  \item \hyperref[enum:neg_m]{\texttt{neg2}}: Stage A floating-point negation on \texttt{src2} (flips IEEE sign bit of FMA addend after swizzle/widen and \texttt{abs2}). \emph{Example:} {\scriptsize\ttfamily \mnem{FMA.f32} \dst{r0}, \src{r1}, \src{r2}, \src{r3}\imod{.neg}} (\textbf{SPIR-V:} \texttt{OpFNegate}; \textbf{GLSL:} \texttt{float d = fma(a, b, -c);})
  \item \hyperref[enum:abs_m]{\texttt{abs2}}: Stage A floating-point absolute value on \texttt{src2} (clears IEEE sign bit of FMA addend after swizzle/widen). \emph{Example:} {\scriptsize\ttfamily \mnem{FMA.f32} \dst{r0}, \src{r1}, \src{r2}, \src{r3}\imod{.abs}} (\textbf{SPIR-V:} \texttt{GLSL.std.450 FAbs}; \textbf{GLSL:} \texttt{float d = fma(a, b, abs(c));})
  \item \hyperref[enum:neg_m]{\texttt{neg1}}: Stage A floating-point negation on \texttt{src1} (flips IEEE sign bit after swizzle/widen and \texttt{abs1}). \emph{Example:} {\scriptsize\ttfamily \mnem{FMA.f32} \dst{r0}, \src{r1}, \src{r2}\imod{.neg}, \src{r3}} (\textbf{SPIR-V:} \texttt{OpFNegate}; \textbf{GLSL:} \texttt{float d = fma(a, -b, c);})
  \item \hyperref[enum:abs_m]{\texttt{abs1}}: Stage A floating-point absolute value on \texttt{src1} (clears IEEE sign bit after swizzle/widen, before \texttt{neg1}). \emph{Example:} {\scriptsize\ttfamily \mnem{FMA.f32} \dst{r0}, \src{r1}, \src{r2}\imod{.abs}, \src{r3}} (\textbf{SPIR-V:} \texttt{GLSL.std.450 FAbs}; \textbf{GLSL:} \texttt{float d = fma(a, abs(b), c);})
  \item \hyperref[enum:neg_m]{\texttt{neg0}}: Stage A floating-point negation on \texttt{src0} (flips IEEE sign bit after swizzle/widen and \texttt{abs0}). \emph{Example:} {\scriptsize\ttfamily \mnem{FADD.f32} \dst{r0}, \src{r1}\imod{.neg}, \src{r2}} (\textbf{SPIR-V:} \texttt{OpFNegate}; \textbf{GLSL:} \texttt{float c = -a + b;})
  \item \hyperref[enum:abs_m]{\texttt{abs0}}: Stage A floating-point absolute value on \texttt{src0} (clears IEEE sign bit after swizzle/widen, before \texttt{neg0}). \emph{Example:} {\scriptsize\ttfamily \mnem{FADD.f32} \dst{r0}, \src{r1}\imod{.abs}, \src{r2}} (\textbf{SPIR-V:} \texttt{GLSL.std.450 FAbs}; \textbf{GLSL:} \texttt{float c = abs(a) + b;})
  \item \texttt{dest\_register\_index}: 6-bit destination GPR index (\texttt{r0..r63}) receiving the Stage C writeback result (even-aligned pair \texttt{[rN:rN+1]} for 64-bit ops). \emph{Example:} {\scriptsize\ttfamily \mnem{FMA.f32} \dst{r0}, \src{r1}, \src{r2}, \src{r3}} (\textbf{SPIR-V:} \texttt{SSA Result ID}; \textbf{GLSL:} \texttt{float r0 = fma(r1, r2, r3);})
  \item \hyperref[enum:dest_width_m]{\texttt{dest\_width}}: Stage C destination lane write-mask (\texttt{none} = full 32-bit, \texttt{h0} = low 16-bit, \texttt{h1} = high 16-bit without bit-shifting). \emph{Example:} {\scriptsize\ttfamily \mnem{FADD.v2f16} \dst{r0}\imod{.h1}, \src{r1}, \src{r2}} (\textbf{SPIR-V:} \texttt{OpVectorInsertDynamic / WriteMask}; \textbf{GLSL:} \texttt{dst.y = a.y + b.y;})
  \item \texttt{opcode1}: Primary sub-opcode within the instruction class uniquely identifying the operation family. \emph{Example:} {\scriptsize\ttfamily \mnem{IADD.s32} \dst{r0}, \src{r1}, \src{r2}} (\textbf{SPIR-V:} \texttt{OpIAdd (operation selector)}; \textbf{GLSL:} \texttt{// Primary opcode bits})
  \item \texttt{instruction\_class}: 4-bit major instruction class selector (\texttt{bits 48:51}) partitioning the 64-bit Valhall opcode space (e.g., \texttt{FMA}, \texttt{CVT}, \texttt{SFU}, \texttt{LS}, \texttt{TEX}, \texttt{BRANCH}). \emph{Example:} {\scriptsize\ttfamily \mnem{FADD.f32} \dst{r0}, \src{r1}, \src{r2}} (\textbf{SPIR-V:} \texttt{N/A (encoding class)}; \textbf{GLSL:} \texttt{// Hardware execution pipe / decode class})
  \item \texttt{fau\_page\_index}: Selects the FAU (Fast Access Uniform) page (\texttt{0..3}) when a source operand references hardware small constants, thread IDs, or lane constants. \emph{Example:} {\scriptsize\ttfamily \mnem{FADD.f32} \dst{r0}, \src{r1}, \src{k1f}} (\textbf{SPIR-V:} \texttt{OpConstant / BuiltIn}; \textbf{GLSL:} \texttt{float y = x + 1.0;})
  \item \hyperref[enum:flow_control_m]{\texttt{flow}}: 4-bit embedded flow-control \& scoreboard wait modifier (\texttt{bits 59:62}: \texttt{none}, \texttt{wait0..wait0126}, \texttt{wait}, \texttt{reconverge}, \texttt{discard}, \texttt{end}). \emph{Example:} {\scriptsize\ttfamily \mnem{FADD.f32}\imod{.wait0.reconverge} \dst{r0}, \src{r1}, \src{r2}} (\textbf{SPIR-V:} \texttt{OpControlBarrier / OpReturn}; \textbf{GLSL:} \texttt{// Scoreboard wait + warp reconvergence})
\end{tightitemize}
\clearpage

\subsection*{CUBE\_SSEL}
\addcontentsline{toc}{subsection}{CUBE\_SSEL}
Selects signed cubemap S coordinate from \texttt{(src0=Z, src1=X, src2=CUBEFACE2)}: \texttt{face 0 -\textgreater{} -Z}, \texttt{1 -\textgreater{} +Z}, \texttt{2/3/4 -\textgreater{} +X}, \texttt{5 -\textgreater{} -X}; when \texttt{is\_inf} (bit 16 of \texttt{src2}) is set, normalizes \texttt{$\pm$Inf -\textgreater{} $\pm$1.0} and finite \texttt{-\textgreater{} $\pm$0.0}.
\par\textbf{Note:} Verified in \texttt{libGLES\_mali.so} (\texttt{0x1ecbc3c}) and \texttt{malioc} G710.
\codebox{\mnem{CUBE\_SSEL}\imod{\{.flow\}} \dst{dest}\imod{\{.dest\_width\}}, \src{src0}\imod{\{.widen0\}}\imod{\{.abs0\}}\imod{\{.neg0\}}, \src{src1}\imod{\{.widen1\}}\imod{\{.abs1\}}\imod{\{.neg1\}}, \src{src2}}
\textbf{Operands:}
\begin{tightitemize}\footnotesize
  \item \texttt{\dst{dest}} (\texttt{dest\_register\_index}, \texttt{dest\_width}): FP32 destination GPR receiving the signed horizontal (\texttt{S}) numerator for the selected face.
  \item \texttt{\src{src0}}: FP32 cubemap direction \texttt{z} component.
  \item \texttt{\src{src1}}: FP32 cubemap direction \texttt{x} component.
  \item \texttt{\src{src2}}: Face word from \texttt{CUBEFACE2} (uses bits \texttt{[31:29]}).
\end{tightitemize}
\par\textbf{Closest SPIR-V Op:} \texttt{None (lowered from OpImageSample* Cube, v9--v10)}
\par\textbf{Pseudo-GLSL (Operational Semantics):}
\codebox{uint face = src2 \textgreater{}\textgreater{} 29; // src0 = z, src1 = x\\dest = (face == 0u) ? -src0 : (face == 1u) ? src0 : (face == 5u) ? -src1 : src1;}
\textbf{Execution Context:}
\begin{tightitemize}\footnotesize
  \item \textbf{Compute stage preloads (\hyperref[sec:csf-spd]{\texttt{SPD\_k}}, \texttt{r16..r23}):} Hardware pre-fills \texttt{r55..r57} with \texttt{WorkGroupID.xyz + JOB\_OFFSET} (from \hyperref[sec:csf-regs]{\texttt{COMPUTESR::JOB\_OFFSET\_X/Y/Z}} \texttt{r34..r36}) and \texttt{r60..r62} with \texttt{LocalInvocationID.xyz} or \texttt{GlobalInvocationID.xyz} (from \texttt{WG\_SIZE} \texttt{r33} \& \texttt{JOB\_SIZE} \texttt{r37..r39}, selected by \texttt{ShaderProgram::preload}) so compute shaders can derive thread/element indices without memory loads.
  \item \textbf{Vertex stage preloads (\hyperref[sec:csf-spd]{\texttt{SPD\_0}/\texttt{SPD\_1}}, \texttt{r16..r19}):} The hardware Index Fetcher reads \hyperref[sec:csf-regs]{\texttt{IDVSSR::INDEX\_BUFFER}} (\texttt{r54:r55}) and pre-fills \texttt{r59, r60} with \texttt{VertexIndex} (\texttt{+ VERTEX\_OFFSET + GLOBAL\_ATTR\_OFF}) and \texttt{r61} with \texttt{InstanceIndex} (\texttt{+ INSTANCE\_OFFSET}) so Position and Varying shaders can pull vertex attributes directly.
  \item \textbf{Fragment stage preloads \& FP modes (\hyperref[sec:csf-spd]{\texttt{SPD\_2}}, \texttt{r20:r21}):} Hardware pre-fills \texttt{r59} with \texttt{floatBitsToUint(FragCoord.z)} and \texttt{r60} with the MSAA sample coverage mask (\texttt{gl\_SampleMaskIn[0]}), while \texttt{ShaderProgram} word~0 latches 32/64 GPR allocation and warp-wide \texttt{flush\_to\_zero\_mode}, \texttt{suppress\_na\_n}, and \texttt{suppress\_inf}.
  \item \textbf{Inline FAU \& Special Operands (\hyperref[sec:csf-tsd-fau]{\texttt{FAU\_k}}, \texttt{r8..r15}):} Any 8-bit source with \texttt{mode=2} reads zero-latency uniform \texttt{u0..u31} from \texttt{FAU\_k} (\texttt{slots64 = bits[63:56]}); \texttt{mode=3} reads \texttt{SmallConstant} (Table~1.2) or \texttt{FauSpecial} (\texttt{pc}, \texttt{lane\_id}, \texttt{lane\_width}, \texttt{tls\_ptr}, \texttt{wls\_ptr}, \texttt{blend\_descriptor\_0..7}).
\end{tightitemize}
\par\textbf{Worked Examples:}
\par\textbf{Example 1 (`+X` face horizontal coordinate (`-Z`)):}
\codebox{\mnem{CUBE\_SSEL} \dst{r0}, \src{r3}, \src{r1}, \src{r4}}
Operands \texttt{(Z, X, face)}: with \texttt{face = 0} (\texttt{+X}, \texttt{r4 = 0x00000000}), selects and negates \texttt{Z} -\textgreater{} \texttt{r0 = -r3}.
\par\textbf{Example 2 (`-Z` face horizontal coordinate (`-X`)):}
\codebox{\mnem{CUBE\_SSEL} \dst{r0}, \src{r3}, \src{r1}, \src{r4}}
With \texttt{face = 5} (\texttt{-Z}, \texttt{r4 = 0xA0000000}): selects and negates \texttt{X} -\textgreater{} \texttt{r0 = -r1}.
\par\textbf{Mined SPIR-V $\to$ Valhall ISA Example (Selecting the signed S numerator for the active cube face --- Fragment shader, Mali-G710 (v10)):}
\codebox{\textbf{; SPIR-V (input)}\\\%color = OpImageSampleImplicitLod \%v4float \%sampled\_cube \%dir\_xyz None\\[3pt]\textbf{; Valhall ISA (malioc output)}\\CUBEFACE2 r6, r5, r4, r3          ; face from (x=r5, y=r4, z=r3)\\CUBE\_SSEL r7, r3, r5, r6          ; src0=z, src1=x, src2=face\\FMA.f32.clamp\_0\_1 r7, r5, r7\textasciicircum{}, k0fr2f.h1}
Notice the operand order: \texttt{CUBE\_SSEL} takes \texttt{z} in \texttt{src0}, \texttt{x} in \texttt{src1}, and the \texttt{CUBEFACE2} output in \texttt{src2}, selecting and negating the horizontal face coordinate (\texttt{-z} for \texttt{+X}, \texttt{+z} for \texttt{-X}, \texttt{+x} for \texttt{$\pm$Y/+Z}, \texttt{-x} for \texttt{-Z}).
\begin{center}\footnotesize
\begin{tabularx}{\linewidth}{@{}lllX@{}}
\toprule
Variant & Arch & Class & Opcode Fields \\
\midrule
\texttt{(none)} & v9--v14 & \texttt{R3} & \texttt{opcode1=CUBE} \\
\bottomrule
\end{tabularx}
\end{center}
\begin{center}\footnotesize
\begin{tabularx}{0.88\linewidth}{@{}l l X@{}}
\toprule
Field Name & Bits & Type / Fixed Value \\
\midrule
\texttt{src0} & \texttt{0:7} & \texttt{SourceEncoding} \\
\texttt{src1} & \texttt{8:15} & \texttt{SourceEncoding} \\
\texttt{src2} & \texttt{16:23} & \texttt{SourceEncoding} \\
\texttt{coord\_select} & \texttt{24} & \texttt{exact=s} \\
\textit{Reserved} & \texttt{25} & \texttt{zero} \\
\hyperref[enum:widen_m]{\texttt{widen1}} & \texttt{26:27} & \texttt{widen\_m (def none)} \\
\hyperref[enum:widen_m]{\texttt{widen0}} & \texttt{28:29} & \texttt{widen\_m (def none)} \\
\textit{Reserved} & \texttt{30:35} & \texttt{zero} \\
\hyperref[enum:neg_m]{\texttt{neg1}} & \texttt{36} & \texttt{neg\_m (def none)} \\
\hyperref[enum:abs_m]{\texttt{abs1}} & \texttt{37} & \texttt{abs\_m (def none)} \\
\hyperref[enum:neg_m]{\texttt{neg0}} & \texttt{38} & \texttt{neg\_m (def none)} \\
\hyperref[enum:abs_m]{\texttt{abs0}} & \texttt{39} & \texttt{abs\_m (def none)} \\
\texttt{dest\_register\_index} & \texttt{40:45} & \texttt{uint} \\
\hyperref[enum:dest_width_m]{\texttt{dest\_width}} & \texttt{46:47} & \texttt{dest\_width\_m} \\
\texttt{opcode1} & \texttt{48:51} & \texttt{exact=CUBE} \\
\texttt{instruction\_class} & \texttt{52:56} & \texttt{exact=R3} \\
\texttt{fau\_page\_index} & \texttt{57:58} & \texttt{uint} \\
\hyperref[enum:flow_control_m]{\texttt{flow}} & \texttt{59:62} & \texttt{flow\_control\_m} \\
\textit{Reserved} & \texttt{63} & \texttt{zero} \\
\bottomrule
\end{tabularx}
\end{center}
\textbf{Fields Explained:}
\begin{tightitemize}\footnotesize
  \item \texttt{src0}: Primary source operand (\texttt{r0..r63}, \texttt{u0..u63} uniform, \texttt{k*} FAU constant, or \texttt{t0..t2} temporary; \texttt{\textasciicircum{}} marks discard/last-use). \emph{Example:} {\scriptsize\ttfamily \mnem{FADD.f32} \dst{r0}, \src{r1}, \src{r2}} (\textbf{SPIR-V:} \texttt{OpFAdd (Operand 0)}; \textbf{GLSL:} \texttt{float c = a + b;})
  \item \texttt{src1}: Secondary source operand (register, uniform, or FAU constant; in shifts, supplies the shift count byte). \emph{Example:} {\scriptsize\ttfamily \mnem{LSHIFT.i32} \dst{r0}, \src{r1}, \src{r2}\imod{.b0}} (\textbf{SPIR-V:} \texttt{OpShiftLeftLogical (Shift)}; \textbf{GLSL:} \texttt{uint y = a \textless{}\textless{} s;})
  \item \texttt{src2}: Third source operand (FMA addend, \texttt{MUX}/\texttt{CSEL} selector, \texttt{IADDC} carry-in, or bitwise mask in fused shift-logic). \emph{Example:} {\scriptsize\ttfamily \mnem{FMA.f32} \dst{r0}, \src{r1}, \src{r2}, \src{r3}} (\textbf{SPIR-V:} \texttt{GLSL.std.450 Fma (Addend)}; \textbf{GLSL:} \texttt{float d = fma(a, b, c);})
  \item \texttt{coord\_select}: Selects which coordinate pair is read from staging registers. \emph{Example:} {\scriptsize\ttfamily \mnem{TEX\_GRADIENT}\imod{.tex2d} \dst{@r0:r1}, \src{@r4:r5}} (\textbf{SPIR-V:} \texttt{OpImageQueryLod}; \textbf{GLSL:} \texttt{vec2 lod = textureQueryLod(tex, uv);})
  \item \hyperref[enum:widen_m]{\texttt{widen1}}: Stage A widening extractor on \texttt{src1} (promotes narrow byte/half/word lane before Stage B). \emph{Example:} {\scriptsize\ttfamily \mnem{IMUL.u64} \dst{[r0:r1]}, \src{[r2:r3]}\imod{.w0}, \src{[r4:r5]}\imod{.w0}} (\textbf{SPIR-V:} \texttt{OpUConvert / OpUMulExtended}; \textbf{GLSL:} \texttt{uint64\_t p = uint64\_t(a) * uint64\_t(b);})
  \item \hyperref[enum:widen_m]{\texttt{widen0}}: Stage A widening extractor on \texttt{src0} (promotes FP16 \texttt{h0}/\texttt{h1} to FP32, or sign/zero-extends \texttt{b0..b3}/\texttt{h0..h1}/\texttt{w0..w1} to 32/64-bit before \texttt{abs}/\texttt{neg}). \emph{Example:} {\scriptsize\ttfamily \mnem{IADD.s32} \dst{r0}, \src{r1}\imod{.h0}, \src{r2}} (\textbf{SPIR-V:} \texttt{OpSConvert / OpFConvert}; \textbf{GLSL:} \texttt{int c = int(int16\_t(a)) + b;})
  \item \hyperref[enum:neg_m]{\texttt{neg1}}: Stage A floating-point negation on \texttt{src1} (flips IEEE sign bit after swizzle/widen and \texttt{abs1}). \emph{Example:} {\scriptsize\ttfamily \mnem{FMA.f32} \dst{r0}, \src{r1}, \src{r2}\imod{.neg}, \src{r3}} (\textbf{SPIR-V:} \texttt{OpFNegate}; \textbf{GLSL:} \texttt{float d = fma(a, -b, c);})
  \item \hyperref[enum:abs_m]{\texttt{abs1}}: Stage A floating-point absolute value on \texttt{src1} (clears IEEE sign bit after swizzle/widen, before \texttt{neg1}). \emph{Example:} {\scriptsize\ttfamily \mnem{FMA.f32} \dst{r0}, \src{r1}, \src{r2}\imod{.abs}, \src{r3}} (\textbf{SPIR-V:} \texttt{GLSL.std.450 FAbs}; \textbf{GLSL:} \texttt{float d = fma(a, abs(b), c);})
  \item \hyperref[enum:neg_m]{\texttt{neg0}}: Stage A floating-point negation on \texttt{src0} (flips IEEE sign bit after swizzle/widen and \texttt{abs0}). \emph{Example:} {\scriptsize\ttfamily \mnem{FADD.f32} \dst{r0}, \src{r1}\imod{.neg}, \src{r2}} (\textbf{SPIR-V:} \texttt{OpFNegate}; \textbf{GLSL:} \texttt{float c = -a + b;})
  \item \hyperref[enum:abs_m]{\texttt{abs0}}: Stage A floating-point absolute value on \texttt{src0} (clears IEEE sign bit after swizzle/widen, before \texttt{neg0}). \emph{Example:} {\scriptsize\ttfamily \mnem{FADD.f32} \dst{r0}, \src{r1}\imod{.abs}, \src{r2}} (\textbf{SPIR-V:} \texttt{GLSL.std.450 FAbs}; \textbf{GLSL:} \texttt{float c = abs(a) + b;})
  \item \texttt{dest\_register\_index}: 6-bit destination GPR index (\texttt{r0..r63}) receiving the Stage C writeback result (even-aligned pair \texttt{[rN:rN+1]} for 64-bit ops). \emph{Example:} {\scriptsize\ttfamily \mnem{FMA.f32} \dst{r0}, \src{r1}, \src{r2}, \src{r3}} (\textbf{SPIR-V:} \texttt{SSA Result ID}; \textbf{GLSL:} \texttt{float r0 = fma(r1, r2, r3);})
  \item \hyperref[enum:dest_width_m]{\texttt{dest\_width}}: Stage C destination lane write-mask (\texttt{none} = full 32-bit, \texttt{h0} = low 16-bit, \texttt{h1} = high 16-bit without bit-shifting). \emph{Example:} {\scriptsize\ttfamily \mnem{FADD.v2f16} \dst{r0}\imod{.h1}, \src{r1}, \src{r2}} (\textbf{SPIR-V:} \texttt{OpVectorInsertDynamic / WriteMask}; \textbf{GLSL:} \texttt{dst.y = a.y + b.y;})
  \item \texttt{opcode1}: Primary sub-opcode within the instruction class uniquely identifying the operation family. \emph{Example:} {\scriptsize\ttfamily \mnem{IADD.s32} \dst{r0}, \src{r1}, \src{r2}} (\textbf{SPIR-V:} \texttt{OpIAdd (operation selector)}; \textbf{GLSL:} \texttt{// Primary opcode bits})
  \item \texttt{instruction\_class}: 4-bit major instruction class selector (\texttt{bits 48:51}) partitioning the 64-bit Valhall opcode space (e.g., \texttt{FMA}, \texttt{CVT}, \texttt{SFU}, \texttt{LS}, \texttt{TEX}, \texttt{BRANCH}). \emph{Example:} {\scriptsize\ttfamily \mnem{FADD.f32} \dst{r0}, \src{r1}, \src{r2}} (\textbf{SPIR-V:} \texttt{N/A (encoding class)}; \textbf{GLSL:} \texttt{// Hardware execution pipe / decode class})
  \item \texttt{fau\_page\_index}: Selects the FAU (Fast Access Uniform) page (\texttt{0..3}) when a source operand references hardware small constants, thread IDs, or lane constants. \emph{Example:} {\scriptsize\ttfamily \mnem{FADD.f32} \dst{r0}, \src{r1}, \src{k1f}} (\textbf{SPIR-V:} \texttt{OpConstant / BuiltIn}; \textbf{GLSL:} \texttt{float y = x + 1.0;})
  \item \hyperref[enum:flow_control_m]{\texttt{flow}}: 4-bit embedded flow-control \& scoreboard wait modifier (\texttt{bits 59:62}: \texttt{none}, \texttt{wait0..wait0126}, \texttt{wait}, \texttt{reconverge}, \texttt{discard}, \texttt{end}). \emph{Example:} {\scriptsize\ttfamily \mnem{FADD.f32}\imod{.wait0.reconverge} \dst{r0}, \src{r1}, \src{r2}} (\textbf{SPIR-V:} \texttt{OpControlBarrier / OpReturn}; \textbf{GLSL:} \texttt{// Scoreboard wait + warp reconvergence})
\end{tightitemize}
\clearpage

\subsection*{CUBE\_TSEL}
\addcontentsline{toc}{subsection}{CUBE\_TSEL}
Selects signed cubemap T coordinate from \texttt{(src0=Y, src1=Z, src2=CUBEFACE2)}: \texttt{face 2 -\textgreater{} +Z}, \texttt{3 -\textgreater{} -Z}, \texttt{0/1/4/5 -\textgreater{} -Y}; when \texttt{is\_inf} (bit 16 of \texttt{src2}) is set, normalizes \texttt{$\pm$Inf -\textgreater{} $\pm$1.0} and finite \texttt{-\textgreater{} $\pm$0.0}.
\par\textbf{Note:} Verified in \texttt{libGLES\_mali.so} (\texttt{0x1ecbc98}) and \texttt{malioc} G710.
\codebox{\mnem{CUBE\_TSEL}\imod{\{.flow\}} \dst{dest}\imod{\{.dest\_width\}}, \src{src0}\imod{\{.widen0\}}\imod{\{.abs0\}}\imod{\{.neg0\}}, \src{src1}\imod{\{.widen1\}}\imod{\{.abs1\}}\imod{\{.neg1\}}, \src{src2}}
\textbf{Operands:}
\begin{tightitemize}\footnotesize
  \item \texttt{\dst{dest}} (\texttt{dest\_register\_index}, \texttt{dest\_width}): FP32 destination GPR receiving the signed vertical (\texttt{T}) numerator for the selected face.
  \item \texttt{\src{src0}}: FP32 cubemap direction \texttt{y} component.
  \item \texttt{\src{src1}}: FP32 cubemap direction \texttt{z} component.
  \item \texttt{\src{src2}}: Face word from \texttt{CUBEFACE2} (uses bits \texttt{[31:29]}).
\end{tightitemize}
\par\textbf{Closest SPIR-V Op:} \texttt{None (lowered from OpImageSample* Cube, v9--v10)}
\par\textbf{Pseudo-GLSL (Operational Semantics):}
\codebox{uint face = src2 \textgreater{}\textgreater{} 29; // src0 = y, src1 = z\\dest = (face == 2u) ? src1 : (face == 3u) ? -src1 : -src0;}
\textbf{Execution Context:}
\begin{tightitemize}\footnotesize
  \item \textbf{Compute stage preloads (\hyperref[sec:csf-spd]{\texttt{SPD\_k}}, \texttt{r16..r23}):} Hardware pre-fills \texttt{r55..r57} with \texttt{WorkGroupID.xyz + JOB\_OFFSET} (from \hyperref[sec:csf-regs]{\texttt{COMPUTESR::JOB\_OFFSET\_X/Y/Z}} \texttt{r34..r36}) and \texttt{r60..r62} with \texttt{LocalInvocationID.xyz} or \texttt{GlobalInvocationID.xyz} (from \texttt{WG\_SIZE} \texttt{r33} \& \texttt{JOB\_SIZE} \texttt{r37..r39}, selected by \texttt{ShaderProgram::preload}) so compute shaders can derive thread/element indices without memory loads.
  \item \textbf{Vertex stage preloads (\hyperref[sec:csf-spd]{\texttt{SPD\_0}/\texttt{SPD\_1}}, \texttt{r16..r19}):} The hardware Index Fetcher reads \hyperref[sec:csf-regs]{\texttt{IDVSSR::INDEX\_BUFFER}} (\texttt{r54:r55}) and pre-fills \texttt{r59, r60} with \texttt{VertexIndex} (\texttt{+ VERTEX\_OFFSET + GLOBAL\_ATTR\_OFF}) and \texttt{r61} with \texttt{InstanceIndex} (\texttt{+ INSTANCE\_OFFSET}) so Position and Varying shaders can pull vertex attributes directly.
  \item \textbf{Fragment stage preloads \& FP modes (\hyperref[sec:csf-spd]{\texttt{SPD\_2}}, \texttt{r20:r21}):} Hardware pre-fills \texttt{r59} with \texttt{floatBitsToUint(FragCoord.z)} and \texttt{r60} with the MSAA sample coverage mask (\texttt{gl\_SampleMaskIn[0]}), while \texttt{ShaderProgram} word~0 latches 32/64 GPR allocation and warp-wide \texttt{flush\_to\_zero\_mode}, \texttt{suppress\_na\_n}, and \texttt{suppress\_inf}.
  \item \textbf{Inline FAU \& Special Operands (\hyperref[sec:csf-tsd-fau]{\texttt{FAU\_k}}, \texttt{r8..r15}):} Any 8-bit source with \texttt{mode=2} reads zero-latency uniform \texttt{u0..u31} from \texttt{FAU\_k} (\texttt{slots64 = bits[63:56]}); \texttt{mode=3} reads \texttt{SmallConstant} (Table~1.2) or \texttt{FauSpecial} (\texttt{pc}, \texttt{lane\_id}, \texttt{lane\_width}, \texttt{tls\_ptr}, \texttt{wls\_ptr}, \texttt{blend\_descriptor\_0..7}).
\end{tightitemize}
\par\textbf{Worked Examples:}
\par\textbf{Example 1 (`+Y` face vertical coordinate (`+Z`)):}
\codebox{\mnem{CUBE\_TSEL} \dst{r0}, \src{r2}, \src{r3}, \src{r4}}
Operands \texttt{(Y, Z, face)}: with \texttt{face = 2} (\texttt{+Y}, \texttt{r4 = 0x40000000}), returns \texttt{+Z} (\texttt{r0 = r3}).
\par\textbf{Example 2 (`+X` face vertical coordinate (`-Y`)):}
\codebox{\mnem{CUBE\_TSEL} \dst{r0}, \src{r2}, \src{r3}, \src{r4}}
With \texttt{face = 0} (\texttt{+X}, \texttt{r4 = 0x00000000}): returns \texttt{-Y} (\texttt{r0 = -r2}).
\par\textbf{Mined SPIR-V $\to$ Valhall ISA Example (Selecting the signed T numerator for the active cube face --- Fragment shader, Mali-G710 (v10)):}
\codebox{\textbf{; SPIR-V (input)}\\\%color = OpImageSampleImplicitLod \%v4float \%sampled\_cube \%dir\_xyz None\\[3pt]\textbf{; Valhall ISA (malioc output)}\\CUBE\_TSEL r3, r4\textasciicircum{}, r3\textasciicircum{}, r6\textasciicircum{}       ; src0=y, src1=z, src2=face\\FMA.f32.clamp\_0\_1 r9, r5\textasciicircum{}, r3\textasciicircum{}, k0fr2f.h1\\TEX\_SINGLE.slot0.reconverge.skip.cube.dst32.rgba.wide.computed @r4:r5:r6:r7, @r8:r9, u2}
\texttt{CUBE\_TSEL} takes \texttt{y} in \texttt{src0}, \texttt{z} in \texttt{src1}, and the \texttt{CUBEFACE2} face word in \texttt{src2}, outputting \texttt{+z} for \texttt{+Y}, \texttt{-z} for \texttt{-Y}, and \texttt{-y} for \texttt{$\pm$X/$\pm$Z}.
\begin{center}\footnotesize
\begin{tabularx}{\linewidth}{@{}lllX@{}}
\toprule
Variant & Arch & Class & Opcode Fields \\
\midrule
\texttt{(none)} & v9--v14 & \texttt{R3} & \texttt{opcode1=CUBE} \\
\bottomrule
\end{tabularx}
\end{center}
\begin{center}\footnotesize
\begin{tabularx}{0.88\linewidth}{@{}l l X@{}}
\toprule
Field Name & Bits & Type / Fixed Value \\
\midrule
\texttt{src0} & \texttt{0:7} & \texttt{SourceEncoding} \\
\texttt{src1} & \texttt{8:15} & \texttt{SourceEncoding} \\
\texttt{src2} & \texttt{16:23} & \texttt{SourceEncoding} \\
\texttt{coord\_select} & \texttt{24} & \texttt{exact=t} \\
\textit{Reserved} & \texttt{25} & \texttt{zero} \\
\hyperref[enum:widen_m]{\texttt{widen1}} & \texttt{26:27} & \texttt{widen\_m (def none)} \\
\hyperref[enum:widen_m]{\texttt{widen0}} & \texttt{28:29} & \texttt{widen\_m (def none)} \\
\textit{Reserved} & \texttt{30:35} & \texttt{zero} \\
\hyperref[enum:neg_m]{\texttt{neg1}} & \texttt{36} & \texttt{neg\_m (def none)} \\
\hyperref[enum:abs_m]{\texttt{abs1}} & \texttt{37} & \texttt{abs\_m (def none)} \\
\hyperref[enum:neg_m]{\texttt{neg0}} & \texttt{38} & \texttt{neg\_m (def none)} \\
\hyperref[enum:abs_m]{\texttt{abs0}} & \texttt{39} & \texttt{abs\_m (def none)} \\
\texttt{dest\_register\_index} & \texttt{40:45} & \texttt{uint} \\
\hyperref[enum:dest_width_m]{\texttt{dest\_width}} & \texttt{46:47} & \texttt{dest\_width\_m} \\
\texttt{opcode1} & \texttt{48:51} & \texttt{exact=CUBE} \\
\texttt{instruction\_class} & \texttt{52:56} & \texttt{exact=R3} \\
\texttt{fau\_page\_index} & \texttt{57:58} & \texttt{uint} \\
\hyperref[enum:flow_control_m]{\texttt{flow}} & \texttt{59:62} & \texttt{flow\_control\_m} \\
\textit{Reserved} & \texttt{63} & \texttt{zero} \\
\bottomrule
\end{tabularx}
\end{center}
\textbf{Fields Explained:}
\begin{tightitemize}\footnotesize
  \item \texttt{src0}: Primary source operand (\texttt{r0..r63}, \texttt{u0..u63} uniform, \texttt{k*} FAU constant, or \texttt{t0..t2} temporary; \texttt{\textasciicircum{}} marks discard/last-use). \emph{Example:} {\scriptsize\ttfamily \mnem{FADD.f32} \dst{r0}, \src{r1}, \src{r2}} (\textbf{SPIR-V:} \texttt{OpFAdd (Operand 0)}; \textbf{GLSL:} \texttt{float c = a + b;})
  \item \texttt{src1}: Secondary source operand (register, uniform, or FAU constant; in shifts, supplies the shift count byte). \emph{Example:} {\scriptsize\ttfamily \mnem{LSHIFT.i32} \dst{r0}, \src{r1}, \src{r2}\imod{.b0}} (\textbf{SPIR-V:} \texttt{OpShiftLeftLogical (Shift)}; \textbf{GLSL:} \texttt{uint y = a \textless{}\textless{} s;})
  \item \texttt{src2}: Third source operand (FMA addend, \texttt{MUX}/\texttt{CSEL} selector, \texttt{IADDC} carry-in, or bitwise mask in fused shift-logic). \emph{Example:} {\scriptsize\ttfamily \mnem{FMA.f32} \dst{r0}, \src{r1}, \src{r2}, \src{r3}} (\textbf{SPIR-V:} \texttt{GLSL.std.450 Fma (Addend)}; \textbf{GLSL:} \texttt{float d = fma(a, b, c);})
  \item \texttt{coord\_select}: Selects which coordinate pair is read from staging registers. \emph{Example:} {\scriptsize\ttfamily \mnem{TEX\_GRADIENT}\imod{.tex2d} \dst{@r0:r1}, \src{@r4:r5}} (\textbf{SPIR-V:} \texttt{OpImageQueryLod}; \textbf{GLSL:} \texttt{vec2 lod = textureQueryLod(tex, uv);})
  \item \hyperref[enum:widen_m]{\texttt{widen1}}: Stage A widening extractor on \texttt{src1} (promotes narrow byte/half/word lane before Stage B). \emph{Example:} {\scriptsize\ttfamily \mnem{IMUL.u64} \dst{[r0:r1]}, \src{[r2:r3]}\imod{.w0}, \src{[r4:r5]}\imod{.w0}} (\textbf{SPIR-V:} \texttt{OpUConvert / OpUMulExtended}; \textbf{GLSL:} \texttt{uint64\_t p = uint64\_t(a) * uint64\_t(b);})
  \item \hyperref[enum:widen_m]{\texttt{widen0}}: Stage A widening extractor on \texttt{src0} (promotes FP16 \texttt{h0}/\texttt{h1} to FP32, or sign/zero-extends \texttt{b0..b3}/\texttt{h0..h1}/\texttt{w0..w1} to 32/64-bit before \texttt{abs}/\texttt{neg}). \emph{Example:} {\scriptsize\ttfamily \mnem{IADD.s32} \dst{r0}, \src{r1}\imod{.h0}, \src{r2}} (\textbf{SPIR-V:} \texttt{OpSConvert / OpFConvert}; \textbf{GLSL:} \texttt{int c = int(int16\_t(a)) + b;})
  \item \hyperref[enum:neg_m]{\texttt{neg1}}: Stage A floating-point negation on \texttt{src1} (flips IEEE sign bit after swizzle/widen and \texttt{abs1}). \emph{Example:} {\scriptsize\ttfamily \mnem{FMA.f32} \dst{r0}, \src{r1}, \src{r2}\imod{.neg}, \src{r3}} (\textbf{SPIR-V:} \texttt{OpFNegate}; \textbf{GLSL:} \texttt{float d = fma(a, -b, c);})
  \item \hyperref[enum:abs_m]{\texttt{abs1}}: Stage A floating-point absolute value on \texttt{src1} (clears IEEE sign bit after swizzle/widen, before \texttt{neg1}). \emph{Example:} {\scriptsize\ttfamily \mnem{FMA.f32} \dst{r0}, \src{r1}, \src{r2}\imod{.abs}, \src{r3}} (\textbf{SPIR-V:} \texttt{GLSL.std.450 FAbs}; \textbf{GLSL:} \texttt{float d = fma(a, abs(b), c);})
  \item \hyperref[enum:neg_m]{\texttt{neg0}}: Stage A floating-point negation on \texttt{src0} (flips IEEE sign bit after swizzle/widen and \texttt{abs0}). \emph{Example:} {\scriptsize\ttfamily \mnem{FADD.f32} \dst{r0}, \src{r1}\imod{.neg}, \src{r2}} (\textbf{SPIR-V:} \texttt{OpFNegate}; \textbf{GLSL:} \texttt{float c = -a + b;})
  \item \hyperref[enum:abs_m]{\texttt{abs0}}: Stage A floating-point absolute value on \texttt{src0} (clears IEEE sign bit after swizzle/widen, before \texttt{neg0}). \emph{Example:} {\scriptsize\ttfamily \mnem{FADD.f32} \dst{r0}, \src{r1}\imod{.abs}, \src{r2}} (\textbf{SPIR-V:} \texttt{GLSL.std.450 FAbs}; \textbf{GLSL:} \texttt{float c = abs(a) + b;})
  \item \texttt{dest\_register\_index}: 6-bit destination GPR index (\texttt{r0..r63}) receiving the Stage C writeback result (even-aligned pair \texttt{[rN:rN+1]} for 64-bit ops). \emph{Example:} {\scriptsize\ttfamily \mnem{FMA.f32} \dst{r0}, \src{r1}, \src{r2}, \src{r3}} (\textbf{SPIR-V:} \texttt{SSA Result ID}; \textbf{GLSL:} \texttt{float r0 = fma(r1, r2, r3);})
  \item \hyperref[enum:dest_width_m]{\texttt{dest\_width}}: Stage C destination lane write-mask (\texttt{none} = full 32-bit, \texttt{h0} = low 16-bit, \texttt{h1} = high 16-bit without bit-shifting). \emph{Example:} {\scriptsize\ttfamily \mnem{FADD.v2f16} \dst{r0}\imod{.h1}, \src{r1}, \src{r2}} (\textbf{SPIR-V:} \texttt{OpVectorInsertDynamic / WriteMask}; \textbf{GLSL:} \texttt{dst.y = a.y + b.y;})
  \item \texttt{opcode1}: Primary sub-opcode within the instruction class uniquely identifying the operation family. \emph{Example:} {\scriptsize\ttfamily \mnem{IADD.s32} \dst{r0}, \src{r1}, \src{r2}} (\textbf{SPIR-V:} \texttt{OpIAdd (operation selector)}; \textbf{GLSL:} \texttt{// Primary opcode bits})
  \item \texttt{instruction\_class}: 4-bit major instruction class selector (\texttt{bits 48:51}) partitioning the 64-bit Valhall opcode space (e.g., \texttt{FMA}, \texttt{CVT}, \texttt{SFU}, \texttt{LS}, \texttt{TEX}, \texttt{BRANCH}). \emph{Example:} {\scriptsize\ttfamily \mnem{FADD.f32} \dst{r0}, \src{r1}, \src{r2}} (\textbf{SPIR-V:} \texttt{N/A (encoding class)}; \textbf{GLSL:} \texttt{// Hardware execution pipe / decode class})
  \item \texttt{fau\_page\_index}: Selects the FAU (Fast Access Uniform) page (\texttt{0..3}) when a source operand references hardware small constants, thread IDs, or lane constants. \emph{Example:} {\scriptsize\ttfamily \mnem{FADD.f32} \dst{r0}, \src{r1}, \src{k1f}} (\textbf{SPIR-V:} \texttt{OpConstant / BuiltIn}; \textbf{GLSL:} \texttt{float y = x + 1.0;})
  \item \hyperref[enum:flow_control_m]{\texttt{flow}}: 4-bit embedded flow-control \& scoreboard wait modifier (\texttt{bits 59:62}: \texttt{none}, \texttt{wait0..wait0126}, \texttt{wait}, \texttt{reconverge}, \texttt{discard}, \texttt{end}). \emph{Example:} {\scriptsize\ttfamily \mnem{FADD.f32}\imod{.wait0.reconverge} \dst{r0}, \src{r1}, \src{r2}} (\textbf{SPIR-V:} \texttt{OpControlBarrier / OpReturn}; \textbf{GLSL:} \texttt{// Scoreboard wait + warp reconvergence})
\end{tightitemize}
\clearpage

\section{Compare, Select \& Mux (9 mnemonics)}
\subsection*{CSEL}
\addcontentsline{toc}{subsection}{CSEL}
Compares \texttt{src0[i]} and \texttt{src1[i]} using predicate \texttt{cmpf} independently for each SIMD element $i \in \{0, \dots, N-1\}$ ($N=1$ for \texttt{f32}/\texttt{s32}/\texttt{u32}, $N=2$ for \texttt{v2f16}/\texttt{v2s16}/\texttt{v2u16}) and writes \texttt{src2[i]} to \texttt{dest[i]} if true, or \texttt{src3[i]} if false.
\par\textbf{Notes:}
\begin{tightitemize}
  \item \texttt{src2}/\texttt{src3} are raw bit payloads (\texttt{VNIN}, \texttt{ops.rs:859}).
  \item Exact match verified in \texttt{hw\_tests.rs:1391-1430}.
  \item \textbf{Source references:} \texttt{kraid/ops.rs:851-894}, \texttt{encode\_v9.rs:1026-1050}, \texttt{nir.rs:1222-1239}, \texttt{hw\_tests.rs:1391-1430}.
\end{tightitemize}
\codebox{\mnem{CSEL}\imod{\{.flow\}}\imod{\{.cmpf\}} \dst{dest}\imod{\{.dest\_width\}}, \src{src0}, \src{src1}, \src{src2}, \src{src3}}
\textbf{Operands:}
\begin{tightitemize}\footnotesize
  \item \texttt{\dst{dest}} (\texttt{dest\_register\_index}, \texttt{dest\_width}): Destination GPR receiving \texttt{src2} when \texttt{cmpf(src0, src1)} is true, else \texttt{src3}.
  \item \texttt{\src{src0}}: Left-hand comparison operand.
  \item \texttt{\src{src1}}: Right-hand comparison operand.
  \item \texttt{\src{src2}}: Value selected when the comparison is true.
  \item \texttt{\src{src3}}: Value selected when the comparison is false.
\end{tightitemize}
\par\textbf{Closest SPIR-V Op:} \texttt{OpSelect + OpFOrd*/OpS*/OpU*}
\par\textbf{Pseudo-GLSL (Operational Semantics):}
\codebox{for (int i = 0; i \textless{} N; ++i)\\    dest[i] = cmpf(src0[i], src1[i]) ? src2[i] : src3[i];}
\textbf{Execution Context:}
\begin{tightitemize}\footnotesize
  \item \textbf{Compute stage preloads (\hyperref[sec:csf-spd]{\texttt{SPD\_k}}, \texttt{r16..r23}):} Hardware pre-fills \texttt{r55..r57} with \texttt{WorkGroupID.xyz + JOB\_OFFSET} (from \hyperref[sec:csf-regs]{\texttt{COMPUTESR::JOB\_OFFSET\_X/Y/Z}} \texttt{r34..r36}) and \texttt{r60..r62} with \texttt{LocalInvocationID.xyz} or \texttt{GlobalInvocationID.xyz} (from \texttt{WG\_SIZE} \texttt{r33} \& \texttt{JOB\_SIZE} \texttt{r37..r39}, selected by \texttt{ShaderProgram::preload}) so compute shaders can derive thread/element indices without memory loads.
  \item \textbf{Vertex stage preloads (\hyperref[sec:csf-spd]{\texttt{SPD\_0}/\texttt{SPD\_1}}, \texttt{r16..r19}):} The hardware Index Fetcher reads \hyperref[sec:csf-regs]{\texttt{IDVSSR::INDEX\_BUFFER}} (\texttt{r54:r55}) and pre-fills \texttt{r59, r60} with \texttt{VertexIndex} (\texttt{+ VERTEX\_OFFSET + GLOBAL\_ATTR\_OFF}) and \texttt{r61} with \texttt{InstanceIndex} (\texttt{+ INSTANCE\_OFFSET}) so Position and Varying shaders can pull vertex attributes directly.
  \item \textbf{Inline FAU \& Special Operands (\hyperref[sec:csf-tsd-fau]{\texttt{FAU\_k}}, \texttt{r8..r15}):} Any 8-bit source with \texttt{mode=2} reads zero-latency uniform \texttt{u0..u31} from \texttt{FAU\_k} (\texttt{slots64 = bits[63:56]}); \texttt{mode=3} reads \texttt{SmallConstant} (Table~1.2) or \texttt{FauSpecial} (\texttt{pc}, \texttt{lane\_id}, \texttt{lane\_width}, \texttt{tls\_ptr}, \texttt{wls\_ptr}, \texttt{blend\_descriptor\_0..7}).
\end{tightitemize}
\par\textbf{Worked Examples:}
\par\textbf{Example 1 (Floating-point conditional select (true branch)):}
\codebox{\mnem{CSEL.f32}\imod{.lt} \dst{r0}, \src{r1}, \src{r2}, \src{r3}, \src{r4}}
With \texttt{r1 = 1.0f, r2 = 2.0f, r3 = 0xAAAA, r4 = 0xBBBB}: \texttt{1.0 \textless{} 2.0} is true -\textgreater{} selects \texttt{r0 = 0xAAAA}.
\par\textbf{Example 2 (NaN comparison falls through to false branch):}
\codebox{\mnem{CSEL.f32}\imod{.lt} \dst{r0}, \src{r1}, \src{r2}, \src{r3}, \src{r4}}
With \texttt{r1 = NaN, r2 = 2.0f}: ordered \texttt{.lt} evaluates to false -\textgreater{} selects \texttt{r0 = r4} (\texttt{0xBBBB}).
\par\textbf{Mined SPIR-V $\to$ Valhall ISA Example (NaN-to-zero sanitization of texture coordinates via CSEL.f32.ne --- Fragment shader, Mali-G710 (v10)):}
\codebox{\textbf{; SPIR-V (input)}\\\%is\_nan = OpIsNan \%bool \%u\\\%safe\_u = OpSelect \%float \%is\_nan \%float\_0 \%u\\[3pt]\textbf{; Valhall ISA (malioc output)}\\CSEL.f32.ne r1, r1\textasciicircum{}, r1\textasciicircum{}, k0, r1\textasciicircum{}\\CSEL.f32.discard.ne r0, r0\textasciicircum{}, r0\textasciicircum{}, k0, r0\textasciicircum{}}
Self-comparing \texttt{r1 != r1} (\texttt{CSEL.f32.ne}) is true if and only if \texttt{r1} is NaN; hardware selects \texttt{k0} (\texttt{0.0f}) when NaN and preserves \texttt{r1\textasciicircum{}} otherwise in a single cycle.
\begin{center}\footnotesize
\begin{tabularx}{\linewidth}{@{}lllX@{}}
\toprule
Variant & Arch & Class & Opcode Fields \\
\midrule
\texttt{f32} & v9--v14 & \texttt{R4\_CSEL} & \texttt{--} \\
\texttt{s32} & v9--v14 & \texttt{R4\_CSEL} & \texttt{--} \\
\texttt{u32} & v9--v14 & \texttt{R4\_CSEL} & \texttt{--} \\
\texttt{v2f16} & v9--v14 & \texttt{R4\_CSEL} & \texttt{--} \\
\texttt{v2s16} & v9--v14 & \texttt{R4\_CSEL} & \texttt{--} \\
\texttt{v2u16} & v9--v14 & \texttt{R4\_CSEL} & \texttt{--} \\
\bottomrule
\end{tabularx}
\end{center}
\begin{center}\footnotesize
\begin{tabularx}{0.88\linewidth}{@{}l l X@{}}
\toprule
Field Name & Bits & Type / Fixed Value \\
\midrule
\texttt{src0} & \texttt{0:7} & \texttt{SourceEncoding} \\
\texttt{src1} & \texttt{8:15} & \texttt{SourceEncoding} \\
\texttt{src2} & \texttt{16:23} & \texttt{SourceEncoding} \\
\texttt{src3} & \texttt{24:31} & \texttt{SourceEncoding} \\
\hyperref[enum:cmpf_m]{\texttt{cmpf}} & \texttt{32:34} & \texttt{cmpf\_m} \\
\textit{Reserved} & \texttt{35:39} & \texttt{zero} \\
\texttt{dest\_register\_index} & \texttt{40:45} & \texttt{uint} \\
\hyperref[enum:dest_width_m]{\texttt{dest\_width}} & \texttt{46:47} & \texttt{dest\_width\_m} \\
\texttt{datatype} & \texttt{48:51} & \texttt{exact=f32} \\
\texttt{instruction\_class} & \texttt{52:56} & \texttt{exact=R4\_CSEL} \\
\texttt{fau\_page\_index} & \texttt{57:58} & \texttt{uint} \\
\hyperref[enum:flow_control_m]{\texttt{flow}} & \texttt{59:62} & \texttt{flow\_control\_m} \\
\textit{Reserved} & \texttt{63} & \texttt{zero} \\
\bottomrule
\end{tabularx}
\end{center}
\textbf{Fields Explained:}
\begin{tightitemize}\footnotesize
  \item \texttt{src0}: Primary source operand (\texttt{r0..r63}, \texttt{u0..u63} uniform, \texttt{k*} FAU constant, or \texttt{t0..t2} temporary; \texttt{\textasciicircum{}} marks discard/last-use). \emph{Example:} {\scriptsize\ttfamily \mnem{FADD.f32} \dst{r0}, \src{r1}, \src{r2}} (\textbf{SPIR-V:} \texttt{OpFAdd (Operand 0)}; \textbf{GLSL:} \texttt{float c = a + b;})
  \item \texttt{src1}: Secondary source operand (register, uniform, or FAU constant; in shifts, supplies the shift count byte). \emph{Example:} {\scriptsize\ttfamily \mnem{LSHIFT.i32} \dst{r0}, \src{r1}, \src{r2}\imod{.b0}} (\textbf{SPIR-V:} \texttt{OpShiftLeftLogical (Shift)}; \textbf{GLSL:} \texttt{uint y = a \textless{}\textless{} s;})
  \item \texttt{src2}: Third source operand (FMA addend, \texttt{MUX}/\texttt{CSEL} selector, \texttt{IADDC} carry-in, or bitwise mask in fused shift-logic). \emph{Example:} {\scriptsize\ttfamily \mnem{FMA.f32} \dst{r0}, \src{r1}, \src{r2}, \src{r3}} (\textbf{SPIR-V:} \texttt{GLSL.std.450 Fma (Addend)}; \textbf{GLSL:} \texttt{float d = fma(a, b, c);})
  \item \texttt{src3}: Fourth source operand (false-branch value in \texttt{CSEL}, exponent scale in \texttt{FMA\_RSCALE}, or fourth matrix/select input). \emph{Example:} {\scriptsize\ttfamily \mnem{CSEL.f32}\imod{.lt} \dst{r0}, \src{r1}, \src{r2}, \src{r3}, \src{r4}} (\textbf{SPIR-V:} \texttt{OpSelect (Object 2)}; \textbf{GLSL:} \texttt{float r = (a \textless{} b) ? x : y;})
  \item \hyperref[enum:cmpf_m]{\texttt{cmpf}}: Comparison predicate (\texttt{eq}, \texttt{ne}, \texttt{lt}, \texttt{le}, \texttt{gt}, \texttt{ge}, \texttt{gtlt}, \texttt{total}). \emph{Example:} {\scriptsize\ttfamily \mnem{CSEL.f32}\imod{.lt} \dst{r0}, \src{r1}, \src{r2}, \src{r3}, \src{r4}} (\textbf{SPIR-V:} \texttt{OpFOrdLessThan + OpSelect}; \textbf{GLSL:} \texttt{float y = (a \textless{} b) ? c : d;})
  \item \texttt{dest\_register\_index}: 6-bit destination GPR index (\texttt{r0..r63}) receiving the Stage C writeback result (even-aligned pair \texttt{[rN:rN+1]} for 64-bit ops). \emph{Example:} {\scriptsize\ttfamily \mnem{FMA.f32} \dst{r0}, \src{r1}, \src{r2}, \src{r3}} (\textbf{SPIR-V:} \texttt{SSA Result ID}; \textbf{GLSL:} \texttt{float r0 = fma(r1, r2, r3);})
  \item \hyperref[enum:dest_width_m]{\texttt{dest\_width}}: Stage C destination lane write-mask (\texttt{none} = full 32-bit, \texttt{h0} = low 16-bit, \texttt{h1} = high 16-bit without bit-shifting). \emph{Example:} {\scriptsize\ttfamily \mnem{FADD.v2f16} \dst{r0}\imod{.h1}, \src{r1}, \src{r2}} (\textbf{SPIR-V:} \texttt{OpVectorInsertDynamic / WriteMask}; \textbf{GLSL:} \texttt{dst.y = a.y + b.y;})
  \item \texttt{datatype}: Selects the lane data type (\texttt{f32}, \texttt{f16}, \texttt{v2f16}, \texttt{s32}, \texttt{u32}, \texttt{v2s16}, \texttt{v2u16}, \texttt{v4s8}, \texttt{v4u8}, \texttt{s64}, \texttt{u64}, etc.). \emph{Example:} {\scriptsize\ttfamily \mnem{IADD.v2s16} \dst{r0}, \src{r1}, \src{r2}} (\textbf{SPIR-V:} \texttt{OpTypeInt / OpTypeFloat}; \textbf{GLSL:} \texttt{i16vec2 c = a + b;})
  \item \texttt{instruction\_class}: 4-bit major instruction class selector (\texttt{bits 48:51}) partitioning the 64-bit Valhall opcode space (e.g., \texttt{FMA}, \texttt{CVT}, \texttt{SFU}, \texttt{LS}, \texttt{TEX}, \texttt{BRANCH}). \emph{Example:} {\scriptsize\ttfamily \mnem{FADD.f32} \dst{r0}, \src{r1}, \src{r2}} (\textbf{SPIR-V:} \texttt{N/A (encoding class)}; \textbf{GLSL:} \texttt{// Hardware execution pipe / decode class})
  \item \texttt{fau\_page\_index}: Selects the FAU (Fast Access Uniform) page (\texttt{0..3}) when a source operand references hardware small constants, thread IDs, or lane constants. \emph{Example:} {\scriptsize\ttfamily \mnem{FADD.f32} \dst{r0}, \src{r1}, \src{k1f}} (\textbf{SPIR-V:} \texttt{OpConstant / BuiltIn}; \textbf{GLSL:} \texttt{float y = x + 1.0;})
  \item \hyperref[enum:flow_control_m]{\texttt{flow}}: 4-bit embedded flow-control \& scoreboard wait modifier (\texttt{bits 59:62}: \texttt{none}, \texttt{wait0..wait0126}, \texttt{wait}, \texttt{reconverge}, \texttt{discard}, \texttt{end}). \emph{Example:} {\scriptsize\ttfamily \mnem{FADD.f32}\imod{.wait0.reconverge} \dst{r0}, \src{r1}, \src{r2}} (\textbf{SPIR-V:} \texttt{OpControlBarrier / OpReturn}; \textbf{GLSL:} \texttt{// Scoreboard wait + warp reconvergence})
\end{tightitemize}
\clearpage

\subsection*{FCMP}
\addcontentsline{toc}{subsection}{FCMP}
Compares floating-point \texttt{src0[i]} and \texttt{src1[i]} using predicate \texttt{cmpf} independently for each SIMD element $i \in \{0, \dots, N-1\}$ ($N=1$ \texttt{f32}, $N=2$ \texttt{v2f16}) and encodes the boolean into \texttt{dest[i]} via \texttt{result\_type} (\texttt{i1} $\to$ \texttt{1}, \texttt{m1} $\to$ all-ones, \texttt{f1} $\to$ \texttt{1.0}, false $\to$ \texttt{0}).
\par\textbf{Notes:}
\begin{tightitemize}
  \item \texttt{ne} is unordered (\texttt{True} on NaN).
  \item \texttt{gtlt} is ordered not-equal (\texttt{False} on NaN).
  \item \texttt{total} is IEEE-754 \texttt{total\_cmp(a, b) \textless{}= 0} (\texttt{ops.rs:1489-1521}).
  \item \textbf{Source references:} \texttt{kraid/ops.rs:741-790, 1386-1522}, \texttt{encode\_v9.rs:933-1024}, \texttt{hw\_tests.rs:1318-1388}.
\end{tightitemize}
\codebox{\mnem{FCMP}\imod{\{.flow\}}\imod{\{.result\_type\}}\imod{\{.cmpf\}} \dst{dest}\imod{\{.dest\_width\}}, \src{src0}\imod{\{.widen0\}}\imod{\{.abs0\}}\imod{\{.neg0\}}, \src{src1}\imod{\{.widen1\}}\imod{\{.abs1\}}\imod{\{.neg1\}}\\\mnem{FCMP}\imod{\{.flow\}}\imod{\{.result\_type\}}\imod{\{.cmpf\}} \dst{dest}\imod{\{.dest\_width\}}, \src{src0}\imod{\{.swz0\}}\imod{\{.abs0\}}\imod{\{.neg0\}}, \src{src1}\imod{\{.swz1\}}\imod{\{.abs1\}}\imod{\{.neg1\}}}
\textbf{Operands:}
\begin{tightitemize}\footnotesize
  \item \texttt{\dst{dest}} (\texttt{dest\_register\_index}, \texttt{dest\_width}): Destination GPR receiving comparison result formatted per \texttt{result\_type} (\texttt{.i1}, \texttt{.m1}, or \texttt{.f1}).
  \item \texttt{\src{src0}}: Left-hand floating-point comparison operand.
  \item \texttt{\src{src1}}: Right-hand floating-point comparison operand.
\end{tightitemize}
\par\textbf{Closest SPIR-V Op:} \texttt{OpFOrdEqual/LessThan/... / OpFUnordNotEqual}
\par\textbf{Pseudo-GLSL (Operational Semantics):}
\codebox{for (int i = 0; i \textless{} N; ++i) \{\\    bool c = cmpf(mod(src0)[i], mod(src1)[i]);\\    dest[i] = (result\_type == m1) ? (c ? \textasciitilde{}0u : 0u) : (result\_type == f1 ? (c ? 1.0 : 0.0) : (c ? 1u : 0u));\\\}}
\textbf{Execution Context:}
\begin{tightitemize}\footnotesize
  \item \textbf{Compute stage preloads (\hyperref[sec:csf-spd]{\texttt{SPD\_k}}, \texttt{r16..r23}):} Hardware pre-fills \texttt{r55..r57} with \texttt{WorkGroupID.xyz + JOB\_OFFSET} (from \hyperref[sec:csf-regs]{\texttt{COMPUTESR::JOB\_OFFSET\_X/Y/Z}} \texttt{r34..r36}) and \texttt{r60..r62} with \texttt{LocalInvocationID.xyz} or \texttt{GlobalInvocationID.xyz} (from \texttt{WG\_SIZE} \texttt{r33} \& \texttt{JOB\_SIZE} \texttt{r37..r39}, selected by \texttt{ShaderProgram::preload}) so compute shaders can derive thread/element indices without memory loads.
  \item \textbf{Vertex stage preloads (\hyperref[sec:csf-spd]{\texttt{SPD\_0}/\texttt{SPD\_1}}, \texttt{r16..r19}):} The hardware Index Fetcher reads \hyperref[sec:csf-regs]{\texttt{IDVSSR::INDEX\_BUFFER}} (\texttt{r54:r55}) and pre-fills \texttt{r59, r60} with \texttt{VertexIndex} (\texttt{+ VERTEX\_OFFSET + GLOBAL\_ATTR\_OFF}) and \texttt{r61} with \texttt{InstanceIndex} (\texttt{+ INSTANCE\_OFFSET}) so Position and Varying shaders can pull vertex attributes directly.
  \item \textbf{Inline FAU \& Special Operands (\hyperref[sec:csf-tsd-fau]{\texttt{FAU\_k}}, \texttt{r8..r15}):} Any 8-bit source with \texttt{mode=2} reads zero-latency uniform \texttt{u0..u31} from \texttt{FAU\_k} (\texttt{slots64 = bits[63:56]}); \texttt{mode=3} reads \texttt{SmallConstant} (Table~1.2) or \texttt{FauSpecial} (\texttt{pc}, \texttt{lane\_id}, \texttt{lane\_width}, \texttt{tls\_ptr}, \texttt{wls\_ptr}, \texttt{blend\_descriptor\_0..7}).
\end{tightitemize}
\par\textbf{Worked Examples:}
\par\textbf{Example 1 (Unordered `.ne` on NaN with `.m1` bitmask):}
\codebox{\mnem{FCMP.f32}\imod{.ne.m1} \dst{r0}, \src{r1}, \src{r2}}
With \texttt{r1 = NaN, r2 = 1.0f}: \texttt{.ne} is unordered (True on NaN), \texttt{.m1} formats True as \texttt{r0 = 0xFFFFFFFF}.
\par\textbf{Example 2 (Ordered `.gtlt` on NaN with `.f1` float boolean):}
\codebox{\mnem{FCMP.f32}\imod{.gtlt.f1} \dst{r0}, \src{r1}, \src{r2}}
With \texttt{r1 = NaN, r2 = 1.0f}: \texttt{.gtlt} is ordered (False on NaN) -\textgreater{} \texttt{r0 = 0.0f}; for \texttt{(2.0f, 1.0f)} -\textgreater{} \texttt{r0 = 1.0f} (\texttt{0x3F800000}).

\begin{center}\footnotesize
\begin{tabularx}{\linewidth}{@{}lllX@{}}
\toprule
Variant & Arch & Class & Opcode Fields \\
\midrule
\texttt{f32} & v9--v14 & \texttt{R3\_CMP} & \texttt{--} \\
\texttt{v2f16} & v9--v14 & \texttt{R3\_CMP} & \texttt{--} \\
\bottomrule
\end{tabularx}
\end{center}
\begin{center}\footnotesize
\begin{tabularx}{0.88\linewidth}{@{}l l X@{}}
\toprule
Field Name & Bits & Type / Fixed Value \\
\midrule
\texttt{src0} & \texttt{0:7} & \texttt{SourceEncoding} \\
\texttt{src1} & \texttt{8:15} & \texttt{SourceEncoding} \\
\texttt{src2} & \texttt{16:23} & \texttt{exact=192} \\
\texttt{combine} & \texttt{24:25} & \texttt{exact=or} \\
\hyperref[enum:widen_m]{\texttt{widen1}} & \texttt{26:27} & \texttt{widen\_m (def none)} \\
\hyperref[enum:widen_m]{\texttt{widen0}} & \texttt{28:29} & \texttt{widen\_m (def none)} \\
\hyperref[enum:cmp_result_type_m]{\texttt{result\_type}} & \texttt{30:31} & \texttt{cmp\_result\_type\_m (def i1)} \\
\hyperref[enum:cmpf_m]{\texttt{cmpf}} & \texttt{32:34} & \texttt{cmpf\_m} \\
\textit{Reserved} & \texttt{35} & \texttt{zero} \\
\hyperref[enum:neg_m]{\texttt{neg1}} & \texttt{36} & \texttt{neg\_m (def none)} \\
\hyperref[enum:abs_m]{\texttt{abs1}} & \texttt{37} & \texttt{abs\_m (def none)} \\
\hyperref[enum:neg_m]{\texttt{neg0}} & \texttt{38} & \texttt{neg\_m (def none)} \\
\hyperref[enum:abs_m]{\texttt{abs0}} & \texttt{39} & \texttt{abs\_m (def none)} \\
\texttt{dest\_register\_index} & \texttt{40:45} & \texttt{uint} \\
\hyperref[enum:dest_width_m]{\texttt{dest\_width}} & \texttt{46:47} & \texttt{dest\_width\_m} \\
\texttt{datatype} & \texttt{48:51} & \texttt{exact=f32} \\
\texttt{instruction\_class} & \texttt{52:56} & \texttt{exact=R3\_CMP} \\
\texttt{fau\_page\_index} & \texttt{57:58} & \texttt{uint} \\
\hyperref[enum:flow_control_m]{\texttt{flow}} & \texttt{59:62} & \texttt{flow\_control\_m} \\
\textit{Reserved} & \texttt{63} & \texttt{zero} \\
\bottomrule
\end{tabularx}
\end{center}
\textbf{Fields Explained:}
\begin{tightitemize}\footnotesize
  \item \texttt{src0}: Primary source operand (\texttt{r0..r63}, \texttt{u0..u63} uniform, \texttt{k*} FAU constant, or \texttt{t0..t2} temporary; \texttt{\textasciicircum{}} marks discard/last-use). \emph{Example:} {\scriptsize\ttfamily \mnem{FADD.f32} \dst{r0}, \src{r1}, \src{r2}} (\textbf{SPIR-V:} \texttt{OpFAdd (Operand 0)}; \textbf{GLSL:} \texttt{float c = a + b;})
  \item \texttt{src1}: Secondary source operand (register, uniform, or FAU constant; in shifts, supplies the shift count byte). \emph{Example:} {\scriptsize\ttfamily \mnem{LSHIFT.i32} \dst{r0}, \src{r1}, \src{r2}\imod{.b0}} (\textbf{SPIR-V:} \texttt{OpShiftLeftLogical (Shift)}; \textbf{GLSL:} \texttt{uint y = a \textless{}\textless{} s;})
  \item \texttt{src2}: Third source operand (FMA addend, \texttt{MUX}/\texttt{CSEL} selector, \texttt{IADDC} carry-in, or bitwise mask in fused shift-logic). \emph{Example:} {\scriptsize\ttfamily \mnem{FMA.f32} \dst{r0}, \src{r1}, \src{r2}, \src{r3}} (\textbf{SPIR-V:} \texttt{GLSL.std.450 Fma (Addend)}; \textbf{GLSL:} \texttt{float d = fma(a, b, c);})
  \item \texttt{combine}: Combines 16-bit halfword predicates (\texttt{h0}, \texttt{h1}, \texttt{and}, \texttt{lowbits}) before evaluating a branch. \emph{Example:} {\scriptsize\ttfamily \mnem{BRANCH}\imod{.and} \src{r0}, \imm{\#+8}} (\textbf{SPIR-V:} \texttt{OpBranchConditional (packed predicate)}; \textbf{GLSL:} \texttt{if (condA \&\& condB) \{ ... \}})
  \item \hyperref[enum:widen_m]{\texttt{widen1}}: Stage A widening extractor on \texttt{src1} (promotes narrow byte/half/word lane before Stage B). \emph{Example:} {\scriptsize\ttfamily \mnem{IMUL.u64} \dst{[r0:r1]}, \src{[r2:r3]}\imod{.w0}, \src{[r4:r5]}\imod{.w0}} (\textbf{SPIR-V:} \texttt{OpUConvert / OpUMulExtended}; \textbf{GLSL:} \texttt{uint64\_t p = uint64\_t(a) * uint64\_t(b);})
  \item \hyperref[enum:widen_m]{\texttt{widen0}}: Stage A widening extractor on \texttt{src0} (promotes FP16 \texttt{h0}/\texttt{h1} to FP32, or sign/zero-extends \texttt{b0..b3}/\texttt{h0..h1}/\texttt{w0..w1} to 32/64-bit before \texttt{abs}/\texttt{neg}). \emph{Example:} {\scriptsize\ttfamily \mnem{IADD.s32} \dst{r0}, \src{r1}\imod{.h0}, \src{r2}} (\textbf{SPIR-V:} \texttt{OpSConvert / OpFConvert}; \textbf{GLSL:} \texttt{int c = int(int16\_t(a)) + b;})
  \item \hyperref[enum:cmp_result_type_m]{\texttt{result\_type}}: Encodes boolean output representation (\texttt{i1} = 0/1, \texttt{f1} = 0.0/1.0, \texttt{m1} = 0/-1, \texttt{c} = carry). \emph{Example:} {\scriptsize\ttfamily \mnem{FCMP.f32}\imod{.lt.m1} \dst{r0}, \src{r1}, \src{r2}} (\textbf{SPIR-V:} \texttt{OpFOrdLessThan + OpSelect}; \textbf{GLSL:} \texttt{int mask = (a \textless{} b) ? -1 : 0;})
  \item \hyperref[enum:cmpf_m]{\texttt{cmpf}}: Comparison predicate (\texttt{eq}, \texttt{ne}, \texttt{lt}, \texttt{le}, \texttt{gt}, \texttt{ge}, \texttt{gtlt}, \texttt{total}). \emph{Example:} {\scriptsize\ttfamily \mnem{CSEL.f32}\imod{.lt} \dst{r0}, \src{r1}, \src{r2}, \src{r3}, \src{r4}} (\textbf{SPIR-V:} \texttt{OpFOrdLessThan + OpSelect}; \textbf{GLSL:} \texttt{float y = (a \textless{} b) ? c : d;})
  \item \hyperref[enum:neg_m]{\texttt{neg1}}: Stage A floating-point negation on \texttt{src1} (flips IEEE sign bit after swizzle/widen and \texttt{abs1}). \emph{Example:} {\scriptsize\ttfamily \mnem{FMA.f32} \dst{r0}, \src{r1}, \src{r2}\imod{.neg}, \src{r3}} (\textbf{SPIR-V:} \texttt{OpFNegate}; \textbf{GLSL:} \texttt{float d = fma(a, -b, c);})
  \item \hyperref[enum:abs_m]{\texttt{abs1}}: Stage A floating-point absolute value on \texttt{src1} (clears IEEE sign bit after swizzle/widen, before \texttt{neg1}). \emph{Example:} {\scriptsize\ttfamily \mnem{FMA.f32} \dst{r0}, \src{r1}, \src{r2}\imod{.abs}, \src{r3}} (\textbf{SPIR-V:} \texttt{GLSL.std.450 FAbs}; \textbf{GLSL:} \texttt{float d = fma(a, abs(b), c);})
  \item \hyperref[enum:neg_m]{\texttt{neg0}}: Stage A floating-point negation on \texttt{src0} (flips IEEE sign bit after swizzle/widen and \texttt{abs0}). \emph{Example:} {\scriptsize\ttfamily \mnem{FADD.f32} \dst{r0}, \src{r1}\imod{.neg}, \src{r2}} (\textbf{SPIR-V:} \texttt{OpFNegate}; \textbf{GLSL:} \texttt{float c = -a + b;})
  \item \hyperref[enum:abs_m]{\texttt{abs0}}: Stage A floating-point absolute value on \texttt{src0} (clears IEEE sign bit after swizzle/widen, before \texttt{neg0}). \emph{Example:} {\scriptsize\ttfamily \mnem{FADD.f32} \dst{r0}, \src{r1}\imod{.abs}, \src{r2}} (\textbf{SPIR-V:} \texttt{GLSL.std.450 FAbs}; \textbf{GLSL:} \texttt{float c = abs(a) + b;})
  \item \texttt{dest\_register\_index}: 6-bit destination GPR index (\texttt{r0..r63}) receiving the Stage C writeback result (even-aligned pair \texttt{[rN:rN+1]} for 64-bit ops). \emph{Example:} {\scriptsize\ttfamily \mnem{FMA.f32} \dst{r0}, \src{r1}, \src{r2}, \src{r3}} (\textbf{SPIR-V:} \texttt{SSA Result ID}; \textbf{GLSL:} \texttt{float r0 = fma(r1, r2, r3);})
  \item \hyperref[enum:dest_width_m]{\texttt{dest\_width}}: Stage C destination lane write-mask (\texttt{none} = full 32-bit, \texttt{h0} = low 16-bit, \texttt{h1} = high 16-bit without bit-shifting). \emph{Example:} {\scriptsize\ttfamily \mnem{FADD.v2f16} \dst{r0}\imod{.h1}, \src{r1}, \src{r2}} (\textbf{SPIR-V:} \texttt{OpVectorInsertDynamic / WriteMask}; \textbf{GLSL:} \texttt{dst.y = a.y + b.y;})
  \item \texttt{datatype}: Selects the lane data type (\texttt{f32}, \texttt{f16}, \texttt{v2f16}, \texttt{s32}, \texttt{u32}, \texttt{v2s16}, \texttt{v2u16}, \texttt{v4s8}, \texttt{v4u8}, \texttt{s64}, \texttt{u64}, etc.). \emph{Example:} {\scriptsize\ttfamily \mnem{IADD.v2s16} \dst{r0}, \src{r1}, \src{r2}} (\textbf{SPIR-V:} \texttt{OpTypeInt / OpTypeFloat}; \textbf{GLSL:} \texttt{i16vec2 c = a + b;})
  \item \texttt{instruction\_class}: 4-bit major instruction class selector (\texttt{bits 48:51}) partitioning the 64-bit Valhall opcode space (e.g., \texttt{FMA}, \texttt{CVT}, \texttt{SFU}, \texttt{LS}, \texttt{TEX}, \texttt{BRANCH}). \emph{Example:} {\scriptsize\ttfamily \mnem{FADD.f32} \dst{r0}, \src{r1}, \src{r2}} (\textbf{SPIR-V:} \texttt{N/A (encoding class)}; \textbf{GLSL:} \texttt{// Hardware execution pipe / decode class})
  \item \texttt{fau\_page\_index}: Selects the FAU (Fast Access Uniform) page (\texttt{0..3}) when a source operand references hardware small constants, thread IDs, or lane constants. \emph{Example:} {\scriptsize\ttfamily \mnem{FADD.f32} \dst{r0}, \src{r1}, \src{k1f}} (\textbf{SPIR-V:} \texttt{OpConstant / BuiltIn}; \textbf{GLSL:} \texttt{float y = x + 1.0;})
  \item \hyperref[enum:flow_control_m]{\texttt{flow}}: 4-bit embedded flow-control \& scoreboard wait modifier (\texttt{bits 59:62}: \texttt{none}, \texttt{wait0..wait0126}, \texttt{wait}, \texttt{reconverge}, \texttt{discard}, \texttt{end}). \emph{Example:} {\scriptsize\ttfamily \mnem{FADD.f32}\imod{.wait0.reconverge} \dst{r0}, \src{r1}, \src{r2}} (\textbf{SPIR-V:} \texttt{OpControlBarrier / OpReturn}; \textbf{GLSL:} \texttt{// Scoreboard wait + warp reconvergence})
\end{tightitemize}
\clearpage

\subsection*{FCMP\_AND}
\addcontentsline{toc}{subsection}{FCMP\_AND}
Compares floating-point \texttt{src0[i]} and \texttt{src1[i]} via \texttt{cmpf} for each SIMD element $i \in \{0, \dots, N-1\}$ ($N=1$ \texttt{f32}, $N=2$ \texttt{v2f16}), combines the predicate with \texttt{(src2[i] != 0)} using logical AND, and formats the boolean into \texttt{dest[i]} according to \texttt{result\_type}.
\par\textbf{Notes:}
\begin{tightitemize}
  \item \texttt{src2[i]} tested as \texttt{!= 0} per lane before formatting via \texttt{result\_type} (\texttt{ops.rs:785-786}).
  \item \textbf{Source references:} \texttt{kraid/ops.rs:741-790}, \texttt{encode\_v9.rs:974-982}, \texttt{opt\_fuse.rs:188-238}, \texttt{hw\_tests.rs:1318-1388}.
\end{tightitemize}
\codebox{\mnem{FCMP\_AND}\imod{\{.flow\}}\imod{\{.result\_type\}}\imod{\{.cmpf\}} \dst{dest}\imod{\{.dest\_width\}}, \src{src0}\imod{\{.widen0\}}\imod{\{.abs0\}}\imod{\{.neg0\}}, \src{src1}\imod{\{.widen1\}}\imod{\{.abs1\}}\imod{\{.neg1\}}, \src{src2}\\\mnem{FCMP\_AND}\imod{\{.flow\}}\imod{\{.result\_type\}}\imod{\{.cmpf\}} \dst{dest}\imod{\{.dest\_width\}}, \src{src0}\imod{\{.swz0\}}\imod{\{.abs0\}}\imod{\{.neg0\}}, \src{src1}\imod{\{.swz1\}}\imod{\{.abs1\}}\imod{\{.neg1\}}, \src{src2}}
\textbf{Operands:}
\begin{tightitemize}\footnotesize
  \item \texttt{\dst{dest}} (\texttt{dest\_register\_index}, \texttt{dest\_width}): Destination GPR receiving \texttt{cmpf(src0, src1) \&\& (src2 != 0)} formatted per \texttt{result\_type}.
  \item \texttt{\src{src0}}: Left-hand floating-point comparison operand.
  \item \texttt{\src{src1}}: Right-hand floating-point comparison operand.
  \item \texttt{\src{src2}}: Existing boolean/mask predicate combined with logical AND.
\end{tightitemize}
\par\textbf{Closest SPIR-V Op:} \texttt{OpLogicalAnd + OpFOrd*}
\par\textbf{Pseudo-GLSL (Operational Semantics):}
\codebox{for (int i = 0; i \textless{} N; ++i) \{\\    bool c = cmpf(mod(src0)[i], mod(src1)[i]) \&\& (src2[i] != 0);\\    dest[i] = formatResult(result\_type, c); // .i1: 1/0, .m1: \textasciitilde{}0/0, .f1: 1.0/0.0\\\}}
\textbf{Execution Context:}
\begin{tightitemize}\footnotesize
  \item \textbf{Compute stage preloads (\hyperref[sec:csf-spd]{\texttt{SPD\_k}}, \texttt{r16..r23}):} Hardware pre-fills \texttt{r55..r57} with \texttt{WorkGroupID.xyz + JOB\_OFFSET} (from \hyperref[sec:csf-regs]{\texttt{COMPUTESR::JOB\_OFFSET\_X/Y/Z}} \texttt{r34..r36}) and \texttt{r60..r62} with \texttt{LocalInvocationID.xyz} or \texttt{GlobalInvocationID.xyz} (from \texttt{WG\_SIZE} \texttt{r33} \& \texttt{JOB\_SIZE} \texttt{r37..r39}, selected by \texttt{ShaderProgram::preload}) so compute shaders can derive thread/element indices without memory loads.
  \item \textbf{Vertex stage preloads (\hyperref[sec:csf-spd]{\texttt{SPD\_0}/\texttt{SPD\_1}}, \texttt{r16..r19}):} The hardware Index Fetcher reads \hyperref[sec:csf-regs]{\texttt{IDVSSR::INDEX\_BUFFER}} (\texttt{r54:r55}) and pre-fills \texttt{r59, r60} with \texttt{VertexIndex} (\texttt{+ VERTEX\_OFFSET + GLOBAL\_ATTR\_OFF}) and \texttt{r61} with \texttt{InstanceIndex} (\texttt{+ INSTANCE\_OFFSET}) so Position and Varying shaders can pull vertex attributes directly.
  \item \textbf{Inline FAU \& Special Operands (\hyperref[sec:csf-tsd-fau]{\texttt{FAU\_k}}, \texttt{r8..r15}):} Any 8-bit source with \texttt{mode=2} reads zero-latency uniform \texttt{u0..u31} from \texttt{FAU\_k} (\texttt{slots64 = bits[63:56]}); \texttt{mode=3} reads \texttt{SmallConstant} (Table~1.2) or \texttt{FauSpecial} (\texttt{pc}, \texttt{lane\_id}, \texttt{lane\_width}, \texttt{tls\_ptr}, \texttt{wls\_ptr}, \texttt{blend\_descriptor\_0..7}).
\end{tightitemize}
\par\textbf{Worked Examples:}
\par\textbf{Example 1 (Short-circuit chain with false accumulator):}
\codebox{\mnem{FCMP\_AND.f32}\imod{.lt.i1} \dst{r0}, \src{r1}, \src{r2}, \src{r3}}
With \texttt{r1 = 1.0f, r2 = 2.0f, r3 = 0}: \texttt{(1.0 \textless{} 2.0) \&\& (0 != 0)} is False -\textgreater{} \texttt{r0 = 0}.
\par\textbf{Example 2 (Both predicates true with `.m1` mask):}
\codebox{\mnem{FCMP\_AND.f32}\imod{.le.m1} \dst{r0}, \src{r1}, \src{r2}, \src{r3}}
With \texttt{r1 = 2.0f, r2 = 2.0f, r3 = 0xFFFFFFFF}: \texttt{(2.0 \textless{}= 2.0) \&\& true} is True -\textgreater{} \texttt{r0 = 0xFFFFFFFF}.
\par\textbf{Mined SPIR-V $\to$ Valhall ISA Example (Fused float comparison AND existing boolean predicate --- Vertex shader, Mali-G710 (v10)):}
\codebox{\textbf{; SPIR-V (input)}\\\%gt = OpFOrdGreaterThan \%bool \%v \%thresh\\\%both = OpLogicalAnd \%bool \%gt \%ubo\_flag\\\%sel = OpSelect \%float \%both \%val\_a \%val\_b\\[3pt]\textbf{; Valhall ISA (malioc output)}\\FCMP\_AND.f32.m1.gt r1, r0, u4.w0, u4.w1\\CSEL.u32.eq r1, r1\textasciicircum{}, k0, r2\textasciicircum{}, u3.w0}
\texttt{FCMP\_AND} folds \texttt{(r0 \textgreater{} u4.w0) \&\& (u4.w1 != 0)} into one instruction, writing an \texttt{.m1} mask to \texttt{r1} that drives \texttt{CSEL.u32.eq}.
\begin{center}\footnotesize
\begin{tabularx}{\linewidth}{@{}lllX@{}}
\toprule
Variant & Arch & Class & Opcode Fields \\
\midrule
\texttt{f32} & v9--v14 & \texttt{R3\_CMP} & \texttt{--} \\
\texttt{v2f16} & v9--v14 & \texttt{R3\_CMP} & \texttt{--} \\
\bottomrule
\end{tabularx}
\end{center}
\begin{center}\footnotesize
\begin{tabularx}{0.88\linewidth}{@{}l l X@{}}
\toprule
Field Name & Bits & Type / Fixed Value \\
\midrule
\texttt{src0} & \texttt{0:7} & \texttt{SourceEncoding} \\
\texttt{src1} & \texttt{8:15} & \texttt{SourceEncoding} \\
\texttt{src2} & \texttt{16:23} & \texttt{SourceEncoding} \\
\texttt{combine} & \texttt{24:25} & \texttt{exact=and} \\
\hyperref[enum:widen_m]{\texttt{widen1}} & \texttt{26:27} & \texttt{widen\_m (def none)} \\
\hyperref[enum:widen_m]{\texttt{widen0}} & \texttt{28:29} & \texttt{widen\_m (def none)} \\
\hyperref[enum:cmp_result_type_m]{\texttt{result\_type}} & \texttt{30:31} & \texttt{cmp\_result\_type\_m (def i1)} \\
\hyperref[enum:cmpf_m]{\texttt{cmpf}} & \texttt{32:34} & \texttt{cmpf\_m} \\
\textit{Reserved} & \texttt{35} & \texttt{zero} \\
\hyperref[enum:neg_m]{\texttt{neg1}} & \texttt{36} & \texttt{neg\_m (def none)} \\
\hyperref[enum:abs_m]{\texttt{abs1}} & \texttt{37} & \texttt{abs\_m (def none)} \\
\hyperref[enum:neg_m]{\texttt{neg0}} & \texttt{38} & \texttt{neg\_m (def none)} \\
\hyperref[enum:abs_m]{\texttt{abs0}} & \texttt{39} & \texttt{abs\_m (def none)} \\
\texttt{dest\_register\_index} & \texttt{40:45} & \texttt{uint} \\
\hyperref[enum:dest_width_m]{\texttt{dest\_width}} & \texttt{46:47} & \texttt{dest\_width\_m} \\
\texttt{datatype} & \texttt{48:51} & \texttt{exact=f32} \\
\texttt{instruction\_class} & \texttt{52:56} & \texttt{exact=R3\_CMP} \\
\texttt{fau\_page\_index} & \texttt{57:58} & \texttt{uint} \\
\hyperref[enum:flow_control_m]{\texttt{flow}} & \texttt{59:62} & \texttt{flow\_control\_m} \\
\textit{Reserved} & \texttt{63} & \texttt{zero} \\
\bottomrule
\end{tabularx}
\end{center}
\textbf{Fields Explained:}
\begin{tightitemize}\footnotesize
  \item \texttt{src0}: Primary source operand (\texttt{r0..r63}, \texttt{u0..u63} uniform, \texttt{k*} FAU constant, or \texttt{t0..t2} temporary; \texttt{\textasciicircum{}} marks discard/last-use). \emph{Example:} {\scriptsize\ttfamily \mnem{FADD.f32} \dst{r0}, \src{r1}, \src{r2}} (\textbf{SPIR-V:} \texttt{OpFAdd (Operand 0)}; \textbf{GLSL:} \texttt{float c = a + b;})
  \item \texttt{src1}: Secondary source operand (register, uniform, or FAU constant; in shifts, supplies the shift count byte). \emph{Example:} {\scriptsize\ttfamily \mnem{LSHIFT.i32} \dst{r0}, \src{r1}, \src{r2}\imod{.b0}} (\textbf{SPIR-V:} \texttt{OpShiftLeftLogical (Shift)}; \textbf{GLSL:} \texttt{uint y = a \textless{}\textless{} s;})
  \item \texttt{src2}: Third source operand (FMA addend, \texttt{MUX}/\texttt{CSEL} selector, \texttt{IADDC} carry-in, or bitwise mask in fused shift-logic). \emph{Example:} {\scriptsize\ttfamily \mnem{FMA.f32} \dst{r0}, \src{r1}, \src{r2}, \src{r3}} (\textbf{SPIR-V:} \texttt{GLSL.std.450 Fma (Addend)}; \textbf{GLSL:} \texttt{float d = fma(a, b, c);})
  \item \texttt{combine}: Combines 16-bit halfword predicates (\texttt{h0}, \texttt{h1}, \texttt{and}, \texttt{lowbits}) before evaluating a branch. \emph{Example:} {\scriptsize\ttfamily \mnem{BRANCH}\imod{.and} \src{r0}, \imm{\#+8}} (\textbf{SPIR-V:} \texttt{OpBranchConditional (packed predicate)}; \textbf{GLSL:} \texttt{if (condA \&\& condB) \{ ... \}})
  \item \hyperref[enum:widen_m]{\texttt{widen1}}: Stage A widening extractor on \texttt{src1} (promotes narrow byte/half/word lane before Stage B). \emph{Example:} {\scriptsize\ttfamily \mnem{IMUL.u64} \dst{[r0:r1]}, \src{[r2:r3]}\imod{.w0}, \src{[r4:r5]}\imod{.w0}} (\textbf{SPIR-V:} \texttt{OpUConvert / OpUMulExtended}; \textbf{GLSL:} \texttt{uint64\_t p = uint64\_t(a) * uint64\_t(b);})
  \item \hyperref[enum:widen_m]{\texttt{widen0}}: Stage A widening extractor on \texttt{src0} (promotes FP16 \texttt{h0}/\texttt{h1} to FP32, or sign/zero-extends \texttt{b0..b3}/\texttt{h0..h1}/\texttt{w0..w1} to 32/64-bit before \texttt{abs}/\texttt{neg}). \emph{Example:} {\scriptsize\ttfamily \mnem{IADD.s32} \dst{r0}, \src{r1}\imod{.h0}, \src{r2}} (\textbf{SPIR-V:} \texttt{OpSConvert / OpFConvert}; \textbf{GLSL:} \texttt{int c = int(int16\_t(a)) + b;})
  \item \hyperref[enum:cmp_result_type_m]{\texttt{result\_type}}: Encodes boolean output representation (\texttt{i1} = 0/1, \texttt{f1} = 0.0/1.0, \texttt{m1} = 0/-1, \texttt{c} = carry). \emph{Example:} {\scriptsize\ttfamily \mnem{FCMP.f32}\imod{.lt.m1} \dst{r0}, \src{r1}, \src{r2}} (\textbf{SPIR-V:} \texttt{OpFOrdLessThan + OpSelect}; \textbf{GLSL:} \texttt{int mask = (a \textless{} b) ? -1 : 0;})
  \item \hyperref[enum:cmpf_m]{\texttt{cmpf}}: Comparison predicate (\texttt{eq}, \texttt{ne}, \texttt{lt}, \texttt{le}, \texttt{gt}, \texttt{ge}, \texttt{gtlt}, \texttt{total}). \emph{Example:} {\scriptsize\ttfamily \mnem{CSEL.f32}\imod{.lt} \dst{r0}, \src{r1}, \src{r2}, \src{r3}, \src{r4}} (\textbf{SPIR-V:} \texttt{OpFOrdLessThan + OpSelect}; \textbf{GLSL:} \texttt{float y = (a \textless{} b) ? c : d;})
  \item \hyperref[enum:neg_m]{\texttt{neg1}}: Stage A floating-point negation on \texttt{src1} (flips IEEE sign bit after swizzle/widen and \texttt{abs1}). \emph{Example:} {\scriptsize\ttfamily \mnem{FMA.f32} \dst{r0}, \src{r1}, \src{r2}\imod{.neg}, \src{r3}} (\textbf{SPIR-V:} \texttt{OpFNegate}; \textbf{GLSL:} \texttt{float d = fma(a, -b, c);})
  \item \hyperref[enum:abs_m]{\texttt{abs1}}: Stage A floating-point absolute value on \texttt{src1} (clears IEEE sign bit after swizzle/widen, before \texttt{neg1}). \emph{Example:} {\scriptsize\ttfamily \mnem{FMA.f32} \dst{r0}, \src{r1}, \src{r2}\imod{.abs}, \src{r3}} (\textbf{SPIR-V:} \texttt{GLSL.std.450 FAbs}; \textbf{GLSL:} \texttt{float d = fma(a, abs(b), c);})
  \item \hyperref[enum:neg_m]{\texttt{neg0}}: Stage A floating-point negation on \texttt{src0} (flips IEEE sign bit after swizzle/widen and \texttt{abs0}). \emph{Example:} {\scriptsize\ttfamily \mnem{FADD.f32} \dst{r0}, \src{r1}\imod{.neg}, \src{r2}} (\textbf{SPIR-V:} \texttt{OpFNegate}; \textbf{GLSL:} \texttt{float c = -a + b;})
  \item \hyperref[enum:abs_m]{\texttt{abs0}}: Stage A floating-point absolute value on \texttt{src0} (clears IEEE sign bit after swizzle/widen, before \texttt{neg0}). \emph{Example:} {\scriptsize\ttfamily \mnem{FADD.f32} \dst{r0}, \src{r1}\imod{.abs}, \src{r2}} (\textbf{SPIR-V:} \texttt{GLSL.std.450 FAbs}; \textbf{GLSL:} \texttt{float c = abs(a) + b;})
  \item \texttt{dest\_register\_index}: 6-bit destination GPR index (\texttt{r0..r63}) receiving the Stage C writeback result (even-aligned pair \texttt{[rN:rN+1]} for 64-bit ops). \emph{Example:} {\scriptsize\ttfamily \mnem{FMA.f32} \dst{r0}, \src{r1}, \src{r2}, \src{r3}} (\textbf{SPIR-V:} \texttt{SSA Result ID}; \textbf{GLSL:} \texttt{float r0 = fma(r1, r2, r3);})
  \item \hyperref[enum:dest_width_m]{\texttt{dest\_width}}: Stage C destination lane write-mask (\texttt{none} = full 32-bit, \texttt{h0} = low 16-bit, \texttt{h1} = high 16-bit without bit-shifting). \emph{Example:} {\scriptsize\ttfamily \mnem{FADD.v2f16} \dst{r0}\imod{.h1}, \src{r1}, \src{r2}} (\textbf{SPIR-V:} \texttt{OpVectorInsertDynamic / WriteMask}; \textbf{GLSL:} \texttt{dst.y = a.y + b.y;})
  \item \texttt{datatype}: Selects the lane data type (\texttt{f32}, \texttt{f16}, \texttt{v2f16}, \texttt{s32}, \texttt{u32}, \texttt{v2s16}, \texttt{v2u16}, \texttt{v4s8}, \texttt{v4u8}, \texttt{s64}, \texttt{u64}, etc.). \emph{Example:} {\scriptsize\ttfamily \mnem{IADD.v2s16} \dst{r0}, \src{r1}, \src{r2}} (\textbf{SPIR-V:} \texttt{OpTypeInt / OpTypeFloat}; \textbf{GLSL:} \texttt{i16vec2 c = a + b;})
  \item \texttt{instruction\_class}: 4-bit major instruction class selector (\texttt{bits 48:51}) partitioning the 64-bit Valhall opcode space (e.g., \texttt{FMA}, \texttt{CVT}, \texttt{SFU}, \texttt{LS}, \texttt{TEX}, \texttt{BRANCH}). \emph{Example:} {\scriptsize\ttfamily \mnem{FADD.f32} \dst{r0}, \src{r1}, \src{r2}} (\textbf{SPIR-V:} \texttt{N/A (encoding class)}; \textbf{GLSL:} \texttt{// Hardware execution pipe / decode class})
  \item \texttt{fau\_page\_index}: Selects the FAU (Fast Access Uniform) page (\texttt{0..3}) when a source operand references hardware small constants, thread IDs, or lane constants. \emph{Example:} {\scriptsize\ttfamily \mnem{FADD.f32} \dst{r0}, \src{r1}, \src{k1f}} (\textbf{SPIR-V:} \texttt{OpConstant / BuiltIn}; \textbf{GLSL:} \texttt{float y = x + 1.0;})
  \item \hyperref[enum:flow_control_m]{\texttt{flow}}: 4-bit embedded flow-control \& scoreboard wait modifier (\texttt{bits 59:62}: \texttt{none}, \texttt{wait0..wait0126}, \texttt{wait}, \texttt{reconverge}, \texttt{discard}, \texttt{end}). \emph{Example:} {\scriptsize\ttfamily \mnem{FADD.f32}\imod{.wait0.reconverge} \dst{r0}, \src{r1}, \src{r2}} (\textbf{SPIR-V:} \texttt{OpControlBarrier / OpReturn}; \textbf{GLSL:} \texttt{// Scoreboard wait + warp reconvergence})
\end{tightitemize}
\clearpage

\subsection*{FCMP\_OR}
\addcontentsline{toc}{subsection}{FCMP\_OR}
Compares floating-point \texttt{src0[i]} and \texttt{src1[i]} via \texttt{cmpf} for each SIMD element $i \in \{0, \dots, N-1\}$ ($N=1$ \texttt{f32}, $N=2$ \texttt{v2f16}), combines the predicate with \texttt{(src2[i] != 0)} using logical OR, and formats the boolean into \texttt{dest[i]} according to \texttt{result\_type}.
\par\textbf{Notes:}
\begin{tightitemize}
  \item \texttt{src2[i]} tested as \texttt{!= 0} per lane before formatting via \texttt{result\_type}.
  \item \textbf{Source references:} \texttt{kraid/ops.rs:741-790}, \texttt{encode\_v9.rs:983-991}, \texttt{opt\_fuse.rs:188-238}, \texttt{hw\_tests.rs:1318-1388}.
\end{tightitemize}
\codebox{\mnem{FCMP\_OR}\imod{\{.flow\}}\imod{\{.result\_type\}}\imod{\{.cmpf\}} \dst{dest}\imod{\{.dest\_width\}}, \src{src0}\imod{\{.widen0\}}\imod{\{.abs0\}}\imod{\{.neg0\}}, \src{src1}\imod{\{.widen1\}}\imod{\{.abs1\}}\imod{\{.neg1\}}, \src{src2}\\\mnem{FCMP\_OR}\imod{\{.flow\}}\imod{\{.result\_type\}}\imod{\{.cmpf\}} \dst{dest}\imod{\{.dest\_width\}}, \src{src0}\imod{\{.swz0\}}\imod{\{.abs0\}}\imod{\{.neg0\}}, \src{src1}\imod{\{.swz1\}}\imod{\{.abs1\}}\imod{\{.neg1\}}, \src{src2}}
\textbf{Operands:}
\begin{tightitemize}\footnotesize
  \item \texttt{\dst{dest}} (\texttt{dest\_register\_index}, \texttt{dest\_width}): Destination GPR receiving \texttt{cmpf(src0, src1) \textbar{}\textbar{} (src2 != 0)} formatted per \texttt{result\_type}.
  \item \texttt{\src{src0}}: Left-hand floating-point comparison operand.
  \item \texttt{\src{src1}}: Right-hand floating-point comparison operand.
  \item \texttt{\src{src2}}: Existing boolean/mask predicate combined with logical OR (\texttt{k0} encodes standalone \texttt{FCMP}).
\end{tightitemize}
\par\textbf{Closest SPIR-V Op:} \texttt{OpLogicalOr + OpFOrd*}
\par\textbf{Pseudo-GLSL (Operational Semantics):}
\codebox{for (int i = 0; i \textless{} N; ++i) \{\\    bool c = cmpf(mod(src0)[i], mod(src1)[i]) \textbar{}\textbar{} (src2[i] != 0);\\    dest[i] = formatResult(result\_type, c); // .i1: 1/0, .m1: \textasciitilde{}0/0, .f1: 1.0/0.0\\\}}
\textbf{Execution Context:}
\begin{tightitemize}\footnotesize
  \item \textbf{Compute stage preloads (\hyperref[sec:csf-spd]{\texttt{SPD\_k}}, \texttt{r16..r23}):} Hardware pre-fills \texttt{r55..r57} with \texttt{WorkGroupID.xyz + JOB\_OFFSET} (from \hyperref[sec:csf-regs]{\texttt{COMPUTESR::JOB\_OFFSET\_X/Y/Z}} \texttt{r34..r36}) and \texttt{r60..r62} with \texttt{LocalInvocationID.xyz} or \texttt{GlobalInvocationID.xyz} (from \texttt{WG\_SIZE} \texttt{r33} \& \texttt{JOB\_SIZE} \texttt{r37..r39}, selected by \texttt{ShaderProgram::preload}) so compute shaders can derive thread/element indices without memory loads.
  \item \textbf{Vertex stage preloads (\hyperref[sec:csf-spd]{\texttt{SPD\_0}/\texttt{SPD\_1}}, \texttt{r16..r19}):} The hardware Index Fetcher reads \hyperref[sec:csf-regs]{\texttt{IDVSSR::INDEX\_BUFFER}} (\texttt{r54:r55}) and pre-fills \texttt{r59, r60} with \texttt{VertexIndex} (\texttt{+ VERTEX\_OFFSET + GLOBAL\_ATTR\_OFF}) and \texttt{r61} with \texttt{InstanceIndex} (\texttt{+ INSTANCE\_OFFSET}) so Position and Varying shaders can pull vertex attributes directly.
  \item \textbf{Inline FAU \& Special Operands (\hyperref[sec:csf-tsd-fau]{\texttt{FAU\_k}}, \texttt{r8..r15}):} Any 8-bit source with \texttt{mode=2} reads zero-latency uniform \texttt{u0..u31} from \texttt{FAU\_k} (\texttt{slots64 = bits[63:56]}); \texttt{mode=3} reads \texttt{SmallConstant} (Table~1.2) or \texttt{FauSpecial} (\texttt{pc}, \texttt{lane\_id}, \texttt{lane\_width}, \texttt{tls\_ptr}, \texttt{wls\_ptr}, \texttt{blend\_descriptor\_0..7}).
\end{tightitemize}
\par\textbf{Worked Examples:}
\par\textbf{Example 1 (Previous condition true overrides false comparison):}
\codebox{\mnem{FCMP\_OR.f32}\imod{.gt.i1} \dst{r0}, \src{r1}, \src{r2}, \src{r3}}
With \texttt{r1 = 1.0f, r2 = 2.0f, r3 = 1}: \texttt{(1.0 \textgreater{} 2.0) \textbar{}\textbar{} (1 != 0)} is True -\textgreater{} \texttt{r0 = 1}.
\par\textbf{Example 2 (Chained range check (both false)):}
\codebox{\mnem{FCMP\_OR.f32}\imod{.gt.i1} \dst{r0}, \src{r1}, \src{r2}, \src{r3}}
With \texttt{r1 = 1.0f, r2 = 2.0f, r3 = 0}: \texttt{(1.0 \textgreater{} 2.0) \textbar{}\textbar{} false} is False -\textgreater{} \texttt{r0 = 0}.
\par\textbf{Mined SPIR-V $\to$ Valhall ISA Example (Testing float non-zero or NaN for branch/discard --- Fragment shader, Mali-G710 (v10)):}
\codebox{\textbf{; SPIR-V (input)}\\\%nz  = OpFOrdNotEqual \%bool \%x \%float\_0\\\%nan = OpFUnordNotEqual \%bool \%x \%x\\[3pt]\textbf{; Valhall ISA (malioc output)}\\FCMP.f32.discard.m1.gtlt r0, r2, k0\\FCMP.f32.m1.ne r1, r2, r2}
Standalone \texttt{FCMP} is encoded as \texttt{FCMP\_OR} with \texttt{src2 = k0}. Note the distinction: \texttt{.gtlt} tests ordered not-equal (\texttt{x \textless{} 0 \textbar{}\textbar{} x \textgreater{} 0}, false on NaN), whereas \texttt{.ne} tests unordered not-equal (true on NaN).
\begin{center}\footnotesize
\begin{tabularx}{\linewidth}{@{}lllX@{}}
\toprule
Variant & Arch & Class & Opcode Fields \\
\midrule
\texttt{f32} & v9--v14 & \texttt{R3\_CMP} & \texttt{--} \\
\texttt{v2f16} & v9--v14 & \texttt{R3\_CMP} & \texttt{--} \\
\bottomrule
\end{tabularx}
\end{center}
\begin{center}\footnotesize
\begin{tabularx}{0.88\linewidth}{@{}l l X@{}}
\toprule
Field Name & Bits & Type / Fixed Value \\
\midrule
\texttt{src0} & \texttt{0:7} & \texttt{SourceEncoding} \\
\texttt{src1} & \texttt{8:15} & \texttt{SourceEncoding} \\
\texttt{src2} & \texttt{16:23} & \texttt{SourceEncoding} \\
\texttt{combine} & \texttt{24:25} & \texttt{exact=or} \\
\hyperref[enum:widen_m]{\texttt{widen1}} & \texttt{26:27} & \texttt{widen\_m (def none)} \\
\hyperref[enum:widen_m]{\texttt{widen0}} & \texttt{28:29} & \texttt{widen\_m (def none)} \\
\hyperref[enum:cmp_result_type_m]{\texttt{result\_type}} & \texttt{30:31} & \texttt{cmp\_result\_type\_m (def i1)} \\
\hyperref[enum:cmpf_m]{\texttt{cmpf}} & \texttt{32:34} & \texttt{cmpf\_m} \\
\textit{Reserved} & \texttt{35} & \texttt{zero} \\
\hyperref[enum:neg_m]{\texttt{neg1}} & \texttt{36} & \texttt{neg\_m (def none)} \\
\hyperref[enum:abs_m]{\texttt{abs1}} & \texttt{37} & \texttt{abs\_m (def none)} \\
\hyperref[enum:neg_m]{\texttt{neg0}} & \texttt{38} & \texttt{neg\_m (def none)} \\
\hyperref[enum:abs_m]{\texttt{abs0}} & \texttt{39} & \texttt{abs\_m (def none)} \\
\texttt{dest\_register\_index} & \texttt{40:45} & \texttt{uint} \\
\hyperref[enum:dest_width_m]{\texttt{dest\_width}} & \texttt{46:47} & \texttt{dest\_width\_m} \\
\texttt{datatype} & \texttt{48:51} & \texttt{exact=f32} \\
\texttt{instruction\_class} & \texttt{52:56} & \texttt{exact=R3\_CMP} \\
\texttt{fau\_page\_index} & \texttt{57:58} & \texttt{uint} \\
\hyperref[enum:flow_control_m]{\texttt{flow}} & \texttt{59:62} & \texttt{flow\_control\_m} \\
\textit{Reserved} & \texttt{63} & \texttt{zero} \\
\bottomrule
\end{tabularx}
\end{center}
\textbf{Fields Explained:}
\begin{tightitemize}\footnotesize
  \item \texttt{src0}: Primary source operand (\texttt{r0..r63}, \texttt{u0..u63} uniform, \texttt{k*} FAU constant, or \texttt{t0..t2} temporary; \texttt{\textasciicircum{}} marks discard/last-use). \emph{Example:} {\scriptsize\ttfamily \mnem{FADD.f32} \dst{r0}, \src{r1}, \src{r2}} (\textbf{SPIR-V:} \texttt{OpFAdd (Operand 0)}; \textbf{GLSL:} \texttt{float c = a + b;})
  \item \texttt{src1}: Secondary source operand (register, uniform, or FAU constant; in shifts, supplies the shift count byte). \emph{Example:} {\scriptsize\ttfamily \mnem{LSHIFT.i32} \dst{r0}, \src{r1}, \src{r2}\imod{.b0}} (\textbf{SPIR-V:} \texttt{OpShiftLeftLogical (Shift)}; \textbf{GLSL:} \texttt{uint y = a \textless{}\textless{} s;})
  \item \texttt{src2}: Third source operand (FMA addend, \texttt{MUX}/\texttt{CSEL} selector, \texttt{IADDC} carry-in, or bitwise mask in fused shift-logic). \emph{Example:} {\scriptsize\ttfamily \mnem{FMA.f32} \dst{r0}, \src{r1}, \src{r2}, \src{r3}} (\textbf{SPIR-V:} \texttt{GLSL.std.450 Fma (Addend)}; \textbf{GLSL:} \texttt{float d = fma(a, b, c);})
  \item \texttt{combine}: Combines 16-bit halfword predicates (\texttt{h0}, \texttt{h1}, \texttt{and}, \texttt{lowbits}) before evaluating a branch. \emph{Example:} {\scriptsize\ttfamily \mnem{BRANCH}\imod{.and} \src{r0}, \imm{\#+8}} (\textbf{SPIR-V:} \texttt{OpBranchConditional (packed predicate)}; \textbf{GLSL:} \texttt{if (condA \&\& condB) \{ ... \}})
  \item \hyperref[enum:widen_m]{\texttt{widen1}}: Stage A widening extractor on \texttt{src1} (promotes narrow byte/half/word lane before Stage B). \emph{Example:} {\scriptsize\ttfamily \mnem{IMUL.u64} \dst{[r0:r1]}, \src{[r2:r3]}\imod{.w0}, \src{[r4:r5]}\imod{.w0}} (\textbf{SPIR-V:} \texttt{OpUConvert / OpUMulExtended}; \textbf{GLSL:} \texttt{uint64\_t p = uint64\_t(a) * uint64\_t(b);})
  \item \hyperref[enum:widen_m]{\texttt{widen0}}: Stage A widening extractor on \texttt{src0} (promotes FP16 \texttt{h0}/\texttt{h1} to FP32, or sign/zero-extends \texttt{b0..b3}/\texttt{h0..h1}/\texttt{w0..w1} to 32/64-bit before \texttt{abs}/\texttt{neg}). \emph{Example:} {\scriptsize\ttfamily \mnem{IADD.s32} \dst{r0}, \src{r1}\imod{.h0}, \src{r2}} (\textbf{SPIR-V:} \texttt{OpSConvert / OpFConvert}; \textbf{GLSL:} \texttt{int c = int(int16\_t(a)) + b;})
  \item \hyperref[enum:cmp_result_type_m]{\texttt{result\_type}}: Encodes boolean output representation (\texttt{i1} = 0/1, \texttt{f1} = 0.0/1.0, \texttt{m1} = 0/-1, \texttt{c} = carry). \emph{Example:} {\scriptsize\ttfamily \mnem{FCMP.f32}\imod{.lt.m1} \dst{r0}, \src{r1}, \src{r2}} (\textbf{SPIR-V:} \texttt{OpFOrdLessThan + OpSelect}; \textbf{GLSL:} \texttt{int mask = (a \textless{} b) ? -1 : 0;})
  \item \hyperref[enum:cmpf_m]{\texttt{cmpf}}: Comparison predicate (\texttt{eq}, \texttt{ne}, \texttt{lt}, \texttt{le}, \texttt{gt}, \texttt{ge}, \texttt{gtlt}, \texttt{total}). \emph{Example:} {\scriptsize\ttfamily \mnem{CSEL.f32}\imod{.lt} \dst{r0}, \src{r1}, \src{r2}, \src{r3}, \src{r4}} (\textbf{SPIR-V:} \texttt{OpFOrdLessThan + OpSelect}; \textbf{GLSL:} \texttt{float y = (a \textless{} b) ? c : d;})
  \item \hyperref[enum:neg_m]{\texttt{neg1}}: Stage A floating-point negation on \texttt{src1} (flips IEEE sign bit after swizzle/widen and \texttt{abs1}). \emph{Example:} {\scriptsize\ttfamily \mnem{FMA.f32} \dst{r0}, \src{r1}, \src{r2}\imod{.neg}, \src{r3}} (\textbf{SPIR-V:} \texttt{OpFNegate}; \textbf{GLSL:} \texttt{float d = fma(a, -b, c);})
  \item \hyperref[enum:abs_m]{\texttt{abs1}}: Stage A floating-point absolute value on \texttt{src1} (clears IEEE sign bit after swizzle/widen, before \texttt{neg1}). \emph{Example:} {\scriptsize\ttfamily \mnem{FMA.f32} \dst{r0}, \src{r1}, \src{r2}\imod{.abs}, \src{r3}} (\textbf{SPIR-V:} \texttt{GLSL.std.450 FAbs}; \textbf{GLSL:} \texttt{float d = fma(a, abs(b), c);})
  \item \hyperref[enum:neg_m]{\texttt{neg0}}: Stage A floating-point negation on \texttt{src0} (flips IEEE sign bit after swizzle/widen and \texttt{abs0}). \emph{Example:} {\scriptsize\ttfamily \mnem{FADD.f32} \dst{r0}, \src{r1}\imod{.neg}, \src{r2}} (\textbf{SPIR-V:} \texttt{OpFNegate}; \textbf{GLSL:} \texttt{float c = -a + b;})
  \item \hyperref[enum:abs_m]{\texttt{abs0}}: Stage A floating-point absolute value on \texttt{src0} (clears IEEE sign bit after swizzle/widen, before \texttt{neg0}). \emph{Example:} {\scriptsize\ttfamily \mnem{FADD.f32} \dst{r0}, \src{r1}\imod{.abs}, \src{r2}} (\textbf{SPIR-V:} \texttt{GLSL.std.450 FAbs}; \textbf{GLSL:} \texttt{float c = abs(a) + b;})
  \item \texttt{dest\_register\_index}: 6-bit destination GPR index (\texttt{r0..r63}) receiving the Stage C writeback result (even-aligned pair \texttt{[rN:rN+1]} for 64-bit ops). \emph{Example:} {\scriptsize\ttfamily \mnem{FMA.f32} \dst{r0}, \src{r1}, \src{r2}, \src{r3}} (\textbf{SPIR-V:} \texttt{SSA Result ID}; \textbf{GLSL:} \texttt{float r0 = fma(r1, r2, r3);})
  \item \hyperref[enum:dest_width_m]{\texttt{dest\_width}}: Stage C destination lane write-mask (\texttt{none} = full 32-bit, \texttt{h0} = low 16-bit, \texttt{h1} = high 16-bit without bit-shifting). \emph{Example:} {\scriptsize\ttfamily \mnem{FADD.v2f16} \dst{r0}\imod{.h1}, \src{r1}, \src{r2}} (\textbf{SPIR-V:} \texttt{OpVectorInsertDynamic / WriteMask}; \textbf{GLSL:} \texttt{dst.y = a.y + b.y;})
  \item \texttt{datatype}: Selects the lane data type (\texttt{f32}, \texttt{f16}, \texttt{v2f16}, \texttt{s32}, \texttt{u32}, \texttt{v2s16}, \texttt{v2u16}, \texttt{v4s8}, \texttt{v4u8}, \texttt{s64}, \texttt{u64}, etc.). \emph{Example:} {\scriptsize\ttfamily \mnem{IADD.v2s16} \dst{r0}, \src{r1}, \src{r2}} (\textbf{SPIR-V:} \texttt{OpTypeInt / OpTypeFloat}; \textbf{GLSL:} \texttt{i16vec2 c = a + b;})
  \item \texttt{instruction\_class}: 4-bit major instruction class selector (\texttt{bits 48:51}) partitioning the 64-bit Valhall opcode space (e.g., \texttt{FMA}, \texttt{CVT}, \texttt{SFU}, \texttt{LS}, \texttt{TEX}, \texttt{BRANCH}). \emph{Example:} {\scriptsize\ttfamily \mnem{FADD.f32} \dst{r0}, \src{r1}, \src{r2}} (\textbf{SPIR-V:} \texttt{N/A (encoding class)}; \textbf{GLSL:} \texttt{// Hardware execution pipe / decode class})
  \item \texttt{fau\_page\_index}: Selects the FAU (Fast Access Uniform) page (\texttt{0..3}) when a source operand references hardware small constants, thread IDs, or lane constants. \emph{Example:} {\scriptsize\ttfamily \mnem{FADD.f32} \dst{r0}, \src{r1}, \src{k1f}} (\textbf{SPIR-V:} \texttt{OpConstant / BuiltIn}; \textbf{GLSL:} \texttt{float y = x + 1.0;})
  \item \hyperref[enum:flow_control_m]{\texttt{flow}}: 4-bit embedded flow-control \& scoreboard wait modifier (\texttt{bits 59:62}: \texttt{none}, \texttt{wait0..wait0126}, \texttt{wait}, \texttt{reconverge}, \texttt{discard}, \texttt{end}). \emph{Example:} {\scriptsize\ttfamily \mnem{FADD.f32}\imod{.wait0.reconverge} \dst{r0}, \src{r1}, \src{r2}} (\textbf{SPIR-V:} \texttt{OpControlBarrier / OpReturn}; \textbf{GLSL:} \texttt{// Scoreboard wait + warp reconvergence})
\end{tightitemize}
\clearpage

\subsection*{ICMP}
\addcontentsline{toc}{subsection}{ICMP}
Compares Stage A widened/swizzled integer lanes \texttt{src0[i]} and \texttt{src1[i]} via \texttt{cmpf} independently for each SIMD element $i \in \{0, \dots, N-1\}$ ($N=1$ for \texttt{s32}/\texttt{u32}, $N=2$ for \texttt{v2s16}/\texttt{v2u16}, $N=4$ for \texttt{v4s8}/\texttt{v4u8} on v9--v10) and encodes the boolean into \texttt{dest[i]} via \texttt{result\_type}.
\par\textbf{Notes:}
\begin{tightitemize}
  \item \texttt{v4s8}/\texttt{v4u8} only arch \textless{}= v10 (\texttt{hw\_tests.rs:1351}).
  \item \texttt{gtlt}/\texttt{total} invalid for ints.
  \item \textbf{Source references:} \texttt{kraid/ops.rs:741-790, 1463-1522}, \texttt{encode\_v9.rs:994-1002}, \texttt{hw\_tests.rs:1318-1388}.
\end{tightitemize}
\codebox{\mnem{ICMP}\imod{\{.flow\}}\imod{\{.result\_type\}}\imod{\{.cmpf\}} \dst{dest}\imod{\{.dest\_width\}}, \src{src0}\imod{\{.widen0\}}, \src{src1}\imod{\{.widen1\}}\\\mnem{ICMP}\imod{\{.flow\}}\imod{\{.result\_type\}}\imod{\{.cmpf\}} \dst{dest}\imod{\{.dest\_width\}}, \src{src0}\imod{\{.swz0\}}, \src{src1}\imod{\{.swz1\}}\\\mnem{ICMP}\imod{\{.flow\}}\imod{\{.result\_type\}}\imod{\{.cmpf\}} \dst{dest}\imod{\{.dest\_width\}}, \src{src0}\imod{\{.lanes0\}}, \src{src1}\imod{\{.lanes1\}}}
\textbf{Operands:}
\begin{tightitemize}\footnotesize
  \item \texttt{\dst{dest}} (\texttt{dest\_register\_index}, \texttt{dest\_width}): Destination GPR receiving integer comparison result formatted per \texttt{result\_type}.
  \item \texttt{\src{src0}}: Left-hand integer comparison operand.
  \item \texttt{\src{src1}}: Right-hand integer comparison operand.
\end{tightitemize}
\par\textbf{Closest SPIR-V Op:} \texttt{OpIEqual / OpSLessThan / OpULessThan / ...}
\par\textbf{Pseudo-GLSL (Operational Semantics):}
\codebox{for (int i = 0; i \textless{} N; ++i) \{\\    bool c = cmpf(widen0(src0)[i], widen1(src1)[i]);\\    dest[i] = formatResult(result\_type, c);\\\}}
\textbf{Execution Context:}
\begin{tightitemize}\footnotesize
  \item \textbf{Compute stage preloads (\hyperref[sec:csf-spd]{\texttt{SPD\_k}}, \texttt{r16..r23}):} Hardware pre-fills \texttt{r55..r57} with \texttt{WorkGroupID.xyz + JOB\_OFFSET} (from \hyperref[sec:csf-regs]{\texttt{COMPUTESR::JOB\_OFFSET\_X/Y/Z}} \texttt{r34..r36}) and \texttt{r60..r62} with \texttt{LocalInvocationID.xyz} or \texttt{GlobalInvocationID.xyz} (from \texttt{WG\_SIZE} \texttt{r33} \& \texttt{JOB\_SIZE} \texttt{r37..r39}, selected by \texttt{ShaderProgram::preload}) so compute shaders can derive thread/element indices without memory loads.
  \item \textbf{Vertex stage preloads (\hyperref[sec:csf-spd]{\texttt{SPD\_0}/\texttt{SPD\_1}}, \texttt{r16..r19}):} The hardware Index Fetcher reads \hyperref[sec:csf-regs]{\texttt{IDVSSR::INDEX\_BUFFER}} (\texttt{r54:r55}) and pre-fills \texttt{r59, r60} with \texttt{VertexIndex} (\texttt{+ VERTEX\_OFFSET + GLOBAL\_ATTR\_OFF}) and \texttt{r61} with \texttt{InstanceIndex} (\texttt{+ INSTANCE\_OFFSET}) so Position and Varying shaders can pull vertex attributes directly.
  \item \textbf{Inline FAU \& Special Operands (\hyperref[sec:csf-tsd-fau]{\texttt{FAU\_k}}, \texttt{r8..r15}):} Any 8-bit source with \texttt{mode=2} reads zero-latency uniform \texttt{u0..u31} from \texttt{FAU\_k} (\texttt{slots64 = bits[63:56]}); \texttt{mode=3} reads \texttt{SmallConstant} (Table~1.2) or \texttt{FauSpecial} (\texttt{pc}, \texttt{lane\_id}, \texttt{lane\_width}, \texttt{tls\_ptr}, \texttt{wls\_ptr}, \texttt{blend\_descriptor\_0..7}).
\end{tightitemize}
\par\textbf{Worked Examples:}
\par\textbf{Example 1 (Signed comparison with negative value):}
\codebox{\mnem{ICMP.s32}\imod{.ge.i1} \dst{r0}, \src{r1}, \src{r2}}
With \texttt{r1 = 0xFFFFFFFF} (\texttt{-1}), \texttt{r2 = 0}: signed \texttt{-1 \textgreater{}= 0} is False -\textgreater{} \texttt{r0 = 0}.
\par\textbf{Example 2 (Unsigned comparison on same bit pattern):}
\codebox{\mnem{ICMP.u32}\imod{.ge.m1} \dst{r0}, \src{r1}, \src{r2}}
With \texttt{r1 = 0xFFFFFFFF} (\texttt{4294967295u}), \texttt{r2 = 0}: unsigned \texttt{\textgreater{}=} is True -\textgreater{} \texttt{r0 = 0xFFFFFFFF} under \texttt{.m1}.

\begin{center}\footnotesize
\begin{tabularx}{\linewidth}{@{}lllX@{}}
\toprule
Variant & Arch & Class & Opcode Fields \\
\midrule
\texttt{s32} & v9--v14 & \texttt{R3\_CMP} & \texttt{--} \\
\texttt{u32} & v9--v14 & \texttt{R3\_CMP} & \texttt{--} \\
\texttt{v2s16} & v9--v14 & \texttt{R3\_CMP} & \texttt{--} \\
\texttt{v2u16} & v9--v14 & \texttt{R3\_CMP} & \texttt{--} \\
\texttt{v4s8} & v9--v14 & \texttt{R3\_CMP} & \texttt{--} \\
\texttt{v4u8} & v9--v14 & \texttt{R3\_CMP} & \texttt{--} \\
\bottomrule
\end{tabularx}
\end{center}
\begin{center}\footnotesize
\begin{tabularx}{0.88\linewidth}{@{}l l X@{}}
\toprule
Field Name & Bits & Type / Fixed Value \\
\midrule
\texttt{src0} & \texttt{0:7} & \texttt{SourceEncoding} \\
\texttt{src1} & \texttt{8:15} & \texttt{SourceEncoding} \\
\texttt{src2} & \texttt{16:23} & \texttt{exact=192} \\
\texttt{combine} & \texttt{24:25} & \texttt{exact=or} \\
\hyperref[enum:widen_int_m]{\texttt{widen1}} & \texttt{26:29} & \texttt{widen\_int\_m (def none)} \\
\hyperref[enum:cmp_result_type_m]{\texttt{result\_type}} & \texttt{30:31} & \texttt{cmp\_result\_type\_m (def i1)} \\
\hyperref[enum:cmpf_m]{\texttt{cmpf}} & \texttt{32:34} & \texttt{cmpf\_m} \\
\textit{Reserved} & \texttt{35} & \texttt{zero} \\
\hyperref[enum:widen_int_m]{\texttt{widen0}} & \texttt{36:39} & \texttt{widen\_int\_m (def none)} \\
\texttt{dest\_register\_index} & \texttt{40:45} & \texttt{uint} \\
\hyperref[enum:dest_width_m]{\texttt{dest\_width}} & \texttt{46:47} & \texttt{dest\_width\_m} \\
\texttt{datatype} & \texttt{48:51} & \texttt{exact=s32} \\
\texttt{instruction\_class} & \texttt{52:56} & \texttt{exact=R3\_CMP} \\
\texttt{fau\_page\_index} & \texttt{57:58} & \texttt{uint} \\
\hyperref[enum:flow_control_m]{\texttt{flow}} & \texttt{59:62} & \texttt{flow\_control\_m} \\
\textit{Reserved} & \texttt{63} & \texttt{zero} \\
\bottomrule
\end{tabularx}
\end{center}
\textbf{Fields Explained:}
\begin{tightitemize}\footnotesize
  \item \texttt{src0}: Primary source operand (\texttt{r0..r63}, \texttt{u0..u63} uniform, \texttt{k*} FAU constant, or \texttt{t0..t2} temporary; \texttt{\textasciicircum{}} marks discard/last-use). \emph{Example:} {\scriptsize\ttfamily \mnem{FADD.f32} \dst{r0}, \src{r1}, \src{r2}} (\textbf{SPIR-V:} \texttt{OpFAdd (Operand 0)}; \textbf{GLSL:} \texttt{float c = a + b;})
  \item \texttt{src1}: Secondary source operand (register, uniform, or FAU constant; in shifts, supplies the shift count byte). \emph{Example:} {\scriptsize\ttfamily \mnem{LSHIFT.i32} \dst{r0}, \src{r1}, \src{r2}\imod{.b0}} (\textbf{SPIR-V:} \texttt{OpShiftLeftLogical (Shift)}; \textbf{GLSL:} \texttt{uint y = a \textless{}\textless{} s;})
  \item \texttt{src2}: Third source operand (FMA addend, \texttt{MUX}/\texttt{CSEL} selector, \texttt{IADDC} carry-in, or bitwise mask in fused shift-logic). \emph{Example:} {\scriptsize\ttfamily \mnem{FMA.f32} \dst{r0}, \src{r1}, \src{r2}, \src{r3}} (\textbf{SPIR-V:} \texttt{GLSL.std.450 Fma (Addend)}; \textbf{GLSL:} \texttt{float d = fma(a, b, c);})
  \item \texttt{combine}: Combines 16-bit halfword predicates (\texttt{h0}, \texttt{h1}, \texttt{and}, \texttt{lowbits}) before evaluating a branch. \emph{Example:} {\scriptsize\ttfamily \mnem{BRANCH}\imod{.and} \src{r0}, \imm{\#+8}} (\textbf{SPIR-V:} \texttt{OpBranchConditional (packed predicate)}; \textbf{GLSL:} \texttt{if (condA \&\& condB) \{ ... \}})
  \item \hyperref[enum:widen_int_m]{\texttt{widen1}}: Stage A widening extractor on \texttt{src1} (promotes narrow byte/half/word lane before Stage B). \emph{Example:} {\scriptsize\ttfamily \mnem{IMUL.u64} \dst{[r0:r1]}, \src{[r2:r3]}\imod{.w0}, \src{[r4:r5]}\imod{.w0}} (\textbf{SPIR-V:} \texttt{OpUConvert / OpUMulExtended}; \textbf{GLSL:} \texttt{uint64\_t p = uint64\_t(a) * uint64\_t(b);})
  \item \hyperref[enum:cmp_result_type_m]{\texttt{result\_type}}: Encodes boolean output representation (\texttt{i1} = 0/1, \texttt{f1} = 0.0/1.0, \texttt{m1} = 0/-1, \texttt{c} = carry). \emph{Example:} {\scriptsize\ttfamily \mnem{FCMP.f32}\imod{.lt.m1} \dst{r0}, \src{r1}, \src{r2}} (\textbf{SPIR-V:} \texttt{OpFOrdLessThan + OpSelect}; \textbf{GLSL:} \texttt{int mask = (a \textless{} b) ? -1 : 0;})
  \item \hyperref[enum:cmpf_m]{\texttt{cmpf}}: Comparison predicate (\texttt{eq}, \texttt{ne}, \texttt{lt}, \texttt{le}, \texttt{gt}, \texttt{ge}, \texttt{gtlt}, \texttt{total}). \emph{Example:} {\scriptsize\ttfamily \mnem{CSEL.f32}\imod{.lt} \dst{r0}, \src{r1}, \src{r2}, \src{r3}, \src{r4}} (\textbf{SPIR-V:} \texttt{OpFOrdLessThan + OpSelect}; \textbf{GLSL:} \texttt{float y = (a \textless{} b) ? c : d;})
  \item \hyperref[enum:widen_int_m]{\texttt{widen0}}: Stage A widening extractor on \texttt{src0} (promotes FP16 \texttt{h0}/\texttt{h1} to FP32, or sign/zero-extends \texttt{b0..b3}/\texttt{h0..h1}/\texttt{w0..w1} to 32/64-bit before \texttt{abs}/\texttt{neg}). \emph{Example:} {\scriptsize\ttfamily \mnem{IADD.s32} \dst{r0}, \src{r1}\imod{.h0}, \src{r2}} (\textbf{SPIR-V:} \texttt{OpSConvert / OpFConvert}; \textbf{GLSL:} \texttt{int c = int(int16\_t(a)) + b;})
  \item \texttt{dest\_register\_index}: 6-bit destination GPR index (\texttt{r0..r63}) receiving the Stage C writeback result (even-aligned pair \texttt{[rN:rN+1]} for 64-bit ops). \emph{Example:} {\scriptsize\ttfamily \mnem{FMA.f32} \dst{r0}, \src{r1}, \src{r2}, \src{r3}} (\textbf{SPIR-V:} \texttt{SSA Result ID}; \textbf{GLSL:} \texttt{float r0 = fma(r1, r2, r3);})
  \item \hyperref[enum:dest_width_m]{\texttt{dest\_width}}: Stage C destination lane write-mask (\texttt{none} = full 32-bit, \texttt{h0} = low 16-bit, \texttt{h1} = high 16-bit without bit-shifting). \emph{Example:} {\scriptsize\ttfamily \mnem{FADD.v2f16} \dst{r0}\imod{.h1}, \src{r1}, \src{r2}} (\textbf{SPIR-V:} \texttt{OpVectorInsertDynamic / WriteMask}; \textbf{GLSL:} \texttt{dst.y = a.y + b.y;})
  \item \texttt{datatype}: Selects the lane data type (\texttt{f32}, \texttt{f16}, \texttt{v2f16}, \texttt{s32}, \texttt{u32}, \texttt{v2s16}, \texttt{v2u16}, \texttt{v4s8}, \texttt{v4u8}, \texttt{s64}, \texttt{u64}, etc.). \emph{Example:} {\scriptsize\ttfamily \mnem{IADD.v2s16} \dst{r0}, \src{r1}, \src{r2}} (\textbf{SPIR-V:} \texttt{OpTypeInt / OpTypeFloat}; \textbf{GLSL:} \texttt{i16vec2 c = a + b;})
  \item \texttt{instruction\_class}: 4-bit major instruction class selector (\texttt{bits 48:51}) partitioning the 64-bit Valhall opcode space (e.g., \texttt{FMA}, \texttt{CVT}, \texttt{SFU}, \texttt{LS}, \texttt{TEX}, \texttt{BRANCH}). \emph{Example:} {\scriptsize\ttfamily \mnem{FADD.f32} \dst{r0}, \src{r1}, \src{r2}} (\textbf{SPIR-V:} \texttt{N/A (encoding class)}; \textbf{GLSL:} \texttt{// Hardware execution pipe / decode class})
  \item \texttt{fau\_page\_index}: Selects the FAU (Fast Access Uniform) page (\texttt{0..3}) when a source operand references hardware small constants, thread IDs, or lane constants. \emph{Example:} {\scriptsize\ttfamily \mnem{FADD.f32} \dst{r0}, \src{r1}, \src{k1f}} (\textbf{SPIR-V:} \texttt{OpConstant / BuiltIn}; \textbf{GLSL:} \texttt{float y = x + 1.0;})
  \item \hyperref[enum:flow_control_m]{\texttt{flow}}: 4-bit embedded flow-control \& scoreboard wait modifier (\texttt{bits 59:62}: \texttt{none}, \texttt{wait0..wait0126}, \texttt{wait}, \texttt{reconverge}, \texttt{discard}, \texttt{end}). \emph{Example:} {\scriptsize\ttfamily \mnem{FADD.f32}\imod{.wait0.reconverge} \dst{r0}, \src{r1}, \src{r2}} (\textbf{SPIR-V:} \texttt{OpControlBarrier / OpReturn}; \textbf{GLSL:} \texttt{// Scoreboard wait + warp reconvergence})
\end{tightitemize}
\clearpage

\subsection*{ICMP\_AND}
\addcontentsline{toc}{subsection}{ICMP\_AND}
For each SIMD element $i \in \{0, \dots, N-1\}$ ($N=1$ for \texttt{s32}/\texttt{u32}, $N=2$ for \texttt{v2s16}/\texttt{v2u16}, $N=4$ for \texttt{v4s8}/\texttt{v4u8} on v9--v10), compares Stage A widened/swizzled \texttt{src0[i]} and \texttt{src1[i]} using predicate \texttt{cmpf}, combines the boolean with \texttt{(src2[i] != 0)} via logical AND, and formats the result into \texttt{dest[i]} according to \texttt{result\_type} (\texttt{i1} $\to$ \texttt{1}, \texttt{m1} $\to$ all-ones, \texttt{f1} $\to$ \texttt{1.0}, false $\to$ \texttt{0}).
\par\textbf{Notes:}
\begin{tightitemize}
  \item \texttt{v4s8}/\texttt{v4u8} only arch \textless{}= v10.
  \item \textbf{Source references:} \texttt{kraid/ops.rs:741-790}, \texttt{encode\_v9.rs:1003-1011}, \texttt{hw\_tests.rs:1318-1388}.
\end{tightitemize}
\codebox{\mnem{ICMP\_AND}\imod{\{.flow\}}\imod{\{.result\_type\}}\imod{\{.cmpf\}} \dst{dest}\imod{\{.dest\_width\}}, \src{src0}\imod{\{.widen0\}}, \src{src1}\imod{\{.widen1\}}, \src{src2}\\\mnem{ICMP\_AND}\imod{\{.flow\}}\imod{\{.result\_type\}}\imod{\{.cmpf\}} \dst{dest}\imod{\{.dest\_width\}}, \src{src0}\imod{\{.swz0\}}, \src{src1}\imod{\{.swz1\}}, \src{src2}\\\mnem{ICMP\_AND}\imod{\{.flow\}}\imod{\{.result\_type\}}\imod{\{.cmpf\}} \dst{dest}\imod{\{.dest\_width\}}, \src{src0}\imod{\{.lanes0\}}, \src{src1}\imod{\{.lanes1\}}, \src{src2}}
\textbf{Operands:}
\begin{tightitemize}\footnotesize
  \item \texttt{\dst{dest}} (\texttt{dest\_register\_index}, \texttt{dest\_width}): Destination GPR receiving \texttt{cmpf(src0, src1) \&\& (src2 != 0)} formatted per \texttt{result\_type}.
  \item \texttt{\src{src0}}: Left-hand integer comparison operand.
  \item \texttt{\src{src1}}: Right-hand integer comparison operand.
  \item \texttt{\src{src2}}: Existing boolean/mask predicate combined with logical AND.
\end{tightitemize}
\par\textbf{Closest SPIR-V Op:} \texttt{OpLogicalAnd + OpI*/OpS*/OpU*}
\par\textbf{Pseudo-GLSL (Operational Semantics):}
\codebox{for (int i = 0; i \textless{} N; ++i) \{\\    bool c = cmpf(widen0(src0)[i], widen1(src1)[i]) \&\& (src2[i] != 0);\\    dest[i] = formatResult(result\_type, c);\\\}}
\textbf{Execution Context:}
\begin{tightitemize}\footnotesize
  \item \textbf{Compute stage preloads (\hyperref[sec:csf-spd]{\texttt{SPD\_k}}, \texttt{r16..r23}):} Hardware pre-fills \texttt{r55..r57} with \texttt{WorkGroupID.xyz + JOB\_OFFSET} (from \hyperref[sec:csf-regs]{\texttt{COMPUTESR::JOB\_OFFSET\_X/Y/Z}} \texttt{r34..r36}) and \texttt{r60..r62} with \texttt{LocalInvocationID.xyz} or \texttt{GlobalInvocationID.xyz} (from \texttt{WG\_SIZE} \texttt{r33} \& \texttt{JOB\_SIZE} \texttt{r37..r39}, selected by \texttt{ShaderProgram::preload}) so compute shaders can derive thread/element indices without memory loads.
  \item \textbf{Vertex stage preloads (\hyperref[sec:csf-spd]{\texttt{SPD\_0}/\texttt{SPD\_1}}, \texttt{r16..r19}):} The hardware Index Fetcher reads \hyperref[sec:csf-regs]{\texttt{IDVSSR::INDEX\_BUFFER}} (\texttt{r54:r55}) and pre-fills \texttt{r59, r60} with \texttt{VertexIndex} (\texttt{+ VERTEX\_OFFSET + GLOBAL\_ATTR\_OFF}) and \texttt{r61} with \texttt{InstanceIndex} (\texttt{+ INSTANCE\_OFFSET}) so Position and Varying shaders can pull vertex attributes directly.
  \item \textbf{Inline FAU \& Special Operands (\hyperref[sec:csf-tsd-fau]{\texttt{FAU\_k}}, \texttt{r8..r15}):} Any 8-bit source with \texttt{mode=2} reads zero-latency uniform \texttt{u0..u31} from \texttt{FAU\_k} (\texttt{slots64 = bits[63:56]}); \texttt{mode=3} reads \texttt{SmallConstant} (Table~1.2) or \texttt{FauSpecial} (\texttt{pc}, \texttt{lane\_id}, \texttt{lane\_width}, \texttt{tls\_ptr}, \texttt{wls\_ptr}, \texttt{blend\_descriptor\_0..7}).
\end{tightitemize}
\par\textbf{Worked Examples:}
\par\textbf{Example 1 (Chained equality (both true)):}
\codebox{\mnem{ICMP\_AND.u32}\imod{.eq.m1} \dst{r0}, \src{r1}, \src{r2}, \src{r3}}
With \texttt{r1 = 5, r2 = 5, r3 = 1}: \texttt{(5 == 5) \&\& (1 != 0)} is True -\textgreater{} \texttt{r0 = 0xFFFFFFFF}.
\par\textbf{Example 2 (64-bit equality high-word step (false)):}
\codebox{\mnem{ICMP\_AND.u32}\imod{.eq.i1} \dst{r0}, \src{r1}, \src{r2}, \src{r3}}
With \texttt{r1 = 0, r2 = 1, r3 = 1} (low words matched, high words differ): returns \texttt{r0 = 0}.
\par\textbf{Mined SPIR-V $\to$ Valhall ISA Example (Chained bit-test and uniform-flag guard --- Vertex shader, Mali-G710 (v10)):}
\codebox{\textbf{; SPIR-V (input)}\\\%b = OpBitwiseAnd \%int \%x \%int\_2\\\%c = OpLogicalAnd \%bool (OpIEqual \%b \%int\_0) \%flag\\[3pt]\textbf{; Valhall ISA (malioc output)}\\AND.i32 r1, k03020100.b2, r0\\ICMP\_AND.u32.m1.eq r2, r1, k0, u3.w1}
\texttt{ICMP\_AND.u32.m1.eq} evaluates \texttt{(r1 == 0) \&\& (u3.w1 != 0)} and outputs an all-ones/zero mask directly into \texttt{r2}.
\begin{center}\footnotesize
\begin{tabularx}{\linewidth}{@{}lllX@{}}
\toprule
Variant & Arch & Class & Opcode Fields \\
\midrule
\texttt{s32} & v9--v14 & \texttt{R3\_CMP} & \texttt{--} \\
\texttt{u32} & v9--v14 & \texttt{R3\_CMP} & \texttt{--} \\
\texttt{v2s16} & v9--v14 & \texttt{R3\_CMP} & \texttt{--} \\
\texttt{v2u16} & v9--v14 & \texttt{R3\_CMP} & \texttt{--} \\
\texttt{v4s8} & v9--v14 & \texttt{R3\_CMP} & \texttt{--} \\
\texttt{v4u8} & v9--v14 & \texttt{R3\_CMP} & \texttt{--} \\
\bottomrule
\end{tabularx}
\end{center}
\begin{center}\footnotesize
\begin{tabularx}{0.88\linewidth}{@{}l l X@{}}
\toprule
Field Name & Bits & Type / Fixed Value \\
\midrule
\texttt{src0} & \texttt{0:7} & \texttt{SourceEncoding} \\
\texttt{src1} & \texttt{8:15} & \texttt{SourceEncoding} \\
\texttt{src2} & \texttt{16:23} & \texttt{SourceEncoding} \\
\texttt{combine} & \texttt{24:25} & \texttt{exact=and} \\
\hyperref[enum:widen_int_m]{\texttt{widen1}} & \texttt{26:29} & \texttt{widen\_int\_m (def none)} \\
\hyperref[enum:cmp_result_type_m]{\texttt{result\_type}} & \texttt{30:31} & \texttt{cmp\_result\_type\_m (def i1)} \\
\hyperref[enum:cmpf_m]{\texttt{cmpf}} & \texttt{32:34} & \texttt{cmpf\_m} \\
\textit{Reserved} & \texttt{35} & \texttt{zero} \\
\hyperref[enum:widen_int_m]{\texttt{widen0}} & \texttt{36:39} & \texttt{widen\_int\_m (def none)} \\
\texttt{dest\_register\_index} & \texttt{40:45} & \texttt{uint} \\
\hyperref[enum:dest_width_m]{\texttt{dest\_width}} & \texttt{46:47} & \texttt{dest\_width\_m} \\
\texttt{datatype} & \texttt{48:51} & \texttt{exact=s32} \\
\texttt{instruction\_class} & \texttt{52:56} & \texttt{exact=R3\_CMP} \\
\texttt{fau\_page\_index} & \texttt{57:58} & \texttt{uint} \\
\hyperref[enum:flow_control_m]{\texttt{flow}} & \texttt{59:62} & \texttt{flow\_control\_m} \\
\textit{Reserved} & \texttt{63} & \texttt{zero} \\
\bottomrule
\end{tabularx}
\end{center}
\textbf{Fields Explained:}
\begin{tightitemize}\footnotesize
  \item \texttt{src0}: Primary source operand (\texttt{r0..r63}, \texttt{u0..u63} uniform, \texttt{k*} FAU constant, or \texttt{t0..t2} temporary; \texttt{\textasciicircum{}} marks discard/last-use). \emph{Example:} {\scriptsize\ttfamily \mnem{FADD.f32} \dst{r0}, \src{r1}, \src{r2}} (\textbf{SPIR-V:} \texttt{OpFAdd (Operand 0)}; \textbf{GLSL:} \texttt{float c = a + b;})
  \item \texttt{src1}: Secondary source operand (register, uniform, or FAU constant; in shifts, supplies the shift count byte). \emph{Example:} {\scriptsize\ttfamily \mnem{LSHIFT.i32} \dst{r0}, \src{r1}, \src{r2}\imod{.b0}} (\textbf{SPIR-V:} \texttt{OpShiftLeftLogical (Shift)}; \textbf{GLSL:} \texttt{uint y = a \textless{}\textless{} s;})
  \item \texttt{src2}: Third source operand (FMA addend, \texttt{MUX}/\texttt{CSEL} selector, \texttt{IADDC} carry-in, or bitwise mask in fused shift-logic). \emph{Example:} {\scriptsize\ttfamily \mnem{FMA.f32} \dst{r0}, \src{r1}, \src{r2}, \src{r3}} (\textbf{SPIR-V:} \texttt{GLSL.std.450 Fma (Addend)}; \textbf{GLSL:} \texttt{float d = fma(a, b, c);})
  \item \texttt{combine}: Combines 16-bit halfword predicates (\texttt{h0}, \texttt{h1}, \texttt{and}, \texttt{lowbits}) before evaluating a branch. \emph{Example:} {\scriptsize\ttfamily \mnem{BRANCH}\imod{.and} \src{r0}, \imm{\#+8}} (\textbf{SPIR-V:} \texttt{OpBranchConditional (packed predicate)}; \textbf{GLSL:} \texttt{if (condA \&\& condB) \{ ... \}})
  \item \hyperref[enum:widen_int_m]{\texttt{widen1}}: Stage A widening extractor on \texttt{src1} (promotes narrow byte/half/word lane before Stage B). \emph{Example:} {\scriptsize\ttfamily \mnem{IMUL.u64} \dst{[r0:r1]}, \src{[r2:r3]}\imod{.w0}, \src{[r4:r5]}\imod{.w0}} (\textbf{SPIR-V:} \texttt{OpUConvert / OpUMulExtended}; \textbf{GLSL:} \texttt{uint64\_t p = uint64\_t(a) * uint64\_t(b);})
  \item \hyperref[enum:cmp_result_type_m]{\texttt{result\_type}}: Encodes boolean output representation (\texttt{i1} = 0/1, \texttt{f1} = 0.0/1.0, \texttt{m1} = 0/-1, \texttt{c} = carry). \emph{Example:} {\scriptsize\ttfamily \mnem{FCMP.f32}\imod{.lt.m1} \dst{r0}, \src{r1}, \src{r2}} (\textbf{SPIR-V:} \texttt{OpFOrdLessThan + OpSelect}; \textbf{GLSL:} \texttt{int mask = (a \textless{} b) ? -1 : 0;})
  \item \hyperref[enum:cmpf_m]{\texttt{cmpf}}: Comparison predicate (\texttt{eq}, \texttt{ne}, \texttt{lt}, \texttt{le}, \texttt{gt}, \texttt{ge}, \texttt{gtlt}, \texttt{total}). \emph{Example:} {\scriptsize\ttfamily \mnem{CSEL.f32}\imod{.lt} \dst{r0}, \src{r1}, \src{r2}, \src{r3}, \src{r4}} (\textbf{SPIR-V:} \texttt{OpFOrdLessThan + OpSelect}; \textbf{GLSL:} \texttt{float y = (a \textless{} b) ? c : d;})
  \item \hyperref[enum:widen_int_m]{\texttt{widen0}}: Stage A widening extractor on \texttt{src0} (promotes FP16 \texttt{h0}/\texttt{h1} to FP32, or sign/zero-extends \texttt{b0..b3}/\texttt{h0..h1}/\texttt{w0..w1} to 32/64-bit before \texttt{abs}/\texttt{neg}). \emph{Example:} {\scriptsize\ttfamily \mnem{IADD.s32} \dst{r0}, \src{r1}\imod{.h0}, \src{r2}} (\textbf{SPIR-V:} \texttt{OpSConvert / OpFConvert}; \textbf{GLSL:} \texttt{int c = int(int16\_t(a)) + b;})
  \item \texttt{dest\_register\_index}: 6-bit destination GPR index (\texttt{r0..r63}) receiving the Stage C writeback result (even-aligned pair \texttt{[rN:rN+1]} for 64-bit ops). \emph{Example:} {\scriptsize\ttfamily \mnem{FMA.f32} \dst{r0}, \src{r1}, \src{r2}, \src{r3}} (\textbf{SPIR-V:} \texttt{SSA Result ID}; \textbf{GLSL:} \texttt{float r0 = fma(r1, r2, r3);})
  \item \hyperref[enum:dest_width_m]{\texttt{dest\_width}}: Stage C destination lane write-mask (\texttt{none} = full 32-bit, \texttt{h0} = low 16-bit, \texttt{h1} = high 16-bit without bit-shifting). \emph{Example:} {\scriptsize\ttfamily \mnem{FADD.v2f16} \dst{r0}\imod{.h1}, \src{r1}, \src{r2}} (\textbf{SPIR-V:} \texttt{OpVectorInsertDynamic / WriteMask}; \textbf{GLSL:} \texttt{dst.y = a.y + b.y;})
  \item \texttt{datatype}: Selects the lane data type (\texttt{f32}, \texttt{f16}, \texttt{v2f16}, \texttt{s32}, \texttt{u32}, \texttt{v2s16}, \texttt{v2u16}, \texttt{v4s8}, \texttt{v4u8}, \texttt{s64}, \texttt{u64}, etc.). \emph{Example:} {\scriptsize\ttfamily \mnem{IADD.v2s16} \dst{r0}, \src{r1}, \src{r2}} (\textbf{SPIR-V:} \texttt{OpTypeInt / OpTypeFloat}; \textbf{GLSL:} \texttt{i16vec2 c = a + b;})
  \item \texttt{instruction\_class}: 4-bit major instruction class selector (\texttt{bits 48:51}) partitioning the 64-bit Valhall opcode space (e.g., \texttt{FMA}, \texttt{CVT}, \texttt{SFU}, \texttt{LS}, \texttt{TEX}, \texttt{BRANCH}). \emph{Example:} {\scriptsize\ttfamily \mnem{FADD.f32} \dst{r0}, \src{r1}, \src{r2}} (\textbf{SPIR-V:} \texttt{N/A (encoding class)}; \textbf{GLSL:} \texttt{// Hardware execution pipe / decode class})
  \item \texttt{fau\_page\_index}: Selects the FAU (Fast Access Uniform) page (\texttt{0..3}) when a source operand references hardware small constants, thread IDs, or lane constants. \emph{Example:} {\scriptsize\ttfamily \mnem{FADD.f32} \dst{r0}, \src{r1}, \src{k1f}} (\textbf{SPIR-V:} \texttt{OpConstant / BuiltIn}; \textbf{GLSL:} \texttt{float y = x + 1.0;})
  \item \hyperref[enum:flow_control_m]{\texttt{flow}}: 4-bit embedded flow-control \& scoreboard wait modifier (\texttt{bits 59:62}: \texttt{none}, \texttt{wait0..wait0126}, \texttt{wait}, \texttt{reconverge}, \texttt{discard}, \texttt{end}). \emph{Example:} {\scriptsize\ttfamily \mnem{FADD.f32}\imod{.wait0.reconverge} \dst{r0}, \src{r1}, \src{r2}} (\textbf{SPIR-V:} \texttt{OpControlBarrier / OpReturn}; \textbf{GLSL:} \texttt{// Scoreboard wait + warp reconvergence})
\end{tightitemize}
\clearpage

\subsection*{ICMP\_MULTI}
\addcontentsline{toc}{subsection}{ICMP\_MULTI}
Multi-word chained integer compare (\texttt{s32}/\texttt{u32}): \texttt{ord = cmp(src0, src1).then((src2 as i32).cmp(\&0))}; if \texttt{result\_type == c} -\textgreater{} \texttt{ord as i8 as u32} (\texttt{-1=0xFFFFFFFF}, \texttt{0}, \texttt{1}), else \texttt{res\_type(cmpf(ord))}.
\par\textbf{Notes:}
\begin{tightitemize}
  \item \texttt{ICMP\_MULTI.s32.c} does NOT exist (\texttt{kraid/nir.rs:1192}, \texttt{hw\_tests.rs:1928}).
  \item Used for 64-bit comparisons (\texttt{nir.rs:1186-1209}).
  \item \textbf{Source references:} \texttt{kraid/ops.rs:2303-2349}, \texttt{encode\_v9.rs:1795-1818}, \texttt{nir.rs:1186-1209}, \texttt{hw\_tests.rs:1905-1949}.
\end{tightitemize}
\codebox{\mnem{ICMP\_MULTI}\imod{\{.flow\}}\imod{\{.result\_type\}}\imod{\{.cmpf\}} \dst{dest}\imod{\{.dest\_width\}}, \src{src0}\imod{\{.widen0\}}, \src{src1}\imod{\{.widen1\}}, \src{src2}}
\textbf{Operands:}
\begin{tightitemize}\footnotesize
  \item \texttt{\dst{dest}} (\texttt{dest\_register\_index}, \texttt{dest\_width}): Destination GPR receiving 3-way carry (\texttt{-1}/\texttt{0}/\texttt{+1} with \texttt{.c}) or chained multi-word comparison result.
  \item \texttt{\src{src0}}: Left-hand 32-bit word of the multi-word integer.
  \item \texttt{\src{src1}}: Right-hand 32-bit word of the multi-word integer.
  \item \texttt{\src{src2}}: Incoming 3-way comparison carry from the less-significant word (\texttt{0} for lowest word).
\end{tightitemize}
\par\textbf{Closest SPIR-V Op:} \texttt{OpULessThan / OpSLessThan (64-bit chained)}
\par\textbf{Pseudo-GLSL (Operational Semantics):}
\codebox{int ord = (src0 \textless{} src1) ? -1 : ((src0 \textgreater{} src1) ? +1 : int(src2));\\if (result\_type == c) dest = uint(ord); // 3-way carry to next word\\else                  dest = formatResult(result\_type, matchesCmpf(cmpf, ord));}
\textbf{Execution Context:}
\begin{tightitemize}\footnotesize
  \item \textbf{Compute stage preloads (\hyperref[sec:csf-spd]{\texttt{SPD\_k}}, \texttt{r16..r23}):} Hardware pre-fills \texttt{r55..r57} with \texttt{WorkGroupID.xyz + JOB\_OFFSET} (from \hyperref[sec:csf-regs]{\texttt{COMPUTESR::JOB\_OFFSET\_X/Y/Z}} \texttt{r34..r36}) and \texttt{r60..r62} with \texttt{LocalInvocationID.xyz} or \texttt{GlobalInvocationID.xyz} (from \texttt{WG\_SIZE} \texttt{r33} \& \texttt{JOB\_SIZE} \texttt{r37..r39}, selected by \texttt{ShaderProgram::preload}) so compute shaders can derive thread/element indices without memory loads.
  \item \textbf{Vertex stage preloads (\hyperref[sec:csf-spd]{\texttt{SPD\_0}/\texttt{SPD\_1}}, \texttt{r16..r19}):} The hardware Index Fetcher reads \hyperref[sec:csf-regs]{\texttt{IDVSSR::INDEX\_BUFFER}} (\texttt{r54:r55}) and pre-fills \texttt{r59, r60} with \texttt{VertexIndex} (\texttt{+ VERTEX\_OFFSET + GLOBAL\_ATTR\_OFF}) and \texttt{r61} with \texttt{InstanceIndex} (\texttt{+ INSTANCE\_OFFSET}) so Position and Varying shaders can pull vertex attributes directly.
  \item \textbf{Inline FAU \& Special Operands (\hyperref[sec:csf-tsd-fau]{\texttt{FAU\_k}}, \texttt{r8..r15}):} Any 8-bit source with \texttt{mode=2} reads zero-latency uniform \texttt{u0..u31} from \texttt{FAU\_k} (\texttt{slots64 = bits[63:56]}); \texttt{mode=3} reads \texttt{SmallConstant} (Table~1.2) or \texttt{FauSpecial} (\texttt{pc}, \texttt{lane\_id}, \texttt{lane\_width}, \texttt{tls\_ptr}, \texttt{wls\_ptr}, \texttt{blend\_descriptor\_0..7}).
\end{tightitemize}
\par\textbf{Worked Examples:}
\par\textbf{Example 1 (64-bit unsigned `\textless{}` low-word step (`.c` 3-way carry)):}
\codebox{\mnem{ICMP\_MULTI.u32}\imod{.lt.c} \dst{r0}, \src{r1}, \src{r2}, \src{k0}}
With \texttt{r1 = 3, r2 = 5, acc = 0}: \texttt{3 \textless{} 5} outputs 3-way carry \texttt{r0 = -1} (\texttt{0xFFFFFFFF}) for the high-word step.
\par\textbf{Example 2 (64-bit unsigned `\textless{}` high-word tie-breaker):}
\codebox{\mnem{ICMP\_MULTI.u32}\imod{.lt.i1} \dst{r0}, \src{r3}, \src{r4}, \src{r0}}
With equal high words \texttt{r3 = 7, r4 = 7} and low-word carry \texttt{r0 = -1}: tie-breaks using carry \texttt{\textless{} 0} -\textgreater{} returns \texttt{r0 = 1} (\texttt{true}).

\begin{center}\footnotesize
\begin{tabularx}{\linewidth}{@{}lllX@{}}
\toprule
Variant & Arch & Class & Opcode Fields \\
\midrule
\texttt{s32} & v9--v14 & \texttt{R3\_CMP} & \texttt{--} \\
\texttt{u32} & v9--v14 & \texttt{R3\_CMP} & \texttt{--} \\
\bottomrule
\end{tabularx}
\end{center}
\begin{center}\footnotesize
\begin{tabularx}{0.88\linewidth}{@{}l l X@{}}
\toprule
Field Name & Bits & Type / Fixed Value \\
\midrule
\texttt{src0} & \texttt{0:7} & \texttt{SourceEncoding} \\
\texttt{src1} & \texttt{8:15} & \texttt{SourceEncoding} \\
\texttt{src2} & \texttt{16:23} & \texttt{SourceEncoding} \\
\texttt{combine} & \texttt{24:25} & \texttt{exact=multi} \\
\hyperref[enum:widen_int_m]{\texttt{widen1}} & \texttt{26:29} & \texttt{widen\_int\_m (def none)} \\
\hyperref[enum:cmp_result_type_m]{\texttt{result\_type}} & \texttt{30:31} & \texttt{cmp\_result\_type\_m (def i1)} \\
\hyperref[enum:cmpf_m]{\texttt{cmpf}} & \texttt{32:34} & \texttt{cmpf\_m} \\
\textit{Reserved} & \texttt{35} & \texttt{zero} \\
\hyperref[enum:widen_int_m]{\texttt{widen0}} & \texttt{36:39} & \texttt{widen\_int\_m (def none)} \\
\texttt{dest\_register\_index} & \texttt{40:45} & \texttt{uint} \\
\hyperref[enum:dest_width_m]{\texttt{dest\_width}} & \texttt{46:47} & \texttt{dest\_width\_m} \\
\texttt{datatype} & \texttt{48:51} & \texttt{exact=s32} \\
\texttt{instruction\_class} & \texttt{52:56} & \texttt{exact=R3\_CMP} \\
\texttt{fau\_page\_index} & \texttt{57:58} & \texttt{uint} \\
\hyperref[enum:flow_control_m]{\texttt{flow}} & \texttt{59:62} & \texttt{flow\_control\_m} \\
\textit{Reserved} & \texttt{63} & \texttt{zero} \\
\bottomrule
\end{tabularx}
\end{center}
\textbf{Fields Explained:}
\begin{tightitemize}\footnotesize
  \item \texttt{src0}: Primary source operand (\texttt{r0..r63}, \texttt{u0..u63} uniform, \texttt{k*} FAU constant, or \texttt{t0..t2} temporary; \texttt{\textasciicircum{}} marks discard/last-use). \emph{Example:} {\scriptsize\ttfamily \mnem{FADD.f32} \dst{r0}, \src{r1}, \src{r2}} (\textbf{SPIR-V:} \texttt{OpFAdd (Operand 0)}; \textbf{GLSL:} \texttt{float c = a + b;})
  \item \texttt{src1}: Secondary source operand (register, uniform, or FAU constant; in shifts, supplies the shift count byte). \emph{Example:} {\scriptsize\ttfamily \mnem{LSHIFT.i32} \dst{r0}, \src{r1}, \src{r2}\imod{.b0}} (\textbf{SPIR-V:} \texttt{OpShiftLeftLogical (Shift)}; \textbf{GLSL:} \texttt{uint y = a \textless{}\textless{} s;})
  \item \texttt{src2}: Third source operand (FMA addend, \texttt{MUX}/\texttt{CSEL} selector, \texttt{IADDC} carry-in, or bitwise mask in fused shift-logic). \emph{Example:} {\scriptsize\ttfamily \mnem{FMA.f32} \dst{r0}, \src{r1}, \src{r2}, \src{r3}} (\textbf{SPIR-V:} \texttt{GLSL.std.450 Fma (Addend)}; \textbf{GLSL:} \texttt{float d = fma(a, b, c);})
  \item \texttt{combine}: Combines 16-bit halfword predicates (\texttt{h0}, \texttt{h1}, \texttt{and}, \texttt{lowbits}) before evaluating a branch. \emph{Example:} {\scriptsize\ttfamily \mnem{BRANCH}\imod{.and} \src{r0}, \imm{\#+8}} (\textbf{SPIR-V:} \texttt{OpBranchConditional (packed predicate)}; \textbf{GLSL:} \texttt{if (condA \&\& condB) \{ ... \}})
  \item \hyperref[enum:widen_int_m]{\texttt{widen1}}: Stage A widening extractor on \texttt{src1} (promotes narrow byte/half/word lane before Stage B). \emph{Example:} {\scriptsize\ttfamily \mnem{IMUL.u64} \dst{[r0:r1]}, \src{[r2:r3]}\imod{.w0}, \src{[r4:r5]}\imod{.w0}} (\textbf{SPIR-V:} \texttt{OpUConvert / OpUMulExtended}; \textbf{GLSL:} \texttt{uint64\_t p = uint64\_t(a) * uint64\_t(b);})
  \item \hyperref[enum:cmp_result_type_m]{\texttt{result\_type}}: Encodes boolean output representation (\texttt{i1} = 0/1, \texttt{f1} = 0.0/1.0, \texttt{m1} = 0/-1, \texttt{c} = carry). \emph{Example:} {\scriptsize\ttfamily \mnem{FCMP.f32}\imod{.lt.m1} \dst{r0}, \src{r1}, \src{r2}} (\textbf{SPIR-V:} \texttt{OpFOrdLessThan + OpSelect}; \textbf{GLSL:} \texttt{int mask = (a \textless{} b) ? -1 : 0;})
  \item \hyperref[enum:cmpf_m]{\texttt{cmpf}}: Comparison predicate (\texttt{eq}, \texttt{ne}, \texttt{lt}, \texttt{le}, \texttt{gt}, \texttt{ge}, \texttt{gtlt}, \texttt{total}). \emph{Example:} {\scriptsize\ttfamily \mnem{CSEL.f32}\imod{.lt} \dst{r0}, \src{r1}, \src{r2}, \src{r3}, \src{r4}} (\textbf{SPIR-V:} \texttt{OpFOrdLessThan + OpSelect}; \textbf{GLSL:} \texttt{float y = (a \textless{} b) ? c : d;})
  \item \hyperref[enum:widen_int_m]{\texttt{widen0}}: Stage A widening extractor on \texttt{src0} (promotes FP16 \texttt{h0}/\texttt{h1} to FP32, or sign/zero-extends \texttt{b0..b3}/\texttt{h0..h1}/\texttt{w0..w1} to 32/64-bit before \texttt{abs}/\texttt{neg}). \emph{Example:} {\scriptsize\ttfamily \mnem{IADD.s32} \dst{r0}, \src{r1}\imod{.h0}, \src{r2}} (\textbf{SPIR-V:} \texttt{OpSConvert / OpFConvert}; \textbf{GLSL:} \texttt{int c = int(int16\_t(a)) + b;})
  \item \texttt{dest\_register\_index}: 6-bit destination GPR index (\texttt{r0..r63}) receiving the Stage C writeback result (even-aligned pair \texttt{[rN:rN+1]} for 64-bit ops). \emph{Example:} {\scriptsize\ttfamily \mnem{FMA.f32} \dst{r0}, \src{r1}, \src{r2}, \src{r3}} (\textbf{SPIR-V:} \texttt{SSA Result ID}; \textbf{GLSL:} \texttt{float r0 = fma(r1, r2, r3);})
  \item \hyperref[enum:dest_width_m]{\texttt{dest\_width}}: Stage C destination lane write-mask (\texttt{none} = full 32-bit, \texttt{h0} = low 16-bit, \texttt{h1} = high 16-bit without bit-shifting). \emph{Example:} {\scriptsize\ttfamily \mnem{FADD.v2f16} \dst{r0}\imod{.h1}, \src{r1}, \src{r2}} (\textbf{SPIR-V:} \texttt{OpVectorInsertDynamic / WriteMask}; \textbf{GLSL:} \texttt{dst.y = a.y + b.y;})
  \item \texttt{datatype}: Selects the lane data type (\texttt{f32}, \texttt{f16}, \texttt{v2f16}, \texttt{s32}, \texttt{u32}, \texttt{v2s16}, \texttt{v2u16}, \texttt{v4s8}, \texttt{v4u8}, \texttt{s64}, \texttt{u64}, etc.). \emph{Example:} {\scriptsize\ttfamily \mnem{IADD.v2s16} \dst{r0}, \src{r1}, \src{r2}} (\textbf{SPIR-V:} \texttt{OpTypeInt / OpTypeFloat}; \textbf{GLSL:} \texttt{i16vec2 c = a + b;})
  \item \texttt{instruction\_class}: 4-bit major instruction class selector (\texttt{bits 48:51}) partitioning the 64-bit Valhall opcode space (e.g., \texttt{FMA}, \texttt{CVT}, \texttt{SFU}, \texttt{LS}, \texttt{TEX}, \texttt{BRANCH}). \emph{Example:} {\scriptsize\ttfamily \mnem{FADD.f32} \dst{r0}, \src{r1}, \src{r2}} (\textbf{SPIR-V:} \texttt{N/A (encoding class)}; \textbf{GLSL:} \texttt{// Hardware execution pipe / decode class})
  \item \texttt{fau\_page\_index}: Selects the FAU (Fast Access Uniform) page (\texttt{0..3}) when a source operand references hardware small constants, thread IDs, or lane constants. \emph{Example:} {\scriptsize\ttfamily \mnem{FADD.f32} \dst{r0}, \src{r1}, \src{k1f}} (\textbf{SPIR-V:} \texttt{OpConstant / BuiltIn}; \textbf{GLSL:} \texttt{float y = x + 1.0;})
  \item \hyperref[enum:flow_control_m]{\texttt{flow}}: 4-bit embedded flow-control \& scoreboard wait modifier (\texttt{bits 59:62}: \texttt{none}, \texttt{wait0..wait0126}, \texttt{wait}, \texttt{reconverge}, \texttt{discard}, \texttt{end}). \emph{Example:} {\scriptsize\ttfamily \mnem{FADD.f32}\imod{.wait0.reconverge} \dst{r0}, \src{r1}, \src{r2}} (\textbf{SPIR-V:} \texttt{OpControlBarrier / OpReturn}; \textbf{GLSL:} \texttt{// Scoreboard wait + warp reconvergence})
\end{tightitemize}
\clearpage

\subsection*{ICMP\_OR}
\addcontentsline{toc}{subsection}{ICMP\_OR}
For each SIMD element $i \in \{0, \dots, N-1\}$ ($N=1$ for \texttt{s32}/\texttt{u32}, $N=2$ for \texttt{v2s16}/\texttt{v2u16}, $N=4$ for \texttt{v4s8}/\texttt{v4u8} on v9--v10), compares Stage A widened/swizzled \texttt{src0[i]} and \texttt{src1[i]} using predicate \texttt{cmpf}, combines the boolean with \texttt{(src2[i] != 0)} via logical OR, and formats the result into \texttt{dest[i]} according to \texttt{result\_type} (\texttt{i1} $\to$ \texttt{1}, \texttt{m1} $\to$ all-ones, \texttt{f1} $\to$ \texttt{1.0}, false $\to$ \texttt{0}).
\par\textbf{Notes:}
\begin{tightitemize}
  \item \texttt{v4s8}/\texttt{v4u8} only arch \textless{}= v10.
  \item \textbf{Source references:} \texttt{kraid/ops.rs:741-790}, \texttt{encode\_v9.rs:1012-1020}, \texttt{hw\_tests.rs:1318-1388}.
\end{tightitemize}
\codebox{\mnem{ICMP\_OR}\imod{\{.flow\}}\imod{\{.result\_type\}}\imod{\{.cmpf\}} \dst{dest}\imod{\{.dest\_width\}}, \src{src0}\imod{\{.widen0\}}, \src{src1}\imod{\{.widen1\}}, \src{src2}\\\mnem{ICMP\_OR}\imod{\{.flow\}}\imod{\{.result\_type\}}\imod{\{.cmpf\}} \dst{dest}\imod{\{.dest\_width\}}, \src{src0}\imod{\{.swz0\}}, \src{src1}\imod{\{.swz1\}}, \src{src2}\\\mnem{ICMP\_OR}\imod{\{.flow\}}\imod{\{.result\_type\}}\imod{\{.cmpf\}} \dst{dest}\imod{\{.dest\_width\}}, \src{src0}\imod{\{.lanes0\}}, \src{src1}\imod{\{.lanes1\}}, \src{src2}}
\textbf{Operands:}
\begin{tightitemize}\footnotesize
  \item \texttt{\dst{dest}} (\texttt{dest\_register\_index}, \texttt{dest\_width}): Destination GPR receiving \texttt{cmpf(src0, src1) \textbar{}\textbar{} (src2 != 0)} formatted per \texttt{result\_type}.
  \item \texttt{\src{src0}}: Left-hand integer comparison operand.
  \item \texttt{\src{src1}}: Right-hand integer comparison operand.
  \item \texttt{\src{src2}}: Existing boolean/mask predicate combined with logical OR (\texttt{k0} encodes standalone \texttt{ICMP}).
\end{tightitemize}
\par\textbf{Closest SPIR-V Op:} \texttt{OpLogicalOr + OpI*/OpS*/OpU*}
\par\textbf{Pseudo-GLSL (Operational Semantics):}
\codebox{for (int i = 0; i \textless{} N; ++i) \{\\    bool c = cmpf(widen0(src0)[i], widen1(src1)[i]) \textbar{}\textbar{} (src2[i] != 0);\\    dest[i] = formatResult(result\_type, c);\\\}}
\textbf{Execution Context:}
\begin{tightitemize}\footnotesize
  \item \textbf{Compute stage preloads (\hyperref[sec:csf-spd]{\texttt{SPD\_k}}, \texttt{r16..r23}):} Hardware pre-fills \texttt{r55..r57} with \texttt{WorkGroupID.xyz + JOB\_OFFSET} (from \hyperref[sec:csf-regs]{\texttt{COMPUTESR::JOB\_OFFSET\_X/Y/Z}} \texttt{r34..r36}) and \texttt{r60..r62} with \texttt{LocalInvocationID.xyz} or \texttt{GlobalInvocationID.xyz} (from \texttt{WG\_SIZE} \texttt{r33} \& \texttt{JOB\_SIZE} \texttt{r37..r39}, selected by \texttt{ShaderProgram::preload}) so compute shaders can derive thread/element indices without memory loads.
  \item \textbf{Vertex stage preloads (\hyperref[sec:csf-spd]{\texttt{SPD\_0}/\texttt{SPD\_1}}, \texttt{r16..r19}):} The hardware Index Fetcher reads \hyperref[sec:csf-regs]{\texttt{IDVSSR::INDEX\_BUFFER}} (\texttt{r54:r55}) and pre-fills \texttt{r59, r60} with \texttt{VertexIndex} (\texttt{+ VERTEX\_OFFSET + GLOBAL\_ATTR\_OFF}) and \texttt{r61} with \texttt{InstanceIndex} (\texttt{+ INSTANCE\_OFFSET}) so Position and Varying shaders can pull vertex attributes directly.
  \item \textbf{Inline FAU \& Special Operands (\hyperref[sec:csf-tsd-fau]{\texttt{FAU\_k}}, \texttt{r8..r15}):} Any 8-bit source with \texttt{mode=2} reads zero-latency uniform \texttt{u0..u31} from \texttt{FAU\_k} (\texttt{slots64 = bits[63:56]}); \texttt{mode=3} reads \texttt{SmallConstant} (Table~1.2) or \texttt{FauSpecial} (\texttt{pc}, \texttt{lane\_id}, \texttt{lane\_width}, \texttt{tls\_ptr}, \texttt{wls\_ptr}, \texttt{blend\_descriptor\_0..7}).
\end{tightitemize}
\par\textbf{Worked Examples:}
\par\textbf{Example 1 (Both false):}
\codebox{\mnem{ICMP\_OR.u32}\imod{.ne.m1} \dst{r0}, \src{r1}, \src{r2}, \src{r3}}
With \texttt{r1 = 4, r2 = 4, r3 = 0}: \texttt{(4 != 4) \textbar{}\textbar{} (0 != 0)} is False -\textgreater{} \texttt{r0 = 0}.
\par\textbf{Example 2 (64-bit inequality (`!=`) high-word step):}
\codebox{\mnem{ICMP\_OR.u32}\imod{.ne.i1} \dst{r0}, \src{r1}, \src{r2}, \src{r3}}
With \texttt{r1 = 0, r2 = 0, r3 = 1} (low words differed): \texttt{(0 != 0) \textbar{}\textbar{} (1 != 0)} is True -\textgreater{} \texttt{r0 = 1}.
\par\textbf{Mined SPIR-V $\to$ Valhall ISA Example (Short-circuit OR of two integer equality tests --- Vertex shader, Mali-G710 (v10)):}
\codebox{\textbf{; SPIR-V (input)}\\\%eq0 = OpIEqual \%bool \%bit0 \%int\_0\\\%any = OpLogicalOr \%bool \%eq0 \%prev\_cond\\[3pt]\textbf{; Valhall ISA (malioc output)}\\ICMP\_OR.u32.m1.eq r0, r0\textasciicircum{}, k0, r2\\BRANCH r1, 0x78}
\texttt{ICMP\_OR.u32.m1.eq} fuses \texttt{(r0 == 0) \textbar{}\textbar{} (r2 != 0)} into a single instruction, avoiding a separate \texttt{OR.i32}.
\begin{center}\footnotesize
\begin{tabularx}{\linewidth}{@{}lllX@{}}
\toprule
Variant & Arch & Class & Opcode Fields \\
\midrule
\texttt{s32} & v9--v14 & \texttt{R3\_CMP} & \texttt{--} \\
\texttt{u32} & v9--v14 & \texttt{R3\_CMP} & \texttt{--} \\
\texttt{v2s16} & v9--v14 & \texttt{R3\_CMP} & \texttt{--} \\
\texttt{v2u16} & v9--v14 & \texttt{R3\_CMP} & \texttt{--} \\
\texttt{v4s8} & v9--v14 & \texttt{R3\_CMP} & \texttt{--} \\
\texttt{v4u8} & v9--v14 & \texttt{R3\_CMP} & \texttt{--} \\
\bottomrule
\end{tabularx}
\end{center}
\begin{center}\footnotesize
\begin{tabularx}{0.88\linewidth}{@{}l l X@{}}
\toprule
Field Name & Bits & Type / Fixed Value \\
\midrule
\texttt{src0} & \texttt{0:7} & \texttt{SourceEncoding} \\
\texttt{src1} & \texttt{8:15} & \texttt{SourceEncoding} \\
\texttt{src2} & \texttt{16:23} & \texttt{SourceEncoding} \\
\texttt{combine} & \texttt{24:25} & \texttt{exact=or} \\
\hyperref[enum:widen_int_m]{\texttt{widen1}} & \texttt{26:29} & \texttt{widen\_int\_m (def none)} \\
\hyperref[enum:cmp_result_type_m]{\texttt{result\_type}} & \texttt{30:31} & \texttt{cmp\_result\_type\_m (def i1)} \\
\hyperref[enum:cmpf_m]{\texttt{cmpf}} & \texttt{32:34} & \texttt{cmpf\_m} \\
\textit{Reserved} & \texttt{35} & \texttt{zero} \\
\hyperref[enum:widen_int_m]{\texttt{widen0}} & \texttt{36:39} & \texttt{widen\_int\_m (def none)} \\
\texttt{dest\_register\_index} & \texttt{40:45} & \texttt{uint} \\
\hyperref[enum:dest_width_m]{\texttt{dest\_width}} & \texttt{46:47} & \texttt{dest\_width\_m} \\
\texttt{datatype} & \texttt{48:51} & \texttt{exact=s32} \\
\texttt{instruction\_class} & \texttt{52:56} & \texttt{exact=R3\_CMP} \\
\texttt{fau\_page\_index} & \texttt{57:58} & \texttt{uint} \\
\hyperref[enum:flow_control_m]{\texttt{flow}} & \texttt{59:62} & \texttt{flow\_control\_m} \\
\textit{Reserved} & \texttt{63} & \texttt{zero} \\
\bottomrule
\end{tabularx}
\end{center}
\textbf{Fields Explained:}
\begin{tightitemize}\footnotesize
  \item \texttt{src0}: Primary source operand (\texttt{r0..r63}, \texttt{u0..u63} uniform, \texttt{k*} FAU constant, or \texttt{t0..t2} temporary; \texttt{\textasciicircum{}} marks discard/last-use). \emph{Example:} {\scriptsize\ttfamily \mnem{FADD.f32} \dst{r0}, \src{r1}, \src{r2}} (\textbf{SPIR-V:} \texttt{OpFAdd (Operand 0)}; \textbf{GLSL:} \texttt{float c = a + b;})
  \item \texttt{src1}: Secondary source operand (register, uniform, or FAU constant; in shifts, supplies the shift count byte). \emph{Example:} {\scriptsize\ttfamily \mnem{LSHIFT.i32} \dst{r0}, \src{r1}, \src{r2}\imod{.b0}} (\textbf{SPIR-V:} \texttt{OpShiftLeftLogical (Shift)}; \textbf{GLSL:} \texttt{uint y = a \textless{}\textless{} s;})
  \item \texttt{src2}: Third source operand (FMA addend, \texttt{MUX}/\texttt{CSEL} selector, \texttt{IADDC} carry-in, or bitwise mask in fused shift-logic). \emph{Example:} {\scriptsize\ttfamily \mnem{FMA.f32} \dst{r0}, \src{r1}, \src{r2}, \src{r3}} (\textbf{SPIR-V:} \texttt{GLSL.std.450 Fma (Addend)}; \textbf{GLSL:} \texttt{float d = fma(a, b, c);})
  \item \texttt{combine}: Combines 16-bit halfword predicates (\texttt{h0}, \texttt{h1}, \texttt{and}, \texttt{lowbits}) before evaluating a branch. \emph{Example:} {\scriptsize\ttfamily \mnem{BRANCH}\imod{.and} \src{r0}, \imm{\#+8}} (\textbf{SPIR-V:} \texttt{OpBranchConditional (packed predicate)}; \textbf{GLSL:} \texttt{if (condA \&\& condB) \{ ... \}})
  \item \hyperref[enum:widen_int_m]{\texttt{widen1}}: Stage A widening extractor on \texttt{src1} (promotes narrow byte/half/word lane before Stage B). \emph{Example:} {\scriptsize\ttfamily \mnem{IMUL.u64} \dst{[r0:r1]}, \src{[r2:r3]}\imod{.w0}, \src{[r4:r5]}\imod{.w0}} (\textbf{SPIR-V:} \texttt{OpUConvert / OpUMulExtended}; \textbf{GLSL:} \texttt{uint64\_t p = uint64\_t(a) * uint64\_t(b);})
  \item \hyperref[enum:cmp_result_type_m]{\texttt{result\_type}}: Encodes boolean output representation (\texttt{i1} = 0/1, \texttt{f1} = 0.0/1.0, \texttt{m1} = 0/-1, \texttt{c} = carry). \emph{Example:} {\scriptsize\ttfamily \mnem{FCMP.f32}\imod{.lt.m1} \dst{r0}, \src{r1}, \src{r2}} (\textbf{SPIR-V:} \texttt{OpFOrdLessThan + OpSelect}; \textbf{GLSL:} \texttt{int mask = (a \textless{} b) ? -1 : 0;})
  \item \hyperref[enum:cmpf_m]{\texttt{cmpf}}: Comparison predicate (\texttt{eq}, \texttt{ne}, \texttt{lt}, \texttt{le}, \texttt{gt}, \texttt{ge}, \texttt{gtlt}, \texttt{total}). \emph{Example:} {\scriptsize\ttfamily \mnem{CSEL.f32}\imod{.lt} \dst{r0}, \src{r1}, \src{r2}, \src{r3}, \src{r4}} (\textbf{SPIR-V:} \texttt{OpFOrdLessThan + OpSelect}; \textbf{GLSL:} \texttt{float y = (a \textless{} b) ? c : d;})
  \item \hyperref[enum:widen_int_m]{\texttt{widen0}}: Stage A widening extractor on \texttt{src0} (promotes FP16 \texttt{h0}/\texttt{h1} to FP32, or sign/zero-extends \texttt{b0..b3}/\texttt{h0..h1}/\texttt{w0..w1} to 32/64-bit before \texttt{abs}/\texttt{neg}). \emph{Example:} {\scriptsize\ttfamily \mnem{IADD.s32} \dst{r0}, \src{r1}\imod{.h0}, \src{r2}} (\textbf{SPIR-V:} \texttt{OpSConvert / OpFConvert}; \textbf{GLSL:} \texttt{int c = int(int16\_t(a)) + b;})
  \item \texttt{dest\_register\_index}: 6-bit destination GPR index (\texttt{r0..r63}) receiving the Stage C writeback result (even-aligned pair \texttt{[rN:rN+1]} for 64-bit ops). \emph{Example:} {\scriptsize\ttfamily \mnem{FMA.f32} \dst{r0}, \src{r1}, \src{r2}, \src{r3}} (\textbf{SPIR-V:} \texttt{SSA Result ID}; \textbf{GLSL:} \texttt{float r0 = fma(r1, r2, r3);})
  \item \hyperref[enum:dest_width_m]{\texttt{dest\_width}}: Stage C destination lane write-mask (\texttt{none} = full 32-bit, \texttt{h0} = low 16-bit, \texttt{h1} = high 16-bit without bit-shifting). \emph{Example:} {\scriptsize\ttfamily \mnem{FADD.v2f16} \dst{r0}\imod{.h1}, \src{r1}, \src{r2}} (\textbf{SPIR-V:} \texttt{OpVectorInsertDynamic / WriteMask}; \textbf{GLSL:} \texttt{dst.y = a.y + b.y;})
  \item \texttt{datatype}: Selects the lane data type (\texttt{f32}, \texttt{f16}, \texttt{v2f16}, \texttt{s32}, \texttt{u32}, \texttt{v2s16}, \texttt{v2u16}, \texttt{v4s8}, \texttt{v4u8}, \texttt{s64}, \texttt{u64}, etc.). \emph{Example:} {\scriptsize\ttfamily \mnem{IADD.v2s16} \dst{r0}, \src{r1}, \src{r2}} (\textbf{SPIR-V:} \texttt{OpTypeInt / OpTypeFloat}; \textbf{GLSL:} \texttt{i16vec2 c = a + b;})
  \item \texttt{instruction\_class}: 4-bit major instruction class selector (\texttt{bits 48:51}) partitioning the 64-bit Valhall opcode space (e.g., \texttt{FMA}, \texttt{CVT}, \texttt{SFU}, \texttt{LS}, \texttt{TEX}, \texttt{BRANCH}). \emph{Example:} {\scriptsize\ttfamily \mnem{FADD.f32} \dst{r0}, \src{r1}, \src{r2}} (\textbf{SPIR-V:} \texttt{N/A (encoding class)}; \textbf{GLSL:} \texttt{// Hardware execution pipe / decode class})
  \item \texttt{fau\_page\_index}: Selects the FAU (Fast Access Uniform) page (\texttt{0..3}) when a source operand references hardware small constants, thread IDs, or lane constants. \emph{Example:} {\scriptsize\ttfamily \mnem{FADD.f32} \dst{r0}, \src{r1}, \src{k1f}} (\textbf{SPIR-V:} \texttt{OpConstant / BuiltIn}; \textbf{GLSL:} \texttt{float y = x + 1.0;})
  \item \hyperref[enum:flow_control_m]{\texttt{flow}}: 4-bit embedded flow-control \& scoreboard wait modifier (\texttt{bits 59:62}: \texttt{none}, \texttt{wait0..wait0126}, \texttt{wait}, \texttt{reconverge}, \texttt{discard}, \texttt{end}). \emph{Example:} {\scriptsize\ttfamily \mnem{FADD.f32}\imod{.wait0.reconverge} \dst{r0}, \src{r1}, \src{r2}} (\textbf{SPIR-V:} \texttt{OpControlBarrier / OpReturn}; \textbf{GLSL:} \texttt{// Scoreboard wait + warp reconvergence})
\end{tightitemize}
\clearpage

\subsection*{MUX}
\addcontentsline{toc}{subsection}{MUX}
Selects between \texttt{src0[i]} and \texttt{src1[i]} using selector \texttt{src2[i]} independently for each SIMD element $i \in \{0, \dots, N-1\}$ ($N=1$ \texttt{i32}, $N=2$ \texttt{v2i16}, $N=4$ \texttt{v4i8}): \texttt{.neg} selects \texttt{src0[i]} if the sign bit of \texttt{src2[i]} is 1; \texttt{.int\_zero} (default) if \texttt{src2[i] == 0}; \texttt{.fp\_zero} (\texttt{i32} only) if \texttt{f32(src2) == $\pm$0.0}; \texttt{.bit} blends bitwise \texttt{(src0[i] \& src2[i]) \textbar{} (src1[i] \& \textasciitilde{}src2[i])}.
\par\textbf{Notes:}
\begin{tightitemize}
  \item \texttt{fp\_zero} is \texttt{i32}-only (\texttt{ops.rs:3639}, \texttt{hw\_tests.rs:2121}) and treats both \texttt{+0.0} and \texttt{-0.0} as true.
  \item Exact match verified in \texttt{hw\_tests.rs:2113-2135}.
  \item \textbf{Source references:} \texttt{kraid/ops.rs:3573-3647}, \texttt{encode\_v9.rs:2858-2885}, \texttt{opt\_fuse.rs:150-180}, \texttt{hw\_tests.rs:2113-2135}.
\end{tightitemize}
\codebox{\mnem{MUX}\imod{\{.flow\}}\imod{\{.muxf\}} \dst{dest}\imod{\{.dest\_width\}}, \src{src0}, \src{src1}, \src{src2}\\\mnem{MUX}\imod{\{.flow\}}\imod{\{.muxf\}} \dst{dest}\imod{\{.dest\_width\}}, \src{src0}\imod{\{.swz0\}}, \src{src1}\imod{\{.swz1\}}, \src{src2}\imod{\{.swz2\}}}
\textbf{Operands:}
\begin{tightitemize}\footnotesize
  \item \texttt{\dst{dest}} (\texttt{dest\_register\_index}, \texttt{dest\_width}): Destination GPR receiving lane/bit-selected output from \texttt{src0} and \texttt{src1}.
  \item \texttt{\src{src0}}: Value selected when the \texttt{mux} test succeeds (or where mask bits are \texttt{1} in \texttt{.bit} mode).
  \item \texttt{\src{src1}}: Value selected when the \texttt{mux} test fails (or where mask bits are \texttt{0} in \texttt{.bit} mode).
  \item \texttt{\src{src2}}: Selector value or bitwise mask evaluated according to \texttt{mux} mode (\texttt{.int\_zero}, \texttt{.neg}, \texttt{.fp\_zero}, \texttt{.bit}).
\end{tightitemize}
\par\textbf{Closest SPIR-V Op:} \texttt{OpSelect / OpenCL `bitselect`}
\par\textbf{Pseudo-GLSL (Operational Semantics):}
\codebox{for (int i = 0; i \textless{} N; ++i) \{\\    if (mux == bit)          dest[i] = (src0[i] \& src2[i]) \textbar{} (src1[i] \& \textasciitilde{}src2[i]);\\    else if (mux == neg)     dest[i] = (int(src2[i]) \textless{} 0) ? src0[i] : src1[i];\\    else if (mux == fp\_zero) dest[i] = (asFloat(src2[i]) == 0.0) ? src0[i] : src1[i];\\    else                     dest[i] = (src2[i] == 0u) ? src0[i] : src1[i]; // .int\_zero\\\}}
\textbf{Execution Context:}
\begin{tightitemize}\footnotesize
  \item \textbf{Compute stage preloads (\hyperref[sec:csf-spd]{\texttt{SPD\_k}}, \texttt{r16..r23}):} Hardware pre-fills \texttt{r55..r57} with \texttt{WorkGroupID.xyz + JOB\_OFFSET} (from \hyperref[sec:csf-regs]{\texttt{COMPUTESR::JOB\_OFFSET\_X/Y/Z}} \texttt{r34..r36}) and \texttt{r60..r62} with \texttt{LocalInvocationID.xyz} or \texttt{GlobalInvocationID.xyz} (from \texttt{WG\_SIZE} \texttt{r33} \& \texttt{JOB\_SIZE} \texttt{r37..r39}, selected by \texttt{ShaderProgram::preload}) so compute shaders can derive thread/element indices without memory loads.
  \item \textbf{Vertex stage preloads (\hyperref[sec:csf-spd]{\texttt{SPD\_0}/\texttt{SPD\_1}}, \texttt{r16..r19}):} The hardware Index Fetcher reads \hyperref[sec:csf-regs]{\texttt{IDVSSR::INDEX\_BUFFER}} (\texttt{r54:r55}) and pre-fills \texttt{r59, r60} with \texttt{VertexIndex} (\texttt{+ VERTEX\_OFFSET + GLOBAL\_ATTR\_OFF}) and \texttt{r61} with \texttt{InstanceIndex} (\texttt{+ INSTANCE\_OFFSET}) so Position and Varying shaders can pull vertex attributes directly.
  \item \textbf{Inline FAU \& Special Operands (\hyperref[sec:csf-tsd-fau]{\texttt{FAU\_k}}, \texttt{r8..r15}):} Any 8-bit source with \texttt{mode=2} reads zero-latency uniform \texttt{u0..u31} from \texttt{FAU\_k} (\texttt{slots64 = bits[63:56]}); \texttt{mode=3} reads \texttt{SmallConstant} (Table~1.2) or \texttt{FauSpecial} (\texttt{pc}, \texttt{lane\_id}, \texttt{lane\_width}, \texttt{tls\_ptr}, \texttt{wls\_ptr}, \texttt{blend\_descriptor\_0..7}).
\end{tightitemize}
\par\textbf{Worked Examples:}
\par\textbf{Example 1 (Per-bit blend (`.bit`)):}
\codebox{\mnem{MUX.i32}\imod{.bit} \dst{r0}, \src{r1}, \src{r2}, \src{r3}}
With \texttt{r1 = 0x12345678, r2 = 0xA0000000, r3 = 0x1FFFFFFF}: takes bits \texttt{[28:0]} from \texttt{r1} and \texttt{[31:29]} from \texttt{r2} -\textgreater{} \texttt{r0 = 0xB2345678}.
\par\textbf{Example 2 (Default `.int\_zero` selector (encoding 1)):}
\codebox{\mnem{MUX.i32}\imod{.int\_zero} \dst{r0}, \src{r1}, \src{r2}, \src{r3}}
With \texttt{r1 = 0x1111, r2 = 0x2222, r3 = 0}: \texttt{r3 == 0} selects \texttt{src0} -\textgreater{} \texttt{r0 = 0x1111} (for \texttt{r3 != 0}, selects \texttt{src1 = 0x2222}).
\par\textbf{Mined SPIR-V $\to$ Valhall ISA Example (Bitwise merge of cubemap face index into S coordinate --- Fragment shader, Mali-G710 (v10)):}
\codebox{\textbf{; SPIR-V (input)}\\\%color = OpImageSampleImplicitLod \%v4float \%sampled\_cube \%dir\_xyz None\\[3pt]\textbf{; Valhall ISA (malioc output)}\\CUBEFACE2 r6, r5, r4, r3\\FMA.f32.clamp\_0\_1 r7, r5, r7\textasciicircum{}, k0fr2f.h1\\MUX.i32.bit r8, r7\textasciicircum{}, r6, u1.w1}
\texttt{MUX.i32.bit} computes \texttt{(r7 \& u1.w1) \textbar{} (r6 \& \textasciitilde{}u1.w1)} bit-by-bit, splicing the 3-bit cube face ID from \texttt{r6[31:29]} into the projected \texttt{s} coordinate \texttt{r7} in one cycle.
\begin{center}\footnotesize
\begin{tabularx}{\linewidth}{@{}lllX@{}}
\toprule
Variant & Arch & Class & Opcode Fields \\
\midrule
\texttt{(none)} & v9--v14 & \texttt{R3} & \texttt{opcode1=MUX\_I32} \\
\texttt{(none)} & v9--v14 & \texttt{R3} & \texttt{opcode1=MUX\_V2I16} \\
\texttt{(none)} & v9--v14 & \texttt{R3} & \texttt{opcode1=MUX\_V4I8} \\
\bottomrule
\end{tabularx}
\end{center}
\begin{center}\footnotesize
\begin{tabularx}{0.88\linewidth}{@{}l l X@{}}
\toprule
Field Name & Bits & Type / Fixed Value \\
\midrule
\texttt{src0} & \texttt{0:7} & \texttt{SourceEncoding} \\
\texttt{src1} & \texttt{8:15} & \texttt{SourceEncoding} \\
\texttt{src2} & \texttt{16:23} & \texttt{SourceEncoding} \\
\textit{Reserved} & \texttt{24:31} & \texttt{zero} \\
\hyperref[enum:mux_cmpf_m]{\texttt{muxf}} & \texttt{32:33} & \texttt{mux\_cmpf\_m (def int\_zero)} \\
\textit{Reserved} & \texttt{34:39} & \texttt{zero} \\
\texttt{dest\_register\_index} & \texttt{40:45} & \texttt{uint} \\
\hyperref[enum:dest_width_m]{\texttt{dest\_width}} & \texttt{46:47} & \texttt{dest\_width\_m} \\
\texttt{opcode1} & \texttt{48:51} & \texttt{exact=MUX\_I32} \\
\texttt{instruction\_class} & \texttt{52:56} & \texttt{exact=R3} \\
\texttt{fau\_page\_index} & \texttt{57:58} & \texttt{uint} \\
\hyperref[enum:flow_control_m]{\texttt{flow}} & \texttt{59:62} & \texttt{flow\_control\_m} \\
\textit{Reserved} & \texttt{63} & \texttt{zero} \\
\bottomrule
\end{tabularx}
\end{center}
\textbf{Fields Explained:}
\begin{tightitemize}\footnotesize
  \item \texttt{src0}: Primary source operand (\texttt{r0..r63}, \texttt{u0..u63} uniform, \texttt{k*} FAU constant, or \texttt{t0..t2} temporary; \texttt{\textasciicircum{}} marks discard/last-use). \emph{Example:} {\scriptsize\ttfamily \mnem{FADD.f32} \dst{r0}, \src{r1}, \src{r2}} (\textbf{SPIR-V:} \texttt{OpFAdd (Operand 0)}; \textbf{GLSL:} \texttt{float c = a + b;})
  \item \texttt{src1}: Secondary source operand (register, uniform, or FAU constant; in shifts, supplies the shift count byte). \emph{Example:} {\scriptsize\ttfamily \mnem{LSHIFT.i32} \dst{r0}, \src{r1}, \src{r2}\imod{.b0}} (\textbf{SPIR-V:} \texttt{OpShiftLeftLogical (Shift)}; \textbf{GLSL:} \texttt{uint y = a \textless{}\textless{} s;})
  \item \texttt{src2}: Third source operand (FMA addend, \texttt{MUX}/\texttt{CSEL} selector, \texttt{IADDC} carry-in, or bitwise mask in fused shift-logic). \emph{Example:} {\scriptsize\ttfamily \mnem{FMA.f32} \dst{r0}, \src{r1}, \src{r2}, \src{r3}} (\textbf{SPIR-V:} \texttt{GLSL.std.450 Fma (Addend)}; \textbf{GLSL:} \texttt{float d = fma(a, b, c);})
  \item \hyperref[enum:mux_cmpf_m]{\texttt{muxf}}: Selects the \texttt{MUX} comparison predicate (\texttt{neg}, \texttt{int\_zero}, \texttt{fp\_zero}, \texttt{bit}). \emph{Example:} {\scriptsize\ttfamily \mnem{MUX.i32}\imod{.bit} \dst{r0}, \src{r1}, \src{r2}, \src{r3}} (\textbf{SPIR-V:} \texttt{OpBitwiseOr / OpBitwiseAnd}; \textbf{GLSL:} \texttt{uint z = (a \& mask) \textbar{} (b \& \textasciitilde{}mask);})
  \item \texttt{dest\_register\_index}: 6-bit destination GPR index (\texttt{r0..r63}) receiving the Stage C writeback result (even-aligned pair \texttt{[rN:rN+1]} for 64-bit ops). \emph{Example:} {\scriptsize\ttfamily \mnem{FMA.f32} \dst{r0}, \src{r1}, \src{r2}, \src{r3}} (\textbf{SPIR-V:} \texttt{SSA Result ID}; \textbf{GLSL:} \texttt{float r0 = fma(r1, r2, r3);})
  \item \hyperref[enum:dest_width_m]{\texttt{dest\_width}}: Stage C destination lane write-mask (\texttt{none} = full 32-bit, \texttt{h0} = low 16-bit, \texttt{h1} = high 16-bit without bit-shifting). \emph{Example:} {\scriptsize\ttfamily \mnem{FADD.v2f16} \dst{r0}\imod{.h1}, \src{r1}, \src{r2}} (\textbf{SPIR-V:} \texttt{OpVectorInsertDynamic / WriteMask}; \textbf{GLSL:} \texttt{dst.y = a.y + b.y;})
  \item \texttt{opcode1}: Primary sub-opcode within the instruction class uniquely identifying the operation family. \emph{Example:} {\scriptsize\ttfamily \mnem{IADD.s32} \dst{r0}, \src{r1}, \src{r2}} (\textbf{SPIR-V:} \texttt{OpIAdd (operation selector)}; \textbf{GLSL:} \texttt{// Primary opcode bits})
  \item \texttt{instruction\_class}: 4-bit major instruction class selector (\texttt{bits 48:51}) partitioning the 64-bit Valhall opcode space (e.g., \texttt{FMA}, \texttt{CVT}, \texttt{SFU}, \texttt{LS}, \texttt{TEX}, \texttt{BRANCH}). \emph{Example:} {\scriptsize\ttfamily \mnem{FADD.f32} \dst{r0}, \src{r1}, \src{r2}} (\textbf{SPIR-V:} \texttt{N/A (encoding class)}; \textbf{GLSL:} \texttt{// Hardware execution pipe / decode class})
  \item \texttt{fau\_page\_index}: Selects the FAU (Fast Access Uniform) page (\texttt{0..3}) when a source operand references hardware small constants, thread IDs, or lane constants. \emph{Example:} {\scriptsize\ttfamily \mnem{FADD.f32} \dst{r0}, \src{r1}, \src{k1f}} (\textbf{SPIR-V:} \texttt{OpConstant / BuiltIn}; \textbf{GLSL:} \texttt{float y = x + 1.0;})
  \item \hyperref[enum:flow_control_m]{\texttt{flow}}: 4-bit embedded flow-control \& scoreboard wait modifier (\texttt{bits 59:62}: \texttt{none}, \texttt{wait0..wait0126}, \texttt{wait}, \texttt{reconverge}, \texttt{discard}, \texttt{end}). \emph{Example:} {\scriptsize\ttfamily \mnem{FADD.f32}\imod{.wait0.reconverge} \dst{r0}, \src{r1}, \src{r2}} (\textbf{SPIR-V:} \texttt{OpControlBarrier / OpReturn}; \textbf{GLSL:} \texttt{// Scoreboard wait + warp reconvergence})
\end{tightitemize}
\clearpage

\section{Conversions (25 mnemonics)}
\subsection*{F16\_TO\_F32}
\addcontentsline{toc}{subsection}{F16\_TO\_F32}
Widens the FP16 halfword selected by \texttt{lane0} (\texttt{.h0} or \texttt{.h1}) of \texttt{src0} losslessly to IEEE-754 single-precision FP32 and writes the result to \texttt{dest} (exact conversion with no rounding or clamping).
\par\textbf{Notes:}
\begin{tightitemize}
  \item Exact IEEE-754 f16-\textgreater{}f32 widening (\texttt{hw\_tests.rs:1453-1459}).
  \item \textbf{Source references:} \texttt{kraid/ops.rs:1064-1088}, \texttt{encode\_v9.rs:1150-1166}, \texttt{hw\_tests.rs:1453-1459}.
\end{tightitemize}
\codebox{\mnem{F16\_TO\_F32}\imod{\{.flow\}} \dst{dest}\imod{\{.dest\_width\}}, \src{src0}\imod{\{.lane0\}}\imod{\{.abs0\}}\imod{\{.neg0\}}}
\textbf{Operands:}
\begin{tightitemize}\footnotesize
  \item \texttt{\dst{dest}} (\texttt{dest\_register\_index}, \texttt{dest\_width}): 32-bit destination GPR receiving the converted value.
  \item \texttt{\src{src0}}: FP16 input lane (\texttt{lane0}) promoted losslessly to FP32.
\end{tightitemize}
\par\textbf{Closest SPIR-V Op:} \texttt{OpFConvert (f16 -\textgreater{} f32) / UnpackHalf2x16}
\par\textbf{Pseudo-GLSL (Operational Semantics):}
\codebox{dest = unpackHalf2x16(src0)[lane0]; // exact f16 -\textgreater{} f32 widening}
\textbf{Execution Context:}
\begin{tightitemize}\footnotesize
  \item \textbf{Compute stage preloads (\hyperref[sec:csf-spd]{\texttt{SPD\_k}}, \texttt{r16..r23}):} Hardware pre-fills \texttt{r55..r57} with \texttt{WorkGroupID.xyz + JOB\_OFFSET} (from \hyperref[sec:csf-regs]{\texttt{COMPUTESR::JOB\_OFFSET\_X/Y/Z}} \texttt{r34..r36}) and \texttt{r60..r62} with \texttt{LocalInvocationID.xyz} or \texttt{GlobalInvocationID.xyz} (from \texttt{WG\_SIZE} \texttt{r33} \& \texttt{JOB\_SIZE} \texttt{r37..r39}, selected by \texttt{ShaderProgram::preload}) so compute shaders can derive thread/element indices without memory loads.
  \item \textbf{Vertex stage preloads (\hyperref[sec:csf-spd]{\texttt{SPD\_0}/\texttt{SPD\_1}}, \texttt{r16..r19}):} The hardware Index Fetcher reads \hyperref[sec:csf-regs]{\texttt{IDVSSR::INDEX\_BUFFER}} (\texttt{r54:r55}) and pre-fills \texttt{r59, r60} with \texttt{VertexIndex} (\texttt{+ VERTEX\_OFFSET + GLOBAL\_ATTR\_OFF}) and \texttt{r61} with \texttt{InstanceIndex} (\texttt{+ INSTANCE\_OFFSET}) so Position and Varying shaders can pull vertex attributes directly.
  \item \textbf{Inline FAU \& Special Operands (\hyperref[sec:csf-tsd-fau]{\texttt{FAU\_k}}, \texttt{r8..r15}):} Any 8-bit source with \texttt{mode=2} reads zero-latency uniform \texttt{u0..u31} from \texttt{FAU\_k} (\texttt{slots64 = bits[63:56]}); \texttt{mode=3} reads \texttt{SmallConstant} (Table~1.2) or \texttt{FauSpecial} (\texttt{pc}, \texttt{lane\_id}, \texttt{lane\_width}, \texttt{tls\_ptr}, \texttt{wls\_ptr}, \texttt{blend\_descriptor\_0..7}).
\end{tightitemize}
\par\textbf{Worked Examples:}
\par\textbf{Example 1 (Upper halfword FP16 -\textgreater{} FP32):}
\codebox{\mnem{F16\_TO\_F32} \dst{r0}, \src{r1}\imod{.h1}}
With \texttt{r1 = 0x3C000000}: extracts \texttt{h1 = 0x3C00} (\texttt{1.0h}) -\textgreater{} exact \texttt{r0 = 0x3F800000} (\texttt{1.0f}).
\par\textbf{Example 2 (FP16 subnormal exact promotion):}
\codebox{\mnem{F16\_TO\_F32} \dst{r0}, \src{r1}\imod{.h0}}
With \texttt{r1.h0 = 0x0001} (\texttt{2\textasciicircum{}-24}, smallest FP16 subnormal): normalizes exactly in FP32 -\textgreater{} \texttt{r0 = 0x33800000} (\texttt{2\textasciicircum{}-24f}).
\par\textbf{Mined SPIR-V $\to$ Valhall ISA Example (Exact FP16-to-FP32 lane promotion --- Vertex shader, Mali-G710 (v10)):}
\codebox{\textbf{; SPIR-V (input)}\\\%v = OpExtInst \%v2float \%glsl UnpackHalf2x16 \%packed\\[3pt]\textbf{; Valhall ISA (malioc output)}\\F16\_TO\_F32 r14, r47\textasciicircum{}.h0}
Extracts the low 16-bit half-float lane (\texttt{.h0}) from \texttt{r47} and promotes it losslessly to IEEE-754 FP32 in \texttt{r14}.
\begin{center}\footnotesize
\begin{tabularx}{\linewidth}{@{}lllX@{}}
\toprule
Variant & Arch & Class & Opcode Fields \\
\midrule
\texttt{(none)} & v9--v14 & \texttt{R1} & \texttt{opcode2=F16\_TO\_F32, opcode1=CONVERT} \\
\texttt{(none)} & v9--v14 & \texttt{R1} & \texttt{opcode2=F16\_TO\_F32, opcode1=CONVERT} \\
\bottomrule
\end{tabularx}
\end{center}
\begin{center}\footnotesize
\begin{tabularx}{0.88\linewidth}{@{}l l X@{}}
\toprule
Field Name & Bits & Type / Fixed Value \\
\midrule
\texttt{src0} & \texttt{0:7} & \texttt{SourceEncoding} \\
\textit{Reserved} & \texttt{8:15} & \texttt{zero} \\
\texttt{opcode2} & \texttt{16:20} & \texttt{exact=F16\_TO\_F32} \\
\textit{Reserved} & \texttt{21:27} & \texttt{zero} \\
\hyperref[enum:lane16_m]{\texttt{lane0}} & \texttt{28} & \texttt{lane16\_m (def h0)} \\
\textit{Reserved} & \texttt{29:37} & \texttt{zero} \\
\hyperref[enum:neg_m]{\texttt{neg0}} & \texttt{38} & \texttt{neg\_m (def none)} \\
\hyperref[enum:abs_m]{\texttt{abs0}} & \texttt{39} & \texttt{abs\_m (def none)} \\
\texttt{dest\_register\_index} & \texttt{40:45} & \texttt{uint} \\
\hyperref[enum:dest_width_m]{\texttt{dest\_width}} & \texttt{46:47} & \texttt{dest\_width\_m} \\
\texttt{opcode1} & \texttt{48:51} & \texttt{exact=CONVERT} \\
\texttt{instruction\_class} & \texttt{52:56} & \texttt{exact=R1} \\
\texttt{fau\_page\_index} & \texttt{57:58} & \texttt{uint} \\
\hyperref[enum:flow_control_m]{\texttt{flow}} & \texttt{59:62} & \texttt{flow\_control\_m} \\
\textit{Reserved} & \texttt{63} & \texttt{zero} \\
\bottomrule
\end{tabularx}
\end{center}
\textbf{Fields Explained:}
\begin{tightitemize}\footnotesize
  \item \texttt{src0}: Primary source operand (\texttt{r0..r63}, \texttt{u0..u63} uniform, \texttt{k*} FAU constant, or \texttt{t0..t2} temporary; \texttt{\textasciicircum{}} marks discard/last-use). \emph{Example:} {\scriptsize\ttfamily \mnem{FADD.f32} \dst{r0}, \src{r1}, \src{r2}} (\textbf{SPIR-V:} \texttt{OpFAdd (Operand 0)}; \textbf{GLSL:} \texttt{float c = a + b;})
  \item \texttt{opcode2}: Secondary sub-opcode further specializing instructions sharing \texttt{instruction\_class} and \texttt{opcode1} (e.g., conversions, SFU tables, moves). \emph{Example:} {\scriptsize\ttfamily \mnem{F32\_TO\_S32}\imod{.round\_zero} \dst{r0}, \src{r1}} (\textbf{SPIR-V:} \texttt{OpConvertFToS}; \textbf{GLSL:} \texttt{int y = int(x);})
  \item \hyperref[enum:lane16_m]{\texttt{lane0}}: Stage A single-lane extractor on \texttt{src0} (\texttt{b0..b3} for 8-bit or \texttt{h0}/\texttt{h1} for 16-bit narrow conversions and \texttt{MKVEC}). \emph{Example:} {\scriptsize\ttfamily \mnem{F16\_TO\_F32} \dst{r0}, \src{r1}\imod{.h1}} (\textbf{SPIR-V:} \texttt{OpCompositeExtract}; \textbf{GLSL:} \texttt{float y = float(x.y);})
  \item \hyperref[enum:neg_m]{\texttt{neg0}}: Stage A floating-point negation on \texttt{src0} (flips IEEE sign bit after swizzle/widen and \texttt{abs0}). \emph{Example:} {\scriptsize\ttfamily \mnem{FADD.f32} \dst{r0}, \src{r1}\imod{.neg}, \src{r2}} (\textbf{SPIR-V:} \texttt{OpFNegate}; \textbf{GLSL:} \texttt{float c = -a + b;})
  \item \hyperref[enum:abs_m]{\texttt{abs0}}: Stage A floating-point absolute value on \texttt{src0} (clears IEEE sign bit after swizzle/widen, before \texttt{neg0}). \emph{Example:} {\scriptsize\ttfamily \mnem{FADD.f32} \dst{r0}, \src{r1}\imod{.abs}, \src{r2}} (\textbf{SPIR-V:} \texttt{GLSL.std.450 FAbs}; \textbf{GLSL:} \texttt{float c = abs(a) + b;})
  \item \texttt{dest\_register\_index}: 6-bit destination GPR index (\texttt{r0..r63}) receiving the Stage C writeback result (even-aligned pair \texttt{[rN:rN+1]} for 64-bit ops). \emph{Example:} {\scriptsize\ttfamily \mnem{FMA.f32} \dst{r0}, \src{r1}, \src{r2}, \src{r3}} (\textbf{SPIR-V:} \texttt{SSA Result ID}; \textbf{GLSL:} \texttt{float r0 = fma(r1, r2, r3);})
  \item \hyperref[enum:dest_width_m]{\texttt{dest\_width}}: Stage C destination lane write-mask (\texttt{none} = full 32-bit, \texttt{h0} = low 16-bit, \texttt{h1} = high 16-bit without bit-shifting). \emph{Example:} {\scriptsize\ttfamily \mnem{FADD.v2f16} \dst{r0}\imod{.h1}, \src{r1}, \src{r2}} (\textbf{SPIR-V:} \texttt{OpVectorInsertDynamic / WriteMask}; \textbf{GLSL:} \texttt{dst.y = a.y + b.y;})
  \item \texttt{opcode1}: Primary sub-opcode within the instruction class uniquely identifying the operation family. \emph{Example:} {\scriptsize\ttfamily \mnem{IADD.s32} \dst{r0}, \src{r1}, \src{r2}} (\textbf{SPIR-V:} \texttt{OpIAdd (operation selector)}; \textbf{GLSL:} \texttt{// Primary opcode bits})
  \item \texttt{instruction\_class}: 4-bit major instruction class selector (\texttt{bits 48:51}) partitioning the 64-bit Valhall opcode space (e.g., \texttt{FMA}, \texttt{CVT}, \texttt{SFU}, \texttt{LS}, \texttt{TEX}, \texttt{BRANCH}). \emph{Example:} {\scriptsize\ttfamily \mnem{FADD.f32} \dst{r0}, \src{r1}, \src{r2}} (\textbf{SPIR-V:} \texttt{N/A (encoding class)}; \textbf{GLSL:} \texttt{// Hardware execution pipe / decode class})
  \item \texttt{fau\_page\_index}: Selects the FAU (Fast Access Uniform) page (\texttt{0..3}) when a source operand references hardware small constants, thread IDs, or lane constants. \emph{Example:} {\scriptsize\ttfamily \mnem{FADD.f32} \dst{r0}, \src{r1}, \src{k1f}} (\textbf{SPIR-V:} \texttt{OpConstant / BuiltIn}; \textbf{GLSL:} \texttt{float y = x + 1.0;})
  \item \hyperref[enum:flow_control_m]{\texttt{flow}}: 4-bit embedded flow-control \& scoreboard wait modifier (\texttt{bits 59:62}: \texttt{none}, \texttt{wait0..wait0126}, \texttt{wait}, \texttt{reconverge}, \texttt{discard}, \texttt{end}). \emph{Example:} {\scriptsize\ttfamily \mnem{FADD.f32}\imod{.wait0.reconverge} \dst{r0}, \src{r1}, \src{r2}} (\textbf{SPIR-V:} \texttt{OpControlBarrier / OpReturn}; \textbf{GLSL:} \texttt{// Scoreboard wait + warp reconvergence})
\end{tightitemize}
\clearpage

\subsection*{F16\_TO\_S32}
\addcontentsline{toc}{subsection}{F16\_TO\_S32}
Saturating f16-\textgreater{}s32 (arch \textless{}= v10): \texttt{x = round\_rm(f16(src0.lane0))}; NaN-\textgreater{}0, \texttt{-Inf -\textgreater{} INT32\_MIN}, \texttt{+Inf -\textgreater{} INT32\_MAX}, else \texttt{x as i32}.
\par\textbf{Notes:}
\begin{tightitemize}
  \item Arch v9-v10 only (\texttt{kraid/ops.rs:1196}, \texttt{nir.rs:727}).
  \item \textbf{Source references:} \texttt{kraid/ops.rs:1196-1221}, \texttt{encode\_v9.rs:1168-1198}, \texttt{nir.rs:726-734}.
\end{tightitemize}
\codebox{\mnem{F16\_TO\_S32}\imod{\{.flow\}}\imod{\{.round\}} \dst{dest}\imod{\{.dest\_width\}}, \src{src0}\imod{\{.lane0\}}\imod{\{.abs0\}}\imod{\{.neg0\}}}
\textbf{Operands:}
\begin{tightitemize}\footnotesize
  \item \texttt{\dst{dest}} (\texttt{dest\_register\_index}, \texttt{dest\_width}): 32-bit destination GPR receiving the converted value.
  \item \texttt{\src{src0}}: FP16 input lane (\texttt{lane0}) rounded and saturated to signed 32-bit integer.
\end{tightitemize}
\par\textbf{Closest SPIR-V Op:} \texttt{OpConvertFToS (f16 -\textgreater{} s32, v9--v10)}
\par\textbf{Pseudo-GLSL (Operational Semantics):}
\codebox{float h = unpackHalf2x16(src0)[lane0];\\dest = isnan(h) ? 0 : int(clamp(round\_rm(h), float(INT32\_MIN), float(INT32\_MAX)));}
\textbf{Execution Context:}
\begin{tightitemize}\footnotesize
  \item \textbf{Compute stage preloads (\hyperref[sec:csf-spd]{\texttt{SPD\_k}}, \texttt{r16..r23}):} Hardware pre-fills \texttt{r55..r57} with \texttt{WorkGroupID.xyz + JOB\_OFFSET} (from \hyperref[sec:csf-regs]{\texttt{COMPUTESR::JOB\_OFFSET\_X/Y/Z}} \texttt{r34..r36}) and \texttt{r60..r62} with \texttt{LocalInvocationID.xyz} or \texttt{GlobalInvocationID.xyz} (from \texttt{WG\_SIZE} \texttt{r33} \& \texttt{JOB\_SIZE} \texttt{r37..r39}, selected by \texttt{ShaderProgram::preload}) so compute shaders can derive thread/element indices without memory loads.
  \item \textbf{Vertex stage preloads (\hyperref[sec:csf-spd]{\texttt{SPD\_0}/\texttt{SPD\_1}}, \texttt{r16..r19}):} The hardware Index Fetcher reads \hyperref[sec:csf-regs]{\texttt{IDVSSR::INDEX\_BUFFER}} (\texttt{r54:r55}) and pre-fills \texttt{r59, r60} with \texttt{VertexIndex} (\texttt{+ VERTEX\_OFFSET + GLOBAL\_ATTR\_OFF}) and \texttt{r61} with \texttt{InstanceIndex} (\texttt{+ INSTANCE\_OFFSET}) so Position and Varying shaders can pull vertex attributes directly.
  \item \textbf{Inline FAU \& Special Operands (\hyperref[sec:csf-tsd-fau]{\texttt{FAU\_k}}, \texttt{r8..r15}):} Any 8-bit source with \texttt{mode=2} reads zero-latency uniform \texttt{u0..u31} from \texttt{FAU\_k} (\texttt{slots64 = bits[63:56]}); \texttt{mode=3} reads \texttt{SmallConstant} (Table~1.2) or \texttt{FauSpecial} (\texttt{pc}, \texttt{lane\_id}, \texttt{lane\_width}, \texttt{tls\_ptr}, \texttt{wls\_ptr}, \texttt{blend\_descriptor\_0..7}).
\end{tightitemize}
\par\textbf{Worked Examples:}
\par\textbf{Example 1 (Negative half-float truncation (v9--v10)):}
\codebox{\mnem{F16\_TO\_S32}\imod{.round\_zero} \dst{r0}, \src{r1}\imod{.h0}}
With \texttt{r1.h0 = 0xC280} (\texttt{-3.25h}): \texttt{.round\_zero} truncates toward zero -\textgreater{} \texttt{r0 = -3} (\texttt{0xFFFFFFFD}).
\par\textbf{Example 2 (NaN and infinity saturation):}
\codebox{\mnem{F16\_TO\_S32}\imod{.round\_zero} \dst{r0}, \src{r1}\imod{.h0}}
With \texttt{r1.h0 = 0x7E00} (\texttt{NaN}) -\textgreater{} \texttt{r0 = 0}; with \texttt{r1.h0 = 0xFC00} (\texttt{-Inf}) -\textgreater{} \texttt{r0 = 0x80000000} (\texttt{INT32\_MIN}).

\begin{center}\footnotesize
\begin{tabularx}{\linewidth}{@{}lllX@{}}
\toprule
Variant & Arch & Class & Opcode Fields \\
\midrule
\texttt{(none)} & v9--v14 & \texttt{R1} & \texttt{opcode2=F16\_TO\_S32, opcode1=CONVERT} \\
\bottomrule
\end{tabularx}
\end{center}
\begin{center}\footnotesize
\begin{tabularx}{0.88\linewidth}{@{}l l X@{}}
\toprule
Field Name & Bits & Type / Fixed Value \\
\midrule
\texttt{src0} & \texttt{0:7} & \texttt{SourceEncoding} \\
\textit{Reserved} & \texttt{8:15} & \texttt{zero} \\
\texttt{opcode2} & \texttt{16:20} & \texttt{exact=F16\_TO\_S32} \\
\textit{Reserved} & \texttt{21:27} & \texttt{zero} \\
\hyperref[enum:lane16_m]{\texttt{lane0}} & \texttt{28} & \texttt{lane16\_m (def h0)} \\
\textit{Reserved} & \texttt{29} & \texttt{zero} \\
\hyperref[enum:round_integer_m]{\texttt{round}} & \texttt{30:32} & \texttt{round\_integer\_m (def none)} \\
\textit{Reserved} & \texttt{33} & \texttt{zero} \\
\textit{Reserved} & \texttt{34:35} & \texttt{ignore} \\
\textit{Reserved} & \texttt{36:37} & \texttt{zero} \\
\hyperref[enum:neg_m]{\texttt{neg0}} & \texttt{38} & \texttt{neg\_m (def none)} \\
\hyperref[enum:abs_m]{\texttt{abs0}} & \texttt{39} & \texttt{abs\_m (def none)} \\
\texttt{dest\_register\_index} & \texttt{40:45} & \texttt{uint} \\
\hyperref[enum:dest_width_m]{\texttt{dest\_width}} & \texttt{46:47} & \texttt{dest\_width\_m} \\
\texttt{opcode1} & \texttt{48:51} & \texttt{exact=CONVERT} \\
\texttt{instruction\_class} & \texttt{52:56} & \texttt{exact=R1} \\
\texttt{fau\_page\_index} & \texttt{57:58} & \texttt{uint} \\
\hyperref[enum:flow_control_m]{\texttt{flow}} & \texttt{59:62} & \texttt{flow\_control\_m} \\
\textit{Reserved} & \texttt{63} & \texttt{zero} \\
\bottomrule
\end{tabularx}
\end{center}
\textbf{Fields Explained:}
\begin{tightitemize}\footnotesize
  \item \texttt{src0}: Primary source operand (\texttt{r0..r63}, \texttt{u0..u63} uniform, \texttt{k*} FAU constant, or \texttt{t0..t2} temporary; \texttt{\textasciicircum{}} marks discard/last-use). \emph{Example:} {\scriptsize\ttfamily \mnem{FADD.f32} \dst{r0}, \src{r1}, \src{r2}} (\textbf{SPIR-V:} \texttt{OpFAdd (Operand 0)}; \textbf{GLSL:} \texttt{float c = a + b;})
  \item \texttt{opcode2}: Secondary sub-opcode further specializing instructions sharing \texttt{instruction\_class} and \texttt{opcode1} (e.g., conversions, SFU tables, moves). \emph{Example:} {\scriptsize\ttfamily \mnem{F32\_TO\_S32}\imod{.round\_zero} \dst{r0}, \src{r1}} (\textbf{SPIR-V:} \texttt{OpConvertFToS}; \textbf{GLSL:} \texttt{int y = int(x);})
  \item \hyperref[enum:lane16_m]{\texttt{lane0}}: Stage A single-lane extractor on \texttt{src0} (\texttt{b0..b3} for 8-bit or \texttt{h0}/\texttt{h1} for 16-bit narrow conversions and \texttt{MKVEC}). \emph{Example:} {\scriptsize\ttfamily \mnem{F16\_TO\_F32} \dst{r0}, \src{r1}\imod{.h1}} (\textbf{SPIR-V:} \texttt{OpCompositeExtract}; \textbf{GLSL:} \texttt{float y = float(x.y);})
  \item \hyperref[enum:round_integer_m]{\texttt{round}}: Rounding mode (\texttt{none}/RNE, \texttt{round\_up}, \texttt{round\_down}, \texttt{round\_zero}, \texttt{round\_na}) applied in FP ALU, conversions, \texttt{HADD}, \texttt{LDEXP}, and Stage C narrowing. \emph{Example:} {\scriptsize\ttfamily \mnem{FROUND.f32}\imod{.round\_down} \dst{r0}, \src{r1}} (\textbf{SPIR-V:} \texttt{GLSL.std.450 Floor / Ceil / Trunc}; \textbf{GLSL:} \texttt{float y = floor(x);})
  \item \hyperref[enum:neg_m]{\texttt{neg0}}: Stage A floating-point negation on \texttt{src0} (flips IEEE sign bit after swizzle/widen and \texttt{abs0}). \emph{Example:} {\scriptsize\ttfamily \mnem{FADD.f32} \dst{r0}, \src{r1}\imod{.neg}, \src{r2}} (\textbf{SPIR-V:} \texttt{OpFNegate}; \textbf{GLSL:} \texttt{float c = -a + b;})
  \item \hyperref[enum:abs_m]{\texttt{abs0}}: Stage A floating-point absolute value on \texttt{src0} (clears IEEE sign bit after swizzle/widen, before \texttt{neg0}). \emph{Example:} {\scriptsize\ttfamily \mnem{FADD.f32} \dst{r0}, \src{r1}\imod{.abs}, \src{r2}} (\textbf{SPIR-V:} \texttt{GLSL.std.450 FAbs}; \textbf{GLSL:} \texttt{float c = abs(a) + b;})
  \item \texttt{dest\_register\_index}: 6-bit destination GPR index (\texttt{r0..r63}) receiving the Stage C writeback result (even-aligned pair \texttt{[rN:rN+1]} for 64-bit ops). \emph{Example:} {\scriptsize\ttfamily \mnem{FMA.f32} \dst{r0}, \src{r1}, \src{r2}, \src{r3}} (\textbf{SPIR-V:} \texttt{SSA Result ID}; \textbf{GLSL:} \texttt{float r0 = fma(r1, r2, r3);})
  \item \hyperref[enum:dest_width_m]{\texttt{dest\_width}}: Stage C destination lane write-mask (\texttt{none} = full 32-bit, \texttt{h0} = low 16-bit, \texttt{h1} = high 16-bit without bit-shifting). \emph{Example:} {\scriptsize\ttfamily \mnem{FADD.v2f16} \dst{r0}\imod{.h1}, \src{r1}, \src{r2}} (\textbf{SPIR-V:} \texttt{OpVectorInsertDynamic / WriteMask}; \textbf{GLSL:} \texttt{dst.y = a.y + b.y;})
  \item \texttt{opcode1}: Primary sub-opcode within the instruction class uniquely identifying the operation family. \emph{Example:} {\scriptsize\ttfamily \mnem{IADD.s32} \dst{r0}, \src{r1}, \src{r2}} (\textbf{SPIR-V:} \texttt{OpIAdd (operation selector)}; \textbf{GLSL:} \texttt{// Primary opcode bits})
  \item \texttt{instruction\_class}: 4-bit major instruction class selector (\texttt{bits 48:51}) partitioning the 64-bit Valhall opcode space (e.g., \texttt{FMA}, \texttt{CVT}, \texttt{SFU}, \texttt{LS}, \texttt{TEX}, \texttt{BRANCH}). \emph{Example:} {\scriptsize\ttfamily \mnem{FADD.f32} \dst{r0}, \src{r1}, \src{r2}} (\textbf{SPIR-V:} \texttt{N/A (encoding class)}; \textbf{GLSL:} \texttt{// Hardware execution pipe / decode class})
  \item \texttt{fau\_page\_index}: Selects the FAU (Fast Access Uniform) page (\texttt{0..3}) when a source operand references hardware small constants, thread IDs, or lane constants. \emph{Example:} {\scriptsize\ttfamily \mnem{FADD.f32} \dst{r0}, \src{r1}, \src{k1f}} (\textbf{SPIR-V:} \texttt{OpConstant / BuiltIn}; \textbf{GLSL:} \texttt{float y = x + 1.0;})
  \item \hyperref[enum:flow_control_m]{\texttt{flow}}: 4-bit embedded flow-control \& scoreboard wait modifier (\texttt{bits 59:62}: \texttt{none}, \texttt{wait0..wait0126}, \texttt{wait}, \texttt{reconverge}, \texttt{discard}, \texttt{end}). \emph{Example:} {\scriptsize\ttfamily \mnem{FADD.f32}\imod{.wait0.reconverge} \dst{r0}, \src{r1}, \src{r2}} (\textbf{SPIR-V:} \texttt{OpControlBarrier / OpReturn}; \textbf{GLSL:} \texttt{// Scoreboard wait + warp reconvergence})
\end{tightitemize}
\clearpage

\subsection*{F16\_TO\_U32}
\addcontentsline{toc}{subsection}{F16\_TO\_U32}
Converts the FP16 halfword selected by \texttt{lane0} of \texttt{src0} to a saturating unsigned 32-bit integer in \texttt{dest} (arch \textless{}= v10): rounds per \texttt{round}, returning \texttt{0} for \texttt{NaN} or \texttt{\textless{}= 0.0}, \texttt{UINT32\_MAX} for \texttt{+Inf}, and \texttt{uint32(x)} otherwise.
\par\textbf{Notes:}
\begin{tightitemize}
  \item Arch v9-v10 only (\texttt{kraid/ops.rs:1196}, \texttt{nir.rs:727}).
  \item \textbf{Source references:} \texttt{kraid/ops.rs:1196-1221}, \texttt{encode\_v9.rs:1168-1198}, \texttt{nir.rs:726-734}.
\end{tightitemize}
\codebox{\mnem{F16\_TO\_U32}\imod{\{.flow\}}\imod{\{.round\}} \dst{dest}\imod{\{.dest\_width\}}, \src{src0}\imod{\{.lane0\}}\imod{\{.abs0\}}\imod{\{.neg0\}}}
\textbf{Operands:}
\begin{tightitemize}\footnotesize
  \item \texttt{\dst{dest}} (\texttt{dest\_register\_index}, \texttt{dest\_width}): 32-bit destination GPR receiving the converted value.
  \item \texttt{\src{src0}}: FP16 input lane (\texttt{lane0}) rounded and saturated to unsigned 32-bit integer.
\end{tightitemize}
\par\textbf{Closest SPIR-V Op:} \texttt{OpConvertFToU (f16 -\textgreater{} u32, v9--v10)}
\par\textbf{Pseudo-GLSL (Operational Semantics):}
\codebox{float h = unpackHalf2x16(src0)[lane0];\\dest = (isnan(h) \textbar{}\textbar{} h \textless{}= 0.0) ? 0u : uint(clamp(round\_rm(h), 0.0, float(UINT32\_MAX)));}
\textbf{Execution Context:}
\begin{tightitemize}\footnotesize
  \item \textbf{Compute stage preloads (\hyperref[sec:csf-spd]{\texttt{SPD\_k}}, \texttt{r16..r23}):} Hardware pre-fills \texttt{r55..r57} with \texttt{WorkGroupID.xyz + JOB\_OFFSET} (from \hyperref[sec:csf-regs]{\texttt{COMPUTESR::JOB\_OFFSET\_X/Y/Z}} \texttt{r34..r36}) and \texttt{r60..r62} with \texttt{LocalInvocationID.xyz} or \texttt{GlobalInvocationID.xyz} (from \texttt{WG\_SIZE} \texttt{r33} \& \texttt{JOB\_SIZE} \texttt{r37..r39}, selected by \texttt{ShaderProgram::preload}) so compute shaders can derive thread/element indices without memory loads.
  \item \textbf{Vertex stage preloads (\hyperref[sec:csf-spd]{\texttt{SPD\_0}/\texttt{SPD\_1}}, \texttt{r16..r19}):} The hardware Index Fetcher reads \hyperref[sec:csf-regs]{\texttt{IDVSSR::INDEX\_BUFFER}} (\texttt{r54:r55}) and pre-fills \texttt{r59, r60} with \texttt{VertexIndex} (\texttt{+ VERTEX\_OFFSET + GLOBAL\_ATTR\_OFF}) and \texttt{r61} with \texttt{InstanceIndex} (\texttt{+ INSTANCE\_OFFSET}) so Position and Varying shaders can pull vertex attributes directly.
  \item \textbf{Inline FAU \& Special Operands (\hyperref[sec:csf-tsd-fau]{\texttt{FAU\_k}}, \texttt{r8..r15}):} Any 8-bit source with \texttt{mode=2} reads zero-latency uniform \texttt{u0..u31} from \texttt{FAU\_k} (\texttt{slots64 = bits[63:56]}); \texttt{mode=3} reads \texttt{SmallConstant} (Table~1.2) or \texttt{FauSpecial} (\texttt{pc}, \texttt{lane\_id}, \texttt{lane\_width}, \texttt{tls\_ptr}, \texttt{wls\_ptr}, \texttt{blend\_descriptor\_0..7}).
\end{tightitemize}
\par\textbf{Worked Examples:}
\par\textbf{Example 1 (Positive half-float conversion (v9--v10)):}
\codebox{\mnem{F16\_TO\_U32}\imod{.round\_zero} \dst{r0}, \src{r1}\imod{.h0}}
With \texttt{r1.h0 = 0x4580} (\texttt{5.5h}): truncates to \texttt{r0 = 5}.
\par\textbf{Example 2 (Negative / NaN clamp to 0):}
\codebox{\mnem{F16\_TO\_U32}\imod{.round\_zero} \dst{r0}, \src{r1}\imod{.h0}}
With \texttt{r1.h0 = 0xC000} (\texttt{-2.0h}) or \texttt{NaN}: clamps to \texttt{r0 = 0}; \texttt{+Inf} saturates to \texttt{0xFFFFFFFF}.

\begin{center}\footnotesize
\begin{tabularx}{\linewidth}{@{}lllX@{}}
\toprule
Variant & Arch & Class & Opcode Fields \\
\midrule
\texttt{(none)} & v9--v14 & \texttt{R1} & \texttt{opcode2=F16\_TO\_U32, opcode1=CONVERT} \\
\bottomrule
\end{tabularx}
\end{center}
\begin{center}\footnotesize
\begin{tabularx}{0.88\linewidth}{@{}l l X@{}}
\toprule
Field Name & Bits & Type / Fixed Value \\
\midrule
\texttt{src0} & \texttt{0:7} & \texttt{SourceEncoding} \\
\textit{Reserved} & \texttt{8:15} & \texttt{zero} \\
\texttt{opcode2} & \texttt{16:20} & \texttt{exact=F16\_TO\_U32} \\
\textit{Reserved} & \texttt{21:27} & \texttt{zero} \\
\hyperref[enum:lane16_m]{\texttt{lane0}} & \texttt{28} & \texttt{lane16\_m (def h0)} \\
\textit{Reserved} & \texttt{29} & \texttt{zero} \\
\hyperref[enum:round_integer_m]{\texttt{round}} & \texttt{30:32} & \texttt{round\_integer\_m (def none)} \\
\textit{Reserved} & \texttt{33} & \texttt{zero} \\
\textit{Reserved} & \texttt{34:35} & \texttt{ignore} \\
\textit{Reserved} & \texttt{36:37} & \texttt{zero} \\
\hyperref[enum:neg_m]{\texttt{neg0}} & \texttt{38} & \texttt{neg\_m (def none)} \\
\hyperref[enum:abs_m]{\texttt{abs0}} & \texttt{39} & \texttt{abs\_m (def none)} \\
\texttt{dest\_register\_index} & \texttt{40:45} & \texttt{uint} \\
\hyperref[enum:dest_width_m]{\texttt{dest\_width}} & \texttt{46:47} & \texttt{dest\_width\_m} \\
\texttt{opcode1} & \texttt{48:51} & \texttt{exact=CONVERT} \\
\texttt{instruction\_class} & \texttt{52:56} & \texttt{exact=R1} \\
\texttt{fau\_page\_index} & \texttt{57:58} & \texttt{uint} \\
\hyperref[enum:flow_control_m]{\texttt{flow}} & \texttt{59:62} & \texttt{flow\_control\_m} \\
\textit{Reserved} & \texttt{63} & \texttt{zero} \\
\bottomrule
\end{tabularx}
\end{center}
\textbf{Fields Explained:}
\begin{tightitemize}\footnotesize
  \item \texttt{src0}: Primary source operand (\texttt{r0..r63}, \texttt{u0..u63} uniform, \texttt{k*} FAU constant, or \texttt{t0..t2} temporary; \texttt{\textasciicircum{}} marks discard/last-use). \emph{Example:} {\scriptsize\ttfamily \mnem{FADD.f32} \dst{r0}, \src{r1}, \src{r2}} (\textbf{SPIR-V:} \texttt{OpFAdd (Operand 0)}; \textbf{GLSL:} \texttt{float c = a + b;})
  \item \texttt{opcode2}: Secondary sub-opcode further specializing instructions sharing \texttt{instruction\_class} and \texttt{opcode1} (e.g., conversions, SFU tables, moves). \emph{Example:} {\scriptsize\ttfamily \mnem{F32\_TO\_S32}\imod{.round\_zero} \dst{r0}, \src{r1}} (\textbf{SPIR-V:} \texttt{OpConvertFToS}; \textbf{GLSL:} \texttt{int y = int(x);})
  \item \hyperref[enum:lane16_m]{\texttt{lane0}}: Stage A single-lane extractor on \texttt{src0} (\texttt{b0..b3} for 8-bit or \texttt{h0}/\texttt{h1} for 16-bit narrow conversions and \texttt{MKVEC}). \emph{Example:} {\scriptsize\ttfamily \mnem{F16\_TO\_F32} \dst{r0}, \src{r1}\imod{.h1}} (\textbf{SPIR-V:} \texttt{OpCompositeExtract}; \textbf{GLSL:} \texttt{float y = float(x.y);})
  \item \hyperref[enum:round_integer_m]{\texttt{round}}: Rounding mode (\texttt{none}/RNE, \texttt{round\_up}, \texttt{round\_down}, \texttt{round\_zero}, \texttt{round\_na}) applied in FP ALU, conversions, \texttt{HADD}, \texttt{LDEXP}, and Stage C narrowing. \emph{Example:} {\scriptsize\ttfamily \mnem{FROUND.f32}\imod{.round\_down} \dst{r0}, \src{r1}} (\textbf{SPIR-V:} \texttt{GLSL.std.450 Floor / Ceil / Trunc}; \textbf{GLSL:} \texttt{float y = floor(x);})
  \item \hyperref[enum:neg_m]{\texttt{neg0}}: Stage A floating-point negation on \texttt{src0} (flips IEEE sign bit after swizzle/widen and \texttt{abs0}). \emph{Example:} {\scriptsize\ttfamily \mnem{FADD.f32} \dst{r0}, \src{r1}\imod{.neg}, \src{r2}} (\textbf{SPIR-V:} \texttt{OpFNegate}; \textbf{GLSL:} \texttt{float c = -a + b;})
  \item \hyperref[enum:abs_m]{\texttt{abs0}}: Stage A floating-point absolute value on \texttt{src0} (clears IEEE sign bit after swizzle/widen, before \texttt{neg0}). \emph{Example:} {\scriptsize\ttfamily \mnem{FADD.f32} \dst{r0}, \src{r1}\imod{.abs}, \src{r2}} (\textbf{SPIR-V:} \texttt{GLSL.std.450 FAbs}; \textbf{GLSL:} \texttt{float c = abs(a) + b;})
  \item \texttt{dest\_register\_index}: 6-bit destination GPR index (\texttt{r0..r63}) receiving the Stage C writeback result (even-aligned pair \texttt{[rN:rN+1]} for 64-bit ops). \emph{Example:} {\scriptsize\ttfamily \mnem{FMA.f32} \dst{r0}, \src{r1}, \src{r2}, \src{r3}} (\textbf{SPIR-V:} \texttt{SSA Result ID}; \textbf{GLSL:} \texttt{float r0 = fma(r1, r2, r3);})
  \item \hyperref[enum:dest_width_m]{\texttt{dest\_width}}: Stage C destination lane write-mask (\texttt{none} = full 32-bit, \texttt{h0} = low 16-bit, \texttt{h1} = high 16-bit without bit-shifting). \emph{Example:} {\scriptsize\ttfamily \mnem{FADD.v2f16} \dst{r0}\imod{.h1}, \src{r1}, \src{r2}} (\textbf{SPIR-V:} \texttt{OpVectorInsertDynamic / WriteMask}; \textbf{GLSL:} \texttt{dst.y = a.y + b.y;})
  \item \texttt{opcode1}: Primary sub-opcode within the instruction class uniquely identifying the operation family. \emph{Example:} {\scriptsize\ttfamily \mnem{IADD.s32} \dst{r0}, \src{r1}, \src{r2}} (\textbf{SPIR-V:} \texttt{OpIAdd (operation selector)}; \textbf{GLSL:} \texttt{// Primary opcode bits})
  \item \texttt{instruction\_class}: 4-bit major instruction class selector (\texttt{bits 48:51}) partitioning the 64-bit Valhall opcode space (e.g., \texttt{FMA}, \texttt{CVT}, \texttt{SFU}, \texttt{LS}, \texttt{TEX}, \texttt{BRANCH}). \emph{Example:} {\scriptsize\ttfamily \mnem{FADD.f32} \dst{r0}, \src{r1}, \src{r2}} (\textbf{SPIR-V:} \texttt{N/A (encoding class)}; \textbf{GLSL:} \texttt{// Hardware execution pipe / decode class})
  \item \texttt{fau\_page\_index}: Selects the FAU (Fast Access Uniform) page (\texttt{0..3}) when a source operand references hardware small constants, thread IDs, or lane constants. \emph{Example:} {\scriptsize\ttfamily \mnem{FADD.f32} \dst{r0}, \src{r1}, \src{k1f}} (\textbf{SPIR-V:} \texttt{OpConstant / BuiltIn}; \textbf{GLSL:} \texttt{float y = x + 1.0;})
  \item \hyperref[enum:flow_control_m]{\texttt{flow}}: 4-bit embedded flow-control \& scoreboard wait modifier (\texttt{bits 59:62}: \texttt{none}, \texttt{wait0..wait0126}, \texttt{wait}, \texttt{reconverge}, \texttt{discard}, \texttt{end}). \emph{Example:} {\scriptsize\ttfamily \mnem{FADD.f32}\imod{.wait0.reconverge} \dst{r0}, \src{r1}, \src{r2}} (\textbf{SPIR-V:} \texttt{OpControlBarrier / OpReturn}; \textbf{GLSL:} \texttt{// Scoreboard wait + warp reconvergence})
\end{tightitemize}
\clearpage

\subsection*{F32\_TO\_F16}
\addcontentsline{toc}{subsection}{F32\_TO\_F16}
Narrows FP32 \texttt{src0} to IEEE-754 FP16 after applying optional \texttt{clamp} and \texttt{round} modifiers, and writes the 16-bit result into \texttt{dest.h0}, \texttt{dest.h1} (preserving the other halfword), or replicated into both halves (\texttt{dest.h01}) per \texttt{dest\_width\_narrow} (v11--v14 alias of \texttt{FADD.f32(src0, -0.0)}).
\par\textbf{Notes:}
\begin{tightitemize}
  \item Arch v11-v14 only (v9-v10 uses \texttt{V2F32\_TO\_V2F16} or \texttt{FADD.f32.hf0/hf1}).
  \item Single source (\texttt{src1} fixed to \texttt{0.0} with \texttt{neg1=neg} in \texttt{malioc} v12 to preserve \texttt{-0.0}; explicit \texttt{neg1=none} adds \texttt{+0.0} flushing \texttt{-0.0 -\textgreater{} +0.0}).
  \item \textbf{Source references:} \texttt{kraid/builder.rs:50-65}, \texttt{nir.rs:642-663}, \texttt{isa-v9-v14.xml}.
\end{tightitemize}
\codebox{\mnem{F32\_TO\_F16}\imod{\{.flow\}}\imod{\{.round\}}\imod{\{.clamp\}} \dst{dest}\imod{\{.dest\_width\_narrow\}}, \src{src0}\imod{\{.abs0\}}\imod{\{.neg0\}}}
\textbf{Operands:}
\begin{tightitemize}\footnotesize
  \item \texttt{\dst{dest}} (\texttt{dest\_register\_index}, \texttt{dest\_width\_narrow}): 32-bit destination GPR receiving the converted value.
  \item \texttt{\src{src0}}: FP32 input value narrowed to FP16 into \texttt{dest.h0}, \texttt{dest.h1}, or replicated \texttt{dest.h01}.
\end{tightitemize}
\par\textbf{Closest SPIR-V Op:} \texttt{OpFConvert (f32 -\textgreater{} f16, v11+)}
\par\textbf{Pseudo-GLSL (Operational Semantics):}
\codebox{uint h = floatToHalfBits(round\_rm(clamp\_mode(src0)));\\if (dest\_width\_narrow == h01) dest = h \textbar{} (h \textless{}\textless{} 16);\\else                          dest = writeHalf(dest, dest\_width\_narrow, h);}
\textbf{Execution Context:}
\begin{tightitemize}\footnotesize
  \item \textbf{Compute stage preloads (\hyperref[sec:csf-spd]{\texttt{SPD\_k}}, \texttt{r16..r23}):} Hardware pre-fills \texttt{r55..r57} with \texttt{WorkGroupID.xyz + JOB\_OFFSET} (from \hyperref[sec:csf-regs]{\texttt{COMPUTESR::JOB\_OFFSET\_X/Y/Z}} \texttt{r34..r36}) and \texttt{r60..r62} with \texttt{LocalInvocationID.xyz} or \texttt{GlobalInvocationID.xyz} (from \texttt{WG\_SIZE} \texttt{r33} \& \texttt{JOB\_SIZE} \texttt{r37..r39}, selected by \texttt{ShaderProgram::preload}) so compute shaders can derive thread/element indices without memory loads.
  \item \textbf{Vertex stage preloads (\hyperref[sec:csf-spd]{\texttt{SPD\_0}/\texttt{SPD\_1}}, \texttt{r16..r19}):} The hardware Index Fetcher reads \hyperref[sec:csf-regs]{\texttt{IDVSSR::INDEX\_BUFFER}} (\texttt{r54:r55}) and pre-fills \texttt{r59, r60} with \texttt{VertexIndex} (\texttt{+ VERTEX\_OFFSET + GLOBAL\_ATTR\_OFF}) and \texttt{r61} with \texttt{InstanceIndex} (\texttt{+ INSTANCE\_OFFSET}) so Position and Varying shaders can pull vertex attributes directly.
  \item \textbf{Inline FAU \& Special Operands (\hyperref[sec:csf-tsd-fau]{\texttt{FAU\_k}}, \texttt{r8..r15}):} Any 8-bit source with \texttt{mode=2} reads zero-latency uniform \texttt{u0..u31} from \texttt{FAU\_k} (\texttt{slots64 = bits[63:56]}); \texttt{mode=3} reads \texttt{SmallConstant} (Table~1.2) or \texttt{FauSpecial} (\texttt{pc}, \texttt{lane\_id}, \texttt{lane\_width}, \texttt{tls\_ptr}, \texttt{wls\_ptr}, \texttt{blend\_descriptor\_0..7}).
\end{tightitemize}
\par\textbf{Worked Examples:}
\par\textbf{Example 1 (Upper-half write preserving low half (v11+)):}
\codebox{\mnem{F32\_TO\_F16} \dst{r0}\imod{.h1}, \src{r1}}
With \texttt{r1 = 1.0f}, \texttt{r0\_old = 0x0000BEEF}: converts \texttt{1.0f -\textgreater{} 0x3C00} in upper half -\textgreater{} \texttt{r0 = 0x3C00BEEF}.
\par\textbf{Example 2 (Pre-narrowing clamp + broadcast (`.clamp\_0\_1` + `.h01`)):}
\codebox{\mnem{F32\_TO\_F16}\imod{.clamp\_0\_1} \dst{r0}\imod{.h01}, \src{r1}}
With \texttt{r1 = 100000.0f} (\texttt{\textgreater{} FP16 MAX 65504}): \texttt{.clamp\_0\_1} runs BEFORE narrowing (clamping to \texttt{1.0f}, not \texttt{+Inf}) -\textgreater{} \texttt{r0 = 0x3C003C00}.

\begin{center}\footnotesize
\begin{tabularx}{\linewidth}{@{}lllX@{}}
\toprule
Variant & Arch & Class & Opcode Fields \\
\midrule
\texttt{f32} & v9--v14 & \texttt{R2} & \texttt{opcode2=ADD} \\
\bottomrule
\end{tabularx}
\end{center}
\begin{center}\footnotesize
\begin{tabularx}{0.88\linewidth}{@{}l l X@{}}
\toprule
Field Name & Bits & Type / Fixed Value \\
\midrule
\texttt{src0} & \texttt{0:7} & \texttt{SourceEncoding} \\
\texttt{src1} & \texttt{8:15} & \texttt{exact=192} \\
\texttt{opcode2} & \texttt{16:19} & \texttt{exact=ADD} \\
\textit{Reserved} & \texttt{20:25} & \texttt{zero} \\
\hyperref[enum:widen_m]{\texttt{widen1}} & \texttt{26:27} & \texttt{exact=none} \\
\hyperref[enum:widen_m]{\texttt{widen0}} & \texttt{28:29} & \texttt{exact=none} \\
\hyperref[enum:round_m]{\texttt{round}} & \texttt{30:31} & \texttt{round\_m (def none)} \\
\hyperref[enum:clamp_m]{\texttt{clamp}} & \texttt{32:33} & \texttt{clamp\_m (def none)} \\
\hyperref[enum:fadd_sticky_m]{\texttt{sticky}} & \texttt{34} & \texttt{exact=none} \\
\textit{Reserved} & \texttt{35} & \texttt{zero} \\
\hyperref[enum:neg_m]{\texttt{neg1}} & \texttt{36} & \texttt{neg\_m} \\
\hyperref[enum:abs_m]{\texttt{abs2}} & \texttt{37} & \texttt{exact=none} \\
\hyperref[enum:neg_m]{\texttt{neg0}} & \texttt{38} & \texttt{neg\_m (def none)} \\
\hyperref[enum:abs_m]{\texttt{abs0}} & \texttt{39} & \texttt{abs\_m (def none)} \\
\texttt{dest\_register\_index} & \texttt{40:45} & \texttt{uint} \\
\hyperref[enum:dest_width_narrow_m]{\texttt{dest\_width\_narrow}} & \texttt{46:47} & \texttt{dest\_width\_narrow\_m} \\
\texttt{datatype} & \texttt{48:51} & \texttt{exact=f32} \\
\texttt{instruction\_class} & \texttt{52:56} & \texttt{exact=R2} \\
\texttt{fau\_page\_index} & \texttt{57:58} & \texttt{uint} \\
\hyperref[enum:flow_control_m]{\texttt{flow}} & \texttt{59:62} & \texttt{flow\_control\_m} \\
\textit{Reserved} & \texttt{63} & \texttt{zero} \\
\bottomrule
\end{tabularx}
\end{center}
\textbf{Fields Explained:}
\begin{tightitemize}\footnotesize
  \item \texttt{src0}: Primary source operand (\texttt{r0..r63}, \texttt{u0..u63} uniform, \texttt{k*} FAU constant, or \texttt{t0..t2} temporary; \texttt{\textasciicircum{}} marks discard/last-use). \emph{Example:} {\scriptsize\ttfamily \mnem{FADD.f32} \dst{r0}, \src{r1}, \src{r2}} (\textbf{SPIR-V:} \texttt{OpFAdd (Operand 0)}; \textbf{GLSL:} \texttt{float c = a + b;})
  \item \texttt{src1}: Secondary source operand (register, uniform, or FAU constant; in shifts, supplies the shift count byte). \emph{Example:} {\scriptsize\ttfamily \mnem{LSHIFT.i32} \dst{r0}, \src{r1}, \src{r2}\imod{.b0}} (\textbf{SPIR-V:} \texttt{OpShiftLeftLogical (Shift)}; \textbf{GLSL:} \texttt{uint y = a \textless{}\textless{} s;})
  \item \texttt{opcode2}: Secondary sub-opcode further specializing instructions sharing \texttt{instruction\_class} and \texttt{opcode1} (e.g., conversions, SFU tables, moves). \emph{Example:} {\scriptsize\ttfamily \mnem{F32\_TO\_S32}\imod{.round\_zero} \dst{r0}, \src{r1}} (\textbf{SPIR-V:} \texttt{OpConvertFToS}; \textbf{GLSL:} \texttt{int y = int(x);})
  \item \hyperref[enum:widen_m]{\texttt{widen1}}: Stage A widening extractor on \texttt{src1} (promotes narrow byte/half/word lane before Stage B). \emph{Example:} {\scriptsize\ttfamily \mnem{IMUL.u64} \dst{[r0:r1]}, \src{[r2:r3]}\imod{.w0}, \src{[r4:r5]}\imod{.w0}} (\textbf{SPIR-V:} \texttt{OpUConvert / OpUMulExtended}; \textbf{GLSL:} \texttt{uint64\_t p = uint64\_t(a) * uint64\_t(b);})
  \item \hyperref[enum:widen_m]{\texttt{widen0}}: Stage A widening extractor on \texttt{src0} (promotes FP16 \texttt{h0}/\texttt{h1} to FP32, or sign/zero-extends \texttt{b0..b3}/\texttt{h0..h1}/\texttt{w0..w1} to 32/64-bit before \texttt{abs}/\texttt{neg}). \emph{Example:} {\scriptsize\ttfamily \mnem{IADD.s32} \dst{r0}, \src{r1}\imod{.h0}, \src{r2}} (\textbf{SPIR-V:} \texttt{OpSConvert / OpFConvert}; \textbf{GLSL:} \texttt{int c = int(int16\_t(a)) + b;})
  \item \hyperref[enum:round_m]{\texttt{round}}: Rounding mode (\texttt{none}/RNE, \texttt{round\_up}, \texttt{round\_down}, \texttt{round\_zero}, \texttt{round\_na}) applied in FP ALU, conversions, \texttt{HADD}, \texttt{LDEXP}, and Stage C narrowing. \emph{Example:} {\scriptsize\ttfamily \mnem{FROUND.f32}\imod{.round\_down} \dst{r0}, \src{r1}} (\textbf{SPIR-V:} \texttt{GLSL.std.450 Floor / Ceil / Trunc}; \textbf{GLSL:} \texttt{float y = floor(x);})
  \item \hyperref[enum:clamp_m]{\texttt{clamp}}: Stage B post-op floating-point clamp (\texttt{none}, \texttt{clamp\_0\_inf}, \texttt{clamp\_m1\_1}, \texttt{clamp\_0\_1}) with \texttt{-0.0 \textless{} +0.0} and NaN suppression (\texttt{NaN -\textgreater{} lo}). \emph{Example:} {\scriptsize\ttfamily \mnem{FMA.f32}\imod{.clamp\_0\_1} \dst{r0}, \src{r1}, \src{r2}, \src{r3}} (\textbf{SPIR-V:} \texttt{GLSL.std.450 FClamp}; \textbf{GLSL:} \texttt{float y = clamp(fma(a, b, c), 0.0, 1.0);})
  \item \hyperref[enum:fadd_sticky_m]{\texttt{sticky}}: OR-reduces shifted-out rounding bits into LSB (round-to-odd for double-precision emulation). \emph{Example:} {\scriptsize\ttfamily \mnem{FADD.f32}\imod{.sticky} \dst{r0}, \src{r1}, \src{r2}} (\textbf{SPIR-V:} \texttt{OpFAdd (emulated FP64)}; \textbf{GLSL:} \texttt{double c = a + b; // lowered})
  \item \hyperref[enum:neg_m]{\texttt{neg1}}: Stage A floating-point negation on \texttt{src1} (flips IEEE sign bit after swizzle/widen and \texttt{abs1}). \emph{Example:} {\scriptsize\ttfamily \mnem{FMA.f32} \dst{r0}, \src{r1}, \src{r2}\imod{.neg}, \src{r3}} (\textbf{SPIR-V:} \texttt{OpFNegate}; \textbf{GLSL:} \texttt{float d = fma(a, -b, c);})
  \item \hyperref[enum:abs_m]{\texttt{abs2}}: Stage A floating-point absolute value on \texttt{src2} (clears IEEE sign bit of FMA addend after swizzle/widen). \emph{Example:} {\scriptsize\ttfamily \mnem{FMA.f32} \dst{r0}, \src{r1}, \src{r2}, \src{r3}\imod{.abs}} (\textbf{SPIR-V:} \texttt{GLSL.std.450 FAbs}; \textbf{GLSL:} \texttt{float d = fma(a, b, abs(c));})
  \item \hyperref[enum:neg_m]{\texttt{neg0}}: Stage A floating-point negation on \texttt{src0} (flips IEEE sign bit after swizzle/widen and \texttt{abs0}). \emph{Example:} {\scriptsize\ttfamily \mnem{FADD.f32} \dst{r0}, \src{r1}\imod{.neg}, \src{r2}} (\textbf{SPIR-V:} \texttt{OpFNegate}; \textbf{GLSL:} \texttt{float c = -a + b;})
  \item \hyperref[enum:abs_m]{\texttt{abs0}}: Stage A floating-point absolute value on \texttt{src0} (clears IEEE sign bit after swizzle/widen, before \texttt{neg0}). \emph{Example:} {\scriptsize\ttfamily \mnem{FADD.f32} \dst{r0}, \src{r1}\imod{.abs}, \src{r2}} (\textbf{SPIR-V:} \texttt{GLSL.std.450 FAbs}; \textbf{GLSL:} \texttt{float c = abs(a) + b;})
  \item \texttt{dest\_register\_index}: 6-bit destination GPR index (\texttt{r0..r63}) receiving the Stage C writeback result (even-aligned pair \texttt{[rN:rN+1]} for 64-bit ops). \emph{Example:} {\scriptsize\ttfamily \mnem{FMA.f32} \dst{r0}, \src{r1}, \src{r2}, \src{r3}} (\textbf{SPIR-V:} \texttt{SSA Result ID}; \textbf{GLSL:} \texttt{float r0 = fma(r1, r2, r3);})
  \item \hyperref[enum:dest_width_narrow_m]{\texttt{dest\_width\_narrow}}: Destination halfword lane-mask / replication mode (\texttt{h0}, \texttt{h1}, \texttt{h01}, \texttt{none}). \emph{Example:} {\scriptsize\ttfamily \mnem{F32\_TO\_F16} \dst{r0}\imod{.h01}, \src{r1}} (\textbf{SPIR-V:} \texttt{OpFConvert (f32 -\textgreater{} f16)}; \textbf{GLSL:} \texttt{float16\_t h = float16\_t(x);})
  \item \texttt{datatype}: Selects the lane data type (\texttt{f32}, \texttt{f16}, \texttt{v2f16}, \texttt{s32}, \texttt{u32}, \texttt{v2s16}, \texttt{v2u16}, \texttt{v4s8}, \texttt{v4u8}, \texttt{s64}, \texttt{u64}, etc.). \emph{Example:} {\scriptsize\ttfamily \mnem{IADD.v2s16} \dst{r0}, \src{r1}, \src{r2}} (\textbf{SPIR-V:} \texttt{OpTypeInt / OpTypeFloat}; \textbf{GLSL:} \texttt{i16vec2 c = a + b;})
  \item \texttt{instruction\_class}: 4-bit major instruction class selector (\texttt{bits 48:51}) partitioning the 64-bit Valhall opcode space (e.g., \texttt{FMA}, \texttt{CVT}, \texttt{SFU}, \texttt{LS}, \texttt{TEX}, \texttt{BRANCH}). \emph{Example:} {\scriptsize\ttfamily \mnem{FADD.f32} \dst{r0}, \src{r1}, \src{r2}} (\textbf{SPIR-V:} \texttt{N/A (encoding class)}; \textbf{GLSL:} \texttt{// Hardware execution pipe / decode class})
  \item \texttt{fau\_page\_index}: Selects the FAU (Fast Access Uniform) page (\texttt{0..3}) when a source operand references hardware small constants, thread IDs, or lane constants. \emph{Example:} {\scriptsize\ttfamily \mnem{FADD.f32} \dst{r0}, \src{r1}, \src{k1f}} (\textbf{SPIR-V:} \texttt{OpConstant / BuiltIn}; \textbf{GLSL:} \texttt{float y = x + 1.0;})
  \item \hyperref[enum:flow_control_m]{\texttt{flow}}: 4-bit embedded flow-control \& scoreboard wait modifier (\texttt{bits 59:62}: \texttt{none}, \texttt{wait0..wait0126}, \texttt{wait}, \texttt{reconverge}, \texttt{discard}, \texttt{end}). \emph{Example:} {\scriptsize\ttfamily \mnem{FADD.f32}\imod{.wait0.reconverge} \dst{r0}, \src{r1}, \src{r2}} (\textbf{SPIR-V:} \texttt{OpControlBarrier / OpReturn}; \textbf{GLSL:} \texttt{// Scoreboard wait + warp reconvergence})
\end{tightitemize}
\clearpage

\subsection*{F32\_TO\_S32}
\addcontentsline{toc}{subsection}{F32\_TO\_S32}
Converts FP32 \texttt{src0} to a saturating signed 32-bit integer in \texttt{dest}: rounds per \texttt{round} (default RNE; \texttt{.round\_zero} for truncation), returning \texttt{0} on \texttt{NaN}, \texttt{INT32\_MIN} for values \texttt{\textless{}= -2\textasciicircum{}31}, \texttt{INT32\_MAX} for values \texttt{\textgreater{}= 2\textasciicircum{}31 - 1}, and \texttt{int32(x)} otherwise.
\par\textbf{Notes:}
\begin{tightitemize}
  \item Matches Rust \texttt{f32 as i32} saturation (\texttt{ops.rs:1249-1262}, \texttt{hw\_tests.rs:1462-1483}).
  \item \textbf{Source references:} \texttt{kraid/ops.rs:1226-1262}, \texttt{encode\_v9.rs:1212-1246}, \texttt{hw\_tests.rs:1462-1483}.
\end{tightitemize}
\codebox{\mnem{F32\_TO\_S32}\imod{\{.flow\}}\imod{\{.round\}} \dst{dest}\imod{\{.dest\_width\}}, \src{src0}\imod{\{.abs0\}}\imod{\{.neg0\}}\\\mnem{F32\_TO\_S32}\imod{\{.flow\}}\imod{\{.round\}} \dst{dest}\imod{\{.dest\_width\}}, \src{src0}\imod{\{.widen0\}}\imod{\{.abs0\}}\imod{\{.neg0\}}}
\textbf{Operands:}
\begin{tightitemize}\footnotesize
  \item \texttt{\dst{dest}} (\texttt{dest\_register\_index}, \texttt{dest\_width}): 32-bit destination GPR receiving the converted value.
  \item \texttt{\src{src0}}: FP32 input rounded per \texttt{round} and saturated to \texttt{[-2\textasciicircum{}31, 2\textasciicircum{}31-1]} (NaN -\textgreater{} \texttt{0}).
\end{tightitemize}
\par\textbf{Closest SPIR-V Op:} \texttt{OpConvertFToS}
\par\textbf{Pseudo-GLSL (Operational Semantics):}
\codebox{if (isnan(src0)) dest = 0;\\else             dest = int(clamp(round\_rm(src0), double(INT32\_MIN), double(INT32\_MAX)));}
\textbf{Execution Context:}
\begin{tightitemize}\footnotesize
  \item \textbf{Compute stage preloads (\hyperref[sec:csf-spd]{\texttt{SPD\_k}}, \texttt{r16..r23}):} Hardware pre-fills \texttt{r55..r57} with \texttt{WorkGroupID.xyz + JOB\_OFFSET} (from \hyperref[sec:csf-regs]{\texttt{COMPUTESR::JOB\_OFFSET\_X/Y/Z}} \texttt{r34..r36}) and \texttt{r60..r62} with \texttt{LocalInvocationID.xyz} or \texttt{GlobalInvocationID.xyz} (from \texttt{WG\_SIZE} \texttt{r33} \& \texttt{JOB\_SIZE} \texttt{r37..r39}, selected by \texttt{ShaderProgram::preload}) so compute shaders can derive thread/element indices without memory loads.
  \item \textbf{Vertex stage preloads (\hyperref[sec:csf-spd]{\texttt{SPD\_0}/\texttt{SPD\_1}}, \texttt{r16..r19}):} The hardware Index Fetcher reads \hyperref[sec:csf-regs]{\texttt{IDVSSR::INDEX\_BUFFER}} (\texttt{r54:r55}) and pre-fills \texttt{r59, r60} with \texttt{VertexIndex} (\texttt{+ VERTEX\_OFFSET + GLOBAL\_ATTR\_OFF}) and \texttt{r61} with \texttt{InstanceIndex} (\texttt{+ INSTANCE\_OFFSET}) so Position and Varying shaders can pull vertex attributes directly.
  \item \textbf{Inline FAU \& Special Operands (\hyperref[sec:csf-tsd-fau]{\texttt{FAU\_k}}, \texttt{r8..r15}):} Any 8-bit source with \texttt{mode=2} reads zero-latency uniform \texttt{u0..u31} from \texttt{FAU\_k} (\texttt{slots64 = bits[63:56]}); \texttt{mode=3} reads \texttt{SmallConstant} (Table~1.2) or \texttt{FauSpecial} (\texttt{pc}, \texttt{lane\_id}, \texttt{lane\_width}, \texttt{tls\_ptr}, \texttt{wls\_ptr}, \texttt{blend\_descriptor\_0..7}).
\end{tightitemize}
\par\textbf{Worked Examples:}
\par\textbf{Example 1 (Truncate toward zero (GLSL `int(x)`)):}
\codebox{\mnem{F32\_TO\_S32}\imod{.round\_zero} \dst{r0}, \src{r1}}
With \texttt{r1 = -3.9f}: \texttt{.round\_zero} truncates to \texttt{r0 = -3} (\texttt{0xFFFFFFFD}).
\par\textbf{Example 2 (NaN \& out-of-range saturation):}
\codebox{\mnem{F32\_TO\_S32}\imod{.round\_zero} \dst{r0}, \src{r1}}
With \texttt{r1 = NaN} -\textgreater{} \texttt{r0 = 0}; with \texttt{r1 = +Inf} (or \texttt{\textgreater{}= 2\textasciicircum{}31}) -\textgreater{} \texttt{r0 = 0x7FFFFFFF}; with \texttt{r1 = -Inf} -\textgreater{} \texttt{r0 = 0x80000000}.
\par\textbf{Mined SPIR-V $\to$ Valhall ISA Example (Truncating float-to-signed-int conversion (.round\_zero) --- Fragment shader, Mali-G710 (v10)):}
\codebox{\textbf{; SPIR-V (input)}\\\%scaled = OpFMul \%float \%v \%float\_256\\\%clamped = OpExtInst \%float \%glsl FMin \%scaled \%max\_val\\\%ival = OpConvertFToS \%int \%clamped\\[3pt]\textbf{; Valhall ISA (malioc output)}\\FMUL.f32 r3, r3\textasciicircum{}, k255f256f.h1\\FMIN.f32.nan\_propagate.clamp\_0\_inf r3, r3\textasciicircum{}, u3.w1\\F32\_TO\_S32.round\_zero r3, r3\textasciicircum{}}
SPIR-V \texttt{OpConvertFToS} requires truncation toward zero, which \texttt{malioc} encodes explicitly via \texttt{.round\_zero} (default without \texttt{.round\_zero} is round-to-nearest-even, as used in \texttt{FEXP} Q8.24 conversion).
\begin{center}\footnotesize
\begin{tabularx}{\linewidth}{@{}lllX@{}}
\toprule
Variant & Arch & Class & Opcode Fields \\
\midrule
\texttt{(none)} & v9--v14 & \texttt{R1} & \texttt{opcode2=F32\_TO\_S32, opcode1=CONVERT} \\
\texttt{(none)} & v9--v14 & \texttt{R1} & \texttt{opcode2=F32\_TO\_S32, opcode1=CONVERT} \\
\texttt{(none)} & v9--v14 & \texttt{R1} & \texttt{opcode2=F32\_TO\_I32, opcode1=CONVERT} \\
\bottomrule
\end{tabularx}
\end{center}
\begin{center}\footnotesize
\begin{tabularx}{0.88\linewidth}{@{}l l X@{}}
\toprule
Field Name & Bits & Type / Fixed Value \\
\midrule
\texttt{src0} & \texttt{0:7} & \texttt{SourceEncoding} \\
\textit{Reserved} & \texttt{8:15} & \texttt{zero} \\
\texttt{opcode2} & \texttt{16:20} & \texttt{exact=F32\_TO\_S32} \\
\textit{Reserved} & \texttt{21:29} & \texttt{zero} \\
\hyperref[enum:round_integer_m]{\texttt{round}} & \texttt{30:32} & \texttt{round\_integer\_m (def none)} \\
\textit{Reserved} & \texttt{33} & \texttt{zero} \\
\textit{Reserved} & \texttt{34:35} & \texttt{ignore} \\
\textit{Reserved} & \texttt{36:37} & \texttt{zero} \\
\hyperref[enum:neg_m]{\texttt{neg0}} & \texttt{38} & \texttt{neg\_m (def none)} \\
\hyperref[enum:abs_m]{\texttt{abs0}} & \texttt{39} & \texttt{abs\_m (def none)} \\
\texttt{dest\_register\_index} & \texttt{40:45} & \texttt{uint} \\
\hyperref[enum:dest_width_m]{\texttt{dest\_width}} & \texttt{46:47} & \texttt{dest\_width\_m} \\
\texttt{opcode1} & \texttt{48:51} & \texttt{exact=CONVERT} \\
\texttt{instruction\_class} & \texttt{52:56} & \texttt{exact=R1} \\
\texttt{fau\_page\_index} & \texttt{57:58} & \texttt{uint} \\
\hyperref[enum:flow_control_m]{\texttt{flow}} & \texttt{59:62} & \texttt{flow\_control\_m} \\
\textit{Reserved} & \texttt{63} & \texttt{zero} \\
\bottomrule
\end{tabularx}
\end{center}
\textbf{Fields Explained:}
\begin{tightitemize}\footnotesize
  \item \texttt{src0}: Primary source operand (\texttt{r0..r63}, \texttt{u0..u63} uniform, \texttt{k*} FAU constant, or \texttt{t0..t2} temporary; \texttt{\textasciicircum{}} marks discard/last-use). \emph{Example:} {\scriptsize\ttfamily \mnem{FADD.f32} \dst{r0}, \src{r1}, \src{r2}} (\textbf{SPIR-V:} \texttt{OpFAdd (Operand 0)}; \textbf{GLSL:} \texttt{float c = a + b;})
  \item \texttt{opcode2}: Secondary sub-opcode further specializing instructions sharing \texttt{instruction\_class} and \texttt{opcode1} (e.g., conversions, SFU tables, moves). \emph{Example:} {\scriptsize\ttfamily \mnem{F32\_TO\_S32}\imod{.round\_zero} \dst{r0}, \src{r1}} (\textbf{SPIR-V:} \texttt{OpConvertFToS}; \textbf{GLSL:} \texttt{int y = int(x);})
  \item \hyperref[enum:round_integer_m]{\texttt{round}}: Rounding mode (\texttt{none}/RNE, \texttt{round\_up}, \texttt{round\_down}, \texttt{round\_zero}, \texttt{round\_na}) applied in FP ALU, conversions, \texttt{HADD}, \texttt{LDEXP}, and Stage C narrowing. \emph{Example:} {\scriptsize\ttfamily \mnem{FROUND.f32}\imod{.round\_down} \dst{r0}, \src{r1}} (\textbf{SPIR-V:} \texttt{GLSL.std.450 Floor / Ceil / Trunc}; \textbf{GLSL:} \texttt{float y = floor(x);})
  \item \hyperref[enum:neg_m]{\texttt{neg0}}: Stage A floating-point negation on \texttt{src0} (flips IEEE sign bit after swizzle/widen and \texttt{abs0}). \emph{Example:} {\scriptsize\ttfamily \mnem{FADD.f32} \dst{r0}, \src{r1}\imod{.neg}, \src{r2}} (\textbf{SPIR-V:} \texttt{OpFNegate}; \textbf{GLSL:} \texttt{float c = -a + b;})
  \item \hyperref[enum:abs_m]{\texttt{abs0}}: Stage A floating-point absolute value on \texttt{src0} (clears IEEE sign bit after swizzle/widen, before \texttt{neg0}). \emph{Example:} {\scriptsize\ttfamily \mnem{FADD.f32} \dst{r0}, \src{r1}\imod{.abs}, \src{r2}} (\textbf{SPIR-V:} \texttt{GLSL.std.450 FAbs}; \textbf{GLSL:} \texttt{float c = abs(a) + b;})
  \item \texttt{dest\_register\_index}: 6-bit destination GPR index (\texttt{r0..r63}) receiving the Stage C writeback result (even-aligned pair \texttt{[rN:rN+1]} for 64-bit ops). \emph{Example:} {\scriptsize\ttfamily \mnem{FMA.f32} \dst{r0}, \src{r1}, \src{r2}, \src{r3}} (\textbf{SPIR-V:} \texttt{SSA Result ID}; \textbf{GLSL:} \texttt{float r0 = fma(r1, r2, r3);})
  \item \hyperref[enum:dest_width_m]{\texttt{dest\_width}}: Stage C destination lane write-mask (\texttt{none} = full 32-bit, \texttt{h0} = low 16-bit, \texttt{h1} = high 16-bit without bit-shifting). \emph{Example:} {\scriptsize\ttfamily \mnem{FADD.v2f16} \dst{r0}\imod{.h1}, \src{r1}, \src{r2}} (\textbf{SPIR-V:} \texttt{OpVectorInsertDynamic / WriteMask}; \textbf{GLSL:} \texttt{dst.y = a.y + b.y;})
  \item \texttt{opcode1}: Primary sub-opcode within the instruction class uniquely identifying the operation family. \emph{Example:} {\scriptsize\ttfamily \mnem{IADD.s32} \dst{r0}, \src{r1}, \src{r2}} (\textbf{SPIR-V:} \texttt{OpIAdd (operation selector)}; \textbf{GLSL:} \texttt{// Primary opcode bits})
  \item \texttt{instruction\_class}: 4-bit major instruction class selector (\texttt{bits 48:51}) partitioning the 64-bit Valhall opcode space (e.g., \texttt{FMA}, \texttt{CVT}, \texttt{SFU}, \texttt{LS}, \texttt{TEX}, \texttt{BRANCH}). \emph{Example:} {\scriptsize\ttfamily \mnem{FADD.f32} \dst{r0}, \src{r1}, \src{r2}} (\textbf{SPIR-V:} \texttt{N/A (encoding class)}; \textbf{GLSL:} \texttt{// Hardware execution pipe / decode class})
  \item \texttt{fau\_page\_index}: Selects the FAU (Fast Access Uniform) page (\texttt{0..3}) when a source operand references hardware small constants, thread IDs, or lane constants. \emph{Example:} {\scriptsize\ttfamily \mnem{FADD.f32} \dst{r0}, \src{r1}, \src{k1f}} (\textbf{SPIR-V:} \texttt{OpConstant / BuiltIn}; \textbf{GLSL:} \texttt{float y = x + 1.0;})
  \item \hyperref[enum:flow_control_m]{\texttt{flow}}: 4-bit embedded flow-control \& scoreboard wait modifier (\texttt{bits 59:62}: \texttt{none}, \texttt{wait0..wait0126}, \texttt{wait}, \texttt{reconverge}, \texttt{discard}, \texttt{end}). \emph{Example:} {\scriptsize\ttfamily \mnem{FADD.f32}\imod{.wait0.reconverge} \dst{r0}, \src{r1}, \src{r2}} (\textbf{SPIR-V:} \texttt{OpControlBarrier / OpReturn}; \textbf{GLSL:} \texttt{// Scoreboard wait + warp reconvergence})
\end{tightitemize}
\clearpage

\subsection*{F32\_TO\_U32}
\addcontentsline{toc}{subsection}{F32\_TO\_U32}
Converts FP32 \texttt{src0} to a saturating unsigned 32-bit integer in \texttt{dest}: rounds per \texttt{round} (default RNE; \texttt{.round\_zero} for truncation), returning \texttt{0} on \texttt{NaN} or values \texttt{\textless{}= 0.0}, \texttt{UINT32\_MAX} for values \texttt{\textgreater{}= 2\textasciicircum{}32 - 1}, and \texttt{uint32(x)} otherwise.
\par\textbf{Notes:}
\begin{tightitemize}
  \item Matches Rust \texttt{f32 as u32} saturation (\texttt{ops.rs:1249-1262}, \texttt{hw\_tests.rs:1462-1483}).
  \item \textbf{Source references:} \texttt{kraid/ops.rs:1226-1262}, \texttt{encode\_v9.rs:1212-1246}, \texttt{hw\_tests.rs:1462-1483}.
\end{tightitemize}
\codebox{\mnem{F32\_TO\_U32}\imod{\{.flow\}}\imod{\{.round\}} \dst{dest}\imod{\{.dest\_width\}}, \src{src0}\imod{\{.abs0\}}\imod{\{.neg0\}}\\\mnem{F32\_TO\_U32}\imod{\{.flow\}}\imod{\{.round\}} \dst{dest}\imod{\{.dest\_width\}}, \src{src0}\imod{\{.widen0\}}\imod{\{.abs0\}}\imod{\{.neg0\}}}
\textbf{Operands:}
\begin{tightitemize}\footnotesize
  \item \texttt{\dst{dest}} (\texttt{dest\_register\_index}, \texttt{dest\_width}): 32-bit destination GPR receiving the converted value.
  \item \texttt{\src{src0}}: FP32 input rounded per \texttt{round} and saturated to \texttt{[0, 2\textasciicircum{}32-1]} (NaN/\textless{}0 -\textgreater{} \texttt{0}).
\end{tightitemize}
\par\textbf{Closest SPIR-V Op:} \texttt{OpConvertFToU}
\par\textbf{Pseudo-GLSL (Operational Semantics):}
\codebox{if (isnan(src0) \textbar{}\textbar{} src0 \textless{}= 0.0) dest = 0u;\\else                            dest = uint(clamp(round\_rm(src0), 0.0, double(UINT32\_MAX)));}
\textbf{Execution Context:}
\begin{tightitemize}\footnotesize
  \item \textbf{Compute stage preloads (\hyperref[sec:csf-spd]{\texttt{SPD\_k}}, \texttt{r16..r23}):} Hardware pre-fills \texttt{r55..r57} with \texttt{WorkGroupID.xyz + JOB\_OFFSET} (from \hyperref[sec:csf-regs]{\texttt{COMPUTESR::JOB\_OFFSET\_X/Y/Z}} \texttt{r34..r36}) and \texttt{r60..r62} with \texttt{LocalInvocationID.xyz} or \texttt{GlobalInvocationID.xyz} (from \texttt{WG\_SIZE} \texttt{r33} \& \texttt{JOB\_SIZE} \texttt{r37..r39}, selected by \texttt{ShaderProgram::preload}) so compute shaders can derive thread/element indices without memory loads.
  \item \textbf{Vertex stage preloads (\hyperref[sec:csf-spd]{\texttt{SPD\_0}/\texttt{SPD\_1}}, \texttt{r16..r19}):} The hardware Index Fetcher reads \hyperref[sec:csf-regs]{\texttt{IDVSSR::INDEX\_BUFFER}} (\texttt{r54:r55}) and pre-fills \texttt{r59, r60} with \texttt{VertexIndex} (\texttt{+ VERTEX\_OFFSET + GLOBAL\_ATTR\_OFF}) and \texttt{r61} with \texttt{InstanceIndex} (\texttt{+ INSTANCE\_OFFSET}) so Position and Varying shaders can pull vertex attributes directly.
  \item \textbf{Inline FAU \& Special Operands (\hyperref[sec:csf-tsd-fau]{\texttt{FAU\_k}}, \texttt{r8..r15}):} Any 8-bit source with \texttt{mode=2} reads zero-latency uniform \texttt{u0..u31} from \texttt{FAU\_k} (\texttt{slots64 = bits[63:56]}); \texttt{mode=3} reads \texttt{SmallConstant} (Table~1.2) or \texttt{FauSpecial} (\texttt{pc}, \texttt{lane\_id}, \texttt{lane\_width}, \texttt{tls\_ptr}, \texttt{wls\_ptr}, \texttt{blend\_descriptor\_0..7}).
\end{tightitemize}
\par\textbf{Worked Examples:}
\par\textbf{Example 1 (Positive float truncation):}
\codebox{\mnem{F32\_TO\_U32}\imod{.round\_zero} \dst{r0}, \src{r1}}
With \texttt{r1 = 42.8f}: truncates to \texttt{r0 = 42} (\texttt{0x0000002A}).
\par\textbf{Example 2 (Negative, NaN, and overflow saturation):}
\codebox{\mnem{F32\_TO\_U32}\imod{.round\_zero} \dst{r0}, \src{r1}}
With \texttt{r1 = -1.5f} or \texttt{NaN} -\textgreater{} \texttt{r0 = 0}; with \texttt{r1 = 1e20f} or \texttt{+Inf} -\textgreater{} \texttt{r0 = 0xFFFFFFFF}.
\par\textbf{Mined SPIR-V $\to$ Valhall ISA Example (Clamped float-to-unsigned-int conversion --- Fragment shader, Mali-G710 (v10)):}
\codebox{\textbf{; SPIR-V (input)}\\\%pos = OpExtInst \%float \%glsl FMax \%v \%float\_0\\\%u   = OpConvertFToU \%uint \%pos\\[3pt]\textbf{; Valhall ISA (malioc output)}\\FADD.f32.clamp\_0\_inf r2, r2\textasciicircum{}, k0.neg\\F32\_TO\_U32 r2, r2\textasciicircum{}}
\texttt{F32\_TO\_U32} saturates negative/NaN inputs to \texttt{0} and values $\ge 2^{32}-1$ to \texttt{0xFFFFFFFF}.
\begin{center}\footnotesize
\begin{tabularx}{\linewidth}{@{}lllX@{}}
\toprule
Variant & Arch & Class & Opcode Fields \\
\midrule
\texttt{(none)} & v9--v14 & \texttt{R1} & \texttt{opcode2=F32\_TO\_U32, opcode1=CONVERT} \\
\texttt{(none)} & v9--v14 & \texttt{R1} & \texttt{opcode2=F32\_TO\_U32, opcode1=CONVERT} \\
\texttt{(none)} & v9--v14 & \texttt{R1} & \texttt{opcode2=F32\_TO\_I32, opcode1=CONVERT} \\
\bottomrule
\end{tabularx}
\end{center}
\begin{center}\footnotesize
\begin{tabularx}{0.88\linewidth}{@{}l l X@{}}
\toprule
Field Name & Bits & Type / Fixed Value \\
\midrule
\texttt{src0} & \texttt{0:7} & \texttt{SourceEncoding} \\
\textit{Reserved} & \texttt{8:15} & \texttt{zero} \\
\texttt{opcode2} & \texttt{16:20} & \texttt{exact=F32\_TO\_U32} \\
\textit{Reserved} & \texttt{21:29} & \texttt{zero} \\
\hyperref[enum:round_integer_m]{\texttt{round}} & \texttt{30:32} & \texttt{round\_integer\_m (def none)} \\
\textit{Reserved} & \texttt{33} & \texttt{zero} \\
\textit{Reserved} & \texttt{34:35} & \texttt{ignore} \\
\textit{Reserved} & \texttt{36:37} & \texttt{zero} \\
\hyperref[enum:neg_m]{\texttt{neg0}} & \texttt{38} & \texttt{neg\_m (def none)} \\
\hyperref[enum:abs_m]{\texttt{abs0}} & \texttt{39} & \texttt{abs\_m (def none)} \\
\texttt{dest\_register\_index} & \texttt{40:45} & \texttt{uint} \\
\hyperref[enum:dest_width_m]{\texttt{dest\_width}} & \texttt{46:47} & \texttt{dest\_width\_m} \\
\texttt{opcode1} & \texttt{48:51} & \texttt{exact=CONVERT} \\
\texttt{instruction\_class} & \texttt{52:56} & \texttt{exact=R1} \\
\texttt{fau\_page\_index} & \texttt{57:58} & \texttt{uint} \\
\hyperref[enum:flow_control_m]{\texttt{flow}} & \texttt{59:62} & \texttt{flow\_control\_m} \\
\textit{Reserved} & \texttt{63} & \texttt{zero} \\
\bottomrule
\end{tabularx}
\end{center}
\textbf{Fields Explained:}
\begin{tightitemize}\footnotesize
  \item \texttt{src0}: Primary source operand (\texttt{r0..r63}, \texttt{u0..u63} uniform, \texttt{k*} FAU constant, or \texttt{t0..t2} temporary; \texttt{\textasciicircum{}} marks discard/last-use). \emph{Example:} {\scriptsize\ttfamily \mnem{FADD.f32} \dst{r0}, \src{r1}, \src{r2}} (\textbf{SPIR-V:} \texttt{OpFAdd (Operand 0)}; \textbf{GLSL:} \texttt{float c = a + b;})
  \item \texttt{opcode2}: Secondary sub-opcode further specializing instructions sharing \texttt{instruction\_class} and \texttt{opcode1} (e.g., conversions, SFU tables, moves). \emph{Example:} {\scriptsize\ttfamily \mnem{F32\_TO\_S32}\imod{.round\_zero} \dst{r0}, \src{r1}} (\textbf{SPIR-V:} \texttt{OpConvertFToS}; \textbf{GLSL:} \texttt{int y = int(x);})
  \item \hyperref[enum:round_integer_m]{\texttt{round}}: Rounding mode (\texttt{none}/RNE, \texttt{round\_up}, \texttt{round\_down}, \texttt{round\_zero}, \texttt{round\_na}) applied in FP ALU, conversions, \texttt{HADD}, \texttt{LDEXP}, and Stage C narrowing. \emph{Example:} {\scriptsize\ttfamily \mnem{FROUND.f32}\imod{.round\_down} \dst{r0}, \src{r1}} (\textbf{SPIR-V:} \texttt{GLSL.std.450 Floor / Ceil / Trunc}; \textbf{GLSL:} \texttt{float y = floor(x);})
  \item \hyperref[enum:neg_m]{\texttt{neg0}}: Stage A floating-point negation on \texttt{src0} (flips IEEE sign bit after swizzle/widen and \texttt{abs0}). \emph{Example:} {\scriptsize\ttfamily \mnem{FADD.f32} \dst{r0}, \src{r1}\imod{.neg}, \src{r2}} (\textbf{SPIR-V:} \texttt{OpFNegate}; \textbf{GLSL:} \texttt{float c = -a + b;})
  \item \hyperref[enum:abs_m]{\texttt{abs0}}: Stage A floating-point absolute value on \texttt{src0} (clears IEEE sign bit after swizzle/widen, before \texttt{neg0}). \emph{Example:} {\scriptsize\ttfamily \mnem{FADD.f32} \dst{r0}, \src{r1}\imod{.abs}, \src{r2}} (\textbf{SPIR-V:} \texttt{GLSL.std.450 FAbs}; \textbf{GLSL:} \texttt{float c = abs(a) + b;})
  \item \texttt{dest\_register\_index}: 6-bit destination GPR index (\texttt{r0..r63}) receiving the Stage C writeback result (even-aligned pair \texttt{[rN:rN+1]} for 64-bit ops). \emph{Example:} {\scriptsize\ttfamily \mnem{FMA.f32} \dst{r0}, \src{r1}, \src{r2}, \src{r3}} (\textbf{SPIR-V:} \texttt{SSA Result ID}; \textbf{GLSL:} \texttt{float r0 = fma(r1, r2, r3);})
  \item \hyperref[enum:dest_width_m]{\texttt{dest\_width}}: Stage C destination lane write-mask (\texttt{none} = full 32-bit, \texttt{h0} = low 16-bit, \texttt{h1} = high 16-bit without bit-shifting). \emph{Example:} {\scriptsize\ttfamily \mnem{FADD.v2f16} \dst{r0}\imod{.h1}, \src{r1}, \src{r2}} (\textbf{SPIR-V:} \texttt{OpVectorInsertDynamic / WriteMask}; \textbf{GLSL:} \texttt{dst.y = a.y + b.y;})
  \item \texttt{opcode1}: Primary sub-opcode within the instruction class uniquely identifying the operation family. \emph{Example:} {\scriptsize\ttfamily \mnem{IADD.s32} \dst{r0}, \src{r1}, \src{r2}} (\textbf{SPIR-V:} \texttt{OpIAdd (operation selector)}; \textbf{GLSL:} \texttt{// Primary opcode bits})
  \item \texttt{instruction\_class}: 4-bit major instruction class selector (\texttt{bits 48:51}) partitioning the 64-bit Valhall opcode space (e.g., \texttt{FMA}, \texttt{CVT}, \texttt{SFU}, \texttt{LS}, \texttt{TEX}, \texttt{BRANCH}). \emph{Example:} {\scriptsize\ttfamily \mnem{FADD.f32} \dst{r0}, \src{r1}, \src{r2}} (\textbf{SPIR-V:} \texttt{N/A (encoding class)}; \textbf{GLSL:} \texttt{// Hardware execution pipe / decode class})
  \item \texttt{fau\_page\_index}: Selects the FAU (Fast Access Uniform) page (\texttt{0..3}) when a source operand references hardware small constants, thread IDs, or lane constants. \emph{Example:} {\scriptsize\ttfamily \mnem{FADD.f32} \dst{r0}, \src{r1}, \src{k1f}} (\textbf{SPIR-V:} \texttt{OpConstant / BuiltIn}; \textbf{GLSL:} \texttt{float y = x + 1.0;})
  \item \hyperref[enum:flow_control_m]{\texttt{flow}}: 4-bit embedded flow-control \& scoreboard wait modifier (\texttt{bits 59:62}: \texttt{none}, \texttt{wait0..wait0126}, \texttt{wait}, \texttt{reconverge}, \texttt{discard}, \texttt{end}). \emph{Example:} {\scriptsize\ttfamily \mnem{FADD.f32}\imod{.wait0.reconverge} \dst{r0}, \src{r1}, \src{r2}} (\textbf{SPIR-V:} \texttt{OpControlBarrier / OpReturn}; \textbf{GLSL:} \texttt{// Scoreboard wait + warp reconvergence})
\end{tightitemize}
\clearpage

\subsection*{S16\_TO\_F32}
\addcontentsline{toc}{subsection}{S16\_TO\_F32}
Converts the signed 16-bit integer selected by \texttt{lane0} (\texttt{.h0} or \texttt{.h1}) of \texttt{src0} exactly to IEEE-754 FP32 and writes the result to \texttt{dest} (arch \textless{}= v10).
\par\textbf{Notes:}
\begin{tightitemize}
  \item Arch v9-v10 only (\texttt{kraid/nir.rs:740}).
  \item On v11+ use \texttt{S32\_TO\_F32} with \texttt{widen0=h0}.
  \item \textbf{Source references:} \texttt{kraid/ops.rs:2575-2600}, \texttt{encode\_v9.rs:1992-1995}, \texttt{nir.rs:738-749}.
\end{tightitemize}
\codebox{\mnem{S16\_TO\_F32}\imod{\{.flow\}} \dst{dest}\imod{\{.dest\_width\}}, \src{src0}\imod{\{.lane0\}}}
\textbf{Operands:}
\begin{tightitemize}\footnotesize
  \item \texttt{\dst{dest}} (\texttt{dest\_register\_index}, \texttt{dest\_width}): 32-bit destination GPR receiving the converted value.
  \item \texttt{\src{src0}}: Signed 16-bit lane (\texttt{lane0}) converted exactly to FP32.
\end{tightitemize}
\par\textbf{Closest SPIR-V Op:} \texttt{OpConvertSToF (s16 -\textgreater{} f32, v9--v10)}
\par\textbf{Pseudo-GLSL (Operational Semantics):}
\codebox{dest = float(bitfieldExtract(int(src0), lane0 * 16, 16));}
\textbf{Execution Context:}
\begin{tightitemize}\footnotesize
  \item \textbf{Compute stage preloads (\hyperref[sec:csf-spd]{\texttt{SPD\_k}}, \texttt{r16..r23}):} Hardware pre-fills \texttt{r55..r57} with \texttt{WorkGroupID.xyz + JOB\_OFFSET} (from \hyperref[sec:csf-regs]{\texttt{COMPUTESR::JOB\_OFFSET\_X/Y/Z}} \texttt{r34..r36}) and \texttt{r60..r62} with \texttt{LocalInvocationID.xyz} or \texttt{GlobalInvocationID.xyz} (from \texttt{WG\_SIZE} \texttt{r33} \& \texttt{JOB\_SIZE} \texttt{r37..r39}, selected by \texttt{ShaderProgram::preload}) so compute shaders can derive thread/element indices without memory loads.
  \item \textbf{Vertex stage preloads (\hyperref[sec:csf-spd]{\texttt{SPD\_0}/\texttt{SPD\_1}}, \texttt{r16..r19}):} The hardware Index Fetcher reads \hyperref[sec:csf-regs]{\texttt{IDVSSR::INDEX\_BUFFER}} (\texttt{r54:r55}) and pre-fills \texttt{r59, r60} with \texttt{VertexIndex} (\texttt{+ VERTEX\_OFFSET + GLOBAL\_ATTR\_OFF}) and \texttt{r61} with \texttt{InstanceIndex} (\texttt{+ INSTANCE\_OFFSET}) so Position and Varying shaders can pull vertex attributes directly.
  \item \textbf{Inline FAU \& Special Operands (\hyperref[sec:csf-tsd-fau]{\texttt{FAU\_k}}, \texttt{r8..r15}):} Any 8-bit source with \texttt{mode=2} reads zero-latency uniform \texttt{u0..u31} from \texttt{FAU\_k} (\texttt{slots64 = bits[63:56]}); \texttt{mode=3} reads \texttt{SmallConstant} (Table~1.2) or \texttt{FauSpecial} (\texttt{pc}, \texttt{lane\_id}, \texttt{lane\_width}, \texttt{tls\_ptr}, \texttt{wls\_ptr}, \texttt{blend\_descriptor\_0..7}).
\end{tightitemize}
\par\textbf{Worked Examples:}
\par\textbf{Example 1 (Negative 16-bit integer (v9--v10)):}
\codebox{\mnem{S16\_TO\_F32} \dst{r0}, \src{r1}\imod{.h0}}
With \texttt{r1.h0 = 0xFFFE} (\texttt{-2}): exact conversion -\textgreater{} \texttt{r0 = -2.0f} (\texttt{0xC0000000}).
\par\textbf{Example 2 (Upper halfword `INT16\_MIN`):}
\codebox{\mnem{S16\_TO\_F32} \dst{r0}, \src{r1}\imod{.h1}}
With \texttt{r1 = 0x80000000} (\texttt{h1 = -32768}): exact conversion -\textgreater{} \texttt{r0 = -32768.0f} (\texttt{0xC7000000}).
\par\textbf{Mined SPIR-V $\to$ Valhall ISA Example (Converting fixed-point 8.8 LOD result from TEX\_GRADIENT to float --- Fragment shader, Mali-G710 (v10)):}
\codebox{\textbf{; SPIR-V (input)}\\\%lod = OpImageQueryLod \%v2float \%sampler \%uv\\[3pt]\textbf{; Valhall ISA (malioc output)}\\TEX\_GRADIENT.slot0.discard.skip.tex2d.dst32.r.wide.disable\_lod\_clamp @r2, @r0:r1, u2\\S16\_TO\_F32 r2, r2\textasciicircum{}.h0\\FMUL.f32 r2, r2\textasciicircum{}, u3.w0           ; multiply by 1/256}
\texttt{TEX\_GRADIENT.disable\_lod\_clamp} returns the raw LOD as a signed 8.8 fixed-point 16-bit integer in \texttt{r2.h0}; \texttt{S16\_TO\_F32} converts \texttt{r2\textasciicircum{}.h0} exactly to FP32 before scaling by $1/256$.
\begin{center}\footnotesize
\begin{tabularx}{\linewidth}{@{}lllX@{}}
\toprule
Variant & Arch & Class & Opcode Fields \\
\midrule
\texttt{(none)} & v9--v14 & \texttt{R1} & \texttt{opcode2=S16\_TO\_F32, opcode1=CONVERT} \\
\bottomrule
\end{tabularx}
\end{center}
\begin{center}\footnotesize
\begin{tabularx}{0.88\linewidth}{@{}l l X@{}}
\toprule
Field Name & Bits & Type / Fixed Value \\
\midrule
\texttt{src0} & \texttt{0:7} & \texttt{SourceEncoding} \\
\textit{Reserved} & \texttt{8:15} & \texttt{zero} \\
\texttt{opcode2} & \texttt{16:20} & \texttt{exact=S16\_TO\_F32} \\
\textit{Reserved} & \texttt{21:27} & \texttt{zero} \\
\hyperref[enum:lane16_m]{\texttt{lane0}} & \texttt{28} & \texttt{lane16\_m (def h0)} \\
\textit{Reserved} & \texttt{29:39} & \texttt{zero} \\
\texttt{dest\_register\_index} & \texttt{40:45} & \texttt{uint} \\
\hyperref[enum:dest_width_m]{\texttt{dest\_width}} & \texttt{46:47} & \texttt{dest\_width\_m} \\
\texttt{opcode1} & \texttt{48:51} & \texttt{exact=CONVERT} \\
\texttt{instruction\_class} & \texttt{52:56} & \texttt{exact=R1} \\
\texttt{fau\_page\_index} & \texttt{57:58} & \texttt{uint} \\
\hyperref[enum:flow_control_m]{\texttt{flow}} & \texttt{59:62} & \texttt{flow\_control\_m} \\
\textit{Reserved} & \texttt{63} & \texttt{zero} \\
\bottomrule
\end{tabularx}
\end{center}
\textbf{Fields Explained:}
\begin{tightitemize}\footnotesize
  \item \texttt{src0}: Primary source operand (\texttt{r0..r63}, \texttt{u0..u63} uniform, \texttt{k*} FAU constant, or \texttt{t0..t2} temporary; \texttt{\textasciicircum{}} marks discard/last-use). \emph{Example:} {\scriptsize\ttfamily \mnem{FADD.f32} \dst{r0}, \src{r1}, \src{r2}} (\textbf{SPIR-V:} \texttt{OpFAdd (Operand 0)}; \textbf{GLSL:} \texttt{float c = a + b;})
  \item \texttt{opcode2}: Secondary sub-opcode further specializing instructions sharing \texttt{instruction\_class} and \texttt{opcode1} (e.g., conversions, SFU tables, moves). \emph{Example:} {\scriptsize\ttfamily \mnem{F32\_TO\_S32}\imod{.round\_zero} \dst{r0}, \src{r1}} (\textbf{SPIR-V:} \texttt{OpConvertFToS}; \textbf{GLSL:} \texttt{int y = int(x);})
  \item \hyperref[enum:lane16_m]{\texttt{lane0}}: Stage A single-lane extractor on \texttt{src0} (\texttt{b0..b3} for 8-bit or \texttt{h0}/\texttt{h1} for 16-bit narrow conversions and \texttt{MKVEC}). \emph{Example:} {\scriptsize\ttfamily \mnem{F16\_TO\_F32} \dst{r0}, \src{r1}\imod{.h1}} (\textbf{SPIR-V:} \texttt{OpCompositeExtract}; \textbf{GLSL:} \texttt{float y = float(x.y);})
  \item \texttt{dest\_register\_index}: 6-bit destination GPR index (\texttt{r0..r63}) receiving the Stage C writeback result (even-aligned pair \texttt{[rN:rN+1]} for 64-bit ops). \emph{Example:} {\scriptsize\ttfamily \mnem{FMA.f32} \dst{r0}, \src{r1}, \src{r2}, \src{r3}} (\textbf{SPIR-V:} \texttt{SSA Result ID}; \textbf{GLSL:} \texttt{float r0 = fma(r1, r2, r3);})
  \item \hyperref[enum:dest_width_m]{\texttt{dest\_width}}: Stage C destination lane write-mask (\texttt{none} = full 32-bit, \texttt{h0} = low 16-bit, \texttt{h1} = high 16-bit without bit-shifting). \emph{Example:} {\scriptsize\ttfamily \mnem{FADD.v2f16} \dst{r0}\imod{.h1}, \src{r1}, \src{r2}} (\textbf{SPIR-V:} \texttt{OpVectorInsertDynamic / WriteMask}; \textbf{GLSL:} \texttt{dst.y = a.y + b.y;})
  \item \texttt{opcode1}: Primary sub-opcode within the instruction class uniquely identifying the operation family. \emph{Example:} {\scriptsize\ttfamily \mnem{IADD.s32} \dst{r0}, \src{r1}, \src{r2}} (\textbf{SPIR-V:} \texttt{OpIAdd (operation selector)}; \textbf{GLSL:} \texttt{// Primary opcode bits})
  \item \texttt{instruction\_class}: 4-bit major instruction class selector (\texttt{bits 48:51}) partitioning the 64-bit Valhall opcode space (e.g., \texttt{FMA}, \texttt{CVT}, \texttt{SFU}, \texttt{LS}, \texttt{TEX}, \texttt{BRANCH}). \emph{Example:} {\scriptsize\ttfamily \mnem{FADD.f32} \dst{r0}, \src{r1}, \src{r2}} (\textbf{SPIR-V:} \texttt{N/A (encoding class)}; \textbf{GLSL:} \texttt{// Hardware execution pipe / decode class})
  \item \texttt{fau\_page\_index}: Selects the FAU (Fast Access Uniform) page (\texttt{0..3}) when a source operand references hardware small constants, thread IDs, or lane constants. \emph{Example:} {\scriptsize\ttfamily \mnem{FADD.f32} \dst{r0}, \src{r1}, \src{k1f}} (\textbf{SPIR-V:} \texttt{OpConstant / BuiltIn}; \textbf{GLSL:} \texttt{float y = x + 1.0;})
  \item \hyperref[enum:flow_control_m]{\texttt{flow}}: 4-bit embedded flow-control \& scoreboard wait modifier (\texttt{bits 59:62}: \texttt{none}, \texttt{wait0..wait0126}, \texttt{wait}, \texttt{reconverge}, \texttt{discard}, \texttt{end}). \emph{Example:} {\scriptsize\ttfamily \mnem{FADD.f32}\imod{.wait0.reconverge} \dst{r0}, \src{r1}, \src{r2}} (\textbf{SPIR-V:} \texttt{OpControlBarrier / OpReturn}; \textbf{GLSL:} \texttt{// Scoreboard wait + warp reconvergence})
\end{tightitemize}
\clearpage

\subsection*{S16\_TO\_S32}
\addcontentsline{toc}{subsection}{S16\_TO\_S32}
Extracts the signed 16-bit halfword selected by \texttt{lane0} (\texttt{.h0} or \texttt{.h1}) of \texttt{src0}, sign-extends it to 32 bits, and writes the result to \texttt{dest}.
\par\textbf{Notes:}
\begin{tightitemize}
  \item Exact sign extension.
  \item \textbf{Source references:} \texttt{isa-v9-v14.xml}, \texttt{bifrost/valhall/ISA.xml}.
\end{tightitemize}
\codebox{\mnem{S16\_TO\_S32}\imod{\{.flow\}} \dst{dest}\imod{\{.dest\_width\}}, \src{src0}\imod{\{.lane0\}}}
\textbf{Operands:}
\begin{tightitemize}\footnotesize
  \item \texttt{\dst{dest}} (\texttt{dest\_register\_index}, \texttt{dest\_width}): 32-bit destination GPR receiving the converted value.
  \item \texttt{\src{src0}}: Signed 16-bit lane (\texttt{lane0}) sign-extended to 32-bit integer.
\end{tightitemize}
\par\textbf{Closest SPIR-V Op:} \texttt{OpSConvert (s16 -\textgreater{} s32)}
\par\textbf{Pseudo-GLSL (Operational Semantics):}
\codebox{dest = bitfieldExtract(int(src0), lane0 * 16, 16);}
\textbf{Execution Context:}
\begin{tightitemize}\footnotesize
  \item \textbf{Compute stage preloads (\hyperref[sec:csf-spd]{\texttt{SPD\_k}}, \texttt{r16..r23}):} Hardware pre-fills \texttt{r55..r57} with \texttt{WorkGroupID.xyz + JOB\_OFFSET} (from \hyperref[sec:csf-regs]{\texttt{COMPUTESR::JOB\_OFFSET\_X/Y/Z}} \texttt{r34..r36}) and \texttt{r60..r62} with \texttt{LocalInvocationID.xyz} or \texttt{GlobalInvocationID.xyz} (from \texttt{WG\_SIZE} \texttt{r33} \& \texttt{JOB\_SIZE} \texttt{r37..r39}, selected by \texttt{ShaderProgram::preload}) so compute shaders can derive thread/element indices without memory loads.
  \item \textbf{Vertex stage preloads (\hyperref[sec:csf-spd]{\texttt{SPD\_0}/\texttt{SPD\_1}}, \texttt{r16..r19}):} The hardware Index Fetcher reads \hyperref[sec:csf-regs]{\texttt{IDVSSR::INDEX\_BUFFER}} (\texttt{r54:r55}) and pre-fills \texttt{r59, r60} with \texttt{VertexIndex} (\texttt{+ VERTEX\_OFFSET + GLOBAL\_ATTR\_OFF}) and \texttt{r61} with \texttt{InstanceIndex} (\texttt{+ INSTANCE\_OFFSET}) so Position and Varying shaders can pull vertex attributes directly.
  \item \textbf{Inline FAU \& Special Operands (\hyperref[sec:csf-tsd-fau]{\texttt{FAU\_k}}, \texttt{r8..r15}):} Any 8-bit source with \texttt{mode=2} reads zero-latency uniform \texttt{u0..u31} from \texttt{FAU\_k} (\texttt{slots64 = bits[63:56]}); \texttt{mode=3} reads \texttt{SmallConstant} (Table~1.2) or \texttt{FauSpecial} (\texttt{pc}, \texttt{lane\_id}, \texttt{lane\_width}, \texttt{tls\_ptr}, \texttt{wls\_ptr}, \texttt{blend\_descriptor\_0..7}).
\end{tightitemize}
\par\textbf{Worked Examples:}
\par\textbf{Example 1 (Sign-extend negative upper halfword):}
\codebox{\mnem{S16\_TO\_S32} \dst{r0}, \src{r1}\imod{.h1}}
With \texttt{r1 = 0x80010000}: sign-extends \texttt{0x8001} (\texttt{-32767}) -\textgreater{} \texttt{r0 = 0xFFFF8001}.
\par\textbf{Example 2 (Positive low halfword):}
\codebox{\mnem{S16\_TO\_S32} \dst{r0}, \src{r1}\imod{.h0}}
With \texttt{r1 = 0xFFFF7FFF}: sign-extends \texttt{h0 = 0x7FFF} (\texttt{+32767}) -\textgreater{} \texttt{r0 = 0x00007FFF}.
\par\textbf{Mined SPIR-V $\to$ Valhall ISA Example (Sign-extending a 16-bit uniform halfword for discard test --- Fragment shader, Mali-G710 (v10)):}
\codebox{\textbf{; SPIR-V (input)}\\\%flag16 = OpBitFieldSExtract \%int \%ubo\_word \%int\_0 \%int\_16\\OpBranchConditional (OpINotEqual \%flag16 \%int\_0) \%kill \%ok\\[3pt]\textbf{; Valhall ISA (malioc output)}\\S16\_TO\_S32.discard r0, u0.w0.h0\\DISCARD.f32.wait0126.ne r0\textasciicircum{}, k0}
\texttt{S16\_TO\_S32} reads the low 16 bits of uniform \texttt{u0.w0} (\texttt{.h0}) and sign-extends them to 32 bits in \texttt{r0}.
\begin{center}\footnotesize
\begin{tabularx}{\linewidth}{@{}lllX@{}}
\toprule
Variant & Arch & Class & Opcode Fields \\
\midrule
\texttt{(none)} & v9--v14 & \texttt{R1} & \texttt{opcode2=S16\_TO\_S32, opcode1=CONVERT} \\
\texttt{(none)} & v9--v14 & \texttt{R1} & \texttt{opcode2=I16\_TO\_I32, opcode1=CONVERT} \\
\bottomrule
\end{tabularx}
\end{center}
\begin{center}\footnotesize
\begin{tabularx}{0.88\linewidth}{@{}l l X@{}}
\toprule
Field Name & Bits & Type / Fixed Value \\
\midrule
\texttt{src0} & \texttt{0:7} & \texttt{SourceEncoding} \\
\textit{Reserved} & \texttt{8:15} & \texttt{zero} \\
\texttt{opcode2} & \texttt{16:20} & \texttt{exact=S16\_TO\_S32} \\
\textit{Reserved} & \texttt{21:27} & \texttt{zero} \\
\hyperref[enum:lane16_m]{\texttt{lane0}} & \texttt{28} & \texttt{lane16\_m (def h0)} \\
\textit{Reserved} & \texttt{29:39} & \texttt{zero} \\
\texttt{dest\_register\_index} & \texttt{40:45} & \texttt{uint} \\
\hyperref[enum:dest_width_m]{\texttt{dest\_width}} & \texttt{46:47} & \texttt{dest\_width\_m} \\
\texttt{opcode1} & \texttt{48:51} & \texttt{exact=CONVERT} \\
\texttt{instruction\_class} & \texttt{52:56} & \texttt{exact=R1} \\
\texttt{fau\_page\_index} & \texttt{57:58} & \texttt{uint} \\
\hyperref[enum:flow_control_m]{\texttt{flow}} & \texttt{59:62} & \texttt{flow\_control\_m} \\
\textit{Reserved} & \texttt{63} & \texttt{zero} \\
\bottomrule
\end{tabularx}
\end{center}
\textbf{Fields Explained:}
\begin{tightitemize}\footnotesize
  \item \texttt{src0}: Primary source operand (\texttt{r0..r63}, \texttt{u0..u63} uniform, \texttt{k*} FAU constant, or \texttt{t0..t2} temporary; \texttt{\textasciicircum{}} marks discard/last-use). \emph{Example:} {\scriptsize\ttfamily \mnem{FADD.f32} \dst{r0}, \src{r1}, \src{r2}} (\textbf{SPIR-V:} \texttt{OpFAdd (Operand 0)}; \textbf{GLSL:} \texttt{float c = a + b;})
  \item \texttt{opcode2}: Secondary sub-opcode further specializing instructions sharing \texttt{instruction\_class} and \texttt{opcode1} (e.g., conversions, SFU tables, moves). \emph{Example:} {\scriptsize\ttfamily \mnem{F32\_TO\_S32}\imod{.round\_zero} \dst{r0}, \src{r1}} (\textbf{SPIR-V:} \texttt{OpConvertFToS}; \textbf{GLSL:} \texttt{int y = int(x);})
  \item \hyperref[enum:lane16_m]{\texttt{lane0}}: Stage A single-lane extractor on \texttt{src0} (\texttt{b0..b3} for 8-bit or \texttt{h0}/\texttt{h1} for 16-bit narrow conversions and \texttt{MKVEC}). \emph{Example:} {\scriptsize\ttfamily \mnem{F16\_TO\_F32} \dst{r0}, \src{r1}\imod{.h1}} (\textbf{SPIR-V:} \texttt{OpCompositeExtract}; \textbf{GLSL:} \texttt{float y = float(x.y);})
  \item \texttt{dest\_register\_index}: 6-bit destination GPR index (\texttt{r0..r63}) receiving the Stage C writeback result (even-aligned pair \texttt{[rN:rN+1]} for 64-bit ops). \emph{Example:} {\scriptsize\ttfamily \mnem{FMA.f32} \dst{r0}, \src{r1}, \src{r2}, \src{r3}} (\textbf{SPIR-V:} \texttt{SSA Result ID}; \textbf{GLSL:} \texttt{float r0 = fma(r1, r2, r3);})
  \item \hyperref[enum:dest_width_m]{\texttt{dest\_width}}: Stage C destination lane write-mask (\texttt{none} = full 32-bit, \texttt{h0} = low 16-bit, \texttt{h1} = high 16-bit without bit-shifting). \emph{Example:} {\scriptsize\ttfamily \mnem{FADD.v2f16} \dst{r0}\imod{.h1}, \src{r1}, \src{r2}} (\textbf{SPIR-V:} \texttt{OpVectorInsertDynamic / WriteMask}; \textbf{GLSL:} \texttt{dst.y = a.y + b.y;})
  \item \texttt{opcode1}: Primary sub-opcode within the instruction class uniquely identifying the operation family. \emph{Example:} {\scriptsize\ttfamily \mnem{IADD.s32} \dst{r0}, \src{r1}, \src{r2}} (\textbf{SPIR-V:} \texttt{OpIAdd (operation selector)}; \textbf{GLSL:} \texttt{// Primary opcode bits})
  \item \texttt{instruction\_class}: 4-bit major instruction class selector (\texttt{bits 48:51}) partitioning the 64-bit Valhall opcode space (e.g., \texttt{FMA}, \texttt{CVT}, \texttt{SFU}, \texttt{LS}, \texttt{TEX}, \texttt{BRANCH}). \emph{Example:} {\scriptsize\ttfamily \mnem{FADD.f32} \dst{r0}, \src{r1}, \src{r2}} (\textbf{SPIR-V:} \texttt{N/A (encoding class)}; \textbf{GLSL:} \texttt{// Hardware execution pipe / decode class})
  \item \texttt{fau\_page\_index}: Selects the FAU (Fast Access Uniform) page (\texttt{0..3}) when a source operand references hardware small constants, thread IDs, or lane constants. \emph{Example:} {\scriptsize\ttfamily \mnem{FADD.f32} \dst{r0}, \src{r1}, \src{k1f}} (\textbf{SPIR-V:} \texttt{OpConstant / BuiltIn}; \textbf{GLSL:} \texttt{float y = x + 1.0;})
  \item \hyperref[enum:flow_control_m]{\texttt{flow}}: 4-bit embedded flow-control \& scoreboard wait modifier (\texttt{bits 59:62}: \texttt{none}, \texttt{wait0..wait0126}, \texttt{wait}, \texttt{reconverge}, \texttt{discard}, \texttt{end}). \emph{Example:} {\scriptsize\ttfamily \mnem{FADD.f32}\imod{.wait0.reconverge} \dst{r0}, \src{r1}, \src{r2}} (\textbf{SPIR-V:} \texttt{OpControlBarrier / OpReturn}; \textbf{GLSL:} \texttt{// Scoreboard wait + warp reconvergence})
\end{tightitemize}
\clearpage

\subsection*{S32\_TO\_F32}
\addcontentsline{toc}{subsection}{S32\_TO\_F32}
Converts signed 32-bit integer \texttt{src0} (with optional Stage A widening) to IEEE-754 FP32 in \texttt{dest}, applying the selected \texttt{round} mode (default RNE; no clamp modifier).
\par\textbf{Notes:}
\begin{tightitemize}
  \item Supports all 5 rounding modes (\texttt{round\_extended\_m}).
  \item \textbf{Source references:} \texttt{kraid/ops.rs:2575-2600}, \texttt{encode\_v9.rs:1960-2010}, \texttt{nir.rs:738-749}.
\end{tightitemize}
\codebox{\mnem{S32\_TO\_F32}\imod{\{.flow\}}\imod{\{.round\}} \dst{dest}\imod{\{.dest\_width\}}, \src{src0}\\\mnem{S32\_TO\_F32}\imod{\{.flow\}}\imod{\{.round\}} \dst{dest}\imod{\{.dest\_width\}}, \src{src0}\imod{\{.widen0\}}}
\textbf{Operands:}
\begin{tightitemize}\footnotesize
  \item \texttt{\dst{dest}} (\texttt{dest\_register\_index}, \texttt{dest\_width}): 32-bit destination GPR receiving the converted value.
  \item \texttt{\src{src0}}: Signed 32-bit integer converted to FP32 (rounded per \texttt{round}).
\end{tightitemize}
\par\textbf{Closest SPIR-V Op:} \texttt{OpConvertSToF}
\par\textbf{Pseudo-GLSL (Operational Semantics):}
\codebox{dest = round\_rm(float(int(widen0(src0))));}
\textbf{Execution Context:}
\begin{tightitemize}\footnotesize
  \item \textbf{Compute stage preloads (\hyperref[sec:csf-spd]{\texttt{SPD\_k}}, \texttt{r16..r23}):} Hardware pre-fills \texttt{r55..r57} with \texttt{WorkGroupID.xyz + JOB\_OFFSET} (from \hyperref[sec:csf-regs]{\texttt{COMPUTESR::JOB\_OFFSET\_X/Y/Z}} \texttt{r34..r36}) and \texttt{r60..r62} with \texttt{LocalInvocationID.xyz} or \texttt{GlobalInvocationID.xyz} (from \texttt{WG\_SIZE} \texttt{r33} \& \texttt{JOB\_SIZE} \texttt{r37..r39}, selected by \texttt{ShaderProgram::preload}) so compute shaders can derive thread/element indices without memory loads.
  \item \textbf{Vertex stage preloads (\hyperref[sec:csf-spd]{\texttt{SPD\_0}/\texttt{SPD\_1}}, \texttt{r16..r19}):} The hardware Index Fetcher reads \hyperref[sec:csf-regs]{\texttt{IDVSSR::INDEX\_BUFFER}} (\texttt{r54:r55}) and pre-fills \texttt{r59, r60} with \texttt{VertexIndex} (\texttt{+ VERTEX\_OFFSET + GLOBAL\_ATTR\_OFF}) and \texttt{r61} with \texttt{InstanceIndex} (\texttt{+ INSTANCE\_OFFSET}) so Position and Varying shaders can pull vertex attributes directly.
  \item \textbf{Inline FAU \& Special Operands (\hyperref[sec:csf-tsd-fau]{\texttt{FAU\_k}}, \texttt{r8..r15}):} Any 8-bit source with \texttt{mode=2} reads zero-latency uniform \texttt{u0..u31} from \texttt{FAU\_k} (\texttt{slots64 = bits[63:56]}); \texttt{mode=3} reads \texttt{SmallConstant} (Table~1.2) or \texttt{FauSpecial} (\texttt{pc}, \texttt{lane\_id}, \texttt{lane\_width}, \texttt{tls\_ptr}, \texttt{wls\_ptr}, \texttt{blend\_descriptor\_0..7}).
\end{tightitemize}
\par\textbf{Worked Examples:}
\par\textbf{Example 1 (Exact small negative integer):}
\codebox{\mnem{S32\_TO\_F32} \dst{r0}, \src{r1}}
With \texttt{r1 = -10} (\texttt{0xFFFFFFF6}): exact conversion -\textgreater{} \texttt{r0 = -10.0f} (\texttt{0xC1200000}).
\par\textbf{Example 2 (Inexact 32-bit integer rounding):}
\codebox{\mnem{S32\_TO\_F32}\imod{.round\_zero} \dst{r0}, \src{r1}}
With \texttt{r1 = 0x01000001} (\texttt{2\textasciicircum{}24 + 1}, exceeds 24-bit mantissa): \texttt{.round\_zero} rounds down to \texttt{16777216.0f} (\texttt{0x4B800000}).
\par\textbf{Mined SPIR-V $\to$ Valhall ISA Example (Converting signed vertex index to float --- Vertex shader, Mali-G710 (v10)):}
\codebox{\textbf{; SPIR-V (input)}\\\%vid = OpLoad \%int \%gl\_VertexIndex\\\%f   = OpConvertSToF \%float \%vid\\[3pt]\textbf{; Valhall ISA (malioc output)}\\S32\_TO\_F32 r0, r60\textasciicircum{}}
\texttt{OpConvertSToF} maps directly to \texttt{S32\_TO\_F32} (default RNE rounding) consuming the preloaded \texttt{VertexIndex} in \texttt{r60}.
\begin{center}\footnotesize
\begin{tabularx}{\linewidth}{@{}lllX@{}}
\toprule
Variant & Arch & Class & Opcode Fields \\
\midrule
\texttt{(none)} & v9--v14 & \texttt{R1} & \texttt{opcode2=S32\_TO\_F32, opcode1=CONVERT} \\
\texttt{(none)} & v9--v14 & \texttt{R1} & \texttt{opcode2=S32\_TO\_F32, opcode1=CONVERT} \\
\texttt{(none)} & v9--v14 & \texttt{R1} & \texttt{opcode2=I32\_TO\_F32, opcode1=CONVERT} \\
\bottomrule
\end{tabularx}
\end{center}
\begin{center}\footnotesize
\begin{tabularx}{0.88\linewidth}{@{}l l X@{}}
\toprule
Field Name & Bits & Type / Fixed Value \\
\midrule
\texttt{src0} & \texttt{0:7} & \texttt{SourceEncoding} \\
\textit{Reserved} & \texttt{8:15} & \texttt{zero} \\
\texttt{opcode2} & \texttt{16:20} & \texttt{exact=S32\_TO\_F32} \\
\textit{Reserved} & \texttt{21:29} & \texttt{zero} \\
\hyperref[enum:round_extended_m]{\texttt{round}} & \texttt{30:32} & \texttt{round\_extended\_m (def none)} \\
\textit{Reserved} & \texttt{33} & \texttt{zero} \\
\textit{Reserved} & \texttt{34:35} & \texttt{ignore} \\
\hyperref[enum:widen_int_m]{\texttt{widen0}} & \texttt{36:39} & \texttt{exact=none} \\
\texttt{dest\_register\_index} & \texttt{40:45} & \texttt{uint} \\
\hyperref[enum:dest_width_m]{\texttt{dest\_width}} & \texttt{46:47} & \texttt{dest\_width\_m} \\
\texttt{opcode1} & \texttt{48:51} & \texttt{exact=CONVERT} \\
\texttt{instruction\_class} & \texttt{52:56} & \texttt{exact=R1} \\
\texttt{fau\_page\_index} & \texttt{57:58} & \texttt{uint} \\
\hyperref[enum:flow_control_m]{\texttt{flow}} & \texttt{59:62} & \texttt{flow\_control\_m} \\
\textit{Reserved} & \texttt{63} & \texttt{zero} \\
\bottomrule
\end{tabularx}
\end{center}
\textbf{Fields Explained:}
\begin{tightitemize}\footnotesize
  \item \texttt{src0}: Primary source operand (\texttt{r0..r63}, \texttt{u0..u63} uniform, \texttt{k*} FAU constant, or \texttt{t0..t2} temporary; \texttt{\textasciicircum{}} marks discard/last-use). \emph{Example:} {\scriptsize\ttfamily \mnem{FADD.f32} \dst{r0}, \src{r1}, \src{r2}} (\textbf{SPIR-V:} \texttt{OpFAdd (Operand 0)}; \textbf{GLSL:} \texttt{float c = a + b;})
  \item \texttt{opcode2}: Secondary sub-opcode further specializing instructions sharing \texttt{instruction\_class} and \texttt{opcode1} (e.g., conversions, SFU tables, moves). \emph{Example:} {\scriptsize\ttfamily \mnem{F32\_TO\_S32}\imod{.round\_zero} \dst{r0}, \src{r1}} (\textbf{SPIR-V:} \texttt{OpConvertFToS}; \textbf{GLSL:} \texttt{int y = int(x);})
  \item \hyperref[enum:round_extended_m]{\texttt{round}}: Rounding mode (\texttt{none}/RNE, \texttt{round\_up}, \texttt{round\_down}, \texttt{round\_zero}, \texttt{round\_na}) applied in FP ALU, conversions, \texttt{HADD}, \texttt{LDEXP}, and Stage C narrowing. \emph{Example:} {\scriptsize\ttfamily \mnem{FROUND.f32}\imod{.round\_down} \dst{r0}, \src{r1}} (\textbf{SPIR-V:} \texttt{GLSL.std.450 Floor / Ceil / Trunc}; \textbf{GLSL:} \texttt{float y = floor(x);})
  \item \hyperref[enum:widen_int_m]{\texttt{widen0}}: Stage A widening extractor on \texttt{src0} (promotes FP16 \texttt{h0}/\texttt{h1} to FP32, or sign/zero-extends \texttt{b0..b3}/\texttt{h0..h1}/\texttt{w0..w1} to 32/64-bit before \texttt{abs}/\texttt{neg}). \emph{Example:} {\scriptsize\ttfamily \mnem{IADD.s32} \dst{r0}, \src{r1}\imod{.h0}, \src{r2}} (\textbf{SPIR-V:} \texttt{OpSConvert / OpFConvert}; \textbf{GLSL:} \texttt{int c = int(int16\_t(a)) + b;})
  \item \texttt{dest\_register\_index}: 6-bit destination GPR index (\texttt{r0..r63}) receiving the Stage C writeback result (even-aligned pair \texttt{[rN:rN+1]} for 64-bit ops). \emph{Example:} {\scriptsize\ttfamily \mnem{FMA.f32} \dst{r0}, \src{r1}, \src{r2}, \src{r3}} (\textbf{SPIR-V:} \texttt{SSA Result ID}; \textbf{GLSL:} \texttt{float r0 = fma(r1, r2, r3);})
  \item \hyperref[enum:dest_width_m]{\texttt{dest\_width}}: Stage C destination lane write-mask (\texttt{none} = full 32-bit, \texttt{h0} = low 16-bit, \texttt{h1} = high 16-bit without bit-shifting). \emph{Example:} {\scriptsize\ttfamily \mnem{FADD.v2f16} \dst{r0}\imod{.h1}, \src{r1}, \src{r2}} (\textbf{SPIR-V:} \texttt{OpVectorInsertDynamic / WriteMask}; \textbf{GLSL:} \texttt{dst.y = a.y + b.y;})
  \item \texttt{opcode1}: Primary sub-opcode within the instruction class uniquely identifying the operation family. \emph{Example:} {\scriptsize\ttfamily \mnem{IADD.s32} \dst{r0}, \src{r1}, \src{r2}} (\textbf{SPIR-V:} \texttt{OpIAdd (operation selector)}; \textbf{GLSL:} \texttt{// Primary opcode bits})
  \item \texttt{instruction\_class}: 4-bit major instruction class selector (\texttt{bits 48:51}) partitioning the 64-bit Valhall opcode space (e.g., \texttt{FMA}, \texttt{CVT}, \texttt{SFU}, \texttt{LS}, \texttt{TEX}, \texttt{BRANCH}). \emph{Example:} {\scriptsize\ttfamily \mnem{FADD.f32} \dst{r0}, \src{r1}, \src{r2}} (\textbf{SPIR-V:} \texttt{N/A (encoding class)}; \textbf{GLSL:} \texttt{// Hardware execution pipe / decode class})
  \item \texttt{fau\_page\_index}: Selects the FAU (Fast Access Uniform) page (\texttt{0..3}) when a source operand references hardware small constants, thread IDs, or lane constants. \emph{Example:} {\scriptsize\ttfamily \mnem{FADD.f32} \dst{r0}, \src{r1}, \src{k1f}} (\textbf{SPIR-V:} \texttt{OpConstant / BuiltIn}; \textbf{GLSL:} \texttt{float y = x + 1.0;})
  \item \hyperref[enum:flow_control_m]{\texttt{flow}}: 4-bit embedded flow-control \& scoreboard wait modifier (\texttt{bits 59:62}: \texttt{none}, \texttt{wait0..wait0126}, \texttt{wait}, \texttt{reconverge}, \texttt{discard}, \texttt{end}). \emph{Example:} {\scriptsize\ttfamily \mnem{FADD.f32}\imod{.wait0.reconverge} \dst{r0}, \src{r1}, \src{r2}} (\textbf{SPIR-V:} \texttt{OpControlBarrier / OpReturn}; \textbf{GLSL:} \texttt{// Scoreboard wait + warp reconvergence})
\end{tightitemize}
\clearpage

\subsection*{S8\_TO\_F32}
\addcontentsline{toc}{subsection}{S8\_TO\_F32}
Converts the signed 8-bit byte selected by \texttt{lane0} (\texttt{.b0}..\texttt{.b3}) of \texttt{src0} exactly to IEEE-754 FP32 and writes the result to \texttt{dest} (arch \textless{}= v10).
\par\textbf{Notes:}
\begin{tightitemize}
  \item Arch v9-v10 only (\texttt{kraid/nir.rs:740}).
  \item \textbf{Source references:} \texttt{kraid/ops.rs:2575-2600}, \texttt{encode\_v9.rs:1984-1987}, \texttt{nir.rs:738-749}.
\end{tightitemize}
\codebox{\mnem{S8\_TO\_F32}\imod{\{.flow\}} \dst{dest}\imod{\{.dest\_width\}}, \src{src0}\imod{\{.lane0\}}}
\textbf{Operands:}
\begin{tightitemize}\footnotesize
  \item \texttt{\dst{dest}} (\texttt{dest\_register\_index}, \texttt{dest\_width}): 32-bit destination GPR receiving the converted value.
  \item \texttt{\src{src0}}: Signed 8-bit byte lane (\texttt{lane0}) converted exactly to FP32.
\end{tightitemize}
\par\textbf{Closest SPIR-V Op:} \texttt{OpConvertSToF (s8 -\textgreater{} f32, v9--v10)}
\par\textbf{Pseudo-GLSL (Operational Semantics):}
\codebox{dest = float(bitfieldExtract(int(src0), lane0 * 8, 8));}
\textbf{Execution Context:}
\begin{tightitemize}\footnotesize
  \item \textbf{Compute stage preloads (\hyperref[sec:csf-spd]{\texttt{SPD\_k}}, \texttt{r16..r23}):} Hardware pre-fills \texttt{r55..r57} with \texttt{WorkGroupID.xyz + JOB\_OFFSET} (from \hyperref[sec:csf-regs]{\texttt{COMPUTESR::JOB\_OFFSET\_X/Y/Z}} \texttt{r34..r36}) and \texttt{r60..r62} with \texttt{LocalInvocationID.xyz} or \texttt{GlobalInvocationID.xyz} (from \texttt{WG\_SIZE} \texttt{r33} \& \texttt{JOB\_SIZE} \texttt{r37..r39}, selected by \texttt{ShaderProgram::preload}) so compute shaders can derive thread/element indices without memory loads.
  \item \textbf{Vertex stage preloads (\hyperref[sec:csf-spd]{\texttt{SPD\_0}/\texttt{SPD\_1}}, \texttt{r16..r19}):} The hardware Index Fetcher reads \hyperref[sec:csf-regs]{\texttt{IDVSSR::INDEX\_BUFFER}} (\texttt{r54:r55}) and pre-fills \texttt{r59, r60} with \texttt{VertexIndex} (\texttt{+ VERTEX\_OFFSET + GLOBAL\_ATTR\_OFF}) and \texttt{r61} with \texttt{InstanceIndex} (\texttt{+ INSTANCE\_OFFSET}) so Position and Varying shaders can pull vertex attributes directly.
  \item \textbf{Inline FAU \& Special Operands (\hyperref[sec:csf-tsd-fau]{\texttt{FAU\_k}}, \texttt{r8..r15}):} Any 8-bit source with \texttt{mode=2} reads zero-latency uniform \texttt{u0..u31} from \texttt{FAU\_k} (\texttt{slots64 = bits[63:56]}); \texttt{mode=3} reads \texttt{SmallConstant} (Table~1.2) or \texttt{FauSpecial} (\texttt{pc}, \texttt{lane\_id}, \texttt{lane\_width}, \texttt{tls\_ptr}, \texttt{wls\_ptr}, \texttt{blend\_descriptor\_0..7}).
\end{tightitemize}
\par\textbf{Worked Examples:}
\par\textbf{Example 1 (Negative byte 2 extraction (v9--v10)):}
\codebox{\mnem{S8\_TO\_F32} \dst{r0}, \src{r1}\imod{.b2}}
With \texttt{r1 = 0x00FC0000} (\texttt{b2 = -4}): exact conversion -\textgreater{} \texttt{r0 = -4.0f} (\texttt{0xC0800000}).
\par\textbf{Example 2 (Positive byte 0 extraction):}
\codebox{\mnem{S8\_TO\_F32} \dst{r0}, \src{r1}\imod{.b0}}
With \texttt{r1 = 0x0000007F} (\texttt{b0 = +127}): exact conversion -\textgreater{} \texttt{r0 = 127.0f} (\texttt{0x42FE0000}).

\begin{center}\footnotesize
\begin{tabularx}{\linewidth}{@{}lllX@{}}
\toprule
Variant & Arch & Class & Opcode Fields \\
\midrule
\texttt{(none)} & v9--v14 & \texttt{R1} & \texttt{opcode2=S8\_TO\_F32, opcode1=CONVERT} \\
\bottomrule
\end{tabularx}
\end{center}
\begin{center}\footnotesize
\begin{tabularx}{0.88\linewidth}{@{}l l X@{}}
\toprule
Field Name & Bits & Type / Fixed Value \\
\midrule
\texttt{src0} & \texttt{0:7} & \texttt{SourceEncoding} \\
\textit{Reserved} & \texttt{8:15} & \texttt{zero} \\
\texttt{opcode2} & \texttt{16:20} & \texttt{exact=S8\_TO\_F32} \\
\textit{Reserved} & \texttt{21:27} & \texttt{zero} \\
\hyperref[enum:lane8_m]{\texttt{lane0}} & \texttt{28:29} & \texttt{lane8\_m (def b0)} \\
\textit{Reserved} & \texttt{30:39} & \texttt{zero} \\
\texttt{dest\_register\_index} & \texttt{40:45} & \texttt{uint} \\
\hyperref[enum:dest_width_m]{\texttt{dest\_width}} & \texttt{46:47} & \texttt{dest\_width\_m} \\
\texttt{opcode1} & \texttt{48:51} & \texttt{exact=CONVERT} \\
\texttt{instruction\_class} & \texttt{52:56} & \texttt{exact=R1} \\
\texttt{fau\_page\_index} & \texttt{57:58} & \texttt{uint} \\
\hyperref[enum:flow_control_m]{\texttt{flow}} & \texttt{59:62} & \texttt{flow\_control\_m} \\
\textit{Reserved} & \texttt{63} & \texttt{zero} \\
\bottomrule
\end{tabularx}
\end{center}
\textbf{Fields Explained:}
\begin{tightitemize}\footnotesize
  \item \texttt{src0}: Primary source operand (\texttt{r0..r63}, \texttt{u0..u63} uniform, \texttt{k*} FAU constant, or \texttt{t0..t2} temporary; \texttt{\textasciicircum{}} marks discard/last-use). \emph{Example:} {\scriptsize\ttfamily \mnem{FADD.f32} \dst{r0}, \src{r1}, \src{r2}} (\textbf{SPIR-V:} \texttt{OpFAdd (Operand 0)}; \textbf{GLSL:} \texttt{float c = a + b;})
  \item \texttt{opcode2}: Secondary sub-opcode further specializing instructions sharing \texttt{instruction\_class} and \texttt{opcode1} (e.g., conversions, SFU tables, moves). \emph{Example:} {\scriptsize\ttfamily \mnem{F32\_TO\_S32}\imod{.round\_zero} \dst{r0}, \src{r1}} (\textbf{SPIR-V:} \texttt{OpConvertFToS}; \textbf{GLSL:} \texttt{int y = int(x);})
  \item \hyperref[enum:lane8_m]{\texttt{lane0}}: Stage A single-lane extractor on \texttt{src0} (\texttt{b0..b3} for 8-bit or \texttt{h0}/\texttt{h1} for 16-bit narrow conversions and \texttt{MKVEC}). \emph{Example:} {\scriptsize\ttfamily \mnem{F16\_TO\_F32} \dst{r0}, \src{r1}\imod{.h1}} (\textbf{SPIR-V:} \texttt{OpCompositeExtract}; \textbf{GLSL:} \texttt{float y = float(x.y);})
  \item \texttt{dest\_register\_index}: 6-bit destination GPR index (\texttt{r0..r63}) receiving the Stage C writeback result (even-aligned pair \texttt{[rN:rN+1]} for 64-bit ops). \emph{Example:} {\scriptsize\ttfamily \mnem{FMA.f32} \dst{r0}, \src{r1}, \src{r2}, \src{r3}} (\textbf{SPIR-V:} \texttt{SSA Result ID}; \textbf{GLSL:} \texttt{float r0 = fma(r1, r2, r3);})
  \item \hyperref[enum:dest_width_m]{\texttt{dest\_width}}: Stage C destination lane write-mask (\texttt{none} = full 32-bit, \texttt{h0} = low 16-bit, \texttt{h1} = high 16-bit without bit-shifting). \emph{Example:} {\scriptsize\ttfamily \mnem{FADD.v2f16} \dst{r0}\imod{.h1}, \src{r1}, \src{r2}} (\textbf{SPIR-V:} \texttt{OpVectorInsertDynamic / WriteMask}; \textbf{GLSL:} \texttt{dst.y = a.y + b.y;})
  \item \texttt{opcode1}: Primary sub-opcode within the instruction class uniquely identifying the operation family. \emph{Example:} {\scriptsize\ttfamily \mnem{IADD.s32} \dst{r0}, \src{r1}, \src{r2}} (\textbf{SPIR-V:} \texttt{OpIAdd (operation selector)}; \textbf{GLSL:} \texttt{// Primary opcode bits})
  \item \texttt{instruction\_class}: 4-bit major instruction class selector (\texttt{bits 48:51}) partitioning the 64-bit Valhall opcode space (e.g., \texttt{FMA}, \texttt{CVT}, \texttt{SFU}, \texttt{LS}, \texttt{TEX}, \texttt{BRANCH}). \emph{Example:} {\scriptsize\ttfamily \mnem{FADD.f32} \dst{r0}, \src{r1}, \src{r2}} (\textbf{SPIR-V:} \texttt{N/A (encoding class)}; \textbf{GLSL:} \texttt{// Hardware execution pipe / decode class})
  \item \texttt{fau\_page\_index}: Selects the FAU (Fast Access Uniform) page (\texttt{0..3}) when a source operand references hardware small constants, thread IDs, or lane constants. \emph{Example:} {\scriptsize\ttfamily \mnem{FADD.f32} \dst{r0}, \src{r1}, \src{k1f}} (\textbf{SPIR-V:} \texttt{OpConstant / BuiltIn}; \textbf{GLSL:} \texttt{float y = x + 1.0;})
  \item \hyperref[enum:flow_control_m]{\texttt{flow}}: 4-bit embedded flow-control \& scoreboard wait modifier (\texttt{bits 59:62}: \texttt{none}, \texttt{wait0..wait0126}, \texttt{wait}, \texttt{reconverge}, \texttt{discard}, \texttt{end}). \emph{Example:} {\scriptsize\ttfamily \mnem{FADD.f32}\imod{.wait0.reconverge} \dst{r0}, \src{r1}, \src{r2}} (\textbf{SPIR-V:} \texttt{OpControlBarrier / OpReturn}; \textbf{GLSL:} \texttt{// Scoreboard wait + warp reconvergence})
\end{tightitemize}
\clearpage

\subsection*{S8\_TO\_S32}
\addcontentsline{toc}{subsection}{S8\_TO\_S32}
Extracts the signed 8-bit byte selected by \texttt{lane0} (\texttt{.b0}..\texttt{.b3}) of \texttt{src0}, sign-extends it to 32 bits, and writes the result to \texttt{dest}.
\par\textbf{Notes:}
\begin{tightitemize}
  \item Exact sign extension.
  \item \textbf{Source references:} \texttt{isa-v9-v14.xml}, \texttt{bifrost/valhall/ISA.xml}.
\end{tightitemize}
\codebox{\mnem{S8\_TO\_S32}\imod{\{.flow\}} \dst{dest}\imod{\{.dest\_width\}}, \src{src0}\imod{\{.lane0\}}}
\textbf{Operands:}
\begin{tightitemize}\footnotesize
  \item \texttt{\dst{dest}} (\texttt{dest\_register\_index}, \texttt{dest\_width}): 32-bit destination GPR receiving the converted value.
  \item \texttt{\src{src0}}: Signed 8-bit byte lane (\texttt{lane0}) sign-extended to 32-bit integer.
\end{tightitemize}
\par\textbf{Closest SPIR-V Op:} \texttt{OpSConvert (s8 -\textgreater{} s32)}
\par\textbf{Pseudo-GLSL (Operational Semantics):}
\codebox{dest = bitfieldExtract(int(src0), lane0 * 8, 8);}
\textbf{Execution Context:}
\begin{tightitemize}\footnotesize
  \item \textbf{Compute stage preloads (\hyperref[sec:csf-spd]{\texttt{SPD\_k}}, \texttt{r16..r23}):} Hardware pre-fills \texttt{r55..r57} with \texttt{WorkGroupID.xyz + JOB\_OFFSET} (from \hyperref[sec:csf-regs]{\texttt{COMPUTESR::JOB\_OFFSET\_X/Y/Z}} \texttt{r34..r36}) and \texttt{r60..r62} with \texttt{LocalInvocationID.xyz} or \texttt{GlobalInvocationID.xyz} (from \texttt{WG\_SIZE} \texttt{r33} \& \texttt{JOB\_SIZE} \texttt{r37..r39}, selected by \texttt{ShaderProgram::preload}) so compute shaders can derive thread/element indices without memory loads.
  \item \textbf{Vertex stage preloads (\hyperref[sec:csf-spd]{\texttt{SPD\_0}/\texttt{SPD\_1}}, \texttt{r16..r19}):} The hardware Index Fetcher reads \hyperref[sec:csf-regs]{\texttt{IDVSSR::INDEX\_BUFFER}} (\texttt{r54:r55}) and pre-fills \texttt{r59, r60} with \texttt{VertexIndex} (\texttt{+ VERTEX\_OFFSET + GLOBAL\_ATTR\_OFF}) and \texttt{r61} with \texttt{InstanceIndex} (\texttt{+ INSTANCE\_OFFSET}) so Position and Varying shaders can pull vertex attributes directly.
  \item \textbf{Inline FAU \& Special Operands (\hyperref[sec:csf-tsd-fau]{\texttt{FAU\_k}}, \texttt{r8..r15}):} Any 8-bit source with \texttt{mode=2} reads zero-latency uniform \texttt{u0..u31} from \texttt{FAU\_k} (\texttt{slots64 = bits[63:56]}); \texttt{mode=3} reads \texttt{SmallConstant} (Table~1.2) or \texttt{FauSpecial} (\texttt{pc}, \texttt{lane\_id}, \texttt{lane\_width}, \texttt{tls\_ptr}, \texttt{wls\_ptr}, \texttt{blend\_descriptor\_0..7}).
\end{tightitemize}
\par\textbf{Worked Examples:}
\par\textbf{Example 1 (Negative byte sign extension):}
\codebox{\mnem{S8\_TO\_S32} \dst{r0}, \src{r1}\imod{.b0}}
With \texttt{r1.b0 = 0xFE} (\texttt{-2}): sign-extends bit 7 across \texttt{[31:8]} -\textgreater{} \texttt{r0 = 0xFFFFFFFE}.
\par\textbf{Example 2 (Byte 3 positive extension):}
\codebox{\mnem{S8\_TO\_S32} \dst{r0}, \src{r1}\imod{.b3}}
With \texttt{r1 = 0x7F000000}: sign-extends \texttt{0x7F} (\texttt{+127}) -\textgreater{} \texttt{r0 = 0x0000007F}.

\begin{center}\footnotesize
\begin{tabularx}{\linewidth}{@{}lllX@{}}
\toprule
Variant & Arch & Class & Opcode Fields \\
\midrule
\texttt{(none)} & v9--v14 & \texttt{R1} & \texttt{opcode2=S8\_TO\_S32, opcode1=CONVERT} \\
\texttt{(none)} & v9--v14 & \texttt{R1} & \texttt{opcode2=I8\_TO\_I32, opcode1=CONVERT} \\
\bottomrule
\end{tabularx}
\end{center}
\begin{center}\footnotesize
\begin{tabularx}{0.88\linewidth}{@{}l l X@{}}
\toprule
Field Name & Bits & Type / Fixed Value \\
\midrule
\texttt{src0} & \texttt{0:7} & \texttt{SourceEncoding} \\
\textit{Reserved} & \texttt{8:15} & \texttt{zero} \\
\texttt{opcode2} & \texttt{16:20} & \texttt{exact=S8\_TO\_S32} \\
\textit{Reserved} & \texttt{21:27} & \texttt{zero} \\
\hyperref[enum:lane8_m]{\texttt{lane0}} & \texttt{28:29} & \texttt{lane8\_m (def b0)} \\
\textit{Reserved} & \texttt{30:39} & \texttt{zero} \\
\texttt{dest\_register\_index} & \texttt{40:45} & \texttt{uint} \\
\hyperref[enum:dest_width_m]{\texttt{dest\_width}} & \texttt{46:47} & \texttt{dest\_width\_m} \\
\texttt{opcode1} & \texttt{48:51} & \texttt{exact=CONVERT} \\
\texttt{instruction\_class} & \texttt{52:56} & \texttt{exact=R1} \\
\texttt{fau\_page\_index} & \texttt{57:58} & \texttt{uint} \\
\hyperref[enum:flow_control_m]{\texttt{flow}} & \texttt{59:62} & \texttt{flow\_control\_m} \\
\textit{Reserved} & \texttt{63} & \texttt{zero} \\
\bottomrule
\end{tabularx}
\end{center}
\textbf{Fields Explained:}
\begin{tightitemize}\footnotesize
  \item \texttt{src0}: Primary source operand (\texttt{r0..r63}, \texttt{u0..u63} uniform, \texttt{k*} FAU constant, or \texttt{t0..t2} temporary; \texttt{\textasciicircum{}} marks discard/last-use). \emph{Example:} {\scriptsize\ttfamily \mnem{FADD.f32} \dst{r0}, \src{r1}, \src{r2}} (\textbf{SPIR-V:} \texttt{OpFAdd (Operand 0)}; \textbf{GLSL:} \texttt{float c = a + b;})
  \item \texttt{opcode2}: Secondary sub-opcode further specializing instructions sharing \texttt{instruction\_class} and \texttt{opcode1} (e.g., conversions, SFU tables, moves). \emph{Example:} {\scriptsize\ttfamily \mnem{F32\_TO\_S32}\imod{.round\_zero} \dst{r0}, \src{r1}} (\textbf{SPIR-V:} \texttt{OpConvertFToS}; \textbf{GLSL:} \texttt{int y = int(x);})
  \item \hyperref[enum:lane8_m]{\texttt{lane0}}: Stage A single-lane extractor on \texttt{src0} (\texttt{b0..b3} for 8-bit or \texttt{h0}/\texttt{h1} for 16-bit narrow conversions and \texttt{MKVEC}). \emph{Example:} {\scriptsize\ttfamily \mnem{F16\_TO\_F32} \dst{r0}, \src{r1}\imod{.h1}} (\textbf{SPIR-V:} \texttt{OpCompositeExtract}; \textbf{GLSL:} \texttt{float y = float(x.y);})
  \item \texttt{dest\_register\_index}: 6-bit destination GPR index (\texttt{r0..r63}) receiving the Stage C writeback result (even-aligned pair \texttt{[rN:rN+1]} for 64-bit ops). \emph{Example:} {\scriptsize\ttfamily \mnem{FMA.f32} \dst{r0}, \src{r1}, \src{r2}, \src{r3}} (\textbf{SPIR-V:} \texttt{SSA Result ID}; \textbf{GLSL:} \texttt{float r0 = fma(r1, r2, r3);})
  \item \hyperref[enum:dest_width_m]{\texttt{dest\_width}}: Stage C destination lane write-mask (\texttt{none} = full 32-bit, \texttt{h0} = low 16-bit, \texttt{h1} = high 16-bit without bit-shifting). \emph{Example:} {\scriptsize\ttfamily \mnem{FADD.v2f16} \dst{r0}\imod{.h1}, \src{r1}, \src{r2}} (\textbf{SPIR-V:} \texttt{OpVectorInsertDynamic / WriteMask}; \textbf{GLSL:} \texttt{dst.y = a.y + b.y;})
  \item \texttt{opcode1}: Primary sub-opcode within the instruction class uniquely identifying the operation family. \emph{Example:} {\scriptsize\ttfamily \mnem{IADD.s32} \dst{r0}, \src{r1}, \src{r2}} (\textbf{SPIR-V:} \texttt{OpIAdd (operation selector)}; \textbf{GLSL:} \texttt{// Primary opcode bits})
  \item \texttt{instruction\_class}: 4-bit major instruction class selector (\texttt{bits 48:51}) partitioning the 64-bit Valhall opcode space (e.g., \texttt{FMA}, \texttt{CVT}, \texttt{SFU}, \texttt{LS}, \texttt{TEX}, \texttt{BRANCH}). \emph{Example:} {\scriptsize\ttfamily \mnem{FADD.f32} \dst{r0}, \src{r1}, \src{r2}} (\textbf{SPIR-V:} \texttt{N/A (encoding class)}; \textbf{GLSL:} \texttt{// Hardware execution pipe / decode class})
  \item \texttt{fau\_page\_index}: Selects the FAU (Fast Access Uniform) page (\texttt{0..3}) when a source operand references hardware small constants, thread IDs, or lane constants. \emph{Example:} {\scriptsize\ttfamily \mnem{FADD.f32} \dst{r0}, \src{r1}, \src{k1f}} (\textbf{SPIR-V:} \texttt{OpConstant / BuiltIn}; \textbf{GLSL:} \texttt{float y = x + 1.0;})
  \item \hyperref[enum:flow_control_m]{\texttt{flow}}: 4-bit embedded flow-control \& scoreboard wait modifier (\texttt{bits 59:62}: \texttt{none}, \texttt{wait0..wait0126}, \texttt{wait}, \texttt{reconverge}, \texttt{discard}, \texttt{end}). \emph{Example:} {\scriptsize\ttfamily \mnem{FADD.f32}\imod{.wait0.reconverge} \dst{r0}, \src{r1}, \src{r2}} (\textbf{SPIR-V:} \texttt{OpControlBarrier / OpReturn}; \textbf{GLSL:} \texttt{// Scoreboard wait + warp reconvergence})
\end{tightitemize}
\clearpage

\subsection*{U16\_TO\_F32}
\addcontentsline{toc}{subsection}{U16\_TO\_F32}
Converts the unsigned 16-bit integer selected by \texttt{lane0} (\texttt{.h0} or \texttt{.h1}) of \texttt{src0} exactly to IEEE-754 FP32 and writes the result to \texttt{dest} (arch \textless{}= v10).
\par\textbf{Notes:}
\begin{tightitemize}
  \item Arch v9-v10 only (\texttt{kraid/nir.rs:692}).
  \item On v11+ use \texttt{U32\_TO\_F32} with \texttt{widen0=h0}.
  \item \textbf{Source references:} \texttt{kraid/ops.rs:2575-2600}, \texttt{encode\_v9.rs:1996-1999}, \texttt{nir.rs:690-701}.
\end{tightitemize}
\codebox{\mnem{U16\_TO\_F32}\imod{\{.flow\}} \dst{dest}\imod{\{.dest\_width\}}, \src{src0}\imod{\{.lane0\}}}
\textbf{Operands:}
\begin{tightitemize}\footnotesize
  \item \texttt{\dst{dest}} (\texttt{dest\_register\_index}, \texttt{dest\_width}): 32-bit destination GPR receiving the converted value.
  \item \texttt{\src{src0}}: Unsigned 16-bit lane (\texttt{lane0}) converted exactly to FP32.
\end{tightitemize}
\par\textbf{Closest SPIR-V Op:} \texttt{OpConvertUToF (u16 -\textgreater{} f32, v9--v10)}
\par\textbf{Pseudo-GLSL (Operational Semantics):}
\codebox{dest = float(bitfieldExtract(uint(src0), lane0 * 16, 16));}
\textbf{Execution Context:}
\begin{tightitemize}\footnotesize
  \item \textbf{Compute stage preloads (\hyperref[sec:csf-spd]{\texttt{SPD\_k}}, \texttt{r16..r23}):} Hardware pre-fills \texttt{r55..r57} with \texttt{WorkGroupID.xyz + JOB\_OFFSET} (from \hyperref[sec:csf-regs]{\texttt{COMPUTESR::JOB\_OFFSET\_X/Y/Z}} \texttt{r34..r36}) and \texttt{r60..r62} with \texttt{LocalInvocationID.xyz} or \texttt{GlobalInvocationID.xyz} (from \texttt{WG\_SIZE} \texttt{r33} \& \texttt{JOB\_SIZE} \texttt{r37..r39}, selected by \texttt{ShaderProgram::preload}) so compute shaders can derive thread/element indices without memory loads.
  \item \textbf{Vertex stage preloads (\hyperref[sec:csf-spd]{\texttt{SPD\_0}/\texttt{SPD\_1}}, \texttt{r16..r19}):} The hardware Index Fetcher reads \hyperref[sec:csf-regs]{\texttt{IDVSSR::INDEX\_BUFFER}} (\texttt{r54:r55}) and pre-fills \texttt{r59, r60} with \texttt{VertexIndex} (\texttt{+ VERTEX\_OFFSET + GLOBAL\_ATTR\_OFF}) and \texttt{r61} with \texttt{InstanceIndex} (\texttt{+ INSTANCE\_OFFSET}) so Position and Varying shaders can pull vertex attributes directly.
  \item \textbf{Inline FAU \& Special Operands (\hyperref[sec:csf-tsd-fau]{\texttt{FAU\_k}}, \texttt{r8..r15}):} Any 8-bit source with \texttt{mode=2} reads zero-latency uniform \texttt{u0..u31} from \texttt{FAU\_k} (\texttt{slots64 = bits[63:56]}); \texttt{mode=3} reads \texttt{SmallConstant} (Table~1.2) or \texttt{FauSpecial} (\texttt{pc}, \texttt{lane\_id}, \texttt{lane\_width}, \texttt{tls\_ptr}, \texttt{wls\_ptr}, \texttt{blend\_descriptor\_0..7}).
\end{tightitemize}
\par\textbf{Worked Examples:}
\par\textbf{Example 1 (Low halfword unsigned -\textgreater{} FP32 (v9--v10)):}
\codebox{\mnem{U16\_TO\_F32} \dst{r0}, \src{r1}\imod{.h0}}
With \texttt{r1.h0 = 0x0100} (\texttt{256}): exact conversion -\textgreater{} \texttt{r0 = 256.0f} (\texttt{0x43800000}).
\par\textbf{Example 2 (Max uint16 (`65535`)):}
\codebox{\mnem{U16\_TO\_F32} \dst{r0}, \src{r1}\imod{.h1}}
With \texttt{r1 = 0xFFFF0000}: exact conversion (fits in 24-bit mantissa) -\textgreater{} \texttt{r0 = 65535.0f} (\texttt{0x477FFF00}).
\par\textbf{Mined SPIR-V $\to$ Valhall ISA Example (Unpacking packed 16-bit pixel coordinates to float --- Fragment shader, Mali-G710 (v10)):}
\codebox{\textbf{; SPIR-V (input)}\\\%fc = OpLoad \%v4float \%gl\_FragCoord\\[3pt]\textbf{; Valhall ISA (malioc output)}\\U16\_TO\_F32 r3, r59.h1\\U16\_TO\_F32 r2, r59\textasciicircum{}.h0}
Hardware packs integer fragment \texttt{(x, y)} coordinates into halfwords \texttt{.h0} and \texttt{.h1} of \texttt{r59}; \texttt{U16\_TO\_F32} converts each halfword losslessly to FP32.
\begin{center}\footnotesize
\begin{tabularx}{\linewidth}{@{}lllX@{}}
\toprule
Variant & Arch & Class & Opcode Fields \\
\midrule
\texttt{(none)} & v9--v14 & \texttt{R1} & \texttt{opcode2=U16\_TO\_F32, opcode1=CONVERT} \\
\bottomrule
\end{tabularx}
\end{center}
\begin{center}\footnotesize
\begin{tabularx}{0.88\linewidth}{@{}l l X@{}}
\toprule
Field Name & Bits & Type / Fixed Value \\
\midrule
\texttt{src0} & \texttt{0:7} & \texttt{SourceEncoding} \\
\textit{Reserved} & \texttt{8:15} & \texttt{zero} \\
\texttt{opcode2} & \texttt{16:20} & \texttt{exact=U16\_TO\_F32} \\
\textit{Reserved} & \texttt{21:27} & \texttt{zero} \\
\hyperref[enum:lane16_m]{\texttt{lane0}} & \texttt{28} & \texttt{lane16\_m (def h0)} \\
\textit{Reserved} & \texttt{29:39} & \texttt{zero} \\
\texttt{dest\_register\_index} & \texttt{40:45} & \texttt{uint} \\
\hyperref[enum:dest_width_m]{\texttt{dest\_width}} & \texttt{46:47} & \texttt{dest\_width\_m} \\
\texttt{opcode1} & \texttt{48:51} & \texttt{exact=CONVERT} \\
\texttt{instruction\_class} & \texttt{52:56} & \texttt{exact=R1} \\
\texttt{fau\_page\_index} & \texttt{57:58} & \texttt{uint} \\
\hyperref[enum:flow_control_m]{\texttt{flow}} & \texttt{59:62} & \texttt{flow\_control\_m} \\
\textit{Reserved} & \texttt{63} & \texttt{zero} \\
\bottomrule
\end{tabularx}
\end{center}
\textbf{Fields Explained:}
\begin{tightitemize}\footnotesize
  \item \texttt{src0}: Primary source operand (\texttt{r0..r63}, \texttt{u0..u63} uniform, \texttt{k*} FAU constant, or \texttt{t0..t2} temporary; \texttt{\textasciicircum{}} marks discard/last-use). \emph{Example:} {\scriptsize\ttfamily \mnem{FADD.f32} \dst{r0}, \src{r1}, \src{r2}} (\textbf{SPIR-V:} \texttt{OpFAdd (Operand 0)}; \textbf{GLSL:} \texttt{float c = a + b;})
  \item \texttt{opcode2}: Secondary sub-opcode further specializing instructions sharing \texttt{instruction\_class} and \texttt{opcode1} (e.g., conversions, SFU tables, moves). \emph{Example:} {\scriptsize\ttfamily \mnem{F32\_TO\_S32}\imod{.round\_zero} \dst{r0}, \src{r1}} (\textbf{SPIR-V:} \texttt{OpConvertFToS}; \textbf{GLSL:} \texttt{int y = int(x);})
  \item \hyperref[enum:lane16_m]{\texttt{lane0}}: Stage A single-lane extractor on \texttt{src0} (\texttt{b0..b3} for 8-bit or \texttt{h0}/\texttt{h1} for 16-bit narrow conversions and \texttt{MKVEC}). \emph{Example:} {\scriptsize\ttfamily \mnem{F16\_TO\_F32} \dst{r0}, \src{r1}\imod{.h1}} (\textbf{SPIR-V:} \texttt{OpCompositeExtract}; \textbf{GLSL:} \texttt{float y = float(x.y);})
  \item \texttt{dest\_register\_index}: 6-bit destination GPR index (\texttt{r0..r63}) receiving the Stage C writeback result (even-aligned pair \texttt{[rN:rN+1]} for 64-bit ops). \emph{Example:} {\scriptsize\ttfamily \mnem{FMA.f32} \dst{r0}, \src{r1}, \src{r2}, \src{r3}} (\textbf{SPIR-V:} \texttt{SSA Result ID}; \textbf{GLSL:} \texttt{float r0 = fma(r1, r2, r3);})
  \item \hyperref[enum:dest_width_m]{\texttt{dest\_width}}: Stage C destination lane write-mask (\texttt{none} = full 32-bit, \texttt{h0} = low 16-bit, \texttt{h1} = high 16-bit without bit-shifting). \emph{Example:} {\scriptsize\ttfamily \mnem{FADD.v2f16} \dst{r0}\imod{.h1}, \src{r1}, \src{r2}} (\textbf{SPIR-V:} \texttt{OpVectorInsertDynamic / WriteMask}; \textbf{GLSL:} \texttt{dst.y = a.y + b.y;})
  \item \texttt{opcode1}: Primary sub-opcode within the instruction class uniquely identifying the operation family. \emph{Example:} {\scriptsize\ttfamily \mnem{IADD.s32} \dst{r0}, \src{r1}, \src{r2}} (\textbf{SPIR-V:} \texttt{OpIAdd (operation selector)}; \textbf{GLSL:} \texttt{// Primary opcode bits})
  \item \texttt{instruction\_class}: 4-bit major instruction class selector (\texttt{bits 48:51}) partitioning the 64-bit Valhall opcode space (e.g., \texttt{FMA}, \texttt{CVT}, \texttt{SFU}, \texttt{LS}, \texttt{TEX}, \texttt{BRANCH}). \emph{Example:} {\scriptsize\ttfamily \mnem{FADD.f32} \dst{r0}, \src{r1}, \src{r2}} (\textbf{SPIR-V:} \texttt{N/A (encoding class)}; \textbf{GLSL:} \texttt{// Hardware execution pipe / decode class})
  \item \texttt{fau\_page\_index}: Selects the FAU (Fast Access Uniform) page (\texttt{0..3}) when a source operand references hardware small constants, thread IDs, or lane constants. \emph{Example:} {\scriptsize\ttfamily \mnem{FADD.f32} \dst{r0}, \src{r1}, \src{k1f}} (\textbf{SPIR-V:} \texttt{OpConstant / BuiltIn}; \textbf{GLSL:} \texttt{float y = x + 1.0;})
  \item \hyperref[enum:flow_control_m]{\texttt{flow}}: 4-bit embedded flow-control \& scoreboard wait modifier (\texttt{bits 59:62}: \texttt{none}, \texttt{wait0..wait0126}, \texttt{wait}, \texttt{reconverge}, \texttt{discard}, \texttt{end}). \emph{Example:} {\scriptsize\ttfamily \mnem{FADD.f32}\imod{.wait0.reconverge} \dst{r0}, \src{r1}, \src{r2}} (\textbf{SPIR-V:} \texttt{OpControlBarrier / OpReturn}; \textbf{GLSL:} \texttt{// Scoreboard wait + warp reconvergence})
\end{tightitemize}
\clearpage

\subsection*{U16\_TO\_U32}
\addcontentsline{toc}{subsection}{U16\_TO\_U32}
Extracts the unsigned 16-bit halfword selected by \texttt{lane0} (\texttt{.h0} or \texttt{.h1}) of \texttt{src0}, zero-extends it to 32 bits, and writes the result to \texttt{dest}.
\par\textbf{Notes:}
\begin{tightitemize}
  \item Exact zero extension.
  \item \textbf{Source references:} \texttt{isa-v9-v14.xml}, \texttt{bifrost/valhall/ISA.xml}.
\end{tightitemize}
\codebox{\mnem{U16\_TO\_U32}\imod{\{.flow\}} \dst{dest}\imod{\{.dest\_width\}}, \src{src0}\imod{\{.lane0\}}}
\textbf{Operands:}
\begin{tightitemize}\footnotesize
  \item \texttt{\dst{dest}} (\texttt{dest\_register\_index}, \texttt{dest\_width}): 32-bit destination GPR receiving the converted value.
  \item \texttt{\src{src0}}: Unsigned 16-bit lane (\texttt{lane0}) zero-extended to 32-bit integer.
\end{tightitemize}
\par\textbf{Closest SPIR-V Op:} \texttt{OpUConvert (u16 -\textgreater{} u32)}
\par\textbf{Pseudo-GLSL (Operational Semantics):}
\codebox{dest = bitfieldExtract(uint(src0), lane0 * 16, 16);}
\textbf{Execution Context:}
\begin{tightitemize}\footnotesize
  \item \textbf{Compute stage preloads (\hyperref[sec:csf-spd]{\texttt{SPD\_k}}, \texttt{r16..r23}):} Hardware pre-fills \texttt{r55..r57} with \texttt{WorkGroupID.xyz + JOB\_OFFSET} (from \hyperref[sec:csf-regs]{\texttt{COMPUTESR::JOB\_OFFSET\_X/Y/Z}} \texttt{r34..r36}) and \texttt{r60..r62} with \texttt{LocalInvocationID.xyz} or \texttt{GlobalInvocationID.xyz} (from \texttt{WG\_SIZE} \texttt{r33} \& \texttt{JOB\_SIZE} \texttt{r37..r39}, selected by \texttt{ShaderProgram::preload}) so compute shaders can derive thread/element indices without memory loads.
  \item \textbf{Vertex stage preloads (\hyperref[sec:csf-spd]{\texttt{SPD\_0}/\texttt{SPD\_1}}, \texttt{r16..r19}):} The hardware Index Fetcher reads \hyperref[sec:csf-regs]{\texttt{IDVSSR::INDEX\_BUFFER}} (\texttt{r54:r55}) and pre-fills \texttt{r59, r60} with \texttt{VertexIndex} (\texttt{+ VERTEX\_OFFSET + GLOBAL\_ATTR\_OFF}) and \texttt{r61} with \texttt{InstanceIndex} (\texttt{+ INSTANCE\_OFFSET}) so Position and Varying shaders can pull vertex attributes directly.
  \item \textbf{Inline FAU \& Special Operands (\hyperref[sec:csf-tsd-fau]{\texttt{FAU\_k}}, \texttt{r8..r15}):} Any 8-bit source with \texttt{mode=2} reads zero-latency uniform \texttt{u0..u31} from \texttt{FAU\_k} (\texttt{slots64 = bits[63:56]}); \texttt{mode=3} reads \texttt{SmallConstant} (Table~1.2) or \texttt{FauSpecial} (\texttt{pc}, \texttt{lane\_id}, \texttt{lane\_width}, \texttt{tls\_ptr}, \texttt{wls\_ptr}, \texttt{blend\_descriptor\_0..7}).
\end{tightitemize}
\par\textbf{Worked Examples:}
\par\textbf{Example 1 (Zero-extend upper halfword):}
\codebox{\mnem{U16\_TO\_U32} \dst{r0}, \src{r1}\imod{.h1}}
With \texttt{r1 = 0xABCD1234}: extracts \texttt{h1} and zero-extends -\textgreater{} \texttt{r0 = 0x0000ABCD}.
\par\textbf{Example 2 (Zero-extend low halfword (masks upper bits)):}
\codebox{\mnem{U16\_TO\_U32} \dst{r0}, \src{r1}\imod{.h0}}
With \texttt{r1 = 0xFFFF8000}: zero-extends \texttt{0x8000} -\textgreater{} \texttt{r0 = 0x00008000} (not sign-extended).

\begin{center}\footnotesize
\begin{tabularx}{\linewidth}{@{}lllX@{}}
\toprule
Variant & Arch & Class & Opcode Fields \\
\midrule
\texttt{(none)} & v9--v14 & \texttt{R1} & \texttt{opcode2=U16\_TO\_U32, opcode1=CONVERT} \\
\texttt{(none)} & v9--v14 & \texttt{R1} & \texttt{opcode2=I16\_TO\_I32, opcode1=CONVERT} \\
\bottomrule
\end{tabularx}
\end{center}
\begin{center}\footnotesize
\begin{tabularx}{0.88\linewidth}{@{}l l X@{}}
\toprule
Field Name & Bits & Type / Fixed Value \\
\midrule
\texttt{src0} & \texttt{0:7} & \texttt{SourceEncoding} \\
\textit{Reserved} & \texttt{8:15} & \texttt{zero} \\
\texttt{opcode2} & \texttt{16:20} & \texttt{exact=U16\_TO\_U32} \\
\textit{Reserved} & \texttt{21:27} & \texttt{zero} \\
\hyperref[enum:lane16_m]{\texttt{lane0}} & \texttt{28} & \texttt{lane16\_m (def h0)} \\
\textit{Reserved} & \texttt{29:39} & \texttt{zero} \\
\texttt{dest\_register\_index} & \texttt{40:45} & \texttt{uint} \\
\hyperref[enum:dest_width_m]{\texttt{dest\_width}} & \texttt{46:47} & \texttt{dest\_width\_m} \\
\texttt{opcode1} & \texttt{48:51} & \texttt{exact=CONVERT} \\
\texttt{instruction\_class} & \texttt{52:56} & \texttt{exact=R1} \\
\texttt{fau\_page\_index} & \texttt{57:58} & \texttt{uint} \\
\hyperref[enum:flow_control_m]{\texttt{flow}} & \texttt{59:62} & \texttt{flow\_control\_m} \\
\textit{Reserved} & \texttt{63} & \texttt{zero} \\
\bottomrule
\end{tabularx}
\end{center}
\textbf{Fields Explained:}
\begin{tightitemize}\footnotesize
  \item \texttt{src0}: Primary source operand (\texttt{r0..r63}, \texttt{u0..u63} uniform, \texttt{k*} FAU constant, or \texttt{t0..t2} temporary; \texttt{\textasciicircum{}} marks discard/last-use). \emph{Example:} {\scriptsize\ttfamily \mnem{FADD.f32} \dst{r0}, \src{r1}, \src{r2}} (\textbf{SPIR-V:} \texttt{OpFAdd (Operand 0)}; \textbf{GLSL:} \texttt{float c = a + b;})
  \item \texttt{opcode2}: Secondary sub-opcode further specializing instructions sharing \texttt{instruction\_class} and \texttt{opcode1} (e.g., conversions, SFU tables, moves). \emph{Example:} {\scriptsize\ttfamily \mnem{F32\_TO\_S32}\imod{.round\_zero} \dst{r0}, \src{r1}} (\textbf{SPIR-V:} \texttt{OpConvertFToS}; \textbf{GLSL:} \texttt{int y = int(x);})
  \item \hyperref[enum:lane16_m]{\texttt{lane0}}: Stage A single-lane extractor on \texttt{src0} (\texttt{b0..b3} for 8-bit or \texttt{h0}/\texttt{h1} for 16-bit narrow conversions and \texttt{MKVEC}). \emph{Example:} {\scriptsize\ttfamily \mnem{F16\_TO\_F32} \dst{r0}, \src{r1}\imod{.h1}} (\textbf{SPIR-V:} \texttt{OpCompositeExtract}; \textbf{GLSL:} \texttt{float y = float(x.y);})
  \item \texttt{dest\_register\_index}: 6-bit destination GPR index (\texttt{r0..r63}) receiving the Stage C writeback result (even-aligned pair \texttt{[rN:rN+1]} for 64-bit ops). \emph{Example:} {\scriptsize\ttfamily \mnem{FMA.f32} \dst{r0}, \src{r1}, \src{r2}, \src{r3}} (\textbf{SPIR-V:} \texttt{SSA Result ID}; \textbf{GLSL:} \texttt{float r0 = fma(r1, r2, r3);})
  \item \hyperref[enum:dest_width_m]{\texttt{dest\_width}}: Stage C destination lane write-mask (\texttt{none} = full 32-bit, \texttt{h0} = low 16-bit, \texttt{h1} = high 16-bit without bit-shifting). \emph{Example:} {\scriptsize\ttfamily \mnem{FADD.v2f16} \dst{r0}\imod{.h1}, \src{r1}, \src{r2}} (\textbf{SPIR-V:} \texttt{OpVectorInsertDynamic / WriteMask}; \textbf{GLSL:} \texttt{dst.y = a.y + b.y;})
  \item \texttt{opcode1}: Primary sub-opcode within the instruction class uniquely identifying the operation family. \emph{Example:} {\scriptsize\ttfamily \mnem{IADD.s32} \dst{r0}, \src{r1}, \src{r2}} (\textbf{SPIR-V:} \texttt{OpIAdd (operation selector)}; \textbf{GLSL:} \texttt{// Primary opcode bits})
  \item \texttt{instruction\_class}: 4-bit major instruction class selector (\texttt{bits 48:51}) partitioning the 64-bit Valhall opcode space (e.g., \texttt{FMA}, \texttt{CVT}, \texttt{SFU}, \texttt{LS}, \texttt{TEX}, \texttt{BRANCH}). \emph{Example:} {\scriptsize\ttfamily \mnem{FADD.f32} \dst{r0}, \src{r1}, \src{r2}} (\textbf{SPIR-V:} \texttt{N/A (encoding class)}; \textbf{GLSL:} \texttt{// Hardware execution pipe / decode class})
  \item \texttt{fau\_page\_index}: Selects the FAU (Fast Access Uniform) page (\texttt{0..3}) when a source operand references hardware small constants, thread IDs, or lane constants. \emph{Example:} {\scriptsize\ttfamily \mnem{FADD.f32} \dst{r0}, \src{r1}, \src{k1f}} (\textbf{SPIR-V:} \texttt{OpConstant / BuiltIn}; \textbf{GLSL:} \texttt{float y = x + 1.0;})
  \item \hyperref[enum:flow_control_m]{\texttt{flow}}: 4-bit embedded flow-control \& scoreboard wait modifier (\texttt{bits 59:62}: \texttt{none}, \texttt{wait0..wait0126}, \texttt{wait}, \texttt{reconverge}, \texttt{discard}, \texttt{end}). \emph{Example:} {\scriptsize\ttfamily \mnem{FADD.f32}\imod{.wait0.reconverge} \dst{r0}, \src{r1}, \src{r2}} (\textbf{SPIR-V:} \texttt{OpControlBarrier / OpReturn}; \textbf{GLSL:} \texttt{// Scoreboard wait + warp reconvergence})
\end{tightitemize}
\clearpage

\subsection*{U32\_TO\_F32}
\addcontentsline{toc}{subsection}{U32\_TO\_F32}
Converts unsigned 32-bit integer \texttt{src0} (with optional Stage A widening) to IEEE-754 FP32 in \texttt{dest}, applying the selected \texttt{round} mode (default RNE; no clamp modifier).
\par\textbf{Notes:}
\begin{tightitemize}
  \item Supports all 5 rounding modes (\texttt{round\_extended\_m}).
  \item \textbf{Source references:} \texttt{kraid/ops.rs:2575-2600}, \texttt{encode\_v9.rs:1960-2010}, \texttt{nir.rs:690-701}.
\end{tightitemize}
\codebox{\mnem{U32\_TO\_F32}\imod{\{.flow\}}\imod{\{.round\}} \dst{dest}\imod{\{.dest\_width\}}, \src{src0}\\\mnem{U32\_TO\_F32}\imod{\{.flow\}}\imod{\{.round\}} \dst{dest}\imod{\{.dest\_width\}}, \src{src0}\imod{\{.widen0\}}}
\textbf{Operands:}
\begin{tightitemize}\footnotesize
  \item \texttt{\dst{dest}} (\texttt{dest\_register\_index}, \texttt{dest\_width}): 32-bit destination GPR receiving the converted value.
  \item \texttt{\src{src0}}: Unsigned 32-bit integer converted to FP32 (rounded per \texttt{round}).
\end{tightitemize}
\par\textbf{Closest SPIR-V Op:} \texttt{OpConvertUToF}
\par\textbf{Pseudo-GLSL (Operational Semantics):}
\codebox{dest = round\_rm(float(uint(widen0(src0))));}
\textbf{Execution Context:}
\begin{tightitemize}\footnotesize
  \item \textbf{Compute stage preloads (\hyperref[sec:csf-spd]{\texttt{SPD\_k}}, \texttt{r16..r23}):} Hardware pre-fills \texttt{r55..r57} with \texttt{WorkGroupID.xyz + JOB\_OFFSET} (from \hyperref[sec:csf-regs]{\texttt{COMPUTESR::JOB\_OFFSET\_X/Y/Z}} \texttt{r34..r36}) and \texttt{r60..r62} with \texttt{LocalInvocationID.xyz} or \texttt{GlobalInvocationID.xyz} (from \texttt{WG\_SIZE} \texttt{r33} \& \texttt{JOB\_SIZE} \texttt{r37..r39}, selected by \texttt{ShaderProgram::preload}) so compute shaders can derive thread/element indices without memory loads.
  \item \textbf{Vertex stage preloads (\hyperref[sec:csf-spd]{\texttt{SPD\_0}/\texttt{SPD\_1}}, \texttt{r16..r19}):} The hardware Index Fetcher reads \hyperref[sec:csf-regs]{\texttt{IDVSSR::INDEX\_BUFFER}} (\texttt{r54:r55}) and pre-fills \texttt{r59, r60} with \texttt{VertexIndex} (\texttt{+ VERTEX\_OFFSET + GLOBAL\_ATTR\_OFF}) and \texttt{r61} with \texttt{InstanceIndex} (\texttt{+ INSTANCE\_OFFSET}) so Position and Varying shaders can pull vertex attributes directly.
  \item \textbf{Inline FAU \& Special Operands (\hyperref[sec:csf-tsd-fau]{\texttt{FAU\_k}}, \texttt{r8..r15}):} Any 8-bit source with \texttt{mode=2} reads zero-latency uniform \texttt{u0..u31} from \texttt{FAU\_k} (\texttt{slots64 = bits[63:56]}); \texttt{mode=3} reads \texttt{SmallConstant} (Table~1.2) or \texttt{FauSpecial} (\texttt{pc}, \texttt{lane\_id}, \texttt{lane\_width}, \texttt{tls\_ptr}, \texttt{wls\_ptr}, \texttt{blend\_descriptor\_0..7}).
\end{tightitemize}
\par\textbf{Worked Examples:}
\par\textbf{Example 1 (Exact unsigned integer):}
\codebox{\mnem{U32\_TO\_F32} \dst{r0}, \src{r1}}
With \texttt{r1 = 100}: returns \texttt{100.0f} (\texttt{0x42C80000}).
\par\textbf{Example 2 (`UINT32\_MAX` rounding under RNE vs RTZ):}
\codebox{\mnem{U32\_TO\_F32}\imod{.round\_zero} \dst{r0}, \src{r1}}
With \texttt{r1 = 0xFFFFFFFF}: default RNE rounds up to \texttt{2\textasciicircum{}32} (\texttt{0x4F800000}), whereas \texttt{.round\_zero} rounds down to \texttt{0x4F7FFFFF}.
\par\textbf{Mined SPIR-V $\to$ Valhall ISA Example (Converting unsigned integer to float --- Fragment shader, Mali-G710 (v10)):}
\codebox{\textbf{; SPIR-V (input)}\\\%f = OpConvertUToF \%float \%uval\\[3pt]\textbf{; Valhall ISA (malioc output)}\\F32\_TO\_S32.round\_zero r2, r2\textasciicircum{}\\U32\_TO\_F32 r2, r2\textasciicircum{}\\FMUL.f32 r2, r2\textasciicircum{}, u3.w0}
\texttt{OpConvertUToF} lowers 1:1 to \texttt{U32\_TO\_F32} with default round-to-nearest-even.
\begin{center}\footnotesize
\begin{tabularx}{\linewidth}{@{}lllX@{}}
\toprule
Variant & Arch & Class & Opcode Fields \\
\midrule
\texttt{(none)} & v9--v14 & \texttt{R1} & \texttt{opcode2=U32\_TO\_F32, opcode1=CONVERT} \\
\texttt{(none)} & v9--v14 & \texttt{R1} & \texttt{opcode2=U32\_TO\_F32, opcode1=CONVERT} \\
\texttt{(none)} & v9--v14 & \texttt{R1} & \texttt{opcode2=I32\_TO\_F32, opcode1=CONVERT} \\
\bottomrule
\end{tabularx}
\end{center}
\begin{center}\footnotesize
\begin{tabularx}{0.88\linewidth}{@{}l l X@{}}
\toprule
Field Name & Bits & Type / Fixed Value \\
\midrule
\texttt{src0} & \texttt{0:7} & \texttt{SourceEncoding} \\
\textit{Reserved} & \texttt{8:15} & \texttt{zero} \\
\texttt{opcode2} & \texttt{16:20} & \texttt{exact=U32\_TO\_F32} \\
\textit{Reserved} & \texttt{21:29} & \texttt{zero} \\
\hyperref[enum:round_extended_m]{\texttt{round}} & \texttt{30:32} & \texttt{round\_extended\_m (def none)} \\
\textit{Reserved} & \texttt{33} & \texttt{zero} \\
\textit{Reserved} & \texttt{34:35} & \texttt{ignore} \\
\hyperref[enum:widen_int_m]{\texttt{widen0}} & \texttt{36:39} & \texttt{exact=none} \\
\texttt{dest\_register\_index} & \texttt{40:45} & \texttt{uint} \\
\hyperref[enum:dest_width_m]{\texttt{dest\_width}} & \texttt{46:47} & \texttt{dest\_width\_m} \\
\texttt{opcode1} & \texttt{48:51} & \texttt{exact=CONVERT} \\
\texttt{instruction\_class} & \texttt{52:56} & \texttt{exact=R1} \\
\texttt{fau\_page\_index} & \texttt{57:58} & \texttt{uint} \\
\hyperref[enum:flow_control_m]{\texttt{flow}} & \texttt{59:62} & \texttt{flow\_control\_m} \\
\textit{Reserved} & \texttt{63} & \texttt{zero} \\
\bottomrule
\end{tabularx}
\end{center}
\textbf{Fields Explained:}
\begin{tightitemize}\footnotesize
  \item \texttt{src0}: Primary source operand (\texttt{r0..r63}, \texttt{u0..u63} uniform, \texttt{k*} FAU constant, or \texttt{t0..t2} temporary; \texttt{\textasciicircum{}} marks discard/last-use). \emph{Example:} {\scriptsize\ttfamily \mnem{FADD.f32} \dst{r0}, \src{r1}, \src{r2}} (\textbf{SPIR-V:} \texttt{OpFAdd (Operand 0)}; \textbf{GLSL:} \texttt{float c = a + b;})
  \item \texttt{opcode2}: Secondary sub-opcode further specializing instructions sharing \texttt{instruction\_class} and \texttt{opcode1} (e.g., conversions, SFU tables, moves). \emph{Example:} {\scriptsize\ttfamily \mnem{F32\_TO\_S32}\imod{.round\_zero} \dst{r0}, \src{r1}} (\textbf{SPIR-V:} \texttt{OpConvertFToS}; \textbf{GLSL:} \texttt{int y = int(x);})
  \item \hyperref[enum:round_extended_m]{\texttt{round}}: Rounding mode (\texttt{none}/RNE, \texttt{round\_up}, \texttt{round\_down}, \texttt{round\_zero}, \texttt{round\_na}) applied in FP ALU, conversions, \texttt{HADD}, \texttt{LDEXP}, and Stage C narrowing. \emph{Example:} {\scriptsize\ttfamily \mnem{FROUND.f32}\imod{.round\_down} \dst{r0}, \src{r1}} (\textbf{SPIR-V:} \texttt{GLSL.std.450 Floor / Ceil / Trunc}; \textbf{GLSL:} \texttt{float y = floor(x);})
  \item \hyperref[enum:widen_int_m]{\texttt{widen0}}: Stage A widening extractor on \texttt{src0} (promotes FP16 \texttt{h0}/\texttt{h1} to FP32, or sign/zero-extends \texttt{b0..b3}/\texttt{h0..h1}/\texttt{w0..w1} to 32/64-bit before \texttt{abs}/\texttt{neg}). \emph{Example:} {\scriptsize\ttfamily \mnem{IADD.s32} \dst{r0}, \src{r1}\imod{.h0}, \src{r2}} (\textbf{SPIR-V:} \texttt{OpSConvert / OpFConvert}; \textbf{GLSL:} \texttt{int c = int(int16\_t(a)) + b;})
  \item \texttt{dest\_register\_index}: 6-bit destination GPR index (\texttt{r0..r63}) receiving the Stage C writeback result (even-aligned pair \texttt{[rN:rN+1]} for 64-bit ops). \emph{Example:} {\scriptsize\ttfamily \mnem{FMA.f32} \dst{r0}, \src{r1}, \src{r2}, \src{r3}} (\textbf{SPIR-V:} \texttt{SSA Result ID}; \textbf{GLSL:} \texttt{float r0 = fma(r1, r2, r3);})
  \item \hyperref[enum:dest_width_m]{\texttt{dest\_width}}: Stage C destination lane write-mask (\texttt{none} = full 32-bit, \texttt{h0} = low 16-bit, \texttt{h1} = high 16-bit without bit-shifting). \emph{Example:} {\scriptsize\ttfamily \mnem{FADD.v2f16} \dst{r0}\imod{.h1}, \src{r1}, \src{r2}} (\textbf{SPIR-V:} \texttt{OpVectorInsertDynamic / WriteMask}; \textbf{GLSL:} \texttt{dst.y = a.y + b.y;})
  \item \texttt{opcode1}: Primary sub-opcode within the instruction class uniquely identifying the operation family. \emph{Example:} {\scriptsize\ttfamily \mnem{IADD.s32} \dst{r0}, \src{r1}, \src{r2}} (\textbf{SPIR-V:} \texttt{OpIAdd (operation selector)}; \textbf{GLSL:} \texttt{// Primary opcode bits})
  \item \texttt{instruction\_class}: 4-bit major instruction class selector (\texttt{bits 48:51}) partitioning the 64-bit Valhall opcode space (e.g., \texttt{FMA}, \texttt{CVT}, \texttt{SFU}, \texttt{LS}, \texttt{TEX}, \texttt{BRANCH}). \emph{Example:} {\scriptsize\ttfamily \mnem{FADD.f32} \dst{r0}, \src{r1}, \src{r2}} (\textbf{SPIR-V:} \texttt{N/A (encoding class)}; \textbf{GLSL:} \texttt{// Hardware execution pipe / decode class})
  \item \texttt{fau\_page\_index}: Selects the FAU (Fast Access Uniform) page (\texttt{0..3}) when a source operand references hardware small constants, thread IDs, or lane constants. \emph{Example:} {\scriptsize\ttfamily \mnem{FADD.f32} \dst{r0}, \src{r1}, \src{k1f}} (\textbf{SPIR-V:} \texttt{OpConstant / BuiltIn}; \textbf{GLSL:} \texttt{float y = x + 1.0;})
  \item \hyperref[enum:flow_control_m]{\texttt{flow}}: 4-bit embedded flow-control \& scoreboard wait modifier (\texttt{bits 59:62}: \texttt{none}, \texttt{wait0..wait0126}, \texttt{wait}, \texttt{reconverge}, \texttt{discard}, \texttt{end}). \emph{Example:} {\scriptsize\ttfamily \mnem{FADD.f32}\imod{.wait0.reconverge} \dst{r0}, \src{r1}, \src{r2}} (\textbf{SPIR-V:} \texttt{OpControlBarrier / OpReturn}; \textbf{GLSL:} \texttt{// Scoreboard wait + warp reconvergence})
\end{tightitemize}
\clearpage

\subsection*{U8\_TO\_F32}
\addcontentsline{toc}{subsection}{U8\_TO\_F32}
Converts the unsigned 8-bit byte selected by \texttt{lane0} (\texttt{.b0}..\texttt{.b3}) of \texttt{src0} exactly to IEEE-754 FP32 and writes the result to \texttt{dest} (arch \textless{}= v10).
\par\textbf{Notes:}
\begin{tightitemize}
  \item Arch v9-v10 only (\texttt{kraid/nir.rs:692}).
  \item \textbf{Source references:} \texttt{kraid/ops.rs:2575-2600}, \texttt{encode\_v9.rs:1988-1991}, \texttt{nir.rs:690-701}.
\end{tightitemize}
\codebox{\mnem{U8\_TO\_F32}\imod{\{.flow\}} \dst{dest}\imod{\{.dest\_width\}}, \src{src0}\imod{\{.lane0\}}}
\textbf{Operands:}
\begin{tightitemize}\footnotesize
  \item \texttt{\dst{dest}} (\texttt{dest\_register\_index}, \texttt{dest\_width}): 32-bit destination GPR receiving the converted value.
  \item \texttt{\src{src0}}: Unsigned 8-bit byte lane (\texttt{lane0}) converted exactly to FP32.
\end{tightitemize}
\par\textbf{Closest SPIR-V Op:} \texttt{OpConvertUToF / UnpackUnorm4x8 (v9--v10)}
\par\textbf{Pseudo-GLSL (Operational Semantics):}
\codebox{dest = float(bitfieldExtract(uint(src0), lane0 * 8, 8));}
\textbf{Execution Context:}
\begin{tightitemize}\footnotesize
  \item \textbf{Compute stage preloads (\hyperref[sec:csf-spd]{\texttt{SPD\_k}}, \texttt{r16..r23}):} Hardware pre-fills \texttt{r55..r57} with \texttt{WorkGroupID.xyz + JOB\_OFFSET} (from \hyperref[sec:csf-regs]{\texttt{COMPUTESR::JOB\_OFFSET\_X/Y/Z}} \texttt{r34..r36}) and \texttt{r60..r62} with \texttt{LocalInvocationID.xyz} or \texttt{GlobalInvocationID.xyz} (from \texttt{WG\_SIZE} \texttt{r33} \& \texttt{JOB\_SIZE} \texttt{r37..r39}, selected by \texttt{ShaderProgram::preload}) so compute shaders can derive thread/element indices without memory loads.
  \item \textbf{Vertex stage preloads (\hyperref[sec:csf-spd]{\texttt{SPD\_0}/\texttt{SPD\_1}}, \texttt{r16..r19}):} The hardware Index Fetcher reads \hyperref[sec:csf-regs]{\texttt{IDVSSR::INDEX\_BUFFER}} (\texttt{r54:r55}) and pre-fills \texttt{r59, r60} with \texttt{VertexIndex} (\texttt{+ VERTEX\_OFFSET + GLOBAL\_ATTR\_OFF}) and \texttt{r61} with \texttt{InstanceIndex} (\texttt{+ INSTANCE\_OFFSET}) so Position and Varying shaders can pull vertex attributes directly.
  \item \textbf{Inline FAU \& Special Operands (\hyperref[sec:csf-tsd-fau]{\texttt{FAU\_k}}, \texttt{r8..r15}):} Any 8-bit source with \texttt{mode=2} reads zero-latency uniform \texttt{u0..u31} from \texttt{FAU\_k} (\texttt{slots64 = bits[63:56]}); \texttt{mode=3} reads \texttt{SmallConstant} (Table~1.2) or \texttt{FauSpecial} (\texttt{pc}, \texttt{lane\_id}, \texttt{lane\_width}, \texttt{tls\_ptr}, \texttt{wls\_ptr}, \texttt{blend\_descriptor\_0..7}).
\end{tightitemize}
\par\textbf{Worked Examples:}
\par\textbf{Example 1 (Byte 3 unsigned -\textgreater{} FP32 (v9--v10)):}
\codebox{\mnem{U8\_TO\_F32} \dst{r0}, \src{r1}\imod{.b3}}
With \texttt{r1 = 0xFF000000}: converts byte 3 (\texttt{255u}) -\textgreater{} \texttt{r0 = 255.0f} (\texttt{0x437F0000}).
\par\textbf{Example 2 (Byte 0 zero conversion):}
\codebox{\mnem{U8\_TO\_F32} \dst{r0}, \src{r1}\imod{.b0}}
With \texttt{r1 = 0xFFFFFF00}: converts byte 0 (\texttt{0u}) -\textgreater{} \texttt{r0 = +0.0f} (\texttt{0x00000000}).
\par\textbf{Mined SPIR-V $\to$ Valhall ISA Example (Unpacking a RGBA8\_UNORM uint32 word into 4 float channels --- Vertex shader, Mali-G710 (v10)):}
\codebox{\textbf{; SPIR-V (input)}\\\%b0 = OpBitFieldUExtract \%uint \%packed \%uint\_0 \%uint\_8\\\%r  = OpFMul \%float (OpConvertUToF \%float \%b0) \%inv\_255\\\%b1 = OpBitFieldUExtract \%uint \%packed \%uint\_8 \%uint\_8\\\%g  = OpFMul \%float (OpConvertUToF \%float \%b1) \%inv\_255\\[3pt]\textbf{; Valhall ISA (malioc output)}\\LD\_PKA.i32.slot0.wait0 @r0, r0\textasciicircum{}, u12.w1\\U8\_TO\_F32 r1, r0.b0\\U8\_TO\_F32 r62, r0.b1\\U8\_TO\_F32 r63, r0.b2\\U8\_TO\_F32 r61, r0\textasciicircum{}.b3\\FMUL.f32 r7, r1\textasciicircum{}, u13.w0          ; multiply by 1/255}
\texttt{malioc} fuses each \texttt{OpBitFieldUExtract(..., 8)} + \texttt{OpConvertUToF} into a single \texttt{U8\_TO\_F32} selecting bytes \texttt{.b0} through \texttt{.b3} of the loaded 32-bit word.
\begin{center}\footnotesize
\begin{tabularx}{\linewidth}{@{}lllX@{}}
\toprule
Variant & Arch & Class & Opcode Fields \\
\midrule
\texttt{(none)} & v9--v14 & \texttt{R1} & \texttt{opcode2=U8\_TO\_F32, opcode1=CONVERT} \\
\bottomrule
\end{tabularx}
\end{center}
\begin{center}\footnotesize
\begin{tabularx}{0.88\linewidth}{@{}l l X@{}}
\toprule
Field Name & Bits & Type / Fixed Value \\
\midrule
\texttt{src0} & \texttt{0:7} & \texttt{SourceEncoding} \\
\textit{Reserved} & \texttt{8:15} & \texttt{zero} \\
\texttt{opcode2} & \texttt{16:20} & \texttt{exact=U8\_TO\_F32} \\
\textit{Reserved} & \texttt{21:27} & \texttt{zero} \\
\hyperref[enum:lane8_m]{\texttt{lane0}} & \texttt{28:29} & \texttt{lane8\_m (def b0)} \\
\textit{Reserved} & \texttt{30:39} & \texttt{zero} \\
\texttt{dest\_register\_index} & \texttt{40:45} & \texttt{uint} \\
\hyperref[enum:dest_width_m]{\texttt{dest\_width}} & \texttt{46:47} & \texttt{dest\_width\_m} \\
\texttt{opcode1} & \texttt{48:51} & \texttt{exact=CONVERT} \\
\texttt{instruction\_class} & \texttt{52:56} & \texttt{exact=R1} \\
\texttt{fau\_page\_index} & \texttt{57:58} & \texttt{uint} \\
\hyperref[enum:flow_control_m]{\texttt{flow}} & \texttt{59:62} & \texttt{flow\_control\_m} \\
\textit{Reserved} & \texttt{63} & \texttt{zero} \\
\bottomrule
\end{tabularx}
\end{center}
\textbf{Fields Explained:}
\begin{tightitemize}\footnotesize
  \item \texttt{src0}: Primary source operand (\texttt{r0..r63}, \texttt{u0..u63} uniform, \texttt{k*} FAU constant, or \texttt{t0..t2} temporary; \texttt{\textasciicircum{}} marks discard/last-use). \emph{Example:} {\scriptsize\ttfamily \mnem{FADD.f32} \dst{r0}, \src{r1}, \src{r2}} (\textbf{SPIR-V:} \texttt{OpFAdd (Operand 0)}; \textbf{GLSL:} \texttt{float c = a + b;})
  \item \texttt{opcode2}: Secondary sub-opcode further specializing instructions sharing \texttt{instruction\_class} and \texttt{opcode1} (e.g., conversions, SFU tables, moves). \emph{Example:} {\scriptsize\ttfamily \mnem{F32\_TO\_S32}\imod{.round\_zero} \dst{r0}, \src{r1}} (\textbf{SPIR-V:} \texttt{OpConvertFToS}; \textbf{GLSL:} \texttt{int y = int(x);})
  \item \hyperref[enum:lane8_m]{\texttt{lane0}}: Stage A single-lane extractor on \texttt{src0} (\texttt{b0..b3} for 8-bit or \texttt{h0}/\texttt{h1} for 16-bit narrow conversions and \texttt{MKVEC}). \emph{Example:} {\scriptsize\ttfamily \mnem{F16\_TO\_F32} \dst{r0}, \src{r1}\imod{.h1}} (\textbf{SPIR-V:} \texttt{OpCompositeExtract}; \textbf{GLSL:} \texttt{float y = float(x.y);})
  \item \texttt{dest\_register\_index}: 6-bit destination GPR index (\texttt{r0..r63}) receiving the Stage C writeback result (even-aligned pair \texttt{[rN:rN+1]} for 64-bit ops). \emph{Example:} {\scriptsize\ttfamily \mnem{FMA.f32} \dst{r0}, \src{r1}, \src{r2}, \src{r3}} (\textbf{SPIR-V:} \texttt{SSA Result ID}; \textbf{GLSL:} \texttt{float r0 = fma(r1, r2, r3);})
  \item \hyperref[enum:dest_width_m]{\texttt{dest\_width}}: Stage C destination lane write-mask (\texttt{none} = full 32-bit, \texttt{h0} = low 16-bit, \texttt{h1} = high 16-bit without bit-shifting). \emph{Example:} {\scriptsize\ttfamily \mnem{FADD.v2f16} \dst{r0}\imod{.h1}, \src{r1}, \src{r2}} (\textbf{SPIR-V:} \texttt{OpVectorInsertDynamic / WriteMask}; \textbf{GLSL:} \texttt{dst.y = a.y + b.y;})
  \item \texttt{opcode1}: Primary sub-opcode within the instruction class uniquely identifying the operation family. \emph{Example:} {\scriptsize\ttfamily \mnem{IADD.s32} \dst{r0}, \src{r1}, \src{r2}} (\textbf{SPIR-V:} \texttt{OpIAdd (operation selector)}; \textbf{GLSL:} \texttt{// Primary opcode bits})
  \item \texttt{instruction\_class}: 4-bit major instruction class selector (\texttt{bits 48:51}) partitioning the 64-bit Valhall opcode space (e.g., \texttt{FMA}, \texttt{CVT}, \texttt{SFU}, \texttt{LS}, \texttt{TEX}, \texttt{BRANCH}). \emph{Example:} {\scriptsize\ttfamily \mnem{FADD.f32} \dst{r0}, \src{r1}, \src{r2}} (\textbf{SPIR-V:} \texttt{N/A (encoding class)}; \textbf{GLSL:} \texttt{// Hardware execution pipe / decode class})
  \item \texttt{fau\_page\_index}: Selects the FAU (Fast Access Uniform) page (\texttt{0..3}) when a source operand references hardware small constants, thread IDs, or lane constants. \emph{Example:} {\scriptsize\ttfamily \mnem{FADD.f32} \dst{r0}, \src{r1}, \src{k1f}} (\textbf{SPIR-V:} \texttt{OpConstant / BuiltIn}; \textbf{GLSL:} \texttt{float y = x + 1.0;})
  \item \hyperref[enum:flow_control_m]{\texttt{flow}}: 4-bit embedded flow-control \& scoreboard wait modifier (\texttt{bits 59:62}: \texttt{none}, \texttt{wait0..wait0126}, \texttt{wait}, \texttt{reconverge}, \texttt{discard}, \texttt{end}). \emph{Example:} {\scriptsize\ttfamily \mnem{FADD.f32}\imod{.wait0.reconverge} \dst{r0}, \src{r1}, \src{r2}} (\textbf{SPIR-V:} \texttt{OpControlBarrier / OpReturn}; \textbf{GLSL:} \texttt{// Scoreboard wait + warp reconvergence})
\end{tightitemize}
\clearpage

\subsection*{U8\_TO\_U32}
\addcontentsline{toc}{subsection}{U8\_TO\_U32}
Extracts the unsigned 8-bit byte selected by \texttt{lane0} (\texttt{.b0}..\texttt{.b3}) of \texttt{src0}, zero-extends it to 32 bits, and writes the result to \texttt{dest}.
\par\textbf{Notes:}
\begin{tightitemize}
  \item Exact zero extension.
  \item \textbf{Source references:} \texttt{isa-v9-v14.xml}, \texttt{bifrost/valhall/ISA.xml}.
\end{tightitemize}
\codebox{\mnem{U8\_TO\_U32}\imod{\{.flow\}} \dst{dest}\imod{\{.dest\_width\}}, \src{src0}\imod{\{.lane0\}}}
\textbf{Operands:}
\begin{tightitemize}\footnotesize
  \item \texttt{\dst{dest}} (\texttt{dest\_register\_index}, \texttt{dest\_width}): 32-bit destination GPR receiving the converted value.
  \item \texttt{\src{src0}}: Unsigned 8-bit byte lane (\texttt{lane0}) zero-extended to 32-bit integer.
\end{tightitemize}
\par\textbf{Closest SPIR-V Op:} \texttt{OpUConvert (u8 -\textgreater{} u32)}
\par\textbf{Pseudo-GLSL (Operational Semantics):}
\codebox{dest = bitfieldExtract(uint(src0), lane0 * 8, 8);}
\textbf{Execution Context:}
\begin{tightitemize}\footnotesize
  \item \textbf{Compute stage preloads (\hyperref[sec:csf-spd]{\texttt{SPD\_k}}, \texttt{r16..r23}):} Hardware pre-fills \texttt{r55..r57} with \texttt{WorkGroupID.xyz + JOB\_OFFSET} (from \hyperref[sec:csf-regs]{\texttt{COMPUTESR::JOB\_OFFSET\_X/Y/Z}} \texttt{r34..r36}) and \texttt{r60..r62} with \texttt{LocalInvocationID.xyz} or \texttt{GlobalInvocationID.xyz} (from \texttt{WG\_SIZE} \texttt{r33} \& \texttt{JOB\_SIZE} \texttt{r37..r39}, selected by \texttt{ShaderProgram::preload}) so compute shaders can derive thread/element indices without memory loads.
  \item \textbf{Vertex stage preloads (\hyperref[sec:csf-spd]{\texttt{SPD\_0}/\texttt{SPD\_1}}, \texttt{r16..r19}):} The hardware Index Fetcher reads \hyperref[sec:csf-regs]{\texttt{IDVSSR::INDEX\_BUFFER}} (\texttt{r54:r55}) and pre-fills \texttt{r59, r60} with \texttt{VertexIndex} (\texttt{+ VERTEX\_OFFSET + GLOBAL\_ATTR\_OFF}) and \texttt{r61} with \texttt{InstanceIndex} (\texttt{+ INSTANCE\_OFFSET}) so Position and Varying shaders can pull vertex attributes directly.
  \item \textbf{Inline FAU \& Special Operands (\hyperref[sec:csf-tsd-fau]{\texttt{FAU\_k}}, \texttt{r8..r15}):} Any 8-bit source with \texttt{mode=2} reads zero-latency uniform \texttt{u0..u31} from \texttt{FAU\_k} (\texttt{slots64 = bits[63:56]}); \texttt{mode=3} reads \texttt{SmallConstant} (Table~1.2) or \texttt{FauSpecial} (\texttt{pc}, \texttt{lane\_id}, \texttt{lane\_width}, \texttt{tls\_ptr}, \texttt{wls\_ptr}, \texttt{blend\_descriptor\_0..7}).
\end{tightitemize}
\par\textbf{Worked Examples:}
\par\textbf{Example 1 (Zero-extend byte 1):}
\codebox{\mnem{U8\_TO\_U32} \dst{r0}, \src{r1}\imod{.b1}}
With \texttt{r1 = 0x0000AB00}: zero-extends byte 1 (\texttt{0xAB}) -\textgreater{} \texttt{r0 = 0x000000AB}.
\par\textbf{Example 2 (Zero-extend byte 3 (MSB set)):}
\codebox{\mnem{U8\_TO\_U32} \dst{r0}, \src{r1}\imod{.b3}}
With \texttt{r1 = 0xFE000000}: zero-extends \texttt{0xFE} -\textgreater{} \texttt{r0 = 0x000000FE} (\texttt{254}, unlike \texttt{S8\_TO\_S32}).

\begin{center}\footnotesize
\begin{tabularx}{\linewidth}{@{}lllX@{}}
\toprule
Variant & Arch & Class & Opcode Fields \\
\midrule
\texttt{(none)} & v9--v14 & \texttt{R1} & \texttt{opcode2=U8\_TO\_U32, opcode1=CONVERT} \\
\texttt{(none)} & v9--v14 & \texttt{R1} & \texttt{opcode2=I8\_TO\_I32, opcode1=CONVERT} \\
\bottomrule
\end{tabularx}
\end{center}
\begin{center}\footnotesize
\begin{tabularx}{0.88\linewidth}{@{}l l X@{}}
\toprule
Field Name & Bits & Type / Fixed Value \\
\midrule
\texttt{src0} & \texttt{0:7} & \texttt{SourceEncoding} \\
\textit{Reserved} & \texttt{8:15} & \texttt{zero} \\
\texttt{opcode2} & \texttt{16:20} & \texttt{exact=U8\_TO\_U32} \\
\textit{Reserved} & \texttt{21:27} & \texttt{zero} \\
\hyperref[enum:lane8_m]{\texttt{lane0}} & \texttt{28:29} & \texttt{lane8\_m (def b0)} \\
\textit{Reserved} & \texttt{30:39} & \texttt{zero} \\
\texttt{dest\_register\_index} & \texttt{40:45} & \texttt{uint} \\
\hyperref[enum:dest_width_m]{\texttt{dest\_width}} & \texttt{46:47} & \texttt{dest\_width\_m} \\
\texttt{opcode1} & \texttt{48:51} & \texttt{exact=CONVERT} \\
\texttt{instruction\_class} & \texttt{52:56} & \texttt{exact=R1} \\
\texttt{fau\_page\_index} & \texttt{57:58} & \texttt{uint} \\
\hyperref[enum:flow_control_m]{\texttt{flow}} & \texttt{59:62} & \texttt{flow\_control\_m} \\
\textit{Reserved} & \texttt{63} & \texttt{zero} \\
\bottomrule
\end{tabularx}
\end{center}
\textbf{Fields Explained:}
\begin{tightitemize}\footnotesize
  \item \texttt{src0}: Primary source operand (\texttt{r0..r63}, \texttt{u0..u63} uniform, \texttt{k*} FAU constant, or \texttt{t0..t2} temporary; \texttt{\textasciicircum{}} marks discard/last-use). \emph{Example:} {\scriptsize\ttfamily \mnem{FADD.f32} \dst{r0}, \src{r1}, \src{r2}} (\textbf{SPIR-V:} \texttt{OpFAdd (Operand 0)}; \textbf{GLSL:} \texttt{float c = a + b;})
  \item \texttt{opcode2}: Secondary sub-opcode further specializing instructions sharing \texttt{instruction\_class} and \texttt{opcode1} (e.g., conversions, SFU tables, moves). \emph{Example:} {\scriptsize\ttfamily \mnem{F32\_TO\_S32}\imod{.round\_zero} \dst{r0}, \src{r1}} (\textbf{SPIR-V:} \texttt{OpConvertFToS}; \textbf{GLSL:} \texttt{int y = int(x);})
  \item \hyperref[enum:lane8_m]{\texttt{lane0}}: Stage A single-lane extractor on \texttt{src0} (\texttt{b0..b3} for 8-bit or \texttt{h0}/\texttt{h1} for 16-bit narrow conversions and \texttt{MKVEC}). \emph{Example:} {\scriptsize\ttfamily \mnem{F16\_TO\_F32} \dst{r0}, \src{r1}\imod{.h1}} (\textbf{SPIR-V:} \texttt{OpCompositeExtract}; \textbf{GLSL:} \texttt{float y = float(x.y);})
  \item \texttt{dest\_register\_index}: 6-bit destination GPR index (\texttt{r0..r63}) receiving the Stage C writeback result (even-aligned pair \texttt{[rN:rN+1]} for 64-bit ops). \emph{Example:} {\scriptsize\ttfamily \mnem{FMA.f32} \dst{r0}, \src{r1}, \src{r2}, \src{r3}} (\textbf{SPIR-V:} \texttt{SSA Result ID}; \textbf{GLSL:} \texttt{float r0 = fma(r1, r2, r3);})
  \item \hyperref[enum:dest_width_m]{\texttt{dest\_width}}: Stage C destination lane write-mask (\texttt{none} = full 32-bit, \texttt{h0} = low 16-bit, \texttt{h1} = high 16-bit without bit-shifting). \emph{Example:} {\scriptsize\ttfamily \mnem{FADD.v2f16} \dst{r0}\imod{.h1}, \src{r1}, \src{r2}} (\textbf{SPIR-V:} \texttt{OpVectorInsertDynamic / WriteMask}; \textbf{GLSL:} \texttt{dst.y = a.y + b.y;})
  \item \texttt{opcode1}: Primary sub-opcode within the instruction class uniquely identifying the operation family. \emph{Example:} {\scriptsize\ttfamily \mnem{IADD.s32} \dst{r0}, \src{r1}, \src{r2}} (\textbf{SPIR-V:} \texttt{OpIAdd (operation selector)}; \textbf{GLSL:} \texttt{// Primary opcode bits})
  \item \texttt{instruction\_class}: 4-bit major instruction class selector (\texttt{bits 48:51}) partitioning the 64-bit Valhall opcode space (e.g., \texttt{FMA}, \texttt{CVT}, \texttt{SFU}, \texttt{LS}, \texttt{TEX}, \texttt{BRANCH}). \emph{Example:} {\scriptsize\ttfamily \mnem{FADD.f32} \dst{r0}, \src{r1}, \src{r2}} (\textbf{SPIR-V:} \texttt{N/A (encoding class)}; \textbf{GLSL:} \texttt{// Hardware execution pipe / decode class})
  \item \texttt{fau\_page\_index}: Selects the FAU (Fast Access Uniform) page (\texttt{0..3}) when a source operand references hardware small constants, thread IDs, or lane constants. \emph{Example:} {\scriptsize\ttfamily \mnem{FADD.f32} \dst{r0}, \src{r1}, \src{k1f}} (\textbf{SPIR-V:} \texttt{OpConstant / BuiltIn}; \textbf{GLSL:} \texttt{float y = x + 1.0;})
  \item \hyperref[enum:flow_control_m]{\texttt{flow}}: 4-bit embedded flow-control \& scoreboard wait modifier (\texttt{bits 59:62}: \texttt{none}, \texttt{wait0..wait0126}, \texttt{wait}, \texttt{reconverge}, \texttt{discard}, \texttt{end}). \emph{Example:} {\scriptsize\ttfamily \mnem{FADD.f32}\imod{.wait0.reconverge} \dst{r0}, \src{r1}, \src{r2}} (\textbf{SPIR-V:} \texttt{OpControlBarrier / OpReturn}; \textbf{GLSL:} \texttt{// Scoreboard wait + warp reconvergence})
\end{tightitemize}
\clearpage

\subsection*{V2F16\_TO\_V2S16}
\addcontentsline{toc}{subsection}{V2F16\_TO\_V2S16}
Converts each half-precision lane $i \in \{0, 1\}$ of swizzled \texttt{src0} to a saturating signed 16-bit integer \texttt{dest[i] = sat\_i16(round\_rm(src0[i]))} (NaN maps to \texttt{0}; arch \textless{}= v10).
\par\textbf{Notes:}
\begin{tightitemize}
  \item Arch v9-v10 only.
  \item \textbf{Source references:} \texttt{isa-v9-v14.xml}, \texttt{bifrost/valhall/ISA.xml}.
\end{tightitemize}
\codebox{\mnem{V2F16\_TO\_V2S16}\imod{\{.flow\}}\imod{\{.round\}} \dst{dest}\imod{\{.dest\_width\}}, \src{src0}\imod{\{.swz0\}}\imod{\{.abs0\}}\imod{\{.neg0\}}}
\textbf{Operands:}
\begin{tightitemize}\footnotesize
  \item \texttt{\dst{dest}} (\texttt{dest\_register\_index}, \texttt{dest\_width}): 32-bit destination GPR receiving the converted value.
  \item \texttt{\src{src0}}: Packed \texttt{v2f16} input converted per lane to saturating \texttt{v2s16}.
\end{tightitemize}
\par\textbf{Closest SPIR-V Op:} \texttt{OpConvertFToS (v2f16 -\textgreater{} v2s16, v9--v10)}
\par\textbf{Pseudo-GLSL (Operational Semantics):}
\codebox{vec2 h = unpackHalf2x16(src0);\\for (int i = 0; i \textless{} 2; ++i) dest[i] = isnan(h[i]) ? 0 : int16\_t(clamp(round\_rm(h[i]), -32768.0, 32767.0));}
\textbf{Execution Context:}
\begin{tightitemize}\footnotesize
  \item \textbf{Compute stage preloads (\hyperref[sec:csf-spd]{\texttt{SPD\_k}}, \texttt{r16..r23}):} Hardware pre-fills \texttt{r55..r57} with \texttt{WorkGroupID.xyz + JOB\_OFFSET} (from \hyperref[sec:csf-regs]{\texttt{COMPUTESR::JOB\_OFFSET\_X/Y/Z}} \texttt{r34..r36}) and \texttt{r60..r62} with \texttt{LocalInvocationID.xyz} or \texttt{GlobalInvocationID.xyz} (from \texttt{WG\_SIZE} \texttt{r33} \& \texttt{JOB\_SIZE} \texttt{r37..r39}, selected by \texttt{ShaderProgram::preload}) so compute shaders can derive thread/element indices without memory loads.
  \item \textbf{Vertex stage preloads (\hyperref[sec:csf-spd]{\texttt{SPD\_0}/\texttt{SPD\_1}}, \texttt{r16..r19}):} The hardware Index Fetcher reads \hyperref[sec:csf-regs]{\texttt{IDVSSR::INDEX\_BUFFER}} (\texttt{r54:r55}) and pre-fills \texttt{r59, r60} with \texttt{VertexIndex} (\texttt{+ VERTEX\_OFFSET + GLOBAL\_ATTR\_OFF}) and \texttt{r61} with \texttt{InstanceIndex} (\texttt{+ INSTANCE\_OFFSET}) so Position and Varying shaders can pull vertex attributes directly.
  \item \textbf{Inline FAU \& Special Operands (\hyperref[sec:csf-tsd-fau]{\texttt{FAU\_k}}, \texttt{r8..r15}):} Any 8-bit source with \texttt{mode=2} reads zero-latency uniform \texttt{u0..u31} from \texttt{FAU\_k} (\texttt{slots64 = bits[63:56]}); \texttt{mode=3} reads \texttt{SmallConstant} (Table~1.2) or \texttt{FauSpecial} (\texttt{pc}, \texttt{lane\_id}, \texttt{lane\_width}, \texttt{tls\_ptr}, \texttt{wls\_ptr}, \texttt{blend\_descriptor\_0..7}).
\end{tightitemize}
\par\textbf{Worked Examples:}
\par\textbf{Example 1 (Per-lane truncation toward zero (v9--v10)):}
\codebox{\mnem{V2F16\_TO\_V2S16}\imod{.round\_zero} \dst{r0}, \src{r1}\imod{.h01}}
With \texttt{r1 = [2.5h, -1.5h]}: truncates each halfword -\textgreater{} \texttt{r0 = [2, -1]} (\texttt{0x0002FFFF}).
\par\textbf{Example 2 (Overflow and NaN saturation):}
\codebox{\mnem{V2F16\_TO\_V2S16}\imod{.round\_zero} \dst{r0}, \src{r1}\imod{.h01}}
With \texttt{r1 = [+Inf, NaN]}: \texttt{+Inf} saturates to \texttt{INT16\_MAX} (\texttt{0x7FFF}), \texttt{NaN} becomes \texttt{0} -\textgreater{} \texttt{r0 = 0x7FFF0000}.

\begin{center}\footnotesize
\begin{tabularx}{\linewidth}{@{}lllX@{}}
\toprule
Variant & Arch & Class & Opcode Fields \\
\midrule
\texttt{(none)} & v9--v14 & \texttt{R1} & \texttt{opcode2=V2F16\_TO\_V2S16, opcode1=CONVERT} \\
\bottomrule
\end{tabularx}
\end{center}
\begin{center}\footnotesize
\begin{tabularx}{0.88\linewidth}{@{}l l X@{}}
\toprule
Field Name & Bits & Type / Fixed Value \\
\midrule
\texttt{src0} & \texttt{0:7} & \texttt{SourceEncoding} \\
\textit{Reserved} & \texttt{8:15} & \texttt{zero} \\
\texttt{opcode2} & \texttt{16:20} & \texttt{exact=V2F16\_TO\_V2S16} \\
\textit{Reserved} & \texttt{21:27} & \texttt{zero} \\
\hyperref[enum:swiz_m]{\texttt{swz0}} & \texttt{28:29} & \texttt{swiz\_m (def h01)} \\
\hyperref[enum:round_integer_m]{\texttt{round}} & \texttt{30:32} & \texttt{round\_integer\_m (def none)} \\
\textit{Reserved} & \texttt{33} & \texttt{zero} \\
\textit{Reserved} & \texttt{34:35} & \texttt{ignore} \\
\textit{Reserved} & \texttt{36:37} & \texttt{zero} \\
\hyperref[enum:neg_m]{\texttt{neg0}} & \texttt{38} & \texttt{neg\_m (def none)} \\
\hyperref[enum:abs_m]{\texttt{abs0}} & \texttt{39} & \texttt{abs\_m (def none)} \\
\texttt{dest\_register\_index} & \texttt{40:45} & \texttt{uint} \\
\hyperref[enum:dest_width_m]{\texttt{dest\_width}} & \texttt{46:47} & \texttt{dest\_width\_m} \\
\texttt{opcode1} & \texttt{48:51} & \texttt{exact=CONVERT} \\
\texttt{instruction\_class} & \texttt{52:56} & \texttt{exact=R1} \\
\texttt{fau\_page\_index} & \texttt{57:58} & \texttt{uint} \\
\hyperref[enum:flow_control_m]{\texttt{flow}} & \texttt{59:62} & \texttt{flow\_control\_m} \\
\textit{Reserved} & \texttt{63} & \texttt{zero} \\
\bottomrule
\end{tabularx}
\end{center}
\textbf{Fields Explained:}
\begin{tightitemize}\footnotesize
  \item \texttt{src0}: Primary source operand (\texttt{r0..r63}, \texttt{u0..u63} uniform, \texttt{k*} FAU constant, or \texttt{t0..t2} temporary; \texttt{\textasciicircum{}} marks discard/last-use). \emph{Example:} {\scriptsize\ttfamily \mnem{FADD.f32} \dst{r0}, \src{r1}, \src{r2}} (\textbf{SPIR-V:} \texttt{OpFAdd (Operand 0)}; \textbf{GLSL:} \texttt{float c = a + b;})
  \item \texttt{opcode2}: Secondary sub-opcode further specializing instructions sharing \texttt{instruction\_class} and \texttt{opcode1} (e.g., conversions, SFU tables, moves). \emph{Example:} {\scriptsize\ttfamily \mnem{F32\_TO\_S32}\imod{.round\_zero} \dst{r0}, \src{r1}} (\textbf{SPIR-V:} \texttt{OpConvertFToS}; \textbf{GLSL:} \texttt{int y = int(x);})
  \item \hyperref[enum:swiz_m]{\texttt{swz0}}: Stage A lane swizzle on \texttt{src0} (\texttt{h00}/\texttt{h01}/\texttt{h10}/\texttt{h11} for 16-bit vectors, or \texttt{bXY} byte-to-halfword widening). \emph{Example:} {\scriptsize\ttfamily \mnem{FADD.v2f16} \dst{r0}, \src{r1}\imod{.h10}, \src{r2}} (\textbf{SPIR-V:} \texttt{OpVectorShuffle}; \textbf{GLSL:} \texttt{f16vec2 c = a.yx + b;})
  \item \hyperref[enum:round_integer_m]{\texttt{round}}: Rounding mode (\texttt{none}/RNE, \texttt{round\_up}, \texttt{round\_down}, \texttt{round\_zero}, \texttt{round\_na}) applied in FP ALU, conversions, \texttt{HADD}, \texttt{LDEXP}, and Stage C narrowing. \emph{Example:} {\scriptsize\ttfamily \mnem{FROUND.f32}\imod{.round\_down} \dst{r0}, \src{r1}} (\textbf{SPIR-V:} \texttt{GLSL.std.450 Floor / Ceil / Trunc}; \textbf{GLSL:} \texttt{float y = floor(x);})
  \item \hyperref[enum:neg_m]{\texttt{neg0}}: Stage A floating-point negation on \texttt{src0} (flips IEEE sign bit after swizzle/widen and \texttt{abs0}). \emph{Example:} {\scriptsize\ttfamily \mnem{FADD.f32} \dst{r0}, \src{r1}\imod{.neg}, \src{r2}} (\textbf{SPIR-V:} \texttt{OpFNegate}; \textbf{GLSL:} \texttt{float c = -a + b;})
  \item \hyperref[enum:abs_m]{\texttt{abs0}}: Stage A floating-point absolute value on \texttt{src0} (clears IEEE sign bit after swizzle/widen, before \texttt{neg0}). \emph{Example:} {\scriptsize\ttfamily \mnem{FADD.f32} \dst{r0}, \src{r1}\imod{.abs}, \src{r2}} (\textbf{SPIR-V:} \texttt{GLSL.std.450 FAbs}; \textbf{GLSL:} \texttt{float c = abs(a) + b;})
  \item \texttt{dest\_register\_index}: 6-bit destination GPR index (\texttt{r0..r63}) receiving the Stage C writeback result (even-aligned pair \texttt{[rN:rN+1]} for 64-bit ops). \emph{Example:} {\scriptsize\ttfamily \mnem{FMA.f32} \dst{r0}, \src{r1}, \src{r2}, \src{r3}} (\textbf{SPIR-V:} \texttt{SSA Result ID}; \textbf{GLSL:} \texttt{float r0 = fma(r1, r2, r3);})
  \item \hyperref[enum:dest_width_m]{\texttt{dest\_width}}: Stage C destination lane write-mask (\texttt{none} = full 32-bit, \texttt{h0} = low 16-bit, \texttt{h1} = high 16-bit without bit-shifting). \emph{Example:} {\scriptsize\ttfamily \mnem{FADD.v2f16} \dst{r0}\imod{.h1}, \src{r1}, \src{r2}} (\textbf{SPIR-V:} \texttt{OpVectorInsertDynamic / WriteMask}; \textbf{GLSL:} \texttt{dst.y = a.y + b.y;})
  \item \texttt{opcode1}: Primary sub-opcode within the instruction class uniquely identifying the operation family. \emph{Example:} {\scriptsize\ttfamily \mnem{IADD.s32} \dst{r0}, \src{r1}, \src{r2}} (\textbf{SPIR-V:} \texttt{OpIAdd (operation selector)}; \textbf{GLSL:} \texttt{// Primary opcode bits})
  \item \texttt{instruction\_class}: 4-bit major instruction class selector (\texttt{bits 48:51}) partitioning the 64-bit Valhall opcode space (e.g., \texttt{FMA}, \texttt{CVT}, \texttt{SFU}, \texttt{LS}, \texttt{TEX}, \texttt{BRANCH}). \emph{Example:} {\scriptsize\ttfamily \mnem{FADD.f32} \dst{r0}, \src{r1}, \src{r2}} (\textbf{SPIR-V:} \texttt{N/A (encoding class)}; \textbf{GLSL:} \texttt{// Hardware execution pipe / decode class})
  \item \texttt{fau\_page\_index}: Selects the FAU (Fast Access Uniform) page (\texttt{0..3}) when a source operand references hardware small constants, thread IDs, or lane constants. \emph{Example:} {\scriptsize\ttfamily \mnem{FADD.f32} \dst{r0}, \src{r1}, \src{k1f}} (\textbf{SPIR-V:} \texttt{OpConstant / BuiltIn}; \textbf{GLSL:} \texttt{float y = x + 1.0;})
  \item \hyperref[enum:flow_control_m]{\texttt{flow}}: 4-bit embedded flow-control \& scoreboard wait modifier (\texttt{bits 59:62}: \texttt{none}, \texttt{wait0..wait0126}, \texttt{wait}, \texttt{reconverge}, \texttt{discard}, \texttt{end}). \emph{Example:} {\scriptsize\ttfamily \mnem{FADD.f32}\imod{.wait0.reconverge} \dst{r0}, \src{r1}, \src{r2}} (\textbf{SPIR-V:} \texttt{OpControlBarrier / OpReturn}; \textbf{GLSL:} \texttt{// Scoreboard wait + warp reconvergence})
\end{tightitemize}
\clearpage

\subsection*{V2F16\_TO\_V2U16}
\addcontentsline{toc}{subsection}{V2F16\_TO\_V2U16}
Converts each half-precision lane $i \in \{0, 1\}$ of swizzled \texttt{src0} to a saturating unsigned 16-bit integer \texttt{dest[i] = sat\_u16(round\_rm(src0[i]))} (NaN or \texttt{\textless{}= 0} maps to \texttt{0}; arch \textless{}= v10).
\par\textbf{Notes:}
\begin{tightitemize}
  \item Arch v9-v10 only.
  \item \textbf{Source references:} \texttt{isa-v9-v14.xml}, \texttt{bifrost/valhall/ISA.xml}.
\end{tightitemize}
\codebox{\mnem{V2F16\_TO\_V2U16}\imod{\{.flow\}}\imod{\{.round\}} \dst{dest}\imod{\{.dest\_width\}}, \src{src0}\imod{\{.swz0\}}\imod{\{.abs0\}}\imod{\{.neg0\}}}
\textbf{Operands:}
\begin{tightitemize}\footnotesize
  \item \texttt{\dst{dest}} (\texttt{dest\_register\_index}, \texttt{dest\_width}): 32-bit destination GPR receiving the converted value.
  \item \texttt{\src{src0}}: Packed \texttt{v2f16} input converted per lane to saturating \texttt{v2u16}.
\end{tightitemize}
\par\textbf{Closest SPIR-V Op:} \texttt{OpConvertFToU (v2f16 -\textgreater{} v2u16, v9--v10)}
\par\textbf{Pseudo-GLSL (Operational Semantics):}
\codebox{vec2 h = unpackHalf2x16(src0);\\for (int i = 0; i \textless{} 2; ++i) dest[i] = (isnan(h[i]) \textbar{}\textbar{} h[i] \textless{}= 0.0) ? 0u : uint16\_t(clamp(round\_rm(h[i]), 0.0, 65535.0));}
\textbf{Execution Context:}
\begin{tightitemize}\footnotesize
  \item \textbf{Compute stage preloads (\hyperref[sec:csf-spd]{\texttt{SPD\_k}}, \texttt{r16..r23}):} Hardware pre-fills \texttt{r55..r57} with \texttt{WorkGroupID.xyz + JOB\_OFFSET} (from \hyperref[sec:csf-regs]{\texttt{COMPUTESR::JOB\_OFFSET\_X/Y/Z}} \texttt{r34..r36}) and \texttt{r60..r62} with \texttt{LocalInvocationID.xyz} or \texttt{GlobalInvocationID.xyz} (from \texttt{WG\_SIZE} \texttt{r33} \& \texttt{JOB\_SIZE} \texttt{r37..r39}, selected by \texttt{ShaderProgram::preload}) so compute shaders can derive thread/element indices without memory loads.
  \item \textbf{Vertex stage preloads (\hyperref[sec:csf-spd]{\texttt{SPD\_0}/\texttt{SPD\_1}}, \texttt{r16..r19}):} The hardware Index Fetcher reads \hyperref[sec:csf-regs]{\texttt{IDVSSR::INDEX\_BUFFER}} (\texttt{r54:r55}) and pre-fills \texttt{r59, r60} with \texttt{VertexIndex} (\texttt{+ VERTEX\_OFFSET + GLOBAL\_ATTR\_OFF}) and \texttt{r61} with \texttt{InstanceIndex} (\texttt{+ INSTANCE\_OFFSET}) so Position and Varying shaders can pull vertex attributes directly.
  \item \textbf{Inline FAU \& Special Operands (\hyperref[sec:csf-tsd-fau]{\texttt{FAU\_k}}, \texttt{r8..r15}):} Any 8-bit source with \texttt{mode=2} reads zero-latency uniform \texttt{u0..u31} from \texttt{FAU\_k} (\texttt{slots64 = bits[63:56]}); \texttt{mode=3} reads \texttt{SmallConstant} (Table~1.2) or \texttt{FauSpecial} (\texttt{pc}, \texttt{lane\_id}, \texttt{lane\_width}, \texttt{tls\_ptr}, \texttt{wls\_ptr}, \texttt{blend\_descriptor\_0..7}).
\end{tightitemize}
\par\textbf{Worked Examples:}
\par\textbf{Example 1 (Negative lane clamps to 0 (v9--v10)):}
\codebox{\mnem{V2F16\_TO\_V2U16}\imod{.round\_zero} \dst{r0}, \src{r1}\imod{.h01}}
With \texttt{r1 = [3.0h, -2.0h]}: saturates negative lane to 0 -\textgreater{} \texttt{r0 = [3, 0]} (\texttt{0x00030000}).
\par\textbf{Example 2 (Swizzled lanes + `+Inf` saturation):}
\codebox{\mnem{V2F16\_TO\_V2U16}\imod{.round\_zero} \dst{r0}, \src{r1}\imod{.h10}}
With \texttt{r1 = [1.0h, +Inf]}: \texttt{h10} swaps lanes, saturating \texttt{+Inf} to \texttt{0xFFFF} -\textgreater{} \texttt{r0 = 0x0001FFFF}.

\begin{center}\footnotesize
\begin{tabularx}{\linewidth}{@{}lllX@{}}
\toprule
Variant & Arch & Class & Opcode Fields \\
\midrule
\texttt{(none)} & v9--v14 & \texttt{R1} & \texttt{opcode2=V2F16\_TO\_V2U16, opcode1=CONVERT} \\
\bottomrule
\end{tabularx}
\end{center}
\begin{center}\footnotesize
\begin{tabularx}{0.88\linewidth}{@{}l l X@{}}
\toprule
Field Name & Bits & Type / Fixed Value \\
\midrule
\texttt{src0} & \texttt{0:7} & \texttt{SourceEncoding} \\
\textit{Reserved} & \texttt{8:15} & \texttt{zero} \\
\texttt{opcode2} & \texttt{16:20} & \texttt{exact=V2F16\_TO\_V2U16} \\
\textit{Reserved} & \texttt{21:27} & \texttt{zero} \\
\hyperref[enum:swiz_m]{\texttt{swz0}} & \texttt{28:29} & \texttt{swiz\_m (def h01)} \\
\hyperref[enum:round_integer_m]{\texttt{round}} & \texttt{30:32} & \texttt{round\_integer\_m (def none)} \\
\textit{Reserved} & \texttt{33} & \texttt{zero} \\
\textit{Reserved} & \texttt{34:35} & \texttt{ignore} \\
\textit{Reserved} & \texttt{36:37} & \texttt{zero} \\
\hyperref[enum:neg_m]{\texttt{neg0}} & \texttt{38} & \texttt{neg\_m (def none)} \\
\hyperref[enum:abs_m]{\texttt{abs0}} & \texttt{39} & \texttt{abs\_m (def none)} \\
\texttt{dest\_register\_index} & \texttt{40:45} & \texttt{uint} \\
\hyperref[enum:dest_width_m]{\texttt{dest\_width}} & \texttt{46:47} & \texttt{dest\_width\_m} \\
\texttt{opcode1} & \texttt{48:51} & \texttt{exact=CONVERT} \\
\texttt{instruction\_class} & \texttt{52:56} & \texttt{exact=R1} \\
\texttt{fau\_page\_index} & \texttt{57:58} & \texttt{uint} \\
\hyperref[enum:flow_control_m]{\texttt{flow}} & \texttt{59:62} & \texttt{flow\_control\_m} \\
\textit{Reserved} & \texttt{63} & \texttt{zero} \\
\bottomrule
\end{tabularx}
\end{center}
\textbf{Fields Explained:}
\begin{tightitemize}\footnotesize
  \item \texttt{src0}: Primary source operand (\texttt{r0..r63}, \texttt{u0..u63} uniform, \texttt{k*} FAU constant, or \texttt{t0..t2} temporary; \texttt{\textasciicircum{}} marks discard/last-use). \emph{Example:} {\scriptsize\ttfamily \mnem{FADD.f32} \dst{r0}, \src{r1}, \src{r2}} (\textbf{SPIR-V:} \texttt{OpFAdd (Operand 0)}; \textbf{GLSL:} \texttt{float c = a + b;})
  \item \texttt{opcode2}: Secondary sub-opcode further specializing instructions sharing \texttt{instruction\_class} and \texttt{opcode1} (e.g., conversions, SFU tables, moves). \emph{Example:} {\scriptsize\ttfamily \mnem{F32\_TO\_S32}\imod{.round\_zero} \dst{r0}, \src{r1}} (\textbf{SPIR-V:} \texttt{OpConvertFToS}; \textbf{GLSL:} \texttt{int y = int(x);})
  \item \hyperref[enum:swiz_m]{\texttt{swz0}}: Stage A lane swizzle on \texttt{src0} (\texttt{h00}/\texttt{h01}/\texttt{h10}/\texttt{h11} for 16-bit vectors, or \texttt{bXY} byte-to-halfword widening). \emph{Example:} {\scriptsize\ttfamily \mnem{FADD.v2f16} \dst{r0}, \src{r1}\imod{.h10}, \src{r2}} (\textbf{SPIR-V:} \texttt{OpVectorShuffle}; \textbf{GLSL:} \texttt{f16vec2 c = a.yx + b;})
  \item \hyperref[enum:round_integer_m]{\texttt{round}}: Rounding mode (\texttt{none}/RNE, \texttt{round\_up}, \texttt{round\_down}, \texttt{round\_zero}, \texttt{round\_na}) applied in FP ALU, conversions, \texttt{HADD}, \texttt{LDEXP}, and Stage C narrowing. \emph{Example:} {\scriptsize\ttfamily \mnem{FROUND.f32}\imod{.round\_down} \dst{r0}, \src{r1}} (\textbf{SPIR-V:} \texttt{GLSL.std.450 Floor / Ceil / Trunc}; \textbf{GLSL:} \texttt{float y = floor(x);})
  \item \hyperref[enum:neg_m]{\texttt{neg0}}: Stage A floating-point negation on \texttt{src0} (flips IEEE sign bit after swizzle/widen and \texttt{abs0}). \emph{Example:} {\scriptsize\ttfamily \mnem{FADD.f32} \dst{r0}, \src{r1}\imod{.neg}, \src{r2}} (\textbf{SPIR-V:} \texttt{OpFNegate}; \textbf{GLSL:} \texttt{float c = -a + b;})
  \item \hyperref[enum:abs_m]{\texttt{abs0}}: Stage A floating-point absolute value on \texttt{src0} (clears IEEE sign bit after swizzle/widen, before \texttt{neg0}). \emph{Example:} {\scriptsize\ttfamily \mnem{FADD.f32} \dst{r0}, \src{r1}\imod{.abs}, \src{r2}} (\textbf{SPIR-V:} \texttt{GLSL.std.450 FAbs}; \textbf{GLSL:} \texttt{float c = abs(a) + b;})
  \item \texttt{dest\_register\_index}: 6-bit destination GPR index (\texttt{r0..r63}) receiving the Stage C writeback result (even-aligned pair \texttt{[rN:rN+1]} for 64-bit ops). \emph{Example:} {\scriptsize\ttfamily \mnem{FMA.f32} \dst{r0}, \src{r1}, \src{r2}, \src{r3}} (\textbf{SPIR-V:} \texttt{SSA Result ID}; \textbf{GLSL:} \texttt{float r0 = fma(r1, r2, r3);})
  \item \hyperref[enum:dest_width_m]{\texttt{dest\_width}}: Stage C destination lane write-mask (\texttt{none} = full 32-bit, \texttt{h0} = low 16-bit, \texttt{h1} = high 16-bit without bit-shifting). \emph{Example:} {\scriptsize\ttfamily \mnem{FADD.v2f16} \dst{r0}\imod{.h1}, \src{r1}, \src{r2}} (\textbf{SPIR-V:} \texttt{OpVectorInsertDynamic / WriteMask}; \textbf{GLSL:} \texttt{dst.y = a.y + b.y;})
  \item \texttt{opcode1}: Primary sub-opcode within the instruction class uniquely identifying the operation family. \emph{Example:} {\scriptsize\ttfamily \mnem{IADD.s32} \dst{r0}, \src{r1}, \src{r2}} (\textbf{SPIR-V:} \texttt{OpIAdd (operation selector)}; \textbf{GLSL:} \texttt{// Primary opcode bits})
  \item \texttt{instruction\_class}: 4-bit major instruction class selector (\texttt{bits 48:51}) partitioning the 64-bit Valhall opcode space (e.g., \texttt{FMA}, \texttt{CVT}, \texttt{SFU}, \texttt{LS}, \texttt{TEX}, \texttt{BRANCH}). \emph{Example:} {\scriptsize\ttfamily \mnem{FADD.f32} \dst{r0}, \src{r1}, \src{r2}} (\textbf{SPIR-V:} \texttt{N/A (encoding class)}; \textbf{GLSL:} \texttt{// Hardware execution pipe / decode class})
  \item \texttt{fau\_page\_index}: Selects the FAU (Fast Access Uniform) page (\texttt{0..3}) when a source operand references hardware small constants, thread IDs, or lane constants. \emph{Example:} {\scriptsize\ttfamily \mnem{FADD.f32} \dst{r0}, \src{r1}, \src{k1f}} (\textbf{SPIR-V:} \texttt{OpConstant / BuiltIn}; \textbf{GLSL:} \texttt{float y = x + 1.0;})
  \item \hyperref[enum:flow_control_m]{\texttt{flow}}: 4-bit embedded flow-control \& scoreboard wait modifier (\texttt{bits 59:62}: \texttt{none}, \texttt{wait0..wait0126}, \texttt{wait}, \texttt{reconverge}, \texttt{discard}, \texttt{end}). \emph{Example:} {\scriptsize\ttfamily \mnem{FADD.f32}\imod{.wait0.reconverge} \dst{r0}, \src{r1}, \src{r2}} (\textbf{SPIR-V:} \texttt{OpControlBarrier / OpReturn}; \textbf{GLSL:} \texttt{// Scoreboard wait + warp reconvergence})
\end{tightitemize}
\clearpage

\subsection*{V2F32\_TO\_V2F16}
\addcontentsline{toc}{subsection}{V2F32\_TO\_V2F16}
Two-source FP32-to-FP16 packing conversion (v9--v10): clamps and rounds \texttt{src0} to FP16 into \texttt{dest.h0} and clamps and rounds \texttt{src1} to FP16 into \texttt{dest.h1} in a single instruction.
\par\textbf{Notes:}
\begin{tightitemize}
  \item Arch v9-v10 only (\texttt{kraid/nir.rs:653, 676}).
  \item Clamp is applied BEFORE f32-\textgreater{}f16 rounding (\texttt{ops.rs:4629-4638}).
  \item \textbf{Source references:} \texttt{kraid/ops.rs:4597-4645}, \texttt{encode\_v9.rs:3462-3482}, \texttt{nir.rs:652-689}.
\end{tightitemize}
\codebox{\mnem{V2F32\_TO\_V2F16}\imod{\{.flow\}}\imod{\{.round\}}\imod{\{.clamp\}} \dst{dest}\imod{\{.dest\_width\}}, \src{src0}\imod{\{.abs0\}}\imod{\{.neg0\}}, \src{src1}\imod{\{.abs1\}}\imod{\{.neg1\}}}
\textbf{Operands:}
\begin{tightitemize}\footnotesize
  \item \texttt{\dst{dest}} (\texttt{dest\_register\_index}, \texttt{dest\_width}): 32-bit destination GPR receiving packed \texttt{v2f16(src0, src1)}.
  \item \texttt{\src{src0}}: FP32 input converted into the low halfword (\texttt{dest.h0}).
  \item \texttt{\src{src1}}: FP32 input converted into the high halfword (\texttt{dest.h1}).
\end{tightitemize}
\par\textbf{Closest SPIR-V Op:} \texttt{GLSL.std.450 PackHalf2x16 (v9--v10)}
\par\textbf{Pseudo-GLSL (Operational Semantics):}
\codebox{dest = packHalf2x16(vec2(round\_rm(clamp\_mode(src0)), round\_rm(clamp\_mode(src1))));}
\textbf{Execution Context:}
\begin{tightitemize}\footnotesize
  \item \textbf{Compute stage preloads (\hyperref[sec:csf-spd]{\texttt{SPD\_k}}, \texttt{r16..r23}):} Hardware pre-fills \texttt{r55..r57} with \texttt{WorkGroupID.xyz + JOB\_OFFSET} (from \hyperref[sec:csf-regs]{\texttt{COMPUTESR::JOB\_OFFSET\_X/Y/Z}} \texttt{r34..r36}) and \texttt{r60..r62} with \texttt{LocalInvocationID.xyz} or \texttt{GlobalInvocationID.xyz} (from \texttt{WG\_SIZE} \texttt{r33} \& \texttt{JOB\_SIZE} \texttt{r37..r39}, selected by \texttt{ShaderProgram::preload}) so compute shaders can derive thread/element indices without memory loads.
  \item \textbf{Vertex stage preloads (\hyperref[sec:csf-spd]{\texttt{SPD\_0}/\texttt{SPD\_1}}, \texttt{r16..r19}):} The hardware Index Fetcher reads \hyperref[sec:csf-regs]{\texttt{IDVSSR::INDEX\_BUFFER}} (\texttt{r54:r55}) and pre-fills \texttt{r59, r60} with \texttt{VertexIndex} (\texttt{+ VERTEX\_OFFSET + GLOBAL\_ATTR\_OFF}) and \texttt{r61} with \texttt{InstanceIndex} (\texttt{+ INSTANCE\_OFFSET}) so Position and Varying shaders can pull vertex attributes directly.
  \item \textbf{Inline FAU \& Special Operands (\hyperref[sec:csf-tsd-fau]{\texttt{FAU\_k}}, \texttt{r8..r15}):} Any 8-bit source with \texttt{mode=2} reads zero-latency uniform \texttt{u0..u31} from \texttt{FAU\_k} (\texttt{slots64 = bits[63:56]}); \texttt{mode=3} reads \texttt{SmallConstant} (Table~1.2) or \texttt{FauSpecial} (\texttt{pc}, \texttt{lane\_id}, \texttt{lane\_width}, \texttt{tls\_ptr}, \texttt{wls\_ptr}, \texttt{blend\_descriptor\_0..7}).
\end{tightitemize}
\par\textbf{Worked Examples:}
\par\textbf{Example 1 (Pre-narrowing `.clamp\_0\_1` packing (v9--v10)):}
\codebox{\mnem{V2F32\_TO\_V2F16}\imod{.clamp\_0\_1} \dst{r0}, \src{r1}, \src{r2}}
With \texttt{r1 = 1.0f, r2 = 2.0f}: clamps BEFORE narrowing (\texttt{1.0f, 1.0f}) -\textgreater{} packs \texttt{r0 = 0x3C003C00}.
\par\textbf{Example 2 (Rounding mode control on pack):}
\codebox{\mnem{V2F32\_TO\_V2F16}\imod{.round\_zero} \dst{r0}, \src{r1}, \src{r2}}
With \texttt{r1 = 65536.0f} (\texttt{\textgreater{} FP16 MAX 65504}), \texttt{r2 = -0.5f}: \texttt{.round\_zero} yields \texttt{MAX\_FP16} (\texttt{0x7BFF}) in \texttt{h0} and \texttt{-0.5h} (\texttt{0xB800}) in \texttt{h1} -\textgreater{} \texttt{r0 = 0xB8007BFF}.
\par\textbf{Mined SPIR-V $\to$ Valhall ISA Example (Packing two FP32 values into packed FP16 (PackHalf2x16) --- Vertex shader, Mali-G710 (v10)):}
\codebox{\textbf{; SPIR-V (input)}\\\%packed = OpExtInst \%uint \%glsl PackHalf2x16 \%v2f\\[3pt]\textbf{; Valhall ISA (malioc output)}\\V2F32\_TO\_V2F16.wait1 r57.h0, r57\textasciicircum{}, r57\textasciicircum{}}
\texttt{GLSL.std.450 PackHalf2x16} lowers directly to \texttt{V2F32\_TO\_V2F16}, narrowing two FP32 sources into \texttt{dest.h0} and \texttt{dest.h1} in a single instruction (v9--v10).
\begin{center}\footnotesize
\begin{tabularx}{\linewidth}{@{}lllX@{}}
\toprule
Variant & Arch & Class & Opcode Fields \\
\midrule
\texttt{v2f16} & v9--v14 & \texttt{R2} & \texttt{opcode2=CONVERT} \\
\bottomrule
\end{tabularx}
\end{center}
\begin{center}\footnotesize
\begin{tabularx}{0.88\linewidth}{@{}l l X@{}}
\toprule
Field Name & Bits & Type / Fixed Value \\
\midrule
\texttt{src0} & \texttt{0:7} & \texttt{SourceEncoding} \\
\texttt{src1} & \texttt{8:15} & \texttt{SourceEncoding} \\
\texttt{opcode2} & \texttt{16:19} & \texttt{exact=CONVERT} \\
\textit{Reserved} & \texttt{20:29} & \texttt{zero} \\
\hyperref[enum:round_m]{\texttt{round}} & \texttt{30:31} & \texttt{round\_m (def none)} \\
\hyperref[enum:clamp_m]{\texttt{clamp}} & \texttt{32:33} & \texttt{clamp\_m (def none)} \\
\textit{Reserved} & \texttt{34:35} & \texttt{ignore} \\
\hyperref[enum:neg_m]{\texttt{neg1}} & \texttt{36} & \texttt{neg\_m (def none)} \\
\hyperref[enum:abs_m]{\texttt{abs1}} & \texttt{37} & \texttt{abs\_m (def none)} \\
\hyperref[enum:neg_m]{\texttt{neg0}} & \texttt{38} & \texttt{neg\_m (def none)} \\
\hyperref[enum:abs_m]{\texttt{abs0}} & \texttt{39} & \texttt{abs\_m (def none)} \\
\texttt{dest\_register\_index} & \texttt{40:45} & \texttt{uint} \\
\hyperref[enum:dest_width_m]{\texttt{dest\_width}} & \texttt{46:47} & \texttt{dest\_width\_m} \\
\texttt{datatype} & \texttt{48:51} & \texttt{exact=v2f16} \\
\texttt{instruction\_class} & \texttt{52:56} & \texttt{exact=R2} \\
\texttt{fau\_page\_index} & \texttt{57:58} & \texttt{uint} \\
\hyperref[enum:flow_control_m]{\texttt{flow}} & \texttt{59:62} & \texttt{flow\_control\_m} \\
\textit{Reserved} & \texttt{63} & \texttt{zero} \\
\bottomrule
\end{tabularx}
\end{center}
\textbf{Fields Explained:}
\begin{tightitemize}\footnotesize
  \item \texttt{src0}: Primary source operand (\texttt{r0..r63}, \texttt{u0..u63} uniform, \texttt{k*} FAU constant, or \texttt{t0..t2} temporary; \texttt{\textasciicircum{}} marks discard/last-use). \emph{Example:} {\scriptsize\ttfamily \mnem{FADD.f32} \dst{r0}, \src{r1}, \src{r2}} (\textbf{SPIR-V:} \texttt{OpFAdd (Operand 0)}; \textbf{GLSL:} \texttt{float c = a + b;})
  \item \texttt{src1}: Secondary source operand (register, uniform, or FAU constant; in shifts, supplies the shift count byte). \emph{Example:} {\scriptsize\ttfamily \mnem{LSHIFT.i32} \dst{r0}, \src{r1}, \src{r2}\imod{.b0}} (\textbf{SPIR-V:} \texttt{OpShiftLeftLogical (Shift)}; \textbf{GLSL:} \texttt{uint y = a \textless{}\textless{} s;})
  \item \texttt{opcode2}: Secondary sub-opcode further specializing instructions sharing \texttt{instruction\_class} and \texttt{opcode1} (e.g., conversions, SFU tables, moves). \emph{Example:} {\scriptsize\ttfamily \mnem{F32\_TO\_S32}\imod{.round\_zero} \dst{r0}, \src{r1}} (\textbf{SPIR-V:} \texttt{OpConvertFToS}; \textbf{GLSL:} \texttt{int y = int(x);})
  \item \hyperref[enum:round_m]{\texttt{round}}: Rounding mode (\texttt{none}/RNE, \texttt{round\_up}, \texttt{round\_down}, \texttt{round\_zero}, \texttt{round\_na}) applied in FP ALU, conversions, \texttt{HADD}, \texttt{LDEXP}, and Stage C narrowing. \emph{Example:} {\scriptsize\ttfamily \mnem{FROUND.f32}\imod{.round\_down} \dst{r0}, \src{r1}} (\textbf{SPIR-V:} \texttt{GLSL.std.450 Floor / Ceil / Trunc}; \textbf{GLSL:} \texttt{float y = floor(x);})
  \item \hyperref[enum:clamp_m]{\texttt{clamp}}: Stage B post-op floating-point clamp (\texttt{none}, \texttt{clamp\_0\_inf}, \texttt{clamp\_m1\_1}, \texttt{clamp\_0\_1}) with \texttt{-0.0 \textless{} +0.0} and NaN suppression (\texttt{NaN -\textgreater{} lo}). \emph{Example:} {\scriptsize\ttfamily \mnem{FMA.f32}\imod{.clamp\_0\_1} \dst{r0}, \src{r1}, \src{r2}, \src{r3}} (\textbf{SPIR-V:} \texttt{GLSL.std.450 FClamp}; \textbf{GLSL:} \texttt{float y = clamp(fma(a, b, c), 0.0, 1.0);})
  \item \hyperref[enum:neg_m]{\texttt{neg1}}: Stage A floating-point negation on \texttt{src1} (flips IEEE sign bit after swizzle/widen and \texttt{abs1}). \emph{Example:} {\scriptsize\ttfamily \mnem{FMA.f32} \dst{r0}, \src{r1}, \src{r2}\imod{.neg}, \src{r3}} (\textbf{SPIR-V:} \texttt{OpFNegate}; \textbf{GLSL:} \texttt{float d = fma(a, -b, c);})
  \item \hyperref[enum:abs_m]{\texttt{abs1}}: Stage A floating-point absolute value on \texttt{src1} (clears IEEE sign bit after swizzle/widen, before \texttt{neg1}). \emph{Example:} {\scriptsize\ttfamily \mnem{FMA.f32} \dst{r0}, \src{r1}, \src{r2}\imod{.abs}, \src{r3}} (\textbf{SPIR-V:} \texttt{GLSL.std.450 FAbs}; \textbf{GLSL:} \texttt{float d = fma(a, abs(b), c);})
  \item \hyperref[enum:neg_m]{\texttt{neg0}}: Stage A floating-point negation on \texttt{src0} (flips IEEE sign bit after swizzle/widen and \texttt{abs0}). \emph{Example:} {\scriptsize\ttfamily \mnem{FADD.f32} \dst{r0}, \src{r1}\imod{.neg}, \src{r2}} (\textbf{SPIR-V:} \texttt{OpFNegate}; \textbf{GLSL:} \texttt{float c = -a + b;})
  \item \hyperref[enum:abs_m]{\texttt{abs0}}: Stage A floating-point absolute value on \texttt{src0} (clears IEEE sign bit after swizzle/widen, before \texttt{neg0}). \emph{Example:} {\scriptsize\ttfamily \mnem{FADD.f32} \dst{r0}, \src{r1}\imod{.abs}, \src{r2}} (\textbf{SPIR-V:} \texttt{GLSL.std.450 FAbs}; \textbf{GLSL:} \texttt{float c = abs(a) + b;})
  \item \texttt{dest\_register\_index}: 6-bit destination GPR index (\texttt{r0..r63}) receiving the Stage C writeback result (even-aligned pair \texttt{[rN:rN+1]} for 64-bit ops). \emph{Example:} {\scriptsize\ttfamily \mnem{FMA.f32} \dst{r0}, \src{r1}, \src{r2}, \src{r3}} (\textbf{SPIR-V:} \texttt{SSA Result ID}; \textbf{GLSL:} \texttt{float r0 = fma(r1, r2, r3);})
  \item \hyperref[enum:dest_width_m]{\texttt{dest\_width}}: Stage C destination lane write-mask (\texttt{none} = full 32-bit, \texttt{h0} = low 16-bit, \texttt{h1} = high 16-bit without bit-shifting). \emph{Example:} {\scriptsize\ttfamily \mnem{FADD.v2f16} \dst{r0}\imod{.h1}, \src{r1}, \src{r2}} (\textbf{SPIR-V:} \texttt{OpVectorInsertDynamic / WriteMask}; \textbf{GLSL:} \texttt{dst.y = a.y + b.y;})
  \item \texttt{datatype}: Selects the lane data type (\texttt{f32}, \texttt{f16}, \texttt{v2f16}, \texttt{s32}, \texttt{u32}, \texttt{v2s16}, \texttt{v2u16}, \texttt{v4s8}, \texttt{v4u8}, \texttt{s64}, \texttt{u64}, etc.). \emph{Example:} {\scriptsize\ttfamily \mnem{IADD.v2s16} \dst{r0}, \src{r1}, \src{r2}} (\textbf{SPIR-V:} \texttt{OpTypeInt / OpTypeFloat}; \textbf{GLSL:} \texttt{i16vec2 c = a + b;})
  \item \texttt{instruction\_class}: 4-bit major instruction class selector (\texttt{bits 48:51}) partitioning the 64-bit Valhall opcode space (e.g., \texttt{FMA}, \texttt{CVT}, \texttt{SFU}, \texttt{LS}, \texttt{TEX}, \texttt{BRANCH}). \emph{Example:} {\scriptsize\ttfamily \mnem{FADD.f32} \dst{r0}, \src{r1}, \src{r2}} (\textbf{SPIR-V:} \texttt{N/A (encoding class)}; \textbf{GLSL:} \texttt{// Hardware execution pipe / decode class})
  \item \texttt{fau\_page\_index}: Selects the FAU (Fast Access Uniform) page (\texttt{0..3}) when a source operand references hardware small constants, thread IDs, or lane constants. \emph{Example:} {\scriptsize\ttfamily \mnem{FADD.f32} \dst{r0}, \src{r1}, \src{k1f}} (\textbf{SPIR-V:} \texttt{OpConstant / BuiltIn}; \textbf{GLSL:} \texttt{float y = x + 1.0;})
  \item \hyperref[enum:flow_control_m]{\texttt{flow}}: 4-bit embedded flow-control \& scoreboard wait modifier (\texttt{bits 59:62}: \texttt{none}, \texttt{wait0..wait0126}, \texttt{wait}, \texttt{reconverge}, \texttt{discard}, \texttt{end}). \emph{Example:} {\scriptsize\ttfamily \mnem{FADD.f32}\imod{.wait0.reconverge} \dst{r0}, \src{r1}, \src{r2}} (\textbf{SPIR-V:} \texttt{OpControlBarrier / OpReturn}; \textbf{GLSL:} \texttt{// Scoreboard wait + warp reconvergence})
\end{tightitemize}
\clearpage

\subsection*{V2S16\_TO\_V2F16}
\addcontentsline{toc}{subsection}{V2S16\_TO\_V2F16}
Converts each signed 16-bit integer lane $i \in \{0, 1\}$ of swizzled \texttt{src0} to IEEE-754 FP16 \texttt{dest[i] = round\_rm(f16(src0[i]))} (arch \textless{}= v10).
\par\textbf{Notes:}
\begin{tightitemize}
  \item Arch v9-v10 only (\texttt{kraid/ops.rs:2542}).
  \item \textbf{Source references:} \texttt{kraid/ops.rs:2542-2570}, \texttt{encode\_v9.rs:1945-1949}, \texttt{nir.rs:750-765}.
\end{tightitemize}
\codebox{\mnem{V2S16\_TO\_V2F16}\imod{\{.flow\}}\imod{\{.round\}} \dst{dest}\imod{\{.dest\_width\}}, \src{src0}\imod{\{.swz0\}}}
\textbf{Operands:}
\begin{tightitemize}\footnotesize
  \item \texttt{\dst{dest}} (\texttt{dest\_register\_index}, \texttt{dest\_width}): 32-bit destination GPR receiving the converted value.
  \item \texttt{\src{src0}}: Packed \texttt{v2s16} input converted per lane to \texttt{v2f16}.
\end{tightitemize}
\par\textbf{Closest SPIR-V Op:} \texttt{OpConvertSToF (v2s16 -\textgreater{} v2f16, v9--v10)}
\par\textbf{Pseudo-GLSL (Operational Semantics):}
\codebox{for (int i = 0; i \textless{} 2; ++i) dest[i] = round\_rm(float16\_t(int16\_t(swz0(src0)[i])));}
\textbf{Execution Context:}
\begin{tightitemize}\footnotesize
  \item \textbf{Compute stage preloads (\hyperref[sec:csf-spd]{\texttt{SPD\_k}}, \texttt{r16..r23}):} Hardware pre-fills \texttt{r55..r57} with \texttt{WorkGroupID.xyz + JOB\_OFFSET} (from \hyperref[sec:csf-regs]{\texttt{COMPUTESR::JOB\_OFFSET\_X/Y/Z}} \texttt{r34..r36}) and \texttt{r60..r62} with \texttt{LocalInvocationID.xyz} or \texttt{GlobalInvocationID.xyz} (from \texttt{WG\_SIZE} \texttt{r33} \& \texttt{JOB\_SIZE} \texttt{r37..r39}, selected by \texttt{ShaderProgram::preload}) so compute shaders can derive thread/element indices without memory loads.
  \item \textbf{Vertex stage preloads (\hyperref[sec:csf-spd]{\texttt{SPD\_0}/\texttt{SPD\_1}}, \texttt{r16..r19}):} The hardware Index Fetcher reads \hyperref[sec:csf-regs]{\texttt{IDVSSR::INDEX\_BUFFER}} (\texttt{r54:r55}) and pre-fills \texttt{r59, r60} with \texttt{VertexIndex} (\texttt{+ VERTEX\_OFFSET + GLOBAL\_ATTR\_OFF}) and \texttt{r61} with \texttt{InstanceIndex} (\texttt{+ INSTANCE\_OFFSET}) so Position and Varying shaders can pull vertex attributes directly.
  \item \textbf{Inline FAU \& Special Operands (\hyperref[sec:csf-tsd-fau]{\texttt{FAU\_k}}, \texttt{r8..r15}):} Any 8-bit source with \texttt{mode=2} reads zero-latency uniform \texttt{u0..u31} from \texttt{FAU\_k} (\texttt{slots64 = bits[63:56]}); \texttt{mode=3} reads \texttt{SmallConstant} (Table~1.2) or \texttt{FauSpecial} (\texttt{pc}, \texttt{lane\_id}, \texttt{lane\_width}, \texttt{tls\_ptr}, \texttt{wls\_ptr}, \texttt{blend\_descriptor\_0..7}).
\end{tightitemize}
\par\textbf{Worked Examples:}
\par\textbf{Example 1 (Per-lane signed int16 -\textgreater{} FP16 (v9--v10)):}
\codebox{\mnem{V2S16\_TO\_V2F16} \dst{r0}, \src{r1}\imod{.h01}}
With \texttt{r1 = 0x0002FFFF} (\texttt{[2, -1]}): converts per-lane -\textgreater{} \texttt{r0 = 0x4000BC00} (\texttt{[2.0h, -1.0h]}).
\par\textbf{Example 2 (Inexact large int16 rounding):}
\codebox{\mnem{V2S16\_TO\_V2F16}\imod{.round\_zero} \dst{r0}, \src{r1}\imod{.h01}}
With \texttt{r1 = [2049, 0]} (exceeds 11-bit FP16 significand): \texttt{.round\_zero} rounds \texttt{2049} down to \texttt{2048.0h} (\texttt{0x6800}).

\begin{center}\footnotesize
\begin{tabularx}{\linewidth}{@{}lllX@{}}
\toprule
Variant & Arch & Class & Opcode Fields \\
\midrule
\texttt{(none)} & v9--v14 & \texttt{R1} & \texttt{opcode2=V2S16\_TO\_V2F16, opcode1=CONVERT} \\
\bottomrule
\end{tabularx}
\end{center}
\begin{center}\footnotesize
\begin{tabularx}{0.88\linewidth}{@{}l l X@{}}
\toprule
Field Name & Bits & Type / Fixed Value \\
\midrule
\texttt{src0} & \texttt{0:7} & \texttt{SourceEncoding} \\
\textit{Reserved} & \texttt{8:15} & \texttt{zero} \\
\texttt{opcode2} & \texttt{16:20} & \texttt{exact=V2S16\_TO\_V2F16} \\
\textit{Reserved} & \texttt{21:29} & \texttt{zero} \\
\hyperref[enum:round_extended_m]{\texttt{round}} & \texttt{30:32} & \texttt{round\_extended\_m (def none)} \\
\textit{Reserved} & \texttt{33} & \texttt{zero} \\
\textit{Reserved} & \texttt{34:35} & \texttt{ignore} \\
\hyperref[enum:swiz_int_m]{\texttt{swz0}} & \texttt{36:39} & \texttt{swiz\_int\_m (def h01)} \\
\texttt{dest\_register\_index} & \texttt{40:45} & \texttt{uint} \\
\hyperref[enum:dest_width_m]{\texttt{dest\_width}} & \texttt{46:47} & \texttt{dest\_width\_m} \\
\texttt{opcode1} & \texttt{48:51} & \texttt{exact=CONVERT} \\
\texttt{instruction\_class} & \texttt{52:56} & \texttt{exact=R1} \\
\texttt{fau\_page\_index} & \texttt{57:58} & \texttt{uint} \\
\hyperref[enum:flow_control_m]{\texttt{flow}} & \texttt{59:62} & \texttt{flow\_control\_m} \\
\textit{Reserved} & \texttt{63} & \texttt{zero} \\
\bottomrule
\end{tabularx}
\end{center}
\textbf{Fields Explained:}
\begin{tightitemize}\footnotesize
  \item \texttt{src0}: Primary source operand (\texttt{r0..r63}, \texttt{u0..u63} uniform, \texttt{k*} FAU constant, or \texttt{t0..t2} temporary; \texttt{\textasciicircum{}} marks discard/last-use). \emph{Example:} {\scriptsize\ttfamily \mnem{FADD.f32} \dst{r0}, \src{r1}, \src{r2}} (\textbf{SPIR-V:} \texttt{OpFAdd (Operand 0)}; \textbf{GLSL:} \texttt{float c = a + b;})
  \item \texttt{opcode2}: Secondary sub-opcode further specializing instructions sharing \texttt{instruction\_class} and \texttt{opcode1} (e.g., conversions, SFU tables, moves). \emph{Example:} {\scriptsize\ttfamily \mnem{F32\_TO\_S32}\imod{.round\_zero} \dst{r0}, \src{r1}} (\textbf{SPIR-V:} \texttt{OpConvertFToS}; \textbf{GLSL:} \texttt{int y = int(x);})
  \item \hyperref[enum:round_extended_m]{\texttt{round}}: Rounding mode (\texttt{none}/RNE, \texttt{round\_up}, \texttt{round\_down}, \texttt{round\_zero}, \texttt{round\_na}) applied in FP ALU, conversions, \texttt{HADD}, \texttt{LDEXP}, and Stage C narrowing. \emph{Example:} {\scriptsize\ttfamily \mnem{FROUND.f32}\imod{.round\_down} \dst{r0}, \src{r1}} (\textbf{SPIR-V:} \texttt{GLSL.std.450 Floor / Ceil / Trunc}; \textbf{GLSL:} \texttt{float y = floor(x);})
  \item \hyperref[enum:swiz_int_m]{\texttt{swz0}}: Stage A lane swizzle on \texttt{src0} (\texttt{h00}/\texttt{h01}/\texttt{h10}/\texttt{h11} for 16-bit vectors, or \texttt{bXY} byte-to-halfword widening). \emph{Example:} {\scriptsize\ttfamily \mnem{FADD.v2f16} \dst{r0}, \src{r1}\imod{.h10}, \src{r2}} (\textbf{SPIR-V:} \texttt{OpVectorShuffle}; \textbf{GLSL:} \texttt{f16vec2 c = a.yx + b;})
  \item \texttt{dest\_register\_index}: 6-bit destination GPR index (\texttt{r0..r63}) receiving the Stage C writeback result (even-aligned pair \texttt{[rN:rN+1]} for 64-bit ops). \emph{Example:} {\scriptsize\ttfamily \mnem{FMA.f32} \dst{r0}, \src{r1}, \src{r2}, \src{r3}} (\textbf{SPIR-V:} \texttt{SSA Result ID}; \textbf{GLSL:} \texttt{float r0 = fma(r1, r2, r3);})
  \item \hyperref[enum:dest_width_m]{\texttt{dest\_width}}: Stage C destination lane write-mask (\texttt{none} = full 32-bit, \texttt{h0} = low 16-bit, \texttt{h1} = high 16-bit without bit-shifting). \emph{Example:} {\scriptsize\ttfamily \mnem{FADD.v2f16} \dst{r0}\imod{.h1}, \src{r1}, \src{r2}} (\textbf{SPIR-V:} \texttt{OpVectorInsertDynamic / WriteMask}; \textbf{GLSL:} \texttt{dst.y = a.y + b.y;})
  \item \texttt{opcode1}: Primary sub-opcode within the instruction class uniquely identifying the operation family. \emph{Example:} {\scriptsize\ttfamily \mnem{IADD.s32} \dst{r0}, \src{r1}, \src{r2}} (\textbf{SPIR-V:} \texttt{OpIAdd (operation selector)}; \textbf{GLSL:} \texttt{// Primary opcode bits})
  \item \texttt{instruction\_class}: 4-bit major instruction class selector (\texttt{bits 48:51}) partitioning the 64-bit Valhall opcode space (e.g., \texttt{FMA}, \texttt{CVT}, \texttt{SFU}, \texttt{LS}, \texttt{TEX}, \texttt{BRANCH}). \emph{Example:} {\scriptsize\ttfamily \mnem{FADD.f32} \dst{r0}, \src{r1}, \src{r2}} (\textbf{SPIR-V:} \texttt{N/A (encoding class)}; \textbf{GLSL:} \texttt{// Hardware execution pipe / decode class})
  \item \texttt{fau\_page\_index}: Selects the FAU (Fast Access Uniform) page (\texttt{0..3}) when a source operand references hardware small constants, thread IDs, or lane constants. \emph{Example:} {\scriptsize\ttfamily \mnem{FADD.f32} \dst{r0}, \src{r1}, \src{k1f}} (\textbf{SPIR-V:} \texttt{OpConstant / BuiltIn}; \textbf{GLSL:} \texttt{float y = x + 1.0;})
  \item \hyperref[enum:flow_control_m]{\texttt{flow}}: 4-bit embedded flow-control \& scoreboard wait modifier (\texttt{bits 59:62}: \texttt{none}, \texttt{wait0..wait0126}, \texttt{wait}, \texttt{reconverge}, \texttt{discard}, \texttt{end}). \emph{Example:} {\scriptsize\ttfamily \mnem{FADD.f32}\imod{.wait0.reconverge} \dst{r0}, \src{r1}, \src{r2}} (\textbf{SPIR-V:} \texttt{OpControlBarrier / OpReturn}; \textbf{GLSL:} \texttt{// Scoreboard wait + warp reconvergence})
\end{tightitemize}
\clearpage

\subsection*{V2S8\_TO\_V2F16}
\addcontentsline{toc}{subsection}{V2S8\_TO\_V2F16}
Extracts two signed 8-bit bytes selected by \texttt{swz0} (\texttt{bXY}, default \texttt{b02}) from \texttt{src0} and converts them exactly to packed FP16 lanes \texttt{dest.h0 = f16(s8(src0.b[X]))} and \texttt{dest.h1 = f16(s8(src0.b[Y]))} (arch \textless{}= v10).
\par\textbf{Notes:}
\begin{tightitemize}
  \item Arch v9-v10 only (\texttt{kraid/ops.rs:2542}, \texttt{encode\_v9.rs:1933}).
  \item \textbf{Source references:} \texttt{kraid/ops.rs:2542-2570}, \texttt{encode\_v9.rs:1937-1940}.
\end{tightitemize}
\codebox{\mnem{V2S8\_TO\_V2F16}\imod{\{.flow\}} \dst{dest}\imod{\{.dest\_width\}}, \src{src0}\imod{\{.swz0\}}}
\textbf{Operands:}
\begin{tightitemize}\footnotesize
  \item \texttt{\dst{dest}} (\texttt{dest\_register\_index}, \texttt{dest\_width}): 32-bit destination GPR receiving the converted value.
  \item \texttt{\src{src0}}: Packed \texttt{v4s8} input (two bytes selected by \texttt{swz0}, default \texttt{b02}) converted to \texttt{v2f16}.
\end{tightitemize}
\par\textbf{Closest SPIR-V Op:} \texttt{OpConvertSToF (v2s8 -\textgreater{} v2f16, v9--v10)}
\par\textbf{Pseudo-GLSL (Operational Semantics):}
\codebox{dest = f16vec2(float16\_t(s8(src0, swz0.x)), float16\_t(s8(src0, swz0.y)));}
\textbf{Execution Context:}
\begin{tightitemize}\footnotesize
  \item \textbf{Compute stage preloads (\hyperref[sec:csf-spd]{\texttt{SPD\_k}}, \texttt{r16..r23}):} Hardware pre-fills \texttt{r55..r57} with \texttt{WorkGroupID.xyz + JOB\_OFFSET} (from \hyperref[sec:csf-regs]{\texttt{COMPUTESR::JOB\_OFFSET\_X/Y/Z}} \texttt{r34..r36}) and \texttt{r60..r62} with \texttt{LocalInvocationID.xyz} or \texttt{GlobalInvocationID.xyz} (from \texttt{WG\_SIZE} \texttt{r33} \& \texttt{JOB\_SIZE} \texttt{r37..r39}, selected by \texttt{ShaderProgram::preload}) so compute shaders can derive thread/element indices without memory loads.
  \item \textbf{Vertex stage preloads (\hyperref[sec:csf-spd]{\texttt{SPD\_0}/\texttt{SPD\_1}}, \texttt{r16..r19}):} The hardware Index Fetcher reads \hyperref[sec:csf-regs]{\texttt{IDVSSR::INDEX\_BUFFER}} (\texttt{r54:r55}) and pre-fills \texttt{r59, r60} with \texttt{VertexIndex} (\texttt{+ VERTEX\_OFFSET + GLOBAL\_ATTR\_OFF}) and \texttt{r61} with \texttt{InstanceIndex} (\texttt{+ INSTANCE\_OFFSET}) so Position and Varying shaders can pull vertex attributes directly.
  \item \textbf{Inline FAU \& Special Operands (\hyperref[sec:csf-tsd-fau]{\texttt{FAU\_k}}, \texttt{r8..r15}):} Any 8-bit source with \texttt{mode=2} reads zero-latency uniform \texttt{u0..u31} from \texttt{FAU\_k} (\texttt{slots64 = bits[63:56]}); \texttt{mode=3} reads \texttt{SmallConstant} (Table~1.2) or \texttt{FauSpecial} (\texttt{pc}, \texttt{lane\_id}, \texttt{lane\_width}, \texttt{tls\_ptr}, \texttt{wls\_ptr}, \texttt{blend\_descriptor\_0..7}).
\end{tightitemize}
\par\textbf{Worked Examples:}
\par\textbf{Example 1 (Default `b02` stride extraction (v9--v10)):}
\codebox{\mnem{V2S8\_TO\_V2F16} \dst{r0}, \src{r1}\imod{.b02}}
Default \texttt{b02} extracts bytes 0 \& 2: with \texttt{r1 = 0x00FE0001} (\texttt{b2=-2, b0=+1}), yields \texttt{r0 = 0xC0003C00} (\texttt{[-2.0h, 1.0h]}).
\par\textbf{Example 2 (Adjacent bytes `b01` extraction):}
\codebox{\mnem{V2S8\_TO\_V2F16} \dst{r0}, \src{r1}\imod{.b01}}
With \texttt{r1 = 0x000004FC} (\texttt{b1=+4, b0=-4}): yields \texttt{r0 = 0x4400C400} (\texttt{[+4.0h, -4.0h]}).

\begin{center}\footnotesize
\begin{tabularx}{\linewidth}{@{}lllX@{}}
\toprule
Variant & Arch & Class & Opcode Fields \\
\midrule
\texttt{(none)} & v9--v14 & \texttt{R1} & \texttt{opcode2=V2S8\_TO\_V2F16, opcode1=CONVERT} \\
\bottomrule
\end{tabularx}
\end{center}
\begin{center}\footnotesize
\begin{tabularx}{0.88\linewidth}{@{}l l X@{}}
\toprule
Field Name & Bits & Type / Fixed Value \\
\midrule
\texttt{src0} & \texttt{0:7} & \texttt{SourceEncoding} \\
\textit{Reserved} & \texttt{8:15} & \texttt{zero} \\
\texttt{opcode2} & \texttt{16:20} & \texttt{exact=V2S8\_TO\_V2F16} \\
\textit{Reserved} & \texttt{21:35} & \texttt{zero} \\
\hyperref[enum:swiz8_2_full_m]{\texttt{swz0}} & \texttt{36:39} & \texttt{swiz8\_2\_full\_m (def b02)} \\
\texttt{dest\_register\_index} & \texttt{40:45} & \texttt{uint} \\
\hyperref[enum:dest_width_m]{\texttt{dest\_width}} & \texttt{46:47} & \texttt{dest\_width\_m} \\
\texttt{opcode1} & \texttt{48:51} & \texttt{exact=CONVERT} \\
\texttt{instruction\_class} & \texttt{52:56} & \texttt{exact=R1} \\
\texttt{fau\_page\_index} & \texttt{57:58} & \texttt{uint} \\
\hyperref[enum:flow_control_m]{\texttt{flow}} & \texttt{59:62} & \texttt{flow\_control\_m} \\
\textit{Reserved} & \texttt{63} & \texttt{zero} \\
\bottomrule
\end{tabularx}
\end{center}
\textbf{Fields Explained:}
\begin{tightitemize}\footnotesize
  \item \texttt{src0}: Primary source operand (\texttt{r0..r63}, \texttt{u0..u63} uniform, \texttt{k*} FAU constant, or \texttt{t0..t2} temporary; \texttt{\textasciicircum{}} marks discard/last-use). \emph{Example:} {\scriptsize\ttfamily \mnem{FADD.f32} \dst{r0}, \src{r1}, \src{r2}} (\textbf{SPIR-V:} \texttt{OpFAdd (Operand 0)}; \textbf{GLSL:} \texttt{float c = a + b;})
  \item \texttt{opcode2}: Secondary sub-opcode further specializing instructions sharing \texttt{instruction\_class} and \texttt{opcode1} (e.g., conversions, SFU tables, moves). \emph{Example:} {\scriptsize\ttfamily \mnem{F32\_TO\_S32}\imod{.round\_zero} \dst{r0}, \src{r1}} (\textbf{SPIR-V:} \texttt{OpConvertFToS}; \textbf{GLSL:} \texttt{int y = int(x);})
  \item \hyperref[enum:swiz8_2_full_m]{\texttt{swz0}}: Stage A lane swizzle on \texttt{src0} (\texttt{h00}/\texttt{h01}/\texttt{h10}/\texttt{h11} for 16-bit vectors, or \texttt{bXY} byte-to-halfword widening). \emph{Example:} {\scriptsize\ttfamily \mnem{FADD.v2f16} \dst{r0}, \src{r1}\imod{.h10}, \src{r2}} (\textbf{SPIR-V:} \texttt{OpVectorShuffle}; \textbf{GLSL:} \texttt{f16vec2 c = a.yx + b;})
  \item \texttt{dest\_register\_index}: 6-bit destination GPR index (\texttt{r0..r63}) receiving the Stage C writeback result (even-aligned pair \texttt{[rN:rN+1]} for 64-bit ops). \emph{Example:} {\scriptsize\ttfamily \mnem{FMA.f32} \dst{r0}, \src{r1}, \src{r2}, \src{r3}} (\textbf{SPIR-V:} \texttt{SSA Result ID}; \textbf{GLSL:} \texttt{float r0 = fma(r1, r2, r3);})
  \item \hyperref[enum:dest_width_m]{\texttt{dest\_width}}: Stage C destination lane write-mask (\texttt{none} = full 32-bit, \texttt{h0} = low 16-bit, \texttt{h1} = high 16-bit without bit-shifting). \emph{Example:} {\scriptsize\ttfamily \mnem{FADD.v2f16} \dst{r0}\imod{.h1}, \src{r1}, \src{r2}} (\textbf{SPIR-V:} \texttt{OpVectorInsertDynamic / WriteMask}; \textbf{GLSL:} \texttt{dst.y = a.y + b.y;})
  \item \texttt{opcode1}: Primary sub-opcode within the instruction class uniquely identifying the operation family. \emph{Example:} {\scriptsize\ttfamily \mnem{IADD.s32} \dst{r0}, \src{r1}, \src{r2}} (\textbf{SPIR-V:} \texttt{OpIAdd (operation selector)}; \textbf{GLSL:} \texttt{// Primary opcode bits})
  \item \texttt{instruction\_class}: 4-bit major instruction class selector (\texttt{bits 48:51}) partitioning the 64-bit Valhall opcode space (e.g., \texttt{FMA}, \texttt{CVT}, \texttt{SFU}, \texttt{LS}, \texttt{TEX}, \texttt{BRANCH}). \emph{Example:} {\scriptsize\ttfamily \mnem{FADD.f32} \dst{r0}, \src{r1}, \src{r2}} (\textbf{SPIR-V:} \texttt{N/A (encoding class)}; \textbf{GLSL:} \texttt{// Hardware execution pipe / decode class})
  \item \texttt{fau\_page\_index}: Selects the FAU (Fast Access Uniform) page (\texttt{0..3}) when a source operand references hardware small constants, thread IDs, or lane constants. \emph{Example:} {\scriptsize\ttfamily \mnem{FADD.f32} \dst{r0}, \src{r1}, \src{k1f}} (\textbf{SPIR-V:} \texttt{OpConstant / BuiltIn}; \textbf{GLSL:} \texttt{float y = x + 1.0;})
  \item \hyperref[enum:flow_control_m]{\texttt{flow}}: 4-bit embedded flow-control \& scoreboard wait modifier (\texttt{bits 59:62}: \texttt{none}, \texttt{wait0..wait0126}, \texttt{wait}, \texttt{reconverge}, \texttt{discard}, \texttt{end}). \emph{Example:} {\scriptsize\ttfamily \mnem{FADD.f32}\imod{.wait0.reconverge} \dst{r0}, \src{r1}, \src{r2}} (\textbf{SPIR-V:} \texttt{OpControlBarrier / OpReturn}; \textbf{GLSL:} \texttt{// Scoreboard wait + warp reconvergence})
\end{tightitemize}
\clearpage

\subsection*{V2S8\_TO\_V2S16}
\addcontentsline{toc}{subsection}{V2S8\_TO\_V2S16}
Selects two signed 8-bit bytes from \texttt{src0} via \texttt{swz0} (default \texttt{b02}: bytes 0 and 2), sign-extends each to 16 bits, and writes the packed \texttt{v2s16} pair to \texttt{dest}.
\par\textbf{Notes:}
\begin{tightitemize}
  \item Exact per-lane sign extension.
  \item \textbf{Source references:} \texttt{isa-v9-v14.xml}, \texttt{bifrost/valhall/ISA.xml}.
\end{tightitemize}
\codebox{\mnem{V2S8\_TO\_V2S16}\imod{\{.flow\}} \dst{dest}\imod{\{.dest\_width\}}, \src{src0}\imod{\{.swz0\}}}
\textbf{Operands:}
\begin{tightitemize}\footnotesize
  \item \texttt{\dst{dest}} (\texttt{dest\_register\_index}, \texttt{dest\_width}): 32-bit destination GPR receiving the converted value.
  \item \texttt{\src{src0}}: Packed \texttt{v4s8} input (two bytes selected by \texttt{swz0}, default \texttt{b02}) sign-extended to \texttt{v2s16}.
\end{tightitemize}
\par\textbf{Closest SPIR-V Op:} \texttt{OpSConvert (v2s8 -\textgreater{} v2s16)}
\par\textbf{Pseudo-GLSL (Operational Semantics):}
\codebox{dest = i16vec2(int16\_t(s8(src0, swz0.x)), int16\_t(s8(src0, swz0.y)));}
\textbf{Execution Context:}
\begin{tightitemize}\footnotesize
  \item \textbf{Compute stage preloads (\hyperref[sec:csf-spd]{\texttt{SPD\_k}}, \texttt{r16..r23}):} Hardware pre-fills \texttt{r55..r57} with \texttt{WorkGroupID.xyz + JOB\_OFFSET} (from \hyperref[sec:csf-regs]{\texttt{COMPUTESR::JOB\_OFFSET\_X/Y/Z}} \texttt{r34..r36}) and \texttt{r60..r62} with \texttt{LocalInvocationID.xyz} or \texttt{GlobalInvocationID.xyz} (from \texttt{WG\_SIZE} \texttt{r33} \& \texttt{JOB\_SIZE} \texttt{r37..r39}, selected by \texttt{ShaderProgram::preload}) so compute shaders can derive thread/element indices without memory loads.
  \item \textbf{Vertex stage preloads (\hyperref[sec:csf-spd]{\texttt{SPD\_0}/\texttt{SPD\_1}}, \texttt{r16..r19}):} The hardware Index Fetcher reads \hyperref[sec:csf-regs]{\texttt{IDVSSR::INDEX\_BUFFER}} (\texttt{r54:r55}) and pre-fills \texttt{r59, r60} with \texttt{VertexIndex} (\texttt{+ VERTEX\_OFFSET + GLOBAL\_ATTR\_OFF}) and \texttt{r61} with \texttt{InstanceIndex} (\texttt{+ INSTANCE\_OFFSET}) so Position and Varying shaders can pull vertex attributes directly.
  \item \textbf{Inline FAU \& Special Operands (\hyperref[sec:csf-tsd-fau]{\texttt{FAU\_k}}, \texttt{r8..r15}):} Any 8-bit source with \texttt{mode=2} reads zero-latency uniform \texttt{u0..u31} from \texttt{FAU\_k} (\texttt{slots64 = bits[63:56]}); \texttt{mode=3} reads \texttt{SmallConstant} (Table~1.2) or \texttt{FauSpecial} (\texttt{pc}, \texttt{lane\_id}, \texttt{lane\_width}, \texttt{tls\_ptr}, \texttt{wls\_ptr}, \texttt{blend\_descriptor\_0..7}).
\end{tightitemize}
\par\textbf{Worked Examples:}
\par\textbf{Example 1 (Default `b02` sign extension):}
\codebox{\mnem{V2S8\_TO\_V2S16} \dst{r0}, \src{r1}\imod{.b02}}
With \texttt{r1 = 0x00FE00FF} (\texttt{b2 = -2, b0 = -1}): sign-extends bytes 0 \& 2 -\textgreater{} \texttt{r0 = 0xFFFEFFFF}.
\par\textbf{Example 2 (Broadcast single byte `b33`):}
\codebox{\mnem{V2S8\_TO\_V2S16} \dst{r0}, \src{r1}\imod{.b33}}
With \texttt{r1 = 0x80000000} (\texttt{b3 = -128}): sign-extends byte 3 into both halfwords -\textgreater{} \texttt{r0 = 0xFF80FF80}.

\begin{center}\footnotesize
\begin{tabularx}{\linewidth}{@{}lllX@{}}
\toprule
Variant & Arch & Class & Opcode Fields \\
\midrule
\texttt{(none)} & v9--v14 & \texttt{R1} & \texttt{opcode2=V2S8\_TO\_V2S16, opcode1=CONVERT} \\
\texttt{(none)} & v9--v14 & \texttt{R1} & \texttt{opcode2=V2I8\_TO\_V2I16, opcode1=CONVERT} \\
\bottomrule
\end{tabularx}
\end{center}
\begin{center}\footnotesize
\begin{tabularx}{0.88\linewidth}{@{}l l X@{}}
\toprule
Field Name & Bits & Type / Fixed Value \\
\midrule
\texttt{src0} & \texttt{0:7} & \texttt{SourceEncoding} \\
\textit{Reserved} & \texttt{8:15} & \texttt{zero} \\
\texttt{opcode2} & \texttt{16:20} & \texttt{exact=V2S8\_TO\_V2S16} \\
\textit{Reserved} & \texttt{21:35} & \texttt{zero} \\
\hyperref[enum:swiz8_2_full_m]{\texttt{swz0}} & \texttt{36:39} & \texttt{swiz8\_2\_full\_m (def b02)} \\
\texttt{dest\_register\_index} & \texttt{40:45} & \texttt{uint} \\
\hyperref[enum:dest_width_m]{\texttt{dest\_width}} & \texttt{46:47} & \texttt{dest\_width\_m} \\
\texttt{opcode1} & \texttt{48:51} & \texttt{exact=CONVERT} \\
\texttt{instruction\_class} & \texttt{52:56} & \texttt{exact=R1} \\
\texttt{fau\_page\_index} & \texttt{57:58} & \texttt{uint} \\
\hyperref[enum:flow_control_m]{\texttt{flow}} & \texttt{59:62} & \texttt{flow\_control\_m} \\
\textit{Reserved} & \texttt{63} & \texttt{zero} \\
\bottomrule
\end{tabularx}
\end{center}
\textbf{Fields Explained:}
\begin{tightitemize}\footnotesize
  \item \texttt{src0}: Primary source operand (\texttt{r0..r63}, \texttt{u0..u63} uniform, \texttt{k*} FAU constant, or \texttt{t0..t2} temporary; \texttt{\textasciicircum{}} marks discard/last-use). \emph{Example:} {\scriptsize\ttfamily \mnem{FADD.f32} \dst{r0}, \src{r1}, \src{r2}} (\textbf{SPIR-V:} \texttt{OpFAdd (Operand 0)}; \textbf{GLSL:} \texttt{float c = a + b;})
  \item \texttt{opcode2}: Secondary sub-opcode further specializing instructions sharing \texttt{instruction\_class} and \texttt{opcode1} (e.g., conversions, SFU tables, moves). \emph{Example:} {\scriptsize\ttfamily \mnem{F32\_TO\_S32}\imod{.round\_zero} \dst{r0}, \src{r1}} (\textbf{SPIR-V:} \texttt{OpConvertFToS}; \textbf{GLSL:} \texttt{int y = int(x);})
  \item \hyperref[enum:swiz8_2_full_m]{\texttt{swz0}}: Stage A lane swizzle on \texttt{src0} (\texttt{h00}/\texttt{h01}/\texttt{h10}/\texttt{h11} for 16-bit vectors, or \texttt{bXY} byte-to-halfword widening). \emph{Example:} {\scriptsize\ttfamily \mnem{FADD.v2f16} \dst{r0}, \src{r1}\imod{.h10}, \src{r2}} (\textbf{SPIR-V:} \texttt{OpVectorShuffle}; \textbf{GLSL:} \texttt{f16vec2 c = a.yx + b;})
  \item \texttt{dest\_register\_index}: 6-bit destination GPR index (\texttt{r0..r63}) receiving the Stage C writeback result (even-aligned pair \texttt{[rN:rN+1]} for 64-bit ops). \emph{Example:} {\scriptsize\ttfamily \mnem{FMA.f32} \dst{r0}, \src{r1}, \src{r2}, \src{r3}} (\textbf{SPIR-V:} \texttt{SSA Result ID}; \textbf{GLSL:} \texttt{float r0 = fma(r1, r2, r3);})
  \item \hyperref[enum:dest_width_m]{\texttt{dest\_width}}: Stage C destination lane write-mask (\texttt{none} = full 32-bit, \texttt{h0} = low 16-bit, \texttt{h1} = high 16-bit without bit-shifting). \emph{Example:} {\scriptsize\ttfamily \mnem{FADD.v2f16} \dst{r0}\imod{.h1}, \src{r1}, \src{r2}} (\textbf{SPIR-V:} \texttt{OpVectorInsertDynamic / WriteMask}; \textbf{GLSL:} \texttt{dst.y = a.y + b.y;})
  \item \texttt{opcode1}: Primary sub-opcode within the instruction class uniquely identifying the operation family. \emph{Example:} {\scriptsize\ttfamily \mnem{IADD.s32} \dst{r0}, \src{r1}, \src{r2}} (\textbf{SPIR-V:} \texttt{OpIAdd (operation selector)}; \textbf{GLSL:} \texttt{// Primary opcode bits})
  \item \texttt{instruction\_class}: 4-bit major instruction class selector (\texttt{bits 48:51}) partitioning the 64-bit Valhall opcode space (e.g., \texttt{FMA}, \texttt{CVT}, \texttt{SFU}, \texttt{LS}, \texttt{TEX}, \texttt{BRANCH}). \emph{Example:} {\scriptsize\ttfamily \mnem{FADD.f32} \dst{r0}, \src{r1}, \src{r2}} (\textbf{SPIR-V:} \texttt{N/A (encoding class)}; \textbf{GLSL:} \texttt{// Hardware execution pipe / decode class})
  \item \texttt{fau\_page\_index}: Selects the FAU (Fast Access Uniform) page (\texttt{0..3}) when a source operand references hardware small constants, thread IDs, or lane constants. \emph{Example:} {\scriptsize\ttfamily \mnem{FADD.f32} \dst{r0}, \src{r1}, \src{k1f}} (\textbf{SPIR-V:} \texttt{OpConstant / BuiltIn}; \textbf{GLSL:} \texttt{float y = x + 1.0;})
  \item \hyperref[enum:flow_control_m]{\texttt{flow}}: 4-bit embedded flow-control \& scoreboard wait modifier (\texttt{bits 59:62}: \texttt{none}, \texttt{wait0..wait0126}, \texttt{wait}, \texttt{reconverge}, \texttt{discard}, \texttt{end}). \emph{Example:} {\scriptsize\ttfamily \mnem{FADD.f32}\imod{.wait0.reconverge} \dst{r0}, \src{r1}, \src{r2}} (\textbf{SPIR-V:} \texttt{OpControlBarrier / OpReturn}; \textbf{GLSL:} \texttt{// Scoreboard wait + warp reconvergence})
\end{tightitemize}
\clearpage

\subsection*{V2U16\_TO\_V2F16}
\addcontentsline{toc}{subsection}{V2U16\_TO\_V2F16}
Converts each unsigned 16-bit integer lane $i \in \{0, 1\}$ of swizzled \texttt{src0} to IEEE-754 FP16 \texttt{dest[i] = round\_rm(f16(src0[i]))} (arch \textless{}= v10).
\par\textbf{Notes:}
\begin{tightitemize}
  \item Arch v9-v10 only.
  \item Values \texttt{\textgreater{}= 65520} overflow to \texttt{+Inf} in RNE.
  \item \textbf{Source references:} \texttt{kraid/ops.rs:2542-2570}, \texttt{encode\_v9.rs:1950-1954}, \texttt{nir.rs:750-765}.
\end{tightitemize}
\codebox{\mnem{V2U16\_TO\_V2F16}\imod{\{.flow\}}\imod{\{.round\}} \dst{dest}\imod{\{.dest\_width\}}, \src{src0}\imod{\{.swz0\}}}
\textbf{Operands:}
\begin{tightitemize}\footnotesize
  \item \texttt{\dst{dest}} (\texttt{dest\_register\_index}, \texttt{dest\_width}): 32-bit destination GPR receiving the converted value.
  \item \texttt{\src{src0}}: Packed \texttt{v2u16} input converted per lane to \texttt{v2f16}.
\end{tightitemize}
\par\textbf{Closest SPIR-V Op:} \texttt{OpConvertUToF (v2u16 -\textgreater{} v2f16, v9--v10)}
\par\textbf{Pseudo-GLSL (Operational Semantics):}
\codebox{for (int i = 0; i \textless{} 2; ++i) dest[i] = round\_rm(float16\_t(uint16\_t(swz0(src0)[i])));}
\textbf{Execution Context:}
\begin{tightitemize}\footnotesize
  \item \textbf{Compute stage preloads (\hyperref[sec:csf-spd]{\texttt{SPD\_k}}, \texttt{r16..r23}):} Hardware pre-fills \texttt{r55..r57} with \texttt{WorkGroupID.xyz + JOB\_OFFSET} (from \hyperref[sec:csf-regs]{\texttt{COMPUTESR::JOB\_OFFSET\_X/Y/Z}} \texttt{r34..r36}) and \texttt{r60..r62} with \texttt{LocalInvocationID.xyz} or \texttt{GlobalInvocationID.xyz} (from \texttt{WG\_SIZE} \texttt{r33} \& \texttt{JOB\_SIZE} \texttt{r37..r39}, selected by \texttt{ShaderProgram::preload}) so compute shaders can derive thread/element indices without memory loads.
  \item \textbf{Vertex stage preloads (\hyperref[sec:csf-spd]{\texttt{SPD\_0}/\texttt{SPD\_1}}, \texttt{r16..r19}):} The hardware Index Fetcher reads \hyperref[sec:csf-regs]{\texttt{IDVSSR::INDEX\_BUFFER}} (\texttt{r54:r55}) and pre-fills \texttt{r59, r60} with \texttt{VertexIndex} (\texttt{+ VERTEX\_OFFSET + GLOBAL\_ATTR\_OFF}) and \texttt{r61} with \texttt{InstanceIndex} (\texttt{+ INSTANCE\_OFFSET}) so Position and Varying shaders can pull vertex attributes directly.
  \item \textbf{Inline FAU \& Special Operands (\hyperref[sec:csf-tsd-fau]{\texttt{FAU\_k}}, \texttt{r8..r15}):} Any 8-bit source with \texttt{mode=2} reads zero-latency uniform \texttt{u0..u31} from \texttt{FAU\_k} (\texttt{slots64 = bits[63:56]}); \texttt{mode=3} reads \texttt{SmallConstant} (Table~1.2) or \texttt{FauSpecial} (\texttt{pc}, \texttt{lane\_id}, \texttt{lane\_width}, \texttt{tls\_ptr}, \texttt{wls\_ptr}, \texttt{blend\_descriptor\_0..7}).
\end{tightitemize}
\par\textbf{Worked Examples:}
\par\textbf{Example 1 (Normal uint16 pair (v9--v10)):}
\codebox{\mnem{V2U16\_TO\_V2F16} \dst{r0}, \src{r1}\imod{.h01}}
With \texttt{r1 = 0x00040001} (\texttt{[4, 1]}): yields \texttt{r0 = 0x44003C00} (\texttt{[4.0h, 1.0h]}).
\par\textbf{Example 2 (Overflow to `+Inf` above `65519`):}
\codebox{\mnem{V2U16\_TO\_V2F16} \dst{r0}, \src{r1}\imod{.h01}}
With \texttt{r1 = 0xFFFF0000} (\texttt{[65535, 0]}): under RNE, \texttt{65535 \textgreater{}= 65520} overflows FP16 to \texttt{+Inf} (\texttt{0x7C00}) -\textgreater{} \texttt{r0 = 0x7C000000}.

\begin{center}\footnotesize
\begin{tabularx}{\linewidth}{@{}lllX@{}}
\toprule
Variant & Arch & Class & Opcode Fields \\
\midrule
\texttt{(none)} & v9--v14 & \texttt{R1} & \texttt{opcode2=V2U16\_TO\_V2F16, opcode1=CONVERT} \\
\bottomrule
\end{tabularx}
\end{center}
\begin{center}\footnotesize
\begin{tabularx}{0.88\linewidth}{@{}l l X@{}}
\toprule
Field Name & Bits & Type / Fixed Value \\
\midrule
\texttt{src0} & \texttt{0:7} & \texttt{SourceEncoding} \\
\textit{Reserved} & \texttt{8:15} & \texttt{zero} \\
\texttt{opcode2} & \texttt{16:20} & \texttt{exact=V2U16\_TO\_V2F16} \\
\textit{Reserved} & \texttt{21:29} & \texttt{zero} \\
\hyperref[enum:round_extended_m]{\texttt{round}} & \texttt{30:32} & \texttt{round\_extended\_m (def none)} \\
\textit{Reserved} & \texttt{33} & \texttt{zero} \\
\textit{Reserved} & \texttt{34:35} & \texttt{ignore} \\
\hyperref[enum:swiz_int_m]{\texttt{swz0}} & \texttt{36:39} & \texttt{swiz\_int\_m (def h01)} \\
\texttt{dest\_register\_index} & \texttt{40:45} & \texttt{uint} \\
\hyperref[enum:dest_width_m]{\texttt{dest\_width}} & \texttt{46:47} & \texttt{dest\_width\_m} \\
\texttt{opcode1} & \texttt{48:51} & \texttt{exact=CONVERT} \\
\texttt{instruction\_class} & \texttt{52:56} & \texttt{exact=R1} \\
\texttt{fau\_page\_index} & \texttt{57:58} & \texttt{uint} \\
\hyperref[enum:flow_control_m]{\texttt{flow}} & \texttt{59:62} & \texttt{flow\_control\_m} \\
\textit{Reserved} & \texttt{63} & \texttt{zero} \\
\bottomrule
\end{tabularx}
\end{center}
\textbf{Fields Explained:}
\begin{tightitemize}\footnotesize
  \item \texttt{src0}: Primary source operand (\texttt{r0..r63}, \texttt{u0..u63} uniform, \texttt{k*} FAU constant, or \texttt{t0..t2} temporary; \texttt{\textasciicircum{}} marks discard/last-use). \emph{Example:} {\scriptsize\ttfamily \mnem{FADD.f32} \dst{r0}, \src{r1}, \src{r2}} (\textbf{SPIR-V:} \texttt{OpFAdd (Operand 0)}; \textbf{GLSL:} \texttt{float c = a + b;})
  \item \texttt{opcode2}: Secondary sub-opcode further specializing instructions sharing \texttt{instruction\_class} and \texttt{opcode1} (e.g., conversions, SFU tables, moves). \emph{Example:} {\scriptsize\ttfamily \mnem{F32\_TO\_S32}\imod{.round\_zero} \dst{r0}, \src{r1}} (\textbf{SPIR-V:} \texttt{OpConvertFToS}; \textbf{GLSL:} \texttt{int y = int(x);})
  \item \hyperref[enum:round_extended_m]{\texttt{round}}: Rounding mode (\texttt{none}/RNE, \texttt{round\_up}, \texttt{round\_down}, \texttt{round\_zero}, \texttt{round\_na}) applied in FP ALU, conversions, \texttt{HADD}, \texttt{LDEXP}, and Stage C narrowing. \emph{Example:} {\scriptsize\ttfamily \mnem{FROUND.f32}\imod{.round\_down} \dst{r0}, \src{r1}} (\textbf{SPIR-V:} \texttt{GLSL.std.450 Floor / Ceil / Trunc}; \textbf{GLSL:} \texttt{float y = floor(x);})
  \item \hyperref[enum:swiz_int_m]{\texttt{swz0}}: Stage A lane swizzle on \texttt{src0} (\texttt{h00}/\texttt{h01}/\texttt{h10}/\texttt{h11} for 16-bit vectors, or \texttt{bXY} byte-to-halfword widening). \emph{Example:} {\scriptsize\ttfamily \mnem{FADD.v2f16} \dst{r0}, \src{r1}\imod{.h10}, \src{r2}} (\textbf{SPIR-V:} \texttt{OpVectorShuffle}; \textbf{GLSL:} \texttt{f16vec2 c = a.yx + b;})
  \item \texttt{dest\_register\_index}: 6-bit destination GPR index (\texttt{r0..r63}) receiving the Stage C writeback result (even-aligned pair \texttt{[rN:rN+1]} for 64-bit ops). \emph{Example:} {\scriptsize\ttfamily \mnem{FMA.f32} \dst{r0}, \src{r1}, \src{r2}, \src{r3}} (\textbf{SPIR-V:} \texttt{SSA Result ID}; \textbf{GLSL:} \texttt{float r0 = fma(r1, r2, r3);})
  \item \hyperref[enum:dest_width_m]{\texttt{dest\_width}}: Stage C destination lane write-mask (\texttt{none} = full 32-bit, \texttt{h0} = low 16-bit, \texttt{h1} = high 16-bit without bit-shifting). \emph{Example:} {\scriptsize\ttfamily \mnem{FADD.v2f16} \dst{r0}\imod{.h1}, \src{r1}, \src{r2}} (\textbf{SPIR-V:} \texttt{OpVectorInsertDynamic / WriteMask}; \textbf{GLSL:} \texttt{dst.y = a.y + b.y;})
  \item \texttt{opcode1}: Primary sub-opcode within the instruction class uniquely identifying the operation family. \emph{Example:} {\scriptsize\ttfamily \mnem{IADD.s32} \dst{r0}, \src{r1}, \src{r2}} (\textbf{SPIR-V:} \texttt{OpIAdd (operation selector)}; \textbf{GLSL:} \texttt{// Primary opcode bits})
  \item \texttt{instruction\_class}: 4-bit major instruction class selector (\texttt{bits 48:51}) partitioning the 64-bit Valhall opcode space (e.g., \texttt{FMA}, \texttt{CVT}, \texttt{SFU}, \texttt{LS}, \texttt{TEX}, \texttt{BRANCH}). \emph{Example:} {\scriptsize\ttfamily \mnem{FADD.f32} \dst{r0}, \src{r1}, \src{r2}} (\textbf{SPIR-V:} \texttt{N/A (encoding class)}; \textbf{GLSL:} \texttt{// Hardware execution pipe / decode class})
  \item \texttt{fau\_page\_index}: Selects the FAU (Fast Access Uniform) page (\texttt{0..3}) when a source operand references hardware small constants, thread IDs, or lane constants. \emph{Example:} {\scriptsize\ttfamily \mnem{FADD.f32} \dst{r0}, \src{r1}, \src{k1f}} (\textbf{SPIR-V:} \texttt{OpConstant / BuiltIn}; \textbf{GLSL:} \texttt{float y = x + 1.0;})
  \item \hyperref[enum:flow_control_m]{\texttt{flow}}: 4-bit embedded flow-control \& scoreboard wait modifier (\texttt{bits 59:62}: \texttt{none}, \texttt{wait0..wait0126}, \texttt{wait}, \texttt{reconverge}, \texttt{discard}, \texttt{end}). \emph{Example:} {\scriptsize\ttfamily \mnem{FADD.f32}\imod{.wait0.reconverge} \dst{r0}, \src{r1}, \src{r2}} (\textbf{SPIR-V:} \texttt{OpControlBarrier / OpReturn}; \textbf{GLSL:} \texttt{// Scoreboard wait + warp reconvergence})
\end{tightitemize}
\clearpage

\subsection*{V2U8\_TO\_V2F16}
\addcontentsline{toc}{subsection}{V2U8\_TO\_V2F16}
Extracts two unsigned 8-bit bytes selected by \texttt{swz0} (\texttt{bXY}, default \texttt{b02}) from \texttt{src0} and converts them exactly to packed FP16 lanes \texttt{dest.h0 = f16(u8(src0.b[X]))} and \texttt{dest.h1 = f16(u8(src0.b[Y]))} (arch \textless{}= v10).
\par\textbf{Notes:}
\begin{tightitemize}
  \item Arch v9-v10 only.
  \item \textbf{Source references:} \texttt{kraid/ops.rs:2542-2570}, \texttt{encode\_v9.rs:1941-1944}.
\end{tightitemize}
\codebox{\mnem{V2U8\_TO\_V2F16}\imod{\{.flow\}} \dst{dest}\imod{\{.dest\_width\}}, \src{src0}\imod{\{.swz0\}}}
\textbf{Operands:}
\begin{tightitemize}\footnotesize
  \item \texttt{\dst{dest}} (\texttt{dest\_register\_index}, \texttt{dest\_width}): 32-bit destination GPR receiving the converted value.
  \item \texttt{\src{src0}}: Packed \texttt{v4u8} input (two bytes selected by \texttt{swz0}, default \texttt{b02}) converted to \texttt{v2f16}.
\end{tightitemize}
\par\textbf{Closest SPIR-V Op:} \texttt{OpConvertUToF (v2u8 -\textgreater{} v2f16, v9--v10)}
\par\textbf{Pseudo-GLSL (Operational Semantics):}
\codebox{dest = f16vec2(float16\_t(u8(src0, swz0.x)), float16\_t(u8(src0, swz0.y)));}
\textbf{Execution Context:}
\begin{tightitemize}\footnotesize
  \item \textbf{Compute stage preloads (\hyperref[sec:csf-spd]{\texttt{SPD\_k}}, \texttt{r16..r23}):} Hardware pre-fills \texttt{r55..r57} with \texttt{WorkGroupID.xyz + JOB\_OFFSET} (from \hyperref[sec:csf-regs]{\texttt{COMPUTESR::JOB\_OFFSET\_X/Y/Z}} \texttt{r34..r36}) and \texttt{r60..r62} with \texttt{LocalInvocationID.xyz} or \texttt{GlobalInvocationID.xyz} (from \texttt{WG\_SIZE} \texttt{r33} \& \texttt{JOB\_SIZE} \texttt{r37..r39}, selected by \texttt{ShaderProgram::preload}) so compute shaders can derive thread/element indices without memory loads.
  \item \textbf{Vertex stage preloads (\hyperref[sec:csf-spd]{\texttt{SPD\_0}/\texttt{SPD\_1}}, \texttt{r16..r19}):} The hardware Index Fetcher reads \hyperref[sec:csf-regs]{\texttt{IDVSSR::INDEX\_BUFFER}} (\texttt{r54:r55}) and pre-fills \texttt{r59, r60} with \texttt{VertexIndex} (\texttt{+ VERTEX\_OFFSET + GLOBAL\_ATTR\_OFF}) and \texttt{r61} with \texttt{InstanceIndex} (\texttt{+ INSTANCE\_OFFSET}) so Position and Varying shaders can pull vertex attributes directly.
  \item \textbf{Inline FAU \& Special Operands (\hyperref[sec:csf-tsd-fau]{\texttt{FAU\_k}}, \texttt{r8..r15}):} Any 8-bit source with \texttt{mode=2} reads zero-latency uniform \texttt{u0..u31} from \texttt{FAU\_k} (\texttt{slots64 = bits[63:56]}); \texttt{mode=3} reads \texttt{SmallConstant} (Table~1.2) or \texttt{FauSpecial} (\texttt{pc}, \texttt{lane\_id}, \texttt{lane\_width}, \texttt{tls\_ptr}, \texttt{wls\_ptr}, \texttt{blend\_descriptor\_0..7}).
\end{tightitemize}
\par\textbf{Worked Examples:}
\par\textbf{Example 1 (Default `b02` unsigned byte pair (v9--v10)):}
\codebox{\mnem{V2U8\_TO\_V2F16} \dst{r0}, \src{r1}\imod{.b02}}
With \texttt{r1 = 0x00020001} (\texttt{b2 = 2, b0 = 1}): yields \texttt{r0 = 0x40003C00} (\texttt{[2.0h, 1.0h]}).
\par\textbf{Example 2 (High-bit byte (`0xFF = 255.0h`)):}
\codebox{\mnem{V2U8\_TO\_V2F16} \dst{r0}, \src{r1}\imod{.b13}}
With \texttt{r1 = 0xFF000000} (\texttt{b3 = 255, b1 = 0}): yields \texttt{r0 = 0x5BF80000} (\texttt{[255.0h, 0.0h]}).

\begin{center}\footnotesize
\begin{tabularx}{\linewidth}{@{}lllX@{}}
\toprule
Variant & Arch & Class & Opcode Fields \\
\midrule
\texttt{(none)} & v9--v14 & \texttt{R1} & \texttt{opcode2=V2U8\_TO\_V2F16, opcode1=CONVERT} \\
\bottomrule
\end{tabularx}
\end{center}
\begin{center}\footnotesize
\begin{tabularx}{0.88\linewidth}{@{}l l X@{}}
\toprule
Field Name & Bits & Type / Fixed Value \\
\midrule
\texttt{src0} & \texttt{0:7} & \texttt{SourceEncoding} \\
\textit{Reserved} & \texttt{8:15} & \texttt{zero} \\
\texttt{opcode2} & \texttt{16:20} & \texttt{exact=V2U8\_TO\_V2F16} \\
\textit{Reserved} & \texttt{21:35} & \texttt{zero} \\
\hyperref[enum:swiz8_2_full_m]{\texttt{swz0}} & \texttt{36:39} & \texttt{swiz8\_2\_full\_m (def b02)} \\
\texttt{dest\_register\_index} & \texttt{40:45} & \texttt{uint} \\
\hyperref[enum:dest_width_m]{\texttt{dest\_width}} & \texttt{46:47} & \texttt{dest\_width\_m} \\
\texttt{opcode1} & \texttt{48:51} & \texttt{exact=CONVERT} \\
\texttt{instruction\_class} & \texttt{52:56} & \texttt{exact=R1} \\
\texttt{fau\_page\_index} & \texttt{57:58} & \texttt{uint} \\
\hyperref[enum:flow_control_m]{\texttt{flow}} & \texttt{59:62} & \texttt{flow\_control\_m} \\
\textit{Reserved} & \texttt{63} & \texttt{zero} \\
\bottomrule
\end{tabularx}
\end{center}
\textbf{Fields Explained:}
\begin{tightitemize}\footnotesize
  \item \texttt{src0}: Primary source operand (\texttt{r0..r63}, \texttt{u0..u63} uniform, \texttt{k*} FAU constant, or \texttt{t0..t2} temporary; \texttt{\textasciicircum{}} marks discard/last-use). \emph{Example:} {\scriptsize\ttfamily \mnem{FADD.f32} \dst{r0}, \src{r1}, \src{r2}} (\textbf{SPIR-V:} \texttt{OpFAdd (Operand 0)}; \textbf{GLSL:} \texttt{float c = a + b;})
  \item \texttt{opcode2}: Secondary sub-opcode further specializing instructions sharing \texttt{instruction\_class} and \texttt{opcode1} (e.g., conversions, SFU tables, moves). \emph{Example:} {\scriptsize\ttfamily \mnem{F32\_TO\_S32}\imod{.round\_zero} \dst{r0}, \src{r1}} (\textbf{SPIR-V:} \texttt{OpConvertFToS}; \textbf{GLSL:} \texttt{int y = int(x);})
  \item \hyperref[enum:swiz8_2_full_m]{\texttt{swz0}}: Stage A lane swizzle on \texttt{src0} (\texttt{h00}/\texttt{h01}/\texttt{h10}/\texttt{h11} for 16-bit vectors, or \texttt{bXY} byte-to-halfword widening). \emph{Example:} {\scriptsize\ttfamily \mnem{FADD.v2f16} \dst{r0}, \src{r1}\imod{.h10}, \src{r2}} (\textbf{SPIR-V:} \texttt{OpVectorShuffle}; \textbf{GLSL:} \texttt{f16vec2 c = a.yx + b;})
  \item \texttt{dest\_register\_index}: 6-bit destination GPR index (\texttt{r0..r63}) receiving the Stage C writeback result (even-aligned pair \texttt{[rN:rN+1]} for 64-bit ops). \emph{Example:} {\scriptsize\ttfamily \mnem{FMA.f32} \dst{r0}, \src{r1}, \src{r2}, \src{r3}} (\textbf{SPIR-V:} \texttt{SSA Result ID}; \textbf{GLSL:} \texttt{float r0 = fma(r1, r2, r3);})
  \item \hyperref[enum:dest_width_m]{\texttt{dest\_width}}: Stage C destination lane write-mask (\texttt{none} = full 32-bit, \texttt{h0} = low 16-bit, \texttt{h1} = high 16-bit without bit-shifting). \emph{Example:} {\scriptsize\ttfamily \mnem{FADD.v2f16} \dst{r0}\imod{.h1}, \src{r1}, \src{r2}} (\textbf{SPIR-V:} \texttt{OpVectorInsertDynamic / WriteMask}; \textbf{GLSL:} \texttt{dst.y = a.y + b.y;})
  \item \texttt{opcode1}: Primary sub-opcode within the instruction class uniquely identifying the operation family. \emph{Example:} {\scriptsize\ttfamily \mnem{IADD.s32} \dst{r0}, \src{r1}, \src{r2}} (\textbf{SPIR-V:} \texttt{OpIAdd (operation selector)}; \textbf{GLSL:} \texttt{// Primary opcode bits})
  \item \texttt{instruction\_class}: 4-bit major instruction class selector (\texttt{bits 48:51}) partitioning the 64-bit Valhall opcode space (e.g., \texttt{FMA}, \texttt{CVT}, \texttt{SFU}, \texttt{LS}, \texttt{TEX}, \texttt{BRANCH}). \emph{Example:} {\scriptsize\ttfamily \mnem{FADD.f32} \dst{r0}, \src{r1}, \src{r2}} (\textbf{SPIR-V:} \texttt{N/A (encoding class)}; \textbf{GLSL:} \texttt{// Hardware execution pipe / decode class})
  \item \texttt{fau\_page\_index}: Selects the FAU (Fast Access Uniform) page (\texttt{0..3}) when a source operand references hardware small constants, thread IDs, or lane constants. \emph{Example:} {\scriptsize\ttfamily \mnem{FADD.f32} \dst{r0}, \src{r1}, \src{k1f}} (\textbf{SPIR-V:} \texttt{OpConstant / BuiltIn}; \textbf{GLSL:} \texttt{float y = x + 1.0;})
  \item \hyperref[enum:flow_control_m]{\texttt{flow}}: 4-bit embedded flow-control \& scoreboard wait modifier (\texttt{bits 59:62}: \texttt{none}, \texttt{wait0..wait0126}, \texttt{wait}, \texttt{reconverge}, \texttt{discard}, \texttt{end}). \emph{Example:} {\scriptsize\ttfamily \mnem{FADD.f32}\imod{.wait0.reconverge} \dst{r0}, \src{r1}, \src{r2}} (\textbf{SPIR-V:} \texttt{OpControlBarrier / OpReturn}; \textbf{GLSL:} \texttt{// Scoreboard wait + warp reconvergence})
\end{tightitemize}
\clearpage

\subsection*{V2U8\_TO\_V2U16}
\addcontentsline{toc}{subsection}{V2U8\_TO\_V2U16}
Selects two unsigned 8-bit bytes from \texttt{src0} via \texttt{swz0} (default \texttt{b02}: bytes 0 and 2), zero-extends each to 16 bits, and writes the packed \texttt{v2u16} pair to \texttt{dest} (respecting \texttt{dest\_width}).
\par\textbf{Notes:}
\begin{tightitemize}
  \item Exact per-lane zero extension.
  \item \textbf{Source references:} \texttt{isa-v9-v14.xml}, \texttt{bifrost/valhall/ISA.xml}.
\end{tightitemize}
\codebox{\mnem{V2U8\_TO\_V2U16}\imod{\{.flow\}} \dst{dest}\imod{\{.dest\_width\}}, \src{src0}\imod{\{.swz0\}}}
\textbf{Operands:}
\begin{tightitemize}\footnotesize
  \item \texttt{\dst{dest}} (\texttt{dest\_register\_index}, \texttt{dest\_width}): 32-bit destination GPR receiving the converted value.
  \item \texttt{\src{src0}}: Packed \texttt{v4u8} input (two bytes selected by \texttt{swz0}, default \texttt{b02}) zero-extended to \texttt{v2u16}.
\end{tightitemize}
\par\textbf{Closest SPIR-V Op:} \texttt{OpUConvert (v2u8 -\textgreater{} v2u16)}
\par\textbf{Pseudo-GLSL (Operational Semantics):}
\codebox{dest = maskWrite(dest\_width, u16vec2(u8(src0, swz0.x), u8(src0, swz0.y)));}
\textbf{Execution Context:}
\begin{tightitemize}\footnotesize
  \item \textbf{Compute stage preloads (\hyperref[sec:csf-spd]{\texttt{SPD\_k}}, \texttt{r16..r23}):} Hardware pre-fills \texttt{r55..r57} with \texttt{WorkGroupID.xyz + JOB\_OFFSET} (from \hyperref[sec:csf-regs]{\texttt{COMPUTESR::JOB\_OFFSET\_X/Y/Z}} \texttt{r34..r36}) and \texttt{r60..r62} with \texttt{LocalInvocationID.xyz} or \texttt{GlobalInvocationID.xyz} (from \texttt{WG\_SIZE} \texttt{r33} \& \texttt{JOB\_SIZE} \texttt{r37..r39}, selected by \texttt{ShaderProgram::preload}) so compute shaders can derive thread/element indices without memory loads.
  \item \textbf{Vertex stage preloads (\hyperref[sec:csf-spd]{\texttt{SPD\_0}/\texttt{SPD\_1}}, \texttt{r16..r19}):} The hardware Index Fetcher reads \hyperref[sec:csf-regs]{\texttt{IDVSSR::INDEX\_BUFFER}} (\texttt{r54:r55}) and pre-fills \texttt{r59, r60} with \texttt{VertexIndex} (\texttt{+ VERTEX\_OFFSET + GLOBAL\_ATTR\_OFF}) and \texttt{r61} with \texttt{InstanceIndex} (\texttt{+ INSTANCE\_OFFSET}) so Position and Varying shaders can pull vertex attributes directly.
  \item \textbf{Inline FAU \& Special Operands (\hyperref[sec:csf-tsd-fau]{\texttt{FAU\_k}}, \texttt{r8..r15}):} Any 8-bit source with \texttt{mode=2} reads zero-latency uniform \texttt{u0..u31} from \texttt{FAU\_k} (\texttt{slots64 = bits[63:56]}); \texttt{mode=3} reads \texttt{SmallConstant} (Table~1.2) or \texttt{FauSpecial} (\texttt{pc}, \texttt{lane\_id}, \texttt{lane\_width}, \texttt{tls\_ptr}, \texttt{wls\_ptr}, \texttt{blend\_descriptor\_0..7}).
\end{tightitemize}
\par\textbf{Worked Examples:}
\par\textbf{Example 1 (ASTC swizzle + lane-mask idiom):}
\codebox{\mnem{V2U8\_TO\_V2U16} \dst{r0}\imod{.h1}, \src{r1}\imod{.b20}}
With \texttt{r1 = 0x000500AB}, \texttt{r0\_old = 0x00001234}: \texttt{b20} puts \texttt{b2=5} in lane 0 and \texttt{b0=0xAB} in lane 1; \texttt{.h1} masks lane 1 -\textgreater{} \texttt{r0 = 0x00AB1234}.
\par\textbf{Example 2 (Full 2-lane unpack (`b01`)):}
\codebox{\mnem{V2U8\_TO\_V2U16} \dst{r0}, \src{r1}\imod{.b01}}
With \texttt{r1 = 0x0000FEDC}: zero-extends \texttt{b0=0xDC} and \texttt{b1=0xFE} -\textgreater{} \texttt{r0 = 0x00FE00DC}.

\begin{center}\footnotesize
\begin{tabularx}{\linewidth}{@{}lllX@{}}
\toprule
Variant & Arch & Class & Opcode Fields \\
\midrule
\texttt{(none)} & v9--v14 & \texttt{R1} & \texttt{opcode2=V2U8\_TO\_V2U16, opcode1=CONVERT} \\
\texttt{(none)} & v9--v14 & \texttt{R1} & \texttt{opcode2=V2I8\_TO\_V2I16, opcode1=CONVERT} \\
\bottomrule
\end{tabularx}
\end{center}
\begin{center}\footnotesize
\begin{tabularx}{0.88\linewidth}{@{}l l X@{}}
\toprule
Field Name & Bits & Type / Fixed Value \\
\midrule
\texttt{src0} & \texttt{0:7} & \texttt{SourceEncoding} \\
\textit{Reserved} & \texttt{8:15} & \texttt{zero} \\
\texttt{opcode2} & \texttt{16:20} & \texttt{exact=V2U8\_TO\_V2U16} \\
\textit{Reserved} & \texttt{21:35} & \texttt{zero} \\
\hyperref[enum:swiz8_2_full_m]{\texttt{swz0}} & \texttt{36:39} & \texttt{swiz8\_2\_full\_m (def b02)} \\
\texttt{dest\_register\_index} & \texttt{40:45} & \texttt{uint} \\
\hyperref[enum:dest_width_m]{\texttt{dest\_width}} & \texttt{46:47} & \texttt{dest\_width\_m} \\
\texttt{opcode1} & \texttt{48:51} & \texttt{exact=CONVERT} \\
\texttt{instruction\_class} & \texttt{52:56} & \texttt{exact=R1} \\
\texttt{fau\_page\_index} & \texttt{57:58} & \texttt{uint} \\
\hyperref[enum:flow_control_m]{\texttt{flow}} & \texttt{59:62} & \texttt{flow\_control\_m} \\
\textit{Reserved} & \texttt{63} & \texttt{zero} \\
\bottomrule
\end{tabularx}
\end{center}
\textbf{Fields Explained:}
\begin{tightitemize}\footnotesize
  \item \texttt{src0}: Primary source operand (\texttt{r0..r63}, \texttt{u0..u63} uniform, \texttt{k*} FAU constant, or \texttt{t0..t2} temporary; \texttt{\textasciicircum{}} marks discard/last-use). \emph{Example:} {\scriptsize\ttfamily \mnem{FADD.f32} \dst{r0}, \src{r1}, \src{r2}} (\textbf{SPIR-V:} \texttt{OpFAdd (Operand 0)}; \textbf{GLSL:} \texttt{float c = a + b;})
  \item \texttt{opcode2}: Secondary sub-opcode further specializing instructions sharing \texttt{instruction\_class} and \texttt{opcode1} (e.g., conversions, SFU tables, moves). \emph{Example:} {\scriptsize\ttfamily \mnem{F32\_TO\_S32}\imod{.round\_zero} \dst{r0}, \src{r1}} (\textbf{SPIR-V:} \texttt{OpConvertFToS}; \textbf{GLSL:} \texttt{int y = int(x);})
  \item \hyperref[enum:swiz8_2_full_m]{\texttt{swz0}}: Stage A lane swizzle on \texttt{src0} (\texttt{h00}/\texttt{h01}/\texttt{h10}/\texttt{h11} for 16-bit vectors, or \texttt{bXY} byte-to-halfword widening). \emph{Example:} {\scriptsize\ttfamily \mnem{FADD.v2f16} \dst{r0}, \src{r1}\imod{.h10}, \src{r2}} (\textbf{SPIR-V:} \texttt{OpVectorShuffle}; \textbf{GLSL:} \texttt{f16vec2 c = a.yx + b;})
  \item \texttt{dest\_register\_index}: 6-bit destination GPR index (\texttt{r0..r63}) receiving the Stage C writeback result (even-aligned pair \texttt{[rN:rN+1]} for 64-bit ops). \emph{Example:} {\scriptsize\ttfamily \mnem{FMA.f32} \dst{r0}, \src{r1}, \src{r2}, \src{r3}} (\textbf{SPIR-V:} \texttt{SSA Result ID}; \textbf{GLSL:} \texttt{float r0 = fma(r1, r2, r3);})
  \item \hyperref[enum:dest_width_m]{\texttt{dest\_width}}: Stage C destination lane write-mask (\texttt{none} = full 32-bit, \texttt{h0} = low 16-bit, \texttt{h1} = high 16-bit without bit-shifting). \emph{Example:} {\scriptsize\ttfamily \mnem{FADD.v2f16} \dst{r0}\imod{.h1}, \src{r1}, \src{r2}} (\textbf{SPIR-V:} \texttt{OpVectorInsertDynamic / WriteMask}; \textbf{GLSL:} \texttt{dst.y = a.y + b.y;})
  \item \texttt{opcode1}: Primary sub-opcode within the instruction class uniquely identifying the operation family. \emph{Example:} {\scriptsize\ttfamily \mnem{IADD.s32} \dst{r0}, \src{r1}, \src{r2}} (\textbf{SPIR-V:} \texttt{OpIAdd (operation selector)}; \textbf{GLSL:} \texttt{// Primary opcode bits})
  \item \texttt{instruction\_class}: 4-bit major instruction class selector (\texttt{bits 48:51}) partitioning the 64-bit Valhall opcode space (e.g., \texttt{FMA}, \texttt{CVT}, \texttt{SFU}, \texttt{LS}, \texttt{TEX}, \texttt{BRANCH}). \emph{Example:} {\scriptsize\ttfamily \mnem{FADD.f32} \dst{r0}, \src{r1}, \src{r2}} (\textbf{SPIR-V:} \texttt{N/A (encoding class)}; \textbf{GLSL:} \texttt{// Hardware execution pipe / decode class})
  \item \texttt{fau\_page\_index}: Selects the FAU (Fast Access Uniform) page (\texttt{0..3}) when a source operand references hardware small constants, thread IDs, or lane constants. \emph{Example:} {\scriptsize\ttfamily \mnem{FADD.f32} \dst{r0}, \src{r1}, \src{k1f}} (\textbf{SPIR-V:} \texttt{OpConstant / BuiltIn}; \textbf{GLSL:} \texttt{float y = x + 1.0;})
  \item \hyperref[enum:flow_control_m]{\texttt{flow}}: 4-bit embedded flow-control \& scoreboard wait modifier (\texttt{bits 59:62}: \texttt{none}, \texttt{wait0..wait0126}, \texttt{wait}, \texttt{reconverge}, \texttt{discard}, \texttt{end}). \emph{Example:} {\scriptsize\ttfamily \mnem{FADD.f32}\imod{.wait0.reconverge} \dst{r0}, \src{r1}, \src{r2}} (\textbf{SPIR-V:} \texttt{OpControlBarrier / OpReturn}; \textbf{GLSL:} \texttt{// Scoreboard wait + warp reconvergence})
\end{tightitemize}
\clearpage

\section{Move, Pack \& Swizzle (5 mnemonics)}
\subsection*{ADR}
\addcontentsline{toc}{subsection}{ADR}
Computes a 32-bit PC-relative shader code address by adding the signed inline offset \texttt{\#imm1w} to the address of the following instruction (\texttt{PC + 8}) and writes \texttt{(PC + 8 + \#imm1w) \& 0xFFFFFFFF} to \texttt{dest}.
\par\textbf{Notes:}
\begin{tightitemize}
  \item Destination is \textbf{32-bit} (\texttt{\#[dst\_type(I32)]} in \texttt{kraid/ops.rs:127}).
  \item \texttt{imm1w} is byte offset from \texttt{PC + 8} (\texttt{encode\_v9.rs:162-164}).
  \item \textbf{Source references:} \texttt{kraid/ops.rs:126-141}, \texttt{encode\_v9.rs:162-164, 582-593}.
\end{tightitemize}
\codebox{\mnem{ADR}\imod{\{.flow\}} \dst{dest}\imod{\{.dest\_width\}}, \imm{\#imm1w}}
\textbf{Operands:}
\begin{tightitemize}\footnotesize
  \item \texttt{\dst{dest}} (\texttt{dest\_register\_index}, \texttt{dest\_width}): 32-bit destination GPR receiving \texttt{(PC + 8 + \#imm1w) \& 0xFFFFFFFF}.
  \item \texttt{\imm{\#imm1w}} (\texttt{imm1w}): Signed 32-bit PC-relative byte offset from the next instruction (\texttt{PC + 8}).
\end{tightitemize}
\par\textbf{Closest SPIR-V Op:} \texttt{None (PC-relative address for blend/call)}
\par\textbf{Pseudo-GLSL (Operational Semantics):}
\codebox{dest = uint((PC + 8ul + int64\_t(\#imm1w)) \& 0xFFFFFFFFul);}
\textbf{Execution Context:}
\begin{tightitemize}\footnotesize
  \item \textbf{Shader BO base (\hyperref[sec:csf-spd]{\texttt{SPD\_k::binary}}, \texttt{r16..r23}):} Computes PC-relative addresses inside the shader binary object (e.g., embedded constant LUTs or blend-shader return labels).
  \item \textbf{Inline FAU \& Special Operands (\hyperref[sec:csf-tsd-fau]{\texttt{FAU\_k}}, \texttt{r8..r15}):} Any 8-bit source with \texttt{mode=2} reads zero-latency uniform \texttt{u0..u31} from \texttt{FAU\_k} (\texttt{slots64 = bits[63:56]}); \texttt{mode=3} reads \texttt{SmallConstant} (Table~1.2) or \texttt{FauSpecial} (\texttt{pc}, \texttt{lane\_id}, \texttt{lane\_width}, \texttt{tls\_ptr}, \texttt{wls\_ptr}, \texttt{blend\_descriptor\_0..7}).
\end{tightitemize}
\par\textbf{Worked Examples:}
\par\textbf{Example 1 (Forward return-address calculation):}
\codebox{\mnem{ADR} \dst{r54}, \imm{\#0x10}}
At \texttt{PC = 0x1000}: computes 32-bit \texttt{r54 = (0x1000 + 8 + 0x10) \& 0xFFFFFFFF = 0x1018}.
\par\textbf{Example 2 (Negative backward offset):}
\codebox{\mnem{ADR} \dst{r54}, \imm{\#-0x18}}
At \texttt{PC = 0x1020}: computes \texttt{0x1020 + 8 - 0x18 = 0x1010} (PC-relative label before current instruction).

\begin{center}\footnotesize
\begin{tabularx}{\linewidth}{@{}lllX@{}}
\toprule
Variant & Arch & Class & Opcode Fields \\
\midrule
\texttt{(none)} & v9--v14 & \texttt{R1I} & \texttt{opcode1=ADD\_IMM\_I32} \\
\bottomrule
\end{tabularx}
\end{center}
\begin{center}\footnotesize
\begin{tabularx}{0.88\linewidth}{@{}l l X@{}}
\toprule
Field Name & Bits & Type / Fixed Value \\
\midrule
\texttt{src0} & \texttt{0:7} & \texttt{exact=254} \\
\texttt{imm1w} & \texttt{8:39} & \texttt{pc\_rel\_label\_signed} \\
\texttt{dest\_register\_index} & \texttt{40:45} & \texttt{uint} \\
\hyperref[enum:dest_width_m]{\texttt{dest\_width}} & \texttt{46:47} & \texttt{exact=none} \\
\texttt{opcode1} & \texttt{48:51} & \texttt{exact=ADD\_IMM\_I32} \\
\texttt{instruction\_class} & \texttt{52:56} & \texttt{exact=R1I} \\
\texttt{fau\_page\_index} & \texttt{57:58} & \texttt{exact=3} \\
\hyperref[enum:flow_control_m]{\texttt{flow}} & \texttt{59:62} & \texttt{flow\_control\_m} \\
\textit{Reserved} & \texttt{63} & \texttt{zero} \\
\bottomrule
\end{tabularx}
\end{center}
\textbf{Fields Explained:}
\begin{tightitemize}\footnotesize
  \item \texttt{src0}: Primary source operand (\texttt{r0..r63}, \texttt{u0..u63} uniform, \texttt{k*} FAU constant, or \texttt{t0..t2} temporary; \texttt{\textasciicircum{}} marks discard/last-use). \emph{Example:} {\scriptsize\ttfamily \mnem{FADD.f32} \dst{r0}, \src{r1}, \src{r2}} (\textbf{SPIR-V:} \texttt{OpFAdd (Operand 0)}; \textbf{GLSL:} \texttt{float c = a + b;})
  \item \texttt{imm1w}: 32-bit inline immediate operand encoded directly in the instruction word (bits 8:39). \emph{Example:} {\scriptsize\ttfamily \mnem{IADD\_IMM.i32} \dst{r0}, \src{r1}, \imm{\#0x10}} (\textbf{SPIR-V:} \texttt{OpIAdd (constant)}; \textbf{GLSL:} \texttt{int y = x + 16;})
  \item \texttt{dest\_register\_index}: 6-bit destination GPR index (\texttt{r0..r63}) receiving the Stage C writeback result (even-aligned pair \texttt{[rN:rN+1]} for 64-bit ops). \emph{Example:} {\scriptsize\ttfamily \mnem{FMA.f32} \dst{r0}, \src{r1}, \src{r2}, \src{r3}} (\textbf{SPIR-V:} \texttt{SSA Result ID}; \textbf{GLSL:} \texttt{float r0 = fma(r1, r2, r3);})
  \item \hyperref[enum:dest_width_m]{\texttt{dest\_width}}: Stage C destination lane write-mask (\texttt{none} = full 32-bit, \texttt{h0} = low 16-bit, \texttt{h1} = high 16-bit without bit-shifting). \emph{Example:} {\scriptsize\ttfamily \mnem{FADD.v2f16} \dst{r0}\imod{.h1}, \src{r1}, \src{r2}} (\textbf{SPIR-V:} \texttt{OpVectorInsertDynamic / WriteMask}; \textbf{GLSL:} \texttt{dst.y = a.y + b.y;})
  \item \texttt{opcode1}: Primary sub-opcode within the instruction class uniquely identifying the operation family. \emph{Example:} {\scriptsize\ttfamily \mnem{IADD.s32} \dst{r0}, \src{r1}, \src{r2}} (\textbf{SPIR-V:} \texttt{OpIAdd (operation selector)}; \textbf{GLSL:} \texttt{// Primary opcode bits})
  \item \texttt{instruction\_class}: 4-bit major instruction class selector (\texttt{bits 48:51}) partitioning the 64-bit Valhall opcode space (e.g., \texttt{FMA}, \texttt{CVT}, \texttt{SFU}, \texttt{LS}, \texttt{TEX}, \texttt{BRANCH}). \emph{Example:} {\scriptsize\ttfamily \mnem{FADD.f32} \dst{r0}, \src{r1}, \src{r2}} (\textbf{SPIR-V:} \texttt{N/A (encoding class)}; \textbf{GLSL:} \texttt{// Hardware execution pipe / decode class})
  \item \texttt{fau\_page\_index}: Selects the FAU (Fast Access Uniform) page (\texttt{0..3}) when a source operand references hardware small constants, thread IDs, or lane constants. \emph{Example:} {\scriptsize\ttfamily \mnem{FADD.f32} \dst{r0}, \src{r1}, \src{k1f}} (\textbf{SPIR-V:} \texttt{OpConstant / BuiltIn}; \textbf{GLSL:} \texttt{float y = x + 1.0;})
  \item \hyperref[enum:flow_control_m]{\texttt{flow}}: 4-bit embedded flow-control \& scoreboard wait modifier (\texttt{bits 59:62}: \texttt{none}, \texttt{wait0..wait0126}, \texttt{wait}, \texttt{reconverge}, \texttt{discard}, \texttt{end}). \emph{Example:} {\scriptsize\ttfamily \mnem{FADD.f32}\imod{.wait0.reconverge} \dst{r0}, \src{r1}, \src{r2}} (\textbf{SPIR-V:} \texttt{OpControlBarrier / OpReturn}; \textbf{GLSL:} \texttt{// Scoreboard wait + warp reconvergence})
\end{tightitemize}
\clearpage

\subsection*{MKVEC}
\addcontentsline{toc}{subsection}{MKVEC}
\texttt{MKVEC.v2i16}: \texttt{Concat(src1.lane1, src0.lane0)}. \texttt{MKVEC.v2i8}: \texttt{Concat(src2[15:0], src1.lane1, src0.lane0)} (bits \texttt{[35:24]} reserved zero; virtual \texttt{lanes2=b01} in XML and fixed \texttt{b01} in \texttt{libGLES\_mali.so} \texttt{ModInfo}).
\par\textbf{Note:} All 3 workstreams reconciled (\texttt{isa-v9-v14.xml}, \texttt{libGLES\_mali.so} \texttt{ModInfo}, 313 ASTC/ETC2/BC6 \texttt{malioc} emissions).
\codebox{\mnem{MKVEC}\imod{\{.flow\}} \dst{dest}\imod{\{.dest\_width\}}, \src{src0}\imod{\{.lane0\}}, \src{src1}\imod{\{.lane1\}}\\\mnem{MKVEC}\imod{\{.flow\}} \dst{dest}\imod{\{.dest\_width\}}, \src{src0}\imod{\{.lane0\}}, \src{src1}\imod{\{.lane1\}}, \src{src2}\imod{\{.lanes2\}}}
\textbf{Operands:}
\begin{tightitemize}\footnotesize
  \item \texttt{\dst{dest}} (\texttt{dest\_register\_index}, \texttt{dest\_width}): 32-bit destination GPR receiving the packed halfwords or bytes.
  \item \texttt{\src{src0}}: Source supplying lane 0 (selected via \texttt{lane0}).
  \item \texttt{\src{src1}}: Source supplying lane 1 (selected via \texttt{lane1}).
  \item \texttt{\src{src2}}: Source supplying upper 16 bits \texttt{[31:16]} in \texttt{MKVEC.v2i8} mode (unused in \texttt{v2i16}).
\end{tightitemize}
\par\textbf{Closest SPIR-V Op:} \texttt{OpCompositeConstruct / PackUnorm4x8}
\par\textbf{Pseudo-GLSL (Operational Semantics):}
\codebox{if (is\_v2i8) dest = u8(src0, lane0) \textbar{} (u8(src1, lane1) \textless{}\textless{} 8) \textbar{} ((src2 \& 0xFFFFu) \textless{}\textless{} 16);\\else         dest = u16(src0, lane0) \textbar{} (u16(src1, lane1) \textless{}\textless{} 16); // v2i16}
\textbf{Execution Context:}
\begin{tightitemize}\footnotesize
  \item \textbf{Compute stage preloads (\hyperref[sec:csf-spd]{\texttt{SPD\_k}}, \texttt{r16..r23}):} Hardware pre-fills \texttt{r55..r57} with \texttt{WorkGroupID.xyz + JOB\_OFFSET} (from \hyperref[sec:csf-regs]{\texttt{COMPUTESR::JOB\_OFFSET\_X/Y/Z}} \texttt{r34..r36}) and \texttt{r60..r62} with \texttt{LocalInvocationID.xyz} or \texttt{GlobalInvocationID.xyz} (from \texttt{WG\_SIZE} \texttt{r33} \& \texttt{JOB\_SIZE} \texttt{r37..r39}, selected by \texttt{ShaderProgram::preload}) so compute shaders can derive thread/element indices without memory loads.
  \item \textbf{Vertex stage preloads (\hyperref[sec:csf-spd]{\texttt{SPD\_0}/\texttt{SPD\_1}}, \texttt{r16..r19}):} The hardware Index Fetcher reads \hyperref[sec:csf-regs]{\texttt{IDVSSR::INDEX\_BUFFER}} (\texttt{r54:r55}) and pre-fills \texttt{r59, r60} with \texttt{VertexIndex} (\texttt{+ VERTEX\_OFFSET + GLOBAL\_ATTR\_OFF}) and \texttt{r61} with \texttt{InstanceIndex} (\texttt{+ INSTANCE\_OFFSET}) so Position and Varying shaders can pull vertex attributes directly.
  \item \textbf{Inline FAU \& Special Operands (\hyperref[sec:csf-tsd-fau]{\texttt{FAU\_k}}, \texttt{r8..r15}):} Any 8-bit source with \texttt{mode=2} reads zero-latency uniform \texttt{u0..u31} from \texttt{FAU\_k} (\texttt{slots64 = bits[63:56]}); \texttt{mode=3} reads \texttt{SmallConstant} (Table~1.2) or \texttt{FauSpecial} (\texttt{pc}, \texttt{lane\_id}, \texttt{lane\_width}, \texttt{tls\_ptr}, \texttt{wls\_ptr}, \texttt{blend\_descriptor\_0..7}).
\end{tightitemize}
\par\textbf{Worked Examples:}
\par\textbf{Example 1 (4-byte pack with fixed virtual `b01` on `src2`):}
\codebox{\mnem{MKVEC.v2i8} \dst{r0}, \src{r1}\imod{.b0}, \src{r2}\imod{.b0}, \src{r3}\imod{.b01}}
With \texttt{r1 = 0x11, r2 = 0x22, r3 = 0xAABB4433}: packs \texttt{Concat(r3[15:0], r2.b0, r1.b0) = 0x44332211}.
\par\textbf{Example 2 (Halfword vector pack (`MKVEC.v2i16`)):}
\codebox{\mnem{MKVEC.v2i16} \dst{r0}, \src{r1}\imod{.h1}, \src{r2}\imod{.h0}}
With \texttt{r1 = 0xAAAA0000}, \texttt{r2 = 0x0000BBBB}: places \texttt{r1.h1} in low half and \texttt{r2.h0} in high half -\textgreater{} \texttt{r0 = 0xBBBBAAAA}.
\par\textbf{Mined SPIR-V $\to$ Valhall ISA Example (Packing two FP16 constants into a v2f16 register for BLEND --- Fragment shader, Mali-G710 (v10)):}
\codebox{\textbf{; SPIR-V (input)}\\OpStore \%frag\_out \%vec4\_1\_1\_1\_1\\[3pt]\textbf{; Valhall ISA (malioc output)}\\MOV\_IMM.i32.discard r62, \#0x3c003c00 ; (1.0h, 1.0h)\\MKVEC.v2i16.wait r63, r62.h0, r62.h1\\BLEND.slot0.end.v4.f16 @r62:r63, blend\_descriptor\_0, 0x68, r60\textasciicircum{}}
\texttt{MKVEC.v2i16} combines two 16-bit halfwords (\texttt{r62.h0} and \texttt{r62.h1}) into a single 32-bit register \texttt{r63} to form the 64-bit \texttt{v4.f16} RGBA output for \texttt{BLEND}.
\begin{center}\footnotesize
\begin{tabularx}{\linewidth}{@{}lllX@{}}
\toprule
Variant & Arch & Class & Opcode Fields \\
\midrule
\texttt{v2u16} & v9--v14 & \texttt{R2} & \texttt{opcode2=MKVEC} \\
\texttt{(none)} & v9--v14 & \texttt{R3} & \texttt{opcode1=MKVEC\_V2I8} \\
\bottomrule
\end{tabularx}
\end{center}
\begin{center}\footnotesize
\begin{tabularx}{0.88\linewidth}{@{}l l X@{}}
\toprule
Field Name & Bits & Type / Fixed Value \\
\midrule
\texttt{src0} & \texttt{0:7} & \texttt{SourceEncoding} \\
\texttt{src1} & \texttt{8:15} & \texttt{SourceEncoding} \\
\texttt{opcode2} & \texttt{16:19} & \texttt{exact=MKVEC} \\
\textit{Reserved} & \texttt{20:25} & \texttt{zero} \\
\hyperref[enum:lane16_m]{\texttt{lane1}} & \texttt{26} & \texttt{lane16\_m (def h0)} \\
\textit{Reserved} & \texttt{27} & \texttt{zero} \\
\hyperref[enum:lane16_m]{\texttt{lane0}} & \texttt{28} & \texttt{lane16\_m (def h0)} \\
\textit{Reserved} & \texttt{29:39} & \texttt{zero} \\
\texttt{dest\_register\_index} & \texttt{40:45} & \texttt{uint} \\
\hyperref[enum:dest_width_m]{\texttt{dest\_width}} & \texttt{46:47} & \texttt{dest\_width\_m} \\
\texttt{datatype} & \texttt{48:51} & \texttt{exact=v2u16} \\
\texttt{instruction\_class} & \texttt{52:56} & \texttt{exact=R2} \\
\texttt{fau\_page\_index} & \texttt{57:58} & \texttt{uint} \\
\hyperref[enum:flow_control_m]{\texttt{flow}} & \texttt{59:62} & \texttt{flow\_control\_m} \\
\textit{Reserved} & \texttt{63} & \texttt{zero} \\
\bottomrule
\end{tabularx}
\end{center}
\textbf{Fields Explained:}
\begin{tightitemize}\footnotesize
  \item \texttt{src0}: Primary source operand (\texttt{r0..r63}, \texttt{u0..u63} uniform, \texttt{k*} FAU constant, or \texttt{t0..t2} temporary; \texttt{\textasciicircum{}} marks discard/last-use). \emph{Example:} {\scriptsize\ttfamily \mnem{FADD.f32} \dst{r0}, \src{r1}, \src{r2}} (\textbf{SPIR-V:} \texttt{OpFAdd (Operand 0)}; \textbf{GLSL:} \texttt{float c = a + b;})
  \item \texttt{src1}: Secondary source operand (register, uniform, or FAU constant; in shifts, supplies the shift count byte). \emph{Example:} {\scriptsize\ttfamily \mnem{LSHIFT.i32} \dst{r0}, \src{r1}, \src{r2}\imod{.b0}} (\textbf{SPIR-V:} \texttt{OpShiftLeftLogical (Shift)}; \textbf{GLSL:} \texttt{uint y = a \textless{}\textless{} s;})
  \item \texttt{opcode2}: Secondary sub-opcode further specializing instructions sharing \texttt{instruction\_class} and \texttt{opcode1} (e.g., conversions, SFU tables, moves). \emph{Example:} {\scriptsize\ttfamily \mnem{F32\_TO\_S32}\imod{.round\_zero} \dst{r0}, \src{r1}} (\textbf{SPIR-V:} \texttt{OpConvertFToS}; \textbf{GLSL:} \texttt{int y = int(x);})
  \item \hyperref[enum:lane16_m]{\texttt{lane1}}: Stage A single-lane extractor on \texttt{src1} (\texttt{b0..b3} or \texttt{h0}/\texttt{h1}) for \texttt{MKVEC} packing and narrow binary ops. \emph{Example:} {\scriptsize\ttfamily \mnem{MKVEC.v2i8} \dst{r0}, \src{r1}\imod{.b0}, \src{r2}\imod{.b1}, \src{r3}\imod{.b01}} (\textbf{SPIR-V:} \texttt{OpCompositeConstruct}; \textbf{GLSL:} \texttt{u8vec4 v = u8vec4(a.x, b.y, c.x, c.y);})
  \item \hyperref[enum:lane16_m]{\texttt{lane0}}: Stage A single-lane extractor on \texttt{src0} (\texttt{b0..b3} for 8-bit or \texttt{h0}/\texttt{h1} for 16-bit narrow conversions and \texttt{MKVEC}). \emph{Example:} {\scriptsize\ttfamily \mnem{F16\_TO\_F32} \dst{r0}, \src{r1}\imod{.h1}} (\textbf{SPIR-V:} \texttt{OpCompositeExtract}; \textbf{GLSL:} \texttt{float y = float(x.y);})
  \item \texttt{dest\_register\_index}: 6-bit destination GPR index (\texttt{r0..r63}) receiving the Stage C writeback result (even-aligned pair \texttt{[rN:rN+1]} for 64-bit ops). \emph{Example:} {\scriptsize\ttfamily \mnem{FMA.f32} \dst{r0}, \src{r1}, \src{r2}, \src{r3}} (\textbf{SPIR-V:} \texttt{SSA Result ID}; \textbf{GLSL:} \texttt{float r0 = fma(r1, r2, r3);})
  \item \hyperref[enum:dest_width_m]{\texttt{dest\_width}}: Stage C destination lane write-mask (\texttt{none} = full 32-bit, \texttt{h0} = low 16-bit, \texttt{h1} = high 16-bit without bit-shifting). \emph{Example:} {\scriptsize\ttfamily \mnem{FADD.v2f16} \dst{r0}\imod{.h1}, \src{r1}, \src{r2}} (\textbf{SPIR-V:} \texttt{OpVectorInsertDynamic / WriteMask}; \textbf{GLSL:} \texttt{dst.y = a.y + b.y;})
  \item \texttt{datatype}: Selects the lane data type (\texttt{f32}, \texttt{f16}, \texttt{v2f16}, \texttt{s32}, \texttt{u32}, \texttt{v2s16}, \texttt{v2u16}, \texttt{v4s8}, \texttt{v4u8}, \texttt{s64}, \texttt{u64}, etc.). \emph{Example:} {\scriptsize\ttfamily \mnem{IADD.v2s16} \dst{r0}, \src{r1}, \src{r2}} (\textbf{SPIR-V:} \texttt{OpTypeInt / OpTypeFloat}; \textbf{GLSL:} \texttt{i16vec2 c = a + b;})
  \item \texttt{instruction\_class}: 4-bit major instruction class selector (\texttt{bits 48:51}) partitioning the 64-bit Valhall opcode space (e.g., \texttt{FMA}, \texttt{CVT}, \texttt{SFU}, \texttt{LS}, \texttt{TEX}, \texttt{BRANCH}). \emph{Example:} {\scriptsize\ttfamily \mnem{FADD.f32} \dst{r0}, \src{r1}, \src{r2}} (\textbf{SPIR-V:} \texttt{N/A (encoding class)}; \textbf{GLSL:} \texttt{// Hardware execution pipe / decode class})
  \item \texttt{fau\_page\_index}: Selects the FAU (Fast Access Uniform) page (\texttt{0..3}) when a source operand references hardware small constants, thread IDs, or lane constants. \emph{Example:} {\scriptsize\ttfamily \mnem{FADD.f32} \dst{r0}, \src{r1}, \src{k1f}} (\textbf{SPIR-V:} \texttt{OpConstant / BuiltIn}; \textbf{GLSL:} \texttt{float y = x + 1.0;})
  \item \hyperref[enum:flow_control_m]{\texttt{flow}}: 4-bit embedded flow-control \& scoreboard wait modifier (\texttt{bits 59:62}: \texttt{none}, \texttt{wait0..wait0126}, \texttt{wait}, \texttt{reconverge}, \texttt{discard}, \texttt{end}). \emph{Example:} {\scriptsize\ttfamily \mnem{FADD.f32}\imod{.wait0.reconverge} \dst{r0}, \src{r1}, \src{r2}} (\textbf{SPIR-V:} \texttt{OpControlBarrier / OpReturn}; \textbf{GLSL:} \texttt{// Scoreboard wait + warp reconvergence})
\end{tightitemize}
\clearpage

\subsection*{MOV}
\addcontentsline{toc}{subsection}{MOV}
Copies the 32-bit value \texttt{src0} (GPR, FAU uniform, or inline constant) into \texttt{dest}, optionally masking the write to \texttt{.h0} or \texttt{.h1} via \texttt{dest\_width}.
\par\textbf{Notes:}
\begin{tightitemize}
  \item No swizzle or source modifiers (\texttt{encode\_v9.rs:2821-2856}).
  \item \textbf{Source references:} \texttt{kraid/ops.rs:3555-3569}, \texttt{encode\_v9.rs:2821-2856}.
\end{tightitemize}
\codebox{\mnem{MOV}\imod{\{.flow\}} \dst{dest}\imod{\{.dest\_width\}}, \src{src0}}
\textbf{Operands:}
\begin{tightitemize}\footnotesize
  \item \texttt{\dst{dest}} (\texttt{dest\_register\_index}, \texttt{dest\_width}): Destination GPR receiving \texttt{src0} (optionally halfword-masked via \texttt{dest\_width}).
  \item \texttt{\src{src0}}: 32-bit source register, uniform, or inline constant.
\end{tightitemize}
\par\textbf{Closest SPIR-V Op:} \texttt{OpCopyObject}
\par\textbf{Pseudo-GLSL (Operational Semantics):}
\codebox{dest = maskWrite(dest\_width, src0);}
\textbf{Execution Context:}
\begin{tightitemize}\footnotesize
  \item \textbf{Compute stage preloads (\hyperref[sec:csf-spd]{\texttt{SPD\_k}}, \texttt{r16..r23}):} Hardware pre-fills \texttt{r55..r57} with \texttt{WorkGroupID.xyz + JOB\_OFFSET} (from \hyperref[sec:csf-regs]{\texttt{COMPUTESR::JOB\_OFFSET\_X/Y/Z}} \texttt{r34..r36}) and \texttt{r60..r62} with \texttt{LocalInvocationID.xyz} or \texttt{GlobalInvocationID.xyz} (from \texttt{WG\_SIZE} \texttt{r33} \& \texttt{JOB\_SIZE} \texttt{r37..r39}, selected by \texttt{ShaderProgram::preload}) so compute shaders can derive thread/element indices without memory loads.
  \item \textbf{Vertex stage preloads (\hyperref[sec:csf-spd]{\texttt{SPD\_0}/\texttt{SPD\_1}}, \texttt{r16..r19}):} The hardware Index Fetcher reads \hyperref[sec:csf-regs]{\texttt{IDVSSR::INDEX\_BUFFER}} (\texttt{r54:r55}) and pre-fills \texttt{r59, r60} with \texttt{VertexIndex} (\texttt{+ VERTEX\_OFFSET + GLOBAL\_ATTR\_OFF}) and \texttt{r61} with \texttt{InstanceIndex} (\texttt{+ INSTANCE\_OFFSET}) so Position and Varying shaders can pull vertex attributes directly.
  \item \textbf{Inline FAU \& Special Operands (\hyperref[sec:csf-tsd-fau]{\texttt{FAU\_k}}, \texttt{r8..r15}):} Any 8-bit source with \texttt{mode=2} reads zero-latency uniform \texttt{u0..u31} from \texttt{FAU\_k} (\texttt{slots64 = bits[63:56]}); \texttt{mode=3} reads \texttt{SmallConstant} (Table~1.2) or \texttt{FauSpecial} (\texttt{pc}, \texttt{lane\_id}, \texttt{lane\_width}, \texttt{tls\_ptr}, \texttt{wls\_ptr}, \texttt{blend\_descriptor\_0..7}).
\end{tightitemize}
\par\textbf{Worked Examples:}
\par\textbf{Example 1 (Full 32-bit register copy):}
\codebox{\mnem{MOV.i32} \dst{r0}, \src{r1}}
With \texttt{r1 = 0xDEADBEEF}: copies all 32 bits -\textgreater{} \texttt{r0 = 0xDEADBEEF} (also loads FAU constants like \texttt{k1f}).
\par\textbf{Example 2 (Halfword lane-masked copy (`r0.h0`)):}
\codebox{\mnem{MOV.i32} \dst{r0}\imod{.h0}, \src{r1}}
With \texttt{r1 = 0x12345678}, \texttt{r0\_old = 0xAAAA0000}: \texttt{.h0} writes only \texttt{r1[15:0]} -\textgreater{} \texttt{r0 = 0xAAAA5678}.
\par\textbf{Mined SPIR-V $\to$ Valhall ISA Example (Copying FAU uniforms into contiguous staging registers for BLEND --- Fragment shader, Mali-G710 (v10)):}
\codebox{\textbf{; SPIR-V (input)}\\OpStore \%out\_color (OpCompositeConstruct \%v4float \%u\_r \%u\_g \%u\_b \%var\_a)\\[3pt]\textbf{; Valhall ISA (malioc output)}\\LD\_VAR\_IMM.slot0.discard.f32.sample.store @r3, r61\textasciicircum{}, \#0x0, \#0xb\\MOV.i32 r2, u0.w0\\MOV.i32 r1, u0.w1\\MOV.i32.wait r0, u1.w0\\BLEND.slot0.end.v4.f32 @r0:r1:r2:r3, blend\_descriptor\_0, 0x28, r60\textasciicircum{}}
\texttt{BLEND} requires its RGBA inputs in a contiguous staging-register range (\texttt{@r0:r1:r2:r3}), so \texttt{malioc} uses \texttt{MOV.i32} to stage uniform channels from \texttt{u0..u1} alongside the varying alpha in \texttt{r3}.
\begin{center}\footnotesize
\begin{tabularx}{\linewidth}{@{}lllX@{}}
\toprule
Variant & Arch & Class & Opcode Fields \\
\midrule
\texttt{(none)} & v9--v14 & \texttt{R1} & \texttt{opcode2=MOV, opcode1=MOV} \\
\bottomrule
\end{tabularx}
\end{center}
\begin{center}\footnotesize
\begin{tabularx}{0.88\linewidth}{@{}l l X@{}}
\toprule
Field Name & Bits & Type / Fixed Value \\
\midrule
\texttt{src0} & \texttt{0:7} & \texttt{SourceEncoding} \\
\textit{Reserved} & \texttt{8:15} & \texttt{zero} \\
\texttt{opcode2} & \texttt{16:20} & \texttt{exact=MOV} \\
\textit{Reserved} & \texttt{21:39} & \texttt{zero} \\
\texttt{dest\_register\_index} & \texttt{40:45} & \texttt{uint} \\
\hyperref[enum:dest_width_m]{\texttt{dest\_width}} & \texttt{46:47} & \texttt{dest\_width\_m} \\
\texttt{opcode1} & \texttt{48:51} & \texttt{exact=MOV} \\
\texttt{instruction\_class} & \texttt{52:56} & \texttt{exact=R1} \\
\texttt{fau\_page\_index} & \texttt{57:58} & \texttt{uint} \\
\hyperref[enum:flow_control_m]{\texttt{flow}} & \texttt{59:62} & \texttt{flow\_control\_m} \\
\textit{Reserved} & \texttt{63} & \texttt{zero} \\
\bottomrule
\end{tabularx}
\end{center}
\textbf{Fields Explained:}
\begin{tightitemize}\footnotesize
  \item \texttt{src0}: Primary source operand (\texttt{r0..r63}, \texttt{u0..u63} uniform, \texttt{k*} FAU constant, or \texttt{t0..t2} temporary; \texttt{\textasciicircum{}} marks discard/last-use). \emph{Example:} {\scriptsize\ttfamily \mnem{FADD.f32} \dst{r0}, \src{r1}, \src{r2}} (\textbf{SPIR-V:} \texttt{OpFAdd (Operand 0)}; \textbf{GLSL:} \texttt{float c = a + b;})
  \item \texttt{opcode2}: Secondary sub-opcode further specializing instructions sharing \texttt{instruction\_class} and \texttt{opcode1} (e.g., conversions, SFU tables, moves). \emph{Example:} {\scriptsize\ttfamily \mnem{F32\_TO\_S32}\imod{.round\_zero} \dst{r0}, \src{r1}} (\textbf{SPIR-V:} \texttt{OpConvertFToS}; \textbf{GLSL:} \texttt{int y = int(x);})
  \item \texttt{dest\_register\_index}: 6-bit destination GPR index (\texttt{r0..r63}) receiving the Stage C writeback result (even-aligned pair \texttt{[rN:rN+1]} for 64-bit ops). \emph{Example:} {\scriptsize\ttfamily \mnem{FMA.f32} \dst{r0}, \src{r1}, \src{r2}, \src{r3}} (\textbf{SPIR-V:} \texttt{SSA Result ID}; \textbf{GLSL:} \texttt{float r0 = fma(r1, r2, r3);})
  \item \hyperref[enum:dest_width_m]{\texttt{dest\_width}}: Stage C destination lane write-mask (\texttt{none} = full 32-bit, \texttt{h0} = low 16-bit, \texttt{h1} = high 16-bit without bit-shifting). \emph{Example:} {\scriptsize\ttfamily \mnem{FADD.v2f16} \dst{r0}\imod{.h1}, \src{r1}, \src{r2}} (\textbf{SPIR-V:} \texttt{OpVectorInsertDynamic / WriteMask}; \textbf{GLSL:} \texttt{dst.y = a.y + b.y;})
  \item \texttt{opcode1}: Primary sub-opcode within the instruction class uniquely identifying the operation family. \emph{Example:} {\scriptsize\ttfamily \mnem{IADD.s32} \dst{r0}, \src{r1}, \src{r2}} (\textbf{SPIR-V:} \texttt{OpIAdd (operation selector)}; \textbf{GLSL:} \texttt{// Primary opcode bits})
  \item \texttt{instruction\_class}: 4-bit major instruction class selector (\texttt{bits 48:51}) partitioning the 64-bit Valhall opcode space (e.g., \texttt{FMA}, \texttt{CVT}, \texttt{SFU}, \texttt{LS}, \texttt{TEX}, \texttt{BRANCH}). \emph{Example:} {\scriptsize\ttfamily \mnem{FADD.f32} \dst{r0}, \src{r1}, \src{r2}} (\textbf{SPIR-V:} \texttt{N/A (encoding class)}; \textbf{GLSL:} \texttt{// Hardware execution pipe / decode class})
  \item \texttt{fau\_page\_index}: Selects the FAU (Fast Access Uniform) page (\texttt{0..3}) when a source operand references hardware small constants, thread IDs, or lane constants. \emph{Example:} {\scriptsize\ttfamily \mnem{FADD.f32} \dst{r0}, \src{r1}, \src{k1f}} (\textbf{SPIR-V:} \texttt{OpConstant / BuiltIn}; \textbf{GLSL:} \texttt{float y = x + 1.0;})
  \item \hyperref[enum:flow_control_m]{\texttt{flow}}: 4-bit embedded flow-control \& scoreboard wait modifier (\texttt{bits 59:62}: \texttt{none}, \texttt{wait0..wait0126}, \texttt{wait}, \texttt{reconverge}, \texttt{discard}, \texttt{end}). \emph{Example:} {\scriptsize\ttfamily \mnem{FADD.f32}\imod{.wait0.reconverge} \dst{r0}, \src{r1}, \src{r2}} (\textbf{SPIR-V:} \texttt{OpControlBarrier / OpReturn}; \textbf{GLSL:} \texttt{// Scoreboard wait + warp reconvergence})
\end{tightitemize}
\clearpage

\subsection*{MOV\_IMM}
\addcontentsline{toc}{subsection}{MOV\_IMM}
Loads the inline 32-bit immediate \texttt{\#imm1w} into \texttt{dest} (interpreted as \texttt{i32}, \texttt{v2i16}, or \texttt{v4i8}), optionally masking the write to \texttt{.h0} or \texttt{.h1} via \texttt{dest\_width}.
\par\textbf{Notes:}
\begin{tightitemize}
  \item Alias of \texttt{IADD\_IMM} with \texttt{src0 = 0}.
  \item \textbf{Source references:} \texttt{kraid/encode\_v9.rs:2842-2847}, \texttt{isa-v9-v14.xml}.
\end{tightitemize}
\codebox{\mnem{MOV\_IMM}\imod{\{.flow\}} \dst{dest}\imod{\{.dest\_width\}}, \imm{\#imm1w}}
\textbf{Operands:}
\begin{tightitemize}\footnotesize
  \item \texttt{\dst{dest}} (\texttt{dest\_register\_index}, \texttt{dest\_width}): Destination GPR receiving \texttt{\#imm1w} (optionally halfword-masked via \texttt{dest\_width}).
  \item \texttt{\imm{\#imm1w}} (\texttt{imm1w}): Inline 32-bit immediate value.
\end{tightitemize}
\par\textbf{Closest SPIR-V Op:} \texttt{OpConstant}
\par\textbf{Pseudo-GLSL (Operational Semantics):}
\codebox{dest = maskWrite(dest\_width, \#imm1w);}
\textbf{Execution Context:}
\begin{tightitemize}\footnotesize
  \item \textbf{Compute stage preloads (\hyperref[sec:csf-spd]{\texttt{SPD\_k}}, \texttt{r16..r23}):} Hardware pre-fills \texttt{r55..r57} with \texttt{WorkGroupID.xyz + JOB\_OFFSET} (from \hyperref[sec:csf-regs]{\texttt{COMPUTESR::JOB\_OFFSET\_X/Y/Z}} \texttt{r34..r36}) and \texttt{r60..r62} with \texttt{LocalInvocationID.xyz} or \texttt{GlobalInvocationID.xyz} (from \texttt{WG\_SIZE} \texttt{r33} \& \texttt{JOB\_SIZE} \texttt{r37..r39}, selected by \texttt{ShaderProgram::preload}) so compute shaders can derive thread/element indices without memory loads.
  \item \textbf{Vertex stage preloads (\hyperref[sec:csf-spd]{\texttt{SPD\_0}/\texttt{SPD\_1}}, \texttt{r16..r19}):} The hardware Index Fetcher reads \hyperref[sec:csf-regs]{\texttt{IDVSSR::INDEX\_BUFFER}} (\texttt{r54:r55}) and pre-fills \texttt{r59, r60} with \texttt{VertexIndex} (\texttt{+ VERTEX\_OFFSET + GLOBAL\_ATTR\_OFF}) and \texttt{r61} with \texttt{InstanceIndex} (\texttt{+ INSTANCE\_OFFSET}) so Position and Varying shaders can pull vertex attributes directly.
  \item \textbf{Inline FAU \& Special Operands (\hyperref[sec:csf-tsd-fau]{\texttt{FAU\_k}}, \texttt{r8..r15}):} Any 8-bit source with \texttt{mode=2} reads zero-latency uniform \texttt{u0..u31} from \texttt{FAU\_k} (\texttt{slots64 = bits[63:56]}); \texttt{mode=3} reads \texttt{SmallConstant} (Table~1.2) or \texttt{FauSpecial} (\texttt{pc}, \texttt{lane\_id}, \texttt{lane\_width}, \texttt{tls\_ptr}, \texttt{wls\_ptr}, \texttt{blend\_descriptor\_0..7}).
\end{tightitemize}
\par\textbf{Worked Examples:}
\par\textbf{Example 1 (ASTC halfword-masked immediate load):}
\codebox{\mnem{MOV\_IMM.i32} \dst{r4}\imod{.h1}, \imm{\#0xE3D0E3D0}}
With \texttt{r4\_old = 0x00001234}: writes upper half \texttt{0xE3D0} into \texttt{r4[31:16]}, preserving \texttt{r4[15:0]} -\textgreater{} \texttt{r4 = 0xE3D01234}.
\par\textbf{Example 2 (Full 32-bit immediate constant load):}
\codebox{\mnem{MOV\_IMM.i32} \dst{r0}, \imm{\#0x3F800000}}
Loads arbitrary 32-bit constant \texttt{0x3F800000} (\texttt{1.0f}) directly into \texttt{r0} without FAU traffic.
\par\textbf{Mined SPIR-V $\to$ Valhall ISA Example (Loading packed FP16 constant colors for BLEND --- Fragment shader, Mali-G710 (v10)):}
\codebox{\textbf{; SPIR-V (input)}\\OpStore \%out\_color \%const\_v4f16\_1\_1\_1\_0\\[3pt]\textbf{; Valhall ISA (malioc output)}\\MOV\_IMM.i32.discard r1.h1, \#0x0\\MOV\_IMM.i32 r0, \#0x3c003c00\\MOV\_IMM.i32.wait r1.h0, \#0x3c003c00\\BLEND.slot0.end.v4.f16 @r0:r1, blend\_descriptor\_0, 0x20, r60\textasciicircum{}}
\texttt{MOV\_IMM.i32} loads 32-bit bit patterns (\texttt{0x3c003c00} = two FP16 \texttt{1.0h} lanes) and uses halfword write masks (\texttt{r1.h1}, \texttt{r1.h0}) to construct \texttt{(1.0h, 1.0h, 1.0h, 0.0h)}.
\begin{center}\footnotesize
\begin{tabularx}{\linewidth}{@{}lllX@{}}
\toprule
Variant & Arch & Class & Opcode Fields \\
\midrule
\texttt{(none)} & v9--v14 & \texttt{R1I} & \texttt{opcode1=ADD\_IMM\_I32} \\
\texttt{(none)} & v9--v14 & \texttt{R1I} & \texttt{opcode1=ADD\_IMM\_V2I16} \\
\texttt{(none)} & v9--v14 & \texttt{R1I} & \texttt{opcode1=ADD\_IMM\_V4I8} \\
\bottomrule
\end{tabularx}
\end{center}
\begin{center}\footnotesize
\begin{tabularx}{0.88\linewidth}{@{}l l X@{}}
\toprule
Field Name & Bits & Type / Fixed Value \\
\midrule
\texttt{src0} & \texttt{0:7} & \texttt{exact=192} \\
\texttt{imm1w} & \texttt{8:39} & \texttt{uint} \\
\texttt{dest\_register\_index} & \texttt{40:45} & \texttt{uint} \\
\hyperref[enum:dest_width_m]{\texttt{dest\_width}} & \texttt{46:47} & \texttt{dest\_width\_m} \\
\texttt{opcode1} & \texttt{48:51} & \texttt{exact=ADD\_IMM\_I32} \\
\texttt{instruction\_class} & \texttt{52:56} & \texttt{exact=R1I} \\
\texttt{fau\_page\_index} & \texttt{57:58} & \texttt{exact=0} \\
\hyperref[enum:flow_control_m]{\texttt{flow}} & \texttt{59:62} & \texttt{flow\_control\_m} \\
\textit{Reserved} & \texttt{63} & \texttt{zero} \\
\bottomrule
\end{tabularx}
\end{center}
\textbf{Fields Explained:}
\begin{tightitemize}\footnotesize
  \item \texttt{src0}: Primary source operand (\texttt{r0..r63}, \texttt{u0..u63} uniform, \texttt{k*} FAU constant, or \texttt{t0..t2} temporary; \texttt{\textasciicircum{}} marks discard/last-use). \emph{Example:} {\scriptsize\ttfamily \mnem{FADD.f32} \dst{r0}, \src{r1}, \src{r2}} (\textbf{SPIR-V:} \texttt{OpFAdd (Operand 0)}; \textbf{GLSL:} \texttt{float c = a + b;})
  \item \texttt{imm1w}: 32-bit inline immediate operand encoded directly in the instruction word (bits 8:39). \emph{Example:} {\scriptsize\ttfamily \mnem{IADD\_IMM.i32} \dst{r0}, \src{r1}, \imm{\#0x10}} (\textbf{SPIR-V:} \texttt{OpIAdd (constant)}; \textbf{GLSL:} \texttt{int y = x + 16;})
  \item \texttt{dest\_register\_index}: 6-bit destination GPR index (\texttt{r0..r63}) receiving the Stage C writeback result (even-aligned pair \texttt{[rN:rN+1]} for 64-bit ops). \emph{Example:} {\scriptsize\ttfamily \mnem{FMA.f32} \dst{r0}, \src{r1}, \src{r2}, \src{r3}} (\textbf{SPIR-V:} \texttt{SSA Result ID}; \textbf{GLSL:} \texttt{float r0 = fma(r1, r2, r3);})
  \item \hyperref[enum:dest_width_m]{\texttt{dest\_width}}: Stage C destination lane write-mask (\texttt{none} = full 32-bit, \texttt{h0} = low 16-bit, \texttt{h1} = high 16-bit without bit-shifting). \emph{Example:} {\scriptsize\ttfamily \mnem{FADD.v2f16} \dst{r0}\imod{.h1}, \src{r1}, \src{r2}} (\textbf{SPIR-V:} \texttt{OpVectorInsertDynamic / WriteMask}; \textbf{GLSL:} \texttt{dst.y = a.y + b.y;})
  \item \texttt{opcode1}: Primary sub-opcode within the instruction class uniquely identifying the operation family. \emph{Example:} {\scriptsize\ttfamily \mnem{IADD.s32} \dst{r0}, \src{r1}, \src{r2}} (\textbf{SPIR-V:} \texttt{OpIAdd (operation selector)}; \textbf{GLSL:} \texttt{// Primary opcode bits})
  \item \texttt{instruction\_class}: 4-bit major instruction class selector (\texttt{bits 48:51}) partitioning the 64-bit Valhall opcode space (e.g., \texttt{FMA}, \texttt{CVT}, \texttt{SFU}, \texttt{LS}, \texttt{TEX}, \texttt{BRANCH}). \emph{Example:} {\scriptsize\ttfamily \mnem{FADD.f32} \dst{r0}, \src{r1}, \src{r2}} (\textbf{SPIR-V:} \texttt{N/A (encoding class)}; \textbf{GLSL:} \texttt{// Hardware execution pipe / decode class})
  \item \texttt{fau\_page\_index}: Selects the FAU (Fast Access Uniform) page (\texttt{0..3}) when a source operand references hardware small constants, thread IDs, or lane constants. \emph{Example:} {\scriptsize\ttfamily \mnem{FADD.f32} \dst{r0}, \src{r1}, \src{k1f}} (\textbf{SPIR-V:} \texttt{OpConstant / BuiltIn}; \textbf{GLSL:} \texttt{float y = x + 1.0;})
  \item \hyperref[enum:flow_control_m]{\texttt{flow}}: 4-bit embedded flow-control \& scoreboard wait modifier (\texttt{bits 59:62}: \texttt{none}, \texttt{wait0..wait0126}, \texttt{wait}, \texttt{reconverge}, \texttt{discard}, \texttt{end}). \emph{Example:} {\scriptsize\ttfamily \mnem{FADD.f32}\imod{.wait0.reconverge} \dst{r0}, \src{r1}, \src{r2}} (\textbf{SPIR-V:} \texttt{OpControlBarrier / OpReturn}; \textbf{GLSL:} \texttt{// Scoreboard wait + warp reconvergence})
\end{tightitemize}
\clearpage

\subsection*{SWZ}
\addcontentsline{toc}{subsection}{SWZ}
Permutes packed 16-bit (\texttt{v2i16}) or 8-bit (\texttt{v4i8}, arch \textless{}= v10) lanes of \texttt{src0} according to \texttt{lanes0} and writes the swizzled vector to \texttt{dest}. Encoded as an alias of \texttt{IADD} with \texttt{src1 = 0}.
\par\textbf{Notes:}
\begin{tightitemize}
  \item In kraid, \texttt{OpSwz} is a virtual op lowered via copy/\texttt{IADD}/\texttt{ShiftLop}/\texttt{ARShift}/\texttt{MKVEC} (\texttt{lower\_mkvec\_swz.rs}).
  \item \textbf{Source references:} \texttt{kraid/ops.rs:4159-4204}, \texttt{lower\_mkvec\_swz.rs:137-499}, \texttt{hw\_tests.rs:2239-2289}.
\end{tightitemize}
\codebox{\mnem{SWZ}\imod{\{.flow\}} \dst{dest}\imod{\{.dest\_width\}}, \src{src0}\imod{\{.swz0\}}\\\mnem{SWZ}\imod{\{.flow\}} \dst{dest}\imod{\{.dest\_width\}}, \src{src0}\imod{\{.lanes0\}}}
\textbf{Operands:}
\begin{tightitemize}\footnotesize
  \item \texttt{\dst{dest}} (\texttt{dest\_register\_index}, \texttt{dest\_width}): 32-bit destination GPR receiving swizzled 16-bit or 8-bit lanes of \texttt{src0}.
  \item \texttt{\src{src0}}: Packed \texttt{v2i16} or \texttt{v4i8} source register.
\end{tightitemize}
\par\textbf{Closest SPIR-V Op:} \texttt{OpVectorShuffle}
\par\textbf{Pseudo-GLSL (Operational Semantics):}
\codebox{for (int i = 0; i \textless{} N; ++i) dest[i] = src0[lanes0[i]];}
\textbf{Execution Context:}
\begin{tightitemize}\footnotesize
  \item \textbf{Compute stage preloads (\hyperref[sec:csf-spd]{\texttt{SPD\_k}}, \texttt{r16..r23}):} Hardware pre-fills \texttt{r55..r57} with \texttt{WorkGroupID.xyz + JOB\_OFFSET} (from \hyperref[sec:csf-regs]{\texttt{COMPUTESR::JOB\_OFFSET\_X/Y/Z}} \texttt{r34..r36}) and \texttt{r60..r62} with \texttt{LocalInvocationID.xyz} or \texttt{GlobalInvocationID.xyz} (from \texttt{WG\_SIZE} \texttt{r33} \& \texttt{JOB\_SIZE} \texttt{r37..r39}, selected by \texttt{ShaderProgram::preload}) so compute shaders can derive thread/element indices without memory loads.
  \item \textbf{Vertex stage preloads (\hyperref[sec:csf-spd]{\texttt{SPD\_0}/\texttt{SPD\_1}}, \texttt{r16..r19}):} The hardware Index Fetcher reads \hyperref[sec:csf-regs]{\texttt{IDVSSR::INDEX\_BUFFER}} (\texttt{r54:r55}) and pre-fills \texttt{r59, r60} with \texttt{VertexIndex} (\texttt{+ VERTEX\_OFFSET + GLOBAL\_ATTR\_OFF}) and \texttt{r61} with \texttt{InstanceIndex} (\texttt{+ INSTANCE\_OFFSET}) so Position and Varying shaders can pull vertex attributes directly.
  \item \textbf{Inline FAU \& Special Operands (\hyperref[sec:csf-tsd-fau]{\texttt{FAU\_k}}, \texttt{r8..r15}):} Any 8-bit source with \texttt{mode=2} reads zero-latency uniform \texttt{u0..u31} from \texttt{FAU\_k} (\texttt{slots64 = bits[63:56]}); \texttt{mode=3} reads \texttt{SmallConstant} (Table~1.2) or \texttt{FauSpecial} (\texttt{pc}, \texttt{lane\_id}, \texttt{lane\_width}, \texttt{tls\_ptr}, \texttt{wls\_ptr}, \texttt{blend\_descriptor\_0..7}).
\end{tightitemize}
\par\textbf{Worked Examples:}
\par\textbf{Example 1 (Byte-reverse endianness swap (ETC2 idiom)):}
\codebox{\mnem{SWZ.v4i8} \dst{r0}, \src{r1}\imod{.b3210}}
With \texttt{r1 = 0x11223344}: reverses all 4 bytes -\textgreater{} \texttt{r0 = 0x44332211}.
\par\textbf{Example 2 (Splat single byte across all 4 lanes):}
\codebox{\mnem{SWZ.v4i8} \dst{r0}, \src{r1}\imod{.b1111}}
With \texttt{r1 = 0x11223344}: replicates byte 1 (\texttt{0x33}) to all 4 lanes -\textgreater{} \texttt{r0 = 0x33333333}.

\begin{center}\footnotesize
\begin{tabularx}{\linewidth}{@{}lllX@{}}
\toprule
Variant & Arch & Class & Opcode Fields \\
\midrule
\texttt{v2u16} & v9--v14 & \texttt{R2} & \texttt{opcode2=ADD} \\
\texttt{v4u8} & v9--v14 & \texttt{R2} & \texttt{opcode2=ADD} \\
\bottomrule
\end{tabularx}
\end{center}
\begin{center}\footnotesize
\begin{tabularx}{0.88\linewidth}{@{}l l X@{}}
\toprule
Field Name & Bits & Type / Fixed Value \\
\midrule
\texttt{src0} & \texttt{0:7} & \texttt{SourceEncoding} \\
\texttt{src1} & \texttt{8:15} & \texttt{exact=192} \\
\texttt{opcode2} & \texttt{16:19} & \texttt{exact=ADD} \\
\textit{Reserved} & \texttt{20:25} & \texttt{zero} \\
\hyperref[enum:swiz_int_m]{\texttt{swz1}} & \texttt{26:29} & \texttt{exact=h01} \\
\hyperref[enum:saturate_m]{\texttt{saturate}} & \texttt{30} & \texttt{exact=none} \\
\textit{Reserved} & \texttt{31:35} & \texttt{zero} \\
\hyperref[enum:swiz_int_m]{\texttt{swz0}} & \texttt{36:39} & \texttt{swiz\_int\_m (def h01)} \\
\texttt{dest\_register\_index} & \texttt{40:45} & \texttt{uint} \\
\hyperref[enum:dest_width_m]{\texttt{dest\_width}} & \texttt{46:47} & \texttt{dest\_width\_m} \\
\texttt{datatype} & \texttt{48:51} & \texttt{exact=v2u16} \\
\texttt{instruction\_class} & \texttt{52:56} & \texttt{exact=R2} \\
\texttt{fau\_page\_index} & \texttt{57:58} & \texttt{uint} \\
\hyperref[enum:flow_control_m]{\texttt{flow}} & \texttt{59:62} & \texttt{flow\_control\_m} \\
\textit{Reserved} & \texttt{63} & \texttt{zero} \\
\bottomrule
\end{tabularx}
\end{center}
\textbf{Fields Explained:}
\begin{tightitemize}\footnotesize
  \item \texttt{src0}: Primary source operand (\texttt{r0..r63}, \texttt{u0..u63} uniform, \texttt{k*} FAU constant, or \texttt{t0..t2} temporary; \texttt{\textasciicircum{}} marks discard/last-use). \emph{Example:} {\scriptsize\ttfamily \mnem{FADD.f32} \dst{r0}, \src{r1}, \src{r2}} (\textbf{SPIR-V:} \texttt{OpFAdd (Operand 0)}; \textbf{GLSL:} \texttt{float c = a + b;})
  \item \texttt{src1}: Secondary source operand (register, uniform, or FAU constant; in shifts, supplies the shift count byte). \emph{Example:} {\scriptsize\ttfamily \mnem{LSHIFT.i32} \dst{r0}, \src{r1}, \src{r2}\imod{.b0}} (\textbf{SPIR-V:} \texttt{OpShiftLeftLogical (Shift)}; \textbf{GLSL:} \texttt{uint y = a \textless{}\textless{} s;})
  \item \texttt{opcode2}: Secondary sub-opcode further specializing instructions sharing \texttt{instruction\_class} and \texttt{opcode1} (e.g., conversions, SFU tables, moves). \emph{Example:} {\scriptsize\ttfamily \mnem{F32\_TO\_S32}\imod{.round\_zero} \dst{r0}, \src{r1}} (\textbf{SPIR-V:} \texttt{OpConvertFToS}; \textbf{GLSL:} \texttt{int y = int(x);})
  \item \hyperref[enum:swiz_int_m]{\texttt{swz1}}: Stage A lane swizzle on \texttt{src1} (\texttt{h00}/\texttt{h01}/\texttt{h10}/\texttt{h11} or \texttt{bXY} byte-to-halfword widening). \emph{Example:} {\scriptsize\ttfamily \mnem{FADD.v2f16} \dst{r0}, \src{r1}, \src{r2}\imod{.h00}} (\textbf{SPIR-V:} \texttt{OpVectorShuffle}; \textbf{GLSL:} \texttt{f16vec2 c = a + b.xx;})
  \item \hyperref[enum:saturate_m]{\texttt{saturate}}: Stage B integer saturation flag (\texttt{.sat}): clamps signed/unsigned overflow or underflow to \texttt{[MIN\_N, MAX\_N]} / \texttt{[0, 2\textasciicircum{}N - 1]} instead of wrapping. \emph{Example:} {\scriptsize\ttfamily \mnem{IADD.s32}\imod{.sat} \dst{r0}, \src{r1}, \src{r2}} (\textbf{SPIR-V:} \texttt{OpSDotAccSat / saturating add}; \textbf{GLSL:} \texttt{int c = clamp(int64\_t(a) + int64\_t(b), INT\_MIN, INT\_MAX);})
  \item \hyperref[enum:swiz_int_m]{\texttt{swz0}}: Stage A lane swizzle on \texttt{src0} (\texttt{h00}/\texttt{h01}/\texttt{h10}/\texttt{h11} for 16-bit vectors, or \texttt{bXY} byte-to-halfword widening). \emph{Example:} {\scriptsize\ttfamily \mnem{FADD.v2f16} \dst{r0}, \src{r1}\imod{.h10}, \src{r2}} (\textbf{SPIR-V:} \texttt{OpVectorShuffle}; \textbf{GLSL:} \texttt{f16vec2 c = a.yx + b;})
  \item \texttt{dest\_register\_index}: 6-bit destination GPR index (\texttt{r0..r63}) receiving the Stage C writeback result (even-aligned pair \texttt{[rN:rN+1]} for 64-bit ops). \emph{Example:} {\scriptsize\ttfamily \mnem{FMA.f32} \dst{r0}, \src{r1}, \src{r2}, \src{r3}} (\textbf{SPIR-V:} \texttt{SSA Result ID}; \textbf{GLSL:} \texttt{float r0 = fma(r1, r2, r3);})
  \item \hyperref[enum:dest_width_m]{\texttt{dest\_width}}: Stage C destination lane write-mask (\texttt{none} = full 32-bit, \texttt{h0} = low 16-bit, \texttt{h1} = high 16-bit without bit-shifting). \emph{Example:} {\scriptsize\ttfamily \mnem{FADD.v2f16} \dst{r0}\imod{.h1}, \src{r1}, \src{r2}} (\textbf{SPIR-V:} \texttt{OpVectorInsertDynamic / WriteMask}; \textbf{GLSL:} \texttt{dst.y = a.y + b.y;})
  \item \texttt{datatype}: Selects the lane data type (\texttt{f32}, \texttt{f16}, \texttt{v2f16}, \texttt{s32}, \texttt{u32}, \texttt{v2s16}, \texttt{v2u16}, \texttt{v4s8}, \texttt{v4u8}, \texttt{s64}, \texttt{u64}, etc.). \emph{Example:} {\scriptsize\ttfamily \mnem{IADD.v2s16} \dst{r0}, \src{r1}, \src{r2}} (\textbf{SPIR-V:} \texttt{OpTypeInt / OpTypeFloat}; \textbf{GLSL:} \texttt{i16vec2 c = a + b;})
  \item \texttt{instruction\_class}: 4-bit major instruction class selector (\texttt{bits 48:51}) partitioning the 64-bit Valhall opcode space (e.g., \texttt{FMA}, \texttt{CVT}, \texttt{SFU}, \texttt{LS}, \texttt{TEX}, \texttt{BRANCH}). \emph{Example:} {\scriptsize\ttfamily \mnem{FADD.f32} \dst{r0}, \src{r1}, \src{r2}} (\textbf{SPIR-V:} \texttt{N/A (encoding class)}; \textbf{GLSL:} \texttt{// Hardware execution pipe / decode class})
  \item \texttt{fau\_page\_index}: Selects the FAU (Fast Access Uniform) page (\texttt{0..3}) when a source operand references hardware small constants, thread IDs, or lane constants. \emph{Example:} {\scriptsize\ttfamily \mnem{FADD.f32} \dst{r0}, \src{r1}, \src{k1f}} (\textbf{SPIR-V:} \texttt{OpConstant / BuiltIn}; \textbf{GLSL:} \texttt{float y = x + 1.0;})
  \item \hyperref[enum:flow_control_m]{\texttt{flow}}: 4-bit embedded flow-control \& scoreboard wait modifier (\texttt{bits 59:62}: \texttt{none}, \texttt{wait0..wait0126}, \texttt{wait}, \texttt{reconverge}, \texttt{discard}, \texttt{end}). \emph{Example:} {\scriptsize\ttfamily \mnem{FADD.f32}\imod{.wait0.reconverge} \dst{r0}, \src{r1}, \src{r2}} (\textbf{SPIR-V:} \texttt{OpControlBarrier / OpReturn}; \textbf{GLSL:} \texttt{// Scoreboard wait + warp reconvergence})
\end{tightitemize}
\clearpage

\section{Subgroup \& Cross-Lane (3 mnemonics)}
\subsection*{CLPER}
\addcontentsline{toc}{subsection}{CLPER}
Cross-lane permute: reads \texttt{src0} from target lane computed from \texttt{lane\_id} and \texttt{src1} (8-bit via \texttt{lanes1}) using \texttt{lane\_op} (\texttt{none}/\texttt{xor}/\texttt{accumulate}/\texttt{shift}/\texttt{rotate}/\texttt{low}/\texttt{low\_alt}/\texttt{prefix}) within \texttt{subgroup}; returns \texttt{inactive\_result} constant if target lane is inactive/OOB.
\par\textbf{Notes:}
\begin{tightitemize}
  \item Used for \texttt{ddx}/\texttt{ddy} (\texttt{builder.rs:263-288}) and subgroup shuffles/reductions.
  \item \textbf{Source references:} \texttt{kraid/ops.rs:537-690}, \texttt{encode\_v9.rs:830-900}, \texttt{builder.rs:260-289}, \texttt{bifrost/valhall/ISA.xml:23}.
\end{tightitemize}
\codebox{\mnem{CLPER}\imod{\{.flow\}}\imod{\{.subgroup\}}\imod{\{.lane\_op\}}\imod{\{.inactive\_result\}} \dst{dest}\imod{\{.dest\_width\}}, \src{src0}, \src{src1}\imod{\{.lanes1\}}}
\textbf{Operands:}
\begin{tightitemize}\footnotesize
  \item \texttt{\dst{dest}} (\texttt{dest\_register\_index}, \texttt{dest\_width}): Destination GPR receiving the value read from the target peer lane (or \texttt{inactive\_result} if inactive).
  \item \texttt{\src{src0}}: Per-lane data value shared across the subgroup.
  \item \texttt{\src{src1}}: Byte 0 (\texttt{src1.b0}) supplies the target lane index, XOR mask, or shift delta per \texttt{lane\_op}.
\end{tightitemize}
\par\textbf{Closest SPIR-V Op:} \texttt{OpGroupNonUniformShuffle* / OpDPdx / OpDPdy}
\par\textbf{Pseudo-GLSL (Operational Semantics):}
\codebox{uint tgt = laneOp(lane\_id, src1 \& 0xFFu, subgroup\_size); // none/xor/accumulate/shift/rotate/low/prefix\\dest = laneActive(tgt) ? readLane(src0, tgt) : inactiveConst(inactive\_result);}
\textbf{Execution Context:}
\begin{tightitemize}\footnotesize
  \item \textbf{Compute stage preloads (\hyperref[sec:csf-spd]{\texttt{SPD\_k}}, \texttt{r16..r23}):} Hardware pre-fills \texttt{r55..r57} with \texttt{WorkGroupID.xyz + JOB\_OFFSET} (from \hyperref[sec:csf-regs]{\texttt{COMPUTESR::JOB\_OFFSET\_X/Y/Z}} \texttt{r34..r36}) and \texttt{r60..r62} with \texttt{LocalInvocationID.xyz} or \texttt{GlobalInvocationID.xyz} (from \texttt{WG\_SIZE} \texttt{r33} \& \texttt{JOB\_SIZE} \texttt{r37..r39}, selected by \texttt{ShaderProgram::preload}) so compute shaders can derive thread/element indices without memory loads.
  \item \textbf{Vertex stage preloads (\hyperref[sec:csf-spd]{\texttt{SPD\_0}/\texttt{SPD\_1}}, \texttt{r16..r19}):} The hardware Index Fetcher reads \hyperref[sec:csf-regs]{\texttt{IDVSSR::INDEX\_BUFFER}} (\texttt{r54:r55}) and pre-fills \texttt{r59, r60} with \texttt{VertexIndex} (\texttt{+ VERTEX\_OFFSET + GLOBAL\_ATTR\_OFF}) and \texttt{r61} with \texttt{InstanceIndex} (\texttt{+ INSTANCE\_OFFSET}) so Position and Varying shaders can pull vertex attributes directly.
  \item \textbf{Inline FAU \& Special Operands (\hyperref[sec:csf-tsd-fau]{\texttt{FAU\_k}}, \texttt{r8..r15}):} Any 8-bit source with \texttt{mode=2} reads zero-latency uniform \texttt{u0..u31} from \texttt{FAU\_k} (\texttt{slots64 = bits[63:56]}); \texttt{mode=3} reads \texttt{SmallConstant} (Table~1.2) or \texttt{FauSpecial} (\texttt{pc}, \texttt{lane\_id}, \texttt{lane\_width}, \texttt{tls\_ptr}, \texttt{wls\_ptr}, \texttt{blend\_descriptor\_0..7}).
\end{tightitemize}
\par\textbf{Worked Examples:}
\par\textbf{Example 1 (16-lane butterfly XOR shuffle):}
\codebox{\mnem{CLPER}\imod{.subgroup16.xor.zero} \dst{r0}, \src{r1}, \src{r2}\imod{.b0}}
Each lane \texttt{L} reads \texttt{r1} from lane \texttt{L \textasciicircum{} (r2 \& 0xF)}: with \texttt{r2 = 1}, lanes \texttt{(0,1), (2,3), ...} swap values; inactive lanes return \texttt{0} (\texttt{.zero}).
\par\textbf{Example 2 (Subgroup min reduction identity (`.inf`)):}
\codebox{\mnem{CLPER}\imod{.subgroup16.xor.inf} \dst{r0}, \src{r1}, \src{r2}\imod{.b0}}
When partner lane \texttt{L \textasciicircum{} r2} is inactive under divergence, \texttt{.inf} substitutes \texttt{+Inf} (\texttt{0x7F800000}) so subsequent \texttt{FMIN} is unaffected.
\par\textbf{Mined SPIR-V $\to$ Valhall ISA Example (Computing fragment screen-space derivative OpDPdx via 2x2 quad shuffle --- Fragment shader, Mali-G710 (v10)):}
\codebox{\textbf{; SPIR-V (input)}\\\%dx = OpDPdx \%float \%var\_x\\[3pt]\textbf{; Valhall ISA (malioc output)}\\LD\_VAR\_IMM.slot0.wait0.v4.f32.sample.store @r12:r13:r14:r15, r61\textasciicircum{}, \#0x0, \#0xb\\CLPER.subgroup4.zero r51, r12, k01000000.b3 ; read lane 1 of 2x2 quad\\CLPER.subgroup4.zero r1, r12, k0.b0         ; read lane 0 of 2x2 quad\\; followed by FADD.f32 dest, r51, r1.neg   -\textgreater{} lane1 - lane0}
Fragment shaders execute in 2x2 pixel quads (\texttt{subgroup4}). \texttt{OpDPdx} lowers to two \texttt{CLPER.subgroup4.zero} instructions fetching quad-lane 1 (\texttt{k01000000.b3 = 1}) and quad-lane 0 (\texttt{k0.b0 = 0}), then subtracting them.
\begin{center}\footnotesize
\begin{tabularx}{\linewidth}{@{}lllX@{}}
\toprule
Variant & Arch & Class & Opcode Fields \\
\midrule
\texttt{u32} & v9--v14 & \texttt{R2} & \texttt{opcode2=CLPER} \\
\texttt{u32} & v9--v14 & \texttt{R2} & \texttt{opcode2=CLPER} \\
\bottomrule
\end{tabularx}
\end{center}
\begin{center}\footnotesize
\begin{tabularx}{0.88\linewidth}{@{}l l X@{}}
\toprule
Field Name & Bits & Type / Fixed Value \\
\midrule
\texttt{src0} & \texttt{0:7} & \texttt{SourceEncoding} \\
\texttt{src1} & \texttt{8:15} & \texttt{SourceEncoding} \\
\texttt{opcode2} & \texttt{16:19} & \texttt{exact=CLPER} \\
\textit{Reserved} & \texttt{20:21} & \texttt{zero} \\
\hyperref[enum:clper_inactive_result_m]{\texttt{inactive\_result}} & \texttt{22:25} & \texttt{clper\_inactive\_result\_m (def zero)} \\
\hyperref[enum:shift_bitwise_i32_lanes_m]{\texttt{lanes1}} & \texttt{26:29} & \texttt{shift\_bitwise\_i32\_lanes\_m (def b0)} \\
\textit{Reserved} & \texttt{30:31} & \texttt{zero} \\
\hyperref[enum:clper_lane_op_m]{\texttt{lane\_op}} & \texttt{32:35} & \texttt{clper\_lane\_op\_m} \\
\hyperref[enum:clper_subgroup_size_m]{\texttt{subgroup}} & \texttt{36:39} & \texttt{clper\_subgroup\_size\_m} \\
\texttt{dest\_register\_index} & \texttt{40:45} & \texttt{uint} \\
\hyperref[enum:dest_width_m]{\texttt{dest\_width}} & \texttt{46:47} & \texttt{dest\_width\_m} \\
\texttt{datatype} & \texttt{48:51} & \texttt{exact=u32} \\
\texttt{instruction\_class} & \texttt{52:56} & \texttt{exact=R2} \\
\texttt{fau\_page\_index} & \texttt{57:58} & \texttt{uint} \\
\hyperref[enum:flow_control_m]{\texttt{flow}} & \texttt{59:62} & \texttt{flow\_control\_m} \\
\textit{Reserved} & \texttt{63} & \texttt{zero} \\
\bottomrule
\end{tabularx}
\end{center}
\textbf{Fields Explained:}
\begin{tightitemize}\footnotesize
  \item \texttt{src0}: Primary source operand (\texttt{r0..r63}, \texttt{u0..u63} uniform, \texttt{k*} FAU constant, or \texttt{t0..t2} temporary; \texttt{\textasciicircum{}} marks discard/last-use). \emph{Example:} {\scriptsize\ttfamily \mnem{FADD.f32} \dst{r0}, \src{r1}, \src{r2}} (\textbf{SPIR-V:} \texttt{OpFAdd (Operand 0)}; \textbf{GLSL:} \texttt{float c = a + b;})
  \item \texttt{src1}: Secondary source operand (register, uniform, or FAU constant; in shifts, supplies the shift count byte). \emph{Example:} {\scriptsize\ttfamily \mnem{LSHIFT.i32} \dst{r0}, \src{r1}, \src{r2}\imod{.b0}} (\textbf{SPIR-V:} \texttt{OpShiftLeftLogical (Shift)}; \textbf{GLSL:} \texttt{uint y = a \textless{}\textless{} s;})
  \item \texttt{opcode2}: Secondary sub-opcode further specializing instructions sharing \texttt{instruction\_class} and \texttt{opcode1} (e.g., conversions, SFU tables, moves). \emph{Example:} {\scriptsize\ttfamily \mnem{F32\_TO\_S32}\imod{.round\_zero} \dst{r0}, \src{r1}} (\textbf{SPIR-V:} \texttt{OpConvertFToS}; \textbf{GLSL:} \texttt{int y = int(x);})
  \item \hyperref[enum:clper_inactive_result_m]{\texttt{inactive\_result}}: Identity value substituted by \texttt{CLPER} when the source lane in the subgroup is inactive. \emph{Example:} {\scriptsize\ttfamily \mnem{CLPER}\imod{.subgroup16.xor.inf} \dst{r0}, \src{r1}, \src{r2}} (\textbf{SPIR-V:} \texttt{OpGroupNonUniformFMin}; \textbf{GLSL:} \texttt{float m = subgroupMin(x);})
  \item \hyperref[enum:shift_bitwise_i32_lanes_m]{\texttt{lanes1}}: Stage A byte-lane selector on \texttt{src1} (selects shift-count byte \texttt{b0..b3} on 32/64-bit shifts, \texttt{bXY} on \texttt{v2i16} shifts, or 4-byte swizzle on \texttt{v4*8}). \emph{Example:} {\scriptsize\ttfamily \mnem{LSHIFT.i32} \dst{r0}, \src{r1}, \src{r2}\imod{.b1}} (\textbf{SPIR-V:} \texttt{OpBitFieldUExtract + OpShiftLeftLogical}; \textbf{GLSL:} \texttt{uint y = a \textless{}\textless{} ((s \textgreater{}\textgreater{} 8) \& 0xFFu);})
  \item \hyperref[enum:clper_lane_op_m]{\texttt{lane\_op}}: Subgroup lane-index transformation in \texttt{CLPER} (\texttt{none}, \texttt{xor}, \texttt{shift}, \texttt{rotate}, \texttt{low}, \texttt{prefix}). \emph{Example:} {\scriptsize\ttfamily \mnem{CLPER}\imod{.subgroup16.xor.zero} \dst{r0}, \src{r1}, \src{r2}} (\textbf{SPIR-V:} \texttt{OpGroupNonUniformShuffleXor}; \textbf{GLSL:} \texttt{uint v = subgroupShuffleXor(x, 1u);})
  \item \hyperref[enum:clper_subgroup_size_m]{\texttt{subgroup}}: Subgroup partition width (\texttt{subgroup2}, \texttt{subgroup4}, \texttt{subgroup8}, \texttt{subgroup16}). \emph{Example:} {\scriptsize\ttfamily \mnem{CLPER}\imod{.subgroup8.rotate.zero} \dst{r0}, \src{r1}, \src{r2}} (\textbf{SPIR-V:} \texttt{OpGroupNonUniformRotateKHR}; \textbf{GLSL:} \texttt{uint v = subgroupRotate(x, 1);})
  \item \texttt{dest\_register\_index}: 6-bit destination GPR index (\texttt{r0..r63}) receiving the Stage C writeback result (even-aligned pair \texttt{[rN:rN+1]} for 64-bit ops). \emph{Example:} {\scriptsize\ttfamily \mnem{FMA.f32} \dst{r0}, \src{r1}, \src{r2}, \src{r3}} (\textbf{SPIR-V:} \texttt{SSA Result ID}; \textbf{GLSL:} \texttt{float r0 = fma(r1, r2, r3);})
  \item \hyperref[enum:dest_width_m]{\texttt{dest\_width}}: Stage C destination lane write-mask (\texttt{none} = full 32-bit, \texttt{h0} = low 16-bit, \texttt{h1} = high 16-bit without bit-shifting). \emph{Example:} {\scriptsize\ttfamily \mnem{FADD.v2f16} \dst{r0}\imod{.h1}, \src{r1}, \src{r2}} (\textbf{SPIR-V:} \texttt{OpVectorInsertDynamic / WriteMask}; \textbf{GLSL:} \texttt{dst.y = a.y + b.y;})
  \item \texttt{datatype}: Selects the lane data type (\texttt{f32}, \texttt{f16}, \texttt{v2f16}, \texttt{s32}, \texttt{u32}, \texttt{v2s16}, \texttt{v2u16}, \texttt{v4s8}, \texttt{v4u8}, \texttt{s64}, \texttt{u64}, etc.). \emph{Example:} {\scriptsize\ttfamily \mnem{IADD.v2s16} \dst{r0}, \src{r1}, \src{r2}} (\textbf{SPIR-V:} \texttt{OpTypeInt / OpTypeFloat}; \textbf{GLSL:} \texttt{i16vec2 c = a + b;})
  \item \texttt{instruction\_class}: 4-bit major instruction class selector (\texttt{bits 48:51}) partitioning the 64-bit Valhall opcode space (e.g., \texttt{FMA}, \texttt{CVT}, \texttt{SFU}, \texttt{LS}, \texttt{TEX}, \texttt{BRANCH}). \emph{Example:} {\scriptsize\ttfamily \mnem{FADD.f32} \dst{r0}, \src{r1}, \src{r2}} (\textbf{SPIR-V:} \texttt{N/A (encoding class)}; \textbf{GLSL:} \texttt{// Hardware execution pipe / decode class})
  \item \texttt{fau\_page\_index}: Selects the FAU (Fast Access Uniform) page (\texttt{0..3}) when a source operand references hardware small constants, thread IDs, or lane constants. \emph{Example:} {\scriptsize\ttfamily \mnem{FADD.f32} \dst{r0}, \src{r1}, \src{k1f}} (\textbf{SPIR-V:} \texttt{OpConstant / BuiltIn}; \textbf{GLSL:} \texttt{float y = x + 1.0;})
  \item \hyperref[enum:flow_control_m]{\texttt{flow}}: 4-bit embedded flow-control \& scoreboard wait modifier (\texttt{bits 59:62}: \texttt{none}, \texttt{wait0..wait0126}, \texttt{wait}, \texttt{reconverge}, \texttt{discard}, \texttt{end}). \emph{Example:} {\scriptsize\ttfamily \mnem{FADD.f32}\imod{.wait0.reconverge} \dst{r0}, \src{r1}, \src{r2}} (\textbf{SPIR-V:} \texttt{OpControlBarrier / OpReturn}; \textbf{GLSL:} \texttt{// Scoreboard wait + warp reconvergence})
\end{tightitemize}
\clearpage

\subsection*{CLPER\_OR}
\addcontentsline{toc}{subsection}{CLPER\_OR}
Cross-lane permute + bitwise OR (v14+): permutes \texttt{src0} across \texttt{subgroup} using \texttt{lane\_select} (\texttt{lanes04}..\texttt{lanesBF}) and \texttt{src1} (\texttt{lanes1}), then ORs with \texttt{src2}.
\par\textbf{Notes:}
\begin{tightitemize}
  \item Has \texttt{lane\_select} (NOT \texttt{lane\_op} or \texttt{inactive\_result}).
  \item \textbf{Source references:} \texttt{isa-v9-v14.xml}.
\end{tightitemize}
\codebox{\mnem{CLPER\_OR}\imod{\{.flow\}}\imod{\{.subgroup\}}\imod{\{.lane\_select\}} \dst{dest}\imod{\{.dest\_width\}}, \src{src0}, \src{src1}\imod{\{.lanes1\}}, \src{src2}}
\textbf{Operands:}
\begin{tightitemize}\footnotesize
  \item \texttt{\dst{dest}} (\texttt{dest\_register\_index}, \texttt{dest\_width}): Destination GPR receiving \texttt{subgroupPermute(src0, src1.b0) \textbar{} src2} (v14+).
  \item \texttt{\src{src0}}: Per-lane data value shared across the subgroup.
  \item \texttt{\src{src1}}: Byte 0 supplies the peer-lane index or shuffle operand.
  \item \texttt{\src{src2}}: 32-bit value bitwise-ORed into the shuffled result.
\end{tightitemize}
\par\textbf{Closest SPIR-V Op:} \texttt{OpGroupNonUniformShuffle* + OpBitwiseOr (v14+)}
\par\textbf{Pseudo-GLSL (Operational Semantics):}
\codebox{uint tgt = laneOp(lane\_id, src1 \& 0xFFu, subgroup\_size);\\dest = (laneActive(tgt) ? readLane(src0, tgt) : inactiveConst(inactive\_result)) \textbar{} src2;}
\textbf{Execution Context:}
\begin{tightitemize}\footnotesize
  \item \textbf{Compute stage preloads (\hyperref[sec:csf-spd]{\texttt{SPD\_k}}, \texttt{r16..r23}):} Hardware pre-fills \texttt{r55..r57} with \texttt{WorkGroupID.xyz + JOB\_OFFSET} (from \hyperref[sec:csf-regs]{\texttt{COMPUTESR::JOB\_OFFSET\_X/Y/Z}} \texttt{r34..r36}) and \texttt{r60..r62} with \texttt{LocalInvocationID.xyz} or \texttt{GlobalInvocationID.xyz} (from \texttt{WG\_SIZE} \texttt{r33} \& \texttt{JOB\_SIZE} \texttt{r37..r39}, selected by \texttt{ShaderProgram::preload}) so compute shaders can derive thread/element indices without memory loads.
  \item \textbf{Vertex stage preloads (\hyperref[sec:csf-spd]{\texttt{SPD\_0}/\texttt{SPD\_1}}, \texttt{r16..r19}):} The hardware Index Fetcher reads \hyperref[sec:csf-regs]{\texttt{IDVSSR::INDEX\_BUFFER}} (\texttt{r54:r55}) and pre-fills \texttt{r59, r60} with \texttt{VertexIndex} (\texttt{+ VERTEX\_OFFSET + GLOBAL\_ATTR\_OFF}) and \texttt{r61} with \texttt{InstanceIndex} (\texttt{+ INSTANCE\_OFFSET}) so Position and Varying shaders can pull vertex attributes directly.
  \item \textbf{Inline FAU \& Special Operands (\hyperref[sec:csf-tsd-fau]{\texttt{FAU\_k}}, \texttt{r8..r15}):} Any 8-bit source with \texttt{mode=2} reads zero-latency uniform \texttt{u0..u31} from \texttt{FAU\_k} (\texttt{slots64 = bits[63:56]}); \texttt{mode=3} reads \texttt{SmallConstant} (Table~1.2) or \texttt{FauSpecial} (\texttt{pc}, \texttt{lane\_id}, \texttt{lane\_width}, \texttt{tls\_ptr}, \texttt{wls\_ptr}, \texttt{blend\_descriptor\_0..7}).
\end{tightitemize}
\par\textbf{Worked Examples:}
\par\textbf{Example 1 (Fused subgroup permute + bitwise OR (v14+)):}
\codebox{\mnem{CLPER\_OR}\imod{.subgroup16.lanes04} \dst{r0}, \src{r1}, \src{r2}\imod{.b0}, \src{r3}}
Permutes \texttt{r1} from the lane pair selected by \texttt{lanes04} + \texttt{r2.b0}, then computes \texttt{r0 = permuted \textbar{} r3} in one instruction.
\par\textbf{Example 2 (Quad bitmask accumulation):}
\codebox{\mnem{CLPER\_OR}\imod{.subgroup4.lanes15} \dst{r0}, \src{r1}, \src{r2}\imod{.b0}, \src{r3}}
Accumulates cross-lane bitflags across a 4-lane quad without needing a separate \texttt{OR.i32} instruction.

\begin{center}\footnotesize
\begin{tabularx}{\linewidth}{@{}lllX@{}}
\toprule
Variant & Arch & Class & Opcode Fields \\
\midrule
\texttt{(none)} & v9--v14 & \texttt{R3} & \texttt{opcode1=CLPER\_OR} \\
\bottomrule
\end{tabularx}
\end{center}
\begin{center}\footnotesize
\begin{tabularx}{0.88\linewidth}{@{}l l X@{}}
\toprule
Field Name & Bits & Type / Fixed Value \\
\midrule
\texttt{src0} & \texttt{0:7} & \texttt{SourceEncoding} \\
\texttt{src1} & \texttt{8:15} & \texttt{SourceEncoding} \\
\texttt{src2} & \texttt{16:23} & \texttt{SourceEncoding} \\
\textit{Reserved} & \texttt{24:25} & \texttt{zero} \\
\hyperref[enum:shift_bitwise_i32_lanes_m]{\texttt{lanes1}} & \texttt{26:29} & \texttt{shift\_bitwise\_i32\_lanes\_m (def b0)} \\
\textit{Reserved} & \texttt{30:32} & \texttt{zero} \\
\hyperref[enum:clper_or_lane_select_m]{\texttt{lane\_select}} & \texttt{33:35} & \texttt{clper\_or\_lane\_select\_m} \\
\hyperref[enum:clper_subgroup_size_m]{\texttt{subgroup}} & \texttt{36:39} & \texttt{clper\_subgroup\_size\_m} \\
\texttt{dest\_register\_index} & \texttt{40:45} & \texttt{uint} \\
\hyperref[enum:dest_width_m]{\texttt{dest\_width}} & \texttt{46:47} & \texttt{dest\_width\_m} \\
\texttt{opcode1} & \texttt{48:51} & \texttt{exact=CLPER\_OR} \\
\texttt{instruction\_class} & \texttt{52:56} & \texttt{exact=R3} \\
\texttt{fau\_page\_index} & \texttt{57:58} & \texttt{uint} \\
\hyperref[enum:flow_control_m]{\texttt{flow}} & \texttt{59:62} & \texttt{flow\_control\_m} \\
\textit{Reserved} & \texttt{63} & \texttt{zero} \\
\bottomrule
\end{tabularx}
\end{center}
\textbf{Fields Explained:}
\begin{tightitemize}\footnotesize
  \item \texttt{src0}: Primary source operand (\texttt{r0..r63}, \texttt{u0..u63} uniform, \texttt{k*} FAU constant, or \texttt{t0..t2} temporary; \texttt{\textasciicircum{}} marks discard/last-use). \emph{Example:} {\scriptsize\ttfamily \mnem{FADD.f32} \dst{r0}, \src{r1}, \src{r2}} (\textbf{SPIR-V:} \texttt{OpFAdd (Operand 0)}; \textbf{GLSL:} \texttt{float c = a + b;})
  \item \texttt{src1}: Secondary source operand (register, uniform, or FAU constant; in shifts, supplies the shift count byte). \emph{Example:} {\scriptsize\ttfamily \mnem{LSHIFT.i32} \dst{r0}, \src{r1}, \src{r2}\imod{.b0}} (\textbf{SPIR-V:} \texttt{OpShiftLeftLogical (Shift)}; \textbf{GLSL:} \texttt{uint y = a \textless{}\textless{} s;})
  \item \texttt{src2}: Third source operand (FMA addend, \texttt{MUX}/\texttt{CSEL} selector, \texttt{IADDC} carry-in, or bitwise mask in fused shift-logic). \emph{Example:} {\scriptsize\ttfamily \mnem{FMA.f32} \dst{r0}, \src{r1}, \src{r2}, \src{r3}} (\textbf{SPIR-V:} \texttt{GLSL.std.450 Fma (Addend)}; \textbf{GLSL:} \texttt{float d = fma(a, b, c);})
  \item \hyperref[enum:shift_bitwise_i32_lanes_m]{\texttt{lanes1}}: Stage A byte-lane selector on \texttt{src1} (selects shift-count byte \texttt{b0..b3} on 32/64-bit shifts, \texttt{bXY} on \texttt{v2i16} shifts, or 4-byte swizzle on \texttt{v4*8}). \emph{Example:} {\scriptsize\ttfamily \mnem{LSHIFT.i32} \dst{r0}, \src{r1}, \src{r2}\imod{.b1}} (\textbf{SPIR-V:} \texttt{OpBitFieldUExtract + OpShiftLeftLogical}; \textbf{GLSL:} \texttt{uint y = a \textless{}\textless{} ((s \textgreater{}\textgreater{} 8) \& 0xFFu);})
  \item \hyperref[enum:clper_or_lane_select_m]{\texttt{lane\_select}}: Selects which pair of warp lanes (\texttt{lanes04}..\texttt{lanesBF}) is OR-combined by \texttt{CLPER\_OR}. \emph{Example:} {\scriptsize\ttfamily \mnem{CLPER\_OR}\imod{.subgroup16.lanes04} \dst{r0}, \src{r1}, \src{r2}} (\textbf{SPIR-V:} \texttt{OpGroupNonUniformBitwiseOr}; \textbf{GLSL:} \texttt{uint v = subgroupOr(x);})
  \item \hyperref[enum:clper_subgroup_size_m]{\texttt{subgroup}}: Subgroup partition width (\texttt{subgroup2}, \texttt{subgroup4}, \texttt{subgroup8}, \texttt{subgroup16}). \emph{Example:} {\scriptsize\ttfamily \mnem{CLPER}\imod{.subgroup8.rotate.zero} \dst{r0}, \src{r1}, \src{r2}} (\textbf{SPIR-V:} \texttt{OpGroupNonUniformRotateKHR}; \textbf{GLSL:} \texttt{uint v = subgroupRotate(x, 1);})
  \item \texttt{dest\_register\_index}: 6-bit destination GPR index (\texttt{r0..r63}) receiving the Stage C writeback result (even-aligned pair \texttt{[rN:rN+1]} for 64-bit ops). \emph{Example:} {\scriptsize\ttfamily \mnem{FMA.f32} \dst{r0}, \src{r1}, \src{r2}, \src{r3}} (\textbf{SPIR-V:} \texttt{SSA Result ID}; \textbf{GLSL:} \texttt{float r0 = fma(r1, r2, r3);})
  \item \hyperref[enum:dest_width_m]{\texttt{dest\_width}}: Stage C destination lane write-mask (\texttt{none} = full 32-bit, \texttt{h0} = low 16-bit, \texttt{h1} = high 16-bit without bit-shifting). \emph{Example:} {\scriptsize\ttfamily \mnem{FADD.v2f16} \dst{r0}\imod{.h1}, \src{r1}, \src{r2}} (\textbf{SPIR-V:} \texttt{OpVectorInsertDynamic / WriteMask}; \textbf{GLSL:} \texttt{dst.y = a.y + b.y;})
  \item \texttt{opcode1}: Primary sub-opcode within the instruction class uniquely identifying the operation family. \emph{Example:} {\scriptsize\ttfamily \mnem{IADD.s32} \dst{r0}, \src{r1}, \src{r2}} (\textbf{SPIR-V:} \texttt{OpIAdd (operation selector)}; \textbf{GLSL:} \texttt{// Primary opcode bits})
  \item \texttt{instruction\_class}: 4-bit major instruction class selector (\texttt{bits 48:51}) partitioning the 64-bit Valhall opcode space (e.g., \texttt{FMA}, \texttt{CVT}, \texttt{SFU}, \texttt{LS}, \texttt{TEX}, \texttt{BRANCH}). \emph{Example:} {\scriptsize\ttfamily \mnem{FADD.f32} \dst{r0}, \src{r1}, \src{r2}} (\textbf{SPIR-V:} \texttt{N/A (encoding class)}; \textbf{GLSL:} \texttt{// Hardware execution pipe / decode class})
  \item \texttt{fau\_page\_index}: Selects the FAU (Fast Access Uniform) page (\texttt{0..3}) when a source operand references hardware small constants, thread IDs, or lane constants. \emph{Example:} {\scriptsize\ttfamily \mnem{FADD.f32} \dst{r0}, \src{r1}, \src{k1f}} (\textbf{SPIR-V:} \texttt{OpConstant / BuiltIn}; \textbf{GLSL:} \texttt{float y = x + 1.0;})
  \item \hyperref[enum:flow_control_m]{\texttt{flow}}: 4-bit embedded flow-control \& scoreboard wait modifier (\texttt{bits 59:62}: \texttt{none}, \texttt{wait0..wait0126}, \texttt{wait}, \texttt{reconverge}, \texttt{discard}, \texttt{end}). \emph{Example:} {\scriptsize\ttfamily \mnem{FADD.f32}\imod{.wait0.reconverge} \dst{r0}, \src{r1}, \src{r2}} (\textbf{SPIR-V:} \texttt{OpControlBarrier / OpReturn}; \textbf{GLSL:} \texttt{// Scoreboard wait + warp reconvergence})
\end{tightitemize}
\clearpage

\subsection*{WMASK}
\addcontentsline{toc}{subsection}{WMASK}
Computes a subgroup ballot bitmask over active lanes within the hardware subgroup (\texttt{subgroup2}, \texttt{subgroup4}, \texttt{subgroup8}, or \texttt{subgroup16}): bit $k$ of \texttt{dest} is \texttt{1} if lane $k$ is active and its \texttt{src0 != 0}, and \texttt{0} otherwise.
\par\textbf{Notes:}
\begin{tightitemize}
  \item Cross-lane subgroup operation.
  \item NOT an identity move! Implements \texttt{nir\_intrinsic\_ballot}.
  \item \textbf{Source references:} \texttt{kraid/ops.rs:5087-5103}, \texttt{encode\_v9.rs:3831-3848}, \texttt{nir.rs:1885-1893}, \texttt{bifrost/valhall/ISA.xml:2172-2188}.
\end{tightitemize}
\codebox{\mnem{WMASK}\imod{\{.flow\}}\imod{\{.subgroup\}} \dst{dest}\imod{\{.dest\_width\}}, \src{src0}}
\textbf{Operands:}
\begin{tightitemize}\footnotesize
  \item \texttt{\dst{dest}} (\texttt{dest\_register\_index}, \texttt{dest\_width}): 32-bit destination GPR receiving the active-lane ballot bitmask within \texttt{subgroup}.
  \item \texttt{\src{src0}}: Per-lane predicate (non-zero sets this lane's bit in the ballot mask).
\end{tightitemize}
\par\textbf{Closest SPIR-V Op:} \texttt{OpGroupNonUniformBallot / All / Any}
\par\textbf{Pseudo-GLSL (Operational Semantics):}
\codebox{uint mask = 0u;\\for (uint k = 0u; k \textless{} subgroup\_size; ++k) if (laneActive(k) \&\& readLane(src0, k) != 0u) mask \textbar{}= (1u \textless{}\textless{} k);\\dest = mask;}
\textbf{Execution Context:}
\begin{tightitemize}\footnotesize
  \item \textbf{Compute stage preloads (\hyperref[sec:csf-spd]{\texttt{SPD\_k}}, \texttt{r16..r23}):} Hardware pre-fills \texttt{r55..r57} with \texttt{WorkGroupID.xyz + JOB\_OFFSET} (from \hyperref[sec:csf-regs]{\texttt{COMPUTESR::JOB\_OFFSET\_X/Y/Z}} \texttt{r34..r36}) and \texttt{r60..r62} with \texttt{LocalInvocationID.xyz} or \texttt{GlobalInvocationID.xyz} (from \texttt{WG\_SIZE} \texttt{r33} \& \texttt{JOB\_SIZE} \texttt{r37..r39}, selected by \texttt{ShaderProgram::preload}) so compute shaders can derive thread/element indices without memory loads.
  \item \textbf{Vertex stage preloads (\hyperref[sec:csf-spd]{\texttt{SPD\_0}/\texttt{SPD\_1}}, \texttt{r16..r19}):} The hardware Index Fetcher reads \hyperref[sec:csf-regs]{\texttt{IDVSSR::INDEX\_BUFFER}} (\texttt{r54:r55}) and pre-fills \texttt{r59, r60} with \texttt{VertexIndex} (\texttt{+ VERTEX\_OFFSET + GLOBAL\_ATTR\_OFF}) and \texttt{r61} with \texttt{InstanceIndex} (\texttt{+ INSTANCE\_OFFSET}) so Position and Varying shaders can pull vertex attributes directly.
  \item \textbf{Inline FAU \& Special Operands (\hyperref[sec:csf-tsd-fau]{\texttt{FAU\_k}}, \texttt{r8..r15}):} Any 8-bit source with \texttt{mode=2} reads zero-latency uniform \texttt{u0..u31} from \texttt{FAU\_k} (\texttt{slots64 = bits[63:56]}); \texttt{mode=3} reads \texttt{SmallConstant} (Table~1.2) or \texttt{FauSpecial} (\texttt{pc}, \texttt{lane\_id}, \texttt{lane\_width}, \texttt{tls\_ptr}, \texttt{wls\_ptr}, \texttt{blend\_descriptor\_0..7}).
\end{tightitemize}
\par\textbf{Worked Examples:}
\par\textbf{Example 1 (16-lane warp ballot):}
\codebox{\mnem{WMASK}\imod{.subgroup16} \dst{r0}, \src{r1}}
If active lanes \texttt{0, 1, 4} have \texttt{r1 != 0} and all other lanes have \texttt{r1 == 0}: returns \texttt{r0 = 0x00000013} (\texttt{(1\textless{}\textless{}0)\textbar{}(1\textless{}\textless{}1)\textbar{}(1\textless{}\textless{}4)}).
\par\textbf{Example 2 (4-lane quad ballot):}
\codebox{\mnem{WMASK}\imod{.subgroup4} \dst{r0}, \src{r1}}
Evaluates non-zero predicate independently within each 4-lane fragment quad (bits \texttt{0..3} per quad).

\begin{center}\footnotesize
\begin{tabularx}{\linewidth}{@{}lllX@{}}
\toprule
Variant & Arch & Class & Opcode Fields \\
\midrule
\texttt{(none)} & v9--v14 & \texttt{R1} & \texttt{opcode2=WMASK, opcode1=WMASK} \\
\bottomrule
\end{tabularx}
\end{center}
\begin{center}\footnotesize
\begin{tabularx}{0.88\linewidth}{@{}l l X@{}}
\toprule
Field Name & Bits & Type / Fixed Value \\
\midrule
\texttt{src0} & \texttt{0:7} & \texttt{SourceEncoding} \\
\textit{Reserved} & \texttt{8:15} & \texttt{zero} \\
\texttt{opcode2} & \texttt{16:20} & \texttt{exact=WMASK} \\
\textit{Reserved} & \texttt{21:35} & \texttt{zero} \\
\hyperref[enum:clper_subgroup_size_m]{\texttt{subgroup}} & \texttt{36:39} & \texttt{clper\_subgroup\_size\_m} \\
\texttt{dest\_register\_index} & \texttt{40:45} & \texttt{uint} \\
\hyperref[enum:dest_width_m]{\texttt{dest\_width}} & \texttt{46:47} & \texttt{dest\_width\_m} \\
\texttt{opcode1} & \texttt{48:51} & \texttt{exact=WMASK} \\
\texttt{instruction\_class} & \texttt{52:56} & \texttt{exact=R1} \\
\texttt{fau\_page\_index} & \texttt{57:58} & \texttt{uint} \\
\hyperref[enum:flow_control_m]{\texttt{flow}} & \texttt{59:62} & \texttt{flow\_control\_m} \\
\textit{Reserved} & \texttt{63} & \texttt{zero} \\
\bottomrule
\end{tabularx}
\end{center}
\textbf{Fields Explained:}
\begin{tightitemize}\footnotesize
  \item \texttt{src0}: Primary source operand (\texttt{r0..r63}, \texttt{u0..u63} uniform, \texttt{k*} FAU constant, or \texttt{t0..t2} temporary; \texttt{\textasciicircum{}} marks discard/last-use). \emph{Example:} {\scriptsize\ttfamily \mnem{FADD.f32} \dst{r0}, \src{r1}, \src{r2}} (\textbf{SPIR-V:} \texttt{OpFAdd (Operand 0)}; \textbf{GLSL:} \texttt{float c = a + b;})
  \item \texttt{opcode2}: Secondary sub-opcode further specializing instructions sharing \texttt{instruction\_class} and \texttt{opcode1} (e.g., conversions, SFU tables, moves). \emph{Example:} {\scriptsize\ttfamily \mnem{F32\_TO\_S32}\imod{.round\_zero} \dst{r0}, \src{r1}} (\textbf{SPIR-V:} \texttt{OpConvertFToS}; \textbf{GLSL:} \texttt{int y = int(x);})
  \item \hyperref[enum:clper_subgroup_size_m]{\texttt{subgroup}}: Subgroup partition width (\texttt{subgroup2}, \texttt{subgroup4}, \texttt{subgroup8}, \texttt{subgroup16}). \emph{Example:} {\scriptsize\ttfamily \mnem{CLPER}\imod{.subgroup8.rotate.zero} \dst{r0}, \src{r1}, \src{r2}} (\textbf{SPIR-V:} \texttt{OpGroupNonUniformRotateKHR}; \textbf{GLSL:} \texttt{uint v = subgroupRotate(x, 1);})
  \item \texttt{dest\_register\_index}: 6-bit destination GPR index (\texttt{r0..r63}) receiving the Stage C writeback result (even-aligned pair \texttt{[rN:rN+1]} for 64-bit ops). \emph{Example:} {\scriptsize\ttfamily \mnem{FMA.f32} \dst{r0}, \src{r1}, \src{r2}, \src{r3}} (\textbf{SPIR-V:} \texttt{SSA Result ID}; \textbf{GLSL:} \texttt{float r0 = fma(r1, r2, r3);})
  \item \hyperref[enum:dest_width_m]{\texttt{dest\_width}}: Stage C destination lane write-mask (\texttt{none} = full 32-bit, \texttt{h0} = low 16-bit, \texttt{h1} = high 16-bit without bit-shifting). \emph{Example:} {\scriptsize\ttfamily \mnem{FADD.v2f16} \dst{r0}\imod{.h1}, \src{r1}, \src{r2}} (\textbf{SPIR-V:} \texttt{OpVectorInsertDynamic / WriteMask}; \textbf{GLSL:} \texttt{dst.y = a.y + b.y;})
  \item \texttt{opcode1}: Primary sub-opcode within the instruction class uniquely identifying the operation family. \emph{Example:} {\scriptsize\ttfamily \mnem{IADD.s32} \dst{r0}, \src{r1}, \src{r2}} (\textbf{SPIR-V:} \texttt{OpIAdd (operation selector)}; \textbf{GLSL:} \texttt{// Primary opcode bits})
  \item \texttt{instruction\_class}: 4-bit major instruction class selector (\texttt{bits 48:51}) partitioning the 64-bit Valhall opcode space (e.g., \texttt{FMA}, \texttt{CVT}, \texttt{SFU}, \texttt{LS}, \texttt{TEX}, \texttt{BRANCH}). \emph{Example:} {\scriptsize\ttfamily \mnem{FADD.f32} \dst{r0}, \src{r1}, \src{r2}} (\textbf{SPIR-V:} \texttt{N/A (encoding class)}; \textbf{GLSL:} \texttt{// Hardware execution pipe / decode class})
  \item \texttt{fau\_page\_index}: Selects the FAU (Fast Access Uniform) page (\texttt{0..3}) when a source operand references hardware small constants, thread IDs, or lane constants. \emph{Example:} {\scriptsize\ttfamily \mnem{FADD.f32} \dst{r0}, \src{r1}, \src{k1f}} (\textbf{SPIR-V:} \texttt{OpConstant / BuiltIn}; \textbf{GLSL:} \texttt{float y = x + 1.0;})
  \item \hyperref[enum:flow_control_m]{\texttt{flow}}: 4-bit embedded flow-control \& scoreboard wait modifier (\texttt{bits 59:62}: \texttt{none}, \texttt{wait0..wait0126}, \texttt{wait}, \texttt{reconverge}, \texttt{discard}, \texttt{end}). \emph{Example:} {\scriptsize\ttfamily \mnem{FADD.f32}\imod{.wait0.reconverge} \dst{r0}, \src{r1}, \src{r2}} (\textbf{SPIR-V:} \texttt{OpControlBarrier / OpReturn}; \textbf{GLSL:} \texttt{// Scoreboard wait + warp reconvergence})
\end{tightitemize}
\clearpage

\section{Control Flow, Barriers \& System (9 mnemonics)}
\subsection*{BARRIER}
\addcontentsline{toc}{subsection}{BARRIER}
Synchronizes execution and workgroup-shared (WLS) memory visibility across all warps in the current compute workgroup; advances the hardware barrier epoch before any warp proceeds past the barrier.
\par\textbf{Notes:}
\begin{tightitemize}
  \item Implements \texttt{nir\_intrinsic\_barrier} with \texttt{SCOPE\_WORKGROUP} (\texttt{kraid/nir.rs:1898-1904}).
  \item \textbf{Source references:} \texttt{kraid/ops.rs:317-327}, \texttt{encode\_v9.rs:738-746}, \texttt{nir.rs:1894-1906}, \texttt{bifrost/valhall/ISA.xml:13}.
\end{tightitemize}
\codebox{\mnem{BARRIER}\imod{\{.message\_slot\_index\}}\imod{\{.flow\}}}
\textbf{Operands:} \textit{None} (control/system instruction with no register or immediate operands; configured via modifiers and \texttt{flow}).
\par\textbf{Closest SPIR-V Op:} \texttt{OpControlBarrier}
\par\textbf{Pseudo-GLSL (Operational Semantics):}
\codebox{barrier();\\memoryBarrierShared(); // synchronize all warps in WG\_SIZE workgroup}
\textbf{Execution Context:}
\begin{tightitemize}\footnotesize
  \item \textbf{Workgroup synchronization (\hyperref[sec:csf-regs]{\texttt{COMPUTESR::WG\_SIZE}}, \texttt{r33}):} Synchronizes all warps in the \texttt{(x, y, z)} workgroup and orders \hyperref[sec:csf-tsd-fau]{\texttt{TSD\_k}} WLS shared-memory visibility. Must encode \texttt{slot7} + \texttt{.wait}.
  \item \textbf{Compute stage preloads (\hyperref[sec:csf-spd]{\texttt{SPD\_k}}, \texttt{r16..r23}):} Hardware pre-fills \texttt{r55..r57} with \texttt{WorkGroupID.xyz + JOB\_OFFSET} (from \hyperref[sec:csf-regs]{\texttt{COMPUTESR::JOB\_OFFSET\_X/Y/Z}} \texttt{r34..r36}) and \texttt{r60..r62} with \texttt{LocalInvocationID.xyz} or \texttt{GlobalInvocationID.xyz} (from \texttt{WG\_SIZE} \texttt{r33} \& \texttt{JOB\_SIZE} \texttt{r37..r39}, selected by \texttt{ShaderProgram::preload}) so compute shaders can derive thread/element indices without memory loads.
\end{tightitemize}
\par\textbf{Worked Examples:}
\par\textbf{Example 1 (Compute workgroup barrier):}
\codebox{\mnem{BARRIER}\imod{.slot7.wait}}
Synchronizes all warps in the compute workgroup and increments the barrier epoch (requires \texttt{slot7} + \texttt{.wait}).
\par\textbf{Example 2 (Shared-memory publish + barrier):}
\codebox{\mnem{BARRIER}\imod{.slot7.wait}}
Placed after \texttt{STORE.i32.slot0} to WLS/shared memory so all warps observe prior writes before subsequent \texttt{LOAD}s.

\begin{center}\footnotesize
\begin{tabularx}{\linewidth}{@{}lllX@{}}
\toprule
Variant & Arch & Class & Opcode Fields \\
\midrule
\texttt{(none)} & v9--v14 & \texttt{S0} & \texttt{opcode1=BARRIER} \\
\bottomrule
\end{tabularx}
\end{center}
\begin{center}\footnotesize
\begin{tabularx}{0.88\linewidth}{@{}l l X@{}}
\toprule
Field Name & Bits & Type / Fixed Value \\
\midrule
\textit{Reserved} & \texttt{0:29} & \texttt{zero} \\
\hyperref[enum:message_slot_index_m]{\texttt{message\_slot\_index}} & \texttt{30:32} & \texttt{exact=slot7} \\
\texttt{sr\_count} & \texttt{33:35} & \texttt{exact=count0} \\
\textit{Reserved} & \texttt{36:39} & \texttt{zero} \\
\texttt{staging\_register\_index} & \texttt{40:45} & \texttt{exact=0} \\
\hyperref[enum:dest_width_m]{\texttt{dest\_width}} & \texttt{46:47} & \texttt{exact=none} \\
\texttt{opcode1} & \texttt{48:51} & \texttt{exact=BARRIER} \\
\texttt{instruction\_class} & \texttt{52:56} & \texttt{exact=S0} \\
\texttt{fau\_page\_index} & \texttt{57:58} & \texttt{exact=0} \\
\hyperref[enum:flow_control_m]{\texttt{flow}} & \texttt{59:62} & \texttt{flow\_control\_m} \\
\textit{Reserved} & \texttt{63} & \texttt{zero} \\
\bottomrule
\end{tabularx}
\end{center}
\textbf{Fields Explained:}
\begin{tightitemize}\footnotesize
  \item \hyperref[enum:message_slot_index_m]{\texttt{message\_slot\_index}}: Asynchronous scoreboard slot (\texttt{slot0}--\texttt{slot2}, \texttt{slot7}) tracking completion of this message. \emph{Example:} {\scriptsize\ttfamily \mnem{LOAD.i32}\imod{.slot0} \dst{@r0}, \src{r2}, \imm{\#0}} (\textbf{SPIR-V:} \texttt{OpLoad}; \textbf{GLSL:} \texttt{uint v = ssbo.arr[i];})
  \item \texttt{sr\_count}: Number of consecutive 32-bit staging registers transferred by a load/store/texture message. \emph{Example:} {\scriptsize\ttfamily \mnem{LD\_MULTI.i128} \dst{@r0:r3}, \src{r4}, \imm{\#0}} (\textbf{SPIR-V:} \texttt{OpLoad (wide vector)}; \textbf{GLSL:} \texttt{uvec4 x = buf[i];})
  \item \texttt{staging\_register\_index}: Base GPR index of the contiguous staging register window (\texttt{@rN:rM}) for message payload/return. \emph{Example:} {\scriptsize\ttfamily \mnem{STORE.i64}\imod{.slot7} \src{@r0:r1}, \src{r2}, \imm{\#8}} (\textbf{SPIR-V:} \texttt{OpStore (uvec2)}; \textbf{GLSL:} \texttt{ssbo.pairs[i] = uvec2(a, b);})
  \item \hyperref[enum:dest_width_m]{\texttt{dest\_width}}: Stage C destination lane write-mask (\texttt{none} = full 32-bit, \texttt{h0} = low 16-bit, \texttt{h1} = high 16-bit without bit-shifting). \emph{Example:} {\scriptsize\ttfamily \mnem{FADD.v2f16} \dst{r0}\imod{.h1}, \src{r1}, \src{r2}} (\textbf{SPIR-V:} \texttt{OpVectorInsertDynamic / WriteMask}; \textbf{GLSL:} \texttt{dst.y = a.y + b.y;})
  \item \texttt{opcode1}: Primary sub-opcode within the instruction class uniquely identifying the operation family. \emph{Example:} {\scriptsize\ttfamily \mnem{IADD.s32} \dst{r0}, \src{r1}, \src{r2}} (\textbf{SPIR-V:} \texttt{OpIAdd (operation selector)}; \textbf{GLSL:} \texttt{// Primary opcode bits})
  \item \texttt{instruction\_class}: 4-bit major instruction class selector (\texttt{bits 48:51}) partitioning the 64-bit Valhall opcode space (e.g., \texttt{FMA}, \texttt{CVT}, \texttt{SFU}, \texttt{LS}, \texttt{TEX}, \texttt{BRANCH}). \emph{Example:} {\scriptsize\ttfamily \mnem{FADD.f32} \dst{r0}, \src{r1}, \src{r2}} (\textbf{SPIR-V:} \texttt{N/A (encoding class)}; \textbf{GLSL:} \texttt{// Hardware execution pipe / decode class})
  \item \texttt{fau\_page\_index}: Selects the FAU (Fast Access Uniform) page (\texttt{0..3}) when a source operand references hardware small constants, thread IDs, or lane constants. \emph{Example:} {\scriptsize\ttfamily \mnem{FADD.f32} \dst{r0}, \src{r1}, \src{k1f}} (\textbf{SPIR-V:} \texttt{OpConstant / BuiltIn}; \textbf{GLSL:} \texttt{float y = x + 1.0;})
  \item \hyperref[enum:flow_control_m]{\texttt{flow}}: 4-bit embedded flow-control \& scoreboard wait modifier (\texttt{bits 59:62}: \texttt{none}, \texttt{wait0..wait0126}, \texttt{wait}, \texttt{reconverge}, \texttt{discard}, \texttt{end}). \emph{Example:} {\scriptsize\ttfamily \mnem{FADD.f32}\imod{.wait0.reconverge} \dst{r0}, \src{r1}, \src{r2}} (\textbf{SPIR-V:} \texttt{OpControlBarrier / OpReturn}; \textbf{GLSL:} \texttt{// Scoreboard wait + warp reconvergence})
\end{tightitemize}
\clearpage

\subsection*{BRANCH}
\addcontentsline{toc}{subsection}{BRANCH}
Evaluates \texttt{base = combine(src0)} (\texttt{none}: \texttt{!=0}, \texttt{h0}: \texttt{lo16!=0}, \texttt{h1}: \texttt{hi16!=0}, \texttt{and}: \texttt{==0xFFFFFFFF}, \texttt{lowbits}: \texttt{bits[20:3] not in \{0, 0x3FFFF\}}); \texttt{taken = not ? !base : base}; \texttt{pc = taken ? pc + 8 + offset : pc + 8}.
\par\textbf{Notes:}
\begin{tightitemize}
  \item Unconditional branch is encoded as \texttt{BRANCH.not 0, target} (\texttt{kraid/ops.rs:513-515}).
  \item \texttt{not} is a boolean branch-condition negation, NOT a bitwise NOT on \texttt{src0}!
  \item \textbf{Source references:} \texttt{kraid/ops.rs:476-533}, \texttt{encode\_v9.rs:772-792}, \texttt{opt\_fuse.rs:240-290}, \texttt{bifrost/valhall/ISA.xml:19}.
\end{tightitemize}
\codebox{\mnem{BRANCH}\imod{\{.flow\}}\imod{\{.conservative\}}\imod{\{.not\}} \src{src0}\imod{\{.combine\}}, \imm{\#offset}}
\textbf{Operands:}
\begin{tightitemize}\footnotesize
  \item \texttt{\src{src0}}: Condition register tested via \texttt{combine} (and inverted if \texttt{.not}).
  \item \texttt{\imm{\#offset}} (\texttt{offset}): Signed PC-relative byte offset added to \texttt{PC + 8} when the branch is taken.
\end{tightitemize}
\par\textbf{Closest SPIR-V Op:} \texttt{OpBranch / OpBranchConditional}
\par\textbf{Pseudo-GLSL (Operational Semantics):}
\codebox{bool cond = combine(src0); // none: !=0, h0/h1, and: both!=0, lowbits\\if (flag\_not ? !cond : cond) PC = PC + 8 + \#offset;}
\textbf{Execution Context:}
\begin{tightitemize}\footnotesize
  \item \textbf{Compute stage preloads (\hyperref[sec:csf-spd]{\texttt{SPD\_k}}, \texttt{r16..r23}):} Hardware pre-fills \texttt{r55..r57} with \texttt{WorkGroupID.xyz + JOB\_OFFSET} (from \hyperref[sec:csf-regs]{\texttt{COMPUTESR::JOB\_OFFSET\_X/Y/Z}} \texttt{r34..r36}) and \texttt{r60..r62} with \texttt{LocalInvocationID.xyz} or \texttt{GlobalInvocationID.xyz} (from \texttt{WG\_SIZE} \texttt{r33} \& \texttt{JOB\_SIZE} \texttt{r37..r39}, selected by \texttt{ShaderProgram::preload}) so compute shaders can derive thread/element indices without memory loads.
  \item \textbf{Vertex stage preloads (\hyperref[sec:csf-spd]{\texttt{SPD\_0}/\texttt{SPD\_1}}, \texttt{r16..r19}):} The hardware Index Fetcher reads \hyperref[sec:csf-regs]{\texttt{IDVSSR::INDEX\_BUFFER}} (\texttt{r54:r55}) and pre-fills \texttt{r59, r60} with \texttt{VertexIndex} (\texttt{+ VERTEX\_OFFSET + GLOBAL\_ATTR\_OFF}) and \texttt{r61} with \texttt{InstanceIndex} (\texttt{+ INSTANCE\_OFFSET}) so Position and Varying shaders can pull vertex attributes directly.
  \item \textbf{Inline FAU \& Special Operands (\hyperref[sec:csf-tsd-fau]{\texttt{FAU\_k}}, \texttt{r8..r15}):} Any 8-bit source with \texttt{mode=2} reads zero-latency uniform \texttt{u0..u31} from \texttt{FAU\_k} (\texttt{slots64 = bits[63:56]}); \texttt{mode=3} reads \texttt{SmallConstant} (Table~1.2) or \texttt{FauSpecial} (\texttt{pc}, \texttt{lane\_id}, \texttt{lane\_width}, \texttt{tls\_ptr}, \texttt{wls\_ptr}, \texttt{blend\_descriptor\_0..7}).
\end{tightitemize}
\par\textbf{Worked Examples:}
\par\textbf{Example 1 (Panfrost canonical unconditional branch):}
\codebox{\mnem{BRANCH}\imod{.not} \src{k0}, \imm{\#0x10}}
\texttt{combine=none} tests \texttt{k0 != 0} (\texttt{False}); \texttt{.not} negates boolean to \texttt{True} -\textgreater{} \texttt{PC = PC + 8 + 0x10}.
\par\textbf{Example 2 (Conditional halfword-AND branch):}
\codebox{\mnem{BRANCH} \src{r4}\imod{.and}, \imm{\#0x20}}
With \texttt{r4 = 0x00010001}: both halfwords non-zero -\textgreater{} taken (\texttt{PC += 8 + 0x20}). With \texttt{r4 = 0x00010000}: falls through (\texttt{PC += 8}).
\par\textbf{Mined SPIR-V $\to$ Valhall ISA Example (Conditional forward skip after fragment discard check --- Fragment shader, Mali-G710 (v10)):}
\codebox{\textbf{; SPIR-V (input)}\\\%le = OpFOrdLessThanEqual \%bool \%alpha \%thresh\\OpSelectionMerge \%merge None\\OpBranchConditional \%le \%merge \%shaded\\[3pt]\textbf{; Valhall ISA (malioc output)}\\FCMP.f32.discard.m1.le r8, r7, u0.w0\\MOV\_IMM.i32 r0, \#0x0\\BRANCH r8\textasciicircum{}, 0x88                  ; skip shading block if r8 != 0\\... ; shaded block ending with .reconverge}
\texttt{BRANCH} tests \texttt{r8\textasciicircum{} != 0} and jumps forward by \texttt{+0x88} bytes when true, reconverging at the merge block (marked with \texttt{.reconverge} on the last instruction of the taken/fallthrough body).
\begin{center}\footnotesize
\begin{tabularx}{\linewidth}{@{}lllX@{}}
\toprule
Variant & Arch & Class & Opcode Fields \\
\midrule
\texttt{(none)} & v9--v14 & \texttt{N1} & \texttt{opcode1=BRANCH} \\
\bottomrule
\end{tabularx}
\end{center}
\begin{center}\footnotesize
\begin{tabularx}{0.88\linewidth}{@{}l l X@{}}
\toprule
Field Name & Bits & Type / Fixed Value \\
\midrule
\texttt{src0} & \texttt{0:7} & \texttt{SourceEncoding} \\
\texttt{offset} & \texttt{8:34} & \texttt{pc\_rel\_label\_signed} \\
\hyperref[enum:branch_conservative_m]{\texttt{conservative}} & \texttt{35} & \texttt{branch\_conservative\_m (def none)} \\
\hyperref[enum:branch_negate_m]{\texttt{not}} & \texttt{36} & \texttt{branch\_negate\_m (def none)} \\
\hyperref[enum:branch_combine_m]{\texttt{combine}} & \texttt{37:39} & \texttt{branch\_combine\_m (def none)} \\
\textit{Reserved} & \texttt{40:45} & \texttt{zero} \\
\hyperref[enum:dest_width_m]{\texttt{dest\_width}} & \texttt{46:47} & \texttt{exact=none} \\
\texttt{opcode1} & \texttt{48:51} & \texttt{exact=BRANCH} \\
\texttt{instruction\_class} & \texttt{52:56} & \texttt{exact=N1} \\
\texttt{fau\_page\_index} & \texttt{57:58} & \texttt{uint} \\
\hyperref[enum:flow_control_m]{\texttt{flow}} & \texttt{59:62} & \texttt{flow\_control\_m} \\
\textit{Reserved} & \texttt{63} & \texttt{zero} \\
\bottomrule
\end{tabularx}
\end{center}
\textbf{Fields Explained:}
\begin{tightitemize}\footnotesize
  \item \texttt{src0}: Primary source operand (\texttt{r0..r63}, \texttt{u0..u63} uniform, \texttt{k*} FAU constant, or \texttt{t0..t2} temporary; \texttt{\textasciicircum{}} marks discard/last-use). \emph{Example:} {\scriptsize\ttfamily \mnem{FADD.f32} \dst{r0}, \src{r1}, \src{r2}} (\textbf{SPIR-V:} \texttt{OpFAdd (Operand 0)}; \textbf{GLSL:} \texttt{float c = a + b;})
  \item \texttt{offset}: Signed byte offset added to base address register in memory load/store/atomic instructions. \emph{Example:} {\scriptsize\ttfamily \mnem{LOAD.i32} \dst{@r0}, \src{[r2:r3]}, \imm{\#16}} (\textbf{SPIR-V:} \texttt{OpAccessChain + OpLoad}; \textbf{GLSL:} \texttt{uint y = ssbo.arr[4];})
  \item \hyperref[enum:branch_conservative_m]{\texttt{conservative}}: Conservative warp-execution hint on \texttt{BRANCH}/\texttt{JUMP} when reconvergence is non-trivial. \emph{Example:} {\scriptsize\ttfamily \mnem{BRANCH}\imod{.conservative} \src{r0}, \imm{\#+12}} (\textbf{SPIR-V:} \texttt{OpBranchConditional}; \textbf{GLSL:} \texttt{while (cond) \{ ... \}})
  \item \hyperref[enum:branch_negate_m]{\texttt{not}}: Boolean negation of the evaluated branch condition on \texttt{BRANCH}/\texttt{JUMP}/\texttt{JUMP\_EX} (inverts the boolean predicate, NOT bitwise \texttt{\textasciitilde{}src0}). \emph{Example:} {\scriptsize\ttfamily \mnem{BRANCH}\imod{.not} \src{k0}, \imm{\#0x10}} (\textbf{SPIR-V:} \texttt{OpLogicalNot + OpBranchConditional}; \textbf{GLSL:} \texttt{if (!cond) goto target;})
  \item \hyperref[enum:branch_combine_m]{\texttt{combine}}: Combines 16-bit halfword predicates (\texttt{h0}, \texttt{h1}, \texttt{and}, \texttt{lowbits}) before evaluating a branch. \emph{Example:} {\scriptsize\ttfamily \mnem{BRANCH}\imod{.and} \src{r0}, \imm{\#+8}} (\textbf{SPIR-V:} \texttt{OpBranchConditional (packed predicate)}; \textbf{GLSL:} \texttt{if (condA \&\& condB) \{ ... \}})
  \item \hyperref[enum:dest_width_m]{\texttt{dest\_width}}: Stage C destination lane write-mask (\texttt{none} = full 32-bit, \texttt{h0} = low 16-bit, \texttt{h1} = high 16-bit without bit-shifting). \emph{Example:} {\scriptsize\ttfamily \mnem{FADD.v2f16} \dst{r0}\imod{.h1}, \src{r1}, \src{r2}} (\textbf{SPIR-V:} \texttt{OpVectorInsertDynamic / WriteMask}; \textbf{GLSL:} \texttt{dst.y = a.y + b.y;})
  \item \texttt{opcode1}: Primary sub-opcode within the instruction class uniquely identifying the operation family. \emph{Example:} {\scriptsize\ttfamily \mnem{IADD.s32} \dst{r0}, \src{r1}, \src{r2}} (\textbf{SPIR-V:} \texttt{OpIAdd (operation selector)}; \textbf{GLSL:} \texttt{// Primary opcode bits})
  \item \texttt{instruction\_class}: 4-bit major instruction class selector (\texttt{bits 48:51}) partitioning the 64-bit Valhall opcode space (e.g., \texttt{FMA}, \texttt{CVT}, \texttt{SFU}, \texttt{LS}, \texttt{TEX}, \texttt{BRANCH}). \emph{Example:} {\scriptsize\ttfamily \mnem{FADD.f32} \dst{r0}, \src{r1}, \src{r2}} (\textbf{SPIR-V:} \texttt{N/A (encoding class)}; \textbf{GLSL:} \texttt{// Hardware execution pipe / decode class})
  \item \texttt{fau\_page\_index}: Selects the FAU (Fast Access Uniform) page (\texttt{0..3}) when a source operand references hardware small constants, thread IDs, or lane constants. \emph{Example:} {\scriptsize\ttfamily \mnem{FADD.f32} \dst{r0}, \src{r1}, \src{k1f}} (\textbf{SPIR-V:} \texttt{OpConstant / BuiltIn}; \textbf{GLSL:} \texttt{float y = x + 1.0;})
  \item \hyperref[enum:flow_control_m]{\texttt{flow}}: 4-bit embedded flow-control \& scoreboard wait modifier (\texttt{bits 59:62}: \texttt{none}, \texttt{wait0..wait0126}, \texttt{wait}, \texttt{reconverge}, \texttt{discard}, \texttt{end}). \emph{Example:} {\scriptsize\ttfamily \mnem{FADD.f32}\imod{.wait0.reconverge} \dst{r0}, \src{r1}, \src{r2}} (\textbf{SPIR-V:} \texttt{OpControlBarrier / OpReturn}; \textbf{GLSL:} \texttt{// Scoreboard wait + warp reconvergence})
\end{tightitemize}
\clearpage

\subsection*{DISCARD}
\addcontentsline{toc}{subsection}{DISCARD}
Fragment discard: evaluates \texttt{FCMP} condition \texttt{cmpf(src0, src1)} (both \texttt{f32}, with \texttt{widen0/1}, \texttt{abs0/1}, \texttt{neg0/1}); if true, discards fragment (\texttt{discarded = true}) and terminates thread.
\par\textbf{Notes:}
\begin{tightitemize}
  \item Takes \textbf{2 FP32 sources} and a \textbf{`cmpf`} modifier (\texttt{ops.rs:1039-1060}, \texttt{encode\_v9.rs:1129-1148}, \texttt{opt\_fuse.rs:98-142}).
  \item Only valid in fragment shaders.
  \item \textbf{Source references:} \texttt{kraid/ops.rs:1039-1060}, \texttt{encode\_v9.rs:1129-1148}, \texttt{opt\_fuse.rs:98-142}, \texttt{bifrost/valhall/ISA.xml:37}.
\end{tightitemize}
\codebox{\mnem{DISCARD}\imod{\{.flow\}}\imod{\{.cmpf\}} \src{src0}\imod{\{.widen0\}}\imod{\{.abs0\}}\imod{\{.neg0\}}, \src{src1}\imod{\{.widen1\}}\imod{\{.abs1\}}\imod{\{.neg1\}}}
\textbf{Operands:}
\begin{tightitemize}\footnotesize
  \item \texttt{\src{src0}}: Left-hand comparison operand; lanes matching \texttt{cmpf(src0, src1)} are marked discarded/helper.
  \item \texttt{\src{src1}}: Right-hand comparison operand.
\end{tightitemize}
\par\textbf{Closest SPIR-V Op:} \texttt{OpKill / OpDemoteToHelperInvocation}
\par\textbf{Pseudo-GLSL (Operational Semantics):}
\codebox{if (cmpf(src0, src1)) discard; // mark active fragment lane discarded/helper}
\textbf{Execution Context:}
\begin{tightitemize}\footnotesize
  \item \textbf{Depth/Stencil \& coverage state (\hyperref[sec:csf-zsd-blend]{\texttt{IDVSSR::ZSD}}, \texttt{r52:r53}):} Interacts with 32B \texttt{Depth/stencil} (\texttt{depth\_func}, \texttt{stencil\_from\_shader}), \hyperref[sec:csf-tiler]{\texttt{DCD2}} (\texttt{r38}: \texttt{no\_shader\_depth\_read} Early-Z vs Late-Z), \texttt{DCD0}/\texttt{DCD1} (\texttt{r57:r58}: \texttt{shader\_modifies\_coverage}, \texttt{alpha\_to\_coverage}), and \texttt{FauSpecial::atest\_datum}.
  \item \textbf{Fragment stage preloads \& FP modes (\hyperref[sec:csf-spd]{\texttt{SPD\_2}}, \texttt{r20:r21}):} Hardware pre-fills \texttt{r59} with \texttt{floatBitsToUint(FragCoord.z)} and \texttt{r60} with the MSAA sample coverage mask (\texttt{gl\_SampleMaskIn[0]}), while \texttt{ShaderProgram} word~0 latches 32/64 GPR allocation and warp-wide \texttt{flush\_to\_zero\_mode}, \texttt{suppress\_na\_n}, and \texttt{suppress\_inf}.
\end{tightitemize}
\par\textbf{Worked Examples:}
\par\textbf{Example 1 (Alpha-test fragment discard):}
\codebox{\mnem{DISCARD.f32}\imod{.lt} \src{r3}, \src{k0fr2f}\imod{.h1}}
Evaluates \texttt{r3 \textless{} 0.5f}; active lanes where true are marked discarded/helper and suppress color/depth/SSBO writes.
\par\textbf{Example 2 (Unconditional discard via `k0 == k0`):}
\codebox{\mnem{DISCARD.f32}\imod{.eq} \src{k0}, \src{k0}}
Always evaluates true (\texttt{0.0 == 0.0}), discarding all currently active lanes in the divergent branch.
\par\textbf{Mined SPIR-V $\to$ Valhall ISA Example (Lowering fragment discard (OpKill / OpDemoteToHelperInvocation) --- Fragment shader, Mali-G710 (v10)):}
\codebox{\textbf{; SPIR-V (input)}\\\%kill = OpFOrdLessThan \%bool \%val \%neg\_thresh\\OpBranchConditional \%kill \%discard\_blk \%cont\_blk\\\%discard\_blk = OpLabel ; OpKill\\[3pt]\textbf{; Valhall ISA (malioc output)}\\DISCARD.f32.wait0126.lt r0\textasciicircum{}, u0.w1.neg\\ATEST.slot0.wait0 @r60, r60, k0, atest\_datum.w0}
\texttt{DISCARD.f32.lt} directly evaluates \texttt{r0 \textless{} -u0.w1} and marks failing fragment lanes as helper/discarded without branching, followed by \texttt{ATEST} to commit updated coverage to \texttt{r60}.
\begin{center}\footnotesize
\begin{tabularx}{\linewidth}{@{}lllX@{}}
\toprule
Variant & Arch & Class & Opcode Fields \\
\midrule
\texttt{(none)} & v9--v14 & \texttt{N2} & \texttt{opcode1=DISCARD} \\
\bottomrule
\end{tabularx}
\end{center}
\begin{center}\footnotesize
\begin{tabularx}{0.88\linewidth}{@{}l l X@{}}
\toprule
Field Name & Bits & Type / Fixed Value \\
\midrule
\texttt{src0} & \texttt{0:7} & \texttt{SourceEncoding} \\
\texttt{src1} & \texttt{8:15} & \texttt{SourceEncoding} \\
\textit{Reserved} & \texttt{16:25} & \texttt{zero} \\
\hyperref[enum:widen_m]{\texttt{widen1}} & \texttt{26:27} & \texttt{widen\_m (def none)} \\
\hyperref[enum:widen_m]{\texttt{widen0}} & \texttt{28:29} & \texttt{widen\_m (def none)} \\
\textit{Reserved} & \texttt{30:31} & \texttt{zero} \\
\hyperref[enum:cmpf_m]{\texttt{cmpf}} & \texttt{32:34} & \texttt{cmpf\_m} \\
\textit{Reserved} & \texttt{35} & \texttt{zero} \\
\hyperref[enum:neg_m]{\texttt{neg1}} & \texttt{36} & \texttt{neg\_m (def none)} \\
\hyperref[enum:abs_m]{\texttt{abs1}} & \texttt{37} & \texttt{abs\_m (def none)} \\
\hyperref[enum:neg_m]{\texttt{neg0}} & \texttt{38} & \texttt{neg\_m (def none)} \\
\hyperref[enum:abs_m]{\texttt{abs0}} & \texttt{39} & \texttt{abs\_m (def none)} \\
\textit{Reserved} & \texttt{40:45} & \texttt{zero} \\
\hyperref[enum:dest_width_m]{\texttt{dest\_width}} & \texttt{46:47} & \texttt{exact=none} \\
\texttt{opcode1} & \texttt{48:51} & \texttt{exact=DISCARD} \\
\texttt{instruction\_class} & \texttt{52:56} & \texttt{exact=N2} \\
\texttt{fau\_page\_index} & \texttt{57:58} & \texttt{uint} \\
\hyperref[enum:flow_control_m]{\texttt{flow}} & \texttt{59:62} & \texttt{flow\_control\_m} \\
\textit{Reserved} & \texttt{63} & \texttt{zero} \\
\bottomrule
\end{tabularx}
\end{center}
\textbf{Fields Explained:}
\begin{tightitemize}\footnotesize
  \item \texttt{src0}: Primary source operand (\texttt{r0..r63}, \texttt{u0..u63} uniform, \texttt{k*} FAU constant, or \texttt{t0..t2} temporary; \texttt{\textasciicircum{}} marks discard/last-use). \emph{Example:} {\scriptsize\ttfamily \mnem{FADD.f32} \dst{r0}, \src{r1}, \src{r2}} (\textbf{SPIR-V:} \texttt{OpFAdd (Operand 0)}; \textbf{GLSL:} \texttt{float c = a + b;})
  \item \texttt{src1}: Secondary source operand (register, uniform, or FAU constant; in shifts, supplies the shift count byte). \emph{Example:} {\scriptsize\ttfamily \mnem{LSHIFT.i32} \dst{r0}, \src{r1}, \src{r2}\imod{.b0}} (\textbf{SPIR-V:} \texttt{OpShiftLeftLogical (Shift)}; \textbf{GLSL:} \texttt{uint y = a \textless{}\textless{} s;})
  \item \hyperref[enum:widen_m]{\texttt{widen1}}: Stage A widening extractor on \texttt{src1} (promotes narrow byte/half/word lane before Stage B). \emph{Example:} {\scriptsize\ttfamily \mnem{IMUL.u64} \dst{[r0:r1]}, \src{[r2:r3]}\imod{.w0}, \src{[r4:r5]}\imod{.w0}} (\textbf{SPIR-V:} \texttt{OpUConvert / OpUMulExtended}; \textbf{GLSL:} \texttt{uint64\_t p = uint64\_t(a) * uint64\_t(b);})
  \item \hyperref[enum:widen_m]{\texttt{widen0}}: Stage A widening extractor on \texttt{src0} (promotes FP16 \texttt{h0}/\texttt{h1} to FP32, or sign/zero-extends \texttt{b0..b3}/\texttt{h0..h1}/\texttt{w0..w1} to 32/64-bit before \texttt{abs}/\texttt{neg}). \emph{Example:} {\scriptsize\ttfamily \mnem{IADD.s32} \dst{r0}, \src{r1}\imod{.h0}, \src{r2}} (\textbf{SPIR-V:} \texttt{OpSConvert / OpFConvert}; \textbf{GLSL:} \texttt{int c = int(int16\_t(a)) + b;})
  \item \hyperref[enum:cmpf_m]{\texttt{cmpf}}: Comparison predicate (\texttt{eq}, \texttt{ne}, \texttt{lt}, \texttt{le}, \texttt{gt}, \texttt{ge}, \texttt{gtlt}, \texttt{total}). \emph{Example:} {\scriptsize\ttfamily \mnem{CSEL.f32}\imod{.lt} \dst{r0}, \src{r1}, \src{r2}, \src{r3}, \src{r4}} (\textbf{SPIR-V:} \texttt{OpFOrdLessThan + OpSelect}; \textbf{GLSL:} \texttt{float y = (a \textless{} b) ? c : d;})
  \item \hyperref[enum:neg_m]{\texttt{neg1}}: Stage A floating-point negation on \texttt{src1} (flips IEEE sign bit after swizzle/widen and \texttt{abs1}). \emph{Example:} {\scriptsize\ttfamily \mnem{FMA.f32} \dst{r0}, \src{r1}, \src{r2}\imod{.neg}, \src{r3}} (\textbf{SPIR-V:} \texttt{OpFNegate}; \textbf{GLSL:} \texttt{float d = fma(a, -b, c);})
  \item \hyperref[enum:abs_m]{\texttt{abs1}}: Stage A floating-point absolute value on \texttt{src1} (clears IEEE sign bit after swizzle/widen, before \texttt{neg1}). \emph{Example:} {\scriptsize\ttfamily \mnem{FMA.f32} \dst{r0}, \src{r1}, \src{r2}\imod{.abs}, \src{r3}} (\textbf{SPIR-V:} \texttt{GLSL.std.450 FAbs}; \textbf{GLSL:} \texttt{float d = fma(a, abs(b), c);})
  \item \hyperref[enum:neg_m]{\texttt{neg0}}: Stage A floating-point negation on \texttt{src0} (flips IEEE sign bit after swizzle/widen and \texttt{abs0}). \emph{Example:} {\scriptsize\ttfamily \mnem{FADD.f32} \dst{r0}, \src{r1}\imod{.neg}, \src{r2}} (\textbf{SPIR-V:} \texttt{OpFNegate}; \textbf{GLSL:} \texttt{float c = -a + b;})
  \item \hyperref[enum:abs_m]{\texttt{abs0}}: Stage A floating-point absolute value on \texttt{src0} (clears IEEE sign bit after swizzle/widen, before \texttt{neg0}). \emph{Example:} {\scriptsize\ttfamily \mnem{FADD.f32} \dst{r0}, \src{r1}\imod{.abs}, \src{r2}} (\textbf{SPIR-V:} \texttt{GLSL.std.450 FAbs}; \textbf{GLSL:} \texttt{float c = abs(a) + b;})
  \item \hyperref[enum:dest_width_m]{\texttt{dest\_width}}: Stage C destination lane write-mask (\texttt{none} = full 32-bit, \texttt{h0} = low 16-bit, \texttt{h1} = high 16-bit without bit-shifting). \emph{Example:} {\scriptsize\ttfamily \mnem{FADD.v2f16} \dst{r0}\imod{.h1}, \src{r1}, \src{r2}} (\textbf{SPIR-V:} \texttt{OpVectorInsertDynamic / WriteMask}; \textbf{GLSL:} \texttt{dst.y = a.y + b.y;})
  \item \texttt{opcode1}: Primary sub-opcode within the instruction class uniquely identifying the operation family. \emph{Example:} {\scriptsize\ttfamily \mnem{IADD.s32} \dst{r0}, \src{r1}, \src{r2}} (\textbf{SPIR-V:} \texttt{OpIAdd (operation selector)}; \textbf{GLSL:} \texttt{// Primary opcode bits})
  \item \texttt{instruction\_class}: 4-bit major instruction class selector (\texttt{bits 48:51}) partitioning the 64-bit Valhall opcode space (e.g., \texttt{FMA}, \texttt{CVT}, \texttt{SFU}, \texttt{LS}, \texttt{TEX}, \texttt{BRANCH}). \emph{Example:} {\scriptsize\ttfamily \mnem{FADD.f32} \dst{r0}, \src{r1}, \src{r2}} (\textbf{SPIR-V:} \texttt{N/A (encoding class)}; \textbf{GLSL:} \texttt{// Hardware execution pipe / decode class})
  \item \texttt{fau\_page\_index}: Selects the FAU (Fast Access Uniform) page (\texttt{0..3}) when a source operand references hardware small constants, thread IDs, or lane constants. \emph{Example:} {\scriptsize\ttfamily \mnem{FADD.f32} \dst{r0}, \src{r1}, \src{k1f}} (\textbf{SPIR-V:} \texttt{OpConstant / BuiltIn}; \textbf{GLSL:} \texttt{float y = x + 1.0;})
  \item \hyperref[enum:flow_control_m]{\texttt{flow}}: 4-bit embedded flow-control \& scoreboard wait modifier (\texttt{bits 59:62}: \texttt{none}, \texttt{wait0..wait0126}, \texttt{wait}, \texttt{reconverge}, \texttt{discard}, \texttt{end}). \emph{Example:} {\scriptsize\ttfamily \mnem{FADD.f32}\imod{.wait0.reconverge} \dst{r0}, \src{r1}, \src{r2}} (\textbf{SPIR-V:} \texttt{OpControlBarrier / OpReturn}; \textbf{GLSL:} \texttt{// Scoreboard wait + warp reconvergence})
\end{tightitemize}
\clearpage

\subsection*{DOORBELL}
\addcontentsline{toc}{subsection}{DOORBELL}
Issues a privileged hardware CSF doorbell message to the command-stream firmware/job manager (system/firmware use only; no user-shader register operands).
\par\textbf{Notes:}
\begin{tightitemize}
  \item Not emitted by graphics/compute compilers.
  \item System message side effect.
  \item \textbf{Source references:} \texttt{bifrost/bifrost/ISA.xml:3851}, \texttt{isa-v9-v14.xml}.
\end{tightitemize}
\codebox{\mnem{DOORBELL}\imod{\{.flow\}}}
\textbf{Operands:} \textit{None} (control/system instruction with no register or immediate operands; configured via modifiers and \texttt{flow}).
\par\textbf{Closest SPIR-V Op:} \texttt{None (firmware/job-management)}
\par\textbf{Pseudo-GLSL (Operational Semantics):}
\codebox{csf\_doorbell(); // signal privileged CSF firmware doorbell}
\textbf{Execution Context:}
\begin{tightitemize}\footnotesize
  \item \textbf{Compute stage preloads (\hyperref[sec:csf-spd]{\texttt{SPD\_k}}, \texttt{r16..r23}):} Hardware pre-fills \texttt{r55..r57} with \texttt{WorkGroupID.xyz + JOB\_OFFSET} (from \hyperref[sec:csf-regs]{\texttt{COMPUTESR::JOB\_OFFSET\_X/Y/Z}} \texttt{r34..r36}) and \texttt{r60..r62} with \texttt{LocalInvocationID.xyz} or \texttt{GlobalInvocationID.xyz} (from \texttt{WG\_SIZE} \texttt{r33} \& \texttt{JOB\_SIZE} \texttt{r37..r39}, selected by \texttt{ShaderProgram::preload}) so compute shaders can derive thread/element indices without memory loads.
  \item \textbf{Vertex stage preloads (\hyperref[sec:csf-spd]{\texttt{SPD\_0}/\texttt{SPD\_1}}, \texttt{r16..r19}):} The hardware Index Fetcher reads \hyperref[sec:csf-regs]{\texttt{IDVSSR::INDEX\_BUFFER}} (\texttt{r54:r55}) and pre-fills \texttt{r59, r60} with \texttt{VertexIndex} (\texttt{+ VERTEX\_OFFSET + GLOBAL\_ATTR\_OFF}) and \texttt{r61} with \texttt{InstanceIndex} (\texttt{+ INSTANCE\_OFFSET}) so Position and Varying shaders can pull vertex attributes directly.
  \item \textbf{Inline FAU \& Special Operands (\hyperref[sec:csf-tsd-fau]{\texttt{FAU\_k}}, \texttt{r8..r15}):} Any 8-bit source with \texttt{mode=2} reads zero-latency uniform \texttt{u0..u31} from \texttt{FAU\_k} (\texttt{slots64 = bits[63:56]}); \texttt{mode=3} reads \texttt{SmallConstant} (Table~1.2) or \texttt{FauSpecial} (\texttt{pc}, \texttt{lane\_id}, \texttt{lane\_width}, \texttt{tls\_ptr}, \texttt{wls\_ptr}, \texttt{blend\_descriptor\_0..7}).
\end{tightitemize}
\par\textbf{Worked Examples:}
\par\textbf{Example 1 (Firmware doorbell notification):}
\codebox{\mnem{DOORBELL}}
Signals the CSF/job-management firmware doorbell from privileged microcode (never emitted by Vulkan/OpenGL compilers).
\par\textbf{Example 2 (With scoreboard drain):}
\codebox{\mnem{DOORBELL}\imod{.wait}}
Drains all in-flight asynchronous slots before ringing the hardware doorbell.

\begin{center}\footnotesize
\begin{tabularx}{\linewidth}{@{}lllX@{}}
\toprule
Variant & Arch & Class & Opcode Fields \\
\midrule
\texttt{(none)} & v9--v14 & \texttt{N0} & \texttt{opcode1=DOORBELL} \\
\bottomrule
\end{tabularx}
\end{center}
\begin{center}\footnotesize
\begin{tabularx}{0.88\linewidth}{@{}l l X@{}}
\toprule
Field Name & Bits & Type / Fixed Value \\
\midrule
\textit{Reserved} & \texttt{0:32} & \texttt{zero} \\
\texttt{sr\_count} & \texttt{33:35} & \texttt{exact=count0} \\
\textit{Reserved} & \texttt{36:39} & \texttt{zero} \\
\texttt{staging\_register\_index} & \texttt{40:45} & \texttt{exact=0} \\
\hyperref[enum:dest_width_m]{\texttt{dest\_width}} & \texttt{46:47} & \texttt{exact=none} \\
\texttt{opcode1} & \texttt{48:51} & \texttt{exact=DOORBELL} \\
\texttt{instruction\_class} & \texttt{52:56} & \texttt{exact=N0} \\
\texttt{fau\_page\_index} & \texttt{57:58} & \texttt{exact=0} \\
\hyperref[enum:flow_control_m]{\texttt{flow}} & \texttt{59:62} & \texttt{flow\_control\_m} \\
\textit{Reserved} & \texttt{63} & \texttt{zero} \\
\bottomrule
\end{tabularx}
\end{center}
\textbf{Fields Explained:}
\begin{tightitemize}\footnotesize
  \item \texttt{sr\_count}: Number of consecutive 32-bit staging registers transferred by a load/store/texture message. \emph{Example:} {\scriptsize\ttfamily \mnem{LD\_MULTI.i128} \dst{@r0:r3}, \src{r4}, \imm{\#0}} (\textbf{SPIR-V:} \texttt{OpLoad (wide vector)}; \textbf{GLSL:} \texttt{uvec4 x = buf[i];})
  \item \texttt{staging\_register\_index}: Base GPR index of the contiguous staging register window (\texttt{@rN:rM}) for message payload/return. \emph{Example:} {\scriptsize\ttfamily \mnem{STORE.i64}\imod{.slot7} \src{@r0:r1}, \src{r2}, \imm{\#8}} (\textbf{SPIR-V:} \texttt{OpStore (uvec2)}; \textbf{GLSL:} \texttt{ssbo.pairs[i] = uvec2(a, b);})
  \item \hyperref[enum:dest_width_m]{\texttt{dest\_width}}: Stage C destination lane write-mask (\texttt{none} = full 32-bit, \texttt{h0} = low 16-bit, \texttt{h1} = high 16-bit without bit-shifting). \emph{Example:} {\scriptsize\ttfamily \mnem{FADD.v2f16} \dst{r0}\imod{.h1}, \src{r1}, \src{r2}} (\textbf{SPIR-V:} \texttt{OpVectorInsertDynamic / WriteMask}; \textbf{GLSL:} \texttt{dst.y = a.y + b.y;})
  \item \texttt{opcode1}: Primary sub-opcode within the instruction class uniquely identifying the operation family. \emph{Example:} {\scriptsize\ttfamily \mnem{IADD.s32} \dst{r0}, \src{r1}, \src{r2}} (\textbf{SPIR-V:} \texttt{OpIAdd (operation selector)}; \textbf{GLSL:} \texttt{// Primary opcode bits})
  \item \texttt{instruction\_class}: 4-bit major instruction class selector (\texttt{bits 48:51}) partitioning the 64-bit Valhall opcode space (e.g., \texttt{FMA}, \texttt{CVT}, \texttt{SFU}, \texttt{LS}, \texttt{TEX}, \texttt{BRANCH}). \emph{Example:} {\scriptsize\ttfamily \mnem{FADD.f32} \dst{r0}, \src{r1}, \src{r2}} (\textbf{SPIR-V:} \texttt{N/A (encoding class)}; \textbf{GLSL:} \texttt{// Hardware execution pipe / decode class})
  \item \texttt{fau\_page\_index}: Selects the FAU (Fast Access Uniform) page (\texttt{0..3}) when a source operand references hardware small constants, thread IDs, or lane constants. \emph{Example:} {\scriptsize\ttfamily \mnem{FADD.f32} \dst{r0}, \src{r1}, \src{k1f}} (\textbf{SPIR-V:} \texttt{OpConstant / BuiltIn}; \textbf{GLSL:} \texttt{float y = x + 1.0;})
  \item \hyperref[enum:flow_control_m]{\texttt{flow}}: 4-bit embedded flow-control \& scoreboard wait modifier (\texttt{bits 59:62}: \texttt{none}, \texttt{wait0..wait0126}, \texttt{wait}, \texttt{reconverge}, \texttt{discard}, \texttt{end}). \emph{Example:} {\scriptsize\ttfamily \mnem{FADD.f32}\imod{.wait0.reconverge} \dst{r0}, \src{r1}, \src{r2}} (\textbf{SPIR-V:} \texttt{OpControlBarrier / OpReturn}; \textbf{GLSL:} \texttt{// Scoreboard wait + warp reconvergence})
\end{tightitemize}
\clearpage

\subsection*{EUREKA}
\addcontentsline{toc}{subsection}{EUREKA}
Signals a privileged hardware job-management event notification to the CSF firmware (system/firmware use only; no user-shader register operands).
\par\textbf{Notes:}
\begin{tightitemize}
  \item System/job-management message.
  \item \textbf{Source references:} \texttt{bifrost/bifrost/ISA.xml:3855}, \texttt{isa-v9-v14.xml}.
\end{tightitemize}
\codebox{\mnem{EUREKA}\imod{\{.flow\}}}
\textbf{Operands:} \textit{None} (control/system instruction with no register or immediate operands; configured via modifiers and \texttt{flow}).
\par\textbf{Closest SPIR-V Op:} \texttt{None (firmware/job-management)}
\par\textbf{Pseudo-GLSL (Operational Semantics):}
\codebox{csf\_signal\_event(); // signal privileged CSF job-management event}
\textbf{Execution Context:}
\begin{tightitemize}\footnotesize
  \item \textbf{Compute stage preloads (\hyperref[sec:csf-spd]{\texttt{SPD\_k}}, \texttt{r16..r23}):} Hardware pre-fills \texttt{r55..r57} with \texttt{WorkGroupID.xyz + JOB\_OFFSET} (from \hyperref[sec:csf-regs]{\texttt{COMPUTESR::JOB\_OFFSET\_X/Y/Z}} \texttt{r34..r36}) and \texttt{r60..r62} with \texttt{LocalInvocationID.xyz} or \texttt{GlobalInvocationID.xyz} (from \texttt{WG\_SIZE} \texttt{r33} \& \texttt{JOB\_SIZE} \texttt{r37..r39}, selected by \texttt{ShaderProgram::preload}) so compute shaders can derive thread/element indices without memory loads.
  \item \textbf{Vertex stage preloads (\hyperref[sec:csf-spd]{\texttt{SPD\_0}/\texttt{SPD\_1}}, \texttt{r16..r19}):} The hardware Index Fetcher reads \hyperref[sec:csf-regs]{\texttt{IDVSSR::INDEX\_BUFFER}} (\texttt{r54:r55}) and pre-fills \texttt{r59, r60} with \texttt{VertexIndex} (\texttt{+ VERTEX\_OFFSET + GLOBAL\_ATTR\_OFF}) and \texttt{r61} with \texttt{InstanceIndex} (\texttt{+ INSTANCE\_OFFSET}) so Position and Varying shaders can pull vertex attributes directly.
  \item \textbf{Inline FAU \& Special Operands (\hyperref[sec:csf-tsd-fau]{\texttt{FAU\_k}}, \texttt{r8..r15}):} Any 8-bit source with \texttt{mode=2} reads zero-latency uniform \texttt{u0..u31} from \texttt{FAU\_k} (\texttt{slots64 = bits[63:56]}); \texttt{mode=3} reads \texttt{SmallConstant} (Table~1.2) or \texttt{FauSpecial} (\texttt{pc}, \texttt{lane\_id}, \texttt{lane\_width}, \texttt{tls\_ptr}, \texttt{wls\_ptr}, \texttt{blend\_descriptor\_0..7}).
\end{tightitemize}
\par\textbf{Worked Examples:}
\par\textbf{Example 1 (Hardware event signal):}
\codebox{\mnem{EUREKA}}
Emits a hardware job-management event notification from system shaders.
\par\textbf{Example 2 (Synchronized event signal):}
\codebox{\mnem{EUREKA}\imod{.wait}}
Waits for all outstanding scoreboard slots before asserting the hardware event.

\begin{center}\footnotesize
\begin{tabularx}{\linewidth}{@{}lllX@{}}
\toprule
Variant & Arch & Class & Opcode Fields \\
\midrule
\texttt{(none)} & v9--v14 & \texttt{N0} & \texttt{opcode1=EUREKA} \\
\bottomrule
\end{tabularx}
\end{center}
\begin{center}\footnotesize
\begin{tabularx}{0.88\linewidth}{@{}l l X@{}}
\toprule
Field Name & Bits & Type / Fixed Value \\
\midrule
\textit{Reserved} & \texttt{0:32} & \texttt{zero} \\
\texttt{sr\_count} & \texttt{33:35} & \texttt{exact=count0} \\
\textit{Reserved} & \texttt{36:39} & \texttt{zero} \\
\texttt{staging\_register\_index} & \texttt{40:45} & \texttt{exact=0} \\
\hyperref[enum:dest_width_m]{\texttt{dest\_width}} & \texttt{46:47} & \texttt{exact=none} \\
\texttt{opcode1} & \texttt{48:51} & \texttt{exact=EUREKA} \\
\texttt{instruction\_class} & \texttt{52:56} & \texttt{exact=N0} \\
\texttt{fau\_page\_index} & \texttt{57:58} & \texttt{exact=0} \\
\hyperref[enum:flow_control_m]{\texttt{flow}} & \texttt{59:62} & \texttt{flow\_control\_m} \\
\textit{Reserved} & \texttt{63} & \texttt{zero} \\
\bottomrule
\end{tabularx}
\end{center}
\textbf{Fields Explained:}
\begin{tightitemize}\footnotesize
  \item \texttt{sr\_count}: Number of consecutive 32-bit staging registers transferred by a load/store/texture message. \emph{Example:} {\scriptsize\ttfamily \mnem{LD\_MULTI.i128} \dst{@r0:r3}, \src{r4}, \imm{\#0}} (\textbf{SPIR-V:} \texttt{OpLoad (wide vector)}; \textbf{GLSL:} \texttt{uvec4 x = buf[i];})
  \item \texttt{staging\_register\_index}: Base GPR index of the contiguous staging register window (\texttt{@rN:rM}) for message payload/return. \emph{Example:} {\scriptsize\ttfamily \mnem{STORE.i64}\imod{.slot7} \src{@r0:r1}, \src{r2}, \imm{\#8}} (\textbf{SPIR-V:} \texttt{OpStore (uvec2)}; \textbf{GLSL:} \texttt{ssbo.pairs[i] = uvec2(a, b);})
  \item \hyperref[enum:dest_width_m]{\texttt{dest\_width}}: Stage C destination lane write-mask (\texttt{none} = full 32-bit, \texttt{h0} = low 16-bit, \texttt{h1} = high 16-bit without bit-shifting). \emph{Example:} {\scriptsize\ttfamily \mnem{FADD.v2f16} \dst{r0}\imod{.h1}, \src{r1}, \src{r2}} (\textbf{SPIR-V:} \texttt{OpVectorInsertDynamic / WriteMask}; \textbf{GLSL:} \texttt{dst.y = a.y + b.y;})
  \item \texttt{opcode1}: Primary sub-opcode within the instruction class uniquely identifying the operation family. \emph{Example:} {\scriptsize\ttfamily \mnem{IADD.s32} \dst{r0}, \src{r1}, \src{r2}} (\textbf{SPIR-V:} \texttt{OpIAdd (operation selector)}; \textbf{GLSL:} \texttt{// Primary opcode bits})
  \item \texttt{instruction\_class}: 4-bit major instruction class selector (\texttt{bits 48:51}) partitioning the 64-bit Valhall opcode space (e.g., \texttt{FMA}, \texttt{CVT}, \texttt{SFU}, \texttt{LS}, \texttt{TEX}, \texttt{BRANCH}). \emph{Example:} {\scriptsize\ttfamily \mnem{FADD.f32} \dst{r0}, \src{r1}, \src{r2}} (\textbf{SPIR-V:} \texttt{N/A (encoding class)}; \textbf{GLSL:} \texttt{// Hardware execution pipe / decode class})
  \item \texttt{fau\_page\_index}: Selects the FAU (Fast Access Uniform) page (\texttt{0..3}) when a source operand references hardware small constants, thread IDs, or lane constants. \emph{Example:} {\scriptsize\ttfamily \mnem{FADD.f32} \dst{r0}, \src{r1}, \src{k1f}} (\textbf{SPIR-V:} \texttt{OpConstant / BuiltIn}; \textbf{GLSL:} \texttt{float y = x + 1.0;})
  \item \hyperref[enum:flow_control_m]{\texttt{flow}}: 4-bit embedded flow-control \& scoreboard wait modifier (\texttt{bits 59:62}: \texttt{none}, \texttt{wait0..wait0126}, \texttt{wait}, \texttt{reconverge}, \texttt{discard}, \texttt{end}). \emph{Example:} {\scriptsize\ttfamily \mnem{FADD.f32}\imod{.wait0.reconverge} \dst{r0}, \src{r1}, \src{r2}} (\textbf{SPIR-V:} \texttt{OpControlBarrier / OpReturn}; \textbf{GLSL:} \texttt{// Scoreboard wait + warp reconvergence})
\end{tightitemize}
\clearpage

\subsection*{JUMP}
\addcontentsline{toc}{subsection}{JUMP}
Conditional indirect control transfer: evaluates \texttt{taken = not ? !combine(src0) : combine(src0)}. If \texttt{taken}, jumps to the shader address in \texttt{src1} (absolute or PC-relative per \texttt{address\_mode}; jumping to address \texttt{0} terminates shader execution); otherwise falls through to \texttt{PC + 8}.
\par\textbf{Notes:}
\begin{tightitemize}
  \item \texttt{src1} is 32-bit address (\texttt{OpJump.address: S32}, \texttt{ops.rs:2609}).
  \item Used for blend shader calls (\texttt{lower\_blend\_call.rs:51-57}).
  \item \textbf{Source references:} \texttt{kraid/ops.rs:2605-2630}, \texttt{encode\_v9.rs:2015-2045}, \texttt{lower\_blend\_call.rs:31-58}.
\end{tightitemize}
\codebox{\mnem{JUMP}\imod{\{.flow\}}\imod{\{.conservative\}}\imod{\{.not\}}\imod{\{.address\_mode\}} \src{src0}\imod{\{.combine\}}, \src{src1}}
\textbf{Operands:}
\begin{tightitemize}\footnotesize
  \item \texttt{\src{src0}}: Condition register tested via \texttt{combine} (and inverted if \texttt{.not}).
  \item \texttt{\src{src1}}: Target shader address (\texttt{0} terminates the shader; absolute or PC-relative per \texttt{address\_mode}).
\end{tightitemize}
\par\textbf{Closest SPIR-V Op:} \texttt{OpFunctionCall / OpReturn (indirect)}
\par\textbf{Pseudo-GLSL (Operational Semantics):}
\codebox{bool cond = combine(src0);\\if (flag\_not ? !cond : cond) \{\\    if (src1 == 0u) return; // terminate shader warp\\    PC = (address\_mode == absolute) ? src1 : (PC + 8 + src1);\\\}}
\textbf{Execution Context:}
\begin{tightitemize}\footnotesize
  \item \textbf{Indirect / blend-shader dispatch (\hyperref[sec:csf-zsd-blend]{\texttt{IDVSSR::BLEND\_DESC}}):} Used for programmable blend-shader calls/returns (\texttt{Blend[rt]::internal.shader.pc}); jumping to \texttt{0} terminates warp execution.
  \item \textbf{Inline FAU \& Special Operands (\hyperref[sec:csf-tsd-fau]{\texttt{FAU\_k}}, \texttt{r8..r15}):} Any 8-bit source with \texttt{mode=2} reads zero-latency uniform \texttt{u0..u31} from \texttt{FAU\_k} (\texttt{slots64 = bits[63:56]}); \texttt{mode=3} reads \texttt{SmallConstant} (Table~1.2) or \texttt{FauSpecial} (\texttt{pc}, \texttt{lane\_id}, \texttt{lane\_width}, \texttt{tls\_ptr}, \texttt{wls\_ptr}, \texttt{blend\_descriptor\_0..7}).
\end{tightitemize}
\par\textbf{Worked Examples:}
\par\textbf{Example 1 (Unconditional absolute indirect jump (blend shader)):}
\codebox{\mnem{JUMP}\imod{.absolute} \src{k1n}, \src{r54}}
Since \texttt{k1n != 0}, jumps to 32-bit absolute shader PC in \texttt{r54} (\texttt{r54 == 0} terminates execution cleanly).
\par\textbf{Example 2 (Predicated indirect return):}
\codebox{\mnem{JUMP}\imod{.absolute} \src{r0}, \src{r54}}
Jumps to \texttt{r54} only if predicate \texttt{r0 != 0}; otherwise falls through to \texttt{PC + 8}.

\begin{center}\footnotesize
\begin{tabularx}{\linewidth}{@{}lllX@{}}
\toprule
Variant & Arch & Class & Opcode Fields \\
\midrule
\texttt{(none)} & v9--v14 & \texttt{N2} & \texttt{opcode1=JUMP} \\
\texttt{(none)} & v9--v14 & \texttt{N2} & \texttt{opcode1=JUMP} \\
\bottomrule
\end{tabularx}
\end{center}
\begin{center}\footnotesize
\begin{tabularx}{0.88\linewidth}{@{}l l X@{}}
\toprule
Field Name & Bits & Type / Fixed Value \\
\midrule
\texttt{src0} & \texttt{0:7} & \texttt{SourceEncoding} \\
\texttt{src1} & \texttt{8:15} & \texttt{SourceEncoding} \\
\textit{Reserved} & \texttt{16:34} & \texttt{zero} \\
\hyperref[enum:branch_conservative_m]{\texttt{conservative}} & \texttt{35} & \texttt{branch\_conservative\_m (def none)} \\
\hyperref[enum:branch_negate_m]{\texttt{not}} & \texttt{36} & \texttt{branch\_negate\_m (def none)} \\
\hyperref[enum:branch_combine_m]{\texttt{combine}} & \texttt{37:39} & \texttt{branch\_combine\_m (def none)} \\
\hyperref[enum:jump_address_mode_m]{\texttt{address\_mode}} & \texttt{40} & \texttt{jump\_address\_mode\_m} \\
\textit{Reserved} & \texttt{41:45} & \texttt{zero} \\
\hyperref[enum:dest_width_m]{\texttt{dest\_width}} & \texttt{46:47} & \texttt{exact=none} \\
\texttt{opcode1} & \texttt{48:51} & \texttt{exact=JUMP} \\
\texttt{instruction\_class} & \texttt{52:56} & \texttt{exact=N2} \\
\texttt{fau\_page\_index} & \texttt{57:58} & \texttt{uint} \\
\hyperref[enum:flow_control_m]{\texttt{flow}} & \texttt{59:62} & \texttt{flow\_control\_m} \\
\textit{Reserved} & \texttt{63} & \texttt{zero} \\
\bottomrule
\end{tabularx}
\end{center}
\textbf{Fields Explained:}
\begin{tightitemize}\footnotesize
  \item \texttt{src0}: Primary source operand (\texttt{r0..r63}, \texttt{u0..u63} uniform, \texttt{k*} FAU constant, or \texttt{t0..t2} temporary; \texttt{\textasciicircum{}} marks discard/last-use). \emph{Example:} {\scriptsize\ttfamily \mnem{FADD.f32} \dst{r0}, \src{r1}, \src{r2}} (\textbf{SPIR-V:} \texttt{OpFAdd (Operand 0)}; \textbf{GLSL:} \texttt{float c = a + b;})
  \item \texttt{src1}: Secondary source operand (register, uniform, or FAU constant; in shifts, supplies the shift count byte). \emph{Example:} {\scriptsize\ttfamily \mnem{LSHIFT.i32} \dst{r0}, \src{r1}, \src{r2}\imod{.b0}} (\textbf{SPIR-V:} \texttt{OpShiftLeftLogical (Shift)}; \textbf{GLSL:} \texttt{uint y = a \textless{}\textless{} s;})
  \item \hyperref[enum:branch_conservative_m]{\texttt{conservative}}: Conservative warp-execution hint on \texttt{BRANCH}/\texttt{JUMP} when reconvergence is non-trivial. \emph{Example:} {\scriptsize\ttfamily \mnem{BRANCH}\imod{.conservative} \src{r0}, \imm{\#+12}} (\textbf{SPIR-V:} \texttt{OpBranchConditional}; \textbf{GLSL:} \texttt{while (cond) \{ ... \}})
  \item \hyperref[enum:branch_negate_m]{\texttt{not}}: Boolean negation of the evaluated branch condition on \texttt{BRANCH}/\texttt{JUMP}/\texttt{JUMP\_EX} (inverts the boolean predicate, NOT bitwise \texttt{\textasciitilde{}src0}). \emph{Example:} {\scriptsize\ttfamily \mnem{BRANCH}\imod{.not} \src{k0}, \imm{\#0x10}} (\textbf{SPIR-V:} \texttt{OpLogicalNot + OpBranchConditional}; \textbf{GLSL:} \texttt{if (!cond) goto target;})
  \item \hyperref[enum:branch_combine_m]{\texttt{combine}}: Combines 16-bit halfword predicates (\texttt{h0}, \texttt{h1}, \texttt{and}, \texttt{lowbits}) before evaluating a branch. \emph{Example:} {\scriptsize\ttfamily \mnem{BRANCH}\imod{.and} \src{r0}, \imm{\#+8}} (\textbf{SPIR-V:} \texttt{OpBranchConditional (packed predicate)}; \textbf{GLSL:} \texttt{if (condA \&\& condB) \{ ... \}})
  \item \hyperref[enum:jump_address_mode_m]{\texttt{address\_mode}}: Selects PC-relative vs absolute 64-bit jump address calculation. \emph{Example:} {\scriptsize\ttfamily \mnem{JUMP}\imod{.absolute} \src{r0}, \src{[r2:r3]}} (\textbf{SPIR-V:} \texttt{OpSwitch / indirect function call}; \textbf{GLSL:} \texttt{// Jump table dispatch})
  \item \hyperref[enum:dest_width_m]{\texttt{dest\_width}}: Stage C destination lane write-mask (\texttt{none} = full 32-bit, \texttt{h0} = low 16-bit, \texttt{h1} = high 16-bit without bit-shifting). \emph{Example:} {\scriptsize\ttfamily \mnem{FADD.v2f16} \dst{r0}\imod{.h1}, \src{r1}, \src{r2}} (\textbf{SPIR-V:} \texttt{OpVectorInsertDynamic / WriteMask}; \textbf{GLSL:} \texttt{dst.y = a.y + b.y;})
  \item \texttt{opcode1}: Primary sub-opcode within the instruction class uniquely identifying the operation family. \emph{Example:} {\scriptsize\ttfamily \mnem{IADD.s32} \dst{r0}, \src{r1}, \src{r2}} (\textbf{SPIR-V:} \texttt{OpIAdd (operation selector)}; \textbf{GLSL:} \texttt{// Primary opcode bits})
  \item \texttt{instruction\_class}: 4-bit major instruction class selector (\texttt{bits 48:51}) partitioning the 64-bit Valhall opcode space (e.g., \texttt{FMA}, \texttt{CVT}, \texttt{SFU}, \texttt{LS}, \texttt{TEX}, \texttt{BRANCH}). \emph{Example:} {\scriptsize\ttfamily \mnem{FADD.f32} \dst{r0}, \src{r1}, \src{r2}} (\textbf{SPIR-V:} \texttt{N/A (encoding class)}; \textbf{GLSL:} \texttt{// Hardware execution pipe / decode class})
  \item \texttt{fau\_page\_index}: Selects the FAU (Fast Access Uniform) page (\texttt{0..3}) when a source operand references hardware small constants, thread IDs, or lane constants. \emph{Example:} {\scriptsize\ttfamily \mnem{FADD.f32} \dst{r0}, \src{r1}, \src{k1f}} (\textbf{SPIR-V:} \texttt{OpConstant / BuiltIn}; \textbf{GLSL:} \texttt{float y = x + 1.0;})
  \item \hyperref[enum:flow_control_m]{\texttt{flow}}: 4-bit embedded flow-control \& scoreboard wait modifier (\texttt{bits 59:62}: \texttt{none}, \texttt{wait0..wait0126}, \texttt{wait}, \texttt{reconverge}, \texttt{discard}, \texttt{end}). \emph{Example:} {\scriptsize\ttfamily \mnem{FADD.f32}\imod{.wait0.reconverge} \dst{r0}, \src{r1}, \src{r2}} (\textbf{SPIR-V:} \texttt{OpControlBarrier / OpReturn}; \textbf{GLSL:} \texttt{// Scoreboard wait + warp reconvergence})
\end{tightitemize}
\clearpage

\subsection*{JUMP\_EX}
\addcontentsline{toc}{subsection}{JUMP\_EX}
Extended indirect jump/call/return: conditional jump with \texttt{not}/\texttt{combine}/\texttt{conservative}, divergence hint (\texttt{diverge}), and call-stack operation (\texttt{stack\_mode}: \texttt{none}/\texttt{call}/\texttt{return}/\texttt{replace}) using \texttt{src1} and \texttt{src2}.
\par\textbf{Notes:}
\begin{tightitemize}
  \item Supports hardware call/return stack management.
  \item \textbf{Source references:} \texttt{isa-v9-v14.xml}, \texttt{bifrost/compiler.h:715}, \texttt{bifrost/bifrost/ISA.xml:1621}.
\end{tightitemize}
\codebox{\mnem{JUMP\_EX}\imod{\{.flow\}}\imod{\{.conservative\}}\imod{\{.not\}}\imod{\{.stack\_mode\}}\imod{\{.diverge\}} \src{src0}\imod{\{.combine\}}, \src{src1}, \src{src2}}
\textbf{Operands:}
\begin{tightitemize}\footnotesize
  \item \texttt{\src{src0}}: Condition register guarding the call/return/tail-call.
  \item \texttt{\src{src1}}: Target shader address (for \texttt{.call} / \texttt{.replace}).
  \item \texttt{\src{src2}}: Auxiliary call/stack operand.
\end{tightitemize}
\par\textbf{Closest SPIR-V Op:} \texttt{OpFunctionCall / OpReturn (stack)}
\par\textbf{Pseudo-GLSL (Operational Semantics):}
\codebox{if (stack\_mode == call)        \{ hw\_stack.push(PC + 8); PC = resolveTarget(src1, src2); \}\\else if (stack\_mode == ret)    \{ PC = hw\_stack.pop(); \}\\else /* replace (tail-call) */ \{ PC = resolveTarget(src1, src2); \}}
\textbf{Execution Context:}
\begin{tightitemize}\footnotesize
  \item \textbf{Hardware call-stack requirement (\hyperref[sec:csf-spd]{\texttt{SPD\_k}} \& \hyperref[sec:csf-tsd-fau]{\texttt{TSD\_k}}):} Requires \texttt{ShaderProgram::shader\_contains\_jump\_ex = 1} (word~0 bit~29) and per-thread stack backing in \texttt{TSD\_k}.
  \item \textbf{Inline FAU \& Special Operands (\hyperref[sec:csf-tsd-fau]{\texttt{FAU\_k}}, \texttt{r8..r15}):} Any 8-bit source with \texttt{mode=2} reads zero-latency uniform \texttt{u0..u31} from \texttt{FAU\_k} (\texttt{slots64 = bits[63:56]}); \texttt{mode=3} reads \texttt{SmallConstant} (Table~1.2) or \texttt{FauSpecial} (\texttt{pc}, \texttt{lane\_id}, \texttt{lane\_width}, \texttt{tls\_ptr}, \texttt{wls\_ptr}, \texttt{blend\_descriptor\_0..7}).
\end{tightitemize}
\par\textbf{Worked Examples:}
\par\textbf{Example 1 (Hardware call-stack push (`.call`)):}
\codebox{\mnem{JUMP\_EX}\imod{.call} \src{k1n}, \src{r1}, \src{r2}}
Pushes return address \texttt{PC + 8} onto the hardware call stack and jumps to target address.
\par\textbf{Example 2 (Hardware call-stack return (`.return`)):}
\codebox{\mnem{JUMP\_EX}\imod{.return} \src{k1n}, \src{k0}, \src{k0}}
Pops the saved return PC from the hardware call stack and resumes the caller.

\begin{center}\footnotesize
\begin{tabularx}{\linewidth}{@{}lllX@{}}
\toprule
Variant & Arch & Class & Opcode Fields \\
\midrule
\texttt{(none)} & v9--v14 & \texttt{N3} & \texttt{opcode1=JUMP\_EX} \\
\bottomrule
\end{tabularx}
\end{center}
\begin{center}\footnotesize
\begin{tabularx}{0.88\linewidth}{@{}l l X@{}}
\toprule
Field Name & Bits & Type / Fixed Value \\
\midrule
\texttt{src0} & \texttt{0:7} & \texttt{SourceEncoding} \\
\texttt{src1} & \texttt{8:15} & \texttt{SourceEncoding} \\
\texttt{src2} & \texttt{16:23} & \texttt{SourceEncoding} \\
\textit{Reserved} & \texttt{24:31} & \texttt{zero} \\
\hyperref[enum:jump_divergence_mode_m]{\texttt{diverge}} & \texttt{32} & \texttt{jump\_divergence\_mode\_m (def none)} \\
\hyperref[enum:jump_stack_mode_m]{\texttt{stack\_mode}} & \texttt{33:34} & \texttt{jump\_stack\_mode\_m (def none)} \\
\hyperref[enum:branch_conservative_m]{\texttt{conservative}} & \texttt{35} & \texttt{branch\_conservative\_m (def none)} \\
\hyperref[enum:branch_negate_m]{\texttt{not}} & \texttt{36} & \texttt{branch\_negate\_m (def none)} \\
\hyperref[enum:branch_combine_m]{\texttt{combine}} & \texttt{37:39} & \texttt{branch\_combine\_m (def none)} \\
\textit{Reserved} & \texttt{40:45} & \texttt{zero} \\
\hyperref[enum:dest_width_m]{\texttt{dest\_width}} & \texttt{46:47} & \texttt{exact=none} \\
\texttt{opcode1} & \texttt{48:51} & \texttt{exact=JUMP\_EX} \\
\texttt{instruction\_class} & \texttt{52:56} & \texttt{exact=N3} \\
\texttt{fau\_page\_index} & \texttt{57:58} & \texttt{uint} \\
\hyperref[enum:flow_control_m]{\texttt{flow}} & \texttt{59:62} & \texttt{flow\_control\_m} \\
\textit{Reserved} & \texttt{63} & \texttt{zero} \\
\bottomrule
\end{tabularx}
\end{center}
\textbf{Fields Explained:}
\begin{tightitemize}\footnotesize
  \item \texttt{src0}: Primary source operand (\texttt{r0..r63}, \texttt{u0..u63} uniform, \texttt{k*} FAU constant, or \texttt{t0..t2} temporary; \texttt{\textasciicircum{}} marks discard/last-use). \emph{Example:} {\scriptsize\ttfamily \mnem{FADD.f32} \dst{r0}, \src{r1}, \src{r2}} (\textbf{SPIR-V:} \texttt{OpFAdd (Operand 0)}; \textbf{GLSL:} \texttt{float c = a + b;})
  \item \texttt{src1}: Secondary source operand (register, uniform, or FAU constant; in shifts, supplies the shift count byte). \emph{Example:} {\scriptsize\ttfamily \mnem{LSHIFT.i32} \dst{r0}, \src{r1}, \src{r2}\imod{.b0}} (\textbf{SPIR-V:} \texttt{OpShiftLeftLogical (Shift)}; \textbf{GLSL:} \texttt{uint y = a \textless{}\textless{} s;})
  \item \texttt{src2}: Third source operand (FMA addend, \texttt{MUX}/\texttt{CSEL} selector, \texttt{IADDC} carry-in, or bitwise mask in fused shift-logic). \emph{Example:} {\scriptsize\ttfamily \mnem{FMA.f32} \dst{r0}, \src{r1}, \src{r2}, \src{r3}} (\textbf{SPIR-V:} \texttt{GLSL.std.450 Fma (Addend)}; \textbf{GLSL:} \texttt{float d = fma(a, b, c);})
  \item \hyperref[enum:jump_divergence_mode_m]{\texttt{diverge}}: Signals potentially divergent indirect jump across warp lanes in \texttt{JUMP\_EX}. \emph{Example:} {\scriptsize\ttfamily \mnem{JUMP\_EX}\imod{.diverge} \src{r0}, \src{[r2:r3]}} (\textbf{SPIR-V:} \texttt{OpFunctionCall (divergent)}; \textbf{GLSL:} \texttt{if (gl\_LocalInvocationID.x \textless{} k) f();})
  \item \hyperref[enum:jump_stack_mode_m]{\texttt{stack\_mode}}: Controls hardware call-stack push/pop/replace on \texttt{JUMP\_EX} (\texttt{call}, \texttt{return}, \texttt{replace}, \texttt{none}). \emph{Example:} {\scriptsize\ttfamily \mnem{JUMP\_EX}\imod{.call} \src{r0}, \imm{\#+32}} (\textbf{SPIR-V:} \texttt{OpFunctionCall}; \textbf{GLSL:} \texttt{float y = helperFunc(x);})
  \item \hyperref[enum:branch_conservative_m]{\texttt{conservative}}: Conservative warp-execution hint on \texttt{BRANCH}/\texttt{JUMP} when reconvergence is non-trivial. \emph{Example:} {\scriptsize\ttfamily \mnem{BRANCH}\imod{.conservative} \src{r0}, \imm{\#+12}} (\textbf{SPIR-V:} \texttt{OpBranchConditional}; \textbf{GLSL:} \texttt{while (cond) \{ ... \}})
  \item \hyperref[enum:branch_negate_m]{\texttt{not}}: Boolean negation of the evaluated branch condition on \texttt{BRANCH}/\texttt{JUMP}/\texttt{JUMP\_EX} (inverts the boolean predicate, NOT bitwise \texttt{\textasciitilde{}src0}). \emph{Example:} {\scriptsize\ttfamily \mnem{BRANCH}\imod{.not} \src{k0}, \imm{\#0x10}} (\textbf{SPIR-V:} \texttt{OpLogicalNot + OpBranchConditional}; \textbf{GLSL:} \texttt{if (!cond) goto target;})
  \item \hyperref[enum:branch_combine_m]{\texttt{combine}}: Combines 16-bit halfword predicates (\texttt{h0}, \texttt{h1}, \texttt{and}, \texttt{lowbits}) before evaluating a branch. \emph{Example:} {\scriptsize\ttfamily \mnem{BRANCH}\imod{.and} \src{r0}, \imm{\#+8}} (\textbf{SPIR-V:} \texttt{OpBranchConditional (packed predicate)}; \textbf{GLSL:} \texttt{if (condA \&\& condB) \{ ... \}})
  \item \hyperref[enum:dest_width_m]{\texttt{dest\_width}}: Stage C destination lane write-mask (\texttt{none} = full 32-bit, \texttt{h0} = low 16-bit, \texttt{h1} = high 16-bit without bit-shifting). \emph{Example:} {\scriptsize\ttfamily \mnem{FADD.v2f16} \dst{r0}\imod{.h1}, \src{r1}, \src{r2}} (\textbf{SPIR-V:} \texttt{OpVectorInsertDynamic / WriteMask}; \textbf{GLSL:} \texttt{dst.y = a.y + b.y;})
  \item \texttt{opcode1}: Primary sub-opcode within the instruction class uniquely identifying the operation family. \emph{Example:} {\scriptsize\ttfamily \mnem{IADD.s32} \dst{r0}, \src{r1}, \src{r2}} (\textbf{SPIR-V:} \texttt{OpIAdd (operation selector)}; \textbf{GLSL:} \texttt{// Primary opcode bits})
  \item \texttt{instruction\_class}: 4-bit major instruction class selector (\texttt{bits 48:51}) partitioning the 64-bit Valhall opcode space (e.g., \texttt{FMA}, \texttt{CVT}, \texttt{SFU}, \texttt{LS}, \texttt{TEX}, \texttt{BRANCH}). \emph{Example:} {\scriptsize\ttfamily \mnem{FADD.f32} \dst{r0}, \src{r1}, \src{r2}} (\textbf{SPIR-V:} \texttt{N/A (encoding class)}; \textbf{GLSL:} \texttt{// Hardware execution pipe / decode class})
  \item \texttt{fau\_page\_index}: Selects the FAU (Fast Access Uniform) page (\texttt{0..3}) when a source operand references hardware small constants, thread IDs, or lane constants. \emph{Example:} {\scriptsize\ttfamily \mnem{FADD.f32} \dst{r0}, \src{r1}, \src{k1f}} (\textbf{SPIR-V:} \texttt{OpConstant / BuiltIn}; \textbf{GLSL:} \texttt{float y = x + 1.0;})
  \item \hyperref[enum:flow_control_m]{\texttt{flow}}: 4-bit embedded flow-control \& scoreboard wait modifier (\texttt{bits 59:62}: \texttt{none}, \texttt{wait0..wait0126}, \texttt{wait}, \texttt{reconverge}, \texttt{discard}, \texttt{end}). \emph{Example:} {\scriptsize\ttfamily \mnem{FADD.f32}\imod{.wait0.reconverge} \dst{r0}, \src{r1}, \src{r2}} (\textbf{SPIR-V:} \texttt{OpControlBarrier / OpReturn}; \textbf{GLSL:} \texttt{// Scoreboard wait + warp reconvergence})
\end{tightitemize}
\clearpage

\subsection*{KABOOM}
\addcontentsline{toc}{subsection}{KABOOM}
Triggers an immediate hardware shader fault/trap message to the CSF job manager; used as an unreachable-code or fatal-assertion sentinel.
\par\textbf{Notes:}
\begin{tightitemize}
  \item Triggers hardware fault/abort.
  \item Not a benign no-op.
  \item \textbf{Source references:} \texttt{bifrost/bifrost/ISA.xml:6225}, \texttt{isa-v9-v14.xml}.
\end{tightitemize}
\codebox{\mnem{KABOOM}\imod{\{.flow\}}}
\textbf{Operands:} \textit{None} (control/system instruction with no register or immediate operands; configured via modifiers and \texttt{flow}).
\par\textbf{Closest SPIR-V Op:} \texttt{OpUnreachable / hardware fault}
\par\textbf{Pseudo-GLSL (Operational Semantics):}
\codebox{gpu\_fault\_trap(KABOOM); // raise immediate shader fault}
\textbf{Execution Context:}
\begin{tightitemize}\footnotesize
  \item \textbf{Compute stage preloads (\hyperref[sec:csf-spd]{\texttt{SPD\_k}}, \texttt{r16..r23}):} Hardware pre-fills \texttt{r55..r57} with \texttt{WorkGroupID.xyz + JOB\_OFFSET} (from \hyperref[sec:csf-regs]{\texttt{COMPUTESR::JOB\_OFFSET\_X/Y/Z}} \texttt{r34..r36}) and \texttt{r60..r62} with \texttt{LocalInvocationID.xyz} or \texttt{GlobalInvocationID.xyz} (from \texttt{WG\_SIZE} \texttt{r33} \& \texttt{JOB\_SIZE} \texttt{r37..r39}, selected by \texttt{ShaderProgram::preload}) so compute shaders can derive thread/element indices without memory loads.
  \item \textbf{Vertex stage preloads (\hyperref[sec:csf-spd]{\texttt{SPD\_0}/\texttt{SPD\_1}}, \texttt{r16..r19}):} The hardware Index Fetcher reads \hyperref[sec:csf-regs]{\texttt{IDVSSR::INDEX\_BUFFER}} (\texttt{r54:r55}) and pre-fills \texttt{r59, r60} with \texttt{VertexIndex} (\texttt{+ VERTEX\_OFFSET + GLOBAL\_ATTR\_OFF}) and \texttt{r61} with \texttt{InstanceIndex} (\texttt{+ INSTANCE\_OFFSET}) so Position and Varying shaders can pull vertex attributes directly.
  \item \textbf{Inline FAU \& Special Operands (\hyperref[sec:csf-tsd-fau]{\texttt{FAU\_k}}, \texttt{r8..r15}):} Any 8-bit source with \texttt{mode=2} reads zero-latency uniform \texttt{u0..u31} from \texttt{FAU\_k} (\texttt{slots64 = bits[63:56]}); \texttt{mode=3} reads \texttt{SmallConstant} (Table~1.2) or \texttt{FauSpecial} (\texttt{pc}, \texttt{lane\_id}, \texttt{lane\_width}, \texttt{tls\_ptr}, \texttt{wls\_ptr}, \texttt{blend\_descriptor\_0..7}).
\end{tightitemize}
\par\textbf{Worked Examples:}
\par\textbf{Example 1 (Unreachable trap):}
\codebox{\mnem{KABOOM}}
Immediately raises a hardware shader fault (\texttt{hw\_trap(KABOOM)}); used as a poison/unreachable sentinel.
\par\textbf{Example 2 (Divergent abort guard):}
\codebox{\mnem{KABOOM}\imod{.end}}
Faults the shader job if execution ever reaches an impossible control-flow path.

\begin{center}\footnotesize
\begin{tabularx}{\linewidth}{@{}lllX@{}}
\toprule
Variant & Arch & Class & Opcode Fields \\
\midrule
\texttt{(none)} & v9--v14 & \texttt{N0} & \texttt{opcode1=KABOOM} \\
\bottomrule
\end{tabularx}
\end{center}
\begin{center}\footnotesize
\begin{tabularx}{0.88\linewidth}{@{}l l X@{}}
\toprule
Field Name & Bits & Type / Fixed Value \\
\midrule
\textit{Reserved} & \texttt{0:32} & \texttt{zero} \\
\texttt{sr\_count} & \texttt{33:35} & \texttt{exact=count0} \\
\textit{Reserved} & \texttt{36:39} & \texttt{zero} \\
\texttt{staging\_register\_index} & \texttt{40:45} & \texttt{exact=0} \\
\hyperref[enum:dest_width_m]{\texttt{dest\_width}} & \texttt{46:47} & \texttt{exact=none} \\
\texttt{opcode1} & \texttt{48:51} & \texttt{exact=KABOOM} \\
\texttt{instruction\_class} & \texttt{52:56} & \texttt{exact=N0} \\
\texttt{fau\_page\_index} & \texttt{57:58} & \texttt{exact=0} \\
\hyperref[enum:flow_control_m]{\texttt{flow}} & \texttt{59:62} & \texttt{flow\_control\_m} \\
\textit{Reserved} & \texttt{63} & \texttt{zero} \\
\bottomrule
\end{tabularx}
\end{center}
\textbf{Fields Explained:}
\begin{tightitemize}\footnotesize
  \item \texttt{sr\_count}: Number of consecutive 32-bit staging registers transferred by a load/store/texture message. \emph{Example:} {\scriptsize\ttfamily \mnem{LD\_MULTI.i128} \dst{@r0:r3}, \src{r4}, \imm{\#0}} (\textbf{SPIR-V:} \texttt{OpLoad (wide vector)}; \textbf{GLSL:} \texttt{uvec4 x = buf[i];})
  \item \texttt{staging\_register\_index}: Base GPR index of the contiguous staging register window (\texttt{@rN:rM}) for message payload/return. \emph{Example:} {\scriptsize\ttfamily \mnem{STORE.i64}\imod{.slot7} \src{@r0:r1}, \src{r2}, \imm{\#8}} (\textbf{SPIR-V:} \texttt{OpStore (uvec2)}; \textbf{GLSL:} \texttt{ssbo.pairs[i] = uvec2(a, b);})
  \item \hyperref[enum:dest_width_m]{\texttt{dest\_width}}: Stage C destination lane write-mask (\texttt{none} = full 32-bit, \texttt{h0} = low 16-bit, \texttt{h1} = high 16-bit without bit-shifting). \emph{Example:} {\scriptsize\ttfamily \mnem{FADD.v2f16} \dst{r0}\imod{.h1}, \src{r1}, \src{r2}} (\textbf{SPIR-V:} \texttt{OpVectorInsertDynamic / WriteMask}; \textbf{GLSL:} \texttt{dst.y = a.y + b.y;})
  \item \texttt{opcode1}: Primary sub-opcode within the instruction class uniquely identifying the operation family. \emph{Example:} {\scriptsize\ttfamily \mnem{IADD.s32} \dst{r0}, \src{r1}, \src{r2}} (\textbf{SPIR-V:} \texttt{OpIAdd (operation selector)}; \textbf{GLSL:} \texttt{// Primary opcode bits})
  \item \texttt{instruction\_class}: 4-bit major instruction class selector (\texttt{bits 48:51}) partitioning the 64-bit Valhall opcode space (e.g., \texttt{FMA}, \texttt{CVT}, \texttt{SFU}, \texttt{LS}, \texttt{TEX}, \texttt{BRANCH}). \emph{Example:} {\scriptsize\ttfamily \mnem{FADD.f32} \dst{r0}, \src{r1}, \src{r2}} (\textbf{SPIR-V:} \texttt{N/A (encoding class)}; \textbf{GLSL:} \texttt{// Hardware execution pipe / decode class})
  \item \texttt{fau\_page\_index}: Selects the FAU (Fast Access Uniform) page (\texttt{0..3}) when a source operand references hardware small constants, thread IDs, or lane constants. \emph{Example:} {\scriptsize\ttfamily \mnem{FADD.f32} \dst{r0}, \src{r1}, \src{k1f}} (\textbf{SPIR-V:} \texttt{OpConstant / BuiltIn}; \textbf{GLSL:} \texttt{float y = x + 1.0;})
  \item \hyperref[enum:flow_control_m]{\texttt{flow}}: 4-bit embedded flow-control \& scoreboard wait modifier (\texttt{bits 59:62}: \texttt{none}, \texttt{wait0..wait0126}, \texttt{wait}, \texttt{reconverge}, \texttt{discard}, \texttt{end}). \emph{Example:} {\scriptsize\ttfamily \mnem{FADD.f32}\imod{.wait0.reconverge} \dst{r0}, \src{r1}, \src{r2}} (\textbf{SPIR-V:} \texttt{OpControlBarrier / OpReturn}; \textbf{GLSL:} \texttt{// Scoreboard wait + warp reconvergence})
\end{tightitemize}
\clearpage

\subsection*{NOP}
\addcontentsline{toc}{subsection}{NOP}
Architectural no-operation that performs no arithmetic or memory access, used to carry \texttt{flow} control actions such as scoreboard waits (\texttt{.wait}, \texttt{.wait0126}), \texttt{.reconverge}, or warp termination (\texttt{.end}).
\par\textbf{Notes:}
\begin{tightitemize}
  \item No register or memory side effects.
  \item \textbf{Source references:} \texttt{kraid/ops.rs:3651-3661}, \texttt{encode\_v9.rs:2887-2895}, \texttt{bifrost/valhall/ISA.xml:135}.
\end{tightitemize}
\codebox{\mnem{NOP}\imod{\{.flow\}}}
\textbf{Operands:} \textit{None} (control/system instruction with no register or immediate operands; configured via modifiers and \texttt{flow}).
\par\textbf{Closest SPIR-V Op:} \texttt{OpNop}
\par\textbf{Pseudo-GLSL (Operational Semantics):}
\codebox{if (flow\_wait\_mask != 0u) scoreboardWait(flow\_wait\_mask);\\if (flow == end)          return; // terminate shader warp}
\textbf{Execution Context:}
\begin{tightitemize}\footnotesize
  \item \textbf{Compute stage preloads (\hyperref[sec:csf-spd]{\texttt{SPD\_k}}, \texttt{r16..r23}):} Hardware pre-fills \texttt{r55..r57} with \texttt{WorkGroupID.xyz + JOB\_OFFSET} (from \hyperref[sec:csf-regs]{\texttt{COMPUTESR::JOB\_OFFSET\_X/Y/Z}} \texttt{r34..r36}) and \texttt{r60..r62} with \texttt{LocalInvocationID.xyz} or \texttt{GlobalInvocationID.xyz} (from \texttt{WG\_SIZE} \texttt{r33} \& \texttt{JOB\_SIZE} \texttt{r37..r39}, selected by \texttt{ShaderProgram::preload}) so compute shaders can derive thread/element indices without memory loads.
  \item \textbf{Vertex stage preloads (\hyperref[sec:csf-spd]{\texttt{SPD\_0}/\texttt{SPD\_1}}, \texttt{r16..r19}):} The hardware Index Fetcher reads \hyperref[sec:csf-regs]{\texttt{IDVSSR::INDEX\_BUFFER}} (\texttt{r54:r55}) and pre-fills \texttt{r59, r60} with \texttt{VertexIndex} (\texttt{+ VERTEX\_OFFSET + GLOBAL\_ATTR\_OFF}) and \texttt{r61} with \texttt{InstanceIndex} (\texttt{+ INSTANCE\_OFFSET}) so Position and Varying shaders can pull vertex attributes directly.
  \item \textbf{Inline FAU \& Special Operands (\hyperref[sec:csf-tsd-fau]{\texttt{FAU\_k}}, \texttt{r8..r15}):} Any 8-bit source with \texttt{mode=2} reads zero-latency uniform \texttt{u0..u31} from \texttt{FAU\_k} (\texttt{slots64 = bits[63:56]}); \texttt{mode=3} reads \texttt{SmallConstant} (Table~1.2) or \texttt{FauSpecial} (\texttt{pc}, \texttt{lane\_id}, \texttt{lane\_width}, \texttt{tls\_ptr}, \texttt{wls\_ptr}, \texttt{blend\_descriptor\_0..7}).
\end{tightitemize}
\par\textbf{Worked Examples:}
\par\textbf{Example 1 (Multi-slot scoreboard wait):}
\codebox{\mnem{NOP}\imod{.wait012}}
Performs no register or memory update; stalls the warp until scoreboard slots 0, 1, and 2 all drain.
\par\textbf{Example 2 (Shader termination carrier (`.end`)):}
\codebox{\mnem{NOP}\imod{.wait.end}}
Drains all pending async stores/messages (\texttt{.wait}) and terminates the shader warp (\texttt{.end}).
\par\textbf{Mined SPIR-V $\to$ Valhall ISA Example (Explicit scoreboard barrier before ATEST --- Fragment shader, Mali-G710 (v10)):}
\codebox{\textbf{; SPIR-V (input)}\\; Fragment prologue before alpha/depth test\\[3pt]\textbf{; Valhall ISA (malioc output)}\\NOP.wait0126\\ATEST.slot0 @r60, r60, k0, atest\_datum.w0}
\texttt{NOP.wait0126} performs no computation but stalls warp execution until asynchronous scoreboard slots 0, 1, 2, and 6 drain --- required before issuing \texttt{ATEST} or \texttt{DISCARD}.
\begin{center}\footnotesize
\begin{tabularx}{\linewidth}{@{}lllX@{}}
\toprule
Variant & Arch & Class & Opcode Fields \\
\midrule
\texttt{(none)} & v9--v14 & \texttt{N0} & \texttt{opcode1=NOP} \\
\bottomrule
\end{tabularx}
\end{center}
\begin{center}\footnotesize
\begin{tabularx}{0.88\linewidth}{@{}l l X@{}}
\toprule
Field Name & Bits & Type / Fixed Value \\
\midrule
\textit{Reserved} & \texttt{0:45} & \texttt{zero} \\
\hyperref[enum:dest_width_m]{\texttt{dest\_width}} & \texttt{46:47} & \texttt{exact=none} \\
\texttt{opcode1} & \texttt{48:51} & \texttt{exact=NOP} \\
\texttt{instruction\_class} & \texttt{52:56} & \texttt{exact=N0} \\
\texttt{fau\_page\_index} & \texttt{57:58} & \texttt{exact=0} \\
\hyperref[enum:flow_control_m]{\texttt{flow}} & \texttt{59:62} & \texttt{flow\_control\_m} \\
\textit{Reserved} & \texttt{63} & \texttt{zero} \\
\bottomrule
\end{tabularx}
\end{center}
\textbf{Fields Explained:}
\begin{tightitemize}\footnotesize
  \item \hyperref[enum:dest_width_m]{\texttt{dest\_width}}: Stage C destination lane write-mask (\texttt{none} = full 32-bit, \texttt{h0} = low 16-bit, \texttt{h1} = high 16-bit without bit-shifting). \emph{Example:} {\scriptsize\ttfamily \mnem{FADD.v2f16} \dst{r0}\imod{.h1}, \src{r1}, \src{r2}} (\textbf{SPIR-V:} \texttt{OpVectorInsertDynamic / WriteMask}; \textbf{GLSL:} \texttt{dst.y = a.y + b.y;})
  \item \texttt{opcode1}: Primary sub-opcode within the instruction class uniquely identifying the operation family. \emph{Example:} {\scriptsize\ttfamily \mnem{IADD.s32} \dst{r0}, \src{r1}, \src{r2}} (\textbf{SPIR-V:} \texttt{OpIAdd (operation selector)}; \textbf{GLSL:} \texttt{// Primary opcode bits})
  \item \texttt{instruction\_class}: 4-bit major instruction class selector (\texttt{bits 48:51}) partitioning the 64-bit Valhall opcode space (e.g., \texttt{FMA}, \texttt{CVT}, \texttt{SFU}, \texttt{LS}, \texttt{TEX}, \texttt{BRANCH}). \emph{Example:} {\scriptsize\ttfamily \mnem{FADD.f32} \dst{r0}, \src{r1}, \src{r2}} (\textbf{SPIR-V:} \texttt{N/A (encoding class)}; \textbf{GLSL:} \texttt{// Hardware execution pipe / decode class})
  \item \texttt{fau\_page\_index}: Selects the FAU (Fast Access Uniform) page (\texttt{0..3}) when a source operand references hardware small constants, thread IDs, or lane constants. \emph{Example:} {\scriptsize\ttfamily \mnem{FADD.f32} \dst{r0}, \src{r1}, \src{k1f}} (\textbf{SPIR-V:} \texttt{OpConstant / BuiltIn}; \textbf{GLSL:} \texttt{float y = x + 1.0;})
  \item \hyperref[enum:flow_control_m]{\texttt{flow}}: 4-bit embedded flow-control \& scoreboard wait modifier (\texttt{bits 59:62}: \texttt{none}, \texttt{wait0..wait0126}, \texttt{wait}, \texttt{reconverge}, \texttt{discard}, \texttt{end}). \emph{Example:} {\scriptsize\ttfamily \mnem{FADD.f32}\imod{.wait0.reconverge} \dst{r0}, \src{r1}, \src{r2}} (\textbf{SPIR-V:} \texttt{OpControlBarrier / OpReturn}; \textbf{GLSL:} \texttt{// Scoreboard wait + warp reconvergence})
\end{tightitemize}
\clearpage

\section{Memory Load / Store \& LEA (19 mnemonics)}
\subsection*{LD\_CHECKED}
\addcontentsline{toc}{subsection}{LD\_CHECKED}
Bounds-checked load (v12+): loads \texttt{i8}..\texttt{i128} from base \texttt{src0} + index/offset \texttt{src1} checked against bound \texttt{src2}; on OOB, returns 0 (or \texttt{oob\_mode} behavior on v14).
\par\textbf{Notes:}
\begin{tightitemize}
  \item Robust buffer access hardware instruction.
  \item \textbf{Source references:} \texttt{isa-v9-v14.xml}.
\end{tightitemize}
\codebox{\mnem{LD\_CHECKED}\imod{\{.message\_slot\_index\}}\imod{\{.flow\}}\imod{\{.access\}} \dst{@sw}, \src{src0}, \src{src1}, \src{src2}\\\mnem{LD\_CHECKED}\imod{\{.message\_slot\_index\}}\imod{\{.flow\}}\imod{\{.oob\_mode\}}\imod{\{.access\}} \dst{@sw}, \src{src0}, \src{src1}, \src{src2}\\\mnem{LD\_CHECKED}\imod{\{.message\_slot\_index\}}\imod{\{.flow\}}\imod{\{.access\}}\imod{\{.lane\}} \dst{@sw}, \src{src0}, \src{src1}, \src{src2}\\\mnem{LD\_CHECKED}\imod{\{.message\_slot\_index\}}\imod{\{.flow\}}\imod{\{.oob\_mode\}}\imod{\{.access\}}\imod{\{.lane\}} \dst{@sw}, \src{src0}, \src{src1}, \src{src2}}
\textbf{Operands:}
\begin{tightitemize}\footnotesize
  \item \texttt{\dst{@sw}} (\texttt{staging\_register\_index} + \texttt{sr\_count}): Destination staging register range receiving loaded data (or \texttt{0} / unchanged on OOB).
  \item \texttt{\src{src0}}: 64-bit base GPU virtual address.
  \item \texttt{\src{src1}}: Dynamic unsigned 32-bit byte offset.
  \item \texttt{\src{src2}}: 32-bit buffer byte-size upper bound (\texttt{src1 + access\_bytes \textless{}= src2}).
\end{tightitemize}
\par\textbf{Closest SPIR-V Op:} \texttt{OpLoad (robustBufferAccess, v12+)}
\par\textbf{Pseudo-GLSL (Operational Semantics):}
\codebox{if (uint64\_t(src1) + access\_bytes \textless{}= uint64\_t(src2)) @sw = mem\_read(src0 + src1, width);\\else if (!flag\_oob\_ignore)                           @sw = 0;}
\textbf{Execution Context:}
\begin{tightitemize}\footnotesize
  \item \textbf{Hardware bounds check (v12+):} Chases \hyperref[sec:csf-srt]{\texttt{SRT\_k}} / \hyperref[sec:csf-tsd-fau]{\texttt{FAU\_k}} base VA and enforces \texttt{offset + width <= Buffer::size} (or bound register); OOB loads return \texttt{0} (unless v14 \texttt{.oob\_ignore}) and OOB stores drop silently.
  \item \textbf{Inline FAU \& Special Operands (\hyperref[sec:csf-tsd-fau]{\texttt{FAU\_k}}, \texttt{r8..r15}):} Any 8-bit source with \texttt{mode=2} reads zero-latency uniform \texttt{u0..u31} from \texttt{FAU\_k} (\texttt{slots64 = bits[63:56]}); \texttt{mode=3} reads \texttt{SmallConstant} (Table~1.2) or \texttt{FauSpecial} (\texttt{pc}, \texttt{lane\_id}, \texttt{lane\_width}, \texttt{tls\_ptr}, \texttt{wls\_ptr}, \texttt{blend\_descriptor\_0..7}).
\end{tightitemize}
\par\textbf{Worked Examples:}
\par\textbf{Example 1 (In-bounds hardware-checked load (v12+)):}
\codebox{\mnem{LD\_CHECKED.i32}\imod{.slot0} \dst{@r12}, \src{u3}, \src{r47\textasciicircum{}}, \src{k1n}}
Checks byte offset \texttt{r47 \textless{} bound}: when in-bounds, loads 4 bytes from base \texttt{u3 + r47} into \texttt{@r12}.
\par\textbf{Example 2 (Out-of-bounds zero return):}
\codebox{\mnem{LD\_CHECKED.i32}\imod{.slot0} \dst{@r12}, \src{u3}, \src{r47}, \src{r48}}
When \texttt{r47 \textgreater{}= r48} (with default \texttt{oob\_mode=none}): suppresses the memory transaction and writes \texttt{@r12 = 0}.
\par\textbf{Mined SPIR-V $\to$ Valhall ISA Example (Robust bounds-checked SSBO load on Immortalis-G925 (v14) --- Compute shader, Immortalis-G925 (v14)):}
\codebox{\textbf{; SPIR-V (input)}\\\%idx = OpIAdd \%uint (OpShiftLeftLogical \%uint \%wg\_x \%uint\_5) \%local\_x\\\%ptr = OpAccessChain \%\_ptr\_Uniform\_float \%ssbo \%int\_0 \%idx\\\%val = OpLoad \%float \%ptr\\[3pt]\textbf{; Valhall ISA (malioc output)}\\LSHIFT.i32 r0, r57\textasciicircum{}, k0B0A0908.b1\\IADD.u32 r6, r0, k0F0E0D0C.b0\\LD\_CHECKED.i32.slot0.wait0 @r7, u7, r6, k1n}
On Valhall v12+ (Immortalis-G925), \texttt{malioc} lowers robust Vulkan SSBO loads to \texttt{LD\_CHECKED.i32}, which checks the byte offset \texttt{r6} against the buffer bound in hardware and returns \texttt{0} on out-of-bounds access without a branch.
\begin{center}\footnotesize
\begin{tabularx}{\linewidth}{@{}lllX@{}}
\toprule
Variant & Arch & Class & Opcode Fields \\
\midrule
\texttt{(none)} & v9--v14 & \texttt{S4} & \texttt{opcode1=LD\_CHECKED} \\
\texttt{(none)} & v9--v14 & \texttt{S4} & \texttt{opcode1=LD\_CHECKED} \\
\texttt{(none)} & v9--v14 & \texttt{S4} & \texttt{opcode1=LD\_CHECKED} \\
\texttt{(none)} & v9--v14 & \texttt{S4} & \texttt{opcode1=LD\_CHECKED} \\
\texttt{(none)} & v9--v14 & \texttt{S4} & \texttt{opcode1=LD\_CHECKED} \\
\texttt{(none)} & v9--v14 & \texttt{S4} & \texttt{opcode1=LD\_CHECKED} \\
\texttt{(none)} & v9--v14 & \texttt{S4} & \texttt{opcode1=LD\_CHECKED} \\
\texttt{(none)} & v9--v14 & \texttt{S4} & \texttt{opcode1=LD\_CHECKED} \\
\texttt{(none)} & v9--v14 & \texttt{S4} & \texttt{opcode1=LD\_CHECKED} \\
\texttt{(none)} & v9--v14 & \texttt{S4} & \texttt{opcode1=LD\_CHECKED} \\
\texttt{(none)} & v9--v14 & \texttt{S4} & \texttt{opcode1=LD\_CHECKED} \\
\texttt{(none)} & v9--v14 & \texttt{S4} & \texttt{opcode1=LD\_CHECKED} \\
\texttt{(none)} & v9--v14 & \texttt{S4} & \texttt{opcode1=LD\_CHECKED} \\
\texttt{(none)} & v9--v14 & \texttt{S4} & \texttt{opcode1=LD\_CHECKED} \\
\texttt{(none)} & v9--v14 & \texttt{S4} & \texttt{opcode1=LD\_CHECKED} \\
\texttt{(none)} & v9--v14 & \texttt{S4} & \texttt{opcode1=LD\_CHECKED} \\
\bottomrule
\end{tabularx}
\end{center}
\begin{center}\footnotesize
\begin{tabularx}{0.88\linewidth}{@{}l l X@{}}
\toprule
Field Name & Bits & Type / Fixed Value \\
\midrule
\texttt{src0} & \texttt{0:7} & \texttt{SourceEncoding64} \\
\texttt{src1} & \texttt{8:15} & \texttt{SourceEncoding} \\
\texttt{src2} & \texttt{16:23} & \texttt{SourceEncoding} \\
\hyperref[enum:access_load_m]{\texttt{access}} & \texttt{24:26} & \texttt{access\_load\_m (def none)} \\
\texttt{memory\_width} & \texttt{27:29} & \texttt{exact=i128} \\
\hyperref[enum:message_slot_index_m]{\texttt{message\_slot\_index}} & \texttt{30:32} & \texttt{message\_slot\_index\_m} \\
\texttt{sr\_count} & \texttt{33:35} & \texttt{exact=count4} \\
\textit{Reserved} & \texttt{36:39} & \texttt{zero} \\
\texttt{staging\_register\_index} & \texttt{40:45} & \texttt{uint} \\
\texttt{staging\_register\_control\_32} & \texttt{46:47} & \texttt{exact=write} \\
\texttt{opcode1} & \texttt{48:51} & \texttt{exact=LD\_CHECKED} \\
\texttt{instruction\_class} & \texttt{52:56} & \texttt{exact=S4} \\
\texttt{fau\_page\_index} & \texttt{57:58} & \texttt{uint} \\
\hyperref[enum:flow_control_m]{\texttt{flow}} & \texttt{59:62} & \texttt{flow\_control\_m} \\
\textit{Reserved} & \texttt{63} & \texttt{zero} \\
\bottomrule
\end{tabularx}
\end{center}
\textbf{Fields Explained:}
\begin{tightitemize}\footnotesize
  \item \texttt{src0}: Primary source operand (\texttt{r0..r63}, \texttt{u0..u63} uniform, \texttt{k*} FAU constant, or \texttt{t0..t2} temporary; \texttt{\textasciicircum{}} marks discard/last-use). \emph{Example:} {\scriptsize\ttfamily \mnem{FADD.f32} \dst{r0}, \src{r1}, \src{r2}} (\textbf{SPIR-V:} \texttt{OpFAdd (Operand 0)}; \textbf{GLSL:} \texttt{float c = a + b;})
  \item \texttt{src1}: Secondary source operand (register, uniform, or FAU constant; in shifts, supplies the shift count byte). \emph{Example:} {\scriptsize\ttfamily \mnem{LSHIFT.i32} \dst{r0}, \src{r1}, \src{r2}\imod{.b0}} (\textbf{SPIR-V:} \texttt{OpShiftLeftLogical (Shift)}; \textbf{GLSL:} \texttt{uint y = a \textless{}\textless{} s;})
  \item \texttt{src2}: Third source operand (FMA addend, \texttt{MUX}/\texttt{CSEL} selector, \texttt{IADDC} carry-in, or bitwise mask in fused shift-logic). \emph{Example:} {\scriptsize\ttfamily \mnem{FMA.f32} \dst{r0}, \src{r1}, \src{r2}, \src{r3}} (\textbf{SPIR-V:} \texttt{GLSL.std.450 Fma (Addend)}; \textbf{GLSL:} \texttt{float d = fma(a, b, c);})
  \item \hyperref[enum:access_load_m]{\texttt{access}}: Memory cache streaming/coherence policy (\texttt{none}, \texttt{istream}, \texttt{estream}, \texttt{force}). \emph{Example:} {\scriptsize\ttfamily \mnem{STORE.i32}\imod{.force} \src{@r0}, \src{r2}, \imm{\#0}} (\textbf{SPIR-V:} \texttt{OpStore (MakePointerVisible / Coherent)}; \textbf{GLSL:} \texttt{coherent buffer B \{ uint x; \}; x = 1u;})
  \item \texttt{memory\_width}: Memory transfer width (\texttt{i8}, \texttt{i16}, \texttt{i24}, \texttt{i32}, \texttt{i48}, \texttt{i64}, \texttt{i96}, \texttt{i128}). \emph{Example:} {\scriptsize\ttfamily \mnem{LOAD.i64} \dst{@r0:r1}, \src{[r2:r3]}, \imm{\#0}} (\textbf{SPIR-V:} \texttt{OpLoad (uvec2)}; \textbf{GLSL:} \texttt{uvec2 v = ssbo.pairs[i];})
  \item \hyperref[enum:message_slot_index_m]{\texttt{message\_slot\_index}}: Asynchronous scoreboard slot (\texttt{slot0}--\texttt{slot2}, \texttt{slot7}) tracking completion of this message. \emph{Example:} {\scriptsize\ttfamily \mnem{LOAD.i32}\imod{.slot0} \dst{@r0}, \src{r2}, \imm{\#0}} (\textbf{SPIR-V:} \texttt{OpLoad}; \textbf{GLSL:} \texttt{uint v = ssbo.arr[i];})
  \item \texttt{sr\_count}: Number of consecutive 32-bit staging registers transferred by a load/store/texture message. \emph{Example:} {\scriptsize\ttfamily \mnem{LD\_MULTI.i128} \dst{@r0:r3}, \src{r4}, \imm{\#0}} (\textbf{SPIR-V:} \texttt{OpLoad (wide vector)}; \textbf{GLSL:} \texttt{uvec4 x = buf[i];})
  \item \texttt{staging\_register\_index}: Base GPR index of the contiguous staging register window (\texttt{@rN:rM}) for message payload/return. \emph{Example:} {\scriptsize\ttfamily \mnem{STORE.i64}\imod{.slot7} \src{@r0:r1}, \src{r2}, \imm{\#8}} (\textbf{SPIR-V:} \texttt{OpStore (uvec2)}; \textbf{GLSL:} \texttt{ssbo.pairs[i] = uvec2(a, b);})
  \item \texttt{staging\_register\_control\_32}: Controls the read/write staging register bank and count for async message instructions. \emph{Example:} {\scriptsize\ttfamily \mnem{LOAD.i128}\imod{.slot0} \dst{@r0:r3}, \src{r4}, \imm{\#0}} (\textbf{SPIR-V:} \texttt{OpLoad (vec4)}; \textbf{GLSL:} \texttt{vec4 v = ssbo.data[idx];})
  \item \texttt{opcode1}: Primary sub-opcode within the instruction class uniquely identifying the operation family. \emph{Example:} {\scriptsize\ttfamily \mnem{IADD.s32} \dst{r0}, \src{r1}, \src{r2}} (\textbf{SPIR-V:} \texttt{OpIAdd (operation selector)}; \textbf{GLSL:} \texttt{// Primary opcode bits})
  \item \texttt{instruction\_class}: 4-bit major instruction class selector (\texttt{bits 48:51}) partitioning the 64-bit Valhall opcode space (e.g., \texttt{FMA}, \texttt{CVT}, \texttt{SFU}, \texttt{LS}, \texttt{TEX}, \texttt{BRANCH}). \emph{Example:} {\scriptsize\ttfamily \mnem{FADD.f32} \dst{r0}, \src{r1}, \src{r2}} (\textbf{SPIR-V:} \texttt{N/A (encoding class)}; \textbf{GLSL:} \texttt{// Hardware execution pipe / decode class})
  \item \texttt{fau\_page\_index}: Selects the FAU (Fast Access Uniform) page (\texttt{0..3}) when a source operand references hardware small constants, thread IDs, or lane constants. \emph{Example:} {\scriptsize\ttfamily \mnem{FADD.f32} \dst{r0}, \src{r1}, \src{k1f}} (\textbf{SPIR-V:} \texttt{OpConstant / BuiltIn}; \textbf{GLSL:} \texttt{float y = x + 1.0;})
  \item \hyperref[enum:flow_control_m]{\texttt{flow}}: 4-bit embedded flow-control \& scoreboard wait modifier (\texttt{bits 59:62}: \texttt{none}, \texttt{wait0..wait0126}, \texttt{wait}, \texttt{reconverge}, \texttt{discard}, \texttt{end}). \emph{Example:} {\scriptsize\ttfamily \mnem{FADD.f32}\imod{.wait0.reconverge} \dst{r0}, \src{r1}, \src{r2}} (\textbf{SPIR-V:} \texttt{OpControlBarrier / OpReturn}; \textbf{GLSL:} \texttt{// Scoreboard wait + warp reconvergence})
\end{tightitemize}
\clearpage

\subsection*{LD\_CHECKED\_IMM}
\addcontentsline{toc}{subsection}{LD\_CHECKED\_IMM}
Bounds-checked memory load (v12+) at 64-bit base address \texttt{src0} plus immediate unsigned byte offset \texttt{\#offset} (\texttt{offset\_lo} + \texttt{offset\_hi}): if \texttt{\#offset + access\_bytes \textless{}= src2}, loads into \texttt{@sw}; otherwise returns \texttt{0} (or leaves \texttt{@sw} unchanged when \texttt{.oob\_ignore} is set on v14).
\par\textbf{Notes:}
\begin{tightitemize}
  \item Robust buffer access with 12-bit unsigned split immediate offset \texttt{((offset\_hi \& 0xF) \textless{}\textless{} 8) \textbar{} (offset\_lo \& 0xFF)} and full-width bounds check \texttt{offset + access\_bytes \textless{}= bound}.
  \item \textbf{Source references:} \texttt{isa-v9-v14.xml}.
\end{tightitemize}
\codebox{\mnem{LD\_CHECKED\_IMM}\imod{\{.message\_slot\_index\}}\imod{\{.flow\}}\imod{\{.access\}} \dst{@sw}, \src{src0}, \imm{\#offset}, \src{src2}\\\mnem{LD\_CHECKED\_IMM}\imod{\{.message\_slot\_index\}}\imod{\{.flow\}}\imod{\{.oob\_mode\}}\imod{\{.access\}} \dst{@sw}, \src{src0}, \imm{\#offset}, \src{src2}\\\mnem{LD\_CHECKED\_IMM}\imod{\{.message\_slot\_index\}}\imod{\{.flow\}}\imod{\{.access\}}\imod{\{.lane\}} \dst{@sw}, \src{src0}, \imm{\#offset}, \src{src2}\\\mnem{LD\_CHECKED\_IMM}\imod{\{.message\_slot\_index\}}\imod{\{.flow\}}\imod{\{.oob\_mode\}}\imod{\{.access\}}\imod{\{.lane\}} \dst{@sw}, \src{src0}, \imm{\#offset}, \src{src2}}
\textbf{Operands:}
\begin{tightitemize}\footnotesize
  \item \texttt{\dst{@sw}} (\texttt{staging\_register\_index} + \texttt{sr\_count}): Destination staging register range receiving loaded data (or \texttt{0} / unchanged on OOB).
  \item \texttt{\src{src0}}: 64-bit base GPU virtual address.
  \item \texttt{\imm{\#offset}} (\texttt{offset\_lo} + \texttt{offset\_hi}): Immediate unsigned byte offset (assembled from \texttt{offset\_lo} and \texttt{offset\_hi}).
  \item \texttt{\src{src2}}: 32-bit buffer byte-size upper bound (\texttt{\#offset + access\_bytes \textless{}= src2}).
\end{tightitemize}
\par\textbf{Closest SPIR-V Op:} \texttt{OpLoad (robustBufferAccess imm, v12+)}
\par\textbf{Pseudo-GLSL (Operational Semantics):}
\codebox{if (uint64\_t(\#offset) + access\_bytes \textless{}= uint64\_t(src2)) @sw = mem\_read(src0 + \#offset, width);\\else if (!flag\_oob\_ignore)                              @sw = 0;}
\textbf{Execution Context:}
\begin{tightitemize}\footnotesize
  \item \textbf{Hardware bounds check (v12+):} Chases \hyperref[sec:csf-srt]{\texttt{SRT\_k}} / \hyperref[sec:csf-tsd-fau]{\texttt{FAU\_k}} base VA and enforces \texttt{offset + width <= Buffer::size} (or bound register); OOB loads return \texttt{0} (unless v14 \texttt{.oob\_ignore}) and OOB stores drop silently.
  \item \textbf{Inline FAU \& Special Operands (\hyperref[sec:csf-tsd-fau]{\texttt{FAU\_k}}, \texttt{r8..r15}):} Any 8-bit source with \texttt{mode=2} reads zero-latency uniform \texttt{u0..u31} from \texttt{FAU\_k} (\texttt{slots64 = bits[63:56]}); \texttt{mode=3} reads \texttt{SmallConstant} (Table~1.2) or \texttt{FauSpecial} (\texttt{pc}, \texttt{lane\_id}, \texttt{lane\_width}, \texttt{tls\_ptr}, \texttt{wls\_ptr}, \texttt{blend\_descriptor\_0..7}).
\end{tightitemize}
\par\textbf{Worked Examples:}
\par\textbf{Example 1 (Immediate-offset bounds-checked 64-bit load):}
\codebox{\mnem{LD\_CHECKED\_IMM.i64}\imod{.slot1} \dst{@r62:r63}, \src{[r8:r9]}, \src{r10}, \imm{\#0x20}}
Checks \texttt{0x20 + 8 \textless{}= r10}: if in-bounds, loads 8 bytes at \texttt{[r8:r9] + 0x20} into \texttt{@r62:r63} on \texttt{slot1}.
\par\textbf{Example 2 (Out-of-bounds zero fill):}
\codebox{\mnem{LD\_CHECKED\_IMM.i64}\imod{.slot1} \dst{@r62:r63}, \src{[r8:r9]}, \src{r10}, \imm{\#0x20}}
If buffer size \texttt{r10 = 0x10} (\texttt{\textless{} 0x20}): suppresses load and zeroes both \texttt{@r62 = 0, @r63 = 0}.

\begin{center}\footnotesize
\begin{tabularx}{\linewidth}{@{}lllX@{}}
\toprule
Variant & Arch & Class & Opcode Fields \\
\midrule
\texttt{(none)} & v9--v14 & \texttt{S3} & \texttt{opcode1=LD\_CHECKED\_IMM} \\
\texttt{(none)} & v9--v14 & \texttt{S3} & \texttt{opcode1=LD\_CHECKED\_IMM} \\
\texttt{(none)} & v9--v14 & \texttt{S3} & \texttt{opcode1=LD\_CHECKED\_IMM} \\
\texttt{(none)} & v9--v14 & \texttt{S3} & \texttt{opcode1=LD\_CHECKED\_IMM} \\
\texttt{(none)} & v9--v14 & \texttt{S3} & \texttt{opcode1=LD\_CHECKED\_IMM} \\
\texttt{(none)} & v9--v14 & \texttt{S3} & \texttt{opcode1=LD\_CHECKED\_IMM} \\
\texttt{(none)} & v9--v14 & \texttt{S3} & \texttt{opcode1=LD\_CHECKED\_IMM} \\
\texttt{(none)} & v9--v14 & \texttt{S3} & \texttt{opcode1=LD\_CHECKED\_IMM} \\
\texttt{(none)} & v9--v14 & \texttt{S3} & \texttt{opcode1=LD\_CHECKED\_IMM} \\
\texttt{(none)} & v9--v14 & \texttt{S3} & \texttt{opcode1=LD\_CHECKED\_IMM} \\
\texttt{(none)} & v9--v14 & \texttt{S3} & \texttt{opcode1=LD\_CHECKED\_IMM} \\
\texttt{(none)} & v9--v14 & \texttt{S3} & \texttt{opcode1=LD\_CHECKED\_IMM} \\
\texttt{(none)} & v9--v14 & \texttt{S3} & \texttt{opcode1=LD\_CHECKED\_IMM} \\
\texttt{(none)} & v9--v14 & \texttt{S3} & \texttt{opcode1=LD\_CHECKED\_IMM} \\
\texttt{(none)} & v9--v14 & \texttt{S3} & \texttt{opcode1=LD\_CHECKED\_IMM} \\
\texttt{(none)} & v9--v14 & \texttt{S3} & \texttt{opcode1=LD\_CHECKED\_IMM} \\
\bottomrule
\end{tabularx}
\end{center}
\begin{center}\footnotesize
\begin{tabularx}{0.88\linewidth}{@{}l l X@{}}
\toprule
Field Name & Bits & Type / Fixed Value \\
\midrule
\texttt{src0} & \texttt{0:7} & \texttt{SourceEncoding64} \\
\texttt{offset\_lo} & \texttt{8:15} & \texttt{uint} \\
\texttt{src2} & \texttt{16:23} & \texttt{SourceEncoding} \\
\hyperref[enum:access_load_m]{\texttt{access}} & \texttt{24:26} & \texttt{access\_load\_m (def none)} \\
\texttt{memory\_width} & \texttt{27:29} & \texttt{exact=i128} \\
\hyperref[enum:message_slot_index_m]{\texttt{message\_slot\_index}} & \texttt{30:32} & \texttt{message\_slot\_index\_m} \\
\texttt{sr\_count} & \texttt{33:35} & \texttt{exact=count4} \\
\texttt{offset\_hi} & \texttt{36:39} & \texttt{uint} \\
\texttt{staging\_register\_index} & \texttt{40:45} & \texttt{uint} \\
\texttt{staging\_register\_control\_32} & \texttt{46:47} & \texttt{exact=write} \\
\texttt{opcode1} & \texttt{48:51} & \texttt{exact=LD\_CHECKED\_IMM} \\
\texttt{instruction\_class} & \texttt{52:56} & \texttt{exact=S3} \\
\texttt{fau\_page\_index} & \texttt{57:58} & \texttt{uint} \\
\hyperref[enum:flow_control_m]{\texttt{flow}} & \texttt{59:62} & \texttt{flow\_control\_m} \\
\textit{Reserved} & \texttt{63} & \texttt{zero} \\
\bottomrule
\end{tabularx}
\end{center}
\textbf{Fields Explained:}
\begin{tightitemize}\footnotesize
  \item \texttt{src0}: Primary source operand (\texttt{r0..r63}, \texttt{u0..u63} uniform, \texttt{k*} FAU constant, or \texttt{t0..t2} temporary; \texttt{\textasciicircum{}} marks discard/last-use). \emph{Example:} {\scriptsize\ttfamily \mnem{FADD.f32} \dst{r0}, \src{r1}, \src{r2}} (\textbf{SPIR-V:} \texttt{OpFAdd (Operand 0)}; \textbf{GLSL:} \texttt{float c = a + b;})
  \item \texttt{offset\_lo}: Signed byte offset added to base address register in memory load/store/atomic instructions. \emph{Example:} {\scriptsize\ttfamily \mnem{LOAD.i32} \dst{@r0}, \src{[r2:r3]}, \imm{\#16}} (\textbf{SPIR-V:} \texttt{OpAccessChain + OpLoad}; \textbf{GLSL:} \texttt{uint y = ssbo.arr[4];})
  \item \texttt{src2}: Third source operand (FMA addend, \texttt{MUX}/\texttt{CSEL} selector, \texttt{IADDC} carry-in, or bitwise mask in fused shift-logic). \emph{Example:} {\scriptsize\ttfamily \mnem{FMA.f32} \dst{r0}, \src{r1}, \src{r2}, \src{r3}} (\textbf{SPIR-V:} \texttt{GLSL.std.450 Fma (Addend)}; \textbf{GLSL:} \texttt{float d = fma(a, b, c);})
  \item \hyperref[enum:access_load_m]{\texttt{access}}: Memory cache streaming/coherence policy (\texttt{none}, \texttt{istream}, \texttt{estream}, \texttt{force}). \emph{Example:} {\scriptsize\ttfamily \mnem{STORE.i32}\imod{.force} \src{@r0}, \src{r2}, \imm{\#0}} (\textbf{SPIR-V:} \texttt{OpStore (MakePointerVisible / Coherent)}; \textbf{GLSL:} \texttt{coherent buffer B \{ uint x; \}; x = 1u;})
  \item \texttt{memory\_width}: Memory transfer width (\texttt{i8}, \texttt{i16}, \texttt{i24}, \texttt{i32}, \texttt{i48}, \texttt{i64}, \texttt{i96}, \texttt{i128}). \emph{Example:} {\scriptsize\ttfamily \mnem{LOAD.i64} \dst{@r0:r1}, \src{[r2:r3]}, \imm{\#0}} (\textbf{SPIR-V:} \texttt{OpLoad (uvec2)}; \textbf{GLSL:} \texttt{uvec2 v = ssbo.pairs[i];})
  \item \hyperref[enum:message_slot_index_m]{\texttt{message\_slot\_index}}: Asynchronous scoreboard slot (\texttt{slot0}--\texttt{slot2}, \texttt{slot7}) tracking completion of this message. \emph{Example:} {\scriptsize\ttfamily \mnem{LOAD.i32}\imod{.slot0} \dst{@r0}, \src{r2}, \imm{\#0}} (\textbf{SPIR-V:} \texttt{OpLoad}; \textbf{GLSL:} \texttt{uint v = ssbo.arr[i];})
  \item \texttt{sr\_count}: Number of consecutive 32-bit staging registers transferred by a load/store/texture message. \emph{Example:} {\scriptsize\ttfamily \mnem{LD\_MULTI.i128} \dst{@r0:r3}, \src{r4}, \imm{\#0}} (\textbf{SPIR-V:} \texttt{OpLoad (wide vector)}; \textbf{GLSL:} \texttt{uvec4 x = buf[i];})
  \item \texttt{offset\_hi}: Signed byte offset added to base address register in memory load/store/atomic instructions. \emph{Example:} {\scriptsize\ttfamily \mnem{LOAD.i32} \dst{@r0}, \src{[r2:r3]}, \imm{\#16}} (\textbf{SPIR-V:} \texttt{OpAccessChain + OpLoad}; \textbf{GLSL:} \texttt{uint y = ssbo.arr[4];})
  \item \texttt{staging\_register\_index}: Base GPR index of the contiguous staging register window (\texttt{@rN:rM}) for message payload/return. \emph{Example:} {\scriptsize\ttfamily \mnem{STORE.i64}\imod{.slot7} \src{@r0:r1}, \src{r2}, \imm{\#8}} (\textbf{SPIR-V:} \texttt{OpStore (uvec2)}; \textbf{GLSL:} \texttt{ssbo.pairs[i] = uvec2(a, b);})
  \item \texttt{staging\_register\_control\_32}: Controls the read/write staging register bank and count for async message instructions. \emph{Example:} {\scriptsize\ttfamily \mnem{LOAD.i128}\imod{.slot0} \dst{@r0:r3}, \src{r4}, \imm{\#0}} (\textbf{SPIR-V:} \texttt{OpLoad (vec4)}; \textbf{GLSL:} \texttt{vec4 v = ssbo.data[idx];})
  \item \texttt{opcode1}: Primary sub-opcode within the instruction class uniquely identifying the operation family. \emph{Example:} {\scriptsize\ttfamily \mnem{IADD.s32} \dst{r0}, \src{r1}, \src{r2}} (\textbf{SPIR-V:} \texttt{OpIAdd (operation selector)}; \textbf{GLSL:} \texttt{// Primary opcode bits})
  \item \texttt{instruction\_class}: 4-bit major instruction class selector (\texttt{bits 48:51}) partitioning the 64-bit Valhall opcode space (e.g., \texttt{FMA}, \texttt{CVT}, \texttt{SFU}, \texttt{LS}, \texttt{TEX}, \texttt{BRANCH}). \emph{Example:} {\scriptsize\ttfamily \mnem{FADD.f32} \dst{r0}, \src{r1}, \src{r2}} (\textbf{SPIR-V:} \texttt{N/A (encoding class)}; \textbf{GLSL:} \texttt{// Hardware execution pipe / decode class})
  \item \texttt{fau\_page\_index}: Selects the FAU (Fast Access Uniform) page (\texttt{0..3}) when a source operand references hardware small constants, thread IDs, or lane constants. \emph{Example:} {\scriptsize\ttfamily \mnem{FADD.f32} \dst{r0}, \src{r1}, \src{k1f}} (\textbf{SPIR-V:} \texttt{OpConstant / BuiltIn}; \textbf{GLSL:} \texttt{float y = x + 1.0;})
  \item \hyperref[enum:flow_control_m]{\texttt{flow}}: 4-bit embedded flow-control \& scoreboard wait modifier (\texttt{bits 59:62}: \texttt{none}, \texttt{wait0..wait0126}, \texttt{wait}, \texttt{reconverge}, \texttt{discard}, \texttt{end}). \emph{Example:} {\scriptsize\ttfamily \mnem{FADD.f32}\imod{.wait0.reconverge} \dst{r0}, \src{r1}, \src{r2}} (\textbf{SPIR-V:} \texttt{OpControlBarrier / OpReturn}; \textbf{GLSL:} \texttt{// Scoreboard wait + warp reconvergence})
\end{tightitemize}
\clearpage

\subsection*{LD\_CVT}
\addcontentsline{toc}{subsection}{LD\_CVT}
Load with format conversion: reads from \texttt{src0 + offset} and converts memory pixel/vertex format using internal conversion descriptor \texttt{src2} into \texttt{vecsize} components of \texttt{register\_format}.
\par\textbf{Notes:}
\begin{tightitemize}
  \item Paired with \texttt{LEA\_TEX}/\texttt{LEA\_ATTR} (\texttt{ISA.xml:70}).
  \item NOT a raw byte load!
  \item \textbf{Source references:} \texttt{kraid/ops.rs:2707-2734}, \texttt{encode\_v9.rs:2065-2085}, \texttt{bifrost/valhall/ISA.xml:70}.
\end{tightitemize}
\codebox{\mnem{LD\_CVT}\imod{\{.message\_slot\_index\}}\imod{\{.flow\}}\imod{\{.access\}}\imod{\{.vecsize\}}\imod{\{.register\_format\}} \dst{@sw}, \src{src0}, \imm{\#offset}, \src{src2}}
\textbf{Operands:}
\begin{tightitemize}\footnotesize
  \item \texttt{\dst{@sw}} (\texttt{staging\_register\_index} + \texttt{sr\_count}): Destination staging register range receiving format-converted texel/attribute components.
  \item \texttt{\src{src0}}: 64-bit base GPU virtual address (typically from \texttt{LEA\_TEX}/\texttt{LEA\_ATTR}).
  \item \texttt{\imm{\#offset}} (\texttt{offset}): Unsigned byte offset added to \texttt{src0}.
  \item \texttt{\src{src2}}: 32-bit hardware format conversion descriptor (typically \texttt{r2} from \texttt{LEA\_TEX}/\texttt{LEA\_ATTR}).
\end{tightitemize}
\par\textbf{Closest SPIR-V Op:} \texttt{OpImageRead (paired with LEA\_TEX)}
\par\textbf{Pseudo-GLSL (Operational Semantics):}
\codebox{@sw = decodePixelFormat(mem\_read(src0 + \#offset), cvt\_desc = src2, register\_format, vecsize);}
\textbf{Execution Context:}
\begin{tightitemize}\footnotesize
  \item \textbf{Texture \& Sampler chain (\hyperref[sec:csf-srt]{\texttt{SRT\_2}} Frag / \texttt{SRT\_0} Comp/Vert):} 64B-aligned VA (\texttt{\& \textasciitilde{}0x3F}) $\to$ \texttt{Resource[table]} $\to$ \hyperref[sec:csf-tex-samp]{\texttt{Texture[tex\_idx]}} (32B, \texttt{type=2}) $\to$ \texttt{GenericPlane[mip]} (32B, \texttt{type=10}: \texttt{pointer}, \texttt{row\_stride}, Morton 16$\times$16 / AFBC) + \texttt{Sampler[samp\_idx]} (32B, \texttt{type=1}: filter, wrap, LOD clamp/bias, compare).
  \item \textbf{Inline FAU \& Special Operands (\hyperref[sec:csf-tsd-fau]{\texttt{FAU\_k}}, \texttt{r8..r15}):} Any 8-bit source with \texttt{mode=2} reads zero-latency uniform \texttt{u0..u31} from \texttt{FAU\_k} (\texttt{slots64 = bits[63:56]}); \texttt{mode=3} reads \texttt{SmallConstant} (Table~1.2) or \texttt{FauSpecial} (\texttt{pc}, \texttt{lane\_id}, \texttt{lane\_width}, \texttt{tls\_ptr}, \texttt{wls\_ptr}, \texttt{blend\_descriptor\_0..7}).
\end{tightitemize}
\par\textbf{Worked Examples:}
\par\textbf{Example 1 (Converted RGBA8 -\textgreater{} FP32 texel load):}
\codebox{\mnem{LD\_CVT}\imod{.f32.v4.slot0} \dst{@r0:r3}, \src{[r4:r5]}, \src{r6}, \imm{\#0}}
Loads texel at address \texttt{[r4:r5]} (from \texttt{LEA\_TEX}) and converts pixel format via descriptor \texttt{r6} into 4x FP32 \texttt{@r0:r3}.
\par\textbf{Example 2 (Relaxed-precision FP16 image load):}
\codebox{\mnem{LD\_CVT}\imod{.f16.v4.slot0} \dst{@r0:r1}, \src{[r4:r5]}, \src{r6}, \imm{\#0}}
Converts the same texel directly into 4x packed FP16 values in just 2 registers \texttt{@r0:r1}.

\begin{center}\footnotesize
\begin{tabularx}{\linewidth}{@{}lllX@{}}
\toprule
Variant & Arch & Class & Opcode Fields \\
\midrule
\texttt{(none)} & v9--v14 & \texttt{S3} & \texttt{opcode1=LD\_CVT} \\
\bottomrule
\end{tabularx}
\end{center}
\begin{center}\footnotesize
\begin{tabularx}{0.88\linewidth}{@{}l l X@{}}
\toprule
Field Name & Bits & Type / Fixed Value \\
\midrule
\texttt{src0} & \texttt{0:7} & \texttt{SourceEncoding64} \\
\texttt{offset} & \texttt{8:15} & \texttt{uint} \\
\texttt{src2} & \texttt{16:23} & \texttt{SourceEncoding} \\
\hyperref[enum:register_file_format_general_m]{\texttt{register\_format}} & \texttt{24:26} & \texttt{register\_file\_format\_general\_m} \\
\textit{Reserved} & \texttt{27} & \texttt{zero} \\
\hyperref[enum:vecsize_varying_m]{\texttt{vecsize}} & \texttt{28:29} & \texttt{vecsize\_varying\_m (def none)} \\
\hyperref[enum:message_slot_index_m]{\texttt{message\_slot\_index}} & \texttt{30:32} & \texttt{message\_slot\_index\_m} \\
\texttt{sr\_count} & \texttt{33:35} & \texttt{sr\_count\_t} \\
\textit{Reserved} & \texttt{36} & \texttt{zero} \\
\hyperref[enum:access_load_m]{\texttt{access}} & \texttt{37:39} & \texttt{access\_load\_m (def none)} \\
\texttt{staging\_register\_index} & \texttt{40:45} & \texttt{uint} \\
\texttt{staging\_register\_control\_32} & \texttt{46:47} & \texttt{exact=write} \\
\texttt{opcode1} & \texttt{48:51} & \texttt{exact=LD\_CVT} \\
\texttt{instruction\_class} & \texttt{52:56} & \texttt{exact=S3} \\
\texttt{fau\_page\_index} & \texttt{57:58} & \texttt{uint} \\
\hyperref[enum:flow_control_m]{\texttt{flow}} & \texttt{59:62} & \texttt{flow\_control\_m} \\
\textit{Reserved} & \texttt{63} & \texttt{zero} \\
\bottomrule
\end{tabularx}
\end{center}
\textbf{Fields Explained:}
\begin{tightitemize}\footnotesize
  \item \texttt{src0}: Primary source operand (\texttt{r0..r63}, \texttt{u0..u63} uniform, \texttt{k*} FAU constant, or \texttt{t0..t2} temporary; \texttt{\textasciicircum{}} marks discard/last-use). \emph{Example:} {\scriptsize\ttfamily \mnem{FADD.f32} \dst{r0}, \src{r1}, \src{r2}} (\textbf{SPIR-V:} \texttt{OpFAdd (Operand 0)}; \textbf{GLSL:} \texttt{float c = a + b;})
  \item \texttt{offset}: Signed byte offset added to base address register in memory load/store/atomic instructions. \emph{Example:} {\scriptsize\ttfamily \mnem{LOAD.i32} \dst{@r0}, \src{[r2:r3]}, \imm{\#16}} (\textbf{SPIR-V:} \texttt{OpAccessChain + OpLoad}; \textbf{GLSL:} \texttt{uint y = ssbo.arr[4];})
  \item \texttt{src2}: Third source operand (FMA addend, \texttt{MUX}/\texttt{CSEL} selector, \texttt{IADDC} carry-in, or bitwise mask in fused shift-logic). \emph{Example:} {\scriptsize\ttfamily \mnem{FMA.f32} \dst{r0}, \src{r1}, \src{r2}, \src{r3}} (\textbf{SPIR-V:} \texttt{GLSL.std.450 Fma (Addend)}; \textbf{GLSL:} \texttt{float d = fma(a, b, c);})
  \item \hyperref[enum:register_file_format_general_m]{\texttt{register\_format}}: Target GPR data conversion format (\texttt{auto32}, \texttt{f32}, \texttt{f16}, \texttt{s32}, \texttt{s16}, \texttt{u32}, \texttt{u16}). \emph{Example:} {\scriptsize\ttfamily \mnem{LD\_TILE}\imod{.f16.v4} \dst{@r0:r1}, \src{r2}} (\textbf{SPIR-V:} \texttt{OpImageRead (SubpassData)}; \textbf{GLSL:} \texttt{f16vec4 c = subpassLoad(inputAtt);})
  \item \hyperref[enum:vecsize_varying_m]{\texttt{vecsize}}: Component count (\texttt{none}=1, \texttt{v2}, \texttt{v3}, \texttt{v4}) for attribute, varying, or blend transfers. \emph{Example:} {\scriptsize\ttfamily \mnem{LD\_ATTR}\imod{.f32.v3} \dst{@r0:r2}, \src{r4}} (\textbf{SPIR-V:} \texttt{OpLoad (vec3 Input)}; \textbf{GLSL:} \texttt{in vec3 a\_normal;})
  \item \hyperref[enum:message_slot_index_m]{\texttt{message\_slot\_index}}: Asynchronous scoreboard slot (\texttt{slot0}--\texttt{slot2}, \texttt{slot7}) tracking completion of this message. \emph{Example:} {\scriptsize\ttfamily \mnem{LOAD.i32}\imod{.slot0} \dst{@r0}, \src{r2}, \imm{\#0}} (\textbf{SPIR-V:} \texttt{OpLoad}; \textbf{GLSL:} \texttt{uint v = ssbo.arr[i];})
  \item \texttt{sr\_count}: Number of consecutive 32-bit staging registers transferred by a load/store/texture message. \emph{Example:} {\scriptsize\ttfamily \mnem{LD\_MULTI.i128} \dst{@r0:r3}, \src{r4}, \imm{\#0}} (\textbf{SPIR-V:} \texttt{OpLoad (wide vector)}; \textbf{GLSL:} \texttt{uvec4 x = buf[i];})
  \item \hyperref[enum:access_load_m]{\texttt{access}}: Memory cache streaming/coherence policy (\texttt{none}, \texttt{istream}, \texttt{estream}, \texttt{force}). \emph{Example:} {\scriptsize\ttfamily \mnem{STORE.i32}\imod{.force} \src{@r0}, \src{r2}, \imm{\#0}} (\textbf{SPIR-V:} \texttt{OpStore (MakePointerVisible / Coherent)}; \textbf{GLSL:} \texttt{coherent buffer B \{ uint x; \}; x = 1u;})
  \item \texttt{staging\_register\_index}: Base GPR index of the contiguous staging register window (\texttt{@rN:rM}) for message payload/return. \emph{Example:} {\scriptsize\ttfamily \mnem{STORE.i64}\imod{.slot7} \src{@r0:r1}, \src{r2}, \imm{\#8}} (\textbf{SPIR-V:} \texttt{OpStore (uvec2)}; \textbf{GLSL:} \texttt{ssbo.pairs[i] = uvec2(a, b);})
  \item \texttt{staging\_register\_control\_32}: Controls the read/write staging register bank and count for async message instructions. \emph{Example:} {\scriptsize\ttfamily \mnem{LOAD.i128}\imod{.slot0} \dst{@r0:r3}, \src{r4}, \imm{\#0}} (\textbf{SPIR-V:} \texttt{OpLoad (vec4)}; \textbf{GLSL:} \texttt{vec4 v = ssbo.data[idx];})
  \item \texttt{opcode1}: Primary sub-opcode within the instruction class uniquely identifying the operation family. \emph{Example:} {\scriptsize\ttfamily \mnem{IADD.s32} \dst{r0}, \src{r1}, \src{r2}} (\textbf{SPIR-V:} \texttt{OpIAdd (operation selector)}; \textbf{GLSL:} \texttt{// Primary opcode bits})
  \item \texttt{instruction\_class}: 4-bit major instruction class selector (\texttt{bits 48:51}) partitioning the 64-bit Valhall opcode space (e.g., \texttt{FMA}, \texttt{CVT}, \texttt{SFU}, \texttt{LS}, \texttt{TEX}, \texttt{BRANCH}). \emph{Example:} {\scriptsize\ttfamily \mnem{FADD.f32} \dst{r0}, \src{r1}, \src{r2}} (\textbf{SPIR-V:} \texttt{N/A (encoding class)}; \textbf{GLSL:} \texttt{// Hardware execution pipe / decode class})
  \item \texttt{fau\_page\_index}: Selects the FAU (Fast Access Uniform) page (\texttt{0..3}) when a source operand references hardware small constants, thread IDs, or lane constants. \emph{Example:} {\scriptsize\ttfamily \mnem{FADD.f32} \dst{r0}, \src{r1}, \src{k1f}} (\textbf{SPIR-V:} \texttt{OpConstant / BuiltIn}; \textbf{GLSL:} \texttt{float y = x + 1.0;})
  \item \hyperref[enum:flow_control_m]{\texttt{flow}}: 4-bit embedded flow-control \& scoreboard wait modifier (\texttt{bits 59:62}: \texttt{none}, \texttt{wait0..wait0126}, \texttt{wait}, \texttt{reconverge}, \texttt{discard}, \texttt{end}). \emph{Example:} {\scriptsize\ttfamily \mnem{FADD.f32}\imod{.wait0.reconverge} \dst{r0}, \src{r1}, \src{r2}} (\textbf{SPIR-V:} \texttt{OpControlBarrier / OpReturn}; \textbf{GLSL:} \texttt{// Scoreboard wait + warp reconvergence})
\end{tightitemize}
\clearpage

\subsection*{LD\_GCLK}
\addcontentsline{toc}{subsection}{LD\_GCLK}
Reads the 64-bit hardware GPU clock (\texttt{cycle\_counter} for per-core shader cycles or \texttt{system\_timestamp} for the cross-core global timestamp) asynchronously into the 2-register staging pair \texttt{@sw[0:1]}.
\par\textbf{Notes:}
\begin{tightitemize}
  \item Non-deterministic architectural clock read.
  \item \textbf{Source references:} \texttt{kraid/ops.rs:2767-2797}, \texttt{bifrost/valhall/ISA.xml:72}.
\end{tightitemize}
\codebox{\mnem{LD\_GCLK}\imod{\{.message\_slot\_index\}}\imod{\{.flow\}}\imod{\{.source\}} \dst{@sw}}
\textbf{Operands:}
\begin{tightitemize}\footnotesize
  \item \texttt{\dst{@sw}} (\texttt{staging\_register\_index} + \texttt{sr\_count}): 2-register destination staging range (\texttt{@rN:rN+1}) receiving the 64-bit GPU clock value.
\end{tightitemize}
\par\textbf{Closest SPIR-V Op:} \texttt{OpReadClockKHR}
\par\textbf{Pseudo-GLSL (Operational Semantics):}
\codebox{if (clock\_source == cycle\_counter) @sw[0:1] = readPerCoreCycleCounter64();\\else                               @sw[0:1] = readSystemTimestamp64();}
\textbf{Execution Context:}
\begin{tightitemize}\footnotesize
  \item \textbf{Hardware clock source:} Reads per-core GPU cycle counter or cross-core system timestamp directly (no \texttt{SRT}/\texttt{FAU}/\texttt{TSD} descriptor dereference; tracked on \texttt{slot0..2}).
\end{tightitemize}
\par\textbf{Worked Examples:}
\par\textbf{Example 1 (GPU shader-core cycle counter):}
\codebox{\mnem{LD\_GCLK}\imod{.cycle\_counter.slot0} \dst{@r0:r1}}
Reads the 64-bit GPU cycle counter into \texttt{@r0:r1} (\texttt{r0} = low 32 bits, \texttt{r1} = high 32 bits) for fine-grained profiling.
\par\textbf{Example 2 (Global system timestamp):}
\codebox{\mnem{LD\_GCLK}\imod{.system\_timestamp.slot0} \dst{@r0:r1}}
Reads the monotonically increasing 64-bit system-wide timestamp synchronized across cores.

\begin{center}\footnotesize
\begin{tabularx}{\linewidth}{@{}lllX@{}}
\toprule
Variant & Arch & Class & Opcode Fields \\
\midrule
\texttt{(none)} & v9--v14 & \texttt{S0} & \texttt{opcode1=LD\_GCLK} \\
\bottomrule
\end{tabularx}
\end{center}
\begin{center}\footnotesize
\begin{tabularx}{0.88\linewidth}{@{}l l X@{}}
\toprule
Field Name & Bits & Type / Fixed Value \\
\midrule
\hyperref[enum:ld_gclk_source_m]{\texttt{source}} & \texttt{0:2} & \texttt{ld\_gclk\_source\_m} \\
\textit{Reserved} & \texttt{3:29} & \texttt{zero} \\
\hyperref[enum:message_slot_index_m]{\texttt{message\_slot\_index}} & \texttt{30:32} & \texttt{message\_slot\_index\_m} \\
\texttt{sr\_count} & \texttt{33:35} & \texttt{exact=count2} \\
\textit{Reserved} & \texttt{36:39} & \texttt{zero} \\
\texttt{staging\_register\_index} & \texttt{40:45} & \texttt{uint} \\
\texttt{staging\_register\_control\_32} & \texttt{46:47} & \texttt{exact=write} \\
\texttt{opcode1} & \texttt{48:51} & \texttt{exact=LD\_GCLK} \\
\texttt{instruction\_class} & \texttt{52:56} & \texttt{exact=S0} \\
\texttt{fau\_page\_index} & \texttt{57:58} & \texttt{exact=0} \\
\hyperref[enum:flow_control_m]{\texttt{flow}} & \texttt{59:62} & \texttt{flow\_control\_m} \\
\textit{Reserved} & \texttt{63} & \texttt{zero} \\
\bottomrule
\end{tabularx}
\end{center}
\textbf{Fields Explained:}
\begin{tightitemize}\footnotesize
  \item \hyperref[enum:ld_gclk_source_m]{\texttt{source}}: Selects global clock source in \texttt{LD\_GCLK} (\texttt{cycle\_counter} vs \texttt{system\_timestamp}). \emph{Example:} {\scriptsize\ttfamily \mnem{LD\_GCLK}\imod{.system\_timestamp} \dst{[r0:r1]}} (\textbf{SPIR-V:} \texttt{OpReadClockKHR}; \textbf{GLSL:} \texttt{uint64\_t t = clockARB();})
  \item \hyperref[enum:message_slot_index_m]{\texttt{message\_slot\_index}}: Asynchronous scoreboard slot (\texttt{slot0}--\texttt{slot2}, \texttt{slot7}) tracking completion of this message. \emph{Example:} {\scriptsize\ttfamily \mnem{LOAD.i32}\imod{.slot0} \dst{@r0}, \src{r2}, \imm{\#0}} (\textbf{SPIR-V:} \texttt{OpLoad}; \textbf{GLSL:} \texttt{uint v = ssbo.arr[i];})
  \item \texttt{sr\_count}: Number of consecutive 32-bit staging registers transferred by a load/store/texture message. \emph{Example:} {\scriptsize\ttfamily \mnem{LD\_MULTI.i128} \dst{@r0:r3}, \src{r4}, \imm{\#0}} (\textbf{SPIR-V:} \texttt{OpLoad (wide vector)}; \textbf{GLSL:} \texttt{uvec4 x = buf[i];})
  \item \texttt{staging\_register\_index}: Base GPR index of the contiguous staging register window (\texttt{@rN:rM}) for message payload/return. \emph{Example:} {\scriptsize\ttfamily \mnem{STORE.i64}\imod{.slot7} \src{@r0:r1}, \src{r2}, \imm{\#8}} (\textbf{SPIR-V:} \texttt{OpStore (uvec2)}; \textbf{GLSL:} \texttt{ssbo.pairs[i] = uvec2(a, b);})
  \item \texttt{staging\_register\_control\_32}: Controls the read/write staging register bank and count for async message instructions. \emph{Example:} {\scriptsize\ttfamily \mnem{LOAD.i128}\imod{.slot0} \dst{@r0:r3}, \src{r4}, \imm{\#0}} (\textbf{SPIR-V:} \texttt{OpLoad (vec4)}; \textbf{GLSL:} \texttt{vec4 v = ssbo.data[idx];})
  \item \texttt{opcode1}: Primary sub-opcode within the instruction class uniquely identifying the operation family. \emph{Example:} {\scriptsize\ttfamily \mnem{IADD.s32} \dst{r0}, \src{r1}, \src{r2}} (\textbf{SPIR-V:} \texttt{OpIAdd (operation selector)}; \textbf{GLSL:} \texttt{// Primary opcode bits})
  \item \texttt{instruction\_class}: 4-bit major instruction class selector (\texttt{bits 48:51}) partitioning the 64-bit Valhall opcode space (e.g., \texttt{FMA}, \texttt{CVT}, \texttt{SFU}, \texttt{LS}, \texttt{TEX}, \texttt{BRANCH}). \emph{Example:} {\scriptsize\ttfamily \mnem{FADD.f32} \dst{r0}, \src{r1}, \src{r2}} (\textbf{SPIR-V:} \texttt{N/A (encoding class)}; \textbf{GLSL:} \texttt{// Hardware execution pipe / decode class})
  \item \texttt{fau\_page\_index}: Selects the FAU (Fast Access Uniform) page (\texttt{0..3}) when a source operand references hardware small constants, thread IDs, or lane constants. \emph{Example:} {\scriptsize\ttfamily \mnem{FADD.f32} \dst{r0}, \src{r1}, \src{k1f}} (\textbf{SPIR-V:} \texttt{OpConstant / BuiltIn}; \textbf{GLSL:} \texttt{float y = x + 1.0;})
  \item \hyperref[enum:flow_control_m]{\texttt{flow}}: 4-bit embedded flow-control \& scoreboard wait modifier (\texttt{bits 59:62}: \texttt{none}, \texttt{wait0..wait0126}, \texttt{wait}, \texttt{reconverge}, \texttt{discard}, \texttt{end}). \emph{Example:} {\scriptsize\ttfamily \mnem{FADD.f32}\imod{.wait0.reconverge} \dst{r0}, \src{r1}, \src{r2}} (\textbf{SPIR-V:} \texttt{OpControlBarrier / OpReturn}; \textbf{GLSL:} \texttt{// Scoreboard wait + warp reconvergence})
\end{tightitemize}
\clearpage

\subsection*{LD\_MULTI}
\addcontentsline{toc}{subsection}{LD\_MULTI}
Multi-register load (v14+): loads \texttt{sr\_count} (\texttt{ls\_multi\_sr\_count\_m}) consecutive 32-bit staging registers from \texttt{src0 + sext64(offset)}.
\par\textbf{Notes:}
\begin{tightitemize}
  \item Wide bulk load beyond \texttt{LOAD.i128}.
  \item \textbf{Source references:} \texttt{isa-v9-v14.xml}.
\end{tightitemize}
\codebox{\mnem{LD\_MULTI}\imod{\{.sr\_count\}}\imod{\{.message\_slot\_index\}}\imod{\{.flow\}}\imod{\{.access\}} \dst{@sw}, \src{src0}, \imm{\#offset}}
\textbf{Operands:}
\begin{tightitemize}\footnotesize
  \item \texttt{\dst{@sw}} (\texttt{staging\_register\_index} + \texttt{sr\_count}): Destination staging register range (\texttt{sr\_count} 32-bit words) receiving bulk-loaded data.
  \item \texttt{\src{src0}}: 64-bit base GPU virtual address.
  \item \texttt{\imm{\#offset}} (\texttt{offset}): Signed byte offset added to \texttt{src0}.
\end{tightitemize}
\par\textbf{Closest SPIR-V Op:} \texttt{OpLoad (bulk \textgreater{}128-bit, v14+)}
\par\textbf{Pseudo-GLSL (Operational Semantics):}
\codebox{for (int i = 0; i \textless{} sr\_count; ++i) @sw[i] = mem\_u32[src0 + \#offset + i * 4];}
\textbf{Execution Context:}
\begin{tightitemize}\footnotesize
  \item \textbf{Direct 64-bit GPU VA memory/atomics:} Base addresses typically come from \hyperref[sec:csf-tsd-fau]{\texttt{FAU\_k}} (\texttt{u0..u31} populated by a Pilot shader chasing \texttt{SRT\_0} $\to$ \texttt{Resource[]} $\to$ \texttt{Buffer[]}), \texttt{FauSpecial::tls\_ptr}/\texttt{wls\_ptr} backed by \hyperref[sec:csf-tsd-fau]{\texttt{TSD\_k}} (\texttt{r24..r31}), or \texttt{FauSpecial::pc} for embedded constant LUTs.
  \item \textbf{Inline FAU \& Special Operands (\hyperref[sec:csf-tsd-fau]{\texttt{FAU\_k}}, \texttt{r8..r15}):} Any 8-bit source with \texttt{mode=2} reads zero-latency uniform \texttt{u0..u31} from \texttt{FAU\_k} (\texttt{slots64 = bits[63:56]}); \texttt{mode=3} reads \texttt{SmallConstant} (Table~1.2) or \texttt{FauSpecial} (\texttt{pc}, \texttt{lane\_id}, \texttt{lane\_width}, \texttt{tls\_ptr}, \texttt{wls\_ptr}, \texttt{blend\_descriptor\_0..7}).
\end{tightitemize}
\par\textbf{Worked Examples:}
\par\textbf{Example 1 (256-bit (8-GPR) bulk load (v14+)):}
\codebox{\mnem{LD\_MULTI}\imod{.i256.slot0} \dst{@r0:r7}, \src{[r8:r9]}, \imm{\#0}}
Loads 8 consecutive 32-bit registers (32 bytes) from \texttt{[r8:r9]} in a single v14 memory message.
\par\textbf{Example 2 (Offset 160-bit (5-GPR) load):}
\codebox{\mnem{LD\_MULTI}\imod{.i160.slot1} \dst{@r0:r4}, \src{[r8:r9]}, \imm{\#0x20}}
Loads 20 bytes from \texttt{[r8:r9] + 0x20} into \texttt{@r0:r4} on scoreboard \texttt{slot1}.

\begin{center}\footnotesize
\begin{tabularx}{\linewidth}{@{}lllX@{}}
\toprule
Variant & Arch & Class & Opcode Fields \\
\midrule
\texttt{(none)} & v9--v14 & \texttt{S2} & \texttt{opcode1=LD\_MULTI} \\
\bottomrule
\end{tabularx}
\end{center}
\begin{center}\footnotesize
\begin{tabularx}{0.88\linewidth}{@{}l l X@{}}
\toprule
Field Name & Bits & Type / Fixed Value \\
\midrule
\texttt{src0} & \texttt{0:7} & \texttt{SourceEncoding64} \\
\texttt{offset} & \texttt{8:23} & \texttt{int} \\
\hyperref[enum:access_load_m]{\texttt{access}} & \texttt{24:26} & \texttt{access\_load\_m (def none)} \\
\textit{Reserved} & \texttt{27:29} & \texttt{zero} \\
\hyperref[enum:message_slot_index_m]{\texttt{message\_slot\_index}} & \texttt{30:32} & \texttt{message\_slot\_index\_m} \\
\hyperref[enum:ls_multi_sr_count_m]{\texttt{sr\_count}} & \texttt{33:38} & \texttt{ls\_multi\_sr\_count\_m} \\
\textit{Reserved} & \texttt{39} & \texttt{zero} \\
\texttt{staging\_register\_index} & \texttt{40:45} & \texttt{uint} \\
\texttt{staging\_register\_control\_32} & \texttt{46:47} & \texttt{exact=write} \\
\texttt{opcode1} & \texttt{48:51} & \texttt{exact=LD\_MULTI} \\
\texttt{instruction\_class} & \texttt{52:56} & \texttt{exact=S2} \\
\texttt{fau\_page\_index} & \texttt{57:58} & \texttt{uint} \\
\hyperref[enum:flow_control_m]{\texttt{flow}} & \texttt{59:62} & \texttt{flow\_control\_m} \\
\textit{Reserved} & \texttt{63} & \texttt{zero} \\
\bottomrule
\end{tabularx}
\end{center}
\textbf{Fields Explained:}
\begin{tightitemize}\footnotesize
  \item \texttt{src0}: Primary source operand (\texttt{r0..r63}, \texttt{u0..u63} uniform, \texttt{k*} FAU constant, or \texttt{t0..t2} temporary; \texttt{\textasciicircum{}} marks discard/last-use). \emph{Example:} {\scriptsize\ttfamily \mnem{FADD.f32} \dst{r0}, \src{r1}, \src{r2}} (\textbf{SPIR-V:} \texttt{OpFAdd (Operand 0)}; \textbf{GLSL:} \texttt{float c = a + b;})
  \item \texttt{offset}: Signed byte offset added to base address register in memory load/store/atomic instructions. \emph{Example:} {\scriptsize\ttfamily \mnem{LOAD.i32} \dst{@r0}, \src{[r2:r3]}, \imm{\#16}} (\textbf{SPIR-V:} \texttt{OpAccessChain + OpLoad}; \textbf{GLSL:} \texttt{uint y = ssbo.arr[4];})
  \item \hyperref[enum:access_load_m]{\texttt{access}}: Memory cache streaming/coherence policy (\texttt{none}, \texttt{istream}, \texttt{estream}, \texttt{force}). \emph{Example:} {\scriptsize\ttfamily \mnem{STORE.i32}\imod{.force} \src{@r0}, \src{r2}, \imm{\#0}} (\textbf{SPIR-V:} \texttt{OpStore (MakePointerVisible / Coherent)}; \textbf{GLSL:} \texttt{coherent buffer B \{ uint x; \}; x = 1u;})
  \item \hyperref[enum:message_slot_index_m]{\texttt{message\_slot\_index}}: Asynchronous scoreboard slot (\texttt{slot0}--\texttt{slot2}, \texttt{slot7}) tracking completion of this message. \emph{Example:} {\scriptsize\ttfamily \mnem{LOAD.i32}\imod{.slot0} \dst{@r0}, \src{r2}, \imm{\#0}} (\textbf{SPIR-V:} \texttt{OpLoad}; \textbf{GLSL:} \texttt{uint v = ssbo.arr[i];})
  \item \hyperref[enum:ls_multi_sr_count_m]{\texttt{sr\_count}}: Number of consecutive 32-bit staging registers transferred by a load/store/texture message. \emph{Example:} {\scriptsize\ttfamily \mnem{LD\_MULTI.i128} \dst{@r0:r3}, \src{r4}, \imm{\#0}} (\textbf{SPIR-V:} \texttt{OpLoad (wide vector)}; \textbf{GLSL:} \texttt{uvec4 x = buf[i];})
  \item \texttt{staging\_register\_index}: Base GPR index of the contiguous staging register window (\texttt{@rN:rM}) for message payload/return. \emph{Example:} {\scriptsize\ttfamily \mnem{STORE.i64}\imod{.slot7} \src{@r0:r1}, \src{r2}, \imm{\#8}} (\textbf{SPIR-V:} \texttt{OpStore (uvec2)}; \textbf{GLSL:} \texttt{ssbo.pairs[i] = uvec2(a, b);})
  \item \texttt{staging\_register\_control\_32}: Controls the read/write staging register bank and count for async message instructions. \emph{Example:} {\scriptsize\ttfamily \mnem{LOAD.i128}\imod{.slot0} \dst{@r0:r3}, \src{r4}, \imm{\#0}} (\textbf{SPIR-V:} \texttt{OpLoad (vec4)}; \textbf{GLSL:} \texttt{vec4 v = ssbo.data[idx];})
  \item \texttt{opcode1}: Primary sub-opcode within the instruction class uniquely identifying the operation family. \emph{Example:} {\scriptsize\ttfamily \mnem{IADD.s32} \dst{r0}, \src{r1}, \src{r2}} (\textbf{SPIR-V:} \texttt{OpIAdd (operation selector)}; \textbf{GLSL:} \texttt{// Primary opcode bits})
  \item \texttt{instruction\_class}: 4-bit major instruction class selector (\texttt{bits 48:51}) partitioning the 64-bit Valhall opcode space (e.g., \texttt{FMA}, \texttt{CVT}, \texttt{SFU}, \texttt{LS}, \texttt{TEX}, \texttt{BRANCH}). \emph{Example:} {\scriptsize\ttfamily \mnem{FADD.f32} \dst{r0}, \src{r1}, \src{r2}} (\textbf{SPIR-V:} \texttt{N/A (encoding class)}; \textbf{GLSL:} \texttt{// Hardware execution pipe / decode class})
  \item \texttt{fau\_page\_index}: Selects the FAU (Fast Access Uniform) page (\texttt{0..3}) when a source operand references hardware small constants, thread IDs, or lane constants. \emph{Example:} {\scriptsize\ttfamily \mnem{FADD.f32} \dst{r0}, \src{r1}, \src{k1f}} (\textbf{SPIR-V:} \texttt{OpConstant / BuiltIn}; \textbf{GLSL:} \texttt{float y = x + 1.0;})
  \item \hyperref[enum:flow_control_m]{\texttt{flow}}: 4-bit embedded flow-control \& scoreboard wait modifier (\texttt{bits 59:62}: \texttt{none}, \texttt{wait0..wait0126}, \texttt{wait}, \texttt{reconverge}, \texttt{discard}, \texttt{end}). \emph{Example:} {\scriptsize\ttfamily \mnem{FADD.f32}\imod{.wait0.reconverge} \dst{r0}, \src{r1}, \src{r2}} (\textbf{SPIR-V:} \texttt{OpControlBarrier / OpReturn}; \textbf{GLSL:} \texttt{// Scoreboard wait + warp reconvergence})
\end{tightitemize}
\clearpage

\subsection*{LD\_PKA}
\addcontentsline{toc}{subsection}{LD\_PKA}
Loads \texttt{i8} through \texttt{i128} bytes from the buffer descriptor identified by handle \texttt{src1} at byte offset \texttt{src0} into staging register range \texttt{@sw}, with hardware descriptor resolution and bounds checking.
\par\textbf{Notes:}
\begin{tightitemize}
  \item If \texttt{handle[31:25] == 0x7F}, loads descriptor.
  \item If \texttt{0}, loads buffer data.
  \item \textbf{Source references:} \texttt{kraid/ops.rs:2802-2828}, \texttt{bifrost/valhall/ISA.xml:74-81}.
\end{tightitemize}
\codebox{\mnem{LD\_PKA}\imod{\{.message\_slot\_index\}}\imod{\{.flow\}}\imod{\{.access\}} \dst{@sw}, \src{src0}, \src{src1}\\\mnem{LD\_PKA}\imod{\{.message\_slot\_index\}}\imod{\{.flow\}}\imod{\{.access\}}\imod{\{.lane\}}\imod{\{.extend\}} \dst{@sw}, \src{src0}, \src{src1}}
\textbf{Operands:}
\begin{tightitemize}\footnotesize
  \item \texttt{\dst{@sw}} (\texttt{staging\_register\_index} + \texttt{sr\_count}): Destination staging register range receiving the loaded buffer data.
  \item \texttt{\src{src0}}: 32-bit byte offset within the buffer.
  \item \texttt{\src{src1}}: 32-bit buffer descriptor handle (table + index).
\end{tightitemize}
\par\textbf{Closest SPIR-V Op:} \texttt{OpLoad (UBO/SSBO descriptor table)}
\par\textbf{Pseudo-GLSL (Operational Semantics):}
\codebox{BufferDesc b = resolveBufferDesc(src1);\\@sw = (src0 + access\_bytes \textless{}= b.size) ? mem\_read(b.address + src0, width) : 0;}
\textbf{Execution Context:}
\begin{tightitemize}\footnotesize
  \item \textbf{Buffer descriptor chain (\hyperref[sec:csf-srt]{\texttt{SRT\_k}}, \texttt{r0:r1..r6:r7}):} Masks 64B-aligned VA (\texttt{\& \textasciitilde{}0x3F}, \texttt{count=bits[5:0]}) $\to$ \texttt{Resource[table]} (16B, \texttt{contains\_descriptors=1}) $\to$ \hyperref[sec:csf-buf-attr]{\texttt{Buffer[index]}} (32B, \texttt{type=9}: 64-bit \texttt{address}, \texttt{size}). Precondition: \texttt{table < count} and \texttt{index*32 < Resource[table].size}.
  \item \textbf{Inline FAU \& Special Operands (\hyperref[sec:csf-tsd-fau]{\texttt{FAU\_k}}, \texttt{r8..r15}):} Any 8-bit source with \texttt{mode=2} reads zero-latency uniform \texttt{u0..u31} from \texttt{FAU\_k} (\texttt{slots64 = bits[63:56]}); \texttt{mode=3} reads \texttt{SmallConstant} (Table~1.2) or \texttt{FauSpecial} (\texttt{pc}, \texttt{lane\_id}, \texttt{lane\_width}, \texttt{tls\_ptr}, \texttt{wls\_ptr}, \texttt{blend\_descriptor\_0..7}).
\end{tightitemize}
\par\textbf{Worked Examples:}
\par\textbf{Example 1 (128-bit descriptor-indexed buffer load):}
\codebox{\mnem{LD\_PKA.i128}\imod{.slot0} \dst{@r0:r3}, \src{r4}, \src{u0}}
Loads 16 bytes from PKA buffer descriptor \texttt{u0} at byte offset \texttt{r4} into \texttt{@r0:r3}.
\par\textbf{Example 2 (32-bit scalar descriptor load):}
\codebox{\mnem{LD\_PKA.i32}\imod{.slot0} \dst{@r0}, \src{r4}, \src{u0}}
Loads 4 bytes from descriptor \texttt{u0} at dynamic offset \texttt{r4} into \texttt{@r0}.
\par\textbf{Mined SPIR-V $\to$ Valhall ISA Example (Loading from a uniform/storage buffer via descriptor handle --- Vertex shader, Mali-G710 (v10)):}
\codebox{\textbf{; SPIR-V (input)}\\\%ptr = OpAccessChain \%\_ptr\_Uniform\_uint \%ubo \%int\_0 \%idx\\\%val = OpLoad \%uint \%ptr\\[3pt]\textbf{; Valhall ISA (malioc output)}\\IMUL.u32 r0, r0\textasciicircum{}, k13121110.b0\\LD\_PKA.i32.slot0.wait0 @r0, r0\textasciicircum{}, u12.w1}
\texttt{LD\_PKA.i32} loads 32 bits at byte offset \texttt{r0\textasciicircum{}} from the buffer descriptor identified by handle \texttt{u12.w1}, performing hardware descriptor resolution and bounds checking in one message.
\begin{center}\footnotesize
\begin{tabularx}{\linewidth}{@{}lllX@{}}
\toprule
Variant & Arch & Class & Opcode Fields \\
\midrule
\texttt{(none)} & v9--v14 & \texttt{S2} & \texttt{opcode1=LD\_PKA} \\
\texttt{(none)} & v9--v14 & \texttt{S2} & \texttt{opcode1=LD\_PKA} \\
\texttt{(none)} & v9--v14 & \texttt{S2} & \texttt{opcode1=LD\_PKA} \\
\texttt{(none)} & v9--v14 & \texttt{S2} & \texttt{opcode1=LD\_PKA} \\
\texttt{(none)} & v9--v14 & \texttt{S2} & \texttt{opcode1=LD\_PKA} \\
\texttt{(none)} & v9--v14 & \texttt{S2} & \texttt{opcode1=LD\_PKA} \\
\texttt{(none)} & v9--v14 & \texttt{S2} & \texttt{opcode1=LD\_PKA} \\
\texttt{(none)} & v9--v14 & \texttt{S2} & \texttt{opcode1=LD\_PKA} \\
\bottomrule
\end{tabularx}
\end{center}
\begin{center}\footnotesize
\begin{tabularx}{0.88\linewidth}{@{}l l X@{}}
\toprule
Field Name & Bits & Type / Fixed Value \\
\midrule
\texttt{src0} & \texttt{0:7} & \texttt{SourceEncoding} \\
\texttt{src1} & \texttt{8:15} & \texttt{SourceEncoding} \\
\textit{Reserved} & \texttt{16:23} & \texttt{zero} \\
\hyperref[enum:access_load_m]{\texttt{access}} & \texttt{24:26} & \texttt{access\_load\_m (def none)} \\
\texttt{memory\_width} & \texttt{27:29} & \texttt{exact=i128} \\
\hyperref[enum:message_slot_index_m]{\texttt{message\_slot\_index}} & \texttt{30:32} & \texttt{message\_slot\_index\_m} \\
\texttt{sr\_count} & \texttt{33:35} & \texttt{exact=count4} \\
\hyperref[enum:lane_all_m]{\texttt{lane}} & \texttt{36:38} & \texttt{exact=d0} \\
\hyperref[enum:sign_extend_or_none_compressed_encoding_m]{\texttt{extend}} & \texttt{39} & \texttt{exact=zext\_or\_none} \\
\texttt{staging\_register\_index} & \texttt{40:45} & \texttt{uint} \\
\texttt{staging\_register\_control\_32} & \texttt{46:47} & \texttt{exact=write} \\
\texttt{opcode1} & \texttt{48:51} & \texttt{exact=LD\_PKA} \\
\texttt{instruction\_class} & \texttt{52:56} & \texttt{exact=S2} \\
\texttt{fau\_page\_index} & \texttt{57:58} & \texttt{uint} \\
\hyperref[enum:flow_control_m]{\texttt{flow}} & \texttt{59:62} & \texttt{flow\_control\_m} \\
\textit{Reserved} & \texttt{63} & \texttt{zero} \\
\bottomrule
\end{tabularx}
\end{center}
\textbf{Fields Explained:}
\begin{tightitemize}\footnotesize
  \item \texttt{src0}: Primary source operand (\texttt{r0..r63}, \texttt{u0..u63} uniform, \texttt{k*} FAU constant, or \texttt{t0..t2} temporary; \texttt{\textasciicircum{}} marks discard/last-use). \emph{Example:} {\scriptsize\ttfamily \mnem{FADD.f32} \dst{r0}, \src{r1}, \src{r2}} (\textbf{SPIR-V:} \texttt{OpFAdd (Operand 0)}; \textbf{GLSL:} \texttt{float c = a + b;})
  \item \texttt{src1}: Secondary source operand (register, uniform, or FAU constant; in shifts, supplies the shift count byte). \emph{Example:} {\scriptsize\ttfamily \mnem{LSHIFT.i32} \dst{r0}, \src{r1}, \src{r2}\imod{.b0}} (\textbf{SPIR-V:} \texttt{OpShiftLeftLogical (Shift)}; \textbf{GLSL:} \texttt{uint y = a \textless{}\textless{} s;})
  \item \hyperref[enum:access_load_m]{\texttt{access}}: Memory cache streaming/coherence policy (\texttt{none}, \texttt{istream}, \texttt{estream}, \texttt{force}). \emph{Example:} {\scriptsize\ttfamily \mnem{STORE.i32}\imod{.force} \src{@r0}, \src{r2}, \imm{\#0}} (\textbf{SPIR-V:} \texttt{OpStore (MakePointerVisible / Coherent)}; \textbf{GLSL:} \texttt{coherent buffer B \{ uint x; \}; x = 1u;})
  \item \texttt{memory\_width}: Memory transfer width (\texttt{i8}, \texttt{i16}, \texttt{i24}, \texttt{i32}, \texttt{i48}, \texttt{i64}, \texttt{i96}, \texttt{i128}). \emph{Example:} {\scriptsize\ttfamily \mnem{LOAD.i64} \dst{@r0:r1}, \src{[r2:r3]}, \imm{\#0}} (\textbf{SPIR-V:} \texttt{OpLoad (uvec2)}; \textbf{GLSL:} \texttt{uvec2 v = ssbo.pairs[i];})
  \item \hyperref[enum:message_slot_index_m]{\texttt{message\_slot\_index}}: Asynchronous scoreboard slot (\texttt{slot0}--\texttt{slot2}, \texttt{slot7}) tracking completion of this message. \emph{Example:} {\scriptsize\ttfamily \mnem{LOAD.i32}\imod{.slot0} \dst{@r0}, \src{r2}, \imm{\#0}} (\textbf{SPIR-V:} \texttt{OpLoad}; \textbf{GLSL:} \texttt{uint v = ssbo.arr[i];})
  \item \texttt{sr\_count}: Number of consecutive 32-bit staging registers transferred by a load/store/texture message. \emph{Example:} {\scriptsize\ttfamily \mnem{LD\_MULTI.i128} \dst{@r0:r3}, \src{r4}, \imm{\#0}} (\textbf{SPIR-V:} \texttt{OpLoad (wide vector)}; \textbf{GLSL:} \texttt{uvec4 x = buf[i];})
  \item \hyperref[enum:lane_all_m]{\texttt{lane}}: Sub-word lane selector (\texttt{b0}..\texttt{b3}, \texttt{h0}, \texttt{h1}, \texttt{w0}, \texttt{d0}) for narrow checked/unpacked memory transfers. \emph{Example:} {\scriptsize\ttfamily \mnem{LOAD.i16}\imod{.h1} \dst{@r0}, \src{r2}, \imm{\#0}} (\textbf{SPIR-V:} \texttt{OpLoad (uint16\_t)}; \textbf{GLSL:} \texttt{uint16\_t x = ssbo.u16[i];})
  \item \hyperref[enum:sign_extend_or_none_compressed_encoding_m]{\texttt{extend}}: Virtual extension mode (\texttt{sext}, \texttt{zext}, \texttt{none}) decoded from \texttt{encoded\_extend} and \texttt{memory\_width}. \emph{Example:} {\scriptsize\ttfamily \mnem{LOAD.i8}\imod{.sext} \dst{@r0}, \src{r2}, \imm{\#0}} (\textbf{SPIR-V:} \texttt{OpLoad + OpSConvert}; \textbf{GLSL:} \texttt{int x = int(ssbo.i8\_arr[i]);})
  \item \texttt{staging\_register\_index}: Base GPR index of the contiguous staging register window (\texttt{@rN:rM}) for message payload/return. \emph{Example:} {\scriptsize\ttfamily \mnem{STORE.i64}\imod{.slot7} \src{@r0:r1}, \src{r2}, \imm{\#8}} (\textbf{SPIR-V:} \texttt{OpStore (uvec2)}; \textbf{GLSL:} \texttt{ssbo.pairs[i] = uvec2(a, b);})
  \item \texttt{staging\_register\_control\_32}: Controls the read/write staging register bank and count for async message instructions. \emph{Example:} {\scriptsize\ttfamily \mnem{LOAD.i128}\imod{.slot0} \dst{@r0:r3}, \src{r4}, \imm{\#0}} (\textbf{SPIR-V:} \texttt{OpLoad (vec4)}; \textbf{GLSL:} \texttt{vec4 v = ssbo.data[idx];})
  \item \texttt{opcode1}: Primary sub-opcode within the instruction class uniquely identifying the operation family. \emph{Example:} {\scriptsize\ttfamily \mnem{IADD.s32} \dst{r0}, \src{r1}, \src{r2}} (\textbf{SPIR-V:} \texttt{OpIAdd (operation selector)}; \textbf{GLSL:} \texttt{// Primary opcode bits})
  \item \texttt{instruction\_class}: 4-bit major instruction class selector (\texttt{bits 48:51}) partitioning the 64-bit Valhall opcode space (e.g., \texttt{FMA}, \texttt{CVT}, \texttt{SFU}, \texttt{LS}, \texttt{TEX}, \texttt{BRANCH}). \emph{Example:} {\scriptsize\ttfamily \mnem{FADD.f32} \dst{r0}, \src{r1}, \src{r2}} (\textbf{SPIR-V:} \texttt{N/A (encoding class)}; \textbf{GLSL:} \texttt{// Hardware execution pipe / decode class})
  \item \texttt{fau\_page\_index}: Selects the FAU (Fast Access Uniform) page (\texttt{0..3}) when a source operand references hardware small constants, thread IDs, or lane constants. \emph{Example:} {\scriptsize\ttfamily \mnem{FADD.f32} \dst{r0}, \src{r1}, \src{k1f}} (\textbf{SPIR-V:} \texttt{OpConstant / BuiltIn}; \textbf{GLSL:} \texttt{float y = x + 1.0;})
  \item \hyperref[enum:flow_control_m]{\texttt{flow}}: 4-bit embedded flow-control \& scoreboard wait modifier (\texttt{bits 59:62}: \texttt{none}, \texttt{wait0..wait0126}, \texttt{wait}, \texttt{reconverge}, \texttt{discard}, \texttt{end}). \emph{Example:} {\scriptsize\ttfamily \mnem{FADD.f32}\imod{.wait0.reconverge} \dst{r0}, \src{r1}, \src{r2}} (\textbf{SPIR-V:} \texttt{OpControlBarrier / OpReturn}; \textbf{GLSL:} \texttt{// Scoreboard wait + warp reconvergence})
\end{tightitemize}
\clearpage

\subsection*{LEA\_ATTR}
\addcontentsline{toc}{subsection}{LEA\_ATTR}
Computes effective address + conversion descriptor (3 staging regs: \texttt{addr\_lo}, \texttt{addr\_hi}, \texttt{cvt\_desc}) for vertex attribute at dynamic handle.
\par\textbf{Notes:}
\begin{tightitemize}
  \item Returns 3 staging registers (\texttt{ISA.xml:103}).
  \item \textbf{Source references:} \texttt{bifrost/valhall/ISA.xml:103}, \texttt{isa-v9-v14.xml}.
\end{tightitemize}
\codebox{\mnem{LEA\_ATTR}\imod{\{.message\_slot\_index\}}\imod{\{.flow\}} \dst{@sw}, \src{src0}, \src{src1}, \src{src2}}
\textbf{Operands:}
\begin{tightitemize}\footnotesize
  \item \texttt{\dst{@sw}} (\texttt{staging\_register\_index} + \texttt{sr\_count}): 3-register destination range: \texttt{@sw[0:1]} = 64-bit attribute VA, \texttt{@sw[2]} = 32-bit format conversion descriptor.
  \item \texttt{\src{src0}}: Vertex index (\texttt{VertexIndex}).
  \item \texttt{\src{src1}}: Instance index (\texttt{InstanceIndex}).
  \item \texttt{\src{src2}}: Dynamic attribute descriptor handle.
\end{tightitemize}
\par\textbf{Closest SPIR-V Op:} \texttt{None (attribute address + cvt desc)}
\par\textbf{Pseudo-GLSL (Operational Semantics):}
\codebox{@sw[0:1] = attrAddress(handle = src2, vert = src0, inst = src1);\\@sw[2]   = attrFormatDesc(handle = src2);}
\textbf{Execution Context:}
\begin{tightitemize}\footnotesize
  \item \textbf{Attribute descriptor chain (\hyperref[sec:csf-srt]{\texttt{IDVSSR::SRT\_0}}, \texttt{r0:r1}):} Masks 64B-aligned VA (\texttt{\& \textasciitilde{}0x3F}, \texttt{count=bits[5:0]}) $\to$ \texttt{Resource[table]} (16B) $\to$ \hyperref[sec:csf-buf-attr]{\texttt{Attribute[index]}} (32B: \texttt{table}, \texttt{buffer\_index}, \texttt{offset}, \texttt{stride}, \texttt{format}) $\to$ \texttt{Resource[Attribute::table]} $\to$ \texttt{Buffer[Attribute::buffer\_index]} (32B: \texttt{address}, \texttt{size}).
  \item \textbf{Vertex stage preloads (\hyperref[sec:csf-spd]{\texttt{SPD\_0}/\texttt{SPD\_1}}, \texttt{r16..r19}):} The hardware Index Fetcher reads \hyperref[sec:csf-regs]{\texttt{IDVSSR::INDEX\_BUFFER}} (\texttt{r54:r55}) and pre-fills \texttt{r59, r60} with \texttt{VertexIndex} (\texttt{+ VERTEX\_OFFSET + GLOBAL\_ATTR\_OFF}) and \texttt{r61} with \texttt{InstanceIndex} (\texttt{+ INSTANCE\_OFFSET}) so Position and Varying shaders can pull vertex attributes directly.
\end{tightitemize}
\par\textbf{Worked Examples:}
\par\textbf{Example 1 (Dynamic vertex attribute address lookup):}
\codebox{\mnem{LEA\_ATTR}\imod{.slot0} \dst{@r0:r2}, \src{r60}, \src{r61}, \src{r4}}
Computes \texttt{(addr\_lo, addr\_hi, cvt\_desc)} in \texttt{@r0:r2} for \texttt{(vertex\_id=r60, instance\_id=r61)} and dynamic attribute handle \texttt{r4}.
\par\textbf{Example 2 (Paired with `LD\_CVT` / `ST\_CVT`):}
\codebox{\mnem{LEA\_ATTR}\imod{.slot1} \dst{@r4:r6}, \src{r60}, \src{r61}, \src{r8}}
Produces the 64-bit address \texttt{[r4:r5]} and format descriptor \texttt{r6} consumed by a subsequent \texttt{ST\_CVT} in vertex-output shaders.

\begin{center}\footnotesize
\begin{tabularx}{\linewidth}{@{}lllX@{}}
\toprule
Variant & Arch & Class & Opcode Fields \\
\midrule
\texttt{(none)} & v9--v14 & \texttt{S3} & \texttt{opcode1=LEA\_ATTR} \\
\bottomrule
\end{tabularx}
\end{center}
\begin{center}\footnotesize
\begin{tabularx}{0.88\linewidth}{@{}l l X@{}}
\toprule
Field Name & Bits & Type / Fixed Value \\
\midrule
\texttt{src0} & \texttt{0:7} & \texttt{SourceEncoding} \\
\texttt{src1} & \texttt{8:15} & \texttt{SourceEncoding} \\
\texttt{src2} & \texttt{16:23} & \texttt{SourceEncoding} \\
\textit{Reserved} & \texttt{24:29} & \texttt{zero} \\
\hyperref[enum:message_slot_index_m]{\texttt{message\_slot\_index}} & \texttt{30:32} & \texttt{message\_slot\_index\_m} \\
\texttt{sr\_count} & \texttt{33:35} & \texttt{exact=count3} \\
\textit{Reserved} & \texttt{36} & \texttt{zero} \\
\texttt{descriptor\_type} & \texttt{37:38} & \texttt{exact=attribute} \\
\textit{Reserved} & \texttt{39} & \texttt{zero} \\
\texttt{staging\_register\_index} & \texttt{40:45} & \texttt{uint} \\
\texttt{staging\_register\_control\_32} & \texttt{46:47} & \texttt{exact=write} \\
\texttt{opcode1} & \texttt{48:51} & \texttt{exact=LEA\_ATTR} \\
\texttt{instruction\_class} & \texttt{52:56} & \texttt{exact=S3} \\
\texttt{fau\_page\_index} & \texttt{57:58} & \texttt{uint} \\
\hyperref[enum:flow_control_m]{\texttt{flow}} & \texttt{59:62} & \texttt{flow\_control\_m} \\
\textit{Reserved} & \texttt{63} & \texttt{zero} \\
\bottomrule
\end{tabularx}
\end{center}
\textbf{Fields Explained:}
\begin{tightitemize}\footnotesize
  \item \texttt{src0}: Primary source operand (\texttt{r0..r63}, \texttt{u0..u63} uniform, \texttt{k*} FAU constant, or \texttt{t0..t2} temporary; \texttt{\textasciicircum{}} marks discard/last-use). \emph{Example:} {\scriptsize\ttfamily \mnem{FADD.f32} \dst{r0}, \src{r1}, \src{r2}} (\textbf{SPIR-V:} \texttt{OpFAdd (Operand 0)}; \textbf{GLSL:} \texttt{float c = a + b;})
  \item \texttt{src1}: Secondary source operand (register, uniform, or FAU constant; in shifts, supplies the shift count byte). \emph{Example:} {\scriptsize\ttfamily \mnem{LSHIFT.i32} \dst{r0}, \src{r1}, \src{r2}\imod{.b0}} (\textbf{SPIR-V:} \texttt{OpShiftLeftLogical (Shift)}; \textbf{GLSL:} \texttt{uint y = a \textless{}\textless{} s;})
  \item \texttt{src2}: Third source operand (FMA addend, \texttt{MUX}/\texttt{CSEL} selector, \texttt{IADDC} carry-in, or bitwise mask in fused shift-logic). \emph{Example:} {\scriptsize\ttfamily \mnem{FMA.f32} \dst{r0}, \src{r1}, \src{r2}, \src{r3}} (\textbf{SPIR-V:} \texttt{GLSL.std.450 Fma (Addend)}; \textbf{GLSL:} \texttt{float d = fma(a, b, c);})
  \item \hyperref[enum:message_slot_index_m]{\texttt{message\_slot\_index}}: Asynchronous scoreboard slot (\texttt{slot0}--\texttt{slot2}, \texttt{slot7}) tracking completion of this message. \emph{Example:} {\scriptsize\ttfamily \mnem{LOAD.i32}\imod{.slot0} \dst{@r0}, \src{r2}, \imm{\#0}} (\textbf{SPIR-V:} \texttt{OpLoad}; \textbf{GLSL:} \texttt{uint v = ssbo.arr[i];})
  \item \texttt{sr\_count}: Number of consecutive 32-bit staging registers transferred by a load/store/texture message. \emph{Example:} {\scriptsize\ttfamily \mnem{LD\_MULTI.i128} \dst{@r0:r3}, \src{r4}, \imm{\#0}} (\textbf{SPIR-V:} \texttt{OpLoad (wide vector)}; \textbf{GLSL:} \texttt{uvec4 x = buf[i];})
  \item \texttt{descriptor\_type}: Selects descriptor table binding class (attribute, buffer, texture, varying). \emph{Example:} {\scriptsize\ttfamily \mnem{LD\_BUFFER.i32} \dst{@r0}, \src{r2}, \imm{\#0}} (\textbf{SPIR-V:} \texttt{OpLoad (StorageBuffer)}; \textbf{GLSL:} \texttt{uint v = ssbo.arr[i];})
  \item \texttt{staging\_register\_index}: Base GPR index of the contiguous staging register window (\texttt{@rN:rM}) for message payload/return. \emph{Example:} {\scriptsize\ttfamily \mnem{STORE.i64}\imod{.slot7} \src{@r0:r1}, \src{r2}, \imm{\#8}} (\textbf{SPIR-V:} \texttt{OpStore (uvec2)}; \textbf{GLSL:} \texttt{ssbo.pairs[i] = uvec2(a, b);})
  \item \texttt{staging\_register\_control\_32}: Controls the read/write staging register bank and count for async message instructions. \emph{Example:} {\scriptsize\ttfamily \mnem{LOAD.i128}\imod{.slot0} \dst{@r0:r3}, \src{r4}, \imm{\#0}} (\textbf{SPIR-V:} \texttt{OpLoad (vec4)}; \textbf{GLSL:} \texttt{vec4 v = ssbo.data[idx];})
  \item \texttt{opcode1}: Primary sub-opcode within the instruction class uniquely identifying the operation family. \emph{Example:} {\scriptsize\ttfamily \mnem{IADD.s32} \dst{r0}, \src{r1}, \src{r2}} (\textbf{SPIR-V:} \texttt{OpIAdd (operation selector)}; \textbf{GLSL:} \texttt{// Primary opcode bits})
  \item \texttt{instruction\_class}: 4-bit major instruction class selector (\texttt{bits 48:51}) partitioning the 64-bit Valhall opcode space (e.g., \texttt{FMA}, \texttt{CVT}, \texttt{SFU}, \texttt{LS}, \texttt{TEX}, \texttt{BRANCH}). \emph{Example:} {\scriptsize\ttfamily \mnem{FADD.f32} \dst{r0}, \src{r1}, \src{r2}} (\textbf{SPIR-V:} \texttt{N/A (encoding class)}; \textbf{GLSL:} \texttt{// Hardware execution pipe / decode class})
  \item \texttt{fau\_page\_index}: Selects the FAU (Fast Access Uniform) page (\texttt{0..3}) when a source operand references hardware small constants, thread IDs, or lane constants. \emph{Example:} {\scriptsize\ttfamily \mnem{FADD.f32} \dst{r0}, \src{r1}, \src{k1f}} (\textbf{SPIR-V:} \texttt{OpConstant / BuiltIn}; \textbf{GLSL:} \texttt{float y = x + 1.0;})
  \item \hyperref[enum:flow_control_m]{\texttt{flow}}: 4-bit embedded flow-control \& scoreboard wait modifier (\texttt{bits 59:62}: \texttt{none}, \texttt{wait0..wait0126}, \texttt{wait}, \texttt{reconverge}, \texttt{discard}, \texttt{end}). \emph{Example:} {\scriptsize\ttfamily \mnem{FADD.f32}\imod{.wait0.reconverge} \dst{r0}, \src{r1}, \src{r2}} (\textbf{SPIR-V:} \texttt{OpControlBarrier / OpReturn}; \textbf{GLSL:} \texttt{// Scoreboard wait + warp reconvergence})
\end{tightitemize}
\clearpage

\subsection*{LEA\_ATTR\_IMM}
\addcontentsline{toc}{subsection}{LEA\_ATTR\_IMM}
Resolves the Attribute descriptor at immediate \texttt{(table=\#attribute\_table\_index, index=\#attribute\_index)} for \texttt{(vertex=src0, instance=src1)} and writes 3 staging registers to \texttt{@sw}: \texttt{@sw[0:1]} holds the 64-bit attribute GPU VA and \texttt{@sw[2]} holds the 32-bit format conversion descriptor for \texttt{LD\_CVT}/\texttt{ST\_CVT}.
\par\textbf{Notes:}
\begin{tightitemize}
  \item Returns 3 staging registers (\texttt{ISA.xml:105}).
  \item \textbf{Source references:} \texttt{bifrost/valhall/ISA.xml:105}, \texttt{isa-v9-v14.xml}.
\end{tightitemize}
\codebox{\mnem{LEA\_ATTR\_IMM}\imod{\{.message\_slot\_index\}}\imod{\{.flow\}} \dst{@sw}, \src{src0}, \src{src1}, \imm{\#attribute\_index}, \imm{\#attribute\_table\_index}}
\textbf{Operands:}
\begin{tightitemize}\footnotesize
  \item \texttt{\dst{@sw}} (\texttt{staging\_register\_index} + \texttt{sr\_count}): 3-register destination range: \texttt{@sw[0:1]} = 64-bit attribute VA, \texttt{@sw[2]} = 32-bit format conversion descriptor.
  \item \texttt{\src{src0}}: Vertex index (\texttt{VertexIndex}).
  \item \texttt{\src{src1}}: Instance index (\texttt{InstanceIndex}).
  \item \texttt{\imm{\#attribute\_index}} (\texttt{attribute\_index}): Immediate index of the Attribute descriptor within the table.
  \item \texttt{\imm{\#attribute\_table\_index}} (\texttt{encoded\_attribute\_table\_index}): Immediate SRT table index holding the Attribute descriptor array.
\end{tightitemize}
\par\textbf{Closest SPIR-V Op:} \texttt{None (attribute address + cvt desc)}
\par\textbf{Pseudo-GLSL (Operational Semantics):}
\codebox{AttrDesc a = SRT[\#attribute\_table\_index].Attr[\#attribute\_index];\\@sw[0:1] = attrAddress(a, vert = src0, inst = src1); @sw[2] = a.format\_desc;}
\textbf{Execution Context:}
\begin{tightitemize}\footnotesize
  \item \textbf{Attribute descriptor chain (\hyperref[sec:csf-srt]{\texttt{IDVSSR::SRT\_0}}, \texttt{r0:r1}):} Masks 64B-aligned VA (\texttt{\& \textasciitilde{}0x3F}, \texttt{count=bits[5:0]}) $\to$ \texttt{Resource[table]} (16B) $\to$ \hyperref[sec:csf-buf-attr]{\texttt{Attribute[index]}} (32B: \texttt{table}, \texttt{buffer\_index}, \texttt{offset}, \texttt{stride}, \texttt{format}) $\to$ \texttt{Resource[Attribute::table]} $\to$ \texttt{Buffer[Attribute::buffer\_index]} (32B: \texttt{address}, \texttt{size}).
  \item \textbf{Vertex stage preloads (\hyperref[sec:csf-spd]{\texttt{SPD\_0}/\texttt{SPD\_1}}, \texttt{r16..r19}):} The hardware Index Fetcher reads \hyperref[sec:csf-regs]{\texttt{IDVSSR::INDEX\_BUFFER}} (\texttt{r54:r55}) and pre-fills \texttt{r59, r60} with \texttt{VertexIndex} (\texttt{+ VERTEX\_OFFSET + GLOBAL\_ATTR\_OFF}) and \texttt{r61} with \texttt{InstanceIndex} (\texttt{+ INSTANCE\_OFFSET}) so Position and Varying shaders can pull vertex attributes directly.
\end{tightitemize}
\par\textbf{Worked Examples:}
\par\textbf{Example 1 (Static varying output address (table 0, index 1)):}
\codebox{\mnem{LEA\_ATTR\_IMM}\imod{.slot0} \dst{@r0:r2}, \src{r59}, \imm{\#0}, \imm{\#1}}
Computes 64-bit output varying address \texttt{[r0:r1]} + conversion descriptor \texttt{r2} for immediate \texttt{(table=0, index=1)}.
\par\textbf{Example 2 (Vertex shader position output (table 0, index 0)):}
\codebox{\mnem{LEA\_ATTR\_IMM}\imod{.slot0} \dst{@r0:r2}, \src{r59}, \imm{\#0}, \imm{\#0}}
Resolves the per-vertex position buffer slot before \texttt{ST\_CVT.f32.v4}.

\begin{center}\footnotesize
\begin{tabularx}{\linewidth}{@{}lllX@{}}
\toprule
Variant & Arch & Class & Opcode Fields \\
\midrule
\texttt{(none)} & v9--v14 & \texttt{S2} & \texttt{opcode1=LEA\_ATTR\_IMM} \\
\bottomrule
\end{tabularx}
\end{center}
\begin{center}\footnotesize
\begin{tabularx}{0.88\linewidth}{@{}l l X@{}}
\toprule
Field Name & Bits & Type / Fixed Value \\
\midrule
\texttt{src0} & \texttt{0:7} & \texttt{SourceEncoding} \\
\texttt{src1} & \texttt{8:15} & \texttt{SourceEncoding} \\
\texttt{encoded\_attribute\_table\_index} & \texttt{16:19} & \texttt{uint} \\
\texttt{attribute\_index} & \texttt{20:23} & \texttt{uint} \\
\textit{Reserved} & \texttt{24:29} & \texttt{zero} \\
\hyperref[enum:message_slot_index_m]{\texttt{message\_slot\_index}} & \texttt{30:32} & \texttt{message\_slot\_index\_m} \\
\texttt{sr\_count} & \texttt{33:35} & \texttt{exact=count3} \\
\textit{Reserved} & \texttt{36} & \texttt{zero} \\
\texttt{descriptor\_type} & \texttt{37:38} & \texttt{exact=attribute} \\
\textit{Reserved} & \texttt{39} & \texttt{zero} \\
\texttt{staging\_register\_index} & \texttt{40:45} & \texttt{uint} \\
\texttt{staging\_register\_control\_32} & \texttt{46:47} & \texttt{exact=write} \\
\texttt{opcode1} & \texttt{48:51} & \texttt{exact=LEA\_ATTR\_IMM} \\
\texttt{instruction\_class} & \texttt{52:56} & \texttt{exact=S2} \\
\texttt{fau\_page\_index} & \texttt{57:58} & \texttt{uint} \\
\hyperref[enum:flow_control_m]{\texttt{flow}} & \texttt{59:62} & \texttt{flow\_control\_m} \\
\textit{Reserved} & \texttt{63} & \texttt{zero} \\
\bottomrule
\end{tabularx}
\end{center}
\textbf{Fields Explained:}
\begin{tightitemize}\footnotesize
  \item \texttt{src0}: Primary source operand (\texttt{r0..r63}, \texttt{u0..u63} uniform, \texttt{k*} FAU constant, or \texttt{t0..t2} temporary; \texttt{\textasciicircum{}} marks discard/last-use). \emph{Example:} {\scriptsize\ttfamily \mnem{FADD.f32} \dst{r0}, \src{r1}, \src{r2}} (\textbf{SPIR-V:} \texttt{OpFAdd (Operand 0)}; \textbf{GLSL:} \texttt{float c = a + b;})
  \item \texttt{src1}: Secondary source operand (register, uniform, or FAU constant; in shifts, supplies the shift count byte). \emph{Example:} {\scriptsize\ttfamily \mnem{LSHIFT.i32} \dst{r0}, \src{r1}, \src{r2}\imod{.b0}} (\textbf{SPIR-V:} \texttt{OpShiftLeftLogical (Shift)}; \textbf{GLSL:} \texttt{uint y = a \textless{}\textless{} s;})
  \item \texttt{encoded\_attribute\_table\_index}: Encoded descriptor table index for \textbackslash\{\}texttt\{encoded\textbackslash\{\}\_attribute\textbackslash\{\}\_table\textbackslash\{\}\_index\}. \emph{Example:} {\scriptsize\ttfamily \mnem{LD\_BUFFER.i32} \dst{@r0}, \src{r2}, \imm{\#0}} (\textbf{SPIR-V:} \texttt{OpAccessChain (DescriptorSet)}; \textbf{GLSL:} \texttt{uint v = buf[i];})
  \item \texttt{attribute\_index}: Resource/slot index field \textbackslash\{\}texttt\{attribute\textbackslash\{\}\_index\} selecting the target descriptor or slot. \emph{Example:} {\scriptsize\ttfamily \mnem{LD\_VAR\_BUF\_IMM.f32} \dst{@r0:r3}, \imm{\#0}} (\textbf{SPIR-V:} \texttt{OpLoad (Descriptor)}; \textbf{GLSL:} \texttt{vec4 v = v\_in;})
  \item \hyperref[enum:message_slot_index_m]{\texttt{message\_slot\_index}}: Asynchronous scoreboard slot (\texttt{slot0}--\texttt{slot2}, \texttt{slot7}) tracking completion of this message. \emph{Example:} {\scriptsize\ttfamily \mnem{LOAD.i32}\imod{.slot0} \dst{@r0}, \src{r2}, \imm{\#0}} (\textbf{SPIR-V:} \texttt{OpLoad}; \textbf{GLSL:} \texttt{uint v = ssbo.arr[i];})
  \item \texttt{sr\_count}: Number of consecutive 32-bit staging registers transferred by a load/store/texture message. \emph{Example:} {\scriptsize\ttfamily \mnem{LD\_MULTI.i128} \dst{@r0:r3}, \src{r4}, \imm{\#0}} (\textbf{SPIR-V:} \texttt{OpLoad (wide vector)}; \textbf{GLSL:} \texttt{uvec4 x = buf[i];})
  \item \texttt{descriptor\_type}: Selects descriptor table binding class (attribute, buffer, texture, varying). \emph{Example:} {\scriptsize\ttfamily \mnem{LD\_BUFFER.i32} \dst{@r0}, \src{r2}, \imm{\#0}} (\textbf{SPIR-V:} \texttt{OpLoad (StorageBuffer)}; \textbf{GLSL:} \texttt{uint v = ssbo.arr[i];})
  \item \texttt{staging\_register\_index}: Base GPR index of the contiguous staging register window (\texttt{@rN:rM}) for message payload/return. \emph{Example:} {\scriptsize\ttfamily \mnem{STORE.i64}\imod{.slot7} \src{@r0:r1}, \src{r2}, \imm{\#8}} (\textbf{SPIR-V:} \texttt{OpStore (uvec2)}; \textbf{GLSL:} \texttt{ssbo.pairs[i] = uvec2(a, b);})
  \item \texttt{staging\_register\_control\_32}: Controls the read/write staging register bank and count for async message instructions. \emph{Example:} {\scriptsize\ttfamily \mnem{LOAD.i128}\imod{.slot0} \dst{@r0:r3}, \src{r4}, \imm{\#0}} (\textbf{SPIR-V:} \texttt{OpLoad (vec4)}; \textbf{GLSL:} \texttt{vec4 v = ssbo.data[idx];})
  \item \texttt{opcode1}: Primary sub-opcode within the instruction class uniquely identifying the operation family. \emph{Example:} {\scriptsize\ttfamily \mnem{IADD.s32} \dst{r0}, \src{r1}, \src{r2}} (\textbf{SPIR-V:} \texttt{OpIAdd (operation selector)}; \textbf{GLSL:} \texttt{// Primary opcode bits})
  \item \texttt{instruction\_class}: 4-bit major instruction class selector (\texttt{bits 48:51}) partitioning the 64-bit Valhall opcode space (e.g., \texttt{FMA}, \texttt{CVT}, \texttt{SFU}, \texttt{LS}, \texttt{TEX}, \texttt{BRANCH}). \emph{Example:} {\scriptsize\ttfamily \mnem{FADD.f32} \dst{r0}, \src{r1}, \src{r2}} (\textbf{SPIR-V:} \texttt{N/A (encoding class)}; \textbf{GLSL:} \texttt{// Hardware execution pipe / decode class})
  \item \texttt{fau\_page\_index}: Selects the FAU (Fast Access Uniform) page (\texttt{0..3}) when a source operand references hardware small constants, thread IDs, or lane constants. \emph{Example:} {\scriptsize\ttfamily \mnem{FADD.f32} \dst{r0}, \src{r1}, \src{k1f}} (\textbf{SPIR-V:} \texttt{OpConstant / BuiltIn}; \textbf{GLSL:} \texttt{float y = x + 1.0;})
  \item \hyperref[enum:flow_control_m]{\texttt{flow}}: 4-bit embedded flow-control \& scoreboard wait modifier (\texttt{bits 59:62}: \texttt{none}, \texttt{wait0..wait0126}, \texttt{wait}, \texttt{reconverge}, \texttt{discard}, \texttt{end}). \emph{Example:} {\scriptsize\ttfamily \mnem{FADD.f32}\imod{.wait0.reconverge} \dst{r0}, \src{r1}, \src{r2}} (\textbf{SPIR-V:} \texttt{OpControlBarrier / OpReturn}; \textbf{GLSL:} \texttt{// Scoreboard wait + warp reconvergence})
\end{tightitemize}
\clearpage

\subsection*{LEA\_BUF}
\addcontentsline{toc}{subsection}{LEA\_BUF}
Resolves the Buffer descriptor identified by dynamic handle \texttt{src1}, adds byte offset (or vertex index) \texttt{src0}, and writes the resulting 64-bit GPU virtual address into the 2-register staging pair \texttt{@sw[0:1]}.
\par\textbf{Notes:}
\begin{tightitemize}
  \item Returns 64-bit address in 2 staging registers.
  \item \textbf{Source references:} \texttt{kraid/ops.rs:3161-3185}, \texttt{bifrost/valhall/ISA.xml:107}.
\end{tightitemize}
\codebox{\mnem{LEA\_BUF}\imod{\{.message\_slot\_index\}}\imod{\{.flow\}} \dst{@sw}, \src{src0}, \src{src1}}
\textbf{Operands:}
\begin{tightitemize}\footnotesize
  \item \texttt{\dst{@sw}} (\texttt{staging\_register\_index} + \texttt{sr\_count}): 2-register destination range (\texttt{@sw[0:1]}) receiving the 64-bit effective buffer GPU VA.
  \item \texttt{\src{src0}}: 32-bit byte offset (or vertex index for IDVS output buffers).
  \item \texttt{\src{src1}}: Dynamic buffer descriptor handle.
\end{tightitemize}
\par\textbf{Closest SPIR-V Op:} \texttt{OpAccessChain (dynamic SSBO base)}
\par\textbf{Pseudo-GLSL (Operational Semantics):}
\codebox{@sw[0:1] = resolveBufferDesc(src1).address + uint64\_t(src0);}
\textbf{Execution Context:}
\begin{tightitemize}\footnotesize
  \item \textbf{Buffer descriptor chain (\hyperref[sec:csf-srt]{\texttt{SRT\_k}}, \texttt{r0:r1..r6:r7}):} Masks 64B-aligned VA (\texttt{\& \textasciitilde{}0x3F}, \texttt{count=bits[5:0]}) $\to$ \texttt{Resource[table]} (16B, \texttt{contains\_descriptors=1}) $\to$ \hyperref[sec:csf-buf-attr]{\texttt{Buffer[index]}} (32B, \texttt{type=9}: 64-bit \texttt{address}, \texttt{size}). Precondition: \texttt{table < count} and \texttt{index*32 < Resource[table].size}.
  \item \textbf{Inline FAU \& Special Operands (\hyperref[sec:csf-tsd-fau]{\texttt{FAU\_k}}, \texttt{r8..r15}):} Any 8-bit source with \texttt{mode=2} reads zero-latency uniform \texttt{u0..u31} from \texttt{FAU\_k} (\texttt{slots64 = bits[63:56]}); \texttt{mode=3} reads \texttt{SmallConstant} (Table~1.2) or \texttt{FauSpecial} (\texttt{pc}, \texttt{lane\_id}, \texttt{lane\_width}, \texttt{tls\_ptr}, \texttt{wls\_ptr}, \texttt{blend\_descriptor\_0..7}).
\end{tightitemize}
\par\textbf{Worked Examples:}
\par\textbf{Example 1 (Dynamic SSBO address resolution):}
\codebox{\mnem{LEA\_BUF}\imod{.slot0} \dst{@r0:r1}, \src{r2}, \src{r3}}
Computes 64-bit GPU VA \texttt{@r0:r1} for dynamic buffer descriptor handle \texttt{r3} plus byte offset \texttt{r2}.
\par\textbf{Example 2 (Bindless SSBO array indexing):}
\codebox{\mnem{LEA\_BUF}\imod{.slot1} \dst{@r4:r5}, \src{r6}, \src{r7}}
Resolves address of \texttt{ssbos[nonuniformEXT(i)]} for subsequent \texttt{LOAD}/\texttt{STORE}/\texttt{ATOM}.

\begin{center}\footnotesize
\begin{tabularx}{\linewidth}{@{}lllX@{}}
\toprule
Variant & Arch & Class & Opcode Fields \\
\midrule
\texttt{(none)} & v9--v14 & \texttt{S2} & \texttt{opcode1=LEA\_BUF} \\
\bottomrule
\end{tabularx}
\end{center}
\begin{center}\footnotesize
\begin{tabularx}{0.88\linewidth}{@{}l l X@{}}
\toprule
Field Name & Bits & Type / Fixed Value \\
\midrule
\texttt{src0} & \texttt{0:7} & \texttt{SourceEncoding} \\
\texttt{src1} & \texttt{8:15} & \texttt{SourceEncoding} \\
\textit{Reserved} & \texttt{16:29} & \texttt{zero} \\
\hyperref[enum:message_slot_index_m]{\texttt{message\_slot\_index}} & \texttt{30:32} & \texttt{message\_slot\_index\_m} \\
\texttt{sr\_count} & \texttt{33:35} & \texttt{exact=count2} \\
\textit{Reserved} & \texttt{36:39} & \texttt{zero} \\
\texttt{staging\_register\_index} & \texttt{40:45} & \texttt{uint} \\
\texttt{staging\_register\_control\_32} & \texttt{46:47} & \texttt{exact=write} \\
\texttt{opcode1} & \texttt{48:51} & \texttt{exact=LEA\_BUF} \\
\texttt{instruction\_class} & \texttt{52:56} & \texttt{exact=S2} \\
\texttt{fau\_page\_index} & \texttt{57:58} & \texttt{uint} \\
\hyperref[enum:flow_control_m]{\texttt{flow}} & \texttt{59:62} & \texttt{flow\_control\_m} \\
\textit{Reserved} & \texttt{63} & \texttt{zero} \\
\bottomrule
\end{tabularx}
\end{center}
\textbf{Fields Explained:}
\begin{tightitemize}\footnotesize
  \item \texttt{src0}: Primary source operand (\texttt{r0..r63}, \texttt{u0..u63} uniform, \texttt{k*} FAU constant, or \texttt{t0..t2} temporary; \texttt{\textasciicircum{}} marks discard/last-use). \emph{Example:} {\scriptsize\ttfamily \mnem{FADD.f32} \dst{r0}, \src{r1}, \src{r2}} (\textbf{SPIR-V:} \texttt{OpFAdd (Operand 0)}; \textbf{GLSL:} \texttt{float c = a + b;})
  \item \texttt{src1}: Secondary source operand (register, uniform, or FAU constant; in shifts, supplies the shift count byte). \emph{Example:} {\scriptsize\ttfamily \mnem{LSHIFT.i32} \dst{r0}, \src{r1}, \src{r2}\imod{.b0}} (\textbf{SPIR-V:} \texttt{OpShiftLeftLogical (Shift)}; \textbf{GLSL:} \texttt{uint y = a \textless{}\textless{} s;})
  \item \hyperref[enum:message_slot_index_m]{\texttt{message\_slot\_index}}: Asynchronous scoreboard slot (\texttt{slot0}--\texttt{slot2}, \texttt{slot7}) tracking completion of this message. \emph{Example:} {\scriptsize\ttfamily \mnem{LOAD.i32}\imod{.slot0} \dst{@r0}, \src{r2}, \imm{\#0}} (\textbf{SPIR-V:} \texttt{OpLoad}; \textbf{GLSL:} \texttt{uint v = ssbo.arr[i];})
  \item \texttt{sr\_count}: Number of consecutive 32-bit staging registers transferred by a load/store/texture message. \emph{Example:} {\scriptsize\ttfamily \mnem{LD\_MULTI.i128} \dst{@r0:r3}, \src{r4}, \imm{\#0}} (\textbf{SPIR-V:} \texttt{OpLoad (wide vector)}; \textbf{GLSL:} \texttt{uvec4 x = buf[i];})
  \item \texttt{staging\_register\_index}: Base GPR index of the contiguous staging register window (\texttt{@rN:rM}) for message payload/return. \emph{Example:} {\scriptsize\ttfamily \mnem{STORE.i64}\imod{.slot7} \src{@r0:r1}, \src{r2}, \imm{\#8}} (\textbf{SPIR-V:} \texttt{OpStore (uvec2)}; \textbf{GLSL:} \texttt{ssbo.pairs[i] = uvec2(a, b);})
  \item \texttt{staging\_register\_control\_32}: Controls the read/write staging register bank and count for async message instructions. \emph{Example:} {\scriptsize\ttfamily \mnem{LOAD.i128}\imod{.slot0} \dst{@r0:r3}, \src{r4}, \imm{\#0}} (\textbf{SPIR-V:} \texttt{OpLoad (vec4)}; \textbf{GLSL:} \texttt{vec4 v = ssbo.data[idx];})
  \item \texttt{opcode1}: Primary sub-opcode within the instruction class uniquely identifying the operation family. \emph{Example:} {\scriptsize\ttfamily \mnem{IADD.s32} \dst{r0}, \src{r1}, \src{r2}} (\textbf{SPIR-V:} \texttt{OpIAdd (operation selector)}; \textbf{GLSL:} \texttt{// Primary opcode bits})
  \item \texttt{instruction\_class}: 4-bit major instruction class selector (\texttt{bits 48:51}) partitioning the 64-bit Valhall opcode space (e.g., \texttt{FMA}, \texttt{CVT}, \texttt{SFU}, \texttt{LS}, \texttt{TEX}, \texttt{BRANCH}). \emph{Example:} {\scriptsize\ttfamily \mnem{FADD.f32} \dst{r0}, \src{r1}, \src{r2}} (\textbf{SPIR-V:} \texttt{N/A (encoding class)}; \textbf{GLSL:} \texttt{// Hardware execution pipe / decode class})
  \item \texttt{fau\_page\_index}: Selects the FAU (Fast Access Uniform) page (\texttt{0..3}) when a source operand references hardware small constants, thread IDs, or lane constants. \emph{Example:} {\scriptsize\ttfamily \mnem{FADD.f32} \dst{r0}, \src{r1}, \src{k1f}} (\textbf{SPIR-V:} \texttt{OpConstant / BuiltIn}; \textbf{GLSL:} \texttt{float y = x + 1.0;})
  \item \hyperref[enum:flow_control_m]{\texttt{flow}}: 4-bit embedded flow-control \& scoreboard wait modifier (\texttt{bits 59:62}: \texttt{none}, \texttt{wait0..wait0126}, \texttt{wait}, \texttt{reconverge}, \texttt{discard}, \texttt{end}). \emph{Example:} {\scriptsize\ttfamily \mnem{FADD.f32}\imod{.wait0.reconverge} \dst{r0}, \src{r1}, \src{r2}} (\textbf{SPIR-V:} \texttt{OpControlBarrier / OpReturn}; \textbf{GLSL:} \texttt{// Scoreboard wait + warp reconvergence})
\end{tightitemize}
\clearpage

\subsection*{LEA\_BUF\_IMM}
\addcontentsline{toc}{subsection}{LEA\_BUF\_IMM}
Resolves the Buffer descriptor at immediate \texttt{(table=\#buffer\_table\_index, index=\#buffer\_index)}, adds byte offset (or vertex index) \texttt{src0}, and writes the resulting 64-bit GPU virtual address into \texttt{@sw[0:1]}.
\par\textbf{Notes:}
\begin{tightitemize}
  \item Immediate-handle form of \texttt{LEA\_BUF}.
  \item \textbf{Source references:} \texttt{bifrost/valhall/ISA.xml:109}, \texttt{isa-v9-v14.xml}.
\end{tightitemize}
\codebox{\mnem{LEA\_BUF\_IMM}\imod{\{.message\_slot\_index\}}\imod{\{.flow\}} \dst{@sw}, \src{src0}, \imm{\#buffer\_index}, \imm{\#buffer\_table\_index}}
\textbf{Operands:}
\begin{tightitemize}\footnotesize
  \item \texttt{\dst{@sw}} (\texttt{staging\_register\_index} + \texttt{sr\_count}): 2-register destination range (\texttt{@sw[0:1]}) receiving the 64-bit effective buffer GPU VA.
  \item \texttt{\src{src0}}: 32-bit byte offset (or vertex index for IDVS output buffers).
  \item \texttt{\imm{\#buffer\_index}} (\texttt{buffer\_index}): Immediate index of the Buffer descriptor within the table.
  \item \texttt{\imm{\#buffer\_table\_index}} (\texttt{encoded\_buffer\_table\_index}): Immediate SRT table index holding the Buffer descriptor array.
\end{tightitemize}
\par\textbf{Closest SPIR-V Op:} \texttt{OpAccessChain (static SSBO base)}
\par\textbf{Pseudo-GLSL (Operational Semantics):}
\codebox{@sw[0:1] = SRT[\#buffer\_table\_index].Buffer[\#buffer\_index].address + uint64\_t(src0);}
\textbf{Execution Context:}
\begin{tightitemize}\footnotesize
  \item \textbf{Buffer descriptor chain (\hyperref[sec:csf-srt]{\texttt{SRT\_k}}, \texttt{r0:r1..r6:r7}):} Masks 64B-aligned VA (\texttt{\& \textasciitilde{}0x3F}, \texttt{count=bits[5:0]}) $\to$ \texttt{Resource[table]} (16B, \texttt{contains\_descriptors=1}) $\to$ \hyperref[sec:csf-buf-attr]{\texttt{Buffer[index]}} (32B, \texttt{type=9}: 64-bit \texttt{address}, \texttt{size}). Precondition: \texttt{table < count} and \texttt{index*32 < Resource[table].size}.
  \item \textbf{Inline FAU \& Special Operands (\hyperref[sec:csf-tsd-fau]{\texttt{FAU\_k}}, \texttt{r8..r15}):} Any 8-bit source with \texttt{mode=2} reads zero-latency uniform \texttt{u0..u31} from \texttt{FAU\_k} (\texttt{slots64 = bits[63:56]}); \texttt{mode=3} reads \texttt{SmallConstant} (Table~1.2) or \texttt{FauSpecial} (\texttt{pc}, \texttt{lane\_id}, \texttt{lane\_width}, \texttt{tls\_ptr}, \texttt{wls\_ptr}, \texttt{blend\_descriptor\_0..7}).
\end{tightitemize}
\par\textbf{Worked Examples:}
\par\textbf{Example 1 (Static SSBO base + dynamic byte offset):}
\codebox{\mnem{LEA\_BUF\_IMM}\imod{.slot0} \dst{@r0:r1}, \src{r2}, \imm{\#0}, \imm{\#0}}
Computes 64-bit effective address of immediate buffer \texttt{(table=0, index=0)} + byte offset \texttt{r2} into \texttt{@r0:r1}.
\par\textbf{Example 2 (Second descriptor binding (`table=0, index=2`)):}
\codebox{\mnem{LEA\_BUF\_IMM}\imod{.slot0} \dst{@r2:r3}, \src{k0}, \imm{\#0}, \imm{\#2}}
With offset \texttt{k0 = 0}: retrieves the base 64-bit GPU pointer of buffer binding 2.
\par\textbf{Mined SPIR-V $\to$ Valhall ISA Example (Computing IDVS varying-buffer output address in a Vertex shader --- Vertex shader, Mali-G710 (v10)):}
\codebox{\textbf{; SPIR-V (input)}\\OpStore \%out\_uv \%attr\_uv\\[3pt]\textbf{; Valhall ISA (malioc output)}\\LD\_ATTR\_IMM.slot0.v2.f32 @r0:r1, r60\textasciicircum{}, r61\textasciicircum{}, \#0x1, \#0x1\\LEA\_BUF\_IMM.slot0 @r4:r5, r59\textasciicircum{}, \#0x0, \#0x3d\\STORE.i128.slot0.end.estream @r0:r1:r2:r3, [r4\textasciicircum{}:r5\textasciicircum{}], \#0x0}
In Mali IDVS vertex/varying shaders, \texttt{LEA\_BUF\_IMM} resolves buffer descriptor \texttt{(table=\#0x3d, index=\#0x0)} for vertex index \texttt{r59\textasciicircum{}} into a 64-bit GPU VA in \texttt{@r4:r5}, which \texttt{STORE.i128.estream} writes to.
\begin{center}\footnotesize
\begin{tabularx}{\linewidth}{@{}lllX@{}}
\toprule
Variant & Arch & Class & Opcode Fields \\
\midrule
\texttt{(none)} & v9--v14 & \texttt{S1} & \texttt{opcode1=LEA\_BUF\_IMM} \\
\bottomrule
\end{tabularx}
\end{center}
\begin{center}\footnotesize
\begin{tabularx}{0.88\linewidth}{@{}l l X@{}}
\toprule
Field Name & Bits & Type / Fixed Value \\
\midrule
\texttt{src0} & \texttt{0:7} & \texttt{SourceEncoding} \\
\texttt{encoded\_buffer\_table\_index} & \texttt{8:11} & \texttt{uint} \\
\texttt{buffer\_index} & \texttt{12:19} & \texttt{uint} \\
\textit{Reserved} & \texttt{20:29} & \texttt{zero} \\
\hyperref[enum:message_slot_index_m]{\texttt{message\_slot\_index}} & \texttt{30:32} & \texttt{message\_slot\_index\_m} \\
\texttt{sr\_count} & \texttt{33:35} & \texttt{exact=count2} \\
\textit{Reserved} & \texttt{36:39} & \texttt{zero} \\
\texttt{staging\_register\_index} & \texttt{40:45} & \texttt{uint} \\
\texttt{staging\_register\_control\_32} & \texttt{46:47} & \texttt{exact=write} \\
\texttt{opcode1} & \texttt{48:51} & \texttt{exact=LEA\_BUF\_IMM} \\
\texttt{instruction\_class} & \texttt{52:56} & \texttt{exact=S1} \\
\texttt{fau\_page\_index} & \texttt{57:58} & \texttt{uint} \\
\hyperref[enum:flow_control_m]{\texttt{flow}} & \texttt{59:62} & \texttt{flow\_control\_m} \\
\textit{Reserved} & \texttt{63} & \texttt{zero} \\
\bottomrule
\end{tabularx}
\end{center}
\textbf{Fields Explained:}
\begin{tightitemize}\footnotesize
  \item \texttt{src0}: Primary source operand (\texttt{r0..r63}, \texttt{u0..u63} uniform, \texttt{k*} FAU constant, or \texttt{t0..t2} temporary; \texttt{\textasciicircum{}} marks discard/last-use). \emph{Example:} {\scriptsize\ttfamily \mnem{FADD.f32} \dst{r0}, \src{r1}, \src{r2}} (\textbf{SPIR-V:} \texttt{OpFAdd (Operand 0)}; \textbf{GLSL:} \texttt{float c = a + b;})
  \item \texttt{encoded\_buffer\_table\_index}: Encoded descriptor table index for \textbackslash\{\}texttt\{encoded\textbackslash\{\}\_buffer\textbackslash\{\}\_table\textbackslash\{\}\_index\}. \emph{Example:} {\scriptsize\ttfamily \mnem{LD\_BUFFER.i32} \dst{@r0}, \src{r2}, \imm{\#0}} (\textbf{SPIR-V:} \texttt{OpAccessChain (DescriptorSet)}; \textbf{GLSL:} \texttt{uint v = buf[i];})
  \item \texttt{buffer\_index}: Resource/slot index field \textbackslash\{\}texttt\{buffer\textbackslash\{\}\_index\} selecting the target descriptor or slot. \emph{Example:} {\scriptsize\ttfamily \mnem{LD\_VAR\_BUF\_IMM.f32} \dst{@r0:r3}, \imm{\#0}} (\textbf{SPIR-V:} \texttt{OpLoad (Descriptor)}; \textbf{GLSL:} \texttt{vec4 v = v\_in;})
  \item \hyperref[enum:message_slot_index_m]{\texttt{message\_slot\_index}}: Asynchronous scoreboard slot (\texttt{slot0}--\texttt{slot2}, \texttt{slot7}) tracking completion of this message. \emph{Example:} {\scriptsize\ttfamily \mnem{LOAD.i32}\imod{.slot0} \dst{@r0}, \src{r2}, \imm{\#0}} (\textbf{SPIR-V:} \texttt{OpLoad}; \textbf{GLSL:} \texttt{uint v = ssbo.arr[i];})
  \item \texttt{sr\_count}: Number of consecutive 32-bit staging registers transferred by a load/store/texture message. \emph{Example:} {\scriptsize\ttfamily \mnem{LD\_MULTI.i128} \dst{@r0:r3}, \src{r4}, \imm{\#0}} (\textbf{SPIR-V:} \texttt{OpLoad (wide vector)}; \textbf{GLSL:} \texttt{uvec4 x = buf[i];})
  \item \texttt{staging\_register\_index}: Base GPR index of the contiguous staging register window (\texttt{@rN:rM}) for message payload/return. \emph{Example:} {\scriptsize\ttfamily \mnem{STORE.i64}\imod{.slot7} \src{@r0:r1}, \src{r2}, \imm{\#8}} (\textbf{SPIR-V:} \texttt{OpStore (uvec2)}; \textbf{GLSL:} \texttt{ssbo.pairs[i] = uvec2(a, b);})
  \item \texttt{staging\_register\_control\_32}: Controls the read/write staging register bank and count for async message instructions. \emph{Example:} {\scriptsize\ttfamily \mnem{LOAD.i128}\imod{.slot0} \dst{@r0:r3}, \src{r4}, \imm{\#0}} (\textbf{SPIR-V:} \texttt{OpLoad (vec4)}; \textbf{GLSL:} \texttt{vec4 v = ssbo.data[idx];})
  \item \texttt{opcode1}: Primary sub-opcode within the instruction class uniquely identifying the operation family. \emph{Example:} {\scriptsize\ttfamily \mnem{IADD.s32} \dst{r0}, \src{r1}, \src{r2}} (\textbf{SPIR-V:} \texttt{OpIAdd (operation selector)}; \textbf{GLSL:} \texttt{// Primary opcode bits})
  \item \texttt{instruction\_class}: 4-bit major instruction class selector (\texttt{bits 48:51}) partitioning the 64-bit Valhall opcode space (e.g., \texttt{FMA}, \texttt{CVT}, \texttt{SFU}, \texttt{LS}, \texttt{TEX}, \texttt{BRANCH}). \emph{Example:} {\scriptsize\ttfamily \mnem{FADD.f32} \dst{r0}, \src{r1}, \src{r2}} (\textbf{SPIR-V:} \texttt{N/A (encoding class)}; \textbf{GLSL:} \texttt{// Hardware execution pipe / decode class})
  \item \texttt{fau\_page\_index}: Selects the FAU (Fast Access Uniform) page (\texttt{0..3}) when a source operand references hardware small constants, thread IDs, or lane constants. \emph{Example:} {\scriptsize\ttfamily \mnem{FADD.f32} \dst{r0}, \src{r1}, \src{k1f}} (\textbf{SPIR-V:} \texttt{OpConstant / BuiltIn}; \textbf{GLSL:} \texttt{float y = x + 1.0;})
  \item \hyperref[enum:flow_control_m]{\texttt{flow}}: 4-bit embedded flow-control \& scoreboard wait modifier (\texttt{bits 59:62}: \texttt{none}, \texttt{wait0..wait0126}, \texttt{wait}, \texttt{reconverge}, \texttt{discard}, \texttt{end}). \emph{Example:} {\scriptsize\ttfamily \mnem{FADD.f32}\imod{.wait0.reconverge} \dst{r0}, \src{r1}, \src{r2}} (\textbf{SPIR-V:} \texttt{OpControlBarrier / OpReturn}; \textbf{GLSL:} \texttt{// Scoreboard wait + warp reconvergence})
\end{tightitemize}
\clearpage

\subsection*{LEA\_PKA}
\addcontentsline{toc}{subsection}{LEA\_PKA}
Resolves the simple/PKA buffer descriptor identified by handle \texttt{src1}, adds byte offset \texttt{src0}, and writes the resulting 64-bit GPU virtual address into the 2-register staging pair \texttt{@sw[0:1]}.
\par\textbf{Notes:}
\begin{tightitemize}
  \item Address-generation companion to \texttt{LD\_PKA}.
  \item \textbf{Source references:} \texttt{kraid/ops.rs:3187-3211}, \texttt{bifrost/valhall/ISA.xml:111}.
\end{tightitemize}
\codebox{\mnem{LEA\_PKA}\imod{\{.message\_slot\_index\}}\imod{\{.flow\}} \dst{@sw}, \src{src0}, \src{src1}}
\textbf{Operands:}
\begin{tightitemize}\footnotesize
  \item \texttt{\dst{@sw}} (\texttt{staging\_register\_index} + \texttt{sr\_count}): 2-register destination range (\texttt{@sw[0:1]}) receiving the 64-bit effective PKA buffer GPU VA.
  \item \texttt{\src{src0}}: 32-bit byte offset within the buffer.
  \item \texttt{\src{src1}}: 32-bit PKA buffer descriptor handle.
\end{tightitemize}
\par\textbf{Closest SPIR-V Op:} \texttt{OpAccessChain (PKA buffer address)}
\par\textbf{Pseudo-GLSL (Operational Semantics):}
\codebox{@sw[0:1] = resolvePkaBuffer(src1).address + uint64\_t(src0);}
\textbf{Execution Context:}
\begin{tightitemize}\footnotesize
  \item \textbf{Buffer descriptor chain (\hyperref[sec:csf-srt]{\texttt{SRT\_k}}, \texttt{r0:r1..r6:r7}):} Masks 64B-aligned VA (\texttt{\& \textasciitilde{}0x3F}, \texttt{count=bits[5:0]}) $\to$ \texttt{Resource[table]} (16B, \texttt{contains\_descriptors=1}) $\to$ \hyperref[sec:csf-buf-attr]{\texttt{Buffer[index]}} (32B, \texttt{type=9}: 64-bit \texttt{address}, \texttt{size}). Precondition: \texttt{table < count} and \texttt{index*32 < Resource[table].size}.
  \item \textbf{Inline FAU \& Special Operands (\hyperref[sec:csf-tsd-fau]{\texttt{FAU\_k}}, \texttt{r8..r15}):} Any 8-bit source with \texttt{mode=2} reads zero-latency uniform \texttt{u0..u31} from \texttt{FAU\_k} (\texttt{slots64 = bits[63:56]}); \texttt{mode=3} reads \texttt{SmallConstant} (Table~1.2) or \texttt{FauSpecial} (\texttt{pc}, \texttt{lane\_id}, \texttt{lane\_width}, \texttt{tls\_ptr}, \texttt{wls\_ptr}, \texttt{blend\_descriptor\_0..7}).
\end{tightitemize}
\par\textbf{Worked Examples:}
\par\textbf{Example 1 (PKA buffer address resolution):}
\codebox{\mnem{LEA\_PKA}\imod{.slot0} \dst{@r0:r1}, \src{r2}, \src{r3}}
Computes 64-bit effective address \texttt{@r0:r1} for PKA buffer handle \texttt{r3} + byte offset \texttt{r2}.
\par\textbf{Example 2 (Zero-offset base pointer query):}
\codebox{\mnem{LEA\_PKA}\imod{.slot0} \dst{@r0:r1}, \src{k0}, \src{u0}}
Extracts the 64-bit base pointer of uniform-supplied PKA handle \texttt{u0} into \texttt{@r0:r1}.

\begin{center}\footnotesize
\begin{tabularx}{\linewidth}{@{}lllX@{}}
\toprule
Variant & Arch & Class & Opcode Fields \\
\midrule
\texttt{(none)} & v9--v14 & \texttt{S2} & \texttt{opcode1=LEA\_PKA} \\
\bottomrule
\end{tabularx}
\end{center}
\begin{center}\footnotesize
\begin{tabularx}{0.88\linewidth}{@{}l l X@{}}
\toprule
Field Name & Bits & Type / Fixed Value \\
\midrule
\texttt{src0} & \texttt{0:7} & \texttt{SourceEncoding} \\
\texttt{src1} & \texttt{8:15} & \texttt{SourceEncoding} \\
\textit{Reserved} & \texttt{16:29} & \texttt{zero} \\
\hyperref[enum:message_slot_index_m]{\texttt{message\_slot\_index}} & \texttt{30:32} & \texttt{message\_slot\_index\_m} \\
\texttt{sr\_count} & \texttt{33:35} & \texttt{exact=count2} \\
\textit{Reserved} & \texttt{36:39} & \texttt{zero} \\
\texttt{staging\_register\_index} & \texttt{40:45} & \texttt{uint} \\
\texttt{staging\_register\_control\_32} & \texttt{46:47} & \texttt{exact=write} \\
\texttt{opcode1} & \texttt{48:51} & \texttt{exact=LEA\_PKA} \\
\texttt{instruction\_class} & \texttt{52:56} & \texttt{exact=S2} \\
\texttt{fau\_page\_index} & \texttt{57:58} & \texttt{uint} \\
\hyperref[enum:flow_control_m]{\texttt{flow}} & \texttt{59:62} & \texttt{flow\_control\_m} \\
\textit{Reserved} & \texttt{63} & \texttt{zero} \\
\bottomrule
\end{tabularx}
\end{center}
\textbf{Fields Explained:}
\begin{tightitemize}\footnotesize
  \item \texttt{src0}: Primary source operand (\texttt{r0..r63}, \texttt{u0..u63} uniform, \texttt{k*} FAU constant, or \texttt{t0..t2} temporary; \texttt{\textasciicircum{}} marks discard/last-use). \emph{Example:} {\scriptsize\ttfamily \mnem{FADD.f32} \dst{r0}, \src{r1}, \src{r2}} (\textbf{SPIR-V:} \texttt{OpFAdd (Operand 0)}; \textbf{GLSL:} \texttt{float c = a + b;})
  \item \texttt{src1}: Secondary source operand (register, uniform, or FAU constant; in shifts, supplies the shift count byte). \emph{Example:} {\scriptsize\ttfamily \mnem{LSHIFT.i32} \dst{r0}, \src{r1}, \src{r2}\imod{.b0}} (\textbf{SPIR-V:} \texttt{OpShiftLeftLogical (Shift)}; \textbf{GLSL:} \texttt{uint y = a \textless{}\textless{} s;})
  \item \hyperref[enum:message_slot_index_m]{\texttt{message\_slot\_index}}: Asynchronous scoreboard slot (\texttt{slot0}--\texttt{slot2}, \texttt{slot7}) tracking completion of this message. \emph{Example:} {\scriptsize\ttfamily \mnem{LOAD.i32}\imod{.slot0} \dst{@r0}, \src{r2}, \imm{\#0}} (\textbf{SPIR-V:} \texttt{OpLoad}; \textbf{GLSL:} \texttt{uint v = ssbo.arr[i];})
  \item \texttt{sr\_count}: Number of consecutive 32-bit staging registers transferred by a load/store/texture message. \emph{Example:} {\scriptsize\ttfamily \mnem{LD\_MULTI.i128} \dst{@r0:r3}, \src{r4}, \imm{\#0}} (\textbf{SPIR-V:} \texttt{OpLoad (wide vector)}; \textbf{GLSL:} \texttt{uvec4 x = buf[i];})
  \item \texttt{staging\_register\_index}: Base GPR index of the contiguous staging register window (\texttt{@rN:rM}) for message payload/return. \emph{Example:} {\scriptsize\ttfamily \mnem{STORE.i64}\imod{.slot7} \src{@r0:r1}, \src{r2}, \imm{\#8}} (\textbf{SPIR-V:} \texttt{OpStore (uvec2)}; \textbf{GLSL:} \texttt{ssbo.pairs[i] = uvec2(a, b);})
  \item \texttt{staging\_register\_control\_32}: Controls the read/write staging register bank and count for async message instructions. \emph{Example:} {\scriptsize\ttfamily \mnem{LOAD.i128}\imod{.slot0} \dst{@r0:r3}, \src{r4}, \imm{\#0}} (\textbf{SPIR-V:} \texttt{OpLoad (vec4)}; \textbf{GLSL:} \texttt{vec4 v = ssbo.data[idx];})
  \item \texttt{opcode1}: Primary sub-opcode within the instruction class uniquely identifying the operation family. \emph{Example:} {\scriptsize\ttfamily \mnem{IADD.s32} \dst{r0}, \src{r1}, \src{r2}} (\textbf{SPIR-V:} \texttt{OpIAdd (operation selector)}; \textbf{GLSL:} \texttt{// Primary opcode bits})
  \item \texttt{instruction\_class}: 4-bit major instruction class selector (\texttt{bits 48:51}) partitioning the 64-bit Valhall opcode space (e.g., \texttt{FMA}, \texttt{CVT}, \texttt{SFU}, \texttt{LS}, \texttt{TEX}, \texttt{BRANCH}). \emph{Example:} {\scriptsize\ttfamily \mnem{FADD.f32} \dst{r0}, \src{r1}, \src{r2}} (\textbf{SPIR-V:} \texttt{N/A (encoding class)}; \textbf{GLSL:} \texttt{// Hardware execution pipe / decode class})
  \item \texttt{fau\_page\_index}: Selects the FAU (Fast Access Uniform) page (\texttt{0..3}) when a source operand references hardware small constants, thread IDs, or lane constants. \emph{Example:} {\scriptsize\ttfamily \mnem{FADD.f32} \dst{r0}, \src{r1}, \src{k1f}} (\textbf{SPIR-V:} \texttt{OpConstant / BuiltIn}; \textbf{GLSL:} \texttt{float y = x + 1.0;})
  \item \hyperref[enum:flow_control_m]{\texttt{flow}}: 4-bit embedded flow-control \& scoreboard wait modifier (\texttt{bits 59:62}: \texttt{none}, \texttt{wait0..wait0126}, \texttt{wait}, \texttt{reconverge}, \texttt{discard}, \texttt{end}). \emph{Example:} {\scriptsize\ttfamily \mnem{FADD.f32}\imod{.wait0.reconverge} \dst{r0}, \src{r1}, \src{r2}} (\textbf{SPIR-V:} \texttt{OpControlBarrier / OpReturn}; \textbf{GLSL:} \texttt{// Scoreboard wait + warp reconvergence})
\end{tightitemize}
\clearpage

\subsection*{LEA\_TEX}
\addcontentsline{toc}{subsection}{LEA\_TEX}
Computes effective address + conversion descriptor (3 staging regs) for image texel at packed 16-bit coordinates and dynamic handle.
\par\textbf{Notes:}
\begin{tightitemize}
  \item Returns \texttt{(addr\_lo, addr\_hi, cvt\_desc)} for \texttt{LD\_CVT}/\texttt{ST\_CVT} (\texttt{ISA.xml:113}).
  \item \textbf{Source references:} \texttt{kraid/ops.rs:3213-3239}, \texttt{bifrost/valhall/ISA.xml:113}.
\end{tightitemize}
\codebox{\mnem{LEA\_TEX}\imod{\{.message\_slot\_index\}}\imod{\{.flow\}} \dst{@sw}, \src{src0}, \src{src1}, \src{src2}}
\textbf{Operands:}
\begin{tightitemize}\footnotesize
  \item \texttt{\dst{@sw}} (\texttt{staging\_register\_index} + \texttt{sr\_count}): 3-register destination range: \texttt{@sw[0:1]} = 64-bit texel VA, \texttt{@sw[2]} = 32-bit format conversion descriptor.
  \item \texttt{\src{src0}}: Packed \texttt{u16vec2(x, y)} texel coordinates.
  \item \texttt{\src{src1}}: 3D slice / array layer \texttt{z} coordinate (and optional MSAA sample index).
  \item \texttt{\src{src2}}: Dynamic image/texture descriptor handle.
\end{tightitemize}
\par\textbf{Closest SPIR-V Op:} \texttt{None (image texel address + cvt desc)}
\par\textbf{Pseudo-GLSL (Operational Semantics):}
\codebox{@sw[0:1] = texelAddress(handle = src2, xy = src0, z = src1);\\@sw[2]   = texFormatDesc(handle = src2);}
\textbf{Execution Context:}
\begin{tightitemize}\footnotesize
  \item \textbf{Texture \& Sampler chain (\hyperref[sec:csf-srt]{\texttt{SRT\_2}} Frag / \texttt{SRT\_0} Comp/Vert):} 64B-aligned VA (\texttt{\& \textasciitilde{}0x3F}) $\to$ \texttt{Resource[table]} $\to$ \hyperref[sec:csf-tex-samp]{\texttt{Texture[tex\_idx]}} (32B, \texttt{type=2}) $\to$ \texttt{GenericPlane[mip]} (32B, \texttt{type=10}: \texttt{pointer}, \texttt{row\_stride}, Morton 16$\times$16 / AFBC) + \texttt{Sampler[samp\_idx]} (32B, \texttt{type=1}: filter, wrap, LOD clamp/bias, compare).
  \item \textbf{Inline FAU \& Special Operands (\hyperref[sec:csf-tsd-fau]{\texttt{FAU\_k}}, \texttt{r8..r15}):} Any 8-bit source with \texttt{mode=2} reads zero-latency uniform \texttt{u0..u31} from \texttt{FAU\_k} (\texttt{slots64 = bits[63:56]}); \texttt{mode=3} reads \texttt{SmallConstant} (Table~1.2) or \texttt{FauSpecial} (\texttt{pc}, \texttt{lane\_id}, \texttt{lane\_width}, \texttt{tls\_ptr}, \texttt{wls\_ptr}, \texttt{blend\_descriptor\_0..7}).
\end{tightitemize}
\par\textbf{Worked Examples:}
\par\textbf{Example 1 (Dynamic storage image texel address):}
\codebox{\mnem{LEA\_TEX}\imod{.slot0} \dst{@r0:r2}, \src{r3}, \src{r4}, \src{r5}}
Computes \texttt{(addr\_lo, addr\_hi, cvt\_desc)} in \texttt{@r0:r2} for dynamic image handle \texttt{r5} at packed 16-bit coords \texttt{(r3, r4)}.
\par\textbf{Example 2 (Image atomic address setup):}
\codebox{\mnem{LEA\_TEX}\imod{.slot0} \dst{@r2:r4}, \src{r0}, \src{k0}, \src{r5}}
Resolves the exact 64-bit texel pointer \texttt{[r2:r3]} for \texttt{imageAtomicAdd} via subsequent \texttt{ATOM\_RETURN.i32}.

\begin{center}\footnotesize
\begin{tabularx}{\linewidth}{@{}lllX@{}}
\toprule
Variant & Arch & Class & Opcode Fields \\
\midrule
\texttt{(none)} & v9--v14 & \texttt{S3} & \texttt{opcode1=LEA\_ATTR} \\
\bottomrule
\end{tabularx}
\end{center}
\begin{center}\footnotesize
\begin{tabularx}{0.88\linewidth}{@{}l l X@{}}
\toprule
Field Name & Bits & Type / Fixed Value \\
\midrule
\texttt{src0} & \texttt{0:7} & \texttt{SourceEncoding} \\
\texttt{src1} & \texttt{8:15} & \texttt{SourceEncoding} \\
\texttt{src2} & \texttt{16:23} & \texttt{SourceEncoding} \\
\textit{Reserved} & \texttt{24:29} & \texttt{zero} \\
\hyperref[enum:message_slot_index_m]{\texttt{message\_slot\_index}} & \texttt{30:32} & \texttt{message\_slot\_index\_m} \\
\texttt{sr\_count} & \texttt{33:35} & \texttt{exact=count3} \\
\textit{Reserved} & \texttt{36} & \texttt{zero} \\
\texttt{descriptor\_type} & \texttt{37:38} & \texttt{exact=texture\_u16} \\
\textit{Reserved} & \texttt{39} & \texttt{zero} \\
\texttt{staging\_register\_index} & \texttt{40:45} & \texttt{uint} \\
\texttt{staging\_register\_control\_32} & \texttt{46:47} & \texttt{exact=write} \\
\texttt{opcode1} & \texttt{48:51} & \texttt{exact=LEA\_ATTR} \\
\texttt{instruction\_class} & \texttt{52:56} & \texttt{exact=S3} \\
\texttt{fau\_page\_index} & \texttt{57:58} & \texttt{uint} \\
\hyperref[enum:flow_control_m]{\texttt{flow}} & \texttt{59:62} & \texttt{flow\_control\_m} \\
\textit{Reserved} & \texttt{63} & \texttt{zero} \\
\bottomrule
\end{tabularx}
\end{center}
\textbf{Fields Explained:}
\begin{tightitemize}\footnotesize
  \item \texttt{src0}: Primary source operand (\texttt{r0..r63}, \texttt{u0..u63} uniform, \texttt{k*} FAU constant, or \texttt{t0..t2} temporary; \texttt{\textasciicircum{}} marks discard/last-use). \emph{Example:} {\scriptsize\ttfamily \mnem{FADD.f32} \dst{r0}, \src{r1}, \src{r2}} (\textbf{SPIR-V:} \texttt{OpFAdd (Operand 0)}; \textbf{GLSL:} \texttt{float c = a + b;})
  \item \texttt{src1}: Secondary source operand (register, uniform, or FAU constant; in shifts, supplies the shift count byte). \emph{Example:} {\scriptsize\ttfamily \mnem{LSHIFT.i32} \dst{r0}, \src{r1}, \src{r2}\imod{.b0}} (\textbf{SPIR-V:} \texttt{OpShiftLeftLogical (Shift)}; \textbf{GLSL:} \texttt{uint y = a \textless{}\textless{} s;})
  \item \texttt{src2}: Third source operand (FMA addend, \texttt{MUX}/\texttt{CSEL} selector, \texttt{IADDC} carry-in, or bitwise mask in fused shift-logic). \emph{Example:} {\scriptsize\ttfamily \mnem{FMA.f32} \dst{r0}, \src{r1}, \src{r2}, \src{r3}} (\textbf{SPIR-V:} \texttt{GLSL.std.450 Fma (Addend)}; \textbf{GLSL:} \texttt{float d = fma(a, b, c);})
  \item \hyperref[enum:message_slot_index_m]{\texttt{message\_slot\_index}}: Asynchronous scoreboard slot (\texttt{slot0}--\texttt{slot2}, \texttt{slot7}) tracking completion of this message. \emph{Example:} {\scriptsize\ttfamily \mnem{LOAD.i32}\imod{.slot0} \dst{@r0}, \src{r2}, \imm{\#0}} (\textbf{SPIR-V:} \texttt{OpLoad}; \textbf{GLSL:} \texttt{uint v = ssbo.arr[i];})
  \item \texttt{sr\_count}: Number of consecutive 32-bit staging registers transferred by a load/store/texture message. \emph{Example:} {\scriptsize\ttfamily \mnem{LD\_MULTI.i128} \dst{@r0:r3}, \src{r4}, \imm{\#0}} (\textbf{SPIR-V:} \texttt{OpLoad (wide vector)}; \textbf{GLSL:} \texttt{uvec4 x = buf[i];})
  \item \texttt{descriptor\_type}: Selects descriptor table binding class (attribute, buffer, texture, varying). \emph{Example:} {\scriptsize\ttfamily \mnem{LD\_BUFFER.i32} \dst{@r0}, \src{r2}, \imm{\#0}} (\textbf{SPIR-V:} \texttt{OpLoad (StorageBuffer)}; \textbf{GLSL:} \texttt{uint v = ssbo.arr[i];})
  \item \texttt{staging\_register\_index}: Base GPR index of the contiguous staging register window (\texttt{@rN:rM}) for message payload/return. \emph{Example:} {\scriptsize\ttfamily \mnem{STORE.i64}\imod{.slot7} \src{@r0:r1}, \src{r2}, \imm{\#8}} (\textbf{SPIR-V:} \texttt{OpStore (uvec2)}; \textbf{GLSL:} \texttt{ssbo.pairs[i] = uvec2(a, b);})
  \item \texttt{staging\_register\_control\_32}: Controls the read/write staging register bank and count for async message instructions. \emph{Example:} {\scriptsize\ttfamily \mnem{LOAD.i128}\imod{.slot0} \dst{@r0:r3}, \src{r4}, \imm{\#0}} (\textbf{SPIR-V:} \texttt{OpLoad (vec4)}; \textbf{GLSL:} \texttt{vec4 v = ssbo.data[idx];})
  \item \texttt{opcode1}: Primary sub-opcode within the instruction class uniquely identifying the operation family. \emph{Example:} {\scriptsize\ttfamily \mnem{IADD.s32} \dst{r0}, \src{r1}, \src{r2}} (\textbf{SPIR-V:} \texttt{OpIAdd (operation selector)}; \textbf{GLSL:} \texttt{// Primary opcode bits})
  \item \texttt{instruction\_class}: 4-bit major instruction class selector (\texttt{bits 48:51}) partitioning the 64-bit Valhall opcode space (e.g., \texttt{FMA}, \texttt{CVT}, \texttt{SFU}, \texttt{LS}, \texttt{TEX}, \texttt{BRANCH}). \emph{Example:} {\scriptsize\ttfamily \mnem{FADD.f32} \dst{r0}, \src{r1}, \src{r2}} (\textbf{SPIR-V:} \texttt{N/A (encoding class)}; \textbf{GLSL:} \texttt{// Hardware execution pipe / decode class})
  \item \texttt{fau\_page\_index}: Selects the FAU (Fast Access Uniform) page (\texttt{0..3}) when a source operand references hardware small constants, thread IDs, or lane constants. \emph{Example:} {\scriptsize\ttfamily \mnem{FADD.f32} \dst{r0}, \src{r1}, \src{k1f}} (\textbf{SPIR-V:} \texttt{OpConstant / BuiltIn}; \textbf{GLSL:} \texttt{float y = x + 1.0;})
  \item \hyperref[enum:flow_control_m]{\texttt{flow}}: 4-bit embedded flow-control \& scoreboard wait modifier (\texttt{bits 59:62}: \texttt{none}, \texttt{wait0..wait0126}, \texttt{wait}, \texttt{reconverge}, \texttt{discard}, \texttt{end}). \emph{Example:} {\scriptsize\ttfamily \mnem{FADD.f32}\imod{.wait0.reconverge} \dst{r0}, \src{r1}, \src{r2}} (\textbf{SPIR-V:} \texttt{OpControlBarrier / OpReturn}; \textbf{GLSL:} \texttt{// Scoreboard wait + warp reconvergence})
\end{tightitemize}
\clearpage

\subsection*{LEA\_TEX\_IMM}
\addcontentsline{toc}{subsection}{LEA\_TEX\_IMM}
Resolves the Image/Texture descriptor at immediate \texttt{(table=\#texture\_table\_index, index=\#texture\_index)} for coordinates \texttt{(xy=src0, z=src1)} and writes 3 staging registers to \texttt{@sw}: \texttt{@sw[0:1]} holds the 64-bit texel GPU VA and \texttt{@sw[2]} holds the 32-bit format conversion descriptor for \texttt{LD\_CVT}/\texttt{ST\_CVT}.
\par\textbf{Notes:}
\begin{tightitemize}
  \item Immediate-handle form of \texttt{LEA\_TEX} (\texttt{ISA.xml:115}).
  \item \textbf{Source references:} \texttt{bifrost/valhall/ISA.xml:115}, \texttt{isa-v9-v14.xml}.
\end{tightitemize}
\codebox{\mnem{LEA\_TEX\_IMM}\imod{\{.message\_slot\_index\}}\imod{\{.flow\}} \dst{@sw}, \src{src0}, \src{src1}, \imm{\#texture\_index}, \imm{\#texture\_table\_index}}
\textbf{Operands:}
\begin{tightitemize}\footnotesize
  \item \texttt{\dst{@sw}} (\texttt{staging\_register\_index} + \texttt{sr\_count}): 3-register destination range: \texttt{@sw[0:1]} = 64-bit texel VA, \texttt{@sw[2]} = 32-bit format conversion descriptor.
  \item \texttt{\src{src0}}: Packed \texttt{u16vec2(x, y)} texel coordinates.
  \item \texttt{\src{src1}}: 3D slice / array layer \texttt{z} coordinate.
  \item \texttt{\imm{\#texture\_index}} (\texttt{texture\_index}): Immediate index of the Image/Texture descriptor within the table.
  \item \texttt{\imm{\#texture\_table\_index}} (\texttt{encoded\_texture\_table\_index}): Immediate SRT table index holding the Image/Texture descriptor array.
\end{tightitemize}
\par\textbf{Closest SPIR-V Op:} \texttt{None (image texel address + cvt desc)}
\par\textbf{Pseudo-GLSL (Operational Semantics):}
\codebox{TexDesc t = SRT[\#texture\_table\_index].Tex[\#texture\_index];\\@sw[0:1] = texelAddress(t, xy = src0, z = src1); @sw[2] = t.format\_desc;}
\textbf{Execution Context:}
\begin{tightitemize}\footnotesize
  \item \textbf{Texture \& Sampler chain (\hyperref[sec:csf-srt]{\texttt{SRT\_2}} Frag / \texttt{SRT\_0} Comp/Vert):} 64B-aligned VA (\texttt{\& \textasciitilde{}0x3F}) $\to$ \texttt{Resource[table]} $\to$ \hyperref[sec:csf-tex-samp]{\texttt{Texture[tex\_idx]}} (32B, \texttt{type=2}) $\to$ \texttt{GenericPlane[mip]} (32B, \texttt{type=10}: \texttt{pointer}, \texttt{row\_stride}, Morton 16$\times$16 / AFBC) + \texttt{Sampler[samp\_idx]} (32B, \texttt{type=1}: filter, wrap, LOD clamp/bias, compare).
  \item \textbf{Inline FAU \& Special Operands (\hyperref[sec:csf-tsd-fau]{\texttt{FAU\_k}}, \texttt{r8..r15}):} Any 8-bit source with \texttt{mode=2} reads zero-latency uniform \texttt{u0..u31} from \texttt{FAU\_k} (\texttt{slots64 = bits[63:56]}); \texttt{mode=3} reads \texttt{SmallConstant} (Table~1.2) or \texttt{FauSpecial} (\texttt{pc}, \texttt{lane\_id}, \texttt{lane\_width}, \texttt{tls\_ptr}, \texttt{wls\_ptr}, \texttt{blend\_descriptor\_0..7}).
\end{tightitemize}
\par\textbf{Worked Examples:}
\par\textbf{Example 1 (Static 2D storage image texel address):}
\codebox{\mnem{LEA\_TEX\_IMM}\imod{.slot0} \dst{@r0:r2}, \src{r3}, \src{r4}, \imm{\#0}, \imm{\#0}}
Computes \texttt{(addr\_lo, addr\_hi, cvt\_desc)} for static image \texttt{(table=0, index=0)} at coords \texttt{(r3, r4)} for \texttt{ST\_CVT}/\texttt{LD\_CVT}.
\par\textbf{Example 2 (Static 3D/array image layer address):}
\codebox{\mnem{LEA\_TEX\_IMM}\imod{.slot0} \dst{@r0:r2}, \src{r3}, \src{r4}, \imm{\#0}, \imm{\#1}}
Passes packed \texttt{(x, y)} in \texttt{r3} and \texttt{(z/layer, msaa\_sample)} in \texttt{r4} for binding \texttt{(0, 1)}.
\par\textbf{Mined SPIR-V $\to$ Valhall ISA Example (Computing storage-image texel address + conversion descriptor for OpImageWrite --- Compute shader, Mali-G710 (v10)):}
\codebox{\textbf{; SPIR-V (input)}\\OpImageWrite \%storage\_img \%coord\_xy \%rgba\_val None\\[3pt]\textbf{; Valhall ISA (malioc output)}\\SWZ.v2i16 r0.h1, r1\textasciicircum{}.h10          ; pack (x, y) into u16vec2\\LEA\_TEX\_IMM.slot0.wait02 @r0:r1:r2, r0\textasciicircum{}, k0, \#0x6, \#0x4\\ST\_CVT.slot0.reconverge.estream.v4.f32 @r8:r9:r10:r11, [r0\textasciicircum{}:r1\textasciicircum{}], \#0x0, r2\textasciicircum{}}
\texttt{OpImageWrite} is a two-message sequence on Valhall: \texttt{LEA\_TEX\_IMM} resolves image descriptor \texttt{(table=\#0x4, index=\#0x6)} at \texttt{(x, y)} into 3 staging registers (\texttt{@r0:r1} = 64-bit texel address, \texttt{r2} = 32-bit format conversion descriptor), and \texttt{ST\_CVT} converts and stores \texttt{@r8:r9:r10:r11}.
\begin{center}\footnotesize
\begin{tabularx}{\linewidth}{@{}lllX@{}}
\toprule
Variant & Arch & Class & Opcode Fields \\
\midrule
\texttt{(none)} & v9--v14 & \texttt{S2} & \texttt{opcode1=LEA\_ATTR\_IMM} \\
\bottomrule
\end{tabularx}
\end{center}
\begin{center}\footnotesize
\begin{tabularx}{0.88\linewidth}{@{}l l X@{}}
\toprule
Field Name & Bits & Type / Fixed Value \\
\midrule
\texttt{src0} & \texttt{0:7} & \texttt{SourceEncoding} \\
\texttt{src1} & \texttt{8:15} & \texttt{SourceEncoding} \\
\texttt{encoded\_texture\_table\_index} & \texttt{16:19} & \texttt{uint} \\
\texttt{texture\_index} & \texttt{20:23} & \texttt{uint} \\
\textit{Reserved} & \texttt{24:29} & \texttt{zero} \\
\hyperref[enum:message_slot_index_m]{\texttt{message\_slot\_index}} & \texttt{30:32} & \texttt{message\_slot\_index\_m} \\
\texttt{sr\_count} & \texttt{33:35} & \texttt{exact=count3} \\
\textit{Reserved} & \texttt{36} & \texttt{zero} \\
\texttt{descriptor\_type} & \texttt{37:38} & \texttt{exact=texture\_u16} \\
\textit{Reserved} & \texttt{39} & \texttt{zero} \\
\texttt{staging\_register\_index} & \texttt{40:45} & \texttt{uint} \\
\texttt{staging\_register\_control\_32} & \texttt{46:47} & \texttt{exact=write} \\
\texttt{opcode1} & \texttt{48:51} & \texttt{exact=LEA\_ATTR\_IMM} \\
\texttt{instruction\_class} & \texttt{52:56} & \texttt{exact=S2} \\
\texttt{fau\_page\_index} & \texttt{57:58} & \texttt{uint} \\
\hyperref[enum:flow_control_m]{\texttt{flow}} & \texttt{59:62} & \texttt{flow\_control\_m} \\
\textit{Reserved} & \texttt{63} & \texttt{zero} \\
\bottomrule
\end{tabularx}
\end{center}
\textbf{Fields Explained:}
\begin{tightitemize}\footnotesize
  \item \texttt{src0}: Primary source operand (\texttt{r0..r63}, \texttt{u0..u63} uniform, \texttt{k*} FAU constant, or \texttt{t0..t2} temporary; \texttt{\textasciicircum{}} marks discard/last-use). \emph{Example:} {\scriptsize\ttfamily \mnem{FADD.f32} \dst{r0}, \src{r1}, \src{r2}} (\textbf{SPIR-V:} \texttt{OpFAdd (Operand 0)}; \textbf{GLSL:} \texttt{float c = a + b;})
  \item \texttt{src1}: Secondary source operand (register, uniform, or FAU constant; in shifts, supplies the shift count byte). \emph{Example:} {\scriptsize\ttfamily \mnem{LSHIFT.i32} \dst{r0}, \src{r1}, \src{r2}\imod{.b0}} (\textbf{SPIR-V:} \texttt{OpShiftLeftLogical (Shift)}; \textbf{GLSL:} \texttt{uint y = a \textless{}\textless{} s;})
  \item \texttt{encoded\_texture\_table\_index}: Encoded descriptor table index for \textbackslash\{\}texttt\{encoded\textbackslash\{\}\_texture\textbackslash\{\}\_table\textbackslash\{\}\_index\}. \emph{Example:} {\scriptsize\ttfamily \mnem{LD\_BUFFER.i32} \dst{@r0}, \src{r2}, \imm{\#0}} (\textbf{SPIR-V:} \texttt{OpAccessChain (DescriptorSet)}; \textbf{GLSL:} \texttt{uint v = buf[i];})
  \item \texttt{texture\_index}: Resource/slot index field \textbackslash\{\}texttt\{texture\textbackslash\{\}\_index\} selecting the target descriptor or slot. \emph{Example:} {\scriptsize\ttfamily \mnem{LD\_VAR\_BUF\_IMM.f32} \dst{@r0:r3}, \imm{\#0}} (\textbf{SPIR-V:} \texttt{OpLoad (Descriptor)}; \textbf{GLSL:} \texttt{vec4 v = v\_in;})
  \item \hyperref[enum:message_slot_index_m]{\texttt{message\_slot\_index}}: Asynchronous scoreboard slot (\texttt{slot0}--\texttt{slot2}, \texttt{slot7}) tracking completion of this message. \emph{Example:} {\scriptsize\ttfamily \mnem{LOAD.i32}\imod{.slot0} \dst{@r0}, \src{r2}, \imm{\#0}} (\textbf{SPIR-V:} \texttt{OpLoad}; \textbf{GLSL:} \texttt{uint v = ssbo.arr[i];})
  \item \texttt{sr\_count}: Number of consecutive 32-bit staging registers transferred by a load/store/texture message. \emph{Example:} {\scriptsize\ttfamily \mnem{LD\_MULTI.i128} \dst{@r0:r3}, \src{r4}, \imm{\#0}} (\textbf{SPIR-V:} \texttt{OpLoad (wide vector)}; \textbf{GLSL:} \texttt{uvec4 x = buf[i];})
  \item \texttt{descriptor\_type}: Selects descriptor table binding class (attribute, buffer, texture, varying). \emph{Example:} {\scriptsize\ttfamily \mnem{LD\_BUFFER.i32} \dst{@r0}, \src{r2}, \imm{\#0}} (\textbf{SPIR-V:} \texttt{OpLoad (StorageBuffer)}; \textbf{GLSL:} \texttt{uint v = ssbo.arr[i];})
  \item \texttt{staging\_register\_index}: Base GPR index of the contiguous staging register window (\texttt{@rN:rM}) for message payload/return. \emph{Example:} {\scriptsize\ttfamily \mnem{STORE.i64}\imod{.slot7} \src{@r0:r1}, \src{r2}, \imm{\#8}} (\textbf{SPIR-V:} \texttt{OpStore (uvec2)}; \textbf{GLSL:} \texttt{ssbo.pairs[i] = uvec2(a, b);})
  \item \texttt{staging\_register\_control\_32}: Controls the read/write staging register bank and count for async message instructions. \emph{Example:} {\scriptsize\ttfamily \mnem{LOAD.i128}\imod{.slot0} \dst{@r0:r3}, \src{r4}, \imm{\#0}} (\textbf{SPIR-V:} \texttt{OpLoad (vec4)}; \textbf{GLSL:} \texttt{vec4 v = ssbo.data[idx];})
  \item \texttt{opcode1}: Primary sub-opcode within the instruction class uniquely identifying the operation family. \emph{Example:} {\scriptsize\ttfamily \mnem{IADD.s32} \dst{r0}, \src{r1}, \src{r2}} (\textbf{SPIR-V:} \texttt{OpIAdd (operation selector)}; \textbf{GLSL:} \texttt{// Primary opcode bits})
  \item \texttt{instruction\_class}: 4-bit major instruction class selector (\texttt{bits 48:51}) partitioning the 64-bit Valhall opcode space (e.g., \texttt{FMA}, \texttt{CVT}, \texttt{SFU}, \texttt{LS}, \texttt{TEX}, \texttt{BRANCH}). \emph{Example:} {\scriptsize\ttfamily \mnem{FADD.f32} \dst{r0}, \src{r1}, \src{r2}} (\textbf{SPIR-V:} \texttt{N/A (encoding class)}; \textbf{GLSL:} \texttt{// Hardware execution pipe / decode class})
  \item \texttt{fau\_page\_index}: Selects the FAU (Fast Access Uniform) page (\texttt{0..3}) when a source operand references hardware small constants, thread IDs, or lane constants. \emph{Example:} {\scriptsize\ttfamily \mnem{FADD.f32} \dst{r0}, \src{r1}, \src{k1f}} (\textbf{SPIR-V:} \texttt{OpConstant / BuiltIn}; \textbf{GLSL:} \texttt{float y = x + 1.0;})
  \item \hyperref[enum:flow_control_m]{\texttt{flow}}: 4-bit embedded flow-control \& scoreboard wait modifier (\texttt{bits 59:62}: \texttt{none}, \texttt{wait0..wait0126}, \texttt{wait}, \texttt{reconverge}, \texttt{discard}, \texttt{end}). \emph{Example:} {\scriptsize\ttfamily \mnem{FADD.f32}\imod{.wait0.reconverge} \dst{r0}, \src{r1}, \src{r2}} (\textbf{SPIR-V:} \texttt{OpControlBarrier / OpReturn}; \textbf{GLSL:} \texttt{// Scoreboard wait + warp reconvergence})
\end{tightitemize}
\clearpage

\subsection*{LOAD}
\addcontentsline{toc}{subsection}{LOAD}
Loads \texttt{1..16} bytes (\texttt{i8}/\texttt{i16}/\texttt{i24}/\texttt{i32}/\texttt{i48}/\texttt{i64}/\texttt{i96}/\texttt{i128}) from \texttt{src0 + sext64(offset)} into staging register(s) \texttt{sr\_dst}; applies \texttt{extend} (\texttt{sext} vs \texttt{zext\_or\_none}) and sub-register \texttt{lane} (\texttt{b0..b3}/\texttt{h0..h1}) for narrow loads.
\par\textbf{Notes:}
\begin{tightitemize}
  \item Little-endian memory.
  \item Supports up to 128-bit (4 staging regs, even-aligned when \texttt{count \textgreater{} 1}).
  \item \textbf{Source references:} \texttt{kraid/ops.rs:3241-3269}, \texttt{encode\_v9.rs:2642-2663}, \texttt{bifrost/valhall/ISA.xml:117-124}.
\end{tightitemize}
\codebox{\mnem{LOAD}\imod{\{.message\_slot\_index\}}\imod{\{.flow\}}\imod{\{.access\}} \dst{@sw}, \src{src0}, \imm{\#offset}\\\mnem{LOAD}\imod{\{.message\_slot\_index\}}\imod{\{.flow\}}\imod{\{.access\}}\imod{\{.lane\}}\imod{\{.extend\}} \dst{@sw}, \src{src0}, \imm{\#offset}}
\textbf{Operands:}
\begin{tightitemize}\footnotesize
  \item \texttt{\dst{@sw}} (\texttt{staging\_register\_index} + \texttt{sr\_count}): Contiguous destination staging register range receiving the loaded data.
  \item \texttt{\src{src0}}: 64-bit base GPU virtual address (\texttt{[rN:rN+1]} or 64-bit FAU uniform).
  \item \texttt{\imm{\#offset}} (\texttt{offset}): Signed byte offset added to \texttt{src0}.
\end{tightitemize}
\par\textbf{Closest SPIR-V Op:} \texttt{OpLoad (PhysicalStorageBuffer / SSBO / Shared)}
\par\textbf{Pseudo-GLSL (Operational Semantics):}
\codebox{uint64\_t addr = src0 + int64\_t(\#offset);\\if (width \textless{} 32) @sw = (extend == sext) ? sext32(mem[addr]) : zext32(mem[addr]);\\else            @sw = mem\_read(addr, width, access\_hint);}
\textbf{Execution Context:}
\begin{tightitemize}\footnotesize
  \item \textbf{Direct 64-bit GPU VA memory/atomics:} Base addresses typically come from \hyperref[sec:csf-tsd-fau]{\texttt{FAU\_k}} (\texttt{u0..u31} populated by a Pilot shader chasing \texttt{SRT\_0} $\to$ \texttt{Resource[]} $\to$ \texttt{Buffer[]}), \texttt{FauSpecial::tls\_ptr}/\texttt{wls\_ptr} backed by \hyperref[sec:csf-tsd-fau]{\texttt{TSD\_k}} (\texttt{r24..r31}), or \texttt{FauSpecial::pc} for embedded constant LUTs.
  \item \textbf{Inline FAU \& Special Operands (\hyperref[sec:csf-tsd-fau]{\texttt{FAU\_k}}, \texttt{r8..r15}):} Any 8-bit source with \texttt{mode=2} reads zero-latency uniform \texttt{u0..u31} from \texttt{FAU\_k} (\texttt{slots64 = bits[63:56]}); \texttt{mode=3} reads \texttt{SmallConstant} (Table~1.2) or \texttt{FauSpecial} (\texttt{pc}, \texttt{lane\_id}, \texttt{lane\_width}, \texttt{tls\_ptr}, \texttt{wls\_ptr}, \texttt{blend\_descriptor\_0..7}).
\end{tightitemize}
\par\textbf{Worked Examples:}
\par\textbf{Example 1 (Sign-extended 8-bit load (`.sext`)):}
\codebox{\mnem{LOAD.i8}\imod{.slot0.sext} \dst{@r1}, \src{[r2:r3]}, \imm{\#0x4}}
Reads 1 byte at \texttt{[r2:r3] + 4}: if memory holds \texttt{0xFE}, \texttt{.sext} sign-extends to \texttt{r1 = 0xFFFFFFFE} (\texttt{-2}).
\par\textbf{Example 2 (128-bit streaming vector load):}
\codebox{\mnem{LOAD.i128}\imod{.estream.slot0} \dst{@r0:r3}, \src{[r4:r5]}, \imm{\#0x10}}
Loads 16 bytes little-endian from \texttt{[r4:r5] + 16} into \texttt{@r0:r3} with outer-streaming cache hint (\texttt{.estream}).
\par\textbf{Mined SPIR-V $\to$ Valhall ISA Example (Direct 64-bit virtual-address load with immediate struct offset --- Vertex shader, Mali-G710 (v10)):}
\codebox{\textbf{; SPIR-V (input)}\\\%ptr = OpAccessChain \%\_ptr\_Uniform\_uint \%ssbo \%int\_0 \%idx \%int\_113 ; offset 0x1c4\\\%val = OpLoad \%uint \%ptr\\[3pt]\textbf{; Valhall ISA (malioc output)}\\IMUL.u32 r1, r0, k13121110.b0\\IADD.s64 [r2:r3], u8, [r0:r1].w1\\LOAD.i32.slot0.wait0 @r1, [r2\textasciicircum{}:r3\textasciicircum{}], \#0x1c4}
Once the 64-bit base pointer \texttt{[r2:r3]} is computed by adding the strided element offset to base VA \texttt{u8}, \texttt{LOAD.i32} folds the constant struct member offset \texttt{\#0x1c4} directly into the load instruction.
\begin{center}\footnotesize
\begin{tabularx}{\linewidth}{@{}lllX@{}}
\toprule
Variant & Arch & Class & Opcode Fields \\
\midrule
\texttt{(none)} & v9--v14 & \texttt{S2} & \texttt{opcode1=LOAD} \\
\texttt{(none)} & v9--v14 & \texttt{S2} & \texttt{opcode1=LOAD} \\
\texttt{(none)} & v9--v14 & \texttt{S2} & \texttt{opcode1=LOAD} \\
\texttt{(none)} & v9--v14 & \texttt{S2} & \texttt{opcode1=LOAD} \\
\texttt{(none)} & v9--v14 & \texttt{S2} & \texttt{opcode1=LOAD} \\
\texttt{(none)} & v9--v14 & \texttt{S2} & \texttt{opcode1=LOAD} \\
\texttt{(none)} & v9--v14 & \texttt{S2} & \texttt{opcode1=LOAD} \\
\texttt{(none)} & v9--v14 & \texttt{S2} & \texttt{opcode1=LOAD} \\
\bottomrule
\end{tabularx}
\end{center}
\begin{center}\footnotesize
\begin{tabularx}{0.88\linewidth}{@{}l l X@{}}
\toprule
Field Name & Bits & Type / Fixed Value \\
\midrule
\texttt{src0} & \texttt{0:7} & \texttt{SourceEncoding64} \\
\texttt{offset} & \texttt{8:23} & \texttt{int} \\
\hyperref[enum:access_load_m]{\texttt{access}} & \texttt{24:26} & \texttt{access\_load\_m (def none)} \\
\texttt{memory\_width} & \texttt{27:29} & \texttt{exact=i128} \\
\hyperref[enum:message_slot_index_m]{\texttt{message\_slot\_index}} & \texttt{30:32} & \texttt{message\_slot\_index\_m} \\
\texttt{sr\_count} & \texttt{33:35} & \texttt{exact=count4} \\
\hyperref[enum:lane_all_m]{\texttt{lane}} & \texttt{36:38} & \texttt{exact=d0} \\
\hyperref[enum:sign_extend_or_none_compressed_encoding_m]{\texttt{extend}} & \texttt{39} & \texttt{exact=zext\_or\_none} \\
\texttt{staging\_register\_index} & \texttt{40:45} & \texttt{uint} \\
\texttt{staging\_register\_control\_32} & \texttt{46:47} & \texttt{exact=write} \\
\texttt{opcode1} & \texttt{48:51} & \texttt{exact=LOAD} \\
\texttt{instruction\_class} & \texttt{52:56} & \texttt{exact=S2} \\
\texttt{fau\_page\_index} & \texttt{57:58} & \texttt{uint} \\
\hyperref[enum:flow_control_m]{\texttt{flow}} & \texttt{59:62} & \texttt{flow\_control\_m} \\
\textit{Reserved} & \texttt{63} & \texttt{zero} \\
\bottomrule
\end{tabularx}
\end{center}
\textbf{Fields Explained:}
\begin{tightitemize}\footnotesize
  \item \texttt{src0}: Primary source operand (\texttt{r0..r63}, \texttt{u0..u63} uniform, \texttt{k*} FAU constant, or \texttt{t0..t2} temporary; \texttt{\textasciicircum{}} marks discard/last-use). \emph{Example:} {\scriptsize\ttfamily \mnem{FADD.f32} \dst{r0}, \src{r1}, \src{r2}} (\textbf{SPIR-V:} \texttt{OpFAdd (Operand 0)}; \textbf{GLSL:} \texttt{float c = a + b;})
  \item \texttt{offset}: Signed byte offset added to base address register in memory load/store/atomic instructions. \emph{Example:} {\scriptsize\ttfamily \mnem{LOAD.i32} \dst{@r0}, \src{[r2:r3]}, \imm{\#16}} (\textbf{SPIR-V:} \texttt{OpAccessChain + OpLoad}; \textbf{GLSL:} \texttt{uint y = ssbo.arr[4];})
  \item \hyperref[enum:access_load_m]{\texttt{access}}: Memory cache streaming/coherence policy (\texttt{none}, \texttt{istream}, \texttt{estream}, \texttt{force}). \emph{Example:} {\scriptsize\ttfamily \mnem{STORE.i32}\imod{.force} \src{@r0}, \src{r2}, \imm{\#0}} (\textbf{SPIR-V:} \texttt{OpStore (MakePointerVisible / Coherent)}; \textbf{GLSL:} \texttt{coherent buffer B \{ uint x; \}; x = 1u;})
  \item \texttt{memory\_width}: Memory transfer width (\texttt{i8}, \texttt{i16}, \texttt{i24}, \texttt{i32}, \texttt{i48}, \texttt{i64}, \texttt{i96}, \texttt{i128}). \emph{Example:} {\scriptsize\ttfamily \mnem{LOAD.i64} \dst{@r0:r1}, \src{[r2:r3]}, \imm{\#0}} (\textbf{SPIR-V:} \texttt{OpLoad (uvec2)}; \textbf{GLSL:} \texttt{uvec2 v = ssbo.pairs[i];})
  \item \hyperref[enum:message_slot_index_m]{\texttt{message\_slot\_index}}: Asynchronous scoreboard slot (\texttt{slot0}--\texttt{slot2}, \texttt{slot7}) tracking completion of this message. \emph{Example:} {\scriptsize\ttfamily \mnem{LOAD.i32}\imod{.slot0} \dst{@r0}, \src{r2}, \imm{\#0}} (\textbf{SPIR-V:} \texttt{OpLoad}; \textbf{GLSL:} \texttt{uint v = ssbo.arr[i];})
  \item \texttt{sr\_count}: Number of consecutive 32-bit staging registers transferred by a load/store/texture message. \emph{Example:} {\scriptsize\ttfamily \mnem{LD\_MULTI.i128} \dst{@r0:r3}, \src{r4}, \imm{\#0}} (\textbf{SPIR-V:} \texttt{OpLoad (wide vector)}; \textbf{GLSL:} \texttt{uvec4 x = buf[i];})
  \item \hyperref[enum:lane_all_m]{\texttt{lane}}: Sub-word lane selector (\texttt{b0}..\texttt{b3}, \texttt{h0}, \texttt{h1}, \texttt{w0}, \texttt{d0}) for narrow checked/unpacked memory transfers. \emph{Example:} {\scriptsize\ttfamily \mnem{LOAD.i16}\imod{.h1} \dst{@r0}, \src{r2}, \imm{\#0}} (\textbf{SPIR-V:} \texttt{OpLoad (uint16\_t)}; \textbf{GLSL:} \texttt{uint16\_t x = ssbo.u16[i];})
  \item \hyperref[enum:sign_extend_or_none_compressed_encoding_m]{\texttt{extend}}: Virtual extension mode (\texttt{sext}, \texttt{zext}, \texttt{none}) decoded from \texttt{encoded\_extend} and \texttt{memory\_width}. \emph{Example:} {\scriptsize\ttfamily \mnem{LOAD.i8}\imod{.sext} \dst{@r0}, \src{r2}, \imm{\#0}} (\textbf{SPIR-V:} \texttt{OpLoad + OpSConvert}; \textbf{GLSL:} \texttt{int x = int(ssbo.i8\_arr[i]);})
  \item \texttt{staging\_register\_index}: Base GPR index of the contiguous staging register window (\texttt{@rN:rM}) for message payload/return. \emph{Example:} {\scriptsize\ttfamily \mnem{STORE.i64}\imod{.slot7} \src{@r0:r1}, \src{r2}, \imm{\#8}} (\textbf{SPIR-V:} \texttt{OpStore (uvec2)}; \textbf{GLSL:} \texttt{ssbo.pairs[i] = uvec2(a, b);})
  \item \texttt{staging\_register\_control\_32}: Controls the read/write staging register bank and count for async message instructions. \emph{Example:} {\scriptsize\ttfamily \mnem{LOAD.i128}\imod{.slot0} \dst{@r0:r3}, \src{r4}, \imm{\#0}} (\textbf{SPIR-V:} \texttt{OpLoad (vec4)}; \textbf{GLSL:} \texttt{vec4 v = ssbo.data[idx];})
  \item \texttt{opcode1}: Primary sub-opcode within the instruction class uniquely identifying the operation family. \emph{Example:} {\scriptsize\ttfamily \mnem{IADD.s32} \dst{r0}, \src{r1}, \src{r2}} (\textbf{SPIR-V:} \texttt{OpIAdd (operation selector)}; \textbf{GLSL:} \texttt{// Primary opcode bits})
  \item \texttt{instruction\_class}: 4-bit major instruction class selector (\texttt{bits 48:51}) partitioning the 64-bit Valhall opcode space (e.g., \texttt{FMA}, \texttt{CVT}, \texttt{SFU}, \texttt{LS}, \texttt{TEX}, \texttt{BRANCH}). \emph{Example:} {\scriptsize\ttfamily \mnem{FADD.f32} \dst{r0}, \src{r1}, \src{r2}} (\textbf{SPIR-V:} \texttt{N/A (encoding class)}; \textbf{GLSL:} \texttt{// Hardware execution pipe / decode class})
  \item \texttt{fau\_page\_index}: Selects the FAU (Fast Access Uniform) page (\texttt{0..3}) when a source operand references hardware small constants, thread IDs, or lane constants. \emph{Example:} {\scriptsize\ttfamily \mnem{FADD.f32} \dst{r0}, \src{r1}, \src{k1f}} (\textbf{SPIR-V:} \texttt{OpConstant / BuiltIn}; \textbf{GLSL:} \texttt{float y = x + 1.0;})
  \item \hyperref[enum:flow_control_m]{\texttt{flow}}: 4-bit embedded flow-control \& scoreboard wait modifier (\texttt{bits 59:62}: \texttt{none}, \texttt{wait0..wait0126}, \texttt{wait}, \texttt{reconverge}, \texttt{discard}, \texttt{end}). \emph{Example:} {\scriptsize\ttfamily \mnem{FADD.f32}\imod{.wait0.reconverge} \dst{r0}, \src{r1}, \src{r2}} (\textbf{SPIR-V:} \texttt{OpControlBarrier / OpReturn}; \textbf{GLSL:} \texttt{// Scoreboard wait + warp reconvergence})
\end{tightitemize}
\clearpage

\subsection*{STORE}
\addcontentsline{toc}{subsection}{STORE}
Stores \texttt{1..16} bytes (\texttt{i8}..\texttt{i128}) from staging register(s) \texttt{sr\_src} (with optional \texttt{lane} for \texttt{i8}/\texttt{i16}) to \texttt{src0 + sext64(offset)}.
\par\textbf{Notes:}
\begin{tightitemize}
  \item Little-endian memory.
  \item Supports up to 128-bit.
  \item \textbf{Source references:} \texttt{kraid/ops.rs:4072-4105}, \texttt{encode\_v9.rs:3119-3140}, \texttt{isa-v9-v14.xml}.
\end{tightitemize}
\codebox{\mnem{STORE}\imod{\{.message\_slot\_index\}}\imod{\{.flow\}}\imod{\{.access\}} \src{@sr}, \src{src0}, \imm{\#offset}\\\mnem{STORE}\imod{\{.message\_slot\_index\}}\imod{\{.flow\}}\imod{\{.access\}}\imod{\{.lane\}} \src{@sr}, \src{src0}, \imm{\#offset}}
\textbf{Operands:}
\begin{tightitemize}\footnotesize
  \item \texttt{\src{@sr}} (\texttt{staging\_register\_index} + \texttt{sr\_count}): Contiguous source staging register range holding the data to write.
  \item \texttt{\src{src0}}: 64-bit base GPU virtual address (\texttt{[rN:rN+1]} or 64-bit FAU uniform).
  \item \texttt{\imm{\#offset}} (\texttt{offset}): Signed byte offset added to \texttt{src0}.
\end{tightitemize}
\par\textbf{Closest SPIR-V Op:} \texttt{OpStore (PhysicalStorageBuffer / SSBO / Shared)}
\par\textbf{Pseudo-GLSL (Operational Semantics):}
\codebox{uint64\_t addr = src0 + int64\_t(\#offset);\\mem\_write(addr, @sr, width, access\_hint); // .istream/.estream for IDVS outputs}
\textbf{Execution Context:}
\begin{tightitemize}\footnotesize
  \item \textbf{Direct 64-bit GPU VA memory/atomics:} Base addresses typically come from \hyperref[sec:csf-tsd-fau]{\texttt{FAU\_k}} (\texttt{u0..u31} populated by a Pilot shader chasing \texttt{SRT\_0} $\to$ \texttt{Resource[]} $\to$ \texttt{Buffer[]}), \texttt{FauSpecial::tls\_ptr}/\texttt{wls\_ptr} backed by \hyperref[sec:csf-tsd-fau]{\texttt{TSD\_k}} (\texttt{r24..r31}), or \texttt{FauSpecial::pc} for embedded constant LUTs.
  \item \textbf{Graphics fixed-function streaming (\texttt{STORE.istream}/\texttt{.estream}):} In IDVS Position shaders, \texttt{STORE.istream} writes 16B \texttt{gl\_Position} into \hyperref[sec:csf-tiler]{\texttt{TILER\_CTX::geometry\_buffer}} (\texttt{r40:r41}); in Varying shaders, \texttt{STORE.estream} writes per-vertex varyings into \hyperref[sec:csf-tiler]{\texttt{IDVSSR::VARY\_SIZE}} (\texttt{r48:r49}).
  \item \textbf{Inline FAU \& Special Operands (\hyperref[sec:csf-tsd-fau]{\texttt{FAU\_k}}, \texttt{r8..r15}):} Any 8-bit source with \texttt{mode=2} reads zero-latency uniform \texttt{u0..u31} from \texttt{FAU\_k} (\texttt{slots64 = bits[63:56]}); \texttt{mode=3} reads \texttt{SmallConstant} (Table~1.2) or \texttt{FauSpecial} (\texttt{pc}, \texttt{lane\_id}, \texttt{lane\_width}, \texttt{tls\_ptr}, \texttt{wls\_ptr}, \texttt{blend\_descriptor\_0..7}).
\end{tightitemize}
\par\textbf{Worked Examples:}
\par\textbf{Example 1 (64-bit little-endian store):}
\codebox{\mnem{STORE.i64}\imod{.slot0} \dst{@r0:r1}, \src{[r2:r3]}, \imm{\#0x8}}
Stores 8 bytes from staging pair \texttt{@r0:r1} little-endian to address \texttt{[r2:r3] + 8} on scoreboard slot 0.
\par\textbf{Example 2 (Coherent 32-bit fire-and-forget store (`.force.slot7`)):}
\codebox{\mnem{STORE.i32}\imod{.force.slot7} \dst{@r0}, \src{[r2:r3]}, \imm{\#0}}
Writes \texttt{r0} with cache-coherent visibility (\texttt{.force}) tracked on non-blocking \texttt{slot7}.
\par\textbf{Mined SPIR-V $\to$ Valhall ISA Example (Coalesced 128-bit compute SSBO store with .end --- Compute shader, Mali-G710 (v10)):}
\codebox{\textbf{; SPIR-V (input)}\\OpStore \%out\_buf\_0 \%uint\_1000\\OpStore \%out\_buf\_1 \%uint\_4294966296\\OpStore \%out\_buf\_2 \%uint\_0\\OpStore \%out\_buf\_3 \%uint\_0\\[3pt]\textbf{; Valhall ISA (malioc output)}\\STORE.i128.slot0.end @r0:r1:r2:r3, u0, \#0x0}
\texttt{STORE.i128.slot0.end} writes 16 bytes from \texttt{@r0:r1:r2:r3} to the 64-bit base VA in \texttt{u0} and simultaneously terminates the shader warp via \texttt{.end}.
\begin{center}\footnotesize
\begin{tabularx}{\linewidth}{@{}lllX@{}}
\toprule
Variant & Arch & Class & Opcode Fields \\
\midrule
\texttt{(none)} & v9--v14 & \texttt{S2} & \texttt{opcode1=STORE} \\
\texttt{(none)} & v9--v14 & \texttt{S2} & \texttt{opcode1=STORE} \\
\texttt{(none)} & v9--v14 & \texttt{S2} & \texttt{opcode1=STORE} \\
\texttt{(none)} & v9--v14 & \texttt{S2} & \texttt{opcode1=STORE} \\
\texttt{(none)} & v9--v14 & \texttt{S2} & \texttt{opcode1=STORE} \\
\texttt{(none)} & v9--v14 & \texttt{S2} & \texttt{opcode1=STORE} \\
\texttt{(none)} & v9--v14 & \texttt{S2} & \texttt{opcode1=STORE} \\
\texttt{(none)} & v9--v14 & \texttt{S2} & \texttt{opcode1=STORE} \\
\bottomrule
\end{tabularx}
\end{center}
\begin{center}\footnotesize
\begin{tabularx}{0.88\linewidth}{@{}l l X@{}}
\toprule
Field Name & Bits & Type / Fixed Value \\
\midrule
\texttt{src0} & \texttt{0:7} & \texttt{SourceEncoding64} \\
\texttt{offset} & \texttt{8:23} & \texttt{int} \\
\hyperref[enum:access_store_m]{\texttt{access}} & \texttt{24:26} & \texttt{access\_store\_m (def none)} \\
\texttt{memory\_width} & \texttt{27:29} & \texttt{exact=i128} \\
\hyperref[enum:message_slot_index_m]{\texttt{message\_slot\_index}} & \texttt{30:32} & \texttt{message\_slot\_index\_m} \\
\texttt{sr\_count} & \texttt{33:35} & \texttt{exact=count4} \\
\textit{Reserved} & \texttt{36:39} & \texttt{zero} \\
\texttt{staging\_register\_index} & \texttt{40:45} & \texttt{uint} \\
\texttt{staging\_register\_control\_32} & \texttt{46:47} & \texttt{exact=read} \\
\texttt{opcode1} & \texttt{48:51} & \texttt{exact=STORE} \\
\texttt{instruction\_class} & \texttt{52:56} & \texttt{exact=S2} \\
\texttt{fau\_page\_index} & \texttt{57:58} & \texttt{uint} \\
\hyperref[enum:flow_control_m]{\texttt{flow}} & \texttt{59:62} & \texttt{flow\_control\_m} \\
\textit{Reserved} & \texttt{63} & \texttt{zero} \\
\bottomrule
\end{tabularx}
\end{center}
\textbf{Fields Explained:}
\begin{tightitemize}\footnotesize
  \item \texttt{src0}: Primary source operand (\texttt{r0..r63}, \texttt{u0..u63} uniform, \texttt{k*} FAU constant, or \texttt{t0..t2} temporary; \texttt{\textasciicircum{}} marks discard/last-use). \emph{Example:} {\scriptsize\ttfamily \mnem{FADD.f32} \dst{r0}, \src{r1}, \src{r2}} (\textbf{SPIR-V:} \texttt{OpFAdd (Operand 0)}; \textbf{GLSL:} \texttt{float c = a + b;})
  \item \texttt{offset}: Signed byte offset added to base address register in memory load/store/atomic instructions. \emph{Example:} {\scriptsize\ttfamily \mnem{LOAD.i32} \dst{@r0}, \src{[r2:r3]}, \imm{\#16}} (\textbf{SPIR-V:} \texttt{OpAccessChain + OpLoad}; \textbf{GLSL:} \texttt{uint y = ssbo.arr[4];})
  \item \hyperref[enum:access_store_m]{\texttt{access}}: Memory cache streaming/coherence policy (\texttt{none}, \texttt{istream}, \texttt{estream}, \texttt{force}). \emph{Example:} {\scriptsize\ttfamily \mnem{STORE.i32}\imod{.force} \src{@r0}, \src{r2}, \imm{\#0}} (\textbf{SPIR-V:} \texttt{OpStore (MakePointerVisible / Coherent)}; \textbf{GLSL:} \texttt{coherent buffer B \{ uint x; \}; x = 1u;})
  \item \texttt{memory\_width}: Memory transfer width (\texttt{i8}, \texttt{i16}, \texttt{i24}, \texttt{i32}, \texttt{i48}, \texttt{i64}, \texttt{i96}, \texttt{i128}). \emph{Example:} {\scriptsize\ttfamily \mnem{LOAD.i64} \dst{@r0:r1}, \src{[r2:r3]}, \imm{\#0}} (\textbf{SPIR-V:} \texttt{OpLoad (uvec2)}; \textbf{GLSL:} \texttt{uvec2 v = ssbo.pairs[i];})
  \item \hyperref[enum:message_slot_index_m]{\texttt{message\_slot\_index}}: Asynchronous scoreboard slot (\texttt{slot0}--\texttt{slot2}, \texttt{slot7}) tracking completion of this message. \emph{Example:} {\scriptsize\ttfamily \mnem{LOAD.i32}\imod{.slot0} \dst{@r0}, \src{r2}, \imm{\#0}} (\textbf{SPIR-V:} \texttt{OpLoad}; \textbf{GLSL:} \texttt{uint v = ssbo.arr[i];})
  \item \texttt{sr\_count}: Number of consecutive 32-bit staging registers transferred by a load/store/texture message. \emph{Example:} {\scriptsize\ttfamily \mnem{LD\_MULTI.i128} \dst{@r0:r3}, \src{r4}, \imm{\#0}} (\textbf{SPIR-V:} \texttt{OpLoad (wide vector)}; \textbf{GLSL:} \texttt{uvec4 x = buf[i];})
  \item \texttt{staging\_register\_index}: Base GPR index of the contiguous staging register window (\texttt{@rN:rM}) for message payload/return. \emph{Example:} {\scriptsize\ttfamily \mnem{STORE.i64}\imod{.slot7} \src{@r0:r1}, \src{r2}, \imm{\#8}} (\textbf{SPIR-V:} \texttt{OpStore (uvec2)}; \textbf{GLSL:} \texttt{ssbo.pairs[i] = uvec2(a, b);})
  \item \texttt{staging\_register\_control\_32}: Controls the read/write staging register bank and count for async message instructions. \emph{Example:} {\scriptsize\ttfamily \mnem{LOAD.i128}\imod{.slot0} \dst{@r0:r3}, \src{r4}, \imm{\#0}} (\textbf{SPIR-V:} \texttt{OpLoad (vec4)}; \textbf{GLSL:} \texttt{vec4 v = ssbo.data[idx];})
  \item \texttt{opcode1}: Primary sub-opcode within the instruction class uniquely identifying the operation family. \emph{Example:} {\scriptsize\ttfamily \mnem{IADD.s32} \dst{r0}, \src{r1}, \src{r2}} (\textbf{SPIR-V:} \texttt{OpIAdd (operation selector)}; \textbf{GLSL:} \texttt{// Primary opcode bits})
  \item \texttt{instruction\_class}: 4-bit major instruction class selector (\texttt{bits 48:51}) partitioning the 64-bit Valhall opcode space (e.g., \texttt{FMA}, \texttt{CVT}, \texttt{SFU}, \texttt{LS}, \texttt{TEX}, \texttt{BRANCH}). \emph{Example:} {\scriptsize\ttfamily \mnem{FADD.f32} \dst{r0}, \src{r1}, \src{r2}} (\textbf{SPIR-V:} \texttt{N/A (encoding class)}; \textbf{GLSL:} \texttt{// Hardware execution pipe / decode class})
  \item \texttt{fau\_page\_index}: Selects the FAU (Fast Access Uniform) page (\texttt{0..3}) when a source operand references hardware small constants, thread IDs, or lane constants. \emph{Example:} {\scriptsize\ttfamily \mnem{FADD.f32} \dst{r0}, \src{r1}, \src{k1f}} (\textbf{SPIR-V:} \texttt{OpConstant / BuiltIn}; \textbf{GLSL:} \texttt{float y = x + 1.0;})
  \item \hyperref[enum:flow_control_m]{\texttt{flow}}: 4-bit embedded flow-control \& scoreboard wait modifier (\texttt{bits 59:62}: \texttt{none}, \texttt{wait0..wait0126}, \texttt{wait}, \texttt{reconverge}, \texttt{discard}, \texttt{end}). \emph{Example:} {\scriptsize\ttfamily \mnem{FADD.f32}\imod{.wait0.reconverge} \dst{r0}, \src{r1}, \src{r2}} (\textbf{SPIR-V:} \texttt{OpControlBarrier / OpReturn}; \textbf{GLSL:} \texttt{// Scoreboard wait + warp reconvergence})
\end{tightitemize}
\clearpage

\subsection*{ST\_CHECKED}
\addcontentsline{toc}{subsection}{ST\_CHECKED}
Bounds-checked store (v12+): stores \texttt{i8}..\texttt{i128} to base \texttt{src0} + offset \texttt{src1} checked against bound \texttt{src2}; OOB stores are dropped.
\par\textbf{Notes:}
\begin{tightitemize}
  \item Robust buffer access store.
  \item \textbf{Source references:} \texttt{isa-v9-v14.xml}.
\end{tightitemize}
\codebox{\mnem{ST\_CHECKED}\imod{\{.message\_slot\_index\}}\imod{\{.flow\}}\imod{\{.access\}} \src{@sr}, \src{src0}, \src{src1}, \src{src2}\\\mnem{ST\_CHECKED}\imod{\{.message\_slot\_index\}}\imod{\{.flow\}}\imod{\{.access\}}\imod{\{.lane\}} \src{@sr}, \src{src0}, \src{src1}, \src{src2}}
\textbf{Operands:}
\begin{tightitemize}\footnotesize
  \item \texttt{\src{@sr}} (\texttt{staging\_register\_index} + \texttt{sr\_count}): Source staging register range holding the data to store.
  \item \texttt{\src{src0}}: 64-bit base GPU virtual address.
  \item \texttt{\src{src1}}: Dynamic unsigned 32-bit byte offset.
  \item \texttt{\src{src2}}: 32-bit buffer byte-size upper bound (OOB stores are silently dropped).
\end{tightitemize}
\par\textbf{Closest SPIR-V Op:} \texttt{OpStore (robustBufferAccess, v12+)}
\par\textbf{Pseudo-GLSL (Operational Semantics):}
\codebox{if (uint64\_t(src1) + access\_bytes \textless{}= uint64\_t(src2)) mem\_write(src0 + src1, @sr, width);}
\textbf{Execution Context:}
\begin{tightitemize}\footnotesize
  \item \textbf{Hardware bounds check (v12+):} Chases \hyperref[sec:csf-srt]{\texttt{SRT\_k}} / \hyperref[sec:csf-tsd-fau]{\texttt{FAU\_k}} base VA and enforces \texttt{offset + width <= Buffer::size} (or bound register); OOB loads return \texttt{0} (unless v14 \texttt{.oob\_ignore}) and OOB stores drop silently.
  \item \textbf{Inline FAU \& Special Operands (\hyperref[sec:csf-tsd-fau]{\texttt{FAU\_k}}, \texttt{r8..r15}):} Any 8-bit source with \texttt{mode=2} reads zero-latency uniform \texttt{u0..u31} from \texttt{FAU\_k} (\texttt{slots64 = bits[63:56]}); \texttt{mode=3} reads \texttt{SmallConstant} (Table~1.2) or \texttt{FauSpecial} (\texttt{pc}, \texttt{lane\_id}, \texttt{lane\_width}, \texttt{tls\_ptr}, \texttt{wls\_ptr}, \texttt{blend\_descriptor\_0..7}).
\end{tightitemize}
\par\textbf{Worked Examples:}
\par\textbf{Example 1 (In-bounds checked store):}
\codebox{\mnem{ST\_CHECKED.i32}\imod{.slot0} \dst{@r0}, \src{u3}, \src{r1}, \src{r2}}
Checks \texttt{r1 \textless{} r2}: when in-bounds, writes \texttt{@r0} to \texttt{u3 + r1}.
\par\textbf{Example 2 (Out-of-bounds silent drop):}
\codebox{\mnem{ST\_CHECKED.i32}\imod{.slot0} \dst{@r0}, \src{u3}, \src{r1}, \src{r2}}
When \texttt{r1 \textgreater{}= r2}: silently drops the store without faulting or corrupting adjacent memory.
\par\textbf{Mined SPIR-V $\to$ Valhall ISA Example (Robust bounds-checked SSBO store on Immortalis-G925 (v14) --- Compute shader, Immortalis-G925 (v14)):}
\codebox{\textbf{; SPIR-V (input)}\\\%ptr = OpAccessChain \%\_ptr\_Uniform\_float \%ssbo \%int\_0 \%idx\\OpStore \%ptr \%val\\[3pt]\textbf{; Valhall ISA (malioc output)}\\LSHIFT.i32.wait012 r2, r0, k07060504.b0 ; byte offset = idx \textless{}\textless{} 2\\ST\_CHECKED.i32.slot0 @r1, u2, r2, k1n}
\texttt{ST\_CHECKED.i32} writes \texttt{r1} to base \texttt{u2} at byte offset \texttt{r2}, silently dropping the store if out-of-bounds under Vulkan robust buffer access.
\begin{center}\footnotesize
\begin{tabularx}{\linewidth}{@{}lllX@{}}
\toprule
Variant & Arch & Class & Opcode Fields \\
\midrule
\texttt{(none)} & v9--v14 & \texttt{S4} & \texttt{opcode1=ST\_CHECKED} \\
\texttt{(none)} & v9--v14 & \texttt{S4} & \texttt{opcode1=ST\_CHECKED} \\
\texttt{(none)} & v9--v14 & \texttt{S4} & \texttt{opcode1=ST\_CHECKED} \\
\texttt{(none)} & v9--v14 & \texttt{S4} & \texttt{opcode1=ST\_CHECKED} \\
\texttt{(none)} & v9--v14 & \texttt{S4} & \texttt{opcode1=ST\_CHECKED} \\
\texttt{(none)} & v9--v14 & \texttt{S4} & \texttt{opcode1=ST\_CHECKED} \\
\texttt{(none)} & v9--v14 & \texttt{S4} & \texttt{opcode1=ST\_CHECKED} \\
\texttt{(none)} & v9--v14 & \texttt{S4} & \texttt{opcode1=ST\_CHECKED} \\
\texttt{(none)} & v9--v14 & \texttt{S4} & \texttt{opcode1=ST\_CHECKED} \\
\texttt{(none)} & v9--v14 & \texttt{S4} & \texttt{opcode1=ST\_CHECKED} \\
\texttt{(none)} & v9--v14 & \texttt{S4} & \texttt{opcode1=ST\_CHECKED} \\
\texttt{(none)} & v9--v14 & \texttt{S4} & \texttt{opcode1=ST\_CHECKED} \\
\texttt{(none)} & v9--v14 & \texttt{S4} & \texttt{opcode1=ST\_CHECKED} \\
\texttt{(none)} & v9--v14 & \texttt{S4} & \texttt{opcode1=ST\_CHECKED} \\
\texttt{(none)} & v9--v14 & \texttt{S4} & \texttt{opcode1=ST\_CHECKED} \\
\texttt{(none)} & v9--v14 & \texttt{S4} & \texttt{opcode1=ST\_CHECKED} \\
\bottomrule
\end{tabularx}
\end{center}
\begin{center}\footnotesize
\begin{tabularx}{0.88\linewidth}{@{}l l X@{}}
\toprule
Field Name & Bits & Type / Fixed Value \\
\midrule
\texttt{src0} & \texttt{0:7} & \texttt{SourceEncoding64} \\
\texttt{src1} & \texttt{8:15} & \texttt{SourceEncoding} \\
\texttt{src2} & \texttt{16:23} & \texttt{SourceEncoding} \\
\hyperref[enum:access_store_m]{\texttt{access}} & \texttt{24:26} & \texttt{access\_store\_m (def none)} \\
\texttt{memory\_width} & \texttt{27:29} & \texttt{exact=i128} \\
\hyperref[enum:message_slot_index_m]{\texttt{message\_slot\_index}} & \texttt{30:32} & \texttt{message\_slot\_index\_m} \\
\texttt{sr\_count} & \texttt{33:35} & \texttt{exact=count4} \\
\textit{Reserved} & \texttt{36:39} & \texttt{zero} \\
\texttt{staging\_register\_index} & \texttt{40:45} & \texttt{uint} \\
\texttt{staging\_register\_control\_32} & \texttt{46:47} & \texttt{exact=read} \\
\texttt{opcode1} & \texttt{48:51} & \texttt{exact=ST\_CHECKED} \\
\texttt{instruction\_class} & \texttt{52:56} & \texttt{exact=S4} \\
\texttt{fau\_page\_index} & \texttt{57:58} & \texttt{uint} \\
\hyperref[enum:flow_control_m]{\texttt{flow}} & \texttt{59:62} & \texttt{flow\_control\_m} \\
\textit{Reserved} & \texttt{63} & \texttt{zero} \\
\bottomrule
\end{tabularx}
\end{center}
\textbf{Fields Explained:}
\begin{tightitemize}\footnotesize
  \item \texttt{src0}: Primary source operand (\texttt{r0..r63}, \texttt{u0..u63} uniform, \texttt{k*} FAU constant, or \texttt{t0..t2} temporary; \texttt{\textasciicircum{}} marks discard/last-use). \emph{Example:} {\scriptsize\ttfamily \mnem{FADD.f32} \dst{r0}, \src{r1}, \src{r2}} (\textbf{SPIR-V:} \texttt{OpFAdd (Operand 0)}; \textbf{GLSL:} \texttt{float c = a + b;})
  \item \texttt{src1}: Secondary source operand (register, uniform, or FAU constant; in shifts, supplies the shift count byte). \emph{Example:} {\scriptsize\ttfamily \mnem{LSHIFT.i32} \dst{r0}, \src{r1}, \src{r2}\imod{.b0}} (\textbf{SPIR-V:} \texttt{OpShiftLeftLogical (Shift)}; \textbf{GLSL:} \texttt{uint y = a \textless{}\textless{} s;})
  \item \texttt{src2}: Third source operand (FMA addend, \texttt{MUX}/\texttt{CSEL} selector, \texttt{IADDC} carry-in, or bitwise mask in fused shift-logic). \emph{Example:} {\scriptsize\ttfamily \mnem{FMA.f32} \dst{r0}, \src{r1}, \src{r2}, \src{r3}} (\textbf{SPIR-V:} \texttt{GLSL.std.450 Fma (Addend)}; \textbf{GLSL:} \texttt{float d = fma(a, b, c);})
  \item \hyperref[enum:access_store_m]{\texttt{access}}: Memory cache streaming/coherence policy (\texttt{none}, \texttt{istream}, \texttt{estream}, \texttt{force}). \emph{Example:} {\scriptsize\ttfamily \mnem{STORE.i32}\imod{.force} \src{@r0}, \src{r2}, \imm{\#0}} (\textbf{SPIR-V:} \texttt{OpStore (MakePointerVisible / Coherent)}; \textbf{GLSL:} \texttt{coherent buffer B \{ uint x; \}; x = 1u;})
  \item \texttt{memory\_width}: Memory transfer width (\texttt{i8}, \texttt{i16}, \texttt{i24}, \texttt{i32}, \texttt{i48}, \texttt{i64}, \texttt{i96}, \texttt{i128}). \emph{Example:} {\scriptsize\ttfamily \mnem{LOAD.i64} \dst{@r0:r1}, \src{[r2:r3]}, \imm{\#0}} (\textbf{SPIR-V:} \texttt{OpLoad (uvec2)}; \textbf{GLSL:} \texttt{uvec2 v = ssbo.pairs[i];})
  \item \hyperref[enum:message_slot_index_m]{\texttt{message\_slot\_index}}: Asynchronous scoreboard slot (\texttt{slot0}--\texttt{slot2}, \texttt{slot7}) tracking completion of this message. \emph{Example:} {\scriptsize\ttfamily \mnem{LOAD.i32}\imod{.slot0} \dst{@r0}, \src{r2}, \imm{\#0}} (\textbf{SPIR-V:} \texttt{OpLoad}; \textbf{GLSL:} \texttt{uint v = ssbo.arr[i];})
  \item \texttt{sr\_count}: Number of consecutive 32-bit staging registers transferred by a load/store/texture message. \emph{Example:} {\scriptsize\ttfamily \mnem{LD\_MULTI.i128} \dst{@r0:r3}, \src{r4}, \imm{\#0}} (\textbf{SPIR-V:} \texttt{OpLoad (wide vector)}; \textbf{GLSL:} \texttt{uvec4 x = buf[i];})
  \item \texttt{staging\_register\_index}: Base GPR index of the contiguous staging register window (\texttt{@rN:rM}) for message payload/return. \emph{Example:} {\scriptsize\ttfamily \mnem{STORE.i64}\imod{.slot7} \src{@r0:r1}, \src{r2}, \imm{\#8}} (\textbf{SPIR-V:} \texttt{OpStore (uvec2)}; \textbf{GLSL:} \texttt{ssbo.pairs[i] = uvec2(a, b);})
  \item \texttt{staging\_register\_control\_32}: Controls the read/write staging register bank and count for async message instructions. \emph{Example:} {\scriptsize\ttfamily \mnem{LOAD.i128}\imod{.slot0} \dst{@r0:r3}, \src{r4}, \imm{\#0}} (\textbf{SPIR-V:} \texttt{OpLoad (vec4)}; \textbf{GLSL:} \texttt{vec4 v = ssbo.data[idx];})
  \item \texttt{opcode1}: Primary sub-opcode within the instruction class uniquely identifying the operation family. \emph{Example:} {\scriptsize\ttfamily \mnem{IADD.s32} \dst{r0}, \src{r1}, \src{r2}} (\textbf{SPIR-V:} \texttt{OpIAdd (operation selector)}; \textbf{GLSL:} \texttt{// Primary opcode bits})
  \item \texttt{instruction\_class}: 4-bit major instruction class selector (\texttt{bits 48:51}) partitioning the 64-bit Valhall opcode space (e.g., \texttt{FMA}, \texttt{CVT}, \texttt{SFU}, \texttt{LS}, \texttt{TEX}, \texttt{BRANCH}). \emph{Example:} {\scriptsize\ttfamily \mnem{FADD.f32} \dst{r0}, \src{r1}, \src{r2}} (\textbf{SPIR-V:} \texttt{N/A (encoding class)}; \textbf{GLSL:} \texttt{// Hardware execution pipe / decode class})
  \item \texttt{fau\_page\_index}: Selects the FAU (Fast Access Uniform) page (\texttt{0..3}) when a source operand references hardware small constants, thread IDs, or lane constants. \emph{Example:} {\scriptsize\ttfamily \mnem{FADD.f32} \dst{r0}, \src{r1}, \src{k1f}} (\textbf{SPIR-V:} \texttt{OpConstant / BuiltIn}; \textbf{GLSL:} \texttt{float y = x + 1.0;})
  \item \hyperref[enum:flow_control_m]{\texttt{flow}}: 4-bit embedded flow-control \& scoreboard wait modifier (\texttt{bits 59:62}: \texttt{none}, \texttt{wait0..wait0126}, \texttt{wait}, \texttt{reconverge}, \texttt{discard}, \texttt{end}). \emph{Example:} {\scriptsize\ttfamily \mnem{FADD.f32}\imod{.wait0.reconverge} \dst{r0}, \src{r1}, \src{r2}} (\textbf{SPIR-V:} \texttt{OpControlBarrier / OpReturn}; \textbf{GLSL:} \texttt{// Scoreboard wait + warp reconvergence})
\end{tightitemize}
\clearpage

\subsection*{ST\_CHECKED\_IMM}
\addcontentsline{toc}{subsection}{ST\_CHECKED\_IMM}
Bounds-checked memory store (v12+) at 64-bit base address \texttt{src0} plus immediate unsigned byte offset \texttt{\#offset} (\texttt{offset\_lo} + \texttt{offset\_hi}): if \texttt{\#offset + access\_bytes \textless{}= src2}, writes \texttt{@sr} to memory; out-of-bounds stores are silently dropped.
\par\textbf{Notes:}
\begin{tightitemize}
  \item Robust buffer access store with 12-bit unsigned split immediate offset \texttt{((offset\_hi \& 0xF) \textless{}\textless{} 8) \textbar{} (offset\_lo \& 0xFF)} and full-width bounds check \texttt{offset + access\_bytes \textless{}= bound}.
  \item \textbf{Source references:} \texttt{isa-v9-v14.xml}.
\end{tightitemize}
\codebox{\mnem{ST\_CHECKED\_IMM}\imod{\{.message\_slot\_index\}}\imod{\{.flow\}}\imod{\{.access\}} \src{@sr}, \src{src0}, \imm{\#offset}, \src{src2}\\\mnem{ST\_CHECKED\_IMM}\imod{\{.message\_slot\_index\}}\imod{\{.flow\}}\imod{\{.access\}}\imod{\{.lane\}} \src{@sr}, \src{src0}, \imm{\#offset}, \src{src2}}
\textbf{Operands:}
\begin{tightitemize}\footnotesize
  \item \texttt{\src{@sr}} (\texttt{staging\_register\_index} + \texttt{sr\_count}): Source staging register range holding the data to store.
  \item \texttt{\src{src0}}: 64-bit base GPU virtual address.
  \item \texttt{\imm{\#offset}} (\texttt{offset\_lo} + \texttt{offset\_hi}): Immediate unsigned byte offset (assembled from \texttt{offset\_lo} and \texttt{offset\_hi}).
  \item \texttt{\src{src2}}: 32-bit buffer byte-size upper bound (OOB stores are silently dropped).
\end{tightitemize}
\par\textbf{Closest SPIR-V Op:} \texttt{OpStore (robustBufferAccess imm, v12+)}
\par\textbf{Pseudo-GLSL (Operational Semantics):}
\codebox{if (uint64\_t(\#offset) + access\_bytes \textless{}= uint64\_t(src2)) mem\_write(src0 + \#offset, @sr, width);}
\textbf{Execution Context:}
\begin{tightitemize}\footnotesize
  \item \textbf{Hardware bounds check (v12+):} Chases \hyperref[sec:csf-srt]{\texttt{SRT\_k}} / \hyperref[sec:csf-tsd-fau]{\texttt{FAU\_k}} base VA and enforces \texttt{offset + width <= Buffer::size} (or bound register); OOB loads return \texttt{0} (unless v14 \texttt{.oob\_ignore}) and OOB stores drop silently.
  \item \textbf{Inline FAU \& Special Operands (\hyperref[sec:csf-tsd-fau]{\texttt{FAU\_k}}, \texttt{r8..r15}):} Any 8-bit source with \texttt{mode=2} reads zero-latency uniform \texttt{u0..u31} from \texttt{FAU\_k} (\texttt{slots64 = bits[63:56]}); \texttt{mode=3} reads \texttt{SmallConstant} (Table~1.2) or \texttt{FauSpecial} (\texttt{pc}, \texttt{lane\_id}, \texttt{lane\_width}, \texttt{tls\_ptr}, \texttt{wls\_ptr}, \texttt{blend\_descriptor\_0..7}).
\end{tightitemize}
\par\textbf{Worked Examples:}
\par\textbf{Example 1 (Immediate-offset checked store):}
\codebox{\mnem{ST\_CHECKED\_IMM.i32}\imod{.slot0} \dst{@r0}, \src{u3}, \src{r2}, \imm{\#0x10}}
Checks \texttt{0x10 \textless{} r2}: if in-bounds, writes 4 bytes from \texttt{@r0} to \texttt{u3 + 0x10}.
\par\textbf{Example 2 (Out-of-bounds silent drop):}
\codebox{\mnem{ST\_CHECKED\_IMM.i32}\imod{.slot0} \dst{@r0}, \src{u3}, \src{r2}, \imm{\#0x10}}
If \texttt{r2 \textless{}= 0x10}: suppresses the write cleanly with no memory side effect.

\begin{center}\footnotesize
\begin{tabularx}{\linewidth}{@{}lllX@{}}
\toprule
Variant & Arch & Class & Opcode Fields \\
\midrule
\texttt{(none)} & v9--v14 & \texttt{S3} & \texttt{opcode1=ST\_CHECKED\_IMM} \\
\texttt{(none)} & v9--v14 & \texttt{S3} & \texttt{opcode1=ST\_CHECKED\_IMM} \\
\texttt{(none)} & v9--v14 & \texttt{S3} & \texttt{opcode1=ST\_CHECKED\_IMM} \\
\texttt{(none)} & v9--v14 & \texttt{S3} & \texttt{opcode1=ST\_CHECKED\_IMM} \\
\texttt{(none)} & v9--v14 & \texttt{S3} & \texttt{opcode1=ST\_CHECKED\_IMM} \\
\texttt{(none)} & v9--v14 & \texttt{S3} & \texttt{opcode1=ST\_CHECKED\_IMM} \\
\texttt{(none)} & v9--v14 & \texttt{S3} & \texttt{opcode1=ST\_CHECKED\_IMM} \\
\texttt{(none)} & v9--v14 & \texttt{S3} & \texttt{opcode1=ST\_CHECKED\_IMM} \\
\texttt{(none)} & v9--v14 & \texttt{S3} & \texttt{opcode1=ST\_CHECKED\_IMM} \\
\texttt{(none)} & v9--v14 & \texttt{S3} & \texttt{opcode1=ST\_CHECKED\_IMM} \\
\texttt{(none)} & v9--v14 & \texttt{S3} & \texttt{opcode1=ST\_CHECKED\_IMM} \\
\texttt{(none)} & v9--v14 & \texttt{S3} & \texttt{opcode1=ST\_CHECKED\_IMM} \\
\texttt{(none)} & v9--v14 & \texttt{S3} & \texttt{opcode1=ST\_CHECKED\_IMM} \\
\texttt{(none)} & v9--v14 & \texttt{S3} & \texttt{opcode1=ST\_CHECKED\_IMM} \\
\texttt{(none)} & v9--v14 & \texttt{S3} & \texttt{opcode1=ST\_CHECKED\_IMM} \\
\texttt{(none)} & v9--v14 & \texttt{S3} & \texttt{opcode1=ST\_CHECKED\_IMM} \\
\bottomrule
\end{tabularx}
\end{center}
\begin{center}\footnotesize
\begin{tabularx}{0.88\linewidth}{@{}l l X@{}}
\toprule
Field Name & Bits & Type / Fixed Value \\
\midrule
\texttt{src0} & \texttt{0:7} & \texttt{SourceEncoding64} \\
\texttt{offset\_lo} & \texttt{8:15} & \texttt{uint} \\
\texttt{src2} & \texttt{16:23} & \texttt{SourceEncoding} \\
\hyperref[enum:access_store_m]{\texttt{access}} & \texttt{24:26} & \texttt{access\_store\_m (def none)} \\
\texttt{memory\_width} & \texttt{27:29} & \texttt{exact=i128} \\
\hyperref[enum:message_slot_index_m]{\texttt{message\_slot\_index}} & \texttt{30:32} & \texttt{message\_slot\_index\_m} \\
\texttt{sr\_count} & \texttt{33:35} & \texttt{exact=count4} \\
\texttt{offset\_hi} & \texttt{36:39} & \texttt{uint} \\
\texttt{staging\_register\_index} & \texttt{40:45} & \texttt{uint} \\
\texttt{staging\_register\_control\_32} & \texttt{46:47} & \texttt{exact=read} \\
\texttt{opcode1} & \texttt{48:51} & \texttt{exact=ST\_CHECKED\_IMM} \\
\texttt{instruction\_class} & \texttt{52:56} & \texttt{exact=S3} \\
\texttt{fau\_page\_index} & \texttt{57:58} & \texttt{uint} \\
\hyperref[enum:flow_control_m]{\texttt{flow}} & \texttt{59:62} & \texttt{flow\_control\_m} \\
\textit{Reserved} & \texttt{63} & \texttt{zero} \\
\bottomrule
\end{tabularx}
\end{center}
\textbf{Fields Explained:}
\begin{tightitemize}\footnotesize
  \item \texttt{src0}: Primary source operand (\texttt{r0..r63}, \texttt{u0..u63} uniform, \texttt{k*} FAU constant, or \texttt{t0..t2} temporary; \texttt{\textasciicircum{}} marks discard/last-use). \emph{Example:} {\scriptsize\ttfamily \mnem{FADD.f32} \dst{r0}, \src{r1}, \src{r2}} (\textbf{SPIR-V:} \texttt{OpFAdd (Operand 0)}; \textbf{GLSL:} \texttt{float c = a + b;})
  \item \texttt{offset\_lo}: Signed byte offset added to base address register in memory load/store/atomic instructions. \emph{Example:} {\scriptsize\ttfamily \mnem{LOAD.i32} \dst{@r0}, \src{[r2:r3]}, \imm{\#16}} (\textbf{SPIR-V:} \texttt{OpAccessChain + OpLoad}; \textbf{GLSL:} \texttt{uint y = ssbo.arr[4];})
  \item \texttt{src2}: Third source operand (FMA addend, \texttt{MUX}/\texttt{CSEL} selector, \texttt{IADDC} carry-in, or bitwise mask in fused shift-logic). \emph{Example:} {\scriptsize\ttfamily \mnem{FMA.f32} \dst{r0}, \src{r1}, \src{r2}, \src{r3}} (\textbf{SPIR-V:} \texttt{GLSL.std.450 Fma (Addend)}; \textbf{GLSL:} \texttt{float d = fma(a, b, c);})
  \item \hyperref[enum:access_store_m]{\texttt{access}}: Memory cache streaming/coherence policy (\texttt{none}, \texttt{istream}, \texttt{estream}, \texttt{force}). \emph{Example:} {\scriptsize\ttfamily \mnem{STORE.i32}\imod{.force} \src{@r0}, \src{r2}, \imm{\#0}} (\textbf{SPIR-V:} \texttt{OpStore (MakePointerVisible / Coherent)}; \textbf{GLSL:} \texttt{coherent buffer B \{ uint x; \}; x = 1u;})
  \item \texttt{memory\_width}: Memory transfer width (\texttt{i8}, \texttt{i16}, \texttt{i24}, \texttt{i32}, \texttt{i48}, \texttt{i64}, \texttt{i96}, \texttt{i128}). \emph{Example:} {\scriptsize\ttfamily \mnem{LOAD.i64} \dst{@r0:r1}, \src{[r2:r3]}, \imm{\#0}} (\textbf{SPIR-V:} \texttt{OpLoad (uvec2)}; \textbf{GLSL:} \texttt{uvec2 v = ssbo.pairs[i];})
  \item \hyperref[enum:message_slot_index_m]{\texttt{message\_slot\_index}}: Asynchronous scoreboard slot (\texttt{slot0}--\texttt{slot2}, \texttt{slot7}) tracking completion of this message. \emph{Example:} {\scriptsize\ttfamily \mnem{LOAD.i32}\imod{.slot0} \dst{@r0}, \src{r2}, \imm{\#0}} (\textbf{SPIR-V:} \texttt{OpLoad}; \textbf{GLSL:} \texttt{uint v = ssbo.arr[i];})
  \item \texttt{sr\_count}: Number of consecutive 32-bit staging registers transferred by a load/store/texture message. \emph{Example:} {\scriptsize\ttfamily \mnem{LD\_MULTI.i128} \dst{@r0:r3}, \src{r4}, \imm{\#0}} (\textbf{SPIR-V:} \texttt{OpLoad (wide vector)}; \textbf{GLSL:} \texttt{uvec4 x = buf[i];})
  \item \texttt{offset\_hi}: Signed byte offset added to base address register in memory load/store/atomic instructions. \emph{Example:} {\scriptsize\ttfamily \mnem{LOAD.i32} \dst{@r0}, \src{[r2:r3]}, \imm{\#16}} (\textbf{SPIR-V:} \texttt{OpAccessChain + OpLoad}; \textbf{GLSL:} \texttt{uint y = ssbo.arr[4];})
  \item \texttt{staging\_register\_index}: Base GPR index of the contiguous staging register window (\texttt{@rN:rM}) for message payload/return. \emph{Example:} {\scriptsize\ttfamily \mnem{STORE.i64}\imod{.slot7} \src{@r0:r1}, \src{r2}, \imm{\#8}} (\textbf{SPIR-V:} \texttt{OpStore (uvec2)}; \textbf{GLSL:} \texttt{ssbo.pairs[i] = uvec2(a, b);})
  \item \texttt{staging\_register\_control\_32}: Controls the read/write staging register bank and count for async message instructions. \emph{Example:} {\scriptsize\ttfamily \mnem{LOAD.i128}\imod{.slot0} \dst{@r0:r3}, \src{r4}, \imm{\#0}} (\textbf{SPIR-V:} \texttt{OpLoad (vec4)}; \textbf{GLSL:} \texttt{vec4 v = ssbo.data[idx];})
  \item \texttt{opcode1}: Primary sub-opcode within the instruction class uniquely identifying the operation family. \emph{Example:} {\scriptsize\ttfamily \mnem{IADD.s32} \dst{r0}, \src{r1}, \src{r2}} (\textbf{SPIR-V:} \texttt{OpIAdd (operation selector)}; \textbf{GLSL:} \texttt{// Primary opcode bits})
  \item \texttt{instruction\_class}: 4-bit major instruction class selector (\texttt{bits 48:51}) partitioning the 64-bit Valhall opcode space (e.g., \texttt{FMA}, \texttt{CVT}, \texttt{SFU}, \texttt{LS}, \texttt{TEX}, \texttt{BRANCH}). \emph{Example:} {\scriptsize\ttfamily \mnem{FADD.f32} \dst{r0}, \src{r1}, \src{r2}} (\textbf{SPIR-V:} \texttt{N/A (encoding class)}; \textbf{GLSL:} \texttt{// Hardware execution pipe / decode class})
  \item \texttt{fau\_page\_index}: Selects the FAU (Fast Access Uniform) page (\texttt{0..3}) when a source operand references hardware small constants, thread IDs, or lane constants. \emph{Example:} {\scriptsize\ttfamily \mnem{FADD.f32} \dst{r0}, \src{r1}, \src{k1f}} (\textbf{SPIR-V:} \texttt{OpConstant / BuiltIn}; \textbf{GLSL:} \texttt{float y = x + 1.0;})
  \item \hyperref[enum:flow_control_m]{\texttt{flow}}: 4-bit embedded flow-control \& scoreboard wait modifier (\texttt{bits 59:62}: \texttt{none}, \texttt{wait0..wait0126}, \texttt{wait}, \texttt{reconverge}, \texttt{discard}, \texttt{end}). \emph{Example:} {\scriptsize\ttfamily \mnem{FADD.f32}\imod{.wait0.reconverge} \dst{r0}, \src{r1}, \src{r2}} (\textbf{SPIR-V:} \texttt{OpControlBarrier / OpReturn}; \textbf{GLSL:} \texttt{// Scoreboard wait + warp reconvergence})
\end{tightitemize}
\clearpage

\subsection*{ST\_CVT}
\addcontentsline{toc}{subsection}{ST\_CVT}
Store with format conversion: converts \texttt{vecsize} components of \texttt{register\_format} from \texttt{sr\_src} using conversion descriptor \texttt{src2} and writes to \texttt{src0 + offset}.
\par\textbf{Notes:}
\begin{tightitemize}
  \item Used with \texttt{LEA\_TEX\_IMM} for image stores (\texttt{ISA.xml:142}).
  \item NOT a raw byte store!
  \item \textbf{Source references:} \texttt{kraid/ops.rs:4038-4067}, \texttt{encode\_v9.rs:3095-3117}, \texttt{bifrost/valhall/ISA.xml:142}.
\end{tightitemize}
\codebox{\mnem{ST\_CVT}\imod{\{.message\_slot\_index\}}\imod{\{.flow\}}\imod{\{.access\}}\imod{\{.vecsize\}}\imod{\{.register\_format\}} \src{@sr}, \src{src0}, \imm{\#offset}, \src{src2}}
\textbf{Operands:}
\begin{tightitemize}\footnotesize
  \item \texttt{\src{@sr}} (\texttt{staging\_register\_index} + \texttt{sr\_count}): Source staging register range holding register-format components to encode and store.
  \item \texttt{\src{src0}}: 64-bit base GPU virtual address (typically from \texttt{LEA\_TEX}/\texttt{LEA\_ATTR}).
  \item \texttt{\imm{\#offset}} (\texttt{offset}): Unsigned byte offset added to \texttt{src0}.
  \item \texttt{\src{src2}}: 32-bit hardware format conversion descriptor (typically \texttt{r2} from \texttt{LEA\_TEX}/\texttt{LEA\_ATTR}).
\end{tightitemize}
\par\textbf{Closest SPIR-V Op:} \texttt{OpImageWrite (paired with LEA\_TEX)}
\par\textbf{Pseudo-GLSL (Operational Semantics):}
\codebox{mem\_write(src0 + \#offset, encodePixelFormat(@sr, cvt\_desc = src2, register\_format, vecsize));}
\textbf{Execution Context:}
\begin{tightitemize}\footnotesize
  \item \textbf{Texture \& Sampler chain (\hyperref[sec:csf-srt]{\texttt{SRT\_2}} Frag / \texttt{SRT\_0} Comp/Vert):} 64B-aligned VA (\texttt{\& \textasciitilde{}0x3F}) $\to$ \texttt{Resource[table]} $\to$ \hyperref[sec:csf-tex-samp]{\texttt{Texture[tex\_idx]}} (32B, \texttt{type=2}) $\to$ \texttt{GenericPlane[mip]} (32B, \texttt{type=10}: \texttt{pointer}, \texttt{row\_stride}, Morton 16$\times$16 / AFBC) + \texttt{Sampler[samp\_idx]} (32B, \texttt{type=1}: filter, wrap, LOD clamp/bias, compare).
  \item \textbf{Inline FAU \& Special Operands (\hyperref[sec:csf-tsd-fau]{\texttt{FAU\_k}}, \texttt{r8..r15}):} Any 8-bit source with \texttt{mode=2} reads zero-latency uniform \texttt{u0..u31} from \texttt{FAU\_k} (\texttt{slots64 = bits[63:56]}); \texttt{mode=3} reads \texttt{SmallConstant} (Table~1.2) or \texttt{FauSpecial} (\texttt{pc}, \texttt{lane\_id}, \texttt{lane\_width}, \texttt{tls\_ptr}, \texttt{wls\_ptr}, \texttt{blend\_descriptor\_0..7}).
\end{tightitemize}
\par\textbf{Worked Examples:}
\par\textbf{Example 1 (Converted FP32 -\textgreater{} RGBA8 texel store):}
\codebox{\mnem{ST\_CVT}\imod{.f32.v4.slot0} \dst{@r0:r3}, \src{[r4:r5]}, \src{r6}, \imm{\#0}}
Converts 4x FP32 \texttt{@r0:r3} using conversion descriptor \texttt{r6} (e.g., UNORM8 quantization) and stores to \texttt{[r4:r5]}.
\par\textbf{Example 2 (Integer storage image store):}
\codebox{\mnem{ST\_CVT}\imod{.u32.v4.slot0} \dst{@r0:r3}, \src{[r4:r5]}, \src{r6}, \imm{\#0}}
Writes 4x \texttt{u32} channels into an integer \texttt{uimage2D} texel address returned by \texttt{LEA\_TEX}.
\par\textbf{Mined SPIR-V $\to$ Valhall ISA Example (Format-converting storage-image store paired with LEA\_TEX\_IMM --- Compute shader, Mali-G710 (v10)):}
\codebox{\textbf{; SPIR-V (input)}\\OpImageWrite \%storage\_img \%coord\_xy \%rgba\_val None\\[3pt]\textbf{; Valhall ISA (malioc output)}\\LEA\_TEX\_IMM.slot0.wait02 @r0:r1:r2, r0\textasciicircum{}, k0, \#0x6, \#0x4\\ST\_CVT.slot0.reconverge.estream.v4.f32 @r8:r9:r10:r11, [r0\textasciicircum{}:r1\textasciicircum{}], \#0x0, r2\textasciicircum{}}
\texttt{ST\_CVT} encodes the four FP32 channels in \texttt{@r8:r9:r10:r11} into the storage image's memory format using conversion descriptor \texttt{r2\textasciicircum{}} produced by \texttt{LEA\_TEX\_IMM}.
\begin{center}\footnotesize
\begin{tabularx}{\linewidth}{@{}lllX@{}}
\toprule
Variant & Arch & Class & Opcode Fields \\
\midrule
\texttt{(none)} & v9--v14 & \texttt{S3} & \texttt{opcode1=ST\_CVT} \\
\bottomrule
\end{tabularx}
\end{center}
\begin{center}\footnotesize
\begin{tabularx}{0.88\linewidth}{@{}l l X@{}}
\toprule
Field Name & Bits & Type / Fixed Value \\
\midrule
\texttt{src0} & \texttt{0:7} & \texttt{SourceEncoding64} \\
\texttt{offset} & \texttt{8:15} & \texttt{uint} \\
\texttt{src2} & \texttt{16:23} & \texttt{SourceEncoding} \\
\hyperref[enum:register_file_format_general_m]{\texttt{register\_format}} & \texttt{24:26} & \texttt{register\_file\_format\_general\_m} \\
\textit{Reserved} & \texttt{27} & \texttt{zero} \\
\hyperref[enum:vecsize_varying_m]{\texttt{vecsize}} & \texttt{28:29} & \texttt{vecsize\_varying\_m (def none)} \\
\hyperref[enum:message_slot_index_m]{\texttt{message\_slot\_index}} & \texttt{30:32} & \texttt{message\_slot\_index\_m} \\
\texttt{sr\_count} & \texttt{33:35} & \texttt{sr\_count\_t} \\
\textit{Reserved} & \texttt{36} & \texttt{zero} \\
\hyperref[enum:access_store_m]{\texttt{access}} & \texttt{37:39} & \texttt{access\_store\_m (def none)} \\
\texttt{staging\_register\_index} & \texttt{40:45} & \texttt{uint} \\
\texttt{staging\_register\_control\_32} & \texttt{46:47} & \texttt{exact=read} \\
\texttt{opcode1} & \texttt{48:51} & \texttt{exact=ST\_CVT} \\
\texttt{instruction\_class} & \texttt{52:56} & \texttt{exact=S3} \\
\texttt{fau\_page\_index} & \texttt{57:58} & \texttt{uint} \\
\hyperref[enum:flow_control_m]{\texttt{flow}} & \texttt{59:62} & \texttt{flow\_control\_m} \\
\textit{Reserved} & \texttt{63} & \texttt{zero} \\
\bottomrule
\end{tabularx}
\end{center}
\textbf{Fields Explained:}
\begin{tightitemize}\footnotesize
  \item \texttt{src0}: Primary source operand (\texttt{r0..r63}, \texttt{u0..u63} uniform, \texttt{k*} FAU constant, or \texttt{t0..t2} temporary; \texttt{\textasciicircum{}} marks discard/last-use). \emph{Example:} {\scriptsize\ttfamily \mnem{FADD.f32} \dst{r0}, \src{r1}, \src{r2}} (\textbf{SPIR-V:} \texttt{OpFAdd (Operand 0)}; \textbf{GLSL:} \texttt{float c = a + b;})
  \item \texttt{offset}: Signed byte offset added to base address register in memory load/store/atomic instructions. \emph{Example:} {\scriptsize\ttfamily \mnem{LOAD.i32} \dst{@r0}, \src{[r2:r3]}, \imm{\#16}} (\textbf{SPIR-V:} \texttt{OpAccessChain + OpLoad}; \textbf{GLSL:} \texttt{uint y = ssbo.arr[4];})
  \item \texttt{src2}: Third source operand (FMA addend, \texttt{MUX}/\texttt{CSEL} selector, \texttt{IADDC} carry-in, or bitwise mask in fused shift-logic). \emph{Example:} {\scriptsize\ttfamily \mnem{FMA.f32} \dst{r0}, \src{r1}, \src{r2}, \src{r3}} (\textbf{SPIR-V:} \texttt{GLSL.std.450 Fma (Addend)}; \textbf{GLSL:} \texttt{float d = fma(a, b, c);})
  \item \hyperref[enum:register_file_format_general_m]{\texttt{register\_format}}: Target GPR data conversion format (\texttt{auto32}, \texttt{f32}, \texttt{f16}, \texttt{s32}, \texttt{s16}, \texttt{u32}, \texttt{u16}). \emph{Example:} {\scriptsize\ttfamily \mnem{LD\_TILE}\imod{.f16.v4} \dst{@r0:r1}, \src{r2}} (\textbf{SPIR-V:} \texttt{OpImageRead (SubpassData)}; \textbf{GLSL:} \texttt{f16vec4 c = subpassLoad(inputAtt);})
  \item \hyperref[enum:vecsize_varying_m]{\texttt{vecsize}}: Component count (\texttt{none}=1, \texttt{v2}, \texttt{v3}, \texttt{v4}) for attribute, varying, or blend transfers. \emph{Example:} {\scriptsize\ttfamily \mnem{LD\_ATTR}\imod{.f32.v3} \dst{@r0:r2}, \src{r4}} (\textbf{SPIR-V:} \texttt{OpLoad (vec3 Input)}; \textbf{GLSL:} \texttt{in vec3 a\_normal;})
  \item \hyperref[enum:message_slot_index_m]{\texttt{message\_slot\_index}}: Asynchronous scoreboard slot (\texttt{slot0}--\texttt{slot2}, \texttt{slot7}) tracking completion of this message. \emph{Example:} {\scriptsize\ttfamily \mnem{LOAD.i32}\imod{.slot0} \dst{@r0}, \src{r2}, \imm{\#0}} (\textbf{SPIR-V:} \texttt{OpLoad}; \textbf{GLSL:} \texttt{uint v = ssbo.arr[i];})
  \item \texttt{sr\_count}: Number of consecutive 32-bit staging registers transferred by a load/store/texture message. \emph{Example:} {\scriptsize\ttfamily \mnem{LD\_MULTI.i128} \dst{@r0:r3}, \src{r4}, \imm{\#0}} (\textbf{SPIR-V:} \texttt{OpLoad (wide vector)}; \textbf{GLSL:} \texttt{uvec4 x = buf[i];})
  \item \hyperref[enum:access_store_m]{\texttt{access}}: Memory cache streaming/coherence policy (\texttt{none}, \texttt{istream}, \texttt{estream}, \texttt{force}). \emph{Example:} {\scriptsize\ttfamily \mnem{STORE.i32}\imod{.force} \src{@r0}, \src{r2}, \imm{\#0}} (\textbf{SPIR-V:} \texttt{OpStore (MakePointerVisible / Coherent)}; \textbf{GLSL:} \texttt{coherent buffer B \{ uint x; \}; x = 1u;})
  \item \texttt{staging\_register\_index}: Base GPR index of the contiguous staging register window (\texttt{@rN:rM}) for message payload/return. \emph{Example:} {\scriptsize\ttfamily \mnem{STORE.i64}\imod{.slot7} \src{@r0:r1}, \src{r2}, \imm{\#8}} (\textbf{SPIR-V:} \texttt{OpStore (uvec2)}; \textbf{GLSL:} \texttt{ssbo.pairs[i] = uvec2(a, b);})
  \item \texttt{staging\_register\_control\_32}: Controls the read/write staging register bank and count for async message instructions. \emph{Example:} {\scriptsize\ttfamily \mnem{LOAD.i128}\imod{.slot0} \dst{@r0:r3}, \src{r4}, \imm{\#0}} (\textbf{SPIR-V:} \texttt{OpLoad (vec4)}; \textbf{GLSL:} \texttt{vec4 v = ssbo.data[idx];})
  \item \texttt{opcode1}: Primary sub-opcode within the instruction class uniquely identifying the operation family. \emph{Example:} {\scriptsize\ttfamily \mnem{IADD.s32} \dst{r0}, \src{r1}, \src{r2}} (\textbf{SPIR-V:} \texttt{OpIAdd (operation selector)}; \textbf{GLSL:} \texttt{// Primary opcode bits})
  \item \texttt{instruction\_class}: 4-bit major instruction class selector (\texttt{bits 48:51}) partitioning the 64-bit Valhall opcode space (e.g., \texttt{FMA}, \texttt{CVT}, \texttt{SFU}, \texttt{LS}, \texttt{TEX}, \texttt{BRANCH}). \emph{Example:} {\scriptsize\ttfamily \mnem{FADD.f32} \dst{r0}, \src{r1}, \src{r2}} (\textbf{SPIR-V:} \texttt{N/A (encoding class)}; \textbf{GLSL:} \texttt{// Hardware execution pipe / decode class})
  \item \texttt{fau\_page\_index}: Selects the FAU (Fast Access Uniform) page (\texttt{0..3}) when a source operand references hardware small constants, thread IDs, or lane constants. \emph{Example:} {\scriptsize\ttfamily \mnem{FADD.f32} \dst{r0}, \src{r1}, \src{k1f}} (\textbf{SPIR-V:} \texttt{OpConstant / BuiltIn}; \textbf{GLSL:} \texttt{float y = x + 1.0;})
  \item \hyperref[enum:flow_control_m]{\texttt{flow}}: 4-bit embedded flow-control \& scoreboard wait modifier (\texttt{bits 59:62}: \texttt{none}, \texttt{wait0..wait0126}, \texttt{wait}, \texttt{reconverge}, \texttt{discard}, \texttt{end}). \emph{Example:} {\scriptsize\ttfamily \mnem{FADD.f32}\imod{.wait0.reconverge} \dst{r0}, \src{r1}, \src{r2}} (\textbf{SPIR-V:} \texttt{OpControlBarrier / OpReturn}; \textbf{GLSL:} \texttt{// Scoreboard wait + warp reconvergence})
\end{tightitemize}
\clearpage

\subsection*{ST\_MULTI}
\addcontentsline{toc}{subsection}{ST\_MULTI}
Multi-register store (v14+): stores \texttt{sr\_count} (\texttt{ls\_multi\_sr\_count\_m}) consecutive 32-bit staging registers to \texttt{src0 + sext64(offset)}.
\par\textbf{Notes:}
\begin{tightitemize}
  \item Wide bulk store beyond \texttt{STORE.i128}.
  \item \textbf{Source references:} \texttt{isa-v9-v14.xml}.
\end{tightitemize}
\codebox{\mnem{ST\_MULTI}\imod{\{.sr\_count\}}\imod{\{.message\_slot\_index\}}\imod{\{.flow\}}\imod{\{.access\}} \src{@sr}, \src{src0}, \imm{\#offset}}
\textbf{Operands:}
\begin{tightitemize}\footnotesize
  \item \texttt{\src{@sr}} (\texttt{staging\_register\_index} + \texttt{sr\_count}): Source staging register range (\texttt{sr\_count} 32-bit words) written sequentially to memory.
  \item \texttt{\src{src0}}: 64-bit base GPU virtual address.
  \item \texttt{\imm{\#offset}} (\texttt{offset}): Signed byte offset added to \texttt{src0}.
\end{tightitemize}
\par\textbf{Closest SPIR-V Op:} \texttt{OpStore (bulk \textgreater{}128-bit, v14+)}
\par\textbf{Pseudo-GLSL (Operational Semantics):}
\codebox{for (int i = 0; i \textless{} sr\_count; ++i) mem\_u32[src0 + \#offset + i * 4] = @sr[i];}
\textbf{Execution Context:}
\begin{tightitemize}\footnotesize
  \item \textbf{Direct 64-bit GPU VA memory/atomics:} Base addresses typically come from \hyperref[sec:csf-tsd-fau]{\texttt{FAU\_k}} (\texttt{u0..u31} populated by a Pilot shader chasing \texttt{SRT\_0} $\to$ \texttt{Resource[]} $\to$ \texttt{Buffer[]}), \texttt{FauSpecial::tls\_ptr}/\texttt{wls\_ptr} backed by \hyperref[sec:csf-tsd-fau]{\texttt{TSD\_k}} (\texttt{r24..r31}), or \texttt{FauSpecial::pc} for embedded constant LUTs.
  \item \textbf{Inline FAU \& Special Operands (\hyperref[sec:csf-tsd-fau]{\texttt{FAU\_k}}, \texttt{r8..r15}):} Any 8-bit source with \texttt{mode=2} reads zero-latency uniform \texttt{u0..u31} from \texttt{FAU\_k} (\texttt{slots64 = bits[63:56]}); \texttt{mode=3} reads \texttt{SmallConstant} (Table~1.2) or \texttt{FauSpecial} (\texttt{pc}, \texttt{lane\_id}, \texttt{lane\_width}, \texttt{tls\_ptr}, \texttt{wls\_ptr}, \texttt{blend\_descriptor\_0..7}).
\end{tightitemize}
\par\textbf{Worked Examples:}
\par\textbf{Example 1 (256-bit (8-GPR) bulk store (v14+)):}
\codebox{\mnem{ST\_MULTI}\imod{.i256.slot0} \dst{@r0:r7}, \src{[r8:r9]}, \imm{\#0}}
Stores 8 consecutive 32-bit registers (32 bytes) to \texttt{[r8:r9]} in a single v14 store message.
\par\textbf{Example 2 (Streaming 192-bit (6-GPR) store):}
\codebox{\mnem{ST\_MULTI}\imod{.i192.estream.slot7} \dst{@r0:r5}, \src{[r8:r9]}, \imm{\#0x10}}
Streams 24 bytes from \texttt{@r0:r5} to \texttt{[r8:r9] + 0x10} on \texttt{slot7}.

\begin{center}\footnotesize
\begin{tabularx}{\linewidth}{@{}lllX@{}}
\toprule
Variant & Arch & Class & Opcode Fields \\
\midrule
\texttt{(none)} & v9--v14 & \texttt{S2} & \texttt{opcode1=ST\_MULTI} \\
\bottomrule
\end{tabularx}
\end{center}
\begin{center}\footnotesize
\begin{tabularx}{0.88\linewidth}{@{}l l X@{}}
\toprule
Field Name & Bits & Type / Fixed Value \\
\midrule
\texttt{src0} & \texttt{0:7} & \texttt{SourceEncoding64} \\
\texttt{offset} & \texttt{8:23} & \texttt{int} \\
\hyperref[enum:access_store_m]{\texttt{access}} & \texttt{24:26} & \texttt{access\_store\_m (def none)} \\
\textit{Reserved} & \texttt{27:29} & \texttt{zero} \\
\hyperref[enum:message_slot_index_m]{\texttt{message\_slot\_index}} & \texttt{30:32} & \texttt{message\_slot\_index\_m} \\
\hyperref[enum:ls_multi_sr_count_m]{\texttt{sr\_count}} & \texttt{33:38} & \texttt{ls\_multi\_sr\_count\_m} \\
\textit{Reserved} & \texttt{39} & \texttt{zero} \\
\texttt{staging\_register\_index} & \texttt{40:45} & \texttt{uint} \\
\texttt{staging\_register\_control\_32} & \texttt{46:47} & \texttt{exact=read} \\
\texttt{opcode1} & \texttt{48:51} & \texttt{exact=ST\_MULTI} \\
\texttt{instruction\_class} & \texttt{52:56} & \texttt{exact=S2} \\
\texttt{fau\_page\_index} & \texttt{57:58} & \texttt{uint} \\
\hyperref[enum:flow_control_m]{\texttt{flow}} & \texttt{59:62} & \texttt{flow\_control\_m} \\
\textit{Reserved} & \texttt{63} & \texttt{zero} \\
\bottomrule
\end{tabularx}
\end{center}
\textbf{Fields Explained:}
\begin{tightitemize}\footnotesize
  \item \texttt{src0}: Primary source operand (\texttt{r0..r63}, \texttt{u0..u63} uniform, \texttt{k*} FAU constant, or \texttt{t0..t2} temporary; \texttt{\textasciicircum{}} marks discard/last-use). \emph{Example:} {\scriptsize\ttfamily \mnem{FADD.f32} \dst{r0}, \src{r1}, \src{r2}} (\textbf{SPIR-V:} \texttt{OpFAdd (Operand 0)}; \textbf{GLSL:} \texttt{float c = a + b;})
  \item \texttt{offset}: Signed byte offset added to base address register in memory load/store/atomic instructions. \emph{Example:} {\scriptsize\ttfamily \mnem{LOAD.i32} \dst{@r0}, \src{[r2:r3]}, \imm{\#16}} (\textbf{SPIR-V:} \texttt{OpAccessChain + OpLoad}; \textbf{GLSL:} \texttt{uint y = ssbo.arr[4];})
  \item \hyperref[enum:access_store_m]{\texttt{access}}: Memory cache streaming/coherence policy (\texttt{none}, \texttt{istream}, \texttt{estream}, \texttt{force}). \emph{Example:} {\scriptsize\ttfamily \mnem{STORE.i32}\imod{.force} \src{@r0}, \src{r2}, \imm{\#0}} (\textbf{SPIR-V:} \texttt{OpStore (MakePointerVisible / Coherent)}; \textbf{GLSL:} \texttt{coherent buffer B \{ uint x; \}; x = 1u;})
  \item \hyperref[enum:message_slot_index_m]{\texttt{message\_slot\_index}}: Asynchronous scoreboard slot (\texttt{slot0}--\texttt{slot2}, \texttt{slot7}) tracking completion of this message. \emph{Example:} {\scriptsize\ttfamily \mnem{LOAD.i32}\imod{.slot0} \dst{@r0}, \src{r2}, \imm{\#0}} (\textbf{SPIR-V:} \texttt{OpLoad}; \textbf{GLSL:} \texttt{uint v = ssbo.arr[i];})
  \item \hyperref[enum:ls_multi_sr_count_m]{\texttt{sr\_count}}: Number of consecutive 32-bit staging registers transferred by a load/store/texture message. \emph{Example:} {\scriptsize\ttfamily \mnem{LD\_MULTI.i128} \dst{@r0:r3}, \src{r4}, \imm{\#0}} (\textbf{SPIR-V:} \texttt{OpLoad (wide vector)}; \textbf{GLSL:} \texttt{uvec4 x = buf[i];})
  \item \texttt{staging\_register\_index}: Base GPR index of the contiguous staging register window (\texttt{@rN:rM}) for message payload/return. \emph{Example:} {\scriptsize\ttfamily \mnem{STORE.i64}\imod{.slot7} \src{@r0:r1}, \src{r2}, \imm{\#8}} (\textbf{SPIR-V:} \texttt{OpStore (uvec2)}; \textbf{GLSL:} \texttt{ssbo.pairs[i] = uvec2(a, b);})
  \item \texttt{staging\_register\_control\_32}: Controls the read/write staging register bank and count for async message instructions. \emph{Example:} {\scriptsize\ttfamily \mnem{LOAD.i128}\imod{.slot0} \dst{@r0:r3}, \src{r4}, \imm{\#0}} (\textbf{SPIR-V:} \texttt{OpLoad (vec4)}; \textbf{GLSL:} \texttt{vec4 v = ssbo.data[idx];})
  \item \texttt{opcode1}: Primary sub-opcode within the instruction class uniquely identifying the operation family. \emph{Example:} {\scriptsize\ttfamily \mnem{IADD.s32} \dst{r0}, \src{r1}, \src{r2}} (\textbf{SPIR-V:} \texttt{OpIAdd (operation selector)}; \textbf{GLSL:} \texttt{// Primary opcode bits})
  \item \texttt{instruction\_class}: 4-bit major instruction class selector (\texttt{bits 48:51}) partitioning the 64-bit Valhall opcode space (e.g., \texttt{FMA}, \texttt{CVT}, \texttt{SFU}, \texttt{LS}, \texttt{TEX}, \texttt{BRANCH}). \emph{Example:} {\scriptsize\ttfamily \mnem{FADD.f32} \dst{r0}, \src{r1}, \src{r2}} (\textbf{SPIR-V:} \texttt{N/A (encoding class)}; \textbf{GLSL:} \texttt{// Hardware execution pipe / decode class})
  \item \texttt{fau\_page\_index}: Selects the FAU (Fast Access Uniform) page (\texttt{0..3}) when a source operand references hardware small constants, thread IDs, or lane constants. \emph{Example:} {\scriptsize\ttfamily \mnem{FADD.f32} \dst{r0}, \src{r1}, \src{k1f}} (\textbf{SPIR-V:} \texttt{OpConstant / BuiltIn}; \textbf{GLSL:} \texttt{float y = x + 1.0;})
  \item \hyperref[enum:flow_control_m]{\texttt{flow}}: 4-bit embedded flow-control \& scoreboard wait modifier (\texttt{bits 59:62}: \texttt{none}, \texttt{wait0..wait0126}, \texttt{wait}, \texttt{reconverge}, \texttt{discard}, \texttt{end}). \emph{Example:} {\scriptsize\ttfamily \mnem{FADD.f32}\imod{.wait0.reconverge} \dst{r0}, \src{r1}, \src{r2}} (\textbf{SPIR-V:} \texttt{OpControlBarrier / OpReturn}; \textbf{GLSL:} \texttt{// Scoreboard wait + warp reconvergence})
\end{tightitemize}
\clearpage

\section{Atomics (6 mnemonics)}
\subsection*{ACMPXCHG}
\addcontentsline{toc}{subsection}{ACMPXCHG}
Atomic compare-and-swap (\texttt{i32}/\texttt{i64}) at \texttt{src0 + offset}: \texttt{sr\_src} holds \texttt{(compare, swap)} pair (\texttt{V2I32}/\texttt{V2I64}); if \texttt{*addr == compare}, \texttt{*addr = swap}; returns old \texttt{*addr} in \texttt{sr\_dst}.
\par\textbf{Notes:}
\begin{tightitemize}
  \item \texttt{sr\_src} is 2x width (\texttt{V2IN}, \texttt{ops.rs:100}).
  \item \textbf{Source references:} \texttt{kraid/ops.rs:96-122}, \texttt{encode\_v9.rs:565-580}.
\end{tightitemize}
\codebox{\mnem{ACMPXCHG}\imod{\{.message\_slot\_index\}}\imod{\{.flow\}} \dst{@sw}, \src{@sr}, \src{src0}, \imm{\#offset}}
\textbf{Operands:}
\begin{tightitemize}\footnotesize
  \item \texttt{\dst{@sw}} (\texttt{staging\_register\_write\_index} + \texttt{sr\_write\_count}): Destination staging register(s) receiving the pre-modification memory value.
  \item \texttt{\src{@sr}} (\texttt{staging\_register\_index} + \texttt{sr\_count}): Source staging register pair holding \texttt{(compare\_value, swap\_value)}.
  \item \texttt{\src{src0}}: 64-bit base GPU virtual address.
  \item \texttt{\imm{\#offset}} (\texttt{offset}): Unsigned byte offset added to \texttt{src0}.
\end{tightitemize}
\par\textbf{Closest SPIR-V Op:} \texttt{OpAtomicCompareExchange}
\par\textbf{Pseudo-GLSL (Operational Semantics):}
\codebox{uint64\_t addr = src0 + \#offset;\\@sw = atomicCompSwap(mem[addr], @sr[0] /*cmp*/, @sr[1] /*swap*/);}
\textbf{Execution Context:}
\begin{tightitemize}\footnotesize
  \item \textbf{Direct 64-bit GPU VA memory/atomics:} Base addresses typically come from \hyperref[sec:csf-tsd-fau]{\texttt{FAU\_k}} (\texttt{u0..u31} populated by a Pilot shader chasing \texttt{SRT\_0} $\to$ \texttt{Resource[]} $\to$ \texttt{Buffer[]}), \texttt{FauSpecial::tls\_ptr}/\texttt{wls\_ptr} backed by \hyperref[sec:csf-tsd-fau]{\texttt{TSD\_k}} (\texttt{r24..r31}), or \texttt{FauSpecial::pc} for embedded constant LUTs.
  \item \textbf{Inline FAU \& Special Operands (\hyperref[sec:csf-tsd-fau]{\texttt{FAU\_k}}, \texttt{r8..r15}):} Any 8-bit source with \texttt{mode=2} reads zero-latency uniform \texttt{u0..u31} from \texttt{FAU\_k} (\texttt{slots64 = bits[63:56]}); \texttt{mode=3} reads \texttt{SmallConstant} (Table~1.2) or \texttt{FauSpecial} (\texttt{pc}, \texttt{lane\_id}, \texttt{lane\_width}, \texttt{tls\_ptr}, \texttt{wls\_ptr}, \texttt{blend\_descriptor\_0..7}).
\end{tightitemize}
\par\textbf{Worked Examples:}
\par\textbf{Example 1 (Successful 32-bit CAS):}
\codebox{\mnem{ACMPXCHG.i32}\imod{.slot0} \dst{@r0}, \src{[r2:r3]}, \imm{\#0}}
With \texttt{*addr == 10}, compare \texttt{r0 = 10}, swap \texttt{r1 = 42}: updates \texttt{*addr = 42} and returns old value \texttt{r0 = 10}.
\par\textbf{Example 2 (Failed 32-bit CAS (memory unchanged)):}
\codebox{\mnem{ACMPXCHG.i32}\imod{.slot0} \dst{@r0}, \src{[r2:r3]}, \imm{\#0}}
With \texttt{*addr == 99}, compare \texttt{r0 = 10}, swap \texttt{r1 = 42}: leaves \texttt{*addr = 99} untouched and returns observed \texttt{r0 = 99}.

\begin{center}\footnotesize
\begin{tabularx}{\linewidth}{@{}lllX@{}}
\toprule
Variant & Arch & Class & Opcode Fields \\
\midrule
\texttt{(none)} & v9--v14 & \texttt{S2RW} & \texttt{opcode1=ATOM\_RETURN} \\
\texttt{(none)} & v9--v14 & \texttt{S2RW} & \texttt{opcode1=ATOM\_RETURN} \\
\bottomrule
\end{tabularx}
\end{center}
\begin{center}\footnotesize
\begin{tabularx}{0.88\linewidth}{@{}l l X@{}}
\toprule
Field Name & Bits & Type / Fixed Value \\
\midrule
\texttt{src0} & \texttt{0:7} & \texttt{SourceEncoding64} \\
\texttt{offset} & \texttt{8:15} & \texttt{uint} \\
\texttt{staging\_register\_write\_index} & \texttt{16:21} & \texttt{uint} \\
\hyperref[enum:atom_operation_m]{\texttt{atom\_opc}} & \texttt{22:25} & \texttt{exact=axchg} \\
\texttt{compare} & \texttt{26} & \texttt{exact=compare} \\
\texttt{memory\_width} & \texttt{27:29} & \texttt{exact=i32} \\
\hyperref[enum:message_slot_index_m]{\texttt{message\_slot\_index}} & \texttt{30:32} & \texttt{message\_slot\_index\_m} \\
\texttt{sr\_count} & \texttt{33:35} & \texttt{exact=count2} \\
\texttt{sr\_write\_count} & \texttt{36:38} & \texttt{exact=write1} \\
\textit{Reserved} & \texttt{39} & \texttt{zero} \\
\texttt{staging\_register\_index} & \texttt{40:45} & \texttt{uint} \\
\texttt{staging\_register\_control\_32} & \texttt{46:47} & \texttt{exact=read\_write} \\
\texttt{opcode1} & \texttt{48:51} & \texttt{exact=ATOM\_RETURN} \\
\texttt{instruction\_class} & \texttt{52:56} & \texttt{exact=S2RW} \\
\texttt{fau\_page\_index} & \texttt{57:58} & \texttt{uint} \\
\hyperref[enum:flow_control_m]{\texttt{flow}} & \texttt{59:62} & \texttt{flow\_control\_m} \\
\textit{Reserved} & \texttt{63} & \texttt{zero} \\
\bottomrule
\end{tabularx}
\end{center}
\textbf{Fields Explained:}
\begin{tightitemize}\footnotesize
  \item \texttt{src0}: Primary source operand (\texttt{r0..r63}, \texttt{u0..u63} uniform, \texttt{k*} FAU constant, or \texttt{t0..t2} temporary; \texttt{\textasciicircum{}} marks discard/last-use). \emph{Example:} {\scriptsize\ttfamily \mnem{FADD.f32} \dst{r0}, \src{r1}, \src{r2}} (\textbf{SPIR-V:} \texttt{OpFAdd (Operand 0)}; \textbf{GLSL:} \texttt{float c = a + b;})
  \item \texttt{offset}: Signed byte offset added to base address register in memory load/store/atomic instructions. \emph{Example:} {\scriptsize\ttfamily \mnem{LOAD.i32} \dst{@r0}, \src{[r2:r3]}, \imm{\#16}} (\textbf{SPIR-V:} \texttt{OpAccessChain + OpLoad}; \textbf{GLSL:} \texttt{uint y = ssbo.arr[4];})
  \item \texttt{staging\_register\_write\_index}: Separate destination staging register window base when read and write staging windows differ. \emph{Example:} {\scriptsize\ttfamily \mnem{ATOM\_RETURN.i32}\imod{.aadd} \dst{@r0}, \src{@r1}, \src{r2}, \imm{\#0}} (\textbf{SPIR-V:} \texttt{OpAtomicIAdd}; \textbf{GLSL:} \texttt{uint old = atomicAdd(ssbo.cnt, 1u);})
  \item \hyperref[enum:atom_operation_m]{\texttt{atom\_opc}}: Selects the atomic read-modify-write ALU operation (\texttt{aadd}, \texttt{asmin}, \texttt{aumax}, \texttt{aand}, \texttt{aor}, \texttt{axor}, \texttt{axchg}). \emph{Example:} {\scriptsize\ttfamily \mnem{ATOM.i32}\imod{.aumax} \src{@r0}, \src{r2}, \imm{\#0}} (\textbf{SPIR-V:} \texttt{OpAtomicUMax}; \textbf{GLSL:} \texttt{atomicMax(ssbo.val, x);})
  \item \texttt{compare}: Enables depth comparison predicate in texture/shadow lookup. \emph{Example:} {\scriptsize\ttfamily \mnem{TEX\_SINGLE}\imod{.tex2d.compare} \dst{@r0}, \src{@r4:r6}} (\textbf{SPIR-V:} \texttt{OpImageSampleDrefImplicitLod}; \textbf{GLSL:} \texttt{float s = texture(shadow, vec3(uv, z));})
  \item \texttt{memory\_width}: Memory transfer width (\texttt{i8}, \texttt{i16}, \texttt{i24}, \texttt{i32}, \texttt{i48}, \texttt{i64}, \texttt{i96}, \texttt{i128}). \emph{Example:} {\scriptsize\ttfamily \mnem{LOAD.i64} \dst{@r0:r1}, \src{[r2:r3]}, \imm{\#0}} (\textbf{SPIR-V:} \texttt{OpLoad (uvec2)}; \textbf{GLSL:} \texttt{uvec2 v = ssbo.pairs[i];})
  \item \hyperref[enum:message_slot_index_m]{\texttt{message\_slot\_index}}: Asynchronous scoreboard slot (\texttt{slot0}--\texttt{slot2}, \texttt{slot7}) tracking completion of this message. \emph{Example:} {\scriptsize\ttfamily \mnem{LOAD.i32}\imod{.slot0} \dst{@r0}, \src{r2}, \imm{\#0}} (\textbf{SPIR-V:} \texttt{OpLoad}; \textbf{GLSL:} \texttt{uint v = ssbo.arr[i];})
  \item \texttt{sr\_count}: Number of consecutive 32-bit staging registers transferred by a load/store/texture message. \emph{Example:} {\scriptsize\ttfamily \mnem{LD\_MULTI.i128} \dst{@r0:r3}, \src{r4}, \imm{\#0}} (\textbf{SPIR-V:} \texttt{OpLoad (wide vector)}; \textbf{GLSL:} \texttt{uvec4 x = buf[i];})
  \item \texttt{sr\_write\_count}: Number of consecutive 32-bit staging registers written back by a dual-staging message (e.g. ATOM\textbackslash\{\}\_RETURN / TEX). \emph{Example:} {\scriptsize\ttfamily \mnem{TEX\_SINGLE}\imod{.tex2d.f32.rgba} \dst{@r0:r3}, \src{@r4:r5}} (\textbf{SPIR-V:} \texttt{OpImageSampleImplicitLod}; \textbf{GLSL:} \texttt{vec4 c = texture(tex, uv);})
  \item \texttt{staging\_register\_index}: Base GPR index of the contiguous staging register window (\texttt{@rN:rM}) for message payload/return. \emph{Example:} {\scriptsize\ttfamily \mnem{STORE.i64}\imod{.slot7} \src{@r0:r1}, \src{r2}, \imm{\#8}} (\textbf{SPIR-V:} \texttt{OpStore (uvec2)}; \textbf{GLSL:} \texttt{ssbo.pairs[i] = uvec2(a, b);})
  \item \texttt{staging\_register\_control\_32}: Controls the read/write staging register bank and count for async message instructions. \emph{Example:} {\scriptsize\ttfamily \mnem{LOAD.i128}\imod{.slot0} \dst{@r0:r3}, \src{r4}, \imm{\#0}} (\textbf{SPIR-V:} \texttt{OpLoad (vec4)}; \textbf{GLSL:} \texttt{vec4 v = ssbo.data[idx];})
  \item \texttt{opcode1}: Primary sub-opcode within the instruction class uniquely identifying the operation family. \emph{Example:} {\scriptsize\ttfamily \mnem{IADD.s32} \dst{r0}, \src{r1}, \src{r2}} (\textbf{SPIR-V:} \texttt{OpIAdd (operation selector)}; \textbf{GLSL:} \texttt{// Primary opcode bits})
  \item \texttt{instruction\_class}: 4-bit major instruction class selector (\texttt{bits 48:51}) partitioning the 64-bit Valhall opcode space (e.g., \texttt{FMA}, \texttt{CVT}, \texttt{SFU}, \texttt{LS}, \texttt{TEX}, \texttt{BRANCH}). \emph{Example:} {\scriptsize\ttfamily \mnem{FADD.f32} \dst{r0}, \src{r1}, \src{r2}} (\textbf{SPIR-V:} \texttt{N/A (encoding class)}; \textbf{GLSL:} \texttt{// Hardware execution pipe / decode class})
  \item \texttt{fau\_page\_index}: Selects the FAU (Fast Access Uniform) page (\texttt{0..3}) when a source operand references hardware small constants, thread IDs, or lane constants. \emph{Example:} {\scriptsize\ttfamily \mnem{FADD.f32} \dst{r0}, \src{r1}, \src{k1f}} (\textbf{SPIR-V:} \texttt{OpConstant / BuiltIn}; \textbf{GLSL:} \texttt{float y = x + 1.0;})
  \item \hyperref[enum:flow_control_m]{\texttt{flow}}: 4-bit embedded flow-control \& scoreboard wait modifier (\texttt{bits 59:62}: \texttt{none}, \texttt{wait0..wait0126}, \texttt{wait}, \texttt{reconverge}, \texttt{discard}, \texttt{end}). \emph{Example:} {\scriptsize\ttfamily \mnem{FADD.f32}\imod{.wait0.reconverge} \dst{r0}, \src{r1}, \src{r2}} (\textbf{SPIR-V:} \texttt{OpControlBarrier / OpReturn}; \textbf{GLSL:} \texttt{// Scoreboard wait + warp reconvergence})
\end{tightitemize}
\clearpage

\subsection*{ATOM}
\addcontentsline{toc}{subsection}{ATOM}
Fire-and-forget atomic RMW (\texttt{i32}/\texttt{i64}) at \texttt{src0 + offset} with operand \texttt{sr\_src} and op \texttt{atom\_opc} (\texttt{aadd}/\texttt{aaddu}/\texttt{aadds}/\texttt{asmin}/\texttt{asmax}/\texttt{aumin}/\texttt{aumax}/\texttt{aand}/\texttt{aor}/\texttt{axor}/\texttt{axchg}). No return value.
\par\textbf{Notes:}
\begin{tightitemize}
  \item \texttt{atom\_opc} selects the exact RMW operation.
  \item \texttt{aaddu}/\texttt{aadds} are 64-bit atomics with 32-bit widened operand.
  \item \textbf{Source references:} \texttt{kraid/ops.rs:176-276}, \texttt{encode\_v9.rs:620-700}.
\end{tightitemize}
\codebox{\mnem{ATOM}\imod{\{.message\_slot\_index\}}\imod{\{.flow\}}\imod{\{.atom\_opc\}} \src{@sr}, \src{src0}, \imm{\#offset}}
\textbf{Operands:}
\begin{tightitemize}\footnotesize
  \item \texttt{\src{@sr}} (\texttt{staging\_register\_index} + \texttt{sr\_count}): Source staging register(s) holding the atomic operand value (no return value written).
  \item \texttt{\src{src0}}: 64-bit base GPU virtual address.
  \item \texttt{\imm{\#offset}} (\texttt{offset}): Unsigned byte offset added to \texttt{src0}.
\end{tightitemize}
\par\textbf{Closest SPIR-V Op:} \texttt{OpAtomicIAdd/Min/Max/And/Or/Xor (unused result)}
\par\textbf{Pseudo-GLSL (Operational Semantics):}
\codebox{atomicOp(atom\_opc, mem[src0 + \#offset], @sr); // fire-and-forget RMW}
\textbf{Execution Context:}
\begin{tightitemize}\footnotesize
  \item \textbf{Direct 64-bit GPU VA memory/atomics:} Base addresses typically come from \hyperref[sec:csf-tsd-fau]{\texttt{FAU\_k}} (\texttt{u0..u31} populated by a Pilot shader chasing \texttt{SRT\_0} $\to$ \texttt{Resource[]} $\to$ \texttt{Buffer[]}), \texttt{FauSpecial::tls\_ptr}/\texttt{wls\_ptr} backed by \hyperref[sec:csf-tsd-fau]{\texttt{TSD\_k}} (\texttt{r24..r31}), or \texttt{FauSpecial::pc} for embedded constant LUTs.
  \item \textbf{Inline FAU \& Special Operands (\hyperref[sec:csf-tsd-fau]{\texttt{FAU\_k}}, \texttt{r8..r15}):} Any 8-bit source with \texttt{mode=2} reads zero-latency uniform \texttt{u0..u31} from \texttt{FAU\_k} (\texttt{slots64 = bits[63:56]}); \texttt{mode=3} reads \texttt{SmallConstant} (Table~1.2) or \texttt{FauSpecial} (\texttt{pc}, \texttt{lane\_id}, \texttt{lane\_width}, \texttt{tls\_ptr}, \texttt{wls\_ptr}, \texttt{blend\_descriptor\_0..7}).
\end{tightitemize}
\par\textbf{Worked Examples:}
\par\textbf{Example 1 (Fire-and-forget atomic add):}
\codebox{\mnem{ATOM.i32}\imod{.aadd.slot0} \dst{@r0}, \src{[r2:r3]}, \imm{\#0}}
With \texttt{*addr = 100}, \texttt{r0 = 5}: atomically updates \texttt{*addr = 105} without returning the old value.
\par\textbf{Example 2 (Fire-and-forget atomic bitwise OR):}
\codebox{\mnem{ATOM.i32}\imod{.aor.slot7} \dst{@r0}, \src{[r2:r3]}, \imm{\#0x4}}
Atomically sets bits \texttt{r0} in \texttt{*([r2:r3] + 4)} on non-blocking \texttt{slot7}.
\par\textbf{Mined SPIR-V $\to$ Valhall ISA Example (Fire-and-forget unsigned atomic maximum (return value unused) --- Fragment shader, Mali-G710 (v10)):}
\codebox{\textbf{; SPIR-V (input)}\\\%old = OpAtomicUMax \%uint \%ptr \%uint\_1 \%uint\_0 \%val ; \%old unused\\[3pt]\textbf{; Valhall ISA (malioc output)}\\IADD.u64.wait0 [r50:r51], u22, [r48:r49].w1\\ATOM.i32.slot0.aumax @r48, [r50\textasciicircum{}:r51\textasciicircum{}], \#0x0}
Because the return value of \texttt{OpAtomicUMax} is unused in SPIR-V, \texttt{malioc} selects non-returning \texttt{ATOM.i32.aumax} instead of \texttt{ATOM\_RETURN}, avoiding a destination staging register writeback and scoreboard stall.
\begin{center}\footnotesize
\begin{tabularx}{\linewidth}{@{}lllX@{}}
\toprule
Variant & Arch & Class & Opcode Fields \\
\midrule
\texttt{(none)} & v9--v14 & \texttt{S2} & \texttt{opcode1=ATOM} \\
\texttt{(none)} & v9--v14 & \texttt{S2} & \texttt{opcode1=ATOM} \\
\bottomrule
\end{tabularx}
\end{center}
\begin{center}\footnotesize
\begin{tabularx}{0.88\linewidth}{@{}l l X@{}}
\toprule
Field Name & Bits & Type / Fixed Value \\
\midrule
\texttt{src0} & \texttt{0:7} & \texttt{SourceEncoding64} \\
\texttt{offset} & \texttt{8:15} & \texttt{uint} \\
\textit{Reserved} & \texttt{16:21} & \texttt{zero} \\
\hyperref[enum:atom_operation_m]{\texttt{atom\_opc}} & \texttt{22:25} & \texttt{atom\_operation\_m} \\
\texttt{compare} & \texttt{26} & \texttt{exact=no\_compare} \\
\texttt{memory\_width} & \texttt{27:29} & \texttt{exact=i32} \\
\hyperref[enum:message_slot_index_m]{\texttt{message\_slot\_index}} & \texttt{30:32} & \texttt{message\_slot\_index\_m} \\
\texttt{sr\_count} & \texttt{33:35} & \texttt{exact=count1} \\
\textit{Reserved} & \texttt{36:39} & \texttt{zero} \\
\texttt{staging\_register\_index} & \texttt{40:45} & \texttt{uint} \\
\texttt{staging\_register\_control\_32} & \texttt{46:47} & \texttt{exact=read} \\
\texttt{opcode1} & \texttt{48:51} & \texttt{exact=ATOM} \\
\texttt{instruction\_class} & \texttt{52:56} & \texttt{exact=S2} \\
\texttt{fau\_page\_index} & \texttt{57:58} & \texttt{uint} \\
\hyperref[enum:flow_control_m]{\texttt{flow}} & \texttt{59:62} & \texttt{flow\_control\_m} \\
\textit{Reserved} & \texttt{63} & \texttt{zero} \\
\bottomrule
\end{tabularx}
\end{center}
\textbf{Fields Explained:}
\begin{tightitemize}\footnotesize
  \item \texttt{src0}: Primary source operand (\texttt{r0..r63}, \texttt{u0..u63} uniform, \texttt{k*} FAU constant, or \texttt{t0..t2} temporary; \texttt{\textasciicircum{}} marks discard/last-use). \emph{Example:} {\scriptsize\ttfamily \mnem{FADD.f32} \dst{r0}, \src{r1}, \src{r2}} (\textbf{SPIR-V:} \texttt{OpFAdd (Operand 0)}; \textbf{GLSL:} \texttt{float c = a + b;})
  \item \texttt{offset}: Signed byte offset added to base address register in memory load/store/atomic instructions. \emph{Example:} {\scriptsize\ttfamily \mnem{LOAD.i32} \dst{@r0}, \src{[r2:r3]}, \imm{\#16}} (\textbf{SPIR-V:} \texttt{OpAccessChain + OpLoad}; \textbf{GLSL:} \texttt{uint y = ssbo.arr[4];})
  \item \hyperref[enum:atom_operation_m]{\texttt{atom\_opc}}: Selects the atomic read-modify-write ALU operation (\texttt{aadd}, \texttt{asmin}, \texttt{aumax}, \texttt{aand}, \texttt{aor}, \texttt{axor}, \texttt{axchg}). \emph{Example:} {\scriptsize\ttfamily \mnem{ATOM.i32}\imod{.aumax} \src{@r0}, \src{r2}, \imm{\#0}} (\textbf{SPIR-V:} \texttt{OpAtomicUMax}; \textbf{GLSL:} \texttt{atomicMax(ssbo.val, x);})
  \item \texttt{compare}: Enables depth comparison predicate in texture/shadow lookup. \emph{Example:} {\scriptsize\ttfamily \mnem{TEX\_SINGLE}\imod{.tex2d.compare} \dst{@r0}, \src{@r4:r6}} (\textbf{SPIR-V:} \texttt{OpImageSampleDrefImplicitLod}; \textbf{GLSL:} \texttt{float s = texture(shadow, vec3(uv, z));})
  \item \texttt{memory\_width}: Memory transfer width (\texttt{i8}, \texttt{i16}, \texttt{i24}, \texttt{i32}, \texttt{i48}, \texttt{i64}, \texttt{i96}, \texttt{i128}). \emph{Example:} {\scriptsize\ttfamily \mnem{LOAD.i64} \dst{@r0:r1}, \src{[r2:r3]}, \imm{\#0}} (\textbf{SPIR-V:} \texttt{OpLoad (uvec2)}; \textbf{GLSL:} \texttt{uvec2 v = ssbo.pairs[i];})
  \item \hyperref[enum:message_slot_index_m]{\texttt{message\_slot\_index}}: Asynchronous scoreboard slot (\texttt{slot0}--\texttt{slot2}, \texttt{slot7}) tracking completion of this message. \emph{Example:} {\scriptsize\ttfamily \mnem{LOAD.i32}\imod{.slot0} \dst{@r0}, \src{r2}, \imm{\#0}} (\textbf{SPIR-V:} \texttt{OpLoad}; \textbf{GLSL:} \texttt{uint v = ssbo.arr[i];})
  \item \texttt{sr\_count}: Number of consecutive 32-bit staging registers transferred by a load/store/texture message. \emph{Example:} {\scriptsize\ttfamily \mnem{LD\_MULTI.i128} \dst{@r0:r3}, \src{r4}, \imm{\#0}} (\textbf{SPIR-V:} \texttt{OpLoad (wide vector)}; \textbf{GLSL:} \texttt{uvec4 x = buf[i];})
  \item \texttt{staging\_register\_index}: Base GPR index of the contiguous staging register window (\texttt{@rN:rM}) for message payload/return. \emph{Example:} {\scriptsize\ttfamily \mnem{STORE.i64}\imod{.slot7} \src{@r0:r1}, \src{r2}, \imm{\#8}} (\textbf{SPIR-V:} \texttt{OpStore (uvec2)}; \textbf{GLSL:} \texttt{ssbo.pairs[i] = uvec2(a, b);})
  \item \texttt{staging\_register\_control\_32}: Controls the read/write staging register bank and count for async message instructions. \emph{Example:} {\scriptsize\ttfamily \mnem{LOAD.i128}\imod{.slot0} \dst{@r0:r3}, \src{r4}, \imm{\#0}} (\textbf{SPIR-V:} \texttt{OpLoad (vec4)}; \textbf{GLSL:} \texttt{vec4 v = ssbo.data[idx];})
  \item \texttt{opcode1}: Primary sub-opcode within the instruction class uniquely identifying the operation family. \emph{Example:} {\scriptsize\ttfamily \mnem{IADD.s32} \dst{r0}, \src{r1}, \src{r2}} (\textbf{SPIR-V:} \texttt{OpIAdd (operation selector)}; \textbf{GLSL:} \texttt{// Primary opcode bits})
  \item \texttt{instruction\_class}: 4-bit major instruction class selector (\texttt{bits 48:51}) partitioning the 64-bit Valhall opcode space (e.g., \texttt{FMA}, \texttt{CVT}, \texttt{SFU}, \texttt{LS}, \texttt{TEX}, \texttt{BRANCH}). \emph{Example:} {\scriptsize\ttfamily \mnem{FADD.f32} \dst{r0}, \src{r1}, \src{r2}} (\textbf{SPIR-V:} \texttt{N/A (encoding class)}; \textbf{GLSL:} \texttt{// Hardware execution pipe / decode class})
  \item \texttt{fau\_page\_index}: Selects the FAU (Fast Access Uniform) page (\texttt{0..3}) when a source operand references hardware small constants, thread IDs, or lane constants. \emph{Example:} {\scriptsize\ttfamily \mnem{FADD.f32} \dst{r0}, \src{r1}, \src{k1f}} (\textbf{SPIR-V:} \texttt{OpConstant / BuiltIn}; \textbf{GLSL:} \texttt{float y = x + 1.0;})
  \item \hyperref[enum:flow_control_m]{\texttt{flow}}: 4-bit embedded flow-control \& scoreboard wait modifier (\texttt{bits 59:62}: \texttt{none}, \texttt{wait0..wait0126}, \texttt{wait}, \texttt{reconverge}, \texttt{discard}, \texttt{end}). \emph{Example:} {\scriptsize\ttfamily \mnem{FADD.f32}\imod{.wait0.reconverge} \dst{r0}, \src{r1}, \src{r2}} (\textbf{SPIR-V:} \texttt{OpControlBarrier / OpReturn}; \textbf{GLSL:} \texttt{// Scoreboard wait + warp reconvergence})
\end{tightitemize}
\clearpage

\subsection*{ATOM1}
\addcontentsline{toc}{subsection}{ATOM1}
Fire-and-forget unary atomic RMW (\texttt{i32}/\texttt{i64}) at \texttt{src0 + offset} with implicit operand \texttt{1}/\texttt{-1}: \texttt{ainc} (\texttt{+1}), \texttt{adec} (\texttt{-1}), \texttt{aor1} (\texttt{\textbar{}1}), \texttt{asmax1} (\texttt{smax 1}), \texttt{aumax1} (\texttt{umax 1}).
\par\textbf{Notes:}
\begin{tightitemize}
  \item No \texttt{sr\_src} operand.
  \item Constant \texttt{1}/\texttt{-1} implicit (\texttt{ops.rs:190-199}).
  \item \textbf{Source references:} \texttt{kraid/ops.rs:189-237, 281-313}, \texttt{encode\_v9.rs:702-736}.
\end{tightitemize}
\codebox{\mnem{ATOM1}\imod{\{.message\_slot\_index\}}\imod{\{.flow\}}\imod{\{.atom1\_opc\}} \src{src0}, \imm{\#offset}}
\textbf{Operands:}
\begin{tightitemize}\footnotesize
  \item \texttt{\src{src0}}: 64-bit base GPU virtual address for unary atomic RMW (\texttt{ainc}/\texttt{adec}/etc.).
  \item \texttt{\imm{\#offset}} (\texttt{offset}): Unsigned byte offset added to \texttt{src0}.
\end{tightitemize}
\par\textbf{Closest SPIR-V Op:} \texttt{OpAtomicIIncrement / IDecrement (unused result)}
\par\textbf{Pseudo-GLSL (Operational Semantics):}
\codebox{atomicUnary(atom1\_opc, mem[src0 + \#offset]); // ainc/adec/asmax1/aumax1/aor1}
\textbf{Execution Context:}
\begin{tightitemize}\footnotesize
  \item \textbf{Direct 64-bit GPU VA memory/atomics:} Base addresses typically come from \hyperref[sec:csf-tsd-fau]{\texttt{FAU\_k}} (\texttt{u0..u31} populated by a Pilot shader chasing \texttt{SRT\_0} $\to$ \texttt{Resource[]} $\to$ \texttt{Buffer[]}), \texttt{FauSpecial::tls\_ptr}/\texttt{wls\_ptr} backed by \hyperref[sec:csf-tsd-fau]{\texttt{TSD\_k}} (\texttt{r24..r31}), or \texttt{FauSpecial::pc} for embedded constant LUTs.
  \item \textbf{Inline FAU \& Special Operands (\hyperref[sec:csf-tsd-fau]{\texttt{FAU\_k}}, \texttt{r8..r15}):} Any 8-bit source with \texttt{mode=2} reads zero-latency uniform \texttt{u0..u31} from \texttt{FAU\_k} (\texttt{slots64 = bits[63:56]}); \texttt{mode=3} reads \texttt{SmallConstant} (Table~1.2) or \texttt{FauSpecial} (\texttt{pc}, \texttt{lane\_id}, \texttt{lane\_width}, \texttt{tls\_ptr}, \texttt{wls\_ptr}, \texttt{blend\_descriptor\_0..7}).
\end{tightitemize}
\par\textbf{Worked Examples:}
\par\textbf{Example 1 (Fire-and-forget atomic increment):}
\codebox{\mnem{ATOM1.i32}\imod{.ainc.slot0} \src{[r2:r3]}, \imm{\#0}}
With \texttt{*addr = 7}: atomically increments \texttt{*addr} to \texttt{8} without needing a value staging register.
\par\textbf{Example 2 (Fire-and-forget atomic decrement):}
\codebox{\mnem{ATOM1.i32}\imod{.adec.slot7} \src{[r2:r3]}, \imm{\#0x8}}
Atomically decrements \texttt{*([r2:r3] + 8)} by \texttt{1}.

\begin{center}\footnotesize
\begin{tabularx}{\linewidth}{@{}lllX@{}}
\toprule
Variant & Arch & Class & Opcode Fields \\
\midrule
\texttt{(none)} & v9--v14 & \texttt{S2} & \texttt{opcode1=ATOM1} \\
\texttt{(none)} & v9--v14 & \texttt{S2} & \texttt{opcode1=ATOM1} \\
\bottomrule
\end{tabularx}
\end{center}
\begin{center}\footnotesize
\begin{tabularx}{0.88\linewidth}{@{}l l X@{}}
\toprule
Field Name & Bits & Type / Fixed Value \\
\midrule
\texttt{src0} & \texttt{0:7} & \texttt{SourceEncoding64} \\
\texttt{offset} & \texttt{8:15} & \texttt{uint} \\
\textit{Reserved} & \texttt{16:21} & \texttt{zero} \\
\hyperref[enum:atom1_operation_m]{\texttt{atom1\_opc}} & \texttt{22:24} & \texttt{atom1\_operation\_m} \\
\textit{Reserved} & \texttt{25} & \texttt{zero} \\
\texttt{compare} & \texttt{26} & \texttt{exact=no\_compare} \\
\texttt{memory\_width} & \texttt{27:29} & \texttt{exact=i32} \\
\hyperref[enum:message_slot_index_m]{\texttt{message\_slot\_index}} & \texttt{30:32} & \texttt{message\_slot\_index\_m} \\
\texttt{sr\_count} & \texttt{33:35} & \texttt{exact=count0} \\
\textit{Reserved} & \texttt{36:39} & \texttt{zero} \\
\texttt{staging\_register\_index} & \texttt{40:45} & \texttt{exact=0} \\
\texttt{staging\_register\_control\_32} & \texttt{46:47} & \texttt{exact=read} \\
\texttt{opcode1} & \texttt{48:51} & \texttt{exact=ATOM1} \\
\texttt{instruction\_class} & \texttt{52:56} & \texttt{exact=S2} \\
\texttt{fau\_page\_index} & \texttt{57:58} & \texttt{uint} \\
\hyperref[enum:flow_control_m]{\texttt{flow}} & \texttt{59:62} & \texttt{flow\_control\_m} \\
\textit{Reserved} & \texttt{63} & \texttt{zero} \\
\bottomrule
\end{tabularx}
\end{center}
\textbf{Fields Explained:}
\begin{tightitemize}\footnotesize
  \item \texttt{src0}: Primary source operand (\texttt{r0..r63}, \texttt{u0..u63} uniform, \texttt{k*} FAU constant, or \texttt{t0..t2} temporary; \texttt{\textasciicircum{}} marks discard/last-use). \emph{Example:} {\scriptsize\ttfamily \mnem{FADD.f32} \dst{r0}, \src{r1}, \src{r2}} (\textbf{SPIR-V:} \texttt{OpFAdd (Operand 0)}; \textbf{GLSL:} \texttt{float c = a + b;})
  \item \texttt{offset}: Signed byte offset added to base address register in memory load/store/atomic instructions. \emph{Example:} {\scriptsize\ttfamily \mnem{LOAD.i32} \dst{@r0}, \src{[r2:r3]}, \imm{\#16}} (\textbf{SPIR-V:} \texttt{OpAccessChain + OpLoad}; \textbf{GLSL:} \texttt{uint y = ssbo.arr[4];})
  \item \hyperref[enum:atom1_operation_m]{\texttt{atom1\_opc}}: Selects unary atomic operation (\texttt{ainc}, \texttt{adec}, \texttt{aumax1}, \texttt{asmax1}, \texttt{aor1}) with implicit operand 1. \emph{Example:} {\scriptsize\ttfamily \mnem{ATOM1.i32}\imod{.ainc} \src{r2}, \imm{\#0}} (\textbf{SPIR-V:} \texttt{OpAtomicIIncrement}; \textbf{GLSL:} \texttt{atomicAdd(ssbo.cnt, 1u);})
  \item \texttt{compare}: Enables depth comparison predicate in texture/shadow lookup. \emph{Example:} {\scriptsize\ttfamily \mnem{TEX\_SINGLE}\imod{.tex2d.compare} \dst{@r0}, \src{@r4:r6}} (\textbf{SPIR-V:} \texttt{OpImageSampleDrefImplicitLod}; \textbf{GLSL:} \texttt{float s = texture(shadow, vec3(uv, z));})
  \item \texttt{memory\_width}: Memory transfer width (\texttt{i8}, \texttt{i16}, \texttt{i24}, \texttt{i32}, \texttt{i48}, \texttt{i64}, \texttt{i96}, \texttt{i128}). \emph{Example:} {\scriptsize\ttfamily \mnem{LOAD.i64} \dst{@r0:r1}, \src{[r2:r3]}, \imm{\#0}} (\textbf{SPIR-V:} \texttt{OpLoad (uvec2)}; \textbf{GLSL:} \texttt{uvec2 v = ssbo.pairs[i];})
  \item \hyperref[enum:message_slot_index_m]{\texttt{message\_slot\_index}}: Asynchronous scoreboard slot (\texttt{slot0}--\texttt{slot2}, \texttt{slot7}) tracking completion of this message. \emph{Example:} {\scriptsize\ttfamily \mnem{LOAD.i32}\imod{.slot0} \dst{@r0}, \src{r2}, \imm{\#0}} (\textbf{SPIR-V:} \texttt{OpLoad}; \textbf{GLSL:} \texttt{uint v = ssbo.arr[i];})
  \item \texttt{sr\_count}: Number of consecutive 32-bit staging registers transferred by a load/store/texture message. \emph{Example:} {\scriptsize\ttfamily \mnem{LD\_MULTI.i128} \dst{@r0:r3}, \src{r4}, \imm{\#0}} (\textbf{SPIR-V:} \texttt{OpLoad (wide vector)}; \textbf{GLSL:} \texttt{uvec4 x = buf[i];})
  \item \texttt{staging\_register\_index}: Base GPR index of the contiguous staging register window (\texttt{@rN:rM}) for message payload/return. \emph{Example:} {\scriptsize\ttfamily \mnem{STORE.i64}\imod{.slot7} \src{@r0:r1}, \src{r2}, \imm{\#8}} (\textbf{SPIR-V:} \texttt{OpStore (uvec2)}; \textbf{GLSL:} \texttt{ssbo.pairs[i] = uvec2(a, b);})
  \item \texttt{staging\_register\_control\_32}: Controls the read/write staging register bank and count for async message instructions. \emph{Example:} {\scriptsize\ttfamily \mnem{LOAD.i128}\imod{.slot0} \dst{@r0:r3}, \src{r4}, \imm{\#0}} (\textbf{SPIR-V:} \texttt{OpLoad (vec4)}; \textbf{GLSL:} \texttt{vec4 v = ssbo.data[idx];})
  \item \texttt{opcode1}: Primary sub-opcode within the instruction class uniquely identifying the operation family. \emph{Example:} {\scriptsize\ttfamily \mnem{IADD.s32} \dst{r0}, \src{r1}, \src{r2}} (\textbf{SPIR-V:} \texttt{OpIAdd (operation selector)}; \textbf{GLSL:} \texttt{// Primary opcode bits})
  \item \texttt{instruction\_class}: 4-bit major instruction class selector (\texttt{bits 48:51}) partitioning the 64-bit Valhall opcode space (e.g., \texttt{FMA}, \texttt{CVT}, \texttt{SFU}, \texttt{LS}, \texttt{TEX}, \texttt{BRANCH}). \emph{Example:} {\scriptsize\ttfamily \mnem{FADD.f32} \dst{r0}, \src{r1}, \src{r2}} (\textbf{SPIR-V:} \texttt{N/A (encoding class)}; \textbf{GLSL:} \texttt{// Hardware execution pipe / decode class})
  \item \texttt{fau\_page\_index}: Selects the FAU (Fast Access Uniform) page (\texttt{0..3}) when a source operand references hardware small constants, thread IDs, or lane constants. \emph{Example:} {\scriptsize\ttfamily \mnem{FADD.f32} \dst{r0}, \src{r1}, \src{k1f}} (\textbf{SPIR-V:} \texttt{OpConstant / BuiltIn}; \textbf{GLSL:} \texttt{float y = x + 1.0;})
  \item \hyperref[enum:flow_control_m]{\texttt{flow}}: 4-bit embedded flow-control \& scoreboard wait modifier (\texttt{bits 59:62}: \texttt{none}, \texttt{wait0..wait0126}, \texttt{wait}, \texttt{reconverge}, \texttt{discard}, \texttt{end}). \emph{Example:} {\scriptsize\ttfamily \mnem{FADD.f32}\imod{.wait0.reconverge} \dst{r0}, \src{r1}, \src{r2}} (\textbf{SPIR-V:} \texttt{OpControlBarrier / OpReturn}; \textbf{GLSL:} \texttt{// Scoreboard wait + warp reconvergence})
\end{tightitemize}
\clearpage

\subsection*{ATOM1\_RETURN}
\addcontentsline{toc}{subsection}{ATOM1\_RETURN}
Executes a unary atomic read-modify-write (\texttt{i32} or \texttt{i64}, selected by \texttt{atom1\_opc}: \texttt{ainc}, \texttt{adec}, \texttt{asmax1}, \texttt{aumax1}, \texttt{aor1}) at GPU VA \texttt{src0 + \#offset} and returns the pre-modification memory value in \texttt{@sw}.
\par\textbf{Notes:}
\begin{tightitemize}
  \item Returns old value in \texttt{sr\_dst}.
  \item \textbf{Source references:} \texttt{kraid/ops.rs:281-313}, \texttt{encode\_v9.rs:725-735}.
\end{tightitemize}
\codebox{\mnem{ATOM1\_RETURN}\imod{\{.message\_slot\_index\}}\imod{\{.flow\}}\imod{\{.atom1\_opc\}} \dst{@sw}, \src{src0}, \imm{\#offset}}
\textbf{Operands:}
\begin{tightitemize}\footnotesize
  \item \texttt{\dst{@sw}} (\texttt{staging\_register\_index} + \texttt{sr\_count}): Destination staging register(s) receiving the pre-modification memory value.
  \item \texttt{\src{src0}}: 64-bit base GPU virtual address.
  \item \texttt{\imm{\#offset}} (\texttt{offset}): Unsigned byte offset added to \texttt{src0}.
\end{tightitemize}
\par\textbf{Closest SPIR-V Op:} \texttt{OpAtomicIIncrement / IDecrement}
\par\textbf{Pseudo-GLSL (Operational Semantics):}
\codebox{@sw = atomicUnaryReturn(atom1\_opc, mem[src0 + \#offset]); // returns old value}
\textbf{Execution Context:}
\begin{tightitemize}\footnotesize
  \item \textbf{Direct 64-bit GPU VA memory/atomics:} Base addresses typically come from \hyperref[sec:csf-tsd-fau]{\texttt{FAU\_k}} (\texttt{u0..u31} populated by a Pilot shader chasing \texttt{SRT\_0} $\to$ \texttt{Resource[]} $\to$ \texttt{Buffer[]}), \texttt{FauSpecial::tls\_ptr}/\texttt{wls\_ptr} backed by \hyperref[sec:csf-tsd-fau]{\texttt{TSD\_k}} (\texttt{r24..r31}), or \texttt{FauSpecial::pc} for embedded constant LUTs.
  \item \textbf{Inline FAU \& Special Operands (\hyperref[sec:csf-tsd-fau]{\texttt{FAU\_k}}, \texttt{r8..r15}):} Any 8-bit source with \texttt{mode=2} reads zero-latency uniform \texttt{u0..u31} from \texttt{FAU\_k} (\texttt{slots64 = bits[63:56]}); \texttt{mode=3} reads \texttt{SmallConstant} (Table~1.2) or \texttt{FauSpecial} (\texttt{pc}, \texttt{lane\_id}, \texttt{lane\_width}, \texttt{tls\_ptr}, \texttt{wls\_ptr}, \texttt{blend\_descriptor\_0..7}).
\end{tightitemize}
\par\textbf{Worked Examples:}
\par\textbf{Example 1 (Atomic counter allocation (`.ainc`)):}
\codebox{\mnem{ATOM1\_RETURN.i32}\imod{.ainc.slot0} \dst{@r0}, \src{[r2:r3]}, \imm{\#0}}
With \texttt{*addr = 12}: atomically sets \texttt{*addr = 13} and returns pre-increment ticket \texttt{@r0 = 12}.
\par\textbf{Example 2 (Atomic reference-count decrement (`.adec`)):}
\codebox{\mnem{ATOM1\_RETURN.i32}\imod{.adec.slot0} \dst{@r0}, \src{[r2:r3]}, \imm{\#0}}
With \texttt{*addr = 1}: atomically sets \texttt{*addr = 0} and returns old count \texttt{@r0 = 1}.

\begin{center}\footnotesize
\begin{tabularx}{\linewidth}{@{}lllX@{}}
\toprule
Variant & Arch & Class & Opcode Fields \\
\midrule
\texttt{(none)} & v9--v14 & \texttt{S2} & \texttt{opcode1=ATOM1} \\
\texttt{(none)} & v9--v14 & \texttt{S2} & \texttt{opcode1=ATOM1} \\
\bottomrule
\end{tabularx}
\end{center}
\begin{center}\footnotesize
\begin{tabularx}{0.88\linewidth}{@{}l l X@{}}
\toprule
Field Name & Bits & Type / Fixed Value \\
\midrule
\texttt{src0} & \texttt{0:7} & \texttt{SourceEncoding64} \\
\texttt{offset} & \texttt{8:15} & \texttt{uint} \\
\textit{Reserved} & \texttt{16:21} & \texttt{zero} \\
\hyperref[enum:atom1_operation_m]{\texttt{atom1\_opc}} & \texttt{22:24} & \texttt{atom1\_operation\_m} \\
\textit{Reserved} & \texttt{25} & \texttt{zero} \\
\texttt{compare} & \texttt{26} & \texttt{exact=no\_compare} \\
\texttt{memory\_width} & \texttt{27:29} & \texttt{exact=i32} \\
\hyperref[enum:message_slot_index_m]{\texttt{message\_slot\_index}} & \texttt{30:32} & \texttt{message\_slot\_index\_m} \\
\texttt{sr\_count} & \texttt{33:35} & \texttt{exact=count1} \\
\textit{Reserved} & \texttt{36:39} & \texttt{zero} \\
\texttt{staging\_register\_index} & \texttt{40:45} & \texttt{uint} \\
\texttt{staging\_register\_control\_32} & \texttt{46:47} & \texttt{exact=write} \\
\texttt{opcode1} & \texttt{48:51} & \texttt{exact=ATOM1} \\
\texttt{instruction\_class} & \texttt{52:56} & \texttt{exact=S2} \\
\texttt{fau\_page\_index} & \texttt{57:58} & \texttt{uint} \\
\hyperref[enum:flow_control_m]{\texttt{flow}} & \texttt{59:62} & \texttt{flow\_control\_m} \\
\textit{Reserved} & \texttt{63} & \texttt{zero} \\
\bottomrule
\end{tabularx}
\end{center}
\textbf{Fields Explained:}
\begin{tightitemize}\footnotesize
  \item \texttt{src0}: Primary source operand (\texttt{r0..r63}, \texttt{u0..u63} uniform, \texttt{k*} FAU constant, or \texttt{t0..t2} temporary; \texttt{\textasciicircum{}} marks discard/last-use). \emph{Example:} {\scriptsize\ttfamily \mnem{FADD.f32} \dst{r0}, \src{r1}, \src{r2}} (\textbf{SPIR-V:} \texttt{OpFAdd (Operand 0)}; \textbf{GLSL:} \texttt{float c = a + b;})
  \item \texttt{offset}: Signed byte offset added to base address register in memory load/store/atomic instructions. \emph{Example:} {\scriptsize\ttfamily \mnem{LOAD.i32} \dst{@r0}, \src{[r2:r3]}, \imm{\#16}} (\textbf{SPIR-V:} \texttt{OpAccessChain + OpLoad}; \textbf{GLSL:} \texttt{uint y = ssbo.arr[4];})
  \item \hyperref[enum:atom1_operation_m]{\texttt{atom1\_opc}}: Selects unary atomic operation (\texttt{ainc}, \texttt{adec}, \texttt{aumax1}, \texttt{asmax1}, \texttt{aor1}) with implicit operand 1. \emph{Example:} {\scriptsize\ttfamily \mnem{ATOM1.i32}\imod{.ainc} \src{r2}, \imm{\#0}} (\textbf{SPIR-V:} \texttt{OpAtomicIIncrement}; \textbf{GLSL:} \texttt{atomicAdd(ssbo.cnt, 1u);})
  \item \texttt{compare}: Enables depth comparison predicate in texture/shadow lookup. \emph{Example:} {\scriptsize\ttfamily \mnem{TEX\_SINGLE}\imod{.tex2d.compare} \dst{@r0}, \src{@r4:r6}} (\textbf{SPIR-V:} \texttt{OpImageSampleDrefImplicitLod}; \textbf{GLSL:} \texttt{float s = texture(shadow, vec3(uv, z));})
  \item \texttt{memory\_width}: Memory transfer width (\texttt{i8}, \texttt{i16}, \texttt{i24}, \texttt{i32}, \texttt{i48}, \texttt{i64}, \texttt{i96}, \texttt{i128}). \emph{Example:} {\scriptsize\ttfamily \mnem{LOAD.i64} \dst{@r0:r1}, \src{[r2:r3]}, \imm{\#0}} (\textbf{SPIR-V:} \texttt{OpLoad (uvec2)}; \textbf{GLSL:} \texttt{uvec2 v = ssbo.pairs[i];})
  \item \hyperref[enum:message_slot_index_m]{\texttt{message\_slot\_index}}: Asynchronous scoreboard slot (\texttt{slot0}--\texttt{slot2}, \texttt{slot7}) tracking completion of this message. \emph{Example:} {\scriptsize\ttfamily \mnem{LOAD.i32}\imod{.slot0} \dst{@r0}, \src{r2}, \imm{\#0}} (\textbf{SPIR-V:} \texttt{OpLoad}; \textbf{GLSL:} \texttt{uint v = ssbo.arr[i];})
  \item \texttt{sr\_count}: Number of consecutive 32-bit staging registers transferred by a load/store/texture message. \emph{Example:} {\scriptsize\ttfamily \mnem{LD\_MULTI.i128} \dst{@r0:r3}, \src{r4}, \imm{\#0}} (\textbf{SPIR-V:} \texttt{OpLoad (wide vector)}; \textbf{GLSL:} \texttt{uvec4 x = buf[i];})
  \item \texttt{staging\_register\_index}: Base GPR index of the contiguous staging register window (\texttt{@rN:rM}) for message payload/return. \emph{Example:} {\scriptsize\ttfamily \mnem{STORE.i64}\imod{.slot7} \src{@r0:r1}, \src{r2}, \imm{\#8}} (\textbf{SPIR-V:} \texttt{OpStore (uvec2)}; \textbf{GLSL:} \texttt{ssbo.pairs[i] = uvec2(a, b);})
  \item \texttt{staging\_register\_control\_32}: Controls the read/write staging register bank and count for async message instructions. \emph{Example:} {\scriptsize\ttfamily \mnem{LOAD.i128}\imod{.slot0} \dst{@r0:r3}, \src{r4}, \imm{\#0}} (\textbf{SPIR-V:} \texttt{OpLoad (vec4)}; \textbf{GLSL:} \texttt{vec4 v = ssbo.data[idx];})
  \item \texttt{opcode1}: Primary sub-opcode within the instruction class uniquely identifying the operation family. \emph{Example:} {\scriptsize\ttfamily \mnem{IADD.s32} \dst{r0}, \src{r1}, \src{r2}} (\textbf{SPIR-V:} \texttt{OpIAdd (operation selector)}; \textbf{GLSL:} \texttt{// Primary opcode bits})
  \item \texttt{instruction\_class}: 4-bit major instruction class selector (\texttt{bits 48:51}) partitioning the 64-bit Valhall opcode space (e.g., \texttt{FMA}, \texttt{CVT}, \texttt{SFU}, \texttt{LS}, \texttt{TEX}, \texttt{BRANCH}). \emph{Example:} {\scriptsize\ttfamily \mnem{FADD.f32} \dst{r0}, \src{r1}, \src{r2}} (\textbf{SPIR-V:} \texttt{N/A (encoding class)}; \textbf{GLSL:} \texttt{// Hardware execution pipe / decode class})
  \item \texttt{fau\_page\_index}: Selects the FAU (Fast Access Uniform) page (\texttt{0..3}) when a source operand references hardware small constants, thread IDs, or lane constants. \emph{Example:} {\scriptsize\ttfamily \mnem{FADD.f32} \dst{r0}, \src{r1}, \src{k1f}} (\textbf{SPIR-V:} \texttt{OpConstant / BuiltIn}; \textbf{GLSL:} \texttt{float y = x + 1.0;})
  \item \hyperref[enum:flow_control_m]{\texttt{flow}}: 4-bit embedded flow-control \& scoreboard wait modifier (\texttt{bits 59:62}: \texttt{none}, \texttt{wait0..wait0126}, \texttt{wait}, \texttt{reconverge}, \texttt{discard}, \texttt{end}). \emph{Example:} {\scriptsize\ttfamily \mnem{FADD.f32}\imod{.wait0.reconverge} \dst{r0}, \src{r1}, \src{r2}} (\textbf{SPIR-V:} \texttt{OpControlBarrier / OpReturn}; \textbf{GLSL:} \texttt{// Scoreboard wait + warp reconvergence})
\end{tightitemize}
\clearpage

\subsection*{ATOM\_RETURN}
\addcontentsline{toc}{subsection}{ATOM\_RETURN}
Atomic RMW (\texttt{i32}/\texttt{i64}) at \texttt{src0 + offset} with operand \texttt{sr\_src} and op \texttt{atom\_opc}; returns pre-modification memory value in \texttt{sr\_dst}.
\par\textbf{Notes:}
\begin{tightitemize}
  \item Returns old value in \texttt{sr\_dst}.
  \item \textbf{Source references:} \texttt{kraid/ops.rs:242-276}, \texttt{encode\_v9.rs:689-698}.
\end{tightitemize}
\codebox{\mnem{ATOM\_RETURN}\imod{\{.message\_slot\_index\}}\imod{\{.flow\}}\imod{\{.atom\_opc\}} \dst{@sw}, \src{@sr}, \src{src0}, \imm{\#offset}}
\textbf{Operands:}
\begin{tightitemize}\footnotesize
  \item \texttt{\dst{@sw}} (\texttt{staging\_register\_write\_index} + \texttt{sr\_write\_count}): Destination staging register(s) receiving the pre-modification memory value.
  \item \texttt{\src{@sr}} (\texttt{staging\_register\_index} + \texttt{sr\_count}): Source staging register(s) holding the atomic operand value.
  \item \texttt{\src{src0}}: 64-bit base GPU virtual address.
  \item \texttt{\imm{\#offset}} (\texttt{offset}): Unsigned byte offset added to \texttt{src0}.
\end{tightitemize}
\par\textbf{Closest SPIR-V Op:} \texttt{OpAtomicIAdd / SMin / UMax / And / Or / Xor}
\par\textbf{Pseudo-GLSL (Operational Semantics):}
\codebox{@sw = atomicOpReturn(atom\_opc, mem[src0 + \#offset], @sr); // returns old value}
\textbf{Execution Context:}
\begin{tightitemize}\footnotesize
  \item \textbf{Direct 64-bit GPU VA memory/atomics:} Base addresses typically come from \hyperref[sec:csf-tsd-fau]{\texttt{FAU\_k}} (\texttt{u0..u31} populated by a Pilot shader chasing \texttt{SRT\_0} $\to$ \texttt{Resource[]} $\to$ \texttt{Buffer[]}), \texttt{FauSpecial::tls\_ptr}/\texttt{wls\_ptr} backed by \hyperref[sec:csf-tsd-fau]{\texttt{TSD\_k}} (\texttt{r24..r31}), or \texttt{FauSpecial::pc} for embedded constant LUTs.
  \item \textbf{Inline FAU \& Special Operands (\hyperref[sec:csf-tsd-fau]{\texttt{FAU\_k}}, \texttt{r8..r15}):} Any 8-bit source with \texttt{mode=2} reads zero-latency uniform \texttt{u0..u31} from \texttt{FAU\_k} (\texttt{slots64 = bits[63:56]}); \texttt{mode=3} reads \texttt{SmallConstant} (Table~1.2) or \texttt{FauSpecial} (\texttt{pc}, \texttt{lane\_id}, \texttt{lane\_width}, \texttt{tls\_ptr}, \texttt{wls\_ptr}, \texttt{blend\_descriptor\_0..7}).
\end{tightitemize}
\par\textbf{Worked Examples:}
\par\textbf{Example 1 (Atomic unsigned maximum (`.aumax`)):}
\codebox{\mnem{ATOM\_RETURN.i32}\imod{.aumax.slot0} \dst{@r0}, \src{[r2:r3]}, \imm{\#0}}
With \texttt{*addr = 10}, \texttt{r0 = 25}: updates \texttt{*addr = 25} and returns previous value \texttt{@r0 = 10}.
\par\textbf{Example 2 (Atomic signed minimum (`.asmin`)):}
\codebox{\mnem{ATOM\_RETURN.i32}\imod{.asmin.slot0} \dst{@r0}, \src{[r2:r3]}, \imm{\#0x4}}
With \texttt{*addr = 5}, \texttt{r0 = -3}: updates \texttt{*addr = -3} and returns previous value \texttt{@r0 = 5}.

\begin{center}\footnotesize
\begin{tabularx}{\linewidth}{@{}lllX@{}}
\toprule
Variant & Arch & Class & Opcode Fields \\
\midrule
\texttt{(none)} & v9--v14 & \texttt{S2RW} & \texttt{opcode1=ATOM\_RETURN} \\
\texttt{(none)} & v9--v14 & \texttt{S2RW} & \texttt{opcode1=ATOM\_RETURN} \\
\bottomrule
\end{tabularx}
\end{center}
\begin{center}\footnotesize
\begin{tabularx}{0.88\linewidth}{@{}l l X@{}}
\toprule
Field Name & Bits & Type / Fixed Value \\
\midrule
\texttt{src0} & \texttt{0:7} & \texttt{SourceEncoding64} \\
\texttt{offset} & \texttt{8:15} & \texttt{uint} \\
\texttt{staging\_register\_write\_index} & \texttt{16:21} & \texttt{uint} \\
\hyperref[enum:atom_operation_m]{\texttt{atom\_opc}} & \texttt{22:25} & \texttt{atom\_operation\_m} \\
\texttt{compare} & \texttt{26} & \texttt{exact=no\_compare} \\
\texttt{memory\_width} & \texttt{27:29} & \texttt{exact=i32} \\
\hyperref[enum:message_slot_index_m]{\texttt{message\_slot\_index}} & \texttt{30:32} & \texttt{message\_slot\_index\_m} \\
\texttt{sr\_count} & \texttt{33:35} & \texttt{exact=count1} \\
\texttt{sr\_write\_count} & \texttt{36:38} & \texttt{exact=write1} \\
\textit{Reserved} & \texttt{39} & \texttt{zero} \\
\texttt{staging\_register\_index} & \texttt{40:45} & \texttt{uint} \\
\texttt{staging\_register\_control\_32} & \texttt{46:47} & \texttt{exact=read\_write} \\
\texttt{opcode1} & \texttt{48:51} & \texttt{exact=ATOM\_RETURN} \\
\texttt{instruction\_class} & \texttt{52:56} & \texttt{exact=S2RW} \\
\texttt{fau\_page\_index} & \texttt{57:58} & \texttt{uint} \\
\hyperref[enum:flow_control_m]{\texttt{flow}} & \texttt{59:62} & \texttt{flow\_control\_m} \\
\textit{Reserved} & \texttt{63} & \texttt{zero} \\
\bottomrule
\end{tabularx}
\end{center}
\textbf{Fields Explained:}
\begin{tightitemize}\footnotesize
  \item \texttt{src0}: Primary source operand (\texttt{r0..r63}, \texttt{u0..u63} uniform, \texttt{k*} FAU constant, or \texttt{t0..t2} temporary; \texttt{\textasciicircum{}} marks discard/last-use). \emph{Example:} {\scriptsize\ttfamily \mnem{FADD.f32} \dst{r0}, \src{r1}, \src{r2}} (\textbf{SPIR-V:} \texttt{OpFAdd (Operand 0)}; \textbf{GLSL:} \texttt{float c = a + b;})
  \item \texttt{offset}: Signed byte offset added to base address register in memory load/store/atomic instructions. \emph{Example:} {\scriptsize\ttfamily \mnem{LOAD.i32} \dst{@r0}, \src{[r2:r3]}, \imm{\#16}} (\textbf{SPIR-V:} \texttt{OpAccessChain + OpLoad}; \textbf{GLSL:} \texttt{uint y = ssbo.arr[4];})
  \item \texttt{staging\_register\_write\_index}: Separate destination staging register window base when read and write staging windows differ. \emph{Example:} {\scriptsize\ttfamily \mnem{ATOM\_RETURN.i32}\imod{.aadd} \dst{@r0}, \src{@r1}, \src{r2}, \imm{\#0}} (\textbf{SPIR-V:} \texttt{OpAtomicIAdd}; \textbf{GLSL:} \texttt{uint old = atomicAdd(ssbo.cnt, 1u);})
  \item \hyperref[enum:atom_operation_m]{\texttt{atom\_opc}}: Selects the atomic read-modify-write ALU operation (\texttt{aadd}, \texttt{asmin}, \texttt{aumax}, \texttt{aand}, \texttt{aor}, \texttt{axor}, \texttt{axchg}). \emph{Example:} {\scriptsize\ttfamily \mnem{ATOM.i32}\imod{.aumax} \src{@r0}, \src{r2}, \imm{\#0}} (\textbf{SPIR-V:} \texttt{OpAtomicUMax}; \textbf{GLSL:} \texttt{atomicMax(ssbo.val, x);})
  \item \texttt{compare}: Enables depth comparison predicate in texture/shadow lookup. \emph{Example:} {\scriptsize\ttfamily \mnem{TEX\_SINGLE}\imod{.tex2d.compare} \dst{@r0}, \src{@r4:r6}} (\textbf{SPIR-V:} \texttt{OpImageSampleDrefImplicitLod}; \textbf{GLSL:} \texttt{float s = texture(shadow, vec3(uv, z));})
  \item \texttt{memory\_width}: Memory transfer width (\texttt{i8}, \texttt{i16}, \texttt{i24}, \texttt{i32}, \texttt{i48}, \texttt{i64}, \texttt{i96}, \texttt{i128}). \emph{Example:} {\scriptsize\ttfamily \mnem{LOAD.i64} \dst{@r0:r1}, \src{[r2:r3]}, \imm{\#0}} (\textbf{SPIR-V:} \texttt{OpLoad (uvec2)}; \textbf{GLSL:} \texttt{uvec2 v = ssbo.pairs[i];})
  \item \hyperref[enum:message_slot_index_m]{\texttt{message\_slot\_index}}: Asynchronous scoreboard slot (\texttt{slot0}--\texttt{slot2}, \texttt{slot7}) tracking completion of this message. \emph{Example:} {\scriptsize\ttfamily \mnem{LOAD.i32}\imod{.slot0} \dst{@r0}, \src{r2}, \imm{\#0}} (\textbf{SPIR-V:} \texttt{OpLoad}; \textbf{GLSL:} \texttt{uint v = ssbo.arr[i];})
  \item \texttt{sr\_count}: Number of consecutive 32-bit staging registers transferred by a load/store/texture message. \emph{Example:} {\scriptsize\ttfamily \mnem{LD\_MULTI.i128} \dst{@r0:r3}, \src{r4}, \imm{\#0}} (\textbf{SPIR-V:} \texttt{OpLoad (wide vector)}; \textbf{GLSL:} \texttt{uvec4 x = buf[i];})
  \item \texttt{sr\_write\_count}: Number of consecutive 32-bit staging registers written back by a dual-staging message (e.g. ATOM\textbackslash\{\}\_RETURN / TEX). \emph{Example:} {\scriptsize\ttfamily \mnem{TEX\_SINGLE}\imod{.tex2d.f32.rgba} \dst{@r0:r3}, \src{@r4:r5}} (\textbf{SPIR-V:} \texttt{OpImageSampleImplicitLod}; \textbf{GLSL:} \texttt{vec4 c = texture(tex, uv);})
  \item \texttt{staging\_register\_index}: Base GPR index of the contiguous staging register window (\texttt{@rN:rM}) for message payload/return. \emph{Example:} {\scriptsize\ttfamily \mnem{STORE.i64}\imod{.slot7} \src{@r0:r1}, \src{r2}, \imm{\#8}} (\textbf{SPIR-V:} \texttt{OpStore (uvec2)}; \textbf{GLSL:} \texttt{ssbo.pairs[i] = uvec2(a, b);})
  \item \texttt{staging\_register\_control\_32}: Controls the read/write staging register bank and count for async message instructions. \emph{Example:} {\scriptsize\ttfamily \mnem{LOAD.i128}\imod{.slot0} \dst{@r0:r3}, \src{r4}, \imm{\#0}} (\textbf{SPIR-V:} \texttt{OpLoad (vec4)}; \textbf{GLSL:} \texttt{vec4 v = ssbo.data[idx];})
  \item \texttt{opcode1}: Primary sub-opcode within the instruction class uniquely identifying the operation family. \emph{Example:} {\scriptsize\ttfamily \mnem{IADD.s32} \dst{r0}, \src{r1}, \src{r2}} (\textbf{SPIR-V:} \texttt{OpIAdd (operation selector)}; \textbf{GLSL:} \texttt{// Primary opcode bits})
  \item \texttt{instruction\_class}: 4-bit major instruction class selector (\texttt{bits 48:51}) partitioning the 64-bit Valhall opcode space (e.g., \texttt{FMA}, \texttt{CVT}, \texttt{SFU}, \texttt{LS}, \texttt{TEX}, \texttt{BRANCH}). \emph{Example:} {\scriptsize\ttfamily \mnem{FADD.f32} \dst{r0}, \src{r1}, \src{r2}} (\textbf{SPIR-V:} \texttt{N/A (encoding class)}; \textbf{GLSL:} \texttt{// Hardware execution pipe / decode class})
  \item \texttt{fau\_page\_index}: Selects the FAU (Fast Access Uniform) page (\texttt{0..3}) when a source operand references hardware small constants, thread IDs, or lane constants. \emph{Example:} {\scriptsize\ttfamily \mnem{FADD.f32} \dst{r0}, \src{r1}, \src{k1f}} (\textbf{SPIR-V:} \texttt{OpConstant / BuiltIn}; \textbf{GLSL:} \texttt{float y = x + 1.0;})
  \item \hyperref[enum:flow_control_m]{\texttt{flow}}: 4-bit embedded flow-control \& scoreboard wait modifier (\texttt{bits 59:62}: \texttt{none}, \texttt{wait0..wait0126}, \texttt{wait}, \texttt{reconverge}, \texttt{discard}, \texttt{end}). \emph{Example:} {\scriptsize\ttfamily \mnem{FADD.f32}\imod{.wait0.reconverge} \dst{r0}, \src{r1}, \src{r2}} (\textbf{SPIR-V:} \texttt{OpControlBarrier / OpReturn}; \textbf{GLSL:} \texttt{// Scoreboard wait + warp reconvergence})
\end{tightitemize}
\clearpage

\subsection*{AXCHG}
\addcontentsline{toc}{subsection}{AXCHG}
Atomically swaps the 32-bit or 64-bit memory location at GPU VA \texttt{src0 + \#offset} with the replacement value in \texttt{@sr} and returns the previous memory value in \texttt{@sw}.
\par\textbf{Notes:}
\begin{tightitemize}
  \item Encoded from \texttt{OpAtom \{ atom\_op: AtomOp::Xchg \}} (\texttt{encode\_v9.rs:679-688}).
  \item \textbf{Source references:} \texttt{kraid/ops.rs:186, 242-276}, \texttt{encode\_v9.rs:679-688}.
\end{tightitemize}
\codebox{\mnem{AXCHG}\imod{\{.message\_slot\_index\}}\imod{\{.flow\}} \dst{@sw}, \src{@sr}, \src{src0}, \imm{\#offset}}
\textbf{Operands:}
\begin{tightitemize}\footnotesize
  \item \texttt{\dst{@sw}} (\texttt{staging\_register\_write\_index} + \texttt{sr\_write\_count}): Destination staging register(s) receiving the previous memory value.
  \item \texttt{\src{@sr}} (\texttt{staging\_register\_index} + \texttt{sr\_count}): Source staging register(s) holding the replacement value to swap into memory.
  \item \texttt{\src{src0}}: 64-bit base GPU virtual address.
  \item \texttt{\imm{\#offset}} (\texttt{offset}): Unsigned byte offset added to \texttt{src0}.
\end{tightitemize}
\par\textbf{Closest SPIR-V Op:} \texttt{OpAtomicExchange}
\par\textbf{Pseudo-GLSL (Operational Semantics):}
\codebox{@sw = atomicExchange(mem[src0 + \#offset], @sr); // swaps \& returns old value}
\textbf{Execution Context:}
\begin{tightitemize}\footnotesize
  \item \textbf{Direct 64-bit GPU VA memory/atomics:} Base addresses typically come from \hyperref[sec:csf-tsd-fau]{\texttt{FAU\_k}} (\texttt{u0..u31} populated by a Pilot shader chasing \texttt{SRT\_0} $\to$ \texttt{Resource[]} $\to$ \texttt{Buffer[]}), \texttt{FauSpecial::tls\_ptr}/\texttt{wls\_ptr} backed by \hyperref[sec:csf-tsd-fau]{\texttt{TSD\_k}} (\texttt{r24..r31}), or \texttt{FauSpecial::pc} for embedded constant LUTs.
  \item \textbf{Inline FAU \& Special Operands (\hyperref[sec:csf-tsd-fau]{\texttt{FAU\_k}}, \texttt{r8..r15}):} Any 8-bit source with \texttt{mode=2} reads zero-latency uniform \texttt{u0..u31} from \texttt{FAU\_k} (\texttt{slots64 = bits[63:56]}); \texttt{mode=3} reads \texttt{SmallConstant} (Table~1.2) or \texttt{FauSpecial} (\texttt{pc}, \texttt{lane\_id}, \texttt{lane\_width}, \texttt{tls\_ptr}, \texttt{wls\_ptr}, \texttt{blend\_descriptor\_0..7}).
\end{tightitemize}
\par\textbf{Worked Examples:}
\par\textbf{Example 1 (32-bit atomic swap):}
\codebox{\mnem{AXCHG.i32}\imod{.slot0} \dst{@r0}, \src{[r2:r3]}, \imm{\#0}}
With \texttt{*addr = 0x1111}, \texttt{r0 = 0x2222}: stores \texttt{0x2222} to \texttt{*addr} and returns \texttt{@r0 = 0x1111}.
\par\textbf{Example 2 (64-bit spinlock/pointer exchange):}
\codebox{\mnem{AXCHG.i64}\imod{.slot0} \dst{@r0:r1}, \src{[r2:r3]}, \imm{\#0}}
Atomically swaps 64-bit value at \texttt{[r2:r3]} with \texttt{@r0:r1}, returning the prior 64-bit word in \texttt{@r0:r1}.

\begin{center}\footnotesize
\begin{tabularx}{\linewidth}{@{}lllX@{}}
\toprule
Variant & Arch & Class & Opcode Fields \\
\midrule
\texttt{(none)} & v9--v14 & \texttt{S2RW} & \texttt{opcode1=ATOM\_RETURN} \\
\texttt{(none)} & v9--v14 & \texttt{S2RW} & \texttt{opcode1=ATOM\_RETURN} \\
\bottomrule
\end{tabularx}
\end{center}
\begin{center}\footnotesize
\begin{tabularx}{0.88\linewidth}{@{}l l X@{}}
\toprule
Field Name & Bits & Type / Fixed Value \\
\midrule
\texttt{src0} & \texttt{0:7} & \texttt{SourceEncoding64} \\
\texttt{offset} & \texttt{8:15} & \texttt{uint} \\
\texttt{staging\_register\_write\_index} & \texttt{16:21} & \texttt{uint} \\
\hyperref[enum:atom_operation_m]{\texttt{atom\_opc}} & \texttt{22:25} & \texttt{exact=axchg} \\
\texttt{compare} & \texttt{26} & \texttt{exact=no\_compare} \\
\texttt{memory\_width} & \texttt{27:29} & \texttt{exact=i32} \\
\hyperref[enum:message_slot_index_m]{\texttt{message\_slot\_index}} & \texttt{30:32} & \texttt{message\_slot\_index\_m} \\
\texttt{sr\_count} & \texttt{33:35} & \texttt{exact=count1} \\
\texttt{sr\_write\_count} & \texttt{36:38} & \texttt{exact=write1} \\
\textit{Reserved} & \texttt{39} & \texttt{zero} \\
\texttt{staging\_register\_index} & \texttt{40:45} & \texttt{uint} \\
\texttt{staging\_register\_control\_32} & \texttt{46:47} & \texttt{exact=read\_write} \\
\texttt{opcode1} & \texttt{48:51} & \texttt{exact=ATOM\_RETURN} \\
\texttt{instruction\_class} & \texttt{52:56} & \texttt{exact=S2RW} \\
\texttt{fau\_page\_index} & \texttt{57:58} & \texttt{uint} \\
\hyperref[enum:flow_control_m]{\texttt{flow}} & \texttt{59:62} & \texttt{flow\_control\_m} \\
\textit{Reserved} & \texttt{63} & \texttt{zero} \\
\bottomrule
\end{tabularx}
\end{center}
\textbf{Fields Explained:}
\begin{tightitemize}\footnotesize
  \item \texttt{src0}: Primary source operand (\texttt{r0..r63}, \texttt{u0..u63} uniform, \texttt{k*} FAU constant, or \texttt{t0..t2} temporary; \texttt{\textasciicircum{}} marks discard/last-use). \emph{Example:} {\scriptsize\ttfamily \mnem{FADD.f32} \dst{r0}, \src{r1}, \src{r2}} (\textbf{SPIR-V:} \texttt{OpFAdd (Operand 0)}; \textbf{GLSL:} \texttt{float c = a + b;})
  \item \texttt{offset}: Signed byte offset added to base address register in memory load/store/atomic instructions. \emph{Example:} {\scriptsize\ttfamily \mnem{LOAD.i32} \dst{@r0}, \src{[r2:r3]}, \imm{\#16}} (\textbf{SPIR-V:} \texttt{OpAccessChain + OpLoad}; \textbf{GLSL:} \texttt{uint y = ssbo.arr[4];})
  \item \texttt{staging\_register\_write\_index}: Separate destination staging register window base when read and write staging windows differ. \emph{Example:} {\scriptsize\ttfamily \mnem{ATOM\_RETURN.i32}\imod{.aadd} \dst{@r0}, \src{@r1}, \src{r2}, \imm{\#0}} (\textbf{SPIR-V:} \texttt{OpAtomicIAdd}; \textbf{GLSL:} \texttt{uint old = atomicAdd(ssbo.cnt, 1u);})
  \item \hyperref[enum:atom_operation_m]{\texttt{atom\_opc}}: Selects the atomic read-modify-write ALU operation (\texttt{aadd}, \texttt{asmin}, \texttt{aumax}, \texttt{aand}, \texttt{aor}, \texttt{axor}, \texttt{axchg}). \emph{Example:} {\scriptsize\ttfamily \mnem{ATOM.i32}\imod{.aumax} \src{@r0}, \src{r2}, \imm{\#0}} (\textbf{SPIR-V:} \texttt{OpAtomicUMax}; \textbf{GLSL:} \texttt{atomicMax(ssbo.val, x);})
  \item \texttt{compare}: Enables depth comparison predicate in texture/shadow lookup. \emph{Example:} {\scriptsize\ttfamily \mnem{TEX\_SINGLE}\imod{.tex2d.compare} \dst{@r0}, \src{@r4:r6}} (\textbf{SPIR-V:} \texttt{OpImageSampleDrefImplicitLod}; \textbf{GLSL:} \texttt{float s = texture(shadow, vec3(uv, z));})
  \item \texttt{memory\_width}: Memory transfer width (\texttt{i8}, \texttt{i16}, \texttt{i24}, \texttt{i32}, \texttt{i48}, \texttt{i64}, \texttt{i96}, \texttt{i128}). \emph{Example:} {\scriptsize\ttfamily \mnem{LOAD.i64} \dst{@r0:r1}, \src{[r2:r3]}, \imm{\#0}} (\textbf{SPIR-V:} \texttt{OpLoad (uvec2)}; \textbf{GLSL:} \texttt{uvec2 v = ssbo.pairs[i];})
  \item \hyperref[enum:message_slot_index_m]{\texttt{message\_slot\_index}}: Asynchronous scoreboard slot (\texttt{slot0}--\texttt{slot2}, \texttt{slot7}) tracking completion of this message. \emph{Example:} {\scriptsize\ttfamily \mnem{LOAD.i32}\imod{.slot0} \dst{@r0}, \src{r2}, \imm{\#0}} (\textbf{SPIR-V:} \texttt{OpLoad}; \textbf{GLSL:} \texttt{uint v = ssbo.arr[i];})
  \item \texttt{sr\_count}: Number of consecutive 32-bit staging registers transferred by a load/store/texture message. \emph{Example:} {\scriptsize\ttfamily \mnem{LD\_MULTI.i128} \dst{@r0:r3}, \src{r4}, \imm{\#0}} (\textbf{SPIR-V:} \texttt{OpLoad (wide vector)}; \textbf{GLSL:} \texttt{uvec4 x = buf[i];})
  \item \texttt{sr\_write\_count}: Number of consecutive 32-bit staging registers written back by a dual-staging message (e.g. ATOM\textbackslash\{\}\_RETURN / TEX). \emph{Example:} {\scriptsize\ttfamily \mnem{TEX\_SINGLE}\imod{.tex2d.f32.rgba} \dst{@r0:r3}, \src{@r4:r5}} (\textbf{SPIR-V:} \texttt{OpImageSampleImplicitLod}; \textbf{GLSL:} \texttt{vec4 c = texture(tex, uv);})
  \item \texttt{staging\_register\_index}: Base GPR index of the contiguous staging register window (\texttt{@rN:rM}) for message payload/return. \emph{Example:} {\scriptsize\ttfamily \mnem{STORE.i64}\imod{.slot7} \src{@r0:r1}, \src{r2}, \imm{\#8}} (\textbf{SPIR-V:} \texttt{OpStore (uvec2)}; \textbf{GLSL:} \texttt{ssbo.pairs[i] = uvec2(a, b);})
  \item \texttt{staging\_register\_control\_32}: Controls the read/write staging register bank and count for async message instructions. \emph{Example:} {\scriptsize\ttfamily \mnem{LOAD.i128}\imod{.slot0} \dst{@r0:r3}, \src{r4}, \imm{\#0}} (\textbf{SPIR-V:} \texttt{OpLoad (vec4)}; \textbf{GLSL:} \texttt{vec4 v = ssbo.data[idx];})
  \item \texttt{opcode1}: Primary sub-opcode within the instruction class uniquely identifying the operation family. \emph{Example:} {\scriptsize\ttfamily \mnem{IADD.s32} \dst{r0}, \src{r1}, \src{r2}} (\textbf{SPIR-V:} \texttt{OpIAdd (operation selector)}; \textbf{GLSL:} \texttt{// Primary opcode bits})
  \item \texttt{instruction\_class}: 4-bit major instruction class selector (\texttt{bits 48:51}) partitioning the 64-bit Valhall opcode space (e.g., \texttt{FMA}, \texttt{CVT}, \texttt{SFU}, \texttt{LS}, \texttt{TEX}, \texttt{BRANCH}). \emph{Example:} {\scriptsize\ttfamily \mnem{FADD.f32} \dst{r0}, \src{r1}, \src{r2}} (\textbf{SPIR-V:} \texttt{N/A (encoding class)}; \textbf{GLSL:} \texttt{// Hardware execution pipe / decode class})
  \item \texttt{fau\_page\_index}: Selects the FAU (Fast Access Uniform) page (\texttt{0..3}) when a source operand references hardware small constants, thread IDs, or lane constants. \emph{Example:} {\scriptsize\ttfamily \mnem{FADD.f32} \dst{r0}, \src{r1}, \src{k1f}} (\textbf{SPIR-V:} \texttt{OpConstant / BuiltIn}; \textbf{GLSL:} \texttt{float y = x + 1.0;})
  \item \hyperref[enum:flow_control_m]{\texttt{flow}}: 4-bit embedded flow-control \& scoreboard wait modifier (\texttt{bits 59:62}: \texttt{none}, \texttt{wait0..wait0126}, \texttt{wait}, \texttt{reconverge}, \texttt{discard}, \texttt{end}). \emph{Example:} {\scriptsize\ttfamily \mnem{FADD.f32}\imod{.wait0.reconverge} \dst{r0}, \src{r1}, \src{r2}} (\textbf{SPIR-V:} \texttt{OpControlBarrier / OpReturn}; \textbf{GLSL:} \texttt{// Scoreboard wait + warp reconvergence})
\end{tightitemize}
\clearpage

\section{Attributes, Varyings \& Fragment (18 mnemonics)}
\subsection*{ATEST}
\addcontentsline{toc}{subsection}{ATEST}
Performs hardware alpha-to-coverage and early-Z coverage update in fragment shaders: evaluates \texttt{(coverage=src0 [r60], alpha=src1 [f32], datum=src2 [atest\_datum])} and writes the updated MSAA sample coverage mask to \texttt{@sw} (pinned \texttt{r60}). Does not implicitly discard helper lanes.
\par\textbf{Notes:}
\begin{tightitemize}
  \item Does NOT implicitly discard.
  \item Must precede \texttt{ZS\_EMIT}/\texttt{BLEND} (\texttt{ISA.xml:2}).
  \item \textbf{Source references:} \texttt{kraid/ops.rs:145-173}, \texttt{encode\_v9.rs:595-618}, \texttt{bifrost/valhall/ISA.xml:2}.
\end{tightitemize}
\codebox{\mnem{ATEST}\imod{\{.message\_slot\_index\}}\imod{\{.flow\}} \dst{@sw}, \src{src0}, \src{src1}\imod{\{.widen1\}}, \src{src2}}
\textbf{Operands:}
\begin{tightitemize}\footnotesize
  \item \texttt{\dst{@sw}} (\texttt{staging\_register\_index} + \texttt{sr\_count}): Destination staging register (pinned \texttt{r60}) receiving the updated MSAA sample coverage mask.
  \item \texttt{\src{src0}}: Incoming sample coverage mask (pinned \texttt{r60}).
  \item \texttt{\src{src1}}: Fragment alpha value (\texttt{f32}) for alpha-to-coverage.
  \item \texttt{\src{src2}}: Hardware alpha-test datum (\texttt{atest\_datum.w0} FAU special).
\end{tightitemize}
\par\textbf{Closest SPIR-V Op:} \texttt{None (alpha-to-coverage)}
\par\textbf{Pseudo-GLSL (Operational Semantics):}
\codebox{@sw = evalAlphaToCoverageAndEarlyZ(cov = src0, alpha = src1, datum = src2);}
\textbf{Execution Context:}
\begin{tightitemize}\footnotesize
  \item \textbf{Depth/Stencil \& coverage state (\hyperref[sec:csf-zsd-blend]{\texttt{IDVSSR::ZSD}}, \texttt{r52:r53}):} Interacts with 32B \texttt{Depth/stencil} (\texttt{depth\_func}, \texttt{stencil\_from\_shader}), \hyperref[sec:csf-tiler]{\texttt{DCD2}} (\texttt{r38}: \texttt{no\_shader\_depth\_read} Early-Z vs Late-Z), \texttt{DCD0}/\texttt{DCD1} (\texttt{r57:r58}: \texttt{shader\_modifies\_coverage}, \texttt{alpha\_to\_coverage}), and \texttt{FauSpecial::atest\_datum}.
  \item \textbf{Fragment stage preloads \& FP modes (\hyperref[sec:csf-spd]{\texttt{SPD\_2}}, \texttt{r20:r21}):} Hardware pre-fills \texttt{r59} with \texttt{floatBitsToUint(FragCoord.z)} and \texttt{r60} with the MSAA sample coverage mask (\texttt{gl\_SampleMaskIn[0]}), while \texttt{ShaderProgram} word~0 latches 32/64 GPR allocation and warp-wide \texttt{flush\_to\_zero\_mode}, \texttt{suppress\_na\_n}, and \texttt{suppress\_inf}.
\end{tightitemize}
\par\textbf{Worked Examples:}
\par\textbf{Example 1 (Early-Z / alpha-to-coverage update):}
\codebox{\mnem{ATEST}\imod{.slot0} \dst{@r60}, \src{r60}, \src{r3}, \src{k0}}
Combines input sample coverage \texttt{r60} with RT0 alpha \texttt{r3} and writes updated coverage mask back to \texttt{@r60}.
\par\textbf{Example 2 (Discard check after ATEST):}
\codebox{\mnem{ATEST}\imod{.slot0.wait0} \dst{@r60}, \src{r60}, \src{r3}, \src{k0}}
Does NOT implicitly discard the lane; shaders test \texttt{@r60 == 0} afterward if skipping expensive shading.
\par\textbf{Mined SPIR-V $\to$ Valhall ISA Example (Committing fragment alpha/discard coverage before termination --- Fragment shader, Mali-G710 (v10)):}
\codebox{\textbf{; SPIR-V (input)}\\OpKill ; unconditional discard fragment shader\\[3pt]\textbf{; Valhall ISA (malioc output)}\\DISCARD.f32.wait0126.ne k1n, k0\\ATEST.slot0.end @r60, r60, k0, atest\_datum.w0}
Even in a 2-instruction unconditional-discard shader, \texttt{DISCARD} must be followed by \texttt{ATEST} (updating pinned coverage register \texttt{r60} against \texttt{atest\_datum.w0}) before terminating with \texttt{.end}.
\begin{center}\footnotesize
\begin{tabularx}{\linewidth}{@{}lllX@{}}
\toprule
Variant & Arch & Class & Opcode Fields \\
\midrule
\texttt{(none)} & v9--v14 & \texttt{S3} & \texttt{opcode1=ATEST} \\
\bottomrule
\end{tabularx}
\end{center}
\begin{center}\footnotesize
\begin{tabularx}{0.88\linewidth}{@{}l l X@{}}
\toprule
Field Name & Bits & Type / Fixed Value \\
\midrule
\texttt{src0} & \texttt{0:7} & \texttt{SourceEncoding} \\
\texttt{src1} & \texttt{8:15} & \texttt{SourceEncoding} \\
\texttt{src2} & \texttt{16:23} & \texttt{SourceEncoding} \\
\textit{Reserved} & \texttt{24:25} & \texttt{zero} \\
\hyperref[enum:widen_m]{\texttt{widen1}} & \texttt{26:27} & \texttt{widen\_m (def none)} \\
\textit{Reserved} & \texttt{28:29} & \texttt{zero} \\
\hyperref[enum:message_slot_index_m]{\texttt{message\_slot\_index}} & \texttt{30:32} & \texttt{exact=slot0} \\
\texttt{sr\_count} & \texttt{33:35} & \texttt{exact=count1} \\
\textit{Reserved} & \texttt{36:39} & \texttt{zero} \\
\texttt{staging\_register\_index} & \texttt{40:45} & \texttt{uint} \\
\texttt{staging\_register\_control\_32} & \texttt{46:47} & \texttt{exact=write} \\
\texttt{opcode1} & \texttt{48:51} & \texttt{exact=ATEST} \\
\texttt{instruction\_class} & \texttt{52:56} & \texttt{exact=S3} \\
\texttt{fau\_page\_index} & \texttt{57:58} & \texttt{uint} \\
\hyperref[enum:flow_control_m]{\texttt{flow}} & \texttt{59:62} & \texttt{flow\_control\_m} \\
\textit{Reserved} & \texttt{63} & \texttt{zero} \\
\bottomrule
\end{tabularx}
\end{center}
\textbf{Fields Explained:}
\begin{tightitemize}\footnotesize
  \item \texttt{src0}: Primary source operand (\texttt{r0..r63}, \texttt{u0..u63} uniform, \texttt{k*} FAU constant, or \texttt{t0..t2} temporary; \texttt{\textasciicircum{}} marks discard/last-use). \emph{Example:} {\scriptsize\ttfamily \mnem{FADD.f32} \dst{r0}, \src{r1}, \src{r2}} (\textbf{SPIR-V:} \texttt{OpFAdd (Operand 0)}; \textbf{GLSL:} \texttt{float c = a + b;})
  \item \texttt{src1}: Secondary source operand (register, uniform, or FAU constant; in shifts, supplies the shift count byte). \emph{Example:} {\scriptsize\ttfamily \mnem{LSHIFT.i32} \dst{r0}, \src{r1}, \src{r2}\imod{.b0}} (\textbf{SPIR-V:} \texttt{OpShiftLeftLogical (Shift)}; \textbf{GLSL:} \texttt{uint y = a \textless{}\textless{} s;})
  \item \texttt{src2}: Third source operand (FMA addend, \texttt{MUX}/\texttt{CSEL} selector, \texttt{IADDC} carry-in, or bitwise mask in fused shift-logic). \emph{Example:} {\scriptsize\ttfamily \mnem{FMA.f32} \dst{r0}, \src{r1}, \src{r2}, \src{r3}} (\textbf{SPIR-V:} \texttt{GLSL.std.450 Fma (Addend)}; \textbf{GLSL:} \texttt{float d = fma(a, b, c);})
  \item \hyperref[enum:widen_m]{\texttt{widen1}}: Stage A widening extractor on \texttt{src1} (promotes narrow byte/half/word lane before Stage B). \emph{Example:} {\scriptsize\ttfamily \mnem{IMUL.u64} \dst{[r0:r1]}, \src{[r2:r3]}\imod{.w0}, \src{[r4:r5]}\imod{.w0}} (\textbf{SPIR-V:} \texttt{OpUConvert / OpUMulExtended}; \textbf{GLSL:} \texttt{uint64\_t p = uint64\_t(a) * uint64\_t(b);})
  \item \hyperref[enum:message_slot_index_m]{\texttt{message\_slot\_index}}: Asynchronous scoreboard slot (\texttt{slot0}--\texttt{slot2}, \texttt{slot7}) tracking completion of this message. \emph{Example:} {\scriptsize\ttfamily \mnem{LOAD.i32}\imod{.slot0} \dst{@r0}, \src{r2}, \imm{\#0}} (\textbf{SPIR-V:} \texttt{OpLoad}; \textbf{GLSL:} \texttt{uint v = ssbo.arr[i];})
  \item \texttt{sr\_count}: Number of consecutive 32-bit staging registers transferred by a load/store/texture message. \emph{Example:} {\scriptsize\ttfamily \mnem{LD\_MULTI.i128} \dst{@r0:r3}, \src{r4}, \imm{\#0}} (\textbf{SPIR-V:} \texttt{OpLoad (wide vector)}; \textbf{GLSL:} \texttt{uvec4 x = buf[i];})
  \item \texttt{staging\_register\_index}: Base GPR index of the contiguous staging register window (\texttt{@rN:rM}) for message payload/return. \emph{Example:} {\scriptsize\ttfamily \mnem{STORE.i64}\imod{.slot7} \src{@r0:r1}, \src{r2}, \imm{\#8}} (\textbf{SPIR-V:} \texttt{OpStore (uvec2)}; \textbf{GLSL:} \texttt{ssbo.pairs[i] = uvec2(a, b);})
  \item \texttt{staging\_register\_control\_32}: Controls the read/write staging register bank and count for async message instructions. \emph{Example:} {\scriptsize\ttfamily \mnem{LOAD.i128}\imod{.slot0} \dst{@r0:r3}, \src{r4}, \imm{\#0}} (\textbf{SPIR-V:} \texttt{OpLoad (vec4)}; \textbf{GLSL:} \texttt{vec4 v = ssbo.data[idx];})
  \item \texttt{opcode1}: Primary sub-opcode within the instruction class uniquely identifying the operation family. \emph{Example:} {\scriptsize\ttfamily \mnem{IADD.s32} \dst{r0}, \src{r1}, \src{r2}} (\textbf{SPIR-V:} \texttt{OpIAdd (operation selector)}; \textbf{GLSL:} \texttt{// Primary opcode bits})
  \item \texttt{instruction\_class}: 4-bit major instruction class selector (\texttt{bits 48:51}) partitioning the 64-bit Valhall opcode space (e.g., \texttt{FMA}, \texttt{CVT}, \texttt{SFU}, \texttt{LS}, \texttt{TEX}, \texttt{BRANCH}). \emph{Example:} {\scriptsize\ttfamily \mnem{FADD.f32} \dst{r0}, \src{r1}, \src{r2}} (\textbf{SPIR-V:} \texttt{N/A (encoding class)}; \textbf{GLSL:} \texttt{// Hardware execution pipe / decode class})
  \item \texttt{fau\_page\_index}: Selects the FAU (Fast Access Uniform) page (\texttt{0..3}) when a source operand references hardware small constants, thread IDs, or lane constants. \emph{Example:} {\scriptsize\ttfamily \mnem{FADD.f32} \dst{r0}, \src{r1}, \src{k1f}} (\textbf{SPIR-V:} \texttt{OpConstant / BuiltIn}; \textbf{GLSL:} \texttt{float y = x + 1.0;})
  \item \hyperref[enum:flow_control_m]{\texttt{flow}}: 4-bit embedded flow-control \& scoreboard wait modifier (\texttt{bits 59:62}: \texttt{none}, \texttt{wait0..wait0126}, \texttt{wait}, \texttt{reconverge}, \texttt{discard}, \texttt{end}). \emph{Example:} {\scriptsize\ttfamily \mnem{FADD.f32}\imod{.wait0.reconverge} \dst{r0}, \src{r1}, \src{r2}} (\textbf{SPIR-V:} \texttt{OpControlBarrier / OpReturn}; \textbf{GLSL:} \texttt{// Scoreboard wait + warp reconvergence})
\end{tightitemize}
\clearpage

\subsection*{BLEND}
\addcontentsline{toc}{subsection}{BLEND}
Blends color \texttt{sr\_src} (\texttt{r0..}) with render target using descriptor \texttt{src0} (64-bit) and coverage \texttt{src2} (\texttt{r60}). In fixed-function path, writes tilebuffer and either terminates (\texttt{.end}) or branches by \texttt{offset}; in blend-shader path, falls through to next instruction.
\par\textbf{Notes:}
\begin{tightitemize}
  \item Dual-path control-flow semantics (\texttt{ISA.xml:17}, \texttt{lower\_blend\_call.rs:8-58}).
  \item \textbf{Source references:} \texttt{kraid/ops.rs:361-402}, \texttt{encode\_v9.rs:748-770}, \texttt{lower\_blend\_call.rs:8-58}, \texttt{bifrost/valhall/ISA.xml:17}.
\end{tightitemize}
\codebox{\mnem{BLEND}\imod{\{.message\_slot\_index\}}\imod{\{.flow\}}\imod{\{.vecsize\}}\imod{\{.register\_format\}} \src{@sr}, \src{src0}, \imm{\#offset}, \src{src2}}
\textbf{Operands:}
\begin{tightitemize}\footnotesize
  \item \texttt{\src{@sr}} (\texttt{staging\_register\_index} + \texttt{sr\_count}): Contiguous source staging registers holding the fragment RGBA output color.
  \item \texttt{\src{src0}}: 64-bit blend descriptor (\texttt{blend\_descriptor\_0..7} FAU special).
  \item \texttt{\imm{\#offset}} (\texttt{offset}): Unsigned PC-relative byte offset to the blend shader continuation (or return offset).
  \item \texttt{\src{src2}}: Sample coverage mask (typically \texttt{r60\textasciicircum{}}).
\end{tightitemize}
\par\textbf{Closest SPIR-V Op:} \texttt{OpStore (Fragment color output)}
\par\textbf{Pseudo-GLSL (Operational Semantics):}
\codebox{tileSRAM[rt] = blendColor(desc = src0, rgba = @sr, dst = tileSRAM[rt], cov = src2);\\if (flow == end) return;}
\textbf{Execution Context:}
\begin{tightitemize}\footnotesize
  \item \textbf{Blend descriptor chain (\hyperref[sec:csf-zsd-blend]{\texttt{IDVSSR::BLEND\_DESC}}, \texttt{r50:r51}):} 16B-aligned VA (\texttt{\& \textasciitilde{}0xF}, \texttt{count=bits[3:0]}) $\to$ \texttt{Blend[rt]} (16B: fixed-function equation or programmable blend shader PC), selected explicitly via \texttt{src0 = FauSpecial::blend\_descriptor\_0..7}.
  \item \textbf{On-chip 16$\times$16 Tile SRAM (\hyperref[sec:csf-fbd]{\texttt{FRAGMENTSR::FBD\_POINTER}}, \texttt{r40:r41}):} 64B-aligned VA (\texttt{\& \textasciitilde{}0x3F}, \texttt{type[0]|zs\_ext[1]|rt\_cnt-1[5:2]}) $\to$ \texttt{FramebufferParameters} (128B) + optional \texttt{ZSCRCExtension} (64B, for \texttt{.z\_stencil}) + \texttt{RenderTarget[rt]} (64B: \texttt{internal\_format}, \texttt{clear\_color}, Morton 16$\times$16 \texttt{writeback\_buffer}).
  \item \textbf{Fragment stage preloads \& FP modes (\hyperref[sec:csf-spd]{\texttt{SPD\_2}}, \texttt{r20:r21}):} Hardware pre-fills \texttt{r59} with \texttt{floatBitsToUint(FragCoord.z)} and \texttt{r60} with the MSAA sample coverage mask (\texttt{gl\_SampleMaskIn[0]}), while \texttt{ShaderProgram} word~0 latches 32/64 GPR allocation and warp-wide \texttt{flush\_to\_zero\_mode}, \texttt{suppress\_na\_n}, and \texttt{suppress\_inf}.
\end{tightitemize}
\par\textbf{Worked Examples:}
\par\textbf{Example 1 (Fixed-function RT0 RGBA32F blend + shader end):}
\codebox{\mnem{BLEND}\imod{.f32.v4.slot0.end} \dst{@r0:r3}, \src{u0}, \src{r60}, \imm{\#0}}
Blends RGBA \texttt{@r0:r3} with coverage \texttt{r60} using blend descriptor \texttt{u0} for RT0; \texttt{.end} terminates the fragment warp.
\par\textbf{Example 2 (MRT secondary render target (`\#1`, no `.end`)):}
\codebox{\mnem{BLEND}\imod{.f16.v4.slot0} \dst{@r4:r5}, \src{u1}, \src{r60}, \imm{\#1}}
Emits packed FP16 RGBA \texttt{@r4:r5} to render target 1 before issuing the final RT0 \texttt{.end} blend.
\par\textbf{Mined SPIR-V $\to$ Valhall ISA Example (Writing FP16 RGBA output to the tilebuffer and ending the warp --- Fragment shader, Mali-G710 (v10)):}
\codebox{\textbf{; SPIR-V (input)}\\OpStore \%out\_color \%rgba\_const\\[3pt]\textbf{; Valhall ISA (malioc output)}\\MOV\_IMM.i32 r0, \#0x3c003c00\\MOV\_IMM.i32.wait r1.h0, \#0x3c003c00\\BLEND.slot0.end.v4.f16 @r0:r1, blend\_descriptor\_0, 0x20, r60\textasciicircum{}}
\texttt{BLEND} emits the final fragment color (\texttt{@r0:r1} as 4x FP16) to render target 0 via \texttt{blend\_descriptor\_0} masked by sample coverage \texttt{r60\textasciicircum{}}, and terminates the fragment warp with \texttt{.end}.
\begin{center}\footnotesize
\begin{tabularx}{\linewidth}{@{}lllX@{}}
\toprule
Variant & Arch & Class & Opcode Fields \\
\midrule
\texttt{(none)} & v9--v14 & \texttt{S3} & \texttt{opcode1=BLEND} \\
\texttt{(none)} & v9--v14 & \texttt{S3} & \texttt{opcode1=BLEND} \\
\bottomrule
\end{tabularx}
\end{center}
\begin{center}\footnotesize
\begin{tabularx}{0.88\linewidth}{@{}l l X@{}}
\toprule
Field Name & Bits & Type / Fixed Value \\
\midrule
\texttt{src0} & \texttt{0:7} & \texttt{SourceEncoding64} \\
\texttt{offset} & \texttt{8:15} & \texttt{pc\_rel\_label\_unsigned} \\
\texttt{src2} & \texttt{16:23} & \texttt{SourceEncoding} \\
\hyperref[enum:register_file_format_general_m]{\texttt{register\_format}} & \texttt{24:26} & \texttt{register\_file\_format\_general\_m} \\
\textit{Reserved} & \texttt{27} & \texttt{zero} \\
\hyperref[enum:vecsize_varying_m]{\texttt{vecsize}} & \texttt{28:29} & \texttt{exact=v4} \\
\hyperref[enum:message_slot_index_m]{\texttt{message\_slot\_index}} & \texttt{30:32} & \texttt{message\_slot\_index\_m} \\
\texttt{sr\_count} & \texttt{33:35} & \texttt{sr\_count\_t} \\
\textit{Reserved} & \texttt{36:39} & \texttt{zero} \\
\texttt{staging\_register\_index} & \texttt{40:45} & \texttt{uint} \\
\texttt{staging\_register\_control\_32} & \texttt{46:47} & \texttt{exact=read} \\
\texttt{opcode1} & \texttt{48:51} & \texttt{exact=BLEND} \\
\texttt{instruction\_class} & \texttt{52:56} & \texttt{exact=S3} \\
\texttt{fau\_page\_index} & \texttt{57:58} & \texttt{uint} \\
\hyperref[enum:flow_control_m]{\texttt{flow}} & \texttt{59:62} & \texttt{flow\_control\_m} \\
\textit{Reserved} & \texttt{63} & \texttt{zero} \\
\bottomrule
\end{tabularx}
\end{center}
\textbf{Fields Explained:}
\begin{tightitemize}\footnotesize
  \item \texttt{src0}: Primary source operand (\texttt{r0..r63}, \texttt{u0..u63} uniform, \texttt{k*} FAU constant, or \texttt{t0..t2} temporary; \texttt{\textasciicircum{}} marks discard/last-use). \emph{Example:} {\scriptsize\ttfamily \mnem{FADD.f32} \dst{r0}, \src{r1}, \src{r2}} (\textbf{SPIR-V:} \texttt{OpFAdd (Operand 0)}; \textbf{GLSL:} \texttt{float c = a + b;})
  \item \texttt{offset}: Signed byte offset added to base address register in memory load/store/atomic instructions. \emph{Example:} {\scriptsize\ttfamily \mnem{LOAD.i32} \dst{@r0}, \src{[r2:r3]}, \imm{\#16}} (\textbf{SPIR-V:} \texttt{OpAccessChain + OpLoad}; \textbf{GLSL:} \texttt{uint y = ssbo.arr[4];})
  \item \texttt{src2}: Third source operand (FMA addend, \texttt{MUX}/\texttt{CSEL} selector, \texttt{IADDC} carry-in, or bitwise mask in fused shift-logic). \emph{Example:} {\scriptsize\ttfamily \mnem{FMA.f32} \dst{r0}, \src{r1}, \src{r2}, \src{r3}} (\textbf{SPIR-V:} \texttt{GLSL.std.450 Fma (Addend)}; \textbf{GLSL:} \texttt{float d = fma(a, b, c);})
  \item \hyperref[enum:register_file_format_general_m]{\texttt{register\_format}}: Target GPR data conversion format (\texttt{auto32}, \texttt{f32}, \texttt{f16}, \texttt{s32}, \texttt{s16}, \texttt{u32}, \texttt{u16}). \emph{Example:} {\scriptsize\ttfamily \mnem{LD\_TILE}\imod{.f16.v4} \dst{@r0:r1}, \src{r2}} (\textbf{SPIR-V:} \texttt{OpImageRead (SubpassData)}; \textbf{GLSL:} \texttt{f16vec4 c = subpassLoad(inputAtt);})
  \item \hyperref[enum:vecsize_varying_m]{\texttt{vecsize}}: Component count (\texttt{none}=1, \texttt{v2}, \texttt{v3}, \texttt{v4}) for attribute, varying, or blend transfers. \emph{Example:} {\scriptsize\ttfamily \mnem{LD\_ATTR}\imod{.f32.v3} \dst{@r0:r2}, \src{r4}} (\textbf{SPIR-V:} \texttt{OpLoad (vec3 Input)}; \textbf{GLSL:} \texttt{in vec3 a\_normal;})
  \item \hyperref[enum:message_slot_index_m]{\texttt{message\_slot\_index}}: Asynchronous scoreboard slot (\texttt{slot0}--\texttt{slot2}, \texttt{slot7}) tracking completion of this message. \emph{Example:} {\scriptsize\ttfamily \mnem{LOAD.i32}\imod{.slot0} \dst{@r0}, \src{r2}, \imm{\#0}} (\textbf{SPIR-V:} \texttt{OpLoad}; \textbf{GLSL:} \texttt{uint v = ssbo.arr[i];})
  \item \texttt{sr\_count}: Number of consecutive 32-bit staging registers transferred by a load/store/texture message. \emph{Example:} {\scriptsize\ttfamily \mnem{LD\_MULTI.i128} \dst{@r0:r3}, \src{r4}, \imm{\#0}} (\textbf{SPIR-V:} \texttt{OpLoad (wide vector)}; \textbf{GLSL:} \texttt{uvec4 x = buf[i];})
  \item \texttt{staging\_register\_index}: Base GPR index of the contiguous staging register window (\texttt{@rN:rM}) for message payload/return. \emph{Example:} {\scriptsize\ttfamily \mnem{STORE.i64}\imod{.slot7} \src{@r0:r1}, \src{r2}, \imm{\#8}} (\textbf{SPIR-V:} \texttt{OpStore (uvec2)}; \textbf{GLSL:} \texttt{ssbo.pairs[i] = uvec2(a, b);})
  \item \texttt{staging\_register\_control\_32}: Controls the read/write staging register bank and count for async message instructions. \emph{Example:} {\scriptsize\ttfamily \mnem{LOAD.i128}\imod{.slot0} \dst{@r0:r3}, \src{r4}, \imm{\#0}} (\textbf{SPIR-V:} \texttt{OpLoad (vec4)}; \textbf{GLSL:} \texttt{vec4 v = ssbo.data[idx];})
  \item \texttt{opcode1}: Primary sub-opcode within the instruction class uniquely identifying the operation family. \emph{Example:} {\scriptsize\ttfamily \mnem{IADD.s32} \dst{r0}, \src{r1}, \src{r2}} (\textbf{SPIR-V:} \texttt{OpIAdd (operation selector)}; \textbf{GLSL:} \texttt{// Primary opcode bits})
  \item \texttt{instruction\_class}: 4-bit major instruction class selector (\texttt{bits 48:51}) partitioning the 64-bit Valhall opcode space (e.g., \texttt{FMA}, \texttt{CVT}, \texttt{SFU}, \texttt{LS}, \texttt{TEX}, \texttt{BRANCH}). \emph{Example:} {\scriptsize\ttfamily \mnem{FADD.f32} \dst{r0}, \src{r1}, \src{r2}} (\textbf{SPIR-V:} \texttt{N/A (encoding class)}; \textbf{GLSL:} \texttt{// Hardware execution pipe / decode class})
  \item \texttt{fau\_page\_index}: Selects the FAU (Fast Access Uniform) page (\texttt{0..3}) when a source operand references hardware small constants, thread IDs, or lane constants. \emph{Example:} {\scriptsize\ttfamily \mnem{FADD.f32} \dst{r0}, \src{r1}, \src{k1f}} (\textbf{SPIR-V:} \texttt{OpConstant / BuiltIn}; \textbf{GLSL:} \texttt{float y = x + 1.0;})
  \item \hyperref[enum:flow_control_m]{\texttt{flow}}: 4-bit embedded flow-control \& scoreboard wait modifier (\texttt{bits 59:62}: \texttt{none}, \texttt{wait0..wait0126}, \texttt{wait}, \texttt{reconverge}, \texttt{discard}, \texttt{end}). \emph{Example:} {\scriptsize\ttfamily \mnem{FADD.f32}\imod{.wait0.reconverge} \dst{r0}, \src{r1}, \src{r2}} (\textbf{SPIR-V:} \texttt{OpControlBarrier / OpReturn}; \textbf{GLSL:} \texttt{// Scoreboard wait + warp reconvergence})
\end{tightitemize}
\clearpage

\subsection*{LD\_ATTR}
\addcontentsline{toc}{subsection}{LD\_ATTR}
Loads \texttt{vecsize} components of \texttt{register\_format} from vertex attribute descriptor at dynamic handle \texttt{(src2)} for \texttt{(vertex\_id=src0, instance\_id=src1)}.
\par\textbf{Notes:}
\begin{tightitemize}
  \item Handle must be subgroup-uniform (\texttt{ISA.xml:66}).
  \item \textbf{Source references:} \texttt{kraid/ops.rs:2668-2694}, \texttt{bifrost/valhall/ISA.xml:66}.
\end{tightitemize}
\codebox{\mnem{LD\_ATTR}\imod{\{.message\_slot\_index\}}\imod{\{.flow\}}\imod{\{.vecsize\}}\imod{\{.register\_format\}} \dst{@sw}, \src{src0}, \src{src1}, \src{src2}}
\textbf{Operands:}
\begin{tightitemize}\footnotesize
  \item \texttt{\dst{@sw}} (\texttt{staging\_register\_index} + \texttt{sr\_count}): Destination staging register range receiving \texttt{vecsize} converted attribute components.
  \item \texttt{\src{src0}}: Vertex index (\texttt{VertexIndex}, preloaded in \texttt{r60}).
  \item \texttt{\src{src1}}: Instance index (\texttt{InstanceIndex}, preloaded in \texttt{r61}).
  \item \texttt{\src{src2}}: Dynamic attribute descriptor handle.
\end{tightitemize}
\par\textbf{Closest SPIR-V Op:} \texttt{OpLoad (dynamic vertex input)}
\par\textbf{Pseudo-GLSL (Operational Semantics):}
\codebox{@sw = fetchConvertAttr(handle = src2, vert = src0, inst = src1, register\_format, vecsize);}
\textbf{Execution Context:}
\begin{tightitemize}\footnotesize
  \item \textbf{Attribute descriptor chain (\hyperref[sec:csf-srt]{\texttt{IDVSSR::SRT\_0}}, \texttt{r0:r1}):} Masks 64B-aligned VA (\texttt{\& \textasciitilde{}0x3F}, \texttt{count=bits[5:0]}) $\to$ \texttt{Resource[table]} (16B) $\to$ \hyperref[sec:csf-buf-attr]{\texttt{Attribute[index]}} (32B: \texttt{table}, \texttt{buffer\_index}, \texttt{offset}, \texttt{stride}, \texttt{format}) $\to$ \texttt{Resource[Attribute::table]} $\to$ \texttt{Buffer[Attribute::buffer\_index]} (32B: \texttt{address}, \texttt{size}).
  \item \textbf{Vertex stage preloads (\hyperref[sec:csf-spd]{\texttt{SPD\_0}/\texttt{SPD\_1}}, \texttt{r16..r19}):} The hardware Index Fetcher reads \hyperref[sec:csf-regs]{\texttt{IDVSSR::INDEX\_BUFFER}} (\texttt{r54:r55}) and pre-fills \texttt{r59, r60} with \texttt{VertexIndex} (\texttt{+ VERTEX\_OFFSET + GLOBAL\_ATTR\_OFF}) and \texttt{r61} with \texttt{InstanceIndex} (\texttt{+ INSTANCE\_OFFSET}) so Position and Varying shaders can pull vertex attributes directly.
\end{tightitemize}
\par\textbf{Worked Examples:}
\par\textbf{Example 1 (Dynamic vertex attribute fetch (FP32 vec4)):}
\codebox{\mnem{LD\_ATTR}\imod{.f32.v4.slot0} \dst{@r0:r3}, \src{r60}, \src{r61}, \src{r4}}
Loads 4x FP32 vertex attribute for \texttt{(vertex\_id=r60, instance\_id=r61)} at dynamic descriptor handle \texttt{r4} into \texttt{@r0:r3}.
\par\textbf{Example 2 (Integer vertex attribute (u32 vec2)):}
\codebox{\mnem{LD\_ATTR}\imod{.u32.v2.slot0} \dst{@r0:r1}, \src{r60}, \src{r61}, \src{r4}}
Loads 2x unsigned 32-bit joint indices into \texttt{@r0:r1}.

\begin{center}\footnotesize
\begin{tabularx}{\linewidth}{@{}lllX@{}}
\toprule
Variant & Arch & Class & Opcode Fields \\
\midrule
\texttt{(none)} & v9--v14 & \texttt{S3} & \texttt{opcode1=LD\_ATTR} \\
\bottomrule
\end{tabularx}
\end{center}
\begin{center}\footnotesize
\begin{tabularx}{0.88\linewidth}{@{}l l X@{}}
\toprule
Field Name & Bits & Type / Fixed Value \\
\midrule
\texttt{src0} & \texttt{0:7} & \texttt{SourceEncoding} \\
\texttt{src1} & \texttt{8:15} & \texttt{SourceEncoding} \\
\texttt{src2} & \texttt{16:23} & \texttt{SourceEncoding} \\
\hyperref[enum:register_file_format_general_m]{\texttt{register\_format}} & \texttt{24:26} & \texttt{register\_file\_format\_general\_m} \\
\textit{Reserved} & \texttt{27} & \texttt{zero} \\
\hyperref[enum:vecsize_varying_m]{\texttt{vecsize}} & \texttt{28:29} & \texttt{vecsize\_varying\_m (def none)} \\
\hyperref[enum:message_slot_index_m]{\texttt{message\_slot\_index}} & \texttt{30:32} & \texttt{message\_slot\_index\_m} \\
\texttt{sr\_count} & \texttt{33:35} & \texttt{sr\_count\_t} \\
\textit{Reserved} & \texttt{36} & \texttt{zero} \\
\texttt{descriptor\_type} & \texttt{37:38} & \texttt{exact=attribute} \\
\textit{Reserved} & \texttt{39} & \texttt{zero} \\
\texttt{staging\_register\_index} & \texttt{40:45} & \texttt{uint} \\
\texttt{staging\_register\_control\_32} & \texttt{46:47} & \texttt{exact=write} \\
\texttt{opcode1} & \texttt{48:51} & \texttt{exact=LD\_ATTR} \\
\texttt{instruction\_class} & \texttt{52:56} & \texttt{exact=S3} \\
\texttt{fau\_page\_index} & \texttt{57:58} & \texttt{uint} \\
\hyperref[enum:flow_control_m]{\texttt{flow}} & \texttt{59:62} & \texttt{flow\_control\_m} \\
\textit{Reserved} & \texttt{63} & \texttt{zero} \\
\bottomrule
\end{tabularx}
\end{center}
\textbf{Fields Explained:}
\begin{tightitemize}\footnotesize
  \item \texttt{src0}: Primary source operand (\texttt{r0..r63}, \texttt{u0..u63} uniform, \texttt{k*} FAU constant, or \texttt{t0..t2} temporary; \texttt{\textasciicircum{}} marks discard/last-use). \emph{Example:} {\scriptsize\ttfamily \mnem{FADD.f32} \dst{r0}, \src{r1}, \src{r2}} (\textbf{SPIR-V:} \texttt{OpFAdd (Operand 0)}; \textbf{GLSL:} \texttt{float c = a + b;})
  \item \texttt{src1}: Secondary source operand (register, uniform, or FAU constant; in shifts, supplies the shift count byte). \emph{Example:} {\scriptsize\ttfamily \mnem{LSHIFT.i32} \dst{r0}, \src{r1}, \src{r2}\imod{.b0}} (\textbf{SPIR-V:} \texttt{OpShiftLeftLogical (Shift)}; \textbf{GLSL:} \texttt{uint y = a \textless{}\textless{} s;})
  \item \texttt{src2}: Third source operand (FMA addend, \texttt{MUX}/\texttt{CSEL} selector, \texttt{IADDC} carry-in, or bitwise mask in fused shift-logic). \emph{Example:} {\scriptsize\ttfamily \mnem{FMA.f32} \dst{r0}, \src{r1}, \src{r2}, \src{r3}} (\textbf{SPIR-V:} \texttt{GLSL.std.450 Fma (Addend)}; \textbf{GLSL:} \texttt{float d = fma(a, b, c);})
  \item \hyperref[enum:register_file_format_general_m]{\texttt{register\_format}}: Target GPR data conversion format (\texttt{auto32}, \texttt{f32}, \texttt{f16}, \texttt{s32}, \texttt{s16}, \texttt{u32}, \texttt{u16}). \emph{Example:} {\scriptsize\ttfamily \mnem{LD\_TILE}\imod{.f16.v4} \dst{@r0:r1}, \src{r2}} (\textbf{SPIR-V:} \texttt{OpImageRead (SubpassData)}; \textbf{GLSL:} \texttt{f16vec4 c = subpassLoad(inputAtt);})
  \item \hyperref[enum:vecsize_varying_m]{\texttt{vecsize}}: Component count (\texttt{none}=1, \texttt{v2}, \texttt{v3}, \texttt{v4}) for attribute, varying, or blend transfers. \emph{Example:} {\scriptsize\ttfamily \mnem{LD\_ATTR}\imod{.f32.v3} \dst{@r0:r2}, \src{r4}} (\textbf{SPIR-V:} \texttt{OpLoad (vec3 Input)}; \textbf{GLSL:} \texttt{in vec3 a\_normal;})
  \item \hyperref[enum:message_slot_index_m]{\texttt{message\_slot\_index}}: Asynchronous scoreboard slot (\texttt{slot0}--\texttt{slot2}, \texttt{slot7}) tracking completion of this message. \emph{Example:} {\scriptsize\ttfamily \mnem{LOAD.i32}\imod{.slot0} \dst{@r0}, \src{r2}, \imm{\#0}} (\textbf{SPIR-V:} \texttt{OpLoad}; \textbf{GLSL:} \texttt{uint v = ssbo.arr[i];})
  \item \texttt{sr\_count}: Number of consecutive 32-bit staging registers transferred by a load/store/texture message. \emph{Example:} {\scriptsize\ttfamily \mnem{LD\_MULTI.i128} \dst{@r0:r3}, \src{r4}, \imm{\#0}} (\textbf{SPIR-V:} \texttt{OpLoad (wide vector)}; \textbf{GLSL:} \texttt{uvec4 x = buf[i];})
  \item \texttt{descriptor\_type}: Selects descriptor table binding class (attribute, buffer, texture, varying). \emph{Example:} {\scriptsize\ttfamily \mnem{LD\_BUFFER.i32} \dst{@r0}, \src{r2}, \imm{\#0}} (\textbf{SPIR-V:} \texttt{OpLoad (StorageBuffer)}; \textbf{GLSL:} \texttt{uint v = ssbo.arr[i];})
  \item \texttt{staging\_register\_index}: Base GPR index of the contiguous staging register window (\texttt{@rN:rM}) for message payload/return. \emph{Example:} {\scriptsize\ttfamily \mnem{STORE.i64}\imod{.slot7} \src{@r0:r1}, \src{r2}, \imm{\#8}} (\textbf{SPIR-V:} \texttt{OpStore (uvec2)}; \textbf{GLSL:} \texttt{ssbo.pairs[i] = uvec2(a, b);})
  \item \texttt{staging\_register\_control\_32}: Controls the read/write staging register bank and count for async message instructions. \emph{Example:} {\scriptsize\ttfamily \mnem{LOAD.i128}\imod{.slot0} \dst{@r0:r3}, \src{r4}, \imm{\#0}} (\textbf{SPIR-V:} \texttt{OpLoad (vec4)}; \textbf{GLSL:} \texttt{vec4 v = ssbo.data[idx];})
  \item \texttt{opcode1}: Primary sub-opcode within the instruction class uniquely identifying the operation family. \emph{Example:} {\scriptsize\ttfamily \mnem{IADD.s32} \dst{r0}, \src{r1}, \src{r2}} (\textbf{SPIR-V:} \texttt{OpIAdd (operation selector)}; \textbf{GLSL:} \texttt{// Primary opcode bits})
  \item \texttt{instruction\_class}: 4-bit major instruction class selector (\texttt{bits 48:51}) partitioning the 64-bit Valhall opcode space (e.g., \texttt{FMA}, \texttt{CVT}, \texttt{SFU}, \texttt{LS}, \texttt{TEX}, \texttt{BRANCH}). \emph{Example:} {\scriptsize\ttfamily \mnem{FADD.f32} \dst{r0}, \src{r1}, \src{r2}} (\textbf{SPIR-V:} \texttt{N/A (encoding class)}; \textbf{GLSL:} \texttt{// Hardware execution pipe / decode class})
  \item \texttt{fau\_page\_index}: Selects the FAU (Fast Access Uniform) page (\texttt{0..3}) when a source operand references hardware small constants, thread IDs, or lane constants. \emph{Example:} {\scriptsize\ttfamily \mnem{FADD.f32} \dst{r0}, \src{r1}, \src{k1f}} (\textbf{SPIR-V:} \texttt{OpConstant / BuiltIn}; \textbf{GLSL:} \texttt{float y = x + 1.0;})
  \item \hyperref[enum:flow_control_m]{\texttt{flow}}: 4-bit embedded flow-control \& scoreboard wait modifier (\texttt{bits 59:62}: \texttt{none}, \texttt{wait0..wait0126}, \texttt{wait}, \texttt{reconverge}, \texttt{discard}, \texttt{end}). \emph{Example:} {\scriptsize\ttfamily \mnem{FADD.f32}\imod{.wait0.reconverge} \dst{r0}, \src{r1}, \src{r2}} (\textbf{SPIR-V:} \texttt{OpControlBarrier / OpReturn}; \textbf{GLSL:} \texttt{// Scoreboard wait + warp reconvergence})
\end{tightitemize}
\clearpage

\subsection*{LD\_ATTR\_IMM}
\addcontentsline{toc}{subsection}{LD\_ATTR\_IMM}
Fetches \texttt{vecsize} vertex-attribute components for \texttt{(vertex=src0, instance=src1)} from the Attribute descriptor at immediate \texttt{(table=\#attribute\_table\_index, index=\#attribute\_index)}, converts them to \texttt{register\_format}, and writes them into \texttt{@sw}.
\par\textbf{Notes:}
\begin{tightitemize}
  \item Immediate-handle form of \texttt{LD\_ATTR} (\texttt{ISA.xml:68}).
  \item \textbf{Source references:} \texttt{bifrost/valhall/ISA.xml:68}, \texttt{isa-v9-v14.xml}.
\end{tightitemize}
\codebox{\mnem{LD\_ATTR\_IMM}\imod{\{.message\_slot\_index\}}\imod{\{.flow\}}\imod{\{.vecsize\}}\imod{\{.register\_format\}} \dst{@sw}, \src{src0}, \src{src1}, \imm{\#attribute\_index}, \imm{\#attribute\_table\_index}}
\textbf{Operands:}
\begin{tightitemize}\footnotesize
  \item \texttt{\dst{@sw}} (\texttt{staging\_register\_index} + \texttt{sr\_count}): Destination staging register range receiving \texttt{vecsize} converted attribute components.
  \item \texttt{\src{src0}}: Vertex index (\texttt{VertexIndex}, preloaded in \texttt{r60}).
  \item \texttt{\src{src1}}: Instance index (\texttt{InstanceIndex}, preloaded in \texttt{r61}).
  \item \texttt{\imm{\#attribute\_index}} (\texttt{attribute\_index}): Immediate index of the Attribute descriptor within the table.
  \item \texttt{\imm{\#attribute\_table\_index}} (\texttt{encoded\_attribute\_table\_index}): Immediate SRT table index holding the Attribute descriptor array.
\end{tightitemize}
\par\textbf{Closest SPIR-V Op:} \texttt{OpLoad (vertex input)}
\par\textbf{Pseudo-GLSL (Operational Semantics):}
\codebox{AttrDesc a = SRT[\#attribute\_table\_index].Attr[\#attribute\_index];\\@sw = fetchConvertAttr(a, vert = src0, inst = src1, register\_format, vecsize);}
\textbf{Execution Context:}
\begin{tightitemize}\footnotesize
  \item \textbf{Attribute descriptor chain (\hyperref[sec:csf-srt]{\texttt{IDVSSR::SRT\_0}}, \texttt{r0:r1}):} Masks 64B-aligned VA (\texttt{\& \textasciitilde{}0x3F}, \texttt{count=bits[5:0]}) $\to$ \texttt{Resource[table]} (16B) $\to$ \hyperref[sec:csf-buf-attr]{\texttt{Attribute[index]}} (32B: \texttt{table}, \texttt{buffer\_index}, \texttt{offset}, \texttt{stride}, \texttt{format}) $\to$ \texttt{Resource[Attribute::table]} $\to$ \texttt{Buffer[Attribute::buffer\_index]} (32B: \texttt{address}, \texttt{size}).
  \item \textbf{Vertex stage preloads (\hyperref[sec:csf-spd]{\texttt{SPD\_0}/\texttt{SPD\_1}}, \texttt{r16..r19}):} The hardware Index Fetcher reads \hyperref[sec:csf-regs]{\texttt{IDVSSR::INDEX\_BUFFER}} (\texttt{r54:r55}) and pre-fills \texttt{r59, r60} with \texttt{VertexIndex} (\texttt{+ VERTEX\_OFFSET + GLOBAL\_ATTR\_OFF}) and \texttt{r61} with \texttt{InstanceIndex} (\texttt{+ INSTANCE\_OFFSET}) so Position and Varying shaders can pull vertex attributes directly.
\end{tightitemize}
\par\textbf{Worked Examples:}
\par\textbf{Example 1 (Static vertex position attribute (`table=0, index=0`)):}
\codebox{\mnem{LD\_ATTR\_IMM}\imod{.f32.v4.slot0} \dst{@r0:r3}, \src{r60}, \src{r61}, \imm{\#0}, \imm{\#0}}
Loads 4x FP32 vertex attribute from immediate \texttt{(table=0, index=0)} into \texttt{@r0:r3}.
\par\textbf{Example 2 (Relaxed-precision normal (`f16 vec3`)):}
\codebox{\mnem{LD\_ATTR\_IMM}\imod{.f16.v3.slot0} \dst{@r4:r5}, \src{r60}, \src{r61}, \imm{\#0}, \imm{\#1}}
Converts attribute 1 to 3x FP16 packed into \texttt{@r4:r5}.
\par\textbf{Mined SPIR-V $\to$ Valhall ISA Example (Fetching a vec2 vertex input attribute in a Vertex shader --- Vertex shader, Mali-G710 (v10)):}
\codebox{\textbf{; SPIR-V (input)}\\\%uv = OpLoad \%v2float \%in\_location\_1\\[3pt]\textbf{; Valhall ISA (malioc output)}\\LD\_ATTR\_IMM.slot0.v2.f32 @r0:r1, r60\textasciicircum{}, r61\textasciicircum{}, \#0x1, \#0x1}
\texttt{LD\_ATTR\_IMM} fetches 2 FP32 components (\texttt{v2.f32}) into \texttt{@r0:r1} using hardware-preloaded \texttt{VertexIndex} (\texttt{r60\textasciicircum{}}) and \texttt{InstanceIndex} (\texttt{r61\textasciicircum{}}) from attribute descriptor \texttt{(table=\#0x1, index=\#0x1)}.
\begin{center}\footnotesize
\begin{tabularx}{\linewidth}{@{}lllX@{}}
\toprule
Variant & Arch & Class & Opcode Fields \\
\midrule
\texttt{(none)} & v9--v14 & \texttt{S2} & \texttt{opcode1=LD\_ATTR\_IMM} \\
\bottomrule
\end{tabularx}
\end{center}
\begin{center}\footnotesize
\begin{tabularx}{0.88\linewidth}{@{}l l X@{}}
\toprule
Field Name & Bits & Type / Fixed Value \\
\midrule
\texttt{src0} & \texttt{0:7} & \texttt{SourceEncoding} \\
\texttt{src1} & \texttt{8:15} & \texttt{SourceEncoding} \\
\texttt{encoded\_attribute\_table\_index} & \texttt{16:19} & \texttt{uint} \\
\texttt{attribute\_index} & \texttt{20:23} & \texttt{uint} \\
\hyperref[enum:register_file_format_general_m]{\texttt{register\_format}} & \texttt{24:26} & \texttt{register\_file\_format\_general\_m} \\
\textit{Reserved} & \texttt{27} & \texttt{zero} \\
\hyperref[enum:vecsize_varying_m]{\texttt{vecsize}} & \texttt{28:29} & \texttt{vecsize\_varying\_m (def none)} \\
\hyperref[enum:message_slot_index_m]{\texttt{message\_slot\_index}} & \texttt{30:32} & \texttt{message\_slot\_index\_m} \\
\texttt{sr\_count} & \texttt{33:35} & \texttt{sr\_count\_t} \\
\textit{Reserved} & \texttt{36} & \texttt{zero} \\
\texttt{descriptor\_type} & \texttt{37:38} & \texttt{exact=attribute} \\
\textit{Reserved} & \texttt{39} & \texttt{zero} \\
\texttt{staging\_register\_index} & \texttt{40:45} & \texttt{uint} \\
\texttt{staging\_register\_control\_32} & \texttt{46:47} & \texttt{exact=write} \\
\texttt{opcode1} & \texttt{48:51} & \texttt{exact=LD\_ATTR\_IMM} \\
\texttt{instruction\_class} & \texttt{52:56} & \texttt{exact=S2} \\
\texttt{fau\_page\_index} & \texttt{57:58} & \texttt{uint} \\
\hyperref[enum:flow_control_m]{\texttt{flow}} & \texttt{59:62} & \texttt{flow\_control\_m} \\
\textit{Reserved} & \texttt{63} & \texttt{zero} \\
\bottomrule
\end{tabularx}
\end{center}
\textbf{Fields Explained:}
\begin{tightitemize}\footnotesize
  \item \texttt{src0}: Primary source operand (\texttt{r0..r63}, \texttt{u0..u63} uniform, \texttt{k*} FAU constant, or \texttt{t0..t2} temporary; \texttt{\textasciicircum{}} marks discard/last-use). \emph{Example:} {\scriptsize\ttfamily \mnem{FADD.f32} \dst{r0}, \src{r1}, \src{r2}} (\textbf{SPIR-V:} \texttt{OpFAdd (Operand 0)}; \textbf{GLSL:} \texttt{float c = a + b;})
  \item \texttt{src1}: Secondary source operand (register, uniform, or FAU constant; in shifts, supplies the shift count byte). \emph{Example:} {\scriptsize\ttfamily \mnem{LSHIFT.i32} \dst{r0}, \src{r1}, \src{r2}\imod{.b0}} (\textbf{SPIR-V:} \texttt{OpShiftLeftLogical (Shift)}; \textbf{GLSL:} \texttt{uint y = a \textless{}\textless{} s;})
  \item \texttt{encoded\_attribute\_table\_index}: Encoded descriptor table index for \textbackslash\{\}texttt\{encoded\textbackslash\{\}\_attribute\textbackslash\{\}\_table\textbackslash\{\}\_index\}. \emph{Example:} {\scriptsize\ttfamily \mnem{LD\_BUFFER.i32} \dst{@r0}, \src{r2}, \imm{\#0}} (\textbf{SPIR-V:} \texttt{OpAccessChain (DescriptorSet)}; \textbf{GLSL:} \texttt{uint v = buf[i];})
  \item \texttt{attribute\_index}: Resource/slot index field \textbackslash\{\}texttt\{attribute\textbackslash\{\}\_index\} selecting the target descriptor or slot. \emph{Example:} {\scriptsize\ttfamily \mnem{LD\_VAR\_BUF\_IMM.f32} \dst{@r0:r3}, \imm{\#0}} (\textbf{SPIR-V:} \texttt{OpLoad (Descriptor)}; \textbf{GLSL:} \texttt{vec4 v = v\_in;})
  \item \hyperref[enum:register_file_format_general_m]{\texttt{register\_format}}: Target GPR data conversion format (\texttt{auto32}, \texttt{f32}, \texttt{f16}, \texttt{s32}, \texttt{s16}, \texttt{u32}, \texttt{u16}). \emph{Example:} {\scriptsize\ttfamily \mnem{LD\_TILE}\imod{.f16.v4} \dst{@r0:r1}, \src{r2}} (\textbf{SPIR-V:} \texttt{OpImageRead (SubpassData)}; \textbf{GLSL:} \texttt{f16vec4 c = subpassLoad(inputAtt);})
  \item \hyperref[enum:vecsize_varying_m]{\texttt{vecsize}}: Component count (\texttt{none}=1, \texttt{v2}, \texttt{v3}, \texttt{v4}) for attribute, varying, or blend transfers. \emph{Example:} {\scriptsize\ttfamily \mnem{LD\_ATTR}\imod{.f32.v3} \dst{@r0:r2}, \src{r4}} (\textbf{SPIR-V:} \texttt{OpLoad (vec3 Input)}; \textbf{GLSL:} \texttt{in vec3 a\_normal;})
  \item \hyperref[enum:message_slot_index_m]{\texttt{message\_slot\_index}}: Asynchronous scoreboard slot (\texttt{slot0}--\texttt{slot2}, \texttt{slot7}) tracking completion of this message. \emph{Example:} {\scriptsize\ttfamily \mnem{LOAD.i32}\imod{.slot0} \dst{@r0}, \src{r2}, \imm{\#0}} (\textbf{SPIR-V:} \texttt{OpLoad}; \textbf{GLSL:} \texttt{uint v = ssbo.arr[i];})
  \item \texttt{sr\_count}: Number of consecutive 32-bit staging registers transferred by a load/store/texture message. \emph{Example:} {\scriptsize\ttfamily \mnem{LD\_MULTI.i128} \dst{@r0:r3}, \src{r4}, \imm{\#0}} (\textbf{SPIR-V:} \texttt{OpLoad (wide vector)}; \textbf{GLSL:} \texttt{uvec4 x = buf[i];})
  \item \texttt{descriptor\_type}: Selects descriptor table binding class (attribute, buffer, texture, varying). \emph{Example:} {\scriptsize\ttfamily \mnem{LD\_BUFFER.i32} \dst{@r0}, \src{r2}, \imm{\#0}} (\textbf{SPIR-V:} \texttt{OpLoad (StorageBuffer)}; \textbf{GLSL:} \texttt{uint v = ssbo.arr[i];})
  \item \texttt{staging\_register\_index}: Base GPR index of the contiguous staging register window (\texttt{@rN:rM}) for message payload/return. \emph{Example:} {\scriptsize\ttfamily \mnem{STORE.i64}\imod{.slot7} \src{@r0:r1}, \src{r2}, \imm{\#8}} (\textbf{SPIR-V:} \texttt{OpStore (uvec2)}; \textbf{GLSL:} \texttt{ssbo.pairs[i] = uvec2(a, b);})
  \item \texttt{staging\_register\_control\_32}: Controls the read/write staging register bank and count for async message instructions. \emph{Example:} {\scriptsize\ttfamily \mnem{LOAD.i128}\imod{.slot0} \dst{@r0:r3}, \src{r4}, \imm{\#0}} (\textbf{SPIR-V:} \texttt{OpLoad (vec4)}; \textbf{GLSL:} \texttt{vec4 v = ssbo.data[idx];})
  \item \texttt{opcode1}: Primary sub-opcode within the instruction class uniquely identifying the operation family. \emph{Example:} {\scriptsize\ttfamily \mnem{IADD.s32} \dst{r0}, \src{r1}, \src{r2}} (\textbf{SPIR-V:} \texttt{OpIAdd (operation selector)}; \textbf{GLSL:} \texttt{// Primary opcode bits})
  \item \texttt{instruction\_class}: 4-bit major instruction class selector (\texttt{bits 48:51}) partitioning the 64-bit Valhall opcode space (e.g., \texttt{FMA}, \texttt{CVT}, \texttt{SFU}, \texttt{LS}, \texttt{TEX}, \texttt{BRANCH}). \emph{Example:} {\scriptsize\ttfamily \mnem{FADD.f32} \dst{r0}, \src{r1}, \src{r2}} (\textbf{SPIR-V:} \texttt{N/A (encoding class)}; \textbf{GLSL:} \texttt{// Hardware execution pipe / decode class})
  \item \texttt{fau\_page\_index}: Selects the FAU (Fast Access Uniform) page (\texttt{0..3}) when a source operand references hardware small constants, thread IDs, or lane constants. \emph{Example:} {\scriptsize\ttfamily \mnem{FADD.f32} \dst{r0}, \src{r1}, \src{k1f}} (\textbf{SPIR-V:} \texttt{OpConstant / BuiltIn}; \textbf{GLSL:} \texttt{float y = x + 1.0;})
  \item \hyperref[enum:flow_control_m]{\texttt{flow}}: 4-bit embedded flow-control \& scoreboard wait modifier (\texttt{bits 59:62}: \texttt{none}, \texttt{wait0..wait0126}, \texttt{wait}, \texttt{reconverge}, \texttt{discard}, \texttt{end}). \emph{Example:} {\scriptsize\ttfamily \mnem{FADD.f32}\imod{.wait0.reconverge} \dst{r0}, \src{r1}, \src{r2}} (\textbf{SPIR-V:} \texttt{OpControlBarrier / OpReturn}; \textbf{GLSL:} \texttt{// Scoreboard wait + warp reconvergence})
\end{tightitemize}
\clearpage

\subsection*{LD\_TEX}
\addcontentsline{toc}{subsection}{LD\_TEX}
Loads \texttt{vecsize} components of \texttt{register\_format} from image/texture descriptor at dynamic handle (\texttt{src2}) at integer coords (\texttt{src0}, \texttt{src1}).
\par\textbf{Notes:}
\begin{tightitemize}
  \item Direct image texel load (\texttt{ISA.xml:83}).
  \item \textbf{Source references:} \texttt{kraid/ops.rs:2841-2864}, \texttt{bifrost/valhall/ISA.xml:83}.
\end{tightitemize}
\codebox{\mnem{LD\_TEX}\imod{\{.message\_slot\_index\}}\imod{\{.flow\}}\imod{\{.vecsize\}}\imod{\{.register\_format\}} \dst{@sw}, \src{src0}, \src{src1}, \src{src2}}
\textbf{Operands:}
\begin{tightitemize}\footnotesize
  \item \texttt{\dst{@sw}} (\texttt{staging\_register\_index} + \texttt{sr\_count}): Destination staging register range receiving the loaded texel components.
  \item \texttt{\src{src0}}: Packed \texttt{u16vec2(x, y)} integer texel coordinates.
  \item \texttt{\src{src1}}: 3D slice / array layer \texttt{z} coordinate.
  \item \texttt{\src{src2}}: Dynamic image/texture descriptor handle.
\end{tightitemize}
\par\textbf{Closest SPIR-V Op:} \texttt{OpImageRead (dynamic handle)}
\par\textbf{Pseudo-GLSL (Operational Semantics):}
\codebox{@sw = imageLoad(handle = src2, ivec2(src0.h0, src0.h1), z = src1, register\_format, vecsize);}
\textbf{Execution Context:}
\begin{tightitemize}\footnotesize
  \item \textbf{Texture \& Sampler chain (\hyperref[sec:csf-srt]{\texttt{SRT\_2}} Frag / \texttt{SRT\_0} Comp/Vert):} 64B-aligned VA (\texttt{\& \textasciitilde{}0x3F}) $\to$ \texttt{Resource[table]} $\to$ \hyperref[sec:csf-tex-samp]{\texttt{Texture[tex\_idx]}} (32B, \texttt{type=2}) $\to$ \texttt{GenericPlane[mip]} (32B, \texttt{type=10}: \texttt{pointer}, \texttt{row\_stride}, Morton 16$\times$16 / AFBC) + \texttt{Sampler[samp\_idx]} (32B, \texttt{type=1}: filter, wrap, LOD clamp/bias, compare).
  \item \textbf{Inline FAU \& Special Operands (\hyperref[sec:csf-tsd-fau]{\texttt{FAU\_k}}, \texttt{r8..r15}):} Any 8-bit source with \texttt{mode=2} reads zero-latency uniform \texttt{u0..u31} from \texttt{FAU\_k} (\texttt{slots64 = bits[63:56]}); \texttt{mode=3} reads \texttt{SmallConstant} (Table~1.2) or \texttt{FauSpecial} (\texttt{pc}, \texttt{lane\_id}, \texttt{lane\_width}, \texttt{tls\_ptr}, \texttt{wls\_ptr}, \texttt{blend\_descriptor\_0..7}).
\end{tightitemize}
\par\textbf{Worked Examples:}
\par\textbf{Example 1 (Direct texel load from bindless image handle):}
\codebox{\mnem{LD\_TEX}\imod{.f32.v4.slot0} \dst{@r0:r3}, \src{r4}, \src{r5}, \src{r6}}
Reads unfiltered texel from dynamic image handle \texttt{r6} at integer coordinates \texttt{(r4, r5)} into 4x FP32 \texttt{@r0:r3}.
\par\textbf{Example 2 (Unsigned integer image read):}
\codebox{\mnem{LD\_TEX}\imod{.u32.v4.slot0} \dst{@r0:r3}, \src{r4}, \src{r5}, \src{r6}}
Reads raw \texttt{uvec4} texel values from \texttt{uimage2D} without float normalization.

\begin{center}\footnotesize
\begin{tabularx}{\linewidth}{@{}lllX@{}}
\toprule
Variant & Arch & Class & Opcode Fields \\
\midrule
\texttt{(none)} & v9--v14 & \texttt{S3} & \texttt{opcode1=LD\_ATTR} \\
\bottomrule
\end{tabularx}
\end{center}
\begin{center}\footnotesize
\begin{tabularx}{0.88\linewidth}{@{}l l X@{}}
\toprule
Field Name & Bits & Type / Fixed Value \\
\midrule
\texttt{src0} & \texttt{0:7} & \texttt{SourceEncoding} \\
\texttt{src1} & \texttt{8:15} & \texttt{SourceEncoding} \\
\texttt{src2} & \texttt{16:23} & \texttt{SourceEncoding} \\
\hyperref[enum:register_file_format_general_m]{\texttt{register\_format}} & \texttt{24:26} & \texttt{register\_file\_format\_general\_m} \\
\textit{Reserved} & \texttt{27} & \texttt{zero} \\
\hyperref[enum:vecsize_varying_m]{\texttt{vecsize}} & \texttt{28:29} & \texttt{vecsize\_varying\_m (def none)} \\
\hyperref[enum:message_slot_index_m]{\texttt{message\_slot\_index}} & \texttt{30:32} & \texttt{message\_slot\_index\_m} \\
\texttt{sr\_count} & \texttt{33:35} & \texttt{sr\_count\_t} \\
\textit{Reserved} & \texttt{36} & \texttt{zero} \\
\texttt{descriptor\_type} & \texttt{37:38} & \texttt{exact=texture\_u16} \\
\textit{Reserved} & \texttt{39} & \texttt{zero} \\
\texttt{staging\_register\_index} & \texttt{40:45} & \texttt{uint} \\
\texttt{staging\_register\_control\_32} & \texttt{46:47} & \texttt{exact=write} \\
\texttt{opcode1} & \texttt{48:51} & \texttt{exact=LD\_ATTR} \\
\texttt{instruction\_class} & \texttt{52:56} & \texttt{exact=S3} \\
\texttt{fau\_page\_index} & \texttt{57:58} & \texttt{uint} \\
\hyperref[enum:flow_control_m]{\texttt{flow}} & \texttt{59:62} & \texttt{flow\_control\_m} \\
\textit{Reserved} & \texttt{63} & \texttt{zero} \\
\bottomrule
\end{tabularx}
\end{center}
\textbf{Fields Explained:}
\begin{tightitemize}\footnotesize
  \item \texttt{src0}: Primary source operand (\texttt{r0..r63}, \texttt{u0..u63} uniform, \texttt{k*} FAU constant, or \texttt{t0..t2} temporary; \texttt{\textasciicircum{}} marks discard/last-use). \emph{Example:} {\scriptsize\ttfamily \mnem{FADD.f32} \dst{r0}, \src{r1}, \src{r2}} (\textbf{SPIR-V:} \texttt{OpFAdd (Operand 0)}; \textbf{GLSL:} \texttt{float c = a + b;})
  \item \texttt{src1}: Secondary source operand (register, uniform, or FAU constant; in shifts, supplies the shift count byte). \emph{Example:} {\scriptsize\ttfamily \mnem{LSHIFT.i32} \dst{r0}, \src{r1}, \src{r2}\imod{.b0}} (\textbf{SPIR-V:} \texttt{OpShiftLeftLogical (Shift)}; \textbf{GLSL:} \texttt{uint y = a \textless{}\textless{} s;})
  \item \texttt{src2}: Third source operand (FMA addend, \texttt{MUX}/\texttt{CSEL} selector, \texttt{IADDC} carry-in, or bitwise mask in fused shift-logic). \emph{Example:} {\scriptsize\ttfamily \mnem{FMA.f32} \dst{r0}, \src{r1}, \src{r2}, \src{r3}} (\textbf{SPIR-V:} \texttt{GLSL.std.450 Fma (Addend)}; \textbf{GLSL:} \texttt{float d = fma(a, b, c);})
  \item \hyperref[enum:register_file_format_general_m]{\texttt{register\_format}}: Target GPR data conversion format (\texttt{auto32}, \texttt{f32}, \texttt{f16}, \texttt{s32}, \texttt{s16}, \texttt{u32}, \texttt{u16}). \emph{Example:} {\scriptsize\ttfamily \mnem{LD\_TILE}\imod{.f16.v4} \dst{@r0:r1}, \src{r2}} (\textbf{SPIR-V:} \texttt{OpImageRead (SubpassData)}; \textbf{GLSL:} \texttt{f16vec4 c = subpassLoad(inputAtt);})
  \item \hyperref[enum:vecsize_varying_m]{\texttt{vecsize}}: Component count (\texttt{none}=1, \texttt{v2}, \texttt{v3}, \texttt{v4}) for attribute, varying, or blend transfers. \emph{Example:} {\scriptsize\ttfamily \mnem{LD\_ATTR}\imod{.f32.v3} \dst{@r0:r2}, \src{r4}} (\textbf{SPIR-V:} \texttt{OpLoad (vec3 Input)}; \textbf{GLSL:} \texttt{in vec3 a\_normal;})
  \item \hyperref[enum:message_slot_index_m]{\texttt{message\_slot\_index}}: Asynchronous scoreboard slot (\texttt{slot0}--\texttt{slot2}, \texttt{slot7}) tracking completion of this message. \emph{Example:} {\scriptsize\ttfamily \mnem{LOAD.i32}\imod{.slot0} \dst{@r0}, \src{r2}, \imm{\#0}} (\textbf{SPIR-V:} \texttt{OpLoad}; \textbf{GLSL:} \texttt{uint v = ssbo.arr[i];})
  \item \texttt{sr\_count}: Number of consecutive 32-bit staging registers transferred by a load/store/texture message. \emph{Example:} {\scriptsize\ttfamily \mnem{LD\_MULTI.i128} \dst{@r0:r3}, \src{r4}, \imm{\#0}} (\textbf{SPIR-V:} \texttt{OpLoad (wide vector)}; \textbf{GLSL:} \texttt{uvec4 x = buf[i];})
  \item \texttt{descriptor\_type}: Selects descriptor table binding class (attribute, buffer, texture, varying). \emph{Example:} {\scriptsize\ttfamily \mnem{LD\_BUFFER.i32} \dst{@r0}, \src{r2}, \imm{\#0}} (\textbf{SPIR-V:} \texttt{OpLoad (StorageBuffer)}; \textbf{GLSL:} \texttt{uint v = ssbo.arr[i];})
  \item \texttt{staging\_register\_index}: Base GPR index of the contiguous staging register window (\texttt{@rN:rM}) for message payload/return. \emph{Example:} {\scriptsize\ttfamily \mnem{STORE.i64}\imod{.slot7} \src{@r0:r1}, \src{r2}, \imm{\#8}} (\textbf{SPIR-V:} \texttt{OpStore (uvec2)}; \textbf{GLSL:} \texttt{ssbo.pairs[i] = uvec2(a, b);})
  \item \texttt{staging\_register\_control\_32}: Controls the read/write staging register bank and count for async message instructions. \emph{Example:} {\scriptsize\ttfamily \mnem{LOAD.i128}\imod{.slot0} \dst{@r0:r3}, \src{r4}, \imm{\#0}} (\textbf{SPIR-V:} \texttt{OpLoad (vec4)}; \textbf{GLSL:} \texttt{vec4 v = ssbo.data[idx];})
  \item \texttt{opcode1}: Primary sub-opcode within the instruction class uniquely identifying the operation family. \emph{Example:} {\scriptsize\ttfamily \mnem{IADD.s32} \dst{r0}, \src{r1}, \src{r2}} (\textbf{SPIR-V:} \texttt{OpIAdd (operation selector)}; \textbf{GLSL:} \texttt{// Primary opcode bits})
  \item \texttt{instruction\_class}: 4-bit major instruction class selector (\texttt{bits 48:51}) partitioning the 64-bit Valhall opcode space (e.g., \texttt{FMA}, \texttt{CVT}, \texttt{SFU}, \texttt{LS}, \texttt{TEX}, \texttt{BRANCH}). \emph{Example:} {\scriptsize\ttfamily \mnem{FADD.f32} \dst{r0}, \src{r1}, \src{r2}} (\textbf{SPIR-V:} \texttt{N/A (encoding class)}; \textbf{GLSL:} \texttt{// Hardware execution pipe / decode class})
  \item \texttt{fau\_page\_index}: Selects the FAU (Fast Access Uniform) page (\texttt{0..3}) when a source operand references hardware small constants, thread IDs, or lane constants. \emph{Example:} {\scriptsize\ttfamily \mnem{FADD.f32} \dst{r0}, \src{r1}, \src{k1f}} (\textbf{SPIR-V:} \texttt{OpConstant / BuiltIn}; \textbf{GLSL:} \texttt{float y = x + 1.0;})
  \item \hyperref[enum:flow_control_m]{\texttt{flow}}: 4-bit embedded flow-control \& scoreboard wait modifier (\texttt{bits 59:62}: \texttt{none}, \texttt{wait0..wait0126}, \texttt{wait}, \texttt{reconverge}, \texttt{discard}, \texttt{end}). \emph{Example:} {\scriptsize\ttfamily \mnem{FADD.f32}\imod{.wait0.reconverge} \dst{r0}, \src{r1}, \src{r2}} (\textbf{SPIR-V:} \texttt{OpControlBarrier / OpReturn}; \textbf{GLSL:} \texttt{// Scoreboard wait + warp reconvergence})
\end{tightitemize}
\clearpage

\subsection*{LD\_TEX\_IMM}
\addcontentsline{toc}{subsection}{LD\_TEX\_IMM}
Reads \texttt{vecsize} texel components at integer coordinates \texttt{(xy=src0, z=src1)} from the Image/Texture descriptor at immediate \texttt{(table=\#texture\_table\_index, index=\#texture\_index)}, converts them to \texttt{register\_format}, and writes them into \texttt{@sw}.
\par\textbf{Notes:}
\begin{tightitemize}
  \item Immediate-handle form of \texttt{LD\_TEX} (\texttt{ISA.xml:85}).
  \item \textbf{Source references:} \texttt{bifrost/valhall/ISA.xml:85}, \texttt{isa-v9-v14.xml}.
\end{tightitemize}
\codebox{\mnem{LD\_TEX\_IMM}\imod{\{.message\_slot\_index\}}\imod{\{.flow\}}\imod{\{.vecsize\}}\imod{\{.register\_format\}} \dst{@sw}, \src{src0}, \src{src1}, \imm{\#texture\_index}, \imm{\#texture\_table\_index}}
\textbf{Operands:}
\begin{tightitemize}\footnotesize
  \item \texttt{\dst{@sw}} (\texttt{staging\_register\_index} + \texttt{sr\_count}): Destination staging register range receiving the loaded texel components.
  \item \texttt{\src{src0}}: Packed \texttt{u16vec2(x, y)} integer texel coordinates.
  \item \texttt{\src{src1}}: 3D slice / array layer \texttt{z} coordinate.
  \item \texttt{\imm{\#texture\_index}} (\texttt{texture\_index}): Immediate index of the Image/Texture descriptor within the table.
  \item \texttt{\imm{\#texture\_table\_index}} (\texttt{encoded\_texture\_table\_index}): Immediate SRT table index holding the Image/Texture descriptor array.
\end{tightitemize}
\par\textbf{Closest SPIR-V Op:} \texttt{OpImageRead}
\par\textbf{Pseudo-GLSL (Operational Semantics):}
\codebox{TexDesc t = SRT[\#texture\_table\_index].Tex[\#texture\_index];\\@sw = imageLoad(t, ivec2(src0.h0, src0.h1), z = src1, register\_format, vecsize);}
\textbf{Execution Context:}
\begin{tightitemize}\footnotesize
  \item \textbf{Texture \& Sampler chain (\hyperref[sec:csf-srt]{\texttt{SRT\_2}} Frag / \texttt{SRT\_0} Comp/Vert):} 64B-aligned VA (\texttt{\& \textasciitilde{}0x3F}) $\to$ \texttt{Resource[table]} $\to$ \hyperref[sec:csf-tex-samp]{\texttt{Texture[tex\_idx]}} (32B, \texttt{type=2}) $\to$ \texttt{GenericPlane[mip]} (32B, \texttt{type=10}: \texttt{pointer}, \texttt{row\_stride}, Morton 16$\times$16 / AFBC) + \texttt{Sampler[samp\_idx]} (32B, \texttt{type=1}: filter, wrap, LOD clamp/bias, compare).
  \item \textbf{Inline FAU \& Special Operands (\hyperref[sec:csf-tsd-fau]{\texttt{FAU\_k}}, \texttt{r8..r15}):} Any 8-bit source with \texttt{mode=2} reads zero-latency uniform \texttt{u0..u31} from \texttt{FAU\_k} (\texttt{slots64 = bits[63:56]}); \texttt{mode=3} reads \texttt{SmallConstant} (Table~1.2) or \texttt{FauSpecial} (\texttt{pc}, \texttt{lane\_id}, \texttt{lane\_width}, \texttt{tls\_ptr}, \texttt{wls\_ptr}, \texttt{blend\_descriptor\_0..7}).
\end{tightitemize}
\par\textbf{Worked Examples:}
\par\textbf{Example 1 (Static `uimage2D` texel read):}
\codebox{\mnem{LD\_TEX\_IMM}\imod{.u32.v4.slot0} \dst{@r0:r3}, \src{r4}, \src{r5}, \imm{\#0}, \imm{\#0}}
Reads texel \texttt{(r4.h0, r4.h1)} from static image \texttt{(table=0, index=0)} into \texttt{@r0:r3}.
\par\textbf{Example 2 (FP16 storage image read (`table=0, index=2`)):}
\codebox{\mnem{LD\_TEX\_IMM}\imod{.f16.v4.slot0} \dst{@r0:r1}, \src{r4}, \src{r5}, \imm{\#0}, \imm{\#2}}
Reads 4x FP16 channels from binding 2 into \texttt{@r0:r1}.
\par\textbf{Mined SPIR-V $\to$ Valhall ISA Example (Direct storage-image texel read (OpImageRead) --- Compute shader, Mali-G710 (v10)):}
\codebox{\textbf{; SPIR-V (input)}\\\%texel = OpImageRead \%v4uint \%img \%coord None\\[3pt]\textbf{; Valhall ISA (malioc output)}\\LD\_TEX\_IMM.slot0.u32 @r1, k0, k0, \#0x4, \#0x4}
\texttt{OpImageRead} lowers to \texttt{LD\_TEX\_IMM}, reading unsigned 32-bit texel data from descriptor \texttt{(table=\#0x4, index=\#0x4)} at integer coordinates \texttt{(src0, src1)}.
\begin{center}\footnotesize
\begin{tabularx}{\linewidth}{@{}lllX@{}}
\toprule
Variant & Arch & Class & Opcode Fields \\
\midrule
\texttt{(none)} & v9--v14 & \texttt{S2} & \texttt{opcode1=LD\_ATTR\_IMM} \\
\bottomrule
\end{tabularx}
\end{center}
\begin{center}\footnotesize
\begin{tabularx}{0.88\linewidth}{@{}l l X@{}}
\toprule
Field Name & Bits & Type / Fixed Value \\
\midrule
\texttt{src0} & \texttt{0:7} & \texttt{SourceEncoding} \\
\texttt{src1} & \texttt{8:15} & \texttt{SourceEncoding} \\
\texttt{encoded\_texture\_table\_index} & \texttt{16:19} & \texttt{uint} \\
\texttt{texture\_index} & \texttt{20:23} & \texttt{uint} \\
\hyperref[enum:register_file_format_general_m]{\texttt{register\_format}} & \texttt{24:26} & \texttt{register\_file\_format\_general\_m} \\
\textit{Reserved} & \texttt{27} & \texttt{zero} \\
\hyperref[enum:vecsize_varying_m]{\texttt{vecsize}} & \texttt{28:29} & \texttt{vecsize\_varying\_m (def none)} \\
\hyperref[enum:message_slot_index_m]{\texttt{message\_slot\_index}} & \texttt{30:32} & \texttt{message\_slot\_index\_m} \\
\texttt{sr\_count} & \texttt{33:35} & \texttt{sr\_count\_t} \\
\textit{Reserved} & \texttt{36} & \texttt{zero} \\
\texttt{descriptor\_type} & \texttt{37:38} & \texttt{exact=texture\_u16} \\
\textit{Reserved} & \texttt{39} & \texttt{zero} \\
\texttt{staging\_register\_index} & \texttt{40:45} & \texttt{uint} \\
\texttt{staging\_register\_control\_32} & \texttt{46:47} & \texttt{exact=write} \\
\texttt{opcode1} & \texttt{48:51} & \texttt{exact=LD\_ATTR\_IMM} \\
\texttt{instruction\_class} & \texttt{52:56} & \texttt{exact=S2} \\
\texttt{fau\_page\_index} & \texttt{57:58} & \texttt{uint} \\
\hyperref[enum:flow_control_m]{\texttt{flow}} & \texttt{59:62} & \texttt{flow\_control\_m} \\
\textit{Reserved} & \texttt{63} & \texttt{zero} \\
\bottomrule
\end{tabularx}
\end{center}
\textbf{Fields Explained:}
\begin{tightitemize}\footnotesize
  \item \texttt{src0}: Primary source operand (\texttt{r0..r63}, \texttt{u0..u63} uniform, \texttt{k*} FAU constant, or \texttt{t0..t2} temporary; \texttt{\textasciicircum{}} marks discard/last-use). \emph{Example:} {\scriptsize\ttfamily \mnem{FADD.f32} \dst{r0}, \src{r1}, \src{r2}} (\textbf{SPIR-V:} \texttt{OpFAdd (Operand 0)}; \textbf{GLSL:} \texttt{float c = a + b;})
  \item \texttt{src1}: Secondary source operand (register, uniform, or FAU constant; in shifts, supplies the shift count byte). \emph{Example:} {\scriptsize\ttfamily \mnem{LSHIFT.i32} \dst{r0}, \src{r1}, \src{r2}\imod{.b0}} (\textbf{SPIR-V:} \texttt{OpShiftLeftLogical (Shift)}; \textbf{GLSL:} \texttt{uint y = a \textless{}\textless{} s;})
  \item \texttt{encoded\_texture\_table\_index}: Encoded descriptor table index for \textbackslash\{\}texttt\{encoded\textbackslash\{\}\_texture\textbackslash\{\}\_table\textbackslash\{\}\_index\}. \emph{Example:} {\scriptsize\ttfamily \mnem{LD\_BUFFER.i32} \dst{@r0}, \src{r2}, \imm{\#0}} (\textbf{SPIR-V:} \texttt{OpAccessChain (DescriptorSet)}; \textbf{GLSL:} \texttt{uint v = buf[i];})
  \item \texttt{texture\_index}: Resource/slot index field \textbackslash\{\}texttt\{texture\textbackslash\{\}\_index\} selecting the target descriptor or slot. \emph{Example:} {\scriptsize\ttfamily \mnem{LD\_VAR\_BUF\_IMM.f32} \dst{@r0:r3}, \imm{\#0}} (\textbf{SPIR-V:} \texttt{OpLoad (Descriptor)}; \textbf{GLSL:} \texttt{vec4 v = v\_in;})
  \item \hyperref[enum:register_file_format_general_m]{\texttt{register\_format}}: Target GPR data conversion format (\texttt{auto32}, \texttt{f32}, \texttt{f16}, \texttt{s32}, \texttt{s16}, \texttt{u32}, \texttt{u16}). \emph{Example:} {\scriptsize\ttfamily \mnem{LD\_TILE}\imod{.f16.v4} \dst{@r0:r1}, \src{r2}} (\textbf{SPIR-V:} \texttt{OpImageRead (SubpassData)}; \textbf{GLSL:} \texttt{f16vec4 c = subpassLoad(inputAtt);})
  \item \hyperref[enum:vecsize_varying_m]{\texttt{vecsize}}: Component count (\texttt{none}=1, \texttt{v2}, \texttt{v3}, \texttt{v4}) for attribute, varying, or blend transfers. \emph{Example:} {\scriptsize\ttfamily \mnem{LD\_ATTR}\imod{.f32.v3} \dst{@r0:r2}, \src{r4}} (\textbf{SPIR-V:} \texttt{OpLoad (vec3 Input)}; \textbf{GLSL:} \texttt{in vec3 a\_normal;})
  \item \hyperref[enum:message_slot_index_m]{\texttt{message\_slot\_index}}: Asynchronous scoreboard slot (\texttt{slot0}--\texttt{slot2}, \texttt{slot7}) tracking completion of this message. \emph{Example:} {\scriptsize\ttfamily \mnem{LOAD.i32}\imod{.slot0} \dst{@r0}, \src{r2}, \imm{\#0}} (\textbf{SPIR-V:} \texttt{OpLoad}; \textbf{GLSL:} \texttt{uint v = ssbo.arr[i];})
  \item \texttt{sr\_count}: Number of consecutive 32-bit staging registers transferred by a load/store/texture message. \emph{Example:} {\scriptsize\ttfamily \mnem{LD\_MULTI.i128} \dst{@r0:r3}, \src{r4}, \imm{\#0}} (\textbf{SPIR-V:} \texttt{OpLoad (wide vector)}; \textbf{GLSL:} \texttt{uvec4 x = buf[i];})
  \item \texttt{descriptor\_type}: Selects descriptor table binding class (attribute, buffer, texture, varying). \emph{Example:} {\scriptsize\ttfamily \mnem{LD\_BUFFER.i32} \dst{@r0}, \src{r2}, \imm{\#0}} (\textbf{SPIR-V:} \texttt{OpLoad (StorageBuffer)}; \textbf{GLSL:} \texttt{uint v = ssbo.arr[i];})
  \item \texttt{staging\_register\_index}: Base GPR index of the contiguous staging register window (\texttt{@rN:rM}) for message payload/return. \emph{Example:} {\scriptsize\ttfamily \mnem{STORE.i64}\imod{.slot7} \src{@r0:r1}, \src{r2}, \imm{\#8}} (\textbf{SPIR-V:} \texttt{OpStore (uvec2)}; \textbf{GLSL:} \texttt{ssbo.pairs[i] = uvec2(a, b);})
  \item \texttt{staging\_register\_control\_32}: Controls the read/write staging register bank and count for async message instructions. \emph{Example:} {\scriptsize\ttfamily \mnem{LOAD.i128}\imod{.slot0} \dst{@r0:r3}, \src{r4}, \imm{\#0}} (\textbf{SPIR-V:} \texttt{OpLoad (vec4)}; \textbf{GLSL:} \texttt{vec4 v = ssbo.data[idx];})
  \item \texttt{opcode1}: Primary sub-opcode within the instruction class uniquely identifying the operation family. \emph{Example:} {\scriptsize\ttfamily \mnem{IADD.s32} \dst{r0}, \src{r1}, \src{r2}} (\textbf{SPIR-V:} \texttt{OpIAdd (operation selector)}; \textbf{GLSL:} \texttt{// Primary opcode bits})
  \item \texttt{instruction\_class}: 4-bit major instruction class selector (\texttt{bits 48:51}) partitioning the 64-bit Valhall opcode space (e.g., \texttt{FMA}, \texttt{CVT}, \texttt{SFU}, \texttt{LS}, \texttt{TEX}, \texttt{BRANCH}). \emph{Example:} {\scriptsize\ttfamily \mnem{FADD.f32} \dst{r0}, \src{r1}, \src{r2}} (\textbf{SPIR-V:} \texttt{N/A (encoding class)}; \textbf{GLSL:} \texttt{// Hardware execution pipe / decode class})
  \item \texttt{fau\_page\_index}: Selects the FAU (Fast Access Uniform) page (\texttt{0..3}) when a source operand references hardware small constants, thread IDs, or lane constants. \emph{Example:} {\scriptsize\ttfamily \mnem{FADD.f32} \dst{r0}, \src{r1}, \src{k1f}} (\textbf{SPIR-V:} \texttt{OpConstant / BuiltIn}; \textbf{GLSL:} \texttt{float y = x + 1.0;})
  \item \hyperref[enum:flow_control_m]{\texttt{flow}}: 4-bit embedded flow-control \& scoreboard wait modifier (\texttt{bits 59:62}: \texttt{none}, \texttt{wait0..wait0126}, \texttt{wait}, \texttt{reconverge}, \texttt{discard}, \texttt{end}). \emph{Example:} {\scriptsize\ttfamily \mnem{FADD.f32}\imod{.wait0.reconverge} \dst{r0}, \src{r1}, \src{r2}} (\textbf{SPIR-V:} \texttt{OpControlBarrier / OpReturn}; \textbf{GLSL:} \texttt{// Scoreboard wait + warp reconvergence})
\end{tightitemize}
\clearpage

\subsection*{LD\_TILE}
\addcontentsline{toc}{subsection}{LD\_TILE}
Loads pixel data from tilebuffer at \texttt{(pixel=src0, coverage=src1, conversion=src2)} into \texttt{sr\_dst}; \texttt{z\_stencil} loads depth/stencil instead of color.
\par\textbf{Notes:}
\begin{tightitemize}
  \item Used for framebuffer fetch and blend shaders (\texttt{ISA.xml:87}).
  \item \textbf{Source references:} \texttt{kraid/ops.rs:2877-2911}, \texttt{bifrost/valhall/ISA.xml:87}.
\end{tightitemize}
\codebox{\mnem{LD\_TILE}\imod{\{.message\_slot\_index\}}\imod{\{.flow\}}\imod{\{.z\_stencil\}}\imod{\{.vecsize\}}\imod{\{.register\_format\}} \dst{@sw}, \src{src0}, \src{src1}, \src{src2}\\\mnem{LD\_TILE}\imod{\{.message\_slot\_index\}}\imod{\{.flow\}}\imod{\{.z\_stencil\}}\imod{\{.z\_last\_read\}}\imod{\{.vecsize\}}\imod{\{.register\_format\}} \dst{@sw}, \src{src0}, \src{src1}, \src{src2}}
\textbf{Operands:}
\begin{tightitemize}\footnotesize
  \item \texttt{\dst{@sw}} (\texttt{staging\_register\_index} + \texttt{sr\_count}): Destination staging register range receiving pixel data read from on-chip tile SRAM.
  \item \texttt{\src{src0}}: Pixel coordinate / sample selector.
  \item \texttt{\src{src1}}: Sample coverage mask.
  \item \texttt{\src{src2}}: Tilebuffer conversion descriptor / render-target selector.
\end{tightitemize}
\par\textbf{Closest SPIR-V Op:} \texttt{OpImageRead (subpassInput / framebuffer fetch)}
\par\textbf{Pseudo-GLSL (Operational Semantics):}
\codebox{if (flag\_z\_stencil) @sw = readTileDepthStencilSRAM(pixel = src0, cov = src1);\\else                @sw = decodeTileColorSRAM(cvt = src2, pixel = src0, cov = src1);}
\textbf{Execution Context:}
\begin{tightitemize}\footnotesize
  \item \textbf{On-chip 16$\times$16 Tile SRAM (\hyperref[sec:csf-fbd]{\texttt{FRAGMENTSR::FBD\_POINTER}}, \texttt{r40:r41}):} 64B-aligned VA (\texttt{\& \textasciitilde{}0x3F}, \texttt{type[0]|zs\_ext[1]|rt\_cnt-1[5:2]}) $\to$ \texttt{FramebufferParameters} (128B) + optional \texttt{ZSCRCExtension} (64B, for \texttt{.z\_stencil}) + \texttt{RenderTarget[rt]} (64B: \texttt{internal\_format}, \texttt{clear\_color}, Morton 16$\times$16 \texttt{writeback\_buffer}).
  \item \textbf{Fragment stage preloads \& FP modes (\hyperref[sec:csf-spd]{\texttt{SPD\_2}}, \texttt{r20:r21}):} Hardware pre-fills \texttt{r59} with \texttt{floatBitsToUint(FragCoord.z)} and \texttt{r60} with the MSAA sample coverage mask (\texttt{gl\_SampleMaskIn[0]}), while \texttt{ShaderProgram} word~0 latches 32/64 GPR allocation and warp-wide \texttt{flush\_to\_zero\_mode}, \texttt{suppress\_na\_n}, and \texttt{suppress\_inf}.
\end{tightitemize}
\par\textbf{Worked Examples:}
\par\textbf{Example 1 (Color subpass / framebuffer fetch):}
\codebox{\mnem{LD\_TILE}\imod{.f32.v4.slot0} \dst{@r0:r3}, \src{r59}, \src{r60}, \src{u0}}
Reads on-chip tilebuffer color at \texttt{(pixel=r59, coverage=r60, conv\_desc=u0)} into 4x FP32 \texttt{@r0:r3}.
\par\textbf{Example 2 (Depth/stencil tilebuffer read (`.z\_stencil.last\_read`)):}
\codebox{\mnem{LD\_TILE}\imod{.f32.z\_stencil.last\_read.slot0} \dst{@r0}, \src{r59}, \src{r60}, \src{u1}}
Reads depth from the Z/S tilebuffer and marks it \texttt{.last\_read} so tile storage can be reclaimed.

\begin{center}\footnotesize
\begin{tabularx}{\linewidth}{@{}lllX@{}}
\toprule
Variant & Arch & Class & Opcode Fields \\
\midrule
\texttt{(none)} & v9--v14 & \texttt{S3} & \texttt{opcode1=LD\_TILE} \\
\texttt{(none)} & v9--v14 & \texttt{S3} & \texttt{opcode1=LD\_TILE} \\
\bottomrule
\end{tabularx}
\end{center}
\begin{center}\footnotesize
\begin{tabularx}{0.88\linewidth}{@{}l l X@{}}
\toprule
Field Name & Bits & Type / Fixed Value \\
\midrule
\texttt{src0} & \texttt{0:7} & \texttt{SourceEncoding} \\
\texttt{src1} & \texttt{8:15} & \texttt{SourceEncoding} \\
\texttt{src2} & \texttt{16:23} & \texttt{SourceEncoding} \\
\hyperref[enum:register_file_format_general_m]{\texttt{register\_format}} & \texttt{24:26} & \texttt{register\_file\_format\_general\_m} \\
\textit{Reserved} & \texttt{27} & \texttt{zero} \\
\hyperref[enum:vecsize_varying_m]{\texttt{vecsize}} & \texttt{28:29} & \texttt{vecsize\_varying\_m (def none)} \\
\hyperref[enum:message_slot_index_m]{\texttt{message\_slot\_index}} & \texttt{30:32} & \texttt{message\_slot\_index\_m} \\
\texttt{sr\_count} & \texttt{33:35} & \texttt{sr\_count\_t} \\
\hyperref[enum:ls_z_stencil_m]{\texttt{z\_stencil}} & \texttt{36} & \texttt{ls\_z\_stencil\_m} \\
\textit{Reserved} & \texttt{37:39} & \texttt{zero} \\
\texttt{staging\_register\_index} & \texttt{40:45} & \texttt{uint} \\
\texttt{staging\_register\_control\_32} & \texttt{46:47} & \texttt{exact=write} \\
\texttt{opcode1} & \texttt{48:51} & \texttt{exact=LD\_TILE} \\
\texttt{instruction\_class} & \texttt{52:56} & \texttt{exact=S3} \\
\texttt{fau\_page\_index} & \texttt{57:58} & \texttt{uint} \\
\hyperref[enum:flow_control_m]{\texttt{flow}} & \texttt{59:62} & \texttt{flow\_control\_m} \\
\textit{Reserved} & \texttt{63} & \texttt{zero} \\
\bottomrule
\end{tabularx}
\end{center}
\textbf{Fields Explained:}
\begin{tightitemize}\footnotesize
  \item \texttt{src0}: Primary source operand (\texttt{r0..r63}, \texttt{u0..u63} uniform, \texttt{k*} FAU constant, or \texttt{t0..t2} temporary; \texttt{\textasciicircum{}} marks discard/last-use). \emph{Example:} {\scriptsize\ttfamily \mnem{FADD.f32} \dst{r0}, \src{r1}, \src{r2}} (\textbf{SPIR-V:} \texttt{OpFAdd (Operand 0)}; \textbf{GLSL:} \texttt{float c = a + b;})
  \item \texttt{src1}: Secondary source operand (register, uniform, or FAU constant; in shifts, supplies the shift count byte). \emph{Example:} {\scriptsize\ttfamily \mnem{LSHIFT.i32} \dst{r0}, \src{r1}, \src{r2}\imod{.b0}} (\textbf{SPIR-V:} \texttt{OpShiftLeftLogical (Shift)}; \textbf{GLSL:} \texttt{uint y = a \textless{}\textless{} s;})
  \item \texttt{src2}: Third source operand (FMA addend, \texttt{MUX}/\texttt{CSEL} selector, \texttt{IADDC} carry-in, or bitwise mask in fused shift-logic). \emph{Example:} {\scriptsize\ttfamily \mnem{FMA.f32} \dst{r0}, \src{r1}, \src{r2}, \src{r3}} (\textbf{SPIR-V:} \texttt{GLSL.std.450 Fma (Addend)}; \textbf{GLSL:} \texttt{float d = fma(a, b, c);})
  \item \hyperref[enum:register_file_format_general_m]{\texttt{register\_format}}: Target GPR data conversion format (\texttt{auto32}, \texttt{f32}, \texttt{f16}, \texttt{s32}, \texttt{s16}, \texttt{u32}, \texttt{u16}). \emph{Example:} {\scriptsize\ttfamily \mnem{LD\_TILE}\imod{.f16.v4} \dst{@r0:r1}, \src{r2}} (\textbf{SPIR-V:} \texttt{OpImageRead (SubpassData)}; \textbf{GLSL:} \texttt{f16vec4 c = subpassLoad(inputAtt);})
  \item \hyperref[enum:vecsize_varying_m]{\texttt{vecsize}}: Component count (\texttt{none}=1, \texttt{v2}, \texttt{v3}, \texttt{v4}) for attribute, varying, or blend transfers. \emph{Example:} {\scriptsize\ttfamily \mnem{LD\_ATTR}\imod{.f32.v3} \dst{@r0:r2}, \src{r4}} (\textbf{SPIR-V:} \texttt{OpLoad (vec3 Input)}; \textbf{GLSL:} \texttt{in vec3 a\_normal;})
  \item \hyperref[enum:message_slot_index_m]{\texttt{message\_slot\_index}}: Asynchronous scoreboard slot (\texttt{slot0}--\texttt{slot2}, \texttt{slot7}) tracking completion of this message. \emph{Example:} {\scriptsize\ttfamily \mnem{LOAD.i32}\imod{.slot0} \dst{@r0}, \src{r2}, \imm{\#0}} (\textbf{SPIR-V:} \texttt{OpLoad}; \textbf{GLSL:} \texttt{uint v = ssbo.arr[i];})
  \item \texttt{sr\_count}: Number of consecutive 32-bit staging registers transferred by a load/store/texture message. \emph{Example:} {\scriptsize\ttfamily \mnem{LD\_MULTI.i128} \dst{@r0:r3}, \src{r4}, \imm{\#0}} (\textbf{SPIR-V:} \texttt{OpLoad (wide vector)}; \textbf{GLSL:} \texttt{uvec4 x = buf[i];})
  \item \hyperref[enum:ls_z_stencil_m]{\texttt{z\_stencil}}: Targets depth/stencil tile buffer instead of color attachment in \texttt{LD\_TILE}. \emph{Example:} {\scriptsize\ttfamily \mnem{LD\_TILE}\imod{.z\_stencil} \dst{@r0}, \src{r4}} (\textbf{SPIR-V:} \texttt{OpImageRead (Depth/Stencil subpass)}; \textbf{GLSL:} \texttt{float d = subpassLoad(depthAtt).x;})
  \item \texttt{staging\_register\_index}: Base GPR index of the contiguous staging register window (\texttt{@rN:rM}) for message payload/return. \emph{Example:} {\scriptsize\ttfamily \mnem{STORE.i64}\imod{.slot7} \src{@r0:r1}, \src{r2}, \imm{\#8}} (\textbf{SPIR-V:} \texttt{OpStore (uvec2)}; \textbf{GLSL:} \texttt{ssbo.pairs[i] = uvec2(a, b);})
  \item \texttt{staging\_register\_control\_32}: Controls the read/write staging register bank and count for async message instructions. \emph{Example:} {\scriptsize\ttfamily \mnem{LOAD.i128}\imod{.slot0} \dst{@r0:r3}, \src{r4}, \imm{\#0}} (\textbf{SPIR-V:} \texttt{OpLoad (vec4)}; \textbf{GLSL:} \texttt{vec4 v = ssbo.data[idx];})
  \item \texttt{opcode1}: Primary sub-opcode within the instruction class uniquely identifying the operation family. \emph{Example:} {\scriptsize\ttfamily \mnem{IADD.s32} \dst{r0}, \src{r1}, \src{r2}} (\textbf{SPIR-V:} \texttt{OpIAdd (operation selector)}; \textbf{GLSL:} \texttt{// Primary opcode bits})
  \item \texttt{instruction\_class}: 4-bit major instruction class selector (\texttt{bits 48:51}) partitioning the 64-bit Valhall opcode space (e.g., \texttt{FMA}, \texttt{CVT}, \texttt{SFU}, \texttt{LS}, \texttt{TEX}, \texttt{BRANCH}). \emph{Example:} {\scriptsize\ttfamily \mnem{FADD.f32} \dst{r0}, \src{r1}, \src{r2}} (\textbf{SPIR-V:} \texttt{N/A (encoding class)}; \textbf{GLSL:} \texttt{// Hardware execution pipe / decode class})
  \item \texttt{fau\_page\_index}: Selects the FAU (Fast Access Uniform) page (\texttt{0..3}) when a source operand references hardware small constants, thread IDs, or lane constants. \emph{Example:} {\scriptsize\ttfamily \mnem{FADD.f32} \dst{r0}, \src{r1}, \src{k1f}} (\textbf{SPIR-V:} \texttt{OpConstant / BuiltIn}; \textbf{GLSL:} \texttt{float y = x + 1.0;})
  \item \hyperref[enum:flow_control_m]{\texttt{flow}}: 4-bit embedded flow-control \& scoreboard wait modifier (\texttt{bits 59:62}: \texttt{none}, \texttt{wait0..wait0126}, \texttt{wait}, \texttt{reconverge}, \texttt{discard}, \texttt{end}). \emph{Example:} {\scriptsize\ttfamily \mnem{FADD.f32}\imod{.wait0.reconverge} \dst{r0}, \src{r1}, \src{r2}} (\textbf{SPIR-V:} \texttt{OpControlBarrier / OpReturn}; \textbf{GLSL:} \texttt{// Scoreboard wait + warp reconvergence})
\end{tightitemize}
\clearpage

\subsection*{LD\_VAR}
\addcontentsline{toc}{subsection}{LD\_VAR}
Interpolates varying via descriptor at dynamic handle (\texttt{src1}) at \texttt{sample\_position} (\texttt{center}/\texttt{centroid}/\texttt{sample}/\texttt{explicit}), updating hidden barycentric cache per \texttt{update} (\texttt{store}/\texttt{retrieve}/\texttt{clobber}/\texttt{none}).
\par\textbf{Notes:}
\begin{tightitemize}
  \item Software varying buffer interpolation (\texttt{ops.rs:2934-2999}).
  \item \textbf{Source references:} \texttt{kraid/ops.rs:2914-2999}, \texttt{bifrost/valhall/ISA.xml:89}.
\end{tightitemize}
\codebox{\mnem{LD\_VAR}\imod{\{.message\_slot\_index\}}\imod{\{.flow\}}\imod{\{.vecsize\}}\imod{\{.register\_format\}}\imod{\{.sample\_position\}}\imod{\{.update\}} \dst{@sw}, \src{src0}, \src{src1}}
\textbf{Operands:}
\begin{tightitemize}\footnotesize
  \item \texttt{\dst{@sw}} (\texttt{staging\_register\_index} + \texttt{sr\_count}): Destination staging register range receiving interpolated varying components.
  \item \texttt{\src{src0}}: Fragment sample/pixel coordinate (\texttt{r61}, or explicit offset for \texttt{interpolateAtOffset}).
  \item \texttt{\src{src1}}: Dynamic varying descriptor handle.
\end{tightitemize}
\par\textbf{Closest SPIR-V Op:} \texttt{OpLoad (interpolated fragment input)}
\par\textbf{Pseudo-GLSL (Operational Semantics):}
\codebox{@sw = interpolateVarying(handle = src1, pixel = src0, sample\_position, update\_mode);}
\textbf{Execution Context:}
\begin{tightitemize}\footnotesize
  \item \textbf{Varying buffer \& barycentric cache (\hyperref[sec:csf-tiler]{\texttt{IDVSSR::VARY\_SIZE}}, \texttt{r48:r49}):} Reads per-vertex varying RAM written by Varying EBIN \texttt{STORE.estream}; descriptor forms (\texttt{LD\_VAR*}) also chase \hyperref[sec:csf-srt]{\texttt{SRT\_2}} $\to$ \texttt{Resource[table]} $\to$ \texttt{Attribute[index]} and reuse cached barycentric weights via \texttt{.store}/\texttt{.retrieve}.
  \item \textbf{Fragment stage preloads \& FP modes (\hyperref[sec:csf-spd]{\texttt{SPD\_2}}, \texttt{r20:r21}):} Hardware pre-fills \texttt{r59} with \texttt{floatBitsToUint(FragCoord.z)} and \texttt{r60} with the MSAA sample coverage mask (\texttt{gl\_SampleMaskIn[0]}), while \texttt{ShaderProgram} word~0 latches 32/64 GPR allocation and warp-wide \texttt{flush\_to\_zero\_mode}, \texttt{suppress\_na\_n}, and \texttt{suppress\_inf}.
\end{tightitemize}
\par\textbf{Worked Examples:}
\par\textbf{Example 1 (Center interpolation saving barycentrics (`.store`)):}
\codebox{\mnem{LD\_VAR}\imod{.f32.v4.center.store.slot0} \dst{@r0:r3}, \src{r59}, \src{r4}}
Interpolates varying handle \texttt{r4} at pixel center and caches barycentric weights in hardware (\texttt{.store}).
\par\textbf{Example 2 (Centroid interpolation reusing cache (`.retrieve`)):}
\codebox{\mnem{LD\_VAR}\imod{.f32.v4.centroid.retrieve.slot0} \dst{@r4:r7}, \src{r59}, \src{r5}}
Interpolates a second varying at centroid reusing cached barycentric weights (\texttt{.retrieve}).

\begin{center}\footnotesize
\begin{tabularx}{\linewidth}{@{}lllX@{}}
\toprule
Variant & Arch & Class & Opcode Fields \\
\midrule
\texttt{(none)} & v9--v14 & \texttt{S2} & \texttt{opcode1=LD\_VAR} \\
\bottomrule
\end{tabularx}
\end{center}
\begin{center}\footnotesize
\begin{tabularx}{0.88\linewidth}{@{}l l X@{}}
\toprule
Field Name & Bits & Type / Fixed Value \\
\midrule
\texttt{src0} & \texttt{0:7} & \texttt{SourceEncoding} \\
\texttt{src1} & \texttt{8:15} & \texttt{SourceEncoding} \\
\textit{Reserved} & \texttt{16:23} & \texttt{zero} \\
\hyperref[enum:register_file_format_general_m]{\texttt{register\_format}} & \texttt{24:26} & \texttt{register\_file\_format\_general\_m} \\
\textit{Reserved} & \texttt{27} & \texttt{zero} \\
\hyperref[enum:vecsize_varying_m]{\texttt{vecsize}} & \texttt{28:29} & \texttt{vecsize\_varying\_m (def none)} \\
\hyperref[enum:message_slot_index_m]{\texttt{message\_slot\_index}} & \texttt{30:32} & \texttt{message\_slot\_index\_m} \\
\texttt{sr\_count} & \texttt{33:35} & \texttt{sr\_count\_t} \\
\hyperref[enum:update_mode_m]{\texttt{update}} & \texttt{36:37} & \texttt{update\_mode\_m} \\
\hyperref[enum:encoded_sample_position_none_m]{\texttt{encoded\_sample}} & \texttt{38:39} & \texttt{encoded\_sample\_position\_none\_m} \\
\texttt{staging\_register\_index} & \texttt{40:45} & \texttt{uint} \\
\texttt{staging\_register\_control\_32} & \texttt{46:47} & \texttt{exact=write} \\
\texttt{opcode1} & \texttt{48:51} & \texttt{exact=LD\_VAR} \\
\texttt{instruction\_class} & \texttt{52:56} & \texttt{exact=S2} \\
\texttt{fau\_page\_index} & \texttt{57:58} & \texttt{uint} \\
\hyperref[enum:flow_control_m]{\texttt{flow}} & \texttt{59:62} & \texttt{flow\_control\_m} \\
\textit{Reserved} & \texttt{63} & \texttt{zero} \\
\bottomrule
\end{tabularx}
\end{center}
\textbf{Fields Explained:}
\begin{tightitemize}\footnotesize
  \item \texttt{src0}: Primary source operand (\texttt{r0..r63}, \texttt{u0..u63} uniform, \texttt{k*} FAU constant, or \texttt{t0..t2} temporary; \texttt{\textasciicircum{}} marks discard/last-use). \emph{Example:} {\scriptsize\ttfamily \mnem{FADD.f32} \dst{r0}, \src{r1}, \src{r2}} (\textbf{SPIR-V:} \texttt{OpFAdd (Operand 0)}; \textbf{GLSL:} \texttt{float c = a + b;})
  \item \texttt{src1}: Secondary source operand (register, uniform, or FAU constant; in shifts, supplies the shift count byte). \emph{Example:} {\scriptsize\ttfamily \mnem{LSHIFT.i32} \dst{r0}, \src{r1}, \src{r2}\imod{.b0}} (\textbf{SPIR-V:} \texttt{OpShiftLeftLogical (Shift)}; \textbf{GLSL:} \texttt{uint y = a \textless{}\textless{} s;})
  \item \hyperref[enum:register_file_format_general_m]{\texttt{register\_format}}: Target GPR data conversion format (\texttt{auto32}, \texttt{f32}, \texttt{f16}, \texttt{s32}, \texttt{s16}, \texttt{u32}, \texttt{u16}). \emph{Example:} {\scriptsize\ttfamily \mnem{LD\_TILE}\imod{.f16.v4} \dst{@r0:r1}, \src{r2}} (\textbf{SPIR-V:} \texttt{OpImageRead (SubpassData)}; \textbf{GLSL:} \texttt{f16vec4 c = subpassLoad(inputAtt);})
  \item \hyperref[enum:vecsize_varying_m]{\texttt{vecsize}}: Component count (\texttt{none}=1, \texttt{v2}, \texttt{v3}, \texttt{v4}) for attribute, varying, or blend transfers. \emph{Example:} {\scriptsize\ttfamily \mnem{LD\_ATTR}\imod{.f32.v3} \dst{@r0:r2}, \src{r4}} (\textbf{SPIR-V:} \texttt{OpLoad (vec3 Input)}; \textbf{GLSL:} \texttt{in vec3 a\_normal;})
  \item \hyperref[enum:message_slot_index_m]{\texttt{message\_slot\_index}}: Asynchronous scoreboard slot (\texttt{slot0}--\texttt{slot2}, \texttt{slot7}) tracking completion of this message. \emph{Example:} {\scriptsize\ttfamily \mnem{LOAD.i32}\imod{.slot0} \dst{@r0}, \src{r2}, \imm{\#0}} (\textbf{SPIR-V:} \texttt{OpLoad}; \textbf{GLSL:} \texttt{uint v = ssbo.arr[i];})
  \item \texttt{sr\_count}: Number of consecutive 32-bit staging registers transferred by a load/store/texture message. \emph{Example:} {\scriptsize\ttfamily \mnem{LD\_MULTI.i128} \dst{@r0:r3}, \src{r4}, \imm{\#0}} (\textbf{SPIR-V:} \texttt{OpLoad (wide vector)}; \textbf{GLSL:} \texttt{uvec4 x = buf[i];})
  \item \hyperref[enum:update_mode_m]{\texttt{update}}: Controls hardware varying-cache lifecycle (\texttt{store}, \texttt{retrieve}, \texttt{clobber}, \texttt{none}) to avoid redundant barycentric math. \emph{Example:} {\scriptsize\ttfamily \mnem{LD\_VAR\_BUF.f32}\imod{.retrieve.v4} \dst{@r0:r3}, \src{r4}} (\textbf{SPIR-V:} \texttt{OpLoad (Input varying)}; \textbf{GLSL:} \texttt{vec4 a = v\_pos; vec4 b = v\_norm;})
  \item \hyperref[enum:encoded_sample_position_none_m]{\texttt{encoded\_sample}}: Compressed encoding of MSAA sample interpolation position (\texttt{center\_or\_none}, \texttt{centroid}, \texttt{sample}, \texttt{explicit}). \emph{Example:} {\scriptsize\ttfamily \mnem{LD\_VAR.f32}\imod{.sample.v4} \dst{@r0:r3}, \src{r4}} (\textbf{SPIR-V:} \texttt{OpInterpolateAtSample}; \textbf{GLSL:} \texttt{vec4 v = interpolateAtSample(v\_col, 0);})
  \item \texttt{staging\_register\_index}: Base GPR index of the contiguous staging register window (\texttt{@rN:rM}) for message payload/return. \emph{Example:} {\scriptsize\ttfamily \mnem{STORE.i64}\imod{.slot7} \src{@r0:r1}, \src{r2}, \imm{\#8}} (\textbf{SPIR-V:} \texttt{OpStore (uvec2)}; \textbf{GLSL:} \texttt{ssbo.pairs[i] = uvec2(a, b);})
  \item \texttt{staging\_register\_control\_32}: Controls the read/write staging register bank and count for async message instructions. \emph{Example:} {\scriptsize\ttfamily \mnem{LOAD.i128}\imod{.slot0} \dst{@r0:r3}, \src{r4}, \imm{\#0}} (\textbf{SPIR-V:} \texttt{OpLoad (vec4)}; \textbf{GLSL:} \texttt{vec4 v = ssbo.data[idx];})
  \item \texttt{opcode1}: Primary sub-opcode within the instruction class uniquely identifying the operation family. \emph{Example:} {\scriptsize\ttfamily \mnem{IADD.s32} \dst{r0}, \src{r1}, \src{r2}} (\textbf{SPIR-V:} \texttt{OpIAdd (operation selector)}; \textbf{GLSL:} \texttt{// Primary opcode bits})
  \item \texttt{instruction\_class}: 4-bit major instruction class selector (\texttt{bits 48:51}) partitioning the 64-bit Valhall opcode space (e.g., \texttt{FMA}, \texttt{CVT}, \texttt{SFU}, \texttt{LS}, \texttt{TEX}, \texttt{BRANCH}). \emph{Example:} {\scriptsize\ttfamily \mnem{FADD.f32} \dst{r0}, \src{r1}, \src{r2}} (\textbf{SPIR-V:} \texttt{N/A (encoding class)}; \textbf{GLSL:} \texttt{// Hardware execution pipe / decode class})
  \item \texttt{fau\_page\_index}: Selects the FAU (Fast Access Uniform) page (\texttt{0..3}) when a source operand references hardware small constants, thread IDs, or lane constants. \emph{Example:} {\scriptsize\ttfamily \mnem{FADD.f32} \dst{r0}, \src{r1}, \src{k1f}} (\textbf{SPIR-V:} \texttt{OpConstant / BuiltIn}; \textbf{GLSL:} \texttt{float y = x + 1.0;})
  \item \hyperref[enum:flow_control_m]{\texttt{flow}}: 4-bit embedded flow-control \& scoreboard wait modifier (\texttt{bits 59:62}: \texttt{none}, \texttt{wait0..wait0126}, \texttt{wait}, \texttt{reconverge}, \texttt{discard}, \texttt{end}). \emph{Example:} {\scriptsize\ttfamily \mnem{FADD.f32}\imod{.wait0.reconverge} \dst{r0}, \src{r1}, \src{r2}} (\textbf{SPIR-V:} \texttt{OpControlBarrier / OpReturn}; \textbf{GLSL:} \texttt{// Scoreboard wait + warp reconvergence})
\end{tightitemize}
\clearpage

\subsection*{LD\_VAR\_BUF}
\addcontentsline{toc}{subsection}{LD\_VAR\_BUF}
Interpolates varying directly from hardware varying buffer at dynamic byte offset (\texttt{src1}) with \texttt{source\_format} (\texttt{f32}/\texttt{f16}, v10+) conversion to \texttt{register\_format}.
\par\textbf{Notes:}
\begin{tightitemize}
  \item Descriptor-less varying interpolation (\texttt{ops.rs:3001-3042}).
  \item \textbf{Source references:} \texttt{kraid/ops.rs:3001-3042}, \texttt{isa-v9-v14.xml}.
\end{tightitemize}
\codebox{\mnem{LD\_VAR\_BUF}\imod{\{.message\_slot\_index\}}\imod{\{.flow\}}\imod{\{.vecsize\}}\imod{\{.source\_format\}}\imod{\{.sample\_position\}}\imod{\{.update\}} \dst{@sw}, \src{src0}, \src{src1}}
\textbf{Operands:}
\begin{tightitemize}\footnotesize
  \item \texttt{\dst{@sw}} (\texttt{staging\_register\_index} + \texttt{sr\_count}): Destination staging register range receiving interpolated varying components from the varying buffer.
  \item \texttt{\src{src0}}: Fragment sample/pixel coordinate (\texttt{r61}).
  \item \texttt{\src{src1}}: Dynamic byte offset within the descriptor-less varying buffer.
\end{tightitemize}
\par\textbf{Closest SPIR-V Op:} \texttt{OpLoad (hardware varying buffer)}
\par\textbf{Pseudo-GLSL (Operational Semantics):}
\codebox{@sw = interpolateVaryingBuf(byte\_off = src1, pixel = src0, source\_format, sample\_position, update\_mode);}
\textbf{Execution Context:}
\begin{tightitemize}\footnotesize
  \item \textbf{Varying buffer \& barycentric cache (\hyperref[sec:csf-tiler]{\texttt{IDVSSR::VARY\_SIZE}}, \texttt{r48:r49}):} Reads per-vertex varying RAM written by Varying EBIN \texttt{STORE.estream}; descriptor forms (\texttt{LD\_VAR*}) also chase \hyperref[sec:csf-srt]{\texttt{SRT\_2}} $\to$ \texttt{Resource[table]} $\to$ \texttt{Attribute[index]} and reuse cached barycentric weights via \texttt{.store}/\texttt{.retrieve}.
  \item \textbf{Fragment stage preloads \& FP modes (\hyperref[sec:csf-spd]{\texttt{SPD\_2}}, \texttt{r20:r21}):} Hardware pre-fills \texttt{r59} with \texttt{floatBitsToUint(FragCoord.z)} and \texttt{r60} with the MSAA sample coverage mask (\texttt{gl\_SampleMaskIn[0]}), while \texttt{ShaderProgram} word~0 latches 32/64 GPR allocation and warp-wide \texttt{flush\_to\_zero\_mode}, \texttt{suppress\_na\_n}, and \texttt{suppress\_inf}.
\end{tightitemize}
\par\textbf{Worked Examples:}
\par\textbf{Example 1 (Descriptor-less varying buffer interpolation):}
\codebox{\mnem{LD\_VAR\_BUF.f32}\imod{.v4.center.store.slot0} \dst{@r0:r3}, \src{r59}, \src{r4}}
Interpolates 4x FP32 varying at dynamic byte offset \texttt{r4} in the hardware varying buffer and stores barycentrics.
\par\textbf{Example 2 (FP16-stored varying widened to FP32 (`.src\_f16`)):}
\codebox{\mnem{LD\_VAR\_BUF.f32}\imod{.src\_f16.v2.center.retrieve.slot0} \dst{@r0:r1}, \src{r59}, \src{r4}}
Reads \texttt{mediump} FP16 varying storage (\texttt{.src\_f16}) and delivers interpolated FP32 \texttt{vec2} in \texttt{@r0:r1}.

\begin{center}\footnotesize
\begin{tabularx}{\linewidth}{@{}lllX@{}}
\toprule
Variant & Arch & Class & Opcode Fields \\
\midrule
\texttt{(none)} & v9--v14 & \texttt{S2} & \texttt{opcode1=LD\_VAR\_BUF\_F16} \\
\texttt{(none)} & v9--v14 & \texttt{S2} & \texttt{opcode1=LD\_VAR\_BUF\_F16} \\
\texttt{(none)} & v9--v14 & \texttt{S2} & \texttt{opcode1=LD\_VAR\_BUF\_F32} \\
\texttt{(none)} & v9--v14 & \texttt{S2} & \texttt{opcode1=LD\_VAR\_BUF\_F32} \\
\bottomrule
\end{tabularx}
\end{center}
\begin{center}\footnotesize
\begin{tabularx}{0.88\linewidth}{@{}l l X@{}}
\toprule
Field Name & Bits & Type / Fixed Value \\
\midrule
\texttt{src0} & \texttt{0:7} & \texttt{SourceEncoding} \\
\texttt{src1} & \texttt{8:15} & \texttt{SourceEncoding} \\
\textit{Reserved} & \texttt{16:23} & \texttt{zero} \\
\hyperref[enum:ld_var_buf_source_format_m]{\texttt{source\_format}} & \texttt{24:25} & \texttt{ld\_var\_buf\_source\_format\_m} \\
\textit{Reserved} & \texttt{26:27} & \texttt{zero} \\
\hyperref[enum:vecsize_varying_m]{\texttt{vecsize}} & \texttt{28:29} & \texttt{vecsize\_varying\_m (def none)} \\
\hyperref[enum:message_slot_index_m]{\texttt{message\_slot\_index}} & \texttt{30:32} & \texttt{message\_slot\_index\_m} \\
\texttt{sr\_count} & \texttt{33:35} & \texttt{sr\_count\_t} \\
\hyperref[enum:ld_var_buf_encoded_update_mode_none_m]{\texttt{encoded\_update}} & \texttt{36:37} & \texttt{ld\_var\_buf\_encoded\_update\_mode\_none\_m} \\
\hyperref[enum:encoded_sample_position_none_m]{\texttt{encoded\_sample}} & \texttt{38:39} & \texttt{encoded\_sample\_position\_none\_m} \\
\texttt{staging\_register\_index} & \texttt{40:45} & \texttt{uint} \\
\texttt{staging\_register\_control\_32} & \texttt{46:47} & \texttt{exact=write} \\
\texttt{opcode1} & \texttt{48:51} & \texttt{exact=LD\_VAR\_BUF\_F16} \\
\texttt{instruction\_class} & \texttt{52:56} & \texttt{exact=S2} \\
\texttt{fau\_page\_index} & \texttt{57:58} & \texttt{uint} \\
\hyperref[enum:flow_control_m]{\texttt{flow}} & \texttt{59:62} & \texttt{flow\_control\_m} \\
\textit{Reserved} & \texttt{63} & \texttt{zero} \\
\bottomrule
\end{tabularx}
\end{center}
\textbf{Fields Explained:}
\begin{tightitemize}\footnotesize
  \item \texttt{src0}: Primary source operand (\texttt{r0..r63}, \texttt{u0..u63} uniform, \texttt{k*} FAU constant, or \texttt{t0..t2} temporary; \texttt{\textasciicircum{}} marks discard/last-use). \emph{Example:} {\scriptsize\ttfamily \mnem{FADD.f32} \dst{r0}, \src{r1}, \src{r2}} (\textbf{SPIR-V:} \texttt{OpFAdd (Operand 0)}; \textbf{GLSL:} \texttt{float c = a + b;})
  \item \texttt{src1}: Secondary source operand (register, uniform, or FAU constant; in shifts, supplies the shift count byte). \emph{Example:} {\scriptsize\ttfamily \mnem{LSHIFT.i32} \dst{r0}, \src{r1}, \src{r2}\imod{.b0}} (\textbf{SPIR-V:} \texttt{OpShiftLeftLogical (Shift)}; \textbf{GLSL:} \texttt{uint y = a \textless{}\textless{} s;})
  \item \hyperref[enum:ld_var_buf_source_format_m]{\texttt{source\_format}}: Declares the storage precision of the varying buffer (\texttt{src\_f32}, \texttt{src\_f16}, \texttt{src\_flat32}, \texttt{src\_flat16}). \emph{Example:} {\scriptsize\ttfamily \mnem{LD\_VAR\_BUF.f32}\imod{.src\_f16.v2} \dst{@r0:r1}, \src{r2}} (\textbf{SPIR-V:} \texttt{OpLoad (RelaxedPrecision Input)}; \textbf{GLSL:} \texttt{mediump in vec2 v\_uv;})
  \item \hyperref[enum:vecsize_varying_m]{\texttt{vecsize}}: Component count (\texttt{none}=1, \texttt{v2}, \texttt{v3}, \texttt{v4}) for attribute, varying, or blend transfers. \emph{Example:} {\scriptsize\ttfamily \mnem{LD\_ATTR}\imod{.f32.v3} \dst{@r0:r2}, \src{r4}} (\textbf{SPIR-V:} \texttt{OpLoad (vec3 Input)}; \textbf{GLSL:} \texttt{in vec3 a\_normal;})
  \item \hyperref[enum:message_slot_index_m]{\texttt{message\_slot\_index}}: Asynchronous scoreboard slot (\texttt{slot0}--\texttt{slot2}, \texttt{slot7}) tracking completion of this message. \emph{Example:} {\scriptsize\ttfamily \mnem{LOAD.i32}\imod{.slot0} \dst{@r0}, \src{r2}, \imm{\#0}} (\textbf{SPIR-V:} \texttt{OpLoad}; \textbf{GLSL:} \texttt{uint v = ssbo.arr[i];})
  \item \texttt{sr\_count}: Number of consecutive 32-bit staging registers transferred by a load/store/texture message. \emph{Example:} {\scriptsize\ttfamily \mnem{LD\_MULTI.i128} \dst{@r0:r3}, \src{r4}, \imm{\#0}} (\textbf{SPIR-V:} \texttt{OpLoad (wide vector)}; \textbf{GLSL:} \texttt{uvec4 x = buf[i];})
  \item \hyperref[enum:ld_var_buf_encoded_update_mode_none_m]{\texttt{encoded\_update}}: Compressed bitfield encoding for varying-cache update mode (\texttt{store}, \texttt{retrieve}, \texttt{clobber}). \emph{Example:} {\scriptsize\ttfamily \mnem{LD\_VAR\_SPECIAL}\imod{.bary.store} \dst{@r0:r1}, \src{r2}} (\textbf{SPIR-V:} \texttt{OpLoad (BuiltIn BaryCoordKHR)}; \textbf{GLSL:} \texttt{vec3 b = gl\_BaryCoordEXT;})
  \item \hyperref[enum:encoded_sample_position_none_m]{\texttt{encoded\_sample}}: Compressed encoding of MSAA sample interpolation position (\texttt{center\_or\_none}, \texttt{centroid}, \texttt{sample}, \texttt{explicit}). \emph{Example:} {\scriptsize\ttfamily \mnem{LD\_VAR.f32}\imod{.sample.v4} \dst{@r0:r3}, \src{r4}} (\textbf{SPIR-V:} \texttt{OpInterpolateAtSample}; \textbf{GLSL:} \texttt{vec4 v = interpolateAtSample(v\_col, 0);})
  \item \texttt{staging\_register\_index}: Base GPR index of the contiguous staging register window (\texttt{@rN:rM}) for message payload/return. \emph{Example:} {\scriptsize\ttfamily \mnem{STORE.i64}\imod{.slot7} \src{@r0:r1}, \src{r2}, \imm{\#8}} (\textbf{SPIR-V:} \texttt{OpStore (uvec2)}; \textbf{GLSL:} \texttt{ssbo.pairs[i] = uvec2(a, b);})
  \item \texttt{staging\_register\_control\_32}: Controls the read/write staging register bank and count for async message instructions. \emph{Example:} {\scriptsize\ttfamily \mnem{LOAD.i128}\imod{.slot0} \dst{@r0:r3}, \src{r4}, \imm{\#0}} (\textbf{SPIR-V:} \texttt{OpLoad (vec4)}; \textbf{GLSL:} \texttt{vec4 v = ssbo.data[idx];})
  \item \texttt{opcode1}: Primary sub-opcode within the instruction class uniquely identifying the operation family. \emph{Example:} {\scriptsize\ttfamily \mnem{IADD.s32} \dst{r0}, \src{r1}, \src{r2}} (\textbf{SPIR-V:} \texttt{OpIAdd (operation selector)}; \textbf{GLSL:} \texttt{// Primary opcode bits})
  \item \texttt{instruction\_class}: 4-bit major instruction class selector (\texttt{bits 48:51}) partitioning the 64-bit Valhall opcode space (e.g., \texttt{FMA}, \texttt{CVT}, \texttt{SFU}, \texttt{LS}, \texttt{TEX}, \texttt{BRANCH}). \emph{Example:} {\scriptsize\ttfamily \mnem{FADD.f32} \dst{r0}, \src{r1}, \src{r2}} (\textbf{SPIR-V:} \texttt{N/A (encoding class)}; \textbf{GLSL:} \texttt{// Hardware execution pipe / decode class})
  \item \texttt{fau\_page\_index}: Selects the FAU (Fast Access Uniform) page (\texttt{0..3}) when a source operand references hardware small constants, thread IDs, or lane constants. \emph{Example:} {\scriptsize\ttfamily \mnem{FADD.f32} \dst{r0}, \src{r1}, \src{k1f}} (\textbf{SPIR-V:} \texttt{OpConstant / BuiltIn}; \textbf{GLSL:} \texttt{float y = x + 1.0;})
  \item \hyperref[enum:flow_control_m]{\texttt{flow}}: 4-bit embedded flow-control \& scoreboard wait modifier (\texttt{bits 59:62}: \texttt{none}, \texttt{wait0..wait0126}, \texttt{wait}, \texttt{reconverge}, \texttt{discard}, \texttt{end}). \emph{Example:} {\scriptsize\ttfamily \mnem{FADD.f32}\imod{.wait0.reconverge} \dst{r0}, \src{r1}, \src{r2}} (\textbf{SPIR-V:} \texttt{OpControlBarrier / OpReturn}; \textbf{GLSL:} \texttt{// Scoreboard wait + warp reconvergence})
\end{tightitemize}
\clearpage

\subsection*{LD\_VAR\_BUF\_FLAT}
\addcontentsline{toc}{subsection}{LD\_VAR\_BUF\_FLAT}
Fetches flat varying directly from hardware varying buffer at dynamic offset without conversion (v14+ native; v9-v13 lowers to \texttt{LD\_VAR\_BUF}).
\par\textbf{Notes:}
\begin{tightitemize}
  \item Descriptor-less flat varying fetch (\texttt{ops.rs:3044-3071}).
  \item \textbf{Source references:} \texttt{kraid/ops.rs:3044-3071}, \texttt{bifrost/valhall/ISA.xml:91}.
\end{tightitemize}
\codebox{\mnem{LD\_VAR\_BUF\_FLAT}\imod{\{.message\_slot\_index\}}\imod{\{.flow\}}\imod{\{.vecsize\}}\imod{\{.register\_format\}}\imod{\{.vertex\_type\}} \dst{@sw}, \src{src0}}
\textbf{Operands:}
\begin{tightitemize}\footnotesize
  \item \texttt{\dst{@sw}} (\texttt{staging\_register\_index} + \texttt{sr\_count}): Destination staging register range receiving uninterpolated vertex data from the varying buffer.
  \item \texttt{\src{src0}}: Dynamic byte offset (and vertex index when \texttt{.explicit\_vertex} is set).
\end{tightitemize}
\par\textbf{Closest SPIR-V Op:} \texttt{OpLoad (flat varying buffer, v14+)}
\par\textbf{Pseudo-GLSL (Operational Semantics):}
\codebox{@sw = readVaryingBufFlat(off\_and\_vtx = src0, flag\_explicit\_vertex);}
\textbf{Execution Context:}
\begin{tightitemize}\footnotesize
  \item \textbf{Varying buffer \& barycentric cache (\hyperref[sec:csf-tiler]{\texttt{IDVSSR::VARY\_SIZE}}, \texttt{r48:r49}):} Reads per-vertex varying RAM written by Varying EBIN \texttt{STORE.estream}; descriptor forms (\texttt{LD\_VAR*}) also chase \hyperref[sec:csf-srt]{\texttt{SRT\_2}} $\to$ \texttt{Resource[table]} $\to$ \texttt{Attribute[index]} and reuse cached barycentric weights via \texttt{.store}/\texttt{.retrieve}.
  \item \textbf{Fragment stage preloads \& FP modes (\hyperref[sec:csf-spd]{\texttt{SPD\_2}}, \texttt{r20:r21}):} Hardware pre-fills \texttt{r59} with \texttt{floatBitsToUint(FragCoord.z)} and \texttt{r60} with the MSAA sample coverage mask (\texttt{gl\_SampleMaskIn[0]}), while \texttt{ShaderProgram} word~0 latches 32/64 GPR allocation and warp-wide \texttt{flush\_to\_zero\_mode}, \texttt{suppress\_na\_n}, and \texttt{suppress\_inf}.
\end{tightitemize}
\par\textbf{Worked Examples:}
\par\textbf{Example 1 (Provoking-vertex flat varying load (v14+)):}
\codebox{\mnem{LD\_VAR\_BUF\_FLAT}\imod{.u32.slot0} \dst{@r0}, \src{r4}}
Fetches uninterpolated 32-bit flat varying from the varying buffer at dynamic offset \texttt{r4}.
\par\textbf{Example 2 (Per-vertex explicit index (`.explicit\_vertex`)):}
\codebox{\mnem{LD\_VAR\_BUF\_FLAT}\imod{.f32.v4.explicit\_vertex.slot0} \dst{@r0:r3}, \src{r4}}
Implements \texttt{GL\_EXT\_fragment\_shader\_barycentric} \texttt{pervertexEXT} by selecting triangle vertex \texttt{0..2} dynamically.

\begin{center}\footnotesize
\begin{tabularx}{\linewidth}{@{}lllX@{}}
\toprule
Variant & Arch & Class & Opcode Fields \\
\midrule
\texttt{(none)} & v9--v14 & \texttt{S1} & \texttt{opcode1=LD\_VAR\_BUF\_FLAT} \\
\bottomrule
\end{tabularx}
\end{center}
\begin{center}\footnotesize
\begin{tabularx}{0.88\linewidth}{@{}l l X@{}}
\toprule
Field Name & Bits & Type / Fixed Value \\
\midrule
\texttt{src0} & \texttt{0:7} & \texttt{SourceEncoding} \\
\textit{Reserved} & \texttt{8:23} & \texttt{zero} \\
\hyperref[enum:ld_var_buf_flat_format_m]{\texttt{register\_format}} & \texttt{24:25} & \texttt{ld\_var\_buf\_flat\_format\_m} \\
\textit{Reserved} & \texttt{26:27} & \texttt{zero} \\
\hyperref[enum:vecsize_varying_m]{\texttt{vecsize}} & \texttt{28:29} & \texttt{vecsize\_varying\_m (def none)} \\
\hyperref[enum:message_slot_index_m]{\texttt{message\_slot\_index}} & \texttt{30:32} & \texttt{message\_slot\_index\_m} \\
\texttt{sr\_count} & \texttt{33:35} & \texttt{sr\_count\_t} \\
\hyperref[enum:vertex_type_m]{\texttt{vertex\_type}} & \texttt{36} & \texttt{vertex\_type\_m} \\
\textit{Reserved} & \texttt{37:39} & \texttt{zero} \\
\texttt{staging\_register\_index} & \texttt{40:45} & \texttt{uint} \\
\texttt{staging\_register\_control\_32} & \texttt{46:47} & \texttt{exact=write} \\
\texttt{opcode1} & \texttt{48:51} & \texttt{exact=LD\_VAR\_BUF\_FLAT} \\
\texttt{instruction\_class} & \texttt{52:56} & \texttt{exact=S1} \\
\texttt{fau\_page\_index} & \texttt{57:58} & \texttt{uint} \\
\hyperref[enum:flow_control_m]{\texttt{flow}} & \texttt{59:62} & \texttt{flow\_control\_m} \\
\textit{Reserved} & \texttt{63} & \texttt{zero} \\
\bottomrule
\end{tabularx}
\end{center}
\textbf{Fields Explained:}
\begin{tightitemize}\footnotesize
  \item \texttt{src0}: Primary source operand (\texttt{r0..r63}, \texttt{u0..u63} uniform, \texttt{k*} FAU constant, or \texttt{t0..t2} temporary; \texttt{\textasciicircum{}} marks discard/last-use). \emph{Example:} {\scriptsize\ttfamily \mnem{FADD.f32} \dst{r0}, \src{r1}, \src{r2}} (\textbf{SPIR-V:} \texttt{OpFAdd (Operand 0)}; \textbf{GLSL:} \texttt{float c = a + b;})
  \item \hyperref[enum:ld_var_buf_flat_format_m]{\texttt{register\_format}}: Target GPR data conversion format (\texttt{auto32}, \texttt{f32}, \texttt{f16}, \texttt{s32}, \texttt{s16}, \texttt{u32}, \texttt{u16}). \emph{Example:} {\scriptsize\ttfamily \mnem{LD\_TILE}\imod{.f16.v4} \dst{@r0:r1}, \src{r2}} (\textbf{SPIR-V:} \texttt{OpImageRead (SubpassData)}; \textbf{GLSL:} \texttt{f16vec4 c = subpassLoad(inputAtt);})
  \item \hyperref[enum:vecsize_varying_m]{\texttt{vecsize}}: Component count (\texttt{none}=1, \texttt{v2}, \texttt{v3}, \texttt{v4}) for attribute, varying, or blend transfers. \emph{Example:} {\scriptsize\ttfamily \mnem{LD\_ATTR}\imod{.f32.v3} \dst{@r0:r2}, \src{r4}} (\textbf{SPIR-V:} \texttt{OpLoad (vec3 Input)}; \textbf{GLSL:} \texttt{in vec3 a\_normal;})
  \item \hyperref[enum:message_slot_index_m]{\texttt{message\_slot\_index}}: Asynchronous scoreboard slot (\texttt{slot0}--\texttt{slot2}, \texttt{slot7}) tracking completion of this message. \emph{Example:} {\scriptsize\ttfamily \mnem{LOAD.i32}\imod{.slot0} \dst{@r0}, \src{r2}, \imm{\#0}} (\textbf{SPIR-V:} \texttt{OpLoad}; \textbf{GLSL:} \texttt{uint v = ssbo.arr[i];})
  \item \texttt{sr\_count}: Number of consecutive 32-bit staging registers transferred by a load/store/texture message. \emph{Example:} {\scriptsize\ttfamily \mnem{LD\_MULTI.i128} \dst{@r0:r3}, \src{r4}, \imm{\#0}} (\textbf{SPIR-V:} \texttt{OpLoad (wide vector)}; \textbf{GLSL:} \texttt{uvec4 x = buf[i];})
  \item \hyperref[enum:vertex_type_m]{\texttt{vertex\_type}}: Enables explicit register-supplied vertex index on \texttt{LD\_VAR\_BUF\_FLAT}. \emph{Example:} {\scriptsize\ttfamily \mnem{LD\_VAR\_BUF\_FLAT.f32}\imod{.explicit\_vertex} \dst{@r0:r3}, \src{r4}} (\textbf{SPIR-V:} \texttt{OpLoad (PerVertexKHR)}; \textbf{GLSL:} \texttt{pervertexEXT in vec4 v[3]; vec4 x = v[i];})
  \item \texttt{staging\_register\_index}: Base GPR index of the contiguous staging register window (\texttt{@rN:rM}) for message payload/return. \emph{Example:} {\scriptsize\ttfamily \mnem{STORE.i64}\imod{.slot7} \src{@r0:r1}, \src{r2}, \imm{\#8}} (\textbf{SPIR-V:} \texttt{OpStore (uvec2)}; \textbf{GLSL:} \texttt{ssbo.pairs[i] = uvec2(a, b);})
  \item \texttt{staging\_register\_control\_32}: Controls the read/write staging register bank and count for async message instructions. \emph{Example:} {\scriptsize\ttfamily \mnem{LOAD.i128}\imod{.slot0} \dst{@r0:r3}, \src{r4}, \imm{\#0}} (\textbf{SPIR-V:} \texttt{OpLoad (vec4)}; \textbf{GLSL:} \texttt{vec4 v = ssbo.data[idx];})
  \item \texttt{opcode1}: Primary sub-opcode within the instruction class uniquely identifying the operation family. \emph{Example:} {\scriptsize\ttfamily \mnem{IADD.s32} \dst{r0}, \src{r1}, \src{r2}} (\textbf{SPIR-V:} \texttt{OpIAdd (operation selector)}; \textbf{GLSL:} \texttt{// Primary opcode bits})
  \item \texttt{instruction\_class}: 4-bit major instruction class selector (\texttt{bits 48:51}) partitioning the 64-bit Valhall opcode space (e.g., \texttt{FMA}, \texttt{CVT}, \texttt{SFU}, \texttt{LS}, \texttt{TEX}, \texttt{BRANCH}). \emph{Example:} {\scriptsize\ttfamily \mnem{FADD.f32} \dst{r0}, \src{r1}, \src{r2}} (\textbf{SPIR-V:} \texttt{N/A (encoding class)}; \textbf{GLSL:} \texttt{// Hardware execution pipe / decode class})
  \item \texttt{fau\_page\_index}: Selects the FAU (Fast Access Uniform) page (\texttt{0..3}) when a source operand references hardware small constants, thread IDs, or lane constants. \emph{Example:} {\scriptsize\ttfamily \mnem{FADD.f32} \dst{r0}, \src{r1}, \src{k1f}} (\textbf{SPIR-V:} \texttt{OpConstant / BuiltIn}; \textbf{GLSL:} \texttt{float y = x + 1.0;})
  \item \hyperref[enum:flow_control_m]{\texttt{flow}}: 4-bit embedded flow-control \& scoreboard wait modifier (\texttt{bits 59:62}: \texttt{none}, \texttt{wait0..wait0126}, \texttt{wait}, \texttt{reconverge}, \texttt{discard}, \texttt{end}). \emph{Example:} {\scriptsize\ttfamily \mnem{FADD.f32}\imod{.wait0.reconverge} \dst{r0}, \src{r1}, \src{r2}} (\textbf{SPIR-V:} \texttt{OpControlBarrier / OpReturn}; \textbf{GLSL:} \texttt{// Scoreboard wait + warp reconvergence})
\end{tightitemize}
\clearpage

\subsection*{LD\_VAR\_BUF\_FLAT\_IMM}
\addcontentsline{toc}{subsection}{LD\_VAR\_BUF\_FLAT\_IMM}
Fetches an uninterpolated (flat) provoking-vertex varying from the descriptor-less varying buffer at immediate byte offset \texttt{\#offset} into \texttt{@sw} without interpolation or format conversion.
\par\textbf{Notes:}
\begin{tightitemize}
  \item Descriptor-less flat varying fetch (\texttt{ISA.xml:93}).
  \item \textbf{Source references:} \texttt{bifrost/valhall/ISA.xml:93}, \texttt{isa-v9-v14.xml}.
\end{tightitemize}
\codebox{\mnem{LD\_VAR\_BUF\_FLAT\_IMM}\imod{\{.message\_slot\_index\}}\imod{\{.flow\}}\imod{\{.vecsize\}}\imod{\{.register\_format\}}\imod{\{.vertex\}} \dst{@sw}, \imm{\#offset}}
\textbf{Operands:}
\begin{tightitemize}\footnotesize
  \item \texttt{\dst{@sw}} (\texttt{staging\_register\_index} + \texttt{sr\_count}): Destination staging register range receiving uninterpolated provoking-vertex data.
  \item \texttt{\imm{\#offset}} (\texttt{offset}): Immediate unsigned byte offset within the varying buffer.
\end{tightitemize}
\par\textbf{Closest SPIR-V Op:} \texttt{OpLoad (flat varying buffer imm, v14+)}
\par\textbf{Pseudo-GLSL (Operational Semantics):}
\codebox{@sw = readVaryingBufFlat(provoking\_vertex, byte\_off = \#offset);}
\textbf{Execution Context:}
\begin{tightitemize}\footnotesize
  \item \textbf{Varying buffer \& barycentric cache (\hyperref[sec:csf-tiler]{\texttt{IDVSSR::VARY\_SIZE}}, \texttt{r48:r49}):} Reads per-vertex varying RAM written by Varying EBIN \texttt{STORE.estream}; descriptor forms (\texttt{LD\_VAR*}) also chase \hyperref[sec:csf-srt]{\texttt{SRT\_2}} $\to$ \texttt{Resource[table]} $\to$ \texttt{Attribute[index]} and reuse cached barycentric weights via \texttt{.store}/\texttt{.retrieve}.
  \item \textbf{Fragment stage preloads \& FP modes (\hyperref[sec:csf-spd]{\texttt{SPD\_2}}, \texttt{r20:r21}):} Hardware pre-fills \texttt{r59} with \texttt{floatBitsToUint(FragCoord.z)} and \texttt{r60} with the MSAA sample coverage mask (\texttt{gl\_SampleMaskIn[0]}), while \texttt{ShaderProgram} word~0 latches 32/64 GPR allocation and warp-wide \texttt{flush\_to\_zero\_mode}, \texttt{suppress\_na\_n}, and \texttt{suppress\_inf}.
\end{tightitemize}
\par\textbf{Worked Examples:}
\par\textbf{Example 1 (Immediate-offset flat 32-bit varying (v14+)):}
\codebox{\mnem{LD\_VAR\_BUF\_FLAT\_IMM}\imod{.u32.slot0} \dst{@r0}, \imm{\#0x10}}
Fetches uninterpolated flat varying at immediate byte offset \texttt{0x10} into \texttt{@r0}.
\par\textbf{Example 2 (Flat 16-bit varying (`flat16`)):}
\codebox{\mnem{LD\_VAR\_BUF\_FLAT\_IMM}\imod{.flat16.slot0} \dst{@r0}, \imm{\#0x08}}
Reads 16-bit flat varying at immediate offset \texttt{0x08}.

\begin{center}\footnotesize
\begin{tabularx}{\linewidth}{@{}lllX@{}}
\toprule
Variant & Arch & Class & Opcode Fields \\
\midrule
\texttt{(none)} & v9--v14 & \texttt{S0} & \texttt{opcode1=LD\_VAR\_BUF\_FLAT\_IMM} \\
\bottomrule
\end{tabularx}
\end{center}
\begin{center}\footnotesize
\begin{tabularx}{0.88\linewidth}{@{}l l X@{}}
\toprule
Field Name & Bits & Type / Fixed Value \\
\midrule
\textit{Reserved} & \texttt{0:7} & \texttt{zero} \\
\texttt{offset} & \texttt{8:18} & \texttt{uint} \\
\textit{Reserved} & \texttt{19:23} & \texttt{zero} \\
\hyperref[enum:ld_var_buf_flat_format_m]{\texttt{register\_format}} & \texttt{24:25} & \texttt{ld\_var\_buf\_flat\_format\_m} \\
\textit{Reserved} & \texttt{26:27} & \texttt{zero} \\
\hyperref[enum:vecsize_varying_m]{\texttt{vecsize}} & \texttt{28:29} & \texttt{vecsize\_varying\_m (def none)} \\
\hyperref[enum:message_slot_index_m]{\texttt{message\_slot\_index}} & \texttt{30:32} & \texttt{message\_slot\_index\_m} \\
\texttt{sr\_count} & \texttt{33:35} & \texttt{sr\_count\_t} \\
\hyperref[enum:vertex_type_m]{\texttt{vertex\_type}} & \texttt{36} & \texttt{vertex\_type\_m} \\
\texttt{vertex\_id} & \texttt{37:38} & \texttt{uint} \\
\textit{Reserved} & \texttt{39} & \texttt{zero} \\
\texttt{staging\_register\_index} & \texttt{40:45} & \texttt{uint} \\
\texttt{staging\_register\_control\_32} & \texttt{46:47} & \texttt{exact=write} \\
\texttt{opcode1} & \texttt{48:51} & \texttt{exact=LD\_VAR\_BUF\_FLAT\_IMM} \\
\texttt{instruction\_class} & \texttt{52:56} & \texttt{exact=S0} \\
\texttt{fau\_page\_index} & \texttt{57:58} & \texttt{uint} \\
\hyperref[enum:flow_control_m]{\texttt{flow}} & \texttt{59:62} & \texttt{flow\_control\_m} \\
\textit{Reserved} & \texttt{63} & \texttt{zero} \\
\bottomrule
\end{tabularx}
\end{center}
\textbf{Fields Explained:}
\begin{tightitemize}\footnotesize
  \item \texttt{offset}: Signed byte offset added to base address register in memory load/store/atomic instructions. \emph{Example:} {\scriptsize\ttfamily \mnem{LOAD.i32} \dst{@r0}, \src{[r2:r3]}, \imm{\#16}} (\textbf{SPIR-V:} \texttt{OpAccessChain + OpLoad}; \textbf{GLSL:} \texttt{uint y = ssbo.arr[4];})
  \item \hyperref[enum:ld_var_buf_flat_format_m]{\texttt{register\_format}}: Target GPR data conversion format (\texttt{auto32}, \texttt{f32}, \texttt{f16}, \texttt{s32}, \texttt{s16}, \texttt{u32}, \texttt{u16}). \emph{Example:} {\scriptsize\ttfamily \mnem{LD\_TILE}\imod{.f16.v4} \dst{@r0:r1}, \src{r2}} (\textbf{SPIR-V:} \texttt{OpImageRead (SubpassData)}; \textbf{GLSL:} \texttt{f16vec4 c = subpassLoad(inputAtt);})
  \item \hyperref[enum:vecsize_varying_m]{\texttt{vecsize}}: Component count (\texttt{none}=1, \texttt{v2}, \texttt{v3}, \texttt{v4}) for attribute, varying, or blend transfers. \emph{Example:} {\scriptsize\ttfamily \mnem{LD\_ATTR}\imod{.f32.v3} \dst{@r0:r2}, \src{r4}} (\textbf{SPIR-V:} \texttt{OpLoad (vec3 Input)}; \textbf{GLSL:} \texttt{in vec3 a\_normal;})
  \item \hyperref[enum:message_slot_index_m]{\texttt{message\_slot\_index}}: Asynchronous scoreboard slot (\texttt{slot0}--\texttt{slot2}, \texttt{slot7}) tracking completion of this message. \emph{Example:} {\scriptsize\ttfamily \mnem{LOAD.i32}\imod{.slot0} \dst{@r0}, \src{r2}, \imm{\#0}} (\textbf{SPIR-V:} \texttt{OpLoad}; \textbf{GLSL:} \texttt{uint v = ssbo.arr[i];})
  \item \texttt{sr\_count}: Number of consecutive 32-bit staging registers transferred by a load/store/texture message. \emph{Example:} {\scriptsize\ttfamily \mnem{LD\_MULTI.i128} \dst{@r0:r3}, \src{r4}, \imm{\#0}} (\textbf{SPIR-V:} \texttt{OpLoad (wide vector)}; \textbf{GLSL:} \texttt{uvec4 x = buf[i];})
  \item \hyperref[enum:vertex_type_m]{\texttt{vertex\_type}}: Enables explicit register-supplied vertex index on \texttt{LD\_VAR\_BUF\_FLAT}. \emph{Example:} {\scriptsize\ttfamily \mnem{LD\_VAR\_BUF\_FLAT.f32}\imod{.explicit\_vertex} \dst{@r0:r3}, \src{r4}} (\textbf{SPIR-V:} \texttt{OpLoad (PerVertexKHR)}; \textbf{GLSL:} \texttt{pervertexEXT in vec4 v[3]; vec4 x = v[i];})
  \item \texttt{vertex\_id}: 2-bit immediate triangle vertex index (\texttt{0..2}) selecting which provoking vertex to read in flat varying loads. \emph{Example:} {\scriptsize\ttfamily \mnem{LD\_VAR\_FLAT\_IMM.f32} \dst{@r0:r3}, \imm{\#0}, \imm{\#1}} (\textbf{SPIR-V:} \texttt{OpLoad (Flat / PerVertexKHR)}; \textbf{GLSL:} \texttt{flat in vec4 v\_flat;})
  \item \texttt{staging\_register\_index}: Base GPR index of the contiguous staging register window (\texttt{@rN:rM}) for message payload/return. \emph{Example:} {\scriptsize\ttfamily \mnem{STORE.i64}\imod{.slot7} \src{@r0:r1}, \src{r2}, \imm{\#8}} (\textbf{SPIR-V:} \texttt{OpStore (uvec2)}; \textbf{GLSL:} \texttt{ssbo.pairs[i] = uvec2(a, b);})
  \item \texttt{staging\_register\_control\_32}: Controls the read/write staging register bank and count for async message instructions. \emph{Example:} {\scriptsize\ttfamily \mnem{LOAD.i128}\imod{.slot0} \dst{@r0:r3}, \src{r4}, \imm{\#0}} (\textbf{SPIR-V:} \texttt{OpLoad (vec4)}; \textbf{GLSL:} \texttt{vec4 v = ssbo.data[idx];})
  \item \texttt{opcode1}: Primary sub-opcode within the instruction class uniquely identifying the operation family. \emph{Example:} {\scriptsize\ttfamily \mnem{IADD.s32} \dst{r0}, \src{r1}, \src{r2}} (\textbf{SPIR-V:} \texttt{OpIAdd (operation selector)}; \textbf{GLSL:} \texttt{// Primary opcode bits})
  \item \texttt{instruction\_class}: 4-bit major instruction class selector (\texttt{bits 48:51}) partitioning the 64-bit Valhall opcode space (e.g., \texttt{FMA}, \texttt{CVT}, \texttt{SFU}, \texttt{LS}, \texttt{TEX}, \texttt{BRANCH}). \emph{Example:} {\scriptsize\ttfamily \mnem{FADD.f32} \dst{r0}, \src{r1}, \src{r2}} (\textbf{SPIR-V:} \texttt{N/A (encoding class)}; \textbf{GLSL:} \texttt{// Hardware execution pipe / decode class})
  \item \texttt{fau\_page\_index}: Selects the FAU (Fast Access Uniform) page (\texttt{0..3}) when a source operand references hardware small constants, thread IDs, or lane constants. \emph{Example:} {\scriptsize\ttfamily \mnem{FADD.f32} \dst{r0}, \src{r1}, \src{k1f}} (\textbf{SPIR-V:} \texttt{OpConstant / BuiltIn}; \textbf{GLSL:} \texttt{float y = x + 1.0;})
  \item \hyperref[enum:flow_control_m]{\texttt{flow}}: 4-bit embedded flow-control \& scoreboard wait modifier (\texttt{bits 59:62}: \texttt{none}, \texttt{wait0..wait0126}, \texttt{wait}, \texttt{reconverge}, \texttt{discard}, \texttt{end}). \emph{Example:} {\scriptsize\ttfamily \mnem{FADD.f32}\imod{.wait0.reconverge} \dst{r0}, \src{r1}, \src{r2}} (\textbf{SPIR-V:} \texttt{OpControlBarrier / OpReturn}; \textbf{GLSL:} \texttt{// Scoreboard wait + warp reconvergence})
\end{tightitemize}
\clearpage

\subsection*{LD\_VAR\_BUF\_IMM}
\addcontentsline{toc}{subsection}{LD\_VAR\_BUF\_IMM}
Interpolates a fragment varying from the descriptor-less varying buffer at immediate byte offset \texttt{\#offset} using pixel/sample coordinate \texttt{src0}, converting from \texttt{source\_format} to \texttt{register\_format} into \texttt{@sw} and updating the barycentric cache per \texttt{update\_mode}.
\par\textbf{Notes:}
\begin{tightitemize}
  \item Descriptor-less varying interpolation.
  \item \textbf{Source references:} \texttt{isa-v9-v14.xml}.
\end{tightitemize}
\codebox{\mnem{LD\_VAR\_BUF\_IMM}\imod{\{.message\_slot\_index\}}\imod{\{.flow\}}\imod{\{.vecsize\}}\imod{\{.source\_format\}}\imod{\{.sample\_position\}}\imod{\{.update\}} \dst{@sw}, \src{src0}, \imm{\#offset}}
\textbf{Operands:}
\begin{tightitemize}\footnotesize
  \item \texttt{\dst{@sw}} (\texttt{staging\_register\_index} + \texttt{sr\_count}): Destination staging register range receiving interpolated varying components from the varying buffer.
  \item \texttt{\src{src0}}: Fragment sample/pixel coordinate (\texttt{r61}).
  \item \texttt{\imm{\#offset}} (\texttt{offset}): Immediate unsigned byte offset within the varying buffer.
\end{tightitemize}
\par\textbf{Closest SPIR-V Op:} \texttt{OpLoad (varying buffer imm)}
\par\textbf{Pseudo-GLSL (Operational Semantics):}
\codebox{@sw = interpolateVaryingBuf(byte\_off = \#offset, pixel = src0, source\_format, update\_mode);}
\textbf{Execution Context:}
\begin{tightitemize}\footnotesize
  \item \textbf{Varying buffer \& barycentric cache (\hyperref[sec:csf-tiler]{\texttt{IDVSSR::VARY\_SIZE}}, \texttt{r48:r49}):} Reads per-vertex varying RAM written by Varying EBIN \texttt{STORE.estream}; descriptor forms (\texttt{LD\_VAR*}) also chase \hyperref[sec:csf-srt]{\texttt{SRT\_2}} $\to$ \texttt{Resource[table]} $\to$ \texttt{Attribute[index]} and reuse cached barycentric weights via \texttt{.store}/\texttt{.retrieve}.
  \item \textbf{Fragment stage preloads \& FP modes (\hyperref[sec:csf-spd]{\texttt{SPD\_2}}, \texttt{r20:r21}):} Hardware pre-fills \texttt{r59} with \texttt{floatBitsToUint(FragCoord.z)} and \texttt{r60} with the MSAA sample coverage mask (\texttt{gl\_SampleMaskIn[0]}), while \texttt{ShaderProgram} word~0 latches 32/64 GPR allocation and warp-wide \texttt{flush\_to\_zero\_mode}, \texttt{suppress\_na\_n}, and \texttt{suppress\_inf}.
\end{tightitemize}
\par\textbf{Worked Examples:}
\par\textbf{Example 1 (Cached barycentric varying load (`.retrieve`)):}
\codebox{\mnem{LD\_VAR\_BUF\_IMM.f32}\imod{.v2.center.retrieve.slot0} \dst{@r0:r1}, \src{r59}, \imm{\#0x10}}
Interpolates \texttt{vec2} varying at immediate byte offset \texttt{0x10} reusing cached barycentric coordinates.
\par\textbf{Example 2 (Per-sample MSAA interpolation (`.sample`)):}
\codebox{\mnem{LD\_VAR\_BUF\_IMM.f32}\imod{.v4.sample.store.slot0} \dst{@r0:r3}, \src{r59}, \imm{\#0x00}}
Evaluates \texttt{sample in vec4} at the active MSAA sample position rather than pixel center.

\begin{center}\footnotesize
\begin{tabularx}{\linewidth}{@{}lllX@{}}
\toprule
Variant & Arch & Class & Opcode Fields \\
\midrule
\texttt{(none)} & v9--v14 & \texttt{S1} & \texttt{opcode1=LD\_VAR\_BUF\_IMM\_F16} \\
\texttt{(none)} & v9--v14 & \texttt{S1} & \texttt{opcode1=LD\_VAR\_BUF\_IMM\_F16} \\
\texttt{(none)} & v9--v14 & \texttt{S1} & \texttt{opcode1=LD\_VAR\_BUF\_IMM\_F16} \\
\texttt{(none)} & v9--v14 & \texttt{S1} & \texttt{opcode1=LD\_VAR\_BUF\_IMM\_F32} \\
\texttt{(none)} & v9--v14 & \texttt{S1} & \texttt{opcode1=LD\_VAR\_BUF\_IMM\_F32} \\
\texttt{(none)} & v9--v14 & \texttt{S1} & \texttt{opcode1=LD\_VAR\_BUF\_IMM\_F32} \\
\bottomrule
\end{tabularx}
\end{center}
\begin{center}\footnotesize
\begin{tabularx}{0.88\linewidth}{@{}l l X@{}}
\toprule
Field Name & Bits & Type / Fixed Value \\
\midrule
\texttt{src0} & \texttt{0:7} & \texttt{SourceEncoding} \\
\textit{Reserved} & \texttt{8:15} & \texttt{zero} \\
\texttt{offset} & \texttt{16:23} & \texttt{uint} \\
\hyperref[enum:ld_var_buf_source_format_m]{\texttt{source\_format}} & \texttt{24:25} & \texttt{ld\_var\_buf\_source\_format\_m} \\
\textit{Reserved} & \texttt{26:27} & \texttt{zero} \\
\hyperref[enum:vecsize_varying_m]{\texttt{vecsize}} & \texttt{28:29} & \texttt{vecsize\_varying\_m (def none)} \\
\hyperref[enum:message_slot_index_m]{\texttt{message\_slot\_index}} & \texttt{30:32} & \texttt{message\_slot\_index\_m} \\
\texttt{sr\_count} & \texttt{33:35} & \texttt{sr\_count\_t} \\
\hyperref[enum:ld_var_buf_encoded_update_mode_none_m]{\texttt{encoded\_update}} & \texttt{36:37} & \texttt{ld\_var\_buf\_encoded\_update\_mode\_none\_m} \\
\hyperref[enum:encoded_sample_position_none_m]{\texttt{encoded\_sample}} & \texttt{38:39} & \texttt{encoded\_sample\_position\_none\_m} \\
\texttt{staging\_register\_index} & \texttt{40:45} & \texttt{uint} \\
\texttt{staging\_register\_control\_32} & \texttt{46:47} & \texttt{exact=write} \\
\texttt{opcode1} & \texttt{48:51} & \texttt{exact=LD\_VAR\_BUF\_IMM\_F16} \\
\texttt{instruction\_class} & \texttt{52:56} & \texttt{exact=S1} \\
\texttt{fau\_page\_index} & \texttt{57:58} & \texttt{uint} \\
\hyperref[enum:flow_control_m]{\texttt{flow}} & \texttt{59:62} & \texttt{flow\_control\_m} \\
\textit{Reserved} & \texttt{63} & \texttt{zero} \\
\bottomrule
\end{tabularx}
\end{center}
\textbf{Fields Explained:}
\begin{tightitemize}\footnotesize
  \item \texttt{src0}: Primary source operand (\texttt{r0..r63}, \texttt{u0..u63} uniform, \texttt{k*} FAU constant, or \texttt{t0..t2} temporary; \texttt{\textasciicircum{}} marks discard/last-use). \emph{Example:} {\scriptsize\ttfamily \mnem{FADD.f32} \dst{r0}, \src{r1}, \src{r2}} (\textbf{SPIR-V:} \texttt{OpFAdd (Operand 0)}; \textbf{GLSL:} \texttt{float c = a + b;})
  \item \texttt{offset}: Signed byte offset added to base address register in memory load/store/atomic instructions. \emph{Example:} {\scriptsize\ttfamily \mnem{LOAD.i32} \dst{@r0}, \src{[r2:r3]}, \imm{\#16}} (\textbf{SPIR-V:} \texttt{OpAccessChain + OpLoad}; \textbf{GLSL:} \texttt{uint y = ssbo.arr[4];})
  \item \hyperref[enum:ld_var_buf_source_format_m]{\texttt{source\_format}}: Declares the storage precision of the varying buffer (\texttt{src\_f32}, \texttt{src\_f16}, \texttt{src\_flat32}, \texttt{src\_flat16}). \emph{Example:} {\scriptsize\ttfamily \mnem{LD\_VAR\_BUF.f32}\imod{.src\_f16.v2} \dst{@r0:r1}, \src{r2}} (\textbf{SPIR-V:} \texttt{OpLoad (RelaxedPrecision Input)}; \textbf{GLSL:} \texttt{mediump in vec2 v\_uv;})
  \item \hyperref[enum:vecsize_varying_m]{\texttt{vecsize}}: Component count (\texttt{none}=1, \texttt{v2}, \texttt{v3}, \texttt{v4}) for attribute, varying, or blend transfers. \emph{Example:} {\scriptsize\ttfamily \mnem{LD\_ATTR}\imod{.f32.v3} \dst{@r0:r2}, \src{r4}} (\textbf{SPIR-V:} \texttt{OpLoad (vec3 Input)}; \textbf{GLSL:} \texttt{in vec3 a\_normal;})
  \item \hyperref[enum:message_slot_index_m]{\texttt{message\_slot\_index}}: Asynchronous scoreboard slot (\texttt{slot0}--\texttt{slot2}, \texttt{slot7}) tracking completion of this message. \emph{Example:} {\scriptsize\ttfamily \mnem{LOAD.i32}\imod{.slot0} \dst{@r0}, \src{r2}, \imm{\#0}} (\textbf{SPIR-V:} \texttt{OpLoad}; \textbf{GLSL:} \texttt{uint v = ssbo.arr[i];})
  \item \texttt{sr\_count}: Number of consecutive 32-bit staging registers transferred by a load/store/texture message. \emph{Example:} {\scriptsize\ttfamily \mnem{LD\_MULTI.i128} \dst{@r0:r3}, \src{r4}, \imm{\#0}} (\textbf{SPIR-V:} \texttt{OpLoad (wide vector)}; \textbf{GLSL:} \texttt{uvec4 x = buf[i];})
  \item \hyperref[enum:ld_var_buf_encoded_update_mode_none_m]{\texttt{encoded\_update}}: Compressed bitfield encoding for varying-cache update mode (\texttt{store}, \texttt{retrieve}, \texttt{clobber}). \emph{Example:} {\scriptsize\ttfamily \mnem{LD\_VAR\_SPECIAL}\imod{.bary.store} \dst{@r0:r1}, \src{r2}} (\textbf{SPIR-V:} \texttt{OpLoad (BuiltIn BaryCoordKHR)}; \textbf{GLSL:} \texttt{vec3 b = gl\_BaryCoordEXT;})
  \item \hyperref[enum:encoded_sample_position_none_m]{\texttt{encoded\_sample}}: Compressed encoding of MSAA sample interpolation position (\texttt{center\_or\_none}, \texttt{centroid}, \texttt{sample}, \texttt{explicit}). \emph{Example:} {\scriptsize\ttfamily \mnem{LD\_VAR.f32}\imod{.sample.v4} \dst{@r0:r3}, \src{r4}} (\textbf{SPIR-V:} \texttt{OpInterpolateAtSample}; \textbf{GLSL:} \texttt{vec4 v = interpolateAtSample(v\_col, 0);})
  \item \texttt{staging\_register\_index}: Base GPR index of the contiguous staging register window (\texttt{@rN:rM}) for message payload/return. \emph{Example:} {\scriptsize\ttfamily \mnem{STORE.i64}\imod{.slot7} \src{@r0:r1}, \src{r2}, \imm{\#8}} (\textbf{SPIR-V:} \texttt{OpStore (uvec2)}; \textbf{GLSL:} \texttt{ssbo.pairs[i] = uvec2(a, b);})
  \item \texttt{staging\_register\_control\_32}: Controls the read/write staging register bank and count for async message instructions. \emph{Example:} {\scriptsize\ttfamily \mnem{LOAD.i128}\imod{.slot0} \dst{@r0:r3}, \src{r4}, \imm{\#0}} (\textbf{SPIR-V:} \texttt{OpLoad (vec4)}; \textbf{GLSL:} \texttt{vec4 v = ssbo.data[idx];})
  \item \texttt{opcode1}: Primary sub-opcode within the instruction class uniquely identifying the operation family. \emph{Example:} {\scriptsize\ttfamily \mnem{IADD.s32} \dst{r0}, \src{r1}, \src{r2}} (\textbf{SPIR-V:} \texttt{OpIAdd (operation selector)}; \textbf{GLSL:} \texttt{// Primary opcode bits})
  \item \texttt{instruction\_class}: 4-bit major instruction class selector (\texttt{bits 48:51}) partitioning the 64-bit Valhall opcode space (e.g., \texttt{FMA}, \texttt{CVT}, \texttt{SFU}, \texttt{LS}, \texttt{TEX}, \texttt{BRANCH}). \emph{Example:} {\scriptsize\ttfamily \mnem{FADD.f32} \dst{r0}, \src{r1}, \src{r2}} (\textbf{SPIR-V:} \texttt{N/A (encoding class)}; \textbf{GLSL:} \texttt{// Hardware execution pipe / decode class})
  \item \texttt{fau\_page\_index}: Selects the FAU (Fast Access Uniform) page (\texttt{0..3}) when a source operand references hardware small constants, thread IDs, or lane constants. \emph{Example:} {\scriptsize\ttfamily \mnem{FADD.f32} \dst{r0}, \src{r1}, \src{k1f}} (\textbf{SPIR-V:} \texttt{OpConstant / BuiltIn}; \textbf{GLSL:} \texttt{float y = x + 1.0;})
  \item \hyperref[enum:flow_control_m]{\texttt{flow}}: 4-bit embedded flow-control \& scoreboard wait modifier (\texttt{bits 59:62}: \texttt{none}, \texttt{wait0..wait0126}, \texttt{wait}, \texttt{reconverge}, \texttt{discard}, \texttt{end}). \emph{Example:} {\scriptsize\ttfamily \mnem{FADD.f32}\imod{.wait0.reconverge} \dst{r0}, \src{r1}, \src{r2}} (\textbf{SPIR-V:} \texttt{OpControlBarrier / OpReturn}; \textbf{GLSL:} \texttt{// Scoreboard wait + warp reconvergence})
\end{tightitemize}
\clearpage

\subsection*{LD\_VAR\_FLAT}
\addcontentsline{toc}{subsection}{LD\_VAR\_FLAT}
Fetches an uninterpolated (flat) varying from the varying descriptor identified by dynamic handle \texttt{src0} for the selected triangle \texttt{vertex} (\texttt{vertex0}/\texttt{vertex1}/\texttt{vertex2}) into \texttt{@sw} without format conversion.
\par\textbf{Notes:}
\begin{tightitemize}
  \item Software varying buffer flat fetch (\texttt{ops.rs:3073-3100}).
  \item \textbf{Source references:} \texttt{kraid/ops.rs:3073-3100}, \texttt{bifrost/valhall/ISA.xml:95}.
\end{tightitemize}
\codebox{\mnem{LD\_VAR\_FLAT}\imod{\{.message\_slot\_index\}}\imod{\{.flow\}}\imod{\{.vecsize\}}\imod{\{.register\_format\}} \dst{@sw}, \src{src0}\\\mnem{LD\_VAR\_FLAT}\imod{\{.message\_slot\_index\}}\imod{\{.flow\}}\imod{\{.vecsize\}}\imod{\{.register\_format\}}\imod{\{.vertex\_type\}} \dst{@sw}, \src{src0}}
\textbf{Operands:}
\begin{tightitemize}\footnotesize
  \item \texttt{\dst{@sw}} (\texttt{staging\_register\_index} + \texttt{sr\_count}): Destination staging register range receiving uninterpolated flat varying components.
  \item \texttt{\src{src0}}: Dynamic varying descriptor handle.
\end{tightitemize}
\par\textbf{Closest SPIR-V Op:} \texttt{OpLoad (flat varying, dynamic)}
\par\textbf{Pseudo-GLSL (Operational Semantics):}
\codebox{@sw = readFlatVarying(handle = src0, vertex);}
\textbf{Execution Context:}
\begin{tightitemize}\footnotesize
  \item \textbf{Varying buffer \& barycentric cache (\hyperref[sec:csf-tiler]{\texttt{IDVSSR::VARY\_SIZE}}, \texttt{r48:r49}):} Reads per-vertex varying RAM written by Varying EBIN \texttt{STORE.estream}; descriptor forms (\texttt{LD\_VAR*}) also chase \hyperref[sec:csf-srt]{\texttt{SRT\_2}} $\to$ \texttt{Resource[table]} $\to$ \texttt{Attribute[index]} and reuse cached barycentric weights via \texttt{.store}/\texttt{.retrieve}.
  \item \textbf{Fragment stage preloads \& FP modes (\hyperref[sec:csf-spd]{\texttt{SPD\_2}}, \texttt{r20:r21}):} Hardware pre-fills \texttt{r59} with \texttt{floatBitsToUint(FragCoord.z)} and \texttt{r60} with the MSAA sample coverage mask (\texttt{gl\_SampleMaskIn[0]}), while \texttt{ShaderProgram} word~0 latches 32/64 GPR allocation and warp-wide \texttt{flush\_to\_zero\_mode}, \texttt{suppress\_na\_n}, and \texttt{suppress\_inf}.
\end{tightitemize}
\par\textbf{Worked Examples:}
\par\textbf{Example 1 (Dynamic-handle flat varying load):}
\codebox{\mnem{LD\_VAR\_FLAT}\imod{.u32.v4.slot0} \dst{@r0:r3}, \src{r4}}
Fetches 4x \texttt{u32} uninterpolated flat varying via descriptor handle \texttt{r4}.
\par\textbf{Example 2 (Vertex-1 flat read (`.vertex1`)):}
\codebox{\mnem{LD\_VAR\_FLAT}\imod{.f32.v4.vertex1.slot0} \dst{@r0:r3}, \src{r4}}
Reads triangle vertex 1's attribute directly without perspective interpolation.

\begin{center}\footnotesize
\begin{tabularx}{\linewidth}{@{}lllX@{}}
\toprule
Variant & Arch & Class & Opcode Fields \\
\midrule
\texttt{(none)} & v9--v14 & \texttt{S1} & \texttt{opcode1=LD\_VAR\_FLAT} \\
\texttt{(none)} & v9--v14 & \texttt{S1} & \texttt{opcode1=LD\_VAR\_FLAT} \\
\bottomrule
\end{tabularx}
\end{center}
\begin{center}\footnotesize
\begin{tabularx}{0.88\linewidth}{@{}l l X@{}}
\toprule
Field Name & Bits & Type / Fixed Value \\
\midrule
\texttt{src0} & \texttt{0:7} & \texttt{SourceEncoding} \\
\textit{Reserved} & \texttt{8:23} & \texttt{zero} \\
\hyperref[enum:register_file_format_general_m]{\texttt{register\_format}} & \texttt{24:26} & \texttt{register\_file\_format\_general\_m} \\
\textit{Reserved} & \texttt{27} & \texttt{zero} \\
\hyperref[enum:vecsize_varying_m]{\texttt{vecsize}} & \texttt{28:29} & \texttt{vecsize\_varying\_m (def none)} \\
\hyperref[enum:message_slot_index_m]{\texttt{message\_slot\_index}} & \texttt{30:32} & \texttt{message\_slot\_index\_m} \\
\texttt{sr\_count} & \texttt{33:35} & \texttt{sr\_count\_t} \\
\textit{Reserved} & \texttt{36:39} & \texttt{zero} \\
\texttt{staging\_register\_index} & \texttt{40:45} & \texttt{uint} \\
\texttt{staging\_register\_control\_32} & \texttt{46:47} & \texttt{exact=write} \\
\texttt{opcode1} & \texttt{48:51} & \texttt{exact=LD\_VAR\_FLAT} \\
\texttt{instruction\_class} & \texttt{52:56} & \texttt{exact=S1} \\
\texttt{fau\_page\_index} & \texttt{57:58} & \texttt{uint} \\
\hyperref[enum:flow_control_m]{\texttt{flow}} & \texttt{59:62} & \texttt{flow\_control\_m} \\
\textit{Reserved} & \texttt{63} & \texttt{zero} \\
\bottomrule
\end{tabularx}
\end{center}
\textbf{Fields Explained:}
\begin{tightitemize}\footnotesize
  \item \texttt{src0}: Primary source operand (\texttt{r0..r63}, \texttt{u0..u63} uniform, \texttt{k*} FAU constant, or \texttt{t0..t2} temporary; \texttt{\textasciicircum{}} marks discard/last-use). \emph{Example:} {\scriptsize\ttfamily \mnem{FADD.f32} \dst{r0}, \src{r1}, \src{r2}} (\textbf{SPIR-V:} \texttt{OpFAdd (Operand 0)}; \textbf{GLSL:} \texttt{float c = a + b;})
  \item \hyperref[enum:register_file_format_general_m]{\texttt{register\_format}}: Target GPR data conversion format (\texttt{auto32}, \texttt{f32}, \texttt{f16}, \texttt{s32}, \texttt{s16}, \texttt{u32}, \texttt{u16}). \emph{Example:} {\scriptsize\ttfamily \mnem{LD\_TILE}\imod{.f16.v4} \dst{@r0:r1}, \src{r2}} (\textbf{SPIR-V:} \texttt{OpImageRead (SubpassData)}; \textbf{GLSL:} \texttt{f16vec4 c = subpassLoad(inputAtt);})
  \item \hyperref[enum:vecsize_varying_m]{\texttt{vecsize}}: Component count (\texttt{none}=1, \texttt{v2}, \texttt{v3}, \texttt{v4}) for attribute, varying, or blend transfers. \emph{Example:} {\scriptsize\ttfamily \mnem{LD\_ATTR}\imod{.f32.v3} \dst{@r0:r2}, \src{r4}} (\textbf{SPIR-V:} \texttt{OpLoad (vec3 Input)}; \textbf{GLSL:} \texttt{in vec3 a\_normal;})
  \item \hyperref[enum:message_slot_index_m]{\texttt{message\_slot\_index}}: Asynchronous scoreboard slot (\texttt{slot0}--\texttt{slot2}, \texttt{slot7}) tracking completion of this message. \emph{Example:} {\scriptsize\ttfamily \mnem{LOAD.i32}\imod{.slot0} \dst{@r0}, \src{r2}, \imm{\#0}} (\textbf{SPIR-V:} \texttt{OpLoad}; \textbf{GLSL:} \texttt{uint v = ssbo.arr[i];})
  \item \texttt{sr\_count}: Number of consecutive 32-bit staging registers transferred by a load/store/texture message. \emph{Example:} {\scriptsize\ttfamily \mnem{LD\_MULTI.i128} \dst{@r0:r3}, \src{r4}, \imm{\#0}} (\textbf{SPIR-V:} \texttt{OpLoad (wide vector)}; \textbf{GLSL:} \texttt{uvec4 x = buf[i];})
  \item \texttt{staging\_register\_index}: Base GPR index of the contiguous staging register window (\texttt{@rN:rM}) for message payload/return. \emph{Example:} {\scriptsize\ttfamily \mnem{STORE.i64}\imod{.slot7} \src{@r0:r1}, \src{r2}, \imm{\#8}} (\textbf{SPIR-V:} \texttt{OpStore (uvec2)}; \textbf{GLSL:} \texttt{ssbo.pairs[i] = uvec2(a, b);})
  \item \texttt{staging\_register\_control\_32}: Controls the read/write staging register bank and count for async message instructions. \emph{Example:} {\scriptsize\ttfamily \mnem{LOAD.i128}\imod{.slot0} \dst{@r0:r3}, \src{r4}, \imm{\#0}} (\textbf{SPIR-V:} \texttt{OpLoad (vec4)}; \textbf{GLSL:} \texttt{vec4 v = ssbo.data[idx];})
  \item \texttt{opcode1}: Primary sub-opcode within the instruction class uniquely identifying the operation family. \emph{Example:} {\scriptsize\ttfamily \mnem{IADD.s32} \dst{r0}, \src{r1}, \src{r2}} (\textbf{SPIR-V:} \texttt{OpIAdd (operation selector)}; \textbf{GLSL:} \texttt{// Primary opcode bits})
  \item \texttt{instruction\_class}: 4-bit major instruction class selector (\texttt{bits 48:51}) partitioning the 64-bit Valhall opcode space (e.g., \texttt{FMA}, \texttt{CVT}, \texttt{SFU}, \texttt{LS}, \texttt{TEX}, \texttt{BRANCH}). \emph{Example:} {\scriptsize\ttfamily \mnem{FADD.f32} \dst{r0}, \src{r1}, \src{r2}} (\textbf{SPIR-V:} \texttt{N/A (encoding class)}; \textbf{GLSL:} \texttt{// Hardware execution pipe / decode class})
  \item \texttt{fau\_page\_index}: Selects the FAU (Fast Access Uniform) page (\texttt{0..3}) when a source operand references hardware small constants, thread IDs, or lane constants. \emph{Example:} {\scriptsize\ttfamily \mnem{FADD.f32} \dst{r0}, \src{r1}, \src{k1f}} (\textbf{SPIR-V:} \texttt{OpConstant / BuiltIn}; \textbf{GLSL:} \texttt{float y = x + 1.0;})
  \item \hyperref[enum:flow_control_m]{\texttt{flow}}: 4-bit embedded flow-control \& scoreboard wait modifier (\texttt{bits 59:62}: \texttt{none}, \texttt{wait0..wait0126}, \texttt{wait}, \texttt{reconverge}, \texttt{discard}, \texttt{end}). \emph{Example:} {\scriptsize\ttfamily \mnem{FADD.f32}\imod{.wait0.reconverge} \dst{r0}, \src{r1}, \src{r2}} (\textbf{SPIR-V:} \texttt{OpControlBarrier / OpReturn}; \textbf{GLSL:} \texttt{// Scoreboard wait + warp reconvergence})
\end{tightitemize}
\clearpage

\subsection*{LD\_VAR\_FLAT\_IMM}
\addcontentsline{toc}{subsection}{LD\_VAR\_FLAT\_IMM}
Fetches an uninterpolated (flat) varying from the varying descriptor at immediate \texttt{(table=\#varying\_table\_index, index=\#varying\_index)} for the selected triangle \texttt{vertex} into \texttt{@sw} without format conversion.
\par\textbf{Notes:}
\begin{tightitemize}
  \item Software varying buffer flat fetch (\texttt{ISA.xml:97}).
  \item \textbf{Source references:} \texttt{bifrost/valhall/ISA.xml:97}, \texttt{isa-v9-v14.xml}.
\end{tightitemize}
\codebox{\mnem{LD\_VAR\_FLAT\_IMM}\imod{\{.message\_slot\_index\}}\imod{\{.flow\}}\imod{\{.vecsize\}}\imod{\{.register\_format\}} \dst{@sw}, \imm{\#varying\_index}, \imm{\#varying\_table\_index}\\\mnem{LD\_VAR\_FLAT\_IMM}\imod{\{.message\_slot\_index\}}\imod{\{.flow\}}\imod{\{.vecsize\}}\imod{\{.register\_format\}}\imod{\{.vertex\}} \dst{@sw}, \imm{\#varying\_index}, \imm{\#varying\_table\_index}}
\textbf{Operands:}
\begin{tightitemize}\footnotesize
  \item \texttt{\dst{@sw}} (\texttt{staging\_register\_index} + \texttt{sr\_count}): Destination staging register range receiving uninterpolated flat varying components.
  \item \texttt{\imm{\#varying\_index}} (\texttt{varying\_index}): Immediate varying slot index within the varying table.
  \item \texttt{\imm{\#varying\_table\_index}} (\texttt{encoded\_varying\_table\_index}): Immediate SRT table index for varyings (typically \texttt{\#0xb}).
\end{tightitemize}
\par\textbf{Closest SPIR-V Op:} \texttt{OpLoad (flat varying)}
\par\textbf{Pseudo-GLSL (Operational Semantics):}
\codebox{AttrDesc v = SRT[\#varying\_table\_index].Attr[\#varying\_index];\\@sw = readFlatVarying(v, vertex);}
\textbf{Execution Context:}
\begin{tightitemize}\footnotesize
  \item \textbf{Varying buffer \& barycentric cache (\hyperref[sec:csf-tiler]{\texttt{IDVSSR::VARY\_SIZE}}, \texttt{r48:r49}):} Reads per-vertex varying RAM written by Varying EBIN \texttt{STORE.estream}; descriptor forms (\texttt{LD\_VAR*}) also chase \hyperref[sec:csf-srt]{\texttt{SRT\_2}} $\to$ \texttt{Resource[table]} $\to$ \texttt{Attribute[index]} and reuse cached barycentric weights via \texttt{.store}/\texttt{.retrieve}.
  \item \textbf{Fragment stage preloads \& FP modes (\hyperref[sec:csf-spd]{\texttt{SPD\_2}}, \texttt{r20:r21}):} Hardware pre-fills \texttt{r59} with \texttt{floatBitsToUint(FragCoord.z)} and \texttt{r60} with the MSAA sample coverage mask (\texttt{gl\_SampleMaskIn[0]}), while \texttt{ShaderProgram} word~0 latches 32/64 GPR allocation and warp-wide \texttt{flush\_to\_zero\_mode}, \texttt{suppress\_na\_n}, and \texttt{suppress\_inf}.
\end{tightitemize}
\par\textbf{Worked Examples:}
\par\textbf{Example 1 (Immediate descriptor flat load (`table=0, index=1`)):}
\codebox{\mnem{LD\_VAR\_FLAT\_IMM}\imod{.u32.v4.slot0} \dst{@r0:r3}, \imm{\#0}, \imm{\#1}}
Fetches 4x \texttt{u32} uninterpolated flat varying at immediate \texttt{(table=0, index=1)}.
\par\textbf{Example 2 (Explicit provoking vertex 2 (`.vertex2`)):}
\codebox{\mnem{LD\_VAR\_FLAT\_IMM}\imod{.f32.v4.vertex2.slot0} \dst{@r0:r3}, \imm{\#0}, \imm{\#2}}
Reads vertex 2 of the current primitive for \texttt{pervertexEXT}.
\par\textbf{Mined SPIR-V $\to$ Valhall ISA Example (Loading a flat-qualified fragment input varying --- Fragment shader, Mali-G710 (v10)):}
\codebox{\textbf{; SPIR-V (input)}\\OpDecorate \%in\_flat Flat\\\%v = OpLoad \%float \%in\_flat\\[3pt]\textbf{; Valhall ISA (malioc output)}\\LD\_VAR\_FLAT\_IMM.slot1.f32 @r4, \#0x1, \#0xb}
Because \texttt{\%in\_flat} carries \texttt{OpDecorate Flat}, \texttt{malioc} emits \texttt{LD\_VAR\_FLAT\_IMM} (no barycentric coordinate operand needed) reading directly from the provoking vertex at \texttt{(table=\#0xb, index=\#0x1)}.
\begin{center}\footnotesize
\begin{tabularx}{\linewidth}{@{}lllX@{}}
\toprule
Variant & Arch & Class & Opcode Fields \\
\midrule
\texttt{(none)} & v9--v14 & \texttt{S0} & \texttt{opcode1=LD\_VAR\_FLAT\_IMM} \\
\texttt{(none)} & v9--v14 & \texttt{S0} & \texttt{opcode1=LD\_VAR\_FLAT\_IMM} \\
\bottomrule
\end{tabularx}
\end{center}
\begin{center}\footnotesize
\begin{tabularx}{0.88\linewidth}{@{}l l X@{}}
\toprule
Field Name & Bits & Type / Fixed Value \\
\midrule
\textit{Reserved} & \texttt{0:7} & \texttt{zero} \\
\texttt{encoded\_varying\_table\_index} & \texttt{8:11} & \texttt{uint} \\
\texttt{varying\_index} & \texttt{12:19} & \texttt{uint} \\
\textit{Reserved} & \texttt{20:23} & \texttt{zero} \\
\hyperref[enum:register_file_format_general_m]{\texttt{register\_format}} & \texttt{24:26} & \texttt{register\_file\_format\_general\_m} \\
\textit{Reserved} & \texttt{27} & \texttt{zero} \\
\hyperref[enum:vecsize_varying_m]{\texttt{vecsize}} & \texttt{28:29} & \texttt{vecsize\_varying\_m (def none)} \\
\hyperref[enum:message_slot_index_m]{\texttt{message\_slot\_index}} & \texttt{30:32} & \texttt{message\_slot\_index\_m} \\
\texttt{sr\_count} & \texttt{33:35} & \texttt{sr\_count\_t} \\
\textit{Reserved} & \texttt{36:39} & \texttt{zero} \\
\texttt{staging\_register\_index} & \texttt{40:45} & \texttt{uint} \\
\texttt{staging\_register\_control\_32} & \texttt{46:47} & \texttt{exact=write} \\
\texttt{opcode1} & \texttt{48:51} & \texttt{exact=LD\_VAR\_FLAT\_IMM} \\
\texttt{instruction\_class} & \texttt{52:56} & \texttt{exact=S0} \\
\texttt{fau\_page\_index} & \texttt{57:58} & \texttt{exact=0} \\
\hyperref[enum:flow_control_m]{\texttt{flow}} & \texttt{59:62} & \texttt{flow\_control\_m} \\
\textit{Reserved} & \texttt{63} & \texttt{zero} \\
\bottomrule
\end{tabularx}
\end{center}
\textbf{Fields Explained:}
\begin{tightitemize}\footnotesize
  \item \texttt{encoded\_varying\_table\_index}: Encoded descriptor table index for \textbackslash\{\}texttt\{encoded\textbackslash\{\}\_varying\textbackslash\{\}\_table\textbackslash\{\}\_index\}. \emph{Example:} {\scriptsize\ttfamily \mnem{LD\_BUFFER.i32} \dst{@r0}, \src{r2}, \imm{\#0}} (\textbf{SPIR-V:} \texttt{OpAccessChain (DescriptorSet)}; \textbf{GLSL:} \texttt{uint v = buf[i];})
  \item \texttt{varying\_index}: Resource/slot index field \textbackslash\{\}texttt\{varying\textbackslash\{\}\_index\} selecting the target descriptor or slot. \emph{Example:} {\scriptsize\ttfamily \mnem{LD\_VAR\_BUF\_IMM.f32} \dst{@r0:r3}, \imm{\#0}} (\textbf{SPIR-V:} \texttt{OpLoad (Descriptor)}; \textbf{GLSL:} \texttt{vec4 v = v\_in;})
  \item \hyperref[enum:register_file_format_general_m]{\texttt{register\_format}}: Target GPR data conversion format (\texttt{auto32}, \texttt{f32}, \texttt{f16}, \texttt{s32}, \texttt{s16}, \texttt{u32}, \texttt{u16}). \emph{Example:} {\scriptsize\ttfamily \mnem{LD\_TILE}\imod{.f16.v4} \dst{@r0:r1}, \src{r2}} (\textbf{SPIR-V:} \texttt{OpImageRead (SubpassData)}; \textbf{GLSL:} \texttt{f16vec4 c = subpassLoad(inputAtt);})
  \item \hyperref[enum:vecsize_varying_m]{\texttt{vecsize}}: Component count (\texttt{none}=1, \texttt{v2}, \texttt{v3}, \texttt{v4}) for attribute, varying, or blend transfers. \emph{Example:} {\scriptsize\ttfamily \mnem{LD\_ATTR}\imod{.f32.v3} \dst{@r0:r2}, \src{r4}} (\textbf{SPIR-V:} \texttt{OpLoad (vec3 Input)}; \textbf{GLSL:} \texttt{in vec3 a\_normal;})
  \item \hyperref[enum:message_slot_index_m]{\texttt{message\_slot\_index}}: Asynchronous scoreboard slot (\texttt{slot0}--\texttt{slot2}, \texttt{slot7}) tracking completion of this message. \emph{Example:} {\scriptsize\ttfamily \mnem{LOAD.i32}\imod{.slot0} \dst{@r0}, \src{r2}, \imm{\#0}} (\textbf{SPIR-V:} \texttt{OpLoad}; \textbf{GLSL:} \texttt{uint v = ssbo.arr[i];})
  \item \texttt{sr\_count}: Number of consecutive 32-bit staging registers transferred by a load/store/texture message. \emph{Example:} {\scriptsize\ttfamily \mnem{LD\_MULTI.i128} \dst{@r0:r3}, \src{r4}, \imm{\#0}} (\textbf{SPIR-V:} \texttt{OpLoad (wide vector)}; \textbf{GLSL:} \texttt{uvec4 x = buf[i];})
  \item \texttt{staging\_register\_index}: Base GPR index of the contiguous staging register window (\texttt{@rN:rM}) for message payload/return. \emph{Example:} {\scriptsize\ttfamily \mnem{STORE.i64}\imod{.slot7} \src{@r0:r1}, \src{r2}, \imm{\#8}} (\textbf{SPIR-V:} \texttt{OpStore (uvec2)}; \textbf{GLSL:} \texttt{ssbo.pairs[i] = uvec2(a, b);})
  \item \texttt{staging\_register\_control\_32}: Controls the read/write staging register bank and count for async message instructions. \emph{Example:} {\scriptsize\ttfamily \mnem{LOAD.i128}\imod{.slot0} \dst{@r0:r3}, \src{r4}, \imm{\#0}} (\textbf{SPIR-V:} \texttt{OpLoad (vec4)}; \textbf{GLSL:} \texttt{vec4 v = ssbo.data[idx];})
  \item \texttt{opcode1}: Primary sub-opcode within the instruction class uniquely identifying the operation family. \emph{Example:} {\scriptsize\ttfamily \mnem{IADD.s32} \dst{r0}, \src{r1}, \src{r2}} (\textbf{SPIR-V:} \texttt{OpIAdd (operation selector)}; \textbf{GLSL:} \texttt{// Primary opcode bits})
  \item \texttt{instruction\_class}: 4-bit major instruction class selector (\texttt{bits 48:51}) partitioning the 64-bit Valhall opcode space (e.g., \texttt{FMA}, \texttt{CVT}, \texttt{SFU}, \texttt{LS}, \texttt{TEX}, \texttt{BRANCH}). \emph{Example:} {\scriptsize\ttfamily \mnem{FADD.f32} \dst{r0}, \src{r1}, \src{r2}} (\textbf{SPIR-V:} \texttt{N/A (encoding class)}; \textbf{GLSL:} \texttt{// Hardware execution pipe / decode class})
  \item \texttt{fau\_page\_index}: Selects the FAU (Fast Access Uniform) page (\texttt{0..3}) when a source operand references hardware small constants, thread IDs, or lane constants. \emph{Example:} {\scriptsize\ttfamily \mnem{FADD.f32} \dst{r0}, \src{r1}, \src{k1f}} (\textbf{SPIR-V:} \texttt{OpConstant / BuiltIn}; \textbf{GLSL:} \texttt{float y = x + 1.0;})
  \item \hyperref[enum:flow_control_m]{\texttt{flow}}: 4-bit embedded flow-control \& scoreboard wait modifier (\texttt{bits 59:62}: \texttt{none}, \texttt{wait0..wait0126}, \texttt{wait}, \texttt{reconverge}, \texttt{discard}, \texttt{end}). \emph{Example:} {\scriptsize\ttfamily \mnem{FADD.f32}\imod{.wait0.reconverge} \dst{r0}, \src{r1}, \src{r2}} (\textbf{SPIR-V:} \texttt{OpControlBarrier / OpReturn}; \textbf{GLSL:} \texttt{// Scoreboard wait + warp reconvergence})
\end{tightitemize}
\clearpage

\subsection*{LD\_VAR\_IMM}
\addcontentsline{toc}{subsection}{LD\_VAR\_IMM}
Interpolates \texttt{vecsize} varying components from the varying descriptor at immediate \texttt{(table=\#varying\_table\_index, index=\#varying\_index)} at pixel/sample coordinate \texttt{src0}, converting to \texttt{register\_format} into \texttt{@sw} and caching/reusing barycentrics per \texttt{update\_mode}.
\par\textbf{Notes:}
\begin{tightitemize}
  \item Software varying buffer interpolation (\texttt{ISA.xml:99}).
  \item \textbf{Source references:} \texttt{bifrost/valhall/ISA.xml:99}, \texttt{isa-v9-v14.xml}.
\end{tightitemize}
\codebox{\mnem{LD\_VAR\_IMM}\imod{\{.message\_slot\_index\}}\imod{\{.flow\}}\imod{\{.vecsize\}}\imod{\{.register\_format\}}\imod{\{.sample\_position\}}\imod{\{.update\}} \dst{@sw}, \src{src0}, \imm{\#varying\_index}, \imm{\#varying\_table\_index}}
\textbf{Operands:}
\begin{tightitemize}\footnotesize
  \item \texttt{\dst{@sw}} (\texttt{staging\_register\_index} + \texttt{sr\_count}): Destination staging register range receiving interpolated varying components.
  \item \texttt{\src{src0}}: Fragment sample/pixel coordinate (\texttt{r61}, or explicit offset).
  \item \texttt{\imm{\#varying\_index}} (\texttt{varying\_index}): Immediate varying slot index within the varying table.
  \item \texttt{\imm{\#varying\_table\_index}} (\texttt{encoded\_varying\_table\_index}): Immediate SRT table index for varyings (typically \texttt{\#0xb}).
\end{tightitemize}
\par\textbf{Closest SPIR-V Op:} \texttt{OpLoad (interpolated varying)}
\par\textbf{Pseudo-GLSL (Operational Semantics):}
\codebox{AttrDesc v = SRT[\#varying\_table\_index].Attr[\#varying\_index];\\@sw = interpolateVarying(v, pixel = src0, sample\_position, update\_mode);}
\textbf{Execution Context:}
\begin{tightitemize}\footnotesize
  \item \textbf{Varying buffer \& barycentric cache (\hyperref[sec:csf-tiler]{\texttt{IDVSSR::VARY\_SIZE}}, \texttt{r48:r49}):} Reads per-vertex varying RAM written by Varying EBIN \texttt{STORE.estream}; descriptor forms (\texttt{LD\_VAR*}) also chase \hyperref[sec:csf-srt]{\texttt{SRT\_2}} $\to$ \texttt{Resource[table]} $\to$ \texttt{Attribute[index]} and reuse cached barycentric weights via \texttt{.store}/\texttt{.retrieve}.
  \item \textbf{Fragment stage preloads \& FP modes (\hyperref[sec:csf-spd]{\texttt{SPD\_2}}, \texttt{r20:r21}):} Hardware pre-fills \texttt{r59} with \texttt{floatBitsToUint(FragCoord.z)} and \texttt{r60} with the MSAA sample coverage mask (\texttt{gl\_SampleMaskIn[0]}), while \texttt{ShaderProgram} word~0 latches 32/64 GPR allocation and warp-wide \texttt{flush\_to\_zero\_mode}, \texttt{suppress\_na\_n}, and \texttt{suppress\_inf}.
\end{tightitemize}
\par\textbf{Worked Examples:}
\par\textbf{Example 1 (Center-interpolated `vec4` varying):}
\codebox{\mnem{LD\_VAR\_IMM}\imod{.f32.v4.center.store.slot0} \dst{@r0:r3}, \src{r59}, \imm{\#0}, \imm{\#0}}
Interpolates \texttt{vec4} varying at immediate \texttt{(table=0, index=0)} into \texttt{@r0:r3} and caches barycentrics.
\par\textbf{Example 2 (Centroid-interpolated `vec2` varying):}
\codebox{\mnem{LD\_VAR\_IMM}\imod{.f32.v2.centroid.clobber.slot0} \dst{@r4:r5}, \src{r59}, \imm{\#0}, \imm{\#1}}
Interpolates \texttt{centroid in vec2} at \texttt{(table=0, index=1)} without preserving the prior barycentric cache.
\par\textbf{Mined SPIR-V $\to$ Valhall ISA Example (Interpolating a fragment varying and caching barycentrics with .store/.retrieve --- Fragment shader, Mali-G710 (v10)):}
\codebox{\textbf{; SPIR-V (input)}\\\%v0 = OpLoad \%float \%in\_loc\_1\\\%v1 = OpLoad \%v4float \%in\_loc\_0\\[3pt]\textbf{; Valhall ISA (malioc output)}\\LD\_VAR\_IMM.slot1.discard.f32.sample.store @r0, r61, \#0x1, \#0xb\\LD\_VAR\_IMM.slot0.wait0.v4.f32.retrieve @r4:r5:r6:r7, r61\textasciicircum{}, \#0x0, \#0xb}
The first \texttt{LD\_VAR\_IMM} uses \texttt{.sample.store} to compute and cache the perspective barycentric weights in hardware; the second uses \texttt{.retrieve} to reuse those cached weights without recomputing \texttt{1/w}.
\begin{center}\footnotesize
\begin{tabularx}{\linewidth}{@{}lllX@{}}
\toprule
Variant & Arch & Class & Opcode Fields \\
\midrule
\texttt{(none)} & v9--v14 & \texttt{S1} & \texttt{opcode1=LD\_VAR\_IMM} \\
\bottomrule
\end{tabularx}
\end{center}
\begin{center}\footnotesize
\begin{tabularx}{0.88\linewidth}{@{}l l X@{}}
\toprule
Field Name & Bits & Type / Fixed Value \\
\midrule
\texttt{src0} & \texttt{0:7} & \texttt{SourceEncoding} \\
\texttt{encoded\_varying\_table\_index} & \texttt{8:11} & \texttt{uint} \\
\texttt{varying\_index} & \texttt{12:19} & \texttt{uint} \\
\textit{Reserved} & \texttt{20:23} & \texttt{zero} \\
\hyperref[enum:register_file_format_general_m]{\texttt{register\_format}} & \texttt{24:26} & \texttt{register\_file\_format\_general\_m} \\
\textit{Reserved} & \texttt{27} & \texttt{zero} \\
\hyperref[enum:vecsize_varying_m]{\texttt{vecsize}} & \texttt{28:29} & \texttt{vecsize\_varying\_m (def none)} \\
\hyperref[enum:message_slot_index_m]{\texttt{message\_slot\_index}} & \texttt{30:32} & \texttt{message\_slot\_index\_m} \\
\texttt{sr\_count} & \texttt{33:35} & \texttt{sr\_count\_t} \\
\hyperref[enum:update_mode_m]{\texttt{update}} & \texttt{36:37} & \texttt{update\_mode\_m} \\
\hyperref[enum:encoded_sample_position_none_m]{\texttt{encoded\_sample}} & \texttt{38:39} & \texttt{encoded\_sample\_position\_none\_m} \\
\texttt{staging\_register\_index} & \texttt{40:45} & \texttt{uint} \\
\texttt{staging\_register\_control\_32} & \texttt{46:47} & \texttt{exact=write} \\
\texttt{opcode1} & \texttt{48:51} & \texttt{exact=LD\_VAR\_IMM} \\
\texttt{instruction\_class} & \texttt{52:56} & \texttt{exact=S1} \\
\texttt{fau\_page\_index} & \texttt{57:58} & \texttt{uint} \\
\hyperref[enum:flow_control_m]{\texttt{flow}} & \texttt{59:62} & \texttt{flow\_control\_m} \\
\textit{Reserved} & \texttt{63} & \texttt{zero} \\
\bottomrule
\end{tabularx}
\end{center}
\textbf{Fields Explained:}
\begin{tightitemize}\footnotesize
  \item \texttt{src0}: Primary source operand (\texttt{r0..r63}, \texttt{u0..u63} uniform, \texttt{k*} FAU constant, or \texttt{t0..t2} temporary; \texttt{\textasciicircum{}} marks discard/last-use). \emph{Example:} {\scriptsize\ttfamily \mnem{FADD.f32} \dst{r0}, \src{r1}, \src{r2}} (\textbf{SPIR-V:} \texttt{OpFAdd (Operand 0)}; \textbf{GLSL:} \texttt{float c = a + b;})
  \item \texttt{encoded\_varying\_table\_index}: Encoded descriptor table index for \textbackslash\{\}texttt\{encoded\textbackslash\{\}\_varying\textbackslash\{\}\_table\textbackslash\{\}\_index\}. \emph{Example:} {\scriptsize\ttfamily \mnem{LD\_BUFFER.i32} \dst{@r0}, \src{r2}, \imm{\#0}} (\textbf{SPIR-V:} \texttt{OpAccessChain (DescriptorSet)}; \textbf{GLSL:} \texttt{uint v = buf[i];})
  \item \texttt{varying\_index}: Resource/slot index field \textbackslash\{\}texttt\{varying\textbackslash\{\}\_index\} selecting the target descriptor or slot. \emph{Example:} {\scriptsize\ttfamily \mnem{LD\_VAR\_BUF\_IMM.f32} \dst{@r0:r3}, \imm{\#0}} (\textbf{SPIR-V:} \texttt{OpLoad (Descriptor)}; \textbf{GLSL:} \texttt{vec4 v = v\_in;})
  \item \hyperref[enum:register_file_format_general_m]{\texttt{register\_format}}: Target GPR data conversion format (\texttt{auto32}, \texttt{f32}, \texttt{f16}, \texttt{s32}, \texttt{s16}, \texttt{u32}, \texttt{u16}). \emph{Example:} {\scriptsize\ttfamily \mnem{LD\_TILE}\imod{.f16.v4} \dst{@r0:r1}, \src{r2}} (\textbf{SPIR-V:} \texttt{OpImageRead (SubpassData)}; \textbf{GLSL:} \texttt{f16vec4 c = subpassLoad(inputAtt);})
  \item \hyperref[enum:vecsize_varying_m]{\texttt{vecsize}}: Component count (\texttt{none}=1, \texttt{v2}, \texttt{v3}, \texttt{v4}) for attribute, varying, or blend transfers. \emph{Example:} {\scriptsize\ttfamily \mnem{LD\_ATTR}\imod{.f32.v3} \dst{@r0:r2}, \src{r4}} (\textbf{SPIR-V:} \texttt{OpLoad (vec3 Input)}; \textbf{GLSL:} \texttt{in vec3 a\_normal;})
  \item \hyperref[enum:message_slot_index_m]{\texttt{message\_slot\_index}}: Asynchronous scoreboard slot (\texttt{slot0}--\texttt{slot2}, \texttt{slot7}) tracking completion of this message. \emph{Example:} {\scriptsize\ttfamily \mnem{LOAD.i32}\imod{.slot0} \dst{@r0}, \src{r2}, \imm{\#0}} (\textbf{SPIR-V:} \texttt{OpLoad}; \textbf{GLSL:} \texttt{uint v = ssbo.arr[i];})
  \item \texttt{sr\_count}: Number of consecutive 32-bit staging registers transferred by a load/store/texture message. \emph{Example:} {\scriptsize\ttfamily \mnem{LD\_MULTI.i128} \dst{@r0:r3}, \src{r4}, \imm{\#0}} (\textbf{SPIR-V:} \texttt{OpLoad (wide vector)}; \textbf{GLSL:} \texttt{uvec4 x = buf[i];})
  \item \hyperref[enum:update_mode_m]{\texttt{update}}: Controls hardware varying-cache lifecycle (\texttt{store}, \texttt{retrieve}, \texttt{clobber}, \texttt{none}) to avoid redundant barycentric math. \emph{Example:} {\scriptsize\ttfamily \mnem{LD\_VAR\_BUF.f32}\imod{.retrieve.v4} \dst{@r0:r3}, \src{r4}} (\textbf{SPIR-V:} \texttt{OpLoad (Input varying)}; \textbf{GLSL:} \texttt{vec4 a = v\_pos; vec4 b = v\_norm;})
  \item \hyperref[enum:encoded_sample_position_none_m]{\texttt{encoded\_sample}}: Compressed encoding of MSAA sample interpolation position (\texttt{center\_or\_none}, \texttt{centroid}, \texttt{sample}, \texttt{explicit}). \emph{Example:} {\scriptsize\ttfamily \mnem{LD\_VAR.f32}\imod{.sample.v4} \dst{@r0:r3}, \src{r4}} (\textbf{SPIR-V:} \texttt{OpInterpolateAtSample}; \textbf{GLSL:} \texttt{vec4 v = interpolateAtSample(v\_col, 0);})
  \item \texttt{staging\_register\_index}: Base GPR index of the contiguous staging register window (\texttt{@rN:rM}) for message payload/return. \emph{Example:} {\scriptsize\ttfamily \mnem{STORE.i64}\imod{.slot7} \src{@r0:r1}, \src{r2}, \imm{\#8}} (\textbf{SPIR-V:} \texttt{OpStore (uvec2)}; \textbf{GLSL:} \texttt{ssbo.pairs[i] = uvec2(a, b);})
  \item \texttt{staging\_register\_control\_32}: Controls the read/write staging register bank and count for async message instructions. \emph{Example:} {\scriptsize\ttfamily \mnem{LOAD.i128}\imod{.slot0} \dst{@r0:r3}, \src{r4}, \imm{\#0}} (\textbf{SPIR-V:} \texttt{OpLoad (vec4)}; \textbf{GLSL:} \texttt{vec4 v = ssbo.data[idx];})
  \item \texttt{opcode1}: Primary sub-opcode within the instruction class uniquely identifying the operation family. \emph{Example:} {\scriptsize\ttfamily \mnem{IADD.s32} \dst{r0}, \src{r1}, \src{r2}} (\textbf{SPIR-V:} \texttt{OpIAdd (operation selector)}; \textbf{GLSL:} \texttt{// Primary opcode bits})
  \item \texttt{instruction\_class}: 4-bit major instruction class selector (\texttt{bits 48:51}) partitioning the 64-bit Valhall opcode space (e.g., \texttt{FMA}, \texttt{CVT}, \texttt{SFU}, \texttt{LS}, \texttt{TEX}, \texttt{BRANCH}). \emph{Example:} {\scriptsize\ttfamily \mnem{FADD.f32} \dst{r0}, \src{r1}, \src{r2}} (\textbf{SPIR-V:} \texttt{N/A (encoding class)}; \textbf{GLSL:} \texttt{// Hardware execution pipe / decode class})
  \item \texttt{fau\_page\_index}: Selects the FAU (Fast Access Uniform) page (\texttt{0..3}) when a source operand references hardware small constants, thread IDs, or lane constants. \emph{Example:} {\scriptsize\ttfamily \mnem{FADD.f32} \dst{r0}, \src{r1}, \src{k1f}} (\textbf{SPIR-V:} \texttt{OpConstant / BuiltIn}; \textbf{GLSL:} \texttt{float y = x + 1.0;})
  \item \hyperref[enum:flow_control_m]{\texttt{flow}}: 4-bit embedded flow-control \& scoreboard wait modifier (\texttt{bits 59:62}: \texttt{none}, \texttt{wait0..wait0126}, \texttt{wait}, \texttt{reconverge}, \texttt{discard}, \texttt{end}). \emph{Example:} {\scriptsize\ttfamily \mnem{FADD.f32}\imod{.wait0.reconverge} \dst{r0}, \src{r1}, \src{r2}} (\textbf{SPIR-V:} \texttt{OpControlBarrier / OpReturn}; \textbf{GLSL:} \texttt{// Scoreboard wait + warp reconvergence})
\end{tightitemize}
\clearpage

\subsection*{LD\_VAR\_SPECIAL}
\addcontentsline{toc}{subsection}{LD\_VAR\_SPECIAL}
Reads a fixed-function rasterizer varying selected by \texttt{varying\_name} (\texttt{point} for \texttt{gl\_PointCoord}, \texttt{bary} for barycentric coordinates, \texttt{frag\_w} for \texttt{gl\_FragCoord.w}, or \texttt{frag\_z} for \texttt{gl\_FragCoord.z}) at \texttt{sample\_position} (\texttt{src0}) into \texttt{@sw}.
\par\textbf{Notes:}
\begin{tightitemize}
  \item \texttt{frag\_z} does not support \texttt{explicit} sample position (\texttt{ops.rs:3111}).
  \item \textbf{Source references:} \texttt{kraid/ops.rs:3104-3150}, \texttt{isa-v9-v14.xml}.
\end{tightitemize}
\codebox{\mnem{LD\_VAR\_SPECIAL}\imod{\{.message\_slot\_index\}}\imod{\{.flow\}}\imod{\{.varying\_name\}}\imod{\{.vecsize\}}\imod{\{.register\_format\}}\imod{\{.sample\_position\}}\imod{\{.update\}} \dst{@sw}, \src{src0}}
\textbf{Operands:}
\begin{tightitemize}\footnotesize
  \item \texttt{\dst{@sw}} (\texttt{staging\_register\_index} + \texttt{sr\_count}): Destination staging register range receiving \texttt{point}, \texttt{bary}, \texttt{frag\_w}, or \texttt{frag\_z}.
  \item \texttt{\src{src0}}: Sample position / pixel coordinate (\texttt{r61} for \texttt{.sample}, \texttt{k0} for \texttt{.center}).
\end{tightitemize}
\par\textbf{Closest SPIR-V Op:} \texttt{OpLoad (FragCoord / PointCoord / BaryCoord)}
\par\textbf{Pseudo-GLSL (Operational Semantics):}
\codebox{if      (varying\_name == frag\_w) @sw = gl\_FragCoord.w;\\else if (varying\_name == frag\_z) @sw = gl\_FragCoord.z;\\else if (varying\_name == point)  @sw = gl\_PointCoord;\\else /* bary */                  @sw = gl\_BaryCoordEXT.xy;}
\textbf{Execution Context:}
\begin{tightitemize}\footnotesize
  \item \textbf{Varying buffer \& barycentric cache (\hyperref[sec:csf-tiler]{\texttt{IDVSSR::VARY\_SIZE}}, \texttt{r48:r49}):} Reads per-vertex varying RAM written by Varying EBIN \texttt{STORE.estream}; descriptor forms (\texttt{LD\_VAR*}) also chase \hyperref[sec:csf-srt]{\texttt{SRT\_2}} $\to$ \texttt{Resource[table]} $\to$ \texttt{Attribute[index]} and reuse cached barycentric weights via \texttt{.store}/\texttt{.retrieve}.
  \item \textbf{Fragment stage preloads \& FP modes (\hyperref[sec:csf-spd]{\texttt{SPD\_2}}, \texttt{r20:r21}):} Hardware pre-fills \texttt{r59} with \texttt{floatBitsToUint(FragCoord.z)} and \texttt{r60} with the MSAA sample coverage mask (\texttt{gl\_SampleMaskIn[0]}), while \texttt{ShaderProgram} word~0 latches 32/64 GPR allocation and warp-wide \texttt{flush\_to\_zero\_mode}, \texttt{suppress\_na\_n}, and \texttt{suppress\_inf}.
\end{tightitemize}
\par\textbf{Worked Examples:}
\par\textbf{Example 1 (Reciprocal clip-space W (`gl\_FragCoord.w`)):}
\codebox{\mnem{LD\_VAR\_SPECIAL}\imod{.frag\_w.center.clobber.slot0} \dst{@r0}, \src{r59}}
Loads perspective-correct \texttt{1/w} (\texttt{gl\_FragCoord.w}) at pixel center into \texttt{@r0}.
\par\textbf{Example 2 (Barycentric coordinates (`gl\_BaryCoordEXT`)):}
\codebox{\mnem{LD\_VAR\_SPECIAL}\imod{.bary.center.store.slot0} \dst{@r0:r1}, \src{r59}}
Loads \texttt{(i, j)} barycentric weights into \texttt{@r0:r1} (third weight is \texttt{1 - i - j}).
\par\textbf{Mined SPIR-V $\to$ Valhall ISA Example (Loading gl\_FragCoord.w via .frag\_w special varying --- Fragment shader, Mali-G710 (v10)):}
\codebox{\textbf{; SPIR-V (input)}\\\%p\_w = OpAccessChain \%\_ptr\_Input\_float \%gl\_FragCoord \%int\_3\\\%w   = OpLoad \%float \%p\_w\\[3pt]\textbf{; Valhall ISA (malioc output)}\\LD\_VAR\_SPECIAL.slot0.wait0.frag\_w.f32.center.clobber @r2, k0}
\texttt{gl\_FragCoord.w} does not use a user varying slot; \texttt{malioc} fetches it directly from the fixed-function rasterizer via \texttt{LD\_VAR\_SPECIAL.frag\_w}.
\begin{center}\footnotesize
\begin{tabularx}{\linewidth}{@{}lllX@{}}
\toprule
Variant & Arch & Class & Opcode Fields \\
\midrule
\texttt{(none)} & v9--v14 & \texttt{S1} & \texttt{opcode1=LD\_VAR\_SPECIAL} \\
\texttt{(none)} & v9--v14 & \texttt{S1} & \texttt{opcode1=LD\_VAR\_SPECIAL} \\
\bottomrule
\end{tabularx}
\end{center}
\begin{center}\footnotesize
\begin{tabularx}{0.88\linewidth}{@{}l l X@{}}
\toprule
Field Name & Bits & Type / Fixed Value \\
\midrule
\texttt{src0} & \texttt{0:7} & \texttt{SourceEncoding} \\
\textit{Reserved} & \texttt{8:11} & \texttt{zero} \\
\hyperref[enum:varying_name_m]{\texttt{varying\_name}} & \texttt{12:14} & \texttt{varying\_name\_m} \\
\textit{Reserved} & \texttt{15:23} & \texttt{zero} \\
\hyperref[enum:register_file_format_general_m]{\texttt{register\_format}} & \texttt{24:26} & \texttt{exact=f32} \\
\textit{Reserved} & \texttt{27} & \texttt{zero} \\
\hyperref[enum:vecsize_varying_m]{\texttt{vecsize}} & \texttt{28:29} & \texttt{vecsize\_varying\_m (def none)} \\
\hyperref[enum:message_slot_index_m]{\texttt{message\_slot\_index}} & \texttt{30:32} & \texttt{message\_slot\_index\_m} \\
\texttt{sr\_count} & \texttt{33:35} & \texttt{sr\_count\_t} \\
\hyperref[enum:encoded_update_mode_special_m]{\texttt{encoded\_update}} & \texttt{36:37} & \texttt{exact=clobber\_or\_none} \\
\hyperref[enum:sample_position_m]{\texttt{sample\_position}} & \texttt{38:39} & \texttt{sample\_position\_m} \\
\texttt{staging\_register\_index} & \texttt{40:45} & \texttt{uint} \\
\texttt{staging\_register\_control\_32} & \texttt{46:47} & \texttt{exact=write} \\
\texttt{opcode1} & \texttt{48:51} & \texttt{exact=LD\_VAR\_SPECIAL} \\
\texttt{instruction\_class} & \texttt{52:56} & \texttt{exact=S1} \\
\texttt{fau\_page\_index} & \texttt{57:58} & \texttt{uint} \\
\hyperref[enum:flow_control_m]{\texttt{flow}} & \texttt{59:62} & \texttt{flow\_control\_m} \\
\textit{Reserved} & \texttt{63} & \texttt{zero} \\
\bottomrule
\end{tabularx}
\end{center}
\textbf{Fields Explained:}
\begin{tightitemize}\footnotesize
  \item \texttt{src0}: Primary source operand (\texttt{r0..r63}, \texttt{u0..u63} uniform, \texttt{k*} FAU constant, or \texttt{t0..t2} temporary; \texttt{\textasciicircum{}} marks discard/last-use). \emph{Example:} {\scriptsize\ttfamily \mnem{FADD.f32} \dst{r0}, \src{r1}, \src{r2}} (\textbf{SPIR-V:} \texttt{OpFAdd (Operand 0)}; \textbf{GLSL:} \texttt{float c = a + b;})
  \item \hyperref[enum:varying_name_m]{\texttt{varying\_name}}: Selects fixed-function special varying (\texttt{point}, \texttt{bary}, \texttt{frag\_w}, \texttt{frag\_z}). \emph{Example:} {\scriptsize\ttfamily \mnem{LD\_VAR\_SPECIAL}\imod{.frag\_z} \dst{@r0}, \src{r2}} (\textbf{SPIR-V:} \texttt{OpLoad (BuiltIn FragCoord)}; \textbf{GLSL:} \texttt{float z = gl\_FragCoord.z;})
  \item \hyperref[enum:register_file_format_general_m]{\texttt{register\_format}}: Target GPR data conversion format (\texttt{auto32}, \texttt{f32}, \texttt{f16}, \texttt{s32}, \texttt{s16}, \texttt{u32}, \texttt{u16}). \emph{Example:} {\scriptsize\ttfamily \mnem{LD\_TILE}\imod{.f16.v4} \dst{@r0:r1}, \src{r2}} (\textbf{SPIR-V:} \texttt{OpImageRead (SubpassData)}; \textbf{GLSL:} \texttt{f16vec4 c = subpassLoad(inputAtt);})
  \item \hyperref[enum:vecsize_varying_m]{\texttt{vecsize}}: Component count (\texttt{none}=1, \texttt{v2}, \texttt{v3}, \texttt{v4}) for attribute, varying, or blend transfers. \emph{Example:} {\scriptsize\ttfamily \mnem{LD\_ATTR}\imod{.f32.v3} \dst{@r0:r2}, \src{r4}} (\textbf{SPIR-V:} \texttt{OpLoad (vec3 Input)}; \textbf{GLSL:} \texttt{in vec3 a\_normal;})
  \item \hyperref[enum:message_slot_index_m]{\texttt{message\_slot\_index}}: Asynchronous scoreboard slot (\texttt{slot0}--\texttt{slot2}, \texttt{slot7}) tracking completion of this message. \emph{Example:} {\scriptsize\ttfamily \mnem{LOAD.i32}\imod{.slot0} \dst{@r0}, \src{r2}, \imm{\#0}} (\textbf{SPIR-V:} \texttt{OpLoad}; \textbf{GLSL:} \texttt{uint v = ssbo.arr[i];})
  \item \texttt{sr\_count}: Number of consecutive 32-bit staging registers transferred by a load/store/texture message. \emph{Example:} {\scriptsize\ttfamily \mnem{LD\_MULTI.i128} \dst{@r0:r3}, \src{r4}, \imm{\#0}} (\textbf{SPIR-V:} \texttt{OpLoad (wide vector)}; \textbf{GLSL:} \texttt{uvec4 x = buf[i];})
  \item \hyperref[enum:encoded_update_mode_special_m]{\texttt{encoded\_update}}: Compressed bitfield encoding for varying-cache update mode (\texttt{store}, \texttt{retrieve}, \texttt{clobber}). \emph{Example:} {\scriptsize\ttfamily \mnem{LD\_VAR\_SPECIAL}\imod{.bary.store} \dst{@r0:r1}, \src{r2}} (\textbf{SPIR-V:} \texttt{OpLoad (BuiltIn BaryCoordKHR)}; \textbf{GLSL:} \texttt{vec3 b = gl\_BaryCoordEXT;})
  \item \hyperref[enum:sample_position_m]{\texttt{sample\_position}}: Selects MSAA interpolation location (\texttt{center}, \texttt{centroid}, \texttt{sample}, or \texttt{explicit} offset). \emph{Example:} {\scriptsize\ttfamily \mnem{LD\_VAR\_BUF.f32}\imod{.centroid.v2} \dst{@r0:r1}, \src{r2}} (\textbf{SPIR-V:} \texttt{OpInterpolateAtCentroid}; \textbf{GLSL:} \texttt{vec2 p = interpolateAtCentroid(v\_uv);})
  \item \texttt{staging\_register\_index}: Base GPR index of the contiguous staging register window (\texttt{@rN:rM}) for message payload/return. \emph{Example:} {\scriptsize\ttfamily \mnem{STORE.i64}\imod{.slot7} \src{@r0:r1}, \src{r2}, \imm{\#8}} (\textbf{SPIR-V:} \texttt{OpStore (uvec2)}; \textbf{GLSL:} \texttt{ssbo.pairs[i] = uvec2(a, b);})
  \item \texttt{staging\_register\_control\_32}: Controls the read/write staging register bank and count for async message instructions. \emph{Example:} {\scriptsize\ttfamily \mnem{LOAD.i128}\imod{.slot0} \dst{@r0:r3}, \src{r4}, \imm{\#0}} (\textbf{SPIR-V:} \texttt{OpLoad (vec4)}; \textbf{GLSL:} \texttt{vec4 v = ssbo.data[idx];})
  \item \texttt{opcode1}: Primary sub-opcode within the instruction class uniquely identifying the operation family. \emph{Example:} {\scriptsize\ttfamily \mnem{IADD.s32} \dst{r0}, \src{r1}, \src{r2}} (\textbf{SPIR-V:} \texttt{OpIAdd (operation selector)}; \textbf{GLSL:} \texttt{// Primary opcode bits})
  \item \texttt{instruction\_class}: 4-bit major instruction class selector (\texttt{bits 48:51}) partitioning the 64-bit Valhall opcode space (e.g., \texttt{FMA}, \texttt{CVT}, \texttt{SFU}, \texttt{LS}, \texttt{TEX}, \texttt{BRANCH}). \emph{Example:} {\scriptsize\ttfamily \mnem{FADD.f32} \dst{r0}, \src{r1}, \src{r2}} (\textbf{SPIR-V:} \texttt{N/A (encoding class)}; \textbf{GLSL:} \texttt{// Hardware execution pipe / decode class})
  \item \texttt{fau\_page\_index}: Selects the FAU (Fast Access Uniform) page (\texttt{0..3}) when a source operand references hardware small constants, thread IDs, or lane constants. \emph{Example:} {\scriptsize\ttfamily \mnem{FADD.f32} \dst{r0}, \src{r1}, \src{k1f}} (\textbf{SPIR-V:} \texttt{OpConstant / BuiltIn}; \textbf{GLSL:} \texttt{float y = x + 1.0;})
  \item \hyperref[enum:flow_control_m]{\texttt{flow}}: 4-bit embedded flow-control \& scoreboard wait modifier (\texttt{bits 59:62}: \texttt{none}, \texttt{wait0..wait0126}, \texttt{wait}, \texttt{reconverge}, \texttt{discard}, \texttt{end}). \emph{Example:} {\scriptsize\ttfamily \mnem{FADD.f32}\imod{.wait0.reconverge} \dst{r0}, \src{r1}, \src{r2}} (\textbf{SPIR-V:} \texttt{OpControlBarrier / OpReturn}; \textbf{GLSL:} \texttt{// Scoreboard wait + warp reconvergence})
\end{tightitemize}
\clearpage

\subsection*{ST\_TILE}
\addcontentsline{toc}{subsection}{ST\_TILE}
Encodes \texttt{vecsize} color components from \texttt{@sr} using \texttt{register\_format} and the tilebuffer conversion descriptor in \texttt{src2}, and stores the result into on-chip tile color SRAM at pixel \texttt{src0} masked by sample coverage \texttt{src1}.
\par\textbf{Notes:}
\begin{tightitemize}
  \item Used internally in blend shaders (\texttt{ISA.xml:144}).
  \item \textbf{Source references:} \texttt{kraid/ops.rs:4118-4146}, \texttt{encode\_v9.rs:3142-3165}, \texttt{bifrost/valhall/ISA.xml:144}.
\end{tightitemize}
\codebox{\mnem{ST\_TILE}\imod{\{.message\_slot\_index\}}\imod{\{.flow\}}\imod{\{.vecsize\}}\imod{\{.register\_format\}} \src{@sr}, \src{src0}, \src{src1}, \src{src2}}
\textbf{Operands:}
\begin{tightitemize}\footnotesize
  \item \texttt{\src{@sr}} (\texttt{staging\_register\_index} + \texttt{sr\_count}): Source staging register range holding the RGBA pixel value to write into tile SRAM.
  \item \texttt{\src{src0}}: Pixel coordinate / sample selector.
  \item \texttt{\src{src1}}: Sample coverage mask.
  \item \texttt{\src{src2}}: Tilebuffer conversion descriptor / render-target selector.
\end{tightitemize}
\par\textbf{Closest SPIR-V Op:} \texttt{None (blend shader tilebuffer write)}
\par\textbf{Pseudo-GLSL (Operational Semantics):}
\codebox{tileColorSRAM[pixel = src0, cov = src1] = encodeTileColor(@sr, cvt = src2, register\_format, vecsize);}
\textbf{Execution Context:}
\begin{tightitemize}\footnotesize
  \item \textbf{On-chip 16$\times$16 Tile SRAM (\hyperref[sec:csf-fbd]{\texttt{FRAGMENTSR::FBD\_POINTER}}, \texttt{r40:r41}):} 64B-aligned VA (\texttt{\& \textasciitilde{}0x3F}, \texttt{type[0]|zs\_ext[1]|rt\_cnt-1[5:2]}) $\to$ \texttt{FramebufferParameters} (128B) + optional \texttt{ZSCRCExtension} (64B, for \texttt{.z\_stencil}) + \texttt{RenderTarget[rt]} (64B: \texttt{internal\_format}, \texttt{clear\_color}, Morton 16$\times$16 \texttt{writeback\_buffer}).
  \item \textbf{Fragment stage preloads \& FP modes (\hyperref[sec:csf-spd]{\texttt{SPD\_2}}, \texttt{r20:r21}):} Hardware pre-fills \texttt{r59} with \texttt{floatBitsToUint(FragCoord.z)} and \texttt{r60} with the MSAA sample coverage mask (\texttt{gl\_SampleMaskIn[0]}), while \texttt{ShaderProgram} word~0 latches 32/64 GPR allocation and warp-wide \texttt{flush\_to\_zero\_mode}, \texttt{suppress\_na\_n}, and \texttt{suppress\_inf}.
\end{tightitemize}
\par\textbf{Worked Examples:}
\par\textbf{Example 1 (Programmable blend shader color write):}
\codebox{\mnem{ST\_TILE}\imod{.f32.v4.slot0} \dst{@r0:r3}, \src{r59}, \src{r60}, \src{u0}}
Converts and writes blended RGBA \texttt{@r0:r3} directly to on-chip tilebuffer at \texttt{(pixel=r59, cov=r60, desc=u0)}.
\par\textbf{Example 2 (Packed FP16 tilebuffer write):}
\codebox{\mnem{ST\_TILE}\imod{.f16.v4.slot0} \dst{@r0:r1}, \src{r59}, \src{r60}, \src{u0}}
Writes 4x FP16 channels from \texttt{@r0:r1} into an \texttt{R16G16B16A16\_SFLOAT} tile attachment.

\begin{center}\footnotesize
\begin{tabularx}{\linewidth}{@{}lllX@{}}
\toprule
Variant & Arch & Class & Opcode Fields \\
\midrule
\texttt{(none)} & v9--v14 & \texttt{S3} & \texttt{opcode1=ST\_TILE} \\
\bottomrule
\end{tabularx}
\end{center}
\begin{center}\footnotesize
\begin{tabularx}{0.88\linewidth}{@{}l l X@{}}
\toprule
Field Name & Bits & Type / Fixed Value \\
\midrule
\texttt{src0} & \texttt{0:7} & \texttt{SourceEncoding} \\
\texttt{src1} & \texttt{8:15} & \texttt{SourceEncoding} \\
\texttt{src2} & \texttt{16:23} & \texttt{SourceEncoding} \\
\hyperref[enum:register_file_format_general_m]{\texttt{register\_format}} & \texttt{24:26} & \texttt{register\_file\_format\_general\_m} \\
\textit{Reserved} & \texttt{27} & \texttt{zero} \\
\hyperref[enum:vecsize_varying_m]{\texttt{vecsize}} & \texttt{28:29} & \texttt{vecsize\_varying\_m (def none)} \\
\hyperref[enum:message_slot_index_m]{\texttt{message\_slot\_index}} & \texttt{30:32} & \texttt{message\_slot\_index\_m} \\
\texttt{sr\_count} & \texttt{33:35} & \texttt{sr\_count\_t} \\
\textit{Reserved} & \texttt{36:39} & \texttt{zero} \\
\texttt{staging\_register\_index} & \texttt{40:45} & \texttt{uint} \\
\texttt{staging\_register\_control\_32} & \texttt{46:47} & \texttt{exact=read} \\
\texttt{opcode1} & \texttt{48:51} & \texttt{exact=ST\_TILE} \\
\texttt{instruction\_class} & \texttt{52:56} & \texttt{exact=S3} \\
\texttt{fau\_page\_index} & \texttt{57:58} & \texttt{uint} \\
\hyperref[enum:flow_control_m]{\texttt{flow}} & \texttt{59:62} & \texttt{flow\_control\_m} \\
\textit{Reserved} & \texttt{63} & \texttt{zero} \\
\bottomrule
\end{tabularx}
\end{center}
\textbf{Fields Explained:}
\begin{tightitemize}\footnotesize
  \item \texttt{src0}: Primary source operand (\texttt{r0..r63}, \texttt{u0..u63} uniform, \texttt{k*} FAU constant, or \texttt{t0..t2} temporary; \texttt{\textasciicircum{}} marks discard/last-use). \emph{Example:} {\scriptsize\ttfamily \mnem{FADD.f32} \dst{r0}, \src{r1}, \src{r2}} (\textbf{SPIR-V:} \texttt{OpFAdd (Operand 0)}; \textbf{GLSL:} \texttt{float c = a + b;})
  \item \texttt{src1}: Secondary source operand (register, uniform, or FAU constant; in shifts, supplies the shift count byte). \emph{Example:} {\scriptsize\ttfamily \mnem{LSHIFT.i32} \dst{r0}, \src{r1}, \src{r2}\imod{.b0}} (\textbf{SPIR-V:} \texttt{OpShiftLeftLogical (Shift)}; \textbf{GLSL:} \texttt{uint y = a \textless{}\textless{} s;})
  \item \texttt{src2}: Third source operand (FMA addend, \texttt{MUX}/\texttt{CSEL} selector, \texttt{IADDC} carry-in, or bitwise mask in fused shift-logic). \emph{Example:} {\scriptsize\ttfamily \mnem{FMA.f32} \dst{r0}, \src{r1}, \src{r2}, \src{r3}} (\textbf{SPIR-V:} \texttt{GLSL.std.450 Fma (Addend)}; \textbf{GLSL:} \texttt{float d = fma(a, b, c);})
  \item \hyperref[enum:register_file_format_general_m]{\texttt{register\_format}}: Target GPR data conversion format (\texttt{auto32}, \texttt{f32}, \texttt{f16}, \texttt{s32}, \texttt{s16}, \texttt{u32}, \texttt{u16}). \emph{Example:} {\scriptsize\ttfamily \mnem{LD\_TILE}\imod{.f16.v4} \dst{@r0:r1}, \src{r2}} (\textbf{SPIR-V:} \texttt{OpImageRead (SubpassData)}; \textbf{GLSL:} \texttt{f16vec4 c = subpassLoad(inputAtt);})
  \item \hyperref[enum:vecsize_varying_m]{\texttt{vecsize}}: Component count (\texttt{none}=1, \texttt{v2}, \texttt{v3}, \texttt{v4}) for attribute, varying, or blend transfers. \emph{Example:} {\scriptsize\ttfamily \mnem{LD\_ATTR}\imod{.f32.v3} \dst{@r0:r2}, \src{r4}} (\textbf{SPIR-V:} \texttt{OpLoad (vec3 Input)}; \textbf{GLSL:} \texttt{in vec3 a\_normal;})
  \item \hyperref[enum:message_slot_index_m]{\texttt{message\_slot\_index}}: Asynchronous scoreboard slot (\texttt{slot0}--\texttt{slot2}, \texttt{slot7}) tracking completion of this message. \emph{Example:} {\scriptsize\ttfamily \mnem{LOAD.i32}\imod{.slot0} \dst{@r0}, \src{r2}, \imm{\#0}} (\textbf{SPIR-V:} \texttt{OpLoad}; \textbf{GLSL:} \texttt{uint v = ssbo.arr[i];})
  \item \texttt{sr\_count}: Number of consecutive 32-bit staging registers transferred by a load/store/texture message. \emph{Example:} {\scriptsize\ttfamily \mnem{LD\_MULTI.i128} \dst{@r0:r3}, \src{r4}, \imm{\#0}} (\textbf{SPIR-V:} \texttt{OpLoad (wide vector)}; \textbf{GLSL:} \texttt{uvec4 x = buf[i];})
  \item \texttt{staging\_register\_index}: Base GPR index of the contiguous staging register window (\texttt{@rN:rM}) for message payload/return. \emph{Example:} {\scriptsize\ttfamily \mnem{STORE.i64}\imod{.slot7} \src{@r0:r1}, \src{r2}, \imm{\#8}} (\textbf{SPIR-V:} \texttt{OpStore (uvec2)}; \textbf{GLSL:} \texttt{ssbo.pairs[i] = uvec2(a, b);})
  \item \texttt{staging\_register\_control\_32}: Controls the read/write staging register bank and count for async message instructions. \emph{Example:} {\scriptsize\ttfamily \mnem{LOAD.i128}\imod{.slot0} \dst{@r0:r3}, \src{r4}, \imm{\#0}} (\textbf{SPIR-V:} \texttt{OpLoad (vec4)}; \textbf{GLSL:} \texttt{vec4 v = ssbo.data[idx];})
  \item \texttt{opcode1}: Primary sub-opcode within the instruction class uniquely identifying the operation family. \emph{Example:} {\scriptsize\ttfamily \mnem{IADD.s32} \dst{r0}, \src{r1}, \src{r2}} (\textbf{SPIR-V:} \texttt{OpIAdd (operation selector)}; \textbf{GLSL:} \texttt{// Primary opcode bits})
  \item \texttt{instruction\_class}: 4-bit major instruction class selector (\texttt{bits 48:51}) partitioning the 64-bit Valhall opcode space (e.g., \texttt{FMA}, \texttt{CVT}, \texttt{SFU}, \texttt{LS}, \texttt{TEX}, \texttt{BRANCH}). \emph{Example:} {\scriptsize\ttfamily \mnem{FADD.f32} \dst{r0}, \src{r1}, \src{r2}} (\textbf{SPIR-V:} \texttt{N/A (encoding class)}; \textbf{GLSL:} \texttt{// Hardware execution pipe / decode class})
  \item \texttt{fau\_page\_index}: Selects the FAU (Fast Access Uniform) page (\texttt{0..3}) when a source operand references hardware small constants, thread IDs, or lane constants. \emph{Example:} {\scriptsize\ttfamily \mnem{FADD.f32} \dst{r0}, \src{r1}, \src{k1f}} (\textbf{SPIR-V:} \texttt{OpConstant / BuiltIn}; \textbf{GLSL:} \texttt{float y = x + 1.0;})
  \item \hyperref[enum:flow_control_m]{\texttt{flow}}: 4-bit embedded flow-control \& scoreboard wait modifier (\texttt{bits 59:62}: \texttt{none}, \texttt{wait0..wait0126}, \texttt{wait}, \texttt{reconverge}, \texttt{discard}, \texttt{end}). \emph{Example:} {\scriptsize\ttfamily \mnem{FADD.f32}\imod{.wait0.reconverge} \dst{r0}, \src{r1}, \src{r2}} (\textbf{SPIR-V:} \texttt{OpControlBarrier / OpReturn}; \textbf{GLSL:} \texttt{// Scoreboard wait + warp reconvergence})
\end{tightitemize}
\clearpage

\subsection*{ZS\_EMIT}
\addcontentsline{toc}{subsection}{ZS\_EMIT}
Emits programmatic depth (\texttt{src0}, if \texttt{.z}) and/or stencil (\texttt{src1}, if \texttt{.stencil}) with coverage \texttt{src2} (\texttt{r60}), returning updated coverage in \texttt{sr\_dst} (\texttt{r60}).
\par\textbf{Notes:}
\begin{tightitemize}
  \item Must execute after \texttt{DISCARD} and \texttt{ATEST} (\texttt{ops.rs:5106}, \texttt{encode\_v9.rs:3863-3865}).
  \item \textbf{Source references:} \texttt{kraid/ops.rs:5111-5146}, \texttt{encode\_v9.rs:3850-3875}, \texttt{bifrost/valhall/ISA.xml:174}.
\end{tightitemize}
\codebox{\mnem{ZS\_EMIT}\imod{\{.message\_slot\_index\}}\imod{\{.flow\}}\imod{\{.z\}}\imod{\{.stencil\}} \dst{@sw}, \src{src0}, \src{src1}, \src{src2}}
\textbf{Operands:}
\begin{tightitemize}\footnotesize
  \item \texttt{\dst{@sw}} (\texttt{staging\_register\_index} + \texttt{sr\_count}): Destination staging register (pinned \texttt{r60}) receiving the post-depth/stencil coverage mask.
  \item \texttt{\src{src0}}: Shader-computed depth value \texttt{gl\_FragDepth} (\texttt{f32}, used when \texttt{.z} is set).
  \item \texttt{\src{src1}}: Shader-computed 8-bit stencil reference (\texttt{src1 \& 0xFF}, used when \texttt{.stencil} is set).
  \item \texttt{\src{src2}}: Incoming sample coverage mask (pinned \texttt{r60}).
\end{tightitemize}
\par\textbf{Closest SPIR-V Op:} \texttt{OpStore (FragDepth / FragStencilRefEXT)}
\par\textbf{Pseudo-GLSL (Operational Semantics):}
\codebox{if (flag\_z)       gl\_FragDepth = src0;\\if (flag\_stencil) gl\_FragStencilRefARB = int(src1 \& 0xFFu);\\@sw = evalLateDepthStencilTest(cov = src2, z = src0, stencil = src1);}
\textbf{Execution Context:}
\begin{tightitemize}\footnotesize
  \item \textbf{Depth/Stencil \& coverage state (\hyperref[sec:csf-zsd-blend]{\texttt{IDVSSR::ZSD}}, \texttt{r52:r53}):} Interacts with 32B \texttt{Depth/stencil} (\texttt{depth\_func}, \texttt{stencil\_from\_shader}), \hyperref[sec:csf-tiler]{\texttt{DCD2}} (\texttt{r38}: \texttt{no\_shader\_depth\_read} Early-Z vs Late-Z), \texttt{DCD0}/\texttt{DCD1} (\texttt{r57:r58}: \texttt{shader\_modifies\_coverage}, \texttt{alpha\_to\_coverage}), and \texttt{FauSpecial::atest\_datum}.
  \item \textbf{Fragment stage preloads \& FP modes (\hyperref[sec:csf-spd]{\texttt{SPD\_2}}, \texttt{r20:r21}):} Hardware pre-fills \texttt{r59} with \texttt{floatBitsToUint(FragCoord.z)} and \texttt{r60} with the MSAA sample coverage mask (\texttt{gl\_SampleMaskIn[0]}), while \texttt{ShaderProgram} word~0 latches 32/64 GPR allocation and warp-wide \texttt{flush\_to\_zero\_mode}, \texttt{suppress\_na\_n}, and \texttt{suppress\_inf}.
\end{tightitemize}
\par\textbf{Worked Examples:}
\par\textbf{Example 1 (Programmatic fragment depth export (`.z`)):}
\codebox{\mnem{ZS\_EMIT}\imod{.z.slot0} \dst{@r60}, \src{r0}, \src{k0}, \src{r60}}
Emits shader-computed \texttt{gl\_FragDepth = r0} with coverage \texttt{r60}, returning post-depth-test coverage in \texttt{@r60}.
\par\textbf{Example 2 (Combined depth + stencil ref export (`.z.stencil`)):}
\codebox{\mnem{ZS\_EMIT}\imod{.z.stencil.slot0} \dst{@r60}, \src{r0}, \src{r1}, \src{r60}}
Exports both depth \texttt{r0} and 8-bit stencil reference \texttt{r1} simultaneously, updating \texttt{@r60}.
\par\textbf{Mined SPIR-V $\to$ Valhall ISA Example (Exporting shader-computed gl\_FragDepth --- Fragment shader, Mali-G710 (v10)):}
\codebox{\textbf{; SPIR-V (input)}\\OpDecorate \%gl\_FragDepth BuiltIn FragDepth\\OpStore \%gl\_FragDepth \%depth\_val\\[3pt]\textbf{; Valhall ISA (malioc output)}\\TEX\_SINGLE.slot0.wait0.skip.tex2d.dst32.r.wide.computed\_force\_delta @r0, @r0:r1, u1\\ZS\_EMIT.slot0.wait0.z @r60, r0, k0, r60\textasciicircum{}}
Writing \texttt{gl\_FragDepth} lowers to \texttt{ZS\_EMIT.z}, sending the sampled depth \texttt{r0} (with stencil unused = \texttt{k0}) and current coverage \texttt{r60\textasciicircum{}} to the late-ZS unit and receiving the updated post-depth-test coverage mask back into \texttt{@r60}.
\begin{center}\footnotesize
\begin{tabularx}{\linewidth}{@{}lllX@{}}
\toprule
Variant & Arch & Class & Opcode Fields \\
\midrule
\texttt{(none)} & v9--v14 & \texttt{S3} & \texttt{opcode1=ZS\_EMIT} \\
\bottomrule
\end{tabularx}
\end{center}
\begin{center}\footnotesize
\begin{tabularx}{0.88\linewidth}{@{}l l X@{}}
\toprule
Field Name & Bits & Type / Fixed Value \\
\midrule
\texttt{src0} & \texttt{0:7} & \texttt{SourceEncoding} \\
\texttt{src1} & \texttt{8:15} & \texttt{SourceEncoding} \\
\texttt{src2} & \texttt{16:23} & \texttt{SourceEncoding} \\
\hyperref[enum:stencil_m]{\texttt{stencil}} & \texttt{24} & \texttt{stencil\_m (def none)} \\
\hyperref[enum:z_m]{\texttt{z}} & \texttt{25} & \texttt{z\_m (def none)} \\
\textit{Reserved} & \texttt{26:29} & \texttt{zero} \\
\hyperref[enum:message_slot_index_m]{\texttt{message\_slot\_index}} & \texttt{30:32} & \texttt{exact=slot0} \\
\texttt{sr\_count} & \texttt{33:35} & \texttt{exact=count1} \\
\textit{Reserved} & \texttt{36:39} & \texttt{zero} \\
\texttt{staging\_register\_index} & \texttt{40:45} & \texttt{uint} \\
\texttt{staging\_register\_control\_32} & \texttt{46:47} & \texttt{exact=write} \\
\texttt{opcode1} & \texttt{48:51} & \texttt{exact=ZS\_EMIT} \\
\texttt{instruction\_class} & \texttt{52:56} & \texttt{exact=S3} \\
\texttt{fau\_page\_index} & \texttt{57:58} & \texttt{uint} \\
\hyperref[enum:flow_control_m]{\texttt{flow}} & \texttt{59:62} & \texttt{flow\_control\_m} \\
\textit{Reserved} & \texttt{63} & \texttt{zero} \\
\bottomrule
\end{tabularx}
\end{center}
\textbf{Fields Explained:}
\begin{tightitemize}\footnotesize
  \item \texttt{src0}: Primary source operand (\texttt{r0..r63}, \texttt{u0..u63} uniform, \texttt{k*} FAU constant, or \texttt{t0..t2} temporary; \texttt{\textasciicircum{}} marks discard/last-use). \emph{Example:} {\scriptsize\ttfamily \mnem{FADD.f32} \dst{r0}, \src{r1}, \src{r2}} (\textbf{SPIR-V:} \texttt{OpFAdd (Operand 0)}; \textbf{GLSL:} \texttt{float c = a + b;})
  \item \texttt{src1}: Secondary source operand (register, uniform, or FAU constant; in shifts, supplies the shift count byte). \emph{Example:} {\scriptsize\ttfamily \mnem{LSHIFT.i32} \dst{r0}, \src{r1}, \src{r2}\imod{.b0}} (\textbf{SPIR-V:} \texttt{OpShiftLeftLogical (Shift)}; \textbf{GLSL:} \texttt{uint y = a \textless{}\textless{} s;})
  \item \texttt{src2}: Third source operand (FMA addend, \texttt{MUX}/\texttt{CSEL} selector, \texttt{IADDC} carry-in, or bitwise mask in fused shift-logic). \emph{Example:} {\scriptsize\ttfamily \mnem{FMA.f32} \dst{r0}, \src{r1}, \src{r2}, \src{r3}} (\textbf{SPIR-V:} \texttt{GLSL.std.450 Fma (Addend)}; \textbf{GLSL:} \texttt{float d = fma(a, b, c);})
  \item \hyperref[enum:stencil_m]{\texttt{stencil}}: Emits shader-computed fragment stencil ref (\texttt{gl\_FragStencilRefARB}) in \texttt{ZS\_EMIT}. \emph{Example:} {\scriptsize\ttfamily \mnem{ZS\_EMIT}\imod{.stencil} \src{r0}, \src{r1}} (\textbf{SPIR-V:} \texttt{OpStore (BuiltIn FragStencilRefEXT)}; \textbf{GLSL:} \texttt{gl\_FragStencilRefARB = s;})
  \item \hyperref[enum:z_m]{\texttt{z}}: Emits shader-computed fragment depth (\texttt{gl\_FragDepth}) in \texttt{ZS\_EMIT}. \emph{Example:} {\scriptsize\ttfamily \mnem{ZS\_EMIT}\imod{.z} \src{r0}, \src{r1}} (\textbf{SPIR-V:} \texttt{OpStore (BuiltIn FragDepth)}; \textbf{GLSL:} \texttt{gl\_FragDepth = d;})
  \item \hyperref[enum:message_slot_index_m]{\texttt{message\_slot\_index}}: Asynchronous scoreboard slot (\texttt{slot0}--\texttt{slot2}, \texttt{slot7}) tracking completion of this message. \emph{Example:} {\scriptsize\ttfamily \mnem{LOAD.i32}\imod{.slot0} \dst{@r0}, \src{r2}, \imm{\#0}} (\textbf{SPIR-V:} \texttt{OpLoad}; \textbf{GLSL:} \texttt{uint v = ssbo.arr[i];})
  \item \texttt{sr\_count}: Number of consecutive 32-bit staging registers transferred by a load/store/texture message. \emph{Example:} {\scriptsize\ttfamily \mnem{LD\_MULTI.i128} \dst{@r0:r3}, \src{r4}, \imm{\#0}} (\textbf{SPIR-V:} \texttt{OpLoad (wide vector)}; \textbf{GLSL:} \texttt{uvec4 x = buf[i];})
  \item \texttt{staging\_register\_index}: Base GPR index of the contiguous staging register window (\texttt{@rN:rM}) for message payload/return. \emph{Example:} {\scriptsize\ttfamily \mnem{STORE.i64}\imod{.slot7} \src{@r0:r1}, \src{r2}, \imm{\#8}} (\textbf{SPIR-V:} \texttt{OpStore (uvec2)}; \textbf{GLSL:} \texttt{ssbo.pairs[i] = uvec2(a, b);})
  \item \texttt{staging\_register\_control\_32}: Controls the read/write staging register bank and count for async message instructions. \emph{Example:} {\scriptsize\ttfamily \mnem{LOAD.i128}\imod{.slot0} \dst{@r0:r3}, \src{r4}, \imm{\#0}} (\textbf{SPIR-V:} \texttt{OpLoad (vec4)}; \textbf{GLSL:} \texttt{vec4 v = ssbo.data[idx];})
  \item \texttt{opcode1}: Primary sub-opcode within the instruction class uniquely identifying the operation family. \emph{Example:} {\scriptsize\ttfamily \mnem{IADD.s32} \dst{r0}, \src{r1}, \src{r2}} (\textbf{SPIR-V:} \texttt{OpIAdd (operation selector)}; \textbf{GLSL:} \texttt{// Primary opcode bits})
  \item \texttt{instruction\_class}: 4-bit major instruction class selector (\texttt{bits 48:51}) partitioning the 64-bit Valhall opcode space (e.g., \texttt{FMA}, \texttt{CVT}, \texttt{SFU}, \texttt{LS}, \texttt{TEX}, \texttt{BRANCH}). \emph{Example:} {\scriptsize\ttfamily \mnem{FADD.f32} \dst{r0}, \src{r1}, \src{r2}} (\textbf{SPIR-V:} \texttt{N/A (encoding class)}; \textbf{GLSL:} \texttt{// Hardware execution pipe / decode class})
  \item \texttt{fau\_page\_index}: Selects the FAU (Fast Access Uniform) page (\texttt{0..3}) when a source operand references hardware small constants, thread IDs, or lane constants. \emph{Example:} {\scriptsize\ttfamily \mnem{FADD.f32} \dst{r0}, \src{r1}, \src{k1f}} (\textbf{SPIR-V:} \texttt{OpConstant / BuiltIn}; \textbf{GLSL:} \texttt{float y = x + 1.0;})
  \item \hyperref[enum:flow_control_m]{\texttt{flow}}: 4-bit embedded flow-control \& scoreboard wait modifier (\texttt{bits 59:62}: \texttt{none}, \texttt{wait0..wait0126}, \texttt{wait}, \texttt{reconverge}, \texttt{discard}, \texttt{end}). \emph{Example:} {\scriptsize\ttfamily \mnem{FADD.f32}\imod{.wait0.reconverge} \dst{r0}, \src{r1}, \src{r2}} (\textbf{SPIR-V:} \texttt{OpControlBarrier / OpReturn}; \textbf{GLSL:} \texttt{// Scoreboard wait + warp reconvergence})
\end{tightitemize}
\clearpage

\section{Texture \& Varying-Texture (14 mnemonics)}
\subsection*{TEX\_DUAL}
\addcontentsline{toc}{subsection}{TEX\_DUAL}
Paired dual texture sample issuing two samples with shared coordinates and independent \texttt{primary}/\texttt{secondary} widths and write masks.
\par\textbf{Notes:}
\begin{tightitemize}
  \item Combines two texture fetches into one message (\texttt{ISA.xml:146}).
  \item \textbf{Source references:} \texttt{kraid/ops.rs:4343-4391}, \texttt{bifrost/valhall/ISA.xml:146}.
\end{tightitemize}
\codebox{\mnem{TEX\_DUAL}\imod{\{.message\_slot\_index\}}\imod{\{.flow\}}\imod{\{.skip\}}\imod{\{.dimensionality\}}\imod{\{.primary\_width\}}\imod{\{.primary\_write\_mask\}}\imod{\{.secondary\_width\}}\imod{\{.secondary\_write\_mask\}}\imod{\{.array\_enable\}}\imod{\{.texel\_offset\}}\imod{\{.compare\_enable\}}\imod{\{.lod\_mode\}} \dst{@sw}, \src{@sr}, \src{src0}\\\mnem{TEX\_DUAL}\imod{\{.message\_slot\_index\}}\imod{\{.flow\}}\imod{\{.skip\}}\imod{\{.dimensionality\}}\imod{\{.projection\_enable\}}\imod{\{.primary\_width\}}\imod{\{.primary\_write\_mask\}}\imod{\{.secondary\_width\}}\imod{\{.secondary\_write\_mask\}}\imod{\{.array\_enable\}}\imod{\{.texel\_offset\}}\imod{\{.compare\_enable\}}\imod{\{.lod\_mode\}} \dst{@sw}, \src{@sr}, \src{src0}}
\textbf{Operands:}
\begin{tightitemize}\footnotesize
  \item \texttt{\dst{@sw}} (\texttt{staging\_register\_write\_index} + \texttt{sr\_write\_count}): Destination staging register range receiving two consecutive 4-component sample results.
  \item \texttt{\src{@sr}} (\texttt{staging\_register\_index} + \texttt{sr\_count}): Source staging register range holding shared texture coordinates and LOD parameters.
  \item \texttt{\src{src0}}: Dual texture/sampler descriptor handle.
\end{tightitemize}
\par\textbf{Closest SPIR-V Op:} \texttt{2x OpImageSample* (shared coords)}
\par\textbf{Pseudo-GLSL (Operational Semantics):}
\codebox{@sw[0:3] = textureSample(tex0\_of(src0), coords = @sr, lod\_mode);\\@sw[4:7] = textureSample(tex1\_of(src0), coords = @sr, lod\_mode);}
\textbf{Execution Context:}
\begin{tightitemize}\footnotesize
  \item \textbf{Texture \& Sampler chain (\hyperref[sec:csf-srt]{\texttt{SRT\_2}} Frag / \texttt{SRT\_0} Comp/Vert):} 64B-aligned VA (\texttt{\& \textasciitilde{}0x3F}) $\to$ \texttt{Resource[table]} $\to$ \hyperref[sec:csf-tex-samp]{\texttt{Texture[tex\_idx]}} (32B, \texttt{type=2}) $\to$ \texttt{GenericPlane[mip]} (32B, \texttt{type=10}: \texttt{pointer}, \texttt{row\_stride}, Morton 16$\times$16 / AFBC) + \texttt{Sampler[samp\_idx]} (32B, \texttt{type=1}: filter, wrap, LOD clamp/bias, compare).
  \item \textbf{Fragment stage preloads \& FP modes (\hyperref[sec:csf-spd]{\texttt{SPD\_2}}, \texttt{r20:r21}):} Hardware pre-fills \texttt{r59} with \texttt{floatBitsToUint(FragCoord.z)} and \texttt{r60} with the MSAA sample coverage mask (\texttt{gl\_SampleMaskIn[0]}), while \texttt{ShaderProgram} word~0 latches 32/64 GPR allocation and warp-wide \texttt{flush\_to\_zero\_mode}, \texttt{suppress\_na\_n}, and \texttt{suppress\_inf}.
\end{tightitemize}
\par\textbf{Worked Examples:}
\par\textbf{Example 1 (Paired 2D albedo + normal sample):}
\codebox{\mnem{TEX\_DUAL}\imod{.tex2d.computed.slot0} \dst{@r0:r7}, \src{@r8:r9}, \src{u0}}
Issues two filtered 2D samples sharing UV \texttt{@r8:r9}, writing primary RGBA to \texttt{@r0:r3} and secondary RGBA to \texttt{@r4:r7}.
\par\textbf{Example 2 (Mixed-mask dual sample (`.primary\_rgb.secondary\_r`)):}
\codebox{\mnem{TEX\_DUAL}\imod{.tex2d.primary\_rgb.secondary\_r.slot0} \dst{@r0:r3}, \src{@r8:r9}, \src{u0}}
Packs primary \texttt{.rgb} into \texttt{@r0:r2} and secondary \texttt{.r} (e.g., roughness) into \texttt{@r3}, saving 4 staging GPRs.

\begin{center}\footnotesize
\begin{tabularx}{\linewidth}{@{}lllX@{}}
\toprule
Variant & Arch & Class & Opcode Fields \\
\midrule
\texttt{(none)} & v9--v14 & \texttt{S2RW} & \texttt{opcode1=TEX\_DUAL} \\
\texttt{(none)} & v9--v14 & \texttt{S2RW} & \texttt{opcode1=TEX\_DUAL} \\
\texttt{(none)} & v9--v14 & \texttt{S2RW} & \texttt{opcode1=TEX\_DUAL} \\
\bottomrule
\end{tabularx}
\end{center}
\begin{center}\footnotesize
\begin{tabularx}{0.88\linewidth}{@{}l l X@{}}
\toprule
Field Name & Bits & Type / Fixed Value \\
\midrule
\texttt{src0} & \texttt{0:7} & \texttt{SourceEncoding64} \\
\hyperref[enum:tex_wide_indices_enable_m]{\texttt{wide\_indices}} & \texttt{8} & \texttt{exact=none} \\
\textit{Reserved} & \texttt{9} & \texttt{zero} \\
\hyperref[enum:tex_dual_array_enable_m]{\texttt{array\_enable}} & \texttt{10} & \texttt{tex\_dual\_array\_enable\_m (def none)} \\
\hyperref[enum:tex_texel_offset_m]{\texttt{texel\_offset}} & \texttt{11} & \texttt{tex\_texel\_offset\_m (def none)} \\
\hyperref[enum:tex_compare_enable_m]{\texttt{compare\_enable}} & \texttt{12} & \texttt{tex\_compare\_enable\_m (def none)} \\
\hyperref[enum:tex_dual_lod_mode_m]{\texttt{lod\_mode}} & \texttt{13:15} & \texttt{tex\_dual\_lod\_mode\_m} \\
\texttt{staging\_register\_write\_index} & \texttt{16:21} & \texttt{uint} \\
\hyperref[enum:tex_dual_primary_write_mask_m]{\texttt{primary\_write\_mask}} & \texttt{22:25} & \texttt{tex\_dual\_primary\_write\_mask\_m (def primary\_rgba)} \\
\textit{Reserved} & \texttt{26:27} & \texttt{ignore} \\
\hyperref[enum:tex_dimensionality_m]{\texttt{dimensionality}} & \texttt{28:29} & \texttt{tex\_dimensionality\_m} \\
\hyperref[enum:message_slot_index_m]{\texttt{message\_slot\_index}} & \texttt{30:32} & \texttt{message\_slot\_index\_m} \\
\texttt{sr\_count} & \texttt{33:35} & \texttt{sr\_count\_tex\_t} \\
\texttt{sr\_write\_count} & \texttt{36:38} & \texttt{sr\_write\_count\_t} \\
\hyperref[enum:tex_skip_m]{\texttt{skip}} & \texttt{39} & \texttt{tex\_skip\_m (def none)} \\
\texttt{staging\_register\_index} & \texttt{40:45} & \texttt{uint} \\
\hyperref[enum:tex_dual_primary_width_m]{\texttt{primary\_width}} & \texttt{46} & \texttt{tex\_dual\_primary\_width\_m} \\
\hyperref[enum:tex_dual_secondary_width_m]{\texttt{secondary\_width}} & \texttt{47} & \texttt{tex\_dual\_secondary\_width\_m} \\
\texttt{opcode1} & \texttt{48:51} & \texttt{exact=TEX\_DUAL} \\
\texttt{instruction\_class} & \texttt{52:56} & \texttt{exact=S2RW} \\
\texttt{fau\_page\_index} & \texttt{57:58} & \texttt{uint} \\
\hyperref[enum:flow_control_m]{\texttt{flow}} & \texttt{59:62} & \texttt{flow\_control\_m} \\
\textit{Reserved} & \texttt{63} & \texttt{zero} \\
\bottomrule
\end{tabularx}
\end{center}
\textbf{Fields Explained:}
\begin{tightitemize}\footnotesize
  \item \texttt{src0}: Primary source operand (\texttt{r0..r63}, \texttt{u0..u63} uniform, \texttt{k*} FAU constant, or \texttt{t0..t2} temporary; \texttt{\textasciicircum{}} marks discard/last-use). \emph{Example:} {\scriptsize\ttfamily \mnem{FADD.f32} \dst{r0}, \src{r1}, \src{r2}} (\textbf{SPIR-V:} \texttt{OpFAdd (Operand 0)}; \textbf{GLSL:} \texttt{float c = a + b;})
  \item \hyperref[enum:tex_wide_indices_enable_m]{\texttt{wide\_indices}}: Uses 32-bit wide texture/sampler table indices instead of compact 16-bit packed indices. \emph{Example:} {\scriptsize\ttfamily \mnem{TEX\_SINGLE}\imod{.tex2d.wide} \dst{@r0:r3}, \src{@r4:r5}, \src{r6}} (\textbf{SPIR-V:} \texttt{OpImageSampleImplicitLod (bindless)}; \textbf{GLSL:} \texttt{vec4 c = texture(textures[nonuniformEXT(i)], uv);})
  \item \hyperref[enum:tex_dual_array_enable_m]{\texttt{array\_enable}}: Enables array-layer coordinate in texture staging registers (\texttt{sampler2DArray} / \texttt{samplerCubeArray}). \emph{Example:} {\scriptsize\ttfamily \mnem{TEX\_SINGLE}\imod{.tex2d.array} \dst{@r0:r3}, \src{@r4:r6}} (\textbf{SPIR-V:} \texttt{OpImageSampleImplicitLod (Arrayed)}; \textbf{GLSL:} \texttt{vec4 c = texture(texArr, vec3(uv, layer));})
  \item \hyperref[enum:tex_texel_offset_m]{\texttt{texel\_offset}}: Applies a signed 4-bit per-axis texel offset $(dx, dy, dz)$ packed in staging registers. \emph{Example:} {\scriptsize\ttfamily \mnem{TEX\_SINGLE}\imod{.tex2d.texel\_offset} \dst{@r0:r3}, \src{@r4:r6}} (\textbf{SPIR-V:} \texttt{OpImageSampleImplicitLod (ConstOffset)}; \textbf{GLSL:} \texttt{vec4 c = textureOffset(tex, uv, ivec2(1, -1));})
  \item \hyperref[enum:tex_compare_enable_m]{\texttt{compare\_enable}}: Enables depth-reference shadow comparison (\texttt{Dref} in staging registers). \emph{Example:} {\scriptsize\ttfamily \mnem{TEX\_SINGLE}\imod{.tex2d.compare} \dst{@r0}, \src{@r4:r6}} (\textbf{SPIR-V:} \texttt{OpImageSampleDrefImplicitLod}; \textbf{GLSL:} \texttt{float s = texture(shadowMap, vec3(uv, refZ));})
  \item \hyperref[enum:tex_dual_lod_mode_m]{\texttt{lod\_mode}}: Selects mipmap LOD calculation mode (\texttt{none}, \texttt{computed}, \texttt{explicit}, \texttt{computed\_bias}, \texttt{grdesc}). \emph{Example:} {\scriptsize\ttfamily \mnem{TEX\_SINGLE}\imod{.tex2d.explicit} \dst{@r0:r3}, \src{@r4:r6}} (\textbf{SPIR-V:} \texttt{OpImageSampleExplicitLod}; \textbf{GLSL:} \texttt{vec4 c = textureLod(tex, uv, 2.0);})
  \item \texttt{staging\_register\_write\_index}: Separate destination staging register window base when read and write staging windows differ. \emph{Example:} {\scriptsize\ttfamily \mnem{ATOM\_RETURN.i32}\imod{.aadd} \dst{@r0}, \src{@r1}, \src{r2}, \imm{\#0}} (\textbf{SPIR-V:} \texttt{OpAtomicIAdd}; \textbf{GLSL:} \texttt{uint old = atomicAdd(ssbo.cnt, 1u);})
  \item \hyperref[enum:tex_dual_primary_write_mask_m]{\texttt{primary\_write\_mask}}: Channel mask (\texttt{primary\_r}..\texttt{primary\_rgba}) for the first texture fetch in \texttt{TEX\_DUAL}. \emph{Example:} {\scriptsize\ttfamily \mnem{TEX\_DUAL}\imod{.primary\_rgb.secondary\_r} \dst{@r0:r3}, \src{@r4:r5}} (\textbf{SPIR-V:} \texttt{2x OpImageSampleImplicitLod (fused)}; \textbf{GLSL:} \texttt{vec3 a = texture(t0, uv).rgb; float b = texture(t1, uv).r;})
  \item \hyperref[enum:tex_dimensionality_m]{\texttt{dimensionality}}: Texture dimensionality (\texttt{tex1d}, \texttt{tex2d}, \texttt{tex3d}, \texttt{cube}). \emph{Example:} {\scriptsize\ttfamily \mnem{TEX\_SINGLE}\imod{.tex3d.f32.rgba} \dst{@r0:r3}, \src{@r4:r6}} (\textbf{SPIR-V:} \texttt{OpImageSampleImplicitLod (Dim=3D)}; \textbf{GLSL:} \texttt{vec4 c = texture(volTex, uvw);})
  \item \hyperref[enum:message_slot_index_m]{\texttt{message\_slot\_index}}: Asynchronous scoreboard slot (\texttt{slot0}--\texttt{slot2}, \texttt{slot7}) tracking completion of this message. \emph{Example:} {\scriptsize\ttfamily \mnem{LOAD.i32}\imod{.slot0} \dst{@r0}, \src{r2}, \imm{\#0}} (\textbf{SPIR-V:} \texttt{OpLoad}; \textbf{GLSL:} \texttt{uint v = ssbo.arr[i];})
  \item \texttt{sr\_count}: Number of consecutive 32-bit staging registers transferred by a load/store/texture message. \emph{Example:} {\scriptsize\ttfamily \mnem{LD\_MULTI.i128} \dst{@r0:r3}, \src{r4}, \imm{\#0}} (\textbf{SPIR-V:} \texttt{OpLoad (wide vector)}; \textbf{GLSL:} \texttt{uvec4 x = buf[i];})
  \item \texttt{sr\_write\_count}: Number of consecutive 32-bit staging registers written back by a dual-staging message (e.g. ATOM\textbackslash\{\}\_RETURN / TEX). \emph{Example:} {\scriptsize\ttfamily \mnem{TEX\_SINGLE}\imod{.tex2d.f32.rgba} \dst{@r0:r3}, \src{@r4:r5}} (\textbf{SPIR-V:} \texttt{OpImageSampleImplicitLod}; \textbf{GLSL:} \texttt{vec4 c = texture(tex, uv);})
  \item \hyperref[enum:tex_skip_m]{\texttt{skip}}: Skips texture/memory fetch on helper/discarded fragment lanes while preserving 2x2 quad derivatives. \emph{Example:} {\scriptsize\ttfamily \mnem{TEX\_SINGLE}\imod{.tex2d.computed.skip} \dst{@r0:r3}, \src{@r4:r5}} (\textbf{SPIR-V:} \texttt{OpImageSampleImplicitLod (after OpDemoteToHelperInvocation)}; \textbf{GLSL:} \texttt{if (a \textless{} 0.0) discard; vec4 c = texture(s, uv);})
  \item \texttt{staging\_register\_index}: Base GPR index of the contiguous staging register window (\texttt{@rN:rM}) for message payload/return. \emph{Example:} {\scriptsize\ttfamily \mnem{STORE.i64}\imod{.slot7} \src{@r0:r1}, \src{r2}, \imm{\#8}} (\textbf{SPIR-V:} \texttt{OpStore (uvec2)}; \textbf{GLSL:} \texttt{ssbo.pairs[i] = uvec2(a, b);})
  \item \hyperref[enum:tex_dual_primary_width_m]{\texttt{primary\_width}}: Selects 16-bit vs 32-bit GPR element width for the corresponding message return window. \emph{Example:} {\scriptsize\ttfamily \mnem{TEX\_DUAL}\imod{.primary\_f16.secondary\_f32} \dst{@r0:r2}, \src{@r4:r5}} (\textbf{SPIR-V:} \texttt{OpImageSampleImplicitLod (RelaxedPrecision)}; \textbf{GLSL:} \texttt{mediump vec2 a = texture(t0, uv).rg;})
  \item \hyperref[enum:tex_dual_secondary_width_m]{\texttt{secondary\_width}}: Selects 16-bit vs 32-bit GPR element width for the corresponding message return window. \emph{Example:} {\scriptsize\ttfamily \mnem{TEX\_DUAL}\imod{.primary\_f16.secondary\_f32} \dst{@r0:r2}, \src{@r4:r5}} (\textbf{SPIR-V:} \texttt{OpImageSampleImplicitLod (RelaxedPrecision)}; \textbf{GLSL:} \texttt{mediump vec2 a = texture(t0, uv).rg;})
  \item \texttt{opcode1}: Primary sub-opcode within the instruction class uniquely identifying the operation family. \emph{Example:} {\scriptsize\ttfamily \mnem{IADD.s32} \dst{r0}, \src{r1}, \src{r2}} (\textbf{SPIR-V:} \texttt{OpIAdd (operation selector)}; \textbf{GLSL:} \texttt{// Primary opcode bits})
  \item \texttt{instruction\_class}: 4-bit major instruction class selector (\texttt{bits 48:51}) partitioning the 64-bit Valhall opcode space (e.g., \texttt{FMA}, \texttt{CVT}, \texttt{SFU}, \texttt{LS}, \texttt{TEX}, \texttt{BRANCH}). \emph{Example:} {\scriptsize\ttfamily \mnem{FADD.f32} \dst{r0}, \src{r1}, \src{r2}} (\textbf{SPIR-V:} \texttt{N/A (encoding class)}; \textbf{GLSL:} \texttt{// Hardware execution pipe / decode class})
  \item \texttt{fau\_page\_index}: Selects the FAU (Fast Access Uniform) page (\texttt{0..3}) when a source operand references hardware small constants, thread IDs, or lane constants. \emph{Example:} {\scriptsize\ttfamily \mnem{FADD.f32} \dst{r0}, \src{r1}, \src{k1f}} (\textbf{SPIR-V:} \texttt{OpConstant / BuiltIn}; \textbf{GLSL:} \texttt{float y = x + 1.0;})
  \item \hyperref[enum:flow_control_m]{\texttt{flow}}: 4-bit embedded flow-control \& scoreboard wait modifier (\texttt{bits 59:62}: \texttt{none}, \texttt{wait0..wait0126}, \texttt{wait}, \texttt{reconverge}, \texttt{discard}, \texttt{end}). \emph{Example:} {\scriptsize\ttfamily \mnem{FADD.f32}\imod{.wait0.reconverge} \dst{r0}, \src{r1}, \src{r2}} (\textbf{SPIR-V:} \texttt{OpControlBarrier / OpReturn}; \textbf{GLSL:} \texttt{// Scoreboard wait + warp reconvergence})
\end{tightitemize}
\clearpage

\subsection*{TEX\_FETCH}
\addcontentsline{toc}{subsection}{TEX\_FETCH}
Performs an unfiltered integer-coordinate texel fetch (\texttt{texelFetch}) from texture descriptor \texttt{src0} using integer coordinates \texttt{(x, [y], [z/layer], [lod/sample])} in \texttt{@sr}, writing \texttt{write\_mask} channels into \texttt{@sw}.
\par\textbf{Notes:}
\begin{tightitemize}
  \item No sampler filtering (\texttt{ISA.xml:148}).
  \item \textbf{Source references:} \texttt{kraid/ops.rs:4393-4435}, \texttt{bifrost/valhall/ISA.xml:148}.
\end{tightitemize}
\codebox{\mnem{TEX\_FETCH}\imod{\{.message\_slot\_index\}}\imod{\{.flow\}}\imod{\{.skip\}}\imod{\{.dimensionality\}}\imod{\{.register\_width\}}\imod{\{.write\_mask\}}\imod{\{.wide\_indices\}}\imod{\{.array\_enable\}}\imod{\{.texel\_offset\}} \dst{@sw}, \src{@sr}, \src{src0}}
\textbf{Operands:}
\begin{tightitemize}\footnotesize
  \item \texttt{\dst{@sw}} (\texttt{staging\_register\_write\_index} + \texttt{sr\_write\_count}): Destination staging register range receiving unfiltered texel components.
  \item \texttt{\src{@sr}} (\texttt{staging\_register\_index} + \texttt{sr\_count}): Source staging register range holding integer coordinates \texttt{(x, [y], [z/layer], [lod/sample])}.
  \item \texttt{\src{src0}}: Texture descriptor handle.
\end{tightitemize}
\par\textbf{Closest SPIR-V Op:} \texttt{OpImageFetch}
\par\textbf{Pseudo-GLSL (Operational Semantics):}
\codebox{@sw = texelFetch(handle = src0, icoords = @sr, flag\_array, flag\_texel\_offset);}
\textbf{Execution Context:}
\begin{tightitemize}\footnotesize
  \item \textbf{Texture \& Sampler chain (\hyperref[sec:csf-srt]{\texttt{SRT\_2}} Frag / \texttt{SRT\_0} Comp/Vert):} 64B-aligned VA (\texttt{\& \textasciitilde{}0x3F}) $\to$ \texttt{Resource[table]} $\to$ \hyperref[sec:csf-tex-samp]{\texttt{Texture[tex\_idx]}} (32B, \texttt{type=2}) $\to$ \texttt{GenericPlane[mip]} (32B, \texttt{type=10}: \texttt{pointer}, \texttt{row\_stride}, Morton 16$\times$16 / AFBC) + \texttt{Sampler[samp\_idx]} (32B, \texttt{type=1}: filter, wrap, LOD clamp/bias, compare).
  \item \textbf{Inline FAU \& Special Operands (\hyperref[sec:csf-tsd-fau]{\texttt{FAU\_k}}, \texttt{r8..r15}):} Any 8-bit source with \texttt{mode=2} reads zero-latency uniform \texttt{u0..u31} from \texttt{FAU\_k} (\texttt{slots64 = bits[63:56]}); \texttt{mode=3} reads \texttt{SmallConstant} (Table~1.2) or \texttt{FauSpecial} (\texttt{pc}, \texttt{lane\_id}, \texttt{lane\_width}, \texttt{tls\_ptr}, \texttt{wls\_ptr}, \texttt{blend\_descriptor\_0..7}).
\end{tightitemize}
\par\textbf{Worked Examples:}
\par\textbf{Example 1 (Unfiltered integer-coord 2D fetch):}
\codebox{\mnem{TEX\_FETCH}\imod{.tex2d.dst32.rgba.slot0} \dst{@r0:r3}, \src{@r4:r6}, \src{u0}}
Fetches exact texel at integer \texttt{(x=r4, y=r5, lod=r6)} via handle \texttt{u0} into 4x 32-bit \texttt{@r0:r3} without filtering.
\par\textbf{Example 2 (2D array texel fetch (`.array`)):}
\codebox{\mnem{TEX\_FETCH}\imod{.tex2d.array.dst32.rgba.slot0} \dst{@r0:r3}, \src{@r4:r7}, \src{u0}}
Adds integer array layer in \texttt{@r7} for \texttt{texelFetch(sampler2DArray, ivec3(x,y,layer), lod)}.
\par\textbf{Mined SPIR-V $\to$ Valhall ISA Example (Unfiltered integer texelFetch with explicit LOD 0 --- Vertex shader, Mali-G710 (v10)):}
\codebox{\textbf{; SPIR-V (input)}\\\%t = OpImageFetch \%v4float \%img \%ivec2\_xy Lod \%int\_0\\[3pt]\textbf{; Valhall ISA (malioc output)}\\TEX\_FETCH.slot0.tex2d.dst32.rgba @r20:r21:r22:r23, @r10:r11, [r4:r5]}
\texttt{OpImageFetch} lowers directly to \texttt{TEX\_FETCH.tex2d.dst32.rgba}, reading integer coordinates from \texttt{@r10:r11} and the image descriptor handle from \texttt{[r4:r5]} without filtering.
\begin{center}\footnotesize
\begin{tabularx}{\linewidth}{@{}lllX@{}}
\toprule
Variant & Arch & Class & Opcode Fields \\
\midrule
\texttt{(none)} & v9--v14 & \texttt{S2RW} & \texttt{opcode1=TEX\_FETCH} \\
\texttt{(none)} & v9--v14 & \texttt{S2RW} & \texttt{opcode1=TEX\_FETCH} \\
\bottomrule
\end{tabularx}
\end{center}
\begin{center}\footnotesize
\begin{tabularx}{0.88\linewidth}{@{}l l X@{}}
\toprule
Field Name & Bits & Type / Fixed Value \\
\midrule
\texttt{src0} & \texttt{0:7} & \texttt{SourceEncoding64} \\
\hyperref[enum:tex_wide_indices_enable_m]{\texttt{wide\_indices}} & \texttt{8} & \texttt{tex\_wide\_indices\_enable\_m (def none)} \\
\textit{Reserved} & \texttt{9} & \texttt{zero} \\
\hyperref[enum:tex_array_enable_m]{\texttt{array\_enable}} & \texttt{10} & \texttt{tex\_array\_enable\_m (def none)} \\
\hyperref[enum:tex_texel_offset_si_m]{\texttt{texel\_offset}} & \texttt{11} & \texttt{tex\_texel\_offset\_si\_m (def none)} \\
\textit{Reserved} & \texttt{12:15} & \texttt{zero} \\
\texttt{staging\_register\_write\_index} & \texttt{16:21} & \texttt{uint} \\
\hyperref[enum:tex_write_mask_m]{\texttt{write\_mask}} & \texttt{22:25} & \texttt{tex\_write\_mask\_m (def rgba)} \\
\textit{Reserved} & \texttt{26:27} & \texttt{ignore} \\
\hyperref[enum:tex_dimensionality_m]{\texttt{dimensionality}} & \texttt{28:29} & \texttt{tex\_dimensionality\_m} \\
\hyperref[enum:message_slot_index_m]{\texttt{message\_slot\_index}} & \texttt{30:32} & \texttt{message\_slot\_index\_m} \\
\texttt{sr\_count} & \texttt{33:35} & \texttt{sr\_count\_tex\_t} \\
\texttt{sr\_write\_count} & \texttt{36:38} & \texttt{sr\_write\_count\_t} \\
\hyperref[enum:tex_skip_m]{\texttt{skip}} & \texttt{39} & \texttt{tex\_skip\_m (def none)} \\
\texttt{staging\_register\_index} & \texttt{40:45} & \texttt{uint} \\
\hyperref[enum:tex_width_m]{\texttt{register\_width}} & \texttt{46} & \texttt{tex\_width\_m} \\
\texttt{secondary\_width} & \texttt{47} & \texttt{exact=0} \\
\texttt{opcode1} & \texttt{48:51} & \texttt{exact=TEX\_FETCH} \\
\texttt{instruction\_class} & \texttt{52:56} & \texttt{exact=S2RW} \\
\texttt{fau\_page\_index} & \texttt{57:58} & \texttt{uint} \\
\hyperref[enum:flow_control_m]{\texttt{flow}} & \texttt{59:62} & \texttt{flow\_control\_m} \\
\textit{Reserved} & \texttt{63} & \texttt{zero} \\
\bottomrule
\end{tabularx}
\end{center}
\textbf{Fields Explained:}
\begin{tightitemize}\footnotesize
  \item \texttt{src0}: Primary source operand (\texttt{r0..r63}, \texttt{u0..u63} uniform, \texttt{k*} FAU constant, or \texttt{t0..t2} temporary; \texttt{\textasciicircum{}} marks discard/last-use). \emph{Example:} {\scriptsize\ttfamily \mnem{FADD.f32} \dst{r0}, \src{r1}, \src{r2}} (\textbf{SPIR-V:} \texttt{OpFAdd (Operand 0)}; \textbf{GLSL:} \texttt{float c = a + b;})
  \item \hyperref[enum:tex_wide_indices_enable_m]{\texttt{wide\_indices}}: Uses 32-bit wide texture/sampler table indices instead of compact 16-bit packed indices. \emph{Example:} {\scriptsize\ttfamily \mnem{TEX\_SINGLE}\imod{.tex2d.wide} \dst{@r0:r3}, \src{@r4:r5}, \src{r6}} (\textbf{SPIR-V:} \texttt{OpImageSampleImplicitLod (bindless)}; \textbf{GLSL:} \texttt{vec4 c = texture(textures[nonuniformEXT(i)], uv);})
  \item \hyperref[enum:tex_array_enable_m]{\texttt{array\_enable}}: Enables array-layer coordinate in texture staging registers (\texttt{sampler2DArray} / \texttt{samplerCubeArray}). \emph{Example:} {\scriptsize\ttfamily \mnem{TEX\_SINGLE}\imod{.tex2d.array} \dst{@r0:r3}, \src{@r4:r6}} (\textbf{SPIR-V:} \texttt{OpImageSampleImplicitLod (Arrayed)}; \textbf{GLSL:} \texttt{vec4 c = texture(texArr, vec3(uv, layer));})
  \item \hyperref[enum:tex_texel_offset_si_m]{\texttt{texel\_offset}}: Applies a signed 4-bit per-axis texel offset $(dx, dy, dz)$ packed in staging registers. \emph{Example:} {\scriptsize\ttfamily \mnem{TEX\_SINGLE}\imod{.tex2d.texel\_offset} \dst{@r0:r3}, \src{@r4:r6}} (\textbf{SPIR-V:} \texttt{OpImageSampleImplicitLod (ConstOffset)}; \textbf{GLSL:} \texttt{vec4 c = textureOffset(tex, uv, ivec2(1, -1));})
  \item \texttt{staging\_register\_write\_index}: Separate destination staging register window base when read and write staging windows differ. \emph{Example:} {\scriptsize\ttfamily \mnem{ATOM\_RETURN.i32}\imod{.aadd} \dst{@r0}, \src{@r1}, \src{r2}, \imm{\#0}} (\textbf{SPIR-V:} \texttt{OpAtomicIAdd}; \textbf{GLSL:} \texttt{uint old = atomicAdd(ssbo.cnt, 1u);})
  \item \hyperref[enum:tex_write_mask_m]{\texttt{write\_mask}}: Channel enable bitmask (\texttt{r}..\texttt{rgba}) controlling which texture components are written to staging GPRs. \emph{Example:} {\scriptsize\ttfamily \mnem{TEX\_SINGLE}\imod{.tex2d.f32.rg} \dst{@r0:r1}, \src{@r4:r5}} (\textbf{SPIR-V:} \texttt{OpImageSampleImplicitLod + OpVectorShuffle}; \textbf{GLSL:} \texttt{vec2 rg = texture(tex, uv).rg;})
  \item \hyperref[enum:tex_dimensionality_m]{\texttt{dimensionality}}: Texture dimensionality (\texttt{tex1d}, \texttt{tex2d}, \texttt{tex3d}, \texttt{cube}). \emph{Example:} {\scriptsize\ttfamily \mnem{TEX\_SINGLE}\imod{.tex3d.f32.rgba} \dst{@r0:r3}, \src{@r4:r6}} (\textbf{SPIR-V:} \texttt{OpImageSampleImplicitLod (Dim=3D)}; \textbf{GLSL:} \texttt{vec4 c = texture(volTex, uvw);})
  \item \hyperref[enum:message_slot_index_m]{\texttt{message\_slot\_index}}: Asynchronous scoreboard slot (\texttt{slot0}--\texttt{slot2}, \texttt{slot7}) tracking completion of this message. \emph{Example:} {\scriptsize\ttfamily \mnem{LOAD.i32}\imod{.slot0} \dst{@r0}, \src{r2}, \imm{\#0}} (\textbf{SPIR-V:} \texttt{OpLoad}; \textbf{GLSL:} \texttt{uint v = ssbo.arr[i];})
  \item \texttt{sr\_count}: Number of consecutive 32-bit staging registers transferred by a load/store/texture message. \emph{Example:} {\scriptsize\ttfamily \mnem{LD\_MULTI.i128} \dst{@r0:r3}, \src{r4}, \imm{\#0}} (\textbf{SPIR-V:} \texttt{OpLoad (wide vector)}; \textbf{GLSL:} \texttt{uvec4 x = buf[i];})
  \item \texttt{sr\_write\_count}: Number of consecutive 32-bit staging registers written back by a dual-staging message (e.g. ATOM\textbackslash\{\}\_RETURN / TEX). \emph{Example:} {\scriptsize\ttfamily \mnem{TEX\_SINGLE}\imod{.tex2d.f32.rgba} \dst{@r0:r3}, \src{@r4:r5}} (\textbf{SPIR-V:} \texttt{OpImageSampleImplicitLod}; \textbf{GLSL:} \texttt{vec4 c = texture(tex, uv);})
  \item \hyperref[enum:tex_skip_m]{\texttt{skip}}: Skips texture/memory fetch on helper/discarded fragment lanes while preserving 2x2 quad derivatives. \emph{Example:} {\scriptsize\ttfamily \mnem{TEX\_SINGLE}\imod{.tex2d.computed.skip} \dst{@r0:r3}, \src{@r4:r5}} (\textbf{SPIR-V:} \texttt{OpImageSampleImplicitLod (after OpDemoteToHelperInvocation)}; \textbf{GLSL:} \texttt{if (a \textless{} 0.0) discard; vec4 c = texture(s, uv);})
  \item \texttt{staging\_register\_index}: Base GPR index of the contiguous staging register window (\texttt{@rN:rM}) for message payload/return. \emph{Example:} {\scriptsize\ttfamily \mnem{STORE.i64}\imod{.slot7} \src{@r0:r1}, \src{r2}, \imm{\#8}} (\textbf{SPIR-V:} \texttt{OpStore (uvec2)}; \textbf{GLSL:} \texttt{ssbo.pairs[i] = uvec2(a, b);})
  \item \hyperref[enum:tex_width_m]{\texttt{register\_width}}: Selects 16-bit vs 32-bit GPR element width for the corresponding message return window. \emph{Example:} {\scriptsize\ttfamily \mnem{TEX\_DUAL}\imod{.primary\_f16.secondary\_f32} \dst{@r0:r2}, \src{@r4:r5}} (\textbf{SPIR-V:} \texttt{OpImageSampleImplicitLod (RelaxedPrecision)}; \textbf{GLSL:} \texttt{mediump vec2 a = texture(t0, uv).rg;})
  \item \texttt{secondary\_width}: Selects 16-bit vs 32-bit GPR element width for the corresponding message return window. \emph{Example:} {\scriptsize\ttfamily \mnem{TEX\_DUAL}\imod{.primary\_f16.secondary\_f32} \dst{@r0:r2}, \src{@r4:r5}} (\textbf{SPIR-V:} \texttt{OpImageSampleImplicitLod (RelaxedPrecision)}; \textbf{GLSL:} \texttt{mediump vec2 a = texture(t0, uv).rg;})
  \item \texttt{opcode1}: Primary sub-opcode within the instruction class uniquely identifying the operation family. \emph{Example:} {\scriptsize\ttfamily \mnem{IADD.s32} \dst{r0}, \src{r1}, \src{r2}} (\textbf{SPIR-V:} \texttt{OpIAdd (operation selector)}; \textbf{GLSL:} \texttt{// Primary opcode bits})
  \item \texttt{instruction\_class}: 4-bit major instruction class selector (\texttt{bits 48:51}) partitioning the 64-bit Valhall opcode space (e.g., \texttt{FMA}, \texttt{CVT}, \texttt{SFU}, \texttt{LS}, \texttt{TEX}, \texttt{BRANCH}). \emph{Example:} {\scriptsize\ttfamily \mnem{FADD.f32} \dst{r0}, \src{r1}, \src{r2}} (\textbf{SPIR-V:} \texttt{N/A (encoding class)}; \textbf{GLSL:} \texttt{// Hardware execution pipe / decode class})
  \item \texttt{fau\_page\_index}: Selects the FAU (Fast Access Uniform) page (\texttt{0..3}) when a source operand references hardware small constants, thread IDs, or lane constants. \emph{Example:} {\scriptsize\ttfamily \mnem{FADD.f32} \dst{r0}, \src{r1}, \src{k1f}} (\textbf{SPIR-V:} \texttt{OpConstant / BuiltIn}; \textbf{GLSL:} \texttt{float y = x + 1.0;})
  \item \hyperref[enum:flow_control_m]{\texttt{flow}}: 4-bit embedded flow-control \& scoreboard wait modifier (\texttt{bits 59:62}: \texttt{none}, \texttt{wait0..wait0126}, \texttt{wait}, \texttt{reconverge}, \texttt{discard}, \texttt{end}). \emph{Example:} {\scriptsize\ttfamily \mnem{FADD.f32}\imod{.wait0.reconverge} \dst{r0}, \src{r1}, \src{r2}} (\textbf{SPIR-V:} \texttt{OpControlBarrier / OpReturn}; \textbf{GLSL:} \texttt{// Scoreboard wait + warp reconvergence})
\end{tightitemize}
\clearpage

\subsection*{TEX\_GATHER}
\addcontentsline{toc}{subsection}{TEX\_GATHER}
Samples the 2x2 bilinear footprint around floating-point UV coordinates in \texttt{@sr} from texture descriptor \texttt{src0} and returns the 4 values of the selected \texttt{gather\_component} (\texttt{r}/\texttt{g}/\texttt{b}/\texttt{a}, or shadow compare results when \texttt{.compare} is set) into \texttt{@sw}.
\par\textbf{Notes:}
\begin{tightitemize}
  \item 2D/Cube only (\texttt{ISA.xml:150}).
  \item \textbf{Source references:} \texttt{kraid/ops.rs:4437-4482}, \texttt{bifrost/valhall/ISA.xml:150}.
\end{tightitemize}
\codebox{\mnem{TEX\_GATHER}\imod{\{.message\_slot\_index\}}\imod{\{.flow\}}\imod{\{.skip\}}\imod{\{.dimensionality\}}\imod{\{.register\_width\}}\imod{\{.write\_mask\}}\imod{\{.wide\_indices\}}\imod{\{.array\_enable\}}\imod{\{.texel\_offset\}}\imod{\{.compare\_enable\}}\imod{\{.coordinate\_mode\}}\imod{\{.gather\_component\}} \dst{@sw}, \src{@sr}, \src{src0}\\\mnem{TEX\_GATHER}\imod{\{.message\_slot\_index\}}\imod{\{.flow\}}\imod{\{.skip\}}\imod{\{.dimensionality\}}\imod{\{.projection\_enable\}}\imod{\{.register\_width\}}\imod{\{.write\_mask\}}\imod{\{.wide\_indices\}}\imod{\{.array\_enable\}}\imod{\{.texel\_offset\}}\imod{\{.compare\_enable\}}\imod{\{.coordinate\_mode\}}\imod{\{.gather\_component\}} \dst{@sw}, \src{@sr}, \src{src0}}
\textbf{Operands:}
\begin{tightitemize}\footnotesize
  \item \texttt{\dst{@sw}} (\texttt{staging\_register\_write\_index} + \texttt{sr\_write\_count}): Destination staging register range receiving 4 gathered values from the 2x2 texel footprint.
  \item \texttt{\src{@sr}} (\texttt{staging\_register\_index} + \texttt{sr\_count}): Source staging register range holding \texttt{(u, v, [layer], [dref], [offset])}.
  \item \texttt{\src{src0}}: Texture/sampler descriptor handle.
\end{tightitemize}
\par\textbf{Closest SPIR-V Op:} \texttt{OpImageGather / OpImageDrefGather}
\par\textbf{Pseudo-GLSL (Operational Semantics):}
\codebox{if (flag\_compare) @sw = textureGatherCompare(handle = src0, uv\_dref = @sr);\\else              @sw = textureGather(handle = src0, uv = @sr, gather\_component);}
\textbf{Execution Context:}
\begin{tightitemize}\footnotesize
  \item \textbf{Texture \& Sampler chain (\hyperref[sec:csf-srt]{\texttt{SRT\_2}} Frag / \texttt{SRT\_0} Comp/Vert):} 64B-aligned VA (\texttt{\& \textasciitilde{}0x3F}) $\to$ \texttt{Resource[table]} $\to$ \hyperref[sec:csf-tex-samp]{\texttt{Texture[tex\_idx]}} (32B, \texttt{type=2}) $\to$ \texttt{GenericPlane[mip]} (32B, \texttt{type=10}: \texttt{pointer}, \texttt{row\_stride}, Morton 16$\times$16 / AFBC) + \texttt{Sampler[samp\_idx]} (32B, \texttt{type=1}: filter, wrap, LOD clamp/bias, compare).
  \item \textbf{Fragment stage preloads \& FP modes (\hyperref[sec:csf-spd]{\texttt{SPD\_2}}, \texttt{r20:r21}):} Hardware pre-fills \texttt{r59} with \texttt{floatBitsToUint(FragCoord.z)} and \texttt{r60} with the MSAA sample coverage mask (\texttt{gl\_SampleMaskIn[0]}), while \texttt{ShaderProgram} word~0 latches 32/64 GPR allocation and warp-wide \texttt{flush\_to\_zero\_mode}, \texttt{suppress\_na\_n}, and \texttt{suppress\_inf}.
\end{tightitemize}
\par\textbf{Worked Examples:}
\par\textbf{Example 1 (2x2 red-channel gather (`.gather4\_r`)):}
\codebox{\mnem{TEX\_GATHER}\imod{.tex2d.gather4\_r.dst32.rgba.slot0} \dst{@r0:r3}, \src{@r4:r5}, \src{u0}}
Gathers the red channel of the 4 texels in the 2x2 bilinear footprint around \texttt{uv = (@r4, @r5)} into \texttt{@r0:r3}.
\par\textbf{Example 2 (Shadow depth-compare gather (`.compare`)):}
\codebox{\mnem{TEX\_GATHER}\imod{.tex2d.compare.dst32.rgba.slot0} \dst{@r0:r3}, \src{@r4:r6}, \src{u0}}
Compares the 2x2 depth footprint against \texttt{Dref = @r6}, returning 4 binary PCF results (\texttt{0.0} or \texttt{1.0}) in \texttt{@r0:r3}.
\par\textbf{Mined SPIR-V $\to$ Valhall ISA Example (2x2 red-channel textureGather for min/max depth reduction --- Fragment shader, Mali-G710 (v10)):}
\codebox{\textbf{; SPIR-V (input)}\\\%g = OpImageGather \%v4float \%sampled\_img \%uv \%int\_0 None\\[3pt]\textbf{; Valhall ISA (malioc output)}\\TEX\_GATHER.slot0.wait0.skip.tex2d.dst32.rgba.wide.float\_coordinates.gather4\_r @r0:r1:r2:r3, @r0:r1, u0}
\texttt{OpImageGather} with component \texttt{0} (\texttt{R}) lowers to \texttt{TEX\_GATHER} with \texttt{.gather4\_r} and \texttt{.float\_coordinates}, returning the 4 footprint red texels into \texttt{@r0:r1:r2:r3}.
\begin{center}\footnotesize
\begin{tabularx}{\linewidth}{@{}lllX@{}}
\toprule
Variant & Arch & Class & Opcode Fields \\
\midrule
\texttt{(none)} & v9--v14 & \texttt{S2RW} & \texttt{opcode1=TEX\_GATHER} \\
\texttt{(none)} & v9--v14 & \texttt{S2RW} & \texttt{opcode1=TEX\_GATHER} \\
\bottomrule
\end{tabularx}
\end{center}
\begin{center}\footnotesize
\begin{tabularx}{0.88\linewidth}{@{}l l X@{}}
\toprule
Field Name & Bits & Type / Fixed Value \\
\midrule
\texttt{src0} & \texttt{0:7} & \texttt{SourceEncoding64} \\
\hyperref[enum:tex_wide_indices_enable_m]{\texttt{wide\_indices}} & \texttt{8} & \texttt{tex\_wide\_indices\_enable\_m (def none)} \\
\textit{Reserved} & \texttt{9} & \texttt{zero} \\
\hyperref[enum:tex_array_enable_m]{\texttt{array\_enable}} & \texttt{10} & \texttt{tex\_array\_enable\_m (def none)} \\
\hyperref[enum:tex_texel_offset_m]{\texttt{texel\_offset}} & \texttt{11} & \texttt{tex\_texel\_offset\_m (def none)} \\
\hyperref[enum:tex_compare_enable_m]{\texttt{compare\_enable}} & \texttt{12} & \texttt{tex\_compare\_enable\_m (def none)} \\
\hyperref[enum:tex_coordinate_mode_m]{\texttt{coordinate\_mode}} & \texttt{13} & \texttt{tex\_coordinate\_mode\_m (def float\_coordinates)} \\
\hyperref[enum:tex_gather_component_m]{\texttt{gather\_component}} & \texttt{14:15} & \texttt{tex\_gather\_component\_m} \\
\texttt{staging\_register\_write\_index} & \texttt{16:21} & \texttt{uint} \\
\hyperref[enum:tex_write_mask_m]{\texttt{write\_mask}} & \texttt{22:25} & \texttt{tex\_write\_mask\_m (def rgba)} \\
\textit{Reserved} & \texttt{26:27} & \texttt{ignore} \\
\hyperref[enum:tex_dimensionality_m]{\texttt{dimensionality}} & \texttt{28:29} & \texttt{tex\_dimensionality\_m} \\
\hyperref[enum:message_slot_index_m]{\texttt{message\_slot\_index}} & \texttt{30:32} & \texttt{message\_slot\_index\_m} \\
\texttt{sr\_count} & \texttt{33:35} & \texttt{sr\_count\_tex\_t} \\
\texttt{sr\_write\_count} & \texttt{36:38} & \texttt{sr\_write\_count\_t} \\
\hyperref[enum:tex_skip_m]{\texttt{skip}} & \texttt{39} & \texttt{tex\_skip\_m (def none)} \\
\texttt{staging\_register\_index} & \texttt{40:45} & \texttt{uint} \\
\hyperref[enum:tex_width_m]{\texttt{register\_width}} & \texttt{46} & \texttt{tex\_width\_m} \\
\texttt{secondary\_width} & \texttt{47} & \texttt{exact=0} \\
\texttt{opcode1} & \texttt{48:51} & \texttt{exact=TEX\_GATHER} \\
\texttt{instruction\_class} & \texttt{52:56} & \texttt{exact=S2RW} \\
\texttt{fau\_page\_index} & \texttt{57:58} & \texttt{uint} \\
\hyperref[enum:flow_control_m]{\texttt{flow}} & \texttt{59:62} & \texttt{flow\_control\_m} \\
\textit{Reserved} & \texttt{63} & \texttt{zero} \\
\bottomrule
\end{tabularx}
\end{center}
\textbf{Fields Explained:}
\begin{tightitemize}\footnotesize
  \item \texttt{src0}: Primary source operand (\texttt{r0..r63}, \texttt{u0..u63} uniform, \texttt{k*} FAU constant, or \texttt{t0..t2} temporary; \texttt{\textasciicircum{}} marks discard/last-use). \emph{Example:} {\scriptsize\ttfamily \mnem{FADD.f32} \dst{r0}, \src{r1}, \src{r2}} (\textbf{SPIR-V:} \texttt{OpFAdd (Operand 0)}; \textbf{GLSL:} \texttt{float c = a + b;})
  \item \hyperref[enum:tex_wide_indices_enable_m]{\texttt{wide\_indices}}: Uses 32-bit wide texture/sampler table indices instead of compact 16-bit packed indices. \emph{Example:} {\scriptsize\ttfamily \mnem{TEX\_SINGLE}\imod{.tex2d.wide} \dst{@r0:r3}, \src{@r4:r5}, \src{r6}} (\textbf{SPIR-V:} \texttt{OpImageSampleImplicitLod (bindless)}; \textbf{GLSL:} \texttt{vec4 c = texture(textures[nonuniformEXT(i)], uv);})
  \item \hyperref[enum:tex_array_enable_m]{\texttt{array\_enable}}: Enables array-layer coordinate in texture staging registers (\texttt{sampler2DArray} / \texttt{samplerCubeArray}). \emph{Example:} {\scriptsize\ttfamily \mnem{TEX\_SINGLE}\imod{.tex2d.array} \dst{@r0:r3}, \src{@r4:r6}} (\textbf{SPIR-V:} \texttt{OpImageSampleImplicitLod (Arrayed)}; \textbf{GLSL:} \texttt{vec4 c = texture(texArr, vec3(uv, layer));})
  \item \hyperref[enum:tex_texel_offset_m]{\texttt{texel\_offset}}: Applies a signed 4-bit per-axis texel offset $(dx, dy, dz)$ packed in staging registers. \emph{Example:} {\scriptsize\ttfamily \mnem{TEX\_SINGLE}\imod{.tex2d.texel\_offset} \dst{@r0:r3}, \src{@r4:r6}} (\textbf{SPIR-V:} \texttt{OpImageSampleImplicitLod (ConstOffset)}; \textbf{GLSL:} \texttt{vec4 c = textureOffset(tex, uv, ivec2(1, -1));})
  \item \hyperref[enum:tex_compare_enable_m]{\texttt{compare\_enable}}: Enables depth-reference shadow comparison (\texttt{Dref} in staging registers). \emph{Example:} {\scriptsize\ttfamily \mnem{TEX\_SINGLE}\imod{.tex2d.compare} \dst{@r0}, \src{@r4:r6}} (\textbf{SPIR-V:} \texttt{OpImageSampleDrefImplicitLod}; \textbf{GLSL:} \texttt{float s = texture(shadowMap, vec3(uv, refZ));})
  \item \hyperref[enum:tex_coordinate_mode_m]{\texttt{coordinate\_mode}}: Selects normalized float coordinates vs unnormalized integer texel indices. \emph{Example:} {\scriptsize\ttfamily \mnem{TEX\_FETCH}\imod{.tex2d.integer\_coordinates} \dst{@r0:r3}, \src{@r4:r5}} (\textbf{SPIR-V:} \texttt{OpImageFetch}; \textbf{GLSL:} \texttt{vec4 c = texelFetch(tex, ivec2(xy), 0);})
  \item \hyperref[enum:tex_gather_component_m]{\texttt{gather\_component}}: Selects the R, G, B, or A channel gathered across the 2x2 texel footprint. \emph{Example:} {\scriptsize\ttfamily \mnem{TEX\_GATHER}\imod{.tex2d.gather4\_g} \dst{@r0:r3}, \src{@r4:r5}} (\textbf{SPIR-V:} \texttt{OpImageGather}; \textbf{GLSL:} \texttt{vec4 g = textureGather(tex, uv, 1);})
  \item \texttt{staging\_register\_write\_index}: Separate destination staging register window base when read and write staging windows differ. \emph{Example:} {\scriptsize\ttfamily \mnem{ATOM\_RETURN.i32}\imod{.aadd} \dst{@r0}, \src{@r1}, \src{r2}, \imm{\#0}} (\textbf{SPIR-V:} \texttt{OpAtomicIAdd}; \textbf{GLSL:} \texttt{uint old = atomicAdd(ssbo.cnt, 1u);})
  \item \hyperref[enum:tex_write_mask_m]{\texttt{write\_mask}}: Channel enable bitmask (\texttt{r}..\texttt{rgba}) controlling which texture components are written to staging GPRs. \emph{Example:} {\scriptsize\ttfamily \mnem{TEX\_SINGLE}\imod{.tex2d.f32.rg} \dst{@r0:r1}, \src{@r4:r5}} (\textbf{SPIR-V:} \texttt{OpImageSampleImplicitLod + OpVectorShuffle}; \textbf{GLSL:} \texttt{vec2 rg = texture(tex, uv).rg;})
  \item \hyperref[enum:tex_dimensionality_m]{\texttt{dimensionality}}: Texture dimensionality (\texttt{tex1d}, \texttt{tex2d}, \texttt{tex3d}, \texttt{cube}). \emph{Example:} {\scriptsize\ttfamily \mnem{TEX\_SINGLE}\imod{.tex3d.f32.rgba} \dst{@r0:r3}, \src{@r4:r6}} (\textbf{SPIR-V:} \texttt{OpImageSampleImplicitLod (Dim=3D)}; \textbf{GLSL:} \texttt{vec4 c = texture(volTex, uvw);})
  \item \hyperref[enum:message_slot_index_m]{\texttt{message\_slot\_index}}: Asynchronous scoreboard slot (\texttt{slot0}--\texttt{slot2}, \texttt{slot7}) tracking completion of this message. \emph{Example:} {\scriptsize\ttfamily \mnem{LOAD.i32}\imod{.slot0} \dst{@r0}, \src{r2}, \imm{\#0}} (\textbf{SPIR-V:} \texttt{OpLoad}; \textbf{GLSL:} \texttt{uint v = ssbo.arr[i];})
  \item \texttt{sr\_count}: Number of consecutive 32-bit staging registers transferred by a load/store/texture message. \emph{Example:} {\scriptsize\ttfamily \mnem{LD\_MULTI.i128} \dst{@r0:r3}, \src{r4}, \imm{\#0}} (\textbf{SPIR-V:} \texttt{OpLoad (wide vector)}; \textbf{GLSL:} \texttt{uvec4 x = buf[i];})
  \item \texttt{sr\_write\_count}: Number of consecutive 32-bit staging registers written back by a dual-staging message (e.g. ATOM\textbackslash\{\}\_RETURN / TEX). \emph{Example:} {\scriptsize\ttfamily \mnem{TEX\_SINGLE}\imod{.tex2d.f32.rgba} \dst{@r0:r3}, \src{@r4:r5}} (\textbf{SPIR-V:} \texttt{OpImageSampleImplicitLod}; \textbf{GLSL:} \texttt{vec4 c = texture(tex, uv);})
  \item \hyperref[enum:tex_skip_m]{\texttt{skip}}: Skips texture/memory fetch on helper/discarded fragment lanes while preserving 2x2 quad derivatives. \emph{Example:} {\scriptsize\ttfamily \mnem{TEX\_SINGLE}\imod{.tex2d.computed.skip} \dst{@r0:r3}, \src{@r4:r5}} (\textbf{SPIR-V:} \texttt{OpImageSampleImplicitLod (after OpDemoteToHelperInvocation)}; \textbf{GLSL:} \texttt{if (a \textless{} 0.0) discard; vec4 c = texture(s, uv);})
  \item \texttt{staging\_register\_index}: Base GPR index of the contiguous staging register window (\texttt{@rN:rM}) for message payload/return. \emph{Example:} {\scriptsize\ttfamily \mnem{STORE.i64}\imod{.slot7} \src{@r0:r1}, \src{r2}, \imm{\#8}} (\textbf{SPIR-V:} \texttt{OpStore (uvec2)}; \textbf{GLSL:} \texttt{ssbo.pairs[i] = uvec2(a, b);})
  \item \hyperref[enum:tex_width_m]{\texttt{register\_width}}: Selects 16-bit vs 32-bit GPR element width for the corresponding message return window. \emph{Example:} {\scriptsize\ttfamily \mnem{TEX\_DUAL}\imod{.primary\_f16.secondary\_f32} \dst{@r0:r2}, \src{@r4:r5}} (\textbf{SPIR-V:} \texttt{OpImageSampleImplicitLod (RelaxedPrecision)}; \textbf{GLSL:} \texttt{mediump vec2 a = texture(t0, uv).rg;})
  \item \texttt{secondary\_width}: Selects 16-bit vs 32-bit GPR element width for the corresponding message return window. \emph{Example:} {\scriptsize\ttfamily \mnem{TEX\_DUAL}\imod{.primary\_f16.secondary\_f32} \dst{@r0:r2}, \src{@r4:r5}} (\textbf{SPIR-V:} \texttt{OpImageSampleImplicitLod (RelaxedPrecision)}; \textbf{GLSL:} \texttt{mediump vec2 a = texture(t0, uv).rg;})
  \item \texttt{opcode1}: Primary sub-opcode within the instruction class uniquely identifying the operation family. \emph{Example:} {\scriptsize\ttfamily \mnem{IADD.s32} \dst{r0}, \src{r1}, \src{r2}} (\textbf{SPIR-V:} \texttt{OpIAdd (operation selector)}; \textbf{GLSL:} \texttt{// Primary opcode bits})
  \item \texttt{instruction\_class}: 4-bit major instruction class selector (\texttt{bits 48:51}) partitioning the 64-bit Valhall opcode space (e.g., \texttt{FMA}, \texttt{CVT}, \texttt{SFU}, \texttt{LS}, \texttt{TEX}, \texttt{BRANCH}). \emph{Example:} {\scriptsize\ttfamily \mnem{FADD.f32} \dst{r0}, \src{r1}, \src{r2}} (\textbf{SPIR-V:} \texttt{N/A (encoding class)}; \textbf{GLSL:} \texttt{// Hardware execution pipe / decode class})
  \item \texttt{fau\_page\_index}: Selects the FAU (Fast Access Uniform) page (\texttt{0..3}) when a source operand references hardware small constants, thread IDs, or lane constants. \emph{Example:} {\scriptsize\ttfamily \mnem{FADD.f32} \dst{r0}, \src{r1}, \src{k1f}} (\textbf{SPIR-V:} \texttt{OpConstant / BuiltIn}; \textbf{GLSL:} \texttt{float y = x + 1.0;})
  \item \hyperref[enum:flow_control_m]{\texttt{flow}}: 4-bit embedded flow-control \& scoreboard wait modifier (\texttt{bits 59:62}: \texttt{none}, \texttt{wait0..wait0126}, \texttt{wait}, \texttt{reconverge}, \texttt{discard}, \texttt{end}). \emph{Example:} {\scriptsize\ttfamily \mnem{FADD.f32}\imod{.wait0.reconverge} \dst{r0}, \src{r1}, \src{r2}} (\textbf{SPIR-V:} \texttt{OpControlBarrier / OpReturn}; \textbf{GLSL:} \texttt{// Scoreboard wait + warp reconvergence})
\end{tightitemize}
\clearpage

\subsection*{TEX\_GRADIENT}
\addcontentsline{toc}{subsection}{TEX\_GRADIENT}
Samples texture descriptor \texttt{src0} using explicit \texttt{(dPdx, dPdy)} screen-space gradients packed in \texttt{@sr} (\texttt{textureGrad}), or---when \texttt{.disable\_lod\_clamp} / \texttt{.derivative} is set---returns the hardware-computed 8.8 fixed-point mipmap LOD (\texttt{textureQueryLod}) or derivatives into \texttt{@sw}.
\par\textbf{Notes:}
\begin{tightitemize}
  \item Supports explicit gradient sampling (\texttt{ISA.xml:152}).
  \item \textbf{Source references:} \texttt{kraid/ops.rs:4504-4549}, \texttt{bifrost/valhall/ISA.xml:152}.
\end{tightitemize}
\codebox{\mnem{TEX\_GRADIENT}\imod{\{.message\_slot\_index\}}\imod{\{.flow\}}\imod{\{.skip\}}\imod{\{.dimensionality\}}\imod{\{.register\_width\}}\imod{\{.write\_mask\}}\imod{\{.wide\_indices\}}\imod{\{.coordinate\_or\_derivative\}}\imod{\{.lod\_bias\_disable\}}\imod{\{.lod\_clamp\_disable\}} \dst{@sw}, \src{@sr}, \src{src0}\\\mnem{TEX\_GRADIENT}\imod{\{.message\_slot\_index\}}\imod{\{.flow\}}\imod{\{.skip\}}\imod{\{.dimensionality\}}\imod{\{.projection\_enable\}}\imod{\{.register\_width\}}\imod{\{.write\_mask\}}\imod{\{.wide\_indices\}}\imod{\{.coordinate\_or\_derivative\}}\imod{\{.lod\_bias\_disable\}}\imod{\{.lod\_clamp\_disable\}} \dst{@sw}, \src{@sr}, \src{src0}}
\textbf{Operands:}
\begin{tightitemize}\footnotesize
  \item \texttt{\dst{@sw}} (\texttt{staging\_register\_write\_index} + \texttt{sr\_write\_count}): Destination staging register range receiving sampled texels (or fixed-point 8.8 LOD / derivatives).
  \item \texttt{\src{@sr}} (\texttt{staging\_register\_index} + \texttt{sr\_count}): Source staging register range holding coordinates and explicit \texttt{(dPdx, dPdy)} gradients (or UV for LOD query).
  \item \texttt{\src{src0}}: Texture/sampler descriptor handle.
\end{tightitemize}
\par\textbf{Closest SPIR-V Op:} \texttt{OpImageSampleExplicitLod (Grad)}
\par\textbf{Pseudo-GLSL (Operational Semantics):}
\codebox{if (flag\_disable\_lod\_clamp) @sw = packFixed8\_8(textureQueryLod(handle = src0, uv = @sr));\\else if (flag\_derivative)   @sw = evalScreenSpaceDerivatives(uv = @sr);\\else                        @sw = textureGrad(handle = src0, uv\_grads = @sr);}
\textbf{Execution Context:}
\begin{tightitemize}\footnotesize
  \item \textbf{Texture \& Sampler chain (\hyperref[sec:csf-srt]{\texttt{SRT\_2}} Frag / \texttt{SRT\_0} Comp/Vert):} 64B-aligned VA (\texttt{\& \textasciitilde{}0x3F}) $\to$ \texttt{Resource[table]} $\to$ \hyperref[sec:csf-tex-samp]{\texttt{Texture[tex\_idx]}} (32B, \texttt{type=2}) $\to$ \texttt{GenericPlane[mip]} (32B, \texttt{type=10}: \texttt{pointer}, \texttt{row\_stride}, Morton 16$\times$16 / AFBC) + \texttt{Sampler[samp\_idx]} (32B, \texttt{type=1}: filter, wrap, LOD clamp/bias, compare).
  \item \textbf{Fragment stage preloads \& FP modes (\hyperref[sec:csf-spd]{\texttt{SPD\_2}}, \texttt{r20:r21}):} Hardware pre-fills \texttt{r59} with \texttt{floatBitsToUint(FragCoord.z)} and \texttt{r60} with the MSAA sample coverage mask (\texttt{gl\_SampleMaskIn[0]}), while \texttt{ShaderProgram} word~0 latches 32/64 GPR allocation and warp-wide \texttt{flush\_to\_zero\_mode}, \texttt{suppress\_na\_n}, and \texttt{suppress\_inf}.
\end{tightitemize}
\par\textbf{Worked Examples:}
\par\textbf{Example 1 (Explicit anisotropic gradient sample):}
\codebox{\mnem{TEX\_GRADIENT}\imod{.tex2d.dst32.rgba.slot0} \dst{@r0:r3}, \src{@r4:r9}, \src{u0}}
Samples texture at \texttt{uv = (@r4, @r5)} using explicit \texttt{dPdx = (@r6, @r7)} and \texttt{dPdy = (@r8, @r9)}.
\par\textbf{Example 2 (Query raw mipmap LOD (`.disable\_lod\_bias.disable\_lod\_clamp`)):}
\codebox{\mnem{TEX\_GRADIENT}\imod{.tex2d.disable\_lod\_bias.disable\_lod\_clamp.slot0} \dst{@r0:r1}, \src{@r4:r5}, \src{u0}}
Returns \texttt{(clamped\_lod, unclamped\_lod)} for GLSL \texttt{textureQueryLod(samp, uv)}.
\par\textbf{Mined SPIR-V $\to$ Valhall ISA Example (Querying hardware texture LOD via OpImageQueryLod --- Fragment shader, Mali-G710 (v10)):}
\codebox{\textbf{; SPIR-V (input)}\\\%lod = OpImageQueryLod \%v2float \%sampled\_img \%uv\\[3pt]\textbf{; Valhall ISA (malioc output)}\\TEX\_GRADIENT.slot0.discard.skip.tex2d.dst32.r.wide.disable\_lod\_clamp @r2, @r0:r1, u2\\S16\_TO\_F32 r2, r2\textasciicircum{}.h0\\FMUL.f32 r2, r2\textasciicircum{}, u3.w0}
Instead of sampling texels, \texttt{TEX\_GRADIENT} with \texttt{.disable\_lod\_clamp} queries the texture unit's implicit mipmap LOD calculation and returns the signed 8.8 fixed-point LOD in \texttt{r2.h0}.
\begin{center}\footnotesize
\begin{tabularx}{\linewidth}{@{}lllX@{}}
\toprule
Variant & Arch & Class & Opcode Fields \\
\midrule
\texttt{(none)} & v9--v14 & \texttt{S2RW} & \texttt{opcode1=TEX\_GRADIENT} \\
\texttt{(none)} & v9--v14 & \texttt{S2RW} & \texttt{opcode1=TEX\_GRADIENT} \\
\texttt{(none)} & v9--v14 & \texttt{S2RW} & \texttt{opcode1=TEX\_GRADIENT} \\
\bottomrule
\end{tabularx}
\end{center}
\begin{center}\footnotesize
\begin{tabularx}{0.88\linewidth}{@{}l l X@{}}
\toprule
Field Name & Bits & Type / Fixed Value \\
\midrule
\texttt{src0} & \texttt{0:7} & \texttt{SourceEncoding64} \\
\hyperref[enum:tex_wide_indices_enable_m]{\texttt{wide\_indices}} & \texttt{8} & \texttt{tex\_wide\_indices\_enable\_m (def none)} \\
\textit{Reserved} & \texttt{9:12} & \texttt{zero} \\
\hyperref[enum:tex_lod_bias_disable_m]{\texttt{lod\_bias\_disable}} & \texttt{13} & \texttt{tex\_lod\_bias\_disable\_m} \\
\hyperref[enum:tex_lod_clamp_disable_m]{\texttt{lod\_clamp\_disable}} & \texttt{14} & \texttt{tex\_lod\_clamp\_disable\_m} \\
\hyperref[enum:tex_derivative_enable_m]{\texttt{derivative\_enable}} & \texttt{15} & \texttt{tex\_derivative\_enable\_m} \\
\texttt{staging\_register\_write\_index} & \texttt{16:21} & \texttt{uint} \\
\hyperref[enum:tex_gradient_write_mask_m]{\texttt{write\_mask}} & \texttt{22:25} & \texttt{tex\_gradient\_write\_mask\_m (def rgba)} \\
\textit{Reserved} & \texttt{26:27} & \texttt{ignore} \\
\hyperref[enum:tex_dimensionality_m]{\texttt{dimensionality}} & \texttt{28:29} & \texttt{tex\_dimensionality\_m} \\
\hyperref[enum:message_slot_index_m]{\texttt{message\_slot\_index}} & \texttt{30:32} & \texttt{message\_slot\_index\_m} \\
\texttt{sr\_count} & \texttt{33:35} & \texttt{sr\_count\_tex\_t} \\
\texttt{sr\_write\_count} & \texttt{36:38} & \texttt{sr\_write\_count\_t} \\
\hyperref[enum:tex_gradient_skip_m]{\texttt{skip}} & \texttt{39} & \texttt{tex\_gradient\_skip\_m (def none)} \\
\texttt{staging\_register\_index} & \texttt{40:45} & \texttt{uint} \\
\hyperref[enum:tex_gradient_width_m]{\texttt{register\_width}} & \texttt{46} & \texttt{tex\_gradient\_width\_m} \\
\texttt{secondary\_width} & \texttt{47} & \texttt{exact=0} \\
\texttt{opcode1} & \texttt{48:51} & \texttt{exact=TEX\_GRADIENT} \\
\texttt{instruction\_class} & \texttt{52:56} & \texttt{exact=S2RW} \\
\texttt{fau\_page\_index} & \texttt{57:58} & \texttt{uint} \\
\hyperref[enum:flow_control_m]{\texttt{flow}} & \texttt{59:62} & \texttt{flow\_control\_m} \\
\textit{Reserved} & \texttt{63} & \texttt{zero} \\
\bottomrule
\end{tabularx}
\end{center}
\textbf{Fields Explained:}
\begin{tightitemize}\footnotesize
  \item \texttt{src0}: Primary source operand (\texttt{r0..r63}, \texttt{u0..u63} uniform, \texttt{k*} FAU constant, or \texttt{t0..t2} temporary; \texttt{\textasciicircum{}} marks discard/last-use). \emph{Example:} {\scriptsize\ttfamily \mnem{FADD.f32} \dst{r0}, \src{r1}, \src{r2}} (\textbf{SPIR-V:} \texttt{OpFAdd (Operand 0)}; \textbf{GLSL:} \texttt{float c = a + b;})
  \item \hyperref[enum:tex_wide_indices_enable_m]{\texttt{wide\_indices}}: Uses 32-bit wide texture/sampler table indices instead of compact 16-bit packed indices. \emph{Example:} {\scriptsize\ttfamily \mnem{TEX\_SINGLE}\imod{.tex2d.wide} \dst{@r0:r3}, \src{@r4:r5}, \src{r6}} (\textbf{SPIR-V:} \texttt{OpImageSampleImplicitLod (bindless)}; \textbf{GLSL:} \texttt{vec4 c = texture(textures[nonuniformEXT(i)], uv);})
  \item \hyperref[enum:tex_lod_bias_disable_m]{\texttt{lod\_bias\_disable}}: Disables sampler/shader LOD bias when computing explicit gradients (\texttt{TEX\_GRADIENT}). \emph{Example:} {\scriptsize\ttfamily \mnem{TEX\_GRADIENT}\imod{.tex2d.disable\_lod\_bias} \dst{@r0:r1}, \src{@r4:r5}} (\textbf{SPIR-V:} \texttt{OpImageQueryLod}; \textbf{GLSL:} \texttt{vec2 lod = textureQueryLod(tex, uv);})
  \item \hyperref[enum:tex_lod_clamp_disable_m]{\texttt{lod\_clamp\_disable}}: Disables sampler min/max LOD clamp bounds when querying raw mip level. \emph{Example:} {\scriptsize\ttfamily \mnem{TEX\_GRADIENT}\imod{.tex2d.disable\_lod\_clamp} \dst{@r0:r1}, \src{@r4:r5}} (\textbf{SPIR-V:} \texttt{OpImageQueryLod}; \textbf{GLSL:} \texttt{vec2 lod = textureQueryLod(tex, uv);})
  \item \hyperref[enum:tex_derivative_enable_m]{\texttt{derivative\_enable}}: Returns raw screen-space derivatives ($dP/dx, dP/dy$) instead of scalar LOD. \emph{Example:} {\scriptsize\ttfamily \mnem{TEX\_GRADIENT}\imod{.tex2d.derivative} \dst{@r0:r3}, \src{@r4:r5}} (\textbf{SPIR-V:} \texttt{OpDPdx / OpDPdy}; \textbf{GLSL:} \texttt{vec2 dx = dFdx(uv);})
  \item \texttt{staging\_register\_write\_index}: Separate destination staging register window base when read and write staging windows differ. \emph{Example:} {\scriptsize\ttfamily \mnem{ATOM\_RETURN.i32}\imod{.aadd} \dst{@r0}, \src{@r1}, \src{r2}, \imm{\#0}} (\textbf{SPIR-V:} \texttt{OpAtomicIAdd}; \textbf{GLSL:} \texttt{uint old = atomicAdd(ssbo.cnt, 1u);})
  \item \hyperref[enum:tex_gradient_write_mask_m]{\texttt{write\_mask}}: Channel enable bitmask (\texttt{r}..\texttt{rgba}) controlling which texture components are written to staging GPRs. \emph{Example:} {\scriptsize\ttfamily \mnem{TEX\_SINGLE}\imod{.tex2d.f32.rg} \dst{@r0:r1}, \src{@r4:r5}} (\textbf{SPIR-V:} \texttt{OpImageSampleImplicitLod + OpVectorShuffle}; \textbf{GLSL:} \texttt{vec2 rg = texture(tex, uv).rg;})
  \item \hyperref[enum:tex_dimensionality_m]{\texttt{dimensionality}}: Texture dimensionality (\texttt{tex1d}, \texttt{tex2d}, \texttt{tex3d}, \texttt{cube}). \emph{Example:} {\scriptsize\ttfamily \mnem{TEX\_SINGLE}\imod{.tex3d.f32.rgba} \dst{@r0:r3}, \src{@r4:r6}} (\textbf{SPIR-V:} \texttt{OpImageSampleImplicitLod (Dim=3D)}; \textbf{GLSL:} \texttt{vec4 c = texture(volTex, uvw);})
  \item \hyperref[enum:message_slot_index_m]{\texttt{message\_slot\_index}}: Asynchronous scoreboard slot (\texttt{slot0}--\texttt{slot2}, \texttt{slot7}) tracking completion of this message. \emph{Example:} {\scriptsize\ttfamily \mnem{LOAD.i32}\imod{.slot0} \dst{@r0}, \src{r2}, \imm{\#0}} (\textbf{SPIR-V:} \texttt{OpLoad}; \textbf{GLSL:} \texttt{uint v = ssbo.arr[i];})
  \item \texttt{sr\_count}: Number of consecutive 32-bit staging registers transferred by a load/store/texture message. \emph{Example:} {\scriptsize\ttfamily \mnem{LD\_MULTI.i128} \dst{@r0:r3}, \src{r4}, \imm{\#0}} (\textbf{SPIR-V:} \texttt{OpLoad (wide vector)}; \textbf{GLSL:} \texttt{uvec4 x = buf[i];})
  \item \texttt{sr\_write\_count}: Number of consecutive 32-bit staging registers written back by a dual-staging message (e.g. ATOM\textbackslash\{\}\_RETURN / TEX). \emph{Example:} {\scriptsize\ttfamily \mnem{TEX\_SINGLE}\imod{.tex2d.f32.rgba} \dst{@r0:r3}, \src{@r4:r5}} (\textbf{SPIR-V:} \texttt{OpImageSampleImplicitLod}; \textbf{GLSL:} \texttt{vec4 c = texture(tex, uv);})
  \item \hyperref[enum:tex_gradient_skip_m]{\texttt{skip}}: Skips texture/memory fetch on helper/discarded fragment lanes while preserving 2x2 quad derivatives. \emph{Example:} {\scriptsize\ttfamily \mnem{TEX\_SINGLE}\imod{.tex2d.computed.skip} \dst{@r0:r3}, \src{@r4:r5}} (\textbf{SPIR-V:} \texttt{OpImageSampleImplicitLod (after OpDemoteToHelperInvocation)}; \textbf{GLSL:} \texttt{if (a \textless{} 0.0) discard; vec4 c = texture(s, uv);})
  \item \texttt{staging\_register\_index}: Base GPR index of the contiguous staging register window (\texttt{@rN:rM}) for message payload/return. \emph{Example:} {\scriptsize\ttfamily \mnem{STORE.i64}\imod{.slot7} \src{@r0:r1}, \src{r2}, \imm{\#8}} (\textbf{SPIR-V:} \texttt{OpStore (uvec2)}; \textbf{GLSL:} \texttt{ssbo.pairs[i] = uvec2(a, b);})
  \item \hyperref[enum:tex_gradient_width_m]{\texttt{register\_width}}: Selects 16-bit vs 32-bit GPR element width for the corresponding message return window. \emph{Example:} {\scriptsize\ttfamily \mnem{TEX\_DUAL}\imod{.primary\_f16.secondary\_f32} \dst{@r0:r2}, \src{@r4:r5}} (\textbf{SPIR-V:} \texttt{OpImageSampleImplicitLod (RelaxedPrecision)}; \textbf{GLSL:} \texttt{mediump vec2 a = texture(t0, uv).rg;})
  \item \texttt{secondary\_width}: Selects 16-bit vs 32-bit GPR element width for the corresponding message return window. \emph{Example:} {\scriptsize\ttfamily \mnem{TEX\_DUAL}\imod{.primary\_f16.secondary\_f32} \dst{@r0:r2}, \src{@r4:r5}} (\textbf{SPIR-V:} \texttt{OpImageSampleImplicitLod (RelaxedPrecision)}; \textbf{GLSL:} \texttt{mediump vec2 a = texture(t0, uv).rg;})
  \item \texttt{opcode1}: Primary sub-opcode within the instruction class uniquely identifying the operation family. \emph{Example:} {\scriptsize\ttfamily \mnem{IADD.s32} \dst{r0}, \src{r1}, \src{r2}} (\textbf{SPIR-V:} \texttt{OpIAdd (operation selector)}; \textbf{GLSL:} \texttt{// Primary opcode bits})
  \item \texttt{instruction\_class}: 4-bit major instruction class selector (\texttt{bits 48:51}) partitioning the 64-bit Valhall opcode space (e.g., \texttt{FMA}, \texttt{CVT}, \texttt{SFU}, \texttt{LS}, \texttt{TEX}, \texttt{BRANCH}). \emph{Example:} {\scriptsize\ttfamily \mnem{FADD.f32} \dst{r0}, \src{r1}, \src{r2}} (\textbf{SPIR-V:} \texttt{N/A (encoding class)}; \textbf{GLSL:} \texttt{// Hardware execution pipe / decode class})
  \item \texttt{fau\_page\_index}: Selects the FAU (Fast Access Uniform) page (\texttt{0..3}) when a source operand references hardware small constants, thread IDs, or lane constants. \emph{Example:} {\scriptsize\ttfamily \mnem{FADD.f32} \dst{r0}, \src{r1}, \src{k1f}} (\textbf{SPIR-V:} \texttt{OpConstant / BuiltIn}; \textbf{GLSL:} \texttt{float y = x + 1.0;})
  \item \hyperref[enum:flow_control_m]{\texttt{flow}}: 4-bit embedded flow-control \& scoreboard wait modifier (\texttt{bits 59:62}: \texttt{none}, \texttt{wait0..wait0126}, \texttt{wait}, \texttt{reconverge}, \texttt{discard}, \texttt{end}). \emph{Example:} {\scriptsize\ttfamily \mnem{FADD.f32}\imod{.wait0.reconverge} \dst{r0}, \src{r1}, \src{r2}} (\textbf{SPIR-V:} \texttt{OpControlBarrier / OpReturn}; \textbf{GLSL:} \texttt{// Scoreboard wait + warp reconvergence})
\end{tightitemize}
\clearpage

\subsection*{TEX\_SINGLE}
\addcontentsline{toc}{subsection}{TEX\_SINGLE}
Standard filtered texture sample via sampler/texture handle \texttt{src0} and staging coordinates \texttt{sr\_src}, writing masked channels to \texttt{sr\_dst}.
\par\textbf{Notes:}
\begin{tightitemize}
  \item Hardware texture unit message.
  \item \texttt{skip} disables helper-invocation sampling when derivatives not needed.
  \item \textbf{Source references:} \texttt{kraid/ops.rs:4551-4595}, \texttt{encode\_v9.rs:3442-3459}, \texttt{bifrost/valhall/ISA.xml:154}.
\end{tightitemize}
\codebox{\mnem{TEX\_SINGLE}\imod{\{.message\_slot\_index\}}\imod{\{.flow\}}\imod{\{.skip\}}\imod{\{.dimensionality\}}\imod{\{.register\_width\}}\imod{\{.write\_mask\}}\imod{\{.wide\_indices\}}\imod{\{.array\_enable\}}\imod{\{.texel\_offset\}}\imod{\{.compare\_enable\}}\imod{\{.lod\_mode\}} \dst{@sw}, \src{@sr}, \src{src0}\\\mnem{TEX\_SINGLE}\imod{\{.message\_slot\_index\}}\imod{\{.flow\}}\imod{\{.skip\}}\imod{\{.dimensionality\}}\imod{\{.projection\_enable\}}\imod{\{.register\_width\}}\imod{\{.write\_mask\}}\imod{\{.wide\_indices\}}\imod{\{.array\_enable\}}\imod{\{.texel\_offset\}}\imod{\{.compare\_enable\}}\imod{\{.lod\_mode\}} \dst{@sw}, \src{@sr}, \src{src0}}
\textbf{Operands:}
\begin{tightitemize}\footnotesize
  \item \texttt{\dst{@sw}} (\texttt{staging\_register\_write\_index} + \texttt{sr\_write\_count}): Destination staging register range receiving sampled RGBA texels (per \texttt{write\_mask}).
  \item \texttt{\src{@sr}} (\texttt{staging\_register\_index} + \texttt{sr\_count}): Source staging register range holding packed coordinates \texttt{(u, v, [w/layer], [lod/bias], [dref], [offset])}.
  \item \texttt{\src{src0}}: Texture/sampler descriptor handle (32-bit FAU uniform or 64-bit register pair).
\end{tightitemize}
\par\textbf{Closest SPIR-V Op:} \texttt{OpImageSampleImplicitLod / ExplicitLod}
\par\textbf{Pseudo-GLSL (Operational Semantics):}
\codebox{@sw = textureSample(handle = src0, coords = @sr, dim, lod\_mode, flag\_project, flag\_compare);}
\textbf{Execution Context:}
\begin{tightitemize}\footnotesize
  \item \textbf{Texture \& Sampler chain (\hyperref[sec:csf-srt]{\texttt{SRT\_2}} Frag / \texttt{SRT\_0} Comp/Vert):} 64B-aligned VA (\texttt{\& \textasciitilde{}0x3F}) $\to$ \texttt{Resource[table]} $\to$ \hyperref[sec:csf-tex-samp]{\texttt{Texture[tex\_idx]}} (32B, \texttt{type=2}) $\to$ \texttt{GenericPlane[mip]} (32B, \texttt{type=10}: \texttt{pointer}, \texttt{row\_stride}, Morton 16$\times$16 / AFBC) + \texttt{Sampler[samp\_idx]} (32B, \texttt{type=1}: filter, wrap, LOD clamp/bias, compare).
  \item \textbf{Fragment stage preloads \& FP modes (\hyperref[sec:csf-spd]{\texttt{SPD\_2}}, \texttt{r20:r21}):} Hardware pre-fills \texttt{r59} with \texttt{floatBitsToUint(FragCoord.z)} and \texttt{r60} with the MSAA sample coverage mask (\texttt{gl\_SampleMaskIn[0]}), while \texttt{ShaderProgram} word~0 latches 32/64 GPR allocation and warp-wide \texttt{flush\_to\_zero\_mode}, \texttt{suppress\_na\_n}, and \texttt{suppress\_inf}.
\end{tightitemize}
\par\textbf{Worked Examples:}
\par\textbf{Example 1 (Implicit quad-LOD 2D sample with `.skip`):}
\codebox{\mnem{TEX\_SINGLE}\imod{.tex2d.computed.dst32.rgba.skip.slot0} \dst{@r0:r3}, \src{@r4:r5}, \src{u0}}
Computes 2x2 quad LOD from \texttt{@r4:r5}, skips memory fetch on helper lanes (\texttt{.skip}), and returns RGBA in \texttt{@r0:r3}.
\par\textbf{Example 2 (Hardware-projected cubemap sample (v11+ `.cube.project`)):}
\codebox{\mnem{TEX\_SINGLE}\imod{.cube.project.computed.dst32.rgba.slot0} \dst{@r0:r3}, \src{@r4:r6}, \src{u0}}
Takes 3D direction \texttt{(x, y, z)} in \texttt{@r4:r6} directly on v11+, eliminating v9--v10 \texttt{CUBEFACE1/2} + \texttt{CUBE\_SSEL/TSEL}.
\par\textbf{Mined SPIR-V $\to$ Valhall ISA Example (Sampling a 2D texture array with explicit LOD 0 in a Vertex shader --- Vertex shader, Mali-G710 (v10)):}
\codebox{\textbf{; SPIR-V (input)}\\\%color = OpImageSampleExplicitLod \%v4float \%tex2d\_arr \%uv\_layer Lod \%float\_0\\[3pt]\textbf{; Valhall ISA (malioc output)}\\MOV\_IMM.i32 r0, \#0x3f000000       ; u = 0.5\\MOV.i32 r1, r0                    ; v = 0.5\\MOV\_IMM.i32 r2, \#0x40000000       ; layer = 2.0\\TEX\_SINGLE.slot0.wait0.tex2d.dst32.rgba.wide.array @r0:r1:r2:r3, @r0:r1:r2, u3}
\texttt{OpImageSampleExplicitLod} on a \texttt{2DArray} sampler lowers to \texttt{TEX\_SINGLE} with \texttt{.tex2d.array} (default LOD mode = zero/explicit) reading \texttt{(u, v, layer)} from \texttt{@r0:r1:r2} and writing 4 FP32 channels to \texttt{@r0:r1:r2:r3}.
\begin{center}\footnotesize
\begin{tabularx}{\linewidth}{@{}lllX@{}}
\toprule
Variant & Arch & Class & Opcode Fields \\
\midrule
\texttt{(none)} & v9--v14 & \texttt{S2RW} & \texttt{opcode1=TEX\_SINGLE} \\
\texttt{(none)} & v9--v14 & \texttt{S2RW} & \texttt{opcode1=TEX\_SINGLE} \\
\texttt{(none)} & v9--v14 & \texttt{S2RW} & \texttt{opcode1=TEX\_SINGLE} \\
\bottomrule
\end{tabularx}
\end{center}
\begin{center}\footnotesize
\begin{tabularx}{0.88\linewidth}{@{}l l X@{}}
\toprule
Field Name & Bits & Type / Fixed Value \\
\midrule
\texttt{src0} & \texttt{0:7} & \texttt{SourceEncoding64} \\
\hyperref[enum:tex_wide_indices_enable_m]{\texttt{wide\_indices}} & \texttt{8} & \texttt{tex\_wide\_indices\_enable\_m (def none)} \\
\textit{Reserved} & \texttt{9} & \texttt{zero} \\
\hyperref[enum:tex_array_enable_m]{\texttt{array\_enable}} & \texttt{10} & \texttt{tex\_array\_enable\_m (def none)} \\
\hyperref[enum:tex_texel_offset_m]{\texttt{texel\_offset}} & \texttt{11} & \texttt{tex\_texel\_offset\_m (def none)} \\
\hyperref[enum:tex_compare_enable_m]{\texttt{compare\_enable}} & \texttt{12} & \texttt{tex\_compare\_enable\_m (def none)} \\
\hyperref[enum:tex_lod_mode_m]{\texttt{lod\_mode}} & \texttt{13:15} & \texttt{tex\_lod\_mode\_m} \\
\texttt{staging\_register\_write\_index} & \texttt{16:21} & \texttt{uint} \\
\hyperref[enum:tex_write_mask_m]{\texttt{write\_mask}} & \texttt{22:25} & \texttt{tex\_write\_mask\_m (def rgba)} \\
\textit{Reserved} & \texttt{26:27} & \texttt{ignore} \\
\hyperref[enum:tex_dimensionality_m]{\texttt{dimensionality}} & \texttt{28:29} & \texttt{tex\_dimensionality\_m} \\
\hyperref[enum:message_slot_index_m]{\texttt{message\_slot\_index}} & \texttt{30:32} & \texttt{message\_slot\_index\_m} \\
\texttt{sr\_count} & \texttt{33:35} & \texttt{sr\_count\_tex\_t} \\
\texttt{sr\_write\_count} & \texttt{36:38} & \texttt{sr\_write\_count\_t} \\
\hyperref[enum:tex_skip_m]{\texttt{skip}} & \texttt{39} & \texttt{tex\_skip\_m (def none)} \\
\texttt{staging\_register\_index} & \texttt{40:45} & \texttt{uint} \\
\hyperref[enum:tex_width_m]{\texttt{register\_width}} & \texttt{46} & \texttt{tex\_width\_m} \\
\texttt{secondary\_width} & \texttt{47} & \texttt{exact=0} \\
\texttt{opcode1} & \texttt{48:51} & \texttt{exact=TEX\_SINGLE} \\
\texttt{instruction\_class} & \texttt{52:56} & \texttt{exact=S2RW} \\
\texttt{fau\_page\_index} & \texttt{57:58} & \texttt{uint} \\
\hyperref[enum:flow_control_m]{\texttt{flow}} & \texttt{59:62} & \texttt{flow\_control\_m} \\
\textit{Reserved} & \texttt{63} & \texttt{zero} \\
\bottomrule
\end{tabularx}
\end{center}
\textbf{Fields Explained:}
\begin{tightitemize}\footnotesize
  \item \texttt{src0}: Primary source operand (\texttt{r0..r63}, \texttt{u0..u63} uniform, \texttt{k*} FAU constant, or \texttt{t0..t2} temporary; \texttt{\textasciicircum{}} marks discard/last-use). \emph{Example:} {\scriptsize\ttfamily \mnem{FADD.f32} \dst{r0}, \src{r1}, \src{r2}} (\textbf{SPIR-V:} \texttt{OpFAdd (Operand 0)}; \textbf{GLSL:} \texttt{float c = a + b;})
  \item \hyperref[enum:tex_wide_indices_enable_m]{\texttt{wide\_indices}}: Uses 32-bit wide texture/sampler table indices instead of compact 16-bit packed indices. \emph{Example:} {\scriptsize\ttfamily \mnem{TEX\_SINGLE}\imod{.tex2d.wide} \dst{@r0:r3}, \src{@r4:r5}, \src{r6}} (\textbf{SPIR-V:} \texttt{OpImageSampleImplicitLod (bindless)}; \textbf{GLSL:} \texttt{vec4 c = texture(textures[nonuniformEXT(i)], uv);})
  \item \hyperref[enum:tex_array_enable_m]{\texttt{array\_enable}}: Enables array-layer coordinate in texture staging registers (\texttt{sampler2DArray} / \texttt{samplerCubeArray}). \emph{Example:} {\scriptsize\ttfamily \mnem{TEX\_SINGLE}\imod{.tex2d.array} \dst{@r0:r3}, \src{@r4:r6}} (\textbf{SPIR-V:} \texttt{OpImageSampleImplicitLod (Arrayed)}; \textbf{GLSL:} \texttt{vec4 c = texture(texArr, vec3(uv, layer));})
  \item \hyperref[enum:tex_texel_offset_m]{\texttt{texel\_offset}}: Applies a signed 4-bit per-axis texel offset $(dx, dy, dz)$ packed in staging registers. \emph{Example:} {\scriptsize\ttfamily \mnem{TEX\_SINGLE}\imod{.tex2d.texel\_offset} \dst{@r0:r3}, \src{@r4:r6}} (\textbf{SPIR-V:} \texttt{OpImageSampleImplicitLod (ConstOffset)}; \textbf{GLSL:} \texttt{vec4 c = textureOffset(tex, uv, ivec2(1, -1));})
  \item \hyperref[enum:tex_compare_enable_m]{\texttt{compare\_enable}}: Enables depth-reference shadow comparison (\texttt{Dref} in staging registers). \emph{Example:} {\scriptsize\ttfamily \mnem{TEX\_SINGLE}\imod{.tex2d.compare} \dst{@r0}, \src{@r4:r6}} (\textbf{SPIR-V:} \texttt{OpImageSampleDrefImplicitLod}; \textbf{GLSL:} \texttt{float s = texture(shadowMap, vec3(uv, refZ));})
  \item \hyperref[enum:tex_lod_mode_m]{\texttt{lod\_mode}}: Selects mipmap LOD calculation mode (\texttt{none}, \texttt{computed}, \texttt{explicit}, \texttt{computed\_bias}, \texttt{grdesc}). \emph{Example:} {\scriptsize\ttfamily \mnem{TEX\_SINGLE}\imod{.tex2d.explicit} \dst{@r0:r3}, \src{@r4:r6}} (\textbf{SPIR-V:} \texttt{OpImageSampleExplicitLod}; \textbf{GLSL:} \texttt{vec4 c = textureLod(tex, uv, 2.0);})
  \item \texttt{staging\_register\_write\_index}: Separate destination staging register window base when read and write staging windows differ. \emph{Example:} {\scriptsize\ttfamily \mnem{ATOM\_RETURN.i32}\imod{.aadd} \dst{@r0}, \src{@r1}, \src{r2}, \imm{\#0}} (\textbf{SPIR-V:} \texttt{OpAtomicIAdd}; \textbf{GLSL:} \texttt{uint old = atomicAdd(ssbo.cnt, 1u);})
  \item \hyperref[enum:tex_write_mask_m]{\texttt{write\_mask}}: Channel enable bitmask (\texttt{r}..\texttt{rgba}) controlling which texture components are written to staging GPRs. \emph{Example:} {\scriptsize\ttfamily \mnem{TEX\_SINGLE}\imod{.tex2d.f32.rg} \dst{@r0:r1}, \src{@r4:r5}} (\textbf{SPIR-V:} \texttt{OpImageSampleImplicitLod + OpVectorShuffle}; \textbf{GLSL:} \texttt{vec2 rg = texture(tex, uv).rg;})
  \item \hyperref[enum:tex_dimensionality_m]{\texttt{dimensionality}}: Texture dimensionality (\texttt{tex1d}, \texttt{tex2d}, \texttt{tex3d}, \texttt{cube}). \emph{Example:} {\scriptsize\ttfamily \mnem{TEX\_SINGLE}\imod{.tex3d.f32.rgba} \dst{@r0:r3}, \src{@r4:r6}} (\textbf{SPIR-V:} \texttt{OpImageSampleImplicitLod (Dim=3D)}; \textbf{GLSL:} \texttt{vec4 c = texture(volTex, uvw);})
  \item \hyperref[enum:message_slot_index_m]{\texttt{message\_slot\_index}}: Asynchronous scoreboard slot (\texttt{slot0}--\texttt{slot2}, \texttt{slot7}) tracking completion of this message. \emph{Example:} {\scriptsize\ttfamily \mnem{LOAD.i32}\imod{.slot0} \dst{@r0}, \src{r2}, \imm{\#0}} (\textbf{SPIR-V:} \texttt{OpLoad}; \textbf{GLSL:} \texttt{uint v = ssbo.arr[i];})
  \item \texttt{sr\_count}: Number of consecutive 32-bit staging registers transferred by a load/store/texture message. \emph{Example:} {\scriptsize\ttfamily \mnem{LD\_MULTI.i128} \dst{@r0:r3}, \src{r4}, \imm{\#0}} (\textbf{SPIR-V:} \texttt{OpLoad (wide vector)}; \textbf{GLSL:} \texttt{uvec4 x = buf[i];})
  \item \texttt{sr\_write\_count}: Number of consecutive 32-bit staging registers written back by a dual-staging message (e.g. ATOM\textbackslash\{\}\_RETURN / TEX). \emph{Example:} {\scriptsize\ttfamily \mnem{TEX\_SINGLE}\imod{.tex2d.f32.rgba} \dst{@r0:r3}, \src{@r4:r5}} (\textbf{SPIR-V:} \texttt{OpImageSampleImplicitLod}; \textbf{GLSL:} \texttt{vec4 c = texture(tex, uv);})
  \item \hyperref[enum:tex_skip_m]{\texttt{skip}}: Skips texture/memory fetch on helper/discarded fragment lanes while preserving 2x2 quad derivatives. \emph{Example:} {\scriptsize\ttfamily \mnem{TEX\_SINGLE}\imod{.tex2d.computed.skip} \dst{@r0:r3}, \src{@r4:r5}} (\textbf{SPIR-V:} \texttt{OpImageSampleImplicitLod (after OpDemoteToHelperInvocation)}; \textbf{GLSL:} \texttt{if (a \textless{} 0.0) discard; vec4 c = texture(s, uv);})
  \item \texttt{staging\_register\_index}: Base GPR index of the contiguous staging register window (\texttt{@rN:rM}) for message payload/return. \emph{Example:} {\scriptsize\ttfamily \mnem{STORE.i64}\imod{.slot7} \src{@r0:r1}, \src{r2}, \imm{\#8}} (\textbf{SPIR-V:} \texttt{OpStore (uvec2)}; \textbf{GLSL:} \texttt{ssbo.pairs[i] = uvec2(a, b);})
  \item \hyperref[enum:tex_width_m]{\texttt{register\_width}}: Selects 16-bit vs 32-bit GPR element width for the corresponding message return window. \emph{Example:} {\scriptsize\ttfamily \mnem{TEX\_DUAL}\imod{.primary\_f16.secondary\_f32} \dst{@r0:r2}, \src{@r4:r5}} (\textbf{SPIR-V:} \texttt{OpImageSampleImplicitLod (RelaxedPrecision)}; \textbf{GLSL:} \texttt{mediump vec2 a = texture(t0, uv).rg;})
  \item \texttt{secondary\_width}: Selects 16-bit vs 32-bit GPR element width for the corresponding message return window. \emph{Example:} {\scriptsize\ttfamily \mnem{TEX\_DUAL}\imod{.primary\_f16.secondary\_f32} \dst{@r0:r2}, \src{@r4:r5}} (\textbf{SPIR-V:} \texttt{OpImageSampleImplicitLod (RelaxedPrecision)}; \textbf{GLSL:} \texttt{mediump vec2 a = texture(t0, uv).rg;})
  \item \texttt{opcode1}: Primary sub-opcode within the instruction class uniquely identifying the operation family. \emph{Example:} {\scriptsize\ttfamily \mnem{IADD.s32} \dst{r0}, \src{r1}, \src{r2}} (\textbf{SPIR-V:} \texttt{OpIAdd (operation selector)}; \textbf{GLSL:} \texttt{// Primary opcode bits})
  \item \texttt{instruction\_class}: 4-bit major instruction class selector (\texttt{bits 48:51}) partitioning the 64-bit Valhall opcode space (e.g., \texttt{FMA}, \texttt{CVT}, \texttt{SFU}, \texttt{LS}, \texttt{TEX}, \texttt{BRANCH}). \emph{Example:} {\scriptsize\ttfamily \mnem{FADD.f32} \dst{r0}, \src{r1}, \src{r2}} (\textbf{SPIR-V:} \texttt{N/A (encoding class)}; \textbf{GLSL:} \texttt{// Hardware execution pipe / decode class})
  \item \texttt{fau\_page\_index}: Selects the FAU (Fast Access Uniform) page (\texttt{0..3}) when a source operand references hardware small constants, thread IDs, or lane constants. \emph{Example:} {\scriptsize\ttfamily \mnem{FADD.f32} \dst{r0}, \src{r1}, \src{k1f}} (\textbf{SPIR-V:} \texttt{OpConstant / BuiltIn}; \textbf{GLSL:} \texttt{float y = x + 1.0;})
  \item \hyperref[enum:flow_control_m]{\texttt{flow}}: 4-bit embedded flow-control \& scoreboard wait modifier (\texttt{bits 59:62}: \texttt{none}, \texttt{wait0..wait0126}, \texttt{wait}, \texttt{reconverge}, \texttt{discard}, \texttt{end}). \emph{Example:} {\scriptsize\ttfamily \mnem{FADD.f32}\imod{.wait0.reconverge} \dst{r0}, \src{r1}, \src{r2}} (\textbf{SPIR-V:} \texttt{OpControlBarrier / OpReturn}; \textbf{GLSL:} \texttt{// Scoreboard wait + warp reconvergence})
\end{tightitemize}
\clearpage

\subsection*{TEX\_WEIGHTS}
\addcontentsline{toc}{subsection}{TEX\_WEIGHTS}
Texture footprint/filter weights query (arch v9 only): returns bilinear/footprint weights for \texttt{footprint} (\texttt{2x2}, etc.) at coordinates in \texttt{sr\_src}.
\par\textbf{Notes:}
\begin{tightitemize}
  \item Removed after v9.
  \item \textbf{Source references:} \texttt{isa-v9-v14.xml}.
\end{tightitemize}
\codebox{\mnem{TEX\_WEIGHTS}\imod{\{.message\_slot\_index\}}\imod{\{.flow\}}\imod{\{.skip\}}\imod{\{.dimensionality\}}\imod{\{.register\_width\}}\imod{\{.write\_mask\}}\imod{\{.wide\_indices\}}\imod{\{.array\_enable\}}\imod{\{.texel\_offset\}}\imod{\{.coordinate\_mode\}}\imod{\{.footprint\}} \dst{@sw}, \src{@sr}, \src{src0}}
\textbf{Operands:}
\begin{tightitemize}\footnotesize
  \item \texttt{\dst{@sw}} (\texttt{staging\_register\_write\_index} + \texttt{sr\_write\_count}): Destination staging register range receiving the 4 bilinear filter weights (v9 only).
  \item \texttt{\src{@sr}} (\texttt{staging\_register\_index} + \texttt{sr\_count}): Source staging register range holding \texttt{(u, v)} texture coordinates.
  \item \texttt{\src{src0}}: Texture/sampler descriptor handle.
\end{tightitemize}
\par\textbf{Closest SPIR-V Op:} \texttt{None (v9 footprint weights query)}
\par\textbf{Pseudo-GLSL (Operational Semantics):}
\codebox{@sw = bilinearFootprintWeights2x2(handle = src0, uv = @sr); // v9 only}
\textbf{Execution Context:}
\begin{tightitemize}\footnotesize
  \item \textbf{Texture \& Sampler chain (\hyperref[sec:csf-srt]{\texttt{SRT\_2}} Frag / \texttt{SRT\_0} Comp/Vert):} 64B-aligned VA (\texttt{\& \textasciitilde{}0x3F}) $\to$ \texttt{Resource[table]} $\to$ \hyperref[sec:csf-tex-samp]{\texttt{Texture[tex\_idx]}} (32B, \texttt{type=2}) $\to$ \texttt{GenericPlane[mip]} (32B, \texttt{type=10}: \texttt{pointer}, \texttt{row\_stride}, Morton 16$\times$16 / AFBC) + \texttt{Sampler[samp\_idx]} (32B, \texttt{type=1}: filter, wrap, LOD clamp/bias, compare).
  \item \textbf{Fragment stage preloads \& FP modes (\hyperref[sec:csf-spd]{\texttt{SPD\_2}}, \texttt{r20:r21}):} Hardware pre-fills \texttt{r59} with \texttt{floatBitsToUint(FragCoord.z)} and \texttt{r60} with the MSAA sample coverage mask (\texttt{gl\_SampleMaskIn[0]}), while \texttt{ShaderProgram} word~0 latches 32/64 GPR allocation and warp-wide \texttt{flush\_to\_zero\_mode}, \texttt{suppress\_na\_n}, and \texttt{suppress\_inf}.
\end{tightitemize}
\par\textbf{Worked Examples:}
\par\textbf{Example 1 (2x2 bilinear weight query (v9 only)):}
\codebox{\mnem{TEX\_WEIGHTS}\imod{.tex2d.footprint\_2x2.slot0} \dst{@r0:r3}, \src{@r4:r5}, \src{u0}}
Returns the 4 normalized bilinear filter weights for the 2x2 texel footprint at \texttt{@r4:r5} (removed in v10+).
\par\textbf{Example 2 (Sub-texel center check):}
\codebox{\mnem{TEX\_WEIGHTS}\imod{.tex2d.footprint\_2x2.slot0} \dst{@r0:r3}, \src{@r4:r5}, \src{u0}}
At exact texel midpoint \texttt{(frac=0.5, 0.5)}: returns \texttt{[0.25f, 0.25f, 0.25f, 0.25f]} in \texttt{@r0:r3}.

\begin{center}\footnotesize
\begin{tabularx}{\linewidth}{@{}lllX@{}}
\toprule
Variant & Arch & Class & Opcode Fields \\
\midrule
\texttt{(none)} & v9--v14 & \texttt{S2RW} & \texttt{opcode1=TEX\_WEIGHTS} \\
\bottomrule
\end{tabularx}
\end{center}
\begin{center}\footnotesize
\begin{tabularx}{0.88\linewidth}{@{}l l X@{}}
\toprule
Field Name & Bits & Type / Fixed Value \\
\midrule
\texttt{src0} & \texttt{0:7} & \texttt{SourceEncoding64} \\
\hyperref[enum:tex_wide_indices_enable_m]{\texttt{wide\_indices}} & \texttt{8} & \texttt{tex\_wide\_indices\_enable\_m (def none)} \\
\textit{Reserved} & \texttt{9} & \texttt{zero} \\
\hyperref[enum:tex_array_enable_m]{\texttt{array\_enable}} & \texttt{10} & \texttt{tex\_array\_enable\_m (def none)} \\
\hyperref[enum:tex_texel_offset_m]{\texttt{texel\_offset}} & \texttt{11} & \texttt{tex\_texel\_offset\_m (def none)} \\
\textit{Reserved} & \texttt{12} & \texttt{zero} \\
\hyperref[enum:tex_coordinate_mode_m]{\texttt{coordinate\_mode}} & \texttt{13} & \texttt{tex\_coordinate\_mode\_m (def float\_coordinates)} \\
\hyperref[enum:tex_footprint_m]{\texttt{footprint}} & \texttt{14:15} & \texttt{exact=footprint\_2x2} \\
\texttt{staging\_register\_write\_index} & \texttt{16:21} & \texttt{uint} \\
\hyperref[enum:tex_write_mask_m]{\texttt{write\_mask}} & \texttt{22:25} & \texttt{tex\_write\_mask\_m (def rgba)} \\
\textit{Reserved} & \texttt{26:27} & \texttt{ignore} \\
\hyperref[enum:tex_dimensionality_m]{\texttt{dimensionality}} & \texttt{28:29} & \texttt{exact=tex2d} \\
\hyperref[enum:message_slot_index_m]{\texttt{message\_slot\_index}} & \texttt{30:32} & \texttt{message\_slot\_index\_m} \\
\texttt{sr\_count} & \texttt{33:35} & \texttt{sr\_count\_tex\_t} \\
\texttt{sr\_write\_count} & \texttt{36:38} & \texttt{sr\_write\_count\_t} \\
\hyperref[enum:tex_skip_m]{\texttt{skip}} & \texttt{39} & \texttt{tex\_skip\_m (def none)} \\
\texttt{staging\_register\_index} & \texttt{40:45} & \texttt{uint} \\
\hyperref[enum:tex_width_m]{\texttt{register\_width}} & \texttt{46} & \texttt{tex\_width\_m} \\
\texttt{secondary\_width} & \texttt{47} & \texttt{exact=0} \\
\texttt{opcode1} & \texttt{48:51} & \texttt{exact=TEX\_WEIGHTS} \\
\texttt{instruction\_class} & \texttt{52:56} & \texttt{exact=S2RW} \\
\texttt{fau\_page\_index} & \texttt{57:58} & \texttt{uint} \\
\hyperref[enum:flow_control_m]{\texttt{flow}} & \texttt{59:62} & \texttt{flow\_control\_m} \\
\textit{Reserved} & \texttt{63} & \texttt{zero} \\
\bottomrule
\end{tabularx}
\end{center}
\textbf{Fields Explained:}
\begin{tightitemize}\footnotesize
  \item \texttt{src0}: Primary source operand (\texttt{r0..r63}, \texttt{u0..u63} uniform, \texttt{k*} FAU constant, or \texttt{t0..t2} temporary; \texttt{\textasciicircum{}} marks discard/last-use). \emph{Example:} {\scriptsize\ttfamily \mnem{FADD.f32} \dst{r0}, \src{r1}, \src{r2}} (\textbf{SPIR-V:} \texttt{OpFAdd (Operand 0)}; \textbf{GLSL:} \texttt{float c = a + b;})
  \item \hyperref[enum:tex_wide_indices_enable_m]{\texttt{wide\_indices}}: Uses 32-bit wide texture/sampler table indices instead of compact 16-bit packed indices. \emph{Example:} {\scriptsize\ttfamily \mnem{TEX\_SINGLE}\imod{.tex2d.wide} \dst{@r0:r3}, \src{@r4:r5}, \src{r6}} (\textbf{SPIR-V:} \texttt{OpImageSampleImplicitLod (bindless)}; \textbf{GLSL:} \texttt{vec4 c = texture(textures[nonuniformEXT(i)], uv);})
  \item \hyperref[enum:tex_array_enable_m]{\texttt{array\_enable}}: Enables array-layer coordinate in texture staging registers (\texttt{sampler2DArray} / \texttt{samplerCubeArray}). \emph{Example:} {\scriptsize\ttfamily \mnem{TEX\_SINGLE}\imod{.tex2d.array} \dst{@r0:r3}, \src{@r4:r6}} (\textbf{SPIR-V:} \texttt{OpImageSampleImplicitLod (Arrayed)}; \textbf{GLSL:} \texttt{vec4 c = texture(texArr, vec3(uv, layer));})
  \item \hyperref[enum:tex_texel_offset_m]{\texttt{texel\_offset}}: Applies a signed 4-bit per-axis texel offset $(dx, dy, dz)$ packed in staging registers. \emph{Example:} {\scriptsize\ttfamily \mnem{TEX\_SINGLE}\imod{.tex2d.texel\_offset} \dst{@r0:r3}, \src{@r4:r6}} (\textbf{SPIR-V:} \texttt{OpImageSampleImplicitLod (ConstOffset)}; \textbf{GLSL:} \texttt{vec4 c = textureOffset(tex, uv, ivec2(1, -1));})
  \item \hyperref[enum:tex_coordinate_mode_m]{\texttt{coordinate\_mode}}: Selects normalized float coordinates vs unnormalized integer texel indices. \emph{Example:} {\scriptsize\ttfamily \mnem{TEX\_FETCH}\imod{.tex2d.integer\_coordinates} \dst{@r0:r3}, \src{@r4:r5}} (\textbf{SPIR-V:} \texttt{OpImageFetch}; \textbf{GLSL:} \texttt{vec4 c = texelFetch(tex, ivec2(xy), 0);})
  \item \hyperref[enum:tex_footprint_m]{\texttt{footprint}}: Selects 2x2 bilinear footprint in \texttt{TEX\_WEIGHTS}. \emph{Example:} {\scriptsize\ttfamily \mnem{TEX\_WEIGHTS}\imod{.footprint\_2x2} \dst{@r0:r3}, \src{@r4:r5}} (\textbf{SPIR-V:} \texttt{OpImageSampleFootprintNV}; \textbf{GLSL:} \texttt{// Bilinear weight query})
  \item \texttt{staging\_register\_write\_index}: Separate destination staging register window base when read and write staging windows differ. \emph{Example:} {\scriptsize\ttfamily \mnem{ATOM\_RETURN.i32}\imod{.aadd} \dst{@r0}, \src{@r1}, \src{r2}, \imm{\#0}} (\textbf{SPIR-V:} \texttt{OpAtomicIAdd}; \textbf{GLSL:} \texttt{uint old = atomicAdd(ssbo.cnt, 1u);})
  \item \hyperref[enum:tex_write_mask_m]{\texttt{write\_mask}}: Channel enable bitmask (\texttt{r}..\texttt{rgba}) controlling which texture components are written to staging GPRs. \emph{Example:} {\scriptsize\ttfamily \mnem{TEX\_SINGLE}\imod{.tex2d.f32.rg} \dst{@r0:r1}, \src{@r4:r5}} (\textbf{SPIR-V:} \texttt{OpImageSampleImplicitLod + OpVectorShuffle}; \textbf{GLSL:} \texttt{vec2 rg = texture(tex, uv).rg;})
  \item \hyperref[enum:tex_dimensionality_m]{\texttt{dimensionality}}: Texture dimensionality (\texttt{tex1d}, \texttt{tex2d}, \texttt{tex3d}, \texttt{cube}). \emph{Example:} {\scriptsize\ttfamily \mnem{TEX\_SINGLE}\imod{.tex3d.f32.rgba} \dst{@r0:r3}, \src{@r4:r6}} (\textbf{SPIR-V:} \texttt{OpImageSampleImplicitLod (Dim=3D)}; \textbf{GLSL:} \texttt{vec4 c = texture(volTex, uvw);})
  \item \hyperref[enum:message_slot_index_m]{\texttt{message\_slot\_index}}: Asynchronous scoreboard slot (\texttt{slot0}--\texttt{slot2}, \texttt{slot7}) tracking completion of this message. \emph{Example:} {\scriptsize\ttfamily \mnem{LOAD.i32}\imod{.slot0} \dst{@r0}, \src{r2}, \imm{\#0}} (\textbf{SPIR-V:} \texttt{OpLoad}; \textbf{GLSL:} \texttt{uint v = ssbo.arr[i];})
  \item \texttt{sr\_count}: Number of consecutive 32-bit staging registers transferred by a load/store/texture message. \emph{Example:} {\scriptsize\ttfamily \mnem{LD\_MULTI.i128} \dst{@r0:r3}, \src{r4}, \imm{\#0}} (\textbf{SPIR-V:} \texttt{OpLoad (wide vector)}; \textbf{GLSL:} \texttt{uvec4 x = buf[i];})
  \item \texttt{sr\_write\_count}: Number of consecutive 32-bit staging registers written back by a dual-staging message (e.g. ATOM\textbackslash\{\}\_RETURN / TEX). \emph{Example:} {\scriptsize\ttfamily \mnem{TEX\_SINGLE}\imod{.tex2d.f32.rgba} \dst{@r0:r3}, \src{@r4:r5}} (\textbf{SPIR-V:} \texttt{OpImageSampleImplicitLod}; \textbf{GLSL:} \texttt{vec4 c = texture(tex, uv);})
  \item \hyperref[enum:tex_skip_m]{\texttt{skip}}: Skips texture/memory fetch on helper/discarded fragment lanes while preserving 2x2 quad derivatives. \emph{Example:} {\scriptsize\ttfamily \mnem{TEX\_SINGLE}\imod{.tex2d.computed.skip} \dst{@r0:r3}, \src{@r4:r5}} (\textbf{SPIR-V:} \texttt{OpImageSampleImplicitLod (after OpDemoteToHelperInvocation)}; \textbf{GLSL:} \texttt{if (a \textless{} 0.0) discard; vec4 c = texture(s, uv);})
  \item \texttt{staging\_register\_index}: Base GPR index of the contiguous staging register window (\texttt{@rN:rM}) for message payload/return. \emph{Example:} {\scriptsize\ttfamily \mnem{STORE.i64}\imod{.slot7} \src{@r0:r1}, \src{r2}, \imm{\#8}} (\textbf{SPIR-V:} \texttt{OpStore (uvec2)}; \textbf{GLSL:} \texttt{ssbo.pairs[i] = uvec2(a, b);})
  \item \hyperref[enum:tex_width_m]{\texttt{register\_width}}: Selects 16-bit vs 32-bit GPR element width for the corresponding message return window. \emph{Example:} {\scriptsize\ttfamily \mnem{TEX\_DUAL}\imod{.primary\_f16.secondary\_f32} \dst{@r0:r2}, \src{@r4:r5}} (\textbf{SPIR-V:} \texttt{OpImageSampleImplicitLod (RelaxedPrecision)}; \textbf{GLSL:} \texttt{mediump vec2 a = texture(t0, uv).rg;})
  \item \texttt{secondary\_width}: Selects 16-bit vs 32-bit GPR element width for the corresponding message return window. \emph{Example:} {\scriptsize\ttfamily \mnem{TEX\_DUAL}\imod{.primary\_f16.secondary\_f32} \dst{@r0:r2}, \src{@r4:r5}} (\textbf{SPIR-V:} \texttt{OpImageSampleImplicitLod (RelaxedPrecision)}; \textbf{GLSL:} \texttt{mediump vec2 a = texture(t0, uv).rg;})
  \item \texttt{opcode1}: Primary sub-opcode within the instruction class uniquely identifying the operation family. \emph{Example:} {\scriptsize\ttfamily \mnem{IADD.s32} \dst{r0}, \src{r1}, \src{r2}} (\textbf{SPIR-V:} \texttt{OpIAdd (operation selector)}; \textbf{GLSL:} \texttt{// Primary opcode bits})
  \item \texttt{instruction\_class}: 4-bit major instruction class selector (\texttt{bits 48:51}) partitioning the 64-bit Valhall opcode space (e.g., \texttt{FMA}, \texttt{CVT}, \texttt{SFU}, \texttt{LS}, \texttt{TEX}, \texttt{BRANCH}). \emph{Example:} {\scriptsize\ttfamily \mnem{FADD.f32} \dst{r0}, \src{r1}, \src{r2}} (\textbf{SPIR-V:} \texttt{N/A (encoding class)}; \textbf{GLSL:} \texttt{// Hardware execution pipe / decode class})
  \item \texttt{fau\_page\_index}: Selects the FAU (Fast Access Uniform) page (\texttt{0..3}) when a source operand references hardware small constants, thread IDs, or lane constants. \emph{Example:} {\scriptsize\ttfamily \mnem{FADD.f32} \dst{r0}, \src{r1}, \src{k1f}} (\textbf{SPIR-V:} \texttt{OpConstant / BuiltIn}; \textbf{GLSL:} \texttt{float y = x + 1.0;})
  \item \hyperref[enum:flow_control_m]{\texttt{flow}}: 4-bit embedded flow-control \& scoreboard wait modifier (\texttt{bits 59:62}: \texttt{none}, \texttt{wait0..wait0126}, \texttt{wait}, \texttt{reconverge}, \texttt{discard}, \texttt{end}). \emph{Example:} {\scriptsize\ttfamily \mnem{FADD.f32}\imod{.wait0.reconverge} \dst{r0}, \src{r1}, \src{r2}} (\textbf{SPIR-V:} \texttt{OpControlBarrier / OpReturn}; \textbf{GLSL:} \texttt{// Scoreboard wait + warp reconvergence})
\end{tightitemize}
\clearpage

\subsection*{VAR\_TEX\_BUF\_DUAL}
\addcontentsline{toc}{subsection}{VAR\_TEX\_BUF\_DUAL}
Fuses descriptor-less varying-buffer UV interpolation (at byte offset \texttt{src2}) with a dual texture sample (\texttt{TEX\_DUAL} via handle \texttt{src0}), writing both 4-component sample results into \texttt{@sw} in a single message.
\par\textbf{Notes:}
\begin{tightitemize}
  \item Fragment-only (\texttt{encode\_v9.rs:3535-3585}).
  \item \textbf{Source references:} \texttt{kraid/ops.rs:4649-4708}, \texttt{encode\_v9.rs:3535-3585}, \texttt{bifrost/valhall/ISA.xml:155}.
\end{tightitemize}
\codebox{\mnem{VAR\_TEX\_BUF\_DUAL}\imod{\{.message\_slot\_index\}}\imod{\{.flow\}}\imod{\{.skip\}}\imod{\{.source\_format\}}\imod{\{.sample\_position\}}\imod{\{.update\}}\imod{\{.dimensionality\}}\imod{\{.primary\_width\}}\imod{\{.primary\_write\_mask\}}\imod{\{.secondary\_width\}}\imod{\{.secondary\_write\_mask\}}\imod{\{.array\_enable\}}\imod{\{.compare\_enable\}}\imod{\{.lod\_mode\}} \dst{@sw}, \src{src0}, \src{src2}\\\mnem{VAR\_TEX\_BUF\_DUAL}\imod{\{.message\_slot\_index\}}\imod{\{.flow\}}\imod{\{.skip\}}\imod{\{.source\_format\}}\imod{\{.sample\_position\}}\imod{\{.update\}}\imod{\{.dimensionality\}}\imod{\{.projection\_enable\}}\imod{\{.primary\_width\}}\imod{\{.primary\_write\_mask\}}\imod{\{.secondary\_width\}}\imod{\{.secondary\_write\_mask\}}\imod{\{.array\_enable\}}\imod{\{.compare\_enable\}}\imod{\{.lod\_mode\}} \dst{@sw}, \src{src0}, \src{src2}}
\textbf{Operands:}
\begin{tightitemize}\footnotesize
  \item \texttt{\dst{@sw}} (\texttt{staging\_register\_index} + \texttt{sr\_write\_count}): Destination staging register range receiving the dual texture sample result.
  \item \texttt{\src{src0}}: Texture/sampler descriptor handle.
  \item \texttt{\src{src2}}: Descriptor-less varying buffer byte offset providing interpolated UV coordinates (encoded in \texttt{src2}).
\end{tightitemize}
\par\textbf{Closest SPIR-V Op:} \texttt{OpLoad (in) + 2x OpImageSample*}
\par\textbf{Pseudo-GLSL (Operational Semantics):}
\codebox{vec2 uv = interpolateVaryingBuf(byte\_off = src2);\\@sw[0:3] = textureSample(tex0\_of(src0), uv); @sw[4:7] = textureSample(tex1\_of(src0), uv);}
\textbf{Execution Context:}
\begin{tightitemize}\footnotesize
  \item \textbf{Varying buffer \& barycentric cache (\hyperref[sec:csf-tiler]{\texttt{IDVSSR::VARY\_SIZE}}, \texttt{r48:r49}):} Reads per-vertex varying RAM written by Varying EBIN \texttt{STORE.estream}; descriptor forms (\texttt{LD\_VAR*}) also chase \hyperref[sec:csf-srt]{\texttt{SRT\_2}} $\to$ \texttt{Resource[table]} $\to$ \texttt{Attribute[index]} and reuse cached barycentric weights via \texttt{.store}/\texttt{.retrieve}.
  \item \textbf{Texture \& Sampler chain (\hyperref[sec:csf-srt]{\texttt{SRT\_2}} Frag / \texttt{SRT\_0} Comp/Vert):} 64B-aligned VA (\texttt{\& \textasciitilde{}0x3F}) $\to$ \texttt{Resource[table]} $\to$ \hyperref[sec:csf-tex-samp]{\texttt{Texture[tex\_idx]}} (32B, \texttt{type=2}) $\to$ \texttt{GenericPlane[mip]} (32B, \texttt{type=10}: \texttt{pointer}, \texttt{row\_stride}, Morton 16$\times$16 / AFBC) + \texttt{Sampler[samp\_idx]} (32B, \texttt{type=1}: filter, wrap, LOD clamp/bias, compare).
  \item \textbf{Fragment stage preloads \& FP modes (\hyperref[sec:csf-spd]{\texttt{SPD\_2}}, \texttt{r20:r21}):} Hardware pre-fills \texttt{r59} with \texttt{floatBitsToUint(FragCoord.z)} and \texttt{r60} with the MSAA sample coverage mask (\texttt{gl\_SampleMaskIn[0]}), while \texttt{ShaderProgram} word~0 latches 32/64 GPR allocation and warp-wide \texttt{flush\_to\_zero\_mode}, \texttt{suppress\_na\_n}, and \texttt{suppress\_inf}.
\end{tightitemize}
\par\textbf{Worked Examples:}
\par\textbf{Example 1 (Fused varying-buffer interpolation + dual sample):}
\codebox{\mnem{VAR\_TEX\_BUF\_DUAL}\imod{.tex2d.computed.slot0} \dst{@r0:r7}, \src{u0}, \src{r59}, \src{r4}}
Interpolates \texttt{vUV} directly from varying buffer offset \texttt{r4} and feeds both texture units without intermediate UV GPRs.
\par\textbf{Example 2 (FP16 varying source (`.src\_f16`)):}
\codebox{\mnem{VAR\_TEX\_BUF\_DUAL}\imod{.tex2d.src\_f16.computed.slot0} \dst{@r0:r7}, \src{u0}, \src{r59}, \src{r4}}
Reads \texttt{mediump} FP16 UVs from the varying buffer and issues paired filtered samples.

\begin{center}\footnotesize
\begin{tabularx}{\linewidth}{@{}lllX@{}}
\toprule
Variant & Arch & Class & Opcode Fields \\
\midrule
\texttt{(none)} & v9--v14 & \texttt{S3A} & \texttt{opcode1=VAR\_TEX\_BUF\_DUAL} \\
\texttt{(none)} & v9--v14 & \texttt{S3A} & \texttt{opcode1=VAR\_TEX\_BUF\_DUAL} \\
\texttt{(none)} & v9--v14 & \texttt{S3A} & \texttt{opcode1=VAR\_TEX\_BUF\_DUAL} \\
\bottomrule
\end{tabularx}
\end{center}
\begin{center}\footnotesize
\begin{tabularx}{0.88\linewidth}{@{}l l X@{}}
\toprule
Field Name & Bits & Type / Fixed Value \\
\midrule
\texttt{src0} & \texttt{0:7} & \texttt{SourceEncoding64} \\
\hyperref[enum:tex_wide_indices_enable_m]{\texttt{wide\_indices}} & \texttt{8} & \texttt{exact=none} \\
\textit{Reserved} & \texttt{9} & \texttt{zero} \\
\hyperref[enum:var_tex_dual_array_enable_m]{\texttt{array\_enable}} & \texttt{10} & \texttt{var\_tex\_dual\_array\_enable\_m (def none)} \\
\hyperref[enum:var_tex_buf_source_format_m]{\texttt{source\_format}} & \texttt{11} & \texttt{exact=src\_f32} \\
\hyperref[enum:var_tex_compare_enable_m]{\texttt{compare\_enable}} & \texttt{12} & \texttt{var\_tex\_compare\_enable\_m (def none)} \\
\hyperref[enum:var_tex_dual_lod_mode_m]{\texttt{lod\_mode}} & \texttt{13:15} & \texttt{var\_tex\_dual\_lod\_mode\_m} \\
\texttt{src2} & \texttt{16:23} & \texttt{SourceEncoding} \\
\hyperref[enum:tex_dual_primary_width_m]{\texttt{primary\_width}} & \texttt{24} & \texttt{tex\_dual\_primary\_width\_m} \\
\hyperref[enum:tex_dual_secondary_width_m]{\texttt{secondary\_width}} & \texttt{25} & \texttt{tex\_dual\_secondary\_width\_m} \\
\textit{Reserved} & \texttt{26:27} & \texttt{ignore} \\
\hyperref[enum:var_tex_dimensionality_m]{\texttt{dimensionality}} & \texttt{28:29} & \texttt{var\_tex\_dimensionality\_m} \\
\hyperref[enum:message_slot_index_m]{\texttt{message\_slot\_index}} & \texttt{30:32} & \texttt{message\_slot\_index\_m} \\
\hyperref[enum:encoded_sample_and_update_m]{\texttt{encoded\_sample\_and\_update}} & \texttt{33:35} & \texttt{encoded\_sample\_and\_update\_m} \\
\texttt{sr\_write\_count} & \texttt{36:38} & \texttt{sr\_write\_count\_t} \\
\hyperref[enum:tex_skip_m]{\texttt{skip}} & \texttt{39} & \texttt{tex\_skip\_m (def none)} \\
\texttt{staging\_register\_index} & \texttt{40:45} & \texttt{uint} \\
\texttt{staging\_register\_control\_32} & \texttt{46:47} & \texttt{exact=write} \\
\texttt{opcode1} & \texttt{48:51} & \texttt{exact=VAR\_TEX\_BUF\_DUAL} \\
\texttt{instruction\_class} & \texttt{52:56} & \texttt{exact=S3A} \\
\texttt{fau\_page\_index} & \texttt{57:58} & \texttt{uint} \\
\hyperref[enum:flow_control_m]{\texttt{flow}} & \texttt{59:62} & \texttt{flow\_control\_m} \\
\textit{Reserved} & \texttt{63} & \texttt{zero} \\
\bottomrule
\end{tabularx}
\end{center}
\textbf{Fields Explained:}
\begin{tightitemize}\footnotesize
  \item \texttt{src0}: Primary source operand (\texttt{r0..r63}, \texttt{u0..u63} uniform, \texttt{k*} FAU constant, or \texttt{t0..t2} temporary; \texttt{\textasciicircum{}} marks discard/last-use). \emph{Example:} {\scriptsize\ttfamily \mnem{FADD.f32} \dst{r0}, \src{r1}, \src{r2}} (\textbf{SPIR-V:} \texttt{OpFAdd (Operand 0)}; \textbf{GLSL:} \texttt{float c = a + b;})
  \item \hyperref[enum:tex_wide_indices_enable_m]{\texttt{wide\_indices}}: Uses 32-bit wide texture/sampler table indices instead of compact 16-bit packed indices. \emph{Example:} {\scriptsize\ttfamily \mnem{TEX\_SINGLE}\imod{.tex2d.wide} \dst{@r0:r3}, \src{@r4:r5}, \src{r6}} (\textbf{SPIR-V:} \texttt{OpImageSampleImplicitLod (bindless)}; \textbf{GLSL:} \texttt{vec4 c = texture(textures[nonuniformEXT(i)], uv);})
  \item \hyperref[enum:var_tex_dual_array_enable_m]{\texttt{array\_enable}}: Enables array-layer coordinate in texture staging registers (\texttt{sampler2DArray} / \texttt{samplerCubeArray}). \emph{Example:} {\scriptsize\ttfamily \mnem{TEX\_SINGLE}\imod{.tex2d.array} \dst{@r0:r3}, \src{@r4:r6}} (\textbf{SPIR-V:} \texttt{OpImageSampleImplicitLod (Arrayed)}; \textbf{GLSL:} \texttt{vec4 c = texture(texArr, vec3(uv, layer));})
  \item \hyperref[enum:var_tex_buf_source_format_m]{\texttt{source\_format}}: Declares the storage precision of the varying buffer (\texttt{src\_f32}, \texttt{src\_f16}, \texttt{src\_flat32}, \texttt{src\_flat16}). \emph{Example:} {\scriptsize\ttfamily \mnem{LD\_VAR\_BUF.f32}\imod{.src\_f16.v2} \dst{@r0:r1}, \src{r2}} (\textbf{SPIR-V:} \texttt{OpLoad (RelaxedPrecision Input)}; \textbf{GLSL:} \texttt{mediump in vec2 v\_uv;})
  \item \hyperref[enum:var_tex_compare_enable_m]{\texttt{compare\_enable}}: Enables depth-reference shadow comparison (\texttt{Dref} in staging registers). \emph{Example:} {\scriptsize\ttfamily \mnem{TEX\_SINGLE}\imod{.tex2d.compare} \dst{@r0}, \src{@r4:r6}} (\textbf{SPIR-V:} \texttt{OpImageSampleDrefImplicitLod}; \textbf{GLSL:} \texttt{float s = texture(shadowMap, vec3(uv, refZ));})
  \item \hyperref[enum:var_tex_dual_lod_mode_m]{\texttt{lod\_mode}}: Selects mipmap LOD calculation mode (\texttt{none}, \texttt{computed}, \texttt{explicit}, \texttt{computed\_bias}, \texttt{grdesc}). \emph{Example:} {\scriptsize\ttfamily \mnem{TEX\_SINGLE}\imod{.tex2d.explicit} \dst{@r0:r3}, \src{@r4:r6}} (\textbf{SPIR-V:} \texttt{OpImageSampleExplicitLod}; \textbf{GLSL:} \texttt{vec4 c = textureLod(tex, uv, 2.0);})
  \item \texttt{src2}: Third source operand (FMA addend, \texttt{MUX}/\texttt{CSEL} selector, \texttt{IADDC} carry-in, or bitwise mask in fused shift-logic). \emph{Example:} {\scriptsize\ttfamily \mnem{FMA.f32} \dst{r0}, \src{r1}, \src{r2}, \src{r3}} (\textbf{SPIR-V:} \texttt{GLSL.std.450 Fma (Addend)}; \textbf{GLSL:} \texttt{float d = fma(a, b, c);})
  \item \hyperref[enum:tex_dual_primary_width_m]{\texttt{primary\_width}}: Selects 16-bit vs 32-bit GPR element width for the corresponding message return window. \emph{Example:} {\scriptsize\ttfamily \mnem{TEX\_DUAL}\imod{.primary\_f16.secondary\_f32} \dst{@r0:r2}, \src{@r4:r5}} (\textbf{SPIR-V:} \texttt{OpImageSampleImplicitLod (RelaxedPrecision)}; \textbf{GLSL:} \texttt{mediump vec2 a = texture(t0, uv).rg;})
  \item \hyperref[enum:tex_dual_secondary_width_m]{\texttt{secondary\_width}}: Selects 16-bit vs 32-bit GPR element width for the corresponding message return window. \emph{Example:} {\scriptsize\ttfamily \mnem{TEX\_DUAL}\imod{.primary\_f16.secondary\_f32} \dst{@r0:r2}, \src{@r4:r5}} (\textbf{SPIR-V:} \texttt{OpImageSampleImplicitLod (RelaxedPrecision)}; \textbf{GLSL:} \texttt{mediump vec2 a = texture(t0, uv).rg;})
  \item \hyperref[enum:var_tex_dimensionality_m]{\texttt{dimensionality}}: Texture dimensionality (\texttt{tex1d}, \texttt{tex2d}, \texttt{tex3d}, \texttt{cube}). \emph{Example:} {\scriptsize\ttfamily \mnem{TEX\_SINGLE}\imod{.tex3d.f32.rgba} \dst{@r0:r3}, \src{@r4:r6}} (\textbf{SPIR-V:} \texttt{OpImageSampleImplicitLod (Dim=3D)}; \textbf{GLSL:} \texttt{vec4 c = texture(volTex, uvw);})
  \item \hyperref[enum:message_slot_index_m]{\texttt{message\_slot\_index}}: Asynchronous scoreboard slot (\texttt{slot0}--\texttt{slot2}, \texttt{slot7}) tracking completion of this message. \emph{Example:} {\scriptsize\ttfamily \mnem{LOAD.i32}\imod{.slot0} \dst{@r0}, \src{r2}, \imm{\#0}} (\textbf{SPIR-V:} \texttt{OpLoad}; \textbf{GLSL:} \texttt{uint v = ssbo.arr[i];})
  \item \hyperref[enum:encoded_sample_and_update_m]{\texttt{encoded\_sample\_and\_update}}: Jointly encodes MSAA sample position and hardware varying-cache update mode (\texttt{center\_store}, \texttt{retrieve}, etc.). \emph{Example:} {\scriptsize\ttfamily \mnem{VAR\_TEX\_BUF\_SINGLE}\imod{.center\_store.tex2d} \dst{@r0:r3}, \src{r4}} (\textbf{SPIR-V:} \texttt{OpImageSampleImplicitLod (fused varying)}; \textbf{GLSL:} \texttt{vec4 c = texture(s, v\_uv);})
  \item \texttt{sr\_write\_count}: Number of consecutive 32-bit staging registers written back by a dual-staging message (e.g. ATOM\textbackslash\{\}\_RETURN / TEX). \emph{Example:} {\scriptsize\ttfamily \mnem{TEX\_SINGLE}\imod{.tex2d.f32.rgba} \dst{@r0:r3}, \src{@r4:r5}} (\textbf{SPIR-V:} \texttt{OpImageSampleImplicitLod}; \textbf{GLSL:} \texttt{vec4 c = texture(tex, uv);})
  \item \hyperref[enum:tex_skip_m]{\texttt{skip}}: Skips texture/memory fetch on helper/discarded fragment lanes while preserving 2x2 quad derivatives. \emph{Example:} {\scriptsize\ttfamily \mnem{TEX\_SINGLE}\imod{.tex2d.computed.skip} \dst{@r0:r3}, \src{@r4:r5}} (\textbf{SPIR-V:} \texttt{OpImageSampleImplicitLod (after OpDemoteToHelperInvocation)}; \textbf{GLSL:} \texttt{if (a \textless{} 0.0) discard; vec4 c = texture(s, uv);})
  \item \texttt{staging\_register\_index}: Base GPR index of the contiguous staging register window (\texttt{@rN:rM}) for message payload/return. \emph{Example:} {\scriptsize\ttfamily \mnem{STORE.i64}\imod{.slot7} \src{@r0:r1}, \src{r2}, \imm{\#8}} (\textbf{SPIR-V:} \texttt{OpStore (uvec2)}; \textbf{GLSL:} \texttt{ssbo.pairs[i] = uvec2(a, b);})
  \item \texttt{staging\_register\_control\_32}: Controls the read/write staging register bank and count for async message instructions. \emph{Example:} {\scriptsize\ttfamily \mnem{LOAD.i128}\imod{.slot0} \dst{@r0:r3}, \src{r4}, \imm{\#0}} (\textbf{SPIR-V:} \texttt{OpLoad (vec4)}; \textbf{GLSL:} \texttt{vec4 v = ssbo.data[idx];})
  \item \texttt{opcode1}: Primary sub-opcode within the instruction class uniquely identifying the operation family. \emph{Example:} {\scriptsize\ttfamily \mnem{IADD.s32} \dst{r0}, \src{r1}, \src{r2}} (\textbf{SPIR-V:} \texttt{OpIAdd (operation selector)}; \textbf{GLSL:} \texttt{// Primary opcode bits})
  \item \texttt{instruction\_class}: 4-bit major instruction class selector (\texttt{bits 48:51}) partitioning the 64-bit Valhall opcode space (e.g., \texttt{FMA}, \texttt{CVT}, \texttt{SFU}, \texttt{LS}, \texttt{TEX}, \texttt{BRANCH}). \emph{Example:} {\scriptsize\ttfamily \mnem{FADD.f32} \dst{r0}, \src{r1}, \src{r2}} (\textbf{SPIR-V:} \texttt{N/A (encoding class)}; \textbf{GLSL:} \texttt{// Hardware execution pipe / decode class})
  \item \texttt{fau\_page\_index}: Selects the FAU (Fast Access Uniform) page (\texttt{0..3}) when a source operand references hardware small constants, thread IDs, or lane constants. \emph{Example:} {\scriptsize\ttfamily \mnem{FADD.f32} \dst{r0}, \src{r1}, \src{k1f}} (\textbf{SPIR-V:} \texttt{OpConstant / BuiltIn}; \textbf{GLSL:} \texttt{float y = x + 1.0;})
  \item \hyperref[enum:flow_control_m]{\texttt{flow}}: 4-bit embedded flow-control \& scoreboard wait modifier (\texttt{bits 59:62}: \texttt{none}, \texttt{wait0..wait0126}, \texttt{wait}, \texttt{reconverge}, \texttt{discard}, \texttt{end}). \emph{Example:} {\scriptsize\ttfamily \mnem{FADD.f32}\imod{.wait0.reconverge} \dst{r0}, \src{r1}, \src{r2}} (\textbf{SPIR-V:} \texttt{OpControlBarrier / OpReturn}; \textbf{GLSL:} \texttt{// Scoreboard wait + warp reconvergence})
\end{tightitemize}
\clearpage

\subsection*{VAR\_TEX\_BUF\_GATHER}
\addcontentsline{toc}{subsection}{VAR\_TEX\_BUF\_GATHER}
Fuses descriptor-less varying-buffer UV interpolation (at byte offset \texttt{src2}) with a 2x2 texture gather (\texttt{TEX\_GATHER} via handle \texttt{src0}), writing the 4 gathered component values into \texttt{@sw} in a single message.
\par\textbf{Notes:}
\begin{tightitemize}
  \item Fragment-only.
  \item 2D (\texttt{encode\_v9.rs:3587-3618}).
  \item \textbf{Source references:} \texttt{kraid/ops.rs:4710-4762}, \texttt{encode\_v9.rs:3587-3618}, \texttt{bifrost/valhall/ISA.xml:157}.
\end{tightitemize}
\codebox{\mnem{VAR\_TEX\_BUF\_GATHER}\imod{\{.message\_slot\_index\}}\imod{\{.flow\}}\imod{\{.skip\}}\imod{\{.source\_format\}}\imod{\{.sample\_position\}}\imod{\{.update\}}\imod{\{.dimensionality\}}\imod{\{.register\_width\}}\imod{\{.write\_mask\}}\imod{\{.wide\_indices\}}\imod{\{.array\_enable\}}\imod{\{.compare\_enable\}}\imod{\{.gather\_component\}} \dst{@sw}, \src{src0}, \src{src2}\\\mnem{VAR\_TEX\_BUF\_GATHER}\imod{\{.message\_slot\_index\}}\imod{\{.flow\}}\imod{\{.skip\}}\imod{\{.source\_format\}}\imod{\{.sample\_position\}}\imod{\{.update\}}\imod{\{.dimensionality\}}\imod{\{.projection\_enable\}}\imod{\{.register\_width\}}\imod{\{.write\_mask\}}\imod{\{.wide\_indices\}}\imod{\{.array\_enable\}}\imod{\{.compare\_enable\}}\imod{\{.gather\_component\}} \dst{@sw}, \src{src0}, \src{src2}}
\textbf{Operands:}
\begin{tightitemize}\footnotesize
  \item \texttt{\dst{@sw}} (\texttt{staging\_register\_index} + \texttt{sr\_write\_count}): Destination staging register range receiving the 2x2 texture gather result.
  \item \texttt{\src{src0}}: Texture/sampler descriptor handle.
  \item \texttt{\src{src2}}: Descriptor-less varying buffer byte offset providing interpolated UV coordinates (encoded in \texttt{src2}).
\end{tightitemize}
\par\textbf{Closest SPIR-V Op:} \texttt{OpLoad (in) + OpImageGather}
\par\textbf{Pseudo-GLSL (Operational Semantics):}
\codebox{@sw = textureGather(handle = src0, interpolateVaryingBuf(byte\_off = src2), gather\_component);}
\textbf{Execution Context:}
\begin{tightitemize}\footnotesize
  \item \textbf{Varying buffer \& barycentric cache (\hyperref[sec:csf-tiler]{\texttt{IDVSSR::VARY\_SIZE}}, \texttt{r48:r49}):} Reads per-vertex varying RAM written by Varying EBIN \texttt{STORE.estream}; descriptor forms (\texttt{LD\_VAR*}) also chase \hyperref[sec:csf-srt]{\texttt{SRT\_2}} $\to$ \texttt{Resource[table]} $\to$ \texttt{Attribute[index]} and reuse cached barycentric weights via \texttt{.store}/\texttt{.retrieve}.
  \item \textbf{Texture \& Sampler chain (\hyperref[sec:csf-srt]{\texttt{SRT\_2}} Frag / \texttt{SRT\_0} Comp/Vert):} 64B-aligned VA (\texttt{\& \textasciitilde{}0x3F}) $\to$ \texttt{Resource[table]} $\to$ \hyperref[sec:csf-tex-samp]{\texttt{Texture[tex\_idx]}} (32B, \texttt{type=2}) $\to$ \texttt{GenericPlane[mip]} (32B, \texttt{type=10}: \texttt{pointer}, \texttt{row\_stride}, Morton 16$\times$16 / AFBC) + \texttt{Sampler[samp\_idx]} (32B, \texttt{type=1}: filter, wrap, LOD clamp/bias, compare).
  \item \textbf{Fragment stage preloads \& FP modes (\hyperref[sec:csf-spd]{\texttt{SPD\_2}}, \texttt{r20:r21}):} Hardware pre-fills \texttt{r59} with \texttt{floatBitsToUint(FragCoord.z)} and \texttt{r60} with the MSAA sample coverage mask (\texttt{gl\_SampleMaskIn[0]}), while \texttt{ShaderProgram} word~0 latches 32/64 GPR allocation and warp-wide \texttt{flush\_to\_zero\_mode}, \texttt{suppress\_na\_n}, and \texttt{suppress\_inf}.
\end{tightitemize}
\par\textbf{Worked Examples:}
\par\textbf{Example 1 (Fused varying-buffer gather (red channel)):}
\codebox{\mnem{VAR\_TEX\_BUF\_GATHER}\imod{.tex2d.gather4\_r.slot0} \dst{@r0:r3}, \src{u0}, \src{r59}, \src{r4}}
Interpolates \texttt{vUV} from varying buffer offset \texttt{r4} and gathers the 2x2 red-channel footprint in one message.
\par\textbf{Example 2 (Fused shadow-compare gather (`.compare`)):}
\codebox{\mnem{VAR\_TEX\_BUF\_GATHER}\imod{.tex2d.compare.slot0} \dst{@r0:r3}, \src{u0}, \src{r59}, \src{r4}}
Interpolates \texttt{(u, v, refZ)} from varying buffer and returns 4 shadow-comparison results.

\begin{center}\footnotesize
\begin{tabularx}{\linewidth}{@{}lllX@{}}
\toprule
Variant & Arch & Class & Opcode Fields \\
\midrule
\texttt{(none)} & v9--v14 & \texttt{S3A} & \texttt{opcode1=VAR\_TEX\_BUF\_GATHER} \\
\texttt{(none)} & v9--v14 & \texttt{S3A} & \texttt{opcode1=VAR\_TEX\_BUF\_GATHER} \\
\texttt{(none)} & v9--v14 & \texttt{S3A} & \texttt{opcode1=VAR\_TEX\_BUF\_GATHER} \\
\bottomrule
\end{tabularx}
\end{center}
\begin{center}\footnotesize
\begin{tabularx}{0.88\linewidth}{@{}l l X@{}}
\toprule
Field Name & Bits & Type / Fixed Value \\
\midrule
\texttt{src0} & \texttt{0:7} & \texttt{SourceEncoding64} \\
\hyperref[enum:tex_wide_indices_enable_m]{\texttt{wide\_indices}} & \texttt{8} & \texttt{tex\_wide\_indices\_enable\_m (def none)} \\
\textit{Reserved} & \texttt{9} & \texttt{zero} \\
\hyperref[enum:var_tex_array_enable_m]{\texttt{array\_enable}} & \texttt{10} & \texttt{var\_tex\_array\_enable\_m (def none)} \\
\hyperref[enum:var_tex_buf_source_format_m]{\texttt{source\_format}} & \texttt{11} & \texttt{exact=src\_f32} \\
\hyperref[enum:var_tex_compare_enable_m]{\texttt{compare\_enable}} & \texttt{12} & \texttt{var\_tex\_compare\_enable\_m (def none)} \\
\hyperref[enum:tex_coordinate_mode_m]{\texttt{coordinate\_mode}} & \texttt{13} & \texttt{exact=float\_coordinates} \\
\hyperref[enum:tex_gather_component_m]{\texttt{gather\_component}} & \texttt{14:15} & \texttt{tex\_gather\_component\_m} \\
\texttt{src2} & \texttt{16:23} & \texttt{SourceEncoding} \\
\hyperref[enum:tex_width_m]{\texttt{register\_width}} & \texttt{24} & \texttt{tex\_width\_m} \\
\texttt{secondary\_width} & \texttt{25} & \texttt{exact=0} \\
\textit{Reserved} & \texttt{26:27} & \texttt{ignore} \\
\hyperref[enum:var_tex_dimensionality_m]{\texttt{dimensionality}} & \texttt{28:29} & \texttt{exact=tex2d} \\
\hyperref[enum:message_slot_index_m]{\texttt{message\_slot\_index}} & \texttt{30:32} & \texttt{message\_slot\_index\_m} \\
\hyperref[enum:encoded_sample_and_update_m]{\texttt{encoded\_sample\_and\_update}} & \texttt{33:35} & \texttt{encoded\_sample\_and\_update\_m} \\
\texttt{sr\_write\_count} & \texttt{36:38} & \texttt{sr\_write\_count\_t} \\
\hyperref[enum:tex_skip_m]{\texttt{skip}} & \texttt{39} & \texttt{tex\_skip\_m (def none)} \\
\texttt{staging\_register\_index} & \texttt{40:45} & \texttt{uint} \\
\texttt{staging\_register\_control\_32} & \texttt{46:47} & \texttt{exact=write} \\
\texttt{opcode1} & \texttt{48:51} & \texttt{exact=VAR\_TEX\_BUF\_GATHER} \\
\texttt{instruction\_class} & \texttt{52:56} & \texttt{exact=S3A} \\
\texttt{fau\_page\_index} & \texttt{57:58} & \texttt{uint} \\
\hyperref[enum:flow_control_m]{\texttt{flow}} & \texttt{59:62} & \texttt{flow\_control\_m} \\
\textit{Reserved} & \texttt{63} & \texttt{zero} \\
\bottomrule
\end{tabularx}
\end{center}
\textbf{Fields Explained:}
\begin{tightitemize}\footnotesize
  \item \texttt{src0}: Primary source operand (\texttt{r0..r63}, \texttt{u0..u63} uniform, \texttt{k*} FAU constant, or \texttt{t0..t2} temporary; \texttt{\textasciicircum{}} marks discard/last-use). \emph{Example:} {\scriptsize\ttfamily \mnem{FADD.f32} \dst{r0}, \src{r1}, \src{r2}} (\textbf{SPIR-V:} \texttt{OpFAdd (Operand 0)}; \textbf{GLSL:} \texttt{float c = a + b;})
  \item \hyperref[enum:tex_wide_indices_enable_m]{\texttt{wide\_indices}}: Uses 32-bit wide texture/sampler table indices instead of compact 16-bit packed indices. \emph{Example:} {\scriptsize\ttfamily \mnem{TEX\_SINGLE}\imod{.tex2d.wide} \dst{@r0:r3}, \src{@r4:r5}, \src{r6}} (\textbf{SPIR-V:} \texttt{OpImageSampleImplicitLod (bindless)}; \textbf{GLSL:} \texttt{vec4 c = texture(textures[nonuniformEXT(i)], uv);})
  \item \hyperref[enum:var_tex_array_enable_m]{\texttt{array\_enable}}: Enables array-layer coordinate in texture staging registers (\texttt{sampler2DArray} / \texttt{samplerCubeArray}). \emph{Example:} {\scriptsize\ttfamily \mnem{TEX\_SINGLE}\imod{.tex2d.array} \dst{@r0:r3}, \src{@r4:r6}} (\textbf{SPIR-V:} \texttt{OpImageSampleImplicitLod (Arrayed)}; \textbf{GLSL:} \texttt{vec4 c = texture(texArr, vec3(uv, layer));})
  \item \hyperref[enum:var_tex_buf_source_format_m]{\texttt{source\_format}}: Declares the storage precision of the varying buffer (\texttt{src\_f32}, \texttt{src\_f16}, \texttt{src\_flat32}, \texttt{src\_flat16}). \emph{Example:} {\scriptsize\ttfamily \mnem{LD\_VAR\_BUF.f32}\imod{.src\_f16.v2} \dst{@r0:r1}, \src{r2}} (\textbf{SPIR-V:} \texttt{OpLoad (RelaxedPrecision Input)}; \textbf{GLSL:} \texttt{mediump in vec2 v\_uv;})
  \item \hyperref[enum:var_tex_compare_enable_m]{\texttt{compare\_enable}}: Enables depth-reference shadow comparison (\texttt{Dref} in staging registers). \emph{Example:} {\scriptsize\ttfamily \mnem{TEX\_SINGLE}\imod{.tex2d.compare} \dst{@r0}, \src{@r4:r6}} (\textbf{SPIR-V:} \texttt{OpImageSampleDrefImplicitLod}; \textbf{GLSL:} \texttt{float s = texture(shadowMap, vec3(uv, refZ));})
  \item \hyperref[enum:tex_coordinate_mode_m]{\texttt{coordinate\_mode}}: Selects normalized float coordinates vs unnormalized integer texel indices. \emph{Example:} {\scriptsize\ttfamily \mnem{TEX\_FETCH}\imod{.tex2d.integer\_coordinates} \dst{@r0:r3}, \src{@r4:r5}} (\textbf{SPIR-V:} \texttt{OpImageFetch}; \textbf{GLSL:} \texttt{vec4 c = texelFetch(tex, ivec2(xy), 0);})
  \item \hyperref[enum:tex_gather_component_m]{\texttt{gather\_component}}: Selects the R, G, B, or A channel gathered across the 2x2 texel footprint. \emph{Example:} {\scriptsize\ttfamily \mnem{TEX\_GATHER}\imod{.tex2d.gather4\_g} \dst{@r0:r3}, \src{@r4:r5}} (\textbf{SPIR-V:} \texttt{OpImageGather}; \textbf{GLSL:} \texttt{vec4 g = textureGather(tex, uv, 1);})
  \item \texttt{src2}: Third source operand (FMA addend, \texttt{MUX}/\texttt{CSEL} selector, \texttt{IADDC} carry-in, or bitwise mask in fused shift-logic). \emph{Example:} {\scriptsize\ttfamily \mnem{FMA.f32} \dst{r0}, \src{r1}, \src{r2}, \src{r3}} (\textbf{SPIR-V:} \texttt{GLSL.std.450 Fma (Addend)}; \textbf{GLSL:} \texttt{float d = fma(a, b, c);})
  \item \hyperref[enum:tex_width_m]{\texttt{register\_width}}: Selects 16-bit vs 32-bit GPR element width for the corresponding message return window. \emph{Example:} {\scriptsize\ttfamily \mnem{TEX\_DUAL}\imod{.primary\_f16.secondary\_f32} \dst{@r0:r2}, \src{@r4:r5}} (\textbf{SPIR-V:} \texttt{OpImageSampleImplicitLod (RelaxedPrecision)}; \textbf{GLSL:} \texttt{mediump vec2 a = texture(t0, uv).rg;})
  \item \texttt{secondary\_width}: Selects 16-bit vs 32-bit GPR element width for the corresponding message return window. \emph{Example:} {\scriptsize\ttfamily \mnem{TEX\_DUAL}\imod{.primary\_f16.secondary\_f32} \dst{@r0:r2}, \src{@r4:r5}} (\textbf{SPIR-V:} \texttt{OpImageSampleImplicitLod (RelaxedPrecision)}; \textbf{GLSL:} \texttt{mediump vec2 a = texture(t0, uv).rg;})
  \item \hyperref[enum:var_tex_dimensionality_m]{\texttt{dimensionality}}: Texture dimensionality (\texttt{tex1d}, \texttt{tex2d}, \texttt{tex3d}, \texttt{cube}). \emph{Example:} {\scriptsize\ttfamily \mnem{TEX\_SINGLE}\imod{.tex3d.f32.rgba} \dst{@r0:r3}, \src{@r4:r6}} (\textbf{SPIR-V:} \texttt{OpImageSampleImplicitLod (Dim=3D)}; \textbf{GLSL:} \texttt{vec4 c = texture(volTex, uvw);})
  \item \hyperref[enum:message_slot_index_m]{\texttt{message\_slot\_index}}: Asynchronous scoreboard slot (\texttt{slot0}--\texttt{slot2}, \texttt{slot7}) tracking completion of this message. \emph{Example:} {\scriptsize\ttfamily \mnem{LOAD.i32}\imod{.slot0} \dst{@r0}, \src{r2}, \imm{\#0}} (\textbf{SPIR-V:} \texttt{OpLoad}; \textbf{GLSL:} \texttt{uint v = ssbo.arr[i];})
  \item \hyperref[enum:encoded_sample_and_update_m]{\texttt{encoded\_sample\_and\_update}}: Jointly encodes MSAA sample position and hardware varying-cache update mode (\texttt{center\_store}, \texttt{retrieve}, etc.). \emph{Example:} {\scriptsize\ttfamily \mnem{VAR\_TEX\_BUF\_SINGLE}\imod{.center\_store.tex2d} \dst{@r0:r3}, \src{r4}} (\textbf{SPIR-V:} \texttt{OpImageSampleImplicitLod (fused varying)}; \textbf{GLSL:} \texttt{vec4 c = texture(s, v\_uv);})
  \item \texttt{sr\_write\_count}: Number of consecutive 32-bit staging registers written back by a dual-staging message (e.g. ATOM\textbackslash\{\}\_RETURN / TEX). \emph{Example:} {\scriptsize\ttfamily \mnem{TEX\_SINGLE}\imod{.tex2d.f32.rgba} \dst{@r0:r3}, \src{@r4:r5}} (\textbf{SPIR-V:} \texttt{OpImageSampleImplicitLod}; \textbf{GLSL:} \texttt{vec4 c = texture(tex, uv);})
  \item \hyperref[enum:tex_skip_m]{\texttt{skip}}: Skips texture/memory fetch on helper/discarded fragment lanes while preserving 2x2 quad derivatives. \emph{Example:} {\scriptsize\ttfamily \mnem{TEX\_SINGLE}\imod{.tex2d.computed.skip} \dst{@r0:r3}, \src{@r4:r5}} (\textbf{SPIR-V:} \texttt{OpImageSampleImplicitLod (after OpDemoteToHelperInvocation)}; \textbf{GLSL:} \texttt{if (a \textless{} 0.0) discard; vec4 c = texture(s, uv);})
  \item \texttt{staging\_register\_index}: Base GPR index of the contiguous staging register window (\texttt{@rN:rM}) for message payload/return. \emph{Example:} {\scriptsize\ttfamily \mnem{STORE.i64}\imod{.slot7} \src{@r0:r1}, \src{r2}, \imm{\#8}} (\textbf{SPIR-V:} \texttt{OpStore (uvec2)}; \textbf{GLSL:} \texttt{ssbo.pairs[i] = uvec2(a, b);})
  \item \texttt{staging\_register\_control\_32}: Controls the read/write staging register bank and count for async message instructions. \emph{Example:} {\scriptsize\ttfamily \mnem{LOAD.i128}\imod{.slot0} \dst{@r0:r3}, \src{r4}, \imm{\#0}} (\textbf{SPIR-V:} \texttt{OpLoad (vec4)}; \textbf{GLSL:} \texttt{vec4 v = ssbo.data[idx];})
  \item \texttt{opcode1}: Primary sub-opcode within the instruction class uniquely identifying the operation family. \emph{Example:} {\scriptsize\ttfamily \mnem{IADD.s32} \dst{r0}, \src{r1}, \src{r2}} (\textbf{SPIR-V:} \texttt{OpIAdd (operation selector)}; \textbf{GLSL:} \texttt{// Primary opcode bits})
  \item \texttt{instruction\_class}: 4-bit major instruction class selector (\texttt{bits 48:51}) partitioning the 64-bit Valhall opcode space (e.g., \texttt{FMA}, \texttt{CVT}, \texttt{SFU}, \texttt{LS}, \texttt{TEX}, \texttt{BRANCH}). \emph{Example:} {\scriptsize\ttfamily \mnem{FADD.f32} \dst{r0}, \src{r1}, \src{r2}} (\textbf{SPIR-V:} \texttt{N/A (encoding class)}; \textbf{GLSL:} \texttt{// Hardware execution pipe / decode class})
  \item \texttt{fau\_page\_index}: Selects the FAU (Fast Access Uniform) page (\texttt{0..3}) when a source operand references hardware small constants, thread IDs, or lane constants. \emph{Example:} {\scriptsize\ttfamily \mnem{FADD.f32} \dst{r0}, \src{r1}, \src{k1f}} (\textbf{SPIR-V:} \texttt{OpConstant / BuiltIn}; \textbf{GLSL:} \texttt{float y = x + 1.0;})
  \item \hyperref[enum:flow_control_m]{\texttt{flow}}: 4-bit embedded flow-control \& scoreboard wait modifier (\texttt{bits 59:62}: \texttt{none}, \texttt{wait0..wait0126}, \texttt{wait}, \texttt{reconverge}, \texttt{discard}, \texttt{end}). \emph{Example:} {\scriptsize\ttfamily \mnem{FADD.f32}\imod{.wait0.reconverge} \dst{r0}, \src{r1}, \src{r2}} (\textbf{SPIR-V:} \texttt{OpControlBarrier / OpReturn}; \textbf{GLSL:} \texttt{// Scoreboard wait + warp reconvergence})
\end{tightitemize}
\clearpage

\subsection*{VAR\_TEX\_BUF\_GRADIENT}
\addcontentsline{toc}{subsection}{VAR\_TEX\_BUF\_GRADIENT}
Fuses descriptor-less varying-buffer UV interpolation (at byte offset \texttt{src2}) with a texture gradient/LOD operation (\texttt{TEX\_GRADIENT} via handle \texttt{src0}), writing the result into \texttt{@sw} in a single message.
\par\textbf{Notes:}
\begin{tightitemize}
  \item Fragment-only (\texttt{encode\_v9.rs:3620-3650}).
  \item \textbf{Source references:} \texttt{kraid/ops.rs:4764-4817}, \texttt{encode\_v9.rs:3620-3650}, \texttt{bifrost/valhall/ISA.xml:159}.
\end{tightitemize}
\codebox{\mnem{VAR\_TEX\_BUF\_GRADIENT}\imod{\{.message\_slot\_index\}}\imod{\{.flow\}}\imod{\{.skip\}}\imod{\{.source\_format\}}\imod{\{.sample\_position\}}\imod{\{.update\}}\imod{\{.dimensionality\}}\imod{\{.register\_width\}}\imod{\{.write\_mask\}}\imod{\{.wide\_indices\}}\imod{\{.lod\_bias\_disable\}}\imod{\{.lod\_clamp\_disable\}} \dst{@sw}, \src{src0}, \src{src2}\\\mnem{VAR\_TEX\_BUF\_GRADIENT}\imod{\{.message\_slot\_index\}}\imod{\{.flow\}}\imod{\{.skip\}}\imod{\{.source\_format\}}\imod{\{.sample\_position\}}\imod{\{.update\}}\imod{\{.dimensionality\}}\imod{\{.projection\_enable\}}\imod{\{.register\_width\}}\imod{\{.write\_mask\}}\imod{\{.wide\_indices\}}\imod{\{.lod\_bias\_disable\}}\imod{\{.lod\_clamp\_disable\}} \dst{@sw}, \src{src0}, \src{src2}}
\textbf{Operands:}
\begin{tightitemize}\footnotesize
  \item \texttt{\dst{@sw}} (\texttt{staging\_register\_index} + \texttt{sr\_write\_count}): Destination staging register range receiving the gradient/LOD texture operation result.
  \item \texttt{\src{src0}}: Texture/sampler descriptor handle.
  \item \texttt{\src{src2}}: Descriptor-less varying buffer byte offset providing interpolated UV coordinates (encoded in \texttt{src2}).
\end{tightitemize}
\par\textbf{Closest SPIR-V Op:} \texttt{OpLoad (in) + OpImageSampleExplicitLod}
\par\textbf{Pseudo-GLSL (Operational Semantics):}
\codebox{@sw = evalTexGradOrLod(handle = src0, interpolateVaryingBuf(byte\_off = src2));}
\textbf{Execution Context:}
\begin{tightitemize}\footnotesize
  \item \textbf{Varying buffer \& barycentric cache (\hyperref[sec:csf-tiler]{\texttt{IDVSSR::VARY\_SIZE}}, \texttt{r48:r49}):} Reads per-vertex varying RAM written by Varying EBIN \texttt{STORE.estream}; descriptor forms (\texttt{LD\_VAR*}) also chase \hyperref[sec:csf-srt]{\texttt{SRT\_2}} $\to$ \texttt{Resource[table]} $\to$ \texttt{Attribute[index]} and reuse cached barycentric weights via \texttt{.store}/\texttt{.retrieve}.
  \item \textbf{Texture \& Sampler chain (\hyperref[sec:csf-srt]{\texttt{SRT\_2}} Frag / \texttt{SRT\_0} Comp/Vert):} 64B-aligned VA (\texttt{\& \textasciitilde{}0x3F}) $\to$ \texttt{Resource[table]} $\to$ \hyperref[sec:csf-tex-samp]{\texttt{Texture[tex\_idx]}} (32B, \texttt{type=2}) $\to$ \texttt{GenericPlane[mip]} (32B, \texttt{type=10}: \texttt{pointer}, \texttt{row\_stride}, Morton 16$\times$16 / AFBC) + \texttt{Sampler[samp\_idx]} (32B, \texttt{type=1}: filter, wrap, LOD clamp/bias, compare).
  \item \textbf{Fragment stage preloads \& FP modes (\hyperref[sec:csf-spd]{\texttt{SPD\_2}}, \texttt{r20:r21}):} Hardware pre-fills \texttt{r59} with \texttt{floatBitsToUint(FragCoord.z)} and \texttt{r60} with the MSAA sample coverage mask (\texttt{gl\_SampleMaskIn[0]}), while \texttt{ShaderProgram} word~0 latches 32/64 GPR allocation and warp-wide \texttt{flush\_to\_zero\_mode}, \texttt{suppress\_na\_n}, and \texttt{suppress\_inf}.
\end{tightitemize}
\par\textbf{Worked Examples:}
\par\textbf{Example 1 (Fused varying-buffer gradient query):}
\codebox{\mnem{VAR\_TEX\_BUF\_GRADIENT}\imod{.tex2d.slot0} \dst{@r0:r3}, \src{u0}, \src{r59}, \src{r4}}
Interpolates \texttt{vUV} from varying buffer and computes screen-space texture gradients / LOD in one fused message.
\par\textbf{Example 2 (Forced quad derivatives (`.force\_delta`)):}
\codebox{\mnem{VAR\_TEX\_BUF\_GRADIENT}\imod{.tex2d.force\_delta.slot0} \dst{@r0:r1}, \src{u0}, \src{r59}, \src{r4}}
Evaluates 2x2 quad derivatives across helper lanes for accurate mip selection.

\begin{center}\footnotesize
\begin{tabularx}{\linewidth}{@{}lllX@{}}
\toprule
Variant & Arch & Class & Opcode Fields \\
\midrule
\texttt{(none)} & v9--v14 & \texttt{S3A} & \texttt{opcode1=VAR\_TEX\_BUF\_GRADIENT} \\
\texttt{(none)} & v9--v14 & \texttt{S3A} & \texttt{opcode1=VAR\_TEX\_BUF\_GRADIENT} \\
\texttt{(none)} & v9--v14 & \texttt{S3A} & \texttt{opcode1=VAR\_TEX\_BUF\_GRADIENT} \\
\bottomrule
\end{tabularx}
\end{center}
\begin{center}\footnotesize
\begin{tabularx}{0.88\linewidth}{@{}l l X@{}}
\toprule
Field Name & Bits & Type / Fixed Value \\
\midrule
\texttt{src0} & \texttt{0:7} & \texttt{SourceEncoding64} \\
\hyperref[enum:tex_wide_indices_enable_m]{\texttt{wide\_indices}} & \texttt{8} & \texttt{tex\_wide\_indices\_enable\_m (def none)} \\
\textit{Reserved} & \texttt{9:10} & \texttt{zero} \\
\hyperref[enum:var_tex_buf_source_format_m]{\texttt{source\_format}} & \texttt{11} & \texttt{exact=src\_f32} \\
\textit{Reserved} & \texttt{12} & \texttt{zero} \\
\hyperref[enum:tex_lod_bias_disable_m]{\texttt{lod\_bias\_disable}} & \texttt{13} & \texttt{tex\_lod\_bias\_disable\_m} \\
\hyperref[enum:tex_lod_clamp_disable_m]{\texttt{lod\_clamp\_disable}} & \texttt{14} & \texttt{tex\_lod\_clamp\_disable\_m} \\
\hyperref[enum:tex_derivative_enable_m]{\texttt{derivative\_enable}} & \texttt{15} & \texttt{exact=none} \\
\texttt{src2} & \texttt{16:23} & \texttt{SourceEncoding} \\
\hyperref[enum:tex_gradient_width_m]{\texttt{register\_width}} & \texttt{24} & \texttt{tex\_gradient\_width\_m} \\
\texttt{secondary\_width} & \texttt{25} & \texttt{exact=0} \\
\textit{Reserved} & \texttt{26:27} & \texttt{ignore} \\
\hyperref[enum:var_tex_dimensionality_m]{\texttt{dimensionality}} & \texttt{28:29} & \texttt{var\_tex\_dimensionality\_m} \\
\hyperref[enum:message_slot_index_m]{\texttt{message\_slot\_index}} & \texttt{30:32} & \texttt{message\_slot\_index\_m} \\
\hyperref[enum:encoded_sample_and_update_m]{\texttt{encoded\_sample\_and\_update}} & \texttt{33:35} & \texttt{encoded\_sample\_and\_update\_m} \\
\texttt{sr\_write\_count} & \texttt{36:38} & \texttt{sr\_write\_count\_t} \\
\hyperref[enum:tex_gradient_skip_m]{\texttt{skip}} & \texttt{39} & \texttt{tex\_gradient\_skip\_m (def none)} \\
\texttt{staging\_register\_index} & \texttt{40:45} & \texttt{uint} \\
\texttt{staging\_register\_control\_32} & \texttt{46:47} & \texttt{exact=write} \\
\texttt{opcode1} & \texttt{48:51} & \texttt{exact=VAR\_TEX\_BUF\_GRADIENT} \\
\texttt{instruction\_class} & \texttt{52:56} & \texttt{exact=S3A} \\
\texttt{fau\_page\_index} & \texttt{57:58} & \texttt{uint} \\
\hyperref[enum:flow_control_m]{\texttt{flow}} & \texttt{59:62} & \texttt{flow\_control\_m} \\
\textit{Reserved} & \texttt{63} & \texttt{zero} \\
\bottomrule
\end{tabularx}
\end{center}
\textbf{Fields Explained:}
\begin{tightitemize}\footnotesize
  \item \texttt{src0}: Primary source operand (\texttt{r0..r63}, \texttt{u0..u63} uniform, \texttt{k*} FAU constant, or \texttt{t0..t2} temporary; \texttt{\textasciicircum{}} marks discard/last-use). \emph{Example:} {\scriptsize\ttfamily \mnem{FADD.f32} \dst{r0}, \src{r1}, \src{r2}} (\textbf{SPIR-V:} \texttt{OpFAdd (Operand 0)}; \textbf{GLSL:} \texttt{float c = a + b;})
  \item \hyperref[enum:tex_wide_indices_enable_m]{\texttt{wide\_indices}}: Uses 32-bit wide texture/sampler table indices instead of compact 16-bit packed indices. \emph{Example:} {\scriptsize\ttfamily \mnem{TEX\_SINGLE}\imod{.tex2d.wide} \dst{@r0:r3}, \src{@r4:r5}, \src{r6}} (\textbf{SPIR-V:} \texttt{OpImageSampleImplicitLod (bindless)}; \textbf{GLSL:} \texttt{vec4 c = texture(textures[nonuniformEXT(i)], uv);})
  \item \hyperref[enum:var_tex_buf_source_format_m]{\texttt{source\_format}}: Declares the storage precision of the varying buffer (\texttt{src\_f32}, \texttt{src\_f16}, \texttt{src\_flat32}, \texttt{src\_flat16}). \emph{Example:} {\scriptsize\ttfamily \mnem{LD\_VAR\_BUF.f32}\imod{.src\_f16.v2} \dst{@r0:r1}, \src{r2}} (\textbf{SPIR-V:} \texttt{OpLoad (RelaxedPrecision Input)}; \textbf{GLSL:} \texttt{mediump in vec2 v\_uv;})
  \item \hyperref[enum:tex_lod_bias_disable_m]{\texttt{lod\_bias\_disable}}: Disables sampler/shader LOD bias when computing explicit gradients (\texttt{TEX\_GRADIENT}). \emph{Example:} {\scriptsize\ttfamily \mnem{TEX\_GRADIENT}\imod{.tex2d.disable\_lod\_bias} \dst{@r0:r1}, \src{@r4:r5}} (\textbf{SPIR-V:} \texttt{OpImageQueryLod}; \textbf{GLSL:} \texttt{vec2 lod = textureQueryLod(tex, uv);})
  \item \hyperref[enum:tex_lod_clamp_disable_m]{\texttt{lod\_clamp\_disable}}: Disables sampler min/max LOD clamp bounds when querying raw mip level. \emph{Example:} {\scriptsize\ttfamily \mnem{TEX\_GRADIENT}\imod{.tex2d.disable\_lod\_clamp} \dst{@r0:r1}, \src{@r4:r5}} (\textbf{SPIR-V:} \texttt{OpImageQueryLod}; \textbf{GLSL:} \texttt{vec2 lod = textureQueryLod(tex, uv);})
  \item \hyperref[enum:tex_derivative_enable_m]{\texttt{derivative\_enable}}: Returns raw screen-space derivatives ($dP/dx, dP/dy$) instead of scalar LOD. \emph{Example:} {\scriptsize\ttfamily \mnem{TEX\_GRADIENT}\imod{.tex2d.derivative} \dst{@r0:r3}, \src{@r4:r5}} (\textbf{SPIR-V:} \texttt{OpDPdx / OpDPdy}; \textbf{GLSL:} \texttt{vec2 dx = dFdx(uv);})
  \item \texttt{src2}: Third source operand (FMA addend, \texttt{MUX}/\texttt{CSEL} selector, \texttt{IADDC} carry-in, or bitwise mask in fused shift-logic). \emph{Example:} {\scriptsize\ttfamily \mnem{FMA.f32} \dst{r0}, \src{r1}, \src{r2}, \src{r3}} (\textbf{SPIR-V:} \texttt{GLSL.std.450 Fma (Addend)}; \textbf{GLSL:} \texttt{float d = fma(a, b, c);})
  \item \hyperref[enum:tex_gradient_width_m]{\texttt{register\_width}}: Selects 16-bit vs 32-bit GPR element width for the corresponding message return window. \emph{Example:} {\scriptsize\ttfamily \mnem{TEX\_DUAL}\imod{.primary\_f16.secondary\_f32} \dst{@r0:r2}, \src{@r4:r5}} (\textbf{SPIR-V:} \texttt{OpImageSampleImplicitLod (RelaxedPrecision)}; \textbf{GLSL:} \texttt{mediump vec2 a = texture(t0, uv).rg;})
  \item \texttt{secondary\_width}: Selects 16-bit vs 32-bit GPR element width for the corresponding message return window. \emph{Example:} {\scriptsize\ttfamily \mnem{TEX\_DUAL}\imod{.primary\_f16.secondary\_f32} \dst{@r0:r2}, \src{@r4:r5}} (\textbf{SPIR-V:} \texttt{OpImageSampleImplicitLod (RelaxedPrecision)}; \textbf{GLSL:} \texttt{mediump vec2 a = texture(t0, uv).rg;})
  \item \hyperref[enum:var_tex_dimensionality_m]{\texttt{dimensionality}}: Texture dimensionality (\texttt{tex1d}, \texttt{tex2d}, \texttt{tex3d}, \texttt{cube}). \emph{Example:} {\scriptsize\ttfamily \mnem{TEX\_SINGLE}\imod{.tex3d.f32.rgba} \dst{@r0:r3}, \src{@r4:r6}} (\textbf{SPIR-V:} \texttt{OpImageSampleImplicitLod (Dim=3D)}; \textbf{GLSL:} \texttt{vec4 c = texture(volTex, uvw);})
  \item \hyperref[enum:message_slot_index_m]{\texttt{message\_slot\_index}}: Asynchronous scoreboard slot (\texttt{slot0}--\texttt{slot2}, \texttt{slot7}) tracking completion of this message. \emph{Example:} {\scriptsize\ttfamily \mnem{LOAD.i32}\imod{.slot0} \dst{@r0}, \src{r2}, \imm{\#0}} (\textbf{SPIR-V:} \texttt{OpLoad}; \textbf{GLSL:} \texttt{uint v = ssbo.arr[i];})
  \item \hyperref[enum:encoded_sample_and_update_m]{\texttt{encoded\_sample\_and\_update}}: Jointly encodes MSAA sample position and hardware varying-cache update mode (\texttt{center\_store}, \texttt{retrieve}, etc.). \emph{Example:} {\scriptsize\ttfamily \mnem{VAR\_TEX\_BUF\_SINGLE}\imod{.center\_store.tex2d} \dst{@r0:r3}, \src{r4}} (\textbf{SPIR-V:} \texttt{OpImageSampleImplicitLod (fused varying)}; \textbf{GLSL:} \texttt{vec4 c = texture(s, v\_uv);})
  \item \texttt{sr\_write\_count}: Number of consecutive 32-bit staging registers written back by a dual-staging message (e.g. ATOM\textbackslash\{\}\_RETURN / TEX). \emph{Example:} {\scriptsize\ttfamily \mnem{TEX\_SINGLE}\imod{.tex2d.f32.rgba} \dst{@r0:r3}, \src{@r4:r5}} (\textbf{SPIR-V:} \texttt{OpImageSampleImplicitLod}; \textbf{GLSL:} \texttt{vec4 c = texture(tex, uv);})
  \item \hyperref[enum:tex_gradient_skip_m]{\texttt{skip}}: Skips texture/memory fetch on helper/discarded fragment lanes while preserving 2x2 quad derivatives. \emph{Example:} {\scriptsize\ttfamily \mnem{TEX\_SINGLE}\imod{.tex2d.computed.skip} \dst{@r0:r3}, \src{@r4:r5}} (\textbf{SPIR-V:} \texttt{OpImageSampleImplicitLod (after OpDemoteToHelperInvocation)}; \textbf{GLSL:} \texttt{if (a \textless{} 0.0) discard; vec4 c = texture(s, uv);})
  \item \texttt{staging\_register\_index}: Base GPR index of the contiguous staging register window (\texttt{@rN:rM}) for message payload/return. \emph{Example:} {\scriptsize\ttfamily \mnem{STORE.i64}\imod{.slot7} \src{@r0:r1}, \src{r2}, \imm{\#8}} (\textbf{SPIR-V:} \texttt{OpStore (uvec2)}; \textbf{GLSL:} \texttt{ssbo.pairs[i] = uvec2(a, b);})
  \item \texttt{staging\_register\_control\_32}: Controls the read/write staging register bank and count for async message instructions. \emph{Example:} {\scriptsize\ttfamily \mnem{LOAD.i128}\imod{.slot0} \dst{@r0:r3}, \src{r4}, \imm{\#0}} (\textbf{SPIR-V:} \texttt{OpLoad (vec4)}; \textbf{GLSL:} \texttt{vec4 v = ssbo.data[idx];})
  \item \texttt{opcode1}: Primary sub-opcode within the instruction class uniquely identifying the operation family. \emph{Example:} {\scriptsize\ttfamily \mnem{IADD.s32} \dst{r0}, \src{r1}, \src{r2}} (\textbf{SPIR-V:} \texttt{OpIAdd (operation selector)}; \textbf{GLSL:} \texttt{// Primary opcode bits})
  \item \texttt{instruction\_class}: 4-bit major instruction class selector (\texttt{bits 48:51}) partitioning the 64-bit Valhall opcode space (e.g., \texttt{FMA}, \texttt{CVT}, \texttt{SFU}, \texttt{LS}, \texttt{TEX}, \texttt{BRANCH}). \emph{Example:} {\scriptsize\ttfamily \mnem{FADD.f32} \dst{r0}, \src{r1}, \src{r2}} (\textbf{SPIR-V:} \texttt{N/A (encoding class)}; \textbf{GLSL:} \texttt{// Hardware execution pipe / decode class})
  \item \texttt{fau\_page\_index}: Selects the FAU (Fast Access Uniform) page (\texttt{0..3}) when a source operand references hardware small constants, thread IDs, or lane constants. \emph{Example:} {\scriptsize\ttfamily \mnem{FADD.f32} \dst{r0}, \src{r1}, \src{k1f}} (\textbf{SPIR-V:} \texttt{OpConstant / BuiltIn}; \textbf{GLSL:} \texttt{float y = x + 1.0;})
  \item \hyperref[enum:flow_control_m]{\texttt{flow}}: 4-bit embedded flow-control \& scoreboard wait modifier (\texttt{bits 59:62}: \texttt{none}, \texttt{wait0..wait0126}, \texttt{wait}, \texttt{reconverge}, \texttt{discard}, \texttt{end}). \emph{Example:} {\scriptsize\ttfamily \mnem{FADD.f32}\imod{.wait0.reconverge} \dst{r0}, \src{r1}, \src{r2}} (\textbf{SPIR-V:} \texttt{OpControlBarrier / OpReturn}; \textbf{GLSL:} \texttt{// Scoreboard wait + warp reconvergence})
\end{tightitemize}
\clearpage

\subsection*{VAR\_TEX\_BUF\_SINGLE}
\addcontentsline{toc}{subsection}{VAR\_TEX\_BUF\_SINGLE}
Fuses descriptor-less varying-buffer UV interpolation (at byte offset \texttt{src2}, \texttt{source\_format} f32/f16 on v10+) with a filtered texture sample (\texttt{TEX\_SINGLE} via handle \texttt{src0}), writing the sampled texel into \texttt{@sw} in a single message.
\par\textbf{Notes:}
\begin{tightitemize}
  \item Fragment-only (\texttt{encode\_v9.rs:3519-3533, 3652-3683}).
  \item \textbf{Source references:} \texttt{kraid/ops.rs:4819-4870}, \texttt{encode\_v9.rs:3652-3683}, \texttt{bifrost/valhall/ISA.xml:161}.
\end{tightitemize}
\codebox{\mnem{VAR\_TEX\_BUF\_SINGLE}\imod{\{.message\_slot\_index\}}\imod{\{.flow\}}\imod{\{.skip\}}\imod{\{.source\_format\}}\imod{\{.sample\_position\}}\imod{\{.update\}}\imod{\{.dimensionality\}}\imod{\{.register\_width\}}\imod{\{.write\_mask\}}\imod{\{.wide\_indices\}}\imod{\{.array\_enable\}}\imod{\{.compare\_enable\}}\imod{\{.lod\_mode\}} \dst{@sw}, \src{src0}, \src{src2}\\\mnem{VAR\_TEX\_BUF\_SINGLE}\imod{\{.message\_slot\_index\}}\imod{\{.flow\}}\imod{\{.skip\}}\imod{\{.source\_format\}}\imod{\{.sample\_position\}}\imod{\{.update\}}\imod{\{.dimensionality\}}\imod{\{.projection\_enable\}}\imod{\{.register\_width\}}\imod{\{.write\_mask\}}\imod{\{.wide\_indices\}}\imod{\{.array\_enable\}}\imod{\{.compare\_enable\}}\imod{\{.lod\_mode\}} \dst{@sw}, \src{src0}, \src{src2}}
\textbf{Operands:}
\begin{tightitemize}\footnotesize
  \item \texttt{\dst{@sw}} (\texttt{staging\_register\_index} + \texttt{sr\_write\_count}): Destination staging register range receiving the single texture sample result.
  \item \texttt{\src{src0}}: Texture/sampler descriptor handle.
  \item \texttt{\src{src2}}: Descriptor-less varying buffer byte offset providing interpolated UV coordinates (encoded in \texttt{src2}).
\end{tightitemize}
\par\textbf{Closest SPIR-V Op:} \texttt{OpLoad (in) + OpImageSampleImplicitLod}
\par\textbf{Pseudo-GLSL (Operational Semantics):}
\codebox{@sw = textureSample(handle = src0, interpolateVaryingBuf(byte\_off = src2), dim, flag\_project, flag\_compare);}
\textbf{Execution Context:}
\begin{tightitemize}\footnotesize
  \item \textbf{Varying buffer \& barycentric cache (\hyperref[sec:csf-tiler]{\texttt{IDVSSR::VARY\_SIZE}}, \texttt{r48:r49}):} Reads per-vertex varying RAM written by Varying EBIN \texttt{STORE.estream}; descriptor forms (\texttt{LD\_VAR*}) also chase \hyperref[sec:csf-srt]{\texttt{SRT\_2}} $\to$ \texttt{Resource[table]} $\to$ \texttt{Attribute[index]} and reuse cached barycentric weights via \texttt{.store}/\texttt{.retrieve}.
  \item \textbf{Texture \& Sampler chain (\hyperref[sec:csf-srt]{\texttt{SRT\_2}} Frag / \texttt{SRT\_0} Comp/Vert):} 64B-aligned VA (\texttt{\& \textasciitilde{}0x3F}) $\to$ \texttt{Resource[table]} $\to$ \hyperref[sec:csf-tex-samp]{\texttt{Texture[tex\_idx]}} (32B, \texttt{type=2}) $\to$ \texttt{GenericPlane[mip]} (32B, \texttt{type=10}: \texttt{pointer}, \texttt{row\_stride}, Morton 16$\times$16 / AFBC) + \texttt{Sampler[samp\_idx]} (32B, \texttt{type=1}: filter, wrap, LOD clamp/bias, compare).
  \item \textbf{Fragment stage preloads \& FP modes (\hyperref[sec:csf-spd]{\texttt{SPD\_2}}, \texttt{r20:r21}):} Hardware pre-fills \texttt{r59} with \texttt{floatBitsToUint(FragCoord.z)} and \texttt{r60} with the MSAA sample coverage mask (\texttt{gl\_SampleMaskIn[0]}), while \texttt{ShaderProgram} word~0 latches 32/64 GPR allocation and warp-wide \texttt{flush\_to\_zero\_mode}, \texttt{suppress\_na\_n}, and \texttt{suppress\_inf}.
\end{tightitemize}
\par\textbf{Worked Examples:}
\par\textbf{Example 1 (Fused varying-buffer 2D texture sample):}
\codebox{\mnem{VAR\_TEX\_BUF\_SINGLE}\imod{.tex2d.computed.slot0} \dst{@r0:r3}, \src{u0}, \src{r59}, \src{r4}}
Interpolates \texttt{vUV} from varying buffer offset \texttt{r4} directly into \texttt{TEX\_SINGLE}, saving an entire \texttt{LD\_VAR\_BUF} + scoreboard wait.
\par\textbf{Example 2 (Fused projective texture sample (`.project`)):}
\codebox{\mnem{VAR\_TEX\_BUF\_SINGLE}\imod{.tex2d.project.computed.slot0} \dst{@r0:r3}, \src{u0}, \src{r59}, \src{r4}}
Interpolates homogeneous \texttt{(u, v, q)} and performs hardware \texttt{textureProj(s, vUVQ)} in one instruction.

\begin{center}\footnotesize
\begin{tabularx}{\linewidth}{@{}lllX@{}}
\toprule
Variant & Arch & Class & Opcode Fields \\
\midrule
\texttt{(none)} & v9--v14 & \texttt{S3A} & \texttt{opcode1=VAR\_TEX\_BUF\_SINGLE} \\
\texttt{(none)} & v9--v14 & \texttt{S3A} & \texttt{opcode1=VAR\_TEX\_BUF\_SINGLE} \\
\texttt{(none)} & v9--v14 & \texttt{S3A} & \texttt{opcode1=VAR\_TEX\_BUF\_SINGLE} \\
\bottomrule
\end{tabularx}
\end{center}
\begin{center}\footnotesize
\begin{tabularx}{0.88\linewidth}{@{}l l X@{}}
\toprule
Field Name & Bits & Type / Fixed Value \\
\midrule
\texttt{src0} & \texttt{0:7} & \texttt{SourceEncoding64} \\
\hyperref[enum:tex_wide_indices_enable_m]{\texttt{wide\_indices}} & \texttt{8} & \texttt{tex\_wide\_indices\_enable\_m (def none)} \\
\textit{Reserved} & \texttt{9} & \texttt{zero} \\
\hyperref[enum:var_tex_array_enable_m]{\texttt{array\_enable}} & \texttt{10} & \texttt{var\_tex\_array\_enable\_m (def none)} \\
\hyperref[enum:var_tex_buf_source_format_m]{\texttt{source\_format}} & \texttt{11} & \texttt{exact=src\_f32} \\
\hyperref[enum:var_tex_compare_enable_m]{\texttt{compare\_enable}} & \texttt{12} & \texttt{var\_tex\_compare\_enable\_m (def none)} \\
\hyperref[enum:var_tex_lod_mode_m]{\texttt{lod\_mode}} & \texttt{13:15} & \texttt{var\_tex\_lod\_mode\_m} \\
\texttt{src2} & \texttt{16:23} & \texttt{SourceEncoding} \\
\hyperref[enum:tex_width_m]{\texttt{register\_width}} & \texttt{24} & \texttt{tex\_width\_m} \\
\texttt{secondary\_width} & \texttt{25} & \texttt{exact=0} \\
\textit{Reserved} & \texttt{26:27} & \texttt{ignore} \\
\hyperref[enum:var_tex_dimensionality_m]{\texttt{dimensionality}} & \texttt{28:29} & \texttt{var\_tex\_dimensionality\_m} \\
\hyperref[enum:message_slot_index_m]{\texttt{message\_slot\_index}} & \texttt{30:32} & \texttt{message\_slot\_index\_m} \\
\hyperref[enum:encoded_sample_and_update_m]{\texttt{encoded\_sample\_and\_update}} & \texttt{33:35} & \texttt{encoded\_sample\_and\_update\_m} \\
\texttt{sr\_write\_count} & \texttt{36:38} & \texttt{sr\_write\_count\_t} \\
\hyperref[enum:tex_skip_m]{\texttt{skip}} & \texttt{39} & \texttt{tex\_skip\_m (def none)} \\
\texttt{staging\_register\_index} & \texttt{40:45} & \texttt{uint} \\
\texttt{staging\_register\_control\_32} & \texttt{46:47} & \texttt{exact=write} \\
\texttt{opcode1} & \texttt{48:51} & \texttt{exact=VAR\_TEX\_BUF\_SINGLE} \\
\texttt{instruction\_class} & \texttt{52:56} & \texttt{exact=S3A} \\
\texttt{fau\_page\_index} & \texttt{57:58} & \texttt{uint} \\
\hyperref[enum:flow_control_m]{\texttt{flow}} & \texttt{59:62} & \texttt{flow\_control\_m} \\
\textit{Reserved} & \texttt{63} & \texttt{zero} \\
\bottomrule
\end{tabularx}
\end{center}
\textbf{Fields Explained:}
\begin{tightitemize}\footnotesize
  \item \texttt{src0}: Primary source operand (\texttt{r0..r63}, \texttt{u0..u63} uniform, \texttt{k*} FAU constant, or \texttt{t0..t2} temporary; \texttt{\textasciicircum{}} marks discard/last-use). \emph{Example:} {\scriptsize\ttfamily \mnem{FADD.f32} \dst{r0}, \src{r1}, \src{r2}} (\textbf{SPIR-V:} \texttt{OpFAdd (Operand 0)}; \textbf{GLSL:} \texttt{float c = a + b;})
  \item \hyperref[enum:tex_wide_indices_enable_m]{\texttt{wide\_indices}}: Uses 32-bit wide texture/sampler table indices instead of compact 16-bit packed indices. \emph{Example:} {\scriptsize\ttfamily \mnem{TEX\_SINGLE}\imod{.tex2d.wide} \dst{@r0:r3}, \src{@r4:r5}, \src{r6}} (\textbf{SPIR-V:} \texttt{OpImageSampleImplicitLod (bindless)}; \textbf{GLSL:} \texttt{vec4 c = texture(textures[nonuniformEXT(i)], uv);})
  \item \hyperref[enum:var_tex_array_enable_m]{\texttt{array\_enable}}: Enables array-layer coordinate in texture staging registers (\texttt{sampler2DArray} / \texttt{samplerCubeArray}). \emph{Example:} {\scriptsize\ttfamily \mnem{TEX\_SINGLE}\imod{.tex2d.array} \dst{@r0:r3}, \src{@r4:r6}} (\textbf{SPIR-V:} \texttt{OpImageSampleImplicitLod (Arrayed)}; \textbf{GLSL:} \texttt{vec4 c = texture(texArr, vec3(uv, layer));})
  \item \hyperref[enum:var_tex_buf_source_format_m]{\texttt{source\_format}}: Declares the storage precision of the varying buffer (\texttt{src\_f32}, \texttt{src\_f16}, \texttt{src\_flat32}, \texttt{src\_flat16}). \emph{Example:} {\scriptsize\ttfamily \mnem{LD\_VAR\_BUF.f32}\imod{.src\_f16.v2} \dst{@r0:r1}, \src{r2}} (\textbf{SPIR-V:} \texttt{OpLoad (RelaxedPrecision Input)}; \textbf{GLSL:} \texttt{mediump in vec2 v\_uv;})
  \item \hyperref[enum:var_tex_compare_enable_m]{\texttt{compare\_enable}}: Enables depth-reference shadow comparison (\texttt{Dref} in staging registers). \emph{Example:} {\scriptsize\ttfamily \mnem{TEX\_SINGLE}\imod{.tex2d.compare} \dst{@r0}, \src{@r4:r6}} (\textbf{SPIR-V:} \texttt{OpImageSampleDrefImplicitLod}; \textbf{GLSL:} \texttt{float s = texture(shadowMap, vec3(uv, refZ));})
  \item \hyperref[enum:var_tex_lod_mode_m]{\texttt{lod\_mode}}: Selects mipmap LOD calculation mode (\texttt{none}, \texttt{computed}, \texttt{explicit}, \texttt{computed\_bias}, \texttt{grdesc}). \emph{Example:} {\scriptsize\ttfamily \mnem{TEX\_SINGLE}\imod{.tex2d.explicit} \dst{@r0:r3}, \src{@r4:r6}} (\textbf{SPIR-V:} \texttt{OpImageSampleExplicitLod}; \textbf{GLSL:} \texttt{vec4 c = textureLod(tex, uv, 2.0);})
  \item \texttt{src2}: Third source operand (FMA addend, \texttt{MUX}/\texttt{CSEL} selector, \texttt{IADDC} carry-in, or bitwise mask in fused shift-logic). \emph{Example:} {\scriptsize\ttfamily \mnem{FMA.f32} \dst{r0}, \src{r1}, \src{r2}, \src{r3}} (\textbf{SPIR-V:} \texttt{GLSL.std.450 Fma (Addend)}; \textbf{GLSL:} \texttt{float d = fma(a, b, c);})
  \item \hyperref[enum:tex_width_m]{\texttt{register\_width}}: Selects 16-bit vs 32-bit GPR element width for the corresponding message return window. \emph{Example:} {\scriptsize\ttfamily \mnem{TEX\_DUAL}\imod{.primary\_f16.secondary\_f32} \dst{@r0:r2}, \src{@r4:r5}} (\textbf{SPIR-V:} \texttt{OpImageSampleImplicitLod (RelaxedPrecision)}; \textbf{GLSL:} \texttt{mediump vec2 a = texture(t0, uv).rg;})
  \item \texttt{secondary\_width}: Selects 16-bit vs 32-bit GPR element width for the corresponding message return window. \emph{Example:} {\scriptsize\ttfamily \mnem{TEX\_DUAL}\imod{.primary\_f16.secondary\_f32} \dst{@r0:r2}, \src{@r4:r5}} (\textbf{SPIR-V:} \texttt{OpImageSampleImplicitLod (RelaxedPrecision)}; \textbf{GLSL:} \texttt{mediump vec2 a = texture(t0, uv).rg;})
  \item \hyperref[enum:var_tex_dimensionality_m]{\texttt{dimensionality}}: Texture dimensionality (\texttt{tex1d}, \texttt{tex2d}, \texttt{tex3d}, \texttt{cube}). \emph{Example:} {\scriptsize\ttfamily \mnem{TEX\_SINGLE}\imod{.tex3d.f32.rgba} \dst{@r0:r3}, \src{@r4:r6}} (\textbf{SPIR-V:} \texttt{OpImageSampleImplicitLod (Dim=3D)}; \textbf{GLSL:} \texttt{vec4 c = texture(volTex, uvw);})
  \item \hyperref[enum:message_slot_index_m]{\texttt{message\_slot\_index}}: Asynchronous scoreboard slot (\texttt{slot0}--\texttt{slot2}, \texttt{slot7}) tracking completion of this message. \emph{Example:} {\scriptsize\ttfamily \mnem{LOAD.i32}\imod{.slot0} \dst{@r0}, \src{r2}, \imm{\#0}} (\textbf{SPIR-V:} \texttt{OpLoad}; \textbf{GLSL:} \texttt{uint v = ssbo.arr[i];})
  \item \hyperref[enum:encoded_sample_and_update_m]{\texttt{encoded\_sample\_and\_update}}: Jointly encodes MSAA sample position and hardware varying-cache update mode (\texttt{center\_store}, \texttt{retrieve}, etc.). \emph{Example:} {\scriptsize\ttfamily \mnem{VAR\_TEX\_BUF\_SINGLE}\imod{.center\_store.tex2d} \dst{@r0:r3}, \src{r4}} (\textbf{SPIR-V:} \texttt{OpImageSampleImplicitLod (fused varying)}; \textbf{GLSL:} \texttt{vec4 c = texture(s, v\_uv);})
  \item \texttt{sr\_write\_count}: Number of consecutive 32-bit staging registers written back by a dual-staging message (e.g. ATOM\textbackslash\{\}\_RETURN / TEX). \emph{Example:} {\scriptsize\ttfamily \mnem{TEX\_SINGLE}\imod{.tex2d.f32.rgba} \dst{@r0:r3}, \src{@r4:r5}} (\textbf{SPIR-V:} \texttt{OpImageSampleImplicitLod}; \textbf{GLSL:} \texttt{vec4 c = texture(tex, uv);})
  \item \hyperref[enum:tex_skip_m]{\texttt{skip}}: Skips texture/memory fetch on helper/discarded fragment lanes while preserving 2x2 quad derivatives. \emph{Example:} {\scriptsize\ttfamily \mnem{TEX\_SINGLE}\imod{.tex2d.computed.skip} \dst{@r0:r3}, \src{@r4:r5}} (\textbf{SPIR-V:} \texttt{OpImageSampleImplicitLod (after OpDemoteToHelperInvocation)}; \textbf{GLSL:} \texttt{if (a \textless{} 0.0) discard; vec4 c = texture(s, uv);})
  \item \texttt{staging\_register\_index}: Base GPR index of the contiguous staging register window (\texttt{@rN:rM}) for message payload/return. \emph{Example:} {\scriptsize\ttfamily \mnem{STORE.i64}\imod{.slot7} \src{@r0:r1}, \src{r2}, \imm{\#8}} (\textbf{SPIR-V:} \texttt{OpStore (uvec2)}; \textbf{GLSL:} \texttt{ssbo.pairs[i] = uvec2(a, b);})
  \item \texttt{staging\_register\_control\_32}: Controls the read/write staging register bank and count for async message instructions. \emph{Example:} {\scriptsize\ttfamily \mnem{LOAD.i128}\imod{.slot0} \dst{@r0:r3}, \src{r4}, \imm{\#0}} (\textbf{SPIR-V:} \texttt{OpLoad (vec4)}; \textbf{GLSL:} \texttt{vec4 v = ssbo.data[idx];})
  \item \texttt{opcode1}: Primary sub-opcode within the instruction class uniquely identifying the operation family. \emph{Example:} {\scriptsize\ttfamily \mnem{IADD.s32} \dst{r0}, \src{r1}, \src{r2}} (\textbf{SPIR-V:} \texttt{OpIAdd (operation selector)}; \textbf{GLSL:} \texttt{// Primary opcode bits})
  \item \texttt{instruction\_class}: 4-bit major instruction class selector (\texttt{bits 48:51}) partitioning the 64-bit Valhall opcode space (e.g., \texttt{FMA}, \texttt{CVT}, \texttt{SFU}, \texttt{LS}, \texttt{TEX}, \texttt{BRANCH}). \emph{Example:} {\scriptsize\ttfamily \mnem{FADD.f32} \dst{r0}, \src{r1}, \src{r2}} (\textbf{SPIR-V:} \texttt{N/A (encoding class)}; \textbf{GLSL:} \texttt{// Hardware execution pipe / decode class})
  \item \texttt{fau\_page\_index}: Selects the FAU (Fast Access Uniform) page (\texttt{0..3}) when a source operand references hardware small constants, thread IDs, or lane constants. \emph{Example:} {\scriptsize\ttfamily \mnem{FADD.f32} \dst{r0}, \src{r1}, \src{k1f}} (\textbf{SPIR-V:} \texttt{OpConstant / BuiltIn}; \textbf{GLSL:} \texttt{float y = x + 1.0;})
  \item \hyperref[enum:flow_control_m]{\texttt{flow}}: 4-bit embedded flow-control \& scoreboard wait modifier (\texttt{bits 59:62}: \texttt{none}, \texttt{wait0..wait0126}, \texttt{wait}, \texttt{reconverge}, \texttt{discard}, \texttt{end}). \emph{Example:} {\scriptsize\ttfamily \mnem{FADD.f32}\imod{.wait0.reconverge} \dst{r0}, \src{r1}, \src{r2}} (\textbf{SPIR-V:} \texttt{OpControlBarrier / OpReturn}; \textbf{GLSL:} \texttt{// Scoreboard wait + warp reconvergence})
\end{tightitemize}
\clearpage

\subsection*{VAR\_TEX\_DUAL}
\addcontentsline{toc}{subsection}{VAR\_TEX\_DUAL}
Fuses descriptor-based varying UV interpolation (via varying handle \texttt{src2}) with a dual texture sample (\texttt{TEX\_DUAL} via texture handle \texttt{src0}), writing both sample results into \texttt{@sw} in a single message.
\par\textbf{Notes:}
\begin{tightitemize}
  \item Fragment-only.
  \item FP32 varyings (\texttt{ISA.xml:163}).
  \item \textbf{Source references:} \texttt{kraid/ops.rs:4872-4929}, \texttt{encode\_v9.rs:3685-3734}, \texttt{bifrost/valhall/ISA.xml:163}.
\end{tightitemize}
\codebox{\mnem{VAR\_TEX\_DUAL}\imod{\{.message\_slot\_index\}}\imod{\{.flow\}}\imod{\{.skip\}}\imod{\{.sample\_position\}}\imod{\{.update\}}\imod{\{.dimensionality\}}\imod{\{.primary\_width\}}\imod{\{.primary\_write\_mask\}}\imod{\{.secondary\_width\}}\imod{\{.secondary\_write\_mask\}}\imod{\{.array\_enable\}}\imod{\{.compare\_enable\}}\imod{\{.lod\_mode\}} \dst{@sw}, \src{src0}, \src{src2}\\\mnem{VAR\_TEX\_DUAL}\imod{\{.message\_slot\_index\}}\imod{\{.flow\}}\imod{\{.skip\}}\imod{\{.sample\_position\}}\imod{\{.update\}}\imod{\{.dimensionality\}}\imod{\{.projection\_enable\}}\imod{\{.primary\_width\}}\imod{\{.primary\_write\_mask\}}\imod{\{.secondary\_width\}}\imod{\{.secondary\_write\_mask\}}\imod{\{.array\_enable\}}\imod{\{.compare\_enable\}}\imod{\{.lod\_mode\}} \dst{@sw}, \src{src0}, \src{src2}}
\textbf{Operands:}
\begin{tightitemize}\footnotesize
  \item \texttt{\dst{@sw}} (\texttt{staging\_register\_index} + \texttt{sr\_write\_count}): Destination staging register range receiving the dual texture sample result.
  \item \texttt{\src{src0}}: Texture/sampler descriptor handle.
  \item \texttt{\src{src2}}: Varying descriptor handle providing interpolated UV coordinates (encoded in \texttt{src2}).
\end{tightitemize}
\par\textbf{Closest SPIR-V Op:} \texttt{OpLoad (in) + 2x OpImageSample*}
\par\textbf{Pseudo-GLSL (Operational Semantics):}
\codebox{vec2 uv = interpolateVarying(handle = src2);\\@sw[0:3] = textureSample(tex0\_of(src0), uv); @sw[4:7] = textureSample(tex1\_of(src0), uv);}
\textbf{Execution Context:}
\begin{tightitemize}\footnotesize
  \item \textbf{Varying buffer \& barycentric cache (\hyperref[sec:csf-tiler]{\texttt{IDVSSR::VARY\_SIZE}}, \texttt{r48:r49}):} Reads per-vertex varying RAM written by Varying EBIN \texttt{STORE.estream}; descriptor forms (\texttt{LD\_VAR*}) also chase \hyperref[sec:csf-srt]{\texttt{SRT\_2}} $\to$ \texttt{Resource[table]} $\to$ \texttt{Attribute[index]} and reuse cached barycentric weights via \texttt{.store}/\texttt{.retrieve}.
  \item \textbf{Texture \& Sampler chain (\hyperref[sec:csf-srt]{\texttt{SRT\_2}} Frag / \texttt{SRT\_0} Comp/Vert):} 64B-aligned VA (\texttt{\& \textasciitilde{}0x3F}) $\to$ \texttt{Resource[table]} $\to$ \hyperref[sec:csf-tex-samp]{\texttt{Texture[tex\_idx]}} (32B, \texttt{type=2}) $\to$ \texttt{GenericPlane[mip]} (32B, \texttt{type=10}: \texttt{pointer}, \texttt{row\_stride}, Morton 16$\times$16 / AFBC) + \texttt{Sampler[samp\_idx]} (32B, \texttt{type=1}: filter, wrap, LOD clamp/bias, compare).
  \item \textbf{Fragment stage preloads \& FP modes (\hyperref[sec:csf-spd]{\texttt{SPD\_2}}, \texttt{r20:r21}):} Hardware pre-fills \texttt{r59} with \texttt{floatBitsToUint(FragCoord.z)} and \texttt{r60} with the MSAA sample coverage mask (\texttt{gl\_SampleMaskIn[0]}), while \texttt{ShaderProgram} word~0 latches 32/64 GPR allocation and warp-wide \texttt{flush\_to\_zero\_mode}, \texttt{suppress\_na\_n}, and \texttt{suppress\_inf}.
\end{tightitemize}
\par\textbf{Worked Examples:}
\par\textbf{Example 1 (Descriptor varying + dual 2D texture sample):}
\codebox{\mnem{VAR\_TEX\_DUAL}\imod{.tex2d.center.slot0} \dst{@r0:r7}, \src{u0}, \src{r59}, \src{u1}}
Interpolates descriptor-based varying \texttt{u1} at pixel center and issues two paired samples via \texttt{u0}.
\par\textbf{Example 2 (Centroid-interpolated dual sample):}
\codebox{\mnem{VAR\_TEX\_DUAL}\imod{.tex2d.centroid.slot0} \dst{@r0:r7}, \src{u0}, \src{r59}, \src{u1}}
Uses MSAA centroid interpolation for \texttt{vUV} to prevent out-of-primitive UV extrapolation on both textures.

\begin{center}\footnotesize
\begin{tabularx}{\linewidth}{@{}lllX@{}}
\toprule
Variant & Arch & Class & Opcode Fields \\
\midrule
\texttt{(none)} & v9--v14 & \texttt{S3A} & \texttt{opcode1=VAR\_TEX\_DUAL} \\
\texttt{(none)} & v9--v14 & \texttt{S3A} & \texttt{opcode1=VAR\_TEX\_DUAL} \\
\bottomrule
\end{tabularx}
\end{center}
\begin{center}\footnotesize
\begin{tabularx}{0.88\linewidth}{@{}l l X@{}}
\toprule
Field Name & Bits & Type / Fixed Value \\
\midrule
\texttt{src0} & \texttt{0:7} & \texttt{SourceEncoding64} \\
\hyperref[enum:tex_wide_indices_enable_m]{\texttt{wide\_indices}} & \texttt{8} & \texttt{exact=none} \\
\textit{Reserved} & \texttt{9} & \texttt{zero} \\
\hyperref[enum:var_tex_dual_array_enable_m]{\texttt{array\_enable}} & \texttt{10} & \texttt{var\_tex\_dual\_array\_enable\_m (def none)} \\
\hyperref[enum:tex_texel_offset_m]{\texttt{texel\_offset}} & \texttt{11} & \texttt{exact=none} \\
\hyperref[enum:var_tex_compare_enable_m]{\texttt{compare\_enable}} & \texttt{12} & \texttt{var\_tex\_compare\_enable\_m (def none)} \\
\hyperref[enum:var_tex_dual_lod_mode_m]{\texttt{lod\_mode}} & \texttt{13:15} & \texttt{var\_tex\_dual\_lod\_mode\_m} \\
\texttt{src2} & \texttt{16:23} & \texttt{SourceEncoding} \\
\hyperref[enum:tex_dual_primary_width_m]{\texttt{primary\_width}} & \texttt{24} & \texttt{tex\_dual\_primary\_width\_m} \\
\hyperref[enum:tex_dual_secondary_width_m]{\texttt{secondary\_width}} & \texttt{25} & \texttt{tex\_dual\_secondary\_width\_m} \\
\textit{Reserved} & \texttt{26:27} & \texttt{ignore} \\
\hyperref[enum:var_tex_dimensionality_m]{\texttt{dimensionality}} & \texttt{28:29} & \texttt{var\_tex\_dimensionality\_m} \\
\hyperref[enum:message_slot_index_m]{\texttt{message\_slot\_index}} & \texttt{30:32} & \texttt{message\_slot\_index\_m} \\
\hyperref[enum:encoded_sample_and_update_m]{\texttt{encoded\_sample\_and\_update}} & \texttt{33:35} & \texttt{encoded\_sample\_and\_update\_m} \\
\texttt{sr\_write\_count} & \texttt{36:38} & \texttt{sr\_write\_count\_t} \\
\hyperref[enum:tex_skip_m]{\texttt{skip}} & \texttt{39} & \texttt{tex\_skip\_m (def none)} \\
\texttt{staging\_register\_index} & \texttt{40:45} & \texttt{uint} \\
\texttt{staging\_register\_control\_32} & \texttt{46:47} & \texttt{exact=write} \\
\texttt{opcode1} & \texttt{48:51} & \texttt{exact=VAR\_TEX\_DUAL} \\
\texttt{instruction\_class} & \texttt{52:56} & \texttt{exact=S3A} \\
\texttt{fau\_page\_index} & \texttt{57:58} & \texttt{uint} \\
\hyperref[enum:flow_control_m]{\texttt{flow}} & \texttt{59:62} & \texttt{flow\_control\_m} \\
\textit{Reserved} & \texttt{63} & \texttt{zero} \\
\bottomrule
\end{tabularx}
\end{center}
\textbf{Fields Explained:}
\begin{tightitemize}\footnotesize
  \item \texttt{src0}: Primary source operand (\texttt{r0..r63}, \texttt{u0..u63} uniform, \texttt{k*} FAU constant, or \texttt{t0..t2} temporary; \texttt{\textasciicircum{}} marks discard/last-use). \emph{Example:} {\scriptsize\ttfamily \mnem{FADD.f32} \dst{r0}, \src{r1}, \src{r2}} (\textbf{SPIR-V:} \texttt{OpFAdd (Operand 0)}; \textbf{GLSL:} \texttt{float c = a + b;})
  \item \hyperref[enum:tex_wide_indices_enable_m]{\texttt{wide\_indices}}: Uses 32-bit wide texture/sampler table indices instead of compact 16-bit packed indices. \emph{Example:} {\scriptsize\ttfamily \mnem{TEX\_SINGLE}\imod{.tex2d.wide} \dst{@r0:r3}, \src{@r4:r5}, \src{r6}} (\textbf{SPIR-V:} \texttt{OpImageSampleImplicitLod (bindless)}; \textbf{GLSL:} \texttt{vec4 c = texture(textures[nonuniformEXT(i)], uv);})
  \item \hyperref[enum:var_tex_dual_array_enable_m]{\texttt{array\_enable}}: Enables array-layer coordinate in texture staging registers (\texttt{sampler2DArray} / \texttt{samplerCubeArray}). \emph{Example:} {\scriptsize\ttfamily \mnem{TEX\_SINGLE}\imod{.tex2d.array} \dst{@r0:r3}, \src{@r4:r6}} (\textbf{SPIR-V:} \texttt{OpImageSampleImplicitLod (Arrayed)}; \textbf{GLSL:} \texttt{vec4 c = texture(texArr, vec3(uv, layer));})
  \item \hyperref[enum:tex_texel_offset_m]{\texttt{texel\_offset}}: Applies a signed 4-bit per-axis texel offset $(dx, dy, dz)$ packed in staging registers. \emph{Example:} {\scriptsize\ttfamily \mnem{TEX\_SINGLE}\imod{.tex2d.texel\_offset} \dst{@r0:r3}, \src{@r4:r6}} (\textbf{SPIR-V:} \texttt{OpImageSampleImplicitLod (ConstOffset)}; \textbf{GLSL:} \texttt{vec4 c = textureOffset(tex, uv, ivec2(1, -1));})
  \item \hyperref[enum:var_tex_compare_enable_m]{\texttt{compare\_enable}}: Enables depth-reference shadow comparison (\texttt{Dref} in staging registers). \emph{Example:} {\scriptsize\ttfamily \mnem{TEX\_SINGLE}\imod{.tex2d.compare} \dst{@r0}, \src{@r4:r6}} (\textbf{SPIR-V:} \texttt{OpImageSampleDrefImplicitLod}; \textbf{GLSL:} \texttt{float s = texture(shadowMap, vec3(uv, refZ));})
  \item \hyperref[enum:var_tex_dual_lod_mode_m]{\texttt{lod\_mode}}: Selects mipmap LOD calculation mode (\texttt{none}, \texttt{computed}, \texttt{explicit}, \texttt{computed\_bias}, \texttt{grdesc}). \emph{Example:} {\scriptsize\ttfamily \mnem{TEX\_SINGLE}\imod{.tex2d.explicit} \dst{@r0:r3}, \src{@r4:r6}} (\textbf{SPIR-V:} \texttt{OpImageSampleExplicitLod}; \textbf{GLSL:} \texttt{vec4 c = textureLod(tex, uv, 2.0);})
  \item \texttt{src2}: Third source operand (FMA addend, \texttt{MUX}/\texttt{CSEL} selector, \texttt{IADDC} carry-in, or bitwise mask in fused shift-logic). \emph{Example:} {\scriptsize\ttfamily \mnem{FMA.f32} \dst{r0}, \src{r1}, \src{r2}, \src{r3}} (\textbf{SPIR-V:} \texttt{GLSL.std.450 Fma (Addend)}; \textbf{GLSL:} \texttt{float d = fma(a, b, c);})
  \item \hyperref[enum:tex_dual_primary_width_m]{\texttt{primary\_width}}: Selects 16-bit vs 32-bit GPR element width for the corresponding message return window. \emph{Example:} {\scriptsize\ttfamily \mnem{TEX\_DUAL}\imod{.primary\_f16.secondary\_f32} \dst{@r0:r2}, \src{@r4:r5}} (\textbf{SPIR-V:} \texttt{OpImageSampleImplicitLod (RelaxedPrecision)}; \textbf{GLSL:} \texttt{mediump vec2 a = texture(t0, uv).rg;})
  \item \hyperref[enum:tex_dual_secondary_width_m]{\texttt{secondary\_width}}: Selects 16-bit vs 32-bit GPR element width for the corresponding message return window. \emph{Example:} {\scriptsize\ttfamily \mnem{TEX\_DUAL}\imod{.primary\_f16.secondary\_f32} \dst{@r0:r2}, \src{@r4:r5}} (\textbf{SPIR-V:} \texttt{OpImageSampleImplicitLod (RelaxedPrecision)}; \textbf{GLSL:} \texttt{mediump vec2 a = texture(t0, uv).rg;})
  \item \hyperref[enum:var_tex_dimensionality_m]{\texttt{dimensionality}}: Texture dimensionality (\texttt{tex1d}, \texttt{tex2d}, \texttt{tex3d}, \texttt{cube}). \emph{Example:} {\scriptsize\ttfamily \mnem{TEX\_SINGLE}\imod{.tex3d.f32.rgba} \dst{@r0:r3}, \src{@r4:r6}} (\textbf{SPIR-V:} \texttt{OpImageSampleImplicitLod (Dim=3D)}; \textbf{GLSL:} \texttt{vec4 c = texture(volTex, uvw);})
  \item \hyperref[enum:message_slot_index_m]{\texttt{message\_slot\_index}}: Asynchronous scoreboard slot (\texttt{slot0}--\texttt{slot2}, \texttt{slot7}) tracking completion of this message. \emph{Example:} {\scriptsize\ttfamily \mnem{LOAD.i32}\imod{.slot0} \dst{@r0}, \src{r2}, \imm{\#0}} (\textbf{SPIR-V:} \texttt{OpLoad}; \textbf{GLSL:} \texttt{uint v = ssbo.arr[i];})
  \item \hyperref[enum:encoded_sample_and_update_m]{\texttt{encoded\_sample\_and\_update}}: Jointly encodes MSAA sample position and hardware varying-cache update mode (\texttt{center\_store}, \texttt{retrieve}, etc.). \emph{Example:} {\scriptsize\ttfamily \mnem{VAR\_TEX\_BUF\_SINGLE}\imod{.center\_store.tex2d} \dst{@r0:r3}, \src{r4}} (\textbf{SPIR-V:} \texttt{OpImageSampleImplicitLod (fused varying)}; \textbf{GLSL:} \texttt{vec4 c = texture(s, v\_uv);})
  \item \texttt{sr\_write\_count}: Number of consecutive 32-bit staging registers written back by a dual-staging message (e.g. ATOM\textbackslash\{\}\_RETURN / TEX). \emph{Example:} {\scriptsize\ttfamily \mnem{TEX\_SINGLE}\imod{.tex2d.f32.rgba} \dst{@r0:r3}, \src{@r4:r5}} (\textbf{SPIR-V:} \texttt{OpImageSampleImplicitLod}; \textbf{GLSL:} \texttt{vec4 c = texture(tex, uv);})
  \item \hyperref[enum:tex_skip_m]{\texttt{skip}}: Skips texture/memory fetch on helper/discarded fragment lanes while preserving 2x2 quad derivatives. \emph{Example:} {\scriptsize\ttfamily \mnem{TEX\_SINGLE}\imod{.tex2d.computed.skip} \dst{@r0:r3}, \src{@r4:r5}} (\textbf{SPIR-V:} \texttt{OpImageSampleImplicitLod (after OpDemoteToHelperInvocation)}; \textbf{GLSL:} \texttt{if (a \textless{} 0.0) discard; vec4 c = texture(s, uv);})
  \item \texttt{staging\_register\_index}: Base GPR index of the contiguous staging register window (\texttt{@rN:rM}) for message payload/return. \emph{Example:} {\scriptsize\ttfamily \mnem{STORE.i64}\imod{.slot7} \src{@r0:r1}, \src{r2}, \imm{\#8}} (\textbf{SPIR-V:} \texttt{OpStore (uvec2)}; \textbf{GLSL:} \texttt{ssbo.pairs[i] = uvec2(a, b);})
  \item \texttt{staging\_register\_control\_32}: Controls the read/write staging register bank and count for async message instructions. \emph{Example:} {\scriptsize\ttfamily \mnem{LOAD.i128}\imod{.slot0} \dst{@r0:r3}, \src{r4}, \imm{\#0}} (\textbf{SPIR-V:} \texttt{OpLoad (vec4)}; \textbf{GLSL:} \texttt{vec4 v = ssbo.data[idx];})
  \item \texttt{opcode1}: Primary sub-opcode within the instruction class uniquely identifying the operation family. \emph{Example:} {\scriptsize\ttfamily \mnem{IADD.s32} \dst{r0}, \src{r1}, \src{r2}} (\textbf{SPIR-V:} \texttt{OpIAdd (operation selector)}; \textbf{GLSL:} \texttt{// Primary opcode bits})
  \item \texttt{instruction\_class}: 4-bit major instruction class selector (\texttt{bits 48:51}) partitioning the 64-bit Valhall opcode space (e.g., \texttt{FMA}, \texttt{CVT}, \texttt{SFU}, \texttt{LS}, \texttt{TEX}, \texttt{BRANCH}). \emph{Example:} {\scriptsize\ttfamily \mnem{FADD.f32} \dst{r0}, \src{r1}, \src{r2}} (\textbf{SPIR-V:} \texttt{N/A (encoding class)}; \textbf{GLSL:} \texttt{// Hardware execution pipe / decode class})
  \item \texttt{fau\_page\_index}: Selects the FAU (Fast Access Uniform) page (\texttt{0..3}) when a source operand references hardware small constants, thread IDs, or lane constants. \emph{Example:} {\scriptsize\ttfamily \mnem{FADD.f32} \dst{r0}, \src{r1}, \src{k1f}} (\textbf{SPIR-V:} \texttt{OpConstant / BuiltIn}; \textbf{GLSL:} \texttt{float y = x + 1.0;})
  \item \hyperref[enum:flow_control_m]{\texttt{flow}}: 4-bit embedded flow-control \& scoreboard wait modifier (\texttt{bits 59:62}: \texttt{none}, \texttt{wait0..wait0126}, \texttt{wait}, \texttt{reconverge}, \texttt{discard}, \texttt{end}). \emph{Example:} {\scriptsize\ttfamily \mnem{FADD.f32}\imod{.wait0.reconverge} \dst{r0}, \src{r1}, \src{r2}} (\textbf{SPIR-V:} \texttt{OpControlBarrier / OpReturn}; \textbf{GLSL:} \texttt{// Scoreboard wait + warp reconvergence})
\end{tightitemize}
\clearpage

\subsection*{VAR\_TEX\_GATHER}
\addcontentsline{toc}{subsection}{VAR\_TEX\_GATHER}
Fuses descriptor-based varying UV interpolation (via varying handle \texttt{src2}) with a 2x2 texture gather (\texttt{TEX\_GATHER} via texture handle \texttt{src0}), writing the 4 gathered component values into \texttt{@sw} in a single message.
\par\textbf{Notes:}
\begin{tightitemize}
  \item Fragment-only.
  \item 2D (\texttt{encode\_v9.rs:3749}).
  \item \textbf{Source references:} \texttt{kraid/ops.rs:4931-4981}, \texttt{encode\_v9.rs:3736-3766}, \texttt{bifrost/valhall/ISA.xml:165}.
\end{tightitemize}
\codebox{\mnem{VAR\_TEX\_GATHER}\imod{\{.message\_slot\_index\}}\imod{\{.flow\}}\imod{\{.skip\}}\imod{\{.sample\_position\}}\imod{\{.update\}}\imod{\{.dimensionality\}}\imod{\{.register\_width\}}\imod{\{.write\_mask\}}\imod{\{.wide\_indices\}}\imod{\{.array\_enable\}}\imod{\{.compare\_enable\}}\imod{\{.gather\_component\}} \dst{@sw}, \src{src0}, \src{src2}\\\mnem{VAR\_TEX\_GATHER}\imod{\{.message\_slot\_index\}}\imod{\{.flow\}}\imod{\{.skip\}}\imod{\{.sample\_position\}}\imod{\{.update\}}\imod{\{.dimensionality\}}\imod{\{.projection\_enable\}}\imod{\{.register\_width\}}\imod{\{.write\_mask\}}\imod{\{.wide\_indices\}}\imod{\{.array\_enable\}}\imod{\{.compare\_enable\}}\imod{\{.gather\_component\}} \dst{@sw}, \src{src0}, \src{src2}}
\textbf{Operands:}
\begin{tightitemize}\footnotesize
  \item \texttt{\dst{@sw}} (\texttt{staging\_register\_index} + \texttt{sr\_write\_count}): Destination staging register range receiving the 2x2 texture gather result.
  \item \texttt{\src{src0}}: Texture/sampler descriptor handle.
  \item \texttt{\src{src2}}: Varying descriptor handle providing interpolated UV coordinates (encoded in \texttt{src2}).
\end{tightitemize}
\par\textbf{Closest SPIR-V Op:} \texttt{OpLoad (in) + OpImageGather}
\par\textbf{Pseudo-GLSL (Operational Semantics):}
\codebox{@sw = textureGather(handle = src0, interpolateVarying(handle = src2), gather\_component);}
\textbf{Execution Context:}
\begin{tightitemize}\footnotesize
  \item \textbf{Varying buffer \& barycentric cache (\hyperref[sec:csf-tiler]{\texttt{IDVSSR::VARY\_SIZE}}, \texttt{r48:r49}):} Reads per-vertex varying RAM written by Varying EBIN \texttt{STORE.estream}; descriptor forms (\texttt{LD\_VAR*}) also chase \hyperref[sec:csf-srt]{\texttt{SRT\_2}} $\to$ \texttt{Resource[table]} $\to$ \texttt{Attribute[index]} and reuse cached barycentric weights via \texttt{.store}/\texttt{.retrieve}.
  \item \textbf{Texture \& Sampler chain (\hyperref[sec:csf-srt]{\texttt{SRT\_2}} Frag / \texttt{SRT\_0} Comp/Vert):} 64B-aligned VA (\texttt{\& \textasciitilde{}0x3F}) $\to$ \texttt{Resource[table]} $\to$ \hyperref[sec:csf-tex-samp]{\texttt{Texture[tex\_idx]}} (32B, \texttt{type=2}) $\to$ \texttt{GenericPlane[mip]} (32B, \texttt{type=10}: \texttt{pointer}, \texttt{row\_stride}, Morton 16$\times$16 / AFBC) + \texttt{Sampler[samp\_idx]} (32B, \texttt{type=1}: filter, wrap, LOD clamp/bias, compare).
  \item \textbf{Fragment stage preloads \& FP modes (\hyperref[sec:csf-spd]{\texttt{SPD\_2}}, \texttt{r20:r21}):} Hardware pre-fills \texttt{r59} with \texttt{floatBitsToUint(FragCoord.z)} and \texttt{r60} with the MSAA sample coverage mask (\texttt{gl\_SampleMaskIn[0]}), while \texttt{ShaderProgram} word~0 latches 32/64 GPR allocation and warp-wide \texttt{flush\_to\_zero\_mode}, \texttt{suppress\_na\_n}, and \texttt{suppress\_inf}.
\end{tightitemize}
\par\textbf{Worked Examples:}
\par\textbf{Example 1 (Descriptor varying + 2x2 red gather):}
\codebox{\mnem{VAR\_TEX\_GATHER}\imod{.tex2d.gather4\_r.center.slot0} \dst{@r0:r3}, \src{u0}, \src{r59}, \src{u1}}
Fuses descriptor-based \texttt{vUV} interpolation (\texttt{u1}) with \texttt{textureGather(..., 0)} into \texttt{@r0:r3}.
\par\textbf{Example 2 (Green-channel gather (`.gather4\_g`)):}
\codebox{\mnem{VAR\_TEX\_GATHER}\imod{.tex2d.gather4\_g.center.slot0} \dst{@r0:r3}, \src{u0}, \src{r59}, \src{u1}}
Gathers the 2x2 green channel footprint directly from interpolated \texttt{vUV}.

\begin{center}\footnotesize
\begin{tabularx}{\linewidth}{@{}lllX@{}}
\toprule
Variant & Arch & Class & Opcode Fields \\
\midrule
\texttt{(none)} & v9--v14 & \texttt{S3A} & \texttt{opcode1=VAR\_TEX\_GATHER} \\
\texttt{(none)} & v9--v14 & \texttt{S3A} & \texttt{opcode1=VAR\_TEX\_GATHER} \\
\bottomrule
\end{tabularx}
\end{center}
\begin{center}\footnotesize
\begin{tabularx}{0.88\linewidth}{@{}l l X@{}}
\toprule
Field Name & Bits & Type / Fixed Value \\
\midrule
\texttt{src0} & \texttt{0:7} & \texttt{SourceEncoding64} \\
\hyperref[enum:tex_wide_indices_enable_m]{\texttt{wide\_indices}} & \texttt{8} & \texttt{tex\_wide\_indices\_enable\_m (def none)} \\
\textit{Reserved} & \texttt{9} & \texttt{zero} \\
\hyperref[enum:var_tex_array_enable_m]{\texttt{array\_enable}} & \texttt{10} & \texttt{var\_tex\_array\_enable\_m (def none)} \\
\hyperref[enum:tex_texel_offset_m]{\texttt{texel\_offset}} & \texttt{11} & \texttt{exact=none} \\
\hyperref[enum:var_tex_compare_enable_m]{\texttt{compare\_enable}} & \texttt{12} & \texttt{var\_tex\_compare\_enable\_m (def none)} \\
\hyperref[enum:tex_coordinate_mode_m]{\texttt{coordinate\_mode}} & \texttt{13} & \texttt{exact=float\_coordinates} \\
\hyperref[enum:tex_gather_component_m]{\texttt{gather\_component}} & \texttt{14:15} & \texttt{tex\_gather\_component\_m} \\
\texttt{src2} & \texttt{16:23} & \texttt{SourceEncoding} \\
\hyperref[enum:tex_width_m]{\texttt{register\_width}} & \texttt{24} & \texttt{tex\_width\_m} \\
\texttt{secondary\_width} & \texttt{25} & \texttt{exact=0} \\
\textit{Reserved} & \texttt{26:27} & \texttt{ignore} \\
\hyperref[enum:var_tex_dimensionality_m]{\texttt{dimensionality}} & \texttt{28:29} & \texttt{exact=tex2d} \\
\hyperref[enum:message_slot_index_m]{\texttt{message\_slot\_index}} & \texttt{30:32} & \texttt{message\_slot\_index\_m} \\
\hyperref[enum:encoded_sample_and_update_m]{\texttt{encoded\_sample\_and\_update}} & \texttt{33:35} & \texttt{encoded\_sample\_and\_update\_m} \\
\texttt{sr\_write\_count} & \texttt{36:38} & \texttt{sr\_write\_count\_t} \\
\hyperref[enum:tex_skip_m]{\texttt{skip}} & \texttt{39} & \texttt{tex\_skip\_m (def none)} \\
\texttt{staging\_register\_index} & \texttt{40:45} & \texttt{uint} \\
\texttt{staging\_register\_control\_32} & \texttt{46:47} & \texttt{exact=write} \\
\texttt{opcode1} & \texttt{48:51} & \texttt{exact=VAR\_TEX\_GATHER} \\
\texttt{instruction\_class} & \texttt{52:56} & \texttt{exact=S3A} \\
\texttt{fau\_page\_index} & \texttt{57:58} & \texttt{uint} \\
\hyperref[enum:flow_control_m]{\texttt{flow}} & \texttt{59:62} & \texttt{flow\_control\_m} \\
\textit{Reserved} & \texttt{63} & \texttt{zero} \\
\bottomrule
\end{tabularx}
\end{center}
\textbf{Fields Explained:}
\begin{tightitemize}\footnotesize
  \item \texttt{src0}: Primary source operand (\texttt{r0..r63}, \texttt{u0..u63} uniform, \texttt{k*} FAU constant, or \texttt{t0..t2} temporary; \texttt{\textasciicircum{}} marks discard/last-use). \emph{Example:} {\scriptsize\ttfamily \mnem{FADD.f32} \dst{r0}, \src{r1}, \src{r2}} (\textbf{SPIR-V:} \texttt{OpFAdd (Operand 0)}; \textbf{GLSL:} \texttt{float c = a + b;})
  \item \hyperref[enum:tex_wide_indices_enable_m]{\texttt{wide\_indices}}: Uses 32-bit wide texture/sampler table indices instead of compact 16-bit packed indices. \emph{Example:} {\scriptsize\ttfamily \mnem{TEX\_SINGLE}\imod{.tex2d.wide} \dst{@r0:r3}, \src{@r4:r5}, \src{r6}} (\textbf{SPIR-V:} \texttt{OpImageSampleImplicitLod (bindless)}; \textbf{GLSL:} \texttt{vec4 c = texture(textures[nonuniformEXT(i)], uv);})
  \item \hyperref[enum:var_tex_array_enable_m]{\texttt{array\_enable}}: Enables array-layer coordinate in texture staging registers (\texttt{sampler2DArray} / \texttt{samplerCubeArray}). \emph{Example:} {\scriptsize\ttfamily \mnem{TEX\_SINGLE}\imod{.tex2d.array} \dst{@r0:r3}, \src{@r4:r6}} (\textbf{SPIR-V:} \texttt{OpImageSampleImplicitLod (Arrayed)}; \textbf{GLSL:} \texttt{vec4 c = texture(texArr, vec3(uv, layer));})
  \item \hyperref[enum:tex_texel_offset_m]{\texttt{texel\_offset}}: Applies a signed 4-bit per-axis texel offset $(dx, dy, dz)$ packed in staging registers. \emph{Example:} {\scriptsize\ttfamily \mnem{TEX\_SINGLE}\imod{.tex2d.texel\_offset} \dst{@r0:r3}, \src{@r4:r6}} (\textbf{SPIR-V:} \texttt{OpImageSampleImplicitLod (ConstOffset)}; \textbf{GLSL:} \texttt{vec4 c = textureOffset(tex, uv, ivec2(1, -1));})
  \item \hyperref[enum:var_tex_compare_enable_m]{\texttt{compare\_enable}}: Enables depth-reference shadow comparison (\texttt{Dref} in staging registers). \emph{Example:} {\scriptsize\ttfamily \mnem{TEX\_SINGLE}\imod{.tex2d.compare} \dst{@r0}, \src{@r4:r6}} (\textbf{SPIR-V:} \texttt{OpImageSampleDrefImplicitLod}; \textbf{GLSL:} \texttt{float s = texture(shadowMap, vec3(uv, refZ));})
  \item \hyperref[enum:tex_coordinate_mode_m]{\texttt{coordinate\_mode}}: Selects normalized float coordinates vs unnormalized integer texel indices. \emph{Example:} {\scriptsize\ttfamily \mnem{TEX\_FETCH}\imod{.tex2d.integer\_coordinates} \dst{@r0:r3}, \src{@r4:r5}} (\textbf{SPIR-V:} \texttt{OpImageFetch}; \textbf{GLSL:} \texttt{vec4 c = texelFetch(tex, ivec2(xy), 0);})
  \item \hyperref[enum:tex_gather_component_m]{\texttt{gather\_component}}: Selects the R, G, B, or A channel gathered across the 2x2 texel footprint. \emph{Example:} {\scriptsize\ttfamily \mnem{TEX\_GATHER}\imod{.tex2d.gather4\_g} \dst{@r0:r3}, \src{@r4:r5}} (\textbf{SPIR-V:} \texttt{OpImageGather}; \textbf{GLSL:} \texttt{vec4 g = textureGather(tex, uv, 1);})
  \item \texttt{src2}: Third source operand (FMA addend, \texttt{MUX}/\texttt{CSEL} selector, \texttt{IADDC} carry-in, or bitwise mask in fused shift-logic). \emph{Example:} {\scriptsize\ttfamily \mnem{FMA.f32} \dst{r0}, \src{r1}, \src{r2}, \src{r3}} (\textbf{SPIR-V:} \texttt{GLSL.std.450 Fma (Addend)}; \textbf{GLSL:} \texttt{float d = fma(a, b, c);})
  \item \hyperref[enum:tex_width_m]{\texttt{register\_width}}: Selects 16-bit vs 32-bit GPR element width for the corresponding message return window. \emph{Example:} {\scriptsize\ttfamily \mnem{TEX\_DUAL}\imod{.primary\_f16.secondary\_f32} \dst{@r0:r2}, \src{@r4:r5}} (\textbf{SPIR-V:} \texttt{OpImageSampleImplicitLod (RelaxedPrecision)}; \textbf{GLSL:} \texttt{mediump vec2 a = texture(t0, uv).rg;})
  \item \texttt{secondary\_width}: Selects 16-bit vs 32-bit GPR element width for the corresponding message return window. \emph{Example:} {\scriptsize\ttfamily \mnem{TEX\_DUAL}\imod{.primary\_f16.secondary\_f32} \dst{@r0:r2}, \src{@r4:r5}} (\textbf{SPIR-V:} \texttt{OpImageSampleImplicitLod (RelaxedPrecision)}; \textbf{GLSL:} \texttt{mediump vec2 a = texture(t0, uv).rg;})
  \item \hyperref[enum:var_tex_dimensionality_m]{\texttt{dimensionality}}: Texture dimensionality (\texttt{tex1d}, \texttt{tex2d}, \texttt{tex3d}, \texttt{cube}). \emph{Example:} {\scriptsize\ttfamily \mnem{TEX\_SINGLE}\imod{.tex3d.f32.rgba} \dst{@r0:r3}, \src{@r4:r6}} (\textbf{SPIR-V:} \texttt{OpImageSampleImplicitLod (Dim=3D)}; \textbf{GLSL:} \texttt{vec4 c = texture(volTex, uvw);})
  \item \hyperref[enum:message_slot_index_m]{\texttt{message\_slot\_index}}: Asynchronous scoreboard slot (\texttt{slot0}--\texttt{slot2}, \texttt{slot7}) tracking completion of this message. \emph{Example:} {\scriptsize\ttfamily \mnem{LOAD.i32}\imod{.slot0} \dst{@r0}, \src{r2}, \imm{\#0}} (\textbf{SPIR-V:} \texttt{OpLoad}; \textbf{GLSL:} \texttt{uint v = ssbo.arr[i];})
  \item \hyperref[enum:encoded_sample_and_update_m]{\texttt{encoded\_sample\_and\_update}}: Jointly encodes MSAA sample position and hardware varying-cache update mode (\texttt{center\_store}, \texttt{retrieve}, etc.). \emph{Example:} {\scriptsize\ttfamily \mnem{VAR\_TEX\_BUF\_SINGLE}\imod{.center\_store.tex2d} \dst{@r0:r3}, \src{r4}} (\textbf{SPIR-V:} \texttt{OpImageSampleImplicitLod (fused varying)}; \textbf{GLSL:} \texttt{vec4 c = texture(s, v\_uv);})
  \item \texttt{sr\_write\_count}: Number of consecutive 32-bit staging registers written back by a dual-staging message (e.g. ATOM\textbackslash\{\}\_RETURN / TEX). \emph{Example:} {\scriptsize\ttfamily \mnem{TEX\_SINGLE}\imod{.tex2d.f32.rgba} \dst{@r0:r3}, \src{@r4:r5}} (\textbf{SPIR-V:} \texttt{OpImageSampleImplicitLod}; \textbf{GLSL:} \texttt{vec4 c = texture(tex, uv);})
  \item \hyperref[enum:tex_skip_m]{\texttt{skip}}: Skips texture/memory fetch on helper/discarded fragment lanes while preserving 2x2 quad derivatives. \emph{Example:} {\scriptsize\ttfamily \mnem{TEX\_SINGLE}\imod{.tex2d.computed.skip} \dst{@r0:r3}, \src{@r4:r5}} (\textbf{SPIR-V:} \texttt{OpImageSampleImplicitLod (after OpDemoteToHelperInvocation)}; \textbf{GLSL:} \texttt{if (a \textless{} 0.0) discard; vec4 c = texture(s, uv);})
  \item \texttt{staging\_register\_index}: Base GPR index of the contiguous staging register window (\texttt{@rN:rM}) for message payload/return. \emph{Example:} {\scriptsize\ttfamily \mnem{STORE.i64}\imod{.slot7} \src{@r0:r1}, \src{r2}, \imm{\#8}} (\textbf{SPIR-V:} \texttt{OpStore (uvec2)}; \textbf{GLSL:} \texttt{ssbo.pairs[i] = uvec2(a, b);})
  \item \texttt{staging\_register\_control\_32}: Controls the read/write staging register bank and count for async message instructions. \emph{Example:} {\scriptsize\ttfamily \mnem{LOAD.i128}\imod{.slot0} \dst{@r0:r3}, \src{r4}, \imm{\#0}} (\textbf{SPIR-V:} \texttt{OpLoad (vec4)}; \textbf{GLSL:} \texttt{vec4 v = ssbo.data[idx];})
  \item \texttt{opcode1}: Primary sub-opcode within the instruction class uniquely identifying the operation family. \emph{Example:} {\scriptsize\ttfamily \mnem{IADD.s32} \dst{r0}, \src{r1}, \src{r2}} (\textbf{SPIR-V:} \texttt{OpIAdd (operation selector)}; \textbf{GLSL:} \texttt{// Primary opcode bits})
  \item \texttt{instruction\_class}: 4-bit major instruction class selector (\texttt{bits 48:51}) partitioning the 64-bit Valhall opcode space (e.g., \texttt{FMA}, \texttt{CVT}, \texttt{SFU}, \texttt{LS}, \texttt{TEX}, \texttt{BRANCH}). \emph{Example:} {\scriptsize\ttfamily \mnem{FADD.f32} \dst{r0}, \src{r1}, \src{r2}} (\textbf{SPIR-V:} \texttt{N/A (encoding class)}; \textbf{GLSL:} \texttt{// Hardware execution pipe / decode class})
  \item \texttt{fau\_page\_index}: Selects the FAU (Fast Access Uniform) page (\texttt{0..3}) when a source operand references hardware small constants, thread IDs, or lane constants. \emph{Example:} {\scriptsize\ttfamily \mnem{FADD.f32} \dst{r0}, \src{r1}, \src{k1f}} (\textbf{SPIR-V:} \texttt{OpConstant / BuiltIn}; \textbf{GLSL:} \texttt{float y = x + 1.0;})
  \item \hyperref[enum:flow_control_m]{\texttt{flow}}: 4-bit embedded flow-control \& scoreboard wait modifier (\texttt{bits 59:62}: \texttt{none}, \texttt{wait0..wait0126}, \texttt{wait}, \texttt{reconverge}, \texttt{discard}, \texttt{end}). \emph{Example:} {\scriptsize\ttfamily \mnem{FADD.f32}\imod{.wait0.reconverge} \dst{r0}, \src{r1}, \src{r2}} (\textbf{SPIR-V:} \texttt{OpControlBarrier / OpReturn}; \textbf{GLSL:} \texttt{// Scoreboard wait + warp reconvergence})
\end{tightitemize}
\clearpage

\subsection*{VAR\_TEX\_GRADIENT}
\addcontentsline{toc}{subsection}{VAR\_TEX\_GRADIENT}
Fuses descriptor-based varying UV interpolation (via varying handle \texttt{src2}) with a texture gradient/LOD operation (\texttt{TEX\_GRADIENT} via texture handle \texttt{src0}), writing the result into \texttt{@sw} in a single message.
\par\textbf{Notes:}
\begin{tightitemize}
  \item Fragment-only.
  \item FP32 varyings (\texttt{ISA.xml:167}).
  \item \textbf{Source references:} \texttt{kraid/ops.rs:4983-5034}, \texttt{encode\_v9.rs:3768-3797}, \texttt{bifrost/valhall/ISA.xml:167}.
\end{tightitemize}
\codebox{\mnem{VAR\_TEX\_GRADIENT}\imod{\{.message\_slot\_index\}}\imod{\{.flow\}}\imod{\{.skip\}}\imod{\{.sample\_position\}}\imod{\{.update\}}\imod{\{.dimensionality\}}\imod{\{.register\_width\}}\imod{\{.write\_mask\}}\imod{\{.wide\_indices\}}\imod{\{.lod\_bias\_disable\}}\imod{\{.lod\_clamp\_disable\}} \dst{@sw}, \src{src0}, \src{src2}\\\mnem{VAR\_TEX\_GRADIENT}\imod{\{.message\_slot\_index\}}\imod{\{.flow\}}\imod{\{.skip\}}\imod{\{.sample\_position\}}\imod{\{.update\}}\imod{\{.dimensionality\}}\imod{\{.projection\_enable\}}\imod{\{.register\_width\}}\imod{\{.write\_mask\}}\imod{\{.wide\_indices\}}\imod{\{.lod\_bias\_disable\}}\imod{\{.lod\_clamp\_disable\}} \dst{@sw}, \src{src0}, \src{src2}}
\textbf{Operands:}
\begin{tightitemize}\footnotesize
  \item \texttt{\dst{@sw}} (\texttt{staging\_register\_index} + \texttt{sr\_write\_count}): Destination staging register range receiving the gradient/LOD texture operation result.
  \item \texttt{\src{src0}}: Texture/sampler descriptor handle.
  \item \texttt{\src{src2}}: Varying descriptor handle providing interpolated UV coordinates (encoded in \texttt{src2}).
\end{tightitemize}
\par\textbf{Closest SPIR-V Op:} \texttt{OpLoad (in) + OpImageSampleExplicitLod}
\par\textbf{Pseudo-GLSL (Operational Semantics):}
\codebox{@sw = evalTexGradOrLod(handle = src0, interpolateVarying(handle = src2));}
\textbf{Execution Context:}
\begin{tightitemize}\footnotesize
  \item \textbf{Varying buffer \& barycentric cache (\hyperref[sec:csf-tiler]{\texttt{IDVSSR::VARY\_SIZE}}, \texttt{r48:r49}):} Reads per-vertex varying RAM written by Varying EBIN \texttt{STORE.estream}; descriptor forms (\texttt{LD\_VAR*}) also chase \hyperref[sec:csf-srt]{\texttt{SRT\_2}} $\to$ \texttt{Resource[table]} $\to$ \texttt{Attribute[index]} and reuse cached barycentric weights via \texttt{.store}/\texttt{.retrieve}.
  \item \textbf{Texture \& Sampler chain (\hyperref[sec:csf-srt]{\texttt{SRT\_2}} Frag / \texttt{SRT\_0} Comp/Vert):} 64B-aligned VA (\texttt{\& \textasciitilde{}0x3F}) $\to$ \texttt{Resource[table]} $\to$ \hyperref[sec:csf-tex-samp]{\texttt{Texture[tex\_idx]}} (32B, \texttt{type=2}) $\to$ \texttt{GenericPlane[mip]} (32B, \texttt{type=10}: \texttt{pointer}, \texttt{row\_stride}, Morton 16$\times$16 / AFBC) + \texttt{Sampler[samp\_idx]} (32B, \texttt{type=1}: filter, wrap, LOD clamp/bias, compare).
  \item \textbf{Fragment stage preloads \& FP modes (\hyperref[sec:csf-spd]{\texttt{SPD\_2}}, \texttt{r20:r21}):} Hardware pre-fills \texttt{r59} with \texttt{floatBitsToUint(FragCoord.z)} and \texttt{r60} with the MSAA sample coverage mask (\texttt{gl\_SampleMaskIn[0]}), while \texttt{ShaderProgram} word~0 latches 32/64 GPR allocation and warp-wide \texttt{flush\_to\_zero\_mode}, \texttt{suppress\_na\_n}, and \texttt{suppress\_inf}.
\end{tightitemize}
\par\textbf{Worked Examples:}
\par\textbf{Example 1 (Descriptor varying + explicit gradient/LOD):}
\codebox{\mnem{VAR\_TEX\_GRADIENT}\imod{.tex2d.center.slot0} \dst{@r0:r3}, \src{u0}, \src{r59}, \src{u1}}
Fuses descriptor-based \texttt{vUV} interpolation (\texttt{u1}) with hardware texture gradient evaluation.
\par\textbf{Example 2 (Raw derivative output (`.derivative`)):}
\codebox{\mnem{VAR\_TEX\_GRADIENT}\imod{.tex2d.derivative.center.slot0} \dst{@r0:r3}, \src{u0}, \src{r59}, \src{u1}}
Outputs raw \texttt{(dU/dx, dV/dx, dU/dy, dV/dy)} derivatives directly into \texttt{@r0:r3}.

\begin{center}\footnotesize
\begin{tabularx}{\linewidth}{@{}lllX@{}}
\toprule
Variant & Arch & Class & Opcode Fields \\
\midrule
\texttt{(none)} & v9--v14 & \texttt{S3A} & \texttt{opcode1=VAR\_TEX\_GRADIENT} \\
\texttt{(none)} & v9--v14 & \texttt{S3A} & \texttt{opcode1=VAR\_TEX\_GRADIENT} \\
\texttt{(none)} & v9--v14 & \texttt{S3A} & \texttt{opcode1=VAR\_TEX\_GRADIENT} \\
\bottomrule
\end{tabularx}
\end{center}
\begin{center}\footnotesize
\begin{tabularx}{0.88\linewidth}{@{}l l X@{}}
\toprule
Field Name & Bits & Type / Fixed Value \\
\midrule
\texttt{src0} & \texttt{0:7} & \texttt{SourceEncoding64} \\
\hyperref[enum:tex_wide_indices_enable_m]{\texttt{wide\_indices}} & \texttt{8} & \texttt{tex\_wide\_indices\_enable\_m (def none)} \\
\textit{Reserved} & \texttt{9:12} & \texttt{zero} \\
\hyperref[enum:tex_lod_bias_disable_m]{\texttt{lod\_bias\_disable}} & \texttt{13} & \texttt{tex\_lod\_bias\_disable\_m} \\
\hyperref[enum:tex_lod_clamp_disable_m]{\texttt{lod\_clamp\_disable}} & \texttt{14} & \texttt{tex\_lod\_clamp\_disable\_m} \\
\hyperref[enum:tex_derivative_enable_m]{\texttt{derivative\_enable}} & \texttt{15} & \texttt{exact=none} \\
\texttt{src2} & \texttt{16:23} & \texttt{SourceEncoding} \\
\hyperref[enum:tex_gradient_width_m]{\texttt{register\_width}} & \texttt{24} & \texttt{tex\_gradient\_width\_m} \\
\texttt{secondary\_width} & \texttt{25} & \texttt{exact=0} \\
\textit{Reserved} & \texttt{26:27} & \texttt{ignore} \\
\hyperref[enum:var_tex_dimensionality_m]{\texttt{dimensionality}} & \texttt{28:29} & \texttt{var\_tex\_dimensionality\_m} \\
\hyperref[enum:message_slot_index_m]{\texttt{message\_slot\_index}} & \texttt{30:32} & \texttt{message\_slot\_index\_m} \\
\hyperref[enum:encoded_sample_and_update_m]{\texttt{encoded\_sample\_and\_update}} & \texttt{33:35} & \texttt{encoded\_sample\_and\_update\_m} \\
\texttt{sr\_write\_count} & \texttt{36:38} & \texttt{sr\_write\_count\_t} \\
\hyperref[enum:tex_gradient_skip_m]{\texttt{skip}} & \texttt{39} & \texttt{tex\_gradient\_skip\_m (def none)} \\
\texttt{staging\_register\_index} & \texttt{40:45} & \texttt{uint} \\
\texttt{staging\_register\_control\_32} & \texttt{46:47} & \texttt{exact=write} \\
\texttt{opcode1} & \texttt{48:51} & \texttt{exact=VAR\_TEX\_GRADIENT} \\
\texttt{instruction\_class} & \texttt{52:56} & \texttt{exact=S3A} \\
\texttt{fau\_page\_index} & \texttt{57:58} & \texttt{uint} \\
\hyperref[enum:flow_control_m]{\texttt{flow}} & \texttt{59:62} & \texttt{flow\_control\_m} \\
\textit{Reserved} & \texttt{63} & \texttt{zero} \\
\bottomrule
\end{tabularx}
\end{center}
\textbf{Fields Explained:}
\begin{tightitemize}\footnotesize
  \item \texttt{src0}: Primary source operand (\texttt{r0..r63}, \texttt{u0..u63} uniform, \texttt{k*} FAU constant, or \texttt{t0..t2} temporary; \texttt{\textasciicircum{}} marks discard/last-use). \emph{Example:} {\scriptsize\ttfamily \mnem{FADD.f32} \dst{r0}, \src{r1}, \src{r2}} (\textbf{SPIR-V:} \texttt{OpFAdd (Operand 0)}; \textbf{GLSL:} \texttt{float c = a + b;})
  \item \hyperref[enum:tex_wide_indices_enable_m]{\texttt{wide\_indices}}: Uses 32-bit wide texture/sampler table indices instead of compact 16-bit packed indices. \emph{Example:} {\scriptsize\ttfamily \mnem{TEX\_SINGLE}\imod{.tex2d.wide} \dst{@r0:r3}, \src{@r4:r5}, \src{r6}} (\textbf{SPIR-V:} \texttt{OpImageSampleImplicitLod (bindless)}; \textbf{GLSL:} \texttt{vec4 c = texture(textures[nonuniformEXT(i)], uv);})
  \item \hyperref[enum:tex_lod_bias_disable_m]{\texttt{lod\_bias\_disable}}: Disables sampler/shader LOD bias when computing explicit gradients (\texttt{TEX\_GRADIENT}). \emph{Example:} {\scriptsize\ttfamily \mnem{TEX\_GRADIENT}\imod{.tex2d.disable\_lod\_bias} \dst{@r0:r1}, \src{@r4:r5}} (\textbf{SPIR-V:} \texttt{OpImageQueryLod}; \textbf{GLSL:} \texttt{vec2 lod = textureQueryLod(tex, uv);})
  \item \hyperref[enum:tex_lod_clamp_disable_m]{\texttt{lod\_clamp\_disable}}: Disables sampler min/max LOD clamp bounds when querying raw mip level. \emph{Example:} {\scriptsize\ttfamily \mnem{TEX\_GRADIENT}\imod{.tex2d.disable\_lod\_clamp} \dst{@r0:r1}, \src{@r4:r5}} (\textbf{SPIR-V:} \texttt{OpImageQueryLod}; \textbf{GLSL:} \texttt{vec2 lod = textureQueryLod(tex, uv);})
  \item \hyperref[enum:tex_derivative_enable_m]{\texttt{derivative\_enable}}: Returns raw screen-space derivatives ($dP/dx, dP/dy$) instead of scalar LOD. \emph{Example:} {\scriptsize\ttfamily \mnem{TEX\_GRADIENT}\imod{.tex2d.derivative} \dst{@r0:r3}, \src{@r4:r5}} (\textbf{SPIR-V:} \texttt{OpDPdx / OpDPdy}; \textbf{GLSL:} \texttt{vec2 dx = dFdx(uv);})
  \item \texttt{src2}: Third source operand (FMA addend, \texttt{MUX}/\texttt{CSEL} selector, \texttt{IADDC} carry-in, or bitwise mask in fused shift-logic). \emph{Example:} {\scriptsize\ttfamily \mnem{FMA.f32} \dst{r0}, \src{r1}, \src{r2}, \src{r3}} (\textbf{SPIR-V:} \texttt{GLSL.std.450 Fma (Addend)}; \textbf{GLSL:} \texttt{float d = fma(a, b, c);})
  \item \hyperref[enum:tex_gradient_width_m]{\texttt{register\_width}}: Selects 16-bit vs 32-bit GPR element width for the corresponding message return window. \emph{Example:} {\scriptsize\ttfamily \mnem{TEX\_DUAL}\imod{.primary\_f16.secondary\_f32} \dst{@r0:r2}, \src{@r4:r5}} (\textbf{SPIR-V:} \texttt{OpImageSampleImplicitLod (RelaxedPrecision)}; \textbf{GLSL:} \texttt{mediump vec2 a = texture(t0, uv).rg;})
  \item \texttt{secondary\_width}: Selects 16-bit vs 32-bit GPR element width for the corresponding message return window. \emph{Example:} {\scriptsize\ttfamily \mnem{TEX\_DUAL}\imod{.primary\_f16.secondary\_f32} \dst{@r0:r2}, \src{@r4:r5}} (\textbf{SPIR-V:} \texttt{OpImageSampleImplicitLod (RelaxedPrecision)}; \textbf{GLSL:} \texttt{mediump vec2 a = texture(t0, uv).rg;})
  \item \hyperref[enum:var_tex_dimensionality_m]{\texttt{dimensionality}}: Texture dimensionality (\texttt{tex1d}, \texttt{tex2d}, \texttt{tex3d}, \texttt{cube}). \emph{Example:} {\scriptsize\ttfamily \mnem{TEX\_SINGLE}\imod{.tex3d.f32.rgba} \dst{@r0:r3}, \src{@r4:r6}} (\textbf{SPIR-V:} \texttt{OpImageSampleImplicitLod (Dim=3D)}; \textbf{GLSL:} \texttt{vec4 c = texture(volTex, uvw);})
  \item \hyperref[enum:message_slot_index_m]{\texttt{message\_slot\_index}}: Asynchronous scoreboard slot (\texttt{slot0}--\texttt{slot2}, \texttt{slot7}) tracking completion of this message. \emph{Example:} {\scriptsize\ttfamily \mnem{LOAD.i32}\imod{.slot0} \dst{@r0}, \src{r2}, \imm{\#0}} (\textbf{SPIR-V:} \texttt{OpLoad}; \textbf{GLSL:} \texttt{uint v = ssbo.arr[i];})
  \item \hyperref[enum:encoded_sample_and_update_m]{\texttt{encoded\_sample\_and\_update}}: Jointly encodes MSAA sample position and hardware varying-cache update mode (\texttt{center\_store}, \texttt{retrieve}, etc.). \emph{Example:} {\scriptsize\ttfamily \mnem{VAR\_TEX\_BUF\_SINGLE}\imod{.center\_store.tex2d} \dst{@r0:r3}, \src{r4}} (\textbf{SPIR-V:} \texttt{OpImageSampleImplicitLod (fused varying)}; \textbf{GLSL:} \texttt{vec4 c = texture(s, v\_uv);})
  \item \texttt{sr\_write\_count}: Number of consecutive 32-bit staging registers written back by a dual-staging message (e.g. ATOM\textbackslash\{\}\_RETURN / TEX). \emph{Example:} {\scriptsize\ttfamily \mnem{TEX\_SINGLE}\imod{.tex2d.f32.rgba} \dst{@r0:r3}, \src{@r4:r5}} (\textbf{SPIR-V:} \texttt{OpImageSampleImplicitLod}; \textbf{GLSL:} \texttt{vec4 c = texture(tex, uv);})
  \item \hyperref[enum:tex_gradient_skip_m]{\texttt{skip}}: Skips texture/memory fetch on helper/discarded fragment lanes while preserving 2x2 quad derivatives. \emph{Example:} {\scriptsize\ttfamily \mnem{TEX\_SINGLE}\imod{.tex2d.computed.skip} \dst{@r0:r3}, \src{@r4:r5}} (\textbf{SPIR-V:} \texttt{OpImageSampleImplicitLod (after OpDemoteToHelperInvocation)}; \textbf{GLSL:} \texttt{if (a \textless{} 0.0) discard; vec4 c = texture(s, uv);})
  \item \texttt{staging\_register\_index}: Base GPR index of the contiguous staging register window (\texttt{@rN:rM}) for message payload/return. \emph{Example:} {\scriptsize\ttfamily \mnem{STORE.i64}\imod{.slot7} \src{@r0:r1}, \src{r2}, \imm{\#8}} (\textbf{SPIR-V:} \texttt{OpStore (uvec2)}; \textbf{GLSL:} \texttt{ssbo.pairs[i] = uvec2(a, b);})
  \item \texttt{staging\_register\_control\_32}: Controls the read/write staging register bank and count for async message instructions. \emph{Example:} {\scriptsize\ttfamily \mnem{LOAD.i128}\imod{.slot0} \dst{@r0:r3}, \src{r4}, \imm{\#0}} (\textbf{SPIR-V:} \texttt{OpLoad (vec4)}; \textbf{GLSL:} \texttt{vec4 v = ssbo.data[idx];})
  \item \texttt{opcode1}: Primary sub-opcode within the instruction class uniquely identifying the operation family. \emph{Example:} {\scriptsize\ttfamily \mnem{IADD.s32} \dst{r0}, \src{r1}, \src{r2}} (\textbf{SPIR-V:} \texttt{OpIAdd (operation selector)}; \textbf{GLSL:} \texttt{// Primary opcode bits})
  \item \texttt{instruction\_class}: 4-bit major instruction class selector (\texttt{bits 48:51}) partitioning the 64-bit Valhall opcode space (e.g., \texttt{FMA}, \texttt{CVT}, \texttt{SFU}, \texttt{LS}, \texttt{TEX}, \texttt{BRANCH}). \emph{Example:} {\scriptsize\ttfamily \mnem{FADD.f32} \dst{r0}, \src{r1}, \src{r2}} (\textbf{SPIR-V:} \texttt{N/A (encoding class)}; \textbf{GLSL:} \texttt{// Hardware execution pipe / decode class})
  \item \texttt{fau\_page\_index}: Selects the FAU (Fast Access Uniform) page (\texttt{0..3}) when a source operand references hardware small constants, thread IDs, or lane constants. \emph{Example:} {\scriptsize\ttfamily \mnem{FADD.f32} \dst{r0}, \src{r1}, \src{k1f}} (\textbf{SPIR-V:} \texttt{OpConstant / BuiltIn}; \textbf{GLSL:} \texttt{float y = x + 1.0;})
  \item \hyperref[enum:flow_control_m]{\texttt{flow}}: 4-bit embedded flow-control \& scoreboard wait modifier (\texttt{bits 59:62}: \texttt{none}, \texttt{wait0..wait0126}, \texttt{wait}, \texttt{reconverge}, \texttt{discard}, \texttt{end}). \emph{Example:} {\scriptsize\ttfamily \mnem{FADD.f32}\imod{.wait0.reconverge} \dst{r0}, \src{r1}, \src{r2}} (\textbf{SPIR-V:} \texttt{OpControlBarrier / OpReturn}; \textbf{GLSL:} \texttt{// Scoreboard wait + warp reconvergence})
\end{tightitemize}
\clearpage

\subsection*{VAR\_TEX\_SINGLE}
\addcontentsline{toc}{subsection}{VAR\_TEX\_SINGLE}
Fuses descriptor-based varying UV interpolation (via varying handle \texttt{src2}) with a filtered texture sample (\texttt{TEX\_SINGLE} via texture handle \texttt{src0}), writing the sampled texel into \texttt{@sw} in a single message.
\par\textbf{Notes:}
\begin{tightitemize}
  \item Fragment-only.
  \item FP32 varyings (\texttt{ISA.xml:169}).
  \item \textbf{Source references:} \texttt{kraid/ops.rs:5036-5083}, \texttt{encode\_v9.rs:3799-3829}, \texttt{bifrost/valhall/ISA.xml:169}.
\end{tightitemize}
\codebox{\mnem{VAR\_TEX\_SINGLE}\imod{\{.message\_slot\_index\}}\imod{\{.flow\}}\imod{\{.skip\}}\imod{\{.sample\_position\}}\imod{\{.update\}}\imod{\{.dimensionality\}}\imod{\{.register\_width\}}\imod{\{.write\_mask\}}\imod{\{.wide\_indices\}}\imod{\{.array\_enable\}}\imod{\{.compare\_enable\}}\imod{\{.lod\_mode\}} \dst{@sw}, \src{src0}, \src{src2}\\\mnem{VAR\_TEX\_SINGLE}\imod{\{.message\_slot\_index\}}\imod{\{.flow\}}\imod{\{.skip\}}\imod{\{.sample\_position\}}\imod{\{.update\}}\imod{\{.dimensionality\}}\imod{\{.projection\_enable\}}\imod{\{.register\_width\}}\imod{\{.write\_mask\}}\imod{\{.wide\_indices\}}\imod{\{.array\_enable\}}\imod{\{.compare\_enable\}}\imod{\{.lod\_mode\}} \dst{@sw}, \src{src0}, \src{src2}}
\textbf{Operands:}
\begin{tightitemize}\footnotesize
  \item \texttt{\dst{@sw}} (\texttt{staging\_register\_index} + \texttt{sr\_write\_count}): Destination staging register range receiving the single texture sample result.
  \item \texttt{\src{src0}}: Texture/sampler descriptor handle.
  \item \texttt{\src{src2}}: Varying descriptor handle providing interpolated UV coordinates (encoded in \texttt{src2}).
\end{tightitemize}
\par\textbf{Closest SPIR-V Op:} \texttt{OpLoad (in) + OpImageSampleImplicitLod}
\par\textbf{Pseudo-GLSL (Operational Semantics):}
\codebox{@sw = textureSample(handle = src0, interpolateVarying(handle = src2), dim, flag\_project, flag\_compare);}
\textbf{Execution Context:}
\begin{tightitemize}\footnotesize
  \item \textbf{Varying buffer \& barycentric cache (\hyperref[sec:csf-tiler]{\texttt{IDVSSR::VARY\_SIZE}}, \texttt{r48:r49}):} Reads per-vertex varying RAM written by Varying EBIN \texttt{STORE.estream}; descriptor forms (\texttt{LD\_VAR*}) also chase \hyperref[sec:csf-srt]{\texttt{SRT\_2}} $\to$ \texttt{Resource[table]} $\to$ \texttt{Attribute[index]} and reuse cached barycentric weights via \texttt{.store}/\texttt{.retrieve}.
  \item \textbf{Texture \& Sampler chain (\hyperref[sec:csf-srt]{\texttt{SRT\_2}} Frag / \texttt{SRT\_0} Comp/Vert):} 64B-aligned VA (\texttt{\& \textasciitilde{}0x3F}) $\to$ \texttt{Resource[table]} $\to$ \hyperref[sec:csf-tex-samp]{\texttt{Texture[tex\_idx]}} (32B, \texttt{type=2}) $\to$ \texttt{GenericPlane[mip]} (32B, \texttt{type=10}: \texttt{pointer}, \texttt{row\_stride}, Morton 16$\times$16 / AFBC) + \texttt{Sampler[samp\_idx]} (32B, \texttt{type=1}: filter, wrap, LOD clamp/bias, compare).
  \item \textbf{Fragment stage preloads \& FP modes (\hyperref[sec:csf-spd]{\texttt{SPD\_2}}, \texttt{r20:r21}):} Hardware pre-fills \texttt{r59} with \texttt{floatBitsToUint(FragCoord.z)} and \texttt{r60} with the MSAA sample coverage mask (\texttt{gl\_SampleMaskIn[0]}), while \texttt{ShaderProgram} word~0 latches 32/64 GPR allocation and warp-wide \texttt{flush\_to\_zero\_mode}, \texttt{suppress\_na\_n}, and \texttt{suppress\_inf}.
\end{tightitemize}
\par\textbf{Worked Examples:}
\par\textbf{Example 1 (Descriptor varying + implicit-LOD 2D sample):}
\codebox{\mnem{VAR\_TEX\_SINGLE}\imod{.tex2d.computed.center.slot0} \dst{@r0:r3}, \src{u0}, \src{r59}, \src{u1}}
Fuses descriptor-based \texttt{vUV} interpolation (\texttt{u1}) directly into \texttt{TEX\_SINGLE} (\texttt{u0}), returning RGBA in \texttt{@r0:r3}.
\par\textbf{Example 2 (Arrayed cubemap / 2D-array fused sample (`.array`)):}
\codebox{\mnem{VAR\_TEX\_SINGLE}\imod{.tex2d.array.computed.center.slot0} \dst{@r0:r3}, \src{u0}, \src{r59}, \src{u1}}
Interpolates \texttt{(u, v, layer)} from varying \texttt{u1} and samples a \texttt{sampler2DArray} in a single message.

\begin{center}\footnotesize
\begin{tabularx}{\linewidth}{@{}lllX@{}}
\toprule
Variant & Arch & Class & Opcode Fields \\
\midrule
\texttt{(none)} & v9--v14 & \texttt{S3A} & \texttt{opcode1=VAR\_TEX\_SINGLE} \\
\texttt{(none)} & v9--v14 & \texttt{S3A} & \texttt{opcode1=VAR\_TEX\_SINGLE} \\
\bottomrule
\end{tabularx}
\end{center}
\begin{center}\footnotesize
\begin{tabularx}{0.88\linewidth}{@{}l l X@{}}
\toprule
Field Name & Bits & Type / Fixed Value \\
\midrule
\texttt{src0} & \texttt{0:7} & \texttt{SourceEncoding64} \\
\hyperref[enum:tex_wide_indices_enable_m]{\texttt{wide\_indices}} & \texttt{8} & \texttt{tex\_wide\_indices\_enable\_m (def none)} \\
\textit{Reserved} & \texttt{9} & \texttt{zero} \\
\hyperref[enum:var_tex_array_enable_m]{\texttt{array\_enable}} & \texttt{10} & \texttt{var\_tex\_array\_enable\_m (def none)} \\
\hyperref[enum:tex_texel_offset_m]{\texttt{texel\_offset}} & \texttt{11} & \texttt{exact=none} \\
\hyperref[enum:var_tex_compare_enable_m]{\texttt{compare\_enable}} & \texttt{12} & \texttt{var\_tex\_compare\_enable\_m (def none)} \\
\hyperref[enum:var_tex_lod_mode_m]{\texttt{lod\_mode}} & \texttt{13:15} & \texttt{var\_tex\_lod\_mode\_m} \\
\texttt{src2} & \texttt{16:23} & \texttt{SourceEncoding} \\
\hyperref[enum:tex_width_m]{\texttt{register\_width}} & \texttt{24} & \texttt{tex\_width\_m} \\
\texttt{secondary\_width} & \texttt{25} & \texttt{exact=0} \\
\textit{Reserved} & \texttt{26:27} & \texttt{ignore} \\
\hyperref[enum:var_tex_dimensionality_m]{\texttt{dimensionality}} & \texttt{28:29} & \texttt{var\_tex\_dimensionality\_m} \\
\hyperref[enum:message_slot_index_m]{\texttt{message\_slot\_index}} & \texttt{30:32} & \texttt{message\_slot\_index\_m} \\
\hyperref[enum:encoded_sample_and_update_m]{\texttt{encoded\_sample\_and\_update}} & \texttt{33:35} & \texttt{encoded\_sample\_and\_update\_m} \\
\texttt{sr\_write\_count} & \texttt{36:38} & \texttt{sr\_write\_count\_t} \\
\hyperref[enum:tex_skip_m]{\texttt{skip}} & \texttt{39} & \texttt{tex\_skip\_m (def none)} \\
\texttt{staging\_register\_index} & \texttt{40:45} & \texttt{uint} \\
\texttt{staging\_register\_control\_32} & \texttt{46:47} & \texttt{exact=write} \\
\texttt{opcode1} & \texttt{48:51} & \texttt{exact=VAR\_TEX\_SINGLE} \\
\texttt{instruction\_class} & \texttt{52:56} & \texttt{exact=S3A} \\
\texttt{fau\_page\_index} & \texttt{57:58} & \texttt{uint} \\
\hyperref[enum:flow_control_m]{\texttt{flow}} & \texttt{59:62} & \texttt{flow\_control\_m} \\
\textit{Reserved} & \texttt{63} & \texttt{zero} \\
\bottomrule
\end{tabularx}
\end{center}
\textbf{Fields Explained:}
\begin{tightitemize}\footnotesize
  \item \texttt{src0}: Primary source operand (\texttt{r0..r63}, \texttt{u0..u63} uniform, \texttt{k*} FAU constant, or \texttt{t0..t2} temporary; \texttt{\textasciicircum{}} marks discard/last-use). \emph{Example:} {\scriptsize\ttfamily \mnem{FADD.f32} \dst{r0}, \src{r1}, \src{r2}} (\textbf{SPIR-V:} \texttt{OpFAdd (Operand 0)}; \textbf{GLSL:} \texttt{float c = a + b;})
  \item \hyperref[enum:tex_wide_indices_enable_m]{\texttt{wide\_indices}}: Uses 32-bit wide texture/sampler table indices instead of compact 16-bit packed indices. \emph{Example:} {\scriptsize\ttfamily \mnem{TEX\_SINGLE}\imod{.tex2d.wide} \dst{@r0:r3}, \src{@r4:r5}, \src{r6}} (\textbf{SPIR-V:} \texttt{OpImageSampleImplicitLod (bindless)}; \textbf{GLSL:} \texttt{vec4 c = texture(textures[nonuniformEXT(i)], uv);})
  \item \hyperref[enum:var_tex_array_enable_m]{\texttt{array\_enable}}: Enables array-layer coordinate in texture staging registers (\texttt{sampler2DArray} / \texttt{samplerCubeArray}). \emph{Example:} {\scriptsize\ttfamily \mnem{TEX\_SINGLE}\imod{.tex2d.array} \dst{@r0:r3}, \src{@r4:r6}} (\textbf{SPIR-V:} \texttt{OpImageSampleImplicitLod (Arrayed)}; \textbf{GLSL:} \texttt{vec4 c = texture(texArr, vec3(uv, layer));})
  \item \hyperref[enum:tex_texel_offset_m]{\texttt{texel\_offset}}: Applies a signed 4-bit per-axis texel offset $(dx, dy, dz)$ packed in staging registers. \emph{Example:} {\scriptsize\ttfamily \mnem{TEX\_SINGLE}\imod{.tex2d.texel\_offset} \dst{@r0:r3}, \src{@r4:r6}} (\textbf{SPIR-V:} \texttt{OpImageSampleImplicitLod (ConstOffset)}; \textbf{GLSL:} \texttt{vec4 c = textureOffset(tex, uv, ivec2(1, -1));})
  \item \hyperref[enum:var_tex_compare_enable_m]{\texttt{compare\_enable}}: Enables depth-reference shadow comparison (\texttt{Dref} in staging registers). \emph{Example:} {\scriptsize\ttfamily \mnem{TEX\_SINGLE}\imod{.tex2d.compare} \dst{@r0}, \src{@r4:r6}} (\textbf{SPIR-V:} \texttt{OpImageSampleDrefImplicitLod}; \textbf{GLSL:} \texttt{float s = texture(shadowMap, vec3(uv, refZ));})
  \item \hyperref[enum:var_tex_lod_mode_m]{\texttt{lod\_mode}}: Selects mipmap LOD calculation mode (\texttt{none}, \texttt{computed}, \texttt{explicit}, \texttt{computed\_bias}, \texttt{grdesc}). \emph{Example:} {\scriptsize\ttfamily \mnem{TEX\_SINGLE}\imod{.tex2d.explicit} \dst{@r0:r3}, \src{@r4:r6}} (\textbf{SPIR-V:} \texttt{OpImageSampleExplicitLod}; \textbf{GLSL:} \texttt{vec4 c = textureLod(tex, uv, 2.0);})
  \item \texttt{src2}: Third source operand (FMA addend, \texttt{MUX}/\texttt{CSEL} selector, \texttt{IADDC} carry-in, or bitwise mask in fused shift-logic). \emph{Example:} {\scriptsize\ttfamily \mnem{FMA.f32} \dst{r0}, \src{r1}, \src{r2}, \src{r3}} (\textbf{SPIR-V:} \texttt{GLSL.std.450 Fma (Addend)}; \textbf{GLSL:} \texttt{float d = fma(a, b, c);})
  \item \hyperref[enum:tex_width_m]{\texttt{register\_width}}: Selects 16-bit vs 32-bit GPR element width for the corresponding message return window. \emph{Example:} {\scriptsize\ttfamily \mnem{TEX\_DUAL}\imod{.primary\_f16.secondary\_f32} \dst{@r0:r2}, \src{@r4:r5}} (\textbf{SPIR-V:} \texttt{OpImageSampleImplicitLod (RelaxedPrecision)}; \textbf{GLSL:} \texttt{mediump vec2 a = texture(t0, uv).rg;})
  \item \texttt{secondary\_width}: Selects 16-bit vs 32-bit GPR element width for the corresponding message return window. \emph{Example:} {\scriptsize\ttfamily \mnem{TEX\_DUAL}\imod{.primary\_f16.secondary\_f32} \dst{@r0:r2}, \src{@r4:r5}} (\textbf{SPIR-V:} \texttt{OpImageSampleImplicitLod (RelaxedPrecision)}; \textbf{GLSL:} \texttt{mediump vec2 a = texture(t0, uv).rg;})
  \item \hyperref[enum:var_tex_dimensionality_m]{\texttt{dimensionality}}: Texture dimensionality (\texttt{tex1d}, \texttt{tex2d}, \texttt{tex3d}, \texttt{cube}). \emph{Example:} {\scriptsize\ttfamily \mnem{TEX\_SINGLE}\imod{.tex3d.f32.rgba} \dst{@r0:r3}, \src{@r4:r6}} (\textbf{SPIR-V:} \texttt{OpImageSampleImplicitLod (Dim=3D)}; \textbf{GLSL:} \texttt{vec4 c = texture(volTex, uvw);})
  \item \hyperref[enum:message_slot_index_m]{\texttt{message\_slot\_index}}: Asynchronous scoreboard slot (\texttt{slot0}--\texttt{slot2}, \texttt{slot7}) tracking completion of this message. \emph{Example:} {\scriptsize\ttfamily \mnem{LOAD.i32}\imod{.slot0} \dst{@r0}, \src{r2}, \imm{\#0}} (\textbf{SPIR-V:} \texttt{OpLoad}; \textbf{GLSL:} \texttt{uint v = ssbo.arr[i];})
  \item \hyperref[enum:encoded_sample_and_update_m]{\texttt{encoded\_sample\_and\_update}}: Jointly encodes MSAA sample position and hardware varying-cache update mode (\texttt{center\_store}, \texttt{retrieve}, etc.). \emph{Example:} {\scriptsize\ttfamily \mnem{VAR\_TEX\_BUF\_SINGLE}\imod{.center\_store.tex2d} \dst{@r0:r3}, \src{r4}} (\textbf{SPIR-V:} \texttt{OpImageSampleImplicitLod (fused varying)}; \textbf{GLSL:} \texttt{vec4 c = texture(s, v\_uv);})
  \item \texttt{sr\_write\_count}: Number of consecutive 32-bit staging registers written back by a dual-staging message (e.g. ATOM\textbackslash\{\}\_RETURN / TEX). \emph{Example:} {\scriptsize\ttfamily \mnem{TEX\_SINGLE}\imod{.tex2d.f32.rgba} \dst{@r0:r3}, \src{@r4:r5}} (\textbf{SPIR-V:} \texttt{OpImageSampleImplicitLod}; \textbf{GLSL:} \texttt{vec4 c = texture(tex, uv);})
  \item \hyperref[enum:tex_skip_m]{\texttt{skip}}: Skips texture/memory fetch on helper/discarded fragment lanes while preserving 2x2 quad derivatives. \emph{Example:} {\scriptsize\ttfamily \mnem{TEX\_SINGLE}\imod{.tex2d.computed.skip} \dst{@r0:r3}, \src{@r4:r5}} (\textbf{SPIR-V:} \texttt{OpImageSampleImplicitLod (after OpDemoteToHelperInvocation)}; \textbf{GLSL:} \texttt{if (a \textless{} 0.0) discard; vec4 c = texture(s, uv);})
  \item \texttt{staging\_register\_index}: Base GPR index of the contiguous staging register window (\texttt{@rN:rM}) for message payload/return. \emph{Example:} {\scriptsize\ttfamily \mnem{STORE.i64}\imod{.slot7} \src{@r0:r1}, \src{r2}, \imm{\#8}} (\textbf{SPIR-V:} \texttt{OpStore (uvec2)}; \textbf{GLSL:} \texttt{ssbo.pairs[i] = uvec2(a, b);})
  \item \texttt{staging\_register\_control\_32}: Controls the read/write staging register bank and count for async message instructions. \emph{Example:} {\scriptsize\ttfamily \mnem{LOAD.i128}\imod{.slot0} \dst{@r0:r3}, \src{r4}, \imm{\#0}} (\textbf{SPIR-V:} \texttt{OpLoad (vec4)}; \textbf{GLSL:} \texttt{vec4 v = ssbo.data[idx];})
  \item \texttt{opcode1}: Primary sub-opcode within the instruction class uniquely identifying the operation family. \emph{Example:} {\scriptsize\ttfamily \mnem{IADD.s32} \dst{r0}, \src{r1}, \src{r2}} (\textbf{SPIR-V:} \texttt{OpIAdd (operation selector)}; \textbf{GLSL:} \texttt{// Primary opcode bits})
  \item \texttt{instruction\_class}: 4-bit major instruction class selector (\texttt{bits 48:51}) partitioning the 64-bit Valhall opcode space (e.g., \texttt{FMA}, \texttt{CVT}, \texttt{SFU}, \texttt{LS}, \texttt{TEX}, \texttt{BRANCH}). \emph{Example:} {\scriptsize\ttfamily \mnem{FADD.f32} \dst{r0}, \src{r1}, \src{r2}} (\textbf{SPIR-V:} \texttt{N/A (encoding class)}; \textbf{GLSL:} \texttt{// Hardware execution pipe / decode class})
  \item \texttt{fau\_page\_index}: Selects the FAU (Fast Access Uniform) page (\texttt{0..3}) when a source operand references hardware small constants, thread IDs, or lane constants. \emph{Example:} {\scriptsize\ttfamily \mnem{FADD.f32} \dst{r0}, \src{r1}, \src{k1f}} (\textbf{SPIR-V:} \texttt{OpConstant / BuiltIn}; \textbf{GLSL:} \texttt{float y = x + 1.0;})
  \item \hyperref[enum:flow_control_m]{\texttt{flow}}: 4-bit embedded flow-control \& scoreboard wait modifier (\texttt{bits 59:62}: \texttt{none}, \texttt{wait0..wait0126}, \texttt{wait}, \texttt{reconverge}, \texttt{discard}, \texttt{end}). \emph{Example:} {\scriptsize\ttfamily \mnem{FADD.f32}\imod{.wait0.reconverge} \dst{r0}, \src{r1}, \src{r2}} (\textbf{SPIR-V:} \texttt{OpControlBarrier / OpReturn}; \textbf{GLSL:} \texttt{// Scoreboard wait + warp reconvergence})
\end{tightitemize}
\clearpage



% ============================================================================
% Appendix A: Concrete Decompiled libGLES_mali.so vs Mesa Panfrost Differences
% ============================================================================
\appendix
\chapter[DDK vs Panfrost Differences]{Arm \texttt{libGLES\_mali.so} vs Mesa Panfrost: Concrete Decompiled Differences}

Comparing Arm's proprietary LLVM Valhall backend in \texttt{libGLES\_mali.so} (Pixel 8 r54p3, 2,265 \texttt{XMLInstrInfo} entries) against Mesa Panfrost (\texttt{kraid}/\texttt{bifrost}, \texttt{isa-v9-v14.xml} with 627 specs) shows 100\% Valhall v9--v11 parity (374/374 specs), 45 v12--v14 Mesa-only specs, 12 real hardware DDK-only opcodes, and 12 concrete behavioral/modifier differences. Every entry below provides exact AArch64 disassembly from \texttt{libGLES\_mali.so} \emph{and} C/pseudocode decompilation alongside Mesa Panfrost.

\section{Hardware Opcodes in \texttt{libGLES\_mali.so} Absent from \texttt{isa-v9-v14.xml}}

\begin{table}[!htbp]
\centering\footnotesize
\begin{tabularx}{\linewidth}{@{}l>{\raggedright\arraybackslash}X@{}}
\toprule
Group & Mnemonics in \texttt{libGLES\_mali.so} (\texttt{ValhallInsts} \texttt{0x300fd18}) \\
\midrule
Ray Tracing (3) & \texttt{RT\_RAY\_BOX} (\texttt{0x319}), \texttt{RT\_RAY\_TRI} (\texttt{0x31a}), \texttt{RT\_RAY\_TRI\_CO.i\{32,64,96,128\}\_i288} (\texttt{0x805..0x808}) \\
Geometry / Viewport (3) & \texttt{VIEWPORT\_TRANSFORM} (\texttt{0xaa5}), \texttt{VIEWPORT\_TRANSFORM\_W1} (\texttt{0xaa6}), \texttt{VIEWPORT\_TRANSFORM\_Z0W1} (\texttt{0xaa7}) \\
System / Sync / FAU (5) & \texttt{FAU\_STORE.i\{8..128\}} (\texttt{0x450..0x459}), \texttt{FENCE} (\texttt{0x46f}), \texttt{CONDITIONAL\_WAIT} (\texttt{0x40f}), \texttt{ENTER\_SIMD\_STATE} (\texttt{0x210}), \texttt{ENTER\_SIMD\_STATE\_RQ\_LOOP} (\texttt{0x42f}), \texttt{EXIT\_SIMD\_STATE} (\texttt{0x212}) \\
Bitwise Variant (1) & \texttt{AND\_IMM.i16} (\texttt{0x35c}) \\
\bottomrule
\end{tabularx}
\end{table}

\subsubsection*{A.1.1 \texttt{RT\_RAY\_BOX} (\texttt{0x319}, \texttt{xml=360}, \texttt{ops=6, defs=1})}
\textbf{Usage \& Induced Semantics:} Hardware Ray--AABB slab intersection unit (v11+ Immortalis ray-query/ray-tracing). Takes a staging register window containing ray origin, inverse direction, $t_{\min}/t_{\max}$, and AABB bounds (with \texttt{S1.zext} on the descriptor/node offset) and asynchronously writes hit mask / sorted child node distances to destination staging registers on \texttt{slot0..slot2}.\\
\textbf{AArch64 Disassembly (\texttt{ValhallHoistLargeRegSequence::runOnMachineFunction} @ \texttt{0x1ea9a5c}):}
\codebox{%
0x01ea9a5c: ldr   x8, [x25, \#0x18] ; MI->getDesc()\\
0x01ea9a60: ldrh  w8, [x8]          ; Opcode = MI->getOpcode()\\
0x01ea9a64: cmp   w8, \#0x319       ; RT\_RAY\_BOX\\
0x01ea9a68: b.ne  \#0x1eab08c\\
0x01ea9a6c: ldr   x8, [x25, \#0x28] ; MI->getOperand(1) (staging RegSequence)\\
0x01ea9a70: ldr   x0, [x19, \#0x100]\\
0x01ea9a74: ldr   w1, [x8, \#0x24]  ; src staging vreg\\
0x01ea9a78: bl    \#0x27c6a30        ; MRI->getVRegDef(vreg) for hoisting%
}
\textbf{C/Pseudocode Decompilation (Compiler \& HW Semantics):}
\codebox{%
// Compiler pass (ValhallHoistLargeRegSequence):\\
if (MI.getOpcode() == Valhall::RT\_RAY\_BOX) \{\\
~~MachineInstr *Def = MRI.getVRegDef(MI.getOperand(1).getReg());\\
~~hoistStagingRegSequence(Def, MI); // minimize live-range of wide ray payload\\
\}\\
// Induced HW Semantics (async on slot0..2):\\
vec3 t0 = (box\_min - ray\_orig) * ray\_inv\_dir;\\
vec3 t1 = (box\_max - ray\_orig) * ray\_inv\_dir;\\
float t\_enter = max(ray\_tmin, max(min(t0.x,t1.x), max(min(t0.y,t1.y), min(t0.z,t1.z))));\\
float t\_exit  = min(ray\_tmax, min(max(t0.x,t1.x), min(max(t0.y,t1.y), max(t0.z,t1.z))));\\
dst\_sr = (t\_enter <= t\_exit) ? pack\_hit(t\_enter) : MISS\_SENTINEL;%
}

\subsubsection*{A.1.2 \texttt{RT\_RAY\_TRI} (\texttt{0x31a}, \texttt{xml=361}, \texttt{ops=8, defs=1})}
\textbf{Usage \& Induced Semantics:} Hardware watertight Ray--Triangle intersection unit. \texttt{XMLInstrInfo} \texttt{ModInfo} exposes \texttt{S4 = [slot0..slot2, t, uv, flow\_control]}, selecting whether the async return window delivers hit distance $t$ (\texttt{.t}) or barycentric coordinates $(u, v)$ (\texttt{.uv}).\\
\textbf{AArch64 Disassembly (\texttt{XMLInstrInfo} \texttt{0x4b1a28 + 361*0x48} \& \texttt{ValhallShrinkRayTri} @ \texttt{0x1f40ed0}):}
\codebox{%
// XMLInstrInfo[361]: q1=0x000000003c000000, S1=[zext], S4=[slot0..2, t, uv, flow]\\
0x01f40ecc: ldr   x8, [x28, \#0x18] ; MI->getDesc()\\
0x01f40ed0: ldrh  w8, [x8]          ; inspect RT\_RAY\_TRI / RT\_RAY\_TRI\_CO opcode\\
0x01f40ed4: cmp   w8, \#0x805       ; RT\_RAY\_TRI\_CO.i128\_i288%
}
\textbf{C/Pseudocode Decompilation:}
\codebox{%
// Induced HW Semantics (async on slot0..2, Watertight Moller-Trumbore):\\
vec3 e1 = v1 - v0, e2 = v2 - v0, p = cross(ray\_dir, e2);\\
float det = dot(e1, p);\\
float u = dot(ray\_orig - v0, p) / det;\\
float v = dot(ray\_dir, cross(ray\_orig - v0, e1)) / det;\\
float t = dot(e2, cross(ray\_orig - v0, e1)) / det;\\
bool hit = (u >= 0 \&\& v >= 0 \&\& u + v <= 1 \&\& t >= ray\_tmin \&\& t <= ray\_tmax);\\
dst\_sr = mode\_uv ? pack(hit, u, v) : pack(hit, t);%
}

\subsubsection*{A.1.3 \texttt{RT\_RAY\_TRI\_CO.i\{32,64,96,128\}\_i288} (\texttt{0x805..0x808}, \texttt{ops=8, defs=1})}
\textbf{Usage \& Induced Semantics:} Compact-output Ray--Triangle intersection taking a 288-bit (9-GPR) ray+triangle input window and writing a shrunken 32/64/96/128-bit result window. \texttt{ValhallShrinkRayTri::runOnMachineFunction} (\texttt{0x1f40d60}) analyzes live destination subregisters of \texttt{RT\_RAY\_TRI\_CO.i128\_i288} (\texttt{0x805}) and shrinks the opcode to \texttt{i96}, \texttt{i64}, or \texttt{i32} when upper output components are dead.\\
\textbf{AArch64 Disassembly (\texttt{ValhallShrinkRayTri::runOnMachineFunction} @ \texttt{0x1f40ecc}):}
\codebox{%
0x01f40ecc: ldr   x8, [x28, \#0x18] ; MI->getDesc()\\
0x01f40ed0: ldrh  w8, [x8]          ; MI->getOpcode()\\
0x01f40ed4: cmp   w8, \#0x805       ; RT\_RAY\_TRI\_CO.i128\_i288\\
0x01f40ed8: b.ne  \#0x1f41844        ; skip if not 128-bit output variant\\
0x01f40edc: ldr   x8, [x28, \#0x28] ; operand array\\
0x01f40ee8: ldr   w23, [x8, \#4]    ; dst staging super-register (4x32b)%
}
\textbf{C/Pseudocode Decompilation:}
\codebox{%
if (MI.getOpcode() == Valhall::RT\_RAY\_TRI\_CO\_i128\_i288) \{\\
~~unsigned dstReg = MI.getOperand(0).getReg();\\
~~unsigned liveWords = countLiveSubregs(MRI, dstReg); // 1..4 GPRs\\
~~static const unsigned shrunk[4] = \{0x806, 0x807, 0x808, 0x805\}; // i32..i128\\
~~if (liveWords < 4) MI.setDesc(TII.get(shrunk[liveWords - 1]));\\
\}%
}

\subsubsection*{A.1.4 \texttt{VIEWPORT\_TRANSFORM}, \texttt{\_W1}, \texttt{\_Z0W1} (\texttt{0xaa5..0xaa7}, \texttt{ops=7, defs=0})}
\textbf{Usage \& Induced Semantics:} Post-RA vertex-position viewport transform instructions lowered by \texttt{ValhallExpandViewportTransform::runOnMachineFunction} (\texttt{0x1e66970}) and recognized in \texttt{ValhallTTIImpl::instCombineIntrinsic} (\texttt{0x1f67150}). Transforms clip-space $(x, y, z, w)$ to window-space coordinates via perspective divide $1/w$ and viewport scale/bias; \texttt{\_W1} specializes $w = 1.0$ (elides \texttt{FRCP.f32}), and \texttt{\_Z0W1} specializes $z = 0.0, w = 1.0$ (2D/UI fast path).\\
\textbf{AArch64 Disassembly (\texttt{ValhallExpandViewportTransform} @ \texttt{0x1e669c4} \& \texttt{instCombineIntrinsic} @ \texttt{0x1f67150}):}
\codebox{%
0x01e669c4: ldrb  w8, [x8, \#0x17] ; subtarget/shader stage flags\\
0x01e669c8: tbnz  w8, \#5, \#0x1e66a50 ; skip if viewport transform disabled\\
0x01e669cc: ldr   w8, [x20, \#0x10] ; ShaderType == VERTEX (1)\\
0x01e669d0: cmp   w8, \#1 ; b.ne \#0x1e66a50\\
// instCombineIntrinsic @ 0x1f67150:\\
0x01f67150: mov   w9, \#0xaa7      ; VIEWPORT\_TRANSFORM\_Z0W1\\
0x01f6715c: str   w9, [sp]         ; record specialized intrinsic opcode%
}
\textbf{C/Pseudocode Decompilation:}
\codebox{%
// Induced Semantics (expanded post-RA for Vertex/Position shaders):\\
float inv\_w = (opc == VIEWPORT\_TRANSFORM) ? frcp\_f32(clip.w) : 1.0f;\\
float ndc\_x = clip.x * inv\_w, ndc\_y = clip.y * inv\_w;\\
float ndc\_z = (opc == VIEWPORT\_TRANSFORM\_Z0W1) ? 0.0f : (clip.z * inv\_w);\\
win\_pos = vec4(fma(ndc\_x, vp\_scale.x, vp\_bias.x), fma(ndc\_y, vp\_scale.y, vp\_bias.y),\\
~~~~~~~~~~~~~~~fma(ndc\_z, vp\_scale.z, vp\_bias.z), inv\_w);%
}

\subsubsection*{A.1.5 \texttt{FAU\_STORE.i\{8,16,24,32,48,64,96,128\}} (\texttt{0x450..0x459}, \texttt{ops=7, defs=0})}
\textbf{Usage \& Induced Semantics:} Asynchronous store into the per-workgroup/pilot-shader FAU (Fast Access Uniform) RAM buffer. Emitted by \texttt{selectFauStore} (\texttt{0x1ed04c0}) during pilot-shader preloading so subsequent main-shader instructions can read computed uniforms directly via FAU operands without GPR pressure. Allows \texttt{slot0..slot2, slot7}.\\
\textbf{AArch64 Disassembly (\texttt{selectFauStore} @ \texttt{0x1ed04d4..0x1ed050c}):}
\codebox{%
0x01ed04d8: ldr   x21, [x1, \#0x28] ; MI->operands()\\
0x01ed04dc: mov   w1, w2            ; target FAU\_STORE.iN opcode (0x450..0x459)\\
0x01ed04e4: bl    \#0x1e65620        ; ValhallMIBuilder::build(Opc)\\
0x01ed04e8: add   x1, x21, \#0x20   ; src staging value operand\\
0x01ed04ec: bl    \#0x1e641c0        ; addOperand(val)\\
0x01ed04f0: add   x1, x21, \#0x40   ; FAU base/offset operand\\
0x01ed04f4: bl    \#0x1e641c0        ; addOperand(offset)\\
0x01ed04f8: mov   x1, x21 ; bl \#0x1e641c0 ; addOperand(base)%
}
\textbf{C/Pseudocode Decompilation:}
\codebox{%
MachineInstr *selectFauStore(ValhallMIBuilder \&B, MachineInstr \&MI, unsigned Opc) \{\\
~~auto NewMI = B.build(Opc).add(MI.getOperand(1)).add(MI.getOperand(2)).add(MI.getOperand(0));\\
~~return NewMI; // HW: FAU\_RAM[base + offset : +N\_bytes] = src\_staging\_regs (async slot)\\
\}%
}

\subsubsection*{A.1.6 \texttt{FENCE} (\texttt{0x46f}, \texttt{xml=833}, \texttt{ops=1, defs=0}) \& \texttt{CONDITIONAL\_WAIT} (\texttt{0x40f}, \texttt{xml=710}, \texttt{ops=5, defs=0})}
\textbf{Usage \& Induced Semantics:} \texttt{FENCE} is a standalone memory/pipeline ordering fence (accepting only \texttt{S4 = flow\_control}). \texttt{CONDITIONAL\_WAIT} is a predicated scoreboard wait (\texttt{S0, S1, S4 = flow\_control}) that stalls on the specified scoreboard slots only if its runtime predicate register is non-zero, avoiding unnecessary warp stalls on skipped async paths. Both are finalized together in \texttt{ValhallFinalize::runOnMachineFunction} (\texttt{0x1e71350}).\\
\textbf{AArch64 Disassembly (\texttt{ValhallFinalize::runOnMachineFunction} @ \texttt{0x1e71350}):}
\codebox{%
0x01e71350: ldr   x8, [x22, \#0x18] ; MI->getDesc()\\
0x01e71354: ldrh  w9, [x8]          ; MI->getOpcode()\\
0x01e71358: cmp   w9, \#0x46f       ; FENCE\\
0x01e7135c: b.eq  \#0x1e713c0        ; handle standalone fence\\
0x01e71360: cmp   w9, \#0x40f       ; CONDITIONAL\_WAIT\\
0x01e71364: b.ne  \#0x1e713f0%
}
\textbf{C/Pseudocode Decompilation:}
\codebox{%
if (MI.getOpcode() == Valhall::FENCE) \{\\
~~emitPipelineMemoryFence(MI.getFlowControl()); // HW: order prior loads/stores\\
\} else if (MI.getOpcode() == Valhall::CONDITIONAL\_WAIT) \{\\
~~// HW: if (any\_active\_lane(src0 != 0)) scoreboard\_wait(MI.wait\_mask);\\
~~finalizeConditionalWait(MI, MI.getOperand(0), MI.getOperand(1));\\
\}%
}

\subsubsection*{A.1.7 \texttt{ENTER\_SIMD\_STATE} (\texttt{0x210}), \texttt{\_RQ\_LOOP} (\texttt{0x42f}), \& \texttt{EXIT\_SIMD\_STATE} (\texttt{0x212})}
\textbf{Usage \& Induced Semantics:} Hardware SIMD execution-state management instructions used for call-stack / ray-query loop state entry and exit. \texttt{ModInfo} restricts \texttt{S4} strictly to \texttt{['reconverge']}. \texttt{ValhallSplitEnterSIMD} (\texttt{0x1f4db40}) splits ray-query loop regions around \texttt{ENTER\_SIMD\_STATE\_RQ\_LOOP} (\texttt{0x42f}), while \texttt{EmitInstrWithCustomInserter} (\texttt{0x1eccb8c}, \texttt{0x1ecebf0}) emits \texttt{EXIT\_SIMD\_STATE} (\texttt{0x212}) and \texttt{ENTER\_SIMD\_STATE\_RQ\_LOOP} (\texttt{0x42f}).\\
\textbf{AArch64 Disassembly (\texttt{EmitInstrWithCustomInserter} @ \texttt{0x1eccb88}, \texttt{0x1ecebf0} \& \texttt{ValhallSplitEnterSIMD} @ \texttt{0x1f4dbf0}):}
\codebox{%
0x01eccb88: add   x0, sp, \#0x280 ; ValhallMIBuilder\\
0x01eccb8c: mov   w1, \#0x212     ; EXIT\_SIMD\_STATE\\
0x01eccb90: bl    \#0x1e65620      ; build(EXIT\_SIMD\_STATE)\\
0x01ecebf0: mov   w1, \#0x42f     ; ENTER\_SIMD\_STATE\_RQ\_LOOP\\
0x01ecebf4: mov   w2, \#0x1e1     ; target register class\\
0x01ecebf8: bl    \#0x1e65f80      ; buildDef(ENTER\_SIMD\_STATE\_RQ\_LOOP)\\
0x01f4dbf4: cmp   w11, \#0x42f    ; ValhallSplitEnterSIMD detects RQ\_LOOP%
}
\textbf{C/Pseudocode Decompilation:}
\codebox{%
// HW Semantics (Ray-Query / Divergent SIMD State Stack):\\
uint32\_t saved\_mask = ENTER\_SIMD\_STATE(active\_lane\_mask); // push SIMD exec state\\
// ENTER\_SIMD\_STATE\_RQ\_LOOP: initialize HW ray-query traversal loop state\\
// ... divergent ray-query traversal body ...\\
EXIT\_SIMD\_STATE(saved\_mask); // pop/restore prior warp exec mask + reconverge%
}

\subsubsection*{A.1.8 \texttt{AND\_IMM.i16} (\texttt{0x35c}, \texttt{xml=449}, \texttt{ops=6, defs=1})}
\textbf{Usage \& Induced Semantics:} 16-bit bitwise AND with inline 16-bit immediate, synthesized by \texttt{ValhallMachineOptimizeNode::optimizeInstruction} (\texttt{0x1f12bb0}) when folding a 16-bit mask onto a halfword operand without allocating an FAU constant.\\
\textbf{AArch64 Disassembly (\texttt{ValhallMachineOptimizeNode::optimizeInstruction} @ \texttt{0x1f12bac}):}
\codebox{%
0x01f12bac: ldr   x8, [x20, \#0x18] ; MI->getDesc()\\
0x01f12bb0: mov   w23, \#0x35c      ; Valhall::AND\_IMM\_i16\\
0x01f12bb4: ldrh  w2, [x8]          ; original opcode\\
0x01f12bbc: add   x0, sp, \#0x110   ; builder context\\
0x01f12bc4: mov   w3, \#3           ; imm operand index%
}
\textbf{C/Pseudocode Decompilation:}
\codebox{%
// Compiler optimization fold (ValhallMachineOptimizeNode):\\
if (isMaskConstant16(MI, imm16)) \{\\
~~replaceWith(MI, Valhall::AND\_IMM\_i16, MI.getOperand(0), MI.getOperand(1), imm16);\\
\}\\
// Induced HW Semantics: dst.h0 = (uint16\_t)(src0.h0 \& imm16);%
}

\section{Concrete Decompiled Comparisons (DDK vs Panfrost)}

Table~\ref{tab:a2-ddk-panfrost-summary} summarizes all 12 concrete behavioral, encoding, and modifier differences identified between Arm's \texttt{libGLES\_mali.so} (r54p3) and Mesa Panfrost (\texttt{kraid}/\texttt{isa-v9-v14.xml}) before the decompiled deep dives below.

\begin{table}[!htbp]
\centering\footnotesize
\caption{Summary of Concrete DDK (\texttt{libGLES\_mali.so}) vs Mesa Panfrost Differences}
\label{tab:a2-ddk-panfrost-summary}
\begin{tabularx}{\linewidth}{@{}l l >{\raggedright\arraybackslash}p{0.23\linewidth} >{\raggedright\arraybackslash}p{0.23\linewidth} >{\raggedright\arraybackslash}X@{}}
\toprule
Ref & Instruction / Feature & Arm \texttt{libGLES\_mali.so} & Mesa Panfrost & Resolution / Impact \\
\midrule
A.2.1 & \texttt{FMIN}/\texttt{FMAX} & Deterministic $-0.0 < +0.0$; payload-preserving quietened NaN & Canonical \texttt{0x7FC00000} on \texttt{nan\_propagate} & Z3 model adopts DDK hardware semantics \\
A.2.2 & \texttt{FRCP}/\texttt{FRSQ.f32} & Shared input + output FTZ tail on both & Input+output FTZ on \texttt{FRCP}; input-only on \texttt{FRSQ} & Mathematically identical on normal inputs \\
A.2.3 & \texttt{*\_RSCALE} \texttt{.n\_add} & \texttt{FMA}/\texttt{FADD\_RSCALE} only (\texttt{ModId 0xe2}) & XML lists on all 3; \texttt{kraid} hardcodes \texttt{None} & Track allow-list per opcode \\
A.2.4 & \texttt{FREXPM.f32.log($\pm\infty$)} & Returns $\pm 1.0$ (\texttt{x \& 0xbf800000}) & Not constant-folded in \texttt{kraid} & Z3 model implements exact DDK formula \\
A.2.5 & \texttt{F32\_TO\_F16} dispatch & \texttt{V2F32\_TO\_V2F16} (v9--v10) vs \texttt{F32\_TO\_F16} (v11+) & Matches 100\% on arch $\le 10$ vs $\ge 11$ & Full parity verified \\
A.2.6 & \texttt{MKVEC.v2i8} \texttt{src2} & \texttt{ModInfo} allows only \texttt{b01} (rejects \texttt{b02}) & Reserved zero + virtual \texttt{exact="b01"} & Identical \texttt{Concat(src2[15:0], ...)} \\
A.2.7 & \texttt{LDEXP.f32} \texttt{widen0} & Permits \texttt{h0}/\texttt{h1} (\texttt{F136}) & Omitted in XML on v9--v13 & DDK supports FP16 widen on \texttt{src0} \\
A.2.7 & \texttt{MMUL} \texttt{sub0/1} \& \texttt{.sat} & \texttt{F77} on \texttt{f32}; \texttt{.sat} on v9--v11 \texttt{v4u8} & \texttt{sub0/1} on \texttt{v2f16/f16}; \texttt{.sat} v14-only & Minor table-gating delta \\
A.2.7 & \texttt{JUMP} \texttt{address\_mode} & Dynamic relative/absolute (\texttt{F37}) & Fixes \texttt{absolute(1)} in v11--v14 & Equivalent for absolute jumps \\
A.2.7 & \texttt{LOAD.extend} / \texttt{BLEND} & \texttt{sext/zext/none}; \texttt{BLEND} \texttt{v2..v4/none} & \texttt{encoded\_extend}; \texttt{BLEND} fixed \texttt{v4} & Compressed encoding vs full enum \\
A.2.7 & \texttt{LD\_VAR\_BUF*} / \texttt{TEX\_WEIGHTS} & \texttt{F49} flat mode; \texttt{project}, \texttt{4x1/1x4} & Split \texttt{LD\_VAR\_BUF\_FLAT*}; \texttt{2x2} only & Mesa splits flat op in v14 \\
A.2.7 & Unconditional \texttt{BRANCH} & \texttt{BRANCH k1n, \#off} & \texttt{BRANCH.not k0, \#off} & Identical always-taken semantics \\
\bottomrule
\end{tabularx}
\end{table}

\subsubsection*{A.2.1 \texttt{FMIN}/\texttt{FMAX} Signed-Zero Ordering \& \texttt{nan\_propagate} Payload Preservation}
\textbf{Arm \texttt{libGLES\_mali.so} (\texttt{\_mali\_max\_nan\_propagate\_sf32} @ \texttt{0x2e28300}) --- Disassembly:}
\codebox{%
0x02e28320: asr  w8, w0, \#0x1f ; sign-magnitude -> signed int key\\
0x02e28328: eor  w8, w0, w8, lsr \#1 ; -0.0 (0x80000000) -> -1 < +0.0 (0)\\
0x02e28330: cmp  w8, w9 ; csel w8, w0, w1, gt\\
0x02e28340: orr  w9, w1, \#0x400000 ; quieten winning NaN preserving sign \& payload!%
}
\textbf{Arm \texttt{libGLES\_mali.so} --- C/Pseudocode Decompilation:}
\codebox{%
uint32\_t \_mali\_max\_nan\_propagate\_sf32(uint32\_t a, uint32\_t b) \{\\
~~int32\_t ka = a \^ ((int32\_t)a >> 31 >> 1), kb = b \^ ((int32\_t)b >> 31 >> 1);\\
~~if (is\_nan(a) \&\& is\_nan(b)) return ((ka > kb) ? a : b) | 0x00400000u;\\
~~if (is\_nan(a)) return a | 0x00400000u;\\
~~if (is\_nan(b)) return b | 0x00400000u;\\
~~return (ka > kb) ? a : b; // -0.0 (-1) < +0.0 (0) deterministically\\
\}%
}
\textbf{Mesa Panfrost (\texttt{kraid/ops.rs:1892-1896, 1937-1941}):}
\codebox{%
if self.propagate\_nan \&\& (a.is\_nan() || b.is\_nan()) \{ f32::NAN \} // canonical 0x7FC00000%
}
\textbf{Impact:} \texttt{panfrost\_z3\_model.py} adopts DDK's exact hardware behavior: deterministic $-0.0 < +0.0$ and payload-preserving quietened NaN on \texttt{nan\_propagate}.

\subsubsection*{A.2.2 \texttt{FRCP.f32} / \texttt{FRSQ.f32} Output Subnormal Flushing}
\textbf{Arm \texttt{libGLES\_mali.so} (\texttt{ValhallInstrEval::compute} @ \texttt{0x2e24ba0..0x2e24c28}) --- Disassembly:}
\codebox{%
0x02e24ba8: and  w9, w8, \#0x80000000 ; tst w8, \#0x7f800000 ; csel w0, w9, w8, eq\\
0x02e24bb4: bl   eval\_rsqrt (0x2e2ae80) ; b 0x2e24c20\\
0x02e24c20: and  w8, w0, \#0x80000000 ; tst w0, \#0x7f800000 ; csel w0, w8, w0, eq%
}
\textbf{Arm \texttt{libGLES\_mali.so} --- C/Pseudocode Decompilation:}
\codebox{%
uint32\_t ftz(uint32\_t x) \{ return ((x \& 0x7f800000u) == 0) ? (x \& 0x80000000u) : x; \}\\
uint32\_t eval\_frsq\_or\_frcp(uint32\_t x, bool is\_rsq) \{\\
~~uint32\_t in = ftz(x);\\
~~uint32\_t r  = is\_rsq ? \_mali\_rsqrt\_sf32(in) : \_mali\_rcp\_sf32(in);\\
~~return ftz(r); // shared output-FTZ tail for both FRCP and FRSQ\\
\}%
}
\textbf{Mesa Panfrost (\texttt{kraid/ops.rs:1971-1974, 2101}):} Flushes subnormals on input+output for \texttt{FRCP.f32}, and input-only for \texttt{FRSQ.f32} (mathematically equivalent since $1/\sqrt{\text{normal}}$ never underflows FP32).

\subsubsection*{A.2.3 \texttt{FMA\_RSCALE} / \texttt{FADD\_RSCALE} \texttt{.special=n\_add} (\texttt{ModId 0xe2})}
\textbf{Arm \texttt{libGLES\_mali.so} (\texttt{insert\_FEXPM1\_32} @ \texttt{0x1ed6b30} \& \texttt{ModInfo} @ \texttt{0x4f42d8}) --- Disassembly:}
\codebox{%
0x01ed6b30: mov  w1, \#0xe2 ; ModId 0xe2 = n\_add\\
0x01ed6b34: bl   0x1e6abf0 ; attaches .n\_add to FADD\_RSCALE.f32%
}
\textbf{Arm \texttt{libGLES\_mali.so} --- C/Pseudocode Decompilation:}
\codebox{%
// In insert\_FEXPM1\_32():\\
MachineInstr *MI = build(Valhall::FADD\_RSCALE\_f32, dst, a, b, exp);\\
MI->addModifier(ModId::n\_add /* 0xe2 */);\\
// ModInfo allow-list: FMA\_RSCALE(0x234): {n, left, n\_add};\\
// FADD\_RSCALE(0x220): {n\_add} ONLY; FMUL\_RSCALE(0x23a): {n, left} ONLY.%
}
\textbf{Mesa Panfrost (\texttt{kraid/encode\_v9.rs:1518}):} \texttt{isa-v9-v14.xml} lists \texttt{n\_add} on all three opcodes, while \texttt{kraid} hardcodes \texttt{special: FmaRscaleSpecial32M::None}.

\subsubsection*{A.2.4 \texttt{FREXPM.f32.log($\pm\infty$)} Behavior}
\textbf{Arm \texttt{libGLES\_mali.so} (\texttt{\_mali\_frexpm\_log\_sf32} @ \texttt{0x2e290fc..0x2e29114}) --- Disassembly:}
\codebox{%
0x02e290fc: mov  w8, \#0x7f800000 ; cmp w0, w8 ; b.eq 0x2e2910c ; cbnz w0, 0x2e29118\\
0x02e2910c: mov  w8, \#-0x40800000 ; // 0xBF800000\\
0x02e29110: and  w0, w0, w8 ; ret%
}
\textbf{Arm \texttt{libGLES\_mali.so} --- C/Pseudocode Decompilation:}
\codebox{%
uint32\_t \_mali\_frexpm\_log\_sf32(uint32\_t x) \{\\
~~uint32\_t ax = x \& 0x7fffffffu;\\
~~if (ax > 0x7f800000u) return x | 0x00400000u; // NaN -> quiet NaN\\
~~if (ax == 0 || ax == 0x7f800000u) return x \& 0xbf800000u; // +-0->+-0, +-Inf->+-1.0!\\
~~return normalize\_mantissa\_log(x); // [0.75, 1.5)\\
\}%
}
\textbf{Mesa Panfrost (\texttt{kraid/ops.rs:2031-2047}):} \texttt{OpFrexpM} is not constant-folded in \texttt{kraid}; \texttt{panfrost\_z3\_model.py} implements the exact DDK bit-vector formula.

\subsubsection*{A.2.5 \texttt{F32\_TO\_F16} vs \texttt{V2F32\_TO\_V2F16} Architecture Dispatch}
\textbf{Arm \texttt{libGLES\_mali.so} (\texttt{Valhall::expandF32ToF16} @ \texttt{0x1e81040}) --- Disassembly:}
\codebox{%
0x01e81054: ldrb w8, [x6, \#0x155] ; feature bit 33 (1 on v9--v10, 0 on v11+)\\
0x01e8106c: cmp  w8, \#1 ; b.ne 0x1e81100%
}
\textbf{Arm \texttt{libGLES\_mali.so} --- C/Pseudocode Decompilation:}
\codebox{%
void Valhall::expandF32ToF16(MachineInstr \&MI, const ValhallSubtarget \&ST) \{\\
~~if (ST.hasFeatureBit33() /* v9--v10: tNAx..tVAx */)\\
~~~~emit(Valhall::V2F32\_TO\_V2F16 /* 0x33d */, dst, src0, src1);\\
~~else /* v11--v14: tTUx..tMAx */\\
~~~~emit(Valhall::F32\_TO\_F16 /* 0x216 */, dst, src0);\\
\}%
}
\textbf{Mesa Panfrost (\texttt{kraid/nir.rs:642-689}):} Matches 100\%: emits \texttt{V2F32\_TO\_V2F16} on arch $\le 10$ and \texttt{F32\_TO\_F16} on arch $\ge 11$.

\subsubsection*{A.2.6 \texttt{MKVEC.v2i8} \texttt{src2} Modifier Encoding}
\textbf{Arm \texttt{libGLES\_mali.so} (\texttt{OperandModifiers::getDefaultModId} @ \texttt{0x23e37d0} \& \texttt{ModInfo} @ \texttt{0x4f42d8}) --- Disassembly:}
\codebox{%
0x023e37f4: bl   OperandModifiers::find (0x23e2e20) ; binary search ModInfo[3004..3005]\\
// XMLInstrInfo[MKVEC.v2i8 (0x2e0)] S2: fams=F79(def=b02), mods=2@3004 = [\_\_last\_use, b01]%
}
\textbf{Arm \texttt{libGLES\_mali.so} --- C/Pseudocode Decompilation:}
\codebox{%
uint16\_t getDefaultModId(const XMLInstrInfo \&Info, unsigned slot, uint8\_t famId) \{\\
~~uint16\_t cand = ModifierFamily[famId].defaultModId(); // b02 (0x1f) for F79\\
~~return Info.modInfo(slot).contains(cand) ? cand : 0;  // fails -> only b01 is legal!\\
\}%
}
\textbf{Mesa Panfrost (\texttt{isa-v9-v14.xml}):} Bits \texttt{24:35} are \texttt{<reserved type="zero"/>} with \texttt{<virtual name="lanes2" type="swiz8\_2\_full\_m" exact="b01"/>}. Both implementations agree on \texttt{Concat(src2[15:0], src1.lane1, src0.lane0)}.

\subsubsection*{A.2.7 Additional \texttt{XMLInstrInfo} vs \texttt{isa-v9-v14.xml} Table Deltas (6 Items with Pseudocode)}
\begin{itemize}
  \item \textbf{\texttt{LDEXP.f32} \texttt{widen0}:} DDK \texttt{XMLInstrInfo[0x294].S0.F136} permits \texttt{h0}/\texttt{h1} (\texttt{widen\_f16\_to\_f32(src0)} before scaling); \texttt{isa-v9-v14.xml} omits \texttt{widen0} on v9--v13.
  \item \textbf{\texttt{MMUL.f32} \texttt{sub0/1} \& \texttt{MMUL.v4u8.sat}:} DDK \texttt{0x2e1}/\texttt{0x2e4} decode \texttt{F77} (\texttt{f0/f1/none}) on \texttt{MMUL.f32} and \texttt{.sat} on v9--v11 \texttt{MMUL.v4u8}; Mesa restricts \texttt{sub0/1} to \texttt{v2f16}/\texttt{f16} and gates \texttt{.sat} to v14.
  \item \textbf{\texttt{JUMP} \texttt{address\_mode}:} DDK \texttt{0x291.S4.F37} executes \texttt{pc = (mode == 0) ? pc + off : target}; Mesa fixes \texttt{absolute(1)} in v11--v14.
  \item \textbf{\texttt{LOAD} \texttt{extend} \& \texttt{BLEND} \texttt{vecsize}:} DDK \texttt{0x2b8.S4.F75} decodes \texttt{sext(0)/zext(1)/none(2)} (Mesa: \texttt{encoded\_extend} \texttt{sext/zext\_or\_none}); DDK \texttt{BLEND} \texttt{S4.F132} allows \texttt{v2/v3/v4/none} (Mesa fixes \texttt{v4}).
  \item \textbf{\texttt{LD\_VAR\_BUF*} Flat Varyings \& \texttt{TEX\_WEIGHTS}:} DDK encodes \texttt{src\_flat16}/\texttt{src\_flat32} in \texttt{F49} on \texttt{LD\_VAR\_BUF*} (Mesa splits to \texttt{LD\_VAR\_BUF\_FLAT*} in v14); DDK \texttt{TEX\_WEIGHTS} adds \texttt{project} and \texttt{4x1}/\texttt{1x4} footprints (Mesa: \texttt{2x2} only).
  \item \textbf{Unconditional Branch Idiom:} DDK emits \texttt{\mnem{BRANCH} \src{k1n}, \imm{\#off}} (\texttt{if (0xFFFFFFFFu != 0) pc += off;}); Mesa emits \texttt{\mnem{BRANCH}\imod{.not} \src{k0}, \imm{\#off}} (\texttt{if (!(0u != 0)) pc += off;}). Both evaluate identically to always-taken.
\end{itemize}
\end{document}
